\documentclass[11pt,letterpaper,reqno]{amsart}
\usepackage{amsmath,amssymb,amsthm,mathrsfs}
\usepackage[T1]{fontenc}
\usepackage{lmodern}
\usepackage[expansion=false]{microtype}
\usepackage[margin=1in,headheight=15pt]{geometry}
\usepackage{enumitem}
\usepackage{xcolor}
\definecolor{manuscriptred}{rgb}{0.6,0.1,0.1}
\definecolor{manuscriptblue}{rgb}{0.1,0.3,0.5}
\definecolor{manuscriptmagenta}{HTML}{60445C}
\usepackage[hypertexnames=false,colorlinks,linkcolor=manuscriptred,citecolor=manuscriptblue,urlcolor=manuscriptmagenta]{hyperref}
\usepackage[nameinlink,noabbrev]{cleveref}
\setcounter{tocdepth}{1}
\setlength{\emergencystretch}{3em}
\allowdisplaybreaks[2]

\numberwithin{equation}{section}
\renewcommand{\theequation}{\thesection.\roman{equation}}
\renewcommand{\eqref}[1]{\textup{eq.~\ref{#1}}}
\crefname{equation}{eq.}{eqs.}
\Crefname{equation}{Eq.}{Eqs.}
\creflabelformat{equation}{#2#1#3}
\hypersetup{
  pdftitle={Scalar continuum limits of three-dimensional lattice Yang--Mills theory},
  pdfauthor={}
}
\theoremstyle{definition}
\newtheorem{bodydefinition}{Definition}[section]
\newtheorem{bodylemma}[bodydefinition]{Lemma}
\newtheorem{bodyproposition}[bodydefinition]{Proposition}
\newtheorem{bodytheorem}[bodydefinition]{Theorem}
\newtheorem{bodycorollary}[bodydefinition]{Corollary}
\newtheorem{precisetheorem}{Theorem}
\renewcommand{\theprecisetheorem}{\Alph{precisetheorem}}
\newtheorem*{introductorytheorem}{Theorem 0}
\newcommand{\minorend}{\pushQED{\qed}\renewcommand{\qedsymbol}{$\blacksquare$}}
\newcommand{\majorend}{\pushQED{\qed}\renewcommand{\qedsymbol}{$\blacksquare\square$}}
\newenvironment{definition}[1][]{\minorend\bodydefinition[#1]}{\popQED\endbodydefinition}
\newenvironment{lemma}[1][]{\minorend\bodylemma[#1]}{\popQED\endbodylemma}
\newenvironment{proposition}[1][]{\minorend\bodyproposition[#1]}{\popQED\endbodyproposition}
\newenvironment{theorem}[1][]{\majorend\bodytheorem[#1]}{\popQED\endbodytheorem}
\newenvironment{corollary}[1][]{\majorend\bodycorollary[#1]}{\popQED\endbodycorollary}
\newenvironment{maintheorem}[1][]{\majorend\precisetheorem[#1]}{\popQED\endprecisetheorem}
\newenvironment{theoremzero}{\majorend\introductorytheorem}{\popQED\endintroductorytheorem}
\theoremstyle{remark}
\newtheorem{remark}[bodydefinition]{Remark}
\crefname{bodydefinition}{Def.}{Defs.}
\Crefname{bodydefinition}{Def.}{Defs.}
\crefname{bodylemma}{Lem.}{Lems.}
\Crefname{bodylemma}{Lem.}{Lems.}
\crefname{bodyproposition}{Prop.}{Props.}
\Crefname{bodyproposition}{Prop.}{Props.}
\crefname{bodytheorem}{Thm.}{Thms.}
\Crefname{bodytheorem}{Thm.}{Thms.}
\crefname{precisetheorem}{Thm.}{Thms.}
\Crefname{precisetheorem}{Thm.}{Thms.}
\crefname{bodycorollary}{Cor.}{Cors.}
\Crefname{bodycorollary}{Cor.}{Cors.}
\crefname{remark}{Rem.}{Rems.}
\Crefname{remark}{Rem.}{Rems.}

\usepackage{graphicx}
\newsavebox{\WideFormula}
\newcommand{\FitDisplay}[2]{%
  \sbox{\WideFormula}{$\displaystyle #2$}%
  \ifdim\wd\WideFormula>#1\relax
    \resizebox{#1}{!}{\usebox{\WideFormula}}%
  \else
    \usebox{\WideFormula}%
  \fi}
\begin{document}
\title[Scalar continuum limits of lattice Yang--Mills theory]{Scalar continuum limits of three-dimensional lattice Yang--Mills theory}

\begin{abstract}
We study the action density of the three-dimensional $\operatorname{SU}(2)$
Wilson measure at lattice spacing $L^{-K}$ and inverse coupling
$L^K\mathsf b$, for $\mathsf b>0$. A local probability kernel draws the
coarse links. Contraction and source estimates for the resulting
renormalisation give full-sequence convergence of separated cumulants
of order at least two on fixed tori and subsequential limits when the side tends to infinity.
For translation-invariant Gibbs states the limits have full Euclidean
symmetry and reflection positivity. Their canonical clustering
components reconstruct local scalar Wightman fields with unique
invariant vacua. A reflected block-field test gives nontriviality and
a two-point lower bound at short distances in each component.
\end{abstract}
\maketitle
\tableofcontents


\section*{Introduction}

Let $U$ have the $\operatorname{SU}(2)$ Wilson law on a three-dimensional
periodic lattice of spacing $a_K=L^{-K}$. We study the action density
smeared against $f$,
\[
 \mathcal O_K(f)=a_K^{-1}\sum_p f(x(p))
                   \left(1-\tfrac12\operatorname{tr}\partial U(p)\right).
\]
Here $x(p)$ is the plaquette barycentre and $\partial U(p)$ is the
plaquette holonomy. We take the inverse lattice coupling to be
$\beta=L^K\mathsf b$, with positive label $\mathsf b$. Subtracting the
mean gives the centered observable. Cumulants of order at least two
are independent of this subtraction.

\begin{theoremzero}
Fix a sufficiently large admissible odd blocking factor $L$. For each
positive label $\mathsf b$, the action-density cumulants of order at
least two converge on real continuous tests with pairwise disjoint supports as $K\to\infty$ on a fixed torus of
side $2L^m$. On $\mathbb R^3$, taking the cutoff and volume limits in
either order or jointly gives subsequential limits with a common
subsequence for every separated tuple. These limits have cubic
symmetry and coordinate reflection positivity.

For sequences of translation-invariant infinite-lattice Gibbs states
whose labels tend to $\mathsf b$, the continuum limits are invariant
under all Euclidean isometries and reflection positive at every plane.
Their canonical clustering components yield local scalar Wightman
fields with unique invariant vacua. The connected two-point kernel
in each component is bounded below by a positive multiple of
$\mathsf b^{-2}r^{-6}$ at sufficiently small distances $r$.
\end{theoremzero}
The precise statement, including the conditions on the blocking
factor, the three limiting procedures and the test spaces, is
\Cref{Y-00-Y}.

The construction is in the style of OpenAI's nonlinear sigma-model
paper \emph{The canonical massive continuum limit of the two-dimensional
$\operatorname{O}(3)$ model}, from result family \#215 \cite{OpenAI215}. Given the fine
links $U$, draw each coarse link independently from a normalized
$\operatorname{SU}(2)$ density around a gauge-covariant local average
$\mathsf M(U)(c)$. Here $x$ is the reference inverse coupling at the
fine scale and $a>0$ is the smoothing parameter of the kernel. If
$\rho$ is the fine density, the coarse density is
\[
 \rho'(V')=\int \rho(U)\prod_c k_{ax}(V'(c),\mathsf M(U)(c))\,dU .
\]
The kernel $k_{ax}$ has total mass one in its first argument. This
preserves the partition integral at every step and lets us compare
cutoffs by comparing effective interactions in a common class of
densities. The Gaussian rows describe the reference quadratic form.
After matching the rows, a weighted contraction controls the difference
of the interactions. Its complex-source version controls separated
cumulants by differentiation. A chessboard estimate moves the volume
factor into a bound per occupied cell. This makes the continuum
estimates uniform in the side.

The construction and its Gaussian estimates occupy
\S\ref{sec:blocking-and-free}--\S\ref{sec:compare-cutoffs}. The source
analysis and the limits are in \S\ref{sec:single-source-step}--\S\ref{sec:correlation-limits}.
Growth, symmetry and nontriviality are proved next. In
\S\ref{Y-P5-1} we pass the lower bound to each clustering component and
apply the scalar reconstruction theorem of Appendix~\ref{app:scalar-reconstruction}.

\section{The setting and the main theorem}\label{Y-00-1}

\subsection{Wilson measures, sources and dilation}\label{sec:wilson-conventions}

Write $G=\operatorname{SU}(2)$. On the torus
$T_\ell=\mathbb R^3/\ell\mathbb Z^3$, with $\ell=2L^m$ and
$a_K=L^{-K}$, take the lattice $\Lambda_{K,\ell}=a_K\mathbb Z^3/\ell\mathbb Z^3$.
We allow integers $K,m$ with $K+m\ge1$ in this convention. Each
positively oriented coordinate bond has one independent variable in
$G$, with the inverse variable on the reversed bond. The plaquettes
are counted once, in the three positive coordinate planes, and the
trace of the ordered plaquette product is independent of the choice
of its base point and orientation. Write $dU$ for the product of
normalized Haar probabilities, $a_p(U)=1-\tfrac12\operatorname{tr}\partial U(p)$
and $A(U)=\sum_p a_p(U)$. The Wilson probability is
$Z_{K,\ell,\beta}(0)^{-1}e^{-\beta A(U)}dU$. The signed
permutations $W_3$ and lattice translations act on bonds and
plaquettes. They preserve this probability, since they permute
bond variables with possible inversions and preserve $A$.

\begin{definition}[Action-density sources]\label{support:150-2}
For a bounded $f$ on $T_\ell$ define
\[\mathcal O_K(f)=a_K^{-1}\sum_p f(x(p))a_p(U),\]
where $x(p)$ is the
plaquette barycentre. For a tuple $f=(f_1,\ldots,f_n)$ put
\[
 Z_{K,\ell,\beta}(z)=\int\exp\left\{-\beta A(U)
                          +\sum_i z_i\mathcal O_K(f_i)\right\}\,dU.
\]
The joint cumulant $S^T_{n;K,\ell,\beta}(f)$ is the derivative
$\partial_{z_1}\cdots\partial_{z_n}\log(Z(z)/Z(0))|_{z=0}$.
The ratio is entire by boundedness of the observables. It equals $1$
at zero, so the logarithm is holomorphic on a neighborhood of zero.
\end{definition}

\begin{lemma}[Centered moments and cumulants]
Let $X_i$ be bounded random variables and $Y_i=X_i-\mathbb E X_i$.
Then
\begin{equation}\label{support:9-14}
\mathbb E\prod_{i=1}^nY_i
=\sum_{\pi\in\mathrm{Part}_{\ge2}(n)}
    \prod_{B\in\pi}\kappa(X_i:i\in B).
\end{equation}
\end{lemma}
\begin{proof}
In the commutative algebra generated by $z_i$ subject to $z_i^2=0$,
write $z_B=\prod_{i\in B}z_i$. The coefficient of $z_B$ in
$\log\mathbb E\prod_i(1+z_iX_i)$ is $\kappa(X_i:i\in B)$.
Centering subtracts $\sum_i z_i\mathbb E X_i$ from this logarithm,
so its coefficients in degrees at least two are unchanged and its
linear coefficients vanish. Exponentiating gives the partition sum.
For each partition the factorial in the exponential cancels the
number of orders of its blocks. The empty partition has weight $1$.
\end{proof}

\begin{lemma}[Exact change of length units]\label{lem:dilation}
Let $j\in\mathbb Z$ and $f^{[j]}(x)=f(L^{-j}x)$. For each tuple
$f$ of length $n\ge1$ one has
\begin{equation}\label{support:150-42}
 S^T_{n;K,\ell,\beta}(f)=L^{jn}
       S^T_{n;K-j,L^j\ell,\beta}(f^{[j]}).
\end{equation}
The full-cutoff limits at $(\ell,\mathsf b)$ and
$(L^j\ell,L^j\mathsf b)$ exist together. When they exist,
\begin{equation}\label{support:150-43}
 S^{(\ell,\mathsf b)}_n(f)=L^{jn}
         S^{(L^j\ell,L^j\mathsf b)}_n(f^{[j]}).
\end{equation}
\end{lemma}
\begin{proof}
Dilation maps the lattice bijectively onto the lattice with spacing
$L^{j-K}$. It preserves the graph, plaquette products, Haar law and
$\beta$. The observable on the rescaled lattice is
$\mathcal O_{K-j}(f^{[j]})=L^{-j}\mathcal O_K(f)$. Multilinearity
of the joint cumulant gives the finite-cutoff identity. For the
limit use $\beta=L^K\mathsf b=L^{K-j}(L^j\mathsf b)$.
\end{proof}


\subsection{Parameters and limits}
The parameters below ensure that the coarse interaction stays in the
class on which the step and the comparison estimates are proved.

\begin{definition}[\textup{(S1)}: the parameter table]\label{Y-00-S1} The gauge group is $G=\operatorname{SU}(2)$, and the dimension is 3. We work with one table, \textup{(O0)} to \textup{(O13)}, with all its floors: the table of item~(1) of \S\ref{Y-H0-1}, repeated in item~1 of \S\ref{Y-H9-1}. The number $L$ is odd and is the $L$ of the table, so that $L\ge L_*^{\rm su}\ge10^{2048.9}$ (see the remark following \Cref{Y-H9-T}), and it satisfies \textup{(L3su)} and \textup{(L5*)}. Further, $M=L^{m'}$ with $m'\ge1$, subject to \textup{(M2)}. We fix $\kappa\ge6$ and put $\kappa_1=\kappa+1$.

We set $p_0(x)=A_0(\log x)^{\mathfrak p_0}$ with $\mathfrak p_0>2$, and $\varepsilon(x)=p_0(x)x^{-1/2}$; further $c_a(x)=E_ap_0(x)^3x^{-1/2}$ and $c_h(x)=e^{-c_hp_0(x)^2}$, where $c_h$ without an argument, as in the exponent, is the number of \textup{(O11)}, while $c_h(x)$, with an argument, is the function defined here; and $r(x)=\lceil20c_R\log x/M\rceil$, $s=8r$, $\check r=r+3$, $\rho(x)=\lceil c_R\log x\rceil$, which is the radius denoted $\varrho(x)$ in \S\ref{Y-G0-1} and in item~1 of \S\ref{Y-H9-1}. The number $x_s\ge e^2$ is the last of the floors.

The constants are fixed in this order: first the table, then the label, then the numbers of \textup{(O14)}. Every number of item 2 but one belongs to \textup{(O14)}, so no floor depends on it; the exception is the threshold $N^*(K,\mathsf b')$ of \Cref{Y-P4-Th4}\,\textup{(4.9)}\,\textup{(i)}. \end{definition}

\begin{definition}[\textup{(S2)}: the data of a torus]\label{Y-00-S2} The data are those of \Cref{Y-H9-T}, which agree with those of item~1 of \S\ref{Y-H6-1}. They consist of a label $\mathsf b'\ge x_s$; of a side $N=2L^{m_l}$ with $m_l\ge m'$, subject to \textup{(Tor)}, that is, $N\ge2Ms(L\mathsf b')$; and of a cutoff $K\ge1$, with $a_K=L^{-K}$ and $\beta=L^K\mathsf b'$. By $\mu_{K,N,\beta}$ we denote Wilson's measure on $a_K\mathbb Z^3/N\mathbb Z^3$. The label, the side and the cutoff, with $a_K$ and $\beta$, are what we call the data. \end{definition}

\begin{definition}[\textup{(S3)}: the field, cumulants, the centred functional, distances]\label{Y-00-S3} For a function $f$ we put $\mathcal O_Kf=a_K^{-1}\sum_pf(x(p))a_p(U)$, the sum running over the plaquettes $p$, where $a_p=1-\tfrac12\operatorname{tr}\partial U(p)\in[0,2]$ and $x(p)$ is the barycentre of $p$. On a torus we write $S^T_n$, and in a state $\rho$ we write $S^T_n[\rho]$, for the joint cumulants of the $\mathcal O_Kf_i$; and we put $\sigma_{K,N}(f_1,\dots,f_n)=\langle\prod(\mathcal O_Kf_i-\langle \mathcal O_Kf_i\rangle)\rangle$.

For a family $S=(S_n)_{n\ge2}$, the centred functional $\sigma_S$ is that of \eqref{support:9-14}: \[ \sigma_S(X)=\sum\prod_JS_{|J|}(X_J), \] where the sum runs over the partitions of $X$ into blocks $J$ of size $\ge2$, and $\sigma_S(\emptyset)=1$ for the empty tuple $\emptyset$.

Distances are sup distances unless we say ``Euclidean''; a tuple whose Euclidean type is $(d;R)$ has the sup type $(d/\sqrt3;R)$ (item~1 of \S\ref{Y-H9-1}). We put $\check\delta_n:=\min\{d_{\sup},1/(4n)\}$, where $d_{\sup}$ is a lower bound for separation in the sup distance; we take $d_{\sup}=d$ in Sections~\ref{Y-P2-1} and~\ref{Y-P3-1}, and $d_{\sup}=d/\sqrt3$ in Sections~\ref{Y-P4-1} and~\ref{Y-P5-1}. Further, $\|f\|_\sharp=\max\{\|f\|_\infty,40M\,\mathrm{Lip}f\}$, where $\mathrm{Lip}f$ is taken with respect to the Euclidean distance, as in item~1 of \S\ref{Y-H9-1}; $f^{[z]}(x)=f(L^{-z}x)$, and $W_3$ denotes the cubic group.

In the limit statements, a separated tuple consists of real
continuous functions with compact, pairwise disjoint supports;
on a fixed torus the functions are continuous on that torus.
For a volume limit, a tuple on $\mathbb R^3$ is periodised once
the side exceeds twice the diameter of the union of its supports.
The definition of $\mathrm{Lim}_G$ in \textup{(S6)} uses the
$C^1_c$ subspace. Write $\mathfrak S(O)$ for the separated
$C^1_c$ tuples in $O$, including the empty tuple, and
$\mathbb C\mathfrak S(O)$ for their finite complex linear combinations.
The finite-cutoff estimates also allow the
bounded sources specified in their statements. For a tuple of
distinct points, write $d(x)=\min_{i\ne j}|x_i-x_j|_2$, with
$d(x)=\infty$ when $n=1$.
\end{definition}

\begin{definition}[\textup{(S4)}: the drawn block fields and the subtractions]\label{Y-00-S4} The fields $(V_0=U,V_1,\dots,V_K)$ are distributed by the joint law of the $K$ drawn steps (see item~2 of \S\ref{Y-J-1}, under ``The step'' and ``Runs'', and Step~11\,\textup{(0)} of the proof of \Cref{Y-H9-T}). The level $k$ is the lattice $L^{-(K-k)}\mathbb Z^3/N\mathbb Z^3$, whose depth is $K-k$, and $x_k=L^{K-k}\mathsf b'$. We write $\rho_k^z$ for the $k$-fold drawn density of Wilson's weight multiplied by $\exp\{\sum_tz_t\mathcal O_Kg_t\}$. For $0\le s\le K-1$ the subtractions are $c^{(s)}_K(f):=\frac{d}{dz}\log\rho_{K-s}^{[zf]}(1)$ at $z=0$, where $\rho_k^{[zf]}$ is the density $\rho_k^z$ above for the single source $f$ with the variable $z$, that is, the $k$-fold drawn density of Wilson's weight multiplied by $\exp\{z\mathcal O_Kf\}$, and the argument $1$ is the trivial field of the level $k$; and $\hat c_K:=c^{(0)}_K$. These numbers depend on $K$ and $N$. \end{definition}

\begin{definition}[\textup{(S5)}: sources, the setting of \Cref{Y-P1-DATAo}]\label{Y-00-S5} The number $Q\ge1$ is arbitrary, and $\check\delta,\check\rho>0$. The sources $g_1,\dots,g_Q$ are bounded functions on $\mathbb R^3/N\mathbb Z^3$, complex in \S\ref{Y-P1-1} and real from \S\ref{Y-P2-1} on. The sets $\{g_t\ne0\}$ are pairwise $\ge\check\delta$ apart, and each of them lies within $\check\rho$ of a single point. \end{definition}

\begin{definition}[\textup{(S6)}: states and limits]\label{Y-00-S6} For any $\mathsf b'>0$ we let $\mathbb G_K(\mathsf b')$ consist of those Gibbs states $\nu$ of Wilson's action at $\beta=L^K\mathsf b'$ on $a_K\mathbb Z^3$ (\textup{\ref{support:14-9}}\,\textup{(2)}) for which $(\tau_v)_*\nu=\nu$ holds for every $v\in a_K\mathbb Z^3$. This set lies in $\mathcal G^{\rm per}$ (\textup{\ref{support:14-9}}\,\textup{(3)}) and contains $\mathcal G^{\rm tor}$, the set of all limit points of tori (\textup{\ref{support:14-9}}\,\textup{(4)}, \textup{\ref{support:14-16}}\,\textup{(c)}).

For $\mathsf b>0$, $\mathrm{Lim}_G(\mathsf b)$ is the set of the families $S=(S_n)_{n\ge2}$ for which $S_n(X)=\lim_jS^T_n[\nu_j](X)$ holds on all separated $C^1_c$ tuples over $\mathbb R^3$, for some $K_j\to\infty$, $\nu_j\in\mathbb G_{K_j}(\mathsf b_j)$ and $\mathsf b_j\to \mathsf b$. The subset $\mathrm{Lim}_{\rm tor}(\mathsf b)$ consists of the families obtained in this way with $\nu_j\in \mathcal G^{\rm tor}$.

By the limits of \textup{(R5)} we mean the limits of \Cref{Y-P2-ThR}\,\textup{(R5)} in any of its three orders: \textup{(a)} the continuum first, \textup{(b)} the volume first, \textup{(c)} both jointly. Every limit of \textup{(R5)}\,\textup{(b)} belongs to $\mathrm{Lim}_{\rm tor}(\mathsf b)$, by item~1 of \S\ref{Y-P4-1}. We call the $\sigma_\xi$ of a canonical decomposition $(\mathsf X,\lambda,(\sigma_\xi))$ of $\sigma_S$ (\textup{\ref{support:35-34}}) the components of $S$. \end{definition}

\begin{definition}[\textup{(S7)}: labels; $z_0(\mathsf b)$, $z(\mathsf b)$, $U(\mathsf b)$; normal form]\label{Y-00-S7} We use $\mathsf b>0$ for the label of a limit and $\mathsf b'$ for that of an approximant, that is, of a torus or of a state. In a dilation, the letter $z$ written without an argument, as opposed to $z_0(\mathsf b)$ and $z(\mathsf b)$ below, is a bound variable: it stands for any integer that the clause admits.

We put $z_0(\mathsf b):=$ the least integer $z\ge0$ with $L^z\mathsf b\ge x_s$; it brings one label into the data. We put $z(\mathsf b):=$ the least integer $z\ge0$ with $L^z\mathsf b>x_s$; it brings a whole neighbourhood of the label into the data. Thus $z(\mathsf b)=z_0(\mathsf b)$, except when $L^{z_0(\mathsf b)}\mathsf b=x_s$, in which case $z(\mathsf b)=z_0(\mathsf b)+1$.

Let $U(\mathsf b):=[\mathsf b/2,2\mathsf b]\cap[L^{-z(\mathsf b)}x_s,\infty)$. This is a compact neighbourhood of $\mathsf b$, determined by $\mathsf b$ alone; it satisfies $L^{z(\mathsf b)}U(\mathsf b)=U(L^{z(\mathsf b)}\mathsf b)\subset[x_s,\infty)$; and for a sequence $\mathsf b_j\to \mathsf b$, at most finitely many of the labels lie outside $U(\mathsf b)$. We leave out of the sequence the $j$ with $\mathsf b_j\notin U(\mathsf b)$. This changes no convergence and no limit.

The label of a limit is in normal form if $\mathsf b>x_s$; then $z(\mathsf b)=0$ and $U(\mathsf b)=[\mathsf b/2,2\mathsf b]\cap[x_s,\infty)$. Any other $\mathsf b>0$ is brought to normal form by the dilation $L^{z(\mathsf b)}$. \end{definition}

\begin{definition}[\textup{(S8)}: numbers of \textup{(O14)} at a label]\label{Y-00-S8} Let $X$ be any number of \textup{(O14)}, regarded as a function of the label. The suprema below are taken over its original formula. \textup{(n1)} Except in \textup{(n2)}, if the label belongs to the data (a torus, a state at $\mathsf b'\ge x_s$, or a limit at a fixed label $\mathsf b\ge x_s$), then $X(\mathsf b')$ is simply the value of its formula.

\textup{(n2)} In Sections~\ref{Y-P4-1} and~\ref{Y-P5-1}, in a clause on $\mathrm{Lim}_G(\mathsf b)$ or on a frame, and with the label in normal form, we put $X(\mathsf b):=\sup_{\mathsf b'\in U(\mathsf b)}X(\mathsf b')$ and $j^\circ_U(\mathsf b,\check\delta):=\sup_{U(\mathsf b)}j^\circ(\cdot,\check\delta)$. These suprema are finite, by \textup{(Lab1)} in item~1 of \S\ref{Y-P3-1} and by item~1 of \S\ref{Y-P4-1}. For $\mathsf b\le x_s$, a number taken ``at $L^z\mathsf b$'', with $z=z(\mathsf b)$, means its supremum over $L^zU(\mathsf b)$.

\textup{(n3)} In \S\ref{Y-P3-1}, at a limit of \textup{(R5)} whose labels move, that is, a limit taken along approximants with labels $\mathsf b_j\to\mathsf b$ rather than at the fixed label of \textup{(n1)}, the number is that of \textup{(Lab3)}, namely $\limsup_jX(\mathsf b_j)$ along any subsequence; for a label in normal form, this is $\le$ the number of \textup{(n2)}. \end{definition}

\begin{remark}[The one label $\mathsf b=x_s$] Let $\mathsf b=x_s$. For every sequence, the numbers that enter \Cref{Y-P4-Th4}\,\textup{(4.6)} and \textup{(4.9)}\,\textup{(ii)} are then given by \textup{(4.8)} at $z=1$; and when $\mathsf b_j\ge x_s$ along the sequence, the suprema over $[x_s,2x_s]$ may be used as well, as the remark in item~1 of \S\ref{Y-P4-1} explains. \end{remark}

\begin{definition}[\textup{(S9)}: boards]\label{Y-00-S9} We use a single family of boards $B_i(o)$, for $-m_l\le i\le K-1$ and $o\in a_K\mathbb Z^3$: the walls of $B_i(o)$ are $o_\mu+L^{-i}(\mathbb Z+\tfrac12)$, and its cells have side $L^{-i}$. We write $B_i:=B_i(0)$; its cell $0$ is centred at $0$. In the notation of \textup{\ref{support:4-1}}\,\textup{(2)} the board $B_i(o)$ has the parameters $(a_K,\lambda=L^{K-i},\mathsf n_\mu=2L^{m_l+i},o)$, and in that of \textup{\ref{support:18-2}} the parameters $(i,o)$.

Without arguments, $B_0$ is the unit board. As to offsets, \S\ref{Y-P2-1} takes $o\in\mathbb Z^3$ or $o\in L^{-s}\mathbb Z^3$, where $s$, with $0\le s\le K-1$, is a depth as in \textup{(S4)}; \S\ref{Y-P4-1} takes any $o\in a_K\mathbb Z^3$ and $i\ge0$; and in Sections~\ref{Y-P3-1}, \ref{Y-P5-1} and~\ref{Y-P6-1} the boards carry no offset. \end{definition}

\begin{definition}[Item 2: constants in the main theorem]
The constants are those of the results cited below, evaluated at the
label according to \textup{(S8)}. Except for $N^*(K,\mathsf b')$,
they belong to \textup{(O14)} and are fixed after the table and the
label; no floor depends on them.

The source constants $j^\circ$, $r^\circ$, $e_h$ and
$G_1(\mathsf b,\check\delta,\check\rho)$ are defined in
\Cref{Y-P1-D1}; in particular,
$\mathrm{sep}^\circ=8\log(8L^{1/2})$.
Their depth-$s$ versions $r_s^\circ$, $e_h^{(s)}$ and $G_{1,s}$
are defined in item~2 of \S\ref{Y-P2-1}. That item also sets
\[
\begin{gathered}
c_{\rm vol}(\mathsf b)=225/M^3+\tfrac32\mathsf N_Zp_0(\mathsf b)^2
                  +3\log(384/\varepsilon(\mathsf b)^3),\\
1+0.182C_{\rm vol}(N,\mathsf b)
       \le e^{c_{\rm vol}(\mathsf b)N^3}.
\end{gathered}
\]

With $A_4$ from item~1 of \S\ref{Y-P3-1}, put
$A_5=64(4e)^5A_4$, $A_6=3^{5/2}A_5$ and
$C_E=(10^{53})^6M\,8A_6$. For $0<\mathsf b<x_s$ use the
dilated constants of \Cref{Y-P3-ThG}\,\textup{(G2)(iv)} and
\textup{(G4)(e)}. The insertion constant
$\check\Gamma_*(\mathsf b,\eta)$ and norm
$|\!|\!|h|\!|\!|_{\eta,y}$ are those of
\Cref{Y-P3-ThG}\,\textup{(G6)} and \textup{(G6+)}.
Set
\[
\begin{gathered}
C_n^\infty(d;R)=1.45^{n+1}n!\,A_6^n(n!)^5
                  \max\{1,d^{-5n}\}(1+R)^{3n},\\
C_1(R)=A_6(1+R)^3,
\quad C^{\mathsf b}=(C_1,(C_n^\infty)).
\end{gathered}
\]
The constants $\hat G$, $\Gamma(\mathsf b,\mathsf l)$,
$C_\varphi$ and $C_n^\sigma$ are those of item~2 of
\S\ref{Y-P4-1} and \Cref{Y-P4-Th4}\,\textup{(4.6)}; in particular,
$(c_{\rm I},C_{\rm I},\mathsf p_{\rm I})
=(1/(48L),24Le^{c_{\rm vol}}\hat G,1)$.
The threshold $N^*(K,\mathsf b')$ is the one in
\Cref{Y-P4-Th4}\,\textup{(4.9)(i)}.
The definition ``New numbers'' of \S\ref{Y-P5-1} gives
$\bar c_*=3A_6/(4\pi)$ and $C_{\rm pt}$; set
$C_\wp=32C_{\rm pt}$, $\wp=5$, $c_0=1/2$ and
$C_*=C_\wp+\bar c_*$.

For the lower bounds, \Cref{Y-P6-N} defines $J(\mathsf b)$,
$x^\#$, $x^{\#\#}$ and the pure constants
$D_0=10^{140}$,
$c_F\in[5.4\cdot10^{-5},4.0\cdot10^{-4}]D_0^{-5}$ and
$c_*=c_F^2/7\ge10^{-1410}$.
For $\mathsf b\ge x_s$, let
\[
\begin{gathered}
y(\mathsf b)=\mathsf b\max\{c_{\rm vol}(\mathsf b)^{1/3},L^{1-m'}\},\\
j_*(\mathsf b)=\min\{j\ge1:L^j\mathsf b
                  \ge\max\{x^\#,y(\mathsf b)^{4/3}\}\},
\end{gathered}
\]
and let $j^+(\mathsf b)$ be the supremum of $j_*$ over
$[\mathsf b/2,2\mathsf b]\cap[x_s,\infty)$.
For $0<\mathsf b<x_s$, set
$j_*(\mathsf b)=z_0(\mathsf b)+j_*(L^{z_0(\mathsf b)}\mathsf b)$.
Finally, $c_{**}=(4/L)^6c_*$ and, with $z=z(\mathsf b)$,
$r_*(\mathsf b)=L^{-z}D_0L^{-j^+(L^z\mathsf b)}$,
as in \Cref{Y-P5-ThC}\,\textup{(C4)}.
\end{definition}

\begin{maintheorem}[Bounds and continuum limits]\label{Y-00-Y} Assume \textup{(S1)}. Then \textup{(Y0)} to \textup{(Y7)}, \textup{(Y7$'$)}, \textup{(Y8)} and \textup{(Y9)} below hold, and so does the corollary with which they close. \end{maintheorem}

\textup{(Y0)} \emph{Subtraction.} For real $f$ the number $c^{(s)}_K(f)$ is real, and it is linear in $f$. One has $c^{(s)}_K(f\circ R)=c^{(s)}_K(f)$ for the reflections $R:y_\mu\mapsto c-y_\mu$ with any $c\in L^{-s}\mathbb Z$, and for the translations by $L^{-s}\mathbb Z^3$; for $\hat c_K$ every integer $c$ is admitted, and the reflections with odd $c$ are the mirrors in the walls of $B_0$. Moreover, for given data and $0\le s\le K-1$, \[ c^{(s)}_K(f)=\frac{d}{dz}\Big|_{z=0}\sum_{1\le j\le K-s}V_j(z)+\sum_X\delta^h[f](K-s)(X,1), \] where $X$ runs over $\mathcal D$ of the level $K-s$ and the $V_j(z)$ are given by \Cref{Y-H4-U4}\,\textup{(h)} along the run whose single source is $f$, with the variable $z$; with the argument $z$, the $V_j(z)$ are numbers, not the fields $V_j$ of \textup{(S4)}. If $\{f\ne0\}$ lies within $\check\rho$ of a single point, as each source of \textup{(S5)} does, then the last sum has modulus $\le e_h^{(s)}\|f\|_\infty/r^\circ_s$, at any $\check\delta>0$.

\begin{proof} Clause \textup{(Y0)} is \Cref{Y-P2-ThR}\,\textup{(R0)}, with Lems.~\ref{Y-P2-L2} and~\ref{Y-P2-L3}. At $s=0$ the formula for $c^{(0)}_K(f)$ and the bound on its last sum come from \Cref{Y-P1-Jo}\,\textup{(J7$^\circ$)}\,\textup{(vii)} and \textup{(J6$^\circ$)}\,\textup{(e)}. \end{proof}

\textup{(Y1)} \emph{Per unit of volume.} In \textup{(Y1)}, in the bound on cells and in lines 1 and 2 of \textup{(Y2)}, and in \textup{(Y7)}\,\textup{(D)}, \textup{(N0)} to \textup{(N3)} and the first assertion of \textup{(N4)}, $\mathsf b$ is the label of the data, written $\mathsf b'$ in \textup{(S2)}, so that $\mathsf b\ge x_s$; in the rest of \textup{(N4)} and in \textup{(N5)} it is the label of the limit, that is, a fixed label or the limit of the sequence of labels. Throughout \textup{(Y7)}, except in \textup{(N6$'$)}, we assume $\mathsf b\ge x_s$, and for the parts \textup{(N3)}, \textup{(N4)} and \textup{(N5)} of \textup{(Y7)}, every other $\mathsf b>0$ is treated in \textup{(Y8)}. In \textup{(N6$'$)} the label is any $\mathsf b>0$. For given data, with sources as in \textup{(S5)}, and for all $K\ge1$ and all $Q\ge1$, \[ |\langle\prod_t(\mathcal O_Kg_t-\hat c_Kg_t)\rangle|\le G_1(\mathsf b,\check\delta,\check\rho)^Q(1+0.182C_{\rm vol}(N,\mathsf b))\prod\|g_t\|_\infty\le G_1^Qe^{c_{\rm vol}(\mathsf b)N^3}\prod\|g_t\|_\infty. \] Here $C_{\rm vol}$ enters to the first power only. For every bounded Borel function $F$ of $V_K$, with the same $K$, $Q$ and sources, one has $|\langle F(V_K)\prod_t(\mathcal O_Kg_t-\hat c_Kg_t)\rangle|\le\sup|F|\,G_1^Q(1+0.182C_{\rm vol}(N,\mathsf b))\prod\|g_t\|_\infty\le\sup|F|\,G_1^Qe^{c_{\rm vol}(\mathsf b)N^3}\prod\|g_t\|_\infty$.

\begin{proof} The proof is that of \Cref{Y-P1-Cor}\,\textup{(a)}, \textup{(b)}, which Proof~11 of \S\ref{Y-P1-1} gives. At a depth $s\ge0$ the bound is \Cref{Y-P2-ThR}\,\textup{(R8)}\,\textup{(d)}\,\textup{(i)}, with Lems.~\ref{Y-P2-L3} to~\ref{Y-P2-L5}. \end{proof}

\textup{(Y2)} \emph{Bounds free of the side, of the cutoff and of the places.} \emph{Cells.} Take the board $B_i(o)$ with $o\in\mathbb Z^3$ and $-m_l\le i\le0$, and let $h=L^{-i}$. If every support lies in one open cell, at distance $\ge\check\delta/2$ from the walls, and $q$ cells are occupied, then the moment of \textup{(Y1)} is $\le G_1^Qe^{qc_{\rm vol}(\mathsf b)h^3}\prod\|g_t\|_\infty$.

\emph{Lines 1, 2.} Let $f=(f_1,\dots,f_n)$ be a tuple of real bounded functions on $\mathbb R^3/N\mathbb Z^3$ whose sets $\{f_i\ne0\}$ are pairwise $\ge d$ apart, each lying in some sup-ball of radius $R_f$. Wherever these balls lie, for all $K\ge1$, all $N$ of the data, all $n\ge2$ and all $d,R_f>0$, \[ |\sigma_{K,N}(f)|\le A_5(\mathsf b)^n(n!)^5\max\{1,d^{-5n}\}(1+R_f)^{3n}\prod\|f_i\|_\infty, \] and $|S^T_n(f)|\le1.45^{n+1}n!$ times the right side of this bound. The limits of \Cref{Y-H9-T}\,\textup{(T3)} obey these two bounds too.

\emph{Line 3.} For $\mathsf b'\ge x_s$, all $\rho,\nu\in\mathbb G_K(\mathsf b')$, all $K\ge1$ and $n\ge1$, and tuples of the Euclidean type $(d;R)$, \[ |\rho(\prod(\mathcal O_Kf_i-\nu(\mathcal O_Kf_i)))|\le A_6(\mathsf b')^n(n!)^5\max\{1,d^{-5n}\}(1+R)^{3n}\prod\|f_i\|_\infty; \] and for $n\ge2$, $|S^T_n[\rho]|\le C_n^\infty(d;R)\prod\|f_i\|_\infty$. In line 3 the label is $\mathsf b'$.

\emph{Line 4.} Let $K_j\to\infty$ and $\mathsf b_j\to\mathsf b>x_s$. After discarding the finitely many labels outside $U(\mathsf b)$, the family $(K_j,\mathbb G_{K_j}(\mathsf b_j))$ is a class system that obeys $C^{\mathsf b}$, in the sense of \textup{(Cls1)}, \textup{(Cls2)}, \textup{(Bd1)} and \textup{(Bd2)} of \textup{\ref{support:35-2}}.

\begin{proof} The bound on cells is \Cref{Y-P2-ThR}\,\textup{(R2)}, with \Cref{Y-P2-L9}. \Cref{Y-P3-ThG}\,\textup{(G3)}, with \Cref{Y-P3-L3}, is a bound of the same kind. It covers more than \textup{(R2)} in that each occupied cell $t$ has data $(\check\delta_t,\check\rho_t)$ of its own, and less in that \textup{(G3)} is stated on the board $B_0$ alone, for $\check\delta_t\in(0,1]$ and $K\ge\max_tj^\circ(\mathsf b,\check\delta_t)$. Bounds of the same kind, as seminorms with free centres, are \Cref{Y-P4-Th4}\,\textup{(4.1)}, \textup{(4.2)}, by Steps~1 and~2 of its proof.

Lines 1 and 2 are \Cref{Y-P3-ThG}\,\textup{(G2)}\,\textup{(i)}, \textup{(ii)}, proved in \Cref{Y-P3-L4}; a form with $C_n^{(2)}$ is \Cref{Y-P2-ThR}\,\textup{(R3)}, with \Cref{Y-P2-L10}. For the limits of \Cref{Y-H9-T}\,\textup{(T3)}, \Cref{Y-P2-ThR}\,\textup{(R3)} gives the bounds with $C_n^{(2)}$; with $A_5$ they follow from \Cref{Y-P3-ThG}\,\textup{(G2)}\,\textup{(i)}, \textup{(ii)} at fixed $N$, as $K\to\infty$.

Line 3 is proved in \Cref{Y-P4-Th4}\,\textup{(4.3)}\,\textup{(a)}, \textup{(b)}, by Step~3 of its proof, and also in \Cref{Y-P5-ThE}\,\textup{(E0)}, by Step~1 of \S\ref{Y-P5-1}. Line 4 is \Cref{Y-P5-ThE}\,\textup{(E0)}. \end{proof}

\textup{(Y3)} \emph{Limits, for every $\mathsf b>0$.} \textup{(a)} Continuum first: for each side, the limit as $K\to\infty$ exists along the whole sequence of cutoffs, without passing to a subsequence; every sequence of sides $l_j=2L^{m_j}$ with $l_j\to\infty$ has a subsequence converging on all tuples; and the same holds for $\mathsf b_j\to \mathsf b$. \textup{(b)} Volume first: along one subsequence of sides $l_j\to\infty$, common to every fixed $K$, the measures $\mu_{K,l_j,L^K\mathsf b}$ tend weakly to a state $\nu_K\in\mathbb G_K(\mathsf b)$ that is invariant under $W_3$ and reflection positive at planes of sites and at mid planes; after this one passes to a subsequence of $K$. \textup{(c)} Jointly: every sequence $(K_j,l_j)$ with $K_j\to\infty$ and $l_j\to\infty$ has a convergent subsequence, and so does every such sequence at labels $\mathsf b_j\to \mathsf b$. \textup{(d)} If $K_j/l_j^3\to\infty$, then the limit points in \textup{(c)} coincide with those in \textup{(a)} along the same sides and labels.

Every sequence $(K_j,\rho_j,\mathsf b_j)$ with $\rho_j\in\mathbb G_{K_j}(\mathsf b_j)$, $K_j\to\infty$ and $\mathsf b_j\to\mathsf b$ has a convergent subsequence; in particular $\mathrm{Lim}_{\rm tor}(\mathsf b)$ and $\mathrm{Lim}_G(\mathsf b)$ are not empty. For any $\mathsf b>0$, the limit points of joint sequences $(K_j,N_j)$, in which the $N_j$ are the sides themselves, at labels $\mathsf b_j\to \mathsf b$ with $N_j\ge N^*(K_j,\mathsf b_j)$ are functionals of members of $\mathrm{Lim}_{\rm tor}(\mathsf b)$.

\begin{proof} The limits are \Cref{Y-P2-ThR}\,\textup{(R5)}, with \Cref{Y-P2-L12}, together with \Cref{Y-P4-Th4}\,\textup{(4.3)}\,\textup{(c)}, \textup{(4.8)} and \textup{(4.9)}\,\textup{(i)}. The assertion on joint sequences with $N_j\ge N^*(K_j,\mathsf b_j)$ is \Cref{Y-P4-Th4}\,\textup{(4.9)}\,\textup{(i)}, proved in Step~8 of its proof; \Cref{Y-P5-ThE}\,\textup{(E6)} is the same statement. \end{proof}

\textup{(Y4)} \emph{Every limit $S$ of \textup{(R5)}.} Each $S_n$ is real, symmetric and $n$-linear. The cumulants $S_n$ and centered moments $\sigma_S$ obey the bounds of lines~1 and~2 of \textup{(Y2)}, with the label conventions of \textup{(S8)}. The family $S$ is invariant under all translations and under $W_3$, and $\sigma_S$ is reflection positive at every coordinate mirror. On the set of points of $(\mathbb R^3)^n$ whose entries are pairwise different, there is exactly one distribution that takes the value $S_n(f_1,\dots,f_n)$ at $f_1\otimes\cdots\otimes f_n$ for every smooth separated tuple $(f_1,\dots,f_n)$, and its order on compact subsets of that set is $\le3n+1$. For every integer $z$, one has $S^{\mathsf b}_n(f)=L^{zn}S^{L^z\mathsf b}_n(f^{[z]})$ along matched subsequences: each side and cutoff $(l,K)$ is changed to $(L^zl,K-z)$, with the same order of limits. For moving labels, the matching is $\mathsf b_j\mapsto L^z\mathsf b_j$ at each member of the sequence. The symmetry group of $S$ is $W_3$ or $\operatorname{O}(3)$.

The family $S$ determines $\mathcal H(S)$, $\Omega$, $e^{-tH}$ with $H\ge0$, the momenta, $W_2$, and the class $[X]$ of continuous entries. Its centered functional $\sigma_S$ obeys the bounds of \textup{(Y2)}, with $A_6$ for the Euclidean type, and the growth condition \Cref{support:26-33} with $(13,1,C_E(\mathsf b),13)$. The extensions of its moments to $\mathcal S_{\ne}$, the Schwartz tests flat on the diagonals, satisfy
\[
 |S_n^{\sigma_S}(\varphi)|
 \le (10^{36}A_5)^n(n!)^5|\varphi|_{12n+1}.
\]
For the extensions of the cumulants $S_n$, the same estimate has the further factor $1.45^{n+1}n!$. At finite $(K,N)$ of the data, the centered lattice distributions obey the displayed bound for smooth tests flat on the diagonals and supported in $(-N/4+1,N/4-1)^{3n}$, as in \Cref{Y-P3-ThG}\,\textup{(G4)(d)}. The extensions of $\sigma_S$ are invariant under translations and $W_3$, satisfy the positivity \eqref{support:26-42} at coordinate mirrors, and are symmetric.

For $\eta\in(0,1]$, $F,F'\in\mathbb C\mathfrak S(\{x_3>\eta\})$, real $h\in C^1_c(\{|x_3|<\eta/4\})$ and $y\in\mathbb R^3$, a field between two vectors obeys a bound free of $n$, \[ |\sigma_S(\theta\bar F\cdot(h)\cdot F')|\le\check\Gamma_*(\mathsf b,\eta)\,|\!|\!|h|\!|\!|_{\eta,y}\,\|[\tau_{-\eta e_3}F]\|\,\|[\tau_{-\eta e_3}F']\|, \] in which $\theta$ is the map $x_3\mapsto-x_3$; the bound holds in the restricted form, that is, at a fixed label $\mathsf b\ge x_s$, or with $\mathsf b_j\ge x_s$ for infinitely many $j$. At every $\mathsf b>0$ a bound of the same form holds on a thinner slab, with its numbers at the dilated data, by the last bullet of \Cref{Y-P3-ThG}\,\textup{(G6)}. On the whole slab, \textup{(G6+)} gives bounds in two cases: in \textup{(m)}, for labels of the restricted form and any $\eta>0$, a bound of the same form with the numbers of \textup{(G6+)}\,\textup{(m)}; in \textup{(n)}, for every $\mathsf b>0$ and $\eta\in(0,1]$, a bound with its numbers at the dilated data.

The centered functional $\sigma_S$ has \textup{(R0)} to \textup{(R3)}, \textup{(R5)} and \textup{(R7)} of \Cref{support:32-25}. For $S$ the following three are equivalent: the symmetry group $\operatorname{O}(3)$; reflection positivity at the one mirror $((e_1+e_3)/\sqrt2,0)$; and \textup{(R4)} together with \textup{(R6)} of \ref{support:32-25}.

\begin{proof} The assertions of \textup{(Y4)} from reality up to the symmetry group are \Cref{Y-P2-ThR}\,\textup{(R6)}\,\textup{(i)} to \textup{(vii)}, with \Cref{Y-P2-L13}. The objects $\mathcal H(S)$, $\Omega$, $e^{-tH}$, the momenta, $W_2$ and $[X]$ come from \Cref{Y-P2-ThR}\,\textup{(R7)}, with \Cref{Y-P2-L14}. The bounds with $A_6$, the growth condition, the sharper bound and the properties on $\mathcal S_{\ne}$ are \Cref{Y-P3-ThG}\,\textup{(G2)}\,\textup{(iii)}, \textup{(G4)}, \textup{(G5)}, proved in Lems.~\ref{Y-P3-L4} to~\ref{Y-P3-L6}. The bound on a field between two vectors is \Cref{Y-P3-ThG}\,\textup{(G6)}, \textup{(G6+)}, with \Cref{Y-P3-L7}: the main clause of \textup{(G6)} gives it in the restricted form, the last bullet of \textup{(G6)} gives its form at the dilated data for every $\mathsf b>0$, and \textup{(G6+)} gives it on the whole slab. The properties of \Cref{support:32-25} and the equivalence of the three conditions are \Cref{Y-P4-Th4}\,\textup{(4.9)}\,\textup{(ii)}, proved in Step~8 of its proof. \end{proof}

\textup{(Y5)} \emph{Every $S\in\mathrm{Lim}_G(\mathsf b)$.} In \textup{(Y5)} the label is in normal form, $\mathsf b>x_s$, and every other $\mathsf b>0$ is treated in \textup{(Y8)}. On the lattice the insertion bound holds with $2\,\Gamma(\mathsf b',\mathsf l)\|f\|_\infty$, and the plates are $\le C_\varphi t^{-7/2}$. The functional $\sigma_S$ has \textup{(R0)} to \textup{(R7)} of \Cref{support:32-25}, with $C^\sigma_n$, $(c_{\rm I},C_{\rm I},\mathsf p_{\rm I})$ and $\mathsf a_\Gamma=7/2<4$. The family $S$ is invariant under $\mathbb R^3\rtimes \operatorname{O}(3)$, and $\sigma_S$ is reflection positive at every plane. Moreover $H\ge|P_k|$ for $k=1,2$, the joint spectrum lies in $\{E\ge|q|\}$, and $\operatorname{O}(2)$ acts on $\mathcal H(\sigma_S)$.

\begin{proof} Clause \textup{(Y5)} is \Cref{Y-P4-Th4}\,\textup{(4.4)}, \textup{(4.5)}, \textup{(4.6)}, \textup{(4.7)}, proved in Steps~4 to~7 of its proof. \end{proof}

\textup{(Y6)} \emph{Components of $S\in\mathrm{Lim}_G(\mathsf b)$.} In \textup{(Y6)} the label is in normal form, $\mathsf b>x_s$, and every other $\mathsf b>0$ is treated in \textup{(Y8)}. \textup{(E1)} Every sequence $(K_j,\nu_j,\mathsf b_j)$ with $K_j\to\infty$, $\nu_j\in\mathbb G_{K_j}(\mathsf b_j)$ and $\mathsf b_j\to \mathsf b$ has a convergent subsequence. A member $S$ of $\mathrm{Lim}_G(\mathsf b)$ is obtained along such a sequence if and only if $(K_j,\nu_j)$, which by \Cref{support:35-5} is a frame in the class system $(K_j,\mathbb G_{K_j}(\mathsf b_j))$, converges in the sense of \eqref{support:35-6}. \textup{(E2)} The canonical decomposition exists and is unique up to one isomorphism; $\ker H=L^2(\mathsf X,\lambda)$ and $\#\mathsf X=\dim\ker H$ when $\mathsf X$ is finite, and $\ker H$ is infinite-dimensional when $\mathsf X$ is infinite.

\textup{(E3)} Every limit $\sigma'$ in the frame, hence every $\sigma_\xi$, coincides, up to the shift $\mathrm{sh}_c$, with the centred functional of some member of $\mathrm{Lim}_G(\mathsf b)$. It is invariant under all isometries and reflection positive at every mirror, it has order $\le4n+1$, its first member satisfies $S_1=c\int f$ with $|c|\le\bar c_*$, it obeys the bound of line 3 of \textup{(Y2)}, and it satisfies the growth condition \textup{\ref{support:26-33}}, with $(13,1,C_E(\mathsf b),13)$ as its quadruple. Moreover $\ker H_\xi=\mathbb C\,\Omega$, and $\sigma_\xi$ clusters.

\textup{(E4)} Every $\sigma_\xi$ satisfies \textup{(E0$'$)}, \textup{(E1)} for $\mathbb R^3\rtimes \operatorname{O}(3)$, \textup{(E2)}, \textup{(E3)} and \textup{(E4)} of \textup{\ref{support:28-35}}, as well as \textup{(PtB)} and \textup{(Rec1)} to \textup{(Rec5)}. Its point functions are real-analytic, with $|\mathsf S_n^\xi(x)|\le[C_*\max\{1,(2n^2/d(x))^5\}]^n$. It yields a Wightman theory of a single Hermitian scalar field on $\mathbb R^{1+2}$.

\begin{proof} Clause \textup{(Y6)} is \Cref{Y-P5-ThE}\,\textup{(E1)} to \textup{(E4)} of \S\ref{Y-P5-1}. \end{proof}

\textup{(Y7)} \emph{Non-triviality on tori and for every limit of \textup{(R5)}.} \textup{(D)} The number $x^\#$ is finite, $[j_*(\mathsf b),\infty)$ lies in $J(\mathsf b)$, and $j_*$ is bounded on compact sets. The clauses \textup{(N0)} to \textup{(N3)} hold for every $j\in J(\mathsf b)$, hence for all $j\ge j_*(\mathsf b)$, for all $K\ge j+1$, and for all $N$ of the data.

\textup{(N0)} $\mathrm{Cov}(G_-,\mathcal O_Kf)^2\le\mathrm{Var}(G_-)\,S^T_2(\theta f,f)$ for each real bounded $f$ with $f(x(p))=0$ whenever some site of the plaquette $p$ lies outside the region $\frac12<y_1<\frac12+NL^j/2$; here the unit of $y$ is that of the level $k=K-j$, namely $L^{-j}$, and $y_1$ is represented in the fundamental domain $\frac12-NL^j/2<y_1\le\frac12+NL^j/2$, the unit and the domain being as in ``Geometry'' and \textup{(N0)} of \Cref{Y-P6-N}. Throughout \textup{(Y7)}, $\theta$ is the reflection $y_1\to1-y_1$ in units of the level $k$, and $G_-$ and $f_*$ are those of ``Geometry'' of \Cref{Y-P6-N}. \textup{(N1)} $\mathsf b\,\mathrm{Cov}(G_-,\mathcal O_Kf_*)\ge c_F$. \textup{(N2)} $\mathrm{Var}(G_-)\le7$. \textup{(N3)} $\mathsf b^2S^T_{2;K,N,L^K\mathsf b}(\theta f_*,f_*)\ge c_*$, where $c_*$ is a pure number.

\textup{(N4)} The same bound holds for $S_2^{(N,\mathsf b)}$ of \Cref{Y-H9-T}\,\textup{(T3)}; it holds for each limit of \textup{(R5)}\,\textup{(a)}, \textup{(b)}, \textup{(c)} taken at the fixed label $\mathsf b$; and it holds for \textup{(R5)}\,\textup{(a)}, \textup{(c)} along labels $\mathsf b_n\to \mathsf b$ with $\mathsf b_n\ge x_s$, at every $j\ge j^+(\mathsf b)$.

\textup{(N5)} Let $\tau_s$ be the translation by the distance $s$ perpendicular to the mirror of $\theta$, towards the side of that mirror on which $f_*$ is supported. Every limit $S$ of \textup{(R5)} satisfies $S_2(\theta\tau_sf_*,\tau_sf_*)>0$ for all $s\ge0$ and has $\dim \mathcal H(S)\ge2$; for a limit along labels $\mathsf b_n\to\mathsf b$, in place of a fixed label, this holds at the $j$ of \textup{(N4)} or at the $i$ of \textup{(N6$'$)}. \textup{(N6$'$)} Let $f_0$ and the mirror $\theta_i$ be those of \Cref{Y-P6-N}. For every $\mathsf b>0$ and every integer $i$ with $L^i\mathsf b\ge x^{\#\#}$, the inequality $\mathsf b^2S_2(\theta_if_0(L^i\cdot),f_0(L^i\cdot))\ge c_*$ holds for each limit of \textup{(R5)} taken at the fixed label $\mathsf b$; on a torus at the label $\mathsf b$ it holds provided that $K\ge i+1$ and that, for some integer $z$ with $L^z\mathsf b\in[x_s,L^2x_s]$, the side of the torus meets the conditions that \Cref{Y-P2-ThR}\,\textup{(R4)} imposes on a side at that $z$. With $Lx^{\#\#}$ in place of $x^{\#\#}$, limits along labels $\mathsf b_n\to\mathsf b$ satisfy it as well.

\begin{proof} Clause \textup{(Y7)} is \Cref{Y-P6-N}, proved with Lems.~\ref{Y-P6-L1} to~\ref{Y-P6-L16}; its part \textup{(D)} is \Cref{Y-P6-L16}. \end{proof}

\textup{(Y7$'$)} \emph{Every periodic state, every component.} Throughout \textup{(Y7$'$)}, $g_f$, $V$, $P_\rho$, $\theta_*$, $G_-$, $f_*$ and $s_\xi$ are those of Thms.~\ref{Y-P5-ThE}, \ref{Y-P5-ThP} and~\ref{Y-P5-ThC}; $P_0$ is the projection on $\ker H$; and $\theta$ is the map $x_3\mapsto-x_3$, as in \textup{(Y4)}, not the $\theta$ of \textup{(Y7)}. \Cref{Y-P5-ThE}\,\textup{(E7)} gives, for every $f\in C^1_c(\{x_3>0\})$, the identity $S_2(\theta f,f)=\int g_f\,d\lambda+V(\int f)^2$, in which $V$ is a number. For a given $\xi$, it also states four conditions on the component $\sigma_\xi$ that are equivalent to one another. A component is trivial if it satisfies them, and non-trivial if it does not.

Let $\mathsf b'\ge x_s$, $j\ge j_*(\mathsf b')$ and $K\ge j+1$; these are the data of \Cref{Y-P5-ThP}. For every $\rho\in\mathbb G_K(\mathsf b')$ that theorem gives the following.
\begin{enumerate}
\item[\textup{(P0)}] In product form, $\mathrm{Cov}_{P_\rho}(G_-,\mathcal O_Kf)^2\le\mathrm{Var}_{P_\rho}(G_-)\,S^T_2[\rho](\theta_*f,f)$ for every real bounded local $f$ carried by the open right half.
\item[\textup{(P1)}] $\mathsf b'\,\mathrm{Cov}_{P_\rho}(G_-,\mathcal O_Kf_*)\ge c_F$.
\item[\textup{(P2)}] $\mathrm{Var}_{P_\rho}(G_-)\le7$.
\item[\textup{(P3)}] $\mathsf b'^2S^T_2[\rho](\theta_*f_*,f_*)\ge c_*$.
\item[\textup{(P4)}] Conditionally on the invariant sets, the same holds.
\end{enumerate}

In normal form, $\mathsf b>x_s$, \Cref{Y-P5-ThC} gives the following for every $\xi$. At every other $\mathsf b>0$, the same conclusions hold with the test and radius of its clause \textup{(C4)}, as stated in \textup{(Y8)}.
\begin{enumerate}
\item[\textup{(C1)}] For every $j\ge j^+(\mathsf b)$, with the mirror carried to $(e_3,0)$, we have $\mathsf b^2g_{f_*}(\xi)\ge c_*$.
\item[\textup{(C2)}] On $0<r\le r_*(\mathsf b)$ we have $s_\xi(r)\ge c_{**}\mathsf b^{-2}r^{-6}$; moreover $H_\xi$ is unbounded, and $\dim \mathcal H(\sigma_\xi)$ is infinite.
\item[\textup{(C3)}] $\|(1-P_0)[(f_*)]\|^2\ge c_*/\mathsf b^2$.
\end{enumerate}

\begin{proof}
Clause \textup{(Y7$'$)} follows from \Cref{Y-P5-ThE}\,\textup{(E7)}, from Thms.~\ref{Y-P5-ThP} and~\ref{Y-P5-ThC}, and from Steps~9 to~13 of \S\ref{Y-P5-1}.
\end{proof}

\textup{(Y8)} \emph{Every $\mathsf b>0$ and every side $2L^m$.} By the dilation $L^z$, the following hold, in their dilated forms, at every $\mathsf b>0$ and every side $2L^m$: the two bounds of lines 1 and 2 of \textup{(Y2)}, on $\sigma_{K,N}(f)$ and on $S^T_n(f)$; lines 3 and 4 of \textup{(Y2)}; the clauses \textup{(Y4)}, \textup{(Y5)} and \textup{(Y6)}; the parts \textup{(N3)}, \textup{(N4)} and \textup{(N5)} of \textup{(Y7)}; and clause \textup{(Y7$'$)}. The exact finite-torus dilation allows every integer $z$ satisfying the four conditions of \Cref{Y-P2-ThR}\,\textup{(R4)}, including $K-z\ge1$. The simplified constants in \Cref{Y-P3-ThG}\,\textup{(G2)(iv)} use $z\ge0$. For finite-cutoff Gibbs-state bounds, take $K\ge z+1$ as well. For the clauses on $\mathrm{Lim}_G(\mathsf b)$, frames and components, take $z=z(\mathsf b)$, so that $L^z\mathsf b>x_s$, and use the constants of \Cref{Y-P5-ThE}\,\textup{(E5)} and the tests and radius of \Cref{Y-P5-ThC}\,\textup{(C4)}. These choices apply to fixed and convergent labels, including $\mathsf b=x_s$, and give \textup{(Y5)}, \textup{(Y6)} and \textup{(Y7$'$)} at every $\mathsf b>0$.

\begin{proof}
Clause \textup{(Y8)} follows from \Cref{Y-P2-ThR}\,\textup{(R4)} and \textup{(R6)\,(vi)}; from \Cref{Y-P3-ThG}\,\textup{(G2)\,(iv)} and \textup{(G4)\,(e)}, from the last bullet of \textup{(G6)} and from \textup{(G6+)\,(n)}; from \Cref{Y-P4-Th4}\,\textup{(4.8)} together with \Cref{Y-P4-LZ}; from \Cref{Y-P5-ThE}\,\textup{(E5)} and \Cref{Y-P5-ThC}\,\textup{(C4)}; and from \Cref{Y-P6-N}\,\textup{(N6)} and \textup{(N6$'$)}.
\end{proof}

\textup{(Y9)} \emph{Block fields at every depth.} Let $s\ge0$ be any depth and let $K\ge s+1$. The block fields at depth $s$ have the following properties.
\begin{enumerate}
\item[\textup{(a)}] Locality, conditional independence and covariance hold.
\item[\textup{(b)}] Reflection positivity holds without a collar at every $c\in L^{-s}(\mathbb Z+\tfrac12)$.
\item[\textup{(c)}] First half: along a further subsequence of sides, the laws of \Cref{Y-J-CorC}\,\textup{(b)} converge locally in total variation; this half uses \Cref{Y-H9-T}\,\textup{(T7)}, which supplies these laws at each side. Second half: every sequence $(K_j,N_j)$ with $N_j\to\infty$ has a subsequence along which the same holds for all $s\le\liminf K_j-1$.
\item[\textup{(d)}] Mixed bounds hold, both without cells and with cells.
\end{enumerate}

\begin{proof}
Clause \textup{(Y9)} follows from \Cref{Y-P2-ThR}\,\textup{(R8)} and from Lems.~\ref{Y-P2-L5}, \ref{Y-P2-L6}, \ref{Y-P2-L9} and~\ref{Y-P2-L15}.
\end{proof}

The closing corollary is \Cref{Y-P5-Cor}, and \S\ref{Y-P5-1} proves it. It states that $\mathrm{Lim}_G(\mathsf b)$ is not empty for every $\mathsf b>0$, and that every component $\sigma_\xi$ of every member of $\mathrm{Lim}_G(\mathsf b)$, and hence of every limit taken volume first, satisfies \textup{(E0$'$)} to \textup{(E4)} of \textup{\ref{support:28-35}} with full Euclidean invariance and has $S_2^{\sigma_\xi}\ne0$.



\begin{remark}[Two orders of distributions]
The bound $\le3n+1$ of \textup{(Y4)}, on the order of the distribution whose value at $f_1\otimes\cdots\otimes f_n$ is $S_n(f_1,\dots,f_n)$, holds for the limits of \Cref{Y-P2-ThR}\,\textup{(R5)}; it is \Cref{Y-P2-ThR}\,\textup{(R6)\,(v)}. The bound $\le4n+1$ of \textup{(E3)} in \textup{(Y6)}, on the order of the limit $\sigma'$ in the frame, holds for every limit in a frame of $\mathrm{Lim}_G(\mathsf b)$; it is \Cref{Y-P5-ThE}\,\textup{(E3)\,(i)}, which comes from \textup{\ref{support:36-12}~(e)}. These are two upper bounds, on two objects and for two classes.

A limit taken volume first lies in both classes, and so the bound $\le3n+1$ holds for the distribution of its $S_n$.
\end{remark}



\section{Blocking setup and free kernels}\label{sec:blocking-and-free}

\subsection{Blocking geometry and scales}\label{Y-G0-1}

The step takes a density at coupling $x$ on a fine torus to a density
at coupling $x/L$ on a coarse torus. We specify the class of densities,
the geometry of a step and the order of the constants here. \Cref{Y-G-ThA}
shows that one step preserves the class, and \Cref{Y-GB-ThB}
iterates this result from Wilson's weight.

The bounds on the free kernels in \S\ref{Y-A-1} determine the
radii used to localize the nonlinear estimates. The parameters in
\textup{(O0)} to \textup{(O13)} are fixed before the lower bounds on
$x$ in \textup{(A0)} to \textup{(A11)}. We write $c_h$ for the constant
in the exponential decay and $c_h(x)$ for the resulting bound at $x$.
In a tower the coupling is $x_k=L^{-k}\beta$. A subscript on a fixed
threshold, as in $x_9$ of \Cref{Y-H-C9}, is a label for that
threshold.

\subsubsection*{Item 1. Setting}

Throughout, $G=\operatorname{SU}(2)$ and $d=3$.

\begin{definition}[Lattice]
The full definitions are those of items 3.1, 3.3 and 3.5 of \S\ref{Y-E-1}. The radii are those of item~2 below.

The number $L$ is odd, and $\mathsf h_L=(L-1)/2$; further $M=L^{m'}$ with $m'\ge1$; and $N$ is even, with $M\mid N$.

The fine torus is $\mathbb T=\mathbb Z^3/NL\mathbb Z^3$, and the coarse torus is $\mathbb T'=L\mathbb Z^3/NL\mathbb Z^3$.

Fine bonds and fine plaquettes are denoted $b$ and $p$; coarse ones are denoted $c$ and $P$.

For an object $i$, $\pi(i)$ is its base site, and $y(i)$ is the coarse site whose block $B(y)$ contains $\pi(i)$; the distance ${\rm dist}(i,j)=L^{-1}|y(i)-y(j)|_\infty$ is measured in coarse units.

The combs are those of item 3.1 of \S\ref{Y-E-1}. $I^\circ$ is the set of the fine bonds off the combs, of which a block has at most $3L^3$, and $\mathfrak L_{\mathfrak g}$ is the space of $\mathfrak g$-valued functions on $I^\circ$.

We also use $\mathrm{d}$, $Q_1$, $E$, $T(y,y'')$, $K^V$, $\ell^V$, $\mathsf B^V=1+\ell^{V*}\ell^V$, $q=q^V_\ell$ and $\mathsf g^V=\Pi_{\mathfrak L}\mathrm{d}^{V*}\mathsf B^VE$. Here $\Pi_{\mathfrak L}$ is the orthogonal projection from $\Omega^1$ onto the $\mathfrak L$ of \Cref{Y-A-Lifts}.

Norms: $|\cdot|$ is the operator norm of a $2\times2$ matrix, $\|Y\|^2=\frac12\operatorname{tr}Y^*Y$ and $Y\cdot Y'=\frac12\operatorname{tr}YY'$. Then $\|Y\|\le|Y|\le\sqrt2\|Y\|$, and $|Y|=\|Y\|$ on $\mathfrak g$.
\end{definition}

\begin{definition}[Cubes]
At each level the cubes consist of $M^3$ sites of a fixed grid, and a coarse cube is called a cell $\square$. For a cell, $\hat\square$ collects the fine sites that lie in the blocks of the sites of $\square$; for a set $A$ of cells, $\hat A$ collects the fine objects based there.

The grids fit: if the grid of cells has the origin $o'$, then the grid of fine cubes has the origin $o=Lo'-\mathsf h_L\bmod M$.

Then $\hat\square$ is the box of side $LM$ whose lower corner is $L(o'+Mn)-\mathsf h_L=o+LMn$, where $n$ is a point of $\mathbb Z^3$; it is made up of exactly $L^3$ fine cubes.

An object belongs to the cube of its base site, and to the cell of $y(i)$.

Distances between cubes, and between cells, are sup distances, measured in cubes and in cells respectively.

$\mathcal D$ is the family of the non-empty sets of cubes that are connected through faces, edges or corners.

\emph{Count 1.} In a graph of degree at most $26$, at most
$(26e)^{n-1}$ connected sets of $n$ vertices contain a prescribed vertex.
This applies to fine cubes and to coarse cells with the same adjacency.
\begin{proof}[Proof of Count 1]
Order the neighbours of each vertex. A breadth-first tree from the
prescribed vertex encodes a connected set as a rooted subtree of the
tree with $26$ ordered children per vertex. This encoding is injective,
since following the recorded neighbour indices recovers the set.
If $T(z)$ counts these subtrees by size, then $T=z(1+T)^{26}$.
For completeness, substituting $z=u/(1+u)^{26}$ in the formal residue
for $[z^n]T$ gives
\[
[z^n]T=[u^{n-1}](1+u)^{26n-1}(1-25u)
       =\frac1n\binom{26n}{n-1}.
\]
The last number is at most $26^{n-1}n^{n-1}/n!\le(26e)^{n-1}$,
using $\log n!\ge\int_1^n\log t\,\mathrm{d}t$.
\end{proof}

$\mathcal P$ is the family of the non-empty finite sets of cubes.

For $X\in\mathcal P$, ${\rm hull}(X)$ is a member of $\mathcal D$ that contains $X$ with the least number of cubes, chosen by a fixed order. We put $|X|_c=\#{\rm hull}(X)$; for $X\in\mathcal D$, ${\rm hull}(X)=X$.

\emph{Count 2.} Given a cube, at most $(52e)^n$ members of $\mathcal P$ with $|X|_c=n$ contain it.
\begin{proof}[Proof of Count 2]
Such an $X$ is a subset of ${\rm hull}(X)$, and by Count 1 the hull is one of at most $(26e)^{n-1}$ members of $\mathcal D$ with $n$ cubes that contain the given cube; as a set of $n$ cubes has $2^n$ subsets, multiplying the two numbers gives the bound.
\end{proof}

For a set $X$ of cubes, $X^{+j}$ is the set of the cubes at distance at most $j$ from $X$, and for a set $X$ of cells it is the set of the cells at distance at most $j$ from $X$, the distances being the sup distances set above; and we put $X^{\sim}=X^{+1}$.

For a set $A$ of fine cubes, ${\rm cl}(A)$ is the set of the cells $\square$ with $\hat\square\cap A\ne\emptyset$.

The bonds of $A$ are the bonds with both ends in $A$, and the paths and the plaquettes of $A$ are those with all their sites in $A$. This is the convention of item 3.3 of \S\ref{Y-E-1} and of item~C.II~1.4 of \S\ref{Y-C-1}, and the definitions of $\mathcal U_X(x)$, \textup{(ca)}, \textup{(ch)}, \textup{(c$\Xi$)}, as well as \Cref{Y-H-Th53}, are understood in this sense.

A plaquette $p$ is in $A$ when its base site lies in $A$; this is the sense of ``in'' in $1_A$ and in $D(V)\cap X$, and of the words ``holds'' and ``based in''.

We say that a field given on the bonds of $A$ is tame there when $\phi_p\le\varepsilon$ holds at every plaquette of $A$.

The convention on ``of $A$'' and ``in $A$'' is in force throughout Sections~\ref{Y-G0-1}, \ref{Y-G-1}, \ref{Y-GB-1}.
\end{definition}

\begin{definition}[The cells $Y_0(X)$]
For a set $X$ of fine cubes let $n=\#X$ if $X\in\mathcal D$, and $n=|X|_c$ if $X\in\mathcal P$. For $X\in\mathcal D$ both choices give the same $n$, since then ${\rm hull}(X)=X$. We call $X$ small if $n\le L/2-2$. The cubes and the cells are ordered once and for all, and ``first'' refers to these orders.

If $X$ is small, $Y_0(X)$ is the one cell that contains the first cube of $X$, and ${\rm hull}(X)^{\sim}$ lies within $M/2$ coarse units of it.

If $X$ is not small, $Y_0(X)$ is the set of the cells that meet ${\rm hull}(X)^{\sim}$; this set is connected, and $\#Y_0\le8+432n/L$.

In both cases ${\rm cl}(X)\subset Y_0^{+1}$.
\begin{proof}[Proof for small $X$]
The set ${\rm hull}(X)$ is connected, has $n$ cubes and contains the first cube of $X$; hence every cube of ${\rm hull}(X)^{\sim}$ is within $n\le L/2-2$ cubes of that first cube. Put $\square:=Y_0(X)$. Then the sites of ${\rm hull}(X)^{\sim}$ are within $(L/2-2)M$ fine sites of $\hat\square$, and their blocks lie at most $(L/2-2)M/L+1\le M/2$ coarse units away from the sites of $\square$, as $M\ge L$. Since $M/2<M$, the cells that meet ${\rm hull}(X)^{\sim}$, those of ${\rm cl}(X)$ among them, lie in $\square^{+1}$.
\end{proof}
\begin{proof}[Proof for $X$ not small]
The set ${\rm hull}(X)^{\sim}$ is connected and has $m\le27n$ cubes; two neighbouring cubes lie in one cell or in two neighbouring cells, so $Y_0$ is connected. Walking twice along a spanning tree of ${\rm hull}(X)^{\sim}$ we obtain a closed walk through all of its cubes with at most $2m$ steps, which we cut into at most $2m/L+1$ pieces of at most $L$ steps each. Each piece lies in a box of side $L+1$ cubes, while a cell has side $L$ cubes because the grids fit; hence a box meets at most 2 cells along each coordinate direction, so at most 8 cells altogether, and $\#Y_0\le8(2m/L+1)\le8+432n/L$. Finally ${\rm cl}(X)\subset Y_0$, since every cell of ${\rm cl}(X)$ meets $X\subset{\rm hull}(X)^{\sim}$.
\end{proof}
\end{definition}

\begin{definition}[At $x$]
The coupling of the step is $x$, with $x'=x/L$, and a prime marks a quantity evaluated at $x'$.

We set $p_0=A_0(\log x)^{\mathfrak p_0}$ and $\varepsilon=p_0x^{-1/2}$; then $\varepsilon'\le L^{1/2}\varepsilon$.

We set $R=\varrho(x)=\lceil c_R\log x\rceil$, $r(x)=\lceil20c_R\log x/M\rceil$, $s(x)=8r(x)$ and $\ell_w(x)=6Mr(x)+12$.

For sets of fine cubes, $X^{\approx}=X^{+r(x)}$ and $X^{\approx\approx}=X^{+s(x)}$; we write $r'=r(x')$ and $s'=s(x')$; for sets of cells, $Y^{\approx}=Y^{+r'}$; and $v=(LM)^3$.

The caps are $c_a(x)=E_ap_0^3x^{-1/2}$ and $c_h(x)=e^{-c_hp_0^2}$.
\end{definition}

\begin{definition}[Fields, step, pair, energy]
\emph{Fields.} We set $\phi_p=\operatorname{ang}\partial U(p)$, $F_p=\mathsf L(\partial U(p))$ and $D(U)=\{\phi_p>\varepsilon\}$, and on the coarse lattice we use the angles $\phi'_P=\operatorname{ang}\partial V'(P)$ and $D'=\{\phi'_P>\varepsilon'\}$. Here $\operatorname{ang}$ and $\mathsf L$, with $\mathsf L(-1)=0$, are those of item~C.I~4 of \S\ref{Y-C-1}, and $F_p$, $\phi_p$, $D(U)$ are as in item 3.5 of \S\ref{Y-E-1}.

\emph{Step.} The step is given by \[\rho'(V')=\int\rho(U)\prod_ck_{ax}(V'(c),\mathsf M(U)(c))\,dU.\] Here $k_{ax}$ is the kernel of the step and $\mathsf M(U)$ is the average, both as in item 3.4 of \S\ref{Y-E-1}.

\emph{Pair.} The second entry of a pair is $e\in\mathfrak E=\prod_PS^2$, where $S^2$ is the unit sphere of $(\mathfrak g,\|\cdot\|)$, one copy for each coarse plaquette $P$; see also item 3.7 of \S\ref{Y-E-1} and E2 of \S\ref{Y-D-1}.

On each of these spheres we take the probability measure invariant under rotations, and $de$ denotes the product of these measures over all $P$.

Property 1: $\int_{\mathfrak E}de=1$. Property 2: under $de$ the $e_P$ are independent.

Property 3: for every $G$-valued $u$, the map $e\mapsto\operatorname{Ad}_ue$, $(\operatorname{Ad}_ue)_P=\operatorname{Ad}_{u(y_P)}e_P$, where $y_P$ is the base site of $P$, preserves $de$, since every $\operatorname{Ad}_{u(y_P)}$ is a rotation of the three-dimensional Euclidean space $(\mathfrak g,\|\cdot\|)$.

The mean over $e$ is the integral against $de$.

We put $X'_P(V',e)=\mathsf L(\partial V'(P))+\pi e_P1\{\partial V'(P)=-1\}$. At every $P$ one has $|X'_P|=\phi'_P$, and $e_P$ enters only at $P\in D'$.

We write $1_A(V)=1\{\phi_p\le\varepsilon\text{ for all }p\text{ in }A\}$, and $V$ is tame on $A$ when this indicator is $1$.

\emph{Energy.} Following E2 of \S\ref{Y-F-1}, \[\mathrm{En}_\ell(U)=\sum_pS_\varepsilon(\phi_p)+\frac12\sum_ym_y(U)|\ell^U_y\zeta|^2.\]

The energy $\mathrm{En}'(V')$ is the same expression at $x'$, defined with the rows $\ell'^{\rm c}$, the cut of $\ell'$ at $\varrho(x')$.
\end{definition}

\begin{definition}[Domains]
We let $\mathcal U_X(x)$ be the set of the $\operatorname{SL}(2,\mathbb C)$ fields $V$ on the bonds of $X^{\approx}$, with $V(\bar b)=V(b)^{-1}$, that satisfy \textup{(u1)}, \textup{(u2)} and \textup{(u3)}:
\begin{enumerate}
\item[\textup{(u1)}] $|V(b)|<1+\vartheta(x)$, where $\vartheta(x)=\vartheta_0^\flat(1-1/\log x)$;

\item[\textup{(u2)}] $|V(\pi)|<2-1/\log x$ for the paths $\pi$ of $X^{\approx}$ with at most $\ell_w(x)$ bonds;

\item[\textup{(u3)}] $|\partial V(p)-1|<2\varepsilon$ at every plaquette $p$ of $X^{\approx}$.
\end{enumerate}

The domain $U^\tau(B;x')$ is that of item 3.3 of \S\ref{Y-E-1}, given by \textup{(e1)}, \textup{(e2)}, \textup{(e3)}.

$\mathfrak V$ is the set of the $\operatorname{SL}(2,\mathbb C)$ fields on $\mathbb T'$ with $|V'(c)|\le1+\vartheta_0^\flat$, and $\mathfrak V^\circ$ is its interior.

We extend any field that is given on a set of bonds by setting it equal to $1$ on all other bonds.

Convention. Let a complex field be given on the bonds of a set $B$, and let it satisfy \textup{(u3)} or \textup{(e3)}. We then put $X'_P:=-i\log\partial V'(P)$ for those $P$ whose four sites all lie in $B$, and $X'_P:=0$ for every other $P$; this is the convention of item 3.7 of \S\ref{Y-E-1}.

The convention does no harm. Indeed, $a_{0,b}$ depends on $X'$ within $R_3$ of $y(b)$, by the last clause of \textup{($\Gamma3$)} of Step~0 of the proof of \Cref{Y-G-ThA}; and a letter with $J_\alpha\subset Y$ depends on the field in $Y^{+(r'-2)}$, by item~5 of Step~7 of that proof. Hence none of the values in use depends on the convention.
\end{definition}

\begin{definition}[Members]
For $A\ge1$, $M_A^{\rm su}(x)$ is the set of the triples $u=(E^a,E^h,w)$ that satisfy \textup{(ca$_A$)}, \textup{(ch$_A$)} and \textup{(c$\Xi_A$)}:
\begin{enumerate}
\item[\textup{(ca$_A$)}] for $X\in\mathcal D$, the function $E^a(X,\cdot)$ is holomorphic on $\mathcal U_X(x)$, is invariant under $G$-valued gauge transformations, and vanishes at the pure gauges, by which we mean here and below the fields $1^u$ with $u$ $G$-valued; and $\sup|E^a(X,\cdot)|\le Ac_a(x)e^{-\kappa(\#X-1)}$, the supremum being taken over $\mathcal U_X(x)$;

\item[\textup{(ch$_A$)}] $E^h(X,\cdot)$ is a complex function, Borel and gauge invariant, of those $G$-valued fields on the bonds of $X^{\approx}$ which are tame there, and $\sup|E^h(X,\cdot)|\le Ac_h(x)e^{-\kappa(\#X-1)}$, where the supremum runs over these fields;

\item[\textup{(c$\Xi_A$)}] for $X\in\mathcal P$, $w(X,\cdot)$ is a complex function of $V$ on the bonds of $X^{\approx\approx}$, Borel and gauge invariant, and $|w(X,V)|\le Ae^{-\kappa(|X|_c-1)}$ at every $V$; this bound is the weak cap.
\end{enumerate}

The strong cap is the bound \[|w(X,V)|\le\exp\{-c_\Xi x\varepsilon\sum_{D(V)\cap X}\phi_p-\kappa(|X|_c-1)\}.\]

We call $u$ real if $E^a$ is real at $G$-valued fields and $E^h$, $w$ are real.

The density of $u$ with the rows $\ell$ is \[\rho[u;\ell]=\exp\{-x\mathrm{En}_\ell+\sum_{X\in\mathcal D}1_{X^{\approx}}[E^a+E^h](X,\cdot)\}\,\Xi_w,\] where $\Xi_w(V)=\sum_{\{X_i\}}\prod_iw(X_i,V)$; the sum runs over the finite families $\{X_i\}$ of sets in $\mathcal P$ at pairwise distance $>s(x)$ such that each $X_i$ holds a plaquette of $D(V)$ and $D(V)$ lies inside their union. The empty family, with the product $1$, is admitted if and only if $D(V)=\emptyset$.
\end{definition}

\begin{definition}[The class]
The class of densities at $x$ with the rows $\ell$ is \[C^{\rm su}_\flat(x;\ell)=\{e^{\mathsf V}\rho[u;\ell]>0:\ \mathsf V\text{ real},\ u\in M_1^{\rm su}(x)\text{ real, strong cap}\}.\]

The clause ``$\mathsf V$ real'' is called \textup{(c0)}; see item~6 of Step~13 of the proof of \Cref{Y-G-ThA}.
\end{definition}

\begin{definition}[The data \textup{(SD)}]
The data \textup{(SD)} of a step at $x$ consist of the side $N$, the grids and the rows $\ell$, subject to the following conditions. First, $N$ is even, $M\mid N$, \textup{(Tor)} $N\ge2Ms(x)$ holds, and the grids fit.

The rows of \textup{(SD)} are $\ell\in \mathcal R(b_*,K_*,\lambda_*,\frac19)$ on $\mathbb T$, of range $\le\varrho(x)$.

Empty rows are admitted. The family $\ell=\emptyset$, with $\mathsf H_y=\{0\}$ for all $y$, lies in $\mathcal R(1,0,\lambda_*,0)\subset \mathcal R(b_*,K_*,\lambda_*,\frac19)$, by \Cref{Y-A-SetsOfRows}, with its display and its last sentence; its range is $0$.

At $\ell=\emptyset$ we have $\mathrm{En}_\ell=\sum_pS_\varepsilon(\phi_p)$, $\mathsf B^V=1$ and $q=|\mathrm{d}^V\cdot|^2+a|Q_1^V\cdot|^2$; and the numbers $b$, $\mathsf N_{\mathsf B}$, $\gamma$, $\varrho_L$ and all the radii keep the values that belong to $(K_*,\lambda_*,\varrho(x))$.

\textup{(SD)} involves no level and no depth.

The last clause of \textup{(SD)} is a condition on the rows: the pair $(\ell,\ell')$, where $\ell'$ denotes the rows built from $\ell$ by \Cref{Y-A-P5} (with $\mathsf B'=1+\ell'^{\mathsf T}\ell'$, tensored with $1_{\mathfrak g}$; at $\ell=\emptyset$ the group \textup{(5)} of the list of that proposition is void), satisfies the four conclusions of \Cref{Y-A-Th36} listed below.
\begin{enumerate}
\item[\textup{(1)}] \textup{(Dec)}$(C_D,\theta_D)$;

\item[\textup{(2)}] \textup{(N$'$)} with \eqref{eq:row-cut-norms} at every cut;

\item[\textup{(3)}] $\mathsf t_\rho\le\frac18(\frac12-d_0-c_\Xi)$;

\item[\textup{(4)}] \textup{(L1)}, \textup{(L2)}, taken as \textup{(L1$^{\rm su}$)}, \textup{(L2$^{\rm su}$)}.
\end{enumerate}

Here $(C_D,\theta_D)$ are those of \textup{(O6)}.

The number \[\mathsf t_\rho:=\|\ell'^{\rm c}\|_{\rm r}\sup_{P_0}\sum_{y':\,{\rm dist}(y',P_0)>\rho_0}|\ell'^{\rm c}_{y'}(P_0)|,\] as in E5 of \S\ref{Y-F-1}, is a quantity of the rows.

The part ``Conclusion'' of \Cref{Y-A-Fr} contains \textup{(Fr0)}, \textup{(Fr2)}, \textup{(Fr3$'$)} and \textup{(Fr4)}. At the $(C_D,\theta_D)$ of \textup{(O6)} and for the pair $(\ell,\ell')$ of that theorem, the conclusions \textup{(1)}, \textup{(2)}, \textup{(3)} are \textup{(Fr2)} for the sixteen kernels, \textup{(Fr3$'$)} and \textup{(Fr4)}, respectively, and \textup{(Fr0)} gives the two conditions of \textup{(SD)} on the rows, that $\ell\in\mathcal R(b_*,K_*,\lambda_*,\frac19)$ and that $\ell$ is of range $\le\varrho(x)$.

The four conclusions of \textup{(SD)} imply \textup{(Fr2)}, \textup{(Fr3)} of item~C.II~2.2 of \S\ref{Y-C-1}, and \textup{(Fr1)} is \Cref{Y-A-L7}\,(c), which is algebra and holds for all rows.
\end{definition}

\subsubsection*{Item 2. The radii}

\begin{definition}[The radii]
The radii set here are not letters of the class $C^{\rm su}_\flat$. In coarse units we set \[\varrho_L=\lceil R/L\rceil,\quad r_w=\lfloor R/2\rfloor,\quad r_G=r_w+\varrho_L+1,\quad R_2=2R+2,\quad R_3=2R+\varrho_L+5,\] \[R_\chi=R+\varrho_L+2,\quad R_v=2R+2\varrho_L+7,\quad R_6=4R+6\varrho_L+11,\] \[D_w=\lceil\ell_w(x)/L\rceil+1,\quad\omega=M+\lceil Mr(x)/L\rceil+D_w+R_3+2.\]

In cells we set $r_l=\lceil(\varrho_L+2)/M\rceil$, $r_1=\lceil R_2/M\rceil+3r_l+5$ (the $r_1$ of E8 of \S\ref{Y-F-1}), $s_0=2r'$ and $\varrho_*=2r_1+\lceil s(x)/L\rceil+2$.

Identically, $r_1\ge9$ and $r_1M\ge2R+3\varrho_L+8+5M$; hence $r_1M\ge R_v+1$, $r_1M\ge R_\chi$ and $r_1M+1>2(R_\chi+1)$; and $\lceil R_2/M\rceil\le r_1-8$.
\begin{proof}
Since $\lceil R_2/M\rceil\ge1$ and $r_l\ge1$, we find $r_1\ge9$ and $\lceil R_2/M\rceil=r_1-3r_l-5\le r_1-8$. Since $M\lceil R_2/M\rceil\ge R_2$ and $Mr_l\ge\varrho_L+2$, we find $r_1M\ge R_2+3(\varrho_L+2)+5M=2R+3\varrho_L+8+5M$. Finally $R_v+1=2R+2\varrho_L+8$ and $2(R_\chi+1)=2R+2\varrho_L+6$, and $M\ge1$; together with the last bound these give the three inequalities for $r_1M$.
\end{proof}
\end{definition}

\subsubsection*{Item 3. The letters, in the order of the table}

\begin{definition}[The table]
We choose the constants of Sections~\ref{Y-A-1} to~\ref{Y-J-1} in the order in which they are listed below. Wherever one constant is subject to two conditions, we impose the stricter of them: at \textup{(O11)} it is the smaller bound, at \textup{(O12)} and \textup{(O13)} the larger value.

The entries are as follows.
\begin{enumerate}
\item[\textup{(O0)}] The pure numbers $b_*,\gamma_*,\lambda_*,K_*,L_A,L_B,c_1^*,c_2^*,\lambda_2^*,\theta_J$, and $\mathsf g_*=\mathsf g(\lambda_*)$. All of them but $L_B$ are those of \S\ref{Y-A-1}; $L_B$ is a pure number set in \S\ref{Y-I-1}, \textup{(O0)}, and it enters here only through $L\ge L_B$ in \textup{(O4)}.

\item[\textup{(O1)}] $d_0\in(0,\frac12)$.

\item[\textup{(O2)}] $0<c_\Xi\le\frac12(\frac12-d_0)$.

\item[\textup{(O3)}] The integer $\rho_0$, fixed by \eqref{eq:row-mask-radius} of \S\ref{Y-A-1}.

Then the following functions of $L$: $\mathsf l_L=1+\log L$, $\mathsf N_Z=\mathsf N_{\mathsf B}=1+3K_*^2\mathsf g_*^2$, $\mathsf P_\rho=1+3(2+3(2\rho_0+1)^3)K_*^2\mathsf g_*^2$ (this is \eqref{eq:row-cut-norms} of \S\ref{Y-A-1}), and \[\mathsf N_\sharp=c_1^*\mathsf l_L\mathsf N_Z\mathsf P_\rho,\quad \mathsf N_\Xi=c_2^*\mathsf N_Z,\quad \mathsf N_\nu=\mathsf N_\sharp+(3+4c_2^*L)\mathsf N_Z,\quad \mathsf N_\sigma=3K_*\mathsf g_*c_1^*\mathsf l_L\mathsf N_Z,\] \[\mathsf N_\nu^\flat=\mathsf N_Z+2L\mathsf N_\Xi+(3+3(2\rho_0+1)^3)K_*\mathsf g_*\mathsf N_\sigma,\qquad \mathsf N_{\rm c}=17(\mathsf N_\sharp+1).\]

\item[\textup{(O4)}] $L$ is odd and $L\ge L_*^{\rm su}$, where $L_*^{\rm su}$ is the least integer $\ge\max\{L_A,L_B\}$ such that every $L\ge L_*^{\rm su}$ satisfies the three conditions below: \begin{equation}\tag{L1$^{\rm su}$} [\mathsf N_\nu/L+2200\mathsf N_Z+40\mathsf N_\Xi]L^{-1/2}\le d_0/4,\end{equation} \begin{equation}\tag{L1$^\flat$} [\mathsf N_\nu^\flat/L+2200\mathsf N_Z+40\mathsf N_\Xi]L^{-1/2}\le\frac3{16}d_0,\end{equation} \begin{equation}\tag{L2$^{\rm su}$} L^{1/2}\ge1.2\cdot10^4\mathsf N_{\rm c}^2.\end{equation}

Further, $c_L=(L^2-1)/(6L^2)$ and $a\ge3/(c_LL)$, so that $aL\ge18$.

Of these three conditions, the proof of \Cref{Y-G-ThA} itself depends on \textup{(L2$^{\rm su}$)} alone. It gives $L\ge1.2\cdot10^{13}$, as $\mathsf N_{\rm c}\ge17$, and $T:=L^{3/2}/(19\mathsf N_{\rm c})>1$.

Indeed, $T\ge(1.2\cdot10^4\mathsf N_{\rm c}^2)^3/(19\mathsf N_{\rm c})\ge1.7\cdot10^{12}\cdot17^5/19>10^{17}$, so that $T\ge2$ as well.

\item[\textup{(O5)}] $\kappa\ge6$, and $\kappa^\sharp=10(\kappa+4)$.

\item[\textup{(O6)}] Formed from $(L,a,K_*,\lambda_*)$ as in item~C.I~1.2 of \S\ref{Y-C-1}: $\gamma=(73L^4)^{-1}$, $b=16\mathsf N_{\mathsf B}+a/L$, $A^\flat$, $\theta^\flat$, $\theta_{\mathbb C}=\theta^\flat/6$, $\vartheta_0^\flat=\theta^\flat/24$, $\theta_2$, $\Lambda=2+A^\flat/b$, $\mu=\gamma\theta_2/(4b\log\Lambda)$, and $\mathsf G_0\le C_\varphi^\flat$.

Next, $\theta_D=\frac12\min\{\lambda_2^*,\theta_{\mathbb C},\theta_J\}$.

The number $C_D$ is at least the expression that \Cref{Y-A-Fr} gives for $C_D$ in item \textup{(6)} of its part ``The order of the constants'', evaluated at this $\theta_D$; that expression is $\ge A_q\ge32e^2$, so $C_D\ge4$.

The number $C_D$ belongs to the table and is fixed here, before any rows are chosen: for every $C$ at least that expression and every $\ell\in \mathcal R(b_*,K_*,\lambda_*,\frac19)$, \textup{(Dec)}$(C,\theta_D)$ holds, by Step~H13 of the proof of \Cref{Y-A-Fr}, which is \Cref{Y-A-P35} with $\theta_{\mathbb C}$ in place of $\theta_q$. Every later constant rises with $C_D$, by \textup{(Fr2)} of item~C.II~2.2 of \S\ref{Y-C-1}.

Then $c_R=64/\theta_D$, so that $c_R\theta_{\mathbb C}\ge128$.

The condition \textup{(R$^\flat$)} $c_R\mu\ge1.35$ follows, by a chain of four links. First, $\lambda_*L\ge12$ by $L\ge L_A$, see item~C.II~5 of \S\ref{Y-C-1}, under ``Conditions on the letters''; hence $\theta_2=\frac12$ and $\theta^\flat=\gamma/(2A^\flat)$.

Second, $c_R\ge768/\theta^\flat=1536A^\flat/\gamma$, and $c_R\mu\ge192\,t/\log(2+t)$ with $t=A^\flat/b\ge14$.

Third, $A^\flat\ge14b$, by the ``Conditions on the letters'' in item~C.II~5 of \S\ref{Y-C-1}: there $2e^2\ge14$ and $384e^8\mathsf g_2(\lambda/2)^2\ge14\cdot48\mathsf g(\lambda)^2$, and a quotient of two sums is at least the least quotient of their terms.

Fourth, $t/\log(2+t)$ increases with $t$; hence $c_R\mu\ge969$.

After these come $\mathsf K_{\rm c}^0$ and $\mathsf a_{\rm c}$, from item~C.II~4 of \S\ref{Y-C-1}, under ``Constants''.

\item[\textup{(O7)}] $M=L^{m'}$ with \textup{(M2)} $M\ge1920(\kappa+4)c_R$ and \textup{(M4$^\flat$)} $M\ge(54/\mu)(\kappa^\sharp+3)$. The condition \textup{(M4$^\sharp$)} is a consequence of \textup{(M4$^\flat$)}. Further \textup{(M5r)} $M\ge\rho_0$, for Thms.~\ref{Y-I-ThC} and \ref{Y-J-ThA}.

After these come $\kappa_+=M/(40c_R)$ and $\mathsf K_1^\flat$, the number of item~C.II~5 of \S\ref{Y-C-1}.

\item[\textup{(O8)}] $\mathfrak p_0>2$, strictly, and $A_0>0$.

\item[\textup{(O9)}] $c_\xi=(6.5C_\varphi^\flat)^{-1}$, and $\mathsf n=60$.

\item[\textup{(O10)}] First, the numbers that come from the cited theorems. $\mathsf K_{\rm c}:=\sqrt2\mathsf K_{\rm c}^0$ is used for \textup{(Cen)} alone, that is \Cref{Y-C-CIICor72}\,(a); there it is the corollary's own number, for each $0<\hat\varepsilon\le\frac1{40}$.

We take the numbers $\mathsf K_1^\flat$ and $1.5\mathsf K_1^\flat$ in the two bounds on $\tilde{\mathsf D}=-\mathsf D_k/2$, of \Cref{Y-C-CIICor75}; further $K_{\rm q}:=K_{{\rm q},k}$, of \Cref{Y-D-C5}, and $K_v:=K_v^\flat$, of \Cref{Y-E-C6}.

From item F of \S\ref{Y-F-1}: $h_*$, $u_b$, $u_t$, $u_1$, $\delta_3=\min\{d_0/1536,\ 3c_\xi^2/(256h_*^2)\}$, $e_1^\flat$, ${\rm Res}$ and $e^{o}$, with the primed $u_t'$, $u_1'$, $\delta_3'=\min\{d_0/2304,\ 3c_\xi^2/(256h_*^2)\}$ and $e^{o\prime}$.

Our own numbers are $a_\infty:=4\pi \mathsf N_{\mathsf B}(C_\varphi^\flat)^2$ and $n_a:=(3v/2)a_\infty^2$.

As to dependence: first, $\mathsf K_{\rm c},\mathsf K_1^\flat,K_{\rm q},K_v,h_*,\delta_3,\delta_3',a_\infty,n_a$ are numbers; none of them varies with $c_h$, $E_a$, $x$, $N$ or a member, and of $\ell$ only the pure numbers $b_*$, $K_*$, $\lambda_*$ and the range enter them.

Second, $u_b,u_t,u_1,u_t',u_1',e_1^\flat$ are functions of $x$ built from those numbers and $(A_0,\mathfrak p_0)$, and $u_t$, $u_t'$ depend in addition on $r_1(x)$.

Third, ${\rm Res}$, $e^{o}$, $e^{o\prime}$ depend in addition on the $\phi'_P$ of $D'$ and on the $\mathsf t_\rho$ of the rows; no floor contains them.

Fourth, not one of all these quantities depends on $c_h$, $E_a$, $N$ or a member.

\item[\textup{(O11)}] $0<c_h\le\delta_3/2$. The table takes $c_h\le\delta_3'/2$: Thms.~\ref{Y-B-C1} and \ref{Y-G-ThA} require $c_h\le\delta_3/2$, and $\delta_3'\le\delta_3$ by \textup{(O10)}.

\item[\textup{(O12)}] \[E_a\ge E_a^0:=\max\{32L^{-1/2}K_v,\ 32K_{\rm q}\}\]. The table takes the larger value $\max\{32K_{\rm q},\ 1500L^{-1/2}K_v\}$; \Cref{Y-B-C1} requires $E_a\ge E_a^0$ alone.

\item[\textup{(O13)}] Last: $x_\flat\ge x_\flat^0=x_\flat^0(c_h,E_a):=1+\inf S(c_h,E_a)$, where $S(c_h,E_a)$ is the set of those $y\ge e^2$ for which \textup{(A0)} to \textup{(A11)} hold at all $x\ge Ly$.

The set $S$ is an upper half line, and it is not empty, by Step~14 of the proof of \Cref{Y-G-ThA}.

The set $S$ need not contain its infimum. The floors contain ceilings, namely $r'$, $r_1$, $\lceil s(x)/L\rceil$, $R$, $\varrho_L$, which are continuous from the left and jump upwards on the right; so a least such number need not exist.

But $1+\inf S\in S$, and it depends on what $S$ depends on.

The table takes $x_\flat\ge\max\{x_\flat^{(1)},x_\flat^0,x_9/L\}$, together with \textup{(X12)}.

Here $x_\flat^{(1)}$ is the number of \Cref{Y-B-C1}, namely $x_\flat^0$ evaluated at $(c_h,E_a)=(\delta_3/2,E_a^0)$, and hence computed from \textup{(O0)} to \textup{(O10)} alone.

And $x_9$ is that of \textup{(O13)} of \S\ref{Y-H-1}, determined by \textup{(O0)} to \textup{(O12)} alone; so there is no circle.

\S\ref{Y-J-1} adds \textup{(X13)} to \textup{(X18)} under its \textup{(O13)}.
\end{enumerate}
\end{definition}

\subsubsection*{Item 4. The floors}

\begin{definition}[The numbers of the floors]
We set $r_\eta:=(0.22+L^{1/2}\mathsf a_{\rm c})\varepsilon$, $r_T:=\sqrt2C_\varphi^\flat c_\xi\varepsilon+17T\mathsf a_{\rm c}\varepsilon'$, $\varepsilon'':=x^{-1/4}$ and $\upsilon:=e^{-c_\Xi p_0^2}$.

We set $r_\Xi:=\frac12-d_0-\mathsf t_\rho-3e_1^\flat-c_\Xi$; it depends on the rows through $\mathsf t_\rho$.

We set $r_\Xi^\flat:=\frac78(\frac12-d_0-c_\Xi)-3e_1^\flat$; it depends on $x$ and the letters alone.

By \textup{(SD)}, $\mathsf t_\rho\le\frac18(\frac12-d_0-c_\Xi)$; hence $r_\Xi\ge r_\Xi^\flat$.
\end{definition}

\begin{definition}[The floors \textup{(A0)} to \textup{(A11)}]
The floors are the following conditions on $x$.
\begin{enumerate}
\item[\textup{(A0)}] \emph{Part 1.} The floor $x\ge x_{\rm fl}$ of item G of \S\ref{Y-F-1}, which contains \textup{(W)}; \textup{(F$_{\rm law}$)} of \Cref{Y-C-CITh66}\,(k); $x\ge L^2$, $R\ge6$, \textup{(P1)} and \textup{(P3)} of item~C.II~4 of \S\ref{Y-C-1}; and \textup{(S1)}, \textup{(S2)} of item~C.II~5 of \S\ref{Y-C-1}, with $r_u=\lceil R_2/M\rceil$.

\emph{Part 2.} The floors of \Cref{Y-D-C5}, that is \textup{(G1$_k$)}, \textup{(F8$_k$)}, \textup{(F9$_k$)}, and those of \Cref{Y-E-C6}, that is \textup{(W0)} to \textup{(W5)}.

\emph{On \textup{(P0)}.} Of \textup{(P0)}, which \textup{(W0)} names, only the two clauses $x\ge L^2$ and $R\ge6$ are floors. In item~C.II~4 of \S\ref{Y-C-1}, under ``Floors'', \textup{(P0)} carries a third clause, \textup{(Tor)}. It involves $N$: for fixed $N$ it is violated by every large enough $x$. So \textup{(Tor)} is a datum of \textup{(SD)} and no floor, and the same is true in the other sections from which \textup{(A0)} takes its floors, for instance in Step~14 of the proof of \Cref{Y-E-C6}, in item G of \S\ref{Y-F-1}, and in item G of \S\ref{Y-D-1}. Where a statement taken from another section and listed in Step~0 of the proof of \Cref{Y-G-ThA} uses \textup{(Tor)}, \textup{(SD)} supplies it, as for the statements \textup{($\Gamma3$)} and \textup{($\Gamma4$)} of that step. Hence none of the conditions that define $S(c_h,E_a)$ involves $N$.

\emph{Part 3.} $3D_w\le\ell_w(x')$; $17T\varepsilon'\le\frac1{40}$; and \textup{(D2)}, \textup{(D3)} of \Cref{Y-H-Th53} at $(\tau,r)=(1,r_\eta)$ and at $(17T,r_T)$.

Here \textup{(D2)}$(\tau,r)$ reads $\ell_w(x)(7.1r+18D_w\tau\varepsilon')\le \mathsf g_x/2$, and \textup{(D3)}$(\tau,r)$ reads $(4.1r+2.1\tau\varepsilon')^2\le0.1\varepsilon$.

\emph{Part 4.} \textup{(C0)}, \textup{(C1)}, \textup{(C2)} of \Cref{Y-H-Th57}, at this $\omega$ and this $T$.

These are: \textup{(C0)} $18\omega\varepsilon'\le\frac1{10}$ and $160T\omega\varepsilon'\le1$; \textup{(C1)} $2640\,\ell_w(x')(T+1)\omega\varepsilon'\le \mathsf g_x$; \textup{(C2)} $102400\,T\omega^2\varepsilon'\le1$.

\emph{Part 5.} $\mathsf K_{\rm c}(17T\varepsilon')^2\le0.12\varepsilon$.

\emph{Notes.} Here $\mathsf g_x=\vartheta_0^\flat\log L/(4(\log x)^2)$. The condition \textup{(D0)}$(\tau)$ is: $2D_w+1<N$, $3D_w\le\ell_w(x')$, $\tau\varepsilon'\le\frac1{40}$; its first clause is no floor, since it follows from \textup{(Tor)}, see Step~0 of the proof of \Cref{Y-G-ThA}, at the clauses on $N$.

\item[\textup{(A1)}] $\varepsilon\le\frac1{40}$; $\rho_0\le\varrho(x')$; $\varrho_L+2\le\lceil R/2\rceil$; and $0.72(4.1r_\eta+2.1\varepsilon')^2+\mathsf K_{\rm c}\varepsilon'^2\le0.1\varepsilon$.

\item[\textup{(A2)}] $5r_1+2\lceil s(x)/L\rceil+6\le r'$, $\omega\le Mr'$ and $R_6+1\le Mr'$.

\item[\textup{(A3)}] $50x^{-1/8}(2(s_0+\varrho_*+3r_1+2)+1)^3\le1$.

\item[\textup{(A4)}] $(c_a+c_h)(x)\le\frac12$.

\item[\textup{(A5)}] \[x^{1/4}(2s'+1)^3\cdot50e^{(\kappa+1)s'}\bigl[5e^{\kappa+3}L^3(c_a+c_h)(x)+2e^{\kappa+3}K_vp_0^3x^{-1/2}+x^{-20}+e^{\kappa+3}u_b+u_1\bigr]\le x^{-1/8}.\]

\item[\textup{(A6)}] \[\bigl[u_te^{\kappa+5}+220(2r_1+1)^3L^3e^{\kappa+5}\upsilon+6((2r_1+1)L+2r(x))^3e^{\kappa+5}(c_a+c_h)(x)\bigr]e^{(\kappa+4)s'}\le\frac12\varepsilon''.\]

\item[\textup{(A7)}] $3e_1^\flat\le\frac18(\frac12-d_0-c_\Xi)$; and, with \[\upsilon_f:=\exp\{-r_\Xi^\flat p_0'^2+1.01\mathsf K_1^\flat(2r'+1)^3+4+(\kappa+4)(1+2s')\},\] defined with $r_\Xi^\flat$ and not with $r_\Xi$ so that this floor depends on no rows, $e^{(\kappa+1)s'}\upsilon_f^{1/2}\le\frac14\varepsilon''$.

\item[\textup{(A8)}] $(p_0'/p_0)^4\ge\frac12$; $25.5\mathsf K_1^\flat\le E_ap_0'^2$; $111\cdot4.86L^3x^{-1/4}\le0.0026L^{1/2}$; $111x^{-1/4}(2K_v/E_a+1)\le0.003L^{1/2}$; $x^{-20}\le c_a(x)$; and $2x^{-20}\le0.001c_a(x')$.

\item[\textup{(A9)}] $10^3L^3e^{-c_h(p_0^2-p_0'^2)}\le\frac12$; $e^{\delta_3p_0^2/2}\ge200$; and $250v(c_\xi p_0+1/(c_\xi p_0))e^{-0.49c_\xi^2p_0^2}\le\frac14$.

\item[\textup{(A10)}] $4.04x^{-20}\le(\gamma/2)x^{-2}$; $2.02x^{-19}n_a+38.8v\gamma^{-1}x^{-20}(2r_1+1)^3\le x^{-1/8}$; and $\varsigma':=4.04x^{-20}\cdot3L^3(R+1)^3(C_\varphi^\flat)^2\le\frac12$ together with $9v\varsigma'(2r_1+1)^3+2.02x^{-19}n_a\le x^{-1/8}$.

\item[\textup{(A11)}] $\mathsf K_1^\flat e^{-\kappa^\sharp r'}\le1$, and $(2r'+1)^3\cdot2.01\mathsf K_1^\flat e^{\kappa+1}e^{-(\kappa+26.74)r'}\le1$.
\end{enumerate}
\end{definition}

\begin{remark}[Independence of the floors]
No floor depends on $N$, since \textup{(Tor)} is a datum of \textup{(SD)} and no floor, by \textup{(A0)}; and no floor depends on $\ell$ beyond the pure numbers $b_*$, $K_*$, $\lambda_*$ and the range, nor on a member, since $\mathsf t_\rho$ enters none of them, \textup{(A7)} being defined with $r_\Xi^\flat$.

Hence one and the same number $x_\flat^0$ applies to every choice of rows in \textup{(SD)}, the empty rows included, every $N$ and every level $k$.
\end{remark}

\subsection{Free kernels and row estimates}\label{Y-A-1}

The quadratic part of one blocking step is a constrained minimization.
The fine quadratic form is $1+\ell^{\mathsf T}\ell$, with rows $\ell$
indexed by sites. We first identify the minimizer and show that its
coarse form again has a sum-of-squares representation.

To iterate the construction, we need estimates on the new rows that
are uniform in the number of steps. The exponential bounds below give
this control and quantify the error from cutting off distant entries.
The sharper sup bounds control the centre of the Gaussian integration.
\Cref{Y-A-Fr} assembles these estimates at the parameters used
in the nonlinear argument.

\begin{remark}[Letters with several senses]\label{Y-A-Letters} The letter $K$ is the bound of the rows in $\mathcal R(b,K,\lambda,\tau)$, the number of steps in $L^{K-k}$, and a matrix in $\|K\|_\theta$.

On the base line of a formula, $\mathsf T$ is the kernel of \Cref{Y-A-XT}, while a raised ${}^{\mathsf T}$ is the transpose.

The letter $N$, when it denotes the side of a torus, has two senses, and both occur in \Cref{Y-A-Th29}: where a single step is treated, as in part~\textup{(a)} of that theorem, $N$ is the side of the coarse torus of that step; where a tower is treated, as in part~\textup{(b)}, $N$ is the side of the last torus $\mathbb T_K$, so that step $k$ has a coarse torus of side $N'=NL^{K-k-1}$. \end{remark}

\begin{remark}[Standing conditions with $a>0$]\label{Y-A-Standing2}
In the next five subsections, $L\ge3$ is odd, $N\ge2$ is even and
$a>0$. We state the stronger bounds on $a$ where they first enter.
The row iteration will use $aL\ge3/c_L$.
\end{remark}

\subsubsection{Lattices, forms, averages and distances}\label{Y-A-SubLattices}

\begin{definition}[Lattices, forms and the four averages]\label{Y-A-Lattices} Let $L\ge3$ be odd, $\mathsf h_L=\frac12(L-1)$ and $c_L=\frac{L^2-1}{6L^2}<\frac16$. Let $N\ge2$ be even. The fine and the coarse torus are $\mathbb T=\mathbb Z^3/NL\mathbb Z^3$ and $\mathbb T'=L\mathbb Z^3/NL\mathbb Z^3$, and $\Omega^i,\Omega'^i$ denote the real $i$-cochains on them. We write $\mathrm{d}$ and $\delta=\mathrm{d}^*$ for the fine coboundary and its adjoint, and $\mathrm{d}'$, $\delta'$ for the coarse ones; on $2$-forms, $\mathrm{d}'^*$ denotes the adjoint of $\mathrm{d}'$ on $1$-forms, whereas $\delta'$ on $3$-forms is the adjoint of $\mathrm{d}'$ on $2$-forms. Fine bonds, plaquettes and cubes are denoted $b,p,q$, and coarse ones $c,P,q'$. The block $B(y)$ is centred at $y$; $W_c$ is the double block of $c$, and $W_P$ is made of the four blocks at the corners of $P$.

The four averages are $(Q_0\lambda)(y)=L^{-3}\sum_{x\in B(y)}\lambda(x)$ and $(Q_1A)(c)=L^{-3}\sum_{x\in B(c_-)}A([x,x'])$, where, for $c=(y,\mu)$, $c_-=y$, $x'=x+Le_\mu$ and $[x,x']$ is the straight path from $x$ to $x'$; then $W=Q_2$, given by \[(WF)(y;\mu\nu)=L^{-3}\sum_{x\in B(y)}\sum_{0\le i,j<L}F(x+ie_\mu+je_\nu;\mu\nu),\] and $Q_3$, given by \[(Q_3\Theta)(y)=L^{-3}\sum_{x\in B(y)}\sum_{0\le i,j,k<L}\Theta(x+ie_1+je_2+ke_3).\]\end{definition}

\begin{lemma}[Facts on the averages]\label{Y-A-Facts} Let $t(s)=L-|s-\mathsf h_L|$ for $-\mathsf h_L\le s\le3\mathsf h_L$, and $t(s)=0$ for all other $s$. In the kernel formulas below, $t$ and
the indicator of $[-\mathsf h_L,\mathsf h_L]$ are periodized modulo $NL$.
Thus, for example, $t(z_\mu-y_\mu)$ in a kernel means
$\sum_{m\in\mathbb Z}t(z_\mu-y_\mu+mNL)$. For $p=(z;\mu\nu)$ and $P=(y;\mu\nu)$ we have, with $\lambda$ the third direction, normal to $p$, \[W(P,p)=L^{-3}t(z_\mu-y_\mu)t(z_\nu-y_\nu)1\{|z_\lambda-y_\lambda|\le\mathsf h_L\},\] and, for the coarse cube $q'$ at $y$ and the fine cube $q$ at $z$, \[Q_3(q',q)=L^{-3}\prod_\mu t(z_\mu-y_\mu)=:\psi_{q'}(q).\] The function $t$ satisfies $\sum_{y\in L\mathbb Z}t(s-y)=L$, $\sum_st(s)=L^2$ and $t\le L$.

The averages intertwine the coboundaries: $\mathrm{d}'Q_i=Q_{i+1}\mathrm{d}$. For $i=1$ this follows from $(\mathrm{d}'Q_1A)(P)=L^{-3}\sum_xA(\partial S_x)$, where $S_x$ is the $L\times L$ square at $x$. For $i=0,2$ the corresponding identities are that $(\mathrm{d}'Q_0\lambda)(c)$ equals $L^{-3}$ times the sum over $x\in B(c_-)$ of $\lambda$ on the boundary of the path $[x,x']$, and that $(\mathrm{d}'WF)(q')$ equals $L^{-3}$ times the sum over $x$ of $F$ on the boundary of the $L\times L\times L$ cube at $x$.

The kernel of $W$ satisfies $W(P,p)\ge0$, $\sum_PW(P,p)=L^{-1}$, $\sum_pW(P,p)=L^2$ and $W(P,p)\le L^{-1}$, and $W(P,p)\ne0$ only for $p\subset W_P$. Further, $0\le\psi_{q'}\le1$, $\sum_{q'}\psi_{q'}=1$ and $\sum_q\psi_{q'}(q)=L^3$. By Schur's test, $\|Q_1\|^2\le L^{-1}$, $\|W\|^2\le L$ and $\|Q_3\|^2\le L^3$. Moreover, $\sum_st(s)^2=\frac13L(2L^2+1)$ and $\sum_st(s)t(s-L)=\frac16L(L^2-1)$.

Let $\Delta'_\mu=2-\mathsf S_\mu-\mathsf S_\mu^*$, where $\mathsf S_\mu$ is the coarse shift. Then $Q_3Q_3^*=L^3\hat{\mathsf J}$, with $\hat{\mathsf J}=\prod_\mu(1-c_L\Delta'_\mu)\ge(1-4c_L)^3\ge\frac1{27}$, and on the component $\mu\nu$, $L^{-1}WW^*=(1-c_L\Delta'_\mu)(1-c_L\Delta'_\nu)$. Furthermore, $\mathsf P:=1-L^{-1}WW^*-(aL)^{-1}\mathrm{d}'\mathrm{d}'^*\ge0$ if $aL\ge3/c_L$. Indeed, $\mathrm{d}'\mathrm{d}'^*\le2\,{\rm diag}(\Delta'_\mu+\Delta'_\nu)$; on the component $\mu\nu$ the operator $1-L^{-1}WW^*$ equals $c_L(\Delta'_\mu+\Delta'_\nu)-c_L^2\Delta'_\mu\Delta'_\nu$, and as $\Delta'_\mu\Delta'_\nu\le2(\Delta'_\mu+\Delta'_\nu)$, \[c_L(\Delta'_\mu+\Delta'_\nu)-c_L^2\Delta'_\mu\Delta'_\nu\ge(c_L-2c_L^2)(\Delta'_\mu+\Delta'_\nu)\ge\frac23c_L(\Delta'_\mu+\Delta'_\nu);\] lastly, $\frac2{aL}\le\frac23c_L$. The condition $aL\ge3/c_L$ also gives $aL\ge18$.\end{lemma}

\begin{definition}[Base sites and the two distances]\label{Y-A-Distances} Every cell $i$ has a base site $\pi(i)\in\mathbb T$, which for a coarse cell is a point of $\mathbb T'\subset\mathbb T$. The fine distance, in block units, is $\mathsf d(i,j):=L^{-1}|\pi(i)-\pi(j)|_\infty$; it is a pseudo-metric on all cells of the step. Let $y(i)$ be the coarse site whose block contains $\pi(i)$; the coarse distance is ${\rm dist}(i,j):=L^{-1}|y(i)-y(j)|_\infty$. Since $|\pi(i)-y(i)|_\infty\le\mathsf h_L$, we have $|\mathsf d-{\rm dist}|<1$, and $\mathsf d={\rm dist}$ for two coarse cells. Finally, $\mathsf g(\lambda):=\sum_{n\in\mathbb Z^3}e^{-\lambda|n|_\infty}$.\end{definition}

\begin{definition}[The weighted norm of a matrix between cells]\label{Y-A-Norm} For $\theta\ge0$ and a matrix $K$ between two kinds of cells we set \[\|K\|_\theta:=\max\{\sup_i\sum_j|K(i,j)|e^{\theta{\rm dist}(i,j)},\ \sup_j\sum_i|K(i,j)|e^{\theta{\rm dist}(i,j)}\}.\] The triangle inequality gives $\|K_1K_2\|_\theta\le\|K_1\|_\theta\|K_2\|_\theta$. By Schur's test, the norm $\|K\|_\infty$ of $K$ as a map $\ell^\infty\to\ell^\infty$ and its operator norm $\|K\|$ on $\ell^2$ are both $\le\|K\|_0$. Moreover, $\sup_i\sum_{{\rm dist}(i,j)>R}|K(i,j)|\le e^{-\theta R}\|K\|_\theta$.\end{definition}

\subsubsection{Edge lifts and the slice}\label{Y-A-SubLifts}

\begin{definition}[Edge lifts, combs and the slice]\label{Y-A-Lifts} The crossing bonds of a coarse bond $c=(y,\mu)$ are the fine bonds $b=(z,\mu)$ with $z_\mu=y_\mu+\mathsf h_L$ and $|z_\nu-y_\nu|\le\mathsf h_L$ for $\nu\ne\mu$; among them, $b_c$ is the one with $z=y+\mathsf h_Le_\mu$. The edge plaquettes under a coarse plaquette $P=(y;\mu\nu)$ are the plaquettes $p=(z;\mu\nu)$ with $z_\mu=y_\mu+\mathsf h_L$, $z_\nu=y_\nu+\mathsf h_L$ and $|z_\lambda-y_\lambda|\le\mathsf h_L$, where $\lambda$ is the third direction. There are $L$ of them. Under the coarse cube at $y$ lies its corner cube, the fine cube at $z=y+\mathsf h_L(1,1,1)$. The lifts are given by $(\mathsf JA')_b=A'_c$ on the crossing bonds of $c$, with $\mathsf JA'=0$ on the bonds inside blocks; by $(EX')_p=X'_P$ on the edge plaquettes under $P$, with $EX'=0$ elsewhere; and by $(E_3\Theta)_q=\Theta_{q'}$ on the corner cube under $q'$, with $E_3\Theta=0$ elsewhere.

For $x\in B(y)$, the comb $\Gamma(y,x)$ is the path from $y$ to $x$ that runs first along $e_1$, then along $e_2$ and then along $e_3$; the bonds of all combs form a set $\mathcal T$, which is a forest with one tree to a block. The slice is $\mathfrak L=\{\varphi\in\Omega^1:\varphi=0\text{ on }\mathcal T\}$.\end{definition}

\subsubsection{Rows and the mask}

\begin{definition}[Rows, their norms, and the mask]\label{Y-A-Rows} Rows of a level are given by a Euclidean space $\mathsf H_y$ of finite dimension for each site $y$, together with a linear map $\ell_y:\Omega^2\to\mathsf H_y$ whose columns are $\ell_y(p)\in\mathsf H_y$; we write $\ell=(\ell_y)_y:\Omega^2\to\mathsf H=\bigoplus\mathsf H_y$. The norm $|\cdot|$ on $\mathsf H_y$ is the Euclidean one, and for $\sigma\in\mathsf H$ we set $|\sigma|_{2,\infty}=\sup_y|\sigma_y|$. In $\|\cdot\|_\theta$, the entry of $\ell$ at $(y,p)$ is taken to be $|\ell_y(p)|$. We put \[\|\ell\|_{\rm r}=\sup_y\sum_p|\ell_y(p)|,\qquad\|\ell\|_{\rm c}=\sup_p\sum_y|\ell_y(p)|,\qquad\mathsf h_\ell^2=\sup_y\sum_p|\ell_y(p)|^2\ (\le\|\ell\|_{\rm r}^2).\] Nothing that follows depends on the dimension of $\mathsf H_y$; all of it is invariant under orthogonal maps of each $\mathsf H_y$. The rows have range $\varrho$ if $\ell_y(p)=0$ whenever, in units of that level, the distance between $y$ and $p$ is $>\varrho$. We set $\mathsf B=1+\ell^{\mathsf T}\ell$, $\mathsf R=\mathsf B^{-1}\le1$ and $b_{\mathsf B}=\|\mathsf B\|\le1+\|\ell\|_0^2$. At depth $0$, $\mathsf H=0$.

The mask is $m_y(V)=1\{\text{no plaquette of }D(V)\text{ within the distance }\rho_0\text{ of }y\}$, the distance being measured in units of the level of $V$ and $y$; a single plaquette of $D(V)$ switches off at most $(2\rho_0+1)^3$ sites. The mask of the class has a second clause, in item~3.5 of \S\ref{Y-E-1}: in addition, no $p$ with $\phi_p>1$ lies within $\varrho(x)+1$ of $y$. The arguments of \S\ref{Y-A-1} do not depend on the clauses of a mask: \Cref{Y-A-P8}, \textup{(N$'$)} and \textup{(Fr3$'$)} are stated for \emph{every} map $m$ from the sites to $\{0,1\}$.\end{definition}

\subsubsection{The lifts, the form on the slice, and an inverse}

\begin{lemma}[The lifts invert the averages and intertwine the coboundaries]\label{Y-A-L1} We have $Q_1\mathsf J=1$, $\mathrm{d}\mathsf J=E\mathrm{d}'$, $WE=1$, $\mathrm{d}E=E_3\mathrm{d}'$ and $Q_3E_3=1$.\end{lemma}

\begin{proof} For $x\in B(c_-)$, the path $[x,x']$ contains exactly one crossing bond, and this bond belongs to $c$; hence $Q_1\mathsf J=1$.

Let $p=(z;\mu\nu)$, and say that $z_\mu$ is at the top of its block when $z_\mu=y_\mu+\mathsf h_L$. If neither $z_\mu$ nor $z_\nu$ is at the top, the four bonds of $p$ lie in one block. If exactly one of them is at the top, say $z_\mu$, then $(z,\mu)$ and $(z+e_\nu,\mu)$ are crossing bonds of one and the same $c$ and cancel, while the bonds of $p$ in the direction $\nu$ lie in blocks. If both are at the top, the four bonds of $p$ are crossing bonds for the four bonds of $P$, with the signs of $\mathrm{d}'$. So $\mathrm{d}\mathsf J=E\mathrm{d}'$.

Since $t(\mathsf h_L)=L$ and $t(\mathsf h_L\pm L)=0$, we have $W(P,p)=L^{-1}$ at the $L$ edge plaquettes under $P$ and $W(P,p)=0$ at those under any other $P'$; hence $WE=1$. The same two values of $t$, inserted in $\psi_{q'}(q)=L^{-3}\prod_\mu t(z_\mu-y_\mu)$, give $\psi_{q'}=1$ at the corner cube under $q'$ and $\psi_{q'}=0$ at the corner cubes under the other coarse cubes; hence $Q_3E_3=1$.

Let $q$ be the cube at $z$. If $z_\mu,z_\nu$ are at the top and $z_\lambda$ is not, the faces $(z;\mu\nu)$ and $(z+e_\lambda;\mu\nu)$ lie under one $P$ and cancel, and no other face is an edge plaquette. The six faces of a corner cube lie under the six faces of the coarse cube. Hence $\mathrm{d}E=E_3\mathrm{d}'$.\end{proof}

For $\varphi\in\Omega^1$ let $f:=\mathrm{d}\varphi$, and let $f^\circ$ be $f$ with the edge plaquettes of Subsection~\ref{Y-A-SubLifts} left out. On $\mathfrak L$ we use the form $q(\varphi):=\langle \mathrm{d}\varphi,\mathsf B\mathrm{d}\varphi\rangle+a|Q_1\varphi|^2$, that is, $q=\Pi_{\mathfrak L}(\mathrm{d}^*\mathsf B\mathrm{d}+aQ_1^*Q_1)\Pi_{\mathfrak L}$, as in the proof of \Cref{Y-A-P35}.

\begin{lemma}[Bounds on the form on the slice]\label{Y-A-L2} Let $\varphi$ vanish on the combs. The local estimates are valid for
cochains with values in a real Hilbert space. For a block $B(y)$ and
a coarse bond $c$ they read, with $f=\mathrm{d}\varphi$,
\begin{align}
 \sum_{b\subset B(y)}|\varphi_b|^2
     &\le \tfrac12 L^3\sum_{p\subset B(y)}|f_p|^2,
       \label{eq:comb-local-poincare}\\
 \sum_{b\text{ crossing }c}|\varphi_b|^2
     &\le 6L^2|(Q_1\varphi)_c|^2+12L^4\sum_{p\subset W_c}|f_p|^2.
       \label{eq:crossing-local-poincare}
\end{align}
Summing gives $|\varphi|^2\le73L^4|f^\circ|^2+6L^2|Q_1\varphi|^2$. Hence $\gamma_q|\varphi|^2\le|f^\circ|^2+a|Q_1\varphi|^2\le q(\varphi)\le b_q|\varphi|^2$, the first inequality for $a\ge6/(73L^2)$ and the other two for every $a>0$, where $\gamma_q:=(73L^4)^{-1}$ and $b_q:=16b_{\mathsf B}+a/L$. If $aL\ge3/c_L$, then $a\ge6/(73L^2)$.\end{lemma}

\begin{proof}\emph{Step 1: bonds inside a block.} Let $b=(z,\nu)\subset B(y)$. If $\nu=3$, then $b\in\mathcal T$. If $\nu=2$, the closed path $\Gamma(y,z)+b-\Gamma(y,z+e_2)$ bounds a ladder of $|z_3-y_3|$ plaquettes $\{23\}$. If $\nu=1$, the closed path $\Gamma(y,z)+b-\Gamma(y,z+e_1)$ bounds a ladder of $|z_3-y_3|$ plaquettes $\{13\}$ together with a ladder of $|z_2-y_2|$ plaquettes $\{12\}$ at the height $y_3$. As $\varphi$ vanishes on the combs, in each case $\varphi_b=\sum_{p\in\sigma_b}\pm f_p$ with $\#\sigma_b<L$ and $\sigma_b\subset B(y)$, and a plaquette belongs to at most $\mathsf h_LL\le\frac12L^2$ of the sets $\sigma_b$. Hence, by the Cauchy--Schwarz inequality, \[\sum_{b\subset B(y)}|\varphi_b|^2\le \tfrac12 L^3\sum_{p\subset B(y)}|f_p|^2\] in every block $B(y)$.

\emph{Step 2: crossing bonds.} On the combs, $\varphi=0$. For a coarse bond $c=(y,y')$ of direction $\mu$, for $x\in B(y)$ and $x'=x+Le_\mu$, let $\gamma_{c,x}=\Gamma(y,x)+[x,x']-\Gamma(y',x')$. Then $\gamma_{c,x}-\gamma_{c,y}=\partial n_{c,x}$, where $n_{c,x}$ consists of fewer than $L^2$ plaquettes of $W_c$, each occurring once. To see this, we move the bonds of $\Gamma(y,x)$ that are not parallel to $e_\mu$ by $L$ unit steps in the direction $e_\mu$; they sweep out at most $2\mathsf h_L$ strips of $L$ plaquettes each, these strips are disjoint, and their boundaries telescope. On every torus with $N\ge2$ even, the identity is the image of the identity of chains on $\mathbb Z^3$, and still each plaquette occurs once: the $L$ steps of a strip are pairwise distinct mod $NL$, and the other two coordinates of a strip stay within one block. So $\varphi([x,x'])-\varphi([y,y'])=\sum_{n_{c,x}}\pm f_p$.

Let $\mathsf F_c=\sum_{p\subset W_c}|f_p|^2$. Averaging over $x$, we get $|\varphi([y,y'])|^2\le2|(Q_1\varphi)_c|^2+2L^2\mathsf F_c$. For a crossing bond $b=(z,\mu)$ of $c$ we take $x=z$; then $\varphi_b=\varphi([y,y'])+\sum_{n_{c,x}}\pm f_p-\beta_b$, where $\beta_b$ is the sum of $\varphi$ over the $L-1$ other bonds of $[x,x']$, which lie inside $B(y')$. The paths of different $b$ are disjoint, so the Cauchy--Schwarz inequality and Step~1, applied in $B(y')$, give $\sum_b|\beta_b|^2\le L\cdot L^3\mathsf F_c$. Therefore \[\sum_{b\text{ of }c}|\varphi_b|^2\le3L^2[2|(Q_1\varphi)_c|^2+3L^2\mathsf F_c]+3L^4\mathsf F_c=6L^2|(Q_1\varphi)_c|^2+12L^4\mathsf F_c\] for every coarse bond $c$.

A plaquette lies in at most six of the sets $W_c$, and an edge plaquette in none. Summing over the blocks and the coarse bonds, we obtain $|\varphi|^2\le(L^3+72L^4)|f^\circ|^2+6L^2|Q_1\varphi|^2$. Since $L^3\le L^4$, this proves the summed estimate. If $a\ge6/(73L^2)$, then $6L^2\le73L^4a$, giving the lower bound with $\gamma_q$. The second inequality, $|f^\circ|^2+a|Q_1\varphi|^2\le q(\varphi)$, holds because $\mathsf B\ge1$ and $|f^\circ|\le|f|$. The last inequality follows from $\langle f,\mathsf Bf\rangle\le b_{\mathsf B}|f|^2$, from $|\mathrm{d}\varphi|^2\le16|\varphi|^2$ and from $\|Q_1\|^2\le L^{-1}$ of \Cref{Y-A-Facts}. Since $c_L<\frac16$, the condition $aL\ge3/c_L$ gives $aL>18$, and so $a>18/L>6/(73L^2)$.\end{proof}

\begin{lemma}[Decay of an inverse; a ratio of determinants]\label{Y-A-L3} Let $I$ be a finite set with a pseudo-metric ${\rm dist}$, let $\mathsf A$ be real symmetric on $\ell^2(I)$ with $0<\gamma\le\mathsf A$, and let $\varpi_{\mathsf A}(\lambda):=\sup_i\sum_j|\mathsf A(i,j)|(e^{\lambda{\rm dist}(i,j)}-1)$.\begin{enumerate}\item[\textup{(a)}] If $\varpi_{\mathsf A}(\lambda)\le\frac\gamma2$, then $|\mathsf A^{-1}(i,j)|\le\frac2\gamma e^{-\lambda{\rm dist}(i,j)}$.

\item[\textup{(b)}] If $\mathsf A_1,\mathsf A_2>0$, $\operatorname{rank}(\mathsf A_1-\mathsf A_2)\le r$ and $\|\mathsf A_1\|\|\mathsf A_2^{-1}\|\ge1$, then $\det\mathsf A_1/\det\mathsf A_2\le(\|\mathsf A_1\|\|\mathsf A_2^{-1}\|)^r$.\end{enumerate}\end{lemma}

\begin{proof}\emph{Step (a): the decay of the inverse.} For $\lambda\le0$, the bound $\|\mathsf A^{-1}\|\le\gamma^{-1}$ proves the assertion directly. Hence assume $\lambda>0$. Fix $j_0$, and let $\rho={\rm dist}(\cdot,j_0)$ and $\mathsf A_\lambda=e^{\lambda\rho}\mathsf Ae^{-\lambda\rho}$. By the triangle inequality for ${\rm dist}$, $|\mathsf A_\lambda(i,j)-\mathsf A(i,j)|\le|\mathsf A(i,j)|(e^{\lambda{\rm dist}(i,j)}-1)$, and Schur's test gives $\|\mathsf A_\lambda-\mathsf A\|\le\varpi_{\mathsf A}(\lambda)$. So $\|\mathsf A_\lambda u\|\ge\frac\gamma2\|u\|$, and $\mathsf A^{-1}(i,j_0)=e^{-\lambda\rho(i)}\mathsf A_\lambda^{-1}(i,j_0)$; since $|\mathsf A_\lambda^{-1}(i,j_0)|\le\frac2\gamma$, this proves \textup{(a)}.

\emph{Step (b): the ratio of determinants.} Write $\mathsf A_2^{-1/2}\mathsf A_1\mathsf A_2^{-1/2}=1+K$ with $\operatorname{rank}K\le r$. All but $r$ eigenvalues of $1+K$ are $1$, and the others are $\le\|\mathsf A_1\|\|\mathsf A_2^{-1}\|$. As $\det\mathsf A_1/\det\mathsf A_2=\det(1+K)$ and $\|\mathsf A_1\|\|\mathsf A_2^{-1}\|\ge1$, this proves \textup{(b)}.\end{proof}

\subsubsection{Pairs and the minimiser of one step}

\begin{definition}[Pairs, the move, and $R_\sharp$]\label{Y-A-Pairs} Let $\mathcal H=\Omega^2\oplus\Omega'^1$, with $\langle(F,\mathsf n),(F',\mathsf n')\rangle_{\mathcal H}=\langle F,\mathsf BF'\rangle+a\langle\mathsf n,\mathsf n'\rangle$ and the norm $\|\cdot\|_{\mathcal H}$. For $Z\in\Omega'^2$ and $\Xi\in\Omega^3$ let $Y=\mathbb Y(Z,\Xi):=W^*Z+\delta\Xi$, ${\rm pair}(Z,\Xi):=(\nu,\omega)=(\mathsf RY,a^{-1}\mathrm{d}'^*Z)$ and $\sigma=\ell\nu$. Then $\nu+\ell^{\mathsf T}\sigma=Y$, and the identity \[\|{\rm pair}(Z,\Xi)\|_{\mathcal H}^2=\langle Y,\mathsf RY\rangle+a^{-1}|\mathrm{d}'^*Z|^2=|\nu|^2+|\sigma|^2+a|\omega|^2\] holds.

We set $\mathcal A:=\{{\rm pair}(Z,\Xi):Z\in\Omega'^2,\ \Xi\in\Omega^3\}$ and call its elements admissible; further, $\mathcal K=\{(\mathrm{d}\varphi,-Q_1\varphi):\varphi\in\Omega^1\}$. For all $\varphi$, by $\mathrm{d}'Q_1=W\mathrm{d}$ of \Cref{Y-A-Facts} and $\mathrm{d}\mathrm{d}=0$, \begin{equation}\label{eq:hodge-orthogonality}\langle{\rm pair}(Z,\Xi),(\mathrm{d}\varphi,-Q_1\varphi)\rangle_{\mathcal H}=\langle Z,W\mathrm{d}\varphi-\mathrm{d}'Q_1\varphi\rangle+\langle\Xi,\mathrm{d}\mathrm{d}\varphi\rangle=0:\qquad\mathcal A\perp\mathcal K.\end{equation}

The adjoint of $Q_3\mathrm{d}=\mathrm{d}'W$ is $\delta Q_3^*=W^*\delta'$; together with $\mathrm{d}'^*\delta'=0$ it gives ${\rm pair}(-\delta'\Theta,Q_3^*\Theta)=0$ for all $\Theta\in\Omega'^3$. This identity is called the move. For $X'\in\Omega'^2$ we set $g(Z,\Xi;X')=\langle Z,X'\rangle-\frac12\|{\rm pair}(Z,\Xi)\|^2_{\mathcal H}$.

Let $\Omega_0^3=\ker Q_3$, and let $\Xi_Z\in\Omega_0^3$ be the minimiser of $\langle Y,\mathsf RY\rangle$ over $\Omega_0^3$. It is unique and linear in $Z$: for $\Xi\in\Omega_0^3$, $\delta\Xi=0$ makes $\Xi$ a constant $\kappa$, and $Q_3\kappa=L^3\kappa$ then forces $\kappa=0$. We set $\mathsf T_\sharp Z=W^*Z+\delta\Xi_Z$ and $R_\sharp=\mathsf T_\sharp^*\mathsf R\mathsf T_\sharp+a^{-1}\mathrm{d}'\mathrm{d}'^*$; thus $\langle Z,R_\sharp Z\rangle=\min_{\Xi\in\Omega^3_0}\|{\rm pair}(Z,\Xi)\|^2_{\mathcal H}$.\end{definition}

\begin{proposition}[The minimiser of one step]\label{Y-A-P4}\begin{enumerate}\item[\textup{(a)}] $R_\sharp>0$. Let $\mathsf B':=LR_\sharp^{-1}$, $\hat R=\mathsf B'^{-1}$, $H^\sharp:=L^{-1}\mathsf R\mathsf T_\sharp\mathsf B'$ and $\mathsf N^\sharp:=(aL)^{-1}\mathrm{d}'^*\mathsf B'$. For every $X'\in\Omega'^2$, the element $(\bar F,\bar{\mathsf n})=(H^\sharp X',\mathsf N^\sharp X')={\rm pair}(\bar Z,\Xi_{\bar Z})$, where $\bar Z=L^{-1}\mathsf B'X'$, satisfies $W\bar F+\mathrm{d}'\bar{\mathsf n}=X'$ and $\|(\bar F,\bar{\mathsf n})\|_{\mathcal H}^2=L^{-1}\langle X',\mathsf B'X'\rangle$.

\item[\textup{(b)}] $\hat R\le1-\mathsf P\le1$; the first inequality holds for every $a>0$, and the second, which is the inequality $\mathsf P\ge0$ of \Cref{Y-A-Facts}, holds if $aL\ge3/c_L$.

\item[\textup{(c)}] $\mathcal A=\mathcal K^\perp$, and $\mathcal K=\{(\mathrm{d}\varphi,-Q_1\varphi):\varphi\in\mathfrak L\}$. Let $\mathsf g=\Pi_{\mathfrak L}\mathrm{d}^*\mathsf BE$, $H_\circ:=E-\mathrm{d}q^{-1}\mathsf g$ and $\mathsf N_\circ:=Q_1q^{-1}\mathsf g$. Here $q>0$ on $\mathfrak L$ for every $a>0$, by the summed estimate and the bound $|f^\circ|^2+a|Q_1\varphi|^2\le q(\varphi)$ in \Cref{Y-A-L2}, which give $|\varphi|^2\le\max\{73L^4,6L^2/a\}\,q(\varphi)$ for $\varphi\in\mathfrak L$. For every $X'\in\Omega'^2$, $(H_\circ X',\mathsf N_\circ X')$ is the orthogonal projection of $(EX',0)$ on $\mathcal A$.\end{enumerate}\end{proposition}

\begin{proof}\emph{Step (a): positivity, and the minimiser.} To prove $R_\sharp>0$, let $\mathsf T_\sharp Z=0$ and $\mathrm{d}'^*Z=0$. We pair $\mathsf T_\sharp Z$ with $EX'$ and use \Cref{Y-A-L1}: $0=\langle Z,WEX'\rangle+\langle\Xi_Z,\mathrm{d}EX'\rangle=\langle Z+\delta'E_3^{\mathsf T}\Xi_Z,X'\rangle$ for all $X'$, so $Z=\delta'\Theta$ with $\Theta=-E_3^{\mathsf T}\Xi_Z$. Then, by $\delta Q_3^*=W^*\delta'$, $0=W^*\delta'\Theta+\delta\Xi_Z=\delta(Q_3^*\Theta+\Xi_Z)$; hence $Q_3^*\Theta+\Xi_Z=\kappa$ is a constant, and $\kappa=Q_3^*\kappa$, where on the right $\kappa$ is the coarse cochain with the same constant value. Thus $\Xi_Z=Q_3^*(\kappa-\Theta)$, and $\Xi_Z\in\ker Q_3\cap{\rm ran}\,Q_3^*=0$; so $Q_3Q_3^*(\Theta-\kappa)=0$, which gives $\Theta=\kappa$ because $Q_3Q_3^*=L^3\hat{\mathsf J}$ is invertible, and $Z=0$ because $\delta'\kappa=0$.

Next, $\sup_{\Xi\in\Omega_0^3}g=\langle Z,X'\rangle-\frac12\langle Z,R_\sharp Z\rangle$, which is maximal at $\bar Z=R_\sharp^{-1}X'=L^{-1}\mathsf B'X'$. As the maximum over both variables is attained at $(\bar Z,\Xi_{\bar Z})$, the function $g$ is stationary in $Z$ at fixed $\Xi=\Xi_{\bar Z}$; this gives $W\mathsf R\bar Y+a^{-1}\mathrm{d}'\mathrm{d}'^*\bar Z=X'$, with $\bar Y=\mathbb Y(\bar Z,\Xi_{\bar Z})$, which is $W\bar F+\mathrm{d}'\bar{\mathsf n}=X'$. Finally, $\|(\bar F,\bar{\mathsf n})\|_{\mathcal H}^2=\langle\bar Z,R_\sharp\bar Z\rangle=L^{-1}\langle X',\mathsf B'X'\rangle$.

\emph{Step (b): the bound on $\hat R$.} With the competitor $\Xi=0$ and $\mathsf R\le1$, $\langle Z,R_\sharp Z\rangle\le|W^*Z|^2+a^{-1}|\mathrm{d}'^*Z|^2=L\langle Z,(1-\mathsf P)Z\rangle$; since $\hat R=L^{-1}R_\sharp$, this is the first inequality of \textup{(b)}.

\emph{Step (c): admissible pairs, and the projection.} The relation $\mathcal A\perp\mathcal K$ in the display of \Cref{Y-A-Pairs} gives $\mathcal A\subset\mathcal K^\perp$. For the converse inclusion $\mathcal K^\perp\subset\mathcal A$, an element of $\mathcal K^\perp$ has the form $(\mathsf RG,\mathsf n)$ with $\langle G,\mathrm{d}\varphi\rangle=a\langle\mathsf n,Q_1\varphi\rangle$ for all $\varphi\in\Omega^1$, and we show that $(\mathsf RG,\mathsf n)\in\mathcal A$. The range of the map $(Z,\Xi)\mapsto(W^*Z+\delta\Xi,\mathrm{d}'^*Z)$ is the orthogonal complement of $\{(F,n):WF+\mathrm{d}'n=0,\ \mathrm{d}F=0\}$. For such $(F,n)$ we have $F=\mathrm{d}\varphi+h$ with $h$ constant. Indeed,
on a nonzero Fourier mode the coboundary is exterior multiplication by
a nonzero vector $v$. Contraction by $\bar v/|v|^2$ gives a contracting
homotopy, since its anticommutator with exterior multiplication by $v$
is the identity. Thus closed forms are exact on nonzero modes. The zero
mode is constant, and an exact form has zero constant mode. Since
$WF=\mathrm{d}'Q_1\varphi+L^2h$ and $WF+\mathrm{d}'n=0$, the constant mode gives $h=0$.
Then $\mathrm{d}'(n+Q_1\varphi)=0$, so the same argument on the coarse torus gives
$n=-Q_1\varphi+\mathrm{d}'\lambda'+h'$ with $h'$ constant. Then $\langle G,F\rangle+a\langle\mathsf n,n\rangle=a\langle\mathsf n,\mathrm{d}'\lambda'+h'\rangle=0$, the last equality by the tests $\varphi=\mathsf J\mathrm{d}'\lambda'$ and $\varphi=\mathsf Jh'$, for which $\mathrm{d}\varphi=E\mathrm{d}'(\cdot)=0$, while $Q_1\varphi=\mathrm{d}'\lambda'$ for the first test and $Q_1\varphi=h'$ for the second. Hence $(G,a\mathsf n)$ lies in that range, and $(\mathsf RG,\mathsf n)\in\mathcal A$.

For the slice, every $\varphi\in\Omega^1$ is of the form $\varphi=\varphi_{\mathfrak L}+\mathrm{d}\lambda$, with $\lambda(x)=\varphi(\Gamma(y,x))+\kappa_y$ on $B(y)$: a constant $\kappa_y$ on a block alters $\mathrm{d}\lambda$ on crossing bonds only, so $\varphi-\mathrm{d}\lambda\in\mathfrak L$ for every choice of the $\kappa_y$. We choose them with $Q_0\lambda=0$; then $Q_1\mathrm{d}\lambda=\mathrm{d}'Q_0\lambda=0$ and $\mathrm{d}\mathrm{d}\lambda=0$, so that $(\mathrm{d}\varphi,-Q_1\varphi)=(\mathrm{d}\varphi_{\mathfrak L},-Q_1\varphi_{\mathfrak L})$, and $\varphi\in\mathfrak L$ suffices to describe $\mathcal K$.

For the projection, $\bar\eta=-q^{-1}\mathsf gX'$ minimises \[\frac12\|(EX'+\mathrm{d}\eta,-Q_1\eta)\|^2_{\mathcal H}=\frac12\langle EX',\mathsf BEX'\rangle+\langle\mathsf gX',\eta\rangle+\frac12q(\eta)\] over $\eta\in\mathfrak L$; hence the orthogonal projection of $(EX',0)$ on $\mathcal K^\perp=\mathcal A$ is $(EX'+\mathrm{d}\bar\eta,-Q_1\bar\eta)=(H_\circ X',\mathsf N_\circ X')$.\end{proof}

\begin{proposition}[$\mathsf B'$ as one plus a sum of squares of the next rows]\label{Y-A-P5}
Suppose that $aL\ge3/c_L$, and put $\alpha=c_L-\frac2{aL}$, $\beta=\frac12c_L^2$ and $u,v=\frac12(\sqrt\alpha\pm\sqrt{\alpha-4\beta})$. Then $\alpha>4\beta>0$, and therefore $u$ and $v$ are real. Moreover $\mathsf B'=1+\ell'^{\mathsf T}\ell'$, where $\ell'_{y'}\zeta\in\mathsf H'_{y'}$ is the list whose entries are those of \textup{(1)} to \textup{(5)}:
\begin{enumerate}
\item[\textup{(1)}] the six numbers $((u+v\mathsf S_{\rho'})\partial'_\rho\zeta_{\mu\nu})(y')$, for $\mu<\nu$ and $(\rho,\rho')\in\{(\mu,\nu),(\nu,\mu)\}$;
\item[\textup{(2)}] the three numbers $(aL)^{-1/2}(\mathrm{d}^\dagger\zeta)(y',\cdot)$;
\item[\textup{(3)}] the numbers $L^{1/2}((H^\sharp-L^{-1}W^*)\zeta)_p$, one for each $p$ based in $B(y')$;
\item[\textup{(4)}] the numbers $(aL)^{1/2}((\mathsf N^\sharp-(aL)^{-1}\mathrm{d}'^*)\zeta)_c$, one for each $c$ based at $y'$;
\item[\textup{(5)}] the entries $L^{1/2}\ell_yH^\sharp\zeta$, one for each $y\in B(y')$.
\end{enumerate}
Here $\mathrm{d}^\dagger$ is obtained from $\mathrm{d}'^*$ by reversing the sign of the second of its two terms, and $(\partial'_\rho\zeta_{\mu\nu})(y')=\zeta_{\mu\nu}(y'+Le_\rho)-\zeta_{\mu\nu}(y')$; the same formula defines $\partial'_\rho$ on every function on $\mathbb T'$, and $\partial'^*_\rho\partial'_\rho=\partial'_\rho\partial'^*_\rho=\Delta'_\rho$.
\end{proposition}

\begin{proof}
Let $(\bar F,\bar{\mathsf n})$ be the pair of \Cref{Y-A-P4}\,(a) at $X'=\zeta$. Expanding the squares, we find \[L|\bar F-L^{-1}W^*\zeta|^2+aL|\bar{\mathsf n}-(aL)^{-1}\mathrm{d}'^*\zeta|^2=L|\bar F|^2+aL|\bar{\mathsf n}|^2-2\langle\zeta,W\bar F+\mathrm{d}'\bar{\mathsf n}\rangle+\langle\zeta,(1-\mathsf P)\zeta\rangle.\]

To both members of this identity we add $|\zeta|^2+\langle\zeta,\mathsf P\zeta\rangle+L|\ell\bar F|^2$; the resulting sum is $L\|(\bar F,\bar{\mathsf n})\|^2_{\mathcal H}=\langle\zeta,\mathsf B'\zeta\rangle$.

We now turn to $\mathsf P$, for which we use the identity \[1-(1-c_L\Delta'_\mu)(1-c_L\Delta'_\nu)=c_L\Delta'_\mu(1-\frac{c_L}2\Delta'_\nu)+c_L\Delta'_\nu(1-\frac{c_L}2\Delta'_\mu).\] Every component of $\mathrm{d}'^*\zeta$ has the form $s_1A+s_2B$, where $s_1,s_2\in\{+1,-1\}$ and $A,B$ are two of the six functions $\partial'^*_\rho\zeta_{\mu\nu}$, and each of the six appears in exactly one component. From $(s_1A+s_2B)^2+(s_1A-s_2B)^2=2A^2+2B^2$ we get \[|\mathrm{d}'^*\zeta|^2=2\sum\langle\zeta_{\mu\nu},(\Delta'_\mu+\Delta'_\nu)\zeta_{\mu\nu}\rangle-|\mathrm{d}^\dagger\zeta|^2,\] and consequently \[\mathsf P=\sum\partial'^*_\rho(\alpha-\beta\Delta'_{\rho'})\partial'_\rho+\frac1{aL}\mathrm{d}^{\dagger\mathsf T}\mathrm{d}^\dagger.\] Moreover, $(u+v\mathsf S)^*(u+v\mathsf S)=u^2+v^2+uv(\mathsf S+\mathsf S^*)=(u+v)^2-uv\Delta'=\alpha-\beta\Delta'$.

Finally, $\alpha>4\beta$ holds because $\frac2{aL}\le\frac23c_L<c_L(1-2c_L)$.
\end{proof}

\begin{definition}[The map to the next form]\label{Y-A-MapPhi} The form $\mathsf B'$ of \Cref{Y-A-P4} depends on the rows $\ell$ only through $\mathsf B$; we write $\Phi(\mathsf B):=\mathsf B'$. \end{definition}

\begin{definition}[The auxiliary form]\label{Y-A-AuxForm}
We put $\mathsf S_0:=L^{-5}$. On $\Omega'^2\oplus\Omega^3$ we define $\mathfrak R(Z,\Xi)=\|{\rm pair}(Z,\Xi)\|^2_{\mathcal H}+\mathsf S_0|Q_3\Xi|^2$, and we write $\iota X'=(X',0)$. Lems.~\ref{Y-A-L5p} and~\ref{Y-A-L7} are valid for every $\mathsf S_0>0$, with $\mathsf S_0$ kept in the formulas, except for the last member $\mathsf B'+\delta'\hat{\mathsf J}^{-2}\mathrm{d}'$ in the formula for $L\iota^{\mathsf T}\mathfrak R^{-1}\iota$, which is valid at $\mathsf S_0=L^{-5}$. The kernels $H^\sharp$ and $\mathsf B'$, as well as the kernels $\mathsf T$, $\mathsf X$, $\mathfrak G$, $\mathsf K_1$, $\mathsf K_2$ introduced below, are independent of $\mathsf S_0$; by contrast, $\mathsf G_Z$ and $\mathsf G_\Xi$ depend on it, through the move.
\end{definition}

\begin{lemma}[The penalty does not move the pair, at every depth]\label{Y-A-L5p}
For $X'\in\Omega'^2$ we have \begin{equation}\label{eq:constrained-response-formula}\FitDisplay{\dimexpr\linewidth-36pt\relax}{\mathfrak R^{-1}\iota X'=(\bar Z-\delta'\Theta,\ \Xi_{\bar Z}+Q_3^*\Theta),\ \ \Theta=-\mathsf S_0^{-1}(Q_3Q_3^*)^{-2}\mathrm{d}'X';\qquad H^\sharp=\mathsf R\mathbb Y\mathfrak R^{-1}\iota,\quad H_\circ=\mathsf R\mathbb Y\mathfrak R^{-1}\bigl(\iota+(0,E_3\mathrm{d}')\bigr);}\end{equation} moreover $L\iota^{\mathsf T}\mathfrak R^{-1}\iota=\mathsf B'+L\mathsf S_0^{-1}\delta'(Q_3Q_3^*)^{-2}\mathrm{d}'=\mathsf B'+\delta'\hat{\mathsf J}^{-2}\mathrm{d}'$, and the pair of $\mathfrak R^{-1}(X',E_3\mathrm{d}'X')$ is $(H_\circ X',\mathsf N_\circ X')$.
\end{lemma}

\begin{proof}
Given $(Z,\Xi)$, we write $\Xi=\Xi_0+Q_3^*\Theta$ with $\Xi_0\in\Omega_0^3$, the two summands being orthogonal, and we put $\tilde Z=Z+\delta'\Theta$. The move gives ${\rm pair}(Z,\Xi)={\rm pair}(\tilde Z,\Xi_0)$, while $Q_3\Xi=Q_3Q_3^*\Theta$. Therefore \[\langle Z,X'\rangle-\frac12\mathfrak R(Z,\Xi)=g(\tilde Z,\Xi_0;X')-\langle\Theta,\mathrm{d}'X'\rangle-\frac{\mathsf S_0}2|Q_3Q_3^*\Theta|^2.\] The maximiser is $\tilde Z=\bar Z$, $\Xi_0=\Xi_{\bar Z}$, with $\Theta$ as in the statement, and $\mathfrak R>0$ follows.

For the source $(X',E_3\mathrm{d}'X')$ there is one more term, $\langle\Xi,E_3\mathrm{d}'X'\rangle=\langle\Xi_0,E_3\mathrm{d}'X'\rangle+\langle\Theta,\mathrm{d}'X'\rangle$, because $Q_3E_3=1$; as a result, the maximiser has $\Theta=0$. Then $(\tilde Z,\Xi_0)$ maximises \[\langle\tilde Z,X'\rangle+\langle\Xi_0,E_3\mathrm{d}'X'\rangle-\frac12\|{\rm pair}\|^2=\langle{\rm pair}(\tilde Z,\Xi_0),(EX',0)\rangle_{\mathcal H}-\frac12\|{\rm pair}\|^2,\] the equality being a consequence of $\langle\mathsf RY,\mathsf BEX'\rangle=\langle Z,WEX'\rangle+\langle\Xi,\mathrm{d}EX'\rangle$; this identity holds by $\mathsf R\mathsf B=1$, $Y=W^*Z+\delta\Xi$ and adjointness, and by $WE=1$ and $\mathrm{d}E=E_3\mathrm{d}'$ of \Cref{Y-A-L1} its right side is $\langle Z,X'\rangle+\langle\Xi,E_3\mathrm{d}'X'\rangle$. As every pair is of the form ${\rm pair}(\tilde Z,\Xi_0)$, the maximiser is the projection of $(EX',0)$ on $\mathcal A$, as in \Cref{Y-A-P4}\,(c).
\end{proof}

\begin{definition}[The numbers $w_*$, $w_1$, $\mathsf r_\sharp$, $\gamma_1$, $\lambda_1$; the natural norm]\label{Y-A-Cst}
We set $w_*:=3^{-6}\le w_1:=(1-4c_L)^6$, and further $\mathsf r_\sharp:=w_*/(1.3\cdot10^7)$, $\gamma_1:=\mathsf r_\sharp/3074$ and $\lambda_1:=\frac1{40}\gamma_1^{1/2}$. The numbers $w_*$, $\mathsf r_\sharp$, $\gamma_1$ and $\lambda_1$ are free of $L$, $a$ and $N$, while $w_1$ depends on $L$ through $c_L$. The natural norm is given by $\mathsf D={\rm diag}(L\text{ on }\Omega'^2,\ L^{-2}\text{ on }\Omega^3)$, with $N(Z,\Xi)=L|Z|^2+L^{-2}|\Xi|^2$.
\end{definition}

\begin{proposition}[Poincar\'e bounds for empty rows and a lower bound for the auxiliary form]\label{Y-A-P6}
Let $\ell=\emptyset$, let $L\ge3$ be odd and $N\ge2$ even, and let $a>0$. Then the following hold.
\begin{enumerate}
\item[\textup{(a)}] For $u\in\Omega^3$, with $\bar u:=L^{-3}Q_3u$, we have $|u|^2\le768L^2|\delta u|^2+2L^{-3}|Q_3u|^2$.
\item[\textup{(b)}] $|\delta'\bar u|^2\le1.4\cdot10^5L^{-1}|\delta u|^2$.
\item[\textup{(c)}] For $Z\in\Omega'^2$ and $\Xi_0\in\Omega^3_0$ we have $|W^*Z+\delta\Xi_0|^2\ge\mathsf r_\sharp L|Z|^2$; consequently $\hat R\ge\mathsf r_\sharp$ and $\mathsf B'\le\mathsf r_\sharp^{-1}$.
\item[\textup{(d)}] $\mathfrak R\ge\gamma_1\mathsf D$.
\end{enumerate}
\end{proposition}

\begin{proof}
We prove \textup{(a)}. First, \[|u|^2=\sum_{q'}\sum_q\psi_{q'}(q)u_q^2\le2\sum_{q'}\sum_q\psi_{q'}(u_q-\bar u_{q'})^2+2L^3|\bar u|^2,\] and for every $q'$ \[\sum_q\psi_{q'}(u_q-\bar u_{q'})^2=\frac1{2L^3}\sum_{q,\tilde q}\psi_{q'}(q)\psi_{q'}(\tilde q)(u_q-u_{\tilde q})^2.\] Inside the box of side $<2L$ we join $q$ to $\tilde q$ by the coordinate path that moves first in the direction $e_1$, then in the direction $e_2$ and last in the direction $e_3$; the number of its steps is $<6L$. A given dual bond is crossed by at most $(2L)^4$ of these paths. Indeed, a bond along $e_1$ fixes $q_2,q_3$, and the free data are $q_1$ and $\tilde q$; a bond along $e_2$ fixes $\tilde q_1$ and $q_3$, and leaves $q_1,q_2,\tilde q_2,\tilde q_3$ free; a bond along $e_3$ fixes $\tilde q_1$ and $\tilde q_2$, and leaves $q$ and $\tilde q_3$ free. The double sum is therefore $\le\frac1{2L^3}6L(2L)^4\sum_{\rm box}(\delta u)^2=48L^2\sum_{\rm box}(\delta u)^2$, and each dual bond belongs to at most $8$ boxes.

We prove \textup{(b)}. For neighbouring $q',q''$, Jensen's inequality gives \[(\bar u_{q'}-\bar u_{q''})^2\le L^{-6}\sum_{q,\tilde q}\psi_{q'}(q)\psi_{q''}(\tilde q)(u_q-u_{\tilde q})^2.\]
For $N\ge4$, the two supports lie in a box of side less than $3L$
that embeds in the fine torus. The coordinate paths within the box have
length at most $9L$, and a dual bond belongs to at most $(3L)^4$ paths,
by the count in part~(a). Thus the last sum with its factor $L^{-6}$ is
bounded by $729L^{-1}$ times the sum of $(\delta u)^2$ in the box. Each dual
bond belongs to at most $3\cdot4^3=192$ such boxes. Thus the total
coefficient is at most $729\cdot192L^{-1}\le1.4\cdot10^5L^{-1}$.

For $N=2$, use the Jensen sum over the fine torus of side $2L$.
Join its two endpoints by positive coordinate steps in the order
$1,2,3$. The paths have length less than $6L$. A given dual bond
belongs to at most $(2L)^4$ paths, by the same count of free coordinates
as in part~(a). Since $\psi_{q'}\le1$, each coarse plaquette contributes
at most $6L(2L)^4L^{-6}|\delta u|^2=96L^{-1}|\delta u|^2$. There are
$24$ coarse plaquettes, so $2304L^{-1}$ suffices in this case. Part~(a)
applies unchanged, because its boxes have side $2L-1$.

We prove \textup{(c)}. We decompose $Z=Z_{\rm cl}+\delta'\Theta_1$ into two orthogonal parts, where $\mathrm{d}'Z_{\rm cl}=0$; then $Y=W^*Z+\delta\Xi_0=W^*Z_{\rm cl}+\delta u$, where $u=Q_3^*\Theta_1+\Xi_0$.

The closed part is bounded by means of a test field. For $z=y+r$ with $y\in L\mathbb Z^3$ and $0\le r_\rho<L$ we put \[(\mathsf W_2Z)(z;\mu\nu)=L^{-2}[(1-\frac{r_\lambda}L)Z(y;\mu\nu)+\frac{r_\lambda}LZ(y+Le_\lambda;\mu\nu)]\] and $(\mathsf W_3\Theta)(z)=L^{-3}\Theta(y)$. Then $\mathrm{d}\mathsf W_2=\mathsf W_3\mathrm{d}'$, because the difference along $e_\lambda$ equals $L^{-3}\partial'_\lambda Z_{\mu\nu}(y)$ for each value of $r_\lambda$, the value $r_\lambda=L-1$ included. By convexity, $|\mathsf W_2Z|^2\le L^{-1}|Z|^2$, the weight being $L^{-4}\cdot L^2\sum_{r_\lambda}((1-\frac{r_\lambda}L)+\frac{r_\lambda}L)$.

Next, $W\mathsf W_2=\prod_\rho(1-c''\Delta'_\rho)=:\mathsf K$, with $c''=\frac{L^2-1}{8L^2}$. We look at one coordinate at a time. Along $\mu$, for a function $g$ on the coarse sites and with $\lfloor z\rfloor$ denoting the coarse site $y'$ for which $z=y'+r$ and $0\le r<L$ in the coordinate at hand, the sum $\sum_zL^{-1}t(z-y)L^{-1}g(\lfloor z\rfloor)$ assigns to $g(y\pm L)$ the weight $L^{-2}(1+\dots+\mathsf h_L)=\frac{\mathsf h_L(\mathsf h_L+1)}{2L^2}=c''$. Along $\lambda$, the average $L^{-1}\sum_{|s|\le\mathsf h_L}$ of the linear interpolation assigns to $g(y\pm L)$ the weight $L^{-1}\sum_{s=1}^{\mathsf h_L}\frac sL=c''$. The weights sum to $1$. We have $\mathsf K\ge(1-4c'')^3\ge\frac18$; moreover $\mathsf K$ acts in the same manner on every component and commutes with translations, hence with $\mathrm{d}'$ and $\delta'$.

Hence $\langle Y,\mathsf W_2Z_{\rm cl}\rangle=\langle Z,\mathsf KZ_{\rm cl}\rangle+\langle\Xi_0,\mathsf W_3\mathrm{d}'Z_{\rm cl}\rangle=\langle Z_{\rm cl},\mathsf KZ_{\rm cl}\rangle\ge\frac18|Z_{\rm cl}|^2$, and therefore $|Y|^2\ge\frac L{64}|Z_{\rm cl}|^2$.

For the co-exact part, $|\delta u|\le|Y|+L^{1/2}|Z_{\rm cl}|\le9|Y|$. Further, $\bar u=L^{-3}Q_3Q_3^*\Theta_1=\hat{\mathsf J}\Theta_1$, and $\hat{\mathsf J}\ge w_1^{1/2}$ commutes with $\delta'$. Using \textup{(b)}, we get $w_1|\delta'\Theta_1|^2\le|\delta'\bar u|^2\le1.4\cdot10^5L^{-1}\cdot81|Y|^2$.

Adding the bounds for the two parts, we obtain $|Z|^2\le(64+1.134\cdot10^7/w_*)L^{-1}|Y|^2\le\mathsf r_\sharp^{-1}L^{-1}|Y|^2$. For the consequences concerning $\hat R$ and $\mathsf B'$ we note that $\langle Z,R_\sharp Z\rangle\ge|\mathsf T_\sharp Z|^2$.

We prove \textup{(d)}. Let $\Xi=\Xi_0+Q_3^*\Theta$, $\tilde Z=Z+\delta'\Theta$ and $Y=W^*\tilde Z+\delta\Xi_0$, as in the proof of \Cref{Y-A-L5p}; then $L|\tilde Z|^2\le\mathsf r_\sharp^{-1}|Y|^2$. Further, $|\delta\Xi_0|\le|Y|+L^{1/2}|\tilde Z|\le2\mathsf r_\sharp^{-1/2}|Y|$, so that \textup{(a)} gives $|\Xi_0|^2\le3072L^2\mathsf r_\sharp^{-1}|Y|^2$. As $\|\delta'\|^2\le12$, we have $L|Z|^2\le2L|\tilde Z|^2+24L|\Theta|^2$; in addition $|\Xi|^2\le|\Xi_0|^2+L^3|\Theta|^2$ and $|\Theta|^2\le|Q_3\Xi|^2/(L^6w_1)$. Altogether, \[N(Z,\Xi)\le3074\,\mathsf r_\sharp^{-1}|Y|^2+25L\cdot L^{-6}w_*^{-1}|Q_3\Xi|^2=3074\,\mathsf r_\sharp^{-1}|Y|^2+25w_*^{-1}\mathsf S_0|Q_3\Xi|^2\le\gamma_1^{-1}\mathfrak R(Z,\Xi).\]
\end{proof}

\begin{remark}[The torus of side two]\label{Y-A-C13}
\Cref{Y-A-P6} includes $N=2$. Its proof uses the whole fine
torus for the path estimate in part~(b) at that side.
\end{remark}

\begin{definition}[The kernels $\mathsf X$ and $\mathsf T$]\label{Y-A-XT}
For $v\in\Omega'^2$ we set $\mathsf Xv:=\Xi_{L^{-1}v}$ and $\mathsf Tv:=\mathsf R\mathsf T_\sharp L^{-1}v$. Thus ${\rm pair}(L^{-1}v,\mathsf Xv)=(\mathsf Tv,(aL)^{-1}\mathrm{d}'^*v)$, and $H^\sharp=\mathsf T\mathsf B'$.
\end{definition}

\begin{remark}[Decay of $\hat{\mathsf J}^{-2}$; the rate $\theta_J$]\label{Y-A-Jinv}
We have $\hat{\mathsf J}^2\ge w_*$; moreover $\hat{\mathsf J}^2$ has range $2$, and its absolute row sums are $\le(1+4c_L)^6\le25$. Applying \Cref{Y-A-L3}\,(a) to $\mathsf A=\hat{\mathsf J}^2$, for which $\varpi_{\hat{\mathsf J}^2}(\lambda)\le25(e^{2\lambda}-1)$, we obtain $|\hat{\mathsf J}^{-2}(i,j)|\le\frac2{w_*}e^{-\theta_J{\rm dist}}$, where $\theta_J:=\frac12\log(1+\frac{w_*}{50})$ is chosen so that $25(e^{2\theta_J}-1)=\frac12w_*$.
\end{remark}

\begin{definition}[The kernels with a source on coarse cubes]\label{Y-A-GG}
For $\theta\in\Omega'^3$ we define $(\mathsf G_Z\theta,\mathsf G_\Xi\theta):=\mathfrak R^{-1}(0,E_3\theta)$ and $\mathfrak G\theta:={\rm pair}(\mathsf G_Z\theta,\mathsf G_\Xi\theta)\in\mathcal A$. We denote by $\mathsf G\theta$ the first component of $\mathfrak G\theta$; its second component is given by $\mathsf G_{\mathsf n}:=a^{-1}\mathrm{d}'^*\mathsf G_Z$.
\end{definition}

\begin{lemma}[The projected lift: a minimiser plus a term in $\mathrm{d}'X'$]\label{Y-A-L7}
\begin{enumerate}
\item[\textup{(a)}] For all $(Z,\Xi)$ and all $\theta$, $\langle{\rm pair}(Z,\Xi),\mathfrak G\theta\rangle_{\mathcal H}=\langle E_3^{\mathsf T}\Xi,\theta\rangle-\langle Q_3\Xi,(Q_3Q_3^*)^{-1}\theta\rangle$.
\item[\textup{(b)}] $\|\mathfrak G\theta\|^2_{\mathcal H}=\langle\theta,\mathsf K_2\theta\rangle$, where $\mathsf K_2:=E_3^{\mathsf T}\mathsf G_\Xi-\mathsf S_0^{-1}(Q_3Q_3^*)^{-2}$.
\item[\textup{(c)}] For every $X'\in\Omega'^2$, $(H_\circ X',\mathsf N_\circ X')=(H^\sharp X',\mathsf N^\sharp X')+\mathfrak G\mathrm{d}'X'$; in particular $H_\circ=H^\sharp+\mathsf G\mathrm{d}'$. Moreover, with $\mathsf m_\circ(X'):=\frac12\|(H_\circ X',\mathsf N_\circ X')\|^2_{\mathcal H}=\frac12\langle EX',\mathsf BEX'\rangle-\frac12\langle\mathsf gX',q^{-1}\mathsf gX'\rangle$ and $\mathsf K_1:=E_3^{\mathsf T}\mathsf X\mathsf B'$, we have \[\mathsf m_\circ(X')=\tfrac1{2L}\langle X',\mathsf B'X'\rangle+\langle\mathsf K_1X',\mathrm{d}'X'\rangle+\tfrac12\langle \mathrm{d}'X',\mathsf K_2\mathrm{d}'X'\rangle .\]
\item[\textup{(d)}] For $v\in\Omega'^2$, $\|{\rm pair}(L^{-1}v,\mathsf Xv)\|_{\mathcal H}^2=L^{-1}\langle v,\hat Rv\rangle$.
\end{enumerate}
\end{lemma}

\begin{proof}
We prove \textup{(a)}. The vector $\mathfrak R^{-1}(0,E_3\theta)$ is the stationary point of $\langle\Xi,E_3\theta\rangle-\frac12\mathfrak R(Z,\Xi)$; this means that for all $(Z,\Xi)$ \[\langle\Xi,E_3\theta\rangle=\langle{\rm pair}(Z,\Xi),\mathfrak G\theta\rangle_{\mathcal H}+\mathsf S_0\langle Q_3\Xi,Q_3\mathsf G_\Xi\theta\rangle.\] At $(Z,\Xi)=(-\delta'\Theta,Q_3^*\Theta)$ the pair vanishes, and $Q_3E_3=1$; hence $\langle\Theta,\theta\rangle=\mathsf S_0\langle Q_3Q_3^*\Theta,Q_3\mathsf G_\Xi\theta\rangle$ for every $\Theta$, which gives $Q_3\mathsf G_\Xi\theta=\mathsf S_0^{-1}(Q_3Q_3^*)^{-1}\theta$.

Part \textup{(b)} is \textup{(a)} taken at $(Z,\Xi)=(\mathsf G_Z\theta,\mathsf G_\Xi\theta)$.

In \textup{(c)}, the first claim follows from \Cref{Y-A-L5p}, since $\mathfrak R^{-1}$ and ${\rm pair}$ are linear. To obtain the formula for $\mathsf m_\circ$ we expand the square: on the one hand $\|(H^\sharp X',\mathsf N^\sharp X')\|^2=L^{-1}\langle X',\mathsf B'X'\rangle$; on the other hand the cross term is given by \textup{(a)} at $(\bar Z,\Xi_{\bar Z})$, where $Q_3\Xi_{\bar Z}=0$ and $\Xi_{\bar Z}=\mathsf X\mathsf B'X'$. The second expression for $\mathsf m_\circ$ is the minimum in the proof of \Cref{Y-A-P4}\,(c).

Part \textup{(d)} is $\langle Z,R_\sharp Z\rangle$ evaluated at $Z=L^{-1}v$.
\end{proof}

\begin{definition}[Cut rows, masks, the number $\mathsf t_\rho$]\label{Y-A-CutRows}
The cut rows are $\ell'^{\rm c}_{y'}(P):=\ell'_{y'}(P)1\{{\rm dist}(y',\allowbreak P)\le\varrho(x')\}$, and we put $\ell'^>=\ell'-\ell'^{\rm c}$. Further, $\varsigma_>:=(\|\ell'^>\|_{\rm r}\|\ell'^>\|_{\rm c})^{1/2}\ge\|\ell'^>\|$ and $\delta_{\rm c}:=3\varsigma_>$. A mask $m$ is a map from the coarse sites to $\{0,1\}$; for a mask we set $\mathsf B'^{\rm c}_m:=1+\ell'^{{\rm c}\mathsf T}m\,\allowbreak \ell'^{\rm c}$, and $\mathsf B'^{\rm c}=\mathsf B'^{\rm c}_1$. Finally, \[\mathsf t_\rho:=\|\ell'^{\rm c}\|_{\rm r}\ \sup_{P_0}\sum_{y':\,{\rm dist}(y',P_0)>\rho_0}|\ell'^{\rm c}_{y'}(P_0)| .\]
\end{definition}

\begin{proposition}[Cutting and masking cost a factor $1+\delta_{\rm c}$]\label{Y-A-P8}
Suppose that $\varsigma_>\le1$.
For every mask, $\mathsf B'^{\rm c}_m\le(1+\delta_{\rm c})\mathsf B'$; consequently $\hat R\le(1+\delta_{\rm c})(\mathsf B_m'^{\rm c})^{-1}$.
\end{proposition}



\begin{proof}
If $\varsigma_>=0$, the claim reduces to $m\le1$. Otherwise, for $\delta>0$ we have $|m^{1/2}\ell'^{\rm c}\zeta|^2\le(1+\delta)|m^{1/2}\ell'\zeta|^2+(1+\delta^{-1})|\ell'^>\zeta|^2$. The choice $\delta=\varsigma_>$ yields $\mathsf B_m'^{\rm c}\le1+(1+\varsigma_>)\ell'^{\mathsf T}\ell'+(\varsigma_>^2+\varsigma_>)\le(1+3\varsigma_>)\mathsf B'$. Passing to inverses reverses the order, which gives the bound on $\hat R$.
\end{proof}

\begin{definition}[Uniform exponential decay of the sixteen kernels]\label{Y-A-Dec}
Condition \textup{(Dec)}$(C_D,\theta_D)$ holds if $C_D\ge4$, $\theta_D\in(0,1]$, $\varpi_q(\theta_D)\le\frac12\gamma_q$, and each of the kernels \[q,\ \ \mathsf g,\ \ \ell,\ \ \mathsf T,\ \ \ell\mathsf T,\ \ \mathsf X,\ \ \mathsf G,\ \ \ell\mathsf G,\ \ \mathsf G_Z,\ \ \mathsf G_\Xi,\ \ H^\sharp,\ \ \mathsf N^\sharp,\ \ \mathsf B',\ \ \ell',\ \ \mathsf K_1,\ \ \mathsf K_2\] satisfies $\|\cdot\|_{\theta_D}\le C_D$ on every torus of \Cref{Y-A-Lattices}. In that case we set $c_R:=64/\theta_D$.
\end{definition}

\begin{definition}[Sup bounds on the masked minimiser; two conditions on $L$]\label{Y-A-Np}
We are given rows $\ell$ at some level, the next rows $\ell'^{\rm c}$ and a radius $\rho_0$; the mask $m$ ranges over all masks, and $\zeta'$ over the elements of $\Omega'^2$ with $|\zeta'|_\infty\le1$. We put $Z=L^{-1}\mathsf B_m'^{\rm c}\zeta'$, $\Xi=\Xi_Z$, $(\nu,\omega)={\rm pair}(Z,\Xi)$ and $\sigma=\ell\nu$. Condition \textup{(N$'$)} holds if there exist numbers $\mathsf N_\sharp,\mathsf N_Z,\mathsf N_\Xi$, free of $N$ and $x$, such that \[L^2\bigl[|\nu|_\infty+\bigl(2\|\ell\|_{\rm c}+(2\rho_0+1)^3\|\ell\|_{\rm r}\bigr)|\sigma|_{2,\infty}\bigr]\le\mathsf N_\sharp,\qquad L|Z|_\infty\le\mathsf N_Z,\qquad L|\Xi|_\infty\le\mathsf N_\Xi .\] The two conditions on $L$ read \[\textbf{(L1)}\quad\mathsf c_{\rm loss}:=\bigl(\mathsf N_\sharp L^{-1}+60\,\mathsf N_Z+4\,\mathsf N_\Xi\bigr)L^{-1/2}\le\tfrac14d_0;\qquad\textbf{(L2)}\quad L\ge115\,\mathsf N_\sharp^2 .\] Since $\mathsf N_Z\ge1$, as the mask $m\equiv0$ gives $LZ=\zeta'$, condition \textup{(L1)} requires $L>(240/d_0)^2>2.3\cdot10^5$. For the mask $m\equiv1$ one has $\nu=H^\sharp\zeta'-\mathsf T(\mathsf B'-\mathsf B'^{\rm c})\zeta'$.
\end{definition}

\begin{definition}[Distances in the units of a level]\label{Y-A-OwnUnits} For a site $y$ and plaquettes $p,p'$ of one and the same level we put $|y-p|:=|y-\pi(p)|_\infty$ and $|p-p'|:=|\pi(p)-\pi(p')|_\infty$; these distances are periodic and measured in the units of that level. We write $\beta(p)\in\{1,2,3\}$ for the direction normal to $p$. For $F\in\Omega^2$, $F_\beta(z)$ denotes the value of $F$ at the plaquette based at $z$ with normal $e_\beta$. \end{definition}

\begin{definition}[The sums $\mathsf g_1$ and $\mathsf g_2$]\label{Y-A-Sums} We put \[\mathsf g_1(\lambda):=\sum_{u\in\mathbb Z^3}|u|_\infty e^{-\lambda|u|_\infty},\qquad\mathsf g_2(\lambda):=\sum_{u\in\mathbb Z^3}(1+|u|_\infty)^2e^{-\lambda|u|_\infty}.\] For $r\ge1$ there are $24r^2+2$ points $u\in\mathbb Z^3$ with $|u|_\infty=r$. For $0<\lambda\le1$ we have $\mathsf g(\lambda)\le51\lambda^{-3}$ and $\mathsf g_1(\lambda)\le182\lambda^{-4}$. Moreover, $\mathsf g(\lambda)\le\mathsf g_2(\lambda)$ and $\mathsf g_1(\lambda)\le\mathsf g_2(\lambda)$. \end{definition}

\begin{proof}[Proof of the bounds] Put $t=\lambda/2\le\frac12$; all sums below run over $r\ge1$. Since $e^\lambda-1\ge\lambda$ and $\sinh t\ge t$, we have $\sum e^{-\lambda r}=(e^\lambda-1)^{-1}\le\lambda^{-1}$ and $\sum re^{-\lambda r}=(4\sinh^2t)^{-1}\le\lambda^{-2}$. Next, $\sum r^2e^{-\lambda r}=\cosh t/(4\sinh^3t)\le(4t^3)^{-1}=2\lambda^{-3}$, where the inequality is equivalent to $\cosh t\le(\sinh t/t)^3$. This holds because \[\cosh t\le1+\frac{t^2}2+\frac{t^4}{23}\le1+\frac{t^2}2+\frac{13t^4}{120}\le(1+\frac{t^2}6+\frac{t^4}{120})^3\le(\sinh t/t)^3,\] the first and the last step by the power series of $\cosh t$ and $\sinh t$ for $t\le\frac12$, the third by expanding the cube.

Further, \[\sum r^3e^{-\lambda r}=(2\cosh2t+4)/(16\sinh^4t)\le[6(1+\frac{2t^2}3)+1.5t^4]/[16t^4(1+\frac{2t^2}3)]\le6\lambda^{-4}+0.1,\] where the middle step bounds $2\cosh2t+4$ from above and $\sinh^4t$ from below by power series, and the last step uses $16t^4=\lambda^4$. By the count of points, and since $\lambda\le1$, \[\mathsf g(\lambda)=1+24\sum r^2e^{-\lambda r}+2\sum e^{-\lambda r}\le1+48\lambda^{-3}+2\lambda^{-1}\le51\lambda^{-3}\] and \[\mathsf g_1(\lambda)=24\sum r^3e^{-\lambda r}+2\sum re^{-\lambda r}\le144\lambda^{-4}+2.4+2\lambda^{-2}\le182\lambda^{-4}.\] \end{proof}

\begin{definition}[Differences, gradient and Laplacian]\label{Y-A-Differences} For a cochain $F$ of any degree and a direction $\alpha$ we put $(\partial_\alpha F)_p=F_{p+e_\alpha}-F_p$, where $e_\alpha$ is the unit step of the level of $F$, and $|\nabla F|^2=\sum_\alpha|\partial_\alpha F|^2$. The operator $\nabla^*\nabla=\sum_\alpha\partial_\alpha^*\partial_\alpha$ acts on each component as the scalar Laplacian, whose symbol is $\hat\Delta(\kappa)=\sum_\alpha|e^{i\kappa_\alpha}-1|^2\le12$. We write $\nabla'$ for the same objects on $\mathbb T'$. \end{definition}

\begin{definition}[Translation covariant rows]\label{Y-A-CovRows} Rows $\ell$ are called translation covariant if for each translation $\mathsf t$ of the lattice and each site $y$ there is an orthogonal map $U_{\mathsf t,y}:\mathsf H_y\to\mathsf H_{y+\mathsf t}$ such that $\ell_{y+\mathsf t}(p+\mathsf t)=U_{\mathsf t,y}\ell_y(p)$ for all $p$. For such rows $\mathsf B$ and $\mathsf R$ commute with the translations, $|\ell_y(p)|$ does not change when $y$ and $p$ are moved by the same translation, and $|\ell_yh|$, for a constant form $h$, does not change when $y$ is moved by a translation. \end{definition}

\begin{definition}[The defect on constant forms]\label{Y-A-Defect} The defect of rows $\ell$ on constant forms is \[\tau(\ell):=\sup_y\sup\{|\ell_yh|:h\in\Omega^2\text{ constant},|h|_\infty\le1\}.\] \end{definition}

\begin{definition}[The sets of rows]\label{Y-A-SetsOfRows} Let $b\ge1$, $K\ge0$, $0<\lambda\le1$ and $\tau\ge0$. We denote by $\mathcal R(b,K,\lambda,\tau)$ the set of translation covariant rows $\ell$ on a torus with \[b_{\mathsf B}\le b,\qquad|\ell_y(p)|\le Ke^{-\lambda|y-p|}\ \text{ for all }y,p,\qquad\tau(\ell)\le\tau .\]

Rows in this set satisfy $\|\ell\|_{\rm r}\le3K\mathsf g(\lambda)$, $\|\ell\|_{\rm c}\le K\mathsf g(\lambda)$ and $|(\mathsf B-1)(p,p')|\le K^2\mathsf g(\frac\lambda2)e^{-\frac\lambda2|p-p'|}$. For the third bound we write $(\mathsf B-1)(p,p')=\sum_y\langle\ell_y(p),\ell_y(p')\rangle$ and use $\sum_ye^{-\lambda(|y-p|+|y-p'|)}\le e^{-\frac\lambda2|p-p'|}\sum_ye^{-\frac\lambda2|y-p|}$, which follows from the triangle inequality. The empty rows belong to the family: $\emptyset\in\mathcal R(1,0,\lambda,0)$. \end{definition}

\begin{definition}[Towers of tori and chains of rows]\label{Y-A-Towers} Let $N$ be even and $K\ge1$. The tower of tori consists of $\mathbb T_k:=\mathbb Z^3/NL^{K-k}\mathbb Z^3$, $0\le k\le K$, each in its own units, and step $k$ of the tower is the block step with $\mathbb T=\mathbb T_k$ and $\mathbb T'\cong\mathbb T_{k+1}$. Given radii $\varrho_1,\dots,\varrho_K\in\mathbb N\cup\{\infty\}$, the chain of rows starts with $\ell_0=\emptyset$ on $\mathbb T_0$ and continues with $\ell_{k+1}=((\ell_k)')^{\rm c}$ for $0\le k<K$, that is, with the rows of \Cref{Y-A-P5} built from $\ell_k$ and cut at $\varrho_{k+1}$ as in \Cref{Y-A-CutRows}; its forms are $\mathsf B_k=1+\ell_k^{\mathsf T}\ell_k$. The chain is uncut when $\varrho_k=\infty$ for all $k$. \end{definition}

\begin{remark}[Standing conditions]\label{Y-A-Standing} Throughout the part of the section that runs from the subsection ``Four tools'', which opens with \Cref{Y-A-L24}, to the end of the proof of \Cref{Y-A-Th36}, we assume that $L\ge3$ is odd, $aL\ge3/c_L$, $N\ge2$ is even, and $\ell$ are rows on $\mathbb T$. \Cref{Y-A-L31} has hypotheses of its own, and they take precedence over the standing conditions: in that lemma the side of the torus is an integer $\ge2$, $L$ is a real number with $1\le L\le$ half that side, and $\mathsf R$ is not tied to rows. \end{remark}

\subsubsection{Four tools}

\begin{lemma}[Weighted decay of a column of an inverse]\label{Y-A-L24} Let $I$ and $J$ be finite sets, each point of which is placed in a pseudo-metric space with distance $\mathsf d$. Let $\mathbb A$ be a real $J\times I$ matrix and $\mathfrak R=\mathbb A^{\mathsf T}\mathbb A$; let $\mathsf D$ be a positive diagonal matrix on $I$, with $N(\phi)=\langle\phi,\mathsf D\phi\rangle$; and let $\gamma>0$ be such that $\mathfrak R\ge\gamma\mathsf D$. Let $\lambda>0$ and $0\le\varkappa\le\frac14\gamma^{1/2}$ be such that \begin{equation}\tag{$\sinh$}\sum_{r\in J}\Bigl(\sum_{i\in I}|\mathbb A(r,i)|\sinh\bigl(\lambda\mathsf d(r,i)\bigr)|\phi_i|\Bigr)^2\le\varkappa^2N(\phi)\quad\text{for all }\phi.\end{equation} Then for every $j\in I$ the vector $\Phi:=\mathfrak R^{-1}\delta_j$ satisfies \[\text{(a)}\ \sum_{i\in I}\mathsf D_ie^{2\lambda\mathsf d(i,j)}|\Phi_i|^2\le4\gamma^{-2}\mathsf D_j^{-1};\qquad\text{(b)}\ \sum_{r\in J}e^{2\lambda\mathsf d(r,j)}|(\mathbb A\Phi)_r|^2\le4\gamma^{-1}\mathsf D_j^{-1}.\] If \textup{($\sinh$)} holds at $\lambda$, it holds with the same $\varkappa$ at every smaller $\lambda$. \end{lemma}

\begin{proof} Put $w:=\lambda\mathsf d(\cdot,j)$ on $I\cup J$. By the triangle inequality, $|w_i-w_r|\le\lambda\mathsf d(r,i)$ for $i\in I$ and $r\in J$. We let $e^{\pm w}$ act on $I$ as diagonal matrices and put $\mathfrak R_w:=e^w\mathfrak Re^{-w}$.

For real $\phi$ put $\mathsf C_r=\sum_i\mathbb A(r,i)\cosh(w_i-w_r)\phi_i$ and $\mathsf S_r=\sum_i\mathbb A(r,i)\sinh(w_i-w_r)\phi_i$, and let $\mathsf C$ and $\mathsf S$ be the vectors with these entries. Writing $e^{\pm(w_i-w_r)}=\cosh(w_i-w_r)\pm\sinh(w_i-w_r)$, we find $e^{-w_r}(\mathbb Ae^w\phi)_r=\mathsf C_r+\mathsf S_r$ and $e^{w_r}(\mathbb Ae^{-w}\phi)_r=\mathsf C_r-\mathsf S_r$. Consequently $\langle\phi,\mathfrak R_w\phi\rangle=\langle\mathbb Ae^w\phi,\mathbb Ae^{-w}\phi\rangle=|\mathsf C|^2-|\mathsf S|^2$.

Since $|w_i-w_r|\le\lambda\mathsf d(r,i)$ and $\cosh t-1\le\sinh|t|$, both $|\sinh(w_i-w_r)|$ and $\cosh(w_i-w_r)-1$ are at most $\sinh(\lambda\mathsf d(r,i))$; so \textup{($\sinh$)} yields $|\mathsf S|,|\mathsf C-\mathbb A\phi|\le\varkappa N(\phi)^{1/2}$. From $\mathfrak R\ge\gamma\mathsf D$ and $\varkappa\le\frac14\gamma^{1/2}$ we get $|\mathbb A\phi|\ge\gamma^{1/2}N(\phi)^{1/2}\ge4\varkappa N(\phi)^{1/2}$. Therefore \[\langle\phi,\mathfrak R_w\phi\rangle\ge(|\mathbb A\phi|-\varkappa N(\phi)^{1/2})^2-\varkappa^2N(\phi)\ge\frac12|\mathbb A\phi|^2\ge\frac\gamma2N(\phi),\] where the first step uses $|\mathsf C|\ge|\mathbb A\phi|-\varkappa N(\phi)^{1/2}\ge0$ and the second uses $\varkappa N(\phi)^{1/2}\le\frac14|\mathbb A\phi|$.

We now choose $\phi=\mathfrak R_w^{-1}\delta_j$. Since $w_j=0$, we have $e^{-w}\delta_j=\delta_j$, and therefore $\phi=e^w\mathfrak R^{-1}\delta_j=e^w\Phi$.

For (a): as $\mathfrak R_w\phi=\delta_j$, the lower bound just proved gives $\frac\gamma2N(\phi)\le\langle\phi,\delta_j\rangle\le N(\phi)^{1/2}\mathsf D_j^{-1/2}$, the second step because $\mathsf D_j\phi_j^2\le N(\phi)$; hence $N(\phi)\le4\gamma^{-2}\mathsf D_j^{-1}$. Since $\phi_i=e^{\lambda\mathsf d(i,j)}\Phi_i$, the number $N(\phi)$ is the left side of (a).

For (b): now $e^{w_r}(\mathbb A\Phi)_r=\mathsf C_r-\mathsf S_r$, so the left side of (b) equals $|\mathsf C-\mathsf S|^2\le(|\mathsf C|+|\mathsf S|)^2$. By (a) and $\varkappa^2\le\gamma/16$, \[|\mathsf S|^2\le\varkappa^2N(\phi)\le\frac\gamma{16}\cdot\frac4{\gamma^2}\mathsf D_j^{-1}=\frac1{4\gamma}\mathsf D_j^{-1}.\] Moreover $|\mathsf C|^2-|\mathsf S|^2=\langle\phi,\mathfrak R_w\phi\rangle=\langle\phi,\delta_j\rangle$, and $\langle\phi,\delta_j\rangle\le N(\phi)^{1/2}\mathsf D_j^{-1/2}\le\frac2\gamma\mathsf D_j^{-1}$ by the bounds for (a); so \[|\mathsf C|^2=\langle\phi,\delta_j\rangle+|\mathsf S|^2\le(\frac2\gamma+\frac1{4\gamma})\mathsf D_j^{-1}.\] Hence $|\mathsf C|+|\mathsf S|\le(\frac32+\frac12)\gamma^{-1/2}\mathsf D_j^{-1/2}$, and squaring gives (b). \end{proof}

In the rest of this section, $\mathfrak R$ denotes the form of \Cref{Y-A-AuxForm} on $\Omega'^2\oplus\Omega^3$, and $\mathsf D$ the matrix of \Cref{Y-A-Cst}, outside \Cref{Y-A-L24}, where both matrices are general.

\begin{lemma}[Unfolding $\mathsf R$ by one more variable]\label{Y-A-L25} On the space $\Omega'^2\oplus\Omega^3\oplus\mathsf H$ we define \[\FitDisplay{\dimexpr\linewidth-1pt\relax}{\mathbb A^{\rm e}(Z,\Xi,\sigma):=\bigl(W^*Z+\delta\Xi-\ell^{\mathsf T}\sigma,\ \ \sigma,\ \ a^{-1/2}\mathrm{d}'^*Z,\ \ \mathsf S_0^{1/2}Q_3\Xi\bigr),\qquad\mathfrak R^{\rm e}:=\mathbb A^{{\rm e}\mathsf T}\mathbb A^{\rm e},\qquad\mathsf D^{\rm e}:={\rm diag}(L,L^{-2},1).}\] The rows of $\mathbb A^{\rm e}$ are indexed by the plaquettes $p$, by the components of the spaces $\mathsf H_y$, each component being placed at its site $y$, by the coarse bonds $c$, and by the coarse cubes $q'$.

\begin{enumerate} \item[\textup{(a)}] Let $\psi\in\Omega'^2\oplus\Omega^3$, and let $(Z,\Xi)=\mathfrak R^{-1}\psi$ and $\nu=\mathsf R\mathbb Y(Z,\Xi)$. Then $(\mathfrak R^{\rm e})^{-1}(\psi,0)=(Z,\Xi,\ell\nu)$ and $\mathbb A^{\rm e}(Z,\Xi,\ell\nu)=(\nu,\ell\nu,a^{-1/2}\mathrm{d}'^*Z,\mathsf S_0^{1/2}Q_3\Xi)$. In the case $\psi=\iota\delta_P$ one has $\nu=H^\sharp\delta_P$ and $a^{-1}\mathrm{d}'^*Z=\mathsf N^\sharp\delta_P$.

\item[\textup{(b)}] The form $\mathfrak R^{\rm e}$ satisfies $\mathfrak R^{\rm e}\ge\gamma_{\rm e}\mathsf D^{\rm e}$, where $\gamma_{\rm e}:=\gamma_1/2b_{\mathsf B}$; this stays true if any $b\ge b_{\mathsf B}$ takes the place of $b_{\mathsf B}$.

\item[\textup{(c)}] Let $|\ell_y(p)|\le Ke^{-\lambda|y-p|}$ and $b_{\mathsf B}\le b$, let $\gamma_{\rm e}=\gamma_1/2b$, and put \[\lambda':=\min\Bigl\{\frac{\gamma_{\rm e}^{1/2}}{44},\ \frac{\lambda L}2,\ \frac{\lambda L\,\gamma_{\rm e}^{1/2}}{12K\mathsf g(\lambda/2)}\Bigr\}\quad(\text{the first alone if }K=0).\] Then \textup{($\sinh$)} holds for $\mathbb A^{\rm e}$ and $\mathsf D^{\rm e}$, with the distance $\mathsf d$ of \Cref{Y-A-Distances}, at the rate $\lambda'$ and with a constant $\varkappa\le\frac14\gamma_{\rm e}^{1/2}$. \end{enumerate} \end{lemma}

\begin{proof} \textup{(a)} Fix $Y$. As a function of $\sigma\in\mathsf H$, the quantity $|Y-\ell^{\mathsf T}\sigma|^2+|\sigma|^2$ takes its least value at the point where $\sigma=\ell(Y-\ell^{\mathsf T}\sigma)$. We write $\nu:=Y-\ell^{\mathsf T}\sigma$ at this point. Then $\sigma=\ell\nu$, hence $\nu+\ell^{\mathsf T}\ell\nu=Y$, which means $\nu=\mathsf RY$; and the least value is $|\nu|^2+|\ell\nu|^2=\langle\nu,\mathsf B\nu\rangle=\langle Y,\mathsf RY\rangle$.

Now let $Y=\mathbb Y(Z,\Xi)$. By the definition of $\mathbb A^{\rm e}$, the value $\mathfrak R^{\rm e}(Z,\Xi,\sigma)$ is the sum of $|Y-\ell^{\mathsf T}\sigma|^2+|\sigma|^2$ and of $a^{-1}|\mathrm{d}'^*Z|^2+\mathsf S_0|Q_3\Xi|^2$, which does not involve $\sigma$; by \Cref{Y-A-Pairs} of ${\rm pair}$ and \Cref{Y-A-AuxForm} of $\mathfrak R$, the value $\mathfrak R(Z,\Xi)$ is the sum of $\langle Y,\mathsf RY\rangle$ and of the same two terms. Therefore $\min_\sigma\mathfrak R^{\rm e}(Z,\Xi,\sigma)=\mathfrak R(Z,\Xi)$. The vector $(\mathfrak R^{\rm e})^{-1}(\psi,0)$ maximises $\langle\psi,(Z,\Xi)\rangle-\frac12\mathfrak R^{\rm e}$. We may maximise first over $\sigma$, which yields $\sigma=\ell\nu$, and then maximise $\langle\psi,(Z,\Xi)\rangle-\frac12\mathfrak R(Z,\Xi)$, which yields $(Z,\Xi)=\mathfrak R^{-1}\psi$; so the maximiser is $(Z,\Xi,\ell\nu)$ as in \textup{(a)}. The first component of $\mathbb A^{\rm e}(Z,\Xi,\ell\nu)$ is $Y-\ell^{\mathsf T}\ell\nu=\nu$, and its other components are read off from the definition.

For the case $\psi=\iota\delta_P$, \Cref{Y-A-L5p} gives $H^\sharp=\mathsf R\mathbb Y\mathfrak R^{-1}\iota$, so that $\nu=\mathsf R\mathbb Y\mathfrak R^{-1}\iota\delta_P=H^\sharp\delta_P$; the same lemma gives $Z=\bar Z-\delta'\Theta$, with $\bar Z$ and $\Theta$ evaluated at $X'=\delta_P$. As $\mathrm{d}'^*\delta'=0$ and $\bar Z=L^{-1}\mathsf B'\delta_P$, we obtain $a^{-1}\mathrm{d}'^*Z=a^{-1}\mathrm{d}'^*\bar Z=\mathsf N^\sharp\delta_P$.

\textup{(b)} The minimum over $\sigma$ found above shows that $\mathfrak R^{\rm e}\ge\mathfrak R$ as forms in $(Z,\Xi,\sigma)$. Let $\mathfrak R_{\ell=\emptyset}$ be the form $\mathfrak R$ at $\ell=\emptyset$, where $\mathsf R=1$. We apply the bound $\mathsf R\ge b_{\mathsf B}^{-1}$ to the term $\langle Y,\mathsf RY\rangle$, and $b_{\mathsf B}\ge1$ to the two remaining terms; together with \Cref{Y-A-P6}\,(d) this gives \[\mathfrak R^{\rm e}\ge\mathfrak R\ge b_{\mathsf B}^{-1}\mathfrak R_{\ell=\emptyset}\ge b_{\mathsf B}^{-1}\gamma_1\mathsf D.\] Here $\gamma_1$ and $\mathsf D$ are those of \Cref{Y-A-Cst}, with $\mathsf D={\rm diag}(L\text{ on }\Omega'^2,\ L^{-2}\text{ on }\Omega^3)$. On the other hand $\mathfrak R^{\rm e}\ge|\sigma|^2$, since $\sigma$ is the second component of $\mathbb A^{\rm e}(Z,\Xi,\sigma)$. Half the sum of these two lower bounds is at least $\gamma_{\rm e}\mathsf D^{\rm e}$, because $\gamma_1/b_{\mathsf B}\le1$.

\textup{(c)} We begin by bounding, in each row of $\mathbb A^{\rm e}$, the distance from the row to every entry that occurs in it; the facts needed come from \Cref{Y-A-Facts} and \Cref{Y-A-Distances}. Write $p=(z;\mu\nu)$ and $P=(y;\mu\nu)$, and let $\lambda$ denote the direction normal to $p$. For the kernel supports, take the representatives of the periodic displacements from \Cref{Y-A-Facts}. We have $W(P,p)\ne0$ only when both $z_\mu-y_\mu$ and $z_\nu-y_\nu$ belong to $[-\mathsf h_L,3\mathsf h_L]$ while $|z_\lambda-y_\lambda|\le\mathsf h_L$; moreover $3\mathsf h_L/L<\frac32$. The two cubes of $p$ have the base points $z$ and $z-e_\lambda$. Every row $(y,\alpha)$ has a single entry, namely $\sigma_{y,\alpha}$, which is located at $y$; in these rows of $\sigma$, and in $\mathsf d(p,y)$ below, $y$ is a fine site. The four $P\ni c$ have their base points at $c_-$ or at $c_--Le_\nu$. Lastly, if a fine cube $q$ has its base point at $z$ and a coarse cube $q'$ at $y$, as in \Cref{Y-A-Facts}, then $\psi_{q'}(q)\ne0$ only if $z-y\in[-\mathsf h_L,3\mathsf h_L]^3$. It follows that $\mathsf d(p,P)\le\frac32$ whenever $W(P,p)\ne0$, that the distance between $p$ and each of its two cubes is $\le\frac1L$, that $\mathsf d(p,y)=|y-p|/L$, that in the rows of $\sigma$ the distance is $0$, that the distance between $c$ and each of the four $P\ni c$ is $\le1$, and that the distance is $\le\frac32$ whenever $\psi_{q'}(q)\ne0$.

Put $\mathsf s:=\sinh(\frac32\lambda')$; then $\mathsf s\le1.8\lambda'$, because $\lambda'\le\gamma_{\rm e}^{1/2}/44\le\frac1{44}$. Since $\sinh$ is convex on $[0,\infty)$ and vanishes at $0$, we have $\sinh(\theta t)\le\theta\sinh t$ for $0\le\theta\le1$ and $t\ge0$. With $\theta=2/(3L)$ and $t=\frac32\lambda'$, this bound gives $\sinh(\lambda'/L)\le\frac2{3L}\mathsf s\le\mathsf s/L$. We put $\theta_L:=2\lambda'/\lambda L$ for every $K$. If $K>0$, then $\theta_L\le1$ by the second entry of the minimum that defines $\lambda'$, and the bound $\sinh(\theta t)\le\theta\sinh t$ at $\theta=\theta_L$ gives $\sinh(\lambda'n/L)\le\theta_L\sinh(\frac\lambda2n)\le\frac12\theta_Le^{\lambda n/2}$ for $n\ge0$. If $K=0$, then $\ell=0$, the rows $p$ contain no entry of $\sigma$, and this bound is not used.

Consider now the row of a plaquette $p$. The components of $\sigma_y$ occur in it with the coefficients of $\ell_y(p)$, and they are controlled through $\sum_\alpha|\ell_y(p)_\alpha||\sigma_{y,\alpha}|\le|\ell_y(p)||\sigma_y|$. Let $x_1$, $x_2$ and $x_3$ be the parts of the sum of this row in \textup{($\sinh$)} that come from the entries of $Z$, of $\Xi$ and of $\sigma$; we use $(x_1+x_2+x_3)^2\le3\sum x_i^2$. Since $\mathsf d(p,P)\le\frac32$ where $W(P,p)\ne0$, and since $\sum_PW(P,p)=L^{-1}$, the Cauchy--Schwarz inequality gives $x_1^2\le\mathsf s^2L^{-1}\sum_PW(P,p)Z_P^2$. The entries of $\Xi$ in this row belong to the two cubes $q$ and $\tilde q$ of $p$, carry the coefficients $\pm1$ and lie at distance at most $\frac1L$, whence $x_2^2\le2\mathsf s^2L^{-2}(\Xi_q^2+\Xi_{\tilde q}^2)$. Finally, $x_3\le\sum_y|\ell_y(p)|\sinh(\lambda'|y-p|/L)|\sigma_y|$, by the Cauchy--Schwarz bound over the components $\alpha$ and because $\mathsf d(p,y)=|y-p|/L$.

We now sum over $p$. For the first term we use $\sum_pW(P,p)=L^2$, and for the second the fact that every cube has six faces. For the third we apply Schur's test: by the hypothesis on $\ell$ and the bound obtained from convexity, the sums of its kernel are $\le\frac12\theta_LK\mathsf g(\frac\lambda2)$ over $y$ and $\le\frac32\theta_LK\mathsf g(\frac\lambda2)$ over $p$. Altogether the rows $p$ contribute at most \[3\mathsf s^2L|Z|^2+36\mathsf s^2L^{-2}|\Xi|^2+\frac94\theta_L^2K^2\mathsf g(\frac\lambda2)^2|\sigma|^2.\]

The row of a coarse bond $c$ contains $a^{-1/2}$ times the entries $\pm Z_P$ of the four $P\ni c$, at distance at most $1$. By the Cauchy--Schwarz inequality, and as every $P$ has four bonds, the rows $c$ contribute at most $16a^{-1}\mathsf s^2|Z|^2\le\mathsf s^2L|Z|^2$; the last step holds because $aL\ge3/c_L>18\ge16$ by the standing conditions of \Cref{Y-A-Standing}. The row of a coarse cube $q'$ contains the entries $\mathsf S_0^{1/2}\psi_{q'}(q)\Xi_q$, at distance at most $\frac32$; by the Cauchy--Schwarz inequality and $\sum_q\psi_{q'}(q)=L^3$, the square of its sum is at most $\mathsf S_0\mathsf s^2L^3\sum\psi_{q'}\Xi^2$, and the total over all $q'$, by $\sum_{q'}\psi_{q'}=1$ and $\mathsf S_0=L^{-5}$, is $\mathsf s^2L^{-2}|\Xi|^2$. The rows of $\sigma$ contribute nothing, as their distances vanish.

Adding the contributions of all rows and comparing with $\mathsf D^{\rm e}={\rm diag}(L,L^{-2},1)$, we find that \textup{($\sinh$)} holds with $\varkappa=\max\{\sqrt{37}\,\mathsf s,\frac32\theta_LK\mathsf g(\frac\lambda2)\}$. For the first member, $\sqrt{37}\,\mathsf s\le\sqrt{37}\cdot1.8\lambda'\le10.95\lambda'\le\frac14\gamma_{\rm e}^{1/2}$, by the first entry of the minimum that defines $\lambda'$. For the second member, $\frac32\theta_LK\mathsf g(\frac\lambda2)=3\lambda'K\mathsf g(\frac\lambda2)/\lambda L\le\frac14\gamma_{\rm e}^{1/2}$; for $K>0$ this follows from the third entry of that minimum, and for $K=0$ the second member is $0$, so that $\varkappa=\sqrt{37}\,\mathsf s$. \end{proof}

\begin{lemma}[Admitted pairs and the minimiser]\label{Y-A-L26} Let $X'\in\Omega'^2$. We call a pair $(F,\mathsf n)\in\Omega^2\oplus\Omega'^1$ \emph{admitted at $X'$} when $WF+\mathrm{d}'\mathsf n=X'$ and, for every $\Xi\in\ker Q_3$, $\langle\Xi,\mathrm{d}F\rangle=0$. The notion is independent of $\ell$.

\begin{enumerate} \item[\textup{(a)}] The pair $(\bar F,\bar{\mathsf n})=(H^\sharp X',\mathsf N^\sharp X')$ is admitted at $X'$, and $L^{-1}\langle X',\mathsf B'X'\rangle=\|(\bar F,\bar{\mathsf n})\|^2_{\mathcal H}\le\|(F,\mathsf n)\|_{\mathcal H}^2$ holds for every pair $(F,\mathsf n)$ admitted at $X'$. Every such pair satisfies, moreover, $\|(F,\mathsf n)-(\bar F,\bar{\mathsf n})\|^2_{\mathcal H}=\|(F,\mathsf n)\|^2_{\mathcal H}-\|(\bar F,\bar{\mathsf n})\|^2_{\mathcal H}$.

\item[\textup{(b)}] If $\mathsf B_1\le\mathsf B_2$, then $\Phi(\mathsf B_1)\le\Phi(\mathsf B_2)$.

\item[\textup{(c)}] Let $\mathsf B_1,\mathsf B_2\ge1$, put $\mathsf B_s=\mathsf B_1+s(\mathsf B_2-\mathsf B_1)$ for $0\le s\le1$, and let $H_s^\sharp$ be defined with $\mathsf B_s$. Then \[\Phi(\mathsf B_2)-\Phi(\mathsf B_1)=L\int_0^1H_s^{\sharp\mathsf T}(\mathsf B_2-\mathsf B_1)H_s^\sharp\,ds.\] \end{enumerate} \end{lemma}

\begin{proof} \textup{(a)} By \Cref{Y-A-P4}\,(a), $W\bar F+\mathrm{d}'\bar{\mathsf n}=X'$. Since $\Xi_{\bar Z}$ minimises $\langle Y,\mathsf RY\rangle$ over $\ker Q_3$, the first variation at $\Xi_{\bar Z}$ vanishes: $\langle\delta\Xi,\mathsf R\bar Y\rangle=0$ for $\Xi\in\ker Q_3$. As $\mathsf R\bar Y=\bar F$ and $\delta=\mathrm{d}^*$, this is the condition $\langle\Xi,\mathrm{d}\bar F\rangle=0$ for $\Xi\in\ker Q_3$, and so $(\bar F,\bar{\mathsf n})$ is admitted at $X'$. Now let $(F,\mathsf n)$ be any pair admitted at $X'$. Using first $\bar F=\mathsf R\bar Y$ and $a\bar{\mathsf n}=\mathrm{d}'^*\bar Z$, then $\bar Y=W^*\bar Z+\delta\Xi_{\bar Z}$, and finally the two conditions of admission together with $\Xi_{\bar Z}\in\ker Q_3$, we get \[\langle(\bar F,\bar{\mathsf n}),(F,\mathsf n)\rangle_{\mathcal H}=\langle\bar Y,F\rangle+\langle \mathrm{d}'^*\bar Z,\mathsf n\rangle=\langle\bar Z,WF+\mathrm{d}'\mathsf n\rangle+\langle\Xi_{\bar Z},\mathrm{d}F\rangle=\langle\bar Z,X'\rangle.\] This number is the same for all pairs admitted at $X'$. Taking $(F,\mathsf n)=(\bar F,\bar{\mathsf n})$ we see that it equals $\|(\bar F,\bar{\mathsf n})\|^2_{\mathcal H}$, and since $\bar Z=L^{-1}\mathsf B'X'$ it also equals $L^{-1}\langle X',\mathsf B'X'\rangle$. Expanding the square in $\|(F,\mathsf n)-(\bar F,\bar{\mathsf n})\|^2_{\mathcal H}$ yields the last identity of \textup{(a)}; as its left side is $\ge0$, the inequality follows too.

\textup{(b)} Let $(\bar F_2,\bar{\mathsf n}_2)$ be the minimiser for $\mathsf B_2$. By \textup{(a)} it is admitted at $X'$, and admission is independent of $\ell$. Hence \textup{(a)} applied with $\mathsf B_1$, then $\mathsf B_1\le\mathsf B_2$, and \textup{(a)} applied with $\mathsf B_2$ give \[L^{-1}\langle X',\Phi(\mathsf B_1)X'\rangle\le\langle\bar F_2,\mathsf B_1\bar F_2\rangle+a|\bar{\mathsf n}_2|^2\le L^{-1}\langle X',\Phi(\mathsf B_2)X'\rangle\] for every $X'$, which is \textup{(b)}.

\textup{(c)} Put $m(s):=L^{-1}\langle X',\Phi(\mathsf B_s)X'\rangle$, $D=\mathsf B_2-\mathsf B_1$ and $\bar F_s=H^\sharp_sX'$, and let $s,s'\in[0,1]$. The minimiser for $\mathsf B_s$ is admitted at $X'$, and the square of its norm in the space $\mathcal H$ defined with $\mathsf B_{s'}=\mathsf B_s+(s'-s)D$ equals $m(s)+(s'-s)\langle\bar F_s,D\bar F_s\rangle$; so \textup{(a)} applied with $\mathsf B_{s'}$ gives $m(s')\le m(s)+(s'-s)\langle\bar F_s,D\bar F_s\rangle$. Using this as it stands and with $s$ and $s'$ exchanged, we obtain for $s'>s$ \[\langle\bar F_{s'},D\bar F_{s'}\rangle\le\frac{m(s')-m(s)}{s'-s}\le\langle\bar F_s,D\bar F_s\rangle.\] The operator $H_s^\sharp$ depends continuously on $s$: it is a rational function of the entries of $\mathsf B_s$ that has no poles for $0\le s\le1$, because $R_\sharp>0$. Letting $s'$ tend to $s$ from above in this display, and from below in the display with $s$ and $s'$ exchanged, we conclude that $m'(s)=\langle\bar F_s,D\bar F_s\rangle$. Integrating from $0$ to $1$ we obtain the identity of \textup{(c)} between the quadratic forms of its two sides at $X'$, and polarisation gives \textup{(c)}. \end{proof}

\begin{lemma}[The gradient of the minimiser]\label{Y-A-L27} Let $\ell$ be translation covariant, and let $b_{\mathsf B}\le b$ and $\|\mathsf B'\|\le b'$. Then for every $X'\in\Omega'^2$ \[|\nabla H^\sharp X'|\le23\,b\,b'\,L^{-1}|X'|.\] \end{lemma}

The natural size of $|\nabla H^\sharp X'|/|X'|$ is $L^{-3/2}$; the bound above is weaker because $W^*Z$ jumps across the seams of the blocks.

\begin{proof} \emph{Step (i): the gradient of $W^*Z$.} We show that $|\nabla W^*Z|^2\le(4+\frac{48}L)|Z|^2\le20|Z|^2$ for $Z\in\Omega'^2$. Take the component $\mu\nu$ and a direction $\alpha\in\{\mu,\nu\}$. We use the product form of $W(P,p)$ in \Cref{Y-A-Facts}, with the tent function $t$. When $p$ moves by $e_\alpha$, one tent factor of $W(P,p)$ changes by at most $1$ while the other one stays at most $L$, so $|W(P,p+e_\alpha)-W(P,p)|\le L^{-3}\cdot1\cdot L$. For a given $p$ this difference is non-zero for at most $3\cdot2\cdot1=6$ plaquettes $P$, and for a given $P$ for at most $2L\cdot2L\cdot L$ plaquettes $p$. By Schur's test, $|\partial_\alpha W^*Z|^2\le6L^{-2}\cdot4L\,|Z|^2$, and there are two such directions.

In the normal direction $\alpha=\lambda$, the number $W(P,p)$ depends on $z_\lambda$ only through the layer of blocks in which $z_\lambda$ lies. Hence $\partial_\lambda W^*Z$ is zero unless $z_\lambda$ lies at the top of its block, and at such $p$ it equals $\sum_PW(P,p)(Z_{P+Le_\lambda}-Z_P)$. By the Cauchy--Schwarz inequality and $\sum_PW(P,p)=L^{-1}$, its square is at most $L^{-1}\sum_PW(P,p)(Z_{P+Le_\lambda}-Z_P)^2$; moreover $\sum_{p:\,z_\lambda\text{ fixed}}W(P,p)=L$. Therefore $|\partial_\lambda W^*Z|^2\le\sum_P(Z_{P+Le_\lambda}-Z_P)^2\le4|Z|^2$. Adding the two tangential directions and the normal one over all components gives the claim of Step~(i); its last inequality holds as $L\ge3$.

\emph{Step (ii): the bound on $\nabla H^\sharp X'$.} Put $\nu:=H^\sharp X'=\mathsf RY$, where $Y=W^*\bar Z+\delta\Xi$ and $\Xi=\Xi_{\bar Z}$. Since $\bar Z=L^{-1}\mathsf B'X'$, we have $|\bar Z|\le L^{-1}b'|X'|$. As $\mathsf R$ commutes with translations and $\|\mathsf R\|\le1$, we get $|\nabla\nu|\le|\nabla Y|\le|\nabla W^*\bar Z|+|\nabla\delta\Xi|$. Since the difference operators commute, \[|\nabla\delta\Xi|^2=\sum_{\alpha,\beta}|\partial_\alpha\partial_\beta^*\Xi|^2=\langle\Xi,(\sum_\alpha\partial_\alpha^*\partial_\alpha)(\sum_\beta\partial_\beta\partial_\beta^*)\Xi\rangle=|\mathrm{d}\delta\Xi|^2.\]

We now use Fourier series on $\mathbb T$. Let $\hat u(\kappa)\in\mathbb C^3$ be the symbol of $\delta$ on $\Omega^3$, whose entries are $\pm(1-e^{-i\kappa_\beta})$, and let $\hat{\mathsf R}(\kappa)$ be the $3\times3$ symbol of $\mathsf R$. Then $\langle\hat u,\hat{\mathsf R}\hat u\rangle$ is the symbol of the scalar operator $\mathrm{d}\mathsf R\delta$, and $\langle\hat u,\hat{\mathsf R}\hat u\rangle\ge b^{-1}|\hat u|^2$ as $\mathsf R\ge b^{-1}$, where $|\hat u|^2=\hat\Delta$ is the symbol of $\mathrm{d}\delta$. Since $\mathsf R\delta\Xi=\nu-\mathsf RW^*\bar Z$, this gives $|\mathrm{d}\delta\Xi|\le b|\mathrm{d}\mathsf R\delta\Xi|\le b(|\mathrm{d}\nu|+|\mathrm{d}\mathsf RW^*\bar Z|)$. For $F\in\Omega^2$ one has $(\mathrm{d}\mathsf RF)(z)=\sum_\alpha\pm(\mathsf R\partial_\alpha F)_\alpha(z)$, hence $|\mathrm{d}\mathsf RF|^2\le3|\nabla F|^2$, again by $\|\mathsf R\|\le1$.

By \Cref{Y-A-L26}\,(a) the pair $(\nu,\bar{\mathsf n})$ is admitted at $X'$. So $\mathrm{d}\nu$ is orthogonal to $\ker Q_3$, that is $\mathrm{d}\nu\in{\rm ran}\,Q_3^*$, and $Q_3\mathrm{d}\nu=\mathrm{d}'W\nu=\mathrm{d}'(X'-\mathrm{d}'\bar{\mathsf n})=\mathrm{d}'X'$, the first equality by $Q_3\mathrm{d}=\mathrm{d}'W$ from \Cref{Y-A-Facts}. With $Q_3Q_3^*=L^3\hat{\mathsf J}$, $\hat{\mathsf J}\ge\frac1{27}$, again from \Cref{Y-A-Facts}, and with $|\mathrm{d}'X'|^2\le12|X'|^2$, it follows that $|\mathrm{d}\nu|^2=\langle \mathrm{d}'X',(Q_3Q_3^*)^{-1}\mathrm{d}'X'\rangle\le27L^{-3}\cdot12|X'|^2$. We collect the bounds, using Step~(i) for $\bar Z$: \[|\nabla\nu|\le\sqrt{20}(1+\sqrt3b)b'L^{-1}|X'|+18bL^{-3/2}|X'|\le(4.48+7.75+10.4)\,bb'L^{-1}|X'|,\] where the second inequality uses $b,b'\ge1$ and $L\ge3$. As $4.48+7.75+10.4<23$, the lemma follows. \end{proof}

\subsubsection{One step on rows}

\begin{proposition}[One step on a set of rows]\label{Y-A-P28} Let $\ell\in\mathcal R(b,K,\lambda,\tau)$ on $\mathbb T$, let $\gamma_{\rm e}=\gamma_1/2b$, let $\lambda'$ be the number of \Cref{Y-A-L25}\,(c), and put $K_\natural:=3\gamma_{\rm e}^{-1/2}$.

\begin{enumerate} \item[\textup{(a)}] The rows $\ell'$ of \Cref{Y-A-P5} are translation covariant on $\mathbb T'$, and $|\ell'_{y'}(P)|\le K_\natural e^{-\lambda'{\rm dist}(y',P)}$ for all $y',P$. Furthermore, for all cells $i,j$ of $\Omega'^2\oplus\Omega^3$ and every $0<\lambda''\le\lambda'$ we have \[|\mathfrak R^{-1}(i,j)|\le\frac2{\gamma_{\rm e}}(\mathsf D_i\mathsf D_j)^{-1/2}e^{-\lambda''\mathsf d(i,j)},\] and, for $(Z,\Xi)=\mathfrak R^{-1}\iota\delta_P$ and $\nu=H^\sharp\delta_P$, \[\sum_qe^{2\lambda''\mathsf d(q,P)}\Xi_q^2\le4\gamma_{\rm e}^{-2}L,\qquad\sum_pe^{2\lambda''\mathsf d(p,P)}\nu_p^2\le4\gamma_{\rm e}^{-1}L^{-1}.\]

\item[\textup{(b)}] The form $\mathsf B'$ of the rows $\ell'$ satisfies \[1\le\mathsf B'\le\mathsf r_\sharp^{-1}+9522\,\mathsf r_\sharp^{-2}K^2\mathsf g_1(\lambda)^2L^{-1}+2\tau^2\mathsf r_\sharp^{-1}=:b'.\]

\item[\textup{(c)}] We have $\tau(\ell')\le\tau$. If $\tau=0$, then $H^\sharp h'=L^{-2}h$ for every constant form $h'$ on $\mathbb T'$, where $h$ denotes the constant form on $\mathbb T$ with the components of $h'$.

\item[\textup{(d)}] Let $\varrho$ be a cut radius with $\mathsf t:=3K_\natural\mathsf g(\frac{\lambda'}2)e^{-\lambda'\varrho/2}\le\frac13$. Then $\|\ell'^>\|_{\rm r},\varsigma_>\le\mathsf t$, $\delta_{\rm c}\le3\mathsf t$, and $\ell'^{\rm c}\in\mathcal R\bigl((1+3\mathsf t)b',K_\natural,\lambda',\tau+\mathsf t\bigr)$. \end{enumerate} \end{proposition}

\begin{proof} \textup{(a)} We begin with the covariance. Let $\mathsf t'\in L\mathbb Z^3$. The maps $W,Q_3,\mathrm{d},\delta,\mathrm{d}'$ commute with the shift by $\mathsf t'$, and so does $\mathsf R$. Hence the shift maps $\ker Q_3$ onto itself and carries $\Xi_Z$ to the minimiser that belongs to the shifted $Z$, and it commutes with $R_\sharp$, $\mathsf B'$, $H^\sharp$ and $\mathsf N^\sharp$. It follows that the groups (1) to (4) of $\ell'_{y'+\mathsf t'}$, evaluated at the shifted form, consist of the numbers of $\ell'_{y'}$ with shifted indices. For the group (5) one has $\ell_{y+\mathsf t'}H^\sharp(\text{shifted }\zeta)=U_{\mathsf t',y}\ell_yH^\sharp\zeta$. Thus $\ell'$ is translation covariant on $\mathbb T'$.

Next we prove the bound. Fix $P$, put $\Phi=(\mathfrak R^{\rm e})^{-1}(\iota\delta_P,0)$, and note that $\mathsf D^{\rm e}_P=L$. By \Cref{Y-A-L25}\,(a), $\Phi=(Z,\Xi,\ell\nu)$ with $(Z,\Xi)=\mathfrak R^{-1}\iota\delta_P$ and $\nu=H^\sharp\delta_P$, and the column $\ell'_{y'}(P)$ splits into a sum of two vectors. The first of them carries the numbers $L^{1/2}\nu_p$ for the $p$ based in $B(y')$, the numbers $L^{1/2}a^{-1/2}(\mathrm{d}'^*Z)_c=(aL)^{1/2}\mathsf N^\sharp(c,P)$ for the $c$ based at $y'$, the vectors $L^{1/2}(\ell_y\nu)$ for $y\in B(y')$, and zeros in the groups (1) and (2). Apart from these zeros, the entries of the first vector are the entries $(\mathbb A^{\rm e}\Phi)_r$ multiplied by $L^{1/2}$, for the rows $r$ of three kinds: the $p$ based in $B(y')$, the $c$ based at $y'$, and the $(y,\alpha)$ with $y\in B(y')$. For each of these rows, ${\rm dist}(r,P)={\rm dist}(y',P)$; hence $\mathsf d(r,P)>{\rm dist}(y',P)-1$, because $|\mathsf d-{\rm dist}|<1$ by \Cref{Y-A-Distances}. By \Cref{Y-A-L25}\,(b) and (c) and by \Cref{Y-A-L24}\,(b), the norm of the first vector is therefore at most $L^{1/2}\cdot2\gamma_{\rm e}^{-1/2}L^{-1/2}e^{\lambda'}e^{-\lambda'{\rm dist}(y',P)}$.

The second vector is explicit, and it vanishes for ${\rm dist}(y',P)>1$. In the group (1) of \Cref{Y-A-P5}, where $\zeta=\delta_P$ with $P=(y;\mu\nu)$, the four of its six numbers that involve the two other components of $\zeta$ vanish, so that at most the two that involve $\zeta_{\mu\nu}$ differ from $0$. Each of these two is the value of $u(\partial'_\rho\delta_P)$ at $y'$ plus the value of $v(\partial'_\rho\delta_P)$ at the neighbour of $y'$ along $\rho'$. Four sites enter here: $y'$, the neighbours of $y'$ along $\rho$ and along $\rho'$, and the site reached from $y'$ by one step along $\rho$ and one along $\rho'$; they are distinct, even for $N=2$. Hence each of the two numbers has at most one non-zero term and lies in $\{\pm u,\pm v,0\}$. So this group has norm at most $\sqrt2\cdot0.41$, since $0<v\le u\le\sqrt\alpha<\sqrt{c_L}<6^{-1/2}<0.41$ for the number $\alpha$ of \Cref{Y-A-P5}. The group (2) has norm at most $18^{-1/2}\sqrt3$. The norm of the group (3) is $(\sum_pL^{-1}W(P,p)^2)^{1/2}\le(L^{-2}\sum_pW(P,p))^{1/2}=1$, since $W(P,p)\le L^{-1}$ and $\sum_pW(P,p)=L^2$. The group (4) has norm at most $18^{-1/2}\sqrt3$. Altogether the second vector has norm at most $2.4$. Where it does not vanish, $e^{\lambda'}e^{-\lambda'{\rm dist}(y',P)}\ge1$; so the two bounds add up to $|\ell'_{y'}(P)|\le(2\gamma_{\rm e}^{-1/2}+2.4)e^{\lambda'}e^{-\lambda'{\rm dist}(y',P)}$ for all $y',P$, and $(2\gamma_{\rm e}^{-1/2}+2.4)e^{\lambda'}\le3\gamma_{\rm e}^{-1/2}$ because $\lambda'\le\frac1{44}$ and $\gamma_{\rm e}=\gamma_1/2b\le\gamma_1/2\le1.72\cdot10^{-14}\le10^{-13}$. This estimate of $\gamma_{\rm e}$ rests on $b\ge1$ and on the numbers of \Cref{Y-A-Cst}.

The remaining claims of \textup{(a)} follow from \Cref{Y-A-L24}\,(a) and (b), applied with the rate $\lambda''$, because by \Cref{Y-A-L25}\,(a) the block of $(\mathfrak R^{\rm e})^{-1}$ in the variables $(Z,\Xi)$ is $\mathfrak R^{-1}$.

\textup{(b)} The lower bound $\mathsf B'\ge1$ is \Cref{Y-A-P4}\,(b). For the upper bound, let $X'\in\Omega'^2$, and let $(\bar F_0,\bar{\mathsf n}_0)$ denote the minimiser for $\ell=\emptyset$ on the same two lattices. \Cref{Y-A-L26}\,(a) shows that this pair is admitted at $X'$, and admission is independent of $\ell$; so \Cref{Y-A-L26}\,(a) for the rows $\ell$, together with \Cref{Y-A-P6}\,(c), gives \[L^{-1}\langle X',\mathsf B'X'\rangle\le|\bar F_0|^2+a|\bar{\mathsf n}_0|^2+|\ell\bar F_0|^2\le(L\mathsf r_\sharp)^{-1}|X'|^2+|\ell\bar F_0|^2.\]

We claim that every $F\in\Omega^2$ satisfies $|\ell F|\le3K\mathsf g_1(\lambda)|\nabla F|+\tau|F|$. For the proof, let $h^{(y)}$ be the constant form whose components are the numbers $F_\beta(y)$. Splitting $\ell_yF=\ell_y(F-h^{(y)})+\ell_yh^{(y)}$ and using the decay of the entries $\ell_y(p)$, we get $|\ell_yF|\le\sum_{\beta,u}Ke^{-\lambda|u|_\infty}|F_\beta(y+u)-F_\beta(y)|+|\ell_yh^{(y)}|$. We take the norm in $\ell^2(y)$ and use Minkowski's inequality. The norm in $\ell^2(y)$ of the function $y\mapsto F_\beta(y+u)-F_\beta(y)$ is at most $\sum_i|u_i||\partial_iF_\beta|$, and $\sum_{\beta,i}|\partial_iF_\beta|\le3|\nabla F|$; hence the first sum contributes at most $3K\mathsf g_1(\lambda)|\nabla F|$. For the last term, $|\ell_yh^{(y)}|\le\tau|h^{(y)}|_\infty\le\tau(\sum_\beta F_\beta(y)^2)^{1/2}$, and the norm of the right side in $\ell^2(y)$ is $\tau|F|$. This proves the claim.

\Cref{Y-A-L27} applied with $\ell=\emptyset$, where $b=1$ and $b'=\mathsf r_\sharp^{-1}$, gives $|\nabla\bar F_0|\le23\mathsf r_\sharp^{-1}L^{-1}|X'|$, and the display above gives $|\bar F_0|\le(L\mathsf r_\sharp)^{-1/2}|X'|$. By the claim, $|\ell\bar F_0|$ is at most $|X'|$ times $69K\mathsf g_1(\lambda)\mathsf r_\sharp^{-1}L^{-1}+\tau(L\mathsf r_\sharp)^{-1/2}$; since the square of a sum of two terms is at most twice the sum of their squares, \[L(69K\mathsf g_1(\lambda)\mathsf r_\sharp^{-1}L^{-1}+\tau(L\mathsf r_\sharp)^{-1/2})^2\le9522K^2\mathsf g_1(\lambda)^2\mathsf r_\sharp^{-2}L^{-1}+2\tau^2\mathsf r_\sharp^{-1}.\] Multiplying the first display of the proof of \textup{(b)} by $L$ and inserting the bound just found for $L|\ell\bar F_0|^2/|X'|^2$, we obtain \textup{(b)}.

\textup{(c)} Let $h'$ be a constant form on $\mathbb T'$, and let $h$ be the constant form on $\mathbb T$ with the components of $h'$. The pair $(L^{-2}h,0)$ is admitted at $h'$, because $\sum_pW(P,p)=L^2$ and $dh=0$. \Cref{Y-A-L26}\,(a) therefore gives $L^{-1}\langle h',\mathsf B'h'\rangle\le L^{-4}(|h|^2+|\ell h|^2)$, where $|h|^2=L^3|h'|^2$ and $|\ell h|^2\le(NL)^3\tau^2|h|_\infty^2$. Hence \[\sum_{y'}|\ell'_{y'}h'|^2=\langle h',(\mathsf B'-1)h'\rangle\le N^3\tau^2|h|^2_\infty,\] and the $N^3$ terms on the left are equal by covariance; this gives $\tau(\ell')\le\tau$. If $\tau=0$, the bound reads $L^{-1}\langle h',\mathsf B'h'\rangle\le L^{-1}|h'|^2$; as $\mathsf B'\ge1$, the trial pair $(L^{-2}h,0)$ attains the minimum, and the last identity of \Cref{Y-A-L26}\,(a) then shows that it is the minimiser $(H^\sharp h',\mathsf N^\sharp h')$.

\textup{(d)} By \textup{(a)}, $\sum_{{\rm dist}>\varrho}|\ell'_{y'}(P)|\le3K_\natural e^{-\lambda'\varrho/2}\mathsf g(\frac{\lambda'}2)$ for the sum over $P$ at fixed $y'$, and the sum over $y'$ at fixed $P$ is at most a third of this bound. Hence $\|\ell'^>\|_{\rm r}\le\mathsf t$ and $\varsigma_>\le\mathsf t$, so that $\delta_{\rm c}\le3\mathsf t$. The cut is decided by ${\rm dist}(y',P)$, which is invariant under a common translation of $y'$ and $P$, so the cut rows $\ell'^{\rm c}$ inherit the translation covariance of $\ell'$. For the form of the cut rows we apply \Cref{Y-A-P8} with $m\equiv1$. Finally, for a constant form $h'$ we have $|\ell'^{\rm c}_{y'}h'|\le|\ell'_{y'}h'|+\|\ell'^>\|_{\rm r}|h'|_\infty$, which bounds the defect of $\ell'^{\rm c}$ by $\tau+\mathsf t$, by \textup{(c)}. \end{proof}

\subsubsection{The absorbing set of rows}

\begin{definition}[Pure numbers of the absorbing set]\label{Y-A-PureNumbers} A number is called \emph{pure} if it depends neither on $L,a,N,k$, nor on the cuts, nor on $x$. We put \[\FitDisplay{\dimexpr\linewidth-1pt\relax}{b_*:=\frac2{\mathsf r_\sharp},\qquad\gamma_*:=\frac{\gamma_1}{2b_*}=\frac{\mathsf r_\sharp^2}{12296},\qquad\lambda_*:=\frac{\gamma_*^{1/2}}{44},\qquad K_*:=3\gamma_*^{-1/2},\qquad\mathsf t_*(\varrho):=3K_*\mathsf g(\tfrac{\lambda_*}2)e^{-\lambda_*\varrho/2},}\] and \[\FitDisplay{\dimexpr\linewidth-1pt\relax}{L_A:=\max\Bigl\{12K_*\mathsf g(\tfrac{\lambda_*}2)\gamma_*^{-1/2},\ 38088\,K_*^2\mathsf g_1(\lambda_*)^2\mathsf r_\sharp^{-1}\Bigr\}.}\] \end{definition}

\begin{theorem}[One step, and the whole chain]\label{Y-A-Th29} Let $L\ge L_A$ be odd, let $aL\ge3/c_L$, and let $N\ge2$ be even.

\begin{enumerate} \item[\textup{(a)}] \emph{One step.} Let $\ell\in\mathcal R(b_*,K,\lambda_*,\tau)$, where $\tau\le\frac19$ and where $K$ satisfies $12K\mathsf g(\frac{\lambda_*}2)\gamma_*^{-1/2}\le L$ and $38088K^2\mathsf g_1(\lambda_*)^2\mathsf r_\sharp^{-1}\le L$. Then $|\ell'_{y'}(P)|\le K_*e^{-\lambda_*{\rm dist}(y',P)}$, $1\le\mathsf B'\le\frac32\mathsf r_\sharp^{-1}$ and $\tau(\ell')\le\tau$. Moreover, $\ell'^{\rm c}\in\mathcal R(b_*,K_*,\lambda_*,\tau+\mathsf t_*(\varrho))$ for every $\varrho$ with $\mathsf t_*(\varrho)\le\frac19$.

\item[\textup{(b)}] \emph{The chain.} The condition on the radii of the cuts is \begin{equation}\tag{Cut}\sum_{k=1}^K\mathsf t_*(\varrho_k)\le\frac19.\end{equation} Let the radii satisfy \textup{(Cut)}. Then, for every tower, $\ell_k\in\mathcal R(b_*,K_*,\lambda_*,\frac19)$ for $0\le k\le K$, and $(\ell_k)'$ obeys the bounds of \textup{(a)} for $0\le k<K$. For the uncut chain, $\tau(\ell_k)=0$. \end{enumerate} \end{theorem}

\begin{proof} \textup{(a)} We apply \Cref{Y-A-P28} with $b=b_*$ and $\lambda=\lambda_*$, so that $\gamma_{\rm e}=\gamma_*$. As $L\ge3$, and by the first condition on $K$, we get $\lambda'=\min\{\lambda_*,\frac12\lambda_*L,\lambda_*\cdot L\gamma_*^{1/2}/12K\mathsf g(\frac{\lambda_*}2)\}=\lambda_*$; moreover $K_\natural=K_*$. \Cref{Y-A-P28}\,(b) gives $b'\le\mathsf r_\sharp^{-1}(1+\frac14+\frac2{81})\le\frac32\mathsf r_\sharp^{-1}$, where the term $\frac14$ comes from the second condition on $K$ and the term $\frac2{81}$ from $\tau\le\frac19$. With \Cref{Y-A-P28}\,(a) and (c), this proves the bounds on $\ell'$, $\mathsf B'$ and $\tau(\ell')$ in \textup{(a)}. If $\mathsf t_*(\varrho)\le\frac19$, then \Cref{Y-A-P28}\,(d) applies with $\mathsf t=\mathsf t_*(\varrho)$, and $(1+3\mathsf t_*(\varrho))b'\le\frac43\cdot\frac32\mathsf r_\sharp^{-1}=b_*$; this proves the claim on $\ell'^{\rm c}$.

\textup{(b)} We argue by induction on $k$, starting from $\emptyset\in\mathcal R(1,0,\lambda_*,0)$ and carrying the bound $\tau(\ell_k)\le\sum_{j\le k}\mathsf t_*(\varrho_j)$, which is at most $\frac19$ by \textup{(Cut)}. The sets $\mathcal R(\cdot)$ become larger when $b$, $K$ or $\tau$ is raised and $\lambda$ is held fixed, so that $\mathcal R(1,0,\lambda_*,0)\subset\mathcal R(b_*,K_*,\lambda_*,0)$. At every step $K_*$ satisfies the two conditions on $K$ in \textup{(a)}, by the definition of $L_A$. At step $k$ the two tori are $\mathbb T_k$ and $\mathbb T_{k+1}$, and the side $NL^{K-k-1}$ of the coarse one is even. We apply \textup{(a)} to $\ell_k$, with $K_*$ as the bound $K$ of the rows in \textup{(a)}, with $\tau=\tau(\ell_k)$ and with $\varrho=\varrho_{k+1}$; this choice of $\varrho$ is allowed in \textup{(a)}, since \textup{(Cut)} also bounds each single term, $\mathsf t_*(\varrho_{k+1})\le\frac19$. It yields the bounds on $(\ell_k)'$ and shows that $\ell_{k+1}$ lies in $\mathcal R(b_*,K_*,\lambda_*,\tau(\ell_k)+\mathsf t_*(\varrho_{k+1}))$. For the uncut chain, $\varrho_j=\infty$ gives $\mathsf t_*(\varrho_j)=0$, and the carried bound yields $\tau(\ell_k)=0$. \end{proof}

\subsubsection{Kernels on cubes and their decay}

\begin{lemma}[The kernel of $(\mathrm{d}\mathsf R\delta+L^{-2})^{-1}$ on cubes]\label{Y-A-L31} Let $\mathcal N\ge2$ be an integer, and let $\mathsf R$ be a real symmetric operator on $\Omega^2(\mathbb Z^3/\mathcal N\mathbb Z^3)$, commuting with translations, with $b^{-1}\le\mathsf R\le1$ and $|\mathsf R(p,p')|\le Ce^{-\theta|p-p'|}$, where $b\ge1$, $C\ge1$ and $0<\theta\le1$. Put \[\FitDisplay{\dimexpr\linewidth-1pt\relax}{\eta_0:=\min\Bigl\{\tfrac12,\tfrac\theta2,\frac1{24bC\mathsf g_1(\theta/2)}\Bigr\},\quad\omega:=\min\bigl\{(156b)^{-1/2},\tfrac14\eta_0\bigr\},\quad n_0:=18\pi bC\mathsf g_1(\theta),\quad c^{\mathbb Z}:=4b\,(12\eta_0^{-2}+900b^2),}\] \[\FitDisplay{\dimexpr\linewidth-1pt\relax}{\omega_{\mathcal G}:=\tfrac12\omega,\qquad c_{\mathcal G}:=\max\bigl\{c^{\mathbb Z}\mathsf g(\tfrac\omega2),\ n_0^2(1+n_0)^2e^{n_0}\bigr\}.}\] Then, for each real $L$ in the range $1\le L\le\frac12\mathcal N$, the kernel $\mathcal G=(\mathrm{d}\mathsf R\delta+L^{-2})^{-1}$ on $\Omega^3$ obeys, with the periodic distance $r=|x-y|_\infty$ and with $x^+$ a neighbour of $x$, \begin{equation}\tag{G$_{\mathsf R}$}|\mathcal G(x,y)|\le c_{\mathcal G}(1+r)^{-1}e^{-\omega_{\mathcal G}r/L},\qquad|\mathcal G(x^+,y)-\mathcal G(x,y)|\le c_{\mathcal G}(1+r)^{-2}e^{-\omega_{\mathcal G}r/L}.\end{equation} If $\mathcal N\ge n_0$, these bounds hold with the first member of the maximum in place of $c_{\mathcal G}$, and $c^{\mathbb Z}\mathsf g(\frac\omega2)\le2.2\cdot10^{30}\,b^6C^5\theta^{-20}$. \end{lemma}

\begin{proof} We write $n:=\mathcal N$ and $m^2:=L^{-2}$. Because $\mathsf R$ commutes with translations, the entry of $\mathsf R$ that links the plaquettes with normals $e_\alpha,e_\beta$ based at $z,z'$ is a function of $z-z'$ alone, which we denote by $\mathsf R_{\alpha\beta}(z-z')$. The symmetry of $\mathsf R$ gives $\mathsf R_{\alpha\beta}(u)=\mathsf R_{\beta\alpha}(-u)$.

\emph{Step (1): a lift to $\mathbb Z^3$.} For $u\in\mathbb Z^3$ let $[u]$ be its class modulo $n$, and put $\mathsf R^{\mathbb Z}(u):=2^{-\#\{i:|u_i|=n/2\}}\mathsf R([u])$ if $|u|_\infty\le\frac n2$, and $\mathsf R^{\mathbb Z}(u):=0$ otherwise. Then $\sum_{j\in\mathbb Z^3}\mathsf R^{\mathbb Z}(u+nj)=\mathsf R([u])$, $\mathsf R^{\mathbb Z}_{\alpha\beta}(u)=\mathsf R^{\mathbb Z}_{\beta\alpha}(-u)$ and $|\mathsf R^{\mathbb Z}_{\alpha\beta}(u)|\le Ce^{-\theta|u|_\infty}$. We set $\tilde{\mathsf R}(k):=\sum_u\mathsf R^{\mathbb Z}(u)e^{-ik\cdot u}$. This is a $3\times3$ matrix of trigonometric polynomials, Hermitian for real $k$; at the points $k\in\frac{2\pi}n\mathbb Z^3$ it coincides with the symbol of $\mathsf R$, whose spectrum lies in $[b^{-1},1]$.

\emph{Step (2): real $k$.} From this step to the end of Step~(6) we assume $n\ge n_0$; the case $n<n_0$ is treated in Step~(7). For real $k$, let $k_0$ be a point of the grid $\frac{2\pi}n\mathbb Z^3$ with $|k-k_0|_\infty\le\frac\pi n$. The norm of a $3\times3$ matrix is at most $3\max|{\rm entry}|$, and $|e^{-i(k-k_0)u}-1|\le\frac\pi n3|u|_\infty$. Hence $\|\tilde{\mathsf R}(k)-\tilde{\mathsf R}(k_0)\|\le\frac{9\pi C}n\mathsf g_1(\theta)\le\frac1{2b}$, the second bound by $n\ge n_0$. Since the spectrum of $\tilde{\mathsf R}(k_0)$ lies in $[b^{-1},1]$ and $b\ge1$, we conclude that $\frac1{2b}\le\tilde{\mathsf R}(k)\le\frac32$.

\emph{Step (3): the strip.} We have $(\delta\Xi)_\beta(z)=s_\beta(\Xi(z)-\Xi(z-e_\beta))$ with signs $s_\beta=\pm1$, so $\delta$ and $\mathrm{d}$ have the symbols $u_\beta(k)=s_\beta(1-e^{-ik_\beta})$ and $\tilde u_\beta(k)=s_\beta(1-e^{ik_\beta})$. The function $E(k):=\sum_{\alpha\beta}\tilde u_\alpha(k)\tilde{\mathsf R}_{\alpha\beta}(k)u_\beta(k)$ is entire and $2\pi$-periodic. Let $k=\kappa+i\eta e_i$ with $\kappa$ real and $|\eta|\le\eta_0$, and put $v:=u(\kappa)$, so that $|v|^2=\hat\Delta(\kappa)=\sum_\beta2(1-\cos\kappa_\beta)$. We write $u(k)=v+\mathsf a$ and $\tilde u(k)=\bar v+\mathsf a'$; as $|\eta|\le\frac12$, we have $|\mathsf a|,|\mathsf a'|\le e^{|\eta|}-1\le1.3|\eta|$. Further, we write $\tilde{\mathsf R}(k)=\tilde{\mathsf R}(\kappa)+\varrho$; by the definition of $\eta_0$, $\|\varrho\|\le3C\sum_ue^{-\theta|u|_\infty}(e^{|\eta||u|_\infty}-1)\le3C|\eta|\mathsf g_1(\frac\theta2)\le\frac1{8b}$.

We expand $E(k)$ in these terms. The main term, with $\tilde{\mathsf R}=\tilde{\mathsf R}(\kappa)$, satisfies $\bar v^{\mathsf T}\tilde{\mathsf R}v\ge\frac1{2b}|v|^2$ by Step~(2). Next, $|\bar v^{\mathsf T}\varrho v|\le\frac1{8b}|v|^2$. The four terms with one factor $\mathsf a$ or $\mathsf a'$ are together at most $2(\frac32+\frac18)1.3|\eta||v|=4.225|\eta||v|\le\frac1{8b}|v|^2+35.8b\eta^2$, and the two terms with both $\mathsf a'$ and $\mathsf a$ are at most $(\frac32+\frac18)1.69\eta^2\le2.8\eta^2$. As $b\ge1$, it follows that \begin{equation}\label{eq:free-symbol-coercivity}\operatorname{Re}E(\kappa+i\eta e_i)\ \ge\ \tfrac1{4b}\hat\Delta(\kappa)-39b\,\eta^2 .\end{equation}

\emph{Step (4): the heat kernel on $\mathbb Z^3$.} For $t\ge0$ and $x\in\mathbb Z^3$ we define \[P_t(x):=(2\pi)^{-3}\int_{[-\pi,\pi]^3}e^{ik\cdot x-tE(k)}dk.\] Because $E\ge0$ at real $k$, the kernel $\mathcal G^{\mathbb Z}$ is an integral of $P_t$ over $t$: \[\mathcal G^{\mathbb Z}(x):=(2\pi)^{-3}\int\frac{e^{ik\cdot x}}{E(k)+m^2}dk=\int_0^\infty e^{-tm^2}P_t(x)dt.\] Let $r=|x|_\infty=|x_i|$, where $i$ is a direction in which the maximum is attained. As a function of $k_i$ the integrand of $P_t(x)$ is entire and periodic, so the shift $k_i\to k_i+i\eta\operatorname{sgn}x_i$, $0\le\eta\le\eta_0$, of the path of integration leaves $P_t(x)$ unchanged. We write $s:=t/4b$ and $A:=156b^2$.

We use \eqref{eq:free-symbol-coercivity} and the inequality $1-\cos\theta\ge2\theta^2/\pi^2$, valid for $|\theta|\le\pi$. Since $\hat\Delta(\kappa)=\sum_\alpha2(1-\cos\kappa_\alpha)$, the factor $e^{-tE}$ of the shifted integrand has modulus at most $e^{A\eta^2s}\prod_\alpha e^{-4s\kappa_\alpha^2/\pi^2}$, and the factor $e^{ik\cdot x}$ has modulus $e^{-\eta r}$. We put $\mathsf a_s:=\min\{1,0.45s^{-1/2}\}\ge\frac1{2\pi}\int e^{-4s\theta^2/\pi^2}d\theta$ and $\mathsf b_s:=\min\{2,0.4s^{-1}\}\ge\frac1{2\pi}\int|\theta|e^{-4s\theta^2/\pi^2}d\theta$, with integrals over $[-\pi,\pi]$; the bounds $0.45s^{-1/2}$ and $0.4s^{-1}$ follow by integrating
on $\mathbb R$, where the exact coefficients are $\sqrt\pi/4$ and
$\pi/8$. The bounds $1$ and $2$ follow by dropping the exponential
on $[-\pi,\pi]$. In a difference $P_t(x\pm e_j)-P_t(x)$ the integrand has the additional factor $e^{\pm ik_j}-1$, and after the shift $|e^{\pm i(\kappa_j+i\eta_j)}-1|\le|\kappa_j|+1.3\eta$, where $\eta_j$, of modulus at most $\eta$, is the imaginary part of the shifted coordinate $k_j$. Integrating in each coordinate separately, we obtain \[|P_t(x)|\le e^{-\eta r+A\eta^2s}\mathsf a_s^3,\qquad|P_t(x\pm e_j)-P_t(x)|\le e^{-\eta r+A\eta^2s}(\mathsf b_s+1.3\eta\mathsf a_s)\mathsf a_s^2 .\] Here only the bound splits into a product over the three coordinates; $P_t$ itself does not.

We choose $\eta=\min\{\eta_0,r/2As\}$. In region II, that is, for $s\ge r/2A\eta_0$, the exponent $-\eta r+A\eta^2s$ then equals $-r^2/4As$; in region I, the smaller values of $s$, we have $\eta=\eta_0$ and the exponent is $\le-\frac12\eta_0r$.

\emph{Step (5): the integrals over $s$.} The substitution $t=4bs$ turns the formula of Step~(4) into $\mathcal G^{\mathbb Z}=4b\int_0^\infty e^{-\mu s}P_{4bs}ds$, with $\mu=4b/L^2$. We estimate the integral $\int_0^\infty e^{-\mu s}|P_{4bs}(x)|\,ds$, and the same integral with $P_{4bs}(x\pm e_j)-P_{4bs}(x)$ in place of $P_{4bs}(x)$, separately over the two regions; the factor $4b$ is restored at the end. Splitting the integrals at the breakpoints $0.2$ and $0.2025$ gives
$\int_0^\infty\mathsf a_s^3\,\mathrm{d}s=0.6075$ and
$\int_0^\infty\mathsf b_s\mathsf a_s^2\,\mathrm{d}s
=0.8+0.4\log(0.2025/0.2)<0.805$.

Consider first $\mathcal G^{\mathbb Z}$ with $r\ge1$. Region I contributes at most $0.61e^{-\eta_0r/2}$, and since $(1+r)e^{-\eta_0r/4}\le\frac4{\eta_0}e^{-7/8}$, this is at most $\frac{1.02}{\eta_0}(1+r)^{-1}e^{-\eta_0r/4}$. In region II, $\mathsf a_s^3\le0.0912s^{-3/2}$, and we write $e^{-r^2/4As}$ as the product of two factors $e^{-r^2/8As}$. One of them, together with $e^{-\mu s}$, is at most $e^{-r/(L\sqrt{78b})}$, because $\frac{r^2}{8As}+\mu s\ge r(\frac\mu{2A})^{1/2}=\frac r{L\sqrt{78b}}$; the other is integrated by means of $\int_0^\infty s^{-3/2}e^{-r^2/8As}ds=\sqrt{8\pi A}/r\le62.7b/r$. Hence region II contributes at most $5.72\,b\,r^{-1}e^{-r/(L\sqrt{78b})}\le11.5b(1+r)^{-1}e^{-\omega r/L}$.

Consider next the differences, again with $r\ge1$. In region I we have $1.3\eta=1.3\eta_0\le0.65$, so this region contributes at most $(0.805+0.65\cdot0.6075)e^{-\eta_0r/2}\le1.2e^{-\eta_0r/2}$; since $(1+r)^2e^{-\eta_0r/4}\le\frac{64}{\eta_0^2}e^{-15/8}$, this is at most $\frac{11.8}{\eta_0^2}(1+r)^{-2}e^{-\eta_0r/4}$. In region II, \[1.3\eta\mathsf a_se^{-r^2/8As}\le1.3\cdot0.45\cdot\frac1{2\sqrt A}\sup_y(ye^{-y^2/8})s^{-1}\le0.0285s^{-1},\] where $y=r/\sqrt{As}$; the supremum equals $2e^{-1/2}$, and $\sqrt A=\sqrt{156}\,b\ge12.48$, which yields the value $0.02844$. With $\mathsf b_s\le0.4s^{-1}$ and $\mathsf a_s^2\le0.2025s^{-1}$, the integrand is therefore at most $0.4285\cdot0.2025\,s^{-2}e^{-r^2/8As-\mu s}$. Splitting $e^{-r^2/8As}$ once more into two equal factors, we use $\frac{r^2}{16As}+\mu s\ge\frac r{L\sqrt{156b}}$ and $\int_0^\infty s^{-2}e^{-r^2/16As}ds=16A/r^2$; so region II contributes at most $216.6\,b^2r^{-2}e^{-r/(L\sqrt{156b})}\le867b^2(1+r)^{-2}e^{-\omega r/L}$.

For $r=0$ we take $\eta=0$, and the two integrals are at most $0.61$ and $0.81$. As $L\ge1$ and $\omega\le\frac14\eta_0$, we have $e^{-\eta_0r/4}\le e^{-\omega r/L}$. We add the contributions of the two regions and multiply by $4b$. Since $\eta_0\le\frac12$ and $b\ge1$, the constants so obtained are at most $4b(12\eta_0^{-2}+900b^2)=c^{\mathbb Z}$, and for all $x\in\mathbb Z^3$ \[|\mathcal G^{\mathbb Z}(x)|\le c^{\mathbb Z}(1+r)^{-1}e^{-\omega r/L},\qquad|\mathcal G^{\mathbb Z}(x\pm e_j)-\mathcal G^{\mathbb Z}(x)|\le c^{\mathbb Z}(1+r)^{-2}e^{-\omega r/L}.\]

\emph{Step (6): the torus.} The function $f(k)=(E(k)+m^2)^{-1}$ has the Fourier coefficients $\mathcal G^{\mathbb Z}$, which by Step~(5) are absolutely summable. Consequently, for $x\in\mathbb Z^3$, \[n^{-3}\sum_{k\in\frac{2\pi}n\mathbb Z_n^3}f(k)e^{ik\cdot x}=\sum_y\mathcal G^{\mathbb Z}(y)1\{y\equiv x\},\] where $y\equiv x$ is congruence modulo $n$. On the grid $\frac{2\pi}n\mathbb Z^3$ the function $E$ is the symbol of $\mathrm{d}\mathsf R\delta$, so the left side is $\mathcal G(x,0)$, and $\mathcal G(x,0)=\sum_j\mathcal G^{\mathbb Z}(x+nj)$, where $j$ runs over $\mathbb Z^3$. Let $|x|_\infty=r\le\frac n2$. For $j\ne0$ we have $|x+nj|_\infty\ge\max\{r,\frac n2|j|_\infty\}\ge\frac r2+\frac n4|j|_\infty$, and since $n\ge2L$, the weight satisfies $e^{-\omega|\cdot|_\infty/L}\le e^{-\omega r/2L}e^{-\omega|j|_\infty/2}$ at the point $x+nj$. As moreover $(1+|x+nj|_\infty)^{-1}\le(1+r)^{-1}$, the bounds of Step~(5), summed over $j$, give the bounds of \Cref{Y-A-L31} for $n\ge n_0$, with the rate $\omega_{\mathcal G}=\frac12\omega$ and the constant $c^{\mathbb Z}\mathsf g(\frac\omega2)\le c_{\mathcal G}$.

\emph{Step (7): the case $n<n_0$.} Since $L\le\frac12n$, we have $0<\mathcal G\le L^2\le\frac14n_0^2$ as an operator; so the entries of $\mathcal G$ and the differences of its entries are bounded by $\frac12n_0^2$ in absolute value. As $r\le\frac12n_0$, $L\ge1$ and $\omega_{\mathcal G}\le1$, we have $(1+r)^{-2}e^{-\omega_{\mathcal G}r/L}\ge(1+n_0)^{-2}e^{-n_0}$, and since $c_{\mathcal G}\ge n_0^2(1+n_0)^2e^{n_0}$, the bounds of \Cref{Y-A-L31} hold in this case too.

\emph{Step (8): bounds by powers of $b$, $C$ and $\theta^{-1}$.} By the bounds $\mathsf g_1(\lambda)\le182\lambda^{-4}$ and $\mathsf g(\lambda)\le51\lambda^{-3}$ of \Cref{Y-A-Sums}, valid for $0<\lambda\le1$, the first applied to $\mathsf g_1(\frac\theta2)$, we have $\eta_0^{-1}\le24\cdot182\cdot16\,bC\theta^{-4}\le7\cdot10^4bC\theta^{-4}$. It follows that $c^{\mathbb Z}\le2.4\cdot10^{11}b^3C^2\theta^{-8}$ and $\omega^{-1}\le2.8\cdot10^5bC\theta^{-4}$. The second of the two bounds of \Cref{Y-A-Sums} used above, $\mathsf g(\lambda)\le51\lambda^{-3}$, applied with $\lambda=\frac\omega2$, gives $\mathsf g(\frac\omega2)\le408\omega^{-3}$. We insert the bound on $\omega^{-1}$ into $\mathsf g(\frac\omega2)\le408\omega^{-3}$ and multiply the result by the bound on $c^{\mathbb Z}$. In the product, the powers are $b^{3+3}$, $C^{2+3}$ and $\theta^{-8-12}$, and the numerical factor is $2.4\cdot10^{11}\cdot408\cdot(2.8\cdot10^5)^3\le2.2\cdot10^{30}$; hence $c^{\mathbb Z}\mathsf g(\frac\omega2)\le2.2\cdot10^{30}\,b^6C^5\theta^{-20}$, which is the last bound of \Cref{Y-A-L31}. \end{proof}

\begin{lemma}[Sums over all cubes and over seam cubes]\label{Y-A-L32} Let $0<\varpi\le1$ and $L\ge3$, let $K_i(r):=(1+r)^{-i}e^{-\varpi r/L}$ for $i=1,2$ and $r\ge0$, and let $\mathsf l_L:=1+\log L$. A seam cube of direction $\beta$ is a cube $z$ whose coordinate $z_\beta$ lies at the top of its block. Then, for every cube $x$ of $\mathbb T$, with $r=|x-y|_\infty$ and with sums over the cubes $y$ of $\mathbb T$, \[\FitDisplay{\dimexpr\linewidth-1pt\relax}{\sum_yK_1\le\frac{25L^2}{\varpi^2},\quad\sum_yK_2\le\frac{25L}\varpi,\quad\sum_\beta\sum_{y\ \rm seam\ of\ \beta}K_1\le\frac{216L}{\varpi^2},\quad\sum_\beta\sum_{y\ \rm seam\ of\ \beta}K_2\le\frac{144}\varpi\bigl(1+\log\tfrac3\varpi\bigr)\mathsf l_L .}\] \end{lemma}

\begin{proof} Let $\alpha:=\varpi/L$; then $\alpha\le\frac13$. For $r\ge1$ we have $\frac{24r^2+2}{1+r}\le24r$ and $\frac{24r^2+2}{(1+r)^2}\le24$, and also $\sum_{r\ge1}re^{-\alpha r}\le\alpha^{-2}$ and $\sum_{r\ge1}e^{-\alpha r}\le\alpha^{-1}$. At most $24r^2+2$ cubes $y$ lie at distance $r\ge1$ from $x$, and the cube $y=x$ contributes $1$; as $1\le\alpha^{-1}\le\alpha^{-2}$, this gives $\sum_yK_1\le1+24\alpha^{-2}\le25\alpha^{-2}$ and $\sum_yK_2\le1+24\alpha^{-1}\le25\alpha^{-1}$, which are the first two bounds.

We estimate the two sums over seam cubes one plane at a time. In one plane $y_\beta=s$ let $\mathsf a=|s-x_\beta|$, and let $\tau$ be the transverse sup distance from $y$ to $x$; the plane contains at most $8\tau$ cubes with transverse distance $\tau\ge1$ and one with $\tau=0$, and $r=\max\{\mathsf a,\tau\}$. For $K_2$, the cubes with $\tau\le\mathsf a$, of which there are at most $(2\mathsf a+1)^2\le4(1+\mathsf a)^2$, contribute at most $4e^{-\alpha\mathsf a}$; the cubes with $\tau>\mathsf a$ contribute at most $8e^{-\alpha\mathsf a/2}\sum_{\tau\ge1}\frac{e^{-\alpha\tau/2}}{1+\tau}$. Put $\beta_0=\frac\alpha2\le\frac16$. In the last sum, the terms with $\tau\le1/\beta_0$ add up to at most $\sum_{\tau\le1/\beta_0}\frac1{1+\tau}\le\log(1+\frac{2L}\varpi)\le\log L+\log\frac3\varpi$, and the terms with $\tau>1/\beta_0$ to at most $\beta_0\frac{e^{-1}}{1-e^{-\beta_0}}\le0.4$. So for $K_2$ the plane contributes at most $e^{-\alpha\mathsf a/2}[4+8(\mathsf l_L+\log\frac3\varpi-0.6)]\le8e^{-\alpha\mathsf a/2}(\mathsf l_L+\log\frac3\varpi)$.

For $K_1$ the plane contributes at most $4(1+\mathsf a)e^{-\alpha\mathsf a}+\frac8\alpha e^{-\alpha\mathsf a}$; since $(1+\mathsf a)e^{-\alpha\mathsf a/2}\le\frac2\alpha e^{-5/6}$, this is at most $\frac{12}\alpha e^{-\alpha\mathsf a/2}$.

The seam cubes of direction $\beta$ fill the planes $y_\beta=s$ with $s$ at the top of a block. For each integer $j\ge0$ at most two of these planes have $\mathsf a\in[jL,(j+1)L)$, so that over the seam planes of one direction $\sum e^{-\alpha\mathsf a/2}\le\frac2{1-e^{-\varpi/2}}\le\frac6\varpi$. There are three directions $\beta$; hence the sum over seam cubes is at most $3\cdot\frac{12}\alpha\cdot\frac6\varpi=\frac{216L}{\varpi^2}$ for $K_1$, and at most $3\cdot8\cdot\frac6\varpi(\mathsf l_L+\log\frac3\varpi)\le\frac{144}\varpi(1+\log\frac3\varpi)\mathsf l_L$ for $K_2$, the last step because $\mathsf l_L\ge1$. \end{proof}

\begin{theorem}[Decay of the kernels of one step]\label{Y-A-Th33} Let $\ell\in\mathcal R(b,K,\lambda,\tau)$ be rows on $\mathbb T$, and let $L$, $a$ and $N$ satisfy the standing conditions of \Cref{Y-A-Standing}; none of the constants and bounds below depends on $\tau$. Let $\gamma_{\rm e}=\gamma_1/2b$, let $\lambda'$ be the rate of \Cref{Y-A-L25}\,(c), let $K_\natural=3\gamma_{\rm e}^{-1/2}$, and let $0<\lambda^\circ\le\lambda'$. We define the following constants, in the order in which they are listed. First, $\theta_R:=\min\{\frac\lambda4,[6K^2\mathsf g(\frac\lambda2)\mathsf g_1(\frac\lambda4)]^{-1}\}$, read as $\frac\lambda4$ if $K=0$. Next, $(c_{\mathcal G},\omega_{\mathcal G})$ are the constants of \Cref{Y-A-L31} at $(b,C,\theta)=(b,2,\theta_R)$; then $\varpi:=\frac12\omega_{\mathcal G}$ and $\lambda_B:=\min\{\lambda^\circ,\varpi\}$. Further, $\theta_J:=\frac12\log(1+\frac{w_*}{50})$, a pure number; it is the rate in the bound $|\hat{\mathsf J}^{-2}(i,j)|\le\frac2{w_*}e^{-\theta_J{\rm dist}}$ of \Cref{Y-A-Jinv}. Finally, \[c_\Xi^\natural:=\frac{6000\,c_{\mathcal G}\,\mathsf g_2(\theta_R/4)}{\gamma_{\rm e}\varpi^2},\qquad c_H^\natural:=\frac{180\,\mathsf g(\theta_R/2)\,c_{\mathcal G}\,c_\Xi^\natural}\varpi,\qquad\lambda_2^\natural:=\min\Bigl\{\frac{\lambda^\circ}4,\frac1{6K_\natural^2\mathsf g(\lambda^\circ/2)\mathsf g_1(\lambda^\circ/4)}\Bigr\},\] and \[c_1^\natural:=19\,c_H^\natural\,\mathsf g(\lambda_B)\mathsf g(\lambda_2^\natural),\qquad c_2^\natural:=19\,c_\Xi^\natural\,\mathsf g(\lambda_B)\mathsf g(\lambda_2^\natural)+72\,w_*^{-1}\mathsf g(\theta_J)\mathsf g(\lambda_2^\natural).\] We denote by ${\rm pr}_\Xi$ the projection of $\Omega'^2\oplus\Omega^3$ onto $\Omega^3$, that is, ${\rm pr}_\Xi(Z,\Xi)=\Xi$. Then the following bounds hold for all cells.

\begin{enumerate} \item[\textup{(a)}] $|\mathsf R(p,p')|\le2e^{-\theta_R|p-p'|}$.

\item[\textup{(b)}] $|({\rm pr}_\Xi\mathfrak R^{-1}\iota)(q,P)|\le c_\Xi^\natural L^{-1}e^{-\lambda_B\mathsf d(q,P)}$, $\quad|H^\sharp(p,P)|\le c_H^\natural\,\mathsf l_L\,L^{-2}e^{-\lambda_B\mathsf d(p,P)}$.

\item[\textup{(c)}] $|(\mathsf B'-1)(P,P')|\le K_\natural^2\mathsf g(\frac{\lambda^\circ}2)e^{-\frac12\lambda^\circ{\rm dist}(P,P')}$, $\quad|\hat R(P,P')|\le2e^{-\lambda_2^\natural{\rm dist}(P,P')}$.

\item[\textup{(d)}] $L^2\|\mathsf T\|_\infty\le c_1^\natural\mathsf l_L$, $\quad L\|\mathsf X\|_\infty\le c_2^\natural$, $\quad L^2\sup_{|v|_\infty\le1}|\ell\mathsf Tv|_{2,\infty}\le3K\mathsf g(\lambda)c_1^\natural\mathsf l_L$. \end{enumerate}

Each of these constants is a function of $(b,K,\lambda,\lambda^\circ)$ alone; none of them depends on $L$, $a$ or $N$. \end{theorem}

\begin{proof} \emph{Step (a): the kernel $\mathsf R$.} We have $\mathsf B\ge1$. For $\theta\le\frac\lambda4$, the third bound in the second paragraph of \Cref{Y-A-SetsOfRows}, $|(\mathsf B-1)(p,p')|\le K^2\mathsf g(\frac\lambda2)e^{-\frac\lambda2|p-p'|}$, and the inequality $e^t-1\le te^t$ give \[\varpi_{\mathsf B}(\theta)\le3K^2\mathsf g(\frac\lambda2)\sum_ue^{-\frac\lambda2|u|_\infty}(e^{\theta|u|_\infty}-1)\le3K^2\mathsf g(\frac\lambda2)\theta\mathsf g_1(\frac\lambda4).\] At $\theta=\theta_R$ the right side is at most $\frac12$, and \Cref{Y-A-L3}\,(a), with the distance $|p-p'|$, gives the bound \textup{(a)}.

\emph{Step (b): the rates, and the equation for $\Xi$.} By the definitions of $\omega$ and $\eta_0$ in \Cref{Y-A-L31}, at $\theta=\theta_R$, we have $\omega_{\mathcal G}=\frac12\omega\le\frac18\eta_0\le\frac1{16}\theta_R$. Hence $\lambda_B\le\varpi\le\frac1{32}\theta_R$; moreover $\varpi\le\frac1{128}$, because $\theta_R\le\frac\lambda4\le\frac14$, and $\lambda_B+\varpi\le\omega_{\mathcal G}$.

We fix a coarse plaquette $P$ and put $(Z,\Xi)=\mathfrak R^{-1}\iota\delta_P$. \Cref{Y-A-P28}\,(a), applied with $\lambda''=\lambda_B$, gives the two bounds \textup{(e1)} $|Z_{P'}|\le\frac2{\gamma_{\rm e}}L^{-1}e^{-\lambda_B\mathsf d(P',P)}$ for all $P'$, and \textup{(e2)} $\sum_qe^{2\lambda_B\mathsf d(q,P)}\Xi_q^2\le4\gamma_{\rm e}^{-2}L$. The $\Xi$-component of the equation $\mathfrak R(Z,\Xi)=\iota\delta_P$ reads $\mathrm{d}\mathsf R(W^*Z+\delta\Xi)+\mathsf S_0Q_3^*Q_3\Xi=0$; in other words, \[(\mathrm{d}\mathsf R\delta+L^{-2})\Xi=s_1+s_2+s_3,\qquad s_1=-\mathrm{d}\mathsf RW^*Z,\quad s_2=-\mathsf S_0Q_3^*Q_3\Xi,\quad s_3=L^{-2}\Xi ,\] and $\mathcal G=(\mathrm{d}\mathsf R\delta+L^{-2})^{-1}$ is the kernel of \Cref{Y-A-L31} with $\mathcal N=NL\ge2L$. We put $\mathsf z:=\frac{2.1}{\gamma_{\rm e}}L^{-1}$; the numeral $2.1$ rests on $2e^{3\lambda_B}\le2e^{3/128}\le2.0475\le2.1$.

\emph{Step (b): the source $s_1$.} Let $Y=W^*Z$. Then $(\mathrm{d}\mathsf RY)(z)=\sum_\alpha\pm(\mathsf R\partial_\alpha Y)_\alpha(z)$. By Step~(i) of the proof of \Cref{Y-A-L27}, for $\alpha\ne\beta$ we have $|(\partial_\alpha Y)_\beta(z')|\le6L^{-2}\max|Z_{P'}|$. For $\alpha=\beta$ the value $(\partial_\beta Y)_\beta(z')$ is zero unless $z'$, as a base site, is a seam cube of $\beta$, and there it is at most $2L^{-1}\max|Z_{P'}|$. In both bounds the maximum is taken over the $P'$ with $\mathsf d(P',z')\le\frac52$, and by \textup{(e1)} it is therefore at most $\mathsf ze^{-\lambda_B\mathsf d(z',P)}$. Using \textup{(a)} and $e^{\lambda_B\mathsf d(z,P)}\le e^{\lambda_B\mathsf d(z',P)}e^{\frac12\theta_R|z-z'|}$, we arrive at \[e^{\lambda_B\mathsf d(z,P)}|s_1(z)|\le\mathsf z\Bigl[72\,\mathsf g(\tfrac{\theta_R}2)L^{-2}+4L^{-1}\mathcal S(z)\Bigr],\qquad\mathcal S(z):=\sum_\beta\sum_{z'\ \rm seam\ of\ \beta}e^{-\frac12\theta_R|z-z'|}\] in which $72=6\cdot2\cdot6$ comes from the six pairs $\alpha\ne\beta$ and $4=2\cdot2$ from the three pairs $\alpha=\beta$.

\emph{Step (b): the source $s_2$.} By the Cauchy--Schwarz inequality and \textup{(e2)}, \[|(Q_3\Xi)_{q'}|\le(\sum\psi_{q'}^2)^{1/2}(\sum_{\psi_{q'}\ne0}\Xi^2)^{1/2}\le L^{3/2}\frac2{\gamma_{\rm e}}L^{1/2}e^{-\lambda_B(\mathsf d(q',P)-\frac32)}.\] Since $\sum_{q'}\psi_{q'}=1$, it follows that $e^{\lambda_B\mathsf d(q,P)}|s_2(q)|\le\mathsf zL^{-2}$.

\emph{Step (b): the bound on $\Xi$.} We have $e^{\lambda_B\mathsf d(x,P)}\le e^{\lambda_B\mathsf d(y,P)}e^{\lambda_Br/L}$, with $r=|x-y|_\infty$, and $\omega_{\mathcal G}-\lambda_B\ge\varpi$; so the kernels that take the weighted sources to the weighted field are $c_{\mathcal G}K_i$, with the functions $K_i$ of \Cref{Y-A-L32} at the present $\varpi$. For the seam term of $s_1$ we use a smearing, which compares the kernel $K_i$ at the cube $z$ with $K_i$ at the seam cube $z'$: $K_i(|x-z|)\le K_i(|x-z'|)(1+|z-z'|)^ie^{\varpi|z-z'|/L}$, and therefore $\sum_zK_i(|x-z|)e^{-\frac12\theta_R|z-z'|}\le\mathsf g_2(\frac{\theta_R}4)K_i(|x-z'|)$. For $s_3$ we apply the Cauchy--Schwarz inequality together with \textup{(e2)}, and $\sum_yK_1(r)^2\le\sum_yK_2\le25L/\varpi$. With \Cref{Y-A-L32} this gives \[e^{\lambda_B\mathsf d(x,P)}|\Xi_x|\le c_{\mathcal G}\mathsf z\bigl[72\mathsf g\cdot\frac{25}{\varpi^2}+4\mathsf g_2\cdot\frac{216}{\varpi^2}+\frac{25}{\varpi^2}\bigr]+L^{-2}c_{\mathcal G}(\frac{25L}\varpi)^{1/2}\frac2{\gamma_{\rm e}}L^{1/2}\] \[\le\frac{c_{\mathcal G}}{\gamma_{\rm e}\varpi^2}L^{-1}\bigl[2.1(1800+864+25)+10\bigr]\mathsf g_2(\tfrac{\theta_R}4)\le c_\Xi^\natural L^{-1},\] where $\mathsf g$ and $\mathsf g_2$ stand for $\mathsf g(\frac{\theta_R}2)$ and $\mathsf g_2(\frac{\theta_R}4)$; in the second step we apply $\mathsf g(\frac{\theta_R}2)\le\mathsf g_2(\frac{\theta_R}4)$, $1\le\mathsf g_2$ and $\varpi^{-1/2}\le\varpi^{-2}$. Since $\Xi={\rm pr}_\Xi\mathfrak R^{-1}\iota\delta_P$, this is the first bound in \textup{(b)}.

\emph{Step (b): the bounds on $\delta\Xi$ and on $H^\sharp$.} Let the plaquette $p$ lie between the cubes $x$ and $x^+$. We bound $(\delta\Xi)_p$ through the same three sources, now with the differences of $\mathcal G$: that is, with $c_{\mathcal G}K_2$ in place of $c_{\mathcal G}K_1$ and with the second and fourth sums of \Cref{Y-A-L32} in place of the first and third; and $s_3$ is now bounded by the bound on $\Xi$ just proved, $|s_3|\le L^{-3}c_\Xi^\natural e^{-\lambda_B\mathsf d}$. We have $(\delta\Xi)_p=\pm(\Xi_{x^+}-\Xi_x)=\pm\sum_y[\mathcal G(x^+,y)-\mathcal G(x,y)](s_1+s_2+s_3)(y)$. The base sites of $p$ and of $x$ are at most one fine step apart, so that $\mathsf d(p,P)\le\mathsf d(x,P)+\frac1L$; as $\lambda_B\le\frac1{128}$ and $L\ge3$, we have $e^{\lambda_B/L}\le e^{1/384}\le1.1$, which is the factor $1.1$ of the following bound. With $\mathsf g$ and $\mathsf g_2$ abbreviated as in the bound on $\Xi$, we obtain \[e^{\lambda_B\mathsf d(p,P)}|(\delta\Xi)_p|\le1.1c_{\mathcal G}\bigl\{\mathsf z\bigl[72\mathsf gL^{-2}\frac{25L}\varpi+4L^{-1}\mathsf g_2\frac{144}\varpi(1+\log\frac3\varpi)\mathsf l_L+L^{-2}\frac{25L}\varpi\bigr]+L^{-3}c_\Xi^\natural\frac{25L}\varpi\bigr\}.\] By $\log\frac3\varpi\le\frac3\varpi$ and $\varpi\le\frac12$, \[2.1[1800+576(1+\frac3\varpi)+25]\le\frac{6150}\varpi\le\frac{6000}{\varpi^2}.\] We use $c_{\mathcal G}\ge1$, which follows from $c_{\mathcal G}\ge c^{\mathbb Z}\ge4(48+900)$, as well as $\mathsf l_L\ge1$ and $\mathsf g(\frac{\theta_R}2)\le\mathsf g_2(\frac{\theta_R}4)$. So the expression in braces in the bound on $e^{\lambda_B\mathsf d(p,P)}|(\delta\Xi)_p|$ is at most $L^{-2}\mathsf l_L\frac{26}\varpi c_\Xi^\natural$, and \[|(\delta\Xi)_p|\le\frac{29}\varpi c_{\mathcal G}c_\Xi^\natural\mathsf l_LL^{-2}e^{-\lambda_B\mathsf d(p,P)}.\]

Moreover $|(W^*Z)_p|\le L^{-1}\max|Z|\le\mathsf zL^{-1}e^{-\lambda_B\mathsf d(p,P)}$ and $\frac{2.1}{\gamma_{\rm e}}\le c_\Xi^\natural$, so that $|(W^*Z+\delta\Xi)_p|\le\frac{30}\varpi c_{\mathcal G}c_\Xi^\natural\mathsf l_LL^{-2}e^{-\lambda_B\mathsf d(p,P)}$. By \Cref{Y-A-L5p}, $H^\sharp\delta_P=\mathsf R(W^*Z+\delta\Xi)$. We apply \textup{(a)} with $e^{\lambda_B\mathsf d(p,P)}\le e^{\lambda_B\mathsf d(p',P)}e^{\frac12\theta_R|p-p'|}$; as $\sum_{p'}2e^{-\theta_R|p-p'|}e^{\frac12\theta_R|p-p'|}\le6\mathsf g(\frac{\theta_R}2)$ and $6\cdot30=180$, we obtain the bound on $H^\sharp$ in \textup{(b)}.

\emph{Step (c): the kernels $\mathsf B'$ and $\hat R$.} We have $\mathsf B'-1=\ell'^{\mathsf T}\ell'$. \Cref{Y-A-P28}\,(a), whose bound holds with the rate $\lambda^\circ$ in place of $\lambda'$ because $\lambda^\circ\le\lambda'$, bounds the rows $\ell'$ on $\mathbb T'$ with the constant $K_\natural$ at the rate $\lambda^\circ$; the third bound in the second paragraph of \Cref{Y-A-SetsOfRows}, applied to the rows $\ell'$ on $\mathbb T'$ with $K_\natural$, $\lambda^\circ$ and ${\rm dist}(P,P')$ in place of $K$, $\lambda$ and $|p-p'|$, then gives the first bound in \textup{(c)}. For the second bound we carry out Step~(a) on $\mathbb T'$, with $\mathsf B'$, $K_\natural$, $\lambda^\circ$, $\lambda_2^\natural$, ${\rm dist}(P,P')$ and $\hat R$ in place of $\mathsf B$, $K$, $\lambda$, $\theta_R$, $|p-p'|$ and $\mathsf R$: we have $\mathsf B'\ge1$, by the identity just stated; for $\theta\le\frac{\lambda^\circ}4$, the same bound of \Cref{Y-A-SetsOfRows} for the rows $\ell'$ and $e^t-1\le te^t$ give $\varpi_{\mathsf B'}(\theta)\le3K_\natural^2\mathsf g(\frac{\lambda^\circ}2)\theta\mathsf g_1(\frac{\lambda^\circ}4)$, which is at most $\frac12$ at $\theta=\lambda_2^\natural$; and \Cref{Y-A-L3}\,(a), with the distance ${\rm dist}(P,P')$, gives the bound on $\hat R$.

\emph{Step (d): the norms of $\mathsf T$, $\mathsf X$ and $\ell\mathsf T$.} We have $\sum_Pe^{-\lambda_B\mathsf d(p,P)}\le3e^{\lambda_B}\mathsf g(\lambda_B)\le3.1\mathsf g(\lambda_B)$, and by \textup{(c)}, $\|\hat R\|_\infty\le6\mathsf g(\lambda_2^\natural)$. Since $\mathsf T=H^\sharp\hat R$, the bound on $H^\sharp$ in \textup{(b)}, together with $3.1\cdot6\le19$, gives the first bound in \textup{(d)}. By \Cref{Y-A-L5p}, applied with $X'=\hat Rv$, where $\bar Z=L^{-1}v$, \[\mathsf X={\rm pr}_\Xi\mathfrak R^{-1}\iota\hat R+\mathsf S_0^{-1}Q_3^*(Q_3Q_3^*)^{-2}\mathrm{d}'\hat R={\rm pr}_\Xi\mathfrak R^{-1}\iota\hat R+L^{-1}Q_3^*\hat{\mathsf J}^{-2}\mathrm{d}'\hat R.\] Here $\|Q_3^*\|_\infty=1$, $\|\hat{\mathsf J}^{-2}\|_\infty\le\frac2{w_*}\mathsf g(\theta_J)$ and $\|\mathrm{d}'\|_\infty=6$; with the first bound in \textup{(b)} for the first term, this gives the second bound in \textup{(d)}. Finally $|\ell_y\mathsf Tv|\le\|\ell\|_{\rm r}|\mathsf Tv|_\infty$, and $\|\ell\|_{\rm r}\le3K\mathsf g(\lambda)$ by the first bound in the second paragraph of \Cref{Y-A-SetsOfRows}; with the first bound in \textup{(d)} this gives the third. \end{proof}

\begin{definition}[Constants at the pure numbers]\label{Y-A-Stars} We denote by $\theta_R^*,c_{\mathcal G}^*,\varpi^*,\lambda_B^*,c_\Xi^*,c_H^*,\lambda_2^*,c_1^*,c_2^*$ the constants of \Cref{Y-A-Th33} at $(b,K,\lambda,\lambda^\circ)=(b_*,K_*,\lambda_*,\lambda_*)$. They are pure numbers. \end{definition}

At these values $K_\natural=K_*$, since then $\gamma_{\rm e}=\gamma_*$; hence the formula for $\theta_R$ and the formula for $\lambda_2^\natural$ in \Cref{Y-A-Th33} coincide, and $\lambda_2^*=\theta_R^*$.

\subsubsection{(Dec) on the absorbing set}

\begin{proposition}[The condition \textup{(Dec)} on the absorbing set]\label{Y-A-P35} Let $L\ge L_A$ be odd, $aL\ge3/c_L$, $N\ge2$ even, and $\ell\in\mathcal R(b_*,K_*,\lambda_*,\frac19)$. With $\gamma_q=(73L^4)^{-1}$, let \[\FitDisplay{\dimexpr\linewidth-1pt\relax}{A_q:=2e^2(16+a)+48e^2K_*^2\mathsf g(\tfrac{\lambda_*}2)\mathsf g_2(\tfrac{\lambda_*}4),\quad\theta_q:=\min\Bigl\{\frac{\lambda_*}4,\frac{\gamma_q}{2A_q}\Bigr\},\quad\theta_D:=\tfrac12\min\{\lambda_2^*,\theta_q,\theta_J\},}\] \[\FitDisplay{\dimexpr\linewidth-1pt\relax}{\mathsf C_0:=4L^3\mathsf g(\theta_D)\frac{2.1L^2}{\gamma_*},\quad\mathsf C_R:=6e\,\mathsf g(\tfrac12\theta_R^*),\quad\mathsf C_\ell:=3eK_*\mathsf g(\tfrac{\lambda_*}2),\quad C_D:=\max\bigl\{(4L^3\mathsf C_0)^4\mathsf C_R\mathsf C_\ell,\ (4L^3)^2(1+\mathsf C_\ell^2),\ A_q\bigr\}.}\] Then \textup{(Dec)}$(C_D,\theta_D)$ holds for the pair $(\ell,\ell')$. The numbers $C_D$ and $\theta_D$ depend on $L$, $a$ and pure numbers alone, and on none of $\ell$, $N$, the cuts, $x$ and $k$. Here $\theta_D$ is defined with $\theta_q$; in the number $\theta_D$ of \Cref{Y-A-Fr} the number $\theta_{\mathbb C}$ takes the place of $\theta_q$, as in Steps~H13 and H16\,(ii) of its proof. \end{proposition}

\begin{proof} \emph{The form $q$.} For cells of the fine level, ${\rm dist}<|\cdot|/L+1$. We have $q=\Pi_{\mathfrak L}(\mathrm{d}^*\mathrm{d}+aQ_1^*Q_1+\mathrm{d}^*(\mathsf B-1)\mathrm{d})\Pi_{\mathfrak L}$, where, throughout this section, $\Pi_{\mathfrak L}$ denotes the orthogonal projection of $\Omega^1$ onto $\mathfrak L$, with $\mathfrak L$ and $\mathcal T$ as in \Cref{Y-A-Lifts}: for $\varphi\in\Omega^1$, $(\Pi_{\mathfrak L}\varphi)_b=\varphi_b$ if $b\notin\mathcal T$ and $(\Pi_{\mathfrak L}\varphi)_b=0$ if $b\in\mathcal T$. The first two terms of $q$ have coarse range at most $2$ and absolute row sums at most $16+a$, for the following reasons. First, $(\mathrm{d}^*\mathrm{d})(b,b')=\sum_p\mathrm{d}(p,b)\mathrm{d}(p,b')$, where a bond lies in four plaquettes $p$ and a plaquette has four bonds $b'$, which gives $16$; and any two bonds of one plaquette are based within one fine step, at coarse distance at most $1$. Second, $(Q_1^*Q_1)(b,b')=\sum_cQ_1(c,b)Q_1(c,b')$ with $Q_1(c,b)=L^{-3}\#\{x\in B(c_-):b\in[x,x']\}\ge0$; here $\sum_cQ_1(c,b)=L^{-2}$, since there are $L$ points $x$, each in one block, and $\sum_{b'}Q_1(c,b')=L$, so that the row sums are $a/L\le a$; and $b,b'$ are based in $W_c$, at coarse distance at most $1$. Third, the factors $\Pi_{\mathfrak L}$ on the two sides set every row and every column with an index $b\in\mathcal T$ to zero and leave the other entries unchanged. These facts show that the contribution of the first two terms of $q$ to $\varpi_q(\theta)$ is at most $(16+a)(e^{2\theta}-1)\le2e^2(16+a)\theta$ when $\theta\le1$.

For the third term, $|(\cdot)(b,b')|\le16K_*^2\mathsf g(\frac{\lambda_*}2)e^{\lambda_*}e^{-\frac12\lambda_*n}$, where $n$ is the distance of the base sites of $b$ and $b'$: each bond lies in four plaquettes, and the plaquettes at $b$ lie at distance $\ge n-2$ from those at $b'$. Since $e^{\theta{\rm dist}}-1\le\theta(1+n)e^{\theta(1+n)}$ and three bonds are based at each site, the third term contributes, for $\theta\le\frac{\lambda_*}4$, at most $48e^{\lambda_*+\theta}K_*^2\mathsf g(\frac{\lambda_*}2)\mathsf g_2(\frac{\lambda_*}4)\theta$. So $\varpi_q(\theta)\le A_q\theta$, and $\varpi_q(\theta_D)\le\frac12\gamma_q$.

Further, $\|q\|_{\theta_D}\le\|q\|_0+\varpi_q(\theta_D)\le\frac12A_q+1\le A_q$. The first inequality rests on $e^{\theta{\rm dist}}=1+(e^{\theta{\rm dist}}-1)$. For $\|q\|_0\le\frac12A_q$, note that the first two terms of $q$ have absolute row sums $\le16+a\le e^2(16+a)$, while the third term, by the bound on its entries and because three bonds are based at each site, has absolute row sums \[\le48e^{\lambda_*}K_*^2\mathsf g(\frac{\lambda_*}2)\sum_{u\in\mathbb Z^3}e^{-\frac12\lambda_*|u|_\infty}=48e^{\lambda_*}K_*^2\mathsf g(\frac{\lambda_*}2)^2\le24e^2K_*^2\mathsf g(\frac{\lambda_*}2)\mathsf g_2(\frac{\lambda_*}4),\] because $2e^{\lambda_*}\le e^2$ and $\mathsf g(\frac{\lambda_*}2)\le\mathsf g(\frac{\lambda_*}4)\le\mathsf g_2(\frac{\lambda_*}4)$ by \Cref{Y-A-Sums}, whose bound $\mathsf g(\lambda)\le\mathsf g_2(\lambda)$ we apply at $\lambda=\frac{\lambda_*}4$; and $q$ is symmetric, so its column sums equal its row sums. Finally $\varpi_q(\theta_D)\le\frac12\gamma_q\le1\le\frac12A_q$, because $A_q\ge32e^2$.

\emph{Basic kernels; the kernel $\mathsf R$, the rows $\ell$, the local maps.} The basic kernels $\mathfrak R^{-1},\mathsf B',\hat R,\hat{\mathsf J}^{-2}$ are bounded entrywise by $\frac{2.1L^2}{\gamma_*}e^{-2\theta_D{\rm dist}}$. Since $\ell\in\mathcal R(b_*,K_*,\lambda_*,\frac19)$ and $L\ge L_A$, we have $\lambda'=\lambda_*$ by part \textup{(a)} of the proof of \Cref{Y-A-Th29}. Hence the bound of \Cref{Y-A-P28}\,(a) holds with $\lambda''=\lambda_*$, and \Cref{Y-A-Th33} applies with $\lambda^\circ=\lambda_*$. For $\mathfrak R^{-1}$ the entrywise bound follows from \Cref{Y-A-P28}\,(a) with $\lambda''=\lambda_*$, together with $\mathsf D_i\ge L^{-2}$, $\mathsf d>{\rm dist}-1$ and $2e^{\lambda_*}\le2.1$. For $\mathsf B'$ and $\hat R$ it follows from \Cref{Y-A-Th33}\,(c), with $1+K_*^2\mathsf g(\frac{\lambda_*}2)=1+9\mathsf g(\frac{\lambda_*}2)/\gamma_*\le L^2/\gamma_*$, by $L\ge L_A\ge36\mathsf g(\frac{\lambda_*}2)/\gamma_*$. For $\hat{\mathsf J}^{-2}$ it follows from the bound on $\hat{\mathsf J}^{-2}$ of \Cref{Y-A-Jinv}. In each of the four bounds just cited the rate is at least $2\theta_D$, since $2\theta_D\le\lambda_2^*\le\frac14\lambda_*$ and $2\theta_D\le\theta_J$. A block holds at most $4L^3$ cells of one kind, so that $\|\cdot\|_{\theta_D}\le\mathsf C_0$ for each basic kernel.

For $\mathsf R$ let $n=|p-p'|$, so that ${\rm dist}<n/L+1$. Since three plaquettes are based at each site, we have $\sum_{p'}2e^{-\theta_R^*n}e^{\theta_D(n/L+1)}\le2\cdot3\cdot e\,\mathsf g(\theta_R^*-\theta_D)\le6e\,\mathsf g(\frac12\theta_R^*)=\mathsf C_R$. In the first inequality we also use that $\theta_D\le1$ and $L\ge1$ give $e^{\theta_D(n/L+1)}\le e\cdot e^{\theta_Dn}$, and that the sum over the torus is at most the sum over $\mathbb Z^3$; the second inequality follows from $\theta_D\le\frac12\lambda_2^*=\frac12\theta_R^*$. For $\ell$ we have $\|\ell\|_{\theta_D}\le\mathsf C_\ell$, and $\|\ell K\|_\theta\le\|\ell\|_\theta\|K\|_\theta$ for every kernel $K$ whose first index is a plaquette, because $|(\ell K)_y(j)|\le\sum_p|\ell_y(p)||K(p,j)|$. The local maps $W,W^*,\delta,\mathrm{d},\mathrm{d}^*,\mathrm{d}',\delta',\mathrm{d}'^*,Q_3,Q_3^*,E,E_3,\iota$ have coarse range at most $1$ and row and column sums at most $L^3$, so $\|\cdot\|_{\theta_D}\le4L^3$ for each of them.

\emph{Composites.} The composite kernels are $H^\sharp=\mathsf R\mathbb Y\mathfrak R^{-1}\iota$, $\mathsf T=H^\sharp\hat R$ and $\ell\mathsf T$; $\mathsf X={\rm pr}_\Xi\mathfrak R^{-1}\iota\hat R+L^{-1}Q_3^*\hat{\mathsf J}^{-2}\mathrm{d}'\hat R$; $\mathsf N^\sharp=(aL)^{-1}\mathrm{d}'^*\mathsf B'$; $(\mathsf G_Z,\mathsf G_\Xi)=\mathfrak R^{-1}(0,E_3\cdot)$, $\mathsf G=\mathsf R\mathbb Y\mathfrak R^{-1}(0,E_3\cdot)$ and $\ell\mathsf G$; and $\mathsf K_1=E_3^{\mathsf T}\mathsf X\mathsf B'$, $\mathsf K_2=E_3^{\mathsf T}\mathsf G_\Xi-L^{-1}\hat{\mathsf J}^{-2}$. Each of them is a sum of at most two products, each of at most three basic kernels, four local maps, one $\mathsf R$ and one $\ell$; hence each has norm $\|\cdot\|_{\theta_D}\le2\mathsf C_0^3(4L^3)^4\mathsf C_R\mathsf C_\ell$. The kernel $\mathsf g=\Pi_{\mathfrak L}\mathrm{d}^*\mathsf BE$ has norm at most $(4L^3)^2(1+\mathsf C_\ell^2)$. For the rows $\ell'$, $\|\ell'\|_{\theta_D}\le3K_*\mathsf g(\frac{\lambda_*}2)\le\mathsf C_\ell$ by \Cref{Y-A-Th29}\,(a). Since $\mathsf C_0\ge2$ and $\mathsf C_0,\mathsf C_R,\mathsf C_\ell,4L^3\ge1$, we have $2\mathsf C_0^3(4L^3)^4\mathsf C_R\mathsf C_\ell\le(4L^3\mathsf C_0)^4\mathsf C_R\mathsf C_\ell\le C_D$, and the same holds for products with fewer factors. Moreover $\|\mathsf B'\|_{\theta_D}\le\mathsf C_0\le C_D$, $\|\ell\|_{\theta_D},\|\ell'\|_{\theta_D}\le\mathsf C_\ell\le C_D$, $\|q\|_{\theta_D}\le A_q\le C_D$ and $\|\mathsf g\|_{\theta_D}\le(4L^3)^2(1+\mathsf C_\ell^2)\le C_D$. The composite kernels, $\mathsf B'$, $\ell$, $\ell'$, $q$ and $\mathsf g$ are the sixteen kernels of \Cref{Y-A-Dec}, so each of them has norm $\|\cdot\|_{\theta_D}$ at most $C_D$. The bound $\varpi_q(\theta_D)\le\frac12\gamma_q$ was shown above, and we also have $C_D\ge4$ and $\theta_D\le1$. \end{proof}

\subsubsection{(S-chain); $K$ steps}

\begin{definition}[The order of the letters; the choice of $\rho_0$, $L_*$ and $x_\flat$]\label{Y-A-OrderLetters} The letters are chosen in this order: $d_0$; $c_\Xi$; $\rho_0$; $L,a$; $\kappa$; $(C_D,\theta_D)$, which fix $c_R=64/\theta_D$; $M$; $\mathfrak p_0$, $A_0$; $c_\xi$; $c_h$; $E_a$; and, last of all, $x_\flat$.

Given $0<d_0<\frac12$ and $0<c_\Xi\le\frac12(\frac12-d_0)$, we define \begin{equation}\label{eq:row-mask-radius}\rho_0:=\text{the least integer}\ge0\text{ with}\quad3K_*^2\,\mathsf g(\lambda_*)\,\mathsf g(\tfrac{\lambda_*}2)\,e^{-\lambda_*\rho_0/2}\le\tfrac18\bigl(\tfrac12-d_0-c_\Xi\bigr);\end{equation} and, with this $\rho_0$, the four numbers \begin{equation}\label{eq:row-cut-norms}\mathsf N_Z:=1+3K_*^2\mathsf g(\lambda_*)^2,\quad\mathsf P_\rho:=1+3\bigl(2+3(2\rho_0+1)^3\bigr)K_*^2\mathsf g(\lambda_*)^2,\quad\mathsf N_\sharp:=c_1^*\mathsf l_L\mathsf N_Z\mathsf P_\rho,\quad\mathsf N_\Xi:=c_2^*\mathsf N_Z.\end{equation} The number $L_*(d_0,c_\Xi)$, which is defined next, and the conditions \textup{(L1)}, \textup{(L2)} are those of this section. In item~3 of \S\ref{Y-G0-1}, their place is taken by $L_*^{\rm su}$, \textup{(L1$^{\rm su}$)} and \textup{(L2$^{\rm su}$)} of \textup{(O4)}. As to the size of $L$, \Cref{Y-A-Fr} requires $L\ge L_A$ alone.

We let $L_*(d_0,c_\Xi)$ be the least integer $\ge L_A$ for which the two inequalities \[\bigl(c_1^*\mathsf l_L\mathsf P_\rho L^{-1}+60+4c_2^*\bigr)\mathsf N_Z\le\tfrac14d_0L^{1/2},\qquad115\,(c_1^*\mathsf l_L\mathsf N_Z\mathsf P_\rho)^2\le L\] hold for every real $L\ge L_*$; such an integer exists, because $\mathsf l_L^2/L\to0$. The floor condition on $x_\flat$ is \textup{(X12)}: $3K_*\mathsf g(\frac{\lambda_*}2)\,x_\flat^{-e_R}/(1-L^{-e_R})\le\frac19$, where $e_R:=\frac12\lambda_*c_R$.\end{definition}

\begin{theorem}[Decay and bounds along the cut chain]\label{Y-A-Th36} Let $0<d_0<\frac12$ and $0<c_\Xi\le\frac12(\frac12-d_0)$, let $\rho_0$ be given by \eqref{eq:row-mask-radius}, let $L\ge L_*(d_0,c_\Xi)$ be odd and $a\ge3/(c_LL)$, let $(C_D,\theta_D)$ be the constants of \Cref{Y-A-P35}, $c_R=64/\theta_D$ and $\varrho(x)=\lceil c_R\log x\rceil$, and let $x_\flat>1$ satisfy \textup{(X12)}. For every $x_0$, consider the chain $x_k=L^{-k}x_0$, $\ell_0=\emptyset$, $\ell_{k+1}=((\ell_k)')^{\rm c}$, in which the cut is taken at $\varrho(x_{k+1})$. Then, for every $k$ with $x_{k+1}\ge x_\flat$ and every tower of $K\ge k+1$ steps with $N\ge2$ even, the pair $(\ell_k,(\ell_k)')$ satisfies \textup{(Dec)}$(C_D,\theta_D)$, and it satisfies \textup{(N$'$)} with the numbers \eqref{eq:row-cut-norms} and with this $\rho_0$, whatever the radius at which $(\ell_k)'$ is cut; moreover $\mathsf t_\rho\le\frac18(\frac12-d_0-c_\Xi)$, and \textup{(L1)}, \textup{(L2)} hold; we call the assertion of this theorem, as a whole, \textup{(S-chain)}.\end{theorem}

\begin{proof}\emph{Step (Cut): the condition \textup{(Cut)} along the chain.} Since $\varrho(x)\ge c_R\log x$ and $e_R=\frac12\lambda_*c_R$, we have $\mathsf t_*(\varrho(x))\le3K_*\mathsf g(\frac{\lambda_*}2)x^{-e_R}$. The inequality $x_j\ge L^{k+1-j}x_\flat$ holds for every $j\le k+1$, as $x_{k+1}\ge x_\flat$; summing the resulting geometric series in $j$, we find that $\sum_{j\le k+1}\mathsf t_*(\varrho(x_j))$ is at most the left side of \textup{(X12)}, hence at most $\frac19$. This choice of the radii is not circular: the number $L_A$, \Cref{Y-A-Th29} and \Cref{Y-A-P35} depend on no cut radius, the number $c_R$ is defined after them, and the floor \textup{(X12)} is chosen last. We may therefore apply \Cref{Y-A-Th29}\,(b) to the tower of $K'=k+1$ steps and side $N'=NL^{K-k-1}$, an even number. That tower has the same tori $\mathbb T_0$ to $\mathbb T_{k+1}$ and the same rows $\ell_0$ to $\ell_{k+1}$, because these rows are determined by the radii $\varrho(x_1)$ to $\varrho(x_{k+1})$ alone; and its condition \textup{(Cut)}, which reads $\sum_{j\le k+1}\mathsf t_*(\varrho(x_j))\le\frac19$, is the bound just proved. The argument is independent of the radii $\varrho(x_j)$, $j>k+1$, of a tower with $K>k+1$; for such $j$, $x_j$ can be smaller than $x_\flat$. \Cref{Y-A-Th29}\,(b) thus yields $\ell_k\in\mathcal R(b_*,K_*,\lambda_*,\frac19)$ and $|(\ell_k)'_{y'}(P)|\le K_*e^{-\lambda_*{\rm dist}(y',P)}$.

\emph{The condition \textup{(Dec)} and the bound on $\mathsf t_\rho$.} \Cref{Y-A-P35} gives \textup{(Dec)}$(C_D,\theta_D)$. Next, the chain of inequalities \[\mathsf t_\rho\le3K_*\mathsf g(\lambda_*)\cdot K_*\sum_{|u|_\infty>\rho_0}e^{-\lambda_*|u|_\infty}\le3K_*^2\mathsf g(\lambda_*)\mathsf g(\frac{\lambda_*}2)e^{-\lambda_*\rho_0/2}\] holds, the second inequality because $e^{-\lambda_*|u|_\infty}\le e^{-\lambda_*\rho_0/2}e^{-\lambda_*|u|_\infty/2}$ for $|u|_\infty>\rho_0$; by \eqref{eq:row-mask-radius}, its right side is at most $\frac18(\frac12-d_0-c_\Xi)$.

\emph{The conditions \textup{(N$'$)}, \textup{(L1)} and \textup{(L2)}.} For a mask $m$ and $|\zeta'|_\infty\le1$ we have $|LZ|_\infty=|\mathsf B'^{\rm c}_m\zeta'|_\infty\le1+\|\ell'^{\rm c}\|_{\rm c}\|\ell'^{\rm c}\|_{\rm r}\le\mathsf N_Z$. The two inequalities hold because the bound $|(\ell_k)'_{y'}(P)|\le K_*e^{-\lambda_*{\rm dist}(y',P)}$ at the end of Step~(Cut) gives $\|\ell'^{\rm c}\|_{\rm r}\le3K_*\mathsf g(\lambda_*)$ and $\|\ell'^{\rm c}\|_{\rm c}\le K_*\mathsf g(\lambda_*)$ for every radius of the cut, and because a mask raises neither norm. Here $Z=L^{-1}\mathsf B'^{\rm c}_m\zeta'$, $\Xi$, $(\nu,\omega)={\rm pair}(Z,\Xi)$ and $\sigma$ are the objects of \textup{(N$'$)} in \Cref{Y-A-Np}, with $\ell=\ell_k$; by that definition and by \Cref{Y-A-XT}, taken at $v=LZ$, they satisfy the identities $\nu=\mathsf R\mathsf T_\sharp Z=\mathsf T(LZ)$, $\Xi=\Xi_Z=\mathsf X(LZ)$ and $\sigma=\ell_k\nu$. The first two bounds in the second paragraph of \Cref{Y-A-SetsOfRows} are $\|\ell\|_{\rm r}\le3K\mathsf g(\lambda)$ and $\|\ell\|_{\rm c}\le K\mathsf g(\lambda)$. They hold for the rows $\ell_k$ of $\mathcal R(b_*,K_*,\lambda_*,\frac19)$ and give $|\sigma|_{2,\infty}\le3K_*\mathsf g(\lambda_*)|\nu|_\infty$ and $2\|\ell_k\|_{\rm c}+(2\rho_0+1)^3\|\ell_k\|_{\rm r}\le(2+3(2\rho_0+1)^3)K_*\mathsf g(\lambda_*)$. \Cref{Y-A-Th33}\,(d), applied, as in \Cref{Y-A-Stars}, with $(b,K,\lambda,\lambda^\circ)=(b_*,K_*,\lambda_*,\lambda_*)$, where $\lambda'=\lambda_*$ because $L\ge L_A$, now gives \textup{(N$'$)} with the numbers \eqref{eq:row-cut-norms}. Finally, \textup{(L1)} and \textup{(L2)} follow from the definition of $L_*$ in \Cref{Y-A-OrderLetters}.\end{proof}

\begin{theorem}[Uniform bounds in the cut tower]\label{Y-A-Fr}
Use the one-step geometry of \Cref{Y-A-Lattices} on the tower
of \Cref{Y-A-Towers}. In that geometry the coarse side at
step $k$ is $N'=NL^{K-k-1}$. The kernels are the scalar kernels defined
above, and $\mathsf l_L=1+\log L$.

\emph{The chain.} We put $x_j=L^{-j}x_0$; $\ell_0$ is empty, and $\ell_{j+1}=((\ell_j)')^{\rm c}$, with the cut taken at $\varrho(x_{j+1})$, where $\varrho(x)=\lceil c_R\log x\rceil$.

\emph{The order of the constants.} The constants are fixed in the order \textup{(0)} to \textup{(8)}.
\begin{enumerate}
\item[\textup{(0)}] The pure numbers are defined exactly by
\[
w_*=3^{-6},\quad
\mathsf r_\sharp=\frac{w_*}{1.3\cdot10^7},\quad
\gamma_1=\frac{\mathsf r_\sharp}{3074},\quad
b_*=\frac2{\mathsf r_\sharp}=1.8954\cdot10^{10},\quad
\gamma_*=\frac{\gamma_1}{2b_*}=\frac{\mathsf r_\sharp^2}{12296},
\]
\[
\lambda_*=\frac{\gamma_*^{1/2}}{44},\qquad
K_*=3\gamma_*^{-1/2},\qquad
\theta_J=\tfrac12\log(1+w_*/50).
\]
The following outward intervals, rather than rounded equalities, may be used in numerical bounds:
\[
\begin{aligned}
\mathsf r_\sharp&\in[1.05518,1.05519]\cdot10^{-10},&
\gamma_1&\in[3.43261,3.43262]\cdot10^{-14},\\
\gamma_*&\in[9.05512,9.05513]\cdot10^{-25},&
\lambda_*&\in[2.16269,2.16270]\cdot10^{-14},\\
K_*&\in[3.15263,3.15264]\cdot10^{12},&
\theta_J&\in[1.37172,1.37173]\cdot10^{-5}.
\end{aligned}
\]
The number $L_A=\max\{12K_*\mathsf g(\lambda_*/2)\gamma_*^{-1/2},\ 38088K_*^2\mathsf g_1(\lambda_*)^2/\mathsf r_\sharp\}$ is at most $2.5\cdot10^{153}$, and $\lambda_2^*=\theta_R^*=\min\{\lambda_*/4,\ [6K_*^2\mathsf g(\lambda_*/2)\mathsf g_1(\lambda_*/4)]^{-1}\}$ is at least $1.9\cdot10^{-129}$. The pair $(c_{\mathcal G}^*,\omega_{\mathcal G}^*)$ is the one given by \Cref{Y-A-L31} at $(b,C,\theta)=(b_*,2,\theta_R^*)$; we put $\varpi^*=\omega_{\mathcal G}^*/2$ and $\lambda_B^*=\min\{\lambda_*,\varpi^*\}$. Then $c_\Xi^*=6000c_{\mathcal G}^*\mathsf g_2(\theta_R^*/4)/(\gamma_*\varpi^{*2})$, $c_H^*=180\mathsf g(\theta_R^*/2)c_{\mathcal G}^*c_\Xi^*/\varpi^*$ and $c_1^*=19c_H^*\mathsf g(\lambda_B^*)\mathsf g(\lambda_2^*)$; moreover $c_2^*=19c_\Xi^*\mathsf g(\lambda_B^*)\mathsf g(\lambda_2^*)+72w_*^{-1}\mathsf g(\theta_J)\mathsf g(\lambda_2^*)$, which is the $c_2^\natural$ of \Cref{Y-A-Th33}, taken at the starred values, and enters \textup{(Fr3$'$)}. Further, $\mathsf N_Z=1+3K_*^2\mathsf g(\lambda_*)^2$.
\item[\textup{(1)}] $d_0\in(0,\frac12)$.
\item[\textup{(2)}] $0<c_\Xi\le\frac12(\frac12-d_0)$.
\item[\textup{(3)}] The integer $\rho_0$ is the least integer $\ge0$ with $3K_*^2\mathsf g(\lambda_*)\mathsf g(\lambda_*/2)e^{-\lambda_*\rho_0/2}\le\frac18(\frac12-d_0-c_\Xi)$, and $\mathsf P_\rho=1+3(2+3(2\rho_0+1)^3)K_*^2\mathsf g(\lambda_*)^2$.
\item[\textup{(4)}] $L$ is odd, $L\ge L_A$, and $a\ge3/(c_LL)$; we put $\mathsf N_\sharp=c_1^*\mathsf l_L\mathsf N_Z\mathsf P_\rho$. Where the later sections apply the theorem, at the values of item~3 of \S\ref{Y-G0-1}, \textup{(O4)} of that item gives $L\ge L_*^{\rm su}(d_0,c_\Xi)\ge L_A$; here $L_*^{\rm su}$ is the number of \textup{(O4)}, built from the $\mathsf N_\nu$ of \textup{(O3)} in that item. This bound is not among the hypotheses of the theorem: the lower bound on $L$ that this theorem uses is $L\ge L_A$.
\item[\textup{(5)}] The letter $\kappa$.
\item[\textup{(6)}] The rate and the bound, defined in the following order: $\gamma=(73L^4)^{-1}$; $A^\flat=2e^2(16+a)+384e^8K_*^2\mathsf g_2(\lambda_*/2)^2$; $\theta_1=\min\{1,\lambda_*L/12\}$; $\theta_{\mathbb C}=\frac16\min\{\theta_1,\gamma/2A^\flat\}$; $\theta_D=\frac12\min\{\lambda_2^*,\theta_{\mathbb C},\theta_J\}$. After $\theta_D$ come $\mathsf C_0=4L^3\mathsf g(\theta_D)\,2.1L^2/\gamma_*$, $\mathsf C_R=6e\,\mathsf g(\theta_R^*/2)$, $\mathsf C_\ell=3eK_*\mathsf g(\lambda_*/2)$ and $A_q=2e^2(16+a)+48e^2K_*^2\mathsf g(\lambda_*/2)\mathsf g_2(\lambda_*/4)$; then $C_D=\max\{(4L^3\mathsf C_0)^4\mathsf C_R\mathsf C_\ell,\ (4L^3)^2(1+\mathsf C_\ell^2),\ A_q\}$, and $c_R=64/\theta_D$.
\item[\textup{(7)}] The letters $M$; $\mathfrak p_0$, $A_0$; $c_\xi$; $c_h$; $E_a$.
\item[\textup{(8)}] Last, a number $x_\flat>1$ satisfying \textup{(X12)}, that is, $3K_*\mathsf g(\lambda_*/2)x_\flat^{-e_R}/(1-L^{-e_R})\le\frac19$, where $e_R=\lambda_*c_R/2$.
\end{enumerate}
This order is not circular: $L_A$, \Cref{Y-A-Th29} and \Cref{Y-A-P35} are independent of the cut; $c_R$ is defined after them; and \textup{(X12)} is a lower bound on $x_\flat$.

\emph{Conclusion.} Let $x_0$ be arbitrary, let $k\ge0$ satisfy $x_{k+1}\ge x_\flat$, take any tower with $N\ge2$ even, and let $(\ell,\ell')=(\ell_k,(\ell_k)')$. Then the following hold.
\begin{enumerate}
\item[\textup{(Fr0)}] $\ell_k\in\mathcal R(b_*,K_*,\lambda_*,\frac19)$, of range $\le\varrho(x_k)$; here $\ell_0$ is empty, and for $k\ge1$ the range bound comes from the cut. Moreover $|(\ell_k)'_{y'}(P)|\le K_*e^{-\lambda_*{\rm dist}(y',P)}$ and $1\le\mathsf B'\le\frac32/\mathsf r_\sharp$.
\item[\textup{(Fr2)}] $\|K\|_{\theta_D}\le C_D$ for $K\in\{\mathsf G,\mathsf G_Z,H^\sharp,\mathsf N^\sharp,\mathsf T,\ell'\}$. The bound $\|\cdot\|_{\theta_D}\le C_D$ also holds for each of the sixteen kernels of \textup{(Dec)}, and $\varpi_q(\theta_D)\le\gamma/24\le\gamma_q/2$.
\item[\textup{(Fr3)}] At every cut radius, $L^2\|H^\sharp-\mathsf T(\mathsf B'-\mathsf B'^{\rm c})\|_\infty=L^2\|\mathsf T\mathsf B'^{\rm c}\|_\infty\le c_1^*\mathsf l_L\mathsf N_Z\le\mathsf N_\sharp$, and $\|\mathsf B'^{\rm c}\|_\infty\le\mathsf N_Z$.
\item[\textup{(Fr3$'$)}] Condition \textup{(N$'$)} of \Cref{Y-A-Np} holds with the numbers \eqref{eq:row-cut-norms}. Explicitly, let a cut radius, a mask $m$ and $\zeta'$ with $|\zeta'|_\infty\le1$ be arbitrary, and let $\mathsf B'^{\rm c}_m=1+\ell'^{{\rm c}\mathsf T}m\,\ell'^{\rm c}$, $Z=L^{-1}\mathsf B'^{\rm c}_m\zeta'$, $\Xi=\Xi_Z$, $(\nu,\omega)={\rm pair}(Z,\Xi)$ and $\sigma=\ell\nu$; then $L^2[|\nu|_\infty+(2\|\ell\|_{\rm c}+(2\rho_0+1)^3\|\ell\|_{\rm r})|\sigma|_{2,\infty}]\le\mathsf N_\sharp$, $L|Z|_\infty\le\mathsf N_Z$ and $L|\Xi|_\infty\le\mathsf N_\Xi:=c_2^*\mathsf N_Z$; here $\mathsf N_\Xi$ is a pure number.
\item[\textup{(Fr4)}] At every cut radius, $\mathsf t_\rho\le\frac18(\frac12-d_0-c_\Xi)$, with $\mathsf t_\rho$ the number of \Cref{Y-A-CutRows}.
\end{enumerate}
Concerning dependence: $C_D$ and $\theta_D$ are determined by $(L,a)$ together with pure numbers; $\mathsf N_Z$ is a pure number; and $\mathsf N_\sharp$ depends on $L$ only by way of $\mathsf l_L$, and on $(d_0,c_\Xi)$ only by way of $\rho_0$. Each of these numbers is independent of $k$, $N$, $K$, $x_0$, of the cuts and of the rows.
\end{theorem}

\begin{remark}[The four hypotheses]
The proof of \Cref{Y-A-Fr} rests on four hypotheses: $L\ge L_A$ odd; $aL\ge3/c_L$, which is stronger than $aL\ge18$; $N$ even; and $x_{k+1}\ge x_\flat$, with $x_\flat$ satisfying \textup{(X12)}. In \S\ref{Y-G0-1}, item~1 requires $N$ to be even; \textup{(O4)} of item~3 requires $L$ to be odd, with $L\ge L_*^{\rm su}\ge L_A$, and $a\ge3/(c_LL)$; and \textup{(O13)} of item~3 requires \textup{(X12)}. By Step~H16\,(v) of the proof of \Cref{Y-A-Fr}, condition \textup{(X12)} holds automatically for every $x_\flat\ge2$. Part \textup{(Fr0)} is proved in Step~H9, from \Cref{Y-A-Th29}\,(a) and~(b); parts \textup{(Fr3$'$)} and \textup{(Fr4)} are proved in Step~H14.
\end{remark}

\begin{proof}[Proof of \Cref{Y-A-Fr}]
We apply the one-step estimates to the cut tower and check the chosen
rate for the Gaussian window. The estimates for a single pair of rows
are proved in Props.~\ref{Y-A-P28}, \ref{Y-A-P35} and
\Cref{Y-A-Th33}.

\emph{Step H9. The cut tower.}
For $j\le k+1$, we have $x_j\ge L^{k+1-j}x_\flat$. The definition of
$\varrho$ and \textup{(X12)} give
\[
 \sum_{j=1}^{k+1}\mathsf t_*(\varrho(x_j))
 \le \frac{3K_*\mathsf g(\lambda_*/2)x_\flat^{-e_R}}{1-L^{-e_R}}
 \le \frac19 .
\]
Thus \Cref{Y-A-Th29}(b) applies to the first $k+1$ steps, viewed
as a tower with final side $N'=NL^{K-k-1}$. It gives \textup{(Fr0)}.
Only radii through step $k+1$ enter this argument. The numerical bounds
on $L_A$ and $\lambda_2^*$ in item~(0) follow by inserting
$\mathsf g(\lambda)\le51\lambda^{-3}$ and
$\mathsf g_1(\lambda)\le182\lambda^{-4}$ in their defining formulas.

\emph{Step H13. The decay rate for the window.}
The proof of \Cref{Y-A-P35} applies at any rate $\theta_D$
with $\theta_D\le1$, $2\theta_D\le\min\{\lambda_2^*,\theta_J\}$ and
$\varpi_q(\theta_D)\le\gamma_q/2$. Its constants are evaluated at that
rate. The choice in item~(6) satisfies the first two conditions and
gives $\theta_D\le\lambda_*/4$, since $\lambda_2^*\le\lambda_*/4$.
The bound $\varpi_q(\theta)\le A_q\theta$ in that proof
and $A_q\le A^\flat$ from Step~H16(ii) give
$\varpi_q(\theta_D)\le A_q\gamma/(24A^\flat)\le\gamma/24$.
By Step~H9 the rows are in the required class, so these estimates give
\textup{(Fr2)}.

As to the group, tensoring a scalar kernel with $1_{\mathfrak g}$
preserves the norms in \Cref{Y-A-Norm}, since the norm of
each matrix entry becomes the operator norm of that scalar times the
identity.

\emph{Step H14: \Cref{Y-A-Th36}; \textup{(Cut)}, \textup{(N$'$)}, and hence \textup{(Fr3)}.} For every $k$ with $x_{k+1}\ge x_\flat$, the chain cut at $\varrho(x)=\lceil c_R\log x\rceil$ satisfies: $\ell_k\in\mathcal R(b_*,K_*,\lambda_*,\frac19)$; \textup{(Dec)} holds; \textup{(N$'$)} holds with \eqref{eq:row-cut-norms}, at every cut; and $\mathsf t_\rho\le\frac18(\frac12-d_0-c_\Xi)$.

Condition \textup{(Cut)} holds by Step~H9, and \Cref{Y-A-Th29}\,(b) then applies; by Step~H13, \textup{(Dec)}$(C_D,\theta_D)$, with $C_D$ and $\theta_D$ as in item~\textup{(6)} of the theorem, holds for $(\ell_k,(\ell_k)')$.

For \textup{(N$'$)}, we have $|\mathsf B'^{\rm c}_m\zeta'|_\infty\le1+\|\ell'^{\rm c}\|_{\rm c}\|\ell'^{\rm c}\|_{\rm r}\le1+K_*\mathsf g_*\cdot3K_*\mathsf g_*=\mathsf N_Z$, where $\mathsf g_*=\mathsf g(\lambda_*)$; here cutting and masking lower both norms. As $\nu=\mathsf T(LZ)$, \Cref{Y-A-Th33}\,(d) at the starred values of \Cref{Y-A-Stars} gives $L^2|\nu|_\infty\le c_1^*\mathsf l_L\mathsf N_Z$. Further, $|\sigma|_{2,\infty}\le3K_*\mathsf g_*|\nu|_\infty$ and $2\|\ell\|_{\rm c}+(2\rho_0+1)^3\|\ell\|_{\rm r}\le(2+3(2\rho_0+1)^3)K_*\mathsf g_*$, and together these give the factor $\mathsf P_\rho$. As $\Xi=\mathsf X(LZ)$, \Cref{Y-A-Th33}\,(d) at the starred values gives $L|\Xi|_\infty\le c_2^*\mathsf N_Z=\mathsf N_\Xi$. This is \textup{(Fr3$'$)} at every mask and at every cut. The bound on $\mathsf t_\rho$ at the end of this step, with the choice \eqref{eq:row-mask-radius} of $\rho_0$, is \textup{(Fr4)}; it holds at every cut, because cutting lowers both factors of $\mathsf t_\rho$. The step rests on $L\ge L_A$, $aL\ge3/c_L$, the evenness of $N$ and $x_{k+1}\ge x_\flat$, and it depends on no other condition on $L$.

For \textup{(Fr3)}, the identity $H^\sharp=\mathsf T\mathsf B'$ gives $H^\sharp-\mathsf T(\mathsf B'-\mathsf B'^{\rm c})=\mathsf T\mathsf B'^{\rm c}$ exactly. Hence $L^2\|\cdot\|_\infty\le L^2\|\mathsf T\|_\infty\|\mathsf B'^{\rm c}\|_\infty\le c_1^*\mathsf l_L\mathsf N_Z\le\mathsf N_\sharp$, since $\mathsf P_\rho\ge1$; this bound does not depend on the cut radius. Finally $\mathsf t_\rho\le3K_*\mathsf g(\lambda_*)K_*\sum_{|u|>\rho_0}e^{-\lambda_*|u|}\le3K_*^2\mathsf g(\lambda_*)\mathsf g(\lambda_*/2)e^{-\lambda_*\rho_0/2}$, and by \eqref{eq:row-mask-radius} the last member is $\le\frac18(\frac12-d_0-c_\Xi)$.

\emph{Step H16: the passage to $c_R=64/\theta_D$ with $\theta_D\le\theta_{\mathbb C}/2$, and \textup{(X12)}.} We prove five assertions. \textup{(i)} For $\theta'\le\theta$, \textup{(Dec)}$(C,\theta)$ implies \textup{(Dec)}$(C,\theta')$. \textup{(ii)} $A_q\le A^\flat$, and $\theta_{\mathbb C}\le\theta_q/6$, with $\theta_q$ that of \Cref{Y-A-P35}; hence, writing $\theta_D^{\rm sc}$ for the $\theta_D$ given by the display of \Cref{Y-A-P35}, the smaller of $\theta_D^{\rm sc}$ and $\theta_{\mathbb C}/2$ is the $\theta_D$ of \Cref{Y-A-Fr}. \textup{(iii)} When the chain is cut with the larger $c_R$, it remains in the absorbing set $\mathcal R(b_*,K_*,\lambda_*,\frac19)$. \textup{(iv)} \textup{(N$'$)} holds at every cut. \textup{(v)} \textup{(X12)} holds. In \Cref{Y-A-Fr} and in its proof, $\theta_D$, $C_D$ and $c_R$ are those of its item~\textup{(6)}, and $e_R=\lambda_*c_R/2$ is that of its item~\textup{(8)}, whereas the display of \Cref{Y-A-P35} yields $\theta_D^{\rm sc}$ together with the $C_D$ evaluated at $\theta_D^{\rm sc}$.

For \textup{(i)}: $e^{\theta{\rm dist}}$ and $e^{\theta{\rm dist}}-1$ increase with $\theta$.

For \textup{(ii)}, we first treat $A_q\le A^\flat$, for which it is enough that $\mathsf g(\mu)\mathsf g_2(\mu/2)\le8e^6\mathsf g_2(\mu)^2$ at $\mu=\lambda_*/2$. To begin with, $\mathsf g(\mu)\le\mathsf g_2(\mu)$. Writing $u=2v+\varepsilon$ with $\varepsilon\in\{0,1\}^3$, we get $(1+|u|)^2\le4(1+|v|)^2$ and $e^{-\mu|u|/2}\le e^{\mu/2}e^{-\mu|v|}$. Since $\mu=\lambda_*/2$ with $\lambda_*=\gamma_*^{1/2}/44\le\frac1{44}$, we have $\mu\le\frac12$ and hence $e^{\mu/2}\le e^{1/4}$. So $\mathsf g_2(\mu/2)\le32e^{1/4}\mathsf g_2(\mu)\le41.1\,\mathsf g_2(\mu)\le3227\,\mathsf g_2(\mu)$, where $32=8\cdot4$ accounts for the eight values of $\varepsilon$ and the factor $4$. Since $8e^6=3227.4\dots$, we have $3227\le8e^6$, and together with $\mathsf g(\mu)\le\mathsf g_2(\mu)$ this gives $\mathsf g(\mu)\mathsf g_2(\mu/2)\le8e^6\mathsf g_2(\mu)^2$.

Hence $\theta_{\mathbb C}\le\gamma/(12A^\flat)\le\frac16\gamma_q/(2A_q)$. Also $\theta_{\mathbb C}\le1/(12\cdot384e^8\cdot9/\gamma_*)\le\gamma_*/(1.236\cdot10^8)\le\gamma_*^{1/2}/1056=\lambda_*/24$, since $\gamma_*^{1/2}\le1$ and $44\cdot24=1056$. Combining this with the preceding bound, $\theta_{\mathbb C}\le\frac16\min\{\lambda_*/4,\gamma_q/(2A_q)\}=\theta_q/6$. Consequently $\min\{\theta_D^{\rm sc},\theta_{\mathbb C}/2\}=\frac12\min\{\lambda_2^*,\theta_{\mathbb C},\theta_J\}$. We define $C_D$ with $\mathsf g(\theta_D)$ at this $\theta_D$, following item~\textup{(6)} of \Cref{Y-A-Fr}; by \textup{(i)}, the $C_D$ evaluated at $\theta_D^{\rm sc}$, which is the smaller one, would suffice equally well.

For \textup{(iii)}: \Cref{Y-A-Th29}\,(a) and \Cref{Y-A-P35} hold for every $\ell$ in the set, and neither of these depends on a radius. The radii enter through \textup{(Cut)} alone, and $\mathsf t_*(\varrho)$ decreases as $\varrho$ increases. For \textup{(iv)}: by Step~H14, $\|\ell'^{\rm c}\|_{{\rm r},{\rm c}}\le\|\ell'\|_{{\rm r},{\rm c}}$ at every radius.

For \textup{(v)}: $x^{-e_R}/(1-L^{-e_R})$ decreases as $e_R=\lambda_*c_R/2$ increases, so a larger $c_R$ helps. Moreover $\theta_D\le\lambda_2^*/2\le\lambda_*/8$ gives $e_R=32\lambda_*/\theta_D\ge256$, and \[27K_*\mathsf g(\lambda_*/2)\le27\cdot3.153\cdot10^{12}\cdot4.034\cdot10^{43}\le3.44\cdot10^{57}<1.157\cdot10^{77}\le2^{256}(1-3^{-256}).\] Here the bound $\mathsf g(\lambda)\le51\lambda^{-3}$ of \Cref{Y-A-Sums} gives $\mathsf g(\lambda_*/2)\le51(\lambda_*/2)^{-3}\le4.034\cdot10^{43}$. So \textup{(X12)} holds for every $x_\flat\ge2$ with $c_R=64/\theta_D$, the pair $(\theta_D,c_R)$ being that of item~\textup{(6)} of \Cref{Y-A-Fr}, and also with $c_R=64/\theta_D^{\rm sc}$: both $\theta_D$ and $\theta_D^{\rm sc}$ are $\le\lambda_2^*/2$, so that $e_R\ge256$ in both cases.

In addition, $e_R\ge64\lambda_*/\lambda_2^*$, because $\theta_D\le\lambda_2^*/2$. The step depends on $e_R$ only through $e_R\ge256$. The requirement $x_\flat\ge e^2$ is part of \textup{(O13)}, in item~3 of \S\ref{Y-G0-1} and in item~1 of \S\ref{Y-J-1}; so \textup{(X12)} imposes no lower bound of its own on $x_\flat$. \end{proof}

\section{Local Gaussian integration}\label{Y-C-1}

In this section a block field $V'$ is given. On the fine bonds off the combs we study a quadratic form $q^V_\ell$, built with rows and dressed with transports of $V'$, and from it we construct a Gaussian law of finite range. Part C.I shows that $\gamma\le q\le b$ at every $V'$ and proves a weighted bound; it then forms a matrix $G$ row by row from windows, establishes the properties \textup{(a)} to \textup{(k)} of the law $\varphi=G^V\xi$, whose range is at most $R$, and gives an identity that writes the step as a mean over the white noise $\xi$. Part C.II proves, for this law, the bounds on the centre, \textup{(S-cen$_k$)}, and the locality in $V'$ of $\log\det q^V_\ell$ and of the remainder, \textup{(S-loc$_k$)}. Item C.II 4, in which \textup{(S-cen$_k$)} is established, uses the free scalar model \textup{(Fr1)} to \textup{(Fr3)}, which is stated in item C.II 2.2 and proved in \S\ref{Y-A-1}, whereas \textup{(S-loc$_k$)} is independent of \textup{(Fr1)} to \textup{(Fr3)}.

\begin{remark}[The level and the rows] The level is of arbitrary depth. The rows are those of item C.I 1.1 and of item C.II 1.2, with $K\ge0$; the empty rows, for which $\mathsf H_y=\{0\}$ for every $y$, are therefore admitted.\end{remark}

\subsection*{Part C.I. \textup{(Law$_k$)}: the law of windows}

\subsubsection*{C.I 1. Objects}

\begin{remark}[Notation] In Part C.I some letters have two senses, and the context makes clear which one is meant: $b$ is the number of item C.I 1.2 or a bond; $G$ is the group or the matrix of the law of item C.I 3; $Y$ is a subset of $I^\circ$ or a matrix in $\mathfrak g_{\mathbb C}$; $x$ is a fine site or the parameter of the law; $\mathsf g$ denotes either the sums $\mathsf g(\lambda)$ and $\mathsf g_2(\lambda)$ or the kernel $\mathsf g^V$.\end{remark}

\subsubsection*{C.I 1.1. Group, lattices, transports, form}

Let $G=\operatorname{SU}(2)$, let $\mathfrak g$ be the space of hermitian traceless $2\times2$ matrices, and write $\operatorname{Ad}_gY=gYg^{-1}$. A field $V$ is a map from the oriented bonds to $G$, or to $\operatorname{SL}(2,\mathbb C)$, such that $V(\bar b)=V(b)^{-1}$, where $\bar b$ is the bond $b$ with the opposite orientation. For a path $\pi$ from $\pi_-$ to $\pi_+$, $V(\pi)$ is the product of the values of $V$ on the bonds of $\pi$, taken in the order in which $\pi$ runs through them. For a gauge $u$ we set $V^u(b)=u(b_-)V(b)u(b_+)^{-1}$, where $b_-$ and $b_+$ are the initial and the final site of $b$. For a plaquette $p=(z;\mu\nu)$ with $\mu<\nu$ we put \[\partial V(p)=V(z,\mu)V(z+e_\mu,\nu)V(z+e_\nu,\mu)^{-1}V(z,\nu)^{-1}.\] The letters $V'$ and $\mathsf V'$ denote fields on the coarse lattice, and $\mathcal V'$ is the set of those with values in $G$.

The integer $L$ is odd, and $\mathsf h_L=\frac12(L-1)$. The integer $N$ is even; it is the coarse side of the present step, that is, the number $N'=NL^{K-k-1}$ of \S\ref{Y-A-1}, in which $N$ denotes the side of the last torus. We set $\mathbb T=\mathbb Z^3/NL\mathbb Z^3$ and $\mathbb T'=L\mathbb Z^3/NL\mathbb Z^3$. Fine bonds and fine plaquettes are written $b$ and $p$; coarse bonds, coarse plaquettes and coarse cubes are written $c$, $P$ and $q'$. Each cell $i$ has a base site $\pi(i)$, namely its least corner, and $y(i)$ is the coarse site whose block $B(y)=\{x:|x-y|_\infty\le\mathsf h_L\}$ contains $\pi(i)$. The distance of two cells is ${\rm dist}(i,j)=L^{-1}|y(i)-y(j)|_\infty$, taken periodically. For $x\in B(y)$, the comb $\Gamma(y,x)$ is the path from $y$ to $x$ that runs first along $e_1$, then along $e_2$, then along $e_3$; the set of the bonds of all combs is $\mathcal T$.

\emph{The set $I^\circ$.} We denote by $I^\circ$ the set of the fine bonds off $\mathcal T$, that is, of the fine bonds that lie on no comb; at most $3L^3$ bonds of $I^\circ$ are based in a block.

We write $\mathfrak L_{\mathfrak g}$ for the space of the $\mathfrak g$-valued functions on $I^\circ$. A coarse bond $c=(y,y')$ of direction $\mu$ has $c_-=y$ and $c_+=y'=y+Le_\mu$, and we set $W_c=B(y)\cup B(y')$. For $x\in B(y)$ we put $x'=x+Le_\mu$, which is a fine site. We write $[x,x']$ for the straight path from $x$ to $x'$, and $A([x,x'])$ for the sum of $A$ over its $L$ bonds. A crossing bond of $c$ is a bond $(z,\mu)$ with $z\in B(y)$ and $z+e_\mu\in B(y')$. The coboundary $\mathrm{d}$ acts by $(\mathrm{d}\eta)_p=\sum_b\mathrm{d}(p,b)\eta_b$, and the kernel $Q_1$ by \[(Q_1A)(c)=L^{-3}\sum_{x\in B(c_-)}A([x,x']).\] The edge lift $E$ is given by $(EX')_p=X'_P$ when $p$ is one of the $L$ edge plaquettes under $P$, and by $0$ at every other plaquette. The edge plaquettes under $P=(y;\mu\nu)$ are the plaquettes $p=(z;\mu\nu)$ with $z_\mu=y_\mu+\mathsf h_L$, $z_\nu=y_\nu+\mathsf h_L$ and $|z_\lambda-y_\lambda|\le\mathsf h_L$ for the direction $\lambda$ normal to $P$, the one other than $\mu$ and $\nu$. They are based in $B(y_P)$, where $y_P=\pi(P)$.

For coarse sites $y$ and $y''$ at sup distance less than half the side of the torus, $\Gamma(y,y'')$ is the coordinate path from $y$ to $y''$, along $e_1$, then $e_2$, then $e_3$, each time in the shorter direction, and the transport between them is $T(y,y''):=\operatorname{Ad}_{V'(\Gamma(y,y''))}$. A scalar kernel $K$ is dressed with transports by $K^V(i,j):=K(i,j)T(y(i),y(j))$ and $(K^V)^*(j,i):=K(i,j)T(y(i),y(j))^{-1}$. The rows act on a $\mathfrak g$-valued function $F$ on the fine plaquettes by $(\ell^VF)_y=\sum_p\ell_y(p)\otimes T(y(y),y(p))F_p$. In particular,
\[
 (\mathrm{d}^V\eta)_p=\sum_b \mathrm{d}(p,b)T(y(p),y(b))\eta_b,\qquad
 (Q_1^V\eta)_c=\sum_b Q_1(c,b)T(c_-,y(b))\eta_b.
\]
The displayed adjoints are transposes for the bilinear form on complex
cochains, with inverse transports. With $\mathsf B^V=1+\ell^{V*}\ell^V$
we define the form and the kernel \[q^V_\ell(\eta)=\langle \mathrm{d}^V\eta,\mathsf B^V \mathrm{d}^V\eta\rangle+a\langle Q_1^V\eta,Q_1^V\eta\rangle,\qquad\mathsf g^V=\Pi_{\mathfrak L}\mathrm{d}^{V*}\mathsf B^VE ,\] where $a>0$ is the parameter of the form and $\Pi_{\mathfrak L}$ is the restriction to $I^\circ$.

\emph{Norms.} The operator norm of a $2\times2$ matrix is written $|\cdot|$. On $\mathfrak g_{\mathbb C}$, the traceless complex $2\times2$ matrices, we use $\|Y\|^2=\frac12\operatorname{tr}Y^*Y$ and $Y\!\cdot\!Y'=\frac12\operatorname{tr}YY'$; then $\|Y\|\le|Y|\le\sqrt2\|Y\|$ and $\|AYB\|\le|A|\|Y\||B|$. For $g\in\operatorname{SL}(2,\mathbb C)$ we have $|g^{-1}|=|g|\ge1$ and $\|\operatorname{Ad}_g\|\le|g|^2$. In the real basis $\sigma_1,\sigma_2,\sigma_3$ of $\mathfrak g$, the form $Y\!\cdot\!Y'$ reads $\sum y_iy_i'$ and the square of the norm reads $\sum|y_i|^2$. The transpose ${}^{\mathsf T}$ is taken with respect to $Y\!\cdot\!Y'$; it preserves operator norms, and $\operatorname{Ad}_g^{\mathsf T}=\operatorname{Ad}_g^{-1}$. The blocks of a matrix on $\ell^2(I^\circ;\mathfrak g_{\mathbb C})$ are measured in the operator norm for $\|\cdot\|$, the norm $\|\cdot\|$ of the matrix itself is its operator norm on $\ell^2$, and $\|A\|_\infty:=\sup_b\sum_{b'}\|A(b,b')\|$.

For matrices $Y_k$ with $S=\sum|Y_k|\le1$ we have \begin{equation}\tag{E}|e^{iY_1}\cdots e^{iY_n}-\mathbf 1-i\sum Y_k|\le e^S-1-S\le0.72S^2.\end{equation}

\begin{proof}[Proof of \textup{(E)}] We expand every factor in its power series and multiply out: the product $e^{iY_1}\cdots e^{iY_n}$ becomes an absolutely convergent sum, over all $(m_1,\dots,m_n)$ with integers $m_k\ge0$, of the products $\prod_k(iY_k)^{m_k}/m_k!$, whose factors keep their order. The terms with $\sum m_k=0$ add up to $\mathbf 1$, and those with $\sum m_k=1$ add up to $i\sum Y_k$. By the triangle inequality and the submultiplicativity of the operator norm, the norm of the remaining terms is therefore at most the sum, over the $(m_k)$ with $\sum m_k\ge2$, of the numbers $\prod_k|Y_k|^{m_k}/m_k!$. This sum equals $e^S-1-S$: summed over all $(m_k)$, these numbers give $\prod_ke^{|Y_k|}=e^S$, and those with $\sum m_k\le1$ give $1+S$. The first inequality of \textup{(E)} thus holds for every $S\ge0$; this is its form in item C.II 1.1, and it is applied with $S\le1.236$ in \Cref{Y-C-L67}\,(d). For the second inequality we observe that $(e^S-1-S)/S^2=\sum_{j\ge2}S^{j-2}/j!$ increases with $S$ and takes the value $e-2<0.7183$ at $S=1$; hence this quotient is at most $0.72$ whenever $S\le1$.\end{proof}

We also write \[\mathsf g(\lambda)=\sum_{n\in\mathbb Z^3}e^{-\lambda|n|_\infty},\qquad\mathsf g_2(\lambda)=\sum_n(1+|n|_\infty)^2e^{-\lambda|n|_\infty}.\]

\emph{Rows.} The rows form a family $\ell=(\ell_y)$, in which $y$ runs over the sites of $\mathbb T$ and each $\ell_y(p)$ is a vector of a Euclidean space $\mathsf H_y$ with inner product $\langle\cdot,\cdot\rangle$. They satisfy $|\ell_y(p)|\le Ke^{-\lambda|y-p|}$, where $|y-p|=|y-\pi(p)|_\infty$ is measured in fine units, $K\ge0$ and $0<\lambda\le1$; and their range is at most $\varrho$, with $\frac\varrho L+2<\frac N2$. It follows that $\|\ell\|_{\rm r}\le3K\mathsf g(\lambda)$ and $\|\ell\|_{\rm c}\le K\mathsf g(\lambda)$, where $\|\ell\|_{\rm r}=\sup_y\sum_p|\ell_y(p)|$ and $\|\ell\|_{\rm c}=\sup_p\sum_y|\ell_y(p)|$; the factor $3$ comes from the three plaquettes based at each site.

\emph{Condition \textup{(u1$^\flat$)}.} A field $\mathsf V'$ satisfies \textup{(u1$^\flat$)} if it is an $\operatorname{SL}(2,\mathbb C)$ field with $\mathsf V'(\bar c)=\mathsf V'(c)^{-1}$ and $|\mathsf V'(c)|\le1+\vartheta_0^\flat$ at every bond $c$. A field that is given only on a set of bonds is extended by $\mathbf 1$ to the remaining bonds. At a field $\mathsf V'$ of this kind, $T$, $\mathrm{d}^{\mathsf V}$, $\ell^{\mathsf V}$, $q^{\mathsf V}_\ell$ and $\mathsf g^{\mathsf V}$ are given by the same formulas as above, with ${}^*$ defined as for $(K^V)^*$; consequently $q^{\mathsf V}_\ell$ is symmetric for $Y\!\cdot\!Y'$ rather than hermitian.

\subsubsection*{C.I 1.2. Numbers}

We set $\gamma=(73L^4)^{-1}$ and \[\FitDisplay{\dimexpr\linewidth-1pt\relax}{b=16\bigl(1+3K^2\mathsf g(\lambda)^2\bigr)+\tfrac aL,\quad A^\flat=2e^2(16+a)+384e^8K^2\mathsf g_2(\tfrac\lambda2)^2,\quad\theta_1=\min\{1,\tfrac{\lambda L}{12}\},\quad\theta^\flat=\min\{\theta_1,\tfrac\gamma{2A^\flat}\},\quad\theta_{\mathbb C}=\tfrac16\theta^\flat,\quad\vartheta_0^\flat=\tfrac1{24}\theta^\flat,}\] \[\FitDisplay{\dimexpr\linewidth-1pt\relax}{\varkappa=1-\tfrac\gamma{2b},\qquad C_\varphi^\flat=\tfrac{2b^{1/2}}\gamma\,\frac{3L^3\mathsf g(\theta_{\mathbb C})}{1-e^{-\theta_{\mathbb C}}},\qquad c_k=4^{-k}\tbinom{2k}k .}\] These numbers satisfy $\vartheta_0^\flat\le\frac\gamma{48A^\flat}\le10^{-7}$; the second of the two bounds rests on $L\ge3$, a hypothesis of every statement, which gives $\gamma\le\frac1{5913}$ and $A^\flat\ge32e^2$.

In the chain these numbers are evaluated at $(K,\lambda)=(K_*,\lambda_*)$ at every level, the level of the empty rows included (item~C.II~1.3; item~3, \textup{(O6)}, of \S\ref{Y-G0-1}); they thus depend on the letters of the bound and are independent of the rows.

We further put $\theta_2=\frac12\theta_1$ and $\Lambda=2+A^\flat/b$, and, for $R\in\mathbb N$, \[r_w:=\lfloor R/2\rfloor,\quad\varrho_L:=\lceil\varrho/L\rceil,\quad r_q:=2\varrho_L+2,\quad\mathsf R_{\mathsf E}:=2r_w+3\varrho_L+3,\]\[d_R(\theta):=9L^3\,b\,\gamma^{-3/2}(R+3)^3\,\mathsf g(\theta)\,e^{-\theta R/2},\quad d_R:=d_R(\theta_{\mathbb C}),\] \[\mathsf G_0:=(\tfrac2\gamma)^{1/2}3L^3\mathsf g(\theta_{\mathbb C})\ (\le C_\varphi^\flat),\qquad\mathsf K_L:=\frac{153\,A^\flat}{\gamma^{3/2}\theta^\flat}.\]

In what follows we use the identities $\sum_kc_k\kappa^k=(1-\kappa)^{-1/2}$ and $\sum_kkc_k\kappa^{k-1}=\frac12(1-\kappa)^{-3/2}$; at $\kappa=\varkappa$ they give $b^{-1/2}\sum_kc_k\varkappa^k=(\frac2\gamma)^{1/2}$ and $b^{-1/2}\sum_kkc_k\varkappa^{k-1}=\sqrt2\,b\gamma^{-3/2}$. We also use $\sum_k|\binom{1/2}k|=2$ and the inequality $(1+t)e^{-\delta t}\le\frac1\delta$ for $t\ge0$ and $0<\delta\le1$.

\subsubsection*{C.I 1.3. Steps, sub-forms}

The matrix of the form is a sum over steps, \[q^V_\ell(b,b')=\sum_{\mathsf s:b\to b'}q_{\mathsf s}\mathsf T_{\mathsf s},\] in which a step $\mathsf s$ from $b$ to $b'$ is of one of the three kinds \textup{(P)}, \textup{(Q)}, \textup{(R)}, with the following weight $q_{\mathsf s}$ and transport $\mathsf T_{\mathsf s}$.

\begin{enumerate}\item[\textup{(P)}] A step $\mathsf s=(p)$ with $p\ni b,b'$ has $q_{\mathsf s}=\mathrm{d}(p,b)\mathrm{d}(p,b')$ and $\mathsf T_{\mathsf s}=T(y(p),y(b))^{-1}T(y(p),y(b'))$.

\item[\textup{(Q)}] A step $\mathsf s=(c)$ has $q_{\mathsf s}=aQ_1(c,b)Q_1(c,b')$ and $\mathsf T_{\mathsf s}=T(c_-,y(b))^{-1}T(c_-,y(b'))$.

\item[\textup{(R)}] A step $\mathsf s=(y,p,p')$ with $p\ni b$ and $p'\ni b'$ has $q_{\mathsf s}=\mathrm{d}(p,b)\mathrm{d}(p',b')\langle\ell_y(p),\ell_y(p')\rangle$ and \[\mathsf T_{\mathsf s}=T(y(p),y(b))^{-1}T(y(y),y(p))^{-1}T(y(y),y(p'))T(y(p'),y(b')).\]\end{enumerate}

For each kind, $\mathsf T_{\mathsf s}=\operatorname{Ad}_{V'(\pi_{\mathsf s})}$, where $\pi_{\mathsf s}$ is the coarse path from $y(b)$ to $y(b')$ that the product of transports traces, taken as it stands. We write $|\mathsf s|$ for the number of its bonds, backtracks included, and $\bar{\mathsf s}:b'\to b$ for the reversed step. The weight $q_{\mathsf s}$ is independent of the field. For the steps $(y,p,p')$ of kind \textup{(R)} only the sites $y$ with $|y-p|\le\varrho$ and $|y-p'|\le\varrho$ are admitted: for every other $y$ the range of the rows gives $q_{\mathsf s}=0$, and so the transport $T(y(y),y(p))$ is required only for admitted $y$. If $\mathsf H_y=\{0\}$, every step $(y,p,p')$ of kind \textup{(R)} has $q_{\mathsf s}=0$. Every sum over steps below carries the factor $q_{\mathsf s}$; hence the steps with $q_{\mathsf s}=0$ change no sum and no bound.

\emph{Sub-forms.} Let $\mathsf S$ be a set of coarse sites, and let $I_{\mathsf S}:=\{b\in I^\circ:y(b)\in\mathsf S\}$. A step $\mathsf s:b\to b'$ lies \emph{in} $\mathsf S$ if $b,b'\in I_{\mathsf S}$ and if, for a step of kind \textup{(R)}, all sites of $\Gamma(y(y),y(p))$ and of $\Gamma(y(y),y(p'))$ belong to $\mathsf S$. For $\zeta\in\mathbb C$ and $Y\subset I_{\mathsf S}$ we define \[q_{\mathsf S}(\zeta):=\sum_{\mathsf s\text{ in }\mathsf S}\zeta^{|\mathsf s|}q_{\mathsf s}\mathsf T_{\mathsf s},\qquad q_{\mathsf S}:=q_{\mathsf S}(1),\qquad K_{\mathsf S,Y}(\zeta):=\mathbf 1_Y-b^{-1}\mathbf 1_Yq_{\mathsf S}(\zeta)\mathbf 1_Y\ \text{ on }\ell^2(Y;\mathfrak g_{\mathbb C}),\] and wherever a product with another matrix is taken, $K_{\mathsf S,Y}(\zeta)$ is extended by $0$ to $\ell^2(I^\circ)$. We also write $K_{\mathsf S,Y}:=K_{\mathsf S,Y}(1)$ and $K_{\mathsf S}:=K_{\mathsf S,I_{\mathsf S}}$. When $\mathsf S$ is the set of all sites, the index $\mathsf S$ is omitted: then $q=q^V_\ell$, $K=1-q/b$, and $K_Y$ stands for $K_{\mathsf S,Y}$. Finally, for a polynomial $A(\zeta)$, $[A(\zeta)]_j$ denotes its coefficient of $\zeta^j$.

\subsubsection*{C.I 2. The form; one weighted bound}

\begin{lemma}[The form and its sub-forms]\label{Y-C-CIL64} In the setting of item~C.I~1.1, with rows as there, let $L\ge3$ be odd and let $aL\ge18$. Then the following hold.

\begin{enumerate}

\item[\textup{(a)}] At every $V'\in\mathcal V'$, for every $\mathsf S$ and every $\eta\in\mathfrak L_{\mathfrak g}$ supported in $I_{\mathsf S}$, we have $\gamma|\eta|^2\le\langle\eta,q_{\mathsf S}\eta\rangle\le b|\eta|^2$. In consequence, $0\le K_{\mathsf S,Y}\le1-\frac\gamma b$ for all $Y\subset I_{\mathsf S}$.

\item[\textup{(b)}] Let $0\le\beta\le\theta_1$, and write $\sum^{(b)}:=\sup_b\sum_{b'}\sum_{\mathsf s:b\to b'}$. Then we have the first bound $\sum^{(b)}|q_{\mathsf s}|\le b$, the second bound $\sum^{(b)}|q_{\mathsf s}|(e^{\beta|\mathsf s|}-1)\le A^\flat\beta$, and the third bound $\sum^{(b)}|q_{\mathsf s}|\,|\mathsf s|(1+|\mathsf s|)e^{\beta|\mathsf s|}\le9A^\flat$. These three bounds also hold for the sums over $b$ at fixed $b'$.

\item[\textup{(c)}] For every step $\mathsf s:b\to b'$ we have ${\rm dist}(b,b')\le\min\{|\mathsf s|,r_q\}$, and each site of $\pi_{\mathsf s}$ lies within $\varrho_L+1$ of $y(b)$ or of $y(b')$.

\item[\textup{(d)}] For every step $\mathsf s$ and the reversed step $\bar{\mathsf s}$ we have $\mathsf T_{\bar{\mathsf s}}=\mathsf T_{\mathsf s}^{-1}=\mathsf T_{\mathsf s}^{\mathsf T}$, $q_{\bar{\mathsf s}}=q_{\mathsf s}$ and $|\bar{\mathsf s}|=|\mathsf s|$. Moreover, for each map $u$ from the coarse sites to $\operatorname{SL}(2,\mathbb C)$ and each step $\mathsf s:b\to b'$, $\mathsf T_{\mathsf s}[\mathsf V'^u]=\operatorname{Ad}_{u(y(b))}\mathsf T_{\mathsf s}[\mathsf V']\operatorname{Ad}_{u(y(b'))}^{-1}$.

\item[\textup{(e)}] The entries of $\mathsf g^V$ are $\mathsf g^V(b,P)=\sum_{\mathsf t:b\to P}g_{\mathsf t}\operatorname{Ad}_{V'(\pi_{\mathsf t})}$, the sum running over steps of two kinds:

\begin{enumerate}

\item[\textup{(E)}] $\mathsf t=(p)$, where $p\ni b$ is an edge plaquette under $P$, with $g_{\mathsf t}=\mathrm{d}(p,b)$;

\item[\textup{(RE)}] $\mathsf t=(y,p,p')$, where $p\ni b$ and $p'$ is an edge plaquette under $P$, with $g_{\mathsf t}=\mathrm{d}(p,b)\langle\ell_y(p),\ell_y(p')\rangle$; as in item~C.I~1.3, the site $y$ is restricted to $|y-p|\le\varrho$ and $|y-p'|\le\varrho$, and $g_{\mathsf t}=0$ for every other $y$.

\end{enumerate}

For each of these steps, $\pi_{\mathsf t}$ is a coarse path from $y(b)$ to $y_P$, all of whose sites lie within $2\varrho_L+1$ of $y(b)$. We write $|\mathsf t|$ for the number of bonds of $\pi_{\mathsf t}$, backtracks included, just as $|\mathsf s|$ is defined in item~C.I~1.3. Moreover $\sup_b\sum_P\sum_{\mathsf t}|g_{\mathsf t}|\le\frac b4$ and $\sup_b\sum_P\sum_{\mathsf t}|g_{\mathsf t}|e^{2\vartheta_0^\flat|\mathsf t|}\le b$, and for $\beta\le\theta_1$ also $\sup_b\sum_P\sum_{\mathsf t}|g_{\mathsf t}||\mathsf t|(1+|\mathsf t|)e^{\beta|\mathsf t|}\le9A^\flat$.

\end{enumerate}

\end{lemma}

\begin{proof} \emph{Step (a): the lower bound for the form without rows.} Let $\eta\in\mathfrak L_{\mathfrak g}$, and put $f=\mathrm{d}^V\eta$. We define the form without rows, which is the form $q^V_\ell$ at the empty rows, by $q^V_\emptyset(\eta):=|\mathrm{d}^V\eta|^2+a|Q_1^V\eta|^2$, that is, $q^V_\emptyset(\eta)=|f|^2+a|Q_1^V\eta|^2$; by the display that defines $q^V_\ell$ in item~C.I~1.1, $q^V_\ell(\eta)=q^V_\emptyset(\eta)+|\ell^V\mathrm{d}^V\eta|^2$.

In each block all transports are the identity, so
\eqref{eq:comb-local-poincare} applies. For a double block $W_c$ of
$c=(y,y')$, keep the values of $\eta$ based in $B(y)$ and apply
$\operatorname{Ad}_{V'(c)}$ to those based in $B(y')$. On a plaquette
of $W_c$ the ordinary coboundary of the transformed cochain is
$\mathrm{d}^V\eta$ expressed in the frame at $y$. Across the seam this follows
from $T(y,y')=\operatorname{Ad}_{V'(c)}$, and inside either block it
follows because all its bonds are in the same frame. The transformed
cochain vanishes on the combs and its average on $c$ is $(Q_1^V\eta)_c$.
The transport is an isometry, so \eqref{eq:crossing-local-poincare}
gives the same inequality with $f=\mathrm{d}^V\eta$ and $Q_1^V\eta$.
Summing, as in \Cref{Y-A-L2}, gives
$|\eta|^2\le73L^4|f|^2+6L^2|Q_1^V\eta|^2\le73L^4q^V_\emptyset(\eta)$.
The last inequality uses $aL\ge18$.

\emph{Step (a): the sub-forms.} We put $\ell^{[\mathsf S]}_y(p):=\ell_y(p)1\{$all sites of $\Gamma(y(y),y(p))$ lie in $\mathsf S\}$. Then $q_{\mathsf S}=\mathbf 1_{I_{\mathsf S}}q^V_\emptyset\mathbf 1_{I_{\mathsf S}}+(\ell^{[\mathsf S]V}\mathrm{d}^V\mathbf 1_{I_{\mathsf S}})^*(\ell^{[\mathsf S]V}\mathrm{d}^V\mathbf 1_{I_{\mathsf S}})$: when the right side is multiplied out, its first term yields the steps \textup{(P)} and \textup{(Q)} between bonds of $I_{\mathsf S}$, and its second term the steps \textup{(R)} in $\mathsf S$. Thus $q_{\mathsf S}$ is the sum of a compression of $q^V_\emptyset$ and of a square, where $q^V_\emptyset\ge\gamma$ by the local estimates just proved; this gives the lower bound of \textup{(a)}.

\emph{Step (a): the upper bound.} The upper bound of \textup{(a)} follows from Schur's test. Indeed, every transport has $\|\mathsf T_{\mathsf s}\|=1$, and the sums over the steps that make up $q_{\mathsf S}$ are sub-sums of the sums in the first bound of \textup{(b)}, for rows and for columns.

\emph{Step (b): the three bounds.} For the first bound, the three kinds of steps contribute as follows.

\begin{enumerate}

\item[\textup{(P)}] There are four plaquettes $p\ni b$, each with four bonds $b'$, which gives $16$; as $\Gamma(y(p),y(b))$ is empty or consists of one bond, $|\mathsf s|\le2$.

\item[\textup{(Q)}] As $\sum_cQ_1(c,b)=L^{-2}$ and $\sum_{b'}Q_1(c,b')=L$, the contribution is $\frac aL\le a$; as $y(b)\in\{c_-,c_+\}$, again $|\mathsf s|\le2$.

\item[\textup{(R)}] The contribution is $\le4\cdot4\,\|\ell\|_{\rm c}\|\ell\|_{\rm r}\le48K^2\mathsf g(\lambda)^2$.

\end{enumerate}

These three contributions add up to at most $b$, which proves the first bound. For the other two bounds we consider a step of kind \textup{(R)} with the site $y$ and the plaquettes $p$ and $p'$, and we put $n_1=|y-p|$ and $n_2=|y-p'|$. Then $L^{-1}|y(y)-y(p)|_\infty<\frac{n_1}L+1$, and so $\Gamma(y(y),y(p))$ has at most $3(\frac{n_1}L+1)$ bonds; the path from $y(y)$ to $y(p')$ obeys the same bound with $n_2$ in place of $n_1$. Consequently $m:=|\mathsf s|\le8+\frac3L(n_1+n_2)\le8(1+n_1)(1+n_2)$, hence $1+m\le9(1+n_1)(1+n_2)$, and also $e^{\beta m}\le e^8e^{\frac\lambda4(n_1+n_2)}$. In addition we use $e^{\beta m}-1\le\beta me^{\beta m}$. The contributions to the second and to the third sum are then bounded as follows.

\begin{enumerate}

\item[\textup{(P), (Q)}] Since $|\mathsf s|\le2$, these two kinds contribute at most $(16+a)2\beta e^2$ to the second sum and at most $(16+a)6e^2$ to the third.

\item[\textup{(R)}] As there are three plaquettes to each site, the contribution to the second sum is at most $16\cdot8e^8K^2\beta\cdot3[\sum_u(1+|u|)e^{-\frac34\lambda|u|}]^2\le384e^8K^2\mathsf g_2(\frac\lambda2)^2\beta$, where $u$ runs over $\mathbb Z^3$ and $|u|=|u|_\infty$, here and in the next sum over $u$, and the contribution to the third sum is at most $16\cdot72e^8K^2\cdot3[\sum_u(1+|u|)^2e^{-\frac34\lambda|u|}]^2\le3456e^8K^2\mathsf g_2(\frac\lambda2)^2$.

\end{enumerate}

The contributions to the second sum add up to at most $(16+a)2\beta e^2+384e^8K^2\mathsf g_2(\frac\lambda2)^2\beta=A^\flat\beta$, and those to the third sum add up to at most $(16+a)6e^2+3456e^8K^2\mathsf g_2(\frac\lambda2)^2\le9A^\flat$; in both cases we use the definition $A^\flat=2e^2(16+a)+384e^8K^2\mathsf g_2(\tfrac\lambda2)^2$ in item~C.I~1.2. This proves the second and the third bound of \textup{(b)}. For the columns, that is, for the sums over $b$ at fixed $b'$, the same bounds follow from \textup{(d)}.

\emph{Step (c): distances and paths.} The path $\pi_{\mathsf s}$ joins $y(b)$ to $y(b')$. For the plaquette $p\ni b$ of a step, $y(p)$ is within $1$ of $y(b)$. For a step through a site $y$ we have $L^{-1}|y(y)-y(p)|_\infty<\frac\varrho L+1$, and as this number is an integer, it is at most $\varrho_L$. Finally, a coordinate path does not leave the box spanned by its two ends; this proves \textup{(c)}.

\emph{Step (d): reversal and gauge transformations.} The reversed step runs along the reversed path, and $\operatorname{Ad}_g^{\mathsf T}=\operatorname{Ad}_g^{-1}$; this gives the first part of \textup{(d)}. The second part follows from $\mathsf V'^u(\pi)=u(\pi_-)\mathsf V'(\pi)u(\pi_+)^{-1}$, where $\pi_-$ and $\pi_+$ are the initial and the final site of the path $\pi$.

\emph{Step (e): the matrix $\mathsf g^V$.} The two kinds of steps arise on multiplying out $\Pi_{\mathfrak L}\mathrm{d}^{V*}(1+\ell^{V*}\ell^V)E$; here an edge plaquette under $P$ is based in $B(y_P)$. The sums in \textup{(e)} are sub-sums of those in \textup{(b)}, without the last factor $4$, because no $b'\in\partial p'$ occurs, and with $|\mathsf t|\le7+\frac3L(n_1+n_2)$; in particular, the third bound of \textup{(e)} follows from the third bound of \textup{(b)}. The first sum in \textup{(e)} is therefore at most $4+4\cdot3K^2\mathsf g(\lambda)^2\le\frac b4$, and, by the second bound of \textup{(b)}, the second sum is at most $\frac b4+2A^\flat\vartheta_0^\flat\le b$.

\end{proof}

\begin{lemma}[One weighted bound for all sub-forms]\label{Y-C-CIL65} Under the hypotheses of \Cref{Y-C-CIL64}, that is, with rows as in item~C.I~1.1, with $L\ge3$ odd and with $aL\ge18$, consider a family of retained steps on a set $Y\subset I^\circ$, closed
under reversal. Write $q_{\mathrm{ret}}(\zeta)$ for its step sum as in
item~C.I~1.3 and $K_{\mathrm{ret}}(\zeta)=1-b^{-1}q_{\mathrm{ret}}(\zeta)$.
When $\zeta$ is omitted, it is set to $1$. Assume that
$q_{\mathrm{ret}}\ge\gamma$ at every $G$-valued field. The families
defining $K_{\mathsf S,Y}$ satisfy this assumption by
\Cref{Y-C-CIL64}(a). Let $w$ be real on $Y$ with
$|w(b_1)-w(b_2)|\le\theta\,{\rm dist}(b_1,b_2)$. Then the following
bounds hold for this family.

\begin{enumerate}

\item[\textup{(r)}] At every $V'\in\mathcal V'$, if $\theta\le\theta^\flat$, then $\|e^wK_{\mathrm{ret}}e^{-w}\|\le\varkappa$.

\item[\textup{(c)}] At every $\mathsf V'$ satisfying \textup{(u1$^\flat$)}, if $\theta\le\theta_{\mathbb C}$, then $\|e^wK_{\mathrm{ret}}e^{-w}\|\le1-\frac{5\gamma}{8b}\le\varkappa$.

\item[\textup{(z)}] At every $\mathsf V'$ satisfying \textup{(u1$^\flat$)}, $\|K_{\mathrm{ret}}(\zeta)\|\le\Lambda$ for $|\zeta|\le e^{\theta_2}$; consequently $\|[K_{\mathrm{ret}}(\zeta)^k]_j\|\le\Lambda^ke^{-\theta_2j}$.

\end{enumerate}

\end{lemma}

\begin{proof} At a real field the lower bound on the form is the
assumption. Its upper bound by $b$ follows from Schur's test and the
first bound of \Cref{Y-C-CIL64}(b), since the transports are
orthogonal. Thus $0\le K_{\mathrm{ret}}\le1-\gamma/b$.

\emph{Step (r): real fields.} The entry $(b_1,b_2)$ of $e^wK_{\mathrm{ret}}e^{-w}-K_{\mathrm{ret}}$ equals \[\FitDisplay{\dimexpr\linewidth-1pt\relax}{-b^{-1}\sum_{\mathsf s}q_{\mathsf s}\mathsf T_{\mathsf s}(e^{w(b_1)-w(b_2)}-1)}\]. For real $t$ one has $|e^t-1|\le e^{|t|}-1$, and $|w(b_1)-w(b_2)|\le\theta|\mathsf s|$ by \Cref{Y-C-CIL64}\,(c). Schur's test and \Cref{Y-C-CIL64}\,(b) thus bound the norm of the difference by $b^{-1}A^\flat\theta^\flat\le\frac\gamma{2b}$. Combined with $0\le K_{\mathrm{ret}}\le1-\gamma/b$ at a real field, this proves \textup{(r)}.

\emph{Step (c): fields satisfying \textup{(u1$^\flat$)}.} We decompose every $\mathsf V'(c)$ in polar form, $\mathsf V'(c)=\mathsf p_ck_c$, where $k_c\in G$ and $\mathsf p_c>0$ has the eigenvalues $s^{\pm1}$ with $1\le s\le1+\vartheta_0^\flat$; thus $|\mathsf p_c^{\pm1}-\mathbf 1|\le\vartheta_0^\flat$. Then $V_0:=(k_c)$ is a field with values in $G$. Consider a path of $m$ bonds, and let $g=\mathsf V'(\pi)$ and $g_0=V_0(\pi)$. Every factor of $g$ is the sum of a unitary matrix and a matrix of norm $\le\vartheta_0^\flat$, so that $|g|\le(1+\vartheta_0^\flat)^m$ and $|g-g_0|\le(1+\vartheta_0^\flat)^m-1$. From the identity $gYg^{-1}-g_0Yg_0^{-1}=(g-g_0)Yg^{-1}+g_0Yg^{-1}(g_0-g)g_0^{-1}$ we deduce $\|\operatorname{Ad}_g-\operatorname{Ad}_{g_0}\|\le2|g||g-g_0|\le2(e^{2\vartheta_0^\flat m}-1)$. The entries of $e^wq_{\mathrm{ret}}^{\mathsf V}e^{-w}-q_{\mathrm{ret}}^{V_0}$ are therefore, in norm, $\le\sum|q_{\mathsf s}|[(e^{\theta m}-1)+2e^{\theta m}(e^{2\vartheta_0^\flat m}-1)]\le3\sum|q_{\mathsf s}|(e^{(\theta+2\vartheta_0^\flat)m}-1)$, by the elementary inequality $2e^{A+B}-e^A-1\le3e^{A+B}-3$ for $A,B\ge0$. Here $\theta_{\mathbb C}+2\vartheta_0^\flat=\frac14\theta^\flat$. The row sums and the column sums share one majorant, and Schur's test, with the second bound of \Cref{Y-C-CIL64}\,(b), bounds the norm of this difference by $\frac34A^\flat\theta^\flat\le\frac38\gamma$. Using $0\le K_{\mathrm{ret}}[V_0]\le1-\gamma/b$, we conclude that $\|e^wK^{\mathsf V}_{\mathrm{ret}}e^{-w}\|\le1-\frac\gamma b+\frac{3\gamma}{8b}\le\varkappa$. The coercivity assumption is used at the polar field $V_0$.

\emph{Step (z): the variable $\zeta$.} We have $\|\mathsf T_{\mathsf s}\|\le|\mathsf V'(\pi_{\mathsf s})|^2\le e^{2\vartheta_0^\flat|\mathsf s|}$ and $\theta_2+2\vartheta_0^\flat\le\theta_1$; Schur's test and \Cref{Y-C-CIL64}\,(b) give $\|q_{\mathrm{ret}}(\zeta)\|\le b+A^\flat\theta_1$, and from this the first assertion of \textup{(z)} follows. The second assertion is Cauchy's estimate for the polynomial $K(\zeta)^k$ on the circle $|\zeta|=e^{\theta_2}$.

\end{proof}

\subsubsection*{C.I 3. \textup{(Law$_k$)}}

\begin{definition}[Windows, the matrix of the law, the law]\label{Y-C-CIDef4} Let $R\in\mathbb N$. For $b\in I^\circ$, the window of $b$ is $Y_b:=\{b'\in I^\circ:{\rm dist}(b,b')\le r_w\}$; we let $K_b:=K_{Y_b}$, taken as in item~C.I~1.3 with $\mathsf S$ the set of all sites, and we write $q_b=\mathbf 1_{Y_b}q\mathbf 1_{Y_b}=b(\mathbf 1_{Y_b}-K_b)$. The matrix of the law, $G^{\mathsf V}$, and the matrices $C$, $\mathsf E$, $F$ and $\Delta$ are defined by \[G^{\mathsf V}(b,b'):=b^{-1/2}\sum_{k\ge0}c_k\,(K_b)^k(b,b')\ \ (b'\in Y_b),\quad0\text{ for all other }b'\in I^\circ;\]\[C:=GG^{\mathsf T},\quad\mathsf E:=G^{\mathsf T}qG-1,\quad F:=b^{-1/2}\sum_kc_kK^k,\quad\Delta:=F-G .\] Let $x>1$ be the parameter of the law, and let $\xi\sim\mathcal N(0,x^{-1}\mathbf 1)$ on $\mathfrak L_{\mathfrak g}$; this is a product measure, and it depends on no field. The law is the image of $\mathcal N(0,x^{-1}\mathbf 1)$ under the map $\xi\mapsto\varphi:=G^V\xi$. Under the hypotheses of \Cref{Y-C-CITh66}, at every $V'\in\mathcal V'$ the matrix $G^V$ is real, as \Cref{Y-C-CITh66}\,(a) asserts, and the law at such a field is a probability measure on $\mathfrak L_{\mathfrak g}$.

\end{definition}

\begin{theorem}[The law of windows, \textup{(Law$_k$)}]\label{Y-C-CITh66} In the setting of item~C.I~1.1, with rows as there, let $L\ge3$ be odd, let $aL\ge18$ and let $R\in\mathbb N$, and let $G$, $C$, $\mathsf E$, $F$ and $\Delta$ be as in \Cref{Y-C-CIDef4}. Each clause below holds at every $V'\in\mathcal V'$ with $\theta=\theta^\flat$, and at every $\mathsf V'$ with \textup{(u1$^\flat$)} with $\theta=\theta_{\mathbb C}$.

\begin{enumerate}

\item[\textup{(a)}] \emph{Well defined.} The series converge in norm, $F^2q=1$, and $F^{\mathsf T}=F$. At $V'\in\mathcal V'$, moreover, $F=(q^V_\ell)^{-1/2}$, $G^V(b,\cdot)$ is the row $b$ of $q_b^{-1/2}$, and $G^V$ is real.

\item[\textup{(b)}] \emph{Range and dependence.} We have $G(b,b')=0$ unless ${\rm dist}(b,b')\le r_w\le\frac12R$. The matrix $C$ is symmetric, and $C(b,b')=0$ unless ${\rm dist}(b,b')\le2r_w\le R$. The component $(G\xi)_b$ is a function of the values of $\xi$ on $Y_b$. The row $G(b,\cdot)$ depends only on $\operatorname{Ad}\mathsf V'$ on the coarse bonds whose two ends lie within $r_w+\varrho_L+1$ of $y(b)$, and the row $C(b,\cdot)$ only on $\operatorname{Ad}\mathsf V'$ on those whose two ends lie within $3r_w+\varrho_L+1$ of $y(b)$.

\item[\textup{(c)}] \emph{Covariance.} For $G$-valued $u$, $G[\mathsf V'^u]=\operatorname{Ad}_uG[\mathsf V']\operatorname{Ad}_u^{-1}$, where $(\operatorname{Ad}_u\eta)_b=\operatorname{Ad}_{u(y(b))}\eta_b$; the same holds for $C,\mathsf E,F$.

\item[\textup{(c$'$)}] The identities of \textup{(c)} hold also for $\operatorname{SL}(2,\mathbb C)$-valued $u$, provided both $\mathsf V'$ and $\mathsf V'^u$ satisfy \textup{(u1$^\flat$)}. For a single row $b$ it suffices to impose the bound
\textup{(u1$^\flat$)} and the gauge relation on the union of the paths
of the steps between bonds of $Y_b$. These paths are contained in the
region of part~(b).

\item[\textup{(d)}] \emph{Decay in the sup norm.} We have $\|G(b,b')\|,\|F(b,b')\|\le(\frac2\gamma)^{1/2}e^{-\theta\,{\rm dist}(b,b')}$. Consequently $\|G\|_\infty,\|G^{\mathsf T}\|_\infty\le\mathsf G_0$ and $\|(G\xi)_b\|\le\mathsf G_0\sup_{Y_b}\|\xi_{b'}\|\le C_\varphi^\flat|\xi|_\infty$, where $|\xi|_\infty:=\sup_{b'\in I^\circ}\|\xi_{b'}\|$.

\item[\textup{(e)}] \emph{Closeness.} We have $\|\Delta(b,b')\|\le\sqrt2\,b\,\gamma^{-3/2}\,e^{-\theta\max\{r_w+1,\ {\rm dist}(b,b')\}}$; consequently $\|\Delta\|,\|\Delta\|_\infty,\|\Delta^{\mathsf T}\|_\infty\le d_R(\theta)$.

\item[\textup{(f)}] \emph{Inverse.} If $4b^{1/2}d_R(\theta)\le1$, then $G$ is invertible and $\|G^{-1}\|\le4b^{1/2}$; under the same condition, $S:=C^{-1}-q$ is symmetric, $S=-G^{-\mathsf T}\mathsf EG^{-1}$, and $\|S\|\le48\,b^{3/2}d_R(\theta)$. Whatever the value of $d_R(\theta)$, $\|q^{-1}-C\|_\infty\le2\mathsf G_0d_R(\theta)+d_R(\theta)^2$.

\item[\textup{(g)}] \emph{The ratio in white noise.} The matrix $\mathsf E$ is symmetric, and $\mathsf E(b,b')=0$ unless ${\rm dist}(b,b')\le2r_w+r_q\le\mathsf R_{\mathsf E}$; the row $\mathsf E(b,\cdot)$ depends only on $\operatorname{Ad}\mathsf V'$ on the bonds within $\mathsf R_{\mathsf E}$ of $y(b)$. Further, $\mathsf E=-F^{-1}\Delta-\Delta^{\mathsf T}F^{-1}+\Delta^{\mathsf T}q\Delta$, and if $4b^{1/2}d_R(\theta)\le1$, then $\|\mathsf E\|\le5\,b^{1/2}d_R(\theta)$ and $\|\mathsf E\|_\infty\le3L^3(2\mathsf R_{\mathsf E}+1)^3\|\mathsf E\|$.

\item[\textup{(h)}] \emph{Determinant.} We have $(\det G)^2\det q=\det(1+\mathsf E)$. If $4b^{1/2}d_R(\theta)\le1$, then at $V'\in\mathcal V'$ we have $\det G^V>0$ and $\log\det G^V=-\frac12\log\det q^V_\ell+\frac12\operatorname{tr}\log(1+\mathsf E^V)$; the logarithm here is that of the real symmetric matrix $1+\mathsf E^V=(1-\mathsf X)^{\mathsf T}(1-\mathsf X)$, where $\mathsf X:=F^{-1}\Delta$, and the spectrum of this matrix lies in $[\frac14,\frac94]$. If $5b^{1/2}d_R(\theta)<1$, this logarithm is the sum of its power series.

\item[\textup{(i)}] \emph{Holomorphy.} On the set of the fields satisfying \textup{(u1$^\flat$)}, each entry of $G,C,\mathsf E,F$ is the uniform limit of polynomials in the entries of those matrices $\mathsf V'(c)$ on which the entry depends; hence it is continuous on this set and holomorphic in its interior. This holds independently of any bound on products along paths and of any bound on the curvature.

\item[\textup{(j)}] \emph{Lipschitz bound in the field, for differences growing linearly with the distance.} Let $\mathsf V_1,\mathsf V_2$ be fields satisfying \textup{(u1$^\flat$)}, let $b_0\in I^\circ$ and $\epsilon_1\ge0$, and let $\|\operatorname{Ad}^{\pm1}_{\mathsf V_1(c)}-\operatorname{Ad}^{\pm1}_{\mathsf V_2(c)}\|\le\epsilon_1\bigl(1+L^{-1}|w-y(b_0)|_\infty\bigr)$ for each coarse bond $c=(w,\mu)$ on a transport path of a step
between bonds of $Y_{b_0}$. Then for all $b'$ \[\bigl\|(G^{\mathsf V_1}-G^{\mathsf V_2})(b_0,b')\bigr\|\le\mathsf K_L\,\epsilon_1\,e^{-\frac12\theta_{\mathbb C}{\rm dist}(b_0,b')} .\] The constant $\mathsf K_L$ depends neither on $R$ nor on $\varrho$ nor on $N$.

\item[\textup{(k)}] \emph{Numbers.} Let $x>1$ and $R=\lceil c_R\log x\rceil$, where $c_R\theta_{\mathbb C}\ge128$, and let the floor condition \begin{equation}\tag{F$_{\rm law}$}432\,L^3b^{5/2}\gamma^{-3/2}\mathsf g(\theta_{\mathbb C})(R+3)^3\le x^4\end{equation} be satisfied. Then $4b^{1/2}d_R\le1$, $\|\mathsf E\|\le x^{-60}$ and $\|S\|\le x^{-60}$.

\end{enumerate}

\end{theorem}

\begin{proof} \emph{Preliminary step: the bound \textup{(W)}.} Let $w$ be a real function on $I^\circ$ with $|w(b_1)-w(b_2)|\le\theta\,{\rm dist}(b_1,b_2)$, and let $Y$ be arbitrary. By \Cref{Y-C-CIL65}\,(r) at the fields of $\mathcal V'$, and by \Cref{Y-C-CIL65}\,(c) at the fields satisfying \textup{(u1$^\flat$)}, \begin{equation}\tag{W}\|e^wK_Ye^{-w}\|\le\varkappa .\end{equation} For $w=0$ this gives $\|K_Y\|\le\varkappa$.

\emph{Step (a): the series.} The coefficients satisfy $c_k\ge0$ and $\sum_kc_k\varkappa^k<\infty$, so by \textup{(W)} at $w=0$ all the series converge in norm. As power series, $(\sum_kc_k\kappa^k)^2(1-\kappa)=1$, and the series converge absolutely at $\varkappa$; hence $F^2b(1-K)=1$, which is the identity $F^2q=1$. By \Cref{Y-C-CIL64}\,(d), $K^{\mathsf T}=K$, and so $F^{\mathsf T}=F$. At real fields the matrices $K$ and $K_b$ are real and symmetric with $0\le K,K_b\le1-\frac\gamma b$, and the series are the functions $(b(1-\kappa))^{-1/2}$ of these matrices; this gives the remaining assertions of \textup{(a)}.

\emph{Step (b): range and dependence.} The first assertion of \textup{(b)} holds by definition, and so does the assertion on $(G\xi)_b$. The formula $C(b,b')=\sum_{b''}G(b,b'')G(b',b'')^{\mathsf T}$ gives the assertions on $C$. The row $G(b,\cdot)$ is a function of the entries $q(b_1,b_2)$ with $b_i\in Y_b$, and \Cref{Y-C-CIL64}\,(c) bounds the region on which these entries depend. The row $C(b,\cdot)$ is computed from the rows $G(b',\cdot)$ with ${\rm dist}(b,b')\le2r_w$.

\emph{Steps (c) and (c$'$): covariance.} Let $\mathsf U:={\rm diag}(\operatorname{Ad}_{u(y(b))})$. This matrix is bounded and has a bounded inverse, and it commutes with every $\mathbf 1_Y$. By \Cref{Y-C-CIL64}\,(d), $K_Y[\mathsf V'^u]=\mathsf UK_Y[\mathsf V']\mathsf U^{-1}$, and therefore $\sum_kc_kK_Y[\mathsf V'^u]^k=\mathsf U(\sum_kc_kK_Y[\mathsf V']^k)\mathsf U^{-1}$. As $\mathsf U^{\mathsf T}=\mathsf U^{-1}$, the same follows for $G^{\mathsf T}$, $C$ and $\mathsf E$. For a single row, extend each field by the identity off the transport
paths of the window. The entries of $K_b$ are related by the same
similarity, since only the products along those paths enter them.

\emph{Step (d): decay.} Let $w:=\theta\,{\rm dist}(\cdot,b')$. Since $w(b')=0$, we have \[(K_b^k)(b,b')=e^{-w(b)}(e^wK_be^{-w})^k(b,b')\]. By \textup{(W)}, and because $b^{-1/2}\sum_kc_k\varkappa^k=(\frac2\gamma)^{1/2}$, this gives the bound on $\|G(b,b')\|$. The same computation with $K$ in place of $K_b$ gives the bound on $\|F(b,b')\|$. There are at most $3L^3$ bonds of $I^\circ$ to a block, and $(\frac2\gamma)^{1/2}\le\frac{2b^{1/2}}\gamma$; with these two facts, summing the bounds on the entries yields the remaining assertions of \textup{(d)}.

\emph{Step (e): closeness.} Fix $b$, let $Y=Y_b$, and let $\hat K:=K_Y$, extended by $0$, with the convention $\hat K^0:=1$. As $b\in Y$, we have $G(b,\cdot)=b^{-1/2}\sum_kc_k\delta_b^{\mathsf T}\hat K^k$. Moreover $K^k-\hat K^k=\sum_{j<k}\hat K^j(K-\hat K)K^{k-1-j}$, and every row vector $u^{\mathsf T}$ supported in $Y$ satisfies $u^{\mathsf T}(K-\hat K)=u^{\mathsf T}K\mathbf 1_{Y^c}$. It follows that $\Delta(b,b')=b^{-1/2}\sum_kc_k\sum_{j<k}\delta_b^{\mathsf T}\mathsf A\,\mathbf 1_{Y^c}\,\mathsf B\,\delta_{b'}$, with $\mathsf A=\hat K^jK$ and $\mathsf B=K^{k-1-j}$. Let $w_1=\theta\,{\rm dist}(b,\cdot)$ and $w_2=\theta\,{\rm dist}(\cdot,b')$. Then $\mathsf A(b,b_2)=\tilde{\mathsf A}(b,b_2)e^{-w_1(b_2)}$, where $\tilde{\mathsf A}=e^{-w_1}\mathsf Ae^{w_1}$ satisfies $\|\tilde{\mathsf A}\|\le\varkappa^{j+1}$, and $\mathsf B(b_2,b')=e^{-w_2(b_2)}\tilde{\mathsf B}(b_2,b')$, where $\tilde{\mathsf B}=e^{w_2}\mathsf Be^{-w_2}$ satisfies $\|\tilde{\mathsf B}\|\le\varkappa^{k-1-j}$; both norm bounds come from \textup{(W)}. Hence the term with the indices $k$ and $j$ equals $(\delta_b^{\mathsf T}\tilde{\mathsf A})\,\mathsf M\,(\tilde{\mathsf B}\delta_{b'})$, where $\mathsf M$ is the diagonal matrix $\mathsf M(b_2)=\mathbf 1_{Y^c}(b_2)e^{-\theta({\rm dist}(b,b_2)+{\rm dist}(b_2,b'))}$. Off $Y$, ${\rm dist}(b,b_2)\ge r_w+1$; together with the triangle inequality this yields $\|\mathsf M\|\le e^{-\theta\max\{r_w+1,{\rm dist}(b,b')\}}$. As $b^{-1/2}\sum_kkc_k\varkappa^k\le\sqrt2b\gamma^{-3/2}$, the bound on $\|\Delta(b,b')\|$ follows.

\emph{Step (e), continued: the sums.} The bound on $\|\Delta(b,b')\|$ is symmetric in $(b,b')$. Put $r=r_w+1$. At most $3L^3(2r+1)^3$ bonds $b'$ of $I^\circ$ satisfy ${\rm dist}(b,b')\le r$, for a given bond $b$. Moreover $\sum_{|n|_\infty>r}e^{-\theta|n|_\infty}=\sum_{m\ge1}(24(m+r)^2+2)e^{-\theta(m+r)}\le(1+r)^2e^{-\theta r}\mathsf g(\theta)$, since $(m+r)^2\le m^2(1+r)^2$. So the sums of $\|\Delta(b,b')\|$ over a row and over a column are $\le\sqrt2b\gamma^{-3/2}3L^3[(2r+1)^3+(1+r)^2\mathsf g(\theta)]e^{-\theta r}\le\sqrt2\cdot3\cdot2\,L^3b\gamma^{-3/2}(R+3)^3\mathsf g(\theta)e^{-\theta R/2}$, because $2r+1\le R+3$ and $r\ge\frac12R$. The bound on $\|\Delta\|$ follows by Schur's test.

\emph{Step (f): the inverse.} Let $F^{-1}:=b^{1/2}\sum_k\binom{1/2}k(-K)^k$. Then $FF^{-1}=1$, $F^{-2}=q$, and $\|F^{-1}\|\le2b^{1/2}$. We write $G=F(1-F^{-1}\Delta)$; under the condition of \textup{(f)}, $\|F^{-1}\Delta\|\le\frac12$, so that $\|G^{-1}\|\le4b^{1/2}$. Next, $G^{-1}-F^{-1}=G^{-1}\Delta F^{-1}$, and this difference has norm $\le8b\,d_R(\theta)$; further, $C^{-1}-F^{-2}=G^{-\mathsf T}(G^{-1}-F^{-1})+(G^{-1}-F^{-1})^{\mathsf T}F^{-1}$, of norm $\le6b^{1/2}\cdot8b\,d_R(\theta)$, and since $F^{-2}=q$ this is the bound on $\|S\|$. The identity $G^{-\mathsf T}(1+\mathsf E)G^{-1}=q$ gives the formula for $S$. Finally, $q^{-1}-C=F^2-(F-\Delta)(F-\Delta^{\mathsf T})=F\Delta^{\mathsf T}+\Delta F-\Delta\Delta^{\mathsf T}$, and \textup{(d)} and \textup{(e)} yield the last bound of \textup{(f)}. In this identity no inverse of $G$ occurs, and $\|\cdot\|_\infty$ is submultiplicative; so the last bound holds whatever the value of $d_R(\theta)$.

\emph{Step (g): the ratio.} We have $\mathsf E(b,b')+\delta_{bb'}=\sum_{b_1,b_2}G(b_1,b)^{\mathsf T}q(b_1,b_2)G(b_2,b')$, and only terms with ${\rm dist}(b,b_1)\le r_w$, ${\rm dist}(b_1,b_2)\le r_q$ and ${\rm dist}(b_2,b')\le r_w$ contribute; this gives the range of $\mathsf E$. As to the dependence, measured from $y(b)$: the row $G(b_1,\cdot)$ depends on bonds within $2r_w+\varrho_L+1$, the entry $q(b_1,b_2)$ on bonds within $r_w+r_q+\varrho_L+1$, and the row $G(b_2,\cdot)$ on bonds within $(r_w+r_q)+(r_w+\varrho_L+1)=\mathsf R_{\mathsf E}$. The formula for $\mathsf E$ results from expanding $(F-\Delta^{\mathsf T})q(F-\Delta)$ and using $FqF=1$ and $Fq=qF=F^{-1}$. Since $\|q\|\le b(1+\varkappa)\le2b$, under the condition of \textup{(g)} we obtain $\|\mathsf E\|\le4b^{1/2}d_R(\theta)+2b\,d_R(\theta)^2\le4.5\,b^{1/2}d_R(\theta)$. Lastly, each entry of $\mathsf E$ has norm at most $\|\mathsf E\|$, and the bound on $\|\mathsf E\|_\infty$ follows.

\emph{Step (h): the determinant.} The first identity follows from $\det G^{\mathsf T}=\det G$. At real fields $\det F>0$; moreover, the eigenvalues of the real matrix $1-F^{-1}\Delta$ lie in the disc $|z-1|\le\frac12$, and those which are not real come in conjugate pairs, so that its determinant is $>0$ and $\det G>0$. With $\mathsf X=F^{-1}\Delta$ and $\|\mathsf X\|\le\frac12$, the relations $G=F(1-\mathsf X)$, $F^{\mathsf T}=F$ and $FqF=1$ give $1+\mathsf E=G^{\mathsf T}qG=(1-\mathsf X)^{\mathsf T}(1-\mathsf X)$. At real fields this matrix is real and symmetric, while the singular values of $1-\mathsf X$ lie in $[\frac12,\frac32]$. By the spectral theorem, then, $\log(1+\mathsf E)$ exists, and $\det(1+\mathsf E)=\exp\operatorname{tr}\log(1+\mathsf E)$. As $\det G>0$ and $\det q>0$, we may take logarithms in the first identity. If $\|\mathsf E\|<1$, the power series of the logarithm converges, and its sum is the logarithm just described; by \textup{(g)}, $\|\mathsf E\|<1$ whenever $5b^{1/2}d_R(\theta)<1$.

\emph{Step (i): holomorphy.} We have $\operatorname{Ad}_gY=gYg^{-1}$, and $g^{-1}$ is the adjugate of $g$ because $\det g=1$; thus the entries of $\operatorname{Ad}_g^{\pm1}$ are polynomials in the entries of $g$. The partial sums of the series are therefore polynomials, and the bounds $\varkappa^k$ make the convergence uniform.

\emph{Step (j): the Lipschitz bound.} Let $Y=Y_{b_0}$, $K_i:=K_Y[\mathsf V_i]$, $D:=K_1-K_2$, $t(b):={\rm dist}(b_0,b)$, $w=\theta_{\mathbb C}t$ and $w'=\frac12\theta_{\mathbb C}t$. For a step $\mathsf s:b_1\to b_2$ whose path is made of the $m$ bonds $c_1,\dots,c_m$, \[\operatorname{Ad}_{\mathsf V_1(\pi)}-\operatorname{Ad}_{\mathsf V_2(\pi)}=\sum_i\operatorname{Ad}_{\mathsf V_1(c_1\cdots c_{i-1})}[\operatorname{Ad}_{\mathsf V_1(c_i)}-\operatorname{Ad}_{\mathsf V_2(c_i)}]\operatorname{Ad}_{\mathsf V_2(c_{i+1}\cdots c_m)}.\] The base of $c_i$ lies within $t(b_1)+m$ of $y(b_0)$, and $\|\operatorname{Ad}_{\mathsf V(c)}^{\pm1}\|\le e^{2\vartheta_0^\flat}$ for both fields. By the hypothesis of \textup{(j)}, therefore, $\|\mathsf T^{(1)}_{\mathsf s}-\mathsf T_{\mathsf s}^{(2)}\|\le me^{2\vartheta_0^\flat m}\epsilon_1(1+t(b_1)+m)\le\epsilon_1(1+t(b_1))\,m(1+m)e^{2\vartheta_0^\flat m}$. In $e^{-w}De^{w'}$ the entry $(b_1,b_2)$ carries the factor $e^{-\theta_{\mathbb C}t(b_1)+\frac12\theta_{\mathbb C}t(b_2)}\le e^{-\frac12\theta_{\mathbb C}t(b_1)}e^{\frac12\theta_{\mathbb C}m}$. Also $(1+t)e^{-\frac12\theta_{\mathbb C}t}\le\frac2{\theta_{\mathbb C}}=\frac{12}{\theta^\flat}$ and $2\vartheta_0^\flat+\frac12\theta_{\mathbb C}=\frac16\theta^\flat\le\theta_1$. Schur's test, with the third bound of \Cref{Y-C-CIL64}\,(b) for rows and columns alike, now gives $\|e^{-w}De^{w'}\|\le\frac{108A^\flat}{b\,\theta^\flat}\epsilon_1$. Because $w(b_0)=w'(b_0)=0$, we have $\delta_{b_0}^{\mathsf T}K_1^lDK_2^me^{w'}=\delta_{b_0}^{\mathsf T}(e^{-w}K_1e^w)^l(e^{-w}De^{w'})(e^{-w'}K_2e^{w'})^m$, and by \textup{(W)} the norm of this row is at most $\varkappa^{l+m}\frac{108A^\flat}{b\theta^\flat}\epsilon_1$. Further, $K_1^k-K_2^k=\sum_{l+m=k-1}K_1^lDK_2^m$. Finally, $b^{-1/2}\sum_kkc_k\varkappa^{k-1}=\sqrt2b\gamma^{-3/2}$ and $108\sqrt2\le153$; summing first over $l+m=k-1$ and then over $k$, we obtain the bound of \textup{(j)}.

\emph{Step (k): the bound on $d_R$.} By the choice of $R$ and of $c_R$ we have $e^{-\theta_{\mathbb C}R/2}\le x^{-64}$. With the definition of $d_R$ and the floor \textup{(F$_{\rm law}$)} this gives $48b^{3/2}d_R\le432L^3b^{5/2}\gamma^{-3/2}\mathsf g(\theta_{\mathbb C})(R+3)^3x^{-64}\le x^{-60}$. Since $b\ge16$, we obtain $4b^{1/2}d_R\le5b^{1/2}d_R\le48b^{3/2}d_R\le x^{-60}\le1$. By clause \textup{(f)} we have $\|S\|\le48b^{3/2}d_R$, and by clause \textup{(g)} we have $\|\mathsf E\|\le5b^{1/2}d_R\le48b^{3/2}d_R$. Hence the three bounds of clause \textup{(k)} hold: $4b^{1/2}d_R\le1$, $\|\mathsf E\|\le x^{-60}$ and $\|S\|\le x^{-60}$.
\end{proof}

\begin{remark}[On \Cref{Y-C-CITh66}]
\begin{enumerate}
\item[\textup{(1)}] Neither \Cref{Y-C-CITh66} nor its proof depends on a symmetry of $G$. The covariance of $G\xi$ is $x^{-1}GG^{\mathsf T}=x^{-1}C$.

\item[\textup{(2)}] The theorem is free of $M$; of $N$ it uses only $\frac\varrho L+2<\frac N2$; and its clauses \textup{(a)} to \textup{(j)} inclusive are free of $x$. Along the chain, the condition $c_R\theta_{\mathbb C}\ge128$ comes from $c_R=64/\theta_D$ together with $\theta_D\le\frac12\theta_{\mathbb C}$, as in \textup{(Fr2)} of item C.II 2.2.

\item[\textup{(3)}] If $\varrho\le R$, then $3r_w+\varrho_L+1\le2R+1$. Indeed, $\varrho_L\le\lceil R/3\rceil$, and $\frac32R+\frac13(R+2)+1\le2R+1$ for $R\ge4$; for $R=0,1,2,3$ the quantity $3r_w+\varrho_L+1$ is at most $1,2,5,5$, respectively, while $2R+1$ takes the values $1,3,5,7$.

\item[\textup{(4)}] Under the hypotheses of clause \textup{(k)}, $\|C\|\le\frac2\gamma$ at $G$-valued fields. Indeed, there $q\ge\gamma$ and $S$ is real symmetric, so that $C^{-1}=q+S\ge\gamma-x^{-60}\ge\frac\gamma2$, since \textup{(F$_{\rm law}$)} gives $x^{60}\ge x^4\ge432\gamma^{-3/2}\ge\frac2\gamma$. At the complex fields on which \textup{(u1$^\flat$)} holds, \Cref{Y-C-L66}\,(a) gives instead $\|C\|\le\frac4\gamma$. The hypotheses stated in that lemma are the clauses \textup{($\lambda$0)} and \textup{($\rho$5)} of item~C.II~2.1, the bound $\frac2\gamma x^{-60}\le\frac12$, and \textup{(u1$^\flat$)}. The proof of $\|C\|\le\frac4\gamma$ in \Cref{Y-C-L66}\,(a) uses, besides these hypotheses, one fact from the setting of the lemma, namely $\|q^{-1}\|\le\frac2\gamma$, which holds for the rows of item C.I~1.1 by \Cref{Y-C-CIL65}\,(c) at $w=0$. With exponent $60$, \textup{($\lambda$0)} and \textup{($\rho$5)} follow from \Cref{Y-C-CITh66}\,(a), (f), (k), and $x^{60}\ge x^4\ge432\gamma^{-3/2}\ge\frac4\gamma$ gives $\frac2\gamma x^{-60}\le\frac12$.

\item[\textup{(5)}] The floor \textup{(F$_{\rm law}$)} holds on a half line of values of $x$. Indeed, $R+3\le c_R\log x+4$, and the function $(c_R\log x+4)^3x^{-4}$ decreases for $\log x\ge\frac34$, because there $3c_R<4c_R\log x+16$.
\end{enumerate}
\end{remark}

\subsubsection*{C.I 4. White noise: the exact formula and the letters of the ratio}

\begin{definition}[The step, the chart]
We put $\operatorname{ang}(g):=\arccos\frac12\operatorname{tr}g\in[0,\pi]$. For $g\ne-\mathbf 1$ we let $\mathsf L(g)\in\mathfrak g$ be the unique $Y$ with $e^{iY}=g$ and $|Y|<\pi$, and we set $\mathsf L(-\mathbf 1):=0$.

For a $G$-valued fine field $U$ we write $\mathsf t_U(x):=U(\Gamma(y(x),x))$. For $c=(y,y')$ and $x\in B(y)$ we put $h_{c,x}:=\mathsf t_U(x)\,U([x,x'])\,\mathsf t_U(x')^{-1}$, then $\bar\tau_c:=L^{-3}\sum_{x\in B(y)}\mathsf L(h_{c,x}h_{c,y}^{-1})$ and $\mathsf M(U)(c):=e^{i\bar\tau_c}h_{c,y}$. Consequently $\mathsf M(U^u)(c)=u(y)\mathsf M(U)(c)u(y')^{-1}$.

Let $z(s):=\int_Ge^{-s(1-\cos\operatorname{ang}g)}dg$. We put $k_s(g,h)=z(s)^{-1}e^{-s(1-\cos\operatorname{ang}(gh^{-1}))}$ with $s=ax$. The step is \[\rho'(V')=\int\rho(U)\prod_ck_s\bigl(V'(c),\mathsf M(U)(c)\bigr)\,dU,\] where $dU$ denotes the Haar measure. The chart is the field $\mathsf U(\eta;V')$ defined by $\mathsf U(\eta;V')_b:=\mathbf 1$ on $\mathcal T$, by the value $e^{i\eta_b}$ at the other bonds inside blocks, and by the value $e^{i\eta_b}V'(c)$ at a crossing bond of $c$.

Further, $\mathfrak D=\{|\eta_b|<\pi\ \forall b\}$, $\mathsf j(\eta_b)=(\sin|\eta_b|/|\eta_b|)^2$, $\mathsf n_{\mathfrak L}=\#I^\circ$ and $\mathsf c_0=(2\pi^2)^{-\mathsf n_{\mathfrak L}}$. Finally, $\mathsf n_c=\operatorname{ang}(V'(c)\mathsf M(U)(c)^{-1})$.
\end{definition}

\begin{lemma}[The step in the coordinates of the chart]\label{Y-C-L51}
For every bounded, Borel, gauge invariant $\rho$, \[\rho'(V')=\mathsf c_0\int_{\mathfrak D}\rho(\mathsf U(\eta;V'))\prod_ck_s(V'(c),\mathsf M(\mathsf U(\eta;V'))(c))\prod_{b\in I^\circ}\mathsf j(\eta_b)\,d\eta.\]
\end{lemma}

\begin{proof}
Fix the variables on $\mathcal T$. For each block let $u(y)=\mathbf 1$
at its centre and set $u(x)=U(\Gamma(y,x))$ at its other sites. The
transformation $U\mapsto U^u$ sends the fixed tree variables to the
identity. On every bond outside the forest it is left and right
multiplication by fixed elements of $G$, and therefore preserves the
conditional product Haar measure. The integrand is invariant under
this transformation, since $u$ is the identity at the centres.
Consequently the conditional integral is the same for every choice of
tree variables. Their product Haar probability integrates to $1$.

On crossing bonds use right invariance of Haar measure to pass from
$U_b$ to $U_bV'(c)^{-1}$. In Pauli coordinates,
$\operatorname{SU}(2)$ is the unit sphere $S^3$ with normalized surface
measure, since multiplication preserves the Euclidean norm of the
four real coordinates. The exponential of $iY$, for
$Y=r\omega\cdot\sigma$ and $|\omega|=1$, has coordinates
$(\cos r,\sin r\,\omega)$. Thus $dg=(2\pi^2)^{-1}\sin^2 r\,
\mathrm{d}r\,\mathrm{d}\omega=(2\pi^2)^{-1}\mathsf j(Y)\,\mathrm{d}^3Y$
for $r<\pi$. The omitted element $-\mathbf 1$ has Haar measure zero.
\end{proof}

\begin{proposition}[The step in white noise]\label{Y-C-P67}
In the setting of \Cref{Y-C-CITh66} and under its hypotheses, let $4b^{1/2}d_R\le1$, and let $\mathsf V'$ satisfy \textup{(u1$^\flat$)}. Let $X'$ be an arbitrary $\mathfrak g_{\mathbb C}$-valued function on the coarse plaquettes, let $\xi$ be a function on $I^\circ$, and put \[\FitDisplay{\dimexpr\linewidth-1pt\relax}{h:=G^{\mathsf T}\mathsf gX',\quad a_0:=Gh=C\mathsf gX',\quad\eta:=G\xi-a_0,\quad F^\flat:=EX'+\mathrm{d}^{\mathsf V}\eta,\quad\bar F:=EX'-\mathrm{d}^{\mathsf V}a_0,\quad\mathcal Q(\eta):=\tfrac12F^\flat\!\cdot\!\mathsf BF^\flat+\tfrac a2(Q_1^{\mathsf V}\eta)\!\cdot\!(Q_1^{\mathsf V}\eta) .}\]
\begin{enumerate}
\item[\textup{(a)}] \emph{Algebra.} The function $\mathcal Q$ satisfies \[\mathcal Q(\eta)=\mathcal Q(-a_0)+\tfrac12\,\xi\!\cdot\!\xi+\tfrac12\,\xi\!\cdot\!\mathsf E\xi-\xi\!\cdot\!\mathsf Eh,\qquad\mathcal Q(-a_0)=\tfrac12\bar F\!\cdot\!\mathsf B\bar F+\tfrac a2(Q_1^{\mathsf V}a_0)\!\cdot\!(Q_1^{\mathsf V}a_0)=\mathsf m_C+\tfrac12h\!\cdot\!\mathsf Eh,\] where $\mathsf m_C=\frac12EX'\!\cdot\!\mathsf BEX'-\frac12\mathsf gX'\!\cdot\!C\mathsf gX'$; moreover $\eta\!\cdot\!S\eta=-(\xi-h)\!\cdot\!\mathsf E(\xi-h)$.

\item[\textup{(b)}] \emph{Exact formula, at every $V'\in\mathcal V'$.} Let $\rho$ be bounded, Borel and gauge invariant; let $n^\flat=-Q_1^V\eta$ and $U=\mathsf U(\eta;V')$; and let $X'=\mathsf L(\partial V')$, or else any other $\mathfrak g$-valued $X'$ on the coarse plaquettes that does not depend on $\eta$, whatever its dependence on $V'$ and on other variables; in particular $X'$ can be the $X'(V',e)$ of the pairs of item~E2 of \Cref{Y-D-C5}. Let $E_\xi$ denote the mean over $\xi\sim\mathcal N(0,x^{-1}\mathbf 1)$ on $\mathfrak L_{\mathfrak g}$. Then \[\FitDisplay{\dimexpr\linewidth-1pt\relax}{\rho'(V')=\mathsf c_0\,z(s)^{-3\#\mathbb T'}\bigl(\tfrac{2\pi}x\bigr)^{\frac32\mathsf n_{\mathfrak L}}\det G^V\ e^{-x\mathcal Q(-a_0)}\ E_\xi\Bigl[e^{-\frac x2\langle\xi,\mathsf E\xi\rangle+x\langle\xi,\mathsf Eh\rangle}\ 1_{\mathfrak D}(\eta)\prod_b\mathsf j(\eta_b)\ \rho(U)\,e^{\frac x2\langle F^\flat,\mathsf B^VF^\flat\rangle}\prod_ce^{\frac s2|n^\flat_c|^2-s(1-\cos\mathsf n_c)}\Bigr].}\]

\item[\textup{(c)}] \emph{The letters of the ratio.} Let $M\in\mathbb N$ with $M\mid N$. The cells are the cubes of $M^3$ coarse sites that tile $\mathbb T'$, and $b\in\hat\square$ means $y(b)\in\square$. For cells $\square,\square'$ we put \[T_{\square\square'}(\xi,\mathsf V'):=-\frac x2\sum_{b\in\hat\square,\,b'\in\hat\square'}\xi_b\!\cdot\!\mathsf E(b,b')\xi_{b'}+\delta_{\square\square'}\,x\sum_{b\in\hat\square}\xi_b\!\cdot\!(\mathsf Eh)_b.\] Then the ratio, by which we mean the factor $e^{-\frac x2\langle\xi,\mathsf E\xi\rangle+x\langle\xi,\mathsf Eh\rangle}$ under $E_\xi$ in the formula of clause \textup{(b)}, has $\sum_{\square,\square'}T_{\square\square'}$ as its exponent, and $T_{\square\square'}=0$ unless the coarse distance from $\square$ to $\square'$ is at most $\mathsf R_{\mathsf E}$. In the following three assertions, on the dependence, the holomorphy and the covariance of $T_{\square\square'}$, we take $X'$ to be given by $X'_P=-i\log\partial\mathsf V'(P)$, at fields with $|\partial\mathsf V'(P)-\mathbf 1|<1$. The function $T_{\square\square'}$ depends only on the values of $\xi$ on $\hat\square\cup\hat\square'$ and on those of $\mathsf V'$ on the bonds within distance $R_t:=4r_w+5\varrho_L+5$ of $\square$. It is holomorphic in $\mathsf V'$ on the interior of the region where \textup{(u1$^\flat$)} holds and $|\partial\mathsf V'(P)-\mathbf 1|<1$. For $G$-valued $u$ we have $T_{\square\square'}(\operatorname{Ad}_u\xi,\mathsf V'^u)=T_{\square\square'}(\xi,\mathsf V')$. Finally, let $\xi$ be real and put $|\xi_{\hat\square}|^2:=\sum_{b\in\hat\square}|\xi_b|^2$. Then, with the supremum taken over the $P$ on which $(\mathsf Eh)_b$, $b\in\hat\square$, depends, \[|T_{\square\square'}|\le\tfrac x4\|\mathsf E\|\bigl(|\xi_{\hat\square}|^2+|\xi_{\hat\square'}|^2\bigr)+\delta_{\square\square'}\,x\,\|\mathsf E\|_\infty\,\mathsf G_0\,b\,\sup\|X'_P\|\sum_{b\in\hat\square}|\xi_b| .\]
\end{enumerate}
\end{proposition}

\begin{proof}
\emph{Step (a): the algebra.} Since $\mathsf B$ and $q$ are symmetric for $Y\!\cdot\!Y'$, we have $\mathcal Q(\eta)=\frac12EX'\!\cdot\!\mathsf BEX'+\mathsf gX'\!\cdot\!\eta+\frac12\eta\!\cdot\!q\eta$. As $\eta=G(\xi-h)$, it follows that $\eta\!\cdot\!q\eta=(\xi-h)\!\cdot\!(1+\mathsf E)(\xi-h)$ and $\mathsf gX'\!\cdot\!\eta=h\!\cdot\!(\xi-h)$. Next, $\frac12(\xi-h)\!\cdot\!(\xi-h)+h\!\cdot\!(\xi-h)=\frac12\xi\!\cdot\!\xi-\frac12h\!\cdot\!h$, where $h\!\cdot\!h=\mathsf gX'\!\cdot\!C\mathsf gX'$. Therefore $\mathcal Q(\eta)=\mathsf m_C+\frac12\xi\!\cdot\!\xi+\frac12(\xi-h)\!\cdot\!\mathsf E(\xi-h)$, and at $\xi=0$ we have $\eta=-a_0$. Finally, $S=-G^{-\mathsf T}\mathsf EG^{-1}$.

\emph{Step (b): the exact formula.} In the formula of \Cref{Y-C-L51} we multiply and divide by $e^{-x\mathcal Q(\eta)}$, noting that $x\frac a2|Q_1^V\eta|^2=\frac s2|n^\flat|^2$. We then substitute $\eta=G\xi-a_0$, which is a bijection of $\mathfrak L_{\mathfrak g}$ with Jacobian $\det G^V>0$ (\Cref{Y-C-CITh66}\,(f), (h)), and we use $\int e^{-\frac x2|\xi|^2}d\xi=(2\pi/x)^{\frac32\mathsf n_{\mathfrak L}}$.

The argument depends on $X'$ through two facts alone. First, $a_0=C\mathsf gX'$ is a fixed vector of $\mathfrak L_{\mathfrak g}$, because at $V'\in\mathcal V'$ the matrices $C$ and $\mathsf g^V$ are real (\Cref{Y-C-CITh66}\,(a)) while $X'$ is $\mathfrak g$-valued and free of $\eta$; hence $\xi\mapsto G\xi-a_0$ is the bijection just used. Second, the argument uses clause \textup{(a)}, which is stated for every $X'$.

\emph{Step (c): the letters of the ratio.} The range of $\mathsf E$, and the variables on which $\mathsf E$ depends, are given by \Cref{Y-C-CITh66}\,(g). We have $h_b=\sum_{b_1}G(b_1,b)^{\mathsf T}(\mathsf gX')_{b_1}$ with ${\rm dist}(b,b_1)\le r_w$. Here $G(b_1,\cdot)$ depends on the field within $2r_w+\varrho_L+1$ of $y(b)$, and $(\mathsf gX')_{b_1}$ depends on the field within $r_w+2\varrho_L+2$ of $y(b)$, by \Cref{Y-C-CIL64}\,(e) and since $X'_P=-i\log\partial\mathsf V'(P)$ depends on $\mathsf V'$ through the bonds of $P$. Hence $(\mathsf Eh)_b$ depends on the field within $\mathsf R_{\mathsf E}+2r_w+2\varrho_L+2=R_t$ of $y(b)$. Holomorphy and covariance follow from \Cref{Y-C-CITh66}\,(i), (c). The bound follows from the Cauchy--Schwarz inequality; in it, $\|h_b\|\le\mathsf G_0\sup\|(\mathsf gX')_{b_1}\|\le\mathsf G_0b\sup\|X'_P\|$, by \Cref{Y-C-CIL64}\,(e) and $\|\operatorname{Ad}_{\mathsf V'(\pi)}\|\le e^{2\vartheta_0^\flat|\pi|}$.
\end{proof}

\subsection*{Part C.II. \textup{(S-cen$_k$)} and \textup{(S-loc$_k$)} at the law of windows}

\subsubsection*{C.II 1. Objects}

\begin{remark}
The $k$ of $\mathsf D_k(Y,\cdot)$, $S_k(Y;\cdot)$, \textup{(S-loc$_k$)}, \textup{(S-cen$_k$)}, \textup{(Law$_k$)} is part of the name of each of them and is not the number of a level. In Part C.II the letter $b$ denotes either a bond or the number given in item C.II~1.3, the letter $G$ denotes either the matrix of the law or the group, and the context shows which of the two senses is meant.
\end{remark}

\begin{definition}[Item C.II~1.1: lattices, transports and norms]
The objects of this item are those of item C.I~1.1. Throughout Part C.II, $L\ge3$ is odd, the number $\mathsf h_L$ is computed from it as in item C.I~1.1, and $aL\ge18$. We fix $M=L^{m'}$ with $m'\ge1$; then $M\ge L\ge3$, and $M$ is odd, so that every cell has a centre site. The number $N$ is even, and $M\mid N$.

The letter $N$ stands for the coarse side of the present step, and the condition \textup{(Tor)} below bears on this number. In \S\ref{Y-A-1} the same number is written $N'=NL^{K-k-1}$, and there $N$ denotes the side of the last torus.

We use the tori, base sites, distance, combs, transport convention and
norms of item~C.I~1.1. The coarse coboundary is $\mathrm{d}'$.
With $\mathsf g_I(\lambda):=3L^3\mathsf g(\lambda)$, the count of bonds
per block gives
\[
 \sum_{j\in I^\circ}{\rm dist}(i,j)e^{-\lambda{\rm dist}(i,j)}
 \le \lambda^{-1}\mathsf g_I(\lambda/2),
\]
since $t e^{-\lambda t/2}\le 1/\lambda$ for $t\ge0$.

We also use
$\|\operatorname{Ad}_g^{\pm1}-1\|\le2|g||g-\mathbf 1|$
for $g\in\operatorname{SL}(2,\mathbb C)$. For the positive sign, write
$gYg^{-1}-Y=[(g-\mathbf 1)Y-Y(g-\mathbf 1)]g^{-1}$.
For the negative sign, write
$g^{-1}Yg-Y=g^{-1}[Y(g-\mathbf 1)-(g-\mathbf 1)Y]$.
The norm bound follows from $|g^{-1}|=|g|$.
The map $\operatorname{Ad}_g$ preserves trace and the bilinear form
$Y\!\cdot\!Y'$.

The first bound in \textup{(E)} of item~C.I~1.1 holds for every $S\ge0$.
For $S\le1$ we also use its bound $0.72S^2$.

From now on $x>Le$ is a parameter and $x'$ denotes the number $x'=x/L$; the fine site $x'=x+Le_\mu$ of item C.I~1.1 appears from here on only within the symbol $[x,x']$. We put $p_0(x)=A_0(\log x)^{\mathfrak p_0}$, where $A_0>0$ and $\mathfrak p_0>2$, and $\varepsilon=p_0(x)x^{-1/2}$; further $p_0'=p_0(x')$ and $\varepsilon'=p_0'x'^{-1/2}$.

We also set $R=R(x)=\varrho(x)=\lceil c_R\log x\rceil$, $r(x)=\lceil\frac{20c_R}M\log x\rceil$ with $r'=r(x')$, $\ell_w(x)=6Mr(x)+12$, and $\mathsf b(x)=2-\frac1{\log x}$.

A cell $\square$ consists of $M^3$ coarse sites, and $\hat\square$ is the set of fine sites in the blocks of its sites; we assign each object to the cell of its base site. For a set $Y$ of cells, $Y^{+j}$ is the set of those cells whose sup distance to some cell of $Y$, counted in cells, is $\le j$, and we write $Y^\sim=Y^{+1}$ and $Y^\approx=Y^{+r'}$.

The condition \textup{(Tor)} is $N\ge16Mr(x)$. The number $\kappa$ satisfies $\kappa\ge6$, and $\kappa^\sharp=10(\kappa+4)$.

We write $\mathcal V'$ for the set of $G$-valued fields on $\mathbb T'$. For $V'$ in this set we put $\phi'_P=\operatorname{ang}\partial V'(P)$, where $\operatorname{ang}g=\arccos\frac12\operatorname{tr}g$, and we say that $V'$ is \emph{tame on $Y$} if $\phi'_P\le\varepsilon'$ for the plaquettes $P$ of $Y$.

For a set $Y$ of cells, $\mathcal U_Y(x')$ consists of the maps $\mathsf V$ that send the oriented coarse bonds of $Y^\approx$ into $\operatorname{SL}(2,\mathbb C)$, satisfy $\mathsf V(\bar c)=\mathsf V(c)^{-1}$ and obey \textup{(u1)}, \textup{(u2)} and \textup{(u3)} below; a bond, a path or a plaquette of $Y^\approx$ is one whose sites all lie in $Y^\approx$.
\begin{enumerate}
\item[\textup{(u1)}] $|\mathsf V(c)|<1+\vartheta(x')$ at every bond $c$, where $\vartheta$ denotes the function of item C.II~1.3, built with $\vartheta_0^\flat$.
\item[\textup{(u2)}] $|\mathsf V(\pi)|<\mathsf b(x')$ for each path $\pi$ in $Y^\approx$ with at most $\ell_w(x')$ bonds.
\item[\textup{(u3)}] $|\partial\mathsf V(P)-\mathbf 1|<2\varepsilon'$ at each plaquette of $Y^\approx$.
\end{enumerate}

Finally we set \[\varrho_L:=\lceil\varrho(x)/L\rceil,\qquad R_2=2R+2,\qquad R_3=R_2+\varrho_L+3,\qquad r_u:=\lceil R_2/M\rceil .\]
\end{definition}

We use the graph of cells with adjacency at sup distance one.
Count~1 in \S\ref{Y-G0-1} bounds the number of connected sets
of size $n$ through a cell by $(26e)^{n-1}$.

\begin{definition}[Item C.II~1.2: rows, the form and its steps]
The rows, the form and its steps are those of items C.I~1.1 and C.I~1.3, where $\mathsf H_y$, $\langle\cdot,\cdot\rangle$, $|y-p|$ and $\Pi_{\mathfrak L}$ are defined. As in \Cref{Y-A-Rows}, a row $\ell$ is of range $\le\varrho$ if $\ell_y(p)=0$ whenever $|y-p|>\varrho$. Rows $\ell$ on $\mathbb T$ obey $|\ell_y(p)|\le Ke^{-\lambda|y-p|}$, where $0<\lambda\le1$ and $K\ge0$, and they are of range $\le\varrho(x)$; moreover $\varrho_L+2<\frac N2$.

We use the form and the two step expansions of item~C.I~1.3 and
\Cref{Y-C-CIL64}(e). The transport of a step $\mathsf s$ from $b$
stays within $r_q=2\varrho_L+2$ of $y(b)$: its sites lie within
$\varrho_L+1$ of one of the endpoints, and for a step of kind
\textup{(R)} the second leg from $y(y)$ to $y(p')$ lies within
$\varrho_L$ of $y(y)$, which is within $\varrho_L+1$ of $y(b)$.
The last leg has length at most one. Thus the row $q(b,\cdot)$ depends
on the coarse bonds in that neighbourhood. The paths of the
$\mathsf g$-steps from $b$ stay within $2\varrho_L+1$ of $y(b)$ by
\Cref{Y-C-CIL64}(e).
\end{definition}

\begin{definition}[Item C.II~1.3: the numbers]
Use the numbers of item~C.I~1.2. In addition, set
\[
 \mathsf N_{\mathsf B}=1+3K^2\mathsf g(\lambda)^2,\qquad
 \mu=\frac{\gamma\theta_2}{4b\log\Lambda},\qquad
 \kappa_+=\frac M{40c_R}.
\]
We use $\vartheta(x)=\vartheta_0^\flat(1-1/\log x)$ and
$\mathsf g_x=\vartheta_0^\flat\log L/[4(\log x)^2]$.

For $\ell$ in the set $\mathcal R(b_*,K_*,\lambda_*,\frac19)$, the empty rows included, all these numbers are evaluated at $(K_*,\lambda_*)$, that is, with $K=K_*$ and $\lambda=\lambda_*$, and not with the constants of the bound that the rows at hand obey.

We have $\vartheta_0^\flat\le10^{-7}$, as in item C.I~1.2, because $\vartheta_0^\flat\le\frac\gamma{48A^\flat}$, $\gamma\le\frac1{5913}$ and $A^\flat\ge32e^2$; a fortiori $\vartheta_0^\flat\le10^{-4}$. Moreover $\theta_{\mathbb C}\le\frac16\theta_1\le\frac16$.

Throughout, we require $c_R\theta_{\mathbb C}\ge128$. If the numbers of this item are evaluated at $K=K_*$ and $\lambda=\lambda_*$, this inequality follows from the relations $\theta_D\le\frac12\theta_{\mathbb C}$ and $c_R=64/\theta_D$ in \textup{(Fr2)} of item C.II~2.2. For rows of item C.II~1.2 with other values of $K$ and $\lambda$, it is a requirement of this item and not a consequence of these relations. As $R\ge c_R\log x$, it follows that $e^{-\theta_{\mathbb C}R}\le x^{-128}$.
\end{definition}

\begin{definition}[Item C.II~1.4: fields on a box]
The data are the following: a box $\mathsf B$ of coarse sites, that is, a product of three intervals embedded in the torus, whose bonds, paths, plaquettes and cubes are those with all their sites in $\mathsf B$; a site $o\in\mathsf B$; a number $0<\hat\varepsilon\le\frac1{40}$; and an $\operatorname{SL}(2,\mathbb C)$ field $\mathsf V'$ on the bonds of $\mathsf B$.

We say that $\mathsf V'$ has \textup{(h)} at $\hat\varepsilon$ if it satisfies \textup{(h1)}, \textup{(h2)} and \textup{(h3)}.
\begin{enumerate}
\item[\textup{(h1)}] $|\mathsf V'(c)|\le1+\vartheta_0^\flat$ at every bond $c$ of $\mathsf B$.
\item[\textup{(h2)}] $|\mathsf V'(\pi)|\le2$ for each path $\pi$ in $\mathsf B$ with at most $\ell_w(x')$ bonds.
\item[\textup{(h3)}] $|\partial\mathsf V'(P)-\mathbf 1|\le2\hat\varepsilon$ at every plaquette $P$ of $\mathsf B$.
\end{enumerate}

We also use the two strict bounds below.
\begin{enumerate}
\item[\textup{(e1)}] $|\mathsf V'(c)|<1+\vartheta(x')+\mathsf g_x$ at every bond $c$ of $\mathsf B$.
\item[\textup{(e2)}] $|\mathsf V'(\pi)|<\mathsf b(x')+\mathsf g_x$ for each path $\pi$ in $\mathsf B$ with at most $\ell_w(x')$ bonds.
\end{enumerate}

These bounds give \textup{(h1)} and \textup{(h2)}, with the room $\vartheta_0^\flat-\vartheta(x')-\mathsf g_x\ge\frac{3\vartheta_0^\flat}{4\log x}$ and $2-\mathsf b(x')-\mathsf g_x\ge\frac3{4\log x}$. The room follows from the definitions of $\vartheta$ and $\mathsf b$, since $\mathsf g_x\le\frac{\vartheta_0^\flat}{4\log x}$ and $\log x'\le\log x$.

We put $X'_P:=-i\log\partial\mathsf V'(P)$, defined by the power series; under \textup{(h3)} it satisfies $|X'_P|\le2.06\hat\varepsilon$. With $\mathsf k(y):=\mathsf V'(\Gamma(o,y))$ we define $\mathsf W(c):=\mathsf k(c_-)\mathsf V'(c)\mathsf k(c_+)^{-1}$ and $\tilde X_P:=\operatorname{Ad}_{\mathsf k(y_P)}X'_P=T(o,y_P)X'_P$. If $\pi$ is a path in $\mathsf B$ from $y$ to $y''$, then $\mathsf W(\pi)=\mathsf k(y)\mathsf V'(\pi)\mathsf k(y'')^{-1}$.

Every field on $\mathsf B$ is extended by $\mathbf 1$ on the remaining bonds of $\mathbb T'$, and we set $X'_P:=0$ for $P\not\subset\mathsf B$.
\end{definition}

\begin{definition}[Item C.II~1.5: sub-forms]
A \emph{family} $\mathfrak S$ consists of a set $I_{\mathfrak S}\subset I^\circ$ and of steps joining bonds of $I_{\mathfrak S}$, the set of steps being closed under reversal. On $\ell^2(I_{\mathfrak S};\mathbb C^3)$ we set $q_{\mathfrak S}^{\mathsf V}:=\sum_{\mathsf s\in\mathfrak S}q_{\mathsf s}\mathsf T_{\mathsf s}$, and we put $K_{\mathfrak S}^{\mathsf V}(\zeta):=1-b^{-1}\sum_{\mathsf s\in\mathfrak S}\zeta^{|\mathsf s|}q_{\mathsf s}\mathsf T_{\mathsf s}$ and $K_{\mathfrak S}^{\mathsf V}=K_{\mathfrak S}^{\mathsf V}(1)$. We call $\mathfrak S$ \emph{coercive} when $\gamma\le q^{V_0}_{\mathfrak S}$ for all $G$-valued $V_0$.

Families of two kinds occur, \textup{(w)} and \textup{(c)}.
\begin{enumerate}
\item[\textup{(w)}] Windows: $I_{\mathfrak S}=\mathsf Y\subset I^\circ$, with all the steps between bonds of $\mathsf Y$; for them $q_{\mathfrak S}=\mathbf 1_{\mathsf Y}q_\ell\mathbf 1_{\mathsf Y}=:q_{\mathsf Y}$.
\item[\textup{(c)}] Sets of cells: given a set $\mathsf Y$ of cells, we let $I_{\mathsf Y}$ consist of the bonds of $I^\circ$ based in $\hat{\mathsf Y}$. To a step $\mathsf s$ of kind \textup{(P)} or \textup{(Q)} we attach $\mathsf c(\mathsf s):=\{$the cells of $y(b)$, $y(b')\}$, and to a step of kind \textup{(R)} we attach these cells together with the cells met by $\Gamma(y(y),y(p))$ and by $\Gamma(y(y),y(p'))$. The family consists of the steps $\mathsf s$ with $\mathsf c(\mathsf s)\subset\mathsf Y$, and its form is \[q_{[\mathsf Y]}:=q_{\mathfrak S}=\mathbf 1_{\mathsf Y}q_\emptyset\mathbf 1_{\mathsf Y}+\mathbf 1_{\mathsf Y}\mathrm{d}^{V*}\ell^{[\mathsf Y]V*}\ell^{[\mathsf Y]V}\mathrm{d}^V\mathbf 1_{\mathsf Y},\] where $\ell_y^{[\mathsf Y]}(p)=\ell_y(p)1\{$the cells met by $\Gamma(y(y),y(p))$ lie in $\mathsf Y\}$.
\end{enumerate}

Both kinds are coercive by \Cref{Y-C-CIIL64}: a window yields a compression of a form $\ge\gamma$, and a set of cells yields a compression plus a square.

For every step, $\mathsf c(\mathsf s)\subset\{$cells of the sites of $\pi_{\mathsf s}\}$, and $\mathsf c(\mathsf s)$ is connected, since two coarse sites at sup distance $\le1$ always lie in one cell or in two neighbouring cells, whatever $M\ge1$ is.
\end{definition}

\subsubsection*{C.II 2.1. The law: what Part C.II uses of \textup{(Law$_k$)}}

\begin{definition}[The law and the sets of fields]
In Part C.II, $(G^{\mathsf V},C^{\mathsf V})$ is the law of windows of Part C.I, given by \Cref{Y-C-CIDef4}, taken at $R=R(x)$ and for rows of range $\le\varrho(x)=R$, in the setting where $c_R\theta_{\mathbb C}\ge128$ and the floor \textup{(F$_{\rm law}$)} of \Cref{Y-C-CITh66}\,(k) hold.

We write $\mathfrak V$ for the set of $\operatorname{SL}(2,\mathbb C)$ fields on $\mathbb T'$ that satisfy \textup{(u1$^\flat$)} of item C.I~1.1, namely $|\mathsf V'(c)|\le1+\vartheta_0^\flat$ for all bonds $c$, and $\mathfrak V^\circ$ for the set given by the same bounds with strict inequalities. Then $\mathcal V'\subset\mathfrak V$.

In Part C.II, $\varrho_L$ stands for $\lceil\varrho(x)/L\rceil$; it is at least the $\varrho_L$ of Part C.I, and every radius below increases with it.

The proofs of Part C.II rest on the law through the clauses listed in the proposition below and through nothing else. Each of these clauses holds at every field of $\mathfrak V$, complex fields included, by the clause of \Cref{Y-C-CITh66} named at its end, applied with $\theta=\theta_{\mathbb C}$.
\end{definition}

\begin{proposition}[Item C.II~2.1: What Part C.II uses of the law]
The law has the following properties.
\begin{enumerate}
\item[\textup{($\lambda$0)}] $G^{\mathsf V}$ and $C^{\mathsf V}$ are defined at every field of $\mathfrak V$ (by the part of \Cref{Y-C-CITh66} that precedes its clauses, and by its clause (a)).
\item[\textup{($\lambda$1)}] $G^{\mathsf V}(b,b')=0$ whenever ${\rm dist}(b,b')>\frac12R$, and $(G\xi)_b$ depends only on the values of $\xi$ on $Y_b$ (\Cref{Y-C-CITh66}\,(b)).
\item[\textup{($\lambda$2)}] $G^{\mathsf V}(b,\cdot)$ is determined by the restriction of $\mathsf V'$ to the coarse bonds within $r_w+\varrho_L+1$ of $y(b)$, and $r_w+\varrho_L+1\le R_3$ (\Cref{Y-C-CITh66}\,(b)).
\item[\textup{($\lambda$3)}] $G[\mathsf V'^u]=\operatorname{Ad}_uG[\mathsf V']\operatorname{Ad}_u^{-1}$ for $u:\mathbb T'\to G$ (\Cref{Y-C-CITh66}\,(c)).
\item[\textup{($\lambda$4)}] $G^{\mathsf V}$ has entries that are holomorphic on $\mathfrak V^\circ$ and real at $G$-valued fields (\Cref{Y-C-CITh66}\,(i) and (a)).
\item[\textup{($\lambda$5)}] $\|(G^{\mathsf V}\xi)_b\|\le\mathsf G_0\sup_{Y_b}\|\xi_{b'}\|\le C_\varphi^\flat|\xi|_\infty$ (\Cref{Y-C-CITh66}\,(d)).
\item[\textup{($\lambda$6) = ($\rho$5)}] $C^{\mathsf V}$ is invertible, and $\|(C^{\mathsf V})^{-1}-q^{\mathsf V}_\ell\|\le x^{-\mathsf n}$ in the operator norm on $\ell^2(I^\circ;\mathbb C^3)$, with $\mathsf n=60$ (\Cref{Y-C-CITh66}\,(f) and (k)).
\item[\textup{($\lambda$7)}] $C=GG^{\mathsf T}$ (by \Cref{Y-C-CIDef4}).
\item[\textup{($\rho$1)}] $C^{\mathsf V}(b,b')=0$ whenever ${\rm dist}(b,b')>R$ (\Cref{Y-C-CITh66}\,(b), since $2r_w\le R$).
\item[\textup{($\rho$2)}] $C^{\mathsf V}(b,\cdot)$ is determined by the restriction of $\mathsf V'$ to the coarse bonds within $R_2-1$ of $y(b)$ (\Cref{Y-C-CITh66}\,(b), since $3r_w+\varrho_L+1\le2R+1$ by Remark~(3) after \Cref{Y-C-CITh66}).
\item[\textup{($\rho$3)}] $C[\mathsf V'^u]=\operatorname{Ad}_uC[\mathsf V']\operatorname{Ad}_u^{-1}$ for $G$-valued $u$ (\Cref{Y-C-CITh66}\,(c)).
\item[\textup{($\rho$4)}] $C^{\mathsf V}$ has entries that are holomorphic on $\mathfrak V^\circ$ (\Cref{Y-C-CITh66}\,(i)).
\item[\textup{($\rho$6)}] $C^{\mathsf T}=C$, and at $G$-valued fields the matrix $C$ is real with $\det C=(\det G)^2>0$ (\Cref{Y-C-CITh66}\,(b), (a) and (f)).
\end{enumerate}
\end{proposition}

\begin{remark}[Item C.II~2.1: dependence on $\mathsf n$, and notation]
In each statement below, the exponent $\mathsf n$ of \textup{($\rho$5)} is a letter, and a law is any map $\mathsf V'\mapsto(G^{\mathsf V},C^{\mathsf V})$ on $\mathfrak V$ that satisfies the clauses on which the statement rests. These statements hold for every such law whose exponent lies in the range $2\le\mathsf n\le127$ and, in a statement that assumes \textup{(S2)}, also obeys $\mathsf n\ge4$. The law of windows fixed at the start of this item is among these laws, with $\mathsf n=60$. The proofs from \Cref{Y-C-L66} to \Cref{Y-C-CIICor72} depend on $\mathsf n$ only through this range. The passage from the law to a window, in \Cref{Y-C-L66}\,(c), needs $\mathsf n\le128$: there the term of the window carries the factor $e^{-\theta_{\mathbb C}R}\le x^{-128}$ of item C.II~1.3, and $x^{-128}\le x^{-\mathsf n}$. This condition follows from the range. The proofs of \Cref{Y-C-CIITh74} and \Cref{Y-C-CIICor75} need, in addition, $\mathsf n\ge4$, which holds in the setting \textup{(S2)} of item C.II~5, under ``Floors''. Indeed, the geometric sum that occurs in \Cref{Y-C-CIITh74}\,(c) is controlled through $x^{-(\mathsf n-2)}\le\frac12$, and \Cref{Y-C-CIICor75} rests on $x^{-(\mathsf n-3)}\le\varepsilon'$.

We write $S^{\mathsf V}:=(C^{\mathsf V})^{-1}-q^{\mathsf V}_\ell$, $a_0:=C^{\mathsf V}\mathsf g^{\mathsf V}X'$ and $\mathsf u^{\mathsf V}:=1-C^{\mathsf V}q_\ell^{\mathsf V}$.
\end{remark}

\subsubsection*{C.II 2.2. \textup{(Fr)}: the free scalar model at a level with rows}

\begin{proposition}[The free scalar model at a level with rows]
All kernels here are real and scalar, and none of them depends on a field. We write $q=q^{\mathbf 1}_\ell$, $\mathsf g=\Pi_{\mathfrak L}\mathrm{d}^*\mathsf BE$, $H_\circ=E-\mathrm{d}q^{-1}\mathsf g$ and $\mathsf N_\circ=Q_1q^{-1}\mathsf g$, as in \Cref{Y-A-P4}\,(c). The kernels $\mathsf G,\mathsf G_Z,\mathsf G_{\mathsf n},H^\sharp,\mathsf N^\sharp,\mathsf T$, the next rows $\ell'$ with $\mathsf B'$ and the cut rows $\ell'^{\rm c}$ with $\mathsf B'^{\rm c}$ are those of Props.~\ref{Y-A-P4} and \ref{Y-A-P5}, Defs.~\ref{Y-A-XT} and \ref{Y-A-GG}, \Cref{Y-A-L7} and \Cref{Y-A-CutRows}, the weighted norm $\|\cdot\|_\theta$ is that of \Cref{Y-A-Norm}, and the constants $\mathsf N_Z$, $\mathsf N_\sharp$ are those defined in the remark after this proposition. We consider the following hypotheses.
\begin{enumerate}
\item[\textup{(H1)}] $L\ge L_A$, and $L$ is odd.
\item[\textup{(H2)}] $aL\ge3/c_L$, where $c_L=\frac{L^2-1}{6L^2}$; in particular $aL>18$.
\item[\textup{(H3)}] $N$ is even.
\item[\textup{(H4)}] $x'=x_{k+1}\ge x_\flat$. Together with \textup{(X12)} of \S\ref{Y-A-1}, this keeps the rows $\ell_{k+1}=((\ell_k)')^{\rm c}$ of the chain, cut at $\varrho(x_{k+1})$, in $\mathcal R(b_*,K_*,\lambda_*,\frac19)$.
\end{enumerate}

Clause \textup{(Fr1)} below holds for all rows $\ell$. Under \textup{(H1)}, \textup{(H2)} and \textup{(H3)}, clauses \textup{(Fr2)} and \textup{(Fr3)} below hold for every $\ell\in\mathcal R(b_*,K_*,\lambda_*,\frac19)$ and at every cut; under the same three hypotheses, \textup{(H4)} and \textup{(X12)} keep the rows of the chain in this set, and \textup{(H4)} is used for that alone. The relations $C_D\ge4$, $0<\theta_D\le\frac12\theta_{\mathbb C}$ and $c_R=64/\theta_D$ in \textup{(Fr2)} are fixed by the choice of the letters at \textup{(O6)} in item~3 of \S\ref{Y-G0-1}; they are relations between letters of \textup{(O6)}, with $\theta_{\mathbb C}$ taken at $(K_*,\lambda_*)$, and they involve no rows, no torus and none of \textup{(H1)} to \textup{(H3)}.
\begin{enumerate}
\item[\textup{(Fr1)}] $H_\circ=H^\sharp+\mathsf G\mathrm{d}'$ and $\mathsf N_\circ=\mathsf N^\sharp+\mathsf G_{\mathsf n}\mathrm{d}'$, with $\mathsf G_{\mathsf n}=a^{-1}\mathrm{d}'^*\mathsf G_Z$. This is \Cref{Y-A-L7}\,(c), which is formulated and proved there for rows $\ell$.
\item[\textup{(Fr2)}] $C_D\ge4$, $0<\theta_D\le\frac12\theta_{\mathbb C}$ and $c_R=64/\theta_D$, and each of $\mathsf G,\mathsf G_Z,H^\sharp,\mathsf N^\sharp,\mathsf T,\ell'$ satisfies $\|\cdot\|_{\theta_D}\le C_D$ on every torus. This comes from \textup{(Dec)}$(C_D,\theta_D)$ of \Cref{Y-A-Dec}, applied to the pair $(\ell,\ell')$, that is, from \Cref{Y-A-P35} with $\theta_D:=\frac12\min\{\lambda_2^*,\theta_{\mathbb C},\theta_J\}$, which amounts to putting $\theta_{\mathbb C}$ in the place of $\theta_q$. The constant $C_D$ is the letter $C_D$ of \textup{(O6)} in item~3 of \S\ref{Y-G0-1}; it is at least the value at this $\theta_D$ of the expression that \Cref{Y-A-P35} displays for its own $C_D$, a maximum built from $A_q$, $\mathsf C_0$, $\mathsf C_R$ and $\mathsf C_\ell$, and in \Cref{Y-A-Fr} that expression is item (6) of the order of the constants. The bound \textup{(Dec)}$(C,\theta_D)$ holds for every $C$ at least that expression, and every constant below increases with $C_D$. Since $\mathsf G_{\mathsf n}=a^{-1}\mathrm{d}'^*\mathsf G_Z$ and $aL\ge18$, it follows that $\|\mathsf G_{\mathsf n}\|_{\theta_D}\le\frac L{18}\cdot4e\,C_D\le LC_D$; moreover $e^{-\theta_DR(x)}\le x^{-64}$.
\item[\textup{(Fr3)}] $L^2\|H^\sharp-\mathsf T(\mathsf B'-\mathsf B'^{\rm c})\|_\infty\le\mathsf N_\sharp$ and $\|\mathsf B'^{\rm c}\|_\infty\le\mathsf N_Z$, where $\mathsf N^\sharp=(aL)^{-1}\mathrm{d}'^*\mathsf B'$. This rests on \textup{(N$'$)} of \Cref{Y-A-Np}, taken at $m\equiv1$ and with the objects $m$, $\zeta'$, $Z$, $\Xi$, $\nu$, $\omega$, $\sigma$ that are defined there just before \textup{(N$'$)}.
\end{enumerate}
\end{proposition}

\begin{remark}[Source and hypotheses]
The objects of this item include the next rows $\ell'$ and $\mathsf B'=1+\ell'^{\mathsf T}\ell'$, and the cut rows $\ell'^{\rm c}$, namely $\ell'$ cut at $\varrho(x')$, together with $\ell'^>:=\ell'-\ell'^{\rm c}$ and $\mathsf B'^{\rm c}=1+\ell'^{{\rm c}\mathsf T}\ell'^{\rm c}$. The statements and definitions named in the proposition build on Defs.~\ref{Y-A-Pairs} and \ref{Y-A-AuxForm} and on \Cref{Y-A-L5p}.

The weighted norm of \Cref{Y-A-Norm} is \[\|K\|_\theta=\max\{\sup_i\sum_j|K(i,j)|e^{\theta{\rm dist}(i,j)},\ \sup_j\sum_i|K(i,j)|e^{\theta{\rm dist}(i,j)}\}\] and $\|K\|_0$ denotes the same expression at $\theta=0$. The sets $\mathcal R(b,K,\lambda,\tau)$ are those of \S\ref{Y-A-1}, from \Cref{Y-A-CovRows} to \Cref{Y-A-Towers}. On $\mathfrak g$-valued cochains a scalar kernel $K$ operates as $K\otimes1$ and keeps its norms, so these norms are free of the group.

Clause \textup{(Fr1)} is algebra. The proofs of \textup{(Fr2)} and \textup{(Fr3)} for $\ell\in\mathcal R(b_*,K_*,\lambda_*,\frac19)$ are given in \S\ref{Y-A-1}, in Lems.~\ref{Y-A-L24} to \ref{Y-A-L27}, \Cref{Y-A-P28}, \Cref{Y-A-Th29}, Lems.~\ref{Y-A-L31} and \ref{Y-A-L32}, \Cref{Y-A-Th33}, \Cref{Y-A-P35} and \Cref{Y-A-Th36}, and in Props.~\ref{Y-A-P4}, \ref{Y-A-P5}, \ref{Y-A-P6} and \ref{Y-A-P8}.

In \Cref{Y-A-Fr} the conclusion is stated for the pairs $(\ell_k,(\ell_k)')$ of the chain. The two clauses hold, however, for every $\ell$ of the set, as follows. Clause \textup{(Fr2)} is obtained in Step~H13 of the proof of \Cref{Y-A-Fr}, namely from \Cref{Y-A-P35} and its proof, taken with $\theta_{\mathbb C}$ in place of $\theta_q$, together with Step~H16\,(i), (ii) of that proof. Clause \textup{(Fr3)} follows from the argument of Step~H14 of the same proof; this argument uses of $\ell$ only that it belongs to the set, and it goes as follows.

\Cref{Y-A-Th29}\,(a) applied with $K=K_*$ and $\tau\le\frac19$, which $L\ge L_A$ allows, yields $|\ell'_{y'}(P)|\le K_*e^{-\lambda_*{\rm dist}(y',P)}$; consequently, at every cut, $\|\mathsf B'^{\rm c}\|_\infty\le1+K_*\mathsf g(\lambda_*)\cdot3K_*\mathsf g(\lambda_*)=\mathsf N_Z$. Next, \Cref{Y-A-Th33}\,(d) with $(b,K,\lambda,\lambda^\circ)=(b_*,K_*,\lambda_*,\lambda_*)$, as in \Cref{Y-A-Stars}, yields $L^2\|\mathsf T\|_\infty\le c_1^*\mathsf l_L$; here $\lambda'=\lambda_*$, by part \textup{(a)} of the proof of \Cref{Y-A-Th29}. Finally, the identity $H^\sharp=\mathsf T\mathsf B'$ gives $H^\sharp-\mathsf T(\mathsf B'-\mathsf B'^{\rm c})=\mathsf T\mathsf B'^{\rm c}$, and the two bounds just found give, for this kernel, $L^2\|\cdot\|_\infty\le c_1^*\mathsf l_L\mathsf N_Z\le\mathsf N_\sharp$, where $\mathsf N_\sharp=c_1^*\mathsf l_L\mathsf N_Z\mathsf P_\rho$, with $\mathsf P_\rho\ge1$, is the constant that \Cref{Y-A-Fr} fixes in items (3), (4) of the order of the constants.

Each statement of Part C.II that rests on \textup{(Fr2)} or on \textup{(Fr3)} carries with it the hypotheses \textup{(H1)} to \textup{(H4)}.

Where $\ell\in\mathcal R(b_*,K_*,\lambda_*,\frac19)$ is part of the data, as in \textup{(SD)}, the hypotheses \textup{(H1)} to \textup{(H3)} suffice: \Cref{Y-A-P35} requires nothing more, and the argument for \textup{(Fr3)} given above rests on membership alone. For the empty rows membership holds: $\emptyset\in\mathcal R(1,0,\lambda,0)$, by \Cref{Y-A-SetsOfRows}.

The constants $C_D,\theta_D,\mathsf N_\sharp,\mathsf N_Z$ depend neither on $k$ nor on $N$ nor on $x$. Moreover $\theta_D\le\frac12\theta_{\mathbb C}\le\frac18\lambda_*$, a bound on which \Cref{Y-J-ThA} rests; it follows from $\theta_{\mathbb C}\le\frac\gamma{12A^\flat}$, $A^\flat\ge384e^8K_*^2$, $K_*^2=9/\gamma_*$, $\lambda_*=\gamma_*^{1/2}/44$ and $\gamma,\gamma_*\le1$.

The data \textup{(SD)} of item~1 of \S\ref{Y-G0-1}, namely the four conclusions that \Cref{Y-A-Th36} draws for the pair $(\ell,\ell')$, imply \textup{(Fr1)} to \textup{(Fr3)}, but they form a different statement: \textup{(Fr2)} is \textup{(Dec)} for six of the sixteen kernels of \textup{(SD)}, \textup{(Fr3)} is \textup{(N$'$)} at the mask $m\equiv1$, and \textup{(Fr1)} is \Cref{Y-A-L7}\,(c).

Thms.~\ref{Y-C-Th68} and \ref{Y-C-Th70} rest on \textup{(Fr2)} for $\mathsf G,\mathsf G_Z$ alone, and Thms.~\ref{Y-C-CIITh73}, \ref{Y-C-CIITh74} and \Cref{Y-C-CIICor75} are free of \textup{(Fr1)} to \textup{(Fr3)}.
\end{remark}

\begin{definition}[The matrices of \textup{(S-cen$_k$)}] The matrices of \textup{(S-cen$_k$)} are built from the matrix $C^{\mathbf 1}$ of the law, that is, from $C^{\mathsf V}$ at the trivial field. We put $H^{(x)}:=E-\mathrm{d}C^{\mathbf 1}\mathsf g$ and $\mathsf N^{(x)}:=Q_1C^{\mathbf 1}\mathsf g$. For a fine plaquette $i=p$ and for a coarse bond $i=c$, let $\mathcal Q_i$ denote the set of those coarse cubes each face of which lies within $R_2$ of $i$. With these sets we define \[H^\natural(p,P):=H^{(x)}(p,P)-\sum_{q'\in\mathcal Q_p}\mathsf G(p,q')\mathrm{d}'(q',P),\qquad\mathsf N^\natural(c,P):=\mathsf N^{(x)}(c,P)-\sum_{q'\in\mathcal Q_c}\mathsf G_{\mathsf n}(c,q')\mathrm{d}'(q',P).\] \end{definition}

\begin{remark} A priori the entries of these matrices are $3\times3$ blocks. For a constant $G$-valued $u$ one has $\mathbf 1^u=\mathbf 1$; so \textup{($\rho$3)}, taken at this field, shows that each block $C^{\mathbf 1}(b,b')$ commutes with all $\operatorname{Ad}_u$; by \textup{($\rho$6)} each such block is also real, since the field $\mathbf 1$ is $G$-valued, and therefore $C^{\mathbf 1}(b,b')$ is a real number times $1_{\mathfrak g}$. \end{remark}

\subsubsection*{C.II 3. Tools}

\begin{lemma}[Bounds on the quadratic form $q^V_\ell$ and sums over steps]\label{Y-C-CIIL64} In the setting of item~C.II~1.2, with rows as there, let $L\ge3$ be odd and $aL\ge18$. Then the following hold.

\begin{enumerate}\item[\textup{(a)}] At every $V'\in\mathcal V'$ and for every $\eta$, $\gamma|\eta|^2\le q^V_\emptyset(\eta)\le q^V_\ell(\eta)\le b|\eta|^2$.

\item[\textup{(b)}] $\sup_b\sum_{\mathsf s\text{ from }b}|q_{\mathsf s}|\le b$; and $\varpi^\flat(\theta):=\sup_b\sum_{\mathsf s\text{ from }b}|q_{\mathsf s}|(e^{\theta|\mathsf s|}-1)\le A^\flat\theta$ whenever $0\le\theta\le\theta_1$.

\item[\textup{(c)}] $\sup_b\sum_{\mathsf s}|q_{\mathsf s}||\mathsf s|\le A^\flat$, $\sup_b\sum_{\mathsf s}|q_{\mathsf s}||\mathsf s|^2\le2A^\flat/\theta_1$ and $\sup_b\sum_{\mathsf s}|q_{\mathsf s}||\mathsf s|e^{\theta_2|\mathsf s|}\le2A^\flat$.

\item[\textup{(d)}] $\sup_b\sum_{\mathsf t\text{ from }b}|g_{\mathsf t}|\le4\mathsf N_{\mathsf B}$; moreover, for each bond $b$ and each increasing function $f\ge0$, $\sum_{\mathsf t\text{ from }b}|g_{\mathsf t}|f(|\mathsf t|)\le\sum_{\mathsf s\text{ from }b}|q_{\mathsf s}|f(|\mathsf s|)$. \end{enumerate}\end{lemma}

\begin{proof}
Part~(a) is \Cref{Y-C-CIL64}(a) on the full bond set, applied both
to the empty rows and to the given rows. Adding rows adds a nonnegative
square. Part~(b) is the first and second bounds of
\Cref{Y-C-CIL64}(b).

\emph{Step (c): moments of the length.} For $m\ge0$ we have $e^{\theta m}-1\ge\theta m$ and $e^{\theta m}-1\ge\frac12\theta^2m^2$. With $\theta=\theta_1$, these inequalities and \textup{(b)} give the first two bounds of \textup{(c)}. For the third we apply $ue^u\le e^{2u}-1$, valid for $u\ge0$, at $u=\theta_2m$; as $\theta_2=\frac12\theta_1$, this gives $\sum|q_{\mathsf s}||\mathsf s|e^{\theta_2|\mathsf s|}\le\theta_2^{-1}\varpi^\flat(\theta_1)$, and by \textup{(b)} the right side is at most $2A^\flat$.

\emph{Step (d): the $\mathsf g$-steps.} A bond $b$ lies in four plaquettes $p\ni b$; for each of them, the sum of $|g_{\mathsf t}|$ over the $\mathsf g$-steps $\mathsf t$ from $b$ through $p$ runs once over the $\mathsf g$-step $(p)$ and once over each $\mathsf g$-step $(y,p,p'')$, because each plaquette $p''$ lies under a single $P$; this sum is at most $\sum_{p''}\bigl[\delta_{pp''}+\sum_y|\langle\ell_y(p),\ell_y(p'')\rangle|\bigr]$, which by the two row bounds in \Cref{Y-A-SetsOfRows} is $\le1+\sum_y|\ell_y(p)|\sum_{p''}|\ell_y(p'')|\le1+K\mathsf g(\lambda)\cdot3K\mathsf g(\lambda)=\mathsf N_{\mathsf B}$. This gives the first assertion of \textup{(d)}. For the second we map the $\mathsf g$-steps from $b$ into the steps from $b$. A $\mathsf g$-step $(p)$ is sent to the step $(p):b\to b$ of kind \textup{(P)}; then $|\mathsf s|=2|\mathsf t|$. A $\mathsf g$-step $(y,p,p'')$ is sent to the step $(y,p,p''):b\to b''$ of kind \textup{(R)}, where $b''$ is the bond of $p''$ based at $\pi(p'')$ that points in the first direction of $p''$. This $b''$ is a crossing bond and therefore lies in $I^\circ$; moreover $y(b'')=y(p'')$, so that $|\mathsf s|=|\mathsf t|$. The map so defined is injective, with $|q_{\mathsf s}|=|g_{\mathsf t}|$ and $|\mathsf s|\ge|\mathsf t|$; as $f\ge0$ is increasing, the second assertion follows. \end{proof}

\begin{lemma}[A weighted bound for families and a crude bound on the circle]\label{Y-C-CIIL65} Let $\mathfrak S$ be a coercive family, let $\mathsf V'$ be an $\operatorname{SL}(2,\mathbb C)$ field such that $|\mathsf V'(c)|\le1+\vartheta_0^\flat$ for every bond $c$ on which $\mathfrak S$ depends, and let $w$ be a real function on $I_{\mathfrak S}$ with $|w(b)-w(b')|\le\theta_{\mathbb C}{\rm dist}(b,b')$.

\begin{enumerate}\item[\textup{(a)}] $\|e^wK_{\mathfrak S}^{\mathsf V}e^{-w}\|\le\varkappa$; in fact this norm is at most $1-\frac58\frac\gamma b$, so that $\|(q_{\mathfrak S}^{\mathsf V})^{-1}\|\le\frac8{5\gamma}$. Consequently $q_{\mathfrak S}^{\mathsf V}$ is invertible, with $\|e^w(q_{\mathfrak S}^{\mathsf V})^{-1}e^{-w}\|\le\frac2\gamma$, $\|e^wq_{\mathfrak S}^{\mathsf V}e^{-w}\|\le2b$ and $\|(q_{\mathfrak S}^{\mathsf V})^{-1}(b,b')\|\le\frac2\gamma e^{-\theta_{\mathbb C}{\rm dist}(b,b')}$. At $G$-valued fields, $0\le K_{\mathfrak S}\le1-\frac\gamma b$.

\item[\textup{(b)}] For $|\zeta|\le e^{\theta_2}$ we have $\|K_{\mathfrak S}^{\mathsf V}(\zeta)\|\le\Lambda$; hence, in the polynomial $K_{\mathfrak S}^{\mathsf V}(\zeta)^k$, the coefficient of $\zeta^j$ has norm at most $\Lambda^ke^{-\theta_2j}$. \end{enumerate}\end{lemma}

\begin{proof}
The family satisfies the hypotheses of \Cref{Y-C-CIL65}, which
supplies the real bounds, the sharp weighted contraction and the
coefficient estimate of part~(b). Only the bonds in the transport
paths matter, so the field can be extended by the identity off them.

Write $q=q_{\mathfrak S}^{\mathsf V}$ and $K=K_{\mathfrak S}^{\mathsf V}$. Then $q=b(1-K)$, and the Neumann series $q^{-1}=b^{-1}\sum_kK^k$ converges after conjugation with $e^w$, with $\|e^wq^{-1}e^{-w}\|\le b^{-1}(1-\varkappa)^{-1}=\frac2\gamma$; at $w=0$ the sharper bound, used in place of $\varkappa$, gives $\|q^{-1}\|\le\frac8{5\gamma}$. From $q=b(1-K)$ we also obtain $\|e^wqe^{-w}\|\le2b$. For the entries we take $w=\theta_{\mathbb C}{\rm dist}(\cdot,b')$, which satisfies the hypothesis on $w$ by the triangle inequality and vanishes at $b'$; the bound on $e^wq^{-1}e^{-w}$ then gives the bound on the entry at $(b,b')$.

\end{proof}

\begin{lemma}[From the law to a window]\label{Y-C-L66} Let a law satisfy \textup{($\lambda$0)} and \textup{($\rho$5)}, and let $\frac2\gamma x^{-\mathsf n}\le\frac12$ and $\mathsf V'\in\mathfrak V$. Clause \textup{(c)} rests in addition on the standing conditions $2\le\mathsf n\le127$ (item~C.II~2.1) and $c_R\theta_{\mathbb C}\ge128$ (item~C.II~1.3).

\begin{enumerate}\item[\textup{(a)}] $\|C^{\mathsf V}\|\le\frac4\gamma$ and $\|C^{\mathsf V}-(q_\ell^{\mathsf V})^{-1}\|\le8\gamma^{-2}x^{-\mathsf n}$. Moreover $\mathsf u^{\mathsf V}=C^{\mathsf V}S^{\mathsf V}$, with $\|\mathsf u^{\mathsf V}\|\le\frac4\gamma x^{-\mathsf n}$ and $\|\mathsf u^{\mathsf V}\|\le\frac23$, and $S^{\mathsf V}=q^{\mathsf V}_\ell\sum_{m\ge1}(\mathsf u^{\mathsf V})^m$.

\item[\textup{(b)}] Let $\mathsf Y\subset I^\circ$ and $b\in\mathsf Y$ be such that every bond of $I^\circ\setminus\mathsf Y$ lies at distance at least $\rho$ from $b$. Then, for all $b'\in\mathsf Y$, $\|(q_{\mathsf Y}^{\mathsf V})^{-1}(b,b')-(q_\ell^{\mathsf V})^{-1}(b,b')\|\le8b\gamma^{-2}e^{-\theta_{\mathbb C}\rho}$.

\item[\textup{(c)}] Let $o$ be a coarse site, and let $\mathsf Y_o$ be the set of those bonds of $I^\circ$ whose base lies in a block of the box of radius $R+1$ centred at $o$. For a bond $b$ with ${\rm dist}(b,o)\le1$ and for $b'\in\mathsf Y_o$, \[\|C^{\mathsf V}(b,b')-(q_{\mathsf Y_o}^{\mathsf V})^{-1}(b,b')\|\le\tau_1:=8\gamma^{-2}(1+b)\,x^{-\mathsf n}.\] If the law also satisfies \textup{($\rho$1)}, then $C^{\mathsf V}(b,b')=0$ for $b'\notin\mathsf Y_o$. Finally, $q_{\mathsf Y_o}^{\mathsf V}$ depends on $\mathsf V'$ only through its values on the bonds of the box of radius $R+\varrho_L+2$ centred at $o$, other bonds entering at most through products $T^{-1}T=1$. \end{enumerate}\end{lemma}

\begin{proof} \emph{Step (a).} We write $q=q_\ell^{\mathsf V}$, $C=C^{\mathsf V}$, $S=S^{\mathsf V}$ and $\mathsf u=\mathsf u^{\mathsf V}$. The family of all steps is coercive by \Cref{Y-C-CIIL64}\,(a), and $\mathsf V'\in\mathfrak V$; so \Cref{Y-C-CIIL65}\,(a), applied to this family with $w=0$, gives $\|q^{-1}\|\le\frac2\gamma$. By \textup{($\rho$5)} we have $C^{-1}=q(1+\mathsf x)$ with $\mathsf x=q^{-1}S$, and $\|\mathsf x\|\le\frac12$ by the hypothesis $\frac2\gamma x^{-\mathsf n}\le\frac12$. Hence $C=(1+\mathsf x)^{-1}q^{-1}$ and $\|C\|\le\frac4\gamma$. Further, $q^{-1}-C=\mathsf x(1+\mathsf x)^{-1}q^{-1}$, whose norm is at most $2\cdot\frac2\gamma x^{-\mathsf n}\cdot\frac2\gamma$; this is the second bound.

From $C(q+S)=1$ we get $1-Cq=CS$, that is, $\mathsf u=CS$, and $\|C\|\le\frac4\gamma$ gives $\|\mathsf u\|\le\frac4\gamma x^{-\mathsf n}$. Next, $Cq=1-\mathsf u$ yields $C^{-1}=q(1-\mathsf u)^{-1}$; as $S=C^{-1}-q$, this gives the series for $S^{\mathsf V}$ in \textup{(a)}, provided that series converges. It does: the sharper bound of \Cref{Y-C-CIIL65}\,(a) gives $\|K\|\le1-\frac58\frac\gamma b$, hence $\|q^{-1}\|\le\frac8{5\gamma}$ and $\|\mathsf x\|\le\frac8{5\gamma}x^{-\mathsf n}\le\frac25$. Since $\mathsf u=CS=(1+\mathsf x)^{-1}q^{-1}S=(1+\mathsf x)^{-1}\mathsf x$, we obtain $\|\mathsf u\|\le\frac{2/5}{3/5}=\frac23$.

\emph{Step (b).} Multiplying the identity $\mathbf 1_{\mathsf Y}q(\mathbf 1_{\mathsf Y}+\mathbf 1_{\mathsf Y^c})q^{-1}\mathbf 1_{\mathsf Y}=\mathbf 1_{\mathsf Y}$ from the left by $q_{\mathsf Y}^{-1}$, we obtain $q_{\mathsf Y}^{-1}-\mathbf 1_{\mathsf Y}q^{-1}\mathbf 1_{\mathsf Y}=\mathsf M\,\mathbf 1_{\mathsf Y^c}q^{-1}\mathbf 1_{\mathsf Y}$, where $\mathsf M:=q_{\mathsf Y}^{-1}\mathbf 1_{\mathsf Y}q$. Let $w=\theta_{\mathbb C}{\rm dist}(b,\cdot)$. Then $\mathsf M(b,b_2)=e^{-w(b_2)}(e^{-w}\mathsf Me^w)(b,b_2)$, and \Cref{Y-C-CIIL65}\,(a), used twice with the weight $-w$, once for the window $\mathsf Y$ and once for all steps, gives $\|e^{-w}\mathsf Me^w\|\le\frac2\gamma\cdot2b$. Off $\mathsf Y$ we have $w\ge\theta_{\mathbb C}\rho$, so the row $\delta_b^{\mathsf T}\mathsf M\mathbf 1_{\mathsf Y^c}$ has $\ell^2$ norm at most $\frac{4b}\gamma e^{-\theta_{\mathbb C}\rho}$. Multiplied by $\|q^{-1}\|\le\frac2\gamma$, this gives \textup{(b)}.

\emph{Step (c).} A bond off $\mathsf Y_o$ lies at distance at least $R+2$ from $o$, hence at least $R+1$ from $b$. As $e^{-\theta_{\mathbb C}(R+1)}\le x^{-128}\le x^{-\mathsf n}$, the bounds of \textup{(a)} and of \textup{(b)}, the latter taken with $\mathsf Y=\mathsf Y_o$ and $\rho=R+1$, add up to $\tau_1$. For $b'\notin\mathsf Y_o$ we get ${\rm dist}(b,b')>R$, and \textup{($\rho$1)} gives $C^{\mathsf V}(b,b')=0$.

The dependence on $\mathsf V'$ comes from the steps as described in item~C.II~1.2. Let $\mathsf s:b_1\to b_2$ be a step between bonds of $\mathsf Y_o$. Such a step of kind \textup{(P)} or \textup{(Q)} with $y(b_1)=y(b_2)$ has $\mathsf T_{\mathsf s}=1$. Indeed, by item~C.II~1.2, for these two kinds $\mathsf T_{\mathsf s}$ is the inverse of the transport from one site to $y(b_1)$ times the transport from the same site to $y(b_2)$, and at $y(b_1)=y(b_2)$ the two transports are equal. These $\mathsf T_{\mathsf s}=1$ are the products $T^{-1}T=1$ in the last assertion of \textup{(c)}. In every case the sites of $\pi_{\mathsf s}$ are within $\varrho_L+1$ of an end of the step, by item~C.II~1.2, and so within $R+\varrho_L+2$ of $o$. \end{proof}

\begin{lemma}[The comb gauge of a box, paths, and the Bianchi identity]\label{Y-C-L67} Let $\mathsf B$ be a box and $o\in\mathsf B$ a site such that every site of $\mathsf B$ lies within $\rho_{\mathsf B}$ of $o$, where $3\rho_{\mathsf B}+1\le\ell_w(x')$, and let $\mathsf V'$ be a field on $\mathsf B$ that satisfies \textup{(h2)} and \textup{(h3)} at $\hat\varepsilon\le\frac1{40}$. In this lemma and in its proof, coordinates of coarse sites, $|\cdot|_\infty$ and $e_\lambda$ are given in coarse units.

\begin{enumerate}\item[\textup{(a)}] Let $c=(w,\rho)$ be a bond of $\mathsf B$, and let $t=|w-o|_\infty$. Then $|\mathsf W(c)-\mathbf 1|\le e^{18t\hat\varepsilon}-1$ and $|\mathsf W(c)|\le e^{18t\hat\varepsilon}$. Moreover $|\mathsf k(y)|\le2$ and $\|\operatorname{Ad}_{\mathsf k(y)}^{\pm1}\|\le4$ for the sites $y$ of $\mathsf B$.

\item[\textup{(b)}] For a path $\pi$ in $\mathsf B$ of $m\le\ell_w(x')$ bonds whose sites all lie within $\tau$ of $o$, $\|\operatorname{Ad}_{\mathsf W(\pi)}^{\pm1}-1\|\le2340\,m\,\tau\,\hat\varepsilon$.

\item[\textup{(c)}] For every plaquette $P$ of $\mathsf B$, $|\tilde X_P|\le8.24\,\hat\varepsilon$ and $e^{i\tilde X_P}=\partial\mathsf W(P)$.

\item[\textup{(d)}] Let $q'\subset\mathsf B$ be a cube at $z$, let $t=|z-o|_\infty$, and let $(\mathrm{d}'\tilde X)_{q'}=\sum_P\mathrm{d}'(q',P)\tilde X_P$. Then $\|(\mathrm{d}'\tilde X)_{q'}\|\le6\cdot10^4\,(1+t)\,\hat\varepsilon^2$. \end{enumerate}

None of these bounds requires $\rho_{\mathsf B}\hat\varepsilon$ to be small. Of \textup{(h2)}, clause \textup{(a)} needs only the bounds for the combs $\Gamma(o,y)$, $y\in\mathsf B$; since a comb has at most $3\rho_{\mathsf B}$ bonds, \textup{(a)} is valid already under $3\rho_{\mathsf B}\le\ell_w(x')$, and the additional bond is used in \textup{(d)}. \end{lemma}

\begin{proof} \emph{Step (a).} By definition, $\mathsf W(c)=\mathsf V'(\Gamma(o,c_-)\,c\,\Gamma(o,c_+)^{-1})$. If $\rho=3$, the two combs differ by the bond $c$ alone, and $\mathsf W(c)=\mathbf 1$. If $\rho=2$, the loop is a single ladder of $|w_3-o_3|$ plaquettes; if $\rho=1$, it consists of two ladders, with $|w_3-o_3|+|w_2-o_2|\le2t$ plaquettes in all.

A ladder is described by the following data: a path $\sigma$ from $o$; steps $u_0,\dots,u_{n-1}$, each equal to $\pm e_\lambda$, that start from the end of $\sigma$; the paths $\tau_j=u_0\cdots u_{j-1}$; the bond $h_j$ of direction $\rho$ at the end of $\sigma\tau_j$; the step $v_j$, which is $u_j$ shifted by $e_\rho$; and the loop $\mathsf p_j=u_jh_{j+1}v_j^{-1}h_j^{-1}$ along the boundary of one plaquette. For these data \[\tau_nh_nv_{n-1}^{-1}\cdots v_0^{-1}h_0^{-1}=\prod_{j=n-1}^{0}\tau_j\mathsf p_j\tau_j^{-1}.\] Indeed, $\tau_j\mathsf p_j\tau_j^{-1}=\tau_{j+1}h_{j+1}v_j^{-1}h_j^{-1}\tau_j^{-1}$, and the product telescopes.

The identity holds between words and requires of the $u_j,h_j,v_j$ nothing but invertibility. Hence it holds in any group, and also when paths of a fine field take the place of the $h_j$; in this form it enters Step~3, \textup{(3)}, of the proof of \Cref{Y-F-C7}.

Here $\sigma=\Gamma(o,(w_1,w_2,o_3))$ and $\lambda=3$; for $\rho=1$ the second ladder has $\sigma=\Gamma(o,(w_1,o_2,o_3))$ and $\lambda=2$. So every $\sigma\tau_j$ is a comb $\Gamma(o,\cdot)$. Thus each plaquette of a ladder contributes the factor $\mathsf k(y)\,\partial\mathsf V'(P)^{\pm1}\mathsf k(y)^{-1}$, in which $\mathsf k(y)$ is a comb. When the step is $-e_\lambda$, the plaquette $P$ is based at the endpoint of that step, and the comb leading to this endpoint is the conjugating one.

By \textup{(h2)}, $|\mathsf k|\le2$; hence $\|\operatorname{Ad}_{\mathsf k(y)}^{\pm1}\|\le4$ by item~C.II~1.1. For a plaquette factor $\partial=\partial\mathsf V'(P)$ we have $|\partial^{-1}-\mathbf 1|\le1.05\cdot2\hat\varepsilon$: this follows from $|g^{-1}-\mathbf 1|\le|g^{-1}||g-\mathbf 1|$ and $|g^{-1}|=|g|\le1+2\hat\varepsilon\le1.05$ at $g=\partial$, valid on $\operatorname{SL}(2,\mathbb C)$ by item~C.II~1.1 and \textup{(h3)}. With $|\mathsf k|\le2$, each factor is therefore of the form $\mathbf 1+\delta$ with $|\delta|\le8.4\hat\varepsilon$, and $(1+8.4\hat\varepsilon)^{2t}\le e^{16.8t\hat\varepsilon}$; since $16.8\le18$, the two bounds on $\mathsf W(c)$ follow.

\emph{Step (b).} By \textup{(a)}, $|\mathsf W(c)^{-1}-\mathbf 1|\le|\mathsf W(c)||\mathsf W(c)-\mathbf 1|\le e^{36t\hat\varepsilon}-1$. Thus $g=\mathsf W(\pi)$ is a product of $m$ factors $\mathbf 1+h_i$ with $1+|h_i|\le e^{36\tau\hat\varepsilon}$, and with $\varsigma:=36m\tau\hat\varepsilon$ we get $|g|\le e^\varsigma$ and $|g-\mathbf 1|\le e^\varsigma-1$. On the other hand, if $\pi$ runs from $y$ to $y''$, then $|g|=|\mathsf k(y)\mathsf V'(\pi)\mathsf k(y'')^{-1}|\le8$ by \textup{(h2)}. With the bounds on $\operatorname{Ad}_g$ of item~C.II~1.1 we conclude that $\|\operatorname{Ad}_g^{\pm1}-1\|\le\min\{65,\ 2e^\varsigma(e^\varsigma-1)\}\le65\varsigma$. For the last inequality, take first $\varsigma\le1$: the function $2e^\varsigma(e^\varsigma-1)$ is convex, is $0$ at $0$ and is below $9.35$ at $1$, hence it is at most $9.35\varsigma$; for $\varsigma>1$ we have $65\le65\varsigma$. Since $65\cdot36=2340$, this is \textup{(b)}.

\emph{Step (c).} Since $|\mathsf k|\le2$ and $|X'_P|\le2.06\hat\varepsilon$, we have $|\mathsf kX'\mathsf k^{-1}|\le4\cdot2.06\hat\varepsilon$, which is the bound on $\tilde X_P$. The exponential commutes with conjugation, $e^{iX'_P}=\partial\mathsf V'(P)$, and $\partial\mathsf W(P)=\mathsf k(y_P)\partial\mathsf V'(P)\mathsf k(y_P)^{-1}$; hence $e^{i\tilde X_P}=\partial\mathsf W(P)$.

\emph{Step (d).} We use Bianchi's identity, which is an identity of invertible matrices. Take the six monotone paths from $z$ to $z+e_1+e_2+e_3$ in the cyclic order $123,132,312,321,231,213$; any two neighbours in this order differ by one exchange. The quotients of neighbouring paths are then the six matrices $\kappa_1\partial(z+e_1;23)\kappa_1^{-1}$, $\partial(z;13)$, $\kappa_3\partial(z+e_3;12)\kappa_3^{-1}$, $\partial(z;23)^{-1}$, $\kappa_2\partial(z+e_2;13)^{-1}\kappa_2^{-1}$ and $\partial(z;12)^{-1}$, where $\kappa_\lambda$ denotes the bond $(z,\lambda)$; the product of these six is $\mathbf 1$.

For the field $\mathsf W$, by \textup{(c)}, this means $\prod_{k=1}^6e^{iY_k}=\mathbf 1$, where $Y_k=\epsilon_k\operatorname{Ad}_{\kappa_k}\tilde X_{P_k}$ and $\epsilon_k=\mathrm{d}'(q',P_k)$, with $\kappa_k=\mathbf 1$ for the faces based at $z$ and $\kappa_k=\mathsf W(z,\lambda)$ for the face based at $z+e_\lambda$. Now $\operatorname{Ad}_{\mathsf W(z,\lambda)}\tilde X_P=\operatorname{Ad}_{\mathsf k(z)\mathsf V'(z,\lambda)}X'_P$, and $|\mathsf k(z)\mathsf V'(z,\lambda)|\le2$ by \textup{(h2)}, applied to the comb $\Gamma(o,z)$ followed by the bond $(z,\lambda)$, a path of at most $3\rho_{\mathsf B}+1$ bonds. Therefore $|Y_k|\le8.24\hat\varepsilon$ for every $k$, and the sum $S$ of \textup{(E)} satisfies $S\le49.44\hat\varepsilon\le1.236$.

As the product is $\mathbf 1$, \textup{(E)} gives $|\sum Y_k|\le e^S-1-S\le0.7894\,S^2\le1929.6\,\hat\varepsilon^2\le1931\,\hat\varepsilon^2$. Here we used that $(e^S-1-S)/S^2$ increases with $S$ and is $\le0.78931$ at $S=1.236$, and that $0.7894\cdot49.44^2\le1929.55$. Finally, $(\mathrm{d}'\tilde X)_{q'}-\sum Y_k=\sum_\lambda\pm(1-\operatorname{Ad}_{\mathsf W(z,\lambda)})\tilde X_{P_\lambda}$; by \textup{(b)}, with $m=1$ and $\tau=t+1$, and by \textup{(c)}, its norm is at most $3\cdot2340(t+1)\hat\varepsilon\cdot8.24\hat\varepsilon\le57845(t+1)\hat\varepsilon^2$. As $1931+57845\le6\cdot10^4$, the two bounds together give \textup{(d)}. \end{proof}

\subsubsection*{C.II 4. \textup{(S-cen$_k$)}}

\begin{definition}[Constants of the centre] Throughout item C.II 4 we write $\theta=\theta_{\mathbb C}$. We define \[\FitDisplay{\dimexpr\linewidth-1pt\relax}{\mathsf K_A:=1.24\cdot10^6\,\frac{A^\flat}{\gamma\,\theta}\,\mathsf g_I(\tfrac\theta2)\Bigl[\frac{8\mathsf N_{\mathsf B}}\gamma\mathsf g_I(\theta)+1\Bigr],\qquad\mathsf K_{\rm c}^0:=(L+1)\mathsf K_A+6\cdot10^4\,L\,C_D\,\theta_D^{-1},\qquad\mathsf K_{\rm c}^1:=289\,\mathsf K_{\rm c}^0+17,}\] together with \[\FitDisplay{\dimexpr\linewidth-1pt\relax}{\mathsf a_{\rm c}:=\frac{100}\gamma\mathsf N_{\mathsf B}\,\mathsf g_I(\theta),\qquad\tau_{\rm c}:=165\,(L+1)\,\mathsf N_{\mathsf B}\cdot3L^3(2R+3)^3\,\tau_1 .}\] Each of these constants except $\tau_{\rm c}$ is computed from $(L,a,K,\lambda,C_D,\theta_D)$; as $K_*,\lambda_*$ are pure numbers, in the case $K=K_*$ and $\lambda=\lambda_*$ each such constant depends on $(L,a,C_D,\theta_D)$ and on pure numbers alone. \end{definition}

\begin{definition}[Floors on $x$ and on the torus] The hypotheses \textup{(P0)} to \textup{(P2)} separate the coupling
bounds from the torus-size requirements.

\begin{enumerate}\item[\textup{(P0)}] $x\ge L^2$, $R\ge6$, and \textup{(Tor)}.

The condition \textup{(Tor)} of item~C.II~1.1 is a datum of \textup{(SD)} (item~1 of \S\ref{Y-G0-1}). It is no floor on $x$, since for fixed $N$ it is violated once $x$ is large.

From \textup{(P0)} we obtain $3R_3+1\le7R+19\le14c_R\log x'+26\le\ell_w(x')$, $|\mathsf s|,|\mathsf t|\le8+6\varrho/L\le\ell_w(x')$, $r_q\le R+2$ and $2R_3+1<N$.

\item[\textup{(P1)}] $\frac4\gamma x^{-\mathsf n}\le1$, and $165(L+1)\mathsf N_{\mathsf B}\,3L^3(2R+3)^3\,8\gamma^{-2}(1+b)\le x$; by the latter, $\tau_{\rm c}\le x^{-(\mathsf n-1)}$.

\item[\textup{(P2)}] $10^4\,\ell_w(x')\,R_3^2\,\tau_{\rm c}\log x\le\vartheta_0^\flat$, and $2(R_3+\lceil1/\vartheta_0^\flat\rceil+1)+1<N$. \end{enumerate}

With the letters fixed, every clause that is free of $N$ holds once $x$ is large enough. Only \textup{(Tor)} and the second clause of \textup{(P2)} involve $N$; at the given $x$ they are conditions on $N$, and \textup{(Tor)} entails $N\ge320c_R\log x$. \end{definition}

\begin{definition}[The remainders $\Lambda_i$] Let $i=p$ be a fine plaquette or $i=c$ a coarse bond; put $o:=y(p)$, resp.\ $o:=c_-$, and let $\mathsf B(i)$ be the box of radius $R_3$ centred at $o$. The remainders are \[\FitDisplay{\dimexpr\linewidth-1pt\relax}{\Lambda_p(\mathsf V'):=\bigl(EX'-\mathrm{d}^{\mathsf V}C^{\mathsf V}\mathsf g^{\mathsf V}X'\bigr)_p-\sum_PH^\natural(p,P)T(y(p),y_P)X'_P,\qquad\Lambda_c(\mathsf V'):=(Q_1^{\mathsf V}C^{\mathsf V}\mathsf g^{\mathsf V}X')_c-\sum_P\mathsf N^\natural(c,P)T(c_-,y_P)X'_P .}\] \end{definition}

\begin{remark} Each coarse site $o$ is the end $c_-$ of some coarse bond $c$; hence \Cref{Y-C-Th68}, applied with $i=c$ for each coarse bond $c$, covers all coarse sites $o$, and in this form Sections~\ref{Y-E-1} and~\ref{Y-G-1} use its clauses \textup{(a)} and \textup{(c)}. \end{remark}

\begin{theorem}[The remainders: a quadratic bound with a small linear term]\label{Y-C-Th68} Let $L\ge3$ be odd with $aL\ge18$, let the rows be as in item~C.II~1.2, let a law with \textup{($\lambda$0)}, \textup{($\rho$1)}, \textup{($\rho$2)} and \textup{($\rho$5)} be given, let $\mathsf G,\mathsf G_Z$ satisfy \textup{(Fr2)}, and let \textup{(P0)} and \textup{(P1)} hold. Then, for every $i$, every $0<\hat\varepsilon\le\frac1{40}$ and every $\operatorname{SL}(2,\mathbb C)$ field $\mathsf V'$ on $\mathsf B(i)$ with \textup{(h)} at $\hat\varepsilon$, the following hold.
\begin{enumerate}
\item[\textup{(a)}] $\Lambda_i(\mathsf V')$, and $a_{0,b}$ at the bonds with ${\rm dist}(b,o)\le1$, depend on $\mathsf V'$ only through its values on the bonds of $\mathsf B(i)$, and on $X'$ only through its values at the plaquettes of $\mathsf B(i)$. Moreover, independently of \textup{(h2)}, \textup{(h3)} and \textup{(Fr2)}: whenever two pairs, each consisting of a field of $\mathfrak V$ and a $\mathfrak g_{\mathbb C}$-valued function $X'$ on the plaquettes of $\mathbb T'$, coincide on the bonds and on the plaquettes of $\mathsf B(i)$, the formulas of $\Lambda_i$ and of $a_{0,b}=(C^{\mathsf V}\mathsf g^{\mathsf V}X')_b$ take the same values at both pairs. Consequently, when the field is given on a larger set and $X'$ is defined on the plaquettes of that set, the values obtained are those computed on $\mathsf B(i)$.
\item[\textup{(b)}] $\|\Lambda_i(\mathsf V')\|\le\mathsf K_{\rm c}^0\,\hat\varepsilon^2+\tau_{\rm c}\,\hat\varepsilon$.
\item[\textup{(c)}] $|a_{0,b}|\le\sqrt2\|a_{0,b}\|\le\mathsf a_{\rm c}\hat\varepsilon$ at the bonds with ${\rm dist}(b,o)\le1$.
\end{enumerate}

(By item~C.II~1.4, every field with \textup{(e1)}, \textup{(e2)} and $|\partial\mathsf V'(P)-\mathbf 1|\le2\hat\varepsilon$ has \textup{(h)} at $\hat\varepsilon$. A $G$-valued field always satisfies \textup{(h1)} and \textup{(h2)}, and it satisfies \textup{(h3)} if and only if $\sin\frac12\phi'_P\le\hat\varepsilon$ at the plaquettes of $\mathsf B(i)$.)
\end{theorem}

\begin{proof} We take $i=p$ and write $\mathsf B=\mathsf B(p)$. We extend $\mathsf V'$ by $\mathbf 1$; by \textup{(h1)}, the extended field belongs to $\mathfrak V$. Let $\mathsf Y:=\mathsf Y_o$ be the set of \Cref{Y-C-L66}\,(c); its size is $\#\mathsf Y\le3L^3(2R+3)^3$. We put $t_j:={\rm dist}(j,o)$.

\emph{Step (0): on what the objects depend.} Let $b$ be a bond of $p$, or of a path $[x,x']$ with $x\in B(c_-)$. Then $y(b)\in\{o\}\cup\{o+Le_\mu\}$; therefore ${\rm dist}(b,o)\le1$, and $T(o,y(b))=\operatorname{Ad}_{\mathsf k(y(b))}$, because the path $\Gamma(o,y(b))$ is either empty or a single bond. By \textup{($\rho$1)}, $a_{0,b}=\sum_{j\in\mathsf Y}C^{\mathsf V}(b,j)\mathsf g_j$, where $\mathsf g_j:=(\mathsf g^{\mathsf V}X')_j$. By \textup{($\rho$2)}, $C^{\mathsf V}(b,\cdot)$ depends on $\mathsf V'$ only within $R_2\le R_3$ of $o$. For $j\in\mathsf Y$, the term $\mathsf g_j$ depends only on paths $\pi_{\mathsf t}$ and plaquettes $P$ all of whose sites lie within $(R+1)+2\varrho_L+2\le R_3$ of $o$; this inequality holds as $\varrho_L\le R+2$.

By \Cref{Y-C-L71}\,(a), $H^\natural(p,P)\ne0$ only for ${\rm dist}(p,P)\le R_2$. The proof of clause (a) of that lemma rests on \textup{($\rho$1)}, $R\ge6$ and $L\ge3$ alone and is independent of the present theorem. All four sites of such a plaquette $P$ lie within $R_2+1\le R_3$ of $o$, and the path $\Gamma(o,y_P)$ is contained in $\mathsf B$; further, $(EX')_p$ equals $X'_P$ for a plaquette $P$ with $y_P=o$, and $0$ otherwise. This proves (a): none of these objects sees the extension of $\mathsf V'$.

Step (0) also proves the second sentence of (a), because it is independent of the sizes of $\mathsf V'(\pi)$, of $\partial\mathsf V'(P)$ and of $X'$, and of the formula that defines $X'$: the properties \textup{($\lambda$0)}, \textup{($\rho$1)}, \textup{($\rho$2)} hold at every field of $\mathfrak V$, and the ranges of $\mathrm{d}$, $Q_1$, $\mathsf g$, $H^\natural$, $\mathsf N^\natural$ are the same for every field.

\emph{Step (1): sizes in the given gauge.} For $j\in\mathsf Y$ we have $\|\mathsf T_{\mathsf t}\|\le|\mathsf V'(\pi_{\mathsf t})|^2\le4$ by \textup{(h2)}; with \Cref{Y-C-CIIL64}\,(d) this gives $\|\mathsf g_j\|\le4\mathsf N_{\mathsf B}\cdot4\cdot2.06\hat\varepsilon\le33\mathsf N_{\mathsf B}\hat\varepsilon$.

\emph{Step (2): from the law to the window.} Let $\hat a:=(q^{\mathsf V}_{\mathsf Y})^{-1}\mathbf 1_{\mathsf Y}\mathsf g$, where $\mathsf g$ denotes the vector $(\mathsf g_j)_j$ of Step (0). \Cref{Y-C-L66}\,(c) yields $\|a_{0,b}-\hat a_b\|\le\tau_1\#\mathsf Y\cdot33\mathsf N_{\mathsf B}\hat\varepsilon$, and together with \Cref{Y-C-CIIL65}\,(a) we obtain \[\|a_{0,b}\|\le[\frac2\gamma\mathsf g_I(\theta)+\tau_1\#\mathsf Y]33\mathsf N_{\mathsf B}\hat\varepsilon\le\frac{70}\gamma\mathsf N_{\mathsf B}\mathsf g_I(\theta)\hat\varepsilon,\] the second inequality by \textup{(P1)}. Since $\sqrt2\cdot70\le100$, this is (c).

\emph{Step (3): the gauge.} Let $\hat{\mathsf G}:={\rm diag}(\operatorname{Ad}_{\mathsf k(y(j))})_{j\in\mathsf Y}$. For a path $\pi$ in $\mathsf B$ from $y$ to $y''$, the relation $\operatorname{Ad}_{\mathsf k(y)}\operatorname{Ad}_{\mathsf V'(\pi)}\operatorname{Ad}_{\mathsf k(y'')}^{-1}=\operatorname{Ad}_{\mathsf W(\pi)}$ is an identity of matrices. Hence $q^{\mathsf W}:=\hat{\mathsf G}q_{\mathsf Y}^{\mathsf V}\hat{\mathsf G}^{-1}$ has the entries $\sum q_{\mathsf s}\operatorname{Ad}_{\mathsf W(\pi_{\mathsf s})}$, where the steps with $\mathsf T_{\mathsf s}=1$ and $y(b)=y(b')$ keep the factor $1$. Also $\hat{\mathsf G}\mathbf 1_{\mathsf Y}\mathsf g=\mathsf g^{\mathsf W}\tilde X$, with $\mathsf g^{\mathsf W}(j,P):=\sum g_{\mathsf t}\operatorname{Ad}_{\mathsf W(\pi_{\mathsf t})}$. Put $\mathsf a^{\mathsf W}:=(q^{\mathsf W})^{-1}\mathsf g^{\mathsf W}\tilde X=\hat{\mathsf G}\hat a$; then Step (0) gives \[(\mathrm{d}^{\mathsf V}\hat a)_p=\sum_b\mathrm{d}(p,b)\operatorname{Ad}_{\mathsf k(y(b))}\hat a_b=(d\,\mathsf a^{\mathsf W})_p.\]

Since $\mathsf k(o)=\mathbf 1$ and an edge plaquette $p$ sits under some $P$ with $y_P=o$, we have $(EX')_p=(E\tilde X)_p$. From now on, $\mathsf g$ once more denotes the kernel of item~C.II~2.2, that is, $\mathsf g^{\mathsf V}$ at $\mathsf V'=\mathbf 1$. By the definitions of $\tilde X$ and of $H^\natural$, \[\sum_PH^\natural(p,P)T(o,y_P)X'_P=(H^\natural\tilde X)_p=(E\tilde X)_p-(\mathrm{d}C^{\mathbf 1}\mathsf g\tilde X)_p-\sum_{\mathcal Q_p}\mathsf G(p,q')(\mathrm{d}'\tilde X)_{q'}.\] Let $q^{\mathbf 1}:=q_{\mathsf Y}^{\mathbf 1}$ and $\mathsf a^{\mathbf 1}:=(q^{\mathbf 1})^{-1}\mathbf 1_{\mathsf Y}\mathsf g\tilde X$. \Cref{Y-C-L66}\,(c), applied with $\mathbf 1$ as its field $\mathsf V'$ and used with $\|(\mathsf g\tilde X)_j\|\le4\mathsf N_{\mathsf B}\cdot8.24\hat\varepsilon\le33\mathsf N_{\mathsf B}\hat\varepsilon$, gives $\|(C^{\mathbf 1}\mathsf g\tilde X)_b-\mathsf a^{\mathbf 1}_b\|\le\tau_1\#\mathsf Y\,33\mathsf N_{\mathsf B}\hat\varepsilon$.

Inserting all this into the definition of $\Lambda_p$, we find \[\Lambda_p=-\bigl(d[\mathsf a^{\mathsf W}-\mathsf a^{\mathbf 1}]\bigr)_p+\sum_{q'\in\mathcal Q_p}\mathsf G(p,q')(\mathrm{d}'\tilde X)_{q'}+\mathsf{ti},\qquad\|\mathsf{ti}\|\le4\,(4+1)\,33\,\mathsf N_{\mathsf B}\tau_1\#\mathsf Y\,\hat\varepsilon .\] Here the factor $4$ inside $(4+1)$ enters because the error of Step (2) is transported by $\operatorname{Ad}_{\mathsf k}$, whose norm is at most $4$. When $i=c$, the row of $Q_1$ sums to $L$; in either case $\|\mathsf{ti}\|\le\tau_{\rm c}\hat\varepsilon$.

\emph{Step (4): the resolvent identity, row by row.} We write $\mathsf a^{\mathsf W}-\mathsf a^{\mathbf 1}={\rm A}+{\rm B}$, where ${\rm A}:=-(q^{\mathsf W})^{-1}\delta\,(q^{\mathbf 1})^{-1}\mathsf g\tilde X$ with $\delta:=q^{\mathsf W}-q^{\mathbf 1}$, and ${\rm B}:=(q^{\mathsf W})^{-1}(\mathsf g^{\mathsf W}-\mathsf g)\tilde X$. The entries obey $\|(q^{\mathsf W})^{-1}(b,j)\|\le16\cdot\frac2\gamma e^{-\theta{\rm dist}(b,j)}$; this comes from \Cref{Y-C-CIIL65}\,(a) for $q_{\mathsf Y}^{\mathsf V}$, which applies by \textup{(h1)} and independently of any closeness to $q^{\mathbf 1}$, while the conjugation costs the factor $4\cdot4$. Moreover $|(q^{\mathbf 1})^{-1}(j,j')|\le\frac2\gamma e^{-\theta{\rm dist}}$.

All sites of a step from $j$ lie within $t_j+|\mathsf s|$ of $o$. Therefore \Cref{Y-C-L67}\,(b) and \Cref{Y-C-CIIL64}\,(c) give \[\sum_{j'}\|\delta(j,j')\|\le\sum_{\mathsf s\text{ from }j}|q_{\mathsf s}|\,2340\,|\mathsf s|(t_j+|\mathsf s|)\hat\varepsilon\le2340A^\flat\bigl(t_j+\tfrac2{\theta_1}\bigr)\hat\varepsilon .\] \Cref{Y-C-CIIL64}\,(d) gives the same bound for $\sum_P\|(\mathsf g^{\mathsf W}-\mathsf g)(j,P)\|$. If ${\rm dist}(b,o)\le1$, then $t_j\le1+{\rm dist}(b,j)$, and so $\sum_je^{-\theta{\rm dist}(b,j)}(t_j+\frac2{\theta_1})\le\mathsf g_I(\frac\theta2)[1+\frac2{\theta_1}+\frac1\theta]\le\frac2\theta\mathsf g_I(\frac\theta2)$. For the last inequality we use $\theta\le\frac16\theta_1\le\frac16$, by which the bracket is at most $\frac1\theta$ times $\frac16+\frac13+1$.

Consequently \[\|{\rm A}_b\|\le\tfrac{32}\gamma\cdot2340A^\flat\hat\varepsilon\,\tfrac2\theta\mathsf g_I(\tfrac\theta2)\cdot\tfrac2\gamma\mathsf g_I(\theta)\cdot4\mathsf N_{\mathsf B}\,8.24\hat\varepsilon,\qquad\|{\rm B}_b\|\le\tfrac{32}\gamma\cdot2340A^\flat\hat\varepsilon\,\tfrac2\theta\mathsf g_I(\tfrac\theta2)\cdot8.24\hat\varepsilon .\] Since $32\cdot2340\cdot2\cdot8.24=1234022.4$, it follows that $\|{\rm A}_b\|+\|{\rm B}_b\|\le\mathsf K_A\hat\varepsilon^2$. For the first term of $\Lambda_i$ we use that the row of $d$ has four entries, that the row of $Q_1$ sums to $L$, and that $4\le L+1$.

\emph{Step (5): the monopole sum.} The corners of a cube of $\mathcal Q_p$ lie within $R_2+1\le R_3$ of $o$. \Cref{Y-C-L67}\,(d), together with ${\rm dist}(q',o)={\rm dist}(p,q')$, $1+t\le\theta_D^{-1}e^{\theta_Dt}$ and \textup{(Fr2)}, bounds the norm of the monopole sum of Step (3) by $\sum_{q'}|\mathsf G(p,q')|\,6\cdot10^4(1+{\rm dist}(p,q'))\hat\varepsilon^2\le6\cdot10^4C_D\theta_D^{-1}\hat\varepsilon^2$. For $\mathsf G_{\mathsf n}$, which takes the place of $\mathsf G$ when $i=c$, the factor $C_D$ is replaced by $LC_D$. Together with Steps (3) and (4), this proves (b).
\end{proof}

\begin{lemma}[Covariance of a holomorphic law at complex pure gauges]\label{Y-C-L69} Let a law satisfy \textup{($\lambda$0)}, \textup{($\rho$2)}, \textup{($\rho$3)} and \textup{($\rho$4)}, and let \textup{(P0)} and the second clause of \textup{(P2)} hold. Let $\mathsf B$ be the box of radius $R_3$ about $o$, let $u:\mathsf B\to\operatorname{SL}(2,\mathbb C)$ satisfy $|u(y)|\le2$, and $|u(c_-)u(c_+)^{-1}|<1+\vartheta_0^\flat$ on the bonds of $\mathsf B$, and let $\mathsf V_0(c):=u(c_-)u(c_+)^{-1}$ on $\mathsf B$, extended by $\mathbf 1$. Then, for ${\rm dist}(b,o)\le1$ and for all $b'$ based in the blocks of $\mathsf B$, \[C^{\mathsf V_0}(b,b')=\operatorname{Ad}_{u(y(b))}\,C^{\mathbf 1}(b,b')\,\operatorname{Ad}_{u(y(b'))}^{-1}.\]
\end{lemma}

\begin{proof} \emph{Step (1): an extension of the gauge.} We write $u$ in polar form, $u(y)=k_0(y)e^{h(y)}$ with $k_0\in G$ and $h$ hermitian traceless; then $|e^{h(y)}|\le2$. Let $D:=\lceil1/\vartheta_0^\flat\rceil+1$; let $\mathsf B^+$ be the box of radius $R_3+D$ about $o$, which embeds by \textup{(P2)}; let $\varpi:\mathsf B^+\to\mathsf B$ be the map that, in each direction $\mu$, sends the coordinate $y_\mu$ of a site $y$ to the point of the interval of $\mathsf B$ nearest to $y_\mu$, so that $\varpi$ fixes the sites of $\mathsf B$; and let $\chi(y):=1-{\rm dist}(y,\mathsf B)/D$. For $0\le s\le1$ we define $u_s(y):=k_0(\varpi y)e^{s\chi(y)h(\varpi y)}$ on $\mathsf B^+$, and $u_s(y):=\mathbf 1$ off it.

We claim that $\mathbf 1^{u_s}\in\mathfrak V^\circ$. Take a lattice step in $\mathsf B^+$ in the direction $\mu$, from $y$ to $y'$. Either the step leaves $\varpi y$ unchanged, or it moves $\varpi y$ along a bond $c$ of $\mathsf B$ and leaves ${\rm dist}(y,\mathsf B)$ unchanged. The former happens when $y_\mu$ remains outside the interval of $\mathsf B$, when it arrives from outside at an end of this interval, or when it leaves the interval from one of its ends; then $\chi$ changes by at most $\frac1D$, and $|u_s(y)u_s(y')^{-1}|=|e^{s(\chi-\chi')h}|\le|e^h|^{1/D}\le2^{1/D}<1+\vartheta_0^\flat$, by $\frac{\log2}D<0.7\vartheta_0^\flat<\log(1+\vartheta_0^\flat)$. In the latter case, with $h_1$ and $h_2$ the values of $h$ at $c_-$ and at $c_+$, Cordes' inequality bounds the norm of $u_s(y)u_s(y')^{-1}$: $|e^{s\chi h_1}e^{-s\chi h_2}|\le|e^{h_1}e^{-h_2}|^{s\chi}=|u(c_-)u(c_+)^{-1}|^{s\chi}<1+\vartheta_0^\flat$.

Cordes' inequality says that $|A^\alpha B^\alpha|\le|AB|^\alpha$ for positive matrices $A,B$ and $0\le\alpha\le1$. For the proof, consider the set of those $\alpha$ for which $|A^\alpha B^\alpha|\le|AB|^\alpha$. It is closed, it contains $0$ and $1$, and together with two points $\alpha,\beta$ it contains their midpoint $m=\frac{\alpha+\beta}2$, because, with $\mathrm r$ the spectral radius, \[|A^mB^m|^2=\mathrm r(B^mA^{2m}B^m)=\mathrm r(A^{\alpha+\beta}B^{\alpha+\beta})=\mathrm r(B^\alpha A^\alpha A^\beta B^\beta)\le|A^\alpha B^\alpha||A^\beta B^\beta|.\] Hence this set is the whole interval $[0,1]$.

At the rim, where $\chi=0$, the gauge $u_s$ is unitary; so, on every bond from $y$ to $y'$ that leaves $\mathsf B^+$, the value $u_s(y)u_s(y')^{-1}$ of $\mathbf 1^{u_s}$ has norm $1$. This proves the claim.

\emph{Step (2): the derivative vanishes.} For $\mathsf Z\in\mathfrak V^\circ$ and $\lambda:\mathbb T'\to\mathfrak g_{\mathbb C}$ we set \[\mathfrak d(\mathsf Z;\lambda):=\frac d{d\sigma}\bigl[\operatorname{Ad}_{v_\sigma(y(b))}^{-1}C^{\mathsf Z^{v_\sigma}}(b,b')\operatorname{Ad}_{v_\sigma(y(b'))}\bigr]_{\sigma=0},\] where $v_\sigma=e^{i\sigma\lambda}$. This derivative is $\mathbb C$-linear in $\lambda$, since \textup{($\rho$4)} makes the bracket a holomorphic function of $v$ near $\mathbf 1$ and the tangent vector of the curve $v_\sigma$ is $\mathbb C$-linear in $\lambda$. If $\lambda$ is $\mathfrak g$-valued, the bracket is constant by \textup{($\rho$3)}. Every $\lambda$ can be written as $\lambda_1+i\lambda_2$ with $\mathfrak g$-valued $\lambda_1,\lambda_2$; hence $\mathfrak d\equiv0$.

Let $\Psi(s):=\operatorname{Ad}_{u_s(y(b))}^{-1}C^{\mathbf 1^{u_s}}(b,b')\operatorname{Ad}_{u_s(y(b'))}$. We have $\mathbf 1^{u_{s+\epsilon}}=(\mathbf 1^{u_s})^{u_{s+\epsilon}u_s^{-1}}$ and $u_{s+\epsilon}u_s^{-1}=\exp(\epsilon\dot u_su_s^{-1}+o(\epsilon))$, where $\dot u_su_s^{-1}$ is traceless. As $\mathbf 1^{u_s}\in\mathfrak V^\circ$ by Step (1), this gives $\Psi'(s)=\operatorname{Ad}_{u_s}^{-1}\mathfrak d(\mathbf 1^{u_s};-i\dot u_su_s^{-1})\operatorname{Ad}_{u_s}=0$. Moreover $\Psi(0)=C^{\mathbf 1}(b,b')$ by \textup{($\rho$3)}, because $u_0$ is $G$-valued.

\emph{Step (3): comparison on the box.} Since $u_1=u$ on $\mathsf B$, we have $\mathbf 1^{u_1}=\mathsf V_0$ on the bonds of $\mathsf B$. By \textup{($\rho$2)}, $C(b,\cdot)$ depends on the field only within $1+(R_2-1)\le R_3$ of $o$; hence $C^{\mathsf V_0}(b,b')=C^{\mathbf 1^{u_1}}(b,b')$. As $\Psi$ is constant by Step (2), this is the assertion.
\end{proof}

\begin{theorem}[The remainders without the linear term, \textup{(S-cen$_k$)}]\label{Y-C-Th70} Let $L\ge3$ be odd with $aL\ge18$, let the rows be as in item~C.II~1.2, in particular $\ell\in\mathcal R(b_*,K_*,\lambda_*,\frac19)$ of range $\le\varrho(x)$, let a law with \textup{($\lambda$0)}, \textup{($\rho$1)} to \textup{($\rho$5)} be given, let $\mathsf G,\mathsf G_Z$ satisfy \textup{(Fr2)}, and let \textup{(P0)}, \textup{(P1)} and \textup{(P2)} hold. Then, for every fine plaquette $p$ and every coarse bond $c$, for every $0<\hat\varepsilon\le\frac1{40}$, and for every $\operatorname{SL}(2,\mathbb C)$ field $\mathsf V'$ on the box of radius $R_3$ about $y(p)$, resp.\ $c_-$, that satisfies \textup{(e1)}, \textup{(e2)} and $|\partial\mathsf V'(P)-\mathbf 1|\le2\hat\varepsilon$, \[\|\Lambda_p(\mathsf V')\|\le\mathsf K_{\rm c}^1\hat\varepsilon^2,\qquad\|\Lambda_c(\mathsf V')\|\le\mathsf K_{\rm c}^1\hat\varepsilon^2,\qquad\|(C^{\mathsf V}\mathsf g^{\mathsf V}X')_b\|\le\mathsf a_{\rm c}\hat\varepsilon\ \ ({\rm dist}(b,o)\le1).\]
\end{theorem}

\begin{proof} The third bound is \Cref{Y-C-Th68}\,(c). If $\hat\varepsilon\ge\tau_{\rm c}$, the first two bounds follow from \Cref{Y-C-Th68}\,(b), because then $\tau_{\rm c}\hat\varepsilon\le\hat\varepsilon^2$. Now let $\hat\varepsilon<\tau_{\rm c}$, and let $T:=\tau_{\rm c}/\hat\varepsilon>1$. We take $i=p$; for $i=c$ the proof is the same, with $\Lambda_c$ in place of $\Lambda_p$ and the box about $c_-$ in place of the box about $y(p)$.

We first collect the numerical facts that Step (1) uses. The first clause of \textup{(P2)}, together with $\vartheta_0^\flat\le10^{-4}$, $\ell_w(x')\ge12$ and $\log x\ge1$, gives $R_3^2\tau_{\rm c}\le10^{-9}$; besides, we use $R_3\ge17$. With $\hat\varepsilon<\tau_{\rm c}$ and $T\hat\varepsilon=\tau_{\rm c}$, and with the number $\mathsf a$ of Step (1), these facts give: $18R_3\hat\varepsilon\le\frac1{10}$, so that $e^{18R_3\hat\varepsilon}-1\le19R_3\hat\varepsilon\le0.106$; $(T+1)\mathsf a<42R_3\tau_{\rm c}\le1$, so that $4(e^{(T+1)\mathsf a}-1)\le8(T+1)\mathsf a$; $4T\mathsf a=84R_3\tau_{\rm c}\le1$; $56448\,R_3^2\tau_{\rm c}\le2$; and $17\tau_{\rm c}\le\frac1{40}$.

\emph{Step (1): the curve, with the room of item~C.II~1.4.} By \Cref{Y-C-L67}\,(a), $|\mathsf W(c)-\mathbf 1|\le e^{18R_3\hat\varepsilon}-1\le19R_3\hat\varepsilon\le0.106$. Hence $A'_c:=-i\log\mathsf W(c)$ is traceless, with $|A'_c|\le\mathsf a:=21R_3\hat\varepsilon$. We set $\mathsf V'_t(c):=\mathsf k(c_-)^{-1}e^{itA'_c}\mathsf k(c_+)$; this is entire in $t$, and $\mathsf V'_1=\mathsf V'$. For $|t|<T$ we have $(T+1)\mathsf a<42R_3\tau_{\rm c}$, and $\mathsf V'_t(c)=\mathsf V'(c)(\mathbf 1+\delta_c)$ with $\mathbf 1+\delta_c=\mathsf k(c_+)^{-1}e^{i(t-1)A'_c}\mathsf k(c_+)$, where $|\delta_c|,|(\mathbf 1+\delta_c)^{-1}-\mathbf 1|\le4(e^{(T+1)\mathsf a}-1)\le\delta:=336R_3\tau_{\rm c}$.

We now check \textup{(h)} for $\mathsf V'_t$ with $|t|<T$. First, \textup{(h1)} holds strictly: $|\mathsf V'_t(c)|<(1+\delta)(1+\vartheta(x')+\mathsf g_x)\le1+\vartheta_0^\flat$, since $1.01\delta\le\frac{3\vartheta_0^\flat}{4\log x}$ by \textup{(P2)}. Next, \textup{(h2)}: in a path of $n\le\ell_w(x')$ bonds we carry the factors $(\mathbf 1+\delta_c)^{\pm1}$ to the left through the partial products, whose norms are $<2$; we obtain $|\mathsf V'_t(\pi)|\le(1+4\delta)^n|\mathsf V'(\pi)|<e^{\mathsf y}(2-\mathsf r)$ with $\mathsf y=4n\delta$ and $\mathsf r=\frac3{4\log x}$. Now $(2-\mathsf r)(1+2\mathsf y)\le2$ as soon as $4\mathsf y\le\mathsf r$, and the latter holds because $5376\,\ell_w(x')R_3\tau_{\rm c}\le\frac3{4\log x}$ by \textup{(P2)}.

Finally, \textup{(h3)}: $\partial\mathsf V'_t(P)=\mathsf k(y_P)^{-1}\Pi_t\mathsf k(y_P)$ with $\Pi_t=e^{itA'_1}e^{itA'_2}e^{-itA'_3}e^{-itA'_4}$, and \textup{(E)} gives $|\Pi_t-\mathbf 1-it(\mathrm{d}'A')_P|\le(4|t|\mathsf a)^2$, as $4T\mathsf a\le1$. At $t=1$ this yields $|(\mathrm{d}'A')_P|\le8\hat\varepsilon+16\mathsf a^2$. Therefore $|\partial\mathsf V'_t(P)-\mathbf 1|\le32T\hat\varepsilon+128T^2\mathsf a^2\le34T\hat\varepsilon$, by $128\cdot441\,TR_3^2\hat\varepsilon=56448\,R_3^2\tau_{\rm c}\le2$. Thus $\mathsf V'_t$ has \textup{(h)} at $\hat\varepsilon_T:=17T\hat\varepsilon=17\tau_{\rm c}\le\frac1{40}$.

\emph{Step (2): holomorphy and a first bound.} The function $f(t):=\Lambda_p(\mathsf V'_t)$ is holomorphic on $|t|<T$. This follows from \textup{($\rho$4)}, since $t\mapsto\mathsf V'_t$, extended by $\mathbf 1$, is a holomorphic curve in $\mathfrak V^\circ$, since $X'$ is a power series, and since the transports are polynomial. By \Cref{Y-C-Th68}\,(b), $\|f(t)\|\le\mathsf K_{\rm c}^0(17T\hat\varepsilon)^2+\tau_{\rm c}17T\hat\varepsilon$.

\emph{Step (3): the value and the derivative at $t=0$.} We have $f(0)=0$, because $\partial\mathsf V'_0=\mathbf 1$ and hence $X'=0$. We also have $f'(0)=0$. Indeed, only $X'$ is differentiated, and $\dot X_P=\operatorname{Ad}_{\mathsf k(y_P)}^{-1}(\mathrm{d}'A')_P$. On the box $\mathsf B=\mathsf B(p)$ of radius $R_3$ about $y(p)$, $\mathsf V'_0=\mathbf 1^u$ with $u=\mathsf k^{-1}$; here $|u|\le2$, and $|u(c_-)u(c_+)^{-1}|=|\mathsf V'_0(c)|<1+\vartheta_0^\flat$ by Step (1). The transports at $\mathsf V'_0$ are $T_0(y,y'')=\operatorname{Ad}_{\mathsf k(y)}^{-1}\operatorname{Ad}_{\mathsf k(y'')}$; therefore $(\mathsf g^{\mathsf V_0}\dot X)_j=\operatorname{Ad}_{\mathsf k(y(j))}^{-1}(\mathsf g\mathrm{d}'A')_j$. \Cref{Y-C-L69} gives $C^{\mathsf V_0}(b,j)=\operatorname{Ad}_{\mathsf k(y(b))}^{-1}C^{\mathbf 1}(b,j)\operatorname{Ad}_{\mathsf k(y(j))}$. In addition, $T_0(o,y(b))=\operatorname{Ad}_{\mathsf k(y(b))}$ and $T_0(o,y_P)\dot X_P=(\mathrm{d}'A')_P$. Hence \[f'(0)=\bigl((E-\mathrm{d}C^{\mathbf 1}\mathsf g-H^\natural)\mathrm{d}'A'\bigr)_p=\sum_{q'\in\mathcal Q_p}\mathsf G(p,q')(\mathrm{d}'\mathrm{d}'A')_{q'}=0.\]

\emph{Step (4): the maximum principle.} The function $f(t)/t^2$ is holomorphic on $|t|<T$, and the maximum principle for its norm, at $t=1$, gives $\|f(1)\|\le T^{-2}\sup\|f\|\le289\mathsf K_{\rm c}^0\hat\varepsilon^2+17\tau_{\rm c}\hat\varepsilon/T=(289\mathsf K_{\rm c}^0+17)\hat\varepsilon^2$.
\end{proof}

\begin{remark}[The three constants]
\begin{enumerate}
\item[\textup{(1)}] Three constants occur: $\mathsf K_{\rm c}^0$, which appears together with the linear term $\tau_{\rm c}\hat\varepsilon$ in \Cref{Y-C-Th68}\,(b); $\mathsf K_{\rm c}^1=289\mathsf K_{\rm c}^0+17$, in \Cref{Y-C-Th70}, which assumes \textup{(P2)}; and $\sqrt2\mathsf K_{\rm c}^0$, in \textup{(Cen)}, that is, in \Cref{Y-C-CIICor72}\,(a). Each of the three bounds holds for every $0<\hat\varepsilon\le\frac1{40}$.
\item[\textup{(2)}] For $\tau_{\rm c}\le\hat\varepsilon\le\frac1{40}$, in the setting of \Cref{Y-C-Th68} and with its hypotheses alone, $\|\Lambda_i(\mathsf V')\|\le(\mathsf K_{\rm c}^0+1)\hat\varepsilon^2\le\mathsf K_{\rm c}\hat\varepsilon^2$, where $\mathsf K_{\rm c}:=\sqrt2\mathsf K_{\rm c}^0$. Indeed, $\tau_{\rm c}\hat\varepsilon\le\hat\varepsilon^2$, and $\mathsf K_{\rm c}^0\ge6\cdot10^4LC_D\theta_D^{-1}\ge7.2\cdot10^5>(\sqrt2-1)^{-1}$. This $\mathsf K_{\rm c}$ is the letter $\mathsf K_{\rm c}$ of the table of letters, entered at \textup{(O10)} in item~3 of \S\ref{Y-G0-1}.
\item[\textup{(3)}] If \textup{(P1)} holds at $\mathsf n=60$, $p_0'\ge1$ and $\hat\varepsilon=\tau\varepsilon'$ with $\tau\ge1$, then $\hat\varepsilon\ge\tau_{\rm c}$, which is the lower bound on $\hat\varepsilon$ in \textup{(2)}. Indeed, $\varepsilon'\ge x'^{-1/2}>x^{-59}\ge\tau_{\rm c}$, by \textup{(P1)} applied with the value $60$ of its letter $\mathsf n$.
\end{enumerate}
\end{remark}

\begin{definition}[Two numbers and the floor \textup{(P3)}] We set \[\tau_2:=24L^3(2R+1)^3\gamma^{-2}x^{-\mathsf n}+\tfrac2\gamma\mathsf g_I(\tfrac\theta2)x^{-64},\qquad\mathsf{ti}_{\rm c}:=(16+4L)\,\mathsf N_{\mathsf B}\,\tau_2+7e\,(L+1)\,C_D\,x^{-128}.\] The floor \textup{(P3)} is the condition \[\textbf{(P3)}\ \ L^2(1+a)\bigl(\mathsf{ti}_{\rm c}+8C_D^3x'^{-64}\bigr)\le1 .\]
\end{definition}

\begin{lemma}[Range, closeness to the free matrices, and size of $H^\natural$ and $\mathsf N^\natural$]\label{Y-C-L71} Let a law satisfy \textup{($\lambda$0)}, \textup{($\rho$1)} and \textup{($\rho$5)}, and let \textup{(Fr1)} to \textup{(Fr3)}, \textup{(P0)} and \textup{(P1)} hold. Then the following hold.
\begin{enumerate}
\item[\textup{(a)}] $H^\natural(p,P)=0=\mathsf N^\natural(c,P)$ for ${\rm dist}>R_2$. Of the properties of the law, \textup{($\rho$1)} alone suffices for this clause, together with $R\ge6$ and $L\ge3$; \textup{(Fr1)} to \textup{(Fr3)} are not needed for it.
\item[\textup{(b)}] $\|H^\natural-H^\sharp\|_\infty+\|\mathsf N^\natural-\mathsf N^\sharp\|_\infty\le\mathsf{ti}_{\rm c}$.
\item[\textup{(c)}] If, in addition, \textup{(P3)} holds, then $L^2\|H^\natural\|_\infty\le\mathsf N_\sharp+1$ and $aL\|\mathsf N^\natural\|_\infty\le4\mathsf N_Z+1$.
\end{enumerate}
\end{lemma}

\begin{proof} \emph{Step (a): the ranges.} By \textup{($\rho$1)}, $C^{\mathbf 1}$ has range $\le R$; the matrices $d$ and $Q_1$ have range $\le1$, and $\mathsf g$ has range $\le2\varrho_L+1$. Moreover $1+R+2\varrho_L+1\le R_2$, by $R\ge6$ and $L\ge3$, through $2\varrho_L\le\frac23R+2$. No cube of $\mathcal Q_i$ has a face beyond $R_2$. This gives (a).

\emph{Step (b): closeness to the free matrices.} By \textup{(Fr1)}, $H^{(x)}=H_\circ+d(q^{-1}-C^{\mathbf 1})\mathsf g=H^\sharp+\mathsf G\mathrm{d}'+\mathsf{Ti}$, the last term being $\mathsf{Ti}=d(q^{-1}-C^{\mathbf 1})\mathsf g$. We bound the absolute row sums of $q^{-1}-C^{\mathbf 1}$. Its entries at $(b,b')$ with ${\rm dist}(b,b')\le R$ are $\le8\gamma^{-2}x^{-\mathsf n}$ by \Cref{Y-C-L66}\,(a), and a row has $3L^3(2R+1)^3$ of them. For the entries with ${\rm dist}(b,b')>R$, we have $C^{\mathbf 1}=0$ and $|q^{-1}(b,b')|\le\frac2\gamma e^{-\theta{\rm dist}}$, and they contribute at most $\frac2\gamma e^{-\theta R/2}\mathsf g_I(\frac\theta2)$ to a row sum. Altogether, the row sums are $\le\tau_2$, and so $\|\mathsf{Ti}\|_\infty\le4\cdot\tau_2\cdot4\mathsf N_{\mathsf B}$.

For ${\rm dist}(p,P)\le R_2$ one has $(H^\natural-H^\sharp)(p,P)=\mathsf{Ti}(p,P)+\sum_{q'\notin\mathcal Q_p}\mathsf G(p,q')\mathrm{d}'(q',P)$, and beyond $R_2$ the difference is $-H^\sharp(p,P)$. A cube off $\mathcal Q_p$ lies at distance $\ge R_2-1$ and has six faces, so that such cubes add at most $6C_De^{-\theta_D(R_2-1)}$. The tail of $H^\sharp$ contributes $C_De^{-\theta_DR_2}$. Here $e^{-\theta_DR_2}\le x^{-128}$ and $e^{\theta_D}\le e$. For $\mathsf N^\natural$, the matrix $Q_1$, whose rows sum to $L$, replaces $d$, which gives $L\tau_24\mathsf N_{\mathsf B}$, and $LC_D$ replaces $C_D$. Adding these contributions, we obtain (b).

\emph{Step (c): the sizes.} By \textup{(Fr3)}, $\|H^\sharp\|_\infty\le L^{-2}\mathsf N_\sharp+\|\mathsf T\|_0\|\mathsf B'-\mathsf B'^{\rm c}\|_0$. As $\mathsf B'-\mathsf B'^{\rm c}=\ell'^{>\mathsf T}\ell'+\ell'^{{\rm c}\mathsf T}\ell'^>$, we have $\|\mathsf B'-\mathsf B'^{\rm c}\|_0\le2C_D\|\ell'^>\|_0\le2C_D^2e^{-\theta_D\varrho(x')}\le2C_D^2x'^{-64}$. Further, $aL\|\mathsf N^\sharp\|_\infty\le4\|\mathsf B'\|_\infty\le4\mathsf N_Z+8C_D^2x'^{-64}$. With (b) and \textup{(P3)}, these bounds give (c).
\end{proof}

\begin{corollary}[The centre bound and the shrinking of the analytic terms given at $x$]\label{Y-C-CIICor72} Let the hypotheses of \Cref{Y-C-Th68} hold, together with \textup{(Fr1)} to \textup{(Fr3)} and \textup{(P3)}. Then $\tau_{\rm c}\le L^{-2}$, since $\tau_{\rm c}\le x^{-(\mathsf n-1)}\le x^{-1}\le L^{-2}$ by \textup{(P1)} and \textup{(P0)}, and the two assertions below are true.

\begin{enumerate}\item[\textup{(a)}] Condition \textup{(Cen)}$(\mathsf N_{\rm c},\sqrt2\mathsf K_{\rm c}^0,\mathsf a_{\rm c})$ holds with $\mathsf N_{\rm c}=17(\mathsf N_\sharp+1)$. Spelled out, this means the following. Take any fine plaquette $p$, any bond $b$ of $p$, any number $\hat\varepsilon$ in the range $0<\hat\varepsilon\le\frac1{40}$, and any field defined on the box of radius $R_3$ about $y(p)$ that satisfies \textup{(e1)}, \textup{(e2)} and $|\partial\mathsf V'(P)-\mathbf 1|\le2\hat\varepsilon$ at its plaquettes $P$. At such a field the first bound $|a_{0,b}|\le\mathsf a_{\rm c}\hat\varepsilon$ holds, and so does the second bound \[|(EX'-\mathrm{d}^{\mathsf V}a_0)_p|\le\mathsf N_{\rm c}L^{-2}\hat\varepsilon+\sqrt2\mathsf K_{\rm c}^0\hat\varepsilon^2.\]

\item[\textup{(b)}] \emph{At the law of windows, the analytic terms given at $x$ shrink.} We begin with the objects that enter this clause. The fine cubes are the cubes with $M^3$ fine sites, taken from the grid in which each $\hat\square$ is a union of $L^3$ such cubes (see the definition ``Cubes'' in item~1 of \S\ref{Y-G0-1}); the distance between two fine cubes is measured in units of cubes. For a set $X$ of fine cubes we set $X^\approx:=X^{+r(x)}$, and by a bond, a path or a plaquette of $X^\approx$ we mean one whose sites all belong to $X^\approx$.

We write $\mathcal U_X(x)$ for the set of those $\operatorname{SL}(2,\mathbb C)$ fields $\mathsf U$ on the bonds of $X^\approx$ that satisfy the three conditions below.

\begin{enumerate}\item[\textup{(u1)}] $|\mathsf U(b)|<1+\vartheta(x)$ at each bond $b$ of $X^\approx$.

\item[\textup{(u2)}] $|\mathsf U(\pi)|<\mathsf b(x)$ for each path $\pi$ in $X^\approx$ that has at most $\ell_w(x)$ bonds.

\item[\textup{(u3)}] $|\partial\mathsf U(p)-\mathbf 1|<2\varepsilon$ at each plaquette $p$ of $X^\approx$.\end{enumerate}

For a set $\mathsf B$ of coarse sites and a number $\tau\ge1$, the set $U^\tau(\mathsf B;x')$ consists of the $\operatorname{SL}(2,\mathbb C)$ fields on the bonds of $\mathsf B$ that obey the three conditions below.

\begin{enumerate}\item[\textup{(e1)}] Condition \textup{(e1)} of item~C.II~1.4.

\item[\textup{(e2)}] Condition \textup{(e2)} of item~C.II~1.4, for those paths in $\mathsf B$ that have at most $\ell_w(x')$ bonds.

\item[\textup{(e3)}] $|\partial\mathsf V'(P)-\mathbf 1|<2\tau\varepsilon'$ at each plaquette $P$ of $\mathsf B$.\end{enumerate}

We put $c_a(x):=E_ap_0(x)^3x^{-1/2}$ with a number $E_a>0$, and we let $T:=L^{3/2}/(19\mathsf N_{\rm c})$, $c_\xi:=(6.5C_\varphi^\flat)^{-1}$ and $\mathsf K_{\rm c}=\sqrt2\mathsf K_{\rm c}^0$. With these numbers we define $r_T:=\sqrt2C_\varphi^\flat c_\xi\varepsilon+17T\mathsf a_{\rm c}\varepsilon'$. Further, $D_w:=\lceil\ell_w(x)/L\rceil+1$ and $\omega:=M+\lceil Mr(x)/L\rceil+D_w+R_3+2$; these are the numbers $D_w$, $\omega$ of item~2 of \S\ref{Y-G0-1}.

In what follows, $X$ is a connected set of at most $\frac L2-2$ fine cubes (such a set is called small) whose first cube, with respect to a fixed ordering of all cubes, lies in $\hat\square_0$. Let $o$ be the centre site of the cell $\square_0$, and let $\mathsf B_\omega$ be the box of radius $\omega$ about $o$. We define $\chi_X(\xi):=1\{|\xi_b|\le c_\xi\varepsilon$ for all bonds $b$ of $I^\circ$ that lie within $R(x)$ of some block meeting $X^\approx\}$.

Let $E^a(X,\cdot)$ be holomorphic on $\mathcal U_X(x)$ and invariant under the gauge transformations with values in $G$, and let $\sup|E^a(X,\cdot)|\le\mathsf Ac_a(x)e^{-\kappa(\#X-1)}\le\frac12$ hold for a number $\mathsf A\ge1$; for the class of item~1 of \S\ref{Y-G0-1} one has $\mathsf A=1$. We put \[\mathcal G_X(\mathsf V'):=E_\xi\bigl[\chi_X(\xi)\bigl(e^{E^a(X,\mathsf U)}-1\bigr)\bigr]\] with $\xi\sim\mathcal N(0,x^{-1}\mathbf 1)$, where $\mathsf U=\mathsf U(G^{\mathsf V}\xi-a_0;\mathsf V')$ is the chart of item~C.I~4, defined at a complex field by the same expression. Here $E^a(X,\mathsf U)$ stands for the value of $E^a(X,\cdot)$ at the restriction of $\mathsf U$ to the bonds of $X^\approx$. Wherever $\chi_X(\xi)=0$, the integrand is taken to be $0$: for such $\xi$ the restriction need not belong to $\mathcal U_X(x)$, and $E^a(X,\mathsf U)$ need not be defined.

The hypotheses of clause \textup{(b)} are those of clause \textup{(a)}, a law with \textup{($\lambda$0)} to \textup{($\lambda$7)}, the condition \textup{(Tor)}, and nine floors. Four of these floors have no name: $T>1$, $17T\varepsilon'\le\frac1{40}$, $\mathsf K_{\rm c}(17T\varepsilon')^2\le0.12\varepsilon$ and $\omega\le Mr'$. The remaining five, \textup{(D2)}, \textup{(D3)}, \textup{(C0)}, \textup{(C1)} and \textup{(C2)}, are as follows.

\begin{enumerate}\item[\textup{(D2)}] $\ell_w(x)(7.1r_T+18D_w\cdot17T\varepsilon')\le\frac12\mathsf g_x$.

\item[\textup{(D3)}] $(4.1r_T+2.1\cdot17T\varepsilon')^2\le0.1\varepsilon$.

\item[\textup{(C0)}] $18\omega\varepsilon'\le\frac1{10}$ and $160T\omega\varepsilon'\le1$.

\item[\textup{(C1)}] $2640\,\ell_w(x')(T+1)\omega\varepsilon'\le\mathsf g_x$.

\item[\textup{(C2)}] $102400\,T\omega^2\varepsilon'\le1$.\end{enumerate}

Each of the nine floors except $T>1$ appears among \textup{(A0)}, \textup{(A2)} of item~4 of \S\ref{Y-G0-1}, and $T>1$ follows from \textup{(L2$^{\rm su}$)}, which is the condition $L^{1/2}\ge1.2\cdot10^4\mathsf N_{\rm c}^2$ of claim \textup{(3)} below. When the letters have the values fixed in item~3 of \S\ref{Y-G0-1}, all these floors are satisfied for every sufficiently large $x$; this is shown in Step~14 of the proof of \Cref{Y-G-ThA} and in Step~0 of the proof of \Cref{Y-H-C9}. Under the hypotheses of clause \textup{(b)}, the following three claims hold.

\begin{enumerate}\item[\textup{(1)}] On $U^{17T}(\mathsf B_\omega;x')$ the function $\mathcal G_X$ is defined, holomorphic and invariant under the gauge transformations with values in $G$, and there $|\mathcal G_X|\le2\mathsf Ac_a(x)e^{-\kappa(\#X-1)}$.

\item[\textup{(2)}] The bound \[|\mathcal G_X(\mathsf V')-\mathcal G_X(\mathbf 1)|\le4T^{-2}\mathsf Ac_a(x)e^{-\kappa(\#X-1)}\] holds at every field $\mathsf V'$ on $\mathsf B_\omega$ that satisfies \textup{(u1)}, \textup{(u2)}, \textup{(u3)} at $x'$, and hence at the restriction to $\mathsf B_\omega$ of each member of $\mathcal U_{\{\square_0\}}(x')$.

\item[\textup{(3)}] If, in addition, $(p_0'/p_0)^4\ge\frac12$ and \textup{(L2$^{\rm su}$)} $L^{1/2}\ge1.2\cdot10^4\mathsf N_{\rm c}^2$ hold, then the right sides of \textup{(2)}, summed over all small $X$ whose first cube lies in $\hat\square_0$, total at most $\frac14\mathsf Ac_a(x')$.\end{enumerate}\end{enumerate}\end{corollary}

\begin{proof}[Proof of \Cref{Y-C-CIICor72}\,(a)] By the parenthetical remark at the end of \Cref{Y-C-Th68}, the conditions \textup{(e1)}, \textup{(e2)} and $|\partial\mathsf V'(P)-\mathbf 1|\le2\hat\varepsilon$ on the field give \textup{(h)} at $\hat\varepsilon$. By Step~(0) of the proof of \Cref{Y-C-Th68}, every bond $b$ of $p$ has ${\rm dist}(b,o)\le1$ for $o=y(p)$; hence \Cref{Y-C-Th68}\,(c), applied with $i=p$ and $o=y(p)$, gives the first bound $|a_{0,b}|\le\mathsf a_{\rm c}\hat\varepsilon$. The sum $\sum_PH^\natural(p,P)\tilde X_P$ obeys $\|\sum_PH^\natural(p,P)\tilde X_P\|\le\|H^\natural\|_\infty8.24\hat\varepsilon$, since $\|\tilde X_P\|\le|\tilde X_P|\le8.24\hat\varepsilon$ at every plaquette $P$ of the box of radius $R_3$ about $y(p)$: the first inequality is the bound $\|Y\|\le|Y|$ of item~C.II~1.1, and the second is \Cref{Y-C-L67}\,(c), applied to that box with $o=y(p)$ and $\rho_{\mathsf B}=R_3$, which is allowed because $3R_3+1\le\ell_w(x')$ by \textup{(P0)}. Combined with \Cref{Y-C-L71}\,(c) and \Cref{Y-C-Th68}\,(b), this gives the estimate $\|(EX'-\mathrm{d}^{\mathsf V}a_0)_p\|\le[8.24(\mathsf N_\sharp+1)+1]L^{-2}\hat\varepsilon+\mathsf K_{\rm c}^0\hat\varepsilon^2$, where the $+1$ in the bracket accounts for the term $\tau_{\rm c}\hat\varepsilon$ of \Cref{Y-C-Th68}\,(b), since $\tau_{\rm c}\le L^{-2}$ gives $\tau_{\rm c}\hat\varepsilon\le L^{-2}\hat\varepsilon$. The second bound of clause \textup{(a)} now follows from $|\cdot|\le\sqrt2\|\cdot\|$ and $\sqrt2\cdot9.24\le17$.\end{proof}

\begin{proof}[Proof of \Cref{Y-C-CIICor72}\,(b)] \emph{Step (o): dependence on the data, and sizes.} We put $\hat\varepsilon:=17T\varepsilon'$. Let $\mathsf V'\in U^{17T}(\mathsf B_\omega;x')$, let $\xi$ be such that $\chi_X(\xi)=1$, and set $\eta:=G^{\mathsf V}\xi-a_0$ and $F^\flat:=EX'+\mathrm{d}^{\mathsf V}\eta$. Because $\omega\le Mr'$, we have $3D_w+1\le3\omega+1\le\ell_w(x')$. Next, $r_T\ge0.2175\varepsilon$; indeed $r_T\ge\sqrt2C_\varphi^\flat c_\xi\varepsilon=\frac{\sqrt2}{6.5}\varepsilon$, and $\frac{\sqrt2}{6.5}=0.21757\dots$. As $(4.1\cdot0.2175)^2\ge0.7952$, the floor \textup{(D3)} gives $0.79\varepsilon^2\le0.1\varepsilon$, whence $\varepsilon\le0.13$.

We turn to the geometry. Each site $z$ of $X^\approx$ lies at sup distance $<M(\frac L2-2+r(x))$ from some site $z_0$ of the first cube of $X$. Consequently $L^{-1}|y(z)-y(z_0)|_\infty<\frac M2-\frac{2M}L+\frac{Mr(x)}L+1\le\frac M2+\frac{Mr(x)}L-1$, where the second inequality rests on $M\ge L$, which holds as $m'\ge1$ (item~C.II~1.1). The coarse site $y(z_0)\in\square_0$ is within $\frac12(M-1)$ of $o$. It follows that the blocks meeting $X^\approx$ lie within $M+\lceil Mr(x)/L\rceil$ of $o$, and that $\mathsf B_\omega$ contains the box of radius $D_w$ and the box of radius $R_3$ about the coarse site of each of these blocks. Every bond and every plaquette of $X^\approx$ is based in a block of this kind.

The field given on $\mathsf B_\omega$ is extended by $\mathbf 1$ to the remaining coarse bonds; \textup{($\lambda$0)} of the proposition in item~C.II~2.1 applies to the extended field, because \textup{(e1)} gives \textup{(h1)}. By \textup{($\lambda$1)} and \textup{($\lambda$2)}, the variable $(G^{\mathsf V}\xi)_b$ depends on $\xi$ through $Y_b$ alone, hence within $\frac12R$, and on the field within $r_w+\varrho_L+1\le R_3$. By the second and the third sentence of \Cref{Y-C-Th68}\,(a), $a_0$ depends on the field only within $R_3$. In detail: the field is given on $\mathsf B_\omega$, and $X'$ is defined there; hence $a_{0,b}$ and $(EX'-\mathrm{d}^{\mathsf V}a_0)_p$ take the values that are computed on the box of radius $R_3$ about $y(b)$ and about $y(p)$, respectively, and it is to these values that we apply \Cref{Y-C-Th68}\,(c) and clause \textup{(a)} of the corollary, in \textup{($\eta$1)} and \textup{($\eta$2)} of this step. It follows from these dependences and the geometry above that none of the quantities used in this proof depends on the extension of the field by $\mathbf 1$.

The data \textup{($\eta$1)}, \textup{($\eta$2)} of Step~(ii), with $r_T$ in the role of $r$, are obtained as follows.

\begin{enumerate}\item[\textup{($\eta$1)}] By \textup{($\lambda$5)}, $|(G^{\mathsf V}\xi)_b|\le\sqrt2C_\varphi^\flat c_\xi\varepsilon$: the supremum in \textup{($\lambda$5)} runs over $Y_b$, and there $\chi_X(\xi)=1$ bounds $\xi$. \Cref{Y-C-Th68}\,(c), applied with $i=c$ and $c_-=y(b)$, gives $|a_{0,b}|\le\mathsf a_{\rm c}\hat\varepsilon$. Hence $|\eta_b|\le r_T$ at every bond of $X^\approx$.

\item[\textup{($\eta$2)}] We use clause \textup{(a)}, applied with the value $\hat\varepsilon=17T\varepsilon'$ of this step; the inequality $\varepsilon'\le L^{1/2}\varepsilon$, which holds because $p_0'\le p_0$; the third floor $\mathsf K_{\rm c}(17T\varepsilon')^2\le0.12\varepsilon$; and $|\operatorname{Ad}_gY|\le|g|^2|Y|\le1.01|Y|$. Combined, they yield \[|F^\flat_p|\le\mathsf N_{\rm c}L^{-2}\hat\varepsilon+\mathsf K_{\rm c}\hat\varepsilon^2+4\cdot1.01\sqrt2C_\varphi^\flat c_\xi\varepsilon\le1.9\varepsilon\] at each plaquette of $X^\approx$. As to the numerals, $17T\mathsf N_{\rm c}L^{-3/2}=\frac{17}{19}$ and $4.04\sqrt2/6.5\le0.879$, and with $0.12$ the total is at most $1.894$, hence $\le1.9$.\end{enumerate}

In the invariance argument of Step~(iii), a gauge transformation $u$ is extended by $\mathbf 1$ outside $\mathsf B_\omega$. The extended field, once transformed, and the extension of $\mathsf V'^u$ disagree only at the bonds leaving $\mathsf B_\omega$; at these bonds the transformed field takes unitary values, so that it lies in $\mathfrak V$. For a bond $b$ and a plaquette $p$ of $X^\approx$, the dependences and the geometry described earlier in this step show that the row $G^{\mathsf V}(b,\cdot)$ and the values $a_{0,b}$ and $(EX'-\mathrm{d}^{\mathsf V}a_0)_p$ do not involve the field at the bonds leaving $\mathsf B_\omega$. Finally, $\mathsf B_\omega$ is embedded: by \textup{(Tor)} we have $2\omega+1<N$, as $2\omega+1\le2Mr'+1<16Mr(x)$.

\emph{Step (i): an expansion of the curvature of the chart.} The data are: numbers $\hat\varepsilon\le\frac1{40}$ and $r>0$ such that $4.1r+2.1\hat\varepsilon\le1$; a fine plaquette $p=(z;\mu\nu)$, with $y=y(p)$; a field $\mathsf V'$ that satisfies \textup{(h1)} at the coarse bonds and \textup{(h3)} with parameter $\hat\varepsilon$ at the coarse plaquette that occurs below; and $\eta$ with $|\eta_b|\le r$ at the bonds of $p$, where $\eta:=0$ on $\mathcal T$. We claim the bound \[|\partial\mathsf U(p)-\mathbf 1-iF^\flat_p|\le0.72(4.1r+2.1\hat\varepsilon)^2\] for $\mathsf U=\mathsf U(\eta;\mathsf V')$ and $F^\flat=EX'+\mathrm{d}^{\mathsf V}\eta$.

For the proof, let $\eta_1,\dots,\eta_4$ denote the values of $\eta$ at $(z,\mu)$, $(z+e_\mu,\nu)$, $(z+e_\nu,\mu)$, $(z,\nu)$; we say that $z_\mu$ is at the top if $z_\mu=y_\mu+\mathsf h_L$, and the same definition, with $\nu$ in place of $\mu$, applies to $z_\nu$. If neither $z_\mu$ nor $z_\nu$ is at the top, then $\partial\mathsf U(p)=e^{i\eta_1}e^{i\eta_2}e^{-i\eta_3}e^{-i\eta_4}$. If only $z_\mu$ is at the top, let $c=(y,\mu)$: the first and the third bond of $p$ cross $c$, while the second is based in $B(c_+)$, so that \[\partial\mathsf U(p)=e^{i\eta_1}\mathsf V'(c)e^{i\eta_2}\mathsf V'(c)^{-1}e^{-i\eta_3}e^{-i\eta_4}=e^{i\eta_1}e^{i\operatorname{Ad}_{\mathsf V'(c)}\eta_2}e^{-i\eta_3}e^{-i\eta_4}.\] If only $z_\nu$ is at the top, let $c=(y,\nu)$: now the second and the fourth bond cross $c$, while the third is based in $B(c_+)$, and $\partial\mathsf U(p)$ equals $e^{i\eta_1}e^{i\eta_2}e^{-i\operatorname{Ad}_{\mathsf V'(c)}\eta_3}e^{-i\eta_4}$.

If both $z_\mu$ and $z_\nu$ are at the top, then $p$ is an edge plaquette under $P=(y;\mu\nu)$, whose curvature, written through its four coarse bonds, is $\partial\mathsf V'(P)=\mathsf V'(c_1)\mathsf V'(c_2)\mathsf V'(c_3)^{-1}\mathsf V'(c_4)^{-1}$, and \[\partial\mathsf U(p)=e^{i\eta_1}\mathsf V'(c_1)e^{i\eta_2}\mathsf V'(c_2)\mathsf V'(c_3)^{-1}e^{-i\eta_3}\mathsf V'(c_4)^{-1}e^{-i\eta_4}=e^{i\eta_1}e^{i\operatorname{Ad}_{\mathsf V'(c_1)}\eta_2}e^{iX'_P}e^{-i\operatorname{Ad}_{\mathsf V'(c_4)}\eta_3}e^{-i\eta_4}.\] In all four cases the transports that occur are the $T(y(p),y(b))$ of $\mathrm{d}^{\mathsf V}$; hence the sum of the exponents is $iF^\flat_p$. The norms of the exponents add up to at most $(2+2(1+\vartheta_0^\flat)^2)r+2.06\hat\varepsilon\le4.1r+2.1\hat\varepsilon$, and \textup{(E)} of item~C.I~1.1 gives the claim.

\emph{Step (ii): the fine field lies in $\mathcal U_X(x)$.} The data are: a number $\tau$ with $1\le\tau\le17T$ and $\hat\varepsilon=\tau\varepsilon'\le\frac1{40}$; a number $r>0$ such that \textup{(D2)}, \textup{(D3)} hold at $(\tau,r)$, that is, $\ell_w(x)(7.1r+18D_w\hat\varepsilon)\le\frac12\mathsf g_x$ and $(4.1r+2.1\hat\varepsilon)^2\le0.1\varepsilon\le1$; the inequality $3D_w\le\ell_w(x')$; an arbitrary non-empty set $X$ of fine cubes; a set $\mathsf B$ of coarse sites containing the box of radius $D_w$ about each site whose block meets $X^\approx$; a field $\mathsf V'\in U^\tau(\mathsf B;x')$; and $\eta$ with the two properties below.

\begin{enumerate}\item[\textup{($\eta$1)}] $|\eta_b|\le r$ at the bonds of $X^\approx$.

\item[\textup{($\eta$2)}] $|F^\flat_p|\le1.9\varepsilon$ at the plaquettes of $X^\approx$.\end{enumerate}

We claim that $\mathsf U(\eta;\mathsf V')$, restricted to the bonds of $X^\approx$, belongs to $\mathcal U_X(x)$. We prove the three conditions one by one.

\begin{enumerate}\item[\textup{(u3)}] The coarse bonds and the coarse plaquette that enter Step~(i) lie within $1$ of $y(p)$, hence in $\mathsf B$, where \textup{(e1)} and \textup{(e3)} give \textup{(h1)} and \textup{(h3)}. Step~(i) therefore yields $|\partial\mathsf U(p)-\mathbf 1|\le1.9\varepsilon+0.072\varepsilon<2\varepsilon$.

\item[\textup{(u1)}] We have $|\mathsf U(b)|\le e^r(1+\vartheta(x')+\mathsf g_x)<1+\vartheta(x')+4\mathsf g_x\le1+\vartheta(x)$; here $r\le\frac1{14}\mathsf g_x$ by \textup{(D2)}, and $\vartheta(x)-\vartheta(x')=\vartheta_0^\flat\frac{\log L}{\log x\log x'}\ge4\mathsf g_x$.

\item[\textup{(u2)}] Let $\pi$ be a path in $X^\approx$ with $n\le\ell_w(x)$ bonds, say $b_1,\dots,b_n$, which passes through the sites $z_0,\dots,z_n$, and write $y_i:=y(z_i)$ and $o_1:=y_0$. Then $L^{-1}|y_i-o_1|_\infty<\frac nL+1$, so every $y_i$ is within $D_w$ of $o_1$, and the box of radius $D_w$ about $o_1$ is contained in $\mathsf B$. Since by its definition in item~C.II~1.4 a box is embedded, the hypothesis on $\mathsf B$ includes $2D_w+1<N$; in Step~(iii), that inequality holds because $D_w<\omega$ and $2\omega+1<N$ by Step~(o).

Inside this box, $\mathsf k(y):=\mathsf V'(\Gamma(o_1,y))$ obeys $|\mathsf k(y)|<\mathsf b(x')+\mathsf g_x<2$ by \textup{(e2)}, since every comb there has at most $3D_w$ bonds. With $g_i:=\mathsf k(y_{i-1})\mathsf U(b_i)^{\pm1}\mathsf k(y_i)^{-1}$ we have $\mathsf U(\pi)=g_1\cdots g_n\mathsf k(y_n)$. For a bond of $\pi$ that lies inside a block, $g_i=\mathsf ke^{\pm i\eta}\mathsf k^{-1}$, and $|g_i|\le1+4(e^r-1)\le e^{7.1r}$. For a crossing bond of $c$, the factor $g_i$ is one of this kind multiplied by $\mathsf W(c)^{\pm1}$, and $|\mathsf W(c)^{\pm1}|=|\mathsf W(c)|\le e^{18D_w\hat\varepsilon}$ by \Cref{Y-C-L67}\,(a), applied in this box; the proof of that clause invokes \textup{(h2)} only along combs, so that $3D_w\le\ell_w(x')$ suffices for it.

Therefore \[|\mathsf U(\pi)|<e^{\ell_w(x)(7.1r+18D_w\hat\varepsilon)}(\mathsf b(x')+\mathsf g_x)\le e^{\mathsf g_x/2}(\mathsf b(x')+\mathsf g_x)<\mathsf b(x')+4\mathsf g_x\le\mathsf b(x),\] where for the last inequality we use $\mathsf b(x)-\mathsf b(x')=\frac{\log L}{\log x\log x'}\ge4\mathsf g_x$.\end{enumerate}

\emph{Step (iii): proof of claim \textup{(1)}.} By Step~(o), Step~(ii) applies with $(\tau,r)=(17T,r_T)$ and $\mathsf B=\mathsf B_\omega$; thus, wherever $\chi_X(\xi)=1$, the field $\mathsf U$ on the bonds of $X^\approx$ belongs to $\mathcal U_X(x)$. Since $|e^w-1|\le2|w|$ for $|w|\le\frac12$, the bound on $E^a(X,\cdot)$ gives the bound of claim \textup{(1)}. As to holomorphy: by \textup{($\lambda$4)} and \textup{($\lambda$7)} the entries of $G^{\mathsf V}$ and $C^{\mathsf V}$ are holomorphic at the extended field, which belongs to $\mathfrak V^\circ$ by \textup{(e1)}; by \textup{(e3)}, $X'$ is given by a power series; the chart is a polynomial in $\mathsf V'$ and an entire function of $\eta$; and the integrand is bounded and depends on finitely many $\xi_b$ only. As to invariance: \textup{($\lambda$3)} and the covariance of $\mathsf g^{\mathsf V}$ and of $X'$ give $\eta[\mathsf V'^u,\operatorname{Ad}_u\xi]=\operatorname{Ad}_u\eta[\mathsf V',\xi]$; moreover $\mathsf U(\operatorname{Ad}_u\eta;\mathsf V'^u)=\mathsf U(\eta;\mathsf V')^{\tilde u}$ with the fine gauge transformation $\tilde u(z):=u(y(z))$; and $E^a$, $\chi_X$ and the law of $\xi$ are invariant. This proves claim \textup{(1)}.

\emph{Step (iv): the change of a bounded holomorphic invariant function.} The data are: an embedded box $\mathsf B$ all of whose sites lie within $\omega$ of a site $o\in\mathsf B$; the inequalities $3\omega+1\le\ell_w(x')$, $T>1$ and $17T\varepsilon'\le\frac1{40}$, together with \textup{(C0)}, \textup{(C1)}, \textup{(C2)}; the set $\Omega:=U^{17T}(\mathsf B;x')$; and a function $\Phi$ on $\Omega$ that is holomorphic, invariant under the gauge transformations with values in $G$, and bounded by $|\Phi|\le\mathsf S$. The claim is that $|\Phi(\mathsf V')-\Phi(\mathbf 1)|\le2\mathsf ST^{-2}$ at every field $\mathsf V'$ on $\mathsf B$ that satisfies \textup{(u1)}, \textup{(u2)}, \textup{(u3)} at $x'$.

We first build a curve of fields. The field $\mathsf V'$ has \textup{(h2)} and \textup{(h3)} at $\varepsilon'$. By \Cref{Y-C-L67}\,(a) and \textup{(C0)}, $|\mathsf W(c)-\mathbf 1|\le e^{18\omega\varepsilon'}-1\le19\omega\varepsilon'\le0.106$, and $A'_c:=-i\log\mathsf W(c)$ is a traceless matrix with $|A'_c|\le\mathsf a:=21\omega\varepsilon'$. For complex $t$ we put $\mathsf V'_t(c):=\mathsf k(c_-)^{-1}e^{itA'_c}\mathsf k(c_+)=\mathsf V'(c)(\mathbf 1+\delta_c)$, where $\mathbf 1+\delta_c=\mathsf k(c_+)^{-1}e^{i(t-1)A'_c}\mathsf k(c_+)$; this field is entire in $t$, and $\mathsf V'_1=\mathsf V'$. Now let $|t|<T$. Then $(T+1)\mathsf a\le1$ by \textup{(C1)}, and therefore $|\delta_c|,|(\mathbf 1+\delta_c)^{-1}-\mathbf 1|\le4(e^{(T+1)\mathsf a}-1)\le\delta:=168(T+1)\omega\varepsilon'$. We show that $\mathsf V'_t$ has the three properties that define $\Omega$.

\begin{enumerate}\item[\textup{(e1)}] $|\mathsf V'_t(c)|<(1+\delta)(1+\vartheta(x'))\le1+\vartheta(x')+1.01\delta\le1+\vartheta(x')+\mathsf g_x$, where the last inequality comes from \textup{(C1)}.

\item[\textup{(e2)}] Take a path of $n\le\ell_w(x')$ bonds, and shift each factor $(\mathbf 1+\delta_c)^{\pm1}$ to the left past the partial products of $\mathsf V'$, whose norms are $<2$ by \textup{(u2)}. This gives \[|\mathsf V'_t(\pi)|\le(1+4\delta)^n|\mathsf V'(\pi)|<e^{4n\delta}\mathsf b(x')\le\mathsf b(x')+2(e^{4n\delta}-1)<\mathsf b(x')+\mathsf g_x\] with $4n\delta\le\frac{672}{2640}\mathsf g_x$ by \textup{(C1)}.

\item[\textup{(e3)}] The computation of Step~(1) of the proof of \Cref{Y-C-Th70} gives $\partial\mathsf V'_t(P)=\mathsf k(y_P)^{-1}\Pi_t\mathsf k(y_P)$ with $\Pi_t=e^{itA'_1}e^{itA'_2}e^{-itA'_3}e^{-itA'_4}$. By \textup{(E)}, $|\Pi_t-\mathbf 1-it(\mathrm{d}'A')_P|\le(4|t|\mathsf a)^2$, where $4T\mathsf a=84T\omega\varepsilon'\le1$ by \textup{(C0)}. At $t=1$, \textup{(u3)} gives $|\Pi_1-\mathbf 1|<8\varepsilon'$, and therefore $|(\mathrm{d}'A')_P|<8\varepsilon'+16\mathsf a^2$. It follows that $|\partial\mathsf V'_t(P)-\mathbf 1|<32T\varepsilon'+128T^2\mathsf a^2\le34T\varepsilon'$; for the last inequality note that $128\cdot441\,T\omega^2\varepsilon'=56448\,T\omega^2\varepsilon'\le2$ by \textup{(C2)}.\end{enumerate}

Thus $\mathsf V'_t\in\Omega$ for $|t|<T$; on this disc the function $f(t):=\Phi(\mathsf V'_t)$ is holomorphic, and $|f|\le\mathsf S$.

We now bring in complex gauge transformations. The set $\Omega$ is open. Let $\mathsf Z\in\Omega$ and $\lambda:\mathsf B\to\mathfrak g_{\mathbb C}$. The derivative of $\Phi(\mathsf Z^{\exp(i\sigma\lambda)})$ at $\sigma=0$ is $\mathbb C$-linear in $\lambda$, as $\Phi$ is holomorphic, and it vanishes whenever $\lambda$ is $\mathfrak g$-valued, as $\Phi$ is invariant; since every $\lambda$ can be written as $\lambda_1+i\lambda_2$ with $\mathfrak g$-valued summands, the derivative vanishes for every $\lambda$. By the argument of Step~(2) of the proof of \Cref{Y-C-L69}, it follows that $\Phi(\mathsf Z^{u_s})$ is constant along each $C^1$ curve of gauge transformations $u_s:\mathsf B\to\operatorname{SL}(2,\mathbb C)$ for which $\mathsf Z^{u_s}\in\Omega$ throughout.

Write $\mathsf k(y)^{-1}=k_0(y)e^{h(y)}$ in polar form, with $k_0\in G$ and $h$ hermitian and traceless, and set $u_s:=k_0e^{sh}$ for $0\le s\le1$. For a path $\pi$ from $y$ to $y''$, Cordes' inequality, in the form used in Step~(1) of the proof of \Cref{Y-C-L69}, gives \[|\mathbf 1^{u_s}(\pi)|=|e^{sh(y)}e^{-sh(y'')}|\le|e^{h(y)}e^{-h(y'')}|^s=|\mathsf V'_0(\pi)|^s\le|\mathsf V'_0(\pi)|.\] Since $\mathsf V'_0=\mathbf 1^{u_1}$ lies in $\Omega$, each $\mathbf 1^{u_s}$ therefore has \textup{(e1)} and \textup{(e2)}; its curvature is $\mathbf 1$; and so $\mathbf 1^{u_s}\in\Omega$. Consequently $f(0)=\Phi(\mathbf 1^{u_1})=\Phi(\mathbf 1^{u_0})=\Phi(\mathbf 1)$, because $u_0$ takes values in $G$.

Next we show $f'(0)=0$. The set $\{\mathbf 1^{u_s}\}$ is compact and contained in the open set $\Omega$; hence some $t_0>0$ has the property that $\mathsf Z_t^{u_s}\in\Omega$ whenever $|t|<t_0$, for every $s$, where $\mathsf Z_t:=(e^{itA'_c})$. As $\mathsf V'_t=\mathsf Z_t^{u_1}$, we obtain $f(t)=\Phi(\mathsf Z_t^{u_0})=\Phi(\mathsf Z_t)$ for these $t$. The derivative of $\Phi(\mathsf Z_t)$ at $t=0$ is a $\mathbb C$-linear function of $A'$, and it is unchanged under $A'\mapsto\operatorname{Ad}_uA'$ for every constant $u\in G$, because $\mathsf Z_t^u=(e^{it\operatorname{Ad}_uA'_c})$. It therefore takes the same value at $\int_G\operatorname{Ad}_uA'\,du$. Now $\int_G\operatorname{Ad}_uY\,du$ commutes with $G$ and is traceless, so it equals $0$, and the derivative vanishes.

To finish, $(f(t)-f(0))/t^2$ is holomorphic on $|t|<T$, and $|f-f(0)|\le2\mathsf S$ there. The maximum principle on the circles $|t|=\rho$ with $1<\rho<T$, followed by the limit $\rho\to T$, gives $|f(1)-f(0)|\le2\mathsf ST^{-2}$, and this is the claim.

\emph{Step (v): proof of claims \textup{(2)} and \textup{(3)}.} Claim \textup{(2)} is the case $\mathsf B=\mathsf B_\omega$, $\Phi=\mathcal G_X$, $\mathsf S=2\mathsf Ac_a(x)e^{-\kappa(\#X-1)}$ of Step~(iv), the bound $\mathsf S$ being supplied by claim \textup{(1)}; moreover $\mathsf B_\omega\subset\square_0^\approx$, since $\omega\le Mr'$. For claim \textup{(3)}, observe that $\hat\square_0$ contains $L^3$ cubes that can serve as first cubes. The bound of item~C.II~1.1 on the number of connected sets gives, for any given cube, $\sum_{X\ni\text{the given cube}}e^{-\kappa(\#X-1)}\le\sum_{n\ge1}(26e^{1-\kappa})^{n-1}\le(1-26e^{-5})^{-1}\le1.2124$. Moreover $L^3T^{-2}=361\mathsf N_{\rm c}^2$, and $c_a(x')=L^{1/2}(p_0'/p_0)^3c_a(x)\ge2^{-3/4}L^{1/2}c_a(x)$. Hence the sum in claim \textup{(3)} is $\le4\cdot1.2124\cdot361\,\mathsf N_{\rm c}^2\,2^{3/4}L^{-1/2}\mathsf Ac_a(x')\le\frac14\mathsf Ac_a(x')$, where the last inequality follows from \textup{(L2$^{\rm su}$)} because $4\cdot1.2124\cdot361/(0.25\cdot2^{-3/4})\le11777.3\le1.2\cdot10^4$.\end{proof}

\subsubsection*{C.II 5. \textup{(S-loc$_k$)}}

\begin{definition}[Walks and their traces sorted by sets of cells]\label{Y-C-CIIWalks}
Let a family of steps be given. A walk $\omega$ of $k$ steps in this family consists of bonds $b_0,\dots,b_k$ together with, for each $i$, either a step $\mathsf s_i:b_{i-1}\to b_i$ of the family or the trivial step; the trivial step has $b_i=b_{i-1}$, weight $1$ and length $0$, and its set of cells is $\mathsf c=\{$the cell of $b_i\}$. The length of $\omega$ is $|\omega|:=\sum|\mathsf s_i|$, and its set of cells, $\mathsf c(\omega):=\bigcup\mathsf c(\mathsf s_i)$, is a connected set of cells. Whatever the number $k$ of its steps, \begin{equation}\tag{Cnt}\#\mathsf c(\omega)\le27(\frac{|\omega|}M+1).\end{equation} Indeed, $\pi_\omega=\pi_{\mathsf s_1}\cdots\pi_{\mathsf s_k}$ is a coarse lattice path of $|\omega|$ bonds, and $\mathsf c(\omega)\subset\{$cells of its sites$\}$. Along $M$ consecutive sites this path moves by at most $M-1$, so such sites lie in the cube of $27$ cells centred at the cell of the first of them; and the sites of the path, in the order in which it passes through them, split into $\lceil\frac{|\omega|+1}M\rceil\le\frac{|\omega|}M+1$ runs of at most $M$ consecutive sites, each of which therefore meets at most $27$ cells. This proves \textup{(Cnt)}.

For a set $\mathsf Y$ of cells, let $\mathfrak S$ be the family of $\mathsf Y$ (item~C.II~1.5\,(c)), and let $K^{\mathsf V}_{[\mathsf Y]}(\zeta):=K^{\mathsf V}_{\mathfrak S}(\zeta)$, where $K^{\mathsf V}_{\mathfrak S}(\zeta)$ is the operator of item~C.II~1.5. We write $K_{[\mathsf Y]}=K_{[\mathsf Y]}(1)$, and $K_{[\mathsf Y]}[\mathsf V']$ for $K_{[\mathsf Y]}$ at the field $\mathsf V'$. For a non-empty finite set $Y$ of cells, for $k\ge1$ and for $j\ge0$, we define the traces sorted by sets of cells, \[T_k(Y;\mathsf V'):=\sum_{\mathsf Y\subset Y}(-1)^{\#Y\setminus\mathsf Y}\operatorname{tr}(K^{\mathsf V}_{[\mathsf Y]})^k=\sum_jT_{k,j}(Y;\mathsf V'),\qquad T_{k,j}:=\sum_{\mathsf Y\subset Y}(-1)^{\#Y\setminus\mathsf Y}\bigl[\operatorname{tr}K^{\mathsf V}_{[\mathsf Y]}(\zeta)^k\bigr]_j,\] where $[\,\cdot\,]_j$ stands for the coefficient of $\zeta^j$ in the polynomial $\operatorname{tr}K^{\mathsf V}_{[\mathsf Y]}(\zeta)^k$, and we put \[\mathsf D^q(Y,\mathsf V'):=\sum_{k\ge1}\tfrac1k\bigl[T_k(Y;\mathsf V')-T_k(Y;\mathbf 1)\bigr].\]
\end{definition}

\emph{Conditions on the letters.} We use the following conditions.
\begin{enumerate}
\item[\textup{(M4$^\flat$)}] $M\ge\frac{54}\mu(\kappa^\sharp+3)=216\,\frac b\gamma\,\frac{\log\Lambda}{\theta_2}\,(\kappa^\sharp+3)$.
\item[\textup{(R$^\flat$)}] $c_R\,\mu\ge1.35$.
\end{enumerate}

The condition \textup{(M4$^\sharp$)}, that is $M\ge216\frac b\gamma(\kappa^\sharp+3)$, is a consequence of \textup{(M4$^\flat$)}. Indeed, the right side of \textup{(M4$^\flat$)} is that of \textup{(M4$^\sharp$)} multiplied by the factor $\frac{\log\Lambda}{\theta_2}\ge2\log16\ge5.5$; this lower bound holds since $\theta_2\le\frac12$ and $A^\flat\ge14b$, and the latter inequality holds since $2e^2\ge14$ and $384e^8\mathsf g_2(\frac\lambda2)^2\ge14\cdot48\,\mathsf g(\lambda)^2$. The numbers $\Lambda,\theta_2,\mu$ are computed from $(L,a,K,\lambda)$ before $M$ is chosen, and they do not depend on $M$.

If $\lambda L\ge12$, then \textup{(R$^\flat$)} is a consequence of $c_R\theta_{\mathbb C}\ge128$. In that case $\theta_2=\frac12$, and with $t=A^\flat/b$ we get $c_R\ge768/\theta^\flat\ge1536\,t\frac b\gamma\ge10.8\log(2+t)\frac b\gamma$, which is the lower bound on $c_R$ to which \textup{(R$^\flat$)} reduces when $\theta_2=\frac12$. At the starred values the inequality $\lambda_*L\ge12$ holds, because $L\ge L_A\ge36\,\mathsf g(\frac{\lambda_*}2)/\gamma_*$, $\lambda_*=\gamma_*^{1/2}/44$ and $\gamma_*\le1$, and because $\mathsf g(\lambda)\ge1+26e^{-1}+98e^{-2}>16$ for $0<\lambda\le1$, so that $\mathsf g(\frac{\lambda_*}2)\ge16$. The starred values and $L_A$ are those of \S\ref{Y-A-1}.

Further, we use the following condition.
\begin{enumerate}
\item[\textup{(M2)}] $M\ge1920(\kappa+4)c_R=192\kappa^\sharp c_R$.
\end{enumerate}

We define the number \[\mathsf K_1^\flat:=18\,(LM)^3\Bigl[\frac\Lambda{1-e^{-\theta_2}}+\frac{2b}\gamma\Bigr]e^{27(\kappa^\sharp+3+\kappa_+)}.\]

\emph{Floors.} We use the following floors \textup{(S1)} and \textup{(S2)} on $x$.
\begin{enumerate}
\item[\textup{(S1)}] The floor \textup{(S1)} consists of four clauses, in this order: $p_0'\ge1$; $r'\ge1$; $18M(r'+1)\varepsilon'\le\frac1{10}$; $4\bigl(170\,M(r'+1)\log x'/\vartheta_0^\flat\bigr)^2\varepsilon'\le1$.
\item[\textup{(S2)}] The floor \textup{(S2)} has six clauses. The first four are $3L^3(2R_2+1)^3\frac4\gamma\le x$; $(\kappa^\sharp+1)(3r_u+3)\le\log x$; $18(LM)^3\le x$; $4(b+A^\flat)e^{27\kappa^\sharp}\le x$. The last two are $r_u\le r'$ and $x^{-(\mathsf n-3)}\le\varepsilon'$, where $\mathsf n\ge4$.
\end{enumerate}

Under \textup{(M2)} we have $3r_u+3\le\frac{6c_R\log x+12}M+6$ and $(\kappa^\sharp+1)\frac{6c_R}M\le0.04$, so that the second clause of \textup{(S2)} holds as soon as $\log x\ge11(\kappa^\sharp+1)$. Thus, if \textup{(M2)} is assumed, every clause of \textup{(S1)} and of \textup{(S2)} holds for all large $x$.

The next theorem shows that $\log\det q_\ell^V$ is local in $V'$, at a level whose rows may have any depth: the rows are those of item~C.II~1.2 with $K\ge0$, so that $\mathsf H_y=\{0\}$ is admitted.

\begin{theorem}[The determinant of the form is local in the coarse field, at a level with rows]\label{Y-C-CIITh73}
Let $L\ge3$ be odd with $aL\ge18$, let the rows be as in item~C.II~1.2, and let $\kappa\ge6$; we assume \textup{(M4$^\flat$)}, \textup{(R$^\flat$)} and \textup{(Tor)}. Then the following hold.
\begin{enumerate}
\item[\textup{(a)}] At every $V'\in\mathcal V'$ we have $\log\det q^{\mathbf 1}_\ell-\log\det q^V_\ell=\sum_Y\mathsf D^q(Y,V')$. Moreover, $T_{k,j}(Y;\cdot)=0$ unless $Y$ is connected and $\#Y\le27(\frac jM+1)$; in particular $\mathsf D^q(Y,\cdot)=0$ if $Y$ is not connected.
\item[\textup{(b)}] Every gauge transformation $u$ with values in $G$ leaves $\mathsf D^q(Y,\cdot)$ unchanged; so does every $u$ taking its values in $\operatorname{SL}(2,\mathbb C)$, provided that $\mathsf V'$ and $\mathsf V'^u$ both belong to the domain defined at the end of this clause. The function $\mathsf D^q(Y,\cdot)$ is real at $G$-valued fields, and it depends on the field only through $\operatorname{Ad}_{\mathsf V'(c)}$ for the bonds $c$ of $Y^\sim$. It is holomorphic on the domain $\{|\mathsf V'(c)|<1+\vartheta_0^\flat$ on the bonds of $Y^\sim\}$, and it vanishes at those pure gauges $\mathbf 1^u$, with $u$ taking values in $G$ or in $\operatorname{SL}(2,\mathbb C)$, that belong to this domain.
\item[\textup{(c)}] At every $V'\in\mathcal V'$, and at every field of this domain, $|\mathsf D^q(Y,\mathsf V')|\le\mathsf K_1^\flat\,e^{-(\kappa^\sharp+\kappa_+)\#Y}$.
\item[\textup{(d)}] If \textup{(S1)} holds, then on $\mathcal U_Y(x')$, and hence at every $G$-valued field that is tame on $Y^\approx$, we have $|\mathsf D^q(Y,\mathsf V')|\le\mathsf K_1^\flat\,\varepsilon'\,e^{-\kappa^\sharp(\#Y-1)}$.
\end{enumerate}
\end{theorem}

\begin{proof}
Throughout the proof $n:=\#Y$ and $j_0:=\max\{0,M(\frac n{27}-1)\}$, and we put $\mathsf d_Y:=9(LM)^3n$; for every $\mathsf Y\subset Y$ this number bounds the dimension of $\ell^2(I_{\mathsf Y};\mathbb C^3)$.

\emph{Step (a): sums over sets of cells.} By \Cref{Y-C-CIIL65}\,(a), $\|K_{[\mathsf Y]}\|\le\varkappa<1$; consequently $-\log\det(q_{[\mathsf Y]}/b)=\sum_k\frac1k\operatorname{tr}K_{[\mathsf Y]}^k$. For the set $\mathsf Y=\mathfrak c$ of all cells, $q_{[\mathfrak c]}=q_\ell$. M\"obius inversion states that $\sum_{Y\subset\mathfrak c}\sum_{\mathsf Y\subset Y}(-1)^{\#Y\setminus\mathsf Y}f(\mathsf Y)=f(\mathfrak c)$, and we use it for functions with $f(\emptyset)=0$: applied to $-\log\det(q_{[\mathsf Y]}/b)$, at $\mathsf V'$ and at $\mathbf 1$, it yields, through the definition of $\mathsf D^q$, the identity of \textup{(a)}. Next we expand the $k$-th power and write $\operatorname{tr}K_{[\mathsf Y]}(\zeta)^k$ as a sum indexed by the closed walks $\omega$ of $k$ steps with $\mathsf c(\omega)\subset\mathsf Y$, the term of $\omega$ being $\zeta^{|\omega|}$ multiplied by a number and by $\operatorname{tr}_{\mathfrak g}\operatorname{Ad}_{\mathsf V'(\pi_\omega)}$. This rests on the fact that a step lies in the family of $\mathsf Y$ precisely when $\mathsf c(\mathsf s)\subset\mathsf Y$ (item~C.II~1.5\,(c)). For a fixed walk, $\sum_{\mathsf c(\omega)\subset\mathsf Y\subset Y}(-1)^{\#Y\setminus\mathsf Y}=1\{\mathsf c(\omega)=Y\}$. Therefore only the walks with $\mathsf c(\omega)=Y$ and $|\omega|=j$ contribute to $T_{k,j}(Y;\cdot)$; for these, $Y$ is connected, and \textup{(Cnt)} gives $\#Y\le27(\frac jM+1)$.

\emph{Step (b): invariance, dependence and holomorphy.} Let $u$ be a gauge transformation with values in $\operatorname{SL}(2,\mathbb C)$. For a step $\mathsf s:b\to b'$ we have the identity of matrices $\mathsf T_{\mathsf s}[\mathsf V'^u]=\operatorname{Ad}_{u(y(b))}\mathsf T_{\mathsf s}\operatorname{Ad}_{u(y(b'))}^{-1}$, so that $K_{[\mathsf Y]}[\mathsf V'^u]=\operatorname{Ad}_uK_{[\mathsf Y]}[\mathsf V']\operatorname{Ad}_u^{-1}$, and the traces at the two fields coincide for every $k$ and every $\mathsf Y\subset Y$. By \Cref{Y-C-CIIL65}\,(a), the series defining $\mathsf D^q$ converges at each of the two fields $\mathsf V'$, $\mathsf V'^u$ as soon as that field belongs to the domain; this proves the invariance. For the unit field the same identity reads $T_k(Y;\mathbf 1^u)=T_k(Y;\mathbf 1)$, and so $\mathsf D^q(Y,\cdot)$ vanishes at the pure gauges of the domain. At $G$-valued fields the function $\mathsf D^q(Y,\cdot)$ is real: for $g\in G$ the map $\operatorname{Ad}_g$ is real on $\mathfrak g$, being real orthogonal in the real basis $\sigma_1,\sigma_2,\sigma_3$, and $q_{\mathsf s}$ is real, so that $K_{[\mathsf Y]}$ and every trace $\operatorname{tr}K_{[\mathsf Y]}^k$ are real. As to the dependence on the field, a step of kind \textup{(R)} depends only on bonds with both ends in cells of $\mathsf Y$, a step of kind \textup{(P)} on at most two bonds from $y(p)$, each within $1$ of $y(b)$, and a step of kind \textup{(Q)} on one bond from $c_-$ within $1$ of $y(b)$; this gives the dependence on the bonds of $Y^\sim$ alone. Finally, again by \Cref{Y-C-CIIL65}\,(a), the series converges uniformly on the domain, and its terms are polynomials in the entries of the matrices $\operatorname{Ad}^{\pm1}_{\mathsf V'(c)}$; hence $\mathsf D^q(Y,\cdot)$ is holomorphic there.

\emph{Step (c): the bound on the domain.} By \textup{(a)}, $T_k=\sum_{j\ge j_0}T_{k,j}$, and we bound $T_k$ in two ways at every field of the domain. The first bound, $|T_k(Y)|\le2^n\mathsf d_Y\varkappa^k$, comes from \Cref{Y-C-CIIL65}\,(a). For the second, \Cref{Y-C-CIIL65}\,(b) gives $|T_{k,j}(Y)|\le2^n\mathsf d_Y\Lambda^ke^{-\theta_2j}$, and summation over $j\ge j_0$ yields $|T_k(Y)|\le2^n\mathsf d_Y\Lambda^k\frac{e^{-\theta_2j_0}}{1-e^{-\theta_2}}$. Let $k_1:=\lfloor\frac{\theta_2j_0}{2\log\Lambda}\rfloor$. For $k\le k_1$ we use the second bound, together with $\sum_{k\le k_1}\frac1k\Lambda^k\le\Lambda^{k_1+1}\le\Lambda e^{\theta_2j_0/2}$, which holds as $\Lambda\ge2$. For $k>k_1$ we use the first, together with $\sum_{k>k_1}\frac1k\varkappa^k\le\frac{2b}\gamma\varkappa^{k_1+1}\le\frac{2b}\gamma e^{-\frac\gamma{2b}\frac{\theta_2j_0}{2\log\Lambda}}=\frac{2b}\gamma e^{-\mu j_0}$. Since $\mu\le\frac12\theta_2$, we add the bounds for $k\le k_1$ and for $k>k_1$ and obtain $\sum_k\frac1k|T_k(Y)|\le2^n\mathsf d_Y[\frac\Lambda{1-e^{-\theta_2}}+\frac{2b}\gamma]e^{-\mu j_0}$.

Now \textup{(R$^\flat$)} gives $\kappa_+=\frac M{40c_R}\le\frac{\mu M}{54}$, and \textup{(M4$^\flat$)} gives $\frac{\mu M}{54}\ge\kappa^\sharp+3$; therefore $\frac{\mu M}{27}\ge\kappa^\sharp+3+\kappa_+$. If $n\le27$, then $j_0=0$, so $e^{-\mu j_0}=1\le e^{(\kappa^\sharp+3+\kappa_+)(27-n)}$. If $n>27$, then $j_0=\frac M{27}(n-27)$, so $e^{-\mu j_0}=e^{-\frac{\mu M}{27}(n-27)}\le e^{-(\kappa^\sharp+3+\kappa_+)(n-27)}$. In both cases $e^{-\mu j_0}\le e^{27(\kappa^\sharp+3+\kappa_+)}e^{-(\kappa^\sharp+3+\kappa_+)n}$. As $\mathsf d_Y=9(LM)^3n$ and $n2^ne^{-3n}\le1$, the sum $\sum_k\frac1k|T_k(Y)|$ is at most one half of the bound in \textup{(c)}. Since $\mathsf D^q$ is defined in terms of the values of $T_k$ at the two fields $\mathsf V'$ and $\mathbf 1$, both of which belong to the domain, \textup{(c)} follows.

\emph{Step (d): the bound on $\mathcal U_Y(x')$.} Suppose first that $n>r'$. Then $e^{-\kappa_+n}\le e^{-\kappa_+r'}\le e^{-\frac12\log x'}\le\varepsilon'$, where we use $r'\ge\frac{20c_R}M\log x'$ and $p_0'\ge1$. A member of $\mathcal U_Y(x')$ satisfies \textup{(u1)}, and $\vartheta(x')<\vartheta_0^\flat$; hence it belongs to the domain of \textup{(b)}, and \textup{(c)} yields the bound.

Suppose next that $n\le r'$. We may take $Y$ connected, since otherwise $\mathsf D^q(Y,\cdot)=0$ by \textup{(a)}. Fix a cell $\square\in Y$ and a site $o$ of $\square$. Then $Y^\sim\subset\square^{+n}\subset\mathsf B:=$ the sites of $\square^{+r'}\subset Y^\approx$; here $\mathsf B$ is a box, embedded by \textup{(Tor)}, all of whose sites lie within $\omega_{\mathsf B}:=M(r'+1)$ of $o$, and $3\omega_{\mathsf B}+1\le\ell_w(x')$. Let $\Omega:=\{$fields on the bonds of $\mathsf B$ with $|\mathsf V(c)|<1+\vartheta_0^\flat\}$. By \textup{(b)} and \textup{(c)}, the function $\Phi:=\mathsf D^q(Y,\cdot)$ is holomorphic on $\Omega$, and $|\Phi|\le\mathsf S:=\mathsf K_1^\flat e^{-\kappa^\sharp n}$.

Let $\mathsf V'\in\mathcal U_Y(x')$. By \textup{(u2)} and \textup{(u3)}, it satisfies \textup{(h2)} and \textup{(h3)} on $\mathsf B$, with $\hat\varepsilon=\varepsilon'$. The hypothesis $\hat\varepsilon\le\frac1{40}$ of \Cref{Y-C-L67} reads $\varepsilon'\le\frac1{40}$ here; this follows from the third clause of \textup{(S1)}, which together with $M\ge3$ and $r'\ge1$ even gives $\varepsilon'\le\frac1{1080}$. \Cref{Y-C-L67}\,(a) and \textup{(S1)} then give $|\mathsf W(c)-\mathbf 1|\le19\omega_{\mathsf B}\varepsilon'$. We set $A'_c=-i\log\mathsf W(c)$. This matrix satisfies $|A'_c|\le\mathsf a:=21\omega_{\mathsf B}\varepsilon'$, and it is traceless because $\det\mathsf W(c)=1$ and $|\mathsf W(c)-\mathbf 1|<1$. For complex $t$ we consider the field $\mathsf V'_t(c)=\mathsf k(c_-)^{-1}e^{itA'_c}\mathsf k(c_+)=\mathsf V'(c)(\mathbf 1+\delta_c)$, where $\mathbf 1+\delta_c=\mathsf k(c_+)^{-1}e^{i(t-1)A'_c}\mathsf k(c_+)$. By \Cref{Y-C-L67}\,(a) we have $|\mathsf k|\le2$; hence, if $(T+1)\mathsf a\le1$, then $|\delta_c|\le4(e^{(T+1)\mathsf a}-1)\le4(e-1)(T+1)\mathsf a\le8(T+1)\mathsf a$ for $|t|<T$. We choose $T+1:=\vartheta_0^\flat/(170\,\omega_{\mathsf B}\varepsilon'\log x')$. Then $(T+1)\mathsf a=\frac{21}{170}\frac{\vartheta_0^\flat}{\log x'}\le1$, since $\vartheta_0^\flat\le10^{-4}$ (item~C.II~1.3) and $\log x'>1$, the latter because $x>Le$ and $x'=x/L$ (item~C.II~1.1); so the bound on $|\delta_c|$ holds for $|t|<T$, and $1.01\cdot8(T+1)\mathsf a\le\frac{169.7}{170}\frac{\vartheta_0^\flat}{\log x'}$. Hence, by \textup{(u1)}, $|\mathsf V'_t(c)|<(1+|\delta_c|)(1+\vartheta_0^\flat-\frac{\vartheta_0^\flat}{\log x'})\le1+\vartheta_0^\flat$, that is, $\mathsf V'_t\in\Omega$. Membership in $\Omega$ depends neither on the curvature of $\mathsf V'_t$ nor on its values along paths.

The function $f(t):=\Phi(\mathsf V'_t)$ is therefore holomorphic on the disc $|t|<T$, and $|f|\le\mathsf S$ there. By \textup{(b)}, $f(0)=0$, because $\mathsf V'_0=\mathbf 1^{\mathsf k^{-1}}$. For small $|t|$ we also have $\mathsf Z_t:=(e^{itA'_c})\in\Omega$, and $f(t)=\Phi(\mathsf Z_t^{\mathsf k^{-1}})=\Phi(\mathsf Z_t)$ by \textup{(b)}: the gauge $\mathsf k^{-1}$ takes its values in $\operatorname{SL}(2,\mathbb C)$, and $\mathsf Z_t$ and $\mathsf Z_t^{\mathsf k^{-1}}=\mathsf V'_t$ both lie in $\Omega$, hence in the domain of \textup{(b)}. Moreover $\frac d{dt}\Phi(\mathsf Z_t)|_0=0$. Indeed, this derivative, which we denote by $\Lambda(A')$, is $\mathbb C$-linear in $A'$, and, by \textup{(b)}, $\Lambda(\operatorname{Ad}_uA')=\Lambda(A')$ for every constant $u\in G$; hence $\Lambda(A')=\Lambda(\int_G\operatorname{Ad}_uA'\,du)$, whereas the matrix $\int_G\operatorname{Ad}_uA'_c\,du$, for each bond $c$, commutes with $G$ and is traceless, so that it is $0$.

Consequently $f(t)/t^2$ is holomorphic on the disc, and the maximum principle gives $|f(1)|\le\mathsf ST^{-2}$, once $t=1$ is known to lie in the disc. By the last clause of \textup{(S1)}, $(T+1)^{-2}\le\frac14\varepsilon'$, and the third clause of \textup{(S1)} gives $\varepsilon'<1$, so that $T>1$ and $t=1$ lies in the open disc $|t|<T$. Then $T\ge\frac12(T+1)$, and so $T^{-2}\le\varepsilon'$. Since $\mathsf V'_1$ coincides with $\mathsf V'$ on the bonds of $\mathsf B$, the value $f(1)$ equals $\mathsf D^q(Y,\mathsf V')$, and the bound of \textup{(d)} follows.
\end{proof}

\begin{definition}[The $\mathsf u$-steps and the terms $\mathsf D^u$, $S^\circ$, $S_k$, $\mathsf D_k$]\label{Y-C-CIITheRest}
Let $\mathsf u=\mathsf u^{\mathsf V}=1-C^{\mathsf V}q_\ell^{\mathsf V}$. If \textup{($\rho$1)}, \textup{($\rho$2)} and \textup{(P0)} hold, then $\mathsf u(b,b')=0$ for ${\rm dist}(b,b')>R+r_q$, and $R+r_q\le R_2$. Moreover, $\mathsf u(b,\cdot)$ depends on $\mathsf V'$ only through its values on the coarse bonds within $R_2$ of $y(b)$. Indeed, $C(b,\cdot)$ has this property with $R_2-1$ in place of $R_2$, while for ${\rm dist}(b,b'')\le R$ the row $q(b'',\cdot)$ has it with $r_q$ and $y(b'')$ in place of $R_2$ and $y(b)$, by the bound of item~C.II~1.2 that is measured from $y(b'')$ alone: every site of $\pi_{\mathsf s}$, for a step $\mathsf s$ from $b''$, lies within $r_q$ of $y(b'')$.

A $\mathsf u$-step $b\to b'$ carries the set $\mathsf c(b,b'):=$ the cells met by $\Gamma(y(b),y(b'))$; this set is connected and has at most $\mathsf n_u:=3\lceil R_2/M\rceil+3=3r_u+3$ members, since the coordinate path is made of three straight stretches of at most $R_2$ bonds each. Under the hypotheses of \Cref{Y-C-CIITh74} and at every $\mathsf V'\in\mathfrak V$, each row sum and each column sum of $\|\mathsf u(b,b')\|$ is at most $\upsilon_1:=3L^3(2R_2+1)^3\frac4\gamma x^{-\mathsf n}$. This follows from \Cref{Y-C-L66}\,(a): the lemma assumes \textup{($\lambda$0)}, \textup{($\rho$5)}, $\frac2\gamma x^{-\mathsf n}\le\frac12$ and $\mathsf V'\in\mathfrak V$, and \textup{(P1)} gives $\frac4\gamma x^{-\mathsf n}\le1$.

Let $\mathsf V'\in\mathfrak V$, and let $Y$ be a connected set of cells. Under the hypotheses of \Cref{Y-C-CIITh74} we have $\upsilon_1\le\frac12$ by \textup{(S2)}, so that the series in the following definitions converge absolutely: \[\mathsf D^u(Y,\mathsf V'):=-\sum_{m\ge1}\frac1m\sum_{\substack{b_0,\dots,b_{m-1};\ b_m=b_0\\\bigcup\mathsf c(b_i,b_{i+1})=Y}}\operatorname{tr}_{\mathfrak g}\prod_{i=0}^{m-1}\mathsf u(b_i,b_{i+1}),\] \[S^\circ(Y;\mathsf V')(b,b'):=\sum_{m\ge1}\sum_{\substack{\mathsf s:b\to b_1;\ b_2,\dots,b_{m+1}=b'\\\mathsf c(\mathsf s)\cup\bigcup\mathsf c(b_i,b_{i+1})=Y}}q_{\mathsf s}\mathsf T_{\mathsf s}\prod_{i=1}^m\mathsf u(b_i,b_{i+1}),\] and \[S_k(Y;\mathsf V'):=\tfrac12\bigl(S^\circ(Y;\mathsf V')+S^\circ(Y;\mathsf V')^{\mathsf T}\bigr),\qquad\mathsf D_k(Y,\mathsf V'):=-\mathsf D^q(Y,\mathsf V')-\bigl[\mathsf D^u(Y,\mathsf V')-\mathsf D^u(Y,\mathbf 1)\bigr].\]
\end{definition}

\begin{theorem}[The terms $\mathsf D^u(Y,\cdot)$ and $S_k(Y;\cdot)$ are local in the coarse field]\label{Y-C-CIITh74}
Let $L\ge3$ be odd with $aL\ge18$, let the rows be as in item~C.II~1.2, and let $\kappa\ge6$; we take a law with \textup{($\lambda$0)}, \textup{($\rho$1)} to \textup{($\rho$6)}, and we assume \textup{(M4$^\flat$)}, \textup{(P0)}, \textup{(P1)} and \textup{(S2)}. Then the following hold at every $\mathsf V'\in\mathfrak V$, and so in particular at every $V'\in\mathcal V'$, with $Y$ ranging over the connected sets of cells.
\begin{enumerate}
\item[\textup{(a)}] At $V'\in\mathcal V'$, $\log\det C^V=-\log\det q^V_\ell+\sum_Y\mathsf D^u(Y,V')$; moreover $S^{\mathsf V}=\sum_YS_k(Y;\mathsf V')$; both sums converge absolutely.
\item[\textup{(b)}] The functions $\mathsf D^u(Y,\cdot)$ and $S_k(Y;\cdot)$ depend on $\mathsf V'$ only through its values on the coarse bonds within $R_2$ of the sites of $Y$, and hence only through its values on bonds of $Y^{+r_u}$. They are holomorphic on $\mathfrak V^\circ$ and real at $G$-valued fields. For $G$-valued $u$, $\mathsf D^u(Y,\mathsf V'^u)=\mathsf D^u(Y,\mathsf V')$ and $S_k(Y;\mathsf V'^u)=\operatorname{Ad}_uS_k(Y;\mathsf V')\operatorname{Ad}_u^{-1}$. Finally, $S_k(Y)^{\mathsf T}=S_k(Y)$, and $S_k(Y)(b,b')=0$ unless $b,b'\in\hat Y$.
\item[\textup{(c)}] $|\mathsf D^u(Y,\mathsf V')|\le x^{-(\mathsf n-3)}e^{-\kappa^\sharp\#Y}$ and $\|S_k(Y;\mathsf V')\|\le x^{-(\mathsf n-3)}e^{-\kappa^\sharp\#Y}$.
\end{enumerate}
\end{theorem}

\begin{proof}
By \textup{(S2)}, $\upsilon_1\le x^{-(\mathsf n-1)}\le\frac12$ and $e^{(\kappa^\sharp+1)\mathsf n_u}\le x$.

\emph{Step (a): the two expansions.} We have $\det C\det q=\det(1-\mathsf u)$ with $\|\mathsf u\|<1$, so that $\log\det(1-\mathsf u)=-\sum\frac1m\operatorname{tr}\mathsf u^m$, a real number by \textup{($\rho$6)} and \Cref{Y-C-CIIL64}\,(a). We write each trace as a sum over closed walks of $\mathsf u$-steps and sort these walks by the union $\bigcup\mathsf c$ of their sets of cells; this gives the first identity. Likewise $S=q\sum_{m\ge1}\mathsf u^m$ by \Cref{Y-C-L66}\,(a); we sort the walks in the same manner, and since $S^{\mathsf T}=S$ by \textup{($\rho$6)}, the second identity follows.

\emph{Step (b): dependence, holomorphy and gauge behaviour.} The dependence on the field is read off from the walks and from the clause \textup{($\rho$2)} of the law. By \textup{($\rho$3)}, $\mathsf u[\mathsf V'^u]=\operatorname{Ad}_u\mathsf u\operatorname{Ad}_u^{-1}$, while the sets $\mathsf c(\cdot)$ depend on no field; this gives the behaviour under gauge transformations. Holomorphy follows from \textup{($\rho$4)} and the uniform bounds of \textup{(c)}. At $G$-valued fields the matrix $C$ is real by \textup{($\rho$6)}, and $q_\ell$ is real, since $q_{\mathsf s}$ is real and $\operatorname{Ad}$ is real on $\mathfrak g$; hence $\mathsf u$ is real, and so are $\mathsf D^u(Y,\cdot)$, $S^\circ(Y;\cdot)$ and $S_k(Y;\cdot)$. The definition of $S_k$ gives $S_k(Y)^{\mathsf T}=S_k(Y)$. Finally, in a walk that contributes to $S^\circ(Y)(b,b')$, the cells of $b$ and $b'$ lie in $\mathsf c(\mathsf s)$ and in $\mathsf c(b_m,b')$, respectively, and so in $Y$.

\emph{Step (c): the bounds.} Let $n=\#Y$. A closed walk of $m$ $\mathsf u$-steps whose cells make up $Y$ has $n\le m\mathsf n_u$. Its first bond satisfies $b_0\in\hat Y$, which leaves at most $3(LM)^3n$ choices; moreover $|\operatorname{tr}_{\mathfrak g}|\le3\|\cdot\|$, and the sum over the remaining bonds is controlled by the row sums. Hence $|\mathsf D^u(Y)|\le9(LM)^3n\sum_{m\ge m_0}\upsilon_1^m\le18(LM)^3n\,\upsilon_1^{m_0}$, where $m_0=\max\{1,\lceil n/\mathsf n_u\rceil\}$. Here $\upsilon_1^{m_0}\le x^{-(\mathsf n-2)}x^{-m_0}$ and $x^{-m_0}\le e^{-(\kappa^\sharp+1)\mathsf n_um_0}\le e^{-(\kappa^\sharp+1)n}$; together with $18(LM)^3ne^{-n}\le x$ this gives the bound on $|\mathsf D^u(Y,\mathsf V')|$ stated in \textup{(c)}.

For $S^\circ$, the set $Y$ of a walk obeys $\#Y\le\#\mathsf c(\mathsf s)+m\mathsf n_u$, where $\#\mathsf c(\mathsf s)\le27(\frac{|\mathsf s|}M+1)$ by \textup{(Cnt)}; moreover $\|\mathsf T_{\mathsf s}\|\le e^{2\vartheta_0^\flat|\mathsf s|}$. By \textup{(M4$^\flat$)}, $\frac{27\kappa^\sharp}M\le\frac18\theta_2$, and so $2\vartheta_0^\flat+\frac{27\kappa^\sharp}M\le\theta_1$; hence \Cref{Y-C-CIIL64}\,(b) yields $\sum_{\mathsf s\text{ from }b}|q_{\mathsf s}|\|\mathsf T_{\mathsf s}\|e^{\kappa^\sharp\#\mathsf c(\mathsf s)}\le e^{27\kappa^\sharp}(b+A^\flat)$. Also $\sum_{m\ge1}(\upsilon_1e^{\kappa^\sharp\mathsf n_u})^m\le2x^{-(\mathsf n-2)}$. Combining these, $e^{\kappa^\sharp\#Y}\sum_{b'}\|S^\circ(Y)(b,b')\|\le2(b+A^\flat)e^{27\kappa^\sharp}x^{-(\mathsf n-2)}\le\frac12x^{-(\mathsf n-3)}$. The column sums obey the same bound; for them one reverses the steps and uses the column sums of $\mathsf u$. The Schur test now bounds the norm of $S^\circ(Y)$, the norm of its transpose $S^\circ(Y)^{\mathsf T}$ obeys the same bound, and the bound on $\|S_k(Y;\mathsf V')\|$ stated in \textup{(c)} follows.
\end{proof}

\begin{corollary}[\textup{(S-loc$_k$)}: the determinant and $S^V$ as sums of local terms]\label{Y-C-CIICor75}
Let the hypotheses of Thms.~\ref{Y-C-CIITh73} and~\ref{Y-C-CIITh74} hold, assume \textup{(S1)} in addition, and consider the families $\mathsf D_k(Y,\cdot)$ and $S_k(Y;\cdot)$, where $Y$ ranges over the connected sets of cells. Then, at every $V'\in\mathcal V'$, \[\log\det C^{\mathbf 1}-\log\det C^V=\sum_Y\mathsf D_k(Y,V'),\qquad S^V=\sum_YS_k(Y;V').\] Each $\mathsf D_k(Y,\cdot)$ is real and gauge invariant, takes the value $0$ at pure gauges, and depends on $V'$ only through the bonds of $Y^{+r_u}\subset Y^\approx$. At every $V'$ it obeys $|\mathsf D_k(Y,V')|\le2\mathsf K_1^\flat e^{-\kappa^\sharp(\#Y-1)}$, and it is continuous on $\mathcal V'$. It continues to a holomorphic function on $\mathcal U_Y(x')$, where $|\mathsf D_k|\le3\mathsf K_1^\flat\,\varepsilon'\,e^{-\kappa^\sharp(\#Y-1)}$, and this continuation is invariant under the $G$-valued gauge transformations on the sites of $Y^\approx$. Finally, $\|S_k(Y;V')\|\le x^{-(\mathsf n-3)}e^{-\kappa^\sharp\#Y}$, and $S_k(Y;\cdot)$ has the properties of \Cref{Y-C-CIITh74}\,(b).
\end{corollary}

\begin{proof}
The two identities follow from \Cref{Y-C-CIITh73}\,(a) and \Cref{Y-C-CIITh74}\,(a), applied at $V'$ and at $\mathbf 1$. At a pure gauge $\mathbf 1^u$ we have $\mathsf D^q=0$ and $\mathsf D^u(Y,\mathbf 1^u)=\mathsf D^u(Y,\mathbf 1)$, so that $\mathsf D_k(Y,\cdot)$ vanishes there. A member of $\mathcal U_Y(x')$, extended by $\mathbf 1$ off $Y^\approx$, is a member of $\mathfrak V^\circ$ by \textup{(u1)}, and by \textup{(S2)} the bonds on whose values $\mathsf D^q(Y,\cdot)$ and $\mathsf D^u(Y,\cdot)$ depend lie in $Y^{+r_u}\subset Y^\approx$; this yields the continuation to $\mathcal U_Y(x')$. As $2x^{-(\mathsf n-3)}\le2\varepsilon'\le\mathsf K_1^\flat\varepsilon'$, the two bounds on $|\mathsf D_k|$ stated in the corollary follow from \Cref{Y-C-CIITh73}\,(c), (d) and \Cref{Y-C-CIITh74}\,(c).

Continuity on $\mathcal V'$ holds because $\mathsf D^q(Y,\cdot)$ is holomorphic on the open set of \Cref{Y-C-CIITh73}\,(b) and $\mathsf D^u(Y,\cdot)$ is holomorphic on $\mathfrak V^\circ$ by \Cref{Y-C-CIITh74}\,(b); both sets contain every $G$-valued field, and $\mathsf D^u(Y,\mathbf 1)$ is a number.

For the invariance on $\mathcal U_Y(x')$, let $\mathsf V\in\mathcal U_Y(x')$ and let $u$ be $G$-valued on the sites of $Y^\approx$. Then $\mathsf V^u\in\mathcal U_Y(x')$, since $|u_1gu_2|=|g|$ and $|u_1gu_1^{-1}-\mathbf 1|=|g-\mathbf 1|$ for unitary $u_1,u_2$. We extend $u$ by $\mathbf 1$. On the bonds that leave $Y^\approx$ the transform of the extended field then takes unitary values, which places it in $\mathfrak V^\circ$. Only on these bonds can it differ from the extension of $\mathsf V^u$; none of them is a bond of $Y^{+r_u}$, and the terms do not depend on them. \Cref{Y-C-CIITh73}\,(b) and \Cref{Y-C-CIITh74}\,(b), both applied with $G$-valued $u$, now give the invariance.
\end{proof}

\begin{remark}[The special case of the law of windows]\label{Y-C-CIIAtTheLawOfWindows}
For the law of windows, where $\mathsf n=60$, we have $\|S_k(Y;\mathsf V')\|\le x^{-57}e^{-\kappa^\sharp\#Y}$ and $|\mathsf D^u(Y,\mathsf V')|\le x^{-57}e^{-\kappa^\sharp\#Y}$ on $\mathfrak V$. Let $\tilde{\mathsf D}(Y,V'):=-\frac12\mathsf D_k(Y,V')$. Then $\det G^V=(\det C^V)^{1/2}=(\det C^{\mathbf 1})^{1/2}\exp\sum_Y\tilde{\mathsf D}(Y,V')$, where $\det G^V>0$ by \Cref{Y-C-CITh66}\,(h); moreover $|\tilde{\mathsf D}|\le\mathsf K_1^\flat e^{-\kappa^\sharp(\#Y-1)}$ at every $V'$, and $|\tilde{\mathsf D}|\le1.5\,\mathsf K_1^\flat\varepsilon'e^{-\kappa^\sharp(\#Y-1)}$ on $\mathcal U_Y(x')$.

\emph{The conditions, whole.} The conditions of Thms.~\ref{Y-C-CIITh73}, \ref{Y-C-CIITh74} and \Cref{Y-C-CIICor75} are $\kappa\ge6$, \textup{(M4$^\flat$)}, \textup{(R$^\flat$)}, \textup{(M2)}, \textup{(Tor)} and the floors \textup{(F$_{\rm law}$)}, \textup{(P0)}, \textup{(P1)}, \textup{(S1)}, \textup{(S2)}. The condition \textup{(M2)} gives $6c_R(\kappa^\sharp+1)/M\le0.04$. Since $2c_R\log x/M\le r_u\le(2c_R\log x+4)/M+1$, the second clause of \textup{(S2)} holds for all large $x$ if $M>6c_R(\kappa^\sharp+1)$, and for no $x$ if $M\le6c_R(\kappa^\sharp+1)$.

The ratio in the formula of \Cref{Y-C-P67}\,(b) is the same function whether we write it in terms of $\varphi$ or of $\xi$; here $\varphi=G\xi$, $\eta=\varphi-a_0$ and $h=G^{\mathsf T}\mathsf gX'$. For the ratio we use the letters $T_Y:=\frac x2\langle\varphi,S_k(Y)\varphi\rangle-x\langle\varphi,S_k(Y)a_0\rangle$. Then $\sum_YT_Y=\frac x2\langle\eta,S\eta\rangle-\frac x2\langle a_0,Sa_0\rangle=-\frac x2\langle\xi,\mathsf E\xi\rangle+x\langle\xi,\mathsf Eh\rangle$, which is the exponent of \Cref{Y-C-P67}\,(b); the reason is \Cref{Y-C-P67}\,(a), by which $\eta\!\cdot\!S\eta=-(\xi-h)\!\cdot\!\mathsf E(\xi-h)$ and $a_0\!\cdot\!Sa_0=-h\!\cdot\!\mathsf Eh$.
\end{remark}

\section{The nonlinear step}\label{sec:nonlinear-step}

\subsection{The local prefactor}\label{Y-D-1}

The Gaussian integral in Prop.~\ref{Y-C-P67} has prefactor
$e^{-x\mathcal Q(-a_0)}$. We separate its exponent into the next-scale
energy and local terms,
\[
 x\mathcal Q(-a_0)=x'\mathrm{En}'(V')+\sum_\square T_\square.
\]
The identity holds for every real coarse field. In tame windows the
$T_\square$ are holomorphic and small enough for the analytic cap.
Near rough plaquettes the same terms enter the weights of the
corresponding clusters.

\begin{remark}[Notation from earlier sections; letters and tags with more than one sense] The letters $\mathsf G$, $\mathsf G_Z$ and $\mathsf G_{\mathsf n}=a^{-1}\mathrm{d}'^*\mathsf G_Z$ are the free kernels of \S\ref{Y-A-1}, introduced in \Cref{Y-A-GG}.

The letter $b$ is the number of \textup{(O6)} or a fine bond (item~C1); $e$ is the second member of a pair $(V',e)$ or the number $2.718\dots$, and $e_\mu$ is a unit vector; $q$ is either a matrix, namely the form $q^V_\ell$ of item~C4, or a fine cube (item~C1), and $q'$ is a coarse cube. In $M=L^{m'}$ of \textup{(O7)}, $m'$ is the exponent, whereas $m'_{y'}$ and the $m'$ in $\ell'^{{\rm c}V*}m'\ell'^{{\rm c}V}$ refer to the mask of item~E2. The tag \textup{(h)}, said of a field at a number $\hat\varepsilon$ or $\varepsilon'$, denotes the conditions \textup{(h1)} to \textup{(h3)} of \textup{(C3)} in item~D, and otherwise the monotonicity clause \textup{(h)} of item~H. In the names $K_{{\rm q},k}$, \textup{(G1$_k$)}, \textup{(F8$_k$)} and \textup{(F9$_k$)}, the subscript $k$ is a fixed label, not an index.

The sign $\partial V'(P)$ denotes the plaquette product $\partial V(p)$ of item~C.I~1.1 of \S\ref{Y-C-1}, evaluated at the coarse field. For $P=(y';\mu\nu)$ with $\mu<\nu$ and the bonds $c_1=(y',\mu)$, $c_2=(y'+Le_\mu,\nu)$, $c_3=(y'+Le_\nu,\mu)$, $c_4=(y',\nu)$, it is \[\partial V'(P)=V'(c_1)V'(c_2)V'(c_3)^{-1}V'(c_4)^{-1}.\] With no argument (Step~6), $\partial V'$ is the family of these products.

We write $V'(\pi)$ for the ordered product along a path $\pi$, and $V'^u(c)=u(c_-)V'(c)u(c_+)^{-1}$; also $\operatorname{ang}g=\arccos(\frac12\operatorname{tr}g)\in[0,\pi]$ (items~C.I~1.1, C.I~4 of \S\ref{Y-C-1}).

\end{remark}

\begin{theorem}[The prefactor at a level with rows]\label{Y-D-C5} The rows may be of any depth, and the empty rows are admitted (item~B and Remark~(4) of item~I). The group is $\operatorname{SU}(2)$, and $d=3$. The step is the drawn step of item~C.I~4 of \S\ref{Y-C-1}, set out there under ``The step, the chart''; it coincides with the step of item~1 of \S\ref{Y-G0-1}, which is given there under ``Fields'' and ``Step''.

For the table of the letters, with \textup{(O0)} to \textup{(O13)} and $L_*^{\rm su}$, we refer to item~3 of \S\ref{Y-G0-1}, and for \textup{(SD)} to item~1 of \S\ref{Y-G0-1}.

\begin{enumerate}

\item[\textup{A}] \emph{Letters.} The letters follow the order \textup{(O0)} to \textup{(O13)} of item~3 of \S\ref{Y-G0-1}.

\begin{enumerate}

\item[\textup{(O0)}] The pure numbers are $b_*$, $K_*$, $\lambda_*$, $L_A$, $\lambda_2^*$, $\theta_J$, as in clause \textup{(0)} of the order of the constants in \Cref{Y-A-Fr}. Further, $\mathsf g(t)=\sum_{n\in\mathbb Z^3}e^{-t|n|_\infty}$, $\mathsf g_2(t)=\sum_n(1+|n|_\infty)^2e^{-t|n|_\infty}$ and $\mathsf g_I(t)=3L^3\mathsf g(t)$.

\item[\textup{(O1)}] The number $d_0\in(0,\frac12)$ enters in two ways: through $S_{\varepsilon'}$, and through the number $\rho_0$ of the display \eqref{eq:row-mask-radius} of \S\ref{Y-A-1}.

\item[\textup{(O3)}] The display \eqref{eq:row-mask-radius} of \S\ref{Y-A-1} gives $\rho_0$; moreover $\mathsf N_{\mathsf B}=1+3K_*^2\mathsf g(\lambda_*)^2$.

\item[\textup{(O4)}] The number $L$ is odd, with $L\ge L_*^{\rm su}$ of item~3 of \S\ref{Y-G0-1}; of this the theorem uses only $L\ge3$ odd and $L\ge L_A$. We set $c_L=(L^2-1)/(6L^2)$. Further, $a\ge3/(c_LL)$, so that $aL\ge18$.

\item[\textup{(O5)}] The letter $\kappa\ge6$ enters the theorem only through the bound of clause \textup{(ca)}.

\item[\textup{(O6)}] These are the numbers of item~C.II~1.3 of \S\ref{Y-C-1}, evaluated at $K_*,\lambda_*$: $\gamma=(73L^4)^{-1}$, $b=16\mathsf N_{\mathsf B}+a/L$, $A^\flat=2e^2(16+a)+384e^8K_*^2\mathsf g_2(\lambda_*/2)^2$, $\theta_1=\min\{1,\lambda_*L/12\}$, $\theta^\flat=\min\{\theta_1,\gamma/(2A^\flat)\}$, $\theta_{\mathbb C}=\theta^\flat/6$ and $\vartheta_0^\flat=\theta^\flat/24$ ($\le10^{-4}$). In addition $\theta_D=\frac12\min\{\lambda_2^*,\theta_{\mathbb C},\theta_J\}$ ($\le1$).

The number $C_D\ge4$ is that of \textup{(O6)} in item~3 of \S\ref{Y-G0-1}. It is not smaller than the maximum by which \Cref{Y-A-P35} defines its own $C_D$, taken at this $\theta_D$; this number also occurs in clause \textup{(6)} of the order of the constants in \Cref{Y-A-Fr}. Moreover \textup{(Dec)}$(C,\theta_D)$ holds for every $C$ at least as large as that maximum. Each $K$ defined below grows with $C_D$.

Finally $c_R=64/\theta_D$, so that $c_R\theta_{\mathbb C}\ge128$.

\item[\textup{(O7)}] Here $M=L^{m'}$ with $m'\ge1$; it serves as the side of a cell.

\item[\textup{(O8)}] The letters are $\mathfrak p_0>2$ and $A_0>0$, and $p_0(x)=A_0(\log x)^{\mathfrak p_0}$, $\varepsilon=p_0x^{-1/2}$, $x'=x/L$; a prime on a letter means that it is taken at $x'$. So $x\varepsilon'^2=Lp_0'^2$.

\item[\textup{(O9)}] By \Cref{Y-C-CITh66}\,(k), $\mathsf n=60$.

\item[\textup{(O10)}] We take from item~C.II~4 (``Constants'') of \S\ref{Y-C-1} the numbers \[\mathsf K_A=1.24\cdot10^6\,\frac{A^\flat}{\gamma\theta_{\mathbb C}}\,\mathsf g_I(\theta_{\mathbb C}/2)\bigl[\frac{8\mathsf N_{\mathsf B}}\gamma \mathsf g_I(\theta_{\mathbb C})+1\bigr]\] and $\mathsf K_{\rm c}^0=(L+1)\mathsf K_A+6\cdot10^4LC_D/\theta_D$, $\mathsf a_{\rm c}=\frac{100}\gamma \mathsf N_{\mathsf B}\mathsf g_I(\theta_{\mathbb C})$.

We put $\mathsf K_{\rm c}:=\sqrt2\mathsf K_{\rm c}^0$. As $\mathsf K_{\rm c}^0\ge6\cdot10^4\cdot3\cdot4=7.2\cdot10^5>2.42$, it follows that $\mathsf K_{\rm c}\ge \mathsf K_{\rm c}^0+1$. The numbers of item~F come next in the order.

\item[\textup{(O12)}] A single condition is imposed on $E_a$, namely $E_a\ge32K_{{\rm q},k}$, and it enters only the last inequality of \textup{(d)}. Since $c_a(x)=E_ap_0^3x^{-1/2}$, we have $c_a(x')=E_ap_0'^2\varepsilon'$.

\item[\textup{(O13)}] The letter $x_\flat\ge e^2$ is the last one, and the floors of item~G hold for every $x\ge Lx_\flat$.

\end{enumerate}

\item[\textup{B}] \emph{Data.} These are the data \textup{(SD)} of item~1 of \S\ref{Y-G0-1}. We take $x\ge Lx_\flat$. The side $N$ is even, $M\mid N$, and \textup{(Tor)} holds: $N\ge2Ms(x)=16Mr(x)$.

Here $N$ is the coarse side of the present step: it is the $N'=NL^{K-k-1}$ of \S\ref{Y-A-1}, where the $N$ of that formula is the side of the last torus. The fine torus is $\mathbb Z^3/NL\mathbb Z^3$, and the coarse torus is $L\mathbb Z^3/NL\mathbb Z^3$.

The rows $\ell\in\mathcal R(b_*,K_*,\lambda_*,\frac19)$ live on the fine torus and have range $\le\varrho(x)$. The empty rows are admitted: $\ell=\emptyset$, with $\mathsf H_y=\{0\}$ for all $y$, belongs to $\mathcal R(1,0,\lambda_*,0)$, a subset of that set (see \Cref{Y-A-SetsOfRows}), and has range $0$. Neither a level nor a depth is among the data; see Remark~(4) of item~I.

In the sense of \Cref{Y-A-Rows}, rows attach a Euclidean space $\mathsf H_y$ to each fine site $y$, together with columns $\ell_y(p)\in \mathsf H_y$. Since $\ell\in\mathcal R(b_*,K_*,\lambda_*,\frac19)$, the second condition in the display of \Cref{Y-A-SetsOfRows}, taken with the distance $|y-p|:=|y-\pi(p)|_\infty$ of \Cref{Y-A-OwnUnits}, gives $|\ell_y(p)|\le K_*e^{-\lambda_*|y-\pi(p)|_\infty}$. The rows are such that the pair $(\ell,\ell')$ satisfies the four conclusions of \Cref{Y-A-Th36}. Of these data on the rows the theorem uses only $\ell\in\mathcal R(b_*,K_*,\lambda_*,\frac19)$, the range, and the bounds that \textup{(Rows)} of item~D asserts.

Finally, $(V',e)$ is any pair (item~E2 below).

\item[\textup{C}] \emph{First objects: lattice, norms, fields, form, radii, domains.}

\begin{enumerate}

\item[\textup{C1}] \emph{Lattice, blocks, distances, averages, lifts.} The objects come from Subsections~\ref{Y-A-SubLattices} and~\ref{Y-A-SubLifts}, including \Cref{Y-A-Facts} and Defs.~\ref{Y-A-Distances}, \ref{Y-A-Norm} and~\ref{Y-A-Lifts}, and from item~C.II~1.1 of \S\ref{Y-C-1}. Let $\mathsf h_L=(L-1)/2$. Fine bonds, plaquettes, cubes and sites are denoted $b,p,q,y$, and the coarse ones $c,P,q',y'$.

Every object $i$ has a base site $\pi(i)$, its least corner; $B(y')=y'+[-\mathsf h_L,\mathsf h_L]^3$; and $y(i)$ denotes the coarse site whose block contains $\pi(i)$. We write $y_P=\pi(P)$, $y_{q'}=\pi(q')$, $y(c)=c_-$. Distances are ${\rm dist}(i,j)=L^{-1}|y(i)-y(j)|_\infty$, measured in coarse units, and ``within $r$ of $o$'' means ${\rm dist}\le r$. A quantity depends on $V'$ within $r$ of $o$ when it is a function of $V'$ restricted to the coarse bonds whose two ends both lie within $r$ of $o$.

Concerning the combs, $I^\circ$ is the set of the fine bonds off the combs, at most $3L^3$ of which are based in a block; $\mathfrak L$ is the space of functions on $I^\circ$ with values in $\mathfrak g$ or $\mathfrak g_{\mathbb C}$; and $\Pi_{\mathfrak L}$ is the restriction to $I^\circ$. The maps $\mathrm{d}$ and $\mathrm{d}'$ are coboundaries. On cubes, $\delta=\mathrm{d}^*$.

For $c=(c_-,\mu)$ we put \[(Q_1A)(c)=L^{-3}\sum_{x\in B(c_-)}A([x,x+Le_\mu])\] and, for a function $\Theta$ on the fine cubes (Subsection~\ref{Y-A-SubLattices}), \[(Q_3\Theta)(y')=L^{-3}\sum_{x\in B(y')}\sum_{0\le i,j,k<L}\Theta(x+ie_1+je_2+ke_3).\]

The kernel $W=Q_2$ (Subsection~\ref{Y-A-SubLattices} and \Cref{Y-A-Facts}) satisfies $W\ge0$, $\sum_PW(P,p)=L^{-1}$ and $\sum_pW(P,p)=L^2$, and $W(P,p)=0$ unless $y(p)$ is a corner of $P$.

The edge plaquettes under $P=(y';\mu\nu)$ are the $p=(z;\mu\nu)$ with $z_\mu=y'_\mu+\mathsf h_L$, $z_\nu=y'_\nu+\mathsf h_L$ and $|z_\lambda-y'_\lambda|\le \mathsf h_L$; there are $L$ of them, all based in $B(y_P)$. The corner cube under $q'$ is the cube at $y_{q'}+\mathsf h_L(1,1,1)$. The lift $E$ is given by $(EX')_p=X'_P$ on the edge plaquettes under $P$ and by $0$ on all other plaquettes, and the lift $E_3$ by $(E_3\Theta)_q=\Theta_{q'}$ on the corner cube under $q'$ and by $0$ on all other cubes. Hence $(E^*F)_P$ is the sum of $F_p$ over the $L$ edge plaquettes under $P$, and $(E_3^*\Xi)_{q'}$ is the value of $\Xi$ at the corner cube under $q'$.

\item[\textup{C2}] \emph{Norms.} We follow item~C.II~1.1 of \S\ref{Y-C-1} and \Cref{Y-A-Norm}. Here $\mathfrak g$ is the space of hermitian traceless $2\times2$ matrices, and $\mathfrak g_{\mathbb C}$ is its complexification. We use the bilinear form $Y\cdot Y'=\frac12\operatorname{tr}YY'$, the pairing $\langle F,F'\rangle=\sum_iF_i\cdot F'_i$ and the norm $\|Y\|^2=\frac12\operatorname{tr}Y^*Y$, and $|\cdot|$ is the operator norm of a $2\times2$ matrix. Then $\|Y\|\le|Y|\le\sqrt2\|Y\|$, where the first two agree on $\mathfrak g$, and $|Y\cdot Y'|\le\|Y\|\,\|Y'\|$.

For $g\in\operatorname{SL}(2,\mathbb C)$ one has $|g^{-1}|=|g|$ and $\|\operatorname{Ad}_g^{\pm1}\|\le|g|^2$, and $\operatorname{Ad}_g^{-1}$ is the transpose of $\operatorname{Ad}_g$ with respect to $Y\cdot Y'$. By $\|A\|$ we mean the operator norm on $\ell^2(I^\circ;\mathfrak g_{\mathbb C})$, and $\|K\|_\infty=\sup_i\sum_j\|K(i,j)\|$. The norm $\|K\|_\theta$ is the larger of $\sup_i\sum_j|K(i,j)|e^{\theta\,{\rm dist}(i,j)}$ and $\sup_j\sum_i|K(i,j)|e^{\theta\,{\rm dist}(i,j)}$; for rows, the entry at $(y,p)$ is $|\ell_y(p)|$. Then $\|K\|_\infty\le\|K\|_0\le\|K\|_\theta$, and $\sup_i\sum_{{\rm dist}>r}|K(i,j)|$, as well as the corresponding sums over columns, is $\le e^{-\theta r}\|K\|_\theta$.

\item[\textup{C3}] \emph{Fields, transports, dressed kernels.} We follow items~C.II~1.1, C.II~1.2 of \S\ref{Y-C-1}. We let $\mathcal V'$ be the set of $\operatorname{SU}(2)$-valued fields on the coarse bonds, with $V'(\bar c)=V'(c)^{-1}$, and $\mathfrak V$ the set of $\operatorname{SL}(2,\mathbb C)$ fields with $|V'(c)|\le1+\vartheta_0^\flat$ at every bond; $\mathfrak V^\circ$ is the set of those with $|V'(c)|<1+\vartheta_0^\flat$ at every bond. We extend by $1$ every field that is given on a set of bonds only.

The coordinate path $\Gamma(y',y'')$ runs first in the direction $e_1$, then in $e_2$, then in $e_3$, each time the short way. It has no more than $3\,{\rm dist}(y',y'')$ bonds, and each of its sites lies within ${\rm dist}(y',y'')$ both of $y'$ and of $y''$, since every coordinate of any site of the path lies on the short arc joining the matching coordinates of $y'$ and $y''$. We set $T(y',y'')=\operatorname{Ad}_{V'(\Gamma(y',y''))}$.

For a free kernel $K$ we set $K^V(i,j)=K(i,j)T(y(i),y(j))$; then $(K^V)^*(j,i)=K(i,j)T(y(i),y(j))^{-1}$, and $(K^V)^*$ is the transpose of $K^V$ for $\langle,\rangle$. So $(\ell^VF)_y=\sum_p\ell_y(p)\otimes T(y(y),y(p))F_p$ and $(\mathrm{d}'^VX')_{q'}=\sum_P\mathrm{d}'(q',P)T(y_{q'},y_P)X'_P$. The lifts $E$ and $E_3$ join objects $i$, $j$ with $y(i)=y(j)$ and carry no transport.

\item[\textup{C4}] \emph{The form and the kernel $\mathsf g^V$.} We follow item~C.II~1.2 of \S\ref{Y-C-1}. With $\mathsf B^V=1+\ell^{V*}\ell^V$ we set \[q=q^V_\ell=\Pi_{\mathfrak L}[\mathrm{d}^{V*}\mathsf B^V\mathrm{d}^V+aQ_1^{V*}Q_1^V]\Pi_{\mathfrak L}\] and $\mathsf g^V=\Pi_{\mathfrak L} \mathrm{d}^{V*}\mathsf B^VE$. Here $\ell_y(p)=0$ unless ${\rm dist}(y,p)\le\varrho_L$. The form satisfies $q(b,b')=0$ beyond the distance $2\varrho_L+2$, and $q(b,\cdot)$ depends on $V'$ within $2\varrho_L+2$ of $y(b)$.

The reason is the following. By item~C3, a coordinate path keeps each of its sites, as seen from either end, within the distance between its two ends. Starting from $y(b)$ and following a step $(y,p,p')$ from $b$ to $b'$, we find $y(p)$ within $1$, the path $\Gamma(y(y),y(p))$ within $\varrho_L+1$, the path $\Gamma(y(y),y(p'))$ within $2\varrho_L+1$, and $y(b')$ within $2\varrho_L+2$. Item~C.II~1.2 of \S\ref{Y-C-1} makes the same estimate.

Moreover $\mathsf g^V(b,P)=0$ unless $y_P$ lies within $2\varrho_L+1$ of $y(b)$, and $\mathsf g^V(b,P)=\sum_{\mathsf t}g_{\mathsf t}\operatorname{Ad}_{V'(\pi_{\mathsf t})}$, the sum running over the $\mathsf g$-steps $\mathsf t$ of item~C.II~1.2 of \S\ref{Y-C-1}, where $\pi_{\mathsf t}$ is a single coarse path from $y(b)$ to $y_P$ with at most $6\varrho_L+8$ bonds, all of whose sites lie within $2\varrho_L+1$ of $y(b)$.

\item[\textup{C5}] \emph{Functions of $x$, radii, cells, the window.} At $x$ we set $R=R(x)=\varrho(x)=\lceil c_R\log x\rceil$, $r(x)=\lceil20c_R\log x/M\rceil$, $s(x)=8r(x)$, $r'=r(x')$, $\ell_w(x')=6Mr'+12$, $\mathsf b(x')=2-1/\log x'$ and $\vartheta(x')=\vartheta_0^\flat(1-1/\log x')$. In coarse units the radii are $\varrho_L=\lceil R/L\rceil$, $r_w=\lfloor R/2\rfloor$, $R_2=2R+2$, $R_3=2R+\varrho_L+5$ and $R_T:=4R+2\varrho_L+6$. In item~E8 of \S\ref{Y-F-1}, $R_5$ stands for $4R+4\varrho_L+8$ and $R_6$ for $4R+6\varrho_L+11$; see \textup{(h)}.

The cells $\square$ are the cubes of the fixed grid that contain $M^3$ coarse sites each. Every object $i$, fine or coarse, is assigned to the cell containing $y(i)$. This assignment defines the relations $P\in\square$, $q'\in\square$ and $b\in\hat\square$. The distance from one cell to another is their sup distance counted in cells; $Y^{+j}$ is the set of the cells within distance $j$ of $Y$, and $\square^\approx=\square^{+r'}$. The window $W_\square$ consists of the coarse sites within $R_T$ of some site of $\square$; it is a box of side $M+2R_T$.

\item[\textup{C6}] \emph{Domains of complex fields; tame fields.} Let $A$ be a set of coarse sites; by the bonds, paths and plaquettes of $A$ we mean those whose sites all lie in $A$. Let a field with values in $\operatorname{SL}(2,\mathbb C)$ be given on the bonds of $A$. It has \textup{(u)} on $A$ if the following hold.

\begin{enumerate}

\item[\textup{(u1)}] $|V'(c)|<1+\vartheta(x')$.

\item[\textup{(u2)}] $|V'(\pi)|<\mathsf b(x')$ for every path in $A$ with at most $\ell_w(x')$ bonds.

\item[\textup{(u3)}] $|\partial V'(P)-1|<2\varepsilon'$.

\end{enumerate}

We let $\mathcal U_{\{\square\}}(x')$ be the set of fields on the bonds of $\square^\approx$ with \textup{(u)} on $\square^\approx$. This set is the domain of clause \textup{(ca)} of the class $C^{\rm su}_\flat$ at $x'$, with $X=\{\square\}$: in item~1 of \S\ref{Y-G0-1}, under ``Domains'' and ``Members'', these are $\mathcal U_X(x)$ and \textup{(ca$_A$)}, taken at $x'$ with $X=\{\square\}$ and $A=1$. In the same way $U^T_\square(x')$ consists of the fields on the bonds of $W_\square$ that have \textup{(u)} on $W_\square$.

For a field with \textup{(u3)} we put $X'_P:=-i\log \partial V'(P)$, the logarithm being its power series, and then $|X'_P|\le2.06\varepsilon'$ by item~C.II~1.4 of \S\ref{Y-C-1}. When a field is given on the bonds of $A$ only and is extended by $1$, this definition is used at the plaquettes of $A$, while $X'_P:=0$ at every plaquette which is not a plaquette of $A$; this is the rule of item~C.II~1.4 of \S\ref{Y-C-1}. For the extended field, $\partial V'(P)$ at such a plaquette is a product of no more than three bond matrices, and there the power series may fail to converge. By this rule $X'_P$ has a value at every coarse plaquette of the torus, and so $X'$, $\mathsf g^VX'$, $a_0=C\mathsf g^VX'$, $v_*$ and $\Xi_*$ are defined at each of their indices.

The rule $X'_P:=0$ at a plaquette which is not a plaquette of $A$ enters neither a value nor a bound of $T_\square$, $P^T_\square$, \textup{(f)} or \textup{(f$'$)} (see Step~2 of the proof for the first two, and Steps~4 and~8 with \textup{(0.1)} for the last two).

We call a field of $\mathcal V'$ tame on $A$ if $\phi'_P\le\varepsilon'$ holds at every plaquette of $A$.

\end{enumerate}

\item[\textup{D}] \emph{Inputs.} The following inputs are statements quoted from earlier sections, and none of them is a hypothesis of the theorem.

\begin{enumerate}

\item[\textup{(C2)}] \emph{The law.} We use \Cref{Y-C-CITh66}\,(a), (b), (c), (f), (i), (k), and only through $C=C^V=G^V(G^V)^{\mathsf T}$, in the following form.

\begin{enumerate}

\item[\textup{(C2.0)}] The matrix $C$ is defined at all fields of $\mathfrak V$.

\item[\textup{(C2.1)}] By \Cref{Y-C-CITh66}\,(b), $C(b,b')=0$ unless ${\rm dist}(b,b')\le2r_w$ ($\le R$), and $C(b,\cdot)$ depends on $V'$ within $3r_w+\varrho_L+1$ ($\le R_2-1$) of $y(b)$.

\item[\textup{(C2.2)}] By \Cref{Y-C-CITh66}\,(c), $C[V'^u]=\operatorname{Ad}_uC[V']\operatorname{Ad}_u^{-1}$ for $\operatorname{SU}(2)$-valued $u$, where $(\operatorname{Ad}_u\eta)_b=\operatorname{Ad}_{u(y(b))}\eta_b$.

\item[\textup{(C2.3)}] By \Cref{Y-C-CITh66}\,(i), (a), each entry of $C$ is continuous on $\mathfrak V$, holomorphic on $\mathfrak V^\circ$ and real at $\operatorname{SU}(2)$-valued fields.

\item[\textup{(C2.4)}] By \Cref{Y-C-CITh66}\,(f), (k), under \textup{(F$_{\rm law}$)} and $c_R\theta_{\mathbb C}\ge128$, the matrix $C$ is invertible on all of $\mathfrak V$, and $S:=C^{-1}-q$ has $\|S\|\le x^{-60}$.

\item[\textup{(C2.5)}] From \textup{(C2.4)} and \textup{(P1)}, \Cref{Y-C-L66}\,(a) gives $\|C\|\le4/\gamma$ and $\|C-q^{-1}\|\le8\gamma^{-2}x^{-60}$ on $\mathfrak V$.

\end{enumerate}

\item[\textup{(C3)}] \emph{The centre.} We use \Cref{Y-C-Th68}\,(a), (b), (c) together with \Cref{Y-C-L71}\,(a), (b), each applied with the law $C$ of \textup{(C2)} in the role of the law in its hypotheses. For the matrices we follow item~C.II~2.2 of \S\ref{Y-C-1} (``The matrices''): $q^1$, $\mathsf g^1$, $C^1$ are $q$, $\mathsf g^V$, $C$ at $V'=1$, $H^{(x)}=E-\mathrm{d}C^1\mathsf g^1$, $\mathsf N^{(x)}=Q_1C^1\mathsf g^1$, and $\mathcal Q_i$ is the set of those coarse cubes all of whose faces lie within $R_2$ of $i$. We set \[H^\natural(p,P)=H^{(x)}(p,P)-\sum_{q'\in \mathcal Q_p}\mathsf G(p,q')\mathrm{d}'(q',P)\] and \[\mathsf N^\natural(c,P)=\mathsf N^{(x)}(c,P)-\sum_{q'\in \mathcal Q_c}\mathsf G_{\mathsf n}(c,q')\mathrm{d}'(q',P).\]

The numbers are $\tau_1=8\gamma^{-2}(1+b)x^{-60}$, $\tau_{\rm c}=165(L+1)\mathsf N_{\mathsf B}\,3L^3(2R+3)^3\tau_1$, \[\tau_2=24L^3(2R+1)^3\gamma^{-2}x^{-60}+(2/\gamma)\mathsf g_I(\theta_{\mathbb C}/2)x^{-64}\] and $\mathsf{ti}_{\rm c}=(16+4L)\mathsf N_{\mathsf B}\tau_2+7e(L+1)C_Dx^{-128}$.

Throughout \textup{(C3)}, either $i=p$ and $o=y(p)$, or $i=c$ and $o=c_-$, and $\mathsf B(i)$ denotes the box of radius $R_3$ about $o$. A field on $\mathsf B(i)$ is said to have \textup{(h)} at $\hat\varepsilon$ if the following hold.

\begin{enumerate}

\item[\textup{(h1)}] $|V'(c)|\le1+\vartheta_0^\flat$.

\item[\textup{(h2)}] $|V'(\pi)|\le2$ for every path in the box with at most $\ell_w(x')$ bonds.

\item[\textup{(h3)}] $|\partial V'(P)-1|\le2\hat\varepsilon$.

\end{enumerate}

We set \[\Lambda_p=(EX'-\mathrm{d}^VC\mathsf g^VX')_p-\sum_PH^\natural(p,P)T(y(p),y_P)X'_P\] and \[\Lambda_c=(Q_1^VC\mathsf g^VX')_c-\sum_P\mathsf N^\natural(c,P)T(c_-,y_P)X'_P.\]

Under \textup{(P0)} and \textup{(P1)}, for every $i$, every $0<\hat\varepsilon\le\frac1{40}$ and every field on $\mathsf B(i)$ with \textup{(h)} at $\hat\varepsilon$, the following hold.

\begin{enumerate}

\item[\textup{(C3.1)}] The quantities $\Lambda_i$ and $a_{0,b}=(C\mathsf g^VX')_b$, the latter at the bonds with ${\rm dist}(b,o)\le1$, depend on $V'$ only through the restriction of $V'$ to the bonds of $\mathsf B(i)$, and on $X'$ only at the plaquettes of $\mathsf B(i)$; \Cref{Y-C-Th68}\,(a) contains both statements.

By the part of \Cref{Y-C-Th68}\,(a) that does not depend on \textup{(h2)}, \textup{(h3)} or \textup{(Fr2)}, this clause holds with no bound on the size of the field: if a field of $\mathfrak V$ is given on a bigger set and $X'$ is defined on all plaquettes of that set (item~C6), then $\Lambda_i$ and $a_{0,b}$ keep the values that they have when computed on $\mathsf B(i)$.

\item[\textup{(C3.2)}] $\|\Lambda_i\|\le \mathsf K_{\rm c}^0\hat\varepsilon^2+\tau_{\rm c}\hat\varepsilon$; so $\|\Lambda_i\|\le \mathsf K_{\rm c}\hat\varepsilon^2$ when $\tau_{\rm c}\le\hat\varepsilon\le\frac1{40}$.

\item[\textup{(C3.3)}] $\|a_{0,b}\|\le \mathsf a_{\rm c}\hat\varepsilon$ at those bonds.

\item[\textup{(C3.4)}] $H^\natural(p,P)=0=\mathsf N^\natural(c,P)$ for ${\rm dist}>R_2$.

\item[\textup{(C3.5)}] $\|H^\natural-H^\sharp\|_\infty+\|\mathsf N^\natural-\mathsf N^\sharp\|_\infty\le\mathsf{ti}_{\rm c}$.

\end{enumerate}

\item[\textup{(Rows)}] \emph{Decay of the free kernels.} For the pair $(\ell,\ell')$, \textup{(Dec)}$(C_D,\theta_D)$ of \Cref{Y-A-Dec} holds: $\|\cdot\|_{\theta_D}\le C_D$ for $\ell$, $\mathsf X$, $H^\sharp$, $\mathsf B'$, $\ell'$, as well as for $\mathsf T$, $\mathsf N^\sharp$ in clause \textup{(f$'$)} of item~H, and for $\mathsf G$, $\mathsf G_Z$ inside \textup{(C3)}. None of $\ell$, $\mathsf X$, $\mathsf B'$ belongs to the six kernels of \textup{(Fr2)} in item~C.II~2.2 of \S\ref{Y-C-1}, but every kernel named here is one of the sixteen kernels of \textup{(Dec)}.

This comes from \Cref{Y-A-P35}, for every $\ell\in\mathcal R(b_*,K_*,\lambda_*,\frac19)$, $L\ge L_A$ odd and $aL\ge3/c_L$, taken with $\theta_{\mathbb C}$ for its $\theta_q$; the proof of \Cref{Y-A-Fr} obtains it in Step~H13 with the help of Step~H16\,(i), (ii). Thus \textup{(Rows)} holds for every $\ell$ of that set; it is not a hypothesis on $\ell$, and it adds no condition to item~B.

\item[\textup{(Tools)}] \Cref{Y-C-CIIL64}\,(d) gives $\sup_b\sum_{\mathsf t\text{ from }b}|g_{\mathsf t}|\le4\mathsf N_{\mathsf B}$, and \Cref{Y-C-CIIL65}\,(a) gives $\|q^{-1}(b,b')\|\le(2/\gamma)e^{-\theta_{\mathbb C}{\rm dist}(b,b')}$ on $\mathfrak V$. We invoke \Cref{Y-C-L67}\,(b), (c) and Bianchi's identity, which the proof of \Cref{Y-C-L67}\,(d) establishes as an identity of invertible matrices, together with its six quotients. We use \Cref{Y-A-L1} as well.

Finally, \textup{(E)} of item~C.I~1.1 of \S\ref{Y-C-1}, where it is proved, states that \[|e^{iY_1}\cdots e^{iY_n}-1-i\sum Y_k|\le e^S-1-S\le0.72S^2\] whenever the matrices $Y_k$ satisfy $S=\sum|Y_k|\le1$.

\end{enumerate}

\item[\textup{E}] \emph{Second objects.}

\begin{enumerate}

\item[\textup{E1}] \emph{The free kernels and their cuts.} These kernels are taken from \S\ref{Y-A-1}, namely from \Cref{Y-A-Pairs}, from \Cref{Y-A-P4} and its proof, from \Cref{Y-A-P5} and from \Cref{Y-A-XT}; each of them is real and depends on no field. We set $\mathsf B=1+\ell^{\mathsf T}\ell$, and, given $\Xi$ on the fine cubes and $Z$ on the coarse plaquettes, $Y=W^*Z+\delta\Xi$.

Let $\Xi_Z\in\ker Q_3$ be the minimiser of $\langle Y,\mathsf B^{-1}Y\rangle$; then $\mathsf T_\sharp Z=W^*Z+\delta\Xi_Z$ and $R_\sharp=\mathsf T_\sharp^*\mathsf B^{-1}\mathsf T_\sharp+a^{-1}\mathrm{d}'\mathrm{d}'^*>0$. We put $\mathsf B'=LR_\sharp^{-1}$, $H^\sharp=L^{-1}\mathsf B^{-1}\mathsf T_\sharp \mathsf B'$ and $\mathsf N^\sharp=(aL)^{-1}\mathrm{d}'^*\mathsf B'$ (\Cref{Y-A-P4}\,(a)), as well as $\mathsf Xv=\Xi_{v/L}$ and $\mathsf Tv=\mathsf B^{-1}\mathsf T_\sharp v/L$. The rows $\ell'$ are those of \Cref{Y-A-P5}, with $\mathsf B'=1+\ell'^{\mathsf T}\ell'$. These definitions give the identities

\begin{enumerate}

\item[\textup{(I1)}] $\mathsf B\mathsf T=L^{-1}W^*+\delta\mathsf X$;

\item[\textup{(I2)}] $H^\sharp=\mathsf T\mathsf B'$, so $\mathsf BH^\sharp=(L^{-1}W^*+\delta\mathsf X)\mathsf B'$.

\end{enumerate}

The cuts are $\ell'^{\rm c}=\ell'\,1\{{\rm dist}\le\varrho(x')\}$, the cuts $\mathsf X^<$, $\mathsf T^<$ at $R(x')$ and the cuts $H^{\sharp<}$, $\mathsf N^{\sharp<}$ at $R_2$; we also set $K^>=K-K^<$ and $\mathsf B'^{\rm c}=1+\ell'^{{\rm c}\mathsf T}\ell'^{\rm c}$.

\item[\textup{E2}] \emph{Pairs, clipped curvature, mask, energy.} The objects are those of items~E2, E7 of \S\ref{Y-F-1} and of item~1 of \S\ref{Y-G0-1}, under ``Pair'' and ``Energy''. Here $\mathfrak E=\prod_PS^2$, a product of unit spheres of $\mathfrak g$.

We set $\phi'_P=\operatorname{ang}\partial V'(P)$ and \[X'_P(V',e)=\mathsf L(\partial V'(P))+\pi e_P1\{\partial V'(P)=-1\},\] where $\mathsf L$ is the logarithm with $|\mathsf L|<\pi$ and $\mathsf L(-1)=0$. So $|X'_P|=\phi'_P$, and $X'$ is a Borel function of $(V',e)$; moreover $X'_P(V'^u,\operatorname{Ad}_ue)=\operatorname{Ad}_{u(y_P)}X'_P(V',e)$, where $(\operatorname{Ad}_ue)_P=\operatorname{Ad}_{u(y_P)}e_P$, and $e_P$ plays a role only where $\phi'_P=\pi$.

Let $D'=\{P:\phi'_P>\varepsilon'\}$, $\zeta'_P=X'_P\min\{1,\varepsilon'/\phi'_P\}$, and $S_{\varepsilon'}(t)=t^2/2$ ($t\le\varepsilon'$), $d_0\varepsilon't$ ($t>\varepsilon'$). Here and below a clipping multiplier $\min\{1,\alpha/t\}$, with $\alpha>0$ and $t\ge0$, has value $1$ at $t=0$. The mask is $m'_{y'}=1$ when no $P$ of $D'$ has ${\rm dist}(y',P)\le\rho_0$ and no $P$ with $\phi'_P>1$ has ${\rm dist}(y',P)\le\varrho(x')+1$, and $0$ in every other case. The energy \[\mathrm{En}'(V')=\sum_PS_{\varepsilon'}(\phi'_P)+\frac12\sum_{y'}m'_{y'}|(\ell'^{{\rm c}V}\zeta')_{y'}|^2\] does not depend on $e$.

\item[\textup{E3}] \emph{The centre.} For a field of $\mathfrak V$ and an arbitrary $\mathfrak g_{\mathbb C}$-valued $X'$ we set $a_0=C\mathsf g^VX'$, $S=C^{-1}-q$ and $u^T:=1-qC$ ($=SC$). Because $C^{\mathsf T}=C$ (\Cref{Y-C-CITh66}\,(b)), $u^T$ equals the transpose of the $\mathsf u$ that item~C.II~2.1 of \S\ref{Y-C-1} defines. Further, $\bar F=EX'-\mathrm{d}^Va_0$ and $\bar n=Q_1^Va_0$, and \[\mathsf m_C:=\frac12\langle EX',\mathsf B^VEX'\rangle-\frac12\langle \mathsf g^VX',C\mathsf g^VX'\rangle\] is the $\mathsf m_C$ of \Cref{Y-C-P67}\,(a).

\item[\textup{E4}] \emph{The multipliers.}

We set $v_*=(1+\ell'^{{\rm c}V*}m'\ell'^{{\rm c}V})\zeta'$, $Z_*=v_*/L$, $\Xi_*=(\mathsf X^<)^Vv_*$, $Y_*=W^{V*}Z_*+\mathrm{d}^{V*}\Xi_*$ and $\Delta_Y=\mathsf B^V\bar F-Y_*$. For clauses \textup{(f$'$)} and \textup{(g)} of item~H we set $\nu^{\rm ref}=(\mathsf T^<)^Vv_*$, $\sigma_*=\ell^V\nu^{\rm ref}$, $\nu_*=Y_*-\ell^{V*}\sigma_*$, $\omega_*=(aL)^{-1}\mathrm{d}'^{V*}v_*$, $\Delta^{\rm ref}=\bar F-\nu^{\rm ref}$, $\Delta_\sigma=\ell^V\bar F-\sigma_*$, $\Delta_F=\bar F-\nu_*$ and $\Delta_n=\bar n-\omega_*$.

At a field with \textup{(u3)}, the tame formulas are $X'_P=-i\log \partial V'(P)$ at every plaquette of the set on which the field is given, with $X'_P:=0$, by convention, at all other plaquettes (item~C6), and $v_*=(1+\ell'^{{\rm c}V*}\ell'^{{\rm c}V})X'$; every other object is computed from these by the same formulas.

\item[\textup{E5}] \emph{The local term $T_\square$ and its holomorphic form.} With $\mathrm{En}'_P=S_{\varepsilon'}(\phi'_P)+\frac12\langle\zeta'_P,\allowbreak (\ell'^{{\rm c}V*}m'\ell'^{{\rm c}V}\zeta')_P\rangle$ we set \[\FitDisplay{\dimexpr\linewidth-1pt\relax}{T_\square:=\sum_{P\in\square}\bigl[\frac{x'}2\langle X'_P,v_{*P}\rangle-x'\mathrm{En}'_P+\frac x2\langle X'_P,(E^*\Delta_Y)_P\rangle\bigr]+\frac x2\sum_{q'\in\square}\langle(\mathrm{d}'^VX')_{q'},(E_3^*\Xi_*)_{q'}\rangle-\frac x2\sum_{b\in\hat\square}\langle a_{0,b},(u^T\mathsf g^VX')_b\rangle.}\]

For $V'\in U^T_\square(x')$, extended by $1$, we define $P^T_\square(V')$ as the sum of the three sums \[\frac x2\sum_{P\in\square}\langle X'_P,(E^*\Delta_Y)_P\rangle+\frac x2\sum_{q'\in\square}\langle(\mathrm{d}'^VX')_{q'},(E_3^*\Xi_*)_{q'}\rangle-\frac x2\sum_{b\in\hat\square}\langle a_{0,b},(u^T\mathsf g^VX')_b\rangle,\] in which all objects are given by the tame formulas of item~E4. For $V'\in \mathcal U_{\{\square\}}(x')$ we set $P(\{\square\},V'):=P^T_\square(V'|W_\square)$, and $P(Y,\cdot):=0$ for $\#Y\ge2$.

\item[\textup{E6}] \emph{Marked cells, class partitions, the weights $w_j$.} We let $\mathrm{mk}'=\mathrm{mk}'(V')$ be the set of the cells that hold a plaquette of $D'$.

By a class partition we mean a partition $\{\mathrm{mk}'_j\}$ of $\mathrm{mk}'$ into sets whose pairwise distance is $>2r'$ cells. We call classes of the chain the classes into which $\mathrm{mk}'$ falls under the relation generated by ``distance $\le2s_0+4r_1+2$'', with $s_0=2r'$, $r_l=\lceil(\varrho_L+2)/M\rceil$ and $r_1=\lceil R_2/M\rceil+3r_l+5$; since their pairwise distance is $>4r'$, they form a class partition.

We set $w_j:=\sum_{\square\in \mathrm{mk}_j'^{+r'}}T_\square$ and $1_{\square^\approx}(V')=1\{\phi'_P\le\varepsilon'\text{ for all }P\in\square^\approx\}$.

In \textup{(e)} and in Step~7, for a field $V'$ of $\mathcal V'$, $P(\{\square\},V')$ stands for $P(\{\square\},V'|\square^\approx)$, where $V'|\square^\approx$ is the restriction to the bonds of $\square^\approx$. By convention, the product $1_{\square^\approx}(V')P(\{\square\},V')$ has the value $0$ wherever $1_{\square^\approx}(V')=0$, whether or not $P(\{\square\},V')$ is defined there. At such fields the restriction $V'|\square^\approx$ may lie outside $\mathcal U_{\{\square\}}(x')$, and then $P$ may be undefined.

Where $1_{\square^\approx}(V')=1$, we have $V'|\square^\approx\in \mathcal U_{\{\square\}}(x')$ by \textup{(0.6)}.

\end{enumerate}

\item[\textup{F}] \emph{Constants.} In the order of the letters, these numbers come after $M$ and before $c_h$, $E_a$. We set \[K_Y:=4\mathsf N_{\mathsf B}\mathsf K_{\rm c}+10^6\theta_D^{-2}(1+C_D^2)(C_D+eL+6eC_D)+1\] and \[K_{{\rm q},k}:=M^3[3.09L^2K_Y+520LC_D(1+4C_D^2)+1].\]

For clause \textup{(f$'$)} of item~H we set $K_{\rm ref}:=\mathsf K_{\rm c}+10^6\theta_D^{-2}C_D(1+C_D^2)+1$, $K_\sigma:=4C_DK_{\rm ref}$, $K_F:=K_Y+16C_D^2K_{\rm ref}$ and $K_n:=\mathsf K_{\rm c}+10^6\theta_D^{-2}(1+C_D^2)+1$. These numbers depend on $(L,a,K_*,\lambda_*,C_D,\theta_D)$, and $K_{{\rm q},k}$ on $M$ as well. None of them depends on $x$, $N$, $c_h$, $E_a$, a member of a class, or on more of $\ell$ than its starred constants and its range.

\item[\textup{G}] \emph{Floors.} Each floor is a condition on $x$ and on the letters \textup{(O0)} to \textup{(O10)} alone; every one of them holds for all large $x$, and all of them enter the choice of $x_\flat$.

\begin{enumerate}

\item[\textup{(F$_{\rm law}$)}] $432L^3b^{5/2}\gamma^{-3/2}\mathsf g(\theta_{\mathbb C})(R+3)^3\le x^4$.

\item[\textup{(P0)}] $x\ge L^2$; $R\ge6$.

The \textup{(P0)} of item~C.II~4 of \S\ref{Y-C-1} is this condition with \textup{(Tor)} as a third clause. Since \textup{(Tor)} depends on $N$, it is a datum here (item~B) and not a floor. So the conditions \textup{(P0)}, \textup{(P1)} as \textup{(C3)} assumes them, together with item~B, make up the whole hypothesis \textup{(P0)}, \textup{(P1)} of \Cref{Y-C-Th68} and of \Cref{Y-C-L71}.

\item[\textup{(P1)}] $(4/\gamma)x^{-60}\le1$; $165(L+1)\mathsf N_{\mathsf B}\,3L^3(2R+3)^3\,8\gamma^{-2}(1+b)\le x$.

\item[\textup{(G1$_k$)}] $Mr'\ge R_T+1$; $\rho_0\le R(x')$.

This floor is a consequence of \textup{(G1$^\flat$)} of item~G of \S\ref{Y-F-1}, which states $Mr'\ge5R+4\varrho_L+12$ and $\rho_0\le R(x')$ and, besides, $R\ge\max\{M,6\}$.

\item[\textup{(F8$_k$)}] $\varepsilon'\le\frac1{40}$; $35100R_T^2\varepsilon'\le1$; $p_0'\ge1$.

\item[\textup{(F9$_k$)}] $33\mathsf N_{\mathsf B}(1+C_D)\mathsf{ti}_{\rm c}+100\mathsf N_{\mathsf B}C_D^4x'^{-64}\le\varepsilon'$; $198L^3\mathsf a_{\rm c}\gamma^{-1}\mathsf N_{\mathsf B}(3L^3)^{1/2}(2R+4\varrho_L+5)^{3/2}\le x^{59}$.

\end{enumerate}

\item[\textup{H}] \emph{Conclusions.} In this setting the following hold.

\begin{enumerate}

\item[\textup{(a)}] At every $(V',e)$, \[x\,\mathsf m_C-\frac x2\langle a_0,Sa_0\rangle=\frac x2\bigl[\langle\bar F,\mathsf B^V\bar F\rangle+a|\bar n|^2\bigr]=x'\mathrm{En}'(V')+\sum_\square T_\square.\]

\item[\textup{(b)}] The term $T_\square$ is real and jointly Borel in $(V',e)$, and $T_\square(V'^u,\operatorname{Ad}_ue)=T_\square(V',e)$ for every $\operatorname{SU}(2)$-valued $u$. It depends on $V'$ only through the restriction of $V'$ to the bonds of $W_\square$, hence within $R_T\le Mr'-1$ of $\square$, and on $e_P$ only at those plaquettes $P$ of $W_\square$ where $\phi'_P=\pi$.

\item[\textup{(c)}] Let $V'$ be tame on $W_\square$; this is the case in particular when $\phi'_P\le\varepsilon'$ at each $P$ within $R_T+1$ of $\square$. Then $T_\square(V',e)=P^T_\square(V'|W_\square)$ and $|T_\square(V',e)|\le K_{{\rm q},k}p_0'^2\varepsilon'$.

\item[\textup{(d)}] $P^T_\square$ is holomorphic on $U^T_\square(x')$, with $\sup|P^T_\square|\le K_{{\rm q},k}p_0'^2\varepsilon'$.

It follows that $P(\{\square\},\cdot)$ is holomorphic on $\mathcal U_{\{\square\}}(x')$, is invariant under $\operatorname{SU}(2)$-valued gauge transformations, is real at $\operatorname{SU}(2)$-valued fields, vanishes at the pure gauges, and has $\sup|P(\{\square\},\cdot)|\le K_{{\rm q},k}p_0'^2\varepsilon'\le\frac1{32}c_a(x')$.

Therefore the family $P$ satisfies clause \textup{(ca)} at $x'$ with $\frac1{32}$ of the cap.

\item[\textup{(e)}] At every pair $(V',e)$ and for every class partition $\{\mathrm{mk}'_j\}$, so in particular for the classes of the chain, \[x\,\mathsf m_C-\frac x2\langle a_0,Sa_0\rangle=x'\mathrm{En}'(V')+\sum_jw_j+\sum_\square1_{\square^\approx}(V')\,P(\{\square\},V').\] Moreover, each $w_j$ depends on $V'$ within $2r'\le s(x')$ cells of $\mathrm{mk}'_j$.

\item[\textup{(f)}] Let $p$ be a fine plaquette, let $o=y(p)$, and let a field with values in $\operatorname{SL}(2,\mathbb C)$ satisfy \textup{(u)} on the box of radius $R_T$ about $o$; the field is extended by $\mathbf 1$, and the objects are defined by the tame formulas. Then $\|\Delta_{Y,p}\|\le K_Y\varepsilon'^2$.

\item[\textup{(f$'$)}] Let $i$ be a fine plaquette $p$, a fine site $y$ or a coarse bond $c$, put $o=y(i)$, and let the field satisfy the hypothesis of \textup{(f)} about this $o$. Then, according to the kind of $i$, $\|\Delta^{\rm ref}_p\|\le K_{\rm ref}\varepsilon'^2$, $\|\Delta_{\sigma,y}\|\le K_\sigma\varepsilon'^2$, $\|\Delta_{F,p}\|\le K_F\varepsilon'^2$ and $\|\Delta_{n,c}\|\le K_n\varepsilon'^2$.

\item[\textup{(g)}] For every field of $\mathfrak V$, every choice of $v_*$ and every $\psi\in\mathfrak L$, \[\langle\nu_*,\mathrm{d}^V\psi\rangle+\langle\sigma_*,\ell^V\mathrm{d}^V\psi\rangle-a\langle\omega_*,Q_1^V\psi\rangle=\langle Z_*,(W^V\mathrm{d}^V-\mathrm{d}'^VQ_1^V)\psi\rangle+\langle\Xi_*,\mathrm{d}^V\mathrm{d}^V\psi\rangle,\] and this identity is exact: it has no remainder. Moreover, at every pair $(V',e)$ and for every $b\in I^\circ$, \[|a_{0,b}|\le(8.4/\gamma)\mathsf N_{\mathsf B}\mathsf g_I(\theta_{\mathbb C})\sup|X'_P|,\] where the supremum runs over the $P$ whose site $y_P$ lies within $R+2\varrho_L+1$ of $y(b)$.

\item[\textup{(h)}] Clauses \textup{(b)} and \textup{(c)} are monotone in the radius: they remain true when $R_T$ is replaced by any integer $R'$ with $R_T\le R'\le Mr'-1$. Here the letter $R'$ denotes a free integer; it is not the radius $R(x')$. For \textup{(b)} this means that $T_\square$ is a function of $V'$ on the bonds, and of the $e_P$ at the plaquettes, of the box consisting of the sites within $R'$ of $\square$. For \textup{(c)} it means that the hypothesis becomes tameness on that box, whereas the value is still $P^T_\square(V'|W_\square)$, with the window $W_\square$ defined with $R_T$; no function $P^T_\square$ is defined with $R'$. In the setting \textup{(G1$^\flat$)} of \Cref{Y-F-C7} (item G of \S\ref{Y-F-1}), the admissible values of $R'$ include $4R+4\varrho_L+8$, which is the radius $R_5$ of \Cref{Y-F-C7}, and $4R+6\varrho_L+11$, which is the radius $R_6$ of Thms.~\ref{Y-F-C7} and~\ref{Y-G-ThA}; both radii $R_5$ and $R_6$ are those of item E8 of \S\ref{Y-F-1}.

\end{enumerate}

\item[\textup{I}] Remarks.

\begin{enumerate}

\item[\textup{(1)}] The middle member of \textup{(a)} is the quantity $x\mathcal Q(-a_0)$ of \Cref{Y-C-P67}\,(a), and $e^{-x\mathcal Q(-a_0)}$ is the prefactor in the exact formula of \Cref{Y-C-P67}\,(b). That proposition is stated for every $\mathfrak g$-valued $X'$, and so it holds at $X'=X'(V',e)$.

\item[\textup{(2)}] Clauses \textup{(a)}, \textup{(b)} and \textup{(e)} hold also when $e$ is set to $0$, that is, with $X'_P=\mathsf L(\partial V'(P))$, and with $X'_P=0$ where $\partial V'(P)=-\mathbf 1$; in this form they hold at every $V'$, and $e$ does not enter them. We keep the pairs $(V',e)$ because \Cref{Y-F-C7} uses $|X'_P|=\phi'_P$.

\item[\textup{(3)}] No clause of the theorem rests on \Cref{Y-B-C1}.

\item[\textup{(4)}] We consider the case $\ell=\emptyset$. Then $\mathsf H_y=\{0\}$, and $\emptyset\in\mathcal R(b_*,K_*,\lambda_*,\frac19)$; this follows from \Cref{Y-A-SetsOfRows} of the sets of rows $\mathcal R(b,K,\lambda,\tau)$, according to which $\emptyset\in\mathcal R(1,0,\lambda,0)$. The conclusion of \Cref{Y-A-Fr} is stated for every $k\ge0$ with $\ell_0$ empty, and item B admits this case. Here $\mathsf B^V=1$, $\Delta_Y=\bar F-Y_*$ and $\sigma_*=0$, and the steps $(y,p,p')$ of item C4 carry the weight $0$. Every clause holds with the same letters, since the constants are computed from $(K_*,\lambda_*)$ and not from the rows.

\end{enumerate}

\end{enumerate}

\end{theorem}

\begin{proof} \emph{Step 0: what the data and the floors give.} We collect the following facts.

\begin{enumerate}

\item[\textup{(0.1)}] For $L\ge3$ and $R\ge6$ we have $2r_w\le R$, $3r_w+\varrho_L+1\le R_2-1$ and $\max\{R_2,R+2\varrho_L+3\}\le R_3$; each of the numbers $R_3+2\varrho_L$, $3R+3$, $2R+2\varrho_L+3$ and $R+4\varrho_L+4$ is $\le R_T$; and $2\varrho_L+R_2\le R_T-1$, $3R+1\le R_T-1$, $2\varrho_L+\max\{R_3,3R+1\}\le R_T$, $R(x')\le R$.

\item[\textup{(0.2)}] We have $\ell_w(x')\ge6R_T+18$. Hence $3R_T+1\le\ell_w(x')$, and a path of at most $6\varrho_L+8$, $6\varrho(x')$, $3R_2$ or $3R_T$ bonds has at most $\ell_w(x')$ bonds.

\item[\textup{(0.3)}] We have $r'\ge1$ and $N\ge16(R_T+1)$. Hence every box of radius $\le R_T$, every window $W_\square$ and every set $\square^\approx$ embeds in the coarse torus, and $4R_T+2\le N$.

\item[\textup{(0.4)}] We have $W_\square\subset\square^\approx$, and no bond of $W_\square$ has an end outside $\square^\approx$ or in the outermost layer of sites of $\square^\approx$.

\item[\textup{(0.5)}] We have $\tau_{\rm c}\le\varepsilon'$. Hence \textup{(C3.2)} holds at $\hat\varepsilon=\varepsilon'$, in the form $\|\Lambda_i\|\le \mathsf K_{\rm c}\varepsilon'^2$.

\item[\textup{(0.6)}] A field with \textup{(u)} on a box satisfies, on that box, the conditions \textup{(h)} at $\hat\varepsilon=\varepsilon'$. A field of $\mathcal V'$ that is tame on a set $A$ satisfies \textup{(u)} on $A$, and on $A$ we have $X'_P(V',e)=-i\log \partial V'(P)$.

\item[\textup{(0.7)}] On a set $A$ where \textup{(u2)} holds, the operators $\operatorname{Ad}^{\pm1}$ of $V'$ along a single path in $A$ of at most $\ell_w(x')$ bonds have norm $\le4$.

\item[\textup{(0.8)}] We have $e^{-\theta_DR_2}\le x^{-128}$, $e^{-\theta_DR(x')}\le x'^{-64}$ and $\|\mathsf B'-\mathsf B'^{\rm c}\|_0\le2C_D^2x'^{-64}$.

\end{enumerate}

We prove these facts in turn.

\begin{enumerate}

\item[\textup{(0.1)}] For $L\ge3$ we have $\varrho_L\le(R+2)/3$; moreover $r_w\le R/2$, which is the first inequality. Hence $3r_w+\varrho_L+1\le11R/6+5/3\le2R+1$ for $R\ge4$, and $2R+1=R_2-1$. The inequality $R+2\varrho_L+3\le2R+\varrho_L+5$, that is $R+2\varrho_L+3\le R_3$, holds if and only if $\varrho_L\le R+2$, which is the case. The inequality $2R+3\varrho_L+5\le R_T$, that is $R_3+2\varrho_L\le R_T$, holds if and only if $\varrho_L\le2R+1$, and $R+4\varrho_L+4\le R_T$ holds if and only if $2\varrho_L\le3R+2$; both conditions are met. The remaining inequalities are read off the definitions of the radii.

\item[\textup{(0.2)}] By \textup{(G1$_k$)}, $\ell_w(x')=6Mr'+12\ge6(R_T+1)+12=6R_T+18$, so that $3R_T+1\le\ell_w(x')$. The numbers $6\varrho_L+8$, $6R$, $3R_2=6R+6$ and $3R_T$ are all $\le6R_T+18$, and $6\varrho(x')\le6R$ by \textup{(0.1)}.

\item[\textup{(0.3)}] By \textup{(G1$_k$)}, $Mr'\ge R_T+1>0$, so that $r'\ge1$. By \textup{(Tor)}, $N\ge16Mr(x)\ge16Mr'\ge16(R_T+1)$, which gives $4R_T+2\le N$. The sides of $W_\square$ and of $\square^\approx$ are $M+2R_T$ and $M(2r'+1)$, and both are $\le3Mr'<N$.

\item[\textup{(0.4)}] A site within $R_T\le Mr'-1$ of a site of $\square$ lies in a cell within $r'$ of $\square$, and its distance from $\square$ is $\le Mr'-1$, whereas the sites of the outermost layer of $\square^\approx$ are at distance $Mr'$ from $\square$.

\item[\textup{(0.5)}] By \textup{(P1)}, $\tau_{\rm c}\le x\cdot x^{-60}=x^{-59}$, while $\varepsilon'=p_0'x'^{-1/2}\ge x'^{-1/2}\ge x^{-59}$, because $p_0'\ge1$ and $x\ge x'\ge1$. Hence $\tau_{\rm c}\le\varepsilon'$, and therefore $\mathsf K_{\rm c}^0\varepsilon'^2+\tau_{\rm c}\varepsilon'\le(\mathsf K_{\rm c}^0+1)\varepsilon'^2\le\sqrt2\mathsf K_{\rm c}^0\varepsilon'^2$, where the last step rests on $\mathsf K_{\rm c}^0\ge2.42$.

\item[\textup{(0.6)}] We have $\vartheta(x')<\vartheta_0^\flat$ and $\mathsf b(x')<2$, and \textup{(u3)} is \textup{(h3)} at $\varepsilon'$; this gives the first assertion. For a real field, $|V'(c)|=|V'(\pi)|=1$, and $1<1+\vartheta(x')$, $1<\mathsf b(x')$, because $\log x'\ge\log x_\flat\ge2$. At a tame plaquette, $|\partial V'(P)-\mathbf 1|=2\sin(\phi'_P/2)\le\varepsilon'<2\varepsilon'$; and for $\phi'_P\le\varepsilon'<\pi$ the logarithm $\mathsf L$ is given by its power series, so that no $e_P$ enters. This gives the second assertion.

\item[\textup{(0.7)}] Let $g$ be the matrix of $\operatorname{SL}(2,\mathbb C)$ that $V'$ assigns to the path. Then $\|\operatorname{Ad}_g^{\pm1}\|\le|g|^2<\mathsf b(x')^2<4$.

\item[\textup{(0.8)}] We have $\theta_Dc_R=64$, $R_2\ge2c_R\log x$ and $R(x')\ge c_R\log x'$, which give the first two bounds. For the third, $\mathsf B'-\mathsf B'^{\rm c}=\ell'^{>\mathsf T}\ell'+\ell'^{{\rm c}\mathsf T}\ell'^>$, and $\|\ell'^>\|_0\le e^{-\theta_D\varrho(x')}\|\ell'\|_{\theta_D}\le C_Dx'^{-64}$.

\end{enumerate}

\emph{Step 1: the identity, and clause \textup{(a)}.} Let a field of $\mathfrak V$ be given, together with an arbitrary $\mathfrak g_{\mathbb C}$-valued $X'$, an arbitrary $Z$ on the coarse plaquettes, an arbitrary $\Xi$ on the fine cubes and an arbitrary symmetric invertible $C$, and put $\Delta:=\mathsf B^V\bar F-W^{V*}Z-\mathrm{d}^{V*}\Xi$. We claim that \[2\mathsf m_C-\langle a_0,Sa_0\rangle=\langle\bar F,\mathsf B^V\bar F\rangle+a\langle\bar n,\bar n\rangle=\langle X',Z\rangle+\langle \mathrm{d}'^VX',E_3^*\Xi\rangle+\langle EX',\Delta\rangle-\langle a_0,u^T\mathsf g^VX'\rangle,\] and that clause \textup{(a)} follows from this identity.

Throughout, the brackets denote the bilinear form; $\mathsf B^V$ and $q$ are symmetric for it, and $a_0\in\mathfrak L$.

\begin{enumerate}

\item[\textup{(i)}] The first equality. Since $a_0=C\mathsf g^VX'$, we have $\langle a_0,C^{-1}a_0\rangle=\langle a_0,\mathsf g^VX'\rangle$, and with the definitions of $\mathsf m_C$ and of $S=C^{-1}-q$ this gives \[2\mathsf m_C-\langle a_0,Sa_0\rangle=\langle EX',\mathsf B^VEX'\rangle-2\langle \mathsf g^VX',a_0\rangle+\langle a_0,qa_0\rangle.\] For $\psi\in\mathfrak L$ we have $\langle \mathsf g^VX',\psi\rangle=\langle \mathsf B^VEX',\mathrm{d}^V\psi\rangle$; moreover $\langle a_0,qa_0\rangle=\langle \mathrm{d}^Va_0,\mathsf B^V\mathrm{d}^Va_0\rangle+a\langle Q_1^Va_0,Q_1^Va_0\rangle$. Inserting these, the first with $\psi=a_0$, we find that the sum of the three terms is $\langle\bar F,\mathsf B^V\bar F\rangle+a\langle\bar n,\bar n\rangle$.

\item[\textup{(ii)}] For $\psi\in\mathfrak L$, \[\langle\bar F,\mathsf B^V\mathrm{d}^V\psi\rangle-a\langle\bar n,Q_1^V\psi\rangle=\langle \mathsf g^VX',\psi\rangle-\langle qa_0,\psi\rangle=\langle u^T\mathsf g^VX',\psi\rangle;\] the first equality follows from the definitions of $\bar F$ and $\bar n$ and from the identities in \textup{(i)}, the second from $u^T=1-qC$ and $a_0=C\mathsf g^VX'$.

\item[\textup{(iii)}] Since $\bar F=EX'-\mathrm{d}^Va_0$, we have $\langle\bar F,\mathsf B^V\bar F\rangle=\langle \mathsf B^V\bar F,EX'\rangle-\langle\bar F,\mathsf B^V\mathrm{d}^Va_0\rangle$. Taking $\psi=a_0$ in \textup{(ii)} and using $\bar n=Q_1^Va_0$, we obtain \[\langle\bar F,\mathsf B^V\bar F\rangle+a\langle\bar n,\bar n\rangle=\langle \mathsf B^V\bar F,EX'\rangle-\langle u^T\mathsf g^VX',a_0\rangle.\]

\item[\textup{(iv)}] We show that $W^VE=1$ and $\mathrm{d}^VE=E_3\mathrm{d}'^V$. By the proof of \Cref{Y-A-L1}, $W(P',p)=L^{-1}\delta_{PP'}$ at the $L$ edge plaquettes $p$ under $P$; these are based in $B(y_P)$, so that $T(y_P,y(p))=1$, and hence $(W^VEX')_P=X'_P$. For the second identity, take a fine cube at $z$ and two directions $\mu,\nu$, and let $\lambda$ be the third direction. The face $(z;\mu\nu)$ of the cube is an edge plaquette if and only if $z_\mu$ and $z_\nu$ are at the top of their block, and the opposite face $(z+e_\lambda;\mu\nu)$ is an edge plaquette under the same condition. If $z_\lambda$ is not at the top of its block, both faces lie under one $P$, in the block of $z$, with transport $1$ and with opposite signs, so that their contributions cancel; and no other face of the cube is an edge plaquette, as that would require $z_\lambda$ to be at the top. If $z_\mu$, $z_\nu$ and $z_\lambda$ are all at the top, the cube is the corner cube under $q'$. Its face at $z$ lies under $(y_{q'};\mu\nu)$, with transport $1$; its face at $z+e_\lambda$ is based in $B(y_{q'}+Le_\lambda)$ and lies under $(y_{q'}+Le_\lambda;\mu\nu)$, with transport $T(y_{q'},y_{q'}+Le_\lambda)=T(y_{q'},y_P)$; and the signs are those of $\mathrm{d}'$. Therefore $(\mathrm{d}^VEX')_q=(\mathrm{d}'^VX')_{q'}$ at the corner cube, and $(\mathrm{d}^VEX')_q=0$ at all other cubes.

\item[\textup{(v)}] By the definition of $\Delta$ and by \textup{(iv)}, \[\langle \mathsf B^V\bar F,EX'\rangle=\langle Z,W^VEX'\rangle+\langle\Xi,\mathrm{d}^VEX'\rangle+\langle\Delta,EX'\rangle=\langle X',Z\rangle+\langle E_3^*\Xi,\mathrm{d}'^VX'\rangle+\langle EX',\Delta\rangle.\] Together with \textup{(iii)}, this proves the claimed identity.

\end{enumerate}

We turn to clause \textup{(a)}. At the pair $(V',e)$ we take $(Z,\Xi)=(Z_*,\Xi_*)$; then $\Delta=\Delta_Y$. We multiply the identity by $x/2$ and use $x\langle X',Z_*\rangle=x'\langle X',v_*\rangle$. Further, $\langle EX',\Delta_Y\rangle=\sum_P\langle X'_P,(E^*\Delta_Y)_P\rangle$. Each $P$, each $q'$ and each $b\in I^\circ$ lies in exactly one cell, because $M\mid N$, so that the sums over these objects split into sums over the cells. Finally we add and subtract $x'\sum_P\mathrm{En}'_P$; this sum equals $x'\mathrm{En}'(V')$, since \[\sum_P\langle\zeta'_P,(\ell'^{{\rm c}V*}m'\ell'^{{\rm c}V}\zeta')_P\rangle=\langle\ell'^{{\rm c}V}\zeta',m'\ell'^{{\rm c}V}\zeta'\rangle=\sum_{y'}m'_{y'}|(\ell'^{{\rm c}V}\zeta')_{y'}|^2.\]

\emph{Step 2: the sets on which the objects depend.} We claim that every quantity entering $T_\square$ or $P^T_\square$, taken at an object $i$ with $o=y(i)$ in $\square$, depends on $V'$ within $R_T$ of $o$, hence only on the bonds of $W_\square$; and that it depends on $e$ only through the $e_P$ at those plaquettes of $W_\square$ where $\phi'_P=\pi$.

We give the radii of dependence about $o=y(i)$, counting all four sites of a plaquette. Recall that the distance of every site of a coordinate path from the start of the path is at most the distance between its two ends (item C3). The radius is $1$ for $X'_P$, $\zeta'_P$ and $\phi'_P$, and $\varrho(x')+1$ for $(\ell'^{{\rm c}V}\zeta')_{y'}$. For $m'_{y'}$ it is $\max\{\rho_0,\varrho(x')+1\}+1\le R(x')+2$, as $\rho_0\le R(x')$. For $v_{*P}$ and $\mathrm{En}'_P$, which involve the $y'$ within $\varrho(x')$ of $y_P$, it is $\varrho(x')+R(x')+2\le2R+2$. For $\Xi_{*q}$, which involves the $P$ within $R(x')$ of $q$, it is $3R+2$, and for $Y_{*p}$, which involves the $P$ and the $q$ within $1$ of $p$, it is $3R+3$. For $(\mathsf g^VX')_b$ it is $2\varrho_L+2$ (item C4), and for $a_{0,b}$ it is $\max\{3r_w+\varrho_L+1,\ R+2\varrho_L+2\}\le R_3-1$, by \textup{(C2.1)}. For $\bar F_p$ and $\bar n_c$, which involve the bonds within $1$ of $o$, it is $\le R_3$; for $(\mathsf B^V\bar F)_p$, which involves the $p'$ within $2\varrho_L$ of $p$ and paths within $2\varrho_L$, it is $R_3+2\varrho_L$; and for $(\mathrm{d}'^VX')_{q'}$ it is $1$.

Finally, \[(u^T\mathsf g^VX')_b=(\mathsf g^VX')_b-\sum_{b''}q(b,b'')\sum_{b'}C(b'',b')(\mathsf g^VX')_{b'}.\] Here $q(b,\cdot)$ has radius $2\varrho_L+2$; the factor $C(b'',\cdot)$ adds $R_2-1$, which leads to $2R+2\varrho_L+3$; and $(\mathsf g^VX')_{b'}$, with $b'$ within $R+2\varrho_L+2$ of $b$, adds $2\varrho_L+2$, which leads to $R+4\varrho_L+4$. By \textup{(0.1)}, all these radii are $\le R_T$.

The edge plaquettes under $P$ and the corner cube under $q'$ are based in $B(y_P)$ and in $B(y_{q'})$ respectively; hence $y(p)=y_P$ and $y(q)=y_{q'}$ for them. The second member $e$ of the pair enters through $X'$ alone. If a field is given on $W_\square$ and extended by $\mathbf 1$, nothing depends on the continuation, and $X'$ enters at no plaquette other than the plaquettes of $W_\square$ (item C6).

\emph{Step 3: telescoping in the comb gauge of a box.} We prove the following lemma.

\begin{lemma}[The difference between dressed and free products of kernels in a box]\label{Y-D-LA} Let $\mathrm{Bx}$ be the box of radius $\rho_B$ about a coarse site $o$, where $3\rho_B+1\le\ell_w(x')$ and $4\rho_B+2\le N$. Let $V'$ be an $\operatorname{SL}(2,\mathbb C)$ field on the bonds of $\mathrm{Bx}$ that satisfies \textup{(h2)} and \textup{(h3)} at some $\hat\varepsilon\le1/40$ with $35100\rho_B^2\hat\varepsilon\le1$. Let $K_1,\dots,K_n$, where $1\le n\le5$, be free kernels of finite coarse ranges whose sum is $\le\rho_B-1$; each of them is scalar or consists of rows, and rows occur either as a neighbouring pair made up of an adjoint and a row, or as the kernel $K_1$ ($k=1$). Let $A_k$ be $K_k^V$ or $(K_k^V)^*$, and let $A_k^0$ be the same object at $V'=\mathbf 1$. Let $i_0$ be an index with $y(i_0)=o$, and let the last index run over the coarse plaquettes. Put $\tilde X_P:=T(o,y_P)X'_P$ at the plaquettes of $\mathrm{Bx}$ and $\tilde X_P=0$ at all other plaquettes, where $X'_P=-i\log \partial V'(P)$ at the plaquettes of $\mathrm{Bx}$ (item~C.II~1.4 of \S\ref{Y-C-1}), and let $0<\theta\le1$. Then \[\|(A_1\cdots A_nX')_{i_0}-(A_1^0\cdots A_n^0\tilde X)_{i_0}\|\le7020\cdot8\cdot1.2^{n-1}\theta^{-2}\prod_k\|K_k\|_\theta\,\hat\varepsilon\sup_P\|\tilde X_P\|\le10^6\theta^{-2}\prod_k\|K_k\|_\theta\,\hat\varepsilon^2,\] and the first inequality holds for any $\mathfrak g_{\mathbb C}$-valued $X'$. If $n=1$ and $A_1=K_1^V$, the left side vanishes.\end{lemma}

\begin{proof} We use the following facts from item~C.II~1.4 of \S\ref{Y-C-1} and from \Cref{Y-C-L67}; they apply because $3\rho_B+1\le\ell_w(x')$ and \textup{(h2)}, \textup{(h3)} hold. The comb is $\mathsf k(y')=V'(\Gamma(o,y'))$, and for a path $\pi$ in $\mathrm{Bx}$ from $y'$ to $y''$ the field in the comb gauge is $\mathsf W(\pi)=\mathsf k(y')V'(\pi)\mathsf k(y'')^{-1}$. By \Cref{Y-C-L67}\,(b), if $\pi$ has $m\le\ell_w(x')$ bonds and all its sites lie within $\tau$ of $o$, then $\|\operatorname{Ad}_{\mathsf W(\pi)}^{\pm1}-1\|\le2340\,m\,\tau\,\hat\varepsilon$. By \Cref{Y-C-L67}\,(c), $|\tilde X_P|\le8.24\hat\varepsilon$.

We pass to the comb gauge. We have $T(y',y'')=\operatorname{Ad}_{\mathsf k(y')}^{-1}\operatorname{Ad}_{\mathsf W(\Gamma(y',y''))}\operatorname{Ad}_{\mathsf k(y'')}$, and an adjoint carries the inverse of this operator. Along a chain of indices $i_0,\dots,i_n$ the factors $\operatorname{Ad}_{\mathsf k}$ therefore cancel; moreover $\mathsf k(o)=\mathbf 1$ and $\operatorname{Ad}_{\mathsf k(y_P)}X'_P=\tilde X_P$. Consequently $(A_1\cdots A_nX')_{i_0}=(A_1^{\mathsf W}\cdots A_n^{\mathsf W}\tilde X)_{i_0}$, where $A_k^{\mathsf W}$ arises from $A_k$ when $\operatorname{Ad}_{\mathsf W(\Gamma)}^{\pm1}$ is put in the place of $T^{\pm1}$. This is an identity of matrices as soon as all paths lie in $\mathrm{Bx}$, which we now verify. Put $d_k={\rm dist}(i_{k-1},i_k)$ and $\ell=\sum_kd_k\le\rho_B-1$. A site of the path between $y(i_{k-1})$ and $y(i_k)$ lies within $d_k$ of each of the two ends, whichever of them is the start (item C3; the short way is unique, since $d_k\le\rho_B-1<N/4$); in an adjoint the path is $\Gamma(y(i_k),y(i_{k-1}))$. As the end $y(i_{k-1})$ lies within $d_1+\cdots+d_{k-1}$ of $o$, the site lies within $d_1+\cdots+d_k\le\ell$ of $o$. For the last plaquette, $y_P$ lies within $\rho_B-1$ of $o$, and so all its sites lie in $\mathrm{Bx}$. Since $4\rho_B+2\le N$, the short way is unique and $\mathrm{Bx}$ is embedded. When $n=1$ and $A_1=K_1^V$, the transport is $T(o,y_P)$ itself, and the left side of the claim is $0$.

Now consider one chain. Its $k$-th path has at most $3d_k\le\ell_w(x')$ bonds and lies within $\ell$ of $o$, so \Cref{Y-C-L67}\,(b) gives $\|\operatorname{Ad}_{\mathsf W}^{\pm1}-1\|\le2340\cdot3d_k\cdot\ell\cdot\hat\varepsilon=:\delta_k$, and $\delta_k\le7020\rho_B^2\hat\varepsilon\le0.2$. Writing $a_k$ for the factor of the $k$-th path, we telescope: $\prod a_k-\prod1=\sum_ka_1\cdots a_{k-1}(a_k-1)$. Since $\|a_j\|\le1.2$, this difference has norm $\le1.2^{n-1}\sum_k\delta_k=1.2^{n-1}\,7020\,\ell^2\hat\varepsilon$; the sum over $k$ produces $\ell^2$, and not $n\ell^2$, because $\sum_kd_k=\ell$.

To sum over the chains, we use $u^2\le2e^u$ (indeed $\sup u^2e^{-u}\le0.5414$) at $u=\theta\ell/2$, which gives $\ell^2\le(8/\theta^2)e^{\theta\ell/2}$. The sum, over the chains that start at $i_0$, of $\prod_k|K_k(i_{k-1},i_k)|e^{\theta d_k}$ is $\le\prod_k\|K_k\|_\theta$; as this norm is the larger of the weighted row sums and the weighted column sums, adjoints are covered too. For a pair of rows we use $|\langle\ell_y(p),\ell_y(p')\rangle|\le|\ell_y(p)||\ell_y(p')|$, and for rows at the place $k=1$ we use $\|\sum_p\ell_y(p)\otimes w_p\|\le\sum_p|\ell_y(p)|\,\|w_p\|$. As to the numerals, $3\cdot2340=7020$, $5\cdot7020=35100$ and $7020\cdot8\cdot1.2^4=116453.376$, and $116453.376\cdot8.24=959575.81824\le10^6$; with $|\tilde X_P|\le8.24\hat\varepsilon$, this yields the second inequality.\end{proof}

We also note three norms, for $\theta\le1$ and range $1$. First, $\|L^{-1}W\|_\theta\le eL$, the sums that enter being $L^{-1}L^2$ and $L^{-1}L^{-1}$. Second, $\|d\|_\theta\le6e$ between cubes and plaquettes, since a cube has six faces and a plaquette lies in two cubes. Third, $\|\mathrm{d}'\|_\theta\le4e$ between coarse plaquettes and bonds.

\emph{Step 4: clause \textup{(f)}.} Let $p$ be a fine plaquette and $o=y(p)$, and let an $\operatorname{SL}(2,\mathbb C)$ field satisfy \textup{(u)} on the box $\mathrm{Bx}$ of radius $R_T$ about $o$; the field is extended by $\mathbf 1$, and the objects are defined by the tame formulas. We claim that $\|(\mathsf B^V\bar F)_p-Y_{*p}\|\le K_Y\varepsilon'^2$; this is clause \textup{(f)}.

The extended field lies in $\mathfrak V^\circ$. By \textup{(0.2)}, \textup{(0.3)}, \textup{(0.6)} and \textup{(F8$_k$)}, \Cref{Y-D-LA} applies with $\rho_B=R_T$ and $\hat\varepsilon=\varepsilon'$. We define $\tilde X$ at $o$. Then $\sup\|\tilde X_P\|\le8.24\varepsilon'$.

\begin{enumerate}

\item[\textup{(1)}] The neighbours. We write $(\mathsf B^V\bar F)_p=\sum_{p'}\mathsf B^V(p,p')\bar F_{p'}$, where \[\mathsf B^V(p,p')=\delta_{pp'}+\sum_y\langle\ell_y(p),\ell_y(p')\rangle T(y(y),y(p))^{-1}T(y(y),y(p')),\] and only $p'$ with $y(p')$ within $2\varrho_L$ of $o$ occur. The transport in this formula is $\operatorname{Ad}$ of $V'$ along a single path of at most $6\varrho_L$ bonds in $\mathrm{Bx}$, so its norm is $\le4$ by \textup{(0.7)}. For the dressed entries we use the absolute row bound
\[
 \sum_{p'}\biggl(\delta_{pp'}+\sum_y|\ell_y(p)|\,|\ell_y(p')|\biggr)
 \le1+3K_*^2\mathsf g(\lambda_*)^2=\mathsf N_{\mathsf B}.
\] The box of radius $R_3$ about $y(p')$ lies in $\mathrm{Bx}$ by \textup{(0.1)}, and by \textup{(0.6)} the field satisfies the conditions \textup{(h)} at $\varepsilon'$ there. By \textup{(C3.1)}, $\bar F_{p'}$ coincides with the one computed in that box, and by \textup{(0.5)}, \[\bar F_{p'}=\sum_PH^\natural(p',P)T(y(p'),y_P)X'_P+\Lambda_{p'},\] with $\|\Lambda_{p'}\|\le \mathsf K_{\rm c}\varepsilon'^2$. Hence $\|(\mathsf B^V\Lambda)_p\|\le4\mathsf N_{\mathsf B}\mathsf K_{\rm c}\varepsilon'^2$.

\item[\textup{(2)}] The split. Put $dH:=H^\natural-H^{\sharp<}$. By \textup{(C3.4)} this kernel vanishes beyond $R_2$, and within $R_2$ it equals $H^\natural-H^\sharp$; so by \textup{(C3.5)} its absolute row sums are $\le\mathsf{ti}_{\rm c}$. The transport $T(y(p'),y_P)$ runs along at most $3R_2$ bonds in $\mathrm{Bx}$ and so has norm $\le4$. Thus \[\|(\mathsf B^V(dH)^VX')_p\|\le4\mathsf N_{\mathsf B}\cdot\mathsf{ti}_{\rm c}\cdot4\cdot2.06\varepsilon'=32.96\mathsf N_{\mathsf B}\mathsf{ti}_{\rm c}\varepsilon'.\] This bound does not depend on a weighted norm of $dH$; indeed $\mathsf{ti}_{\rm c}$ is of the order of $(\log x)^3x^{-60}$, whereas $e^{\theta_DR_2/2}$ is of the order of $x^{64}$.

\item[\textup{(3)}] We apply \Cref{Y-D-LA} to the products $(H^{\sharp<})$ and $(\ell^*,\ell,H^{\sharp<})$. Their ranges are $\le2\varrho_L+R_2\le R_T-1$, and $\|\ell\|_{\theta_D},\|H^\sharp\|_{\theta_D}\le C_D$; this gives \[\|(\mathsf B^V(H^{\sharp<})^VX')_p-(\mathsf BH^{\sharp<}\tilde X)_p\|\le10^6\theta_D^{-2}(1+C_D^2)C_D\varepsilon'^2.\]

\item[\textup{(4)}] We apply \Cref{Y-D-LA} to $Y_*=(L^{-1}W^{V*}+\mathrm{d}^{V*}(\mathsf X^<)^V)(1+\ell'^{{\rm c}V*}\ell'^{{\rm c}V})X'$, which is made of four products of $1$ to $4$ kernels, with ranges $\le1+R(x')+2\varrho(x')\le3R+1\le R_T-1$. With the norms noted at the end of Step 3, this gives \[\|Y_{*p}-((L^{-1}W^*+\delta \mathsf X^<)\mathsf B'^{\rm c}\tilde X)_p\|\le10^6\theta_D^{-2}(eL+6eC_D)(1+C_D^2)\varepsilon'^2.\]

\item[\textup{(5)}] The free identity. By \textup{(I2)}, as an identity of matrices and hence also on $\tilde X$ extended by $0$, \[\mathsf BH^{\sharp<}-(L^{-1}W^*+\delta \mathsf X^<)\mathsf B'^{\rm c}=-\mathsf BH^{\sharp>}+\delta \mathsf X^>\mathsf B'+(L^{-1}W^*+\delta \mathsf X^<)(\mathsf B'-\mathsf B'^{\rm c}).\] By \textup{(0.8)}, the three terms on the right have row sums at most $\mathsf N_{\mathsf B}C_Dx^{-128}$, $2C_Dx'^{-64}\cdot C_D$ and $(L^{-2}+2C_D)\cdot2C_D^2x'^{-64}$. Their sum is $\le(\mathsf N_{\mathsf B}C_D+8C_D^3)x'^{-64}$, because $2C_D^2+(2/9)C_D^2+4C_D^3\le8C_D^3$ for $C_D\ge4$; and it is to be multiplied by $8.24\varepsilon'$.

\item[\textup{(6)}] The sum. The terms of \textup{(2)} and \textup{(5)} are linear in $\varepsilon'$, and together they are at most \[[33\mathsf N_{\mathsf B}\mathsf{ti}_{\rm c}+8.24(\mathsf N_{\mathsf B}C_D+8C_D^3)x'^{-64}]\varepsilon'\le[33\mathsf N_{\mathsf B}(1+C_D)\mathsf{ti}_{\rm c}+100\mathsf N_{\mathsf B}C_D^4x'^{-64}]\varepsilon'\le\varepsilon'^2,\] the last inequality by \textup{(F9$_k$)}. Adding the bounds of \textup{(1)}, \textup{(3)} and \textup{(4)}, we arrive at the total $4\mathsf N_{\mathsf B}\mathsf K_{\rm c}+10^6\theta_D^{-2}(1+C_D^2)(C_D+eL+6eC_D)+1=K_Y$.

\end{enumerate}

\emph{Step 5: the sizes of the three sums on $U^T_\square(x')$, which give the sup bound of \textup{(d)} and the bound of \textup{(c)}.} Let $V'\in U^T_\square(x')$, extended by $\mathbf 1$. We claim that $|P^T_\square(V')|\le K_{{\rm q},k}p_0'^2\varepsilon'$.

The extended field lies in $\mathfrak V^\circ$. For $o$ in $\square$, the box of radius $R_T$ about $o$ lies in $W_\square$, so that \textup{(u)} holds on it. We note that $x\varepsilon'^2=Lp_0'^2$.

\begin{enumerate}

\item[\textup{(1)}] The first sum. There are $3M^3$ plaquettes $P$ in $\square$, and for each of the $L$ edge plaquettes $p$ under such a $P$ we have $y(p)=y_P$ in $\square$. Clause \textup{(f)} is stated for a field that is given on the box of radius $R_T$ about $o=y(p)$ and extended by $\mathbf 1$, while here the field is given on $W_\square$. By Step 2, where the radii are $R_3+2\varrho_L$ for $(\mathsf B^V\bar F)_p$ and $3R+3$ for $Y_{*p}$, both $\le R_T$ about $o$ when the four sites of a plaquette are counted, and by item C6, the quantity $\Delta_{Y,p}$ is a function of $V'$ on the bonds of that box and of $X'$ at its plaquettes, and of nothing else; hence it takes the same value at the two fields. The restriction to the box satisfies \textup{(u)} there, because the paths and plaquettes of the box are paths and plaquettes of $W_\square$. Therefore \textup{(f)} gives \[|\frac x2\langle X'_P,(E^*\Delta_Y)_P\rangle|\le\frac x2\cdot2.06\varepsilon'\cdot L\cdot K_Y\varepsilon'^2=1.03L^2K_Yp_0'^2\varepsilon',\] and the first sum is at most $3.09M^3L^2K_Yp_0'^2\varepsilon'$.

\item[\textup{(2)}] The second sum. We use the Bianchi identity, which the proof of \Cref{Y-C-L67}\,(d) establishes as an identity between invertible matrices. For a coarse cube $q'$ at $z$ it states that $\prod_{k=1}^6\kappa_k\,\partial V'(P_k)^{\varepsilon_k}\kappa_k^{-1}=\mathbf 1$, where $\varepsilon_k=\mathrm{d}'(q',P_k)$, $\kappa_k=\mathbf 1$ at the faces based at $z$, and $\kappa_k=V'(z,\lambda)$ at the face based at $z+Le_\lambda$. With $Y_k=\varepsilon_k\operatorname{Ad}_{\kappa_k}X'_{P_k}$ we get $\prod e^{iY_k}=\mathbf 1$ and $\sum_kY_k=(\mathrm{d}'^VX')_{q'}$, the latter because $T(z,z+Le_\lambda)=\operatorname{Ad}_{V'(z,\lambda)}$. Each $|Y_k|\le(1+\vartheta_0^\flat)^2\cdot2.06\varepsilon'$, so that $S=\sum|Y_k|\le12.3625\varepsilon'\le0.31$. By \textup{(E)} of item~C.I~1.1 of \S\ref{Y-C-1}, which item D lists under \textup{(Tools)}, \[|\sum Y_k|=|\prod e^{iY_k}-\mathbf 1-i\sum Y_k|\le e^S-1-S\le0.72S^2,\] where the ratio of $e^S-1-S$ to $S^2$ is $\le0.556$ at $S=0.31$; and this is $\le110.04\varepsilon'^2\le126\varepsilon'^2$. Next we bound the multipliers. We have $v_*=X'+\ell'^{{\rm c}V*}\ell'^{{\rm c}V}X'$, where the transports run along a single path of at most $6\varrho(x')$ bonds and so have norm $\le4$; since $\|\ell'\|_0\le C_D$, this gives $\|v_{*P}\|\le(1+4C_D^2)\cdot2.06\varepsilon'$. With paths of at most $3R(x')$ bonds, $\|\Xi_{*q}\|\le C_D\cdot4\cdot\|v_*\|_{\sup}=8.24C_D(1+4C_D^2)\varepsilon'$. Summing over the $M^3$ cubes, we find that the second sum is at most \[M^3\frac x2\cdot126\varepsilon'^2\cdot8.24C_D(1+4C_D^2)\varepsilon'=519.12M^3LC_D(1+4C_D^2)p_0'^2\varepsilon'\le520M^3LC_D(1+4C_D^2)p_0'^2\varepsilon'.\]

\item[\textup{(3)}] The third sum. There are at most $3L^3M^3$ bonds $b$ in $\hat\square$. For each of them $\|a_{0,b}\|\le \mathsf a_{\rm c}\varepsilon'$, by \textup{(C3.3)} in the box of radius $R_3$ about $o=y(b)$: we take $i=c$, the coarse bond $(y(b),e_1)$, so that $o=c_-=y(b)$ and ${\rm dist}(b,o)=0\le1$. Since $u^T=SC$, we have $\|u^T\|\le x^{-60}\cdot4/\gamma$ at the extended field, by \textup{(C2.4)} and \textup{(C2.5)}. The row $b$ of $u^T=1-qC$ vanishes beyond $R+2\varrho_L+2$, so that at most $3L^3(2R+4\varrho_L+5)^3$ bonds $b'$ occur. At these $\|(\mathsf g^VX')_{b'}\|\le4\mathsf N_{\mathsf B}\cdot4\cdot2.06\varepsilon'\le33\mathsf N_{\mathsf B}\varepsilon'$, by \Cref{Y-C-CIIL64}\,(d) and by \textup{(0.7)} for one path $\pi_{\mathsf t}$ of at most $6\varrho_L+8$ bonds. Consequently \[\FitDisplay{\dimexpr\linewidth-1pt\relax}{\|(u^T\mathsf g^VX')_b\|\le\|u^T\|\,(\text{number of }b')^{1/2}\sup\|(\mathsf g^VX')_{b'}\|\le(4/\gamma)x^{-60}(3L^3)^{1/2}(2R+4\varrho_L+5)^{3/2}\cdot33\mathsf N_{\mathsf B}\varepsilon'.}\] With the numeral $3\cdot\frac12\cdot4\cdot33=198$, the third sum is at most \[\FitDisplay{\dimexpr\linewidth-1pt\relax}{198L^3M^3\mathsf a_{\rm c}\gamma^{-1}\mathsf N_{\mathsf B}(3L^3)^{1/2}(2R+4\varrho_L+5)^{3/2}x^{-59}\varepsilon'^2\le M^3\varepsilon'^2,}\] the last inequality by \textup{(F9$_k$)}; and $M^3\varepsilon'^2\le M^3p_0'^2\varepsilon'$, since $\varepsilon'\le1/40\le1\le p_0'^2$ by \textup{(F8$_k$)}.

\end{enumerate}

The sum of the bounds in \textup{(1)} to \textup{(3)} is $K_{{\rm q},k}p_0'^2\varepsilon'$.

\emph{Step 6: clauses \textup{(c)}, \textup{(d)} and \textup{(b)}.} We show that $T_\square=P^T_\square$ on tame windows; that $P^T_\square$ is holomorphic; that $P(\{\square\},\cdot)$ is holomorphic, invariant and real, vanishes at pure gauges and satisfies \textup{(ca)}; and that $T_\square$ is real, Borel and invariant.

For \textup{(c)}, let $V'\in\mathcal V'$ be tame on $W_\square$. A plaquette of $W_\square$ has its base within $R_T$ of $\square$; hence the hypothesis that $\phi'_P\le\varepsilon'$ within $R_T+1$ of $\square$ implies tameness on $W_\square$. By \textup{(0.6)}, $V'|W_\square\in U^T_\square(x')$, and $X'(V',e)=-i\log \partial V'$ on $W_\square$. By Step 2, all that $T_\square$ depends on lies in $W_\square$. There $\zeta'=X'$, and we meet neither a plaquette of $D'$ nor one with $\phi'_P>1$; so $m'=1$ at every $y'$ that enters, and the objects are those of the tame formulas. Since $S_{\varepsilon'}(\phi'_P)=\frac12|X'_P|^2$, we obtain, row by row, \[x'\mathrm{En}'_P=\frac{x'}2\langle X'_P,X'_P+(\ell'^{{\rm c}V*}\ell'^{{\rm c}V}X')_P\rangle=\frac{x'}2\langle X'_P,v_{*P}\rangle.\] What is left of $T_\square$ is therefore $P^T_\square(V'|W_\square)$. The operators $C$, $q$ and $u^T$ are evaluated at $V'$ itself, but only entries depending on $V'$ in $W_\square$ occur (Step 2), and these are the entries of $V'|W_\square$ extended by $\mathbf 1$. The bound in \textup{(c)} is that of Step 5.

For \textup{(d)}, we prove holomorphy. The function $P^T_\square$ is a finite sum of products of factors of three kinds. The factors of the first kind are the entries of $\operatorname{Ad}^{\pm1}_{V'(c)}$; they are polynomials in the entries of $V'(c)$, because $V'(c)^{-1}$ is the adjugate. Those of the second kind are the $X'_P$, which are power series in $\partial V'(P)-\mathbf 1$, where $|\partial V'(P)-\mathbf 1|<2\varepsilon'<1$. Those of the third kind are the entries of $C^V$; by \textup{(C2.3)} they are holomorphic on $\mathfrak V^\circ$, and they are composed with the extension by $\mathbf 1$, which maps $U^T_\square(x')$ into $\mathfrak V^\circ$ because $|V'(c)|<1+\vartheta(x')<1+\vartheta_0^\flat$. The restriction $\mathcal U_{\{\square\}}(x')\to U^T_\square(x')$ is well defined by \textup{(0.4)}, and it is a projection of coordinates.

Next, invariance. Let $u$ be $\operatorname{SU}(2)$-valued on $\square^\approx$ and extended by $\mathbf 1$. Then $V'^u\in \mathcal U_{\{\square\}}(x')$, since the norms are unchanged. Now $P(\{\square\},V')=P^T_\square(V'|W_\square)$ is evaluated at $V'|W_\square$ extended by $\mathbf 1$ off $W_\square$. By Step 2 and item C6, the three sums are functions of $V'$ on the bonds of $W_\square$ and of $X'$ at the plaquettes of $W_\square$, and of nothing else; they therefore have the same value at every field of $\mathfrak V$ which coincides with $V'$ on the bonds of $W_\square$. One such field is $V'$ extended by $\mathbf 1$ off $\square^\approx$, and we work with it. This extension of $V'$, transformed by $u$, and the extension of $V'^u$ by $\mathbf 1$ off $\square^\approx$ both belong to $\mathfrak V$, and they differ only on bonds with one end off $\square^\approx$; by \textup{(0.4)}, no such bond is a bond of $W_\square$. By \textup{(C2.2)}, by the covariance of all dressed kernels and of $X'$, and by the orthogonality of $\operatorname{Ad}_u$ for $Y\cdot Y'$, each bracket is unchanged.

Reality at $\operatorname{SU}(2)$-valued fields holds because there the transports are real on $\mathfrak g$, $X'$ takes values in $\mathfrak g$, and $C$ is real by \textup{(C2.3)}. At a pure gauge, given by a $u$ with values in $\operatorname{SU}(2)$ or in $\operatorname{SL}(2,\mathbb C)$ such that the field lies in the domain, we have $\partial V'=\mathbf 1$ and hence $X'=0$; as each of the three sums is bilinear in $X'$, $P(\{\square\},\cdot)$ vanishes there.

The bound on the supremum follows from Step 5 and from $K_{{\rm q},k}p_0'^2\varepsilon'\le\frac1{32}E_ap_0'^2\varepsilon'=\frac1{32}c_a(x')$; for $\#Y\ge2$ we have $P(Y,\cdot)=0$. The properties obtained are the five clauses of \textup{(ca)}, which item~1 of \S\ref{Y-G0-1} states for members as \textup{(ca$_A$)} at $A=1$ together with reality: holomorphy on $\mathcal U_Y(x')$; invariance under $\operatorname{SU}(2)$-valued gauge transformations; reality at $\operatorname{SU}(2)$-valued fields; vanishing at the pure gauges; and $\sup\le c_a(x')e^{-\kappa(\#Y-1)}$.

For \textup{(b)}: the dependence is that of Step 2. Reality follows by the argument given for $P(\{\square\},\cdot)$: at $\operatorname{SU}(2)$-valued fields the transports are real on $\mathfrak g$, $X'$ is $\mathfrak g$-valued, and $C$ is real. The function $T_\square$ is Borel, since $C$ is continuous on $\mathfrak V$ by \textup{(C2.3)}, $X'$ is Borel in $(V',e)$, and $\phi'$, $\zeta'$, $m'$ are Borel. As to invariance under $(V',e)\to(V'^u,\operatorname{Ad}_ue)$: $X'$ and $\zeta'$ are covariant, $\phi'$, $D'$ and $m'$ are invariant, and $C$ and the dressed kernels are covariant.

\emph{Step 7: clauses \textup{(e)} and \textup{(h)}, and remark \textup{(2)} of item I.} We prove the decomposition over an arbitrary class partition, together with the set on which $w_j$ depends; the monotonicity in the radius; and the case $e=0$.

For \textup{(e)}: the set $\square^\approx$ holds a plaquette of $D'$ if and only if a mark lies within $r'$ cells of $\square$. Hence the cells with $1_{\square^\approx}(V')=0$ form exactly the union of the sets $\mathrm{mk}_j'^{+r'}$. These sets are disjoint, as a common cell would put two classes within $2r'$ of each other. In a cell with $1_{\square^\approx}(V')=1$, every plaquette based in $\square^\approx$ is tame; so $V'|\square^\approx\in \mathcal U_{\{\square\}}(x')$ by \textup{(0.6)}, $V'$ is tame on $W_\square$ by \textup{(0.4)}, and \textup{(c)} yields $T_\square=P^T_\square(V'|W_\square)=P(\{\square\},V')$. Summing \textup{(a)} over the cells now gives the identity of \textup{(e)}. As to the dependence, a cell of $\mathrm{mk}_j'^{+r'}$ lies within $r'$ of $\mathrm{mk}'_j$, and $T_\square$ depends on $V'$ within $R_T\le Mr'$ coarse units, that is, within $r'$ cells of $\square$; the two distances add up to $2r'\le8r'=s(x')$. For the classes of the chain, any two classes are farther apart than $2s_0+4r_1+2>4r'$, for every $r_1\ge0$.

For \textup{(h)}: a function of $V'$ on the bonds within $R_T$ of $\square$ is a function of $V'$ on the bonds within $R'$, and tameness within $R'$ implies tameness on $W_\square$. Further, $4R+6\varrho_L+11\le Mr'-1$ follows from $Mr'\ge5R+4\varrho_L+12$ together with $2\varrho_L\le R$. The former is part of \textup{(G1$^\flat$)}, the setting in which \textup{(h)} admits the radii $R_5$ and $R_6$; as to the latter, the proof of \textup{(0.1)} gives $\varrho_L\le(R+2)/3$, so that $2\varrho_L\le(2R+4)/3\le R$ for $R\ge4$. Thus $4R+6\varrho_L+11\le5R+4\varrho_L+11\le Mr'-1$, and $R_T\le R_5\le R_6$.

For remark \textup{(2)}: the identity of Step 1 holds for any $X'$. With $X'_P=\mathsf L(\partial V'(P))$ and $\zeta'_P=X'_P\min\{1,\varepsilon'/\phi'_P\}$, the sum of the $\mathrm{En}'_P$ is still $\mathrm{En}'(V')$: since a row $y'$ with $m'_{y'}=1$ meets no $P$ with $\phi'_P>1$, the number $\mathrm{En}'$ is the same. The proofs of \textup{(b)}, \textup{(c)} and \textup{(e)} depend on $|X'_P|=\phi'_P$ at tame plaquettes only.

\emph{Step 8: clauses \textup{(f$'$)} and \textup{(g)}, for the triple $\boldsymbol\nu=(\nu_*,\sigma_*,\omega_*)$ of item E6 of \S\ref{Y-F-1}.} We show the following: for a field with \textup{(u)} on the box of radius $R_T$ about $o=y(i)$, as in \textup{(f$'$)}, the objects being defined by the tame formulas, $\|\Delta^{\rm ref}_p\|\le K_{\rm ref}\varepsilon'^2$, $\|\Delta_{\sigma,y}\|\le K_\sigma\varepsilon'^2$, $\|\Delta_{F,p}\|\le K_F\varepsilon'^2$ and $\|\Delta_{n,c}\|\le K_n\varepsilon'^2$; the identity of \textup{(g)}; and the sup bound on $a_0$ at every $(V',e)$. For three of these bounds a smaller box suffices: for $\Delta^{\rm ref}_p$, the box of radius $R_{\rm ref}$ about $y(p)$, by \textup{(1)} below; for $\Delta_{\sigma,y}$, the box of radius $\varrho_L+R_{\rm ref}$ about $y(y)$, by \textup{(2)}; and for $\Delta_{n,c}$, the box of radius $R_3$ about $c_-$, by \textup{(3)}. Here $R_{\rm ref}$ is the radius that \textup{(1)} sets. The bound on $\Delta_{F,p}$ uses the bound on $\Delta_{\sigma,y}$ at the $y$ within $\varrho_L$ of $p$, with the smaller box of \textup{(2)}; for $\Delta_{F,p}$ itself the box remains the one of radius $R_T$ about $y(p)$, because of \textup{(f)}.

\begin{enumerate}

\item[\textup{(1)}] The bound on $\Delta^{\rm ref}_p$. Let $o=y(p)$, and let \textup{(u)} hold on the box of radius $R_{\rm ref}:=\max\{R_3,3R+1\}$ about $o$. By \textup{(C3.1)} and \textup{(0.5)}, and since $T(o,y_P)X'_P=\tilde X_P$, we have the exact equality $\bar F_p=(H^\natural\tilde X)_p+\Lambda_p$, with $\|\Lambda_p\|\le \mathsf K_{\rm c}\varepsilon'^2$. We split $(H^\natural\tilde X)_p=(H^{\sharp<}\tilde X)_p+(dH\,\tilde X)_p$, where $\|(dH\,\tilde X)_p\|\le8.24\mathsf{ti}_{\rm c}\varepsilon'$. \Cref{Y-D-LA}, with $\rho_B=R_{\rm ref}\le R_T$, applied to $(\mathsf T^<)$ and to $(\mathsf T^<,\ell'^{{\rm c}*},\ell'^{\rm c})$, whose ranges are $\le3R$, gives $\|\nu^{\rm ref}_p-(\mathsf T^<\mathsf B'^{\rm c}\tilde X)_p\|\le10^6\theta_D^{-2}C_D(1+C_D^2)\varepsilon'^2$. On the free side, \textup{(I2)} gives \[H^{\sharp<}-\mathsf T^<\mathsf B'^{\rm c}=-H^{\sharp>}+\mathsf T^>\mathsf B'+\mathsf T^<(\mathsf B'-\mathsf B'^{\rm c}),\] with row sums $\le(C_D+C_D^2+2C_D^3)x'^{-64}$, which are to be multiplied by $8.24\varepsilon'$. The terms linear in $\varepsilon'$ add up to at most $[8.24\mathsf{ti}_{\rm c}+8.24(C_D+C_D^2+2C_D^3)x'^{-64}]\varepsilon'\le\varepsilon'^2$, by \textup{(F9$_k$)}. The sum of all contributions is $K_{\rm ref}$.

\item[\textup{(2)}] The bounds on $\Delta_\sigma$ and $\Delta_F$. We have $\Delta_\sigma=\ell^V(\bar F-\nu^{\rm ref})=\ell^V\Delta^{\rm ref}$. The $p$ with $\ell_y(p)\ne0$ lie within $\varrho_L$ of $y$, and their boxes of radius $R_{\rm ref}$ lie within $\varrho_L+R_{\rm ref}\le R_T$ of $y(y)$; the transports have norm $\le4$, and $\sum_p|\ell_y(p)|\le C_D$; hence $\|\Delta_{\sigma,y}\|\le4C_DK_{\rm ref}\varepsilon'^2$. Next, $\Delta_F=\bar F-Y_*+\ell^{V*}\sigma_*=\Delta_Y-\ell^{V*}\Delta_\sigma$ identically, because $\mathsf B^V=1+\ell^{V*}\ell^V$. Only the $y$ within $\varrho_L$ of $p$ occur, and $2\varrho_L+R_{\rm ref}\le R_T$ by \textup{(0.1)}; so $\|(\ell^{V*}\Delta_\sigma)_p\|\le C_D\cdot4\cdot K_\sigma\varepsilon'^2=16C_D^2K_{\rm ref}\varepsilon'^2$. Combined with \textup{(f)}, this gives the constant $K_F$.

\item[\textup{(3)}] The bound on $\Delta_{n,c}$. Let $o=c_-$, and consider the box of radius $R_3$ about $o$. Here $\bar n_c=(\mathsf N^\natural\tilde X)_c+\Lambda_c$ exactly. The kernel $dN:=\mathsf N^\natural-\mathsf N^{\sharp<}$ vanishes beyond $R_2$ and has row sums $\le\mathsf{ti}_{\rm c}$, so that $\|(dN\,\tilde X)_c\|\le8.24\mathsf{ti}_{\rm c}\varepsilon'$. We apply \Cref{Y-D-LA} to $((aL)^{-1}\mathrm{d}'^*)$ and to $((aL)^{-1}\mathrm{d}'^*,\ell'^{{\rm c}*},\ell'^{\rm c})$, with ranges $\le1+2\varrho(x')$; the norms contribute $(aL)^{-1}4e(1+C_D^2)\le(1+C_D^2)$, since $4e/18\le0.605$, and therefore $\|\omega_{*c}-((aL)^{-1}\mathrm{d}'^*\mathsf B'^{\rm c}\tilde X)_c\|\le10^6\theta_D^{-2}(1+C_D^2)\varepsilon'^2$. On the free side, \[\mathsf N^{\sharp<}-(aL)^{-1}\mathrm{d}'^*\mathsf B'^{\rm c}=-\mathsf N^{\sharp>}+(aL)^{-1}\mathrm{d}'^*(\mathsf B'-\mathsf B'^{\rm c}),\] with row sums $C_Dx^{-128}+(4/18)2C_D^2x'^{-64}$, to be multiplied by $8.24\varepsilon'$. The terms linear in $\varepsilon'$ are $\le\varepsilon'^2$, by \textup{(F9$_k$)}. The sum of all contributions is $K_n$.

\item[\textup{(4)}] The identity \textup{(g)}. For $\psi\in\mathfrak L$, \[\langle\nu_*,\mathrm{d}^V\psi\rangle+\langle\sigma_*,\ell^V\mathrm{d}^V\psi\rangle=\langle\nu_*+\ell^{V*}\sigma_*,\mathrm{d}^V\psi\rangle=\langle Y_*,\mathrm{d}^V\psi\rangle=\langle Z_*,W^V\mathrm{d}^V\psi\rangle+\langle\Xi_*,\mathrm{d}^V\mathrm{d}^V\psi\rangle,\] and \[a\langle\omega_*,Q_1^V\psi\rangle=L^{-1}\langle \mathrm{d}'^{V*}v_*,Q_1^V\psi\rangle=\langle Z_*,\mathrm{d}'^VQ_1^V\psi\rangle.\] Subtracting the second identity from the first, we obtain \textup{(g)}.

\end{enumerate}

It remains to prove the sup bound on $a_0$; let $V'\in\mathcal V'$. By \textup{(C2.1)}, \Cref{Y-C-CIIL65}\,(a) and \textup{(C2.5)}, \[\sum_{b'}\|C(b,b')\|\le\sum_{{\rm dist}\le R}[\|q^{-1}(b,b')\|+\|(C-q^{-1})(b,b')\|]\le(2/\gamma)\mathsf g_I(\theta_{\mathbb C})+8\gamma^{-2}x^{-60}\,3L^3(2R+1)^3.\] By \textup{(P1)}, the second term is $\le x^{-59}/(165(L+1)\mathsf N_{\mathsf B}(1+b))\le1/(165\cdot4\cdot17)<10^{-4}$, whereas $(0.1/\gamma)\mathsf g_I(\theta_{\mathbb C})\ge0.1\cdot5913\cdot81$. Hence the sum is $\le(2.1/\gamma)\mathsf g_I(\theta_{\mathbb C})$. Moreover $|(\mathsf g^VX')_{b'}|\le4\mathsf N_{\mathsf B}\sup|X'_P|$, with the supremum over the $P$ with $y_P$ within $2\varrho_L+1$ of $y(b')$, since the transports are isometries (\Cref{Y-C-CIIL64}\,(d)). As $2.1\cdot4=8.4$, the bound of \textup{(g)} follows.

\emph{Step 9: the law, the order of the constants, and the floors.} We show that the law of windows has \textup{(C2.0)} to \textup{(C2.5)}; that the hypotheses of \Cref{Y-C-Th68} and of \Cref{Y-C-L71}\,(a), (b) are met; that the letters are fixed without a circle; and that each floor of item G holds for all large $x$ and depends neither on $N$, $\ell$, $c_h$ or $E_a$ nor on a member.

The law of windows is that of item~C.I~3 of \S\ref{Y-C-1}, \Cref{Y-C-CITh66}. The head of \Cref{Y-C-CITh66}, before its clauses, states that they hold at every field with \textup{(u1$^\flat$)}, that is, on $\mathfrak V$; this is \textup{(C2.0)}. \Cref{Y-C-CIL64} and \Cref{Y-C-CITh66} require $L\ge3$ odd, $aL\ge18$ and rows as in item~C.I~1.1 of \S\ref{Y-C-1}; \Cref{Y-C-CIIL64} and \Cref{Y-C-Th68} require rows as in item~C.II~1.2 of \S\ref{Y-C-1}; and Lems.~\ref{Y-C-CIIL65}, \ref{Y-C-L66}, \ref{Y-C-L67} and~\ref{Y-C-L71} are stated under the same requirements. These requirements are met. Indeed, $L\ge3$ odd and $aL\ge18$ is \textup{(O4)}, and $N$ is even by item B. Item B supplies rows of the kind described in item~C.I~1.1 of \S\ref{Y-C-1}, and equally of the kind described in item~C.II~1.2 of \S\ref{Y-C-1}, at $(K,\lambda,\varrho)=(K_*,\lambda_*,\varrho(x))$, since it gives the bound $|\ell_y(p)|\le K_*e^{-\lambda_*|y-\pi(p)|_\infty}$ and the range $\le\varrho(x)$. Here $K_*\ge0$ and $0<\lambda_*\le1$, because these are the pure numbers of \textup{(O0)}, which \Cref{Y-A-PureNumbers} gives as $K_*=3\gamma_*^{-1/2}$ and $\lambda_*=\gamma_*^{1/2}/44$, and because $\gamma_*\le1$, as clause~\textup{(0)} of the order of the constants in \Cref{Y-A-Fr} gives $\gamma_*=\mathsf r_\sharp^2/12296\le9.05513\cdot10^{-25}$. Further, $\varrho(x)/L+2\le\varrho_L+2<N/2$, since $N\ge16(R_T+1)$ by \textup{(0.3)} and $R_T+1=4R+2\varrho_L+7>\varrho_L+2$. At $\ell=\emptyset$ the same requirements read $0\le K_*e^{-\lambda_*|y-\pi(p)|_\infty}$ and range $0\le\varrho(x)$.

\Cref{Y-C-CITh66}\,(b) gives \textup{(C2.1)}: the range is $2r_w$ and the radius of dependence is $3r_w+\varrho_L+1$, and by \textup{(0.1)} these are $\le R$ and $\le R_2-1$, which are \textup{($\rho$1)} and \textup{($\rho$2)} of the proposition ``What Part C.II uses of the law'' in item~C.II~2.1 of \S\ref{Y-C-1}. \Cref{Y-C-CITh66}\,(c) gives \textup{(C2.2)}. \Cref{Y-C-CITh66}\,(i), together with \Cref{Y-C-CITh66}\,(a), by which $G$ is real at $\mathcal V'$ and so $C=GG^{\mathsf T}$ is real, gives \textup{(C2.3)}. \Cref{Y-C-CITh66}\,(f), (k), which apply since \textup{(F$_{\rm law}$)} and $c_R\theta_{\mathbb C}\ge128$ hold, give \textup{(C2.4)} in the norm of $\ell^2(I^\circ;\mathfrak g_{\mathbb C})$, that is, \textup{($\rho$5)} of the same proposition with $\mathsf n=60$. Finally \textup{(C2.5)} is \Cref{Y-C-L66}\,(a): with $\mathsf w:=q^{-1}S$ we have $\|\mathsf w\|\le(2/\gamma)x^{-60}\le1/2$ by \textup{(P1)}, hence $C=(1+\mathsf w)^{-1}q^{-1}$ and $q^{-1}-C=\mathsf w(1+\mathsf w)^{-1}q^{-1}$.

For the input \textup{(C3)}: \Cref{Y-C-Th68} requires \textup{($\lambda$0)}, \textup{($\rho$1)}, \textup{($\rho$2)}, \textup{($\rho$5)}, \textup{(Dec)} for $\mathsf G,\mathsf G_Z$, \textup{(P0)} and \textup{(P1)}. \Cref{Y-C-L71}\,(a) assumes \textup{($\rho$1)} and $R\ge6$. \Cref{Y-C-L71}\,(b) requires \textup{(Fr1)}, \textup{(Dec)} for $\mathsf G,\mathsf G_Z,H^\sharp,\mathsf N^\sharp$, and \textup{($\lambda$0)}, \textup{($\rho$1)}, \textup{($\rho$5)}, \textup{(P0)}, \textup{(P1)}; here \textup{(Fr1)} is \Cref{Y-A-L7}\,(c), and \textup{($\lambda$0)} and \textup{($\rho$1)} are \textup{(C2.0)} and \textup{(C2.1)}. All clauses of \Cref{Y-C-L71} are stated under \textup{(Fr1)} to \textup{(Fr3)}. Of these, \textup{(Fr1)} is algebra and holds for all rows, and \textup{(Fr2)} is \textup{(Rows)}. Condition \textup{(Fr3)} holds as well, although the proof of \Cref{Y-C-L71}\,(b) does not depend on it: by the remark ``Source and hypotheses'' in item~C.II~2.2 of \S\ref{Y-C-1}, it holds for every $\ell$ of $\mathcal R(b_*,K_*,\lambda_*,\frac19)$ in the setting \textup{(H1)} $L\ge L_A$ odd, \textup{(H2)} $aL\ge3/c_L$, \textup{(H3)} $N$ even, which is \textup{(O4)} together with item B. Hence the hypotheses of the lemma are met. Finally, \textup{(Tor)}, which Part C.II of \S\ref{Y-C-1} lists in \textup{(P0)}, belongs to the data. The hypotheses of \Cref{Y-C-Th68} are met as well: the preceding paragraph gives \textup{($\rho$2)} and \textup{($\rho$5)}, \textup{(Dec)} for $\mathsf G,\mathsf G_Z,H^\sharp,\mathsf N^\sharp$ is \textup{(Rows)} of item D, and \textup{(P0)} and \textup{(P1)} are floors of item G.

The letters are fixed in the following order. First come $\gamma,\mathsf N_{\mathsf B},b,A^\flat,\theta_{\mathbb C},\vartheta_0^\flat$, which depend on $(L,a,K_*,\lambda_*)$; then $\theta_D,C_D,c_R$; then $\mathsf K_A,\mathsf K_{\rm c}^0,\mathsf a_{\rm c},\mathsf K_{\rm c}$; then $K_Y$, together with $K_{\rm ref},K_\sigma,K_F,K_n$. After $M$ comes $K_{{\rm q},k}$, and then $c_h$, on which the proof does not depend. Then comes the condition $E_a\ge32K_{{\rm q},k}$, which is a lower bound on $E_a$. It can be met, because \Cref{Y-B-C1} is stated for all $E_a$ above its own two floors, namely for all $E_a\ge E_a^0$, where $E_a^0=\max\{32L^{-1/2}K_v,32K_{\rm q}\}$ is the number of \textup{(O12)} in item~3 of \S\ref{Y-G0-1} and $K_{\rm q}=K_{{\rm q},k}$. No clause of \Cref{Y-D-C5} uses this fact about \Cref{Y-B-C1}. The condition $E_a\ge32K_{{\rm q},k}$ itself enters only the last inequality of \textup{(d)} and, through that inequality, \textup{(ca)} with $\frac1{32}$ of the cap. Last comes $x_\flat$. The constant $\mathsf a_{\rm c}$ enters none of the constants $K$, only \textup{(F9$_k$)}, and nothing earlier in the order depends on $E_a$; hence there is no circle.

Finally, each floor holds for all large $x$. In \textup{(F$_{\rm law}$)}, in \textup{(P1)} and in the second half of \textup{(F9$_k$)}, powers of $\log x$ are compensated by powers of $x$. In \textup{(P0)} we need $x\ge L^2$, and $R\ge6$ holds once $c_R\log x>5$. For \textup{(G1$_k$)}, $Mr'\ge20c_R\log x'$, while \[R_T+1\le4(c_R\log x+1)+2((c_R\log x+1)/3+1)+7=(14/3)c_R\log x+41/3;\] since $\log x=\log x'+\log L$, the floor holds as soon as $(46/3)c_R\log x'\ge(14/3)c_R\log L+41/3$; and $\rho_0\le R(x')$ holds for large $x$. In \textup{(F8$_k$)}, $\varepsilon'=A_0(\log x')^{\mathfrak p_0}x'^{-1/2}$ is compared with $R_T^{-2}$, and $p_0'\ge1$ for large $x$. In the first half of \textup{(F9$_k$)}, $\mathsf{ti}_{\rm c}\le{\rm const}(\log x)^3x^{-60}+{\rm const}\,x^{-64}$ and $x'^{-64}$ are measured relative to $\varepsilon'\ge x'^{-1/2}$. None of the floors depends on $N$, $c_h$, $E_a$ or a member, nor on more of $\ell$ than $(K_*,\lambda_*,\text{range})$.\end{proof}

\subsection{Local vertices}\label{Y-E-1}

We compare the exact integrand in the chart $\eta=G\xi-a_0$
with its quadratic approximation. The cell remainder $V_\square$
collects the differences in the plaquette action, the row energy and
the step kernel, as well as the logarithm of the Haar density. The
corresponding real factor $A_\square$ is defined on the entire chart.
The key estimate is a local holomorphic bound of order
$p_0^3x^{-1/2}$ for $V_\square$ when the test $\chi_\square$ passes.
In that region $A_\square=e^{V_\square}$.

\begin{remark}[Letters with more than one meaning] The letter $b$ is a fine bond where it is an index or an argument, and the number $b=16\mathsf N_{\mathsf B}+a/L$ of \textup{(O6)} where it is the base of a power ($b^{-1}$, $b^{-1/2}$, $b^{1/2}$); the letter $d$ is the coboundary, but in $b\in dp$ the set $dp$ consists of the four bonds of $p$. The letter $y$ is a coarse site in $B(y)$ and in $T(y,y'')$; as the index of a row $\ell_y$ it is a fine site, and then $y(y)$ is its coarse site. The letter $u$, where it acts on a field or on a vector, is a gauge transformation; in $\rho[u;\ell]$ it is the triple $(E^a,E^h,w)$ of members, and both meanings occur in Step~1\,\textup{(3)} of the proof of \Cref{Y-E-C6}.\end{remark}

\begin{theorem}[The vertex letter at a level of any depth]\label{Y-E-C6} The setting is the group $G=\operatorname{SU}(2)$ in dimension $d=3$, with the drawn step of item~C.I~4 (``The step, the chart'') of \S\ref{Y-C-1}.

\emph{Item 1: the letters.} The letters follow the order \textup{(O0)} to \textup{(O13)} of item~3 of \S\ref{Y-G0-1}. The whole list is that of item~3 of \S\ref{Y-G0-1}, and the functions of $L$ are written out in item~A of \S\ref{Y-F-1}. Some signs are only named below and are set elsewhere: $\mathsf l_L$ and $\mathsf P_\rho$ in item~A of \S\ref{Y-F-1}, under \textup{(O3)}; $A^\flat$ in item~C.I~1.2 of \S\ref{Y-C-1}; $c_1^*$, $\lambda_2^*$, $\theta_J$ and $C_D$ in the statement of \Cref{Y-A-Fr}, where $c_1^*$ is set in the order of the constants \textup{(0)}; \textup{(M2)}, \textup{(M4$^\flat$)} and \textup{(M5r)} in item~3 of \S\ref{Y-G0-1}, under \textup{(O7)}; $\mathsf K_{\rm c}^0$ in item~C.II~4 of \S\ref{Y-C-1}, under ``Constants''; $\delta_3'$ in item~F of \S\ref{Y-F-1}; $K_{{\rm q},k}$ in item~F of \S\ref{Y-D-1}; \textup{(F$_{\rm law}$)} in \Cref{Y-C-CITh66}\,\textup{(k)}; \textup{(P0)} and \textup{(P1)} in item~C.II~4 of \S\ref{Y-C-1}, under ``Floors'', and \textup{(P3)} there, before \Cref{Y-C-L71}; $\mathfrak E$ in item~E2 of \S\ref{Y-D-1}; and $\mathcal R(b_*,K_*,\lambda_*,\frac19)$ in Defs.~\ref{Y-A-CovRows}, \ref{Y-A-Defect} and~\ref{Y-A-SetsOfRows}.

\begin{enumerate}\item[\textup{(O0)}] The pure numbers of \S\ref{Y-A-1}, of which the theorem uses $K_*$ and $\lambda_*$, where $0<\lambda_*\le1$. With them come the function $\mathsf g(t)=\sum_{n\in\mathbb Z^3}e^{-t|n|_\infty}$ of \Cref{Y-A-Distances} and the function $\mathsf g_1(t)=\sum_n|n|_\infty e^{-t|n|_\infty}$ of the display of $\mathsf g_1$, $\mathsf g_2$ in \Cref{Y-A-Sums}; we put $\mathsf g_*=\mathsf g(\lambda_*)$, $\mathsf g_1=\mathsf g_1(\lambda_*)$ and $\mathsf g_I(t)=3L^3\mathsf g(t)$.

\item[\textup{(O1)}] A number $d_0\in(0,\frac12)$, which enters through the function $S_\varepsilon$ and, by the display \eqref{eq:row-mask-radius} of \S\ref{Y-A-1}, through $\rho_0$.

\item[\textup{(O3)}] The radius $\rho_0$ of the display \eqref{eq:row-mask-radius} of \S\ref{Y-A-1}, which enters the mask and \textup{(W0)}. Further, $\mathsf N_{\mathsf B}=\mathsf N_Z=1+3K_*^2\mathsf g_*^2$, $\mathsf N_\sharp=c_1^*\mathsf l_L\mathsf N_Z\mathsf P_\rho$ and $\mathsf N_{\rm c}=17(\mathsf N_\sharp+1)$; the last two enter \textup{(d)} alone.

\item[\textup{(O4)}] An odd $L\ge L_*^{\rm su}$, with the floor $L_*^{\rm su}$ of item~3 of \S\ref{Y-G0-1}, and $a\ge3/(c_LL)$. Of this, the theorem uses that $L\ge3$ is odd and that $aL\ge18$. Clause \textup{(d)} depends on the condition \textup{(L2$^{\rm su}$)}, $L^{1/2}\ge1.2\cdot10^4\mathsf N_{\rm c}^2$, only through $\mathsf N_{\rm c}L^{-3/2}\le0.2$.

\item[\textup{(O6)}] From item~C.I~1.2 of \S\ref{Y-C-1}, taken at $(K_*,\lambda_*)$: $\gamma=(73L^4)^{-1}$, $b=16\mathsf N_{\mathsf B}+a/L$, the number $A^\flat$, and $\theta^\flat=\min\{1,\lambda_*L/12,\gamma/(2A^\flat)\}$. Then $\theta_{\mathbb C}=\theta^\flat/6$ and $\vartheta_0^\flat=\theta^\flat/24$, which is at most $10^{-4}$. Further, $\theta_D=\frac12\min\{\lambda_2^*,\theta_{\mathbb C},\theta_J\}$, $C_D\ge4$ and $c_R=64/\theta_D$; hence $c_R\theta_{\mathbb C}\ge128$ and $c_R\ge768$.

\item[\textup{(O7)}] The cell side $M=L^{m'}$, $m'\ge1$, with \textup{(M2)}, \textup{(M4$^\flat$)} and \textup{(M5r)}. The theorem depends on no condition on $M$.

\item[\textup{(O8)}] Numbers $\mathfrak p_0>2$ and $A_0>0$. We put $p_0(x)=A_0(\log x)^{\mathfrak p_0}$, $\varepsilon=p_0x^{-1/2}$ and $x'=x/L$, and a prime means ``at $x'$''. Then $\varepsilon'\le L^{1/2}\varepsilon$; indeed $\varepsilon'=A_0(\log x')^{\mathfrak p_0}L^{1/2}x^{-1/2}$, and $0<\log x'\le\log x$ because $x'=x/L\ge L>1$.

\item[\textup{(O9)}] The numbers $\mathsf G_0=(2/\gamma)^{1/2}\mathsf g_I(\theta_{\mathbb C})\le C_\varphi^\flat=(2b^{1/2}/\gamma)\,\mathsf g_I(\theta_{\mathbb C})/(1-e^{-\theta_{\mathbb C}})$; further $\mathsf n=60$ and $c_\xi=(6.5\,C_\varphi^\flat)^{-1}$.

\item[\textup{(O10)}] From item~C.II~4 of \S\ref{Y-C-1}, under ``Constants'', the number $\mathsf a_{\rm c}=(100/\gamma)\mathsf N_{\mathsf B}\mathsf g_I(\theta_{\mathbb C})$, which exceeds $7$, and the number $\mathsf K_{\rm c}^0$ of that item; we put $\mathsf K_{\rm c}:=\sqrt2\mathsf K_{\rm c}^0$, which is the number of \textup{(Cen)}. The theorem has three numbers of its own. The first is $C_\eta^\flat:=\frac1{6.5}+\max\{\mathsf a_{\rm c},7\}L^{1/2}$, and for $\tau\ge1$ we set $C_\eta^\flat(\tau):=\frac1{6.5}+\max\{\mathsf a_{\rm c},7\}\,\tau L^{1/2}$. The second is $K_\ell:=291K_*^2\mathsf g_*\,[1320\mathsf g_1+216\mathsf g_*]$. The third is $K_v^\flat:=(LM)^3(1157+474a+K_\ell)(C_\eta^\flat)^3$, and $K_v^\flat(\tau)$ is given by the same formula with $C_\eta^\flat(\tau)$ in place of $C_\eta^\flat$. These three numbers are computed from $(L,M,a,K_*,\lambda_*)$ alone, and from $\tau$ where it occurs.

\item[\textup{(O11)}] A number $c_h$ with $0<c_h\le\delta_3'/2$, on which the theorem does not depend.

\item[\textup{(O12)}] A number $E_a\ge\max\{32K_{{\rm q},k},\,1500K_v^\flat L^{-1/2}\}$, whose two entries are those of item~3 of \S\ref{Y-G0-1} under \textup{(O12)}. Although the condition \textup{(O12)} involves $K_v^\flat$, the order of the constants is not circular, because no statement and no proof of this section depends on \textup{(O12)}.

\item[\textup{(O13)}] The letter $x_\flat$ comes last; it is taken so that \textup{(W0)} to \textup{(W5)} of item~6 hold for all $x\ge Lx_\flat$.\end{enumerate}

\emph{Item 2: the data.} The data are the data \textup{(SD)} of item~1 of \S\ref{Y-G0-1}. They consist of a coupling $x\ge L^2$; of an even coarse side $N$ with $M\mid N$ and \textup{(Tor)} $N\ge2Ms(x)=16Mr(x)$, together with the tori $\mathbb T=\mathbb Z^3/NL\mathbb Z^3$ and $\mathbb T'=L\mathbb Z^3/NL\mathbb Z^3$; and of rows $\ell\in \mathcal R(b_*,K_*,\lambda_*,\frac19)$ of range at most $\varrho(x)$ such that the pair $(\ell,\ell')$ satisfies the four conclusions of \Cref{Y-A-Th36}.

The empty rows are admitted. Indeed, $\ell=\emptyset$, with $\mathsf H_y=\{0\}$ for all $y$, has range $0$ and lies in $\mathcal R(1,0,\lambda_*,0)$, which is contained in $\mathcal R(b_*,K_*,\lambda_*,\frac19)$ by the display and the last sentence of \Cref{Y-A-SetsOfRows}. The data refer to no level and to no depth. At $\ell=\emptyset$, the second sum in $V_\square$, the second product in $A_\square$ and the mask are void, and $K_\ell$ is a bound for $0$; all numbers and radii remain those evaluated at $(K_*,\lambda_*,\varrho(x))$.

The four conclusions are stated in full in item~1 of \S\ref{Y-G0-1}, under \textup{(SD)}. The fourth, \textup{(L1)}, \textup{(L2)}, is taken in the form \textup{(L1$^{\rm su}$)}, \textup{(L2$^{\rm su}$)} of item~3 of \S\ref{Y-G0-1}, under \textup{(O4)}, as \textup{(SD)} prescribes. The first, \textup{(Dec)}$(C_D,\theta_D)$, is taken at the $\theta_D$ of \textup{(O6)}, which is defined with $\theta_{\mathbb C}$; in this form, and for all sixteen kernels of \textup{(Dec)}, it is the conclusion of \Cref{Y-A-Fr}, whereas \Cref{Y-A-P35} has $\theta_q$ in the place of $\theta_{\mathbb C}$.

The four conclusions imply \textup{(Fr2)} and \textup{(Fr3)} of item~C.II~2.2 of \S\ref{Y-C-1}, and \textup{(Fr1)} is \Cref{Y-A-L7}\,\textup{(c)}, which is algebra and holds for all rows. Of the rows $\ell$, the proofs of \textup{(a)}, \textup{(b)}, \textup{(c)}, \textup{(e)}, \textup{(f)}, \textup{(g)} and \textup{(h)} depend only on the bound $|\ell_y(p)|\le K_*e^{-\lambda_*|y-p|}$ and on the range, while \textup{(d)} uses the rest, through \Cref{Y-C-CIICor72}\,\textup{(a)}.

\emph{Item 3: the objects.}

\begin{definition}[Item 3.1: the lattice] We use the objects of item~C.I~1.1 of \S\ref{Y-C-1}. There are the blocks $B(y)$; the base site $\pi(i)$ of an object $i$, which is its least corner; the coarse site $y(i)$ whose block contains $\pi(i)$; the site $y_P=\pi(P)$; and the distance ${\rm dist}(i,j)=L^{-1}|y(i)-y(j)|_\infty$. The combs are $\Gamma(y,x)$, $\mathcal T$ is the set of their bonds, $I^\circ$ is the set of fine bonds off $\mathcal T$, and $\mathfrak L_{\mathfrak g}$ is the space of $\mathfrak g$-valued functions on $I^\circ$. Each coarse bond $c$ has its crossing bonds, and each coarse plaquette $P$ has the edge plaquettes under it. Besides the coboundary $d$, we use $(Q_1A)(c)=L^{-3}\sum_{x\in B(c_-)}A([x,x'])$ and the edge lift, $(EX')_p=X'_P$ on the $L$ edge plaquettes under $P$ and $0$ elsewhere. In fine units, $|y-p|=|y-\pi(p)|_\infty$. At a level, $\Gamma(z,z')$ is the coordinate path that runs along $e_1$, $e_2$, $e_3$ in this order, each time the short way.

A cell $\square$ is a cube of $M^3$ coarse sites of a fixed grid, and $\hat\square$ is the set of the fine sites of the blocks of $\square$. We write $p,b,y\in\hat\square$ for a fine plaquette, bond or site that is based there, and $c\in\square$ for a coarse bond with $c_-\in\square$. Let ${\rm dist}(i,\square)=\min_{y\in\square}L^{-1}|y(i)-y|_\infty$; let $\square_{[j]}$ be the box of the coarse sites within $j$ of $\square$; and let $\square^{+k}$ be $\square$ enlarged by $k$ layers of cells; its sites are $\square_{[kM]}$. Finally $r(x)=\lceil20c_R\log x/M\rceil$, $r'=r(x')$, $s(x)=8r(x)$ and $\square^\approx=\square^{+r'}$.\end{definition}

\begin{definition}[Item 3.2: the group and the norms] Here $G=\operatorname{SU}(2)$, and $\mathfrak g$ is the space of hermitian traceless $2\times2$ matrices; we use the angle $\operatorname{ang}$ and the logarithm $\mathsf L(g)$ on $G$, with $\mathsf L(-1)=0$. The operator norm is written $|\cdot|$. On $\mathfrak g_{\mathbb C}=\mathfrak{sl}_2(\mathbb C)$ we use the bilinear form $Y\cdot Y'=\frac12\operatorname{tr}YY'$ and the norm given by $\|Y\|^2=\frac12\operatorname{tr}Y^*Y$; on $\mathfrak g$, $\|Y\|=|Y|$. For $|g-1|<1$ we put $L_{\mathbb C}(g)=-i\sum_{k\ge1}\frac{(-1)^{k-1}}k(g-1)^k$. Further, $\mathsf j(Y)=s(Y\cdot Y)^2$ with $s(t)=\sum_k(-t)^k/(2k+1)!$; on $\mathfrak g$ this equals $(\sin|Y|/|Y|)^2$. The real Euclidean space of the row $y$ is $\mathsf H_y$. On $\mathsf H_y\otimes\mathfrak g_{\mathbb C}$, in an orthonormal basis $e_i$, we write $A=\sum e_i\otimes A_i$ and set $A\cdot A'=\sum A_i\cdot A'_i$ and $\|A\|^2=\sum\|A_i\|^2$.\end{definition}

\begin{lemma}[Complex $2\times2$ matrices: inverse, logarithm, products of exponentials, trace and Haar density]\label{Y-E-L1C} In this lemma $g$, $u$ and $Y_k$ are complex $2\times2$ matrices.

\begin{enumerate}\item[\textup{(a)}] For $g\in\operatorname{SL}(2,\mathbb C)$ we have $|g^{-1}|=|g|\ge1$, $|g^{-1}-1|=|g-1|$, $\|\operatorname{Ad}_g\|\le|g|^2$ and $\|\operatorname{Ad}_g-1\|\le2|g||g-1|$; moreover $\operatorname{tr}Y$ and $Y\cdot Y'$ are invariant under $\operatorname{Ad}_g$.

\item[\textup{(b)}] Let $|g-1|=t<1$. Then the series $L_{\mathbb C}(g)$ converges, $|L_{\mathbb C}(g)|\le t+t^2/(2(1-t))$, $e^{iL_{\mathbb C}(g)}=g$. For every invertible $u$, the same logarithm series converges at $ugu^{-1}$ and equals $uL_{\mathbb C}(g)u^{-1}$. If $\det g=1$, then $\operatorname{tr}L_{\mathbb C}(g)=0$; if $\det g=1$ and $t<\frac12$, then $L_{\mathbb C}(g^{-1})=-L_{\mathbb C}(g)$; and if $g\in\operatorname{SU}(2)$, then $L_{\mathbb C}(g)=\mathsf L(g)$.

\item[\textup{(c)}] Let $S=\sum|Y_k|\le\frac14$ and $g=\prod e^{iY_k}$. Then $|g-1|\le e^S-1$ and $|L_{\mathbb C}(g)-\sum Y_k|\le\frac32S^2$.

\item[\textup{(d)}] Let $Y\in\mathfrak{sl}_2(\mathbb C)$ and $t=Y\cdot Y$ with $|t|\le1$. Then $|1-\frac12\operatorname{tr}e^{iY}-\frac t2|\le|t|^2/23$ and $|\log \mathsf j(Y)|\le|t|/2$, where $\log \mathsf j(Y):=2\log s(t)$ with the principal value of the logarithm. For all $Y$ and $Y'$ in $\mathfrak g_{\mathbb C}$, $|Y\cdot Y'|\le\|Y\|\,\|Y'\|$.\end{enumerate}\end{lemma}

\begin{proof}\emph{Step (a): the inverse and the adjoint action.} Let $g\in\operatorname{SL}(2,\mathbb C)$. Then $g^{-1}$ is the adjugate of $g$. For $2\times2$ matrices the adjugate is linear in the matrix: $\operatorname{adj}A=JA^{\mathsf T}J^{-1}$ with the unitary $J=\bigl(\begin{smallmatrix}0&1\\-1&0\end{smallmatrix}\bigr)$. Since transposition and conjugation by a unitary matrix leave the operator norm unchanged, we have $|\operatorname{adj}A|=|A|$; this gives $|g^{-1}|=|g|$ and, by linearity, $|g^{-1}-1|=|\operatorname{adj}(g-1)|=|g-1|$. Moreover $|g|^2\ge|\det g|=1$. For the two bounds on $\operatorname{Ad}_g$ we refer to the paragraph on norms in item~C.II~1.1 of \S\ref{Y-C-1}. Finally, the trace of a conjugate equals the trace of the matrix, so $\operatorname{tr}(gYg^{-1})=\operatorname{tr}Y$ and $\operatorname{tr}(gYg^{-1}gY'g^{-1})=\operatorname{tr}YY'$; these two identities prove the last assertion of \textup{(a)}, that $\operatorname{tr}Y$ and $Y\cdot Y'$ are invariant under $\operatorname{Ad}_g$.

\emph{Step (b): the logarithm.} Let $|g-1|=t<1$. The $k$-th term of the series $L_{\mathbb C}(g)$ has norm at most $t^k/k$, and $\sum_{k\ge1}t^k/k\le t+(t^2/2)\sum_{k\ge0}t^k$, which is the bound stated in \textup{(b)}; so the series converges and obeys that bound. With $A=g-1$, the identity $e^{iL_{\mathbb C}(g)}=g$ reads $\exp(\log(1+A))=1+A$, an identity of power series that converge absolutely for $|A|<1$. At $ugu^{-1}$ each series term is the conjugate of the corresponding term at $g$. Its norm is bounded by $|u|\,|u^{-1}|t^k/k$, which proves absolute convergence and the covariance formula.

Now let $\det g=1$, and let $z$ and $1/z$ be the eigenvalues of $g$. They are eigenvalues of $1+(g-1)$ and of $1+(g^{-1}-1)$ respectively, and $|g^{-1}-1|=|g-1|=t$ by \textup{(a)}; hence $|z-1|\le t$ and $|1/z-1|\le t$. So both lie in the region $\operatorname{Re}w>0$, in which the series is the principal logarithm and $\log(1/z)=-\log z$. In a basis in which $g$ is triangular, the diagonal of $L_{\mathbb C}(g)$ consists of $-i\log z$ and $-i\log(1/z)$; hence the trace of $L_{\mathbb C}(g)$ is $0$.

If moreover $t<\frac12$, then $|L_{\mathbb C}(g)|\le-\log(1-t)<\log2$. The identity $\log(\exp X)=X$ holds as an identity of power series that converge absolutely for $|X|<\log2$. At $X=-iL_{\mathbb C}(g)$ we have $\exp X=g^{-1}$, and there the identity says $L_{\mathbb C}(g^{-1})=-L_{\mathbb C}(g)$. Finally, let $g\in\operatorname{SU}(2)$ have the eigenvalues $e^{\pm i\vartheta}$ with $0\le\vartheta\le\pi$. Then $|g-1|=2\sin(\vartheta/2)<1$ gives $\vartheta<\pi/3$, and the series produces the matrix that has the eigenvectors of $g$, carries the eigenvalues $\pm\vartheta$ on them and is hermitian and traceless; this matrix is $\mathsf L(g)$.

\emph{Step (c): products of exponentials.} By \textup{(E)} of item~C.I~1.1 of \S\ref{Y-C-1}, $|g-1-i\sum Y_k|\le e^S-1-S$, so $|g-1|\le e^S-1=:w\le0.2841$. The terms of the series with $k\ge2$ give $|L_{\mathbb C}(g)+i(g-1)|\le w^2/(2(1-w))$. We split $L_{\mathbb C}(g)-\sum Y_k=[L_{\mathbb C}(g)+i(g-1)]-i[g-1-i\sum Y_k]$, so $|L_{\mathbb C}(g)-\sum Y_k|$ is at most the sum of the two bounds found. Both bounds, $e^S-1-S$ and $w^2/(2(1-w))$, divided by $S^2$, increase with $S$; at $S=\frac14$ they are at most $0.5445$ and $0.9014$, and their sum is at most $1.4459\le\frac32$. This proves $|L_{\mathbb C}(g)-\sum Y_k|\le\frac32S^2$. Since $S\mapsto e^S-1$ and $S\mapsto\frac32S^2$ are increasing functions of $S$, both estimates of \textup{(c)} hold when any bound $\le\frac14$ for $\sum|Y_k|$ replaces $S$.

\emph{Step (d): trace and Haar density.} Since $\operatorname{tr}Y=0$ and $\operatorname{tr}Y^2=2t$, the Cayley--Hamilton theorem gives $Y^2=t1$. Hence $\frac12\operatorname{tr}e^{iY}=\sum_k(-t)^k/(2k)!$. In $1-\frac12\operatorname{tr}e^{iY}-\frac t2$ the terms with $k\le1$ cancel, and the terms with $k\ge2$ have modulus at most $|t|^2(\cosh1-\frac32)\le0.04309|t|^2\le|t|^2/23$; this is the first bound. Next, $|s(t)-1|\le(\sinh1-1)|t|=:u\le0.17521$, so $|\log \mathsf j(Y)|=2|\log s(t)|\le2u/(1-u)\le0.4249|t|$, which gives the second bound. Finally, the Cauchy--Schwarz inequality gives $|\operatorname{tr}YY'|\le(\operatorname{tr}Y^*Y)^{1/2}(\operatorname{tr}Y'^*Y')^{1/2}$, and after division by $2$ this is $|Y\cdot Y'|\le\|Y\|\,\|Y'\|$.\end{proof}

\begin{definition}[Item 3.3: fields and domains] The set of $G$-valued coarse fields is $\mathcal V'$. We put \[\mathfrak V=\{\operatorname{SL}(2,\mathbb C)\text{ fields on }\mathbb T',\ V'(\bar c)=V'(c)^{-1},\ |V'(c)|\le1+\vartheta_0^\flat\},\] and $\mathfrak V^\circ$ is defined in the same way with $<$. Further, $\vartheta(x)=\vartheta_0^\flat(1-1/\log x)$, $\mathsf b(x)=2-1/\log x$, the room $\mathsf g_x=\vartheta_0^\flat\log L/(4(\log x)^2)$, and $\ell_w(x)=6Mr(x)+12$.

Given a set $B$ of coarse sites, the bonds, the paths and the plaquettes of $B$ are those all of whose sites belong to $B$. For $\tau\ge1$, the domain $U^\tau(B;x')$ is the set of $\operatorname{SL}(2,\mathbb C)$ fields on the bonds of $B$ that satisfy three clauses:

\begin{enumerate}\item[\textup{(e1)}] $|V'(c)|<1+\vartheta(x')+\mathsf g_x$;

\item[\textup{(e2)}] $|V'(\pi)|<\mathsf b(x')+\mathsf g_x$ for paths in $B$ of at most $\ell_w(x')$ bonds;

\item[\textup{(e3)}] $|\partial V'(P)-1|<2\tau\varepsilon'$.\end{enumerate}

The domain $\mathcal U_{\{\square\}}(x')$ is the set of $\operatorname{SL}(2,\mathbb C)$ fields on the bonds of $\square^\approx$ that satisfy

\begin{enumerate}\item[\textup{(u1)}] $|V'(c)|<1+\vartheta(x')$;

\item[\textup{(u2)}] $|V'(\pi)|<\mathsf b(x')$ for paths in $\square^\approx$ of at most $\ell_w(x')$ bonds;

\item[\textup{(u3)}] $|\partial V'(P)-1|<2\varepsilon'$.\end{enumerate}

Hence $\mathcal U_{\{\square\}}(x')\subset U^1(\text{sites of }\square^\approx;x')$. We write $\operatorname{cont}(V')$ for $V'$ extended by $1$ to the other bonds of $\mathbb T'$. Finally, $T(y,y'')=\operatorname{Ad}_{V'(\Gamma(y,y''))}$, and for a scalar kernel $K$ we put $K^V(i,j)=K(i,j)T(y(i),y(j))$; in an adjoint, the transport appears as $T^{-1}$.\end{definition}

\begin{definition}[Item 3.4: the step and the chart] The chart is the fine field $\mathsf U(\eta;V')$ with $\mathsf U(\eta;V')_b=1$ on $\mathcal T$, $e^{i\eta_b}$ at the other bonds inside blocks, and $e^{i\eta_b}V'(c)$ at a crossing bond of $c$. At any $G$-valued fine field $U$, let $h_{c,x}=U(\Gamma(y,x))\,U([x,x'])\,U(\Gamma(y',x'))^{-1}$ for $c=(y,y')$ and $x\in B(y)$. With it, the three formulas below define $\tau_{c,x}$, $\bar\tau_c$ and $\mathsf M(U)(c)$ at every such $U$, as in item~C.I~4 of \S\ref{Y-C-1}, and $\mathsf M$ is a Borel map of $U$ computed from $U$, $L$ and the blocks alone. The formulas are $\tau_{c,x}=\mathsf L(h_{c,x}h_{c,y}^{-1})$, with $L_{\mathbb C}$ in place of $\mathsf L$ at complex fields, $\bar\tau_c=L^{-3}\sum_x\tau_{c,x}$ and $\mathsf M(U)(c)=e^{i\bar\tau_c}h_{c,y}$. In the chart, $h_{c,x}=U([x,x'])$.

With $s=ax$ and with the number $z(s)$ of item~C.I~4 of \S\ref{Y-C-1}, the kernel of the step is $k_s(g,h)=z(s)^{-1}e^{-s(1-\cos\operatorname{ang}(gh^{-1}))}$, and $\mathsf n_c=\operatorname{ang}(V'(c)\mathsf M(U)(c)^{-1})$. Finally, $\mathfrak D=\{|\eta_b|<\pi\text{ for all }b\}$ and $\mathsf c_0=(2\pi^2)^{-\#I^\circ}$.\end{definition}

\begin{definition}[Item 3.5: rows and energy] The dressed rows act by $(\ell^VF)_y=\sum_p\ell_y(p)\otimes T(y(y),y(p))F_p$, and $\mathsf B^V=1+\ell^{V*}\ell^V$. The quadratic form is $q(\eta)=q^V_\ell(\eta)=\langle \mathrm{d}^V\eta,\mathsf B^V\mathrm{d}^V\eta\rangle+a\langle Q_1^V\eta,Q_1^V\eta\rangle$, with the bilinear pairing of item~3.2; at real $V'$ and $\eta$ it equals $|\mathrm{d}^V\eta|^2+|\ell^V\mathrm{d}^V\eta|^2+a|Q_1^V\eta|^2$. The kernel is $\mathsf g^V=\Pi_{\mathfrak L}\mathrm{d}^{V*}\mathsf B^VE$, where $\Pi_{\mathfrak L}$ is the restriction to $I^\circ$. At a fine field $U$ we set $\ell_y^UF=\sum_p\ell_y(p)\otimes\operatorname{Ad}_{U(\Gamma(y,\pi(p)))}F_p$, where $\Gamma$ is now the fine coordinate path.

At a real field $U$ we put $F_p=\mathsf L(\partial U(p))$, $\phi_p=\operatorname{ang}\partial U(p)$, $D(U)=\{\phi_p>\varepsilon\}$ and $\zeta_p=F_p\min\{1,\varepsilon/\phi_p\}$, with multiplier $1$ at $\phi_p=0$. Thus $|F_p|\le\phi_p$, with equality unless $\partial U(p)=-1$; at that point $F_p=\zeta_p=0$ and $\phi_p=\pi$, so the plaquette remains rough. The function $S_\varepsilon$ is $S_\varepsilon(t)=t^2/2$ for $t\le\varepsilon$, and $d_0\varepsilon t$ for $t>\varepsilon$. The mask is given by $m_y(U)=1$ if and only if no $p$ of $D(U)$ has $|y-p|\le\rho_0$ and no $p$ with $\phi_p>1$ has $|y-p|\le\varrho(x)+1$; otherwise it is $0$.

The class $C^{\rm su}_\flat(x;\ell)$ uses exactly this mask, with both of its radii. In \S\ref{Y-A-1}, the mask of \Cref{Y-A-Rows} keeps only the first of the two clauses; the difference is harmless, because \textup{(Fr3$'$)}, in the conclusion of \Cref{Y-A-Fr}, is stated for every mask $m$. The energy is $\mathrm{En}_\ell(U)=\sum_pS_\varepsilon(\phi_p)+\frac12\sum_ym_y(U)|\ell_y^U\zeta|^2$.\end{definition}

\begin{definition}[Item 3.6: the law] The law is that of \Cref{Y-C-CIDef4}, in item~C.I~3 of \S\ref{Y-C-1}, taken at $R=R(x)=\varrho(x)=\lceil c_R\log x\rceil$. Let $r_w=\lfloor R/2\rfloor$ and $Y_b=\{b'\in I^\circ:{\rm dist}(b,b')\le r_w\}$. Put $K_b=1_{Y_b}-b^{-1}1_{Y_b}q1_{Y_b}$ and $c_k=4^{-k}\binom{2k}k$. The matrix of the law is $G(b,b')=b^{-1/2}\sum_kc_k(K_b)^k(b,b')$ on $Y_b$, and $0$ elsewhere. Further, $C=GG^{\mathsf T}$, $\mathsf E=G^{\mathsf T}qG-1$ and $S=C^{-1}-q$, and the field is drawn as $\xi\sim \mathcal N(0,x^{-1}1)$ on $\mathfrak L_{\mathfrak g}$.\end{definition}

\begin{definition}[Item 3.7: the curvature, the centre and the field] At a pair $(V',e)\in\mathcal V'\times\mathfrak E$, as in item~E2 of \S\ref{Y-D-1}, we put $X'_P=X'_P(V',e)=\mathsf L(\partial V'(P))+\pi e_P1\{\partial V'(P)=-1\}$. At a field $V'$ that lives on the bonds of a set $B$ and satisfies \textup{(e3)}, the field in use is $\operatorname{cont}(V')$, and we put $X'_P:=L_{\mathbb C}(\partial V'(P))$ for each $P$ with all four sites in $B$ and $X'_P:=0$ for every other $P$; this is the convention of item~C.II~1.4 of \S\ref{Y-C-1}.

In both cases, for real $\xi\in\mathfrak L_{\mathfrak g}$, we set $h=G^{\mathsf T}\mathsf g^VX'$, $a_0=Gh=C\mathsf g^VX'$, $\eta=G\xi-a_0$, $U=\mathsf U(\eta;V')$, $F^\flat=EX'+\mathrm{d}^V\eta$, $\bar F=EX'-\mathrm{d}^Va_0$ and $n^\flat=-Q_1^V\eta$. Finally $\mathcal Q(\eta)=\frac12F^\flat\cdot \mathsf B^VF^\flat+\frac a2(Q_1^V\eta)\cdot(Q_1^V\eta)$ and $\mathsf m_C=\frac12EX'\cdot \mathsf B^VEX'-\frac12\mathsf g^VX'\cdot C\mathsf g^VX'$.\end{definition}

\begin{definition}[Item 3.8: the vertex term] With $F_p:=L_{\mathbb C}(\partial U(p))$, the vertex term is \[\FitDisplay{\dimexpr\linewidth-1pt\relax}{\begin{aligned}
&V_\square(\xi,V'):=\frac x2\sum_{p\in\hat\square}\bigl[F^\flat_p\cdot F^\flat_p-F_p\cdot F_p\bigr]
\\
&\quad{}+\frac x2\sum_{y\in\hat\square}\bigl[(\ell^VF^\flat)_y\cdot(\ell^VF^\flat)_y-(\ell_y^UF)\cdot(\ell_y^UF)\bigr]
\\
&\quad{}+\sum_{c\in\square}\bigl[\tfrac s2n^\flat_c\cdot n^\flat_c-s\bigl(1-\tfrac12\operatorname{tr}V'(c)\mathsf M(U)(c)^{-1}\bigr)\bigr]
\\
&\quad{}+\sum_{b\in I^\circ\cap\hat\square}\log \mathsf j(\eta_b).
\end{aligned}}\]\end{definition}

\begin{definition}[Item 3.9: the vertex letter] At every real $\xi$ and every pair $(V',e)$, the vertex letter is \[\FitDisplay{\dimexpr\linewidth-1pt\relax}{\begin{aligned}
&A_\square(\xi,V',e):=\prod_{p\in\hat\square}\exp\bigl\{\tfrac x2|F^\flat_p|^2-xS_\varepsilon(\phi_p)\bigr\}
\\
&\quad{}\cdot\prod_{y\in\hat\square}\exp\bigl\{\tfrac x2|(\ell^VF^\flat)_y|^2-\tfrac x2m_y(U)|\ell_y^U\zeta|^2\bigr\}
\\
&\quad{}\cdot\prod_{c\in\square}\exp\bigl\{\tfrac s2|n^\flat_c|^2-s(1-\cos\mathsf n_c)\bigr\}
\\
&\quad{}\cdot\prod_{b\in I^\circ\cap\hat\square}1\{|\eta_b|<\pi\}\,\mathsf j(\eta_b)\ \ge0.
\end{aligned}}\]\end{definition}

\begin{definition}[Item 3.10: the tests] Let $R'\ge R_\chi^{\min}$ be an integer. The test is given by $\chi_\square(\xi)=1$ if and only if $|\xi_b|\le c_\xi\varepsilon$ at every $b\in I^\circ$ with ${\rm dist}(b,\square)\le R'$; otherwise it is $0$. The radius in force is $R'=R_\chi$. Further, $b_\square(\xi)=1$ if and only if $|\xi_b|>c_\xi\varepsilon$ at a bond of $I^\circ$ in $\hat\square$, as in item~E10 of \S\ref{Y-F-1}.\end{definition}

\emph{Item 4: the radii.} The radii form the one set of item~2 of \S\ref{Y-G0-1}; they are not letters of the class, and those from $\varrho_L$ to $r_1$ are computed from $x$ and the letters $L$, $M$ and $c_R$. They are $\varrho_L=\lceil R/L\rceil$, $r_w=\lfloor R/2\rfloor$, $R_2=2R+2$, $R_3=2R+\varrho_L+5$ and $R_v=\varrho_L+2+R_3=2R+2\varrho_L+7$; then $R_\chi^{\min}=r_w+\varrho_L+2$ and $R_\chi=R+\varrho_L+2$; and, as in item~E8 of \S\ref{Y-F-1}, $r_l=\lceil(\varrho_L+2)/M\rceil$ and $r_1=\lceil R_2/M\rceil+3r_l+5$. We also put $r:=C_\eta^\flat\varepsilon$ and $r_\tau:=C_\eta^\flat(\tau)\varepsilon$.

\emph{Item 5: the inputs.} The first group concerns the law and comes from Part C.I of \S\ref{Y-C-1}. Of \Cref{Y-C-CITh66} it contains the following clauses.

\begin{enumerate}\item[\textup{(a)}] The matrix $G$ is real at $\mathcal V'$.

\item[\textup{(b)}] The ranges, and what each object depends on.

\item[\textup{(c)}] Covariance.

\item[\textup{(d)}] The bound $\|(G\xi)_b\|\le \mathsf G_0\sup_{Y_b}\|\xi_{b'}\|$.

\item[\textup{(i)}] Holomorphy on $\mathfrak V^\circ$ and continuity on $\mathfrak V$.\end{enumerate}

This group also contains \Cref{Y-C-CITh66}\,\textup{(f)}, \textup{(h)}, \textup{(k)}: under \textup{(F$_{\rm law}$)}, and with $c_R\theta_{\mathbb C}\ge128$, the matrix $G$ is invertible, $\det G^V>0$ at $\mathcal V'$, and $\|S\|\le x^{-60}$. It contains further \Cref{Y-C-CIL64}\,\textup{(c)}, \textup{(d)}, \textup{(e)} and \Cref{Y-C-P67}\,\textup{(a)}, which is algebra and holds for any $X'$.

The second group concerns the centre and comes from Part C.II of \S\ref{Y-C-1}. It contains \Cref{Y-C-Th68}\,\textup{(a)}, \textup{(c)}: one has $|a_{0,b}|\le\sqrt2\|a_{0,b}\|\le \mathsf a_{\rm c}\hat\varepsilon$ at ${\rm dist}(b,o)\le1$ for every coarse site $o$, every $0<\hat\varepsilon\le\frac1{40}$ and every $\operatorname{SL}(2,\mathbb C)$ field on the box of radius $R_3$ about $o$ that has the three properties

\begin{enumerate}\item[\textup{(h1)}] $|V'(c)|\le1+\vartheta_0^\flat$;

\item[\textup{(h2)}] $|V'(\pi)|\le2$ on paths in the box of at most $\ell_w(x')$ bonds;

\item[\textup{(h3)}] $|\partial V'(P)-1|\le2\hat\varepsilon$.\end{enumerate}

Item~C.II~4 of \S\ref{Y-C-1} states \Cref{Y-C-Th68} for every $i$, with $i=p$ or $i=c$ and with $o=y(p)$ or $o=c_-$ respectively; since every coarse site $o$ is the $c_-$ of the bond $c=(o,1)$, the form given here, at a coarse site $o$, is that theorem taken with $i=c$. The proof of \Cref{Y-C-Th68} obtains \textup{(a)} for $a_0$ and \textup{(c)} in its Steps \textup{(0)} to \textup{(2)}, and these steps rest on \Cref{Y-C-CITh66}\,\textup{(b)}, \textup{(f)}, \textup{(k)}, on \textup{(P0)} and on \textup{(P1)}, and on no condition \textup{(Fr)}. Among its hypotheses, \Cref{Y-C-Th68} lists \textup{(Fr2)} for $\mathsf G$ and $\mathsf G_Z$, the two free kernels of item~C.II~2.2 of \S\ref{Y-C-1}; in the setting \textup{(SD)} this hypothesis holds, by \S\ref{Y-A-1}, so \Cref{Y-C-Th68} applies here as it stands.

For \textup{(d)} the second group contains \Cref{Y-C-CIICor72}\,\textup{(a)}, which is \textup{(Cen)}$(\mathsf N_{\rm c},\mathsf K_{\rm c},\mathsf a_{\rm c})$: under \textup{(Fr1)} to \textup{(Fr3)}, \textup{(P0)}, \textup{(P1)}, \textup{(P3)}, for every $p$, every $0<\hat\varepsilon\le\frac1{40}$ and every field on the box of radius $R_3$ about $y(p)$ that satisfies \textup{(e1)}, \textup{(e2)} and $|\partial V'(P)-1|\le2\hat\varepsilon$, we have $|(EX'-\mathrm{d}^Va_0)_p|\le \mathsf N_{\rm c}L^{-2}\hat\varepsilon+\mathsf K_{\rm c}\hat\varepsilon^2$.

The third group consists of \Cref{Y-E-L1C} of item 3.2 and of the formula of the chart for bounded, Borel, gauge invariant $\rho$, together with its justification, from \Cref{Y-C-L51}.

The fourth group concerns the other sections. Of \S\ref{Y-F-1}, the definitions of its items E8, E10 and E11 enter. Of \S\ref{Y-D-1}, item~E2 (the set $\mathfrak E$, and $X'(V',e)$ with the facts that $X'$ is Borel in $(V',e)$ and that $X'(V'^u,\operatorname{Ad}_ue)=\operatorname{Ad}_uX'(V',e)$) enters. Nothing of \Cref{Y-B-C1} is used by any clause of the theorem or by any step that proves one.

\emph{Item 6: the floors.} Each floor below holds for all large $x$ when the letters and $\tau$ are fixed.

\begin{enumerate}\item[\textup{(W0)}] $x\ge L^2$; $R\ge12$; $\rho_0\le\varrho(x)+1$; the floor \textup{(F$_{\rm law}$)} of \Cref{Y-C-CITh66}\,\textup{(k)}; and, from item~C.II~4 of \S\ref{Y-C-1}, the clauses $x\ge L^2$ and $R\ge6$ of \textup{(P0)}, together with \textup{(P1)}. The third clause of \textup{(P0)} in item~C.II~4 of \S\ref{Y-C-1}, which is \textup{(Tor)}, is not a floor, since it involves $N$ and, for each fixed $N$, fails once $x$ is large; it is a datum of item 2. Hence no clause of \textup{(W0)} to \textup{(W5)} depends on $N$.

\item[\textup{(W1$_\tau$)}] $22LC_\eta^\flat(\tau)\varepsilon\le1$.

\item[\textup{(W2$_\tau$)}] $17.1\,C_\eta^\flat(\tau)\,\varepsilon\,\varrho(x)\le0.1$.

\item[\textup{(W3)}] $p_0(x)\ge1$.

\item[\textup{(W4$_\tau$)}] $K_v^\flat(\tau)p_0^3x^{-1/2}\le1$.

\item[\textup{(W5)}] For \textup{(d)} alone: \textup{(P3)} of \S\ref{Y-C-1}, $\mathsf N_{\rm c}L^{-3/2}\le0.2$, and $(\mathsf K_{\rm c}L+54(C_\eta^\flat)^2)\varepsilon\le0.18$.\end{enumerate}

We write \textup{(W1)}, \textup{(W2)} and \textup{(W4)} for the cases $\tau=1$.

\emph{Item 7: the conclusions.} Assume \textup{(W0)} to \textup{(W4)}, and let $\square$ be a cell. In \textup{(a)}, \textup{(b)} and \textup{(c)}, $V'$ is any field in $\mathcal U_{\{\square\}}(x')$ and $\xi$ is any real field with $\chi_\square(\xi)=1$.

\begin{enumerate}\item[\textup{(a)}] Every series in $V_\square$ converges. As a function of $\xi$, $V_\square$ depends only on the values at the bonds of $I^\circ$ within $r_w+\varrho_L+1$ of $\square$. As a function of $V'$, it depends only on the coarse bonds with both ends in $\square_{[R_v]}$; since $R_v\le Mr'$, these are bonds of $\square^\approx$. It depends on $x$, on $\ell$ and on the letters of the law, and on no member $E^a$, $E^h$, $w$, on no $e$, and on neither $c_h$ nor $E_a$. The map $V'\mapsto V_\square(\xi,V')$ is holomorphic on $\mathcal U_{\{\square\}}(x')$.

\item[\textup{(b)}] For every $G$-valued $u$ on the sites of $\square^\approx$, extended by $1$, and with $(\operatorname{Ad}_u\xi)_b=\operatorname{Ad}_{u(y(b))}\xi_b$, one has $V_\square(\operatorname{Ad}_u\xi,V'^u)=V_\square(\xi,V')$ and $\chi_\square(\operatorname{Ad}_u\xi)=\chi_\square(\xi)$.

\item[\textup{(c)}] One has $\|\eta_b\|\le r$ at every $b$ with ${\rm dist}(b,\square)\le\varrho_L+2$, and $|X'_P|\le2.06\varepsilon'$. At every fine plaquette $p$ with ${\rm dist}(p,\square)\le\varrho_L+1$, $|F^\flat_p|\le6r$, $|F_p|\le8.25r$ and $|F_p-F^\flat_p|\le54r^2$. For $y\in\hat\square$ and $|y-p|\le\varrho(x)$, $\|\operatorname{Ad}_{U(\Gamma(y,\pi(p)))}-T(y(y),y(p))\|\le160\,r\,|y-p|$. Finally $|V_\square|\le K_v^\flat p_0^3x^{-1/2}$ and $|e^{V_\square}-1|\le2K_v^\flat p_0^3x^{-1/2}$.

\item[\textup{(d)}] Assume \textup{(W5)} as well, let $(V',e)$ be any pair on $\mathbb T'$ such that the restriction of $V'$ to the bonds of $\square^\approx$ lies in $\mathcal U_{\{\square\}}(x')$, and let $\chi_\square(\xi)=1$. This condition on $V'$ holds in particular when $\phi'_P\le\varepsilon'$ on $\square^\approx$; that case rests on $\log x'>1$, which is true because $x'\ge L\ge3$. Then $V_\square$ is real, and $|\eta_b|<\pi$ at every $b$ with ${\rm dist}(b,\square)\le\varrho_L+2$; next, $\phi_p\le\varepsilon$ at every fine plaquette $p$ with $|y-p|\le\varrho(x)+1$ for some $y\in\hat\square$; hence $\zeta_p=F_p$ at these plaquettes, $m_y(U)=1$ for $y\in\hat\square$, and $S_\varepsilon(\phi_p)=\frac12|F_p|^2$; and finally $A_\square(\xi,V',e)=e^{V_\square(\xi,V')}$ for every $e$.

\item[\textup{(e)}] Here $\xi$ is any real field and $(V',e)$ any pair; \textup{(e.i)} and \textup{(e.ii)} hold for every $\mathfrak g$-valued $X'$ put in place of $X'(V',e)$.

Of the floors, \textup{(e)}, \textup{(g)} and \textup{(h)} use the following. The letter $A_\square$, and the objects of items 3.1 to 3.7 that enter it, are defined at every real $\xi$ and every pair once $L\ge3$ is odd, $aL\ge18$, and the inequalities on $N$, namely $M+2R_v<N$, $\varrho(x)+1<NL/2-L$ and $\varrho_L+2<N/2$, hold; here \Cref{Y-C-CITh66}\,\textup{(a)} holds with no condition on the size of $R$. These inequalities on $N$ rest on $x\ge L^2$, on $R\ge6$ and on \textup{(Tor)}: the first and the third are data of item 2, and $R\ge6$ follows from $x\ge L^2$, because $c_R\ge768$ and $\log x\ge2\log3$. Beyond this, \textup{(e.i)}, \textup{(e.iv)}, \textup{(e.v)} and \textup{(g)} depend on no floor. The clauses \textup{(e.iii)} and \textup{(h)} use, of \textup{(W0)}, only $\rho_0\le\varrho(x)+1$, for the mask. The clause \textup{(e.ii)} uses only \textup{(F$_{\rm law}$)}, with $c_R\theta_{\mathbb C}\ge128$ of \textup{(O6)}. The floor \textup{(P1)} and the floors \textup{(W1)} to \textup{(W4)} enter \textup{(a)} to \textup{(d)} and \textup{(f)} alone. So \textup{(e.i)}, \textup{(e.iv)}, \textup{(e.v)} and \textup{(g)} hold in the setting of the data of item 2 alone.

\begin{enumerate}\item[\textup{(e.i)}] The product over all cells is \[\prod_\square A_\square=1_{\mathfrak D}(\eta)\prod_{b\in I^\circ}\mathsf j(\eta_b)\;e^{-x\mathrm{En}_\ell(U)}\;e^{\frac x2\langle F^\flat,\mathsf B^VF^\flat\rangle}\prod_ce^{\frac s2|n^\flat_c|^2-s(1-\cos\mathsf n_c)}.\] For every set $\mathfrak Z$ of cells, $\prod_{\square\in\mathfrak Z}A_\square=[A]^\ell_{\mathfrak Z}$, as in item~E11 of \S\ref{Y-F-1}. These identities do not depend on the law.

\item[\textup{(e.ii)}] Let $\rho=e^{\mathsf V}\rho[u;\ell]$ be a density of the class $C^{\rm su}_\flat(x;\ell)$, or any $\rho=e^{\mathsf V}\rho[u;\ell]$ with $\mathsf V\in\mathbb C$ and $u=(E^a,E^h,w)\in M_A^{\rm su}(x)$, $A\ge1$, real or complex. The signs $\rho[u;\ell]$, $\Xi_w$, $E^a$, $E^h$, $w$, $X^\approx$, $M_A^{\rm su}(x)$ and $C^{\rm su}_\flat(x;\ell)$ are those of item~1 of \S\ref{Y-G0-1}, under ``At $x$'', ``Members'' and ``The class''; the sum over $X$ below runs over the connected sets of fine cubes. The step is linear in $\rho$, while nothing is linear in $u$. Then \[\FitDisplay{\dimexpr\linewidth-1pt\relax}{\begin{aligned}
&\rho'(V')=\mathsf c_0\,z(s)^{-3\#\mathbb T'}\Bigl(\frac{2\pi}x\Bigr)^{\frac32\#I^\circ}\det G^V\;e^{\mathsf V-x\mathcal Q(-a_0)}\;
\\
&\quad{}E_\xi\Bigl[e^{-\frac x2\langle\xi,\mathsf E\xi\rangle+x\langle\xi,\mathsf Eh\rangle}
\\
&\quad{}\prod_\square A_\square\;
\\
&\quad{}\exp\Bigl\{\sum_X1_{X^\approx}(U)(E^a+E^h)(X,U)\Bigr\}\,\Xi_w(U)\Bigr],
\end{aligned}}\] where $\mathcal Q(-a_0)=\mathsf m_C-\frac12\langle a_0,Sa_0\rangle$ and $E_\xi$ is the mean over $\xi\sim\mathcal N(0,x^{-1}1)$ on $\mathfrak L_{\mathfrak g}$ of item 3.6. This formula rests on \Cref{Y-C-CITh66}\,\textup{(f)}, \textup{(h)}, \textup{(k)}.

\item[\textup{(e.iii)}] The letter $A_\square$ depends on $\xi$ only through the bonds of $I^\circ$ within $R_\chi^{\min}$ of $\square$, and on $(V',e)$ only through bonds and plaquettes in $\square_{[R_v]}$.

\item[\textup{(e.iv)}] One has $A_\square(\operatorname{Ad}_u\xi,V'^u,\operatorname{Ad}_ue)=A_\square(\xi,V',e)$ for every $u:\mathbb T'\to G$.

\item[\textup{(e.v)}] The letter $A_\square$ is a Borel function of $(\xi,V',e)$.\end{enumerate}

\item[\textup{(f)}] Let $\tau\ge1$, and let $B$ be a set of coarse sites with $B\supset\square_{[R_v]}$. Suppose that \textup{(W0)}, \textup{(W1$_\tau$)}, \textup{(W2$_\tau$)}, \textup{(W3)} and \textup{(W4$_\tau$)} hold. These conditions by themselves give $\tau\varepsilon'\le1/40$, and the clauses \textup{(a)}, \textup{(b)} and \textup{(c)} hold for every $V'\in U^\tau(B;x')$, where $B$ takes the place of $\square^\approx$, $r_\tau$ that of $r$, $\tau\varepsilon'$ that of $\varepsilon'$ in the bound on $X'$, and $K_v^\flat(\tau)$ that of $K_v^\flat$.

\item[\textup{(g)}] The functions $V_\square$ and $A_\square$ are built from the cell $\square$ (and thus from $M$ and the grid of cells) on the torus (and thus from $N$), from $x$, $\ell$, $L$, $a$ and $R(x)$ (and thus from $c_R$), and from the number $b$ of \textup{(O6)} (and thus from $K_*$ and $\lambda_*$). In $A_\square$ there occur, in addition, $d_0$, $\rho_0$, $A_0$ and $\mathfrak p_0$. Neither function is built from anything else. Each of these data is a letter of the table or a datum of \textup{(SD)}, and none of them is a member, $c_h$ or $E_a$. Consequently, any two densities belonging to a single class $C^{\rm su}_\flat(x;\ell)$ on a single torus share their vertex letters. This clause is independent of the reality and of the positivity of the members.

\item[\textup{(h)}] The radii obey the inequalities \[R_\chi^{\min}\le R_\chi,\qquad 2\lceil R_\chi/M\rceil\le r_1-5,\qquad \lceil R_v/M\rceil\le r_1-5.\] Consequently, if $b_{\square'}(\xi)=0$ for every cell $\square'$ of $\square^{+r_1}$, then $\chi_\square(\xi)=1$. Moreover, $A_\square$ depends on $\xi$ only under $\square^{+r_1}$ and on $(V',e)$ only in $\square^{+r_1}$. Finally, if two cells $\square$ and $\tilde\square$ are more than $2\lceil R_\chi^{\min}/M\rceil$ cells apart, then the letters $A_\square$ and $A_{\tilde\square}$ depend on disjoint sets of the variables $\xi$.
\end{enumerate}

\emph{Item 8: what is used.} The proofs of \textup{(a)}, \textup{(b)}, \textup{(c)} and \textup{(f)} rest on the following results: \Cref{Y-C-CITh66}\,(a) to~(d), (f), (i) and~(k); \Cref{Y-C-CIL64}\,(c), (d) and~(e); \Cref{Y-C-Th68}\,(c); \Cref{Y-E-L1C}.

The proof of \textup{(d)} rests, in addition, on \Cref{Y-C-CIICor72}\,(a), and hence, for the pair $(\ell,\ell')$ of the data, on \textup{(Fr1)} to \textup{(Fr3)} of item~C.II~2.2 of \S\ref{Y-C-1}, which hold in the setting \textup{(SD)} of item~2 of this section.
\end{theorem}

\begin{proof}[Proof of \Cref{Y-E-C6}]
\emph{Step 1: the letters as factors of the integrand; clauses \textup{(e.i)} and \textup{(e.ii)}.}\label{Y-E-LA} Let $\xi$ be real and $V'\in\mathcal V'$, and let $X'$ be any $\mathfrak g$-valued function on the coarse plaquettes, for instance $X'(V',e)$. We claim three things. First, \[\prod_\square A_\square=1_{\mathfrak D}(\eta)\prod_{b\in I^\circ}\mathsf j(\eta_b)\,e^{-x\mathrm{En}_\ell(U)}\,e^{(x/2)\langle F^\flat,\mathsf B^VF^\flat\rangle}\prod_ce^{(s/2)|n^\flat_c|^2-s(1-\cos\mathsf n_c)}.\] Second, $\prod_{\square\in\mathfrak Z}A_\square=[A]^\ell_{\mathfrak Z}$. Third, if \textup{(F$_{\rm law}$)} holds, then the display of \textup{(e.ii)} holds.

\emph{\textup{(1)} Splitting.} For real $V'$ the transports satisfy $T^{-1}=T^{\mathsf T}$, so $(\ell^V)^*$ is the adjoint of $\ell^V$, and therefore \[\langle F^\flat,\mathsf B^VF^\flat\rangle=\sum_p|F^\flat_p|^2+\sum_y|(\ell^VF^\flat)_y|^2.\] By item~3.5, $\mathrm{En}_\ell(U)=\sum_pS_\varepsilon(\phi_p)+\frac12\sum_ym_y(U)|\ell_y^U\zeta|^2$; moreover $1_{\mathfrak D}(\eta)=\prod_b1\{|\eta_b|<\pi\}$. Each fine plaquette, each fine site and each bond of $I^\circ$ is based in exactly one $\hat\square$, and for each coarse bond $c$ the site $c_-$ lies in exactly one $\square$. So the right side of the first claim is the product, over the cells, of the four products of item~3.9, which proves the first claim; no estimate enters here, and this part is independent of $G$ and of $X'$. Put $\hat{\mathfrak Z}:=\bigcup_{\square\in\mathfrak Z}\hat\square$. Taken over the set $\hat{\mathfrak Z}$, the four products agree, letter by letter, with those of item~E11 of \S\ref{Y-F-1}, because $ax=s$; this proves the second claim.

\emph{\textup{(2)} Algebra.} This part is \Cref{Y-C-P67}\,(a), which is stated for every $X'$; we give its derivation. We have $\mathcal Q(\eta)=\frac12EX'\cdot\mathsf BEX'+\mathsf g^VX'\cdot\eta+\frac12\eta\cdot q\eta$, where $q$ and $\mathsf B$ are symmetric. With $\eta=G(\xi-h)$ and $h=G^{\mathsf T}\mathsf g^VX'$ we get $\mathsf g^VX'\cdot\eta=h\cdot(\xi-h)$ and $\eta\cdot q\eta=(\xi-h)\cdot(1+\mathsf E)(\xi-h)$, and hence \[\mathcal Q(\eta)=[\frac12EX'\cdot\mathsf BEX'-\frac12h\cdot h+\frac12h\cdot\mathsf Eh]+\frac12\xi\cdot\xi+\frac12\xi\cdot\mathsf E\xi-\xi\cdot\mathsf Eh.\] Since $h\cdot h=\mathsf g^VX'\cdot C\mathsf g^VX'$, the bracket equals $\mathcal Q(-a_0)=\mathsf m_C+\frac12h\cdot\mathsf Eh$. With $S=-G^{-\mathsf T}\mathsf EG^{-1}$, as given by \Cref{Y-C-CITh66}\,(f), we obtain $h\cdot\mathsf Eh=-a_0\cdot Sa_0$.

\emph{\textup{(3)} Change of variables.} The formula of the chart, \Cref{Y-C-L51}, states that for every bounded, Borel, gauge invariant $\rho$ \[\rho'(V')=\mathsf c_0\int_{\mathfrak D}\rho(\mathsf U(\eta;V'))\prod_ck_s(V'(c),\mathsf M(U)(c))\prod_{b\in I^\circ}\mathsf j(\eta_b)\,d\eta.\] As there are $3\#\mathbb T'$ coarse bonds, the kernels contribute the factor $z(s)^{-3\#\mathbb T'}\prod_ce^{-s(1-\cos\mathsf n_c)}$. We multiply and divide the integrand by $e^{-x\mathcal Q(\eta)}$, where $x\mathcal Q(\eta)=(x/2)\langle F^\flat,\mathsf B^VF^\flat\rangle+(s/2)\sum_c|n^\flat_c|^2$. If \textup{(F$_{\rm law}$)} holds, then, as $c_R\theta_{\mathbb C}\ge128$, \Cref{Y-C-CITh66}\,(k) gives $4b^{1/2}d_R\le1$, and by \Cref{Y-C-CITh66}\,(a), (f) and~(h) the matrix $G^V$ is real and invertible with $\det G^V>0$. Therefore $\eta=G\xi-a_0$ is a bijection of $\mathfrak L_{\mathfrak g}$, with $d\eta=\det G^V\,d\xi$, and $a_0$ does not depend on $\xi$. We also have $\int e^{-(x/2)|\xi|^2}d\xi=(2\pi/x)^{(3/2)\#I^\circ}$. Now, in the exponent $-x\mathcal Q(\eta)$, we replace $\mathcal Q(\eta)$ by the expression for it found in part~\textup{(2)}, and for $\rho(U)$ we put the formula of the class from item~1 of \S\ref{Y-G0-1}, where the members are defined, \[\rho[u;\ell]=\exp\{-x\mathrm{En}_\ell+\sum_{X\in\mathcal D}1_{X^\approx}[E^a+E^h](X,\cdot)\}\Xi_w;\] using in addition the splitting of part~\textup{(1)}, we arrive at the display of \textup{(e.ii)}.

It remains to see that $\rho=e^{\mathsf V}\rho[u;\ell]$, with $u\in M_A^{\rm su}(x)$ real or complex, is admissible in this formula: it is gauge invariant, bounded and Borel on the $G$-valued fine fields of a finite torus. As to invariance, $\phi_p$, $D(U)$, $m_y(U)$, $1_{X^\approx}$, $S_\varepsilon(\phi_p)$ and $|\ell_y^U\zeta|$ are invariant, and so is $\mathrm{En}_\ell$. Indeed, $\partial U^u(p)=u(\pi(p))\partial U(p)u(\pi(p))^{-1}$, so that $\zeta_p\to\operatorname{Ad}_{u(\pi(p))}\zeta_p$, and this holds also at $\partial U(p)=-1$, where both sides are $0$; and the transports are conjugated. The functions $E^a$, $E^h$ and $w$ are invariant by their clauses \textup{(ca$_A$)}, \textup{(ch$_A$)} and \textup{(c$\Xi_A$)}.

As to boundedness, $\mathrm{En}_\ell\ge0$ gives $e^{-x\mathrm{En}_\ell}\le1$, and there are only finitely many sets $X$. A $G$-valued field that is tame on $X^\approx$ lies in $\mathcal U_X(x)$ (item~1 of \S\ref{Y-G0-1}, where the domains are defined): we have $|U(b)|=1<1+\vartheta(x)$ and $|U(\pi)|=1<\mathsf b(x)$, both because $\log x>1$, which holds since $x\ge L^2\ge9$; and $|\partial U(p)-1|=2\sin(\phi_p/2)\le\varepsilon<2\varepsilon$. Hence $|1_{X^\approx}(E^a+E^h)(X,\cdot)|\le A(c_a(x)+c_h(x))$ and $|w|\le A$. As to measurability, $\phi_p$ is continuous; $\mathsf L$, the clip, the masks and the indicators are Borel; $E^a$ is continuous on $\mathcal U_X(x)$; $E^h$ and $w$ are Borel by their clauses; and all sums and products are finite.

The left side does not depend on $e$, so the display holds for every $e$. The display is linear in $\rho$; applying it to the real and to the imaginary part of $\rho$, we find that it holds for complex bounded Borel invariant $\rho[u;\ell]$.

\emph{\textup{(4)} The ratio factor through the $S_k(Y)$.} Put $\varphi=G\xi$ and $T_Y=(x/2)\langle\varphi,S_k(Y)\varphi\rangle-x\langle\varphi,S_k(Y)a_0\rangle$, where $S=\sum_YS_k(Y)$ by \Cref{Y-C-CIITh74}\,(a). Here \Cref{Y-C-CIITh74}\,(a) is taken under its own hypotheses, among them $\kappa\ge6$, \textup{(M4$^\flat$)} and its floor \textup{(S2)}; Thms.~\ref{Y-G-ThA} and~\ref{Y-H-C9} suppose these, \Cref{Y-E-C6} does not, and no clause of \Cref{Y-E-C6} depends on part~\textup{(4)}. Then \[\sum_YT_Y=(x/2)\eta\cdot S\eta-(x/2)a_0\cdot Sa_0=-(x/2)\xi\cdot\mathsf E\xi+x\,\xi\cdot\mathsf Eh,\] where we have used $\eta\cdot S\eta=-(\xi-h)\cdot\mathsf E(\xi-h)$. Hence the ratio factor defined with the $T_Y$ is that of \textup{(e.ii)}.

\emph{Step 2: on what the letters depend, and $R_v\le Mr'$; geometry only, with no field.}\label{Y-E-LB} We claim the following.
\begin{enumerate}
\item[\textup{(i)}] $V_\square$ is a function of three things: the $\eta_b$ with $b\in I^\circ$ and ${\rm dist}(b,\square)\le\varrho_L+1$; the values of $V'$ on bonds with both ends in $\square_{[\varrho_L+1]}$; and the $X'_P$ with $P$ in $\square_{[\varrho_L+1]}$. For $A_\square$, and for the angles $\phi_p$ of \textup{(d)}, the same holds with $\varrho_L+2$ in place of $\varrho_L+1$.

\item[\textup{(ii)}] Let $o$ be a coarse site and let ${\rm dist}(b,o)\le1$. Then $(G\xi)_b$ depends on $\xi$ on $Y_b$; and the row $G(b,\cdot)$ and $a_{0,b}$ depend only on the field on those bonds, and on $X'$ on those plaquettes, all of whose sites lie within $R_3$ of $o$.

\item[\textup{(iii)}] Hence two fields of $\mathfrak V$ which, together with their $X'$, agree on the box of radius $R_3$ about $o$ have the same $G(b,\cdot)$ and the same $a_{0,b}$. Statements \textup{(ii)} and \textup{(iii)} hold at every field of $\mathfrak V$ and for every $X'$ on the coarse plaquettes, and therefore at every pair $(V',e)$: their proof uses the ranges alone and is independent of the size of the curvature.

\item[\textup{(iv)}] $R_v\le20c_R\log x'\le Mr'$ and $3\varrho_L\le\ell_w(x')$; the box $\square_{[R_v]}$ has side $M+2R_v<N$, so that this box and every box inside it embed in $\mathbb T'$; and $\varrho(x)+1<NL/2-L$.
\end{enumerate}

\emph{Proof of \textup{(i)}.} In the first and in the fourth of the four sums of item~3.8, $p$ and $b$ lie in $\hat\square$. The four corners of a fine plaquette $p$ lie in blocks within $1$ of $y(p)$, so that its bonds have ${\rm dist}(b,y(p))\le1$. The letters $V'(c)$ in $\partial U(p)$, and the transports along one bond in $(\mathrm{d}^V\eta)_p$, have both ends within $1$ of $y(p)$. An edge plaquette $p$ lies under a $P$ with $y_P=y(p)$, and the sites of $P$ are within $1$ of $y(p)$. In the third sum, the bonds of $[x,x']$ with $x\in B(c_-)$ are based in $B(c_-)$ or in $B(c_+)$, and $V'(c)$ enters; this gives distance at most $1$. In the second sum, $y\in\hat\square$ and $|y-p|$ is at most the range of the rows, which is at most $\varrho(x)=R$; so $L^{-1}|y(y)-y(p)|_\infty<R/L+1$, and since the left side is an integer, it is at most $\varrho_L$ (\Cref{Y-C-CIL64}\,(c)). The bonds of the fine path $\Gamma(y,\pi(p))$ are based in the box spanned by $y$ and $\pi(p)$, and are therefore at distance at most $\varrho_L$. The path $\Gamma(y(y),y(p))$ stays in the box of its ends, inside $\square_{[\varrho_L]}$, and the plaquette $p$ adds $1$. By \textup{(W0)} we have $\rho_0\le\varrho(x)+1$, so the mask $m_y(U)$ uses the plaquettes $p$ with $|y-p|\le\max\{\rho_0,\varrho(x)+1\}=R+1$; clause \textup{(d)} concerns those with $|y-p|\le\varrho(x)+1=R+1$. From $|y-p|\le R+1$ we get $L^{-1}|y(y)-y(p)|_\infty<(R+1)/L+1\le\varrho_L+1+1/L$; as the left side is an integer, it is at most $\varrho_L+1$, and once more the plaquette adds $1$.

\emph{Proof of \textup{(ii)}.} By \Cref{Y-C-CITh66}\,(b), $(G\xi)_b$ depends on $\xi$ on $Y_b$; $G(b,\cdot)$ depends on the bonds with both ends within $r_w+\varrho_L+1$ of $y(b)$; $C(b,b')=0$ unless ${\rm dist}(b,b')\le2r_w$; and $C(b,\cdot)$ depends on bonds within $3r_w+\varrho_L+1$ of $y(b)$. By \Cref{Y-C-CIL64}\,(e), $(\mathsf g^VX')_{b'}$ depends on paths that end at $y_P$ and all of whose sites lie within $2\varrho_L+1$ of $y(b')$, and on $X'_P$. The sites of $P$ are then within $2\varrho_L+2$ of $y(b')$, hence within $2r_w+2\varrho_L+2$ of $y(b)$. We add $1$ for ${\rm dist}(b,o)\le1$. For the rows of $G$ and $C$ this gives $3r_w+\varrho_L+2\le\frac32R+\varrho_L+2\le R_3$; for $X'$ it gives $2r_w+2\varrho_L+3\le R+2\varrho_L+3\le2R+\varrho_L+5=R_3$, since $\varrho_L\le R$. Here the $\varrho_L$ of Part C.I of \S\ref{Y-C-1} (item~1.2 there) is defined with the range, which is at most $R$, so it is not larger than the $\varrho_L$ of item~4 of this section.

\emph{Proof of \textup{(iii)}.} Statement \textup{(iii)} follows from \textup{(ii)}, applied to each of the two fields.

\emph{Proof of \textup{(iv)}.} We have $\varrho_L\le R/3+1$ and $R\le c_R\log x+1$, and $\log x\le2\log x'$ because $x\ge L^2$. Therefore \[R_v=2R+2\varrho_L+7\le\frac83R+9\le\frac{16}3c_R\log x'+11.67\le20c_R\log x',\] where the last inequality holds because $\frac{44}3c_R\log x'\ge\frac{44}3\cdot768\cdot\log3>11.67$. Next, $Mr'\ge20c_R\log x'$, and so $3\varrho_L\le3R_v\le3Mr'\le\ell_w(x')$. By \textup{(Tor)}, $N\ge16Mr(x)\ge320c_R\log x$. So the side $M+2R_v\le M(2r'+1)$ of $\square_{[R_v]}$ is less than $N$, and $NL/2-L\ge L(160c_R\log x-1)>R+1$. Also $N/2\ge160c_R\log x\ge160(R-1)>R/3+3\ge\varrho_L+2\ge R/L+2$, since $R\ge6$ by \textup{(P0)} and $L\ge3$. Rows of range at most $\varrho(x)$ are thus rows as in item~C.I~1.1 of \S\ref{Y-C-1}, where $\varrho/L+2<N/2$ is required, and rows as in item~C.II~1.2 of \S\ref{Y-C-1}, where $\varrho_L+2<N/2$ is required. This is the hypothesis on $N$ under which we apply, in this section, \Cref{Y-C-CIL64}, \Cref{Y-C-CITh66}, \Cref{Y-C-P67}, \Cref{Y-C-Th68} and \Cref{Y-C-CIICor72}.

By \textup{(i)} and \textup{(ii)}, $V_\square$ and $A_\square$ depend on the field and on $X'$ only through the bonds and plaquettes in $\square_{[\varrho_L+2+R_3]}=\square_{[R_v]}$. Only the $\xi_{b'}$ with $b'\in I^\circ$ and ${\rm dist}(b',\square)\le r_w+\varrho_L+1$ occur in $V_\square$, and only those with ${\rm dist}(b',\square)\le r_w+\varrho_L+2=R_\chi^{\min}$ occur in $A_\square$ and in the angles $\phi_p$ of \textup{(d)}.

\emph{Step 3: the size of $\eta$ and of $X'$.} From this step up to Step~10 we keep the following convention: $\tau\ge1$, $B\supset\square_{[R_v]}$, $V'\in U^\tau(B;x')$ and $\chi_\square(\xi)=1$; we put $\hat\varepsilon:=\tau\varepsilon'$; and $r$ stands for $r_\tau$ and $C_\eta$ for $C_\eta^\flat(\tau)$. In the case $\tau=1$, with $B$ the set of sites of $\square^\approx$, the set $U^1(B;x')$ contains $\mathcal U_{\{\square\}}(x')$. We claim: $\hat\varepsilon\le1/462$; $\operatorname{cont}(V')\in\mathfrak V^\circ$; $|V'(\pi)|<2$ on paths in $B$ of at most $\ell_w(x')$ bonds; $|X'_P|\le2.06\hat\varepsilon$ for $P$ in $B$; and $\|\eta_b\|\le r$ at every $b\in I^\circ$ with ${\rm dist}(b,\square)\le\varrho_L+2$.

By \textup{(W1$_\tau$)}, together with $C_\eta\ge7\tau L^{1/2}$, we have $154\tau L^{3/2}\varepsilon\le1$, and hence $\hat\varepsilon\le\tau L^{1/2}\varepsilon\le1/(154L)\le1/462$.

\emph{Room.} We argue as in item~C.II~1.4 of \S\ref{Y-C-1} and give the derivation. Since $\log L\le\log x$, we have $\mathsf g_x\le\vartheta_0^\flat/(4\log x)$; also $\log x'\le\log x$. So $\vartheta_0^\flat-\vartheta(x')-\mathsf g_x\ge\vartheta_0^\flat/\log x-\vartheta_0^\flat/(4\log x)>0$ and $2-\mathsf b(x')-\mathsf g_x\ge1/\log x-1/(4\log x)>0$. Hence \textup{(e1)} gives $|V'(c)|<1+\vartheta_0^\flat$, and as the continuation has norm $1$, $\operatorname{cont}(V')\in\mathfrak V^\circ$; and \textup{(e2)} gives $|V'(\pi)|<2$.

\emph{The field $X'$.} Let $t=|\partial V'(P)-1|$. Then $t<2\hat\varepsilon\le0.05$, and \Cref{Y-E-L1C}\,(b) gives $|X'_P|\le t+t^2/(2(1-t))\le1.0264t\le2.06\hat\varepsilon$ and $e^{iX'_P}=\partial V'(P)$.

\emph{Noise.} Let ${\rm dist}(b,\square)\le\varrho_L+2$. The bonds of $Y_b$ lie within $r_w+\varrho_L+2=R_\chi^{\min}\le R'$ of $\square$, so that $|\xi_{b'}|\le c_\xi\varepsilon$ there. By \Cref{Y-C-CITh66}\,(d), $\|(G\xi)_b\|\le \mathsf G_0\sup_{Y_b}\|\xi_{b'}\|\le C_\varphi^\flat c_\xi\varepsilon=\varepsilon/6.5$.

\emph{Centre.} Let $o=y(b)$. The box of radius $R_3$ about $o$ lies in $\square_{[R_v]}$, and so in $B$. The restriction of $V'$ to this box satisfies the hypotheses of \Cref{Y-C-Th68} on the field, \textup{(h1)}, \textup{(h2)} and \textup{(h3)}. Indeed, \textup{(h1)} follows from \textup{(e1)} and the bound $|V'(c)|<1+\vartheta_0^\flat$ of the paragraph ``Room'', and \textup{(h2)} from \textup{(e2)} and the bound $|V'(\pi)|<2$ on paths in $B$ of at most $\ell_w(x')$ bonds, which that paragraph also gives; and \textup{(h3)}, with $\hat\varepsilon=\tau\varepsilon'$, is \textup{(e3)}. The parameter $\hat\varepsilon$ lies in the range $0<\hat\varepsilon\le\frac1{40}$ of that theorem, because $\hat\varepsilon\le1/462$. By Step~2\,(iii), the $a_{0,b}$ evaluated at $\operatorname{cont}(V')$ with the $X'$ of item~3.7 equals the $a_{0,b}$ evaluated at the restriction extended by $1$, with $X'=0$ off the box. We apply \Cref{Y-C-Th68}\,(c). Its hypotheses \textup{(P0)} and \textup{(P1)} are contained in \textup{(W0)}, with one exception: the third clause of \textup{(P0)} is \textup{(Tor)}, which is a datum of item~2 and is not contained in \textup{(W0)}. Those on the law, \textup{($\lambda$0)}, \textup{($\rho$1)}, \textup{($\rho$2)} and \textup{($\rho$5)}, hold for the law of windows by item~C.II~2.1 of \S\ref{Y-C-1}, which derives them from \Cref{Y-C-CITh66}\,(a), (b), (f), (k) under \textup{(F$_{\rm law}$)} of \textup{(W0)} and $c_R\theta_{\mathbb C}\ge128$ of \textup{(O6)}, namely \textup{($\lambda$0)} from~(a), \textup{($\rho$1)} and \textup{($\rho$2)} from~(b), and \textup{($\rho$5)} from~(f) and~(k). The hypothesis \textup{(Fr2)} of \Cref{Y-C-Th68} holds in the setting \textup{(SD)}. \Cref{Y-C-Th68}\,(c) therefore gives $\|a_{0,b}\|\le \mathsf a_{\rm c}\hat\varepsilon/\sqrt2\le \mathsf a_{\rm c}\tau L^{1/2}\varepsilon$.

\emph{Sum.} Hence $\|\eta_b\|\le(1/6.5+\mathsf a_{\rm c}\tau L^{1/2})\varepsilon\le r$.

\emph{Step 4: four or five exponentials; the first sum.}\label{Y-E-LC} We claim that at every fine plaquette $p$ with ${\rm dist}(p,\square)\le\varrho_L+1$ the following holds: $\partial U(p)=\prod_ke^{iY_k}$ with $4$ or $5$ factors, where $\sum Y_k=F^\flat_p$ and $\sum|Y_k|\le S:=6r\le1/11$; further, $|F^\flat_p|\le6r$, $|F_p-F^\flat_p|\le54r^2$ and $|F_p|\le8.25r$; and finally $(x/2)|F^\flat_p\cdot F^\flat_p-F_p\cdot F_p|\le384.75\,xr^3$.

The argument rests on the identity of matrices \[\partial U(p)=e^{i\eta_1}e^{iT_2\eta_2}[\partial V'(P)]e^{-iT_3\eta_3}e^{-i\eta_4},\qquad T_j=T(y(p),y(b_j))\in\{1,\operatorname{Ad}_{V'(c_1)},\operatorname{Ad}_{V'(c_4)}\},\] in which the bracket is present exactly when $p$ is an edge plaquette under $P=(y(p);\mu\nu)$; in that case $\partial V'(P)=e^{iX'_P}$ by Step~3.

To prove the identity we fix notation: $p=(z;\mu\nu)$, $b_1=(z,\mu)$, $b_2=(z+e_\mu,\nu)$, $b_3=(z+e_\nu,\mu)$ and $b_4=(z,\nu)$, so that $\partial U(p)=U_{b_1}U_{b_2}U_{b_3}^{-1}U_{b_4}^{-1}$; further $y=y(z)$, $c_1=(y,\mu)$, $c_4=(y,\nu)$, $c_2=(y+Le_\mu,\nu)$ and $c_3=(y+Le_\nu,\mu)$; and $\eta_j=\eta_{b_j}$, where we put $0$ on a comb bond. Let $\alpha=1$ if $z+e_\mu\notin B(y)$ and $\alpha=0$ otherwise, and let $\beta=1$ if $z+e_\nu\notin B(y)$ and $\beta=0$ otherwise. If $\alpha=\beta=0$, there are four exponentials. If $\alpha=1$ and $\beta=0$, the bonds $b_1$ and $b_3$ cross $c_1$ and $b_2$ lies in $B(y+Le_\mu)$, so that \[\FitDisplay{\dimexpr\linewidth-1pt\relax}{\partial U(p)=e^{i\eta_1}V'(c_1)e^{i\eta_2}V'(c_1)^{-1}e^{-i\eta_3}e^{-i\eta_4}.}\] If $\alpha=0$ and $\beta=1$, the bonds $b_2$ and $b_4$ cross $c_4$ and $b_3$ lies in $B(y+Le_\nu)$, so that \[\FitDisplay{\dimexpr\linewidth-1pt\relax}{\partial U(p)=e^{i\eta_1}e^{i\eta_2}V'(c_4)e^{-i\eta_3}V'(c_4)^{-1}e^{-i\eta_4}.}\] If $\alpha=\beta=1$, which is the case of the $L$ edge plaquettes under $P=(y;\mu\nu)$, then \[\FitDisplay{\dimexpr\linewidth-1pt\relax}{\partial U(p)=e^{i\eta_1}V'(c_1)e^{i\eta_2}V'(c_2)V'(c_3)^{-1}e^{-i\eta_3}V'(c_4)^{-1}e^{-i\eta_4}=e^{i\eta_1}e^{i\operatorname{Ad}_{V'(c_1)}\eta_2}\,\partial V'(P)\,e^{-i\operatorname{Ad}_{V'(c_4)}\eta_3}e^{-i\eta_4}.}\] Each time we use $ge^{iY}g^{-1}=e^{i\operatorname{Ad}_gY}$, and $\Gamma(y,y(b_j))$ is empty or is the one bond named. The exponents add up to $(\mathrm{d}^V\eta)_p+(EX')_p=F^\flat_p$.

The bonds of $p$ have ${\rm dist}(b,\square)\le\varrho_L+2$, so by Step~3 we have $|\eta_b|\le\sqrt2\|\eta_b\|\le\sqrt2r$ and $|\operatorname{Ad}_{V'(c)}\eta_b|\le|V'(c)|^2\sqrt2r\le(1+10^{-4})^2\sqrt2r\le1.42r$. Hence $\sum|Y_k|\le(2\sqrt2+2.84)r+2.06\hat\varepsilon\le5.6685r+(2.06/7)r\le6r$, because $r\ge7\tau L^{1/2}\varepsilon\ge7\hat\varepsilon$. By \textup{(W1$_\tau$)}, $r\le1/(22L)\le1/66$, so $S\le1/11\le1/4$. Then $|\partial U(p)-1|\le e^S-1<1$, and the series of $F_p$ converges. \Cref{Y-E-L1C}\,(c) gives $|F_p-F^\flat_p|\le\frac32S^2\le54r^2$, and hence $|F_p|\le S+\frac32S^2\le1.375S\le8.25r$. Finally, $F\cdot F-F^\flat\cdot F^\flat=(F-F^\flat)\cdot(F+F^\flat)$ and $|Y\cdot Y'|\le\|Y\|\|Y'\|\le|Y||Y'|$, so the difference is at most $\frac32S^2\cdot2.375S=3.5625S^3$; multiplied by $x/2$, this is $1.78125\cdot216\,xr^3=384.75\,xr^3$.

\emph{Step 5: the fine transport along a coordinate path compared with the coarse transport, at complex fields.}\label{Y-E-LD} Let $y\in\hat\square$, let $p$ be a fine plaquette with $n:=|y-p|\le\varrho(x)$, and put $z=y$ and $z'=\pi(p)$. We claim that $U(\Gamma(z,z'))=\Pi\,V'(\Gamma(y(z),y(z')))$ with $\Pi\in\operatorname{SL}(2,\mathbb C)$, $|\Pi-1|\le18rn$ and $|\Pi|\le1.106$, and that \[\|\operatorname{Ad}_{U(\Gamma(z,z'))}-T(y(z),y(z'))\|\le160rn,\qquad\|\operatorname{Ad}_{U(\Gamma(z,z'))}\|\le4.9,\qquad\|T(y(z),y(z'))\|<4.\]

\emph{\textup{(1)} Geometry, without the field.} Only seams of the direction $\mu$ are crossed by a fine stretch along $\pm e_\mu$, and they are crossed one after the other. The stretch along $e_1$ goes from the block of $y(z)$ to the block whose first coordinate is $y(z')_1$, and so on for the other directions. In each direction the coarse displacement is less than $n+L<NL/2$ (Step~2\,(iv)), and it has the sign of the fine displacement or is $0$; so the crossings take the short way round $\mathbb T'$, as $\Gamma(y(z),y(z'))$ does by its definition. Hence the path of the crossings of $\Gamma(z,z')$ is $\pi':=\Gamma(y(z),y(z'))$.

\emph{\textup{(2)} An identity.} A factor of $U(\Gamma)$ is $1$, $e^{\pm i\eta_b}$, $e^{i\eta_b}V'(c)$ or $V'(c)^{-1}e^{-i\eta_b}$. We move every exponential to the left past the letters $V'^{\pm1}$ and obtain $U(\Gamma)=[\prod_ig_ie^{\pm i\eta_{b_i}}g_i^{-1}]V'(\pi')$ with $g_i=V'(\pi'_{j(i)})$; here $\pi'_j$ consists of the first $j$ bonds of $\pi'$, and $j(i)$ is the number of letters to the left of the factor in question.

\emph{\textup{(3)} The factors.} The path $\pi'_j$ lies in $\square_{[\varrho_L]}$, hence in $B$, and it has at most $3\varrho_L\le\ell_w(x')$ bonds (Step~2). So $|g_i|=|g_i^{-1}|<2$ (Step~3 and \Cref{Y-E-L1C}\,(a)). The bonds $b$ of $\Gamma$ have ${\rm dist}(b,\square)\le\varrho_L$, so $|\eta_b|\le\sqrt2r$. By \textup{(W2$_\tau$)}, together with $\varrho(x)\ge1$, we have $\sqrt2r\le\sqrt2\cdot0.1/17.1\le0.0083\le0.01$. So each bracketed factor is $1+\delta_i$ with $|\delta_i|\le4(e^{\sqrt2r}-1)\le4\cdot1.00502\cdot1.41422\,r\le5.7r$.

\emph{\textup{(4)} The product.} Since $\Gamma$ has at most $3n$ bonds, $|\Pi-1|\le(1+5.7r)^{3n}-1\le e^{17.1rn}-1$. By \textup{(W2$_\tau$)}, $17.1rn\le0.1$, and then the right side is at most $1.05171\cdot17.1rn\le18rn$. Also $|\Pi|\le e^{0.1}\le1.106$; and $\det\Pi=1$, because the $\eta_b$ are traceless, the blocks of $G$ acting on $\mathfrak g_{\mathbb C}$.

\emph{\textup{(5)} The adjoint action.} \Cref{Y-E-L1C}\,(a) gives $\|\operatorname{Ad}_\Pi-1\|\le2|\Pi||\Pi-1|\le2\cdot1.106\cdot18rn\le39.82rn$, and $\|T\|\le|V'(\pi')|^2<4$. As $\operatorname{Ad}_{U(\Gamma)}=\operatorname{Ad}_\Pi T$, the difference is at most $4\cdot39.82rn\le159.3rn\le160rn$, and the norm is at most $e^{0.2}\cdot4\le4.9$.

\emph{Step 6: the second sum, over $y\in\hat\square$, with the rows.}\label{Y-E-LE} We claim that for every $y\in\hat\square$ \[(x/2)|(\ell^VF^\flat)_y\cdot(\ell^VF^\flat)_y-(\ell_y^UF)\cdot(\ell_y^UF)|\le K_\ell\,xr^3,\] where $K_\ell$ does not depend on $\varrho(x)$.

On $\mathsf H_y\otimes \mathfrak g_{\mathbb C}$ we have $|A\cdot A'|\le\|A\|\|A'\|$, by the Cauchy--Schwarz inequality in $\mathbb C^3$ and in $i$, and $\|\sum_p\ell_y(p)\otimes Y_p\|\le\sum_p|\ell_y(p)|\|Y_p\|$. Put $A=\ell_y^UF$, $B=(\ell^VF^\flat)_y$ and $n_p=|y-p|\le\varrho(x)$. In the chart $U=1$ on the combs, so the frame of a fine site is that of its block; by Step~4, $F^\flat_p$ is the sum of the exponents of $\partial U(p)$ in the frame of $\pi(p)$. Therefore \[A-B=\sum_p\ell_y(p)\otimes\{[\operatorname{Ad}_{U(\Gamma(y,\pi(p)))}-T(y(y),y(p))]F_p+T(y(y),y(p))(F_p-F^\flat_p)\}.\] By Steps~4 and~5, and since $\|\cdot\|\le|\cdot|$, we get $\|A-B\|\le\sum_p|\ell_y(p)|[160rn_p\cdot8.25r+4\cdot54r^2]$. There are three plaquettes to a site, and the range embeds in $\mathbb Z^3$ (Step~2\,(iv)); so $\sum_p|\ell_y(p)|n_p\le3K_*\mathsf g_1$ and $\sum_p|\ell_y(p)|\le3K_*\mathsf g_*$, and hence $\|A-B\|\le3K_*[1320\mathsf g_1+216\mathsf g_*]r^2$. Further, $\|A\|+\|B\|\le\sum_p|\ell_y(p)|(4.9\cdot8.25+4\cdot6)r\le64.425r\cdot3K_*\mathsf g_*\le194K_*\mathsf g_*r$. Now $A\cdot A-B\cdot B=(A-B)\cdot(A+B)$, and $(x/2)\cdot3\cdot194=291x$; this gives the bound $291K_*^2\mathsf g_*[1320\mathsf g_1+216\mathsf g_*]\,xr^3$. The decay $e^{-\lambda_*n_p}$ of $|\ell_y(p)|$ outweighs the weight $n_p$ from Step~5, and $\varrho(x)$ does not occur in the inequality $\sum_p|\ell_y(p)|n_p\le3K_*\mathsf g_1$; so no logarithm appears.

\emph{Step 7: the kernel and the Jacobian.}\label{Y-E-LF} We claim: $\varsigma:=1.42Lr\le1/15$; for $c\in\square$ the series of $\tau_{c,x}$, of $\mathsf M(U)(c)$ and of $n_c:=L_{\mathbb C}(V'(c)\mathsf M(U)(c)^{-1})$ converge; $|n_c-n^\flat_c|\le28\varsigma^2$; \[|\frac12n^\flat_c\cdot n^\flat_c-(1-\frac12\operatorname{tr}V'(c)\mathsf M(U)(c)^{-1})|\le55\varsigma^3;\] and $|\log \mathsf j(\eta_b)|\le r^2/2$ for $b\in\hat\square$.

By \textup{(W1$_\tau$)}, $\varsigma\le1.42/22\le0.0646\le1/15$. Put $\tilde\eta_b=\eta_b$ for $b$ based in $B(c_-)$ and $\tilde\eta_b=\operatorname{Ad}_{V'(c)}\eta_b$ for $b$ based in $B(c_+)$; by Step~4, $|\tilde\eta_b|\le1.42r$. The crossing bond of $[x,x']$ is based in $B(c_-)$ and carries $e^{i\eta_b}V'(c)$. Hence $g_x:=h_{c,x}V'(c)^{-1}=\prod_{b\in[x,x']}e^{i\tilde\eta_b}$, a product of $L$ factors. In this step, $g_x$ and $g_y$ denote these matrices, $x$ and $y$ being sites of $B(c_-)$. Let $\Sigma_x:=\sum\tilde\eta_b$. Then $\sum|\tilde\eta_b|\le\varsigma$, and the mean of $\Sigma_x$ over $x\in B(c_-)$ is $(Q_1^V\eta)_c=-n^\flat_c$.

We have $h_{c,x}h_{c,y}^{-1}=g_xg_y^{-1}$, and for this product the sum of the absolute values of the exponents is the $S$ of \Cref{Y-E-L1C}\,(c), not the $S=6r$ of Step~4, and satisfies $S\le2\varsigma\le1/4$. So \Cref{Y-E-L1C}\,(c) gives $|\tau_{c,x}-\Sigma_x+\Sigma_y|\le\frac32(2\varsigma)^2=6\varsigma^2$, and therefore $|\bar\tau_c|\le2\varsigma+6\varsigma^2\le2.4\varsigma$. Next, $\mathsf M(U)(c)V'(c)^{-1}=e^{i\bar\tau_c}g_y$, and for this product $S\le3.4\varsigma\le3.75\varsigma\le1/4$; its distance from $1$ is at most $e^{1/4}-1\le0.2841<1/2$. So $L_{\mathbb C}$ of its inverse is minus its $L_{\mathbb C}$ (\Cref{Y-E-L1C}\,(b)). The hypothesis $\det(e^{i\bar\tau_c}g_y)=1$ of that clause is met: the $\tilde\eta_b$ are traceless, so $\det g_x=1$; then $\operatorname{tr}\tau_{c,x}=0$ by the clause ``$\det g=1\Rightarrow\operatorname{tr}L_{\mathbb C}(g)=0$'' of \Cref{Y-E-L1C}\,(b), and so $\operatorname{tr}\bar\tau_c=0$; and the same clause gives $\operatorname{tr}n_c=0$, which \Cref{Y-E-L1C}\,(d) requires below.

Applying \Cref{Y-E-L1C}\,(c) to $e^{i\bar\tau_c}g_y$, we obtain $|-n_c-\bar\tau_c-\Sigma_y|\le\frac32(3.75\varsigma)^2\le21.1\varsigma^2$. Together with the bound on $\tau_{c,x}$, averaged over $x$, this gives $|n_c+(Q_1^V\eta)_c|\le27.1\varsigma^2\le28\varsigma^2$. Further, $|n^\flat_c|\le\varsigma$ and $|n_c|\le\varsigma+28\varsigma^2\le2.867\varsigma\le2.9\varsigma$, so that $\frac12|n\cdot n-n^\flat\cdot n^\flat|\le\frac12\cdot28\varsigma^2\cdot3.9\varsigma=54.6\varsigma^3$. Moreover $|n\cdot n|\le8.5\varsigma^2\le1$ and $e^{in_c}=V'(c)\mathsf M(U)(c)^{-1}$, and \Cref{Y-E-L1C}\,(d) gives $|1-\frac12\operatorname{tr}e^{in}-\frac12n\cdot n|\le72.25\varsigma^4/23\le0.21\varsigma^3$. As $54.6+0.21\le55$, the bound on the kernel follows. Finally, $|\eta_b\cdot\eta_b|\le r^2\le1$, and \Cref{Y-E-L1C}\,(d) gives $|\log \mathsf j(\eta_b)|\le r^2/2$.

\emph{Step 8: the sum; clause \textup{(c)}, also within \textup{(f)}.} We claim that $|V_\square|\le K_v^\flat(\tau)p_0^3x^{-1/2}$ and $|e^{V_\square}-1|\le2K_v^\flat(\tau)p_0^3x^{-1/2}$.

The set $\hat\square$ consists of the $(LM)^3$ fine sites of the blocks of $\square$. Its $3(LM)^3$ plaquettes give at most $3\cdot384.75=1154.25$ times $(LM)^3xr^3$, by Step~4. Its $(LM)^3$ sites give at most $K_\ell$ times $(LM)^3xr^3$, by Step~6. Each of the $3M^3$ coarse bonds gives at most $55s\varsigma^3=55a(1.42)^3L^3\,xr^3\le158aL^3\,xr^3$, by Step~7, which makes $474a(LM)^3\,xr^3$ in all. The at most $3(LM)^3$ bonds of $I^\circ$ give $\frac32r^2(LM)^3\le\frac32(LM)^3\,xr^3$, because $xr=C_\eta p_0x^{1/2}\ge1$ by \textup{(W3)}, where we also use $C_\eta^\flat\ge7L^{1/2}\ge1$ and $x\ge1$. Adding up, we obtain $(LM)^3(1155.75+474a+K_\ell)\,xr^3\le(LM)^3(1157+474a+K_\ell)C_\eta^3p_0^3x^{-1/2}$. For $|z|\le1$ we have $|e^z-1|\le(e-1)|z|\le2|z|$, and we apply \textup{(W4$_\tau$)}. The other bounds of \textup{(c)} come from Steps~3, 4 and~5.

\emph{Step 9: clause \textup{(a)}: convergence, dependence, holomorphy.} We claim that $V'\mapsto V_\square(\xi,V')$ is holomorphic on $U^\tau(B;x')$, hence on $\mathcal U_{\{\square\}}(x')$, and that its dependence is the one stated in \textup{(a)}.

The dependence is that of Step~2. The formula of item~3.8 contains no member, no $e$, no $c_h$ and no $E_a$. On the domain we have $\partial V'(P)\ne-1$ at every $P$ in $B$, while $e_P$ occurs in $X'_P$ only where $\partial V'(P)=-1$; hence there is no dependence on $e$. The convergence follows from Steps~3, 4 and~7.

As to holomorphy, $U^\tau(B;x')$ is open in $\operatorname{SL}(2,\mathbb C)^{\text{bonds of }B}$, and $V'\mapsto\operatorname{cont}(V')$ maps it holomorphically into $\mathfrak V^\circ$ (Step~3). By \Cref{Y-C-CITh66}\,(i), the entries of $G$ and of $C$ are holomorphic on $\mathfrak V^\circ$. The entries of $\operatorname{Ad}^{\pm1}_{V'(c)}$ are polynomials in those of $V'(c)$, the inverse being the adjugate; so $\mathrm{d}^V$, $Q_1^V$, $\ell^V$ and $\mathsf g^V$ are polynomial. Each $X'_P$ is either a power series in $\partial V'(P)-1$, which has norm $<2\hat\varepsilon<1$, or $0$. So $\eta$, $F^\flat$ and $n^\flat$ are holomorphic. The functions $e^{i\eta_b}$, $\partial U(p)$, $U(\Gamma)$ and $h_{c,x}$ are entire in $(\eta,V')$, and $\operatorname{Ad}_{U(\Gamma)}$ is polynomial in $U(\Gamma)$, because $\det=1$. The functions $F_p$, $\tau_{c,x}$, $\mathsf M(U)(c)$ and $\log \mathsf j(\eta_b)$ are power series whose arguments stay inside the discs of convergence on the whole domain, since there $|\partial U(p)-1|\le e^{1/11}-1$, $|g_xg_y^{-1}-1|\le e^{1/4}-1$ and $|\eta\cdot\eta|\le1$ (Steps~4 and~7, whose bounds are uniform on the domain). All sums are finite. Finally, $\mathcal U_{\{\square\}}(x')$ is an open subset of $U^1(\text{sites of }\square^\approx;x')$.

\emph{Step 10: clause \textup{(b)}: invariance, with the extension by $1$.} We claim that for every $G$-valued $u$ on the sites of $B$ we have $V_\square(\operatorname{Ad}_u\xi,V'^u)=V_\square(\xi,V')$, and that the domain and $\chi_\square$ are invariant.

\textup{(1)} Norms do not change under unitary factors; moreover $V'^u(\pi)=uV'(\pi)u^{-1}$ and $\partial(V'^u)(P)-1=u(\partial V'(P)-1)u^{-1}$. So \textup{(e1)} to \textup{(e3)}, and likewise \textup{(u1)} to \textup{(u3)}, are invariant; and $|\operatorname{Ad}_u\xi_b|=|\xi_b|$.

\textup{(2)} Let $W=\operatorname{cont}(V')$, let $\tilde u=u$ on $B$ and $\tilde u=1$ elsewhere, and let $W_2=\operatorname{cont}(V'^u)$. Then $W^{\tilde u}\in\mathfrak V^\circ$, and $W^{\tilde u}$ agrees with $W_2$ on every bond with both ends in $B$, while the two may differ on bonds that leave $B$. For $P$ in $B$ we have $X'[W_2]_P=\operatorname{Ad}_{u(y_P)}X'[W]_P$ (\Cref{Y-E-L1C}\,(b)), and elsewhere both vanish; so $X'[W_2]=\operatorname{Ad}_{\tilde u}X'[W]$.

\textup{(3)} \Cref{Y-C-CITh66}\,(c), applied to fields of $\mathfrak V$ and to the $G$-valued $\tilde u$ on all of $\mathbb T'$, gives $G[W^{\tilde u}]=\operatorname{Ad}_{\tilde u}G[W]\operatorname{Ad}_{\tilde u}^{-1}$, and the same for $C$. \Cref{Y-C-CIL64}\,(d) and~(e) give $\mathsf g^{W^{\tilde u}}=\operatorname{Ad}_{\tilde u}\mathsf g^W\operatorname{Ad}_{\tilde u}^{-1}$. Hence at the field $W^{\tilde u}$, with $\operatorname{Ad}_{\tilde u}X'$ and $\operatorname{Ad}_{\tilde u}\xi$, we have $\eta=\operatorname{Ad}_{\tilde u}\eta[W]$.

\textup{(4)} By Step~2\,(ii) and~(iii), the rows $G(b,\cdot)$ and the $a_{0,b}$ at ${\rm dist}(b,\square)\le\varrho_L+2$ depend only on the field and on $X'$ in $\square_{[R_v]}$, which lies inside $B$; and there $W^{\tilde u}=W_2$. The $b'$ that enter have $y(b')\in\square_{[R_v]}$, where $\tilde u=u$. So $\eta[W_2,\operatorname{Ad}_u\xi]_b=\operatorname{Ad}_{u(y(b))}\eta[W,\xi]_b$ at these bonds.

\textup{(5)} \emph{The rest.} We have $T^u(y,y'')=\operatorname{Ad}_{u(y)}T\operatorname{Ad}_{u(y'')}^{-1}$. Hence $F^\flat_p\to\operatorname{Ad}_{u(y(p))}F^\flat_p$, $n^\flat_c\to\operatorname{Ad}_{u(c_-)}n^\flat_c$ and $(\ell^VF^\flat)_y\to(1\otimes\operatorname{Ad}_{u(y(y))})(\ell^VF^\flat)_y$. Further, $\mathsf U(\operatorname{Ad}_u\eta;V'^u)=\mathsf U(\eta;V')^{\hat u}$ with $\hat u(x)=u(y(x))$: inside a block $e^{i\operatorname{Ad}_u\eta}=ue^{i\eta}u^{-1}$, and at a crossing bond $e^{i\operatorname{Ad}_{u(c_-)}\eta}u(c_-)V'(c)u(c_+)^{-1}=u(c_-)e^{i\eta}V'(c)u(c_+)^{-1}$. So $F_p\to\operatorname{Ad}_{\hat u(\pi(p))}F_p$, $\operatorname{Ad}_{U(\Gamma(y,\pi(p)))}\to\operatorname{Ad}_{\hat u(y)}\operatorname{Ad}_{U(\Gamma)}\operatorname{Ad}_{\hat u(\pi(p))}^{-1}$ and $\ell_y^UF\to(1\otimes\operatorname{Ad}_{\hat u(y)})\ell_y^UF$. Also $h_{c,x}\to u(c_-)h_{c,x}u(c_+)^{-1}$, and $\mathsf M(U)(c)$ transforms in the same way. Traces and the products $Y\cdot Y'$ are $\operatorname{Ad}$-invariant (\Cref{Y-E-L1C}\,(a)).

\emph{Step 11: clause \textup{(d)}.} Let $(V',e)$ be given on $\mathbb T'$, with the restriction of $V'$ to $\square^\approx$ in $\mathcal U_{\{\square\}}(x')$; let $\chi_\square(\xi)=1$, and let \textup{(W5)} hold. We claim that $\phi_p\le\varepsilon$ at every fine plaquette $p$ with $|y-p|\le\varrho(x)+1$ for some $y\in\hat\square$; that $V_\square$ is real; that $|\eta_b|<\pi$; and that $A_\square(\xi,V',e)=e^{V_\square(\xi,V')}$.

\emph{\textup{(1)} One $\eta$.} On $\square^\approx$ we have $|\partial V'(P)-1|<2\varepsilon'<1$. Hence $\partial V'(P)\ne-1$, and $X'_P(V',e)=\mathsf L(\partial V'(P))=L_{\mathbb C}(\partial V'(P))$ (\Cref{Y-E-L1C}\,(b)). The global pair $(V',X'(V',e))$ agrees on $\square_{[R_v]}$ with the pair consisting of $\operatorname{cont}(V'|\square^\approx)$ and the $X'$ of item~3.7. By Step~2, the $\eta_b$ at ${\rm dist}(b,\square)\le\varrho_L+2$, and everything on which $A_\square$ depends, are the same for the two; and $e$ does not enter.

\emph{\textup{(2)} Reality.} The matrix $G^V$ is real (\Cref{Y-C-CITh66}\,(a)), and $X'$ is $\mathfrak g$-valued. So $\eta_b\in\mathfrak g$ and $|\eta_b|=\|\eta_b\|\le r<\pi$. On $\mathfrak g$ we have $Y\cdot Y=|Y|^2$, and $L_{\mathbb C}=\mathsf L$ where $|g-1|<1$. Therefore $F_p=\mathsf L(\partial U(p))$, $|F_p|=\phi_p$, $\mathsf M(U)$ is the one of item~3.4, and $1-\frac12\operatorname{tr}g=1-\cos\operatorname{ang}g$. So every bracket of $V_\square$ is real.

\emph{\textup{(3)} The angle.} Let $p$ be as in the claim. Then ${\rm dist}(p,\square)\le\varrho_L+1$, its bonds are at distance at most $\varrho_L+2$, and the box of radius $R_3$ about $y(p)$ lies in $\square_{[R_v]}$ (Step~2). We have $F^\flat_p=\bar F_p+(\mathrm{d}^VG\xi)_p$. The restriction of $V'$ to that box satisfies \textup{(e1)} and \textup{(e2)}, which follow from \textup{(u1)} and \textup{(u2)}, and $|\partial V'(P)-1|\le2\varepsilon'$; and its $a_{0,b}$, $b\in dp$, are the ones at hand (Step~2\,(iii)). \Cref{Y-C-CIICor72}\,(a), applied with the value $\hat\varepsilon=\varepsilon'$ of its parameter, gives $|\bar F_p|\le \mathsf N_{\rm c}L^{-2}\varepsilon'+\mathsf K_{\rm c}\varepsilon'^2\le(\mathsf N_{\rm c}L^{-3/2}+\mathsf K_{\rm c}L\varepsilon)\varepsilon$. At real fields $T$ is an isometry, so $|(\mathrm{d}^VG\xi)_p|\le4\sup|(G\xi)_b|\le(4/6.5)\varepsilon\le0.61539\varepsilon$ (Step~3). By Step~4, $|F_p-F^\flat_p|\le54(C_\eta^\flat)^2\varepsilon^2$. Adding the three bounds, we get from \textup{(W5)} that $\phi_p\le[0.61539+0.2+0.18]\varepsilon=0.99539\varepsilon<\varepsilon$. We note that \textup{(L2$^{\rm su}$)} gives $\mathsf N_{\rm c}L^{-3/2}\le(1.2\cdot10^4\mathsf N_{\rm c}L)^{-1}\le10^{-4}$.

\emph{\textup{(4)} Consequences.} Let $y\in\hat\square$. At every $p$ in the range of $\ell_y$ we have $\zeta_p=F_p$. No plaquette of $D(U)$ lies within $\rho_0\le\varrho(x)+1$ of $y$ (by \textup{(W0)}); and no plaquette with $\phi_p>1\ge\varepsilon$ lies within $\varrho(x)+1$ of $y$, where $\varepsilon\le1$ by \textup{(W1)}, because $C_\eta^\flat\ge1$. So $m_y(U)=1$. For $p\in\hat\square$ we have $S_\varepsilon(\phi_p)=\frac12|F_p|^2$. Moreover $1\{|\eta_b|<\pi\}=1$ and $\mathsf j(\eta_b)=e^{\log \mathsf j(\eta_b)}>0$. Therefore the four products of item~3.9 are the exponentials of the four sums of item~3.8.

\begin{remark}[On Step~11]
If the tests were taken at $\varepsilon$ instead of at $c_\xi\varepsilon$, the second term would be bounded only by $4C_\varphi^\flat\varepsilon$, and the step would fail.
\end{remark}

\emph{Step 12: clauses \textup{(e.iii)}, \textup{(e.iv)}, \textup{(e.v)} and \textup{(g)}: $A_\square$ at every real point.} We claim that $A_\square(\xi,V',e)$ has the dependence, the covariance and the measurability that these clauses state, and that no member enters it.

\emph{Proof of \textup{(e.iii)}.} We use Step~2\,(i) with $\varrho_L+2$, and Step~2\,(ii); no field is estimated there. The value $e_P$ enters only through $X'_P$, hence only at $P$ in $\square_{[R_v]}$, and only where $\partial V'(P)=-1$.

\emph{Proof of \textup{(e.iv)}.} Here $u:\mathbb T'\to G$ is global, and no continuation is constructed. We have $X'(V'^u,\operatorname{Ad}_ue)=\operatorname{Ad}_uX'(V',e)$ (item~E2 of \S\ref{Y-D-1}). By \Cref{Y-C-CITh66}\,(c) and \Cref{Y-C-CIL64}\,(d), (e), $\eta\to\operatorname{Ad}_u\eta$. As in Step~10\,(5), $U\to U^{\hat u}$. The quantities $\phi_p$, $D(U)$, $m_y(U)$ and $|\ell_y^U\zeta|$ are invariant under fine gauge transformations, because $\zeta_p\to\operatorname{Ad}_{\hat u(\pi(p))}\zeta_p$ and the transports are conjugated. Moreover $\mathsf L(ugu^{-1})=\operatorname{Ad}_u\mathsf L(g)$, also at $g=-1$; this gives $\tau_{c,x}\to\operatorname{Ad}_{u(c_-)}\tau_{c,x}$ and $\mathsf M(U)(c)\to u(c_-)\mathsf M(U)(c)u(c_+)^{-1}$, so that $\mathsf n_c$ is invariant. Finally, $|F^\flat_p|$, $|(\ell^VF^\flat)_y|$, $|n^\flat_c|$ and $|\eta_b|$ are invariant, since $\operatorname{Ad}_u$ is orthogonal.

\emph{Proof of \textup{(e.v)}.} The matrices $G^V$ and $C^V$ are continuous on $\mathfrak V$ (\Cref{Y-C-CITh66}\,(i)), and $\mathsf g^V$, $\mathrm{d}^V$, $Q_1^V$, $\ell^V$ are polynomial. By item~E2 of \S\ref{Y-D-1}, $X'$ is Borel in $(V',e)$; so $\eta$ is Borel in $(\xi,V',e)$. The field $U$ is continuous in $(\eta,V')$, $\operatorname{ang}$ is continuous and $\mathsf L$ is Borel; and $m_y$, the clip and $S_\varepsilon$ are Borel functions of these.

\emph{Proof of \textup{(g)}.} The formulas of items~3.6 to~3.9 contain $x$, $\ell$, $L$, $a$ and $R(x)$; the number $b$ of \textup{(O6)}, in $K_b$ and in $G$ (item~3.6); the sets $\hat\square$ and $\square$, and with them $M$ and the grid; and the lattice operations of item~3.1, and with them $N$. Besides these, $A_\square$ contains $\varepsilon$, and with it $A_0$ and $\mathfrak p_0$, as well as $d_0$ and $\rho_0$. Nothing else occurs.

\emph{Step 13: clause \textup{(f)}, with \textup{(a)} to \textup{(c)} as its case $\tau=1$.} We claim that \textup{(a)}, \textup{(b)} and \textup{(c)} hold on $U^\tau(B;x')$ with $r_\tau$ and $K_v^\flat(\tau)$, and on $\mathcal U_{\{\square\}}(x')$ with $r$ and $K_v^\flat$.

Steps~3 to~10 are carried out for $\tau\ge1$, $B\supset\square_{[R_v]}$ and $V'\in U^\tau(B;x')$. They depend on the field only through $\operatorname{cont}(V')\in\mathfrak V^\circ$ (from \textup{(e1)}), through $|V'(\pi)|<2$ on paths in $B$ of at most $\ell_w(x')$ bonds (from \textup{(e2)}), and through $|\partial V'(P)-1|<2\tau\varepsilon'$, which is \textup{(e3)}. Of the floors, they depend only on \textup{(W0)}, \textup{(W1$_\tau$)}, \textup{(W2$_\tau$)}, \textup{(W3)} and \textup{(W4$_\tau$)}. The inequality $\tau\varepsilon'\le1/40$ is a consequence of these, by Step~3. At $\tau=1$, with $B$ the set of sites of $\square^\approx$, which contains $\square_{[R_v]}$ by Step~2\,(iv), the conditions \textup{(u1)}, \textup{(u2)} and \textup{(u3)} imply \textup{(e1)}, \textup{(e2)} and \textup{(e3)}. Hence $\mathcal U_{\{\square\}}(x')\subset U^1(B;x')$, with the constants $C_\eta^\flat$ and $K_v^\flat$.

\emph{Step 14: the order of the letters, and the floors; no circle.} We claim that $C_\eta^\flat$, $K_\ell$ and $K_v^\flat$ depend only on earlier entries of the table, and that \textup{(W0)} to \textup{(W5)}, together with \textup{(W1$_\tau$)}, \textup{(W2$_\tau$)} and \textup{(W4$_\tau$)} for every $\tau\ge1$, are conditions on nothing but $x$, the letters \textup{(O0)} to \textup{(O10)} and $\tau$, and that each of them holds for all large $x$ once the letters and $\tau$ are fixed.

\emph{Order.} The numbers $\gamma$, $\mathsf N_{\mathsf B}$, $b$, $\theta_{\mathbb C}$, $\vartheta_0^\flat$, $\mathsf G_0$, $C_\varphi^\flat$, $c_\xi$ and $\mathsf a_{\rm c}$ depend on $(L,a,K_*,\lambda_*)$, and $\mathsf K_{\rm c}^0$ depends, besides, on $C_D$ and $\theta_D$. The constant $C_\eta^\flat$ comes after $\mathsf a_{\rm c}$; $K_\ell$ is computed from $K_*$ and $\lambda_*$; $K_v^\flat$ comes after $M$; $E_a\ge1500K_v^\flat L^{-1/2}$ comes after $K_v^\flat$; and $x_\flat$ comes last. The number $\mathsf N_{\rm c}$ is a function of $L$ through $\mathsf l_L$, and as such it occurs in \textup{(L2$^{\rm su}$)}, a condition to which $L$ is subject. No number of \Cref{Y-E-C6} enters a condition on $L$, $M$ or $c_R$. As $\mathsf a_{\rm c}\ge100\cdot73\cdot3^4>7$, the maximum in $C_\eta^\flat$ is $\mathsf a_{\rm c}$.

\emph{Floors.} The number $\varepsilon=A_0(\log x)^{\mathfrak p_0}x^{-1/2}$ tends to $0$, and on the range $\log x>2\mathfrak p_0$ it is a decreasing function of $x$. In \textup{(W1$_\tau$)} and in the last clause of \textup{(W5)} a number times $\varepsilon$ is bounded; in \textup{(W2$_\tau$)}, a number times $\varepsilon(c_R\log x+1)$; and in \textup{(W4$_\tau$)}, a number times $(\log x)^{3\mathfrak p_0}x^{-1/2}$. Condition \textup{(W3)} amounts to $\log x\ge A_0^{-1/\mathfrak p_0}$. Condition \textup{(W0)} rests on $R\to\infty$. Condition \textup{(F$_{\rm law}$)} compares $(R+3)^3$ with $x^4$; \textup{(P1)} compares $(2R+3)^3$ with $x$, and contains $x^{-60}$; and \textup{(P3)} contains negative powers of $x$ and $x'$. The second clause of \textup{(W5)}, $\mathsf N_{\rm c}L^{-3/2}\le0.2$, involves no $x$: it is a condition on the letters alone, and Step~11\,\textup{(3)} shows that $\mathsf N_{\rm c}L^{-3/2}\le10^{-4}$. The clauses of \textup{(P0)} on $N$ are \textup{(Tor)}, which belongs to the data. None of these conditions depends on $c_h$, $E_a$, $N$, a member, or on more of $\ell$ than $K_*$, $\lambda_*$ and the range. All the conditions are one-sided: a larger $\mathsf K_{\rm c}$ or $\mathsf a_{\rm c}$ only raises the floors and $E_a$.

\emph{Step 15: clause \textup{(h)}, the inequalities that compare the radii, counted in cells, with $r_1$.} We claim that $R_\chi^{\min}\le R_\chi$, $2\lceil R_\chi/M\rceil\le r_1-5$ and $\lceil R_v/M\rceil\le r_1-5$, together with the three consequences stated in \textup{(h)}.

Since $r_w\le R$, we have $R_\chi^{\min}\le R_\chi$. Next, $\lceil(R+\varrho_L+2)/M\rceil\le\lceil R/M\rceil+r_l$; moreover $2\lceil t\rceil\le\lceil2t\rceil+1$, which gives $2\lceil R/M\rceil\le\lceil R_2/M\rceil+1$. Hence $2\lceil R_\chi/M\rceil\le\lceil R_2/M\rceil+1+2r_l=r_1-r_l-4\le r_1-5$, because $r_l\ge1$. Further, $R_v=R_2+2(\varrho_L+2)+1\le M\lceil R_2/M\rceil+2Mr_l+M$, and therefore $\lceil R_v/M\rceil\le\lceil R_2/M\rceil+2r_l+1=r_1-r_l-4\le r_1-5$.

\emph{Consequences.} A bond of $I^\circ$ with ${\rm dist}(b,\square)\le R_\chi$ is based under $\square^{+\lceil R_\chi/M\rceil}$, which lies inside $\square^{+r_1}$; so, if the cell tests there are passed, then $\chi_\square=1$. Clause \textup{(e.iii)} places the $\xi$ on which $A_\square$ depends under $\square^{+\lceil R_\chi^{\min}/M\rceil}$, and the $(V',e)$ on which it depends in $\square^{+\lceil R_v/M\rceil}$. Two cells that are more than $2\lceil R_\chi^{\min}/M\rceil$ cells apart have disjoint sets $\square_{[R_\chi^{\min}]}$.
\end{proof}

\subsection{Gaussian integration on zones}\label{Y-F-1}

We estimate the cell factors after integrating the white noise.
Compatible sets have disjoint noise supports, so their integrals
factor. The main bound suppresses a failed test by charging either
a large fluctuation in a tame block or a rough plaquette. It follows
from a lower bound for the fine energy plus the noise, after
subtracting the coarse energy and controlling the terms at the
boundary of the region.

\begin{remark}[Notation and letters with more than one meaning]
The lift $E_3$ maps a function $\Theta=(\Theta_{q'})$ on the coarse cubes $q'$ to the function $E_3\Theta=((E_3\Theta)_q)$ on the fine cubes $q$: $(E_3\Theta)_q=\Theta_{q'}$ when $q$ is the corner cube under $q'$, that is, the fine cube at $y_{q'}+\mathsf h_L(1,1,1)$ with $\mathsf h_L=(L-1)/2$, and $(E_3\Theta)_q=0$ at the other fine cubes (\S\ref{Y-D-1}, C1); the identity $\mathrm{d}^VE=E_3\mathrm{d}'^V$ is Step~1\,\textup{(iv)} of the proof of \Cref{Y-D-C5}. We write $E_\xi$ for the mean over $\xi$. Moreover, $\mathrm e=\exp(1)$, and $\mathfrak g$ and $\mathfrak L_{\mathfrak g}$ are those of item~C.I~1.1 of \S\ref{Y-C-1}.

The exponent in $M=L^{m'}$ of \textup{(O7)} is $m'$, whereas $m'_{y'}$, and every other $m'$ of this section, denote the mask at $x'$. The letter $G$ denotes the matrix of the law and also the group $\operatorname{SU}(2)$, as in ``$G$-valued'' and ``$g_i\in G$''. The letters $k_F,k_n,k_p$ are the components of $\mathfrak k$, while $k_p$ in Step~3\,\textup{(3)} is a matrix.

The coboundaries are $\mathrm{d},\mathrm{d}'$, while in ``$b\in dp$'' and ``$c\in dP$'', $dp$ and $dP$ are the sets of the four bonds.
\end{remark}

\begin{theorem}[Means of zone and test letters at any depth]\label{Y-F-C7}

\emph{Item A: letters.} The letters are chosen in the order below. This order, from \textup{(O0)} to \textup{(O13)}, is that of the table of the letters in item~3 of \S\ref{Y-G0-1}. After each tag we state what this theorem uses of the entry of the table under that tag.
\begin{enumerate}
\item[\textup{(O0)}] The pure numbers $b_*,K_*,\lambda_*,L_A,c_1^*,c_2^*,\lambda_2^*,\theta_J$ of \Cref{Y-A-Fr}, the order of the constants, \textup{(0)}. We put $K=K_*$ and $\lambda=\lambda_*$, and we set $\mathsf g(s)=\sum_{n\in\mathbb Z^3}e^{-s|n|_\infty}$, $\mathsf g_*=\mathsf g(\lambda)$, $\mathsf g_1=\sum_{n\in\mathbb Z^3}|n|_\infty e^{-\lambda|n|_\infty}$, $l_r:=3K\mathsf g_*$ and $l_c:=K\mathsf g_*$.

\item[\textup{(O1)}] A number $d_0\in(0,1/2)$.

\item[\textup{(O2)}] A number $c_\Xi$ with $0<c_\Xi\le\frac12(\frac12-d_0)$.

\item[\textup{(O3)}] The number $\rho_0$ is the least integer $\ge0$ for which $3K^2\mathsf g(\lambda)\mathsf g(\lambda/2)e^{-\lambda\rho_0/2}\le\frac18(\frac12-d_0-c_\Xi)$; this is display \eqref{eq:row-mask-radius} of \S\ref{Y-A-1}. The functions of $L$ are $\mathsf l_L=1+\log L$ and the four numbers of display \eqref{eq:row-cut-norms} of \S\ref{Y-A-1}, namely $\mathsf N_Z=\mathsf N_{\mathsf B}=1+3K^2\mathsf g_*^2$, $\mathsf P_\rho=1+3(2+3(2\rho_0+1)^3)K^2\mathsf g_*^2$, $\mathsf N_\sharp=c_1^*\mathsf l_L\mathsf N_Z\mathsf P_\rho$ and $\mathsf N_\Xi=c_2^*\mathsf N_Z$. Moreover $\mathsf N_\sigma=3K\mathsf g_*c_1^*\mathsf l_L\mathsf N_Z$, $N_1=\mathsf N_Z+2L\mathsf N_\Xi+l_c\mathsf N_\sigma$, and \[\mathsf N_\nu^\flat=N_1+(2l_c+(2\rho_0+1)^3l_r)\mathsf N_\sigma=\mathsf N_Z+2L\mathsf N_\Xi+(3+3(2\rho_0+1)^3)K\mathsf g_*\mathsf N_\sigma.\]

\item[\textup{(O4)}] An odd integer $L$ with the following three properties.
\begin{enumerate}
\item[\textup{(i)}] $L\ge L_A$.

\item[\textup{(ii)}] \textup{(L1$^\flat$)} $[\mathsf N_\nu^\flat L^{-1}+2200\mathsf N_Z+40\mathsf N_\Xi]L^{-1/2}\le\frac3{16}d_0$.

\item[\textup{(iii)}] $(\mathsf N_\sharp+2)L^{-3/2}\le0.08$.
\end{enumerate}

Properties \textup{(ii)} and \textup{(iii)} hold for all large $L$, because $\mathsf N_Z,\mathsf N_\Xi$ are free of $L$, while each of $\mathsf N_\sharp$, $\mathsf N_\sigma$, $\mathsf N_\nu^\flat/L$ is at most a number free of $L$ times $(1+\log L)$. We set $c_L=(L^2-1)/(6L^2)$, and the letter $a$ satisfies $a\ge3/(c_LL)$, so that $aL\ge18$. The condition $L\ge L_*^{\rm su}$ of item~3 of \S\ref{Y-G0-1}, \textup{(O4)}, implies \textup{(i)}, \textup{(ii)} and \textup{(iii)}: it includes $L\ge L_A$ and \textup{(L1$^\flat$)}, and its \textup{(L2$^{\rm su}$)}, which reads $L^{1/2}\ge1.2\cdot10^4\mathsf N_{\rm c}^2$ with $\mathsf N_{\rm c}=17(\mathsf N_\sharp+1)$, gives $L^{3/2}\ge10^{12}\mathsf N_{\rm c}^6\ge12.5(\mathsf N_\sharp+2)$.

\item[\textup{(O5)}] A number $\kappa\ge6$.

\item[\textup{(O6)}] The numbers of item~C.I~1.2 of \S\ref{Y-C-1}, evaluated at $(K_*,\lambda_*)$: $\gamma=(73L^4)^{-1}$, $b_h=16(1+3K^2\mathsf g_*^2)+a/L$, \[A^\flat=2\mathrm e^2(16+a)+384\mathrm e^8K^2\mathsf g_2(\lambda/2)^2\] with $\mathsf g_2(s)=\sum_{n\in\mathbb Z^3}(1+|n|_\infty)^2e^{-s|n|_\infty}$, then $\theta_1=\min\{1,\lambda L/12\}$, $\theta^\flat=\min\{\theta_1,\gamma/(2A^\flat)\}$, $\theta_{\mathbb C}=\theta^\flat/6$ and $\vartheta_0^\flat=\theta^\flat/24$. We also set $\theta_D=\frac12\min\{\lambda_2^*,\theta_{\mathbb C},\theta_J\}$. The number $C_D\ge4$ is fixed at this stage, before any rows are chosen: it is the number given, at this $\theta_D$, by the expression for $C_D$ in \Cref{Y-A-Fr}, the order of the constants, \textup{(6)}, or else any number above that one, and that expression is $\ge A_q\ge32\mathrm e^2>4$. That \textup{(Dec)} holds at this $C_D$ for the pair $(\ell,\ell')$ is itself one of the data of item~C. A number $C_D$ with \textup{(Dec)}$(C_D,\theta_D)$ comes from \Cref{Y-A-P35}, taken with $\theta_{\mathbb C}$ in the place of $\theta_q$; with this change, Step~H13 of the proof of \Cref{Y-A-Fr} states the assertion of that proposition for every $\ell\in\mathcal R(b_*,K_*,\lambda_*,1/9)$, and clause \textup{(Fr2)} of the part ``Conclusion'' of that theorem states it for all sixteen kernels of \textup{(Dec)} at the pairs $(\ell_k,(\ell_k)')$ of the chain. Finally $c_R=64/\theta_D$, so that $c_R\theta_{\mathbb C}\ge128$.

\item[\textup{(O7)}] The side $M=L^{m'}$ with $m'\ge1$. We put $v:=(LM)^3$.

\item[\textup{(O8)}] Numbers $\mathfrak p_0>2$ and $A_0>0$. We put $p_0=p_0(x)=A_0(\log x)^{\mathfrak p_0}$, $\varepsilon=p_0x^{-1/2}$ and $x'=x/L$, and a prime on a quantity means that it is taken at $x'$. Then $\varepsilon'\le L^{1/2}\varepsilon$, since $\varepsilon'=A_0(\log x')^{\mathfrak p_0}L^{1/2}x^{-1/2}$ and $0<\log x'\le\log x$, the latter because $x'=x/L\ge L>1$.

\item[\textup{(O9)}] The numbers \[\mathsf G_0=(2/\gamma)^{1/2}3L^3\mathsf g(\theta_{\mathbb C})\le C_\varphi^\flat=(2b_h^{1/2}/\gamma)3L^3\mathsf g(\theta_{\mathbb C})/(1-e^{-\theta_{\mathbb C}})\] and $c_\xi=(6.5C_\varphi^\flat)^{-1}$. We use \S\ref{Y-C-1}, Part C.II, with the value $\mathsf n=60$ of its exponent.

\item[\textup{(O10)}] The number $\mathsf K_{\rm c}^0$ of item~C.II~4 of \S\ref{Y-C-1}, ``Constants''; here $\mathsf K_{\rm c}^0\ge6\cdot10^4$. We also define $\mathsf a_{\rm c}=(100/\gamma)\mathsf N_{\mathsf B}3L^3\mathsf g(\theta_{\mathbb C})$ and $\mathsf K_{\rm c}:=\sqrt2\mathsf K_{\rm c}^0\ge \mathsf K_{\rm c}^0+1$. The constants of item~F, $\delta_3$ and $\delta_3'$ among them, are defined in that item, and they come at \textup{(O10)} in the order \textup{(O0)} to \textup{(O13)} of the letters.

\item[\textup{(O11), (O12)}] The letters $c_h$, $E_a$ play no role in this section.

\item[\textup{(O13)}] The number $x_{\rm fl}$, which is defined in item~G. To apply this theorem for every $x\ge Lx_\flat$, one needs $Lx_\flat\ge x_{\rm fl}$.
\end{enumerate}

\emph{Item C: data.} The coupling satisfies $x\ge x_{\rm fl}$. The coarse side $N$ is even, $M\mid N$, and the condition \textup{(Tor)} $N\ge2Ms(x)$ holds. The fine torus is $\mathbb Z^3/NL\mathbb Z^3$ and the coarse torus is $L\mathbb Z^3/NL\mathbb Z^3$, with $N^3$ sites; the grids of cubes are fitted as in item~1 of \S\ref{Y-G0-1}, Cubes, so that a cell fills exactly $L^3$ fine cubes.

The rows are $\ell\in\mathcal R(b_*,K_*,\lambda_*,1/9)$ of range $\le\varrho(x)$, the sets of rows being those of \S\ref{Y-A-1}, from ``Translation covariant rows'' to ``The sets of rows''. The empty rows are admitted: $\ell=\emptyset$, with $\mathsf H_y=\{0\}$ for all $y$, lies in $\mathcal R(1,0,\lambda_*,0)$, which is a part of that set (\Cref{Y-A-SetsOfRows}, with its display and its last sentence), and its range is $0$. At $\ell=\emptyset$ we have $\mathsf B=1$ and $\sigma_*=0$, and the entries at the site rows of a triple of $\mathrm{Hil}$ lie in the spaces $\mathsf H_y=\{0\}$, so that they vanish and \textup{(Ta)}, \textup{(Tb)}, \textup{(Tc)}, $\mathrm{Rim}_{\mathfrak Z}$ and the ``Second'' determinant of Step~8 are void. There, every quantity that $l_r,l_c,K\mathsf g_1,c_{\rm tr}$ bound is $0$, and all numbers and radii are those evaluated at $(K_*,\lambda_*,\varrho(x))$.

We denote by $\ell'$ the family of rows obtained from $\ell$ by \Cref{Y-A-P5}, and we set $\mathsf B'=1+\ell'^{\mathsf T}\ell'$. The pair $(\ell,\ell')$ satisfies the four conclusions of \Cref{Y-A-Th36}: \textup{(Dec)}$(C_D,\theta_D)$; \textup{(N$'$)} with \eqref{eq:row-cut-norms} at every cut; $\mathsf t_\rho\le\frac18(\frac12-d_0-c_\Xi)$; and \textup{(L1)}, \textup{(L2)}. Here \textup{(L1)}, \textup{(L2)} are taken as \textup{(L1$^{\rm su}$)}, \textup{(L2$^{\rm su}$)} of item~3 of \S\ref{Y-G0-1}, \textup{(O4)}, as in item~1 of \S\ref{Y-G0-1}, \textup{(SD)}.

The pair $(V',e)$ is any pair of \S\ref{Y-D-1}, \textup{E2}. The system $w(X,\cdot)$, $X\in\mathcal P$, is any system of complex Borel functions, each of which is a function of $U$ on the bonds of $X^{+s(x)}$. A further datum is a number $A\ge1$, the constant in the two caps that follow. The system obeys either the strong cap \[|w(X,U)|\le A\exp\{-c_\Xi x\varepsilon\sum_{D(U)\cap X}\phi_p-\kappa(|X|_c-1)\}\] or the weak cap \[|w(X,U)|\le Ae^{-\kappa(|X|_c-1)}.\]

\emph{Item D: inputs.}
\begin{enumerate}
\item[\textup{(C2)}] We take \Cref{Y-C-CITh66}\,(a), (b), (d), (f), (k), together with \Cref{Y-C-CIL64}\,(a), (b), (c), (e), for real $V'$; they give the following. The matrix $G$ is real, and $G(b,b')=0$ beyond the distance $r_w$; $(G\xi)_b$ depends on $\xi$ only through its values on $Y_b=\{b'\in I^\circ:{\rm dist}(b,b')\le r_w\}$; $G(b,\cdot)$ depends on $V'$ only within $r_w+\varrho_L+1$ of $y(b)$, and $C(b,\cdot)$ only within $3r_w+\varrho_L+1$. Next, $C=GG^{\mathsf T}$, $C(b,b')=0$ beyond the distance $2r_w\le R$, and $\|G\|_\infty,\|G^{\mathsf T}\|_\infty\le \mathsf G_0$. In the setting \textup{(F$_{\rm law}$)}, with $c_R\theta_{\mathbb C}\ge128$, the matrix $G$ is invertible and $S=C^{-1}-q$ is symmetric with $\|S\|\le x^{-60}$. Further, $\gamma\le q^V_0:=|\mathrm{d}^V\cdot|^2+a|Q_1^V\cdot|^2\le q^V_\ell\le b_h$; the first inequality is \Cref{Y-C-CIL64}\,(a) applied to the empty rows $\ell=\emptyset$, which item~C.I~1.1 of \S\ref{Y-C-1} admits, with $K=0$, and $\gamma$ depends on no row. The absolute row sums of $q^V_\ell$ are at most $b_h$, $q^V_\ell(b,b')=0$ beyond the distance $r_q=2\varrho_L+2$, and $|\mathsf g^VX'|_\infty\le(b_h/4)\sup|X'|$. As a consequence, at real ($G$-valued) fields, $\|C\|\le2/\gamma$, because $C^{-1}=q+S\ge\gamma-x^{-60}\ge\gamma/2$, where $\gamma^{-2}\le x^8$ by \textup{(F9$^\flat$)}.

\item[\textup{(C3)}] We take \Cref{Y-C-Th68}\,(a), (b), (c) for real fields and with $\hat\varepsilon=\varepsilon'$, as in \S\ref{Y-C-1}, Part C.II, Remarks \textup{(2)}, \textup{(3)} after \Cref{Y-C-Th70}; the hypothesis $\tau_{\rm c}\le\varepsilon'$ of these remarks holds by \textup{(P1)} and by $x^{-59}\le\varepsilon'$ of \textup{(FA)}. For $i=p$ or $c$, with $o=y(p)$ or $o=c_-$ respectively, they give $|\Lambda_i|\le \mathsf K_{\rm c}^0\varepsilon'^2+\tau_{\rm c}\varepsilon'\le \mathsf K_{\rm c}\varepsilon'^2$, where $\Lambda_i$ is the remainder given by the display before \Cref{Y-C-Th68} in item~C.II~4 of \S\ref{Y-C-1}, and differs from the $\Lambda_i$ of item~H\,\textup{(a)}; they also give $|a_{0,b}|\le \mathsf a_{\rm c}\varepsilon'$ whenever ${\rm dist}(b,o)\le1$. Both quantities depend on $V'$ only through its values on the bonds in the box of radius $R_3$ about $o$, at fields $V'$ tame there; a $G$-valued field with $\phi'_P\le\varepsilon'$ has \textup{(h)} at $\varepsilon'$, since $\sin(\phi'/2)\le\varepsilon'$ (the passage in parentheses after clause \textup{(c)} of \Cref{Y-C-Th68}). The law of \textup{(C2)} has the clauses \textup{($\lambda$0)}, \textup{($\rho$1)}, \textup{($\rho$2)}, \textup{($\rho$5)} that \Cref{Y-C-Th68} assumes of a law:
\begin{enumerate}
\item[\textup{($\lambda$0), ($\rho$1)}] Clause \textup{($\lambda$0)} follows from clause \textup{(a)} of \Cref{Y-C-CITh66} together with the part of that theorem preceding its clauses, and clause \textup{($\rho$1)} holds by \Cref{Y-C-CITh66}\,(b).

\item[\textup{($\rho$2)}] It holds because $3r_w+\varrho_L+1\le R_2-1$.

\item[\textup{($\rho$5)}] It holds with $\mathsf n=60$ by \Cref{Y-C-CITh66}\,(f), (k).
\end{enumerate}

\Cref{Y-C-L71}\,(b), (c) are inputs as well, with the matrices $H^\natural,\mathsf N^\natural$ of item~C.II~2.2 of \S\ref{Y-C-1}, ``The matrices of \textup{(S-cen$_k$)}''. At $\mathsf n=60$ we have $\tau_1=8\gamma^{-2}(1+b_h)x^{-60}$ (\Cref{Y-C-L66}\,(c)) and $\tau_{\rm c}=165(L+1)\mathsf N_{\mathsf B}3L^3(2R+3)^3\tau_1$ (item~C.II~4 of \S\ref{Y-C-1}, ``Constants''), and $\tau_{\rm c}\le x^{-59}$ by \textup{(P1)}. The display before \Cref{Y-C-L71} gives \[\tau_2=24L^3(2R+1)^3\gamma^{-2}x^{-60}+(2/\gamma)3L^3\mathsf g(\theta_{\mathbb C}/2)x^{-64}\] and \[\mathsf{ti}_{\rm c}=(16+4L)\mathsf N_{\mathsf B}\tau_2+7\mathrm e(L+1)C_Dx^{-128}.\] These inputs rest on the following three conditions \textup{(Fr1)}, \textup{(Fr2)}, \textup{(Fr3)} of item~C.II~2.2 of \S\ref{Y-C-1}, all of them implied by the data.
\begin{enumerate}
\item[\textup{(Fr1)}] This is \Cref{Y-A-L7}\,(c).

\item[\textup{(Fr2)}] This is \textup{(Dec)}, also for the free kernels $\mathsf G,\mathsf G_Z$, which belong to the list of \textup{(Dec)} in \S\ref{Y-A-1}.

\item[\textup{(Fr3)}] This is \textup{(N$'$)} at $m=1$ (\Cref{Y-A-Np}).
\end{enumerate}

\item[\textup{(Rows)}] On the rows we take \Cref{Y-A-Th36}, \Cref{Y-A-Th33}\,(d), \Cref{Y-A-Th29}\,(a) and \Cref{Y-A-P35}, together with \Cref{Y-A-Pairs}, \Cref{Y-A-XT}, \Cref{Y-A-P4}\,(a) and \textup{(Dec)}.

\item[\textup{(C5)}] The present theorem takes from \S\ref{Y-D-1} the lift $E_3$ of item~C1 there, the pairs $(V',e)$ and $X'$ of item~E2 there, and the definitions of $T_\square$ and $w_j$ in items~E5 and~E6 there, which are also those of items~\textup{E7} and~\textup{E9} of this section. It also uses Step~1\,\textup{(iv)} (the identities $W^VE=1$ and $\mathrm{d}^VE=E_3\mathrm{d}'^V$) and Step~5 (Bianchi's identity) of the proof of \Cref{Y-D-C5}.

\item[\textup{(C6)}] The present theorem depends on five passages of \S\ref{Y-E-1}. The first is the identity of matrices in Step~4 of the proof of \Cref{Y-E-C6}, which writes out $\partial U(p)$ in the chart case by case. The second is Step~5\,\textup{(1)}, \textup{(2)} of that proof, a piece of geometry and an identity. The third is the sentence of item~3.4 of \S\ref{Y-E-1} stating that $\mathsf M(U)(c)$ is a Borel map of $U$. The fourth is the passage ``As to invariance'' of Step~1\,\textup{(3)} of that proof. The fifth consists of the inequalities $\varrho_L+2<N/2$, $\varrho(x)/L+2<N/2$ and $\varrho(x)+L<NL/2$ of Step~2 of that proof: the first two appear in the proof of statement \textup{(iv)} of that step, and the third follows from the stronger inequality $\varrho(x)+1<NL/2-L$ of statement \textup{(iv)} itself.

The first two passages are valid at every field, for every fine site $y$, every $\eta\in\mathfrak L_{\mathfrak g}$ and every $p$ with $|y-p|\le\varrho(x)$; at a pair $(V',e)$ we have $\partial V'(P)=e^{iX'_P}$ also where $\partial V'(P)=-1$, because $e^{i\pi e_P}=-1$. The third and the fourth are statements about every $G$-valued fine field $U$. The fifth rests on $x\ge L^2$, $R\ge6$, $L\ge3$ and \textup{(Tor)} alone, which hold here by item~G, \textup{(O4)} and item~C. None of the five depends on any of the floors \textup{(W1)} to \textup{(W5)} of \Cref{Y-E-C6}, on a condition that bounds $|\eta_b|$ from above, on $\chi_\square=1$, on a standing hypothesis of Steps~3 to~10 of its proof, or on anything of the present theorem.
\end{enumerate}

We also take results from earlier sections. Bianchi's identity appears in the proof of \Cref{Y-C-L67}\,(d), and the decay of an inverse and the quotient of two determinants are \Cref{Y-A-L3}\,(a) and~(b). The formula of the chart is given in \Cref{Y-C-L51}.

\begin{lemma}[Angles and logarithms in $\operatorname{SU}(2)$]\label{Y-F-L1} On $G=\operatorname{SU}(2)$, the angle of an element $g$ is fixed by $\cos\operatorname{ang}g=\frac12\operatorname{tr}g$ and $\operatorname{ang}g\in[0,\pi]$, and $|\mathsf L\,g|=\operatorname{ang}g$ for every $g$ with the one exception $\mathsf L(-1)=0$. For finitely many $Y_k\in\mathfrak g$ let $S=\sum|Y_k|$.

\begin{enumerate}\item[\textup{(a)}] For $g,h,u\in G$ and $Y\in\mathfrak g$, we have $\operatorname{ang}(gh)\le\operatorname{ang}g+\operatorname{ang}h$, also $\operatorname{ang}(g^{-1})=\operatorname{ang}(ugu^{-1})=\operatorname{ang}g$, and $\operatorname{ang}e^{iY}\le|Y|$.

\item[\textup{(b)}] For $g\in G$, $\|\operatorname{Ad}_g-1\|\le\min\{2,2\operatorname{ang}g\}$.

\item[\textup{(c)}] If $S<\pi$, then $|\mathsf L(\prod e^{iY_k})-\sum Y_k|\le S^2$.

\item[\textup{(c$'$)}] If $\prod e^{iY_k}=1$, then $|\sum Y_k|\le\frac12S^2$, whatever the value of $S$.\end{enumerate}\end{lemma}

\begin{proof}\emph{Step \textup{(a)}: the arc distance.} We view $\operatorname{SU}(2)$ as the unit sphere $S^3$. For $g,h\in G$, the number $\operatorname{ang}(g^{-1}h)$ is the arc distance between $g$ and $h$ on $S^3=\operatorname{SU}(2)$, and this metric is invariant under left and under right translations. Thus $\operatorname{ang}g$ is the distance from $1$ to $g$; the triangle inequality and the left invariance give $\operatorname{ang}(gh)\le\operatorname{ang}g+\operatorname{ang}h$, and the invariance on both sides gives $\operatorname{ang}(g^{-1})=\operatorname{ang}(ugu^{-1})=\operatorname{ang}g$. The curve $t\mapsto e^{itY}$, $0\le t\le1$, leads from $1$ to $e^{iY}$ with the speed $|Y|$; its length is therefore $|Y|$, and $\operatorname{ang}e^{iY}\le|Y|$.

\emph{Step \textup{(b)}: the adjoint action.} The map $\operatorname{Ad}_g$ is the rotation of $\mathfrak g=\mathbb R^3$ about the axis of $g$ by the angle $2\operatorname{ang}g$. A rotation by an angle $\alpha$ moves each unit vector by at most $2|\sin(\alpha/2)|$, with equality for the unit vectors orthogonal to its axis. Hence $\|\operatorname{Ad}_g-1\|=2\sin(\operatorname{ang}g)$, and this number is at most $2$ and at most $2\operatorname{ang}g$.

\emph{Step \textup{(c)}: the logarithm of a product.} Write $g$ for $\prod e^{iY_k}$. If $S\ge2$, then $|\mathsf L\,g|\le S$ by \textup{(a)}, as $|\mathsf L\,g|$ never exceeds $\operatorname{ang}g$; moreover $|\sum Y_k|\le S$, and the triangle inequality together with $2S\le S^2$ gives the claim. Now let $S<2$. We reach the product $\prod e^{iY_k}$ factor by factor, along a curve $g(s)$ of speed $1$ whose parameter $s$ is the arc length: the interval $[0,S]$ is cut into consecutive subintervals, one of length $|Y_k|$ for each nonzero $Y_k$ in order. On that subinterval put $Y(s)=Y_k/|Y_k|$. The path $g(s)$ is the solution of $g(0)=1$ and $g'=g\,iY(s)$ on $0\le s\le S$. So $g(S)=\prod e^{iY_k}$ and $\int_0^SY=\sum Y_k$. Put $Z=\mathsf L\,g(s)$ and $\check\theta=|Z|=\operatorname{ang}g(s)\le s$; indeed $g(s)$ is a product of elements $e^{iY}$ whose $|Y|$ add up to $s$, so that \textup{(a)} bounds its angle by $s$, and $s<2<\pi$ excludes $g(s)=-1$.

The identity $e^{-A}(e^A)'=((1-e^{-\operatorname{ad}A})/\operatorname{ad}A)A'$ at $A=iZ$ gives $Z'=f(\operatorname{ad}_{iZ})Y$ with $f(z)=z/(1-e^{-z})$. The operator $\operatorname{ad}_{iZ}$ is skew on $\mathfrak g$ and has the eigenvalues $0,\pm2i\,\check\theta$, and $f(0)=1$, $f(\pm2i\,\check\theta)=\check\theta\cot\check\theta\pm i\,\check\theta$. Since $|Y(s)|=1$, it follows that $|Z'-Y|\le m(\check\theta):=|\check\theta\cot\check\theta-1+i\,\check\theta|$. The function $1/t-\cot t$ is increasing on $(0,2]$, since its derivative is $1/\sin^2t-1/t^2\ge0$. At $t=2$ it is less than $0.958<\sqrt3$. Since $m(\check\theta)^2=(1-\check\theta\cot\check\theta)^2+\check\theta^2$, we obtain $m(\check\theta)\le2\,\check\theta$ for $0<\check\theta\le2$. Since $Z(0)=0$ and $\check\theta\le s<2$, integration over $[0,S]$ yields $|Z(S)-\int_0^SY|\le\int_0^S2s\,ds=S^2$. This is \textup{(c)}, because $Z(S)=\mathsf L(\prod e^{iY_k})$ and $\int_0^SY=\sum Y_k$.

\emph{Step \textup{(c$'$)}: a product equal to $1$.} If $S\ge2$, then $|\sum Y_k|\le S\le S^2/2$. Let $S<2$, and define $g(s)$, $Z$ and $\check\theta$ as in Step~\textup{(c)}. Here $g(S)=1$; so $g(s)^{-1}=g(s)^{-1}g(S)$ is a product of elements $e^{iY}$ whose $|Y|$ add up to $S-s$, and \textup{(a)} gives $\check\theta(s)\le\min\{s,S-s\}$. Also $Z(S)=0$, and $Z(0)=0$. The bound $|Z'-Y|\le m(\check\theta)\le2\,\check\theta$, proved in Step~\textup{(c)} for $\check\theta\le2$, now gives $|\sum Y_k|=|Z(S)-\int_0^SY|\le\int_0^S2\min\{s,S-s\}ds=S^2/2$.\end{proof}

\begin{lemma}[A comb gauge on a box]\label{Y-F-L3} Let $\mathrm{Bx}$ be an embedded box with a site $o$, let $V$ be a $G$-valued field on the bonds of $\mathrm{Bx}$, and let $\phi_p=\operatorname{ang}\partial V(p)$. Let $\Gamma$ be the coordinate path, taken in the order $e_1,e_2,e_3$; put $\check\gamma(z)=V(\Gamma(o,z))$ and $V^{\check\gamma}(b)=\check\gamma(b_-)V(b)\check\gamma(b_+)^{-1}$.

\begin{enumerate}\item[\textup{(1)}] Let $b=(z,\cdot)$ be a bond. Then $\operatorname{ang}V^{\check\gamma}(b)\le\sum_p\phi_p$, where $p$ ranges over the plaquettes of the ladders of $b$, which are as follows: for the direction $3$ there is none; for the direction $2$ there is one ladder of $|z_3-o_3|$ plaquettes $\{23\}$; for the direction $1$ there are two, one ladder of $|z_3-o_3|$ plaquettes $\{13\}$ and one of $|z_2-o_2|$ plaquettes $\{12\}$ at the height $o_3$.

\item[\textup{(2)}] Distinct bonds of one straight stretch have disjoint ladders. Hence, along a straight stretch, $\operatorname{ang}V^{\check\gamma}\le\sum_{p\in\mathrm{Bx}}\phi_p$.

\item[\textup{(3)}] Along a closed path in $\mathrm{Bx}$ made of $k$ straight stretches, $\operatorname{ang}V\le k\sum_{p\in\mathrm{Bx}}\phi_p$.\end{enumerate}\end{lemma}

\begin{proof}\emph{Step \textup{(1)}: the closed path of a bond.} The path $\Gamma(o,b_-)+b-\Gamma(o,b_+)$ is closed, and it bounds the ladders of $b$ named in \textup{(1)}; see Step~1 of the proof of \Cref{Y-A-L2}. By the definition of $\check\gamma$, the element $V^{\check\gamma}(b)$ is the product of $V$ along this closed path. Along a closed path that bounds ladders, the product of $V$ is an ordered product of conjugates of the $\partial V(p)^{\pm1}$, in which every plaquette of the ladders enters once; see the proof of \Cref{Y-C-L67}\,(a), the passage ``a ladder''. So \Cref{Y-F-L1}\,(a) bounds $\operatorname{ang}V^{\check\gamma}(b)$ by the sum of the $\phi_p$ over these plaquettes.

\emph{Step \textup{(2)}: disjoint ladders.} Two bonds of a straight stretch have different values of the coordinate in the direction of the stretch, and the plaquettes of their ladders differ in that coordinate too; moreover a plaquette $\{13\}$ is never a plaquette $\{12\}$. So the ladders of distinct bonds of the stretch are disjoint. The product of $V^{\check\gamma}$ along the stretch is the product of the $V^{\check\gamma}(b)$ over its bonds; \Cref{Y-F-L1}\,(a) bounds its angle by the sum of their angles, and \textup{(1)} with the disjointness bounds that sum by $\sum_{p\in\mathrm{Bx}}\phi_p$.

\emph{Step \textup{(3)}: a closed path.} Let $\pi$ be a closed path in $\mathrm{Bx}$ made of $k$ straight stretches, so that $\pi_-=\pi_+$. The gauge factors cancel in pairs along $\pi$, and $V(\pi)=\check\gamma(\pi_-)^{-1}V^{\check\gamma}(\pi)\check\gamma(\pi_+)$; as the path is closed, this is a conjugation. Here $V^{\check\gamma}(\pi)$ is the product of the $V^{\check\gamma}$ of the $k$ stretches, and \Cref{Y-F-L1}\,(a) together with \textup{(2)} gives $\operatorname{ang}V(\pi)\le k\sum_{p\in\mathrm{Bx}}\phi_p$.\end{proof}

\begin{remark}In the chart, $U=1$ on the combs. Therefore, for the field $U$ of the chart, the box $\mathrm{Bx}=B(y)$ and $o=y$, one has $\check\gamma=1$.\end{remark}

\begin{lemma}[Sizes of the vectors $\tau_{c,x}$ and $\bar\tau_c$ in the chart]\label{Y-F-L2} In the chart, let $c=(y,y')$ be of direction $\mu$, and let $x\in B(y)$. Here $h_{c,x}$ is that of item~C.I~4 of \S\ref{Y-C-1}, and, as in item~3.4 of \S\ref{Y-E-1}, $\tau_{c,x}=\mathsf L(h_{c,x}h_{c,y}^{-1})$ and $\bar\tau_c=L^{-3}\sum_{z\in B(y)}\tau_{c,z}$. Then $|\tau_{c,x}|\le\operatorname{ang}(h_{c,x}h_{c,y}^{-1})\le s_c$; consequently $|\bar\tau_c|\le s_c$.\end{lemma}

\begin{proof}Let $v_0=y,\dots,v_m=x$ be the sites of the comb $\Gamma(y,x)$, in their order along it; for a site $v$ write $v'=v+Le_\mu$, and put \[a_k=U(\Gamma(y,v_k))U([v_k,v_k'])U(\Gamma(y',v_k'))^{-1}.\] Then $a_0=h_{c,y}$ and $a_m=h_{c,x}$. If the bond $(v_k,v_{k+1})$ is parallel to $e_\mu$, then this bond followed by $[v_{k+1},v_{k+1}']$, and $[v_k,v_k']$ followed by the bond $(v_k',v_{k+1}')$, are two paths from $v_k$ to $v_{k+1}'$ on one straight line of direction $e_\mu$; since steps forth and back along this line cancel, the products of $U$ along the two paths agree, and $a_{k+1}=a_k$. If it is not, then $a_{k+1}a_k^{-1}$ is conjugate to the product of $U$ along the boundary of a strip of $L$ plaquettes, and that product is an ordered product of $L$ conjugates of $\partial U(p)^{\pm1}$. No two of these strips meet, and all of them lie in $W_c$; see Step~2 of the proof of \Cref{Y-A-L2}. As $h_{c,x}h_{c,y}^{-1}=a_ma_0^{-1}$ is the ordered product of the $a_{k+1}a_k^{-1}$, \Cref{Y-F-L1}\,(a) bounds $\operatorname{ang}(h_{c,x}h_{c,y}^{-1})$ by the sum of the $\operatorname{ang}\partial U(p)$ over $p\in W_c$, which is $s_c$. Finally, $|\mathsf L\,g|\le\operatorname{ang}g$ holds at every $g$, also at $g=-1$; this gives $|\tau_{c,x}|\le\operatorname{ang}(h_{c,x}h_{c,y}^{-1})$, and with it the bound for $|\bar\tau_c|$.\end{proof}

\emph{Item E: objects.}

\begin{enumerate}\item[\textup{E1}] \emph{Lattice, radii at $x$, norms, cell and fine cube of an object.} We take from item~C.I~1.1 of \S\ref{Y-C-1}: the blocks $B(y)$; $\pi(i)$ and $y(i)$; the distance ${\rm dist}(i,j)=L^{-1}|y(i)-y(j)|_\infty$, in coarse units; the combs; the set $I^\circ$ of the fine bonds off the combs, of which at most $3L^3$ are based in one block; $d$, $Q_1$ and $E$; the coarse path $\Gamma(y,y'')$ and $T(y,y'')$; $K^V$ and $(K^V)^*$; $\ell^V$; $\mathsf B^V=1+\ell^{V*}\ell^V$; $q=q^V_\ell$; $\mathsf g^V=\Pi_{\mathfrak L}\mathrm{d}^{V*}\mathsf B^VE$; and $|y-p|=|y-\pi(p)|_\infty$, in fine units.

From \Cref{Y-A-Lattices} and \Cref{Y-A-Facts} we take $W=Q_2$ with the weights $W(P,p)$; the double block $W_c$ of $c$; the union $W_P$ of the four blocks at the corners of $P$, which is a box $2L\times2L\times L$ holding at most $12L^3$ plaquettes; and $\mathrm{d}'$. The cubes, the cells $\square$, $\hat\square$, $\mathcal P$, $|X|_c$ and $X^{+j}$ are those of the definition ``Cubes'' of item~1 of \S\ref{Y-G0-1}.

The radii at $x$ are $R=R(x)=\varrho(x)=\lceil c_R\log x\rceil$, $r(x)=\lceil20c_R\log x/M\rceil$ and $s(x)=8r(x)$. We let $S_\varepsilon(t)=t^2/2$ for $t\le\varepsilon$, and $d_0\varepsilon t$ for $t>\varepsilon$. The norms $\|K\|_\theta,\|K\|_0,\|K\|_\infty$ are those of \Cref{Y-A-Norm}.

An object $i$ belongs to the cell containing $y(i)$ and to the fine cube in which its base site lies. For a set $A$ of cells, $\hat A$ denotes the set of the fine objects that belong to a cell of $A$ (the objects $i$ with $y(i)$ in a cell of $A$), and $\langle,\rangle_A$, $\|\cdot\|_A$ are defined by summing over the objects (rows) of $A$. The quantity $|\xi_A|^2$ is the sum of the $|\xi_b|^2$ over the bonds $b$ of $I^\circ$ in $\hat A$, and $|\varphi_A|^2$ is the same sum with $(G\xi)_b$ in place of $\xi_b$. The distance between two cells, or between two cubes, is the sup distance measured in cells, or in cubes. We say that $V'$ is tame on the box of radius $\rho$ about $i$ if $\phi'_P\le\varepsilon'$ for all $P$ with ${\rm dist}(P,i)\le\rho$.

\item[\textup{E2}] \emph{Fields, masks, energy, noise, defects.} As in \textup{E2} of \S\ref{Y-D-1}, $X'_P=X'_P(V',e)=\mathsf L(\partial V'(P))+\pi e_P1\{\partial V'(P)=-1\}$ and $\phi'_P=|X'_P|=\operatorname{ang}\partial V'(P)$. Let $D'=\{\phi'_P>\varepsilon'\}$ and $\zeta'_P=X'_P\min\{1,\varepsilon'/\phi'_P\}$, and let $\ell'^{\rm c}$ denote the cut of $\ell'$ at $\varrho(x')$. We put $m'_{y'}=1$ exactly when no plaquette of $D'$ is within $\rho_0$ of $y'$ and none with $\phi'_P>1$ is within $\varrho(x')+1$ of it, and $m'_{y'}=0$ otherwise.

The chart $\mathsf U(\eta;V')$ is that of item~C.I~4 of \S\ref{Y-C-1}, and there $\mathsf j(\eta_b)\le1$. For every $G$-valued fine field $U$ we let $F_p=\mathsf L(\partial U(p))$, $\phi_p=\operatorname{ang}\partial U(p)$, $D=D(U)=\{\phi_p>\varepsilon\}$ and $\zeta_p=F_p\min\{1,\varepsilon/\phi_p\}$, with multiplier $1$ at $\phi_p=0$. Below, we also define $\ell^U_y\zeta$, $m_y(U)$, $\mathrm{En}_\ell(U)$ and the noise $n_c$, $\mathsf n_c$ at every such $U$. Where $\eta$ is given, namely in the defects of \textup{E2}, in \textup{E4}, \textup{E11} and \textup{E12}, and in \textup{(a)} to \textup{(e)} of item~H, all these functions of $U$, when written without argument, are taken at $U=\mathsf U(\eta;V')$ for the $\eta$ of \textup{E3}. As in Step~5\,\textup{(1)} of the proof of \Cref{Y-E-C6}, for fine sites $z,z'$ the coordinate path $\Gamma(z,z')$ runs first along $e_1$, then along $e_2$ and last along $e_3$, always by the shorter way. We set $\ell^U_y\zeta=\sum_p\ell_y(p)\otimes\operatorname{Ad}_{U(\Gamma(y,\pi(p)))}\zeta_p$, and $m_y(U)=1$ exactly when there is no $p$ of $D$ with $|y-p|\le\rho_0$ and no $p$ with $\phi_p>1$ and $|y-p|\le\varrho(x)+1$, while $m_y(U)=0$ otherwise. The energy is \[\mathrm{En}_\ell(U):=\sum_pS_\varepsilon(\phi_p)+\frac12\sum_ym_y(U)|\ell^U_y\zeta|^2.\]

The noise is given by $\mathsf n_c=\operatorname{ang}(V'(c)\mathsf M(U)(c)^{-1})\in[0,\pi]$ and $n_c=\mathsf L(V'(c)\mathsf M(U)(c)^{-1})\in\mathfrak g$, where $|n_c|=\mathsf n_c$ if $\mathsf n_c<\pi$. The defects are $\mathrm{de}_P=X'_P-(W^VF)_P-(\mathrm{d}'^Vn)_P$ and $\mathrm{db}_q=(\mathrm{d}^VF)_q$. We put $\Sigma_P=\sum_{p\in W_P}\phi_p$, $\mathrm{na}_P=\sum_{c\in dP}\mathsf n_c$, $s_c=\sum_{p\in W_c}\phi_p$ and $s_B=\sum_{p\in B}\phi_p$; a superscript $D$ on $\Sigma_P$, $s_c$ or $s_B$ restricts the sum to the rough $p$, and a superscript $t$ restricts it to the tame $p$, where a fine plaquette $p$ is rough if $p\in D(U)$ and tame if $p\notin D(U)$.

\item[\textup{E3}] \emph{White noise and centre.} Let $\xi\sim \mathcal N(0,x^{-1}1)$ on $\mathfrak L_{\mathfrak g}$, and put $\varphi=G^V\xi$, $a_0=C\mathsf g^VX'$ and $\eta=\varphi-a_0$; further, $u^T=1-qC$, $F^\flat=EX'+\mathrm{d}^V\eta$ and $n^\flat=-Q_1^V\eta$.

\item[\textup{E4}] \emph{The space of triples.} The space $\mathrm{Hil}$ consists of the triples $(F,\Sigma,n)$ with $F$ on the fine plaquettes, $\Sigma$ on the fine sites and $n$ on the coarse bonds, which together are its rows; it carries $\langle,\rangle_{\mathrm{Hil}}=\langle F,F'\rangle+\langle\Sigma,\Sigma'\rangle+a\langle n,n'\rangle$. Let $\iota(F,n)=(F,\ell^VF,n)$ and $k\psi=\iota(\mathrm{d}^V\psi,-Q_1^V\psi)$; thus $q^V_\ell=\|k\cdot\|^2$, and $k_i$ denotes the row $i$ of $k$. We put $\bar F=\iota(EX',0)-ka_0$, and $\mathfrak k=\iota(k_F,k_n)$ with $k_F=F^\flat-F$ and $k_n=n^\flat-n$. We set $d_1=W^V\mathrm{d}^V-\mathrm{d}'^VQ_1^V$ and $d_2=\mathrm{d}^V\mathrm{d}^V$.

\item[\textup{E5}] \emph{Free kernels and their cuts.} The kernels $\mathsf T$, $\mathsf X$, $\mathsf B=1+\ell^{\mathsf T}\ell$, $\mathsf B'$, $H^\sharp$, $\mathsf N^\sharp$ are those of \Cref{Y-A-Pairs}, of \Cref{Y-A-XT}, and of \Cref{Y-A-P4}\,(a). We write $\mathsf T^<,\mathsf X^<$ for the cuts of $\mathsf T$ and $\mathsf X$ at $R(x')$, and $\mathsf B'^{\rm c}=1+\ell'^{{\rm c}\mathsf T}\ell'^{\rm c}$. With $\|\cdot\|_{\rm r}$ the largest absolute row sum, \[\mathsf t_\rho:=\|\ell'^{\rm c}\|_{\rm r}\sup_{P_0}\sum_{y':{\rm dist}(y',P_0)>\rho_0}|\ell'^{\rm c}_{y'}(P_0)|,\] as in \Cref{Y-A-CutRows}.

\item[\textup{E6}] \emph{Multipliers and the pair map.} We put $v_*=(1+\ell'^{{\rm c}V*}m'\ell'^{{\rm c}V})\zeta'$, $Z_*=v_*/L$, $\Xi_*=(\mathsf X^<)^Vv_*$ and $Y_*=W^{V*}Z_*+\mathrm{d}^{V*}\Xi_*$; further, $\nu^{\rm ref}=(\mathsf T^<)^Vv_*$, $\sigma_*=\ell^V\nu^{\rm ref}$, $\nu_*=Y_*-\ell^{V*}\sigma_*$ and $\omega_*=(aL)^{-1}\mathrm{d}'^{V*}v_*$. The pair map is \[\mathrm{pm}(v)=(W^{V*}v/L+\mathrm{d}^{V*}(\mathsf X^<)^Vv-\ell^{V*}\ell^V(\mathsf T^<)^Vv,\ \ell^V(\mathsf T^<)^Vv,\ (aL)^{-1}\mathrm{d}'^{V*}v),\] and we put $\boldsymbol\nu=\mathrm{pm}(v_*)=(\nu_*,\sigma_*,\omega_*)$ and $\hat R^V=L\,\mathrm{pm}^*\mathrm{pm}$. Let $\Delta=\bar F-\boldsymbol\nu$, with the components $\Delta_p,\Delta_y,\Delta_c$, and let $\Delta_Y=\mathsf B^V\bar F_F-Y_*$, where $\bar F_F$ is the plaquette component of $\bar F$. With $r_*=\mathsf B^V\nu^{\rm ref}-Y_*$ we then have $\nu_*=\nu^{\rm ref}-r_*$.

\item[\textup{E7}] \emph{Coarse energy and the prefactor of a cell.} For a coarse plaquette $P$ let $\mathrm{En}'_P=S_{\varepsilon'}(\phi'_P)+\frac12\langle\zeta'_P,(\ell'^{{\rm c}V*}m'\ell'^{{\rm c}V}\zeta')_P\rangle$. The coarse energy is \[\FitDisplay{\dimexpr\linewidth-1pt\relax}{\mathrm{En}'(V'):=\sum_P\mathrm{En}'_P=\sum_PS_{\varepsilon'}(\phi'_P)+\frac12\sum_{y'}m'_{y'}|\ell'^{{\rm c}V}_{y'}\zeta'|^2,}\] and for a cell $\square$ we put \[\FitDisplay{\dimexpr\linewidth-1pt\relax}{\begin{aligned}
&T_\square:=\sum_{P\in\square}[(x'/2)\langle X'_P,v_{*P}\rangle-x'\mathrm{En}'_P+(x/2)\langle X'_P,(E^*\Delta_Y)_P\rangle]
\\
&\quad{}+(x/2)\sum_{q'\in\square}\langle(\mathrm{d}'^VX')_{q'},(E_3^*\Xi_*)_{q'}\rangle
\\
&\quad{}-(x/2)\sum_{b\in\hat\square}\langle a_{0,b},(u^T\,\mathsf g^VX')_b\rangle.
\end{aligned}}\]

This $T_\square$ coincides, term by term, with the $T_\square$ of \textup{E5} of \S\ref{Y-D-1}, where the $\bar F$ of that section is the $\bar F_F$ of \textup{E6}, the plaquette component of the triple $\bar F$ of \textup{E4}. The $\nu^{\rm ref}$ and $\nu_*$ of \textup{E4} of \S\ref{Y-D-1} are the same objects as in \textup{E6}, given by the same formulas $\nu^{\rm ref}=(\mathsf T^<)^Vv_*$ and $\nu_*=Y_*-\ell^{V*}\sigma_*$; hence the $\Delta_F=\bar F-\nu_*$ of \S\ref{Y-D-1} is the $\Delta_p$ of \textup{E6}. By \textup{E6}, $\mathsf B^V(\bar F_F-\nu^{\rm ref})+r_*=\Delta_Y$; expressed in the components of $\Delta$ from \textup{E6}, $\Delta_Y=(\Delta_p)_p+\ell^{V*}(\Delta_y)_y$.

\item[\textup{E8}] \emph{Radii.} In coarse units, $\varrho_L=\lceil R/L\rceil$, $r_w=\lfloor R/2\rfloor$, $r_q=2\varrho_L+2$, $R_2=2R+2$, $R_3=2R+\varrho_L+5$, $R_5=4R+4\varrho_L+8$ and $R_6=R_5+2\varrho_L+3$. In cells, $r_l=\lceil(\varrho_L+2)/M\rceil$, which is $\ge1$, $r_1=\lceil R_2/M\rceil+3r_l+5$ and $s_0=2r(x')$.

\item[\textup{E9}] \emph{Marks, forced sets, classes.} Let $\mathrm{mk}'$ be the set of the cells that hold a plaquette of $D'$, and $f=\mathrm{mk}'^{+s_0}$. The classes $\mathrm{mk}'_j$ are those of $\mathrm{mk}'$ under the relation generated by distance $\le2s_0+4r_1+2$. We put $f_j=\mathrm{mk}'^{+s_0}_j$; any two of these sets are at a distance $>4r_1+2$. The set $D'_j$ consists of the plaquettes of $D'$ in $\mathrm{mk}'_j$, and $f_J,D'_J$ are the unions over $j\in J$. Further, $w_j=\sum_{\square\in\mathrm{mk}'^{+r(x')}_j}T_\square$.

\item[\textup{E10}] \emph{Tests.} For a cell $\square$, $b_\square(\xi)=1$ if $|\xi_b|>c_\xi\varepsilon$ at some bond $b$ of $I^\circ$ in $\hat\square$, and $b_\square(\xi)=0$ otherwise.

\item[\textup{E11}] \emph{The factor of a set of cells.} For a set $\mathfrak Z$ of cells, $[A]^\ell_{\mathfrak Z}$ is the product of $\exp\{(x/2)|F^\flat_p|^2-xS_\varepsilon(\phi_p)\}$ over $p\in\hat{\mathfrak Z}$, of $\exp\{(x/2)|(\ell^VF^\flat)_y|^2-(x/2)m_y|\ell^U_y\zeta|^2\}$ over $y\in\hat{\mathfrak Z}$, of $\exp\{(ax/2)|n^\flat_c|^2-ax(1-\cos\mathsf n_c)\}$ over $c\in\mathfrak Z$, and of $1\{|\eta_b|<\pi\}\mathsf j(\eta_b)$ over $b\in I^\circ\cap\hat{\mathfrak Z}$. We put $A_\square:=[A]^\ell_{\{\square\}}$; since every object lies in one cell, $[A]^\ell_{\mathfrak Z}=\prod_{\square\in\mathfrak Z}A_\square$.

\item[\textup{E12}] \emph{Cell data and letters at $V'$.} A cell datum $(J,t)$ consists of a set $J$ of classes and a finite set $t$ of cells off $f_J$, such that $\mathfrak b:=f_J\cup t$ is non-empty; here $\mathfrak b$ is any such set, connected or disconnected. The sets of the datum are $\mathfrak Z=\mathfrak b^{+r_1}$, $\mathfrak Z'=\mathfrak b^{+r_1-r_l}$, $\mathfrak S=\mathfrak b^{+r_1-r_l-1}$, $\mathfrak S^-=\mathfrak b^{+r_1-r_l-2}$, $\mathfrak Z_2=\mathfrak b^{+2r_1}$ and the rim $\mathrm{Lr}=\mathfrak b^{+r_1+r_l+1}\setminus\mathfrak b^{+r_1-2r_l-3}$. We write $\chi$ for the restriction to the objects of $\mathfrak Z$, and $\chi_S$ for the restriction to the bonds based in $\hat{\mathfrak S}$.

We put \[\mathcal G_{\mathfrak Z}=\sum_{P\in\mathfrak Z}[\langle v_*,X'\rangle_P-\frac12\langle v_*,\hat R^Vv_*\rangle_P]\] and \[\hat\lambda_{(J,t)}:=e^{-\sum_Jw_j}\prod_tb_\square\prod_{\mathfrak Z_2\setminus\mathfrak b}(1-b_\square)[A]^\ell_{\mathfrak Z}\ge0.\] A cell datum is admitted if \textup{(H17)} holds: each plaquette of $D'$ that lies in a cell at distance at most $4r_1$ from $\mathfrak b$ lies in $D'_j$ for some $j\in J$. Two cell data are compatible if each cell of the one set $\mathfrak b$ is at distance more than $4r_1+2$ from each cell of the other.

In the rest of \textup{E12}, an object is a forced class $f_j$, whose cells are those of $f_j$, or a cell, or a polymer $X\in\mathcal P$. Letters are made of these objects, which are distinct from the objects $i$ of \textup{E1}. Two objects with cells are adjacent if they have cells at distance $\le4r_1+2$; a polymer $X$ and an object with a cell $\square$ are adjacent if and only if $X$ contains a cube of $\widehat{\square^{+r_1}}$. A letter $\check\zeta=(J,t,\mathfrak X)$ is a set of objects that is connected under adjacency and in which $(J,t)$ is a cell datum, each $X$ is adjacent to some object with cells, and any two $X$ have distance $>s(x)$. The letter is bare if $\mathfrak X$ is empty.

Two letters are compatible if each object of the first is different from, and not adjacent to, every object of the second, and if any polymer of the first and any polymer of the second have distance $>s(x)$; the cell data of two compatible letters are then compatible. A letter is admitted if its cell datum is. We put $1_{\rm ind}=1\{D(U)\cap\hat{\mathfrak Z}$ is contained in the union of $\mathfrak X$, and every $X$ of $\mathfrak X$ holds a plaquette of $D(U)\cap\hat{\mathfrak Z}\}$, as well as $\Lambda_{\check\zeta}=\hat\lambda_{(J,t)}1_{\rm ind}$ and $\lambda^o_{\check\zeta}=\Lambda_{\check\zeta}\prod_{\mathfrak X}w(X,U)$. We set $I(\check\zeta)=\mathfrak Z_2\cup\{\text{cells that meet some }X^{+s(x)}\}$.\end{enumerate}

\emph{Item F: constants.} We set $a_\infty^\flat=\mathsf G_0^2b_h/4$ and $c_q^\flat=\pi b_h\mathsf G_0(1+a_\infty^\flat)$. Further, we put $K_D=\mathsf K_{\rm c}+1+6700\theta_D^{-2}C_D^3$, $K_r=1+23000\theta_D^{-2}(C_D^5+LC_D^2)$ and $K_\Delta^\flat=(1+l_r)(K_D+K_r)$, and $K_{\rm q}^\flat=M^3[1.5L^2(\mathsf N_{\mathsf B}K_D+K_r)+9\mathsf N_\Xi+1]$. Next, $a_\nu=\mathrm eL+6\mathrm eC_D+C_D^3$ and $C_{\rm hol}^\flat=162\mathsf g(\theta_D)(1+C_D^2)[La_\nu^2+LC_D^4+7]$; then $c_B=6\mathsf g(\theta_D/(4C_D^2))$ and $e_1^\flat=\frac12\mathsf N_Z^2c_BC_{\rm hol}^\flat\varepsilon'$. We also set $\mathsf{ti}_0^\flat:=8C_D^2La_\nu^2x'^{-64}$ and $c_\eta^\flat=1+L^{1/2}(1+\mathsf a_{\rm c})$.

With these, \[K_s^\flat=200vL^3(1+a)(1+K_\Delta^\flat)(\mathsf N_Z+\mathsf N_\Xi+\mathsf N_\nu^\flat+\mathsf N_\sigma)(1+l_r)(1+K\mathsf g_1)(1+c_\eta^\flat)^2\] and \[\sigma_{\rm tot}^\flat=K_s^\flat p_0^3x^{-1/2}+K_{\rm q}^\flat p_0'^2\varepsilon'+(9e_1^\flat+\frac92\mathsf N_Z^2c_B\mathsf{ti}_0^\flat)M^3p_0'^2.\]

Let $t_0=\varepsilon^{1/2}$ and, in this section alone, $\theta_2=\varepsilon/12$; this $\theta_2$ differs from the $\theta_2=\theta_1/2$ of item~3 of \S\ref{Y-G0-1}. We put \[\vartheta_F^\flat=[(86400\mathsf N_Z+110\mathsf N_\Xi)L^{5/2}+2949K\mathsf g_1\mathsf N_\sigma L^{9/2}]\varepsilon\] and \[\vartheta_n^\flat=[(192\mathsf N_Z+\mathsf N_\Xi)L^{-1/2}+3K\mathsf g_1\mathsf N_\sigma L^{1/2}]\varepsilon,\] and further, $z_b=21\mathsf N_ZL^{-1/2}\varepsilon$, $r_o=x^{-3/8}$ and $\check\theta_x=x^{-1/4}$.

Let $\mathrm{py}=\min\{d_0p_0^2/2,0.05ax^{1/4}\}$ and $\kappa_1=\mathrm{py}/(240\pi^2L^3x)$; $h_*=C_\varphi^\flat(4+12K\mathsf g_*+a^{1/2}L)$, $t_Q=1/(12h_*^2)$ and $\rho_Q=3c_\xi^2p_0^2/(128h_*^2)$; $c_{\rm tr}=(4/\gamma)144K^2\mathsf g(2\lambda)$ and $c_\tau=\max\{6.1,c_{\rm tr}\}/12$. Moreover \[Q_\theta^\flat=[1.3\cdot10^5L^3+4704aL^4+(8.4\cdot10^4l_r^2+1296K^2\mathsf g_1^2)L^3]M^3\] and $u_b=2.5v(c_\xi p_0+(c_\xi p_0)^{-1})e^{-c_\xi^2p_0^2/2}$.

We put $u_t=3v\max\{e^{c_\tau-\rho_Q/2},e^{-\mathrm{py}/(768(2r_1+1)^3)}\}$, $u_1=\max\{e^{-\mathrm{py}/768},e^{-\rho_Q/2}\}$ and $\delta_3=\min\{d_0/1536,3c_\xi^2/(256h_*^2)\}$. The constants $u_t',u_1',\delta_3'$ are given by the same formulas, with $1152$ in place of $768$ and $2304$ in place of $1536$, and in them, as in $e^{\circ\prime}$ and $e^{o\prime}$ below, the prime is not that of \textup{(O8)}, and $p_0$, $\mathrm{py}$, $\rho_Q$, $r_1$ are those at $x$.

Then \[e_2^\flat=9e_1^\flat+\frac92\mathsf N_Z^2c_B\mathsf{ti}_0^\flat+6.3(\mathsf N_\nu^\flat)^2\vartheta_F^\flat+3\mathsf N_Z^2(\theta_2+L\vartheta_n^\flat)\] and \[\mathrm{Res}_{(J,t)}=\exp\{-(\frac12-d_0-\mathsf t_\rho-3e_1^\flat)x'\varepsilon'\sum_J\sum_{D'_j}\phi'_P\}.\] For a cell datum $(J,t)$ let $e^\circ_{(J,t)}=\mathrm{Res}\,u_t^{\#t}u_1^{1[t\ne\emptyset]}$ and $e^{\circ\prime}_{(J,t)}=\mathrm{Res}\,u_t'^{\#t}u_1'^{1[t\ne\emptyset]}$; for a letter $\check\zeta$ let \[e^{o}_{\check\zeta}=e^\circ_{(J,t)}\prod_{\mathfrak X}Ae^{-c_\Xi p_0^2-\kappa(|X|_c-1)}\] and \[e^{o\prime}_{\check\zeta}=e^{\circ\prime}_{(J,t)}\prod_{\mathfrak X}Ae^{-\mathrm{py}/1152-\kappa(|X|_c-1)}.\]

\emph{Item G: floors.} Let $S_{\rm fl}$ denote the set of all $y\ge L^2\mathrm e$ with the property that every floor in the list below holds at each $x\ge y$. If $S_{\rm fl}$ contains a number $y$, then it contains every number larger than $y$, so it is an upper half line; by Step~13 of the proof of \Cref{Y-F-C7} it is not empty.

The set $S_{\rm fl}$ need not contain its infimum. Indeed, $R$, $\varrho_L$, $r(x')$ and $r_1$ are values that the ceiling function $\lceil\cdot\rceil$ takes; as functions of $x$ they are continuous from the left, and their jumps are all upward. A floor with an inequality in which $R$, $\varrho_L$, $r(x')$ or $r_1$ appears on the side asserted to be the larger, as in \textup{(G1$^\flat$)}, can therefore fail at a point $x_0$ and yet hold at every $x>x_0$.

We put $x_{\rm fl}:=1+\inf S_{\rm fl}$. As $S_{\rm fl}$ is an upper half line and $x_{\rm fl}$ lies above its infimum, $x_{\rm fl}\in S_{\rm fl}$; the number $x_{\rm fl}$ depends only on what $S_{\rm fl}$ depends on. Since $x_{\rm fl}\in S_{\rm fl}$, the defining property of $S_{\rm fl}$ shows that all the floors hold at every $x\ge x_{\rm fl}$. This definition of $x_{\rm fl}$ has the same form as the definition of $x_\flat^0$ in \Cref{Y-B-C1}. Each floor is a condition on $x$ and on the letters of \textup{(O0)} to \textup{(O10)} alone. The floors are the following.

\begin{enumerate}\item[\textup{(G1$^\flat$)}] $R\ge\max\{M,6\}$, $Mr(x')\ge5R+4\varrho_L+12$ and $\rho_0\le R(x')$.

\item[\textup{(F$_{\rm law}$)}] $432L^3b_h^{5/2}\gamma^{-3/2}\mathsf g(\theta_{\mathbb C})(R+3)^3\le x^4$ (see \Cref{Y-C-CITh66}\,(k)).

\item[\textup{(P1)}] $(4/\gamma)x^{-60}\le1$ and $165(L+1)\mathsf N_{\mathsf B}\,3L^3(2R+3)^3\,8\gamma^{-2}(1+b_h)\le x$ (see the floors in item~C.II~4 of \S\ref{Y-C-1}).

\item[\textup{(P3)}] $L^2(1+a)(\mathsf{ti}_{\rm c}+8C_D^3x'^{-64})\le1$ (see the display before \Cref{Y-C-L71} in \S\ref{Y-C-1}, C.II).

Among the floors in item~C.II~4 of \S\ref{Y-C-1}, the floor \textup{(P0)} consists of $x\ge L^2$, $R\ge6$ and \textup{(Tor)}. The first two hold by what precedes, because $S_{\rm fl}$ lies above $L^2\mathrm e$ and $R\ge6$ is part of \textup{(G1$^\flat$)}; the third is a datum. The inequality $2R_3+1<N$, which that section derives from \textup{(P0)}, follows as well: $2R_3+1=4R+2\varrho_L+11$, and \[N\ge16Mr(x)\ge320c_R\log x>320(R-1)>4R+2\varrho_L+11,\] where we use \textup{(Tor)} with $s(x)=8r(x)$, the definition of $r(x)$, the inequality $R<c_R\log x+1$, and $R\ge6$ together with $\varrho_L\le R$. Since $N>4R+2\varrho_L+11$ and $\varrho(x)=R$, also $\varrho_L+2<N/2$ and $\varrho(x)/L+2<N/2$ (see Step~2\,(iv) of the proof of \Cref{Y-E-C6}); hence every row of item~C is a row as in 1.1 of \S\ref{Y-C-1}, C.I, and also a row as in 1.2 of C.II there.

\item[\textup{(FA)}] $\mathsf K_{\rm c}L^2\varepsilon'\le1$, $\varepsilon'\le1/40$, $\varepsilon\le\varepsilon'$, $x^{-59}\le\varepsilon'$ and $\mathsf{ti}_{\rm c}+[3C_D^3+(\mathsf N_{\mathsf B}+6)C_D\mathsf N_Z]x'^{-64}\le\varepsilon'$.

\item[\textup{(F0$^\flat$)}] $\mathsf{ti}_0^\flat\le x^{-5}$.

\item[\textup{(F1)}] $12L^3\varepsilon\le\pi/24$, $90L^3\varepsilon\le1$ and $2L^2c_2^*\varepsilon\le1$.

\item[\textup{(F3$^\flat$)}] $4\cdot10^5\mathsf N_Z^2\varepsilon\le d_0$ and $3.2\cdot10^4\mathsf N_Z^2L\varepsilon^2\le1$.

\item[\textup{(F4$^\flat$)}] $\vartheta_F^\flat\le1/16$ and $\vartheta_n^\flat\le a/20$.

\item[\textup{(F5)}] $\mathsf N_ZL^{1/2}\varepsilon^{1/2}\le0.01$.

\item[\textup{(Fm)}] $l_r\mathsf N_\sigma L^{-3/2}(2\varrho(x)+3)^3\varepsilon\le d_0/16$.

\item[\textup{(Fc)}] $127K\mathsf g_1\mathsf N_\sigma L^{3/2}\varepsilon\le d_0/16$.

\item[\textup{(F8)}] $8400Mr(x')\varepsilon'\le\gamma/b_h$, $800(Mr(x'))^2\varepsilon'\le1$, $\varepsilon'\le0.04$ and $x^{-4}\le\varepsilon'$.

\item[\textup{(F9$^\flat$)}] One has \[10^6v^2b_h^3\gamma^{-2}(1+a)(1+l_r)^2(1+La_\infty^\flat)^2(1+\mathsf G_0)^2(2R+4\varrho_L+5)^3\le x^8,\] and $\mathrm{py}\ge1$.

\item[\textup{(U0$^\flat$)}] $60L^3\varepsilon\le r_o\le t_0$, $(2/9)\mathsf N_ZL^{1/2}\varepsilon\le0.09r_o$, $4\vartheta_F^\flat\le\check\theta_x/2$, $\theta_2+4\vartheta_n^\flat/a\le\check\theta_x/2$, $6Lr_o<\pi$, $\kappa_1\le\gamma/2$, $\check\theta_x\le1/6$ and $p_0^2x^{-1/4}\le1$.

\item[\textup{(U1$^\flat$)}] $19L^3\log(2(b_h+1)/\kappa_1)+(c_{\mathrm{tr}}/2)L^3(2\varrho_L+5)^3\le\mathrm{py}/192$.

\item[\textup{(Z0$^\flat$)}] $c_\eta^\flat\varepsilon\le r_o$, $25(c_\eta^\flat)^2\varepsilon\le0.3$, $h_*(1+a^{1/2})[K_\Delta^\flat\varepsilon'^2+8\theta_2\mathsf N_Z(aL)^{-1}\varepsilon']\le c_\xi\varepsilon/4$ and $c_\xi p_0\ge1$.

\item[\textup{(Z1$^\flat$)}] One has \[\FitDisplay{\dimexpr\linewidth-1pt\relax}{\begin{aligned}
&\sigma_{\rm tot}^\flat+5K_s^\flat p_0^3x^{-1/2}
\\
&\quad{}+[6.3(\mathsf N_\nu^\flat)^2\vartheta_F^\flat+3\mathsf N_Z^2(\theta_2+L\vartheta_n^\flat)]M^3p_0'^2
\\
&\quad{}+(Q_\theta^\flat+60c_\tau v)x^{-1/4}
\\
&\quad{}+27M^3e^{-\mathrm{py}/384}+3v((c_q^\flat)^2+c_\xi^2p_0^2+2)x^{-1}+x^{-30}
\\
&\quad{}\le x^{-1/8}.
\end{aligned}}\]

\item[\textup{(W)}] $0.05ax^{1/4}\ge d_0p_0^2/2$.\end{enumerate}

\emph{Item H: conclusions.} Let $n\ge0$. Let $(J_i,t_i)$, $i=1,\dots,n$, be pairwise compatible admitted cell data, or, in the clauses on letters, let $\check\zeta_i$, $i=1,\dots,n$, be pairwise compatible admitted letters; $\mathfrak b_i$ denotes the set $\mathfrak b$ of $(J_i,t_i)$, or of $\check\zeta_i$. Let $J_0$ be a finite set of cells off $\bigcup_i\mathfrak b_i^{+2r_1}$, let $\mathsf Y$ be a finite set of cells, and let $0\le\varsigma_1\le x^{-2}$ and $0\le\varsigma\le(\gamma/2)x^{-2}$. Then the following hold.

\begin{enumerate}\item[\textup{(a)}] For each $i$, the functions $\hat\lambda_i$ and $\Lambda_i=\hat\lambda_i1_{\mathrm{ind},i}$ depend only on the values of $\xi$ at the bonds of $I^\circ$ under $\mathfrak b_i^{+2r_1}$; here the product $[A]^\ell_{\mathfrak Z}1_{\mathrm{ind}}$ depends on $\eta$ only at bonds within the coarse distance $\varrho_L+2$ of $\hat{\mathfrak Z}$. The sets of bonds obtained in this way for the different $i$, and the bonds of $\hat J_0$, form pairwise disjoint sets. With $\mathsf Y_i:=\mathsf Y\cap\mathfrak b_i^{+2r_1}$ and $\mathsf Y_r:=\mathsf Y\setminus(\bigcup_i\mathfrak b_i^{+2r_1}\cup J_0)$, the following identity holds exactly, where for each $i$ the bracket $(1_{\mathrm{ind},i})$ is a factor that is present on both sides or absent from both: \[\FitDisplay{\dimexpr\linewidth-1pt\relax}{E_\xi[\prod_i\hat\lambda_i(1_{\mathrm{ind},i})\prod_{J_0}b_\square\,e^{(x/2)\varsigma_1|\xi_{\mathsf Y}|^2}]=\prod_iE_\xi[\hat\lambda_i(1_{\mathrm{ind},i})e^{(x/2)\varsigma_1|\xi_{\mathsf Y_i}|^2}]\cdot\prod_{\square\in J_0}E_\xi[b_\square e^{(x/2)\varsigma_1|\xi_{\mathsf Y\cap\{\square\}}|^2}]\cdot E_\xi[e^{(x/2)\varsigma_1|\xi_{\mathsf Y_r}|^2}].}\] Pointwise $|\lambda^o|\le\Lambda\prod_{\mathfrak X}(\text{cap})$, and wherever $1_{\mathrm{ind}}=1$ the strong cap is at most the constant $Ae^{-c_\Xi p_0^2-\kappa(|X|_c-1)}$.

\item[\textup{(a$'$)}] The function $\lambda^o_{\check\zeta}$ depends only on the values of $\xi$ under $\mathfrak Z_2$ and within $r_w$ of the bonds based in the cells of $I(\check\zeta)$. If the sets $I(\check\zeta)$ of some letters lie pairwise more than $r_1$ cells apart, then the functions $\lambda^o_{\check\zeta}$ of these letters are independent, and so are the products $\lambda^o_{\check\zeta}e^{(x/2)\varsigma_1|\xi_{\mathsf Y(\check\zeta)}|^2}$ with $\mathsf Y(\check\zeta)\subset I(\check\zeta)$.

\item[\textup{(b)}] For every cell $\square$, $E_\xi[b_\square e^{(x/2)\varsigma_1|\xi_{\{\square\}}|^2}]\le e^{x^{-1/8}}u_b$; moreover $E_\xi[e^{(x/2)\varsigma_1|\xi_{\mathsf Y}|^2}]\le e^{4.8v\varsigma_1\#\mathsf Y}$.

\item[\textup{(c$^\circ$)}] In the bare case, for every admitted cell datum $(J,t)$ and every $\mathsf Y\subset\mathfrak Z_2$, \[E_\xi[\hat\lambda_{(J,t)}e^{(x/2)\varsigma_1|\xi_{\mathsf Y}|^2}]\le e^{x^{-1/8}\#\mathfrak Z_2}e^\circ_{(J,t)}.\] The mean on the left contains no indicator, no polymer and no member.

\item[\textup{(c)}] With the strong cap, for an admitted letter $\check\zeta$ and every $\mathsf Y\subset\mathfrak Z_2$, \[E_\xi[|\lambda^o_{\check\zeta}|e^{(x/2)\varsigma_1|\xi_{\mathsf Y}|^2}]\le e^{x^{-1/8}\#\mathfrak Z_2}e^o_{\check\zeta}.\]

\item[\textup{(c$'$)}] With the weak cap, the bound of \textup{(c)} holds with $e^{o\prime}_{\check\zeta}$ in the place of $e^o_{\check\zeta}$. Even the mean of $\Lambda_{\check\zeta}$ satisfies \[E_\xi[\Lambda_{\check\zeta}e^{(x/2)\varsigma_1|\xi_{\mathsf Y}|^2}]\le e^{x^{-1/8}\#\mathfrak Z_2}e^{\circ\prime}_{(J,t)}e^{-(\mathrm{py}/1152)\#\mathfrak X}.\]

\item[\textup{(c$''$)}] The bounds \textup{(c)} and \textup{(c$'$)} remain true when $1_{\mathrm{ind}}$ is replaced by any Borel function $0\le\mathsf{th}(\xi)\le1$ such that, on its support, every $X$ of $\mathfrak X$ holds a plaquette of $D(U)\cap\hat{\mathfrak Z}$. Since $\hat{\mathfrak b}\subset\hat{\mathfrak Z}$, they remain true in particular with the indicator of the same kind built with $D(U)\cap\hat{\mathfrak b}$ instead of $D(U)\cap\hat{\mathfrak Z}$.

\item[\textup{(d$^\circ$)}] For the cell data $(J_i,t_i)$ above, \[E_\xi[\prod_i\hat\lambda_i\prod_{J_0}b_\square\,e^{(x/2)\varsigma|\varphi_{\mathsf Y}|^2}]\le\exp\{4.8v(2/\gamma)\varsigma(2r_1+1)^3\#\mathsf Y+x^{-1/8}(\sum_i\#\mathfrak b_i^{+2r_1}+\#J_0)\}\,u_b^{\#J_0}\prod_ie^\circ_{(J_i,t_i)}.\]

\item[\textup{(d)}] For the letters $\check\zeta_i$ above, the bound of \textup{(d$^\circ$)} holds with $|\lambda^o_i|$ in the place of $\hat\lambda_i$, and with $e^o_{\check\zeta_i}$ for the strong cap, or $e^{o\prime}_{\check\zeta_i}$ for the weak cap, in the place of $e^\circ_{(J_i,t_i)}$.

\item[\textup{(d$'$)}] The bound of \textup{(d$^\circ$)} also holds with $\Lambda_i$ in the place of $\hat\lambda_i$ and with $e^{\circ\prime}_{(J_i,t_i)}e^{-(\mathrm{py}/1152)\#\mathfrak X_i}$ in the place of $e^\circ_{(J_i,t_i)}$.

\item[\textup{(e)}] One has $u_1=e^{-\delta_3p_0^2}$, $u_1'=e^{-\delta_3'p_0^2}$ and $e^{-\mathrm{py}/1152}=e^{-d_0p_0^2/2304}$, where $\delta_3$ and $\delta_3'$ are computed from $(d_0,L,a,K_*,\lambda_*)$.

\item[\textup{(f)}] For all $U$ and $V'$, \[L[\mathrm{En}_\ell(U)+a\sum_c(1-\cos\mathsf n_c)]\ge\mathrm{En}'(V')+(\frac12-d_0-\mathsf t_\rho-3e_1^\flat)\varepsilon'\sum_{D'}\phi'_P+\frac12Ld_0\varepsilon\sum_D\phi_p-e_2^\flat\varepsilon'^2N^3.\]

\item[\textup{(g*1)}] Let $O$ be a finite set of objects of \textup{E12}, consisting of classes, which form a set $J$, of cells off $f_J$, which form a set $t$ such that $\mathfrak b$ is non-empty, and of polymers, which form a family $\mathfrak X$; any two of these polymers are farther apart than $s(x)$, and each of them is adjacent to an object with cells. The set $O$ need not be connected. Let \textup{(H17)} hold for $(J,t)$ at $V'$; this is so, for instance, when $(J,t)$ is the cell datum of a zone letter that is actual at $V'$ (see \textup{(Act2)} in Step~6 of the proof of \Cref{Y-G-ThA}). Then the adjacency of \textup{E12} splits $O$ into components $\check\zeta^{(m)}=(J_m,t_m,\mathfrak X_m)$, which are admitted letters, any two of them compatible. Moreover $\mathfrak Z$ is the disjoint union of the $\mathfrak Z_m$, and $\mathfrak Z_2$ that of the $\mathfrak Z_{2,m}$; $1_{\mathrm{ind}}=\prod_m1_{\mathrm{ind},m}$; $\hat\lambda_{(J,t)}1_{\mathrm{ind}}\prod_{\mathfrak X}w=\prod_m\lambda^o_{\check\zeta^{(m)}}$; the factors $\mathrm{Res}$ and $u_t^{\#t}$ split into products over the components; and $\prod_mu_1^{1[t_m\ne\emptyset]}\le u_1^{1[t\ne\emptyset]}$. Hence \textup{(d)} applies to the product $\prod_m\lambda^o_{\check\zeta^{(m)}}$.

\item[\textup{(g*2)}] Let the same hypothesis hold for $(J,t)$. Then the components $(J_k,t_k)$ of $(J,t)$ under the adjacency of cells alone are pairwise compatible admitted bare letters, $\hat\lambda_{(J,t)}=\prod_k\hat\lambda_{(J_k,t_k)}$, and $\mathrm{Res}$ and $u_t^{\#t}$ split into products over these components.\end{enumerate}

No conclusion of this theorem depends on the reality or the positivity of a member.\end{theorem}

\begin{proof}\emph{Step 1: the free kernels and the multipliers.}

\begin{lemma}[Norms, identities and tails of the free kernels]\label{Y-F-LR} Let $(\ell,\ell')$ be a pair of item~C of \Cref{Y-F-C7}. Then the following hold.

\begin{enumerate}\item[\textup{(i)}] Each of $\ell,\mathsf T,\mathsf X,\ell',\mathsf B',H^\sharp,\mathsf N^\sharp$, and each of the free kernels $\mathsf G,\mathsf G_Z$, has $\|\cdot\|_{\theta_D}\le C_D$.

\item[\textup{(ii)}] $\mathsf B\mathsf T=L^{-1}W^*+\mathrm{d}^*\mathsf X$, $H^\sharp=\mathsf T\mathsf B'$, $\mathsf N^\sharp=(aL)^{-1}\mathrm{d}'^*\mathsf B'$ and $\mathsf B'(L\mathsf T^*\mathsf B\mathsf T+(aL)^{-1}\mathrm{d}'\mathrm{d}'^*)=1$.

\item[\textup{(iii)}] $L^2\|\mathsf T^<\|_\infty\le c_1^*\mathsf l_L$, $L\|\mathsf X^<\|_\infty\le c_2^*$, $\|\mathsf B'^{\rm c}\|_\infty\le \mathsf N_Z$ and $\mathsf t_\rho\le\frac18(\frac12-d_0-c_\Xi)$.

\item[\textup{(iv)}] For each of the kernels of \textup{(i)}, the tail that the cut at $R(x')$ removes has $\|\cdot\|_0\le C_Dx'^{-64}$; and $\|\mathsf B'-\mathsf B'^{\rm c}\|_0\le2C_D^2x'^{-64}$.\end{enumerate}\end{lemma}

\begin{lemma}[Sizes of the multipliers]\label{Y-F-LM} At every pair $(V',e)$ one has $|v_*|\le \mathsf N_Z\varepsilon'$, $|Z_*|\le \mathsf N_ZL^{-1}\varepsilon'$, $|\Xi_*|\le \mathsf N_\Xi L^{-1}\varepsilon'$, $|\sigma_*|\le \mathsf N_\sigma L^{-2}\varepsilon'$, $|\ell^{V*}\sigma_*|\le l_c\mathsf N_\sigma L^{-2}\varepsilon'$, $|\nu_*|\le N_1L^{-2}\varepsilon'$ and $|\omega_*|\le4\mathsf N_Z(aL)^{-1}\varepsilon'$. The components of $\mathrm{pm}(\chi v_*)$ and of $\mathrm{pm}((1-\chi)v_*)$ obey the same bounds.\end{lemma}

\begin{proof}[Proof of \Cref{Y-F-LR}]\begin{enumerate}\item[\textup{(i)}] The kernels in \textup{(i)} are among those that the condition \textup{(Dec)}$(C_D,\theta_D)$ of \Cref{Y-A-Dec} lists, and \Cref{Y-A-Th36} gives this condition as its first conclusion.

\item[\textup{(ii)}] By \Cref{Y-A-Pairs}, $\nu=\mathsf RY$ satisfies $\nu+\ell^{\mathsf T}\sigma=Y$ with $\sigma=\ell\nu$; as $\mathsf B=1+\ell^{\mathsf T}\ell$, this says $\mathsf B\nu=Y$. By \Cref{Y-A-XT}, ${\rm pair}(v/L,\mathsf Xv)=(\mathsf Tv,(aL)^{-1}\mathrm{d}'^*v)$. For this pair $\nu=\mathsf Tv$ and $Y=L^{-1}W^*v+\mathrm{d}^*\mathsf Xv$, so the identity $\mathsf B\nu=Y$ yields $\mathsf B\mathsf T=L^{-1}W^*+\mathrm{d}^*\mathsf X$. \Cref{Y-A-XT} also gives $H^\sharp=\mathsf T\mathsf B'$, and the identity for $\mathsf N^\sharp$ is its definition $\mathsf N^\sharp:=(aL)^{-1}\mathrm{d}'^*\mathsf B'$ in \Cref{Y-A-P4}\,(a). Next, \Cref{Y-A-P4}\,(a), applied with $v=\mathsf B'X'$, gives $\|{\rm pair}(\mathsf B'X'/L,\mathsf X\mathsf B'X')\|^2=L^{-1}\langle X',\mathsf B'X'\rangle$. Written with the components $\mathsf Tv$ and $(aL)^{-1}\mathrm{d}'^*v$ of the pair, the left side is $\langle \mathsf B'X',(\mathsf T^*\mathsf B\mathsf T+a(aL)^{-2}\mathrm{d}'\mathrm{d}'^*)\mathsf B'X'\rangle$. Both $\mathsf B'$ and $\mathsf T^*\mathsf B\mathsf T+a(aL)^{-2}\mathrm{d}'\mathrm{d}'^*$ are symmetric, $\mathsf B'=1+\ell'^{\mathsf T}\ell'$ is invertible, and the identity holds for every $X'$; hence $\mathsf B'^{-1}=L\mathsf T^*\mathsf B\mathsf T+(aL)^{-1}\mathrm{d}'\mathrm{d}'^*$, because $La(aL)^{-2}=(aL)^{-1}$.

\item[\textup{(iii)}] \Cref{Y-A-Th33}\,(d), taken at the starred parameters, gives the first two bounds, since a cut can only lower absolute row sums. \Cref{Y-A-Th29}\,(a) applies, as $L\ge L_A$ and $\ell\in\mathcal R(b_*,K_*,\lambda_*,\frac19)$, and gives $|\ell'_{y'}(P)|\le K_*e^{-\lambda_*{\rm dist}}$; hence the absolute row sums of $\ell'^{\rm c}$ are at most $l_r$, its absolute column sums are at most $l_c$, and $\|\mathsf B'^{\rm c}\|_\infty\le1+l_cl_r=\mathsf N_Z$. The bound on $\mathsf t_\rho$ is the third conclusion of \Cref{Y-A-Th36}.

\item[\textup{(iv)}] The tail bound in \Cref{Y-A-Norm}, combined with \textup{(i)}, bounds the norm $\|\cdot\|_0$ of each such tail by $C_De^{-\theta_DR(x')}$; and $e^{-\theta_DR(x')}\le x'^{-64}$, since $R(x')\ge c_R\log x'$ and $c_R\theta_D=64$. Finally $\mathsf B'-\mathsf B'^{\rm c}=(\ell'-\ell'^{\rm c})^{\mathsf T}\ell'+\ell'^{{\rm c}\mathsf T}(\ell'-\ell'^{\rm c})$. Each of the two terms is a tail of $\ell'$, or its transpose, times a kernel of norm at most $C_D$, and so $\|\mathsf B'-\mathsf B'^{\rm c}\|_0\le2C_D^2x'^{-64}$.\end{enumerate}\end{proof}

\begin{proof}[Proof of \Cref{Y-F-LM}] At real $V'$ every transport $T(y,y'')$ is an isometry; so the absolute row sums obey the same bounds as in the free model, and a mask can only lower them. Since $|\zeta'|\le\varepsilon'$, we get $|v_*|\le\|\mathsf B'^{\rm c}\|_\infty\varepsilon'\le \mathsf N_Z\varepsilon'$ by \Cref{Y-F-LR}\,(iii), and $|Z_*|=|v_*|/L\le \mathsf N_ZL^{-1}\varepsilon'$. Next, $|\Xi_*|\le\|\mathsf X^<\|_\infty|v_*|$, and by \Cref{Y-F-LR}\,(iii) this is at most $c_2^*L^{-1}\mathsf N_Z\varepsilon'=\mathsf N_\Xi L^{-1}\varepsilon'$. As $\nu^{\mathrm{ref}}=(\mathsf T^<)^Vv_*$, the bound on $L^2\|\mathsf T^<\|_\infty$ in \Cref{Y-F-LR}\,(iii) gives $|\nu^{\mathrm{ref}}|\le c_1^*\mathsf l_LL^{-2}\mathsf N_Z\varepsilon'$; so $|\sigma_*|\le l_r|\nu^{\mathrm{ref}}|\le l_rc_1^*\mathsf l_LL^{-2}\mathsf N_Z\varepsilon'=\mathsf N_\sigma L^{-2}\varepsilon'$.

Since $\sum_PW(P,p)\le1/L$ and a fine plaquette lies in two fine cubes, $|Y_*|\le(\mathsf N_Z+2L\mathsf N_\Xi)L^{-2}\varepsilon'$. The absolute column sums of $\ell$ are at most $l_c$, so $|\ell^{V*}\sigma_*|\le l_c\mathsf N_\sigma L^{-2}\varepsilon'$, and $|\nu_*|\le|Y_*|+l_c|\sigma_*|\le N_1L^{-2}\varepsilon'$ by the definition of $N_1$. A coarse bond lies in four coarse plaquettes $P$, and this gives $|\omega_*|\le4\mathsf N_Z(aL)^{-1}\varepsilon'$. In these arguments $v_*$ enters only through the bound $|v|_\infty\le \mathsf N_Z\varepsilon'$, which $\chi v_*$ and $(1-\chi)v_*$ satisfy as well; so the components of $\mathrm{pm}(\chi v_*)$ and $\mathrm{pm}((1-\chi)v_*)$ obey the same bounds.\end{proof}

\emph{Step 2: the geometry of an admitted cell datum.}

\begin{lemma}[Distances of rough plaquettes, tame boxes, where objects are based, and inequalities between the radii]\label{Y-F-Geok} Let $(J,t)$ be an admitted cell datum, and let $x$ satisfy \textup{(G1$^\flat$)}. Then the following hold.

\begin{enumerate}\item[\textup{(i)}] A plaquette of $D'$ outside $D'_J$ is at coarse distance $>R_6$ from $\mathfrak Z_2$.

\item[\textup{(ii)}] Let an object be based in $\mathfrak Z_2$, and let it either be based off $f_J$ or lie within $r_w+\varrho_L+2$ of a site off $f_J$. Then $V'$ is tame on the box of radius $R_6$ about this object.

\item[\textup{(iii)}] Every plaquette of $D'$ within $4R+4\varrho_L+6$ of $\mathfrak Z$ belongs to $D'_J$, and $D'\cap\mathfrak Z=D'_J$.

\item[\textup{(iv)}] An object within $R(x')+2\varrho_L+3$ both of $\hat{\mathfrak Z}$ and of its complement is based in $\mathfrak Z_2\setminus\mathfrak b$.

\item[\textup{(v)}] Every site within $4R+4\varrho_L+5$ of $D'_J$ lies in a cell of $f_J$; so does every object based at such a site, and every field supported on such objects is supported in $f_J$.

\item[\textup{(vi)}] The bonds based in the rim $\mathrm{Lr}$ are farther than $r_w$ both from $\hat{\mathfrak b}$ and from the complement of $\hat{\mathfrak Z}_2$.

\item[\textup{(vii)}] $r_w+2\varrho_L+5\le r_1M$ and $3r_w+\varrho_L+1\le R_2-1$.

\item[\textup{(viii)}] For the terms $\mathrm{Rim}_{\mathfrak Z}$, $\mathrm{rim}_0$, $\mathrm{bd}_1$, $\mathrm{bd}_2$ and $\mathrm{bd}_3$ of Steps~3 and~7 of the proof of \Cref{Y-F-C7}, the following bonds are based in $\mathrm{Lr}$: for each row $y\in\hat{\mathfrak Z}\setminus\hat{\mathfrak Z}'$ of $\mathrm{Rim}_{\mathfrak Z}$, the bonds of every plaquette $p$ with $|y-p|\le\varrho(x)+1$ and those of the paths $\Gamma(y,\pi(p))$; for each row $p$ at which $k_p$ enters $\mathrm{rim}_0$ or $\mathrm{bd}_1$, all four bonds of $p$; for each row $c$ at which $k_n$ enters $\mathrm{rim}_0$ or $\mathrm{bd}_1$, every bond of $W_c$; and the bonds at which $\varphi$ and $\eta$ enter $\mathrm{bd}_2$ or $\mathrm{bd}_3$, those of $\mathrm{bd}_2$ being the bonds off $\hat{\mathfrak S}$ within $r_l$ cells of a row of $\mathfrak Z$.

\item[\textup{(ix)}] $R_6\le Mr(x')$ and $r_1-r_l-1-\lceil R/M\rceil\ge1$.\end{enumerate}\end{lemma}

\begin{proof}[Proof of \Cref{Y-F-Geok}] We use five facts.

\begin{enumerate}\item[\textup{(A)}] Sites of two cells at cell distance $d\ge1$ are at coarse distance $\ge(d-1)M+1$; and if a site is at coarse distance at most $D$ from a cell, then the cell that contains the site is within $\lceil D/M\rceil$ cells of that cell.

\item[\textup{(B)}] $r_l\ge1$ and $r_lM\ge\varrho_L+2$, by the definition $r_l=\lceil(\varrho_L+2)/M\rceil$.

\item[\textup{(C)}] $r_1M\ge R_2+3(\varrho_L+2)+5M=2R+3\varrho_L+8+5M$, by the definition of $r_1$ and \textup{(B)}.

\item[\textup{(D)}] $s_0M=2Mr(x')\ge10R+8\varrho_L+24$, by \textup{(G1$^\flat$)}.

\item[\textup{(E)}] $2\varrho_L\le R$ and $R(x')\le R$. Indeed $L\ge3$ gives $\varrho_L\le\lceil R/3\rceil$, and $2\lceil R/3\rceil\le2(R+2)/3\le R$ for $R\ge4$, which holds as $R\ge6$; and $R(x')\le R$ because $x'\le x$.\end{enumerate}

With these facts we prove the clauses in turn.

\begin{enumerate}\item[\textup{(i)}] By \textup{(H17)} the cell of such a plaquette is farther than $4r_1$ from $\mathfrak b$, so it is at cell distance $\ge2r_1+1$ from $\mathfrak Z_2=\mathfrak b^{+2r_1}$. By \textup{(A)} and \textup{(C)} its coarse distance from $\mathfrak Z_2$ is then $\ge2r_1M+1\ge4R+6\varrho_L+17+10M$, which exceeds $R_6=4R+6\varrho_L+11$.

\item[\textup{(ii)}] Let $j\in J$. A site off $f_j$ is at cell distance $\ge s_0+1$ from $\mathrm{mk}'_j$, hence, by \textup{(A)}, at coarse distance $\ge s_0M+1$ from $D'_j$. An object within $r_w+\varrho_L+2$ of such a site is therefore at coarse distance at least $s_0M+1-(r_w+\varrho_L+2)\ge9.5R+7\varrho_L+23>R_6$ from $D'_j$, by \textup{(D)} and $r_w\le R/2$. A site off $f_J$ is off every $f_j$ with $j\in J$, so the box of radius $R_6$ about the object holds no plaquette of $D'_J$. The other rough plaquettes are farther than $R_6$ from $\mathfrak Z_2$ by \textup{(i)}, and the object is based in $\mathfrak Z_2$; hence $V'$ is tame on that box.

\item[\textup{(iii)}] Let a plaquette of $D'$ lie within $4R+4\varrho_L+6$ of $\mathfrak Z$. If it were not in $D'_J$, then by \textup{(H17)} its cell would be farther than $4r_1$ from $\mathfrak b$, hence at cell distance $\ge3r_1+1$ from $\mathfrak Z=\mathfrak b^{+r_1}$, and so, by \textup{(A)} and \textup{(C)}, at coarse distance $\ge3r_1M+1>4R+4\varrho_L+6$ from $\mathfrak Z$. A plaquette of $D'$ in $\mathfrak Z$ lies within $r_1\le4r_1$ of $\mathfrak b$, hence in $D'_J$ by \textup{(H17)}; conversely $D'_J\subset\mathrm{mk}'_J\subset f_J\subset\mathfrak Z$. Thus $D'\cap\mathfrak Z=D'_J$.

\item[\textup{(iv)}] By \textup{(E)}, such an object lies within $R+2\varrho_L+3$ of $\hat{\mathfrak Z}$ and of its complement. The complement of $\mathfrak Z$ is at cell distance $\ge r_1+1$ from $\mathfrak b$, and so, by \textup{(A)} and \textup{(C)}, at coarse distance $\ge r_1M+1>R+2\varrho_L+3$ from $\hat{\mathfrak b}$; hence the object is not based in $\mathfrak b$. Since $\lceil(R+2\varrho_L+3)/M\rceil\le r_1$ by \textup{(C)}, the object lies, by \textup{(A)}, in a cell within $r_1$ of $\mathfrak Z$, that is, in $\mathfrak Z_2$.

\item[\textup{(v)}] By \textup{(D)}, $4R+4\varrho_L+5\le s_0M$. So, by \textup{(A)}, a site within $4R+4\varrho_L+5$ of $D'_j\subset\mathrm{mk}'_j$ lies in a cell within $s_0$ of $\mathrm{mk}'_j$, that is, in $f_j$.

\item[\textup{(vi)}] A cell of $\mathrm{Lr}$ is at cell distance $\ge r_1-2r_l-2$ from $\mathfrak b$. By \textup{(A)}, the bonds based in it are at coarse distance $\ge(r_1-2r_l-3)M+1$ from $\hat{\mathfrak b}$, and by the definition of $r_1$ this is $\ge R_2+(r_l+2)M+1>r_w$. Such a cell is also at cell distance $\ge(2r_1+1)-(r_1+r_l+1)=r_1-r_l$ from the complement of $\mathfrak Z_2$, and $(r_1-r_l-1)M+1>r_w$.

\item[\textup{(vii)}] As $r_w\le R/2$, the first inequality follows from \textup{(C)}. For the second, $3r_w\le1.5R$, and $1.5R+\varrho_L+1\le2R+1=R_2-1$ by \textup{(E)}.

\item[\textup{(viii)}] For $\mathrm{Rim}_{\mathfrak Z}$, the summand at a site $y\in\hat{\mathfrak Z}\setminus\hat{\mathfrak Z}'$ depends on the plaquettes $p$ with $|y-p|\le\varrho(x)+1$ and on the paths $\Gamma(y,\pi(p))$; these are within the coarse distance $\varrho_L+2\le r_lM$ of $y$, hence, by \textup{(A)} and \textup{(B)}, within $r_l$ cells of $y$, and their bonds lie one cell further out; all of these lie in $\mathfrak b^{+r_1+r_l+1}\setminus\mathfrak b^{+r_1-2r_l-1}$. For $\mathrm{rim}_0$, the bonds of the plaquettes $p\in\widehat{\mathfrak Z^{+r_l}}\setminus\hat{\mathfrak Z}'$ are based in $\mathfrak b^{+r_1+r_l+1}\setminus\mathfrak b^{+r_1-r_l-1}$. The term $\mathrm{rim}_0$ also lives on the bonds $c\in\mathfrak Z\setminus\mathfrak Z'$, and since such a $c$ is based in a cell of $\mathfrak b^{+r_1}\setminus\mathfrak b^{+r_1-r_l}$ and $W_c=B(c_-)\cup B(c_+)$ lies under that cell or a neighbouring one, every bond of $W_c$ is based in $\mathfrak b^{+r_1+1}\setminus\mathfrak b^{+r_1-r_l-1}$, inside the set obtained above for $\mathrm{rim}_0$. For $\mathrm{bd}_1$, the contributions at the $P$ and $q$ in $\mathfrak Z'^{+1}\setminus\mathfrak S$ depend on $W_P$, and the bonds lie one cell further out; so these bonds are based in $\mathfrak b^{+r_1-r_l+3}\setminus\mathfrak b^{+r_1-r_l-3}$, which lies inside $\mathrm{Lr}$ because $r_l\ge1$. For $\mathrm{bd}_2$, the bonds off $\hat{\mathfrak S}$ that are within reach of a row of $\mathfrak Z$, that is, within $r_l$ cells of it, lie in $\mathfrak b^{+r_1+r_l}\setminus\mathfrak S$. For $\mathrm{bd}_3$, the $P$ and $q$ lie in $\mathfrak S^{+1}\setminus\mathfrak S^-$, and their bonds one cell further out. All these sets lie in $\mathrm{Lr}=\mathfrak b^{+r_1+r_l+1}\setminus\mathfrak b^{+r_1-2r_l-3}$.

\item[\textup{(ix)}] By \textup{(E)}, $R_6=4R+6\varrho_L+11\le5R+4\varrho_L+12$, and $5R+4\varrho_L+12\le Mr(x')$ by \textup{(G1$^\flat$)}. Moreover $r_1\ge\lceil R/M\rceil+3r_l+5$, as $R_2\ge R$; hence $r_1-r_l-1-\lceil R/M\rceil\ge2r_l+4\ge1$.\end{enumerate}\end{proof}

\emph{Step 3: a local lower bound for the fine energy.}

\begin{lemma}[The energy of a set of cells bounded below by a linear term, defects and squares, minus volume and rim terms]\label{Y-F-LB} Let \textup{(L1$^\flat$)}, \textup{(F1)}, \textup{(F3$^\flat$)}, \textup{(F4$^\flat$)}, \textup{(F5)}, \textup{(Fm)}, \textup{(Fc)}, $\rho_0\le\varrho(x)$ and \textup{(Tor)} hold. Let $V'$ be any coarse field, let $\eta$ be any element of $\mathfrak L_{\mathfrak g}$, and let $U=\mathsf U(\eta;V')$. Let $\mathfrak Z,\mathfrak Z',\mathfrak S$ be any sets of cells with $\mathfrak Z'^{+r_l}\subset\mathfrak Z$ and $\mathfrak S^{+1}\subset\mathfrak Z'$. Let $(\nu,\sigma,\omega,\mathsf{zf},\Xi)$ be any fields on fine plaquettes, fine sites, coarse bonds, coarse plaquettes and fine cubes, in this order, that satisfy \begin{equation}\tag{Sup}|\nu|\le N_1L^{-3/2}\varepsilon,\quad|\sigma|\le \mathsf N_\sigma L^{-3/2}\varepsilon,\quad|\mathsf{zf}|\le \mathsf N_ZL^{-1/2}\varepsilon,\quad|\Xi|\le \mathsf N_\Xi L^{-1/2}\varepsilon,\quad|\omega|\le(2/9)\mathsf N_ZL^{1/2}\varepsilon,\end{equation} and write $\boldsymbol\nu=(\nu,\sigma,\omega)$. With \[\mathcal A_{\mathfrak Z}:=\sum_{\hat{\mathfrak Z}}S_\varepsilon(\phi_p)+\frac12\sum_{y\in\hat{\mathfrak Z}}m_y|\ell^U_y\zeta|^2+a\sum_{c\in\mathfrak Z}(1-\cos\mathsf n_c)\] one then has \[\mathcal A_{\mathfrak Z}\ge\langle\boldsymbol\nu,\iota(F,n)\rangle_{\mathfrak Z}-\frac12\|\boldsymbol\nu\|_{\mathfrak Z}^2+|\langle\mathsf{zf},\mathrm{de}\rangle_{\mathfrak S}|+|\langle\Xi,\mathrm{db}\rangle_{\mathfrak S}|+\mathrm{Sq}_{\mathfrak Z}+\mathrm{Pay}_{\mathfrak Z}-\mathrm{vol}_{\mathfrak Z}-\mathrm{Rim}_{\mathfrak Z},\] where the terms on the right are defined as follows.

The squares are \[\mathrm{Sq}_{\mathfrak Z}=\frac12(1-4\vartheta_F^\flat)\sum_{\hat{\mathfrak Z}\setminus D}|F_p-\nu_p|^2+\frac12\sum_{y\in\hat{\mathfrak Z},m_y=1}|\ell^U_y\zeta-\sigma_y|^2+\sum_{c\in\mathfrak Z,\mathsf n_c\le t_0}\mathsf{kc}_c|n_c-\tilde\omega_c|^2,\] where $\tilde\omega=\omega/(1-\theta_2)$; here $\mathsf{kc}_c=(a/2)(1-\theta_2)-2\vartheta_n^\flat$ off $\mathsf{Cb}$ and $\mathsf{kc}_c=a/4$ on $\mathsf{Cb}$, and $\mathsf{Cb}$ is the set of the bonds of those $P$ in $\mathfrak S$ that have $\Sigma_P\ge\pi/12$ or $\mathrm{na}_P\ge\pi/3$. The payments are \[\mathrm{Pay}_{\mathfrak Z}=(a/20)\sum_{c\in\mathfrak Z,\mathsf n_c>t_0}\mathsf n_c^2+\frac12d_0\varepsilon\sum_{D\cap\hat{\mathfrak Z}}\phi_p.\] The volume term is \[\mathrm{vol}_{\mathfrak Z}=2\vartheta_F^\flat|\nu|_\infty^2\,3L^3M^3\#\mathfrak Z+(a\theta_2+2.1\vartheta_n^\flat)|\omega|_\infty^2\,3M^3\#\mathfrak Z,\] and the rim term is \[\mathrm{Rim}_{\mathfrak Z}=\sum_{y\in\hat{\mathfrak Z}\setminus\hat{\mathfrak Z}'}|\sigma_y|[1\{m_y=1\}(|\mathsf{dl}_y|+|(\ell^V(F-\zeta))_y|)+1\{m_y=0\}|(\ell^VF)_y|],\] with $\mathsf{dl}_y=\ell^U_y\zeta-(\ell^V\zeta)_y$.

For the starred fields $(\nu_*,\sigma_*,\omega_*,Z_*,\Xi_*)$, condition \textup{(Sup)} holds by \Cref{Y-F-LM}, by $\varepsilon'\le L^{1/2}\varepsilon$ and by $aL\ge18$; and for these fields \[L\,\mathrm{vol}_{\mathfrak Z}\le[6N_1^2\vartheta_F^\flat+3\mathsf N_Z^2(\theta_2+L\vartheta_n^\flat)]\varepsilon'^2M^3\#\mathfrak Z.\]\end{lemma}

\begin{proof}[Proof of \Cref{Y-F-LB}] As $t_0=\varepsilon^{1/2}$, the bound on $|\omega|$ in \textup{(Sup)} together with \textup{(F5)} gives $|\omega|\le0.0023t_0$. By \textup{(F1)}, $90L^3\varepsilon\le1$, and $L\ge3$; hence $\varepsilon\le1/2430$ and $\theta_2=\varepsilon/12\le3.5\cdot10^{-5}$. It follows that $1/(1-\theta_2)\le1.01$ and $2/(1-\theta_2)^2\le2.1$.

\emph{\textup{(1)} One row.} We treat one row at a time. At a plaquette $p\notin D$ the term $S_\varepsilon(\phi_p)$ equals $\frac12\phi_p^2$, and completing the square gives $\frac12\phi_p^2=\langle\nu_p,F_p\rangle-\frac12|\nu_p|^2+\frac12|F_p-\nu_p|^2$. At a plaquette $p\in D$ the term equals $d_0\varepsilon\phi_p$; since $|F_p|\le\phi_p$, we have $\langle\nu_p,F_p\rangle\le|\nu|_\infty\phi_p$, and hence $d_0\varepsilon\phi_p\ge\langle\nu_p,F_p\rangle-\frac12|\nu_p|^2+(d_0\varepsilon-|\nu|_\infty)\phi_p$.

At a fine site $y$ with $m_y=1$ we set $a=\ell^U_y\zeta=(\ell^VF)_y+\mathsf{dl}_y-(\ell^V(F-\zeta))_y$; in this case only, $a$ denotes a vector and not the letter $a$ of \textup{(O4)}. Then $\frac12|a|^2=\langle\sigma_y,a\rangle-\frac12|\sigma_y|^2+\frac12|a-\sigma_y|^2$. At a fine site $y$ with $m_y=0$, the inequality $\langle\sigma_y,(\ell^VF)_y\rangle\le|\sigma_y||(\ell^VF)_y|$ gives $0\ge\langle\sigma_y,(\ell^VF)_y\rangle-\frac12|\sigma_y|^2-|\sigma_y||(\ell^VF)_y|$. At a coarse bond $c$ we define $\mathsf{et}_c$ by $a(1-\cos\mathsf n_c)=a\langle\omega_c,n_c\rangle-(a/2)|\omega_c|^2+\mathsf{et}_c$.

We sum these relations over the rows of $\mathfrak Z$. The site terms of no definite sign at the sites $y\in\hat{\mathfrak Z}\setminus\hat{\mathfrak Z}'$ make up $\mathrm{Rim}_{\mathfrak Z}$. Those at the sites $y\in\hat{\mathfrak Z}'$ involve the plaquettes within $\varrho(x)+1$ of $y$, the paths $\Gamma(y,\pi(p))$, and their blocks and double blocks; all of these lie within the coarse distance $\varrho_L+2\le r_lM$ of $y$, hence, as $\mathfrak Z'^{+r_l}\subset\mathfrak Z$, in $\hat{\mathfrak Z}$ and in $\mathfrak Z$.

\emph{\textup{(2)} Noise.} Let $c$ be a coarse bond. If $\mathsf n_c\le t_0$, note that $t^2/12\le\varepsilon/12=\theta_2$ for $0\le t\le t_0$, so that $1-\cos t\ge t^2/2-t^4/24$ gives $1-\cos t\ge(1-\theta_2)t^2/2$. With $|n_c|\le\mathsf n_c$ and $\tilde\omega_c=\omega_c/(1-\theta_2)$, completing the square in the definition of $\mathsf{et}_c$ then yields $\mathsf{et}_c\ge(a/2)(1-\theta_2)|n_c-\tilde\omega_c|^2-a\theta_2|\omega_c|^2$. If $\mathsf n_c>t_0$, then $1-\cos\mathsf n_c\ge2\mathsf n_c^2/\pi^2$ and $|\langle\omega_c,n_c\rangle|\le|\omega|\mathsf n_c\le(|\omega|/t_0)\mathsf n_c^2$; so $\mathsf{et}_c\ge a\,\mathsf n_c^2(2/\pi^2-|\omega|/t_0)\ge0.2a\,\mathsf n_c^2$, by the bound $|\omega|\le0.0023t_0$.

\emph{\textup{(3)} The sum of $|\mathrm{de}_P|$ over $\mathfrak S$.} We start from a bound on the plaquette defect at a single coarse plaquette.

\begin{lemma}[Sizes of the plaquette defect and the cube defect]\label{Y-F-L8} Let $\eta\in\mathfrak L_{\mathfrak g}$, let $U=\mathsf U(\eta;V')$, and let $(V',e)$ be any pair. Then \textup{(a)} and \textup{(b)} hold for every coarse plaquette $P$; a third clause \textup{(c)}, on the cube defect $\mathrm{db}_q$ at every fine cube $q$ and for the same $\eta$, $U$ and $(V',e)$, is stated and proved in \textup{(4)} below.

\begin{enumerate}\item[\textup{(a)}] $|\mathrm{de}_P|\le\pi+L^{-1}\Sigma_P+\mathrm{na}_P$.

\item[\textup{(b)}] If $5\Sigma_P+\mathrm{na}_P<3\pi/4$, then $|\mathrm{de}_P|\le6(5\Sigma_P+\mathrm{na}_P)^2$.\end{enumerate}\end{lemma}

\begin{proof}[Proof of \Cref{Y-F-L8}\,(a), (b)]\begin{enumerate}\item[\textup{(a)}] In $\mathrm{de}_P=X'_P-(W^VF)_P-(\mathrm{d}'^Vn)_P$ we have $|X'_P|\le\pi$. By \Cref{Y-A-Facts}, $0\le W(P,p)\le L^{-1}$, and $W(P,p)=0$ unless $p$ lies in $W_P$; moreover $|F_p|\le\phi_p$, $|n_c|\le\mathsf n_c$, and the transports $T$ are isometries. Hence $|(W^VF)_P|\le L^{-1}\Sigma_P$ and $|(\mathrm{d}'^Vn)_P|\le\mathrm{na}_P$, which proves \textup{(a)}.

\item[\textup{(b)}] We write $\mathsf{Sg}=\Sigma_P$ and $\mathrm{na}=\mathrm{na}_P$, and $P=(y;\mu\nu)$, with $y_1=y+Le_\mu$, $y_2=y+Le_\nu$, $y_3=y_1+Le_\nu$ and with the bonds $c_1=(y,\mu)$, $c_2=(y_1,\nu)$, $c_3=(y_2,\mu)$, $c_4=(y,\nu)$ of $P$; the letters $n_j,\bar\tau_j,\tau_{j,x}$ belong to $c_j$.

\emph{1. The gauge.} We transform $(U,V',e)$ by the $u$ that is constant on blocks and has $u(y)=1$, $u(y_1)=V'(c_1)$, $u(y_2)=V'(c_4)$, $u(y_3)=V'(c_1)V'(c_2)$. The combs keep the value $1$, the values $V'(c_1),V'(c_2),V'(c_4)$ become $1$, and $V'(c_3)$ becomes $\partial V'(P)^{-1}$. For $p\in W_P$ the value $F_p$ becomes $T(y,y(p))F_p$, and $n_c$ becomes $T(y,c_-)n_c$, while $X'_P,\phi_p,\mathsf n_c,s_c,s_B,\operatorname{ang}(e^{i\eta_b})$ keep their values. So $\mathrm{de}_P$ keeps its value, and in the new field $\mathrm{de}_P=X'_P-\sum_pW(P,p)F_p-(n_1+n_2-n_3-n_4)$.

\emph{2. The letters.} By \Cref{Y-F-L2}, $\operatorname{ang}(h_{c,x}h_{c,c_-}^{-1})\le s_c\le\mathsf{Sg}<\pi$, the last inequality by the hypothesis of \textup{(b)}, which also gives $\mathsf n_c<\pi$. Combining these gives $h_{c,x}=e^{i\tau_{c,x}}e^{-i\bar\tau_c}\mathsf M(U)(c)$ and $\mathsf M(U)(c)=e^{-in_c}V'(c)$.

\emph{3. One square.} For $x\in B(y)$ let $x_1=x+Le_\mu$ and $x_2=x+Le_\nu$, and put \[H_x:=h_{c_1,x}h_{c_2,x_1}h_{c_3,x_2}^{-1}h_{c_4,x}^{-1};\] this is the product of $U$ along the boundary of the square $S_x$ of $L\times L$ fine plaquettes whose least corner is $x$. Solving for $\partial V'(P)$ with the help of items~1 and~2, we find \[\partial V'(P)=e^{in_2}e^{i\bar\tau_2}e^{-i\tau_{2,x_1}}e^{in_1}e^{i\bar\tau_1}e^{-i\tau_{1,x}}H_xe^{i\tau_{4,x}}e^{-i\bar\tau_4}e^{-in_4}e^{i\tau_{3,x_2}}e^{-i\bar\tau_3}e^{-in_3}.\]

\emph{4. The square as a product.} In a fixed order, column by column and each column as a ladder, \[H_x=\prod_{p\in S_x}e^{i\operatorname{Ad}_{k_p}F_p}.\] We derive it from the ladder identity that appears in the proof of \Cref{Y-C-L67}\,(a). That identity is a telescoping product, valid in an arbitrary group, so it remains true when its $h_j$ are paths. In a first application, $u_j$ and $v_j$ are the lowest and the highest bond of column number $j$, and $h_j$ is the straight path of $L$ bonds that starts at $x+je_\mu$ and runs along $e_\nu$; then $H_x$ is the ordered product of the loops along the boundaries of the columns, each conjugated by $U$ along $e_\mu$ from $x$ to $x+je_\mu$. A second application, inside each column, concerns the inverse of the loop of that column; now the steps go along $e_\nu$ and the rungs point in the direction $\mu$, and the loops that appear are the $\partial U(p)^{-1}$, each conjugated by $U$ along $e_\nu$ from $x+je_\mu$ to $\pi(p)$. Thus $k_p$ is $U$ along the path from $x$ along $e_\mu$, then along $e_\nu$, to $\pi(p)$; and $e^{iF_p}=\partial U(p)$ because $\phi_p\le\mathsf{Sg}<\pi$.

\emph{5. Size.} The sets $W_{c_1}$ and $W_{c_3}$ are disjoint parts of $W_P$, and so are $W_{c_2}$ and $W_{c_4}$; hence $\sum_js_{c_j}\le2\mathsf{Sg}$. By \Cref{Y-F-L2}, the norms of all the exponents in the formula of item~3, with $H_x$ written as in item~4, add up to at most $\mathrm{na}+2\sum_js_{c_j}+\mathsf{Sg}\le5\mathsf{Sg}+\mathrm{na}<3\pi/4$. Therefore $\partial V'(P)\ne-1$ and $X'_P=\mathsf L(\partial V'(P))$, and \Cref{Y-F-L1}\,(c) gives \[X'_P=(n_1+n_2-n_3-n_4)+[\bar\tau_1-\tau_{1,x}+\bar\tau_2-\tau_{2,x_1}-\bar\tau_3+\tau_{3,x_2}-\bar\tau_4+\tau_{4,x}]+\sum_{p\in S_x}\operatorname{Ad}_{k_p}F_p+\mathrm{err}_x,\] with $|\mathrm{err}_x|\le(5\mathsf{Sg}+\mathrm{na})^2$.

\emph{6. Mean over $x\in B(y)$.} As $x$ runs through $B(y)$, the points $x_1$ and $x_2$ run through $B(y_1)$ and $B(y_2)$; so the bracket has mean $0$. Moreover $L^{-3}\sum_x\sum_{p\in S_x}F_p=\sum_pW(P,p)F_p$, by the formula that gives $W=Q_2$ in \Cref{Y-A-Lattices}.

\emph{7. Tails.} The path of $k_p$ runs inside blocks along at most four stretches, and it uses at most two crossing bonds: one of them belongs to $c_1$, the other to $c_2$ or to $c_4$, and none belongs to $c_3$. By \Cref{Y-F-L3}\,(2), applied with $\mathrm{Bx}=B$ and $o$ the centre of $B$, each stretch inside a block $B$ has $\operatorname{ang}\le s_B$, and together the stretches contribute at most $2\mathsf{Sg}$. On a crossing bond $V'(c)=1$ now, so $\operatorname{ang}U_b=\operatorname{ang}(e^{i\eta_b})\le9s_c+\mathsf n_c$ by \textup{(Tc)}\,(i) in \textup{(4$'$)} below; the proof of that bound rests on Lems.~\ref{Y-F-L1}, \ref{Y-F-L2} and~\ref{Y-F-L3} alone. As $s_c\le\mathsf{Sg}$ for each bond $c$ of $P$, we obtain $\operatorname{ang}k_p\le20\mathsf{Sg}+\mathrm{na}$, and \Cref{Y-F-L1}\,(b) gives $\sum_{p\in S_x}\|\operatorname{Ad}_{k_p}-1\|\phi_p\le2(20\mathsf{Sg}+\mathrm{na})\mathsf{Sg}$.

\emph{8. Conclusion.} By items~1 and~6, averaging the identity of item~5 over $x\in B(y)$ shows that $\mathrm{de}_P$ is the mean over $x$ of $\sum_{p\in S_x}(\operatorname{Ad}_{k_p}-1)F_p+\mathrm{err}_x$. With $|F_p|\le\phi_p$ and the bounds of items~5 and~7 we obtain \[|\mathrm{de}_P|\le(5\mathsf{Sg}+\mathrm{na})^2+40\mathsf{Sg}^2+2\mathsf{Sg}\,\mathrm{na}\le2.6(5\mathsf{Sg}+\mathrm{na})^2\le6(5\mathsf{Sg}+\mathrm{na})^2.\]\end{enumerate}\end{proof}

Let $\mathrm{Pl}_1=\{\Sigma_P\ge\pi/12\}$. For $P$ in $\mathrm{Pl}_1$, the tame part of the sum over $W_P$ is at most $12L^3\varepsilon\le\pi/24$; hence $\Sigma^D\ge\pi/24$, and \Cref{Y-F-L8}\,(a), combined with the inequalities $\pi\le24\Sigma^D$ and $\Sigma_P\le2\Sigma^D$, yields $|\mathrm{de}_P|\le(24+2/L)\Sigma^D+\mathrm{na}_P\le25\Sigma^D+\mathrm{na}_P$. Let $\mathrm{Pl}_2$ be the set of those $P$ outside $\mathrm{Pl}_1$ with $\mathrm{na}_P\ge\pi/3$. On $\mathrm{Pl}_2$ we combine \Cref{Y-F-L8}\,(a) with $\Sigma_P<\pi/12$, $L\ge3$ and $\pi\le3\mathrm{na}_P$ and obtain $|\mathrm{de}_P|\le\pi(1+1/36)+\mathrm{na}_P\le5\mathrm{na}_P$. Every remaining $P$ has $5\Sigma+\mathrm{na}<3\pi/4$, because $\Sigma_P<\pi/12$ and $\mathrm{na}_P<\pi/3$; this is the hypothesis of \Cref{Y-F-L8}\,(b), and that clause gives $|\mathrm{de}_P|\le6(5\Sigma+\mathrm{na})^2\le300\Sigma^2+12\mathrm{na}^2$. Here $\Sigma^2\le(\pi/6)\Sigma^D+24L^3\sum_{\rm tame}\phi^2$, with the tame sum taken over $W_P$, and $\mathrm{na}_P^2\le4\sum\mathsf n_c^2$, with the sum taken over the four bonds $c$ of $P$.

A fine plaquette lies in twelve boxes $W_P$, and a bond lies in four plaquettes $P$. Summing the three cases over $\mathfrak S$, we arrive at \[\sum_{\mathfrak S}|\mathrm{de}_P|\le2200\sum_D\phi+86400L^3\sum_{\rm tame}\phi^2+192\sum_{\mathsf n_c<\pi/3}\mathsf n_c^2+20\sum_{\mathsf{Cb}}\mathsf n_c.\] The four coefficients come from the bound $300+600\pi\le2184.96\le2200$ and from the numbers $300\cdot24\cdot12$, $12\cdot4\cdot4$ and $5\cdot4$.

\emph{\textup{(4)} The sum of $|\mathrm{db}_q|$ over $\mathfrak S$.} We show that \Cref{Y-F-L8}\,(c) gives \[\sum_{\mathfrak S}|\mathrm{db}_q|\le40\sum_D\phi+110L^3\sum_{\rm tame}\phi^2+\sum_{\mathsf n_c<\pi/3}\mathsf n_c^2+2L^2\varepsilon\sum_{\mathsf n_c\ge\pi/3}\mathsf n_c.\]

In this bound the rough plaquettes carry at most the factor $6\pi+2+1\le22$, and the tame ones the factor $6+99L^3+L^2\le110L^3$.

The size of the cube defect is the subject of clause \textup{(c)} of \Cref{Y-F-L8}, which reads as follows.

\begin{enumerate}\item[\textup{(c)}] For every $\eta\in\mathfrak L_{\mathfrak g}$, every pair $(V',e)$ and every fine cube $q$, with $U=\mathsf U(\eta;V')$, \[|\mathrm{db}_q|\le\frac12S_q^2+\sum_{\lambda=1,2,3}\phi_{p_\lambda}\min\{2,2(9A+\mathrm{na})\},\] where $S_q=\sum_{p\in dq}\phi_p$ and $p_\lambda$ is the face of $q$ based at $\pi(q)+e_\lambda$; here $(A,\mathrm{na})=(s_B,0)$, with $B=B(y(q))$, if $\pi(q)+e_\lambda$ lies in the block of $\pi(q)$, and otherwise $(A,\mathrm{na})=(s_c,\mathsf n_c)$, with $c=(y(q),\lambda)$.\end{enumerate}

\begin{proof}[Proof of \Cref{Y-F-L8}\,(c)] The transports $T$ are isometries and $|F_p|\le\phi_p$, so $|\mathrm{db}_q|\le S_q$; if $S_q\ge2$, it follows that $|\mathrm{db}_q|\le S_q\le\frac12S_q^2$.

Now let $S_q<2$. Then $e^{iF_p}=\partial U(p)$ at each of the six faces $p$ of $q$.

We apply Bianchi's identity to $U$ at the fine cube $q$, with $z=\pi(q)$; as in the proof of \Cref{Y-C-L67}\,(d), it is an identity of invertible matrices. It states that $\prod_{k=1,\dots,6}e^{iY_k}=1$, where $Y_k=d(q,p_k)\operatorname{Ad}_{\kappa_k}F_{p_k}$, with $\kappa_k=1$ at the faces based at $z$ and $\kappa_k=U_b$, $b=(z,\lambda)$, at the face $p_\lambda$.

By \Cref{Y-F-L1}\,(c$'$), $|\sum_kY_k|\le\frac12S_q^2$. Moreover \[\mathrm{db}_q-\sum_kY_k=\sum_\lambda d(q,p_\lambda)[T(y(q),y(p_\lambda))-\operatorname{Ad}_{U_b}]F_{p_\lambda}.\]

If $b$ lies in the block $B$, then $T=1$ and $\operatorname{ang}U_b\le s_B$ (on a comb, $U_b=1$). If $b$ crosses $c$, then $T=\operatorname{Ad}_{V'(c)}$ and $U_b=e^{i\eta_b}V'(c)$, and $\operatorname{ang}(e^{i\eta_b})\le9s_c+\mathsf n_c$.

These two angle bounds are proved in the cases \textup{(o)} and \textup{(i)} of the family \textup{(Tc)} in \textup{(4$'$)} below. Given them, \Cref{Y-F-L1}\,(b) bounds the norm of the bracket by $\min\{2,2(9A+\mathrm{na})\}$, and clause \textup{(c)} follows.\end{proof}

We now sum over the cubes $q$ of $\mathfrak S$. A cube has six faces and a plaquette is a face of two cubes; therefore \[\sum_q\frac12S_q^2\le3\sum_q\sum_{dq}\phi_p^2=6\sum_p\phi_p^2\le6\pi\sum_D\phi+6\sum_{\rm tame}\phi^2.\]

A plaquette is the face $p_\lambda$ of one cube only. If $p_\lambda$ is rough, its term is at most $2\phi_{p_\lambda}$; if it is tame, its term is at most $\phi_{p_\lambda}(18A+\min\{2,2\mathrm{na}\})$.

The region of $A$ is the set of those fine plaquettes $p$ for which $\phi_p$ is a term of $A$: the block $B$ if $A=s_B$, and $W_c$ if $A=s_c$. It holds at most $6L^3$ plaquettes, and a given plaquette lies in the region of at most $3L^3+6L^2\le5L^3$ faces $p_\lambda$, namely the faces of its block and $L^2$ seam faces for each of six $W_c$.

So the part with $A^D$ is at most $18\varepsilon\,5L^3\sum_D\phi\le\sum_D\phi$, by \textup{(F1)}, and the part with $A^t$ is at most $18\cdot\frac12(6+5)L^3\sum_{\rm tame}\phi^2=99L^3\sum_{\rm tame}\phi^2$.

In the part with $\mathsf n_c$, each bond $c$ has $L^2$ seam faces. If $\mathsf n_c<\pi/3$, then $2\mathsf n_c\sum_{\rm seam}\phi\le\mathsf n_c^2+L^2\sum_{\rm seam}\phi^2$. Otherwise the part is at most $2L^2\varepsilon\le2L^2\varepsilon\,\mathsf n_c$, as $\pi/3>1$.

For $P$ and $q$ in $\mathfrak S$, everything that enters these bounds lies in $\mathfrak S^{+1}\subset\mathfrak Z$, and restricting $P$ and $q$ to $\mathfrak S$ does not increase the multiplicities counted in \textup{(3)} and \textup{(4)}.

\emph{\textup{(4$'$)} The three families of site terms of no definite sign at $y\in\hat{\mathfrak Z}'$.}

\begin{enumerate}\item[\textup{(Ta)}] Where $m_y=1$, we have $|F_p-\zeta_p|\le\phi_p1_D$, and therefore $\sum_y|\sigma_y||(\ell^V(F-\zeta))_y|\le|\sigma|_\infty l_c\sum_D\phi$.

\item[\textup{(Tb)}] Where $m_y=0$, the rough part of $(\ell^VF)_y$ contributes at most $|\sigma|_\infty l_c\sum_D\phi$, while the tame part is at most $l_r\varepsilon$ at each site. The number of these sites obeys $\#\{m_y=0\}\le(2\rho_0+1)^3\#D+(2\varrho(x)+3)^3\#\{\phi_p>1\}$, and $\varepsilon\#D\le\sum_D\phi$, $\varepsilon\#\{\phi_p>1\}\le\varepsilon\sum_D\phi$; so the tame part is in all at most $|\sigma|_\infty l_r[(2\rho_0+1)^3+(2\varrho(x)+3)^3\varepsilon]\sum_D\phi$.

\item[\textup{(Tc)}] We put $\operatorname{ang}_b:=\min\{1,\operatorname{ang}(e^{i\eta_b})\}$ and $z_p:=|\zeta_p|\le\min\{\phi_p,\varepsilon\}$.

Let $\Gamma=\Gamma(y,\pi(p))$. The crossings of $\Gamma$ form the path $\Gamma'=\Gamma(y(y),y(p))$. We claim that, at every $\eta$, \[\operatorname{ang}(U(\Gamma)V'(\Gamma')^{-1})\le\sum\operatorname{ang}(e^{i\eta_b}),\] the sum running over the bonds $b$ of $\Gamma$ in $I^\circ$.

To see this, we use parts \textup{(1)} and \textup{(2)} of Step~5 of the proof of \Cref{Y-E-C6}, which give \[U(\Gamma)=[\prod_ig_ie^{\pm i\eta_{b_i}}g_i^{-1}]V'(\Gamma'),\] with $g_i\in G$ and $b_i$ the bonds of $\Gamma$ in $I^\circ$. Part \textup{(1)} is a statement on the paths alone and does not involve the field, and part \textup{(2)} is an identity between products of group elements. They do not depend on the size of $\eta$; they apply because $\varrho(x)+L<NL/2$ by \textup{(Tor)}.

Since $\operatorname{ang}$ is subadditive and invariant under conjugation and under passage to the inverse (\Cref{Y-F-L1}\,(a)), the claim follows. \Cref{Y-F-L1}\,(b) then gives $\|\operatorname{Ad}_U-T\|\le\min\{2,2\sum\operatorname{ang}\}\le2\sum\operatorname{ang}_b$, and hence \[|\mathsf{dl}_y|\le2\sum_p|\ell_y(p)|z_p\sum_{b\in\Gamma(y,\pi(p))}\operatorname{ang}_b.\]

We now bound $\operatorname{ang}(e^{i\eta_b})$ at every $\eta$; in the chart, $U=1$ on the combs.

\begin{enumerate}\item[\textup{(o)}] Let $b$ lie inside $B=B(y)$. Then $e^{i\eta_b}=U_b$, and $U_b$ is the variable of the bond $b$ when the box $B$ is put in its comb gauge with origin $y$; so its angle is at most $s_B$, by \Cref{Y-F-L3}.

\item[\textup{(i)}] Let $b=(z,\lambda)$ be a crossing bond of $c$, and put $z'=z+Le_\lambda$ and $h_{c,z}=U([z,z'])$. Then \[e^{i\eta_b}=U_bV'(c)^{-1}=[U_bh_{c,z}^{-1}][h_{c,z}\mathsf M(U)(c)^{-1}][\mathsf M(U)(c)V'(c)^{-1}].\]

The first bracket is conjugate to the closed path that leads from $z+e_\lambda$ back along its comb to $c_+$, then along the comb to $z'$, and then straight back. It consists of $3+3+1=7$ straight segments in $B(c_+)$, so the angle of the first bracket is at most $7s_c$ by \Cref{Y-F-L3}.

The second bracket equals $h_{c,z}h_{c,y}^{-1}e^{-i\bar\tau_c}$ with $y=c_-$. Its angle is at most $s_c+|\bar\tau_c|\le2s_c$, because $|\tau_{c,x}|\le\operatorname{ang}(h_{c,x}h_{c,y}^{-1})\le s_c$ for every $x\in B(y)$, by \Cref{Y-F-L2} and the convention $\mathsf L(-1)=0$.

The angle of the third bracket is $\mathsf n_c$. Hence $\operatorname{ang}(e^{i\eta_b})\le9s_c+\mathsf n_c$, and $\min\{1,9s_c+\mathsf n_c\}\le9s_c+\min\{1,\mathsf n_c\}$.\end{enumerate}

As a result, $\operatorname{ang}_b\le s_B^D+s_B^t$ inside a block $B$, and $\operatorname{ang}_b\le9s_c^D+9s_c^t+\min\{1,\mathsf n_c\}$ at a crossing bond of $c$.

For a bond $b$ of direction $\mu$ and a displacement $u$, at most $|u_\mu|$ sites $y$ have $b\in\Gamma(y,y+u)$. With three plaquettes to a site, this gives \[\mu_\ell:=\sup_b\sum_{(y,p):b\in\Gamma}|\ell_y(p)|\le3K\mathsf g_1,\] and likewise $\sum_y|\ell_y(p)||\Gamma|\le K\sum_u|u|_1e^{-\lambda_*|u|_\infty}\le3K\mathsf g_1$.

The part with $D$ is at most $2\varepsilon\mu_\ell(3L^3+9\cdot6L^2)\sum_D\phi\le K\mathsf g_1(18L^3+324L^2)\varepsilon\sum_D\phi$, by $\mu_\ell\le3K\mathsf g_1$.

For the tame part we use $2z_pa\le z_p^2+a^2$, which holds for any number $a\ge0$; we put $a=\operatorname{ang}_b^t$, the tame term in the bound on $\operatorname{ang}_b$, that is, $\operatorname{ang}_b^t:=s_B^t$ at a bond $b$ inside a block $B$ and $\operatorname{ang}_b^t:=9s_c^t$ at a crossing bond $b$ of $c$. We obtain $\sum|\ell||\Gamma|z_p^2\le3K\mathsf g_1[\sum_{\rm tame}\phi^2+\varepsilon\sum_D\phi]$ and \[\mu_\ell\sum_b(\operatorname{ang}_b^t)^2\le3K\mathsf g_1[3L^3\cdot3L^3+81\cdot6L^3\cdot L^2\cdot6]\sum_{\rm tame}\phi^2.\]

For the noise at bonds with $\mathsf n_c<\pi/3$, the inequality $2z_p\min\{1,\mathsf n_c\}\le z_p^2+\mathsf n_c^2$ gives once more $3K\mathsf g_1\sum z_p^2$, together with $3K\mathsf g_1L^2\mathsf n_c^2$ for each $c$.

For the noise at bonds with $\mathsf n_c\ge\pi/3$, where $\mathsf n_c>1$, the term is at most $2\varepsilon\,\mathsf n_c^2$, and we obtain $6K\mathsf g_1L^2\varepsilon\,\mathsf n_c^2\le3K\mathsf g_1L^2\mathsf n_c^2$.

Adding all parts, we find \[\sum_y|\mathsf{dl}_y|\le K\mathsf g_1(18L^3+324L^2+6)\varepsilon\sum_D\phi+3K\mathsf g_1(2+9L^6+2916L^5)\sum_{\rm tame}\phi^2+3K\mathsf g_1L^2\sum_{c\in\mathfrak Z}\mathsf n_c^2,\] and the right side is at most \[127K\mathsf g_1L^3\varepsilon\sum_D\phi+2949K\mathsf g_1L^6\sum_{\rm tame}\phi^2+3K\mathsf g_1L^2\sum\mathsf n_c^2.\]

The last step holds for $L\ge3$, where $18L^3+324L^2+6\le126.23L^3$ and $2+9L^6+2916L^5\le981.01L^6$; moreover $2949=3\cdot983$.\end{enumerate}

\emph{\textup{(5)} Collecting the bounds.} A bond of $\mathfrak S$ with $\mathsf n_c\ge\pi/3$ lies in $\mathsf{Cb}$; hence the sum over these bonds in \textup{(4)} runs over bonds of $\mathsf{Cb}$, and we add it to the sum over $\mathsf{Cb}$ in \textup{(3)}. By \textup{(Sup)}, \textup{(3)}, \textup{(4)} and \textup{(4$'$)}, the sum of the term $|\nu|_\infty\sum_{D\cap\hat{\mathfrak Z}}\phi$ from \textup{(1)}, of the site terms of \textup{(Ta)}, \textup{(Tb)} and \textup{(Tc)} at $y\in\hat{\mathfrak Z}'$ and of $|\langle\mathsf{zf},\mathrm{de}\rangle_{\mathfrak S}|+|\langle\Xi,\mathrm{db}\rangle_{\mathfrak S}|$ is at most \[c_{\mathrm{loss}}\varepsilon\sum_{D\cap\hat{\mathfrak Z}}\phi+\vartheta_F^\flat\sum_{\rm tame}\phi^2+\vartheta_n^\flat\sum_{c\in\mathfrak Z}\mathsf n_c^2+z_b\sum_{\mathsf{Cb}}\mathsf n_c,\] with $(20\mathsf N_Z+2L^2\varepsilon \mathsf N_\Xi)L^{-1/2}\varepsilon\le z_b$ by \textup{(F1)}, whose third clause, together with $\mathsf N_\Xi=c_2^*\mathsf N_Z$ from \textup{(O3)}, gives $2L^2\varepsilon\mathsf N_\Xi\le\mathsf N_Z$, and with \[c_{\mathrm{loss}}=[\mathsf N_\nu^\flat L^{-1}+2200\mathsf N_Z+40\mathsf N_\Xi]L^{-1/2}+l_r\mathsf N_\sigma L^{-3/2}(2\varrho(x)+3)^3\varepsilon+127K\mathsf g_1\mathsf N_\sigma L^{3/2}\varepsilon\le(\frac3{16}+\frac1{16}+\frac1{16})d_0.\] The three terms of $c_{\mathrm{loss}}$ are at most $\frac3{16}d_0$, $d_0/16$ and $d_0/16$, by \textup{(L1$^\flat$)}, \textup{(Fm)} and \textup{(Fc)} in this order. Here $\mathsf N_\nu^\flat$ collects $N_1$, which comes from $|\nu|_\infty$, and $(2l_c+(2\rho_0+1)^3l_r)\mathsf N_\sigma$.

On the tame plaquettes, $\phi_p^2\le2|F_p-\nu_p|^2+2|\nu_p|^2$. On the bonds with $\mathsf n_c\le t_0$, $\mathsf n_c^2\le2|n_c-\tilde\omega_c|^2+2.1|\omega_c|^2$. On the bonds $c\in\mathsf{Cb}$ with $\mathsf n_c\le t_0$ we have in addition $z_b\mathsf n_c\le(a/8)|n_c-\tilde\omega_c|^2+2z_b^2/a+1.01z_b|\omega_c|$. This follows from $z_b\mathsf n_c\le z_b|n_c-\tilde\omega_c|+z_b|\tilde\omega_c|$, which uses $|n_c|=\mathsf n_c$, valid when $\mathsf n_c\le t_0$ because $t_0=\varepsilon^{1/2}<\pi$, and from $|\tilde\omega_c|=|\omega_c|/(1-\theta_2)\le1.01|\omega_c|$, as at the start of the proof. At such a bond, the coefficient of $|n_c-\tilde\omega_c|^2$ is $(a/2)(1-\theta_2)$ in \textup{(2)}, and what remains of it after we subtract $2\vartheta_n^\flat$ and $a/8$ is $(a/2)(1-\theta_2)-2\vartheta_n^\flat-a/8\ge a/4$, by \textup{(F4$^\flat$)} and $\theta_2\le3.5\cdot10^{-5}$. On the bonds with $\mathsf n_c>t_0$, $\vartheta_n^\flat\mathsf n_c^2+z_b\mathsf n_c\le0.05a\,\mathsf n_c^2+0.05a\,\mathsf n_c^2+5z_b^2/a$. The cost of a bond of $\mathsf{Cb}$ consists of the terms without $n_c$ on the right side of the bound above on $z_b\mathsf n_c$ that applies to it: of $2z_b^2/a+1.01z_b|\omega_c|$ if $\mathsf n_c\le t_0$, and of $5z_b^2/a$ if $\mathsf n_c>t_0$. Since $z_b=21\mathsf N_ZL^{-1/2}\varepsilon$ by item~F of \S\ref{Y-F-1}, and $aL\ge18$, we have $5z_b^2/a\le(5\cdot441/18)\mathsf N_Z^2\varepsilon^2=122.5\mathsf N_Z^2\varepsilon^2$ and $2z_b^2/a\le49\mathsf N_Z^2\varepsilon^2$, and \textup{(Sup)} gives $1.01z_b|\omega_c|\le4.72\mathsf N_Z^2\varepsilon^2$. So in both cases the cost is at most $123\mathsf N_Z^2\varepsilon^2$.

The bonds of $\mathrm{Pl}_1$ number at most $4\#\mathrm{Pl}_1\le(4\cdot12\cdot24/\pi)\sum_D\phi\le366.694\sum_D\phi$. Their cost is thus at most $45104\mathsf N_Z^2\varepsilon\cdot\varepsilon\sum_D\phi\le\frac18d_0\varepsilon\sum_D\phi$, by $8\cdot45104\le4\cdot10^5$ and \textup{(F3$^\flat$)}. For $P\in\mathrm{Pl}_2$ we have $\mathrm{na}_P\ge\pi/3$, and $\mathrm{na}_P$ is the sum of $\mathsf n_c$ over the four bonds $c$ of $P$; hence $\mathsf n_c\ge\pi/12$ at one of these bonds, and $\pi/12>t_0$. At that bond, $\mathsf{et}_c\ge0.2a\,\mathsf n_c^2$ by \textup{(2)}; when we subtract from it the terms $0.05a\,\mathsf n_c^2+0.05a\,\mathsf n_c^2$ of the bound above on the bonds with $\mathsf n_c>t_0$, the term $0.1a\,\mathsf n_c^2$ is left, and $0.1a\,\mathsf n_c^2\ge0.1233/L$ as $aL\ge18$. From this term we subtract the costs of the bonds of every $P\in\mathrm{Pl}_2$ of which $c$ is a bond. A bond lies in four plaquettes $P$, and each of them has four bonds, so these are at most sixteen bonds, and as $1968=16\cdot123$, their cost is at most $1968\mathsf N_Z^2\varepsilon^2\le0.0615/L$ by \textup{(F3$^\flat$)}. This cost is at most $0.05a\,\mathsf n_c^2$, because $0.05a\,\mathsf n_c^2\ge0.06168/L$, and so the term $0.05a\,\mathsf n_c^2$ is left at that bond. The coefficient of $\varepsilon\sum_{D\cap\hat{\mathfrak Z}}\phi$ that is left, once the three terms of $c_{\mathrm{loss}}$ and the $\frac18d_0$ for the bonds of $\mathrm{Pl}_1$ are subtracted from the $d_0$ of \textup{(1)}, is at least $d_0(1-\frac3{16}-\frac1{16}-\frac1{16}-\frac18)=\frac9{16}d_0\ge\frac12d_0$, and $\frac12d_0$ is its coefficient in $\mathrm{Pay}_{\mathfrak Z}$. At the sites with $m_y=1$, the squares $\frac12|\ell^U_y\zeta-\sigma_y|^2$ enter $\mathrm{Sq}_{\mathfrak Z}$ with the full coefficient $\frac12$; none of the bounds in \textup{(3)} to \textup{(5)} uses any part of them.

For the starred fields, the contributions are $L\cdot2\cdot3L^3N_1^2L^{-4}=6N_1^2$ and \[L\cdot3\cdot16\mathsf N_Z^2(aL)^{-2}(a\theta_2+2.1\vartheta_n^\flat)\le48\mathsf N_Z^2(\theta_2/18+2.1L\vartheta_n^\flat/324)\le3\mathsf N_Z^2(\theta_2+L\vartheta_n^\flat).\]\end{proof}

\emph{Step 4: tame windows at real fields.}

\begin{lemma}[Tame windows at real fields: centre, remainder terms and prefactor]\label{Y-F-LA} In the setting of \Cref{Y-F-C7}, the following hold.

\begin{enumerate}\item[\textup{(a)}] At every $(V',e)$, $|a_{0,b}|\le a_\infty^\flat\sup|X'|\le\pi a_\infty^\flat$; the matrix $u^T=SC$ has range at most $R+2\varrho_L+2$; and $|(u^T\mathsf g^VX')_b|\le x^{-50}$.

\item[\textup{(b)}] If $V'$ is tame on the box of radius $R_5$ about $i$, then $|\bar F_p-\nu^{\mathrm{ref}}_p|\le K_D\varepsilon'^2$, $|r_{*p}|\le K_r\varepsilon'^2$, $|\Delta_p|\le(K_D+K_r)\varepsilon'^2$, $|\Delta_c|\le K_D\varepsilon'^2$, $|a_{0,b}|\le\mathsf a_{\rm c}\varepsilon'$ and $|\bar F_p|\le(\mathsf N_\sharp+2)L^{-2}\varepsilon'$. If the first of these bounds, $|\bar F_p-\nu^{\mathrm{ref}}_p|\le K_D\varepsilon'^2$, holds at all $p$ within $\varrho(x)$ of $y$, then $|\Delta_y|\le l_rK_D\varepsilon'^2$. Consequently, if $V'$ is tame on the box of radius $R_6$ about $i$, then $|\Delta_i|\le K_\Delta^\flat\varepsilon'^2$ for rows $i$ of all three kinds.

\item[\textup{(c)}] If all $P$ within $R_6$ of $\square$ are tame, then $|T_\square|\le K_{\rm q}^\flat p_0'^2\varepsilon'$.\end{enumerate}\end{lemma}

\begin{proof}[Proof of \Cref{Y-F-LA}] We begin with \textup{(a)}. Since $\|C\|_\infty\le\|G\|_\infty\|G^{\mathsf T}\|_\infty\le \mathsf G_0^2$ and $|\mathsf g^VX'|\le(b_h/4)\sup|X'|$, the centre $a_0=C\mathsf g^VX'$ obeys the first bound of \textup{(a)}. Next, $u^T=1-qC=SC$, and $1-qC$ has range at most $2r_w+r_q\le R+2\varrho_L+2$. A row of $u^T$ contains at most $3L^3(2R+4\varrho_L+5)^3$ blocks, each of norm at most $\|S\|\|C\|\le(2/\gamma)x^{-60}$; hence \[|(u^T\mathsf g^VX')_b|\le3L^3(2R+4\varrho_L+5)^3(2/\gamma)(b_h/4)\pi x^{-60}\le x^{-52}\] by \textup{(F9$^\flat$)}.

For \textup{(b)} we work in the comb gauge at $o=y(i)$. Let $\check\gamma(y)=V'(\Gamma(o,y))$ and $\tilde X_P:=\operatorname{Ad}_{\check\gamma(y_P)}X'_P=T(o,y_P)X'_P$; then $|\tilde X_P|=|X'_P|\le\varepsilon'$.

The radius $\rho$ of the gauge enters, at real fields, only through three facts:

\begin{enumerate}\item[\textup{(i)}] the embedding of the box of radius $\rho$ about $o$ in the torus, which \Cref{Y-F-L3} requires;

\item[\textup{(ii)}] the bound $\|T^{\check\gamma}(y,y'')^{\pm1}-1\|\le130rt\varepsilon'\le0.2$ for $y,y''$ within $t\le\rho$ of $o$ at distance $r$, so that $r\le2t$, where $T^{\check\gamma}$ are the transports of $V'^{\check\gamma}(c)=\check\gamma(c_-)V'(c)\check\gamma(c_+)^{-1}$;

\item[\textup{(iii)}] in \Cref{Y-F-L20} below, the bound $\rho$ on the sum of the ranges.\end{enumerate}

We take $\rho=R_5$; by the hypothesis of \textup{(b)}, $V'$ is tame on the box of this radius about $i$, and $R_5\le Mr(x')$ by \textup{(G1$^\flat$)}. At this radius the three facts hold, as follows.

\begin{enumerate}\item[\textup{(i)}] $2R_5+1\le2Mr(x')+1<16Mr(x)\le N$, by \textup{(Tor)}.

\item[\textup{(ii)}] If both ends of a bond lie within $t$ of $o$, \Cref{Y-F-L3} with origin $o$, applied to two ladders of at most $t$ plaquettes, gives $\operatorname{ang}V'^{\check\gamma}(c)\le2t\varepsilon'$; so $\|\operatorname{Ad}^{\pm1}_{V'^{\check\gamma}(c)}-1\|\le4t\varepsilon'$ by \Cref{Y-F-L1}\,(b). The path $\Gamma(y,y'')$ has at most $3r$ bonds, all within $t$ of $o$; hence the left side of the bound in \textup{(ii)} is at most $(1+4t\varepsilon')^{3r}-1\le e^{12rt\varepsilon'}-1$. By the second clause of \textup{(F8)}, $12rt\varepsilon'\le24(Mr(x'))^2\varepsilon'\le0.03$, so the left side is at most $12.2rt\varepsilon'\le130rt\varepsilon'$ and at most $0.031\le0.2$.

\item[\textup{(iii)}] The sums of the ranges that occur below are at most $4R+2\varrho_L+3\le R_5\le Mr(x')$, by \textup{(G1$^\flat$)}.\end{enumerate}

Moreover $T^{\check\gamma}(o,y)=1$, and every dressed kernel satisfies $K^{V^{\check\gamma}}(i,j)=\operatorname{Ad}_{\check\gamma(y(i))}K^V(i,j)\operatorname{Ad}_{\check\gamma(y(j))}^{-1}$. Indeed, products along paths telescope, $V'^{\check\gamma}(\pi)=\check\gamma(\pi_-)V'(\pi)\check\gamma(\pi_+)^{-1}$, by the argument in the proof of \Cref{Y-C-CIL64}\,(d). At $\pi=\Gamma(o,y)$ this equals $\check\gamma(o)\check\gamma(y)\check\gamma(y)^{-1}=1$, because $\check\gamma(o)=1$; at $\pi=\Gamma(y(i),y(j))$ we take $\operatorname{Ad}$ and multiply by $K(i,j)$.

The norms are $\|\cdot\|_{\theta_D}$. We have $\|\mathsf T^<\|,\|\mathsf X^<\|,\|\ell\|,\|\ell'^{\rm c}\|\le C_D$; further $\|L^{-1}W^*\|\le\mathrm eL$, because $\sum_pW(P,p)=L^2$, $\sum_PW(P,p)=L^{-1}$ and the range is $1$; and $\|\mathrm{d}^*\|\le6\mathrm e$, $\|\mathrm{d}'^*\|\le4\mathrm e$.

For each coarse plaquette $P$ and each coarse site $y'$ occurring in the quantities bounded in \textup{(b)}, we have $\zeta'_P=X'_P$ and $m'_{y'}=1$, so that $v_*=(1+\ell'^{{\rm c}V*}\ell'^{{\rm c}V})X'$.

\emph{Item 1.} By \Cref{Y-C-Th68}\,(b), $|\bar F_p-(H^\natural\tilde X)_p|=|\Lambda_p|\le(\mathsf K_{\rm c}^0+1)\varepsilon'^2$, since $\tau_{\rm c}\le x^{-59}\le\varepsilon'$.

\emph{Item 2.} By \Cref{Y-C-L71}\,(b) and \Cref{Y-F-LR}, $\|H^\natural-\mathsf T^<\mathsf B'^{\rm c}\|_\infty\le\mathsf{ti}_{\rm c}+(C_D^2+2C_D^3)x'^{-64}\le\varepsilon'$, where the last inequality is \textup{(FA)}.

\emph{Item 3.} By \Cref{Y-F-L20}, which we state and prove now, replacing a dressed product of $n$ kernels by the free one costs at most $1040\theta_D^{-2}\varepsilon'\,n\prod1.2\|K_k\|$, times $\sup|\tilde X|$.

\begin{lemma}[Dressed and free products of kernels in the comb gauge]\label{Y-F-L20} In the setting of \Cref{Y-F-C7} and in the comb gauge at $o$ of radius $\rho$ as above, let $K_1,\dots,K_n$ be kernels of the free model, each of finite range, and let the sum of their ranges be at most $\rho$. We assume that the entries of these kernels are scalars, or vectors of $\mathsf H_y$ or $\mathsf H'_{y'}$, and that the vectors occur in the pairs $\ell^*\ell$, $\ell'^{{\rm c}*}\ell'^{\rm c}$. An entry contributes its modulus $|K_k(i,j)|$ to the norms $\|K_k\|_\theta$. Let $A_k$ be $K_k$ or its adjoint, dressed with $V'^{\check\gamma}$, put $A_k^0=K_k\otimes1$ in the first case and $A_k^0=K_k^*\otimes1$ in the second, let $y(i_0)=o$, and let $0<\theta\le1$. Then $\sum_j\|(A_1\cdots A_n-A_1^0\cdots A_n^0)(i_0,j)\|\le130\varepsilon'\,n(8/\theta^2)\prod_k1.2\|K_k\|_\theta$.\end{lemma}



\begin{proof}[Proof of \Cref{Y-F-L20}] We write $\prod A-\prod A^0=\sum_kA_1\cdots A_{k-1}(A_k-A_k^0)A^0_{k+1}\cdots A^0_n$ and expand each term along chains of indices. For a chain $i_0,\dots,i_n$ of length $\mathsf{le}=\sum_k{\rm dist}(i_{k-1},i_k)\le\rho$, each $y(i_k)$ is within $\mathsf{le}$ of $o$. By \textup{(ii)}, therefore, $\|A_k(i,j)\|\le1.2|K_k(i,j)|$ and $\|(A_k-A_k^0)(i_{k-1},i_k)\|\le|K_k(i_{k-1},i_k)|130\mathsf{le}^2\varepsilon'$. From $u^2\le2e^u$ we get $\mathsf{le}^2\le(8/\theta^2)e^{\theta\mathsf{le}/2}$, and the sum over the chains of $\prod|K_k|e^{\theta\mathsf{le}/2}$ is bounded by $\prod\|K_k\|_\theta$. Here, whenever $A_k$ is the dressed adjoint, $|K_k(i,j)|$ denotes, in the two bounds derived from \textup{(ii)} and in this sum, the modulus of the entry $K_k^*(i,j)$ of the adjoint kernel. Since $\|\cdot\|_\theta$ is the larger of the weighted row and column sums (\Cref{Y-A-Norm}), we have $\|K_k^*\|_\theta=\|K_k\|_\theta$. The $n$ terms of the sum over $k$ give the factor $n$.\end{proof}

We return to item 3. By covariance and $\check\gamma(o)=1$, we have $(K^VX')_{i_0}=(K^{V^{\check\gamma}}\tilde X)_{i_0}$. For $(\mathsf T^<)^V$ and $(\mathsf T^<\ell'^{{\rm c}*}\ell'^{\rm c})^V$ the lemma gives $1040\theta_D^{-2}(1.2C_D+5.184C_D^3)\varepsilon'^2\le5470\theta_D^{-2}C_D^3\varepsilon'^2$, for $C_D\ge4$. Adding the bounds of items 1 to 3 gives the constant $K_D$.

\emph{Item 4.} We have $r_*=[(1+\ell^*\ell)\mathsf T^<-L^{-1}W^*-\mathrm{d}^*\mathsf X^<]^V(1+\ell'^{{\rm c}*}\ell'^{\rm c})^VX'$. For the products that occur, the numbers $n\prod1.2\|K_k\|$ are $1.2C_D$ for $\mathsf T^<$, $3\cdot1.728C_D^3$ for $\ell^*\ell\mathsf T^<$, $3\cdot1.728C_D^3$ for $\mathsf T^<\ell'^{{\rm c}*}\ell'^{\rm c}$, $5\cdot2.48832C_D^5$ for $\ell^*\ell\mathsf T^<\ell'^{{\rm c}*}\ell'^{\rm c}$, $2\cdot1.44\cdot6\mathrm eC_D$ for $\mathrm{d}^*\mathsf X^<$, $4\cdot2.0736\cdot6\mathrm eC_D^3$ for $\mathrm{d}^*\mathsf X^<\ell'^{{\rm c}*}\ell'^{\rm c}$, $1.2\mathrm eL$ for $L^{-1}W^*$, and $3\cdot1.728\mathrm eLC_D^2$ for $L^{-1}W^*\ell'^{{\rm c}*}\ell'^{\rm c}$.

Their sum is $r(C_D)C_D^5+\mathrm e(1.2C_D^{-2}+5.184)LC_D^2$, where \[r(C)=[1.2C+10.368C^3+12.4416C^5+6\mathrm e(2.88C+8.2944C^3)]C^{-5}\] is a decreasing function with $r(4)\le21.733$, and where $\mathrm e(1.2/16+5.184)\le14.296$. Multiplied by $1040$, the sum is at most $22603C_D^5+14868LC_D^2\le23000(C_D^5+LC_D^2)$.

By \Cref{Y-F-LR}\,(ii), the free product is $[\mathsf B(\mathsf T^<-\mathsf T)-\mathrm{d}^*(\mathsf X^<-\mathsf X)]\mathsf B'^{\rm c}$. Its sup norm is at most $(\mathsf N_{\mathsf B}+2)C_Dx'^{-64}\mathsf N_Z$, and multiplied by $\varepsilon'$ it is at most $\varepsilon'^2$, by \textup{(FA)}. This yields $K_r$. Every coarse plaquette $P$ on which $r_{*p}$ depends, by way of $X'_P$, $\zeta'_P$ or $m'$, lies within $2(\varrho_L+1)+3R(x')+R+1\le R_5$ of $i$, with $i$ as in the hypothesis of \textup{(b)}.

\emph{Item 5.} The relation $\nu_*=\nu^{\mathrm{ref}}-r_*$ gives the bound on $|\Delta_p|$. Next, $\Delta_c=(Q_1a_0)_c-\omega_{*c}$, and \Cref{Y-C-Th68}\,(b) for $\Lambda_c$ gives \[|(Q_1a_0)_c-\sum_P\mathsf N^\natural(c,P)T(c_-,y_P)X'_P|\le(\mathsf K_{\rm c}^0+1)\varepsilon'^2.\]

Further, $\|\mathsf N^\natural-\mathsf N^\sharp\|_\infty\le\mathsf{ti}_{\rm c}$, and $\|\mathsf N^\sharp-(aL)^{-1}\mathrm{d}'^*\mathsf B'^{\rm c}\|_0\le(4/18)2C_D^2x'^{-64}$; multiplied by $\varepsilon'$, the two are together at most $\varepsilon'^2$, by \textup{(FA)}. \Cref{Y-F-L20} for $(aL)^{-1}\mathrm{d}'^{V*}(1+\ell'^{{\rm c}V*}\ell'^{{\rm c}V})$ gives $1040\theta_D^{-2}(4\mathrm e/18)(1.2+5.184C_D^2)\varepsilon'^2\le3304\theta_D^{-2}C_D^2\varepsilon'^2\le6700\theta_D^{-2}C_D^3\varepsilon'^2$. The sum of these bounds is at most $K_D\varepsilon'^2$.

The site term $\Delta_y=(\ell^V(\bar F_F-\nu^{\mathrm{ref}}))_y$ is bounded by means of the row sums of $\ell^V$. The bound on $|a_{0,b}|$ is given by \Cref{Y-C-Th68}\,(c). Finally, $|\bar F_p|\le\|H^\natural\|_\infty\varepsilon'+\mathsf K_{\rm c}\varepsilon'^2\le(\mathsf N_\sharp+1)L^{-2}\varepsilon'+L^{-2}\varepsilon'$, by \Cref{Y-C-L71}\,(c) and \textup{(FA)}. The site rows use the $p$ within $\varrho_L+1$, and $R_5+\varrho_L+1\le R_6$.

We turn to \textup{(c)}. At a tame $P$ with $m'=1$ on its reach, $(x'/2)\langle X',v_*\rangle_P-x'\mathrm{En}'_P=0$. Moreover \[|\Delta_Y|=|\mathsf B^V(\bar F_F-\nu^{\mathrm{ref}})+r_*|\le(\mathsf N_{\mathsf B}K_D+K_r)\varepsilon'^2,\] where the windows at the $p'$ within $2(\varrho_L+1)$ enter; this accounts for the radius $R_6$. The operator $E^*$ sums $L$ edge plaquettes, a cell holds $3M^3$ plaquettes $P$, and $x\varepsilon'^2=Lp_0'^2$; so this part contributes $1.5M^3L^2(\mathsf N_{\mathsf B}K_D+K_r)p_0'^2\varepsilon'$.

For the corner cubes we use Bianchi's identity in the real case, as in the proof of \Cref{Y-C-L67}\,(d) and in Step~5 of the proof of \Cref{Y-D-C5}: $\prod_{k=1,\dots,6}e^{iY_k}=1$ with $Y_k=\pm\operatorname{Ad}_{\kappa_k}X'_{P_k}$. Here $\kappa_k=1$ at the faces based at $z=\pi(q')$ and $\kappa_k=V'(z,\lambda)$ at the face based at $z+Le_\lambda$, and the signs are $\mathrm{d}'(q',P_k)$.

Thus $\sum_kY_k=(\mathrm{d}'^VX')_{q'}$ with $|Y_k|\le\varepsilon'$, and \Cref{Y-F-L1}\,(c$'$), applied with $S\le6\varepsilon'$, gives $|\sum_kY_k|\le\frac12S^2$, which is at most $18\varepsilon'^2$. A $q'$ has one corner cube, so this part contributes $M^3(x/2)18\varepsilon'^2\mathsf N_\Xi L^{-1}\varepsilon'=9M^3\mathsf N_\Xi p_0'^2\varepsilon'$.

The last sum is at most $(x/2)3v\mathsf a_{\rm c}\varepsilon'x^{-50}$, and $v\mathsf a_{\rm c}\le4.42v\gamma^{-1/2}b_h\mathsf G_0\le x^8$ by \textup{(F9$^\flat$)}. As $\varepsilon\le\varepsilon'$, $c_\xi p_0\ge1$, $c_\xi\le1/6.5$ and $x\ge L^2$, we have $p_0'^2\ge p_0^2/L\ge42x^{-1/2}$; hence the last sum is at most $1.5x^{-41}\varepsilon'\le M^3p_0'^2\varepsilon'$.\end{proof}

\emph{Step 5: a lower bound for $\mathcal G_{\mathfrak Z}$.}

\begin{lemma}[Lower bound for $\mathcal G_{\mathfrak Z}$ in terms of the coarse energy and the rough coarse plaquettes]\label{Y-F-L23k} In the setting of \Cref{Y-F-C7}, let $(J,t)$ be an admitted cell datum. Then \[\mathcal G_{\mathfrak Z}\ge\sum_{P\in\mathfrak Z}\mathrm{En}'_P+(\frac12-d_0-\mathsf t_\rho-3e_1^\flat)\varepsilon'\sum_{D'_J}\phi'_P-(9e_1^\flat+\frac92\mathsf N_Z^2c_B\mathsf{ti}_0^\flat)\varepsilon'^2M^3\#\mathfrak Z.\]\end{lemma}

\begin{proof}[Proof of \Cref{Y-F-L23k}] In this step alone, $\ell$ stands for $\ell'^{{\rm c}V}$, except in $\nu^<$, in $\sigma^<$ and in the uncut $\nu$ and $\sigma$ of \textup{(4)} below, where $\ell$ has the fine rows. We put $B_1^V:=1+\ell^*\ell$, $\mathsf{Hm}:=B_1^V\hat R^V-1$ and $\hat R^<:=L(\mathrm{pm}^<)^*(\mathrm{pm}^<)$, where $\mathrm{pm}^<$ is the pair map at $V'=1$; its three components are $\nu^<=L^{-1}W^*+\mathrm{d}^*\mathsf X^<-\ell^*\ell\mathsf T^<$, $\sigma^<=\ell\mathsf T^<$ and $(aL)^{-1}\mathrm{d}'^*$, and in the first two of them $\ell$ has the fine rows.

The operator $\hat R^V$ is symmetric, since $\hat R^V=L\,\mathrm{pm}^*\mathrm{pm}$. The ranges are at most $R(x')+2\varrho_L+2$ for $\mathrm{pm}$, $2R(x')+4\varrho_L+4$ for $\hat R^V$, and $4R(x')+4\varrho_L+4$ for $\mathsf{Hm}$.

\emph{\textup{(1)} The quadratic term.} Let $\mathsf{ww}:=(1-m')\ell\zeta'$ and $\mathsf{uu}:=\ell^*\mathsf{ww}$. Then $v_*=B_1^V\zeta'-\mathsf{uu}$, and $\hat R^Vv_*=\zeta'+\mathsf{Hm}^*\zeta'-\hat R^V\mathsf{uu}$.

A site with $m'=0$ lies within $R(x')+1$ of $D'$. Hence $\mathsf{uu}$ lives within $R(x')$ of such sites, and $\hat R^V\mathsf{uu}$ within $4R(x')+4\varrho_L+5$ of $D'$.

We split $\mathsf{uu}=\mathsf{uu}^J+\mathsf{uu}^o$ as follows: $\mathsf{ww}^J:=\mathsf{ww}$ at those sites with $m'=0$ that lie within $R(x')+1$ of $D'_J$, and $0$ at every other site; $\mathsf{ww}^o:=\mathsf{ww}-\mathsf{ww}^J$; $\mathsf{uu}^J:=\ell^*\mathsf{ww}^J$ and $\mathsf{uu}^o:=\ell^*\mathsf{ww}^o$. A site at which $\mathsf{ww}^o$ does not vanish has $m'=0$ and is not within $R(x')+1$ of $D'_J$. Since every site with $m'=0$ is within $R(x')+1$ of $D'$, such a site is within $R(x')+1$ of $D'\setminus D'_J$. No site lies within $R(x')+1$ of $D'_J$ and of $D'\setminus D'_J$ at once: a site within $R(x')+1$ of both $D'_J$ and $D'\setminus D'_J$ would place the latter within $2R(x')+2$ of $\mathfrak Z_2$, in contradiction with \Cref{Y-F-Geok}\,(i). By \Cref{Y-F-Geok}\,(iii) and \textup{(v)}, $\hat R^V\mathsf{uu}^o=0$ on $\mathfrak Z$, while $\mathsf{uu}^J$ and $\hat R^V\mathsf{uu}^J$ live in $f_J\subset\mathfrak Z$.

It follows that \[\langle v_*,\hat R^V\mathsf{uu}\rangle_{\mathfrak Z}=\langle\hat R^Vv_*,\mathsf{uu}^J\rangle=\langle\zeta'+\mathsf{Hm}^*\zeta'-\hat R^V\mathsf{uu}^J,\mathsf{uu}^J\rangle,\] and therefore \[\langle v_*,\hat R^Vv_*\rangle_{\mathfrak Z}=\langle v_*,\zeta'\rangle_{\mathfrak Z}-\langle\zeta',\mathsf{uu}^J\rangle+\langle\mathsf{uu}^J,\hat R^V\mathsf{uu}^J\rangle+\langle v_*,\mathsf{Hm}^*\zeta'\rangle_{\mathfrak Z}-\langle\mathsf{Hm}^*\zeta',\mathsf{uu}^J\rangle.\]

Now $\hat R^V=(B_1^V)^{-1}(1+\mathsf{Hm})$, and $\langle\ell^*\mathsf{ww}^J,(1+\ell^*\ell)^{-1}\ell^*\mathsf{ww}^J\rangle\le|\mathsf{ww}^J|^2=\langle\ell\zeta',\mathsf{ww}^J\rangle=\langle\zeta',\mathsf{uu}^J\rangle$, since $(1-m')^2=1-m'$. Hence $\langle v_*,\hat R^Vv_*\rangle_{\mathfrak Z}\le\langle v_*,\zeta'\rangle_{\mathfrak Z}+\mathrm{Hol}$, with the number \[\mathrm{Hol}=|\langle v_*,\mathsf{Hm}^*\zeta'\rangle_{\mathfrak Z}|+|\langle\mathsf{Hm}^*\zeta',\mathsf{uu}^J\rangle|+|\langle(B_1^V)^{-1}\mathsf{uu}^J,\mathsf{Hm}\,\mathsf{uu}^J\rangle|.\]

\emph{\textup{(2)} The inverse of $B_1^V$, and a bound for $\mathrm{Hol}$ through the row and column sums of $\mathsf{Hm}$.} The fields $v_*$ and $\mathsf{uu}$ obey $|v_*|,\ |\mathsf{uu}|\le \mathsf N_Z\varepsilon'$. We show that $\|(B_1^V)^{-1}\|_\infty\le c_B$. To invert $B_1^V$ we apply \Cref{Y-A-L3}\,(a) with $\gamma=1$, in the form of \Cref{Y-F-L5}, which we state first.

\begin{lemma}[Two bounds for symmetric matrices of $3\times3$ blocks]\label{Y-F-L5} The assertions of \Cref{Y-A-L3}\,(a) and~(b) remain valid for symmetric matrices whose entries are real $3\times3$ blocks, where $|\cdot|$ now denotes the operator norm of a block.\end{lemma}

\begin{proof} The proofs of \Cref{Y-A-L3}\,(a) and~(b), given after that lemma in \S\ref{Y-A-1}, rest on Schur's test and on eigenvalues alone. Both arguments go through unchanged for symmetric matrices with real $3\times3$ blocks as entries, once the modulus of each entry is replaced by the operator norm of the block.\end{proof}

The matrix $B_1^V$ is of this kind, and $B_1^V\ge1$. For real $V'$ all transports $T$ are isometries; hence, with $|\cdot|$ the operator norm of a block, \[|(B_1^V-1)(P,P')|\le\sum_{y'}|\ell_{y'}(P)|\,|\ell_{y'}(P')|.\] For $0<s\le\theta_D$ and $d\ge0$, convexity gives $e^{sd}-1\le(s/\theta_D)(e^{\theta_Dd}-1)$. Since $\|\ell'^{\rm c}\|_{\theta_D}\le C_D$, it follows that, for $s=\theta_D/(4C_D^2)$, \[\sup_P\sum_{P'}|(B_1^V-1)(P,P')|\,(e^{s{\rm dist}(P,P')}-1)\le(s/\theta_D)C_D^2=\frac14.\] Because $\frac14\le\frac12=\gamma/2$ at $\gamma=1$, the hypothesis $\varpi_{\mathsf A}(\lambda)\le\frac\gamma2$ of \Cref{Y-A-L3}\,(a) holds. \Cref{Y-F-L5} now yields $|(B_1^V)^{-1}(P,P')|\le2e^{-s{\rm dist}(P,P')}$, and as there are three plaquettes $P$ to a site, the bound $\|(B_1^V)^{-1}\|_\infty\le c_B$ follows. Consequently \[\mathrm{Hol}\le3\mathsf N_Z^2c_B\varepsilon'^2\max\{\sum_{P\in\mathfrak Z}\sum_{P'}|\mathsf{Hm}(P,P')|,\ \sum_{P\in\mathfrak Z}\sum_{P'}|\mathsf{Hm}(P',P)|\}.\]

\emph{\textup{(3)} Row and column sums of $\mathsf{Hm}$.} We split \[\mathsf{Hm}(P,P')=\mathrm{Hol}(P,P')+(\mathsf B'^{\rm c}\hat R^<-1)(P,P')\,T(y_P,y_{P'}).\] Within (3), $\mathrm{Hol}(P,P')$ and $\mathrm{Hol}(i_0,i_n)$ denote kernels, whereas outside (3) $\mathrm{Hol}$ is the number defined in (1). We bound the kernel $\mathrm{Hol}(P,P')$ by the following lemma, used at fixed $i_0$ for the row sums and at fixed $i_n$ for the column sums.

\begin{lemma}[The deviation of a product of dressed kernels from the bare product times one transport]\label{Y-F-L4} Let $V'$ be any $G$-valued field and let $\theta\ge0$. Let $K_1,\dots,K_n$ be kernels of finite range whose entries are scalars, or vectors of $\mathsf H_y$ and $\mathsf H'_{y'}$ that occur in pairs $\ell^*\ell$; only the moduli of the entries enter. For each $k$, choose $\tilde K_k=K_k^V$ or
$\tilde K_k=(K_k^V)^*$ and let $A_k^0=K_k$ or $K_k^*$,
respectively. Assume the chosen factors are composable. Define \[\mathrm{Hol}(i_0,i_n):=(\tilde K_1\cdots\tilde K_n)(i_0,i_n)-(A_1^0\cdots A_n^0)(i_0,i_n)\,T(y(i_0),y(i_n)).\] Assume that the sum of the ranges of $K_1,\dots,K_n$ is $<N/2-1$. Then \[\sum_{i_n}|\mathrm{Hol}(i_0,i_n)|\le6(n+1)\prod_k\|K_k\|_\theta\sum_{P''}e^{-\theta{\rm dist}(i_0,P'')}\phi'_{P''},\] and the column sums $\sum_{i_0}|\mathrm{Hol}(i_0,i_n)|$ obey the same bound with ${\rm dist}(i_n,P'')$ in place of ${\rm dist}(i_0,P'')$. In the bound on the row sums, $P''$ runs through the plaquettes with ${\rm dist}(i_0,P'')$ at most the sum of the ranges of $K_1,\dots,K_n$, and in the bound on the column sums through those with ${\rm dist}(i_n,P'')$ at most the same sum.\end{lemma}

\begin{proof} By the hypothesis on the ranges, the box spanned by the sites $y(i_k)$ embeds into the coarse torus. Expanding both products we obtain \[\mathrm{Hol}=\sum_{i_1,\dots,i_{n-1}}\prod_kA_k^0(i_{k-1},i_k)\,[\operatorname{Ad}_{V'(\text{chain})}-\operatorname{Ad}_{V'(\Gamma(y(i_0),y(i_n)))}].\] Here the chain is the path through $y(i_0),\dots,y(i_n)$ made of the coordinate paths $\Gamma(y(i_{k-1}),y(i_k))$, except that, where the factor $\tilde K_k$ is the adjoint of a dressed kernel, the path $\Gamma(y(i_k),y(i_{k-1}))$, run backwards, takes the place of $\Gamma(y(i_{k-1}),y(i_k))$. Each coordinate path has at most $3$ straight stretches, so that the chain has at most $3n$ of them and $\Gamma(y(i_0),y(i_n))$ at most $3$ more. Together with $\Gamma(y(i_0),y(i_n))$, the chain makes up a closed path in that box, with at most $3(n+1)$ straight stretches. By \Cref{Y-F-L3} and \Cref{Y-F-L1}\,(b), the bracket is therefore at most $6(n+1)\sum_{P''\in\text{box}}\phi'_{P''}$. For every plaquette $P''$ of the box, both ${\rm dist}(i_0,P'')$ and ${\rm dist}(i_n,P'')$ are bounded by $\mathsf{le}:=\sum_k{\rm dist}(i_{k-1},i_k)$. Hence, for fixed $i_0$ and $P''$, \[\sum_{i_1,\dots,i_n}\prod_k|A_k^0|\,\mathbf 1[\mathsf{le}\ge{\rm dist}(i_0,P'')]\le e^{-\theta{\rm dist}(i_0,P'')}\prod_k\|K_k\|_\theta,\] Here $\|A_k^0\|_\theta=\|K_k\|_\theta$, since the norm
controls both row and column sums. Summing over $P''$ gives the
row bound. For fixed $i_n$ and $P''$ the same estimate holds with column sums and with ${\rm dist}(i_n,P'')$, because $\|\cdot\|_\theta$ controls the column sums as well (\Cref{Y-A-Norm}); this gives the bound on the column sums.\end{proof}

Measured in $\|\cdot\|_{\theta_D}$, the $\nu$-map has norm at most $\mathrm eL+6\mathrm eC_D+C_D^3=a_\nu$ and the $\sigma$-map has norm at most $C_D^2$, while $(aL)^{-1}\|\mathrm{d}'\|^2\le16\mathrm e^2/18\le7$. Every product in $B_1^V\hat R^V$ has $n\le2+3+3=8$ factors, so that $6(n+1)\le54$, and the products of the norms add up to at most $(1+C_D^2)[La_\nu^2+LC_D^4+7]$. For each of these products, the sum of the ranges of the factors is $\le4R(x')+4\varrho_L+4<N/2-1$ by \textup{(Tor)}, since $N\ge16Mr(x)\ge320c_R\log x$. With $54\cdot3\mathsf g(\theta_D)=162\mathsf g(\theta_D)$ and the constant $C_{\rm hol}^\flat$ of item~F of \S\ref{Y-F-1}, \Cref{Y-F-L4} gives \[\sum_{P'}|\mathsf{Hm}(P,P')|\le\dfrac{C_{\rm hol}^\flat}{3\mathsf g(\theta_D)}\sum_{P''}e^{-\theta_D{\rm dist}(P,P'')}\phi'_{P''}+\|\mathsf B'^{\rm c}\hat R^<-1\|_0.\] For the column sums, the bound of \Cref{Y-F-L4} with ${\rm dist}(i_n,P'')$ gives the same right-hand side, \[\sum_{P'}|\mathsf{Hm}(P',P)|\le\dfrac{C_{\rm hol}^\flat}{3\mathsf g(\theta_D)}\sum_{P''}e^{-\theta_D{\rm dist}(P,P'')}\phi'_{P''}+\|\mathsf B'^{\rm c}\hat R^<-1\|_0,\] because $\|\cdot\|_0$ is the larger of the largest row sum and the largest column sum (\Cref{Y-A-Norm}).

\emph{\textup{(4)} Cuts.} In (4) alone we write $t:=x'^{-64}$. Without the cuts, $\nu$ and $\sigma$ denote the components of $\mathrm{pm}$ and $\hat R$ denotes $L\,\mathrm{pm}^*\mathrm{pm}$, at $V'=1$ and with $\mathsf T,\mathsf X$ in place of $\mathsf T^<,\mathsf X^<$. In (4), $\ell$ denotes the fine rows at $V'=1$, as in $\sigma^<$, and not $\ell'^{{\rm c}V}$. By definition, the second component of the uncut $\mathrm{pm}$ is $\sigma=\ell\mathsf T$. With $\mathsf B=1+\ell^{\mathsf T}\ell$, the first identity of \Cref{Y-F-LR}\,(ii), $\mathsf B\mathsf T=L^{-1}W^*+\mathrm{d}^*\mathsf X$, gives $\nu=L^{-1}W^*+\mathrm{d}^*\mathsf X-\ell^{\mathsf T}\ell\mathsf T=\mathsf B\mathsf T-\ell^{\mathsf T}\ell\mathsf T=\mathsf T$. Written out, $\hat R=L(\mathsf T^*\mathsf T+\mathsf T^*\ell^{\mathsf T}\ell\mathsf T)+La(aL)^{-2}\mathrm{d}'\mathrm{d}'^*=L\mathsf T^*\mathsf B\mathsf T+(aL)^{-1}\mathrm{d}'\mathrm{d}'^*$, and the last identity of \Cref{Y-F-LR}\,(ii) yields $\mathsf B'\hat R=1$. For the cut maps, $\|\nu^<-\mathsf T\|_0\le(6C_D+C_D^3)t\le a_\nu t$ and $\|\sigma^<-\ell\mathsf T\|_0\le C_D^2t$. Hence $\|\hat R^<-\hat R\|_0\le L[2a_\nu^2+2C_D^4]t\le2.125\,La_\nu^2t$, since $a_\nu^2\ge16C_D^4$. Further, $\|\hat R^<\|_0\le La_\nu^2+LC_D^4+7\le1.07\,La_\nu^2$, $\|\mathsf B'^{\rm c}-\mathsf B'\|_0\le2C_D^2t$ and $\|\mathsf B'\|_0\le C_D$. From $\mathsf B'^{\rm c}\hat R^<-1=(\mathsf B'^{\rm c}-\mathsf B')\hat R^<+\mathsf B'(\hat R^<-\hat R)$ we conclude \[\|\mathsf B'^{\rm c}\hat R^<-1\|_0\le(2.14C_D^2+2.125C_D)La_\nu^2t\le2.7\,C_D^2La_\nu^2t\le\mathsf{ti}_0^\flat.\]

\emph{\textup{(5)} The sum over $P\in\mathfrak Z$.} By \Cref{Y-F-Geok}\,(iii), every plaquette $P''$ that enters either lies in $D'_J$ or is tame. Moreover, $\sum_Pe^{-\theta_D{\rm dist}(P,P'')}\le3\mathsf g(\theta_D)$. Summing the bound of (3) over $P\in\mathfrak Z$ and inserting (4), we find \[\sum_{P\in\mathfrak Z}\sum_{P'}|\mathsf{Hm}(P,P')|\le C_{\rm hol}^\flat[\sum_{D'_J}\phi'+3M^3\#\mathfrak Z\,\varepsilon']+3M^3\#\mathfrak Z\,\mathsf{ti}_0^\flat.\] Summed in the same way, the bound of (3) on the column sums gives the same bound for $\sum_{P\in\mathfrak Z}\sum_{P'}|\mathsf{Hm}(P',P)|$. Hence, by (2), \[\frac12\mathrm{Hol}\le3e_1^\flat\varepsilon'\sum_{D'_J}\phi'+(9e_1^\flat+\frac92\mathsf N_Z^2c_B\mathsf{ti}_0^\flat)\varepsilon'^2M^3\#\mathfrak Z.\]

\emph{\textup{(6)} The main part, row by row.} Let $\mathsf{et}'=X'-\zeta'$; this field is parallel to $\zeta'$. At each plaquette $P$, \[\langle v_{*P},X'_P\rangle-\frac12\langle v_{*P},\zeta'_P\rangle=\frac12|\zeta'_P|^2+\frac12\langle\zeta'_P,(\ell^*m'\ell\zeta')_P\rangle+\varepsilon'|\mathsf{et}'_P|+\langle(\ell^*m'\ell\zeta')_P,\mathsf{et}'_P\rangle.\] For tame $P$ the right-hand side equals $\mathrm{En}'_P$. Now let $P\in D'$; for $P$ in $\mathfrak Z$ this means $P\in D'_J$. The mask $m'$ vanishes at all sites within $\rho_0$ of $P$, and so $|(\ell^*m'\ell\zeta')_P|\le \mathsf t_\rho\varepsilon'$. At such $P$ we also have $|\mathsf{et}'_P|=\phi'_P-\varepsilon'$, and \[\frac12\varepsilon'^2+\varepsilon'(\phi'-\varepsilon')=d_0\varepsilon'\phi'+(\frac12-d_0)\varepsilon'\phi'+\frac12\varepsilon'(\phi'-\varepsilon').\] Summing over $P\in\mathfrak Z$ and combining the result with (1) and (5) proves \Cref{Y-F-L23k}.\end{proof}

\emph{Step 6: the balance and its three rim brackets.}

\begin{lemma}[The balance, with three rim brackets]\label{Y-F-L25k} Let $(J,t)$ be an admitted cell datum, and set \[\mathrm{Bal}:=\frac x2\|\Delta\|_{\mathfrak Z}^2+x\langle\Xi_*,\mathrm{d}^VEX'\rangle_{\mathfrak S}-x\langle Z_*,d_1a_0\rangle_{\mathfrak S^-}-x\langle\Xi_*,d_2a_0\rangle_{\mathfrak S^-}-\sum_Jw_j.\] Then \[\mathrm{Bal}\le-(\frac12-d_0-\mathsf t_\rho-3e_1^\flat)x'\varepsilon'\sum_{D'_J}\phi'_P+\sigma_{\rm tot}^\flat\#\mathfrak Z_2.\]\end{lemma}

\begin{proof} \emph{Identities.} What follows holds at every $V'$ and for all $\mathsf{zf},\Xi,\sigma,F,n,\psi$. Put \[\mathsf{pair}(\mathsf{zf},\Xi;\sigma):=(W^{V*}\mathsf{zf}+\mathrm{d}^{V*}\Xi-\ell^{V*}\sigma,\ \sigma,\ a^{-1}\mathrm{d}'^{V*}\mathsf{zf}).\] Then \begin{equation}\tag{E1$_k$}\langle\mathsf{pair}(\mathsf{zf},\Xi;\sigma),\iota(F,n)\rangle_{\mathrm{Hil}}=\langle\mathsf{zf},W^VF+\mathrm{d}'^Vn\rangle+\langle\Xi,\mathrm{d}^VF\rangle,\end{equation} and at $\iota(F,n)=k\psi$ the right-hand side of \textup{(E1$_k$)} becomes $\langle\mathsf{zf},d_1\psi\rangle+\langle\Xi,d_2\psi\rangle$. Moreover, \begin{equation}\tag{E3$_k$}\langle\bar F,k\psi\rangle_{\mathrm{Hil}}=\langle EX',\mathsf B^V\mathrm{d}^V\psi\rangle-\langle a_0,q\psi\rangle=\langle u^T\,\mathsf g^VX',\psi\rangle.\end{equation} We have $\bar F=\iota(EX'-\mathrm{d}^Va_0,\ Q_1^Va_0)$, and we write $\bar n=Q_1^Va_0$ for its bond part. Because $W^VE=1$ (Step~1\,(iv) of the proof of \Cref{Y-D-C5}), it follows that $W^V\bar F_F+\mathrm{d}'^V\bar n=X'-d_1a_0$; also $\mathrm{d}^V\bar F_F=\mathrm{d}^VEX'-d_2a_0$. Finally, $\mathrm{d}^VE=E_3\mathrm{d}'^V$, by Step~1\,(iv) of the proof of \Cref{Y-D-C5}.

\emph{\textup{(1)} Cut potentials.} Put \[\Delta^\chi:=\iota(E\chi X',0)-k(\chi a_0)-\mathrm{pm}(\chi v_*),\] and $\Xi^\chi=(\mathsf X^<)^V\chi v_*$. Let $\nu^\chi$ be the plaquette component of $\mathrm{pm}(\chi v_*)$, the first of its three components. We bound the components of $\mathbf 1_{\mathfrak Z}\Delta-\Delta^\chi$, beginning with the plaquette rows, where three terms occur. The first term vanishes, as $E\chi X'=\mathbf 1_{\mathfrak Z}EX'$. The second, $\mathbf 1_{\mathfrak Z}\mathrm{d}^Va_0-\mathrm{d}^V\chi a_0$, is at most $4\sup|a_0|$ and is nonzero only at those plaquette rows $p$ that have a bond $b$ with $p\in\hat{\mathfrak Z}$ and $b\notin\hat{\mathfrak Z}$, or with $p\notin\hat{\mathfrak Z}$ and $b\in\hat{\mathfrak Z}$. The third, $\mathbf 1_{\mathfrak Z}\nu_*-\nu^\chi$, is the $\nu$-component of $\mathrm{pm}((1-\chi)v_*)$ inside $\mathfrak Z$ and that of $-\mathrm{pm}(\chi v_*)$ outside; it is at most $N_1L^{-2}\varepsilon'$ and is nonzero only at rows of $\hat{\mathfrak Z}$ that lie within $R(x')+2\varrho_L+3$ of the complement of $\hat{\mathfrak Z}$ and at rows off $\hat{\mathfrak Z}$ that lie within $R(x')+2\varrho_L+3$ of $\hat{\mathfrak Z}$. At the site rows there are again three terms. The first is $(\ell^VE(1-\chi)X')_y$ at $y\in\hat{\mathfrak Z}$ and $-(\ell^VE\chi X')_y$ off it, and is at most $l_r\varepsilon'$. The second, $\mathbf 1_{\mathfrak Z}(y)(\ell^V\mathrm{d}^Va_0)_y-(\ell^V\mathrm{d}^V\chi a_0)_y$, equals $(\ell^V\mathrm{d}^V(1-\chi)a_0)_y$ at $y\in\hat{\mathfrak Z}$ and $-(\ell^V\mathrm{d}^V\chi a_0)_y$ off it, and is at most $4l_r\sup|a_0|$: the absolute row sum of $\ell^V\mathrm{d}^V$ is at most $4l_r$, and a bond $b$ contributes only if it lies within $\varrho_L+2$ of $y$, with $b\notin\hat{\mathfrak Z}$ when $y\in\hat{\mathfrak Z}$ and $b\in\hat{\mathfrak Z}$ when $y\notin\hat{\mathfrak Z}$. The third, the $\sigma$-component, is at most $\mathsf N_\sigma L^{-2}\varepsilon'$. At the bond rows, which carry the weight $a$, the component is at most $L\sup|a_0|+4\mathsf N_Z(aL)^{-1}\varepsilon'$. All these rows satisfy the hypothesis of \Cref{Y-F-Geok}\,(iv); hence they are based in $\mathfrak Z_2\setminus\mathfrak b$, and since $f_J\subset\mathfrak b$ by item~E12 of \S\ref{Y-F-1}, they lie off $f_J$. By \Cref{Y-F-Geok}\,(ii) and \Cref{Y-F-LA}\,(b), at these rows $|\Delta_i|\le K_\Delta^\flat\varepsilon'^2$ and $|a_0|\le \mathsf a_{\rm c}\varepsilon'$. The rows that belong to one cell are those of $3v$ fine plaquettes, $v$ fine sites and $3M^3$ coarse bonds. Hence \[\frac x2\bigl|\,\|\Delta\|_{\mathfrak Z}^2-\langle\Delta,\Delta^\chi\rangle_{\mathrm{Hil}}\bigr|\le\frac x2K_\Delta^\flat\varepsilon'^3[3v(4\mathsf a_{\rm c}+N_1)+v(l_r(4\mathsf a_{\rm c}+1)+\mathsf N_\sigma)+3M^3(aL\mathsf a_{\rm c}+4\mathsf N_Z)]\#\mathfrak Z_2.\] To simplify the bracket, let $c=c_\eta^\flat\ge \mathsf a_{\rm c}$ and $\mathsf{Nn}=\mathsf N_Z+\mathsf N_\nu^\flat+\mathsf N_\sigma\ge1$, and note that $M^3L\le v/9$. Then $12c+3\mathsf{Nn}\le12(1+c)\mathsf{Nn}$, $4l_r(1+c)+\mathsf{Nn}\le4(1+l_r)(1+c)\mathsf{Nn}$ and $\frac13ac+\frac49\mathsf{Nn}\le\frac43(1+a)(1+c)\mathsf{Nn}$. The bracket is therefore at most $17.34\,v(1+a)(1+l_r)(1+c)\mathsf{Nn}$, and the part of the bound above that comes from the objects based in one cell is at most $8.67\,L^{3/2}vK_\Delta^\flat(1+a)(1+l_r)(1+c)\mathsf{Nn}\,p_0^3x^{-1/2}$.

\emph{\textup{(2)} Six pairings.} With the identities above, the six pairings are as follows. First, $\langle\bar F,\iota(E\chi X',0)\rangle=\langle \mathsf B^V\bar F_F,EX'\rangle_{\mathfrak Z}$. Second, $\langle\bar F,k\chi a_0\rangle=\langle u^T\,\mathsf g^VX',\chi a_0\rangle$. Third, \[\FitDisplay{\dimexpr\linewidth-1pt\relax}{\langle\bar F,\mathrm{pm}(\chi v_*)\rangle=\langle\chi Z_*,X'-d_1a_0\rangle+\langle\Xi^\chi,\mathrm{d}^VEX'-d_2a_0\rangle.}\] Fourth, $\langle\boldsymbol\nu,\iota(E\chi X',0)\rangle=\langle Y_*,EX'\rangle_{\mathfrak Z}$. Fifth, $\langle\boldsymbol\nu,k\chi a_0\rangle=\langle Z_*,d_1\chi a_0\rangle+\langle\Xi_*,d_2\chi a_0\rangle$. Sixth, $\langle\boldsymbol\nu,\mathrm{pm}(\chi v_*)\rangle=L^{-1}\langle v_*,\hat R^Vv_*\rangle_{\mathfrak Z}$. Combining the six pairings, we obtain \[\FitDisplay{\dimexpr\linewidth-1pt\relax}{\begin{aligned}
&\langle\Delta,\Delta^\chi\rangle=\langle\Delta_Y,EX'\rangle_{\mathfrak Z}-\langle u^T\,\mathsf g^VX',\chi a_0\rangle
\\
&\quad{}-\langle X',Z_*\rangle_{\mathfrak Z}-\langle \mathrm{d}^VEX',\Xi^\chi\rangle
\\
&\quad{}+\langle\chi Z_*,d_1a_0\rangle+\langle Z_*,d_1\chi a_0\rangle
\\
&\quad{}+\langle\Xi^\chi,d_2a_0\rangle+\langle\Xi_*,d_2\chi a_0\rangle
\\
&\quad{}+L^{-1}\langle v_*,\hat R^Vv_*\rangle_{\mathfrak Z}.
\end{aligned}}\]

\emph{\textup{(3)} The prefactor.} We have $\mathrm{mk}'^{+r(x')}_j\subset f_j\subset\mathfrak Z$. A cell of $\mathfrak Z$ off $\mathrm{mk}'^{+r(x')}_J$ is farther than $Mr(x')\ge R_6$ from $D'_J$, and by \Cref{Y-F-Geok}\,(i) it is also farther than $R_6$ from the other rough plaquettes. \Cref{Y-F-LA}\,(c) therefore gives \[|\sum_Jw_j-\sum_{\square\in\mathfrak Z}T_\square|\le K_{\rm q}^\flat p_0'^2\varepsilon'\#\mathfrak Z.\] By item~E7 of \S\ref{Y-F-1}, \[\sum_{\mathfrak Z}T_\square=\frac x2\langle X',Z_*\rangle_{\mathfrak Z}-x'\sum_{\mathfrak Z}\mathrm{En}'_P+\frac x2\langle EX',\Delta_Y\rangle_{\mathfrak Z}+\frac x2\langle \mathrm{d}^VEX',\Xi_*\rangle_{\mathfrak Z}-\frac x2\langle a_0,u^T\,\mathsf g^VX'\rangle_{\mathfrak Z}.\]

\emph{\textup{(4)} The sum; support and sizes of the rim.} We consider the expression \[\FitDisplay{\dimexpr\linewidth-1pt\relax}{\frac x2\langle\Delta,\Delta^\chi\rangle+x\langle\Xi_*,\mathrm{d}^VEX'\rangle_{\mathfrak S}-x\langle Z_*,d_1a_0\rangle_{\mathfrak S^-}-x\langle\Xi_*,d_2a_0\rangle_{\mathfrak S^-}-\sum_{\mathfrak Z}T_\square.}\] Once (2) and (3) are inserted, the terms with $\Delta_Y$ and those with $u^T$ cancel. The term $-\frac x2\langle X',Z_*\rangle_{\mathfrak Z}$ appears twice and gives $-x'\langle v_*,X'\rangle_{\mathfrak Z}$, and the terms $+\frac{x'}2\langle v_*,\hat R^Vv_*\rangle_{\mathfrak Z}+x'\sum\mathrm{En}'$ remain. So the expression equals $-x'\mathcal G_{\mathfrak Z}+x'\sum_{\mathfrak Z}\mathrm{En}'+\mathrm{rim}$, where \[\FitDisplay{\dimexpr\linewidth-1pt\relax}{\begin{aligned}
&\mathrm{rim}=x[\langle \mathrm{d}^VEX',\Xi_*\rangle_{\mathfrak S}-\frac12\langle \mathrm{d}^VEX',\Xi^\chi\rangle-\frac12\langle \mathrm{d}^VEX',\Xi_*\rangle_{\mathfrak Z}]
\\
&\quad{}+x[\frac12\langle\chi Z_*,d_1a_0\rangle+\frac12\langle Z_*,d_1\chi a_0\rangle-\langle Z_*,d_1a_0\rangle_{\mathfrak S^-}]
\\
&\quad{}+x[\frac12\langle\Xi^\chi,d_2a_0\rangle+\frac12\langle\Xi_*,d_2\chi a_0\rangle-\langle\Xi_*,d_2a_0\rangle_{\mathfrak S^-}].
\end{aligned}}\]

\emph{Support.} The first bracket is a sum over corner cubes $q$. At $q\in\mathfrak S$ farther than $R(x')$ from the complement of $\mathfrak Z$ we have $\Xi^\chi=\Xi_*$, and the coefficients give $1-\frac12-\frac12=0$; outside $\mathfrak Z$ enlarged by $R(x')$ the first bracket vanishes. In the second bracket, $(d_1\psi)_P$ depends only on the bonds based in $B(y_3(P))$, one block away; this is \Cref{Y-F-P17}, which we now state and prove.

\begin{proposition}[$d_1$ and $d_2$ see the curvature of a single coarse plaquette]\label{Y-F-P17} Let $V'$ be any $G$-valued field, let $\psi\in\mathfrak L_{\mathfrak g}$ and $P=(y;\mu\nu)$ with $\mu<\nu$, and let $y_1,y_2,y_3=y_3(P)$ and $c_1$ to $c_4$ be as in Step~3\,(3) of the proof of \Cref{Y-F-C7}. Put $T_a:=\operatorname{Ad}_{V'(c_1)V'(c_2)}=T(y,y_3)$ and $T_b:=\operatorname{Ad}_{V'(c_4)V'(c_3)}$. Then the following hold.

\begin{enumerate}\item[\textup{(1)}] There are numbers $\mathsf{be}_b$ with $\sum_b|\mathsf{be}_b|\le L-1$ such that $(d_1\psi)_P=\sum_{b\text{ based in }B(y_3)}\mathsf{be}_b\,(T_b-T_a)\psi_b$. Hence $|(d_1\psi)_P|\le2\phi'_P(L-1)|\psi|_\infty\le2\phi'_P\,2L|\psi|_\infty$.

\item[\textup{(2)}] Let $q$ be a cube and $z=\pi(q)$. Then $(d_2\psi)_q$ is a sum over those pairs $\mu<\nu$ for which $z_\mu$ and $z_\nu$ lie at the top of the block of $z$. The term of such a pair is $\pm(T_a-T_b)\psi_b$, where $b=(z+e_\mu+e_\nu,\lambda)$ with $\lambda$ the third direction, and $T_a$, $T_b$ are those of $P=(y(q);\mu\nu)$. There are $L$ such cubes to a $P$, and the share of $P$ in a sum $\sum_q|\Xi_q|\,|(d_2\psi)_q|$ is at most $|\Xi|_\infty2\phi'_PL|\psi|_\infty$.\end{enumerate}\end{proposition}

\begin{proof} We have $T_aT_b^{-1}=\operatorname{Ad}_{\partial V'(P)}$, so that $\|T_b-T_a\|\le2\phi'_P$ by \Cref{Y-F-L1}\,(b); with this, the bounds in (1) and (2) follow from the formulas there, which we now prove.

\emph{\textup{(1)}.} As scalar kernels, $W\mathrm{d}=\mathrm{d}'Q_1$ (\Cref{Y-A-Facts}). In $W^V\mathrm{d}^V$ the bond $b$ carries the transport $T(y,y(p))T(y(p),y(b))$, and in $\mathrm{d}'^VQ_1^V$ it carries $T(y,c_-)T(c_-,y(b))$. Each of these is $\operatorname{Ad}$ of $V'$ along a monotone path from $y$ to $y(b)$ whose sites are corners of $P$. To $y$, $y_1$ and $y_2$ there is just one such path; to $y_3$ there are two, and they give $T_a$ and $T_b$. If $T_a$ is put in place of $T_b$, the sum becomes $\sum_b(W\mathrm{d}-\mathrm{d}'Q_1)(P,b)\,T(y,y(b))\psi_b=0$. Hence only the terms carrying $T_b$ remain, with $T_b-T_a$ in place of $T_b$, and these terms are of two kinds. The terms of the first kind are those of the plaquettes $p=(z;\mu\nu)$ based in $B(y_2)$ with $z+e_\mu\in B(y_3)$; a plaquette based in $B(y)$ has no bond based in $B(y_3)$. For these $p$ we have $W(P,p)=L^{-3}L\,t(z_\nu-y_\nu)$, where $t(\cdot)$ denotes the tent function of the weights $W(P,p)$; its values run through $1,\dots,L-1$, at each of $L$ heights, and these weights add up to $(L-1)/2$. The terms of the second kind are those of the coarse bond $c_3$, for which $\sum_{b\text{ based in }B(y_3)}Q_1(c_3,b)=L^{-3}L^2\sum_{k=0}^{L-1}k=(L-1)/2$. The two kinds together give $\sum_b|\mathsf{be}_b|\le L-1$.

\emph{\textup{(2)}.} As scalars, $dd=0$: every bond $b$ of $q$ lies in two faces of $q$, with opposite signs. These faces are based at $z$ or at some $z+e_\rho$, and their bonds lie at most one step further. The two transports from $y(q)$ to $y(b)$ can therefore differ only if $b=(z+e_\mu+e_\nu,\lambda)$ and both steps cross a seam. In that case the two faces are $(z+e_\mu;\nu\lambda)$ and $(z+e_\nu;\mu\lambda)$, and the transports are $T_a$ and $T_b$.\end{proof}

We now return to the second bracket of $\mathrm{rim}$. For $P\in\mathfrak S^-$, the $\chi$ in $\chi Z_*$ is evaluated at $P$ and, by \Cref{Y-F-P17}, the $\chi$ in $d_1\chi a_0$ at the bonds based in $B(y_3(P))$, one block away. Since $\mathfrak S^-=\mathfrak b^{+r_1-r_l-2}$ and $\mathfrak Z=\mathfrak b^{+r_1}$, we have $P\in\mathfrak Z$, and these bonds lie in $\hat{\mathfrak Z}$. Hence both cuts of this bracket change nothing at $P$, and the coefficients give $\frac12+\frac12-1=0$; off $\mathfrak Z^{+1}$ the second bracket is $0$. The third bracket is treated as the first, with $\mathfrak S^-$ in place of $\mathfrak S$: at $q\in\mathfrak S^-$ farther than $R(x')$ from the complement of $\mathfrak Z$ we have $\Xi^\chi=\Xi_*$, and, since by \Cref{Y-F-P17} the $\chi$ in $d_2\chi a_0$ is evaluated at the bonds $(z+e_\mu+e_\nu,\lambda)$ with $z=\pi(q)$, which lie in $\hat{\mathfrak Z}$, this $\chi$ changes nothing at $q$, so that the coefficients give $\frac12+\frac12-1=0$; outside $\mathfrak Z$ enlarged by $R(x')$ the third bracket vanishes. What remains is based in $\mathfrak b^{+r_1+\lceil R/M\rceil+1}\setminus\mathfrak b^{+r_1-r_l-2-\lceil R/M\rceil}$, which by \Cref{Y-F-Geok}\,(ix) lies inside $\mathfrak Z_2\setminus\mathfrak b$ and off $f_J$; its boxes are tame by \Cref{Y-F-Geok}\,(ii).

\emph{Sizes.} There $|(\mathrm{d}^VEX')_q|=|(\mathrm{d}'^VX')_{q'}|\le18\varepsilon'^2$, with $M^3$ cubes $q'$ to a cell, and $|\Xi_*|,\ |\Xi^\chi|\le \mathsf N_\Xi L^{-1}\varepsilon'$, $|Z_*|\le \mathsf N_ZL^{-1}\varepsilon'$. \Cref{Y-F-P17} with $\phi'_P\le\varepsilon'$ gives $|(d_1\psi)_P|\le2\varepsilon'\cdot2L|\psi|_\infty$ and, for the sum over the $L$ cubes $q$ that go with one $P$, $\sum_q|\Xi_q|\,|(d_2\psi)_q|\le|\Xi|_\infty2\varepsilon'L|\psi|_\infty$. Here $|\psi|_\infty\le \mathsf a_{\rm c}\varepsilon'$, also at the bonds based in $B(y_3(P))$, one block further. This rests on the following two facts, which together are the two hypotheses of the second alternative of \Cref{Y-F-Geok}\,(ii). First, these bonds are based in $\mathfrak b^{+r_1+\lceil R/M\rceil+2}$, which is part of $\mathfrak Z_2$ because $r_1\ge\lceil R_2/M\rceil+8$. Second, they lie within $1\le r_w+\varrho_L+2$ of the site $y_P$, which is off $f_J$; so the second alternative of \Cref{Y-F-Geok}\,(ii) holds, giving a tame box of radius $R_6$, and \Cref{Y-F-LA}\,(b) yields the bound $\mathsf a_{\rm c}\varepsilon'$. There are $3M^3$ plaquettes $P$ to a cell, and in each bracket the moduli of the coefficients add up to at most $2$. The parts of the three brackets that come from the objects based in one cell are therefore at most $36\mathsf N_\Xi$, $24\mathsf N_Z\mathsf a_{\rm c}$ and $12\mathsf N_\Xi \mathsf a_{\rm c}$, each times $x\varepsilon'^3M^3$, and \[|\mathrm{rim}|\le36L^{3/2}v(\mathsf N_Z+\mathsf N_\Xi)(1+c_\eta^\flat)p_0^3x^{-1/2}\#\mathfrak Z_2.\]

\emph{\textup{(5)} Conclusion.} We have \[\FitDisplay{\dimexpr\linewidth-1pt\relax}{\mathrm{Bal}=\frac x2[\|\Delta\|_{\mathfrak Z}^2-\langle\Delta,\Delta^\chi\rangle_{\mathrm{Hil}}]+[-x'\mathcal G_{\mathfrak Z}+x'\sum_{\mathfrak Z}\mathrm{En}'+\mathrm{rim}]+[\sum_{\square\in\mathfrak Z}T_\square-\sum_Jw_j],}\] where the three summands are, in this order, the error of (1), the expression of (4) and the error of (3). We treat the third summand below. By (1) and (4), the parts of the first summand and of the term $\mathrm{rim}$ of the second summand that come from the objects based in one cell of $\mathfrak Z_2$ are together at most \[\FitDisplay{\dimexpr\linewidth-1pt\relax}{\begin{aligned}
&L^{3/2}v(1+c_\eta^\flat)\,[\,8.67\,K_\Delta^\flat(1+a)(1+l_r)\mathsf{Nn}+36(\mathsf N_Z+\mathsf N_\Xi)\,]\,p_0^3x^{-1/2}
\\
&\quad{}\le44.67\,L^{3/2}v(1+a)(1+K_\Delta^\flat)(1+l_r)(1+c_\eta^\flat)(\mathsf N_Z+\mathsf N_\Xi+\mathsf N_\nu^\flat+\mathsf N_\sigma)p_0^3x^{-1/2}
\\
&\quad{}\le K_s^\flat p_0^3x^{-1/2}.
\end{aligned}}\] Indeed, the product $K_\Delta^\flat(1+a)(1+l_r)(1+c_\eta^\flat)\mathsf{Nn}$ from (1) and the product $(\mathsf N_Z+\mathsf N_\Xi)(1+c_\eta^\flat)$ from (4) are both at most $(1+a)(1+K_\Delta^\flat)(\mathsf N_Z+\mathsf N_\Xi+\mathsf N_\nu^\flat+\mathsf N_\sigma)(1+l_r)(1+c_\eta^\flat)^2$, while $44.67L^{3/2}\le200L^3$. Besides the two errors of (1) and (4), there is a third, the error $K_{\rm q}^\flat p_0'^2\varepsilon'\#\mathfrak Z$ of (3), which yields the second term of $\sigma_{\rm tot}^\flat$ because $\#\mathfrak Z\le\#\mathfrak Z_2$. The third term, $(9e_1^\flat+\frac92\mathsf N_Z^2c_B\mathsf{ti}_0^\flat)M^3p_0'^2$, is what the last term of \Cref{Y-F-L23k} becomes after multiplication by $x'$, with $x'\varepsilon'^2=p_0'^2$ and again $\#\mathfrak Z\le\#\mathfrak Z_2$. Together with the bound $K_s^\flat p_0^3x^{-1/2}$ above, these are the three terms of $\sigma_{\rm tot}^\flat$; both $K_s^\flat$ and $\sigma_{\rm tot}^\flat$ are defined in item~F of \S\ref{Y-F-1}. The inequality of \Cref{Y-F-L23k}, multiplied by $x'$, gives the main term $-(\frac12-d_0-\mathsf t_\rho-3e_1^\flat)x'\varepsilon'\sum_{D'_J}\phi'_P$. This proves the claim.\end{proof}

\emph{Step 7: the Gaussian weight row by row, the rim, and a summation by parts.} Let $(J,t)$ be an admitted cell datum. We call the set $\mathfrak Z_2\setminus\mathfrak b$ the collar. By item~E10 of \S\ref{Y-F-1}, $b_\square(\xi)=0$ for a cell $\square$ exactly when $|\xi_b|\le c_\xi\varepsilon$ at every bond $b$ of $I^\circ$ in $\hat\square$. We claim that, on the event that $b_\square(\xi)=0$ for every cell $\square$ of the collar and that $|\eta_b|<\pi$ at each bond $b$ of $I^\circ\cap\hat{\mathfrak Z}$, \[\FitDisplay{\dimexpr\linewidth-1pt\relax}{\begin{aligned}
&-\sum_Jw_j-\frac x2\|k\varphi\|_{\mathfrak Z}^2+\frac x2\|\iota(F^\flat,n^\flat)\|_{\mathfrak Z}^2-x\mathcal A_{\mathfrak Z}
\\
&\quad{}\le-(\frac12-d_0-\mathsf t_\rho-3e_1^\flat)x'\varepsilon'\sum_{D'_J}\phi'_P
\\
&\quad{}-x\,\mathrm{Sq}_{\mathfrak Z}-x\,\mathrm{Pay}_{\mathfrak Z}
\\
&\quad{}+[\sigma_{\rm tot}^\flat+5K_s^\flat p_0^3x^{-1/2}+x^{-40}
\\
&\quad{}+(6.3(\mathsf N_\nu^\flat)^2\vartheta_F^\flat+3\mathsf N_Z^2(\theta_2+L\vartheta_n^\flat))M^3p_0'^2]\#\mathfrak Z_2.
\end{aligned}}\] No term linear in $\varphi$ is left on the right-hand side.

\emph{Substep 1: the exponent of the Gaussian weight cancelled row by row.} As $G$ is invertible and $C=GG^{\mathsf T}$, we have $|\xi|^2=\langle\varphi,C^{-1}\varphi\rangle=\|k\varphi\|_{\mathrm{Hil}}^2+\langle\varphi,S\varphi\rangle$, and \[E_\xi[f]=(x/2\pi)^{\frac32\#I^\circ}(\det C)^{-1/2}\int f\,e^{-\frac x2\langle\varphi,C^{-1}\varphi\rangle}d\varphi.\] Moreover, \[[A]^\ell_{\mathfrak Z}=\exp\{\frac x2\|\iota(F^\flat,n^\flat)\|_{\mathfrak Z}^2-x\mathcal A_{\mathfrak Z}\}\prod\mathbf 1\{|\eta_b|<\pi\}\mathsf j(\eta_b),\] with $\mathsf j\le1$. Since $\iota(F^\flat,n^\flat)=\bar F+k\varphi$ and $\iota(F,n)=\bar F+k\varphi-\mathfrak k$, we have, row by row on $\mathfrak Z$, \[-\frac12\|k\varphi\|^2+\frac12\|\bar F+k\varphi\|^2-\langle\boldsymbol\nu,\iota(F,n)\rangle+\frac12\|\boldsymbol\nu\|^2=\frac12\|\Delta\|^2+\langle\Delta,k\varphi\rangle+\langle\boldsymbol\nu,\mathfrak k\rangle.\] With \Cref{Y-F-LB}, applied to the starred fields and the sets of item~E12 of \S\ref{Y-F-1}, the left-hand side of the claim without $-\sum_Jw_j$ is at most \[\frac x2\|\Delta\|_{\mathfrak Z}^2+x\langle\Delta,k\varphi\rangle_{\mathfrak Z}+x\langle\boldsymbol\nu,\mathfrak k\rangle_{\mathfrak Z}-x|\langle Z_*,\mathrm{de}\rangle_{\mathfrak S}|-x|\langle\Xi_*,\mathrm{db}\rangle_{\mathfrak S}|-x\,\mathrm{Sq}-x\,\mathrm{Pay}+x\,\mathrm{vol}+x\,\mathrm{Rim}_{\mathfrak Z}.\]

\emph{Substep 2: the rim, and the one use of the collar tests.} By \Cref{Y-F-Geok}\,(vi), and as $G(b,b')=0$ beyond the distance $r_w$ by \textup{(C2)} of item~D of \S\ref{Y-F-1}, at a bond $b$ based in $\mathrm{Lr}$ the value $(G\xi)_b$ depends only on the values of $\xi$ at the bonds $b'$ of $I^\circ$ in the cells of $\mathfrak Z_2\setminus\mathfrak b$. At these bonds $|\xi_{b'}|\le c_\xi\varepsilon$, so that $|\varphi_b|\le \mathsf G_0c_\xi\varepsilon\le\varepsilon/6.5$. By \Cref{Y-F-Geok}\,(ii) and \textup{(C3)} of item~D of \S\ref{Y-F-1}, $|a_{0,b}|\le \mathsf a_{\rm c}\varepsilon'$. Consequently $|\eta_b|\le r_\eta:=c_\eta^\flat\varepsilon\le r_o$. We now apply the following lemma.

\begin{lemma}[The deviations of curvature and noise from their linear parts, in local form]\label{Y-F-L13} For the chart $U=\mathsf U(\eta;V')$ the following hold; each assertion involves only the $\eta_b$ that it names.

\begin{enumerate}\item[\textup{(i)}] Let $\sigma_p:=\sum_{b\in dp}|\eta_b|$. If $p$ is not an edge plaquette and $\sigma_p<\pi$, then $|F_p-(\mathrm{d}^V\eta)_p|\le\sigma_p^2$. If $p$ is an edge plaquette under $P$ and $\sigma_p+\phi'_P<\pi$, then $|F_p-X'_P-(\mathrm{d}^V\eta)_p|\le(\sigma_p+\phi'_P)^2$; and always $\phi'_P\le\phi_p+\sigma_p$.

\item[\textup{(ii)}] For each $r>0$ with $3Lr<\pi$, if $|\eta_b|\le r$ on the bonds of $W_c$, then $|n_c+(Q_1^V\eta)_c|\le14L^2r^2$.\end{enumerate}\end{lemma}

\begin{proof} \emph{\textup{(i)}.} By Step~4 of the proof of \Cref{Y-E-C6}, \[\partial U(p)=e^{i\eta_1}e^{iT_2\eta_2}[\partial V'(P)]e^{-iT_3\eta_3}e^{-i\eta_4},\] where the factor in brackets is present only for an edge plaquette under $P$. Here $\partial V'(P)=e^{iX'_P}$, and the $T_j$ are isometries. The exponents add up to $(\mathrm{d}^V\eta)_p$ if $p$ is not an edge plaquette, and to $X'_P+(\mathrm{d}^V\eta)_p$ if it is one. \Cref{Y-F-L1}\,(c) now gives the two bounds of (i), and \Cref{Y-F-L1}\,(a) gives its last assertion.

\emph{\textup{(ii)}.} Let $c=(y,y')$ have direction $\mu$, so that $y=c_-$, and let $x\in B(y)$ and $x'=x+Le_\mu$; here $h_{c,x}$, $\bar\tau_c$ and $\mathsf M(U)(c)$ are as in item~C.I~4 of \S\ref{Y-C-1}, and $\tau_{c,x}=\mathsf L(h_{c,x}h_{c,y}^{-1})$. We write $h_{c,x}V'(c)^{-1}=\prod_{b\in[x,x']}e^{i\tilde\eta_b}=:e^{i\mathsf{sg}_x}$, where $\tilde\eta=\operatorname{Ad}_{V'(c)}\eta$ in $B(c_+)$ and $\tilde\eta=\eta$ elsewhere. Then $|\mathsf{sg}_x|\le Lr$ and $|\mathsf{sg}_x-\sum\tilde\eta_b|\le L^2r^2$. Next, $\tau_{c,x}=\mathsf L(e^{i\mathsf{sg}_x}e^{-i\mathsf{sg}_y})$, whence $|\tau_{c,x}-\mathsf{sg}_x+\mathsf{sg}_y|\le4L^2r^2$ and $|\bar\tau_c|\le2Lr$. Further, $-n_c=\mathsf L(e^{i\bar\tau_c}e^{i\mathsf{sg}_y})$, whence $|n_c+\bar\tau_c+\mathsf{sg}_y|\le9L^2r^2$. The mean over $x$ of $\sum\tilde\eta_b$ equals $(Q_1^V\eta)_c$. Adding the three errors, with $1+4+9=14$, we obtain (ii).\end{proof}

The hypotheses of \Cref{Y-F-L13} are met, since $6Lr_o<\pi$ and $\phi'_P\le\varepsilon'$. Therefore $|k_p|\le(4r_\eta+\varepsilon')^2\le25(c_\eta^\flat)^2\varepsilon^2$, because $L^{1/2}\le c_\eta^\flat$, and $|k_n|\le14L^2(c_\eta^\flat)^2\varepsilon^2$. On these rows $\partial U(p)\ne-1$: the identity of Step~4 of the proof of \Cref{Y-E-C6} and \Cref{Y-F-L1}\,(a) give $\operatorname{ang}\partial U(p)\le\sigma_p+\phi'_P\le4r_o+\varepsilon'<\pi$, since $6Lr_o<\pi$ by \textup{(U0$^\flat$)} and $\varepsilon'\le1/40$ by \textup{(FA)}. So $\mathsf L$ is the genuine logarithm there, and $\phi_p=|F_p|$. Here $F_p=\bar F_p+(\mathrm{d}^V\varphi)_p-k_p$, and, as \Cref{Y-F-Geok}\,(ii) gives a tame box, \Cref{Y-F-LA}\,(b) yields $|\bar F_p|\le(\mathsf N_\sharp+2)L^{-2}\varepsilon'$. With $\varepsilon'\le L^{1/2}\varepsilon$, \textup{(O4)\,(iii)} and \textup{(Z0$^\flat$)}, it follows that \[\phi_p\le|\bar F_p|+4\varepsilon/6.5+|k_p|\le[(\mathsf N_\sharp+2)L^{-3/2}+0.6154+0.3]\varepsilon\le0.9954\,\varepsilon.\] As a result, no rough plaquette lies on the rows named in \Cref{Y-F-Geok}\,(viii), and on these rows $\zeta=F$ and $m_y=1$.

\emph{Substep 3: summation by parts, and five rim terms.} (i) We have \[\langle\boldsymbol\nu,\mathfrak k\rangle_{\mathfrak Z}=\sum_{\hat{\mathfrak Z}}\langle\nu_{*p},k_p\rangle+\sum_p\langle(\ell^{V*}\mathbf 1_{\mathfrak Z}\sigma_*)_p,k_p\rangle+a\langle\omega_*,k_n\rangle_{\mathfrak Z}.\] For $p\in\hat{\mathfrak Z}'$, every $y$ with $\ell_y(p)\ne0$ lies in $\hat{\mathfrak Z}$, and so $\nu_{*p}+(\ell^{V*}\mathbf 1_{\mathfrak Z}\sigma_*)_p=Y_{*p}$. It follows that $\langle\boldsymbol\nu,\mathfrak k\rangle_{\mathfrak Z}=\langle Y_*,k_F\rangle_{\mathfrak Z'}+a\langle\omega_*,k_n\rangle_{\mathfrak Z'}+\mathrm{rim}_0$, where $\mathrm{rim}_0$ lives on the plaquettes $p\in\widehat{\mathfrak Z^{+r_l}}\setminus\hat{\mathfrak Z}'$ and on the bonds $c\in\mathfrak Z\setminus\mathfrak Z'$. For the part of $\mathrm{rim}_0$ that comes from the objects based in one cell, \[x|\mathrm{rim}_0|\le[75v(N_1+l_c\mathsf N_\sigma)L^{-2}+168M^3L\mathsf N_Z]L^{1/2}(c_\eta^\flat)^2p_0^3x^{-1/2},\] the numerals arising as $3v\cdot25$ and $3M^3\cdot a\cdot4\mathsf N_Z(aL)^{-1}\cdot14L^2$. Using $Y_*=W^{V*}Z_*+\mathrm{d}^{V*}\Xi_*$ and $a\omega_*=\mathrm{d}'^{V*}Z_*$, we write the main part as \[\langle Z_*,W^V\mathbf 1_{\mathfrak Z'}k_F+\mathrm{d}'^V\mathbf 1_{\mathfrak Z'}k_n\rangle+\langle\Xi_*,\mathrm{d}^V\mathbf 1_{\mathfrak Z'}k_F\rangle.\] At $P,q\in\mathfrak S$ the cut by $\mathbf 1_{\mathfrak Z'}$ changes nothing. Indeed, $(W^V\mathbf 1_{\mathfrak Z'}k_F+\mathrm{d}'^V\mathbf 1_{\mathfrak Z'}k_n)_P$ involves the values of $k_F$ at the $p$ with $W(P,p)\ne0$ and those of $k_n$ at the four bonds of $P$, all of them within one block of $y_P$, and $(\mathrm{d}^V\mathbf 1_{\mathfrak Z'}k_F)_q$ involves the values of $k_F$ at the six faces of $q$. As $\mathfrak S=\mathfrak b^{+r_1-r_l-1}$ and $\mathfrak Z'=\mathfrak S^{+1}$, these plaquettes, bonds and faces all lie in $\mathfrak Z'$. Since $W^VF^\flat+\mathrm{d}'^Vn^\flat=X'+d_1\eta$ and $\mathrm{d}^VF^\flat=\mathrm{d}^VEX'+d_2\eta$, we obtain there $(W^Vk_F+\mathrm{d}'^Vk_n)_P=\mathrm{de}_P+(d_1\eta)_P$ and $(\mathrm{d}^Vk_F)_q=(\mathrm{d}^VEX')_q+(d_2\eta)_q-\mathrm{db}_q$. The remaining $P$ and $q$ lie in $\mathfrak Z'^{+1}\setminus\mathfrak S$, and the term $\mathrm{bd}_1$ is $x$ times their contribution to the main part of $\langle\boldsymbol\nu,\mathfrak k\rangle_{\mathfrak Z}$. At such a $P$ the contribution is at most \[\mathsf N_ZL^{-1}\varepsilon'(L^{-1}12L^3\cdot25+4\cdot14L^2)(c_\eta^\flat)^2\varepsilon^2=\mathsf N_ZL^{-1}\varepsilon'\,356L^2(c_\eta^\flat)^2\varepsilon^2,\] with $3M^3$ plaquettes $P$ to a cell; at such a $q$ it is at most $\mathsf N_\Xi L^{-1}\varepsilon'\cdot6\cdot25(c_\eta^\flat)^2\varepsilon^2$, with $v$ cubes to a cell. As $1068/L^2\le118.7$, this yields, for the part of $\mathrm{bd}_1$ that comes from the objects based in one cell, \[|\mathrm{bd}_1|\le x\varepsilon'\varepsilon^2(c_\eta^\flat)^2[1068\mathsf N_ZLM^3+150\mathsf N_\Xi v/L]\le150L^{1/2}v(\mathsf N_Z+\mathsf N_\Xi)(1+c_\eta^\flat)^2p_0^3x^{-1/2}.\]

(ii) The field $k(\chi_S\varphi)$ lives on rows within $\varrho_L+2$ of $\hat{\mathfrak S}$, all of which belong to $\mathfrak Z$; hence $x\langle\Delta,k\varphi\rangle_{\mathfrak Z}=x\langle\Delta,k\chi_S\varphi\rangle_{\mathrm{Hil}}+\mathrm{bd}_2$. Moreover, $\langle\Delta,k\psi\rangle_{\mathrm{Hil}}=\langle u^T\,\mathsf g^VX',\psi\rangle-\langle Z_*,d_1\psi\rangle-\langle\Xi_*,d_2\psi\rangle$. Here $\mathrm{bd}_2=x\langle\Delta,k(1-\chi_S)\varphi\rangle_{\mathfrak Z}$; the absolute row sums of $k$ are at most $4$, $4l_r$, and $L$ with the weight $a$, while $|\Delta_i|\le K_\Delta^\flat\varepsilon'^2$ and $|\varphi_b|\le\varepsilon/6.5$. Thus, for the part of $\mathrm{bd}_2$ that comes from the objects based in one cell, \[|\mathrm{bd}_2|\le K_\Delta^\flat L(12v+4l_rv+3aLM^3)/6.5\;p_0^3x^{-1/2}.\] On $\hat{\mathfrak Z}$ we have $|\varphi_b|<\pi(1+a_\infty^\flat)$, and therefore \[x|\langle u^T\,\mathsf g^VX',\chi_S\varphi\rangle|\le x\,x^{-50}\pi(1+a_\infty^\flat)3v\#\mathfrak Z\le x^{-40}\#\mathfrak Z.\]

(iii) For $P,q\in\mathfrak S^-$, the bonds on which $d_1$ and $d_2$ depend lie in $\hat{\mathfrak S}$, and $\eta-\varphi=-a_0$. Hence \[x\langle Z_*,d_1\eta\rangle_{\mathfrak S}-x\langle Z_*,d_1\chi_S\varphi\rangle+x\langle\Xi_*,d_2\eta\rangle_{\mathfrak S}-x\langle\Xi_*,d_2\chi_S\varphi\rangle=-x\langle Z_*,d_1a_0\rangle_{\mathfrak S^-}-x\langle\Xi_*,d_2a_0\rangle_{\mathfrak S^-}+\mathrm{bd}_3.\] The term $\mathrm{bd}_3$ lives at $P,q\in\mathfrak S^{+1}\setminus\mathfrak S^-$; at a $P$ it is at most $\mathsf N_ZL^{-1}\varepsilon'\,2\varepsilon'\,2L\cdot2c_\eta^\flat\varepsilon+\mathsf N_\Xi L^{-1}\varepsilon'\,2\varepsilon'L\cdot2c_\eta^\flat\varepsilon$. Therefore, as there are $3M^3$ plaquettes $P$ to a cell, we obtain, for the part of $\mathrm{bd}_3$ that the objects based in one cell contribute, \[|\mathrm{bd}_3|\le x\varepsilon'^2\,2c_\eta^\flat\varepsilon(4\mathsf N_Z+2\mathsf N_\Xi)3M^3\le24Lv(\mathsf N_Z+\mathsf N_\Xi)c_\eta^\flat p_0^3x^{-1/2}.\]

(iv) The two moduli of Substep 1 absorb the terms \[\FitDisplay{\dimexpr\linewidth-1pt\relax}{x\langle Z_*,\allowbreak \mathrm{de}\rangle_{\mathfrak S}-x\langle\Xi_*,\allowbreak \mathrm{db}\rangle_{\mathfrak S}}\]. The term $x\,\mathrm{Rim}_{\mathfrak Z}$ lives in the rim, where $\mathrm{Rim}_{\mathfrak Z}=\sum|\sigma_{*y}|\,\allowbreak |\mathsf{dl}_y|$ and \[|\mathsf{dl}_y|\le2\varepsilon r_\eta\sum_p|\ell_y(p)|\,|\Gamma(y,\pi(p))|\le2\varepsilon r_\eta\,9K\mathsf g_1;\] hence, for the part that comes from the objects based in one cell, \[\FitDisplay{\dimexpr\linewidth-1pt\relax}{x\,\allowbreak \mathrm{Rim}_{\mathfrak Z}\le v\mathsf N_\sigma L^{-3/2}18K\mathsf g_1c_\eta^\flat p_0^3x^{-1/2}}\]. For each of the five terms $x\,\mathrm{rim}_0$, $\mathrm{bd}_1$, $\mathrm{bd}_2$, $\mathrm{bd}_3$ and $x\,\mathrm{Rim}_{\mathfrak Z}$, the share of one cell is at most $K_s^\flat p_0^3x^{-1/2}$, as a comparison factor by factor shows, with $243\le200L^3$. What is left is $\mathrm{Bal}$.

\emph{Substep 4: the balance and the volume term.} We bound $\mathrm{Bal}$ by \Cref{Y-F-L25k}. Moreover, $x\,\mathrm{vol}_{\mathfrak Z}=x'L\,\mathrm{vol}_{\mathfrak Z}$ and $6N_1^2\le6.3(\mathsf N_\nu^\flat)^2$. This proves the claim of Step 7.

\emph{Step 8: a bound for the contribution of one loud set $\mathfrak N$ and one map $\beta$.}

\begin{proposition}[A bound for one loud set and one choice of tested bonds]\label{Y-F-PB}
Let $(J,t)$ be an admitted cell datum, let $\mathsf Y\subset\mathfrak Z_2$ and $0\le\varsigma_1\le x^{-2}$, let $\mathsf{th}$ be a Borel function of $\xi$ with $0\le\mathsf{th}(\xi)\le1$, let $\mathfrak N$ be a set of blocks, and let $\beta$ be a map that assigns to each $\square\in t$ a bond of $I^\circ$ in $\hat\square$. In the mean $E_\xi[\hat\lambda_{(J,t)}\,\mathsf{th}\,e^{\frac x2\varsigma_1|\xi_{\mathsf Y}|^2}]$, bound the product of the tests by $\prod_tb_\square\le\sum_\beta\prod1\{|\xi_{\beta(\square)}|>c_\xi\varepsilon\}$. Then the contribution of the event $\{\text{loud set}=\mathfrak N\}$ and of the map $\beta$ to this mean is \[\le\mathrm{Res}\;e^{\sigma'\#\mathfrak Z_2}\,e^{-(\mathrm{py}/192)\#\mathfrak N}\,(e^{c_\tau-\rho_Q})^{\#\beta_Q},\] where $\sigma'$ is the left side of \textup{(Z1$^\flat$)} without the term $27M^3e^{-\mathrm{py}/384}$. Moreover, $\#\{\text{type L}\}\le(2r_1+1)^3\#\mathfrak N$.
\end{proposition}

The loud set is introduced in Substep~5 of the proof, and the cells of type L and of type Q, together with the set $\beta_Q$, in Substep~6.

\begin{proof}
Since $\mathsf{th}\le1$ and the integrand is nonnegative, we may drop $\mathsf{th}$.

\emph{Substep 5: quiet and loud blocks; the splitting; the payments.} A coarse plaquette $P$ is quiet if $W_P\cap D\cap\hat{\mathfrak Z}$ is empty and if $\mathsf n_c\le r_o$ holds at each of its bonds $c$ that lie in $\mathfrak Z$. A block is quiet if all twelve of its coarse plaquettes $P$ are quiet, and loud otherwise. The set of the loud blocks, the loud set, lies inside $\mathfrak Z^{+1}$; on the event of the proposition it is $\mathfrak N$.

\emph{The splitting over $\mathfrak N$ and $\beta$.} We fix the pair $(V',e)$. Then $\eta=G\xi-a_0$ is affine in $\xi$, the chart $\mathsf U(\eta;V')$ is continuous in $\eta$, and $\phi_p=\operatorname{ang}\partial U(p)$ is continuous in $U$. By item~3.4 of \S\ref{Y-E-1}, $\mathsf M(U)(c)$ is a Borel map of $U$; hence $\mathsf n_c$ is Borel in $\xi$. Consequently, for each set $\mathfrak N$ of blocks of $\mathfrak Z^{+1}$ the set $\{\xi:\text{loud set}=\mathfrak N\}$ is Borel; there are finitely many such sets, they are pairwise disjoint, and together they cover the whole space.

The integrand is nonnegative and Borel. Indeed, at the fixed
pair $(V',e)$ the vectors $F^\flat$, $n^\flat$ and $\eta$ are affine
in $\xi$. The chart and all path products are continuous, the
logarithm with value $0$ at $-1$ is Borel, and the masks, cutoff
factors and tests are indicators of finite combinations of open
or closed sets. The block map $\mathsf M$ is Borel by
\S\ref{Y-E-1}, item~3.4, so its noise angle is Borel as well.
The members and $\mathsf{th}$ are Borel by hypothesis.
Every sum and product defining $\hat\lambda$ is finite.
Thus splitting the integral over the finitely many loud sets
and bounding the tests by the sum over $\beta$ gives the
contributions in the proposition.

Let $\mathsf{Pq}$ be the set of the plaquettes of $\hat{\mathfrak Z}$ whose corners all lie in quiet blocks, and $\mathsf{Cq}$ the set of the bonds of $\mathfrak Z$ whose two ends are quiet. A site row $y$ of $\hat{\mathfrak Z}$ is quiet if every block within $\varrho_L+2$ of $y(y)$ is quiet; $\mathsf{Yq}$ is the set of the quiet site rows, and $\mathsf{qu}$ is the union of $\mathsf{Pq}$, $\mathsf{Cq}$ and $\mathsf{Yq}$. We let $\mathrm{Bd}(\mathfrak N)$ be the set of the bonds of $I^\circ$ that lie inside a loud block, or that cross between two blocks at least one of which is loud; a loud block has $\le3L^3+6L^2\le5L^3$ of them. The rows are normalised by $k_c\to a^{1/2}k_c$. We set $\mathsf{sh}_p:=-\Delta_p$, $\mathsf{sh}_y:=-\Delta_y$, $\mathsf{sh}_c:=a^{1/2}(-\Delta_c+\theta_2(1-\theta_2)^{-1}\omega_{*c})$, and $\mathsf{dv}_i:=k_i\varphi-\mathsf{sh}_i$ for every row $i$.

\begin{enumerate}
\item[\textup{(1)}] On the bonds based in $\mathfrak Z^{+r_l}$ off $\mathrm{Bd}(\mathfrak N)$ we have $|\eta_b|\le2r_o$. For a bond based in $\mathfrak S$, the block of the bond when the bond lies inside a block, or $W_c$ when the bond crosses along $c$, lies in $\hat{\mathfrak Z}$ and is tame, and \textup{(Tc)}\,(o), (i) in Step~3 of the proof of \Cref{Y-F-C7} bounds $\operatorname{ang}(e^{i\eta_b})$ at such a bond $b$. If $b$ lies inside its block $B$, then \textup{(o)} gives $\operatorname{ang}(e^{i\eta_b})\le s_B\le3L^3\varepsilon$; if $b$ is a crossing bond of $c$, then \textup{(i)} gives $\operatorname{ang}(e^{i\eta_b})\le9s_c+\mathsf n_c$ with $s_c\le6L^3\varepsilon$ and $\mathsf n_c\le r_o$, the latter because $c$ lies in $\mathfrak Z$ and is a bond of a quiet $P$. In both cases $\operatorname{ang}(e^{i\eta_b})\le54L^3\varepsilon+r_o$. To pass from the angle to the modulus, note that $b\in I^\circ\cap\hat{\mathfrak Z}$; hence $|\eta_b|<\pi$ by the indicator $1\{|\eta_b|<\pi\}$ in $[A]^\ell_{\mathfrak Z}$ (item~E11 of \S\ref{Y-F-1}), so that $\mathsf L(e^{i\eta_b})=\eta_b$ by the definition of $\mathsf L$ in item~C.I~4 of \S\ref{Y-C-1}, and $|\eta_b|=\operatorname{ang}(e^{i\eta_b})$ by \Cref{Y-F-L1}. Finally, $54L^3\varepsilon+r_o\le2r_o$, since $60L^3\varepsilon\le r_o$ by \textup{(U0$^\flat$)}. For a bond based in $\mathfrak Z^{+r_l}\setminus\mathfrak S$, which lies inside $\mathrm{Lr}$, the bound is that of Substep~2 of Step~7 of the proof of \Cref{Y-F-C7}.

Consequently every $p\in\mathsf{Pq}$ has $p\notin D$ and $|k_p|\le(17r_o)^2$, the latter because $\phi'_P\le\phi_p+\sigma_p\le9r_o$, where $\sigma_p=\sum_{b\in dp}|\eta_b|\le8r_o$ is the number that enters \Cref{Y-F-L13}\,(i). Every $c\in\mathsf{Cq}$ has $c\notin\mathsf{Cb}$, $\mathsf n_c\le r_o\le t_0$ and $|k_n|\le56L^2r_o^2$. Every $y\in\mathsf{Yq}$ has $m_y=1$ and $\zeta=F$ on its reach; moreover $\ell^U_y\zeta-\sigma_{*y}=(k\varphi)_y+\Delta_y-(\ell^Vk_F)_y+\mathsf{dl}_y$, where the sum of the last two terms of the right side, the rest, satisfies $|\text{rest}|\le289l_rr_o^2+36K\mathsf g_1\varepsilon r_o$.

\emph{The reach of a quiet site row, including the part outside $\hat{\mathfrak Z}$.} Take $y\in\mathsf{Yq}$ and a plaquette $p$ with $|y-p|\le\varrho(x)+1$. The corners of $p$ and the sites of $\Gamma(y,\pi(p))$ lie within $\varrho(x)+2$ of $y$, and so, as $L\ge3$, in blocks within $\varrho_L+1$ of $y(y)$. These blocks are quiet, and they lie in $\mathfrak Z^{+r_l}$, since $\varrho_L+2\le r_lM$. Therefore the bonds of $p$ and of $\Gamma(y,\pi(p))$ are based in $\mathfrak Z^{+r_l}$ and lie off $\mathrm{Bd}(\mathfrak N)$, and the first assertion of (1) gives $|\eta_b|\le2r_o$ on them.

If $y\in\hat{\mathfrak Z}'$, then $p\in\hat{\mathfrak Z}$ by item~(1) of Step~3 of the proof of \Cref{Y-F-C7}, and $p=(z;\mu\nu)\in W_P$ for $P=(y(z);\mu\nu)$, which is one of the twelve $P$ of the quiet block $B(y(z))$; hence $p\notin D$. If $y\in\hat{\mathfrak Z}\setminus\hat{\mathfrak Z}'$, then $y$ is a row of $\mathrm{Rim}_{\mathfrak Z}$; on the support of $\hat\lambda$ the tests of $\mathfrak Z_2\setminus\mathfrak b$ pass, and by \Cref{Y-F-Geok}\,(viii) and the last assertion of Substep~2 of Step~7 of the proof of \Cref{Y-F-C7} no rough plaquette lies within $\varrho(x)+1$ of $y$. In either case $\zeta_p=F_p$; and $m_y=1$, because $\rho_0\le\varrho(x)$.

Further, $\sigma_p\le8r_o$ and $\phi'_P\le\phi_p+\sigma_p\le9r_o$. \Cref{Y-F-L13}\,(i), which applies because $17r_o<\pi$ by \textup{(U0$^\flat$)}, then gives $|k_p|\le(17r_o)^2$, also at the plaquettes $p$ based off $\hat{\mathfrak Z}$; hence $|(\ell^Vk_F)_y|\le289l_rr_o^2$. The estimate of item~(iv) of Substep~3 of Step~7 of the proof of \Cref{Y-F-C7}, here with $|\eta_b|\le2r_o$ on the bonds of the paths, gives $|\mathsf{dl}_y|\le2\varepsilon\,2r_o\sum_p|\ell_y(p)|\,|\Gamma(y,\pi(p))|\le36K\mathsf g_1\varepsilon r_o$. When the rows are empty, $\mathsf{Yq}$ is empty.

\item[\textup{(2)}] We have $F_p-\nu_{*p}=\mathsf{dv}_p-k_p$ and $n_c-\tilde\omega_c=a^{-1/2}\mathsf{dv}_c-k_n$. For vectors $u$, $r$ one has $|u+r|^2\ge(1-\check\theta_x)|u|^2-\check\theta_x^{-1}|r|^2$; this and \textup{(U0$^\flat$)} give \[x\,\mathrm{Sq}_{\mathfrak Z}\ge\frac x2(1-2\check\theta_x)\sum_{\mathsf{qu}}|\mathsf{dv}_i|^2+[\text{the squares of the bonds with }r_o<\mathsf n_c\le t_0]-r_\theta,\] with $r_\theta\le Q_\theta^\flat x^{-1/4}\#\mathfrak Z$. The numerals behind this bound are these: for one block, the plaquette rows give $\frac x2x^{1/4}17^4x^{-3/2}3L^3\le1.3\cdot10^5L^3x^{-1/4}$, the bond rows $4704aL^4x^{-1/4}$, and the site rows $L^3(83521l_r^2+1296K^2\mathsf g_1^2p_0^2x^{-1/4})x^{-1/4}$. A cell has $M^3$ blocks, and $p_0^2x^{-1/4}\le1$ by \textup{(U0$^\flat$)}; hence the sum for one cell is at most $Q_\theta^\flat x^{-1/4}$, where $Q_\theta^\flat$ is the constant of item~F of \S\ref{Y-F-1}.

\item[\textup{(3)}] Every rough plaquette of $\hat{\mathfrak Z}$ pays $\ge d_0p_0^2/2$ out of $x\,\mathrm{Pay}$. A bond of $\mathfrak Z$ with $\mathsf n_c>r_o$ has loud ends, does not belong to $\mathsf{Cq}$, and keeps $\ge0.05ax\mathsf n_c^2\ge0.05ax^{1/4}$: if $\mathsf n_c>t_0$, this is taken from $\mathrm{Pay}$, and otherwise $\mathsf{kc}_c|n_c-\tilde\omega_c|^2\ge\frac a4(0.9\mathsf n_c)^2$. One object makes at most $48$ blocks loud ($12$ boxes with $4$ corners each); so the payments are $\ge(\mathrm{py}/48)\#\mathfrak N$.

\item[\textup{(4)}] For $b\in\mathrm{Bd}(\mathfrak N)$ we have $|\eta_b|<\pi$, by the indicator or by Substep~2 of Step~7 of the proof of \Cref{Y-F-C7}, because a bond of $\mathrm{Bd}(\mathfrak N)$ that is not based in $\hat{\mathfrak Z}$ is based in the rim $\mathrm{Lr}$. Indeed, the distance in blocks from any loud block to $\mathfrak Z$ is $\le1$; hence the base of a crossing bond with an end in a loud block has distance in blocks $\le2\le M-1$ to $\mathfrak Z$, and so it lies in $\mathfrak Z^{+1}$; and $\mathfrak Z^{+1}\setminus\mathfrak Z=\mathfrak b^{+r_1+1}\setminus\mathfrak b^{+r_1}\subset\mathrm{Lr}$ because $r_l\ge1$. Hence $1\le e^{\frac x2\kappa_1\pi^2}e^{-\frac x2\kappa_1|\varphi_b-a_{0,b}|^2}$ on these bonds, at the cost $\frac x2\kappa_1\pi^2\,5L^3\#\mathfrak N=(\mathrm{py}/96)\#\mathfrak N$.
\end{enumerate}

\emph{Substep 6: the tested bonds; type L and type Q; the tilt.} As $\xi=G^{-1}\varphi=G^{\mathsf T}(q+S)\varphi$, we have $\xi_b=\sum_ih(b,i)k_i\varphi+(G^{\mathsf T}S\varphi)_b$ with $h(b,i)=(G^{\mathsf T}k_i^*)(b)$, and $h(b,i)$ is $0$ beyond $r_w+\varrho_L+2$. The row sums of $h$ are $\le \mathsf G_0(4+4l_c+a^{1/2}L^{-2})$ and its column sums are $\le \mathsf G_0\max\{4,4l_r,a^{1/2}L\}$; both bounds are $\le h_*$.

A cell $\square\in t$ is of type L if a loud block lies within $r_w+2\varrho_L+5$ ($\le r_1M$) of $\beta(\square)$; this gives $\#\{\text{type L}\}\le(2r_1+1)^3\#\mathfrak N$, the last assertion of the proposition. All other cells of $t$ are of type Q. We write $\beta_Q$ for the bonds of type Q and $\Pi_\beta$ for the projection on them. For a bond $b$ of type Q, every row $i$ with $h\ne0$ is a row of $\mathfrak Z$; such a row is quiet and is based in $\mathfrak Z_2$, within $r_w+\varrho_L+2$ of a site off $f_J$. By \Cref{Y-F-Geok}\,(ii), Lems.~\ref{Y-F-LA} and~\ref{Y-F-LM}, and \textup{(Z0$^\flat$)}, $\sum_i|h(b,i)|\,|\mathsf{sh}_i|\le c_\xi\varepsilon/4$.

Therefore $|\xi_b|>c_\xi\varepsilon$ gives $2|\sum h\,\mathsf{dv}_i|^2+2|(G^{\mathsf T}S\varphi)_b|^2>\frac9{16}c_\xi^2\varepsilon^2$, and $|\sum h\,\mathsf{dv}_i|^2\le h_*\sum|h|\,|\mathsf{dv}_i|^2$. Chebyshev's inequality with the parameter $\frac x2t_Q$ yields \[\prod_{\beta_Q}1\{\dots\}\le\exp\{\frac x2\sum_{\mathsf{qu}}\tau_i|\mathsf{dv}_i|^2+\frac x22t_Q\langle\varphi,SG\Pi_\beta G^{\mathsf T}S\varphi\rangle-\rho_Q\#\beta_Q\},\] where $\tau_i:=2t_Qh_*\sum_{b\in\beta_Q}|h(b,i)|\le\frac16$ for $i\in\mathsf{qu}$, while $\tau_i$ is $0$ for all other rows; moreover $\sum\tau_i\le\#\beta_Q/6$ and $\frac9{32}t_Qc_\xi^2p_0^2=\rho_Q$.

\emph{The tilt.} In the collar, $|\xi_b|\le c_\xi\varepsilon$, which gives one factor $e^{\frac32v\varsigma_1c_\xi^2p_0^2}$ for every cell of $\mathsf Y$ that lies in the collar. For $b\in\hat{\mathfrak b}$, $(G^{\mathsf T}q\varphi)_b$ depends on $\varphi$ within $r_w+r_q\le r_1M$, hence in $\hat{\mathfrak Z}$, where $|\varphi|<\pi(1+a_\infty^\flat)$; so $|(G^{\mathsf T}q\varphi)_b|\le c_q^\flat$ and $|\xi_b|^2\le2(c_q^\flat)^2+2|(G^{\mathsf T}S\varphi)_b|^2$, whence one factor $e^{3vx\varsigma_1(c_q^\flat)^2}$ for each cell of $\mathsf Y$ in $\mathfrak b$. We put $S':=2t_QSG\Pi_\beta G^{\mathsf T}S+2\varsigma_1SG\mathbf 1_{\hat{\mathfrak b}}G^{\mathsf T}S\ge0$; then $\|S'\|\le(2t_Q+2\varsigma_1)(2/\gamma)x^{-120}\le x^{-110}$.

\emph{Substep 7: the Gaussian integral; three quotients of determinants.} Let \[\FitDisplay{\dimexpr\linewidth-1pt\relax}{\sigma'':=\sigma_{\rm tot}^\flat+5K_s^\flat p_0^3x^{-1/2}+x^{-40}+(6.3(\mathsf N_\nu^\flat)^2\vartheta_F^\flat+3\mathsf N_Z^2(\theta_2+L\vartheta_n^\flat))M^3p_0'^2+\,Q_\theta^\flat x^{-1/4}+\,\frac32vc_\xi^2p_0^2x^{-2}+3v(c_q^\flat)^2x^{-1}.}\] The terms up to $M^3p_0'^2$ form the bracket of Step~7 of the proof of \Cref{Y-F-C7}. The term $Q_\theta^\flat x^{-1/4}$ takes care of $r_\theta$ of item~(2) of Substep~5, as $\#\mathfrak Z\le\#\mathfrak Z_2$. The last two terms are the tilt of Substep~6 for $\varsigma_1\le x^{-2}$: a cell of $\mathsf Y$ lies in the collar or in $\mathfrak b$, and the sum covers both cases. The number $\sigma''$ depends only on $x$ and on the letters \textup{(O0)} to \textup{(O10)}.

The exponent $-(\mathrm{py}/96)\#\mathfrak N$ below combines $-(\mathrm{py}/48)\#\mathfrak N$ of item~(3) of Substep~5 with $+(\mathrm{py}/96)\#\mathfrak N$ of item~(4) of Substep~5. The integrand is at most \[\mathrm{Res}\,\exp\{\sigma''\#\mathfrak Z_2-(\mathrm{py}/96)\#\mathfrak N-\rho_Q\#\beta_Q-\frac x2\Phi(\varphi)\},\] where $\Phi=\langle\varphi,(S-S')\varphi\rangle+\Psi$; here $\mathsf{co}_i:=1-2\check\theta_x-\tau_i\ge\frac12$ and \[\Psi(\varphi):=\|k\varphi\|^2\text{ over the rows off }\mathfrak Z+\sum_{\mathsf{qu}}\mathsf{co}_i|k_i\varphi-\mathsf{sh}_i|^2+\kappa_1\sum_{\mathrm{Bd}(\mathfrak N)}|\varphi_b-a_{0,b}|^2,\] and $\Psi_2$ denotes the quadratic part of $\Psi$. We also set $\mathsf{Ez}:=\sum_{\mathsf{qu}}\mathsf{co}_i|\mathsf{sh}_i|^2+\kappa_1\sum_{\mathrm{Bd}(\mathfrak N)}|a_{0,b}|^2$ and $\mathbb A_2:=(S-S')+\Psi_2$.

\emph{Coercivity.} In $\Psi_2$, each plaquette row whose corners all lie in quiet blocks, and each bond row whose two ends are quiet, appears with a coefficient $\ge\frac12$, anywhere on the torus. These rows depend on the bonds off $\mathrm{Bd}(\mathfrak N)$ alone, and on those bonds their sum is $\ge\gamma$, by \Cref{Y-F-L14}, which we state and prove now.

\begin{lemma}[A lower bound for the form away from the loud blocks]\label{Y-F-L14}
For every $\varphi\in\mathfrak L_{\mathfrak g}$, \[\sum_{b\in I^\circ\text{ off }\mathrm{Bd}(\mathfrak N)}|\varphi_b|^2\le73L^4\bigl[\sum_{p\text{ with all corners in quiet blocks}}|(\mathrm{d}^V\varphi)_p|^2+a\sum_{c\text{ with quiet ends}}|(Q_1^V\varphi)_c|^2\bigr].\]
\end{lemma}

\begin{proof}
Apply \eqref{eq:comb-local-poincare} to each quiet block and
\eqref{eq:crossing-local-poincare} to each pair $c$ with quiet ends,
after the blockwise unitary transports used in the proof of
\Cref{Y-C-CIL64}(a). These estimates give, respectively,
\[
 \sum_{b\text{ inside }B(y)}|\varphi_b|^2
 \le\tfrac12L^3\sum_{p\text{ inside }B(y)}|(\mathrm{d}^V\varphi)_p|^2
\]
and
\[
 \sum_{b\text{ crossing }c}|\varphi_b|^2
 \le6L^2|(Q_1^V\varphi)_c|^2
    +12L^4\sum_{p\text{ inside }W_c}|(\mathrm{d}^V\varphi)_p|^2.
\]
Both estimates use only the bonds inside the indicated block or
pair of blocks. A bond of $I^\circ$ off $\mathrm{Bd}(\mathfrak N)$ lies in a quiet block, or it is a crossing bond of some $c$ with two quiet ends. We sum the first bound over the quiet blocks and the second over these $c$; a plaquette lies in at most six of the $W_c$. Finally, $L^3+72L^4\le73L^4$ and $6L^2\le73L^4a$.
\end{proof}

Hence $\Psi_2\ge\min\{\gamma/2,\kappa_1\}=\kappa_1$. From $|u-e|^2\ge\frac12|u|^2-|e|^2$ for vectors $u$, $e$ we obtain $\Psi\ge\frac12\langle\varphi,\Psi_2\varphi\rangle-\mathsf{Ez}$, and so $|\varphi|^2\le(2/\kappa_1)(\Psi+\mathsf{Ez})$. We have $s_S=\|S\|+\|S'\|\le2x^{-60}\le\kappa_1/2$. Indeed, $\kappa_1=\mathrm{py}/(240\pi^2L^3x)\ge(2369L^3x)^{-1}$ because $\mathrm{py}\ge1$ \textup{(F9$^\flat$)}, and $4\cdot2369L^3x\le x^{60}$ because $10^6v^2\le x^8$ \textup{(F9$^\flat$)}. It follows that $\Phi\ge\Psi(1-2s_S/\kappa_1)-2s_S\mathsf{Ez}/\kappa_1\ge-2s_S\mathsf{Ez}/\kappa_1$.

At every $V'$ we have \[|\mathsf{sh}_i|\le(1+a^{1/2})(1+l_r)(\pi(1+(4+L)a_\infty^\flat)+1).\] For this, note that $|\bar F_p|\le\pi+4\pi a_\infty^\flat$, $|(\ell^V\bar F_F)_y|\le l_r(\pi+4\pi a_\infty^\flat)$ and $|(Q_1^Va_0)_c|\le L\pi a_\infty^\flat$ (\Cref{Y-F-LA}\,(a)), and that $\mathsf{sh}_c=a^{1/2}(-(Q_1^Va_0)_c+\omega_{*c}/(1-\theta_2))$. The summand $1$ covers $\nu_*$, $\sigma_*$ and $\omega_*$, by \Cref{Y-F-LM} and $\varepsilon'\le1/40$ \textup{(FA)}: $|\nu_*|\le N_1L^{-2}\varepsilon'\le \mathsf N_\nu^\flat L^{-3/2}\le1$ by \textup{(L1$^\flat$)}; $|\sigma_*|\le \mathsf N_\sigma L^{-2}\varepsilon'\le \mathsf N_\sharp L^{-2}\le1$, because $3K\mathsf g_*\le1+15K^2\mathsf g_*^2\le \mathsf P_\rho$, and by \textup{(O4)}\,(iii); and $|\omega_*|/(1-\theta_2)\le1.03\cdot\frac29\mathsf N_ZL^{1/2}\varepsilon\le1$ \textup{(F5)}.

With $(\dots)=\pi(1+(4+L)a_\infty^\flat)+1$, these components give $|\mathsf{sh}_p|\le(\dots)$, $|\mathsf{sh}_y|\le(1+l_r)(\dots)$ and $|\mathsf{sh}_c|\le a^{1/2}(\dots)$: neither $\mathsf{sh}_p$ nor $\mathsf{sh}_y$ has a factor $a^{1/2}$, which only $\mathsf{sh}_c$ carries. So $|\mathsf{sh}_i|^2\le\max\{1,a\}(1+l_r)^2(\dots)^2\le(1+a)(1+l_r)^2(\dots)^2$. A cell carries at most $4.12v$ rows, and the second sum of $\mathsf{Ez}$ is at most $\kappa_1\,5L^3\,27M^3(\pi a_\infty^\flat)^2$ per cell; this number is at most $(5v-4.12v)(1+a)(1+l_r)^2(\dots)^2$, which is what these rows leave of the bound $5v(1+a)(1+l_r)^2(\dots)^2$ for one cell. For the last inequality, note that $5L^3\cdot27M^3=135v$, that $(\dots)\ge\pi(4+L)a_\infty^\flat\ge7\pi a_\infty^\flat$ as $L\ge3$, and that $\kappa_1\le\gamma/2$ by \textup{(U0$^\flat$)}, where $\gamma=(73L^4)^{-1}$. Hence $135\kappa_1(\pi a_\infty^\flat)^2\le0.88\cdot49(\pi a_\infty^\flat)^2$, and therefore \[\mathsf{Ez}\le5v\#\mathfrak Z\,(1+a)(1+l_r)^2(\pi(1+(4+L)a_\infty^\flat)+1)^2,\] and, by \textup{(F9$^\flat$)}, \[\frac x2\,2s_S\mathsf{Ez}/\kappa_1\le2.2\cdot10^6L^3v(1+a)(1+l_r)^2(1+La_\infty^\flat)^2x^{-58}\#\mathfrak Z\le x^{-40}\#\mathfrak Z_2.\] The numerals: $s_S\le2x^{-60}$; $1/\kappa_1\le2369L^3x$; $4+L\le\frac73L$ gives $\pi(1+(4+L)a_\infty^\flat)+1\le7.34(1+La_\infty^\flat)$; and $2\cdot2369\cdot5\cdot53.9\le1.28\cdot10^6$.

About its minimiser, $\Phi=\min\Phi+\langle\cdot,\mathbb A_2\cdot\rangle$, where $\mathbb A_2\ge\kappa_1-s_S\ge\kappa_1/2>0$; hence the integral is $\le e^{x^{-40}\#\mathfrak Z_2}(\det C^{-1}/\det\mathbb A_2)^{1/2}$.

\emph{The determinants.} Let $\mathbb A_1:=(S-S')+\sum_{i\text{ off }\mathfrak Z}k_i^*k_i+\sum_{i\in\mathfrak Z}\mathsf{co}_ik_i^*k_i$, and let $\mathbb A_{3/2}:=\mathbb A_1-\sum\mathsf{co}_yk_y^*k_y$, where the sum runs over the site rows $y$ of $\hat{\mathfrak Z}$ that are not quiet. Then \[\det C^{-1}/\det\mathbb A_2=(\det C^{-1}/\det\mathbb A_1)(\det\mathbb A_1/\det\mathbb A_{3/2})(\det\mathbb A_{3/2}/\det\mathbb A_2),\] and we bound the three quotients one after the other.

\begin{enumerate}
\item[\textup{First.}] We have $C^{-1}-\mathbb A_1=S'+\sum_{\mathfrak Z}(1-\mathsf{co}_i)k_i^*k_i\ge0$, as $S$ cancels; so the first quotient is $\le\exp\operatorname{tr}[\mathbb A_1^{-1}(C^{-1}-\mathbb A_1)]$. This rests on the following fact: for symmetric $0<A_2\le A_1$, $\det A_1/\det A_2=\det(1+X)\le e^{\operatorname{tr}X}$, where $X=A_2^{-1/2}(A_1-A_2)A_2^{-1/2}\ge0$ and $\operatorname{tr}X=\operatorname{tr}[A_2^{-1}(A_1-A_2)]$; ``Second'' uses it as well. Next, $\mathbb A_1\ge(1-2s_S/\gamma)\sum\mathsf{co}_jk_j^*k_j\ge\gamma/4$, and for the plaquette and bond rows $k_i\mathbb A_1^{-1}k_i^*\le2/(1-2s_S/\gamma)$ on $\mathfrak g$. The lower bound for $\mathbb A_1$ holds because $\mathsf{co}_j=1$ off $\mathfrak Z$ makes $\sum\mathsf{co}_jk_j^*k_j\ge q/2\ge\gamma/2$, while $S-S'\ge-s_S$. The upper bound for the rows holds because $A>0$, $A\ge c\,k^*k$ and $c>0$ imply $kA^{-1}k^*\le1/c$: for $B=kA^{-1/2}$ we have $cB^*B\le1$, hence $cBB^*\le1$; here $c=(1-2s_S/\gamma)\mathsf{co}_i\ge(1-2s_S/\gamma)/2$.

The trace of a plaquette or bond row is $\le6.1$: there are three components, and $6/(1-2s_S/\gamma)\le6.1$ requires $2s_S/\gamma\le1/61$. This holds: $s_S\le2x^{-60}$ and $\gamma^{-1}\le x^4$ by \textup{(F9$^\flat$)}, so that $2s_S/\gamma\le4x^{-56}\le1/61$, because $x\ge L^2\mathrm e\ge24$. The same inequality yields $\mathbb A_1\ge\gamma/4$ and $\mathbb A_{3/2}\ge\gamma/4$. For the site rows we use a bound free of the dimension of $\mathsf H_y$: $\operatorname{tr}\le(4/\gamma)\operatorname{tr}k_yk_y^*\le(4/\gamma)\,3\cdot4\cdot4\sum_p|\ell_y(p)|^2\le c_{\rm tr}$. Both trace bounds are $\le12c_\tau$, and a cell carries at most $5v$ rows; so the square root of the first quotient is $\le e^{c_\tau\#\beta_Q+60c_\tau v\check\theta_x\#\mathfrak Z+x^{-40}\#\mathfrak Z_2}$.

The term $x^{-40}\#\mathfrak Z_2$ arises as follows: $S'\ge0$ has rank $\le3\#\beta_Q+9v\#\mathfrak b\le12v\#\mathfrak Z_2$ and norm $\le x^{-110}$ (Substep~6), and $\mathbb A_1^{-1}\le4/\gamma$; hence $\frac12\operatorname{tr}[\mathbb A_1^{-1}S']\le(2/\gamma)12vx^{-110}\#\mathfrak Z_2\le x^{-40}\#\mathfrak Z_2$, because $24v/\gamma\le x^8$ \textup{(F9$^\flat$)}. A cell thus receives three terms $x^{-40}$, namely the one of Step~7 of the proof of \Cref{Y-F-C7}, the one with $\mathsf{Ez}$, and this one; and $3x^{-40}\le x^{-30}$, the last term of \textup{(Z1$^\flat$)}.

\item[\textup{Second.}] Here $\mathbb A_1-\mathbb A_{3/2}\ge0$, and each plaquette row and each bond row of the torus enters $\mathbb A_{3/2}$ with a coefficient $\ge\frac12$; thus $\mathbb A_{3/2}\ge\frac12q^V_0-s_S\ge\gamma/4$. The second quotient is $\le\exp\{(4/\gamma)\sum\operatorname{tr}k_yk_y^*\}\le\exp\{c_{\rm tr}\#(\text{rows})\}$, with at most $L^3(2\varrho_L+5)^3\#\mathfrak N$ rows; so its square root is $\le e^{(c_{\rm tr}/2)L^3(2\varrho_L+5)^3\#\mathfrak N}$.

\item[\textup{Third.}] Here $\mathbb A_{3/2}-\mathbb A_2=\sum\mathsf{co}_ik_i^*k_i-\kappa_1\mathbf 1_{\mathrm{Bd}(\mathfrak N)}$, the sum running over the plaquette and bond rows of $\mathfrak Z$ off $\mathsf{qu}$, and the rank of this difference is $\le3[7.2L^3+6+5L^3]\#\mathfrak N\le38L^3\#\mathfrak N$. Indeed, a plaquette row off $\mathsf{Pq}$ has a corner in a loud block, so its base site lies in that block enlarged by one layer below: $3(L+1)^3\le7.2L^3$ for $L\ge3$, since $3(4/3)^3\le7.12$. A bond row off $\mathsf{Cq}$ has a loud end, and there are six of these to a block. Moreover $\#\mathrm{Bd}(\mathfrak N)\le5L^3\#\mathfrak N$ (Substep~5). Each row takes values in $\mathfrak g$, whence a factor $3$; and $3(12.2L^3+6)=36.6L^3+18\le38L^3$ for $L\ge3$.

Further, $\|\mathbb A_{3/2}\|\le\|C^{-1}\|\le b_h+1$. We apply \Cref{Y-A-L3}\,(b) with $A_1=\mathbb A_{3/2}$, $A_2=\mathbb A_2\ge\kappa_1/2$, $r\le38L^3\#\mathfrak N$. The resulting bound does not depend on the hypothesis $\|A_1\|\,\|A_2^{-1}\|\ge1$ of that lemma: by the proof of \Cref{Y-A-L3}\,(b), all but $r$ eigenvalues of $A_2^{-1/2}A_1A_2^{-1/2}$ equal $1$, and the others are $\le\|A_1\|\,\|A_2^{-1}\|\le2(b_h+1)/\kappa_1$, a number $\ge1$ because $b_h\ge16$ and $\kappa_1\le\gamma/2$. So $\det A_1/\det A_2\le(2(b_h+1)/\kappa_1)^r$, and the square root of the third quotient is $\le(2(b_h+1)/\kappa_1)^{19L^3\#\mathfrak N}$.
\end{enumerate}

By \textup{(U1$^\flat$)}, ``Second'' and ``Third'' together contribute $\le e^{(\mathrm{py}/192)\#\mathfrak N}$; and $\mathrm{py}/48-\mathrm{py}/96-\mathrm{py}/192=\mathrm{py}/192$.

\emph{The tally for one cell of $\mathfrak Z_2$.} To $\sigma''$, Substep~7 adds $x^{-40}$ from the term with $\mathsf{Ez}$, and $60c_\tau v\check\theta_x=60c_\tau vx^{-1/4}$ and $x^{-40}$ from ``First'', where $\#\mathfrak Z\le\#\mathfrak Z_2$. Further, $\frac32vc_\xi^2p_0^2x^{-2}+3v(c_q^\flat)^2x^{-1}\le3v((c_q^\flat)^2+c_\xi^2p_0^2)x^{-1}$ and $3x^{-40}\le x^{-30}$. Hence \[\FitDisplay{\dimexpr\linewidth-1pt\relax}{\begin{aligned}
&\sigma''+60c_\tau vx^{-1/4}+2x^{-40}
\\
&\quad{}\le\sigma_{\rm tot}^\flat+5K_s^\flat p_0^3x^{-1/2}
\\
&\quad{}+(6.3(\mathsf N_\nu^\flat)^2\vartheta_F^\flat+3\mathsf N_Z^2(\theta_2+L\vartheta_n^\flat))M^3p_0'^2
\\
&\quad{}+(Q_\theta^\flat+60c_\tau v)x^{-1/4}
\\
&\quad{}+3v((c_q^\flat)^2+c_\xi^2p_0^2)x^{-1}+x^{-30}
\\
&\quad{}\le\sigma'.
\end{aligned}}\]
\end{proof}

\emph{Step 9: clauses \textup{(c$^\circ$)}, \textup{(c)}, \textup{(c$'$)}, \textup{(c$''$)}.} Of the clauses of \Cref{Y-F-C7}, as stated in item~H of \S\ref{Y-F-1}, each of Steps~9 to~12 proves those that its head names. Clause \textup{(c$^\circ$)} is not a case of \textup{(c)}: when $\mathfrak X$ is empty, \textup{(c)} bounds $\hat\lambda\,1\{D(U)\cap\hat{\mathfrak Z}\text{ empty}\}$, which is a smaller function.

\begin{enumerate}
\item[\textup{(c$^\circ$)}] We apply \Cref{Y-F-PB} with $\mathsf{th}=1$. There are at most $(3v)^{\#t}$ maps $\beta$. We fix $\mathfrak N$ and $\beta$ and let $n_L$ be the number of cells of type L, so that $\#\beta_Q+n_L=\#t$. We split $\frac1{192}=\frac1{384}+\frac1{768}+\frac1{768}$ and use the parts as follows.
\begin{enumerate}
\item[\textup{(i)}] By the last assertion of \Cref{Y-F-PB}, $e^{-(\mathrm{py}/768)\#\mathfrak N}\le\exp\{-\mathrm{py}\,n_L/(768(2r_1+1)^3)\}$.
\item[\textup{(ii)}] If $n_L\ge1$, then $\mathfrak N$ is not empty, and $e^{-(\mathrm{py}/768)\#\mathfrak N}\le e^{-\mathrm{py}/768}$.
\item[\textup{(iii)}] We write $e^{-\rho_Q\#\beta_Q}=e^{-(\rho_Q/2)\#\beta_Q}e^{-(\rho_Q/2)\#\beta_Q}$; the second factor is $\le e^{-\rho_Q/2}$ if $\#\beta_Q\ge1$.
\end{enumerate}
If $t$ is not empty, then $n_L\ge1$ or $\#\beta_Q\ge1$, and (ii), (iii) give $u_1$. What is left for the tested cells is $m_8^{\#t}$, with $m_8=\max\{e^{c_\tau-\rho_Q/2},e^{-\mathrm{py}/(768(2r_1+1)^3)}\}$, and the sum over $\beta$ gives $(3vm_8)^{\#t}=u_t^{\#t}$. The sum over $\mathfrak N$ obeys \[\sum_{\mathfrak N}e^{-(\mathrm{py}/384)\#\mathfrak N}\le\exp\{e^{-\mathrm{py}/384}M^3\#\mathfrak Z^{+1}\}\le\exp\{27M^3e^{-\mathrm{py}/384}\#\mathfrak Z_2\}.\] All these numbers, each counted for one cell of $\mathfrak Z_2$, add up to the left side of \textup{(Z1$^\flat$)}, which is $\le x^{-1/8}$.

\item[\textup{(c)}] On the support of $1_{\rm ind}$ each $X$ holds a $p$ with $\phi_p>\varepsilon$; since $x\varepsilon\cdot\varepsilon=p_0^2$, sentence~4 of \Cref{Y-F-C7}\,(a) gives $|\lambda^o|\le\hat\lambda$ times constants. Then \textup{(c$^\circ$)} applies.

\item[\textup{(c$'$)}] We apply \Cref{Y-F-PB} with $\mathsf{th}=1_{\rm ind}$. On the support of this indicator each $X$ holds a plaquette $p_X=(z;\mu\nu)\in D\cap\hat{\mathfrak Z}$; as $p_X\in W_P$ for $P=(y(z);\mu\nu)$, this $P$ is loud, and so is $B(y(z))$. Two polymers are more than $s(x)\ge8$ cubes apart, at fine distance $\ge8M+1>L-1$, so their blocks $B(y(z))$ differ; hence only the contributions with $\#\mathfrak N\ge\#\mathfrak X$ are nonzero. Now we split $\frac1{192}=\frac1{384}+\frac3{1152}$: the three factors $e^{-(\mathrm{py}/1152)\#\mathfrak N}$ give $\exp\{-\mathrm{py}\,n_L/(1152(2r_1+1)^3)\}$, then $e^{-\mathrm{py}/1152}$ if $n_L\ge1$, and $e^{-(\mathrm{py}/1152)\#\mathfrak X}$. The rest of the argument is that for \textup{(c$^\circ$)}, with $u_t'$, $u_1'$ in place of $u_t$, $u_1$. The weak cap is the constant $Ae^{-\kappa(|X|_c-1)}$.

\item[\textup{(c$''$)}] Of the indicator, the argument uses only that it is $\le1$ and that on its support each $X$ holds a plaquette of $D\cap\hat{\mathfrak Z}$; moreover $\hat{\mathfrak b}\subset\hat{\mathfrak Z}$.
\end{enumerate}

\emph{Step 10: clauses \textup{(a)}, \textup{(a$'$)}, \textup{(b)}, \textup{(d$^\circ$)}, \textup{(d)}, \textup{(d$'$)}, \textup{(e)}.}

\begin{enumerate}
\item[\textup{(a)}] The tests of $\hat\lambda$ lie in $\mathfrak Z_2$. The factor $[A]^\ell_{\mathfrak Z}1_{\rm ind}$ involves $\eta$ on the bonds of the $p\in\hat{\mathfrak Z}$, of the $W_c$ for $c\in\mathfrak Z$, and of $\hat{\mathfrak Z}$; and, for each $y\in\hat{\mathfrak Z}$, on the bonds of the plaquettes within $\varrho(x)+1$ of $y$ (through the row, through $\zeta$, and through both radii of the mask, where $\rho_0\le\varrho(x)$) and of the paths $\Gamma(y,\pi(p))$. All of these lie at coarse distance $\le\varrho_L+2$ from $\hat{\mathfrak Z}$. Next, $\eta_b=(G\xi)_b-a_{0,b}$ depends on $\xi$ on $Y_b$, that is, within $r_w$, and $\varrho_L+2+r_w\le r_1M$; so everything lies under $\mathfrak Z^{+r_1}=\mathfrak Z_2$. The quantities $a_0$, $\boldsymbol\nu$, $w_j$ do not depend on $\xi$.

For compatible data, cells of $\mathfrak b_i$ and of $\mathfrak b_j$ are farther than $4r_1+2$ apart; hence the $\mathfrak b_i^{+2r_1}$ are disjoint, and $J_0$ lies off them. The law of $\xi$ and the tilt are products over the bonds of $I^\circ$; the factors are Borel in $\xi$ (by the splitting in Substep~5 of Step~8), they are nonnegative, and they depend on disjoint sets of bonds: this is sentence~3. Sentence~4: $|w|\le$ its cap, and on the support of $1_{\rm ind}$ each $X$ holds a $p$ with $\phi_p>\varepsilon$. The second assertion of sentence~4, that the strong cap is at most a constant where $1_{\rm ind}=1$, is needed because, as a function, the strong cap involves $U$ on the whole of $X$, and $X$ may reach beyond $\hat{\mathfrak Z}_2$.

\item[\textup{(a$'$)}] The member $w(X,\cdot)$ depends on $U$ on the bonds of $X^{+s(x)}$; $U_b$ depends on $\eta_b$ and $V'$; and $\eta_b$ depends on $\xi$ within $r_w$. Cells more than $r_1$ apart have their sites at coarse distance $\ge r_1M+1\ge R_2+1>2r_w+2$. The letters are Borel in $\xi$: $\Lambda_{\check\zeta}$ by the splitting in Substep~5 of Step~8, and $w(X,\cdot)$ because it is a Borel function of $U$ (item~C of \S\ref{Y-F-1}) while $U=\mathsf U(\eta;V')$ is continuous in $\xi$. Borel functions of disjoint sets of the $\xi_b$ are independent.

\item[\textup{(b)}] The vector $x^{1/2}\xi_b$ is a standard Gaussian of $\mathbb R^3$. Hence $E_\xi\,e^{\frac x2\varsigma_1|\xi_b|^2}=(1-\varsigma_1)^{-3/2}\le e^{1.6\varsigma_1}$, and a cell has at most $3v$ bonds. Further, $E_\xi[1\{|\xi_b|>c_\xi\varepsilon\}e^{\frac x2\varsigma_1|\xi_b|^2}]=(1-\varsigma_1)^{-3/2}\Pr\{|N_3|>\tau\}$ with $\tau=c_\xi p_0(1-\varsigma_1)^{1/2}$, and $\Pr\{|N_3|>\tau\}\le0.8(\tau+1/\tau)e^{-\tau^2/2}$. Indeed, for a standard Gaussian vector $N_3$ of $\mathbb R^3$, \[\Pr=(2/\pi)^{1/2}\int_\tau^\infty r^2e^{-r^2/2}dr=(2/\pi)^{1/2}[\tau e^{-\tau^2/2}+\int_\tau^\infty e^{-r^2/2}dr]\le(2/\pi)^{1/2}(\tau+1/\tau)e^{-\tau^2/2},\] and $(2/\pi)^{1/2}\le0.798$. Moreover $e^{\varsigma_1c_\xi^2p_0^2/2}(1-\varsigma_1)^{-1/2}\le1.01$ and $3\cdot0.8\cdot1.01\le2.5$. Over all bonds, $e^{4.8v\varsigma_1}\le e^{x^{-1/8}}$, by the term $6vx^{-1}$ of \textup{(Z1$^\flat$)}.

\item[\textup{(d$^\circ$)}] We have $\mathbf 1_{\mathsf Y}G=\mathbf 1_{\mathsf Y}G\mathbf 1_{\mathsf Y^{+r_1}}$, because $r_w\le r_1M$, and $\|G\|^2=\|C\|\le2/\gamma$ at real $V'$; hence $|\varphi_{\mathsf Y}|^2\le(2/\gamma)|\xi_{\mathsf Y^{+r_1}}|^2$, and we take $\varsigma_1=(2/\gamma)\varsigma\le x^{-2}$. We then use sentence~3 of \textup{(a)} for $\mathsf Y^{+r_1}$, and after it \textup{(c$^\circ$)} and \textup{(b)}, with $\#\mathsf Y^{+r_1}\le(2r_1+1)^3\#\mathsf Y$. When $n=0$, the argument rests on \textup{(b)} alone.

\item[\textup{(d), (d$'$)}] On the supports of the $1_{{\rm ind},i}$ the caps are replaced by constants; then the argument for \textup{(d$^\circ$)} applies, with \textup{(c)}, respectively \textup{(c$'$)}, for each factor in place of \textup{(c$^\circ$)}.

\item[\textup{(e)}] By \textup{(W)}, $\mathrm{py}=d_0p_0^2/2$; so $\mathrm{py}/768=d_0p_0^2/1536$ and $\rho_Q/2=3c_\xi^2p_0^2/(256h_*^2)$. The numbers $C_\varphi^\flat$ and $h_*$ depend on $(L,a,K_*,\lambda_*)$.
\end{enumerate}

\emph{Step 11: clause \textup{(f)}.} Fix any $e$. A gauge transformation which is $1$ at the block centres moves every $U$ into the chart; the first paragraph of the proof of \Cref{Y-C-L51}, the formula of the chart, shows this. Thus $U=\mathsf U(\eta;V')$ with an $\eta\in\mathfrak L_{\mathfrak g}$. Both sides of \textup{(f)} are invariant, and \Cref{Y-F-LB} is stated for any element $\eta$ of $\mathfrak L_{\mathfrak g}$.

Such an $\eta$ exists for every $U$ and $V'$, also where $U_b$ or $U_bV'(c)^{-1}$ is $-1$. The combs form a forest in which each tree carries one block centre (item~C.I~4 of \S\ref{Y-C-1}); on them, therefore, the transformation gives $U=1$. The map $Y\mapsto e^{iY}$ sends $\{Y\in\mathfrak g:|Y|\le\pi\}$ onto $G$: one diagonalises, and at $-1$ one takes $\pi$ times any unit vector of $\mathfrak g$. We choose $\eta_b$ with $e^{i\eta_b}=U_b$ at the bonds of $I^\circ$ inside blocks, and with $e^{i\eta_b}=U_bV'(c)^{-1}$ at a crossing bond of $c$. The choice is immaterial: \Cref{Y-F-LB} holds for every $\eta$ of $\mathfrak L_{\mathfrak g}$, and $F$, $\phi$, $\zeta$, $m_y$, $n_c$, $\mathsf n_c$, $\mathrm{de}$, $\mathrm{db}$ are functions of $(U,V',e)$ alone.

The invariance: $\mathrm{En}_\ell$ and $\sum_D\phi_p$ are invariant by the part ``As to invariance'' of item~(3) of Step~1 of the proof of \Cref{Y-E-C6}. The angle $\mathsf n_c$ is invariant since $u=1$ at the block centres and $\mathsf M(U^u)(c)=u(c_-)\mathsf M(U)(c)u(c_+)^{-1}$, by item~C.I~4 of \S\ref{Y-C-1}. The rest of \textup{(f)} does not depend on $U$.

We apply \Cref{Y-F-LB} with $\mathfrak Z=\mathfrak Z'=\mathfrak S=$ all cells and the starred fields; then $\hat{\mathfrak Z}\setminus\hat{\mathfrak Z}'$ is empty and $\mathrm{Rim}=0$. By \textup{(E1$_k$)}, \[\langle\boldsymbol\nu,\iota(F,n)\rangle=\langle Z_*,W^VF+\mathrm{d}'^Vn\rangle+\langle\Xi_*,\mathrm{d}^VF\rangle=\langle Z_*,X'\rangle-\langle Z_*,\mathrm{de}\rangle+\langle\Xi_*,\mathrm{db}\rangle,\] and the defects are absorbed by the two moduli. Moreover $\langle Z_*,X'\rangle-\frac12\|\boldsymbol\nu\|^2=L^{-1}\mathcal G_{\rm all}$, since $\|\mathrm{pm}(v_*)\|^2=L^{-1}\langle v_*,\hat R^Vv_*\rangle$. We apply \Cref{Y-F-L23k} to the cell datum $J=$ all classes, $t=$ all cells off $f$: then $\mathfrak b=\mathfrak Z=\mathfrak Z_2=$ all cells, and \textup{(H17)} holds for this datum, because $J$ consists of all classes and hence every plaquette of $D'$ lies in $D'_j$ for some $j\in J$. When $D'$ is empty, $J$ is empty and $t$ is the set of all cells. We drop $\mathrm{Sq}\ge0$. This gives $L\mathcal A_{\rm all}\ge\mathcal G_{\rm all}+L\,\mathrm{Pay}-L\,\mathrm{vol}$; with $M^3\#\text{cells}=N^3$ and $6N_1^2\le6.3(\mathsf N_\nu^\flat)^2$, this is \textup{(f)} with $e_2^\flat$.

\emph{Step 12: clauses \textup{(g*1)}, \textup{(g*2)}.}



\begin{enumerate}
\item[\textup{(g*1)}]
\begin{enumerate}
\item[\textup{1.}] Each component contains a cell: every polymer is adjacent to some object that has cells, and that object belongs to the same component; so $\mathfrak b_m$ is not empty. The set $t_m$ lies off $f_J$, which contains $f_{J_m}$, and the polymers keep the distance $>s(x)$; so $\check\zeta^{(m)}$ is a letter.
\item[\textup{2.}] Different components have no adjacent objects; so they are compatible, the $\mathfrak b_m$ are pairwise farther than $4r_1+2$ apart, and hence the $\mathfrak b_m^{+r_1}$, and likewise the $\mathfrak b_m^{+2r_1}$, are disjoint.
\item[\textup{3.}] Admitted: let $P\in D'$ lie in a cell $\square_P$ within $4r_1$ of $\mathfrak b_m$. By \textup{(H17)} for $(J,t)$, the class $j$ of $P$ satisfies $j\in J$. The cell $\square_P$ is a mark, hence a cell of $f_j$, and its distance from some cell of an object of $\check\zeta^{(m)}$ is at most $4r_1+2$; thus $f_j$ is adjacent to that object, and $j\in J_m$.
\item[\textup{4.}] The indicator: a $p$ of $D\cap\hat{\mathfrak Z}_m$ is covered by an $X$ containing its cube. Since a cell is made up of exactly $L^3$ fine cubes, that cube lies under $\square^{+r_1}$ for some cell $\square$ of $\mathfrak b_m$; so $X$ is adjacent to an object of $\check\zeta^{(m)}$, and $X\in\mathfrak X_m$. Conversely, an $X$ of $\mathfrak X_m$ that holds a plaquette of $D\cap\hat{\mathfrak Z}_k$ is adjacent to $\check\zeta^{(k)}$, so that $k=m$. Hence $1_{\rm ind}=\prod1_{{\rm ind},m}$.
\item[\textup{5.}] The set $\mathfrak Z_2\setminus\mathfrak b$ is the disjoint union of the $\mathfrak b_m^{+2r_1}\setminus\mathfrak b_m$; $[A]^\ell_{\mathfrak Z}$ is a product over the rows of $\mathfrak Z$; $\sum_Jw_j$, the sum in $\mathrm{Res}$ and $\#t$ split; and $u_1\le1$.
\end{enumerate}
\item[\textup{(g*2)}] The proof is that of \textup{(g*1)}, with cell adjacency alone in place of adjacency and with no polymer; of its items, 1, 2, 3 and 5 enter.
\end{enumerate}

\emph{Step 13: existence of $x_{\rm fl}$, what it depends on, and its place in the order of the constants.} The following claim is proved in items~\textup{(1)} to \textup{(4)} below. Every constant of item~F of \S\ref{Y-F-1} is computed from \textup{(O0)} to \textup{(O10)} and $x$, with the exception of $\mathrm{Res}_{(J,t)}$ and of the four quantities $e^\circ_{(J,t)}$, $e^{\circ\prime}_{(J,t)}$, $e^{o}_{\check\zeta}$, $e^{o\prime}_{\check\zeta}$ that are defined from it. These five also contain $\mathsf t_\rho$, the cell datum $(J,t)$ and the field $V'$, and none of them occurs in a floor of item~G. Every floor of item~G of \S\ref{Y-F-1} is a condition on $x$ and the letters \textup{(O0)} to \textup{(O10)} alone, and it holds for all large $x$. Hence $S_{\rm fl}$ is not empty, and $x_{\rm fl}=1+\inf S_{\rm fl}$ exists; it is free of $N$, $c_h$, $E_a$, $w$, $V'$ and $e$, and likewise of $\ell$: it is one number for all rows $\ell$ of item~C of \S\ref{Y-F-1}. Of a row, only the pure numbers $b_*$, $K_*$, $\lambda_*$ and the bound $\varrho(x)$ of item~C on its range enter the floors of item~G. The three numbers are letters of \textup{(O0)}, and $\varrho(x)=R=\lceil c_R\log x\rceil$ is a function of $x$ given by $c_R$ of \textup{(O6)}. Consequently, $S_{\rm fl}$ and $x_{\rm fl}$ are determined by the letters \textup{(O0)} to \textup{(O10)} alone.

\begin{enumerate}
\item[\textup{(1)}] The constants of item~F depend on $K_*$, $\lambda_*$, $c_1^*$, $c_2^*$, $d_0$, $c_\Xi$, $\kappa$, $\rho_0$, $L$, $a$, $\gamma$, $b_h$, $\theta_{\mathbb C}$, $\theta_D$, $C_D$, $c_R$, $M$, $A_0$, $\mathfrak p_0$, $\mathsf G_0$, $C_\varphi^\flat$, $\mathsf K_{\rm c}^0$, $\mathsf a_{\rm c}$ and $x$. In item~F, $c_\Xi$ occurs only in $e^{o}$, and $\kappa$ only in $e^{o}$ and $e^{o\prime}$. The constant $\mathrm{Res}$ contains $\mathsf t_\rho$, on which no floor depends; it is a quantity of the pair $(\ell,\ell')$ of rows in item~C of \S\ref{Y-F-1}, with $\mathsf t_\rho\le\frac18(\frac12-d_0-c_\Xi)$. The quantity $\mathsf{ti}_0^\flat$ is given in item~F by the explicit expression $8C_D^2La_\nu^2x'^{-64}$ in the letters and $x$; it does not depend on the rows, and $\|\mathsf B'^{\rm c}\hat R^<-1\|_0\le\mathsf{ti}_0^\flat$ by item~(4) of Step~5. With the value $\mathsf n=60$ of \textup{(O9)}, item~D gives explicit formulas for $\tau_{\rm c}$, $\tau_2$ and $\mathsf{ti}_{\rm c}$ in terms of the letters and $x$.

\item[\textup{(2)}] The radii $R$, $\varrho_L$, $r_1$, $r(x')$ and $Mr(x')$ are at most $\log x$ times a number that depends only on the letters \textup{(O0)} to \textup{(O10)}. The left side of each of \textup{(F$_{\rm law}$)}, \textup{(P1)}, \textup{(P3)}, \textup{(FA)}, \textup{(F0$^\flat$)} to \textup{(F9$^\flat$)}, \textup{(U0$^\flat$)}, \textup{(Z0$^\flat$)}, \textup{(Z1$^\flat$)} is the product of such a number, a power of $\log x$, and one of $p_0^3x^{-1/2}$, $p_0^2\varepsilon$, $\varepsilon^{1/2}$, $x^{-1/4}$, $p_0^2x^{-1}$, $e^{-\mathrm{py}/384}$ or a negative power of $x$; it is compared with a constant, with $x^{-1/8}$, with $\varepsilon'$ or with a positive power of $x$. Four kinds of clauses do not have this form:
\begin{enumerate}
\item[\textup{(i)}] The inequality $\varepsilon\le\varepsilon'$ of \textup{(FA)}: both sides carry $x^{-1/2}$. Here $\varepsilon/\varepsilon'=L^{-1/2}(\log x/\log x')^{\mathfrak p_0}$, and $\log x/\log x'=1+\log L/\log x'$ decreases to $1$ as $x$ grows; so $\varepsilon\le\varepsilon'$ holds for exactly those $x$ with $\log x'\ge\log L/(L^{1/(2\mathfrak p_0)}-1)$, and these $x$ make up an upper half line.
\item[\textup{(ii)}] In \textup{(U0$^\flat$)} and \textup{(Z0$^\flat$)} the right sides $r_o=x^{-3/8}$, $0.09r_o$, $t_0=p_0^{1/2}x^{-1/4}$, $\check\theta_x/2=x^{-1/4}/2$ and $c_\xi\varepsilon/4$ decrease too. The left sides are a factor times $\varepsilon$, compared with $r_o$, $0.09r_o$, $\check\theta_x/2$; then $r_o$, compared with $t_0$; and a factor times $\varepsilon'^2+\varepsilon\varepsilon'$, compared with $c_\xi\varepsilon/4$. Each factor before $\varepsilon$, and each factor before $\varepsilon'^2+\varepsilon\varepsilon'$, is a number that depends only on the letters \textup{(O0)} to \textup{(O10)}. In every case the power of $x$ on the left is strictly lower, up to powers of $\log x$.
\item[\textup{(iii)}] The conditions $R\ge\max\{M,6\}$ and $\rho_0\le R(x')$ of \textup{(G1$^\flat$)}, $\mathrm{py}\ge1$ of \textup{(F9$^\flat$)} and $c_\xi p_0\ge1$ of \textup{(Z0$^\flat$)}: in each of them a constant is compared with a function of $x$ that grows without bound.
\item[\textup{(iv)}] The term $e^{-\mathrm{py}/384}$ in \textup{(Z1$^\flat$)}: in the setting \textup{(W)}, $\mathrm{py}=d_0A_0^2(\log x)^{2\mathfrak p_0}/2$ with $2\mathfrak p_0>1$, so $e^{-\mathrm{py}/384}$ decreases faster than any power of $x$.
\end{enumerate}
For the floors \textup{(G1$^\flat$)}, \textup{(U1$^\flat$)} and \textup{(W)} the comparison runs as follows.
\begin{enumerate}
\item[\textup{(G1$^\flat$)}] $Mr(x')\ge20c_R(\log x-\log L)$ is compared with $5R+4\varrho_L+12\le6.34c_R\log x+23$.
\item[\textup{(U1$^\flat$)}] $\log(1/\kappa_1)\le\log x+\text{const}$, and $(\log x)^3$ is compared with $\mathrm{py}$, which by \textup{(W)} is $d_0A_0^2(\log x)^{2\mathfrak p_0}/2$, where $2\mathfrak p_0>4$.
\item[\textup{(W)}] $x^{1/4}$ is compared with a power of $\log x$.
\end{enumerate}

\item[\textup{(3)}] The conditions on $N$ reduce to \textup{(Tor)}, which is a datum; see the paragraph on \textup{(P0)} between \textup{(P3)} and \textup{(FA)} in item~G of \S\ref{Y-F-1}.

\item[\textup{(4)}] The number $x_{\rm fl}$ is defined after \textup{(O10)} and before $c_h$, $E_a$; it enters \textup{(O13)} as $Lx_\flat\ge x_{\rm fl}$; and \textup{(O11)} depends on $\delta_3$, $\delta_3'$. Every condition is one-sided, and there is no circle.
\end{enumerate}
\end{proof}

\subsection{Polymer expansion of the step}\label{Y-G-1}

We assemble the local factors into the density of the coarse
field. Finite-range dependence factors the expectation over separated
sets of letters. A convergent cluster expansion sums the remaining
interactions. The analytic and rough terms of this expansion obey the
caps at $x'=x/L$, giving closure of the real class under the step.

\begin{theorem}[\textup{(S-asm$_k$)}: one step carries the class at $x$ into the class at $x/L$]\label{Y-G-ThA}
Let the letters be chosen in the order \textup{(O0)} to \textup{(O13)} of item~3 of \S\ref{Y-G0-1}. Assume those clauses of the following statements which Step~0 of the proof lists: \Cref{Y-C-CITh66} with \Cref{Y-C-P67}\,(a), (b) and with \Cref{Y-C-CIL64}\,(e); \Cref{Y-C-Th68}, \Cref{Y-C-L71} and \Cref{Y-C-CIICor72}; Thms.~\ref{Y-C-CIITh73}, \ref{Y-C-CIITh74} and \Cref{Y-C-CIICor75}; and further \Cref{Y-D-C5}\,(a) to (e), \Cref{Y-E-C6}\,(a) to (e) and \Cref{Y-F-C7}, which is made up of items~A to H of \S\ref{Y-F-1} (the objects of items~E1 to E12, the constants of item~F, the floors of item~G, the conclusions of item~H). Let $x\ge Lx_\flat$; we work in the setting \textup{(SD)} at $x$, with any data that it admits; and let $\rho=e^{\mathsf V}\rho[u;\ell]\in C^{\rm su}_\flat(x;\ell)$. The setting \textup{(SD)} names no level and no depth: the rows $\ell$ range over the rows of a chain at any level $k\ge0$ and over any other rows with \textup{(SD)}, the empty rows $\ell=\emptyset$ among them.

Then a real $\mathsf V'$ and a real $u'=(E'^a,E'^h,w')\in M_1^{\rm su}(x')$ with the strong cap at $x'$ are defined by \textup{(F1)} to \textup{(F4)}, which Steps~1, 6, 9 and 12 of the proof set out, and by the mean over $e$ of Step~13 for $w'$, which is the last clause of \textup{(F4)}; with them, at every $V'$, \[\rho'(V')=e^{\mathsf V'}\rho[u';\ell'^{\rm c}](V')>0.\] Hence the property \textup{(S-step)} holds: $\rho'\in C^{\rm su}_\flat(x/L;\ell'^{\rm c})$, where $\ell'^{\rm c}$ is the cut at $\varrho(x/L)$, with the same letters $d_0$, $c_\Xi$, $\rho_0$, $\kappa$, $c_R$, $M$, $A_0$, $\mathfrak p_0$, $\vartheta_0^\flat$, $c_h$, $E_a$ and with the caps at $x/L$; and $\rho'$ is a function of $V'$ alone. Moreover $E'^a=E'^{\rm new}+E'^{\rm old}[u]$, where $E'^{\rm new}$ depends on no member, so that it is a single family for all $\rho$. For every cell $\square$, \[\sum_{Y\ni\square}e^{\kappa(\#Y-1)}\sup_{\mathcal U_Y(x')}|E'^{\rm old}(Y,\cdot)|\le0.25\,c_a(x');\] the same sum with $E'^{\rm new}$ in place of $E'^{\rm old}$ is $\le0.31\,c_a(x')$; and the same sum with $E'^h$ in place of $E'^{\rm old}$, its supremum now taken over the tame real fields on $Y^\approx$, is $\le c_h(x')$.
\end{theorem}

\begin{proof}
\emph{Step 0: the inputs, clause by clause.} Below we give the clauses that are used in Steps~1 to 14 of this proof or in Step~15 of the proof of \Cref{Y-GB-ThB} and that are drawn from the statements that \Cref{Y-G-ThA} assumes or from the base. The hypotheses of each such clause, except those of \Cref{Y-A-Th29}\,(b) and \Cref{Y-A-Th36}, are met by the choice of the letters, by the setting \textup{(SD)} and by the floors \textup{(A0)} to \textup{(A11)}. \Cref{Y-A-Th29}\,(b) and \Cref{Y-A-Th36} are used in Step~15 alone, at the $\theta_D$, $C_D$, $c_R$ of \textup{(O6)}; at these values their use rests on the part ``Conclusion'' of \Cref{Y-A-Fr} and on Steps~H9 and H14 of its proof. The two size bounds of \textup{($\Gamma$3)} are applied only at fields that satisfy its hypothesis on the curvature: in Steps~2 and 7, and, in Step~10, within the parts of the proof of \Cref{Y-H-C9} on which this step draws.
\begin{enumerate}
\item[\textup{($\Gamma$2)}] \Cref{Y-C-CITh66}\,(a), (b), (c), (d), (f), (h), (i), (k) and \Cref{Y-C-P67}\,(a), on $\mathfrak V$. The matrix $G$ of the law has $G(b,b')=0$ unless ${\rm dist}(b,b')\le r_w$, and $(G\xi)_b$ depends on $\xi$ only through its restriction to the window $Y_b=\{b'\in I^\circ:{\rm dist}(b,b')\le r_w\}$. The row $G(b,\cdot)$ depends only on the restriction of $V'$ to the coarse bonds whose two ends lie within $r_G$ of $y(b)$. The matrix $C=GG^{\mathsf T}$ is symmetric, $C(b,b')=0$ unless ${\rm dist}(b,b')\le2r_w\le R$, and $C(b,\cdot)$ depends on $V'$ only within $3r_w+\varrho_L+1$ of $y(b)$. By \textup{(A1)} we have $\varrho_L+2\le\lceil R/2\rceil$, and with it $r_G\le R-1$ and $3r_w+\varrho_L+1\le R_2-1$.

For $G$-valued $u$ we have $G[V'^u]=\operatorname{Ad}_uG\,\operatorname{Ad}_u^{-1}$. Every entry is a uniform limit of polynomials in the entries of those $V'(c)$ on which it depends; it is continuous on $\mathfrak V$, holomorphic on $\mathfrak V^\circ$ and real at $G$-valued fields. The bounds $\|G\|_\infty,\ \|G^{\mathsf T}\|_\infty\le \mathsf G_0\le C_\varphi^\flat$ and $\|(G\xi)_b\|\le C_\varphi^\flat\sup_{Y_b}\|\xi_{b'}\|$ hold. Under \textup{(F$_{\rm law}$)}, and with $c_R\theta_{\mathbb C}\ge128$, the matrix $G$ is invertible and $\det G>0$ at real fields (\Cref{Y-C-CITh66}\,(h)); moreover $S=C^{-1}-q$ is symmetric and $\|S\|\le x^{-60}$.

With these we use \Cref{Y-C-CIL64}\,(e). Its hypotheses are $L\ge3$ odd, $aL\ge18$, and rows of the kind fixed in item~C.I~1.1 of \S\ref{Y-C-1}. The rows of \textup{(SD)} are of this kind, and the clause of item~C.I~1.1 on $N$ is treated below, among the clauses on $N$. The lemma gives the finite sum \[(\mathsf g^VX')_{b'}=\sum_P\sum_{\mathsf t}g_{\mathsf t}\operatorname{Ad}_{V'(\pi_{\mathsf t})}X'_P,\] where the numbers $g_{\mathsf t}$ depend on no field and each $\pi_{\mathsf t}$ is a coarse path from $y(b')$ to the base site of $P$, all of whose sites lie within $2\varrho_L+1$ of $y(b')$. So $\mathsf g^V$ is a polynomial in the $\operatorname{Ad}_{V'(c)}$; and at every field of $\mathfrak V$ and for every $X'$, the value $(\mathsf g^VX')_{b'}$ depends only on the restrictions of $V'$ to the bonds, and of $X'$ to the plaquettes, whose sites all lie within $2\varrho_L+2$ of $y(b')$.

\item[\textup{($\Gamma$3)}] \Cref{Y-C-Th68}\,(a), (c) and \Cref{Y-C-CIICor72}\,(a). The data of \Cref{Y-C-Th68}\,(a), (c) are a coarse site $o$, a number $0<\hat\varepsilon\le1/40$, and an $\operatorname{SL}(2,\mathbb C)$ field on the box of radius $R_3$ about $o$ with \textup{(e1)}, \textup{(e2)} and $|\partial V'(P)-1|\le2\hat\varepsilon$. By item~C.II~1.4 of \S\ref{Y-C-1}, the conditions \textup{(e1)}, \textup{(e2)} give \textup{(h1)}, \textup{(h2)}, two of the conditions on the field that \Cref{Y-C-Th68} assumes. For each bond $b$ with ${\rm dist}(b,o)\le1$, the theorem then gives two facts: $a_{0,b}$ is determined by the field restricted to that box, and $|a_{0,b}|\le \mathsf a_{\rm c}\hat\varepsilon$. \Cref{Y-C-CIICor72}\,(a), which is \textup{(Cen)}, adds \[|(EX'-\mathrm{d}^Va_0)_p|\le \mathsf N_{\rm c}L^{-2}\hat\varepsilon+\mathsf K_{\rm c}\hat\varepsilon^2\] for $y(p)=o$, with $\mathsf K_{\rm c}=\sqrt2\mathsf K_{\rm c}^0$. The hypotheses of \Cref{Y-C-Th68} and \Cref{Y-C-CIICor72} are \textup{(Fr1)} to \textup{(Fr3)}, which \textup{(SD)} provides, and $x\ge L^2$, $R\ge6$, \textup{(Tor)}, \textup{(P1)}, \textup{(P3)}; further, $L\ge3$ odd and $aL\ge18$, rows as in item~C.II~1.2, and a law with \textup{($\lambda$0)}, \textup{($\rho$1)}, \textup{($\rho$2)}, \textup{($\rho$5)} as in the proposition of item~C.II~2.1 of \S\ref{Y-C-1}. The conditions $L\ge3$ odd and $aL\ge18$ hold by \textup{(O4)}; the rows of \textup{(SD)} are rows as in item~C.II~1.2, whose clause on $N$ is among the clauses on $N$ below; and the law of \textup{($\Gamma$2)} has \textup{($\lambda$0)}, \textup{($\rho$1)}, \textup{($\rho$2)}, \textup{($\rho$5)}: it is defined on $\mathfrak V$, $C(b,b')=0$ unless ${\rm dist}(b,b')\le2r_w\le R$, $C(b,\cdot)$ depends on $V'$ only within $3r_w+\varrho_L+1\le R_2-1$ of $y(b)$, and, under \textup{(F$_{\rm law}$)} and $c_R\theta_{\mathbb C}\ge128$, the matrix $C$ is invertible and $\|S\|\le x^{-60}$, that is, $\mathsf n=60$.

The last clause of \textup{($\Gamma$3)} is the locality of $a_0$. It holds at every field of $\mathfrak V$ and for every $X'$ whose values lie in $\mathfrak g$ or in its complexification; in particular it holds at every pair $(V',e)$, rough plaquettes included, and, under the convention of item~1 of \S\ref{Y-G0-1}, at a complex field. It says that $a_{0,b}$ depends only on the restrictions of $V'$ to the bonds, and of $X'$ to the plaquettes, whose sites all lie within $R_3$ of $y(b)$, and on nothing else; and that, when ${\rm dist}(b,o)\le1$, the same is true with $o$ in place of $y(b)$. For the proof we write $a_{0,b}=\sum_{b'}C(b,b')(\mathsf g^VX')_{b'}$. By \textup{($\Gamma$2)}, $C(b,b')=0$ unless ${\rm dist}(b,b')\le2r_w\le R$, and $C(b,\cdot)$ depends on $V'$ only within $3r_w+\varrho_L+1\le R_2-1$ of $y(b)$. By \Cref{Y-C-CIL64}\,(e), $(\mathsf g^VX')_{b'}$ depends on $V'$ and on $X'$ only within $2\varrho_L+2$ of $y(b')$. Altogether, $a_{0,b}$ depends on them only within $\max\{R_2-1,\ R+2\varrho_L+2\}$ of $y(b)$, and therefore, when ${\rm dist}(b,o)\le1$, only within $\max\{R_2,\ R+2\varrho_L+3\}\le R_3$ of $o$, because $\varrho_L\le R+2$. The hypothesis $|\partial V'(P)-1|\le2\hat\varepsilon$ of \Cref{Y-C-Th68} is used for the two size bounds alone. Compare Step~(0) of the proof of \Cref{Y-C-Th68}, as well as Step~2~(ii), (iii) of the proof of \Cref{Y-E-C6}.

\item[\textup{($\Gamma$4)}] \Cref{Y-C-CIICor75}, with \Cref{Y-C-CIITh73}\,(b) and \Cref{Y-C-CIITh74}\,(b), (c). We put $\tilde{\mathsf D}:=-\mathsf D_k/2$. At every real $V'$ we have $(\det C^V)^{1/2}=(\det C^1)^{1/2}\exp\sum_Y\tilde{\mathsf D}(Y,V')$ and $S=\sum_YS_k(Y;V')$, where $Y$ runs over the connected sets of cells. Each $\tilde{\mathsf D}(Y,\cdot)$ is real, gauge invariant and $0$ at pure gauges, and $|\tilde{\mathsf D}|\le \mathsf K_1^\flat e^{-\kappa^\sharp(\#Y-1)}$ at every $V'$. It has a holomorphic continuation to $\mathcal U_Y(x')$, on which $|\tilde{\mathsf D}|\le1.5\mathsf K_1^\flat\varepsilon'e^{-\kappa^\sharp(\#Y-1)}$. The function $\mathsf D_k$ is defined in \Cref{Y-C-CIITheRest}, in the display that precedes \Cref{Y-C-CIITh74}: \[\mathsf D_k(Y,V')=-\mathsf D^q(Y,V')-[\mathsf D^u(Y,V')-\mathsf D^u(Y,1)].\] Its continuation is given by the same formulas, taken at the member of $\mathcal U_Y(x')$ extended by $1$, as in the proof of \Cref{Y-C-CIICor75}.

By \Cref{Y-C-CIITh73}\,(b), $\mathsf D^q(Y,\cdot)$ is holomorphic on the open set $\{|V'(c)|<1+\vartheta_0^\flat$ on the bonds of $Y^\sim\}$; by \Cref{Y-C-CIITh74}\,(b), $\mathsf D^u(Y,\cdot)$ is holomorphic on $\mathfrak V^\circ$; and each of these open sets contains every $G$-valued field. Hence $\tilde{\mathsf D}(Y,\cdot)$ is continuous on the $G$-valued fields. On $\mathcal U_Y(x')$ it is invariant under $G$-valued gauge transformations, by \Cref{Y-C-CIITh73}\,(b) and \Cref{Y-C-CIITh74}\,(b); the argument is given in Step~9, Invariance.

On $\mathfrak V$ the matrices $S_k(Y)$ satisfy $S_k(Y)(b,b')=0$ unless $b,b'\in\hat Y$, $S_k(Y)^{\mathsf T}=S_k(Y)$ and $S_k(Y;V'^u)=\operatorname{Ad}_uS_k\operatorname{Ad}_u^{-1}$; they are holomorphic on $\mathfrak V^\circ$ and real at real fields, and $\|S_k(Y)\|\le x^{-57}e^{-\kappa^\sharp\#Y}\le e_Y:=x^{-40}e^{-\kappa^\sharp\#Y}$. Both $\tilde{\mathsf D}(Y,\cdot)$ and $S_k(Y;\cdot)$ depend only on the restriction of $V'$ to the bonds of $Y^{+\lceil R_2/M\rceil}$, and $\lceil R_2/M\rceil\le r_1-8$.

The conditions of these statements are $\kappa\ge6$, which is \textup{(O5)}; \textup{(M4$^\flat$)}; \textup{(R$^\flat$)}; \textup{(M2)}; \textup{(Tor)}; \textup{(F$_{\rm law}$)}; $x\ge L^2$ and $R\ge6$, the two floors of \textup{(P0)}; \textup{(P1)}; \textup{(S1)}; and \textup{(S2)}. In \Cref{Y-C-CIIAtTheLawOfWindows}, they make up the list ``The conditions, whole''. Of them, \textup{(M4$^\flat$)} and \textup{(M2)} are two of the conditions on $M$ in \textup{(O7)}, \textup{(R$^\flat$)} is derived in \textup{(O6)} through four successive steps, the floors are among those of \textup{(A0)}, and \textup{(Tor)} is one of the data of \textup{(SD)}. As under \textup{($\Gamma$3)}, these statements also assume $L\ge3$ odd, $aL\ge18$, and rows as in item~C.II~1.2, its clause on $N$ included; these conditions are met on the grounds given there.

\item[\textup{($\Gamma$5)}] \Cref{Y-D-C5}, with the following parts.
\begin{enumerate}
\item[\textup{(e)}] At every $(V',e)$, \[x\,\mathsf m_C-\frac x2\langle a_0,Sa_0\rangle=x'\mathrm{En}'(V')+\sum_jw_j+\sum_\square1_{\square^\approx}(V')P(\{\square\},V'),\] where $w_j=\sum_{\square\in \mathrm{mk}_j'^{+r'}}T_\square$, for any partition of the marks into classes whose pairwise distance is $>2r'$. Each $\mathrm{mk}'_j$ is a class of the marks $\mathrm{mk}'(V')$ of item~E6 of \S\ref{Y-D-1}; item~E9 of \S\ref{Y-F-1} has the same set of marks, and its classes form one such partition.

\item[\textup{(b)}] The function $T_\square$, given by the formula of item~E7 of \S\ref{Y-F-1}, is real and jointly Borel; it satisfies $T_\square(V'^u,\operatorname{Ad}_ue)=T_\square(V',e)$; and it depends on $V'$ only within $4R+2\varrho_L+6\le R_6$ of $\square$.

\item[\textup{(d)}] The function $P(\{\square\},\cdot)$ is holomorphic on $\mathcal U_{\{\square\}}(x')$, invariant, real at real fields and $0$ at pure gauges, with $\sup|P(\{\square\},\cdot)|\le K_{\rm q}p_0'^2\varepsilon'$, the supremum taken over $\mathcal U_{\{\square\}}(x')$; and we put $P(Y):=0$ for $\#Y\ge2$.
\end{enumerate}

\item[\textup{($\Gamma$6)}] \Cref{Y-E-C6}, applied with the test radius $R_\chi$. The theorem allows any radius $\ge\lfloor R/2\rfloor+\varrho_L+2$. The test is $\chi_\square(\xi)=1$ if and only if $|\xi_b|\le c_\xi\varepsilon$ at all $b\in I^\circ$ within $R_\chi$ of $\hat\square$, and $V_\square$, $A_\square$ are those of items~3.8 and 3.9 of \S\ref{Y-E-1}. In \textup{(a)} to \textup{(d)}, $V'\in \mathcal U_{\{\square\}}(x')$ and $\chi_\square=1$.
\begin{enumerate}
\item[\textup{(a)}] The function $V_\square$ depends on $\xi$ only within $r_w+\varrho_L+1$ of $\hat\square$ and on $V'$ only within $R_v$ of $\square$; it depends on no member and on no $e$; and it is holomorphic in $V'$.

\item[\textup{(b)}] For $G$-valued $u$, $V_\square(\operatorname{Ad}_u\xi,V'^u)=V_\square(\xi,V')$.

\item[\textup{(c)}] The function $V_\square$ satisfies $|e^{V_\square}-1|\le2K_vp_0^3x^{-1/2}$.

\item[\textup{(d)}] At $G$-valued $V'\in \mathcal U_{\{\square\}}(x')$, $V_\square$ is real and $A_\square=e^{V_\square}$.

\item[\textup{(e)}] At every real $\xi$ and every pair $(V',e)$ the following hold.
\begin{enumerate}
\item[\textup{(e.i)}] For $X'(V',e)$, and also for any $\mathfrak g$-valued $X'$ put in its place, \[\prod_\square A_\square=1\{|\eta_b|<\pi\ \text{all }b\}\prod_b\mathsf j(\eta_b)\,e^{-x\mathrm{En}_\ell(U)}e^{\frac x2\langle F^\flat,\mathsf B^VF^\flat\rangle}\prod_ce^{\frac{ax}2|n^\flat_c|^2-ax(1-\cos\mathsf n_c)}.\]

\item[\textup{(e.ii)}] Let $\rho=e^{\mathsf V}\rho[u;\ell]$ with $u\in M_A^{\rm su}(x)$, real or complex; as the step is linear in $\rho$, the real part and the imaginary part of $\rho$ are treated separately, whereas in $u$ nothing is linear. Then \[\FitDisplay{\dimexpr\linewidth-1pt\relax}{\rho'(V')=\mathsf c_0z(ax)^{-3\#\mathbb T'}(2\pi/x)^{\frac32\#I^\circ}\det G^V\,e^{\mathsf V-x\mathcal Q(-a_0)}E_\xi\bigl[e^{-\frac x2\langle\xi,\mathsf E\xi\rangle+x\langle\xi,\mathsf Eh\rangle}\prod_\square A_\square\prod_Xe^{1_{X^\approx}(U)E(X,U)}\Xi_w(U)\bigr],}\] where $\mathsf E=G^{\mathsf T}qG-1$, $h=G^{\mathsf T}\mathsf g^VX'$ and $\mathcal Q(-a_0)=\mathsf m_C-\frac12\langle a_0,Sa_0\rangle$.

\item[\textup{(e.iii)}] The function $A_\square$ depends on $(V',e)$ only through the restriction of the pair to the bonds and plaquettes within $R_v$ of $\square$.

\item[\textup{(e.iv)}] For $G$-valued $u$, $A_\square(\operatorname{Ad}_u\xi,V'^u,\operatorname{Ad}_ue)=A_\square(\xi,V',e)$.

\item[\textup{(e.v)}] The function $A_\square$ is Borel in $(\xi,V',e)$.
\end{enumerate}
\end{enumerate}

\item[\textup{($\Gamma$7)}] \Cref{Y-F-C7}. The definitions that this clause uses are those of items~E9 to E12 of \S\ref{Y-F-1}. Item~E9 gives the marks $\mathrm{mk}'$, the set $f=\mathrm{mk}'^{+s_0}$, the classes $\mathrm{mk}'_j$ under the relation ``distance $\le2s_0+4r_1+2$'', and $f_j$, $D'_j$, $w_j$. Item~E10 gives $b_\square$, and item~E11 gives $[A]^\ell_{\mathfrak Z}=\prod_{\square\in\mathfrak Z}A_\square$. Item~E12 gives the cell datum $(J,t)$, with $\mathfrak b=f_J\cup t$, $\mathfrak Z=\mathfrak b^{+r_1}$ and $\mathfrak Z_2=\mathfrak b^{+2r_1}$; the function $\hat\lambda_{(J,t)}$; \textup{(H17)}; the objects and their adjacency; the letters $(J,t,\mathfrak X)$ and the properties of being compatible and of being admitted; $1_{\rm ind}$; and $\lambda^o$. The constants that occur in this clause are those of item~F of \S\ref{Y-F-1}.
\begin{enumerate}
\item[\textup{(a)}] By the first sentence of \Cref{Y-F-C7}\,(a), $\hat\lambda$ and $\hat\lambda1_{\rm ind}$ depend on $\xi$ only through its restriction to the bonds of $I^\circ$ under $\mathfrak b^{+2r_1}$.

\item[\textup{(d)}] We take this clause for $A=1$ and with the strong cap. Its data are letters $\check\zeta_i$, admitted and pairwise compatible; a finite set $J_0$ of cells lying off every $\mathfrak b_i^{+2r_1}$; a further finite set $\mathsf Y$ of cells; and $0\le\varsigma\le(\gamma/2)x^{-2}$. Then \[E_\xi\bigl[\prod_i|\lambda^o_{\check\zeta_i}|\prod_{J_0}b_\square\,e^{\frac x2\varsigma|\varphi_{\mathsf Y}|^2}\bigr]\le\exp\bigl\{4.8v(2/\gamma)\varsigma(2r_1+1)^3\#\mathsf Y+x^{-1/8}\bigl(\sum_i\#\mathfrak b_i^{+2r_1}+\#J_0\bigr)\bigr\}\,u_b^{\#J_0}\prod_ie^o_{\check\zeta_i},\] the case of no $\check\zeta$ being included.

\item[\textup{(e)}] By this clause, $u_1=e^{-\delta_3p_0^2}$.

\item[\textup{(g*1)}] We use this clause in item~2 of Step~7.
\end{enumerate}

The data of \Cref{Y-F-C7}, which item~C of \S\ref{Y-F-1} sets out, include a system $w$ that is complex Borel, that is a function of $U$ on the bonds of $X^{+s(x)}$, and that obeys the cap. The $w$ of $\rho$ has these properties by \textup{(c$\Xi$)}, with the strong cap at $A=1$; so they hold in \textup{(d)} and in \textup{(g*1)} alike.
\end{enumerate}

The base consists of the following statements: \Cref{Y-H-L51} and \Cref{Y-H-Th53}, in Step~2 of the proof of \Cref{Y-H-C9}; \Cref{Y-H-Th57}, in Step~8, \textup{(8.2)}, of the proof of \Cref{Y-H-C9}; \Cref{Y-GB-P63}, in Step~15, item~2, of the proof of \Cref{Y-GB-ThB}; Lems.~\ref{Y-G-LKP}, \ref{Y-G-L18}, \ref{Y-G-L19}, in Step~8; \Cref{Y-H-L52}, in Step~10~(9) of the proof of \Cref{Y-H-C9}; and Thms.~\ref{Y-A-Th29} and \ref{Y-A-Th36}, with \textup{(X12)}, in \S\ref{Y-A-1}. This proof uses these passages of the proof of \Cref{Y-H-C9}: in Step~2, \Cref{Y-H-L51}, \Cref{Y-H-Th53} and (ii); in Step~8, \textup{(8.2)}, and, in the part of that step headed ``The clauses without a size'', the extension by $1$ of the $G$-valued gauge transformation $u$ given on the sites of $Y^\approx$ or of $\mathsf B_\omega$; and in Step~10~(9), \Cref{Y-H-L52} with its proof and nothing more. These passages use statements of \S\ref{Y-C-1} and the identity of matrices in Step~4 of the proof of \Cref{Y-E-C6}. From Sections~\ref{Y-G0-1}, \ref{Y-G-1} and \ref{Y-GB-1}, they use the definitions in items~1 and 2 and the letters of item~3 of \S\ref{Y-G0-1}, and from Thms.~\ref{Y-G-ThA}, \ref{Y-GB-ThB} and~\ref{Y-B-C1} they use nothing. That identity of matrices is an input of this proof too.

We come to the clauses on $N$. The clause $2D_w+1<N$ of \textup{(D0)} and the clause $2\omega+1<N$ of \Cref{Y-H-Th57} follow from \textup{(Tor)}: indeed $2D_w+1\le(12Mr(x)+24)/L+5<16Mr(x)\le N$, and $2\omega+1\le2Mr(x)+1<N$, since $\omega\le Mr'$ by \textup{(A2)} and $r'=r(x')\le r(x)$. From \textup{(Tor)} we also obtain $\varrho/L+2<N/2$, a clause that item~C.I~1.1 of \S\ref{Y-C-1} imposes on the rows; $\varrho_L+2<N/2$, a clause that item~C.II~1.2 of that section imposes on the rows; and $2R_3+1<N$, which item~C.II~4 of \S\ref{Y-C-1}, in its part on the floors, obtains from \textup{(P0)}. The reason is that $\omega\ge R_3+2=2R+\varrho_L+7$, so that $\varrho/L+2\le\varrho_L+2<\omega<N/2$ for every range $\varrho\le\varrho(x)=R$, and $2R_3+1<2\omega+1<N$. These are the clauses on $N$ in the hypotheses of \Cref{Y-C-CIL64}, of \Cref{Y-C-CITh66} and \Cref{Y-C-P67}, and of \Cref{Y-C-Th68} and \Cref{Y-C-CIICor72}. At the empty rows the range is $0$, while $\varrho_L$ is still determined by the range $\varrho(x)$, not by the range $0$.

\emph{Step 1: \textup{(F1)}, the exact formula.} We claim that at every pair $(V',e)$ \begin{equation}\tag{F1}\rho'(V')=e^{\mathsf V_0}\exp\bigl\{-x'\mathrm{En}'(V')+\sum_Y\tilde{\mathsf D}(Y,V')-\sum_\square1_{\square^\approx}(V')P(\{\square\},V')\bigr\}\,\mathcal Z,\end{equation} where \[\mathcal Z=e^{-\sum_jw_j}E_\xi\bigl[e^{\sum_YT_Y}\prod_\square A_\square\prod_{X\in\mathcal D}e^{1_{X^\approx}(U)E(X,U)}\Xi_w(U)\bigr].\] Here $\xi\sim \mathcal N(0,x^{-1}1)$ on $\mathfrak L_{\mathfrak g}$, $\varphi=G^V\xi$, $a_0=C^V\mathsf g^VX'$, $\eta=\varphi-a_0$, $U=\mathsf U(\eta;V')$ is given by the chart of item~3.4 of \S\ref{Y-E-1}, and $E=E^a+E^h$; further, $T_Y=\frac x2\langle\varphi,S_k(Y)\varphi\rangle-x\langle\varphi,S_k(Y)a_0\rangle$ and $e^{\mathsf V_0}=e^{\mathsf V}\mathsf c_0z(ax)^{-3\#\mathbb T'}(2\pi/x)^{\frac32\#I^\circ}(\det C^1)^{1/2}$.
\begin{enumerate}
\item[\textup{1.}] The function $\rho$ is Borel, gauge invariant and bounded: $-x\mathrm{En}_\ell\le0$, the sum of the members is $\le1.22(c_a+c_h)$ times the number of cubes, and $\Xi_w$ is a finite sum on a finite torus. At a real $U$ with $1_{X^\approx}(U)=1$, the restriction of $U$ to $X^\approx$ lies in $\mathcal U_X(x)$, because
\begin{enumerate}
\item[\textup{(u1)}] $1<1+\vartheta(x)$,
\item[\textup{(u2)}] $1<2-1/\log x$,
\item[\textup{(u3)}] $|\partial U(p)-1|=2\sin(\phi_p/2)\le\varepsilon<2\varepsilon$.
\end{enumerate}
So $E^a(X,U)$ is defined.

\item[\textup{2.}] \Cref{Y-E-C6}\,(e.ii) holds at every pair $(V',e)$. Together with \textup{(e.i)} and with $\rho(U)=e^{\mathsf V}e^{-x\mathrm{En}_\ell(U)}\prod_Xe^{1_{X^\approx}(U)E(X,U)}\Xi_w(U)$, it gives \[\FitDisplay{\dimexpr\linewidth-1pt\relax}{\begin{aligned}
&\rho'(V')=\mathsf c_0z(ax)^{-3\#\mathbb T'}(2\pi/x)^{\frac32\#I^\circ}\det G^V\,e^{-x\mathcal Q(-a_0)}
\\
&\quad{}E_\xi\bigl[e^{-\frac x2\langle\xi,\mathsf E\xi\rangle+x\langle\xi,\mathsf Eh\rangle}
\\
&\quad{}1\{|\eta_b|<\pi\}\prod_b\mathsf j(\eta_b)\,
\\
&\quad{}\rho(U)\,e^{\frac x2\langle F^\flat,\mathsf B^VF^\flat\rangle}
\\
&\quad{}\prod_ce^{\frac{ax}2|n^\flat_c|^2-ax(1-\cos\mathsf n_c)}\bigr],
\end{aligned}}\] with $\mathsf E=G^{\mathsf T}qG-1$ and $h=G^{\mathsf T}\mathsf g^VX'$.

\item[\textup{3.}] By \Cref{Y-C-P67}\,(a), $\eta\!\cdot\!S\eta=-(\xi-h)\!\cdot\!\mathsf E(\xi-h)$, $a_0\!\cdot\!Sa_0=-h\!\cdot\!\mathsf Eh$ and $\mathcal Q(-a_0)=\mathsf m_C-\frac12\langle a_0,Sa_0\rangle$. With $\eta=\varphi-a_0$ and $S=\sum_YS_k(Y)$ this gives \[\sum_YT_Y=\frac x2\langle\eta,S\eta\rangle-\frac x2\langle a_0,Sa_0\rangle=-\frac x2\langle\xi,\mathsf E\xi\rangle+x\langle\xi,\mathsf Eh\rangle.\] Moreover $\det G^V=(\det C^V)^{1/2}$.

\item[\textup{4.}] We insert the expression for $\rho(U)$ given in item~2 and apply \Cref{Y-E-C6}\,(e): the factors that depend on no member are $\prod_\square A_\square$.

\item[\textup{5.}] We apply \textup{($\Gamma$4)} to $(\det C^V)^{1/2}$, and \Cref{Y-D-C5}\,(e) to $x\mathcal Q(-a_0)=x\,\mathsf m_C-\frac x2\langle a_0,Sa_0\rangle$, with the partition of the marks into the classes of item~E9 of \S\ref{Y-F-1}; their pairwise distance is $>2s_0+4r_1+2>2r'$. The factor $e^{-\sum w_j}$ depends on no $\xi$.
\end{enumerate}

Finally, the left side of \textup{(F1)} contains no $e$; hence its right side is the same number at every $e$.

\emph{Step 2: a tame coarse field and passed tests make the fine field tame.}

\begin{lemma}[A tame fine field from passed tests and a tame coarse field]\label{Y-G-LT}
\begin{enumerate}
\item[\textup{(T1)}] Let $p$ be a plaquette. If $V'$ is real, if $\phi'_P\le\varepsilon'$ on the box of radius $R_3$ about $y(p)$, and if every $b'\in I^\circ$ within $r_w$ of the four bonds of $p$ satisfies $|\xi_{b'}|\le c_\xi\varepsilon$, then $\phi_p(U)\le0.72\varepsilon$.

\item[\textup{(T2)}] Let $y$ be a fine site. If the hypotheses of \textup{(T1)} hold at every $p$ with $|y-\pi(p)|_\infty\le\varrho(x)+1$, then $m_y(U)=1$, and $\zeta_p=F_p$ on the range of $\ell_y$.

\item[\textup{(T3)}] Let $X\in\mathcal D$ with $\chi_X(\xi):=1\{|\xi_b|\le c_\xi\varepsilon$ at the bonds of $I^\circ$ within $R$ of the blocks that meet $X^\approx\}=1$, let $Y$ be a set of cells that contains $Y_0(X)$, and let either $V'\in \mathcal U_Y(x')$, or $V'$ be real and tame on $Y^\approx$, extended by $1$. In either case, the restriction of $\mathsf U(\eta;V')$ to the bonds of $X^\approx$ lies in $\mathcal U_X(x)$, and at a real $V'$ it is real and tame on $X^\approx$. The second case is contained in the first. Indeed, a field that is real on the bonds of $Y^\approx$ and tame there belongs to $\mathcal U_Y(x')$: for such a field the three conditions of item~1 of Step~1, taken at $x'$, read as follows.

\begin{enumerate}
\item[\textup{(u1)}] $1<1+\vartheta(x')$.

\item[\textup{(u2)}] $1<2-1/\log x'$.

\item[\textup{(u3)}] $|\partial V'(P)-1|=2\sin(\phi'_P/2)\le\varepsilon'<2\varepsilon'$ at the plaquettes of $Y^\approx$.
\end{enumerate}
\end{enumerate}
\end{lemma}

\begin{proof}[Proof of \Cref{Y-G-LT}]
\begin{enumerate}
\item[\textup{(T1)}] As $V'$ is real, \textup{(e1)} and \textup{(e2)} hold, and $|\partial V'(P)-1|\le\phi'_P\le\varepsilon'$ on the box of radius $R_3$ about $y(p)$. Taken at $\hat\varepsilon=\varepsilon'$ and $o=y(p)$, \textup{($\Gamma$3)} gives $|a_{0,b}|\le \mathsf a_{\rm c}\varepsilon'$ and $|(EX'-\mathrm{d}^Va_0)_p|\le \mathsf N_{\rm c}L^{-2}\varepsilon'+\mathsf K_{\rm c}\varepsilon'^2$. On $\mathfrak g$ the two norms agree, $|\cdot|=\|\cdot\|$; hence $|(G\xi)_b|\le C_\varphi^\flat c_\xi\varepsilon=\varepsilon/6.5$, and since transports are isometries, $|(\mathrm{d}^VG\xi)_p|\le(4/6.5)\varepsilon\le0.6154\varepsilon$. By \textup{(L2$^{\rm su}$)}, $\mathsf N_{\rm c}L^{-2}\varepsilon'\le \mathsf N_{\rm c}L^{-3/2}\varepsilon\le10^{-4}\varepsilon$. Moreover $|\eta_b|\le(1/6.5+L^{1/2}\mathsf a_{\rm c})\varepsilon\le r_\eta$.

We now apply \Cref{Y-H-L51} at $\tau=1$ and $r=r_\eta$; this lemma is stated in Step~2 of the proof of \Cref{Y-H-C9}. Its hypotheses are met: $|X'_P|=\phi'_P\le\varepsilon'$; $4.1r_\eta+2.1\varepsilon'\le1$ by \textup{(A1)}; and $\varepsilon'\le1/40$, since $17T\varepsilon'\le1/40$ by Part~3 of \textup{(A0)} and $T>1$ by \textup{(O4)}. It yields $|\partial U(p)-1-iF^\flat_p|\le0.72(4.1r_\eta+2.1\varepsilon')^2$, where $F^\flat=EX'+\mathrm{d}^V\eta$. Together with \textup{(A1)}, the bounds obtained so far give $|\partial U(p)-1|\le(0.6154+10^{-4}+0.1)\varepsilon=0.7155\varepsilon$. Finally, $|\partial U(p)-1|$ is the chord $2\sin(\phi_p/2)$, and $\varepsilon\le1/40$; therefore $\phi_p\le0.7156\varepsilon\le0.72\varepsilon$.

\item[\textup{(T2)}] The second clause of \textup{(T2)} is void at the empty rows. By \textup{(T1)}, $\phi_p\le0.72\varepsilon$ at every plaquette $p$ with $|y-\pi(p)|_\infty\le\varrho(x)+1$. Hence no plaquette of $D(U)$ lies within $\rho_0\le\varrho(x')\le\varrho(x)+1$ of $y$, and, as $0.72\varepsilon<1$, no plaquette with $\phi_p>1$ lies within $\varrho(x)+1$ of $y$; so $m_y=1$ by its definition in item~E2 of \S\ref{Y-F-1}. Moreover $\zeta_p=F_p\min\{1,\varepsilon/\phi_p\}=F_p$.

\item[\textup{(T3)}] We first locate the sets involved. By the part of item~1 of \S\ref{Y-G0-1} on $Y_0(X)$, every plaquette of $X^\sim$ has a cell of $Y_0(X)$ within $M/2$ coarse units, and $X^\approx$ reaches $Mr(x)$ fine sites farther out. So the blocks that meet $X^\approx$ lie within $M+\lceil Mr(x)/L\rceil$ of the sites of $Y_0(X)$. Their $D_w$-neighbourhood, and the boxes of radius $R_3$ about them, lie within $\omega$ of these sites, where $\omega\le Mr'$ by \textup{(A2)}; so they lie in $Y_0(X)^{+r'}$, which is inside $Y^\approx$. There, \textup{(u1)}, \textup{(u2)}, \textup{(u3)} at $x'$ give \textup{(e1)}, \textup{(e2)}, \textup{(e3)} at $\tau=1$. At a bond $b$ of $X^\approx$, the value $(G\xi)_b$ involves $V'$ only on the coarse bonds within $r_G\le R_3$ of $y(b)$, and $a_{0,b}$ only on the bonds and plaquettes within $R_3$ of $y(b)$; neither depends on the continuation.

We apply \Cref{Y-H-Th53}, which is stated in Step~2 of the proof of \Cref{Y-H-C9}, at $\tau=1$, with $B$ the set of sites within $\omega$ of $Y_0(X)$. Its hypotheses hold, as we now show.

\begin{enumerate}
\item[\textup{(h1)}] At the bonds of $X^\approx$, $|(G\xi)_b|\le\sqrt2\|(G\xi)_b\|\le(\sqrt2/6.5)\varepsilon\le0.2176\varepsilon$; for these bonds are based in blocks that meet $X^\approx$, or in blocks next to one, and $r_w+1\le R$. Moreover, by \Cref{Y-C-Th68}\,(c) at $o=y(b)$, $|a_{0,b}|\le \mathsf a_{\rm c}\varepsilon'\le L^{1/2}\mathsf a_{\rm c}\varepsilon$; therefore $|\eta_b|\le r_\eta$.

\item[\textup{(h2)}] We have $|(EX'-\mathrm{d}^Va_0)_p|\le10^{-4}\varepsilon+\mathsf K_{\rm c}\varepsilon'^2\le(10^{-4}+0.1)\varepsilon$. For the transports, $|\operatorname{Ad}_gY|\le|g|^2|Y|\le1.01|Y|$. Here $g$ is $1$ or $V'(c)$ for a single coarse bond $c$, because $T_j\in\{1,\operatorname{Ad}_{V'(c)}\}$ by Step~4 of the proof of \Cref{Y-E-C6}; so $|g|<1+\vartheta_0^\flat$ by \textup{(u1)} at $x'$, whereas \textup{(u2)} alone would give only the factor $4$. Thus $|(\mathrm{d}^VG\xi)_p|\le4\cdot1.01\cdot(\sqrt2/6.5)\varepsilon\le0.879\varepsilon$, and the two terms together are at most $0.98\varepsilon\le1.9\varepsilon$.
\end{enumerate}

The conditions \textup{(D0)}, \textup{(D2)}, \textup{(D3)} at
$(1,r_\eta)$ hold by \textup{(A0)} and Step~0.

At a real $V'\in\mathcal U_Y(x')$, the field $U$ is real.
To prove its tameness, use $|\partial V'(P)-1|<2\varepsilon'$
on the same boxes and apply \textup{($\Gamma$3)} with
$\hat\varepsilon=\varepsilon'$. Since the transports are now
isometries, the argument for \textup{(T1)} gives
\[
 |F^\flat_p|\le(0.6154+10^{-4})\varepsilon
                    +\mathsf K_{\rm c}\varepsilon'^2.
\]
Use the first case of \Cref{Y-H-L51}, with the principal
logarithm $X'_P$ and $(\tau,r)=(1,r_\eta)$. By
\textup{(A1)}, the sum of its error and
$\mathsf K_{\rm c}\varepsilon'^2$ is at most $0.1\varepsilon$.
Thus $2\sin(\phi_p/2)\le0.7155\varepsilon$, and the last
estimate in \textup{(T1)} gives $\phi_p\le0.72\varepsilon$.
This uses the curvature bound supplied by the analytic domain;
it does not require $\phi'_P\le\varepsilon'$.
\end{enumerate}
\end{proof}

\emph{Step 3: the normaliser; the numbers $n_j$.} For a given real coarse field, we sort the connected sets of cells.

\begin{definition}
Let $V'$ be real, and let $Y$ be a connected set of cells. We call $Y$ \emph{on} if $Y^{+r'}\cap \mathrm{mk}'$ is empty, and \emph{off} otherwise. An off set $Y$ is \emph{short} if it lies in $f$, and \emph{long} otherwise.
\end{definition}

We claim the following.

\begin{enumerate}
\item[\textup{(i)}] $Y$ is on if and only if $1_{Y^\approx}(V')=1$.

\item[\textup{(ii)}] A short $Y$ lies in exactly one $f_j$. The number $n_j(V'):=\sum_{Y\subset f_j,\ {\rm off}}\tilde{\mathsf D}(Y,V')$ is real, Borel and gauge invariant; it does not depend on $e$, and depends on $V'$ only within $3r'$ cells of $\mathrm{mk}'_j$; and $|n_j|\le1.01\mathsf K_1^\flat(2r'+1)^3\#\mathrm{mk}'_j$.

\item[\textup{(iii)}] A long $Y$ has $\#Y>r'$ and a cell within $r'$ of a mark. Moreover, if \[\FitDisplay{\dimexpr\linewidth-1pt\relax}{\hat w(Y):=2\mathsf K_1^\flat e^{-\kappa^\sharp(\#Y-1)+(\kappa+1)(\#Y+s')}}\], then \[\FitDisplay{\dimexpr\linewidth-1pt\relax}{|e^{\tilde{\mathsf D}(Y)}-1|\,\allowbreak e^{(\kappa+1)(\#Y+s')}\le\hat w(Y)}\] and \[\FitDisplay{\dimexpr\linewidth-1pt\relax}{\sum_{Y\ {\rm long,\allowbreak \ within}\ r'\ {\rm of}\ \mathrm{mk}'_j}\hat w(Y)\le\#\mathrm{mk}'_j}\].

\item[\textup{(iv)}] $\sum_Y\tilde{\mathsf D}=\sum_Y1_{Y^\approx}(V')\tilde{\mathsf D}(Y,V')+\sum_jn_j+\sum_{\rm long}\tilde{\mathsf D}$.
\end{enumerate}

We prove \textup{(i)} to \textup{(iv)}.

\begin{enumerate}
\item[\textup{(i)}] A mark is a cell that holds a plaquette of $D'$, and $1_{Y^\approx}(V')=1$ means that no plaquette based in $Y^{+r'}$ is in $D'$; both conditions say that $Y^{+r'}$ contains no mark.

\item[\textup{(ii)}] For $j\ne k$, ${\rm dist}(f_j,f_k)\ge{\rm dist}(\mathrm{mk}'_j,\mathrm{mk}'_k)-2s_0>4r_1+2\ge2$, while consecutive cells of a connected set are at distance $1$; so a short $Y$ lies in a single $f_j$. An off $Y$ has a cell in $\mathrm{mk}'^{+r'}$. If $Y\subset f_j$, the mark within $r'$ of that cell lies within $3r'\le2s_0+4r_1+2$ of $\mathrm{mk}'_j$, so it belongs to the class $j$; and the cell is one of at most $(2r'+1)^3\#\mathrm{mk}'_j$ cells. For a fixed cell $\square$, $\sum_{Y\ni\square}\mathsf K_1^\flat e^{-\kappa^\sharp(\#Y-1)}\le \mathsf K_1^\flat/(1-26e\cdot e^{-100})\le1.01\mathsf K_1^\flat$, and the bound on $|n_j|$ follows. As to the dependence, $\tilde{\mathsf D}(Y,V')$ depends on $V'$ in $Y^{+\lceil R_2/M\rceil}$, where $Y\subset \mathrm{mk}_j'^{+2r'}$ and $\lceil R_2/M\rceil\le r_1\le r'$. That $n_j$ is real and gauge invariant follows from \textup{($\Gamma$4)}. It is Borel, because on each of the finitely many Borel sets $\{\mathrm{mk}'(V')=\mathsf d\}$ it is a finite sum of continuous functions: by \textup{($\Gamma$4)}, $\tilde{\mathsf D}(Y,\cdot)$ is continuous on the $G$-valued fields, the rough ones included.

\item[\textup{(iii)}] A long $Y$ is connected and has both a cell in $\mathrm{mk}'^{+r'}$ and a cell off $\mathrm{mk}'^{+2r'}$; hence $\#Y>r'$. Then $|\tilde{\mathsf D}(Y)|\le \mathsf K_1^\flat e^{-\kappa^\sharp r'}\le1$ by \textup{(A11)}, and $|e^z-1|\le(e-1)|z|\le2|z|$ for $|z|\le1$. For the sum, each long $Y$ within $r'$ of $\mathrm{mk}'_j$ meets one of $(2r'+1)^3\#\mathrm{mk}'_j$ cells. For $n=\#Y>r'$ we have $s'=8r'\le8(n-1)$, so $(\kappa+1)(n+s')\le(\kappa+1)+9(\kappa+1)(n-1)$. Each of the $n-1$ further cells contributes to the exponent the term $\log(26e)-\kappa^\sharp+9(\kappa+1)\le4.2581-10\kappa-40+9\kappa+9=-(\kappa+26.7419)$, and summing gives $\sum_{n-1\ge r'}e^{-(\kappa+26.7419)(n-1)}\le1.001e^{-(\kappa+26.74)r'}$. The sum in \textup{(iii)} is therefore at most $\#\mathrm{mk}'_j(2r'+1)^3\cdot2.01\mathsf K_1^\flat e^{\kappa+1}e^{-(\kappa+26.74)r'}\le\#\mathrm{mk}'_j$, by \textup{(A11)}.

\item[\textup{(iv)}] The on, the short and the long sets form a partition of the connected sets; so \textup{(iv)} follows from \textup{(i)} and \textup{(ii)}.
\end{enumerate}

\emph{Step 4: the alphabet.} We define the letters \textup{(v)}, \textup{(a)}, \textup{(h)}, \textup{(t)}, \textup{(b)} and the zone letters, at a pair $(V',e)$ or at a field of $\mathfrak V$, in three definitions: on letters, on zone letters, and on compatible and actual letters. These definitions also give the meaning of link, core $J$, set $I$, forced and incompatible. In their terms, we make the following four claims.

\begin{enumerate}
\item[\textup{(L)}] Every link of a core has at most $s'$ cells; hence $\#J_\zeta\le\sum_{o\in\zeta}(\#\mathrm{core}(o)+s')$.

\item[\textup{(In)}] $I_\alpha\subset J_\alpha^{+\varrho_*}$ for every letter $\alpha$ other than a forced zone letter; for a forced $\zeta$, $I_\zeta\subset J_\zeta^{+s_0+\varrho_*}$, with $\mathfrak b_\zeta\subset J_\zeta^{+s_0}$.

\item[\textup{(Rd)}] A letter depends on $(V',e)$ in $I_\alpha^{+r_1}$ only; $w_j$ depends on it only within $2r'$ of $\mathrm{mk}'_j$, and $n_j$ only within $3r'$ of $\mathrm{mk}'_j$.

\item[\textup{(Inc)}] Two incompatible letters satisfy ${\rm dist}(I_\alpha,I_\beta)\le2\le r_1$.
\end{enumerate}

\begin{definition}[Letters]
The letters $\lambda$, each with its core $J$ and its set $I$, are the following.

\begin{enumerate}
\item[\textup{(v)}] At a cell $\square$: $\lambda=\chi_\square(e^{V_\square}-1)$, with $J=I=\{\square\}$.

\item[\textup{(a)}] Of $X\in\mathcal D$: $\chi_X(e^{E^a(X,U)}-1)$.

\item[\textup{(h)}] Of $X$: $\chi_Xe^{E^a(X,U)}(e^{E^h(X,U)}-1)$. For both \textup{(a)} and \textup{(h)}, $J=Y_0(X)$ and $I={\rm cl}(X^\approx)\cup Y_0(X)$.

\item[\textup{(t)}] Of a connected set $Y$ of cells: $e^{T_Y}-1$, with $J=I=Y$.

\item[\textup{(b)}] At a cell $\square$: $-b_\square$, with $J=I=\{\square\}$.
\end{enumerate}

The letters \textup{(v)}, \textup{(a)}, \textup{(h)} are used only at fields that are tame on $J_\alpha^\approx$ or lie in $\mathcal U_{J_\alpha}(x')$; there they are defined by \textup{($\Gamma$6)} and \Cref{Y-G-LT}\,(T3). A product with the factor $\chi_\square$ or $\chi_X$ is taken to be $0$ where that factor is $0$; there $U$ may lie outside $\mathcal U_X(x)$ and may fail to be tame, and $V_\square$, $E^a(X,U)$, $E^h(X,U)$ may have no value. In this way every letter is defined at every $\xi$. For a set $\mathfrak B$ of pairwise compatible letters, $K(\mathfrak B)=E_\xi\prod_{\mathfrak B}\lambda_\alpha$ denotes the mean of the product of its letters.
\end{definition}

\begin{definition}[Zone letters]
A \emph{zone letter} $\zeta$ is a set of objects that is connected under adjacency. The \emph{objects} are those of item~E12 of \S\ref{Y-F-1}, that is, a forced class $f_j$, which carries the cells of $f_j$; a cell; and a set $X\in\mathcal P$, of \emph{kind~p}; to these we add the sets $X\in\mathcal D$, of \emph{kind~p$'$}. \emph{Adjacency} is that of item~E12 of \S\ref{Y-F-1}: two objects with cells are adjacent if they have cells at distance $\le4r_1+2$, and an object with a cell $\square$ is adjacent to an $X$ of kind~p if and only if $X$ holds a cube of $\widehat{\square^{+r_1}}$; in addition, an object with a cell $\square$ is adjacent to an $X$ of kind~p$'$ if and only if $X^\approx$ meets $\widehat{\square^{+r_1}}$.

A zone letter is subject to three conditions: its classes $J$ and its cells $t$ form a cell datum, that is, $t$ lies off $f_J$ and $\mathfrak b=f_J\cup t$ is non-empty; each of its polymers is adjacent to an object with cells; and its polymers of kind~p, which make up the family $\mathfrak X$, lie pairwise at distance $>s(x)$. Its polymers of kind~p$'$ make up the family $\mathfrak Y$. The part $(J,t,\mathfrak X)$ alone may fail to be connected. The value of $\zeta$ is \[\lambda_\zeta:=\bigl[\hat\lambda_{(J,t)}1_{\rm ind}\prod_{\mathfrak X}w(X,U)\bigr]\prod_Je^{n_j}\prod_{\mathfrak Y}\bigl(e^{1_{X^\approx}(U)E(X,U)}-1\bigr),\] where $1_{\rm ind}$ is the indicator of item~E12 of \S\ref{Y-F-1}, on $D(U)\cap\hat{\mathfrak Z}$ with $\mathfrak Z=\mathfrak b^{+r_1}$. We put $J(\zeta):=J$, $f_\zeta:=f_{J(\zeta)}=f_J$ and $\mathfrak b_\zeta:=\mathfrak b=f_J\cup t$. For zone letters $\zeta_i$, $\zeta_k$ we set $J_i:=J(\zeta_i)$, the classes of $\zeta_i$, as distinct from the core $J_{\zeta_i}$. We write $t_i$, $\mathfrak b_i$, $\mathfrak X_i$, $\mathfrak Y_i$, $\mathfrak Z_i$, $1_{{\rm ind},i}$ for the other data of $\zeta_i$, and likewise with the index $k$.

A zone letter is \emph{forced} if $J$ is non-empty. Its \emph{core} $J_\zeta$ is made of the marks of the classes of $\zeta$, the cells of $t$ and the sets $Y_0(X)$ for the polymers $X$ of $\zeta$, connected by diagonal paths that follow a spanning tree, chosen by a fixed order. This tree takes its edges from the adjacencies and, within each class, from the pairs of its marks at distance $\le2s_0+4r_1+2$. The classes of item~E9 of \S\ref{Y-F-1} are those of the relation generated by such pairs, so every class is connected under them. We call these diagonal paths, one for each edge of the tree, the \emph{links} of the core. Thus $\mathrm{core}(\square)=\{\square\}$ and $\mathrm{core}(X)=Y_0(X)$, while $\mathrm{core}(f_j)$ consists of the marks of $f_j$, connected in this way, and has at most $(1+s')\#\mathrm{mk}'_j$ cells. The set of $\zeta$ is $I_\zeta:=\mathfrak b^{+2r_1}\cup{\rm cl}(X^{\approx\approx})\ (X\in\mathfrak X)\cup{\rm cl}(X^\approx)\ (X\in\mathfrak Y)\cup J_\zeta$.

The \emph{root} of a letter is the first cell of $J_\alpha$; for a zone letter with non-empty $t$, it is the first cell of $t$. The \emph{universal alphabet} $\mathfrak A$ consists of the letters \textup{(v)}, \textup{(a)}, \textup{(h)}, \textup{(t)}, \textup{(b)} and of the zone letters without a class; these do not depend on $f$. Finally, $d_\alpha:=(\kappa+1)(\#J_\alpha+s')$, and $\alpha\sim\beta$ if and only if ${\rm dist}(I_\alpha,I_\beta)\le r_1$.
\end{definition}

\begin{definition}[Compatible and actual letters]
Two zone letters are \emph{compatible} if every object of the first differs from every object of the second and is not adjacent to it, and if each polymer of kind~p of the first lies at distance $>s(x)$ from each polymer of kind~p of the second. Besides, the following pairs are \emph{incompatible}: a letter \textup{(b)} at a cell of $\mathfrak b_\zeta^{+2r_1}$ and $\zeta$; a letter \textup{(v)} at a cell of $\mathfrak b_\zeta^{+r_1}$ and $\zeta$; a letter \textup{(a)} or \textup{(h)} of an $X$ for which ${\rm cl}(X^\approx)\cap\mathfrak b_\zeta^{+r_1}$ is non-empty, and $\zeta$; and the letters \textup{(a)} and \textup{(h)} of a single $X$. At $V'$, a letter is \emph{actual} in the following cases: it is a forced zone letter $\zeta$, and $f\setminus f_\zeta$ has no cell within $4r_1+2$ of $\mathfrak b_\zeta$; it is a zone letter of $\mathfrak A$, and $f$ has no cell within $4r_1+2$ of $\mathfrak b$; it is a letter \textup{(b)} at a cell that lies off $f^{+2r_1}$; it is a letter \textup{(v)} at a cell that lies off $f^{+r_1}$; it is a letter \textup{(a)} or \textup{(h)} of an $X$ whose set ${\rm cl}(X^\approx)$ lies off $f^{+r_1}$; or it is a letter \textup{(t)}.
\end{definition}

We prove \textup{(L)}, \textup{(In)}, \textup{(Rd)} and \textup{(Inc)}.

\begin{enumerate}
\item[\textup{(L)}] A link bridges at most the following distances: from a cell to a cell, $4r_1+2$; from a cell to a mark of an adjacent class, $s_0+4r_1+2$; between marks of a single class, $2s_0+4r_1+2$. Two classes are never adjacent, by Step~3\,\textup{(ii)}. From a cell $\square$ to $Y_0(X)$ for an $X$ of kind~p the distance is $r_1+1$: the cube that $X$ holds lies in a cell within $r_1$ of $\square$, and ${\rm cl}(X)\subset Y_0^{+1}$. For an $X$ of kind~p$'$, a cube of $X$ lies within $r(x)$ cubes of $\widehat{\square^{+r_1}}$, which gives $r_1+\lceil r(x)/L\rceil+2$. From a mark in place of a cell, $s_0$ is added. By \textup{(A2)}, $4r_1+2\le r'$ and $r_1+\lceil r(x)/L\rceil+2\le r'$; so all these distances are at most $4r'+4r_1+2\le5r'\le s'$. A diagonal path between cells at distance $d$ adds $d-1$ cells.

\item[\textup{(In)}] We have ${\rm cl}(X^\approx)\subset Y_0^{+\lceil r(x)/L\rceil+2}$ and ${\rm cl}(X^{\approx\approx})\subset Y_0^{+\lceil s(x)/L\rceil+2}$; if $\zeta$ is universal, then $\mathfrak b=t\subset J_\zeta$; and $f_j=\mathrm{mk}_j'^{+s_0}$.

\item[\textup{(Rd)}] By Step~0, $\eta_b$ depends on $(V',e)$ only within $R_3$ of $y(b)$, and this holds at every pair, hence also for the forced letters and for the letters \textup{(t)} near marks, where $\phi'_P>\varepsilon'$. For $(G\xi)_b$ this follows from \textup{($\Gamma$2)}, as $r_G\le R-1\le R_3$; for $a_{0,b}$ it follows from the locality of $a_0$, which is the last clause of \textup{($\Gamma$3)} and holds at every pair $(V',e)$ whatever the curvature, rough plaquettes included; and $e$ enters through $X'$ alone. Moreover $R_3\le r_1M$. The letters depend on the field as follows.

\begin{enumerate}
\item[\textup{(v)}] By \textup{($\Gamma$6)\,(a)}, $V_\square$ does not depend on $e$, and depends on $V'$ only within $R_v$ of $\square$; moreover $R_v\le r_1M$.

\item[\textup{(a), (h)}] These letters depend on $U$ on $X^\approx$, so on $V'$ within $R_3$ of ${\rm cl}(X^\approx)$.

\item[\textup{(t)}] Such a letter depends on $\varphi$ and $a_0$ on $\hat Y$ and on $S_k(Y)$. The first paragraph of the proof of \textup{(Rd)} shows that $\varphi_b$ and $a_{0,b}$, at $b\in\hat Y$, depend on the field only within $R_3$ of $y(b)$, and that $R_3\le r_1M$. By \textup{($\Gamma$4)}, $S_k(Y)$ is a function of $V'$ on the bonds of $Y^{+\lceil R_2/M\rceil}$, and $\lceil R_2/M\rceil\le r_1$.
\end{enumerate}

For a zone letter, $[A]^\ell_{\mathfrak Z}$ depends on $\eta$ within $\varrho_L+2$ of $\hat{\mathfrak Z}$, hence on $V'$ within $R_v$ of $\mathfrak Z$, that is, in $\mathfrak b^{+2r_1}$; and $w(X,U)$ and $E(X,U)$ depend on $V'$ within $R_3$ of ${\rm cl}(X^{\approx\approx})$ and of ${\rm cl}(X^\approx)$, respectively. Finally, $w_j$ depends on $V'$ within $r'+\lceil R_6/M\rceil\le2r'$ of $\mathrm{mk}'_j$, by \textup{(A2)}, and the claim on $n_j$ is part of Step~3.

\item[\textup{(Inc)}] We go through the incompatible pairs.

\begin{enumerate}
\item[\textup{(b), (v)}] For a zone letter $\zeta$ and a letter \textup{(b)} or \textup{(v)} that is incompatible with $\zeta$, the cell of this letter lies in $\mathfrak b_\zeta^{+2r_1}$, which is a part of $I_\zeta$; the distance is $0$.

\item[\textup{(a), (h)}] For a letter of $X$ and a zone letter $\zeta$, ${\rm cl}(X^\approx)$ meets $\mathfrak b^{+r_1}$; the distance is $0$. For the two letters of a single $X$ it is $0$ as well.
\end{enumerate}

For two zone letters that share an object the distance is $0$; if they have cells within $4r_1+2$ of each other, their sets $\mathfrak b^{+2r_1}$ are within $2$. If a cell $\square$ of $\zeta$ and a polymer $X$ of $\zeta'$ are adjacent, we take a cube of $\widehat{\square^{+r_1}}$ that $X$ holds when $X$ is of kind~p, and a cube of $X^\approx$ in $\widehat{\square^{+r_1}}$ when $X$ is of kind~p$'$. The cell of this cube lies in $\mathfrak b_\zeta^{+r_1}$, and also in ${\rm cl}(X)$ for kind~p and in ${\rm cl}(X^\approx)$ for kind~p$'$; the distance is $0$. If two polymers of kind~p lie within $s(x)$ of each other, a cube of one lies in $X^{\approx\approx}$ of the other; again the distance is $0$.
\end{enumerate}

\emph{Step 5: independence of letters at a distance.}

\begin{lemma}[Letters whose sets lie far apart are independent]\label{Y-G-LI}
For a set $I$ of cells, let $\mathcal N(I)$ be the set of bonds $b$ of $I^\circ$ for which $y(b)$ lies within the coarse distance $R_\chi+1$ of a site of $I$.

\begin{enumerate}
\item[\textup{(i)}] Every letter $\lambda_\alpha$, at every $(V',e)$, or at every field, at which it is used, is a Borel function of $(\xi_b)$, $b\in\mathcal N(I_\alpha)$.

\item[\textup{(ii)}] If ${\rm dist}(I_\alpha,I_\beta)>r_1$, then $\mathcal N(I_\alpha)$ and $\mathcal N(I_\beta)$ are disjoint.
\end{enumerate}

Consequently, families of letters whose sets $I$ lie pairwise at distance greater than $r_1$ are independent. This is the hypothesis $K(\mathfrak B\cup\mathfrak B')=K(\mathfrak B)K(\mathfrak B')$ at $\mathsf r=r_1$ of \Cref{Y-G-L18}\,(4), which is stated in Step~8.
\end{lemma}

\begin{proof}[Proof of \Cref{Y-G-LI}]
\begin{enumerate}
\item[\textup{(i)}] The field $U_b$ depends on $\eta_b$ and on $V'$. Here $\eta_b=(G\xi)_b-a_{0,b}$, where $a_0$ does not depend on $\xi$, and $(G\xi)_b$ depends on $\xi$ only within $r_w$ of $b$, by \textup{($\Gamma$2)}. The letters depend on $\xi$ as follows.

\begin{enumerate}
\item[\textup{(b)}] Such a letter at a cell $\square$ depends on the $\xi$ at the bonds of $I^\circ$ in $\hat\square$.

\item[\textup{(v)}] The factor $\chi_\square$ depends on $\xi$ within $R_\chi$ of $\hat\square$, and $V_\square$ within $r_w+\varrho_L+1\le R_\chi$, by \Cref{Y-E-C6}\,(a).

\item[\textup{(a), (h)}] The factor $\chi_X$ depends on $\xi$ within $R$ of the blocks that meet $X^\approx$, and these blocks belong to cells of ${\rm cl}(X^\approx)$. By \textup{(ca)} and \textup{(ch)}, $E^a(X,\cdot)$ and $E^h(X,\cdot)$ are functions of $U$ on the bonds of $X^\approx$, hence of $\xi$ within $r_w+1$.

\item[\textup{(t)}] Since $S_k(Y)$ is carried by $\hat Y\times\hat Y$, only $\varphi$ on $\hat Y$ enters, and with it $\xi$ within $r_w$.
\end{enumerate}

For a zone letter, $\hat\lambda1_{\rm ind}$ depends only on the values of $\xi$ at the bonds of $I^\circ$ in $\widehat{\mathfrak b^{+2r_1}}$, by the first sentence of \Cref{Y-F-C7}\,(a); and by \textup{(c$\Xi$)}, $w(X,\cdot)$ is a function of $U$ on the bonds of $X^{\approx\approx}$, hence of $\xi$ within $r_w+1$ of ${\rm cl}(X^{\approx\approx})$. The factors $1_{X^\approx}(U)E(X,U)$ with $X\in\mathfrak Y$ depend on $\xi$ within $r_w+1$ of ${\rm cl}(X^\approx)$, while $e^{-w_j}$ and $e^{n_j}$ do not depend on $\xi$.

All these radii are at most $R_\chi+1$; here a bond of a set is a bond with both ends in it, as in item~1 of \S\ref{Y-G0-1} and in \textup{(ca)}, \textup{(ch)}, \textup{(c$\Xi$)}. The letters are Borel, since the members are Borel and the kernels are continuous.

\item[\textup{(ii)}] Cells at cell distance $\ge r_1+1$ have their sites at coarse distance $\ge r_1M+1>2(R_\chi+1)$, by item~2 of \S\ref{Y-G0-1}; so $\mathcal N(I_\alpha)$ and $\mathcal N(I_\beta)$ are disjoint. As the law of $\xi$ is a product over the bonds of $I^\circ$, the families are independent.
\end{enumerate}
\end{proof}

\begin{remark}
The locality of $E^a$, $E^h$ and $w$ is a clause of the definition of the class, and so a part of the hypothesis of \Cref{Y-G-ThA} (and of $M_A^{\rm su}(x)$).
\end{remark}

\emph{Step 6: the identity \textup{(F2)} and the bijection.} We claim the following.

\begin{enumerate}
\item[\textup{(F2)}] At every $(V',e)$, \[e^{\sum_jn_j}\mathcal Z=\sum_{\mathfrak B}E_\xi\prod_{\alpha\in\mathfrak B}\lambda_\alpha,\] where $\mathfrak B$ ranges over all sets of letters that are pairwise compatible and hold every forced class.

\item[\textup{(Act1)}] All letters of such a set $\mathfrak B$ are actual.

\item[\textup{(Act1$'$)}] Where an actual letter \textup{(v)}, \textup{(a)} or \textup{(h)} is used, $V'$ is tame on $J_\alpha^\approx$.

\item[\textup{(Act2)}] The cell datum of an actual zone letter has \textup{(H17)}.
\end{enumerate}

We begin with the resolution of the integrand. At every $\xi$, \[1=\sum_{\mathfrak b\supset f}\prod_{\mathfrak b\setminus f}b_\square\prod_{\square\notin\mathfrak b}(1-b_\square),\] the product, over the cells off $f$, of $b_\square+(1-b_\square)$. On the support of the term $\mathfrak b$ the following holds.

\begin{enumerate}
\item[\textup{(1)}] Let $\square$ be a cell off $\mathfrak b^{+r_1}$. The bonds within $R_\chi\le r_1M$ of $\hat\square$ lie in cells within $r_1$ of $\square$, hence off $\mathfrak b$; so $\chi_\square=1$. Moreover ${\rm dist}(\square,\mathrm{mk}')>s_0+r_1>r'$, so $V'$ is tame on $\square^\approx$, and \Cref{Y-E-C6}\,(d) gives $A_\square=e^{V_\square}=1+\lambda^v_\square$, where $\lambda^v_\square$ is the value $\chi_\square(e^{V_\square}-1)$ of the letter \textup{(v)} at $\square$.

\item[\textup{(2)}] $D(U)\subset\widehat{\mathfrak b^{+r_1}}$. Indeed, take a plaquette $p$ whose base site lies in a cell $\square$ off $\mathfrak b^{+r_1}$. Every bond $b'\in I^\circ$ within $r_w$ of a bond of $p$ lies within $r_w+1\le R_\chi$ of $\hat\square$, and hence satisfies $|\xi_{b'}|\le c_\xi\varepsilon$. Each plaquette of $D'$ has its base site in a mark, at coarse distance $>(s_0+r_1)M\ge R_3+1$ from $y(p)$. \Cref{Y-G-LT}\,(T1) now gives $\phi_p\le0.72\varepsilon$, and so $p$ does not belong to $D(U)$.

\item[\textup{(3)}] Let $X\in\mathcal D$ be such that ${\rm cl}(X^\approx)\cap\mathfrak b^{+r_1}$ is empty. Every bond within $R\le r_1M$ of a block meeting $X^\approx$ then lies off $\mathfrak b$, which gives $\chi_X=1$; moreover $1_{X^\approx}(U)=1$ by \textup{(2)}. Hence $e^{1_{X^\approx}(U)E(X,U)}=1+\lambda^a_X+\lambda^h_X$, where $\lambda^a_X=\chi_X(e^{E^a(X,U)}-1)$ and $\lambda^h_X=\chi_Xe^{E^a(X,U)}(e^{E^h(X,U)}-1)$ are the values of the letters \textup{(a)} and \textup{(h)} of $X$.

\item[\textup{(4)}] For the other $X\in\mathcal D$, \[\prod e^{1_{X^\approx}(U)E(X,U)}=\sum_{\mathfrak Y}\prod_{\mathfrak Y}(e^{1_{X^\approx}(U)E(X,U)}-1),\] where $\mathfrak Y$ runs over the subsets of $\{X:{\rm cl}(X^\approx)\cap\mathfrak b^{+r_1}$ non-empty$\}$.

\item[\textup{(5)}] By \textup{(2)}, $D(U)=D(U)\cap\hat{\mathfrak Z}$ with $\mathfrak Z=\mathfrak b^{+r_1}$. Hence $\Xi_w(U)=\sum_{\mathfrak X}\prod_{\mathfrak X}w(X,U)1_{\rm ind}$, the sum running over the families in $\mathcal P$ at pairwise distance $>s(x)$, and $1_{\rm ind}=0$ unless each $X$ holds a cube of $\hat{\mathfrak Z}$.
\end{enumerate}

We split three more factors. First, \[\prod_{\square\notin\mathfrak b}(1-b_\square)=\prod_{\mathfrak b^{+2r_1}\setminus\mathfrak b}(1-b_\square)\sum_{J_b}\prod_{J_b}(-b_\square),\] where $J_b$ runs over the sets of cells off $\mathfrak b^{+2r_1}$. Second, $e^{\sum T_Y}=\prod_Y(1+(e^{T_Y}-1))$. Third, $\prod_\square A_\square=[A]^\ell_{\mathfrak Z}\prod_{\square\ {\rm off}\ \mathfrak Z}A_\square$. Off the support of the term $\mathfrak b$, its first factors vanish. So at every $\xi$ the integrand of $\mathcal Z$ is \begin{multline*}\sum_{\mathfrak b\supset f}\prod_{\mathfrak b\setminus f}b_\square\prod_{\mathfrak b^{+2r_1}\setminus\mathfrak b}(1-b_\square)\,[A]^\ell_{\mathfrak Z}\sum_{\mathfrak X}\prod_{\mathfrak X}w(X,U)1_{\rm ind}\sum_{\mathfrak Y}\prod_{\mathfrak Y}(e^{1_{X^\approx}(U)E(X,U)}-1)\\ \times\sum_{J_b}\prod_{J_b}(-b_\square)\prod_{\square\ {\rm off}\ \mathfrak Z}(1+\lambda^v_\square)\prod_{X\ {\rm of}\ (3)}(1+\lambda^a_X+\lambda^h_X)\prod_Y(1+(e^{T_Y}-1)).\end{multline*} All these sums are finite.

We turn to the bijection. Let $(\mathfrak b,\mathfrak X,\mathfrak Y)$ be given with non-zero indicator. Then the components under adjacency of the collection made of all classes, the cells of $\mathfrak b\setminus f$, the polymers of $\mathfrak X$ and those of $\mathfrak Y$ are zone letters; they are pairwise compatible, and together they hold every class once. They are zone letters because for each $X$ of $\mathfrak X$ there is a cell $\square$ of $\mathfrak b$ such that $X$ holds a cube of $\widehat{\square^{+r_1}}$, for each $X$ of $\mathfrak Y$ there is a cell $\square$ of $\mathfrak b$ such that $X^\approx$ meets $\widehat{\square^{+r_1}}$, and every cell of $\mathfrak b$ belongs to an object.

Conversely, suppose that zone letters $\zeta_i$ are pairwise compatible and that together they hold every class. A cell object lying inside some $f_k$ would be at distance $0$ from the object $f_k$; the cell datum rules this out when $k\in J(\zeta_i)$, and compatibility rules it out otherwise. So the cell objects are distinct cells off $f$; we put $\mathfrak b:=\bigcup\mathfrak b_i\supset f$, and the $\zeta_i$ are the components.

As for the factors, ${\rm dist}(\mathfrak b_i,\mathfrak b_k)>4r_1+2$, so the sets $\mathfrak b_i^{+2r_1}$ are disjoint, and the first three factors split into products. As for the indicator, $\hat{\mathfrak Z}$ is the disjoint union of the $\hat{\mathfrak Z}_i$. A plaquette $p$ of $D(U)\cap\hat{\mathfrak Z}_i$ is covered by an $X$ that holds its cube, and this cube lies in $\widehat{\square^{+r_1}}$ for some $\square\in\mathfrak b_i$; so $X$ is adjacent to an object of $\zeta_i$, and $X\in\mathfrak X_i$. If an $X$ of $\mathfrak X_i$ holds a plaquette of $D(U)\cap\hat{\mathfrak Z}_k$, then it is adjacent to $\zeta_k$, so $k=i$. Hence $1_{\rm ind}=\prod_i1_{{\rm ind},i}$.

The sets $J_b$, the cells off $\mathfrak Z$, the sets $X$ of \textup{(3)}, with at most one of \textup{(a)}, \textup{(h)} to each $X$, and the sets $Y$ are exactly the letters \textup{(b)}, \textup{(v)}, \textup{(a)}, \textup{(h)}, \textup{(t)} that are compatible with all the $\zeta_i$. The factors $e^{-w_j}$ and $e^{n_j}$ are given to the letter of the class $j$.

We turn to the clauses on actual letters.

\begin{enumerate}
\item[\textup{(Act1)}] Let $\zeta$ be a zone letter. If some cell of $f_k$ lies within $4r_1+2$ of a cell of $\mathfrak b_\zeta$, then $f_k$ is adjacent to an object of $\zeta$; hence either $f_k$ belongs to $\zeta$, or $\zeta$ and the letter $\zeta_k$ that holds $f_k$ are incompatible. For the letters \textup{(b)}, \textup{(v)}, \textup{(a)}, \textup{(h)}, note that $\mathfrak b_{\zeta_k}\supset f_k$; the incompatibilities of Step~4 exclude the cells of $f_k^{+2r_1}$, the cells of $f_k^{+r_1}$, and the sets $X$ for which ${\rm cl}(X^\approx)\cap f_k^{+r_1}$ is non-empty.

\item[\textup{(Act1$'$)}] We take the letters in turn.

\begin{enumerate}
\item[\textup{(v)}] ${\rm dist}(\square,\mathrm{mk}')>s_0+r_1>r'$.

\item[\textup{(a), (h)}] In both cases of the definition of $Y_0(X)$ in item~1 of \S\ref{Y-G0-1}, $Y_0(X)\subset{\rm cl}(X^\sim)$, a part of ${\rm cl}(X^\approx)$ (for an $X$ that is not small, the inclusion rests on ${\rm hull}(X)=X$, which holds by that item since $X\in\mathcal D$). Hence ${\rm dist}(Y_0(X),\mathrm{mk}')>s_0+r_1>r'$, and \Cref{Y-G-LT}\,(T3) applies with $Y=Y_0(X)$.
\end{enumerate}

\item[\textup{(Act2)}] Suppose that $P\in D'$ lies in a cell at distance at most $4r_1$ from $\mathfrak b_\zeta$. Such a cell is a mark and so belongs to $f$; actuality places it in $f_\zeta$, and for $\zeta\in\mathfrak A$ no such cell exists. The marks of a class $k$ keep a distance greater than $2s_0+4r_1+2$ from those of each class $j\ne k$; such a mark therefore lies in no $f_j$ with $j\ne k$, and its class belongs to $J(\zeta)$.
\end{enumerate}

\emph{Step 7: the product bound.}

\begin{definition}
For a letter $\alpha$ we put $e_\alpha:=e^\circ_\alpha e^{2x^{-1/8}\#I_\alpha}$, where $e^\circ_\alpha$ is as follows.

\begin{enumerate}
\item[\textup{(v)}] $2K_vp_0^3x^{-1/2}$.

\item[\textup{(a)}] $2c_a(x)e^{-\kappa(\#X-1)}$.

\item[\textup{(h)}] $4c_h(x)e^{-\kappa(\#X-1)}$.

\item[\textup{(t)}] $e_Y^{1/2}=x^{-20}e^{-\kappa^\sharp\#Y/2}$.

\item[\textup{(b)}] $u_b$.
\end{enumerate}

For a zone letter, $e^\circ_\zeta$ is \[\mathrm{Res}_{(J,t)}\,u_t^{\#t}u_1^{1[t\ \text{non-empty}]}\prod_{\mathfrak X}\upsilon e^{-\kappa(|X|_c-1)}\prod_Je^{1.01\mathsf K_1^\flat(2r'+1)^3\#\mathrm{mk}'_j}\prod_{\mathfrak Y}2(c_a+c_h)(x)e^{-\kappa(\#X-1)}.\]
\end{definition}

We claim that $E_\xi\bigl|\prod_{\mathfrak B}\lambda_\alpha\bigr|\le\prod_{\mathfrak B}e_\alpha$ in each of the following four situations.

\begin{enumerate}
\item[\textup{(i)}] For pairwise compatible letters that are actual at $V'$, at every $e$.

\item[\textup{(ii)}] At a field that is real on the bonds of $Y^\approx$ and tame there, for letters of $\mathfrak A$ that are pairwise compatible and satisfy $J_\alpha\subset Y$.

\item[\textup{(iii)}] For letters \textup{(a)}, \textup{(v)}, \textup{(t)} with $J_\alpha\subset Y$, at $V'\in \mathcal U_Y(x')$.

\item[\textup{(iv)}] For letters \textup{(b)}.
\end{enumerate}

We now prove the claim.

\begin{enumerate}
\item[\textup{1.}] We first bound the letters at every $\xi$, before any mean over $\xi$ is taken.

\begin{enumerate}
\item[\textup{(v)}] The bound is \textup{($\Gamma$6)\,(c)}.

\item[\textup{(a)}] \Cref{Y-G-LT}\,(T3) shows that the restriction of $U$ to $X^\approx$ lies in $\mathcal U_X(x)$ where $\chi_X=1$; we use $|e^z-1|\le2|z|$ for $|z|\le1/2$, and \textup{(A4)}.

\item[\textup{(h)}] $U$ is real and tame on $X^\approx$, and $e^{1/2}\cdot2c_h(x)\le4c_h(x)$.
\end{enumerate}

For a polymer of kind~p$'$, where the indicator $1_{X^\approx}(U)$ is $1$, the restriction of $U$ to $X^\approx$ is real and lies in $\mathcal U_X(x)$, by item~1 of Step~1; so $|e^{1_{X^\approx}(U)E(X,U)}-1|\le2(c_a+c_h)(x)e^{-\kappa(\#X-1)}$ at every $\xi$. The factors $e^{n_j}$ are bounded by Step~3. The hypotheses of \textup{($\Gamma$6)} and of \Cref{Y-G-LT}\,(T3) hold in the situation \textup{(i)} by \textup{(Act1$'$)}, and in the situations \textup{(ii)} and \textup{(iii)} because $J_\alpha^\approx\subset Y^\approx$.

\item[\textup{2.}] What is left of a zone letter $\zeta_i$ is $\hat\lambda_{(J_i,t_i)}1_{\rm ind}\prod_{\mathfrak X_i}w$, and in the situation \textup{(i)} the datum $(J_i,t_i)$ has \textup{(H17)} by \textup{(Act2)}. We apply \Cref{Y-F-C7}\,(g*1) to $O_i=\zeta_i$ without $\mathfrak Y_i$: the components $\check\zeta_i^{(m)}$ of $O_i$ under the adjacency of item~E12 of \S\ref{Y-F-1} are pairwise compatible admitted letters, and $\hat\lambda_{(J_i,t_i)}1_{\rm ind}\prod_{\mathfrak X_i}w=\prod_m\lambda^o_{\check\zeta_i^{(m)}}$. Components of different $\zeta_i$ are compatible, because the $\zeta_i$ are; and the sets $\mathfrak b^{(m)+2r_1}$ are disjoint parts of $\mathfrak b_i^{+2r_1}$. Moreover $\prod_me^o_{\check\zeta^{(m)}}\le\mathrm{Res}_{(J,t)}\,u_t^{\#t}u_1^{1[t\ \text{non-empty}]}\prod_{\mathfrak X}\upsilon e^{-\kappa(|X|_c-1)}$, since $\mathrm{Res}$ and $u_t^{\#t}$ factor over the components and $u_1\le1$. The cells of the letters \textup{(b)} of $\mathfrak B$ lie off all the $\mathfrak b_i^{+2r_1}$.

\item[\textup{3.}] We turn to the letters \textup{(t)} at a real $V'$. Transports are isometries and $|X'_P|\le\pi$; moreover $\|\ell\|_{\rm r}\le3K_*\mathsf g_*$ and $\|\ell\|_{\rm c}\le K_*\mathsf g_*$, by the first two bounds of \Cref{Y-A-SetsOfRows}, $\|\ell\|_{\rm r}\le3K\mathsf g(\lambda)$ and $\|\ell\|_{\rm c}\le K\mathsf g(\lambda)$, taken at $(K_*,\lambda_*)$ with $\mathsf g_*=\mathsf g(\lambda_*)$, and by \textup{(SD)}. Hence $|(\mathsf B^VEX')_p|\le \mathsf N_{\mathsf B}\pi$. A bond lies in $4$ plaquettes, so $|\mathsf g^VX'|_\infty\le4\pi \mathsf N_{\mathsf B}$; and as $a_0=GG^{\mathsf T}\mathsf g^VX'$, we get $|a_0|_\infty\le a_\infty$ at every $(V',e)$. A cell carries at most $3v$ bonds; therefore \[|T_Y|\le\frac x2e_Y|\varphi_{\hat Y}|^2+xe_Y|\varphi_{\hat Y}||a_{0,\hat Y}|\le xe_Y\bigl(|\varphi_{\hat Y}|^2+\frac12|a_{0,\hat Y}|^2\bigr)\le xe_Y\bigl(|\varphi_{\hat Y}|^2+n_a\#Y\bigr).\] Put $z_Y=xe_Y^{1/2}(n_a\#Y+|\varphi_{\hat Y}|^2)$; then $|e^{T_Y}-1|\le|T_Y|e^{|T_Y|}\le e_Y^{1/2}e^{2z_Y}$. Further, $\sup_\square\sum_{Y\ni\square}e_Y^{1/2}\le x^{-20}\sum_n(26e)^{n-1}e^{-50n}\le1.01x^{-20}$. Writing $\mathsf Y$ for the union of the sets $Y$, we obtain \[\prod|e^{T_Y}-1|\le\prod e_Y^{1/2}\cdot\exp\{2.02x^{-19}n_a\sum\#Y\}\cdot e^{\frac x2\varsigma|\varphi_{\mathsf Y}|^2},\] where $\varsigma=4.04x^{-20}\le(\gamma/2)x^{-2}$ by \textup{(A10)}.

\item[\textup{4.}] For the situations \textup{(i)} and \textup{(iv)}, we apply \Cref{Y-F-C7}\,(d) at $A=1$ to the letters $\check\zeta_i^{(m)}$, with $J_0=J_b$, the set of the cells of the letters \textup{(b)}, and to $\mathsf Y$ and $\varsigma$. The harmless factor that remains is \[\exp\bigl\{[2.02x^{-19}n_a+38.8v\gamma^{-1}x^{-20}(2r_1+1)^3]\sum\#Y+x^{-1/8}\bigl(\sum\#\mathfrak b_i^{+2r_1}+\#J_b\bigr)\bigr\}\] (here $4.8\cdot2\cdot4.04=38.784\le38.8$), and by \textup{(A10)} it is at most $\prod e^{2x^{-1/8}\#I_\alpha}$, since $I_\zeta\supset\mathfrak b^{+2r_1}$.

\item[\textup{5.}] For the situation \textup{(ii)} we extend the field by $1$: let $\tilde V'=V'$ on the bonds of $Y^\approx$, and $1$ elsewhere. A plaquette of $D'(\tilde V')$ has a site off $Y^{+r'}$, and its base site is within one site of that one; so all marks of $\tilde V'$ are at cell distance $\ge r'$ from $Y$. By \textup{(A2)}, $2\varrho_*+r_1\le r'-2$. A letter of $\mathfrak A$ with $J_\alpha\subset Y$ depends on the field in $I_\alpha^{+r_1}$, which is a part of $Y^{+\varrho_*+r_1}$ and so of $Y^{+(r'-2)}$; hence its value does not depend on the extension, and it does not depend on $e$, since $\phi'_P\le\varepsilon'<\pi$ there. A zone letter has $\mathfrak b\subset J_\zeta\subset Y$, and $r'>4r_1$; so no plaquette of $D'(\tilde V')$ lies in a cell within $4r_1$ of $\mathfrak b$, and \textup{(H17)} holds with $J$ empty. Items~1 to~4 therefore apply at $(\tilde V',e)$, for any $e$.

\item[\textup{6.}] In the situation \textup{(iii)} there is no zone letter. The extended field lies in $\mathfrak V^\circ$ by \textup{(u1)}. Since $|Y\cdot Y'|\le\|Y\|\,\|Y'\|$, item~3 holds with $\|\cdot\|$ in place of $|\cdot|$. At $b\in\hat Y$ we have $\|a_{0,b}\|\le \mathsf a_{\rm c}\varepsilon'$, because the box of radius $R_3\le r_1M$ about $y(b)$ lies in $Y^\approx$. By \textup{(A1)}, $(4.1\mathsf a_{\rm c}\varepsilon')^2\le0.1\varepsilon/0.72\le0.1/28.8$, so $\mathsf a_{\rm c}\varepsilon'\le0.0144\le4\pi\le a_\infty$. The last inequality follows from $a_\infty=4\pi \mathsf N_{\mathsf B}(C_\varphi^\flat)^2$, $\mathsf N_{\mathsf B}\ge1$ and $C_\varphi^\flat\ge\mathsf G_0=(2/\gamma)^{1/2}3L^3\mathsf g(\theta_{\mathbb C})>1$, by item~C.I~1.2 of \S\ref{Y-C-1}. Next, $\|\varphi_b\|^2\le(C_\varphi^\flat)^2\sum_{b'\in Y_b}|\xi_{b'}|^2$, and a bond $b'$ lies in at most $3L^3(R+1)^3$ windows; so $\|\varphi_{\hat{\mathsf Y}}\|^2\le3L^3(R+1)^3(C_\varphi^\flat)^2|\xi_{\mathsf Y^{+r_1}}|^2$. Finally, $E_\xi e^{\frac x2\varsigma'|\xi_b|^2}=(1-\varsigma')^{-3/2}\le e^{3\varsigma'}$ for $\varsigma'\le1/2$, and there are at most $3v(2r_1+1)^3\#\mathsf Y$ bonds; \textup{(A10)} then gives the bound.
\end{enumerate}

\emph{Step 8: sums over connected sets, and the summability condition.} We attach the following weights to cells and to polymers: a cell $\square$ receives $\hat w(\square)=u_te^{(\kappa+4)(1+s')}$, a polymer $X$ of kind~p receives $\hat w_p(X)=\upsilon e^{-\kappa(|X|_c-1)+(\kappa+4)(\#Y_0(X)+s')}$, and a polymer $X$ of kind~p$'$ receives $\hat w_{p'}(X)=4(c_a+c_h)e^{-\kappa(\#X-1)+(\kappa+4)(\#Y_0(X)+s')}$.

In the setting \textup{(A3)}, \textup{(A6)}, we have \begin{equation}\tag{An}\sum_{\zeta\ni o}\prod_{o'\in\zeta}\hat w(o')\le\hat w(o)e^{a_o},\qquad a_o\le\#\mathrm{core}(o)/50,\end{equation} where $\zeta$ ranges over those connected sets of cells and polymers that contain $o$, and where the numbers $a_o$ are the ones defined below. The first inequality holds for $o$ a cell, or a polymer of kind~p, or a polymer of kind~p$'$; the second holds for all objects $o$, the classes $f_j$ included.

Moreover, with $\mathsf r=r_1$, $t=x^{-1/4}$, $\varepsilon_2=x^{-1/8}$ and $d_\beta=(\kappa+1)(\#J_\beta+s')$, the alphabet $\mathfrak A$ satisfies the condition \textup{(S)} of \Cref{Y-G-L18} below.

We first state and prove four lemmas, and then return to the step.

\begin{lemma}[Sums over connected sets]\label{Y-G-An0} Let a countable set carry a symmetric relation, and call two of its elements \emph{adjacent} when they are related. Attach to each element $\beta$ numbers $v_\beta\ge0$ and $a_\beta\ge0$, and suppose that $\sum_{\beta\text{ adjacent to }\alpha}v_\beta e^{a_\beta}\le a_\alpha$ for all $\alpha$. We call a finite set that is connected with respect to the relation an \emph{animal}, and the hypothesis just made the \emph{condition for animals}. Then for each $\alpha$, with $S$ running over the animals that contain $\alpha$, \begin{equation}\tag{An0}\sum_{S\ni\alpha}\prod_{\beta\in S}v_\beta\le v_\alpha e^{a_\alpha};\end{equation} and if only the sets $S$ with at least two members are kept, the sum is at most $v_\alpha(e^{a_\alpha}-1)$.\end{lemma}

\begin{proof} For $n\ge1$ let $Z_n(\alpha)$ be the sum of $\prod_{\beta\in S}v_\beta$ over the connected sets $S\ni\alpha$ having at most $n$ members; in particular $Z_1(\alpha)=v_\alpha$. Let $S\ni\alpha$ be connected with at most $n+1$ members, and remove $\alpha$ from it. What is left splits into $k\ge0$ connected components; they are pairwise distinct, none of them has more than $n$ members, and each contains some $\beta$ adjacent to $\alpha$. The components determine $S$, and since they are pairwise distinct, each choice of $k$ of them occurs at least $k!$ times among the ordered $k$-tuples that the $k$-th power of $\sum_{\beta\text{ adjacent to }\alpha}Z_n(\beta)$ counts; hence \[Z_{n+1}(\alpha)\le v_\alpha\sum_k\frac1{k!}\bigl(\sum_{\beta\text{ adjacent to }\alpha}Z_n(\beta)\bigr)^k\le v_\alpha\exp\{\sum_{\beta\text{ adjacent to }\alpha}v_\beta e^{a_\beta}\}\le v_\alpha e^{a_\alpha},\] where the second inequality uses $Z_n(\beta)\le v_\beta e^{a_\beta}$ and the third is the hypothesis. Since $v_\alpha\le v_\alpha e^{a_\alpha}$, the bound holds for $n=1$, and induction on $n$ yields $Z_n(\alpha)\le v_\alpha e^{a_\alpha}$ for every $n$, and the first claim follows as $n\to\infty$. In this bound the term $k=0$ comes from the set $S=\{\alpha\}$ and equals $v_\alpha$; without it, the same chain bounds the sum over the sets with at least two members by $v_\alpha(e^{a_\alpha}-1)$.\end{proof}

\begin{lemma}[A convergence criterion for sums over polymers, due to Koteck\'y and Preiss]\label{Y-G-LKP} Let $\mathfrak P$ be a finite set, whose elements we call \emph{polymers}, and let a symmetric relation on $\mathfrak P$ be given; related polymers are said to \emph{touch}, and every polymer touches itself. Attach to each $P$ numbers $W(P)\in\mathbb C$ and $a(P)\ge0$, and suppose that $\sum_{P'\text{ touches }P}|W(P')|e^{a(P')}\le a(P)$ for all $P$. A \emph{cluster} is a finite sequence $C=(P_1,\dots,P_n)$ of polymers, with $n\ge1$, for which the graph on $\{1,\dots,n\}$ that joins $i\ne j$ whenever $P_i$ touches $P_j$ is connected. For a cluster we set $\varphi^T(C):=\frac1{n!}\sum_g(-1)^{\#g}$, where $g$ ranges over the connected spanning subgraphs of that graph and $\#g$ denotes the number of edges of $g$; we write $C\ni P$ when $P$ is one of the $P_i$.

Then $\sum_{C\ni P}|\varphi^T(C)|\prod_i|W(P_i)|\le|W(P)|e^{a(P)}$ for every polymer $P$, and the clusters with $n\ge2$ contribute at most $|W(P)|(e^{a(P)}-1)$ to this sum. Moreover, $\sum_{\{P_i\}}\prod_iW(P_i)=\exp\sum_C\varphi^T(C)\prod_iW(P_i)$, where on the left $\{P_i\}$ runs over the sets of polymers no two of which touch, and the empty set contributes $1$.\end{lemma}

\begin{proof}
Fix a total order on the edges of the cluster graph. For a spanning
 tree $\tau$, let $R(\tau)$ consist of $\tau$ and every edge whose order
 is greater than that of each edge of the path in $\tau$ joining its
 endpoints. Run the greedy tree algorithm in increasing order on a
 connected spanning subgraph $g$. If its output is $\tau$, then every
 edge of $g\setminus\tau$ was rejected after its endpoints were already
 joined by earlier edges of $\tau$. Hence $g\subset R(\tau)$. Conversely,
 if $\tau\subset g\subset R(\tau)$, induction over the edges shows that
 each edge of $\tau$ is accepted and each edge of $g\setminus\tau$ is
 rejected. The subgraphs giving $\tau$ therefore form the interval
 $[\tau,R(\tau)]$. The alternating sum on an interval is $0$ unless
 $R(\tau)=\tau$, when it is $(-1)^{n-1}$. This gives the Penrose
 tree bound
 \[
   \biggl|\sum_g(-1)^{\#g}\biggr|\le\#\{\text{spanning trees}\}.
 \]
 A cluster containing $P$ can be rooted at an occurrence of $P$.
 Counting all such occurrences only increases the sum. Permutation
 invariance then puts the root at the first place at a cost of a factor
 $n$. Thus the first sum of the lemma is at most $|W(P)|T(P)$, where
 \[
 T(P)=\sum_{n\ge1}\frac1{(n-1)!}\sum_{P_2,\ldots,P_n}
         \sum_{\tau}\prod_{i=2}^n|W(P_i)|,
 \]
 with $P_1=P$ and the tree edges required to join touching polymers.

For $m\ge0$ let $T_m(P)$ be the part of $T(P)$ that comes from the trees whose depth, measured from the root $1$, is at most $m$; thus $T_0=1$. Sorting the trees according to the $k$ subtrees that hang at the root, we find \[T_{m+1}(P)=\exp\{\sum_{P'\text{ touches }P}|W(P')|T_m(P')\}\le\exp\{\sum|W(P')|e^{a(P')}\}\le e^{a(P)},\] where the first inequality uses $T_m(P')\le e^{a(P')}$ and the second is the hypothesis. Induction on $m$ gives $T_m(P)\le e^{a(P)}$ for every $m$, hence $T(P)\le e^{a(P)}$, which proves the first claim. In $T(P)$ the term $n=1$ is the summand $1$; it comes from the cluster $(P)$, which contributes $|W(P)|$ to the sum of the first claim, and omitting it on both sides proves the second claim.

For the identity introduce a variable $z$ and replace every activity
 by $zW(P)$. The coefficient of $z^n$ in the partition function is
 \[
 \frac1{n!}\sum_{P_1,\ldots,P_n}\prod_i W(P_i)
    \prod_{i<j}\bigl(1-1\{P_i\text{ touches }P_j\}\bigr).
 \]
 Self-touching excludes repeated polymers. Expand the product at each
 fixed $n$ and partition the graph into connected components. A component
 contributes its connected graph sum. The factor $1/n!$ accounts for the
 allocation of labels to the components and their permutations, so the
 resulting identity of formal power series is the exponential of the
 connected-cluster series. The rooted estimate above, summed over the
 finite set of possible root polymers, gives absolute convergence of
 that series for $|z|\le1$. The formal identity thus gives equality for
 $|z|<1$. Let $z$ tend to $1$ along the real interval and use absolute
 convergence to obtain the claimed identity.
\end{proof}

\begin{lemma}[Clusters of letters under a summability condition]\label{Y-G-L18} Let a finite alphabet be given, and attach to each letter $\alpha$ a finite non-empty set $I_\alpha$ of cells and two numbers $e_\alpha\ge0$, $d_\alpha\ge0$. For a number $\mathsf r\ge0$ we write $\alpha\sim\beta$ if and only if ${\rm dist}(I_\alpha,I_\beta)\le\mathsf r$. Let \emph{compatibility} be a symmetric relation between different letters, such that two letters that are not compatible are always $\sim$. Let $0<t\le1$ and $\varepsilon_2>0$, and suppose that at every cell $\square$ \begin{equation}\tag{S}\sum_{\beta:{\rm dist}(I_\beta,\square)\le\mathsf r}e_\beta e^{3\varepsilon_2\#I_\beta+d_\beta}\le t\varepsilon_2.\end{equation}

A \emph{polymer} is a non-empty set of pairwise compatible letters that is $\sim$-connected, and two polymers \emph{touch} if some letter of the one is $\sim$ to some letter of the other. For this relation, the clusters $C=(\mathfrak B_1,\dots,\mathfrak B_n)$ of polymers and the numbers $\varphi^T(C)$ are those of \Cref{Y-G-LKP}, and we set $\hat\Phi(C):=|\varphi^T(C)|\prod_i\prod_{\alpha\in\mathfrak B_i}e_\alpha$. We denote by $n(C)$ the number of the letters of $C$ and by $d(C)$ the sum of their numbers $d_\alpha$, both counted with multiplicity. We write $C\ni\alpha$ if $\alpha$ belongs to some $\mathfrak B_i$, and $C\sim\alpha$ if some letter of $C$ is $\sim\alpha$; in the relation $C\sim\alpha$ we allow $\alpha$ to be any object that carries a finite non-empty set $I_\alpha$ of cells. Then the following hold.

\begin{enumerate}\item[\textup{(1)}] For every letter $\alpha$, $\sum_{C\ni\alpha}\hat\Phi(C)e^{d(C)}\le e_\alpha e^{2\varepsilon_2\#I_\alpha+d_\alpha}$.

\item[\textup{(2)}] The clusters with $n(C)\ge2$ contribute to the sum in \textup{(1)} at most $t$ times the right side of \textup{(1)}.

\item[\textup{(3)}] For every such object $\alpha$, $\sum_{C\sim\alpha}\hat\Phi(C)e^{d(C)}\le t\varepsilon_2\#I_\alpha$; and $\hat\Phi(C)e^{d(C)}\le t\varepsilon_2$ for every cluster $C$.

\item[\textup{(4)}] Suppose in addition that a number $K(\mathfrak B)\in\mathbb C$ is given for every set $\mathfrak B$ of pairwise compatible letters, with $K(\emptyset)=1$, with $|K(\mathfrak B)|\le\prod_{\mathfrak B}e_\alpha$, and with $K(\mathfrak B\cup\mathfrak B')=K(\mathfrak B)K(\mathfrak B')$ whenever no letter of $\mathfrak B$ is $\sim$ to a letter of $\mathfrak B'$; and put $\Phi(C):=\varphi^T(C)\prod_iK(\mathfrak B_i)$. Then $|\Phi|\le\hat\Phi$, and $\sum_{\mathfrak B}K(\mathfrak B)=\exp\sum_C\Phi(C)$, where $\mathfrak B$ runs over the sets just named and the series converges absolutely.\end{enumerate}

If an alphabet satisfies \textup{(S)}, so does each of its subsets, and the claims therefore hold for each subset. Claims \textup{(1)} to \textup{(3)} do not involve $K$, and so they do not depend on any field.\end{lemma}

\begin{proof} Put $s_\alpha=\varepsilon_2\#I_\alpha$ and $v_\alpha=e_\alpha e^{d_\alpha}$. If $\beta\sim\alpha$, then ${\rm dist}(I_\beta,\square)\le\mathsf r$ for some cell $\square$ of $I_\alpha$; adding up \textup{(S)} over the cells of $I_\alpha$ we therefore obtain \begin{equation}\tag{S$'$}\sum_{\beta\sim\alpha}v_\beta e^{3s_\beta}\le ts_\alpha\end{equation} for every object $\alpha$ that carries a set $I_\alpha$.

We apply \Cref{Y-G-An0} to the letters, with $\sim$ as the adjacency, with the weights $v_\beta e^{ts_\beta}$ and with $a_\beta=ts_\beta$. Its hypothesis becomes $\sum_{\beta\sim\alpha}v_\beta e^{2ts_\beta}\le ts_\alpha$, which follows from \textup{(S$'$)} since $t\le1$. The polymers that hold $\alpha$ are among the connected sets that hold $\alpha$, so the lemma gives \begin{equation}\tag{$\star$}\sum_{\mathfrak B\ni\alpha}\prod_{\beta\in\mathfrak B}v_\beta e^{ts_\beta}\le v_\alpha e^{2ts_\alpha},\end{equation} the sum running over the polymers $\mathfrak B$ that hold $\alpha$.

Next we apply \Cref{Y-G-LKP} to the polymers, which form a finite set, with $W(\mathfrak B)=\prod_{\mathfrak B}v_\beta$ and $a(\mathfrak B)=t\sum_{\mathfrak B}s_\beta$. If $\mathfrak B'$ touches $\mathfrak B$, then some letter $\beta$ of $\mathfrak B'$ is $\sim$ to some letter $\alpha$ of $\mathfrak B$; hence, by \textup{($\star$)} and \textup{(S$'$)}, \[\sum_{\mathfrak B'\text{ touches }\mathfrak B}W(\mathfrak B')e^{a(\mathfrak B')}\le\sum_{\alpha\in\mathfrak B}\sum_{\beta\sim\alpha}\sum_{\mathfrak B'\ni\beta}\prod_{\mathfrak B'}ve^{ts}\le\sum_{\alpha\in\mathfrak B}\sum_{\beta\sim\alpha}v_\beta e^{2ts_\beta}\le a(\mathfrak B),\] and this is the hypothesis of \Cref{Y-G-LKP}. For every cluster, $\hat\Phi(C)e^{d(C)}=|\varphi^T(C)|\prod_iW(\mathfrak B_i)$. We now prove the claims in turn.

\begin{enumerate}\item[\textup{(1)}] With the summand $\hat\Phi(C)e^{d(C)}$ in the first two sums, \Cref{Y-G-LKP} and \textup{($\star$)} give \[\sum_{C\ni\alpha}\le\sum_{\mathfrak B\ni\alpha}\sum_{C\ni\mathfrak B}\le\sum_{\mathfrak B\ni\alpha}W(\mathfrak B)e^{a(\mathfrak B)}\le v_\alpha e^{2ts_\alpha}\le v_\alpha e^{2s_\alpha},\] and $v_\alpha e^{2s_\alpha}=e_\alpha e^{2\varepsilon_2\#I_\alpha+d_\alpha}$.

\item[\textup{(2)}] Exactly one cluster with $n(C)=1$ holds $\alpha$, namely $(\{\alpha\})$, and its term is $v_\alpha$. By the chain in \textup{(1)}, the other clusters that hold $\alpha$ contribute at most $v_\alpha(e^{2ts_\alpha}-1)\le tv_\alpha(e^{2s_\alpha}-1)$; the last inequality holds because $t\mapsto e^{2ts_\alpha}$ is convex and takes the value $1$ at $t=0$. So these clusters contribute at most $t$ times the bound in \textup{(1)}.

\item[\textup{(3)}] A cluster $C\sim\alpha$ holds a letter $\beta\sim\alpha$; hence, by \textup{(1)} and \textup{(S$'$)}, \[\sum_{C\sim\alpha}\le\sum_{\beta\sim\alpha}\sum_{C\ni\beta}\le\sum_{\beta\sim\alpha}v_\beta e^{2s_\beta}\le ts_\alpha,\] which is the first half, as $s_\alpha=\varepsilon_2\#I_\alpha$. For the second half, choose a letter $\beta$ of $C$ and a cell $\square$ of $I_\beta$; then \textup{(1)} and \textup{(S)} at $\square$ give $\hat\Phi(C)e^{d(C)}\le v_\beta e^{2s_\beta}\le t\varepsilon_2$.

\item[\textup{(4)}] A set $\mathfrak B$ of pairwise compatible letters can be written in exactly one way as a union of polymers no two of which touch, namely as the union of its $\sim$-components, and $K(\mathfrak B)$ is the product of the values of $K$ on these components. Conversely, a union of polymers no two of which touch consists of pairwise compatible letters, since incompatible letters are $\sim$. Thus $\sum_{\mathfrak B}K(\mathfrak B)$ is the left side of the identity of \Cref{Y-G-LKP} for the weight $W=K$ on the polymers. Moreover $|K(\mathfrak B)|\le\prod_{\mathfrak B}e_\alpha\le\prod_{\mathfrak B}v_\alpha$; so the hypothesis of \Cref{Y-G-LKP}, established above for the weight $\prod_{\mathfrak B}v_\beta$, holds for $W=K$ as well, the identity of that lemma becomes the identity in \textup{(4)}, with absolute convergence, and the same bound on $K$ yields $|\Phi|\le\hat\Phi$.\end{enumerate}\end{proof}

\begin{lemma}[Sums over letters and clusters attached to given letters]\label{Y-G-L19} Let the data of \Cref{Y-G-L18} be given and satisfy \textup{(S)}, and let $t\le1/2$. Let $\mathsf F$ be a finite set of additional letters $\zeta$; to each of them a finite non-empty set $I_\zeta$ of cells is attached, and $\sim$ is defined for them by the same rule. Let $\mathfrak M$ be a countable set of objects $C$, each of which has a relation $C\sim\alpha$ to the letters, and let numbers $\hat w_C\ge0$ be given such that $\sum_{C\in\mathfrak M,\,C\sim\alpha}\hat w_C\le\varepsilon_2\#I_\alpha$ for every letter $\alpha$, both of the alphabet and of $\mathsf F$. We consider the pairs $(\mathfrak B_o,\mathfrak N)$ consisting of a set $\mathfrak B_o$ of letters of the alphabet, every $\sim$-component of which holds a letter that is $\sim$ to a letter of $\mathsf F$, and of a finite subset $\mathfrak N$ of $\mathfrak M$, every member of which is $\sim$ to some letter of $\mathsf F\cup\mathfrak B_o$. Then summing $\prod_{\beta\in\mathfrak B_o}e_\beta e^{d_\beta}\prod_{C\in\mathfrak N}\hat w_C$ over all these pairs gives at most $\prod_{\zeta\in\mathsf F}e^{2\varepsilon_2\#I_\zeta}$. When $\mathfrak M$ is the set of the clusters and $\hat w_C=2\hat\Phi(C)e^{d(C)}$, the condition on $\mathfrak M$ is satisfied, by \Cref{Y-G-L18}\,(3) and $2t\le1$.\end{lemma}

\begin{proof} We write again $s_\alpha=\varepsilon_2\#I_\alpha$, now also for the letters of $\mathsf F$, and $v_\alpha=e_\alpha e^{d_\alpha}$. Fix $\mathfrak B_o$. The sets $\mathfrak N$ for which $(\mathfrak B_o,\mathfrak N)$ is one of the pairs considered in the lemma are the finite sets of members $C$ of $\mathfrak M$ with $C\sim\alpha$ for some $\alpha\in\mathsf F\cup\mathfrak B_o$; therefore \[\sum_{\mathfrak N}\prod\hat w_C\le\prod_{C\sim\mathsf F\cup\mathfrak B_o}(1+\hat w_C)\le\exp\sum_{\alpha\in\mathsf F\cup\mathfrak B_o}\sum_{C\sim\alpha}\hat w_C\le\prod_{\alpha\in\mathsf F\cup\mathfrak B_o}e^{s_\alpha}.\] It remains to sum over $\mathfrak B_o$. Such a set is the union of its $\sim$-components, which are pairwise distinct $\sim$-connected sets $S$, and each of them holds a letter $\beta\sim\zeta$ for some $\zeta$ of $\mathsf F$. A choice of $k$ such sets occurs at least $k!$ times among the ordered choices, so summing over $k$ we get \[\sum_{\mathfrak B_o}\prod v_\beta e^{s_\beta}\le\exp\{\sum_{\zeta\in\mathsf F}\sum_{\beta\sim\zeta}\sum_{S\ni\beta}\prod_Sve^{s}\}\le\exp\{\sum_\zeta\sum_{\beta\sim\zeta}v_\beta e^{2s_\beta}\}\le\exp\{t\sum_\zeta s_\zeta\}.\] The second inequality is \Cref{Y-G-An0} for $\sim$ with the weights $v_\beta e^{s_\beta}$ and $a_\beta=s_\beta$, whose hypothesis $\sum_{\beta\sim\alpha}v_\beta e^{2s_\beta}\le ts_\alpha\le s_\alpha$ follows from \textup{(S$'$)}; the third inequality is \textup{(S$'$)} at $\zeta$. Combining the two bounds, each $\zeta\in\mathsf F$ receives the factor $e^{(1+t)s_\zeta}\le e^{2s_\zeta}$, and this is the claim.\end{proof}

We now come to the step itself. The application of \Cref{Y-G-An0} to the objects involves the number $\varepsilon''=x^{-1/4}$ that item~4 of \S\ref{Y-G0-1} sets. It enters at two places: in the numbers $a_o:=\varepsilon''\#\operatorname{fp}(o)^{+(4r_1+2)}\le\#\mathrm{core}(o)$, where $\operatorname{fp}(o)$ is the footprint of $o$ defined below, and in the bound $\sum_{\operatorname{fp}(o)\ni\square}\hat w(o)e^{a_o}\le\varepsilon''$ at every cell $\square$.

The \emph{footprints} of the objects are $\operatorname{fp}(\square)=\{\square\}$ for a cell, $\operatorname{fp}(f_j)=f_j$, $\operatorname{fp}_p(X)={\rm cl}(X)^{+r_1}$ for a polymer of kind~p, and $\operatorname{fp}_{p'}(X)={\rm cl}(X^\approx)^{+r_1}$ for a polymer of kind~p$'$. For adjacent objects $o$ and $o'$ we have $\operatorname{fp}(o')\cap\operatorname{fp}(o)^{+(4r_1+2)}\ne\emptyset$: if each of the two is a cell or a class $f_j$, then by the definition of adjacency some cell of the one lies at distance at most $4r_1+2$ from some cell of the other, and the footprint of each is its set of cells; if an object has a cell $\square$ through which it is adjacent to a polymer $X$, then $\square$ lies in the footprint of that object and in $\operatorname{fp}(X)$, whether $X$ is of kind~p or of kind~p$'$; and two polymers are never adjacent.

We bound $a_o$. The set $\operatorname{fp}(o)^{+(4r_1+2)}$ is contained in $\square^{+(4r_1+2)}$ when $o=\square$ is a cell, in $Y_0^{+(5r_1+\lceil r(x)/L\rceil+4)}$ with $Y_0=Y_0(X)$ when $o=X$ is a polymer, and in $\mathrm{mk}_j'^{+(s_0+4r_1+2)}$ when $o=f_j$. Each of these radii is at most $s_0+\varrho_*+3r_1+2$, and so \textup{(A3)}, together with $x^{-1/4}\le x^{-1/8}$, gives $a_o\le\#\mathrm{core}(o)/50$.

Now we fix a cell $\square$ and estimate $\sum_{\operatorname{fp}(o)\ni\square}\hat w(o)e^{a_o}$, where $o$ ranges over the cells and over the polymers of kind~p and of kind~p$'$. The cell $\square$ itself contributes at most $u_te^{\kappa+5}e^{(\kappa+4)s'}$.

Consider first the polymers $X$ of kind~p. If $\square$ lies in $\operatorname{fp}_p(X)$, then $\square\in{\rm cl}(X)^{+r_1}$, and so, by the definition of ${\rm cl}$, the fine set of some cell of $\square^{+r_1}$ meets $X$. As the grids fit (item~1 of \S\ref{Y-G0-1}), the fine set $\hat\square$ of each cell is made up of exactly $L^3$ fine cubes. Hence $X$ holds one of the $L^3(2r_1+1)^3$ fine cubes that make up the fine sets of the $(2r_1+1)^3$ cells of $\square^{+r_1}$. Up to the factor $\upsilon e^{(\kappa+4)s'}$ of their weights, the polymers of kind~p holding a fixed cube contribute at most $\sum_n(52e)^ne^{-\kappa(n-1)}e^{(\kappa+5)\#Y_0}$, where $n=|X|_c$. For the small $X$, which have $\#Y_0=1$, this is at most $e^{\kappa+5}\cdot52e/(1-52e\cdot e^{-6})\le217.59\,e^{\kappa+5}$. For $X$ not small we have $n\ge(L-3)/2\ge6\cdot10^{12}$ (for the second inequality, \textup{(L2$^{\rm su}$)} and $\mathsf N_{\rm c}\ge17$ give $L\ge(1.2\cdot10^4\cdot17^2)^2\ge1.2027\cdot10^{13}$) and $\#Y_0\le8+432n/L$, and the exponent of the summand of index $n$ is at most $9\kappa+40-n[\kappa-4.952-432(\kappa+5)/L]$. For $\kappa\ge6$ one has $(\kappa+5)/(\kappa-4.952)\le10.5$, and together with $L\ge1.2\cdot10^{13}$ this gives $432(\kappa+5)/L\le10^{-9}(\kappa-4.952)$; hence the bracket is at least $1.04$ and at least $0.174\kappa$. The ratio of two consecutive terms is therefore at most $e^{-1.04}$, and the sum is at most $1.55$ times its first term, whose exponent is at most $\kappa(15.7-0.174\cdot6\cdot10^{12})$ and thus below $-10^{12}$. Altogether, up to the same factor, the polymers of kind~p holding a fixed cube contribute at most $220\,e^{\kappa+5}$.

Consider next the polymers $X$ of kind~p$'$. If $\square$ lies in $\operatorname{fp}_{p'}(X)$, then $X$ has a cube in a set of $((2r_1+1)L+2r(x))^3$ cubes that is determined by $\square$ and does not depend on $X$. For a fixed cube, $\sum_{X\ni\text{cube}}e^{-\kappa(\#X-1)}\le(1-26e\cdot e^{-6})^{-1}\le1.2124$; hence, up to the factor $(c_a+c_h)e^{\kappa+5}e^{(\kappa+4)s'}$, the small polymers of kind~p$'$ that hold the cube contribute at most $4\cdot1.2124\le4.85$, and together with those that are not small at most $6$. For the latter, $n=\#X\ge(L-3)/2$ and $\#Y_0\le8+432n/L$, and, relative to $e^{\kappa+5}$, the exponent of the summand of index $n$ is at most $4.2581(n-1)+35+2160n/L-\kappa[n(1-432/L)-8]$. The bracket is positive, so this bound decreases as $\kappa$ increases; at $\kappa=6$ it is at most $-1.74n+79$.

Adding the contribution of the cell and those of the two kinds of polymers, we obtain the left side of \textup{(A6)}; hence $\sum_{\operatorname{fp}(o)\ni\square}\hat w(o)e^{a_o}\le\varepsilon''/2$ at every cell $\square$.

We apply \Cref{Y-G-An0}, by which $\sum_{\beta\text{ adjacent to }\alpha}v_\beta e^{a_\beta}\le a_\alpha$ implies $\sum_{\text{connected sets}\ni\alpha}\prod v\le v_\alpha e^{a_\alpha}$. Every object $o'$ adjacent to $o$ has a footprint that meets $\operatorname{fp}(o)^{+(4r_1+2)}$; applying the last bound at each cell of $\operatorname{fp}(o)^{+(4r_1+2)}$, we obtain $\sum_{o'\text{ adjacent to }o}\hat w(o')e^{a_{o'}}\le\varepsilon''\#\operatorname{fp}(o)^{+(4r_1+2)}=a_o$. So the lemma applies to the set of all cells and of all polymers, of kind~p as well as of kind~p$'$, with the adjacency of objects, with the weights $\hat w$ attached to them at the beginning of the step and with the numbers $a_o$, and its conclusion is \textup{(An)}.

It remains to prove \textup{(S)}. Fix a cell $\square$ and a letter $\beta\in\mathfrak A$ with ${\rm dist}(I_\beta,\square)\le r_1$. By \textup{(In)}, some cell $\square''$ of $J_\beta$ lies within $\varrho_*+r_1\le s'$ of $\square$, and there are $(2s'+1)^3$ such cells $\square''$. Since $J_\beta$ is connected and holds its root $\square_0$, we have $\#J_\beta>{\rm dist}(\square_0,\square'')$, and so, for non-negative summands, \[\sum_{\beta:J_\beta\ni\square''}(\cdot)\le\sum_{\square_0}e^{-{\rm dist}(\square_0,\square'')}\sum_{\beta\text{ rooted at }\square_0}e^{\#J_\beta}(\cdot),\] where the sum of $e^{-{\rm dist}(\square_0,\square'')}$ over $\square_0$ is at most $\sum_{n\in\mathbb Z^3}e^{-|n|_\infty}\le49.98\le50$. Further, $e_\beta e^{3\varepsilon_2\#I_\beta}=e^\circ_\beta e^{5x^{-1/8}\#I_\beta}\le e^\circ_\beta e^{\#J_\beta}$, by $\#I_\beta\le(2\varrho_*+1)^3\#J_\beta$ and \textup{(A3)}, and $e^{d_\beta}=e^{(\kappa+1)s'}e^{(\kappa+1)\#J_\beta}$. What has to be bounded for each root $\square_0$ is thus the sum of $e^\circ_\beta e^{(\kappa+3)\#J_\beta}$ over the letters $\beta$ rooted at $\square_0$; we treat the kinds of letters in turn.

\begin{enumerate}\item[\textup{(a), (h)}] For $X$ small, these letters give at most $L^3\cdot1.2124\cdot(2c_a+4c_h)e^{\kappa+3}\le4.85L^3(c_a+c_h)e^{\kappa+3}$. For $X$ not small, $X$ has a cube in a set of $(L+2)^3$ cubes that is determined by the root $\square_0$ and does not depend on $X$. Moreover, $n\ge(L-3)/2\ge6\cdot10^{12}$ and $\#Y_0\le8+432n/L$; relative to $e^{\kappa+3}$, the exponent of the summand of index $n$ is at most $4.2581(n-1)+21+1296n/L-\kappa[n(1-432/L)-8]$. The bracket is positive, so this bound decreases as $\kappa$ increases. By \textup{(O5)} we have $\kappa\ge6$. For every such $\kappa$ the bound is therefore at most its value at $\kappa=6$, and this value is at most $-1.74n+65$. Altogether the letters \textup{(a)} and \textup{(h)} give at most $5e^{\kappa+3}L^3(c_a+c_h)$.

\item[\textup{(v)}] These letters give at most $2e^{\kappa+3}K_vp_0^3x^{-1/2}$.

\item[\textup{(t)}] These letters give at most $x^{-20}\sum_n(26e)^{n-1}e^{-5(\kappa+4)n+(\kappa+3)n}\le x^{-20}$.

\item[\textup{(b)}] These letters give at most $e^{\kappa+3}u_b$.\end{enumerate}

For a zone letter $\zeta$, \textup{(L)} gives $e^\circ_\zeta e^{(\kappa+3)\#J_\zeta}\le u_1\prod_o\hat w(o)$, the product running over the objects $o$ of $\zeta$. By \textup{(An)}, the sum over the zone letters $\zeta\ni\square_0$ is therefore at most $u_1\hat w(\square_0)e^{a_{\square_0}}$, and by the bound at the cell $\square_0$ this is at most $u_1\varepsilon''/2\le u_1$. With $t^{-1}=x^{1/4}$, the inequality to which these bounds reduce \textup{(S)} is \textup{(A5)}; hence \textup{(S)} holds.

\emph{Step 9: cores of clusters; the new members $E'^{\rm old}$, $E'^{\rm new}$, $E'^h$ and the number $\mathsf V'$, as defined in \textup{(F3)}; those clauses of \textup{(ca)} and \textup{(ch)} that involve no size.} The polymers and clusters of the alphabet $\mathfrak A$, and the numbers $\varphi^T(C)$, are those of \Cref{Y-G-L18}. For a set $\mathfrak B$ of pairwise compatible letters we put $K(\mathfrak B)=E_\xi\prod_{\mathfrak B}\lambda_\alpha$, and for a cluster $C$ we put $\Phi(C)=\varphi^T(C)\prod_iK(\mathfrak B_i)$ and $\hat\Phi(C)=|\varphi^T(C)|\prod e_\alpha$; moreover, $n(C)$ is the number of the letters of $C$, and $d(C)=\sum d_\alpha$.

The \emph{core} $\operatorname{co}(C)$ of a cluster $C$ is obtained from the sets $J_\alpha$ of its letters by joining them with diagonal paths along a spanning tree of $\sim$, the choices being fixed by an order. We call a cluster \emph{constant} when each of its letters is \textup{(b)}, \emph{analytic} when each of its letters is \textup{(a)}, \textup{(v)} or \textup{(t)}, and \emph{hard} otherwise. Every cluster $C$ satisfies \begin{equation}\tag{Clu}e^{\kappa(\#\operatorname{co}(C)-1)}\le e^{d(C)-(\kappa+1)s'-\#\operatorname{co}(C)-\kappa}.\end{equation}

We define \begin{equation}\tag{F3}\begin{gathered}E'^{\rm old}(Y,\cdot):=\sum_{\operatorname{co}(C)=Y,\text{ analytic, with a letter (a)}}[\Phi(C)(\cdot)-\Phi(C)(1)],\\E'^{\rm new}(Y,\cdot):=\sum_{\operatorname{co}(C)=Y,\text{ analytic, no letter (a)}}[\Phi(C)(\cdot)-\Phi(C)(1)]+\tilde{\mathsf D}(Y,\cdot)-P(Y,\cdot),\\E'^a=E'^{\rm old}+E'^{\rm new},\qquad E'^h(Y,\cdot):=\sum_{\operatorname{co}(C)=Y,\text{ hard}}\Phi(C)(\cdot),\\\mathsf V':=\mathsf V_0+\sum_{\rm analytic}\Phi(C)(1)+\sum_{\rm constant}\Phi(C).\end{gathered}\end{equation}

We claim the following. For $Y\in\mathcal D$, the function $E'^a(Y,\cdot)$ is holomorphic on $\mathcal U_Y(x')$ and real at $G$-valued fields; it is invariant under $G$-valued gauge transformations, and it vanishes at the pure gauges. The function $E'^h(Y,\cdot)$ is real, Borel and gauge invariant, its domain consisting of the $G$-valued fields that live on the bonds of $Y^\approx$ and are tame there. The number $\mathsf V'$ is real, and $E'^{\rm new}$ does not depend on the member.

We prove \textup{(Clu)}. If $\alpha\sim\beta$ in $\mathfrak A$, then ${\rm dist}(J_\alpha,J_\beta)\le2\varrho_*+r_1\le r'-2\le s'$, by \textup{(In)} and \textup{(A2)}. Each of the diagonal paths of the core therefore adds at most $s'$ cells, so that $\#\operatorname{co}\le\sum(\#J_\alpha+s')-s'$, and hence $d(C)\ge(\kappa+1)(\#\operatorname{co}+s')$. Since $(\kappa+1)\#\operatorname{co}-\#\operatorname{co}-\kappa=\kappa(\#\operatorname{co}-1)$, this gives \textup{(Clu)}.

Next we determine on what the values depend. By item~5 of Step~7, a letter $\alpha$ with $J_\alpha\subset Y=\operatorname{co}(C)$ is a function of the field in $I_\alpha^{+r_1}$, and this set is contained in $Y^{+(r'-2)}$. Consequently the value $\Phi(C)(\cdot)$ depends only on the field on the bonds of $Y^\approx$, that is, on the bonds with both ends in $Y^\approx$, the field being extended by $1$. The entry $e$ of the pair does not enter: the letters of $C$ are functions of the field in $Y^{+(r'-2)}$, so the values $\Phi(C)(\cdot)$ with $\operatorname{co}(C)=Y$ depend on $X'_P$ only at the plaquettes $P$ based in $Y^{+(r'-2)}$; these plaquettes have all four sites in $Y^{+r'}=Y^\approx$, and at each of them $\phi'_P\le\varepsilon'<\pi$ holds on the domain of $E'^h$ (the tame $G$-valued fields on the bonds of $Y^\approx$) and \textup{(u3)} holds on the domain $\mathcal U_Y(x')$ of $E'^a$.

\emph{Convergence.} By \textup{(ii)}, \textup{(iii)} and \textup{(iv)} of Step~7, $|\Phi|\le\hat\Phi$, and $\hat\Phi$ does not depend on the field. \Cref{Y-G-L18}, applied with the condition \textup{(S)} proved in Step~8 and with \Cref{Y-G-LI}, gives $\sum_{C\ni\alpha}\hat\Phi(C)e^{d(C)}\le e_\alpha e^{2\varepsilon_2\#I_\alpha+d_\alpha}$, and $t$ times this bound for the clusters with $n(C)\ge2$. The torus is finite, hence so is the alphabet, and all the sums in \textup{(F3)}, those of $\mathsf V'$ included, converge absolutely and uniformly in the field.

The sets $Y_0(X)$ and $\{\square\}$, as well as the sets $Y$ of the letters \textup{(t)} and of $\tilde{\mathsf D}$, are connected; hence $\operatorname{co}(C)$ is connected too.

\emph{Holomorphy.} Fix $\xi$. By the locality just proved, the restriction of $\mathsf U(\eta;V')$ to $X^\approx$ is holomorphic in $V'\in\mathcal U_Y(x')$: the entries of $G$ and $C$ are holomorphic by $(\Gamma2)$, $\mathsf g^V$ is a polynomial in $\operatorname{Ad}V'$ by \Cref{Y-C-CIL64}\,(e), and the values $X'=-i\log\partial V'$ that enter are power series by \textup{(u3)}. Where $\chi_X=1$, \Cref{Y-G-LT}\,(T3) places this restriction in $\mathcal U_X(x)$, where $E^a(X,\cdot)$ is holomorphic; where $\chi_X=0$, the letter \textup{(a)} is identically zero. Likewise, \Cref{Y-E-C6}\,(a) makes the letter \textup{(v)} holomorphic where $\chi_\square=1$, and it is zero elsewhere. The tests depend on $\xi$ alone. Finally, $S_k(Y)$ is holomorphic by $(\Gamma4)$. The product of the letters is dominated by the integrable majorant of \textup{(iii)} of Step~7; hence $K(\mathfrak B)$ is holomorphic, and the series converges uniformly. Finally, $\tilde{\mathsf D}$ is holomorphic by $(\Gamma4)$, and $P$ by \Cref{Y-D-C5}\,(d).

\emph{Invariance.} Let $u$ be $G$-valued. Then $\mathsf U(\operatorname{Ad}_u\eta;V'^u)$ arises from $\mathsf U(\eta;V')$ by the gauge transformation whose value at each fine site is that of $u$ at the coarse site whose block contains that fine site; this is a property of the chart $\mathsf U$ in its two arguments $\eta$ and $V'$. Moreover, $G[V'^u]=\operatorname{Ad}_uG\operatorname{Ad}_u^{-1}$, $X'\mapsto\operatorname{Ad}_uX'$, $a_0\mapsto\operatorname{Ad}_ua_0$, and $S_k$ is covariant, so that $\eta(\operatorname{Ad}_u\xi,V'^u)=\operatorname{Ad}_u\eta(\xi,V')$. The two facts together show that $U$ at $(\operatorname{Ad}_u\xi,V'^u)$ is the gauge transform of $U$ at $(\xi,V')$. The functions $E^a$, $E^h$ and $w$ are gauge invariant by \textup{(ca)}, \textup{(ch)} and \textup{(c$\Xi$)}, and so is $V_\square$, by \Cref{Y-E-C6}\,(b); the factors $b_\square$, $\chi_\square$ and $\chi_X$ are functions of the numbers $|\xi_b|$, and $1_{\rm ind}$ and $1_{X^\approx}(U)$ are functions of the set $D(U)$, which a gauge transformation of $U$ leaves unchanged; the factor $A_\square$ satisfies $A_\square(\operatorname{Ad}_u\xi,V'^u,\operatorname{Ad}_ue)=A_\square(\xi,V',e)$, by \Cref{Y-E-C6}\,(e.iv); and the law of $\xi$ is invariant under $\operatorname{Ad}_u$. Therefore $\lambda_\alpha(\operatorname{Ad}_u\xi,V'^u)=\lambda_\alpha(\xi,V')$, and $K$ and $\Phi$ are invariant. If $u$ is given on the sites of $Y^\approx$ alone, we extend it by $1$. The extension of $V'^u$ by $1$ from the bonds of $Y^\approx$ and the field obtained from $V'$ by extending it by $1$ and then transforming it by the extended $u$ disagree only on bonds leaving $Y^\approx$, and no quantity here depends on such bonds.

We now come to the function $\tilde{\mathsf D}$ on $\mathcal U_Y(x')$. Together with $V'$, the field $V'^u$ lies in $\mathcal U_Y(x')$, since conjugation by a unitary matrix preserves the operator norm, and so \textup{(u1)}, \textup{(u2)} and \textup{(u3)} are preserved. By $(\Gamma4)$, $\tilde{\mathsf D}=-\mathsf D_k/2$ and $\mathsf D_k(Y,V')=-\mathsf D^q(Y,V')-[\mathsf D^u(Y,V')-\mathsf D^u(Y,1)]$, and the same formulas give $\tilde{\mathsf D}$ and $\mathsf D_k$ at complex fields. By \Cref{Y-C-CIITh73}\,(b), $\mathsf D^q(Y,V'^u)=\mathsf D^q(Y,V')$ at every field of its domain $\{|V'(c)|<1+\vartheta_0^\flat$ on the bonds of $Y^\sim\}$: in the proof of that theorem $\mathsf D^q$ is made of traces of matrices of \S\ref{Y-C-1}, and when $V'$ is replaced by $V'^u$, each of these matrices is conjugated by $\operatorname{Ad}_u$; this is an identity of matrices, so the traces agree. By \Cref{Y-C-CIITh74}\,(b), $\mathsf D^u(Y,V'^u)=\mathsf D^u(Y,V')$ at every field of $\mathfrak V$.

An element of $\mathcal U_Y(x')$, extended by $1$, lies in $\mathfrak V^\circ$ by \textup{(u1)}. Extend $u$ by $1$ as well, and transform the extended field by it: the result keeps the norms $|V'(c)|$ on the bonds of $Y^\approx$, is unitary on the bonds leaving $Y^\approx$, and equals $1$ on the bonds with no end in $Y^\approx$; so it, too, lies in $\mathfrak V^\circ$. It agrees with the extension of $V'^u$ except on the bonds leaving $Y^\approx$, whereas $\mathsf D_k(Y,\cdot)$ is determined by the field on the bonds of $Y^{+\lceil R_2/M\rceil}$ alone, and $\lceil R_2/M\rceil\le r_1-8<r'$. Hence $\tilde{\mathsf D}(Y,V'^u)=\tilde{\mathsf D}(Y,V')$ on $\mathcal U_Y(x')$.

By \Cref{Y-D-C5}\,(d), $P$ is invariant on $\mathcal U_{\{\square\}}(x')$. So $E'^a(Y,\cdot)$ is invariant on $\mathcal U_Y(x')$, as \textup{(ca)} requires.

At $G$-valued fields every letter is real, by \Cref{Y-E-C6}\,(d), by \textup{(ca)}, \textup{(ch)} and \textup{(c$\Xi$)}, and by $(\Gamma2)$ and $(\Gamma4)$; hence the values $\Phi(C)(1)$ and the values $\Phi(C)$ of the constant clusters are real as well. The number $\mathsf V_0$ is real, because Step~1 displays $e^{\mathsf V_0}$ as a product of numbers that are positive: $\mathsf V$ is real by \textup{(c0)}, and $\det C^1>0$ by $(\Gamma2)$.

At the pure gauges, invariance gives $\Phi(C)(1^u)-\Phi(C)(1)=0$; for $\tilde{\mathsf D}$ the vanishing at $1^u$ follows from $(\Gamma4)$, and for $P$ from \Cref{Y-D-C5}\,(d).

As to the Borel property of $E'^h$: the letters are jointly Borel in $(\xi,V')$, because the members are Borel and the kernels are continuous; Tonelli's theorem then applies to the means, and the series converges uniformly.

Finally, the letters \textup{(v)} and \textup{(t)}, as well as $\tilde{\mathsf D}$ and $P$, do not depend on the member (\Cref{Y-E-C6}\,(a); $T_Y$ depends on $G$, $S_k$ and $a_0$).

\emph{Step 10: the clause \textup{(ca)} at $x'$.} We claim that for every cell $\square$ \[\sum_{Y\ni\square}e^{\kappa(\#Y-1)}\sup_{\mathcal U_Y(x')}|E'^{\rm old}(Y,\cdot)|\le0.24978\,c_a(x'),\] and that the same sum, with $E'^{\rm new}$ in place of $E'^{\rm old}$, is $\le0.30695\,c_a(x')$. The two bounds together give $\sup|E'^a(Y,\cdot)|\le0.56\,c_a(x')e^{-\kappa(\#Y-1)}$.

We first compare the caps at $x$ and at $x'$. By \textup{(A8)}, \[c_a(x')=L^{1/2}(p_0'/p_0)^3c_a(x)\ge2^{-3/4}L^{1/2}c_a(x)\ge0.59460\,L^{1/2}c_a(x).\]

\emph{The rooted bound.} Let $C$ be a cluster with $\operatorname{co}(C)\ni\square$, let $\alpha$ be one of its letters, and let $\square_0$ be the root of $J_\alpha$. Then $\#\operatorname{co}\ge{\rm dist}(\square_0,\square)+1$. By \textup{(Clu)} and \Cref{Y-G-L18} we obtain, for the clusters that hold a letter $\alpha$ of a given kind, \[\sum_{\operatorname{co}\ni\square,\ C\ni\alpha\text{ of the given kind}}\hat\Phi e^{\kappa(\#\operatorname{co}-1)}\le\sum_{\square_0}e^{-{\rm dist}(\square_0,\square)-1}\sum_{\alpha\text{ rooted at }\square_0}[t]\,e^\circ_\alpha e^{4x^{-1/8}\#I_\alpha+(\kappa+1)\#J_\alpha-\kappa},\] where the bracketed factor $t=x^{-1/4}$ is present when the sum runs over the clusters with $n(C)\ge2$ alone.

By \textup{(A3)}, $e^{4x^{-1/8}\#I_\alpha}\le e^{0.08\#J_\alpha}$, and at $\#J_\alpha=1$ we have $e^{-1}e^{(\kappa+1)-\kappa}=1$. Wherever $\Phi(C)(1)$ is subtracted, each of these bounds is doubled.

\begin{enumerate}\item[\textup{(i)}] Let the cluster consist of one letter \textup{(a)} whose set $X$ is small. Then $\Phi(C)(V')=E_\xi[\chi_X(e^{E^a(X,U)}-1)]=:\mathcal G_X(V')$, and $\operatorname{co}=\{\square_0\}$, where $\square_0$ is the cell of the first cube of $X$.

Using $T>1$, $\omega\le Mr'$ from \textup{(A2)}, $c_a(x)\le1$ from \textup{(A4)}, and \textup{(A0)}, \textup{(A8)}, we have \[\sup_{\mathcal U_{\{\square_0\}}(x')}|\mathcal G_X-\mathcal G_X(1)|\le4T^{-2}c_a(x)e^{-\kappa(\#X-1)}.\]

The argument for this bound is that of Step~8, \textup{(8.2)}, of the proof of \Cref{Y-H-C9} and of the passage on $u$ in ``The clauses without a size'' of the same step, together with part~\textup{(ii)} of Step~2 of that proof, taken at $A=1$. These two steps rely on \Cref{Y-H-L51} and Thms.~\ref{Y-H-Th53} and~\ref{Y-H-Th57}. The passage on $u$ extends by $1$ the $G$-valued gauge transformation $u$ given on the sites of $\mathsf B_\omega$. Since the regions of radii $r_G$ and $R_3$ about the blocks that meet $X^\approx$ lie inside $\mathsf B_\omega$, every quantity that enters $\mathcal G_X$ is independent of the field on the bonds that leave $\mathsf B_\omega$.

Of the test $\chi_X$, that argument uses only the bounds $|\xi_{b'}|\le c_\xi\varepsilon$ on the windows $Y_b$ of the bonds $b$ of $X^\approx$. The test $\chi_X$ of \Cref{Y-G-LT}\,(T3) provides these bounds, since $r_w+1\le R$ and every bond of $X^\approx$ is based in a block that meets $X^\approx$.

This test is a function of the norms $|\xi_b|$ alone: it is independent of $V'$, and it is unchanged under $\xi\mapsto\operatorname{Ad}_u\xi$.

Hence $\mathcal G_X$, defined with this test, is holomorphic and $G$-invariant on $U^{17T}(\mathsf B_\omega;x')$ and is bounded there by $2c_a(x)e^{-\kappa(\#X-1)}$, by the argument of Step~8, \textup{(8.2)}, of the proof of \Cref{Y-H-C9}; and the statement of \Cref{Y-H-Th57} covers all $G$-invariant holomorphic functions on that domain that are bounded.

That argument is independent of the condition $\mathsf a_{\rm c}\varepsilon'\le1$ of \textup{(0b)}. The floors on which it rests are the other conditions of \textup{(0b)}, then \textup{(0c)} and \textup{(0d)}, where \textup{(0b)}, \textup{(0c)}, \textup{(0d)} are clauses of \textup{(Y.0)} in item~5 of \S\ref{Y-H-1}, and finally $\omega\le Mr'$ and $c_a(x)\le1/2$; all of these floors are stated in \textup{(A0)}, \textup{(A2)}, \textup{(A4)}, \textup{(A8)}.

The two hypotheses of \Cref{Y-H-Th57} on $\omega$ are satisfied: $3\omega+1\le3Mr'+1\le6Mr'+12=\ell_w(x')$ by \textup{(A2)}, and $2\omega+1<N$ by the clauses on $N$ in Step~0. Moreover, that argument is independent of Thms.~\ref{Y-G-ThA} and~\ref{Y-GB-ThB}.

Since the grids fit, there are $L^3$ first cubes in $\hat\square_0$; for each cube, $\sum_{X\ni\text{cube}}e^{-\kappa(\#X-1)}\le1.2124$ by Step~8; and $4\cdot1.2124\le4.85$. Summing over the small sets $X$ whose first cube lies in $\hat\square_0$, we obtain by \textup{(L2$^{\rm su}$)} \[\FitDisplay{\dimexpr\linewidth-1pt\relax}{\begin{aligned}
&\sum_{X\text{ small, first cube in }\hat\square_0}\sup_{\mathcal U_{\{\square_0\}}(x')}|\mathcal G_X-\mathcal G_X(1)|
\\
&\quad{}\le4.85L^3T^{-2}c_a(x)=4.85\cdot361\mathsf N_{\rm c}^2c_a(x)
\\
&\quad{}\le(4.85\cdot361/(1.2\cdot10^4))L^{1/2}c_a(x)
\\
&\quad{}\le0.24539\,c_a(x').
\end{aligned}}\]

In numbers: $4.85\cdot361/12000\le0.1459042$, and $0.1459042/0.59460\le0.245383$.

\item[\textup{(ii)}] Let the cluster consist of one letter \textup{(a)} whose set $X$ is not small. Then $\operatorname{co}=Y_0(X)\ni\square$, so $X$ holds one of $(L+2)^3$ cubes, and $n=\#X\ge(L-3)/2$. The term belonging to $n$ is $\le4c_a(26e)^{n-1}e^{-\kappa(n-1)}e^{\kappa(7+432n/L)}$.

The ratio of two consecutive terms is $\le e^{-1.74}$; so the sum is $\le1.22$ times its first term, and the exponent of the first term is $\le-((L-5)/2)(\kappa-4.2581)+223\kappa\le-\kappa L/7$. For the last inequality we note that the difference $((L-5)/2)(\kappa-4.2581)-223\kappa-\kappa L/7$ is $\ge L(5\kappa/14-2.12905)-225.5\kappa+10.64$; at $\kappa=6$, $L\ge1.2\cdot10^{13}$ this lower bound is $\ge0.0138L-1343\ge1.6\cdot10^{11}>0$, and its derivative in $\kappa$ is $5L/14-225.5>0$.

Case~\textup{(ii)} therefore contributes $\le16L^3e^{-\kappa L/7}c_a(x)\le10^{-6}c_a(x')$.

\item[\textup{(iii)}] Let $n(C)\ge2$, and let $C$ hold a letter \textup{(a)}; we root the cluster at one of its letters \textup{(a)}. In the rooted sum of the sizes $e^\circ$, the letters with small $X$ give $L^3\cdot1.2124\cdot2c_a\le2.43L^3c_a$; the letters whose $X$ is not small add a negligible amount, by the argument of case~\textup{(ii)}; in all, the rooted sum is $\le4.86L^3c_a$.

So case~\textup{(iii)} contributes $\le2\cdot49.98e^{0.1}x^{-1/4}\cdot4.86L^3c_a(x)\le110.48\cdot4.86L^3x^{-1/4}c_a(x)\le0.0026L^{1/2}c_a(x)\le0.00438\,c_a(x')$, by \textup{(A8)}, since $2\cdot49.98e^{0.1}\le110.48\le111$ and \textup{(A8)} contains the clause $111\cdot4.86L^3x^{-1/4}\le0.0026L^{1/2}$. The factor $49.98$ comes from $\sum_{n\in\mathbb Z^3}e^{-|n|_\infty}\le49.98$, and the sum on the left bounds the sum $\sum_{\square_0}e^{-{\rm dist}(\square_0,\square)}$ in the rooted bound, by Step~8.

Adding the three cases, we find $0.24539+0.000001+0.00438\le0.24978$, the bound for $E'^{\rm old}$.

\item[\textup{(iv)}] Let the cluster consist of one letter \textup{(v)}. We use the bound $e^\circ_\alpha=2K_vp_0^3x^{-1/2}$ on the supremum of this letter $\alpha$, from item~1 of Step~7. This bound does not carry the factor $e^{2x^{-1/8}\#I_\alpha}$ by which $e_\alpha$ exceeds $e^\circ_\alpha$. We obtain $2\cdot2K_vp_0^3x^{-1/2}=(4K_v/E_a)c_a(x)\le(1/8)L^{1/2}c_a(x)\le0.21023\,c_a(x')$, as $E_a\ge32L^{-1/2}K_v$ by \textup{(O12)}.

\item[\textup{(v)}] For the terms $\tilde{\mathsf D}$ we have $\sum_{Y\ni\square}e^{(\kappa-\kappa^\sharp)(\#Y-1)}1.5\mathsf K_1^\flat\varepsilon'\le1.01\cdot1.5\mathsf K_1^\flat\varepsilon'$. Since $\varepsilon'=c_a(x')/(E_ap_0'^2)$, the bound $1.01\cdot1.5\mathsf K_1^\flat\varepsilon'$ is $\le(1.01/17)c_a(x')\le0.05942\,c_a(x')$, by \textup{(A8)}.

\item[\textup{(vi)}] For the terms $P$ we have $K_{\rm q}p_0'^2\varepsilon'=(K_{\rm q}/E_a)c_a(x')\le(1/32)c_a(x')=0.03125\,c_a(x')$, as $E_a\ge32K_{\rm q}$ by \textup{(O12)}.

\item[\textup{(vii)}] Let $n(C)\ge2$, and let $C$ hold no letter \textup{(a)}; we root it at a letter \textup{(v)} or at a letter \textup{(t)}. The contribution is then $\le110.48x^{-1/4}[2K_vp_0^3x^{-1/2}+x^{-20}]\le111x^{-1/4}(2K_v/E_a+1)c_a(x)\le0.003L^{1/2}c_a(x)\le0.00505\,c_a(x')$, by \textup{(A8)}.

\item[\textup{(viii)}] Let the cluster consist of one letter \textup{(t)}. Then $2\sum_{Y\ni\square}e^{\kappa(\#Y-1)}e_Y^{1/2}e^{2x^{-1/8}\#Y}\le2x^{-20}\le0.001\,c_a(x')$, where the last inequality is the last clause of \textup{(A8)}.

Adding the five cases \textup{(iv)} to \textup{(viii)}, we find $0.21023+0.05942+0.03125+0.00505+0.001=0.30695$, the bound for $E'^{\rm new}$.\end{enumerate}

\emph{Step 11: the clause \textup{(ch)} at $x'$.} For every cell $\square$ we have \[\sum_{Y\ni\square}e^{\kappa(\#Y-1)}\sup|E'^h(Y,\cdot)|\le10^3L^3c_h(x)+50(2u_b+u_1)\le c_h(x').\]

Every hard cluster holds a letter $\alpha$ that is a letter \textup{(h)}, a letter \textup{(b)} or a zone letter of $\mathfrak A$. We sum by means of the rooted bound of Step~10, which rests on claim~\textup{(1)} of \Cref{Y-G-L18}.

If $\alpha$ is a letter \textup{(h)}, we obtain $49.98e^{0.08}\cdot L^3\cdot1.2124\cdot4c_h(x)\le262.6L^3c_h(x)$; when the sets $X$ that are not small are added, by the argument of case~\textup{(ii)} of Step~10, the total is $\le10^3L^3c_h(x)$. If $\alpha$ is a letter \textup{(b)}, we obtain $49.98e^{0.08}u_b\le54.2u_b\le100u_b$.

For a zone letter $\zeta$, \textup{(L)} gives \[e^\circ_\zeta e^{4x^{-1/8}\#I_\zeta+(\kappa+1)\#J_\zeta-\kappa}\le e^\circ_\zeta e^{(\kappa+2)\#J_\zeta}\le u_1\prod_o\hat w(o),\] where we use that the set $t$ of tested cells is non-empty for $\zeta\in\mathfrak A$. By \textup{(An)}, the sum over the zone letters $\zeta$ rooted at $\square_0$ is $\le u_1\hat w(\square_0)e^{a_{\square_0}}\le u_1$, and so the zone letters give $49.98u_1\le50u_1$.

The first clause of \textup{(A9)} gives $10^3L^3c_h(x)\le(1/2)e^{-c_hp_0'^2}=(1/2)c_h(x')$.

By \Cref{Y-F-C7}\,(e), $u_1=e^{-\delta_3p_0^2}$; hence $50u_1\le(1/4)e^{-\delta_3p_0^2/2}$, as the second clause of \textup{(A9)} is $e^{\delta_3p_0^2/2}\ge200$. Next, $h_*=C_\varphi^\flat(4+12K_*\mathsf g_*+a^{1/2}L)\ge4C_\varphi^\flat\ge4(2/\gamma)^{1/2}>1$, since $C_\varphi^\flat\ge\mathsf G_0=(2/\gamma)^{1/2}3L^3\mathsf g(\theta_{\mathbb C})\ge(2/\gamma)^{1/2}$ by item~C.I~1.2 of \S\ref{Y-C-1}, where the last step rests on $3L^3\mathsf g(\theta_{\mathbb C})\ge1$, given by the term $n=0$; therefore $\delta_3\le(3/256)c_\xi^2$ and $c_\xi^2/2-\delta_3/2\ge0.494c_\xi^2$. With this, and with $u_b=2.5v(c_\xi p_0+(c_\xi p_0)^{-1})e^{-c_\xi^2p_0^2/2}$ as in item~F of \S\ref{Y-F-1}, we obtain $100u_be^{\delta_3p_0^2/2}\le250v(c_\xi p_0+1/(c_\xi p_0))e^{-0.49c_\xi^2p_0^2}\le1/4$, where the last inequality is the third clause of \textup{(A9)}.

So $50(2u_b+u_1)\le(1/2)e^{-\delta_3p_0^2/2}\le(1/2)e^{-c_hp_0'^2}$, because $c_h\le\delta_3/2$ and $p_0'\le p_0$.

\emph{Step 12: the expansion of $\Xi'$ at a fixed $e$.} In this step we show that at every $(V',e)$ \[\rho'(V')=e^{\mathsf V'}\exp\{-x'\mathrm{En}'(V')+\sum_Y1_{Y^\approx}(V')[E'^a+E'^h](Y,V')\}\,\Xi',\] where $\Xi'=\sum_{\{\mathsf Y_i\}}\prod_iw'(\mathsf Y_i,V',e)$, the sum being taken over the families of the definition of $\Xi_w$ in item~1 of \S\ref{Y-G0-1} applied at $x'$ with cells in place of cubes ($s'$ in place of $s(x)$, $D'$ in place of $D(V)$), and where $w'$ is given by the definition \textup{(F4)} that follows.

\emph{Definition \textup{(F4)}.} Let $\mathsf Y$ be a finite set of cells that is not empty. When some cell of $\mathrm{mk}'(V')$ lies in $\mathsf Y^{+(s'-1)}\setminus\mathsf Y$, we set $w'(\mathsf Y,V',e):=0$. Otherwise $w'(\mathsf Y,V',e)$ is the sum of \[\prod_{\gamma\in\mathsf G}K(\gamma)\prod_{C\in\mathfrak N}(e^{-\Phi^\natural(C)}-1)\prod_{Y\in\mathfrak N^\flat}(e^{\tilde{\mathsf D}(Y,V')}-1)\] over all triples $(\mathsf G,\mathfrak N,\mathfrak N^\flat)$ with the following properties: $\mathsf G$ is a set of rough polymers at $V'$, no two of them $\sim$, and the marks of its forced letters are exactly $\mathrm{mk}'\cap\mathsf Y$, with every class met once; $\mathfrak N$ is a finite set of clusters $\sim\mathsf G$; $\mathfrak N^\flat$ is a set of long $Y$; the cores together make up $\mathsf Y$; and the graph of links is connected. Rough polymers and $K(\gamma)$ are defined in part~\textup{(a)}, $\Phi^\natural$ in part~\textup{(b)}, and the elements, their cores and the links in part~\textup{(c)}.

The definition \textup{(F4)} has one more clause, its last: the third member of $u'$ is the mean $w'(\mathsf Y,V'):=\int_{\mathfrak E}w'(\mathsf Y,V',e)\,de$, which is treated in item~5 of Step~13.

We claim that $w'(\mathsf Y,V',e)$ depends on $V'$ only through the bonds of $\mathsf Y^{+s'}$, and on $e_P$ only at the plaquettes of cells of $\mathsf Y$.

\begin{enumerate}\item[\textup{(a)}] \emph{Rough polymers.} We call a rough polymer at $V'$ any $\sim$-connected set $\gamma$ of letters, all actual at $V'$ and pairwise compatible, that contains a forced letter; for it we put $K(\gamma)=E_\xi\prod_\gamma\lambda_\alpha$.

We split a set $\mathfrak B$ of \textup{(F2)} into its $\sim$-components, and we let $\mathsf G$ be the set of those components that hold a forced letter; by \textup{(Act1)} these are rough polymers, and the union of the sets $I_\alpha$ of their letters contains $f^{+2r_1}$. We put $\mathfrak A_{\mathsf G}:=\{\alpha\in\mathfrak A:\alpha\text{ not}\sim\text{any letter of }\mathsf G\}$.

Every letter of $\mathfrak A_{\mathsf G}$ is actual. For such a letter, $I_\alpha$ lies off $f^{+3r_1}$; hence the letters \textup{(v)} and \textup{(b)} are actual, and so are the letters \textup{(a)} and \textup{(h)}, for which ${\rm cl}(X^\approx)\subset I_\alpha$; for a zone letter, $\mathfrak b^{+2r_1}\subset I_\alpha$, and so ${\rm dist}(\mathfrak b,f)>5r_1\ge4r_1+2$.

The remaining components of $\mathfrak B$ are the $\sim$-connected compatible sets of letters of $\mathfrak A_{\mathsf G}$, and all letters of such a set are compatible with all letters of $\mathsf G$, because incompatibility occurs only between letters that are $\sim$; hence such a set and $\mathsf G$ can occur together among the components of a set of \textup{(F2)}.

Lems.~\ref{Y-G-LI} and~\ref{Y-G-L18}, applied to the alphabet $\mathfrak A_{\mathsf G}$ (we use \textup{(i)} of Step~7, and the condition \textup{(S)} is weaker for a smaller alphabet), give \[e^{\sum n_j}\mathcal Z=\sum_{\mathsf G}\prod_\gamma K(\gamma)\exp\{\sum_{C\text{ not}\sim\mathsf G}\Phi(C)(V')\},\] where $\mathsf G$ ranges over the sets of rough polymers that are pairwise not $\sim$ and contain every forced class once.

\item[\textup{(b)}] \emph{On and off.} A cluster $C$ is on at $V'$ if it is constant or if $V'$ is tame on $\operatorname{co}(C)^\approx$; otherwise it is off. Every cluster $C$ that is not $\sim\mathsf G$ is on: its letters have ${\rm dist}(J_\alpha,\mathrm{mk}')>s_0+3r_1$, and every cell of a joining path is within $r'$ of one of the $J_\alpha$, so that ${\rm dist}(\operatorname{co},\mathrm{mk}')>s_0-r'=r'$. The value of such a cluster at $V'$ is its value at the restriction of $V'$ to $\operatorname{co}^\approx$, by Step~9.

We put $\Phi^\natural(C):=\Phi(C)(V')$ if $C$ is on, $\Phi(C)(1)$ if $C$ is analytic and off, and $0$ if $C$ is hard and off. This quantity is independent of $e$, and $|\Phi^\natural|\le\hat\Phi$.

We go through the types. For an analytic cluster we have $\Phi(V')=\Phi(1)+(\Phi(V')-\Phi(1))$ if it is not $\sim\mathsf G$, $0=\Phi(1)+(\Phi(V')-\Phi(1))-\Phi(V')$ if it is $\sim\mathsf G$ and on, and $0=\Phi(1)-\Phi(1)$ if it is $\sim\mathsf G$ and off. For a hard cluster the same three cases read $\Phi(V')=\Phi(V')$, $0=\Phi(V')-\Phi(V')$ and $0=0-0$. For a constant cluster they read $\Phi=\Phi$ if it is not $\sim\mathsf G$ and $0=\Phi-\Phi$ if it is $\sim\mathsf G$.

Combining these identities gives \[\FitDisplay{\dimexpr\linewidth-1pt\relax}{\sum_{C\text{ not}\sim\mathsf G}\Phi(C)(V')=[\sum_{\rm analytic}\Phi(C)(1)+\sum_{\rm constant}\Phi(C)]+\sum_Y1_{Y^\approx}(V')[E'^a-\tilde{\mathsf D}+P+E'^h](Y,V')-\sum_{C\sim\mathsf G}\Phi^\natural(C),}\] where all sums converge absolutely.

We insert this, together with \textup{(iv)} of Step~3, into \textup{(F1)}. The terms $1_{Y^\approx}\tilde{\mathsf D}$ and $1_{\square^\approx}P$ cancel, and $\sum n_j$ is now carried by the letters. Hence the formula for $\rho'(V')$ holds with \[\Xi'=\sum_{\mathsf G}\prod K(\gamma)e^{-\sum_{C\sim\mathsf G}\Phi^\natural(C)}e^{\sum_{\rm long}\tilde{\mathsf D}}=\sum_{(\mathsf G,\mathfrak N,\mathfrak N^\flat)}\prod K(\gamma)\prod_{\mathfrak N}(e^{-\Phi^\natural}-1)\prod_{\mathfrak N^\flat}(e^{\tilde{\mathsf D}}-1),\] where $\mathfrak N$ ranges over the finite subsets of $\{C\sim\mathsf G\}$ and $\mathfrak N^\flat$ over the sets of long $Y$.

\item[\textup{(c)}] \emph{Elements, links and a bijection.} The elements of a triple are its letters, each with the core $J_\alpha$, the clusters $C$ of $\mathfrak N$, each with the core $\operatorname{co}(C)$, and the sets $Y$ of $\mathfrak N^\flat$, each with the core $Y$. Two elements are linked if their cores lie within $s'$ of each other.

When no cell of $\mathrm{mk}'(V')$ lies in $\mathsf Y^{+(s'-1)}\setminus\mathsf Y$, that is, when the first clause of \textup{(F4)} does not make $w'$ vanish, every mark outside $\mathsf Y$ lies at a distance greater than $s'-1\ge5r'\ge2s_0+4r_1+2$ from $\mathsf Y$; therefore $\mathrm{mk}'\cap\mathsf Y$ is a union of classes.

We now set up a bijection. If a triple as in part~\textup{(b)} is given, each component of its graph of links yields a set $\mathsf Y_i$; these sets lie at pairwise distance $>s'$, and each of them carries a triple of \textup{(F4)}. Indeed, two letters of one $\gamma$ that are $\sim$, and a cluster $C$ together with a letter of $\mathsf G$ to which $C$ is $\sim$, have cores within $2(s_0+\varrho_*)+r_1=4r'+5r_1+2\lceil s(x)/L\rceil+4\le5r'-2$ of each other, by \textup{(A2)}; and every long $Y$ comes within $r'$ of some mark, and that mark belongs to the core of a forced letter. Every $\mathsf Y_i$ therefore holds a forced letter, thus a mark and thus a plaquette of $D'$; since every mark lies in some $\mathsf Y_i$, the set $D'$ lies in their union; and no $\mathsf Y_i$ has a mark in $\mathsf Y_i^{+(s'-1)}\setminus\mathsf Y_i$.

Conversely, let sets $\mathsf Y_i$ be given, at pairwise distance $>s'$, which together hold $\mathrm{mk}'$, each with a triple. Letters of different indices $i$ have ${\rm dist}(I_\alpha,I_\beta)>s'-2(s_0+\varrho_*)=4r'-2\varrho_*\ge r_1$, so they are not $\sim$ and hence compatible. The condition of compatibility on polymers of kind~p holds for them as well, since, by the proof of \textup{(Inc)} in Step~4, two zone letters $\alpha$, $\beta$ with polymers of kind~p within $s(x)$ of each other have ${\rm dist}(I_\alpha,I_\beta)=0$. The mark that makes a set $Y$ of $\mathfrak N_i^\flat$ off lies in $\mathsf Y_i^{+r'}$, hence in $\mathsf Y_i$. Therefore the union of the given triples is a triple as in part~\textup{(b)}, and its components are the $\mathsf Y_i$.

If $D'$ is empty, only the empty triple remains, and $\Xi'=1$.

\item[\textup{(d)}] \emph{On what $w'$ depends.} Measured in cells from $\mathsf Y$, the ingredients of $w'$ depend on $V'$ within the following distances: a forced letter within $s_0+\varrho_*+r_1\le3r'$, by \textup{(In)} and by \textup{(Rd)}, which holds at every pair; the class terms $w_j$ within $2r'$; the numbers $n_j$ within $3r'$; a long $Y$ within $\lceil R_2/M\rceil$; and $\Phi^\natural(C)$ within $r'$.

The marks that decide whether a letter is actual or forced, what the classes are, and whether a cluster or a set is on, off, short or long lie within $2s_0+4r_1+2\le5r'$ of a core, and the marks that decide whether the first clause of \textup{(F4)} applies are those in $\mathsf Y^{+(s'-1)}\setminus\mathsf Y$, within distance $s'-1$ of $\mathsf Y$. All these distances are $\le s'-1$, and the plaquettes based in a cell decide whether that cell is a mark; so $w'$ depends on $V'$ only through the bonds of $\mathsf Y^{+s'}$.

The component $e_P$ of $e$ occurs at $P\in D'$ only; hence $e$ enters only through plaquettes that belong to marks lying within $s'-1$ of $\mathsf Y$, and, when $w'\ne0$, all such marks are cells of $\mathsf Y$; finally, the first clause of \textup{(F4)} is independent of $e$.\end{enumerate}

\emph{Step 13: the strong cap and the mean over $e$.} In this step we show that at every $(V',e)$ \[|w'(\mathsf Y,V',e)|e^{\kappa(|\mathsf Y|_c-1)}\le\exp\{-c_\Xi x'\varepsilon'\sum_{P\in D'\cap\mathsf Y}\phi'_P\},\] that the mean $w'(\mathsf Y,V'):=\int w'(\mathsf Y,V',e)\,de$ satisfies \textup{(c$\Xi$)} at $x'$ with the strong cap, and that $\Xi'(V')=\sum_{\{\mathsf Y_i\}}\prod w'(\mathsf Y_i,V')$ at every $V'$. \Cref{Y-G-ThA} follows from this.

\begin{enumerate}\item[\textup{1.}] Let $w'\ne0$. By \textup{(i)} of Step~7, $|K(\gamma)|\le\prod e_\alpha$. The cores are connected and are linked within $s'$; hence $|\mathsf Y|_c\le\sum_{\rm elements}(\#\mathrm{core}+s')-s'$, and so \[e^{\kappa(|\mathsf Y|_c-1)}\le\prod_{\rm letters}e^{d_\alpha}\prod_{\mathfrak N}e^{d(C)}\prod_{\mathfrak N^\flat}e^{(\kappa+1)(\#Y+s')}.\]

Since $|\Phi^\natural|\le\hat\Phi\le1$, we have $|e^{-\Phi^\natural}-1|\le2\hat\Phi$. By Step~3, $|e^{\tilde{\mathsf D}(Y)}-1|e^{(\kappa+1)(\#Y+s')}\le\hat w(Y)$.

\item[\textup{2.}] We apply \Cref{Y-G-L19}, with $t=x^{-1/4}\le1/2$, taking for $\mathsf F$ the forced letters of $\mathsf G$, for $\mathfrak B_o$ the other letters of $\mathsf G$, and for $\mathfrak M$ the clusters, with $\hat w_C=2\hat\Phi e^{d(C)}$. Each $\sim$-component of $\mathfrak B_o$ holds a letter that is $\sim$ to a letter of $\mathsf F$, since $\gamma$ is $\sim$-connected; and the condition of the lemma holds with $2t\varepsilon_2\le\varepsilon_2$.

In addition to the lemma, we use a separate fact: for a further finite set $\mathfrak M^\flat$ with $\sum_{\mathfrak M^\flat}\hat w\le\sum_{\mathsf F}A_\zeta$ we have $\sum_{\mathfrak N^\flat\subset\mathfrak M^\flat}\prod\hat w\le\prod(1+\hat w)\le e^{\sum A_\zeta}$. We take for $\mathfrak M^\flat$ the long $Y$ within $r'$ of the marks in $\mathsf Y$, and $A_\zeta=\sum_{J(\zeta)}\#\mathrm{mk}'_j$, by \textup{(iii)} of Step~3.

The outcome is \[|w'|e^{\kappa(|\mathsf Y|_c-1)}\le\sum_{\mathsf F}\prod_{\zeta\in\mathsf F}e_\zeta e^{d_\zeta+2\varepsilon_2\#I_\zeta+A_\zeta},\] where $\mathsf F$ runs over those sets of forced letters for which each class with marks in $\mathsf Y$ occurs exactly once among their classes and no other class occurs. The right side is a majorant, since the forced letters in it need not be actual at $V'$.

\item[\textup{3.}] Since $\#I_\zeta\le(2(s_0+\varrho_*)+1)^3\#J_\zeta$, \textup{(A3)} gives $e^{4x^{-1/8}\#I_\zeta}\le e^{\#J_\zeta}$. With \textup{(L)} it follows that $e_\zeta e^{d_\zeta+2\varepsilon_2\#I_\zeta+A_\zeta}\le e^{(\kappa+1)s'}\prod_{o\in\zeta}\hat w(o)$, where $\hat w(\square)$, $\hat w_p$, $\hat w_{p'}$ are the weights of Step~8, and \[\hat w(f_j):=\exp\{-(r_\Xi+c_\Xi)x'\varepsilon'\sum_{D'_j}\phi'_P+[1.01\mathsf K_1^\flat(2r'+1)^3+1]\#D'_j+(\kappa+4)((1+s')\#\mathrm{mk}'_j+s')\}.\]

Into $\hat w(f_j)$ enter the term $\mathrm{Res}$, the factor $e^{|n_j|}$, the number $A_\zeta$, and the inequality $\#\mathrm{mk}'_j\le\#D'_j$.

We have $r_\Xi\ge r_\Xi^\flat\ge(3/4)(1/2-d_0-c_\Xi)>0$, since $\mathsf t_\rho\le(1/8)(1/2-d_0-c_\Xi)$ by \textup{(SD)} and $3e_1^\flat$ obeys the same bound by \textup{(A7)}. Each $P\in D'$ has $x'\varepsilon'\phi'_P>p_0'^2$, so $-r_\Xi x'\varepsilon'\sum_{D'_j}\phi'_P\le-r_\Xi^\flat p_0'^2\#D'_j$, and therefore $\hat w(f_j)e^{2\#\mathrm{mk}'_j}\le e^{-c_\Xi x'\varepsilon'\sum_{D'_j}\phi'_P}\upsilon_f^{\#D'_j}$.

\item[\textup{4.}] For every class we take out the factor $e^{-c_\Xi(\cdot)}\upsilon_f^{\#D'_j/2}$. With $\hat w''(f_j)=\upsilon_f^{\#D'_j/2}e^{-2\#\mathrm{mk}'_j}$ the condition of \Cref{Y-G-An0} holds: at a cell, the cells and polymers contribute $\varepsilon''/2$, by \textup{(A6)}; at most one class has $f_j\ni\square$; and $a_{f_j}\le\#\mathrm{mk}'_j$.

This last bound comes from the definition of $a_o$ in Step~8, independently of the bound $a_o\le\#\mathrm{core}(o)/50$ obtained there, because $\mathrm{core}(f_j)$ may have $(1+s')\#\mathrm{mk}'_j$ cells: we have $f_j^{+(4r_1+2)}\subset \mathrm{mk}_j'^{+(s_0+4r_1+2)}$, so $a_{f_j}\le x^{-1/4}(2(s_0+4r_1+2)+1)^3\#\mathrm{mk}'_j\le\#\mathrm{mk}'_j/50$ by \textup{(A3)}.

Hence $\hat w''(f_j)e^{a_{f_j}}\le\upsilon_f^{1/2}\le\varepsilon''/4$, by \textup{(A7)}, and so $\sum_{\zeta\ni f_j}e^{(\kappa+1)s'}\prod\hat w''\le e^{(\kappa+1)s'}\upsilon_f^{1/2}\le1$. A set $\mathsf F$ is determined by choosing, for each class, either nothing or the letter of which that class is the first class; thus $\sum_{\mathsf F}\prod\le\prod_j2$. Finally $\prod_j2\upsilon_f^{\#D'_j/2}\le1$.

\item[\textup{5.}] We apply \Cref{Y-H-L52}, which is stated and proved in \textup{(9)} of Step~10 of the proof of \Cref{Y-H-C9}, to the function \[\Xi'(V'):=\rho'(V')e^{-\mathsf V'}\exp\{x'\mathrm{En}'-\sum1_{Y^\approx}[E'^a+E'^h]\},\] which is independent of $e$. Its hypotheses \textup{(i)} to \textup{(iv)} hold, as follows.

\begin{enumerate}\item[\textup{(i)}] The curvature $X'$, the masks, the marks and the classes are Borel in $(V',e)$; $T_\square$ is jointly Borel, by \Cref{Y-D-C5}\,(b); the kernels are continuous; $\tilde{\mathsf D}(Y,\cdot)$ is continuous and $n_j$ is Borel, by $(\Gamma4)$ and by \textup{(ii)} of Step~3; $A_\square$ is Borel, by \Cref{Y-E-C6}\,(e.v); and the members are Borel.

Hence every letter is jointly Borel in $(\xi,V',e)$, and every $K(\gamma)$ is Borel in $(V',e)$, by Tonelli's theorem and Step~7. The Borel sets $\{\mathrm{mk}'(V')\cap\mathsf Y^{+(s'-1)}=\mathsf d\}$ are finite in number, and on every one of them the sum of \textup{(F4)} is a series of functions of this kind that converges absolutely; moreover, the bound obtained in items~1 to~4 is independent of $e$.

\item[\textup{(ii)}] The marks, the classes, and the properties actual, on, off, short and long are functions of the $\phi'_P$. Further, $X'(V'^u,\operatorname{Ad}_ue)=\operatorname{Ad}_uX'(V',e)$; hence $a_0$ and $\eta$ are covariant, and $T_\square$, $w_j$ and the letters are invariant together with $\xi\mapsto\operatorname{Ad}_u\xi$, by Step~9; and $\Phi^\natural$, $\tilde{\mathsf D}$, $n_j$ are invariant.

\item[\textup{(iii)}] Hypothesis \textup{(iii)} is the content of part~\textup{(d)} of Step~12.

\item[\textup{(iv)}] Hypothesis \textup{(iv)} is the content of Step~12.\end{enumerate}

The lemma yields that the mean is real, Borel and gauge invariant, that it depends on $V'$ only through the bonds of $\mathsf Y^{+s'}$, that it satisfies the bound, and that $\Xi'=\sum\prod w'(\mathsf Y_i,V')$ at \emph{every} $V'$.

\item[\textup{6.}] We conclude. Clause \textup{(c0)} holds because $\mathsf V'$ is real, by Step~9; clause \textup{(ca)} holds by Steps~9 and~10, clause \textup{(ch)} by Steps~9 and~11, and clause \textup{(c$\Xi$)} by item~5 above. Moreover $\rho'>0$, as a positive kernel integrated against a positive density, and $\ell'^{\rm c}$ has range $\le\varrho(x')$.

The letters of the class $C^{\rm su}_\flat$ listed in the statement of \Cref{Y-G-ThA} are the same at $x'$ as at $x$, and $\rho'$ is a function of $V'$ alone; besides the clauses just verified, these two facts are what that statement asserts of $\rho'$.\end{enumerate}

\emph{Step 14: the floors on half lines and the order of the constants.} In this step we show that each of the floors \textup{(A0)} to \textup{(A11)} holds for all large $x$ once the letters \textup{(O0)} to \textup{(O12)} are fixed; hence the set $S(c_h,E_a)$ of \textup{(O13)} is non-empty, and $x_\flat^0=1+\inf S(c_h,E_a)$ exists. We show moreover that each entry of the table in item~3 of \S\ref{Y-G0-1} depends only on entries before it, and that each condition is one-sided.

We put $\Lambda=c_R\log x/M$. Item~1 of \S\ref{Y-G0-1} defines $s'=s(x')$, where $s(x)=8r(x)$ and $r(x)=\lceil20c_R\log x/M\rceil$. Since $x'\le x$, it follows that $s'\le8+160\Lambda$, and so $e^{(\kappa+4)s'}\le e^{8(\kappa+4)}x^{1/12}$, because $(\kappa+4)160c_R/M\le1/12$ by \textup{(M2)}.

For \textup{(A0)}: the inputs state that their floors hold for all large $x$; they do so in \S\ref{Y-C-1}, Part~C.II, items~4 and~5 (Floors), in item~G (Floors) of \S\ref{Y-D-1}, in item~6 (Floors) of \S\ref{Y-E-1} with Step~14 of the proof of \Cref{Y-E-C6}, and in item~G of \S\ref{Y-F-1} with Step~13 of the proof of \Cref{Y-F-C7}. Of \textup{(P0)}, only the clauses $x\ge L^2$ and $R\ge6$ are among these floors, since \textup{(Tor)} is a datum by \textup{(A0)}.

For the conditions \textup{(D2)}, \textup{(D3)} and \textup{(C0)} to \textup{(C2)} named in \textup{(A0)}, we use $\ell_w$, $D_w$, $\omega=O(\log x)$, $\mathsf g_x\ge c(\log x)^{-2}$, and $r_\eta$, $r_T$, $\varepsilon'=O(x^{-1/2}(\log x)^{\mathfrak p_0})$. In $3D_w\le\ell_w(x')$, the leading term $360\Lambda M/L$ of the left side is smaller than the leading term $120\Lambda M$ of the right side; in $\mathsf K_{\rm c}(17T\varepsilon')^2\le0.12\varepsilon$, the left side is of the order $\varepsilon^2$ and the right side of the larger order $\varepsilon$.

For \textup{(A1)}: $\varepsilon\to0$; $\varrho(x')\to\infty$; $R/L+3\le R/2$; and in the last condition $\varepsilon^2$ on the left is of smaller order than $\varepsilon$ on the right.

For \textup{(A2)}: $r_1\le(2+3/L)\Lambda+{\rm const}$ and $2\lceil s(x)/L\rceil\le320\Lambda/L+{\rm const}$, so the left side is $\le(10+335/L)\Lambda+{\rm const}$, while $r'\ge20\Lambda-20c_R\log L/M$; thus the condition holds for large $x$ if and only if $L\ge34$, as $10+335/L$, rounded to two places, is $20.15$ at $L=33$ and $19.85$ at $L=34$. Further, for large $x$ the upper bound in $\omega\le(2+141/L)R+M+{\rm const}$ lies below the lower bound in $Mr'\ge20R-{\rm const}$, and so does the upper bound in $R_6+1\le(4+6/L)R+{\rm const}$.

For \textup{(A3)}, \textup{(A4)}, \textup{(A8)} and \textup{(A10)}, powers of $x$ dominate powers of $\log x$ for large $x$; moreover $(p_0'/p_0)\to1$, and $n_a$ is independent of $x$.

For \textup{(A5)}: $1/4+1/12-1/2=-1/6<-1/8$, and $u_b$, $u_1$ decay faster than every power.

For \textup{(A6)}: in the third term, $-1/2+1/12=-5/12<-1/4$, and $\upsilon=e^{-c_\Xi p_0^2}$ decays faster than every power. Further, $-\log u_t\ge\min\{\rho_Q/2-c_\tau,\ \mathrm{py}/(768(2r_1+1)^3)\}-\log3v$, which is of the order $(\log x)^{2\mathfrak p_0-3}$ under the floor \textup{(W)} and exceeds $(1/4+1/12)\log x$ for large $x$; it is here that we use $\mathfrak p_0>2$, which gives $2\mathfrak p_0-3>1$.

For \textup{(A7)}: $e_1^\flat=O(\varepsilon')$, and $r_\Xi^\flat p_0'^2$, with $r_\Xi^\flat\to(7/8)(1/2-d_0-c_\Xi)>0$, is of the order $(\log x)^{2\mathfrak p_0}$ and so of larger order than $(\log x)^3$. For \textup{(A9)}: $p_0^2-p_0'^2\ge2\mathfrak p_0A_0^2(\log x')^{2\mathfrak p_0-1}\log L\to\infty$. For \textup{(A11)}: $r'\to\infty$.

The order of the constants is as follows. First \textup{(O0)} to \textup{(O3)} fix $\mathsf N_Z$, $\mathsf N_\Xi$, $\mathsf N_\sharp$, $\mathsf N_\nu$, $\mathsf N_\nu^\flat$, $\mathsf N_{\rm c}$ as functions of $L$; then come $L$ and $a$; then $\kappa$; then $\gamma$, $\mathsf N_{\mathsf B}$, $b$, $A^\flat$, $\theta^\flat$, $\theta_{\mathbb C}$, $\vartheta_0^\flat$, the number $\Lambda=2+A^\flat/b$ of \textup{(O6)}, $\mu$, $\mathsf G_0$, $C_\varphi^\flat$, determined by $(L,a,K_*,\lambda_*)$; then $\theta_D$, $C_D$, $c_R$; then $\mathsf K_{\rm c}^0$ and $\mathsf a_{\rm c}$, which are independent of $M$ and of $x$; then $M$, by \textup{(M2)}, \textup{(M4$^\flat$)}, \textup{(M5r)}; then $\kappa_+$, $\mathsf K_1^\flat$; then $\mathfrak p_0$, $A_0$; then $c_\xi$, $h_*$; then $K_{\rm q}$, $K_v$, the list of item~F of \S\ref{Y-F-1}, and $\delta_3$; then $a_\infty$, $n_a$; then $c_h$; then $E_a$; and last $x_\flat$.

The conditions are lower bounds on $L$, $C_D$, $M$, $E_a$, $x_\flat$ and upper bounds on $\theta_D$, $c_h$; and $c_R$ is defined as $64/\theta_D$.

The letters $E_a$ and $1/c_h$ enter the floors only through $c_a(x)$, $c_h(x)$ and $1/E_a$; hence \Cref{Y-G-ThA} holds on half lines in $c_h$, $E_a$, $x_\flat$. In other words, for each pair $(c_h,E_a)$ taken from the two half lines, some number $x_\flat^0(c_h,E_a)$ exists such that \Cref{Y-G-ThA} holds for every $x_\flat\ge x_\flat^0(c_h,E_a)$.

Finally, $x_{\rm fl}$ is independent of $\ell$ and of $N$, by item~G of \S\ref{Y-F-1}. So the order of the constants has no circle.\end{proof}

\section{Iteration and cutoff comparison}\label{sec:compare-cutoffs}

\subsection{The first Wilson step}\label{Y-B-1}

We start the iteration by writing Wilson's weight as a member of
the class with empty rows. The free-kernel estimates at depth zero give
the data for the first step. Thm.~\ref{Y-B-C1} checks that the thresholds
in these two arguments can be chosen together and uniformly in the
volume.

\emph{Item 0: names.} We take over from items~1 to~4 of \S\ref{Y-G0-1} and from \Cref{Y-G-ThA} the following: the lattice, the cubes and the cells, and the way the grids fit; the families $\mathcal D$ and $\mathcal P$, the number $|X|_c$ and the sets $X^{+j}$; the numbers $p_0$, $\varepsilon$, $R=\varrho(x)$, $r(x)$, $s(x)$; the step $\rho\to\rho'$; the function $S_\varepsilon$ and the energy $\mathrm{En}_\ell$; the domains $\mathcal U_X(x)$, with their clauses \textup{(u1)} to \textup{(u3)} (item~1 of \S\ref{Y-G0-1}); the sets $M_A^{\rm su}(x)$ of members, with the strong cap; the densities $\rho[u;\ell]$ and the sums $\Xi_w$; the class $C^{\rm su}_\flat(x;\ell)$; the data \textup{(SD)}; the letters \textup{(O0)} to \textup{(O13)}; and the floors \textup{(A0)} to \textup{(A11)}.

By \textup{(F1)} to \textup{(F4)} we mean the formulas that \Cref{Y-G-ThA} names: those of Steps~1, 6, 9 and~12 of its proof, and, for $w'$, Step~13 of that proof, which takes the mean over $e$.

Two conventions fixed in \S\ref{Y-G0-1} are in force throughout. The first stems from item~4 and \textup{(A7)}: \textup{(A7)} is defined with $r_\Xi^\flat=\frac78(\frac12-d_0-c_\Xi)-3e_1^\flat$. The second stems from item~1, on cubes: a bond, a path or a plaquette \emph{of} $A$ has all of its sites in $A$, whereas a plaquette $p$ lies \emph{in} $A$ as soon as its base site does.

We write $\emptyset$ for the family of rows with $\mathsf H_y=\{0\}$; for it, $\mathrm{En}_\emptyset(U)=\sum_pS_\varepsilon(\phi_p)$. We write $\emptyset'$ for the family of rows that \Cref{Y-A-P5} builds from $\emptyset$; at the empty rows its group \textup{(5)} is void. We set $\ell_1:=(\emptyset')^{\rm c}$, that is, the cut of $\emptyset'$ at $\varrho(\beta/L)$, tensored with $1_{\mathfrak g}$ and dressed. Finally, $m_\Xi:=\frac12-d_0-c_\Xi$.

The hull ${\rm hull}(X)$ is a member of $\mathcal D$ that contains $X$ and has the least number of cubes among such members, an order fixing the choice; $|X|_c=\#{\rm hull}(X)$. Both are those of item~1 of \S\ref{Y-G0-1}, on cubes.

The letters $\rho$, $\rho'$, $\rho_1=\rho'$, $\rho_\beta$ are densities, and $\rho_\beta=e^{-\beta\sum_p(1-\cos\phi_p)}$ is Wilson's density at $\beta$.

Further, $\ell_w(x)=6Mr(x)+12$ is a number of bonds (item~1 of \S\ref{Y-G0-1}, at $x$). The number $D_w$ and the radius $\omega$ in \textup{(D2)} and \textup{(C0)} to \textup{(C2)} are set in item~2 of \S\ref{Y-G0-1}, and the number $\mathsf g_x$ in \textup{(A0)}.

The $k$ in $K_{{\rm q},k}$ and in \textup{(G1$_k$)}, \textup{(F8$_k$)}, \textup{(F9$_k$)} is a mark and not a level. The letter $K$ has three senses. In $0\le k\le K$, and also in the phrase `the tower $K=1$', it is the number of steps of a tower. It denotes the bound of the rows in $\mathcal R(b,K,\lambda,\tau)$, and hence in $K=K_*$ and in $K\ge0$. In $K\otimes1_{\mathfrak g}$ it is a scalar kernel. The letter $U$ is a fine field in $\mathrm{En}_\emptyset(U)$ and in $\ell_y^U$; in the clause $U=1$ of the first sentence on covariance of \Cref{Y-B-RZ}, it is the orthogonal map $U_{\mathsf t,y}$ of \Cref{Y-A-CovRows}.

\emph{Item 1.}

\begin{definition}[Wilson's system at $x$, for any grid of fine cubes]\label{Y-B-uW}
For $|t|<2$ we put \[h(t):=\sum_{n\ge2}\dfrac{(2t)^n}{n^2\binom{2n}{n}}=\dfrac{t^2}{6}+\dfrac{2t^3}{45}+\dfrac{t^4}{70}+\dots.\]

We call $X\in\mathcal P$ \emph{$s$-connected} when no splitting of $X$ into two non-empty parts separates them by a distance $>s(x)$.

We set $u_W:=(E^a_W,0,w_W)$. Here, for a single cube $Q$, $E^a_W(\{Q\},V):=x\sum_{p\in Q}h\bigl(1-\tfrac12\operatorname{tr}\partial V(p)\bigr)$, while $E^a_W(X,\cdot):=0$ for $\#X\ge2$; and \[\FitDisplay{\dimexpr\linewidth-1pt\relax}{\begin{aligned}
&w_W(X,V):=1\{X\text{ is }s\text{-connected}\}
\\
&\quad{}\,1\{\text{every cube of }X\text{ holds a plaquette of }D(V)\}
\\
&\quad{}\,\exp\Bigl\{x\sum_{Q\in X^{+r(x)}}\ \sum_{p\in Q,\,p\notin D(V)}h(1-\cos\phi_p)
\\
&\quad{}-x\sum_{p\in D(V)\cap X}\bigl[1-\cos\phi_p-d_0\varepsilon\phi_p\bigr]\Bigr\}.
\end{aligned}}\]
\end{definition}

\emph{Item 2.}

\begin{lemma}[Wilson's weight in the class at the empty rows]\label{Y-B-LW}
Let $x>1$ satisfy the following three conditions.
\begin{enumerate}
\item[\textup{(w1)}] $\varepsilon\le\frac1{40}$.
\item[\textup{(w2)}] $2.01M^3\varepsilon\le E_a$.
\item[\textup{(w3)}] \[\bigl(m_\Xi-\frac{\varepsilon^2}{24}\bigr)p_0^2\ge\kappa s(x)+\dfrac{M^3(2r(x)+1)^3p_0^4}{8x}.\]
\end{enumerate}

Conditions \textup{(w1)} and \textup{(w3)} force $x>1.5\cdot10^5$. Indeed, inserted into \textup{(w3)}, the bounds $\kappa\ge6$ of \textup{(O5)}, $s(x)=8r(x)\ge8$ and $m_\Xi<\frac12$ give $p_0^2>96$, and \textup{(w1)} gives $x=p_0^2/\varepsilon^2\ge1600p_0^2$. Consequently $\log x>1$, $\vartheta(x)=\vartheta_0^\flat(1-1/\log x)>0$ and $2-1/\log x>1$, and the domains $\mathcal U_X(x)$ are non-empty: a $G$-valued field that is tame on $Q^{\approx}$ has $|V(b)|=|V(\pi)|=1$, and $|\partial V(p)-1|=2\sin(\phi_p/2)\le\varepsilon$ at the plaquettes of $Q^{\approx}$, so that it belongs to $\mathcal U_{\{Q\}}(x)$.

Then, on every torus $\mathbb Z^3/n\mathbb Z^3$ with $M\mid n$ and for every grid of cubes, the following hold.
\begin{enumerate}
\item[\textup{(a)}] $e^{-x\sum_p(1-\cos\phi_p)}=\rho[u_W;\emptyset](V)$ at every $V$.
\item[\textup{(b)}] $u_W$ is real, lies in $M_1^{\rm su}(x)$ and has the strong cap. Hence Wilson's weight at $\beta=x$ lies in $C^{\rm su}_\flat(x;\emptyset)$, with $\mathsf V=0$ and $E^h=0$.
\item[\textup{(c)}] \[\sup_{\mathcal U_{\{Q\}}(x)}|E^a_W(\{Q\},\cdot)|\le2.01M^3p_0^4/x=(2.01M^3\varepsilon/E_a)\,c_a(x)\le0.0016\,c_a(x),\] the last inequality holding if $E_a\ge32M^3$.
\item[\textup{(d)}] The conditions \textup{(w1)} to \textup{(w3)} follow from \textup{(O12)}, from $x\ge L^2$ of \textup{(A0)}, from \textup{(A1)}, \textup{(A4)}, \textup{(A6)}, \textup{(A7)}, and from $K_{\rm q}\ge M^3$ (item~F of \S\ref{Y-D-1}); so they hold for every $x\ge Lx_\flat^0$.
\end{enumerate}

Moreover, $u_W$ is independent of each of $c_h$, $E_a$, $x_\flat$, $N$, $c_\Xi$, $\kappa$.
\end{lemma}

\emph{Item 3.}

\begin{theorem}[The first step from Wilson's action stays in the class]\label{Y-B-C1}
Assume the clauses and base statements listed in Step~0 of the
proof of \Cref{Y-G-ThA}, together with the conclusion of
\Cref{Y-A-Fr}. Choose \textup{(O0)} to \textup{(O10)} in the order
of item~3 of \S\ref{Y-G0-1}, and define
\[
 E_a^0:=\max\{32L^{-1/2}K_v,32K_{\rm q}\},\qquad
 x_\flat^{(1)}:=x_\flat^0(\delta_3/2,E_a^0).
\]
Here $x_\flat^0(c_h,E_a)=1+\inf S(c_h,E_a)$ is the threshold of
\textup{(O13)}: $S(c_h,E_a)$ consists of those $y\ge e^2$ for which
\textup{(A0)} to \textup{(A11)} hold at every $x\ge Ly$.
Step~14 of the proof of \Cref{Y-G-ThA} shows that this is a
nonempty upper half line, so $1+\inf S$ belongs to it, whether or
not its infimum does. The threshold $x_\flat^{(1)}$ therefore
depends on \textup{(O0)} to \textup{(O10)} alone.

All size conditions entering these floors are recorded in
\textup{(A0)}. In particular, it includes the floors of
\S\ref{Y-F-1} through $x\ge x_{\rm fl}$, the chart conditions
\textup{(D2)}, \textup{(D3)} at $(1,r_\eta)$ and $(17T,r_T)$,
and \textup{(C0)} to \textup{(C2)}. The volume condition
\textup{(Tor)} is a datum of the step and is excluded from these
floors, as stipulated in \textup{(A0)} and \textup{(P0)}.
Thus $x_\flat^{(1)}$ is independent of $N$ and of the grid.

\emph{Conclusion.} In this setting, let $0<c_h\le\delta_3/2$, $E_a\ge E_a^0$ and $x_\flat\ge x_\flat^{(1)}$ be arbitrary, let $\beta\ge Lx_\flat$, let $N$ be any even number with $M\mid N$ and $N\ge2Ms(\beta)$, and let any grid of cells on $\mathbb T'$ be given. Then the image $\rho'$ of $e^{-\beta\sum_p(1-\cos\phi_p)}$ on $\mathbb Z^3/NL\mathbb Z^3$ under the step at $x=\beta$ lies in $C^{\rm su}_\flat(\beta/L;\ell_1)$, with the letters $d_0$, $c_\Xi$, $\kappa$, $\rho_0$, $c_R$, $M$, $A_0$, $\mathfrak p_0$, $\vartheta_0^\flat$, $c_h$, $E_a$ and with $x'=\beta/L$ exactly. The formula of the class holds at every $V'$, and $V'$ is its only argument. One system covers every $(c_h,E_a,x_\flat)$ of the three half lines: $(\mathsf V',u')$ is built by \textup{(F1)} to \textup{(F4)} from $(\mathsf V,u,\ell)=(0,u_W,\emptyset)$, with the fine grid of origin $o=Lo'-\mathsf h_L$ mod $M$, and $w'$ is the mean over $e$ of the $w'(\mathsf Y,V',e)$ of \textup{(F4)} (Step~13 of the proof of \Cref{Y-G-ThA}). The same conclusion holds at any $(c_h,E_a)$ of the half lines for $\beta\ge Lx_\flat^0(c_h,E_a)$.

\emph{Besides.} The following also hold.
\begin{enumerate}
\item[\textup{(i$'$)}] For every cell $\square$, \[\sum_{Y\ni\square}e^{\kappa(\#Y-1)}\sup_{\mathcal U_Y(x')}|E'^a(Y,\cdot)|\le0.56\,(E_a^0/E_a)\,c_a(x').\]
\item[\textup{(ii$'$)}] For every cell $\square$, \[\sum_{Y\ni\square}e^{\kappa(\#Y-1)}\sup|E'^h(Y,\cdot)|\le e^{-(\delta_3/2)p_0(\beta/L)^2}.\]
\item[\textup{(iii$'$)}] The clauses that follow are those of the scalar rows, namely of the cut of $\emptyset'$ at $\varrho(\beta/L)$, not yet tensored with $1_{\mathfrak g}$ or dressed: indeed $K\otimes1_{\mathfrak g}$ has the norms of $K$ (Step~H13 of the proof of \Cref{Y-A-Fr}, the group), and at $G$-valued fields the transports of the dressing are isometries. The rows $\ell_1$ are translation covariant and have range $\le\varrho(\beta/L)$; further, $|\ell_{1,y'}(P)|\le K_*e^{-\lambda_*{\rm dist}(y',P)}$, $1\le\mathsf B'\le\frac32\mathsf r_\sharp^{-1}$, $\|\ell_1\|_{\rm r}\|\ell_1\|_{\rm c}\le3K_*^2\mathsf g(\lambda_*)^2$ and $\mathsf t_\rho\le\frac18m_\Xi$. These facts come from the conclusion of \Cref{Y-A-Fr}, with \textup{(Fr0)} and \textup{(Fr4)}, at $k=0$, $x_0=\beta$; the estimate of the two norms from \Cref{Y-A-SetsOfRows}, at $(K_*,\lambda_*)$; the covariance of $\emptyset'$ from \Cref{Y-A-P28}\,(a); the covariance of the cut from the proof of \Cref{Y-A-P28}\,(d); and the range from the cut.
\item[\textup{(iv$'$)}] $\theta_D\le\theta_{\mathbb C}/2$, $\theta_D\le\lambda_*/8$ and $c_R\theta_{\mathbb C}\ge128$, by \textup{(O6)} and by item~C.II~2.2 of \S\ref{Y-C-1} (source and hypotheses).
\end{enumerate}
\end{theorem}

\begin{remark}[Comparison of half lines of $c_h$]\label{Y-B-Riv}
For every number $\delta>0$ we have $\frac12\min\{\delta,\delta_3'\}\le\frac12\min\{\delta,\delta_3\}\le\delta_3/2$. Hence the half lines $c_h\le\frac12\min\{\delta,\delta_3'\}$ and $c_h\le\frac12\min\{\delta,\delta_3\}$ are contained in the half line $c_h\le\delta_3/2$ of \Cref{Y-B-C1}, whatever $\delta>0$ is. For instance, at \textup{(O11)} of item~3 of \S\ref{Y-G0-1} the table takes $c_h\le\delta_3'/2$, and $\delta_3'/2=\frac12\min\{\delta,\delta_3'\}$ for every $\delta\ge\delta_3'$.
\end{remark}

\begin{remark}[Statements for later use]
For later sections we collect the following.
\begin{enumerate}
\item[\textup{(R0)}] $\rho_0$ itself lies in $C^{\rm su}_\flat(\beta;\emptyset)$, at all $(c_h,E_a)$ of the half lines and for every $\beta\ge Lx_\flat^{(1)}$; here $\rho_0$ is Wilson's density $e^{-\beta\sum_p(1-\cos\phi_p)}$, the $\rho_0$ of \Cref{Y-GB-ThB}, and not the integer of \textup{(O3)}. Indeed, neither \textup{(w1)} nor \textup{(w3)} depends on $c_h$ or on $E_a$; \textup{(w2)} at $E_a^0$ gives \textup{(w2)} at every $E_a\ge E_a^0$; $E^h=0$ meets \textup{(ch)} at every $c_h$; and \Cref{Y-B-LW}\,(b) then applies.
\item[\textup{(R1)}] $\rho_1\in C^{\rm su}_\flat(\beta/L;\ell_1)$ at every $V'$, for every grid, where $\rho_1=\rho'$ is the image of $e^{-\beta\sum_p(1-\cos\phi_p)}$ under the step at $x=\beta$, as in the conclusion of \Cref{Y-B-C1}.
\end{enumerate}

By \textup{(R0)}, the membership $\rho_k\in C^{\rm su}_\flat(x_k;\ell_k)$ of \Cref{Y-GB-ThB} is true also at $k=0$, where $x_0=\beta$ and $\ell_0=\emptyset$; it is therefore true for $0\le k\le K$.
\end{remark}

\begin{proof}[Proof of \Cref{Y-B-LW} and of \Cref{Y-B-C1}]
\emph{Step S0: the hypotheses at Wilson's weight.}
To apply \Cref{Y-G-ThA}, we need the letters and floors, the data
\textup{(SD)}, and a real member with the strong cap representing
Wilson's density. Steps S1--S4 prove the representation and its
analytic and rough caps; Step S5 checks its three size conditions;
and Step S6 proves \textup{(SD)} for $(\emptyset,\emptyset')$.
The remaining conditions hold directly: $\mathsf V=0$, $E^h=0$,
Wilson's density is positive and gauge invariant, and the chosen
fine grid fits the given coarse grid. The hypotheses on $N$ are
those in the theorem.

The empty rows are allowed in items C.I~1.1 and C.II~1.2 of
\S\ref{Y-C-1}, in Remark~(4) of item~I of \S\ref{Y-D-1}, and in
the data of Sections~\ref{Y-E-1}--\ref{Y-G-1}. They are also
members of $\mathcal R(1,0,\lambda,0)$ by
\Cref{Y-A-SetsOfRows}. \Cref{Y-B-RZ} gives an equivalent
formulation using one-dimensional zero rows. The window law is
constructed within the step from $q^V_\emptyset$; it is not an
additional datum of the class.

\emph{Step S1: the function~$h$.} We show that $h$ has the following four properties.

\begin{enumerate}

\item[\textup{(1)}] The function $h$ is holomorphic in $|t|<2$, and its coefficients are $>0$.

\item[\textup{(2)}] For $0\le\phi<\pi$ we have $h(1-\cos\phi)=\phi^2/2-(1-\cos\phi)$.

\item[\textup{(3)}] Moreover $0\le h(1-\cos\phi)\le\phi^4/24$.

\item[\textup{(4)}] In the same disc, $|h(t)|\le\dfrac{|t|^2/6}{1-|t|/2}$.

\end{enumerate}

The proof rests on two facts: the expansion \[\arcsin^2z=\frac12\sum_{n\ge1}\dfrac{(2z)^{2n}}{n^2\binom{2n}{n}},\] and the relation $\arccos(1-t)=2\arcsin\bigl((t/2)^{1/2}\bigr)$.

To prove the first, we put $f=\arcsin^2z$. This function is even and holomorphic in $|z|<1$, and it satisfies $f(0)=0$ and the equation $(1-z^2)f''-zf'=2$, which we obtain by differentiating $(1-z^2)^{1/2}f'=2\arcsin z$.

If we write $f=\sum_{n\ge1}c_nz^{2n}$, this equation says that $2c_1=2$ and that $(2n+2)(2n+1)c_{n+1}=4n^2c_n$.

On the other hand, the numbers $b_n=\frac12\,4^n/\bigl(n^2\binom{2n}{n}\bigr)$ satisfy $b_1=1$ and $b_{n+1}/b_n=4n^2/((2n+2)(2n+1))$, since $\binom{2n+2}{n+1}/\binom{2n}{n}=(2n+2)(2n+1)/(n+1)^2$. Therefore $c_n=b_n$.

The second fact is the identity $\cos2a=1-2\sin^2a$ at $\sin a=(t/2)^{1/2}$, where $0\le t\le2$.

From the two facts together we obtain \[\frac12\arccos(1-t)^2=\sum_{n\ge1}\dfrac{(2t)^n}{n^2\binom{2n}{n}}=t+h(t).\]

In this series, the quotient of two consecutive coefficients equals $n^2/((n+1)(2n+1))$, which is $<\frac12$. Hence the radius of convergence is $2$, and the tail of the series from $n=2$ on is bounded in absolute value by $|t|^2/6$, the absolute value of its first term, times $1/(1-|t|/2)$.

The bound $\phi^4/24$ comes from the inequality $\cos\phi\le1-\phi^2/2+\phi^4/24$, which is valid for every real $\phi$.

\emph{Step S2: clause \textup{(ca)} for $E^a_W$, in the proof of \Cref{Y-B-LW}\,(c).} We show that $E^a_W(\{Q\},\cdot)$ is holomorphic on $\mathcal U_{\{Q\}}(x)$ and real at $G$-valued fields, that it is invariant under $G$-valued gauge transformations and vanishes at pure gauges, and that its supremum is bounded by $2.01M^3p_0^4/x$, a number that is itself $\le c_a(x)$.

Let $p$ be a plaquette based in $Q$. All its sites lie in $Q^{+1}$, and $Q^{+1}$ is contained in $Q^{+r(x)}$ because $r(x)\ge1$. Thus $p$ is a plaquette of $Q^{\approx}$, and \textup{(u3)} yields $|\partial V(p)-1|<2\varepsilon$.

We set $g=\partial V(p)=1+A\in\operatorname{SL}(2,\mathbb C)$. This uses that $\mathcal U_X(x)$ is a set of $\operatorname{SL}(2,\mathbb C)$ fields (item~1 of \S\ref{Y-G0-1}, on domains); for a cell, the domain of item~C6 of \S\ref{Y-D-1} at $x'$ is a set of $\operatorname{SL}(2,\mathbb C)$ fields as well.

Because $\det g=1$, we have $\operatorname{tr}A=-\det A$, and therefore $t_p:=1-\frac12\operatorname{tr}g=\frac12\det A$. Without $\det g=1$, $t_p$ would be of first order in $A$.

Since $|\det A|$ is the product of the two singular values of $A$, it follows that $|t_p|\le\frac12|A|^2<2\varepsilon^2\le\frac1{800}$.

The number $t_p$ is a polynomial in the entries of $V$, so $h(t_p)$ is holomorphic. Under gauge transformations the trace does not change. At $G$-valued fields, $t_p=1-\cos\phi_p$ is real. At $\partial V=1$ we get $t_p=0$, and $h(0)=0$.

Step~S1 gives $|h(t_p)|\le\dfrac{4\varepsilon^4/6}{1-\varepsilon^2}\le0.668\,\varepsilon^4$. As a cube holds $3M^3$ plaquettes, $|E^a_W|\le x\cdot3M^3\cdot0.668\,\varepsilon^4\le2.01M^3p_0^4/x$.

Since $c_a(x)=E_ap_0^3x^{-1/2}$, dividing this bound by $c_a(x)$ gives $2.01M^3\varepsilon/E_a$, and by \textup{(w2)} this quotient is $\le1$. With $E_a\ge32K_{\rm q}\ge32M^3$ and $\varepsilon\le\frac1{40}$, it is even $\le0.0016$.

\emph{Step S3: the identity, in the proof of \Cref{Y-B-LW}\,(a).} We show that at every $G$-valued $V$ \[e^{-x\sum_p(1-\cos\phi_p)}=\exp\Bigl\{-x\sum_pS_\varepsilon(\phi_p)+\sum_Q1_{Q^{\approx}}(V)E^a_W(\{Q\},V)\Bigr\}\,\Xi_{w_W}(V).\]

We denote by $d(V)$ the set of those cubes that hold a plaquette of $D(V)$. In this step and the next we abbreviate $D=D(V)$, $r=r(x)$ and $s=s(x)$, and under a sum over $p$ we write $h$ for $h(1-\cos\phi_p)$.

\begin{enumerate}

\item[\textup{1.}] \emph{The quotient.} If $p\notin D$, then $-x(1-\cos\phi_p)+x\phi_p^2/2=x\,h(1-\cos\phi_p)$. If $p\in D$, the corresponding expression is $-x(1-\cos\phi_p)+xd_0\varepsilon\phi_p$. Dividing the left side by $e^{-x\sum S_\varepsilon}$, we therefore obtain $\exp\bigl\{x\sum_{p\notin D}h-x\sum_{p\in D}[1-\cos\phi_p-d_0\varepsilon\phi_p]\bigr\}$.

\item[\textup{2.}] \emph{The indicator.} The equality $1_{Q^{\approx}}(V)=1$ holds exactly when no plaquette based in $Q^{+r}$ belongs to $D$, that is, exactly when ${\rm dist}(Q,d(V))>r$. Hence $\sum_Q1_{Q^{\approx}}E^a_W=x\sum_{Q\notin d(V)^{+r}}\sum_{p\in Q}h$, and each $p$ occurring in this sum is tame. The factor that remains for $\Xi$ is thus $\exp\bigl\{x\sum_{Q\in d(V)^{+r}}\sum_{p\in Q,\,p\notin D}h-x\sum_{p\in D}[\,\cdot\,]\bigr\}$, where $[\,\cdot\,]$ is the bracket of part~1.

\item[\textup{3.}] \emph{The families.} Any family $\{X_i\}$ in $\Xi_w$ has pairwise distances $>s$, and $D$ lies inside its union, that is, each plaquette of $D$ is in some cube of the union. When its product is not $0$, every cube of every $X_i$ lies in $d(V)$, and hence the union is exactly $d(V)$. Every $X_i$ is $s$-connected, while two different $X_i$ are at distance $>s$; so the $X_i$ are the classes of $d(V)$ under ``distance $\le s$''. There is thus exactly one such family, and when $D$ is empty this is the empty family, whose value is $1$.

\item[\textup{4.}] \emph{The product.} As $s=8r>2r$, the sets $X_i^{+r}$ are pairwise disjoint, and their union is $d(V)^{+r}$; the sets $D\cap X_i$ form a partition of $D$. Consequently $\prod_iw_W(X_i,V)$ equals the factor that remains in part~2.

\end{enumerate}

\emph{Step S4: clause \textup{(c$\Xi$)} with the strong cap, in the proof of \Cref{Y-B-LW}\,(b).} We show that $w_W(X,\cdot)$ is real, Borel and gauge invariant, that it is a function of $V$ on the bonds of $X^{+s(x)}$, and that at every $V$ \[|w_W(X,V)|\le\exp\bigl\{-c_\Xi x\varepsilon\sum_{D(V)\cap X}\phi_p-\kappa(|X|_c-1)\bigr\}.\]

The function $w_W$ depends only on the angles $\phi_p$ at the plaquettes based in $X^{+r}$. The sites of these plaquettes lie in $X^{+r+1}$, and this set is contained in $X^{+s}$ because $r+1\le8r$. Each $\phi_p$ is a continuous class function of $\partial V(p)$.

\emph{The bound.} Suppose $w_W\ne0$, and put $n=\#X$.

\begin{enumerate}

\item[\textup{(1)}] \emph{The rough plaquettes.} The function $(1-\cos\phi)/\phi$ has the derivative $(\phi\sin\phi-1+\cos\phi)/\phi^2$. The numerator $N(\phi)$ of this derivative satisfies $N(0)=0$ and $N'=\phi\cos\phi$; so $N$ increases on $(0,\pi/2)$, then decreases, and $N(\pi)=-2$. Therefore the function first increases and then decreases on $(0,\pi]$, and on $[\varepsilon,\pi]$ it is bounded below by the smaller of its values at the two end points. These values satisfy $(1-\cos\varepsilon)/\varepsilon\ge\varepsilon/2-\varepsilon^3/24$ and $2/\pi\ge\varepsilon/2$. For $p\in D$ we thus get $1-\cos\phi_p\ge(\frac12-\frac{\varepsilon^2}{24})\varepsilon\phi_p$, and hence $1-\cos\phi_p-d_0\varepsilon\phi_p\ge(c_\Xi+m_\Xi-\frac{\varepsilon^2}{24})\varepsilon\phi_p$.

\item[\textup{(2)}] Each cube of $X$ holds some $p\in D$, and for this $p$ we have $x\varepsilon\phi_p>x\varepsilon^2=p_0^2$. Summing, $x\varepsilon\sum_{D\cap X}\phi_p\ge np_0^2$.

\item[\textup{(3)}] We have $\#X^{+r}\le n(2r+1)^3$; a cube holds $3M^3$ plaquettes; and at tame $p$ we have $h\le\varepsilon^4/24$. The positive part of the exponent is therefore at most $n(2r+1)^3\,3M^3\,x\varepsilon^4/24=nM^3(2r+1)^3p_0^4/(8x)$.

\item[\textup{(4)}] \emph{The hull.} For $X\in\mathcal P$, ${\rm hull}(X)$ is a member of $\mathcal D$ that contains $X$ and has the least number of cubes among such members, the choice being fixed by an order; and $|X|_c:=\#{\rm hull}(X)$. Consequently, every connected set of cubes that contains $X$ has at least $|X|_c$ cubes. The set $X$ has a spanning tree whose edges have distance $\le s$, and for two cubes at sup distance $d$ there are $d-1$ further cubes that link them through corners. Taking the $n$ cubes of $X$ together with, for each of the $n-1$ edges of the tree, its $d-1\le s-1$ further cubes, we obtain a connected set of cubes that contains $X$; this is true on a torus as well. Hence $|X|_c\le n+(n-1)(s-1)$, that is, $|X|_c-1\le(n-1)s$.

\item[\textup{(5)}] \emph{Conclusion.} The number $m_\Xi-\varepsilon^2/24$ is $>0$, since the right side of \textup{(w3)} is $>0$ and $p_0^2>0$; thus we may multiply the inequality of \textup{(2)} by this number. We arrive at \[\log w_W\le-c_\Xi x\varepsilon\sum\phi_p-n\Bigl[(m_\Xi-\frac{\varepsilon^2}{24})p_0^2-\frac{M^3(2r+1)^3p_0^4}{8x}\Bigr]\le-c_\Xi x\varepsilon\sum\phi_p-n\kappa s,\] where the second inequality follows from \textup{(w3)}; and $n\kappa s\ge\kappa(|X|_c-1)$.

\end{enumerate}

\emph{Step S5: the floors, in the proof of \Cref{Y-B-LW}\,(d).} We show that, for any $c_h\le\delta_3/2$ and any $E_a\ge E_a^0$, the three conditions \textup{(w1)}, \textup{(w2)}, \textup{(w3)} hold at every $x\ge Lx_\flat^0$.

\begin{enumerate}

\item[\textup{(w1)}] This condition is the first clause of \textup{(A1)}.

\item[\textup{(w2)}] By \textup{(O12)} we have $E_a\ge32K_{\rm q}$, and $K_{\rm q}=K_{{\rm q},k}=M^3[3.09L^2K_Y+520LC_D(1+4C_D^2)+1]\ge M^3$ (item~F of \S\ref{Y-D-1}). Hence $2.01M^3\varepsilon\le M^3\le E_a$.

\item[\textup{(w3)}] We proceed in four parts.

\begin{enumerate}

\item[\textup{(i)}] Condition \textup{(A7)} gives $\upsilon_f^{1/2}\le\frac14x^{-1/4}e^{-(\kappa+1)s'}$. Taking logarithms, we find \[r_\Xi^\flat p_0'^2\ge2(\kappa+1)s'+(\kappa+4)(1+2s')+4+1.01\mathsf K_1^\flat(2r'+1)^3+2\log4\ge(4\kappa+10)s'+\kappa+8.\] Here $r_\Xi^\flat=\frac78m_\Xi-3e_1^\flat\le m_\Xi$. The signs we use are these: $e_1^\flat>0$ and $u_t>0$ (item~F of \S\ref{Y-F-1}); and $\mathsf K_1^\flat>0$ (Part~C.II of \S\ref{Y-C-1}, in the display before \Cref{Y-C-CIITh73}). The sign of $\mathsf K_1^\flat$ enters the displayed line of logarithms, while $u_t>0$, $\upsilon>0$ and $c_h(x)>0$ are used in \textup{(iii)}, where two terms of \textup{(A6)} are dropped. Since moreover $p_0'\le p_0$, it follows that $m_\Xi p_0^2\ge(4\kappa+10)s'+\kappa+8$.

\item[\textup{(ii)}] By \textup{(A0)}, $x\ge L^2$, and this gives $\log x\le2\log x'$. So $r(x)=\lceil20c_R\log x/M\rceil\le\lceil2(20c_R\log x'/M)\rceil\le2r'$, and $s(x)\le2s'$. Hence $\kappa s(x)\le2\kappa s'$.

\item[\textup{(iii)}] From the third term of \textup{(A6)} we have $6((2r_1+1)L+2r(x))^3c_a(x)\le\frac12$, so that $(2r(x)+1)^3c_a(x)\le\frac1{12}$. Also $c_a(x)=E_ap_0^2\varepsilon\ge32M^3p_0^2\varepsilon$. Hence $M^3(2r(x)+1)^3p_0^4/(8x)=M^3(2r(x)+1)^3p_0^2\varepsilon\,(\varepsilon/8)\le\frac1{384}\cdot\frac1{320}$.

\item[\textup{(iv)}] Condition \textup{(A4)} gives $32M^3p_0^2\varepsilon\le\frac12$, and so $\varepsilon^2p_0^2/24\le\frac1{64}\cdot\frac1{40}\cdot\frac1{24}$.

\end{enumerate}

Adding these estimates, we find that the left side of \textup{(w3)} minus its right side is $\ge(2\kappa+10)s'+\kappa+8-10^{-4}>0$.

\end{enumerate}

\emph{Step S6: the data \textup{(SD)} for
$(\emptyset,\emptyset')$.}
Apply \Cref{Y-A-Fr} with $x_0=\beta$, $k=0$ and the tower $K=1$,
whose tori have sides $NL$ and $N$. Its hypotheses hold:
\textup{(O4)} gives odd $L\ge L_A$ and $aL\ge3/c_L$; $N$ is
even; and $x_1=\beta/L\ge x_\flat\ge e^2$. The condition
\textup{(X12)} holds for every $x_\flat\ge2$, by Step~H16 of
that theorem's proof. Its conclusions give
\[
 \emptyset\in\mathcal R(b_*,K_*,\lambda_*,1/9),\qquad
 \operatorname{range}(\emptyset)=0,
\]
\textup{(Dec)}$(C_D,\theta_D)$ for all sixteen kernels,
\textup{(N$'$)} with \eqref{eq:row-cut-norms} at every cut and
mask, and $\mathsf t_\rho\le m_\Xi/8$, through
\textup{(Fr0)}, \textup{(Fr2)}, \textup{(Fr3$'$)} and
\textup{(Fr4)}. The value of $C_D$ in \textup{(O6)} is at least
the value in that theorem, so it preserves these estimates.

The remaining clauses \textup{(L1$^{\rm su}$)} and
\textup{(L2$^{\rm su}$)} are imposed by \textup{(O4)}. They also
imply the older \textup{(L1)}, \textup{(L2)} of \Cref{Y-A-Np}:
since $\mathsf N_\nu\ge\mathsf N_\sharp$, $2200\ge60$ and
$40\ge4$,
\[
 \mathsf c_{\rm loss}
 \le[\mathsf N_\nu/L+2200\mathsf N_Z+40\mathsf N_\Xi]L^{-1/2}
 \le d_0/4.
\]
Also $L\ge(1.2\cdot10^4)^2\mathsf N_{\rm c}^4$, with
$\mathsf N_{\rm c}=17(\mathsf N_\sharp+1)$, implies
$L\ge115\mathsf N_\sharp^2$. Thus \textup{(SD)} holds.

For later use, \textup{(Fr0)} also gives
$|\emptyset'_{y'}(P)|\le K_*e^{-\lambda_*{\rm dist}(y',P)}$
and $1\le\mathsf B'\le\frac32\mathsf r_\sharp^{-1}$.
Translation covariance follows from \Cref{Y-A-P28}(a), and
cutting by distance preserves it by the proof of
\Cref{Y-A-P28}(d). Tensoring a scalar kernel with
$1_{\mathfrak g}$ preserves its norms; dressing at a real field
uses isometries. These are the tensoring and dressing operations
of item C.II~2.2 of \S\ref{Y-C-1} and item~3.5 of
\S\ref{Y-E-1}.

\begin{remark}[Empty rows as zero rows]\label{Y-B-RZ}
Let $0$ denote the rows with $\mathsf H_y=\mathbb R$ and
$\ell_y(p)=0$. This family is translation covariant, with $U=1$;
it has $b_{\mathsf B}=1$, $\tau=0$ and range $0$, and hence
belongs to $\mathcal R(1,0,\lambda_*,0)$.

Adding a zero component to a row does not change its column
norms or its Gram matrix. Thus $0$ and $\emptyset$ give the same
$\mathsf B$, $\mathsf R$, $q^V_\ell$, $\mathsf g^V$ and
$\mathrm{En}_\ell$; their row norms vanish, as do
$\sigma=\ell\nu$, $|\sigma_y|$ and $\ell^{\mathsf T}\sigma$.
The transported row $(\ell^VF)_y$ and every step of kind
\textup{(R)} vanish as well. The auxiliary output spaces
$\bigoplus_y\mathsf H_y$ and $\mathrm{Hil}$ may differ, but the
extra coordinates are zero and contribute no norm or energy.

By group~\textup{(5)} of \Cref{Y-A-P5}, $0'$ is
$\emptyset'$ with $L^3$ additional zero components. Their
column norms and Gram matrices therefore agree, as do
$\mathsf B'$, $H^\sharp$, $\mathsf N^\sharp$, $\mathsf T$,
$\mathsf X$, the cut forms, $\mathsf t_\rho$ and
$\mathrm{En}'$. It follows that \textup{(SD)} holds for
$(0,0')$ exactly when it holds for $(\emptyset,\emptyset')$,
and that the densities and classes at both scales agree.
Consequently either family yields the same hypothesis and
conclusion in \Cref{Y-G-ThA}.
\end{remark}

\emph{Step S7: the proof of \Cref{Y-B-C1}, with \textup{(i$'$)} to \textup{(iv$'$)}.} We now prove the theorem of item~3 of \S\ref{Y-B-1}.

\begin{enumerate}

\item[\textup{1.}] We fix the letters \textup{(O0)} to \textup{(O10)}. We set $(c_h,E_a)=(\delta_3/2,E_a^0)$, both of them computed from \textup{(O0)} to \textup{(O10)}, and $x_\flat^{(1)}=x_\flat^0$ at these two values (Step~14 of the proof of \Cref{Y-G-ThA}). Let $\beta\ge Lx_\flat^{(1)}$, let $N$ be as in the statement, and let $o'$ be the origin of an arbitrary grid of cells.

\item[\textup{2.}] We choose the grid of fine cubes with origin $o=Lo'-\mathsf h_L\bmod M$; with this choice the grids fit.

\item[\textup{3.}] Wilson's density $\rho_\beta=e^{-\beta\sum_p(1-\cos\phi_p)}$ at $\beta$ is the left side of the identity of Step~S3 at $x=\beta$, and \Cref{Y-B-LW}\,(a), taken at $x=\beta$, identifies it with $\rho[u_W;\emptyset]$. Step~S5 and \Cref{Y-B-LW} applied with $x=\beta$ show that $\rho_\beta$ lies in $C^{\rm su}_\flat(\beta;\emptyset)$.

\item[\textup{4.}] Step~S6 shows that \textup{(SD)} holds at $x=\beta$ for $(\emptyset,\emptyset')$.

\item[\textup{5.}] \Cref{Y-G-ThA} then yields the following: the formulas \textup{(F1)} to \textup{(F4)} define a real $\mathsf V'$ and a real $u'\in M_1^{\rm su}(x')$ with the strong cap at $x'=\beta/L$. Moreover $\rho'=e^{\mathsf V'}\rho[u';(\emptyset')^{\rm c}]>0$ at every $V'$, and $\rho'$ has no argument other than $V'$.

\item[\textup{6.}] \emph{Half lines.} Let $c_h\le\delta_3/2$ and $E_a\ge E_a^0$. In the definition of the class, $E_a$ occurs only in $c_a$, and $c_h$ only in $c_h(x)$. Now $c_a(x')$ increases with $E_a$, and $e^{-c_hp_0'^2}$ increases when $c_h$ decreases; hence the same $(\mathsf V',u')$ satisfies \textup{(ca)} and \textup{(ch)} at $(c_h,E_a)$. The clause ``$\mathsf V$ real'' of the class, the clause \textup{(c$\Xi$)} and the formula of the class depend on neither letter. A larger $x_\flat$ merely shrinks the set of $\beta$.

\item[\textup{7.}] \emph{One system.} The system $u_W$ depends on none of $c_h$, $E_a$ and $x_\flat$ (item~1 of \S\ref{Y-B-1}). The formulas \textup{(F1)} to \textup{(F4)} depend on letters, clusters and cores, but not on $c_h$, $E_a$ or $x_\flat$; and neither does the mean over $e$ in Step~13 of the proof of \Cref{Y-G-ThA}.

\item[\textup{8.}] For \textup{(i$'$)} and \textup{(ii$'$)}: we apply the part ``Moreover'' of \Cref{Y-G-ThA} with $(c_h,E_a)=(\delta_3/2,E_a^0)$. For every cell, this part bounds the sum on the left of \textup{(i$'$)} by $(0.25+0.31)\,E_a^0p_0'^3x'^{-1/2}$, which equals $(0.25+0.31)(E_a^0/E_a)c_a(x')$. Its bound for the sum on the left of \textup{(ii$'$)}, at the same cell, is $e^{-(\delta_3/2)p_0'^2}$.

\item[\textup{9.}] For \textup{(iii$'$)} and \textup{(iv$'$)} we use Step~S6, \textup{(O6)}, and item~C.II~2.2 of \S\ref{Y-C-1} (in its remark ``Source and hypotheses'').

\item[\textup{10.}] For \Cref{Y-B-Riv} we recall from \textup{(O10)} that $\delta_3=\min\{d_0/1536,\,3c_\xi^2/(256h_*^2)\}$ and $\delta_3'=\min\{d_0/2304,\,3c_\xi^2/(256h_*^2)\}$. Since $d_0/2304<d_0/1536$, while the other entry is the same in both, $\delta_3'\le\delta_3$. A minimum increases with each of its entries and does not exceed any of them; therefore $\min\{\delta,\delta_3'\}\le\min\{\delta,\delta_3\}\le\delta_3$.

\item[\textup{11.}] The last clause of the theorem, where $\beta\ge Lx_\flat^0(c_h,E_a)$, follows from parts~2 to~5 of this step, carried out with the values $(c_h,E_a)$ themselves.

\end{enumerate}

\emph{Step S8: order of the constants and dependencies.}
The numbers $\delta_3$, $K_v$ and $K_{\rm q}$ are fixed at
\textup{(O10)}. They determine $E_a^0$ and
$x_\flat^{(1)}$ before $c_h$, $E_a$ and $x_\flat$ are chosen
in \textup{(O11)}--\textup{(O13)}.

The proof uses \Cref{Y-A-Fr} and \Cref{Y-G-ThA}. The latter
uses Sections~\ref{Y-C-1}--\ref{Y-F-1} and the base statements
listed in Step~0 of its proof, all independently of the present
theorem. Its use of \S\ref{Y-H-1} is limited to the chart
lemmas \Cref{Y-H-L51} and \Cref{Y-H-Th53}, Step~2(ii),
Step~8(8.2) and its paragraph on $u$, and \Cref{Y-H-L52} in
Step~10(9) of the proof of \Cref{Y-H-C9}. Those passages use
the definitions and constants of \S\ref{Y-G0-1}, the results of
\S\ref{Y-C-1}, and the chart identity from Step~4 of the proof
of \Cref{Y-E-C6}; they do not use \Cref{Y-B-C1} or
\Cref{Y-GB-ThB}. The references in item~2 of \S\ref{Y-H-1}
to Step~S6 and \Cref{Y-B-RZ} concern only empty-row data,
which those two passages derive from \S\ref{Y-A-1} and the
definitions. Thus the dependency order is the free row estimates,
the real step, the present first step, and then the iteration
\Cref{Y-GB-ThB}.
\end{proof}

\subsection{Iteration of the real step}\label{Y-GB-1}

We fix the constants for the initial step and the iteration. The
side condition at the last step implies the conditions at all earlier
steps. The row estimates and the fitting of the grids then permit an
induction from the Wilson weight through all $K$ steps.

\begin{theorem}[The $K$ steps from Wilson's weight stay in the class, \textup{(S-K)}]\label{Y-GB-ThB} In the setting \textup{(O0)} to \textup{(O10)} of item~3 of \S\ref{Y-G0-1}, let $0<c_h\le\delta_3/2$ and $E_a\ge E_a^0$, and let $x_\flat\ge\max\{x_\flat^{(1)},x_\flat^0\}$ satisfy \textup{(X12)}, which is item~(8) in the order of the constants of \Cref{Y-A-Fr}. Here $x_\flat^{(1)}$ denotes the floor of \Cref{Y-B-C1}, namely $x_\flat^0$ evaluated at $(\delta_3/2,E_a^0)$, while $x_\flat^0$ itself is evaluated at the chosen pair $(c_h,E_a)$. Let $K\ge1$ and $\beta\ge L^Kx_\flat$, and let $N$ be even, with $M\mid N$ and $N\ge2Ms(L^{1-K}\beta)$. For $0\le k\le K$ let $\mathbb T_k=\mathbb Z^3/NL^{K-k}\mathbb Z^3$, with the grids $o_K=0$ and $o_k=Lo_{k+1}-\mathsf h_L$, and put $x_k=L^{-k}\beta$. Let $\rho_0=e^{-\beta\sum_p(1-\cos\phi_p)}$ be Wilson's density on $\mathbb T_0$, and for $0\le k<K$ let $\rho_{k+1}=(\rho_k)'$ be the image of $\rho_k$ under the step at $x=x_k$, in the sense of item~1 of \S\ref{Y-G0-1}, with the noise $ax_k$, by which we mean the index $ax$ of the kernel of that step.

Then, given \Cref{Y-B-C1} and \Cref{Y-G-ThA}, we have $\rho_k\in C^{\rm su}_\flat(x_k;\ell_k)$ for $1\le k\le K$, where the rows are $\ell_0=\emptyset$ and $\ell_{k+1}=((\ell_k)')^{\rm c}$, cut at $\varrho(x_{k+1})$; these are the scalar rows of \S\ref{Y-A-1}, tensored with $1_{\mathfrak g}$ and dressed. This holds with the same letters at every $k$, and at every field of $\mathbb T_k$. No argument other than the field of level $k$ enters $\rho_k$.

We make the following additional assertion: for $1\le k\le K$ we have $\rho_k=e^{\mathsf V_k}\rho[u_k;\ell_k]$, where $u_1$ is a system of members given by \Cref{Y-B-C1} and, for $k<K$, $u_{k+1}$ is made from $u_k$ by \textup{(F1)} to \textup{(F4)}. The map of members so obtained is not the map $F_\ell$ of \Cref{Y-H-C9}: the formulas \textup{(F1)} to \textup{(F4)} and the formulas \textup{(Fm1)} to \textup{(Fm4)} of $F_\ell$ are defined with two different tests $\chi_X$, so that their means $K(\mathfrak B)$ need not agree, and neither need the members $E'^a$, $E'^h$, $w'$ built with these means; however, both the map given by \textup{(F1)} to \textup{(F4)} and the map $F_\ell$ yield members that represent the same $\rho'$. The members are not unique. \end{theorem}

\begin{proof} \emph{Step 15: the tower of $K$ steps.} \emph{Item 1: the letters.} \Cref{Y-B-C1} holds for all $0<c_h\le\delta_3/2$, all $E_a\ge E_a^0$ and all $x_\flat\ge x_\flat^{(1)}$, and the number $x_\flat^{(1)}$ depends on none of $N$, $c_h$, $E_a$; this is true for any $L\ge L_*^{\rm su}$ and for any $M$ that satisfies the bounds of \textup{(O7)}. In turn, \Cref{Y-G-ThA} is valid on the half lines of Step~14 of its proof. Intersecting the two ranges gives the choice of \Cref{Y-GB-ThB}, and this is one choice of all the letters.

\emph{Item 2: the condition \textup{(Tor)}.} We begin with a proposition on numbers.

\begin{proposition}[The condition on the side passes to finer levels]\label{Y-GB-P63} If $x\ge e^2$, $L\ge3$ and $N\ge2Ms(x)$, then $NL^j\ge2Ms(L^jx)$ for every integer $j\ge0$. \end{proposition}

\begin{proof}[Proof of \Cref{Y-GB-P63}] Put $\alpha=20c_R/M$ and $t=\alpha\log x$, so that $r(x)=\lceil t\rceil$. Then \[\alpha\log(L^jx)=t(1+j\log L/\log x)\le t(1+(j/2)\log L)\le L^jt\le L^j\lceil t\rceil\] where the first inequality uses $\log x\ge2$ and the second uses $1+v/2\le e^v$ at $v=j\log L$. Since $r(L^jx)=\lceil\alpha\log(L^jx)\rceil$ and $L^j\lceil t\rceil$ is an integer, this gives $r(L^jx)\le L^jr(x)$. Consequently \[2Ms(L^jx)=16Mr(L^jx)\le L^j\cdot2Ms(x)\le L^jN\] which proves the proposition. \end{proof}

We now apply the proposition. The coarse side of step $k$ is $N_k=NL^{K-k-1}$, which is even and a multiple of $M$. Moreover $x_k=L^{K-1-k}x_{K-1}$, and $x_{K-1}=L^{1-K}\beta\ge Lx_\flat\ge e^2$. \Cref{Y-GB-P63}, applied with $x=x_{K-1}$ and $j=K-1-k$, therefore turns \textup{(Tor)} at the last step, which is the condition $N\ge2Ms(L^{1-K}\beta)$ of the data, into \textup{(Tor)} at every step, namely $N_k\ge2Ms(x_k)$.

\emph{Item 3: the grids.} The grids of the data obey $o_k=Lo_{k+1}-\mathsf h_L$, and this relation is the fit of the grids that \textup{(SD)} requires at every step.

\emph{Item 4: the level $k=0$.} \Cref{Y-B-C1}, applied with $\beta\ge Lx_\flat$, gives $\rho_1\in C^{\rm su}_\flat(\beta/L;\ell_1)$. Of that theorem we use here that it holds at every field of $\mathbb T_1$ and for every grid of cells, hence in particular for the grid $o_1$. The condition $N_0\ge2Ms(\beta)$ of that theorem is \textup{(Tor)} of item~2 for $k=0$. The rows are $\ell_1=(\emptyset')^{\rm c}$, cut at $\varrho(\beta/L)$.

\emph{Item 5: from $k$ to $k+1$, for $1\le k<K$.} Here $x_k\ge Lx_\flat$, and $\ell_k$ has range $\le\varrho(x_k)$. We use \Cref{Y-A-Fr} for this $k$, with $x_{k+1}\ge x_\flat$, with \textup{(X12)}, with $N_k$ even, with $L\ge L_A$ odd, and with $aL\ge3/c_L$ from \textup{(O4)}. The conclusion of that theorem gives \textup{(Dec)}$(C_D,\theta_D)$ for all sixteen kernels, with $\theta_D=\frac12\min\{\lambda_2^*,\theta_{\mathbb C},\theta_J\}$ and with the $C_D$ of \textup{(O6)}. Indeed, in that conclusion $C_D$ stands for the expression displayed in item~(6) of its order of the constants; since \textup{(Dec)}$(C,\theta_D)$ holds for every $C$ at least as large as this expression, it holds for the letter $C_D$ of \textup{(O6)}. The same conclusion yields \textup{(Fr3)} at every cut.

Two steps of the proof of \Cref{Y-A-Fr} are used as well: Step~H9, which is the membership \textup{(Fr0)} of its conclusion, and Step~H14, which, with the present letters, applies \Cref{Y-A-Th29}\,(b) and proves \Cref{Y-A-Th36} in the form that \Cref{Y-A-Fr} uses. For the scalar rows $\ell_k$ of \S\ref{Y-A-1}, before they are tensored with $1_{\mathfrak g}$, the two steps give $\ell_k\in \mathcal R(b_*,K_*,\lambda_*,1/9)$, then \textup{(N$'$)} with \eqref{eq:row-cut-norms} at every cut, and finally $\mathsf t_\rho\le\frac18(\frac12-d_0-c_\Xi)$. All of these are statements on scalar kernels. For a scalar kernel $K$, the kernel $K\otimes1_{\mathfrak g}$ has the same norms as $K$, by item C.II 2.2, ``Source and hypotheses'', of \S\ref{Y-C-1}; hence these statements hold for the rows $\ell\otimes1_{\mathfrak g}$ as well. Therefore $\ell_k\in \mathcal R(b_*,K_*,\lambda_*,1/9)$, and the pair $(\ell_k,(\ell_k)')$ has the four conclusions named in \textup{(SD)}, with \textup{(L1)} and \textup{(L2)} in the forms \textup{(L1$^{\rm su}$)} and \textup{(L2$^{\rm su}$)} by \textup{(O4)}. Together with items~2 and~3, this is \textup{(SD)} at step $k$; it rests on the conclusion of \Cref{Y-A-Fr}, namely on \textup{(Fr0)}, \textup{(Fr2)}, \textup{(Fr3$'$)} and \textup{(Fr4)}, at every $k$ with $x_{k+1}\ge x_\flat$. \Cref{Y-G-ThA} now carries $\rho_k\in C^{\rm su}_\flat(x_k;\ell_k)$ over to $\rho_{k+1}\in C^{\rm su}_\flat(x_{k+1};\ell_{k+1})$.

\emph{Item 6: the members.} The proof of \Cref{Y-G-ThA} constructs the members of $\rho_{k+1}$ from those of $\rho_k$ by \textup{(F1)} to \textup{(F4)}. \end{proof}

\subsection{The complex step}\label{Y-H-1}

The cutoff comparison needs a step defined on complex members
with only a weak cap on the rough term. We use a different test
$\chi_X$ from the real step, which gives a different representation of
the same coarse density. Thm.~\ref{Y-H-C9} proves holomorphic dependence
and restores the strong cap at the output. Cauchy's estimate then
contracts the weak distance between inputs into the strong distance
between outputs.

\begin{remark}[Item 0: signs]
The signs $D(U)$ and $D'$ are sets of rough plaquettes. We call a map $G$-valued when its values lie in $\operatorname{SU}(2)$. Throughout Step~0\,(i) of the proof of \Cref{Y-H-C9}, $c>0$ denotes a number determined by the letters \textup{(O0)} to \textup{(O12)} and independent of $x$, and the value of $c$ may change from one occurrence to the next. On $\mathfrak g_{\mathbb C}$ we write $|\cdot|$ for the operator norm and $\|\cdot\|$ for the norm of $\mathbb C^3$; these obey $\|X\|\le|X|\le\sqrt2\,\|X\|$, and $|X|=\|X\|$ on $\mathfrak g$.
\end{remark}

\begin{definition}[Item 1: the letters, in the order of the table]
The letters are those of the table of item~3 of \S\ref{Y-G0-1}, in its order \textup{(O0)} to \textup{(O13)}.
\begin{enumerate}
\item[\textup{(O0)}] The pure numbers are $b_*,K_*,\lambda_*,L_A,L_B,c_1^*,c_2^*,\lambda_2^*,\theta_J$. All of them but $L_B$ are fixed in \Cref{Y-A-Fr}, in item~\textup{(0)} of its part ``The order of the constants'', and $L_B$ is the pure number defined in item~1, \textup{(O0)}, of \S\ref{Y-I-1}. We set $\mathsf g(\lambda)=\sum_{n\in\mathbb Z^3}e^{-\lambda|n|_\infty}$ and $\mathsf g_*=\mathsf g(\lambda_*)$.

\item[\textup{(O1)}] We fix $d_0\in(0,\frac12)$.

\item[\textup{(O2)}] We fix $c_\Xi$ with $0<c_\Xi\le\frac12(\frac12-d_0)$.

\item[\textup{(O3)}] Let $\rho_0$ be the least integer $\ge0$ such that \[t_\rho:=3K_*^2\,\mathsf g(\lambda_*)\,\mathsf g(\lambda_*/2)\,e^{-\lambda_*\rho_0/2}\le\frac18(\frac12-d_0-c_\Xi);\] this is the display \eqref{eq:row-mask-radius} of \S\ref{Y-A-1}.

Item~E5 of \S\ref{Y-F-1}, the data \textup{(SD)} in item~1 of \S\ref{Y-G0-1}, and \S\ref{Y-A-1} denote by $\mathsf t_\rho$ a different object, a quantity of the rows, namely \[\|\ell'^{\rm c}\|_{\rm r}\sup_{P_0}\sum_{y':\,{\rm dist}(y',P_0)>\rho_0}|\ell'^{\rm c}_{y'}(P_0)|.\]

For the rows of \textup{(SD)} that quantity is at most $3K_*\mathsf g(\lambda_*)\cdot K_*\sum_{|u|_\infty>\rho_0}e^{-\lambda_*|u|_\infty}$, and this in turn is at most the number $t_\rho$ above. The comparison rests on two facts. By \Cref{Y-A-Th29}\,(a), $|\ell'_{y'}(P)|\le K_*e^{-\lambda_*{\rm dist}(y',P)}$ for every $\ell\in \mathcal R(b_*,K_*,\lambda_*,\frac19)$, the empty rows included, where $K=K_*$ is allowed because $L\ge L_A$; and the sum is bounded by the estimate of $\mathsf t_\rho$ in the proof of \Cref{Y-A-Th36}.

We put $\mathsf N_Z=\mathsf N_{\mathsf B}=1+3K_*^2\mathsf g_*^2$ and $\mathsf P_\rho=1+3(2+3(2\rho_0+1)^3)K_*^2\mathsf g_*^2$. The number $\mathsf P_\rho$ is fixed by the entries \textup{(O0)} to \textup{(O3)}, while the numbers $\mathsf l_L=1+\log L$, $\mathsf N_\sharp=c_1^*\,\mathsf l_L\,\mathsf N_Z\,\mathsf P_\rho$ and $\mathsf N_{\rm c}=17(\mathsf N_\sharp+1)$ are functions of $L$; since $L$ is chosen in \textup{(O4)} by inequalities in these functions, no circle arises.

\item[\textup{(O4)}] $L$ is odd and $L\ge L_*^{\rm su}$, where $L_*^{\rm su}$ is the least integer $\ge\max\{L_A,L_B\}$ such that \textup{(L1$^{\rm su}$)}, \textup{(L1$^\flat$)} at $\frac3{16}d_0$ and \textup{(L2$^{\rm su}$)} hold for all values of $L$ from $L_*^{\rm su}$ on (item~3, \textup{(O4)}, of \S\ref{Y-G0-1}). The condition \textup{(L2$^{\rm su}$)} is the inequality $L^{1/2}\ge1.2\cdot10^4\mathsf N_{\rm c}^2$. As $\mathsf N_{\rm c}=17(\mathsf N_\sharp+1)\ge17$, it gives $L\ge(1.2\cdot10^4\cdot17^2)^2>1.2\cdot10^{13}$. The number $c_L=(L^2-1)/(6L^2)$ depends on $L$, and $a\ge3/(c_LL)$; as $c_L<\frac16$, this yields $aL\ge18$.

\item[\textup{(O5)}] $\kappa\ge6$, and $\kappa^\sharp=10(\kappa+4)$.

\item[\textup{(O6)}] This entry consists of $\gamma=(73L^4)^{-1}$, of $b=16(1+3K_*^2\mathsf g_*^2)+a/L$, of the numbers $\theta_{\mathbb C}$ and $\vartheta_0^\flat$ of item~C.I~1.2 of \S\ref{Y-C-1}, taken at $K_*,\lambda_*$, of $\theta_D=\frac12\min\{\lambda_2^*,\theta_{\mathbb C},\theta_J\}$, of $C_D$, and of $c_R=64/\theta_D$, so that $c_R\theta_{\mathbb C}\ge128$. As $L\ge L_A$ yields $\lambda_*L\ge12$, the condition \textup{(R$^\flat$)} $c_R\mu\ge1.35$ of \Cref{Y-C-CIITh73}, \Cref{Y-C-CIITh74} and \Cref{Y-C-CIICor75} follows (item~C.II~5 of \S\ref{Y-C-1}, under ``Conditions on the letters''; item~3, \textup{(O6)}, of \S\ref{Y-G0-1}).

\item[\textup{(O7)}] $M=L^{m'}$ with $m'\ge1$, subject to \textup{(M2)} $M\ge1920(\kappa+4)c_R$, to \textup{(M4$^\flat$)} (item~C.II~5 of \S\ref{Y-C-1}, under ``Conditions on the letters'') and to \textup{(M5r)} $M\ge\rho_0$. The condition \textup{(M4$^\sharp$)} is a consequence of \textup{(M4$^\flat$)}. We put $v:=(LM)^3$.

\item[\textup{(O8)}] $\mathfrak p_0>2$; further $A_0>0$. We set $p_0(x)=A_0(\log x)^{\mathfrak p_0}$, $\varepsilon=p_0x^{-1/2}$ and $x'=x/L$; a prime on a letter means that it is taken at $x'$; and $\varepsilon'\le L^{1/2}\varepsilon$.

\item[\textup{(O9)}] $\mathsf G_0\le C_\varphi^\flat$ (item~C.I~1.2 of \S\ref{Y-C-1}). These two numbers are already determined at \textup{(O6)}, since they depend on $(L,a,K_*,\lambda_*)$ alone. Further, we set $\mathsf n=60$ and $c_\xi=(6.5\,C_\varphi^\flat)^{-1}$.

\item[\textup{(O10)}] \Cref{Y-C-Th68} and \Cref{Y-C-CIICor72} supply $\mathsf a_{\rm c}$ (item~C.II~4 of \S\ref{Y-C-1}, under ``Constants'') and $\mathsf K_{\rm c}:=\sqrt2\,\mathsf K_{\rm c}^0$. Here $\mathsf K_{\rm c}$ is the constant of \textup{(Cen)} in \Cref{Y-C-CIICor72}\,(a), where a single constant suffices for every $0<\hat\varepsilon\le\frac1{40}$. \Cref{Y-C-CIITh73}, \Cref{Y-C-CIITh74} and \Cref{Y-C-CIICor75} supply $\mathsf K_1^\flat$; the numbers $K_1^\sharp:=2\mathsf K_1^\flat$ and $K_2^\sharp:=3\mathsf K_1^\flat$ bound $\mathsf D_k$, and $\tilde{\mathsf D}=-\mathsf D_k/2$ obeys the halves of these bounds. \Cref{Y-D-C5} supplies $K_{{\rm q},k}$, \Cref{Y-E-C6} supplies $K_v^\flat$, and \Cref{Y-F-C7}, through item~F of \S\ref{Y-F-1}, supplies $u_b,u_t,u_1,\delta_3,\delta_3',e_1^\flat,{\rm Res}$. Formed in this section are $T:=L^{3/2}/(19\mathsf N_{\rm c})$, $a_\infty:=4\pi \mathsf N_{\mathsf B}(C_\varphi^\flat)^2$ and $n_a:=\frac32\,v\,a_\infty^2$.

\item[\textup{(O11)}] $0<c_h\le\delta_3'/2$. As $\delta_3'\le\delta_3$, it follows that $c_h\le\delta_3/2$.

\item[\textup{(O12)}] $E_a\ge\max\{32K_{{\rm q},k},\ 1500\,K_v^\flat L^{-1/2}\}$.

\item[\textup{(O13)}] $x_9:=1+\inf S_9$, where $S_9$ is the set of the $y\ge x_{\rm fl}$ (item~G of \S\ref{Y-F-1}) such that \textup{(Y.0)} to \textup{(Y.9)} of item~5 hold for all $x\ge y$. The set $S_9$ is an upper half line, and it is not empty by Step~0\,(i) of the proof of \Cref{Y-H-C9}. The floors \textup{(Y.0)} to \textup{(Y.9)} involve numbers rounded up to an integer, such as the radii $R$, $\varrho_L$, $r(x)$, $r'$ and $r_1$ of item~2, and $\lceil s(x)/L\rceil$ and $\lceil r(x)/L\rceil$. As these numbers jump as functions of $x$, $S_9$ need not contain its infimum. But $1+\inf S_9\in S_9$, and this number depends on what $S_9$ depends on. The floor of \Cref{Y-B-C1} and that of item~3, \textup{(O13)}, of \S\ref{Y-G0-1} have the same form. The table requires $L x_\flat\ge x_9$.
\end{enumerate}

The caps are $c_a(x)=E_ap_0^3x^{-1/2}$ and $c_h(x)=e^{-c_hp_0^2}$.

No constant depends on $A$, $u$, $z$ or $N$, nor on more of $\ell$ than the pure numbers $b_*$, $K_*$, $\lambda_*$ and its range.
\end{definition}

\begin{definition}[Item 2: the data \textup{(SD)} of a step, and the radii]
The data \textup{(SD)} of a step begin with the torus: $x\ge x_9$; $N$ is even, $M\mid N$, and \textup{(Tor)} $N\ge2M s(x)$, with $N$ the coarse side; the fine torus is $\mathbb Z^3/NL\mathbb Z^3$.

The lattices, the blocks, the base sites, the coarse distance, the cubes, the cells $\square$ and $\hat\square$ (the fine set under $\square$), and $\mathcal D$, $\mathcal P$, $|X|_c$ and $X^{+j}$ are those of item~1 of \S\ref{Y-G0-1} and of item~3.1 of \S\ref{Y-E-1}. The radii in force are those defined below.

The grids of cubes are fitted as in item~1 of \S\ref{Y-G0-1}, under ``Cubes'', so that the fine set $\hat\square$ under each cell $\square$ consists of exactly $L^3$ fine cubes. An object is attached to the cube, or the cell, of its base site. We say that $X$ \emph{holds} $p$ if $p$ belongs to a cube of $X$.

The bonds, paths and plaquettes \emph{of} a set $A$ of cubes, cells or sites are those all of whose sites lie in $A$ (item~1 of \S\ref{Y-G0-1}; item~3.3 of \S\ref{Y-E-1}). A plaquette $p$ is \emph{in} $A$ if its base site lies in $A$; this is the sense of ``in'' in $1_A$ and in $D(V)\cap X$, and also of ``based in'' and ``belongs to''. A field on the bonds of $A$ is \emph{tame there} if $\phi_p\le\varepsilon$ at the plaquettes of $A$.

For the rows, \textup{(SD)} requires that $\ell\in \mathcal R(b_*,K_*,\lambda_*,\frac19)$ have range $\le\varrho(x)$ and that $(\ell,\ell')$ satisfy the four conclusions of \Cref{Y-A-Th36}. These are listed in \textup{(SD)} in item~1 of \S\ref{Y-G0-1}. Here $\ell'$ is the family of rows that \Cref{Y-A-P5} forms from $\ell$, so that $\mathsf B'=1+\ell'^{\mathsf T}\ell'$, tensored with $1_{\mathfrak g}$; at $\ell=\emptyset$ its group \textup{(5)} is void. Further, $\ell'^{\rm c}$ is the cut of $\ell'$ at $\varrho(x')$.

Empty rows are admitted. The rows $\ell=\emptyset$, with $\mathsf H_y=\{0\}$ for all $y$, lie in $\mathcal R(1,0,\lambda_*,0)\subset \mathcal R(b_*,K_*,\lambda_*,\frac19)$ (by the display and the last sentence of \Cref{Y-A-SetsOfRows}) and have range $0$. For them $\mathrm{En}_\ell=\sum_pS_\varepsilon(\phi_p)$ and $\mathsf B^V=1$, while the numbers and the radii remain the ones evaluated at $(K_*,\lambda_*,\varrho(x))$. Neither a level nor a depth enters these data.

For $\ell=\emptyset$ the hypotheses of \Cref{Y-H-C9}, that is, the data \textup{(SD)} and the clauses \textup{($\Gamma$2)} to \textup{($\Gamma$7)}, are met as follows.
\begin{enumerate}
\item[\textup{(1)}] \Cref{Y-H-C9} is stated for every choice of the data \textup{(SD)}, hence for $(\emptyset,\emptyset')$, which satisfies the four conclusions that \textup{(SD)} requires. Of these, \textup{(1)}, \textup{(2)} and \textup{(3)} follow from \Cref{Y-A-Fr} at $k=0$ and $x_0=x$, namely from \textup{(Fr2)} for the sixteen kernels, \textup{(Fr3$'$)} and \textup{(Fr4)}. At these values, $x/L\ge2$ gives \textup{(X12)}. Conclusion \textup{(4)} consists of \textup{(L1$^{\rm su}$)} and \textup{(L2$^{\rm su}$)} and holds by \textup{(O4)}. Step~S6 of the proof of \Cref{Y-B-C1} treats this case.

\item[\textup{(2)}] Each of the clauses \textup{($\Gamma$2)} to \textup{($\Gamma$7)} admits $\ell=\emptyset$ within the data of the section to which it belongs: in items~C.I~1.1 and C.II~1.2 of \S\ref{Y-C-1}, with $K\ge0$; in \S\ref{Y-D-1}, item~B and item~I, Remark~(4); in \S\ref{Y-E-1}, item~2; and in \S\ref{Y-F-1}, item~C.

\item[\textup{(3)}] Besides the reason \textup{(2)}, a second argument runs as follows. By \Cref{Y-B-RZ} in Step~S6 of the proof of \Cref{Y-B-C1}, the objects that \textup{(SD)} and the classes attach to the rows are the same at the empty rows as at the zero rows whose $\mathsf H_y$ is the real line; and the zero rows fall under the hypothesis of those sections as a datum with rows.

\item[\textup{(4)}] \textup{($\Gamma$3)} is stated in the setting \textup{(Fr1)} to \textup{(Fr3)}. For $(\emptyset,\emptyset')$ these conditions follow from \textup{(SD)}, by the last sentence of \textup{(SD)} in item~1 of \S\ref{Y-G0-1}.
\end{enumerate}

At $\ell=\emptyset$, Steps~0 to 12 of the proof of \Cref{Y-H-C9} use $\ell$ only through \textup{(i)} to \textup{(iv)} below.
\begin{enumerate}
\item[\textup{(i)}] $\mathrm{En}_\ell\ge0$ and $\mathsf B^V$, in items~1, 3 and 4 of Step~1 of that proof. For $\ell=\emptyset$ we have $\mathrm{En}_\ell=\sum_pS_\varepsilon(\phi_p)$, $\mathsf B^V=1$ and $|\ell^VF^\flat|^2=0$, and the identity stays as it is.

\item[\textup{(ii)}] $\zeta_p$, $\ell_y$ and the mask $m_y$, in Step~2\,(iv) of that proof. For $\ell=\emptyset$ there is no such object, and $A_\square=e^{V_\square}$ holds factor by factor.

\item[\textup{(iii)}] The range $\le\varrho(x)$, in the dependence of $[A]^\ell_{\mathfrak Z}$ in Step~3\,(V) of that proof. For $\ell=\emptyset$ the range is $0\le\varrho(x)$, so $[A]^\ell_{\mathfrak Z}$ depends on no more than what Step~3\,(V) states.

\item[\textup{(iv)}] $\|\ell^{V*}\ell^V\|_\infty\le3K_*^2\mathsf g_*^2$, in Step~6, \textup{(t)}, of that proof. For $\ell=\emptyset$ this norm is $0$.
\end{enumerate}

The family $\ell'=\emptyset'$ is not empty. The steps need two things from it: the energy $\mathrm{En}'=\mathrm{En}_{\ell'^{\rm c}}$, through \textup{(5a)}, and the inequality $\mathsf t_\rho\le t_\rho$ for $(\emptyset')^{\rm c}$, in which $t_\rho$ is the number of \textup{(O3)} and which holds by \textup{(O3)}, \Cref{Y-A-Th29}\,(a) and the proof of \Cref{Y-A-Th36}.

For $\ell=\emptyset$ the letters are those of the table: $E_a\ge\max\{32K_{{\rm q},k},1500K_v^\flat L^{-1/2}\}$ of \textup{(O12)}, and not merely $E_a\ge E_a^0$ of \Cref{Y-B-C1}; $c_h\le\delta_3'/2$ of \textup{(O11)}; and $x\ge x_9$, not merely $x\ge L x_\flat^{(1)}$.

For $\ell=\emptyset$ there are two systems. Let $u_W$ be the member of Wilson's weight, so that $e^{-x\sum_p(1-\cos\phi_p)}=\rho[u_W;\emptyset]$ by \Cref{Y-B-LW}\,(a). The member $F_\emptyset(u_W)$ need not coincide with the system $u'$ of \Cref{Y-B-C1}, which \textup{(F1)} to \textup{(F4)} of \Cref{Y-G-ThA} construct; each of the two represents the image of Wilson's weight under the step (see item~8).

A pair $(V',e)$ consists of a coarse field $V'$ and of $e\in\mathfrak E=\prod_PS^2$, where $S^2$ is the unit sphere of $(\mathfrak g,\|\cdot\|)$, one copy for each coarse plaquette $P$. The measure $de$ is the product, taken over all $P$, of the rotation invariant probability measures on the spheres. It has three properties: $\int_{\mathfrak E}de=1$; the $e_P$ are independent under $de$; and for every $G$-valued $u$ the map $e\mapsto\operatorname{Ad}_ue$, given by $(\operatorname{Ad}_ue)_P=\operatorname{Ad}_{u(y_P)}e_P$ with $y_P$ the base site of $P$, leaves $de$ invariant, since every $\operatorname{Ad}_{u(y_P)}$ acts as a rotation on the three-dimensional Euclidean space $(\mathfrak g,\|\cdot\|)$. The mean over $e$ is the integral against $de$.

For a pair we put $X'_P=\mathsf L(\partial V'(P))+\pi e_P1\{\partial V'(P)=-1\}$; then $|X'_P|=\phi'_P\le\pi$.

The step acts on a bounded Borel function $\rho$ by \[\rho'(V')=\int\rho(U)\prod_ck_{ax}(V'(c),\mathsf M(U)(c))\,dU\] (item~C.I~4 of \S\ref{Y-C-1}, under ``The step, the chart''; item~3.4 of \S\ref{Y-E-1}), and on a complex $\rho$ by linearity.

The radii in coarse units are $R=\varrho(x)=\lceil c_R\log x\rceil$, $\varrho_L=\lceil R/L\rceil$, $r_w=\lfloor R/2\rfloor$, $r_G=r_w+\varrho_L+1$, $R_2=2R+2$, $R_3=2R+\varrho_L+5$, $R_\chi=R+\varrho_L+2$ and $R_v=2R+2\varrho_L+7$. Further, we put $\ell_w(x)=6M r(x)+12$, $D_w=\lceil\ell_w(x)/L\rceil+1$ and $\omega=M+\lceil M r(x)/L\rceil+D_w+R_3+2$.

The radii in units of cubes or of cells are $r(x)=\lceil20c_R\log x/M\rceil$ and $s(x)=8r(x)$. We write $X^{\approx}=X^{+r(x)}$ and $X^{\approx\approx}=X^{+s(x)}$, we put $r'=r(x')$, $s'=s(x')$ and $s_0=2r'$, and for sets of cells we write $Y^{\approx}=Y^{+r'}$. Finally $r_l=\lceil(\varrho_L+2)/M\rceil$, $r_1=\lceil R_2/M\rceil+3r_l+5$, $\varrho_*=2r_1+\lceil s(x)/L\rceil+2$, $k_t=\lceil(r_w+1)/M\rceil$ and $k_\chi=\lceil R_\chi/M\rceil$.
\end{definition}

\begin{definition}[Item 3: domains, members, caps, densities and distances]
For $X\in\mathcal D$, the domain $\mathcal U_X(x)$ (item~1 of \S\ref{Y-G0-1}, under ``Domains'') consists of the maps $\mathsf V$ that send the oriented bonds of $X^{\approx}$ into $\operatorname{SL}(2,\mathbb C)$, with $\mathsf V(\bar b)=\mathsf V(b)^{-1}$, such that the following hold.
\begin{enumerate}
\item[\textup{(u1)}] $|\mathsf V(b)|<1+\vartheta(x)$, where $\vartheta(x)=\vartheta_0^\flat(1-1/\log x)$.

\item[\textup{(u2)}] $|\mathsf V(\pi)|<\mathsf b(x)=2-1/\log x$ for every path $\pi$ in $X^{\approx}$ of at most $\ell_w(x)$ bonds.

\item[\textup{(u3)}] $|\partial\mathsf V(p)-1|<2\varepsilon$.
\end{enumerate}

For a set $Y$ of cells, the domain $\mathcal U_Y(x')$ is defined in the same way, with $x'$ in the place of $x$. The fields in $\mathcal U_Y(x')$ are extended off $Y^{\approx}$ by the value $1$, and $X'$ is then defined by the convention given below for $U^\tau(B;x')$.

For $\tau\ge1$ and a set $B$ of coarse sites, $U^\tau(B;x')$ (item~3.3 of \S\ref{Y-E-1}) is the set of the $\operatorname{SL}(2,\mathbb C)$ fields on the bonds of $B$ such that the following hold.
\begin{enumerate}
\item[\textup{(e1)}] $|V'(c)|<1+\vartheta(x')+\mathsf g_x$.

\item[\textup{(e2)}] $|V'(\pi)|<\mathsf b(x')+\mathsf g_x$ for every path $\pi$ in $B$ of at most $\ell_w(x')$ bonds.

\item[\textup{(e3)}] $|\partial V'(P)-1|<2\tau\varepsilon'$.
\end{enumerate}

Here $\mathsf g_x=\vartheta_0^\flat\log L/(4(\log x)^2)$. On such a field we put $X'_P:=-i\log\partial V'(P)$ for each $P$ with all four sites in $B$; the field is extended by $1$, and $X'_P:=0$ for every other $P$, which is the convention of item~3.7 of \S\ref{Y-E-1}.

The quantity $a_{0,b}$ depends on $X'$ within $R_3$ of $y(b)$ (\textup{($\Gamma$3)}, ``The dependence at every pair''). In what follows, $a_{0,b}$ enters only at those bonds $b$ for which the box of radius $R_3$ about $y(b)$ is contained in $B$ (Step~2\,(ii), Step~6\,(iii) and Step~8, \textup{(8.2)}, of the proof of \Cref{Y-H-C9}); hence every value of $a_{0,b}$ that enters is independent of the convention.

We put $1_Y(V)=1\{\phi_p\le\varepsilon\text{ for all }p\text{ in }Y\}$, and we say that $V$ is \emph{tame on} $Y$ if $1_Y(V)=1$. The energy is \[\mathrm{En}_\ell(U)=\sum_pS_\varepsilon(\phi_p)+\frac12\sum_ym_y(U)\,|\ell^U_y\zeta|^2,\] where $S_\varepsilon$, $\zeta$, the dressed rows $\ell^U_y$ and the mask $m_y$, of radii $\rho_0$ and $\varrho(x)+1$, are those of items~E1 and E2 of \S\ref{Y-F-1}.

For $A\ge1$, $M_A^{\rm su}(x)$ is the set of the triples $u=(E^a,E^h,w)$ with the three properties below.
\begin{enumerate}
\item[\textup{(ca$_A$)}] For $X\in\mathcal D$, $E^a(X,\cdot)$ is holomorphic on $\mathcal U_X(x)$, invariant under $G$-valued gauge transformations and $0$ at pure gauges, and $\sup_{\mathcal U_X(x)}|E^a(X,\cdot)|\le A\,c_a(x)\,e^{-\kappa(\#X-1)}$.

\item[\textup{(ch$_A$)}] $E^h(X,\cdot)$ is defined on those $G$-valued fields on the bonds of $X^{\approx}$ which are tame there; it is complex, Borel and gauge invariant, and $\sup|E^h(X,\cdot)|\le A\,c_h(x)\,e^{-\kappa(\#X-1)}$, with the supremum over these fields.

\item[\textup{(c$\Xi_A$)}] For $X\in\mathcal P$, $w(X,\cdot)$ is complex, Borel and gauge invariant, it is a function of $V$ on the bonds of $X^{\approx\approx}$, and $|w(X,V)|\le A\,e^{-\kappa(|X|_c-1)}$ at every $V$; this bound is the \emph{weak cap}.
\end{enumerate}

The \emph{strong cap} is the bound \[|w(X,V)|\le\exp\{-c_\Xi x\varepsilon\sum_{p\in D(V)\cap X}\phi_p-\kappa(|X|_c-1)\}.\] We call $u$ \emph{real} if $E^a$ is real at $G$-valued fields and $E^h$ and $w$ are real.

Consider a triple $u=(E^a,E^h,w)$ with those clauses of \textup{(ca)}, \textup{(ch)} and \textup{(c$\Xi$)} which carry no size. For it, $\|u\|_x$ is the smallest $A\ge0$ with which all three size bounds are true, that is, the largest of the three weighted suprema; it is a norm (item~3 of \S\ref{Y-I-1}), and $M_A^{\rm su}(x)=\{\|u\|_x\le A\}$. The pure gauges are the fields $1^u$ with $G$-valued $u$.

We put $1E(X)(V):=1_{X^{\approx}}(V)\,[E^a(X,V)+E^h(X,V)]$ and \[\rho[u;\ell](V):=\exp\{-x\mathrm{En}_\ell(V)+\sum_{X\in\mathcal D}1E(X)(V)\}\cdot\Xi_w(V),\] where $\Xi_w(V)=\sum_{\{X_i\}}\prod_iw(X_i,V)$, the sum running over the finite families of sets $X_i\in\mathcal P$ of pairwise distance $>s(x)$ such that every $X_i$ holds a plaquette of $D(V)$ and $D(V)$ lies inside their union. The empty family, whose product is $1$, is admitted exactly when $D(V)$ is empty. These families are called the \emph{families of} \textup{(c$\Xi$)}.

We put \[C^{\rm su}_\flat(x;\ell)=\{e^{\mathsf V}\rho[u;\ell]>0:\ \mathsf V\text{ real},\ u\in M_1^{\rm su}(x)\text{ real with the strong cap}\}.\]

For two members $u_1$ and $u_2$ we put $\mathsf D^{\rm in}_x(u_1,u_2):=\|u_1-u_2\|_x$, and we let $\mathsf D_x(u_1,u_2)$ be the largest of the three numbers \[\sup_X\sup_{\mathcal U_X(x)}|E_1^a-E_2^a|\,e^{\kappa(\#X-1)}/c_a(x),\] $\sup_X\sup|E_1^h-E_2^h|\,e^{\kappa(\#X-1)}/c_h(x)$, with the inner supremum over the tame $G$-valued fields, and \[\sup_{X,V}|w_1-w_2|(X,V)\exp\{c_\Xi x\varepsilon\sum_{D(V)\cap X}\phi_p+\kappa(|X|_c-1)\}.\]

Hence $\mathsf D^{\rm in}_x\le\mathsf D_x$: the first two of these numbers are the corresponding weighted suprema of $u_1-u_2$, and in the third the weight of $|w_1-w_2|$ is at least the weight $e^{\kappa(|X|_c-1)}$ of the weak cap.
\end{definition}

\emph{Item 4: what is used of the earlier theorems.} The inputs are the following clauses of earlier theorems, each stated in full.

\begin{enumerate}
\item[\textup{($\Gamma$2)}] \Cref{Y-C-CITh66}\,(a), (b), (c), (d), (f), (i) and (k), together with \Cref{Y-C-CIL64}\,(e). On the set $\mathfrak V=\{\operatorname{SL}(2,\mathbb C)\text{ fields with }|V'(c)|\le1+\vartheta_0^\flat\}$ there is a matrix $G=G^{V'}$ on $I^\circ$, the set of fine bonds off the combs, such that $G$ and the matrices $C$ and $S$ obtained from it have the properties \textup{(2a)} to \textup{(2f)}, while \textup{(2g)} is a property of $\mathsf g^V$.
\begin{enumerate}
\item[\textup{(2a)}] The entry $G(b,b')$ vanishes as soon as $b'$ lies farther than $r_w$ from $b$, and $(G\xi)_b$ depends on $\xi$ only through its values on a window $Y_b$ of bonds within the distance $r_w$ of $b$.

\item[\textup{(2b)}] The row $G(b,\cdot)$ depends on $V'$ only through the coarse bonds within the distance $r_G$ of $y(b)$. The matrix $C:=G G^{\mathsf T}$ is symmetric; its entry $C(b,b')$ vanishes as soon as $b'$ lies farther than $2r_w$ from $b$, where $2r_w\le R$; and its row $C(b,\cdot)$ depends on $V'$ only within the distance $3r_w+\varrho_L+1$ of $y(b)$, where $3r_w+\varrho_L+1\le R_2-1$.

\item[\textup{(2c)}] For every $G$-valued $u$, the relation $G[V'^u]=\operatorname{Ad}_u\,G[V']\,\operatorname{Ad}_u^{-1}$ holds on all of $\mathfrak V$.

\item[\textup{(2d)}] The entries of $G$ are holomorphic in the interior of $\mathfrak V$, and they are real at $G$-valued fields.

\item[\textup{(2e)}] For every $\xi$ and every bond $b$ we have $\|(G\xi)_b\|\le C_\varphi^\flat\sup_{Y_b}\|\xi\|$, and the norms obey $\|G\|_\infty,\|G^{\mathsf T}\|_\infty\le \mathsf G_0\le C_\varphi^\flat$.

\item[\textup{(2f)}] The matrix $G$ is invertible, and $S:=C^{-1}-q^V_\ell$ is symmetric.

\item[\textup{(2g)}] We recall \Cref{Y-C-CIL64}\,(e). It expresses $(\mathsf g^VX')_{b'}=\sum_P\sum_{\mathsf t}g_{\mathsf t}\,\operatorname{Ad}_{V'(\pi_{\mathsf t})}X'_P$ as a finite sum in which the numbers $g_{\mathsf t}$ are independent of every field and each $\pi_{\mathsf t}$ is a coarse path that runs from $y(b')$ to the base site of $P$ and keeps all its sites within the distance $2\varrho_L+1$ of $y(b')$. Consequently, at every field of $\mathfrak V$ and for every $\mathfrak g_{\mathbb C}$-valued $X'$, the value $(\mathsf g^VX')_{b'}$ depends on $V'$ only through its values on bonds, and on $X'$ only through its values on plaquettes, all of whose sites lie within the distance $2\varrho_L+2$ of $y(b')$.
\end{enumerate}

\item[\textup{($\Gamma$3)}] \Cref{Y-C-Th68}\,(a) and (c), together with \Cref{Y-C-CIICor72}\,(a), which is \textup{(Cen)}. We put $a_0:=C \mathsf g^VX'$.

\emph{The two bounds.} Let $p$ be a fine plaquette, let $b$ be a bond of $p$, let $0<\hat\varepsilon\le\frac1{40}$, and take an $\operatorname{SL}(2,\mathbb C)$ field on the box of radius $R_3$ about $y(p)$ that satisfies \textup{(e1)}, \textup{(e2)} and $|\partial V'(P)-1|\le2\hat\varepsilon$. Then \[ |a_{0,b}|\le \mathsf a_{\rm c}\hat\varepsilon\qquad\text{and}\qquad|(EX'-\mathrm{d}^Va_0)_p|\le \mathsf N_{\rm c}L^{-2}\hat\varepsilon+\mathsf K_{\rm c}\hat\varepsilon^2. \]

Both quantities depend on $(V',e)$ only through its values on that box. \Cref{Y-C-Th68}\,(a) (item~C.II~4 of \S\ref{Y-C-1}) is a clause on the dependence only and contains no bound, whereas \Cref{Y-C-CIICor72}\,(a) states the two bounds alone; for $a_{0,b}$, where ${\rm dist}(b,y(p))\le1$, the dependence just stated is what that clause asserts, and for $(EX'-\mathrm{d}^Va_0)_p$ it is derived below from the values $a_{0,b}$ at the four bonds of $p$. The bound $|a_{0,b}|\le\mathsf a_{\rm c}\hat\varepsilon$ is \Cref{Y-C-Th68}\,(c) as well as part of \Cref{Y-C-CIICor72}\,(a); the bound on $|(EX'-\mathrm{d}^Va_0)_p|$ comes from that corollary alone.

Only the clause on the dependence, \Cref{Y-C-Th68}\,(a), is free of the conditions \textup{(h2)}, \textup{(h3)} and \textup{(Fr2)}; the two bounds are not. Suppose that the field is given on a set containing the box and that $X'$ is defined on the plaquettes of that set. Then at each bond $b$ with ${\rm dist}(b,y(p))\le1$, in particular at the four bonds of $p$, the value $a_{0,b}$ coincides with the value computed on the box. The same is true of $(EX'-\mathrm{d}^Va_0)_p$: either $(EX')_p=X'_P$ for a $P$ based at $y(p)$ or $(EX')_p=0$, and $(\mathrm{d}^Va_0)_p$ is built from these four values $a_{0,b}$ by transports along no more than one coarse bond from $y(p)$ (Step~(0) of the proof of \Cref{Y-C-Th68}). For this reason the two bounds can be applied where the field is given on a larger set.

\emph{The dependence at every pair.} At every field of $\mathfrak V$ and for every $\mathfrak g_{\mathbb C}$-valued $X'$, hence at every pair $(V',e)$, rough plaquettes included, the value $a_{0,b}$ depends on $V'$ only through its values on bonds, and on $X'$ only through its values on plaquettes, all of whose sites lie within the distance $R_3$ of $y(b)$, and it depends on nothing else.
\begin{proof}[Proof of the dependence at every pair]
We have $a_{0,b}=\sum_{b'}C(b,b')(\mathsf g^VX')_{b'}$. By \textup{(2b)}, $C(b,b')=0$ as soon as $b'$ lies farther than $2r_w$ from $b$, and $2r_w\le R$; moreover the row $C(b,\cdot)$ depends on $V'$ only within the distance $R_2-1$ of $y(b)$. By \textup{(2g)}, each value $(\mathsf g^VX')_{b'}$ depends on $V'$ and on $X'$ only within the distance $2\varrho_L+2$ of $y(b')$. Combining the two facts, all sites that enter lie within the distance $\max\{R_2-1,\ R+2\varrho_L+2\}\le R_3$ of $y(b)$, where the last inequality holds because $\varrho_L\le R+3$. If, in addition, ${\rm dist}(b,o)\le1$ for a coarse site $o$, as is the case for a bond $b$ of $p$ and $o=y(p)$, then all these sites also lie within the distance $\max\{R_2,\ R+2\varrho_L+3\}\le R_3$ of $o$, the last inequality now holding because $\varrho_L\le R+2$.
\end{proof}

\item[\textup{($\Gamma$4)}] \Cref{Y-C-CIITh73}\,(b), \Cref{Y-C-CIITh74}\,(a), (b) and (c), and \Cref{Y-C-CIICor75}, all taken at $\mathsf n=60$. Summing over the connected sets $Y$ of cells, we have $\log\det C^1-\log\det C^{V'}=\sum_Y\mathsf D_k(Y,V')$ at every $V'$, and $S=\sum_YS_k(Y;V')$ on $\mathfrak V$.
\begin{enumerate}
\item[\textup{(4a)}] The function $\mathsf D_k(Y,\cdot)$ is real and gauge invariant, it vanishes at pure gauges, and $|\mathsf D_k|\le K_1^\sharp e^{-\kappa^\sharp(\#Y-1)}$ at every $V'$; it has a holomorphic continuation to $\mathcal U_Y(x')$, on which $|\mathsf D_k|\le K_2^\sharp\varepsilon'e^{-\kappa^\sharp(\#Y-1)}$.

By definition (\Cref{Y-C-CIITheRest}, the display before \Cref{Y-C-CIITh74}), $\mathsf D_k(Y,V')=-\mathsf D^q(Y,V')-[\mathsf D^u(Y,V')-\mathsf D^u(Y,1)]$; the continuation is given by the same formulas, evaluated at the element of $\mathcal U_Y(x')$ extended by $1$ (proof of \Cref{Y-C-CIICor75}).

The function $\mathsf D^q(Y,\cdot)$ is holomorphic on the open set \[\FitDisplay{\dimexpr\linewidth-1pt\relax}{\{|V'(c)|<1+\vartheta_0^\flat\text{ on the bonds of }Y^{+1}\}}\] (\Cref{Y-C-CIITh73}\,(b)), and $\mathsf D^u(Y,\cdot)$ is holomorphic in the interior of $\mathfrak V$ (\Cref{Y-C-CIITh74}\,(b)); each of these two open sets contains all $G$-valued fields. It follows that $\mathsf D_k(Y,\cdot)$ is continuous on the $G$-valued fields. On $\mathcal U_Y(x')$ the function is invariant under $G$-valued gauge transformations, by \Cref{Y-C-CIITh73}\,(b) and \Cref{Y-C-CIITh74}\,(b); the argument is given in Step~8 of the proof of \Cref{Y-H-C9}, paragraph ``The clauses without a size''.

\item[\textup{(4b)}] We have $S_k(Y)(b,b')=0$ unless both $b$ and $b'$ lie in $\hat Y$. The matrix $S_k(Y)$ is symmetric, covariant as in \textup{(2c)}, holomorphic and real at $G$-valued fields, and $\|S_k(Y;V')\|\le x^{-57}e^{-\kappa^\sharp\#Y}$ at every field of $\mathfrak V$.

\item[\textup{(4c)}] Both $\mathsf D_k(Y,\cdot)$ and $S_k(Y;\cdot)$ depend on $V'$ only through its values on the bonds of $Y^{+\lceil R_2/M\rceil}$. In \S\ref{Y-C-1} the number $\lceil R_2/M\rceil$ is written $r_u$.
\end{enumerate}

We put $\tilde{\mathsf D}:=-\mathsf D_k/2$ and $e_Y:=x^{-40}e^{-\kappa^\sharp\#Y}$.

\item[\textup{($\Gamma$5)}] \Cref{Y-D-C5}\,(b), (d) and (e): below, \textup{(5a)} is its clause (e), \textup{(5b)} its clause (b) and \textup{(5c)} its clause (d).
\begin{enumerate}
\item[\textup{(5a)}] At every pair $(V',e)$, \[ x\,\mathsf m_C-\frac x2\langle a_0,Sa_0\rangle=x'\mathrm{En}'(V')+\sum_jw_j+\sum_\square1_{\square^{\approx}}(V')\,P(\{\square\},V'). \] Here $\mathsf m_C=\frac12\langle EX',\mathsf B^VEX'\rangle-\frac12\langle \mathsf g^VX',C \mathsf g^VX'\rangle$; the term $\mathrm{En}'=\mathrm{En}_{\ell'^{\rm c}}$ is taken at $x'$; and $w_j=\sum_{\square\in \mathrm{mk}_j'^{+r'}}T_\square$, where $T_\square$ is that of item~E7 of \S\ref{Y-F-1} and the classes are those of item~E9 of \S\ref{Y-F-1}.

\item[\textup{(5b)}] The function $T_\square$ is real, jointly Borel and invariant under $(V'^u,\operatorname{Ad}_ue)$, and it depends on $(V',e)$ only within the distance $M r'-1$ of $\square$.

\item[\textup{(5c)}] The function $P(\{\square\},\cdot)$ is holomorphic on $\mathcal U_{\{\square\}}(x')$, invariant, real at $G$-valued fields and zero at pure gauges, and its supremum does not exceed $K_{{\rm q},k}\,p_0'^2\varepsilon'$. For $\#Y\ge2$ we put $P(Y,\cdot):=0$.
\end{enumerate}

\item[\textup{($\Gamma$6)}] \Cref{Y-E-C6}\,(a), (b), (c), (d), (e.iii), (e.iv) and (e.v). Here $V_\square(\xi,V')$ is the vertex term defined by the display of item~3.8 of \S\ref{Y-E-1}, with the law of \Cref{Y-C-CITh66}, and $\chi_\square(\xi)=1$ if and only if $|\xi_b|\le c_\xi\varepsilon$ at all $b\in I^\circ$ within the distance $R_\chi$ of $\hat\square$.

Let $V'\in\mathcal U_{\{\square\}}(x')$, and let $\xi$ be real with $\chi_\square=1$. Then $V_\square$ is holomorphic in $V'$; it depends on $\xi$ only within the distance $r_w+\varrho_L+1$ of $\hat\square$ and on $V'$ only within the distance $R_v$ of $\square$; it is independent of every member and of $e$; it is invariant under $(\operatorname{Ad}_u\xi,V'^u)$; it satisfies $|V_\square|\le K_v^\flat p_0^3x^{-1/2}$; and it is real at $G$-valued $V'$.
\begin{enumerate}
\item[\textup{(6e)}] \Cref{Y-E-C6}\,(e.iii), (e.iv) and (e.v) state the following for the $A_\square$ of item~6 of this section (item~E11 of \S\ref{Y-F-1}; the same product appears in item~3.9 of \S\ref{Y-E-1}). At every real $\xi$ and every pair $(V',e)$, tame or rough, whether the tests are passed or failed, $A_\square$ depends on $\xi$ only through its values on the bonds of $I^\circ$ within the distance $r_w+\varrho_L+2$ of $\square$, and on $(V',e)$ only through its values on bonds and plaquettes within the distance $R_v$ of $\square$; moreover $A_\square(\operatorname{Ad}_u\xi,V'^u,\operatorname{Ad}_ue)=A_\square(\xi,V',e)$ for $G$-valued $u$, and $A_\square$ is Borel in $(\xi,V',e)$.
\end{enumerate}

\item[\textup{($\Gamma$7)}] \Cref{Y-F-C7}. From \S\ref{Y-F-1} we take item~E (the objects), item~F (the numbers) and item~G (the floor $x_{\rm fl}$), and the clauses \textup{(7a)}, \textup{(7d)} and \textup{(7e)}.
\begin{enumerate}
\item[\textup{(7a)}] Of clause (a) of \Cref{Y-F-C7} we use one assertion alone: $\hat\lambda_{(\mathsf J,t)}$ and $\hat\lambda_{(\mathsf J,t)}\cdot1_{\rm ind}$ depend on $\xi$ only through its values on the bonds of $I^\circ$ under $\mathfrak b^{+2r_1}$. The set of classes of a cell datum, written $J$ in Sections~\ref{Y-G0-1}, \ref{Y-F-1} and~\ref{Y-G-1}, is written $\mathsf J$ in this section, where $J$ is the core of a letter.

\item[\textup{(7d)}] We state clause \textup{(d$^\circ$)} of \Cref{Y-F-C7}. Its data consist of cell data $(\mathsf J_i,t_i)$ that are admitted and pairwise compatible; of a set $J_b$ of finitely many cells (in \Cref{Y-F-C7} this set is written $J_0$), disjoint from every $\mathfrak b_i^{+2r_1}$; of a finite set $\mathsf Y$ of cells, with $|\varphi_{\mathsf Y}|^2$ defined as in item~E1 of \S\ref{Y-F-1}, that is, by summing $|(G\xi)_b|^2$ over all bonds $b$ of $I^\circ$ in $\hat{\mathsf Y}$; of a number $\varsigma$ with $0\le\varsigma\le\frac\gamma2x^{-2}$; and of an arbitrary pair $(V',e)$. For such data, \[\FitDisplay{\dimexpr\linewidth-1pt\relax}{ E_\xi\bigl[\prod_i\hat\lambda_{(\mathsf J_i,t_i)}\cdot\prod_{J_b}b_\square\cdot e^{\frac x2\varsigma|\varphi_{\mathsf Y}|^2}\bigr]\le\exp\bigl\{4.8\,v\,\frac2\gamma\,\varsigma\,(2r_1+1)^3\#\mathsf Y+x^{-1/8}\bigl(\sum_i\#\mathfrak b_i^{+2r_1}+\#J_b\bigr)\bigr\}\cdot u_b^{\#J_b}\cdot\prod_ie^\circ_{(\mathsf J_i,t_i)}. }\]

The sizes in this bound are $e^\circ_{(\mathsf J,t)}={\rm Res}_{(\mathsf J,t)}\cdot u_t^{\#t}\cdot u_1^{1[t\text{ non-empty}]}$ and ${\rm Res}=\exp\{-(\frac12-d_0-t_\rho-3e_1^\flat)\,x'\varepsilon'\sum_{\mathsf J}\sum_{D'_j}\phi'_P\}$, defined with the number $t_\rho$ of \textup{(O3)}. The quantity ${\rm Res}$ of \Cref{Y-F-C7}, defined with the $\mathsf t_\rho$ of the rows, is at most this one. Hence \Cref{Y-F-C7}\,(d$^\circ$) implies the bound above in the form given here, and each of $r_\Xi$, $\upsilon_f$, $\upsilon_*$ and $x_9$ depends on the rows only through their range and the numbers $b_*$, $K_*$, $\lambda_*$.

\item[\textup{(7e)}] By clause (e) of \Cref{Y-F-C7}, $u_1=e^{-\delta_3p_0^2}$.
\end{enumerate}
\end{enumerate}

\emph{Tools.} The proof of \Cref{Y-H-C9} uses the following tools: the formula of the chart, \Cref{Y-C-L51}, in item~C.I~4 of \S\ref{Y-C-1}; \Cref{Y-H-L51} and \Cref{Y-H-Th53}, in Step~2 of the proof of \Cref{Y-H-C9}; \Cref{Y-H-Th57}, in Step~8, \textup{(8.2)}, of the proof of \Cref{Y-H-C9}; \Cref{Y-G-L18}, together with \Cref{Y-G-An0}, and \Cref{Y-G-L19}, in Step~8 of the proof of \Cref{Y-G-ThA}; and \Cref{Y-H-L52}, in Step~10\,(9) of the proof of \Cref{Y-H-C9}.

From item~1 of \S\ref{Y-G0-1} the proof takes the set $Y_0(X)$ and the notion ``small'', the grids, and two bounds: a given cube lies in at most $(26e)^{n-1}$ members of $\mathcal D$ with $n$ cubes, and in at most $(52e)^n$ members of $\mathcal P$ with $|X|_c=n$.

\begin{definition}[Item 5: the numbers defined here and the floor \textup{(Y*)}]
We put $t=x^{-1/4}$, $\varepsilon_2=x^{-1/8}$, $\varepsilon''=x^{-1/4}$, $\varsigma''=1.01\,x^{-20}$, $d_\alpha=(\kappa+1)(\#J_\alpha+s')$ and $r_T:=\sqrt2\,C_\varphi^\flat c_\xi\varepsilon+17T \mathsf a_{\rm c}\varepsilon'$. Further, $r_\Xi=\frac12-d_0-t_\rho-3e_1^\flat-c_\Xi$, \[ \upsilon_f=\exp\{-r_\Xi p_0'^2+1.01\,K_1^\sharp(2r'+1)^3+4+(\kappa+4)(1+2s')\}, \] and $\upsilon_*=(\max\{u_t^{1/2},\upsilon_f^{1/4}\})^{1/L^3}$.

The floor \textup{(Y*)} consists of the following conditions \textup{(Y.0)} to \textup{(Y.9)} on $x$.
\begin{enumerate}
\item[\textup{(Y.0)}] This condition is made up of \textup{(0a)} to \textup{(0e)}.
\begin{enumerate}
\item[\textup{(0a)}] It contains the floors \textup{(S1)} and \textup{(S2)} of item~C.II~5 of \S\ref{Y-C-1}, paragraph ``Floors'', with $r_u=\lceil R_2/M\rceil$ as there, and the floors of \Cref{Y-D-C5} and of \Cref{Y-E-C6}, as item~G of \S\ref{Y-D-1} and item~6 of \S\ref{Y-E-1} list them.

The only constraint that Sections~\ref{Y-D-1} and~\ref{Y-E-1} put on their letter $x_\flat$ is that the floors of the section hold for all $x\ge L x_\flat$, and \S\ref{Y-D-1} asks in addition that $x_\flat\ge e^2$. We therefore apply \Cref{Y-D-C5} and \Cref{Y-E-C6} with $x_9/L$ in the role of their $x_\flat$; here $x_9/L\ge x_{\rm fl}/L\ge L\ge e^2$, because $x_{\rm fl}\ge L^2$ (item~G of \S\ref{Y-F-1}).

The condition $x\ge x_{\rm fl}$ comes from \textup{(O13)} of item~1 of this section: there $S_9$ is the set of the $y\ge x_{\rm fl}$, with $x_{\rm fl}$ as in item~G of \S\ref{Y-F-1}, for which \textup{(Y.0)} to \textup{(Y.9)} are satisfied at every $x\ge y$, and $x_9=1+\inf S_9\in S_9$. So, for $x\ge x_9$, the condition $x\ge x_{\rm fl}$ holds in addition to \textup{(Y.0)} to \textup{(Y.9)}. The condition $x\ge x_{\rm fl}$ entails \textup{(G1$^\flat$)}, which gives $R\ge6$. The same condition also entails \textup{(F$_{\rm law}$)}, the clauses $x\ge L^2$ and $R\ge6$ of \textup{(P0)}, the conditions \textup{(P1)} and \textup{(P3)} of item~C.II~4 of \S\ref{Y-C-1}, the bound $\varepsilon'\le\frac1{40}$, and \textup{(W)}.

The third clause of \textup{(P0)}, as item~C.II~4 of \S\ref{Y-C-1}, paragraph ``Floors'', defines it, is the clause \textup{(Tor)}. This clause depends on $N$; it is a datum of \textup{(SD)} and not a floor. In the other sections whose floors enter the definition of $S_9$, by way of \textup{(0a)} or of $x_{\rm fl}$, the clause \textup{(Tor)} is likewise among the data and not among the floors: in Step~14 of the proof of \Cref{Y-E-C6}, in item~G of \S\ref{Y-F-1} and in item~G of \S\ref{Y-D-1}. In the last of these, \textup{(P0)} consists of the two clauses $x\ge L^2$ and $R\ge6$. The condition \textup{(W0)} of \Cref{Y-E-C6} contains both of them and no other clause of \textup{(P0)}. Consequently, none of the conditions in the definition of $S_9$ depends on $N$.

\item[\textup{(0b)}] $17T\varepsilon'\le\frac1{40}$, $\mathsf K_{\rm c}(17T\varepsilon')^2\le0.12\,\varepsilon$, $(p_0'/p_0)^4\ge\frac12$ and $\mathsf a_{\rm c}\varepsilon'\le1$.

\item[\textup{(0c)}] $3D_w\le\ell_w(x')$, together with the two conditions below.
\begin{enumerate}
\item[\textup{(D2)}] $\ell_w(x)\,(7.1\,r_T+18D_w\cdot17T\varepsilon')\le \mathsf g_x/2$.

\item[\textup{(D3)}] $(4.1\,r_T+2.1\cdot17T\varepsilon')^2\le0.1\,\varepsilon$.
\end{enumerate}

The conditions \textup{(D2)} and \textup{(D3)} are the conditions \textup{(D2)}$(\tau,r_T)$ and \textup{(D3)}$(\tau,r_T)$ of \Cref{Y-H-Th53}, at $\tau=17T$. Since both left sides increase with $\tau$, the two conditions hold at every $1\le\tau\le17T$. The condition \textup{(D0)}$(\tau)$ of the same theorem consists of $2D_w+1<N$, which is part of Step~0\,(iv) of the proof of \Cref{Y-H-C9}, of $3D_w\le\ell_w(x')$, and of $\tau\varepsilon'\le\frac1{40}$, which comes from \textup{(0b)}.

\item[\textup{(0d)}] The following three conditions.
\begin{enumerate}
\item[\textup{(C0)}] $18\omega\varepsilon'\le\frac1{10}$ and $160T\omega\varepsilon'\le1$.

\item[\textup{(C1)}] $2640\,\ell_w(x')(T+1)\omega\varepsilon'\le \mathsf g_x$.

\item[\textup{(C2)}] $102400\,T\omega^2\varepsilon'\le1$.
\end{enumerate}

\item[\textup{(0e)}] $3e_1^\flat\le\frac18(\frac12-d_0-c_\Xi)$.
\end{enumerate}

\item[\textup{(Y.1)}] $16(c_a+c_h)(x)\le\frac12$ and $K_v^\flat p_0^3x^{-1/2}\le\frac12$.

\item[\textup{(Y.2)}] $50\,x^{-1/8}(2(s_0+\varrho_*)+1)^3\le1$.

\item[\textup{(Y.3)}] The inequality \[ x^{1/4}(2s'+1)^3\cdot50\,e^{(\kappa+1)s'}\cdot[\,5e^{\kappa+3}L^3\cdot16(c_a+c_h)+2e^{\kappa+3}K_v^\flat p_0^3x^{-1/2}+x^{-20}+e^{\kappa+3}u_b+u_1\,]\le x^{-1/8}. \]

\item[\textup{(Y.4)}] The inequality \[ [\,u_t^{1/2}e^{\kappa+5}+220((2r_1+1)L)^3e^{\kappa+5}\cdot16\upsilon_*+6((2r_1+1)L+2r(x))^3e^{\kappa+5}\cdot16(c_a+c_h)\,]\cdot e^{(\kappa+4)s'}\le\frac12x^{-1/4}, \] together with $e^{(\kappa+1)s'}\upsilon_f^{1/2}\le\frac14x^{-1/4}$.

\item[\textup{(Y.5)}] $16\cdot10^3L^3e^{-c_h(p_0^2-p_0'^2)}\le\frac12$ and $50(2u_b+u_1)\le\frac12e^{-\delta_3p_0^2/2}$.

\item[\textup{(Y.6)}] $10^3e^{\kappa+3}L^{5/2}x^{-1/4}\le10^{-3}$, $8K_2^\sharp\le E_ap_0'^2$ and $x^{-10}\le c_a(x')/32$.

\item[\textup{(Y.7)}] $5r_1+2\lceil s(x)/L\rceil+6\le r'$ and $R_3+D_w+M(\lceil r(x)/L\rceil+2)\le M r'$.

\item[\textup{(Y.8)}] $4.04\,x^{-20}\le\frac\gamma2x^{-2}$, \[ 2x\varsigma''n_a+4.8\cdot4.04\,v\,\frac2\gamma\,x^{-20}(2r_1+1)^3\le x^{-1/8}, \] \[ 2x\varsigma''n_a+36\varsigma''(C_\varphi^\flat)^2\cdot3L^3(R+1)^3\cdot v(2r_1+1)^3\le x^{-1/8}, \] and $4\varsigma''(C_\varphi^\flat)^2\cdot3L^3(R+1)^3\le\frac12$.

\item[\textup{(Y.9)}] $2.001\,K_1^\sharp e^{\kappa+1}(2r'+1)^3e^{-(\kappa+26.7)r'}\le1$ and $K_1^\sharp e^{-\kappa^\sharp r'}\le1$.
\end{enumerate}
\end{definition}

\begin{definition}[Item 6: the letters of the alphabet, the clusters and the formulas for the output]
All objects of this item are defined at a pair $(V',e)$. We let $\xi\sim \mathcal N(0,x^{-1}1)$ on $I^\circ$, and we put $\varphi=G\xi$, $\eta=\varphi-a_0$ and $U=\mathsf U(\eta;V')$, with the chart of item~3.4 of \S\ref{Y-E-1} and of item~C.I~4 of \S\ref{Y-C-1}.

From items~E3 and E9 to E12 of \S\ref{Y-F-1} we take $F^\flat$ and $n^\flat$; the marks $\mathrm{mk}'$ and the set $f=\mathrm{mk}'^{+s_0}$; the classes $\mathrm{mk}'_j$, $f_j$ and $D'_j$; the tests $b_\square$; the factor $[A]^\ell_{\mathfrak Z}$; the cell data $(\mathsf J,t)$, with $\mathfrak b=f_{\mathsf J}\cup t$ and $\mathfrak Z=\mathfrak b^{+r_1}$, and the functions $\hat\lambda_{(\mathsf J,t)}\ge0$; the notion ``admitted'', which is \textup{(H17)}; and the notion of ``compatible'' cell data.

We put $A_\square:=[A]^\ell_{\{\square\}}$ and $T_Y:=\frac x2\langle\varphi,S_k(Y)\varphi\rangle-x\langle\varphi,S_k(Y)a_0\rangle$, and we let $\chi_X(\xi)=1$ if and only if $|\xi_b|\le c_\xi\varepsilon$ at all $b\in I^\circ$ within the distance $r_w+1$ of a bond of $X^{\approx}$.

The set $Y_0(X)$ and the notion ``small'', which means $n\le L/2-2$, with $n=\#X$ for $X\in\mathcal D$ and $n=|X|_c$ for $X\in\mathcal P$, are those of item~1 of \S\ref{Y-G0-1}. The cubes and the cells are ordered once and for all, and the words ``first'' and ``root'' refer to these orders.

We say that a set $Y$ is \emph{on} at $V'$ if $Y^{+r'}$ misses $\mathrm{mk}'$, and \emph{off} otherwise. A set $Y$ that is off is \emph{short} if it lies inside $f$, and \emph{long} otherwise. We put $n_j(V'):=\sum_{Y\subset f_j,\ {\rm off}}\tilde{\mathsf D}(Y,V')$.

\emph{The letters.} Each letter $\alpha$ carries a function, a core $J_\alpha$ and a set $I_\alpha$; there are the following letters.
\begin{enumerate}
\item[\textup{(a)}] For each $X\in\mathcal D$, the function is $\lambda^a_X=\chi_X\cdot1_{X^{\approx}}(U)\,(e^{E^a(X,U)}-1)$, and at complex $V'$ it is $\chi_X(e^{E^a(X,U)}-1)$; the core is $J=Y_0(X)$, and $I$ consists of $J$ and of the cells that meet $X^{\approx}$. Here and in \textup{(h)} and \textup{(v)} we agree that a product containing the factor $\chi_X$, $\chi_\square$ or $1_{X^{\approx}}(U)$ is $0$ wherever that factor is $0$, even where $E^a(X,U)$, $E^h(X,U)$ or $V_\square$ has no value.

\item[\textup{(h)}] For each $X\in\mathcal D$, the function is $\lambda^h_X=\chi_X\cdot1_{X^{\approx}}(U)\,e^{E^a(X,U)}(e^{E^h(X,U)}-1)$, and $J$ and $I$ are as for \textup{(a)}.

\item[\textup{(v)}] For each cell $\square$, the function is $\lambda^v_\square=\chi_\square(A_\square-1)$, and at complex $V'$ it is $\chi_\square(e^{V_\square}-1)$; here $J=I=\{\square\}$.

\item[\textup{(t)}] For each connected set $Y$ of cells, the function is $e^{T_Y}-1$, and $J=I=Y$.

\item[\textup{(b)}] For each cell $\square$, the function is $-b_\square$, and $J=I=\{\square\}$.
\end{enumerate}

\emph{Zone letters.} The objects are of four sorts: a forced class $f_j$, whose cells are those of $f_j$; a cell; a polymer of kind p, that is, some $X\in\mathcal P$; and a polymer of kind p$'$, that is, some $X\in\mathcal D$. Two objects with cells are adjacent when they have cells at distance at most $4r_1+2$. An object with a cell $\square$ is adjacent to a polymer $X$ of kind p if and only if $X$ contains a cube of the fine set under $\square^{+r_1}$, and to a polymer $X$ of kind p$'$ if and only if $X^{\approx}$ meets that fine set.

A \emph{zone letter} $\zeta=(\mathsf J,t,\mathfrak X,\mathfrak Y)$ is a set of objects that is connected for the adjacency and has these properties: $(\mathsf J,t)$ is a cell datum; each polymer in it is adjacent to some object with cells; and the polymers $X$ of $\mathfrak X$, all of kind p, are at mutual distances $>s(x)$. The set $\mathfrak Y$ collects its polymers of kind p$'$. We call $\zeta$ \emph{forced} when $\mathsf J$ is non-empty, and \emph{admitted} when $(\mathsf J,t)$ is admitted. We put \[ 1_{{\rm ind},\zeta}=1\{D(U)\cap\hat{\mathfrak Z}\text{ lies in the union of }\mathfrak X,\text{ and each }X\text{ of }\mathfrak X\text{ holds a plaquette of }D(U)\cap\hat{\mathfrak Z}\} \] and $\Lambda_\zeta:=\hat\lambda_{(\mathsf J,t)}\cdot1_{{\rm ind},\zeta}$, and we give $\zeta$ the function \[ \lambda_\zeta:=\Lambda_\zeta\cdot\prod_{\mathsf J}e^{n_j}\cdot\prod_{\mathfrak X}w(X,U)\cdot\prod_{\mathfrak Y}(e^{1E(X)(U)}-1). \]

The core $J_\zeta$ is built from the marks of the classes of $\zeta$, from its tested cells $t$ and from the sets $Y_0(X)$, linked by diagonal paths that follow a spanning tree, fixed by the order, of the graph whose edges are the adjacencies together with, within each class, the links of the marks of that class. The set $I_\zeta$ is the union of $\mathfrak b^{+2r_1}$, of the cells that meet $X^{\approx\approx}$ for $X$ in $\mathfrak X$ or $X^{\approx}$ for $X$ in $\mathfrak Y$, and of $J_\zeta$.

Two zone letters are \emph{compatible} when each object of the one differs from, and is not adjacent to, every object of the other, and when the polymers of kind p of the two letters stay at distance $>s(x)$ from one another. In addition, the following pairs are incompatible: $\zeta$ and a letter \textup{(b)} at a cell of $\mathfrak b_\zeta^{+2r_1}$; $\zeta$ and a letter \textup{(v)} at a cell of $\mathfrak Z_\zeta$; $\zeta$ and a letter \textup{(a)} or \textup{(h)} of a polymer $X$ for which $X^{\approx}$ meets $\hat{\mathfrak Z}_\zeta$; and the two letters \textup{(a)} and \textup{(h)} of the same $X$.

We write $\alpha\sim\beta$ if ${\rm dist}(I_\alpha,I_\beta)\le r_1$. The universal alphabet $\mathfrak A$ consists of the letters \textup{(a)}, \textup{(h)}, \textup{(v)}, \textup{(t)} and \textup{(b)}, and of the zone letters without a class.

The letters \emph{actual} at $V'$ are the following. A forced $\zeta$ is actual if $f\setminus f_\zeta$ has no cell within the distance $4r_1+2$ of $\mathfrak b_\zeta$, and a zone letter of $\mathfrak A$ is actual if $f$ has no cell within the distance $4r_1+2$ of $\mathfrak b_\zeta$. A letter \textup{(b)} is actual if it lies off $f^{+2r_1}$, and a letter \textup{(v)} if it lies off $f^{+r_1}$. A letter \textup{(a)} or \textup{(h)} of $X$ is actual if $X^{\approx}$ lies off the fine set under $f^{+r_1}$. Every letter \textup{(t)} is actual.

\emph{Clusters.} The clusters, the function $K(\mathfrak B)=E_\xi\prod_{\mathfrak B}\lambda_\alpha$, and $\Phi(C)$, $\hat\Phi(C)$, $n(C)$ and $d(C)$ are those of \Cref{Y-G-L18} (Step~8 of the proof of \Cref{Y-G-ThA}), defined for $\mathfrak A$ with $\mathsf r=r_1$. The core $\operatorname{co}(C)$ of a cluster $C$ is obtained by linking the cores $J_\alpha$ of its letters with diagonal paths that follow a spanning tree of $\sim$ (fixed by an order). A cluster $C$ is \emph{constant} if all its letters are letters \textup{(b)}, and \emph{analytic} if all its letters are letters \textup{(a)}, \textup{(v)} or \textup{(t)}; an analytic cluster is \emph{new} if it has no letter \textup{(a)}, and \emph{old} otherwise. All remaining clusters are \emph{hard}.

\emph{The formulas for the output.}
\begin{enumerate}
\item[\textup{(Fm1)}] On $\mathcal U_Y(x')$ we put \[ E'^{\rm new}(Y,V'):=\sum_{\operatorname{co}(C)=Y,\ {\rm new}}[\Phi(C)(V')-\Phi(C)(1)]+\tilde{\mathsf D}(Y,V')-P(Y,V'), \] \[ E'^{\rm old}[u](Y,V'):=\sum_{\operatorname{co}(C)=Y,\ {\rm old}}[\Phi(C)(V')-\Phi(C)(1)], \] and $E'^a:=E'^{\rm new}+E'^{\rm old}[u]$.

\item[\textup{(Fm2)}] We put $\mathsf V':=\mathsf V_0+\sum_{\rm analytic}\Phi(C)(1)+\sum_{\rm constant}\Phi(C)$ with \[ \mathsf V_0:=\log[\mathsf c_0\,z(ax)^{-3\#\mathbb T'}(2\pi/x)^{\frac32\#I^\circ}(\det C^1)^{1/2}], \] where $\mathsf c_0=(2\pi^2)^{-\#I^\circ}$ (item~C.I~4 of \S\ref{Y-C-1}, paragraph ``The step, the chart'') and $z$ is the normaliser of $k$.

\item[\textup{(Fm3)}] For $G$-valued fields given on the bonds of $Y^{\approx}$ and tame there, where each such field enters through its extension by $1$ off these bonds (Step~6 of the proof of \Cref{Y-H-C9}, paragraph ``The extension by $\mathbf 1$''), we put $E'^h(Y,V'):=\sum_{\operatorname{co}(C)=Y,\ {\rm hard}}\Phi(C)(V')$.

\item[\textup{(Fm4)}] A \emph{rough polymer} at $V'$ is a $\sim$-connected set $\gamma$ of letters that are pairwise compatible and all actual at $V'$, and among which there is a forced letter. A cluster $C$ is \emph{on} at $V'$ if it is constant or if $V'$ is tame on $\operatorname{co}(C)^{\approx}$, and \emph{off} otherwise. We put $\Phi^\natural(C):=\Phi(C)(V')$ if $C$ is on, $\Phi^\natural(C):=\Phi(C)(1)$ if $C$ is analytic and off, and $\Phi^\natural(C):=0$ if $C$ is hard and off.

Let $\mathsf Y$ be a finite set of cells. We put $w'(\mathsf Y,V',e):=0$ if $\mathrm{mk}'(V')$ has a cell in $\mathsf Y^{+(s'-1)}\setminus\mathsf Y$. The radius is $s'-1$ because the plaquettes based in a cell decide whether that cell is a mark, and these plaquettes involve bonds of the next layer; this is the clause of \textup{(F4)} in Step~12 of the proof of \Cref{Y-G-ThA}.

In the remaining case, $w'(\mathsf Y,V',e)$ is the sum of \[ \prod_{\gamma\in\mathfrak G}K(\gamma)\cdot\prod_{C\in\mathfrak N}(e^{-\Phi^\natural(C)}-1)\cdot\prod_{Y\in\mathfrak N^\flat}(e^{\tilde{\mathsf D}(Y,V')}-1) \] over all triples $(\mathfrak G,\mathfrak N,\mathfrak N^\flat)$ with the following properties: $\mathfrak G$ is a set of rough polymers at $V'$, no two of which are related by $\sim$, and the forced letters of $\mathfrak G$ have as their marks exactly $\mathrm{mk}'\cap\mathsf Y$, with each class appearing once; $\mathfrak N$ is a finite set of clusters of $\mathfrak A$ in which every cluster has some letter related by $\sim$ to some letter of $\mathfrak G$; $\mathfrak N^\flat$ is a set of long sets $Y$; the cores $J_\alpha$, $\operatorname{co}(C)$ and $Y$ have the union $\mathsf Y$; and the graph that joins two elements whenever their cores lie within the distance $s'$ of each other is connected.

Finally, $w'(\mathsf Y,V'):=\int_{\mathfrak E}w'(\mathsf Y,V',e)\,de$.
\end{enumerate}

The output is $F_\ell(u):=(E'^a,E'^h,w')$.
\end{definition}

\begin{theorem}[Item 7: one drawn step for complex members within wider caps]\label{Y-H-C9}
In the setting of the letters of item~1, taken in the order \textup{(O0)} to \textup{(O13)}, in the setting \textup{(SD)} of item~2, whatever its data, and in the setting of the inputs \textup{($\Gamma$2)} to \textup{($\Gamma$7)} of item~4, let $x\ge x_9$, let $1\le A\le16$ and let $u\in M_A^{\rm su}(x)$. Then all series in \textup{(Fm1)} to \textup{(Fm4)} converge absolutely, and \textup{(a)} to \textup{(e)} hold.
\begin{enumerate}
\item[\textup{(a)}] The identity $(\rho[u;\ell])'=e^{\mathsf V'}\cdot\rho[F_\ell(u);\ell'^{\rm c}]$ holds at every $V'$, and $V'$ is its only argument.

\item[\textup{(b)}] By \textup{(Fm1)}, $E'^a=E'^{\rm new}+E'^{\rm old}[u]$, and $E'^{\rm new}$ is one and the same for all $u$. For every cell $\square$, \[ \sum_{Y\ni\square}e^{\kappa(\#Y-1)}\sup_{\mathcal U_Y(x')}|E'^{\rm new}(Y,\cdot)|\le0.35\,c_a(x'), \] \[ \sum_{Y\ni\square}e^{\kappa(\#Y-1)}\sup_{\mathcal U_Y(x')}|E'^{\rm old}[u](Y,\cdot)|\le\vartheta_w\,A\,c_a(x'),\qquad\vartheta_w=0.248\le\frac14, \] and \[ \sum_{Y\ni\square}e^{\kappa(\#Y-1)}\sup|E'^h(Y,\cdot)|\le c_h(x'). \] Moreover, at every $V'$, \[ |w'(\mathsf Y,V')|\le\exp\{-c_\Xi x'\varepsilon'\sum_{P\in D'\cap\mathsf Y}\phi'_P-\kappa(|\mathsf Y|_c-1)\}, \] which is the strong cap at $x'$. Of \textup{(ca)}, \textup{(ch)} and \textup{(c$\Xi$)}, the clauses without a size all hold at $x'$. Hence $F_\ell(u)\in M^{\rm su}_{\max\{1,\,0.35+0.248A\}}(x')$ with the strong cap, and at $A=2$ the bound on $E'^a=E'^{\rm new}+E'^{\rm old}[u]$, in units of $c_a(x')$, is $0.35+0.496=0.846\le1$.

\item[\textup{(c)}] Let $z\mapsto u_z\in M_{16}^{\rm su}(x)$ be affine on an open set of $\mathbb C$. Along it, the number $\mathsf V'_z$ and every value $E'^a_z(Y,V')$ with $V'\in\mathcal U_Y(x')$, $E'^h_z(Y,V')$ and $w'_z(\mathsf Y,V')$ depend holomorphically on $z$.

\item[\textup{(d)}] If $u$ is real, then $F_\ell(u)$ and $\mathsf V'$ are real; if, besides, $\rho[u;\ell]>0$, then $(\rho[u;\ell])'>0$.

\item[\textup{(e)}] For $u_1,u_2\in M_2^{\rm su}(x)$, \[ \mathsf D_{x'}(F_\ell(u_1),F_\ell(u_2))\le0.4\,\mathsf D^{\rm in}_x(u_1,u_2)\le0.4\,\mathsf D_x(u_1,u_2); \] the inputs are thus measured in the weak distance and the outputs in the strong one.
\end{enumerate}

Conclusions \textup{(a)}, \textup{(b)}, \textup{(c)} and \textup{(e)} are independent of the reality, of the sign and of the positivity of the member $u$.
\end{theorem}

\begin{remark}[Item 8: uses of \Cref{Y-H-C9} in other theorems]
In Steps~1, 6, 9 and 12 of the proof of \Cref{Y-G-ThA}, the formulas \textup{(F1)} to \textup{(F4)} of that theorem, which use a different test $\chi_X$, define a map of members other than the map $F_\ell$ of \textup{(Fm1)} to \textup{(Fm4)}. Each of the two maps represents $\rho'$.

\Cref{Y-I-ThC}, which is the second half of the change of rows, uses \Cref{Y-H-C9}, applied to inputs in $M_2^{\rm su}$ for which the size of $w$ is bounded by the weak cap alone.
\end{remark}

\begin{proof}[Proof of \Cref{Y-H-C9}]
\emph{Step 0: the floor \textup{(Y*)}, the constants and the radii.}
\begin{enumerate}
\item[\textup{(i)}] Each of the conditions \textup{(Y.0)} to \textup{(Y.9)} holds for all large $x$; hence $S_9$ is not empty, and $x_9=1+\inf S_9$ exists. The number $x_9$ is computed from \textup{(O0)} to \textup{(O12)} alone; it depends on none of $A$, $u$, $z$ and $N$, and it depends on $\ell$ only through the numbers $b_*$, $K_*$, $\lambda_*$ and the range.

\item[\textup{(ii)}] The inequalities \textup{(Geo)}: for $R\ge6$ we have $2k_\chi\le r_1-1$, $k_t\le r_1$, $M r_1\ge R_v\ge R_3$, $3r_w+\varrho_L+1\le R_2-1$ and $\lceil R_2/M\rceil\le r_1-8$.

\item[\textup{(iii)}] The inequalities \textup{(Rad)}, which follow from \textup{(Y.7)}: $4r_1\le r'$, $r_1\le r'/5$, $2\varrho_*+r_1\le r'$, $2(s_0+\varrho_*)+r_1\le5r'\le s'$, $2s_0+4r_1+2\le5r'$, $s_0+r_1+\lceil r(x)/L\rceil+2\le5r'$, $\omega\le M r'$, $3\omega+1\le\ell_w(x')$ and $2(s_0+k_t)+1\le s(x)+1$.

\item[\textup{(iv)}] We have $T>1$, $r_\Xi\ge\frac34(\frac12-d_0-c_\Xi)>0$, $u_t,\upsilon_f,\upsilon_*<1$ and $2\upsilon_f^{1/4}\le1$; moreover $2D_w+1<N$ and $2\omega+1<N$, and $\varrho(x)/L+2\le\varrho_L+2<N/2$ and $2R_3+1<N$.
\end{enumerate}

We prove \textup{(ii)}. First, $M r_1\ge R_2+3(\varrho_L+2)+5M=2R+3\varrho_L+8+5M\ge R_v=R_3+\varrho_L+2$. Next, $2k_\chi\le2R_\chi/M+2=R_2/M+(2\varrho_L+2)/M+2\le\lceil R_2/M\rceil+2r_l+2\le r_1-1$, and $k_t\le k_\chi$. As $r_l\ge1$, we have $\lceil R_2/M\rceil=r_1-3r_l-5\le r_1-8$. Finally, for $R\ge6$, $3r_w+\varrho_L+1\le3R/2+\lceil R/3\rceil+1\le2R+1$.

We prove \textup{(iii)}, with the help of \textup{(Y.7)}. We have $2\varrho_*+r_1=5r_1+2\lceil s(x)/L\rceil+4$, $2(s_0+\varrho_*)+r_1=4r'+(2\varrho_*+r_1)$ and $2s_0+4r_1+2=4r'+4r_1+2$, and moreover $s_0+r_1+\lceil r(x)/L\rceil+2\le3r'$. Because $M\lceil r(x)/L\rceil\ge\lceil M r(x)/L\rceil$ and $2M\ge M+2$, we obtain $\omega\le R_3+D_w+M(\lceil r(x)/L\rceil+2)\le M r'$, and so $3\omega+1\le3M r'+1\le6M r'+12$. Lastly, $2s_0+2k_t+1\le4r'+0.4\,r'+1\le4.4\,r(x)+1\le8r(x)+1$, since $r'\le r(x)$.

We prove \textup{(iv)}. We have $T=L^{3/2}/(19\mathsf N_{\rm c})\ge(1.2\cdot10^4\mathsf N_{\rm c}^2)^3/(19\mathsf N_{\rm c})>10^{17}>1$. By \textup{(O3)}, $t_\rho\le\frac18(\frac12-d_0-c_\Xi)$, and together with \textup{(0e)} this yields $r_\Xi\ge\frac34(\frac12-d_0-c_\Xi)\ge\frac38(\frac12-d_0)$. From \textup{(Y.4)} we get $u_t^{1/2}<1$ and $\upsilon_f^{1/2}\le\frac14$. Moreover, $N\ge2M s(x)=16M r(x)>2M r'+1\ge2\omega+1\ge2D_w+1$.

The same bounds give three further clauses: the clause $\varrho/L+2<N/2$ on the rows (item~C.I~1.1 of \S\ref{Y-C-1}), the clause $\varrho_L+2<N/2$ on the rows (item~C.II~1.2 of \S\ref{Y-C-1}), and the inequality $2R_3+1<N$ in \textup{(P0)} (item~C.II~4 of \S\ref{Y-C-1}, paragraph ``Floors''). Indeed, $\omega\ge R_3+2=2R+\varrho_L+7$; hence $\varrho/L+2\le\varrho_L+2<\omega<N/2$ for every range $\varrho\le\varrho(x)=R$, and $2R_3+1<2\omega+1<N$. These are the clauses on $N$ of the setting in which \S\ref{Y-C-1} states \textup{($\Gamma$2)} and \textup{($\Gamma$3)}. For empty rows the range is $0$, while $\varrho_L$ remains the one belonging to $\varrho(x)$.

We prove \textup{(i)}, condition by condition. As to the orders, $R$, $r(x)$, $r'$, $r_1$, $s'$, $\varrho_*$, $\omega$, $D_w$ and $\ell_w$ are $O(\log x)$; $\varepsilon,\varepsilon'=O((\log x)^{\mathfrak p_0}x^{-1/2})$; and $\mathsf g_x\ge c(\log x)^{-2}$. By \textup{(M2)}, $(\kappa+4)s'\le(\kappa+4)(160c_R\log x/M+8)\le\frac1{12}\log x+8(\kappa+4)$.
\begin{enumerate}
\item[\textup{(0a)}] That the floors of \textup{(0a)} hold for all large $x$ is stated in each of Sections~\ref{Y-C-1}, \ref{Y-D-1} and \ref{Y-E-1}, each section doing so for its own floors, and item~G of \S\ref{Y-F-1} states that $x_{\rm fl}$ exists; of \textup{(P0)}, the clauses $x\ge L^2$ and $R\ge6$ alone are floors, by \textup{(0a)}.

\item[\textup{(0b)}] We have $17T\varepsilon'\to0$, $\mathsf K_{\rm c}(17T\varepsilon')^2/\varepsilon=O(\varepsilon)$, $p_0'/p_0=(1-\log L/\log x)^{\mathfrak p_0}\to1$ and $\mathsf a_{\rm c}\varepsilon'\to0$.

\item[\textup{(0c)}] Here $3D_w\le18M r(x)/L+43$, and $6M r'+12$ is $\ell_w(x')$. For all large $x$ we find that $18M r(x)/L+43$ is at most $\ell_w(x')$, using $L>10^{13}$ and $r(x)\le r'+1+20c_R\log L/M$, and the first condition of \textup{(0c)} follows.

\item[\textup{(D2), (D3), (C0) to (C2)}] In each of these conditions the left side is a number computed from the letters times $(\log x)^{k+\mathfrak p_0}x^{-1/2}$ with $k\le2$ (in \textup{(D3)}, after division by $\varepsilon$, it is $O(\varepsilon)$), and the right side is a number free of $x$ or is at least $c(\log x)^{-2}$, since $\mathsf g_x\ge c(\log x)^{-2}$.

\item[\textup{(0e)}] The number $e_1^\flat$ is $\varepsilon'$ times a number free of $x$.

\item[\textup{(Y.1)}] We have $c_a$, $c_h$, $p_0^3x^{-1/2}\to0$.

\item[\textup{(Y.2)}] The left side is of the order of $(\log x)^3x^{-1/8}$.

\item[\textup{(Y.3)}] Before the bracket of \textup{(Y.3)} there are the three factors $x^{1/4}$, $50(2s'+1)^3=O((\log x)^3)$ and $e^{(\kappa+1)s'}$, and the bound on $(\kappa+4)s'$ given above yields $e^{(\kappa+1)s'}\le e^{8(\kappa+4)}x^{1/12}$. Inside the bracket, the power $x^{-1/2}$ occurs in the term with $c_a$ and in the one with $K_v^\flat$, each of which also contains the factor $p_0^3=O((\log x)^{3\mathfrak p_0})$. For the products of these two terms with the three factors, the exponents of $x$ add up to $\frac14+\frac1{12}-\frac12=-\frac16<-\frac18$, and $x^{-1/6+1/8}=x^{-1/24}$ absorbs the powers of $\log x$. For the term $x^{-20}$ of the bracket, the sum of the exponents is $\frac14+\frac1{12}-20$, and $c_h(x)$, $u_b$ and $u_1$ tend to $0$ as $x\to\infty$ at the rate $e^{-c(\log x)^{2\mathfrak p_0}}$.

\item[\textup{(Y.4)}] In the third term, the power of $x$ is $1/12-1/2=-5/12<-1/4$.

For the first and the second term, \textup{(W)} gives $\mathrm{py}=d_0p_0^2/2$, and therefore \[\log(1/u_t)=\min\{\rho_Q/2-c_\tau,\ \mathrm{py}/(768(2r_1+1)^3)\}-\log3v\ge c(\log x)^{2\mathfrak p_0-3}\] with a constant $c>0$. Next, \[\log(1/\upsilon_f)\ge(3/8)(1/2-d_0)A_0^2(\log x')^{2\mathfrak p_0}-O((\log x)^3).\] Together these two bounds give \[\log(1/\upsilon_*)=L^{-3}\min\{(1/2)\log(1/u_t),(1/4)\log(1/\upsilon_f)\}\ge cL^{-3}(\log x)^{2\mathfrak p_0-3},\] again with a constant $c>0$, where $2\mathfrak p_0-3>1$. This lower bound is to be compared with $(1/4+1/12)\log x+O(\log\log x)$, and since the power of $\log x$ in it exceeds $1$, it is the larger of the two for all large $x$. It is here that the strict inequality $\mathfrak p_0>2$ is used, for it is what gives $2\mathfrak p_0-3>1$.

The second inequality of \textup{(Y.4)} follows from the bound on $\log(1/\upsilon_f)$ obtained above.

\item[\textup{(Y.5)}] We have $p_0^2-p_0'^2\ge2\mathfrak p_0A_0^2\log L\,(\log x')^{2\mathfrak p_0-1}\to\infty$. By \textup{(7e)}, $u_1=e^{-\delta_3p_0^2}$. The rate of $u_b$ is $c_\xi^2/2>\delta_3/2$, because $\delta_3\le3c_\xi^2/(256h_*^2)$ and $h_*\ge1$.

\item[\textup{(Y.6)}] This condition follows from the factor $x^{-1/4}$, from $p_0'\to\infty$, and from the comparison of $x^{-10}$ with $x^{-1/2}$.

\item[\textup{(Y.7)}] For the first inequality, $5r_1+2\lceil s(x)/L\rceil+6\le(10+335/L)c_R\log x/M+{\rm const}$, which is to be compared with $r'\ge20c_R(\log x-\log L)/M$; and $10+335/L<20$ for $L\ge34$. For the second inequality, $R_3+D_w+M(\lceil r(x)/L\rceil+2)\le(2+141/L)c_R\log x+{\rm const}\cdot M$, which is to be compared with $20c_R\log x'$.

\item[\textup{(Y.8)}] This condition follows from the factors $x^{-18}$ and $x^{-19}(\log x)^6$.

\item[\textup{(Y.9)}] This condition follows from $r'\to\infty$.
\end{enumerate}

As to the order of the choices: the conditions depend only on $x$ and on the letters \textup{(O0)} to \textup{(O12)}. The numbers $c_h$ and $E_a$ occur in \textup{(Y.1)} and in \textup{(Y.3)} to \textup{(Y.6)}, and they are fixed before $x_9$. The floor $x_{\rm fl}$ is free of $c_h$, $E_a$, $\ell$ and $N$ (item~G of \S\ref{Y-F-1}).

\emph{Step 1: the formula for the step at every pair.} Let $u\in M_A^{\rm su}(x)$ with $A\le16$. Then $\rho[u;\ell]$ is a bounded, Borel and gauge invariant complex function, and at every pair $(V',e)$ \[(\rho[u;\ell])'(V')=e^{\mathsf V_0}\cdot\exp\{-x'\mathrm{En}'(V')+\sum_Y\tilde{\mathsf D}(Y,V')-\sum_\square1_{\square^\approx}(V')P(\{\square\},V')\}\cdot Z_u,\] where \[Z_u:=e^{-\sum_jw_j}\cdot E_\xi[e^{\sum_YT_Y}\cdot\prod_\square A_\square\cdot\prod_{X\in\mathcal D}e^{1E(X)(U)}\cdot\Xi_w(U)].\]

\begin{proof}[Proof of Step~1]
1. A $G$-valued field that is tame on $X^\approx$ belongs to $\mathcal U_X(x)$: indeed $|U_b|=1<1+\vartheta(x)$, $|U(\pi)|=1<\mathsf b(x)$, and $|\partial U(p)-\mathbf 1|=2\sin(\phi_p/2)\le\varepsilon<2\varepsilon$. Therefore $1E(X)(V)$ is defined at every $V$, and $|1E(X)|\le A(c_a+c_h)e^{-\kappa(\#X-1)}$. The function $\Xi_w$ is a finite sum of bounded terms. The energy $\mathrm{En}_\ell$ is nonnegative and gauge invariant, and the three kinds of members and the indicators are gauge invariant as well. This gives the first assertion.

2. We use the formula of the chart, \Cref{Y-C-L51}. Since the formula is linear in $\rho$, it holds for the real and for the imaginary part of $\rho$ alike, and so \[\rho'(V')=\mathsf c_0\int_{\mathfrak D}\rho(\mathsf U(\eta;V'))\prod_ck_{ax}\prod_{b\in I^\circ}\mathsf j(\eta_b)\,d\eta,\] where $k_{ax}=z(ax)^{-1}e^{-ax(1-\cos\mathsf n_c)}$.

3. Let $X'$ be any $\mathfrak g$-valued function on the coarse plaquettes; in particular, it may be the one of the pair. We put $\mathcal Q(\eta):=(1/2)\langle F^\flat,\mathsf B^VF^\flat\rangle+(a/2)|Q_1^V\eta|^2$, with $F^\flat=EX'+\mathrm{d}^V\eta$ and $\mathsf B^V=1+\ell^{V*}\ell^V$, and we insert $1=e^{-x\mathcal Q(\eta)}e^{x\mathcal Q(\eta)}$. By the definition in item~E11 of \S\ref{Y-F-1}, \[e^{x\mathcal Q(\eta)}\cdot e^{-x\mathrm{En}_\ell(U)}\cdot\prod_ce^{-ax(1-\cos\mathsf n_c)}\cdot1_{\mathfrak D}(\eta)\prod_b\mathsf j(\eta_b)=\prod_\square A_\square,\] because $\langle F^\flat,\mathsf B^VF^\flat\rangle=|F^\flat|^2+|\ell^VF^\flat|^2$ and $n^\flat=-Q_1^V\eta$.

4. Expanded, $\mathcal Q(\eta)=(1/2)\langle\eta,q\eta\rangle+\langle\eta,\mathsf g^VX'\rangle+(1/2)\langle EX',\mathsf B^VEX'\rangle$. With $C^{-1}=q+S$, $\mathsf g^VX'=(q+S)a_0$ and $\eta=\varphi-a_0$ this becomes \[\mathcal Q=(1/2)\langle\varphi,q\varphi\rangle+\langle\varphi,Sa_0\rangle-(1/2)\langle a_0,qa_0\rangle-\langle a_0,Sa_0\rangle+(1/2)\langle EX',\mathsf B^VEX'\rangle\] \[=[\mathsf m_C-(1/2)\langle a_0,Sa_0\rangle]+(1/2)\langle\varphi,C^{-1}\varphi\rangle-(1/2)\langle\varphi,S\varphi\rangle+\langle\varphi,Sa_0\rangle.\]

5. Now let $\varphi=G\xi$. As $C=GG^{\mathsf T}$ and $G$ is invertible by \textup{(2f)}, we get $\langle\varphi,C^{-1}\varphi\rangle=|\xi|^2$ and $d\eta=|\det G|\,d\xi=(\det C)^{1/2}d\xi$; moreover $\int e^{-(x/2)|\xi|^2}d\xi=(2\pi/x)^{(3/2)\#I^\circ}$. By \textup{($\Gamma$4)}, $(x/2)\langle\varphi,S\varphi\rangle-x\langle\varphi,Sa_0\rangle=\sum_YT_Y$.

6. Again by \textup{($\Gamma$4)}, $(\det C^{V'})^{1/2}=(\det C^{\mathbf 1})^{1/2}\exp\{\sum_Y\tilde{\mathsf D}(Y,V')\}$, and \textup{(5a)} rewrites $x\,\mathsf m_C-(x/2)\langle a_0,Sa_0\rangle$. Inserting all of this into the formula of the chart, we obtain the formula of Step~1.

The argument is independent of the reality of $u$ and of its sign.
\end{proof}

\emph{Step 2: tameness of the fine field where the tests are passed and the coarse field is tame, and the domains.} The following assertions hold.
\begin{enumerate}
\item[\textup{(i)}] Let $V'$ be $G$-valued with $\phi'_P\le\varepsilon'$ on the box of radius $R_3$ about $y(p)$, and let $|\xi_b|\le c_\xi\varepsilon$ at the bonds $b$ of $I^\circ$ within $r_w+1$ of $p$. Then $\phi_p(U)\le\varepsilon$.

\item[\textup{(i$'$)}] If $p\in D(U)$, then some bond $\beta$ of $I^\circ$ within $r_w+1$ of $p$ has $|\xi_\beta|>c_\xi\varepsilon$, or some plaquette of $D'$ lies in the box of radius $R_3$ about $y(p)$.

\item[\textup{(ii)}] Let $X\in\mathcal D$ with $\chi_X(\xi)=1$, let $\tau\in\{1,17T\}$, let $B$ be a set of coarse sites that contains the $D_w$-neighbourhood of the sites whose blocks meet $X^\approx$ and the boxes of radius $R_3$ about them, and let $V'\in U^\tau(B;x')$. Then $U$, restricted to $X^\approx$, lies in $\mathcal U_X(x)$ and is holomorphic in $V'$. If $V'\in\mathcal U_Y(x')$ with $Y\supset Y_0(X)$, such a set $B$ can be chosen in $Y^\approx$.

\item[\textup{(iii)}] Let $\mathfrak b\supset f$ and $\mathfrak Z=\mathfrak b^{+r_1}$. On the event that every cell off $\mathfrak b$ passes its tests, $D(U)$ lies in $\hat{\mathfrak Z}$, $\chi_\square=1$ for every $\square$ off $\mathfrak Z$, and $\chi_X=1=1_{X^\approx}(U)$ for every $X$ with $X^\approx$ off $\hat{\mathfrak Z}$.

\item[\textup{(iv)}] If $V'$ is $G$-valued and tame on $\square^\approx$ and $\chi_\square(\xi)=1$, then $A_\square=e^{V_\square}$.
\end{enumerate}

We first prove a lemma and a theorem.

\begin{lemma}[The plaquette of the chart to first order]\label{Y-H-L51}
Let $p$ be a fine plaquette in the plane $\mu\nu$, let $P=(y(p);\mu\nu)$, and let $\tau\ge1$ with $\tau\varepsilon'\le1/40$. For the field there are two cases. In the first case, $V'$ is an $\operatorname{SL}(2,\mathbb C)$ field with $|V'(c)|\le1+\vartheta_0^\flat$ at the two bonds of $P$ at $y(p)$ and with $|\partial V'(P)-\mathbf 1|<2\tau\varepsilon'$, and $X'_P=-i\log\partial V'(P)$. In the second case, $(V',e)$ is a pair with $\phi'_P\le\tau\varepsilon'$. Let $r$ be any number such that $|\eta_b|\le r$ at the bonds of $p$ and $4.1r+2.1\tau\varepsilon'\le1$. Then $|\partial U(p)-\mathbf 1-iF^\flat_p|\le0.72(4.1r+2.1\tau\varepsilon')^2$, where $F^\flat=EX'+\mathrm{d}^V\eta$.
\end{lemma}




\begin{proof}
We use the identity \[\partial U(p)=e^{i\eta_1}e^{iT_2\eta_2}[e^{iX'_P}]e^{-iT_3\eta_3}e^{-i\eta_4},\] with $T_j\in\{1,\operatorname{Ad}_{V'(c)}\}$, which is given in full in Step~4 of the proof of \Cref{Y-E-C6}; the bracket is present if and only if $p$ is an edge plaquette under $P$, and the exponents add up to $F^\flat_p$. Now $|\operatorname{Ad}_gY|\le|g|^2|Y|\le1.01|Y|$. In the first case $|X'_P|\le2.06\tau\varepsilon'$ (item~C.II~1.4 of \S\ref{Y-C-1}); in the second, $|X'_P|=\phi'_P$. Hence the norms of the exponents add up to $S\le4.02r+2.06\tau\varepsilon'\le1$, and \textup{(E)} of item~C.I~1.1 of \S\ref{Y-C-1} gives $|\partial U(p)-\mathbf 1-iF^\flat_p|\le0.72S^2$.
\end{proof}

\begin{theorem}[The chart maps into the fine domain]\label{Y-H-Th53}
Let $\tau\ge1$, $r>0$ and $\varepsilon\le1/40$. Let $X\in\mathcal D$, let $B$ be a set of coarse sites that contains the $D_w$-neighbourhood of the sites whose blocks meet $X^\approx$, let $V'\in U^\tau(B;x')$, and let $\eta$ be a $\mathfrak g_{\mathbb C}$-valued function on $I^\circ$ with the following two properties.
\begin{enumerate}
\item[\textup{(h1)}] $|\eta_b|\le r$ at the bonds of $X^\approx$.
\item[\textup{(h2)}] $|F^\flat_p|\le1.9\varepsilon$ at the plaquettes of $X^\approx$.
\end{enumerate}
Let, moreover, the following floors hold.
\begin{enumerate}
\item[\textup{(D0)}$(\tau)$] $2D_w+1<N$, $3D_w\le\ell_w(x')$, $\tau\varepsilon'\le1/40$.
\item[\textup{(D2)}$(\tau,r)$] $\ell_w(x)(7.1r+18D_w\tau\varepsilon')\le \mathsf g_x/2$.
\item[\textup{(D3)}$(\tau,r)$] $(4.1r+2.1\tau\varepsilon')^2\le0.1\varepsilon$.
\end{enumerate}
Then the restriction of $\mathsf U(\eta;V')$ to the bonds of $X^\approx$ lies in $\mathcal U_X(x)$; it is polynomial in the entries of $V'$ and entire in $\eta$.
\end{theorem}

The condition $\varepsilon\le1/40$ is the floor \textup{(A1)} of item~4 of \S\ref{Y-G0-1}. Here $\varepsilon<\varepsilon'\le1/40$ by the identity $\varepsilon'=L^{1/2}(p_0'/p_0)\varepsilon$ and by \textup{(0b)}.

\begin{proof}
First we note by how much the radii at $x$ exceed those at $x'$: $\vartheta(x)-\vartheta(x')=\vartheta_0^\flat\log L/(\log x\log x')>4\mathsf g_x$ and $\mathsf b(x)-\mathsf b(x')=\log L/(\log x\log x')>4\mathsf g_x$; moreover $\mathsf g_x\le1$ and $\mathsf b(x')+\mathsf g_x<2$ (item~C.II~1.4 of \S\ref{Y-C-1}). By \textup{(D2)}, as $\ell_w(x)\ge18$, we have $r\le \mathsf g_x/255$.

\textup{(u3)}: A plaquette $p$ of $X^\approx$ has $y(p)$ among the sites named in the hypothesis on $B$, and the sites of $P$ lie within $1$ of $y(p)$, hence in $B$. By \textup{(D3)}, $4.1r+2.1\tau\varepsilon'\le1$. So \Cref{Y-H-L51}, together with \textup{(h2)} and \textup{(D3)}, yields $|\partial U(p)-\mathbf 1|\le1.9\varepsilon+0.072\varepsilon<2\varepsilon$.

\textup{(u1)}: Each $U_b$ is $\mathbf 1$, $e^{i\eta_b}$ or $e^{i\eta_b}V'(c)$. Hence \[|U_b|\le e^r(1+\vartheta(x')+\mathsf g_x)<1+\vartheta(x')+\mathsf g_x+2.1r<1+\vartheta(x')+2\mathsf g_x<1+\vartheta(x).\]

\textup{(u2)}: Let $\pi$ be a path in $X^\approx$ of $m\le\ell_w(x)$ bonds $b_1,\dots,b_m$ through the sites $z_0,\dots,z_m$, and put $y_i:=y(z_i)$ and $o_1:=y_0$. Since $|z_i-z_0|_\infty\le m$, we get ${\rm dist}(y_i,o_1)\le\lceil m/L\rceil\le D_w-1$. Let $B_1$ be the box of radius $D_w-1$ about $o_1$. It is embedded, as $2D_w+1<N$; it lies in $B$; and $3(D_w-1)+1\le\ell_w(x')$. By \textup{(e2)} and \textup{(e3)}, the field $V'$ has on $B_1$ the properties \textup{(h2)} and \textup{(h3)} of item~C.II~1.4 of \S\ref{Y-C-1}, with the value $\hat\varepsilon=\tau\varepsilon'$ of their parameter.

We put $\mathsf k(y):=V'(\Gamma(o_1,y))$ and $\mathsf W(c):=\mathsf k(c_-)V'(c)\mathsf k(c_+)^{-1}$. By \Cref{Y-C-L67}\,(a), $|\mathsf k(y)|=|\mathsf k(y)^{-1}|\le2$ and $|\mathsf W(c)|=|\mathsf W(c)^{-1}|\le e^{18D_w\tau\varepsilon'}$. Let $U(b_i)$ be the factor of $b_i$ in $U(\pi)$, taken in its direction, and put $g_i:=\mathsf k(y_{i-1})U(b_i)\mathsf k(y_i)^{-1}$. As $\mathsf k(o_1)=\mathbf 1$, this gives $U(\pi)=g_1\cdots g_m\cdot\mathsf k(y_m)$.

If $b_i$ stays in a block, $g_i$ is $\mathbf 1$ or $\mathsf ke^{\pm i\eta}\mathsf k^{-1}$. If $b_i$ crosses $c$ forwards, $g_i=[\mathsf k(c_-)e^{i\eta}\mathsf k(c_-)^{-1}]\mathsf W(c)$; if it crosses $c$ backwards, $g_i=\mathsf W(c)^{-1}[\mathsf k(c_-)e^{-i\eta}\mathsf k(c_-)^{-1}]$. Consequently $|g_i|\le(1+4(e^r-1))e^{18D_w\tau\varepsilon'}\le e^{7.1r+18D_w\tau\varepsilon'}$. Further, $|\mathsf k(y_m)|<\mathsf b(x')+\mathsf g_x$ by \textup{(e2)}. With \textup{(D2)} we conclude that \[|U(\pi)|<e^{\mathsf g_x/2}(\mathsf b(x')+\mathsf g_x)<\mathsf b(x')+2.3\mathsf g_x<\mathsf b(x).\]

The argument is independent of the rows of $\ell$.
\end{proof}

\begin{proof}[Proof of Step~2]
\begin{enumerate}
\item[\textup{(i)}] Since $\eta=\varphi-a_0$ and $\varphi=G\xi$, we have $F^\flat_p=(EX'-\mathrm{d}^Va_0)_p+(\mathrm{d}^VG\xi)_p$. A $G$-valued field has \textup{(e1)} and \textup{(e2)}, and $|\partial V'(P)-\mathbf 1|=2\sin(\phi'_P/2)\le2\varepsilon'$. We apply \textup{($\Gamma$3)} with $\hat\varepsilon=\varepsilon'\le L^{1/2}\varepsilon$: the first term is at most $\mathsf N_{\rm c}L^{-3/2}\varepsilon+\mathsf K_{\rm c}\varepsilon'^2\le(10^{-4}+0.12)\varepsilon$, by \textup{(L2$^{\rm su}$)} and by \textup{(0b)} with $T>1$. By \textup{(2a)} and \textup{(2e)}, the second term is at most $4C_\varphi^\flat c_\xi\varepsilon\le0.6154\varepsilon$.

At the bonds of $p$, $|\eta_b|\le C_\varphi^\flat c_\xi\varepsilon+\mathsf a_{\rm c}\varepsilon'\le r_T$, where $r_T$ is the number of item~5 of \S\ref{Y-H-1}; so \Cref{Y-H-L51}, with \textup{(D3)}, gives $|\partial U(p)-\mathbf 1-iF^\flat_p|\le0.72\cdot0.1\varepsilon$. Adding up, $2\sin(\phi_p/2)=|\partial U(p)-\mathbf 1|\le0.8075\varepsilon\le2\sin(\varepsilon/2)$, and since $2\sin(\cdot/2)$ increases on $[0,\pi]$, assertion \textup{(i)} follows.

\item[\textup{(i$'$)}] This is the contrapositive of \textup{(i)}.

\item[\textup{(ii)}] We apply \Cref{Y-H-Th53}, using \textup{($\Gamma$3)} with $\hat\varepsilon=\tau\varepsilon'\le17T\varepsilon'$. As to \textup{(h1)}: $|(G\xi)_b|\le\sqrt2\|(G\xi)_b\|\le\sqrt2C_\varphi^\flat c_\xi\varepsilon$ and $|a_{0,b}|\le \mathsf a_{\rm c}\cdot17T\varepsilon'$, so $|\eta_b|\le r_T$. As to \textup{(h2)}: the two terms of \textup{(Cen)} give $17T\cdot \mathsf N_{\rm c}L^{-2}\varepsilon'\le(17/19)\varepsilon$ and $\mathsf K_{\rm c}(17T\varepsilon')^2\le0.12\varepsilon$, while $|(\mathrm{d}^VG\xi)_p|\le4\cdot1.01\cdot\sqrt2\cdot C_\varphi^\flat c_\xi\varepsilon\le0.8790\varepsilon$; the sum is at most $1.8938\varepsilon\le1.9\varepsilon$.

The floors hold: \textup{(D0)} by Step~0\,(iv), \textup{(0c)} and \textup{(0b)}, while \textup{(D2)} and \textup{(D3)} are contained in \textup{(0c)}. Holomorphy follows from \textup{(2d)}, from $X'$ being a power series by \textup{(e3)}, and from the chart being polynomial in $V'$ and entire in $\eta$.

It remains to show that, for $V'\in\mathcal U_Y(x')$ with $Y\supset Y_0(X)$, a set $B$ as in \textup{(ii)} can be chosen in $Y^\approx$. The set ${\rm hull}(X)^{+1}$ lies within $M/2$ of $Y_0(X)$ when $X$ is small, and under $Y_0(X)$ otherwise (item~1 of \S\ref{Y-G0-1}). Passing to $X^\approx$ adds $r(x)$ cubes, which is at most $M\lceil r(x)/L\rceil$ coarse units; then $D_w$ and $R_3$ are added, and by the second inequality of \textup{(Y.7)} the resulting set lies in $Y^\approx$. Moreover, at $x'$ the conditions \textup{(u1)} to \textup{(u3)} imply \textup{(e1)} to \textup{(e3)} with $\tau=1$.

\item[\textup{(iii)}] Let $p$ lie in a cell $\square_p$ off $\mathfrak Z$. The bonds within $r_w+1$ of $p$ lie in cells within $k_t\le r_1$ of $\square_p$, hence off $\mathfrak b$, so they pass. The box of radius $R_3$ about $y(p)$ lies in cells within $\lceil R_3/M\rceil\le r_1$ of $\square_p$, hence off $\mathfrak b\supset \mathrm{mk}'$, so no plaquette of $D'$ lies in it. By \textup{(i)}, $\phi_p\le\varepsilon$. For $\chi_\square$, the bonds within $R_\chi$ of $\hat\square$ lie in cells within $k_\chi\le r_1$ of $\square$; for $\chi_X$, the bonds within $r_w+1$ of $X^\approx$ lie within $k_t$ of cells off $\mathfrak Z$. In both cases these bonds lie in cells off $\mathfrak b$ and pass, and $1_{X^\approx}(U)=1$ because $D(U)$ lies in $\hat{\mathfrak Z}$.

\item[\textup{(iv)}] Let $p$ be a plaquette within $\varrho(x)+1$ fine units of a site of $\hat\square$. Then $y(p)$ lies within $\varrho_L+1$ of $\square$; the bonds within $r_w+1$ of $p$ lie within $r_w+\varrho_L+2\le R_\chi$ of $\hat\square$, and they pass; and the box of radius $R_3$ about $y(p)$ lies within $R_v\le M\,r_1\le M\,r'$ of $\square$, where $V'$ is tame. By \textup{(i)}, $p$ is tame.

Hence, by the definitions (items E2 and E11 of \S\ref{Y-F-1}, item~3.8 of \S\ref{Y-E-1}), we have $S_\varepsilon(\phi_p)=\phi_p^2/2=(1/2)F_p\cdot F_p$; $\zeta_p=F_p$ on the reach of $\ell_y$; for each of the two radii of the mask, no rough plaquette lies within that radius of $y$, so that $m_y=1$; $1-\cos\mathsf n_c=1-(1/2)\operatorname{tr}V'(c)\mathsf M(U)(c)^{-1}$; and $|\eta_b|\le r_T<\pi$. Comparing factor by factor, $A_\square=e^{V_\square}$.
\end{enumerate}
\end{proof}

\emph{Step 3: the alphabet.} The following assertions hold.
\begin{enumerate}
\item[\textup{(I)}] $e^{\sum_jn_j}\cdot Z_u=\sum_{\mathfrak B}E_\xi\prod_{\alpha\in\mathfrak B}\lambda_\alpha$, where the sum runs over the sets $\mathfrak B$ of pairwise compatible letters such that $\mathfrak B$ holds every forced class.

\item[\textup{(II)}] \textup{(Act1)}: All letters of such a set $\mathfrak B$ are actual. An actual zone letter is admitted.

\item[\textup{(III)}] Any two incompatible letters $\alpha$ and $\beta$ satisfy $\alpha\sim\beta$.

\item[\textup{(IV)}] Families of letters are independent as soon as their sets $I$ lie pairwise at distance more than $r_1$.

\item[\textup{(V)}] Except for the factors $e^{-w_j}$ and $e^{n_j}$ of a forced letter, every letter $\alpha$ depends on $(V',e)$ on the cells of $I_\alpha^{+r_1}$ alone. Here $w_j$ is determined by $(V',e)$ within $2r'$ cells of $\mathrm{mk}'_j$, and $n_j$ by $(V',e)$ within $3r'$ cells of $\mathrm{mk}'_j$. For $\alpha\in\mathfrak A$, the set $I_\alpha^{+r_1}$ lies in $J_\alpha^{+r'}$. The entries $e_P$ enter at the plaquettes of $D'$ alone.

\item[\textup{(VI)}] At a $G$-valued $V'$, the two formulas for the letter \textup{(a)} of $X$ agree if $V'$ is tame on $Y_0(X)^\approx$, and those for the letter \textup{(v)} of $\square$ agree if $V'$ is tame on $\square^\approx$.
\end{enumerate}

\begin{proof}[Proof of Step~3]
\begin{enumerate}
\item[\textup{(I)}] At every $\xi$ we have \[1=\sum_{\mathfrak b\supset f}\prod_{\mathfrak b\setminus f}b_\square\cdot\prod_{\square\text{ off }\mathfrak b}(1-b_\square),\] the sum running over the sets $\mathfrak b$ of cells that contain $f$. On the support of the term of $\mathfrak b$, Step~2\,(iii) applies with $\mathfrak Z=\mathfrak b^{+r_1}$.

There $D(U)=D(U)\cap\hat{\mathfrak Z}$, and hence $\Xi_w(U)=\sum_{\mathfrak X}\prod_{\mathfrak X}w(X,U)\cdot1\{D(U)\cap\hat{\mathfrak Z}$ lies in the union of $\mathfrak X$, and each $X$ of $\mathfrak X$ holds a plaquette of $D(U)\cap\hat{\mathfrak Z}\}$, where $\mathfrak X$ runs over the finite families in $\mathcal P$ at pairwise distance $>s(x)$. A family with a polymer that holds no cube under $\mathfrak Z$ has indicator $0$, and we drop it.

Further, \[\prod_Xe^{1E(X)}=\prod_{X^\approx\text{ off }\hat{\mathfrak Z}}(1+\lambda^a_X+\lambda^h_X)\cdot\sum_{\mathfrak Y}\prod_{\mathfrak Y}(e^{1E(X)}-1),\] where $\mathfrak Y$ runs over the subsets of $\{X:X^\approx$ meets $\hat{\mathfrak Z}\}$; this follows from $e^{E^a+E^h}=1+(e^{E^a}-1)+e^{E^a}(e^{E^h}-1)$ and from $\chi_X=1_{X^\approx}(U)=1$ for the $X$ with $X^\approx$ off $\hat{\mathfrak Z}$. Likewise $\prod_\square A_\square=[A]^\ell_{\mathfrak Z}\cdot\prod_{\square\text{ off }\mathfrak Z}(1+\lambda^v_\square)$, $\prod_{\square\text{ off }\mathfrak b^{+2r_1}}(1-b_\square)=\sum_{J_b}\prod_{J_b}(-b_\square)$, and $e^{\sum T_Y}=\prod_Y(1+(e^{T_Y}-1))$.

Off the support of the term of $\mathfrak b$, the product $\prod_{\mathfrak b\setminus f}b_\square\cdot\prod_{\mathfrak b^{+2r_1}\setminus \mathfrak b}(1-b_\square)\cdot\sum_{J_b}\prod(-b_\square)$ vanishes. The expanded identity therefore holds at every $\xi$, and every letter is defined at every $\xi$. All sums are finite.

We now pass to components: the set consisting of all classes, the cells of $\mathfrak b\setminus f$, and the polymers of $\mathfrak X$ and of $\mathfrak Y$ splits into its components $\zeta_i$ under adjacency. Every $X$ of $\mathfrak X$ holds a cube under some $\square^{+r_1}$ with $\square\in \mathfrak b$, and every $X$ of $\mathfrak Y$ has $X^\approx$ meeting such a set. So every polymer is adjacent to an object with cells, and each $\zeta_i$ is a zone letter.

Objects with cells in different components are farther than $4r_1+2$ apart. Hence the sets $\mathfrak b_i^{+2r_1}$ are pairwise disjoint, the sets $\mathfrak Z_i$ are pairwise disjoint, and $\mathfrak Z$ is their union; the products $\prod b_\square$ and $\prod(1-b_\square)$ and the factor $[A]^\ell_{\mathfrak Z}$ factorise.

The indicator factorises as well. A plaquette $p$ of $D(U)\cap\hat{\mathfrak Z}_i$ is covered by a polymer $X$ of $\mathfrak X$ that holds the cube of $p$. By \textup{(SD)} that cube lies under the cell of $p$, and this cell lies in $\square^{+r_1}$ for a cell $\square$ of an object of $\zeta_i$. So $X$ is adjacent to that object, and $X\in\mathfrak X_i$. Conversely, a polymer $X\in\mathfrak X_i$ holds a plaquette of $D(U)$ in some $\hat{\mathfrak Z}_k$; it is then adjacent to an object of $\zeta_k$, so that $k=i$.

For the converse direction, let the $\zeta_i$ be pairwise compatible and hold every class. A cell object of $\zeta_i$ inside some $f_k$ is at distance $0$ from the object $f_k$: if $k\in\mathsf J(\zeta_i)$, the definition of a cell datum excludes it; if not, $\zeta_i$ is incompatible with the letter that holds $f_k$. Let $\mathfrak b:=\bigcup_i\mathfrak b_i$. It follows that $\mathfrak b$ contains $f$, that the $\zeta_i$ are the components, under adjacency, of the set of all their objects, and that the polymers of the union of the $\mathfrak X_i$ are at pairwise distance $>s(x)$. The other letters that are compatible with all $\zeta_i$ are exactly the following: the letters \textup{(b)} off $\mathfrak b^{+2r_1}$; the letters \textup{(v)} off $\mathfrak Z$; for $X$ with $X^\approx$ off $\hat{\mathfrak Z}$, the letter \textup{(a)} or the letter \textup{(h)} of $X$, not both; and any letter \textup{(t)}. The factors $e^{-w_j}$ and $e^{n_j}$ are absorbed into the function of the letter of the class $j$.

\item[\textup{(II)}] A cell of $f_k$ within $4r_1+2$ of a cell of $\mathfrak b_\zeta$ makes $f_k$ adjacent to an object of $\zeta$; so either $f_k$ is in $\zeta$, or $\zeta$ is incompatible with the letter that holds $f_k$. For a letter \textup{(b)}, a cell of $f_k^{+2r_1}$ belongs to $\mathfrak b^{+2r_1}$; for a letter \textup{(v)}, a cell of $f_k^{+r_1}$ belongs to $\mathfrak Z$; and for a letter \textup{(a)} or \textup{(h)}, the fine set under $f_k^{+r_1}$ lies in $\hat{\mathfrak Z}$. In each of these three cases, $\mathfrak b^{+2r_1}$, $\mathfrak Z$ and $\hat{\mathfrak Z}$ are the sets of the letter that holds $f_k$.

It remains to see that an actual zone letter is admitted. Let $P\in D'$ lie in a cell within $4r_1$ of $\mathfrak b_\zeta$; this cell is a mark, and so a cell of $f$. By actuality it lies in $f_\zeta$, and for $\zeta\in\mathfrak A$ there is no such $P$. A mark of the class $k$ is farther than $2s_0+4r_1+2$ from the marks of every $j\ne k$, so it does not lie in $f_j$; hence its class is in $\mathsf J(\zeta)$.

\item[\textup{(III)}] If two zone letters have cells within $4r_1+2$ of each other, then ${\rm dist}(\mathfrak b^{+2r_1},\mathfrak b'^{+2r_1})\le2\le r_1$. If a cell $\square$ of $\zeta$ is adjacent to a polymer $X$ of $\zeta'$, then a cell of $\square^{+r_1}$, which lies in $I_\zeta$, meets $X$ or $X^\approx$, and so it lies in $I_{\zeta'}$. If two zone letters share a polymer, their sets $I$ have a common cell. For polymers of kind p within $s(x)$ of each other, $X^{\approx\approx}$ meets $X'$, where $X'$ denotes the second polymer. For a letter \textup{(b)}, \textup{(v)}, \textup{(a)} or \textup{(h)} that is incompatible with $\zeta$, the cell that makes them incompatible lies in both sets $I$. The letters \textup{(a)} and \textup{(h)} of one $X$ have the same set $I$.

\item[\textup{(IV)}] Each letter depends on $\xi$ only on the bonds of the cells of $I_\alpha^{+k_\chi}$. For \textup{(v)} this comes from the radius $R_\chi$. For \textup{(a)}, \textup{(h)}, $w(X,U)$ and $e^{1E(X)}$: $U_b$ depends on $(G\xi)_b$, hence on $\xi$ within $r_w$, and $\chi_X$ depends on $\xi$ within $r_w+1$. The factor $\Lambda_\zeta$ depends on $\xi$ under $\mathfrak b^{+2r_1}$, by \textup{(7a)}. A letter \textup{(t)} depends on $\xi$ within $r_w$ of $\hat Y$, and a letter \textup{(b)} on $\xi$ in $\hat\square$. Since $2k_\chi\le r_1-1$, for ${\rm dist}(I,I')>r_1$ the sets $I^{+k_\chi}$ and $I'^{+k_\chi}$ are disjoint; and the $\xi_b$ are independent.

\item[\textup{(V)}] The variable $\eta_b$ depends on $(V',e)$ within $R_3$ of $y(b)$. For $G\xi$ this follows from \textup{(2b)}, since $r_G\le R_3$; for $a_0$ it follows from the part of \textup{($\Gamma$3)} by which, at every pair, $a_{0,b}$ depends on $V'$ and $X'$ within $R_3$ of $y(b)$ alone. That part holds at every pair, hence also for forced letters and for letters \textup{(t)} near marks, where the hypothesis of \textup{(Cen)} on the curvature fails. The chart depends on $V'$ at $y(b)$.

For \textup{(a)}, \textup{(h)}, $w$ and $e^{1E}$, this dependence reaches at most $R_3\le M\,r_1$ from the cells of $I_\alpha$; for \textup{(v)} it reaches at most $R_v\le M\,r_1$; for \textup{(t)} the claim follows from \textup{(4c)} and $R_3$. The factors $[A]^\ell_{\mathfrak Z}$ and $1_{\rm ind}$ depend on $U$ at the plaquettes within $\varrho(x)+1$ fine units of $\hat{\mathfrak Z}$, hence on $(V',e)$ within $R_3+\varrho_L+2=R_v\le M\,r_1$ of $\mathfrak Z$.

For $w_j$ the claim is \textup{(5b)}. For $n_j$: each $Y$ in its sum lies in $\mathrm{mk}_j'^{+2r'}$, \textup{(4c)} applies, and whether $Y$ is off depends on the marks within $r'$.

For $\alpha\in\mathfrak A$, the set $I_\alpha$ lies in $J_\alpha^{+\varrho_*}$, because $\mathfrak b$ lies in $J_\zeta$ and the cells that meet $X^{\approx\approx}$ lie in $Y_0(X)^{+\lceil s(x)/L\rceil+2}$; moreover $\varrho_*+r_1\le r'$. Finally, $e$ enters through $X'_P$ alone, at the plaquettes $P$ with $\partial V'(P)=-\mathbf 1$.

\item[\textup{(VI)}] For \textup{(a)}, Step~2\,(i) with \textup{(Y.7)} gives $1_{X^\approx}(U)=1$ where $\chi_X=1$. For \textup{(v)}, we use Step~2\,(iv).
\end{enumerate}
\end{proof}

\emph{Step 4: the skeleton bound.} Let $(V',e)$ be any pair, let $\zeta_i$ be pairwise compatible admitted zone letters, let $J_b$ be a finite set of cells off the union of the $\mathfrak b_i^{+2r_1}$, let $\mathsf Y$ be a finite set of cells, and let $0\le\varsigma\le(\gamma/2)x^{-2}$. Then \[\FitDisplay{\dimexpr\linewidth-1pt\relax}{E_\xi[\prod_i\Lambda_{\zeta_i}\cdot\prod_{J_b}b_\square\cdot e^{(x/2)\varsigma|\varphi_{\mathsf Y}|^2}]\le\exp\{4.8v(2/\gamma)\varsigma(2r_1+1)^3\#\mathsf Y+x^{-1/8}(\sum_i\#\mathfrak b_i^{+2r_1}+\#J_b)\}\cdot u_b^{\#J_b}\cdot\prod_ie^\circ_{(\mathsf J_i,t_i)}.}\]

\begin{proof}[Proof of Step~4]
1. Pointwise, $0\le\Lambda_\zeta=\hat\lambda_{(\mathsf J,t)}\cdot1_{{\rm ind},\zeta}\le\hat\lambda_{(\mathsf J,t)}$. Indeed, $\hat\lambda\ge0$ by its form, since it is built of exponentials, indicators and Jacobians $\mathsf j$ (item~E12 of \S\ref{Y-F-1}), and the indicator is at most $1$. Neither of them depends on the member $u$.

2. The pairs $(\mathsf J_i,t_i)$ are cell data, by the definition of a zone letter. Any two of the zone letters $\zeta_i$ are compatible, so no object with cells of the one is adjacent to an object with cells of the other; that is, no cell of $\mathfrak b_i$ lies within $4r_1+2$ of a cell of $\mathfrak b_k$. The cell data are therefore pairwise compatible in the sense of item~E12 of \S\ref{Y-F-1}. They are admitted by hypothesis. Item~E12 admits cell data that are not connected.

3. By parts 1 and 2 we may apply \textup{(7d)}, which follows from \Cref{Y-F-C7}\,(d$^\circ$).

The functions $\hat\lambda_{(\mathsf J,t)}$ make up the skeleton, which is free of members.
\end{proof}

\emph{Step 5: packing of the fine polymers under the weak cap.}

\begin{lemma}[At most $L^3$ fine polymers for each cell of $t$ and each mark of the classes in $\mathsf J$]\label{Y-H-LP}
Let $(V',e)$ be any pair, and let $\zeta=(\mathsf J,t,\mathfrak X,\mathfrak Y)$ be a zone letter that is admitted at $V'$. If $\Lambda_\zeta(\xi)\ne0$ holds for at least one $\xi$, then $\#\mathfrak X\le L^3\cdot(\#t+\sum_{j\in\mathsf J}\#\mathrm{mk}'_j)$.
\end{lemma}

\begin{corollary}[A factor for each fine polymer]\label{Y-H-LPCor} In the bound of Step~4, each factor $e^\circ_{(\mathsf J_i,t_i)}$ may be replaced by the number $e^*_\zeta$ of the zone letter $\zeta=\zeta_i$, where, for a zone letter $\zeta=(\mathsf J,t,\mathfrak X,\mathfrak Y)$, \[e^*_\zeta:={\rm Res}_{(\mathsf J,t)}\cdot\upsilon_f^{-(1/4)\sum_{\mathsf J}\#D'_j}\cdot u_t^{\#t/2}\cdot u_1^{1[t\text{ non-empty}]}\cdot\upsilon_*^{\#\mathfrak X}.\]\end{corollary}

\begin{proof}[Proof of \Cref{Y-H-LP}] Let $\xi$ be such that $\Lambda_\zeta(\xi)\ne0$. Since $\Lambda_\zeta$ carries the factor $\prod_{\mathfrak b^{+2r_1}\setminus\mathfrak b}(1-b_\square)$ and the indicator $1_{{\rm ind},\zeta}$, every cell of $\mathfrak b^{+2r_1}\setminus\mathfrak b$ passes all of its tests at $\xi$, and each $X\in\mathfrak X$ holds a plaquette $p_X$ of $D(U)\cap\hat{\mathfrak Z}$. We write $q_X$ for the cube of $p_X$; it is a cube of $X$, and by \textup{(SD)} it lies under the cell $\square(p_X)$, which is a cell of $\mathfrak Z$.

\begin{enumerate}\item[\textup{1.}] By Step~2\,(i$'$), at least one of the following two alternatives holds.

\begin{enumerate}\item[\textup{(A)}] Some bond $\beta$ of $I^\circ$ within $r_w+1$ of $p_X$ has $|\xi_\beta|>c_\xi\varepsilon$.

\item[\textup{(B)}] Some plaquette $P$ of $D'$ lies in the box of radius $R_3$ about $y(p_X)$.\end{enumerate}

\item[\textup{2.}] Case \textup{(A)}. The cell $\square_\beta$ of $\beta$ lies within $k_t\le r_1$ of $\square(p_X)$, and $\square(p_X)$ is a cell of $\mathfrak Z=\mathfrak b^{+r_1}$; so $\square_\beta$ belongs to $\mathfrak b^{+2r_1}$. As $|\xi_\beta|>c_\xi\varepsilon$, the cell $\square_\beta$ fails its test, and therefore it lies in $\mathfrak b=f_{\mathsf J}\cup t$. If $\square_\beta\in t$, we charge $X$ to $\square_\beta$, and then $q_X$ lies under $\square_\beta^{+k_t}$. If $\square_\beta\in f_j=\mathrm{mk}_j'^{+s_0}$, we charge $X$ to a mark $\square_m$ of $\mathrm{mk}'_j$ within $s_0$ of $\square_\beta$, and then $q_X$ lies under $\square_m^{+(s_0+k_t)}$.

\item[\textup{3.}] Case \textup{(B)}. The cell that holds $P$ is a mark; it lies within $\lceil R_3/M\rceil\le r_1$ of $\square(p_X)$, hence within $2r_1\le4r_1$ of $\mathfrak b$. Condition \textup{(H17)} therefore applies and shows that $P$ belongs to $D'_j$ for some $j\in\mathsf J$. We charge $X$ to that mark $\square_m$; then $q_X$ lies under $\square_m^{+r_1}$, where $r_1\le s_0+k_t$. If $\mathsf J$ is empty, \textup{(H17)} excludes such a $P$, and case \textup{(B)} does not occur.

\item[\textup{4.}] Counting. Put $k=s_0+k_t$. The polymers of $\mathfrak X$ lie at distance $>s(x)$ from one another and $q_X$ lies in $X$, so the cubes $q_X$ have pairwise sup distance $>s(x)$. By the two cases above, and as $k_t\le k$ and $r_1\le k$, every $X$ charged to a cell $\square$ has its cube $q_X$ in the fine set under $\square^{+k}$. This fine set is a box of $((2k+1)L)^3$ cubes; it is embedded, and distances within it are those of the torus, because $2k+1\le4.4\,r(x)+1\le8r(x)\le N/(2M)$. Here the first inequality is $2k+1=2s_0+2k_t+1\le4r'+0.4\,r'+1\le4.4\,r(x)+1$, since $s_0=2r'$, $k_t\le r_1$ by \textup{(Geo)}, $r_1\le r'/5$ by \textup{(Rad)} and $r'\le r(x)$; the last one, $8r(x)=s(x)\le N/(2M)$, is \textup{(Tor)}, that is, $N\ge2Ms(x)$.

We divide this box into subboxes with $s(x)+1$ cubes along each edge. Any two cubes in a common subbox are at distance at most $s(x)$, and so a subbox contains at most one of the cubes $q_X$. Since $2k+1\le s(x)+1$, there are at most $\lceil(2k+1)L/(s(x)+1)\rceil^3\le L^3$ subboxes, and at most $L^3$ polymers are charged to any one cell.

\item[\textup{5.}] Each $X\in\mathfrak X$ is charged either to a cell of $t$ or to a mark of a class of $\mathsf J$. Summing over these cells and marks, we obtain $\#\mathfrak X\le L^3(\#t+\sum\#\mathrm{mk}'_j)$, the bound of \Cref{Y-H-LP}.\end{enumerate}\end{proof}

\begin{proof}[Proof of \Cref{Y-H-LPCor}] If $\#\mathfrak X_i>L^3(\#t_i+\sum\#\mathrm{mk}'_j)$ for some $i$, then $\Lambda_{\zeta_i}$ vanishes identically by \Cref{Y-H-LP}, and the mean of Step~4 equals $0$.

In the remaining case we use $\upsilon_*<1$, $\upsilon_*^{L^3}\ge u_t^{1/2}$, $\upsilon_*^{L^3}\ge\upsilon_f^{1/4}$ and $\#\mathrm{mk}'_j\le\#D'_j$, together with $\#\mathfrak X\le L^3(\#t+\sum\#\mathrm{mk}'_j)$ for every zone letter; they give \[u_t^{\#t/2}\cdot\upsilon_f^{(1/4)\sum\#D'_j}\le\upsilon_*^{L^3(\#t+\sum\#D'_j)}\le\upsilon_*^{L^3(\#t+\sum\#\mathrm{mk}'_j)}\le\upsilon_*^{\#\mathfrak X}.\]

Consequently $u_t^{\#t}\le u_t^{\#t/2}\cdot\upsilon_f^{-(1/4)\sum\#D'_j}\cdot\upsilon_*^{\#\mathfrak X}$ for each zone letter separately, while the factor $u_1$ is left as it stands. Since ${\rm Res}_{(\mathsf J,t)}$ and $u_1^{1[t\text{ non-empty}]}$ are factors of both $e^\circ_{(\mathsf J,t)}$ and $e^*_\zeta$, this is the replacement that the corollary asserts.\end{proof}

\begin{remark}[The idea of Step~5] A rough fine plaquette under $\mathfrak Z$ can arise only near a failed test or near a rough coarse plaquette, and each such place has already contributed a small factor to the bound: a tested cell with the factor $u_t$, a rough coarse plaquette with the factor $\upsilon_f$. A single place carries at most $L^3$ polymers, because $s(x)$, measured in cubes, exceeds $2(s_0+k_t)$, measured in cells.

The price of Step~5 is the exponent $1/L^3$ in $\upsilon_*$, which makes the floor $x_9$ very large.\end{remark}

\emph{Step 6: the product bound, in four cases.} To each letter $\alpha$ we attach the size $e_\alpha:=e^\circ_\alpha\cdot e^{2x^{-1/8}\#I_\alpha}$, where $e^\circ_\alpha$ is given according to the kind of $\alpha$.

\begin{enumerate}\item[\textup{(a)}] For a letter \textup{(a)} at $X$: $e^\circ({\rm a})=4A\,c_ae^{-\kappa(\#X-1)}$.

\item[\textup{(h)}] For a letter \textup{(h)} at $X$: $e^\circ({\rm h})=4A\,c_he^{-\kappa(\#X-1)}$.

\item[\textup{(v)}] For a letter \textup{(v)} at $\square$: $e^\circ({\rm v})=2K_v^\flat p_0^3x^{-1/2}$.

\item[\textup{(t)}] For a letter \textup{(t)} at $Y$: $e^\circ({\rm t})=e_Y^{1/2}=x^{-20}e^{-\kappa^\sharp\#Y/2}$.

\item[\textup{(b)}] For a letter \textup{(b)} at $\square$: $e^\circ({\rm b})=u_b$.\end{enumerate}

For a zone letter $\zeta$, with $e^*_\zeta$ as in \Cref{Y-H-LPCor}, we put \[e^\circ_\zeta=e^*_\zeta\cdot\prod_{\mathsf J}e^{|n_j|}\cdot\prod_{\mathfrak X}A\,e^{-\kappa(|X|_c-1)}\cdot\prod_{\mathfrak Y}4A(c_a+c_h)e^{-\kappa(\#X-1)}.\]

The step asserts that, for a set $\mathfrak B$ of letters, \[E_\xi|\prod_{\mathfrak B}\lambda_\alpha|\le\prod_{\mathfrak B}e_\alpha\] in each of the following four cases.

\begin{enumerate}\item[\textup{(i)}] The letters of $\mathfrak B$ are pairwise compatible and actual at $V'$; the bound then holds at every $e$.

\item[\textup{(ii)}] The letters of $\mathfrak B$ are pairwise compatible letters of $\mathfrak A$ with $J_\alpha$ inside $Y$, and either $V'$ is $G$-valued and tame on $Y^\approx$, or $V'$ is a $G$-valued field on the bonds of $Y^\approx$, tame there, and the left side is taken at its extension $\tilde V'$ by $\mathbf 1$.

\item[\textup{(iii)}] The letters of $\mathfrak B$ are of the kinds \textup{(a)}, \textup{(v)}, \textup{(t)}, with $J_\alpha$ inside $Y$, and $V'\in\mathcal U_Y(x')$.

\item[\textup{(iv)}] The letters of $\mathfrak B$ are letters \textup{(b)}.\end{enumerate}

\begin{proof}[Proof of Step~6] \emph{Uniform bounds on the members.} By \textup{(Y.1)}, $|z|\le1/2$ for $z=E^a,E^h,1E(X)$, and therefore $|e^z-1|\le|z|e^{|z|}\le2|z|$.

\begin{enumerate}\item[\textup{(a)}] For a letter \textup{(a)}, $|\lambda^a_X|\le2A\,c_ae^{-\kappa(\#X-1)}$. At real fields the factor $1_{X^\approx}(U)$ keeps $U$ in $\mathcal U_X(x)$; at complex fields Step~2\,(ii) does so.

\item[\textup{(h)}] For a letter \textup{(h)}, $|\lambda^h_X|\le e^{1/2}\cdot2A\,c_he^{-\kappa(\#X-1)}$.\end{enumerate}

For the factors that a zone letter carries, $|e^{1E(X)}-1|\le2A(c_a+c_h)e^{-\kappa(\#X-1)}$ and $|w(X,U)|\le A\,e^{-\kappa(|X|_c-1)}$ hold at every $U$.

\begin{enumerate}\item[\textup{(v)}] For a letter \textup{(v)} at $\square$, the restriction of $V'$ to $\square^\approx$ lies in $\mathcal U_{\{\square\}}(x')$: in case \textup{(i)} because ${\rm dist}(\square,\mathrm{mk}')>s_0+r_1>r'$, and in cases \textup{(ii)}, \textup{(iii)} because $\square^\approx$ lies in $Y^\approx$. Step~2\,(iv), \textup{($\Gamma$6)} and \textup{(Y.1)} then give $|\lambda^v_\square|\le e^{|V_\square|}-1\le2K_v^\flat p_0^3x^{-1/2}$.

\item[\textup{(t)}] \emph{Real fields.} As $a_0=C\mathsf g^VX'$ and $C=GG^{\mathsf T}$, we have $a_0=G(G^{\mathsf T}\mathsf g^VX')$. A bond lies in $4$ plaquettes, $|EX'|_\infty\le\pi$, and $\|\ell^{V*}\ell^V\|_\infty\le(K_*\mathsf g_*)(3K_*\mathsf g_*)$; these three facts give $|(\mathsf g^VX')_b|\le4\pi \mathsf N_{\mathsf B}$. Together with \textup{(2e)} this yields $|a_{0,b}|\le(C_\varphi^\flat)^2\cdot4\pi \mathsf N_{\mathsf B}=a_\infty$ at every pair $(V',e)$, and since $\hat Y$ has $3v\cdot\#Y$ bonds, $(1/2)|a_{0,Y}|^2\le(1/2)\cdot3v\cdot\#Y\cdot a_\infty^2=n_a\#Y$; here, as in item~E1 of \S\ref{Y-F-1}, $|a_{0,Y}|^2$ is the sum of the $|a_{0,b}|^2$ over the bonds $b$ of $I^\circ$ in $\hat Y$, and the forms with $\|\cdot\|$ in the place of $|\cdot|$ are defined alike.

By \textup{(4b)}, \[\FitDisplay{\dimexpr\linewidth-1pt\relax}{|T_Y|\le x\,e_Y((1/2)|\varphi_Y|^2+|\varphi_Y||a_{0,Y}|)\le x\,e_Y(|\varphi_Y|^2+n_a\#Y)=e_Y^{1/2}z_Y,\qquad z_Y:=x\,e_Y^{1/2}(n_a\#Y+|\varphi_Y|^2),}\] and then $|e^T-1|\le|T|e^{|T|}\le e_Y^{1/2}e^{2z_Y}$.

Moreover $\sum_{Y\ni\square}e_Y^{1/2}\le x^{-20}\sum_n(26e)^{n-1}e^{-50n}\le\varsigma''$, where $e^{-50n}$ comes from $\kappa^\sharp/2=5\kappa+20\ge50$, by \textup{(O5)}. Let $\mathsf Y$ denote the union of the sets $Y$ of the letters \textup{(t)}. Then \[\prod|e^{T_Y}-1|\le\prod e_Y^{1/2}\cdot\exp\{2x\varsigma''(n_a\#\mathsf Y+|\varphi_{\mathsf Y}|^2)\},\] and the exponential is the tilt of Step~4 at $\varsigma=4\varsigma''=4.04\,x^{-20}\le(\gamma/2)x^{-2}$, by \textup{(Y.8)}.\end{enumerate}

We now treat the cases of the claim.

\begin{enumerate}\item[\textup{(i), (ii), (iv)}] With the members and the letters \textup{(t)} bounded as above, what is left is the mean of Step~4, and we bound it by \Cref{Y-H-LPCor}. By compatibility, the cells of the letters \textup{(b)} lie off the sets $\mathfrak b_i^{+2r_1}$. The zone letters are admitted: in case \textup{(i)} by Step~3\,(II), and in case \textup{(ii)} because $\mathsf J$ is empty, $\mathfrak b_\zeta$ lies in $J_\zeta$, $J_\zeta$ lies in $Y$, and $r'\ge4r_1$. Moreover $|\lambda_\zeta|\le\Lambda_\zeta\cdot\prod e^{|n_j|}$ times the bounds found above for the factors $w(X,U)$ and $e^{1E(X)}-1$.

What remains is the harmless factor \[\exp\{[2x\varsigma''n_a+4.8\cdot4.04\cdot v(2/\gamma)x^{-20}(2r_1+1)^3]\#\mathsf Y+x^{-1/8}(\sum\#\mathfrak b_i^{+2r_1}+\#J_b)\},\] and it is at most $\prod e^{x^{-1/8}\#I_\alpha}$, by the second inequality of \textup{(Y.8)} and because $I_\zeta$ contains $\mathfrak b_\zeta^{+2r_1}$.

\item[\textup{(ii)}] \emph{The extension by $\mathbf 1$.} Let $V'$ be a $G$-valued field on the bonds of $Y^\approx$, tame there, and let $\tilde V'=V'$ on these bonds and $\mathbf 1$ elsewhere. Every plaquette of $D'(\tilde V')$ has a site off $Y^{+r'}$, and its base site lies within one site of that site; hence all marks of $\tilde V'$ lie at cell distance $\ge r'$ from $Y$.

By \textup{(Y.7)}, $2\varrho_*+r_1=5r_1+2\lceil s(x)/L\rceil+4\le r'-2$. By Step~3\,(V), a letter of $\mathfrak A$ with $J_\alpha$ inside $Y$ depends on the field in $I_\alpha^{+r_1}$, which lies inside $Y^{+(\varrho_*+r_1)}$ and thus inside $Y^{+(r'-2)}$. Its value is therefore a function of the field on the bonds of $Y^\approx$ alone, not of the extension, and $e$ does not enter, since $\phi'_P\le\varepsilon'<\pi$ there.

A zone letter has $\mathfrak b_\zeta$ inside $J_\zeta$, and $J_\zeta$ inside $Y$. As $4r_1<r'$, no plaquette of $D'(\tilde V')$ lies in a cell within $4r_1$ of $\mathfrak b_\zeta$, and so \textup{(H17)} holds, since $\mathsf J$ is empty and \textup{(H17)} then requires nothing more.

For a letter \textup{(v)}, $\square^\approx$ lies in $Y^\approx$, so the restriction to $\square^\approx$ lies in $\mathcal U_{\{\square\}}(x')$; the plaquettes that Step~2\,(iv) uses have all their sites within $R_v\le M\,r_1$ of $\square$, and are therefore plaquettes of $Y^\approx$, on which $\tilde V'$ is tame. The letters \textup{(a)} and \textup{(h)} carry the factor $1_{X^\approx}(U)$. The bound of Step~4 holds at every pair.

Hence the bound holds at $(\tilde V',e)$ for every $e$, and the value of the left side is a function of $V'$ on the bonds of $Y^\approx$.

\item[\textup{(iii)}] \emph{Complex fields.} Here no zone letters occur. Let $V'\in\mathcal U_Y(x')$. Extended by $\mathbf 1$, it lies in $\mathfrak V$, by \textup{(u1)}.

We write $Y_t$ for the connected set of cells at which a letter \textup{(t)} of $\mathfrak B$ is placed. It is the set $J_\alpha$ of that letter, and so, in case \textup{(iii)}, it lies inside $Y$, the set with $V'\in\mathcal U_Y(x')$. On all of $\mathfrak V$, $\|S_k(Y_t;V')\|\le e_{Y_t}$ by \textup{(4b)}. The pairings are bilinear, and $|\varphi\cdot S\psi|\le\|S\|\|\varphi\|\|\psi\|$.

Let $b$ be a bond in $\hat Y_t$. It is a bond of a plaquette $p$ whose box of radius $R_3$ lies in $Y^\approx$, because $R_3+1\le M\,r'$. On that box, \textup{(u1)} to \textup{(u3)} give \textup{(e1)}, \textup{(e2)} and $|\partial V'(P)-\mathbf 1|<2\varepsilon'$, so \textup{($\Gamma$3)} applies at $\hat\varepsilon=\varepsilon'$ and gives $\|a_{0,b}\|\le|a_{0,b}|\le \mathsf a_{\rm c}\varepsilon'\le1\le a_\infty$, by \textup{(0b)}. Hence $(1/2)\|a_{0,Y_t}\|^2\le n_a\#Y_t$. By this inequality, the bound on $\|S_k(Y_t;V')\|$ and the estimate of the pairings, the two bounds of the real case hold here as well: $|T_{Y_t}|\le e_{Y_t}^{1/2}z_{Y_t}$ and $|e^{T_{Y_t}}-1|\le e_{Y_t}^{1/2}e^{2z_{Y_t}}$, where $z_{Y_t}$ is the $z_Y$ defined for real fields, taken at $Y=Y_t$ and with $\|\varphi_{Y_t}\|$ in the place of $|\varphi_Y|$.

By \textup{(2a)} and \textup{(2e)}, $\|(G\xi)_b\|^2\le(C_\varphi^\flat)^2\sum_{b'\in Y_b}|\xi_{b'}|^2$. Each bond $b'$ lies in at most $3L^3(2r_w+1)^3\le3L^3(R+1)^3$ windows, and so $\|\varphi_{\mathsf Y}\|^2\le(C_\varphi^\flat)^2\cdot3L^3(R+1)^3\cdot|\xi_{\mathsf Y^{+r_1}}|^2$.

For a single bond, $E_\xi e^{(x/2)\tilde\varsigma|\xi_b|^2}=(1-\tilde\varsigma)^{-3/2}\le e^{3\tilde\varsigma}$ at $\tilde\varsigma=4\varsigma''(C_\varphi^\flat)^2\cdot3L^3(R+1)^3\le1/2$, since $-\log(1-\varsigma)\le1.39\varsigma$ on $[0,1/2]$. At most $3v(2r_1+1)^3\#\mathsf Y$ bonds occur, and the resulting bound is the content of the third and fourth inequalities of \textup{(Y.8)}.\end{enumerate}

\end{proof}

\emph{Step 7: the animals at $\varepsilon''=x^{-1/4}$, and condition \textup{(S)}.} The objects of zone letters receive the weights \[\FitDisplay{\dimexpr\linewidth-1pt\relax}{\hat w(\square)=u_t^{1/2}e^{(\kappa+4)(1+s')},\qquad\hat w_p(X)=16\upsilon_*e^{-\kappa(|X|_c-1)+(\kappa+4)(\#Y_0(X)+s')},\qquad\hat w_{p'}(X)=64(c_a+c_h)e^{-\kappa(\#X-1)+(\kappa+4)(\#Y_0(X)+s')}}\] for a cell $\square$, a polymer $X$ of kind p and a polymer $X$ of kind p$'$, respectively, and the cores $\mathrm{core}(\square)=\{\square\}$ and $\mathrm{core}(X)=Y_0(X)$, while $\mathrm{core}(f_j)$ consists of the marks of the class $j$ together with their links; here $\#\mathrm{core}(f_j)\le(1+s')\#\mathrm{mk}'_j$. With these, the following holds.

\begin{enumerate}\item[\textup{(o)}] Each link of a core costs at most $s'$ further cells, so that $\#J_\zeta\le\sum_{o\in\zeta}(\#\mathrm{core}(o)+s')$. The set $I_\zeta$ lies in $J_\zeta^{+s_0+\varrho_*}$, and in $J_\zeta^{+\varrho_*}$ if $\zeta$ is not forced.

\item[\textup{(i)}] Put $\operatorname{fp}(\square)=\{\square\}$ and $\operatorname{fp}(f_j)=f_j$, and let $\operatorname{fp}_p(X)$, $\operatorname{fp}_{p'}(X)$ be the set of cells $\square$ adjacent to $X$. With $a_o:=\varepsilon''\#\operatorname{fp}(o)^{+(4r_1+2)}$ we have $a_o\le\#\mathrm{core}(o)$, indeed $a_{f_j}\le\#\mathrm{mk}'_j$, and, at every cell $\square$, \[\sum_{o:\operatorname{fp}(o)\ni\square}\hat w(o)e^{a_o}\le\varepsilon''/2,\] where the sum runs over cells and polymers. Consequently $\sum_{\zeta\ni o}\prod_{o'\in\zeta}\hat w(o')\le\hat w(o)e^{a_o}$.

\item[\textup{(ii)}] Condition \textup{(S)} of \Cref{Y-G-L18} holds for $\mathfrak A$ with the sizes of Step~6, at every $A\le16$, with $t=x^{-1/4}$, $\varepsilon_2=x^{-1/8}$ and $\mathsf r=r_1$.\end{enumerate}

\begin{proof}[Proof of Step~7] \begin{enumerate}\item[\textup{(o)}] A link has length at most $4r_1+2$ from a cell to a cell, $s_0+4r_1+2$ from a cell to a mark of an adjacent class, and $2s_0+4r_1+2$ from a mark to a mark. From a cell to $Y_0(X)$ it has length at most $r_1+1$ when $X$ is of kind p, since a small $X$ lies under $Y_0^{+1}$ and any other $X$ under $Y_0$, and at most $s_0+r_1+1$ from a mark; when $X$ is of kind p$'$, the bound is $r_1+\lceil r(x)/L\rceil+2$ from a cell, and $s_0$ more from a mark. All of these are $\le5r'\le s'$, by \textup{(Rad)}.

The set $\mathfrak b$ lies in $J^{+s_0}$, so $\mathfrak b^{+2r_1}$ lies in $J^{+s_0+2r_1}$; and the cells that meet $X^{\approx\approx}$ lie in $Y_0(X)^{+\lceil s(x)/L\rceil+2}$.

\item[\textup{(i)}] The set $\operatorname{fp}_p(X)$ lies in $Y_0^{+r_1+1}$, and $\operatorname{fp}_{p'}(X)$ in $Y_0^{+r_1+\lceil r(x)/L\rceil+2}$. Using $8r_1+5\le2(s_0+\varrho_*)+1$, $5r_1+3\le s_0+\varrho_*$ and $5r_1+\lceil r(x)/L\rceil+4\le s_0+\varrho_*$, we find $\#\operatorname{fp}(o)^{+(4r_1+2)}\le(2(s_0+\varrho_*)+1)^3\#\mathrm{core}(o)$.

For $f_j$, the inequality $2(s_0+4r_1+2)+1\le2(2(s_0+\varrho_*)+1)$ costs a factor $8$. Then \textup{(Y.2)} and $x^{-1/4}\le x^{-1/8}/50$ give $a_o\le\#\mathrm{core}(o)$.

If two objects $o,o'$ are adjacent, then $\operatorname{fp}(o')$ meets $\operatorname{fp}(o)^{+(4r_1+2)}$. Therefore $\sum_{o'\text{ adjacent to }o}\hat w(o')e^{a_{o'}}\le\varepsilon''\#\operatorname{fp}(o)^{+(4r_1+2)}=a_o$, and \textup{(An0)} of Step~8 of the proof of \Cref{Y-G-ThA} applies.

\emph{The sums at a cell $\square$.} A cell contributes $u_t^{1/2}e^{\kappa+5}e^{(\kappa+4)s'}$.

A polymer $X$ of kind p holds one of $((2r_1+1)L)^3$ cubes. For small $X$, \[\FitDisplay{\dimexpr\linewidth-1pt\relax}{\sum_n(52e)^ne^{-\kappa(n-1)}e^{\kappa+5}\le217.59\,\allowbreak e^{\kappa+5}}\]; the number is the value at $\kappa=6$, and the sum decreases as $\kappa$ grows.

For $X$ of kind p that is not small, $n>L/2-2$ and $\#Y_0\le8+432n/L$, and the term $n$ is $(52e)^ne^{-\kappa(n-1)}e^{(\kappa+4)\#Y_0+a_o}$. Since $a_o\le\#Y_0$, the bound on $\#Y_0$ gives $(\kappa+4)\#Y_0+a_o\le(\kappa+5)(8+432n/L)$. The factors of the resulting bound that do not depend on $n$ are $e^{\kappa}$, from $-\kappa(n-1)$, and $e^{8(\kappa+5)}$, from $(\kappa+5)\cdot8$; their product is $e^{9\kappa+40}=e^{\kappa+5}\cdot e^{35+8\kappa}$. Hence, relative to $e^{\kappa+5}$, the exponent of the term $n$ is at most $n(4.952+2160/L)+35-\kappa[n(1-432/L)-8]$. For $n>L/2-2$ the bracket is positive, so this bound decreases as $\kappa$ grows. Since $\kappa\ge6$ is chosen in \textup{(O5)} only after $L$ has been fixed in \textup{(O4)}, this bound on the exponent, with the factor $e^{35+8\kappa}$ included, is needed at that $L$ for all $\kappa\ge6$ and, by the monotonicity just established, does not exceed its value at $\kappa=6$. At $\kappa=6$ the bound is $\le-1.04n+83$, the geometric sum costs a factor $<3$, and the contribution lies below $e^{-L/4}$. Altogether, the polymers of kind p give at most $220e^{\kappa+5}$.

A polymer $X$ of kind p$'$ holds one of $((2r_1+1)L+2r(x))^3$ cubes. We take $16(c_a+c_h)e^{\kappa+5}e^{(\kappa+4)s'}$ as the unit for one cube. A small $X$ has $\#Y_0(X)=1$ and $a_o\le1$, and $\sum_{X\ni{\rm cube}}e^{-\kappa(\#X-1)}\le1.2124$; as $64=4\cdot16$, the small $X$ contribute $1.2124\cdot4\le4.85$ in this unit. For $X$ not small, again $n>L/2-2$ and $\#Y_0\le8+432n/L$, and relative to $e^{\kappa+5}$ the exponent of the term $n$ is at most $4.2581(n-1)+35+2160n/L-\kappa[n(1-432/L)-8]$. The bracket is positive for $n>L/2-2$, so this bound, too, decreases as $\kappa$ grows; at $\kappa=6$ it is $\le-1.74n+79$, and the contribution lies below $e^{-L/4}$. Altogether, in this unit, the small $X$ contribute at most $4.85$ and the $X$ that are not small at most $6-4.85$, so the polymers of kind p$'$ give at most $6$, the coefficient of the third term in the first of the inequalities \textup{(Y.4)}.

The three contributions add up to the left side of the first inequality of \textup{(Y.4)}.

\item[\textup{(ii)}] Let $\beta\in\mathfrak A$ with ${\rm dist}(I_\beta,\square)\le r_1$. Then some cell $\square''$ of $J_\beta$ lies within $\varrho_*+r_1\le s'$ of $\square$, which leaves $(2s'+1)^3$ cells for $\square''$. The set $J_\beta$ is connected and has a root $\square_0$, namely its first cell, or its first tested cell for a zone letter. Hence, for non-negative summands, $\sum_{\beta:J_\beta\ni\square''}(\cdot)\le\sum_{\square_0}e^{-{\rm dist}(\square_0,\square'')}\sum_{\beta\text{ rooted at }\square_0}e^{\#J_\beta}(\cdot)$, where $\sum_{\square_0}e^{-{\rm dist}(\square_0,\square'')}\le \mathsf g(1)\le49.98$. Further, $e^{5x^{-1/8}\#I_\beta}\le e^{\#J_\beta}$ by \textup{(Y.2)}, and $e^{d_\beta}=e^{(\kappa+1)s'}e^{(\kappa+1)\#J_\beta}$.

It remains to bound, kind by kind, the sums of $e^\circ_\beta e^{(\kappa+3)\#J_\beta}$ over the letters $\beta$ rooted at a given cell $\square_0$.

\begin{enumerate}\item[\textup{(a), (h)}] For small $X$, this rooted sum is at most $L^3\cdot1.2124\cdot4\cdot16(c_a+c_h)e^{\kappa+3}$.

For $X$ not small, $X$ holds one of the $(L+2)^3$ cubes about $\hat\square_0$; with $n>L/2-2$ cubes and $\#Y_0\le8+432n/L$, the exponent of the term $n$, relative to $e^{\kappa+3}$, is at most $4.2581(n-1)+21+1296n/L-\kappa[n(1-432/L)-8]$. Since the bracket is positive for $n>L/2-2$, this bound decreases as $\kappa$ grows, as for the polymers of kind p in \textup{(i)}; at $\kappa=6$ it is $\le-1.74n+65$, and the contribution lies below $e^{-L/4}$. Altogether, the letters \textup{(a)}, \textup{(h)} give at most $5e^{\kappa+3}L^3\cdot16(c_a+c_h)$.

\item[\textup{(v)}] The rooted sum is at most $2e^{\kappa+3}K_v^\flat p_0^3x^{-1/2}$.

\item[\textup{(t)}] The rooted sum is at most $x^{-20}\sum_n(26e)^{n-1}e^{-(5\kappa+20)n+(\kappa+3)n}\le x^{-20}$.

\item[\textup{(b)}] The rooted sum is at most $e^{\kappa+3}u_b$.\end{enumerate}

A zone letter of $\mathfrak A$ has a tested cell, and ${\rm Res}=1$ for it. Hence $e^\circ_\zeta e^{(\kappa+3)\#J_\zeta}\le u_1\prod_o\hat w(o)$, and by \textup{(i)} at the object $\square_0$ the rooted sum is $\le u_1\hat w(\square_0)e^{a_{\square_0}}\le u_1$, with $a_{\square_0}$ the number $a_o$ of \textup{(i)} at $o=\square_0$; the last inequality holds as $\hat w(\square_0)e^{a_{\square_0}}\le\varepsilon''/2\le1$, by the sum at a cell in \textup{(i)}.

With $t^{-1}=x^{1/4}$, this is \textup{(Y.3)}; the factor $u_1$ enters it with its full power, as no part of $u_1$ is spent on the sums of this step.\end{enumerate}\end{proof}

\emph{Step 8: the first two bounds of \Cref{Y-H-C9}\,(b), and the clauses of \textup{(ca)} at $x'$ without a size.} The two parts of $E'^a$ in \textup{(Fm1)} are $E'^{\rm new}$ and $E'^{\rm old}[u]$. The functions $E'^{\rm new}(Y,\cdot)$ and $E'^{\rm old}[u](Y,\cdot)$ are holomorphic on $\mathcal U_Y(x')$, invariant under $G$-valued gauge transformations and $0$ at pure gauges, and they are functions of the field on $Y^\approx$; moreover, $E'^{\rm new}$ does not depend on $u$. At every cell $\square$, the sums over $Y\ni\square$ of the suprema of their moduli on $\mathcal U_Y(x')$, weighted by $e^{\kappa(\#Y-1)}$, are $\le0.35\,c_a(x')$ for $E'^{\rm new}$ and $\le0.248\,A\,c_a(x')$ for $E'^{\rm old}[u]$.

\begin{proof}[Proof of Step~8] \begin{enumerate}\item[\textup{(Clu)}] For $\alpha\sim\beta$ in $\mathfrak A$ we have ${\rm dist}(J_\alpha,J_\beta)\le2\varrho_*+r_1\le r'\le s'$. For a cluster $C$, therefore, $\#\operatorname{co}(C)\le\sum(\#J_\alpha+s')-s'$ and $d(C)\ge(\kappa+1)(\#\operatorname{co}+s')$, and it follows that $e^{\kappa(\#\operatorname{co}-1)}\le e^{d(C)-(\kappa+1)s'-\#\operatorname{co}-\kappa}$.

\item[\textup{(R)}] If $\alpha$ is a letter of $C$ with root $\square_0$ and $\operatorname{co}(C)\ni\square$, then $\#\operatorname{co}\ge{\rm dist}(\square_0,\square)+1$. We apply \Cref{Y-G-L18} to the alphabet $\{\alpha\in\mathfrak A:J_\alpha\text{ inside }Y\}$, on which the hypotheses of the lemma hold: by Step~7\,(ii), condition \textup{(S)}; by Step~6\,(ii) to (iv), the bound $|K(\mathfrak B)|\le\prod_{\mathfrak B}e_\alpha$; by Step~3\,(III), the fact that two letters which are not compatible are $\sim$; and by Step~3\,(IV), the factorisation $K(\mathfrak B\cup\mathfrak B')=K(\mathfrak B)K(\mathfrak B')$ whenever $\mathfrak B$ has no letter that is $\sim$ to a letter of $\mathfrak B'$. For each kind of letter this gives \[\FitDisplay{\dimexpr\linewidth-1pt\relax}{\sum_{\operatorname{co}(C)\ni\square,\ C\ni\text{ a letter of that kind}}\hat\Phi(C)e^{\kappa(\#\operatorname{co}-1)}\le\sum_{\square_0}e^{-{\rm dist}(\square_0,\square)-1}\sum_{\alpha\text{ of that kind rooted at }\square_0}e_\alpha e^{2\varepsilon_2\#I_\alpha+(\kappa+1)\#J_\alpha-\kappa}.}\]\end{enumerate}

\emph{The clauses without a size.} For letters \textup{(a)}, \textup{(v)}, \textup{(t)}, the mean $K(\mathfrak B)(V')=E_\xi\prod\lambda_\alpha$ is holomorphic on $\mathcal U_Y(x')$: the letters \textup{(a)} are holomorphic by Step~2\,(ii), the letters \textup{(v)} by \textup{($\Gamma$6)}, and the letters \textup{(t)} by \textup{(2d)} and \textup{(4b)}; $X'$ is a power series; and the integrand is dominated by Step~6\,(iii). The convergence is uniform. The dependence on the field is that of Step~3\,(V).

We have $\lambda_\alpha(\operatorname{Ad}_u\xi,V'^u)=\lambda_\alpha(\xi,V')$, by \textup{(2c)}, \textup{(4b)} and \textup{($\Gamma$6)}, because $\mathsf U(\operatorname{Ad}_u\eta;V'^u)$ is the gauge transform of $\mathsf U(\eta;V')$ by $u$, taken constant on blocks, and $E^a$ is invariant. The indicators $\chi$ and the law of $\xi$ are invariant under $\operatorname{Ad}_u$.

Here $u$ is given on the sites of $Y^\approx$ (in part \textup{(8.2)} below, on those of $\mathsf B_\omega$) and is extended by $\mathbf 1$. The transform of the extended field differs from the extension of $V'^u$ only on the bonds that leave $Y^\approx$ (respectively $\mathsf B_\omega$); on them it is unitary, so it lies in $\mathfrak V$. None of the quantities in use depends on the field on these bonds: $G$ depends on the field within $r_G$, $a_0$ within $R_3$ of the blocks that meet $X^\approx$, $V_\square$ within $R_v$ of $\square$, and $S_k(Y_t)$ within $\lceil R_2/M\rceil$ cells of $Y_t$, and all these regions lie inside $Y^\approx$ (respectively $\mathsf B_\omega$).

Together with $V'$, also $V'^u$ lies in $\mathcal U_Y(x')$, because the operator norm is invariant under multiplication by unitary matrices on the left and on the right, so that \textup{(u1)}, \textup{(u2)} and \textup{(u3)} hold for $V'^u$ as they do for $V'$. By \textup{(4a)}, $\mathsf D_k(Y,V')=-\mathsf D^q(Y,V')-[\mathsf D^u(Y,V')-\mathsf D^u(Y,\mathbf 1)]$, with the same formulas at complex fields.

Here $\mathsf D^q(Y,V'^u)=\mathsf D^q(Y,V')$ at every field of its domain, by \Cref{Y-C-CIITh73}\,(b): in its proof, the matrices of \S\ref{Y-C-1} whose traces make up $\mathsf D^q$ go over into their conjugates by $\operatorname{Ad}_u$, and as this is an identity of matrices, the traces agree. Likewise $\mathsf D^u(Y,V'^u)=\mathsf D^u(Y,V')$ at every field of $\mathfrak V$, by \Cref{Y-C-CIITh74}\,(b).

An element of $\mathcal U_Y(x')$, extended by $\mathbf 1$, lies in the interior of $\mathfrak V$, by \textup{(u1)}. So does its transform by the extended $u$: on the bonds of $Y^\approx$ its norms are the numbers $|V'(c)|$, on the bonds that leave $Y^\approx$ it is unitary, and on all other bonds it equals $\mathbf 1$. This transform differs from the extension of $V'^u$ only on the bonds that leave $Y^\approx$, whereas $\mathsf D_k(Y,\cdot)$ depends on the bonds of $Y^{+\lceil R_2/M\rceil}$ alone, by \textup{(4c)}, and $\lceil R_2/M\rceil\le r_1-8<r'$. Therefore $\tilde{\mathsf D}(Y,V'^u)=\tilde{\mathsf D}(Y,V')$ on $\mathcal U_Y(x')$.

The function $P$ is invariant by \textup{(5c)}, and $\tilde{\mathsf D}$ and $P$ are holomorphic on $\mathcal U_Y(x')$ by \textup{(4a)} and \textup{(5c)}. Thus $E'^{\rm new}$ has the invariance that \textup{(ca)} demands at $x'$.

Both $E'^{\rm new}(Y,\cdot)$ and $E'^{\rm old}[u](Y,\cdot)$ vanish at pure gauges: the cluster terms by the invariance and the subtraction, and $\tilde{\mathsf D}$ and $P$ by \textup{(4a)} and \textup{(5c)}. No letter of a new cluster depends on $u$, that is, on any of $E^a$, $E^h$ and $w$, and neither $\tilde{\mathsf D}$ nor $P$ does; so $E'^{\rm new}$ does not depend on $u$.

\begin{enumerate}\item[\textup{(8.1)}] By \textup{(R)}, the new clusters contribute at most \[49.98\,e^{4x^{-1/8}}[2K_v^\flat p_0^3x^{-1/2}+x^{-20}]\le102K_v^\flat p_0^3x^{-1/2}+x^{-19}.\]

From $r_1\ge1+3+5=9$ we get $\varrho_*\ge2r_1+3\ge21$, and $s_0\ge2$; so \textup{(Y.2)} gives $x^{-1/8}\le(50\cdot47^3)^{-1}<2\cdot10^{-7}$, whence $e^{4x^{-1/8}}<1.000001$, $2\cdot49.98\cdot1.000001\le102$ and $49.98\cdot1.000001\cdot x^{-20}\le x^{-19}$.

The subtraction doubles this bound. Next, \[\FitDisplay{\dimexpr\linewidth-1pt\relax}{c_a(x')=L^{1/2}(p_0'/p_0)^3c_a(x)\ge0.594L^{1/2}E_ap_0^3x^{-1/2}}\], by \textup{(0b)}. As $E_a\ge1500K_v^\flat L^{-1/2}$ by \textup{(O12)}, the first term is at most $204/(1500\cdot0.594)\le0.2290$ in units of $c_a(x')$; and the third inequality of \textup{(Y.6)} gives \[\FitDisplay{\dimexpr\linewidth-1pt\relax}{2x^{-19}=2x^{-9}\cdot x^{-10}\le2x^{-9}c_a(x')/32\le0.02c_a(x')}\]. So, in units of $c_a(x')$, the clusters give at most $0.25$.

For $\tilde{\mathsf D}$, \[(K_2^\sharp/2)\varepsilon'\cdot\sum_n(26e)^{n-1}e^{-(9\kappa+40)(n-1)}\le1.01(K_2^\sharp/2)\varepsilon'\le(1.01/16)c_a(x'),\] because $c_a(x')=E_ap_0'^2\varepsilon'$ and, by the second inequality of \textup{(Y.6)}, $8K_2^\sharp\le E_ap_0'^2$. For $P$, $K_{{\rm q},k}p_0'^2\varepsilon'\le c_a(x')/32$, as $E_a\ge32K_{{\rm q},k}$ by \textup{(O12)}.

Adding up, $0.25+0.06313+0.03125=0.34438\le0.35$.

\item[\textup{(8.2)}] Consider a cluster $C$ of one letter \textup{(a)}, at a small $X$ whose first cube lies in $\hat\square_0$; let $o$ be the centre site of $\square_0$, and $\mathsf B_\omega$ the box of radius $\omega$ about $o$. Then $\Phi(C)=\mathcal G_X(V'):=E_\xi[\chi_X(e^{E^a(X,U)}-1)]$ and $\operatorname{co}(C)=\{\square_0\}$.

The blocks that meet $X^\approx$ lie within $M+\lceil M\,r(x)/L\rceil$ of $o$, so $\mathsf B_\omega$ contains the set $B$ of Step~2\,(ii). We extend the field on $\mathsf B_\omega$ by $\mathbf 1$. No quantity of this item depends on the field outside $\mathsf B_\omega$. By \textup{(e1)}, $|V'(c)|<1+\vartheta(x')+\mathsf g_x<1+\vartheta_0^\flat$, because $\mathsf g_x=\vartheta_0^\flat\log L/(4(\log x)^2)<\vartheta_0^\flat/\log x'$; so the extended field lies in the interior of $\mathfrak V$, and \textup{(2a)} to \textup{(2f)} apply.

By Step~2\,(ii) at $\tau=17T$, the function $\mathcal G_X$ is holomorphic and $G$-invariant on $U^{17T}(\mathsf B_\omega;x')$, with $|\mathcal G_X|\le2A\,c_ae^{-\kappa(\#X-1)}$. The restriction to $\mathsf B_\omega$ of an element of $\mathcal U_{\{\square_0\}}(x')$ satisfies \textup{(u1)} to \textup{(u3)} at $x'$, since $\omega\le M\,r'$.

\Cref{Y-H-Th57} applies, its hypotheses being $T>1$, $17T\varepsilon'\le1/40$, the embedding of $\mathsf B_\omega$, $3\omega+1\le\ell_w(x')$, and \textup{(C0)} to \textup{(C2)}; it yields \[|\mathcal G_X(V')-\mathcal G_X(\mathbf 1)|\le4T^{-2}A\,c_a(x)e^{-\kappa(\#X-1)}.\]

There are $L^3$ choices of the first cube, and $\sum_{X\ni{\rm cube}}e^{-\kappa(\#X-1)}\le1.2124$. The total is therefore $4.85\cdot361\cdot \mathsf N_{\rm c}^2\cdot A\,c_a(x)\le(2947.6\,\mathsf N_{\rm c}^2/L^{1/2})A\,c_a(x')\le0.2457\,A\,c_a(x')$, by \textup{(L2$^{\rm su}$)}.

\begin{theorem}[Bounded invariant holomorphic functions vary little]\label{Y-H-Th57} Let \textup{(C0)}, \textup{(C1)}, \textup{(C2)} of \textup{(0d)} hold, suppose that $T>1$, $17T\varepsilon'\le1/40$, $2\omega+1<N$ and $3\omega+1\le\ell_w(x')$, let $o$ be a coarse site, and let $\mathsf B_\omega$ be the box of radius $\omega$ about $o$. Let $\Psi$ be a bounded holomorphic function on $\mathcal U:=U^{17T}(\mathsf B_\omega;x')$ that is invariant under $G$-valued gauge transformations. Then every $\operatorname{SL}(2,\mathbb C)$ field $V'$ on the bonds of $\mathsf B_\omega$ that satisfies \textup{(u1)}, \textup{(u2)} and \textup{(u3)} at $x'$ obeys \[|\Psi(V')-\Psi(\mathbf 1)|\le2T^{-2}\sup|\Psi|.\] \end{theorem}

\begin{proof}[Proof of \Cref{Y-H-Th57}] \emph{Step 1: the curve.} We follow Step~(1) of the proof of \Cref{Y-C-Th70}, with $\omega$ in place of $R_3$. The field $V'$ has the properties \textup{(h2)} and \textup{(h3)} of item~C.II~1.4 of \S\ref{Y-C-1} with $\hat\varepsilon=\varepsilon'$. Along the comb paths we define the comb gauge: $\mathsf k(y):=V'(\Gamma(o,y))$ for every site $y$ of $\mathsf B_\omega$, and $\mathsf W(c):=\mathsf k(c_-)V'(c)\mathsf k(c_+)^{-1}$ for every bond $c$ of $\mathsf B_\omega$.

By \Cref{Y-C-L67}\,(a) and by \textup{(C0)}, $|\mathsf k(y)|=|\mathsf k(y)^{-1}|\le2$ and $|\mathsf W(c)-\mathbf 1|\le e^{18\omega\varepsilon'}-1\le19\omega\varepsilon'\le0.106$. We put $A_c:=-i\log\mathsf W(c)$; this matrix is traceless, and $|A_c|\le\mathsf a:=21\omega\varepsilon'$.

For complex $z$ let $V'_z(c):=\mathsf k(c_-)^{-1}e^{izA_c}\mathsf k(c_+)$. This field is entire in $z$ and takes its values in $\operatorname{SL}(2,\mathbb C)$, and $V'_1=V'$. Moreover $V'_z(c)=V'(c)(\mathbf 1+\delta_c)$, where $\mathbf 1+\delta_c=\mathsf k(c_+)^{-1}e^{i(z-1)A_c}\mathsf k(c_+)$.

Now let $|z|<T$. From \textup{(C1)}, together with $\ell_w(x')\ge18$ and $\mathsf g_x\le1$, we obtain $(T+1)\mathsf a\le10^{-3}$. Consequently $|\delta_c|$ and $|(\mathbf 1+\delta_c)^{-1}-\mathbf 1|$ are both $\le4(e^{(T+1)\mathsf a}-1)\le\delta:=85(T+1)\omega\varepsilon'$. We show that $V'_z$ satisfies the three conditions that define $\mathcal U$.

\begin{enumerate}

\item[\textup{(e1)}] For every bond $c$ of $\mathsf B_\omega$, $|V'_z(c)|<(1+\vartheta(x'))(1+\delta)\le1+\vartheta(x')+1.01\delta\le1+\vartheta(x')+\mathsf g_x$, the last inequality by \textup{(C1)}.

\item[\textup{(e2)}] Let $\pi$ be a path in $\mathsf B_\omega$ of $n\le\ell_w(x')$ bonds. In the product $V'_z(\pi)$ we move every factor $(\mathbf 1+\delta_c)^{\pm1}$ to the left, past the partial products of $V'$; by \textup{(u2)}, these partial products have norm $<2$. In this way $|V'_z(\pi)|\le(1+4\delta)^n|V'(\pi)|<e^{4n\delta}\mathsf b(x')\le\mathsf b(x')+\mathsf g_x$, since \textup{(C1)} gives $4n\delta\le(340/2640)\mathsf g_x$.

\item[\textup{(e3)}] Let $P$ be a plaquette of $\mathsf B_\omega$ with base site $y_P$, and let $A_1,\dots,A_4$ be the matrices $A_c$ of its four bonds. Then $\partial V'_z(P)=\mathsf k(y_P)^{-1}\Pi_z\mathsf k(y_P)$ with $\Pi_z:=e^{izA_1}e^{izA_2}e^{-izA_3}e^{-izA_4}$. As $4T\mathsf a=84T\omega\varepsilon'\le1$ by \textup{(C0)}, the estimate \textup{(E)} of item~C.I~1.1 of \S\ref{Y-C-1} yields $|\Pi_z-\mathbf 1-iz(\mathrm{d}'A)_P|\le(4|z|\mathsf a)^2$, where $(\mathrm{d}'A)_P$ is the coarse coboundary. At $z=1$ we have $|\Pi_1-\mathbf 1|=|\mathsf k(y_P)(\partial V'(P)-\mathbf 1)\mathsf k(y_P)^{-1}|<8\varepsilon'$, so that $|(\mathrm{d}'A)_P|<8\varepsilon'+16\mathsf a^2$. Therefore $|\partial V'_z(P)-\mathbf 1|\le4|\Pi_z-\mathbf 1|<32T\varepsilon'+128T^2\mathsf a^2=(32+56448T\omega^2\varepsilon')T\varepsilon'<34T\varepsilon'$, by \textup{(C2)}.

\end{enumerate}

Thus $V'_z\in\mathcal U$ for $|z|<T$. Hence $f(z):=\Psi(V'_z)$ is a holomorphic function on $|z|<T$, with $|f|\le\mathsf s:=\sup|\Psi|$ and $f(1)=\Psi(V')$.

\emph{Step 2: complex gauges.} We follow Step~(2) of the proof of \Cref{Y-C-L69}. The set $\mathcal U$ is open, as it is given by finitely many strict inequalities. For $\mathsf Z\in\mathcal U$ and $\lambda:\mathsf B_\omega\to\mathfrak g_{\mathbb C}$ we define $\mathfrak d(\mathsf Z;\lambda):=(d/d\sigma)\Psi(\mathsf Z^{\exp(i\sigma\lambda)})$ at $\sigma=0$. Because $\Psi$ is holomorphic, $\mathfrak d$ is $\mathbb C$-linear in $\lambda$; because $\Psi$ is invariant, $\mathfrak d$ vanishes for $\mathfrak g$-valued $\lambda$. An arbitrary $\lambda$ has the form $\lambda_1+i\lambda_2$ with $\mathfrak g$-valued $\lambda_1$ and $\lambda_2$, and so $\mathfrak d=0$ for every $\lambda$.

It follows that $\Psi(\mathsf Z^{u_s})$ is constant along every $C^1$ curve $s\mapsto u_s$ of $\operatorname{SL}(2,\mathbb C)$-valued gauges for which $\mathsf Z^{u_s}\in\mathcal U$, since its derivative in $s$ is $\mathfrak d(\mathsf Z^{u_s};-i(du_s/ds)u_s^{-1})$, which vanishes.

We write $\mathsf k(y)^{-1}$ in polar form, $\mathsf k(y)^{-1}=k_0(y)e^{h(y)}$, where $k_0$ is $G$-valued and $h$ is hermitian and traceless, and we set $u_s:=k_0e^{sh}$ for $0\le s\le1$. Then $V'_0=\mathbf 1^{u_1}$.

Let $\pi$ be a path from $y$ to $y''$. Cordes' inequality, which is proved in Step~(1) of the proof of \Cref{Y-C-L69}, gives $|\mathbf 1^{u_s}(\pi)|=|e^{sh(y)}e^{-sh(y'')}|\le|e^{h(y)}e^{-h(y'')}|^s=|V'_0(\pi)|^s\le|V'_0(\pi)|$. By Step~1 of the present proof, $V'_0\in\mathcal U$. Hence $\mathbf 1^{u_s}$ satisfies \textup{(e1)} and \textup{(e2)}, and its curvature is $\mathbf 1$; so $\mathbf 1^{u_s}\in\mathcal U$ for every $s$. The constancy just proved, applied with $\mathsf Z=\mathbf 1$, and the invariance of $\Psi$ under the $G$-valued gauge $u_0=k_0$ give $f(0)=\Psi(\mathbf 1^{u_1})=\Psi(\mathbf 1^{u_0})=\Psi(\mathbf 1)$.

\emph{Step 3: the derivative of $f$ at $0$.} We show that $f'(0)=0$. Let $\mathsf Z_z:=(e^{izA_c})_c$, so that $V'_z=\mathsf Z_z^{u_1}$. The set $\{\mathbf 1^{u_s}\}$ is compact and lies in the open set $\mathcal U$; hence there is $z_0>0$ such that $\mathsf Z_z^{u_s}\in\mathcal U$ for $|z|<z_0$ and for all $s$. For these $z$, Step~2 of the present proof and the invariance of $\Psi$ give $f(z)=\Psi(\mathsf Z_z^{u_0})=\Psi(\mathsf Z_z)$.

Let $\Theta(A):=(d/dz)\Psi(\mathsf Z_z)$ at $z=0$; it is $\mathbb C$-linear in $A$. For a constant $v\in G$ we have $\mathsf Z_z^v=(e^{iz\operatorname{Ad}_vA_c})_c$, and therefore $\Theta(\operatorname{Ad}_vA)=\Theta(A)$. For a traceless matrix $Y$, the mean of $\operatorname{Ad}_vY$ over $G$ commutes with $G$ and is traceless, so it is $0$. Averaging over $v$ and using linearity, we find $\Theta(A)=\Theta(\int\operatorname{Ad}_vA\,dv)=0$; thus $f'(0)=0$.

\emph{Step 4: the bound.} Since $f'(0)=0$, the function $(f(z)-f(0))/z^2$ is holomorphic on $|z|<T$, and $|f-f(0)|\le2\mathsf s$. The maximum principle on the circles $|z|=T_1$, with $T_1\to T$, gives at the point $z=1$ the bound $|f(1)-f(0)|\le2\mathsf sT^{-2}$; as $f(1)=\Psi(V')$ and $f(0)=\Psi(\mathbf 1)$, this is the claim of the theorem.

The theorem and its proof are independent of the rows $\ell$. \end{proof}

\item[\textup{(8.3)}] Here the cluster consists of one letter of kind \textup{(a)} at a set $X$ that is not small, so $n=\#X>L/2-2$; its core is $\operatorname{co}(C)=Y_0(X)\ni\square$, and $\#Y_0\le8+432n/L$. For such a cluster, $|\mathcal G_X(V')-\mathcal G_X(\mathbf 1)|\le4A\,c_a(x)e^{-\kappa(n-1)}$. The cell $\square$ meets $X^{+1}$, so $X$ contains one of $(L+2)^3$ cubes, and, for each $n$, at most $(26e)^{n-1}$ sets $X$ contain a given cube.

Summing over these clusters, we obtain at most $4A\,c_a(x)(L+2)^3\sum_{n>L/2-2}\exp\{(4.2581-\kappa)(n-1)+\kappa(7+432n/L)\}$. For $\kappa\ge6$ we have $\kappa-4.2581\ge0.29\kappa$.
Since $L>10^{13}$, the exponent is at most
$-0.289\kappa(n-1)+7.1\kappa$. Summing this geometric majorant gives
\[
 4A\,c_a(x)(L+2)^3
 \frac{\exp\{-\kappa[0.289(L/2-3)-7.1]\}}{1-e^{-0.289\kappa}}.
\]
Here $0.289(L/2-3)-7.1\ge(0.8/6)L$,
$4/(1-e^{-0.289\kappa})<5$, and $5(L+2)^3<e^{0.3L}$.
The contribution is thus at most $e^{-L/2}A\,c_a(x)\le10^{-3}A\,c_a(x')$.

\item[\textup{(8.4)}] Now $n(C)\ge2$, and we root $C$ at a letter of kind \textup{(a)}. The sizes are computed with the number $A$ for which $u\in M_A^{\rm su}(x)$, and condition \textup{(S)} holds for them, by Step~7\,(ii). Claim~(2) of \Cref{Y-G-L18} then yields the bound $t\,e_\alpha e^{2\varepsilon_2\#I_\alpha+d_\alpha}$, in which a single size $e_\alpha$ occurs; it is linear in $A$.

Two applications of \textup{(R)} give $2\cdot49.98\cdot4.85\cdot e^{0.08}\cdot x^{-1/4}L^3A\,c_a(x)\le525.2\,L^3x^{-1/4}A\,c_a(x)$, where the letters at sets $X$ that are not small are bounded by the count of case \textup{(8.3)}. This is $\le500e^{\kappa+3}L^3x^{-1/4}A\,c_a(x)\le842e^{\kappa+3}L^{5/2}x^{-1/4}A\,c_a(x')\le10^{-3}A\,c_a(x')$, the last step by \textup{(Y.6)}.

\end{enumerate}

The three cases \textup{(8.2)}, \textup{(8.3)} and \textup{(8.4)} cover all old clusters: those that consist of one letter of kind \textup{(a)} at a small set $X$, those that consist of one letter of kind \textup{(a)} at a set $X$ that is not small, and those with $n(C)\ge2$. Adding the bounds of these cases, we obtain the second bound of Step~8: in units of $A\,c_a(x')$, the sum $\sum_{Y\ni\square}e^{\kappa(\#Y-1)}\sup_{\mathcal U_Y(x')}|E'^{\rm old}[u](Y,\cdot)|$ is at most $0.2457+0.001+0.001\le0.248$. \end{proof}

\emph{Step 9: $E'^h$; the third bound of \Cref{Y-H-C9}\,(b), and the clauses of \textup{(ch)} at $x'$ without a size.} For every set $Y$ of cells, $E'^h(Y,\cdot)$ is a complex, Borel and gauge invariant function of the $G$-valued fields on the bonds of $Y^\approx$ that are tame there, and it is independent of $e$. Moreover, $\sum_{Y\ni\square}e^{\kappa(\#Y-1)}\sup|E'^h(Y,\cdot)|\le c_h(x')$ for every $A\le16$ with $u\in M_A^{\rm su}(x)$.

\begin{proof}[Proof of Step~9] The function $E'^h(Y,\cdot)$ is well defined, and its series converges absolutely, by \Cref{Y-G-L18} for the alphabet $\{\alpha\in\mathfrak A:J_\alpha\text{ inside }Y\}$, together with Step~6\,(ii), (iv) and Step~7\,(ii). It is Borel, being the sum of a uniformly convergent series of integrals of Borel functions; that the letters are jointly Borel is shown in the verification of hypothesis~(i) of \Cref{Y-H-L52} in Step~10\,(9).

The function is gauge invariant: by the argument of Step~8, $b_\square$, $A_\square$, $w(X,U)$, $E^a$, $E^h$, $1_{X^\approx}(U)$ and $1_{\rm ind}$ are invariant under $(\operatorname{Ad}_u\xi,V'^u)$. For $A_\square$ this rests on \textup{(6e)}, since in a hard cluster $A_\square$ appears also at failed tests, and there it is not $e^{V_\square}$.

For the dependence on the field we use Step~3\,(V) and the extension by $\mathbf 1$ in Step~6: $\Phi(C)(\cdot)$ is evaluated at the extension of the field by $\mathbf 1$, and only the values of the field in $Y^{+(r'-2)}$ enter. On $Y^\approx$ every angle satisfies $\phi'_P\le\varepsilon'<\pi$, so no $e$ enters.

We now bound the size. A hard cluster contains a letter of kind \textup{(h)}, a letter of kind \textup{(b)} or a zone letter, and in each case we use \textup{(R)} of Step~8.

\begin{enumerate}

\item[\textup{(h)}] For a letter of kind \textup{(h)}, the bound is $49.98\cdot L^3\cdot4.85\cdot e^{0.08}\cdot A\,c_h(x)\le262.6\,L^3A\,c_h(x)$; the sets $X$ that are not small add a negligible amount, by the count of case \textup{(8.3)} of Step~8 with $c_h(x)$ in place of $c_a(x)$. Since $A\le16$, this is $\le10^3L^3\cdot16c_h(x)\le c_h(x')/2$, by the first inequality of \textup{(Y.5)}, because $c_h(x)/c_h(x')=e^{-c_h(p_0^2-p_0'^2)}$.

\item[\textup{(b)}] For a letter of kind \textup{(b)}, the bound is $49.98\,e^{4x^{-1/8}}u_b$.

\item[\textup{Zone}] For a zone letter $\zeta$, $e_\zeta e^{2\varepsilon_2\#I_\zeta+(\kappa+1)\#J_\zeta}\le e^\circ_\zeta e^{(\kappa+2)\#J_\zeta}\le u_1\prod_o\hat w(o)$, where $o$ runs over the objects of $\zeta$. Summing over the zone letters rooted at $\square_0$ by means of the bound for animals, we obtain at most $u_1\cdot(1/2)x^{-1/4}$.

\end{enumerate}

The letters of kind \textup{(b)} and the zone letters together give $\le50(2u_b+u_1)\le(1/2)e^{-\delta_3p_0^2/2}\le c_h(x')/2$, by the second inequality of \textup{(Y.5)}, by $c_h\le\delta_3/2$ and by $p_0'\le p_0$; with the bound $c_h(x')/2$ for kind \textup{(h)}, this gives the bound of Step~9. \end{proof}

\emph{Step 10: $\Xi'$ and $w'$; clause \textup{(a)} of \Cref{Y-H-C9}, the strong cap of the output in clause \textup{(b)}, and the mean over $e$.} The following hold.

\begin{enumerate}

\item[\textup{(8a)}] $e^{\sum n_j}Z_u=\sum_{\mathfrak G}\prod_{\gamma\in\mathfrak G}K(\gamma)\cdot\exp\{\sum_{C\not\sim\mathfrak G}\Phi(C)(V')\}$, where $\mathfrak G$ runs over the sets of rough polymers that are pairwise not related by $\sim$ and that contain every forced class exactly once.

\item[\textup{(8b)}] The formula of \Cref{Y-H-C9}\,(a) holds at $(V',e)$, with $\Xi'=\sum_{(\mathfrak G,\mathfrak N,\mathfrak N^\flat)}\prod K(\gamma)\prod(e^{-\Phi^\natural(C)}-1)\prod(e^{\tilde{\mathsf D}(Y)}-1)$.

\item[\textup{(8c)}] $\Xi'=\sum_{\{\mathsf Y_i\}}\prod_iw'(\mathsf Y_i,V',e)$, the sum running over the families of \textup{(c$\Xi$)} at $x'$. Here $w'(\mathsf Y,\cdot,\cdot)$ depends on $V'$ only through the bonds of $\mathsf Y^{+s'}$, and on $e_P$ only at the plaquettes of cells of $\mathsf Y$.

\item[\textup{(8d)}] At every pair $(V',e)$, $|w'(\mathsf Y,V',e)|e^{\kappa(|\mathsf Y|_c-1)}\le\exp\{-c_\Xi x'\varepsilon'\sum_{D'\cap\mathsf Y}\phi'_P\}$.

\item[\textup{(9)}] The mean $w'(\mathsf Y,V')=\int w'(\mathsf Y,V',e)\,de$ is a complex, Borel and gauge invariant function, which depends on $V'$ only through the bonds of $\mathsf Y^{+s'}$ and obeys the same bound; and $\Xi'=\sum_{\{\mathsf Y_i\}}\prod w'(\mathsf Y_i,V')$ at every $V'$.

\end{enumerate}

\begin{proof}[Proof of Step~10] We begin with the normaliser $n_j$, which we treat by \textup{(4a)}. A set $Y$ is on exactly when $1_{Y^\approx}(V')=1$. A short $Y$ is connected and lies in $f$, and the parts $f_j$ of $f$ are farther than $4r_1+2$ apart; hence it lies in exactly one $f_j$. An off $Y$ inside $f_j$ has a cell within $r'$ of a mark, and this mark lies within $r'+s_0$ of $\mathrm{mk}'_j$, so it belongs to class $j$. Consequently $|n_j|\le(2r'+1)^3\#\mathrm{mk}'_j\cdot1.01\cdot K_1^\sharp/2$.

The function $n_j$ is real, Borel and gauge invariant. It is Borel because on each of the finitely many Borel sets $\{\mathrm{mk}'(V')=\mathsf d\}$ it is a finite sum of functions $\tilde{\mathsf D}(Y,\cdot)$, and these are continuous on $G$-valued fields, the rough ones included, by \textup{(4a)}.

A long $Y$ has a cell within $r'$ of a mark and a cell farther than $2r'$ from it, so $\#Y>r'$. By \textup{(Y.9)} we then have $|\tilde{\mathsf D}(Y)|\le1/2$, and therefore $|e^{\tilde{\mathsf D}(Y)}-1|e^{(\kappa+1)(\#Y+s')}\le\hat w(Y):=K_1^\sharp e^{-\kappa^\sharp(\#Y-1)}e^{(\kappa+1)(\#Y+s')}$.

The exponent per cell is $4.2581-\kappa^\sharp+\kappa+1\le-(9\kappa+34.74)$, and $s'=8r'$. Hence $\sum_{Y\text{ long, within }r'\text{ of }\mathrm{mk}'_j}\hat w(Y)\le(2r'+1)^3\#\mathrm{mk}'_j\cdot1.001\cdot K_1^\sharp e^{\kappa+1}e^{-(\kappa+26.74)r'}\le\#\mathrm{mk}'_j$, the last step by \textup{(Y.9)}.

Each set $Y$ is on, short or long, and exactly one of these holds; this partition gives $\sum_Y\tilde{\mathsf D}=\sum_Y1_{Y^\approx}\tilde{\mathsf D}+\sum_jn_j+\sum_{\rm long}\tilde{\mathsf D}$.

\begin{enumerate}

\item[\textup{(8a)}] We split each set $\mathfrak B$ of Step~3\,(I) into its $\sim$-components and let $\mathfrak G$ consist of the components that contain a forced letter; by \textup{(Act1)}, they are rough polymers. The union of the sets $I_\alpha$, where $\alpha$ runs over the letters of the polymers of $\mathfrak G$, contains $f^{+2r_1}$. The remaining components are $\sim$-connected compatible sets of letters of $\mathfrak A_{\mathfrak G}:=\{\alpha\in\mathfrak A:\alpha\not\sim\text{ any letter of }\mathfrak G\}$.

All letters of $\mathfrak A_{\mathfrak G}$ are actual. Indeed, ${\rm dist}(I_\alpha,f^{+2r_1})>r_1$ for each of them; for a zone letter, ${\rm dist}(\mathfrak b,f)>5r_1\ge4r_1+2$, and for a letter of kind \textup{(a)} or \textup{(h)}, $X^\approx$ lies off the fine set under $f^{+r_1}$. The independence of Step~3\,(IV) and \Cref{Y-G-L18} for $\mathfrak A_{\mathfrak G}$, with complex weights and with Step~6\,(i), now give \textup{(8a)}; condition \textup{(S)} for the smaller alphabet is weaker.

\item[\textup{(8b)}] Every cluster $C\not\sim\mathfrak G$ is on. Indeed, its letters satisfy ${\rm dist}(J_\alpha,\mathrm{mk}')>s_0+3r_1$, and each cell of the joining path lies within $r'$ of some $J_\alpha$, so ${\rm dist}(\operatorname{co}(C),\mathrm{mk}')>r'$.

We sort the clusters by type. For an analytic cluster $C$ we write $\Phi(V')=\Phi(\mathbf 1)+(\Phi(V')-\Phi(\mathbf 1))$ if $C\not\sim\mathfrak G$, $0=\Phi(\mathbf 1)+(\Phi(V')-\Phi(\mathbf 1))-\Phi(V')$ if $C\sim\mathfrak G$ is on, and $0=\Phi(\mathbf 1)-\Phi(\mathbf 1)$ if $C\sim\mathfrak G$ is off. Hard and constant clusters are sorted in the same way, with these changes: a hard or constant $C\not\sim\mathfrak G$ keeps its value $\Phi(C)(V')$; for a hard $C\sim\mathfrak G$ that is on, and for a constant $C\sim\mathfrak G$, we write $0=\Phi(C)(V')-\Phi(C)(V')$; and a hard $C\sim\mathfrak G$ that is off contributes $0$, as $\Phi^\natural(C)=0$ there.

Summing over all clusters, we obtain \[\sum_{C\not\sim\mathfrak G}\Phi(C)(V')=[\sum_{\rm analytic}\Phi(C)(\mathbf 1)+\sum_{\rm constant}\Phi(C)]+\sum_Y1_{Y^\approx}(V')[E'^a-\tilde{\mathsf D}+P+E'^h](Y,V')-\sum_{C\sim\mathfrak G}\Phi^\natural(C),\] where all sums converge absolutely and $|\Phi^\natural|\le\hat\Phi$. We insert this identity into the formula of Step~1, together with the partition of the sets $Y$ into on, short and long ones, and we expand $e^{-\sum\Phi^\natural}$ and $e^{\sum_{\rm long}\tilde{\mathsf D}}$; this gives \textup{(8b)}.

\item[\textup{(8c)}] The graph of \textup{(Fm4)} for a triple $(\mathfrak G,\mathfrak N,\mathfrak N^\flat)$ of \textup{(8b)} has the letters of $\mathfrak G$, the clusters of $\mathfrak N$ and the sets of $\mathfrak N^\flat$ as its vertices, and it joins two of them whenever their cores lie within $s'$ of each other. Its components give sets $\mathsf Y_i$ at pairwise distance $>s'$, each with a triple as in \textup{(Fm4)}. Indeed, two letters of one $\gamma$ that are related by $\sim$, and also a cluster $C$ and a letter of $\mathfrak G$ to which $C$ is related by $\sim$, have cores within $2(s_0+\varrho_*)+r_1\le5r'\le s'$ of each other; and a long $Y$ has a cell within $r'$ of a mark, which lies in the core of a forced letter.

Hence each $\mathsf Y_i$ contains a mark, $D'$ lies in the union of the $\mathsf Y_i$, and no $\mathsf Y_i$ has a mark in $\mathsf Y_i^{+s'}\setminus\mathsf Y_i$; in particular, none has a mark in $\mathsf Y_i^{+(s'-1)}\setminus\mathsf Y_i$.

Conversely, two letters, one of the triple of $\mathsf Y_i$ and the other of the triple of $\mathsf Y_{i'}$ with $i\ne i'$, satisfy ${\rm dist}(I_\alpha,I_\beta)>s'-2(s_0+\varrho_*)\ge r_1$; so they are not related by $\sim$, and hence they are compatible. A mark that makes a set $Y$ of $\mathfrak N_i^\flat$ off lies in $\mathsf Y_i^{+r'}$, and as $r'\le s'-1$, it lies in $\mathsf Y_i$. In the remaining case of the definition of $w'$ in \textup{(Fm4)}, when $\mathrm{mk}'(V')$ has no cell in $\mathsf Y^{+(s'-1)}\setminus\mathsf Y$, the marks outside $\mathsf Y$ are farther than $s'-1\ge5r'\ge2s_0+4r_1+2$ from $\mathsf Y$, and so $\mathrm{mk}'\cap\mathsf Y$ is a union of classes.

For the dependence on the field we use Step~3\,(V), which is valid at every pair, since by \textup{($\Gamma$3)} each $a_{0,b}$ depends at every pair only on $V'$ and $X'$ within $R_3$ of $y(b)$; here $I_\zeta^{+r_1}$ lies inside $J_\zeta^{+s_0+\varrho_*+r_1}$. Whether a letter is actual, the classes, and whether a set is on, off, short or long depend only on the marks within $2s_0+4r_1+2$ of the core.

Measured in cells from $\mathsf Y$, the reach of each ingredient is as follows: for a forced letter it is $s_0+\varrho_*+r_1\le3r'$, for $w_j$ it is $2r'$, for $n_j$ it is $3r'$, for a long $Y$ it is $\lceil R_2/M\rceil\le r_1$, for $\Phi^\natural(C)$ it is $r'$, for the marks that decide whether a letter is actual, what the classes are, and whether a set is on, off, short or long, it is $2s_0+4r_1+2\le5r'$, and for the first clause of the definition of $w'$ in \textup{(Fm4)}, which sets $w'(\mathsf Y,V',e):=0$ if $\mathrm{mk}'(V')$ has a cell in $\mathsf Y^{+(s'-1)}\setminus\mathsf Y$, it is $s'-1$. All of these are $\le s'-1$.

Whether a cell is a mark depends on the plaquettes based in it, whose sites lie in the cell or in a neighbour of it; so only bonds with both ends in $\mathsf Y^{+s'}$ enter. The choice $e_P$ enters only at $P\in D'$, hence at plaquettes of marks within $s'-1$ of $\mathsf Y$; where $w'\ne0$, these marks are cells of $\mathsf Y$. The first clause of the definition of $w'$ in \textup{(Fm4)} is independent of $e$.

\item[\textup{(8d)}] By Step~6\,(i), $|K(\gamma)|\le\prod e_\alpha$. As $|\mathsf Y|_c\le\sum_{\rm elements}(\#\mathrm{core}+s')-s'$, we have $e^{\kappa(|\mathsf Y|_c-1)}\le\prod_{\rm elements}e^d$. Moreover $|e^{-\Phi^\natural}-1|\le2\hat\Phi$: indeed, $\hat\Phi\le t\varepsilon_2\le1$ by \textup{(S)}, and $e^s-1\le(e-1)s\le2s$ on $[0,1]$; besides, $t\varepsilon_2=x^{-3/8}\le\log2$.

We apply \Cref{Y-G-L19}, with $t\le1/2$, taking for $\mathsf F$ the forced letters of $\mathfrak G$, for $\mathfrak B_o$ its other letters, and for $\mathfrak M$ the clusters, with $\hat w_C=2\hat\Phi e^{d(C)}$; the condition of the lemma holds, with $2t\varepsilon_2\le\varepsilon_2$. The sets $\mathfrak N^\flat$ contribute $\prod(1+\hat w(Y))\le\exp\sum\hat w(Y)\le\exp\sum_{\mathsf F}A_\zeta$, where $A_\zeta:=\sum_{\mathsf J(\zeta)}\#\mathrm{mk}'_j$.

Therefore $|w'|e^{\kappa(|\mathsf Y|_c-1)}\le\sum_{\mathsf F}\prod_{\zeta\in\mathsf F}e_\zeta e^{d_\zeta+2\varepsilon_2\#I_\zeta+A_\zeta}$, where $\mathsf F$ runs over the sets of forced letters whose classes are the classes with marks in $\mathsf Y$, each taken once. Here the numbers $e_\zeta$ are used for actual letters and for the others alike, which gives a majorant.

By \textup{(Y.2)}, $e^{4x^{-1/8}\#I_\zeta}\le e^{\#J_\zeta}$, and with $u_1\le1$ we obtain $e_\zeta e^{d_\zeta+2\varepsilon_2\#I_\zeta+A_\zeta}\le e^{(\kappa+1)s'}\prod_{o\in\zeta}\hat w(o)$, with the weights $\hat w(\square),\hat w_p,\hat w_{p'}$ of Step~7 and \[\hat w(f_j):=\upsilon_f^{-\#D'_j/4}\exp\{-(r_\Xi+c_\Xi)x'\varepsilon'\sum_{D'_j}\phi'_P+[1.01K_1^\sharp(2r'+1)^3+1]\#D'_j+(\kappa+4)((1+s')\#\mathrm{mk}'_j+s')\}.\]

Since $x'\varepsilon'\phi'_P>p_0'^2$ on $D'$ and $\#\mathrm{mk}'_j\le\#D'_j$, it follows that $\hat w(f_j)e^{2\#\mathrm{mk}'_j}\le e^{-c_\Xi x'\varepsilon'\sum\phi'}\upsilon_f^{(3/4)\#D'_j}$.

We take out the factor $e^{-c_\Xi(\cdot)}\upsilon_f^{\#D'_j/4}$. With $\hat w''(f_j)=\upsilon_f^{\#D'_j/2}e^{-2\#\mathrm{mk}'_j}$, the condition for animals holds: at most one class has $f_j\ni\square$, and $\hat w''e^{a_{f_j}}\le\upsilon_f^{1/2}\le\varepsilon''/4$, while the cells and the polymers contribute $\le\varepsilon''/2$. Hence $\sum_{\zeta\ni f_j}e^{(\kappa+1)s'}\prod\hat w''\le e^{(\kappa+1)s'}\upsilon_f^{1/2}\le1$, by \textup{(Y.4)}.

The set $\mathsf F$ is fixed by choosing, for each class, either nothing or a letter that has this class as its first class; hence $\sum_{\mathsf F}\prod\le\prod_j2$.

As a result, the left side of \textup{(8d)} is $\le e^{-c_\Xi x'\varepsilon'\sum\phi'}\prod_j2\upsilon_f^{\#D'_j/4}\le e^{-c_\Xi x'\varepsilon'\sum\phi'}$. The strong cap of the output thus comes from the factor $\exp\{-(r_\Xi+c_\Xi)x'\varepsilon'\sum_{D'_j}\phi'_P\}$ of $\hat w(f_j)$, and not from the input $w$.

\item[\textup{(9)}] Here $\Xi'(V'):=(\rho[u;\ell])'(V')e^{-\mathsf V'}\exp\{x'\mathrm{En}'-\sum1_{Y^\approx}[E'^a+E'^h]\}$, which is independent of $e$. \Cref{Y-H-L52} below, applied to the complex weights $w'$, gives \textup{(9)}.

\begin{lemma}[Averaging the choice $e$ out of a polymer sum]\label{Y-H-L52} Let $\Xi$ be a complex function of $V'$, and for every finite set $\mathsf Y$ of cells let $w(\mathsf Y,\cdot,\cdot)$ be a complex function of $(V',e)$, such that the following hold.

\begin{enumerate}

\item[\textup{(i)}] The function $w$ is jointly Borel, and $|w(\mathsf Y,V',e)|\le B(\mathsf Y,V')<\infty$ with $B$ independent of $e$.

\item[\textup{(ii)}] $w(\mathsf Y,V'^u,\operatorname{Ad}_ue)=w(\mathsf Y,V',e)$ for all $G$-valued $u$.

\item[\textup{(iii)}] $w(\mathsf Y,V',e)$ depends on $V'$ only through the bonds of $\mathsf Y^{+s'}$, and on $e_P$ only at the plaquettes that belong to cells of $\mathsf Y$.

\item[\textup{(iv)}] At every $(V',e)$, $\Xi(V')=\sum_{\{\mathsf Y_i\}}\prod_iw(\mathsf Y_i,V',e)$, the sum running over the families of \textup{(c$\Xi$)} at $x'$.

\end{enumerate}

Then $\bar w(\mathsf Y,V'):=\int_{\mathfrak E}w(\mathsf Y,V',e)\,de$ is Borel and gauge invariant, it is real if $w$ is real, it depends on $V'$ only through the bonds of $\mathsf Y^{+s'}$, it satisfies $|\bar w|\le B$, and $\Xi(V')=\sum_{\{\mathsf Y_i\}}\prod_i\bar w(\mathsf Y_i,V')$ at every $V'$. \end{lemma}

\begin{proof}[Proof of \Cref{Y-H-L52}] The positive and negative parts of $\operatorname{Re}w$ and $\operatorname{Im}w$ are jointly Borel and $\ge0$; by Tonelli's theorem, their means over $e$ are Borel, and by \textup{(i)} they are finite. The dependence on $V'$ and the bound carry over to the mean.

For $G$-valued $u$ we have $\bar w(\mathsf Y,V'^u)=\int w(\mathsf Y,V'^u,\operatorname{Ad}_ue)\,de=\bar w(\mathsf Y,V')$; the first equality holds because $e\mapsto\operatorname{Ad}_ue$ is a product of rotations of spheres and preserves $de$, and the second is \textup{(ii)}.

Fix $V'$. The index set in \textup{(iv)} depends on $D'(V')$ and not on $e$, and it is finite; the left side of \textup{(iv)} is independent of $e$, and we integrate \textup{(iv)} over $e$. Within one family the sets $\mathsf Y_i$ are pairwise disjoint, so by \textup{(iii)} the factors depend on disjoint sets of independent $e_P$, and the mean of the product is the product of the means. \end{proof}

We now show that $\Xi'$ and $w'$ satisfy the hypotheses of \Cref{Y-H-L52}.

\begin{enumerate}

\item[\textup{(i)}] Every letter is jointly Borel in $(\xi,V',e)$. For $A_\square$ this is \textup{(6e)}; for $T_\square$, and hence for $w_j$, it is \textup{(5b)}. The variable $X'$, the marks and the classes are Borel in $(V',e)$. The entries of $G$, $C$ and $S_k(Y)$ are continuous at $G$-valued fields, which lie in the interior of $\mathfrak V$, by \textup{(2d)} and \textup{(4b)}; and $\mathsf g^V$ is a polynomial in $\operatorname{Ad}_{V'(c)}$, by \textup{(2g)}. So $\varphi,a_0,\eta,T_Y$ are Borel. The chart is continuous, so $U$ is Borel. The function $E^a$ is continuous on $\mathcal U_X(x)$, and $E^h$ and $w$ are Borel, by \textup{(ca)}, \textup{(ch)} and \textup{(c$\Xi$)}. Finally, $b_\square,\chi_\square,\chi_X,1_{X^\approx}(U),1_{\rm ind}$ are indicators of Borel sets.

Every $K(\gamma)$ is jointly Borel in $(V',e)$, by Tonelli's theorem together with the bound of Step~6\,(i). Further, $e^{\tilde{\mathsf D}(Y,\cdot)}-1$ is continuous by \textup{(4a)}, and $n_j$ is Borel, as shown for the normaliser at the start of the proof of Step~10. On each of the finitely many Borel sets $\{\mathrm{mk}'\cap\mathsf Y^{+(s'-1)}=\mathsf d\}$, the sum that defines $w'$ is an absolutely convergent series, and its real and imaginary parts are Borel; the bound \textup{(8d)} is independent of $e$.

\item[\textup{(ii)}] The marks, the classes, whether a letter is actual, and whether a cluster or a set is on or off depend only on the angles $\phi'_P$. Moreover $X'(V'^u,\operatorname{Ad}_ue)=\operatorname{Ad}_uX'(V',e)$; together with $\xi\mapsto\operatorname{Ad}_u\xi$, which preserves the law of $\xi$, this makes $a_0$ and $\eta$ covariant, by \textup{(2c)} and \textup{(2g)}, and it makes $\mathsf U(\operatorname{Ad}_u\eta;V'^u)$ the gauge transform of $\mathsf U(\eta;V')$, as shown for the clauses without a size in Step~8.

The members, $1_{X^\approx}(U)$ and $1_{\rm ind}$ are invariant; $b_\square$ and $\chi$ depend only on $|\xi_b|$; $A_\square$ is invariant by \textup{(6e)}, $T_\square$ by \textup{(5b)}, $T_Y$ by \textup{(4b)} and $n_j$ by \textup{(4a)}. Hence every letter satisfies $\lambda_\alpha(\operatorname{Ad}_u\xi,V'^u,\operatorname{Ad}_ue)=\lambda_\alpha(\xi,V',e)$, and $K(\gamma)(V'^u,\operatorname{Ad}_ue)=K(\gamma)(V',e)$; the map $e\mapsto\operatorname{Ad}_ue$ preserves $de$.

\item[\textup{(iii)}] For $V'$, this is \textup{(8c)}, with the dependence on the field established in the proof of \textup{(8c)}. The choice $e_P$ enters at plaquettes of marks in $\mathsf Y^{+(s'-1)}$, and where $w'\ne0$, these marks are cells of $\mathsf Y$.

\item[\textup{(iv)}] This is \textup{(8c)}. The index set depends on $D'$ and not on $e$; the sets $\mathsf Y_i$ are disjoint, so the bounded factors depend on disjoint sets of $e_P$, and Fubini's theorem applies.

\end{enumerate}

\end{enumerate}

\end{proof}

\emph{Step 11: clauses \textup{(c)} and \textup{(d)} of \Cref{Y-H-C9}.} Let $z\mapsto u_z$ be an affine map from an open set $\Omega\subset\mathbb C$ into $M_{16}^{\rm su}(x)$. Then the number $\mathsf V'_z$ and all the values $E'^a_z(Y,V')$, $E'^h_z(Y,V')$, $w'_z(\mathsf Y,V')$ are holomorphic in $z$. If $u$ is real, then $F_\ell(u)$ and $\mathsf V'$ are real; if moreover $\rho[u;\ell]>0$, then $(\rho[u;\ell])'>0$.

\begin{proof}[Proof of Step~11] At fixed $(\xi,V',e)$ every letter is holomorphic in $z$ on $\Omega$, as a product of affine values $w_z(X,U)$, of exponentials of the affine functions $E^a_z$ and $E^h_z$, and of factors that contain no member.

The alphabet, the adjacency, the compatibility, the cores, $\chi$, $U$, the indicators, whether a letter is actual, whether a cluster or a set is on or off, whether a set is long, and the alternative of \Cref{Y-H-LP}, which depends only on the supports of $\Lambda_\zeta$, are all independent of $z$.

The majorant of Step~6 at $A=16$ is integrable and independent of $z$. Hence $K_z(\mathfrak B)$ is holomorphic, by the theorems of Morera and Fubini.

Each $\Phi_z(C)$ is a polynomial in the $K_z$, and $|\Phi_z|\le\hat\Phi$ at $A=16$. Therefore \textup{(Fm1)} to \textup{(Fm3)} and the sum in \textup{(Fm4)} converge uniformly in $z$, and $e^{-\Phi^\natural_z}-1$ is holomorphic.

The mean over $e$ of a bounded and jointly measurable family that is holomorphic in $z$ is again holomorphic; here the measurability is shown in the verification of hypothesis~(i) of \Cref{Y-H-L52} in Step~10\,(9), together with \textup{(6e)}.

Reality: at real $u$ and $G$-valued $V'$, every letter is real, by \textup{(2d)}, \textup{(4a)}, \textup{(4b)}, \textup{(5b)}, \textup{(5c)} and \textup{($\Gamma$6)}, and $A_\square>0$ or $A_\square=0$. Hence $K,\Phi,\tilde{\mathsf D},P,n_j,w_j,\mathsf V_0$ are real.

The block kernel is strictly positive. Integrating the positive
fine density against its product gives a positive coarse density at
every $V'$. \end{proof}

\emph{Step 12. Contraction along a segment.}
Let $u_1,u_2\in M_B^{\rm su}(x)$ for $B=1$ or $2$, and put
$d=\|u_2-u_1\|_x$. The case $d=0$ is immediate. Otherwise
$0<d\le2B$, and the affine family
\[
 u_z=u_1+\frac{z}{d}(u_2-u_1)
\]
preserves the linear conditions on the members and satisfies
$\|u_z\|_x\le B+|z|$. Its values lie in $M_{16}^{\rm su}(x)$ for
$|z|<16-B$. At each $t\in[0,d]$ a disc of radius $16-3B$ lies in
that domain. For any scalar output $f(z)$ with bound $\mathsf m$ on
$|z|<16-B$, clause \textup{(c)} and Cauchy's estimate give
\[
 |f(d)-f(0)|\le \frac{d\mathsf m}{16-3B}.
\]
For the analytic part take $f=E'^{\rm old}[u_z](Y,V')$, since
$E'^{\rm new}$ is independent of $z$. Its bound is $16\cdot0.248$
times the analytic cap in clause \textup{(b)}. For $E'^h$ and $w'$
the bounds are their caps at $x'$, using the strong cap for $w'$.
After dividing by the three caps and taking the suprema in the
strong distance, we obtain the factors
$3.968/(16-3B)$, $1/(16-3B)$ and $1/(16-3B)$. At $B=2$ the
largest is $0.3968\le0.4$, which proves clause \textup{(e)}.

At $A=1$, clauses \textup{(a)}, \textup{(b)} and \textup{(d)}
also give the class inclusion of \textup{(S-step)} for $x\ge x_9$.
Indeed the analytic cap is multiplied by $0.35+0.248<1$, the output
has the strong cap, and positive real densities retain positivity
and reality. At $A=2$ the analytic factor is $0.846<1$, again with
the strong output cap. The inputs from Sections~\ref{Y-D-1}
through~\ref{Y-F-1} concern the chart, energy and tests of the
step, which are fixed independently of the input members.
\end{proof}

\subsection{Contraction and row perturbations}\label{Y-I-1}

To compare cutoffs we first express two densities over one family
of rows. The difference in their reference energies is assigned to the
analytic term in tame cells and to a rough weight in each remaining
cell. The change is small in the weak distance. The complex step then
restores the strong caps and contracts the difference, so both runs can
continue over the same rows.

\subsubsection*{Item 0: what is used; notation}

\begin{remark}[What is used of other sections] From other sections we use \Cref{Y-H-C9}, with its clauses (a), (b), (c) and (d), its data \textup{(SD)}, its floor $x_9$ and Step~0\,(iv) of its proof, together with \textup{(Y.0)(0b)}, \textup{(Y.4)} and \textup{(Y.7)} of the floors \textup{(Y.0)} to \textup{(Y.9)}; item~8 of \S\ref{Y-H-1}; \Cref{Y-J-Co50}, with its chains, with \textup{(RU)} and with its hypotheses, among which are \textup{(X12)} and \textup{(X13)}; Steps~1, 2\,(d) and item~1 of Step~4 of the proof of \Cref{Y-J-ThA}; Step~S4\,(4) of the proof of \Cref{Y-B-C1}; and \Cref{Y-A-Fr}, \Cref{Y-A-Th29}\,(a), \Cref{Y-A-P4}\,(a), \Cref{Y-A-L26}\,(a), \Cref{Y-A-Th36} and \Cref{Y-A-P5}. \Cref{Y-J-Co50} belongs to the later \S\ref{Y-J-1}, where it is stated and proved in Step~1a; every statement on which a proof in Step~1a rests is stated either in Step~1a itself or in \S\ref{Y-A-1}, and the proof of \Cref{Y-J-Co50} uses nothing of this section.\end{remark}

\begin{definition}[Notation] Throughout this section $G=\operatorname{SU}(2)$. The letter $V$ stands for a $G$-valued field of the level at $x$, and $\mathsf W$ for an $\operatorname{SL}(2,\mathbb C)$-valued field. The letters $\mathsf V_1,\mathsf V_2,\mathsf V',\mathsf V'_1,\mathsf V'_2$ are numbers. We write $D(V)$ for the set of all rough plaquettes, $\mathsf D_x$ and $\mathsf D^{\rm in}_x$ for distances of members, $D_y$ for a difference of kernels of rows, and $\mathcal D$, $\mathcal P$ for families of sets of cubes. The letter $\mathsf b$ is the label of a run, a number $\ge x_\flat$, whereas $\mathsf b(x)$ is the number $2-1/\log x$ of \textup{(u2)}. A prime on a letter means that the letter is taken at $x'=x/L$.\end{definition}

\begin{remark}[Letters with two senses] The letter $u$ is a system of members $(E^a,E^h,w)$, but in $V^u$, $u(y)$, $\mathbf 1^u$ and $\operatorname{Ad}_{u(y)}$ it is a gauge transformation.

Only in Step~9\,(E4), $F$ and $\bar F$ each stand for a real 2-cochain, $X'$ for a coarse 2-form (and not for $X$ at $x'$), and we set $\mathsf B=1+\ell^{\mathsf T}\ell$.

The radius $r(x)$ is written $r$ wherever $x$ is fixed.

In Steps~9 and~10, $\sigma_3=\operatorname{diag}(1,-1)$.

We write \textup{(9a)} to \textup{(9d)} for the clauses (a) to (d) of \Cref{Y-H-C9}.\end{remark}

\subsubsection*{Item 1: letters}

\begin{definition}[Letters] We give the letters in the order \textup{(O0)} to \textup{(O14)} of the complete table.

The complete table appears in item~1 of \S\ref{Y-J-1}. Its letters are defined in that item as well as in ``The order of the constants''\,\textup{(0)} and \textup{(6)} in \S\ref{Y-A-1}, in item~C.I~1.2 of \S\ref{Y-C-1}, in \textup{(O3)}, \textup{(O4)}, \textup{(O6)}, \textup{(O7)}, \textup{(O10)}, \textup{(O13)} in item~3 of \S\ref{Y-G0-1}, and in \textup{(O13)} in item~1 of \S\ref{Y-H-1}.

\begin{enumerate}\item[\textup{(O0)}] The pure numbers (\Cref{Y-A-PureNumbers}) are $b_*$, $\gamma_*$, $\lambda_*=\gamma_*^{1/2}/44$, $K_*=3\gamma_*^{-1/2}$ ($>1$), $L_A$, and the others of the table. Here $\gamma_*=\mathsf r_\sharp^2/12296\le9.05513\cdot10^{-25}<1$ (``The order of the constants''\,\textup{(0)} in \S\ref{Y-A-1}), whence $K_*>3$. We put $c_0:=108\,b_*\gamma_*^{-1/2}\mathsf g(\lambda_*/4)^3$, $c_{\rm fix}:=48\,\mathsf g(\lambda_\flat)C_1c_0^{1/2}$, $L_B:=\max\{L_A,c_{\rm fix}^{10}\}$ and $\vartheta_1:=c_{\rm fix}L^{-1/2}$ (``Rates, constants'', Step~1a of \S\ref{Y-J-1}); all of these are pure numbers except $\vartheta_1$, which depends on $L$. Further, $\mathsf g(\lambda):=\sum_{n\in\mathbb Z^3}e^{-\lambda|n|_\infty}$. The definition ``Rates, constants'' in Step~1a of \S\ref{Y-J-1} fixes $\lambda_\flat:=\lambda_*/16$ and $C_1:=\sqrt3\,K_*\,\mathsf g(\lambda_*/2)^{3/2}$; hence $C_1>1$, because $K_*>3$ and $\mathsf g(\lambda_*/2)\ge1$.

\item[\textup{(O1)}] $d_0\in(0,\frac12)$.

\item[\textup{(O2)}] $0<c_\Xi\le\frac12(\frac12-d_0)$.

\item[\textup{(O3)}] The radius $\rho_0$ is given by \eqref{eq:row-mask-radius} of \S\ref{Y-A-1}.

\item[\textup{(O4)}] $L$ is odd with $L\ge L_*^{\rm su}$ of the table; $c_L=(L^2-1)/(6L^2)$; and $a\ge3/(c_LL)$. In particular, $L\ge L_B$.

\item[\textup{(O5)}] $\kappa\ge6$.

\item[\textup{(O6)}] $\theta^\flat\le1$, $\theta_{\mathbb C}=\theta^\flat/6$ and $\vartheta_0^\flat=\theta^\flat/24$, where $\theta^\flat=\min\{\theta_1,\gamma/(2A^\flat)\}$ and $\theta_1=\min\{1,\lambda_*L/12\}$, as in item~C.I~1.2 of \S\ref{Y-C-1} at $(K,\lambda)=(K_*,\lambda_*)$ and in ``The order of the constants''\,\textup{(6)} in \S\ref{Y-A-1}. Further, $\theta_D=\frac12\min\{\lambda_2^*,\theta_{\mathbb C},\theta_J\}$, $C_D$ is as in the table, and $c_R=64/\theta_D$. We have $c_R\ge768$, since $\theta_D\le\theta_{\mathbb C}/2\le1/12$.

\item[\textup{(O7)}] $M=L^{m'}$ with $m'\ge1$, subject to \textup{(M2)}, \textup{(M4$^\flat$)} and \textup{(M5r)} $M\ge\rho_0$; the condition \textup{(M4$^\sharp$)} is a consequence of \textup{(M4$^\flat$)}.

\item[\textup{(O8)}] $\mathfrak p_0>2$ and $A_0>0$; $p_0(x)=A_0(\log x)^{\mathfrak p_0}$ and $\varepsilon=p_0x^{-1/2}$, so that $x\varepsilon^2=p_0^2$.

\item[\textup{(O10)}] This section defines $c_Q^{\rm su}:=306\,M^3\,\mathsf g(\lambda_\flat)^2$.

\item[\textup{(O11)}] $0<c_h\le\delta_3'/2$.

\item[\textup{(O12)}] $E_a$ is as in the table. In particular, $E_a>0$.

\item[\textup{(O13)}] $x_\flat\ge e^2$ with $L\,x_\flat\ge x_9$ (\Cref{Y-H-C9}). The number $x_9:=1+\inf S_9$ is defined in \textup{(O13)} of item~1 of \S\ref{Y-J-1} and in \textup{(O13)} of item~1 of \S\ref{Y-H-1}, each of which gives the set $S_9$. Of $x_9$ this section uses that \textup{(Y.0)(0b)}, \textup{(Y.4)}, \textup{(Y.7)} and \Cref{Y-H-C9} hold at every $x\ge x_9$, and that $x_9\ge x_{\rm fl}\ge L^2e$. Further, $x_\flat$ is chosen such that for all $x\ge x_\flat$ \begin{enumerate}\item[\textup{(X14)}] $p_0(x)\ge1$; \item[\textup{(X15)}] $c_\Xi p_0(x)^2\ge1$.\end{enumerate} \Cref{Y-I-CorD} uses in addition $16c_a(x)\le\frac12$ for all $x\ge x_\flat$, together with \textup{(X12)} and \textup{(X13)} of \Cref{Y-J-Co50}, which read \begin{enumerate}\item[\textup{(X12)}] $3K_*\mathsf g(\lambda_*/2)\,x_\flat^{-e_R}/(1-L^{-e_R})\le\frac19$, where $e_R:=\lambda_*c_R/2$; \item[\textup{(X13)}] $16\,x_\flat^{-e_R/4}\le1$.\end{enumerate} The full entry, with $x_\flat\ge x_\flat^{(1)}$, $x_\flat\ge x_\flat^0$ and \textup{(X16)} to \textup{(X18)}, is \textup{(O13)} of item~1 of \S\ref{Y-J-1}.

\item[\textup{(O14)}] The numbers of this entry, $\vartheta_w$, $n_w$, $\vartheta_{10}$, $\vartheta$, $c_4$, $\Lambda(\mathsf b)$, $K_0(\mathsf b)$ and $k_0$, are all defined after \textup{(O13)}, and no floor depends on any of them. Here $\vartheta_w=0.248$ and $n_w=0.35$ are those of \Cref{Y-H-C9}\,(b). We put $\vartheta_{10}:=16\vartheta_w/10=0.3968$, and the $\vartheta$ of the table is $\vartheta:=(16\vartheta_w+\frac12)/10=0.4468$. Thus $\vartheta_{10}\le0.4\le0.4318\le\vartheta\le0.45$, so that the factor $0.4$ of \Cref{Y-H-C9}\,(e), the factors $0.4318$ and $0.45$ of \Cref{Y-I-ThA}\,(iv), and $\vartheta$ itself, each as the factor in $\mathsf D_{x'}\le\text{factor}\cdot\mathsf D^{\rm in}_x$ for pairs in $M_2^{\rm su}(x)$, are upper bounds of $\vartheta_{10}$. Finally, $c_4:=15\,c_Q^{\rm su}\,C_1^2/E_a$, $\Lambda(\mathsf b):=\log_L(c_4\mathsf b^{1/2})$, $K_0(\mathsf b):=\frac52\Lambda+\frac94$ and, for the length $K$ of a run, an integer $\ge1$ (the $K$ of \Cref{Y-I-CorD}), $k_0:=\lceil(5K+10\Lambda)/9\rceil$.\end{enumerate}

The caps are $c_a(x)=E_ap_0^3x^{-1/2}$ and $c_h(x)=e^{-c_hp_0^2}$. No number defined here depends on $N$, on a member, or on more of the rows than their range. As their formulas show, $p_0$, $\varepsilon$ and the caps are functions of $x$; apart from these, none depends on $x$.\end{definition}

\subsubsection*{Item 2: the level at $x$}

\begin{definition}[Lattice] Let $N$ be even, with $M\mid N$ and \textup{(Tor)} $N\ge2M\,s(x)$; this $N$ is the coarse side of the step that starts at $x$. Let $\mathbb T=\mathbb Z^3/NL\mathbb Z^3$. Plaquettes are written $p=(z;\mu\nu)$ with $\mu<\nu$, and $\pi(p)=z$ is the base site of $p$; the base site of a bond $(z,\mu)$ is $z$. We set $|y-p|:=|y-\pi(p)|_\infty$, measured periodically (\Cref{Y-A-OwnUnits}).

Cubes are the boxes of $M^3$ sites of one fixed grid of $\mathbb T$. Each bond and each plaquette belongs to the cube of its base site, and ``$p$ in $X$'' or ``$X$ holds $p$'' means that $p$ belongs to a cube of $X$. The distance of cubes is the sup distance counted in cubes, and $X^{+j}$ consists of the cubes within $j$ of $X$. By $\mathcal D$ we denote the family of all non-empty sets of cubes that are connected through faces, edges and corners, and by $\mathcal P$ the family of all non-empty finite sets of cubes; further, $|X|_c=\#{\rm hull}(X)$. Here ${\rm hull}(X)$ is a member of $\mathcal D$ that holds $X$ and has the least number of cubes, singled out by an order (Step~S4\,(4) of the proof of \Cref{Y-B-C1}); thus $|X|_c$ is the least number of cubes of a member of $\mathcal D$ holding $X$, as in item~2 of \S\ref{Y-J-1}.

The bonds of $X$ are those with both ends in cubes of $X$. The paths and plaquettes of a set $A$ of cubes are those all of whose sites lie in $A$, while ``$p$ in $A$'' means that the base site of $p$ lies in $A$. We call a field on the bonds of $A$ tame there when every plaquette of $A$ has $\phi_p\le\varepsilon$ (see item~1 of \S\ref{Y-G0-1} and item~2 of \S\ref{Y-H-1}). If $1_A(V)=1$, the restriction of $V$ to the bonds of $A$ is tame there, since a plaquette of $A$ is based in $A$. The suprema of $E^h$ in $\|\cdot\|_x$, $\mathsf D^{\rm in}_x$ and $\mathsf D_x$ are taken over the fields that are tame there.\end{definition}

\begin{definition}[At $x$] At $x$ we put $\varrho(x)=\lceil c_R\log x\rceil$, $r(x)=\lceil20c_R\log x/M\rceil$ ($\ge1$), $s(x)=8r(x)$ and $\ell_w(x)=6M\,r(x)+12$, and we write $X^{\approx}=X^{+r(x)}$ and $X^{\approx\approx}=X^{+s(x)}$.\end{definition}

\begin{definition}[Group] Let $\mathfrak g$ be the space of hermitian traceless $2\times2$ matrices, with $\langle Y,Y'\rangle=\frac12\operatorname{tr}YY'$ and $|Y|^2=\langle Y,Y\rangle$; this $|Y|$ is the operator norm. We put $\operatorname{ang}(g)=\arccos(\frac12\operatorname{tr}g)\in[0,\pi]$. For $g\ne-1$, $\mathsf L(g)$ is the unique $Y\in\mathfrak g$ with $e^{iY}=g$ and $|Y|<\pi$, and $\mathsf L(-1):=0$. On complex $2\times2$ matrices, $|\cdot|$ denotes the operator norm and $Y\cdot Y':=\frac12\operatorname{tr}YY'$; hence $|Y\cdot Y'|\le|Y||Y'|$, and on $\mathfrak g$ we have $Y\cdot Y'=\langle Y,Y'\rangle$. For $g\in\operatorname{SL}(2,\mathbb C)$ one has $|g^{-1}|=|g|$; with $\operatorname{Ad}_gY=gYg^{-1}$ we get $|\operatorname{Ad}_gY|\le|g|^2|Y|$ and $(\operatorname{Ad}_gY)\cdot(\operatorname{Ad}_gY')=Y\cdot Y'$.\end{definition}

\begin{definition}[Fields] For a field $V$ we have $V(\bar b)=V(b)^{-1}$; $V(\pi)$ is the ordered product along a path $\pi$, and $V^u(b)=u(b_-)V(b)u(b_+)^{-1}$. The curvature is \[\partial V(z;\mu\nu)=V(z,\mu)V(z+e_\mu,\nu)V(z+e_\nu,\mu)^{-1}V(z,\nu)^{-1}.\] We set $\phi_p=\operatorname{ang}\partial V(p)$, $F_p=\mathsf L(\partial V(p))$ and $D(V)=\{p:\phi_p>\varepsilon\}$, and we put $1_Y(V):=1\{\phi_p\le\varepsilon\text{ for all }p\text{ in }Y\}$; when it equals $1$ we say that $V$ is tame on $Y$.\end{definition}

\begin{definition}[Transports] For sites $y,z$ whose sup distance is less than half the side, $\Gamma(y,z)$ is the coordinate path from $y$ to $z$ that moves first in the direction $e_1$, then $e_2$, then $e_3$, each time the short way; and $T(y,z):=\operatorname{Ad}_{V(\Gamma(y,z))}$.\end{definition}

\begin{definition}[Rows] Rows, columns and range are those of \Cref{Y-A-Rows}, with one exception, the mask: the mask there has the one radius $\rho_0$, whereas the mask of this section, defined below, has two radii. Rows $\ell$ on $\mathbb T$ of range $\le\varrho$ consist of a Euclidean space $\mathsf H_y$ of finite dimension for each site $y$ and of columns $\ell_y(p)\in\mathsf H_y$, with $\ell_y(p)=0$ if $|y-p|>\varrho$. We put $\mathsf E_y(p,p'):=\langle\ell_y(p),\ell_y(p')\rangle$, so that $\mathsf E_y=\ell_y^{\mathsf T}\ell_y$ (``Site forms'', Step~1a of \S\ref{Y-J-1}). For a family $D=(D_y)$ of kernels, \[\|D\|_{(\lambda)}:=\sup_{y,p,p'}|D_y(p,p')|\,e^{\lambda(|y-p|+|y-p'|)}\] (``Norms of a family of symmetric operators'', Step~1a of \S\ref{Y-J-1}). The empty rows $\ell=\emptyset$ are those with all $\mathsf H_y=0$.\end{definition}

\begin{definition}[Clip, dressed rows, mask, energy] The clip is $\zeta_p:=F_p\min\{1,\varepsilon/\phi_p\}$, with $\zeta_p=0$ if $\partial V(p)=1$ or $-1$; hence $|\zeta_p|\le\varepsilon$ at every $V$, and $\zeta_p=F_p$ if $\phi_p\le\varepsilon$. Dressing the rows with the transported clip, we put $\ell^V_y\zeta:=\sum_p\ell_y(p)\otimes T(y,\pi(p))\zeta_p\in\mathsf H_y\otimes\mathfrak g$. The mask is $m_y(V):=1$ if no plaquette of $D(V)$ has $|y-p|\le\rho_0$ and no plaquette with $\phi_p>1$ has $|y-p|\le\varrho(x)+1$, and $m_y(V):=0$ otherwise. With $S_\varepsilon(t)=t^2/2$ ($t\le\varepsilon$), $d_0\varepsilon t$ ($t>\varepsilon$), the energy is \[\mathrm{En}_\ell(V):=\sum_pS_\varepsilon(\phi_p)+\frac12\sum_ym_y(V)\,|\ell^V_y\zeta|^2.\] The mask, the clip and the paths depend on $V$, $\varepsilon$, $\rho_0$ and $\varrho(x)$, and on no row.\end{definition}

\begin{proposition}[Invariance]\label{Y-I-P6} Let $u$ be a gauge transformation, and let $V$ be replaced by $V^u$. Then $\phi_p$, $D(V)$, $m_y$ and $1_Y$ are invariant; $\zeta_p$ becomes $\operatorname{Ad}_{u(\pi(p))}\zeta_p$; and $T(y,\pi(p))\zeta_p$ becomes $\operatorname{Ad}_{u(y)}T(y,\pi(p))\zeta_p$.\end{proposition}

\begin{proof} One has $\partial V^u(p)=u(\pi(p))\,\partial V(p)\,u(\pi(p))^{-1}$. As $\operatorname{ang}$ is a class function, $\phi_p$ is unchanged, and so are $D(V)$, $m_y$ and $1_Y$, which are defined through the angles $\phi_p$ alone. Next, $\mathsf L(ugu^{-1})=\operatorname{Ad}_u\mathsf L(g)$, where at $g=-1$ both sides are $0$. Therefore $F_p=\mathsf L(\partial V(p))$ is mapped to its image under $\operatorname{Ad}_{u(\pi(p))}$, and so is $\zeta_p$, the factor $\min\{1,\varepsilon/\phi_p\}$ being invariant. Finally, $V^u(\Gamma(y,z))=u(y)V(\Gamma(y,z))u(z)^{-1}$. For $z=\pi(p)$, in $\operatorname{Ad}_{u(y)V(\Gamma(y,\pi(p)))u(\pi(p))^{-1}}\operatorname{Ad}_{u(\pi(p))}\zeta_p$ the factor $u(\pi(p))^{-1}$ cancels $u(\pi(p))$, so $T(y,\pi(p))\zeta_p$ is mapped to its image under $\operatorname{Ad}_{u(y)}$, which is an isometry of $\mathfrak g$.\end{proof}

\begin{definition}[The step] For bounded Borel $\rho$, \[\rho'(V')=\int\rho(U)\prod_ck_{ax}(V'(c),\mathsf M(U)(c))\,dU,\] and for complex $\rho$ the definition extends by linearity. The map $\rho\mapsto\rho'$ is linear in $\rho$, it depends on $x$, $a$, $L$ and on no row, and $k_{ax}>0$. The kernel $k_s$, its normaliser $z(s)$ and the average $\mathsf M$ are defined in item~2 of \S\ref{Y-J-1}, ``The step''.\end{definition}

\subsubsection*{Item 3: domains, members, distances, the class}

\begin{remark} The definitions of this item are those of item~3 of \S\ref{Y-H-1}.\end{remark}

\begin{definition}[Domain] For $X\in\mathcal D$, let $\mathcal U_X(x)$ be the set of maps $\mathsf W$ from the oriented bonds of $X^{\approx}$ to $\operatorname{SL}(2,\mathbb C)$ with $\mathsf W(\bar b)=\mathsf W(b)^{-1}$ that satisfy \textup{(u1)}, \textup{(u2)} and \textup{(u3)}:

\begin{enumerate}\item[\textup{(u1)}] $|\mathsf W(b)|<1+\vartheta_0^\flat(1-1/\log x)$;

\item[\textup{(u2)}] $|\mathsf W(\pi)|<\mathsf b(x):=2-1/\log x$ for every path $\pi$ in $X^{\approx}$ of at most $\ell_w(x)$ bonds;

\item[\textup{(u3)}] $|\partial\mathsf W(p)-1|<2\varepsilon$ for every plaquette $p$ of $X^{\approx}$.\end{enumerate}

This set is open and invariant under $G$-valued gauge transformations, and it contains every $G$-valued field that is tame on $X^{\approx}$, because the norms of such a field are $1$ and $|\partial V(p)-1|=2\sin(\phi_p/2)\le\varepsilon$.\end{definition}

\begin{definition}[Members] Let $\mathfrak M(x)$ be the complex linear space of all triples $u=(E^a,E^h,w)$ that have $\|u\|_x<\infty$ and satisfy the clauses without a size, which are the three conditions on $E^a$, $E^h$ and $w$ that follow. For $X\in\mathcal D$, $E^a(X,\cdot)$ is holomorphic on $\mathcal U_X(x)$, invariant under $G$-valued gauge transformations, and $0$ at the pure gauges $\mathbf 1^u$. For $X\in\mathcal D$, $E^h(X,\cdot)$ is a gauge invariant complex Borel function defined on those $G$-valued fields on the bonds of $X^{\approx}$ which are tame there. For $X\in\mathcal P$, $w(X,\cdot)$ is a gauge invariant complex Borel function of the restriction of $V$ to the bonds of $X^{\approx\approx}$.

The number $\|u\|_x$ is the largest of the three numbers \[\FitDisplay{\dimexpr\linewidth-1pt\relax}{\sup_X\sup_{\mathcal U_X(x)}|E^a(X,\cdot)|\,e^{\kappa(\#X-1)}/c_a(x),\qquad\sup_X\sup|E^h(X,\cdot)|\,e^{\kappa(\#X-1)}/c_h(x),\qquad\sup_{X,V}|w(X,V)|\,e^{\kappa(|X|_c-1)},}\] and $\|\cdot\|_x$ is a norm on $\mathfrak M(x)$. For $A\ge1$ we put $M_A^{\rm su}(x):=\{u\in\mathfrak M(x):\|u\|_x\le A\}$; these sets are given by \textup{(ca$_A$)}, \textup{(ch$_A$)}, \textup{(c$\Xi_A$)} of item~3 of \S\ref{Y-H-1}, and the cap of $w$ in them is the weak cap. The strong cap is the bound \[\FitDisplay{\dimexpr\linewidth-1pt\relax}{|w(X,V)|\le\exp\{-c_\Xi x\varepsilon\sum_{p\in D(V)\cap X}\phi_p-\kappa(|X|_c-1)\}}\] at every $V$. We call $u$ real if $E^a$ is real at $G$-valued fields and $E^h$, $w$ are real.\end{definition}

\begin{definition}[Density and class] The density of $u$ over the rows $\ell$ is \[\rho[u;\ell](V):=\exp\bigl\{-x\mathrm{En}_\ell(V)+\sum_{X\in\mathcal D}1_{X^{\approx}}(V)[E^a(X,V)+E^h(X,V)]\bigr\}\cdot\Xi_w(V),\] where $\Xi_w(V):=\sum_{\{X_i\}}\prod_iw(X_i,V)$ is summed over the finite families of $X_i\in\mathcal P$ with pairwise distance $>s(x)$, each of which holds a plaquette of $D(V)$, and whose union contains $D(V)$; these are the families of \textup{(c$\Xi$)}. If $D(V)=\emptyset$, the only such family is the empty one, and $\Xi_w=1$. For rows $\ell$ on $\mathbb T$ of range $\le\varrho(x)$, the class $C^{\rm su}_\flat(x;\ell)$ is the set of all densities $e^{\mathsf V_1}\rho[u;\ell]>0$ with $\mathsf V_1$ real and $u\in M_1^{\rm su}(x)$ real with the strong cap.\end{definition}

\begin{definition}[Distances] We put $\mathsf D^{\rm in}_x(u_1,u_2):=\|u_1-u_2\|_x$, the weak distance. The strong distance $\mathsf D_x(u_1,u_2)$ is defined in the same way, except that in the third supremum the weight $e^{\kappa(|X|_c-1)}$ is replaced by $\exp\{c_\Xi x\varepsilon\sum_{D(V)\cap X}\phi_p+\kappa(|X|_c-1)\}$. Hence $\mathsf D^{\rm in}_x\le\mathsf D_x$, and $\mathsf D_x(u_1,u_2)\le2$ when $u_1,u_2$ both lie in $M_1^{\rm su}(x)$ and have the strong cap.\end{definition}

\begin{definition}[The level at $x'$] At $x'$, the definitions of item~2, with two exceptions, and those of item~3 hold on $\mathbb T'=L\mathbb Z^3/NL\mathbb Z^3$, in its own units $\mathbb Z^3/N\mathbb Z^3$, with cells in the place of cubes. The exceptions are the definition ``The step'' and part of the definition ``Lattice'', namely the conditions that $N$ is even, that $M\mid N$ and that \textup{(Tor)} holds, together with the role of $N$ as the coarse side of the step starting at $x$ and the torus $\mathbb T=\mathbb Z^3/NL\mathbb Z^3$.\end{definition}

\begin{remark} Of the conditions on the side $N$ in ``Lattice'', only the following is used at $x'$: the side $N$ of $\mathbb T'$ is even, $M\mid N$, and $N\ge2M\,s(x)\ge2M\,s(x')=16M\,r(x')$, because $r$ is non-decreasing. Nothing at $x'$ depends on $N$ being $L$ times the coarse side of a further step: \textup{(SD)} does not provide this, at the last level of a run it fails in general, and no step starts there. Four facts suffice instead. First, wherever $x'$ is used (in item~4 and in Thms.~\ref{Y-I-ThA} and~\ref{Y-I-ThC}), we have $x\ge x_9\ge x_{\rm fl}\ge L^2e$ (\textup{(O13)} in item~1 of \S\ref{Y-J-1}), and so $x'\ge Le>e$. Second, \textup{(F2)} therefore holds at $x'$, by the argument for it in Step~0\,(Fl), which uses only $x\ge e$ and $c_R\ge768$; that is, $\varrho(x')+1<M\,r(x')\le N/16$. Third, the paths $\Gamma(y,z)$ of $\mathrm{En}_{\ell'^{\rm c}}$ and the mask at $x'$ are therefore defined, since the sup distances involved are less than half the side $N$. Fourth, every box $\square^{+r(x')}$ of cells, whose side is $(2r(x')+1)M\le3M\,r(x')\le3N/16$, embeds. This is the level of item~2 of \S\ref{Y-J-1}, ``Lattice, cubes''.\end{remark}

\subsubsection*{Item 4. Inputs from the complex step}
The tags \textup{(9a)} through \textup{(9d)} refer to the corresponding
clauses of Thm.~\ref{Y-H-C9}. They give the exact image, the caps,
holomorphy along affine families, and reality and positivity. We write
$n_w=0.35$ for the analytic contribution independent of the input and
$\vartheta_w=0.248$ for the coefficient multiplying the input cap.

\subsubsection*{Item 5: the conditions of \Cref{Y-I-LB}}

\begin{definition}[Five conditions] The five conditions of the change of rows are these.

\begin{enumerate}\item[\textup{(M5r)}] $M\ge\rho_0$.

\item[\textup{(X15)}] $c_\Xi p_0(x)^2\ge1$.

\item[\textup{(F1)}] $\varepsilon(x)\le1/40$.

\item[\textup{(F2)}] $\varrho(x)+2\le M\,r(x)$.

\item[\textup{(F3)}] $(2r(x)+1)^3c_a(x)\le34$.\end{enumerate}\end{definition}

\begin{remark} In the setting of the table, at every $x\ge L\,x_\flat$, condition \textup{(M5r)} is \textup{(O7)} and condition \textup{(X15)} is \textup{(O13)}, while \textup{(F1)}, \textup{(F2)} and \textup{(F3)} follow from $x\ge x_9$ and $c_R\ge768$ (Step~0), with $\varepsilon<1/680$ and $(2r(x)+1)^3c_a(x)\le x^{-1/4}/(192e^{11})<10^{-7}$.\end{remark}

\subsubsection*{Item 6: the statements}

\begin{theorem}[One step contracts from the weak to the strong distance]\label{Y-I-ThA} Let $x\ge x_9$, let the data \textup{(SD)} at $x$ be given, and let $u_1,u_2\in M_2^{\rm su}(x)$, complex members with the weak cap. Then \textup{(9b)} and \textup{(9c)} imply the following.

\begin{enumerate}\item[\textup{(i)}] $F_\ell(u_i)$ lies in $M_1^{\rm su}(x')$ and has the strong cap at $x'$; and $\sup_{\mathcal U_Y(x')}|E'^a(Y,\cdot)|\le(n_w+2\vartheta_w)\,c_a(x')\,e^{-\kappa(\#Y-1)}$ for every $Y$, where $n_w+2\vartheta_w=0.846$ at the numbers of \textup{(O14)}.

\item[\textup{(ii)}] We have \[\mathsf D_{x'}(F_\ell(u_1),F_\ell(u_2))\le\vartheta_{10}\,\mathsf D^{\rm in}_x(u_1,u_2)\le\vartheta_{10}\,\mathsf D_x(u_1,u_2),\qquad\vartheta_{10}=0.3968;\] for the three members the factors are $0.3968$, $0.1$ and $0.1$.

\item[\textup{(iii)}] For $u_1,u_2\in M_1^{\rm su}(x)$ the factor is $\max\{16\vartheta_w/13,1/13\}=3.968/13=0.30523\dots\le0.3053$.

\item[\textup{(iv)}] In a robust form: if, for $E'^a$, the two sums of \textup{(9b)} are replaced by the single requirement that $\sup_{\mathcal U_Y(x')}|E'^a(Y,\cdot)|\le(\vartheta_wA+n_w)\,c_a(x')\,e^{-\kappa(\#Y-1)}$ for every $Y$, with numbers $\vartheta_w,n_w\ge0$ and $2\vartheta_w+n_w\le1$, while every other clause of \textup{(9b)} stays in force (the sum for $E'^h$, the strong cap of $w'$, and the clauses without a size), then \textup{(i)} still holds and the factor in \textup{(ii)} becomes $\max\{(16\vartheta_w+n_w)/10,1/10\}$. That factor never exceeds $0.8$; it is $0.4318$ at $(0.248,0.35)$ and $0.45$ at $(\frac14,\frac12)$.\end{enumerate}\end{theorem}

\begin{lemma}[A density over other rows]\label{Y-I-LB} Let $x>e$, let the lattice of item~2 be given, and let $\ell,\tilde\ell$ be rows on $\mathbb T$ of range $\le\varrho(x)$. Let $\tilde{\mathsf E}_y(p,p'):=\langle\tilde\ell_y(p),\tilde\ell_y(p')\rangle$, the kernel $\mathsf E_y$ of ``Rows'' with $\tilde\ell$ in the place of $\ell$, and let $D_y:=\tilde{\mathsf E}_y-\mathsf E_y$, $\delta:=\|D\|_{(\lambda_\flat)}$ and \[\delta_\Delta:=c_Q^{\rm su}p_0^2\delta/c_a(x)=(c_Q^{\rm su}/E_a)\,x^{1/2}\delta/p_0(x).\] Let the five conditions of item~5 hold, and let $\delta_\Delta\le1$. Then each real $\tilde u=(\tilde E^a,\tilde E^h,\tilde w)\in M_1^{\rm su}(x)$ having the strong cap admits a real $u^\flat=(\tilde E^a+E^\Delta,\tilde E^h,w^\flat)\in M_2^{\rm su}(x)$ such that \textup{(B1)} to \textup{(B4)} hold:

\begin{enumerate}\item[\textup{(B1)}] $\rho[\tilde u;\tilde\ell](V)=\rho[u^\flat;\ell](V)$ at every $V$.

\item[\textup{(B2)}] $E^\Delta(X,\cdot)=0$ unless $\#X=1$, and $\sup_{\mathcal U_{\{\square\}}(x)}|E^\Delta(\{\square\},\cdot)|\le\delta_\Delta c_a(x)$; moreover $E^\Delta$ depends on $\ell,\tilde\ell,x$ and not on $\tilde u$.

\item[\textup{(B3)}] $|w^\flat(X,V)|\le e^{-\kappa(|X|_c-1)}$ and $|w^\flat-\tilde w|(X,V)\le\frac12\delta_\Delta e^{-\kappa(|X|_c-1)}$ at every $X$ and $V$. Here \[w^\flat=\tilde w\cdot\prod\{e^{q_\square(V)}:\square\text{ off at }V,\ p_*(\square,V)\text{ in }X\}\] (see Steps~3 and~5 of the proof); the product involves only the values of $V$ on the bonds of $X^{+(2r(x)+1)}$, so that $w^\flat(X,\cdot)$ is, just as $\tilde w(X,\cdot)$, a function of $V$ on the bonds of $X^{\approx\approx}=X^{+s(x)}$.

\item[\textup{(B4)}] $\mathsf D^{\rm in}_x(\tilde u,u^\flat)\le\delta_\Delta$.\end{enumerate}\end{lemma}

\begin{theorem}[Two families of rows in one class]\label{Y-I-ThC} Let $x\ge x_9$, let the data \textup{(SD)} at $x$ be given for $(\ell,\ell')$, and suppose that the conditions of item~5 hold; let $\tilde\ell$ be rows on the same $\mathbb T$ with the same grid, of range $\le\varrho(x)$, and let $\delta_\Delta\le1$. Let $\rho_1=e^{\mathsf V_1}\rho[u_1;\ell]\in C^{\rm su}_\flat(x;\ell)$ and $\rho_2=e^{\mathsf V_2}\rho[u_2;\tilde\ell]\in C^{\rm su}_\flat(x;\tilde\ell)$, and let $u_2^\flat$ be the system that \Cref{Y-I-LB} gives for $\tilde u=u_2$.

Then, by \textup{(9a)} to \textup{(9d)}, both images lie in one and the same class $C^{\rm su}_\flat(x';\ell'^{\rm c})$: \[\rho_1'=e^{\mathsf V'_1}\rho[F_\ell(u_1);\ell'^{\rm c}],\qquad\rho_2'=e^{\mathsf V'_2}\rho[F_\ell(u_2^\flat);\ell'^{\rm c}],\] where $\mathsf V'_1=\mathsf V_1+\mathsf V'[u_1]$ and $\mathsf V'_2=\mathsf V_2+\mathsf V'[u_2^\flat]$ are real; and \[\mathsf D_{x'}(F_\ell(u_1),F_\ell(u_2^\flat))\le\vartheta_{10}\bigl(\mathsf D_x(u_1,u_2)+\delta_\Delta\bigr)\le3\vartheta_{10}\le3\vartheta.\]\end{theorem}

In \Cref{Y-I-ThC} the hypotheses of \textup{(SD)} are made on the pair $(\ell,\ell')$ alone, and none of them on $\tilde\ell$.

\begin{corollary}[The level of the switch]\label{Y-I-CorD} Let $\mathsf b\ge x_\flat$, $K'\ge K\ge1$, $m=K'-K$ and $x_k=L^{K-k}\mathsf b$. Let $\ell_k$ and $\tilde\ell_{m+k}$ be the two chains of \Cref{Y-J-Co50}, under its hypotheses, that is, with $L\ge L_B$ odd, $aL\ge3/c_L$, $\theta_D\le\lambda_*/8$, \textup{(X12)} and \textup{(X13)}; and let the side $N$ of the torus on which both chains end be even. In this corollary and in its proof, $N$ denotes this side, so that the coarse side of the step that starts at $x_k$ is $NL^{K-k-1}$ for $0\le k\le K-1$.

Then for $0\le k\le K$ the number $\delta_\Delta$ at $(x_k;\ell_k,\tilde\ell_{m+k})$ is $\le c_4\mathsf b^{1/2}L^{K/2-9k/10}$, and this bound is $\le1$ if and only if $k\ge(5K+10\Lambda)/9$. Moreover $\Lambda>0$, so that $k_0\ge1$; and $k_0\le K-1$ if and only if $K\ge K_0(\mathsf b)$. Further, $K-k_0\ge\frac49K-\frac{10}9\Lambda-1$. With the numeral $\varphi:=5/9$ of \textup{(O14)} in item~1 of \S\ref{Y-J-1}, one has $k_0<\varphi K+\frac{10}9\Lambda+1$, since $k_0=\lceil(5K+10\Lambda)/9\rceil$; hence $\vartheta^{K-k_0}\le\vartheta^{-\frac{10}9\Lambda-1}(\vartheta^{4/9})^K$, and the rate in $K$ of this bound is $\vartheta^{4/9}\le0.7013$ (and $0.6632$ at $\vartheta_{10}$).\end{corollary}

The reference coupling changes by $x'=x/L$ in the exact identity
\textup{(9a)}. Matching the reference couplings of two runs therefore
requires a change of rows only.

\begin{proposition}[An unbounded strong distance after a change of rows]\label{Y-I-InF}
Let $x\ge x_9$, let the five conditions of item~5 hold, and let $\delta_\Delta\in(0,1]$. Then there are rows $\ell$ of range $0$ and, with $\tilde\ell=\emptyset$, a real $\tilde u\in M_1^{\rm su}(x)$ with the strong cap, with the following property for the system $u^\flat$ that \Cref{Y-I-LB} forms from these data: for every $n$ such that $\mathbb T$ has $n$ sites whose pairwise sup distances are $\ge2M(r(x)+2)$, there exist a field $V$ and a set $X$ such that \[|w^\flat-\tilde w|(X,V)\cdot\exp\{c_\Xi x\varepsilon\sum_{D(V)\cap X}\phi_p+\kappa(|X|_c-1)\}\ \ge\ \exp\{n\delta_\Delta c_a(x)/(8c_Q^{\rm su})\}-1.\] Hence, for this system $u^\flat$ of \Cref{Y-I-LB}, $\mathsf D_x(\tilde u,u^\flat)>\delta_\Delta$ as soon as $n\ge8c_Q^{\rm su}/c_a(x)$, and $\sup_N\mathsf D_x(\tilde u,u^\flat)=\infty$.
\end{proposition}

\begin{remark}[Item~7: weaker contractions implied by \Cref{Y-I-ThA}, the caps in \Cref{Y-I-LB}, the constants, the change of rows, empty rows]
\emph{Weaker contractions implied by \Cref{Y-I-ThA}.} A contraction with the strong distance $\mathsf D_x$ on the right, and with a factor $\vartheta<1$ computed from the letters, follows from \Cref{Y-I-ThA}\,(ii), because $\mathsf D^{\rm in}_x\le\mathsf D_x$. The factor $0.45$ comes from \Cref{Y-I-ThA}\,(iv). The stronger form is \textup{(ii)}, and \Cref{Y-I-ThC} requires it.

\emph{The caps in \Cref{Y-I-LB}.} In \Cref{Y-I-LB}, $\tilde u$ has the caps of the class, the strong cap of $\tilde w$ included, since every point of a class has it; and $u^\flat\in M_2^{\rm su}(x)$, with the weak cap for $w^\flat$.

\emph{The constants.}
\begin{enumerate}
\item[\textup{(1)}] The contraction factor of \Cref{Y-I-ThA}\,(ii) is $\vartheta_{10}=0.3968$, because $E'^{\rm new}$ cancels, as it does in \Cref{Y-H-C9}\,(e). The numbers $0.4$ and $0.4318$ are upper bounds of this factor. The letter $\vartheta$ of the table has the value $0.4468$.
\item[\textup{(2)}] One has $c_Q^{\rm su}=306M^3\mathsf g(\lambda_\flat)^2$.
\item[\textup{(3)}] The constant $c_4$ may be replaced by any larger number.
\item[\textup{(4)}] The condition \textup{(F3)} of item~5 gives $\delta'\le\frac12\delta_\Delta$ in Step~6, and the inequality $c_a(5+2r)^3\le\frac12$ is \textup{(X18)} of \textup{(O13)} in item~1 of \S\ref{Y-J-1}. At every $x\ge x_9$ both inequalities follow from \textup{(Y.4)} of item~5 of \S\ref{Y-H-1}, by Step~0\,\textup{(Fl)}.
\item[\textup{(5)}] The upper bound of $c_h$ is $c_h\le\delta_3'/2$, which is \textup{(O11)}. By \textup{(O10)} in item~3 of \S\ref{Y-G0-1} we have $\delta_3'\le\delta_3$, and hence $c_h\le\delta_3/2$, the bound of \textup{(O11)} in that item.
\end{enumerate}

\emph{The change of rows in \Cref{Y-H-C9}.} By item~8 of \S\ref{Y-H-1}, all that \Cref{Y-H-C9} requires of the change of rows is $u^\flat\in M_2^{\rm su}(x)$, and \textup{(B2)}, \textup{(B3)} give it.

\emph{Empty rows.}
\begin{enumerate}
\item[\textup{(1)}] Of the rows $\ell,\tilde\ell$ themselves, \Cref{Y-I-LB} assumes only that their range is $\le\varrho(x)$; its one further hypothesis on the pair is $\delta_\Delta\le1$, and it assumes no lower bound of a row and no membership in a set $\mathcal R(\cdot)$, and it allows the rows to be empty. With $\ell=\emptyset$, with $\tilde\ell=\emptyset$, or with both, Steps~3 to~6 hold without any change. The reasons are the implication $\mathsf H_y=0\Rightarrow\mathsf E_y=0$, the independence of the mask, the clip and the paths from the rows, and the fact that the argument nowhere divides by $\delta$ or calls on a lower bound of a row. When both are empty, $D=0$, $\delta_\Delta=0$, $E^\Delta=0$, all $q_\square=0$, and $u^\flat=\tilde u$. \Cref{Y-I-InF} is the case $\tilde\ell=\emptyset$.
\item[\textup{(2)}] The condition \textup{(SD)} admits $\ell=\emptyset$: one has $\emptyset\in\mathcal R(1,0,\lambda_*,0)\subset\mathcal R(b_*,K_*,\lambda_*,\frac19)$, by the last sentence and the display of \Cref{Y-A-SetsOfRows}. At $\ell=\emptyset$ the pair $(\ell,\ell')$ also has the four conclusions required by \textup{(SD)}; \Cref{Y-A-Fr} at $k=0$ and $x_0=x$ yields the first three of them through \textup{(Fr2)} for the sixteen kernels, \textup{(Fr3$'$)} and \textup{(Fr4)}. Item~\textup{(8)} of that theorem asks for a number $>1$ subject to \textup{(X12)}, and here we take $x'$ for it, not the $x_\flat$ of \textup{(O13)}. Step~H16\,(v) of the proof of \Cref{Y-A-Fr} shows that every number $\ge2$ satisfies \textup{(X12)}, and $x_0=x\ge x_9$ gives $x_1=x'\ge Le\ge2$; so $x'$ is admissible, and $x_1$ is at least $x'$, as the theorem requires at $k=0$. The fourth conclusion, made up of \textup{(L1$^{\rm su}$)} and \textup{(L2$^{\rm su}$)}, holds by \textup{(O4)}. The proofs of Thms.~\ref{Y-I-ThA} and~\ref{Y-I-ThC} draw on no row: at $\ell=\emptyset$ they consist of the same deductions from the four clauses \textup{(9a)}, \textup{(9b)}, \textup{(9c)} and \textup{(9d)}, which, by \S\ref{Y-H-1}, hold for all data in the setting \textup{(SD)}; and the condition \textup{(SD)} of \S\ref{Y-H-1} admits $\ell=\emptyset$ expressly, in item~2 of that section. So at $\ell=\emptyset$ Thms.~\ref{Y-I-ThA} and~\ref{Y-I-ThC} rest on \Cref{Y-H-C9} at $\ell=\emptyset$, for members that are complex, lie within up to sixteen caps and obey the weak cap, and for $x\ge x_9$.
\end{enumerate}
\end{remark}

\begin{proof}[Preliminaries to the proofs of this section]
\emph{Step 0: geometry, the conditions of item~5, and the order of the constants.} We claim the following seven statements, \textup{(G1)} to \textup{(Ord)}, and prove them in the paragraphs after the list.
\begin{enumerate}
\item[\textup{(G1)}] If $y$ is a site of a cube $\square$ and $z$ is a site with $|z-y|_\infty\le jM$, then $z$ lies in a cube of $\square^{+j}$.
\item[\textup{(G2)}] Let \textup{(F2)} hold, and let $y\in\square$ and $|y-p|\le\varrho(x)$. Then all sites of $\Gamma(y,\pi(p))$ and of $p$ lie within $\varrho(x)+1\le Mr(x)-1$ of $y$. Consequently $\Gamma(y,\pi(p))$ is a path in $\square^\approx$ with no more than $3\varrho(x)\le\ell_w(x)$ bonds, and $p$ is a plaquette of $\square^\approx$ whose four bonds are bonds of $\square^\approx$.
\item[\textup{(G3)}] Let \textup{(F2)} hold. Then every path $\Gamma(y,\pi(p))$ with $|y-p|\le\varrho(x)+1$ is defined, and $\square^\approx$ is an embedded box.
\item[\textup{(G4)}] For each cube, at most $(2r(x)+1)^3$ cubes $\square$ have it in $\square^\approx$.
\item[\textup{(G5)}] For every $y$, $\sum_{p,p'}|D_y(p,p')|\le9\mathsf g(\lambda_\flat)^2\delta$.
\item[\textup{(Fl)}] For $x\ge x_9$: $\varepsilon<\frac1{680}$; \textup{(F2)} holds; $(5+2r(x))^3c_a(x)\le x^{-1/4}/(192e^{11})$, so that \textup{(F3)} holds; and $r(x)\ge6$.
\item[\textup{(Ord)}] The definitions of $c_Q^{\rm su},c_4,\Lambda,K_0,k_0$ form no circle.
\end{enumerate}

For \textup{(G1)}: the cubes are the boxes $o+Mn+\{0,\dots,M-1\}^3$, and moving one coordinate by no more than $jM$ moves the index $n$ by no more than $j$. For \textup{(G2)}: the coordinate path runs inside the box spanned by $y$ and $\pi(p)$, so its sites lie within $\varrho(x)$ of $y$; the sites of $p$ are $\pi(p),\ \pi(p)+e_\mu,\ \pi(p)+e_\nu,\ \pi(p)+e_\mu+e_\nu$, which lie within $\varrho(x)+1$ of $y$. By \textup{(F2)} and \textup{(G1)} all these sites lie in cubes of $\square^{+r(x)}$, and the path has at most $3\varrho(x)\le3Mr(x)\le\ell_w(x)$ bonds.

For \textup{(G3)}: by \textup{(Tor)}, the side of $\mathbb T$ is $NL\ge16Mr(x)L\ge48Mr(x)$; so $\varrho(x)+1\le Mr(x)$ is less than half the side, and the box $\square^\approx$, of side $(2r(x)+1)M\le3Mr(x)$, embeds. For \textup{(G4)}: $\square_*\subset\square^{+r}$ holds if and only if $\square\subset\square_*^{+r}$, so the cubes $\square$ in question are among those of $\square_*^{+r}$. For \textup{(G5)}: $|D_y(p,p')|\le\delta e^{-\lambda_\flat(|y-p|+|y-p'|)}$; each base site carries three plaquettes, and the periodic sum is bounded by the sum over $\mathbb Z^3$, so that $\sum_pe^{-\lambda_\flat|y-p|}\le3\mathsf g(\lambda_\flat)$.

For \textup{(Fl)}, first the bound on $\varepsilon$: the condition \textup{(Y.0)(0b)} of item~5 of \S\ref{Y-H-1} contains $17T\varepsilon'\le\frac1{40}$ and $(p_0'/p_0)^4\ge\frac12$, with the number $T=L^{3/2}/(19\mathsf N_{\rm c})>1$, by \textup{(L2$^{\rm su}$)} (Step~0\,(iv) of the proof of \Cref{Y-H-C9}). Therefore $\varepsilon'<\frac1{680}$, and $\varepsilon=(p_0/p_0')L^{-1/2}\varepsilon'\le2^{1/4}L^{-1/2}\varepsilon'<\varepsilon'$.

Next, \textup{(F2)}: $\varrho(x)+2\le c_R\log x+3\le20c_R\log x\le Mr(x)$, because $19c_R\log x\ge19\cdot768>3$ for $x\ge e$.

Next, \textup{(F3)}: each term of the bracket in \textup{(Y.4)} of item~5 of \S\ref{Y-H-1} is $\ge0$; hence \[6((2r_1+1)L+2r(x))^3e^{\kappa+5}\cdot16c_a(x)\cdot e^{(\kappa+4)s'}\le\frac12x^{-1/4}.\] Here the third term of that bracket contains $16(c_a+c_h)(x)$, and the summand $c_h(x)>0$ is dropped. As to the letters: $r_1=\lceil R_2/M\rceil+3r_l+5\ge9$ and $s'=s(x')$ are radii of item~2 of \S\ref{Y-H-1}, where $r_1\ge9$ comes from item~2 of \S\ref{Y-G0-1}; and $\mathsf N_{\rm c}=17(\mathsf N_\sharp+1)$, which enters $T$, is the number of \textup{(O3)} in item~3 of \S\ref{Y-G0-1}. Now $(2r_1+1)L\ge19L\ge5$, $e^{\kappa+5}\ge e^{11}$ and $e^{(\kappa+4)s'}\ge1$, which gives $96e^{11}(5+2r(x))^3c_a(x)\le\frac12x^{-1/4}$; and $192e^{11}\ge1.1495\cdot10^7$.

Last, $r(x)\ge6$: the condition \textup{(Y.7)} of item~5 of \S\ref{Y-H-1} contains $5r_1+2\lceil s(x)/L\rceil+6\le r'$, and $r(x)\ge r(x')=r'$.

For \textup{(Ord)}: $c_Q^{\rm su}$ depends on $M$ and on the pure number $\lambda_\flat$, and it appears in \textup{(O10)}. No floor depends on it, since $\delta_\Delta\le1$ is a hypothesis on the pair of rows, not on $x_\flat$. Further, $c_4$ depends on $c_Q^{\rm su},C_1,E_a$, and $\Lambda,K_0,k_0$ depend on $c_4,L,\mathsf b,K$; all of them appear in \textup{(O14)}.
\end{proof}

\begin{proof}[Proof of \Cref{Y-I-ThA}]
A sum in \textup{(9b)} bounds the term for any $Y$ containing the
cell in that sum. Hence at $A=2$ the analytic factor is
$0.35+2\cdot0.248=0.846$. The other caps and the conditions without
a size also follow from \textup{(9b)}, and the strong cap implies
the weak cap. This proves \textup{(i)}.

Apply Step~12 of the proof of Thm.~\ref{Y-H-C9} with $B=2$ for \textup{(ii)}
and with $B=1$ for \textup{(iii)}. The weak distance is bounded by the
strong one because the strong weight for $w$ is larger.

Under the replacement in \textup{(iv)}, the analytic factor at $A=2$
is at most $1$. In the segment estimate use the full analytic member
in place of $E'^{\rm old}$ and its bound at $A=16$. This gives the
factor $\max\{16\vartheta_w+n_w,1\}/10$. Since
$16\vartheta_w+n_w=8(2\vartheta_w+n_w)-7n_w\le8$, the factor is at
most $0.8$. Substitution gives the two stated numerical values.
\end{proof}

\begin{proof}[Proof of \Cref{Y-I-LB}]
\emph{Step 3: the identity of the energies.} We claim that at every $V$ one has $-x\mathrm{En}_{\tilde\ell}(V)=-x\mathrm{En}_\ell(V)+\sum_\square q_\square(V)$, where $q_\square:=\sum_{y\in\square}q_y$ and \[q_y(V):=-\frac x2\,m_y(V)\sum_{p,p'}D_y(p,p')\langle T(y,\pi(p))\zeta_p,T(y,\pi(p'))\zeta_{p'}\rangle;\] further, that $q_y$ is real, Borel and gauge invariant, that $|q_y(V)|\le\frac92\mathsf g(\lambda_\flat)^2p_0^2\delta$, and that $q_y$ depends on $V$ only at sites within $\max\{\rho_0,\varrho(x)+1\}+1$ of $y$.

In $\mathsf H_y\otimes\mathfrak g$ we have \[|\ell^V_y\zeta|^2=\sum_{p,p'}\langle\ell_y(p),\ell_y(p')\rangle\langle T(y,\pi(p))\zeta_p,T(y,\pi(p'))\zeta_{p'}\rangle=\sum_{p,p'}\mathsf E_y(p,p')\langle T(y,\pi(p))\zeta_p,T(y,\pi(p'))\zeta_{p'}\rangle,\] so the spaces $\mathsf H_y$ enter through $\mathsf E_y$ alone. The sum runs over $|y-p|,|y-p'|\le\varrho(x)$, where the paths are defined, by \textup{(G3)}. The terms $S_\varepsilon(\phi_p)$, the clip $\zeta$, the mask $m_y$ and the paths coincide for $\ell$ and $\tilde\ell$, since by item~2 none of them depends on a row; in particular the second radius of the mask is $\varrho(x)+1$, which is a function of $x$ and not the range of the rows. Subtracting the two energies and multiplying by $-x$ gives the identity.

As to the size: at a $G$-valued $V$ each $T$ is an isometry of $\mathfrak g$; $|\zeta_p|\le\varepsilon$ at every $V$; $m_y\le1$; and $x\varepsilon^2=p_0^2$. Together with \textup{(G5)} this yields $|q_y|\le\frac x2\varepsilon^2\cdot9\mathsf g(\lambda_\flat)^2\delta$. The function $q_y$ is Borel since $\mathsf L$, $\operatorname{ang}$ and $\min$ are Borel. It is invariant by \Cref{Y-I-P6}: under $V\to V^u$ each of the two entries of $\langle\cdot,\cdot\rangle$, of the form $T(y,\pi(p))\zeta_p$, becomes $\operatorname{Ad}_{u(y)}T(y,\pi(p))\zeta_p$, and $\operatorname{Ad}_{u(y)}$ is an isometry of $\mathfrak g$; and $m_y$ is invariant. As to dependence, the mask depends on $\phi_p$ at $|y-p|\le\max\{\rho_0,\varrho(x)+1\}$, the sites of such a $p$ lie at most one unit farther out, and the paths and the $\zeta_p$ stay within $\varrho(x)+1$.

\emph{Step 4: the tame part at complex fields.} For $\mathsf W\in\mathcal U_{\{\square\}}(x)$ let $F_p:=-i\log\partial\mathsf W(p)$, defined by the power series, and \[E^\Delta(\{\square\},\mathsf W):=-\frac x2\sum_{y\in\square}\sum_{p,p'}D_y(p,p')\,(\operatorname{Ad}_{\mathsf W(\Gamma(y,\pi(p)))}F_p)\cdot(\operatorname{Ad}_{\mathsf W(\Gamma(y,\pi(p')))}F_{p'}),\] and let $E^\Delta(X,\cdot):=0$ for $\#X\ge2$. We claim: $E^\Delta$ has the clauses without a size of $E^a$; it is real at $G$-valued fields; $\sup_{\mathcal U_{\{\square\}}(x)}|E^\Delta(\{\square\},\cdot)|\le c_Q^{\rm su}p_0^2\delta=\delta_\Delta c_a(x)$; and $E^\Delta(\{\square\},V)=q_\square(V)$ at every $G$-valued $V$ that is tame on $\square^\approx$.

Dependence: by \textup{(G2)}, every path and every plaquette that occurs lies in $\square^\approx$ and keeps a margin of one site there. Holomorphy: by \textup{(u3)} and \textup{(F1)}, $|\partial\mathsf W(p)-1|<2\varepsilon\le\frac1{20}$; hence the series $\log(1+A)=\sum_{k\ge1}(-1)^{k+1}A^k/k$ converges, and $|F_p|\le-\log(1-2\varepsilon)\le2.0518\,\varepsilon\le2.06\,\varepsilon$, because $t\mapsto-\log(1-t)/t$ is increasing and its value at $t=\frac1{20}$ is $\le1.02587$. The entries of $\mathsf W(\bar b)=\mathsf W(b)^{-1}$ are those of the adjugate. Thus $E^\Delta$ is built, by finitely many sums and products, from polynomials in the entries and from convergent power series, so it is holomorphic on $\mathcal U_{\{\square\}}(x)$.

As to the size: by \textup{(G2)} and \textup{(u2)}, $|\mathsf W(\Gamma)|<\mathsf b(x)<2$, where the empty path gives $1$; so $|\operatorname{Ad}_{\mathsf W(\Gamma)}F_p|\le|\mathsf W(\Gamma)|^2|F_p|\le4\cdot2.06\,\varepsilon=8.24\,\varepsilon$, and $|Y\cdot Y'|\le|Y||Y'|$. A single term is therefore at most $\frac x2|D_y(p,p')|(8.24\,\varepsilon)^2=33.9488\,p_0^2|D_y(p,p')|$. By \textup{(G5)}, and since $\square$ has $M^3$ sites, the whole sum is at most $33.9488\cdot9\cdot M^3\mathsf g(\lambda_\flat)^2p_0^2\delta\le305.54\,M^3\mathsf g(\lambda_\flat)^2p_0^2\delta\le c_Q^{\rm su}p_0^2\delta$. At $\#X=1$ the weight $e^{\kappa(\#X-1)}$ is $1$.

Invariance under a $G$-valued $u$: $\partial\mathsf W^u(p)=u(\pi(p))\partial\mathsf W(p)u(\pi(p))^{-1}$, so $F_p\to\operatorname{Ad}_{u(\pi(p))}F_p$; also $\mathsf W^u(\Gamma(y,z))=u(y)\mathsf W(\Gamma(y,z))u(z)^{-1}$, so $\operatorname{Ad}_{\mathsf W(\Gamma)}F_p\to\operatorname{Ad}_{u(y)}\operatorname{Ad}_{\mathsf W(\Gamma)}F_p$; and the form $\cdot$ is $\operatorname{Ad}$-invariant. At a pure gauge, $\partial\mathsf W=1$ and $F=0$.

Reality at the $G$-valued fields of the domain: write $\partial V(p)=e^{iY}$ with $Y=\mathsf L(\partial V(p))$, so that $|Y|=\phi_p$ and $2\sin(\phi_p/2)<\frac1{20}$. The series gives $\log e^{iY}=iY$, hence $F_p=\mathsf L(\partial V(p))\in\mathfrak g$. For $g\in G$ we have $\operatorname{Ad}_gY\in\mathfrak g$ for $Y\in\mathfrak g$, and on $\mathfrak g$ one has $\cdot=\langle,\rangle$.

Agreement with $q_\square$: let $V$ be tame on $\square^\approx$; its restriction then lies in $\mathcal U_{\{\square\}}(x)$, by item~3. Take $y\in\square$. A plaquette with $|y-p|\le\rho_0\le M$, by \textup{(M5r)}, lies in $\square^{+1}$, which is inside $\square^\approx$ because $r(x)\ge1$; so it does not belong to $D(V)$. A plaquette with $|y-p|\le\varrho(x)+1\le Mr(x)$ lies in $\square^\approx$, by \textup{(G1)}, and so $\phi_p\le\varepsilon\le1$. Hence $m_y(V)=1$. For $|y-p|\le\varrho(x)$ we have $\phi_p\le\varepsilon$, so $\zeta_p=F_p$, and the two formulas agree term by term.

\emph{Step 5: the part that is off and \textup{(B1)}.} A cube $\square$ is called \emph{on} at $V$ if $1_{\square^\approx}(V)=1$, and \emph{off} otherwise. For an off $\square$, let $p_*(\square,V)$ be the first plaquette of $D(V)$ lying in $\square^\approx$, with respect to an ordering of the plaquettes of $\mathbb T$ fixed once and for all. We set \[w^\flat(X,V):=\tilde w(X,V)\cdot\prod\{e^{q_\square(V)}:\ \square\text{ off at }V,\ p_*(\square,V)\text{ in }X\}.\] We claim: $\Xi_{w^\flat}(V)=\Xi_{\tilde w}(V)\cdot\prod_{\square\text{ off}}e^{q_\square(V)}$ at every $V$; \textup{(B1)} holds; and $w^\flat(X,\cdot)$ is real, Borel and gauge invariant. Moreover the product $\prod\{e^{q_\square(V)}:\square\text{ off at }V,\ p_*(\square,V)\text{ in }X\}$ depends only on the restriction of $V$ to the bonds of $X^{+(2r(x)+1)}$, a set inside $X^{\approx\approx}$; therefore $w^\flat(X,\cdot)$ depends, as $\tilde w(X,\cdot)$ does, only on the restriction of $V$ to the bonds of $X^{\approx\approx}$.

The factorisation: in a family of \textup{(c$\Xi$)}, the sets $X_i$ are pairwise more than $s(x)\ge8$ apart, hence disjoint, and $D(V)$ is contained in their union. Each $p_*(\square,V)$, being a plaquette of $D(V)$, therefore belongs to exactly one $X_i$. Thus $\prod_iw^\flat(X_i,V)=\prod_i\tilde w(X_i,V)\cdot\prod_{\square\text{ off}}e^{q_\square(V)}$, where the last factor does not change from one family to another, and we sum over the families. If $D(V)=\emptyset$, all cubes are on and both $\Xi$ equal $1$.

For \textup{(B1)}: by Steps~3 and~4, $-x\mathrm{En}_{\tilde\ell}=-x\mathrm{En}_\ell+\sum_{\rm on}q_\square+\sum_{\rm off}q_\square$, and $\sum_{\rm on}q_\square(V)=\sum_\square1_{\square^\approx}(V)E^\Delta(\{\square\},V)=\sum_{X\in\mathcal D}1_{X^\approx}(V)E^\Delta(X,V)$. Hence \[\rho[\tilde u;\tilde\ell]=\exp\{-x\mathrm{En}_\ell+\sum_X1_{X^\approx}[\tilde E^a+E^\Delta+\tilde E^h]\}\cdot\prod_{\rm off}e^{q_\square}\cdot\Xi_{\tilde w}=\rho[u^\flat;\ell].\]

The dependence of the product: if $p_*(\square,V)$ is in $X$, the cube of $p_*$ lies both in $X$ and in $\square^\approx$, and so $\square\subset X^{+r}$. For the cubes $\square$ of $X^{+r}$, whether $\square$ is off, and which plaquette is $p_*(\square,V)$, is decided by $\phi_p$ at the plaquettes in $\square^\approx$; these lie inside $X^{+2r}$, and the end points of their bonds lie in $X^{+(2r+1)}$. By Step~3, \textup{(M5r)} and \textup{(F2)}, $q_\square$ depends on $V$ at sites within $\max\{\rho_0,\varrho(x)+1\}+1\le\max\{M+1,Mr\}$ of $\square$, hence, by \textup{(G1)}, in $\square^{+\max\{2,r\}}$, which lies inside $X^{+(2r+1)}$ because $\max\{2,r\}\le r+1$. Moreover $2r+1\le3r<8r=s(x)$. Finally, $\tilde w(X,\cdot)$ depends on $V$ through the bonds of $X^{\approx\approx}$, by its clause without a size, and it may depend on every one of them; so $w^\flat(X,\cdot)$, like $\tilde w(X,\cdot)$, is a function of $V$ on the bonds of $X^{\approx\approx}$.

For the remaining clauses: the product is finite and is taken over a finite Borel partition of the fields according to $D(V)$, so $w^\flat(X,\cdot)$ is Borel; it is real because $\tilde w$ and the $q_\square$ are real; and it is gauge invariant because $D(V)$, the property of being off, $p_*$, $q_y$ and $\tilde w$ are.

\emph{Step 6: the two bounds on $w^\flat$, \textup{(B2)} to \textup{(B4)}, and $u^\flat\in M_2^{\rm su}(x)$.} We claim that $|w^\flat|\le e^{-\kappa(|X|_c-1)}$ and $|w^\flat-\tilde w|\le\frac2e\,\delta'e^{-\kappa(|X|_c-1)}$, where \[\delta':=(2r(x)+1)^3M^3\,\frac92\,\mathsf g(\lambda_\flat)^2p_0^2\delta=\bigl((2r(x)+1)^3c_a(x)/68\bigr)\delta_\Delta\le\frac12\delta_\Delta.\]

First, fix a plaquette $p_*$ of $D(V)$ and let $\square_*$ be the cube that holds it. Each cube $\square$ with $p_*(\square,V)=p_*$ has $\square_*$ in $\square^\approx$; by \textup{(G4)} there are at most $(2r+1)^3$ such cubes, each with $M^3$ sites. By Step~3, $\sum\{|q_\square(V)|:p_*(\square,V)=p_*\}\le\delta'$. The second form of $\delta'$ follows from $c_Q^{\rm su}p_0^2\delta=\delta_\Delta c_a(x)$ and $4.5/306=1/68$; and \textup{(F3)} together with $\delta_\Delta\le1$ gives $\delta'\le\frac12\delta_\Delta\le\frac12$.

Let $n:=\#(D(V)\cap X)$ and $Q:=\sum\{q_\square(V):\square\text{ off},\ p_*(\square,V)\text{ in }X\}$; then $|Q|\le\delta'n$. If $n=0$, the product is empty, $w^\flat=\tilde w$, and $|\tilde w|\le e^{-\kappa(|X|_c-1)}$. If $n\ge1$, then $\phi_p>\varepsilon$ on $D(V)$, so $x\varepsilon\phi_p>p_0^2$ there, and the strong cap of $\tilde w$ together with \textup{(X15)} gives $|\tilde w(X,V)|\le e^{-c_\Xi p_0^2n}e^{-\kappa(|X|_c-1)}\le e^{-n}e^{-\kappa(|X|_c-1)}$. Then $|w^\flat|=|\tilde w|e^Q\le e^{-n+n/2}e^{-\kappa(|X|_c-1)}$, and \[|w^\flat-\tilde w|=|\tilde w||e^Q-1|\le|\tilde w|\cdot\delta'n\cdot e^{\delta'n}\le\delta'\cdot ne^{-n/2}\cdot e^{-\kappa(|X|_c-1)}\le\frac2e\,\delta'e^{-\kappa(|X|_c-1)}.\] Here we use $|e^Q-1|\le|Q|e^{|Q|}$, valid for either sign of $Q$, and $te^{-t/2}\le2/e$.

Hence \textup{(B3)}. \textup{(B2)} is Step~4. For \textup{(B4)}: $u^\flat-\tilde u=(E^\Delta,0,w^\flat-\tilde w)$, and the three numbers of item~3 whose largest is $\|u^\flat-\tilde u\|_x$ are at most $\delta_\Delta$, $0$ and $\frac1e\delta_\Delta$. Finally, $u^\flat\in M_2^{\rm su}(x)$: the member $\tilde E^a+E^\Delta$ lies within $1+\delta_\Delta\le2$ caps, $\tilde E^h$ within one cap, and $w^\flat$ within one weak cap; the clauses without a size hold by Steps~4 and~5 together with linearity; and reality holds by Steps~4 and~5.
\end{proof}

\begin{remark}[On Step~6]
In the setting of the table and at every $x\ge x_9$, $\delta'\le1.3\cdot10^{-9}\delta_\Delta$, by Step~0.
\end{remark}

\begin{proof}[Proof of \Cref{Y-I-ThC}]
\emph{Step 7: the step applied to the two densities.} We prove the two formulas for $\rho_1'$ and $\rho_2'$, that both of them lie in $C^{\rm su}_\flat(x';\ell'^{\rm c})$, and the bound on $\mathsf D_{x'}$. The definition of the class makes $u_2$ real, in $M_1^{\rm su}(x)$, and subject to the strong cap. So \Cref{Y-I-LB} applies and gives $\rho_2=e^{\mathsf V_2}\rho[u_2^\flat;\ell]$ as functions on the fields of $\mathbb T$, with $u_2^\flat$ real and in $M_2^{\rm su}(x)$. Since the step is linear and independent of the rows (item~2), \textup{(9a)} at $A=2$ gives \[\rho_2'=e^{\mathsf V_2}(\rho[u_2^\flat;\ell])'=e^{\mathsf V_2+\mathsf V'[u_2^\flat]}\rho[F_\ell(u_2^\flat);\ell'^{\rm c}],\] and \textup{(9a)} at $A=1$ gives $\rho_1'$ in the same way.

The class: by \Cref{Y-I-ThA}\,(i), $F_\ell(u_1)$ and $F_\ell(u_2^\flat)$ lie in $M_1^{\rm su}(x')$ with the strong cap, and by \textup{(9d)} they are real, as are $\mathsf V'[u_1]$ and $\mathsf V'[u_2^\flat]$. Since $\rho[u_1;\ell]=e^{-\mathsf V_1}\rho_1>0$ and $\rho[u_2^\flat;\ell]=e^{-\mathsf V_2}\rho_2>0$, the images are $>0$ by \textup{(9d)}; and $\ell'^{\rm c}$ has range $\le\varrho(x')$.

The distance: $\mathsf D^{\rm in}_x$ is the distance of a norm, so item~3 and \textup{(B4)} give \[\mathsf D^{\rm in}_x(u_1,u_2^\flat)\le\mathsf D^{\rm in}_x(u_1,u_2)+\mathsf D^{\rm in}_x(u_2,u_2^\flat)\le\mathsf D_x(u_1,u_2)+\delta_\Delta.\] By item~3, $\mathsf D_x(u_1,u_2)\le2$; and \Cref{Y-I-ThA}\,(ii) is applied to $(u_1,u_2^\flat)$, both of which lie in $M_2^{\rm su}(x)$.
\end{proof}

\begin{proof}[Proof of \Cref{Y-I-CorD}]
\emph{Step 8: the level of the switch.} We claim that $\delta_\Delta(k)\le c_4\mathsf b^{1/2}L^{K/2-9k/10}$, that $\Lambda>0$, that $k_0\ge1$, that $k_0\le K-1$ if and only if $K\ge\frac52\Lambda+\frac94$, and that $K-k_0\ge\frac49K-\frac{10}9\Lambda-1$.

\Cref{Y-J-Co50} gives $\delta_k:=\|\tilde{\mathsf E}_{m+k}-\mathsf E_k\|_{(\lambda_\flat)}\le15C_1^2\vartheta_1^k\le15C_1^2L^{-2k/5}$, in the norm of \Cref{Y-I-LB}; by Step~1 of the proof of \Cref{Y-J-ThA}, the hypotheses of that corollary hold for the letters of the table. The last inequality holds because $\vartheta_1=c_{\rm fix}L^{-1/2}\le L^{1/10}L^{-1/2}$, as $L\ge L_B\ge c_{\rm fix}^{10}$ by \textup{(O0)} and \textup{(O4)}. Both chains, at partner levels, live on a single torus and have range $\le\varrho(x_k)$. The inequality $\theta_D\le\lambda_*/8$ follows from the definitions alone: $\theta_D\le\lambda_2^*/2$ by \textup{(O6)}, and $\lambda_2^*=\min\{\lambda_*/4,[6K_*^2\mathsf g(\lambda_*/2)\mathsf g_1(\lambda_*/4)]^{-1}\}\le\lambda_*/4$, by the order of the constants, \textup{(0)}, in \S\ref{Y-A-1}.

By \textup{(X14)}, \[\delta_\Delta(k)=(c_Q^{\rm su}/E_a)x_k^{1/2}\delta_k/p_0(x_k)\le(c_Q^{\rm su}/E_a)\mathsf b^{1/2}L^{(K-k)/2}\cdot15C_1^2L^{-2k/5}=c_4\mathsf b^{1/2}L^{K/2-9k/10},\] and this is $\le1$ if and only if $\frac9{10}k\ge\frac K2+\Lambda$. As to $\Lambda>0$: $16c_a(\mathsf b)\le\frac12$ and \textup{(X14)} give $\mathsf b^{1/2}\ge32E_ap_0(\mathsf b)^3\ge32E_a$, hence $c_4\mathsf b^{1/2}\ge480c_Q^{\rm su}C_1^2>1$, since $c_Q^{\rm su}\ge306$ and $C_1>1$. Put $y:=(5K+10\Lambda)/9>0$. Then $k_0=\lceil y\rceil\ge1$; and since $K-1$ is an integer, $\lceil y\rceil\le K-1$ holds if and only if $y\le K-1$, that is, if and only if $4K\ge10\Lambda+9$. Also $K-k_0\ge K-y-1$.

At $k=k_0\le K-1$ we have $x_{k_0}\ge L\mathsf b\ge Lx_\flat\ge x_9$. So the floor of \Cref{Y-I-ThC} is met there, and, by Step~0\,\textup{(Fl)}, so are the conditions \textup{(F1)}, \textup{(F2)}, \textup{(F3)}; \textup{(M5r)} and \textup{(X15)} are entries of the table; and $\delta_\Delta\le1$ is the bound just obtained. The remaining hypotheses of \Cref{Y-I-ThC} are established elsewhere: \textup{(SD)} at $(x_{k_0},NL^{K-k_0-1},\ell_{k_0})$, for $N$ even with $M\mid N$ and $N\ge2Ms(L\mathsf b)$, by Step~2\,(d) of the proof of \Cref{Y-J-ThA}; and the membership of the two densities in their classes by item~1 of Step~4 there. With these, \Cref{Y-I-ThC} applies at $k_0$. Numerically, $0.45^{4/9}\le0.70125$, $0.4468^{4/9}\le0.69903$ and $0.3968^{4/9}\le0.66312$.
\end{proof}

\begin{proof}[Proof of \Cref{Y-I-InF}]
\emph{Step 10: the lower bound.} The letters and $x\ge x_9$ are fixed. By Step~0\,\textup{(Fl)}, $r(x)\ge6$ and $\varepsilon\le\frac1{40}$; and $M\ge3$ by \textup{(M2)}, the condition $M\ge1920(\kappa+4)c_R$ of \textup{(O7)}.

The rows: $\tilde\ell=\emptyset$; $\mathsf H_y=\mathbb R$, and $\ell_y(p)=t$ if $p=p_y:=(y;12)$, while $\ell_y(p)=0$ otherwise, so $\ell$ has range $0$. Then $D_y(p,p')=-t^2$ at $p=p'=p_y$ and $D_y(p,p')=0$ otherwise, so $\delta=t^2$; we choose $t^2:=\delta_\Delta c_a(x)/(c_Q^{\rm su}p_0^2)$. Since $T(y,y)=1$, $q_y(V)=\frac x2t^2m_y(V)|\zeta_{p_y}|^2\ge0$ at every $V$.

The members: $\tilde u=(0,0,\tilde w)$ with $\tilde w(X,V):=\exp\{-c_\Xi x\varepsilon\sum_{D(V)\cap X}\phi_p-\kappa(|X|_c-1)\}$. This function is real, Borel and gauge invariant, depends on $\phi$ on $X$ only, and meets the strong cap with equality; moreover $\rho[\tilde u;\emptyset]>0$.

The field: take sites $z_1,\dots,z_n$ whose pairwise sup distances are $\ge2M(r(x)+2)$; put $d:=\max\{\rho_0,\varrho(x)+1\}+3$ and $y_i:=z_i+de_1$; and let $V:=1$ except for $V(z_i,3):=e^{2i\varepsilon\sigma_3}$ and $V(y_i,1):=e^{i(\varepsilon/2)\sigma_3}$. The plaquettes containing the bond $(z_i,3)$ are $(z_i;13),\ (z_i-e_1;13),\ (z_i;23),\ (z_i-e_2;23)$; they have $\phi=2\varepsilon$, which is $\le1$, and they are rough. The plaquettes containing $(y_i,1)$ are $(y_i;12),\ (y_i-e_2;12),\ (y_i;13),\ (y_i-e_3;13)$; they have $\phi=\varepsilon/2$ and are tame. Because $d\ge4$, no plaquette contains two of the changed bonds, and every other plaquette has $\phi=0$. Moreover $d\le Mr(x)$, as $\rho_0+3\le M+3\le2M$ and $\varrho(x)+4\le c_R\log x+5\le20c_R\log x$; so, by \textup{(G1)}, the cube of $z_i$ is one of the cubes of $\square^{+r(x)}$ for the cube $\square$ of $y_i$, and that cube $\square$ is therefore off. The four rough plaquettes that contain the bond $(z_i,3)$ satisfy $|y_i-p|\ge d>\max\{\rho_0,\varrho(x)+1\}$, and the rough plaquettes that contain a bond $(z_j,3)$ with $j\ne i$ satisfy $|y_i-p|\ge2M(r+2)-d-1>Mr$; hence $m_{y_i}(V)=1$. The plaquette $p_{y_i}$ contains $(y_i,1)$, so $|\zeta_{p_{y_i}}|=\varepsilon/2$ and $q_{y_i}(V)=t^2p_0^2/8$.

The set: let $X$ consist of the cubes that hold a plaquette of $D(V)$. Each off cube has its $p_*$ in $D(V)$ and hence in $X$; $q_y\ge0$ for every $y$; and the $y_i$ are distinct sites of off cubes. Therefore $w^\flat(X,V)/\tilde w(X,V)=e^Q$ with $Q\ge nt^2p_0^2/8=n\delta_\Delta c_a(x)/(8c_Q^{\rm su})$. Since $\tilde w$ equals the cap, the weighted difference is $e^Q-1$. For $n\ge8c_Q^{\rm su}/c_a(x)$ this gives $Q\ge\delta_\Delta$ and $e^Q-1>\delta_\Delta$. Sites $z_1,\dots,z_n$ whose pairwise sup distances are $\ge2M(r(x)+2)$ exist on $\mathbb T$ as soon as $\lfloor NL/(2M(r(x)+2))\rfloor^3\ge n$. The conditions of item~2 on $N$ are $N$ even, $M\mid N$ and \textup{(Tor)} $N\ge2M\,s(x)$. None of them bounds $N$ from above.
\end{proof}

\subsection{Comparison of two cutoffs}\label{Y-J-1}

Two runs with the same final torus and label have the same
coupling at partner levels. Their site forms approach each other until
one run can be expressed over the rows of the other. The contraction
then applies at every subsequent step. A lower bound for the
normalizing integral converts the resulting comparison into total
variation at every fixed depth. We also apply the comparison to the
response to a fixed abelian flux.

\emph{Item 0: setting, what is used, signs.} Our group is $\operatorname{SU}(2)$, the dimension is $d=3$, and the step is the drawn step. We rely on \Cref{Y-GB-ThB}; on clauses (a), (b), (d), (e) of \Cref{Y-H-C9}, together with the data \textup{(SD)} of item~2 of \S\ref{Y-H-1} and with item~3 of \S\ref{Y-H-1}; and on \Cref{Y-I-LB}.

The label $\mathsf b$ is the coupling at the last level, and $G=\operatorname{SU}(2)$. For laws with densities $f,g$ against Haar measure we define $\|\mu-\nu\|_{\rm TV}:=\int|f-g|\,dV$, without a factor $\frac12$. With an argument, $s(x)=8r(x)$ is a radius; without an argument, $s$ denotes the depth $K-k$; so in $NL^s\ge2M\,s(L^s\mathsf b)$ the exponents are the depth while $s(\cdot)$ is the radius; and in $k_s$ and $z(s)$ the letter $s$ is the noise, which is a third meaning. The number of \textup{(O11)} is $c_h$, and the cap is $c_h(x)$; next to $c_a$ and without an argument, as in $(c_a+c_h)$, the letter $c_h$ denotes the cap at the coupling in question. Laws are written $\mu,\nu$, but $\mu$ in Step~9\,(4) is Haar measure of $G$. By $K_{{\rm q},k}$ and $K_v^\flat$ we mean the $K_{\rm q}$ and $K_v$ of \textup{(O10)} in item~3 of \S\ref{Y-G0-1}. A system of members is written $u$; in $V^u$, $u(y)$, $u'$, $\operatorname{Ad}_{u(y)}$ the letter stands for a gauge transformation. The coupling of the level is $x$, and $x'=x/L$. In item~2, in the paragraph of the definition ``The step'' that defines $\mathsf M$ and in the proof of the three clauses, and only there, $x$ and $x'=x+Le_\mu$ are fine sites and $y(x)$ is a coarse site, as in item~C.I~4 of \S\ref{Y-C-1}; in $k_{ax}$, and where we state on what the step depends, $x$ is the coupling. In \Cref{Y-J-CorD} and in Step~12, $\sigma_3={\rm diag}(1,-1)$. Finally, \textup{(9a)}, \textup{(9b)}, \textup{(9d)}, \textup{(9e)} stand for clauses (a), (b), (d), (e) of \Cref{Y-H-C9}.

A step that we cite without naming a statement belongs to a proof of this section: Steps 0 to 6 to the proof of \Cref{Y-J-ThA}, Steps 7 to 10 to that of \Cref{Y-J-ThB}, Step~11 to that of \Cref{Y-J-CorC}, and Step~12 to that of \Cref{Y-J-CorD}.

\emph{Item 1: letters.} We go through the entries \textup{(O0)} to \textup{(O13)} of the one table, item~3 of \S\ref{Y-G0-1}, and the entry \textup{(O14)}, in their order.

\begin{enumerate} \item[\textup{(O0)}] Pure numbers. \Cref{Y-A-PureNumbers} provides $b_*$, $\gamma_*$, $\lambda_*=\gamma_*^{1/2}/44$ ($\le1$), $K_*=3\gamma_*^{-1/2}$, $L_A$ and $\mathsf t_*(\varrho)=3K_*\mathsf g(\lambda_*/2)e^{-\lambda_*\varrho/2}$. The numbers $c_1^*,c_2^*,\lambda_2^*,\theta_J$ are those of \Cref{Y-A-Fr}, ``The order of the constants''\,(0); $c_1^*,c_2^*,\lambda_2^*$ also appear in \Cref{Y-A-Stars}, and $\theta_J$ in \Cref{Y-A-Jinv}. Here $\mathsf g(\lambda)=\sum_{n\in\mathbb Z^3}e^{-\lambda|n|_\infty}$. We set $C_1=\sqrt3\,K_*\,\mathsf g(\lambda_*/2)^{3/2}$, $c_0:=108\,b_*\gamma_*^{-1/2}\mathsf g(\lambda_*/4)^3$, $c_{\rm fix}:=48\,\mathsf g(\lambda_\flat)C_1c_0^{1/2}$ and $L_B=\max\{L_A,c_{\rm fix}^{10}\}$, where $\lambda_\flat=\lambda_*/16$; and $\vartheta_1:=c_{\rm fix}L^{-1/2}$, a number that depends on $L$.

\item[\textup{(O1)}] $d_0\in(0,\frac12)$.

\item[\textup{(O2)}] $0<c_\Xi\le\frac12(\frac12-d_0)$.

\item[\textup{(O3)}] The radius $\rho_0$ is given by \eqref{eq:row-mask-radius} of \S\ref{Y-A-1}, and $\mathsf N_Z=1+3K_*^2\mathsf g(\lambda_*)^2$.

\item[\textup{(O4)}] $L$ is odd and $L\ge L_*^{\rm su}$, where $L_*^{\rm su}$ is the least integer $\ge\max\{L_A,L_B\}$ such that \textup{(L1$^{\rm su}$)}, \textup{(L1$^\flat$)} with right side $\frac3{16}d_0$ and \textup{(L2$^{\rm su}$)} hold at that integer and at every larger $L$ (\textup{(O4)} in item~3 of \S\ref{Y-G0-1}). The number $c_L=(L^2-1)/(6L^2)$ depends on $L$, and $a\ge3/(c_LL)$.

\item[\textup{(O5)}] $\kappa\ge6$.

\item[\textup{(O6)}] This entry fixes $\gamma=(73L^4)^{-1}$, $A^\flat$, $\theta^\flat$, $\theta_{\mathbb C}=\theta^\flat/6$ and $\vartheta_0^\flat=\theta^\flat/24$. In it, $A^\flat=2e^2(16+a)+384e^8K_*^2\mathsf g_2(\lambda_*/2)^2$, $\theta^\flat=\min\{\theta_1,\gamma/(2A^\flat)\}$ and $\theta_1=\min\{1,\lambda_*L/12\}$, taken from item~C.I~1.2 of \S\ref{Y-C-1}, at $(K,\lambda)=(K_*,\lambda_*)$, and from \Cref{Y-A-Fr}, ``The order of the constants''\,(6).

We further put $\theta_D=\frac12\min\{\lambda_2^*,\theta_{\mathbb C},\theta_J\}$, and we take $C_D\ge4$ and no smaller than the number $C_D$ of \Cref{Y-A-Fr}, ``The order of the constants''\,(6), defined with this $\theta_D$. This choice is legitimate: \textup{(Dec)} of \S\ref{Y-A-1} bounds sixteen kernels from above by $C_D$, so it stays true for every larger $C_D$; and a number of this kind is what Sections~\ref{Y-B-1}, \ref{Y-G-1} and \ref{Y-H-1} mean by $C_D$. Then $c_R=64/\theta_D$.

We have $\theta_D\le\lambda_*/8$, which follows from the choice of the letters (Step~0\,(i)). Hence, as $c_R=64/\theta_D$ and $\lambda_*\le1$, we get $e_R:=\lambda_*c_R/2\ge256$ and $c_R\ge512$.

\item[\textup{(O7)}] $M=L^{m'}$ with $m'\ge1$, subject to three conditions.

\begin{enumerate} \item[\textup{(M2)}] $M\ge1920(\kappa+4)c_R$.

\item[\textup{(M4$^\flat$)}] The condition of this name in \textup{(O7)} of item~3 of \S\ref{Y-G0-1}.

\item[\textup{(M5r)}] $M\ge\rho_0$. \end{enumerate}

The condition \textup{(M4$^\sharp$)} is a consequence of \textup{(M4$^\flat$)}.

\item[\textup{(O8)}] $\mathfrak p_0>2$ and $A_0>0$; moreover $p_0(x)=A_0(\log x)^{\mathfrak p_0}$ and $\varepsilon=\varepsilon(x)=p_0(x)x^{-1/2}$.

\item[\textup{(O9), (O10)}] These two entries contain the numbers of the law of windows and those of \Cref{Y-B-C1} and of Sections~\ref{Y-C-1}, \ref{Y-D-1}, \ref{Y-E-1}, \ref{Y-F-1}. Four of these numbers are $\delta_3,\delta_3',K_{{\rm q},k},K_v^\flat$. From \S\ref{Y-I-1} we take $c_Q^{\rm su}:=306\,M^3\mathsf g(\lambda_\flat)^2$.

\item[\textup{(O11)}] $0<c_h\le\delta_3'/2$.

\item[\textup{(O12)}] $E_a\ge\max\{32K_{{\rm q},k},\,1500K_v^\flat L^{-1/2}\}$. The caps are $c_a(x)=E_ap_0(x)^3x^{-1/2}$ and $c_h(x)=e^{-c_hp_0(x)^2}$.

\item[\textup{(O13)}] This entry is the last of the choices. It is a number $x_\flat\ge e^2$ such that $x_\flat\ge x_\flat^{(1)}$, $x_\flat\ge x_\flat^0$ and $L\,x_\flat\ge x_9$, such that \textup{(X12)} and \textup{(X13)} hold, and such that \textup{(X14)} to \textup{(X18)} hold for all $x\ge x_\flat$.

Here $x_\flat^{(1)}$ is the number of \Cref{Y-B-C1}; it coincides with the $x_\flat^0$ of \textup{(O13)} in item~3 of \S\ref{Y-G0-1} evaluated at $c_h=\delta_3/2$, $E_a=\max\{32L^{-1/2}K_v^\flat,32K_{{\rm q},k}\}$. Next, $x_\flat^0$ is the number of \textup{(O13)} in item~3 of \S\ref{Y-G0-1}, computed from the $c_h,E_a$ of \textup{(O11)}, \textup{(O12)}; and $x_9$ is the number of \textup{(O13)} in item~1 of \S\ref{Y-H-1}.

We recall the definitions of three numbers, in this order: $x_{\rm fl}$ (item~G of \S\ref{Y-F-1}), $x_\flat^0=x_\flat^0(c_h,E_a)$ and $x_9$. First, $x_{\rm fl}:=1+\inf S_{\rm fl}$, where $S_{\rm fl}$ is the set of the $y\ge L^2e$ for which every floor of item~G (Floors) of \S\ref{Y-F-1} is valid at each $x\ge y$; here $e=\exp1$. Second, $x_\flat^0=x_\flat^0(c_h,E_a):=1+\inf S(c_h,E_a)$, where $S(c_h,E_a)$ consists of the $y\ge e^2$ for which \textup{(A0)} to \textup{(A11)} of item~4 of \S\ref{Y-G0-1} are valid at each $x\ge Ly$. Third, $x_9:=1+\inf S_9$, where $S_9$ is the set of the $y\ge x_{\rm fl}$ for which \textup{(Y.0)} to \textup{(Y.9)} of item~5 of \S\ref{Y-H-1} are valid at each $x\ge y$.

Each of these sets is an upper half line: together with $y$ it contains every $\tilde y\ge y$. None of them is empty, because each floor holds for all large $x$, as Step~13 of the proof of \Cref{Y-F-C7}, Step~14 of the proof of \Cref{Y-G-ThA} and Step~0\,(i) of the proof of \Cref{Y-H-C9} show. A set of this kind need not contain its infimum, since the inequalities of the floors involve numbers rounded up to an integer, $\lceil\cdot\rceil$, and these, as functions of $x$, are continuous from the left and jump upwards to the right. The number $1+\inf$ does lie in the set, which is an upper half line, and it depends on what the set depends on.

The definitions of $S_9$ and $S_{\rm fl}$ give $x_9\ge x_{\rm fl}\ge L^2e$.

The conditions named in this entry read as follows.

\begin{enumerate} \item[\textup{(X12)}] $3K_*\mathsf g(\lambda_*/2)\,x_\flat^{-e_R}/(1-L^{-e_R})\le\frac19$.

\item[\textup{(X13)}] $16\,x_\flat^{-e_R/4}\le1$.

\item[\textup{(X14)}] $p_0(x)\ge1$.

\item[\textup{(X15)}] $c_\Xi p_0(x)^2\ge1$.

\item[\textup{(X16)}] $\varepsilon(x)\le\frac1{40}$.

\item[\textup{(X17)}] $16(c_a+c_h)(x)\le\frac12$.

\item[\textup{(X18)}] $c_a(x)(5+2r(x))^3\le\frac12$. \end{enumerate}

The condition \textup{(X13)} follows from the choice of the letters and imposes nothing more: $x_\flat\ge e^2$ and $e_R\ge256$ give $16x_\flat^{-e_R/4}\le16e^{-128}$, which is less than $1$. The condition \textup{(X12)} is a consequence of that choice as well: $3K_*\mathsf g(\lambda_*/2)<4\cdot10^{56}$, $x_\flat^{-e_R}\le e^{-512}<10^{-222}$ and $1-L^{-e_R}\ge\frac12$, and these three bounds give \textup{(X12)}.

The bound $4\cdot10^{56}$ rests on $\mathsf g(\lambda)\le51\lambda^{-3}$ for $0<\lambda\le1$, one of the bounds after the display of $\mathsf g_1,\mathsf g_2$ in \Cref{Y-A-Sums}, and on $K_*\le3.1527\cdot10^{12}$, $\lambda_*\ge2.1626\cdot10^{-14}$, from \Cref{Y-A-Fr}, ``The order of the constants''\,(0). Indeed, these give $3K_*\mathsf g(\lambda_*/2)\le3\cdot3.1527\cdot10^{12}\cdot51\cdot8/(2.1626\cdot10^{-14})^3\le3.82\cdot10^{56}$.

\item[\textup{(O14)}] The numbers of this entry are defined once all the choices \textup{(O0)} to \textup{(O13)} have been made, and no floor depends on them. We have $\vartheta_w=0.248$ (\Cref{Y-H-C9}\,(b)) and $\vartheta:=(16\vartheta_w+\frac12)/10=0.4468$.

We put $c_4:=15\,c_Q^{\rm su}C_1^2/E_a$, and for $\mathsf b\ge x_\flat$ we put $\Lambda:=\log_L(c_4\mathsf b^{1/2})$ and $K_0(\mathsf b):=\frac52\Lambda+\frac94$, which is the least length of a run. For an integer $K$ the level of the switch is $k_0:=\lceil(5K+10\Lambda)/9\rceil$. Moreover $\varphi:=5/9$ and $\theta:=\vartheta^{4/9}$, so that $0.69902<\theta\le0.69903$, and we set $\varepsilon_L:=\frac{10}9\log(1/\vartheta)/\log L$ and $C_{\mathsf b}:=(3/\vartheta)(c_4\mathsf b^{1/2})^{\varepsilon_L}$, where $3/\vartheta\le6.715$.

Finally, with $n=(N/M)^3$, \[C^{\rm su}_{\rm vol}(N,x):=4n\,e^{220n}\,e^{\frac32\mathsf N_Zp_0(x)^2N^3}\,(384/\varepsilon(x)^3)^{3N^3},\] and \[Q^{\rm su}(\mathsf b):=19.52\,M^{-3}\bigl[E_ap_0(\mathsf b)\mathsf b^{1/2}+\mathsf b\,c_h(\mathsf b)p_0(\mathsf b)^{-2}\bigr].\] \end{enumerate}

\emph{Item 2: objects at a level.} We now fix the objects that live at one level.

\begin{definition}[Lattice, cubes] A level is a torus $\mathbb Z^3/N\mathbb Z^3$, with $N$ even and $M\mid N$, that carries a grid of cubes of $M^3$ sites each; $n:=(N/M)^3$ is the number of its cubes. It has $3N^3$ bonds $b$ and $3N^3$ plaquettes $p$.

The base site $\pi(p)$ of a plaquette is its least corner, and $|y-p|:=|y-\pi(p)|_\infty$. Each object is assigned to the cube that contains its base site, and we say that $X$ holds $p$ when $p$ lies in a cube of $X$. The distance between two cubes is their sup distance, counted in cubes.

By $\mathcal D$ we denote the family of the non-empty sets of cubes that are connected through faces, edges and corners. \emph{First count:} a given cube is held by at most $(26e)^{j-1}$ members of $\mathcal D$ with $j$ cubes.

By $\mathcal P$ we denote the family of the non-empty finite sets of cubes, and $|X|_c$ is the smallest number of cubes that a member of $\mathcal D$ holding $X$ can have. \emph{Second count:} a given cube is held by at most $(52e)^j$ members $X$ of $\mathcal P$ with $|X|_c=j$. \end{definition}

\begin{proof}[Proof of the two counts] We regard the cubes as the vertices of a graph of degree at most $\Delta=26$, in which two cubes are neighbours when their sup distance is $1$, and we number the neighbours of each cube by $1,\dots,\Delta$. Let $X$ be a connected set of $j$ cubes that holds $\square$. It has a spanning tree rooted at $\square$, and we fix one by a rule: breadth first, and least numbers first. Reading the numbers along it, we see this tree as a rooted subtree, on $j$ vertices, of the infinite rooted tree whose vertices each have $\Delta$ numbered children. The set $X$ consists of the cubes reached from $\square$ along this subtree, so distinct sets $X$ give distinct subtrees. These subtrees are counted by the Fuss--Catalan number $\binom{\Delta j}j/((\Delta-1)j+1)$.

Now $\binom{\Delta j}j\le(\Delta j)^j/j!\le\Delta^je^{j-1}$, because $j!\ge e(j/e)^j$. The last inequality is an equality at $j=1$, and when $j$ is replaced by $j+1$ its left side gains the factor $j+1$, while its right side gains the factor $(j+1)(1+1/j)^j/e\le j+1$. Moreover $(\Delta-1)j+1\ge\Delta$. Hence the number of the sets $X$ is $\le(e\Delta)^{j-1}$, which is the first count. (Encoding the tree by a walk along its contour gives only $676^{j-1}$, and $676e^{-\kappa}>1$ at $\kappa=6$, the least value that \textup{(O5)} admits.)

For $\mathcal P$, let $X$ hold the given cube, with $|X|_c=j$, and pick a member of $\mathcal D$ with $j$ cubes that holds $X$. That member holds the given cube as well, so it is one of at most $(26e)^{j-1}$ sets, and $X$ is one of its $2^j$ subsets. As $(26e)^{j-1}2^j\le(52e)^j$, the second count follows. \end{proof}

\begin{definition}[Lattice, cubes, continued] For a set $X$ of cubes, $X^{+j}$ is the set of the cubes at distance $\le j$ from $X$. For a set $A$ of cubes, the bonds of $A$ are the bonds with both ends in $A$, while the paths and the plaquettes of $A$ are those all of whose sites lie in $A$; a plaquette $p$ is said to be in $A$ when its base site lies in $A$.

A field on the bonds of $A$ is called tame there if $\phi_p\le\varepsilon$ at the plaquettes of $A$ (item~1 of \S\ref{Y-G0-1}; item~2 of \S\ref{Y-H-1}). If $1_A(V)=1$, with $1_A$ as in the definition ``Group, fields'' below, then $V$, restricted to the bonds of $A$, is tame there, because each plaquette of $A$ is in $A$. In $\|\cdot\|_x$, $\mathsf D^{\rm in}_x$ and $\mathsf D_x$, the suprema of $E^h$ are taken over the $G$-valued fields on the bonds of $X^{\approx}$ which are tame there. \end{definition}

\begin{definition}[At $x$] The radii at $x$ are $\varrho(x)=\lceil c_R\log x\rceil$, $r(x)=\lceil20c_R\log x/M\rceil$ and $s(x)=8r(x)$; further $\ell_w(x)=6M\,r(x)+12$, $X^{\approx}=X^{+r(x)}$ and $X^{\approx\approx}=X^{+s(x)}$. \end{definition}

\begin{definition}[Group, fields] Let $\mathfrak g$ be the space of the hermitian traceless $2\times2$ matrices, with $\langle Y,Y'\rangle=\frac12\operatorname{tr}YY'$ and $|Y|^2=\langle Y,Y\rangle$; for $g\in G$ the map $\operatorname{Ad}_gY=gYg^{-1}$ is an isometry of $\mathfrak g$. The angle of $g$ is $\operatorname{ang}(g)=\arccos(\frac12\operatorname{tr}g)\in[0,\pi]$. For $g\ne-1$ we let $\mathsf L(g)$ be the unique $Y\in\mathfrak g$ with $e^{iY}=g$ and $|Y|<\pi$, and we set $\mathsf L(-1):=0$; then $|\mathsf L(g)|=\operatorname{ang}(g)$ for $g\ne-1$.

A field $V$ assigns to each oriented bond an element of $G$, with $V(\bar b)=V(b)^{-1}$, and $dV$ denotes the product of the Haar probability measures. We write $\partial V(z;\mu\nu)=V(z,\mu)V(z+e_\mu,\nu)V(z+e_\nu,\mu)^{-1}V(z,\nu)^{-1}$, denote by $V(\pi)$ the ordered product along a path $\pi$, and put $V^u(b)=u(b_-)V(b)u(b_+)^{-1}$, as in the definition ``Fields'' in item~2 of \S\ref{Y-I-1}. Finally $\phi_p=\operatorname{ang}\partial V(p)$ and $F_p=\mathsf L(\partial V(p))$; the rough plaquettes form the set $D(V)=\{p:\phi_p>\varepsilon\}$; and $1_Y(V)=1\{\phi_p\le\varepsilon\text{ for all }p\text{ in }Y\}$, where we say that $V$ is tame on $Y$ when $1_Y(V)=1$. \end{definition}

\begin{definition}[Energy] Rows $\ell=(\ell_y)$ take their values $\ell_y(p)$ in a real Hilbert space $\mathsf H_y$ and have range $\le\varrho(x)$, which means $\ell_y(p)=0$ for $|y-p|>\varrho(x)$. Their site forms are $\mathsf E_y(p,p'):=\langle\ell_y(p),\ell_y(p')\rangle$, and for two families of site forms we put \[\|\tilde{\mathsf E}-\mathsf E\|_{(\lambda)}:=\sup_{y,p,p'}|(\tilde{\mathsf E}_y-\mathsf E_y)(p,p')|\,e^{\lambda(|y-p|+|y-p'|)};\] this is the second norm of the definition ``Norms of a family of symmetric operators'' in Step~1a.

The coordinate path $\Gamma(y,y'')$ goes first in the direction $e_1$, then $e_2$, then $e_3$, each time the short way; it is defined when the sup distance is less than half the side, and we put $T(y,y'')=\operatorname{Ad}_{V(\Gamma(y,y''))}$. The clip is $\zeta_p=F_p\min\{1,\varepsilon/\phi_p\}$, with multiplier $1$ at $\phi_p=0$; hence it is $0$ at $\partial V(p)=1$ or $-1$. The dressed rows are $\ell_y^V\zeta=\sum_p\ell_y(p)\otimes T(y,\pi(p))\zeta_p$.

The mask $m_y(V)$ is $1$ when $D(V)$ has no plaquette within $\rho_0$ of $y$ and there is no plaquette with $\phi_p>1$ within $\varrho(x)+1$ of $y$; otherwise the mask is $0$. With $S_\varepsilon(t)=t^2/2$ for $t\le\varepsilon$ and $d_0\varepsilon t$ for $t>\varepsilon$, the energy is \[\mathrm{En}_\ell(V)=\sum_pS_\varepsilon(\phi_p)+\frac12\sum_ym_y(V)|\ell_y^V\zeta|^2.\] The clip, the mask and the paths are functions of $V,\varepsilon,\rho_0,\varrho(x)$, and of no rows. \end{definition}

\begin{definition}[Domain] For $X\in\mathcal D$, the complex domain $\mathcal U_X(x)$ consists of the maps $\mathsf V$ that send the oriented bonds of $X^{\approx}$ to $\operatorname{SL}(2,\mathbb C)$, satisfy $\mathsf V(\bar b)=\mathsf V(b)^{-1}$, and have the following three properties, $|\cdot|$ being the operator norm.

\begin{enumerate} \item[\textup{(u1)}] $|\mathsf V(b)|<1+\vartheta_0^\flat(1-1/\log x)$.

\item[\textup{(u2)}] $|\mathsf V(\pi)|<2-1/\log x$ for every path $\pi$ in $X^{\approx}$ with at most $\ell_w(x)$ bonds.

\item[\textup{(u3)}] $|\partial\mathsf V(p)-1|<2\varepsilon$. \end{enumerate}

For $x>e$ this set contains every $G$-valued field that is tame on $X^{\approx}$. Indeed, for such a field the norms in \textup{(u1)}, \textup{(u2)} are equal to $1$, and $1-1/\log x>0$; further $|\partial V(p)-1|=2\sin(\phi_p/2)\le\varepsilon<2\varepsilon$, since the eigenvalues of $\partial V(p)$ are $e^{\pm i\phi_p}$. \end{definition}

\begin{definition}[Members, class, distances] The members, the class and the distances below are those of item~3 of \S\ref{Y-H-1}. For $A\ge1$, $M_A^{\rm su}(x)$ is the set of the triples $u=(E^a,E^h,w)$ that have the following three properties.

\begin{enumerate} \item[\textup{(ca$_A$)}] For $X\in\mathcal D$, the function $E^a(X,\cdot)$ is holomorphic on $\mathcal U_X(x)$, invariant under $G$-valued gauge transformations and equal to $0$ at pure gauges, and $\sup_{\mathcal U_X(x)}|E^a(X,\cdot)|\le A\,c_a(x)e^{-\kappa(\#X-1)}$.

\item[\textup{(ch$_A$)}] $E^h(X,\cdot)$ is a complex Borel gauge invariant function of those $G$-valued fields on the bonds of $X^{\approx}$ that are tame there, and $\sup|E^h(X,\cdot)|\le A\,c_h(x)e^{-\kappa(\#X-1)}$, where the supremum runs over these fields.

\item[\textup{(c$\Xi_A$)}] For $X\in\mathcal P$, the function $w(X,\cdot)$ is complex, Borel and gauge invariant, it depends on $V$ only through the bonds of $X^{\approx\approx}$, and $|w(X,V)|\le A\,e^{-\kappa(|X|_c-1)}$ at every $V$; this bound is the weak cap. \end{enumerate}

The strong cap is the bound $|w(X,V)|\le\exp\{-c_\Xi x\varepsilon\sum_{p\in D(V)\cap X}\phi_p-\kappa(|X|_c-1)\}$ at every $V$. We call $u$ real when $E^a$ is real at $G$-valued fields and $E^h$, $w$ are real. For a triple that satisfies every requirement of \textup{(ca$_A$)}, \textup{(ch$_A$)}, \textup{(c$\Xi_A$)} except the three bounds in which $A$ occurs, $\|u\|_x$ is the least $A\ge0$ with which the three size bounds hold; this is a norm, as in item~3 of \S\ref{Y-H-1} and in item~3 of \S\ref{Y-I-1}.

We put \[\rho[u;\ell](V):=\exp\{-x\mathrm{En}_\ell(V)+\sum_{X\in\mathcal D}1_{X^{\approx}}(V)[E^a+E^h](X,V)\}\cdot\Xi_w(V),\] where $\Xi_w(V)=\sum_{\{X_i\}}\prod_iw(X_i,V)$, the sum running over the finite families of sets $X_i\in\mathcal P$ at pairwise distance $>s(x)$ such that each $X_i$ holds a plaquette of $D(V)$ and their union contains $D(V)$; the empty family, with the product $1$, is admitted exactly when $D(V)$ is empty. The class is \[C^{\rm su}_\flat(x;\ell):=\{e^{\mathsf V}\rho[u;\ell]>0:\ \mathsf V\text{ real},\ u\in M_1^{\rm su}(x)\text{ real with the strong cap}\},\] and such a $u$ is called a system of members of the density; it is not unique.

We put $\mathsf D^{\rm in}_x(u_1,u_2):=\|u_1-u_2\|_x$. The distance $\mathsf D_x(u_1,u_2)$ is the largest of three numbers: first $\sup_X\sup_{\mathcal U_X(x)}|E_1^a-E_2^a|e^{\kappa(\#X-1)}/c_a(x)$; second $\sup_X\sup|E_1^h-E_2^h|e^{\kappa(\#X-1)}/c_h(x)$, where the inner supremum runs over the tame $G$-valued fields, that is, over the $G$-valued fields on the bonds of $X^{\approx}$ which are tame there; and third $\sup_{X,V}|w_1-w_2|(X,V)\exp\{c_\Xi x\varepsilon\sum_{D(V)\cap X}\phi_p+\kappa(|X|_c-1)\}$. Comparing these three numbers with the three size bounds, we find $\mathsf D^{\rm in}_x\le\mathsf D_x$. \end{definition}

\begin{definition}[The step] The drawn step leads from $\mathbb Z^3/N'L\mathbb Z^3$ to $\mathbb Z^3/N'\mathbb Z^3$; it maps $\rho$ to \[\rho'(V')=\int\rho(U)\prod_ck_{ax}(V'(c),\mathsf M(U)(c))\,dU,\] with the kernel $k_s(g,h)=z(s)^{-1}e^{-s(1-\cos\operatorname{ang}(gh^{-1}))}$, where $z(s)=\int_Ge^{-s(1-\cos\operatorname{ang}g)}dg$.

Here $\mathsf M$ is the average introduced in the definition ``The step, the chart'' in item~C.I~4 of \S\ref{Y-C-1}. For a fine site $x$, the coarse site $y(x)$ is the one whose centred block $B(y)$ holds $x$, and $\mathsf t_U(x):=U(\Gamma(y(x),x))$. Given a coarse bond $c=(y,y')$, with $y'=y+Le_\mu$ in fine units, and a site $x\in B(y)$, we let $x':=x+Le_\mu$, write $[x,x']$ for the straight path, and put $h_{c,x}:=\mathsf t_U(x)U([x,x'])\mathsf t_U(x')^{-1}$, $\bar\tau_c:=L^{-3}\sum_{x\in B(y)}\mathsf L(h_{c,x}h_{c,y}^{-1})$ and $\mathsf M(U)(c):=e^{i\bar\tau_c}h_{c,y}$. Thus $\mathsf M$ is a Borel function of $U$, built only from $U$, $L$, the blocks and the paths.

At every $U$ the kernel is a probability kernel: the right invariance of Haar measure gives $\int k_s(g,h)\,dg=z(s)^{-1}\int e^{-s(1-\cos\operatorname{ang}(gh^{-1}))}dg=1$, and hence $\int\rho'dV'=\int\rho\,dU$, by Fubini. Three clauses hold: the kernel is gauge covariant; if $\rho\ge0$ is not zero almost everywhere, then $\rho'>0$; and if $\rho$ is gauge invariant, so is $\rho'$. The step depends on $x$, $a$ and $L$ only: on no rows, on no grid of cubes, on no members, and not on $K$. \end{definition}

\begin{proof}[Proof of the three clauses] Because $(ugv^{-1})(uhv^{-1})^{-1}=u(gh^{-1})u^{-1}$ and $\operatorname{ang}$ is a class function, $k_s(ugv^{-1},uhv^{-1})=k_s(g,h)$. Next, $\mathsf t_{U^u}(x)=u(y(x))\mathsf t_U(x)u(x)^{-1}$. Therefore $h_{c,x}\mapsto u(y)h_{c,x}u(y')^{-1}$, so that $h_{c,x}h_{c,y}^{-1}$ is carried into its conjugate by $u(y)$ and $\bar\tau_c\mapsto\operatorname{Ad}_{u(y)}\bar\tau_c$; for the last point we use $\mathsf L(ugu^{-1})=\operatorname{Ad}_u\mathsf L(g)$, where both sides are $0$ at $g=-1$. This gives $\mathsf M(U^u)(c)=u(y)\mathsf M(U)(c)u(y')^{-1}$, and with the invariance of $k_s$ just shown, the kernel is gauge covariant.

Given a coarse $u'$, we take the fine $u(x):=u'(y(x))$ and change the variable $U\to U^u$, which leaves Haar measure invariant; when $\rho$ is invariant, this yields $\rho'(V'^{u'})=\rho'(V')$. Finally $k_s>0$, so that $\rho'>0$ as soon as $\rho\ge0$ is not zero almost everywhere. \end{proof}

\begin{definition}[Runs] Let $\mathsf b\ge x_\flat$, let $N$ be even with $M\mid N$, and let $K\ge1$. On $\mathbb Z^3/NL^K\mathbb Z^3$ we put $\rho^{[K]}_0=e^{-L^K\mathsf b\sum_p(1-\cos\phi_p)}$, and we define $\rho^{[K]}_{k+1}=(\rho^{[K]}_k)'$ by the step with the noise $a\,x_k$, where $x_k=L^{K-k}\mathsf b$. Then $\rho^{[K]}_k$ lives on $\mathbb T_k=\mathbb Z^3/NL^{K-k}\mathbb Z^3$, with the grids of origin $o_K=0$, $o_k=L\,o_{k+1}-(L-1)/2$. The rows $\ell_k^{[K]}$ are the chain of \Cref{Y-A-Towers}, $\ell_0=\emptyset$, $\ell_{k+1}=((\ell_k)')^{\rm c}$, built from the rows of \Cref{Y-A-P5} cut at $\varrho(x_{k+1})$, then tensored with $1_{\mathfrak g}$ and dressed.

For $K'=K+m$, the partner of the level $k$ of the run of length $K$ is the level $m+k$ of the run of length $K'$, and the depth of this level $k$ is $s=K-k$. Finally $\mu^{[K]}_{[s]}$ is the law on the fields of $\mathbb Z^3/NL^s\mathbb Z^3$ with the density $\rho^{[K]}_{K-s}/\int\rho^{[K]}_{K-s}dV$. \end{definition}

\emph{Item 3: what is used.} We write out each clause that we use.

\textup{(R8)} \emph{\Cref{Y-GB-ThB}.} The letters are those of item~1. Let $K\ge1$, $\beta\ge L^Kx_\flat$, $N$ even, $M\mid N$ and $N\ge2M\,s(L^{1-K}\beta)$. Then, for $1\le k\le K$, the density $\rho_k$ of the run that starts from $e^{-\beta\sum(1-\cos\phi_p)}$ satisfies $\rho_k\in C^{\rm su}_\flat(L^{-k}\beta;\ell_k)$ at every field, with the rows, tori and grids of the definition ``Runs'' in item~2 and with the same letters at every $k$.

\textup{(R9)} \emph{\Cref{Y-H-C9}.} The data \textup{(SD)} at $x$ (item~2 of \S\ref{Y-H-1}) are as follows. First, $x\ge x_9$; $N'$ is even, $M\mid N'$, and \textup{(Tor)} $N'\ge2M\,s(x)$ holds, where $N'$ is the coarse side; the fine torus is $\mathbb Z^3/N'L\mathbb Z^3$; and the grids fit, that is, $o=L\,o'-(L-1)/2$ mod $M$. Second, $\ell\in\mathcal R(b_*,K_*,\lambda_*,\frac19)$ (see \Cref{Y-A-SetsOfRows}) on the fine torus, of range $\le\varrho(x)$, and the pair $(\ell,\ell')$ has the four conclusions of \Cref{Y-A-Th36}: \textup{(Dec)}$(C_D,\theta_D)$; \textup{(N$'$)} with \eqref{eq:row-cut-norms} at every cut; $\mathsf t_\rho\le\frac18(\frac12-d_0-c_\Xi)$; and \textup{(L1)}, \textup{(L2)} in the forms \textup{(L1$^{\rm su}$)}, \textup{(L2$^{\rm su}$)}. By $\ell'^{\rm c}$ we denote the cut of $\ell'$ at $\varrho(x/L)$.

For all data \textup{(SD)} at $x$, all $1\le A\le16$ and all $u\in M_A^{\rm su}(x)$ there exist a number $\mathsf V'\in\mathbb C$ and a triple $F_\ell(u)=(E'^a,E'^h,w')$, given by the formulas \textup{(Fm1)} to \textup{(Fm4)} of item~6 of \S\ref{Y-H-1}, such that the following clauses hold with $x'=x/L$.

\begin{enumerate} \item[\textup{(9a)}] $(\rho[u;\ell])'=e^{\mathsf V'}\rho[F_\ell(u);\ell'^{\rm c}]$ at every $V'$.

\item[\textup{(9b)}] $F_\ell(u)$ satisfies, at $x'$, all clauses of \textup{(ca)}, \textup{(ch)}, \textup{(c$\Xi$)} other than the size bounds; $e^{\kappa(\#Y-1)}\sup_{\mathcal U_Y(x')}|E'^a(Y,\cdot)|\le(0.35+0.248A)c_a(x')$; $e^{\kappa(\#Y-1)}\sup|E'^h(Y,\cdot)|\le c_h(x')$; and $w'$ has the strong cap at $x'$. Consequently, for $A\le2$, where $0.35+0.496=0.846$, we have $F_\ell(u)\in M_1^{\rm su}(x')$ with the strong cap.

\item[\textup{(9d)}] If $u$ is real, then $F_\ell(u)$ and $\mathsf V'$ are real; if moreover $\rho[u;\ell]>0$, then $(\rho[u;\ell])'>0$.

\item[\textup{(9e)}] For $u_1,u_2\in M_2^{\rm su}(x)$, $\mathsf D_{x'}(F_\ell(u_1),F_\ell(u_2))\le0.4\,\mathsf D^{\rm in}_x(u_1,u_2)$. Here $u_1,u_2$ are measured in the weak distance $\mathsf D^{\rm in}_x$ and $F_\ell(u_1),F_\ell(u_2)$ in the strong distance $\mathsf D_{x'}$. \end{enumerate}

Clause \textup{(9e)} is \Cref{Y-I-ThA}\,(ii), where the factor is $0.3968$.

\textup{(R11)} \emph{\Cref{Y-I-LB}.} Let $x\ge e^2$, and consider a torus $\mathbb Z^3/N_f\mathbb Z^3$ with $N_f=N'L$, where $N'$ is even, $M\mid N'$ and $N'\ge2M\,s(x)$; this is the lattice of item~2 of \S\ref{Y-I-1}, whose $N$ is the present $N'$. The torus carries a grid of cubes, and $\ell,\tilde\ell$ are any rows on it of range $\le\varrho(x)$, with no further condition on them. We put $\delta:=\|\tilde{\mathsf E}-\mathsf E\|_{(\lambda_\flat)}$ and $\delta_\Delta:=c_Q^{\rm su}p_0^2\delta/c_a(x)=(c_Q^{\rm su}/E_a)x^{1/2}\delta/p_0(x)$. Let $\delta_\Delta\le1$, and let \textup{(M5r)}, $\varepsilon(x)\le\frac1{40}$, $\varrho(x)+2\le M\,r(x)$, $c_\Xi p_0(x)^2\ge1$ and $(2r(x)+1)^3c_a(x)\le34$ hold. Then each real $\tilde u=(\tilde E^a,\tilde E^h,\tilde w)\in M_1^{\rm su}(x)$ that has the strong cap admits a real $u^\flat=(\tilde E^a+E^\Delta,\tilde E^h,w^\flat)\in M_2^{\rm su}(x)$ such that $\rho[\tilde u;\tilde\ell]=\rho[u^\flat;\ell]$ at every $V$ and $\mathsf D^{\rm in}_x(\tilde u,u^\flat)\le\delta_\Delta$. For $w^\flat$ the lemma gives the weak cap alone.

\textup{(RU)} \emph{\Cref{Y-J-Co50}.} Let $L\ge L_B$ be odd, $aL\ge3/c_L$, $\theta_D\le\lambda_*/8$ and $c_R=64/\theta_D$; let $x_\flat$ satisfy \textup{(X12)}, \textup{(X13)}; and let $\mathsf b\ge x_\flat$, $K\le K'$, $m=K'-K$, and $N$ even. Let $\ell_k$ be the chain on the tower $\mathbb Z^3/NL^{K-k}\mathbb Z^3$ with the radii $\varrho(L^{K-j}\mathsf b)$, and $\tilde\ell_k$ the chain of length $K'$ ending on $\mathbb Z^3/N\mathbb Z^3$, with the radii $\varrho(L^{K'-j}\mathsf b)$. For $0\le k\le K$ we then have $\|\tilde{\mathsf E}_{m+k}-\mathsf E_k\|_{(\lambda_\flat)}\le15C_1^2\vartheta_1^k$, where $\vartheta_1=c_{\rm fix}L^{-1/2}$.

\textup{(RB)} \emph{Base.} We use \Cref{Y-A-Th29}\,(b).

We use \Cref{Y-A-Th36} with $(C_D,\theta_D)$ of \textup{(O6)} as its pair of constants. Its conclusions hold with this pair by \Cref{Y-A-Fr} (clause \textup{(Fr2)}; Step~H13 of its proof) and by the proof of \Cref{Y-A-Th36}, which amounts to clauses \textup{(Fr3$'$)}, \textup{(Fr4)} of the conclusion of \Cref{Y-A-Fr} (Step~H14 of its proof); Step~2\,(d) explains this.

We use the bounds $\|\ell\|_{\rm r}\le3K_*\mathsf g(\lambda_*)$ and $\|\ell\|_{\rm c}\le K_*\mathsf g(\lambda_*)$ on $\mathcal R(b_*,K_*,\lambda_*,\cdot)$, which are stated after the display of \Cref{Y-A-SetsOfRows}.

The torus condition passes to all finer levels. This is \Cref{Y-GB-P63}: if $x\ge e^2$, $L\ge3$ and $N\ge2M\,s(x)$, then $NL^j\ge2M\,s(L^jx)$ for all $j\ge0$.

The angle satisfies $\operatorname{ang}(gh)\le\operatorname{ang}g+\operatorname{ang}h$ and $\operatorname{ang}g^{-1}=\operatorname{ang}g$ for $g,h\in G$; this is part of \Cref{Y-F-L1}\,(a).

\begin{proposition}[Every point of the class is gauge invariant]\label{Y-J-P6} All points of $C^{\rm su}_\flat(x;\ell)$ are gauge invariant.
\end{proposition}

\begin{proof} By \Cref{Y-I-P6}, the objects $\phi_p$, $D(V)$, $m_y$ and $1_Y$ are invariant, while $\ell_y^V\zeta$ is rotated by $1\otimes\operatorname{Ad}_{u(y)}$, which is an isometry. Therefore $\mathrm{En}_\ell$, the indicators $1_{X^{\approx}}$ and the families of \textup{(c$\Xi$)} are invariant, and each member is invariant by the clauses that define it.
\end{proof}

\emph{Item 4: the statements.}

\begin{theorem}[Two cutoffs meet]\label{Y-J-ThA} In the setting \textup{(O0)} to \textup{(O13)} of item~1, with the letters chosen in this order, let $\mathsf b\ge x_\flat$; let $N$ be even, with $M\mid N$ and \textup{(Tor)} $N\ge2M\,s(L\mathsf b)$; let $K'\ge K\ge K_0(\mathsf b)$ and $m:=K'-K$; and write $x_k=L^{K-k}\mathsf b$ and $\ell_k:=\ell_k^{[K]}$. Then the following hold.
\begin{enumerate}
\item[\textup{(i)}] $\Lambda>0$ and $1\le k_0\le K-1$.
\item[\textup{(ii)}] Fix a system $u_{k_0}$ of members of $\rho^{[K]}_{k_0}$ (there is one by \textup{(R8)}), and put $u_{k+1}:=F_{\ell_k}(u_k)$ for $k_0\le k<K$. Then for every $k$ with $k_0\le k\le K$, $u_k$ is a system of members of $\rho^{[K]}_k$: it is real, it lies in $M_1^{\rm su}(x_k)$, it satisfies the strong cap, and $\rho^{[K]}_k=e^{\mathsf V}\rho[u_k;\ell_k]$ with a real $\mathsf V$.
\item[\textup{(iii)}] For every $k$ with $k_0<k\le K$ there exist a real $\mathfrak n_k\in M_1^{\rm su}(x_k)$ satisfying the strong cap and a real $\mathsf V$ such that $\rho^{[K']}_{m+k}=e^{\mathsf V}\rho[\mathfrak n_k;\ell_k^{[K]}]$ at every field. Thus $\rho^{[K]}_k$ and $\rho^{[K']}_{m+k}$ belong to one and the same class $C^{\rm su}_\flat(x_k;\ell_k^{[K]})$.
\item[\textup{(iv)}] For $k_0<k\le K$ one has $\mathsf D_{x_k}(u_k,\mathfrak n_k)\le3\vartheta^{k-k_0}$, and at $k=K$ one has $\mathsf D_{\mathsf b}(u_K,\mathfrak n_K)\le C_{\mathsf b}\theta^K$. In terms of the depth $s=K-k$, this says: the level $K-s$ is $>k_0$ if and only if $K-s\ge K_0(L^s\mathsf b)=\frac52\Lambda+\frac54s+\frac94$, and in that case $\mathsf D\le3\vartheta^{\frac49K-s-\frac{10}9\Lambda-1}=C_{L^s\mathsf b}\theta^{K-s}$.
\item[\textup{(v)}] In a cruder form, $k_0<\varphi K+K_0(\mathsf b)$ and $k_0\le\lceil\varphi K\rceil+K_0(\mathsf b)$. Consequently, at every level $k$ with $\varphi K+K_0(\mathsf b)<k\le K$, clause \textup{(iii)} holds and $\mathsf D\le3\vartheta^{k-\lceil\varphi K\rceil-K_0(\mathsf b)}$.
\end{enumerate}
Here $\varphi,\vartheta,\theta$ are numerals, and $K_0(\mathsf b)$, $C_{\mathsf b}$ depend on $L,M,E_a,\mathsf b$ and pure numbers; none of these quantities depends on $N$, $K'$, $m$ or a row.
\end{theorem}

\begin{theorem}[From the distance of members to densities]\label{Y-J-ThB} In the setting \textup{(O0)} to \textup{(O12)} of item~1, let $x\ge e^2$; let $N$ be even, with $M\mid N$ and $N\ge2M\,s(x)$; let $\kappa\ge6$, $\varepsilon(x)\le1$ and $(c_a+c_h)(x)\le\frac12$; let $\ell$ be rows on $\mathbb Z^3/N\mathbb Z^3$ of range $\le\varrho(x)$; and let $\rho_i=e^{\mathsf V_i}\rho[u_i;\ell]\in C^{\rm su}_\flat(x;\ell)$ for $i=1,2$. Put $\mathsf D:=\mathsf D^{\rm in}_x(u_1,u_2)$, which is $\le\mathsf D_x(u_1,u_2)$, so that each bound below remains valid with $\mathsf D_x$ in place of $\mathsf D$; and put $n=(N/M)^3$. Then the following hold.
\begin{enumerate}
\item[\textup{(a)}] At every $V$ for which $D(V)$ is empty, \[|\log\rho_1(V)-\log\rho_2(V)-(\mathsf V_1-\mathsf V_2)|\le1.22\,n(c_a+c_h)(x)\,\mathsf D.\]
\item[\textup{(b)}] $\sup_V|\rho_1-e^{\mathsf V_1-\mathsf V_2}\rho_2|\le e^{\mathsf V_1}\cdot2n\,e^{219n}\,\mathsf D$.
\item[\textup{(c)}] If in addition $|\ell_y(p)|\le K_*e^{-\lambda_*|y-p|}$, then \[\int\rho_1dV\ge e^{\mathsf V_1}e^{-0.61n}e^{-\frac32\mathsf N_Zp_0^2N^3}(\varepsilon^3/384)^{3N^3}\] and \[\int|\rho_1/\!\int\!\rho_1-\rho_2/\!\int\!\rho_2|\,dV\le C^{\rm su}_{\rm vol}(N,x)\,\mathsf D.\]
\item[\textup{(d)}] The integral $\int|\rho_1/\!\int\!\rho_1-\rho_2/\!\int\!\rho_2|\,dV$ is $\le2\max_ir_i+e^{2.44n(c_a+c_h)\mathsf D}-1$, where $r_i$ denotes the mass of $\{D(V)\text{ non-empty}\}$ under $\rho_i/\!\int\!\rho_i$.
\end{enumerate}
\end{theorem}

\begin{corollary}[The laws converge in total variation]\label{Y-J-CorC} Let the letters, $\mathsf b$ and $N$ satisfy the hypotheses that \Cref{Y-J-ThA} imposes on them, and let $s\ge0$.
\begin{enumerate}
\item[\textup{(a)}] For all $K'\ge K\ge s+K_0(L^s\mathsf b)$, \[\|\mu^{[K']}_{[s]}-\mu^{[K]}_{[s]}\|_{\rm TV}\le C^{\rm su}_{\rm vol}(NL^s,L^s\mathsf b)\,C_{L^s\mathsf b}\,\theta^{K-s}.\]
\item[\textup{(b)}] As $K\to\infty$, the full sequence converges: $\mu^{[K]}_{[s]}\to\mu^\infty_{[s]}$ in total variation, at the rate of \textup{(a)}. The limit $\mu^\infty_{[s]}$ is a gauge invariant probability law, absolutely continuous with respect to Haar measure.
\item[\textup{(c)}] For $s\ge1$, the law $\mu^\infty_{[s-1]}$ is obtained from $\mu^\infty_{[s]}$ by the drawn step with the noise $aL^s\mathsf b$.
\item[\textup{(d)}] For $K\ge s+K_0(L^s\mathsf b)$, the distance in $L^1(dV)$ between the density of $\mu^\infty_{[s]}$ and the normalised density of the point $\rho^{[K]}_{K-s}$ of $C^{\rm su}_\flat(L^s\mathsf b;\ell^{[K]}_{K-s})$ is at most the bound of \textup{(a)}.
\end{enumerate}
\end{corollary}

\begin{corollary}[The flux response]\label{Y-J-CorD} Let the letters, $\mathsf b$ and $N$ satisfy the hypotheses that \Cref{Y-J-ThA} imposes on them, let moreover $N^2\ge4\pi/\varepsilon(\mathsf b)$, and let $m_0$ be an integer with $\varepsilon(\mathsf b)/2\le\varepsilon_0:=2\pi m_0/N^2\le\varepsilon(\mathsf b)$; such an integer exists. On $\mathbb Z^3/N\mathbb Z^3$, with $0\le z_i<N$, let $V_*$ be given by $V_*(z,e_2)=e^{i\varepsilon_0z_1\sigma_3}$, by $V_*(z,e_1)=e^{-i\varepsilon_0Nz_2\sigma_3}$ if $z_1=N-1$, and by the value $1$ at all other bonds. For $\rho>0$ set \[\mathfrak x(\rho):=\frac2{N^3\varepsilon_0^2}[\log\rho(1)-\log\rho(V_*)].\] Then for all $K'\ge K\ge K_0(\mathsf b)$, \[|\mathfrak x(\rho^{[K']}_{K'})-\mathfrak x(\rho^{[K]}_K)|\le Q^{\rm su}(\mathsf b)\,C_{\mathsf b}\,\theta^K.\]
\end{corollary}

\emph{Item 5: remarks.}

\begin{remark} The system $\mathfrak n_k$ of \Cref{Y-J-ThA}\,(iii) is a system of members over the rows $\ell_k^{[K]}$ of the shorter run, while the run of length $K'$ has members of its own, over its rows $\ell_{m+k}^{[K']}$; each of the two is a system of members of one and the same function $\rho^{[K']}_{m+k}$. Members are not unique, and the proofs use no identity between the two systems.

The proofs depend on $c_Q^{\rm su}$ only through two facts, that \textup{(R11)} holds with it and that $480c_Q^{\rm su}C_1^2>1$, and on $\vartheta$ only through $0.4\le\vartheta<1$. Inserting $\vartheta=0.4$ in all formulas of \textup{(O14)} gives $\theta=0.4^{4/9}<0.66549$.
\end{remark}

\begin{proof}[Proof of \Cref{Y-J-ThA}] \emph{Step 0: the letters; three inequalities that follow from their choice; the existence of a number $x_\flat$ that satisfies all conditions of \textup{(O13)}; no circularity in the choice of the letters.} We claim the following six facts.
\begin{enumerate}
\item[\textup{(i)}] $\theta_D\le\lambda_*/8$, so that $c_R\ge512$ and $e_R\ge256$.
\item[\textup{(ii)}] For $x\ge e^2$, $\varrho(x)+2\le M\,r(x)$ and $2(\varrho(x)+1)<2M\,s(x)$.
\item[\textup{(iii)}] $480\,c_Q^{\rm su}C_1^2>2000$.
\item[\textup{(iv)}] There is a number $x_\flat$ that satisfies all conditions of \textup{(O13)}.
\item[\textup{(v)}] \textup{(R8)}, \textup{(R9)}, \textup{(R11)} and \textup{(RU)} hold for one and the same choice of the letters, and no entry of this choice depends on a later one.
\item[\textup{(vi)}] \textup{(X18)} gives $(2r+1)^3c_a\le34$.
\end{enumerate}

We prove the six facts in turn.
\begin{enumerate}
\item[\textup{(i)}] By \textup{(O6)}, $\theta_D\le\theta_{\mathbb C}/2$. Next, $\theta_{\mathbb C}=\theta^\flat/6\le\gamma/(12A^\flat)$ and $A^\flat\ge384e^8K_*^2$; these follow from \textup{(O6)}, from clause~\textup{(6)} of the order of the constants in \Cref{Y-A-Fr}, and from $\mathsf g_2\ge1$. As $K_*^2=9/\gamma_*$ and $\gamma\le1$, we get $\theta_{\mathbb C}\le\gamma_*/(41472e^8)$, while $\lambda_*/4=\gamma_*^{1/2}/176$. The bound $\theta_{\mathbb C}\le\lambda_*/4$ therefore comes down to $\gamma_*^{1/2}\le41472e^8/176$, which follows from $\gamma_*^{1/2}=44\lambda_*\le44<41472/176$. Thus $\theta_D\le\lambda_*/8$, whence $c_R=64/\theta_D\ge512/\lambda_*\ge512$ and $e_R=\lambda_*c_R/2\ge256$.
\item[\textup{(ii)}] We have $M\,r(x)\ge20c_R\log x$ and $\varrho(x)+2\le c_R\log x+3$, and moreover $19c_R\log x\ge19\cdot512\cdot2>3$; this gives the first inequality. The second follows from $2M\,s(x)=16M\,r(x)\ge320c_R\log x>2(c_R\log x+2)\ge2(\varrho(x)+1)$.
\item[\textup{(iii)}] From $K_*\lambda_*=3/44$ and $\lambda_*\le1$ we obtain $K_*\ge3/44$; since also $\mathsf g(\lambda_*/2)\ge1$, it follows that $C_1^2=3K_*^2\mathsf g(\lambda_*/2)^3\ge3(3/44)^2>0.0139$. We also have $c_Q^{\rm su}\ge306$, and $480\cdot306\cdot0.0139>2000$.
\item[\textup{(iv)}] The conditions $x_\flat\ge x_\flat^{(1)}$, $x_\flat\ge x_\flat^0$, $L\,x_\flat\ge x_9$, \textup{(X12)} and \textup{(X13)} of \textup{(O13)} bound $x_\flat$ from below by finite numbers. Here $x_\flat^{(1)}$, $x_\flat^0$ and $x_9$ are numbers defined in \textup{(O13)} of item~1, while \textup{(X12)} and \textup{(X13)} hold for all large $x_\flat$, because $x_\flat^{-e_R}\to0$ as $x_\flat\to\infty$. For the floors \textup{(X14)} to \textup{(X18)}, the limits of $p_0$, $\varepsilon$, $c_a$ and $c_h$ as $x\to\infty$ follow from the definitions of item~1. By \textup{(O8)}, $p_0(x)=A_0(\log x)^{\mathfrak p_0}$ with $A_0>0$ and $\mathfrak p_0>2$, and $\varepsilon=p_0x^{-1/2}$; hence $p_0\to\infty$ and $\varepsilon=p_0x^{-1/2}\to0$. Under \textup{(O12)}, the caps are $c_a(x)=E_ap_0(x)^3x^{-1/2}$ and $c_h(x)=e^{-c_hp_0(x)^2}$, so that $c_a,c_h\to0$. With $r(x)=\lceil20c_R\log x/M\rceil$, these definitions give \[c_a(5+2r)^3\le E_aA_0^3(\log x)^{3\mathfrak p_0}x^{-1/2}(7+40c_R\log x/M)^3\to0.\] Thus the five floors hold for all large $x$: the limit of $p_0$ yields \textup{(X14)} and \textup{(X15)}, that of $\varepsilon$ yields \textup{(X16)}, those of $c_a$ and $c_h$ yield \textup{(X17)}, and the display yields \textup{(X18)}. The inequality of each floor holds strictly at the limit that yields it; hence each floor is satisfied by all $x$ beyond some number.
\item[\textup{(v)}] The conditions of \textup{(O13)} depend on the letters of \textup{(O0)} to \textup{(O12)} only. \Cref{Y-B-C1} holds for all $c_h\le\delta_3/2$, all $E_a\ge\max\{32L^{-1/2}K_v^\flat,32K_{{\rm q},k}\}$ and all $x_\flat\ge x_\flat^{(1)}$, where $x_\flat^{(1)}$ does not depend on $c_h,E_a,N$. \Cref{Y-GB-ThB} holds for $c_h\le\delta_3/2$, $E_a\ge\max\{32L^{-1/2}K_v^\flat,32K_{{\rm q},k}\}$ and $x_\flat\ge\max\{x_\flat^{(1)},x_\flat^0\}$ with \textup{(X12)}. \Cref{Y-H-C9} holds for $c_h\le\delta_3'/2$, $E_a\ge\max\{32K_{{\rm q},k},1500K_v^\flat L^{-1/2}\}$ and $x\ge x_9$. Since $\delta_3'\le\delta_3$, and since the maximum in \textup{(O12)} is the largest of these maxima, the choice of the table lies on all these half-lines; and the class does not depend on $x_\flat$. The condition \textup{(M5r)} is a lower bound on $M$ in terms of the radius $\rho_0$ that \textup{(O3)} defines, and Thms.~\ref{Y-B-C1}, \ref{Y-G-ThA}, \ref{Y-GB-ThB} and~\ref{Y-H-C9} admit every larger $M=L^{m'}$. Finally, \textup{(O14)} is defined last and enters no floor.
\item[\textup{(vi)}] Indeed $(2r+1)^3\le(5+2r)^3$, and $\frac12\le34$.
\end{enumerate}

\emph{Step 1: \Cref{Y-J-Co50} for the letters of $\operatorname{SU}(2)$.} We claim that, for $\mathsf b\ge x_\flat$, $K'\ge K$, $m=K'-K$, $N$ even and $0\le k\le K$, \[\delta_k:=\|\tilde{\mathsf E}_{m+k}-\mathsf E_k\|_{(\lambda_\flat)}\le15C_1^2L^{-2k/5},\] where $\mathsf E_k$ and $\tilde{\mathsf E}_j$ denote the site forms of the rows $\ell_k^{[K]}$ and $\ell_j^{[K']}$.

By the definition ``Runs'' of item~2 and by \textup{(R8)}, the rows of both runs are the scalar chains of \Cref{Y-A-Towers}, tensored with $1_{\mathfrak g}$, and $\langle\ell_y(p)\otimes1,\ell_y(p')\otimes1\rangle$ enters only through the scalar site form $\mathsf E_y(p,p')=\langle\ell_y(p),\ell_y(p')\rangle$.

We take the hypotheses of \textup{(RU)} one by one. That $L$ is odd, $L\ge L_B$ and $aL\ge3/c_L$ is \textup{(O4)}. The relations $\theta_D\le\lambda_*/8$ and $c_R=64/\theta_D$ come from \textup{(O6)} and Step~0\,(i). Conditions \textup{(X12)} and \textup{(X13)} hold by \textup{(O13)}, and $\mathsf b\ge x_\flat$, $K\le K'$ and the evenness of $N$ belong to the data. As to the radii, the chain of the run of length $K$ is cut at $\varrho(x_j)=\varrho(L^{K-j}\mathsf b)$, the chain of the run of length $K'$ at $\varrho(L^{K'-j}\mathsf b)$, and both chains end on $\mathbb Z^3/N\mathbb Z^3$.

Hence $\delta_k\le15C_1^2\vartheta_1^k$, and $\vartheta_1=c_{\rm fix}L^{-1/2}\le L^{1/10}L^{-1/2}=L^{-2/5}$ because $L\ge L_B\ge c_{\rm fix}^{10}$. \Cref{Y-J-Co50}, together with \Cref{Y-J-Th49}, \Cref{Y-J-P48}, Lems.~\ref{Y-J-L41} and~\ref{Y-J-L42}, \Cref{Y-J-P43} and Lems.~\ref{Y-J-L44}, \ref{Y-J-L46} and~\ref{Y-J-L47} on which it rests, is stated and proved in Step~1a.

\emph{Step 1a: \textup{(RU)}, that is \Cref{Y-J-Co50}, and its proof.} The claim of this step is \textup{(RU)} of item~3, which is the last inequality of \Cref{Y-J-Co50}.

Objects and statements are taken from \S\ref{Y-A-1}, each with the meaning given to it by the passage named in the parentheses after it: the lattices, $W$, $Q_i$ and the facts on them (\Cref{Y-A-Lattices} and \Cref{Y-A-Facts}); $\mathsf d$, ${\rm dist}$ and $\mathsf g(\lambda)$ (\Cref{Y-A-Distances}); own units, $F_\beta(z)$, $\mathsf g_1$, $\partial_\alpha$, $\nabla$ and $\nabla'$ (Defs.~\ref{Y-A-OwnUnits}, \ref{Y-A-Sums} and~\ref{Y-A-Differences}), where $|(\nabla F)_p|^2:=\sum_\alpha|(\partial_\alpha F)_p|^2$; rows, covariance, $\tau(\ell)$ and $\mathcal R(b,K,\lambda,\tau)$ (\Cref{Y-A-Rows}, and Defs.~\ref{Y-A-CovRows} to~\ref{Y-A-SetsOfRows}); pairs, $\Xi_Z$, $\mathsf T_\sharp$ and $R_\sharp$ (\Cref{Y-A-Pairs}); $\mathsf B'$, $H^\sharp$, $\mathsf N^\sharp$ and the statement that defines them (\Cref{Y-A-P4}\,(a)); $\Phi$ (\Cref{Y-A-MapPhi}); the five groups of $\ell'$ (\Cref{Y-A-P5}); $\mathsf S_0$ and $\mathfrak R$ (\Cref{Y-A-AuxForm}); $\mathsf r_\sharp$, $\gamma_1$ and \Cref{Y-A-P6} (\Cref{Y-A-Cst} and \Cref{Y-A-C13}); the cut and $\ell'^>$ (\Cref{Y-A-CutRows}); towers and chains (\Cref{Y-A-Towers}); and the pure numbers and \textup{(Cut)} (\Cref{Y-A-PureNumbers} and \Cref{Y-A-Th29}).

Lems.~\ref{Y-A-L25} and~\ref{Y-A-L26}, \Cref{Y-A-P28} and \Cref{Y-A-Th29} are proved in \S\ref{Y-A-1}. Clauses \textup{(a)} and \textup{(b)} of \Cref{Y-J-L41} are \Cref{Y-A-L24}, and the condition \textup{(X12)} of \S\ref{Y-A-1} is the \textup{(X12)} of \textup{(O13)}.

For the statements of Step~1a and for their proofs, we fix the meaning of several letters. Wherever Step~1a treats the passage from the fine torus $\mathbb T$ to the coarse torus $\mathbb T'$ in a single step, $N$ is the coarse side of the step, with $\mathbb T=\mathbb Z^3/NL\mathbb Z^3$ and $\mathbb T'=L\mathbb Z^3/NL\mathbb Z^3$ as in \Cref{Y-A-Lattices}; it is the $N'$ of this section. Throughout Step~1a we allow every even $N\ge2$ and do not impose the condition $M\mid N$. In \Cref{Y-J-Co50}, by contrast, $N$ is the side of the last torus, as in the definition ``Runs'' of item~2. In \Cref{Y-J-Th49}, $k\le k'$ are the lengths of the two chains and $n$ is the side of the torus on which both end; \Cref{Y-J-Co50} applies that theorem with $(k,k',n)=(k,m+k,NL^{K-k})$. The letters $\mathrm{d},\mathrm{d}',\delta$ are operators of the free model of \S\ref{Y-A-1}. In \Cref{Y-J-L44}\,(b), $\hat{\mathsf B}(0)$ and $\hat{\mathsf B}'(0)$ denote the values taken at the momentum $0$ by the symbols of $\mathsf B,\mathsf B'$; in other words, they are the $3\times3$ matrices $1+\mathsf c(\mathsf E)$ and $1+\mathsf c(\mathsf E')$ by which $\mathsf B,\mathsf B'$ act on the constant forms. In \Cref{Y-J-P48}\,(b) and in Step~\textup{(7)} of its proof, $\varrho\in\mathbb N\cup\{\infty\}$ is a radius in the sense of \Cref{Y-A-Towers}. The distance $r_P={\rm dist}(y',P)$ equals $|y'-P|$ in the units of $\mathbb T'$ (\Cref{Y-A-Distances}, for coarse cells); accordingly \Cref{Y-J-L46} is applied on $\mathbb T'$ with this distance. In the proofs of \Cref{Y-J-P43} and \Cref{Y-J-L44}, $\mathsf R=\mathsf B^{-1}$, as in \Cref{Y-A-Rows}.

The proofs of Step~1a use the exact constants and outward intervals in \Cref{Y-A-Fr}\,\textup{(0)}, together with the following five numerical facts. In them, $C_1$, $c_0$, $c_{\rm fix}$, $L_B$ and $\vartheta_1$ are as in \textup{(O0)} of item~1, $C_0$ and $\mu$ are fixed in the definition ``Rates, constants'' below, and $\mathsf y$ is the number defined in Step~\textup{(6)} of the proof of \Cref{Y-J-P48}.
\begin{enumerate}
\item[\textup{(1)}] The 27 points with $|n|_\infty\le1$ give $\mathsf g(\lambda)\ge27e^{-\lambda}$, hence $\mathsf g(\lambda)\ge10$ for $\lambda\le\lambda_*$; this yields the bound $\mathsf g(\lambda_*/4)^2\ge100$. Moreover $\mathsf g(\lambda)\ge3$ gives $3K_*\mathsf g(\lambda_*/2)\le\sqrt3K_*\mathsf g(\lambda_*/2)^{3/2}=C_1$.
\item[\textup{(2)}] One has $\gamma_*^{-1/2}/(9b_*)=\sqrt{12296}/(9b_*\mathsf r_\sharp)=\sqrt{12296}/18<6.17$; consequently $C_0/c_0=(\gamma_*^{-1/2}/(9b_*))\,\mathsf g(\lambda_*)/\mathsf g(\lambda_*/4)^3\le6.17/100$.
\item[\textup{(3)}] Since $2\mu=\lambda_*/4<10^{-14}$, we have $e^{\lambda_*}\le1.1$, $6e^{2\mu}+6\le13$ and $3+3e^{2\mu}+1\le7.1$.
\item[\textup{(4)}] One has $b_*\ge27$.
\item[\textup{(5)}] One has $c_{\rm fix}\ge48$, because $\mathsf g(\lambda)\ge1$, $C_1>1$ and $c_0>1$. Hence $L\ge L_B\ge48^{10}>32$; further $\vartheta_1\le L^{-2/5}\le48^{-4}$, and $\mathsf y=\vartheta_1/48\le\frac1{100}$.
\end{enumerate}

We now recall statements of \S\ref{Y-A-1}, beginning with \Cref{Y-A-L25}. On $\Omega'^2\oplus\Omega^3\oplus\mathsf H$ let \[\mathbb A^{\rm e}(Z,\Xi,\sigma)=(W^*Z+\delta\Xi-\ell^{\mathsf T}\sigma,\ \sigma,\ a^{-1/2}\mathrm{d}'^*Z,\ \mathsf S_0^{1/2}Q_3\Xi),\] whose rows are indexed, component by component, by $p$, by the components of the spaces $\mathsf H_y$ (each placed at $y$), by $c$, and by $q'$; let $\mathfrak R^{\rm e}=\mathbb A^{{\rm e}\mathsf T}\mathbb A^{\rm e}$ and $\mathsf D^{\rm e}={\rm diag}(L,L^{-2},1)$. \Cref{Y-A-L25} states the following.
\begin{enumerate}
\item[\textup{(a)}] At the source $(\delta_P,0,0)$, $\mathbb A^{\rm e}(\mathfrak R^{\rm e})^{-1}(\delta_P,0,0)=(\nu,\ell\nu,a^{1/2}\mathsf N^\sharp\delta_P,\cdot)$, where $\nu=H^\sharp\delta_P$.
\item[\textup{(b)}] $\mathfrak R^{\rm e}\ge(\gamma_1/2b)\mathsf D^{\rm e}$ for $b\ge b_{\mathsf B}$.
\item[\textup{(c)}] For $\ell\in\mathcal R(b_*,K_*,\lambda_*,\cdot)$ and $L\ge L_A$, the bound \textup{($\sinh$)} of \Cref{Y-J-L41} is valid for $\mathbb A^{\rm e},\mathsf D^{\rm e},\mathsf d$ at $\lambda_*$, with some $\varkappa\le\frac14\gamma_*^{1/2}$. Here \Cref{Y-A-L25}\,(c) has the rate $\lambda'$, a minimum defined in that lemma, and $\lambda'=\lambda_*$ for $\ell\in\mathcal R(b_*,K_*,\lambda_*,\cdot)$ and $L\ge L_A$, because $L$ is at least the first member of the maximum that defines $L_A$.
\end{enumerate}

Next we recall \Cref{Y-A-L26}. A pair $(F,\mathsf n)$ is \emph{admitted at} $X'$ if $WF+\mathrm{d}'\mathsf n=X'$ and $\mathrm{d}F\perp\ker Q_3$; these pairs form the set ${\rm Adm}(X')$, which does not depend on $\ell$. Each entry below bears the tag of the clause of \Cref{Y-A-L26} to which it belongs.
\begin{enumerate}
\item[\textup{(a)}] $(\bar F,\bar{\mathsf n})\in{\rm Adm}(X')$, and $\langle(\bar F,\bar{\mathsf n}),(F,\mathsf n)\rangle_{\mathcal H}=\langle\bar Z,X'\rangle$ for every $(F,\mathsf n)\in{\rm Adm}(X')$; this identity is shown in the proof of \Cref{Y-A-L26}\,(a). Consequently $(\bar F,\bar{\mathsf n})$ minimises $\|\cdot\|_{\mathcal H}$ on ${\rm Adm}(X')$ and is $\mathcal H$-orthogonal to ${\rm Adm}(0)$. Conversely, an element of ${\rm Adm}(X')$ which is $\mathcal H$-orthogonal to ${\rm Adm}(0)$ is the minimiser, since its difference with the minimiser lies in ${\rm Adm}(0)$ and is orthogonal to both of them.
\item[\textup{(b)}] $\Phi$ is monotone.
\end{enumerate}

Finally we recall \Cref{Y-A-P28} and \Cref{Y-A-Th29} for the starred constants. Let $L\ge L_A$ be odd, let $aL\ge3/c_L$, let $N\ge2$ be even, and let $\ell\in\mathcal R(b_*,K_*,\lambda_*,\tau)$ with $\tau\le\frac19$. Then $\ell'$ is translation covariant on $\mathbb T'$, $|\ell'_{y'}(P)|\le K_*e^{-\lambda_*{\rm dist}(y',P)}$, $1\le\mathsf B'\le\frac32\mathsf r_\sharp^{-1}$, and \[\sum_pe^{2\lambda_*\mathsf d(p,P)}(H^\sharp\delta_P)_p^2\le4\gamma_*^{-1}L^{-1}.\] If moreover $\mathsf t_*(\varrho)\le\frac19$, then $\|\ell'^>\|_{\rm r}:=\sup_{y'}\sum_P|\ell'^>_{y'}(P)|\le\mathsf t_*(\varrho)$ and $\ell'^{\rm c}\in\mathcal R(b_*,K_*,\lambda_*,\tau+\mathsf t_*(\varrho))$. Along a chain with \textup{(Cut)}, $\ell_k\in\mathcal R(b_*,K_*,\lambda_*,\sum_{j\le k}\mathsf t_*(\varrho_j))$; this is shown in the proof of \Cref{Y-A-Th29}\,(b). For every $F$, and for rows with $|\ell_y(p)|\le Ke^{-\lambda|y-p|}$, one has $|\ell F|\le3K\mathsf g_1(\lambda)|\nabla F|+\tau(\ell)|F|$, as shown within the proof of \Cref{Y-A-P28}\,(b).

\begin{definition}[Rates, constants] We fix the pure rates $\mu:=\frac18\lambda_*$ and $\lambda_\flat:=\mu_0:=\frac1{16}\lambda_*$, and the pure constants \[C_1:=\sqrt3\,K_*\,\mathsf g(\tfrac{\lambda_*}2)^{3/2},\quad c_0:=108\,b_*\gamma_*^{-1/2}\mathsf g(\tfrac{\lambda_*}4)^3,\quad C_0:=12\,\gamma_*^{-1}\mathsf g(\lambda_*),\] \[c_{\rm fix}:=48\,\mathsf g(\lambda_\flat)\,C_1\,c_0^{1/2},\qquad L_B:=\max\{L_A,c_{\rm fix}^{10}\}.\] We also set $\vartheta_0:=c_0/L$ and $\vartheta_1:=c_{\rm fix}L^{-1/2}$, the latter being $\le L^{-2/5}$ for $L\ge L_B$. Note that $C_0\le c_0$, because $\gamma_*^{-1/2}/9b_*=\sqrt{12296}/18<6.17$, $\mathsf g(\lambda_*)\le \mathsf g(\frac{\lambda_*}4)$ and $\mathsf g(\frac{\lambda_*}4)^2\ge100$.
\end{definition}

\begin{definition}[Site forms] The site form at $y$ is $\mathsf E_y:=\ell_y^{\mathsf T}\ell_y$, a positive operator on $\Omega^2$, and $\sum_y\mathsf E_y=\mathsf B-1$. The cut acts on site forms by $\mathsf E'^{\rm c}_{y'}=\chi_{y'}\mathsf E'_{y'}\chi_{y'}$, where $\chi_{y'}=1\{{\rm dist}(y',\cdot)\le\varrho\}$.
\end{definition}

\begin{definition}[Constant forms, charge] Let the constant $2$-form $e_\beta\in\Omega^2$ take the value $1$ on the plaquettes whose normal is the unit vector $e_\beta$ and the value $0$ on the others, and write $h=\sum h_\beta e_\beta$; let $e'_\beta$ and $h'$ be the same objects on $\mathbb T'$, and let $h\leftrightarrow h'$ pair forms with equal components. For translation covariant rows, put $\mathsf c(\mathsf E)_{\alpha\beta}:=\langle e_\alpha,\mathsf E_ye_\beta\rangle=\langle\ell_ye_\alpha,\ell_ye_\beta\rangle$. This is a positive $3\times3$ matrix that does not depend on $y$, because $\ell_{y+\mathsf t}e_\alpha=U_{\mathsf t,y}\ell_ye_\alpha$; moreover $\mathsf c_{\beta\beta}\le\tau(\ell)^2$ by the definition of $\tau(\ell)$, since $|e_\beta|_\infty=1$.
\end{definition}

\begin{definition}[Seminorms at a site] In own units, for $\eta\ge0$, let \[n_y(F)^2:=\sum_pe^{-2\mu|y-p|}|(\nabla F)_p|^2,\qquad|F|_y^2:=\sum_pe^{-2\mu_0|y-p|}F_p^2,\qquad\hat n_{y,\eta}(F)^2:=n_y(F)^2+\eta^2|F|_y^2.\]
\end{definition}

\begin{definition}[Norms of a family of symmetric operators] For a family $D=(D_y)_y$ of symmetric operators on $\Omega^2$, let \[\|D\|_\eta:=\inf\{\epsilon:\ |\langle F,D_yG\rangle|\le\epsilon\,\hat n_{y,\eta}(F)\,\hat n_{y,\eta}(G)\ \text{for all }y,F,G\}\] and \[\|D\|_{(\lambda)}:=\sup_{y,p,p'}|D_y(p,p')|\,e^{\lambda(|y-p|+|y-p'|)}.\]
\end{definition}

These seminorms and norms satisfy two bounds.
\begin{enumerate}
\item[\textup{(n1)}] For $\eta\le1$, $\hat n_{y,\eta}(G)^2\le14|G|_y^2$. Indeed, $|(\nabla G)_p|^2\le2\sum_\alpha(G_{p+e_\alpha}^2+G_p^2)$, and therefore \[\sum_pe^{-2\mu|y-p|}|(\nabla G)_p|^2\le(6e^{2\mu}+6)\sum_pe^{-2\mu|y-p|}G_p^2\le13|G|_y^2.\]
\item[\textup{(n2)}] $n_y(\delta_p)^2\le(3+3e^{2\mu})e^{-2\mu|y-p|}$; hence, for $\eta\le1$, $\hat n_{y,\eta}(\delta_p)^2\le7.1e^{-2\mu_0|y-p|}$ and $\|D\|_{(\lambda_\flat)}\le7.1\|D\|_\eta$.
\end{enumerate}
Moreover $\|D\|_\eta$ does not increase in $\eta$, and $\|D\|_0<\infty$ only if every $D_y$ kills constants, since $n_y(h)=0$ for each constant form $h$.

The standing conditions for the rest of Step~1a are that $L$ is odd, that $aL\ge3/c_L$, and that $N\ge2$ is even.

\begin{lemma}[Exponentially weighted estimates for $\mathfrak R^{-1}$ and for the range projection of $\mathbb A$]\label{Y-J-L41} Let $I$ and $J$ be finite sets, and let each point of $I$ and of $J$ be assigned a place, that is, a point of a pseudo-metric space with distance $\mathsf d$. In $\mathsf d(r,i)$, $\mathsf d(i,j)$ and $\mathsf d(r,j)$, the distance is taken between the places of the two points. Let $\mathbb A$ be a real $J\times I$ matrix and $\mathfrak R=\mathbb A^{\mathsf T}\mathbb A$; let $\mathsf D$ be a positive diagonal matrix on $I$, with $N(\phi)=\langle\phi,\mathsf D\phi\rangle$; and let $\gamma>0$ satisfy $\mathfrak R\ge\gamma\mathsf D$. Let $\lambda>0$ and $0\le\varkappa\le\frac14\gamma^{1/2}$ be such that \begin{equation}\tag{$\sinh$}\sum_{r\in J}\Bigl(\sum_{i\in I}|\mathbb A(r,i)|\sinh(\lambda\mathsf d(r,i))|\phi_i|\Bigr)^2\le\varkappa^2N(\phi)\end{equation} for all $\phi$. Then the following hold.
\begin{enumerate}
\item[\textup{(a)}] For $j\in I$ and $\Phi=\mathfrak R^{-1}\delta_j$, $\sum_i\mathsf D_ie^{2\lambda\mathsf d(i,j)}\Phi_i^2\le4\gamma^{-2}\mathsf D_j^{-1}$.
\item[\textup{(b)}] For the same $j$ and $\Phi$, $\sum_re^{2\lambda\mathsf d(r,j)}(\mathbb A\Phi)_r^2\le4\gamma^{-1}\mathsf D_j^{-1}$.
\item[\textup{(c)}] Let $\Pi:=\mathbb A\mathfrak R^{-1}\mathbb A^{\mathsf T}$, which is the orthogonal projection of $\mathbb R^J$ onto ${\rm ran}\,\mathbb A$. If $w$ is any real function on $I\cup J$ such that $|w_i-w_r|\le\lambda\mathsf d(r,i)$ for all $r\in J$, $i\in I$, then each $\rho\in\mathbb R^J$ satisfies $|e^w\Pi\rho|\le3|e^w\rho|$ and $|e^w(1-\Pi)\rho|\le4|e^w\rho|$.
\end{enumerate}
\end{lemma}

\begin{proof}
Clauses \textup{(a)} and \textup{(b)} are \Cref{Y-A-L24} under the same hypotheses. To prove \textup{(c)}, $w$ is a real function on $I\cup J$ as in \textup{(c)}, so that $|w_i-w_r|\le\lambda\mathsf d(r,i)$, and we put $\mathfrak R_w:=e^w\mathfrak Re^{-w}$. To a real vector $\phi$ on $I$ we attach the two vectors on $J$ with the entries $\mathsf C_r=\sum_i\mathbb A(r,i)\cosh(w_i-w_r)\phi_i$ and $\mathsf S_r=\sum_i\mathbb A(r,i)\sinh(w_i-w_r)\phi_i$.

Writing $e^{w_i-w_r}=\cosh(w_i-w_r)+\sinh(w_i-w_r)$, we find $e^{-w_r}(\mathbb Ae^w\phi)_r=\mathsf C_r+\mathsf S_r$, and in the same way $e^{w_r}(\mathbb Ae^{-w}\phi)_r=\mathsf C_r-\mathsf S_r$. Since $\mathfrak R=\mathbb A^{\mathsf T}\mathbb A$, it follows that $\langle\phi,\mathfrak R_w\phi\rangle=\langle\mathbb Ae^w\phi,\mathbb Ae^{-w}\phi\rangle=|\mathsf C|^2-|\mathsf S|^2$.

Next we compare both vectors with $\mathbb A\phi$. As $\sinh$ increases on $[0,\infty)$ and $\cosh t-1\le\sinh|t|$, the condition on $w$ and \textup{($\sinh$)} give $|\mathsf S|,|\mathsf C-\mathbb A\phi|\le\varkappa N^{1/2}\le\frac14|\mathbb A\phi|$, where $N=N(\phi)$; the last inequality follows from $\varkappa\le\frac14\gamma^{1/2}$ and $|\mathbb A\phi|\ge\gamma^{1/2}N^{1/2}$, which is a consequence of $\mathfrak R\ge\gamma\mathsf D$.

From these bounds and the identity above we obtain \begin{equation}\tag{$\dagger$}|\mathsf C|-|\mathsf S|\ge\frac12|\mathbb A\phi|,\qquad|\mathsf C|+|\mathsf S|\le\frac32|\mathbb A\phi|,\qquad\langle\phi,\mathfrak R_w\phi\rangle\ge\frac12|\mathbb A\phi|^2\ge\frac\gamma2N(\phi).\end{equation}

For \textup{(c)} we fix $\rho$ on $J$ and put $\Phi:=\mathfrak R^{-1}\mathbb A^{\mathsf T}\rho$ and $\phi:=e^w\Phi$; then $\Pi\rho=\mathbb A\Phi$ and $\mathfrak R_w\phi=e^w\mathbb A^{\mathsf T}\rho$. By the Cauchy--Schwarz inequality, \[\langle\phi,\mathfrak R_w\phi\rangle=\langle e^w\phi,\mathbb A^{\mathsf T}\rho\rangle=\sum_r(\mathbb Ae^w\phi)_r\rho_r=\sum_r(\mathsf C_r+\mathsf S_r)e^{w_r}\rho_r\le(|\mathsf C|+|\mathsf S|)\,|e^w\rho|.\]

On the other hand, the identity for $\langle\phi,\mathfrak R_w\phi\rangle$ shown above gives $\langle\phi,\mathfrak R_w\phi\rangle=(|\mathsf C|-|\mathsf S|)(|\mathsf C|+|\mathsf S|)$. If $|\mathsf C|+|\mathsf S|=0$, then $\phi=0$ by \textup{($\dagger$)}, and so $|\mathsf C|-|\mathsf S|=0$; otherwise we divide by $|\mathsf C|+|\mathsf S|$. In both cases $|\mathsf C|-|\mathsf S|\le|e^w\rho|$.

With \textup{($\dagger$)} this gives $|\mathbb A\phi|\le2|e^w\rho|$, and therefore $|e^w\Pi\rho|=|e^w\mathbb A\Phi|=|\mathsf C-\mathsf S|\le\frac32|\mathbb A\phi|\le3|e^w\rho|$, where the second equality is the identity $e^{w_r}(\mathbb Ae^{-w}\phi)_r=\mathsf C_r-\mathsf S_r$. The bound for $1-\Pi$ follows by the triangle inequality, $|e^w(1-\Pi)\rho|\le|e^w\rho|+|e^w\Pi\rho|$. \end{proof}

\begin{lemma}[A Weitzenb\"ock identity for $2$-forms]\label{Y-J-L42} For $F,G\in\Omega^2(\mathbb Z^3/\mathcal N\mathbb Z^3)$ we have $\sum_\alpha\langle\partial_\alpha F,\partial_\alpha G\rangle=\langle \mathrm{d}F,\mathrm{d}G\rangle+\langle \mathrm{d}^*F,\mathrm{d}^*G\rangle$. In particular $|\mathrm{d}F|,|\mathrm{d}^*F|\le|\nabla F|$. \end{lemma}

\begin{proof}[Proof of \Cref{Y-J-L42}] By polarisation it suffices to treat $G=F$. All operators that occur commute with the translations of the torus, so that by Parseval's identity we may consider one momentum $k$ at a time. At this momentum we write $v_\beta=e^{ik_\beta}-1$, and we let $E\in\mathbb C^3$ be the vector of the components $F_\beta$, where the orientation signs are absorbed into $E$; then $\widehat{\mathrm{d}F}=\sum_\beta v_\beta E_\beta$.

On a $1$-form $A$ the coboundary is $(\mathrm{d}A)_{\mu\nu}=\partial_\mu A_\nu-\partial_\nu A_\mu$, which at the momentum $k$ is the vector product $v\times\hat A$. Taking adjoints, $\widehat{\mathrm{d}^*F}=(v\times)^*E=-\bar v\times E$.

For complex vectors the Lagrange identity reads $|a\times b|^2=|a|^2|b|^2-|\sum a_i\bar b_i|^2$. At $a=\bar v$ and $b=E$ it becomes $|\bar v\times E|^2+|\sum v_\beta E_\beta|^2=|v|^2|E|^2$. The left side is the symbol of $|\mathrm{d}^*F|^2+|\mathrm{d}F|^2$, and the right side is the symbol of $|\nabla F|^2$; this proves the identity. \end{proof}

\begin{proposition}[A gradient bound for the minimiser]\label{Y-J-P43} Let the rows $\ell$ be translation covariant, with $\mathsf B\le b$ and $\mathsf B'\le b'$. Then every $X'$ satisfies $|\nabla H^\sharp X'|\le(b'+\sqrt{27}(bb')^{1/4})L^{-3/2}|\nabla'X'|$. If $b=b_*$ and $b'\le b_*$, this gives $|\nabla H^\sharp X'|\le2b_*L^{-3/2}|\nabla'X'|$, and for the column of a coarse plaquette $P$ it gives $|\nabla H^\sharp\delta_P|\le2\sqrt6\,b_*L^{-3/2}$. \end{proposition}

\begin{proof}[Proof of \Cref{Y-J-P43}] Let $\bar F=H^\sharp X'$. By the definition of $H^\sharp$ we have $\mathsf B\bar F=\bar Y=W^*\bar Z+\delta\Xi_{\bar Z}$, which defines $\bar Y$. As $\ell$ is translation covariant, $\mathsf B\ge1$ commutes with each $\partial_\alpha$, and \Cref{Y-J-L42} gives \[|\nabla\bar F|^2\le\sum_\alpha\langle\partial_\alpha\bar F,\mathsf B\partial_\alpha\bar F\rangle=\sum_\alpha\langle\partial_\alpha\bar F,\partial_\alpha\bar Y\rangle=\langle \mathrm{d}\bar F,\mathrm{d}\bar Y\rangle+\langle \mathrm{d}^*\bar F,\mathrm{d}^*\bar Y\rangle.\] We bound the two terms on the right in turn.

For the first term, $\mathrm{d}\bar F\perp\ker Q_3$, because $(\bar F,\bar{\mathsf n})$ lies in ${\rm Adm}(X')$; so $\mathrm{d}\bar F=Q_3^*\theta$ for a $3$-cochain $\theta$. From $Q_3\mathrm{d}=\mathrm{d}'W$, $W\bar F+\mathrm{d}'\bar{\mathsf n}=X'$ and $\mathrm{d}'\mathrm{d}'=0$ we get $Q_3\mathrm{d}\bar F=\mathrm{d}'W\bar F=\mathrm{d}'(X'-\mathrm{d}'\bar{\mathsf n})=\mathrm{d}'X'$. Hence $\theta=(Q_3Q_3^*)^{-1}\mathrm{d}'X'$, and as $Q_3Q_3^*\ge L^3/27$ and $|\mathrm{d}'X'|\le|\nabla'X'|$ by \Cref{Y-J-L42} on $\mathbb T'$, we have $|\theta|\le27L^{-3}|\nabla'X'|$.

It follows that $\langle \mathrm{d}\bar F,\mathrm{d}\bar Y\rangle=\langle\theta,Q_3\mathrm{d}\bar Y\rangle=\langle\theta,\mathrm{d}'W\bar Y\rangle$. The map $\mathcal O:X'\mapsto W\bar Y$ is linear, and $\mathcal OX'=W\bar Y=W\mathsf BH^\sharp X'$ because $\bar F=H^\sharp X'$ and $\mathsf B\bar F=\bar Y$. Since $\ell$ is translation covariant, $\mathsf B$ commutes with the translations of $\mathbb T$, and the proof of \Cref{Y-A-P28}\,(a) shows that $W$, $H^\sharp$ and $\mathsf B'$ commute with those of $\mathbb T'$; hence so does $\mathcal O$. To bound its norm, note that $|\bar Y|^2\le b\langle\bar Y,\mathsf R\bar Y\rangle\le bL^{-1}\langle X',\mathsf B'X'\rangle$, where the first inequality uses $\mathsf B\le b$ and the second uses $\langle\bar Y,\mathsf R\bar Y\rangle=\langle\bar F,\mathsf B\bar F\rangle\le\|(\bar F,\bar{\mathsf n})\|^2_{\mathcal H}=L^{-1}\langle X',\mathsf B'X'\rangle$; with $\|W\|^2\le L$ and $\mathsf B'\le b'$ we obtain $\|\mathcal O\|\le(bb')^{1/2}$. By \Cref{Y-J-L42} on $\mathbb T'$ and the commutation with translations, $|\mathrm{d}'W\bar Y|\le|\nabla'\mathcal OX'|=|\mathcal O\nabla'X'|\le(bb')^{1/2}|\nabla'X'|$.

For the second term we use $\mathrm{d}^*\delta=0$ and $\mathrm{d}^*W^*=Q_1^*\mathrm{d}'^*$; with $\bar Z=L^{-1}\mathsf B'X'$ they give $\mathrm{d}^*\bar Y=L^{-1}Q_1^*\mathrm{d}'^*\mathsf B'X'$. As we saw in bounding the first term, $\mathsf B'$ commutes with the translations of $\mathbb T'$; hence $\nabla'\mathsf B'X'=\mathsf B'\nabla'X'$. \Cref{Y-J-L42} on $\mathbb T'$, this identity and $\mathsf B'\le b'$ yield $|\mathrm{d}'^*\mathsf B'X'|\le|\mathsf B'\nabla'X'|\le b'|\nabla'X'|$, so that $\|Q_1\|^2\le L^{-1}$ implies $|\mathrm{d}^*\bar Y|\le b'L^{-3/2}|\nabla'X'|$; moreover $|\mathrm{d}^*\bar F|\le|\nabla\bar F|$ by \Cref{Y-J-L42}.

If $\nabla'X'=0$, the bounds above give $\nabla\bar F=0$, and the claim holds. If $\nabla'X'\ne0$, then for $u=|\nabla\bar F|L^{3/2}/|\nabla'X'|$ the bounds for the two terms combine to $u^2\le27(bb')^{1/2}+b'u$, so $u\le b'+\sqrt{27}(bb')^{1/4}$, because every $u>b'+\sqrt{27}(bb')^{1/4}$ satisfies $u(u-b')>27(bb')^{1/2}$. This is the first claim.

Finally, $b_*\ge27$, so that $\sqrt{27}(b_*b')^{1/4}\le\sqrt{27}\,b_*^{1/2}\le b_*$ when $b'\le b_*$; this gives the second claim. The third is the second at $X'=\delta_P$, as $|\nabla'\delta_P|^2=6$. \end{proof}

\begin{lemma}[The step on constant forms]\label{Y-J-L44} Let $L$ be odd, let the rows $\ell$ on $\mathbb T$ be translation covariant and $\tau(\ell)$ arbitrary, let $h'$ be a constant form on $\mathbb T'$, and let $h$ be the constant form on $\mathbb T$ with the same components as $h'$.

\begin{enumerate}\item[\textup{(a)}] $H^\sharp h'=L^{-2}h$ and $\mathsf N^\sharp h'=0$.

\item[\textup{(b)}] For all sites $y'$ of $\mathbb T'$ and $y$ of $\mathbb T$, all constant forms $h_1,h_2$ on $\mathbb T$ and the constant forms $h'_1,h'_2$ on $\mathbb T'$ with the same components, $\langle\ell'_{y'}h'_1,\ell'_{y'}h'_2\rangle=\langle\ell_yh_1,\ell_yh_2\rangle$; in other words, $\mathsf c(\mathsf E')=\mathsf c(\mathsf E)$ as $3\times3$ matrices. Consequently $\tau(\ell')=\tau(\ell)$ and $\hat{\mathsf B}'(0)=\hat{\mathsf B}(0)$.\end{enumerate}\end{lemma}

\begin{proof}[Proof of \Cref{Y-J-L44}] We begin with \textup{(a)}. As $\mathsf B$ commutes with translations, $\mathsf Bh=:g$ is a constant form. Let $g'$ be the constant form on $\mathbb T'$ with the components of $g$, and put $Z:=L^{-1}g'$.

Since $W$ couples only plaquettes of equal orientation and $\sum_PW(P,p)=L^{-1}$, we get $W^*Z=L^{-1}\cdot L^{-1}g$; also $\mathrm{d}'^*Z=0$, as $Z$ is constant. Therefore ${\rm pair}(Z,0)=(L^{-2}\mathsf Rg,0)=(L^{-2}h,0)$. This element lies in ${\rm Adm}(h')$, because $\sum_pW(P,p)=L^2$ gives $W(L^{-2}h)=h'$, and because $dh=0$.

For every $(G,\mathsf m)\in{\rm Adm}(0)$ we have $\langle{\rm pair}(Z,0),(G,\mathsf m)\rangle_{\mathcal H}=\langle Z,WG+\mathrm{d}'\mathsf m\rangle=0$. \Cref{Y-A-L26}\,(a), as recalled in Step~1a of the proof of \Cref{Y-J-ThA}, states that an element of ${\rm Adm}(X')$ that is $\mathcal H$-orthogonal to ${\rm Adm}(0)$ is the minimiser. Hence ${\rm pair}(Z,0)$ is the pair $(H^\sharp h',\mathsf N^\sharp h')$, and \textup{(a)} is proved.

For \textup{(b)} we fix $y'$ and treat the groups of $\ell'_{y'}$ from \Cref{Y-A-P5} one at a time.

\begin{enumerate}\item[\textup{(1), (2)}] These groups consist of differences of $\zeta$ at $y'$, and they vanish on constants.

\item[\textup{(3)}] Here part \textup{(a)} and $W^*h'=L^{-1}h$ give $L^{-2}h-L^{-1}W^*h'=0$, so that this group vanishes at $h'$.

\item[\textup{(4)}] The two terms of this group, the one with $\mathsf N^\sharp$ and the one with $\mathrm{d}'^*$, are $0$ at $h'$, by \textup{(a)} and $\mathrm{d}'^*h'=0$.

\item[\textup{(5)}] By \textup{(a)}, this group contributes $L\sum_{y\in B(y')}\langle\ell_yL^{-2}h_1,\ell_yL^{-2}h_2\rangle=L\cdot L^{-4}\cdot L^3\,h_1^{\mathsf T}\mathsf c(\mathsf E)h_2$ to the inner product, since $B(y')$ has $L^3$ sites and, by the covariance of $\ell$, $\langle\ell_yh_1,\ell_yh_2\rangle=h_1^{\mathsf T}\mathsf c(\mathsf E)h_2$ for every $y$.\end{enumerate}

Adding the contributions of the groups, we find $\langle\ell'_{y'}h'_1,\ell'_{y'}h'_2\rangle=h_1^{\mathsf T}\mathsf c(\mathsf E)h_2$, which is \textup{(b)}; the last two equalities in \textup{(b)} follow, since $\tau$ and $\hat{\mathsf B}(0)$ are determined by the matrix $\mathsf c(\mathsf E)$. \end{proof}

\begin{lemma}[A form compared with its value at a site]\label{Y-J-L46} Let $y$ be a site of $\mathbb Z^3/n\mathbb Z^3$, and let $\mathsf a\ge0$ be a function on the plaquettes with $\mathsf a(p)\le Ae^{-\theta|y-p|}$, where $\theta\ge4\mu$. For $F\in\Omega^2$ let $h^{(y)}$ be the constant form with the components $F_\beta(y)$.

\begin{enumerate}\item[\textup{Claim 1.}] $\sum_p\mathsf a(p)\,|F_p-h^{(y)}_p|\le\sqrt3\,A\,\mathsf g(\tfrac\theta2)^{3/2}\,n_y(F)$.

\item[\textup{Claim 2.}] For rows with $|\ell_y(p)|\le K_*e^{-\lambda_*|y-p|}$, $|\ell_yF|\le C_1n_y(F)+\tau(\ell)|F|_y$.

\item[\textup{Claim 3.}] If moreover $\tau(\ell)\le C_1\eta$, then $|\ell_yF|\le\sqrt2\,C_1\hat n_{y,\eta}(F)$.\end{enumerate}\end{lemma}

\begin{proof}[Proof of \Cref{Y-J-L46}] Each class of $\mathbb Z^3/n\mathbb Z^3$ has exactly one representative $u\in(-\frac n2,\frac n2]^3$, and the sup norm of $u$ is the periodic norm of the class; we work with these representatives and write $|u|$ for this norm. For such $u$, ${\rm path}(u)$ is the lattice path from $0$ to $u$ that runs first along $e_1$, then along $e_2$ and then along $e_3$; it is simple, that is, it visits no point twice, and it stays in the box spanned by $0$ and $u$.

Each step of the path joins some $v$ and $v+e_\alpha$, and we call $v$ its lower end. Along this step $F_\beta(y+\cdot)$ changes by $\pm(\partial_\alpha F)_{(y+v;\beta)}$, where $(y+v;\beta)$ denotes the plaquette based at $y+v$ with normal $e_\beta$. Let $c(u,v)$ be the number of steps of ${\rm path}(u)$ whose lower end is $v$. On a simple path every point lies on at most two steps, and the starting point on at most one; so $c\le2$ and $c(u,0)\le1$. Moreover $c(u,v)\ne0$ only for $v\in{\rm path}(u)$, and then $|v|_\infty\le|u|_\infty$.

Summing along the path and using $|(\partial_\alpha F)_p|\le|(\nabla F)_p|$, we obtain $|F_\beta(y+u)-F_\beta(y)|\le\sum_vc(u,v)|(\nabla F)_{(y+v;\beta)}|$. The plaquette based at $y+u$ with normal $e_\beta$ has distance $|u|$ from $y$, so \[\sum_p\mathsf a(p)\,|F_p-h^{(y)}_p|\le\sum_{\beta,v}\kappa(v)|(\nabla F)_{(y+v;\beta)}|,\qquad\kappa(v):=A\sum_uc(u,v)e^{-\theta|u|}.\] At $v=0$, the bound $c(u,0)\le1$ gives $\kappa(0)\le A\mathsf g(\theta)\le A\mathsf g(\frac\theta2)$.

We define ${\rm path}(u)$ for every $u\in\mathbb Z^3$ in the same way as for the representatives, and we write $|u|$ for the sup norm $|u|_\infty$ of such $u$. Now let $v\ne0$, and let $U_v:=\{u\in\mathbb Z^3:v\in{\rm path}(u)\}$. In the definition of $\kappa(v)$, the summation variable $u$ runs over the representatives, a subset of $\mathbb Z^3$. For $u\in U_v$ each coordinate $v_i$ lies between $0$ and $u_i$, and some $v_i$ is nonzero and has the sign of $u_i$; hence $v$ lies outside the box spanned by $0$ and $-u$. This shows $-U_v\cap U_v=\emptyset$, and both sets are contained in $\{|u|_\infty\ge|v|_\infty\}$.

From $c\le2$, the symmetry $|-u|=|u|$, this disjointness and $e^{-\theta|u|}\le e^{-\frac\theta2|v|}e^{-\frac\theta2|u|}$ for $|u|\ge|v|$ we get \[\sum_uc(u,v)e^{-\theta|u|}\le\sum_{U_v}e^{-\theta|u|}+\sum_{-U_v}e^{-\theta|u|}\le\sum_{|u|\ge|v|}e^{-\theta|u|}\le e^{-\frac\theta2|v|}\mathsf g(\frac\theta2).\] Together with the case $v=0$, this shows $\kappa(v)\le Ae^{-\frac\theta2|v|}\mathsf g(\frac\theta2)$ for all $v$.

We apply the Cauchy--Schwarz inequality to the sum over $\beta$ and $v$, with the weights $e^{\mp\mu|v|}$. The factor with $e^{-\mu|v|}$ is at most $n_y(F)$, since $e^{-2\mu|v|}$ is the weight of the plaquette $(y+v;\beta)$ in $n_y(F)^2$; the other factor is the square root of \[3\sum_ve^{2\mu|v|}\kappa(v)^2\le3A^2\mathsf g(\frac\theta2)^2\sum_ve^{-(\theta-2\mu)|v|}\le3A^2\mathsf g(\frac\theta2)^3,\] where the last step uses $\theta-2\mu\ge\frac\theta2$. This proves Claim~1.

For Claim~2 we write $|\ell_yF|\le\sum_p|\ell_y(p)||F_p-h^{(y)}_p|+|\ell_yh^{(y)}|$ and apply Claim~1 with $\mathsf a(p)=|\ell_y(p)|$, $A=K_*$ and $\theta=\lambda_*=8\mu$, which bounds the sum by $C_1n_y(F)$. For the last term, $|\ell_yh^{(y)}|\le\tau|h^{(y)}|_\infty\le\tau|F|_y$, because each of the three plaquettes based at $y$ enters $|F|_y$ with the weight $1$. This proves Claim~2. If moreover $\tau(\ell)\le C_1\eta$, the right side of Claim~2 is at most $C_1(n_y(F)+\eta|F|_y)\le\sqrt2\,C_1\hat n_{y,\eta}(F)$, which is Claim~3. \end{proof}

\begin{lemma}[The change of the minimiser as a projection]\label{Y-J-L47} Let $\ell,\tilde\ell$ be rows on $\mathbb T$. Objects with a tilde are defined from $\tilde\ell$ as the corresponding objects without a tilde are defined from $\ell$: $\tilde{\mathsf B}=1+\tilde\ell^{\mathsf T}\tilde\ell$; the inner product of $\tilde{\mathcal H}$ is the inner product of \Cref{Y-A-Pairs} with $\mathsf B$ replaced by $\tilde{\mathsf B}$, $\langle(F,\mathsf n),(F',\mathsf n')\rangle_{\tilde{\mathcal H}}=\langle F,\tilde{\mathsf B}F'\rangle+a\langle\mathsf n,\mathsf n'\rangle$; $\tilde H^\sharp$ and $\tilde{\mathsf N}^\sharp$ are the maps $H^\sharp$ and $\mathsf N^\sharp$ of \Cref{Y-A-P4}\,(a), defined with $\tilde{\mathsf B}$ instead of $\mathsf B$; and $\tilde{\mathsf E}_y=\tilde\ell_y^{\mathsf T}\tilde\ell_y$. Put $\mathsf D:=\tilde{\mathsf B}-\mathsf B$, $\Delta H:=\tilde H^\sharp-H^\sharp$ and $\Delta\mathsf N:=\tilde{\mathsf N}^\sharp-\mathsf N^\sharp$. Let $\tilde\Pi$ be the projection of \Cref{Y-J-L41}\,(c) for $\tilde{\mathbb A}^{\rm e}$, the matrix of \Cref{Y-A-L25} with $\tilde\ell$ in place of $\ell$.

\begin{enumerate}\item[\textup{(a)}] $\ker\tilde{\mathbb A}^{{\rm e}\mathsf T}=\{(G,\tilde\ell G,a^{1/2}\mathsf m,0):(G,\mathsf m)\in{\rm Adm}(0)\}$, and the Euclidean inner product of two vectors of this form is the inner product $\langle\cdot,\cdot\rangle_{\tilde{\mathcal H}}$ of the corresponding pairs $(G,\mathsf m)$.

\item[\textup{(b)}] For every $X'$, \[\bigl(\Delta HX',\ \tilde\ell\,\Delta HX',\ a^{1/2}\Delta\mathsf NX',\ 0\bigr)=-(1-\tilde\Pi)\bigl(\mathsf DH^\sharp X',0,0,0\bigr).\]\end{enumerate}\end{lemma}

\begin{proof}[Proof of \Cref{Y-J-L47}] For \textup{(a)}, the definition of $\tilde{\mathbb A}^{\rm e}$ gives $\tilde{\mathbb A}^{{\rm e}\mathsf T}(\rho_1,\rho_2,\rho_3,\rho_4)=(W\rho_1+a^{-1/2}\mathrm{d}'\rho_3,\ \mathrm{d}\rho_1+\mathsf S_0^{1/2}Q_3^*\rho_4,\ \rho_2-\tilde\ell\rho_1)$. Let $(\rho_1,\rho_2,\rho_3,\rho_4)$ lie in the kernel, and put $G=\rho_1$ and $\mathsf m=a^{-1/2}\rho_3$. Then $WG+\mathrm{d}'\mathsf m=0$, $\rho_2=\tilde\ell G$ and $\mathrm{d}G=-\mathsf S_0^{1/2}Q_3^*\rho_4$. Applying $Q_3$ to the last identity, we find on its left side $Q_3\mathrm{d}G=\mathrm{d}'WG=-\mathrm{d}'\mathrm{d}'\mathsf m=0$; so $Q_3Q_3^*\rho_4=0$, hence $\rho_4=0$ because $Q_3Q_3^*\ge L^3/27$, and then $\mathrm{d}G=0$. Thus $(G,\mathsf m)\in{\rm Adm}(0)$, and the vector has the stated form.

Conversely, let $(G,\mathsf m)\in{\rm Adm}(0)$. Then $\mathrm{d}G\in{\rm ran}\,Q_3^*$, since $\mathrm{d}G\perp\ker Q_3$, and $Q_3\mathrm{d}G=0$, because $Q_3\mathrm{d}=\mathrm{d}'W$, $WG=-\mathrm{d}'\mathsf m$ and $\mathrm{d}'\mathrm{d}'=0$; together these give $\mathrm{d}G=0$, and the formula for the transpose shows that $(G,\tilde\ell G,a^{1/2}\mathsf m,0)$ lies in the kernel. The statement on the inner products holds because $\tilde{\mathsf B}=1+\tilde\ell^{\mathsf T}\tilde\ell$.

For \textup{(b)}, note that $(\Delta HX',\Delta\mathsf NX')\in{\rm Adm}(0)$, being the difference of two elements of ${\rm Adm}(X')$. As ${\rm Adm}$ does not depend on the rows, it is the same for $\ell$ and for $\tilde\ell$. \Cref{Y-A-L26}\,(a), as recalled in Step~1a of the proof of \Cref{Y-J-ThA}, shows that $(H^\sharp X',\mathsf N^\sharp X')$ is $\mathcal H$-orthogonal to ${\rm Adm}(0)$, and, with $\tilde\ell$ in place of $\ell$, that $(\tilde H^\sharp X',\tilde{\mathsf N}^\sharp X')$ is $\tilde{\mathcal H}$-orthogonal to ${\rm Adm}(0)$. For $(G,\mathsf m)\in{\rm Adm}(0)$ this gives \[\langle(G,\mathsf m),(\Delta HX',\Delta\mathsf NX')\rangle_{\tilde{\mathcal H}}=0-\langle(G,\mathsf m),(H^\sharp X',\mathsf N^\sharp X')\rangle_{\tilde{\mathcal H}}=-\langle G,\mathsf DH^\sharp X'\rangle-0,\] where in the second step we write $\tilde{\mathsf B}=\mathsf B+\mathsf D$.

Hence, by \textup{(a)}, the left side of \textup{(b)} is a vector of $\ker\tilde{\mathbb A}^{{\rm e}\mathsf T}={\rm ran}(1-\tilde\Pi)$, and, by \textup{(a)} again, its inner product with any vector of that space equals the inner product of $-(\mathsf DH^\sharp X',0,0,0)$ with the same vector. It is therefore the orthogonal projection of $-(\mathsf DH^\sharp X',0,0,0)$ onto ${\rm ran}(1-\tilde\Pi)$, which is \textup{(b)}. \end{proof}

\begin{proposition}[One step contracts the difference of the site forms]\label{Y-J-P48} Let $L\ge L_B$ be odd, $aL\ge3/c_L$ and $N\ge2$ even. Let $\ell,\tilde\ell\in\mathcal R(b_*,K_*,\lambda_*,\tau)$ be rows on $\mathbb T$ with $\tau\le\frac19$, and let $0\le\eta\le1$ with $\tau\le C_1\eta$. Let $D_y:=\tilde{\mathsf E}_y-\mathsf E_y$, with $\epsilon:=\|D\|_\eta\le2C_1^2$, and put $\eta_+:=L^{1/2}\eta$.

\begin{enumerate}\item[\textup{(a)}] For the uncut rows $\ell'$ and $\tilde\ell'$, $\|\tilde{\mathsf E}'-\mathsf E'\|_{\eta_+}\le\vartheta_1\,\epsilon$.

\item[\textup{(b)}] For the rows $\ell'^{\rm c}$ and $\tilde\ell'^{\rm c}$, both cut at the same radius $\varrho$, $\|\tilde{\mathsf E}'^{\rm c}-\mathsf E'^{\rm c}\|_{\eta'}\le\vartheta_1\,\epsilon$ with $\eta'^2:=L\eta^2+6e^{-\lambda_*\varrho/8}$.\end{enumerate}\end{proposition}

\begin{proof}[Proof of \Cref{Y-J-P48}] We fix $y'\in\mathbb T'$ and define $w:=-\lambda_*\mathsf d(\cdot,y')$ on all cells and rows of the step; by the triangle inequality, $|w_i-w_r|\le\lambda_*\mathsf d(r,i)$. We write $H=H^\sharp$ for the map defined from the rows $\ell$. For $F\in\Omega^2$ we set $\beta[F]^2:=L\bigl(|e^w\nabla F|^2+\eta^2|e^wF|^2\bigr)$, and we abbreviate $\mathsf g_\flat:=\mathsf g(\lambda_\flat)$.

\emph{Step (1): sites of a block.} We claim that $L\sum_{y\in B(y')}\hat n_{y,\eta}(F)^2\le1.1\,\mathsf g_\flat\,\beta[F]^2$. Fix a plaquette $p$ and let $\mathsf a:=\min_{y\in B(y')}|y-p|$; as $\mu_0=\lambda_\flat$, we have $\sum_{y\in B(y')}e^{-2\mu_0|y-p|}\le e^{-\mu_0\mathsf a}\mathsf g_\flat$. On the other hand $|p-y'|\le\mathsf a+\mathsf h_L$, so that $2\lambda_*\mathsf d(p,y')\le\frac{2\lambda_*}L\mathsf a+\lambda_*\le\mu_0\mathsf a+\lambda_*$, where we use $L\ge32$; since $e^{\lambda_*}\le1.1$, this gives $e^{-\mu_0\mathsf a}\le1.1\,e^{2w_p}$. The same bounds hold with $\mu\ge\mu_0$ in place of $\mu_0$ in the sum over $y$. Exchanging the sums over $y$ and $p$ in $L\sum_{y\in B(y')}\hat n_{y,\eta}(F)^2$ and inserting these bounds, we obtain the claim.

\emph{Step (2): the source.} We claim that $|e^w\mathsf DF|\le\sqrt{14}\,\mathsf g_\flat\,\epsilon\,L^{-1/2}\beta[F]$, where $\mathsf D=\sum_yD_y$. Writing $\hat n_y=\hat n_{y,\eta}$, we get from the definition of $\epsilon$ and the Cauchy--Schwarz inequality, for every $G$, \[|\langle G,\mathsf DF\rangle|\le\epsilon\sum_y\hat n_y(G)\hat n_y(F)\le\epsilon(\sum_ye^{-2w_y}\hat n_y(G)^2)^{1/2}(\sum_ye^{2w_y}\hat n_y(F)^2)^{1/2}.\] Since $|w_y-w_p|\le\frac{\lambda_*}L|y-p|\le\frac12\mu_0|y-p|$, again by $L\ge32$, we have $\sum_ye^{\mp2w_y}e^{-2\mu_0|y-p|}\le e^{\mp2w_p}\mathsf g_\flat$, and the same with $\mu$ in place of $\mu_0$. With \textup{(n1)}, the first sum is thus at most $14\mathsf g_\flat|e^{-w}G|^2$ and the second at most $\mathsf g_\flat L^{-1}\beta[F]^2$, so that $|\langle G,\mathsf DF\rangle|\le\epsilon\sqrt{14}\,\mathsf g_\flat|e^{-w}G|\,L^{-1/2}\beta[F]$. Taking $G=e^{2w}\mathsf DF$, we obtain the claim.

\emph{Step (3): the minimiser under the weight.} We claim that every $X'$ satisfies $\beta[HX']^2\le\vartheta_0\,\hat n_{y',\eta_+}(X')^2$. For a coarse plaquette $P$ we put $\nu_P:=H\delta_P$ and $r_P:={\rm dist}(y',P)$. \Cref{Y-A-P28}\,(a) for the given rows $\ell\in\mathcal R(b_*,K_*,\lambda_*,\tau)$ with $\tau\le\frac19$, recalled in Step~1a of the proof of \Cref{Y-J-ThA} in the form $\sum_pe^{2\lambda_*\mathsf d(p,P)}(H^\sharp\delta_P)_p^2\le4\gamma_*^{-1}L^{-1}$, together with $w_p\le-\lambda_*r_P+\lambda_*\mathsf d(p,P)$, gives $|e^w\nu_P|\le2(\gamma_*L)^{-1/2}e^{-\lambda_*r_P}$.

As $|w_p-w_{p+e_\alpha}|\le\lambda_*/L$, we have $|e^w\partial_\alpha\nu_P|\le(1+e^{\lambda_*/L})|e^w\nu_P|$, and with $1+e^{\lambda_*/L}\le2.1$ this yields $|e^w\nabla\nu_P|\le\sqrt3\cdot2.1\cdot2(\gamma_*L)^{-1/2}e^{-\lambda_*r_P}\le7.3(\gamma_*L)^{-1/2}e^{-\lambda_*r_P}$. A second bound comes from \Cref{Y-J-P43}, which applies since $\mathsf B'\le\frac32\mathsf r_\sharp^{-1}\le b_*$: as $w\le0$, it gives $|e^w\nabla\nu_P|\le2\sqrt6\,b_*L^{-3/2}$. Taking the geometric mean of the two bounds, we find $|e^w\nabla\nu_P|\le6(b_*\gamma_*^{-1/2})^{1/2}L^{-1}e^{-\frac12\lambda_*r_P}$; here we use $2\sqrt6\cdot7.3<35.8<36$.

By \Cref{Y-J-L44}\,(a), $\nabla Hh'=0$ for the constant form $h'=h'^{(y')}$ of $X'$, that is, for the constant form on $\mathbb T'$ with the components of $X'$ at $y'$. Hence Claim~1 of \Cref{Y-J-L46}, used on $\mathbb T'$ with $\theta=\frac12\lambda_*=4\mu$, gives \[|e^w\nabla HX'|\le\sum_P|e^w\nabla\nu_P||X'_P-h'_P|\le\sqrt3\cdot6(b_*\gamma_*^{-1/2})^{1/2}L^{-1}\mathsf g(\frac{\lambda_*}4)^{3/2}n_{y'}(X'),\] and so $L|e^w\nabla HX'|^2\le\vartheta_0n_{y'}(X')^2$. For the second part of $\beta$, the bound on $|e^w\nu_P|$ and the Cauchy--Schwarz inequality give $L^{1/2}|e^wHX'|\le2\gamma_*^{-1/2}\sum_Pe^{-\lambda_*r_P}|X'_P|\le2\gamma_*^{-1/2}(3\mathsf g(2\lambda_*-2\mu_0))^{1/2}|X'|_{y'}$, hence $L|e^wHX'|^2\le C_0|X'|^2_{y'}$, as $\mathsf g(2\lambda_*-2\mu_0)\le \mathsf g(\lambda_*)$. Since $C_0\eta^2\le c_0\eta^2=\vartheta_0L\eta^2$ and $\eta_+^2=L\eta^2$, the two parts add up to the claim.

\emph{Step (4): the algebra.} Put $G:=H-L^{-1}W^*$ and $\Gamma:=\mathsf N^\sharp-(aL)^{-1}\mathrm{d}'^*$, and let $\tilde G$, $\tilde\Gamma$ and $\tilde H$ be the same operators defined from $\tilde\ell$, that is, with $\tilde H^\sharp$ and $\tilde{\mathsf N}^\sharp$. We set \[\tilde u(X'):=\bigl(L^{1/2}\mathbf 1_{B(y')}\tilde GX',\ (aL)^{1/2}\mathbf 1_{y'}\tilde\Gamma X',\ L^{1/2}(\tilde\ell_y\tilde HX')_{y\in B(y')}\bigr),\] which is the part of $\tilde\ell'_{y'}X'$ consisting of the groups \textup{(3)}, \textup{(4)} and \textup{(5)}, and \[\Delta u(X'):=\bigl(L^{1/2}\mathbf 1_{B(y')}\Delta HX',\ (aL)^{1/2}\mathbf 1_{y'}\Delta\mathsf NX',\ L^{1/2}(\tilde\ell_y\Delta HX')_{y\in B(y')}\bigr),\qquad v:=\tilde u-\Delta u.\] The groups \textup{(1)} and \textup{(2)} are the same for $\ell'_{y'}$ and for $\tilde\ell'_{y'}$.

The vector $v(X')$ consists of the groups \textup{(3)} and \textup{(4)} of $\ell'_{y'}X'$ and of $L^{1/2}(\tilde\ell_yHX')_{y\in B(y')}$. Since $\langle\tilde\ell_yF,\tilde\ell_yF'\rangle=\langle\ell_yF,\ell_yF'\rangle+\langle F,D_yF'\rangle$, we therefore find that $\langle X',\mathsf E'_{y'}X''\rangle$ is the sum of three terms: the contribution of the groups \textup{(1)} and \textup{(2)}, the term $\langle v(X'),v(X'')\rangle$ and the term $-L\sum_{y\in B(y')}\langle HX',D_yHX''\rangle$. On the other hand, $\langle X',\tilde{\mathsf E}'_{y'}X''\rangle$ is the same contribution plus $\langle\tilde u(X'),\tilde u(X'')\rangle$. With $D'_{y'}:=\tilde{\mathsf E}'_{y'}-\mathsf E'_{y'}$ and $v=\tilde u-\Delta u$ we obtain \[\langle X',D'_{y'}X''\rangle=L\!\!\sum_{y\in B(y')}\!\!\langle HX',D_yHX''\rangle+\langle\Delta u(X'),\tilde u(X'')\rangle+\langle\tilde u(X'),\Delta u(X'')\rangle-\langle\Delta u(X'),\Delta u(X'')\rangle .\]

\emph{Step (5): the three bounds.} By the definition of $\epsilon$, the Cauchy--Schwarz inequality and Step~(1), the first term on the right satisfies $\bigl|L\sum_{y\in B(y')}\langle HX',D_yHX''\rangle\bigr|\le\epsilon L\sum_{B(y')}\hat n_y(HX')\hat n_y(HX'')\le1.1\mathsf g_\flat\epsilon\,\beta[HX']\beta[HX'']$.

The vector $\Delta u(X')$ is $L^{1/2}$ times the rows of the vector of \Cref{Y-J-L47}\,(b) that are placed in $B(y')$ or at $y'$, and on these rows $e^w\ge e^{-\lambda_*/2}$. By \Cref{Y-A-L25}\,(b) and~(c), the bound \textup{($\sinh$)} of \Cref{Y-J-L41} holds for $\tilde{\mathbb A}^{\rm e}$ at $\lambda_*$, so that \Cref{Y-J-L41}\,(c) applies with our weight $w$; followed by Step~(2), it gives $|\Delta u(X')|\le L^{1/2}e^{\lambda_*/2}\,4\,|e^w\mathsf DHX'|\le15.1\,\mathsf g_\flat\,\epsilon\,\beta[HX']$. The factor $15.1$ rests on $4\sqrt{14}<14.97$ and on $\lambda_*<10^{-13}$, which follows from the interval for $\lambda_*$ in \Cref{Y-A-Fr}\,\textup{(0)}; together they give $e^{\lambda_*/2}\cdot4\sqrt{14}\le14.97\cdot1.001\le15.1$.

Finally $|\tilde u(X'')|\le|\tilde\ell'_{y'}X''|\le\sqrt2\,C_1\hat n_{y',\eta_+}(X'')$. The first inequality holds because $\tilde u(X'')$ is a part of $\tilde\ell'_{y'}X''$. The second is Claim~3 of \Cref{Y-J-L46} on $\mathbb T'$, which applies by \Cref{Y-A-Th29}\,(a), recalled in Step~1a of the proof of \Cref{Y-J-ThA} in the form $|\tilde\ell'_{y'}(P)|\le K_*e^{-\lambda_*{\rm dist}(y',P)}$, and because $\tau(\tilde\ell')=\tau(\tilde\ell)\le C_1\eta\le C_1\eta_+$, where the equality is \Cref{Y-J-L44}\,(b).

\emph{Step (6): the sum.} Let $\mathsf y:=\mathsf g_\flat C_1\vartheta_0^{1/2}=\vartheta_1/48\le\frac1{100}$. We insert the three bounds of Step~(5) into the identity of Step~(4), and we use Step~(3) and $\epsilon\le2C_1^2$; this gives \[|\langle X',D'_{y'}X''\rangle|\le\epsilon\,[1.1\mathsf g_\flat\vartheta_0+2\cdot15.1\sqrt2\,\mathsf y+15.1^2\cdot2\,\mathsf y^2]\,\hat n_{y',\eta_+}(X')\hat n_{y',\eta_+}(X'').\] The bracket is $\le(0.02+42.8+4.6)\mathsf y\le48\mathsf y=\vartheta_1$, and \textup{(a)} is proved.

\emph{Step (7): the cut.} For the cut families, $D'^{\rm c}_{y'}=\chi D'_{y'}\chi$ with $\chi=\chi_{y'}$. We say that a coarse plaquette is inside if $\chi$ takes the value $1$ on it, and outside otherwise. Then $(\partial_\alpha(\chi X'))_P$ equals $(\partial_\alpha X')_P$ if $P$ and $P+e_\alpha$ are both inside, $0$ if both are outside, and, in the remaining case, where exactly one of the two is inside, $-X'_P$ if this one is $P$ and $X'_{P+e_\alpha}$ if it is $P+e_\alpha$.

In the remaining case, the one of the two plaquettes that is inside has distance exactly $\varrho$ from $y'$, since the periodic sup distance takes integer values and varies by at most $1$ from a plaquette to its neighbour; moreover $|y'-P|\ge\varrho$, and each plaquette that is inside occurs in at most six such terms. Hence \[n_{y'}(\chi X')^2\le n_{y'}(X')^2+6e^{-2\mu\varrho}\sum_{|y'-P|=\varrho}X_P'^2\le n_{y'}(X')^2+6e^{-2(\mu-\mu_0)\varrho}|X'|_{y'}^2,\] and since $|\chi X'|_{y'}\le|X'|_{y'}$, $2(\mu-\mu_0)=\lambda_*/8$ and $\eta_+^2=L\eta^2$, we get $\hat n_{y',\eta_+}(\chi X')\le\hat n_{y',\eta'}(X')$. Part \textup{(a)}, applied to $\chi X'$ and $\chi X''$, now gives \textup{(b)}. \end{proof}

\begin{theorem}[Contraction of the site forms along two chains]\label{Y-J-Th49} Let $L\ge L_B$ be odd, $aL\ge3/c_L$, $n$ even, $0\le k\le k'$ and $m:=k'-k$. Let $(\ell_j)_{j\le k}$ and $(\tilde\ell_j)_{j\le k'}$ be two chains that end on $\mathbb Z^3/n\mathbb Z^3$, with radii $(\varrho_j)$ and $(\tilde\varrho_j)$ that both satisfy the condition \textup{(Cut)}, that is, $\sum_{j\le k}\mathsf t_*(\varrho_j)\le\frac19$ and $\sum_{j\le k'}\mathsf t_*(\tilde\varrho_j)\le\frac19$, and with equal radii at equal distance from the end, $\tilde\varrho_{m+j}=\varrho_j$ for $1\le j\le k$.

Put \[\eta_0:=C_1^{-1}\sum_{j\le m}\mathsf t_*(\tilde\varrho_j),\qquad\eta_i^2:=L\,\eta_{i-1}^2+6\,e^{-\lambda_*\varrho_i/8}\quad(1\le i\le k),\] and let the radii satisfy \begin{equation}\tag{Bud}\eta_k\le1 .\end{equation} Then for $0\le i\le k$, on $\mathbb Z^3/nL^{k-i}\mathbb Z^3$, $\|\tilde{\mathsf E}_{m+i}-\mathsf E_i\|_{\eta_i}\le2C_1^2\,\vartheta_1^{\,i}$. In particular, for all $y,p,p'$, \[\bigl|(\mathsf E_{k,y}-\tilde{\mathsf E}_{k',y})(p,p')\bigr|\le15\,C_1^2\,\vartheta_1^{\,k}\,e^{-\lambda_\flat(|y-p|+|y-p'|)},\qquad\vartheta_1=c_{\rm fix}L^{-1/2}\le L^{-2/5},\quad\lambda_\flat=\tfrac1{16}\lambda_* .\]\end{theorem}

\begin{proof}[Proof of \Cref{Y-J-Th49}] By \Cref{Y-A-Th29}\,(b), together with the bound on $\tau$ established in its proof, all rows of both chains lie in $\mathcal R(b_*,K_*,\lambda_*,\frac19)$, with $\tau(\ell_i)\le\sum_{j\le i}\mathsf t_*(\varrho_j)$ and $\tau(\tilde\ell_{m+i})\le T_i:=C_1\eta_0+\sum_{j\le i}\mathsf t_*(\varrho_j)$; for the second chain we use $\tilde\varrho_{m+j}=\varrho_j$.

We show by induction on $i$ the bound \textup{(I$_i$)}: $T_i\le C_1\eta_i$. For $i=0$ it holds by the definition of $\eta_0$. If \textup{(I$_i$)} holds, then \[(C_1\eta_i+\mathsf t_*(\varrho_{i+1}))^2\le2C_1^2\eta_i^2+2\mathsf t_*^2\le C_1^2(L\eta_i^2+6e^{-\lambda_*\varrho_{i+1}/8}),\] by $L\ge2$ and by $2\mathsf t_*(\varrho)^2=18K_*^2\mathsf g(\frac{\lambda_*}2)^2e^{-\lambda_*\varrho}\le6C_1^2e^{-\lambda_*\varrho/8}$, where $6C_1^2=18K_*^2\mathsf g(\frac{\lambda_*}2)^3$. As $T_{i+1}=T_i+\mathsf t_*(\varrho_{i+1})$, this is \textup{(I$_{i+1}$)}. Since the $\eta_i$ increase with $i$, \textup{(Bud)} gives $\eta_i\le\eta_k\le1$.

Now we prove the bound of the theorem by induction on $i$. At the start, $\mathsf E_0=0$ because $\ell_0=\emptyset$, and Claim~3 of \Cref{Y-J-L46}, which applies as $\tau(\tilde\ell_m)\le T_0=C_1\eta_0$, gives $|\langle F,\tilde{\mathsf E}_{m,y}G\rangle|\le|\tilde\ell_{m,y}F||\tilde\ell_{m,y}G|\le2C_1^2\hat n_{y,\eta_0}(F)\hat n_{y,\eta_0}(G)$; this is the case $i=0$.

For the step from $i$ to $i+1\le k$ we apply \Cref{Y-J-P48}\,(b) to the pair $(\ell_i,\tilde\ell_{m+i})$, with $\eta=\eta_i$, $\varrho=\varrho_{i+1}$ and $\tau=\min\{C_1\eta_i,\frac19\}$; here the letter $N$ of the proposition is $nL^{k-i-1}$. Its hypotheses hold by the first paragraph of the proof, by \textup{(I$_i$)}, by $\eta_i\le1$ and by the bound of the theorem for the index $i$, which is the hypothesis of the induction and gives $\epsilon\le2C_1^2\vartheta_1^{\,i}\le2C_1^2$. Since $\eta'=\eta_{i+1}$ and both families are cut at $\varrho_{i+1}=\tilde\varrho_{m+i+1}$, the proposition yields this bound for the index $i+1$. Finally, as $\eta_k\le1$, we may use \textup{(n2)} with $\eta=\eta_k$; it gives the bound on the entries, because $7.1\cdot2\le15$. \end{proof}

\begin{corollary}[Row families of two runs approach each other]\label{Y-J-Co50}
Let $L\ge L_B$ be odd, let $aL\ge3/c_L$ and $\theta_D\le\frac18\lambda_*$, and put $c_R=64/\theta_D$, $\varrho(x)=\lceil c_R\log x\rceil$ and $e_R=\frac12\lambda_*c_R\ (\ge256)$. Let $x_\flat$ satisfy the condition \textup{(X12)} of the entry \textup{(O13)}, together with \textup{(X13)}: $16x_\flat^{-e_R/4}\le1$. Let $\mathsf b\ge x_\flat$, let $K\le K'$ with $m=K'-K$, and let $N$ be even. Let $\ell_k$ be the chain on the tower $\mathbb Z^3/NL^{K-k}\mathbb Z^3$ with the radii $\varrho_j=\varrho(L^{K-j}\mathsf b)$, and let $\tilde\ell_k$ be the chain of length $K'$ ending on $\mathbb Z^3/N\mathbb Z^3$, with the radii $\tilde\varrho_j=\varrho(L^{K'-j}\mathsf b)$. Then for $0\le k\le K$, with $x_k=L^{K-k}\mathsf b$, the condition \textup{(Bud)} holds with $\eta_k^2\le16x_k^{-e_R/4}$, we have $\|\tilde{\mathsf E}_{m+k}-\mathsf E_k\|_{\eta_k}\le2C_1^2\vartheta_1^{\,k}$, and \begin{equation}\tag{RU}\|\tilde{\mathsf E}_{m+k}-\mathsf E_k\|_{(\lambda_\flat)}\le15C_1^2\vartheta_1^{\,k}.\end{equation}
\end{corollary}

\begin{proof}[Proof of \Cref{Y-J-Co50}]
At partner levels the radii of the two chains agree: $\tilde\varrho_{m+j}=\varrho_j$. The condition \textup{(Cut)} follows from \textup{(X12)}, by the argument that the proof of \Cref{Y-A-Th36} gives for \textup{(Cut)}. Next, $\mathsf t_*(\varrho(x))\le3K_*\mathsf g(\frac{\lambda_*}2)x^{-e_R}$ and $\sum_{j\le m}(L^{K'-j}\mathsf b)^{-e_R}\le(L^K\mathsf b)^{-e_R}/(1-L^{-e_R})$; together with $3K_*\mathsf g(\frac{\lambda_*}2)\le C_1$, these bounds yield $\eta_0\le2(L^K\mathsf b)^{-e_R}$. Further, $e^{-\lambda_*\varrho_j/8}\le x_j^{-e_R/4}$, where $x_j=L^{k-j}x_k$.

Therefore \[\eta_k^2=L^k\eta_0^2+6\sum_{j\le k}L^{k-j}e^{-\lambda_*\varrho_j/8}\le4L^{k-2e_RK}\mathsf b^{-2e_R}+6x_k^{-e_R/4}\sum_iL^{i(1-e_R/4)}\le4x_k^{-e_R/4}+12x_k^{-e_R/4},\] where the last inequality rests on $k+\frac14(K-k)e_R\le2e_RK$ and on $\mathsf b>1$. This is the bound on $\eta_k^2$ in the statement. As $x_k\ge\mathsf b\ge x_\flat$, we have $16x_k^{-e_R/4}\le16x_\flat^{-e_R/4}\le1$ by \textup{(X13)}, so that $\eta_k\le1$, which is \textup{(Bud)} for the two chains cut off at the level $k$. \Cref{Y-J-Th49}, applied to these two chains, now gives the two bounds on $\tilde{\mathsf E}_{m+k}-\mathsf E_k$.
\end{proof}

\emph{Step 2: partners, and the data of a step at each level.} In this step of the proof of \Cref{Y-J-ThA}, under the hypotheses and in the notation of that theorem, we claim that the following clauses \textup{(a)} to \textup{(e)} hold.
\begin{enumerate}
\item[\textup{(a)}] Level $k$ of the run $K$ and level $m+k$ of the run $K'$ share one torus, one number $x_k$, one grid of cubes, one set of the numbers $\varepsilon(x_k),p_0(x_k),\varrho(x_k),r(x_k),s(x_k),\ell_w(x_k)$, the caps and the domains, one mask and one clip, and one step map.
\item[\textup{(b)}] \textup{(R8)} applies to both runs.
\item[\textup{(c)}] For $0\le k\le K$, $\ell_k\in\mathcal R(b_*,K_*,\lambda_*,\frac19)$, and $\ell_k$ has range $\le\varrho(x_k)$.
\item[\textup{(d)}] For $1\le k\le K-1$, \textup{(SD)} holds at $(x_k,\ N'=NL^{K-k-1},\ \ell_k)$.
\item[\textup{(e)}] For $0\le k\le K$, $NL^{K-k}\ge2M\,s(x_k)$.
\end{enumerate}

Taken by themselves, and with $K\ge1$ in place of the hypothesis $K\ge K_0(\mathsf b)$ of \Cref{Y-J-ThA}, the clauses \textup{(c)} and \textup{(e)} are a statement in its own right. Among the letters, which are chosen in the order \textup{(O0)} to \textup{(O13)}, its data are only these: $L\ge L_A$ odd, $aL\ge3/c_L$, $x_\flat\ge e^2$ with \textup{(X12)}; $\mathsf b\ge x_\flat$; $N$ even, $M\mid N$, $N\ge2M\,s(L\mathsf b)$; $K\ge1$; $x_k=L^{K-k}\mathsf b$, and the rows $\ell_k=\ell_k^{[K]}$ of the runs (item~2 of \S\ref{Y-J-1}). With these data, \textup{(c)} and \textup{(e)} hold for $0\le k\le K$. Clause \textup{(a)} is a list of identities between level $k$ of the run $K$ and level $m+k$ of the run $K'=K+m$; it holds for every $K'\ge K\ge1$ as well. The arguments for \textup{(c)} and \textup{(e)} below rest on these data alone: they are independent of $K_0(\mathsf b)$, of $k_0$, of $\Lambda$, of $K'$, of the members, and of Steps 3 to 6.

For \textup{(a)}: the tori agree, $\mathbb Z^3/NL^{K'-(m+k)}\mathbb Z^3=\mathbb Z^3/NL^{K-k}\mathbb Z^3$, and $L^{K'-(m+k)}\mathsf b=L^{K-k}\mathsf b=x_k$. The origins are defined recursively, downwards from the last level $j=K$ of the run, by $o_K=0$ and $o_j=L\,o_{j+1}-(L-1)/2$; hence $o^{[K']}_{m+k}=o^{[K]}_k$. Each number of the list is a function of $x$ and of the letters, and the mask, the clip and the paths depend on $V,\varepsilon,\rho_0,\varrho(x)$ (the definition ``Energy'' in item~2 of \S\ref{Y-J-1}). The step depends on $x_k,a,L$ alone and is independent of the rows, of the grid and of the members (the definition ``The step'' in item~2 of \S\ref{Y-J-1}). So it is one map for both runs, and $(\rho^{[K']}_{m+k})'=\rho^{[K']}_{m+k+1}$.

For \textup{(b)}: $\beta=L^K\mathsf b\ge L^Kx_\flat$ and $L^{1-K}\beta=L\mathsf b$, and the same holds with $K'$ in place of $K$. The condition $N\ge2M\,s(L^{1-K}\beta)$ is thus \textup{(Tor)} of \Cref{Y-J-ThA}, for both runs.

For \textup{(c)} we begin with \textup{(Cut)}. From $\varrho(x)\ge c_R\log x$ we obtain $\mathsf t_*(\varrho(x_j))\le3K_*\mathsf g(\lambda_*/2)x_j^{-e_R}$; as $x_j\ge L^{K-j}x_\flat$, the geometric series and \textup{(X12)} give $\sum_{j=1}^K\mathsf t_*(\varrho(x_j))\le\frac19$. \Cref{Y-A-Th29}\,(b), whose data are $L\ge L_A$ odd, $aL\ge3/c_L$ and $N$ even, then yields $\ell_k\in\mathcal R(b_*,K_*,\lambda_*,\frac19)$ for $0\le k\le K$, and the range is $\le\varrho(x_k)$ because of the cut.

For \textup{(d)}, let $1\le k\le K-1$. Then $x_k\ge L\mathsf b\ge L\,x_\flat\ge x_9$, and $N'=NL^{K-k-1}$ is even and a multiple of $M$. As to \textup{(Tor)}: \Cref{Y-GB-P63}, one of the base clauses \textup{(RB)}, applied with $x=L\mathsf b\ge e^2$ and $j=K-1-k$, turns $N\ge2M\,s(L\mathsf b)$ into $NL^j\ge2M\,s(L^j\cdot L\mathsf b)=2M\,s(x_k)$. The grids fit, by their definition.

The four conclusions for the pair $(\ell_k,(\ell_k)')$ are provided by \Cref{Y-A-Th36} at this $k$, since its two hypotheses on $k$ and on $x_\flat$, namely $x_{k+1}\ge\mathsf b\ge x_\flat$ and \textup{(X12)}, hold. Applied with $(C_D,\theta_D)$ of \textup{(O6)} as its pair of constants, the theorem gives them in the following sense. The bound \textup{(Dec)}$(C_D,\theta_D)$ at $\theta_D=\frac12\min\{\lambda_2^*,\theta_{\mathbb C},\theta_J\}$ is the conclusion \textup{(Fr2)} of \Cref{Y-A-Fr}, which covers all sixteen kernels of \textup{(Dec)}; the hypotheses of that theorem, $L\ge L_A$ odd, $aL\ge3/c_L$, $N$ even and $x_{k+1}\ge x_\flat$ with \textup{(X12)}, are met. The condition \textup{(N$'$)} with \eqref{eq:row-cut-norms} at every cut, and the bound $\mathsf t_\rho\le\frac18(\frac12-d_0-c_\Xi)$, are \textup{(Fr3$'$)} and \textup{(Fr4)} in the conclusion of \Cref{Y-A-Fr}. The properties of $\ell_k$ and $(\ell_k)'$ that these two clauses use are the ones stated in \textup{(Fr0)} of that theorem. In the proof of \Cref{Y-A-Th36}, the condition and the bound are derived from $\ell_k\in\mathcal R(b_*,K_*,\lambda_*,\frac19)$ and $|(\ell_k)'_{y'}(P)|\le K_*e^{-\lambda_*{\rm dist}(y',P)}$, which \Cref{Y-A-Th29}\,(b) gives; from the radius $\rho_0$ fixed by \eqref{eq:row-mask-radius}; and from \Cref{Y-A-Th33}\,(d) with $(b,K,\lambda,\lambda^\circ)=(b_*,K_*,\lambda_*,\lambda_*)$, as in \Cref{Y-A-Stars}, hence with the constants $c_1^*$ and $c_2^*$. This derivation depends on $\theta_D$ only through \textup{(Cut)}, which \textup{(c)} gives when $c_R$ has the value fixed by \textup{(O6)}.

The hypothesis $L\ge L_*(d_0,c_\Xi)$ of \Cref{Y-A-Th36} follows from \textup{(O4)}. With $\mathsf N_\nu=\mathsf N_\sharp+(3+4c_2^*L)\mathsf N_Z$, $\mathsf N_\sharp=c_1^*\mathsf l_L\mathsf N_Z\mathsf P_\rho$ and $\mathsf N_\Xi=c_2^*\mathsf N_Z$, the left side of \textup{(L1$^{\rm su}$)} equals \[[c_1^*\mathsf l_L\mathsf P_\rho/L+3/L+2200+44c_2^*]\mathsf N_Z\,L^{-1/2}\ge(c_1^*\mathsf l_L\mathsf P_\rho/L+60+4c_2^*)\mathsf N_Z\,L^{-1/2}.\] From \textup{(L2$^{\rm su}$)} we get $L\ge1.44\cdot10^8\cdot17^4(\mathsf N_\sharp+1)^4\ge115\mathsf N_\sharp^2$. The definition of $L_*(d_0,c_\Xi)$, in \Cref{Y-A-OrderLetters}, asks for its two inequalities at all real $L\ge L_*(d_0,c_\Xi)$. For each of the two, the left side divided by the right side decreases in the real variable $L>e$: $\mathsf l_L/L$ has the derivative $-\log L/L^2$, $L^{-1/2}$ decreases, and $\mathsf l_L^2/L$ has the derivative $\mathsf l_L(2-\mathsf l_L)/L^2<0$ for $\mathsf l_L>2$. The two bounds deduced from \textup{(L1$^{\rm su}$)} and \textup{(L2$^{\rm su}$)} show that both inequalities hold at the $L$ of \textup{(O4)}, where \textup{(L1$^{\rm su}$)} and \textup{(L2$^{\rm su}$)} hold; by this monotonicity they hold at every real number $\ge L$. As $L$ is an integer $\ge L_A$, it follows that $L_*(d_0,c_\Xi)\le L$. The conditions \textup{(L1)} and \textup{(L2)} hold in the forms \textup{(L1$^{\rm su}$)} and \textup{(L2$^{\rm su}$)}, by \textup{(O4)}.

For \textup{(e)}: as $r(x)=\lceil20c_R\log x/M\rceil$ is the ceiling of an increasing function, the function $x\mapsto s(x)=8r(x)$ is non-decreasing, so $N\ge2M\,s(L\mathsf b)\ge2M\,s(\mathsf b)$, and \Cref{Y-GB-P63}, applied with $x=\mathsf b\ge e^2$ and $j=K-k$, gives $NL^{K-k}\ge2M\,s(x_k)$.

\emph{Step 3: \Cref{Y-J-ThA}\,(i), and the size of $\delta_\Delta$.} We claim that $\Lambda>0$, that $k_0\ge1$, and that $k_0\le K-1$ holds if and only if $K\ge K_0(\mathsf b)$. As $K\ge K_0(\mathsf b)$ is a hypothesis of \Cref{Y-J-ThA}, this gives $1\le k_0\le K-1$, which, together with $\Lambda>0$, makes up clause \textup{(i)} of that theorem. Moreover, at every level $k$ with $k_0\le k\le K$ \[\delta_\Delta(k):=(c_Q^{\rm su}/E_a)\,x_k^{1/2}\delta_k/p_0(x_k)\le L^{\Lambda+K/2-9k/10}\le1.\]

First, $\Lambda>0$. By \textup{(X17)} at $x=\mathsf b$ we have $16c_a(\mathsf b)\le\frac12$, that is, $\mathsf b^{1/2}\ge32E_ap_0(\mathsf b)^3$, and the right side is $\ge32E_a$ by \textup{(X14)}. Hence $c_4\mathsf b^{1/2}\ge480c_Q^{\rm su}C_1^2>1$, by Step~0\,(iii). Next, $k_0$ is a level: as $(5K+10\Lambda)/9>0$, we have $k_0\ge1$. As $K-1$ is an integer, $\lceil y\rceil\le K-1$ holds if and only if $y\le K-1$; and $(5K+10\Lambda)/9\le K-1$ holds if and only if $4K\ge10\Lambda+9$, that is, if and only if $K\ge\frac52\Lambda+\frac94$.

We turn to $\delta_\Delta$. By Step~1 and by \textup{(X14)} at $x_k\ge\mathsf b\ge x_\flat$, \[\delta_\Delta(k)\le(c_Q^{\rm su}/E_a)\mathsf b^{1/2}L^{(K-k)/2}\cdot15C_1^2L^{-2k/5}=c_4\mathsf b^{1/2}L^{K/2-9k/10}=L^{\Lambda+K/2-9k/10}.\] The exponent is $\le0$ if and only if $k\ge(5K+10\Lambda)/9$, and hence for every $k\ge k_0$.

\emph{Step 4: the switch at the step $k_0\to k_0+1$, following the argument of \Cref{Y-I-ThC}.} We show three things. First, for some real number $\mathsf V$ and some real $\mathfrak n_{k_0+1}\in M_1^{\rm su}(x_{k_0+1})$ having the strong cap, the identity $\rho^{[K']}_{m+k_0+1}=e^{\mathsf V}\rho[\mathfrak n_{k_0+1};\ell_{k_0+1}^{[K]}]$ holds at every field. Second, $u_{k_0+1}=F_{\ell_{k_0}}(u_{k_0})$ is a system of members of $\rho^{[K]}_{k_0+1}$. Third, $\mathsf D_{x_{k_0+1}}(u_{k_0+1},\mathfrak n_{k_0+1})\le0.4(2+\delta_\Delta)\le1.2\le3\vartheta$.

We put $x=x_{k_0}$, $\ell=\ell_{k_0}^{[K]}$ and $\tilde\ell=\ell_{m+k_0}^{[K']}$. Both row families live on $\mathbb T_{k_0}$ and have the same grid of cubes (Step~2\,(a)).

\textup{1.} \emph{Members.} We have $1\le k_0\le K$ and $1\le m+k_0\le K'$. By \textup{(R8)} for both runs (Step~2\,(b)), $\rho^{[K]}_{k_0}=e^{\mathsf V_1}\rho[u_{k_0};\ell]$ and $\rho^{[K']}_{m+k_0}=e^{\mathsf V_2}\rho[\tilde u;\tilde\ell]$, where the numbers $\mathsf V_i$ are real and $u_{k_0},\tilde u$ are real elements of $M_1^{\rm su}(x)$ with the strong cap.

\textup{2.} \emph{The rows are changed.} We go through the conditions of \textup{(R11)} at $x$ in turn. First, $x\ge e^2$. The side is $N_f=NL^{K-k_0}=N'L$ with $N'=NL^{K-k_0-1}$, as $k_0\le K-1$ by Step~3; and $N'$ is even, a multiple of $M$, and $\ge2M\,s(x)$, by Step~2\,(d) at $k=k_0$. The families $\ell$ and $\tilde\ell$ have range $\le\varrho(x)$; for $\tilde\ell$ this holds because it is cut at $\varrho(L^{K'-m-k_0}\mathsf b)=\varrho(x)$. Further, $\delta_\Delta\le1$ by Step~3 at $k=k_0$; \textup{(M5r)} holds by \textup{(O7)}; $\varepsilon\le\frac1{40}$ is \textup{(X16)}; $\varrho+2\le M\,r$ is Step~0\,(ii); $c_\Xi p_0^2\ge1$ is \textup{(X15)}; and $(2r+1)^3c_a\le34$ is Step~0\,(vi). Hence \textup{(R11)} provides a real $u^\flat\in M_2^{\rm su}(x)$ with $\rho[\tilde u;\tilde\ell]=\rho[u^\flat;\ell]$ at every $V$ and $\mathsf D^{\rm in}_x(\tilde u,u^\flat)\le\delta_\Delta$. Thus $\rho^{[K']}_{m+k_0}=e^{\mathsf V_2}\rho[u^\flat;\ell]$ as functions.

\textup{3.} \emph{The step.} By Step~2\,(d), as $k_0\le K-1$, \textup{(SD)} holds at $(x,NL^{K-k_0-1},\ell)$. The step is linear in $\rho$, and it is one map for both runs. By \textup{(9a)} at $A=2$, \[\rho^{[K']}_{m+k_0+1}=e^{\mathsf V_2}(\rho[u^\flat;\ell])'=e^{\mathsf V_2+\mathsf V'}\rho[F_\ell(u^\flat);\ell'^{\rm c}],\] and $\ell'^{\rm c}=((\ell_{k_0})')^{\rm c}$, cut at $\varrho(x_{k_0+1})$, is $\ell_{k_0+1}^{[K]}$, the row family of the short run. We put $\mathfrak n_{k_0+1}:=F_\ell(u^\flat)$. By \textup{(9b)} at $A=2$, it lies in $M_1^{\rm su}(x_{k_0+1})$ and has the strong cap; by \textup{(9d)}, it is real, and so is $\mathsf V'$. Also by \textup{(9d)}, $\rho^{[K']}_{m+k_0+1}>0$, because $\rho[u^\flat;\ell]=e^{-\mathsf V_2}\rho^{[K']}_{m+k_0}>0$. So $\rho^{[K']}_{m+k_0+1}$ is a point of $C^{\rm su}_\flat(x_{k_0+1};\ell_{k_0+1}^{[K]})$. The clauses \textup{(9a)}, \textup{(9b)} and \textup{(9d)}, now taken at $A=1$ and applied to the system $u_{k_0}\in M_1^{\rm su}(x)$ of $\rho^{[K]}_{k_0}$, give the claim on $u_{k_0+1}$.

\textup{4.} \emph{The distance.} As $\|\cdot\|_x$ is a norm, \[\mathsf D^{\rm in}_x(u_{k_0},u^\flat)\le\|u_{k_0}\|_x+\|\tilde u\|_x+\mathsf D^{\rm in}_x(\tilde u,u^\flat)\le2+\delta_\Delta\le3.\] Both systems lie in $M_2^{\rm su}(x)$, so \textup{(9e)} applies; it gives $\mathsf D_{x_{k_0+1}}(u_{k_0+1},\mathfrak n_{k_0+1})\le0.4\cdot3=1.2\le3\vartheta=1.3404$.

\begin{remark}
By \textup{(R11)}, the system $u^\flat$ lies in $M_2^{\rm su}(x)$, that is, it satisfies the three size bounds with $A=2$; and $w^\flat$ has the weak cap only. The definition of the class asks for a system of members in $M_1^{\rm su}(x)$ with the strong cap. Hence \textup{(R11)} does not provide $u^\flat$ as such a system, nor does it give the membership of $\rho[u^\flat;\ell]$ in $C^{\rm su}_\flat(x;\ell)$. The property \textup{(S-step)} of \Cref{Y-G-ThA} concerns the images of points of the class alone. We therefore use \Cref{Y-H-C9} at $A=2$, which admits the weak cap, and we obtain reality from \Cref{Y-H-C9}\,(d).
\end{remark}

\emph{Step 5: after the switch, and the last level $k=K$; \Cref{Y-J-ThA}\,(ii), (iii), (iv) at $k$ and at $K$.} We claim that for $k_0<k\le K$ the systems $u_k$ and $\mathfrak n_k$ are real, lie in $M_1^{\rm su}(x_k)$, have the strong cap, and are systems of members of $\rho^{[K]}_k$ and of $\rho^{[K']}_{m+k}$ respectively, both over the same row family $\ell_k^{[K]}$, and that $\mathsf D_{x_k}(u_k,\mathfrak n_k)\le3\vartheta^{k-k_0}$. At $k=K$ this distance is $\le C_{\mathsf b}\theta^K$.

We argue by induction on $k$, starting from $k_0+1$, which is Step~4. Let the claim hold at some $k$ with $k_0<k<K$. By Step~2\,(d), \textup{(SD)} holds at $(x_k,NL^{K-k-1},\ell_k)$. We put $\mathfrak n_{k+1}:=F_{\ell_k}(\mathfrak n_k)$, and we have $u_{k+1}=F_{\ell_k}(u_k)$. By \textup{(9a)} at $A=1$, together with Step~2\,(a), \[\rho^{[K']}_{m+k+1}=(\rho^{[K']}_{m+k})'=e^{\mathsf V+\mathsf V'}\rho[\mathfrak n_{k+1};(\ell_k)'^{\rm c}],\] and $(\ell_k)'^{\rm c}=\ell_{k+1}^{[K]}$. By \textup{(9b)} at $A=1$, the sizes of $\mathfrak n_{k+1}$ are at most $0.598$ of the cap $c_a$ and at most the cap $c_h$; it has the strong cap, and it satisfies at $x_{k+1}$ all clauses of \textup{(ca)}, \textup{(ch)} and \textup{(c$\Xi$)} other than the size bounds. By \textup{(9d)}, $\mathfrak n_{k+1}$ and $\mathsf V'$ are real, and the image is $>0$. Applied to $u_k$ in place of $\mathfrak n_k$, the clauses \textup{(9a)}, \textup{(9b)}, \textup{(9d)} give these properties for $u_{k+1}$ as well. By \textup{(9e)}, \[\mathsf D_{x_{k+1}}(u_{k+1},\mathfrak n_{k+1})\le0.4\,\mathsf D^{\rm in}_{x_k}(u_k,\mathfrak n_k)\le\vartheta\,\mathsf D_{x_k}(u_k,\mathfrak n_k).\] Starting from $\mathsf D_{x_{k_0+1}}\le3\vartheta$, this gives $\mathsf D_{x_k}\le3\vartheta^{k-k_0}$.

\emph{The last level $k=K$.} Since $k_0<(5K+10\Lambda)/9+1$, we have $K-k_0>\frac49K-\frac{10}9\Lambda-1$. As $\vartheta<1$, the distance at the level $K$, where $x_K=\mathsf b$, satisfies \[\mathsf D_{\mathsf b}\le3\vartheta^{\frac49K-\frac{10}9\Lambda-1}=(3/\vartheta)\cdot\vartheta^{-\frac{10}9\Lambda}\cdot(\vartheta^{4/9})^K.\] Here \[\vartheta^{-\frac{10}9\Lambda}=\exp\{\frac{10}9\Lambda\log(1/\vartheta)\}=(L^\Lambda)^{\varepsilon_L}=(c_4\mathsf b^{1/2})^{\varepsilon_L},\] so that $\mathsf D_{\mathsf b}\le C_{\mathsf b}\theta^K$. Numerically, $0.4468^{4/9}\le0.69903$ and $3/0.4468\le6.715$.

\emph{Step 6: \Cref{Y-J-ThA}\,(iv) at the depth $s$, and (v).} We claim the following.
\begin{enumerate}
\item[\textup{(a)}] $K-s>k_0$ if and only if $K-s\ge K_0(L^s\mathsf b)=\frac52\Lambda+\frac54s+\frac94$; in that case $3\vartheta^{K-s-k_0}\le3\vartheta^{\frac49K-s-\frac{10}9\Lambda-1}=C_{L^s\mathsf b}\theta^{K-s}$.
\item[\textup{(b)}] $k_0<\varphi K+K_0(\mathsf b)$ and $k_0\le\lceil\varphi K\rceil+K_0(\mathsf b)$.
\end{enumerate}

For \textup{(a)}: as $K-s-1$ is an integer, $k_0\le K-s-1$ holds if and only if $(5K+10\Lambda)/9\le K-s-1$, that is, if and only if $4K\ge10\Lambda+9s+9$, that is, if and only if $K-s\ge\frac52\Lambda+\frac54s+\frac94$. Moreover $\Lambda(L^s\mathsf b)=\log_L(c_4L^{s/2}\mathsf b^{1/2})=\Lambda+s/2$, so that $K_0(L^s\mathsf b)=\frac52\Lambda+\frac54s+\frac94$. Further, $K-s-k_0>\frac49K-s-\frac{10}9\Lambda-1$; as $\vartheta<1$, this gives the inequality in \textup{(a)}. And \[C_{L^s\mathsf b}\theta^{K-s}=(3/\vartheta)\vartheta^{-\frac{10}9(\Lambda+s/2)}\vartheta^{\frac49(K-s)}=3\vartheta^{\frac49K-\frac49s-\frac59s-\frac{10}9\Lambda-1},\] with the same exponent.

For \textup{(b)}: $k_0<\varphi K+\frac{10}9\Lambda+1\le\varphi K+\frac52\Lambda+\frac94$. The inequality $\lceil\alpha+\beta\rceil\le\lceil\alpha\rceil+\lceil\beta\rceil$ gives $k_0\le\lceil\varphi K\rceil+\lceil\frac{10}9\Lambda\rceil\le\lceil\varphi K\rceil+\frac{10}9\Lambda+1\le\lceil\varphi K\rceil+K_0(\mathsf b)$. So $k>\varphi K+K_0(\mathsf b)$ implies $k>k_0$; and $k-k_0\ge k-\lceil\varphi K\rceil-K_0(\mathsf b)$, together with $\vartheta<1$, gives $3\vartheta^{k-k_0}\le3\vartheta^{k-\lceil\varphi K\rceil-K_0}$. For $K\le\frac94K_0(\mathsf b)$ the range of $k$ in \textup{(v)} is empty, and the claim at $k=K$ is \textup{(iv)}.
\end{proof}

\begin{proof}[Proof of \Cref{Y-J-ThB}]
\emph{Step 7: \Cref{Y-J-ThB}\,(a).} For $i=1,2$ let \[S_i(V):=\sum_{X\in\mathcal D}1_{X^{\approx}}(V)[E_i^a+E_i^h](X,V),\] and let $c_S:=1.22(c_a+c_h)(x)\le0.61$. We claim that $|S_i|\le c_Sn$ and $|S_1-S_2|\le c_Sn\mathsf D$ at every $V$, and that $\Xi_i=1$ where $D(V)$ is empty; together these give \textup{(a)}.

We have $\rho_i=e^{\mathsf V_i}e^{-x\mathrm{En}_\ell}e^{S_i}\Xi_i$, with one and the same $\mathrm{En}_\ell$ for $i=1,2$, as $\ell$ and $x$ are the same. The indicators $1_{X^{\approx}}$ depend on $V$ and $\varepsilon$ only, and so they are the same for $i=1,2$. Let $1_{X^{\approx}}(V)=1$. Then $V$, restricted to the bonds of $X^{\approx}$, takes its values in $G$ and is tame there, so that it belongs to $\mathcal U_X(x)$ (the definition ``Domain'' in item~2 of \S\ref{Y-J-1}; here $x>e$). Accordingly, at such $V$, \[|E_i^a+E_i^h|(X,V)\le(c_a+c_h)e^{-\kappa(\#X-1)}\] and \[|E_1^a-E_2^a|+|E_1^h-E_2^h|\le\mathsf D(c_a+c_h)e^{-\kappa(\#X-1)};\] for these two parts of the members, $\mathsf D^{\rm in}$ and $\mathsf D_x$ are the same quotients. Each $X$ contains a cube, so for the positive terms $e^{-\kappa(\#X-1)}$ we have $\sum_Xe^{-\kappa(\#X-1)}\le\sum_\square\sum_{X\ni\square}e^{-\kappa(\#X-1)}$, and \[\sum_{X\ni\square}e^{-\kappa(\#X-1)}\le\sum_{j\ge1}(26e)^{j-1}e^{-6(j-1)}=1/(1-26e^{-5})\le1.2124\le1.22;\] summed over the $n$ cubes, this gives the two bounds on $S_i$ and on $S_1-S_2$. Each $S_i$ is real, because $u_i$ is real.

If $D(V)$ is empty, no non-empty family is admitted, since each $X_i$ of a family must hold a plaquette of $D(V)$; the empty family gives $\Xi_i=1$. So there $\log\rho_1-\log\rho_2-(\mathsf V_1-\mathsf V_2)=S_1-S_2$.

\emph{Step 8: \Cref{Y-J-ThB}\,(b), the bound on the supremum.} We show that $|\Xi_i|\le e^{218n}$, that $|\Xi_1-\Xi_2|\le\mathsf Dne^{218n}$, and that $\sup_V|\rho_1-e^{\mathsf V_1-\mathsf V_2}\rho_2|\le e^{\mathsf V_1}2ne^{219n}\mathsf D$.

The sets of a family as in \textup{(c$\Xi$)} lie pairwise at a distance $>s(x)\ge8$, and so they are disjoint. We fix an order of the cubes and send each set of the family to its first cube; different sets go to different first cubes. By the strong cap, $|w_i(X,V)|\le e^{-\kappa(|X|_c-1)}$. Hence \[|\Xi_i|\le\prod_\square(1+\sum_{X\ni\square}e^{-\kappa(|X|_c-1)})\le\exp\{n\sum_{j\ge1}(52e)^je^{-6(j-1)}\},\] and the sum equals $52e/(1-52e^{-5})\le217.588\le218$.

By telescoping, \[|\prod_iw_1(X_i)-\prod_iw_2(X_i)|\le\sum_j|w_1-w_2|(X_j)\prod_{i\ne j}e^{-\kappa(|X_i|_c-1)}\le\mathsf D\cdot\#({\rm family})\cdot\prod_ie^{-\kappa(|X_i|_c-1)},\] where we use the bound $|w_1-w_2|(X,V)\le\mathsf De^{-\kappa(|X|_c-1)}$ that the weak distance provides. A family has at most $n$ sets. So $|\Xi_1-\Xi_2|\le\mathsf D\,n\,e^{218n}$.

By Step~7 and by $|e^\alpha-e^\beta|\le e^{\max\{\alpha,\beta\}}|\alpha-\beta|$, \[\FitDisplay{\dimexpr\linewidth-1pt\relax}{|e^{S_1}\Xi_1-e^{S_2}\Xi_2|\le|e^{S_1}-e^{S_2}||\Xi_1|+e^{S_2}|\Xi_1-\Xi_2|\le e^{0.61n}\cdot0.61n\mathsf D\cdot e^{218n}+e^{0.61n}\cdot n\mathsf D\cdot e^{218n}=1.61\,n\mathsf D\,e^{218.61n}\le2n\,e^{219n}\mathsf D.}\] Moreover $e^{-x\mathrm{En}_\ell}\le1$, because $S_\varepsilon\ge0$, $m_y\in\{0,1\}$, and the other terms are squares. Finally $\rho_1-e^{\mathsf V_1-\mathsf V_2}\rho_2=e^{\mathsf V_1}e^{-x\mathrm{En}_\ell}(e^{S_1}\Xi_1-e^{S_2}\Xi_2)$, and the bound on the supremum follows.

\emph{Step 9: \Cref{Y-J-ThB}\,(c), a lower bound on $\int\rho$ over a neighbourhood of $1$.} We claim that \[\int\rho_1dV\ge e^{\mathsf V_1}e^{-0.61n}e^{-\frac32\mathsf N_Zp_0^2N^3}(\varepsilon^3/384)^{3N^3},\] and that this gives the bound in $L^1$ with $C^{\rm su}_{\rm vol}$.

Let $\mathfrak B:=\{V:\operatorname{ang}V(b)\le\varepsilon/4\text{ at all }3N^3\text{ bonds}\}$.

\textup{(1)} \emph{Tame.} At $V\in\mathfrak B$, $\partial V(p)$ is the product of four $V(b)^{\pm1}$, so $\phi_p$ is at most the sum of four angles, which is $\le\varepsilon$; here we use the clause $\operatorname{ang}(gh)\le\operatorname{ang}g+\operatorname{ang}h$, $\operatorname{ang}g^{-1}=\operatorname{ang}g$ of \textup{(RB)}. Hence $D(V)$ is empty, $\Xi_1=1$, and $S_1\ge-0.61n$ by Step~7.

\textup{(2)} \emph{Mask and clip.} No plaquette has $\phi_p>\varepsilon$, and $\varepsilon\le1$. So $m_y=1$ for every $y$, $\zeta_p=F_p$ with $|F_p|=\phi_p\le\varepsilon$, and $S_\varepsilon(\phi_p)=\phi_p^2/2$.

\textup{(3)} \emph{Energy.} The paths $\Gamma(y,\pi(p))$ with $|y-p|\le\varrho(x)$ are defined, because $2(\varrho(x)+1)<2M\,s(x)\le N$ (Step~0\,(ii)). Here the rows are dressed, and the form differs from the scalar $\mathsf B$. At $G$-valued fields, $T$ is an isometry of $\mathfrak g$, and $|h\otimes Y|=|h||Y|$; therefore $|\ell_y^VF|\le\sum_p|\ell_y(p)|\phi_p$. By Schur's test, \[\sum_y(\sum_p|\ell_y(p)|\phi_p)^2\le\sum_y(\sum_p|\ell_y(p)|)(\sum_p|\ell_y(p)|\phi_p^2)\le\|\ell\|_{\rm r}\|\ell\|_{\rm c}\sum_p\phi_p^2\le3K_*^2\mathsf g(\lambda_*)^2\sum_p\phi_p^2;\] here we use that a base site carries three plaquettes, that the sums over the torus are at most the sums over $\mathbb Z^3$, and the bounds on $\|\ell\|_{\rm r},\|\ell\|_{\rm c}$ in \textup{(RB)}, which \Cref{Y-A-SetsOfRows} states in the sentence after its display. Hence \[x\mathrm{En}_\ell\le\frac x2\mathsf N_Z\sum_p\phi_p^2\le\frac x2\mathsf N_Z\varepsilon^2\cdot3N^3=\frac32\mathsf N_Zp_0^2N^3.\]

\textup{(4)} \emph{Haar measure.} The group $G$ is the unit sphere $S^3$, and $\operatorname{ang}g$ is the polar angle from $1$. So the law of $\operatorname{ang}g$ is $(2/\pi)\sin^2\tau\,d\tau$, and $\mu\{\operatorname{ang}g\le t\}=(t-\sin t\cos t)/\pi$. From $\sin u\le u-u^3/6+u^5/120$ for $u\ge0$, taken at $u=2t$, we get $t-\frac12\sin2t\ge\frac23t^3-\frac2{15}t^5\ge\frac8{15}t^3$ on $(0,1]$; and $8/(15\pi)\ge0.1697\ge\frac16$. At $t=\varepsilon/4\le1$ this gives $\mu\ge\varepsilon^3/384$ for each bond; since $dV$ is the product measure, $\int_{\mathfrak B}dV\ge(\varepsilon^3/384)^{3N^3}$.

\textup{(5)} As $\rho_1>0$, we have $\int\rho_1\ge\int_{\mathfrak B}\rho_1$, and by \textup{(1)} to \textup{(4)} the right side is at least the bound claimed.

\textup{(6)} For positive $f,g$, \[f/\!\int\!f-g/\!\int\!g=(f-g)/\!\int\!f+g(\int g-\int f)/(\int f\int g),\] and hence $\|f/\!\int\!f-g/\!\int\!g\|_1\le2\|f-g\|_1/\!\int\!f$. We take $f=\rho_1$ and $g=e^{\mathsf V_1-\mathsf V_2}\rho_2$, so that $g/\!\int\!g=\rho_2/\!\int\!\rho_2$. As $dV$, the product of the Haar probability measures, has total mass $1$, Step~8 gives $\|f-g\|_1\le e^{\mathsf V_1}2n\,e^{219n}\mathsf D$. Since $2\cdot2n\,e^{219n}e^{0.61n}\le4n\,e^{220n}$, division by the lower bound on $\int\rho_1dV$ gives the bound with $C^{\rm su}_{\rm vol}$.

\begin{remark}
The bound is exponential in $N^3$, with the factor $\varepsilon^{-9N^3}$.
\end{remark}

\emph{Step 10: \Cref{Y-J-ThB}\,(d).} Let $\mu_i=\rho_i/\!\int\!\rho_i$, $T=\{D(V)\text{ empty}\}$, $r_i=\mu_i(\text{complement of }T)$ and $\eta=1.22n(c_a+c_h)\mathsf D$. We show that \[\int|\mu_1-\mu_2|dV\le2\max r_i+e^{2\eta}-1.\]

If some $r_i=1$, the right side is $\ge2$, and there is nothing to show. Otherwise $\nu_i:=\mu_i1_T/(1-r_i)$ are two probability densities on $T$. On $T$, by \textup{(a)}, $\rho_1/\rho_2=e^{\mathsf V_1-\mathsf V_2}e^{S_1-S_2}$ with $|S_1-S_2|\le\eta$. So \[\nu_1/\nu_2=e^{S_1-S_2}/\!\int\!e^{S_1-S_2}d\nu_2\in[e^{-2\eta},e^{2\eta}],\] and $\int|\nu_1-\nu_2|=\int|\nu_1/\nu_2-1|d\nu_2\le e^{2\eta}-1$. Let $r_1\ge r_2$. Then \[\int_T|\mu_1-\mu_2|=\int|(1-r_1)(\nu_1-\nu_2)-(r_1-r_2)\nu_2|\le(e^{2\eta}-1)+(r_1-r_2).\] Off $T$ the integral is $\le r_1+r_2$. The sum of the two parts is $e^{2\eta}-1+2r_1$.
\end{proof}

\begin{proof}[Proof of \Cref{Y-J-CorC}]
\emph{Step 11: the laws.} We prove \Cref{Y-J-CorC}\,(a) to (d).

Relative to Haar measure $dV'$, the step is given by a probability kernel (the definition ``The step'' in item~2 of \S\ref{Y-J-1}); hence $\int\rho_{k+1}dV'=\int\rho_kdU$. So $\rho^{[K]}_k/\!\int\!\rho^{[K]}_k$ is the density of the law of the field obtained by $k$ steps from the Wilson measure at $\beta=L^K\mathsf b$ on $\mathbb Z^3/NL^K\mathbb Z^3$.

For \textup{(a)}, let $K'\ge K\ge s+K_0(L^s\mathsf b)$. Then $K\ge K_0(\mathsf b)$, and $k:=K-s$ is a level $>k_0$ (Step~6\,(a)). By \Cref{Y-J-ThA}\,(ii), (iii), (iv) at this $k$, the densities $\rho^{[K]}_{K-s}=e^{\mathsf V_1}\rho[u_k;\ell_k]$ and $\rho^{[K']}_{K'-s}=e^{\mathsf V_2}\rho[\mathfrak n_k;\ell_k]$ are points of a common class at $x=x_k=L^s\mathsf b$ on $\mathbb Z^3/NL^s\mathbb Z^3$, with $\mathsf D\le C_{L^s\mathsf b}\theta^{K-s}$. The hypotheses of \Cref{Y-J-ThB}\,(c) hold there: $x\ge x_\flat\ge e^2$; $NL^s$ is even, a multiple of $M$, and $\ge2M\,s(L^s\mathsf b)$ (Step~2\,(e)); $\kappa\ge6$; $\varepsilon\le\frac1{40}$ by \textup{(X16)}; $c_a+c_h\le\frac1{32}$ by \textup{(X17)}; and $\ell_k$ has range $\le\varrho(x_k)$ and satisfies $|\ell_{k,y}(p)|\le K_*e^{-\lambda_*|y-p|}$ (Step~2\,(c)). Both densities live on the grid with origin $o^{[K]}_{K-s}=o^{[K']}_{K'-s}$ (Step~2\,(a)), and the class is that of this grid. \Cref{Y-J-ThB}\,(c) therefore gives \textup{(a)}.

For \textup{(b)}: by \textup{(a)}, and as $\theta<1$, the densities form a Cauchy sequence in $L^1(dV)$, which is complete. The limit is $\ge0$ and has integral $1$. Since $\rho^{[K]}_k\in C^{\rm su}_\flat(x_k;\ell_k)$ for $1\le k\le K$, by \textup{(R8)} (Step~2\,(b)), and here $k=K-s\ge1$, the density $\rho^{[K]}_k$ is gauge invariant by \Cref{Y-J-P6}; Haar measure is invariant under gauge transformations; and invariance passes to limits in $L^1$. So the limit density is gauge invariant, and with it $\mu^\infty_{[s]}$, as \Cref{Y-J-CorC}\,(b) asserts.

For \textup{(c)}: for every $K$, $\mu^{[K]}_{[s-1]}$ is the image of $\mu^{[K]}_{[s]}$ under the kernel, \[\prod_ck_{aL^s\mathsf b}(V'(c),\mathsf M(U)(c))dV',\] since the noise is $a\,x_{K-s}=aL^s\mathsf b$; this kernel is independent of $K$. A probability kernel does not increase $\|\cdot\|_{\rm TV}$. We let $K\to\infty$ on both sides, using \textup{(b)} at $s$ and at $s-1$.

For \textup{(d)}, we let $K'\to\infty$ in \textup{(a)}.
\end{proof}

\begin{proof}[Proof of \Cref{Y-J-CorD}]
\emph{Step 12: the flux response.} We claim that \[|\mathfrak x(\rho^{[K']}_{K'})-\mathfrak x(\rho^{[K]}_K)|\le Q^{\rm su}(\mathsf b)\,C_{\mathsf b}\,\theta^K\] for $K'\ge K\ge K_0(\mathsf b)$.

The integer $m_0$ exists: the interval $[\varepsilon(\mathsf b)/2,\varepsilon(\mathsf b)]$ has the length $\varepsilon(\mathsf b)/2\ge2\pi/N^2$, and $2\pi/N^2$ is the spacing of the numbers $2\pi m_0/N^2$.

The field $V_*$ is abelian, as all its values are of the form $e^{i\alpha\sigma_3}$; so $\partial V_*(p)=e^{iF\sigma_3}$, where $F$ is the lattice curl of the exponents. At the plaquettes with normal $e_3$ and $z_1<N-1$, the curl is $\varepsilon_0(z_1+1)-\varepsilon_0z_1=\varepsilon_0$. At those with $z_1=N-1$ and $z_2<N-1$, it is $-\varepsilon_0Nz_2+\varepsilon_0N(z_2+1)-\varepsilon_0(N-1)=\varepsilon_0$. At those with $z_1=z_2=N-1$, it is $-\varepsilon_0N(N-1)-\varepsilon_0(N-1)=\varepsilon_0-2\pi m_0$. At all other plaquettes, $F=0$. So $\phi_p$ is $\varepsilon_0$ or $0$, since $0<\varepsilon_0\le\varepsilon(\mathsf b)\le\frac1{40}<\pi$, and $D(V_*)=D(1)=\emptyset$ at $x=\mathsf b$. Here $\varepsilon_0\le\varepsilon(\mathsf b)$ holds by the hypothesis on $m_0$, and $\varepsilon(\mathsf b)\le\frac1{40}$ by \textup{(X16)} at $x=\mathsf b\ge x_\flat$.

By \Cref{Y-J-ThA} at $k=K$, both densities have the form $e^{\mathsf V_i}\rho[\cdot;\ell_K]$ and lie in a common class, with $\mathsf D\le C_{\mathsf b}\theta^K$. We apply \Cref{Y-J-ThB}\,(a) at $V=1$ and at $V=V_*$; its hypotheses at $(\mathsf b,N)$ are those of Step~11, taken at $s=0$. The numbers $\mathsf V_1-\mathsf V_2$ cancel in the difference, and \[|[\log\rho_1(1)-\log\rho_1(V_*)]-[\log\rho_2(1)-\log\rho_2(V_*)]|\le2.44(N/M)^3(c_a+c_h)(\mathsf b)\mathsf D.\] We multiply by $2/(N^3\varepsilon_0^2)\le8/(N^3\varepsilon(\mathsf b)^2)=8\mathsf b/(N^3p_0^2)$. Since $c_a(\mathsf b)\mathsf b/p_0^2=E_ap_0\mathsf b^{1/2}$ and $8\cdot2.44=19.52$, the left side of the claim is at most $Q^{\rm su}(\mathsf b)\mathsf D\le Q^{\rm su}(\mathsf b)\,C_{\mathsf b}\,\theta^K$.
\end{proof}

\begin{remark}
The bound of the factor $2/(N^3\varepsilon_0^2)$ by $8\mathsf b/(N^3p_0^2)$ rests on $\varepsilon_0\ge\varepsilon(\mathsf b)/2$: at $m_0=1$ and $N$ large, the factor exceeds $8\mathsf b/(N^3p_0^2)$.
\end{remark}

\section{Source estimates for one step}\label{sec:single-source-step}

\subsection{Setup for the sourced step}\label{Y-H0-1}

The following sections analyse the RG flow in the presence of sources. The transport of the quadratic source forms is linear, and their charges obey an exact block-averaging identity. The analytic part of the RG step follows this transport up to a controlled remainder. The shared notation and the additional bounds on the starting scale are collected below.

\begin{definition}[Item (0): the shared letters and their senses] Every letter of Sections~\ref{Y-G0-1} to \ref{Y-J-1} has here the sense it has there, with the exceptions that entries 1 to 13 of this item name. One of them is the following: the letter $\kappa$ of those sections appears here as $\kappa_1$, as entry 1 explains.

\emph{Entry 1: the rates $\kappa$.} We put $\kappa_1:=\kappa+1$, and $\tilde\kappa$ ranges over the interval $[\kappa,\kappa_1]$, whose lower end is $\kappa\ge6$. The rate of \textup{(O5)} is $\kappa_1$: the statements of Sections~\ref{Y-C-1} to \ref{Y-J-1} are used with their letter $\kappa$ replaced by $\kappa_1$, and so are the entries \textup{(O5)} to \textup{(O13)} of the table. Further, $\kappa^\sharp=10(\kappa_1+4)$.

In this section and in Sections~\ref{Y-H1-1} to \ref{Y-H9-1} the symbol $\kappa_1$ always stands for $\kappa+1$. The rate $\kappa_k=\kappa+\sum_{k\le j<k_d}\check\kappa_j$ of the level $k$ is set in item~2, \textup{(o4)}, of \S\ref{Y-H6-1}.

In \S\ref{Y-F-1}, $\kappa_1$ is not the rate $\kappa+1$ but the number ${\rm py}/(240\pi^2L^3x)$, which enters the present section only through $x_{\rm fl}$.

The factor $e^{\tilde\kappa+3}$ of \S\ref{Y-H3-1} is taken at $\kappa_1$ where it enters a constant, and at $\kappa$ where it appears on the right side of the inequality of a floor.

\emph{Entry 2: the letters $n$, $q_N$ and $\nu$.} The letter $n$ is the number of sources. We write $q_N:=(N/M)^3$ for the number of cubes in a level of side $N$; in \S\ref{Y-J-1} this number is called $n$. The symbol $\nu_p$ denotes the local coupling of \Cref{Y-H4-U4}.

We write $\nu_a$ for the homogeneous factor of the $a$-channel, measured in the units $u_3$. By entry 9 one has $\nu_a=\vartheta_aL^{1/4}=0.2431$, and by the condition \textup{(L3su)} the factor in the bound of \Cref{Y-H1-U1}\,(h3) is $\le\vartheta_a$.

\emph{Entry 3: floors and levels.} The floors are $x_3$, $x_4$, $x_9^U$, $x_{\rm fl}$ and $x_s$, besides $x_\flat$ and $x_9$ of Sections~\ref{Y-G0-1} to \ref{Y-J-1}. The floor of \Cref{Y-H1-U1} is $x_9^U$; we write $x_3$ for the floor of \Cref{Y-H3-U3}, and $x_4$ is that of \Cref{Y-H4-U4}. The floor of \Cref{Y-H9-T} is $x_s$, which comes last in the order of \textup{(O13)} of item~(1). The index of $x_s$ is a mark and not the depth $s$ of \Cref{Y-H9-T}\,(T7+) or of \Cref{Y-H7-Df}\,(D3). The letters $x_k$, $x_j$, $x_{\check k}$, $x_{k_0}$, $x_{k_d}$, $x_K$, $x_{K-1}$, $x_0$ and $x_1$ are levels.

\emph{Entry 4: the symbols $D$, $\mathbb D$ and $\mathrm{sep}$.} Any family of kernels is written $D=(D_y)$, as in \S\ref{Y-I-1}. The source forms are written $\mathbb D_k[f]$. The separations are written $\mathrm{sep}_k$ and $\mathrm{sep}_*$. Item~2, \textup{(o4)}, of \S\ref{Y-H6-1} sets $\mathrm{sep}^{[i]}$, $\mathrm{sep}_k=\mathrm{sep}^{[K-k]}$ and $\mathrm{sep}_*$, and from them $j^*$ and $k_d$.

\emph{Entry 5: the symbols $s_x[D]$, $\sigma_D$ and $\sigma_M$.} For a cube $\square$ we set \[s_x[D](\{\square\},\mathsf W):=\frac x2\sum_{y\in\square}\sum_{p,p'}D_y(p,p')\,\bigl(\operatorname{Ad}_{\mathsf W(\Gamma(y,\pi(p)))}F_p\bigr)\cdot\bigl(\operatorname{Ad}_{\mathsf W(\Gamma(y,\pi(p')))}F_{p'}\bigr),\] where $F_p=-i\log\partial\mathsf W(p)$, and $s_x[D](X,\cdot)=0$ whenever $\#X\ge2$. One has $E^\Delta[D]=-s_x[D]$. We further set $\sigma_D:=c_Q^{\rm su}\,p_0^2\,\|D\|_{(\lambda_\flat)}$ and $\sigma_M:=160(\kappa_1+1)c_R/M$.

\emph{Entry 6: amplitudes, and the letters $\alpha$ and $\mu$.} The amplitude at $x$ is $\mathrm{amp}(x):=c_Q^{\rm su}K_5x^{1/2}/(\mathsf bE_ap_0(x))$, and we put $\mathrm{amp}_W:=1/(c_\nu\mathsf b\,\varepsilon(\beta))$. The two satisfy the inequality $1.01\,\mathrm{amp}_W+0.01/\mathsf b\le\mathrm{amp}(x_0)$, for the following reasons. Since $c_Q^{\rm su}=306M^3\mathsf g_\flat^2$, we have $\mathrm{amp}(x_0)=51\mathsf g_\flat^2K_5\,\mathrm{amp}_W$, where $\mathsf g_\flat,K_5\ge1$. Next, \textup{(X17)} and \textup{(X14)} of \S\ref{Y-J-1} give $x^{1/2}/(E_ap_0)\ge32p_0^2\ge32$. Consequently $\mathrm{amp}_W=6M^3x_0^{1/2}/(\mathsf bE_ap_0(x_0))\ge192M^3/\mathsf b\ge0.01/\mathsf b$. The left side of the inequality is therefore at most $2.01\,\mathrm{amp}_W\le51\,\mathrm{amp}_W$, and since $\mathsf g_\flat,K_5\ge1$ the identity above shows that this is at most $\mathrm{amp}(x_0)$.

The letter $\alpha$ is a letter of the alphabet of \Cref{Y-H-C9}, as in item~6 of \S\ref{Y-H-1}. Taylor indices are written $\mu$.

\emph{Entry 7: the function $u_3$, the letters $u$ with an index, and $\check\omega_3$.} We set $u_3(x):=C_3x^{1/4}$. The families of the levels are written $u_k(z)$, $\bar u_k$, $u_k^{\mathrm{pol}}$, $u_0(\nu)$ and $u_1(\nu)$.

The symbol $u_k(z)$ has two senses. In the first it is the honest family of \Cref{Y-H4-U4}\,(h) on $\mathrm{Pd}_0$; it has this sense in that clause, in the data of \Cref{Y-H6-J}, and at every place where the family is called honest. In the second it stands for the representatives of item~4 of \S\ref{Y-H6-1}, that is, for $u_k^{\mathrm{pol}}$ when $k\le k_d$ and for their iterates under $F$ beyond $k_d$; it has this sense in \textup{(J6)} to \textup{(J8)} of \Cref{Y-H6-J}, in \Cref{Y-H7-Df}\,(D1), (D3), (D4), (D6). The two share the base point $\bar u_k$ and the multi-affine jets, by \textup{(J1)} and \textup{(J6)}; beyond that, the two families can differ.

We set $\check\omega_3(x):=C_3E_ap_0(x)x^{-1/4}/c_Q^{\rm su}$. This closed form is defined at \textup{(O12b)} and is free of $\mathsf b$, and it satisfies the identity $\check\omega_3(x)=u_3(x)K_5/(\mathsf b\,\mathrm{amp}(x))$.

\emph{Entry 8: radii.} The radii $r(x)$, $s(x)$ and $\varrho(x)$ are those of item~2 of \S\ref{Y-H-1}. The polydiscs have the radii $\mathfrak r(x):=3.5/(n\,\mathrm{amp}(Lx))$, $\mathfrak r_k:=\mathfrak r(x_k)$ and $\mathfrak r_*:=\mathfrak r_{k_d}$; here $\mathfrak r_*$ is set in item~2, \textup{(o5)}, of \S\ref{Y-H6-1}, and $k_d$ in item~2, \textup{(o4)}, of that section. We write $\mathrm{Pd}(r):=\{|z_i|<r_i\}$, and when $r$ is a single number every $r_i$ equals $r$. The remainder of \Cref{Y-H6-J}\,(J3) is $\mathrm{rem}_{k,i}$: the $a$-part of the jet at the level $k$ in the source $i$ alone, minus $s_{x_k}[\mathbb D_k[f_i]/\mathsf b]$. We put $\check r(x):=r(x)+3$.

\emph{Entry 9: the thetas.} From Sections~\ref{Y-G0-1} to \ref{Y-J-1} we take $\vartheta=0.4468$, $\theta=\vartheta^{4/9}\le0.69903$, $\vartheta_w=0.248$, $\vartheta_{10}$, $\vartheta_1=c_{\rm fix}L^{-1/2}$, and $\vartheta(x)$ of \textup{(u1)}. The thetas introduced in Sections~\ref{Y-H1-1} to~\ref{Y-H9-1} are $\vartheta_s:=c_sL^{-15/32}$, $\vartheta_5:=\max\{\vartheta_s,\vartheta_1^{15/16}\}$, $\vartheta_a:=0.2431\,L^{-1/4}$ and $\vartheta[B]:=\max\{16\vartheta_w,1\}/(16-B)$.

\emph{Entry 10: constants.} The number $C_5$ is that of \textup{(O0a)} (item~2 of \S\ref{Y-H2-1}), and $C_5^\#:=L^{15/32}C_5+K_5$. We set $C_3:=1686\,e^{\kappa+4}L^{11/4}c_Q^{\rm su}/E_a$. The constants $C_n$ and $C'_n$ are those of \Cref{Y-H9-T}. The numbers $K_5$ and $C_{\rm ch}$ belong to \textup{(O0a)}. We set $C_g:=263L^3$; it is bounded above by $10^3L^3$, and at the level 0 \Cref{Y-H4-U4}\,(e) gives the sharper value $2.43L^3$. The numbers $k_1,k_2,k_3$ are those of \textup{(O10a)} (item~1 of \S\ref{Y-H3-1}).

\emph{Entry 11: further symbols, and symbols with a local sense.} The threshold of \Cref{Y-H9-T}\,(T2) is $K_T(\mathsf b,n,d)$.

We set $\tau_7=\max\{1,B_0C_{\rm vol}\}$. This is a function of the datum of \Cref{Y-H8-CoF}\,(F0) and not a single number; in the clauses of \Cref{Y-H9-T} the datum is as follows:

\begin{enumerate}\item[\textup{(T1)}] $B_0=0.91$ and $C_{\rm vol}(N,\mathsf b)$;

\item[\textup{(T2)}] $B_0=1$ and $C_{\rm vol}(N,\mathsf b)$;

\item[\textup{(T7+)}] $B_0=1$ and $C_{\rm vol}(NL^s,L^s\mathsf b)$.\end{enumerate}

We write $C_{\rm vol}:=C^{\rm su}_{\rm vol}$. The radius of the balls of the supports is $R_f$. The exponent of the side is $m_l$, and $m=K'-K$.

The letter $l$ is the side, and $T_l=\mathbb R^3/l\mathbb Z^3$ is the torus. The letter $d$ is the separation. In entry 5 of this section, $\pi(p)$ is the base site of $p$.

We set $m_\Xi:=\frac12-d_0-c_\Xi$. The set of the pairs $p$ of \Cref{Y-H1-U1} is $\mathbb P$, and $P_5$ is the number of \Cref{Y-H2-U5}\,(S3). The number $\check\kappa_k$ is that of item~2, \textup{(o4)}, of \S\ref{Y-H6-1}. It satisfies $\check\kappa_k=\kappa_k-\kappa_{k+1}$ for $k<k_d$. The numbers of \Cref{Y-H7-Df}\,(D2) are $a^D_k$, $h^D_k$ and $e_k$, and $\mathrm{Tail}(J)$ and $\mathrm{Ser}_2$ are the series of item~1.3 of \S\ref{Y-H7-1}. We set $\gamma_r:=1/4$.

The symbol $\mathsf V'$ is the number of \textup{(Fm2)}, and $V'$ is the coarse field.

\emph{Entry 12: tags.} The following clause names are used in the estimates for the sourced flow. The clauses \textup{(U)}, \textup{(c+)}, \textup{(f1)} to \textup{(f6)}, \textup{(g1)} to \textup{(g4)}, \textup{(h)}, \textup{(h1)} to \textup{(h3)}, \textup{(h\#)} and \textup{(b\#)} are those of \Cref{Y-H1-U1}; \textup{(G1)} to \textup{(G4)} and \textup{(J0)} to \textup{(J8)} are those of \S\ref{Y-H6-1}; \textup{(B1)} to \textup{(B8)} are those of \Cref{Y-H5-Bc}; \textup{(D0)} to \textup{(D7)} are those of \Cref{Y-H7-Df}; \textup{(a$^C$)} to \textup{(d$^C$)}, \textup{(S)} and \textup{(c1)} to \textup{(c4)} are those of \Cref{Y-H8-ThBC}; \textup{(F0)} to \textup{(F5)} are those of \Cref{Y-H8-CoF}; \textup{(R0)} to \textup{(R3)} are those of \Cref{Y-H8-CoR7}; and \textup{(S0)} to \textup{(S6)} are those of \Cref{Y-H2-U5}.

The floors \textup{(Y.0)} to \textup{(Y.9)} are stated in item~5 of \S\ref{Y-H-1}, and \textup{(Y.10)} in item~0, \textup{(O13)}\,(v), of \S\ref{Y-H1-1}. For the floor $x_9^U$, \textup{(Y.0)} to \textup{(Y.9)} are taken in the form in which item~0, \textup{(O13)}\,(i) to~(iv), of \S\ref{Y-H1-1} lists them. Further, \textup{(X5)} is stated in item~2 of \S\ref{Y-H2-1}, \textup{(Z.0)} to \textup{(Z.4)} in item~1 of \S\ref{Y-H3-1}, and \textup{(W1)} to \textup{(W5)} in item~1 of \S\ref{Y-H4-1}. A second pair \textup{(W2)}, \textup{(W3)} occurs in item~1.3 of \S\ref{Y-H7-1}. Where \textup{(W2)} and \textup{(W3)} are cited by their tags alone, in item~(1) below and in item~1 of \S\ref{Y-H9-1}, they are the two of item~1.3 of \S\ref{Y-H7-1}, not those of item~1 of \S\ref{Y-H4-1}.

\emph{Entry 13: the observable.} The observable is $\mathcal O_K(f):=a_K^{-1}\sum_pf(x(p))\,a_p(U)$, in which $x(p)$ is the barycentre (\Cref{support:150-2}; item~2 of \S\ref{Y-H9-1} writes it out among its data). Moreover $s_p=a_p$.\end{definition}

\begin{definition}[Item (1): the numbers and floors introduced in Sections~\ref{Y-H1-1} to~\ref{Y-H9-1}, each at its place] Each of them is listed at its place in the one table, together with what it depends on.

\begin{enumerate}\item[\textup{(O0a)}] For \S\ref{Y-H3-1}: $\mathsf g_\flat=\mathsf g(\lambda_\flat)$, $g_{1b}$ and $g_b^-$, which depend on $\lambda_*$. For \S\ref{Y-H2-1}, the pure numbers $C_{\rm ch}=4e^{2\lambda_*}/\gamma_*$; $K_5=1.35\,C_{\rm ch}$, which equals $0.60\,e^{2\lambda_*}K_*^2\le0.61K_*^2$; $\mathsf g'=\mathsf g(\lambda_\flat/2)$; $\mathsf g_1'=\mathsf g_1(\lambda_\flat/2)$; $a_G$; $c_s$; $t_5=18\,\mathsf g_\flat \mathsf g'$; $a_H$; $c_I$; $A_1=\max\{1,15C_1^2\}$; $I_5=c_IK_5A_1$; and $C_5=0.3(K_5+1)+1.02\,I_5$.

\item[\textup{(O2a)}] The number $m_\Xi$ (\S\ref{Y-H4-1}), which depends on $d_0$ and $c_\Xi$.

\item[\textup{(O4a)}] Two conditions on $L$.

\begin{enumerate}\item[\textup{(L3su)}] $L^{1/4}\ge10^4\mathsf N_{\rm c}^2$. It depends on $\mathsf N_{\rm c}$ of \textup{(O3)}, which grows like $\log L$. The condition can be met: one has $\mathsf N_{\rm c}=17\bigl(c_1^*(1+\log L)\mathsf N_Z\mathsf P_\rho+1\bigr)$, where $c_1^*$, $\mathsf N_Z$ and $\mathsf P_\rho$ do not depend on $L$, and the quotient $L^{1/4}/(1+\log L)^2$ increases as soon as $\log L\ge7$; hence \textup{(L3su)} holds for every $L$ from some value on, as do the three conditions of \textup{(O4)}, and it is required in addition to them. Squaring \textup{(L3su)} gives $L^{1/2}\ge10^8\mathsf N_{\rm c}^4\ge1.2\cdot10^4\mathsf N_{\rm c}^2$, so that \textup{(L3su)} implies \textup{(L2$^{\rm su}$)}.

\item[\textup{(L5*)}] $c_s\le L^{43/160}$, which depends on \textup{(O0)}. It follows from $L\ge L_B\ge c_{\rm fix}^{10}$, because $\log_{10}c_s\le145.34<204.89\le\log_{10}c_{\rm fix}$ (\Cref{Y-H2-Rem2}). It therefore adds nothing to the conditions on $L$.\end{enumerate}

\item[\textup{(O5)}] The rates $\kappa$, $\kappa_1$ and $\kappa^\sharp$.

\item[\textup{(O7)}] The conditions \textup{(M2)}, \textup{(M4$^\flat$)} and \textup{(M5r)}, taken at $\kappa_1$ and $\kappa^\sharp$.

\item[\textup{(O7a)}] The number $\sigma_M:=160(\kappa_1+1)c_R/M$, which depends on $\kappa_1$, $c_R$ and $M$. By \textup{(M2)}, $\sigma_M\le1/12$.

\item[\textup{(O10a)}] The numbers $k_1$, $k_2$, $k_3$ (\S\ref{Y-H3-1}), which depend on $\mathsf g_\flat$, $g_{1b}$, $g_b^-$ of \textup{(O0a)} and on $L$, $\mathsf a_{\rm c}$, $\mathsf K_{\rm c}$, $C_D$ and $\theta_D$; and $C_g$, which depends on $L$.

\item[\textup{(O11)}] The number $c_h$ of item~3, \textup{(O11)}, of \S\ref{Y-G0-1}; $c_h/2$ is computed from it, and no condition arises.

\item[\textup{(O12a)}] The number $c_\nu:=E_a/(6M^3)$ (\S\ref{Y-H4-1}), and $C_3$, which depends on $\kappa$, $L$ and $E_a$, and on $c_Q^{\rm su}$ of \textup{(O10)}, hence on $M$ and on $\mathsf g_\flat$ of \textup{(O0a)}.

\item[\textup{(O12b)}] Three functions of $x$. They are defined before $x_s$, they are given in closed form, and they are built from \textup{(O2)}, \textup{(O4)}, \textup{(O8)}, \textup{(O10)}, \textup{(O11)}, \textup{(O12)} and \textup{(O12a)} and from nothing else: $\eta_h(x)=e^{-(c_h/2)(p_0(x)^2-p_0(x/L)^2)}$, $\eta_w(x)=e^{-c_\Xi p_0(x)^2}$ and $\check\omega_3(x)=C_3E_ap_0(x)x^{-1/4}/c_Q^{\rm su}$. They are free of $\mathsf b$, $n$ and $d$. The floor $x_s$ depends on them through \textup{(W2)} and \textup{(W3)}, and through $\check\omega_3$ it depends on $C_3$.

\item[\textup{(O13)}] The one-sided floors, which come in the following order. First, $x_9^U\ge x_{\rm fl}$, subject to \textup{(Y.0)} to \textup{(Y.9)}, in the form in which item~0, \textup{(O13)}\,(i) to~(iv), of \S\ref{Y-H1-1} lists them, and to \textup{(Y.10)}, which reads $1800\,L^{7/2}x^{-1/4}\le1$. Each of \textup{(Y.0)} to \textup{(Y.9)} is taken at the pair of $\mathbb P=[\kappa,\kappa_1]\times\{c_h,c_h/2\}$ at which it is strongest.

Second, $x_3\ge x_9^U$, subject to \textup{(Z.0)} to \textup{(Z.4)} of item~1 of \S\ref{Y-H3-1}, with \textup{(Z.4)} taken at $\kappa$. Third, $x_4\ge x_9^U$, subject to \textup{(W1)} to \textup{(W5)} of item~1 of \S\ref{Y-H4-1}, taken at $\kappa_1$.

Last, $x_s\ge e^2$. It is subject to all of \textup{(O13)} of item~1 of \S\ref{Y-J-1}, taken at $\kappa_1$; to $Lx_s\ge\max\{x_9^U,x_3,x_4\}$; and to the three conditions below, of which the last two are required for all $x\ge Lx_s$.

\begin{enumerate}\item[\textup{(X5)}] $1.01\,(4/3)\,t_5C_{\rm ch}L\,x_s^{-16}\le1/100$;

\item[\textup{(W2)}] $\eta_h+e^{1/2}\eta_w\le1/L$;

\item[\textup{(W3)}] $\check\omega_3(x),\ \check\omega_3(x/L)\le10^{-3}$.\end{enumerate}

Four remarks follow, and none of them adds a floor. First, \textup{(Y.10)} and \textup{(L3su)} give $[2430\mathsf N_{\rm c}^2L^{-1/2}+L^{-1}]L^{1/4}\le0.2431$. Second, there are the identities $\vartheta_au_3(x)/u_3(x/L)=0.2431\le0.3452$ and $\check\omega_3(x/L)/\check\omega_3(x)=L^{1/4}p_0'/p_0\le L^{1/4}$. Third, $\vartheta_5\le L^{-1/5}$ is \textup{(L5*)} together with $L\ge L_B$. Fourth, at the last level only $\tilde c=c_h$ is used, and what is required there is \textup{(X17)} of \S\ref{Y-J-1}.

\item[\textup{(O14)}] The numbers on which no floor depends. They are listed in an order in which no entry depends on a later one. They are the entries of \textup{(O14)} of item~1 of \S\ref{Y-J-1}, taken at $\kappa_1$, namely $c_4$, $\Lambda$, $K_0$, $k_0$, $\varepsilon_L$, $C_{\mathsf b}$ and $C_{\rm vol}$; then $\vartheta_s$; $\vartheta_5$; $C_5^\#$; $\vartheta_a$; $\vartheta[B]$; $\mathrm{amp}$; $u_3$; the identity $\check\omega_3=u_3K_5/(\mathsf b\,\mathrm{amp})$; $\mathfrak r$; the functions $\varepsilon_2^\sharp(x)$, $a^\sharp(x)$ and $\tau^\sharp(x)$ of $x$ from item~0 of \S\ref{Y-H1-1}; $\check r$; $\mathrm{sep}_k$; $\mathrm{sep}_*=8\log(4nL^{1/2})$; $j^*$; $k_d$; $\check\kappa_k$; $\kappa_k$; $\mathfrak r_*$; $N_1^{[j]}$; $m^{[j]}$; $K_T$ with \textup{(K1a)} and \textup{(K1b)}, and $\Omega^{\#}$ (item~1.3 of \S\ref{Y-H7-1}); $\mathrm{Tail}$; $\mathrm{Ser}_1$; $\mathrm{Ser}_2$; $P_*$ (item~1 of \S\ref{Y-H9-1}); $\tau_7$; $C_n$; and $C'_n$.\end{enumerate}

\emph{No circle.} There is no circle in the table, and the reasons are these. Every entry depends only on places that come strictly earlier. The conditions \textup{(L3su)} and \textup{(L5*)} tie $L$ to \textup{(O0)} and \textup{(O3)}. No number of \textup{(O10a)} depends on $c_h$, on $E_a$ or on $x$. The number $c_\nu$ depends on \textup{(O7)} and \textup{(O12)}, and $C_3$ depends on \textup{(O0a)}, \textup{(O4)}, \textup{(O5)}, \textup{(O7)}, \textup{(O10)} and \textup{(O12)}.

Every floor has the same form, a family of one-sided inequalities in $x$. Their left sides are a polylogarithmic factor times a negative power of $x$, or $e^{-cp_0^2}$, where the exponent $\mathfrak p_0$ of $p_0$ is $>2$; or else, for $\eta_h$ in \textup{(W2)}, the left side is $e^{-c(p_0(x)^2-p_0(x/L)^2)}$. In the last case the mean value theorem gives $p_0(x)^2-p_0(x/L)^2\ge2\mathfrak p_0A_0^2\log L\,(\log(x/L))^{2\mathfrak p_0-1}$, and this lower bound tends to infinity because $2\mathfrak p_0-1>3$.

No floor depends on $N$, $K$, $\mathsf b$, $n$, $d$, $f$, on a member or on a row. Finally, $x_\flat^{(1)}$ of \Cref{Y-B-C1} is free of $c_h$ and $E_a$, as that theorem states.\end{definition}

The linear source transport and its comparison across two cutoffs are proved in \S\ref{Y-H2-1}. \S\ref{Y-H3-1} identifies this transport as the leading part of the derivative of the analytic RG step. The initial source family is given in \S\ref{Y-H4-1}. The holomorphic estimates of \S\ref{Y-H5-1} allow the flow to be analysed through its squarefree Taylor coefficients.

Sections~\ref{Y-H6-1} and~\ref{Y-H7-1} estimate those coefficients for one and two runs, respectively. In the two-run comparison the free source forms are compared before the two families are evolved at common rows. \S\ref{Y-H8-1} controls the final integral, and \S\ref{Y-H9-1} proves the finite-volume cutoff limit. \S\ref{Y-P1-1} gives the estimates for arbitrarily many separated sources using the number of sources visible to each polymer.

\subsection{Estimates at two decay rates}\label{Y-H1-1}

In this section we extend \Cref{Y-H-C9}, which is stated for one pair of rates, to every pair of rates. Since the formulas \textup{(Fm1)--(Fm4)} contain neither the rate $\tilde\kappa\in[\kappa,\kappa_1]$, nor $\tilde c\in\{c_h,c_h/2\}$, nor $E_a$, nor $A$, they define one map $F_\ell$, one number $\mathsf V'$ and one family $v$, and the clauses \textup{(a)--(e)} hold at every pair, with the same numerals, above a single floor $x_9^U$. New clauses treat holomorphic families in several variables, the members on which an image member depends together with its reach, the number $\mathsf V'$ as a sum of local numbers with bounds in the cell sum, and the derivative of $E'^{\rm old}$ in a direction $\delta$, split into the part of one letter and a rest, with bounds that in turn bound $E'^{\rm old}$; two of the bounds, the first bound on the rest $\mathrm{Rst}$ in \textup{(h1)} and the bound on $r$ in \textup{(g3)}, also hold when $x^{-1/4}$ is replaced by $\tau^\sharp(x)$, but with larger numerals: $291\,L^3\tau^\sharp(x)$ in the place of $217\,L^3x^{-1/4}$ and $146\,L^3\tau^\sharp(x)$ in the place of $109\,L^3x^{-1/4}$.

\begin{remark}[What the theorem is, and where its signs come from] \Cref{Y-H1-U1} extends \Cref{Y-H-C9} from one pair of rates to every pair: it keeps the clauses \textup{(a)--(e)} of that theorem and adds to them the clauses \textup{(U)}, \textup{(b\#)}, \textup{(c+)}, \textup{(f)}, \textup{(g)}, \textup{(h)} and \textup{(h\#)}, and it is accompanied by \Cref{Y-H1-LK} and \Cref{Y-H1-Rcl}.

From items~0 to~6 of \S\ref{Y-H-1} we take the signs, the lattices, the cubes and the cells $\square$, the families $\mathcal D$ and $\mathcal P$, the radii, the domains $\mathcal U_X(x)$ and $U^\tau(B;x')$, the members and the distances, the data \textup{(SD)}, the inputs \textup{($\Gamma$2)--($\Gamma$7)} and the tools, the alphabet $\mathfrak A$ together with $J_\alpha$, $I_\alpha$, the relation $\sim$ and the relation ``compatible'', $\operatorname{co}(C)$, $n(C)$, $\Phi$ and $\hat\Phi$, the formulas \textup{(Fm1)--(Fm4)}, and the floors \textup{(Y.0)--(Y.9)}. From item~1 of \S\ref{Y-G0-1} we take $Y_0(X)$, the word ``small'', ${\rm cl}(X)$ and the fitted grids. From Lems.~\ref{Y-G-LKP} and~\ref{Y-G-L18} in Step~8 of the proof of \Cref{Y-G-ThA} we take the polymers, the relation ``touch'', the clusters, which are finite sequences, and $\varphi^T$.

The letters and the numbers are those of items~(0) and~(1) of \S\ref{Y-H0-1}. Throughout this section, ``Step~$n$'' denotes Step~$n$ of the proof of \Cref{Y-H-C9}, while the steps of the proof given here are named S1 to S13.\end{remark}

\begin{remark}[Letters with more than one sense]

The letter $t$ is the number $x^{-1/4}$, but as a member of a cell datum $(\mathsf J,t)$, and so in $\hat\lambda_{(\mathsf J,t)}$ and in a zone letter $(\mathsf J,t,\mathfrak X,\mathfrak Y)$, it stands for the set of the tested cells; and the families of numbers of item~1 are written $v$, $v(Y;u)$ and $v^0(Y)$, while in Step S7, in $v(2r_1+1)^3$, we have $v=(LM)^3$. Outside the proof of \Cref{Y-H1-LT57}, the bare $A$ denotes the number of caps, whereas in that proof the bare $A$ is the family $(A_c)_c$ of the $A_c$, indexed by the bonds $c$ of $\mathsf B_\omega$ and with values in $\mathfrak g_{\mathbb C}$, and $\operatorname{Ad}_uA:=(\operatorname{Ad}_uA_c)_c$. The tags \textup{(a)}, \textup{(b)} and \textup{(h)} name clauses of \Cref{Y-H-C9} and of \Cref{Y-H1-U1}, and also kinds of letters, as in \textup{(a)}$_X$, \textup{(h)}$_X$ and \textup{(b)}. The two senses of the bare letter $G$ are the group and the matrix of the law, those of the bare letter $P$ are a coarse plaquette and the function $P(Y,\cdot)$ of \textup{(Fm1)}, and the context shows which sense is meant.\end{remark}

\begin{definition}[Item 0: the letters and their place in the table] The places \textup{(O0)--(O4)}, \textup{(O6)}, \textup{(O8)--(O10)} and \textup{(O12)} of the table are as in item~1 of \S\ref{Y-H-1}. The remaining places, and the conditions added to them, read as follows.

\begin{enumerate} \item[\textup{(O4a)}] The place \textup{(O4a)} consists of the conditions of \textup{(O4)} in item~1 of \S\ref{Y-H-1}, together with one more condition, \begin{equation}\tag{L3su}L^{1/4}\ge10^4\mathsf N_{\rm c}^2.\end{equation}

By item~1, \textup{(O3)}, of \S\ref{Y-H-1}, $\mathsf N_{\rm c}=17(c_1^*\mathsf l_L\mathsf N_Z\mathsf P_\rho+1)$ with $\mathsf l_L=1+\log L$, and the numbers $c_1^*$, $\mathsf N_Z$, $\mathsf P_\rho$ are computed from \textup{(O0)--(O3)} and are independent of $L$. For $\log L\ge7$ we have \[\frac{d}{d\log L}\log(L^{1/4}/\mathsf N_{\rm c}^2)\ge\frac14-\frac2{\mathsf l_L}\ge0,\] and $L^{1/4}/\mathsf N_{\rm c}^2\to\infty$. Hence \textup{(L3su)} holds from some $L$ on. Moreover \textup{(L3su)} implies \textup{(L2$^{\rm su}$)}, since it gives $L^{1/2}\ge10^8\mathsf N_{\rm c}^4\ge1.2\cdot10^4\mathsf N_{\rm c}^2$. As \textup{(L3su)} depends on $L$ and on \textup{(O0)--(O3)} alone, we may impose it on $L$ at the same time as \textup{(O4)}, before any of the later letters is chosen.

\item[\textup{(O5)}] The number $\kappa\ge6$ is the lower end of the range $[\kappa,\kappa_1]$ of the rate $\tilde\kappa$. We set $\kappa_1:=\kappa+1$: the table, and the chain of statements of Sections~\ref{Y-C-1} to \ref{Y-J-1}, are taken with the letter $\kappa$ of \S\ref{Y-H-1} equal to $\kappa_1$. We set $\kappa^\sharp:=10(\kappa_1+4)$.

\item[\textup{(O7)}] The side $M$ is chosen by \textup{(M2)}, \textup{(M4$^\flat$)} and \textup{(M5r)} at $\kappa_1$ and $\kappa^\sharp$. Hence every number of \textup{(O6)--(O12)}, the floor $x_{\rm fl}$ and every radius of item~2 of \S\ref{Y-H-1} are defined once, and so are the objects $G$, $C$, $a_0$, $\mathsf D_k(Y)$, $S_k(Y)$, $P$, $T_\square$, $V_\square$, $\hat\lambda_{(\mathsf J,t)}$, $b_\square$.

\item[\textup{(O11)}] The number $c_h$ is as in item~1 of \S\ref{Y-H-1}, and $c_h/2$ is computed from it. As \textup{(O11)} bounds $c_h$ from above, $c_h/2$ obeys it just as $c_h$ does.

\item[\textup{(O13)}] We let $x_9^U$ be $1$ plus the infimum of the numbers $y\ge x_{\rm fl}$ such that, for all $x\ge y$, the following conditions hold. \begin{enumerate} \item[\textup{(i)}] \textup{(Y.0)}, \textup{(Y.2)}, \textup{(Y.7)}, \textup{(Y.8)}, the second inequality of \textup{(Y.5)}, the second and the third inequality of \textup{(Y.6)} and the second inequality of \textup{(Y.9)}, all as in item~5 of \S\ref{Y-H-1}, with $\kappa^\sharp$ as in \textup{(O5)} above.

\item[\textup{(ii)}] \textup{(Y.1)}, \textup{(Y.3)}, \textup{(Y.4)}, together with the $\upsilon_f$ and $\upsilon_*$ that occur in it, and the first inequality of \textup{(Y.6)}, with $\kappa_1$ in the place of the letter $\kappa$ of \S\ref{Y-H-1} and $e^{-(c_h/2)p_0^2}$ in the place of $c_h(x)$.

\item[\textup{(iii)}] The first inequality of \textup{(Y.5)}, with $c_h/2$ in the place of $c_h$.

\item[\textup{(iv)}] The first inequality of \textup{(Y.9)}, with the lower end $\kappa$ in the place of the letter $\kappa$ of \S\ref{Y-H-1}.

\item[\textup{(v)}] The inequality \begin{equation}\tag{Y.10}1800\,L^{7/2}x^{-1/4}\le1.\end{equation}\end{enumerate}

Each inequality of \textup{(i)--(v)} holds for all large $x$, by Step S1, and it depends on $x$ and on \textup{(O0)--(O12)} alone.

The set of these $y$ is an upper half line, and Step S1 shows that it is not empty. It need not contain its infimum, because the radii are step functions of $x$ and some of them occur on the side of an inequality that is asserted to be the larger, as $r'$ does on the right side of \textup{(Y.7)}. The number $x_9^U$ lies in the set and depends on what the set depends on; this is the same form of a floor as in \Cref{Y-B-C1} (item~3 of \S\ref{Y-B-1}).

We have $x_9^U\ge x_9$, where $x_9=1+\inf S_9$ is the floor of item~1, \textup{(O13)}, of \S\ref{Y-H-1}, taken at $\kappa_1$: by Step S1 the conditions \textup{(i)--(iv)} imply \textup{(Y.0)--(Y.9)} at the pair $(\kappa_1,c_h)$, so that the set of these $y$ lies in $S_9$.

In the table of item~(1) of \S\ref{Y-H0-1}, the floor $x_s$ is taken with $Lx_s\ge x_9^U$.\end{enumerate}

\emph{Pairs.} We set $\mathbb P:=[\kappa,\kappa_1]\times\{c_h,c_h/2\}$, and we write its elements as $p=(\tilde\kappa,\tilde c)$. For $p\in\mathbb P$, the class $M_A(x;p)$, the norm $\|\cdot\|_{x;p}$ and the distances $\mathsf D_{x;p}$, $\mathsf D^{\mathrm{in}}_{x;p}$ are the class $M_A^{\rm su}(x)$, the norm $\|\cdot\|_x$ and the distances $\mathsf D_x$, $\mathsf D^{\mathrm{in}}_x$ of item~3 of \S\ref{Y-H-1}, defined with $\tilde\kappa$ in the place of the letter $\kappa$ of that section and with $e^{-\tilde cp_0^2}$ in the place of $c_h(x)$; at $x'$ they are defined with $\tilde\kappa$ and with $e^{-\tilde cp_0'^2}$. Then $M_A(x;\kappa_1,c_h)\subset M_A(x;p)$ for every $p$, and every $M_A(x;p)$ is contained in $M_A(x;\kappa,c_h/2)$.

\emph{Numbers.} \begin{enumerate} \item[\textup{(O7a)}] $\sigma_M:=160(\kappa_1+1)c_R/M$, which is $\le1/12$ by \textup{(M2)}.

\item[\textup{(O10a)}] $C_g:=263L^3$.

\item[\textup{(O14)}] No floor depends on \textup{(O14)}. It sets $\vartheta_a:=0.2431\,L^{-1/4}$ and three functions of $x$, like $t$ and $\varepsilon_2$ of item~5 of \S\ref{Y-H-1} (here $s_0$, $\varrho_*$ and $s'$ depend on $x$): \[\varepsilon_2^\sharp=\varepsilon_2^\sharp(x):=[6(2(s_0+\varrho_*)+1)^3]^{-1},\] \[a^\sharp(x):=50(2s'+1)^3e^{(\kappa_1+1)s'}\,[\,5e^{\kappa_1+3}L^3\cdot16c_a(x)+2e^{\kappa_1+3}K_v^\flat p_0^3x^{-1/2}+x^{-20}\,],\] \[\tau^\sharp(x):=a^\sharp(x)/\varepsilon_2^\sharp.\]\end{enumerate}

As in item~5 of \S\ref{Y-H-1}, $t=x^{-1/4}$ and $\varepsilon_2=x^{-1/8}$.\end{definition}

\begin{definition}[Item 1: the objects] \emph{Cell sum.} For a family $f=(f(Y,\cdot))$ indexed by the connected sets $Y$ of cells we set \[|f|_{\mathrm{cs}}:=\sup_\square\sum_{Y\ni\square}e^{\tilde\kappa(\#Y-1)}\sup|f(Y,\cdot)|,\] where the inner supremum is taken over $\mathcal U_Y(x')$. For a family of numbers, the modulus of each number takes the place of the inner supremum. The cell sum depends on $\tilde\kappa$.

We let $T^a(x)$ be the set of the families $\delta=(\delta(X,\cdot))$, $X\in\mathcal D$, of bounded functions which have the clauses of \textup{(ca)} that carry no size, namely that $\delta(X,\cdot)$ is holomorphic on $\mathcal U_X(x)$, invariant under $G$-valued gauge transformations and $0$ at pure gauges. For $\delta\in T^a(x)$, $B=B_p(\delta)$ is the least number with $\sup_{\mathcal U_X(x)}|\delta(X,\cdot)|\le B\,c_a(x)e^{-\tilde\kappa(\#X-1)}$ for all $X$; here $B$ may be any number $\ge0$. Further, $A_a=A_a(u):=\sup_X\sup_{\mathcal U_X(x)}|E^a(X,\cdot)|\,e^{\tilde\kappa(\#X-1)}/c_a(x)$ is the $a$-part of $\|u\|_{x;p}$, so that $A_a\le A$ on $M_A(x;p)$.

We set \[v(Y;u):=\sum_{\operatorname{co}(C)=Y,\ \mathrm{old}}\Phi(C)(1),\qquad v^0(Y):=\sum_{\operatorname{co}(C)=Y,\ \mathrm{new}}\Phi(C)(1)+\sum_{\operatorname{co}(C)=Y,\ \mathrm{constant}}\Phi(C),\] $V_0^+:=\mathsf V_0+\sum_Yv^0(Y)$ and $A(u):=E'^{\rm old}[u]$, and \[G^\delta_X[u](V'):=E_\xi[\chi_X(\xi)\,e^{E^a(X,U)}\,\delta(X,U)],\qquad U=\mathsf U(G\xi-a_0;V').\]

For an analytic cluster $C$ and $V'\in\mathcal U_{\operatorname{co}(C)}(x')$, we let $D\Phi(C)[\delta](V')$ be the derivative $\frac{d}{dz}\big|_{z=0}$ of $\Phi(C)(V')$ evaluated at $u+z(\delta,0,0)$; by Step S9, $\Phi(C)(V')$ so evaluated is an entire function of $z$. We set \[\FitDisplay{\dimexpr\linewidth-1pt\relax}{DA(u)[\delta](Y,V'):=\sum_{\operatorname{co}(C)=Y,\ \mathrm{old}}[D\Phi(C)[\delta](V')-D\Phi(C)[\delta](1)],\qquad Dv(Y;u)[\delta]:=\sum_{\operatorname{co}(C)=Y,\ \mathrm{old}}D\Phi(C)[\delta](1).}\]

For a cluster $C=(\mathfrak B_1,\dots,\mathfrak B_n)$, a finite sequence of polymers as in item~3, we set $\mathrm{mult}_\alpha(C):=\#\{i:\mathfrak B_i\ \text{holds the letter}\ \alpha\}$. The root of an analytic letter $\alpha$ is the first cell of $J_\alpha$.\end{definition}

\begin{theorem}[Item 2: the step at every pair of rates]\label{Y-H1-U1} In the setting \textup{(SD)} of item~2 of \S\ref{Y-H-1}, take the data of a step at any $x\ge x_9^U$, the empty rows included wherever they obey \textup{(SD)}, and let $p\in\mathbb P$, $1\le A\le16$ and $u\in M_A(x;p)$. Then \textup{(U)}, \textup{(b\#)}, \textup{(c+)}, \textup{(f)}, \textup{(g)}, \textup{(h)} and \textup{(h\#)} hold.

\begin{enumerate} \item[\textup{(U)}] The formulas \textup{(Fm1)--(Fm4)}, with $s'-1$ in \textup{(Fm4)} as in item~6 of \S\ref{Y-H-1}, depend neither on $\tilde\kappa$, nor on $\tilde c$, nor on $E_a$, nor on $A$. They define one map $F_\ell$, one number $\mathsf V'$ and one family $v$ on $M_{16}(x;\kappa,c_h/2)$, and all their series converge absolutely. The clauses \textup{(a)--(e)} of \Cref{Y-H-C9} hold at $p$, with $\tilde\kappa$ and $\tilde c$ in the classes, in the weights $e^{\tilde\kappa(\#Y-1)}$, in the caps and in the distances, and with the numerals $0.35$, $0.248$, $0.4$ of that theorem.

\item[\textup{(b\#)}] \[|E'^{\rm old}[u]|_{\mathrm{cs}}\le[2430\mathsf N_{\rm c}^2L^{-1/2}+L^{-1}]\,A\,c_a(x')\le\vartheta_a\,A\,c_a(x').\]

\item[\textup{(c+)}] Let $\Omega$ be open in $\mathbb C^n$, and let $z\mapsto u_z\in M_{16}(x;p)$ be such that every value $E^a_z(X,V)$ with $V\in\mathcal U_X(x)$, every value $E^h_z(X,V)$ and every value $w_z(X,V)$ is holomorphic on $\Omega$. Then $\mathsf V'_z$, every $v(Y;u_z)$ and every value of $F_\ell(u_z)$ are holomorphic on $\Omega$. Here, as in \textup{(f6)} and in Step S5, the $n$ of $\mathbb C^n$ is any integer $\ge1$, and not the number of sources in entry~2 of item~(0) of \S\ref{Y-H0-1}.

\item[\textup{(f)}] \emph{Locality in the members.} \begin{enumerate} \item[\textup{(f1)}] The objects $E'^{\rm new}$, $v^0$ and $V_0^+$ depend on no member. Both $E'^{\rm old}[u](Y,\cdot)$ and $v(Y;u)$ are functions of $\{E^a(X,\cdot):X\in\mathcal D,\ Y_0(X)\subset Y\}$ alone; hence neither $E'^a$ nor $\mathsf V'$ depends on $E^h$ or on $w$.

\item[\textup{(f2)}] The member $E'^h(Y,\cdot)$ is a function of $\{E^a(X),E^h(X):X\in\mathcal D,\ Y_0(X)\subset Y\}$ and $\{w(X):X\in\mathcal P,\ Y_0(X)\subset Y\}$ alone.

\item[\textup{(f3)}] The member $w'(\mathsf Y,\cdot)$ is a function of $\{E^a(X),E^h(X):X\in\mathcal D,\ Y_0(X)\subset\mathsf Y\}$ and $\{w(X):X\in\mathcal P,\ Y_0(X)\subset\mathsf Y\}$ alone.

\item[\textup{(f4)}] In $F_\ell(u)$ and in $\mathsf V'$, the member $w(X,V)$ enters only at fields $V$ such that $X$ holds a plaquette of $D(V)$, and $E^h(X,V)$ enters only at fields $V$ that are tame on $X^\approx$. Moreover, each of $\mathsf V'$, $E'^h$, $w'$ and $E'^a(Y,V')$, the last for $G$-valued $V'$ tame on $Y^\approx$, depends on $E^a(X,V)$ only at $G$-valued fields $V$ that are tame on $X^\approx$; at the remaining $V'$ of $\mathcal U_Y(x')$, the value $E'^a(Y,V')$ draws on $E^a(X,\cdot)$ on $\mathcal U_X(x)$.

\item[\textup{(f5)}] \emph{Reach.} Let $S$ be a set of fine cubes, $S':={\rm cl}(S)$, $\hat\rho\ge0$ an integer and $\hat\rho':=\lceil\hat\rho/L\rceil+1$. If $u_z(X)$, that is, each of the three kinds of member at $X$, is independent of $z$ for every $X$ with $X^{+\hat\rho}$ off $S$, then $F_\ell(u_z)(Y)$ and $v(Y;u_z)$ are independent of $z$ for every $Y$ with $Y^{+\hat\rho'}$ off $S'$.

\item[\textup{(f6)}] Let $z\mapsto u_z$ be a family as in \textup{(c+)}, and let $S_1,\dots,S_n$ be sets such that $u_z(X)$ is independent of $z_i$ whenever $X^{+\hat\rho}$ is off $S_i$. Then, for every set $I$ of indices, the mixed derivative $\prod_{i\in I}\partial_{z_i}$ of $F_\ell(u_z)(Y)$ and of $v(Y;u_z)$ vanishes, unless $Y^{+\hat\rho'}$ meets $S_i'$ for every $i\in I$.\end{enumerate}

\item[\textup{(g)}] \emph{The number.} \begin{enumerate} \item[\textup{(g1)}] $\mathsf V'(u)=V_0^++\sum_Yv(Y;u)$, where $Y$ runs over the connected sets of cells and the sum converges absolutely; and $v(Y;u)=0$ if $E^a=0$.

\item[\textup{(g2)}] $|v(\cdot;u)|_{\mathrm{cs}}\le C_g\,A\,c_a(x)$; besides, $|v^0|_{\mathrm{cs}}\le55\,[2K_v^\flat p_0^3x^{-1/2}+x^{-20}+u_b]$.

\item[\textup{(g3)}] For $\delta\in T^a(x)$, \[Dv(Y;u)[\delta]=\sum_{X:\,Y_0(X)=Y}G^\delta_X[u](1)+r(Y),\qquad|r|_{\mathrm{cs}}\le109\,L^3x^{-1/4}\,B\,c_a(x).\] If $A_a<16$, and so if $A<16$, then $Dv(Y;u)[\delta]=\frac d{dz}|_{0}\,v(Y;u+z(\delta,0,0))$, and, since the torus is finite, $\frac d{dz}|_0\,\mathsf V'(u+z(\delta,0,0))=\sum_YDv(Y;u)[\delta]$. If $u_\tau:=u_0+\tau(\delta,0,0)\in M_{16}(x;p)$ for $0\le\tau\le1$, and $A_a(u_\tau)<16$ for every $\tau<1$, then $v(Y;u_1)-v(Y;u_0)=\int_0^1Dv(Y;u_\tau)[\delta]\,d\tau$.

\item[\textup{(g4)}] For $u_1,u_2\in M_2(x;p)$, \[|v(\cdot;u_1)-v(\cdot;u_2)|_{\mathrm{cs}}\le421\,L^3c_a(x)\,\|E^a_1-E^a_2\|,\] where $\|E^a_1-E^a_2\|=B_p(E^a_1-E^a_2)$, with $B_p$ as in item~1.\end{enumerate}

\item[\textup{(h)}] \emph{The one-letter formula.} Let $\delta\in T^a(x)$. The series $DA(u)[\delta](Y,V')$ converges absolutely and uniformly on $\mathcal U_Y(x')$; it is linear in $\delta$; it has, at $x'$, the clauses of \textup{(ca)} without a size; of $u$ it depends on $E^a$ alone; along a family as in \textup{(c+)} it is holomorphic in $z$; and if $A_a<16$, and so if $A<16$, it equals $\frac d{dz}|_0\,E'^{\rm old}[u+z(\delta,0,0)](Y,V')$. Moreover \[DA(u)[\delta](Y,V')=\sum_{X:\,Y_0(X)=Y}[G^\delta_X[u](V')-G^\delta_X[u](1)]+\mathrm{Rst}(Y,V'),\] where \[\mathrm{Rst}(Y,V'):=\sum_{\operatorname{co}(C)=Y,\ \mathrm{old},\ n(C)\ge2}[D\Phi(C)[\delta](V')-D\Phi(C)[\delta](1)],\] and the following bounds hold. \begin{enumerate} \item[\textup{(h1)}] $|\mathrm{Rst}|_{\mathrm{cs}}\le217\,L^3x^{-1/4}B\,c_a(x)\le366\,L^{5/2}x^{-1/4}B\,c_a(x')$.

\item[\textup{(h2)}] For every cell $\square_0$, \[\FitDisplay{\dimexpr\linewidth-1pt\relax}{\sum_{X\ \mathrm{small},\ \text{first cube under}\ \square_0}\sup_{\mathcal U_{\{\square_0\}}(x')}|G^\delta_X[u](V')-G^\delta_X[u](1)|\le2.4248\cdot361\,\mathsf N_{\rm c}^2e^{Ac_a(x)}B\,c_a(x)\le2430\,\mathsf N_{\rm c}^2L^{-1/2}B\,c_a(x'),}\] and for every cell $\square$, \[\FitDisplay{\dimexpr\linewidth-1pt\relax}{\sum_{X\ \text{not small},\ Y_0(X)\ni\square}e^{\tilde\kappa(\#Y_0(X)-1)}\sup|G^\delta_X[u](V')-G^\delta_X[u](1)|\le e^{-L/2}B\,c_a(x).}\]

\item[\textup{(h3)}] $|DA(u)[\delta]|_{\mathrm{cs}}\le[2430\mathsf N_{\rm c}^2L^{-1/2}+L^{-1}]\,B\,c_a(x')\le\vartheta_a\,B\,c_a(x')$, where the last inequality rests on \textup{(L3su)}; under \textup{(L2$^{\rm su}$)} alone the bracket is $\le0.2026$.\end{enumerate}

\item[\textup{(h\#)}] \emph{The bounds at $\tau^\sharp(x)$.} We have \[\tau^\sharp(x)\le0.12\,x^{-1/4}\] and \[\tau^\sharp(x)\le300\,(2(s_0+\varrho_*)+1)^3(2s'+1)^3e^{9\kappa_1+11}[81L^3E_a+2K_v^\flat]\,p_0^3\,x^{-1/2+\sigma_M};\] moreover \textup{(h1)} holds with $291\,L^3\tau^\sharp(x)$ in the place of $217\,L^3x^{-1/4}$, and \textup{(g3)} with $146\,L^3\tau^\sharp(x)$ in the place of $109\,L^3x^{-1/4}$.\end{enumerate}\end{theorem}

\begin{lemma}[Rooted cluster sum with multiplicities and a power of $\tau$]\label{Y-H1-Rcl} Assume the hypotheses \textup{(SD)}, $x\ge x_9^U$, $p=(\tilde\kappa,\tilde c)\in\mathbb P$ and $1\le A\le16$ of \Cref{Y-H1-U1}. Let $e_\alpha$ be the sizes of Step~6 at $(\tilde\kappa,A)$, let $\alpha$ be an analytic letter with root $\square_0$, let $\square$ be a cell and let $\tau\in[\tau^\sharp(x),1]$. Then \[\sum_{C\ \mathrm{analytic},\ \operatorname{co}(C)\ni\square}\mathrm{mult}_\alpha(C)\,\hat\Phi(C)\,e^{\tilde\kappa(\#\operatorname{co}(C)-1)}\,\tau^{-(n(C)-1)}\le e^{-{\rm dist}(\square_0,\square)-1}\,e_\alpha\,e^{2\varepsilon_2^\sharp\#I_\alpha+(\tilde\kappa+1)\#J_\alpha-\tilde\kappa}.\]\end{lemma}

\begin{lemma}[Item 3: the cluster bound with powers of $\tau$ and with multiplicities]\label{Y-H1-LK} Let a finite alphabet be given, and to each letter $\alpha$ a non-empty finite set $I_\alpha$ of cells. Let $\mathsf r\ge0$, and let $\alpha\sim\beta$ if and only if ${\rm dist}(I_\alpha,I_\beta)\le\mathsf r$. Let ``compatible'' be a symmetric relation on pairs of distinct letters under which incompatible letters are $\sim$. Let $e_\alpha\ge0$, $d_\alpha\ge0$, $0<\tau\le1$ and $\epsilon>0$ be numbers such that for every cell $\square$ \begin{equation}\tag{S$_\tau$}\sum_{\beta:\ {\rm dist}(I_\beta,\square)\le\mathsf r}\tau^{-1}e_\beta\,e^{3\epsilon\#I_\beta+d_\beta}\le\epsilon.\end{equation} The polymers are the non-empty sets of pairwise compatible letters that are $\sim$-connected, and two polymers touch when some letter of the first is $\sim$ to some letter of the second. The clusters, as in Step~8 of the proof of \Cref{Y-G-ThA} (Lems.~\ref{Y-G-LKP} and~\ref{Y-G-L18}), are the finite sequences $C=(\mathfrak B_1,\dots,\mathfrak B_n)$, $n\ge1$, of polymers such that the graph of touching on $\{1,\dots,n\}$ is connected, and \[\varphi^T(C):=\frac1{n!}\sum_g(-1)^{\#g},\] where $g$ runs over the connected spanning subgraphs of that graph and $\#g$ is the number of the edges of $g$. We set $\hat\Phi(C):=|\varphi^T(C)|\prod e_\alpha$, we let $n(C)$ be the number of the letters of the polymers of $C$, and we set $d(C):=\sum d_\alpha$; here products and sums run over these letters, and both in them and in $n(C)$ every letter is counted with multiplicity.

Then for every letter $\alpha$ \[\sum_C\mathrm{mult}_\alpha(C)\,\hat\Phi(C)\,e^{d(C)}\,\tau^{-n(C)}\le\tau^{-1}e_\alpha\,e^{2\epsilon\#I_\alpha+d_\alpha};\] consequently, for every $m\ge1$, the part of $\sum_C\mathrm{mult}_\alpha(C)\hat\Phi(C)e^{d(C)}$ with $n(C)\ge m$ is $\le\tau^{m-1}e_\alpha\,e^{2\epsilon\#I_\alpha+d_\alpha}$.

Suppose further that a complex number $K(\mathfrak B)$ is attached to each set $\mathfrak B$ of pairwise compatible letters, with $K(\emptyset)=1$, with $K(\mathfrak B\cup\mathfrak B')=K(\mathfrak B)K(\mathfrak B')$ whenever $\mathfrak B$ has no letter that is $\sim$ to a letter of $\mathfrak B'$, and with $|K(\mathfrak B)|\le\prod_{\mathfrak B}e_\alpha$. Then \[\sum_{\mathfrak B}K(\mathfrak B)=\exp\sum_C\Phi(C),\qquad\Phi(C)=\varphi^T(C)\prod K(\mathfrak B_i),\] and the series converge absolutely.\end{lemma}

\begin{remark}[\Cref{Y-H1-LK} beside \Cref{Y-G-L18}] Taken at $(\tau,\epsilon)=(t,\varepsilon_2)$ and with $m=1,2$, the claims of \Cref{Y-H1-LK} are the two claims of \Cref{Y-G-L18} on which the proof of \Cref{Y-G-ThA}, in its Step~9 on convergence, and the proof of \Cref{Y-H-C9}, in Step~8, rely. The definitions in \Cref{Y-H1-LK} are those of Step~8 of the proof of \Cref{Y-G-ThA} (Lems.~\ref{Y-G-LKP} and~\ref{Y-G-L18}). The paragraph ``Clusters'' in item~6 of \S\ref{Y-H-1} takes from \Cref{Y-G-L18} the clusters as finite sequences, $K(\mathfrak B)$, $\Phi$, $\hat\Phi$, $n(C)$ and $d(C)$, and with them the polymers as sets of letters, the relation ``touch'' and $\varphi^T$; these are the ones defined in \Cref{Y-H1-LK}, for the alphabet $\mathfrak A$ with $\mathsf r=r_1$, and the polymers of the kinds p and p$'$ of the zone letters are not among them. The condition \textup{(S)} of Step~7\,(ii) is \textup{(S$_t$)} at $\epsilon=\varepsilon_2$.

Compared with \Cref{Y-G-L18}\,(1), (2), \Cref{Y-H1-LK} brings two additions: the factor $\mathrm{mult}_\alpha(C)$ takes the place of $1\{C\ni\alpha\}$, and the power $\tau^{-n(C)}$ is added; indeed, if \Cref{Y-G-L18} is taken at the sizes $e_\alpha/\tau$, with $t=1$ and $\varepsilon_2=\epsilon$, its condition \textup{(S)} becomes \textup{(S$_\tau$)}. The last claim of \Cref{Y-H1-LK} is \Cref{Y-G-L18}\,(4). The display \textup{(An0)} of \Cref{Y-G-An0}, \Cref{Y-G-L18}\,(3), (4) and \Cref{Y-G-L19} are used in Step~7\,(i) and in Step~10\,(8a) and~(8d), at $p$ as well; \Cref{Y-H1-LK} covers only \Cref{Y-G-L18}\,(1), (2) and~(4).\end{remark}

\begin{remark}[Item 4: what is used] The constants and the floors of \Cref{Y-H1-U1} are free of $A$, $u$, $z$, $\delta$, $N$, $K$, $\mathsf b$ and $f$, and of $\ell$ they involve only the pure numbers $b_*$, $K_*$, $\lambda_*$ and the range. \textup{(maj)} Each majorant behind \textup{(b\#)}, \textup{(g2)}, \textup{(g3)}, \textup{(h1)}, \textup{(h2)}, \textup{(h3)} is a function of $A$ and $B$ only: these bounds stay true with the supremum over $u\in M_A(x;p)$, and over all $\delta$ with $B_p(\delta)\le B$, placed inside the cell sum.

\Cref{Y-H1-U1}, like \Cref{Y-H-C9}, is proved in the setting \textup{($\Gamma$2)--($\Gamma$7)}, with the tools of item~4 of \S\ref{Y-H-1}. Its clause \textup{(U)} is not obtained by applying the statement of \Cref{Y-H-C9} at each pair, since in item~1, \textup{(O5)}, of \S\ref{Y-H-1} the number $\kappa^\sharp$ is tied to the single letter $\kappa$ of that section, since \textup{(O11)} has a single $c_h$, and since that statement speaks of one rate at a time. Clause \textup{(U)} is obtained instead in part~(3) of Step S2, by repeating the proof of \Cref{Y-H-C9} step by step at $p$: the letter $\kappa$ of \S\ref{Y-H-1} and $c_h$ occur only in Steps~0, 1, 6 to~10 and~12 of that proof.

In the proof of \Cref{Y-H1-U1} we also use, in the manner of clauses of a statement, individual steps from the proof of \Cref{Y-H-C9}, namely Step~0\,(iii) and (iv), Step~2\,(i) and (ii), Step~3\,(V) and (VI), Step~6\,(iii) and (iv), Step~7\,(ii), Step~8 with \textup{(Clu)}, with ``The clauses without a size'' and with the setting of \textup{(8.2)}, Step~10\,(8d), and Step~11. The proof of \Cref{Y-H1-U1} thus uses \S\ref{Y-H-1} together with the proof of \Cref{Y-H-C9}, and not the statement of that theorem alone; Step S13 lists the statements that this section and the following ones use in addition to those of \S\ref{Y-H-1} and of Step~8 of the proof of \Cref{Y-G-ThA}.\end{remark}

\begin{proof}[Proof of \Cref{Y-H1-U1}] \emph{Step S1: the floor $x_9^U$, its worst corners, its existence and its place in the table.} We show two things. First, each inequality of \textup{(Y.0)--(Y.9)} holds, at every $p\in\mathbb P$ and every $x\ge x_9^U$, when it is taken with $\tilde\kappa$ in the place of the letter $\kappa$ of \S\ref{Y-H-1} and with $e^{-\tilde cp_0^2}$ in the place of $c_h(x)$, while $\kappa^\sharp$ is kept. Second, $x_9^U$ exists and is determined by \textup{(O0)--(O12)} alone; each of the conditions \textup{(i)--(v)} of \textup{(O13)} that define it holds for all large $x$; and $x_9^U$ is free of $A$, $u$, $z$ and $N$, and of all of $\ell$ but the pure numbers $b_*$, $K_*$, $\lambda_*$ and the range.

\emph{Where the two letters occur.} Apart from its place inside $\kappa^\sharp$, the letter $\kappa$ of \S\ref{Y-H-1} occurs in \textup{(Y.3)}, in $e^{(\kappa+1)s'}$ and $e^{\kappa+3}$; in \textup{(Y.4)}, in $e^{\kappa+5}$, $e^{(\kappa+4)s'}$ and $e^{(\kappa+1)s'}$, and in $(\kappa+4)(1+2s')$ in the exponent of $\upsilon_f$, hence also in $\upsilon_*=(\max\{u_t^{1/2},\upsilon_f^{1/4}\})^{1/L^3}$; in the first inequality of \textup{(Y.6)}, in $e^{\kappa+3}$; and in the first inequality of \textup{(Y.9)}. The number $c_h$ occurs in \textup{(Y.1)}, \textup{(Y.3)} and \textup{(Y.4)} through $c_h(x)$, and in the first inequality of \textup{(Y.5)}. Neither letter occurs in any other inequality, and $u_t$, $u_b$, $u_1$, $K_1^\sharp$, $K_v^\flat$, $\delta_3$ are numbers of the table, defined once.

\emph{Directions.} In \textup{(Y.3)}, \textup{(Y.4)} and the first inequality of \textup{(Y.6)}, each term of the left side is non-decreasing in the letter $\kappa$ of \S\ref{Y-H-1}, and so are $\upsilon_f$ and hence $\upsilon_*$; the worst end is $\kappa_1$. The left side of the first inequality of \textup{(Y.9)} is \[2.001\,K_1^\sharp(2r'+1)^3\exp\{1-26.7r'-\kappa(r'-1)\},\] where $\kappa$ is the letter of \S\ref{Y-H-1}; as $r'\ge1$, it is non-increasing in that letter, and the worst end is the lower end $\kappa$. The factors $e^{-\tilde cp_0^2}$ and $e^{-\tilde c(p_0^2-p_0'^2)}$ grow as $\tilde c$ falls, so the worst end is $c_h/2$. Since every inequality holds at every $p$ once it holds at the worst ends, and since \textup{(i)--(iv)} of \textup{(O13)} take the inequalities there, these give every inequality at every $p$.

\emph{Existence.} The argument on orders of Step~0\,(i) applies word for word and shows that each inequality of \textup{(i)--(v)} holds for all large $x$; at the new places its inputs are the following. By \textup{(M2)} at $\kappa_1$, \[(\kappa_1+4)s'\le(\kappa_1+4)(160c_R\log x/M+8)\le\frac1{12}\log x+8(\kappa_1+4).\] Hence the exponent of $x$ in the terms of \textup{(Y.3)} with $c_a$ and with $K_v^\flat$ is $\frac14+\frac1{12}-\frac12=-\frac16<-\frac18$, and in the third term of \textup{(Y.4)} it is $\frac1{12}-\frac12=-\frac5{12}<-\frac14$. The factor $e^{-(c_h/2)p_0^2}=\exp\{-(c_h/2)A_0^2(\log x)^{2\mathfrak p_0}\}$ still decays faster than every power of $x$. As in Step~0\,(i), $\log(1/\upsilon_f)\ge\frac38(\frac12-d_0)A_0^2(\log x')^{2\mathfrak p_0}-O((\log x)^3)$, because $(\kappa_1+4)(1+2s')=O(\log x)$; so $\log(1/\upsilon_*)\ge cL^{-3}(\log x)^{2\mathfrak p_0-3}$ for a number $c>0$, where $2\mathfrak p_0-3>1$. Further, $(c_h/2)(p_0^2-p_0'^2)\ge c_h\mathfrak p_0A_0^2\log L\,(\log x')^{2\mathfrak p_0-1}\to\infty$. In the first inequality of \textup{(Y.9)}, $r'\to\infty$, and in \textup{(Y.10)}, $x^{-1/4}\to0$.

The inequality \textup{(Y.10)} is an additional condition: the first inequality of \textup{(Y.6)} bounds $1800L^{7/2}x^{-1/4}$ only by $1.8\cdot10^{-3}e^{-(\kappa_1+3)}L$, a bound that is larger than $1$ once $L>556\,e^{\kappa_1+3}$.

\emph{Place.} The conditions depend on $x$ and \textup{(O0)--(O12)}; the floor $x_{\rm fl}$ is free of $c_h$, $E_a$, $\ell$ and $N$ (item~G of \S\ref{Y-F-1}, on floors, and Step~13 of the proof of \Cref{Y-F-C7}); and, besides the data, only the floors $x_3$, $x_4$ and $x_s$ of item~(1) of \S\ref{Y-H0-1} depend on $x_9^U$.

\emph{The rest of Step~0.} \textup{(Geo)}, \textup{(Rad)}, $T>1$, $r_\Xi\ge\frac34(\frac12-d_0-c_\Xi)$ and the clauses on $N$ contain neither letter. Applied at $\kappa_1$, \textup{(Y.4)} gives $u_t,\upsilon_f,\upsilon_*<1$ and $2\upsilon_f^{1/4}\le1$, and these hold for the $\upsilon_f$ at every $\tilde\kappa\le\kappa_1$.

\emph{Step S2: one map, and \Cref{Y-H-C9} at every pair.} We show that \textup{(Fm1)--(Fm4)} are free of $\tilde\kappa$, $\tilde c$, $E_a$ and $A$; that, for every $p\in\mathbb P$, the clauses \textup{(a)--(e)} of \Cref{Y-H-C9} hold on $M_A(x;p)$ with the numerals $0.35$, $0.248$, $0.4$; and that the classes are nested as stated in item~0.

\begin{enumerate} \item[\textup{(1)}] \emph{The formulas.} We examine the formulas in item~6 of \S\ref{Y-H-1}. The functions of the letters \textup{(a)}, \textup{(h)}, \textup{(v)}, \textup{(t)}, \textup{(b)} and of the zone letters are made of $u$, $\chi$, $U$, $A_\square$, $V_\square$, $T_Y$, $b_\square$, $\hat\lambda$, $1_{\mathrm{ind}}$ and $n_j$. The sets $J_\alpha$ and $I_\alpha$, the relations $\sim$, ``adjacent'', ``compatible'' and ``actual'', the words ``on'', ``off'', ``short'' and ``long'', and $\operatorname{co}(C)$, ``small'' and $Y_0$ are determined by $r_1$, $s_0$, $s'$, $r'$, $s(x)$, $L$ and the fixed orders; and $\Phi=\varphi^T\prod K$. The letter $\kappa$ of \S\ref{Y-H-1}, as well as $c_h$, $E_a$ and $A$, occurs only in $d_\alpha$, $e_\alpha$, $\hat\Phi$, $\hat w$, $\upsilon_f$ and $\upsilon_*$, which are majorants and appear in no formula. Hence \textup{(Fm1)--(Fm4)} are the same expressions at every $p$.

\item[\textup{(2)}] \emph{The inputs.} \Cref{Y-H-C9} is derived from the clauses that item~4 of \S\ref{Y-H-1} writes out. These contain the letter $\kappa$ of that section only through $\kappa^\sharp$, namely in \textup{(4a)}, \textup{(4b)} and in $e_Y$, and they contain no $c_h$. The objects are fixed at $\kappa_1$, and $\kappa^\sharp=10(\kappa_1+4)\ge10(\tilde\kappa+4)$. It follows from the clauses of item~4 of \S\ref{Y-H-1} that $\mathsf D_k(Y)$, $S_k(Y)$ and $\hat\lambda$ constitute a single family of objects at fixed $(M,\kappa^\sharp)$. Since $\kappa^\sharp$ is determined by $\kappa_1$, it does not depend on $p$.

\item[\textup{(3)}] \emph{The proof of \Cref{Y-H-C9} at $p$.} We repeat that proof at $p$, with $\tilde\kappa$ in the place of the letter $\kappa$ of \S\ref{Y-H-1} in the classes \textup{(ca)}, \textup{(ch)}, \textup{(c$\Xi$)}, in the weights, in $d_\alpha=(\tilde\kappa+1)(\#J_\alpha+s')$ and in $e^\circ$, $\hat w$, $\upsilon_f$, $\upsilon_*$, with $e^{-\tilde cp_0^2}$ in the place of $c_h(x)$, and keeping $\kappa^\sharp$. We treat its steps in order.

Steps~1, 2 and~3 involve neither the letter $\kappa$ of \S\ref{Y-H-1} nor $c_h$, except that Step~1 uses them in the statement that $\rho[u;\ell]$ is bounded, and there alone. Steps~4 and~5 involve no member, and $\upsilon_f,\upsilon_*<1$ by Step S1. In Step~6, \textup{(Y.1)} enters at $c_h/2$, and the bound $\sum_{Y\ni\square}e_Y^{1/2}\le x^{-20}\sum_n(26e)^{n-1}e^{-50n}$ rests on $\kappa^\sharp/2\ge50$, which holds, since $\kappa^\sharp/2\ge55$.

In Step~7, and throughout the discussion of that step, $\kappa$ denotes the letter of \S\ref{Y-H-1} as a running variable. We have $52e/(1-52e\cdot e^{-6})\le217.59$ and $(1-26e\cdot e^{-6})^{-1}\le1.2124$, and the left sides are the values at $\kappa=6$ of expressions that decrease in $\kappa$. The exponents per cube, $4.952-\kappa+(\kappa+5)432/L\le-1.04$ and $4.2581-\kappa+(\kappa+3)432/L\le-1.7$, decrease in $\kappa$ for $L>432$.

For $X$ not small, the prefactors also contain $\kappa$, namely $e^{9\kappa+40}$ in Step~7\,(i) and $e^{8(\kappa+3)}$ in Step~7\,(ii), and $\kappa$ is chosen after $L$; so we keep the whole exponent. Measured against $e^{\kappa+5}$ for the kinds p and p$'$, and against $e^{\kappa+3}$ for \textup{(a)} and \textup{(h)}, the term with $n>L/2-2$ has an exponent at most \[n(4.952+2160/L)+35-\kappa[n(1-432/L)-8]\quad\text{for the kind p},\] \[4.2581(n-1)+35+2160n/L-\kappa[n(1-432/L)-8]\quad\text{for the kind p$'$},\] \[4.2581(n-1)+21+1296n/L-\kappa[n(1-432/L)-8]\quad\text{for \textup{(a)} and \textup{(h)}}.\] When $n>L/2-2$ and $L>10^{13}$, the bracket is positive; so each exponent decreases in $\kappa$, and its largest value, at $\kappa=6$, is $\le-1.04n+83$. With the geometric factor, which is $\le1.55$, each sum is therefore $\le e^{-0.52L+90}\le e^{-L/4}$, uniformly in $\tilde\kappa\ge6$. In Step~7, the bound for the letters of kind \textup{(t)} uses $\kappa^\sharp/2=5\kappa_1+20\ge5\tilde\kappa+20$, and Step~7\,(i) is \textup{(Y.4)} and Step~7\,(ii) is \textup{(Y.3)}, both at $p$.

In Step~8, \textup{(Clu)} and \textup{(R)} are pure algebra in the letter $\kappa$ of \S\ref{Y-H-1}; $\tilde{\mathsf D}$ rests on $\kappa^\sharp-\tilde\kappa\ge9\tilde\kappa+40$; \textup{(8.3)} rests on $\kappa-4.2581\ge0.29\kappa$ for the running letter $\kappa$, which is true for $\kappa\ge6$, since at $6$ it reads $1.7419\ge1.74$ and the slope of the left side is $1$ against $0.29$; \textup{(8.4)} rests on the first inequality of \textup{(Y.6)}; \textup{(8.1)} rests on \textup{(O12)} and \textup{(Y.6)}; and the numerals $0.34438\le0.35$ and $0.2457+0.001+0.001\le0.248$ are unchanged. In Step~9, the number $c_h$ is used at two places: in the first inequality of \textup{(Y.5)}, which enters at $\tilde c$, and in the inequality $\frac12e^{-\delta_3p_0^2/2}\le e^{-\tilde cp_0'^2}/2$, which requires $\tilde c\le\delta_3/2$ and only becomes easier as $\tilde c$ falls.

In Step~10, the exponent per cell of a long $Y$ is $4.2581-\kappa^\sharp+\tilde\kappa+1\le-(9\tilde\kappa+34.74)$; with $s'=8r'$ and $\#Y\ge r'+1$, the sum is $\le(2r'+1)^3\#\mathrm{mk}'_j\cdot1.001K_1^\sharp e^{\tilde\kappa+1}e^{-(\tilde\kappa+26.74)r'}$, which the first inequality of \textup{(Y.9)} at $\tilde\kappa$ bounds; the second inequality of \textup{(Y.4)} enters at $\tilde\kappa$; and $|\tilde{\mathsf D}(Y)|\le1/2$ by the second inequality of \textup{(Y.9)}. Steps~11 and~12 use no floor, and the Cauchy bounds from \textup{(b)} and \textup{(c)} at $p$ give $16\cdot0.248/10=0.3968$, $0.1$ and $0.1$. Every floor that is used is an inequality of \textup{(Y.0)--(Y.9)} at $p$, and so it holds by Step S1.

\item[\textup{(4)}] \emph{Inclusions.} For $n\ge0$ we have $e^{-\kappa_1n}\le e^{-\tilde\kappa n}\le e^{-\kappa n}$, and $e^{-c_hp_0^2}\le e^{-\tilde cp_0^2}\le e^{-(c_h/2)p_0^2}$; hence $M_A(x;\kappa_1,c_h)\subset M_A(x;p)\subset M_A(x;\kappa,c_h/2)$.

\item[\textup{(5)}] \emph{One map.} Part~(3), taken at $(\kappa,c_h/2)$, shows that all series converge absolutely on $M_{16}(x;\kappa,c_h/2)$; by part~(1), restricting them to $M_{16}(x;p)$ gives the map, the number and the family $v$ at $p$.

\item[\textup{(6)}] \emph{Rows.} The rows enter the proof only through \textup{($\Gamma$2)--($\Gamma$7)}, through $\mathrm{En}_\ell$ and through the upper bound $\|\ell^{V*}\ell^V\|_\infty\le(K_*\mathsf g_*)(3K_*\mathsf g_*)$ of Step~6. No quantity is divided by a norm of the rows, and the argument does not depend on any row being non-zero.\end{enumerate}

\emph{Step S3: item \textup{(VII)} of the alphabet, and \textup{(f1)} to \textup{(f4)}.} We prove the following further item of the alphabet and deduce \textup{(f1)--(f4)} from it. \textup{(VII)} The member $u$ at $X$ enters a letter $\alpha$ only if $Y_0(X)\subset J_\alpha$: the letter \textup{(a)}$_X$ involves $E^a(X,\cdot)$; the letter \textup{(h)}$_X$ involves $E^a(X,\cdot)$ and $E^h(X,\cdot)$; a zone letter $(\mathsf J,t,\mathfrak X,\mathfrak Y)$ involves $w(X,\cdot)$ with $X\in\mathfrak X$, where $1_{\mathrm{ind},\zeta}=1$, and $1E(X)$ with $X\in\mathfrak Y$; while \textup{(v)}, \textup{(t)}, \textup{(b)}, $\Lambda_\zeta$, $n_j$, $\tilde{\mathsf D}$ and $P$ depend on no member.

Item \textup{(VII)} is read off item~6 of \S\ref{Y-H-1}, from $\lambda^a_X$, $\lambda^h_X$ and $\lambda_\zeta=\Lambda_\zeta\prod e^{n_j}\prod_{\mathfrak X}w(X,\allowbreak U)\prod_{\mathfrak Y}(e^{1E(X)(U)}-1)$: for \textup{(a)} and \textup{(h)}, $J=Y_0(X)$, and $J_\zeta$ contains the sets $Y_0(X)$ of its polymers.

\begin{enumerate} \item[\textup{(f1)}] The set $\operatorname{co}(C)$ contains every $J_\alpha$ of $C$, and an analytic cluster contains letters \textup{(a)}, \textup{(v)} and \textup{(t)} only. Therefore an old $C$ with $\operatorname{co}(C)=Y$ depends on the $E^a(X,\cdot)$ with $Y_0(X)\subset Y$ and on no other part of $u$; a new or a constant $C$ depends on no member; and the same holds for $\tilde{\mathsf D}$, $P$ and $\mathsf V_0$. Together with \textup{(Fm1)} and \textup{(Fm2)}, this gives \textup{(f1)}.

\item[\textup{(f2)}] In \textup{(Fm3)}, a hard $C$ with $\operatorname{co}(C)=Y$ is made of letters of $\mathfrak A$ with $J_\alpha\subset Y$, and so \textup{(VII)} gives \textup{(f2)}.

\item[\textup{(f3)}] In \textup{(Fm4)}, the union of the cores, which are the $J_\alpha$ of the letters of the rough polymers, the sets $\operatorname{co}(C)$ and the long $Y$, equals $\mathsf Y$. Moreover $\Phi^\natural(C)$ is one of $\Phi(C)(V')$, $\Phi(C)(1)$ and $0$, and which one it is depends on $V'$ alone, as do the marks and the attributes ``actual'', ``on'', ``off'' and ``long''. The argument rests neither on the linking distance $s'$ nor on the clause with $s'-1$, and the mean over $e$ depends on no member.

\item[\textup{(f4)}] The member $w$ occurs in zone letters alone, with the factor $1_{\mathrm{ind},\zeta}$, which vanishes unless each $X$ of $\mathfrak X$ holds a plaquette of $D(U)$. The member $E^h$ occurs in \textup{(h)} and in $1E(X)$, each time with the factor $1_{X^\approx}(U)$. The member $E^a$ carries the factor $1_{X^\approx}(U)$ wherever it occurs in \textup{(h)}, in $1E(X)$, or in the formula of \textup{(a)} at a pair $(V',e)$; it is these formulas that appear in $K(\gamma)$ of \textup{(Fm4)} and in \textup{(Fm3)}. In $\Phi^\natural(C)$ for analytic $C$, and in $\mathsf V'$, the field either equals $1$ or is tame on $\operatorname{co}(C)^\approx$; in $E'^a(Y,V')$ at a $G$-valued $V'$ tame on $Y^\approx$ the field is tame as well; and there the two formulas of \textup{(a)} coincide by Step~3\,(VI), so that $1_{X^\approx}(U)=1$ on $\chi_X=1$. At the other fields of $\mathcal U_Y(x')$ we use the holomorphic formula $\chi_X(e^{E^a(X,U)}-1)$, with $U$ on $X^\approx$ in $\mathcal U_X(x)$ by Step~2\,(ii); and at those of them that are $G$-valued, the remark that follows shows that $U$ is $G$-valued and tame on $X^\approx$.\end{enumerate}

\begin{remark}[The $G$-valued fields $V'$ of $\mathcal U_Y(x')$] In the proof of Step~2\,(i), the bound ``$\phi'_P\le\varepsilon'$'' enters only through $|\partial V'(P)-1|\le2\varepsilon'$, and this is the hypothesis of \textup{($\Gamma$3)} at $\hat\varepsilon=\varepsilon'$. A $G$-valued $V'$ of $\mathcal U_Y(x')$ satisfies \textup{(e1)} and \textup{(e2)}, and by \textup{(u3)} it has $|\partial V'(P)-1|<2\varepsilon'$ at the plaquettes of $Y^\approx$; moreover, for $Y_0(X)\subset Y$, the boxes of radius $R_3$ about $y(p)$, with $p$ a plaquette of $X^\approx$, lie in $Y^\approx$, by the second inequality of \textup{(Y.7)}, as in Step~2\,(ii). For such a $V'$ we use \Cref{Y-H-L51} in its first case, with the parameter $\tau=1$ of that lemma: $|V'(c)|=1$, $|\partial V'(P)-1|<2\varepsilon'$ by \textup{(u3)}, and $X'_P=-i\log\partial V'(P)$ by the convention of item~3 of \S\ref{Y-H-1}. Together with \textup{(D3)}, this gives $0.072\varepsilon$, as in the proof of Step~2\,(i). On $\chi_X=1$ the numbers of that proof are therefore unchanged, namely $(10^{-4}+0.12)\varepsilon$, $0.6154\varepsilon$ and $0.072\varepsilon$, and $2\sin(\phi_p(U)/2)\le0.8075\varepsilon\le2\sin(\varepsilon/2)$. So, at every $G$-valued $V'$ of $\mathcal U_Y(x')$, the field $U$ is $G$-valued and tame on $X^\approx$.\end{remark}

\emph{Step S4: the reach, \textup{(f5)} and \textup{(f6)}.}

We begin with \textup{(f5)}. Let $Y^{+\hat\rho'}$ be disjoint from $S'$, and let $Y_0(X)\subset Y$, where $X\in\mathcal D$ or $X\in\mathcal P$. By item~1 of \S\ref{Y-G0-1}, ${\rm cl}(X)\subset Y_0(X)^{+1}$, and therefore ${\rm cl}(X)\subset Y^{+1}$.

Suppose that some cube $q$ of $S$ is within distance $\hat\rho$ of a cube $q'$ of $X$. In a given coordinate let $i$ and $j$ be the indices of $q$ and $q'$, so that $|i-j|\le\hat\rho$. Because the grids fit (item~1 of \S\ref{Y-G0-1}), a cell is made of exactly $L^3$ fine cubes; so the cells of $q,q'$ carry the indices $\lfloor i/L\rfloor,\lfloor j/L\rfloor$ in this coordinate, and $|\lfloor i/L\rfloor-\lfloor j/L\rfloor|\le\lceil|i-j|/L\rceil\le\lceil\hat\rho/L\rceil$. On the torus we choose lifts of $q,q'$ with $|i-j|\le\hat\rho$; as the grids fit, the index of a cell is $\lfloor\cdot/L\rfloor$ of the lift, taken modulo the side.

The cell of $q'$ lies in ${\rm cl}(X)$, hence in $Y^{+1}$; consequently the cell of $q$ lies in $Y^{+1+\lceil\hat\rho/L\rceil}=Y^{+\hat\rho'}$. It lies in $S'$ as well, which is impossible, since $Y^{+\hat\rho'}$ does not meet $S'$. Therefore $X^{+\hat\rho}$ does not meet $S$, and $u_z(X)$ is constant. By \textup{(f1)} to \textup{(f3)}, $F_\ell(u_z)(Y)$ and $v(Y;u_z)$ are functions of these values $u_z(X)$ alone, which proves \textup{(f5)}.

For \textup{(f6)} suppose that $Y^{+\hat\rho'}$ does not meet $S_i'$ for some $i\in I$; we fix this $i$ and all other variables. Then \textup{(f5)} for $S_i$ shows that the value does not depend on $z_i$, and by \textup{(c+)} it is holomorphic. So its derivative $\partial_{z_i}$ is $0$, and with it the mixed derivative.

\emph{Step S5: holomorphic families in several variables, \textup{(c+)}.} The argument is that of Step~11, with ``holomorphic on $\Omega$, open in $\mathbb C^n$'' in place of ``affine on an open set of $\mathbb C$''. As in \textup{(c+)}, $n$ is any integer $\ge1$.

\textup{(1)} Fix $(\xi,V',e)$. Each letter is then a product of values $w_z(X,U)$, of exponentials of values $E^a_z(X,U)$ and $E^h_z(X,U)$, and of factors in which no member occurs; hence it is holomorphic on $\Omega$. Indeed, at $V'\in\mathcal U_Y(x')$ and where $\chi_X=1$, the field $U$ on $X^{\approx}$ lies in $\mathcal U_X(x)$, and there the value is holomorphic in $z$ by hypothesis.

\textup{(2)} None of the following depends on $z$: the alphabet, the adjacency, the compatibility, the cores, $\chi$, $U$, the indicators, the notions actual, on, off and long, and the alternative of \Cref{Y-H-LP}, which is decided by the supports of the $\Lambda_\zeta$.

\textup{(3)} For fixed $z$ each letter is a Borel function of $\xi$, and for fixed $\xi$ it is continuous in $z$; being a Carath\'eodory function, it is jointly measurable. The majorant of Step~6, at $A=16$ and at $p$, is integrable and independent of $z$. By the theorems of Fubini and Morera, $K_z(\mathfrak B)$ is holomorphic in each variable $z_i$ separately; it is bounded by $\prod e_\alpha$; and Osgood's lemma says that a bounded function which is holomorphic in each variable separately is holomorphic.

\textup{(4)} $\Phi_z(C)$ is a polynomial in the $K_z$, and $|\Phi_z|\le\hat\Phi$ at $A=16$, a bound independent of $z$. Hence the following converge uniformly on $\Omega$: the series \textup{(Fm1)} to \textup{(Fm3)}; the series of $v(Y;\cdot)$; and the sum in \textup{(Fm4)}, since $e^{-\Phi^{\natural}_z}-1$ is holomorphic and the majorant \textup{(8d)} of Step~10 applies. By Weierstrass' theorem the sums are holomorphic.

\textup{(5)} $w'_z(\mathsf Y,V')$ is the mean over $e$ of a bounded, jointly measurable family that is holomorphic in $z$; by the argument of \textup{(3)}, namely the theorems of Fubini and Morera and Osgood's lemma, it is holomorphic.

\emph{Step S6: polymers and clusters, the proof of \Cref{Y-H1-LK}.}

\textup{(1)} \emph{The gas of polymers.} Every set $\mathfrak B$ of pairwise compatible letters has a unique splitting into $\sim$-components. These are polymers, no two of which touch, and $K$ factorises over them. Conversely, let polymers be given, no two of which touch. A letter of one of them and a letter of another are then not $\sim$; so they are distinct, as $\sim$ is reflexive, and compatible, as incompatible letters are $\sim$. Hence the union is pairwise compatible, and its components are the given polymers. It follows that $\sum_{\mathfrak B}K(\mathfrak B)$ is the partition function of the polymers with the weights $K$ and the hard core ``touch'', where a polymer touches itself.

\textup{(2)} \emph{Animals.} We claim: if $v_\beta,a_\beta\ge0$ satisfy $\sum_{\beta\sim\alpha}v_\beta e^{a_\beta}\le a_\alpha$ for all $\alpha$, then \[\FitDisplay{\dimexpr\linewidth-1pt\relax}{\sum_{\mathfrak B\ni\alpha,\ \sim\text{-connected}}\prod_{\mathfrak B}v\le v_\alpha e^{a_\alpha}}\]. Let $S_N(\alpha)$ be the part of this sum over the sets with at most $N$ letters; we argue by induction on $N$, starting from $S_1(\alpha)=v_\alpha$. A connected set $\mathfrak B\ni\alpha$ with at most $N+1$ letters consists of $\{\alpha\}$ together with the components $\mathfrak B_1,\dots,\mathfrak B_k$ of $\mathfrak B\setminus\{\alpha\}$; these are distinct connected sets of at most $N$ letters, and each contains a letter $\beta_i\sim\alpha$. Every unordered family of this kind arises from at least $k!$ tuples $((\beta_i,\mathfrak B_i))_i$. With the induction hypothesis and the assumption we obtain \[S_{N+1}(\alpha)\le v_\alpha\sum_k\frac1{k!}\bigl(\sum_{\beta\sim\alpha}S_N(\beta)\bigr)^k\le v_\alpha\exp\bigl\{\sum_{\beta\sim\alpha}v_\beta e^{a_\beta}\bigr\}\le v_\alpha e^{a_\alpha}.\]

\textup{(3)} We put $a_\beta:=\epsilon\#I_\beta$ and $v_\beta:=\tau^{-1}e_\beta e^{d_\beta+\epsilon\#I_\beta}$. Summing \textup{(S$_\tau$)} over the cells of $I_\alpha$ we get $\sum_{\beta\sim\alpha}v_\beta e^{2a_\beta}\le\epsilon\#I_\alpha=a_\alpha$. A polymer is a $\sim$-connected set, so \textup{(2)} yields $\sum_{\text{polymers }\mathfrak B\ni\beta}\prod_{\mathfrak B}v\le v_\beta e^{a_\beta}$.

\textup{(4)} \emph{The criterion of Koteck\'y and Preiss.} We use it with one marked polymer, that is, \Cref{Y-G-LKP} of Step~8 of the proof of \Cref{Y-G-ThA} in the form that its proof yields. The claim is: if weights $\hat z(\gamma)\ge0$ and numbers $a(\gamma)\ge0$ satisfy $\sum_{\gamma'\text{ touching }\gamma}\hat z(\gamma')e^{a(\gamma')}\le a(\gamma)$ for every $\gamma$, then for every $\gamma_1$ \[1+\sum_{k\ge2}k\sum_{\gamma_2,\dots,\gamma_k}|\varphi^T(\gamma_1,\dots,\gamma_k)|\prod_{j\ge2}\hat z(\gamma_j)\le e^{a(\gamma_1)};\] here $\varphi^T$ is the function of \Cref{Y-G-LKP}, $\frac1{k!}\sum_g(-1)^{\#g}$. Moreover, for all complex weights $\varpi$ with $|\varpi|\le\hat z$, the partition function equals the exponential of the absolutely convergent sum of $\varphi^T\prod\varpi$.

For the proof we use Penrose's bound from the proof of \Cref{Y-G-LKP}: $k!\,|\varphi^T(\gamma_1,\dots,\gamma_k)|\le N_{\rm tr}(\gamma_1,\dots,\gamma_k)$, the number of trees on $\{1,\dots,k\}$ in which every edge joins two touching polymers. It gives \[k\sum_{\gamma_2,\dots,\gamma_k}|\varphi^T(\gamma_1,\dots,\gamma_k)|\prod_{j\ge2}\hat z(\gamma_j)\le\frac1{(k-1)!}\sum_{\gamma_2,\dots,\gamma_k}N_{\rm tr}(\gamma_1,\dots,\gamma_k)\prod_{j\ge2}\hat z(\gamma_j).\] Summed over $k\ge1$, where the term for $k=1$ is $1$, the right sides give the quantity $T(\gamma_1)$ of that proof at $W=\hat z$. The recursion used there shows $T(\gamma)\le e^{a(\gamma)}$: with $T_0=1$, \[T_{i+1}(\gamma)=\exp\{\sum_{\gamma'\text{ touching }\gamma}\hat z(\gamma')T_i(\gamma')\}\le\exp\{\sum_{\gamma'\text{ touching }\gamma}\hat z(\gamma')e^{a(\gamma')}\}\le e^{a(\gamma)}.\] The formula for the partition function is the last claim of \Cref{Y-G-LKP}, applied with the weights $W=\varpi$.

Let $\mathrm{mult}_\gamma(C):=\#\{j:\gamma_j=\gamma\}$. Since $\varphi^T$ and $\prod\hat z$ are symmetric in the sequence, all $k$ places contribute the same sum, and so \[\FitDisplay{\dimexpr\linewidth-1pt\relax}{\sum_C\mathrm{mult}_\gamma(C)|\varphi^T(C)|\prod\hat z=\sum_k\sum_{(\gamma_1,\dots,\gamma_k)}\#\{j:\gamma_j=\gamma\}|\varphi^T|\prod\hat z=\hat z(\gamma)\bigl[1+\sum_{k\ge2}k\sum_{\gamma_2,\dots,\gamma_k}|\varphi^T(\gamma,\gamma_2,\dots)|\prod_{j\ge2}\hat z(\gamma_j)\bigr]\le\hat z(\gamma)e^{a(\gamma)}.}\]

We choose $\hat z(\mathfrak B):=\prod_{\mathfrak B}\tau^{-1}e_\beta e^{d_\beta}$ and $a(\mathfrak B):=\epsilon\sum_{\mathfrak B}\#I_\beta$, so that $\hat z(\mathfrak B)e^{a(\mathfrak B)}=\prod_{\mathfrak B}v$. With this choice \[\sum_{\mathfrak B'\text{ touching }\mathfrak B}\hat z(\mathfrak B')e^{a(\mathfrak B')}\le\sum_{\alpha\in\mathfrak B}\sum_{\beta\sim\alpha}\sum_{\mathfrak B'\ni\beta}\prod_{\mathfrak B'}v\le\sum_{\alpha\in\mathfrak B}\sum_{\beta\sim\alpha}v_\beta e^{a_\beta}\le\sum_{\alpha\in\mathfrak B}a_\alpha=a(\mathfrak B).\]

\textup{(5)} We have $\mathrm{mult}_\alpha(C)=\sum_{\mathfrak B\ni\alpha}\mathrm{mult}_{\mathfrak B}(C)$ and $\hat\Phi(C)e^{d(C)}\tau^{-n(C)}=|\varphi^T(C)|\prod\hat z$. Hence the left side of the bound claimed in \Cref{Y-H1-LK} is at most $\sum_{\mathfrak B\ni\alpha}\hat z(\mathfrak B)e^{a(\mathfrak B)}=\sum_{\mathfrak B\ni\alpha}\prod_{\mathfrak B}v\le v_\alpha e^{a_\alpha}=\tau^{-1}e_\alpha e^{d_\alpha+2\epsilon\#I_\alpha}$. When $n(C)\ge m$, we use $1\le\tau^m\tau^{-n(C)}$. The identity for $\sum_{\mathfrak B}K(\mathfrak B)$ follows from \textup{(1)} and \textup{(4)}, because $|K(\mathfrak B)|\le\prod e_\beta\le\hat z(\mathfrak B)$.

\emph{Step S7: two bounds and a sum, \textup{(M)}, \textup{(R*)} and \textup{($\Sigma$a)}.}

\begin{enumerate}\item[\textup{(M)}] Let $\mathfrak B$ be a set of letters \textup{(a)}, \textup{(v)}, \textup{(t)} with $J_\beta\subset Y$; let $V'\in\mathcal U_Y(x')$, $\textup{(a)}_X\in\mathfrak B$ and $\delta\in T^a(x)$. Then \[E_\xi\bigl|\chi_Xe^{E^a(X,U)}\delta(X,U)\cdot\prod_{\beta\in\mathfrak B,\ \beta\ne\textup{(a)}_X}\lambda_\beta\bigr|\le m_X\prod_{\beta\ne\textup{(a)}_X}e_\beta,\] where \[\FitDisplay{\dimexpr\linewidth-1pt\relax}{m_X:=e^{Ac_a(x)}\sup_{\mathcal U_X(x)}|\delta(X,\cdot)|\,e^{2x^{-1/8}\#I_{\textup{(a)}_X}}}\]; consequently $m_X/e_{\textup{(a)}_X}\le e^{Ac_a(x)}B/(4A)$.

\item[\textup{(R*)}] Take the sizes $e_\alpha$ of Step~6 at $(\tilde\kappa,A)$, and $d_\alpha=(\tilde\kappa+1)(\#J_\alpha+s')$. Then for every analytic letter $\alpha$ with root $\square_0$, every cell $\square$, and $m=1,2$, \[\sum_{C\text{ analytic},\ \operatorname{co}(C)\ni\square,\ n(C)\ge m}\mathrm{mult}_\alpha(C)\hat\Phi(C)e^{\tilde\kappa(\#\operatorname{co}(C)-1)}\le t^{m-1}e^{-{\rm dist}(\square_0,\square)-1}e_\alpha e^{2\varepsilon_2\#I_\alpha+(\tilde\kappa+1)\#J_\alpha-\tilde\kappa}.\]

\item[\textup{($\Sigma$a)}] For every cell $\square$, \[\sum_{X\in\mathcal D}e^{-{\rm dist}(\square_0(X),\square)-1}e^{-\tilde\kappa(\#X-1)}e^{4x^{-1/8}\#I_{\textup{(a)}_X}+(\tilde\kappa+1)\#Y_0(X)-\tilde\kappa}\le65.642L^3,\] where $\square_0(X)$ denotes the root of $\textup{(a)}_X$.\end{enumerate}

We prove \textup{(M)} by the argument of Step~6\,\textup{(iii)}, in which each factor containing a member is bounded in modulus by a number that does not depend on $\xi$, and these numbers then appear as constant factors in front of $E_\xi$. Where $\chi_X=1$, the field $U$ on $X^{\approx}$ lies in $\mathcal U_X(x)$, by Step~2\,\textup{(ii)}, since $Y\supset Y_0(X)$; on that domain $|E^a(X,\cdot)|\le Ac_a(x)$, and $\delta(X,\cdot)$ is bounded by its supremum. Each remaining letter \textup{(a)} or \textup{(v)} is handled in the same way: its modulus is at most $2Ac_a(x)e^{-\tilde\kappa(\#X'-1)}\le e^{\circ}({\rm a})$ or $2K_v^{\flat}p_0^3x^{-1/2}=e^{\circ}({\rm v})$, respectively, a number independent of $\xi$. What remains is $E_\xi\prod|e^{T_{Y_t}}-1|$, the product running over the letters \textup{(t)}, and Step~6\,\textup{(iii)} bounds this by $\prod e_{Y_t}^{1/2}\cdot\prod e^{x^{-1/8}\#I}$, again over the letters \textup{(t)}.

In fact Step~6\,\textup{(iii)} leaves the factor \[\exp\{[2x\varsigma''n_a+36\varsigma''(C_\varphi^{\flat})^2\cdot3L^3(R+1)^3\cdot v(2r_1+1)^3]\#\mathsf Y\}\le e^{x^{-1/8}\#\mathsf Y},\] with $\mathsf Y$ the union of the $Y_t$; this inequality follows from the third and fourth inequalities of \textup{(Y.8)}, and $\#\mathsf Y\le\sum\#Y_t$.

The factors $e^{2x^{-1/8}\#I}$ in $e_\beta$ and in $m_X$ are at least $1$. For the last assertion of \textup{(M)} we use $e_{\textup{(a)}_X}=4Ac_a(x)e^{-\tilde\kappa(\#X-1)}e^{2x^{-1/8}\#I_{\textup{(a)}_X}}$.

We turn to \textup{(R*)}. The letters \textup{(a)}, \textup{(v)}, \textup{(t)} are pairwise compatible: item~6 of \S\ref{Y-H-1} lists no incompatibility among them. The condition \textup{(S$_t$)} at $\epsilon=\varepsilon_2$, $\mathsf r=r_1$, holds for $\mathfrak A$ at $p$ and every $A\le16$, by Step~7\,\textup{(ii)}, which applies by Step~S2; hence it holds for the smaller alphabet as well. \Cref{Y-H1-LK} bounds the part with $n(C)\ge m$ of $\sum_C\mathrm{mult}_\alpha\hat\Phi e^{d(C)}$ by $t^{m-1}e_\alpha e^{2\varepsilon_2\#I_\alpha+d_\alpha}$. By \textup{(Clu)} of Step~8, $\#\operatorname{co}(C)\le\sum(\#J_\beta+s')-s'$, the sum being taken with multiplicity; thus $e^{\tilde\kappa(\#\operatorname{co}(C)-1)}\le e^{d(C)-(\tilde\kappa+1)s'-\#\operatorname{co}(C)-\tilde\kappa}$. Since $\operatorname{co}(C)$ is connected and contains both $\square_0$ and $\square$, $\#\operatorname{co}(C)\ge{\rm dist}(\square_0,\square)+1$. Finally, $d_\alpha-(\tilde\kappa+1)s'=(\tilde\kappa+1)\#J_\alpha$.

We turn to \textup{($\Sigma$a)}. By Step~3\,\textup{(V)}, $\#I_\alpha\le(2\varrho_*+1)^3\#J_\alpha$; with \textup{(Y.2)} this gives $4x^{-1/8}\#I_\alpha\le0.08\#J_\alpha$. Consider first a small $X$. Then $\#J=1$, and $e^{-1+(\tilde\kappa+1)-\tilde\kappa}=1$. Next, $\sum_{\square_0}e^{-{\rm dist}(\square_0,\square)}\le \mathsf g(1)\le49.9791$; indeed the distance on the torus is the minimum over the lifts, so that the sum over the torus does not exceed the sum over $\mathbb Z^3$. The exact value is $\mathsf g(1)=49.97902\ldots$. The first cube of $X$ is one of the $L^3$ fine cubes that make up $\hat\square_0$, and $\sum_{X\ni\text{cube}}e^{-\tilde\kappa(\#X-1)}\le1.2124$. Altogether the small $X$ contribute at most $49.9791\cdot1.2124\cdot e^{0.08}L^3\le65.6415L^3$.

Now let $X$ be not small, so $n=\#X>L/2-2$. Then $\#Y_0\le8+432n/L$; $\square_0$ meets $X^{+1}$, so $X$ contains one of $(L+2)^3$ cubes; and a given cube lies in at most $(26e)^{n-1}$ sets $X$ with $n$ cubes, where $\log26e\le4.2581$. So these $X$ contribute at most \[49.9791(L+2)^3\sum_{n>L/2-2}\exp\{(4.2581-\tilde\kappa)(n-1)+(\tilde\kappa+1.08)(8+432n/L)-\tilde\kappa-1\}.\] From $\tilde\kappa-4.2581\ge0.29\tilde\kappa$, $1.08\le0.18\tilde\kappa$ and $L>1.2\cdot10^{13}$, each exponent is at most $-0.289\tilde\kappa(n-1)+9.5\tilde\kappa$. Consecutive terms have ratio at most $e^{-1.73}$, and the first exponent is at most $-0.289\tilde\kappa(L/2-36)\le-0.867L+63$. Hence this part is at most $\exp\{-0.867L+70+3\log(L+2)\}\le\frac18e^{-L/2}$, and together with the small $X$ we arrive at \textup{($\Sigma$a)}.

\emph{Step S8: the cluster series, \textup{(g1)} and \textup{(g2)}.} \textup{(g1)} is \textup{(Fm2)} with the clusters sorted by $\operatorname{co}(C)$. The alphabet is finite, and, by \Cref{Y-H1-LK}, the sums of $\hat\Phi$ over the clusters containing a given letter are finite; so all the sums converge absolutely. At $E^a=0$ each $\lambda^a_X$ vanishes identically; therefore $K(\mathfrak B)=0$ for every $\mathfrak B$ containing a letter \textup{(a)}, and $\Phi(C)=0$ for every old $C$.

For the first bound of \textup{(g2)}, the bound on $v(\cdot;u)$, we note that the field $1$ belongs to $\mathcal U_Y(x')$ for every $Y$; so, by Step~6\,\textup{(iii)}, applied with $Y=\operatorname{co}(C)$, $|\Phi(C)(1)|\le\hat\Phi(C)$ for every analytic $C$. Every old $C$ contains a letter \textup{(a)}, and hence \[\sum_{Y\ni\square}e^{\tilde\kappa(\#Y-1)}|v(Y;u)|\le\sum_X\sum_{C\ni\textup{(a)}_X,\ \operatorname{co}(C)\ni\square}\hat\Phi(C)e^{\tilde\kappa(\#\operatorname{co}(C)-1)}.\] We apply \textup{(R*)} for $m=1$, with $e_\alpha e^{2\varepsilon_2\#I_\alpha}=4Ac_a(x)e^{-\tilde\kappa(\#X-1)}e^{4x^{-1/8}\#I_\alpha}$, and then \textup{($\Sigma$a)}: the right side is at most $4\cdot65.642L^3Ac_a(x)\le262.57L^3Ac_a(x)\le C_gAc_a(x)$.

For the second bound of \textup{(g2)}, the bound on $v^0$, we observe that a new $C$ contains a letter \textup{(v)} or \textup{(t)}, and a constant $C$ contains a letter \textup{(b)}; $|\Phi|\le\hat\Phi$ by Step~6\,\textup{(iii)}, \textup{(iv)}; and \Cref{Y-H1-LK} for $\mathfrak A$, together with \textup{(Clu)}, gives \textup{(R*)} for these letters too. For the letters \textup{(v)} and \textup{(b)} we have $\#J=\#I=1$ and $\mathsf g(1)e^{4x^{-1/8}}\le49.9791e^{0.08}\le54.142$, and this number is multiplied by $2K_v^{\flat}p_0^3x^{-1/2}$ for the letters \textup{(v)} and by $u_b$ for the letters \textup{(b)}. For the letters \textup{(t)}, $\sum_n(26e)^{n-1}e^{-(5\tilde\kappa+20)n+(\tilde\kappa+1.08)n-\tilde\kappa-1}\le e^{-40}$, to be multiplied by $49.9791x^{-20}$. Here $e^{\circ}({\rm t})$ carries $\kappa^{\sharp}/2=5\kappa_1+20$; as $\kappa_1\ge\tilde\kappa$, we have $e^{-(5\kappa_1+20)n}\le e^{-(5\tilde\kappa+20)n}$.

\emph{Step S9: the derivative and the one-letter formula, \textup{(h)} up to and including \textup{(h1)}; \textup{(g3)}, \textup{(g4)} and \textup{(maj)}.}

\textup{(0)} \emph{Each term is entire.} Let $\mathfrak B$ be a set of letters \textup{(a)}, \textup{(v)}, \textup{(t)} with $J\subset Y$, and let $V'\in\mathcal U_Y(x')$. When $u$ is replaced by $u+z(\delta,0,0)$, the letter $\textup{(a)}_X$ becomes $\chi_X(e^{E^a(X,U)+z\delta(X,U)}-1)$, an entire function of $z$, bounded on $|z|\le\rho$ by $e^{Ac_a(x)+\rho\sup|\delta(X)|}+1$; the other factors keep the integrable majorant of Step~6\,\textup{(iii)}. Hence $K_z(\mathfrak B)(V')$ is entire, the derivative may be taken under $E_\xi$, and at $z=0$ one has $\frac{d}{dz}\lambda^a_X=\chi_Xe^{E^a(X,U)}\delta(X,U)$. So $\Phi_z(C)=\varphi^T\prod K_z(\mathfrak B_i)$ is entire.

\textup{(1)} \emph{The majorant.} The product rule and \textup{(M)} give \[\FitDisplay{\dimexpr\linewidth-1pt\relax}{|DK(\mathfrak B)[\delta]|\le\sum_{\textup{(a)}_X\in\mathfrak B}(m_X/e_{\textup{(a)}_X})\prod_{\mathfrak B}e_\beta}\]. Consequently, for $V'\in\mathcal U_{\operatorname{co}(C)}(x')$, \[\FitDisplay{\dimexpr\linewidth-1pt\relax}{|D\Phi(C)[\delta](V')|\le\hat\Phi(C)\sum_X\mathrm{mult}_{\textup{(a)}_X}(C)\,m_X/e_{\textup{(a)}_X},}\] a bound in which no field enters. From \textup{(R*)} and \textup{($\Sigma$a)}, with \[\FitDisplay{\dimexpr\linewidth-1pt\relax}{m_Xe^{2\varepsilon_2\#I_{\textup{(a)}_X}}\le e^{Ac_a(x)}Bc_a(x)e^{-\tilde\kappa(\#X-1)}e^{4x^{-1/8}\#I_{\textup{(a)}_X}}}\] and with $e^{Ac_a(x)}\le e^{1/2}$ by \textup{(Y.1)}, we obtain for $m=1,2$ \begin{equation}\tag{$D_m$}\FitDisplay{\dimexpr\linewidth-36pt\relax}{\sum_{C\text{ old},\ \operatorname{co}(C)\ni\square,\ n(C)\ge m}e^{\tilde\kappa(\#\operatorname{co}(C)-1)}\sup|D\Phi(C)[\delta]|\le t^{m-1}e^{1/2}65.642L^3Bc_a(x)\le t^{m-1}108.23L^3Bc_a(x).}\end{equation}

\textup{(2)} Therefore the series of $DA(u)[\delta](Y,\cdot)$, of $Dv(Y;u)[\delta]$, of $\mathrm{Rst}$ and of $r$ converge absolutely and uniformly. Each of them is linear in $\delta$, as $D\Phi(C)[\delta]$ is; and, by \textup{(f1)}, of the components of $u$ only $E^a$ enters them.

\textup{(3)} \emph{The derivative.} Let $A_a<16$. For $|z|B\le16-A_a$ we have $u+z(\delta,0,0)\in M_{16}(x;p)$, because the clauses without a size are linear, $\|\cdot\|_{x;p}$ is the largest of three weighted suprema (item~3 of \S\ref{Y-H-1}), and $z(\delta,0,0)$ changes only the first of them. On this disc $|\Phi_z|\le\hat\Phi$ at $A=16$, so the series of $E'^{\rm old}$ converges uniformly in $z$, and by Weierstrass' theorem its derivative at $0$ is the sum of the derivatives of its terms, which is $DA(u)[\delta]$. For $v$ we use, at the field $1$, the bound $|\Phi_z(C)(1)|\le\hat\Phi(C)$ at $A=16$; then the series of $v$ converges uniformly in $z$, and Weierstrass' theorem shows that $Dv(Y;u)[\delta]$ is $\frac{d}{dz}$ at $0$ of $v(Y;u+z(\delta,0,0))$. By \textup{(g1)} and \textup{(f1)}, and because there are finitely many connected sets $Y$ of cells, $\frac{d}{dz}$ at $0$ of $\mathsf V'(u+z(\delta,0,0))$ equals $\sum_YDv(Y;u)[\delta]$.

\emph{The integral of \textup{(g3)}.} Let $u_\tau$ be as in \textup{(g3)}. By \textup{(0)}, every $\Phi(C)(1)$ defined with $u_\tau$ is entire in $\tau$, its derivative being $D\Phi(C)[\delta](1)$ evaluated at $u_\tau$. On $[0,1]$ the terms are at most $\hat\Phi(C)$ at $A=16$, and their derivatives are bounded in modulus by the majorant of \textup{(1)} at $A=16$; both majorants are independent of $\tau$ and summable, by \textup{(R*)} and \textup{($D_1$)}. We may therefore differentiate term by term: $\tau\mapsto v(Y;u_\tau)$ is continuously differentiable on $[0,1]$, with derivative $Dv(Y;u_\tau)[\delta]$, and we integrate.

\textup{(4)} \emph{Families.} Along the family $u_z$ of \textup{(c+)}, $D\Phi_z(C)[\delta](V')$ is holomorphic in $z$. To see this, note first that, by \textup{(0)} and the product rule, \[DK_z(\mathfrak B)[\delta](V')=\sum_{\textup{(a)}_X\in\mathfrak B}E_\xi\bigl[\chi_Xe^{E^a_z(X,U)}\delta(X,U)\prod_{\beta\in\mathfrak B,\ \beta\ne\textup{(a)}_X}\lambda_{\beta,z}\bigr].\] For fixed $\xi$ the integrand is holomorphic in $z$ on $\Omega$, by \textup{(1)} and \textup{(2)} of Step~S5 and because $\delta$ is independent of $z$; it is also jointly measurable. By the estimate of the proof of \textup{(M)}, applied with $A=16$, the modulus of the integrand is at most $e^{1/2}\sup|\delta(X,\cdot)|$ times the numbers $e^{\circ}$ of the other letters \textup{(a)}, \textup{(v)} times $\prod|e^{T_{Y_t}}-1|$; this is an integrable majorant independent of $z$, also where $\|u_z\|=16$. The theorems of Fubini and Morera and Osgood's lemma, used as in \textup{(3)} of Step~S5, now show that $DK_z(\mathfrak B)[\delta]$ is holomorphic on $\Omega$. Further, $D\Phi_z(C)[\delta]=\varphi^T\sum_iDK_z(\mathfrak B_i)[\delta]\prod_{j\ne i}K_z(\mathfrak B_j)$ is a polynomial in holomorphic functions. Since the majorant of \textup{(1)} at $A=16$ is independent of $z$, \textup{($D_1$)} makes the series of $DA(u_z)[\delta](Y,V')$ converge uniformly in $z$, and Weierstrass' theorem applies.

\textup{(5)} \emph{The split.} An old cluster with $n(C)=1$ is $C=(\{\textup{(a)}_X\})$, the sequence consisting of the single polymer $\{\textup{(a)}_X\}$, and for it $\varphi^T=1$, $\Phi(C)=\mathcal G_X=E_\xi[\chi_X(e^{E^a(X,U)}-1)]$, $\operatorname{co}(C)=Y_0(X)$ and $D\Phi(C)[\delta]=G^\delta_X[u]$. In $DA(u)[\delta]$ the clusters with $n(C)\ge2$ make up $\mathrm{Rst}$, and in $Dv(Y;u)[\delta]$ they make up $r(Y):=\sum_{\operatorname{co}(C)=Y,\ \text{old},\ n(C)\ge2}D\Phi(C)[\delta](1)$.

\textup{(6)} For \textup{(h1)} we apply \textup{($D_2$)} twice, on account of the subtraction; the result is \[216.46L^3x^{-1/4}Bc_a(x)\le217L^3x^{-1/4}Bc_a(x)\]. By \textup{(Y.0)(0b)}, \[\FitDisplay{\dimexpr\linewidth-1pt\relax}{c_a(x')=L^{1/2}(p_0'/p_0)^3c_a(x)\ge0.594L^{1/2}c_a(x)}\], because $2^{-3/4}\ge0.5946$; and $217/0.594\le365.32\le366$. For \textup{(g3)} we apply \textup{($D_2$)} once: $108.23\le109$.

\textup{(7)} For \textup{(g4)} recall that, by \textup{(f1)}, $v$ depends on $E^a$ alone. The family $u_\tau:=(E_1^a+\tau(E_2^a-E_1^a),E_1^h,w_1)$, $0\le\tau\le1$, lies in the convex set $M_2(x;p)$, and $\delta:=E_2^a-E_1^a\in T^a(x)$ with $B=\|E_1^a-E_2^a\|$. By \textup{(3)}, $v(Y;u_2)-v(Y;u_1)=\int_0^1Dv(Y;u_\tau)[\delta]\,d\tau$. Taking moduli and the cell sum under the integral and using \textup{($D_1$)}, we bound this by $108.23L^3c_a(x)B\le421L^3c_a(x)B$.

\textup{(8)} \emph{Clauses without a size.} We treat them as in Step~8. The function $DK(\mathfrak B)[\delta](V')$ is the mean of a product which depends holomorphically on $V'\in\mathcal U_Y(x')$, by Step~2\,\textup{(ii)} and by \textup{($\Gamma$6)}, \textup{(2d)} and \textup{(4b)} of item~4 of \S\ref{Y-H-1}, and which is dominated by \textup{(M)}; the convergence being uniform, holomorphy in $V'$ follows. The dependence on the field on $Y^{\approx}$ alone follows from Step~3\,\textup{(V)}. For the invariance: like $E^a(X,\cdot)$, the function $\delta(X,\cdot)$ is invariant, so the marked letter satisfies $\lambda(\operatorname{Ad}_u\xi,V'^u)=\lambda(\xi,V')$, just as the other letters do; moreover the $\chi$ and the law of $\xi$ are invariant under $\operatorname{Ad}_u$. The invariance and the subtraction give the value $0$ at pure gauges.

\textup{(9)} For \textup{(maj)} we note that the right sides of \textup{(1)}, and the number $S_X$ of Step~S10, depend on $A$ and $B$ alone; hence the supremum over $u$ and $\delta$ may be moved inside every sum above.

\emph{Step S10: the one-letter part, \textup{(h2)}.} Let $X$ be small, with its first cube in $\hat\square_0$; let $o$ be the centre site of $\square_0$, and $\mathsf B_\omega$ the box of radius $\omega$ about $o$. This is the setting of Step~8, \textup{(8.2)}. The blocks meeting $X^{\approx}$ lie within $M+\lceil Mr(x)/L\rceil$ of $o$, so that $\mathsf B_\omega$ contains the set of sites of Step~2\,\textup{(ii)}. A field on $\mathsf B_\omega$ is extended by $1$; it then lies in the interior of $\mathfrak V$, and no quantity of the argument depends on the extension.

Indeed, the values of $\mathsf U$ on $X^{\approx}$ are determined by $V'$ at $y(b)$, by $G$ at distance at most $r_G$ and by $a_0$ at distance at most $R_3$ from the blocks meeting $X^{\approx}$; this follows from \textup{(2b)} and, for $a_0$, from the paragraph ``The dependence at every pair'' of \textup{($\Gamma$3)} of item~4 of \S\ref{Y-H-1}, by which, at every field of $\mathfrak V$, $a_{0,b}$ depends on $V'$ through its values on bonds and on the $\mathfrak g_{\mathbb C}$-valued variables $X'_P$ of plaquettes $P$, every site of these bonds and plaquettes being within $R_3$ of $y(b)$. By the definition of $\omega$, all these sites lie in $\mathsf B_\omega$. Hence, if a member of $\mathcal U_{\{\square\}}(x')$ is extended by $1$ outside the neighbourhood of the cell, $G^\delta_X[u]$ takes at it the same value as at its restriction to $\mathsf B_\omega$, extended by $1$.

\textup{(1)} We apply Step~2\,\textup{(ii)} with $\tau=17T$: for $V'\in U^{17T}(\mathsf B_\omega;x')$ and $\chi_X(\xi)=1$, the field $U$ on $X^{\approx}$ lies in $\mathcal U_X(x)$, and it depends holomorphically on $V'$. Therefore the integrand of $G^\delta_X[u]$ is holomorphic in $V'$ and bounded by $S_X:=e^{Ac_a(x)}Bc_a(x)e^{-\tilde\kappa(\#X-1)}$; so $G^\delta_X[u]$ is holomorphic on $U^{17T}(\mathsf B_\omega;x')$, with $|G^\delta_X[u]|\le S_X$.

\textup{(2)} The function $G^\delta_X[u]$ is invariant under gauge transformations with values in $G$. Indeed, $E^a(X,\cdot)$ and $\delta(X,\cdot)$ are invariant; the gauge transform of $\mathsf U(\eta;V')$ by the blockwise constant $u$ is $\mathsf U(\operatorname{Ad}_u\eta;V'^u)$; $G$ is covariant, by \textup{(2c)} of item~4 of \S\ref{Y-H-1}, and so is $a_0$, by its definition in \textup{($\Gamma$3)} together with \textup{(2c)} and the formula of \textup{(2g)} of that item; $\chi_X$ and the law of $\xi$ are invariant under $\operatorname{Ad}_u$; and a transformation $u$ given on the sites of $\mathsf B_\omega$ is extended by $1$, as in the part of Step~8 on the clauses without a size.

\textup{(3)} The restriction of a member of $\mathcal U_{\{\square_0\}}(x')$ to $\mathsf B_\omega$ has \textup{(u1)} to \textup{(u3)} at $x'$, since $\omega\le Mr'$. We now apply \Cref{Y-H1-LT57} below with $T=L^{3/2}/(19\mathsf N_{\rm c})$, $\Psi=G^\delta_X[u]$ and $S=S_X$. Here $T>1$ and $2\omega+1<N$ hold by Step~0\,\textup{(iv)}, so that $\mathsf B_\omega$ is embedded; $3\omega\le\ell_w(x')$ follows from $3\omega+1\le\ell_w(x')$, which is given by Step~0\,\textup{(iii)} with \textup{(Rad)}; and \textup{(C0)} to \textup{(C2)} are, for this $T$, the conditions \textup{(Y.0)(0d)}, which are satisfied since $x\ge x_9^U$ and \textup{(Y.0)} is one of the conditions of \textup{(O13)}\,\textup{(i)} of item~0 of this section. Thus the hypotheses of the lemma are met; and its conclusion says that any complex $\Psi$ which is holomorphic, $G$-invariant and bounded by $S$ on $U^{17T}(\mathsf B_\omega;x')$ satisfies $|\Psi(V')-\Psi(1)|\le2ST^{-2}$ at every such $V'$. Hence $|G^\delta_X[u](V')-G^\delta_X[u](1)|\le2T^{-2}S_X$.

The conclusion of \Cref{Y-H1-LT57} is that of \Cref{Y-H-Th57}; the theorem assumes in addition $17T\varepsilon'\le1/40$ and $3\omega+1\le\ell_w(x')$.

\begin{lemma}[Invariant holomorphic functions on a box]\label{Y-H1-LT57} Let $T>1$ be arbitrary. Let $\mathsf B_\omega$ be the box of radius $\omega$ about a coarse site $o$, embedded, that is, $2\omega+1<N$, and let $3\omega\le\ell_w(x')$. Let the following conditions hold; for $T=L^{3/2}/(19\mathsf N_{\rm c})$ they are the conditions \textup{(Y.0)(0d)}.

\begin{enumerate}\item[\textup{(C0)}] $18\omega\varepsilon'\le1/10$ and $160T\omega\varepsilon'\le1$;

\item[\textup{(C1)}] $2640\,\ell_w(x')(T+1)\omega\varepsilon'\le \mathsf g_x$;

\item[\textup{(C2)}] $102400\,T\omega^2\varepsilon'\le1$.\end{enumerate}

Let $\Psi$ be a complex function on $U^{17T}(\mathsf B_\omega;x')$ that is holomorphic, bounded by $S$, and invariant under the gauge transformations with values in $G$ at the sites of $\mathsf B_\omega$. Then $|\Psi(V')-\Psi(1)|\le2ST^{-2}$ for every field $V'$ on $\mathsf B_\omega$ that satisfies \textup{(u1)} to \textup{(u3)} at $x'$.\end{lemma}

\begin{proof}[Proof of \Cref{Y-H1-LT57}] We write $|\cdot|$ for the operator norm. In $\operatorname{SL}(2,\mathbb C)$ one has $M^{-1}=(\operatorname{tr}M)1-M$; hence $|M^{-1}|=|M|\ge1$ and $|M^{-1}-1|=|M-1|$. For $k$ on the sites of $\mathsf B_\omega$ with values in $\operatorname{SL}(2,\mathbb C)$ we set $V'^k(c):=k(c_-)V'(c)k(c_+)^{-1}$. The domain $U^{17T}(\mathsf B_\omega;x')$ is open, and it is invariant under $G$-valued $k$.

\textup{(i)} \emph{A tree gauge.} Let $\Gamma(o,y')$ be the path in $\mathsf B_\omega$ from $o$ to $y'$ that follows $e_1$ first, then $e_2$, and finally $e_3$; it has at most $3\omega\le\ell_w(x')$ bonds. We put $\mathsf k(y'):=V'(\Gamma(o,y'))$; by \textup{(u2)}, $|\mathsf k|,|\mathsf k^{-1}|<2$. The field $\mathsf W:=V'^{\mathsf k}$ equals $1$ on the bonds of the paths $\Gamma$, and $\partial \mathsf W(P)=\mathsf k\,\partial V'(P)\mathsf k^{-1}$ with $\mathsf k$ at the base site of $P$; so \textup{(u3)} gives $|\partial \mathsf W(P)^{\pm1}-1|<4\cdot2\varepsilon'=8\varepsilon'$.

Each bond along $e_3$ lies on a path $\Gamma$. Consider a bond $c=(y';e_2)$. The bond $(y'_1,y'_2,0;e_2)$ lies on a path $\Gamma$, and we pass from it to $c$ across $|y'_3|$ plaquettes forming a ladder in the plane $(2,3)$; the sides of this ladder along $e_3$ lie on paths $\Gamma$, and each plaquette contributes one factor $\partial \mathsf W(P)^{\pm1}$ to $\mathsf W$. For a bond $(y';e_1)$ we start from $(y'_1,0,0;e_1)$ and cross $|y'_2|$ plaquettes in the plane $(1,2)$ at height $0$, then $|y'_3|$ plaquettes in the plane $(1,3)$. So $\mathsf W(c)$ is a product of at most $2\omega$ factors $1+a_i$ with $|a_i|<8\varepsilon'$, and $|\mathsf W(c)-1|\le e^{16\omega\varepsilon'}-1\le18\omega\varepsilon'\le1/10$ by \textup{(C0)}. Hence $A_c:=-i\log \mathsf W(c)$ is traceless, and $|A_c|\le1.054\cdot18\omega\varepsilon'\le19\omega\varepsilon'$.

\textup{(ii)} \emph{The value at $\mathsf W$.} We show that $\Psi(V')=\Psi(\mathsf W)$. Let $\mathsf k=hu$ be the polar decomposition, where $h$ is positive of determinant $1$ and $u$ has its values in $G$; then $|h|=|h^{-1}|=|\mathsf k|<2$. Unitary factors leave the norms unchanged, so $V'':=V'^u$ again has \textup{(u1)} to \textup{(u3)}; moreover $\Psi(V'')=\Psi(V')$ by the invariance, and $\mathsf W=V''^h$. For complex $s$ let $V_s:=V''^{h^s}$, where $h^s=e^{s\log h}$. Then $V_s$ is entire in $s$, $V_0=V''$, $V_1=\mathsf W$, and $V_{s+i\sigma}=(V_s)^{h^{i\sigma}}$, with $h^{i\sigma}$ taking its values in $G$. We claim that $V_s$ lies in the domain for $0\le s\le1$. If so, $V_s$ lies in the domain for all $s$ in an open set containing the strip $0\le\operatorname{Re}s\le1$. On the strip itself this holds because $V_{s+i\sigma}=(V_s)^{h^{i\sigma}}$; the set $\{s:V_s\text{ in the domain}\}$ is open, and we take its component containing the strip, which is a connected open set. On this set $s\mapsto\Psi(V_s)$ is holomorphic, and by the invariance it is constant on the imaginary axis; so it is constant, and $\Psi(V')=\Psi(V'')=\Psi(\mathsf W)$.

\emph{Proof of the claim.} For \textup{(e3)}: $|\partial V_s(P)-1|\le|h^s||h^{-s}||\partial V''(P)-1|<8\varepsilon'<34T\varepsilon'$. For \textup{(e1)} and \textup{(e2)}, let $\pi$ be a path in $\mathsf B_\omega$ from $y'$ to $y''$ with no more than $\ell_w(x')$ bonds, a single bond in the case of \textup{(e1)}. From \textup{(i)}, from \textup{(C1)} with $T+1>2$, and from $\mathsf g_x\le1$ we get $|\mathsf W(\pi)^{\pm1}-1|\le e^{19\omega\varepsilon'\ell_w(x')}-1\le \mathsf g_x/250$; here $\mathsf g_x\le1$ because $\mathsf g_x<\vartheta_0^{\flat}/\log x'$, by Step~8, \textup{(8.2)}, and $\vartheta_0^{\flat}\le10^{-4}$, by item~C.II~1.3 of \S\ref{Y-C-1}.

Next, $V''(\pi)=h(y')^{-1}\mathsf W(\pi)h(y'')$; so, with $\theta:=1-s\in[0,1]$, \[V_s(\pi)=h(y')^{-\theta}h(y'')^{\theta}+h(y')^{-\theta}(\mathsf W(\pi)-1)h(y'')^{\theta},\] and the second term has norm at most $4\mathsf g_x/250$. For positive matrices $H_1,H_2$ and $0\le\theta\le1$, $|H_1^\theta H_2^\theta|\le|H_1H_2|^\theta$; this bound is Cordes' inequality, which is a form of the L\"owner--Heinz theorem, and it is proved in Step~(1) of the proof of \Cref{Y-C-L69}. We take $H_1=h(y')^{-1}$ and $H_2=h(y'')$. Then $H_1H_2=V''(\pi)[1+h(y'')^{-1}(\mathsf W(\pi)^{-1}-1)h(y'')]$, so $|H_1H_2|\le|V''(\pi)|(1+4\mathsf g_x/250)$; as the determinant equals $1$, $|H_1H_2|\ge1$ and $|H_1H_2|^\theta\le|H_1H_2|$. Using $|V''(\pi)|<2$, we find $|V_s(\pi)|\le|V''(\pi)|+12\mathsf g_x/250<|V''(\pi)|+\mathsf g_x$, and this is $<\mathsf b(x')+\mathsf g_x$ by \textup{(u2)} in the case of \textup{(e2)}, where $\pi$ has at most $\ell_w(x')$ bonds, or $<1+\vartheta(x')+\mathsf g_x$ by \textup{(u1)} in the case of \textup{(e1)}, where $\pi$ is a single bond.

\textup{(iii)} \emph{A family of fields.} The family $\mathsf W_t:=(e^{itA_c})_c$, $t$ complex, $|t|<T$, lies in the domain. For \textup{(e1)} and \textup{(e2)}: $|\mathsf W_t(\pi)|\le e^{19T\omega\varepsilon'\ell_w(x')}\le e^{\mathsf g_x/138}<1+\mathsf g_x$, by \textup{(C1)}. For \textup{(e3)}: $\partial \mathsf W_t(P)$ is a product of four exponentials $e^{\pm itA_c}$, with $\sigma_t:=\sum|tA_c|\le76T\omega\varepsilon'\le1/2$ by \textup{(C0)}; therefore $|\partial \mathsf W_t(P)-1-it(\mathrm{d}A)_P|\le e^{\sigma_t}-1-\sigma_t\le0.6\sigma_t^2$, where $(\mathrm{d}A)_P$ denotes the signed sum of the four $A_c$. Taking $t=1$ we get $|(\mathrm{d}A)_P|\le|\partial \mathsf W(P)-1|+0.6\cdot76^2\omega^2\varepsilon'^2\le8\varepsilon'+3466\omega^2\varepsilon'^2$. Hence $|\partial \mathsf W_t(P)-1|\le T\varepsilon'[8+3466\omega^2\varepsilon'+3466T\omega^2\varepsilon']\le8.1T\varepsilon'<34T\varepsilon'$, by \textup{(C2)}.

\textup{(iv)} \emph{Schwarz' lemma.} The function $\psi(t):=\Psi(\mathsf W_t)$ is holomorphic on $|t|<T$, with $|\psi|\le S$, $\psi(0)=\Psi(1)$ and $\psi(1)=\Psi(\mathsf W)=\Psi(V')$. Its derivative at the origin is $\psi'(0)=\mathcal L(A)$, where $\mathcal L$, the derivative of $\Psi$ at the field $1$, is complex linear in $A$. For a constant $u\in G$ we have $(\mathsf W_t)^u=(e^{it\operatorname{Ad}_uA_c})_c$, and so $\mathcal L(\operatorname{Ad}_uA)=\mathcal L(A)$. The Haar mean of $\operatorname{Ad}_u$ on $\mathfrak g_{\mathbb C}$ is $0$, since the adjoint representation of $\operatorname{SU}(2)$ is irreducible and not trivial; hence $\mathcal L(A)=0$. Thus $\psi-\psi(0)$ is bounded by $2S$ and vanishes to second order at $0$, and Schwarz' lemma gives $|\psi(1)-\psi(0)|\le2ST^{-2}$.

\end{proof}

\emph{Step S10, continued.} \textup{(4)} There are $L^3$ first cubes, and $\sum_{X\ni\text{cube}}e^{-\tilde\kappa(\#X-1)}\le1.2124$; since $L^3T^{-2}=361\mathsf N_{\rm c}^2$, the sum over the small $X$ is at most $2.4248\cdot361\mathsf N_{\rm c}^2e^{Ac_a(x)}Bc_a(x)$. Using $e^{Ac_a(x)}\le e^{1/2}$ and $c_a(x)\le c_a(x')/(0.594L^{1/2})$, together with $2.4248\cdot361\cdot1.6488/0.594\le2429.8\le2430$, we bound it by $2430\mathsf N_{\rm c}^2L^{-1/2}Bc_a(x')$.

Now let $X$ be not small. On $\mathcal U_{Y_0(X)}(x')$, Step~2\,\textup{(ii)}, applied with $\tau=1$, gives $|G^\delta_X[u](V')-G^\delta_X[u](1)|\le2S_X\le3.3Bc_a(x)e^{-\tilde\kappa(n-1)}$. A cell $\square\in Y_0(X)$ meets $X^{+1}$, so the sum over these $X$ is at most \[3.3Bc_a(x)(L+2)^3\sum_{n>L/2-2}\exp\{(4.2581-\tilde\kappa)(n-1)+\tilde\kappa(7+432n/L)\},\] and the counting in the proof of \textup{($\Sigma$a)} bounds this by $\exp\{-0.867L+70+3\log(L+2)\}Bc_a(x)\le\frac18e^{-L/2}Bc_a(x)$.

\emph{Step S11: summing the parts, \textup{(h3)} and \textup{(b\#)}.} For \textup{(h3)} we arrange the cell sum of the one-letter part according to the sets $Y$: a set $Y=\{\square\}$ receives, with weight $1$, the small $X$ whose first cube lies in $\hat\square$, and an $X$ that is not small, with $Y_0(X)=Y\ni\square$, enters with the weight $e^{\tilde\kappa(\#Y_0-1)}$. Adding \textup{(h2)} and \textup{(h1)}, we obtain the bound $[2430\mathsf N_{\rm c}^2L^{-1/2}+366L^{5/2}x^{-1/4}]Bc_a(x')+e^{-L/2}Bc_a(x)$.

By \textup{(Y.10)}, \[\FitDisplay{\dimexpr\linewidth-1pt\relax}{366L^{5/2}x^{-1/4}\le(366/1800)L^{-1}\le0.2034L^{-1}}\]. Also \[\FitDisplay{\dimexpr\linewidth-1pt\relax}{e^{-L/2}c_a(x)\le e^{-L/2}c_a(x')/(0.594L^{1/2})\le0.1L^{-1}c_a(x')}\]. So the bracket is at most $2430\mathsf N_{\rm c}^2L^{-1/2}+L^{-1}$. In the setting \textup{(L2$^{\rm su}$)}, $2430/(1.2\cdot10^4)=0.2025$, and $L^{-1}<10^{-13}$. In the setting \textup{(L3su)}, $2430\mathsf N_{\rm c}^2L^{-1/2}\le0.243L^{-1/4}$, and \[\FitDisplay{\dimexpr\linewidth-1pt\relax}{L^{-1}=L^{-3/4}L^{-1/4}\le1.6\cdot10^{-10}L^{-1/4}}\]; thus the bracket is at most $0.2431L^{-1/4}=\vartheta_a$.

For \textup{(b\#)} let $u_\tau:=(\tau E^a,E^h,w)$, $0\le\tau\le1$; this family lies in $M_A(x;p)$, and $E^a\in T^a(x)$ with $B_p(E^a)\le A$. By Step~S8, $E'^{\rm old}[u_0]=0$. For $\tau<1$ and $|z|\le16/A-\tau$ we have $u_\tau+z(E^a,0,0)=u_{\tau+z}\in M_{16}(x;p)$. As $A_a(u_\tau)\le\tau A<16$, \textup{(3)} of Step~S9 shows that $\tau\mapsto E'^{\rm old}[u_\tau](Y,V')$ is differentiable on $[0,1)$ with the derivative $DA(u_\tau)[E^a](Y,V')$. Continuity on $[0,1]$ holds as well: the series consists of functions continuous in $\tau$ and is dominated by the majorant $\hat\Phi$ with the sizes $e_\alpha$ for this $A$, valid for the members within $A$ caps, that is, on $M_A(x;p)$, which contains every $u_\tau$ with $0\le\tau\le1$. Hence $E'^{\rm old}[u](Y,V')=\int_0^1DA(u_\tau)[E^a](Y,V')\,d\tau$. We take the supremum over $V'$ and the cell sum under the integral, and we apply \textup{(h3)} at $u_\tau\in M_A(x;p)$, with $B\le A$.

\emph{Step S12: the bounds at $\tau^{\sharp}$, \Cref{Y-H1-Rcl} and \textup{(h\#)}.}

\textup{(1)} \emph{The condition \textup{(S$_\tau$)} for $\{$\textup{(a)}, \textup{(v)}, \textup{(t)}$\}$ at $(\tau^{\sharp},\varepsilon_2^{\sharp})$, with $\mathsf r=r_1$.} We follow the lines of Step~7\,\textup{(ii)}. If ${\rm dist}(I_\beta,\square)\le r_1$, then some cell $\square''$ of $J_\beta$ is within $\varrho_*+r_1\le s'$ of $\square$, and there are $(2s'+1)^3$ such cells. For the sum over the letters $\beta$ with $J_\beta\ni\square''$ we use $\sum_{\beta:J_\beta\ni\square''}\le\sum_{\square_0}e^{-{\rm dist}(\square_0,\square'')}\sum_{\beta\text{ rooted at }\square_0}e^{\#J_\beta}$, together with $\mathsf g(1)\le50$. Further, $e_\beta e^{3\varepsilon_2^{\sharp}\#I_\beta}=e^{\circ}_\beta e^{(2x^{-1/8}+3\varepsilon_2^{\sharp})\#I_\beta}\le e^{\circ}_\beta e^{\#J_\beta}$, because $\#I_\beta\le(2(s_0+\varrho_*)+1)^3\#J_\beta$ gives $3\varepsilon_2^{\sharp}\#I_\beta\le\#J_\beta/2$ and, by \textup{(Y.2)}, $2x^{-1/8}\#I_\beta\le\#J_\beta/25$. Also $e^{d_\beta}=e^{(\tilde\kappa+1)s'}e^{(\tilde\kappa+1)\#J_\beta}$.

The rooted sums of $e^{\circ}_\beta e^{(\tilde\kappa+3)\#J_\beta}$ are at most $5e^{\tilde\kappa+3}L^3\cdot16c_a(x)$, $2e^{\tilde\kappa+3}K_v^{\flat}p_0^3x^{-1/2}$ and $x^{-20}$, for the letters \textup{(a)}, \textup{(v)} and \textup{(t)} respectively. Therefore $\sum_{\beta:{\rm dist}(I_\beta,\square)\le r_1}e_\beta e^{3\varepsilon_2^{\sharp}\#I_\beta+d_\beta}\le a^{\sharp}(x)=\tau^{\sharp}\varepsilon_2^{\sharp}$, and \textup{(S$_\tau$)} holds for every $\tau\ge\tau^{\sharp}$.

\textup{(2)} \emph{The bound $\tau^{\sharp}\le1$.} We take \textup{(Y.3)} with $\kappa_1$ in the place of $\kappa$. Its bracket and that of $a^{\sharp}$ are multiplied by the same factor $50(2s'+1)^3e^{(\kappa_1+1)s'}$. Each of the three terms of the bracket of $a^{\sharp}$ is also a term of the bracket of \textup{(Y.3)}, except that the first has $16c_a(x)$ where \textup{(Y.3)} has $16c_a(x)$ plus a term $\ge0$; the remaining terms of the bracket of \textup{(Y.3)} are $\ge0$. Therefore the bracket of $a^{\sharp}$ does not exceed that of \textup{(Y.3)}, so $x^{1/4}a^{\sharp}\le x^{-1/8}$. By \textup{(Y.2)}, $1/\varepsilon_2^{\sharp}=6(2(s_0+\varrho_*)+1)^3\le0.12x^{1/8}$. Consequently $\tau^{\sharp}\le0.12x^{-1/4}$.

\textup{(3)} \emph{\Cref{Y-H1-Rcl}.} By \textup{(1)}, the alphabet $\{$\textup{(a)}, \textup{(v)}, \textup{(t)}$\}$, with $\mathsf r=r_1$, the sizes $e_\alpha$ at $(\tilde\kappa,A)$ and $d_\alpha=(\tilde\kappa+1)(\#J_\alpha+s')$, satisfies the hypothesis \textup{(S$_\tau$)} of \Cref{Y-H1-LK} at $\epsilon=\varepsilon_2^{\sharp}$ for every $\tau\ge\tau^{\sharp}$. Let $\tau$ with $\tau^{\sharp}\le\tau\le1$ be as in \Cref{Y-H1-Rcl}. \Cref{Y-H1-LK}, applied with this $\tau$ and with $\epsilon=\varepsilon_2^{\sharp}$, and multiplied by $\tau$, gives $\sum_C\mathrm{mult}_\alpha\hat\Phi e^{d(C)}\tau^{-(n(C)-1)}\le e_\alpha e^{2\varepsilon_2^{\sharp}\#I_\alpha+d_\alpha}$. We conclude as in the proof of \textup{(R*)}, with \textup{(Clu)} and with $\#\operatorname{co}(C)\ge{\rm dist}(\square_0,\square)+1$.

\textup{(4)} \emph{The power.} Since $s'=8r'\le160c_R\log x/M+8$, we have $e^{(\kappa_1+1)s'}\le e^{8(\kappa_1+1)}x^{\sigma_M}$. The third inequality of \textup{(Y.6)} gives $x^{-10}\le c_a(x')/32\le L^{1/2}c_a(x)/32$, and \textup{(Y.10)} gives $x^{-10}L^{1/2}\le1$; hence $x^{-20}\le c_a(x)$. So the bracket of $a^{\sharp}$ is at most $e^{\kappa_1+3}[81L^3E_a+2K_v^{\flat}]p_0^3x^{-1/2}$; moreover $6\cdot50=300$ and $8(\kappa_1+1)+\kappa_1+3=9\kappa_1+11$. Finally $\sigma_M\le160(\kappa_1+1)/(1920(\kappa_1+4))\le1/12$.

\textup{(5)} \emph{The bounds \textup{(h1)} and \textup{(g3)} at $\tau^{\sharp}$.} In \textup{(1)} of Step~S9 we use \Cref{Y-H1-Rcl}, applied with $\tau=\tau^{\sharp}$, instead of \textup{(R*)}; the part with $n(C)\ge2$ then has $\tau^{\sharp}$ in place of $t$, and $e^{2\varepsilon_2^{\sharp}\#I}$ in place of $e^{2\varepsilon_2\#I}$. For small $X$, $(2\varepsilon_2^{\sharp}+2x^{-1/8})\#I\le1/3+1/25$, and $49.9791\cdot1.2124\cdot e^{0.3734}\cdot e^{1/2}\le145.13\le146$; twice this gives $290.26\le291$. For $X$ not small, the factor $e^{0.3734\#Y_0}$ replaces $e^{0.08\#Y_0}$; we have $1.3734\le0.23\tilde\kappa$, the exponents are at most $-0.289\tilde\kappa(n-1)+10\tilde\kappa$, and the counting in the proof of \textup{($\Sigma$a)} remains valid. Indeed, the first exponent is at most $-0.289\tilde\kappa(L/2-37.7)\le-0.867L+65.4$. Consecutive terms of the sum over $n$ have ratio at most $e^{-0.289\tilde\kappa}\le e^{-1.734}$, as $\tilde\kappa\ge6$, so this sum is at most its first term times $1/(1-e^{-1.734})\le1.215$, the sum of the geometric series with ratio $e^{-1.734}$. Since $\log49.98+\log1.215\le4.11$, the constant term becomes at most $69.6\le70$.

\emph{Step S13: additional statements used from this section on.} The present section and the sections after it use the statements of \S\ref{Y-H-1}, the proof of \Cref{Y-H-C9} included, and the statements of Step~8 of the proof of \Cref{Y-G-ThA}. They also use the following statements.

Instead of the floor $x_9$ of \textup{(O13)} of item~1 of \S\ref{Y-H-1}, we use in this section and in the later ones the floor $x_9^U$, where $x_9^U$ is defined in \textup{(O13)} of item~0 of this section, with \textup{(Y.10)} among the conditions that define it, and $x_9^U\ge x_9$. Item~\textup{(VII)} of the alphabet, stated and proved in Step~S3, is used besides the items of Step~3. Together with the product bound of Step~6 we use the bound \textup{(M)}; Step~S7 proves it. In addition to item~\textup{(R)} of Step~8, we apply the claim \textup{(R*)}, with $\mathrm{mult}_\alpha$, proved in Step~S7. Step~S10 carries out the argument of \textup{(8.2)} for $G^\delta_X[u]$, and this argument is used alongside \textup{(8.2)} of Step~8. Step~11 concerns affine families on an open set of $\mathbb C$, and together with it we apply clause \textup{(c+)} of \Cref{Y-H1-U1}, which is proved in Step~S5 and yields the conclusion of Step~11 for families that are holomorphic on an open set of $\mathbb C^n$.

\Cref{Y-G-L18} is part of Step~8 of the proof of \Cref{Y-G-ThA}, and besides its claims \textup{(1)} and \textup{(2)} we use the following two claims:

\begin{enumerate}\item[\textup{(1*)}] if \textup{(S)} of \Cref{Y-G-L18} holds with $\tau$ in place of $t$, $0<\tau\le1$, then \[\FitDisplay{\dimexpr\linewidth-1pt\relax}{\sum_C\mathrm{mult}_\alpha(C)\hat\Phi(C)e^{d(C)}\tau^{-(n(C)-1)}\le e_\alpha e^{2\varepsilon_2\#I_\alpha+d_\alpha}}\], where $\mathrm{mult}_\alpha(C)=\#\{i:\alpha\in\mathfrak B_i\}$;

\item[\textup{(2*)}] consequently the claims \textup{(1)} and \textup{(2)} of \Cref{Y-G-L18} hold with the factor $\mathrm{mult}_\alpha(C)$ under the sum.\end{enumerate}

The claims \textup{(1*)} and \textup{(2*)} are \Cref{Y-H1-LK}. That lemma is proved in Step~S6, and \textup{(1*)} is its first claim, taken with $\epsilon=\varepsilon_2$ and multiplied by $\tau$. In Step~8, \textup{(8.2)}, \Cref{Y-H-Th57} is used. In the place of \Cref{Y-H-Th57}, \Cref{Y-H1-LT57}, which gives the conclusion of that theorem for every $T>1$ for which \textup{(C0)} to \textup{(C2)} hold, is applied in \textup{(3)} of Step~S10 and, with another value of $T$, in Step~4 of the proof of \Cref{Y-H3-U3}; the lemma is stated and proved in Step~S10. Besides \textup{(O4)} of item~1 of \S\ref{Y-H-1}, $L$ also satisfies the condition \textup{(L3su)}, which is stated in \textup{(O4a)} of item~0 of this section.\end{proof}

\subsection{Linear flow of quadratic sources}\label{Y-H2-1}

All objects of this section are scalar kernels on tori: no field, no group, no member and no cap enters them, apart from the letters of \textup{(S5)(iii)}, and no source appears in any kernel, law, centre or row. We show that one step $D\mapsto T_\ell D$ on families of kernels contracts by the factor $\vartheta_s=c_sL^{-15/32}$ away from the charges, that is, away from the values on constant forms; that the step averages the charges exactly, whatever the defect $\tau(\ell)$ of the rows may be; that a cut moves the charges by a power of $x$; and that the steps at two sets of rows differ by $c_I\delta^{15/16}$, with charges $0$. From these facts we obtain \Cref{Y-H2-U2}, which is clause~\textup{(S1)} of \Cref{Y-H2-U5}, and clause~\textup{(S3)}. This section uses \S\ref{Y-A-1}; from \S\ref{Y-J-1} it uses \textup{(O0)} in item~1, the definition ``Runs'' in item~2, \textup{(RU)} in item~3, Steps~0\,(i) and~1 of the proof of \Cref{Y-J-ThA}, and \textup{(X12)}, \textup{(X13)}; from \S\ref{Y-I-1} it uses \Cref{Y-I-LB} and \Cref{Y-I-CorD}.

\emph{Item 1: the setting.}

\begin{remark} The letters of this section are those listed in item~\textup{(0)}, and its numbers those listed in item~\textup{(1)}, of \S\ref{Y-H0-1}.

A ``Step~$n$'' cited without the name of a statement is Step~$n$ of the proof of Thms.~\ref{Y-H2-U2} and~\ref{Y-H2-U5}.

The letter $K$ is the number of steps, except in $K=K_*$ (Step~1\,(a)), where it is the bound of the rows of \S\ref{Y-A-1}.\end{remark}

\emph{Letters and data.} The letters and data come from the table, item~3 of \S\ref{Y-G0-1}, in the form given in item~1 of \S\ref{Y-J-1}, which also contains \textup{(O14)}. Its entries read as follows.

\begin{enumerate}\item[\textup{(O0)}] The numbers $b_*,\gamma_*,\lambda_*,K_*,L_A,C_1,c_{\rm fix}$, together with $L_B=\max\{L_A,c_{\rm fix}^{10}\}$ and $\lambda_\flat=\lambda_*/16$; moreover $\mathsf g(\lambda)=\sum_{n\in\mathbb Z^3}e^{-\lambda|n|_\infty}$, and $\mathsf g_1(\lambda)$ is the same sum with the factor $|n|_\infty$ in each term.

\item[\textup{(O4)}] $L$ is odd, $L\ge L_*^{\rm su}\ge\max\{L_A,L_B\}$ and $a\ge3/(c_LL)$; moreover $\vartheta_1=c_{\rm fix}L^{-1/2}$.

\item[\textup{(O6)}] The number $c_R$, with $e_R=\lambda_*c_R/2\ge256$ (Step~0\,(i) of the proof of \Cref{Y-J-ThA}), and $\varrho(x)=\lceil c_R\log x\rceil$.

\item[\textup{(O13)}] Of the conditions on the last floor $x_s$, this section rests on \textup{(X12)} and \textup{(X13)}, so that $\mathsf b\ge x_s$ meets the requirement $\mathsf b\ge x_\flat$ of \S\ref{Y-A-1} and of \textup{(RU)}, and on \textup{(X5)} below; in \textup{(S5)(iii)} alone, through \Cref{Y-I-CorD}, it rests also on \textup{(X14)} $p_0(x)\ge1$.\end{enumerate}

The data are: $\mathsf b\ge x_s$; $N$ even; $K'\ge K\ge1$; $m=K'-K$; $x_k=L^{K-k}\mathsf b$; $a_K=L^{-K}$; and $\mathbb T_k=\mathbb Z^3/NL^{K-k}\mathbb Z^3$ for $0\le k\le K$, each in the units of its own level. In step $k$ we have $\mathbb T=\mathbb T_k$, and $\mathbb T'=L\mathbb Z^3/NL^{K-k}\mathbb Z^3$ is $\mathbb T_{k+1}$.

\begin{definition}[Rows, partners and the kernels of the step] \emph{Rows.} We write $\ell_k=\ell_k^{[K]}$ and $\tilde\ell_j=\ell_j^{[K']}$ for the scalar chains of the runs $K$ and $K'$, as given by the definition ``Runs'' in item~2 of \S\ref{Y-J-1}: $\ell_0=\emptyset$, and $\ell_{k+1}=((\ell_k)')^{\rm c}$ is obtained by cutting the rows of \Cref{Y-A-P5} at $\varrho(x_{k+1})$. This recursion is used for every even $N$, and the condition $M\mid N$ imposed in that definition plays no part in these chains.

\emph{Site forms.} $\mathsf E_{k,y}(p,p')=\langle\ell_{k,y}(p),\ell_{k,y}(p')\rangle$, so that $\mathsf B_k=1+\sum_y\mathsf E_{k,y}$; the forms $\tilde{\mathsf E}_j$ are given by the same formula with the rows $\tilde\ell_j$ in place of the rows $\ell_k$.

\emph{Partners.} Level $m+k$ of the run $K'$ lives on $\mathbb T_k$, and $\tilde x_{m+k}=x_k$; partners are therefore cut at the same radius.

\emph{Kernels of the step.} We let $H_k:=H^\sharp$ be the map of \Cref{Y-A-P4}\,(a) at the rows $\ell_k$, from $\Omega^2(\mathbb T_{k+1})$ to $\Omega^2(\mathbb T_k)$; $\tilde H_j$ is the same map at the rows $\tilde\ell_j$. We write $\delta_P$ and $\delta_p$ for the $2$-forms on the coarse level and on the fine level that are $1$ at $P$ and at $p$, respectively, and $0$ at every other plaquette. The column $P$ of $H^\sharp$ is $\nu^P:=H^\sharp\delta_P$, which is denoted $\nu$ in \S\ref{Y-A-1}.

\emph{Blocks.} $B(y')$ is the centred block. Further, $B^0(y)=\{y\}$ and $B^{k+1}(y')=\bigcup_{y\in B(y')}B^k(y)$; thus $B^k(y)$ consists of the sites $z$ of $\mathbb T_0$ with $|z-L^ky|_\infty\le(L^k-1)/2$.\end{definition}

\begin{definition}[Places] Throughout, $f$ is a complex bounded function on $T_N=\mathbb R^3/N\mathbb Z^3$. The site $y$ of $\mathbb T_k$ is placed at $a_KL^ky$; in the run $K'$, the site $y$ of $\tilde{\mathbb T}_j$ is placed at $a_{K'}L^jy$, so partners have the same places. For a plaquette $p$, $\iota(p)\in\{1,2,3\}$ is the direction normal to $p$, and $(z;\iota)$ is the plaquette based at $z$ with normal $e_\iota$, spanned by $e_\mu,e_\nu$. Its barycentre is $x(p)=x^{[K]}(z;\iota)=a_K\bigl(z+\frac12(e_\mu+e_\nu)\bigr)$ (entry~13 of item~\textup{(0)} of \S\ref{Y-H0-1}). We denote by $h_\iota$ the constant $2$-form whose component $\iota$ equals $1$ and whose other components are $0$. The Lipschitz constant ${\rm Lip}$ is taken with respect to the Euclidean distance, and $\|f\|_\sharp=\max\{\|f\|_\infty,40M\,{\rm Lip}f\}$.\end{definition}

\begin{definition}[Families, norm, charge, step and cut] A family on a level $\mathbb T$ is $D=(D_y)_{y\in\mathbb T}$, where every $D_y$ is a symmetric complex scalar kernel on pairs of plaquettes of $\mathbb T$, and we write $D_y(F,G)=\sum_{p,p'}D_y(p,p')F_pG_{p'}$. The norm of a family is \[\|D\|=\|D\|_{(\lambda_\flat)}:=\sup_{y,p,p'}|D_y(p,p')|e^{\lambda_\flat(|y-p|+|y-p'|)},\] with $|y-p|$ and $|y-p'|$ measured in the units of the level $\mathbb T$ (\Cref{Y-A-OwnUnits}). The charge of $D$ at $y$ is the $3\times3$ matrix $c_y(D)(\iota,\iota'):=D_y(h_\iota,h_{\iota'})$, and $|c(D)|$ is the supremum over $y$ of its operator norm. The step at the rows $\ell$ is $(T_\ell D)_{y'}:=L\sum_{y\in B(y')}H^{\mathsf T}D_yH$, with $H=H^\sharp$ at $\ell$; entry by entry, $(T_\ell D)_{y'}(P,P')=L\sum_{y\in B(y')}D_y(\nu^P,\nu^{P'})$. The cut at radius $\varrho$ of a family $D'$ on the coarse level is \[({\rm cut}_\varrho D')_{y'}(P,P'):=D'_{y'}(P,P')1\{|y'-P|\le\varrho,\ |y'-P'|\le\varrho\}.\]\end{definition}

\begin{definition}[Source forms and Riemann means] The source forms are given by \[\mathbb D_{0,y}[f]:=\sum_{p\text{ based at }y}f(x(p))\,\delta_p\delta_p^{\mathsf T}\] and, recursively, by $\mathbb D_{k+1}[f]:={\rm cut}_{\varrho(x_{k+1})}(T_{\ell_k}\mathbb D_k[f])$. The forms $\tilde{\mathbb D}_j[f]$ of the run $K'$ are given by the same two formulas with $x^{[K']}$, $\tilde\ell_j$ and $\varrho(\tilde x_{j+1})$ in place of $x^{[K]}$, $\ell_k$ and $\varrho(x_{k+1})$. The Riemann means are ${\rm Mn}^{[K]}_{k,y,\iota}(f):=L^{-3k}\sum_{z\in B^k(y)}f(x^{[K]}(z;\iota))$, and we set \[{\rm os}_i:=\sup_{y\in\mathbb T_i,\iota}|{\rm Mn}^{[K']}_{m+i,y,\iota}(f)-{\rm Mn}^{[K]}_{i,y,\iota}(f)|,\] together with ${\rm osc}_K(f):={\rm os}_0$.\end{definition}

\emph{Item 2: the numbers and their places.}

\begin{definition}[Entries of the table]

\begin{enumerate}\item[\textup{(O0a)}] Pure numbers: $\mathsf g_\flat:=\mathsf g(\lambda_\flat)$, $\mathsf g':=\mathsf g(\lambda_\flat/2)$ and $\mathsf g_1':=\mathsf g_1(\lambda_\flat/2)$. Next, $C_{\rm ch}:=4e^{2\lambda_*}/\gamma_*$ and $K_5:=1.35\,C_{\rm ch}$ $(=0.60\,e^{2\lambda_*}K_*^2\le0.61K_*^2)$. Then $a_G:=(23b_*^2)^{15/16}(7\gamma_*^{-1/2})^{1/16}$, $c_s:=63e^{2\lambda_*}\mathsf g_1'\mathsf g'a_G\gamma_*^{-1/2}$ and $t_5:=18\mathsf g_\flat\mathsf g'$. Further, $a_H:=(4\gamma_*^{-1/2})^{1/16}(3\mathsf g_\flat^2b_*^{1/2})^{15/16}$ and $c_I:=12e^{2\lambda_*}\mathsf g'^2a_H\gamma_*^{-1/2}$. Finally, $A_1:=\max\{1,15C_1^2\}$, $I_5:=c_IK_5A_1$ and $C_5:=0.3(K_5+1)+1.02I_5$.

\item[\textup{(O4a)}] The condition \textup{(L5*)} $c_s\le L^{43/160}$. It depends on \textup{(O0)} alone, and it follows from $L\ge L_B$ (\Cref{Y-H2-Rem2}).

\item[\textup{(O13)}] On $x_s$, the condition \textup{(X5)} $1.01\cdot\frac43t_5C_{\rm ch}Lx_s^{-16}\le\frac1{100}$, which depends on \textup{(O0)} and \textup{(O4)}.

\item[\textup{(O14)}] The numbers $\vartheta_s:=c_sL^{-15/32}$, $\vartheta_5:=\max\{\vartheta_s,\vartheta_1^{15/16}\}$ and $C_5^\#:=L^{15/32}C_5+K_5$, on which no floor depends. Then $L^{-15/32}\le\vartheta_s\le\vartheta_5\le L^{-1/5}\le\frac14$.\end{enumerate}

\emph{Numerals.} The numerals come from those of ``The order of the constants'', \textup{(0)}, \S\ref{Y-A-1}, and from the bounds $\mathsf g(\lambda)\le51\lambda^{-3}$ and $\mathsf g_1(\lambda)\le182\lambda^{-4}$ of \Cref{Y-A-Sums}. We have $\log_{10}\mathsf g_\flat\le46.32$, $\log_{10}\mathsf g'\le47.22$, $\log_{10}\mathsf g_1'\le62.95$, $\log_{10}C_{\rm ch}\le24.646$, $\log_{10}K_5\le24.776$, $\log_{10}a_G\le21.352$, $\log_{10}c_s\le145.34$, $\log_{10}t_5\le94.79$, $\log_{10}a_H\le92.90$, $\log_{10}c_I\le200.44$, $\log_{10}A_1\le157.47$, $\log_{10}I_5\le382.68$ and $\log_{10}C_5\le382.69$. Accordingly, \textup{(L5*)} holds for $\log_{10}L\ge540.8$, and \textup{(X5)} holds once $16\log_{10}x_s\ge121.57+\log_{10}L$. None of the constants below depends on $\mathsf b,N,K,K',k,f$.\end{definition}

\emph{Item 3: the free flow of the source forms.}

\begin{theorem}[Bound, range and locality of the source forms]\label{Y-H2-U2} With the letters and data of item~1, and under the conditions \textup{(L5*)} and \textup{(X5)}, the following hold for all $\mathsf b\ge x_s$, all even $N$, all $K\ge1$, all $0\le k\le K$ and all complex bounded $f$.

\begin{enumerate}\item[\textup{(U2.1)}] $\|\mathbb D_k[f]\|_{(\lambda_\flat)}\le K_5\|f\|_\infty$; site by site, this reads $|\mathbb D_{k,y}[f](p,p')|\le K_5\|f\|_\infty e^{-\lambda_\flat(|y-p|+|y-p'|)}$.

\item[\textup{(U2.2)}] $\mathbb D_{k,y}[f]$ is symmetric, linear in $f$, real when $f$ is real, and of range $\le\varrho(x_k)$: it vanishes unless $|y-p|,|y-p'|\le\varrho(x_k)$. At $k=0$ its range is $0$.

\item[\textup{(U2.3)}] $\mathbb D_{k,y}[f]$ depends on $f$ only through the values of $f$ at the barycentres of the plaquettes based in $B^k(y)$; in particular, it is $0$ if $f$ vanishes at all of them.\end{enumerate}\end{theorem}

\emph{Item 4: one step on families of kernels, and two runs.}

\begin{theorem}[One step on families of kernels, and the source forms of two runs]\label{Y-H2-U5} Under \textup{(O0)} to \textup{(O13)}, \textup{(L5*)} and \textup{(X5)}, the following hold for all data of item~1.

\begin{enumerate}\item[\textup{(S0)}] \emph{One step.} Let $\mathbb T=\mathbb Z^3/N'L\mathbb Z^3$ with $N'$ even, let $\mathbb T'$ be its coarse torus, let $D$ be any family on $\mathbb T$, and let $D'$ be any family on $\mathbb T'$.

\begin{enumerate}\item[\textup{(a)}] For every translation covariant $\ell$ on $\mathbb T$, cut or not, and whatever $\tau(\ell)$ is, $H^\sharp h'=L^{-2}h$ and $\mathsf N^\sharp h'=0$ for every constant form $h'$ on $\mathbb T'$, where $h$ is the constant form on $\mathbb T$ with the components of $h'$.

\item[\textup{(b)}] For $\ell\in\mathcal R(b_*,K_*,\lambda_*,\frac19)$, $\|T_\ell D\|\le\vartheta_s\|D\|+C_{\rm ch}|c(D)|$.

\item[\textup{(c)}] For every translation covariant $\ell$, the identity $c_{y'}(T_\ell D)=L^{-3}\sum_{y\in B(y')}c_y(D)$ holds exactly.

\item[\textup{(d)}] $\|{\rm cut}_\varrho D'\|\le\|D'\|$ and $|c_{y'}({\rm cut}_\varrho D')-c_{y'}(D')|\le t_5e^{-\lambda_\flat\varrho/2}\|D'\|$; at $\varrho=\varrho(x)$ with $x\ge1$, the right side is $\le t_5x^{-e_R/16}\|D'\|\le t_5x^{-16}\|D'\|$.

\item[\textup{(e)}] For $\ell,\tilde\ell\in\mathcal R(b_*,K_*,\lambda_*,\frac19)$ on $\mathbb T$, and with $\delta:=\|\tilde{\mathsf E}-\mathsf E\|_{(\lambda_\flat)}$ of any size, $\|(T_{\tilde\ell}-T_\ell)D\|\le c_I\delta^{15/16}\|D\|$, and every charge of $(T_{\tilde\ell}-T_\ell)D$ vanishes.\end{enumerate}

\item[\textup{(S1)}] The assertions of \Cref{Y-H2-U2} hold. Besides, $\mathbb D_{k,y+\mathsf t}[f(\cdot-a_KL^k\mathsf t)](p+\mathsf t,p'+\mathsf t)=\mathbb D_{k,y}[f](p,p')$ for every translation $\mathsf t$ of $\mathbb T_k$.

\item[\textup{(S2)}] \emph{Charges.} For $0\le k\le K$, \[|c_y(\mathbb D_k[f])-{\rm diag}_\iota{\rm Mn}^{[K]}_{k,y,\iota}(f)|\le0.0102\,L^{-1}(x_s/x_k)^{16}\|f\|_\infty.\]

\item[\textup{(S3)}] \emph{Two runs.} For $1\le k\le K$, \[\|\tilde{\mathbb D}_{m+k}[f]-\mathbb D_k[f]\|_{(\lambda_\flat)}\le K_5\,{\rm osc}_K(f)+C_5k\vartheta_5^{k-1}\|f\|_\infty.\] Here ${\rm os}_i$ does not increase with $i$, and ${\rm osc}_K(f)\le\sqrt3a_K{\rm Lip}f\le(\sqrt3/40M)a_K\|f\|_\sharp$; at $k=0$ the left side is $\le(K_5+1)\|f\|_\infty$. \emph{Finer form:} if $\omega_5,P_5\ge0$ and ${\rm os}_i\le\omega_5+P_5\vartheta_5^i\|f\|_\infty$ for all $i<k$, then the right side may be replaced by $K_5\omega_5+(C_5+1.02C_{\rm ch}P_5)k\vartheta_5^{k-1}\|f\|_\infty$.

\item[\textup{(S3c)}] \emph{Corollary.} For $1\le k\le K$ and Lipschitz $f$, $\|\tilde{\mathbb D}_{m+k}[f]-\mathbb D_k[f]\|_{(\lambda_\flat)}\le C_5^\#k\vartheta_5^k\|f\|_\sharp$, with $\vartheta_5\le L^{-1/5}$.

\item[\textup{(S4)}] \emph{The hybrid chain.} For $1\le k_0\le K$, let $\mathbb D^b_{k_0}:=\tilde{\mathbb D}_{m+k_0}[f]$ and $\mathbb D^b_{k+1}:={\rm cut}_{\varrho(x_{k+1})}(T_{\ell_k}\mathbb D^b_k)$. Then $\|\mathbb D^b_k-\mathbb D_k[f]\|\le2.1K_5{\rm osc}_K(f)+1.2C_5k_0\vartheta_5^{k_0-1}\|f\|_\infty$ for $k_0\le k\le K$.

\item[\textup{(S5)}] \emph{Other forms.}

\begin{enumerate}\item[\textup{(i)}] \emph{Totals.} Let $N_k[f]:=\sum_y\mathbb D_{k,y}[f]$. The operator norm of $\tilde N_{m+k}[f]-N_k[f]$ is $\le3\mathsf g_\flat^2$ times the right side of \textup{(S3)}, and its entries are $\le\mathsf g'e^{-(\lambda_\flat/2)|p-p'|}$ times that right side.

\item[\textup{(ii)}] Let $f_1,\dots,f_n$ be Lipschitz with pairwise disjoint sets $\{f_i\ne0\}$, and let $w=\sum z_if_i$ with complex $z_i$. Then $\sum z_i\mathbb D_k[f_i]=\mathbb D_k[w]$, $\|w\|_\infty=\max|z_i|\|f_i\|_\infty$ and ${\rm Lip}w=\max|z_i|{\rm Lip}f_i$; so \textup{(S1)} to \textup{(S4)} hold for $w$ with $\max_i$ and without any factor $n$.

\item[\textup{(iii)}] For $1\le k\le K$ and the family $D:=(\tilde{\mathsf E}_{m+k}-\tilde{\mathbb D}_{m+k}[w]/\mathsf b)-(\mathsf E_k-\mathbb D_k[w]/\mathsf b)$, \[(c_Q^{\rm su}/E_a)x_k^{1/2}\|D\|_{(\lambda_\flat)}/p_0(x_k)\le c_4\mathsf b^{1/2}L^{K/2-9k/10}+{\rm amp}(x_k)\bigl[{\rm osc}_K(w)+(C_5/K_5)k\vartheta_5^{k-1}\|w\|_\infty\bigr],\] where $c_4$ comes from \Cref{Y-I-CorD}, and the left side equals $\delta_\Delta$, the number of \Cref{Y-I-LB}, when that number is evaluated at this $D$.\end{enumerate}

\item[\textup{(S6)}] \emph{Bounded $f$.}

\begin{enumerate}\item[\textup{(i)}] For $K\ge3$ and $m\ge1$, the bound \textup{(S3)} without the term $K_5{\rm osc}_K(f)$ fails for a bounded Borel $f$; the instance is that of Step~10\,(i).

\item[\textup{(ii)}] If $f$ is the indicator of a box, that is, of a product of three arcs, then ${\rm os}_i\le12L^{-i}\le12\vartheta_5^i$; the finer form therefore holds with $\omega_5=0$ and $P_5=12$.

\item[\textup{(iii)}] If $f$ has a modulus of continuity $\omega_f$, then ${\rm osc}_K(f)\le\omega_f(\sqrt3a_K)$.\end{enumerate}\end{enumerate}\end{theorem}

\emph{Item 5: three remarks.}

\begin{remark}[The uncut chain]\label{Y-H2-Rem1} Let $\ell_k$ be the uncut chain, in which all radii are infinite; let $H_k=H^\sharp$ at $\ell_k$ and $\mathbb H_k:=H_0H_1\cdots H_{k-1}$; and let $D^{\rm unc}_k[f]$ be the forms built without any cut. Then \[D^{\rm unc}_{k,y}[f](P,P')=L^k\sum_{p\text{ based in }B^k(y)}f(x(p))\mathbb H_k(p,P)\mathbb H_k(p,P')\] and \[\FitDisplay{\dimexpr\linewidth-1pt\relax}{|D^{\rm unc}_{k,\allowbreak y}[f](P,\allowbreak P')|\le243K_*^2\|f\|_\infty e^{-\lambda_*(|y-P|+|y-P'|)}}\], so that the left side is also $\le250K_*^2\|f\|_\infty e^{-(\lambda_*/4)(|y-P|+|y-P'|)}$. For $f\ge0$ the total obeys $0\le N^{\rm unc}_k[f]\le\|f\|_\infty\mathsf B_k$.\end{remark}

\begin{remark}[\textup{(L5*)} follows from $L\ge L_B$]\label{Y-H2-Rem2} We have $c_s\le10^{145.34}<10^{204.89}\le c_{\rm fix}$; therefore $L\ge L_B\ge c_{\rm fix}^{10}$ gives $L^{43/160}\ge c_{\rm fix}^{43/16}\ge c_s$. Moreover, $L_B\ge c_{\rm fix}^{10}\ge10^{2048.9}$, so this floor on $L$ is larger than $10^{540.8}$, and item~2 shows that \textup{(L5*)} holds for all $L\ge10^{540.8}$.\end{remark}

\begin{remark}[The kernel of the step]\label{Y-H2-Rem3} The step $T_\ell$ and the cut of this section coincide exactly with those of \Cref{Y-H3-U3}, namely $D'_{y'}={\rm cut}_{\varrho(x')}(T_\ell D)_{y'}$ with $H^\sharp$ at the rows $\ell$. Like the rows, the source forms are chosen by hand.\end{remark}

\begin{proof}[Proof of Thms.~\ref{Y-H2-U2} and~\ref{Y-H2-U5} and of Rems.~\ref{Y-H2-Rem1} and~\ref{Y-H2-Rem2}] \emph{Step 1: the columns $\nu^P=H^\sharp\delta_P$; constant forms.} This step proves \textup{(S0)(a)}.

\begin{lemma}[Bounds on the columns of the minimiser, and its value on constants]\label{Y-H2-L51} Let $\ell\in\mathcal R(b_*,K_*,\lambda_*,\frac19)$ be rows on $\mathbb T$, and let $L\ge L_A$. For a set $A$ of sites, let $|F|_A$ be the $\ell^2$ norm of $F$ over the plaquettes based in $A$, and let $|\nabla F|_A^2:=\sum_{\alpha,\iota}\sum_{z\in A}|\partial_\alpha F_\iota(z)|^2$. Let $y'$ be a coarse site, $P$ a coarse plaquette and $v\in\mathbb Z^3$. Then the following hold.

\begin{enumerate}\item[\textup{(a)}] $|\nu^P|_{B(y')+v}\le2\gamma_*^{-1/2}L^{-1/2}e^{\lambda_*(1+|v|/L)}e^{-\lambda_*{\rm dist}(y',P)}$.

\item[\textup{(b)}] $|\nabla\nu^P|\le23b_*^2L^{-1}$.

\item[\textup{(c)}] $|\nabla\nu^P|_{B(y')+v}\le a_GL^{-31/32}e^{\lambda_\flat(1+(|v|+1)/L)}e^{-\lambda_\flat{\rm dist}(y',P)}$.

\item[\textup{(d)}] The assertion of \textup{(S0)(a)} holds, for every translation covariant $\ell$.\end{enumerate}\end{lemma}

\begin{proof}[Proof of \Cref{Y-H2-L51}] We take the four assertions in turn.

\begin{enumerate}\item[\textup{(a)}] At $b=b_*$, $K=K_*$ and $\lambda=\lambda_*$, the rate $\lambda'$ of \Cref{Y-A-L25}\,(c) is $\lambda_*$, since $L\ge L_A$ (proof of \Cref{Y-A-Th29}\,(a)), and $\gamma_{\rm e}=\gamma_*$. \Cref{Y-A-P28}\,(a), taken at $\lambda''=\lambda_*$, thus yields $\sum_pe^{2\lambda_*\mathsf d(p,P)}(\nu^P_p)^2\le4/(\gamma_*L)$. For $p$ based in $B(y')+v$, the triangle inequality for the periodic distance gives $\mathsf d(p,P)\ge{\rm dist}(y',P)-(\mathsf h_L+|v|)/L$, with $\mathsf h_L=(L-1)/2$ from \Cref{Y-A-Lattices}, and the right side is in turn $\ge{\rm dist}(y',P)-\frac12-|v|/L$. Keeping in the sum only the plaquettes based in $B(y')+v$ and inserting this lower bound for $\mathsf d(p,P)$, we obtain \textup{(a)}. The empty rows also belong to the set of the lemma, as they lie in $\mathcal R(1,0,\lambda_*,0)$.

\item[\textup{(b)}] We apply \Cref{Y-A-L27} with $X'=\delta_P$ and $b=b'=b_*$; the choice $b'=b_*$ is allowed, since $\|\mathsf B'\|\le\frac32/\mathsf r_\sharp\le b_*$ (\Cref{Y-A-Th29}\,(a)). Under $b_{\mathsf B}\le b$ and $\|\mathsf B'\|\le b'$, the lemma gives $|\nabla H^\sharp X'|\le23\,b\,b'\,L^{-1}|X'|$, whence the bound $23b_*^2L^{-1}$.

\item[\textup{(c)}] For all $\alpha$, $\iota$ and $z$, $|\partial_\alpha\nu_\iota(z)|^2\le2\nu_\iota(z+e_\alpha)^2+2\nu_\iota(z)^2$. Summing over $\iota$, over $z\in A=B(y')+v$ and over the three directions $\alpha$, and then applying \textup{(a)} to the shifts $v$ and $v+e_\alpha$, each of length at most $|v|+1$, we get $|\nabla\nu|_A^2\le\sum_\alpha(2|\nu|^2_{A+e_\alpha}+2|\nu|^2_A)\le12$ times the square of the bound \textup{(a)} at $|v|+1$. Hence $|\nabla\nu|_A\le\sqrt{48}\gamma_*^{-1/2}L^{-1/2}e^{\lambda_*(1+(|v|+1)/L)}e^{-\lambda_*{\rm dist}}$, where $\sqrt{48}\le6.93\le7$. Let $X$ be the bound \textup{(b)} and $Y$ the bound just found; both bound $|\nabla\nu|_A$, and we use $\min\{X,Y\}\le X^{15/16}Y^{1/16}$. The power of $L$ is then $-\frac{15}{16}-\frac1{32}=-\frac{31}{32}$, the rate is $\lambda_*/16=\lambda_\flat$ and the constant is $a_G$; this is \textup{(c)}.

\item[\textup{(d)}] The pair $(L^{-2}h,0)$ is admitted at $h'$ in the sense of \Cref{Y-A-L26}: indeed $Wh=L^2h'$, because $\sum_pW(P,p)=L^2$ (\Cref{Y-A-Facts}), and $dh=0$. Now let $(F,\mathsf n)$ be any pair admitted at $h'$, and add up the relation $WF+\mathrm{d}'\mathsf n=h'$ over the $N'^3$ plaquettes $P$ of one component $\iota$. Since $\sum_PW(P,p)=1/L$ (\Cref{Y-A-Facts}) and $\sum_P(\mathrm{d}'\mathsf n)_P=0$, each bond variable entering twice and with opposite signs, we find $(1/L)\sum_{p:\iota(p)=\iota}F_p=N'^3h'_\iota$. As there are $(N'L)^3$ such $p$, the mean of $F_\iota$ is $L^{-2}h'_\iota$; that is, $F=L^{-2}h+G$ with $G$ of mean $0$ in each component. The form $\mathsf B$ commutes with all translations of $\mathbb T$ (\Cref{Y-A-CovRows}), so $\mathsf Bh$ is a constant form and $\langle\mathsf Bh,G\rangle=0$. Consequently \[\|(F,\mathsf n)\|_{\mathcal H}^2=L^{-4}\langle h,\mathsf Bh\rangle+\langle G,\mathsf BG\rangle+a|\mathsf n|^2\ge\|(L^{-2}h,0)\|_{\mathcal H}^2.\]

\Cref{Y-A-L26}\,(a) states that $(H^\sharp h',\mathsf N^\sharp h')$ is an admitted pair and that no admitted pair has a smaller norm; the last display, applied to this pair, shows that its norm is at least that of $(L^{-2}h,0)$, and hence the two norms coincide. The last identity of \Cref{Y-A-L26}\,(a), at $(F,\mathsf n)=(L^{-2}h,0)$, then gives $\|(L^{-2}h,0)-(H^\sharp h',\mathsf N^\sharp h')\|_{\mathcal H}^2=0$, and $\mathsf B\ge1$ together with $a>0$ finishes the proof of \textup{(d)}. This argument does not depend on $\tau(\ell)$; in this it sharpens \Cref{Y-A-P28}\,(c).

The statement of clause~\textup{(S0)(a)} is also \Cref{Y-J-L44}\,(a), which appears in Step~1a of the proof of \Cref{Y-J-ThA}, where it has a different proof.\end{enumerate}\end{proof}

\emph{Step 2: one step on a family, away from its charges.} This step proves \textup{(S0)(b)}, \textup{(c)} and \textup{(d)}.

\begin{lemma}[One step contracts a family away from its charges]\label{Y-H2-L52} Let $\ell\in\mathcal R(b_*,K_*,\lambda_*,\frac19)$, and let $D$ be any family on $\mathbb T$. Then \[|(T_\ell D)_{y'}(P,P')|\le(\vartheta_s\|D\|+C_{\rm ch}|c(D)|)e^{-\lambda_\flat(|y'-P|+|y'-P'|)}.\] Moreover, $c_{y'}(T_\ell D)$ is the mean of $c_y(D)$ over $B(y')$. Finally, a cut at radius $\varrho$ never increases the norm, and it changes a charge by at most $18\mathsf g_\flat\mathsf g'e^{-\lambda_\flat\varrho/2}$ times the norm.\end{lemma}

\begin{proof}[Proof of \Cref{Y-H2-L52}] \emph{Freezing at $y$.} We write $\nu=\nu^P$ and $\nu'=\nu^{P'}$, and we let $\nu(y)$ be the constant form whose components are the numbers $\nu_\iota(y)$. By bilinearity, \[D_y(\nu,\nu')=D_y(\nu(y),\nu'(y))+D_y(\nu-\nu(y),\nu')+D_y(\nu(y),\nu'-\nu'(y)),\] and $D_y(\nu(y),\nu'(y))=\sum_{\iota,\iota'}c_y(D)(\iota,\iota')\nu_\iota(y)\nu'_{\iota'}(y)$.

\emph{The charge term.} The first term contributes to $(T_\ell D)_{y'}(P,P')$ an amount of modulus at most \[L\sum_{y\in B(y')}|c_y(D)||\nu(y)|_2|\nu'(y)|_2\le L|c(D)||\nu|_{B(y')}|\nu'|_{B(y')},\] where we use $|c_y(D)|\le|c(D)|$ and then the Cauchy--Schwarz inequality over $y\in B(y')$. By \Cref{Y-H2-L51}\,(a) at $v=0$, the right side is $\le4e^{2\lambda_*}\gamma_*^{-1}|c(D)|e^{-\lambda_*({\rm dist}(y',P)+{\rm dist}(y',P'))}$, and the factor in front of the exponential is $C_{\rm ch}|c(D)|$.

\emph{The cross terms.} Let ${\rm Gr}_y(\nu):=\sum_pe^{-\lambda_\flat|y-p|}|\nu_p-\nu_{\iota(p)}(y)|$ and ${\rm Ab}_y(\nu'):=\sum_pe^{-\lambda_\flat|y-p|}|\nu'_p|$. By the definition of the norm, $|D_y(\nu-\nu(y),\nu')|\le\|D\|{\rm Gr}_y(\nu){\rm Ab}_y(\nu')$, and the sum over $y\in B(y')$ is then estimated by the Cauchy--Schwarz inequality. In the sums over $p$ we write $p=(y+u;\iota)$, where $u\in\mathbb Z^3$ is a representative with $|u|_\infty=|y-p|$.

For ${\rm Gr}$, Minkowski's inequality in $\ell^2(y\in B(y'))$ gives \[(\sum_y{\rm Gr}_y^2)^{1/2}\le\sum_{u,\iota}e^{-\lambda_\flat|u|}(\sum_y|\nu_\iota(y+u)-\nu_\iota(y)|^2)^{1/2}.\] The difference $\nu_\iota(y+u)-\nu_\iota(y)$ is a sum of $|u|_1\le3|u|$ differences $\pm\partial_\alpha\nu_\iota(y+v_s)$ with $|v_s|\le|u|$. So the inner root is $\le3|u|\max_{|v|\le|u|}|\nabla\nu|_{B(y')+v}$, and the sum over $\iota$ multiplies this bound by $3$. \Cref{Y-H2-L51}\,(c), together with $\lambda_\flat(1-1/L)\ge\lambda_\flat/2$ and $1+1/L\le2$, then gives \[(\sum_y{\rm Gr}_y^2)^{1/2}\le9\mathsf g_1'a_Ge^{2\lambda_\flat}L^{-31/32}e^{-\lambda_\flat{\rm dist}(y',P)}.\]

For ${\rm Ab}$, Minkowski's inequality in $\ell^2(y\in B(y'))$, the Cauchy--Schwarz inequality over $\iota$, \Cref{Y-H2-L51}\,(a) and $\lambda_*/L\le\lambda_\flat/2$, which holds since $L\ge32$, give \[\FitDisplay{\dimexpr\linewidth-1pt\relax}{(\sum_y{\rm Ab}_y^2)^{1/2}\le\sum_ue^{-\lambda_\flat|u|}\sum_\iota(\sum_y\nu'_\iota(y+u)^2)^{1/2}\le\sqrt3\sum_ue^{-\lambda_\flat|u|}|\nu'|_{B(y')+u}\le2\sqrt3\mathsf g'\gamma_*^{-1/2}e^{\lambda_*}L^{-1/2}e^{-\lambda_*{\rm dist}(y',P')}.}\]

In the second cross term, the factor ${\rm Gr}_y(\nu')$ obeys the bound found for ${\rm Gr}$ with $P'$ in place of $P$, and the frozen factor is $\sum_pe^{-\lambda_\flat|y-p|}|\nu_{\iota(p)}(y)|\le\mathsf g_\flat\sum_\iota|\nu_\iota(y)|$, whose $\ell^2(y)$ norm is $\le\sqrt3\mathsf g_\flat|\nu|_{B(y')}$; as $\mathsf g_\flat\le\mathsf g'$, the frozen factor obeys the bound found for ${\rm Ab}$, with $P$ in place of $P'$.

\emph{Conclusion.} Apart from their exponential factors, the two cross terms contribute at most \[\FitDisplay{\dimexpr\linewidth-1pt\relax}{L\cdot2\cdot9\mathsf g_1'a_Ge^{2\lambda_\flat}L^{-31/32}\cdot2\sqrt3\mathsf g'\gamma_*^{-1/2}e^{\lambda_*}L^{-1/2}\|D\|=36\sqrt3e^{2\lambda_\flat+\lambda_*}\mathsf g_1'\mathsf g'a_G\gamma_*^{-1/2}L^{-15/32}\|D\|\le c_sL^{-15/32}\|D\|,}\] because $36\sqrt3\le62.36\le63$, $2\lambda_\flat\le\lambda_*$ and $1-\frac{31}{32}-\frac12=-\frac{15}{32}$. For coarse cells, ${\rm dist}(y',P)$ is $|y'-P|$ measured in the units of the level $\mathbb T'$ (\Cref{Y-A-Distances}), and $e^{-\lambda_*}\le e^{-\lambda_\flat}$; so every exponential factor is at most $e^{-\lambda_\flat(|y'-P|+|y'-P'|)}$, which proves the first assertion of the lemma and, with it, \textup{(S0)(b)}.

We turn to \textup{(S0)(c)} and \textup{(S0)(d)}.

\begin{enumerate}\item[\textup{(c)}] By \textup{(S0)(a)}, \[(T_\ell D)_{y'}(h'_\iota,h'_{\iota'})=L\sum_{y\in B(y')}D_y(L^{-2}h_\iota,L^{-2}h_{\iota'})=L^{-3}\sum_{y\in B(y')}c_y(D)(\iota,\iota').\]

\item[\textup{(d)}] The cut sets entries to $0$; hence it does not raise the norm. The operator norm of a $3\times3$ matrix is at most the sum of the moduli of its entries, and the entries that the cut removes have $|y'-P|>\varrho$ or $|y'-P'|>\varrho$; the sum of their moduli is therefore \[\le\|D'\|\cdot2\cdot(3\sum_{|u|>\varrho}e^{-\lambda_\flat|u|})(3\mathsf g_\flat)\le18\mathsf g_\flat\mathsf g'e^{-\lambda_\flat\varrho/2}\|D'\|.\] At $\varrho=\varrho(x)$, finally, $e^{-\lambda_\flat\varrho(x)/2}\le x^{-\lambda_*c_R/32}=x^{-e_R/16}\le x^{-16}$.\end{enumerate}\end{proof}

\emph{Step 3: two sets of rows on one torus.} This step proves \textup{(S0)(e)}.

\begin{lemma}[The step at two sets of rows]\label{Y-H2-L53} Let $\ell,\tilde\ell\in\mathcal R(b_*,K_*,\lambda_*,\frac19)$ be two sets of rows on $\mathbb T$ for one and the same $a$, let the objects defined with $\tilde\ell$ carry a tilde, and let $\delta=\|\tilde{\mathsf E}-\mathsf E\|_{(\lambda_\flat)}$. Then the following hold.
\begin{enumerate}
\item[\textup{(a)}] $\|\tilde{\mathsf B}-\mathsf B\|\le3\mathsf g_\flat^2\delta$.
\item[\textup{(b)}] For every $X'$, $|(\tilde H-H)X'|\le3\mathsf g_\flat^2\delta(b_*/L)^{1/2}|X'|$.
\item[\textup{(c)}] For every $P$, \[(\sum_pe^{2\lambda_\flat\mathsf d(p,P)}|(\tilde H-H)(p,P)|^2)^{1/2}\le a_HL^{-1/2}\delta^{15/16}.\]
\item[\textup{(d)}] The statement \textup{(S0)(e)} holds for the pair $\ell,\tilde\ell$.
\end{enumerate}
\end{lemma}

\begin{proof}[Proof of \Cref{Y-H2-L53}] We take the four clauses in turn.
\begin{enumerate}
\item[\textup{(a)}] The entries of the difference of the two forms are $(\tilde{\mathsf B}-\mathsf B)(p,p')=\sum_y(\tilde{\mathsf E}_y-\mathsf E_y)(p,p')$, and by the definition of $\delta$ each of them has modulus at most $\delta\sum_ye^{-\lambda_\flat(|y-p|+|y-p'|)}$. For a fixed $p$, the sum of these bounds over $p'$ does not exceed $\delta\cdot\mathsf g_\flat\cdot3\mathsf g_\flat$. The difference is a symmetric matrix, so its column sums obey the same bound as its row sums, and Schur's test yields \textup{(a)}.

\item[\textup{(b)}] Fix $X'$, write $\mathsf B_1=\mathsf B$ and $\mathsf B_2=\tilde{\mathsf B}$, and for $i=1,2$ let $(F_i,\mathsf n_i)$ be the minimiser at this $X'$ for the form $\mathsf B_i$, and $\|\cdot\|_i$ the norm $\|\cdot\|_{\mathcal H}$ built with $\mathsf B_i$. Whether a pair is admitted at $X'$ does not depend on the rows; hence both pairs are admitted at $X'$ for either form.

\Cref{Y-A-L26}\,(a), applied in the model with $\mathsf B=\mathsf B_1$ to the admitted pair $(F_2,\mathsf n_2)$, gives $\|(F_2,\mathsf n_2)-(F_1,\mathsf n_1)\|_1^2=\|(F_2,\mathsf n_2)\|_1^2-\|(F_1,\mathsf n_1)\|_1^2$. Applied in the model with $\mathsf B=\mathsf B_2$ to the admitted pair $(F_1,\mathsf n_1)$, that is with the indices $1$ and $2$ exchanged, it gives $\|(F_1,\mathsf n_1)-(F_2,\mathsf n_2)\|_2^2=\|(F_1,\mathsf n_1)\|_2^2-\|(F_2,\mathsf n_2)\|_2^2$.

We add these two identities and write ${\rm d}F=F_2-F_1$, ${\rm d}\mathsf n=\mathsf n_2-\mathsf n_1$, ${\rm d}\mathsf B=\mathsf B_2-\mathsf B_1$. On the right the terms in $\mathsf n_1$ and $\mathsf n_2$ cancel, and since ${\rm d}\mathsf B$ is symmetric we arrive at \[\langle{\rm d}F,(\mathsf B_1+\mathsf B_2){\rm d}F\rangle+2a|{\rm d}\mathsf n|^2=\langle F_1,{\rm d}\mathsf BF_1\rangle-\langle F_2,{\rm d}\mathsf BF_2\rangle=-\langle{\rm d}F,{\rm d}\mathsf B(F_1+F_2)\rangle.\]

As $\mathsf B_i\ge1$, the left side is at least $2|{\rm d}F|^2$, while the modulus of the right side is at most $|{\rm d}F|\,\|{\rm d}\mathsf B\|(|F_1|+|F_2|)$. Thus $2|{\rm d}F|^2\le|{\rm d}F|\,\|{\rm d}\mathsf B\|(|F_1|+|F_2|)$.

Moreover $|F_i|^2\le\|(F_i,\mathsf n_i)\|_i^2=\langle X',\mathsf B_i'X'\rangle/L\le b_*|X'|^2/L$, where the equality is given by \Cref{Y-A-L26}\,(a), and the second inequality by the bound $\|\mathsf B_i'\|\le\frac32/\mathsf r_\sharp\le b_*$ of \Cref{Y-A-Th29}\,(a). We insert this into the last inequality and divide by $2|{\rm d}F|$ where this is not $0$; the result is $|{\rm d}F|\le\|{\rm d}\mathsf B\|(b_*/L)^{1/2}|X'|$. Since ${\rm d}F=(\tilde H-H)X'$, the bound \textup{(a)} on $\|{\rm d}\mathsf B\|$ turns this into \textup{(b)}.

\item[\textup{(c)}] Fix $P$ and let $v=(\tilde H-H)\delta_P$; thus $v=\tilde\nu-\nu$ is the difference of the columns of $\tilde H$ and $H$ at $P$. Write $\mathsf d$ for $\mathsf d(p,P)$. Since $\lambda_\flat=\lambda_*/16$, H\"older's inequality with the exponents $16$ and $16/15$ gives \[\sum e^{2(\lambda_*/16)\mathsf d}v^2\le(\sum e^{2\lambda_*\mathsf d}v^2)^{1/16}(\sum v^2)^{15/16}.\]

To bound the first sum on the right we use $v^2\le2\tilde\nu^2+2\nu^2$ together with \Cref{Y-A-P28}\,(a), applied to each of the two columns with $\lambda''=\lambda_*$ and $\gamma_{\rm e}=\gamma_*$ as in the proof of \Cref{Y-H2-L51}\,(a), where the bound for one column is $4/(\gamma_*L)$; this sum is then at most $16/(\gamma_*L)$. The second sum equals $|(\tilde H-H)\delta_P|^2$, and \textup{(b)} at $X'=\delta_P$ bounds it by $9\mathsf g_\flat^4\delta^2b_*/L$. The square root of the product of the two bounds is $(4\gamma_*^{-1/2})^{1/16}(3\mathsf g_\flat^2b_*^{1/2})^{15/16}\delta^{15/16}L^{-1/2}$, and \textup{(c)} follows.

\item[\textup{(d)}] Let $D$ be a family, and let $\nu,\nu'$ and $\tilde\nu,\tilde\nu'$ denote the columns of $H$ and of $\tilde H$ at $P$ and at $P'$. The definition of the step and the bilinearity of $D_y$ give \[((T_{\tilde\ell}-T_\ell)D)_{y'}(P,P')=L\sum_{y\in B(y')}[D_y(\tilde\nu-\nu,\tilde\nu')+D_y(\nu,\tilde\nu'-\nu')].\] Each of the two sums is at most $L\|D\|(\sum_y{\rm Ab}_y(\cdot)^2)^{1/2}(\sum_y{\rm Ab}_y(\cdot)^2)^{1/2}$, where the two vectors of that sum take the places of the dots.

We bound these factors by the argument of Step~2, now with \textup{(c)} as the input for a difference of columns and with the inequality $\mathsf d(p,P)\ge{\rm dist}(y',P)-\frac12-|u|/L$. This yields \[(\sum_y{\rm Ab}_y(\tilde\nu-\nu)^2)^{1/2}\le\sqrt3\mathsf g'a_He^{\lambda_\flat}L^{-1/2}\delta^{15/16}e^{-\lambda_\flat{\rm dist}(y',P)},\] and \textup{(c)} at $P'$ gives the same bound for $\tilde\nu'-\nu'$, with $P'$ in the place of $P$. For the factors of $\nu$ and of $\tilde\nu'$, the bound for ${\rm Ab}$ of Step~2 gives \[(\sum_y{\rm Ab}_y(\nu)^2)^{1/2}\le2\sqrt3\mathsf g'\gamma_*^{-1/2}e^{\lambda_*}L^{-1/2}e^{-\lambda_*{\rm dist}(y',P)}\] and \[(\sum_y{\rm Ab}_y(\tilde\nu')^2)^{1/2}\le2\sqrt3\mathsf g'\gamma_*^{-1/2}e^{\lambda_*}L^{-1/2}e^{-\lambda_*{\rm dist}(y',P')}.\]

Twice the product of the two factors, times $L$, bounds the norm of $(T_{\tilde\ell}-T_\ell)D$ by $12e^{\lambda_\flat+\lambda_*}\mathsf g'^2a_H\gamma_*^{-1/2}\delta^{15/16}\|D\|\le c_I\delta^{15/16}\|D\|$. The charges of $(T_{\tilde\ell}-T_\ell)D$ vanish, because $\tilde Hh'=Hh'=L^{-2}h$ by \textup{(S0)(a)}. This proves \textup{(d)}.
\end{enumerate}
\end{proof}

\emph{Step 4: bounds uniform in $k$ and $K$, with the cuts.} We prove \textup{(S1)}, which is \Cref{Y-H2-U2}, and \textup{(S2)}. That is, we show for $0\le k\le K$ the bound $\|\mathbb D_k[f]\|\le K_5\|f\|_\infty$ and \[|c_y(\mathbb D_k[f])-{\rm diag}\,{\rm Mn}^{[K]}_{k,y,\iota}(f)|\le0.0102L^{-1}(x_s/x_k)^{16}\|f\|_\infty,\] and further \textup{(U2.2)}, \textup{(U2.3)} and the covariance.

We begin with the rows. Since $\mathsf t_*(\varrho(x))\le3K_*\mathsf g(\lambda_*/2)x^{-e_R}$ and $x_j\ge L^{K-j}x_s$, \textup{(X12)} yields the condition \textup{(Cut)} of \Cref{Y-A-Th29}\,(b); this deduction is the one made in Step~\textup{(Cut)} of the proof of \Cref{Y-A-Th36}. Consequently $\ell_k\in\mathcal R(b_*,K_*,\lambda_*,\frac19)$ for $0\le k\le K$: for $k<K$ by \Cref{Y-A-Fr}\,\textup{(Fr0)}, and for $k=K$ by \Cref{Y-A-Th29}\,(b). So \textup{(S0)} applies at every step.

We normalise $\|f\|_\infty=1$ and put ${\rm Nm}_k:=\|\mathbb D_k[f]\|$, ${\rm Ch}_k:=|c(\mathbb D_k[f])|$ and ${\rm ct}_i:=x_{i+1}^{-16}$. At the start, $\mathbb D_{0,y}$ is diagonal and is carried by the $p$ with $|y-p|=0$; hence ${\rm Nm}_0\le1$. Its charge is $c_y(\mathbb D_0)={\rm diag}_\iota f(x(y;\iota))$, so that ${\rm Ch}_0\le1$.

From \textup{(S0)(b)}, \textup{(S0)(c)} and \textup{(S0)(d)} we obtain the recursion \[{\rm Nm}_{k+1}\le\vartheta_s{\rm Nm}_k+C_{\rm ch}{\rm Ch}_k,\] \[{\rm Ch}_{k+1}\le{\rm Ch}_k+t_5{\rm ct}_k(\vartheta_s{\rm Nm}_k+C_{\rm ch}{\rm Ch}_k);\] here we use that a mean of matrices of norm at most ${\rm Ch}_k$ has norm at most ${\rm Ch}_k$.

To solve it, let ${\rm Ch}^*:=\max\{1,\max_k{\rm Ch}_k\}$. Iterating the first inequality from ${\rm Nm}_0\le1$ and summing the geometric series, we find with $\vartheta_s\le\frac14$ and $C_{\rm ch}\ge1$ that ${\rm Nm}_k\le\vartheta_s^k+\frac43(1-\vartheta_s^k)C_{\rm ch}{\rm Ch}^*\le\frac43C_{\rm ch}{\rm Ch}^*$; it follows that $\vartheta_s{\rm Nm}_k+C_{\rm ch}{\rm Ch}_k\le\frac43C_{\rm ch}{\rm Ch}^*$. The cut errors form a geometric series: \[\sum_{i<k}{\rm ct}_i=\sum_{j=1}^k(x_kL^{k-j})^{-16}\le x_k^{-16}/(1-L^{-16})\le1.01x_k^{-16}.\]

Summing the second inequality of the recursion and using $x_k\ge x_s$, we get ${\rm Ch}^*\le1+1.01t_5\frac43C_{\rm ch}x_s^{-16}{\rm Ch}^*\le1+{\rm Ch}^*/(100L)$, the last step by \textup{(X5)}. Therefore ${\rm Ch}^*\le1/0.99\le1.0102$, and ${\rm Nm}_k\le\frac43\cdot1.0102\,C_{\rm ch}\le1.347C_{\rm ch}\le K_5$, which is the first bound.

Next we follow the charges. The mean ${\rm Mn}_{k+1,y'}$ is the mean of the ${\rm Mn}_{k,y}$ over $y\in B(y')$, since $B^{k+1}(y')$ is the disjoint union of the $B^k(y)$. Put ${\rm ce}_k:=\sup_y|c_y(\mathbb D_k)-{\rm diag}\,{\rm Mn}_{k,y}|$. Then ${\rm ce}_0=0$, and by \textup{(S0)(c)} and \textup{(S0)(d)}, ${\rm ce}_{k+1}\le{\rm ce}_k+t_5{\rm ct}_k\frac43C_{\rm ch}{\rm Ch}^*$. Summing these increments and using \textup{(X5)} once more, \[{\rm ce}_k\le1.01x_k^{-16}t_5\frac43C_{\rm ch}\cdot1.0102\le(1.0102/100L)(x_s/x_k)^{16},\] and since $1.0102/100\le0.0102$ this is the second bound.

The remaining clauses are shown by induction on $k$. Symmetry, linearity in $f$ and reality pass from one level to the next, because $H$ is real and both the step and the cut are linear. The bound on the range comes from the cut. \textup{(U2.3)} holds since the step sums over $y\in B(y')$ while the cut acts on $P,P'$ alone. The covariance follows from three facts: $H_k$ commutes with the translations of $\mathbb T_{k+1}$, as the proof of \Cref{Y-A-P28}\,(a) shows; translations map blocks to blocks; and the cut is defined by a distance. No positivity enters this argument, so $f$ may be complex.

\emph{Step 5: the start of the comparison, Riemann means at barycentres.} On $\mathbb T_k$ let ${\rm d}\mathbb D_k:=\tilde{\mathbb D}_{m+k}[f]-\mathbb D_k[f]$. We show that $|c({\rm d}\mathbb D_0)|\le{\rm osc}_K(f)+0.0102L^{-16K-1}\|f\|_\infty$ and $\|{\rm d}\mathbb D_0\|\le(K_5+1)\|f\|_\infty$, that ${\rm osc}_K(f)\le\sqrt3a_K{\rm Lip}f$, and that ${\rm os}_{i+1}\le{\rm os}_i$.

By \textup{(S2)} in the run $K'$ at its level $m$, where $\tilde x_m=L^K\mathsf b$, the charge $c_y(\tilde{\mathbb D}_m[f])$ coincides with ${\rm diag}\,{\rm Mn}^{[K']}_{m,y,\iota}(f)$ up to an error of at most $0.0102L^{-1}(x_s/(L^K\mathsf b))^{16}\le0.0102L^{-16K-1}$ times $\|f\|_\infty$. In the run $K$ we have exactly $c_y(\mathbb D_0[f])={\rm diag}f(x^{[K]}(y;\iota))={\rm diag}\,{\rm Mn}^{[K]}_{0,y,\iota}(f)$. These two facts give the bound on the charge. The bound on the norm follows from \textup{(S1)} in the run $K'$ and from ${\rm Nm}_0\le1$, that is $\|\mathbb D_0[f]\|\le\|f\|_\infty$.

For the distances, let $z\in B^m(y)$. Then $a_{K'}|z-L^my|_\infty\le a_{K'}(L^m-1)/2=(a_K-a_{K'})/2$. The offsets of the two barycentres $x^{[K']}(z;\iota)$ and $x^{[K]}(y;\iota)$ differ by $(a_K-a_{K'})/2$ in the directions $\mu,\nu$, and by $0$ in the normal one. Hence $|x^{[K']}(z;\iota)-x^{[K]}(y;\iota)|_\infty\le a_K-a_{K'}<a_K$, and the Euclidean distance of the two points is $<\sqrt3a_K$. So each term of the mean differs from $f(x^{[K]}(y;\iota))$ by $\le\sqrt3a_K{\rm Lip}f$, which is the bound on ${\rm osc}_K(f)$; for $f$ with a modulus of continuity $\omega_f$, each term differs from it by $\le\omega_f(\sqrt3a_K)$, so that ${\rm osc}_K(f)\le\omega_f(\sqrt3a_K)$, which is \textup{(S6)}\,(iii).

Finally, $B^{m+i+1}(y')$ and $B^{i+1}(y')$ are the unions over $y\in B(y')$ of the blocks $B^{m+i}(y)$ and $B^i(y)$, respectively. Hence the difference of the two means at level $i+1$ is the mean of the differences at level $i$, and ${\rm os}_{i+1}\le{\rm os}_i$.

\emph{Step 6: the linear system for the difference of the two runs.} We prove \textup{(S3)}. Let $\|f\|_\infty=1$, and now let ${\rm Nm}_k:=\|{\rm d}\mathbb D_k\|$ and ${\rm Ch}_k:=|c({\rm d}\mathbb D_k)|$. We show that \[{\rm Nm}_k\le1.347C_{\rm ch}{\rm Ch}_0+\vartheta_5^{k-1}[(\vartheta_5+0.0102){\rm Nm}_0+1.0102kI_5]\] for $1\le k\le K$, and we derive from this \textup{(S3)} and its finer form.

Partners have one and the same $x$, hence one and the same cut radius; and the step and the cut are linear. Therefore \[{\rm d}\mathbb D_{k+1}={\rm cut}_{\varrho(x_{k+1})}[T_{\ell_k}{\rm d}\mathbb D_k+I_k],\] with $I_k:=(T_{\tilde\ell_{m+k}}-T_{\ell_k})\tilde{\mathbb D}_{m+k}[f]$.

By Step~4, applied to both runs, the rows $\ell_k$ and $\tilde\ell_{m+k}$ lie in $\mathcal R(b_*,K_*,\lambda_*,\frac19)$ on $\mathbb T_k$. We bound $I_k$ with \textup{(S0)(e)}, with \textup{(S1)} in the run $K'$, and with \textup{(RU)} of item~3 of \S\ref{Y-J-1}, which gives $\delta_k:=\|\tilde{\mathsf E}_{m+k}-\mathsf E_k\|_{(\lambda_\flat)}\le15C_1^2\vartheta_1^k$; the hypotheses of \textup{(RU)} are verified in Step~1 of the proof of \Cref{Y-J-ThA}. The result is $\|I_k\|\le c_IK_5(15C_1^2)^{15/16}\vartheta_1^{15k/16}\le I_5\vartheta_5^k$, and the charges of $I_k$ are $0$.

From \textup{(S0)(b)}, \textup{(S0)(c)}, \textup{(S0)(d)} and $\vartheta_s\le\vartheta_5$ we obtain the recursion \[{\rm Nm}_{k+1}\le\vartheta_5{\rm Nm}_k+C_{\rm ch}{\rm Ch}_k+I_5\vartheta_5^k,\] \[{\rm Ch}_{k+1}\le{\rm Ch}_k+t_5{\rm ct}_k(\vartheta_5{\rm Nm}_k+C_{\rm ch}{\rm Ch}_k+I_5\vartheta_5^k).\]

We solve it as follows, with ${\rm Ch}^*_k:=\max_{i\le k}{\rm Ch}_i$.
\begin{enumerate}
\item[\textup{(1)}] Iterating the first inequality and summing the geometric series with $\vartheta_5\le\frac14$, we get ${\rm Nm}_k\le\vartheta_5^k{\rm Nm}_0+kI_5\vartheta_5^{k-1}+\frac43C_{\rm ch}{\rm Ch}^*_{k-1}$.

\item[\textup{(2)}] Hence, as $\frac43\vartheta_5+1\le\frac43$, \[\vartheta_5{\rm Nm}_i+C_{\rm ch}{\rm Ch}_i+I_5\vartheta_5^i\le\vartheta_5^{i+1}{\rm Nm}_0+(i+1)I_5\vartheta_5^i+\frac43C_{\rm ch}{\rm Ch}^*_i.\]

\item[\textup{(3)}] We have ${\rm ct}_i=(x_kL^{k-1-i})^{-16}$ and $L^{-16}/\vartheta_5\le L^{-16+15/32}\le0.0099$; summing a geometric series of ratio at most $0.0099$, we get $\sum_{i<k}{\rm ct}_i\vartheta_5^i\le1.01x_k^{-16}\vartheta_5^{k-1}$.

\item[\textup{(4)}] The sums over $i<k$ that bound ${\rm Ch}_k$ grow with $k$, so they also bound ${\rm Ch}^*_k$. With \textup{(2)}, \textup{(3)} and the bound $1.01x_k^{-16}$ for $\sum_{i<k}{\rm ct}_i$ we obtain \[{\rm Ch}^*_k\le{\rm Ch}_0+1.01t_5x_k^{-16}[\vartheta_5^k{\rm Nm}_0+kI_5\vartheta_5^{k-1}]+1.01t_5x_k^{-16}\frac43C_{\rm ch}{\rm Ch}^*_k.\]

\item[\textup{(5)}] By \textup{(X5)}, $1.01t_5x_k^{-16}\le3/(400C_{\rm ch}L)$, so the last term of \textup{(4)} is at most ${\rm Ch}^*_k/(100L)$; moreover $1/(L\vartheta_5)\le L^{-17/32}\le1$. We move the last term to the left and divide by $1-1/(100L)$, whose inverse is at most $1.0102$; this gives \[{\rm Ch}^*_k\le1.0102{\rm Ch}_0+1.0102(0.0075/C_{\rm ch})({\rm Nm}_0+kI_5)\vartheta_5^k\le1.0102{\rm Ch}_0+(0.0076/C_{\rm ch})({\rm Nm}_0+kI_5)\vartheta_5^k.\]

\item[\textup{(6)}] We insert \textup{(5)}, taken at $k-1$, into \textup{(1)}. Since $\frac43\cdot1.0102\le1.347$ and $\frac43\cdot0.0076\le0.01014\le0.0102$, the outcome is the claim.
\end{enumerate}

We now derive \textup{(S3)}. By Step~5, ${\rm Nm}_0\le K_5+1$ and ${\rm Ch}_0\le{\rm osc}_K(f)+0.0102L^{-16K-1}$. Although the first of these is not small, the contraction does the work, since ${\rm Nm}_0$ appears in the claim with the coefficient $\vartheta_5^{k-1}(\vartheta_5+0.0102)$. Further, we have $1.347C_{\rm ch}\le K_5$ and $L^{-16K-1}\vartheta_5^{1-k}\le L^{-16K-1+15K/32}\le1$. Therefore \[\FitDisplay{\dimexpr\linewidth-1pt\relax}{\begin{aligned}
&{\rm Nm}_k
\\
&\quad{}\le K_5{\rm osc}_K(f)
\\
&\quad{}+\vartheta_5^{k-1}[(0.25+0.0102)(K_5+1)+0.0102K_5+1.0102kI_5]
\\
&\quad{}\le K_5{\rm osc}_K(f)
\\
&\quad{}+\vartheta_5^{k-1}[0.2704(K_5+1)+1.0102kI_5]
\\
&\quad{}\le K_5{\rm osc}_K(f)+C_5k\vartheta_5^{k-1}.
\end{aligned}}\]

For the finer form, let $d_i:=\sup_y|c_y({\rm d}\mathbb D_i)-{\rm diag}({\rm Mn}^{[K']}_{m+i,y}-{\rm Mn}^{[K]}_{i,y})|$, so that ${\rm Ch}_i\le{\rm os}_i+d_i$. As the means average exactly and $I_i$ has charges $0$, we have $d_{i+1}\le d_i+t_5{\rm ct}_i(\vartheta_5{\rm Nm}_i+C_{\rm ch}{\rm Ch}_i+I_5\vartheta_5^i)$, and $d_0\le0.0102L^{-16K-1}$. If ${\rm os}_i\le\omega_5+P_5\vartheta_5^i$, then ${\rm Ch}_i\le(\omega_5+d_i)+P_5\vartheta_5^i$, and the pair $({\rm Nm}_i,\omega_5+d_i)$ obeys the recursion above with $I_5+C_{\rm ch}P_5$ in place of $I_5$ and with $\omega_5+d_0$ in place of ${\rm Ch}_0$, so that the claim, applied to this pair, gives \[{\rm Nm}_k\le1.347C_{\rm ch}(\omega_5+d_0)+\vartheta_5^{k-1}[(\vartheta_5+0.0102){\rm Nm}_0+1.0102k(I_5+C_{\rm ch}P_5)].\] As in the derivation of \textup{(S3)}, we now use $1.347C_{\rm ch}\le K_5$, $d_0\le0.0102L^{-16K-1}$, $L^{-16K-1}\vartheta_5^{1-k}\le1$, ${\rm Nm}_0\le K_5+1$, $\vartheta_5\le\frac14$, $1.0102\le1.02$, and $k\ge1$, and obtain \[{\rm Nm}_k\le K_5\omega_5+\vartheta_5^{k-1}[0.2704(K_5+1)+1.0102kI_5+1.0102kC_{\rm ch}P_5]\le K_5\omega_5+(C_5+1.02C_{\rm ch}P_5)k\vartheta_5^{k-1},\] and the last expression is the right side of the finer form of \textup{(S3)}.

Nothing is lost at the last levels. The errors of the cuts are of order $x_k^{-16}$, which is not small in $K$ when $k$ is near $K$. Since both runs are cut at the same radius, and since \textup{(S0)(a)} eliminates $\tau(\ell_k)$ completely, such an error enters only as a factor of ${\rm Nm}_i$, of ${\rm Ch}_i$ or of $\vartheta_5^i$.

\emph{Step 7: the size of $\vartheta_5$, and the order of the numbers.} In this step we prove \textup{(S3c)}, the bound $\vartheta_5\le L^{-1/5}$ and \Cref{Y-H2-Rem2}. We show that $L^{-15/32}\le\vartheta_5\le L^{-1/5}\le\frac14$, that the right side of \textup{(S3)} is at most $C_5^\#k\vartheta_5^k\|f\|_\sharp$, and that $c_s\le c_{\rm fix}$; and we show that the numbers are fixed in an order without a circle.

First the sizes. Since $c_s\ge63$, we have $\vartheta_5\ge\vartheta_s\ge L^{-15/32}$. By \textup{(L5*)}, $c_sL^{-15/32}\le L^{43/160-75/160}=L^{-1/5}$. Next, $\vartheta_1=c_{\rm fix}L^{-1/2}\le L^{-2/5}$, because $L\ge L_B\ge c_{\rm fix}^{10}$ by Step~1 of the proof of \Cref{Y-J-ThA}; hence $\vartheta_1^{15/16}\le L^{-3/8}\le L^{-1/5}$. Finally $L\ge1024$ gives $L^{-1/5}\le\frac14$.

For \textup{(S3c)} we use $\vartheta_5^{k-1}\le L^{15/32}\vartheta_5^k$, $a_K=L^{-K}\le\vartheta_5^K\le\vartheta_5^k$, $\sqrt3/(40M)\le1$, $\|f\|_\infty\le\|f\|_\sharp$ and $k\ge1$. Combined, they give \[K_5{\rm osc}_K(f)+C_5k\vartheta_5^{k-1}\|f\|_\infty\le(K_5+L^{15/32}C_5)k\vartheta_5^k\|f\|_\sharp.\]

For \Cref{Y-H2-Rem2}, the closed forms in the proof of the bounds of \Cref{Y-A-Sums} give, with $t=\lambda/2\le0.2$, \[\mathsf g(\lambda)\ge24\cosh t/(4\sinh^3t)\ge24/(4(1.01t)^3)\ge46\lambda^{-3}.\] By \textup{(O0)} in item~1 of \S\ref{Y-J-1},
\[
c_0=108b_*\gamma_*^{-1/2}\mathsf g(\lambda_*/4)^3,\qquad
C_1=\sqrt3K_*\mathsf g(\lambda_*/2)^{3/2},\qquad
c_{\rm fix}=48\mathsf g_\flat C_1c_0^{1/2}.
\]
Use the exact definitions and outward intervals in the order of the constants, \textup{(0)}, of \Cref{Y-A-Fr}. Substitution of $\mathsf g(t)\ge46t^{-3}$ into these three formulas gives
$\log_{10}c_{\rm fix}\ge204.892>204.89$.
For the other direction, substitute
$\mathsf g(t)\le51t^{-3}$ and $\mathsf g_1(t)\le182t^{-4}$ from \Cref{Y-A-Sums} into
\[
a_G=(23b_*^2)^{15/16}(7\gamma_*^{-1/2})^{1/16},\qquad
c_s=63e^{2\lambda_*}\mathsf g_1(\lambda_*/32)
             \mathsf g(\lambda_*/32)a_G\gamma_*^{-1/2}.
\]
They give $\log_{10}c_s\le145.332<145.34$.
So $c_s\le c_{\rm fix}\le c_{\rm fix}^{43/16}\le L_B^{43/160}\le L^{43/160}$.

It remains to fix the order of the numbers. The numbers of \textup{(O0a)} depend on $b_*,\gamma_*,\lambda_*,C_1$ alone. \textup{(L5*)} bounds $L$ from below in terms of \textup{(O0)} only, and \textup{(X5)} bounds $x_s$ from below in terms of \textup{(O0)} and \textup{(O4)}. The numbers $\vartheta_s,\vartheta_5,C_5^\#$ are fixed afterwards, and no floor is defined through them; so there is no circle. None of these numbers depends on $\mathsf b,N,K,K',k,f,M,\kappa,p_0,c_h,E_a$.

\emph{Step 8: the hybrid chain.} We prove \textup{(S4)}: for $1\le k_0\le k\le K$, \[\|\mathbb D^b_k-\mathbb D_k[f]\|\le2.1K_5{\rm osc}_K(f)+1.2C_5k_0\vartheta_5^{k_0-1}\|f\|_\infty.\]

Let $\|f\|_\infty=1$. From $k_0$ on, the difference $\mathbb D^b_k-\mathbb D_k$ obeys the homogeneous system of Step~4, since both chains then use one set of rows and one cut. Let ${\rm Nm}_k$ and ${\rm Ch}_k$ be its norm and its charge, and ${\rm Ch}^*:=\max_{k_0\le i\le K}{\rm Ch}_i$. Iterating from $k_0$ with $\vartheta_s\le\frac14$ we find ${\rm Nm}_k\le\vartheta_s^{k-k_0}{\rm Nm}_{k_0}+\frac43C_{\rm ch}{\rm Ch}^*$, and therefore $\vartheta_s{\rm Nm}_i+C_{\rm ch}{\rm Ch}_i\le{\rm Nm}_{k_0}/4+\frac43C_{\rm ch}{\rm Ch}^*$.

Summing the inequality for the charges from $k_0$ on, we get ${\rm Ch}^*\le{\rm Ch}_{k_0}+1.01t_5x_s^{-16}[{\rm Nm}_{k_0}/4+\frac43C_{\rm ch}{\rm Ch}^*]$, and with \textup{(X5)} this becomes ${\rm Ch}^*\le1.0102[{\rm Ch}_{k_0}+(3/(400C_{\rm ch}L)){\rm Nm}_{k_0}]$. Inserted into the bound for ${\rm Nm}_k$, it yields ${\rm Nm}_k\le1.0102{\rm Nm}_{k_0}+1.347C_{\rm ch}{\rm Ch}_{k_0}$.

Into this we insert \textup{(S3)} for ${\rm Nm}_{k_0}$, and \textup{(5)} of Step~6 for ${\rm Ch}_{k_0}$, which gives \[{\rm Ch}_{k_0}\le1.0102({\rm osc}_K+0.0102L^{-16K-1})+(0.0076/C_{\rm ch})(K_5+1+k_0I_5)\vartheta_5^{k_0}.\] Moreover $1.347\cdot1.0102\,C_{\rm ch}\le1.3608C_{\rm ch}=1.008K_5$. So the coefficient of ${\rm osc}_K$ is $(1.0102+1.008)K_5=2.0182K_5\le2.1K_5$.

We measure the remaining terms in units of $C_5k_0\vartheta_5^{k_0-1}$, using $K_5\le C_5/0.3$, $K_5+1+k_0I_5\le k_0C_5/0.3$, $L^{-16K-1}\le\vartheta_5^{k_0-1}$ and $\vartheta_5\le\frac14$. Their sum is at most \[1.0102+1.008\cdot0.0102/0.3+1.347\cdot0.0076/(4\cdot0.3)\le1.0102+0.0343+0.0086=1.0531\le1.2,\] which proves the claim.

\emph{Step 9: totals, several sources, and the number $\delta_\Delta$.} We prove the three parts of \textup{(S5)}.
\begin{enumerate}
\item[\textup{(i)}] Let $D$ be a family. Then \[|\sum_yD_y(p,p')|\le\|D\|\sum_ye^{-\lambda_\flat(|y-p|+|y-p'|)}\le\|D\|\mathsf g'e^{-(\lambda_\flat/2)|p-p'|},\] where the second inequality is obtained by the computation in \Cref{Y-A-SetsOfRows}, in which the triangle inequality gives $\sum_ye^{-\lambda(|y-p|+|y-p'|)}\le e^{-\frac\lambda2|p-p'|}\sum_ye^{-\frac\lambda2|y-p|}$. For fixed $p$, the middle term of the chain in the display, summed over $p'$, is $\|D\|\sum_{p'}\sum_ye^{-\lambda_\flat(|y-p|+|y-p'|)}\le\|D\|\,\mathsf g_\flat\cdot3\mathsf g_\flat$, as in the proof of \Cref{Y-H2-L53}\,(a). Hence the row sums of $\sum_yD_y$ are at most $3\mathsf g_\flat^2\|D\|$, and as $\sum_yD_y$ is symmetric, Schur's test bounds its norm by the same number; for $D={\rm d}\mathbb D_k$ this gives the two bounds of \textup{(S5)}\,(i), since then $\sum_yD_y=\tilde N_{m+k}[f]-N_k[f]$ and $\|D\|$ is at most the right side of \textup{(S3)}.

\item[\textup{(ii)}] The linearity is \textup{(U2.2)}. At each point at most one of the $f_i$ differs from $0$, which gives the claim on the sup norm.

For the upper bound on ${\rm Lip}w$, let $w(x)=z_if_i(x)$ and $w(x')=z_jf_j(x')$ with $i\ne j$ and both values not $0$, and take a shortest segment from $x$ to $x'$. Let $\xi$ be its first point with $f_i=0$, and $\xi'$ its last point with $f_j=0$. These points exist: the $f_j$ are Lipschitz, hence continuous, and $f_i(x')=0$, $f_j(x)=0$, since the sets $\{f_i\ne0\}$ are disjoint.

On $[x,\xi)$ we have $f_i\ne0$, so $f_j=0$ there, whereas $f_j\ne0$ on $(\xi',x']$. Hence $\xi$ lies before $\xi'$ on the segment, or is equal to it, and \[|w(x)-w(x')|\le|z_i|{\rm Lip}f_i|x-\xi|+|z_j|{\rm Lip}f_j|\xi'-x'|\le\max_i|z_i|{\rm Lip}f_i\,|x-x'|.\]

The remaining cases are direct. If $i=j$, then $|w(x)-w(x')|=|z_i||f_i(x)-f_i(x')|\le|z_i|{\rm Lip}f_i|x-x'|$. If one value is $0$, say $w(x')=0$ while $w(x)=z_if_i(x)\ne0$, then $z_i\ne0$ forces $f_i(x')=0$, and again $|w(x)-w(x')|=|z_i||f_i(x)-f_i(x')|\le|z_i|{\rm Lip}f_i|x-x'|$. If both values are $0$, there is nothing to prove.

For the lower bound ${\rm Lip}w\ge|z_i|{\rm Lip}f_i$, we use once more that the $f_j$ are continuous, being Lipschitz, and take two points $x,x'$. If $f_i\ne0$ at both, then $w=z_if_i$ at both, and $|z_i||f_i(x)-f_i(x')|=|w(x)-w(x')|\le{\rm Lip}w|x-x'|$. If $f_i=0$ at both, there is nothing to prove.

Up to exchanging $x$ and $x'$, the last case is $f_i(x)\ne0$ and $f_i(x')=0$. Let $\xi$ be the first point with $f_i(\xi)=0$ on a shortest segment from $x$ to $x'$. On $[x,\xi)$ we have $f_i\ne0$, so the other $f_j$ vanish and $w=z_if_i$ there; by continuity $w(\xi)=0$. Hence \[|z_i||f_i(x)-f_i(x')|=|w(x)-w(\xi)|\le{\rm Lip}w|x-\xi|\le{\rm Lip}w|x-x'|.\]

We remark that ${\rm Lip}w=\max|z_i|{\rm Lip}f_i$ holds as well, the value $\infty$ allowed, for bounded $f_i$ with pairwise disjoint closed supports and all $z_i\ne0$. If all $f_i$ are Lipschitz, this is the case just treated. Let ${\rm Lip}f_i=\infty$ for one $i$. The supports are compact and finitely many, so they lie at a positive least distance from each other, and $w=z_if_i$ at all points nearer to ${\rm supp}f_i$ than that distance.

Take $x_n,x_n'$ with $|f_i(x_n)-f_i(x_n')|/|x_n-x_n'|\to\infty$. As $f_i$ is bounded, $|x_n-x_n'|\to0$. One of the two points lies in ${\rm supp}f_i$, since otherwise both values would be $0$; so from some $n$ on, both points lie nearer to ${\rm supp}f_i$ than the least distance between the supports, and the difference quotients of $w$ are $|z_i|$ times those of $f_i$. Thus ${\rm Lip}w=\infty$.

\item[\textup{(iii)}] For the family $D$ of this part, $\|D\|\le\|\tilde{\mathsf E}_{m+k}-\mathsf E_k\|+\|\tilde{\mathbb D}_{m+k}[w]-\mathbb D_k[w]\|/\mathsf b$. We multiply by $(c_Q^{\rm su}/E_a)x_k^{1/2}/p_0(x_k)={\rm amp}(x_k)\mathsf b/K_5$. Then the first term is bounded by \Cref{Y-I-CorD}, and the second by \textup{(S3)} for $w$.
\end{enumerate}

\emph{Step 10: the class of $f$.} We prove \textup{(S6)}\,(i) and (ii).
\begin{enumerate}
\item[\textup{(i)}] For $K\ge3$ and $m\ge1$, there exists a bounded Borel $f$ for which \textup{(S3)} fails once the term $K_5{\rm osc}_K(f)$ is deleted.
\item[\textup{(ii)}] For the indicator of a box, ${\rm os}_i\le12L^{-i}$.
\end{enumerate}

We first note a lower bound for the norm: for every family $D$, $|c_y(D)|\le\sum_{p,p'}|D_y(p,p')|\le9\mathsf g_\flat^2\|D\|$.

\begin{enumerate}
\item[\textup{(i)}] Let $m\ge1$, and let $\mathsf b,L,N,K$ be any data of item~1 subject to \textup{(O0)} to \textup{(O13)}, with $K\ge3$. Thus $L\ge L_B$, and the range $3\le k\le K$ is not empty. Let $f=1$ at the barycentres $x^{[K]}(p)$ of all plaquettes $p$ of level $0$ in the run $K$, and $f=0$ elsewhere. This $f$ is bounded and Borel, and $\|f\|_\infty=1$. We show that \textup{(S3)} with the term $K_5{\rm osc}_K(f)$ deleted fails at every $k$ with $3\le k\le K$.

Indeed, with $\iota$ normal to $e_\mu,e_\nu$ and $\iota'$ normal to $e_{\mu'},e_{\nu'}$, the equation $x^{[K']}(z;\iota)=x^{[K]}(y;\iota')$ says $z+\frac12(e_\mu+e_\nu)=L^m(y+\frac12(e_{\mu'}+e_{\nu'}))$. As $L^m$ is odd, the half-integer coordinates of the two sides agree only if $\iota=\iota'$, and then $z=L^my+\frac{L^m-1}2(e_\mu+e_\nu)$, a single site of $B^m(y)$.

Hence ${\rm Mn}^{[K']}_{m+i,y,\iota}(f)=L^{-3m}$ and ${\rm Mn}^{[K]}_{i,y,\iota}(f)=1$, so ${\rm os}_i=1-L^{-3m}$ for all $i$. By \textup{(S2)} for both runs and the lower bound above, $\|{\rm d}\mathbb D_k\|\ge(1-L^{-3m}-0.0204/L)/(9\mathsf g_\flat^2)$ for all $k\le K$. Since $\vartheta_5\le L^{-1/5}\le\frac14$ by Step~7, we have $(k+1)\vartheta_5^k\le2\vartheta_5\cdot k\vartheta_5^{k-1}\le\frac12k\vartheta_5^{k-1}$ for $k\ge1$, and hence $C_5k\vartheta_5^{k-1}\le3C_5\vartheta_5^2$ for $k\ge3$. By the numerals of item~2 of \S\ref{Y-H2-1}, $\log_{10}C_5\le382.69$ and $\log_{10}\mathsf g_\flat\le46.32$, and by \Cref{Y-H2-Rem2}, $L\ge L_B\ge10^{2048.9}$. Therefore $3C_5\vartheta_5^2\le3\cdot10^{382.69}L^{-2/5}\le10^{0.48+382.69-819.56}\le10^{-436}$ and $(1-L^{-3m}-0.0204/L)/(9\mathsf g_\flat^2)\ge0.99/(9\cdot10^{92.64})\ge10^{-94}$, so that $\|{\rm d}\mathbb D_k\|\ge10^{-94}>10^{-436}\ge C_5k\vartheta_5^{k-1}$ for $3\le k\le K$. What \Cref{Y-H9-T}\,(T2) has and this instance lacks is the continuity of $f$.

\item[\textup{(ii)}] Let $f$ be the indicator of a box. By Step~5, the points $x^{[K']}(z';\iota)$, $z'\in B^m(z)$, lie at sup distance $<a_K$ from $x^{[K]}(z;\iota)$. Therefore the difference ${\rm Mn}^{[K']}_{m,z,\iota}(f)-f(x^{[K]}(z;\iota))$, of modulus at most $1$, vanishes unless for some $\mu$ an end $c$ of the $\mu$-th arc of the box lies between two $\mu$-coordinates that differ by less than $a_K$, that is, unless $|x^{[K]}(z;\iota)_\mu-c|<a_K$ for one of the six pairs $(\mu,c)$.

For one pair, the coordinate $z_\mu$ then lies in an open interval of length $2$, so it takes at most $2$ values; this leaves at most $2L^{2i}$ sites of $B^i(y)$ for each pair, and $12L^{2i}$ in all. As ${\rm os}_i$ is a supremum of means over $L^{3i}$ sites, ${\rm os}_i\le12L^{-i}$. By Step~7, $\vartheta_5\ge L^{-15/32}$, so that $L^{-1}\le\vartheta_5$ and therefore ${\rm os}_i\le12L^{-i}\le12\vartheta_5^i$, which is the second bound of \textup{(S6)}\,(ii).
\end{enumerate}

\emph{Step 11: the uncut chain.} We prove \Cref{Y-H2-Rem1}: the formula for $D^{\rm unc}_k[f]$, the bound $|D^{\rm unc}_{k,y}[f](P,P')|\le243K_*^2\|f\|_\infty e^{-\lambda_*(|y-P|+|y-P'|)}$, and $0\le N^{\rm unc}_k[f]\le\|f\|_\infty\mathsf B_k$ for $f\ge0$.

The formula follows by induction on $k$ from the definition of $T_\ell$, because $B^{k+1}(y')$ is the disjoint union of the $B^k(y)$.

For the bound we introduce auxiliary rows. Let $\tau_0:=1/243$, $\hat a:=(1+\tau_0)a$ and $\hat\ell_{0,y}\zeta:=\tau_0^{1/2}(\zeta_p)_{p\text{ based at }y}$. These rows are translation covariant, $\hat{\mathsf B}_0=(1+\tau_0)1$, $|\hat\ell_{0,y}(p)|\le\tau_0^{1/2}\le K_*$, and $\hat\ell_{0,y}(p)=0$ unless $|y-p|=0$; moreover $\tau(\hat\ell_0)=(3\tau_0)^{1/2}=\frac19$ and $\hat aL\ge3/c_L$. We set $\hat\ell_{k+1}:=(\hat\ell_k)'$, given by \Cref{Y-A-P5} in the model with $\hat a$, without cut.

We have $\langle F,(1+\tau_0)\mathsf BF\rangle+\hat a|\mathsf n|^2=(1+\tau_0)[\langle F,\mathsf BF\rangle+a|\mathsf n|^2]$, and being admitted depends neither on the rows nor on $a$. Hence \Cref{Y-A-L26}\,(a) shows: if $\hat{\mathsf B}_k=(1+\tau_0)\mathsf B_k$, then the minimisers of the two models agree, $\hat H_k=H_k$, and $\hat{\mathsf B}_{k+1}=(1+\tau_0)\mathsf B_{k+1}$. By \Cref{Y-A-Th29}\,(a), in which the constants do not depend on $a$, and by $\tau(\ell')\le\tau(\ell)$ from \Cref{Y-A-P28}\,(c), induction on $k$ gives $\hat\ell_k\in\mathcal R(b_*,K_*,\lambda_*,\frac19)$ and $|\hat\ell_{k,y}(P)|\le K_*e^{-\lambda_*|y-P|}$ for $k\ge1$.

By group~(5) of \Cref{Y-A-P5}, and by induction, $\hat\ell_{k,y}\zeta$ has among its components the numbers $\tau_0^{1/2}L^{k/2}(\mathbb H_k\zeta)_p$, $p$ based in $B^k(y)$. Therefore $L^k\sum_{p\text{ based in }B^k(y)}\mathbb H_k(p,P)^2\le243|\hat\ell_{k,y}(P)|^2\le243K_*^2e^{-2\lambda_*|y-P|}$, and the Cauchy--Schwarz inequality, applied in the formula for $D^{\rm unc}_{k,y}[f]$, gives the bound.

We set $N^{\rm unc}_k[f]:=\sum_yD^{\rm unc}_{k,y}[f]$, the total of \Cref{Y-H2-Rem1}, and turn to the inequality $0\le N^{\rm unc}_k[f]\le\|f\|_\infty\mathsf B_k$ for $f\ge0$. Since $\mathbb T_0$ is the disjoint union of the blocks $B^k(y)$, $y\in\mathbb T_k$, the formula for $D^{\rm unc}_{k,y}[f]$ gives $N^{\rm unc}_k[f]=L^k\mathbb H_k^{\mathsf T}{\rm diag}(f(x(p)))\mathbb H_k$; so $0\le N^{\rm unc}_k[f]\le\|f\|_\infty N^{\rm unc}_k[1]$ for $f\ge0$. Further, we have $N^{\rm unc}_0[1]=1=\mathsf B_0$, and if $N^{\rm unc}_k[1]\le\mathsf B_k$, then by \Cref{Y-A-P4}\,(a) \[\FitDisplay{\dimexpr\linewidth-1pt\relax}{\langle X',N^{\rm unc}_{k+1}[1]X'\rangle=L\langle H_kX',N^{\rm unc}_k[1]H_kX'\rangle\le L\langle H_kX',\mathsf B_kH_kX'\rangle\le L\|(H_kX',\mathsf N^\sharp X')\|_{\mathcal H}^2=\langle X',\mathsf B_{k+1}X'\rangle.}\] So $N^{\rm unc}_k[1]\le\mathsf B_k$ for every $k$, and the inequality $0\le N^{\rm unc}_k[f]\le\|f\|_\infty\mathsf B_k$ for $f\ge0$ follows.
\end{proof}

\subsection{One-step derivative}\label{Y-H3-1}

This section studies one drawn step in the setting \textup{(SD)}, for $\operatorname{SU}(2)$ in $d=3$, with kernels, law, centre and rows free of sources, and differentiates it in the direction of the quadratic member $s_x[D]$, where $D$ is any complex symmetric family of kernels of range $\le\varrho(x)$. The $a$-part of this derivative is the quadratic member $s_{x'}[D']$, with $D'$ one cut step of the free source flow, plus a remainder $n[u;D]$; the remainder splits as $n^{\rm fix}+n^{\rm mem}$, its cell sum is bounded by $u_3(x')\|D\|_{(\lambda_\flat)}$ caps at $x'$, and further clauses treat analyticity, locality, the dependence on the member, linearity in $D$ and the constants. Further results are the derivative of the number, the estimate \textup{(2$\tau$)} of a trace, and a corollary in which the direction is $s_x[D]+{\rm rem}$ and the remainders ${\rm rem}$, ${\rm rem}'$ obey $t'\le0.2431t+1$.

\begin{remark}[Comparison of the bounds on the remainder and on the leading member]
By \Cref{Y-H3-U3}\,(a), the remainder $n$ is at most $u_3(x')\|D\|_{(\lambda_\flat)}$ caps at $x'$, a bound that is not small by itself. The leading member $s_{x'}[D']$, however, is bounded by $c_Q^{\rm su}p_0'^2\|D'\|_{(\lambda_\flat)}$, that is, by $(c_Q^{\rm su}/E_a)x'^{1/2}p_0'^{-1}\|D'\|_{(\lambda_\flat)}$ caps at $x'$; this follows from Step~1\,(1) of the proof of \Cref{Y-H3-U3}, applied at $x'$. The coefficient of $\|D\|_{(\lambda_\flat)}$ in the first bound, divided by that of $\|D'\|_{(\lambda_\flat)}$ in the second, is $\check\omega_3(x')=C_3E_ap_0'x'^{-1/4}/c_Q^{\rm su}$ (entry~7 of item~(0) of \S\ref{Y-H0-1}), and this quotient is small by \textup{(W3)} in item~(1) of \S\ref{Y-H0-1}.
\end{remark}

\emph{Item 1: the letters.} All letters are those of items~(0) and~(1) of \S\ref{Y-H0-1}. The list below gives entries \textup{(O0a)} to \textup{(O14)} of the table which occur in this section, and its entry \textup{(O13)} also defines $x_3$ and the conditions \textup{(Z.0)} to \textup{(Z.4)}.

\begin{enumerate}
\item[\textup{(O0a)}] $\lambda_\flat=\lambda_*/16$; $\mathsf g(t)=\sum_{m\in\mathbb Z^3}e^{-t|m|_\infty}$; $\mathsf g_\flat=\mathsf g(\lambda_\flat)$; $g_{1b}=\sum_m|m|_\infty e^{-\lambda_\flat|m|_\infty}$; $g_b^-=\mathsf g(8\lambda_\flat/9)$.
\item[\textup{(O4)}] $L$ is odd and satisfies \textup{(L2$^{\rm su}$)}, that is $L^{1/2}\ge1.2\cdot10^4\mathsf N_{\rm c}^2$. Moreover \textup{(L3su)} reads $L^{1/4}\ge10^4\mathsf N_{\rm c}^2$.
\item[\textup{(O5)}] $\kappa\ge6$, $\kappa_1=\kappa+1$ and $\tilde\kappa\in[\kappa,\kappa_1]$; $\mathbb P=[\kappa,\kappa_1]\times\{c_h,c_h/2\}$, with elements $p=(\tilde\kappa,\tilde c)$.
\item[\textup{(O6)}] $\gamma=(73L^4)^{-1}$; $\theta_D\le1$; $C_D\ge4$; $c_R=64/\theta_D$.
\item[\textup{(O7)}] The letter $M$.
\item[\textup{(O8)}] $p_0(x)=A_0(\log x)^{\mathfrak p_0}$, $\varepsilon=p_0x^{-1/2}$ and $x'=x/L$; a prime means ``at $x'$''.
\item[\textup{(O9)}] $C_\varphi^\flat$, and $c_\xi=(6.5C_\varphi^\flat)^{-1}$.
\item[\textup{(O10)}] $\mathsf a_{\rm c}$ and $\mathsf K_{\rm c}$; $C_\eta^\flat=1/6.5+\max\{\mathsf a_{\rm c},7\}L^{1/2}$ (\S\ref{Y-E-1}); $c_Q^{\rm su}=306M^3\mathsf g_\flat^2$ (\S\ref{Y-I-1}).
\item[\textup{(O10a)}] $k_1:=0.9475(C_\eta^\flat)^3(1320g_{1b}/\mathsf g_\flat+216)$; $k_2:=0.1359L(g_b^-/\mathsf g_\flat)C_D[4\mathsf K_{\rm c}+1+1.51\cdot10^6\theta_D^{-2}C_D]$; $k_3:=6730L^4$.
\item[\textup{(O12)}] $E_a$, and $c_a(x)=E_ap_0^3x^{-1/2}$.
\item[\textup{(O12a)}] $C_3:=1686e^{\kappa+4}L^{11/4}c_Q^{\rm su}/E_a$.
\item[\textup{(O13)}] $x_3:=1+\inf S_3$, where $S_3$ is the set of $y\ge x_9^U$ such that the five conditions \textup{(Z.0)} to \textup{(Z.4)} below hold for every $x\ge y$, where $T_s:=x^{1/4}$ and $n_\chi:=3(M(2r(x)+1)+2L(r_w+2)+1)^3$:
\begin{enumerate}
\item[\textup{(Z.0)}] $160T_s\omega\varepsilon'\le1$, $2640\ell_w(x')(T_s+1)\omega\varepsilon'\le \mathsf g_x$ and $102400T_s\omega^2\varepsilon'\le1$;
\item[\textup{(Z.1)}] $c_\xi p_0\ge1$;
\item[\textup{(Z.2)}] $7L^3c_\xi p_0n_\chi e^{-c_\xi^2p_0^2/2}\le x^{-10}$;
\item[\textup{(Z.3)}] $17.9L^4C_D^2x^{-31}\le x^{-20}$;
\item[\textup{(Z.4)}] $18.2L^3c_a(x)+L^3(2k_1\varepsilon+k_2\varepsilon')+k_3L^3p_0^{-2}x^{-1/2}+2x^{-10}\le e^{\kappa+3}L^3x^{-1/4}$, with $\kappa$ the lower end of the range of $\tilde\kappa$.
\end{enumerate}
The table, in item~(1), \textup{(O13)}, of \S\ref{Y-H0-1}, requires the last floor $x_s$ to satisfy $Lx_s\ge x_3$. The floor $x_3$ imposes no condition on $d_0,c_\Xi,\rho_0,L,a,\kappa,M,\mathfrak p_0,c_h,E_a$.
\item[\textup{(O14)}] $u_3(x)=C_3x^{1/4}$; $\vartheta_a=0.2431L^{-1/4}$; $\sigma_D=c_Q^{\rm su}p_0(x)^2\|D\|_{(\lambda_\flat)}$.
\end{enumerate}

The radii and domains $\varrho(x)$, $\varrho_L$, $r_w$, $R_2$, $R_3$, $R_T$, $r(x)$, $\ell_w$, $D_w$, $\omega$, $\mathsf g_x$, $\mathcal U_X(x)$, $\mathcal U_Y(x')$, $U^\tau(B;x')$, the chart $U=\mathsf U(\eta;V')$ with $\eta=G\xi-a_0$ and $\xi\sim \mathcal N(0,x^{-1}1)$, and the symbols $\chi_X$, $Y_0(X)$, $F^\flat$ and $\bar F=EX'-\mathrm{d}^Va_0$ are taken from items~2, 3 and~6 of \S\ref{Y-H-1} and from Sections~\ref{Y-D-1}, \ref{Y-C-1} and~\ref{Y-E-1}. From item~2 of \S\ref{Y-I-1} we take $|y-p|=|y-\pi(p)|_\infty$, the norm \[\|D\|_{(\lambda)}=\sup_{y,p,p'}|D_y(p,p')|e^{\lambda(|y-p|+|y-p'|)},\] the paths $\Gamma(y,z)$, and the form $Y\cdot Y'=\frac12\operatorname{tr}YY'$. The transport $T(y',y'')=\operatorname{Ad}_{V'(\Gamma(y',y''))}$ is that of item C3 of \S\ref{Y-D-1}.

\begin{definition}[The cell sum of a family]\label{Y-H3-CS}
Let $a=(a(Y,\cdot))$ be a family indexed by the connected sets $Y$ of cells. Its cell sum is \[{\rm CS}_{\tilde\kappa}(a):=\sup_{\square}\sum_{Y\ni\square}e^{\tilde\kappa(\#Y-1)}\sup_{\mathcal U_Y(x')}|a(Y,\cdot)|,\] where the first supremum runs over the cells $\square$, and we put $|a|^{{\rm cs},(\tilde\kappa)}_{a,x'}:={\rm CS}_{\tilde\kappa}(a)/c_a(x')$.
\end{definition}

The number ${\rm CS}_{\tilde\kappa}(a)$ is the $|a|_{\rm cs}$ of item~1 of \S\ref{Y-H1-1}, taken at $\tilde\kappa$. The quantity $|a|^{{\rm cs},(\tilde\kappa)}_{a,x'}$ dominates the size $\sup_Ye^{\tilde\kappa(\#Y-1)}\sup|a(Y,\cdot)|/c_a(x')$, which is the size in which \Cref{Y-H1-U1}\,(h) measures a direction.

\begin{remark}[Letters with more than one meaning]
The letters $p,p'$ denote fine plaquettes, and $p$ is the pair $p=(\tilde\kappa,\tilde c)$ in $\mathbb P$ only in $M_A(x;p)$, in $p\in\mathbb P$, in ``at $p$'' and ``all $p$'', and in the list of clause~(f) of \Cref{Y-H3-U3}; the remainders are $n[u;D]$, $n^{\rm fix}$, $n^{\rm mem}$, and $n=n[u;D]$ as \Cref{Y-H3-U3} sets it, also in $n(Y,\cdot)$, whereas $n(C)$, told apart from them by its argument, a cluster $C$, is the number of letters in the polymers of $C$, each counted with multiplicity (\Cref{Y-H1-LK}); and outside \textup{(CD)} the letter $r$ without argument and without prime is the number $C_\eta^\flat\varepsilon$ of Step~0\,(vi) of the proof of \Cref{Y-H3-U3}, while $r(x)$ and $r'$ are the radii, and $r(Y)$, as well as $r$ where Step~10 of that proof uses it for the family $(r(Y))$, is the remainder of \Cref{Y-H1-U1}\,\textup{(g3)}. The symbol $T(y',y'')$ denotes the transport, $T>1$ is the parameter of \Cref{Y-H1-LT57}, and $T_s=x^{1/4}$ of \textup{(O13)} is the value of $T$ with which Step~4 of the proof of \Cref{Y-H3-U3} applies that lemma. The letter $u$ is a member; in $u(y')$, $\operatorname{Ad}_u$ and $V'^u$ the letter $u$ is a gauge transformation.
\end{remark}

\emph{Item 2: the data.} We fix $x\ge x_3$ and work in the setting \textup{(SD)} of item~2 of \S\ref{Y-H-1}, with $x_9^U$ in place of $x_9$: thus $N$ is even, $M\mid N$, \textup{(Tor)} holds, the grids are fitted, and the rows $\ell\in\mathcal R(b_*,K_*,\lambda_*,\frac19)$ have range $\le\varrho(x)$ and the four conclusions named there. The empty rows, those with all $\mathsf H_y=0$, are admitted wherever they satisfy \textup{(SD)}. We fix $p\in\mathbb P$ and a member $u\in M_8(x;p)$, which may be complex and carries the weak cap. Finally, $D=(D_y)$, with $y$ running over the fine torus, is a family of complex kernels on the fine plaquettes such that $D_y(p,p')=D_y(p',p)$, and $D_y(p,p')=0$ unless $|y-p|,|y-p'|\le\varrho(x)$. Nothing further is asked of $D$: no sign, no covariance under translations, no factorisation through the rows and no bound on its size.

\emph{Item 3: the objects.}

\begin{definition}[The quadratic member, the step on kernels, the derivative and the remainder]\label{Y-H3-Obj}
For a cube $\boxdot$ and a field $\mathsf W\in\mathcal U_{\{\boxdot\}}(x)$, the quadratic member is \[s_x[D](\{\boxdot\},\mathsf W):=+\frac x2\sum_{y\in\boxdot}\sum_{p,p'}D_y(p,p')\,\bigl(\operatorname{Ad}_{\mathsf W(\Gamma(y,\pi(p)))}F_p\bigr)\cdot\bigl(\operatorname{Ad}_{\mathsf W(\Gamma(y,\pi(p')))}F_{p'}\bigr),\] where $F_p=-i\log\partial\mathsf W(p)$, and $s_x[D](X,\cdot):=0$ for $\#X\ge2$. This family is $-E^\Delta[D]$. At $x'$ the same formula is used, with cells, coarse sites $y'$, coarse plaquettes $P$ and $X'_P=-i\log\partial V'(P)$.

The step on kernels is \[(T_\ell D)_{y'}(P,P'):=L\sum_{y\in B(y')}\sum_{p,p'}H^\sharp(p,P)D_y(p,p')H^\sharp(p',P'),\] where $H^\sharp$ is the free kernel of item E1 of \S\ref{Y-D-1}, taken at the rows $\ell$; its cut is \[D'_{y'}(P,P'):=1\{|y'-P|,|y'-P'|\le\varrho(x')\}\,(T_\ell D)_{y'}(P,P'),\] the distances being measured in coarse units. Kernel and cut are those of \Cref{Y-H2-U2} (see \Cref{Y-H2-Rem3}): item E1 of \S\ref{Y-D-1} defines $H^\sharp$ by $H^\sharp=L^{-1}\mathsf B^{-1}\mathsf T_\sharp \mathsf B'$ in the form of \Cref{Y-A-P4}\,(a), and item~1 of \S\ref{Y-H2-1} takes the kernel from \S\ref{Y-A-1} in the form of that proposition.

We write $DA_\ell(u)[\delta](Y,V')$ for the derivative $DA(u)[\delta](Y,V')$ that item~1 of \S\ref{Y-H1-1} defines as a sum, term by term, of derivatives, here taken at the rows $\ell$. For $u\in M_8(x;p)$, $\delta\in T^a(x)$ and $V'\in\mathcal U_Y(x')$ it equals $\frac{d}{dz}E'^{\rm old}[u+z(\delta,0,0)](Y,V')$ at $z=0$, by \Cref{Y-H1-U1}\,(h), as $A=8<16$. By \textup{(Fm1)}, $E'^a$ depends on $E^a$ alone and $E'^{\rm new}$ on no member; hence $DA_\ell(u)[\delta]$ is the $a$-part of $DF_\ell(u)$, the derivative of the step at the member $u$, in every direction whose $a$-part is $\delta$.

The remainder is $n[u;D]:=DA_\ell(u)[s_x[D]]-s_{x'}[D']$.
\end{definition}

\emph{Item 4: the theorem.}

\begin{lemma}[Quadratic source members]\label{Y-H3-I4}
Suppose $\varepsilon(x)\le1/40$, $\varrho(x)+2\le Mr(x)$, and the torus has side at least $16Mr(x)$. Let $D=(D_y)$ be a complex family with $D_y(p,p')=0$ unless both $|y-p|$ and $|y-p'|$ are at most $\varrho(x)$. Then $s_x[D]\in T^a(x)$ and
\[
\sup_{\mathcal U_{\{\square\}}(x)}
  |s_x[D](\{\square\},\cdot)|
  \le 306M^3\mathsf g_\flat^2p_0(x)^2\|D\|_{(\lambda_\flat)}
  =\sigma_D .
\]
For real $D$ the member is real on $G$-valued fields.
\end{lemma}

\begin{proof}
For $y\in\square$ and $|y-p|\le\varrho(x)$, the plaquette $p$ and the path $\Gamma(y,\pi(p))$ lie in the embedded box $\square^\approx$. The path has length at most $3\varrho(x)\le\ell_w(x)$, as in \textup{(G2)} and \textup{(G3)} of Step~0 in \S\ref{Y-I-1}. On $\mathcal U_{\{\square\}}(x)$, \textup{(u2)} bounds the norm of the path transport by $2$. Also, \textup{(u3)} and $\varepsilon(x)\le1/40$ give
\[
  |F_p|\le -\log(1-2\varepsilon(x))\le2.06\varepsilon(x).
\]
The norm of each transported $F_p$ is thus at most $8.24\varepsilon(x)$. Since $\sum_{p,p'}|D_y(p,p')|\le9\mathsf g_\flat^2\|D\|_{(\lambda_\flat)}$, the size bound follows from $x\varepsilon(x)^2=p_0(x)^2$ and $(9/2)(8.24)^2\le306$.

The logarithm defining $F_p$ is holomorphic by the same plaquette bound. A gauge transformation conjugates both transported vectors at $y$, preserving their bilinear product. At a pure gauge every $F_p$ vanishes. For $G$-valued fields the transported vectors belong to $\mathfrak g$ and their bilinear product is real. These properties and the fact that $s_x[D]$ vanishes on sets of two or more cubes give the asserted membership.
\end{proof}

\begin{theorem}[The step differentiated in a quadratic direction]\label{Y-H3-U3}
In the setting of item~2, the remainder $n=n[u;D]$ of \Cref{Y-H3-Obj} has the following properties.
\begin{enumerate}
\item[\textup{(a)}] \emph{Bound.} The remainder splits as $n=n^{\rm fix}[D]+n^{\rm mem}[u;D]$, the two parts being defined in Step~1 of the proof. The part $n^{\rm fix}$ depends on no member and vanishes unless $\#Y=1$, and for every cell $\square$ \[\sup_{\mathcal U_{\{\square\}}(x')}|n^{\rm fix}(\{\square\},\cdot)|\le\sigma_D\bigl[L^3(2k_1\varepsilon+k_2\varepsilon')+k_3L^3p_0^{-2}x^{-1/2}+2x^{-10}\bigr].\] The other part satisfies \[{\rm CS}_{\tilde\kappa}(n^{\rm mem})\le\sigma_D\bigl[1000e^{\tilde\kappa+3}L^3x^{-1/4}+18.2L^3c_a(x)\bigr]=:B_{\rm mem}(\tilde\kappa).\] In $n^{\rm mem}$, the clusters with $n(C)\ge2$ make up the family ${\rm Rst}$ of \Cref{Y-H1-U1}\,(h); by \textup{(h1)} applied with $\delta=s_x[D]$, whose size $B$ satisfies $Bc_a(x)\le\sigma_D$, they are bounded by $217L^3x^{-1/4}\sigma_D$, and $217\le1000e^{\tilde\kappa+3}$. Consequently ${\rm CS}_{\tilde\kappa}(n)\le1001e^{\tilde\kappa+3}L^3x^{-1/4}\sigma_D$ and $|n[u;D]|^{{\rm cs},(\tilde\kappa)}_{a,x'}\le u_3(x')\|D\|_{(\lambda_\flat)}$, for every $p=(\tilde\kappa,\tilde c)\in\mathbb P$ and every $u\in M_8(x;p)$, each at its own $\tilde\kappa$.
\item[\textup{(b)}] The functions $n(Y,\cdot)$, $n^{\rm fix}(Y,\cdot)$ and $n^{\rm mem}(Y,\cdot)$ have the clauses of \textup{(ca)} without a size at $x'$: they are holomorphic on $\mathcal U_Y(x')$, invariant under $G$-valued gauge transformations, $0$ at pure gauges, and functions of the field on $Y^\approx$. They are real at $G$-valued fields if $u$ and $D$ are real.
\item[\textup{(c)}] \emph{Locality.} The function $n(Y,\cdot)$ depends on $D_y$ only at the fine sites $y$ under $Y$, and it vanishes unless some $D_y$ with $y$ under $Y$ is not $0$. The kernel $D'_{y'}$ depends only on the $D_y$ with $y\in B(y')$.
\item[\textup{(d)}] \emph{In the member.} The remainder $n$ depends on the member $u$ through $E^a$ alone, and each value $n[u;D](Y,V')$ is holomorphic along every affine line in $M_8(x;p)$. For $u_1,u_2\in M_2(x;p)$, \[{\rm CS}_{\tilde\kappa}(n[u_1;D]-n[u_2;D])\le\frac13B_{\rm mem}(\tilde\kappa)\,d_a,\] where $d_a:=\sup_X\sup_{\mathcal U_X(x)}|E_1^a-E_2^a|e^{\tilde\kappa(\#X-1)}/c_a(x)$ is the $a$-part of $\mathsf D^{\rm in}_x(u_1,u_2)$.
\item[\textup{(e)}] The map $D\mapsto(s_x[D],D',n^{\rm fix}[D],n^{\rm mem}[u;D])$ is complex linear. Hence \textup{(a)} and \textup{(d)} hold for $\tilde D-D$, with $\|\tilde D-D\|_{(\lambda_\flat)}$ in place of $\|D\|_{(\lambda_\flat)}$.
\item[\textup{(f)}] The constants $k_1,k_2,k_3,C_3,x_3$ and all numerals are independent of $N,K,k,\mathsf b,p,u,D$; of the rows $\ell$ they involve only the pure numbers $b_*,K_*,\lambda_*$ and the range.
\item[\textup{(g)}] \emph{The number.} The derivative of the number in the direction $s_x[D]$ is \[D\mathsf V'(u)[s_x[D]]=\sum_y\frac32\sum_{p,p'}D_y(p,p')(\mathrm{d}C^1\mathrm{d}^{\mathsf T})(p,p')+\sum_{\square}e_\square,\] the last sum running over the cells, with $|e_\square|\le\sigma_D\bigl[L^3(k_1\varepsilon+9.1c_a(x))+109L^3x^{-1/4}+x^{-10}\bigr]$; here $C^1$ is $C$ at $V'=1$, and its blocks are scalars.
\end{enumerate}
\end{theorem}

\emph{Item 5: the clause \textup{(2$\tau$)}.} The clause below is a corollary of \Cref{Y-H1-LT57}, which is stated in Step~S10\,(3) of the proof of \Cref{Y-H1-U1} and proved there by Cordes' inequality; Step~4 of the proof of \Cref{Y-H3-U3} proves \textup{(2$\tau$)}.

\begin{corollary}[\textup{(2$\tau$)}: the variation of a dressed trace of the covariance]\label{Y-H3-2tau}
In the setting of item~2, let $y$ be a fine site and $y'=y(y)$, let $p,p'$ be fine plaquettes with $|y-p|,|y-p'|\le\varrho(x)$, let $o$ be a coarse site within $M$ of $y'$, and let $\mathsf B_\omega$ be the box of radius $\omega$ centred at $o$. Put \[\tau(V'):=\operatorname{tr}\bigl[T(y',y(p))\,(\mathrm{d}^VC\mathrm{d}^{V\mathsf T})(p,p')\,T(y',y(p'))^{\mathsf T}\bigr].\] For every $T>1$ at which \textup{(C0)}, \textup{(C1)} and \textup{(C2)} of \textup{(Y.0)(0d)} (item~5 of \S\ref{Y-H-1}) hold, the function $\tau$ is holomorphic and $G$-invariant on $U^{17T}(\mathsf B_\omega;x')$, depends on $V'$ only within $R_2+\varrho_L$ of $y'$, and satisfies $|\tau|\le3134/\gamma$ there; moreover $\tau(1)=3(\mathrm{d}C^1\mathrm{d}^{\mathsf T})(p,p')$. At the fields on $\mathsf B_\omega$ with \textup{(u1)} to \textup{(u3)} at $x'$, \[|\tau(V')-\tau(1)|\le2(3134/\gamma)T^{-2}.\] At $T=T_s$ this bound becomes $6268\gamma^{-1}x^{-1/2}$.
\end{corollary}

\emph{Item 6: the corollary.} Besides \Cref{Y-H3-U3}, the next corollary rests on \Cref{Y-H1-U1}\,(h3) and on \textup{(L3su)}.

\begin{corollary}[A direction with a remainder, and the recursion of the remainders]\label{Y-H3-R3}
In the setting of item~2, let ${\rm rem}\in T^a(x)$ be of size $B_r$ caps at $p$. Then \[DA_\ell(u)[s_x[D]+{\rm rem}]=s_{x'}[D']+{\rm rem}',\] where ${\rm rem}':=DA_\ell(u)[{\rm rem}]+n[u;D]\in T^a(x')$, and $|{\rm rem}'|^{{\rm cs},(\tilde\kappa)}_{a,x'}\le\vartheta_aB_r+u_3(x')\|D\|_{(\lambda_\flat)}$. Let $\bar d>0$ with $\|D\|\le\bar d$; in the units $t:=B_r/(u_3(x)\bar d)$ and $t':=|{\rm rem}'|/(u_3(x')\bar d)$ one has $t'\le0.2431t+1$. Along levels $x_{k+1}=x_k/L$ with $\|D_k\|\le\bar d$, it follows that $t_k\le\max\{t_0,1.3212\}$, which is at most $2$ if $t_0\le2$.
\end{corollary}

\emph{Item 7: four remarks.}

\begin{remark}[$n^{\rm fix}$ in caps]
In caps at $x'$, \[|n^{\rm fix}|_{a,x'}\le(c_Q^{\rm su}/E_a)(L^{5/2}/0.5946)(2k_1+L^{1/2}k_2+k_3+1)\|D\|_{(\lambda_\flat)},\] a bound that carries no power of $x$ and no logarithm. The power $x'^{1/4}$ of $u_3$ comes from ${\rm Rst}$ alone, and $n^{\rm fix}$ cancels in the difference of two members.
\end{remark}

\begin{remark}[The identity over sub-alphabets]
Let $X$ be a cube under $Y$, and let $s_X$ be the member equal to $s_x[D]$ at the cube $X$ and to $0$ at all other sets of cubes. Let $\mathbb A$ be a finite set of letters \textup{(a)}, \textup{(v)}, \textup{(t)} with cores in $Y$ which contains the letter \textup{(a)} at $X$, and let $V'\in\mathcal U_Y(x')$. With $W_{\mathbb A}:=\sum\chi_XE^a(X,U)+\sum\chi_\square V_\square+\sum T_Y$, the sums running over $\mathbb A$, we have \[\sum_{C\ {\rm in}\ \mathbb A}\frac{d}{dz}\big|_{z=0}\Phi(C)[u+zs_X]=E_\xi[\chi_Xs(X,U)e^{W_{\mathbb A}}]/E_\xi[e^{W_{\mathbb A}}].\] The cluster of the single letter \textup{(a)}$_X$ yields $E_\xi[\chi_Xe^{E^a}s]$, and the clusters with $n(C)\ge2$ yield the truncated part. Thus $n$ is, polymer by polymer, the remainder of the family itself, and not a new decomposition.
\end{remark}

\begin{remark}[The sum rule]
Let $u$ be a real member of $M_1$ with the strong cap and with $\rho[u;\ell]>0$, and let $D$ be real. Put \[q_D(V):=\frac x2\sum_ym_y(V)\sum D_y(p,p')\langle T(y,\pi(p))\zeta_p,T(y,\pi(p'))\zeta_{p'}\rangle,\] and let $u_t:=(E^a+ts_x[D],E^h,w_t)$, where \[w_t(X,V):=w(X,V)\prod\{e^{tq_\boxdot(V)}:\ \boxdot\ \text{off at}\ V,\ p_*(\boxdot,V)\text{ in }X\}.\] Here $q_\boxdot$, with the sign $+$, is the part of $q_D$ with $y\in\boxdot$; a cube $\boxdot$ off at $V$ and $p_*(\boxdot,V)$ are as in Step~5 of the proof of \Cref{Y-I-LB}; $m_y$ and $\zeta_p$ are the mask and the clip of item~2 of \S\ref{Y-I-1}; and $q'_{D'}$ is $q_D$ taken at $x'$, with $D'$ in place of $D$. Then $\rho[u_t;\ell]=\rho[u;\ell]e^{tq_D}$ at every $V$, and $\frac{d}{dt}\log(\rho[u_t;\ell])'(V')$ at $0$ is the mean of $q_D(U)$ given $V'$. At a $G$-valued $V'$ without a rough plaquette, this derivative equals \[D\mathsf V'(u)[s_x[D]]+q'_{D'}(V')+\sum_Y\bigl(n[u;D]+DE'^h(u)[du_t/dt]\bigr)(Y,V').\]
\end{remark}

\begin{remark}[The size of $D'$]
A bound on $\|D'\|_{(\lambda_\flat)}$ is the subject of \Cref{Y-H2-U2}.
\end{remark}

\emph{Item 8: what this section rests on.}

\begin{remark}[The results on which this section rests]
This section rests on \S\ref{Y-H0-1}, on \Cref{Y-H1-U1}\,(U) and on \textup{(c+)}, \textup{(h)}, \textup{(h1)}, \textup{(g1)}, \textup{(g3)} and \textup{(h3)} of that theorem, on Lems.~\ref{Y-H1-LT57} and~\ref{Y-H1-LK}, on \Cref{Y-B-RZ}, and on Sections~\ref{Y-C-1}, \ref{Y-D-1}, \ref{Y-E-1}, \ref{Y-G0-1}, \ref{Y-G-1}, \ref{Y-H-1} and~\ref{Y-I-1}.
\end{remark}

\begin{proof}[Proof of \Cref{Y-H3-U3}]
\emph{Step 0: the floor $x_3$, the order of the constants, and the inequalities that all later steps use.} The floor $x_3$ exists, and each new letter depends only on letters introduced before it. For $x\ge x_3$ we have the following.
\begin{enumerate}
\item[\textup{(i)}] $p_0\ge p_0'\ge1$.
\item[\textup{(ii)}] $\varepsilon<\varepsilon'\le L^{1/2}\varepsilon$, and $\varepsilon'\le1/40$.
\item[\textup{(iii)}] $c_a(x)\le c_a(x')/(0.5946L^{1/2})$.
\item[\textup{(iv)}] $\theta_D/L\le\lambda_\flat/9$.
\item[\textup{(v)}] $8c_a(x)\le1/4$.
\item[\textup{(vi)}] With $r:=C_\eta^\flat\varepsilon$, one has $r\ge7\varepsilon'$, $r\le1/(22L)$ and $17.1r\varrho(x)\le0.1$.
\item[\textup{(vii)}] $\omega\le Mr'$, $3\omega\le\ell_w(x')$ and $\omega\ge M+R_2+\varrho_L$, and $\mathsf B_\omega$ is embedded.
\end{enumerate}

\emph{Order.} The numbers $\mathsf g_\flat,g_{1b},g_b^-$ depend on $\lambda_*$ of \textup{(O0)}. The number $k_1$ depends on $g_{1b}/\mathsf g_\flat$ of \textup{(O0a)} and on $C_\eta^\flat$, that is, on $\mathsf a_{\rm c}$ of \textup{(O10)} and on $L$. The number $k_2$ depends on $g_b^-/\mathsf g_\flat$ of \textup{(O0a)}, on $L$, on $C_D$ and $\theta_D$ of \textup{(O6)}, and on $\mathsf K_{\rm c}$ of \textup{(O10)}; and $k_3=6730L^4$ depends on $L$. None of these three numbers depends on $M,\kappa,c_h,E_a,x$; they are introduced at \textup{(O10a)}. The constant $C_3$ depends on $\kappa$, $L$ and $M$, on $\mathsf g_\flat$ of \textup{(O0a)} (through $c_Q^{\rm su}=306M^3\mathsf g_\flat^2$ of \textup{(O10)}), and on $E_a$. The floors of this section are free of $C_3$; the last floor $x_s$ depends on it, by way of $\check\omega_3$ in \textup{(W3)} of item~(1) of \S\ref{Y-H0-1}, and since $C_3$ is fixed at \textup{(O12a)}, ahead of \textup{(O13)}, the order contains no circle. Finally, $x_3$ depends on $x_9^U$ and on \textup{(O0)} to \textup{(O12a)}.

\emph{Existence.} Each of the conditions \textup{(Z.0)} to \textup{(Z.4)} is one-sided and holds for all large $x$. In \textup{(Z.0)}, $T_s\varepsilon'=L^{1/2}p_0'x^{-1/4}$, while $\omega$ and $\ell_w(x')$ are $O(\log x)$ and $1/\mathsf g_x=4(\log x)^2/(\vartheta_0^\flat\log L)$; so each of its three inequalities, divided by its right side, has on the left at most a constant multiple of a fixed power of $\log x$, times $x^{-1/4}$. \textup{(Z.1)} holds since $p_0$ grows. In \textup{(Z.2)}, $n_\chi=O((\log x)^3)$ and $p_0^2=A_0^2(\log x)^{2\mathfrak p_0}$ with $2\mathfrak p_0>4$, so that, as $x$ grows, $e^{-c_\xi^2p_0^2/2}$ tends to $0$ more quickly than any negative power of $x$. \textup{(Z.3)} amounts to $x^{-11}\le(17.9L^4C_D^2)^{-1}$. In \textup{(Z.4)} the left side is at most a constant multiple of a fixed power of $\log x$, times $x^{-1/2}$. Thus $S_3$ is a non-empty upper half line and $x_3=1+\inf S_3$ belongs to it. No condition depends on $N$, $K$, $\mathsf b$, a member, a row or $D$.

\emph{The least floor.} The factors $\omega$, $\ell_w(x')$ and $n_\chi$ are integer valued, do not decrease in $x$, and appear as factors on the left sides; all other factors, and $\mathsf g_x$, are continuous and positive. Suppose that \textup{(Z.0)} to \textup{(Z.4)} hold for all $x>t$. If, at such $x$, we give these three factors their values at $t$, the left sides become no larger; letting $x$ tend to $t$ from above, we obtain the conditions at $t$. Hence the set of the thresholds is a closed half line, and it has a minimum.

\emph{The inequalities.} Since $x\ge x_3\ge x_9^U$, the floors \textup{(Y.0)} to \textup{(Y.9)} of item~5 of \S\ref{Y-H-1} hold, and with them, by \textup{(Y.0)(0a)}, the floors of Sections~\ref{Y-D-1} and~\ref{Y-E-1}. In \textup{(Y.0)(0b)}, and in \textup{(0d)} as \S\ref{Y-H-1} states it, $T$ is the number $L^{3/2}/(19\mathsf N_{\rm c})>1$ of Step~0\,(iv) of the proof of \Cref{Y-H-C9}. \textup{(0b)} contains $\mathsf K_{\rm c}(17T\varepsilon')^2\le0.12\varepsilon$, and since $17T\ge1$ this gives $\mathsf K_{\rm c}\varepsilon'^2\le0.12\varepsilon$.
\begin{enumerate}
\item[\textup{(i)}] By \textup{(F8$_k$)} of item G of \S\ref{Y-D-1}, $p_0'\ge1$; and $p_0$ increases.
\item[\textup{(ii)}] $\varepsilon'/\varepsilon=L^{1/2}p_0'/p_0\le L^{1/2}$. By \textup{(0b)}, $(p_0'/p_0)^4\ge1/2$, hence $\varepsilon/\varepsilon'\le2^{1/4}L^{-1/2}<1$; and $\varepsilon'\le1/40$ by \textup{(F8$_k$)}.
\item[\textup{(iii)}] $c_a(x')/c_a(x)=L^{1/2}(p_0'/p_0)^3\ge2^{-3/4}L^{1/2}$, where $2^{-3/4}\ge0.5946$.
\item[\textup{(iv)}] By \textup{(O6)} in item A of \S\ref{Y-D-1}, $\theta_D\le\theta_{\mathbb C}/2=\theta^\flat/12\le\theta_1/12\le\lambda_*L/144$, and $\lambda_*/144=\lambda_\flat/9$.
\item[\textup{(v)}] This is \textup{(Y.1)}.
\item[\textup{(vi)}] The bound $r\ge7\varepsilon'$ follows from $C_\eta^\flat\ge7L^{1/2}$ and \textup{(ii)}; the other two bounds are \textup{(W1)} and \textup{(W2)} of item~6 of \S\ref{Y-E-1} at $\tau=1$.
\item[\textup{(vii)}] The bound $\omega\le Mr'$ is the second inequality of \textup{(Y.7)}, together with $\lceil Mr(x)/L\rceil\le M\lceil r(x)/L\rceil$ and $M+2\le2M$; then $3\omega\le\ell_w(x')$, because $3Mr'\le\ell_w(x')$. Next, $\omega\ge M+R_3+2=M+2R+\varrho_L+7$. Finally, $2\omega+1<N$ by \textup{(Tor)} (the clauses on $N$ in Step~0 of the proof of \Cref{Y-G-ThA}), so $\mathsf B_\omega$ is embedded.
\end{enumerate}

\emph{Step 1: the member $s_x[D]$, the derivative, and the split into five terms.} We prove the following three facts.
\begin{enumerate}
\item[\textup{(1)}] $s_x[D]\in T^a(x)$ and $\sup_{\mathcal U_{\{\boxdot\}}(x)}|s_x[D](\{\boxdot\},\cdot)|\le\sigma_D$; moreover $s_{x'}[D']\in T^a(x')$.
\item[\textup{(2)}] For $V'\in\mathcal U_Y(x')$, \[DA_\ell(u)[s_x[D]](Y,V')=1\{Y=\{\square_0\}\}\sum_{X\ {\rm cube\ under}\ \square_0}[g_X(V')-g_X(1)]+{\rm Rst}(Y,V'),\] where $g_X:=E_\xi[\chi_Xe^{E^a(X,U)}s(X,U)]$ and $s:=s_x[D]$.
\item[\textup{(3)}] With the objects defined below, \[g_X=M_X+\frac12{\rm Tr}_X+E_\xi[\chi_XR_X]-E_\xi[(1-\chi_X)Q_X]+E_\xi[\chi_X(e^{E^a(X,U)}-1)s(X,U)].\]
\end{enumerate}

\emph{Proof of \textup{(1)}.} We apply \Cref{Y-H3-I4} to $D$ at $x$ and to $D'$ at $x'$. The hypotheses \textup{(F1)} and \textup{(F2)} hold at $x$ by the claim \textup{(Fl)} of Step~0 in \S\ref{Y-I-1}, and at $x'$ by $\varepsilon'\le1/40$ and $\varrho(x')+2\le c_R\log x'+3\le20c_R\log x'\le Mr'$. \textup{(Tor)} gives $N\ge16Mr(x)$ and hence the required lower bounds on the fine side $NL$ and on the coarse side $N\ge16Mr(x')$. On the coarse torus the coordinate paths are defined and $\square^\approx$ embeds, by \textup{(0.3)} of Step~0 of the proof of \Cref{Y-D-C5}. The family $D'$ is complex, symmetric because $D$ is, and of range $\le\varrho(x')$; so the same conclusions hold at $x'$.

\emph{Proof of \textup{(2)}.} If $\sigma_D=0$, everything vanishes. Otherwise $u+z(s,0,0)\in M_{16}(x;p)$ for $|z|\le z_0:=8c_a(x)/\sigma_D$, so by \Cref{Y-H1-U1}\,(U) and \textup{(c+)} the derivative exists and is one and the same object for all $p$. \Cref{Y-H1-U1}\,(h), applied with $\delta=s$, gives the formula with the sum over the $X$ with $Y_0(X)=Y$, and $s(X)=0$ for $\#X\ge2$. Every single cube is small, as $1\le L/2-2$; by Step~8, \textup{(8.2)}, of the proof of \Cref{Y-H-C9}, for such a cube $X$ the equality $Y_0(X)=\{\square_0\}$ then holds, where $\square_0$ is the cell that contains the cube in its fine set. As the grids are fitted (``The grids fit'' in item~1 of \S\ref{Y-G0-1}, with the proof given there), the fine set of $\square_0$ is the union of $L^3$ cubes, and it is also the union of the blocks $B(y')$, $y'\in\square_0$.

\emph{Proof of \textup{(3)}.} Fix a cube $X$ under $\square_0$ and $V'\in\mathcal U_{\{\square_0\}}(x')$, extended by $1$. Put \[Q_X(\xi,V'):=\frac x2\sum_{y\in X}\sum_{p,p'}D_y(p,p')\,\bigl(T(y(y),y(p))F^\flat_p\bigr)\cdot\bigl(T(y(y),y(p'))F^\flat_{p'}\bigr),\] with $F^\flat=\bar F+\mathrm{d}^VG\xi$. This is a polynomial of degree $2$ in the real variable $\xi$, defined for all $\xi$, and its coefficients are holomorphic in $V'$ by \textup{(2d)} of \textup{($\Gamma$2)} in item~4 of \S\ref{Y-H-1}, since $X'_P=-i\log\partial V'(P)$, the logarithm being given by its power series (item C6 of \S\ref{Y-D-1}). On $\chi_X=1$ put $R_X:=s(X,U)-Q_X$. Then $\chi e^{E^a}s=\chi s+\chi(e^{E^a}-1)s$ and $\chi s=Q-(1-\chi)Q+\chi R$. As $E_\xi\xi=0$, the cross term has mean $0$; moreover $E_\xi(G\xi)(G\xi)^{\mathsf T}=x^{-1}GG^{\mathsf T}=C/x$, the transpose being that of the bilinear form, also at complex $V'$; and, written in an orthonormal basis of $\mathfrak g$, the form $Y\cdot Y'$ becomes the standard bilinear form of $\mathbb C^3$. Therefore $E_\xi Q_X=M_X+\frac12{\rm Tr}_X$, with \[M_X:=\frac x2\sum_{y\in X}\sum D_y(p,p')\,(T\bar F_p)\cdot(T\bar F_{p'})\] and \[{\rm Tr}_X:=\sum_{y\in X}\sum D_y(p,p')\,\operatorname{tr}\bigl[T(y(y),y(p))\,(\mathrm{d}^VC\mathrm{d}^{V\mathsf T})(p,p')\,T(y(y),y(p'))^{\mathsf T}\bigr].\] At $V'=1$ we have $X'=0$, $a_0=0$ and $\bar F=0$, so that $M_X(1)=0$.

\emph{The two parts of the remainder.} We define \[\FitDisplay{\dimexpr\linewidth-1pt\relax}{n^{\rm fix}(\{\square_0\},V'):=\sum_X\bigl\{E_\xi[\chi R_X]-E_\xi[(1-\chi)Q_X]+\frac12{\rm Tr}_X\bigr\}(V')-\text{the same at }1+\bigl[\sum_XM_X(V')-s_{x'}[D'](\{\square_0\},V')\bigr],}\] and $n^{\rm fix}(Y):=0$ for $\#Y\ge2$; further \[\FitDisplay{\dimexpr\linewidth-1pt\relax}{n^{\rm mem}(Y,V'):=1\{Y=\{\square_0\}\}\sum_X\bigl\{E_\xi[\chi_X(e^{E^a}-1)s](V')-\text{the same at }1\bigr\}+{\rm Rst}(Y,V').}\] The part $n^{\rm fix}$ contains no member.

\emph{Step 2: the member $s(X,U)$ against $Q_X$ where $\chi_X=1$, at complex $V'$.}

\begin{lemma}[Comparison of the member at the chart with its quadratic form at the centre]\label{Y-H3-LV}
In the setting of item~2 and with the objects of Step~1, let $V'\in\mathcal U_{\{\square_0\}}(x')$, let $X$ be a cube under $\square_0$, and let $\xi$ be real with $\chi_X(\xi)=1$. Then \[|R_X|\le289.92M^3\mathsf g_\flat(1320g_{1b}+216\mathsf g_\flat)\|D\|xr^3\le k_1\varepsilon\sigma_D.\] The part of $n^{\rm fix}$ consisting of the terms $E_\xi[\chi R_X]$ is at most $2L^3k_1\varepsilon\sigma_D$ in modulus.
\end{lemma}

\begin{proof}[Proof of \Cref{Y-H3-LV}]
\textup{(i)} \emph{The size of $\eta$.} Let $y\in X$ and $|y-p|\le\varrho(x)$. Then, by \textup{(G2)} of Step~0 in \S\ref{Y-I-1}, the four bonds of $p$, as well as every bond of $\Gamma(y,\pi(p))$, are bonds of $X^\approx$. For each such bond $b$, its window $Y_b$ lies within $r_w$ of $b$ by \textup{(2a)}; hence $\chi_X$ tests that window, and \textup{(2e)} gives $\|(G\xi)_b\|\le C_\varphi^\flat c_\xi\varepsilon=\varepsilon/6.5$. Such a bond $b=(z,\mu)$ is a bond of a plaquette based at $z$, and the block of $z$ meets $X^\approx$. By the last sentence of Step~2\,(ii) of the proof of \Cref{Y-H-C9}, the boxes of radius $R_3$ centred at the sites whose blocks meet $X^\approx$ lie in $\square_0^\approx$, where $V'$ satisfies \textup{(e1)}, \textup{(e2)} and $|\partial V'(P)-1|<2\varepsilon'$. So \textup{($\Gamma$3)} of item~4 of \S\ref{Y-H-1}, applied with $\hat\varepsilon=\varepsilon'$, gives $\|a_{0,b}\|\le|a_{0,b}|\le \mathsf a_{\rm c}\varepsilon'\le \mathsf a_{\rm c}L^{1/2}\varepsilon$. Since $\eta=0$ on the combs and $\varepsilon/6.5+\mathsf a_{\rm c}L^{1/2}\varepsilon\le C_\eta^\flat\varepsilon$, we obtain $\|\eta_b\|\le r$ on all bonds that occur.

\textup{(ii)} \emph{Steps~4 and~5 of the proof of \Cref{Y-E-C6}.} In clause \textup{(CD)} below, \textup{(C)} is Step~4 and \textup{(D)} is Step~5 of that proof, each stated under the hypotheses on which its proof rests.

\begin{enumerate}
\item[\textup{(CD)}]\label{Y-H3-CD} Let $0<r\le1/(22L)$, let $V'$ be an $\operatorname{SL}(2,\mathbb C)$ field with $|V'(c)|\le1+10^{-4}$, let $\eta$ be a $\mathfrak g_{\mathbb C}$-valued function on $I^\circ$, with $\eta=0$ on the combs, and let $U=\mathsf U(\eta;V')$.
\begin{enumerate}
\item[\textup{(C)}] If $p$ is a fine plaquette with $\|\eta_b\|\le r$ on its four bonds and, in case $p$ is an edge plaquette under $P$, $\partial V'(P)=e^{iX'_P}$ with $|X'_P|\le(2.06/7)r$, then $|F^\flat_p|\le6r$, $|F_p-F^\flat_p|\le54r^2$ and $|F_p|\le8.25r$.
\item[\textup{(D)}] If $y$ is a fine site with $n:=|y-p|\le\varrho(x)$ and $n+L<NL/2$, if $17.1r\varrho(x)\le0.1$, if $\|\eta_b\|\le r$ on the bonds of $\Gamma(y,\pi(p))$, and if $|V'(\pi'_j)|<2$ on the first pieces $\pi'_j$ of $\Gamma(y(y),y(p))$, then $\|\operatorname{Ad}_{U(\Gamma(y,\pi(p)))}-T(y(y),y(p))\|\le160rn$, and the norm of $\operatorname{Ad}_{U(\Gamma(y,\pi(p)))}$ is at most $4.9$, that of $T(y(y),y(p))$ less than $4$.
\end{enumerate}
\end{enumerate}

\begin{proof}[Proof of \textup{(CD)}]
The proof is that of Steps~4 and~5 of the proof of \Cref{Y-E-C6}, word for word, together with \Cref{Y-E-L1C} (item~3.2 of \S\ref{Y-E-1}), which these two steps cite. The hypotheses of \textup{(CD)} take the place of what the two steps invoke, namely Steps~2 and~3 of that proof, $r\ge7\hat\varepsilon$, \textup{(W1$_\tau$)} and \textup{(W2$_\tau$)}; in particular $r\le1/(22L)$ takes the place of \textup{(W1$_\tau$)}, and $17.1r\varrho(x)\le0.1$ that of \textup{(W2$_\tau$)}. Beyond this, the two steps depend neither on $\chi_\square=1$ nor on $y$ being under $\square$, and they depend on no row.
\end{proof}

The two steps are stated under $\chi_\square=1$ for $y$ under $\square$; the hypotheses of \textup{(CD)} hold here. To begin with, $|V'(c)|<1+\vartheta(x')<1+\vartheta_0^\flat$ by \textup{(u1)} at $x'$, and $\vartheta_0^\flat=\theta^\flat/24\le10^{-4}$ by \textup{(O6)} in item A of \S\ref{Y-D-1}. If $p$ is an edge plaquette under $P$, then $y(p)=y_P$; the block of $y(p)$ meets $X^\approx$, because it contains $\pi(p)$ by \textup{(G2)}; and $P$ is a plaquette of the box of radius $R_3$ centred at $y(p)$, which lies in $\square_0^\approx$ by \textup{(i)}. Hence $|\partial V'(P)-1|<2\varepsilon'\le1/20$ by \textup{(u3)}, the logarithm in $X'_P=-i\log\partial V'(P)$ is given by its power series, $e^{iX'_P}=\partial V'(P)$, and $|X'_P|\le2.06\varepsilon'$ by item C6 of \S\ref{Y-D-1}. As for the paths, $3\varrho_L\le\ell_w(x')$ by \textup{(0.2)} of Step~0 of the proof of \Cref{Y-D-C5}, and $\varrho_L\le Mr'$ by \textup{(G1$_k$)} of item G of \S\ref{Y-D-1}. Thus $|V'(\pi'_j)|<2$ on the first pieces of $\Gamma(y(y),y(p))$: these are coarse paths of at most $3\varrho_L\le\ell_w(x')$ bonds within $\varrho_L$ of $y(y)$ in $\square_0$, so they lie in $\square_0^\approx$, where \textup{(u2)} holds. The bound $\|\eta_b\|\le r$ is \textup{(i)}; the inequalities $7\varepsilon'\le r$, $r\le1/(22L)$ and $17.1r\varrho(x)\le0.1$ are Step~0\,(vi), so that $|X'_P|\le2.06\varepsilon'\le(2.06/7)r$; and $|y-p|+L<NL/2$ follows from \textup{(Tor)} with \textup{(F2)} of Step~1\,(1). The two steps, that is \textup{(C)} and \textup{(D)} of \textup{(CD)}, therefore give $|F^\flat_p|\le6r$, $|F_p-F^\flat_p|\le54r^2$, $|F_p|\le8.25r$ and $\|\operatorname{Ad}_{U(\Gamma(y,\pi(p)))}-T(y(y),y(p))\|\le160r|y-p|$, the norm of the first of these two operators being at most $4.9$ and that of the second less than $4$.

\textup{(iii)} \emph{The sum.} Put $A_p:=\operatorname{Ad}_{U(\Gamma(y,\pi(p)))}F_p$ and $B_p:=T(y(y),y(p))F^\flat_p$. Then $A_p-B_p=(\operatorname{Ad}_U-T)F_p+T(F_p-F^\flat_p)$, so that $\|A_p-B_p\|\le(160\cdot8.25|y-p|+4\cdot54)r^2=(1320|y-p|+216)r^2$, while $\|A_p\|+\|B_p\|\le(4.9\cdot8.25+4\cdot6)r=64.425r$. Since $D_y$ and the form are symmetric, the cross terms cancel, and \[\sum D_y(p,p')(A_p\cdot A_{p'}-B_p\cdot B_{p'})=\sum D_y(p,p')(A_p-B_p)\cdot(A_{p'}+B_{p'}).\] Each base site carries three plaquettes, so $\sum_pe^{-\lambda_\flat|y-p|}|y-p|\le3g_{1b}$ and $\sum_pe^{-\lambda_\flat|y-p|}\le3\mathsf g_\flat$. With $|Y\cdot Y'|\le\|Y\|\|Y'\|$, the modulus of the sum is at most $\|D\|\cdot3(1320g_{1b}+216\mathsf g_\flat)r^2\cdot3\mathsf g_\flat\cdot64.425r$. Multiplication by $\frac x2M^3$ brings in the factor $\frac12\cdot9\cdot64.425=289.9125$. Now $xr^3=(C_\eta^\flat)^3p_0^2\varepsilon$, and division by $\sigma_D=306M^3\mathsf g_\flat^2p_0^2\|D\|$ gives $289.9125/306\le0.94743\le0.9475$, which is the bound $k_1\varepsilon\sigma_D$. The weight $|y-p|$ is absorbed by $e^{-\lambda_\flat|y-p|}$, so that no logarithm appears. Finally, $|E_\xi[\chi_XR_X]|$ is at most the supremum; the subtraction at $1$, a field that lies in $\mathcal U_{\{\square_0\}}(x')$, doubles it; and there are $L^3$ cubes. This gives the bound $2L^3k_1\varepsilon\sigma_D$ in $n^{\rm fix}$.
\end{proof}

\emph{Step 3: the tree part.} In this step we compare the form $D$, taken at the centre, with $s_{x'}[D']$.

\begin{lemma}[The tree part at the centre]\label{Y-H3-LTr} Let $\square_0$ and $M_X$ be as in Step~1. Then every $V'\in\mathcal U_{\{\square_0\}}(x')$ satisfies \[\bigl|\sum_{X\ {\rm under}\ \square_0}M_X(V')-s_{x'}[D'](\{\square_0\},V')\bigr|\le\sigma_D\,[L^3k_2\varepsilon'+x^{-20}].\]\end{lemma}

\begin{proof} We fix a site $y'\in\square_0$ and a site $y\in B(y')$, and we write $o:=y'=y(y)$. We also put $\tilde X_P:=T(y',y_P)X'_P$; this is the dressing at $y'$ with which $s_{x'}[D']$ is defined.

\textup{(1)} \emph{The field.} Let $\mathrm{Bx}$ be the box of radius $R_T$ about $o$. By \textup{(0.4)} in Step~0 of the proof of \Cref{Y-D-C5}, the box $\mathrm{Bx}$ lies inside $W_{\square_0}$ and inside $\square_0^\approx$. On $\mathrm{Bx}$ the field $V'$ has \textup{(u)}, and hence, by \textup{(0.6)} of the same Step~0, it has \textup{(h)} at $\hat\varepsilon=\varepsilon'$. For every plaquette $p$ with $|y-p|\le\varrho(x)$ we have ${\rm dist}(y,p)\le\varrho_L$, and the box of radius $R_3$ about $y(p)$ is contained in $\mathrm{Bx}$, since $R_3+2\varrho_L\le R_T$ by \textup{(0.1)}. Now \textup{(C3.1)} and \textup{(C3.2)} of item~D of \S\ref{Y-D-1}, together with \textup{(0.5)}, give \[\bar F_p=\sum_PH^\natural(p,P)\,T(y(p),y_P)X'_P+\Lambda_p\] with a remainder that satisfies $\|\Lambda_p\|\le \mathsf K_{\rm c}\varepsilon'^2$.

\textup{(2)} We put $dH:=H^\natural-H^{\sharp<}$. This kernel vanishes beyond $R_2$ by \textup{(C3.4)}, and its row sums are at most $\mathsf{ti}_{\rm c}$ by \textup{(C3.5)}. In $T(y',y(p))\bar F_p$ the kernel $dH$ gives a sum over $P$ in which $X'_P$ is carried by two transports, each of norm at most $4$ by \textup{(0.2)} and \textup{(0.7)}, while $|X'_P|\le2.06\varepsilon'$. So this contribution is at most $4\cdot4\cdot2.06\,\mathsf{ti}_{\rm c}\varepsilon'=32.96\,\mathsf{ti}_{\rm c}\varepsilon'\le\varepsilon'^2$, the last inequality being \textup{(F9$_k$)} of item~G of \S\ref{Y-D-1}.

\textup{(3)} We apply \Cref{Y-D-LA}, which is stated in Step~3 of the proof of \Cref{Y-D-C5}, with $n=2$ and $i_0=y$. As $K_1$ we take the point kernel at $(y,p)$; it is scalar, its coarse range does not exceed $\varrho_L$, and $\|K_1\|_\theta=e^{\theta\,{\rm dist}(y,p)}$. As the second kernel we take $K_2:=H^{\sharp<}$; its range is $R_2$, and its norm obeys $\|\cdot\|_{\theta_D}\le C_D$ by \textup{(Dec)}. The remaining choices are $\theta=\theta_D\le1$, $\rho_B=R_T$ and $\hat\varepsilon=\varepsilon'$. The hypotheses of the lemma are met: $\varrho_L+R_2\le R_T-1$ holds by \textup{(0.1)}, $3R_T+1\le\ell_w(x')$ by \textup{(0.2)}, $4R_T+2\le N$ by \textup{(0.3)}, and $35100R_T^2\varepsilon'\le1$ by \textup{(F8$_k$)}.

Let $\tilde a_p:=\sum_PH^{\sharp<}(p,P)\tilde X_P$. Since $\sup\|\tilde X_P\|\le8.24\varepsilon'$, the lemma yields \[\bigl\|T(y',y(p))\sum_PH^{\sharp<}(p,P)T(y(p),y_P)X'_P-\tilde a_p\bigr\|\le7020\cdot8\cdot1.2\cdot8.24\,\theta_D^{-2}e^{\theta_D{\rm dist}(y,p)}C_D\varepsilon'^2,\] and the numerical factor is $7020\cdot8\cdot1.2\cdot8.24\le555310.1$. From ${\rm dist}(y,p)\le(|y-p|+L-1)/L$, from $\theta_D\le1$ and from Step~0\,\textup{(iv)} we get $e^{\theta_D{\rm dist}}\le e\cdot e^{\lambda_\flat|y-p|/9}$; moreover $555310.1\,e\le1.50949\cdot10^6$.

\textup{(4)} Collecting \textup{(1)} to \textup{(3)}, we obtain \[b_p:=T(y',y(p))\bar F_p=\tilde a_p+e_p\] with \[\|e_p\|\le[4\mathsf K_{\rm c}+1+1.51\cdot10^6\theta_D^{-2}C_De^{\lambda_\flat|y-p|/9}]\varepsilon'^2;\] here the term $4\mathsf K_{\rm c}$ stems from $\Lambda_p$ and one transport, the term $1$ from \textup{(2)}, and the last term from \textup{(3)}. By \textup{($\Gamma$3)}, together with \textup{(L2$^{\rm su}$)} and with \textup{(0b)} of \textup{(Y.0)} in item~5 of \S\ref{Y-H-1}, we have $|\bar F_p|\le \mathsf N_{\rm c}L^{-2}\varepsilon'+\mathsf K_{\rm c}\varepsilon'^2\le\varepsilon'$, so that $\|b_p\|\le4\varepsilon'$. Further, we have $\|\tilde a_p\|\le8.24C_D\varepsilon'$, and since $C_D\ge4$, also $\|b_p+\tilde a_p\|\le9.24C_D\varepsilon'$.

\textup{(5)} As $b_p=\tilde a_p+e_p$, the difference $b_p\cdot b_{p'}-\tilde a_p\cdot\tilde a_{p'}$ equals $e_p\cdot b_{p'}+\tilde a_p\cdot e_{p'}$. Since $D_y$ and the form $Y\cdot Y'$ are symmetric, exchanging $p$ and $p'$ in the sum $\sum D_y(p,p')\,\tilde a_p\cdot e_{p'}$ turns it into the sum of $D_y(p,p')\,e_p\cdot\tilde a_{p'}$; with this argument, the same as in Step~2\,\textup{(iii)}, we get \[\sum D_y(p,p')(b_p\cdot b_{p'}-\tilde a_p\cdot\tilde a_{p'})=\sum D_y(p,p')\,e_p\cdot(b_{p'}+\tilde a_{p'}).\]

We use $\sum_pe^{-(8/9)\lambda_\flat|y-p|}\le3g_b^-$ and $\mathsf g_\flat\le g_b^-$. By the definition of $\|D\|$, the decay of $D_y(p,p')$ in $p$ absorbs the factor $e^{\lambda_\flat|y-p|/9}$ in the upper bound for $\|e_p\|$ obtained in \textup{(4)}, and the contribution of a single site $y$ to the right side of \textup{(5)} is at most \[\|D\|\cdot3g_b^-[4\mathsf K_{\rm c}+1+1.51\cdot10^6\theta_D^{-2}C_D]\varepsilon'^2\cdot3\mathsf g_\flat\cdot9.24C_D\varepsilon'.\] For a cube, this is multiplied by $(x/2)M^3$. Now $x\varepsilon'^3=Lp_0'^2\varepsilon'\le Lp_0^2\varepsilon'$ and $9\cdot9.24/612\le0.1359$, so that each cube contributes at most $k_2\varepsilon'\sigma_D$; there are $L^3$ cubes.

\textup{(6)} \emph{The free identity.} Let $D^<:=L\sum_{y\in B(y')}H^{\sharp<\mathsf T}D_yH^{\sharp<}$. Then \[(x/2)\sum_{y\in B(y')}\sum_{p,p'}D_y(p,p')\,\tilde a_p\cdot\tilde a_{p'}=(x'/2)\sum_{P,P'}D^<_{y'}(P,P')\,\tilde X_P\cdot\tilde X_{P'},\] because $x=Lx'$. By Step~1, the sites under $\square_0$ are the sites of the blocks of the sites of $\square_0$.

\textup{(7)} \emph{Tails.} Let $\mathrm{ct}$ denote the cut at $\varrho(x')$. Then \[D^<-D'=\{D^<-\mathrm{ct}\,D^<\mathrm{ct}\}-\mathrm{ct}\,(T_\ell D-D^<)\,\mathrm{ct}.\] By \textup{(0.8)}, the row sums of $H^{\sharp>}$ are at most $C_Dx^{-128}$. If $|y'-P|>\varrho(x')$, then ${\rm dist}(p,P)>\varrho(x')-\varrho_L$, and therefore the part of a row sum of $H^{\sharp<}$ that comes from such $P$ is at most $C_Dx'^{-64}e^{\theta_D\varrho_L}$. Here $\theta_D\varrho_L\le64\log x/L+2$; next, $x'^{-64}\le x^{-32}$, because $x\ge L^2$ by \textup{(P0)} of item~G of \S\ref{Y-D-1}; and $x^{64/L}\le x$, because $L\ge64$, which holds since \textup{(L2$^{\rm su}$)} with $\mathsf N_{\rm c}\ge1$ gives $L\ge1.44\cdot10^8$. Altogether, that part of a row sum is at most $e^2C_Dx^{-31}$.

Hence \[\sum_{P,P'}|D^<-D'|\le L^4\cdot9\mathsf g_\flat^2\|D\|C_D^2(2x^{-128}+2e^2x^{-31})\le161L^4\mathsf g_\flat^2C_D^2\|D\|x^{-31},\] because $9(2+2e^2)\le151.01\le161$. The transports in $\tilde X_P$ run along paths of at most $3R_T$ bonds, whence $|\tilde X_P\cdot\tilde X_{P'}|\le(8.24\varepsilon')^2$. We multiply by $(x'/2)(8.24\varepsilon')^2M^3=33.9488\,p_0'^2M^3$ and divide by $\sigma_D$; since $33.9488\cdot161/306\le17.862$, the tails contribute at most $17.9L^4C_D^2x^{-31}\sigma_D\le x^{-20}\sigma_D$, the last inequality being \textup{(Z.3)}.\end{proof}

\emph{Step 4: the loop part and clause \textup{(2$\tau$)}.} We show that for $V'\in\mathcal U_{\{\square_0\}}(x')$ and for every cube $X$ under $\square_0$ \[\frac12|\mathrm{Tr}_X(V')-\mathrm{Tr}_X(1)|\le S_{\rm tr}T_s^{-2},\] where $S_{\rm tr}:=28204\,M^3\mathsf g_\flat^2\|D\|/\gamma\le k_3p_0^{-2}\sigma_D$. Summing over the $L^3$ cubes, we get $k_3L^3p_0^{-2}x^{-1/2}\sigma_D$. The step also yields clause~\textup{(2$\tau$)} of item~5 of \S\ref{Y-H3-1}.

\emph{The gauge-box estimate.} We apply \Cref{Y-H1-LT57} with $T=T_s$. Then \textup{(C0)} to \textup{(C2)} are \textup{(Z.0)}, completed by $18\omega\varepsilon'\le1/10$, which follows from \textup{(0d)} of \textup{(Y.0)} in item~5 of \S\ref{Y-H-1}; the requirements on $\omega$ and on $\mathsf B_\omega$ hold by Step~0\,\textup{(vii)}.

\emph{The domain.} Let $o$ be the centre site of $\square_0$, and extend a field of $U^{17T_s}(\mathsf B_\omega;x')$ by the value $1$ to the bonds outside $\mathsf B_\omega$. By \textup{(e1)} and $\mathsf g_x<\vartheta_0^\flat/\log x'$, the field so extended lies in the interior of $\mathfrak V$; this is the inequality $|V'(c)|<1+\vartheta(x')+\mathsf g_x<1+\vartheta_0^\flat$ of Step~8, \textup{(8.2)}, of the proof of \Cref{Y-H-C9}. Moreover \textup{(e2)} gives $|V'(\pi)|<2$ for every path of at most $\ell_w(x')$ bonds.

\emph{The part of $V'$ on which $\mathrm{Tr}_X$ depends.} The coarse site $y'=y(y)$ lies within $M$ of $o$. The transport $T(y',y(p))$ depends on $V'$ within $\varrho_L$ of $y'$; the factor $\mathrm{d}^V(p,b)$ depends on the bond between $y(p)$ and $y(b)$; by \textup{(2b)} of \textup{($\Gamma$2)} in item~4 of \S\ref{Y-H-1}, $C(b,\cdot)$ depends on $V'$ within $R_2-1$ of $y(b)$; and $y(b)$ lies within $\varrho_L+1$ of $y'$. Altogether, $\mathrm{Tr}_X$ depends on $V'$ only within $R_2+\varrho_L$ of $y'$, a region inside $\mathsf B_\omega$ by Step~0\,\textup{(vii)}.

\emph{Holomorphy and invariance.} The function $\mathrm{Tr}_X$ is holomorphic by \textup{(2d)}, its other factors being polynomials in $V'$. Under a gauge transformation $u$ we have $T(y',y'')\to\operatorname{Ad}_{u(y')}T\operatorname{Ad}_{u(y'')}^{-1}$, and $\mathrm{d}^V(p,b)$ transforms in the same manner; for $C=GG^{\mathsf T}$ we use \textup{(2c)} together with $\operatorname{Ad}_u^{\mathsf T}=\operatorname{Ad}_u^{-1}$. Thus the matrix under the trace is conjugated by $\operatorname{Ad}_{u(y')}$, and the trace is unchanged. A gauge transformation $u$ on $\mathsf B_\omega$ is extended by $1$, in the manner of Step~8 of the proof of \Cref{Y-H-C9}.

\emph{Size.} On $\mathfrak V$ we have $|V'(c)|\le1+\vartheta_0^\flat\le1+10^{-4}$, by \textup{(O6)} of item~A of \S\ref{Y-D-1}; hence $\|\mathrm{d}^V(p,b)\|\le|V'(c)|^2\le1.01$. A plaquette has four bonds, so $\|(\mathrm{d}^VC\mathrm{d}^{V\mathsf T})(p,p')\|\le4.04^2\sup\|C(b,b')\|$. On $\mathfrak V$ we also have $\|C(b,b')\|\le\|C\|\le4/\gamma$ by \textup{(C2.5)} of item~D of \S\ref{Y-D-1}; this is \Cref{Y-C-L66}\,(a), which is stated on all of $\mathfrak V$. With $\|T\|,\|T^{\mathsf T}\|<4$ and $|\operatorname{tr}|\le3\|\cdot\|$ we find $|\tau|\le3\cdot16\cdot16.3216\cdot4/\gamma\le3133.75/\gamma$. As $\sum_{p,p'}|D_y|\le9\mathsf g_\flat^2\|D\|$ and there are $M^3$ sites, $\sup|\mathrm{Tr}_X|\le28203.8\,M^3\mathsf g_\flat^2\|D\|/\gamma\le S_{\rm tr}$. Finally, $S_{\rm tr}/\sigma_D\le92.170/(\gamma p_0^2)=6728.41\,L^4/p_0^2\le k_3/p_0^2$.

\emph{The bound.} Because $\omega\le Mr'$, the restriction to $\mathsf B_\omega$ of a field in $\mathcal U_{\{\square_0\}}(x')$ has \textup{(u1)} to \textup{(u3)}. The lemma therefore gives $|\mathrm{Tr}_X(V')-\mathrm{Tr}_X(1)|\le2S_{\rm tr}x^{-1/2}$.

\emph{Clause \textup{(2$\tau$)}.} Here the lemma is applied to the single function $\tau$ in the place of $\mathrm{Tr}_X$, with any value of $T$ for which \textup{(C0)} to \textup{(C2)} hold, and about any coarse site $o$ within $M$ of $y'$. The fields are extended by $1$ and lie in the interior of $\mathfrak V$ by \textup{(e1)}; the function $\tau$ depends on $V'$ only within $R_2+\varrho_L$ of $y'$, hence inside $\mathsf B_\omega$ by Step~0\,\textup{(vii)}; it is holomorphic and invariant by the two arguments just given; and $|\tau|\le3133.75/\gamma$. For the value $\tau(1)$ let $V'=1$ and let $u$ be constant: then \textup{(2c)} gives $C^1(b,b')\operatorname{Ad}_u=\operatorname{Ad}_uC^1(b,b')$; since $\operatorname{Ad}$ is irreducible on $\mathbb C^3$, each block is a scalar by Schur's lemma, and $\tau(1)=3(\mathrm{d}C^1\mathrm{d}^{\mathsf T})(p,p')$.

\emph{Step 5: the cut-off and the member factor.} We show that for every $V'\in\mathcal U_{\{\square_0\}}(x')$ and every cube $X$ under $\square_0$ \[|E_\xi[(1-\chi_X)Q_X]|\le3.47\,c_\xi p_0n_\chi e^{-c_\xi^2p_0^2/2}\sigma_D\] and \[|E_\xi[\chi_X(e^{E^a(X,U)}-1)s(X,U)]|\le9.09\,c_a(x)\sigma_D.\] When these bounds are doubled and summed over the $L^3$ cubes, they become at most $x^{-10}\sigma_D$ and $18.2L^3c_a(x)\sigma_D$ respectively.

\emph{The cut-off.} Let $\Xi:=\sup\|\xi_b\|$, the supremum running over the bonds tested by $\chi_X$. For a tested bond $b'$ there is a bond $b$ of $X^\approx$ with $|y(b')-y(b)|_\infty\le L(r_w+1)$; hence the base site of $b'$ lies within $L(r_w+1)+L-1<L(r_w+2)$ of the box $X^\approx$, which has side $M(2r(x)+1)$, and the number of tested bonds is at most $n_\chi$. Recall that $\|T\|<4$. By \textup{($\Gamma$3)}, \textup{(L2$^{\rm su}$)} and \textup{(0b)} we have $\|\bar F_p\|\le \mathsf N_{\rm c}L^{-3/2}\varepsilon+\mathsf K_{\rm c}\varepsilon'^2\le(10^{-4}+0.12)\varepsilon\le\varepsilon$. By \textup{(2a)} and \textup{(2e)}, $\|(\mathrm{d}^VG\xi)_p\|\le4\cdot1.01\,C_\varphi^\flat\Xi$; the factor $1.01$ is the bound $\|\mathrm{d}^V(p,b)\|\le|V'(c)|^2\le1.01$ of Step~4, obtained here from \textup{(u1)} at $x'$.

On the event $\chi_X=0$ we have $\varepsilon<\Xi/c_\xi=6.5\,C_\varphi^\flat\Xi$, and therefore $\|F^\flat_p\|\le10.54\,C_\varphi^\flat\Xi$ and \[\FitDisplay{\dimexpr\linewidth-1pt\relax}{|Q_X|\le(x/2)\,16\cdot9M^3\mathsf g_\flat^2\|D\|(10.54\,C_\varphi^\flat)^2\Xi^2.}\] Moreover $\Xi^21\{\Xi>t\}\le\sum_b\|\xi_b\|^21\{\|\xi_b\|>t\}$, the left side being the largest term on the right. Each $\xi_b$ has the form $\xi_b=x^{-1/2}g$ with $g$ a standard Gaussian vector in $\mathbb R^3$, and for $t_g\ge1$ \[\FitDisplay{\dimexpr\linewidth-1pt\relax}{E_\xi|g|^21\{|g|>t_g\}=(2/\pi)^{1/2}\{(t_g^3+3t_g)e^{-t_g^2/2}+3\int_{t_g}^\infty e^{-v^2/2}dv\}\le0.79789\,(t_g^3+3t_g+3/t_g)e^{-t_g^2/2}\le5.6\,t_g^3e^{-t_g^2/2},}\] since $0.79789\cdot7=5.58523$.

We take $t_g=c_\xi p_0$, which corresponds to $t=c_\xi\varepsilon$ because $\varepsilon=p_0x^{-1/2}$, and which satisfies $t_g\ge1$ by \textup{(Z.1)}. Using $C_\varphi^\flat=(6.5c_\xi)^{-1}$ we arrive at \[|E_\xi(1-\chi)Q|\le72(10.54/6.5)^2\cdot5.6\,M^3\mathsf g_\flat^2\|D\|n_\chi c_\xi p_0^3e^{-c_\xi^2p_0^2/2}.\] After division by $\sigma_D$ the constant is $(72/306)(10.54/6.5)^2\cdot5.6\le3.4647$. Doubling and summing over the $L^3$ cubes gives the factor $6.93\le7$, so that \textup{(Z.2)} yields the bound $x^{-10}\sigma_D$.

\emph{The member factor.} On the event $\chi_X=1$, Step~2\,\textup{(ii)} of the proof of \Cref{Y-H-C9}, taken at $\tau=1$, shows that the restriction of $U$ to $X^\approx$ lies in $\mathcal U_X(x)$. Hence $|s(X,U)|\le\sigma_D$ and $|E^a(X,U)|\le8c_a(x)\le1/4$. For $|w|\le1/4$ we have $|e^w-1|\le e^{|w|}-1\le1.13611|w|$, and $8\cdot1.13611\le9.0889$. Doubling and summing over the $L^3$ cubes yields $18.178\le18.2$. The term $E_\xi[(1-\chi_X)Q_X]$ contains no member, and it belongs to $n^{\rm fix}$ by the definitions of Step~1; the term $E_\xi[\chi_X(e^{E^a}-1)s]$ contains $E^a$ and belongs to $n^{\rm mem}$.

\emph{Step 6: the clusters with $n(C)\ge2$.} We show that \[\mathrm{CS}_{\tilde\kappa}(\mathrm{Rst})\le217L^3x^{-1/4}\sigma_D\le1000\,e^{\tilde\kappa+3}L^3x^{-1/4}\sigma_D.\] To this end we apply \Cref{Y-H1-U1}\,(h1) to the member $u\in M_8(x;p)$, with $\delta=s_x[D]$. Since $\delta$ lives on single cubes, Step~1\,\textup{(1)} shows that at every $\tilde\kappa$ its size in caps at $p$ is $B=\sup|s|/c_a(x)\le\sigma_D/c_a(x)$. Then \textup{(h1)} gives the bound $217L^3x^{-1/4}Bc_a(x)$, whence the first inequality. The second follows from $1000\,e^{\tilde\kappa+3}\ge1000\,e^9\ge8.1\cdot10^6$.

\emph{Step 7: assembly, clause (a), and the first of the four remarks of item~7.} We prove the two bounds of \Cref{Y-H3-U3}\,(a), \[\sup_{\mathcal U_{\{\square\}}(x')}|n^{\rm fix}(\{\square\},\cdot)|\le\sigma_D\bigl[L^3(2k_1\varepsilon+k_2\varepsilon')+k_3L^3p_0^{-2}x^{-1/2}+2x^{-10}\bigr]\] and $\mathrm{CS}_{\tilde\kappa}(n^{\rm mem})\le B_{\rm mem}(\tilde\kappa)$, and their two consequences \[\mathrm{CS}_{\tilde\kappa}(n)\le1001\,e^{\tilde\kappa+3}L^3x^{-1/4}\sigma_D\] and \[|n|^{{\rm cs},(\tilde\kappa)}_{a,x'}\le u_3(x')\|D\|_{(\lambda_\flat)},\] together with the first of the four remarks of item~7 of \S\ref{Y-H3-1}.

By Step~1 we have $n=DA_\ell(u)[s]-s_{x'}[D']=n^{\rm fix}+n^{\rm mem}$, since $M_X(1)=0$ and $s_{x'}[D'](\cdot,1)=0$. For $n^{\rm fix}$, Steps~2, 5, 4 and~3 give the bound \[\sigma_D[2L^3k_1\varepsilon+x^{-10}+k_3L^3p_0^{-2}x^{-1/2}+L^3k_2\varepsilon'+x^{-20}].\] For $n^{\rm mem}$ we use Steps~5 and~6; a term on a single cell has weight $1$ in the cell sum. Applying \textup{(Z.4)}, which is taken at the lower end $\kappa\le\tilde\kappa$, we get \[\mathrm{CS}_{\tilde\kappa}(n)\le(1000\,e^{\tilde\kappa+3}+e^{\kappa+3})L^3x^{-1/4}\sigma_D\le1001\,e^{\tilde\kappa+3}L^3x^{-1/4}\sigma_D.\]

\emph{In caps at $x'$.} As $x=Lx'$ and $c_a(x')=E_ap_0'^3x'^{-1/2}$, the quotient of this bound by $c_a(x')$ equals \[1001\,e^{\tilde\kappa+3}L^{3-1/4}(c_Q^{\rm su}/E_a)\,x'^{1/4}(p_0^2/p_0'^3)\|D\|.\] By Step~0\,\textup{(i)} and \textup{(iii)}, $p_0^2/p_0'^3=(p_0/p_0')^3/p_0\le1/0.5946$; furthermore $1001/0.5946\le1683.5\le1686$ and $e^{\tilde\kappa+3}\le e^{\kappa+4}$. The quotient is therefore at most $C_3x'^{1/4}\|D\|$, which is the second of the two consequences.

\emph{Proof of the first of the four remarks of item~7.} We use three relations. First, $\sigma_D\varepsilon=(c_Q^{\rm su}/E_a)c_a(x)\|D\|$, and $\varepsilon'\le L^{1/2}\varepsilon$. Second, \[\sigma_Dp_0^{-2}x^{-1/2}=(c_Q^{\rm su}/E_a)c_a(x)p_0^{-3}\|D\|\le(c_Q^{\rm su}/E_a)c_a(x)\|D\|.\] Third, $2x^{-10}\sigma_D\le L^3(c_Q^{\rm su}/E_a)c_a(x)\|D\|$. We apply these three relations to the bound on $n^{\rm fix}$ and find that each of its terms is at most a multiple of $(c_Q^{\rm su}/E_a)\,c_a(x)\|D\|$. By Step~0\,\textup{(iii)}, that is, by $c_a(x)\le c_a(x')/(0.5946L^{1/2})$, we then pass from $c_a(x)$ to $c_a(x')$, and $L^3/(0.5946L^{1/2})=L^{5/2}/0.5946$ is the factor that the remark carries. This proves the bound of that remark, in caps at $x'$.

\emph{Step 8: clauses (b), (c), (e) and (f).} We prove \Cref{Y-H3-U3}\,(b), (c), (e) and (f).

\emph{Clause (b).} For $0<\rho_z<z_0$, Cauchy's formula gives \[DA_\ell(u)[s](Y,V')=(2\pi i)^{-1}\int_{|z|=\rho_z}E'^{\rm old}[u+zs](Y,V')\,z^{-2}dz.\] At every $z$ the integrand has the clauses without a size, by \Cref{Y-H-C9}\,(b), which applies at $p$ by \Cref{Y-H1-U1}\,(U). That clause is a statement on $E'^a=E'^{\rm new}+E'^{\rm old}$; since $E'^{\rm new}$ does not depend on $z$ and $\int_{|z|=\rho_z}z^{-2}dz=0$, the integral of $E'^{\rm old}[u+zs]z^{-2}$ coincides with that of $E'^a[u+zs]z^{-2}$. The integrand is bounded, and it is holomorphic in $z$ and in $V'$ separately, hence jointly by Osgood's lemma. Consequently the integral is holomorphic on $\mathcal U_Y(x')$, invariant, equal to $0$ at pure gauges, and a function of the field on $Y^\approx$. By Step~1, $s_{x'}[D']$ has these clauses; hence $n$ has them too.

The functions $g_X$ and $E_\xi[\chi_X(e^{E^a}-1)s]$ are holomorphic on $\mathcal U_{\{\square_0\}}(x')$, by Step~2\,\textup{(ii)} of the proof of \Cref{Y-H-C9}, with domination by $\sigma_De^{1/4}$. They are invariant: the transformation $(\xi,V')\to(\operatorname{Ad}_u\xi,V'^u)$ leaves $\chi_X$, $E^a(X,U)$ and $s(X,U)$ unchanged, and $\operatorname{Ad}_u$ preserves the law of $\xi$; see the part on the clauses without a size in Step~8 of the proof of \Cref{Y-H-C9}. They depend on the field in $\square_0^\approx$, and their differences with their values at $1$ vanish at pure gauges. Hence $\mathrm{Rst}=DA_\ell(u)[s]-\sum[g_X-g_X(1)]$ has the clauses; so does $n^{\rm mem}$, and then also $n^{\rm fix}=n-n^{\rm mem}$.

\emph{Reality.} Let $u$ and $D$ be real, and let $z$ be real. Then $u+zs$ is real, and $E'^a$ is real at $G$-valued fields by \Cref{Y-H-C9}\,(d). Since $E'^{\rm new}$ does not depend on $z$, the $z$-derivative of $E'^a$, which equals that of $E'^{\rm old}$, is real. As $H^\sharp$ is real, so is $D'$, and $s_{x'}[D']$ is real at $G$-valued fields by Step~4 of the proof of \Cref{Y-I-LB}; the same holds for $g_X$.

\emph{Clause (c).} By \textup{(Fm1)}, $E'^{\rm old}[u](Y)$ is a sum over the clusters $C$ with $\operatorname{co}(C)=Y$, and $\operatorname{co}(C)$ contains the cores $J_\alpha$ of the letters of $C$. A letter (a) at $X$ has the core $J=Y_0(X)$, and $\Phi(C)$ depends on $E^a$ only through its letters (a). So $E'^{\rm old}[u+zs](Y)$ depends on $s(X)$ only at the cubes $X$ with $Y_0(X)\subset Y$, which are the cubes under $Y$; and $s(\{X\})$ depends on the $D_y$ with $y\in X$. The second term of $n$, namely $s_{x'}[D'](\{\square\})$, depends on the $D'_{y'}$ with $y'\in\square$, and $D'_{y'}$ depends on the $D_y$ with $y\in B(y')$, which lie under $\square$. Therefore both terms are $0$ when all $D_y$ under $Y$ vanish.

\emph{Clause (e).} The following are complex linear in $D$: $s_x[D]$, $T_\ell D$, the cut, the derivative $DA_\ell(u)[\cdot]$, and every one of the five terms of Step~1.

\emph{Clause (f).} As their definitions show, the numbers $k_1,k_2,k_3,C_3$ and the conditions \textup{(Z.0)} to \textup{(Z.4)} are functions of the letters of Step~0 only. Among these letters, the numbers that come from the rows are $C_D,\theta_D$ (for every $\ell\in \mathcal R(b_*,K_*,\lambda_*,1/9)$), $\mathsf K_{\rm c}$, $\mathsf a_{\rm c}$ and $\mathsf{ti}_{\rm c}$.

\emph{Empty rows.} We show that the steps above hold also for $\ell=\emptyset$, a case admitted by item~2 of \S\ref{Y-H3-1}.

By \Cref{Y-B-RZ} in Step~S6 of the proof of \Cref{Y-B-C1}, the zero rows ($\mathsf H_y=\mathbb R$, $\ell_y(p)=0$) yield the same objects as the empty rows; moreover, the zero rows form a member of $\mathcal R(b_*,K_*,\lambda_*,1/9)$, and this member satisfies \textup{(SD)} exactly when the empty rows satisfy it. Hence every statement made for all $\ell$ with \textup{(SD)} holds for the zero rows. The objects that \Cref{Y-B-RZ} lists are those of Sections~\ref{Y-A-1} and~\ref{Y-C-1} to~\ref{Y-G-1}. From the rows, Sections~\ref{Y-H-1} and~\ref{Y-I-1}, \Cref{Y-H1-U1} and the present section build only objects of three kinds: objects on that list; the numbers $|\ell_y(p)|$, $|\ell'_{y'}(P)|$ and $\mathsf E_y(p,p')=\langle\ell_y(p),\ell_y(p')\rangle$; and $T_\ell$, which depends on $H^\sharp$. Each of them is expressed through the columns $\ell_y(p)$, and so takes the same value at $0$ as at $\emptyset$.

\emph{Step 9: clause (d), the member.} We show three facts: $n$ depends on the member $u$ through $E^a$ alone; $n$ is holomorphic along affine lines in $M_8(x;p)$; and, for $u_1,u_2\in M_2(x;p)$, \[\mathrm{CS}_{\tilde\kappa}(n[u_1;D]-n[u_2;D])\le\frac13B_{\rm mem}(\tilde\kappa)\,d_a.\]

Neither $n^{\rm fix}$ nor $s_{x'}[D']$ depends on the member, and $E'^{\rm old}$ depends on $E^a$ alone, by the proof of clause (c) in Step~8.

\emph{Holomorphy.} Let $\zeta\to u_\zeta$ be an affine map with values in $M_8(x;p)$ on an open set. Then the map $(\zeta,z)\to u_\zeta+zs$ is holomorphic value by value and, for $|z|<z_0$, takes its values in $M_{16}(x;p)$. By \Cref{Y-H1-U1}\,(c+), $E'^{\rm old}[u_\zeta+zs](Y,V')$ is therefore holomorphic in $(\zeta,z)$; consequently its $z$-derivative at $0$ is holomorphic in $\zeta$.

\emph{The bound.} As $n$ depends on $E^a$ alone, we may take $E_2^h=E_1^h$ and $w_2=w_1$. If $d_a=0$, both sides of the bound vanish; in any case $d_a\le4$. For $|\zeta|<\zeta_0:=6/d_a$, a radius with $\zeta_0\ge3/2$, the member $u_\zeta:=u_1+\zeta(u_2-u_1)$ has an $a$-part within $2+|\zeta|d_a<8$ caps. Let $f(\zeta):=n^{\rm mem}[u_\zeta;D]-n^{\rm mem}[u_1;D]$; by \Cref{Y-H3-U3}\,(a), $\mathrm{CS}_{\tilde\kappa}(f(\zeta))\le2B_{\rm mem}$ on the disc, and $f(0)=0$. Since the torus is finite, only finitely many $Y$ contain a given cell $\square$. For a choice of points $V'_Y\in\mathcal U_Y(x')$ and of numbers $|c_Y|\le1$, the function \[\phi(\zeta):=\sum_{Y\ni\square}e^{\tilde\kappa(\#Y-1)}c_Yf(\zeta)(Y,V'_Y)\] is a finite sum of holomorphic functions, with $|\phi|\le2B_{\rm mem}$ and $\phi(0)=0$. By Schwarz's lemma, $|\phi(1)|\le2B_{\rm mem}/\zeta_0=\frac13B_{\rm mem}d_a$. The supremum over all such choices is the sum at $\square$. We bound this sum as a whole, and not term by term in $Y$ followed by a summation, because \Cref{Y-H1-U1}\,(h1) bounds the sum at each member.

\emph{Step 10: clause (g), the number $\mathsf V'$.} We show that \[D\mathsf V'(u)[s_x[D]]=\sum_y\frac32\sum_{p,p'}D_y(p,p')(\mathrm{d}C^1\mathrm{d}^{\mathsf T})(p,p')+\sum_\square e_\square,\] with the bound of \Cref{Y-H3-U3}\,(g).

By \Cref{Y-H1-U1}\,(g1), $\mathsf V'(u)=V_0^++\sum_Yv(Y;u)$, a finite sum, and $V_0^+$ does not depend on the member. Let $D\mathsf V'(u)[s]$ be the derivative at $z=0$ of $\mathsf V'(u+z(s,0,0))$; as the sum over $Y$ is finite, it equals $\sum_YDv(Y;u)[s]$. Because here $A=8<16$, the quantity $Dv(Y;u)[s]$ of item~1 of \S\ref{Y-H1-1}, which is a sum of derivatives of terms, equals the derivative $d/dz$ at $0$ of $v(Y;u+z(s,0,0))$. This equality is stated in \Cref{Y-H1-U1}\,(g3), for the case $A_a<16$.

\Cref{Y-H1-U1}\,(g3), applied with $\delta=s$, gives \[Dv(Y;u)[s]=\sum_{Y_0(X)=Y}g_X(1)+r(Y),\] where the weighted cell sum of $r$ is at most $109L^3x^{-1/4}\sigma_D$. By Step~1\,\textup{(3)} at $V'=1$, \[g_X(1)=\frac12\mathrm{Tr}_X(1)+\{E_\xi[\chi R_X]-E_\xi[(1-\chi)Q_X]+E_\xi[\chi(e^{E^a}-1)s]\}(1),\] and by Step~4, \[\frac12\mathrm{Tr}_X(1)=\sum_{y\in X}\frac32\sum_{p,p'}D_y(p,p')(\mathrm{d}C^1\mathrm{d}^{\mathsf T})(p,p').\] We define \[e_\square:=\sum_{X\ {\rm under}\ \square}\{E_\xi[\chi R_X]-E_\xi[(1-\chi)Q_X]+E_\xi[\chi(e^{E^a}-1)s]\}(1)+\sum\{r(Y):\ \square\ \text{the first cell of}\ Y\}.\] By Steps~2 and~5 without the doubling, the first sum in the definition of $e_\square$ has modulus $\le\sigma_D[L^3k_1\varepsilon+\frac12x^{-10}+9.09L^3c_a(x)]$. For the second sum we have $\bigl|\sum\{r(Y):\ \square\ \text{the first cell of}\ Y\}\bigr|\le\sum_{Y\ni\square}|r(Y)|\le109L^3x^{-1/4}\sigma_D$, by the bound on the weighted cell sum of $r$ stated above, whose weights $e^{\tilde\kappa(\#Y-1)}$ are $\ge1$. Since $9.09\le9.1$ and $\frac12x^{-10}\le x^{-10}$, adding the two estimates, we obtain $|e_\square|\le\sigma_D\bigl[L^3(k_1\varepsilon+9.1c_a(x))+109L^3x^{-1/4}+x^{-10}\bigr]$, which is the bound stated in \Cref{Y-H3-U3}\,(g).

\emph{Step 11: proof of \Cref{Y-H3-R3}.} We prove the identity for $\mathrm{rem}'$, its bound, the inequality $t'\le0.2431\,t+1$, and $t_k\le\max\{t_0,1.3212\}$.

Since $DA_\ell(u)[\cdot]$ is linear, \[DA_\ell(u)[s_x[D]+\mathrm{rem}]=s_{x'}[D']+n[u;D]+DA_\ell(u)[\mathrm{rem}].\] \Cref{Y-H1-U1}\,(h3), in the setting \textup{(L3su)}, gives \[\mathrm{CS}_{\tilde\kappa}(DA_\ell(u)[\mathrm{rem}])\le[2430\mathsf N_{\rm c}^2L^{-1/2}+L^{-1}]B_rc_a(x')\le\vartheta_aB_rc_a(x'),\] by $[0.2430+L^{-3/4}]L^{-1/4}\le0.2431L^{-1/4}$. Together with \Cref{Y-H3-U3}\,(a), this yields the upper bound for $\mathrm{rem}'$ that \Cref{Y-H3-R3} asserts. The clauses without a size hold for $\mathrm{rem}'$: for $DA_\ell(u)[\mathrm{rem}]$ by the argument given for clause (b) in Step~8, with $\mathrm{rem}$ in the place of $s$, and for $n$ by \Cref{Y-H3-U3}\,(b). From $u_3(x)/u_3(x')=L^{1/4}$ we get $\vartheta_au_3(x)/u_3(x')=0.2431$. The map $t\to0.2431\,t+1$ has the fixed point $1/0.7569\le1.32118$ and maps $[0,\max\{t_0,1.3212\}]$ into itself.

\emph{Step 12: proofs of the second and the third of the four remarks of item~7.} The second of these remarks is the identity over sub-alphabets, and the third is the sum rule; both are stated in item~7 of \S\ref{Y-H3-1}.

\emph{Proof of the second of the four remarks of item~7.} Item~6 of \S\ref{Y-H-1} calls a pair of letters incompatible only if the pair contains a zone letter or a letter (h): namely two zone letters, a letter (b), (v), (a) or (h) together with a zone letter, or the letters (a) and (h) of one $X$. Hence any two of the letters (a), (v), (t) are compatible, and \[Z_{\mathbb A}:=\sum_{\mathfrak B\subset\mathbb A}E_\xi\prod_{\mathfrak B}\lambda_\alpha=E_\xi\prod_{\mathbb A}(1+\lambda_\alpha).\] Since $\chi$ takes only the values $0$ and $1$, we have $1+\lambda^a_X=e^{\chi_XE^a(X,U)}$ and $1+\lambda^v_\square=e^{\chi_\square V_\square}$; for the function $\lambda^t_Y:=e^{T_Y}-1$ of the letter \textup{(t)} at $Y$ in item~6 of \S\ref{Y-H-1} we have $1+\lambda^t_Y=e^{T_Y}$; so $Z_{\mathbb A}=E_\xi e^{W_{\mathbb A}}$.

We apply the last claim of \Cref{Y-H1-LK} to the alphabet $\mathbb A$, with $K(\mathfrak B):=E_\xi\prod_{\mathfrak B}\lambda_\alpha$. It gives $Z_{\mathbb A}=\exp\sum_{C\ {\rm in}\ \mathbb A}\Phi(C)$, with an absolutely convergent series; in particular $Z_{\mathbb A}\ne0$. The hypotheses of the lemma are satisfied. First, $K(\mathfrak B\cup\mathfrak B')=K(\mathfrak B)K(\mathfrak B')$ whenever no letter of $\mathfrak B$ is $\sim$ to a letter of $\mathfrak B'$, by the independence shown in Step~3\,\textup{(IV)} of the proof of \Cref{Y-H-C9}. Second, $|K(\mathfrak B)|\le\prod e_\alpha$, by Step~6\,\textup{(iii)} of the proof of \Cref{Y-H-C9}, which concerns the letters (a), (v), (t) whose cores lie in $Y$, at $V'\in\mathcal U_Y(x')$. Third, \textup{(S$_\tau$)} holds a fortiori on a part of the alphabet.

We differentiate at $E^a+zs_X$. When $|z|<z_0$, the member stays within 16 caps and the sizes $e_\alpha$ to be used are the ones at $A=16$; hence the convergence of the series is uniform in $z$, and summing the $(d/dz)\Phi(C)$ gives $(dZ_{\mathbb A}/dz)/Z_{\mathbb A}$. The letter (a) at $X$ belongs to $\mathbb A$; were it absent, the left side would be $0$ while the right side would not. The cluster that consists of this one letter is $E_\xi[\chi_X(e^{E^a}-1)]$, and its derivative is $E_\xi[\chi_Xe^{E^a}s]$. Finally, the split by $\operatorname{co}(C)$ is the one of \textup{(Fm1)}.

\emph{Proof of the third of the four remarks of item~7.} The proof is that of Steps~3 to~6 of the proof of \Cref{Y-I-LB}, with changes that we give below. Those steps depend on $D$ only through its range, its reality and \textup{(G5)} of Step~0 of \S\ref{Y-I-1}. In all of them, $D$ is an arbitrary real symmetric family of range $\le\varrho(x)$, which takes the place of $\tilde{\mathsf E}-\mathsf E$, and in $q_y$ the factor $+(x/2)$ replaces $-(x/2)$. In the present section $\square$ denotes a cell, and we write $\boxdot$ for a cube; in \S\ref{Y-I-1} a cube is written $\square$, as in $q_\square$, $e^{q_\square}$ and $p_*(\square,V)$ of that proof. From Step~3 of that proof we take the definitions of $q_y$ and of $q_\boxdot$, with the properties that $q_y$ is real, Borel and gauge invariant and that $|q_y|\le\frac92\mathsf g_\flat^2p_0^2\|D\|_{(\lambda_\flat)}$, but not the identity of two energies of that step. Of Step~4 of that proof we use two parts: its claims, in the form of \Cref{Y-H3-I4}, and its agreement, which here reads $s_x[D](\{\boxdot\},V)=q_\boxdot(V)$ at every $G$-valued $V$ that is tame on $\boxdot^\approx$. That step proves, in \S\ref{Y-I-1}, the agreement $E^\Delta(\{\square\},V)=q_\square(V)$, in which both sides carry the factor $-(x/2)$; since the two sides change sign together, its proof, which proceeds term by term with $m_y=1$ and $\zeta_p=F_p$ and rests on \textup{(M5r)}, applies unchanged. In Step~5 of that proof, $w$ replaces $\tilde w$, $w_t$ replaces $w^\flat$, and $e^{tq_\boxdot}$ replaces $e^{q_\boxdot}$. In Step~6 of that proof, the condition $|t|\delta'\le c_\Xi p_0^2$ replaces \textup{(X15)}, \textup{(F3)} and $\delta_\Delta\le1$.

With $q_\boxdot:=\sum_{y\in\boxdot}q_y$ and $q_D=\sum q_\boxdot$, we have $q_\boxdot(V)=s_x[D](\{\boxdot\},V)$ at every cube that is on, by the agreement proved in Step~4 of that proof. By Step~5 of that proof, each off cube is assigned through $p_*(\boxdot,V)$, and these cubes factor out of $\Xi_w$; hence $\rho[u;\ell]e^{tq_D}=\rho[u_t;\ell]$ at every $V$.

Let $t$ be complex with $|t|\delta'\le c_\Xi p_0^2$, where $\delta'=(2r(x)+1)^3M^3\frac92\mathsf g_\flat^2p_0^2\|D\|$. The counting of Step~6 of that proof, combined with the strong cap of $w$, gives $|w_t|\le e^{-\tilde\kappa(|X|_c-1)}$; we give the details. Let $m:=\#(D(V)\cap X)$, and let $Q$ denote the sum, over the off cubes with $p_*(\boxdot,V)\in X$, of the $q_\boxdot(V)$. Here $D(V)$ is the set of the rough plaquettes of $V$, and $D'(V')$ that of $V'$; written without a field, $D=(D_y)$ and $D'$ are the kernels. The counting of Step~6, which depends on Step~3 of that proof and on \textup{(G4)} alone, gives $|Q|\le\delta'm$. If $m=0$, then $w_t=w$. If $m\ge1$, then $x\varepsilon\phi_p>p_0^2$ on $D(V)$, so that the strong cap gives $|w|\le e^{-c_\Xi p_0^2m}e^{-\tilde\kappa(|X|_c-1)}$, and \[|w_t|=|w||e^{tQ}|\le|w|e^{|t|\delta'm}\le e^{-\tilde\kappa(|X|_c-1)}.\] Finally, \textup{(M5r)} is \textup{(O7)}, and \textup{(F1)}, \textup{(F2)} hold at $x$ by Step~1\,\textup{(1)} of the present proof.

Hence $u_t\in M_2$ for small complex $t$, and the map $t\to u_t$ is holomorphic value by value. As $q_D$ is bounded, we may differentiate under the integral of the step $\rho\to\rho'$; the result is the mean of $q_D$ given $V'$. By \Cref{Y-H-C9}\,(a), $(\rho[u_t;\ell])'=e^{\mathsf V'[u_t]}\rho[F_\ell(u_t);\ell'^{\rm c}]$, and all values are holomorphic in $t$ by \Cref{Y-H1-U1}\,(c+). At a $G$-valued $V'$ with $D'(V')$ empty, that is without a rough plaquette, we have $\Xi'=1$, since the only family is the empty one, and $1_{Y^\approx}(V')=1$ for every $Y$. Hence, for $t$ close to $0$, the left side lies close to $\rho'(V')>0$, and there the logarithm of the right side equals \[\mathsf V'[u_t]-x'\mathrm{En}_{\ell'^{\rm c}}(V')+\sum_Y(E'^a+E'^h)(Y,V').\] By \textup{(Fm2)}, the number $\mathsf V'$ depends on the member through $E^a$ alone. The $t$-derivative of the logarithm is \[D\mathsf V'(u)[s]+\sum_YDA_\ell(u)[s](Y,V')+\sum_YDE'^h,\] where $DA_\ell(u)[s]=s_{x'}[D']+n$, and $\sum_\square s_{x'}[D'](\{\square\},V')=q'_{D'}(V')$ by Step~4 of the proof of \Cref{Y-I-LB}, taken at $x'$.\end{proof}

\section{Iteration with sources}\label{sec:iterate-sources}

\subsection{The Wilson measure with sources}\label{Y-H4-1}

In this section Wilson's weight with a complex coupling $\nu$, which assigns to each plaquette $p$ of $\mathbb T_0$ a value $\nu_p$ of its own and satisfies $|\nu|_\infty\le15.99c_\nu\varepsilon(\beta)$, is expressed through a system of members at the empty rows: at every field, $\exp\{-\beta\sum_p(1-\nu_p)s_p\}=\rho[u_0(\nu);\emptyset]$, where $u_0(\nu)$ stays within 16 caps and carries the strong cap at $\kappa_1$ together with a reserve. We verify \textup{(SD)} at $(x_k;\ell_k)$ for $0\le k<K$, that is, at the levels from which a step is made, among them $(x_0;\ell_0)=(\beta;\emptyset)$ with the empty rows, and we trace the first step $F_\emptyset$ along the family, with bounds on the members of $u_1(\nu)$, a Lipschitz bound, holomorphy, locality in $\nu$, and the vacuum number estimated cell by cell. We then compute the jets in the sources $z$ at level 0 and at level 1, follow the family $u_k(z)$ over all $K$ steps on a polydisc, and deduce in \Cref{Y-H4-Cor} that each $\rho_k$, $0\le k\le K$, lies in the class $C^{\rm su}_\flat(x_k;\ell_k)$ with the system $\bar u_k$.

\emph{Item 0: names.} We work with the group $\operatorname{SU}(2)$ in dimension $d=3$. The letters are fixed in item (0) of \S\ref{Y-H0-1}, and the numbers in item (1) of \S\ref{Y-H0-1}.

We take over from item 1 of \S\ref{Y-G0-1} the following notions and notation: the lattice, the cubes and the cells; the shift $\mathsf h_L$ of the grids; the operations $\hat{\ }$ and ${\rm cl}$; the convention that an object is \emph{of} $A$ when all its sites lie in $A$ and \emph{in} $A$ when its base site lies in $A$; the families $\mathcal D$ and $\mathcal P$; ${\rm hull}$, $|X|_c$, $X^{+j}$ and $Y_0(X)$; the functions $p_0$, $\varepsilon$, $\varrho(x)$, $r(x)$ and $s(x)=8r(x)$; the caps $c_a(x)=E_ap_0^3x^{-1/2}=E_ap_0^2\varepsilon$ and $c_h(x)$; $S_\varepsilon$ and $\mathrm{En}_\ell$; the sets $\mathcal U_X(x)$ with \textup{(u1)} to \textup{(u3)}; the members, with the weak and the strong cap; the densities $\rho[u;\ell]$; $\Xi_w$; the class $C^{\rm su}_\flat(x;\ell)$; and the data \textup{(SD)}.

Throughout this section, every notion and every statement that comes from Sections~\ref{Y-G0-1} to~\ref{Y-J-1} is used with its letter $\kappa$ replaced by $\kappa_1$.

The sets $M_A(x;p)$, $p\in\mathbb P$, are those defined in the part ``Pairs'' of item~0 of \S\ref{Y-H1-1}. We write $M_A(x):=M_A(x;\kappa_1,c_h)$; it is a subset of every $M_A(x;p)$.

The map $F_\ell$, the number $\mathsf V'$, and \textup{(Fm1)} to \textup{(Fm4)} come from item 6 of \S\ref{Y-H-1}, within the range that \Cref{Y-H1-U1}\,(U) allows.

By $\emptyset$, the empty rows, we mean the rows with $\mathsf H_y=\{0\}$. Hence $\mathrm{En}_\emptyset(U)=\sum_pS_\varepsilon(\phi_p)$, where $S_\varepsilon(t)=t^2/2$ for $t\le\varepsilon$ and $S_\varepsilon(t)=d_0\varepsilon t$ for $t>\varepsilon$.

When $Y$ is a set of cells, an object is said to lie under $Y$ if it is based in $\hat Y$.

The hull and the number $|X|_c=\#{\rm hull}(X)$ are those of item 0 of \S\ref{Y-B-1}: ${\rm hull}(X)$ is an element of $\mathcal D$ that contains $X$ and has the least number of cubes among such elements, the choice being fixed by an order. By item 1 of \S\ref{Y-G0-1}, ${\rm hull}(X)=X$ for $X\in\mathcal D$.

\emph{Letters with more than one sense.} The letter $r$ is the radius $r(\beta)$ (item~2), $r_\tau(Y)$ in Step~8 of the proof of \Cref{Y-H4-U4} is the remainder of \Cref{Y-H1-U1}\,(g3) at $u^\tau$, and the $r_i$ of \Cref{Y-H4-U4}\,(g) are radii; $n$ is the number of sources and, in the series for $\mathrm{Arc}$, the index of summation; and $n[u;D]$ is the remainder of \Cref{Y-H3-U3}. For $u$, the argument decides the sense: $u_0(\nu)$, $u_1(\nu)$, $u_k(z)$, $\bar u_k$ and $u^\tau$, whose argument is $\nu$, $z$ or absent, are members indexed by a level or a parameter, and $u_3(x_1)$, whose argument is the level $x_1$, is the number $C_3x_1^{1/4}$ of \S\ref{Y-H3-1}. The symbols $x_3$, $x_4$, $x_9^U$ and $x_s$ are floors, whereas the $x_k=L^{-k}\beta$ of item~2, among them $x_0$ and $x_1$, are levels.

\emph{Item 1: new entries of the table.}

\begin{enumerate}

\item[\textup{(O2a)}] We set $m_\Xi:=\frac12-d_0-c_\Xi$. By \textup{(O2)}, $\frac12(\frac12-d_0)\le m_\Xi<\frac12$.

\item[\textup{(O12a)}] We set $c_\nu:=E_a/(6M^3)$.

\item[\textup{(O13)}] We put $x_4:=1+\inf S_4$, where $S_4$ is the set of $y\ge x_9^U$ for which \textup{(W1)} to \textup{(W5)} below hold at every $x\ge y$:

\begin{enumerate}

\item[\textup{(W1)}] $\varepsilon(x)\le1/40$;

\item[\textup{(W2)}] $201M^3\varepsilon(x)\le E_a$;

\item[\textup{(W3)}] $\varepsilon^2/24+8c_\nu\varepsilon\le m_\Xi/4$;

\item[\textup{(W4)}] $(m_\Xi/4)p_0^2\ge\kappa_1s(x)+(2r(x)+1)^3[M^3p_0^4/(8x)+4c_a(x)]$;

\item[\textup{(W5)}] $e^{-(m_\Xi/2)p_0^2}\le x^{-10}$.

\end{enumerate}

\end{enumerate}

The table moreover gives $Lx_s\ge\max\{x_9^U,x_3,x_4\}$.

\emph{Item 2: data.} We fix an integer $K\ge1$ and a number $\mathsf b\ge x_s$, and we put $\beta=x_0=L^K\mathsf b$ and $x_k=L^{-k}\beta$.

The integer $N$ is even, $M\mid N$, and the condition \textup{(Tor)}, namely $N\ge2Ms(L\mathsf b)$, holds.

The tori are $\mathbb T_k=\mathbb Z^3/NL^{K-k}\mathbb Z^3$, with the grids $o_K=0$ and $o_k=Lo_{k+1}-\mathsf h_L$ mod $M$ for $0\le k<K$.

The rows are $\ell_0=\emptyset$ and $\ell_{k+1}=((\ell_k)')^{\rm c}$ cut at $\varrho(x_{k+1})$, and the steps $\rho\to\rho'$ are made with the noise $ax_k$. Wilson's density on $\mathbb T_0$ is $\rho_0=e^{-\beta\sum_p(1-\cos\phi_p)}$, and $\rho_{k+1}=(\rho_k)'$ for $0\le k<K$. The rows, the steps, the noise and these densities are those of \Cref{Y-GB-ThB}. The kernels, the law, the centre and the rows do not depend on any source.

Without an argument, $\varepsilon$, $p_0$, $r$, $s$ and $c_a$ are taken at $\beta$.

A coupling $\nu$ is a map from the plaquettes of $\mathbb T_0$ to $\mathbb C$, and $|\nu|_\infty$ is the supremum of its modulus.

\emph{Item 3: objects.} For a field $V$ and a plaquette $p$ we set $s_p(V):=1-\frac12\operatorname{tr}\partial V(p)$; this is a polynomial on the fields with values in $\operatorname{SL}(2,\mathbb C)$, and at $G$-valued fields $s_p=1-\cos\phi_p=a_p$.

For $|s|<2$ we write \[\mathrm{Arc}(s):=\arccos(1-s)^2=2\sum_{n\ge1}(2s)^n/(n^2\binom{2n}{n})=2s+s^2/3+4s^3/45+\dots,\] and $h(s):=\mathrm{Arc}(s)/2-s$, which is the function $h$ of \S\ref{Y-B-1}.

We put $q_p(V):=\beta h(s_p(V))$. The plaquettes in $D(V)=\{\phi_p>\varepsilon\}$ are called rough, and all other plaquettes are called tame.

A set $X\in\mathcal P$ is called \emph{$s$-connected} if the graph on its cubes in which two cubes are joined whenever they lie at distance $\le s$ is connected.

For a coupling $\nu$ we define the triple $u_0(\nu):=(E^a_\nu,0,w_\nu)$ as follows. For a single cube $Q$, \[E^a_\nu(\{Q\},V):=\sum_{p\in Q}[q_p(V)+\beta\nu_ps_p(V)],\] and $E^a_\nu(X,\cdot):=0$ whenever $\#X\ge2$.

For $X\in\mathcal P$, \[\FitDisplay{\dimexpr\linewidth-1pt\relax}{\begin{aligned}
&w_\nu(X,V):=1\{X\ s\text{-connected}\}
\\
&\quad{}\,1\{\text{every cube of }X\text{ holds a plaquette of }D(V)\}
\\
&\quad{}\exp\{-\beta\sum_{p\in X,\ {\rm rough}}[(1-\nu_p)s_p-d_0\varepsilon\phi_p]
\\
&\quad{}+\sum_{p\in X^{+r},\ {\rm tame}}[q_p+\beta\nu_ps_p]\}.
\end{aligned}}\]

At $\nu=0$, the triple $u_0(0)$ is Wilson's system $u_W$ of \S\ref{Y-B-1}, taken at $x=\beta$.

\emph{Sources.} Let $f_1,\dots,f_n$ be bounded complex functions for which the sets $\mathrm{supp}_0f_i:=\{p:f_i(x(p))\ne0\}$ are pairwise disjoint, where $x(p)$ denotes the barycentre of $p$.

We denote by $S(f_i)$ the set of cubes of $\mathbb T_0$ that hold a plaquette of $\mathrm{supp}_0f_i$; it coincides with $S_i(0)$ of item 2, \textup{(o3)}, of \S\ref{Y-H6-1}.

For complex numbers $z_i$ we put $\nu(z)_p:=\mathsf b^{-1}\sum_iz_if_i(x(p))$, so that $|\nu(z)|_\infty\le\max_i|z_i|\,\|f_i\|_\infty/\mathsf b$. Since $\beta/\mathsf b=L^K=a_K^{-1}$ and $\mathcal O_K(f)=a_K^{-1}\sum_pf(x(p))a_p$ by \Cref{Y-00-S3}, \[\exp\{-\beta\sum_p(1-\nu(z)_p)s_p\}=\rho_0\exp\{\sum_iz_i\mathcal O_K(f_i)\}.\]

We put $\mathrm{amp}_W:=1/(c_\nu\mathsf b\varepsilon(\beta))=6M^3x_0^{1/2}/(\mathsf bE_ap_0(x_0))$, and we consider the set $\mathrm{Pd}_0:=\{|z_i|\,\|f_i\|_\infty\le1.99/\mathrm{amp}_W\}$ together with the open set $\mathrm{Pd}_0^o:=\{|z_i|\,\|f_i\|_\infty<15.99/\mathrm{amp}_W\}$.

For a non-empty $S\subset\{1,\dots,n\}$ we write $d_S:=\prod_{i\in S}d/dz_i$ at $z=0$.

\emph{Item 4: the theorem and the corollary.}

\begin{theorem}[Wilson's weight with sources through all levels]\label{Y-H4-U4} In the setting of items 0 to 3, the following hold.

\begin{enumerate}

\item[\textup{(a)}] \emph{Level 0.} If $|\nu|_\infty\le15.99c_\nu\varepsilon$, then \[\exp\{-\beta\sum_p(1-\nu_p)s_p(U)\}=\rho[u_0(\nu);\emptyset](U)\] at every $G$-valued field $U$; that is, the identity holds with the factor $e^{\mathsf V}$ at $\mathsf V=0$, where $\mathsf V$ is the exponent of this factor in the form $e^{\mathsf V}\rho[u;\ell]$ of a point of $C^{\rm su}_\flat(x;\ell)$ (item~1 of \S\ref{Y-G0-1}). This value of $\mathsf V$ is the $V_0$ of \textup{(h)}. Moreover \textup{(a1)} to \textup{(a4)} hold.

\begin{enumerate}

\item[\textup{(a1)}] For every cube $Q$, \[\sup_{\mathcal U_{\{Q\}}(\beta)}|E^a_\nu(\{Q\},\cdot)|\le(0.00996+|\nu|_\infty/(c_\nu\varepsilon))c_a\le(0.01+|\nu|_\infty/(c_\nu\varepsilon))c_a\le16c_a,\] with equality in the last inequality exactly when $|\nu|_\infty=15.99c_\nu\varepsilon$. In addition, $(0.00996+|\nu|_\infty/(c_\nu\varepsilon))c_a\le15.99996c_a<16c_a$.

\item[\textup{(a2)}] For every field $V$ and every $X\in\mathcal P$, \[|w_\nu(X,V)|\le\exp\{-(c_\Xi+m_\Xi/2)\beta\varepsilon\sum_{p\in D(V)\cap X}\phi_p-\kappa_1(|X|_c-1)\},\] and wherever $w_\nu(X,V)\ne0$, also $\beta\varepsilon\sum_{D(V)\cap X}\phi_p\ge p_0^2\#X$. In other words, $w_\nu$ has the strong cap at $\kappa_1$ with the reserve $e^{-(m_\Xi/2)p_0^2\#X}\le x_0^{-10\#X}$. Consequently $u_0(\nu)\in M_A(\beta)$, where $A:=\max\{1,0.01+|\nu|_\infty/(c_\nu\varepsilon)\}\le16$, and $u_0(\nu)$ has the strong cap.

\item[\textup{(a3)}] In each of \textup{(ca)}, \textup{(ch)} and \textup{(c$\Xi$)}, the clauses without a size hold for $u_0(\nu)$. For real $\nu$, the triple $u_0(\nu)$ is real and $\rho[u_0(\nu);\emptyset]>0$; if $\nu$ is real with $|\nu|_\infty\le0.99c_\nu\varepsilon$, the weight of \textup{(a)} is a point of $C^{\rm su}_\flat(\beta;\emptyset)$.

\item[\textup{(a4)}] The function $E^a_\nu(\{Q\},V)$ is affine in $\nu$ and depends only on the $\nu_p$ with $p\in Q$; the function $w_\nu(X,V)$ is entire in the $\nu_p$ and depends only on those with $p\in X^{+r}$.

\end{enumerate}

Clause \textup{(a)} and \textup{(a1)} to \textup{(a4)} hold for every $\beta\ge x_4$, on every torus whose side is divisible by $M$, and for every grid.

\item[\textup{(b)}] \textup{(SD)} holds at $(x_k;\ell_k)$ for every $0\le k<K$. For $k=0$ this is $(\beta;\emptyset)$, and there $\ell'^{\rm c}=(\emptyset')^{\rm c}$ cut at $\varrho(\beta/L)$, so that $\ell'^{\rm c}=\ell_1$.

\item[\textup{(c)}] \emph{The first step.} For $|\nu|_\infty\le15.99c_\nu\varepsilon$, let $u_1(\nu):=F_\emptyset(u_0(\nu))$, and let $V_1(\nu)$ be the number $\mathsf V'$ of \textup{(Fm2)}. Then \[(\exp\{-\beta\sum(1-\nu_p)s_p\})'=e^{V_1(\nu)}\rho[u_1(\nu);\ell_1]\] at every coarse field $V'$, and $E'^a=E'^{\rm new}+E'^{\rm old}[u_0(\nu)]$, where $E'^{\rm new}$ does not depend on $\nu$.

Moreover, for each cell and each $\tilde\kappa\in[\kappa,\kappa_1]$, the three sums, over the $Y$ that hold the cell and with the weight $e^{\tilde\kappa(\#Y-1)}$, of $\sup|E'^{\rm new}(Y,\cdot)|$, of $\sup|E'^{\rm old}[u_0(\nu)](Y,\cdot)|$ and of $\sup|E'^h(Y,\cdot)|$ are $\le0.35c_a(x_1)$, $\le0.248Ac_a(x_1)$ and $\le c_h(x_1)$ respectively, and $w'$ has the strong cap at $x_1$. Hence $u_1(\nu)\in M_1(x_1)$ for $|\nu|_\infty\le1.99c_\nu\varepsilon$, and $u_1(\nu)\in M_2(x_1)$ for $|\nu|_\infty\le6.6c_\nu\varepsilon$.

If $|\nu|_\infty,|\tilde\nu|_\infty\le1.99c_\nu\varepsilon$, then $\mathsf D_{x_1}(u_1(\nu),u_1(\tilde\nu))\le0.4|\nu-\tilde\nu|_\infty/(c_\nu\varepsilon)$. For real $\nu$, both $u_1(\nu)$ and $V_1(\nu)$ are real. Finally, the number $V_1(\nu)$ and all values of $u_1(\nu)$ depend holomorphically on $(\nu_p)$ in the region $|\nu|_\infty<15.99c_\nu\varepsilon$; this is \Cref{Y-H1-U1}\,(c+).

\item[\textup{(d)}] \emph{Locality.} The function $E'^{\rm old}[u_0(\nu)](Y,\cdot)$ and the number $v_0(Y;\nu)$ of \textup{(e)} depend only on the $\nu_p$ with $p$ under $Y$. The functions $E'^h(Y,\cdot)$ and $w'(Y,\cdot)$ depend only on the $\nu_p$ with $p$ under $Y^{+\lceil r(\beta)/L\rceil+1}$.

\item[\textup{(e)}] \emph{Vacuum.} One has $V_1(\nu)=V_1^o+\sum_{Y\in\mathcal D}v_0(Y;\nu)$, where $v_0(Y;\nu):=v(Y;u_0(\nu))$ is the number of \Cref{Y-H1-U1}\,(g1) and $V_1^o$ does not depend on $\nu$. For every cell and every $\tilde\kappa\in[\kappa,\kappa_1]$, \[\sum_{Y\ \text{holding the cell}}e^{\tilde\kappa(\#Y-1)}|v_0(Y;\nu)|\le1.30L^3(0.01+|\nu|_\infty/(c_\nu\varepsilon))c_a(\beta)\le2.43L^3Ac_a(\beta)\le C_gAc_a(\beta).\]

\item[\textup{(f)}] \emph{Jets at level 0.} Let $\nu=\nu(z)$ and $\delta_S(0):=d_Su_0(\nu(z))$. For the $a$-part, $\delta_i^a(0)=s_{x_0}[\mathbb D_0[f_i]/\mathsf b]+\mathrm{rem}_{0,i}$, where \[s_{x_0}[\mathbb D_0[f]/\mathsf b](\{Q\},V)=(x_0/2\mathsf b)\sum_{p\in Q}f(x(p))F_p\cdot F_p\] and \[\mathrm{rem}_{0,i}(\{Q\},V)=-(1/\mathsf b)\sum_{p\in Q}f_i(x(p))q_p(V).\] Both terms are carried by single cubes of $S(f_i)$ and satisfy the clauses of \textup{(ca)} without a size. Measured in caps $c_a(\beta)$, $|\delta_i^a(0)|\le\mathrm{amp}_W\|f_i\|_\infty$, $|s_{x_0}[\mathbb D_0[f_i]/\mathsf b]|\le1.01\,\mathrm{amp}_W\|f_i\|_\infty$ and $|\mathrm{rem}_{0,i}|\le0.01\|f_i\|_\infty/\mathsf b$. For $|S|\ge2$ the $a$-part of $\delta_S(0)$ is 0.

The $h$-part of $\delta_S(0)$ is 0. Its $w$-part is \[\delta_S^w(0)(X,V)=w_0(X,V)\prod_{i\in S}(x_0/\mathsf b)\sum_{p\in\mathrm{Pl}(X,V)}f_i(x(p))s_p,\] where $\mathrm{Pl}(X,V)$ is the set consisting of the rough plaquettes in $X$ together with the tame plaquettes in $X^{+r}$. The $w$-part is 0 unless $X^{+r}$ meets $S(f_i)$ for each $i\in S$, and $|\delta_S^w(0)(X,V)|\le e^{-(m_\Xi/2)p_0^2\#X}\prod_{i\in S}(\mathrm{amp}_W\|f_i\|_\infty/15.99)$ times the strong cap at $\kappa_1$. The leading term of $\delta_i(0)$ is the triple $(s_{x_0}[\mathbb D_0[f_i]/\mathsf b],0,0)$, and the whole $w$-jet is counted in the remainder.

\item[\textup{(g)}] \emph{Jets at level 1.} Let $\delta_S(1):=d_Su_1(\nu(z))$. \emph{First half.} Then the functions $\delta_S^a(1)(Y,\cdot)$ and the numbers $d_Sv_0(Y)$ vanish except when $Y$ meets ${\rm cl}\,S(f_i)$ for every $i\in S$, while $\delta_S^h(1)(Y,\cdot)$ and $\delta_S^w(1)(Y,\cdot)$ vanish except when $Y^{+\lceil r(\beta)/L\rceil+1}$ meets ${\rm cl}\,S(f_i)$ for every $i\in S$.

The bounds below are stated for $0<t\le15.98$, with $r_i:=t/(\mathrm{amp}_W\|f_i\|_\infty)$ and $A_t:=\max\{1,t+0.01\}$, and they hold for every cell and every $\tilde\kappa\in[\kappa,\kappa_1]$. If $\|f_i\|_\infty=0$, then the coupling $\nu(z)$ is independent of $z_i$, each $\delta_S(1)$ and each $d_Sv_0$ with $i\in S$ is 0, and we take $r_i^{-1}$ to be 0. The sums over the $Y$ that hold the cell, with the weight $e^{\tilde\kappa(\#Y-1)}$, of $\sup|\delta_S^a(1)(Y,\cdot)|$, of $\sup|\delta_S^h(1)(Y,\cdot)|$ and of $|d_Sv_0(Y)|$ are $\le0.248A_tc_a(x_1)\prod_Sr_i^{-1}$, $\le c_h(x_1)\prod_Sr_i^{-1}$ and $\le2.43L^3A_tc_a(\beta)\prod_Sr_i^{-1}$ respectively, and $|\delta_S^w(1)(Y,V')|\le\prod_Sr_i^{-1}$ times the strong cap at $x_1$.

\emph{Second half.} Moreover, $\delta_i^a(1)=s_{x_1}[\mathbb D_1[f_i]/\mathsf b]+\mathrm{rem}_{1,i}$ with $\mathrm{rem}_{1,i}=n[u_0(0);\mathbb D_0[f_i]/\mathsf b]+DA_\emptyset(u_0(0))[\mathrm{rem}_{0,i}]$, and for all $\tilde\kappa\in[\kappa,\kappa_1]$, \[|\mathrm{rem}_{1,i}|_{a,x_1}^{{\rm cs},(\tilde\kappa)}\le(u_3(x_1)K_5+0.01\,\vartheta_a)\|f_i\|_\infty/\mathsf b.\]

\item[\textup{(h)}] \emph{The honest family.} Put $u_0(z):=u_0(\nu(z))$ and $V_0:=0$, and for $0\le k<K$ put $u_{k+1}(z):=F_{\ell_k}(u_k(z))$ and let $V_{k+1}(z)$ be the number $\mathsf V'$ at $u_k(z)$; we call $u_k(z)$, $V_k(z)$, $0\le k\le K$, the honest family. On $\mathrm{Pd}_0^o$ all of these are defined; there $u_k(z)\in M_{A^{(k)}}(x_k)$ with $A^{(0)}=16$, $A^{(1)}=4.318$, $A^{(2)}=1.421$ and $A^{(k)}=1$ for $k\ge3$, each $u_k(z)$ has the strong cap, and $V_k(z)$ and every value of $u_k(z)$ are holomorphic. On $\mathrm{Pd}_0$, $u_0(z)\in M_2(x_0)$ and $u_k(z)\in M_1(x_k)$ for $1\le k\le K$.

Let $\rho_0^z:=\rho_0\exp\{\sum z_i\mathcal O_K(f_i)\}$ and $\rho_{k+1}^z:=(\rho_k^z)'$. Then $\rho_k^z=\exp\{\sum_{j\le k}V_j(z)\}\rho[u_k(z);\ell_k]$ at every field of $\mathbb T_k$, and \[Z_K(z):=\int\rho_0^z=\exp\{\sum_{j\le K}V_j(z)\}\int\rho[u_K(z);\ell_K].\] When all $f_i$ are real, all these objects are real at each real point $z$ of $\mathrm{Pd}_0^o$, and $\rho[u_k(z);\ell_k]>0$ there; for arbitrary $f_i$ the same is true at $z=0$.

Hence, for $0\le k\le K$, $\bar u_k:=u_k(0)$ is a system of members of $\rho_k$: it is real, it lies in $M_1(x_k)$, where $M_1(x_k)=M_1(x_k;\kappa_1,c_h)$ by item~0, and it has the strong cap. The radius $1.99/\mathrm{amp}_W=1.99c_\nu\mathsf b^{1/2}L^{-K/2}p_0(L^K\mathsf b)$ depends on $K$.

\item[\textup{(i)}] The numbers $c_\nu$, $m_\Xi$, $x_4$, and every numeral above, do not depend on $K$, $N$, $\mathsf b$, $\nu$, $f$ or $n$.

\end{enumerate}

\end{theorem}

\begin{corollary}[Wilson's density and its images lie in the classes]\label{Y-H4-Cor}
With the letters as fixed in the table, at every field, $\rho_0=e^{-\beta\sum(1-\cos\phi_p)}\in C^{\rm su}_\flat(\beta;\emptyset)$ and $\rho_k\in C^{\rm su}_\flat(x_k;\ell_k)$ for $1\le k\le K$, with the systems $\bar u_k$. For $k=0$ and $k=1$ the two memberships are \textup{(R0)} and \textup{(R1)} of \S\ref{Y-B-1}, in the case of the letters as fixed in the table and of the grids $o_k$ of item~2 of \S\ref{Y-H4-1}; at $k=1$ the system is $\bar u_1=F_\emptyset(u_W)$, defined with $F_\emptyset$ of \S\ref{Y-H-1} in the place of the map of \textup{(F1)} to \textup{(F4)} of \S\ref{Y-G-1}, and the two maps give systems of the same $\rho_1$.

In \Cref{Y-GB-ThB}, for the data of item~2 of \S\ref{Y-H4-1}, that is, with the values of the table, for $\mathsf b\ge x_s$ and at the rate $\kappa_1$, \Cref{Y-H4-U4}\,(a), (b) can be used instead of the hypothesis that \Cref{Y-B-C1} is given; the induction then begins at $k=0$. In \Cref{Y-J-ThA}\,(ii), the system that is fixed may be $\bar u_{k_0}$.
\end{corollary}

\begin{proof}[Proof of \Cref{Y-H4-U4} and of \Cref{Y-H4-Cor}]
\emph{Step 1: the floor $x_4$ and its place in the table.} We first show that $x_4$ exists. Only strictly earlier places of the table enter $m_\Xi$, $c_\nu$ and $x_4$; none of $N$, $K$, $\mathsf b$, $\nu$, $f$, $n$ enters them, and neither does a member or a row. Once $x\ge x_4$, we have $c_\nu\varepsilon\le1/64$, and therefore $\mathrm{amp}_W\ge64/\mathsf b$.

By \textup{(O8)}, $p_0^2=A_0^2(\log x)^{2\mathfrak p_0}$ with $\mathfrak p_0>2$, while both $\varepsilon=p_0x^{-1/2}$ and $c_a=E_ap_0^3x^{-1/2}$ tend to $0$ as $x$ grows; hence \textup{(W1)}, \textup{(W2)} and \textup{(W3)} are true from some $x$ on. In \textup{(W4)} the left side grows like $(\log x)^{2\mathfrak p_0}$, where $2\mathfrak p_0>4$; on the right side, $\kappa_1s(x)\le8\kappa_1(20c_R\log x/M+1)$, while $(2r(x)+1)^3$ is of order $(\log x)^3$ and multiplies a factor tending to $0$. Condition \textup{(W5)} states that $(m_\Xi/2)A_0^2(\log x)^{2\mathfrak p_0}\ge10\log x$. So $S_4$ is a non-empty upper half line, and $x_4=1+\inf S_4$ belongs to it.

\emph{The least threshold.} Apart from $r$ and $s$, which are non-decreasing and enter only the right side of \textup{(W4)}, all functions involved are continuous in $x$. Assume that the five conditions hold for every $x>t$. Replacing $r(x)$, $s(x)$ by $r(t)$, $s(t)$ does not increase the right side of \textup{(W4)}; letting $x$ decrease to $t$, we conclude that the five conditions hold at $t$. Hence the set of thresholds is a closed half line, and it has a least point.

\emph{Order.} The number $m_\Xi$ depends on \textup{(O1)} and \textup{(O2)}, the number $c_\nu$ on \textup{(O7)} and \textup{(O12)}, and $x_4$ on these, on $\kappa_1$ \textup{(O5)}, $c_R$ \textup{(O6)}, $M$ \textup{(O7)}, $A_0$ and $\mathfrak p_0$ \textup{(O8)}, $E_a$ \textup{(O12)}, and on $x_9^U$. The floor $x_s$ is chosen after $x_4$. \emph{The last claim.} From \textup{(W3)} and $m_\Xi<1/2$ we obtain $8c_\nu\varepsilon\le1/8$.

\emph{Step 2: five elementary facts.} We need the following five facts.

\begin{enumerate}
\item[\textup{(E1)}] Let $g\in\operatorname{SL}(2,\mathbb C)$. Then $1-\frac12\operatorname{tr}g=\frac12\det(g-1)$, and consequently $|1-\frac12\operatorname{tr}g|\le\frac12|g-1|^2$.
\item[\textup{(E2)}] The coefficients of the series $h$ are $>0$, and its radius of convergence is $2$. Moreover $|h(t)|\le h(|t|)\le(|t|^2/6)/(1-|t|/2)$, and for $0\le\phi<\pi$ the number $h(1-\cos\phi)=\phi^2/2-(1-\cos\phi)$ lies in $[0,\phi^4/24]$.
\item[\textup{(E3)}] Let $g\in\operatorname{SL}(2,\mathbb C)$ with $|g-1|<1/20$, and let $F=-i\log g$. Then $F\cdot F=\mathrm{Arc}(1-\frac12\operatorname{tr}g)$.
\item[\textup{(E4)}] If $0<\varepsilon\le1$ and $\phi\in[\varepsilon,\pi]$, then $1-\cos\phi\ge(\frac12-\varepsilon^2/24)\varepsilon\phi$.
\item[\textup{(E5)}] If $X\in\mathcal P$ is $s$-connected, then $|X|_c-1\le s(\#X-1)$.
\end{enumerate}

\emph{Proof of \textup{(E1)}.} Write $g=1+y$. From $1=\det(1+y)=1+\operatorname{tr}y+\det y$ we get $\operatorname{tr}y=-\det y$, so that $1-\frac12\operatorname{tr}g=-\frac12\operatorname{tr}y=\frac12\det y$. Finally $|\det y|$ is the product of the two singular values of $y$, and is therefore $\le|y|^2$.

\emph{Proof of \textup{(E2)}.} It rests on the classical series \[\arcsin^2z=\frac12\sum_{n\ge1}(2z)^{2n}/(n^2\binom{2n}{n}),\] which we derive first. The function $\mathrm{fa}(z):=\arcsin^2z$ is even and holomorphic on $|z|<1$; it satisfies $\mathrm{fa}(0)=0$ and $(1-z^2)\mathrm{fa}'^2=4\,\mathrm{fa}$. Differentiating this relation and dividing by $2\,\mathrm{fa}'$, we get $(1-z^2)\mathrm{fa}''-z\,\mathrm{fa}'=2$. Written for $\mathrm{fa}=\sum_{n\ge1}\breve\gamma_nz^{2n}$, it means that $\breve\gamma_1=1$ and $(2n+2)(2n+1)\breve\gamma_{n+1}=4n^2\breve\gamma_n$; both are satisfied by $\breve\gamma_n=\frac124^n/(n^2\binom{2n}{n})$, which proves the displayed formula.

Since $\arccos(1-t)=2\arcsin((t/2)^{1/2})$, the series gives \[\frac12\arccos(1-t)^2=\sum_{n\ge1}(2t)^n/(n^2\binom{2n}{n})=t+h(t).\] Two consecutive coefficients have the quotient $n^2/((n+1)(2n+1))<1/2$. Hence the radius is $2$, and the part of the series from $n=2$ on is bounded by its first term $|t|^2/6$ times $1/(1-|t|/2)$. At $t=1-\cos\phi<2$ we have $\arccos(\cos\phi)=\phi$. Finally $\cos\phi\le1-\phi^2/2+\phi^4/24$ for real $\phi$; together with the positivity of the coefficients, this places $h(1-\cos\phi)$ in $[0,\phi^4/24]$.

\emph{Proof of \textup{(E3)}.} We have $\det g=e^{i\operatorname{tr}F}=1$ and $|\operatorname{tr}F|\le2|F|<2\pi$, so $\operatorname{tr}F=0$. For a traceless $2\times2$ matrix, the theorem of Cayley--Hamilton gives $F^2=-\det F\,1=(F\cdot F)1$, as $F\cdot F=\frac12\operatorname{tr}F^2$. Therefore, with $w=F\cdot F$, \[g=e^{iF}=\cos(w^{1/2})1+iF\sin(w^{1/2})/w^{1/2},\] and $1-\frac12\operatorname{tr}g=c(w):=1-\cos w^{1/2}$, which is an entire function of $w$. On $|w|<1$ we have $|c(w)|\le\cosh1-1\le0.5431<2$, so that $\mathrm{Arc}(c(w))$ is holomorphic there; since it coincides with $w$ on $[0,1)$, it coincides with $w$ on the whole disc. It remains to note that $|w|\le|F|^2\le(\log(20/19))^2\le0.0513^2<1$.

\emph{Proofs of \textup{(E4)} and \textup{(E5)}.} The monotonicity argument in Step~S4\,(1) of the proof of \Cref{Y-B-C1} proves \textup{(E4)} for $0<\varepsilon\le1$: its endpoint comparison only requires $\varepsilon/2\le2/\pi$. For \textup{(E5)} apply the spanning-tree construction of Step~S4\,(4) there. It uses only $s$-connectivity and the definition of $|X|_c$ as the least size of a connected set containing $X$, and applies on a torus as well.

\emph{Step 3: \textup{(a1)}, \textup{(a3)} and \textup{(a4)}, the analytic member and the clauses without a size.} We prove \textup{(a1)}; that $E^a_\nu(\{Q\},\cdot)$ is holomorphic on $\mathcal U_{\{Q\}}(\beta)$, does not change under $G$-valued gauge transformations, vanishes at pure gauges, and is real at $G$-valued fields when $\nu$ is real; that $w_\nu(X,\cdot)$ is Borel, gauge invariant, and a function of $V$ on the bonds of $X^{+s}$; and \textup{(a4)}.

The sites of a plaquette in $Q$ lie in $Q^{+1}$, which is a subset of $Q^{+r}$ ($r\ge1$ since $\beta>1$); thus the plaquette is a plaquette \emph{of} $Q^{+r}$, and \textup{(u3)} gives $|\partial V(p)-1|<2\varepsilon$. Now \textup{(E1)} yields $|s_p|<2\varepsilon^2\le1/800$, and \textup{(E2)} yields \[|q_p|\le\beta(4\varepsilon^4/6)/(1-\varepsilon^2)\le0.66709\beta\varepsilon^4=0.66709p_0^4/\beta.\] As a cube holds $3M^3$ plaquettes, \[|E^a_\nu|\le2.00127M^3p_0^4/\beta+6M^3|\nu|_\infty p_0^2.\] In the first term, $p_0^4/\beta=\varepsilon p_0^3\beta^{-1/2}$, and \textup{(W2)} gives $M^3\varepsilon\le E_a/201$; so the first term is $\le(2.00127/201)c_a\le0.00996c_a$. In the second term we use $c_a=E_ap_0^2\varepsilon$ and $6M^3c_\nu=E_a$, which give $6M^3|\nu|_\infty p_0^2=(|\nu|_\infty/(c_\nu\varepsilon))c_a$. For $|\nu|_\infty\le15.99c_\nu\varepsilon$ the sum is therefore $\le15.99996c_a$.

\emph{The clauses.} Since $s_p$ is a polynomial in the entries of $V$ and $|s_p|<2$, the function $h(s_p)$ is holomorphic. A gauge transformation takes $\partial V(p)$ to $u\,\partial V(p)u^{-1}$, which leaves the trace unchanged. At a pure gauge, $\partial V(p)=1$, so $s_p=0$ and $h(0)=0$. At $G$-valued fields, $s_p=1-\cos\phi_p$ is real. Moreover $E^h=0$. The weight $w_\nu(X,\cdot)$ is built from $\phi_p$ and $s_p$ at plaquettes in $X^{+r}$, whose sites lie in $X^{+r+1}$, a subset of $X^{+s}$ ($r+1\le8r$); and, $\phi_p$ being a continuous class function of $\partial V(p)$, the indicators are Borel. For real $\nu$ the weight $w_\nu\ge0$ is real, and by \textup{(a)} (Step~4) $\rho[u_0(\nu);\emptyset]$ equals $\exp\{-\beta\sum_p(1-\nu_p)s_p\}$, the left side of \textup{(a)}, so that $\rho[u_0(\nu);\emptyset]>0$. If moreover $|\nu|_\infty\le0.99c_\nu\varepsilon$, then $0.01+|\nu|_\infty/(c_\nu\varepsilon)\le0.01+0.99=1$, so $A=1$ in \textup{(a2)}, and \textup{(a2)} (Step~5) places $u_0(\nu)$ in $M_1(\beta)$, with the strong cap; as $u_0(\nu)$ is real and $\rho[u_0(\nu);\emptyset]>0$ (both just shown), and as the exponent of the factor in front of $\rho[u_0(\nu);\emptyset]$ is $0$, a real number (\textup{(a)}, Step~4), the weight of \textup{(a)} is then a point of $C^{\rm su}_\flat(\beta;\emptyset)$. Clause \textup{(a4)} follows from the formulas of item~3 of \S\ref{Y-H4-1}, since $\mathbb T_0$ has finitely many plaquettes.

\emph{Step 4: \textup{(a)}, the identity at every field.} We show that, for every complex $\nu$ and every $G$-valued $V$, \[\exp\{-\beta\sum_p(1-\nu_p)s_p\}=\exp\{-\beta\sum_pS_\varepsilon(\phi_p)+\sum_Q1_{Q^{+r}}(V)E^a_\nu(\{Q\},V)\}\,\Xi_{w_\nu}(V).\] Denote by $d(V)$ the set of those cubes that hold a rough plaquette, and by $X_1,\dots,X_j$ its classes under the links at distance $\le s$.

\begin{enumerate}
\item[\textup{1.}] \emph{One family.} Take a family $\{X_i'\}$ of $\Xi_w$; its members are at pairwise distance $>s$, and $D(V)$ lies inside their union. Suppose that its product does not vanish. Then each cube of each $X_i'$ holds a rough plaquette, so that the union equals $d(V)$; since every $X_i'$ is $s$-connected and no link connects two distinct members, the family must be $\{X_1,\dots,X_j\}$. When $D(V)$ is empty, this is the empty family, whose value is $1$.
\item[\textup{2.}] \emph{On and off.} We have $1_{Q^{+r}}(V)=1$ exactly when no plaquette in $Q^{+r}$ is rough, that is, exactly when $Q$ is not in $d(V)^{+r}$. For such $V$ the field on the bonds of $Q^{+r}$ lies in $\mathcal U_{\{Q\}}(\beta)$, because $|V(b)|=|V(\pi)|=1$ and $|\partial V(p)-1|=2\sin(\phi_p/2)\le\varepsilon$. We say that a cube $Q$ is on at $V$ if $1_{Q^{+r}}(V)=1$, and off at $V$ otherwise. Two distinct classes are at distance $\ge s+1=8r+1>2r$; hence the sets $X_i^{+r}$ are pairwise disjoint, their union is $d(V)^{+r}$, and a rough plaquette in $X_i^{+r}$ lies in $X_i$.
\item[\textup{3.}] \emph{Plaquette by plaquette.} At a tame $p$, \textup{(E2)} gives $q_p=\beta\phi_p^2/2-\beta s_p$. A tame $p$ in an on cube carries the factors $e^{-\beta\phi_p^2/2}e^{q_p+\beta\nu_ps_p}$. A tame $p$ in an off cube lies in exactly one $X_i^{+r}$ and carries the same two factors, the second of which now appears inside $w_\nu(X_i,V)$. A rough $p$ lies in exactly one $X_i$ and carries $e^{-\beta d_0\varepsilon\phi_p}e^{-\beta[(1-\nu_p)s_p-d_0\varepsilon\phi_p]}$. In all three cases the product of the factors is $e^{-\beta(1-\nu_p)s_p}$, and the identity follows on multiplying over all plaquettes.
\end{enumerate}

\emph{Step 5: \textup{(a2)}, the strong cap at $\kappa_1$ with a reserve.} We prove \textup{(a2)} for $|\nu|_\infty\le15.99c_\nu\varepsilon$, at every $V$ and for every $X\in\mathcal P$. Assume that $w_\nu(X,V)\ne0$, and set $\check\Sigma:=\beta\varepsilon\sum_{p\in D(V)\cap X}\phi_p$.

\begin{enumerate}
\item[\textup{(1)}] \emph{Rough plaquettes.} Here $\operatorname{Re}[(1-\nu_p)s_p]\ge(1-|\nu|_\infty)s_p$, and \textup{(E4)} gives $s_p\ge(\frac12-\varepsilon^2/24)\varepsilon\phi_p$. By \textup{(W3)}, \[(1-|\nu|_\infty)(\frac12-\varepsilon^2/24)\ge\frac12-\varepsilon^2/24-|\nu|_\infty/2\ge\frac12-(\varepsilon^2/24+8c_\nu\varepsilon)\ge\frac12-m_\Xi/4.\] Hence the rough plaquettes contribute to the real part of the exponent at most $-(\frac12-d_0-m_\Xi/4)\check\Sigma=-(c_\Xi+3m_\Xi/4)\check\Sigma$.
\item[\textup{(2)}] Each cube of $X$ holds a rough $p$, and for it $\beta\varepsilon\phi_p>\beta\varepsilon^2=p_0^2$; therefore $\check\Sigma\ge p_0^2\#X$.
\item[\textup{(3)}] \emph{Tame plaquettes.} We have $\#X^{+r}\le(2r+1)^3\#X$; each cube holds $3M^3$ plaquettes; $0\le q_p\le\beta\varepsilon^4/24=p_0^4/(24\beta)$ by \textup{(E2)}; and $|\beta\nu_ps_p|\le\beta|\nu|_\infty\varepsilon^2/2\le8c_\nu\varepsilon p_0^2$. The tame plaquettes therefore contribute $\le\#X(2r+1)^3[M^3p_0^4/(8\beta)+24M^3c_\nu\varepsilon p_0^2]$, where $24M^3c_\nu\varepsilon p_0^2=4E_ap_0^2\varepsilon=4c_a$.
\item[\textup{(4)}] From \textup{(2)} and \textup{(W4)}, $(m_\Xi/4)\check\Sigma\ge\#X(m_\Xi/4)p_0^2$, and the right side is at least $\kappa_1s\#X$ plus the bound found in \textup{(3)}. By \textup{(E5)}, $\kappa_1s\#X\ge\kappa_1(|X|_c-1)$.
\item[\textup{(5)}] \emph{Conclusion.} Adding up, $\log|w_\nu(X,V)|\le-(c_\Xi+m_\Xi/2)\check\Sigma-\kappa_1(|X|_c-1)$. For the reserve, \textup{(2)} and \textup{(W5)} give $e^{-(m_\Xi/2)\check\Sigma}\le e^{-(m_\Xi/2)p_0^2\#X}\le x_0^{-10\#X}$.
\end{enumerate}

Rough plaquettes in $X^{+r}$ that lie off $X$ occur in no sum that defines $w_\nu$. Therefore \textup{(1)} to \textup{(5)} are valid at \emph{every} $V$, and not only at those $V$ for which $X$ is a class of $d(V)$.

\emph{Step 6: \textup{(b)}, the data \textup{(SD)} at every level, the empty rows included.} Let $0\le k<K$ and $x=x_k$. We show: the coarse side $N_k:=NL^{K-k-1}$ is even, divisible by $M$, and satisfies $N_k\ge2Ms(x_k)$; the grids fit; $\ell_k\in \mathcal R(b_*,K_*,\lambda_*,1/9)$ has range $\le\varrho(x_k)$, and the pair $(\ell_k,(\ell_k)')$ has the four conclusions named in \textup{(SD)}; and $x_k\ge\max\{x_9^U,x_3,x_4\}$.

\emph{Floors.} We have $x_k\ge x_{K-1}=L\mathsf b\ge Lx_s$.

\emph{The side of the torus, \textup{(Tor)}.} We use the following lemma.

\begin{lemma}[Growth of $s$ under scaling]
If $x\ge e^2$ and $j\ge0$ is an integer, then $L^js(x)\ge s(L^jx)$.
\end{lemma}

To prove it, put $\alpha=20c_R/M$. As $e^t-1\ge t$, we have $(L^j-1)\log x\ge2(L^j-1)\ge2j\log L$, which gives $L^j\alpha\log x\ge\alpha\log(L^jx)$. The number $L^j\lceil\alpha\log x\rceil$ is an integer and is $\ge L^j\alpha\log x$; hence it is $\ge\lceil\alpha\log(L^jx)\rceil$. Since $r(x)=\lceil20c_R\log x/M\rceil$ and $s(x)=8r(x)$, this is the lemma.

At $x=L\mathsf b\ge e^2$ and $j=K-1-k$, the lemma gives $N_k=NL^j\ge2ML^js(L\mathsf b)\ge2Ms(x_k)$. Moreover $N_k$ is even, and $M\mid N_k$ holds because it holds for $N$.

\emph{Grids.} The relation $o_k=Lo_{k+1}-\mathsf h_L$ mod $M$ is the fit of item~1 of \S\ref{Y-G0-1}, under which the hat of each cell consists of exactly $L^3$ fine cubes; and $M$ divides the side $NL^{K-k}$ of $\mathbb T_k$. At $k=0$ this fixes the grid of $\mathbb T_0$. The weight $\exp\{-\beta\sum_p(1-\nu_p)s_p\}$ does not depend on that grid, and \textup{(a)} holds for every grid.

\emph{Rows.} \Cref{Y-A-Fr}, with its clauses \textup{(Fr0)}, \textup{(Fr2)}, \textup{(Fr3$'$)} and \textup{(Fr4)}, is valid for every $x_0$, for every $k\ge0$ with $x_{k+1}\ge x_\flat$, for every tower with $N\ge2$ even, and for the pair $(\ell_k,(\ell_k)')$. Here $x_{k+1}\ge\mathsf b\ge x_s$, and $x_s$ meets every demand that \textup{(O13)} makes on $x_\flat$, \textup{(X12)} included (item~(1), \textup{(O13)}, of \S\ref{Y-H0-1}). The statement therefore provides the following.

\begin{enumerate}
\item[\textup{(Fr0)}] $\ell_k\in \mathcal R(b_*,K_*,\lambda_*,1/9)$, of range $\le\varrho(x_k)$; for $\ell_0$, which is empty, the range is $0$.
\item[\textup{(Fr2)}] This is \textup{(Dec)}$(C_D,\theta_D)$ for the sixteen kernels, with $(C_D,\theta_D)$ as in \textup{(O6)}.
\item[\textup{(Fr3$'$)}] This is \textup{(N$'$)} with the numbers \eqref{eq:row-cut-norms} of \S\ref{Y-A-1}, at every cut and every mask.
\item[\textup{(Fr4)}] $\mathsf t_\rho\le\frac18m_\Xi$.
\end{enumerate}

The conditions \textup{(L1$^{\rm su}$)} and \textup{(L2$^{\rm su}$)} are \textup{(O4)}. These statements concern scalar kernels; the kernel $K\otimes1_{\mathfrak g}$ has the same norms as $K$, by item~5 of Step~15 of the proof of \Cref{Y-GB-ThB}; and at $\ell=\emptyset$ tensoring and dressing act on nothing. What we have obtained are the four conclusions of \textup{(SD)} (item~1 of \S\ref{Y-G0-1}, the last clause of \textup{(SD)}). At $k=0$, $\ell'^{\rm c}$ is $(\emptyset')^{\rm c}$ cut at $\varrho(\beta/L)$, and this is the definition of $\ell_1$.

\emph{The empty rows.} The data \textup{(SD)} of \S\ref{Y-C-1}, of item~B of \S\ref{Y-D-1}, of item~2 of \S\ref{Y-E-1}, of item~C of \S\ref{Y-F-1} and of item~1 of \S\ref{Y-G0-1} admit the empty rows; this is shown in Steps S0 and S6 of the proof of \Cref{Y-B-C1}. In \S\ref{Y-C-1}, items~C.I~1.1 and C.II~1.2 require $|\ell_y(p)|\le Ke^{-\lambda|y-p|}$, $K\ge0$; the inequalities $\gamma|\eta|^2\le q_\emptyset\le q_\ell\le b|\eta|^2$ are not demanded by these items but follow from \Cref{Y-C-CIIL64}\,(a) and \Cref{Y-C-CIL64}\,(a), in the case $K=0$ as in the case $K>0$; and Remark~(4) of item~I of \S\ref{Y-D-1} treats the case $\ell=\emptyset$. The data \textup{(SD)} of item~2 of \S\ref{Y-H-1} admit $\ell=\emptyset$ as well.

For item~2 of \S\ref{Y-H-1} this means that \Cref{Y-H-C9} itself, with its clauses \textup{(a)} to \textup{(e)}, applies at $\ell=\emptyset$, at every complex $u$ within 16 caps that has the weak cap.

\emph{Step 7: \textup{(c)}, the first step along the family.} By \textup{(a)}, $u_0(\nu)\in M_A(\beta)$, a subset of $M_A(\beta;p)$ for every $p\in\mathbb P$; by \textup{(b)}, \textup{(SD)} holds at $(\beta;\emptyset)$; and $\beta\ge x_9^U$. Hence \Cref{Y-H1-U1}\,(U) yields clauses \textup{(a)} to \textup{(e)} of \Cref{Y-H-C9} at every $p$, with a single map $F_\emptyset$ and a single number $\mathsf V'$. Combined with \textup{(a)} of \Cref{Y-H4-U4}, \Cref{Y-H-C9}\,(a) gives the formula of \textup{(c)} at every $V'$. \Cref{Y-H-C9}\,(b) shows that $E'^{\rm new}$ is the same for all members, hence free of $\nu$, and it supplies the three sums, the strong cap, and $F_\emptyset(u_0(\nu))\in M_{\max\{1,0.35+0.248A\}}(x_1)$.

\emph{Numerals.} If $|\nu|_\infty\le1.99c_\nu\varepsilon$, then $A\le2$ and $0.35+0.496=0.846\le1$. If $|\nu|_\infty\le6.6c_\nu\varepsilon$, then $A\le6.61$ and $0.35+1.63928=1.98928\le2$. Reality comes from \Cref{Y-H-C9}\,(d).

\emph{The Lipschitz bound.} Let $\check\delta=\tilde\nu-\nu$. Applied on $M_2$, \Cref{Y-H-C9}\,(e) gives $\mathsf D_{x_1}(u_1(\nu),u_1(\tilde\nu))\le0.4\mathsf D^{\rm in}_\beta(u_0(\nu),u_0(\tilde\nu))$. In the $a$-part, $E^a_\nu-E^a_{\tilde\nu}=-\beta\sum_{p\in Q}\check\delta_ps_p$, whose modulus is $\le6M^3|\check\delta|_\infty p_0^2=(|\check\delta|_\infty/(c_\nu\varepsilon))c_a$; the $h$-part is $0$. In the $w$-part, for $\tau\in[0,1]$ and complex $\zeta$ with $|\zeta|\,|\check\delta|_\infty\le14c_\nu\varepsilon$ we have $|\nu+(\tau+\zeta)\check\delta|_\infty\le15.99c_\nu\varepsilon$. On that disc, Cauchy's estimate together with \textup{(a2)} gives \[|(d/d\tau)w_{\nu+\tau\check\delta}(X,V)|\le(|\check\delta|_\infty/(14c_\nu\varepsilon))x_0^{-10}e^{-\kappa_1(|X|_c-1)}.\] Thus $\mathsf D^{\rm in}\le|\check\delta|_\infty/(c_\nu\varepsilon)$.

\emph{Holomorphy.} On the open set $|\nu|_\infty<15.99c_\nu\varepsilon$ of $\mathbb C^{\#\text{plaquettes}}$, the family $\nu\mapsto u_0(\nu)$ takes its values in $M_{16}(\beta;p)$, and by \textup{(a4)} each of its values is affine or entire. So \Cref{Y-H1-U1}\,(c+) applies.

\emph{Step 8: \textup{(d)} and \textup{(e)}, locality in $\nu$ and the vacuum number cell by cell.} For \textup{(d)}, \Cref{Y-H1-U1}\,(f1) states that $E'^{\rm old}[u](Y)$ and $v(Y;u)$ depend on $\{E^a(X):Y_0(X)\subset Y\}$ alone. The function $E^a_\nu$ is carried by single cubes $Q$. Single cubes are small ($1\le L/2-2$), and the cell holding $Q$ is $Y_0(\{Q\})$; by \textup{(a4)}, $E^a_\nu(\{Q\})$ depends on $\nu$ in $Q$. So both quantities depend only on $\nu$ under $Y$. By \Cref{Y-H1-U1}\,(f2), (f3), the quantities $E'^h(Y)$ and $w'(Y)$ depend on the members at those $X$ with $Y_0(X)\subset Y$. Here ${\rm cl}(X)\subset Y_0(X)^{+1}$ (item~1 of \S\ref{Y-G0-1}); since a cell is $L$ cubes wide, $X^{+r}$ lies under $Y_0(X)^{+1+\lceil r/L\rceil}$; and $w_\nu(X)$ depends on $\nu$ in $X^{+r}$. This is \Cref{Y-H1-U1}\,(f5) at the radius $r(\beta)$.

For \textup{(e)}, \Cref{Y-H1-U1}\,(g1) gives $V_1(\nu)=V_0^++\sum_Yv(Y;u_0(\nu))$, and $V_1^o:=V_0^+$ depends on no member, by \Cref{Y-H1-U1}\,(f1).

\emph{The bound.} Put $a_\nu:=0.00996+|\nu|_\infty/(c_\nu\varepsilon)<16$ and $u^\tau:=(\tau E^a_\nu,0,w_\nu)$. For $|\tau|<16/a_\nu$, an open disc containing $[0,1]$, we have $u^\tau\in M_{16}(\beta;p)$ with affine values; so by \Cref{Y-H1-U1}\,(c+) the map $\tau\mapsto v(Y;u^\tau)$ is holomorphic, and its derivative is $Dv(Y;u^\tau)[E^a_\nu]$. For $0\le\tau\le1$ we have $u^\tau\in M_A(\beta;p)$ with $A=\max\{1,\tau a_\nu\}<16$; on this range $v$ may be differentiated term by term, so that its derivative is the $Dv$ defined in item~1 of \S\ref{Y-H1-1} as the sum of the derivatives of the terms. This is part of \Cref{Y-H1-U1}\,(g3).

By the closing clause of \Cref{Y-H1-U1}\,(g1), $v(Y;u)=0$ if $E^a=0$, so that $v(Y;u^0)=0$. Hence, by \Cref{Y-H1-U1}\,(g3) applied with $\delta=E^a_\nu$, which has size $B\le a_\nu$, \[v_0(Y;\nu)=\int_0^1\{\sum_{Q:Y_0(\{Q\})=Y}G^\delta_{\{Q\}}[u^\tau](1)+r_\tau(Y)\}d\tau.\] Clause \textup{(g3)} holds on all of $M_{16}(\beta;p)$, since item~2 of \S\ref{Y-H1-1} has $1\le A\le16$.

The first sum vanishes when $\#Y\ge2$, and for $Y=\{\text{cell}\}$ it has $L^3$ terms. The field $V'=1$ lies in $\mathcal U_{\{\text{cell}\}}(x_1)$. Item~(ii) of Step~2 of the proof of \Cref{Y-H-C9}, applied at $V'=1$ with $X=\{Q\}$, with $Y=Y_0(\{Q\})$ and with the upper index $\tau$ of its domain $U^\tau$ of coarse fields equal to $1$, shows that the restriction of $U$ to $Q^{+r}$ lies in $\mathcal U_{\{Q\}}(\beta)$ on the set where $\chi_Q=1$; and there $|E^a_\nu|\le a_\nu c_a$ by \textup{(a1)}. Each term, being a mean of $\chi e^{\tau E^a}E^a$, therefore has modulus $\le e^{\tau a_\nu c_a}a_\nu c_a$. As $a_\nu c_a\le16c_a\le1/2$ according to \textup{(Y.1)} of item~5 of \S\ref{Y-H-1} (we take this item at $x=\beta$; it is in force there since $\beta\ge x_9^U$, by item~0, \textup{(O13)}\,(ii), of \S\ref{Y-H1-1}), and $(e^{1/2}-1)/(1/2)\le1.2975$, \[\int_0^1e^{\tau a_\nu c_a}a_\nu c_a\,d\tau=e^{a_\nu c_a}-1\le1.2975a_\nu c_a.\] The remainder, with the sum over $Y$ taken under the integral, is $\le109L^3\beta^{-1/4}a_\nu c_a\le0.00011L^3a_\nu c_a$, because the first inequality of \textup{(Y.6)}, in the same item, gives $\beta^{-1/4}\le10^{-6}$. Both parts together are $\le1.2977L^3a_\nu c_a\le1.30L^3a_\nu c_a$. Finally $a_\nu\le0.01+|\nu|_\infty/(c_\nu\varepsilon)\le A$, and $C_g=263L^3$.

\emph{Step 9: \textup{(f)}, the jets at level 0.} \emph{The $a$-part.} The expression \[E^a_{\nu(z)}(\{Q\},V)=\sum_{p\in Q}q_p+(x_0/\mathsf b)\sum_iz_i\sum_{p\in Q}f_i(x(p))s_p\] is affine in $z$. Therefore $\delta_i^a(0)(\{Q\},V)=(x_0/\mathsf b)\sum_{p\in Q}f_i(x(p))s_p$, while the $a$-part of $\delta_S(0)$ vanishes for $|S|\ge2$. Using $|s_p|<2\varepsilon^2$ we find \[|\delta_i^a(0)|\le3M^3(x_0/\mathsf b)2\varepsilon^2\|f_i\|_\infty=6M^3p_0^2\|f_i\|_\infty/\mathsf b=\mathrm{amp}_W\|f_i\|_\infty c_a.\]

Next, $s_p=\mathrm{Arc}(s_p)/2-q_p/x_0$, and on $\mathcal U_{\{Q\}}(\beta)$ we have $\mathrm{Arc}(s_p)=F_p\cdot F_p$ by \textup{(E3)}, since $2\varepsilon\le1/20$. Consequently \[\delta_i^a(0)=(x_0/2\mathsf b)\sum f_i(x(p))F_p\cdot F_p-(1/\mathsf b)\sum f_i(x(p))q_p.\] The first term equals $s_{x_0}[\mathbb D_0[f_i]/\mathsf b]$. Indeed, $\mathbb D_{0,y}[f](p,p')=f(x(p))$ if $p=p'$ is based at $y$, and $0$ otherwise; the path $\Gamma(y,\pi(p))$ is trivial; and the formula of entry~5 of item~(0) of \S\ref{Y-H0-1} gives $(x_0/2)\sum_{y\in Q}\sum_{p\ \text{based at}\ y}(f(x(p))/\mathsf b)F_p\cdot F_p$.

\emph{Sizes.} For $|s|<2\varepsilon^2$ and $\varepsilon\le1/40$ we have $|\mathrm{Arc}(s)|\le\mathrm{Arc}(|s|)=2|s|+2h(|s|)\le4.00084\varepsilon^2$; hence the first term is $\le3M^3(x_0/2\mathsf b)4.00084\varepsilon^2\|f_i\|_\infty=1.00021\,\mathrm{amp}_W\|f_i\|_\infty c_a$. By Step 3, the second term is $\le(\|f_i\|_\infty/\mathsf b)3M^3\,0.66709p_0^4/x_0\le0.00996c_a\|f_i\|_\infty/\mathsf b$. Just as in Step 3, both terms are holomorphic on $\mathcal U_{\{Q\}}(\beta)$, gauge invariant and $0$ at pure gauges; and they vanish unless $Q$ lies in $S(f_i)$.

\emph{The $h$-part.} Here $E^h=0$. \emph{The $w$-part.} We have $w_\nu(X,V)=w_0(X,V)\exp\{x_0\sum_{p\in\mathrm{Pl}(X,V)}\nu_ps_p\}$, where $\mathrm{Pl}(X,V)$ does not depend on $\nu$. One differentiation in each $z_i$, $i\in S$, at $z=0$ produces the product formula. One of its factors is $0$ when no plaquette in $X^{+r}$ lies in $\mathrm{supp}_0f_i$. For the size: if $\|f_i\|_\infty=0$ for some $i\in S$, the jet is $0$. Otherwise consider the torus $|z_i|=15.99/(\mathrm{amp}_W\|f_i\|_\infty)$, $i\in S$, with the other $z_j=0$. As the supports are disjoint, $|\nu(z)|_\infty\le15.99c_\nu\varepsilon$ there, so \textup{(a2)} bounds $|w_{\nu(z)}(X,V)|$ there; and the Cauchy estimate, applied to a derivative of first order in every variable $z_i$, $i\in S$, produces the factor $\prod_S(\mathrm{amp}_W\|f_i\|_\infty/15.99)$. When the supports overlap, \textup{(a)} to \textup{(e)} remain valid as stated in $\nu$, whereas in $z$ the radii shrink by the factor $n$.

\emph{Step 10: \textup{(g)}, the jets at level 1.} \emph{First half.} Clause \textup{(c)} and \Cref{Y-H1-U1}\,(c+) show that each value of $u_1(\nu(z))$, and each number $v_0(Y;\nu(z))$, is holomorphic on $\mathrm{Pd}_0^o$. \emph{Vanishing.} By \textup{(d)}, any value in which $\nu$ enters only off $\mathrm{supp}_0f_i$ is independent of $z_i$. \emph{Bounds.} If $\|f_i\|_\infty=0$ for some $i\in S$, everything vanishes, by the vanishing just shown. Otherwise, on $|z_i|=r_i$ ($i\in S$), with $z_j=0$ for the remaining $j$, we have $|\nu(z)|_\infty\le tc_\nu\varepsilon$, so that $u_0(\nu(z))\in M_{A_t}(\beta)$; this torus lies in $\mathrm{Pd}_0^o$ because $t\le15.98$. As $E'^{\rm new}$ is free of $z$, $\delta_S^a(1)=d_SE'^{\rm old}$. We write $d_S$ as the Cauchy integral over the torus $|z_i|=r_i$, and bring the supremum over $V'$, together with the weighted sum over $Y$, under the integral sign; at each $z$ the integrand is then bounded by \textup{(c)} or by \textup{(e)}, whichever applies to it. For $w'$, the strong cap at $x_1$ does not depend on $z$.

\emph{Second half.} We proceed in four parts.

\begin{enumerate}
\item[\textup{(1)}] By \Cref{Y-H1-U1}\,(f1), $E'^{\rm old}$ depends on $E^a$ alone, and $E^a_{\nu(z)}=E^a_0+\sum z_i\delta_i^a(0)$ is affine. Hence $\delta_i^a(1)=(d/dt)E'^{\rm old}[(E^a_0+t\delta_i^a(0),0,w_0)]$ at $t=0$, and this is $DA_\emptyset(u_0(0))[\delta_i^a(0)]$ in the sense of \Cref{Y-H1-U1}\,(h). The map $DA$ is additive in its argument, since by \Cref{Y-H1-U1}\,(c+) the two-parameter family is holomorphic.
\item[\textup{(2)}] By \textup{(f)}, \[\delta_i^a(1)=DA_\emptyset(u_0(0))[s_{x_0}[D]]+DA_\emptyset(u_0(0))[\mathrm{rem}_{0,i}],\qquad D:=\mathbb D_0[f_i]/\mathsf b.\]
\item[\textup{(3)}] We apply \Cref{Y-H3-U3} at $x=\beta\ge x_3$, with \textup{(SD)} $=(\beta;\emptyset)$ (its data include the empty rows), with $u=u_0(0)\in M_1(\beta)$, a subset of $M_8(\beta;p)$, and with $D$ complex, symmetric (diagonal) and of range $0$. It gives $DA_\emptyset(u_0(0))[s_{x_0}[D]]=s_{x_1}[D']+n[u_0(0);D]$, where $D'=\mathrm{cut}_{\varrho(x_1)}(T_\emptyset D)=\mathbb D_1[f_i]/\mathsf b$ by the definition in item~1 of \S\ref{Y-H2-1} (with $\ell_0=\emptyset$) and by linearity. \Cref{Y-H3-U3}\,(a), with \Cref{Y-H2-U2}\,(U2.1) at $k=0$, gives $|n|_{a,x_1}^{{\rm cs},(\tilde\kappa)}\le u_3(x_1)\|D\|_{(\lambda_\flat)}\le u_3(x_1)K_5\|f_i\|_\infty/\mathsf b$.
\item[\textup{(4)}] The remainder $\mathrm{rem}_{0,i}$ belongs to $T^a(\beta)$ and has size $B\le0.01\|f_i\|_\infty/\mathsf b$ at every $p$, as it lives on single cubes. By \Cref{Y-H1-U1}\,(h3), the weighted cell sum of $DA_\emptyset(u_0(0))[\mathrm{rem}_{0,i}]$ is $\le\vartheta_aBc_a(x_1)$, and a cell sum bounds the weighted supremum. Since $\vartheta_a=0.2431L^{-1/4}$, we get $0.01\,\vartheta_a\le0.00244$.
\end{enumerate}



\emph{Step 11: \textup{(h)}, \textup{(i)} and the corollary.} \emph{Induction on $k$ over $\mathrm{Pd}_0^o$.} For $k=0$, \textup{(a)} gives $u_0(z)\in M_{16}(x_0)$ with the strong cap, and its values are affine or entire in $z$. For the passage $k\to k+1$: by \textup{(b)}, \textup{(SD)} holds at $(x_k;\ell_k)$, $x_k\ge x_9^U$, and $u_k(z)\in M_{A^{(k)}}(x_k)$ with $A^{(k)}\le16$. \Cref{Y-H1-U1}\,(U), together with \Cref{Y-H-C9}\,(b), then gives $u_{k+1}(z)\in M_{A^{(k+1)}}(x_{k+1})$ with $A^{(k+1)}=\max\{1,0.35+0.248A^{(k)}\}$, and the strong cap. This yields $4.318$, then $1.420864\le1.421$, then at most $0.702408<1$, so that the maximum is $1$. By \Cref{Y-H1-U1}\,(c+), $V_{k+1}(z)$ and all values are holomorphic on the open set $\mathrm{Pd}_0^o$. On $\mathrm{Pd}_0$ one has $A^{(0)}=2$, then $0.846\le1$, then $0.598\le1$.

\emph{The densities.} By \textup{(a)}, $\rho_0^z=\rho[u_0(z);\emptyset]$. Suppose that $\rho_k^z=e^c\rho[u_k(z);\ell_k]$ for a complex number $c$. As the step is linear in $\rho$ (item~2 of \S\ref{Y-H-1}), \Cref{Y-H-C9}\,(a) gives, at every field, \[\rho_{k+1}^z=e^c(\rho[u_k(z);\ell_k])'=e^{c+V_{k+1}(z)}\rho[u_{k+1}(z);\ell_{k+1}].\]

\emph{Normalisation.} The kernel $k_{ax}(\cdot,h)$ is a probability density in $V'(c)$, by item~2 of \S\ref{Y-J-1}, which describes the step: there $k_s(g,h)=z(s)^{-1}e^{-s(1-\cos\operatorname{ang}(gh^{-1}))}$ with $z(s)=\int_Ge^{-s(1-\cos\operatorname{ang}g)}dg$, and the right invariance of Haar measure gives $\int k_s(g,h)dg=1$. The function $\rho_k^z$ is bounded and Borel on a compact group carrying its Haar probability, so Fubini's theorem yields $\int\rho_{k+1}^z=\int\rho_k^z$. Hence $Z_K(z)=\int\rho_K^z$.

\emph{Real $z$ and real $f_i$, or $z=0$.} In this case $\nu(z)$ is real; so $u_0(z)$ is real, and $\rho_0^z>0$ by \textup{(a3)}. Applying \Cref{Y-H-C9}\,(d) step by step, we find that $u_k(z)$ and $V_k(z)$ are real. From $\rho_0^z>0$ we get $\rho_k^z>0$, and then $\rho[u_k(z);\ell_k]>0$.

\emph{Clause \textup{(i)}.} It follows from Steps 1 to 10: among the constants of \Cref{Y-H1-U1}, \Cref{Y-H2-U2} and \Cref{Y-H3-U3} that enter, none depends on $K$, $N$, $\mathsf b$, on a member, or on any part of the rows other than the numbers $b_*$, $K_*$, $\lambda_*$ and the range.

\emph{The corollary.} At $z=0$, $\bar u_0=u_0(0)$ is real, lies in $M_1(\beta)$ ($A=1$) and has the strong cap, and $\rho[\bar u_0;\emptyset]>0$; so $\rho_0$ is a point of $C^{\rm su}_\flat(\beta;\emptyset)$. For $k\ge1$, $\bar u_k$ is real, lies in $M_1(x_k)$ and has the strong cap, and $\rho_k=e^{\sum V_j(0)}\rho[\bar u_k;\ell_k]>0$, where the exponent $\sum V_j(0)$ is real. Further, $u_0(0)=u_W$, since the two formulas agree term by term ($q_p=\beta h(s_p)$). Finally, \Cref{Y-J-ThA}\,(ii) asks that a system be fixed, and any system will do.
\end{proof}

\subsection{Holomorphic contraction}\label{Y-H5-1}

In this section we work on the complex space $\mathfrak M(x)$ of members. For bounded maps that are holomorphic along lines we give Cauchy bounds on their derivatives, a Lipschitz bound and holomorphy along families, in every seminorm that is a supremum of finite linear forms, as the weighted suprema and the cell sums are; for families we give the Cauchy bound on every Taylor coefficient at the origin and the Schwarz bound on the values of a family that vanishes at the origin, we show that the Taylor coefficients of a bounded family of members are again members, and we write the coefficient of $z^S$ of a composed map as a sum over the partitions of $S$. We then follow the step map $F_\ell$ on families within $B\le2.6$ caps: their images lie within one cap and obey the strong cap, the factor $\vartheta[B]$ leads from the weak distance of the inputs to the strong distance of the images, and we give the channels of the first and second derivatives, the local vacuum numbers and the reflection $u\mapsto u^*$. Finally, a lemma of eight clauses \textup{(B1)} to \textup{(B8)} gives the properties of an affine map $\mathbb B_D$ that absorbs into the members a complex symmetric kernel $D$ of the density, added to the kernels $\mathsf E_y$ of the rows, and we carry out one step for two families, one of them over a kernel $D(z)$.

\begin{remark} The statements of this section follow from \S\ref{Y-I-1}, from standard facts, and from \Cref{Y-H1-U1}, which asserts \Cref{Y-H-C9} at every pair of rates and adds further clauses. \end{remark}

\emph{Item 0: the letters.}

\begin{enumerate}
\item[\textup{(1)}] \emph{Letters and data kept.} From items~2 and~3 of \S\ref{Y-I-1} we keep the lattice $\mathbb Z^3/NL\mathbb Z^3$ with \textup{(Tor)}, the cubes and the cells, the neighbourhoods $X^{\approx}=X^{+r(x)}$ and $X^{\approx\approx}=X^{+s(x)}$, and the numbers $\varrho(x)$, $r(x)$ and $s(x)=8r(x)$.

From the same items we take the group $G=\operatorname{SU}(2)$ with the operator norm $|\cdot|$ and the form $Y\cdot Y'=\frac12\operatorname{tr}YY'$, the fields $V$ with values in $G$, and the fields $\mathsf W$ with values in $\operatorname{SL}(2,\mathbb C)$.

These items also give us $\phi_p$, $F_p$, $\zeta_p$, the rough set $D(V)$, the indicator $1_Y(V)$ of the property that $V$ is tame on $Y$, the path $\Gamma(y,z)$, the transport $T(y,z)$ and the mask $m_y(V)$.

We keep as well the rows $\ell$ with their kernels $\mathsf E_y=\ell_y^{\mathsf T}\ell_y$, the norm $\|D\|_{(\lambda)}$ of kernels, the energy $\mathrm{En}_\ell(V)$ and the step $\rho\to\rho'$.

We keep the domains $\mathcal U_X(x)$ with \textup{(u1)}, \textup{(u2)} and \textup{(u3)}, and the complex linear space $\mathfrak M(x)$ of those triples $u=(E^a,E^h,w)$ that obey the clauses without a size and have $\|u\|<\infty$ (item~3 of \S\ref{Y-I-1}).

We keep the density $\rho[u;\ell]$, the sum $\Xi_w$ and the families of \textup{(c$\Xi$)}, the class $C^{\rm su}_\flat(x;\ell)$, and the property of a triple $u$ of being real.

From item~1 of \S\ref{Y-I-1} we keep $p_0(x)$ and $\varepsilon=\varepsilon(x)=p_0(x)x^{-1/2}$, for which $x\varepsilon^2=p_0^2$ is \textup{(O8)}; the caps $c_a(x)=E_ap_0(x)^3x^{-1/2}$ and $c_h(x)=e^{-c_hp_0(x)^2}$; and the numbers $c_h$ of \textup{(O11)} and $E_a$ of \textup{(O12)}. None of these depends on a rate.

From \S\ref{Y-H-1} we keep \textup{(SD)}, the map $F_\ell(u)=(E'^a,E'^h,w')$ with $E'^a=E'^{\rm new}+E'^{\rm old}[u]$, and $\mathsf V'[u]$.

From \S\ref{Y-H0-1} we keep $\kappa$ and $\kappa_1$, the set $\mathbb P=[\kappa,\kappa_1]\times\{c_h,c_h/2\}$ of pairs, and the numbers $\vartheta_w=0.248$, $\vartheta_a=0.2431L^{-1/4}$, $\vartheta[B]=\max\{16\vartheta_w,1\}/(16-B)=3.968/(16-B)$, $C_g=263L^3$, $\mathsf g_\flat=\mathsf g(\lambda_\flat)$ and $c_Q^{\rm su}=306M^3\mathsf g_\flat^2$.

From that section we keep, too, $\eta_h(x)=e^{-(c_h/2)(p_0(x)^2-p_0(x/L)^2)}$, $\eta_w(x)=e^{-c_\Xi p_0(x)^2}$, the floor $x_9^U$, the member $s_x[D]$ and the polydisc $\mathrm{Pd}(r)=\{z\in\mathbb C^n:|z_i|<r_i\}$.

From \S\ref{Y-H1-1} we keep $v(Y;u)$, $V_0^+$, $T^a(x)$ and $DA(u)[\delta]$, where $A=E'^{\rm old}$.

In $x'$, $p_0'$, $V'$, $\mathsf W'$, $E'$, $w'$ and $\ell'$ the prime denotes the object at the level $x'=x/L$.

At the level $x'$ the definitions of items~2 and~3 of \S\ref{Y-I-1} are taken on the coarse torus $\mathbb T'=\mathbb Z^3/N\mathbb Z^3$, with cells in the place of cubes; the definition ``The level at $x'$'' in item~3 of \S\ref{Y-I-1} carries them over, and the remark directly after it shows that this rests on three conditions alone, namely that $N$ is even, that $M\mid N$ and that $N\ge16Mr(x')$.

Lems.~\ref{Y-H5-LH}, \ref{Y-H5-LS} and~\ref{Y-H5-LJ} depend on three features of a level, that is, of $x$ with its torus, and on nothing else of it: the index sets are finite, each $\mathcal U_X$ is an open set of a complex manifold, and the weights of items~0\,(2) and~0\,(3) are free of $z$.

Hence the three lemmas hold at $x'$ on $\mathbb T'$ just as at $x$; in the same way they hold for $\mathfrak M(\mathsf b)$ on the last torus $\mathbb T_K$ of \S\ref{Y-H6-1}, which is the torus $\mathbb T'$ of the step at $x=L\mathsf b$.

\item[\textup{(2)}] \emph{Norms at a pair.} Let $p=(\tilde\kappa,\tilde c)\in\mathbb P$ and $\delta=(\delta^a,\delta^h,\delta^w)\in\mathfrak M(x)$. We set \[|\delta|_a:=\sup_X\sup_{\mathcal U_X(x)}|\delta^a(X,\cdot)|\,e^{\tilde\kappa(\#X-1)}/c_a(x),\] \[|\delta|_h:=\sup_X\sup|\delta^h(X,\cdot)|\,e^{\tilde\kappa(\#X-1)}e^{\tilde cp_0(x)^2},\] \[|\delta|_w^{\mathrm{wk}}:=\sup_{X,V}|\delta^w(X,V)|\,e^{\tilde\kappa(|X|_c-1)},\] and we let $|\delta|_w^{\mathrm{st}}$ be the same supremum as $|\delta|_w^{\mathrm{wk}}$, with the weight \[\exp\{c_\Xi x\varepsilon\sum_{p\in D(V)\cap X}\phi_p+\tilde\kappa(|X|_c-1)\}\] in the place of $e^{\tilde\kappa(|X|_c-1)}$. We also set $\|\delta\|:=\max\{|\delta|_a,|\delta|_h,|\delta|_w^{\mathrm{wk}}\}$ and $\|\delta\|^{\mathrm{st}}:=\max\{|\delta|_a,|\delta|_h,|\delta|_w^{\mathrm{st}}\}$.

We write $M_A(x;p)=\{\|u\|\le A\}$, $\mathsf D^{\mathrm{in}}_x(u_1,u_2)=\|u_1-u_2\|$ and $\mathsf D_x(u_1,u_2)=\|u_1-u_2\|^{\mathrm{st}}$. A member obeys the \emph{strong cap} if $|w|_w^{\mathrm{st}}\le1$. On the torus the index sets are finite, so that $\mathfrak M(x)$ is one and the same set for all $p$. Where the context does not make clear which pair is meant, we write an upper index $(p)$ or $(\tilde\kappa)$.

\item[\textup{(3)}] \emph{Cell sums at $x'$.} For a family $e(Y,\cdot)$ we set \[|e|_a^{\Sigma}:=\sup_{\square}\sum_{Y\ni\square}e^{\tilde\kappa(\#Y-1)}\sup_{\mathcal U_Y(x')}|e(Y,\cdot)|/c_a(x'),\] where $\square$ runs over the cells. We define $|e|_h^{\Sigma}$ in the same way, with the supremum over those $G$-valued fields on the bonds of $Y^{\approx}$ that are tame there, and with $e^{-\tilde cp_0(x')^2}$ in the place of $c_a(x')$. For numbers $v=(v(Y))$ we set $|v|^{\Sigma}:=\sup_{\square}\sum_{Y\ni\square}e^{\tilde\kappa(\#Y-1)}|v(Y)|$. Since every $Y$ holds a cell, $|e|_a\le|e|_a^{\Sigma}$ and $|e|_h\le|e|_h^{\Sigma}$.

The quantity $\sup_{\square}\sum_{Y\ni\square}e^{\tilde\kappa(\#Y-1)}\sup|e(Y,\cdot)|$ is the $|e|_{\mathrm{cs}}$ of item~1 of \S\ref{Y-H1-1} and the ${\rm CS}_{\tilde\kappa}(e)$ of item~1 of \S\ref{Y-H3-1}, and neither of the two is divided by a cap. Thus $|e|_a^{\Sigma}=|e|_{\mathrm{cs}}/c_a(x')$, and this is the $|e|^{{\rm cs},(\tilde\kappa)}_{a,x'}$ of Sections~\ref{Y-H3-1} and~\ref{Y-H4-1}.

\item[\textup{(4)}] \emph{Values, admissible seminorms, domination of coordinates.} A \emph{value} of $u$ is one of the numbers $E^a(X,\mathsf W)$, $E^h(X,V)$, $w(X,V)$. Let $I$ be an index set. A seminorm $N$ on a linear subspace of $\mathbb C^I$ is \emph{admissible} if $N(a)=\sup_{\mathfrak l\in\mathfrak L}|\mathfrak l(a)|$ for a set $\mathfrak L$ of linear forms $\mathfrak l(a)=\sum_ic_ia_i$, each with finitely many $c_i\ne0$. Every weighted supremum norm of item~0\,(2) is admissible: each of its forms $\mathfrak l$ is the weight times one value. Every cell sum of item~0\,(3) is admissible, with the forms $\sum_{Y\ni\square}c_Ye^{\tilde\kappa(\#Y-1)}e(Y,\mathsf W_Y)/c_a(x')$ over the cells $\square$, the arguments $\mathsf W_Y$ and the phases $|c_Y|=1$; each such sum is finite.

Each of these seminorms $N$ \emph{dominates each coordinate}: for every $i$ there is a $c_i>0$ with $c_i|a_i|\le N(a)$. Indeed, a weighted value is itself a form of $\mathfrak L$; and in a cell sum every $Y$ holds a cell and the phases are free, so that the sum of the moduli, and each of its terms, is at most $N(a)$. Fix such a representing family $\mathfrak L$ and use the same supremum to define $N$ on all of $\mathbb C^I$, allowing the value $+\infty$. Its finite domain is a linear subspace, and $N$ is a seminorm there. A bound on a coordinatewise derivative or Taylor coefficient also asserts finiteness of this extended seminorm. Membership of Taylor coefficients in $\mathfrak M(x)$ is proved separately in \Cref{Y-H5-LS}\,\textup{(S3)}. An admissible seminorm can fail to dominate: $N(a)=|a_1|$ on $\mathbb C^2$ is admissible and does not dominate the coordinate $a_2$.

\item[\textup{(5)}] \emph{Families, Taylor coefficients, $u^{\mathrm{pol}}$.} Let $n\ge1$. For $\mu\in\mathbb N^n$, where $\mathbb N^n$ is the set of the $n$-tuples of integers $\ge0$, we write $r^{-\mu}=\prod r_i^{-\mu_i}$, and for a non-empty subset $S$ of $\{1,\dots,n\}$ we write $z^S=\prod_{i\in S}z_i$. A \emph{family} at $x$ on an open set $\Omega\subset\mathbb C^n$ within $B$ caps at $p$ is a map $z\mapsto u(z)\in M_B(x;p)$ all of whose values are holomorphic in $z$. For families, $\mathsf D_{x;\Omega}$ and $\mathsf D^{\mathrm{in}}_{x;\Omega}$ denote the suprema over $z\in\Omega$. We write $u_\mu$ for the Taylor coefficients at $0$, defined value by value, and $u_S$ for the coefficient of $z^S$, and we set $u^{\mathrm{pol}}(z):=u(0)+\sum_Sz^Su_S$.

\item[\textup{(6)}] \emph{Holomorphy along lines, derivatives.} A function $\Phi$ on a norm-open set $O$ of $\mathfrak M(x)$ is \emph{holomorphic along lines} if for all $u,g\in\mathfrak M(x)$ the map $s\mapsto\Phi(u+sg)$ is holomorphic on $\{s\in\mathbb C:u+sg\in O\}$. For such a $\Phi$ we set $D^j\Phi(u)[g_1,\dots,g_j]:=\partial_{s_1}\cdots\partial_{s_j}\Phi(u+\sum s_ig_i)$ at $s=0$.

\item[\textup{(7)}] \emph{The reflection.} For a group element $g$ we put $g^{\#}:=(g^\dagger)^{-1}$, for a field $\mathsf W$ we put $\mathsf W^{\#}(b):=\mathsf W(b)^{\#}$, and for a member $u$ we put \[u^*:=(\overline{E^a(X,\mathsf W^{\#})},\overline{E^h(X,V)},\overline{w(X,V)}).\]

\item[\textup{(8)}] \emph{Letters with several senses.}

The letter $p$ is the pair $(\tilde\kappa,\tilde c)$ when a norm, a ball $M_A(x;p)$ or a family is taken at it, and also in $p^\circ$ and in the upper index $(p)$; it is a plaquette when it indexes a quantity of the field or of a kernel, is an argument of such a quantity or is the variable of a sum, as in $\phi_p$, $D_y(p,p')$ and $\sum_{p\in D(V)\cap X}$, and this role tells the two senses apart.

The letter $A$ is a number of caps; in $DA$ and $D^2A$ it is the map $E'^{\rm old}$; and in $\log(1+A)$ (Step~7 of this section) it is a $2\times2$ matrix.

The letter $r$ denotes the radii in $\mathrm{Pd}(r)$ and in $r^{-\mu}$, and it is $r(x)$ in $(2r+1)^3$ and in $X^{+r}$.
\end{enumerate}

\begin{lemma}[Cauchy bounds on a ball of members]\label{Y-H5-LH} Let $\|\cdot\|$ be the norm of item~0\,(2) at one pair $p$, and let $0\le B<A_{\#}\le16$. Let $\Phi=(\Phi_i)_{i\in I}$ be functions on $\{u\in\mathfrak M(x):\|u\|<A_{\#}\}$, each of them holomorphic along lines and bounded on that ball, and let $N$ be an admissible seminorm with $N(\Phi(u))\le\mathsf m$ for all $\|u\|<A_{\#}$.

\begin{enumerate}
\item[\textup{(a)}] For $\|u\|\le B$ and $j\ge1$, every $D^j\Phi_i(u)$ is symmetric and complex multilinear, and \[N(D^j\Phi(u)[g_1,\dots,g_j])\le\mathsf m\,(j/(A_{\#}-B))^j\prod\|g_i\|.\]

\item[\textup{(b)}] For $\|u_1\|,\|u_2\|\le B$ we have $N(\Phi(u_1)-\Phi(u_2))\le\mathsf m\,\|u_1-u_2\|/(A_{\#}-B)$.

\item[\textup{(c)}] For every family $u$ within $B$ caps at $p$, every $\Phi_i(u(z))$ is holomorphic in $z$.

\item[\textup{(a$^\circ$)}] Off the centre: let $u_0\in\mathfrak M(x)$, let $t>0$ be any number, and let $\Phi$ be a function on $\{u\in\mathfrak M(x):\|u-u_0\|<t\}$ that is holomorphic along lines and obeys $|\Phi|\le\mathsf m$ there. Then for every $j\ge1$ the form $D^j\Phi(u_0)$ is symmetric and complex multilinear, and \[|D^j\Phi(u_0)[g_1,\dots,g_j]|\le\mathsf m\,(j/t)^j\prod\|g_i\|.\]
\end{enumerate}
\end{lemma}

\begin{remark}
The boundedness of each $\Phi_i$ follows from $N(\Phi(u))\le\mathsf m$ wherever $N$ dominates each coordinate (item~0\,(4)); so it follows for every norm of item~0\,(2) and for every cell sum of item~0\,(3).

For $\Phi$ a value of $F_\ell$, of $\mathsf V'$ or of $v(Y;\cdot)$, \Cref{Y-H5-LH}\,(c) is \Cref{Y-H1-U1}\,(c+).
\end{remark}

\begin{lemma}[Cauchy and Schwarz for families; coefficients are members]\label{Y-H5-LS} Let $a(z)=(a_i(z))_{i\in I}$, with every $a_i$ holomorphic on $\mathrm{Pd}(r)$, and let $N$ be an admissible seminorm with $N(a(z))\le m$ on $\mathrm{Pd}(r)$.

\begin{enumerate}
\item[\textup{(S1)}] Every Taylor coefficient $a_\mu$ at $0$ obeys $N(a_\mu)\le m\,r^{-\mu}$.

\item[\textup{(S2)}] If $a(0)=0$, then $N(a(z))\le m\max_i|z_i|/r_i$.

\item[\textup{(S3)}] If $g:\mathrm{Pd}(r)\to\mathfrak M(x)$ has holomorphic values and $\sup_z\|g(z)\|<\infty$, then every $g_\mu$ lies in $\mathfrak M(x)$.
\end{enumerate}

Hence \textup{(S1)} and \textup{(S2)} hold for such a $g$ in $\|\cdot\|$, in $\|\cdot\|^{\mathrm{st}}$, in each of the suprema of item~0\,(2), and in the cell sums of item~0\,(3).
\end{lemma}

\begin{lemma}[The coefficient of $z^S$ of a composed map]\label{Y-H5-LJ} Let $u_0\in\mathfrak M(x)$ and $t>0$, and let $\Phi$ be bounded and holomorphic along lines on $\{\|u-u_0\|<t\}$. Let $u$ be a map from a neighbourhood of $0$ in $\mathbb C^n$ to $\mathfrak M(x)$ that is bounded in norm, has holomorphic values and satisfies $u(0)=u_0$. Then $\Phi(u(z))$ is holomorphic near $0$, and its coefficient of $z^S$ is \[\sum_{\pi\in\mathrm{Part}(S)}D^{|\pi|}\Phi(u_0)[u_{S_1},\dots,u_{S_{|\pi|}}],\] where $\pi=\{S_1,\dots,S_{|\pi|}\}$ runs over the partitions of $S$ into non-empty blocks.

So this coefficient depends on $u(0)$ and on the $u_{S'}$, with $S'$ a non-empty subset of $S$, alone, and $\Phi(u(z))$ and $\Phi(u^{\mathrm{pol}}(z))$ have the same coefficient of $z^S$ for every $S$.

In particular, in the setting of \Cref{Y-H5-Af} ($x\ge x_9^U$, \textup{(SD)} at $x$, $p\in\mathbb P$), and for a family $u$ within $B<16$ caps at $p$ near $0$, the assertions above hold, with $t=16-B$, for every value of $F_\ell(u(z))$, for $\mathsf V'[u(z)]$ and for $v(Y;u(z))$.

They hold also for $I[u]:=\int\rho[u;\ell](V)\,dV$, at any $x>e$, for any rows $\ell$ of range $\le\varrho(x)$, any $u_0$ and any $t$. The lattice for $I[u]$ is that of item~2 of \S\ref{Y-I-1} at $x$, or a torus $\mathbb Z^3/N\mathbb Z^3$ with $N$ even, $M\mid N$ and $N\ge16Mr(x)$, on which the definitions of items~2 and~3 of \S\ref{Y-I-1} are set up by the definition ``The level at $x'$'' in item~3 of that section; this torus is the $\mathbb T'$ of that definition, with $x$ in the place of its $x'$.

Thus they hold for $I[u]$ on $\mathbb T_K$ at $x=\mathsf b$, where $N\ge2Ms(L\mathsf b)\ge16Mr(\mathsf b)$.

The claim for $I[u]$ rests on three facts. First, at each $V$ the density $\rho[u;\ell](V)$ is $e^{-x\mathrm{En}_\ell(V)}$ times the exponential of a finite sum of values of $E^a$ and $E^h$, times a polynomial in finitely many values of $w$, and so it is entire along lines; the factor $e^{-x\mathrm{En}_\ell(V)}$ is free of the members and is at most $1$, since $\mathrm{En}_\ell\ge0$, the energy of item~2 of \S\ref{Y-I-1} being a sum of numbers $S_\varepsilon(\phi_p)\ge0$ and of squares. Second, on $\{\|u\|<R\}$ the density is bounded, uniformly in $V$, by \[\exp\{R(c_a(x)+e^{-\tilde cp_0^2})\#\mathcal D\}(1+R)^{\#\mathcal P},\] where, on the finite torus, $\#\mathcal D$ and $\#\mathcal P$ are the numbers of elements of the families $\mathcal D$ and $\mathcal P$ of item~2 of \S\ref{Y-I-1}. Third, by the theorems of Fubini and Morera the integral is holomorphic along lines.

Thus the two hypotheses of the lemma on $\Phi$, boundedness on the ball and holomorphy along lines on it, hold for four kinds of $\Phi$: for a value of $F_\ell$, for $\mathsf V'[\cdot]$, for $v(Y;\cdot)$ and for $I[\cdot]$.
\end{lemma}

\begin{theorem}[The contraction for pairs, for families; family form of \Cref{Y-I-ThA}]\label{Y-H5-Af} Let $x\ge x_9^U$, in the setting \textup{(SD)} at $x$; let $p=(\tilde\kappa,\tilde c)\in\mathbb P$ and $1\le B\le2.6$; and let $u_1$ and $u_2$ be families on $\Omega$ within $B$ caps at $p$, with the weak cap. All norms are taken at $p$.

\begin{enumerate}
\item[\textup{(i)}] Each $F_\ell(u_i(\cdot))$ is a family at $x'$ in $M_1(x';p)$ that obeys the strong cap, and $|E'^a|_a^{\Sigma}\le0.35+0.248B\le0.9948$; moreover $\mathsf V'[u_i(\cdot)]$ and every $v(Y;u_i(\cdot))$ are holomorphic.

\item[\textup{(ii)}] We have $\mathsf D_{x';\Omega}(F_\ell u_1,F_\ell u_2)\le\vartheta[B]\,\mathsf D^{\mathrm{in}}_{x;\Omega}(u_1,u_2)$. Here $\vartheta[1]\le0.26454$, $\vartheta[2]\le0.28343$, $\vartheta[2.455]\le0.29295$ and $\vartheta[2.6]\le0.29612$, by Step~4 of this section. For the three members the factors are $16\vartheta_w/(16-B)$, $1/(16-B)$ and $1/(16-B)$, and the first two of them hold in $|\cdot|_a^{\Sigma}$ and $|\cdot|_h^{\Sigma}$. Besides, $|E'^a[u_1]-E'^a[u_2]|_a^{\Sigma}\le\vartheta_a|E_1^a-E_2^a|_a$ at each $z$.

\item[\textup{(iii)}] On $\Omega=\mathrm{Pd}(r)$, for every $\mu$, \[\|(F_\ell u_1-F_\ell u_2)_\mu\|^{\mathrm{st}}\le\vartheta[B]\,r^{-\mu}\mathsf D^{\mathrm{in}}_{x;\mathrm{Pd}(r)}(u_1,u_2).\]

\item[\textup{(iv)}] \emph{Channels.} Let $\|u\|\le B$ and $g,g_1,g_2\in\mathfrak M(x)$, and put $g^{w\circ}(X,V):=g^w(X,V)\,1\{D(V)\cap X\ne\emptyset\}$. The inputs are measured in caps at $x$ and the outputs in caps at $x'$, the $a$- and $h$-members in the cell sums and $w'$ in $|\cdot|_w^{\mathrm{st}}$. For the first derivatives, the channel from $a$ to $a$ obeys $|DA(u)[g^a]|_a^{\Sigma}\le\vartheta_a|g^a|_a$, and the channels from $h$ and from $w$ to $a$ are $0$. The channels from $a$ to $h$ and $w$ are at most $|g^a|_a/(16-B)$. The channels from $w$ to $h$ and $w$ are at most $|g^{w\circ}|_w^{\mathrm{wk}}/(16-B)$, and $|g^{w\circ}|_w^{\mathrm{wk}}\le\eta_w|g^w|_w^{\mathrm{st}}$, the input being measured in the strong norm. At $\tilde c=c_h$, the channels from $h$ to $h$ and $w$ are at most $\eta_h|g^h|_h/(16-B)$.

For the second derivatives with output $a$, $|D^2A(u)[g_1^a,g_2^a]|_a^{\Sigma}\le16\vartheta_a\,(4/(16-B)^2)\,|g_1^a|_a|g_2^a|_a$, and the second derivative is $0$ as soon as an entry is an $h$- or a $w$-member. Hence on $M_B$ \[|(DA(u_1)-DA(u_2))[g^a]|_a^{\Sigma}\le(64\vartheta_a/(16-B)^2)|E_1^a-E_2^a|_a|g^a|_a,\] and also $|(DA(u_1)-DA(u_2))[g^a]|_a^{\Sigma}\le(\vartheta_a/(16-B))|E_1^a-E_2^a|_a|g^a|_a$; of these two bounds the second is the smaller. For the second derivatives with outputs $h$ and $w$, the bound is $4/(16-B)^2$ times the product of the sizes of the two entries, where an $a$-entry is measured by $|\cdot|_a$, a $w$-entry by $|g^{w\circ}|_w^{\mathrm{wk}}\le\eta_w|g^w|_w^{\mathrm{st}}$, and an $h$-entry, at $\tilde c=c_h$, by $\eta_h|g^h|_h$. An entry can also be left whole, as a mixed triple $g$, and is then measured by $\|g\|$ at $p$; a mixed entry that is split is measured by the sum $|g^a|_a+\eta_h|g^h|_h+|g^{w\circ}|_w^{\mathrm{wk}}$. The factor $4/(16-B)^2$ is $1/56.25$ at $B=1$ and $1/49$ at $B=2$.

\item[\textup{(v)}] At each $z$, \[|v(\cdot;u_1)-v(\cdot;u_2)|^{\Sigma}\le16C_gc_a(x)|E_1^a-E_2^a|_a/(16-B)\le16C_gc_a(x)\mathsf D^{\mathrm{in}}_x(u_1,u_2)/(16-B),\] and at $B=2$ the factor $16C_gc_a(x)/(16-B)$ is at most $300.6L^3c_a$. On $\mathrm{Pd}(r)$, \[|(v(\cdot;u_1(z)))_\mu|^{\Sigma}\le r^{-\mu}C_gBc_a(x)\] and \[|(v(\cdot;u_1(z))-v(\cdot;u_2(z)))_\mu|^{\Sigma}\le r^{-\mu}16C_gc_a(x)\mathsf D^{\mathrm{in}}_{x;\mathrm{Pd}(r)}/(16-B).\]

\item[\textup{(vi)}] \emph{Reflection.} The map $u\mapsto u^*$ is an antilinear involution of $\mathfrak M(x)$, isometric in every norm of item~0\,(2), and its fixed points are real. For $\|u\|<16$ we have $\mathsf V'[u^*]=\overline{\mathsf V'[u]}$, and $E'^h[u^*]=\overline{E'^h[u]}$ and $w'[u^*]=\overline{w'[u]}$ at every argument; further, $E'^a[u^*](Y,\mathsf W')=\overline{E'^a[u](Y,\mathsf W'^{\#})}$ for each $\mathsf W'$ in $\mathcal U_Y(x')$ with values in $G$, and also on each connected component of $\mathcal U_Y(x')$ that contains a field with values in $G$. Hence $\rho[F_\ell(u^*);\ell'^{\rm c}]=\overline{\rho[F_\ell(u);\ell'^{\rm c}]}$.
\end{enumerate}

Clauses \textup{(ii)} to \textup{(v)} hold for every $1\le B<16$; among the hypotheses of the theorem, the bound $B\le2.6$ is used for \textup{(i)} alone.
\end{theorem}

\begin{lemma}[Absorbing into the members a complex kernel added to those of the rows; complex form of \Cref{Y-I-LB}]\label{Y-H5-Bc} Let $x>e$. We work on the lattice of item~2 of \S\ref{Y-I-1} and in the setting of the five conditions of item~5 of \S\ref{Y-I-1}: \textup{(M5r)} $M\ge\rho_0$, \textup{(X15)} $c_\Xi p_0(x)^2\ge1$, \textup{(F1)} $\varepsilon(x)\le1/40$, \textup{(F2)} $\varrho(x)+2\le Mr(x)$ and \textup{(F3)} $(2r(x)+1)^3c_a(x)\le34$. In the setting of the table of item~(1) of \S\ref{Y-H0-1} these five conditions hold at every $x\ge Lx_s$, by \textup{(O7)}, by \textup{(O13)} and by Step~0\,(Fl) of \S\ref{Y-I-1}, which is stated for $x\ge x_9$, the floor of \Cref{Y-H-C9}, evaluated at $\kappa_1$; every $x\ge Lx_s$ is at least $x_9^U$, since $Lx_s\ge x_9^U$ is a clause of \textup{(O13)}. As $x_9^U\ge x_9$ by item~0\,(O13) of \S\ref{Y-H1-1}, we have $x\ge Lx_s\ge x_9^U\ge x_9$. Let $\ell$ be rows of range $\le\varrho(x)$. Let $D=(D_y)$ be complex kernels with $D_y(p,p')=D_y(p',p)$ and with $D_y(p,p')=0$ unless $|y-p|,|y-p'|\le\varrho(x)$, whatever their sign, whether or not they are Gram forms and whether or not they are covariant. Put $\delta:=\|D\|_{(\lambda_\flat)}$ and $\delta_\Delta=\delta_\Delta(D):=c_Q^{\rm su}p_0^2\delta/c_a(x)$, and let $\delta_\Delta\le1$ and $A>0$. The norms of item~0\,(2) are taken at any rate $\tilde\kappa\ge0$ and any $\tilde c>0$.

We set \[q_y^D(V):=-(x/2)\,m_y(V)\sum_{p,p'}D_y(p,p')\langle T(y,\pi(p))\zeta_p,T(y,\pi(p'))\zeta_{p'}\rangle\] and, for a cube $\square$, $q_\square^D:=\sum_{y\in\square}q_y^D$. We set $E^\Delta[D]:=-s_x[D]$, that is, \[E^\Delta[D](\{\square\},\mathsf W)=-(x/2)\sum_{y\in\square}\sum_{p,p'}D_y(p,p')\,(\operatorname{Ad}_{\mathsf W(\Gamma(y,\pi(p)))}F_p)\cdot(\operatorname{Ad}_{\mathsf W(\Gamma(y,\pi(p')))}F_{p'}),\] with $F_p=-i\log\partial\mathsf W(p)$, and $E^\Delta[D]$ is $0$ at $\#X\ge2$. A cube $\square$ is \emph{on} at $V$ if $V$ is tame on $\square^{\approx}$, and \emph{off} otherwise; $p_*(\square,V)$ is the first plaquette of $D(V)$ in $\square^{\approx}$, in a fixed order. We set $Q_D(X,V):=\sum\{q_\square^D(V):\square\text{ off},\ p_*(\square,V)\in X\}$ and $\rho[u;\ell;D]:=\rho[u;\ell]\exp\{\sum_\square q_\square^D\}$; if $D=\tilde{\mathsf E}-\mathsf E$ for rows $\tilde\ell$ of range $\le\varrho(x)$, then $\rho[u;\ell;D]=\rho[u;\tilde\ell]$.

Then for every complex $\tilde u=(\tilde E^a,\tilde E^h,\tilde w)\in\mathfrak M(x)$ with $|\tilde E^a|_a,|\tilde E^h|_h,|\tilde w|_w^{\mathrm{st}}\le A$, the triple $\mathbb B_D(\tilde u):=(\tilde E^a+E^\Delta[D],\tilde E^h,w^\flat)$, where $w^\flat:=\tilde we^{Q_D}$, lies in $M_{A+\delta_\Delta}(x;p)$, and \textup{(B1)} to \textup{(B8)} hold. In them $\mathfrak n:=\#(D(V)\cap X)$, a number determined by $(X,V)$.

\begin{enumerate}
\item[\textup{(B1)}] $\rho[\tilde u;\ell;D]=\rho[\mathbb B_D\tilde u;\ell]$ at every $V$.

\item[\textup{(B2)}] The member $E^\Delta[D]$ is complex linear in $D$, lives on $\#X=1$ and obeys those of the clauses without a size that concern $E^a$; it obeys $\sup_{\mathcal U_{\{\square\}}(x)}|E^\Delta[D](\{\square\},\cdot)|\le\delta_\Delta c_a(x)$; it depends neither on $\ell$ nor on $\tilde u$; and $E^\Delta[D](\{\square\},V)=q_\square^D(V)$ at every $V$ tame on $\square^{\approx}$. Clause \textup{(B2)} holds for every $D$ that satisfies the hypotheses that the lemma makes on $D$ other than $\delta_\Delta\le1$, whatever $\delta_\Delta(D)$ is: both sides of its bound are homogeneous of degree $1$ in $D$. Likewise, whatever $\delta_\Delta(D)$ is, \[|Q_D(X,V)|\le((2r(x)+1)^3c_a(x)/68)\,\delta_\Delta(D)\,\mathfrak n\le(\delta_\Delta(D)/2)\,\mathfrak n,\] by Step~8 of this section together with \textup{(F3)}.

\item[\textup{(B3)}] $|w^\flat(X,V)|\le Ae^{-(c_\Xi p_0^2-1/2)\mathfrak n}e^{-\tilde\kappa(|X|_c-1)}$ and $|w^\flat-\tilde w|(X,V)\le(A/e)\,\delta_\Delta e^{-\tilde\kappa(|X|_c-1)}$; and $e^{Q_D(X,\cdot)}$ depends on $V$ through its values on the bonds of $X^{+(2r(x)+1)}$, which lie inside $X^{\approx\approx}$.

\item[\textup{(B4)}] $\mathsf D^{\mathrm{in}}_x(\tilde u,\mathbb B_D\tilde u)\le\max\{1,A/e\}\,\delta_\Delta$.

\item[\textup{(B5)}] The map $\mathbb B_D$ is affine in $\tilde u$. For a family $\tilde u(z)$ whose values are holomorphic and whose three sizes are $\le A$ at each $z$, and for a constant $D$, the $a$- and $h$-coefficients of $z^\mu$, $\mu\ne0$, are unchanged, and $w^\flat_\mu=\tilde w_\mu e^{Q_D}$, which is $0$ wherever $\tilde w_\mu$ is $0$; moreover $|w^\flat_\mu|_w^{\mathrm{wk}}\le|\tilde w_\mu|_w^{\mathrm{st}}$, and $|(w^\flat_\mu)^\circ|_w^{\mathrm{wk}}\le e^{1/2}\eta_w|\tilde w_\mu|_w^{\mathrm{st}}$ for the part on $D(V)\cap X\ne\emptyset$. If $z\mapsto D(z)$ has holomorphic entries and at each $z$ the kernel $D(z)$ satisfies the hypotheses that the lemma makes on $D$, with $\sup_z\delta_\Delta(D(z))\le1$, then $\mathbb B_{D(z)}\tilde u(z)$ is a family within $A+1$ caps.

\item[\textup{(B6)}] $\mathsf D^{\mathrm{in}}_x(\mathbb B_{D_1}\tilde u_1,\mathbb B_{D_2}\tilde u_2)\le \mathsf D_x(\tilde u_1,\tilde u_2)+\max\{1,A/e\}\,\delta_\Delta(D_1-D_2)$.

\item[\textup{(B7)}] The derivative $(d/dt)\mathbb B_{tD}\tilde u$ at $t=0$ is $(E^\Delta[D],0,\tilde wQ_D)$, and \[|\tilde wQ_D|(X,V)\exp\{(c_\Xi/2)x\varepsilon\sum_{p\in D(V)\cap X}\phi_p+\tilde\kappa(|X|_c-1)\}\le(A/e)\,\delta_\Delta,\] with half of the strong weight. With the whole strong weight instead of its half, the left-hand side is bounded by $A(\delta_\Delta/2)\mathfrak n$. For a lower bound let $x\ge x_9$, the floor of \Cref{Y-H-C9}, evaluated at $\kappa_1$ (so every $x\ge x_9^U$ qualifies, by item~0\,(O13) of \S\ref{Y-H1-1}), and let $A\ge1$; take the rows, the field and the set of Step~10 of \S\ref{Y-I-1}, the proof of \Cref{Y-I-InF}, with the kernel $D=\tilde{\mathsf E}-\mathsf E$ of that step, for which $D_y(p,p')=-t^2$ at $p=p'=p_y$ and $0$ otherwise, and the member $\tilde u=(0,0,\tilde w)$ with \[\tilde w(X,V):=\exp\{-c_\Xi x\varepsilon\sum_{p\in D(V)\cap X}\phi_p-\tilde\kappa(|X|_c-1)\},\] which is the member of that proposition with the rate $\tilde\kappa$ of the norm in the place of its $\kappa$. In $D_y(p,p')=-t^2$ the letter $t$ denotes not the parameter of $\mathbb B_{tD}$ but the entry $\ell_y(p_y)=t$ of the rows of that step. For this kernel, $\delta=t^2$, so that $t^2=\delta_\Delta c_a(x)/(c_Q^{\rm su}p_0^2)$. Then, with the whole strong weight, \[|\tilde wQ_D|(X,V)\exp\{c_\Xi x\varepsilon\sum_{p\in D(V)\cap X}\phi_p+\tilde\kappa(|X|_c-1)\}\] equals $Q_D(X,V)$ and is $\ge\mathfrak n'\delta_\Delta c_a(x)/(8c_Q^{\rm su})$, where $\mathfrak n'$ is the number $n$ of the sites $z_1,\dots,z_n$ of that step, which grows with $N$.

\item[\textup{(B8)}] $E^\Delta[\overline D]=(E^\Delta[D])^*$ on all of $\mathcal U_{\{\square\}}(x)$, $\mathbb B_{\overline D}(\tilde u^*)=(\mathbb B_D\tilde u)^*$, and $\mathbb B_D\tilde u$ is real if $D$ and $\tilde u$ are.
\end{enumerate}

At $A=1$, with $D=\tilde{\mathsf E}-\mathsf E$, with $\tilde u$ real and with the norms taken at the pair $p=(\kappa_1,c_h)$, the membership of $\mathbb B_D(\tilde u)$ in $M_{A+\delta_\Delta}(x;p)$, \textup{(B1)} to \textup{(B4)} and \textup{(B8)} give \Cref{Y-I-LB}, since $1/e\le\frac12$ in \textup{(B3)}.
\end{lemma}

\begin{theorem}[One step for two families, one of them over a kernel $D(z)$; family form of \Cref{Y-I-ThC}]\label{Y-H5-Cf} Let $x\ge x_9^U$, in the setting \textup{(SD)} at $x$ for $(\ell,\ell')$ and in the setting of the conditions of item~5 of \S\ref{Y-I-1}, and let $p\in\mathbb P$. Let $D(z)$, $z\in\Omega$, be kernels as in \Cref{Y-H5-Bc}\,(B5), with $\delta_\Delta:=\sup_z\delta_\Delta(D(z))\le1$. Let $u_1$ be a family on $\Omega$ within $B_1$ caps, and let $\tilde u_2$ be a family with $|\tilde E_2^a|_a,|\tilde E_2^h|_h,|\tilde w_2|_w^{\mathrm{st}}\le B_2$; let $B:=\max\{B_1,B_2+1\}\le2.6$, and put $u_2^\flat(z):=\mathbb B_{D(z)}\tilde u_2(z)$.

\begin{enumerate}
\item[\textup{(C1)}] At every $z$ and $V'$, $(\rho[\tilde u_2;\ell;D])'=e^{\mathsf V'[u_2^\flat]}\rho[F_\ell(u_2^\flat);\ell'^{\rm c}]$.

\item[\textup{(C2)}] $F_\ell(u_1(\cdot))$ and $F_\ell(u_2^\flat(\cdot))$ are families in $M_1(x';p)$ that obey the strong cap, and $\mathsf V'[u_2^\flat(\cdot)]$ is holomorphic.

\item[\textup{(C3)}] $\mathsf D_{x';\Omega}(F_\ell u_1,F_\ell u_2^\flat)\le\vartheta[B]\,(\mathsf D^{\mathrm{in}}_{x;\Omega}(u_1,\tilde u_2)+\delta_\Delta)$.

\item[\textup{(C4)}] At one point $z$, with $D=\tilde{\mathsf E}-\mathsf E$, with $u_1$ and $\tilde u_2$ real, in $M_1$ and obeying the strong cap, with $\rho>0$ and with $p=(\kappa_1,c_h)$, the theorem becomes \Cref{Y-I-ThC}, with $\vartheta[2]\le0.28343\le\vartheta_{10}$.
\end{enumerate}
\end{theorem}

The steps of the proofs in this section are numbered consecutively from Step~1 to Step~9. \begin{proof}[Proof of \Cref{Y-H5-LH}\,(a), (b)] \emph{Step 1: Cauchy's estimate on the ball, in an admissible seminorm.}

\begin{enumerate} \item[\textup{(a)}] We begin with the scalar case, in which $I$ is a single point and $N=|\cdot|$. Fix $u$ with $\|u\|\le B$ and directions $g_1,\dots,g_j$. If $g_i=0$ for some $i$, both sides of the bound are $0$; we may therefore take every $g_i$ to be non-zero.

The set $O:=\{s\in\mathbb C^j:\|u+\sum s_ig_i\|<A_{\#}\}$ is open and convex. On $O$ the function $\Psi(s):=\Phi(u+\sum s_ig_i)$ is bounded by $\mathsf m$, and it is holomorphic in each $s_i$ when the other variables are held fixed, because $s_i\mapsto u+\sum s_ig_i$ then runs along an affine line. By Osgood's lemma, a bounded function that is holomorphic in each variable separately is holomorphic; hence $\Psi$ is holomorphic on $O$. Consequently the mixed derivative of $\Psi$ at $0$, which is $D^j\Phi(u)[g_1,\dots,g_j]$, exists and is symmetric in the directions.

This derivative is linear in $g_1$. Indeed, for $g_1=cg+c'g'$ we apply the same argument to the $j+1$ directions $g,g',g_2,\dots,g_j$, and the chain rule gives $\partial_{s_1}=c\,\partial_s+c'\partial_{s'}$.

For the bound we put $R_i:=(A_{\#}-B)/(j\|g_i\|)$. Where $|s_i|<R_i$ for every $i$, we have $\|u+\sum s_ig_i\|<B+j(A_{\#}-B)/j=A_{\#}$, so this polydisc lies in $O$. Cauchy's estimate on the polydisc with radii $R_i'<R_i$, followed by the limit $R_i'\to R_i$, gives \[|\partial_1\cdots\partial_j\Psi(0)|\le\mathsf m\prod R_i^{-1}=\mathsf m\,(j/(A_{\#}-B))^j\prod\|g_i\|.\]

We now let $N$ be general. Each $\Phi_i$ is bounded on the ball, so the scalar case applies to $\Phi_i$, with the bound of $\Phi_i$ in the place of $\mathsf m$: the derivative $D^j\Phi_i(u)$ exists and is symmetric and complex multilinear. Let $\mathfrak l\in\mathfrak L$. The function $\mathfrak l(\Phi(\cdot))=\sum c_i\Phi_i$ is a finite sum, hence holomorphic along lines, and $|\mathfrak l(\Phi(u))|\le N(\Phi(u))\le\mathsf m$ on the ball; moreover $D^j$ commutes with the finite sum. The scalar case, applied to $\mathfrak l(\Phi(\cdot))$, therefore gives \[|\mathfrak l(D^j\Phi(u)[g_1,\dots,g_j])|\le\mathsf m\,(j/(A_{\#}-B))^j\prod\|g_i\|,\] and we take the supremum over $\mathfrak L$.

For a cell sum this argument amounts to the following: each derivative is written as a Cauchy integral, and the sum over $Y$ is taken inside the integral, together with the suprema over the fields. Applying Cauchy's estimate to each $Y$ separately and summing afterwards would instead give $\sum_Y\sup_s$, and this can be larger than $\sup_s\sum_Y$. For example, for $v_Y(z)=((1+ze^{-2\pi iY/8})/2)^{12}$, $Y=1,\dots,8$, on $|z|=1$ the first quantity is $8$, and the second is the value of the sum $\sum_Y|v_Y(z)|$, which is constant on $|z|=1$ and equals $231/128\le1.805$.

\item[\textup{(b)}] The closed ball is convex, so $u_t:=u_2+t(u_1-u_2)$ satisfies $\|u_t\|\le B$ for $0\le t\le1$. Let $\mathfrak l\in\mathfrak L$. The function $t\mapsto\mathfrak l(\Phi(u_t))$ is $C^1$ with derivative $\mathfrak l(D\Phi(u_t)[u_1-u_2])$, and by \textup{(a)} at $j=1$ this derivative is bounded by $\mathsf m\|u_1-u_2\|/(A_{\#}-B)$. We integrate over $[0,1]$ and take the supremum over $\mathfrak L$. \end{enumerate}

Step~2 of the proof of \Cref{Y-I-ThA} uses the bound $\|u_z\|\le B+|z|$ on the segment as well, which leads to its constant $16-3B=10$ at $B=2$. Here the norm on the segment is $\le B$, and this gives $16-B$. \end{proof}

\begin{proof}[Proof of \Cref{Y-H5-LS}] \emph{Step 2: the estimates of Cauchy and of Schwarz; the coefficients are members.}

\begin{enumerate} \item[\textup{(S1)}] Let $\mathfrak l\in\mathfrak L$ and $h(z):=\mathfrak l(a(z))$. Then $h$ is a finite sum of holomorphic functions, $|h|\le m$ on $\mathrm{Pd}(r)$, and the Taylor coefficient of $h$ of index $\mu$ is $\mathfrak l(a_\mu)$. Cauchy's estimate on the polydisc gives $|\mathfrak l(a_\mu)|\le m\,r^{-\mu}$, and we take the supremum over $\mathfrak L$.

\item[\textup{(S2)}] Now $h(0)=0$. For $z\ne0$ let $t:=\max|z_i|/r_i<1$. The function $c\mapsto h(cz/t)$ is holomorphic on $|c|<1$, bounded by $m$ and equal to $0$ at $0$, so the Schwarz lemma at $c=t$ gives $|h(z)|\le mt$.

\item[\textup{(S3)}] Let $0<r_i'<r_i$. The coefficients are given, at each value, by \[g_\mu(\text{value})=(2\pi i)^{-n}\oint_{|z_i|=r_i'}g(z)(\text{value})\prod dz_i/z_i^{\mu_i+1}.\]

We first treat $E^a$. The set $\mathcal U_X(x)$ is open in the complex manifold $\operatorname{SL}(2,\mathbb C)^{\text{bonds of }X^{\approx}}$. In a chart, the map $(z,\mathsf W)\mapsto E^a_z(X,\mathsf W)$ is bounded and separately holomorphic, so it is holomorphic by Osgood's lemma. Integrating a jointly holomorphic function over a compact torus in $z$ gives a function that is holomorphic in $\mathsf W$.

Invariance under $G$-valued gauge transformations, the value $0$ at the pure gauges, the dependence sets $X^{\approx}$ and $X^{\approx\approx}$, and the domains of $E^h$ and $w$ are linear conditions that hold at each $z$; hence they hold for the integral.

The coefficients are Borel. Indeed, the integrand is continuous in $z$, so the integral is a pointwise limit of Riemann sums, and these are finite combinations of Borel functions of $V$.

Finally $\|g_\mu\|<\infty$ by \textup{(S1)}, and the norms and cell sums that the lemma names are admissible by item~0\,(4) of \S\ref{Y-H5-1}. \end{enumerate} \end{proof}

\begin{proof}[Proof of \Cref{Y-H5-LH}\,(c), (a$^\circ$) and of \Cref{Y-H5-LJ}] \emph{Step 3: one power series.}

\begin{enumerate} \item[\textup{(a$^\circ$)}] Let $|\Phi|\le\mathsf m$ on $\{\|u-u_0\|<t\}$. We apply Step~1 with the centre $u_0$ in place of $0$, with $B=0$ and with $A_{\#}=t$; this is possible because the proof of Step~1 uses only that the ball is open and convex. It gives \[|D^m\Phi(u_0)[g_1,\dots,g_m]|\le\mathsf m\,(m/t)^m\prod\|g_i\|,\] so $D^m\Phi(u_0)$ is multilinear and continuous for the norm.

\item[\textup{(1)}] Let $\|h\|<t$. The function $s\mapsto\Phi(u_0+sh)$ is holomorphic on $|s|<t/\|h\|$, and its Taylor series at $0$, evaluated at $s=1$, is \[\Phi(u_0+h)=\sum_{m\ge0}(1/m!)D^m\Phi(u_0)[h,\dots,h].\] Here the $m$-th derivative at $0$ equals $D^m\Phi(u_0)[h,\dots,h]$, by the chain rule applied to the function of $m$ variables from Step~1 with all $g_i=h$.

\item[\textup{(2)}] Let $\|u(z)\|\le B'$ on $\mathrm{Pd}(q)$. \Cref{Y-H5-LS} gives $u_\mu\in\mathfrak M(x)$ and $\|u_\mu\|\le B'q^{-\mu}$.

Hence \[\varsigma(z):=\sum_{\mu\ne0}|z^\mu|\,\|u_\mu\|\le B'[\prod(1-|z_i|/q_i)^{-1}-1],\] which tends to $0$ with $z$.

Value by value we have $h(z):=u(z)-u_0=\sum_{\mu\ne0}z^\mu u_\mu$. The norm is a weighted supremum of values, so the norm of a tail of this series is at most the sum of the norms of its terms; thus the series converges in norm.

\item[\textup{(3)}] On the set $\{\varsigma(z)<t/e\}$, multilinearity and continuity give \[D^m\Phi(u_0)[h(z),\dots,h(z)]=\sum_{(\mu^1,\dots,\mu^m),\ \text{all}\ne0}z^{\mu^1+\dots+\mu^m}D^m\Phi(u_0)[u_{\mu^1},\dots,u_{\mu^m}],\] and the sum of the moduli on the right is at most $\mathsf m\,(m/t)^m\varsigma^m$; moreover $m^m/m!\le e^m$.

Therefore $\Phi(u(z))$ is the sum of a multiple power series whose sum of moduli is at most $\mathsf m\sum_m(e\varsigma/t)^m<\infty$. So $\Phi(u(z))$ is holomorphic near $0$, and its coefficients are obtained by collecting terms.

\item[\textup{(4)}] The coefficient of $z^S$ collects the tuples with $\mu^1+\dots+\mu^m=1_S$ and all $\mu^i\ne0$. For such a tuple $\mu^i=1_{S_i}$, where $(S_1,\dots,S_m)$ is an ordered partition of $S$ into non-empty blocks.

By symmetry the $m!$ orderings of one partition give equal terms, which cancels the factor $1/m!$; this is the formula. It involves only $u_0$ and the $u_{S'}$ with $S'\subset S$, and $u$ and $u^{\mathrm{pol}}$ have these in common.

\item[\textup{(c)}] We prove \Cref{Y-H5-LH}\,(c). At $z_0\in\Omega$ we take $\Phi=\Phi_i$ and let $\mathsf m$ be the bound of $\Phi_i$ on $\{\|u\|<A_{\#}\}$. This ball contains $\{\|u-u_0\|<t\}$ for $u_0=u(z_0)$ and $t=A_{\#}-B$. We apply \textup{(3)} to the map $z\mapsto u(z_0+z)$ with $B'=B$, and obtain holomorphy near $z_0$. \end{enumerate}

It remains to prove the sentence of \Cref{Y-H5-LJ} that is introduced by the words \emph{in particular}. By \Cref{Y-H1-U1}\,(U), \Cref{Y-H-C9}\,(b) holds at $p$, and so a value of $F_\ell$ is bounded on $M_{16}(x;p)$. By \Cref{Y-H1-U1}\,(g2) we have $|v(Y;u)|\le16C_gc_a(x)$, and by \Cref{Y-H1-U1}\,(g1), (g2) we have $|\mathsf V'[u]|\le|V_0^+|+16C_gc_a(x)$ times the number of cells. All these maps are holomorphic along lines, by \Cref{Y-H1-U1}\,(c+) applied to affine families.

For $I[u]$ the two hypotheses are shown in the statement of \Cref{Y-H5-LJ} itself, by the three facts named there. \end{proof}

\begin{proof}[Proof of \Cref{Y-H5-Af}] \emph{Step 4: the image, the contraction and the coefficients.} In this step we prove clauses \textup{(i)} to \textup{(iii)} of \Cref{Y-H5-Af}.

\begin{enumerate} \item[\textup{(i)}] By \Cref{Y-H1-U1}\,(U), \Cref{Y-H-C9}\,(b) holds at $p$, and we use it at $A=B$. It gives $|E'^{\rm new}|_a^{\Sigma}\le0.35$ and $|E'^{\rm old}[u]|_a^{\Sigma}\le0.248B$, so that \[|E'^a|_a^{\Sigma}\le0.35+0.248\cdot2.6=0.9948\le1.\]

Moreover $|E'^h|_h^{\Sigma}\le1$; the member $w'$ carries the strong cap at $x'$ with the rate $\tilde\kappa$, and hence also the weak cap; and the clauses without a size hold at $x'$.

Since a sum of non-negative terms bounds each of its terms, $F_\ell(u(z))\in M_1(x';p)$. All values, $\mathsf V'$ and $v(Y;\cdot)$ are holomorphic in $z$ by \Cref{Y-H1-U1}\,(c+).

\item[\textup{(ii)}] Fix $z$; then $\|u_i(z)\|\le B$. We apply \Cref{Y-H5-LH}\,(b), with $A_{\#}=16$, to three maps $\Phi$ on $\{\|u\|<16\}$, each of which is holomorphic along lines by \Cref{Y-H1-U1}\,(c+): the map $\Phi_a=(E'^{\rm old}[u](Y,\mathsf W'))_{Y,\mathsf W'}$, with $N=|\cdot|_a^{\Sigma}$ and $\mathsf m=16\vartheta_w=3.968$ by \Cref{Y-H-C9}\,(b) at $A=16$; the map $\Phi_h=(E'^h(Y,V'))$, with $N=|\cdot|_h^{\Sigma}$ and $\mathsf m=1$; and the map $\Phi_w=(w'(Y,V'))$, with $N=|\cdot|_w^{\mathrm{st}}$ at $x'$ and $\mathsf m=1$.

The member $E'^{\rm new}$ is the same for all $u$, so it cancels in $E'^a[u_1]-E'^a[u_2]$.

Hence the three members of $F_\ell(u_1)-F_\ell(u_2)$ are at most $3.968/(16-B)$, $1/(16-B)$ and $1/(16-B)$ times $\|u_1-u_2\|$, and the largest of these factors is $\vartheta[B]$. Here we use $|\cdot|_a\le|\cdot|_a^{\Sigma}$ and $|\cdot|_h\le|\cdot|_h^{\Sigma}$. Finally we take the supremum over $z$.

We turn to the further bound in \textup{(ii)}. By \Cref{Y-H1-U1}\,(f1), $E'^{\rm old}$ depends on $E^a$ alone. Along $u_t=u_2+t(u_1-u_2)$ we have $(d/dt)E'^{\rm old}[u_t]=DA(u_t)[E_1^a-E_2^a]$ by \Cref{Y-H1-U1}\,(h): the series $DA$, defined in \S\ref{Y-H1-1} termwise as a sum of derivatives, coincides with this derivative wherever $\|u_t\|<16$.

The difference $E_1^a-E_2^a$ lies in $T^a(x)$, since its norm is finite and so each of its values is bounded. At each $t$, \Cref{Y-H1-U1}\,(h3) gives the bound $\vartheta_a|E_1^a-E_2^a|_a$ in the cell sum. As in the proof of \textup{(b)} in Step~1, we apply each linear form of the cell sum, integrate over $0\le t\le1$, and take the supremum over these forms.

\item[\textup{(iii)}] Let $g(z):=F_\ell(u_1(z))-F_\ell(u_2(z))$. This map takes $\mathrm{Pd}(r)$ to $\mathfrak M(x')$: by \Cref{Y-H1-U1}\,(U), \Cref{Y-H-C9}\,(b) holds at $p$, and at $A=B<16$ it yields finite sizes together with all clauses without a size at $x'$, namely $F_\ell(u)\in M_{\max\{1,0.35+0.248B\}}(x';p)$.

The values of $g$ are holomorphic by \Cref{Y-H1-U1}\,(c+). Neither of these two facts depends on $B\le2.6$; for $B\le2.6$ they are \textup{(i)}. By \textup{(ii)}, $\|g(z)\|^{\mathrm{st}}\le\vartheta[B]\mathsf D^{\mathrm{in}}_{x;\mathrm{Pd}(r)}$, and we apply \Cref{Y-H5-LS}\,(S1), (S3) in the norm $\|\cdot\|^{\mathrm{st}}$ at $x'$. \end{enumerate}

The contraction factors satisfy $\vartheta[1]\le0.26454$, $\vartheta[2]\le0.28343$, $\vartheta[2.455]\le0.29295$ and $\vartheta[2.6]\le0.29612$. Also $(1-0.35)/0.248>2.6$, so the radius $2.6$ is admitted.

The numbers $0.3968$ of \Cref{Y-I-ThA}\,(ii) and $0.4$ of \Cref{Y-H-C9}\,(e) are upper bounds of $\vartheta[2]$.

\emph{Step 5: the channels.} We prove the bounds on the first and second derivatives in \Cref{Y-H5-Af}\,(iv).

By Step~1, $DF_\ell(u)[g]$ is linear in $g$, so we may split $g=(g^a,0,0)+(0,g^h,0)+(0,0,g^w)$; likewise $D^2$ is bilinear.

For the channel $a\to a$ we use \Cref{Y-H1-U1}\,(h3) with $\delta=g^a\in T^a(x)$, whose size is $|g^a|_a$ caps at $p$, in the setting \textup{(L3su)}.

For the channels $h,w\to a$ we use \Cref{Y-H1-U1}\,(f1): $E'^{\rm new}$ depends on no member and $E'^{\rm old}[u]$ on $E^a$ alone, so $E'^a[u+s(0,g^h,g^w)]$ is constant in $s$.

For the channels $a\to h,w$ we use \Cref{Y-H5-LH}\,(a) with $j=1$ and $A_{\#}=16$, for the maps $\Phi_h$ and $\Phi_w$ of Step~4, with $\mathsf m=1$, in the direction $g=(g^a,0,0)$, for which $\|g\|=|g^a|_a$.

For the channels $w\to h,w$ we use \Cref{Y-H1-U1}\,(f4): $F_\ell$ and $\mathsf V'$ use $w(X,V)$ only at those $V$ for which $X$ holds a plaquette of $D(V)$. Hence $F_\ell(u)=F_\ell(u^\circ)$, where $w^\circ(X,V):=w(X,V)\,1\{D(V)\cap X\ne\emptyset\}$ and $u^\circ:=(E^a,E^h,w^\circ)$.

Moreover $g^\circ:=(0,0,g^{w\circ})\in\mathfrak M(x)$, because the indicator is Borel and gauge invariant, being a function of $\phi_p$ at the plaquettes of $X$.

Since $(u+\sum s_ig_i)^\circ=(u+\sum s_ig_i^\circ)^\circ$, we have $F_\ell(u+\sum s_ig_i)=F_\ell(u+\sum s_ig_i^\circ)$ for small $s$. Every derivative may therefore be taken in the directions $g_i^\circ$, and we apply \Cref{Y-H5-LH}\,(a) with $\|g^\circ\|=|g^{w\circ}|_w^{\mathrm{wk}}$.

Where $D(V)\cap X\ne\emptyset$, some $p$ has $\phi_p>\varepsilon$, and then $x\varepsilon\phi_p>x\varepsilon^2=p_0^2$. There the strong weight is at least $e^{c_\Xi p_0^2}$ times the weak one, and so $|g^{w\circ}|_w^{\mathrm{wk}}\le\eta_w|g^w|_w^{\mathrm{st}}$.

For the channels $h\to h,w$ let $p=(\tilde\kappa,c_h)$, and let $p^\circ:=(\tilde\kappa,c_h/2)$, which is a pair of $\mathbb P$. The norms at $p$ and at $p^\circ$ differ only in the $h$-unit: $|\cdot|_h^{(p^\circ)}=e^{-(c_h/2)p_0(x)^2}|\cdot|_h^{(p)}$.

Hence $\|u\|^{(p^\circ)}\le\|u\|^{(p)}\le B$ and $\|(0,g^h,0)\|^{(p^\circ)}=e^{-(c_h/2)p_0^2}|g^h|_h^{(p)}$.

By \Cref{Y-H1-U1}\,(U) the same map $F_\ell$ satisfies \Cref{Y-H-C9}\,(b) at $p^\circ$ on $M_{16}(x;p^\circ)$: the cell sum of $E'^h$ is at most $e^{-(c_h/2)p_0'^2}$, and $w'$ has the strong cap.

We apply \Cref{Y-H5-LH}\,(a) in the ball of the norm at $p^\circ$. The cell sum of $DE'^h(u)[g^h]$ is then at most \[e^{-(c_h/2)p_0'^2}e^{-(c_h/2)p_0^2}|g^h|_h/(16-B)=c_h(x')\,e^{(c_h/2)p_0'^2-(c_h/2)p_0^2}|g^h|_h/(16-B)=c_h(x')\,\eta_h|g^h|_h/(16-B),\] and $|Dw'(u)[g^h]|_w^{\mathrm{st}}\le e^{-(c_h/2)p_0^2}|g^h|_h/(16-B)\le\eta_h|g^h|_h/(16-B)$.

We turn to the second derivatives and begin with the output $a$. Every point $u$ of the open ball $\{\|u\|<16\}$ belongs to $M_A(x;p)$ for $A=\max\{1,\|u\|\}<16$; hence \Cref{Y-H1-U1}\,(b\#), applied at values $A<16$ alone, gives $|E'^{\rm old}[u]|_a^{\Sigma}\le16\vartheta_a$ on $\{\|u\|<16\}$.

\Cref{Y-H5-LH}\,(a) with $j=2$ then gives the bound $16\vartheta_a(2/(16-B))^2\|g_1\|\|g_2\|$. By \Cref{Y-H1-U1}\,(f1) each $g_i$ may be replaced by $(g_i^a,0,0)$, and the derivative is $0$ if one of the entries has no $a$-part.

For the Lipschitz bound we write $(DA(u_1)-DA(u_2))[g^a]=\int_0^1D^2A(u_t)[u_1-u_2,g^a]\,dt$; with the bound just obtained, the constant is $16\cdot4=64$.

We turn to the second Lipschitz bound of \Cref{Y-H5-Af}\,(iv), namely $|(DA(u_1)-DA(u_2))[g^a]|_a^{\Sigma}\le(\vartheta_a/(16-B))|E_1^a-E_2^a|_a|g^a|_a$. The map $\Phi(u):=DA(u)[g^a]$, with $N=|\cdot|_a^{\Sigma}$, is bounded by $\mathsf m=\vartheta_a|g^a|_a$ on $\{\|u\|<16\}$ by \Cref{Y-H1-U1}\,(h3); by \Cref{Y-H1-U1}\,(h) it is holomorphic along lines, and it depends on $u$ through $E^a$ alone.

We apply \Cref{Y-H5-LH}\,(b) to the two points $u_1$ and $(E_2^a,E_1^h,w_1)$, whose distance is $|E_1^a-E_2^a|_a$; it gives $\vartheta_a|g^a|_a|E_1^a-E_2^a|_a/(16-B)$. Since $1/(16-B)\le64/(16-B)^2$, this second bound is the smaller of the two.

For the second derivatives in the outputs $h$ and $w$ we use \Cref{Y-H5-LH}\,(a) with $j=2$ and $\mathsf m=1$. If no entry is an $h$-member, we work at $p$, with $g^\circ$ in place of each $w$-entry.

If $j_h\ge1$ entries are $h$-members, we work at $p^\circ$. The entries then give the factor $e^{-j_h(c_h/2)p_0^2}$, and the return of $E'^h$ to the unit $c_h(x')$ gives the factor $e^{(c_h/2)p_0'^2}$.

Since $\eta_h^{j_h}=e^{-j_h(c_h/2)p_0^2+j_h(c_h/2)p_0'^2}$ and $j_h\ge1$, we have $e^{(c_h/2)p_0'^2-j_h(c_h/2)p_0^2}\le\eta_h^{j_h}$. The bound for $w'$ carries no factor $e^{(c_h/2)p_0'^2}$, and $e^{-(c_h/2)p_0^2}\le\eta_h$. The $a$-unit and the $w$-unit are the same at $p$ and at $p^\circ$.

An entry may also be left whole. \Cref{Y-H5-LH}\,(a) is applied in the max norm, so $g_1$ may be kept unsplit. At $p^\circ$ we have $\|g_1\|^{(p^\circ)}\le\|g_1\|^{(p)}$, because the $h$-unit at $p^\circ$ is the larger one.

If $g_2$ is a pure entry, the bound is $(4/(16-B)^2)\|g_1\|^{(p)}$ times $|g_2^a|_a$, or times $|g_2^{w\circ}|_w^{\mathrm{wk}}$ (at $p$, with $g_2^\circ$ in place of $g_2$), or times $\eta_h|g_2^h|_h$ (at $p^\circ$, with $j_h=1$ and with the factor $e^{(c_h/2)p_0'^2}$ occurring once), according to the kind of $g_2$. Since $D^2$ is bilinear, a mixed $g_2$ gives the sum of the three.

The numerals are $4/225=1/56.25$, $4/196=1/49$ and $64/225=16/56.25$, which is $0.28444$ to five decimals.

\emph{Step 6: the local vacuum numbers and the reflection.} We prove clauses \textup{(v)} and \textup{(vi)} of \Cref{Y-H5-Af}.

\begin{enumerate} \item[\textup{(v)}] Let $\Phi=(v(Y;u))_Y$ on $\{\|u\|<16\}$, with $N=|\cdot|^{\Sigma}$. It is holomorphic along lines by \Cref{Y-H1-U1}\,(c+), and $N(\Phi(u))\le16C_gc_a(x)$ by \Cref{Y-H1-U1}\,(g2) at $A=16$.

By \Cref{Y-H1-U1}\,(f1), $v(Y;u)$ depends on $E^a$ alone. Hence $v(Y;u_2)=v(Y;u_2^\circ)$ with $u_2^\circ:=(E_2^a,E_1^h,w_1)$, where $\|u_2^\circ\|\le B$ and $\|u_1-u_2^\circ\|=|E_1^a-E_2^a|_a$, and we apply \Cref{Y-H5-LH}\,(b). The numeral is $16\cdot263/14=300.57$ to two decimals.

For the coefficients we use \Cref{Y-H5-LS}\,(S1) in $|\cdot|^{\Sigma}$: for $a(z)=(v(Y;u_1(z)))_Y$, which is bounded by $C_gBc_a(x)$ by \Cref{Y-H1-U1}\,(g2) at $A=B$, and for the difference, which is bounded by the first sentence of \textup{(v)}.

\item[\textup{(vi)}] We begin with the involution. The map $g\mapsto g^{\#}$ is an antiholomorphic automorphism of $\operatorname{SL}(2,\mathbb C)$ that fixes $G$ pointwise; thus $(gh)^{\#}=g^{\#}h^{\#}$ and $(g^{-1})^{\#}=(g^{\#})^{-1}$.

Hence $\mathsf W^{\#}$ is a field, with $\mathsf W^{\#}(\text{path})=\mathsf W(\text{path})^{\#}$ and $\partial\mathsf W^{\#}(p)=(\partial\mathsf W(p))^{\#}$. Since $|g^{\#}|=|g^{-1}|=|g|$ (the definition ``Group'' in item~2 of \S\ref{Y-I-1}), the conditions \textup{(u1)} and \textup{(u2)} are preserved under $\mathsf W\mapsto\mathsf W^{\#}$.

For $h\in \operatorname{SL}(2,\mathbb C)$ we have $h^{-1}=(\operatorname{tr}h)1-h=\operatorname{adj}h$. The map $\operatorname{adj}$ is linear on $2\times2$ matrices, with $\operatorname{adj}1=1$ and $\operatorname{adj}A=JA^{\mathsf T}J^{-1}$ for a unitary $J$. Therefore $|h^{\#}-1|=|(h^{-1}-1)^\dagger|=|\operatorname{adj}(h-1)|=|h-1|$, and \textup{(u3)} is preserved as well.

So $\mathcal U_X(x)$ is mapped onto itself. The map $\mathsf W\mapsto\overline{E^a(X,\mathsf W^{\#})}$ is holomorphic, $(\mathsf W^u)^{\#}=(\mathsf W^{\#})^u$ for $G$-valued $u$, and the pure gauges are fixed.

Hence $u^*\in\mathfrak M(x)$, $u^{**}=u$ and $(cu)^*=\overline c\,u^*$, and every supremum of item~0\,(2) of \S\ref{Y-H5-1} is unchanged under $u\mapsto u^*$.

If $u=u^*$, then $E^a(X,V)=\overline{E^a(X,V)}$ at $G$-valued $V$, and $E^h$ and $w$ are real; thus $u$ is real.

Next we prove the identity. Let $\Phi$ be either $\mathsf V'$ or, at a $G$-valued argument, a value of $F_\ell$, and put $\psi(u):=\overline{\Phi(u^*)}-\Phi(u)$ on $\{\|u\|<16\}$. This function is bounded, and it is holomorphic along lines, since $(u+sg)^*=u^*+\overline s\,g^*$.

If $u=u^*$, then $u$ is real, so $\Phi(u)$ is real by \Cref{Y-H-C9}\,(d) at $p$, which holds by \Cref{Y-H1-U1}\,(U); hence $\psi(u)=0$.

For an arbitrary $u$ we put $u_1:=(u+u^*)/2$ and $u_2:=(u-u^*)/(2i)$. Both are fixed by $*$, and $u=u_1+iu_2$.

The set $O:=\{(s,t)\in\mathbb C^2:\|su_1+tu_2\|<16\}$ is open and convex and contains $(0,0)$ and $(1,i)$. The function $\chi(s,t):=\psi(su_1+tu_2)$ is holomorphic on $O$ by Osgood's lemma, and it vanishes at the real points of $O$.

Hence all derivatives of $\chi$ at $(0,0)$ vanish, so $\chi=0$ near $(0,0)$ and therefore on the connected set $O$. In particular $\chi(1,i)=\psi(u)=0$.

It remains to treat complex arguments. The map $\mathsf W'\mapsto E'^a[u^*](Y,\mathsf W')-\overline{E'^a[u](Y,\mathsf W'^{\#})}$ is holomorphic on $\mathcal U_Y(x')$ and vanishes at its $G$-valued points.

Now $G^{\text{bonds}}$ is a totally real submanifold of $\operatorname{SL}(2,\mathbb C)^{\text{bonds}}$ of full real dimension. A holomorphic function that is zero on an open piece of such a submanifold is zero near that piece, and so on the whole component.

The reason is that in the chart $\mathsf W(b)=V(b)\exp\{i\sum_kc_{b,k}\sigma_k\}$, with complex $c$, the $G$-valued fields are those with real $c$, and the derivatives of the function in all real directions are zero there.

Finally, the density $\rho[\cdot;\ell'^{\rm c}]$ uses the members at $G$-valued fields alone. \end{enumerate} \end{proof}

\begin{proof}[Proof of \Cref{Y-H5-Bc}] \emph{Step 7: the objects at a complex kernel.} In this step we prove the bound $|q_y^D|\le(9/2)\mathsf g_\flat^2p_0^2\delta$, clause \textup{(B2)}, the factorisation $\Xi_{w^\flat}=\Xi_{\tilde w}\prod_{\mathrm{off}}e^{q_\square^D}$, clause \textup{(B1)}, the statement on the dependence set of $e^{Q_D}$, and the identity $\rho[u;\ell;\tilde{\mathsf E}-\mathsf E]=\rho[u;\tilde\ell]$.

We use the facts \textup{(G1)} to \textup{(G5)} on the geometry, from Step~0 of \S\ref{Y-I-1}. The fact \textup{(G5)}, namely $\sum_{p,p'}|D_y(p,p')|\le9\mathsf g_\flat^2\delta$, involves $|D_y(p,p')|$ alone; by its own proof it holds for every family $D=(D_y)$ of complex kernels.

That proof runs as follows. The norm $\|D\|_{(\lambda)}$ is defined in item~2 of \S\ref{Y-I-1} for a family $D=(D_y)$ of kernels, and its definition gives $|D_y(p,p')|\le\delta e^{-\lambda_\flat(|y-p|+|y-p'|)}$. Three plaquettes share a base site, and the periodic sum is at most the sum over $\mathbb Z^3$; hence $\sum_pe^{-\lambda_\flat|y-p|}\le3\mathsf g_\flat$.

We turn to $q_y^D$. At a $G$-valued $V$ each $T(y,\cdot)$ is an isometry of the real space $\mathfrak g$, $|\zeta_p|\le\varepsilon$, $m_y\le1$ and $x\varepsilon^2=p_0^2$; therefore \[|q_y^D|\le(x/2)\varepsilon^2\cdot9\mathsf g_\flat^2\delta=(9/2)\mathsf g_\flat^2p_0^2\delta.\]

The brackets $\langle\cdot,\cdot\rangle$ are real, Borel and gauge invariant (\Cref{Y-I-P6}), and they are functions of $V$ at the sites within $\max\{\rho_0,\varrho(x)+1\}+1$ of $y$; the number $q_y^D$ is a complex linear combination of them.

If $D=\tilde{\mathsf E}-\mathsf E$, Step~3 of the proof of \Cref{Y-I-LB} gives $-x\mathrm{En}_{\tilde\ell}=-x\mathrm{En}_\ell+\sum_\square q_\square^D$, and so $\rho[u;\ell;D]=\rho[u;\tilde\ell]$. For every family $D=(D_y)$ of kernels that satisfies the hypotheses of \Cref{Y-H5-Bc} on $D$, the lemma defines $\rho[u;\ell;D]:=\rho[u;\ell]\exp\{\sum_\square q_\square^D\}$. Only for $D=\tilde{\mathsf E}-\mathsf E$ does the identity of the energies make it the density over the rows $\tilde\ell$.

\begin{enumerate} \item[\textup{(B2)}] By \textup{(G2)}, every path and every plaquette that occurs lies in $\square^{\approx}$.

By \textup{(u3)} and \textup{(F1)}, $|\partial\mathsf W(p)-1|<2\varepsilon\le1/20$. Hence the series $\log(1+A)$ converges, and $|F_p|\le-\log(1-2\varepsilon)\le1.02587\cdot2\varepsilon\le2.06\varepsilon$.

By \textup{(u2)}, $|\mathsf W(\Gamma)|<2$, so $|\operatorname{Ad}_{\mathsf W(\Gamma)}F_p|\le4\cdot2.06\varepsilon=8.24\varepsilon$; moreover $|Y\cdot Y'|\le|Y||Y'|$.

Consequently one term is at most $(x/2)|D_y(p,p')|(8.24\varepsilon)^2=33.9488\,p_0^2|D_y(p,p')|$. Together with \textup{(G5)} and the $M^3$ sites, the sum is at most \[305.5392\,M^3\mathsf g_\flat^2p_0^2\delta\le c_Q^{\rm su}p_0^2\delta=\delta_\Delta c_a(x).\]

Each term is holomorphic on $\mathcal U_{\{\square\}}(x)$; it is invariant under $G$-valued gauge transformations, because such a transformation $u$ takes each of the two entries $Y$ to $\operatorname{Ad}_{u(y)}Y$ and the form is Ad-invariant; and it is $0$ at the pure gauges, where $F=0$. Hence every complex combination of the terms has these properties.

The formula shows the linearity in $D$ and the independence of $\ell$ and of $\tilde u$.

Let $V$ be tame on $\square^{\approx}$ and $y\in\square$. By \textup{(M5r)}, \textup{(F2)} and \textup{(G1)}, no plaquette within $\rho_0\le M$ of $y$ is rough, and none within $\varrho(x)+1\le Mr(x)$ has $\phi_p>1$; so $m_y=1$. Further, $\zeta_p=F_p=\mathsf L(\partial V(p))$ for $|y-p|\le\varrho(x)$, and on $\mathfrak g$ we have $\operatorname{Ad}=T$ and $\cdot=\langle\,,\rangle$. Hence, for any coefficients $D_y(p,p')$, the two formulas coincide term by term.

\item[\textup{(B1)}] We first prove the factorisation. In a family of \textup{(c$\Xi$)} the sets $X_i$ are pairwise farther apart than $s(x)\ge8$, so they are disjoint, and $D(V)$ lies in their union; therefore each $p_*(\square,V)$ belongs to exactly one $X_i$.

So $\prod_iw^\flat(X_i,V)=\prod_i\tilde w(X_i,V)\cdot\prod_{\mathrm{off}}e^{q_\square^D(V)}$, where the last factor is the same for all families, and we sum over the families. In the case $D(V)=\emptyset$, all cubes are on.

Now we prove \textup{(B1)} itself. We split $\sum_\square q_\square^D=\sum_{\mathrm{on}}+\sum_{\mathrm{off}}$, and \[\sum_{\mathrm{on}}q_\square^D(V)=\sum_\square1_{\square^{\approx}}(V)E^\Delta[D](\{\square\},V)=\sum_X1_{X^{\approx}}(V)E^\Delta[D](X,V).\]

Hence \[\rho[\tilde u;\ell]\exp\{\sum q^D\}=\exp\{-x\mathrm{En}_\ell+\sum_X1_{X^{\approx}}[\tilde E^a+E^\Delta[D]+\tilde E^h]\}\cdot\prod_{\mathrm{off}}e^{q^D}\cdot\Xi_{\tilde w}=\rho[\mathbb B_D\tilde u;\ell].\] \end{enumerate}

We turn to the dependence set. If $p_*(\square,V)\in X$, then $\square\subset X^{+r}$. Whether a cube is off, and the plaquette $p_*$, depend on $\phi_p$ on $\square^{\approx}$, which lies inside $X^{+2r}$, and hence on bonds with ends in $X^{+(2r+1)}$.

The number $q_\square^D$ depends on $V$ within $\max\{M+1,Mr\}$ of the cube, hence inside $X^{+(2r+1)}$; measured in cubes, this is within $\max\{2,r\}\le r+1$ of the cube, so the paragraph on the dependence of the product in Step~5 of the proof of \Cref{Y-I-LB} applies. Moreover $2r+1<8r=s$.

The member $w^\flat$ is Borel, being a finite product taken over a finite Borel partition by $D(V)$, and it is gauge invariant, because $D(V)$, the set of cubes that are off, $p_*$, $q_y^D$ and $\tilde w$ are.

The numerals are $8.24^2/2=33.9488$, $33.9488\cdot9=305.5392$ and $-\log(0.95)/0.05\le1.025866$.

\emph{Step 8: the clauses \textup{(B3)} to \textup{(B8)} and the membership.}

Fix $p_*\in D(V)$, in the cube $c_*$. The cubes $\square$ with $p_*(\square,V)=p_*$ have $c_*\subset\square^{\approx}$; by \textup{(G4)} there are at most $(2r+1)^3$ of them, each with $M^3$ sites.

Hence \[\sum\{|q_\square^D|:p_*(\square,V)=p_*\}\le\delta':=(2r+1)^3M^3(9/2)\mathsf g_\flat^2p_0^2\delta=((2r+1)^3c_a(x)/68)\,\delta_\Delta,\] since $4.5/306=1/68$, and by \textup{(F3)} this is $\le\delta_\Delta/2\le1/2$. Consequently $|Q_D(X,V)|\le\delta'\mathfrak n$.

On $D(V)$ we have $x\varepsilon\phi_p>p_0^2$, so $|\tilde w(X,V)|\le Ae^{-c_\Xi p_0^2\mathfrak n}e^{-\tilde\kappa(|X|_c-1)}$, and by \textup{(X15)} this is $\le Ae^{-\mathfrak n}e^{-\tilde\kappa(|X|_c-1)}$.

\begin{enumerate} \item[\textup{(B3)}] Since $|e^Q|\le e^{|Q|}$, we obtain $|w^\flat|\le Ae^{-c_\Xi p_0^2\mathfrak n+\mathfrak n/2}e^{-\tilde\kappa(|X|_c-1)}$.

Since $|e^Q-1|\le|Q|e^{|Q|}$ for complex $Q$, \[|w^\flat-\tilde w|\le Ae^{-\mathfrak n}\delta'\mathfrak ne^{\mathfrak n/2}e^{-\tilde\kappa(|X|_c-1)}\le(2/e)A\delta'e^{-\tilde\kappa(|X|_c-1)}\le(A/e)\delta_\Delta e^{-\tilde\kappa(|X|_c-1)},\] because $te^{-t/2}\le2/e$. At $\mathfrak n=0$ we have $Q=0$. The dependence is treated in Step~7.

We now prove the membership $\mathbb B_D\tilde u\in M_{A+\delta_\Delta}(x;p)$. We have $|\tilde E^a+E^\Delta[D]|_a\le A+\delta_\Delta$, since at $\#X=1$ the weight is $1$; further $|\tilde E^h|_h\le A$, and $|w^\flat|_w^{\mathrm{wk}}\le A$ since $c_\Xi p_0^2\ge1>1/2$. The clauses without a size follow from Step~7 and linearity.

\item[\textup{(B4)}] We have $\mathbb B_D\tilde u-\tilde u=(E^\Delta[D],0,w^\flat-\tilde w)$, whose members have the sizes $\delta_\Delta$, $0$ and $(A/e)\delta_\Delta$.

\item[\textup{(B5)}] The objects $E^\Delta[D]$ and $Q_D$, the set of cubes that are off, and $p_*$ depend on no $\tilde u$. Hence $\mathbb B_D(\tilde u)=(E^\Delta[D],0,0)+(\tilde E^a,\tilde E^h,\tilde we^{Q_D})$ is a constant plus a linear map.

For constant $D$, the coefficient of $z^\mu$ with $\mu\ne0$ is $(\tilde E^a_\mu,\tilde E^h_\mu,\tilde w_\mu e^{Q_D})$, and $\tilde w_\mu$ is a member with finite $|\cdot|_w^{\mathrm{st}}$ by \Cref{Y-H5-LS}.

Moreover $|\tilde w_\mu e^{Q_D}|e^{\tilde\kappa(|X|_c-1)}\le|\tilde w_\mu|_w^{\mathrm{st}}e^{-(c_\Xi p_0^2-1/2)\mathfrak n}$. This is always $\le|\tilde w_\mu|_w^{\mathrm{st}}$, and for $\mathfrak n\ge1$ it is at most $e^{-(c_\Xi p_0^2-1/2)}=e^{1/2}\eta_w$ times $|\tilde w_\mu|_w^{\mathrm{st}}$.

For holomorphic $D(z)$, the objects $q_y^{D(z)}$ and $E^\Delta[D(z)]$ are linear in the entries and so holomorphic in $z$; products and the exponential keep this property, and the caps hold at each $z$.

\item[\textup{(B6)}] In the $a$-part the bound is $|\tilde E_1^a-\tilde E_2^a|_a+\delta_\Delta(D_1-D_2)$, by linearity and \textup{(B2)}, whose bound on the size holds without the condition $\delta_\Delta\le1$. In the $h$-part it is $|\tilde E_1^h-\tilde E_2^h|_h$.

In the $w$-part we put $w_i:=\tilde w_i$ and $Q_i:=Q_{D_i}(X,V)$ for $i=1,2$, and $\mathsf D_x$ stands for $\mathsf D_x(\tilde u_1,\tilde u_2)$. We write $w_1e^{Q_1}-w_2e^{Q_2}=(w_1-w_2)e^{Q_1}+w_2(e^{Q_1}-e^{Q_2})$. Here $|Q_i|\le\mathfrak n/2$, by the bound on $|Q_D(X,V)|$ at the start of this step together with $\delta_\Delta(D_i)\le1$ from the data.

The first term is at most $\mathsf D_xe^{-c_\Xi p_0^2\mathfrak n+\mathfrak n/2}e^{-\tilde\kappa(|X|_c-1)}\le \mathsf D_xe^{-\tilde\kappa(|X|_c-1)}$; this is where the strong distance is used.

For the second term, $e^{Q_1}-e^{Q_2}=(Q_1-Q_2)\int_0^1e^{Q_2+t(Q_1-Q_2)}dt$, so its modulus is at most $|Q_1-Q_2|e^{\max|Q_i|}$. Moreover $Q_1-Q_2=Q_{D_1-D_2}$. We put $\delta'_{12}:=((2r+1)^3c_a(x)/68)\,\delta_\Delta(D_1-D_2)$. The count at the start of this step, which does not need $\delta_\Delta\le1$, shows that the modulus of $Q_{D_1-D_2}$ is at most $\delta'_{12}\mathfrak n$. By \textup{(F3)}, $\delta'_{12}\le\delta_\Delta(D_1-D_2)/2$. So the second term is at most $Ae^{-\mathfrak n}\delta'_{12}\mathfrak ne^{\mathfrak n/2}\le(A/e)\delta_\Delta(D_1-D_2)$.

\item[\textup{(B7)}] We have $\mathbb B_{tD}\tilde u=(\tilde E^a+tE^\Delta[D],\tilde E^h,\tilde we^{tQ_D})$.

Then \[\FitDisplay{\dimexpr\linewidth-1pt\relax}{|\tilde wQ_D|\exp\{(c_\Xi/2)x\varepsilon\sum_{p\in D(V)\cap X}\phi_p+\tilde\kappa(|X|_c-1)\}\le A\exp\{-(c_\Xi/2)x\varepsilon\sum_{p\in D(V)\cap X}\phi_p\}\delta'\mathfrak n\le A\delta'\mathfrak ne^{-\mathfrak n/2}\le(2/e)A\delta'.}\] With the whole weight instead of half of it, the quantity is at most $A\delta'\mathfrak n$.

In Step~10 of the proof of \Cref{Y-I-InF} we use the rate $\tilde\kappa$ of the norm in place of $\kappa$ in its $\tilde w$. Then $\tilde w$ coincides with the strong cap taken at the rate $\tilde\kappa$; the kernel is $D_y=-t^2$ at a single diagonal entry; and the number $Q$ of that step, at its $X$ and $V$, is $Q_D(X,V)$. Since $q_y\ge0$ at every $y$ for this kernel, $Q_D(X,V)\ge0$; moreover $Q_D(X,V)\ge\mathfrak n't^2p_0^2/8=\mathfrak n'\delta_\Delta c_a(x)/(8c_Q^{\rm su})$. So $|\tilde wQ_D|(X,V)$, multiplied by the whole strong weight $\exp\{c_\Xi x\varepsilon\sum_{p\in D(V)\cap X}\phi_p+\tilde\kappa(|X|_c-1)\}$, equals $Q_D(X,V)\ge\mathfrak n'\delta_\Delta c_a(x)/(8c_Q^{\rm su})$.

Step~10 of the proof of \Cref{Y-I-InF} uses $\kappa$ only in $\tilde w$: its rows, its field, its set and the cubes that are off do not depend on the rate, and the identities $m_{y_i}=1$ and $q_{y_i}=t^2p_0^2/8$ hold for every rate. It uses $x\ge x_9$ to obtain $r(x)\ge2$, which it needs for the inequality $d\le Mr(x)$. In that step the member is $\tilde u=(0,0,\tilde w)$, and its sizes are $0$, $0$ and $1$. Hence $\tilde u$ satisfies the hypothesis $|\tilde E^a|_a,|\tilde E^h|_h,|\tilde w|_w^{\mathrm{st}}\le A$ of \Cref{Y-H5-Bc} if and only if $A\ge1$, which is why \textup{(B7)} assumes $A\ge1$.

\item[\textup{(B8)}] We have \[F_p(\mathsf W^{\#})=-i\log((\partial\mathsf W(p))^{\#})=-i(-(\log\partial\mathsf W(p))^\dagger)=(F_p(\mathsf W))^\dagger,\] and also $\operatorname{Ad}_{g^{\#}}Y^\dagger=(\operatorname{Ad}_gY)^\dagger$ and $Y^\dagger\cdot Y'^\dagger=\overline{Y\cdot Y'}$.

Hence $\overline{E^\Delta[D](\{\square\},\mathsf W^{\#})}=E^\Delta[\overline D](\{\square\},\mathsf W)$ holds term by term on all of $\mathcal U_{\{\square\}}(x)$.

At $G$-valued $V$ the brackets in $q_y^D$ are real, so $\overline{q_y^D}=q_y^{\overline D}$ and $\overline{e^{Q_D}}=e^{Q_{\overline D}}$. \end{enumerate}

The numeral is $\max_nne^{-n/2}=2/e\le0.73576$, the maximum being attained at $n=2$. \end{proof}

\begin{proof}[Proof of \Cref{Y-H5-Cf}] \emph{Step 9: one step for two families, one of them over $D(z)$.}

By \Cref{Y-H5-Bc}\,(B5), $u_2^\flat$ is a family on $\Omega$ within $B_2+\delta_\Delta\le B_2+1\le B$ caps at $p$, with the weak cap.

\begin{enumerate} \item[\textup{(C1)}] By \Cref{Y-H5-Bc}\,(B1), $\rho[\tilde u_2(z);\ell;D(z)]=\rho[u_2^\flat(z);\ell]$ as bounded Borel functions of $V$. The step is linear in $\rho$ and free of the rows (the definition ``The step'' in item~2 of \S\ref{Y-I-1}). \Cref{Y-H-C9}\,(a) at $p$, which holds by \Cref{Y-H1-U1}\,(U), applied at $A=B\le2.6$, gives the formula at every $V'$.

\item[\textup{(C2)}] This is \Cref{Y-H5-Af}\,(i) for $u_1$ and for $u_2^\flat$.

\item[\textup{(C3)}] At each $z$, by \Cref{Y-H5-Bc}\,(B4), \[\mathsf D^{\mathrm{in}}_x(u_1,u_2^\flat)\le \mathsf D^{\mathrm{in}}_x(u_1,\tilde u_2)+\mathsf D^{\mathrm{in}}_x(\tilde u_2,u_2^\flat)\le \mathsf D^{\mathrm{in}}_x(u_1,\tilde u_2)+\max\{1,B_2/e\}\delta_\Delta,\] and $B_2\le1.6<e$. Then we apply \Cref{Y-H5-Af}\,(ii).

\item[\textup{(C4)}] Here $B_1=B_2=1$ and $B=2$. By Step~7, $\rho[\tilde u_2;\ell;\tilde{\mathsf E}-\mathsf E]=\rho[\tilde u_2;\tilde\ell]$, and $u_2^\flat$ is real by \Cref{Y-H5-Bc}\,(B8). The images and the numbers $\mathsf V'$ are real by \Cref{Y-H-C9}\,(d), and the images are $>0$; the rows $\ell'^{\rm c}$ have range $\le\varrho(x')$; and $\mathsf D^{\mathrm{in}}_x\le \mathsf D_x$. \end{enumerate}

\end{proof}

\subsection{Source derivatives for one run}\label{Y-H6-1}

This section treats a single run of $K$ drawn steps from Wilson's weight, in which $n$ sources with pairwise separated supports enter through the bare weight alone. Starting at the level $0$, we carry the multi-affine jets of the members at $z=0$ from level to level; the families of members having these jets are each defined for $z$ in a polydisc: at every level at which the distance between the supports, counted in cubes of the level, is large, the radius increases from that level to the next, and all later levels share one and the same polydisc. At the end we write $S^T_n$ as a sum of local vacuum numbers and one last integral, in \Cref{Y-H6-J}\,\textup{(J7)}, and we bound the local vacuum numbers in \textup{(J8)}.

\emph{Item 0: what the run with sources takes from the one without sources, and what the proof of the proposition uses.} Kernels, law, centre and rows are taken from the run without sources. In the sum of local vacuum numbers and the last integral that together give $S^T_n$, nothing is subtracted. The proof of \Cref{Y-H6-J} uses Thms.~\ref{Y-H1-U1}, \ref{Y-H2-U2}, \ref{Y-H3-U3} and~\ref{Y-H4-U4}, Lems.~\ref{Y-H5-LH}, \ref{Y-H5-LS} and~\ref{Y-H5-LJ}, and \Cref{Y-H5-Af}\,\textup{(i)}, \textup{(ii)}, \textup{(iv)}.

\emph{Item 1: the data.} The data are those of \Cref{Y-H9-T}\,\textup{(T1)}. The group is $G=\operatorname{SU}(2)$ and the dimension is $3$; the letters are taken from item~(0) of \S\ref{Y-H0-1}, the numbers and the floors from item~(1) of \S\ref{Y-H0-1}, and the chain runs at $\kappa_1=\kappa+1$ with $\kappa\ge6$. We take $\mathsf b\ge x_s$, and $N=2L^{m_l}$ with $m_l\ge m'$; since $M=L^{m'}$, this gives $M\mid N$. The data include \textup{(Tor)}, which here reads $N\ge2M\,s(L\mathsf b)$.

The number of steps is $K\ge1$. We set $\beta=x_0=L^K\mathsf b$ and $x_k=L^{K-k}\mathsf b$, so that $x_{k+1}=x_k/L$ and $x_K=\mathsf b$, and the spacing is $a_K=L^{-K}$. The lattices are $\mathbb T_k=\mathbb Z^3/NL^{K-k}\mathbb Z^3$, with the fitted grids fixed by $o_K=0$ and $o_k=Lo_{k+1}-\mathsf h_L$. The rows start from $\ell_0=\emptyset$, and $\ell_{k+1}=(\ell_k')^{\rm c}$ cut at $\varrho(x_{k+1})$, following \Cref{Y-GB-ThB} and \Cref{Y-H4-U4}\,\textup{(b)}.

We place $\mathbb T_0$ in $T_l=\mathbb R^3/N\mathbb Z^3$ with the spacing $a_K$, the site $z$ being the point $a_Kz$. In this way $\mathbb T_0$ becomes the lattice $a_K\mathbb Z^3/N\mathbb Z^3$ in $T_l$, which is the $\Lambda_{K,\ell}$ of \Cref{support:150-2} for $\ell=N$.

We fix $n\ge2$, $d>0$, $R_f>0$ and real bounded functions $f_1,\dots,f_n$ on $T_l$ with $\|f_i\|_\infty\le1$, in such a way that any two of the sets $\{f_i\ne0\}$ are at sup distance $\ge d$ from each other and each of these sets lies within sup distance $R_f$ of some point.

For a plaquette $p$ we write $s_p=a_p=1-\frac12\operatorname{tr}\partial U(p)$, and $x(p)$ for its barycentre: if $p=(z;\mu\nu)$ is the plaquette based at the site $z$, then $x(p)=a_K(z+(e_\mu+e_\nu)/2)$. For a function $f$ on $T_l$ let $\mathcal O_K(f)=a_K^{-1}\sum_pf(x(p))s_p$. We define \[Z_K(z)=\int\exp\{-\beta\sum_ps_p+\sum_iz_i\mathcal O_K(f_i)\}\,dU\] and we let \[S^T_n=S^T_{n;K,l,\beta}(f_1,\dots,f_n)=\partial_1\cdots\partial_n\log Z_K(z)\] at $z=0$.

So defined, $S^T_n$ is the cumulant $\langle \mathcal O_K(f_1);\dots;\mathcal O_K(f_n)\rangle=\partial_1\cdots\partial_n{\rm Log}\,\Gamma(0)$, with $\Gamma(z)=Z_K(z)/Z_K(0)$, of \Cref{support:150-2}; here ${\rm Log}$ denotes the principal branch of the complex logarithm, taken on $|w-1|\le1/2$. That definition introduces $x(p)$, $a_p$, the measure, $\mathcal O(f)$ and $S^T_n$ on a torus, and they agree with the objects defined above; the model is the one of \Cref{support:150-2}. In this item, $L$ and $M$ are those of the table, which is item~3 of \S\ref{Y-G0-1}, with $L$ subject to \textup{(O4)} and $M$ to \textup{(O7)}, and they also satisfy the entries \textup{(O4a)} and \textup{(O7)} of item~(1) of \S\ref{Y-H0-1}.

\emph{Item 2: the objects \textup{(o1)} to \textup{(o9)}.} We introduce the following objects.

\begin{enumerate}
\item[\textup{(o1)}] \emph{The honest family.} For $z\in\mathbb C^n$ we put $\nu_p(z):=\mathsf b^{-1}\sum_iz_if_i(x(p))$. The sets $\{f_i\ne0\}$ are pairwise disjoint, so at any point at most one $f_i$ differs from $0$; therefore $|\nu(z)|_\infty\le \mathsf b^{-1}\max_i|z_i|\,\|f_i\|_\infty\le \mathsf b^{-1}\max_i|z_i|$, where the second inequality uses $\|f_i\|_\infty\le1$.

The set ${\rm Pd}_0=\{|z_i|\,\|f_i\|_\infty\le1.99c_\nu \mathsf b\,\varepsilon(\beta)\}$ is the one of item~3, ``Sources'', of \S\ref{Y-H4-1}, and here $1.99c_\nu \mathsf b\,\varepsilon(\beta)=1.99/{\rm amp}_W$. On ${\rm Pd}_0$ we define $u_0(z):=u_0(\nu(z))$, $u_{k+1}(z):=F_{\ell_k}(u_k(z))$ and $V_{k+1}(z):=\mathsf V'_{\ell_k}[u_k(z)]$, as in \Cref{Y-H4-U4}\,\textup{(h)}. All these values are holomorphic on the open set ${\rm Pd}_0^o=\{|z_i|\,\|f_i\|_\infty<15.99/{\rm amp}_W\}$ of \Cref{Y-H4-U4}\,\textup{(h)}, and this open set contains the closed set ${\rm Pd}_0$. Whenever the honest family is regarded as a family in the sense of \textup{(o6)}, we take it on the interior of ${\rm Pd}_0$. Finally, $\bar u_k:=u_k(0)$.

\item[\textup{(o2)}] \emph{Jets.} For a non-empty subset $S$ of $\{1,\dots,n\}$ we write $z^S:=\prod_{i\in S}z_i$ and $\partial_S:=\prod_{i\in S}\partial/\partial z_i$. Value by value we define $\delta_S(k):=\partial_Su_k(z)$ at $z=0$; its parts are $\delta_S(k)=(\delta_S^a,\delta_S^h,\delta_S^w)(k)$, and for a single index we abbreviate $\delta_i:=\delta_{\{i\}}$. With the sum over all such $S$, we put $u_k^{\rm pol}(z):=\bar u_k+\sum_Sz^S\delta_S(k)$.

\item[\textup{(o3)}] \emph{Blocks.} We denote by $S_i(k)$ the set of those cubes of $\mathbb T_k$ for which some plaquette $p$ of $\mathbb T_0$ with $f_i(x(p))\ne0$ has its base site in the bare box of the cube, and we call these cubes the $k$-blocks of ${\rm supp}\,f_i$. The bare box of a cube consists of the sites of $\mathbb T_0$ lying in the $k$-fold blocks of its $M^3$ sites, and it is a box of side $ML^k$. As the grids fit, $S_i(k+1)={\rm cl}(S_i(k))$ is the set of cells containing a cube of $S_i(k)$. For $k=0$, $S_i(0)$ is the set $S(f_i)$ of \Cref{Y-H4-U4}.

\item[\textup{(o4)}] \emph{Reach, separation, rates.} The reach is $\check r(x)=r(x)+3$. For a depth $i\ge0$ we put ${\rm sep}^{[i]}:=\lfloor dL^i/M\rfloor-2\check r(L^i\mathsf b)-2$, and at the level $k$ we write ${\rm sep}_k:={\rm sep}^{[K-k]}$. With ${\rm sep}_*:=8\log(4nL^{1/2})$, we let $j^*$ be the least $j\ge1$ such that ${\rm sep}^{[i]}\ge{\rm sep}_*$ holds for all $i\ge j-1$, and we let $k_d:=\max\{K-j^*,0\}$. For $k<k_d$ we define $\check\kappa_k:=2\log(4nL^{1/2})/{\rm sep}_{k+1}$, and the rate of the level $k$ is $\kappa_k:=\kappa+\sum_{k\le j<k_d}\check\kappa_j$.

\item[\textup{(o5)}] \emph{Amplitudes, radii.} The amplitude is ${\rm amp}(x)=c_Q^{\rm su}K_5x^{1/2}/(\mathsf bE_ap_0(x))$, where $c_Q^{\rm su}=306M^3\mathsf g_\flat^2$ as in \textup{(O10)} of item~1 of \S\ref{Y-I-1}, and we also use ${\rm amp}_W=1/(c_\nu \mathsf b\,\varepsilon(\beta))=6M^3x_0^{1/2}/(\mathsf bE_ap_0(x_0))$. Further, $u_3(x)=C_3x^{1/4}$. The function $\check\omega_3(x)=C_3E_ap_0(x)x^{-1/4}/c_Q^{\rm su}$ is defined in \textup{(O12b)} of item~(1) of \S\ref{Y-H0-1}; it is independent of $\mathsf b$, $n$ and $d$, which is why the condition \textup{(W3)} on $x_s$, given in \textup{(o9)} below, may involve it.

These definitions yield the identity \[\check\omega_3(x)=u_3(x)K_5/(\mathsf b\,{\rm amp}(x))\] for, once $u_3$ and ${\rm amp}$ are inserted, the right side reads $C_3x^{1/4}K_5\,\mathsf bE_ap_0(x)/(\mathsf b\,c_Q^{\rm su}K_5x^{1/2})$, which is $\check\omega_3(x)$ after $\mathsf b$ and $K_5$ cancel.

The radii are $\mathfrak r(x)=3.5/(n\,{\rm amp}(Lx))$, $\mathfrak r_k:=\mathfrak r(x_k)$ and $\mathfrak r_*:=\mathfrak r_{k_d}$, and for $r>0$ we write ${\rm Pd}(r)=\{z\in\mathbb C^n:|z_i|<r\ \text{for all }i\}$.

\item[\textup{(o6)}] \emph{Norms.} Each norm below is taken at a coupling $x$ and at a rate $\tilde\kappa\in[\kappa,\kappa_1]$; it is defined on the triples $\delta=(\delta^a,\delta^h,\delta^w)$ obeying those clauses of \textup{(ca)}, \textup{(ch)}, \textup{(c$\Xi$)} which carry no size and satisfying $\|\delta\|_x<\infty$, where $\|\cdot\|_x$ is the norm of item~3 of \S\ref{Y-I-1}; these triples form the space $\mathfrak M(x)$ of that item. We set \[|\delta^a|_a^{(\tilde\kappa)}:=\sup_X\sup_{\mathcal U_X(x)}|\delta^a(X,\cdot)|\,e^{\tilde\kappa(\#X-1)}/c_a(x)\] and \[|\delta^h|_h^{(\tilde\kappa)}:=\sup_X\sup|\delta^h(X,\cdot)|\,e^{\tilde\kappa(\#X-1)}/c_h(x)\] where the inner supremum runs over the tame $G$-valued fields, and \[|\delta^w|_w^{{\rm st},(\tilde\kappa)}:=\sup_{X,V}|\delta^w(X,V)|\exp\{c_\Xi x\varepsilon\sum_{p\in D(V),\,p\in X}\phi_p+\tilde\kappa(|X|_c-1)\}.\]

The strong norm $\|\delta\|^{{\rm st},(\tilde\kappa)}$ is the largest of these three numbers. The weak norm $\|\delta\|^{(\tilde\kappa)}$ is built in the same way, except that the term with $c_\Xi$ is dropped from the exponential; it never exceeds the strong one. A norm written without a rate is taken at $\kappa_1$.

The space $T^a(x)$ consists of the $a$-members that obey the clauses of \textup{(ca)} without a size and are each bounded, as in item~1 of \S\ref{Y-H1-1}. In Sections~\ref{Y-H3-1} and~\ref{Y-H4-1} the symbol $|\cdot|^{{\rm cs},(\tilde\kappa)}_{a,x'}$ stands for the cell sum divided by $c_a(x')$. As Step~0\,\textup{(c)} of the proof of \Cref{Y-H6-J} shows, it dominates the supremum $|\cdot|_a^{(\tilde\kappa)}$ used in this section. A family on an open set of $\mathbb C^n$ is a map $z\mapsto u(z)$ into some $M_A(x;p)$ with $A\le16$ all of whose values are holomorphic.

\item[\textup{(o7)}] We take $v_k(Y;u)$ and $V_0^+$ from \Cref{Y-H1-U1}\,\textup{(g1)}, applied with the coupling and the rows $(x_k;\ell_k)$; there $Y$ runs over the connected sets of cells, that is, of cubes of $\mathbb T_{k+1}$. We put $I[u]:=\int\rho[u;\ell_K](V)\,dV$, the integral being over the fields of $\mathbb T_K$.

\item[\textup{(o8)}] The source forms $\mathbb D_k[f]$, the map $T_\ell$ and ${\rm cut}$ are those of \Cref{Y-H2-U2}, in item~1 of \S\ref{Y-H2-1}; $s_x[D]$ is that of entry~5 of item~(0) of \S\ref{Y-H0-1}, and $n[u;D]$ that of \Cref{Y-H3-U3}. By $DA_\ell(u)[\delta]$ we mean the series $DA(u)[\delta]$ of item~1 of \S\ref{Y-H1-1} taken at the rows $\ell$, as in item~3 of \S\ref{Y-H3-1}. \Cref{Y-H1-U1}\,\textup{(h)} shows that it is linear in $\delta\in T^a(x)$ and that, wherever $A_a(u)<16$, it equals $\frac d{dz}|_0\,E'^{\rm old}[u+z(\delta,0,0)]$; here $A_a(u)$ is the $a$-part of the norm of $u$, as in item~1 of \S\ref{Y-H1-1}. We recall $\vartheta_a=0.2431L^{-1/4}$ and $C_g=263L^3$. The functions $\eta_h$, $\eta_w$ are those of item~(1) of \S\ref{Y-H0-1}.

\item[\textup{(o9)}] \emph{One symbol, two senses.} The symbol $\kappa_1$ always means $\kappa+1$, as in entry~1 of item~(0) of \S\ref{Y-H0-1}, while $\kappa_k$, $\kappa_{k+1}$, $\kappa_{k_d}$ and $\kappa_0$ are the rates of \textup{(o4)} at the levels; $u_3(x)$ is always the number $C_3x^{1/4}$, whereas $u_k(z)$, $\bar u_k$ and $u_k^{\rm pol}$, whose index is a letter, $0$ or $K$, are the families of the levels. The symbols $x_3$, $x_4$, $x_9^U$, $x_s$ are floors of \textup{(O13)} of item~(1) of \S\ref{Y-H0-1}, the index in $x_s$ being a mark, which is neither a level nor a depth; $x_k$ with a letter index, $x_0$ and $x_K$ are the couplings of the levels. The letter $p$ is a plaquette, except where it denotes an element of $\mathbb P$, namely a pair $(\tilde\kappa,\tilde c)$ consisting of a rate and a constant; this is the case for the last argument of $M_A(x;p)$, for the $p$ of ``at $p$'' and wherever $p$ is written as such a pair.

The tags \textup{(W2)} and \textup{(W3)} denote the two final conditions that \textup{(O13)} of item~(1) of \S\ref{Y-H0-1} imposes on $x_s$; for all $x\ge Lx_s$ they read
\begin{enumerate}
\item[\textup{(W2)}] $\eta_h(x)+e^{1/2}\eta_w(x)\le1/L$,
\item[\textup{(W3)}] $\check\omega_3(x),\ \check\omega_3(x/L)\le10^{-3}$.
\end{enumerate}

We apply \textup{(W2)} at the couplings $x_k$ with $k<K$ and \textup{(W3)} at $x_k$ with $k\le K$; at $k=K$, \textup{(W3)} enters through its second member $\check\omega_3(x/L)$, taken at $x=L\mathsf b$.
\end{enumerate}

\emph{Item 3: the numbers of \textup{(o4)} and \textup{(o5)}.} These numbers have the following properties.

\begin{enumerate}
\item[\textup{(G1)}] For every $i\ge0$ we have ${\rm sep}^{[i+1]}\ge L\,{\rm sep}^{[i]}$. Consequently $j^*$ is the least $j\ge1$ with ${\rm sep}^{[j-1]}\ge{\rm sep}_*$; it is finite, it depends on $\mathsf b$, $n$, $d$, $L$, $M$, $c_R$ alone, and it does not decrease when $\mathsf b$ increases, so that it is locally bounded in $\mathsf b$.

\item[\textup{(G2)}] We have $k_d\le K-1$. For $k<k_d$, ${\rm sep}_{k+1}\ge L\,{\rm sep}_*$ and $\check\kappa_k\le1/(4L)\le1/4$. The rate $\kappa_k$ is non-increasing in $k$, $\kappa_{k_d}=\kappa$, and $\kappa\le\kappa_k\le\kappa+1/(4(L-1))$; hence $\kappa_k$ lies in $[\kappa,\kappa+0.26]$ and therefore in $[\kappa,\kappa_1]$.

\item[\textup{(G3)}] For $x\ge x_s$ we have \[5.01\,{\rm amp}(x)\le2^{-1/4}L^{1/2}{\rm amp}(x)\le{\rm amp}(Lx)\le L^{1/2}{\rm amp}(x)\] and so $\mathfrak r_k\le\mathfrak r_{k+1}\le L^{1/2}\mathfrak r_k$; moreover $\mathfrak r_*\ge\mathfrak r(L^{j^*}\mathsf b)$ for every $K\ge1$, with equality if $K\ge j^*$.

\item[\textup{(G4)}] The sets $S_i(k)^{+\check r(x_k)}$, $i=1,\dots,n$, are pairwise at a distance $\ge{\rm sep}_k+2$ in cubes.
\end{enumerate}

\emph{Item 4: the step for jets of one run.}

\begin{proposition}[The step for jets of one run]\label{Y-H6-J} In the setting of item~1, \textup{(J0)} to \textup{(J5)} hold for $0\le k\le k_d$, and \textup{(J6)} to \textup{(J8)} hold as well.
\begin{enumerate}
\item[\textup{(J0)}] \Cref{Y-H2-U2}, applied with $f=f_i$, shows that $\mathbb D_k[f_i]$ is real, symmetric and linear in $f_i$, that its range is $\le\varrho(x_k)$, that $\|\mathbb D_k[f_i]\|_{(\lambda_\flat)}\le K_5\|f_i\|_\infty$ site by site, and that $\mathbb D_{k,y}[f_i]$ depends only on the values of $f_i$ at the barycentres of the plaquettes based in $B^k(y)$. In addition, $s_{x_k}[\mathbb D_k[f_i]/\mathsf b]$ lies in $T^a(x_k)$ and lives on single cubes of $S_i(k)$, and $|s_{x_k}[\mathbb D_k[f_i]/\mathsf b]|_a\le{\rm amp}(x_k)$ holds in the norm $|\cdot|_a^{(\tilde\kappa)}$ for every rate $\tilde\kappa$ of \textup{(o6)}.

\item[\textup{(J1)}] Let $k<K$. For every $S$, the coefficient of $z^S$ in $F_{\ell_k}(u_k^{\rm pol}(z))$ is $\delta_S(k+1)$, and the coefficient of $z^S$ in $\mathsf V'_{\ell_k}[u_k^{\rm pol}(z)]$ is $\partial_SV_{k+1}(0)$.

\item[\textup{(J2)}] We have $\delta_S(k)(X)=0$ unless $X^{+\check r(x_k)}$ meets $S_i(k)$ for every $i\in S$. Consequently, if $|S|\ge2$ and ${\rm sep}_k\ge1$, then $\delta_S(k)(X)\ne0$ forces $\#X\ge|S|\,{\rm sep}_k/2$ for the $a$- and $h$-parts, $X$ being connected, and $|X|_c\ge|S|\,{\rm sep}_k/2$ for the $w$-part.

\item[\textup{(J3)}] We have $\delta_i^a(k)=s_{x_k}[\mathbb D_k[f_i]/\mathsf b]+{\rm rem}_{k,i}$ with ${\rm rem}_{k,i}\in T^a(x_k)$ and \[|{\rm rem}_{k,i}|_a^{(\kappa_1)}\le2u_3(x_k)K_5/\mathsf b=2\check\omega_3(x_k){\rm amp}(x_k)\le0.002\,{\rm amp}(x_k).\]

\item[\textup{(J4)}] We have $|\delta_i^h(k)|_h^{(\kappa_1)}\le0.12\,{\rm amp}(Lx_k)$ and $|\delta_i^w(k)|_w^{{\rm st},(\kappa_1)}\le0.12\,{\rm amp}(Lx_k)$. The leading term of a single jet $\delta_i(k)$ is $(s_{x_k}[\mathbb D_k[f_i]/\mathsf b],0,0)$, so the second bound covers the whole $w$-jet.

\item[\textup{(J5)}] We have $n\,\mathfrak r_k\max_i\|\delta_i(k)\|^{{\rm st},(\kappa_1)}\le0.7$ and \[\sum_{|S|\ge2}\mathfrak r_k^{|S|}\|\delta_S(k)\|^{{\rm st},(\kappa_k)}\le0.21\] so that $u_k^{\rm pol}$ is a family on ${\rm Pd}(\mathfrak r_k)$ in $M_{1.91}(x_k;\kappa_k,c_h)$ and $\sup_{{\rm Pd}(\mathfrak r_k)}\|u_k^{\rm pol}(z)-\bar u_k\|^{{\rm st},(\kappa_k)}\le0.91$.

\item[\textup{(J6)}] From here on we set $u_k(z):=u_k^{\rm pol}(z)$ on ${\rm Pd}(\mathfrak r_k)$ for $k\le k_d$. By the representatives we mean these $u_k(z)$ and the $u_k(z)$, $k_d<k\le K$, on ${\rm Pd}(\mathfrak r_*)$ defined below. Where the family of \textup{(o1)} itself is meant, and not a representative, we call it the honest family. For $k_d\le k<K$ we put $u_{k+1}(z):=F_{\ell_k}(u_k(z))$ on ${\rm Pd}(\mathfrak r_*)$, beginning with $u_{k_d}(z)=u_{k_d}^{\rm pol}(z)$. When $k_d<k\le K$, the $u_k(z)$ so obtained are families in $M_1(x_k;\kappa,c_h)$, and they carry the strong cap. For $k_d\le k\le K$ we have \[\sup_{{\rm Pd}(\mathfrak r_*)}\|u_k(z)-\bar u_k\|^{{\rm st},(\kappa)}\le0.91\cdot0.4^{k-k_d}\] and $u_k(0)=\bar u_k$, and for every $S$ the derivative $\partial_Su_k(z)$ at $0$ is $\delta_S(k)$, the true multi-affine jet. For $0\le k\le K$, $\bar u_k$ is real, it lies in $M_1(x_k;\kappa_1,c_h)$ with the strong cap, and $\rho[\bar u_k;\ell_k]>0$, by \Cref{Y-H4-U4}\,\textup{(h)}; thus $\bar u_k$ is a system of members of a point of $C^{\rm su}_\flat(x_k;\ell_k)$, and in particular, for every $K\ge1$, $\bar u_K$ is a system of members of a point of $C^{\rm su}_\flat(\mathsf b;\ell_K)$. As $k_d\le K-1$, the family $u_K(z)$ is an output of $F_{\ell_{K-1}}$.

\item[\textup{(J7)}] With the representatives, \[S^T_n=\sum_{0\le k<K}\sum_Y\partial_1\cdots\partial_nv_k(Y;u_k(z))(0)+\partial_1\cdots\partial_n\log I[u_K(z)](0)\] and $I[\bar u_K]>0$. A set $Y$ contributes $0$ unless $Y^{+\check r(x_{k+1})}$ meets $S_i(k+1)$ for every $i=1,\dots,n$.

\item[\textup{(J8)}] For $0\le k<K$, \[|\sum_Y\partial_1\cdots\partial_nv_k(Y;u_k(z))(0)|\le\sum_Y|\partial_1\cdots\partial_nv_k(Y;u_k(z))(0)|\le\mathfrak r_{\min(k,k_d)}^{-n}N_1(k)\,2C_gc_a(x_k)e^{-\kappa(m_k-1)}\] where $N_1(k):=(2R_fL^{K-k-1}/M+2\check r(x_{k+1})+3)^3$, and where $m_k:=\lceil n\,{\rm sep}_{k+1}/2\rceil$ for $k<k_d$ and $m_k:=1$ for $k\ge k_d$.
\end{enumerate}
\end{proposition}

\begin{remark}[Item 5: constants and the range of the clauses] Every numerical constant that appears above is independent of $K$, $k$, $N$, $\mathsf b$, $n$, $d$, $R_f$ and $f$. As for the data, ${\rm amp}$ and $\mathfrak r$ are functions of $(\mathsf b,n)$ and of the depth $K-k$; ${\rm sep}$, $j^*$, $\check\kappa$, $\kappa_k$ and $m_k$ are functions of $(\mathsf b,n,d)$ and of the depth; $N_1$ is a function of $(\mathsf b,R_f)$ and of the depth; and none of them depends on $N$.

The proof given below establishes \textup{(J1)} to \textup{(J4)} of \Cref{Y-H6-J} for every $k\le K-1$. Its arguments also yield \textup{(J0)} for every $k\le K$, and \textup{(J2)}, \textup{(J3)}, \textup{(J4)} at $k=K$ as well. Indeed, Step~2 shows, for $0\le k\le K$, that every $\delta_S(k)$ belongs to the space of members and has a finite norm $\|\delta_S(k)\|^{{\rm st},(\kappa_1)}$; Steps~3, 5 and~6 pass from a level $k<K$ to the level $k+1\le K$; Step~0\,(a), together with the part of its proof that concerns the lattice, and \Cref{Y-H2-U2} apply at $k=K$; and at $k=K$ the condition \textup{(W3)} enters through the second of its two inequalities. At $k=K$ the clause \textup{(J1)} has no meaning, because no step starts at level $K$.

For complex $f_i$ with $\|f_i\|_\infty\le1$ the arguments of Steps~2 to~6 go through without any change, with one exception: the reality asserted in \textup{(J0)} refers to real $f_i$ (\textup{(U2.2)} of \Cref{Y-H2-U2}). The reason is that \Cref{Y-H2-U2} and item~3 of \S\ref{Y-H4-1}, on sources, treat complex bounded $f$; \Cref{Y-H3-U3} treats complex $D$; \Cref{Y-H1-U1}\,(h), (h3) and \Cref{Y-H5-Af}\,(iv) admit complex directions; and the only reality that enters is that of $\bar u_k$.

\end{remark}

\begin{proof}[Proof of \Cref{Y-H6-J} and of \textup{(G1)} to \textup{(G4)} of item~3] \emph{Step 0: tools.} This step provides the following four tools.

\begin{enumerate}

\item[\textup{(a)}] \emph{The size of the leading term.} Let $x\ge x_s$, and let $D=(D_y)$ be complex, symmetric in the sense that $D_y(p,p')=D_y(p',p)$, and of range $\le\varrho(x)$. Then $s_x[D]\in T^a(x)$; it vanishes on every $X$ with $\#X\ge2$; and $\sup_{\mathcal U_{\{\square\}}(x)}|s_x[D](\{\square\},\cdot)|\le\sigma_D=c_Q^{\rm su}p_0^2\|D\|_{(\lambda_\flat)}$. Measured in caps, and at every rate, $|s_x[D]|_a\le c_Q^{\rm su}x^{1/2}\|D\|_{(\lambda_\flat)}/(E_ap_0(x))$.

\item[\textup{(b)}] \emph{Three numbers.} \textup{(G3)} holds; ${\rm amp}_W\le{\rm amp}(x_0)/51$; and $u_3(x)K_5\ge1$ for $x\ge Lx_s$.

\item[\textup{(c)}] \emph{Norms.} We have $M_A(x;\kappa_1,c_h)\subset M_A(x;\tilde\kappa,c_h)$; any norm taken at the rate $\tilde\kappa$ is at most the same norm taken at the rate $\kappa_1$; and a bound for the cell sums $\sum_{Y\ni\text{a cell}}e^{\tilde\kappa(\#Y-1)}\sup|\cdot(Y)|$ is also a bound for the weighted supremum.

\item[\textup{(d)}] \emph{Representatives.} Let $\Phi$ be a value of $F_\ell(u)$, the number $\mathsf V'_\ell[u]$, or a number $v(Y;u)$, each taken for some data \textup{(SD)} at an $x\ge x_9^U$ and some $p\in\mathbb P$; or else let $\Phi$ be the number $I[u]$. Suppose that $u(z)$ and $\tilde u(z)$ are families near $z=0$, of norm $<16$, such that $\partial_{S'}u(0)=\partial_{S'}\tilde u(0)$ holds for each subset $S'$ of $S$, including the empty one. Then $\partial_S\Phi(u(z))(0)=\partial_S\Phi(\tilde u(z))(0)$.

\end{enumerate}

\emph{Proof of \textup{(a)}.} Condition \textup{(X16)} gives $\varepsilon(x)\le1/40$ for $x\ge x_s$. Since $x\ge e^2$ and $c_R\ge768$,
\[
 \varrho(x)+2\le c_R\log x+3\le20c_R\log x\le Mr(x).
\]
The side condition of \Cref{Y-H3-I4} follows from \textup{(Tor)} for $k<K$, when the coarse side $NL^{K-k-1}$ is at least $16Mr(x_k)$. At $k=K$, the initial condition $N\ge2Ms(L\mathsf b)$ gives $N\ge16Mr(\mathsf b)$. Thus the lemma applies on every $\mathbb T_k$, including the last torus. The bound in caps follows by division by $c_a(x)=E_ap_0(x)^3x^{-1/2}$, since a single cube has weight $1$ at every rate.

\emph{Proof of \textup{(b)}.} We begin with \textup{(G3)}. We have ${\rm amp}(Lx)/{\rm amp}(x)=L^{1/2}p_0(x)/p_0(Lx)$, and since $p_0$ is increasing, this ratio does not exceed $L^{1/2}$. On the other hand, \textup{(Y.0)(0b)} of item~5 of \S\ref{Y-H-1}, applied at $Lx\ge Lx_s\ge x_9^U$, gives $(p_0(x)/p_0(Lx))^4\ge1/2$; it is one of the conditions defining the floor $x_9^U$ (\textup{(O13)}\,(i) of item~0 of \S\ref{Y-H1-1}). Hence the ratio is at least $2^{-1/4}L^{1/2}\ge0.8408L^{1/2}\ge5.01$, because $L>1.2\cdot10^{13}$ (\textup{(O4)} of item~1 of \S\ref{Y-H-1}). Moreover $\mathfrak r_{k+1}/\mathfrak r_k={\rm amp}(Lx_k)/{\rm amp}(x_k)$. If $K\ge j^*$, then $x_{k_d}=L^{j^*}\mathsf b$. If $K<j^*$, then $k_d=0$ and $\mathfrak r_*=\mathfrak r(L^K\mathsf b)\ge\mathfrak r(L^{j^*}\mathsf b)$, as $i\mapsto\mathfrak r(L^i\mathsf b)$ is decreasing.

Next, ${\rm amp}_W/{\rm amp}(x_0)=6M^3/(306M^3\mathsf g_\flat^2K_5)=1/(51\mathsf g_\flat^2K_5)\le1/51$, because $\mathsf g_\flat\ge1$ and $K_5=1.35C_{\rm ch}$, where $C_{\rm ch}=4e^{2\lambda_*}/\gamma_*>4$ as $\gamma_*<1$ (\textup{(O0)} of item~1 of \S\ref{Y-I-1}).

Finally, we turn to \textup{(Y.3)} of item~5 of \S\ref{Y-H-1}. Each term inside its bracket is nonnegative, and each factor in front of the bracket is at least $1$. In the definition of $x_9^U$, \textup{(Y.3)} is imposed with $\kappa_1$ in place of $\kappa$ and with $e^{-(c_h/2)p_0^2}$ in place of $c_h(x)$ (\textup{(O13)}\,(ii) of item~0 of \S\ref{Y-H1-1}); our argument involves \textup{(Y.3)} only through its term in $c_a$. In both forms this term is $16c_a$ times $5$ times factors $\ge1$, so that the inequality $x^{1/4}\cdot50\cdot5\cdot16c_a(x)\le x^{-1/8}$ follows from either form. Therefore $E_ap_0^3\le x^{1/8}/4000$, and $x^{1/4}/E_a\ge4000p_0^3x^{1/8}\ge1$. The last inequality uses $p_0\ge1$; this is \textup{(X14)} in \textup{(O13)} of item~1 of \S\ref{Y-I-1}, a floor of $x_s$. Since \textup{(Y.3)} in turn is a floor of $x_9^U\le Lx_s$, we obtain the range $x\ge Lx_s$. It follows that $u_3(x)K_5=1686e^{\kappa+4}L^{11/4}c_Q^{\rm su}K_5x^{1/4}/E_a\ge1$.

\emph{Proof of \textup{(c)}.} As $\tilde\kappa\le\kappa_1$, we have $e^{\tilde\kappa(\#X-1)}\le e^{\kappa_1(\#X-1)}$, which gives the inclusion and the comparison of norms. For the last assertion, observe that every $Y$ contains a cell, and that in a sum of nonnegative terms each term is bounded by the sum.

\emph{Proof of \textup{(d)}.} Let $B<16$ be a common bound of $\|u(z)\|$ and $\|\tilde u(z)\|$ for $z$ near $0$. Then \textup{(d)} follows from \Cref{Y-H5-LJ}, applied at $u_0=u(0)=\tilde u(0)$ with $t=16-B$: the two coefficients are given by one and the same sum over the partitions of $S$. The lemma covers all four kinds of $\Phi$: the first three at $x\ge x_9^U$, and $\Phi=I$ for any $u_0$ and any $t$.

\emph{Step 1: geometry.} We establish the following.

\begin{enumerate}

\item[\textup{(r)}] $r(Lx)\le r(x)+\log L/960+1$.

\item[\textup{(G1)}] The clause \textup{(G1)} of item~3 holds.

\item[\textup{(G2)}] The clause \textup{(G2)} of that item holds.

\item[\textup{(G4)}] The clause \textup{(G4)} of that item holds.

\item[\textup{(R)}] $\lceil\check r(x_k)/L\rceil+1\le\check r(x_{k+1})$.

\item[\textup{(C)}] Let $X$ be a connected set of cubes, and suppose that $X$ meets $j\ge2$ sets which lie pairwise at distance $\ge\sigma\ge1$. Then $\#X\ge j\lceil\sigma/2\rceil$. For $X\in\mathcal P$ meeting $j\ge2$ such sets, the same bound holds for $|X|_c$.

\item[\textup{(A)}] The set $S_1(k+1)^{+\check r(x_{k+1})}$ has at most $N_1(k)$ cells.

\end{enumerate}

All distances are periodic sup distances.

\emph{Proof of \textup{(r)}.} We have $r(x)=\lceil20c_R\log x/M\rceil$ and $\lceil a+a'\rceil\le\lceil a\rceil+a'+1$, which we apply with $a=20c_R\log x/M$ and $a'=20c_R\log L/M$; moreover $20c_R/M\le1/(96(\kappa_1+4))\le1/960$, by \textup{(M2)} $M\ge1920(\kappa+4)c_R$ (\textup{(O7)} of item~1 of \S\ref{Y-H-1}), taken at $\kappa_1$ as in \textup{(O7)} of item~(1) of \S\ref{Y-H0-1}.

\emph{Proof of \textup{(G1)}.} We use $\lfloor La\rfloor\ge L\lfloor a\rfloor$ together with \textup{(r)}; they give \[\FitDisplay{\dimexpr\linewidth-1pt\relax}{{\rm sep}^{[i+1]}\ge L\lfloor dL^i/M\rfloor-2(\check r(L^i\mathsf b)+\log L/960+1)-2=L\,{\rm sep}^{[i]}+(2L-2)\check r(L^i\mathsf b)+2L-4-\log L/480\ge L\,{\rm sep}^{[i]},}\] where the last step uses $\check r\ge3$. Hence, at every depth greater than the first depth at which ${\rm sep}^{[i]}\ge{\rm sep}_*>0$ holds, this inequality is true as well. Such a depth exists, since ${\rm sep}^{[i]}\ge dL^i/M-40c_R\log(L^i\mathsf b)/M-11$, and the right side tends to $\infty$ with $i$. As $r$ does not decrease when $x$ grows, ${\rm sep}^{[i]}$ does not increase with $\mathsf b$, and therefore $j^*$ does not decrease as $\mathsf b$ grows.

\emph{Proof of \textup{(G2)}.} As $j^*\ge1$, we get $k_d\le K-1$. For $k<k_d$, the depth of $k+1$ is $K-k-1\ge j^*$; hence ${\rm sep}_{k+1}\ge L^{K-k-j^*}{\rm sep}^{[j^*-1]}\ge L\,{\rm sep}_*$, and $\check\kappa_k=({\rm sep}_*/4)/{\rm sep}_{k+1}\le1/(4L)$. Summing over the levels, $\sum_{k\le j<k_d}\check\kappa_j\le(1/4)\sum_{t\ge1}L^{-t}=1/(4(L-1))$.

\emph{Proof of \textup{(G4)}.} Let $i\ne j$, and let $p$, $p'$ be plaquettes with $f_i(x(p))\ne0$ and $f_j(x(p'))\ne0$. Write $z$, $z'$ for their base sites in bare units; these lie in the bare boxes of two cubes $Q$, $Q'$ of $\mathbb T_k$, at distance $D$ from each other. Since $x(p)=a_K(z+(e_\mu+e_\nu)/2)$, we have $|x(p)-x(p')|\le a_K(|z-z'|+1/2)$, and therefore $|z-z'|\ge dL^K-1/2$. On the other hand, two sites lying in boxes of side $ML^k$ at distance $D$ satisfy $|z-z'|\le(D+1)ML^k-1$. It follows that $D+1>dL^{K-k}/M$, that is, $D\ge\lfloor dL^{K-k}/M\rfloor$. Fattening both sets $S_i(k)$ and $S_j(k)$ by $\check r(x_k)$ lowers the distance by at most $2\check r(x_k)$, so the fattened sets lie at distance $\ge{\rm sep}_k+2$.

\emph{Proof of \textup{(R)}.} Using \textup{(r)} in the second inequality, we get \[\lceil(r(x_k)+3)/L\rceil+1\le(r(x_k)+3)/L+2\le(r(x_{k+1})+\log L/960+4)/L+2\le r(x_{k+1})+3,\] because $(\log L/960+4)/L\le1$.

\emph{Proof of \textup{(C)}.} For each $i$ we pick a cube $c_i$ lying both in $X$ and in the $i$-th set, and we put $h:=\lceil\sigma/2\rceil-1$. Then $2h<\sigma$, so the balls $\{c:{\rm dist}(c,c_i)\le h\}$ are pairwise disjoint. Since $X$ is connected, it links $c_i$ to some cube whose distance from $c_i$ is $\ge\sigma>h$; along a chain of neighbouring cubes, the distance from $c_i$ changes by at most $1$ from one link to the next. Hence, for each of the values $0,1,\dots,h$, the set $X$ contains a cube at that distance from $c_i$. Counting these cubes in the disjoint balls, we find $\#X\ge j(h+1)=j\lceil\sigma/2\rceil\ge j\sigma/2$.

Now let $X\in\mathcal P$. Here ${\rm hull}(X)$ is a member of $\mathcal D$ that contains $X$ and has the least number of cubes among such members, and $|X|_c=\#{\rm hull}(X)$; for these notions see item~0 of \S\ref{Y-B-1}, on names, and S4\,(4) of the proof of \Cref{Y-B-C1}, on the hull, as well as item~1 of \S\ref{Y-G0-1}, and, for $|X|_c$, item~2 of \S\ref{Y-J-1}. Thus ${\rm hull}(X)$ is connected and meets every set that is met by $X$, and the first half of \textup{(C)}, applied to ${\rm hull}(X)$, gives the second half.

\emph{Proof of \textup{(A)}.} The barycentres at which $f_1\ne0$ lie within $R_f$ of some point, and so the base sites lie within $R_f+a_K/2$ of it; in bare units, each coordinate of these sites ranges over an interval of length $\lambda=2R_fL^K+1$. The integer points of an interval of length $\lambda$ meet no more than $\lambda/s+2$ intervals of a grid of intervals of length $s$. For $s=ML^{k+1}$ this number is at most $2R_fL^{K-k-1}/M+1/(ML^{k+1})+2\le2R_fL^{K-k-1}/M+3$. The fattening by $\check r(x_{k+1})$ adds $2\check r(x_{k+1})$ to it.

\emph{Step 2: the run, \textup{(J0)} and \textup{(J1)}.} We show the following. For $0\le k<K$, \textup{(SD)} holds at $(x_k;\ell_k)$ with $\ell_k'^{\rm c}=\ell_{k+1}$, and $x_k\ge Lx_s\ge\max\{x_9^U,x_3,x_4\}$. For $0\le k\le K$, $\bar u_k$ is real and lies in $M_1(x_k;\kappa_1,c_h)$ with the strong cap, and $\rho[\bar u_k;\ell_k]>0$. For $0\le k\le K$, every $\delta_S(k)$ lies in the space of members $\mathfrak M(x_k)$ of \textup{(o6)} of item~2, and $\|\delta_S(k)\|^{{\rm st},(\kappa_1)}$ is finite. Finally, \textup{(J0)} and \textup{(J1)} hold.

By \Cref{Y-H4-U4}\,(a), for $|\nu|_\infty\le15.99c_\nu\varepsilon(\beta)$ we have $\exp\{-\beta\sum_p(1-\nu_p)s_p\}=\rho[u_0(\nu);\emptyset]$ at every field. By \Cref{Y-H4-U4}\,(b), \textup{(SD)} holds at $(x_k;\ell_k)$ for every $0\le k<K$, with $\ell_k'^{\rm c}=\ell_{k+1}$; for $k=0$ this is $(\beta;\emptyset)$, with $\ell'^{\rm c}=\ell_1$. The empty rows are admitted by \textup{(SD)} of item~2 of \S\ref{Y-H-1}, and also by the setting at the opening of item~2 of \S\ref{Y-H1-1}; and \Cref{Y-H-C9} itself holds at $\ell=\emptyset$, for every complex $u$ within $16$ caps that has the weak cap. At $k=0$, Steps~2 to~10 apply Thms.~\ref{Y-H1-U1}, \ref{Y-H3-U3} and~\ref{Y-H5-Af} at $(\beta;\emptyset)$.

By \Cref{Y-H4-U4}\,(h), on ${\rm Pd}_0$ the honest family $u_k(z)$ takes values in $M_1(x_k;\kappa_1,c_h)$ with the strong cap when $k\ge1$, and in $M_2$ with the strong cap when $k=0$ (by \textup{(a1)}, under which the $a$-part $E^a_\nu$ is at most $0.01+1.99$ caps $c_a(\beta)$ on ${\rm Pd}_0$, and by \textup{(a2)}, both of \Cref{Y-H4-U4}); and it is holomorphic. At $z=0$ it is real (\textup{(a3)} of \Cref{Y-H4-U4}; \Cref{Y-H-C9}\,(d)) and lies within one cap (by \textup{(a1)}, which bounds the same $a$-part by $0.01$ caps at $z=0$); and $\rho[\bar u_k;\ell_k]>0$ for $0\le k\le K$, by the part of \Cref{Y-H4-U4}\,(h) that concerns $z=0$. Further, $x_k\ge L\mathsf b\ge Lx_s$ for $k<K$, and $Lx_s\ge\max\{x_9^U,x_3,x_4\}$ by \textup{(O13)} of item~(1) of \S\ref{Y-H0-1}.

We apply \Cref{Y-H5-LS} on the open polydisc ${\rm Pd}(1.99/{\rm amp}_W)$, with the strong weighted sup norm at $\kappa_1$. Because $\|f_i\|_\infty\le1$, this polydisc is contained in the interior of ${\rm Pd}_0$, inside the open set ${\rm Pd}_0^o$ of \Cref{Y-H4-U4}\,(h), on which every value is holomorphic. The lemma shows that each coefficient $\delta_S(k)$ lies in the space of members and has a finite norm. Consequently $u_k^{\rm pol}(z)$ is a family of norm $<16$ near $z=0$; by its form, $\partial_{S'}u_k^{\rm pol}(0)=\delta_{S'}(k)=\partial_{S'}u_k(0)$ for every $S'$, and $u_k^{\rm pol}(0)=\bar u_k$. Now Step~0\,(d), applied with $\Phi$ a value of $F_{\ell_k}$, yields the first half of \textup{(J1)}; applied with $\Phi=\mathsf V'_{\ell_k}$, it yields the second half.

For \textup{(J0)} we apply \Cref{Y-H2-U2} with $f=f_i$, in the setting \textup{(L5*)}, \textup{(X5)} of the table (item~(1) of \S\ref{Y-H0-1}). Step~0\,(a), applied with $x=x_k$ and $D=\mathbb D_k[f_i]/\mathsf b$, gives the bound $|\cdot|_a\le c_Q^{\rm su}x_k^{1/2}K_5/(\mathsf bE_ap_0(x_k))={\rm amp}(x_k)$. If a cube does not belong to $S_i(k)$, then for no $y$ in that cube does $B^k(y)$ contain the base site of a plaquette $p$ with $f_i(x(p))\ne0$; since $\mathbb D_{k,y}[f_i]$ is linear and depends only on the values of $f_i$ there, it vanishes.

\emph{Step 3: the reach and \textup{(J2)}.} A family $u(z)$ at level $k$ is said to \emph{have $(P_k)$} if, for every $i$, the value $u(z)(X)$ is independent of $z_i$ for all $X$ such that $X^{+\check r(x_k)}$ misses $S_i(k)$. We claim: if $u(z)$ has $(P_k)$, where $k<K$, then $F_{\ell_k}(u(z))$ has $(P_{k+1})$, and moreover $v_k(Y;u(z))$ is independent of $z_i$ for all $Y$ such that $Y^{+\check r(x_{k+1})}$ misses $S_i(k+1)$. From this claim we shall obtain \textup{(J2)}.

\emph{Proof of the claim.} We fix $z_j$ for $j\ne i$. \Cref{Y-H1-U1}\,(f5), applied with $S=S_i(k)$ and $\hat\rho=\check r(x_k)$, shows that $F_{\ell_k}(u(z))(Y)$ and $v_k(Y;u(z))$ are constant in $z_i$ for every $Y$ such that $Y^{+\hat\rho'}$ does not meet ${\rm cl}(S_i(k))=S_i(k+1)$, where $\hat\rho'=\lceil\check r(x_k)/L\rceil+1\le\check r(x_{k+1})$ by Step~1\,(R). If $Y^{+\check r(x_{k+1})}$ does not meet $S_i(k+1)$, then neither does $Y^{+\hat\rho'}$.

\emph{Proof of \textup{(J2)} at $k=0$.} The $a$-part $E^a_\nu(\{\square\})=\sum_{p\in\square}[q_p+\beta\nu_ps_p]$ (item~3 of \S\ref{Y-H4-1}) is affine in $z$. Hence $\delta_i^a(0)(\{\square\})=(\beta/\mathsf b)\sum_{p\in\square}f_i(x(p))s_p$, which vanishes unless the cube lies in $S_i(0)$, and $\delta_S^a(0)=0$ for $|S|\ge2$. The $h$-part is $0$. As for the $w$-part, \Cref{Y-H4-U4}\,(f) shows that it is $0$ unless $X^{+r(x_0)}$ meets $S_i(0)$ for each $i\in S$; and $r\le\check r$.

\emph{Proof of \textup{(J2)} from $k$ to $k+1$.} By \textup{(J2)} at level $k$, the family $u_k^{\rm pol}$ has $(P_k)$: whenever $X^{+\check r(x_k)}$ misses $S_i(k)$, each term $z^S\delta_S(k)(X)$ with $i\in S$ is $0$. By the claim, $F_{\ell_k}(u_k^{\rm pol}(z))(Y)$ is then independent of $z_i$ whenever $Y^{+\check r(x_{k+1})}$ misses $S_i(k+1)$, so that for such $Y$ its coefficient of $z^S$ with $i\in S$ vanishes. By \textup{(J1)}, this coefficient is $\delta_S(k+1)(Y)$.

\emph{The count.} Suppose that $\delta_S(k)(X)\ne0$. Then, for the $a$- and $h$-parts, $X$ meets each of the $|S|$ sets $S_i(k)^{+\check r(x_k)}$, and ${\rm hull}(X)$ does so for the $w$-part; by \textup{(G4)} these sets lie pairwise at distance $\ge{\rm sep}_k+2$. Step~1\,(C) therefore bounds $\#X$ in the first case, and $|X|_c$ in the second, from below by $|S|\lceil({\rm sep}_k+2)/2\rceil\ge|S|\,{\rm sep}_k/2$.

\emph{Step 4: level $0$.} For every $K\ge1$ we treat \textup{(J3)}, \textup{(J4)} and \textup{(J5)} at $k=0$ by proving the following: $t_0:=|{\rm rem}_{0,i}|_a\mathsf b/(u_3(x_0)K_5)\le0.01$; $\delta_i^h(0)=0$; $|\delta_i^w(0)|_w^{\rm st}\le{\rm amp}(x_0)/815$; $\|\delta_i(0)\|^{\rm st}\le0.2\,{\rm amp}(Lx_0)$; and $\sum_{|S|\ge2}\mathfrak r_0^{|S|}\|\delta_S(0)\|^{{\rm st},(\kappa_1)}\le4\cdot10^{-7}$.

By \Cref{Y-H4-U4}\,(f), $\delta_i^a(0)=s_{x_0}[\mathbb D_0[f_i]/\mathsf b]+{\rm rem}_{0,i}$, where $|{\rm rem}_{0,i}|_a\le0.01/\mathsf b$; and ${\rm rem}_{0,i}\in T^a(x_0)$, being the difference of two elements of that space. As $u_3(x_0)K_5\ge1$ by Step~0\,(b), we get $t_0\le0.01$. The $h$-part is $0$. \Cref{Y-H4-U4}\,(f) for $S=\{i\}$ gives $|\delta_i^w(0)|_w^{\rm st}\le{\rm amp}_W/15.99\le{\rm amp}(x_0)/(51\cdot15.99)\le{\rm amp}(x_0)/815$. Combining these bounds with \textup{(J0)} and \textup{(W3)}, we obtain $\|\delta_i(0)\|^{\rm st}\le1.002\,{\rm amp}(x_0)\le(1.002/5.01){\rm amp}(Lx_0)=0.2\,{\rm amp}(Lx_0)$, and $n\,\mathfrak r_0\cdot0.2\,{\rm amp}(Lx_0)=3.5\cdot0.2=0.7$.

For the mixed jets, the $a$- and $h$-parts vanish, and $\|\delta_S(0)\|^{{\rm st},(\kappa_1)}\le({\rm amp}_W/15.99)^{|S|}$, since the factor $e^{-(m_\Xi/2)p_0^2\#X}$ does not exceed $1$. Now $\mathfrak r_0\,{\rm amp}_W/15.99\le3.5/(n\cdot5.01\cdot815)\le0.00086/n=:q_0$, and so $\sum_{|S|\ge2}q_0^{|S|}=(1+q_0)^n-1-nq_0\le e^{0.00086}-1.00086\le4\cdot10^{-7}$. Since $\kappa_0\le\kappa_1$, the norm at $\kappa_0$ is not larger than the norm at $\kappa_1$. The bounds of this step are independent of the separation of the supports.

\emph{Step 5: the $a$-part of the single jets, \textup{(J3)}.} For $k<K$ we show the identity ${\rm rem}_{k+1,i}=DA_{\ell_k}(\bar u_k)[{\rm rem}_{k,i}]+n[\bar u_k;\mathbb D_k[f_i]/\mathsf b]$, and we show that $t_k:=|{\rm rem}_{k,i}|_a^{(\kappa_1)}\mathsf b/(u_3(x_k)K_5)$ satisfies $t_{k+1}\le0.2431t_k+1$; it follows that $t_k\le2$ for every $k$.

By \Cref{Y-H1-U1}\,(f1), $E'^{\rm new}$ does not depend on any member, and $A_\ell(u):=E'^{\rm old}[u]$ depends on $E^a$ only; thus the $a$-channel is autonomous. From \textup{(J1)} for $S=\{i\}$ we derive $\delta_i^a(k+1)=DA_{\ell_k}(\bar u_k)[\delta_i^a(k)]$, as follows. By \textup{(J1)} and \Cref{Y-H5-LJ}, applied with $S=\{i\}$, where a single partition occurs, the coefficient of $z_i$ in $E'^a[u_k^{\rm pol}(z)](Y,\mathsf W')$ equals $\frac d{ds}|_0\,E'^a[\bar u_k+s\delta_i(k)](Y,\mathsf W')$. By \Cref{Y-H1-U1}\,(f1), this is $\frac d{ds}|_0\,E'^{\rm old}[\bar u_k+s(\delta_i^a(k),0,0)](Y,\mathsf W')$. Here $\delta_i^a(k)\in T^a(x_k)$: being the $a$-part of a member whose norm is finite (Step~2), it satisfies the clauses of \textup{(ca)} without a size, and all its values are bounded. Also $A_a(\bar u_k)\le1<16$ at $p=(\kappa_1,c_h)$. Hence, by \Cref{Y-H1-U1}\,(h), at $x_k\ge x_9^U$ and with the data \textup{(SD)} that Step~2 provides at $(x_k;\ell_k)$, this derivative is given by the series $DA(\bar u_k)[\delta_i^a(k)](Y,\mathsf W')$ of item~1 of \S\ref{Y-H1-1}, taken at the rows $\ell_k$; this series is the $DA_{\ell_k}$ of \textup{(o8)} of item~2 of the present section, the object of \Cref{Y-H1-U1}\,(h3) and of \Cref{Y-H3-U3}.

By \Cref{Y-H1-U1}\,(h) this series is linear in $\delta$; we split it according to $\delta_i^a(k)=s_{x_k}[\mathbb D_k[f_i]/\mathsf b]+{\rm rem}_{k,i}$, where both summands lie in $T^a(x_k)$ (by \textup{(J0)}, and by \textup{(J3)} at $k$). Here \textup{(J3)} at level $k$ serves as the induction hypothesis. To the first summand we apply \Cref{Y-H3-U3} at $x=x_k$, with the data \textup{(SD)} that Step~2 provides at $(x_k;\ell_k)$ (the empty rows at $k=0$), with $u=\bar u_k\in M_1(x_k;\kappa_1,c_h)$, which lies inside $M_8$, and with $D=\mathbb D_k[f_i]/\mathsf b$, which is symmetric, of range $\le\varrho(x_k)$ and of norm $\le K_5/\mathsf b$ by \textup{(J0)}. Its $D'={\rm cut}_{\varrho(x_{k+1})}(T_{\ell_k}D)$ is $\mathbb D_{k+1}[f_i]/\mathsf b$, by the definition of the source forms in item~1 of \S\ref{Y-H2-1}. Therefore $DA_{\ell_k}(\bar u_k)[s_{x_k}[D]]=s_{x_{k+1}}[\mathbb D_{k+1}[f_i]/\mathsf b]+n[\bar u_k;D]$, and together with the splitting this gives the identity for ${\rm rem}_{k+1,i}$ stated at the beginning of the step.

Moreover ${\rm rem}_{k+1,i}\in T^a(x_{k+1})$, since it is $\delta_i^a(k+1)$ (Step~2) minus $s_{x_{k+1}}[\cdot]$ (Step~0\,(a), which is valid at $x_{k+1}\ge \mathsf b\ge x_s$, and hence also when $k+1=K$); \Cref{Y-H3-R3} gives the same, ${\rm rem}'\in T^a(x')$. \Cref{Y-H3-U3}\,(a) at $\tilde\kappa=\kappa_1$ gives $|n|_a\le u_3(x_{k+1})K_5/\mathsf b$ at $x_{k+1}$. \Cref{Y-H1-U1}\,(h3), applied with $u=\bar u_k$, $p=(\kappa_1,c_h)$, $\delta={\rm rem}_{k,i}$ and $B=|{\rm rem}_{k,i}|_a^{(\kappa_1)}$, in the setting \textup{(L3su)}, bounds the cell sum of $DA(\bar u_k)[{\rm rem}_{k,i}]$ by $\vartheta_aBc_a(x_{k+1})$; by Step~0\,(c), the norm obeys the same bound. Since $u_3(x_k)=L^{1/4}u_3(x_{k+1})$ and $\vartheta_aL^{1/4}=0.2431$, we conclude that $t_{k+1}\le0.2431t_k+1$. Now $t_0\le0.01$ (Step~4), and $t\le2$ implies $0.2431t+1\le1.4862\le2$; the fixed point is $1/0.7569\le1.3212$. Finally, $2u_3(x_k)K_5/\mathsf b=2\check\omega_3(x_k){\rm amp}(x_k)$ by the identity of \textup{(o5)} of item~2, and $\check\omega_3(x_k)\le10^{-3}$ by \textup{(W3)}. This condition is imposed for all $x\ge Lx_s$ and consists of the two inequalities $\check\omega_3(x),\ \check\omega_3(x/L)\le10^{-3}$; between them they give $\check\omega_3(x)\le10^{-3}$ at every $x\ge x_s$, through the first inequality from $Lx_s$ on and through the second on $[x_s,Lx_s)$. At $x_K=\mathsf b$ the bound comes from the second inequality at $x=L\mathsf b$.

\emph{Step 6: the $h$- and $w$-parts of the single jets, \textup{(J4)}.} Let $H_k:=\max_i\max\{|\delta_i^h(k)|_h^{(\kappa_1)},|\delta_i^w(k)|_w^{{\rm st},(\kappa_1)}\}$. We show that $H_{k+1}\le(1/15)\,1.002\,{\rm amp}(x_k)+((\eta_h+\eta_w)(x_k)/15)H_k$, and hence that $H_k\le0.12\,{\rm amp}(Lx_k)$.

For $k=0$ the bound $H_0\le{\rm amp}(x_0)/815$ follows from the second and third claims of Step~4, namely $\delta_i^h(0)=0$ and $|\delta_i^w(0)|_w^{\rm st}\le{\rm amp}(x_0)/815$. Since ${\rm amp}(x_0)/815\le{\rm amp}(Lx_0)/(815\cdot5.01)\le0.12\,{\rm amp}(Lx_0)$ by \textup{(G3)}, we have $H_0\le0.12\,{\rm amp}(Lx_0)$, which is the estimate of this step at $k=0$. For the passage from $k$ to $k+1$, \textup{(J1)} for $S=\{i\}$ gives $(\delta_i^h,\delta_i^w)(k+1)=D(F^h,F^w)_{\ell_k}(\bar u_k)[\delta_i(k)]$, which is linear in $\delta_i(k)=(\delta^a,0,0)+(0,\delta^h,0)+(0,0,\delta^w)$. Here we set $(F^h,F^w)_\ell(u):=(E'^h,w')$, the $h$- and $w$-members of $F_\ell(u)$, and we write $D$ for the derivative of \Cref{Y-H5-LH}\,(a). \Cref{Y-H5-Af}\,(iv), at the base $\bar u_k$, with $B=1$ and the rate $\kappa_1$, bounds each of the three contributions to the $h$- and $w$-parts by a factor times the size of its input: the contribution of $a$ by $|\delta_i^a(k)|_a/15$, that of $h$ by $\eta_h(x_k)\,|\delta_i^h(k)|_h/15$, and that of $w$, with the input in strong units, by $\eta_w(x_k)\,|\delta_i^w(k)|_w^{\rm st}/15$; the $h$-output is measured in the cell sum, in caps $c_h(x_{k+1})$, and the $w$-output in strong units at $x_{k+1}$. By Step~0\,(c), this cell sum bounds $|\delta_i^h(k+1)|_h^{(\kappa_1)}$.

By \textup{(J0)} and \textup{(J3)}, $|\delta_i^a(k)|_a\le1.002\,{\rm amp}(x_k)$. By \textup{(W2)} at $x_k\ge Lx_s$, $\eta_h+\eta_w\le1/L$; and by \textup{(G3)}, ${\rm amp}(Lx_k)\le L^{1/2}{\rm amp}(x_k)$. So, if $H_k\le0.12\,{\rm amp}(Lx_k)$, then \[\FitDisplay{\dimexpr\linewidth-1pt\relax}{H_{k+1}\le0.0668\,{\rm amp}(x_k)+(0.12/15)L^{-1}L^{1/2}{\rm amp}(x_k)=(0.0668+0.008L^{-1/2}){\rm amp}(x_k)\le0.12\,{\rm amp}(x_k)=0.12\,{\rm amp}(Lx_{k+1}).}\]

\emph{A single bounded complex source.} The single-jet arguments of Steps~2 and~4--6 also apply to a single complex bounded source $f$, with no condition on its support. With $\|f\|_\infty\le1$, \Cref{Y-H4-U4}\,(h) gives the honest family for $|z|<1.99/\mathrm{amp}_W$, so Step~2 applies to its first jet. \Cref{Y-H4-U4}\,(f) gives the level-zero bounds for complex $f$, \Cref{Y-H2-U2} gives its complex symmetric source form, and all derivatives in Steps~5 and~6 are complex linear. The base $\bar u_k$ remains the sourceless run. The arguments therefore give the splitting and estimates of \textup{(J0)}, \textup{(J3)} and \textup{(J4)} for this single jet on every $\mathbb T_k$. In particular,
\[
 \|\delta[f](k)\|^{\mathrm{st},(\kappa_1)}
 \le\max\{1.002\,\mathrm{amp}(x_k),0.12\,\mathrm{amp}(Lx_k)\}
 \le0.2\,\mathrm{amp}(Lx_k).
\]
Linearity gives the factor $\|f\|_\infty$ for a general $f$.

\emph{Step 7: the early levels, \textup{(J5)}.} We show: if \textup{(J5)} holds at some $k<k_d$, then it holds at $k+1$; in fact $\sum_{|S|\ge2}\mathfrak r_{k+1}^{|S|}\|\delta_S(k+1)\|^{{\rm st},(\kappa_{k+1})}\le0.0160$.

At $k=0$, \textup{(J5)} is given by Step~4.

\emph{Single jets at level $k+1$.} By \textup{(J0)}, \textup{(J3)}, \textup{(J4)} and \textup{(G3)}, \[\|\delta_i(k+1)\|^{{\rm st},(\kappa_1)}\le\max\{1.002\,{\rm amp}(x_{k+1}),0.12\,{\rm amp}(x_k)\}\le0.2\,{\rm amp}(x_k)=0.2\,{\rm amp}(Lx_{k+1}),\] where we used $1.002/5.01=0.2$; and $n\,\mathfrak r_{k+1}\cdot0.2\,{\rm amp}(Lx_{k+1})=0.7$.

\emph{The family at level $k$.} We have $\|\bar u_k\|\le1$. On ${\rm Pd}(\mathfrak r_k)$, using Step~0\,(c) for the single jets and for $\bar u_k$, and taking $0.7$ from the first clause of \textup{(J5)}, \[\|u_k^{\rm pol}(z)-\bar u_k\|^{{\rm st},(\kappa_k)}\le n\,\mathfrak r_k\max_i\|\delta_i(k)\|^{{\rm st},(\kappa_1)}+\sum_{|S|\ge2}\mathfrak r_k^{|S|}\|\delta_S(k)\|^{{\rm st},(\kappa_k)}\le0.7+0.21=0.91.\] Hence $u_k^{\rm pol}(z)$ and $\bar u_k$ lie in $M_2(x_k;\kappa_k,c_h)$; and $(\kappa_k,c_h)\in\mathbb P$, by \textup{(G2)}.

\emph{Mixed jets.} \Cref{Y-H-C9}\,(e) at $(\kappa_k,c_h)$ (\Cref{Y-H1-U1}\,(U)) gives, at $x_{k+1}$, \[\sup_{{\rm Pd}(\mathfrak r_k)}\|F_{\ell_k}(u_k^{\rm pol}(z))-F_{\ell_k}(\bar u_k)\|^{{\rm st},(\kappa_k)}\le0.4\cdot0.91=0.364.\] By \Cref{Y-H1-U1}\,(c+), the difference on the left, as a function of $z$, is a family, and by \textup{(J1)} its coefficient of $z^S$ is $\delta_S(k+1)$. \Cref{Y-H5-LS}, that is, Cauchy's estimate in the variables of $S$ on circles of radius $r<\mathfrak r_k$, with the remaining variables put equal to $0$, followed by the limit $r\to\mathfrak r_k$, yields $\|\delta_S(k+1)\|^{{\rm st},(\kappa_k)}\le0.364\,\mathfrak r_k^{-|S|}$. By \textup{(J2)} at $k+1$, $\delta_S(k+1)$ is carried by sets whose $\#X$ (for the $a$- and $h$-parts) or $|X|_c$ (for the $w$-part) is at least $|S|\,{\rm sep}_{k+1}/2$; since $\kappa_{k+1}=\kappa_k-\check\kappa_k$, it follows that \[\|\delta_S(k+1)\|^{{\rm st},(\kappa_{k+1})}\le e^{-\check\kappa_k(|S|{\rm sep}_{k+1}/2-1)}\cdot0.364\,\mathfrak r_k^{-|S|}.\]

By \textup{(G3)}, $\mathfrak r_{k+1}\le L^{1/2}\mathfrak r_k$. Let $q:=L^{1/2}e^{-\check\kappa_k{\rm sep}_{k+1}/2}=L^{1/2}/(4nL^{1/2})=1/(4n)$. Then $\mathfrak r_{k+1}^{|S|}\|\delta_S(k+1)\|^{{\rm st},(\kappa_{k+1})}\le0.364\,e^{\check\kappa_k}q^{|S|}$, and therefore \[\FitDisplay{\dimexpr\linewidth-1pt\relax}{\begin{aligned}
&\sum_{|S|\ge2}\mathfrak r_{k+1}^{|S|}\|\delta_S(k+1)\|^{{\rm st},(\kappa_{k+1})}
\\
&\quad{}\le0.364\,e^{1/4}[(1+q)^n-1-nq]
\\
&\quad{}\le0.364\cdot1.28403\cdot(e^{1/4}-5/4)
\\
&\quad{}\le0.364\cdot1.28403\cdot0.034026
\\
&\quad{}\le0.0160\le0.21.
\end{aligned}}\] Here we used that $(1+1/(4n))^n$ increases to $e^{1/4}$: the supremum of the bracket is $e^{1/4}-5/4\le0.034026$, and no $n$ attains it.

\emph{Step 8: the late levels and the short runs, \textup{(J6)}.} We prove \textup{(J6)} for every $K\ge1$. At $k=k_d$ it is \textup{(J5)}, with $\kappa_{k_d}=\kappa$.

\emph{From $k$ to $k+1$, for $k_d\le k<K$.} We apply \Cref{Y-H5-Af}\,(i), (ii) at $p=(\kappa,c_h)$ to the pair of families consisting of $u_k(z)$ and the constant family $\bar u_k$ on ${\rm Pd}(\mathfrak r_*)$; they lie within $B=1.91$ caps at $k=k_d$, and within $B=1$ at the later levels. The image is a family in $M_1(x_{k+1};\kappa,c_h)$ with the strong cap, because $0.35+0.248\cdot1.91\le0.824\le1$. The strong distance of the images is at most $\vartheta[B]$ times the weak distance of the inputs, where $\vartheta[1.91]=3.968/14.09\le0.2817$ and $\vartheta[1]\le0.2646$; both values are $\le0.4$. The weak distance does not exceed the strong one. Moreover $F_{\ell_k}(\bar u_k)=\bar u_{k+1}$.

\emph{Jets.} If $u_k(z)$ and the honest family have the same $\partial_{S'}$ at $0$ for every $S'$, then Step~0\,(d), applied to each value, shows that their images have the same property; at $k=k_d$ this hypothesis holds by the form of $u^{\rm pol}$.

\emph{Short runs.} If $K<j^*$, then $k_d=0$. In that case \textup{(J5)} at $k=0$ is given by Step~4, whose bounds are independent of the separation; and the families start from $u_0^{\rm pol}$ on ${\rm Pd}(\mathfrak r_0)$, where $\mathfrak r_0=\mathfrak r_*\ge\mathfrak r(L^{j^*}\mathsf b)$ by \textup{(G3)}. There is then no early level. In every case $k_d\le K-1$, by \textup{(G2)}.

\emph{The base points.} For $0\le k\le K$ they are provided by Step~2 (\Cref{Y-H4-U4}\,(h)).

\emph{Step 9: the identity \textup{(J7)}.} Since $\beta\nu_p(z)s_p=L^K\sum_iz_if_i(x(p))s_p$, we have \[\exp\{-\beta\sum_ps_p+\sum_iz_i\mathcal O_K(f_i)\}=\exp\{-\beta\sum_p(1-\nu_p(z))s_p\},\] which is the weight of \Cref{Y-H4-U4}. With the honest family, \Cref{Y-H4-U4}\,(h) gives $Z_K(z)=\exp\{\sum_{1\le j\le K}V_j(z)\}I[u_K(z)]$ on ${\rm Pd}_0$; here the kernels are normalised, and Fubini's theorem is used. Also $\rho[\bar u_K;\ell_K]>0$ (\Cref{Y-H4-U4}\,(a3) at $\nu=0$; \Cref{Y-H-C9}\,(d), applied step by step), so that $I[\bar u_K]>0$. Therefore, near $z=0$, $\log Z_K(z)=\sum_jV_j(z)+\log I[u_K(z)]$ up to an additive constant. By \Cref{Y-H1-U1}\,(g1), $V_{k+1}(z)=V_0^++\sum_Yv_k(Y;u_k(z))$, a finite sum because $\mathbb T_{k+1}$ is finite; and $V_0^+$ does not depend on any member (\Cref{Y-H1-U1}\,(f1)), so $\partial_1$ annihilates it. This proves \textup{(J7)} with the honest family.

\emph{Representatives.} These have the same $\partial_{S'}$ at $0$ as the honest family, for all $S'$ (\textup{(o2)} of item~2; \textup{(J6)}). Applied with $S=\{1,\dots,n\}$ and $\Phi=v_k(Y;\cdot)$, Step~0\,(d) shows that the derivative is unchanged. As to the last term, $\partial_1\cdots\partial_n\log I[u(z)](0)$ is a polynomial in the numbers $\partial_{S'}I[u(z)](0)/I[u(0)]$, $S'\subset\{1,\dots,n\}$, and we apply Step~0\,(d) with $\Phi=I$ for every $S'$.

\emph{Vanishing.} The representatives have the property $(P_k)$ of Step~3: at the levels $k\le k_d$ by \textup{(J2)}, and at the levels $k>k_d$ by the claim of Step~3, applied from one level to the next. By that claim, $v_k(Y;u_k(z))$ is independent of $z_i$ when $Y^{+\check r(x_{k+1})}$ misses $S_i(k+1)$, and then $\partial_1\cdots\partial_n$ annihilates it.

\emph{Step 10: the local vacuum numbers, \textup{(J8)}.} Fix $k<K$; set $\hat{\mathfrak r}_k:=\mathfrak r_{\min(k,k_d)}$, and let $\hat\kappa:=\kappa_k$ for $k<k_d$ and $\hat\kappa:=\kappa$ for $k\ge k_d$. On ${\rm Pd}(\hat{\mathfrak r}_k)$ the representative takes values in $M_2(x_k;\hat\kappa,c_h)$ (\textup{(J5)}, \textup{(J6)}), and $z\mapsto v_k(Y;u_k(z))$ is holomorphic on this polydisc (\Cref{Y-H1-U1}\,(c+)). For $r<\hat{\mathfrak r}_k$ we have \[\partial_1\cdots\partial_nv_k(Y;u_k(z))(0)=(2\pi i)^{-n}\oint_{|z_1|=\dots=|z_n|=r}v_k(Y;u_k(z))\prod_idz_i/z_i^2.\]

Let $\check Y$ be the family of those $Y$ that remain in \textup{(J7)}. We take the sum over $Y$ inside the integral: \[\sum_{Y\in\check Y}|\partial_1\cdots\partial_nv_k(Y;\cdot)(0)|\le r^{-n}\sup_{|z_i|=r}\sum_{Y\in\check Y}|v_k(Y;u_k(z))|.\] Every $Y\in\check Y$ contains a cell of $S_1(k+1)^{+\check r(x_{k+1})}$, a set with at most $N_1(k)$ cells (Step~1\,(A)). Every $Y\in\check Y$ is connected and satisfies $\#Y\ge m_k$: for $k<k_d$, $Y$ meets each of the $n$ sets $S_i(k+1)^{+\check r(x_{k+1})}$, which lie pairwise at distance $\ge{\rm sep}_{k+1}+2$ by \textup{(G4)}, so that Step~1\,(C) gives $\#Y\ge n\lceil{\rm sep}_{k+1}/2\rceil\ge\lceil n\,{\rm sep}_{k+1}/2\rceil$, the number $n\lceil{\rm sep}_{k+1}/2\rceil$ being an integer $\ge n\,{\rm sep}_{k+1}/2$; for $k\ge k_d$ we have $m_k=1$ by the definition of $m_k$ in \textup{(J8)}, and $\#Y\ge1$ because $Y$ contains a cell.

At each $z$, \Cref{Y-H1-U1}\,(g2) at $p=(\hat\kappa,c_h)$ and $A=2$ gives \[\sum_{Y\in\check Y}|v_k(Y;u)|\le N_1(k)e^{-\hat\kappa(m_k-1)}\max_{\rm cells}\sum_{Y\ni\text{the cell}}e^{\hat\kappa(\#Y-1)}|v_k(Y;u)|\le N_1(k)e^{-\kappa(m_k-1)}2C_gc_a(x_k).\] Letting $r\to\hat{\mathfrak r}_k$, we obtain \textup{(J8)}. No entropy factor appears, because the factor $e^{-\kappa(m_k-1)}$ is drawn out of the weight of \Cref{Y-H1-U1}\,(g2). At $k=0$ the step is the one at $(\beta;\emptyset)$, and \Cref{Y-H1-U1}\,(g1), (g2) apply there.\end{proof}

\subsection{Comparison of sourced runs}\label{Y-H7-1}

We compare two runs of the drawn step, of lengths $K\le K'$, which start from Wilson's weight with $n$ separated sources and end on one and the same torus; along each run we follow the base points of the members and their multi-affine jets at $z=0$, on polydiscs whose radii grow with the level $k$ as long as $k\le k_d$, and on one fixed polydisc for $k_d\le k\le K$, namely $\mathrm{Pd}(\mathfrak r_*/2)$ where the two runs are compared and $\mathrm{Pd}(\mathfrak r_*)$ for one run alone. At the level $k_0$ we rewrite the long run over the rows of the short one by the map $\mathbb B_D$, $D=\tilde{\mathsf E}_{m+k_0}-\mathsf E_{k_0}$; from there on a single map $F_{\ell_k}$ carries both runs, and the distances between their base points, between their jets and between their families are bounded by a numeral times $\vartheta^{k-k_0}$ plus a numeral times $\Omega^{\#}$. From this we derive the bound \textup{(D3)} at the end of the runs, the bounds \textup{(D4)} and \textup{(D5)} on the vacuum numbers, the identity \textup{(D6)} for $S^T_{n;K'}-S^T_{n;K}$, and in \textup{(D7)} the local boundedness in $\mathsf b$ of all constants.

\begin{remark}[Content, place and setting]

The tags \textup{(D0)} to \textup{(D7)} denote the clauses of \Cref{Y-H7-Df}.

\Cref{Y-H7-Df} is the fourth of four statements on holomorphic families of complex members; the other three are \Cref{Y-H5-Af}, \Cref{Y-H5-Bc} and \Cref{Y-H5-Cf}, all of them in \S\ref{Y-H5-1}.

Our setting consists of the group $G=\operatorname{SU}(2)$, dimension three, the drawn step, and one fixed torus.

\end{remark}

\emph{Item 1: notation, norms and the table.}

\begin{definition}[Item 1.1: what is kept]

From \S\ref{Y-J-1} we keep the following objects, taken at the rate $\kappa_1=\kappa+1$ of \textup{(O5)}, item~(1) of \S\ref{Y-H0-1}: the lattice, the cubes, the fields, the rows, the site forms, the norm $\|\cdot\|_{(\lambda_\flat)}$, the energy $\mathrm{En}_\ell$, the domain $\mathcal U_X(x)$, the clauses \textup{(ca)}, \textup{(ch)}, \textup{(c$\Xi$)}, the weak and the strong cap, the density $\rho[u;\ell]$, the class, the step, the runs, the partners and the depth. Here $\kappa\ge6$, and $\kappa$ is the lower end of the interval $[\kappa,\kappa_1]$ of the rates $\tilde\kappa$.

We write $x_k=L^{K-k}\mathsf b$, $\ell_k=\ell_k^{[K]}$ and $\tilde\ell_j=\ell_j^{[K']}$, and we denote by $\mathsf E_k$ and $\tilde{\mathsf E}_j$ the site forms of these rows; moreover $m=K'-K$.

For the run $K$ let $S_i(k):=S_i^{[K]}(k)$ be the $k$-blocks of $\mathrm{supp}\,f_i$ as in \textup{(o3)}, item~2 of \S\ref{Y-H6-1}; for the run $K'$ the corresponding sets are $\tilde S_i(k):=S_i^{[K']}(m+k)$. Both are sets of cubes of one torus $\mathbb T_k=\mathbb Z^3/NL^{K-k}\mathbb Z^3$, and the two may differ, since the bare lattice of $K'$ has barycentres that are missing from the bare lattice of $K$.

The numbers $\Lambda=\log_L(c_4\mathsf b^{1/2})$, $K_0(\mathsf b)=(5/2)\Lambda+9/4$, $k_0=\lceil(5K+10\Lambda)/9\rceil$, $\theta=\vartheta^{4/9}\le0.69903$, $\varepsilon_L$ and $C_{\mathsf b}=(3/\vartheta)(c_4\mathsf b^{1/2})^{\varepsilon_L}$ are entries of \textup{(O14)}, item~1 of \S\ref{Y-J-1}, taken at the rate $\kappa_1$, and $\vartheta=0.4468$ is also an entry of this \textup{(O14)}.

The symbols $F_\ell$, $\mathsf V'[\cdot]$, $v(Y;\cdot)$ and $A_\ell:=E'^{\rm old}$ refer to the one map of \Cref{Y-H1-U1}\,(U). Finally, $I[u]:=\int\rho[u;\ell_K]\,dV$.

\end{definition}

\begin{definition}[Item 1.2: norms, families and local letters]

The norms are those of \S\ref{Y-H5-1}. Let $g=(g^a,g^h,g^w)$ be a triple that satisfies at $x$ the clauses without a size, and let $\tilde\kappa\in[\kappa,\kappa_1]$ be a rate. We set \[|g|_a:=\sup_X\sup_{\mathcal U_X(x)}|g^a(X,\cdot)|\,e^{\tilde\kappa(\#X-1)}/c_a(x)\] and \[|g|_h:=\sup_X\sup|g^h(X,\cdot)|\,e^{\tilde\kappa(\#X-1)}/c_h(x),\] where in the second formula the inner supremum is taken over the tame $G$-valued fields. For the third component we put $|g|_w^{\rm wk}:=\sup_{X,V}|g^w(X,V)|\,e^{\tilde\kappa(|X|_c-1)}$, and $|g|_w^{\rm st}$ is defined in the same way with the additional factor $\exp\{c_\Xi x\varepsilon\sum_{p\in D(V)\cap X}\phi_p\}$. Then $\|g\|:=\max\{|g|_a,|g|_h,|g|_w^{\rm wk}\}$, and $\|g\|^{\rm st}$ is obtained likewise with $|g|_w^{\rm st}$ in place of $|g|_w^{\rm wk}$.

The two distances are $\mathsf D^{\rm in}_x(u_1,u_2)=\|u_1-u_2\|$ and $\mathsf D_x(u_1,u_2)=\|u_1-u_2\|^{\rm st}$. An upper index $(\tilde\kappa)$ names the rate; where no index is written, the rate is $\kappa_1$. The norm at a higher rate dominates the norm at a lower rate. We write $M_B(x;\tilde\kappa,c_h)=\{\|u\|^{(\tilde\kappa)}\le B\}$ and $\mathrm{Pd}(r)=\{z\in\mathbb C^n:\ |z_i|<r\}$; for a non-empty subset $S$ of $\{1,\dots,n\}$ we put $z^S=\prod_{i\in S}z_i$.

A family on an open set $\Omega$ within $B$ caps at $p=(\tilde\kappa,\tilde c)$ is a map $z\mapsto u(z)\in M_B(x;p)$ such that all values of $u(z)$ depend holomorphically on $z$; its coefficient of $z^S$ at $0$ is denoted by $u_S$, and these coefficients are its multi-affine jets.

A step cited with no proof named beside it is a step of the proof of \Cref{Y-H7-Df}.

The symbols $r(x)$ and $s(x)$ denote two of the radii at $x$ of item~2 of \S\ref{Y-J-1}, and $\partial_i=\partial/\partial z_i$.

$s_x[D]$ is the map of entry~5 of item~(0) of \S\ref{Y-H0-1}.

The letter $D$ stands for $\tilde{\mathsf E}_{m+k_0}-\mathsf E_{k_0}$ in \textup{(D1)} and in Step~2, and for any family of kernels in $s_x[D]$ and $n[u;D]$; in Step~3 these families are $\mathbb D_k/\mathsf b$ and $(\tilde{\mathbb D}_{m+k_0}-\mathbb D_{k_0})/\mathsf b$. Further, $D(V)$ is the set of the rough plaquettes, and $D$ placed before $A$ or $F$ denotes a derivative.

Without an argument, $c_h$ is the number of \textup{(O11)}, item~1 of \S\ref{Y-J-1}. The caps are $c_h(x)=e^{-c_hp_0(x)^2}$ and $c_a(x)=E_ap_0(x)^3x^{-1/2}$.

$m_l$ and $m'$ are the exponents in $N=2L^{m_l}$ and $M=L^{m'}$.

The symbol $u_k(z)$ has two senses, and in the first it is the given or the honest family, iterated by $u_{k+1}(z)=F_{\ell_k}(u_k(z))$ from its starting level. In \textup{(D0)} and in the part on \textup{(D0)} of Step~2, that level is $k_0$, and the family $u_{k_0}(z)$ is given on $\mathrm{Pd}$. In the part on \textup{(D1)}\,(c) of Step~2 and in Step~8 up to its part ``Representatives'', $u_k(z)$ is, for $0\le k\le K$, the honest family of \textup{(o1)}, item~2 of \S\ref{Y-H6-1}, and on the interior of $\mathrm{Pd}_0$ of the run $K$, that family is an iterate of this kind with starting level $0$. For each run, $\mathrm{Pd}_0$ and $\mathrm{Pd}_0^o$ denote the sets that the part ``Sources'' of item~3 of \S\ref{Y-H4-1} defines for that run. With its argument, $\tilde u_j(z)$ always denotes the given or the honest family of the run $K'$. Everywhere else $u_k(z)$ has the second sense: like the symbol without an argument, it is the representative of \Cref{Y-H6-J}\,(J6), namely $u_k^{\rm pol}$ on $\mathrm{Pd}(\mathfrak r_k)$ for $k\le k_d$ and, for $k_d\le k\le K$, the family on $\mathrm{Pd}(\mathfrak r_*)$ iterated from $u_{k_d}^{\rm pol}$. The families in the two senses share the base point $\bar u_k$ and the multi-affine jets $\delta_S(k)$, by \textup{(o2)} of item~2 of \S\ref{Y-H6-1} and \Cref{Y-H6-J}\,(J6); beyond these they may differ. Without an argument, as in $\mathsf V'_{\ell_k}[u_k]$, $\mathsf V'_{\tilde\ell_j}[\tilde u_j]$ and $\mathsf V'_{\ell_k}[\mathfrak n_k]$ in \textup{(D4)} to \textup{(D6)}, the symbol means the representative of \textup{(D1)}: the polynomial up to $k_d$ (up to $k_d'$ for the run $K'$), and the family after it.



\end{definition}

\begin{definition}[Item 1.3: the table]

Of item~(1) of \S\ref{Y-H0-1} we use the following entries, in the order \textup{(O0)} to \textup{(O14)} of the table. The entry \textup{(O0a)} gives the number $K_5$. From \textup{(O4a)} we take the conditions \textup{(L3su)}, of which this section uses $L^{1/4}\ge10^4$, and \textup{(L5*)}, which enters through $\vartheta_5\le L^{-1/5}$. The entry \textup{(O10a)} gives the number $C_g=263L^3$, and $c_Q^{\rm su}=306M^3\mathsf g(\lambda_\flat)^2$ is the number of \textup{(O10)}, item~1 of \S\ref{Y-I-1}. From \textup{(O12a)} we take the number $C_3$.

After \textup{(O12a)} and before \textup{(O13)} appear three functions of $x$ alone: $\eta_h(x):=e^{-(c_h/2)(p_0(x)^2-p_0(x/L)^2)}$, $\eta_w(x):=e^{-c_\Xi p_0(x)^2}$ and $\check\omega_3(x):=C_3E_ap_0(x)x^{-1/4}/c_Q^{\rm su}$. They contain $x$, the function $p_0$ and the letters $c_h$, $c_\Xi$, $C_3$, $E_a$, $c_Q^{\rm su}$, $L$, all of which are defined before \textup{(O13)}, and none of $\mathsf b$, $n$, $d$.

By \textup{(O13)}, $\mathsf b\ge x_s$, so that $x_k\ge L\mathsf b\ge Lx_s\ge\max\{x_9^U,x_3\}$ for $k<K$. Also by \textup{(O13)}, the following two conditions hold for all $x\ge Lx_s$:

\begin{enumerate}

\item[\textup{(W2)}] $\eta_h(x)+e^{1/2}\eta_w(x)\le1/L$;

\item[\textup{(W3)}] $\check\omega_3(x),\ \check\omega_3(x/L)\le10^{-3}$.

\end{enumerate}

Under \textup{(O13)} there is moreover, from \S\ref{Y-H-1}, the condition \textup{(Y.0)(0b)}: $p_0(x)^4\le2p_0(x/L)^4$ for $x\ge x_9^U$.

From \textup{(O14)} we take first the following functions of $x$ and numbers. We set $\mathrm{amp}(x)=c_Q^{\rm su}K_5x^{1/2}/(\mathsf bE_ap_0(x))$ and $u_3(x)=C_3x^{1/4}$; they obey the identity $u_3(x)K_5/(\mathsf b\,\mathrm{amp}(x))=\check\omega_3(x)$, in which $\mathsf b$ and $K_5$ cancel. Further, $\mathfrak r(x)=3.5/(n\,\mathrm{amp}(Lx))$ and $\mathfrak r_k=\mathfrak r(x_k)$; $\vartheta_a=0.2431L^{-1/4}$ and $\nu_a:=0.2431$; $\vartheta[B]=3.968/(16-B)$, so that $\vartheta[1]\le0.26454$, $\vartheta[2]\le0.28343$ and $\vartheta[2.455]\le0.29295$; finally $\check r(x)=r(x)+3$ and $C_5^{\#}=L^{15/32}C_5+K_5$.

The further entries of \textup{(O14)} are functions of the depth; they depend on $\mathsf b$, $n$, $d$, $R_f$ and on letters, but on none of $K$, $K'$, $N$ and on no row. For $i\ge0$ let $\mathrm{sep}^{[i]}=\lfloor dL^i/M\rfloor-2\,\check r(L^i\mathsf b)-2$, and let $\mathrm{sep}_*=8\log(4nL^{1/2})$. We let $j^*$ be the least $j\ge1$ such that $\mathrm{sep}^{[i]}\ge\mathrm{sep}_*$ for all $i\ge j-1$. For every $j\ge1$ we set $N_1^{[j]}=(2R_fL^{j-1}/M+2\,\check r(L^{j-1}\mathsf b)+3)^3$; moreover $m^{[j]}=\lceil n\,\mathrm{sep}^{[j-1]}/2\rceil$ for $j>j^*$ and $m^{[j]}=1$ for $1\le j\le j^*$.

At the level $k$ of a run of length $K$ we put $\mathrm{sep}_k=\mathrm{sep}^{[K-k]}$ and $k_d=\max\{K-j^*,0\}$; for $k<k_d$ we put $\check\kappa_k=2\log(4nL^{1/2})/\mathrm{sep}_{k+1}$, as in \textup{(o4)}, item~2 of \S\ref{Y-H6-1}. The rates are $\kappa_k=\kappa+\sum_{k\le j<k_d}\check\kappa_j\in[\kappa,\kappa+0.26]$, and $\kappa_k=\kappa$ for $k\ge k_d$. Finally $\mathfrak r_*=\mathfrak r_{k_d}$, $N_1(k)=N_1^{[K-k]}$ and $m_k=m^{[K-k]}$.

The two series of the table are \[\mathrm{Tail}(J)=\sum_{j>J}\mathfrak r(L^{\max(j,j^*)}\mathsf b)^{-n}N_1^{[j]}\,2C_gc_a(L^j\mathsf b)\,e^{-\kappa(m^{[j]}-1)}\] and \[\mathrm{Ser}_2=\sum_{j\ge1}\vartheta^{-j}\bigl(2/\mathfrak r(L^{\max(j,j^*)}\mathsf b)\bigr)^nN_1^{[j]}\,1.2C_gc_a(L^j\mathsf b)\,e^{-\kappa(m^{[j]}-1)}.\]

The threshold $K_T(\mathsf b,n,d)$ is the least integer $K\ge K_0(\mathsf b)$ that satisfies the two conditions

\begin{enumerate}

\item[\textup{(K1a)}] $(4/9)K-(10/9)\Lambda-1\ge j^*+1$,

\item[\textup{(K1b)}] $C_5^{\#}K_5^{-1}K\,0.3683^{(4/9)K-(10/9)\Lambda-1}\le0.1$.

\end{enumerate}

Both conditions then hold for every larger $K$, by Step~1\,(c). We put $\Omega^{\#}:=C_5^{\#}K_5^{-1}K\,0.3683^{K-k_0}$.

The conditions \textup{(W2)} and \textup{(W3)} use $\eta_h$, $\eta_w$ and $\check\omega_3$ in the closed forms above, which are defined before \textup{(O13)}. No floor depends on an entry of \textup{(O14)}, and no entry depends on a later one.



\end{definition}

\begin{definition}[Item 3: the data]

The data are those of \Cref{Y-H9-T}\,(T2). The letters are those of item~(1) of \S\ref{Y-H0-1}, and $\mathsf b\ge x_s$. We take $N=2L^{m_l}$ with $m_l\ge m'$, so that $N$ is even and $M\mid N$, and we require \textup{(Tor)} $N\ge2M\,s(L\mathsf b)$. Further, let $n\ge2$ and $d,R_f>0$. The functions $f_1,\dots,f_n$ are real and Lipschitz on the torus of side $N$, with $\|f_i\|_\sharp=\max\{\|f_i\|_\infty,40M\,\mathrm{Lip}f_i\}\le1$, where $\mathrm{Lip}$ refers to the Euclidean distance. The sets $\{f_i\ne0\}$ lie pairwise at sup distance at least $d$, and each of them lies within sup distance $R_f$ of one point, as in item~1 of \S\ref{Y-H6-1}; these are the distances of item~1 of \S\ref{Y-H9-1}. Finally $K'\ge K\ge K_T(\mathsf b,n,d)$.

\end{definition}

\begin{theorem}[Item 4: two runs of jets meet on growing radii]\label{Y-H7-Df}

In the setting of item~3, the following hold.

\begin{enumerate}

\item[\textup{(D0)}] \emph{A fixed polydisc.} For this clause it suffices that $K\ge K_0(\mathsf b)$. Let $\mathrm{Pd}$ be a polydisc about $0$ on which the two sourced runs have families $u_{k_0}(z)$ and $\tilde u_{m+k_0}(z)$ in $M_1(x_{k_0};\kappa_1,c_h)$ with the strong cap, such that $\rho^{[K],z}_{k_0}=e^{\mathsf V_0(z)}\rho[u_{k_0}(z);\ell_{k_0}]$ and $\rho^{[K'],z}_{m+k_0}=e^{\tilde{\mathsf V}_0(z)}\rho[\tilde u_{m+k_0}(z);\tilde\ell_{m+k_0}]$, where $\mathsf V_0(z)$ and $\tilde{\mathsf V}_0(z)$ are numbers that depend holomorphically on $z$; by \Cref{Y-H4-U4}\,(h), this is so on the interior of $\mathrm{Pd}_0$ of the run $K'$. Here $\rho^{[K],z}_k$ is the $\rho_k^z$ of \Cref{Y-H4-U4}\,(h) taken for the run $K$, that is, $\rho^{[K],z}_0=\rho^{[K]}_0\exp\{\sum_iz_i\mathcal O_K(f_i)\}$ and $\rho^{[K],z}_{k+1}=(\rho^{[K],z}_k)'$, and $\rho^{[K'],z}_j$ is defined in the same way for the run $K'$. Put $\mathfrak n^{\circ}_{k_0}(z):=\mathbb B_{\tilde{\mathsf E}_{m+k_0}-\mathsf E_{k_0}}(\tilde u_{m+k_0}(z))$, and recursively $\mathfrak n^{\circ}_{k+1}(z):=F_{\ell_k}(\mathfrak n^{\circ}_k(z))$ and $u_{k+1}(z):=F_{\ell_k}(u_k(z))$. Then, for $k_0<k\le K$ and every $z$ of $\mathrm{Pd}$, three assertions hold. First, both $u_k(z)$ and $\mathfrak n^{\circ}_k(z)$ are families on $\mathrm{Pd}$ in $M_1(x_k;\kappa_1,c_h)$ with the strong cap. Second, at every field, \[\rho^{[K'],z}_{m+k}=\exp\{\tilde{\mathsf V}_0(z)+\sum_{k_0\le j<k}\mathsf V'_{\ell_j}[\mathfrak n^{\circ}_j(z)]\}\,\rho[\mathfrak n^{\circ}_k(z);\ell_k],\] where the numbers in the exponent are holomorphic. Third, $\sup_{\mathrm{Pd}}\mathsf D_{x_k}(u_k(z),\mathfrak n^{\circ}_k(z))\le3\vartheta^{k-k_0}$. If moreover, at the real $z$ of $\mathrm{Pd}$, the members $u_{k_0}(z)$ and $\tilde u_{m+k_0}(z)$ are real and $\rho[u_{k_0}(z);\ell_{k_0}]$, $\rho[\tilde u_{m+k_0}(z);\tilde\ell_{m+k_0}]>0$ (the honest families have this property by \Cref{Y-H4-U4}\,(h), since the $f_i$ are real by item~3), then at these $z$ all $u_k(z)$ and $\mathfrak n^{\circ}_k(z)$ and all numbers $\mathsf V'_{\ell_j}[u_j(z)]$, $\mathsf V'_{\ell_j}[\mathfrak n^{\circ}_j(z)]$ are real, and $\rho[u_k(z);\ell_k]$, $\rho[\mathfrak n^{\circ}_k(z);\ell_k]>0$.

\item[\textup{(D1)}] \emph{Objects.} The objects $\bar u_k=u_k(0)$, $\delta_S(k)$, $u_k^{\rm pol}$ for $0\le k\le k_d$, and $u_k(z)$ for $k_d\le k\le K$ on $\mathrm{Pd}(\mathfrak r_*)$, are those of \Cref{Y-H6-J} for the run $K$; the objects $\bar{\tilde u}_j$, $\tilde\delta^{\rm own}_S(j)$, $\tilde u_j^{\rm pol}$ are the corresponding ones for the run $K'$, defined over its own rows. Let $Q:=Q_D$ be as in \Cref{Y-H5-Bc}, taken at $D=\tilde{\mathsf E}_{m+k_0}-\mathsf E_{k_0}$ and $x=x_{k_0}$. On $\mathrm{Pd}(\mathfrak r_{k_0}/2)$ put $\mathfrak n_{k_0}(z):=\mathbb B_D(\tilde u^{\rm pol}_{m+k_0}(z))$. For $k_0\le k\le k_d$ put $\mathfrak n_k^{\rm pol}(z):=\bar{\mathfrak n}_k+\sum_Sz^S\tilde\delta_S(k)$, with $\mathfrak n_{k_0}^{\rm pol}=\mathfrak n_{k_0}$. For $k_0\le k<k_d$ let $\bar{\mathfrak n}_{k+1}:=F_{\ell_k}(\bar{\mathfrak n}_k)$, and let $\tilde\delta_S(k+1)$ be the coefficient of $z^S$ of $F_{\ell_k}(\mathfrak n_k^{\rm pol}(z))$. For $k_d\le k<K$ put $\mathfrak n_{k_d}(z):=\mathfrak n_{k_d}^{\rm pol}(z)$ and $\mathfrak n_{k+1}(z):=F_{\ell_k}(\mathfrak n_k(z))$ on $\mathrm{Pd}(\mathfrak r_*/2)$, and $\bar{\mathfrak n}_{k+1}:=\mathfrak n_{k+1}(0)=F_{\ell_k}(\bar{\mathfrak n}_k)$. Then:

\begin{enumerate}

\item[\textup{(a)}] $k_0+2\le k_d=K-j^*$, and the run $K'$ has $k_d'=m+k_d$; partners share $x_k$, $\mathfrak r_k$, $\mathrm{sep}_k$, $\check\kappa_k$, $\kappa_k$, $\check r$, the cubes and the $L$-blocks, while the sets $S_i(k)$, $\tilde S_i(k)$ of item~1.1 are not among the shared objects.

\item[\textup{(b)}] $\mathfrak n_{k_0}(z)$ is a multi-affine family in $M_{2.455}(x_{k_0};\kappa_{k_0},c_h)$, with the weak cap, and $\tilde\delta_S(k_0)=(\tilde\delta^{{\rm own},a}_S,\tilde\delta^{{\rm own},h}_S,\tilde\delta^{{\rm own},w}_Se^Q)(m+k_0)$.

\item[\textup{(c)}] \emph{True jets.} $\bar{\mathfrak n}_k$ is the base point, and the $\tilde\delta_S(k)$ are the multi-affine jets, of the family $\mathfrak n^{\circ}_k(z)$ that \textup{(D0)} provides on the interior of $\mathrm{Pd}_0$ of the run $K'$, that is, of a system of members of the sourced long run, taken over the rows $\ell_k$ of the short run. For $k>k_0$, $\bar{\mathfrak n}_k$ is an $\mathfrak n_k$ as in \Cref{Y-J-ThA}\,(iii); that clause asserts existence, and $\bar{\mathfrak n}_k$ is the system that Steps~4 and~5 of the proof of \Cref{Y-J-ThA} construct from the systems $u_{k_0}=\bar u_{k_0}$ and $\tilde u=\bar{\tilde u}_{m+k_0}$. Of this, only \textup{(D0)} at $z=0$ enters: $\bar{\mathfrak n}_k$ is real, lies in $M_1(x_k;\kappa_1,c_h)$ and carries the strong cap; for some real number $\mathsf V$ one has $\rho^{[K']}_{m+k}=e^{\mathsf V}\rho[\bar{\mathfrak n}_k;\ell_k]>0$ at every field; and $\mathsf D_{x_k}(\bar u_k,\bar{\mathfrak n}_k)\le3\vartheta^{k-k_0}$.

\item[\textup{(d)}] \emph{Reach.} $\tilde\delta_S(k)(X)=0$ unless $X^{+\check r(x_k)}$ meets $\tilde S_i(k)$ for every $i\in S$; so, where $\tilde\delta_S(k)(X)\ne0$ and $|S|\ge2$, we have $\#X\ge|S|\,\mathrm{sep}_k/2$ for the components $a$ and $h$, and $|X|_c\ge|S|\,\mathrm{sep}_k/2$ for the component $w$.

\item[\textup{(e)}] For $k_0<k\le k_d$, $u_k^{\rm pol}$ and $\mathfrak n_k^{\rm pol}$ are families on $\mathrm{Pd}(\mathfrak r_k/2)$ in $M_2(x_k;\kappa_k,c_h)$, with $\|u_k^{\rm pol}(z)-\bar u_k\|^{{\rm st},(\kappa_k)}\le0.455$ and $\|\mathfrak n_k^{\rm pol}(z)-\bar{\mathfrak n}_k\|^{{\rm st},(\kappa_k)}\le0.455+0.34$.

\item[\textup{(f)}] For $k_d<k\le K$, $u_k(z)$ and $\mathfrak n_k(z)$ are families on $\mathrm{Pd}(\mathfrak r_*/2)$ in $M_1(x_k;\kappa,c_h)$ with the strong cap; they are outputs of $F$ and carry the true multi-affine jets, and $\|u_k(z)-\bar u_k\|^{{\rm st},(\kappa)}\le0.455$ and $\|\mathfrak n_k(z)-\bar{\mathfrak n}_k\|^{{\rm st},(\kappa)}\le0.795$, each right side being multiplied by $\vartheta[2]^{k-k_d}$.

\end{enumerate}

\item[\textup{(D2)}] \emph{Pair bound.} Put $e_k:=\mathsf D_{x_k}(\bar u_k,\bar{\mathfrak n}_k)$ for $k>k_0$ and $e_{k_0}:=3$. Put $a^D_k:=\max_i|\delta_i^a(k)-\tilde\delta_i^a(k)|_a$ and, for $k>k_0$, $h^D_k:=\max_i\max\{|\delta_i^h-\tilde\delta_i^h|_h,|\delta_i^w-\tilde\delta_i^w|_w^{\rm st}\}(k)$. For $k_0<k\le k_d$ let \[\Delta_k:=e_k+(n\mathfrak r_k/2)\max\{a^D_k,h^D_k\}+\sum_{|S|\ge2}(\mathfrak r_k/2)^{|S|}\|\delta_S(k)-\tilde\delta_S(k)\|^{{\rm st},(\kappa_k)},\] and let $\Delta_{k_0}:=3.91$. Then $\Delta_k\ge\sup_{\mathrm{Pd}(\mathfrak r_k/2)}\mathsf D^{(\kappa_k)}_{x_k}(u_k^{\rm pol},\mathfrak n_k^{\rm pol})$, where at $k_0$ the distance is $\mathsf D^{\rm in}$. Moreover, for $k_0<k\le k_d$, with \[\Omega_0:=\max_i\|\tilde{\mathbb D}_{m+k_0}[f_i]-\mathbb D_{k_0}[f_i]\|_{(\lambda_\flat)}/(K_5\check\omega_3(x_{k_0})),\] the bounds \textup{(i)} to \textup{(iii)} hold:

\begin{enumerate}

\item[\textup{(i)}] $a^D_k\le(u_3(x_k)K_5/\mathsf b)\,[7.6\vartheta^{k-k_0}+\nu_a^{k-k_0}\Omega_0]$;

\item[\textup{(ii)}] $n\mathfrak r_ka^D_k\le0.0054\vartheta^{k-k_0}+0.7\Omega^{\#}$ and $n\mathfrak r_kh^D_k\le0.49\vartheta^{k-k_0}+0.25\Omega^{\#}$;

\item[\textup{(iii)}] $\Delta_k\le3.36\vartheta^{k-k_0}+0.36\Omega^{\#}\le3.45\vartheta^{k-k_0}+1.36\Omega^{\#}$, where $\Omega^{\#}\le0.1$, and $\Delta_k-e_k\le0.195\le0.34$.

\end{enumerate}

\item[\textup{(D3)}] \emph{End.} \[\sup_{\mathrm{Pd}(\mathfrak r_*/2)}\mathsf D_{\mathsf b}^{(\kappa)}(u_K(z),\mathfrak n_K(z))\le(1.15+0.46C_5^{\#}K_5^{-1}K)\,C_{\mathsf b}\theta^K.\] At a depth $s\ge0$ with $K-s>k_0$, that is, with $K-s\ge K_0(L^s\mathsf b)$, we have $\sup \mathsf D^{(\kappa_{K-s})}_{L^s\mathsf b}\le(1.15+0.46C_5^{\#}K_5^{-1}K)\,C_{L^s\mathsf b}\theta^{K-s}$; there the supremum is that of the distance between the two families over $\mathrm{Pd}(\mathfrak r_*/2)$ if $s<j^*$, and that of the distance between the two polynomials over $\mathrm{Pd}(\mathfrak r(L^s\mathsf b)/2)$ if $s\ge j^*$.

\item[\textup{(D4)}] \emph{Vacuum, pairs.} For $k_d<k\le K$ put $\Delta_k:=\sup_{\mathrm{Pd}(\mathfrak r_*/2)}\mathsf D^{(\kappa)}_{x_k}(u_k(z),\mathfrak n_k(z))$. For $k_0\le k<K$, \[|\partial_1\cdots\partial_n(\mathsf V'_{\ell_k}[\mathfrak n_k]-\mathsf V'_{\ell_k}[u_k])(0)|\le(2/\mathfrak r_{\min(k,k_d)})^nN_1(k)\,1.2C_gc_a(x_k)\,e^{-\kappa(m_k-1)}\Delta_k,\] with the polynomials for $k\le k_d$ and the families after; the sum of the left sides over $k$ is at most $\mathrm{Ser}_2\,(4+1.36C_5^{\#}K_5^{-1}K)\,\vartheta^{K-k_0}$.

\item[\textup{(D5)}] \emph{Vacuum, early levels.} \[\sum_{0\le k<k_0}|\partial_1\cdots\partial_n\mathsf V'_{\ell_k}[u_k](0)|\le\mathrm{Tail}(K-k_0)\] and \[\sum_{0\le j<m+k_0}|\partial_1\cdots\partial_n\mathsf V'_{\tilde\ell_j}[\tilde u_j](0)|\le\mathrm{Tail}(K-k_0).\] For every $J\ge j^*$, \[\mathrm{Tail}(J)\le2\,\mathfrak r(L^{J+1}\mathsf b)^{-n}N_1^{[J+1]}\,2C_gc_a(L^{J+1}\mathsf b)\,e^{\kappa}(4nL^{1/2})^{-4\kappa nL^{J+1-j^*}}.\]

\item[\textup{(D6)}] \emph{Identity.} Put $S^T_{n;K}:=S^T_{n;K,l,L^K\mathsf b}(f_1,\dots,f_n)$ and $S^T_{n;K'}:=S^T_{n;K',l,L^{K'}\mathsf b}(f_1,\dots,f_n)$, where the right sides are those of item~1 of \S\ref{Y-H6-1}, taken with $l=N$ and with the $f_i$ of item~3. Then \[\FitDisplay{\dimexpr\linewidth-1pt\relax}{\begin{aligned}
&S^T_{n;K'}-S^T_{n;K}=\sum_{0\le j<m+k_0}\partial_1\cdots\partial_n\mathsf V'_{\tilde\ell_j}[\tilde u_j](0)
\\
&\quad{}-\sum_{0\le k<k_0}\partial_1\cdots\partial_n\mathsf V'_{\ell_k}[u_k](0)
\\
&\quad{}+\sum_{k_0\le k<K}\partial_1\cdots\partial_n(\mathsf V'_{\ell_k}[\mathfrak n_k]-\mathsf V'_{\ell_k}[u_k])(0)
\\
&\quad{}+\partial_1\cdots\partial_n(\log I[\mathfrak n_K(z)]-\log I[u_K(z)])(0).
\end{aligned}}\] Here the step $0\to1$ of each run is the term $0$ of its early sum.

\item[\textup{(D7)}] \emph{What the constants depend on.} The numbers $j^*$, $K_T$, $\mathfrak r_*=\mathfrak r(L^{j^*}\mathsf b)$, $\mathrm{Ser}_2$, $\mathrm{Tail}(J)$, $C_{\mathsf b}$ and $\Lambda$ depend only on $L$, $M$, $E_a$, $\kappa$, $c_R$, on the letters $A_0$, $\mathfrak p_0$ of \textup{(O8)} (through $p_0(x)=A_0(\log x)^{\mathfrak p_0}$, which enters $\mathrm{amp}$ and $c_a$), on pure numbers, and on $\mathsf b$, $n$, $d$, $R_f$. In detail, $j^*$ depends on $L$, $M$, $c_R$ and on $\mathsf b$, $n$, $d$; $\Lambda$ and $C_{\mathsf b}$ depend on $L$, $M$, $E_a$, pure numbers and $\mathsf b$; $K_T$ depends on all of these; $\mathfrak r_*$ depends in addition on $A_0$, $\mathfrak p_0$, through $\mathrm{amp}$; and $\mathrm{Ser}_2$, $\mathrm{Tail}(J)$ depend in addition on $\kappa$, $R_f$, and once more on $p_0$, through $c_a$. None of them depends on $N$, $K'$, $m$, a row or an $f_i$. Each of them, as a function of $\mathsf b$, is bounded on every compact subset of the range $[x_s,\infty)$ of $\mathsf b$, and $\mathfrak r_*$ is bounded there from below. All numerals of \textup{(D1)} to \textup{(D4)} are free of every letter.

\end{enumerate}

\end{theorem}

\begin{proof}[Proof of \Cref{Y-H7-Df}]

\emph{Step 1: the arithmetic of the table.} We first collect the following facts on the entries of item~1.3.

\begin{enumerate}

\item[\textup{(a)}] Five relations hold:

\begin{enumerate}

\item[\textup{(t1)}] $\vartheta_au_3(x)/u_3(x/L)=\nu_a=0.2431$;

\item[\textup{(t2)}] $u_3(x)K_5/\mathsf b=\check\omega_3(x)\,\mathrm{amp}(x)$;

\item[\textup{(t3)}] $5\le2^{-1/4}L^{1/2}\le\mathrm{amp}(Lx)/\mathrm{amp}(x)\le L^{1/2}$ for $Lx\ge x_9^U$, and $\mathfrak r_{k+1}/\mathfrak r_k=\mathrm{amp}(Lx_k)/\mathrm{amp}(x_k)$, $n\mathfrak r_{k+1}=3.5/\mathrm{amp}(x_k)$;

\item[\textup{(t4)}] $\check\omega_3(x_k)/\check\omega_3(x_{k_0})\le L^{(k-k_0)/4}$ for $k\ge k_0$;

\item[\textup{(t5)}] $c_a(Lx)\le c_a(x)$ for $Lx\ge x_9^U$.

\end{enumerate}

\item[\textup{(b)}] $1\le k_0\le K-1$; $K-k_0>\mathrm{Ex}_K:=(4/9)K-(10/9)\Lambda-1$; $k_0+2\le k_d=K-j^*$; $k_d'=m+k_d$; and partners share every number of the table.

\item[\textup{(c)}] \textup{(K1a)} and \textup{(K1b)} hold at every $K\ge K_T$, and $\Omega^{\#}\le0.1$.

\item[\textup{(d)}] $(K-k_0)/4-k_0/5\le-\Lambda/2<0$.

\item[\textup{(e)}] $\mathrm{sep}^{[i+1]}\ge L\,\mathrm{sep}^{[i]}$; hence $j^*$ is the least $j\ge1$ with $\mathrm{sep}^{[j-1]}\ge\mathrm{sep}_*$, it is finite, and $\mathrm{sep}^{[i]}\ge L^{i+1-j^*}\mathrm{sep}_*$ for $i\ge j^*-1$; moreover $\check\kappa_k\le1/4$ for $k<k_d$.

\item[\textup{(f)}] $N_1^{[j+1]}\le L^3N_1^{[j]}$.

\end{enumerate}

We prove these facts in turn.

\begin{enumerate}

\item[\textup{(a)}] We take the five relations one by one.

\begin{enumerate}

\item[\textup{(t1)}] Since $u_3(x)=C_3x^{1/4}$, the quotient $u_3(x)/u_3(x/L)$ equals $L^{1/4}$, and this factor is cancelled by $\vartheta_a=0.2431L^{-1/4}$.

\item[\textup{(t2)}] We insert the formulas for $\mathrm{amp}$ and for $u_3$; then $\mathsf b$ and $K_5$ cancel.

\item[\textup{(t3)}] We have $\mathrm{amp}(Lx)/\mathrm{amp}(x)=L^{1/2}p_0(x)/p_0(Lx)$. Since $p_0$ is increasing, $p_0(x)\le p_0(Lx)$; hence $p_0(x)/p_0(Lx)\le1$, and the upper bound $L^{1/2}$ follows. For the lower bound we apply \textup{(Y.0)(0b)} at the point $Lx$, which is admissible as $Lx\ge x_9^U$; it gives $p_0(x)/p_0(Lx)\ge2^{-1/4}$. Since $L^{1/2}\ge10^8$, also $2^{-1/4}L^{1/2}\ge5$. The two formulas for $\mathfrak r$ follow from $\mathfrak r(x)=3.5/(n\,\mathrm{amp}(Lx))$ and $Lx_{k+1}=x_k$.

\item[\textup{(t4)}] Here $\check\omega_3(x)=C_3E_ap_0(x)x^{-1/4}/c_Q^{\rm su}$ and $x_k=L^{-(k-k_0)}x_{k_0}$; so the quotient $\check\omega_3(x_k)/\check\omega_3(x_{k_0})$ is $L^{(k-k_0)/4}$ times $p_0(x_k)/p_0(x_{k_0})$, and the latter factor is at most $1$ because $p_0(x_k)\le p_0(x_{k_0})$.

\item[\textup{(t5)}] By the form of $c_a$, $c_a(Lx)/c_a(x)=(p_0(Lx)/p_0(x))^3L^{-1/2}\le2^{3/4}10^{-8}$, with the two bounds used for \textup{(t3)}; this number is below $1$.

\end{enumerate}

\item[\textup{(b)}] Since $K\ge K_0(\mathsf b)$, \Cref{Y-J-ThA}\,(i) gives $1\le k_0\le K-1$. From $k_0<(5K+10\Lambda)/9+1$ we obtain $K-k_0>\mathrm{Ex}_K$. Together with \textup{(K1a)} this yields $K-k_0>j^*+1$, and as both sides are integers, $k_0\le K-j^*-2$. In particular $K>j^*$, so that $k_d=K-j^*$; and $K'\ge K$ gives $k_d'=K'-j^*=m+k_d$. Every entry of item~1.3 that is taken at the level $k$ depends on $K$ and $k$ only through the depth $K-k$, and the level $m+k$ of $K'$ has the same depth $K-k$. The tori, grids, cubes and $L$-blocks, the $x_k$, the caps and the domains are the same for partners, by Step~2\,(a) of the proof of \Cref{Y-J-ThA}, which involves neither a source nor a support. So partners share every number of the table.

\item[\textup{(c)}] The condition \textup{(K1a)} is monotone in $K$. The left side of \textup{(K1b)} is a constant times $Kc^K$, where $c=0.3683^{4/9}\le0.64151$; passing from $K$ to $K+1$ multiplies it by $(1+1/K)c\le(4/3)0.64151\le0.8554<1$ when $K\ge3$, and every integer $K\ge K_0(\mathsf b)>9/4$ satisfies $K\ge3$. Hence each of the two conditions, once it holds, keeps holding for all larger $K$. Since $0.3683<1$ and $K-k_0>\mathrm{Ex}_K$, we get $\Omega^{\#}\le C_5^{\#}K_5^{-1}K\,0.3683^{\mathrm{Ex}_K}\le0.1$, the last inequality being \textup{(K1b)}.

\item[\textup{(d)}] Since $k_0\ge(5K+10\Lambda)/9$, \[(K-k_0)/4-k_0/5=K/4-(9/20)k_0\le K/4-(9/20)(5K+10\Lambda)/9=-\Lambda/2,\] and $\Lambda>0$ by \Cref{Y-J-ThA}\,(i).

\item[\textup{(e)}] By \textup{(G1)} of item~3 of \S\ref{Y-H6-1}, established in Step~1 of the proof of \Cref{Y-H6-J}, we have $\mathrm{sep}^{[i+1]}\ge L\,\mathrm{sep}^{[i]}$ at each $i\ge0$, and $j^*$ is a finite number. Consequently the inequality $\mathrm{sep}^{[i]}\ge\mathrm{sep}_*>0$, once it holds at $i$, holds at $i+1$ as well, with a gain of the factor $L$; this gives both the description of $j^*$ and the lower bound $\mathrm{sep}^{[i]}\ge L^{i+1-j^*}\mathrm{sep}_*$ for $i\ge j^*-1$. By \textup{(G2)} of item~3 of \S\ref{Y-H6-1}, $\check\kappa_k\le1/(4L)\le1/4$ for $k<k_d$.

\item[\textup{(f)}] By Step~1\,(r) of the proof of \Cref{Y-H6-J}, $\check r(Lx)\le\check r(x)+\log L/960+1\le L\,\check r(x)$, the second inequality because $\check r\ge3$. Therefore $2R_fL^j/M+2\,\check r(L^j\mathsf b)+3\le L\,(2R_fL^{j-1}/M+2\,\check r(L^{j-1}\mathsf b)+3)$, and since both sides are positive, raising them to the third power gives \textup{(f)}.

\end{enumerate}

\emph{Step 2: the switch for families, the true jets, the reach and the level $k_0$.} In this step we prove \textup{(D0)} and the parts \textup{(a)}, \textup{(b)}, \textup{(c)}, \textup{(d)} of \textup{(D1)}, together with \textup{(Sep)}: the sets $(S_i(k)\cup\tilde S_i(k))^{+\check r(x_k)}$, $i=1,\dots,n$, lie pairwise at distance at least $\mathrm{sep}_k+2$, counted in cubes, and $(S_1(k+1)\cup\tilde S_1(k+1))^{+\check r(x_{k+1})}$ has at most $N_1(k)$ cells. We also show that $\Delta_{k_0}=3.91$ bounds $\sup_{\mathrm{Pd}(\mathfrak r_{k_0}/2)}\mathsf D^{{\rm in},(\kappa_{k_0})}_{x_{k_0}}(u_{k_0}^{\rm pol},\mathfrak n_{k_0})$.

\emph{Notation.} We put $x=x_{k_0}$, $\ell=\ell_{k_0}$, $\tilde\ell=\tilde\ell_{m+k_0}$ and $D=\tilde{\mathsf E}_{m+k_0}-\mathsf E_{k_0}$.

\emph{The conditions of \Cref{Y-H5-Bc} at $(x,D)$.} Both the requirements of \Cref{Y-H5-Bc} on the lattice and the five conditions of item~5 of \S\ref{Y-I-1} are satisfied, as item~2 of Step~4 of the proof of \Cref{Y-J-ThA} shows. The family $D$ is real and symmetric, and its range is at most $\varrho(x)$, because both families of rows are cut at $\varrho(x)$. Finally $\delta_\Delta\le1$, by Step~3 of the proof of \Cref{Y-J-ThA} taken at $k=k_0$.

\emph{The rate.} The map $\mathbb B_D(\tilde u)=(\tilde E^a+E^\Delta[D],\tilde E^h,\tilde we^{Q})$ changes no index set: the term of $\mathbb B_D(\tilde u)$ at $X$ is obtained from the term of $\tilde u$ at the same $X$; $E^\Delta$ lives on $\#X=1$, and $Q$ depends on $V$ and $D$ alone. In \textup{(B3)} and \textup{(B5)} the weight $e^{\tilde\kappa(|X|_c-1)}$ is a common factor of the two sides. Hence \Cref{Y-H5-Bc} holds, with the same numerals, at every rate $\tilde\kappa$ of $[\kappa,\kappa_1]$; indeed its data admit any rate $\tilde\kappa\ge0$.

\emph{Proof of \textup{(D0)}.} At each $z$ of $\mathrm{Pd}$ we apply \Cref{Y-H5-Bc} at $A=1$: the member $\mathfrak n^{\circ}_{k_0}(z)$ lies in $M_2(x;\kappa_1,c_h)$, by \textup{(B1)} we have $\rho[\tilde u_{m+k_0}(z);\tilde\ell]=\rho[\mathfrak n^{\circ}_{k_0}(z);\ell]$ at every field, and by \textup{(B5)} $\mathfrak n^{\circ}_{k_0}(z)$ is a family. For $k_0\le k\le K-1$, \textup{(SD)} holds at $(x_k,NL^{K-k-1},\ell_k)$, by Step~2\,(d) of the proof of \Cref{Y-J-ThA}. By Step~2\,(a) of the same proof, the step acts linearly on $\rho$, involves no row, and is the same map for two partners. Applied first at $A=2$ and then at $A=1$, clause \textup{(a)} of \Cref{Y-H-C9}, which is available through \Cref{Y-H1-U1}\,(U) at $p=(\kappa_1,c_h)$, gives \[\rho^{[K'],z}_{m+k+1}=(\rho^{[K'],z}_{m+k})'=e^{\,\cdots+\mathsf V'_{\ell_k}[\mathfrak n^{\circ}_k(z)]}\rho[\mathfrak n^{\circ}_{k+1}(z);\ell_{k+1}],\] since $(\ell_k)'^{\rm c}=\ell_{k+1}$. By \Cref{Y-H5-Af}\,(i) the images are families in $M_1$ with the strong cap, and the numbers are holomorphic. As for the distance, $\|u_{k_0}(z)\|\le1$ and $\|\mathfrak n^{\circ}_{k_0}(z)\|\le2$, so that the two inputs are at distance at most $3$; clause \textup{(e)} of \Cref{Y-H-C9}, used through \Cref{Y-H1-U1}\,(U) at $p=(\kappa_1,c_h)$ and at each $z$, then gives $0.4\cdot3=1.2\le3\vartheta$ at the first level and a further factor $0.4\le\vartheta$ at each later level.

\emph{Real $z$.} Now let the further hypothesis of \textup{(D0)} hold, and let $z$ be real. Then $D=\tilde{\mathsf E}-\mathsf E$ is real, since the rows lie in real Hilbert spaces; by \textup{(B8)} the member $\mathfrak n^{\circ}_{k_0}(z)$ is real, and by \textup{(B1)} $\rho[\mathfrak n^{\circ}_{k_0}(z);\ell]=\rho[\tilde u_{m+k_0}(z);\tilde\ell]>0$. From there on, the reality of the members and of the numbers $\mathsf V'$ passes from each level to the next by clause \textup{(d)} of \Cref{Y-H-C9}, which is available through \Cref{Y-H1-U1}\,(U); the same clause, together with \textup{(a)}, passes on the positivity of the density; this holds for both runs. For the honest families the further hypothesis follows from \Cref{Y-H4-U4}\,(a3), (h), the $f_i$ being real.

These arguments are those of Steps~4 and~5 of the proof of \Cref{Y-J-ThA}, carried out at each $z$, and none of the estimates in those steps depends on the sign of a quantity. The items of \S\ref{Y-J-1} on members are stated for real members, by way of \textup{(R8)} and \textup{(R11)} of its item~3; at a complex $z$ the argument is independent of them. At a complex $z$ we draw on \Cref{Y-H5-Bc}, on \Cref{Y-H1-U1}\,(U), on \Cref{Y-H5-Af}\,(i), and on those items of the proof of \Cref{Y-J-ThA} that involve no member, namely Step~2\,(a) and~(d), Step~3, and item~2 of Step~4.

Clause \textup{(D1)}\,(a) is proved in Step~1\,(b).

For \textup{(D1)}\,(b) we apply \Cref{Y-H6-J}\,(J5) to the run $K'$ at its level $m+k_0\le k_d'$; it gives $n\,\mathfrak r_{k_0}\max_i\|\tilde\delta^{\mathrm{own}}_i\|^{\mathrm{st}}\le0.7$ and $\sum_{|S|\ge2}\mathfrak r_{k_0}^{|S|}\|\tilde\delta^{\mathrm{own}}_S\|^{\mathrm{st},(\kappa_{k_0})}\le0.21$. Let $|z_i|<\mathfrak r_{k_0}/2$ for every $i$. Then the singles contribute at most $0.7/2$ and the sets with $|S|\ge2$ at most $0.21/4$, so that $\|\sum_Sz^S\tilde\delta^{\mathrm{own}}_S\|^{\mathrm{st},(\kappa_{k_0})}\le0.7/2+0.21/4=0.4025\le0.455$. Since moreover $\|\bar{\tilde u}_{m+k_0}\|^{\mathrm{st}}\le1$ by \Cref{Y-H4-U4}\,(h), each of the three members of $\tilde u^{\mathrm{pol}}_{m+k_0}(z)$ has strong size at most $1.455$. \Cref{Y-H5-Bc}, taken with $A=1.455$ and at the rate $\kappa_{k_0}$, now yields $\mathfrak n_{k_0}(z)\in M_{2.455}(x;\kappa_{k_0},c_h)$. The map $\mathbb B_D$ is affine in $\tilde u$, and its shift $E^\Delta[D]$ as well as its factor $e^Q$ are free of $z$; consequently $\mathfrak n_{k_0}(z)$ is multi-affine in $z$, and its coefficients are those named in \textup{(D1)}\,(b). Its base point is $\bar{\mathfrak n}_{k_0}=\mathbb B_D(\bar{\tilde u}_{m+k_0})$, an element of $M_2(x;\kappa_1,c_h)$.

For \textup{(D1)}\,(c) we start from the honest family $\tilde u_{m+k_0}(z)$ of \Cref{Y-H4-U4}\,(h): restricted to the interior of $\mathrm{Pd}_0$ of the run $K'$, it satisfies the hypothesis of \textup{(D0)}. The interior is taken because $\mathrm{Pd}_0$ is closed, while a family is defined on an open set. \Cref{Y-H4-U4}\,(h) provides membership in $M_1$ with the strong cap on $\mathrm{Pd}_0$ at the levels $\ge1$, and holomorphy on the larger open set $\mathrm{Pd}_0^o$; from then on only germs at $0$ are used. The family $u_{k_0}(z)$ satisfies the hypothesis on the same set, because $\mathrm{Pd}_0$ of the run $K'$ is contained in $\mathrm{Pd}_0$ of the run $K$. To see this, note that $\varepsilon(x)=p_0(x)x^{-1/2}=c_Q^{\rm su}K_5/(\mathsf bE_a\,\mathrm{amp}(x))$, so that \textup{(t3)} gives $\varepsilon(Lx)\le\varepsilon(x)/5$ whenever $Lx\ge x_9^U$. Applying this $m$ times, starting from $x=L^K\mathsf b$, we obtain $\varepsilon(L^{K'}\mathsf b)\le5^{-m}\varepsilon(L^K\mathsf b)$; in this iteration only the points $L^j\mathsf b$ occur. The radii of $\mathrm{Pd}_0$ are $1.99c_\nu\mathsf b\,\varepsilon(L^K\mathsf b)/\|f_i\|_\infty$ (item~3 of \S\ref{Y-H4-1}).

By definition, the base point and the multi-affine jets of $\tilde u_{m+k_0}(z)$ are $\bar{\tilde u}_{m+k_0}$ and the $\tilde\delta^{\mathrm{own}}_S$ (\textup{(o1)}, \textup{(o2)} in item~2 of \S\ref{Y-H6-1}, for the run $K'$). Since $\mathbb B_D$ is affine, $\mathfrak n^{\circ}_{k_0}(z)$ and $\mathfrak n_{k_0}(z)$ have the same base point and the same multi-affine jets. By \Cref{Y-H5-LJ}, the coefficient of $z^S$ in $F_\ell(u(z))$ depends only on $u(0)$ and on the $u_{S'}$ with $S'\subset S$. An induction on $k$ therefore shows that $\mathfrak n^{\circ}_k(z)$, $\mathfrak n_k^{\mathrm{pol}}(z)$ for $k\le k_d$, and $\mathfrak n_k(z)$ for $k\ge k_d$ share the same base point and the same multi-affine jets.

We now identify the base points. The base points $\bar u_{k_0}$ and $\bar{\tilde u}_{m+k_0}$, that is, the values at $z=0$, are real systems of members of the two densities without sources; they lie in $M_1$ and obey the strong cap (\Cref{Y-H4-U4}\,(a3), (h); \Cref{Y-H-C9}\,(d), through \Cref{Y-H1-U1}\,(U)). For the short run, \Cref{Y-J-ThA}\,(ii) admits every such system. For the long run, \Cref{Y-J-ThA}\,(iii) asserts that a system $\mathfrak n_k$ exists, without naming one: item~1 of Step~4 of its proof takes $\tilde u$ from \textup{(R8)} in item~3 of \S\ref{Y-J-1}, which is \Cref{Y-GB-ThB}. Yet Steps~4 and~5 of that proof use only three properties of $\tilde u$: it is a real system of members of $\rho^{[K']}_{m+k_0}$, it lies in $M_1$, and it obeys the strong cap; and \textup{(R11)} of the same item, which is \Cref{Y-I-LB}, is stated for every such $\tilde u$. We choose $\tilde u:=\bar{\tilde u}_{m+k_0}$. The $u^\flat$ of \Cref{Y-I-LB} is then $(\tilde E^a+E^\Delta,\tilde E^h,\tilde we^{Q_D})=\mathbb B_D(\bar{\tilde u}_{m+k_0})$, by the formula in \Cref{Y-I-LB}\,(B3) that expresses $w^\flat$ in terms of $\tilde w$, by the definition $\mathbb B_D(\tilde u)=(\tilde E^a+E^\Delta[D],\tilde E^h,w^\flat)$ with $w^\flat:=\tilde we^{Q_D}$ in \Cref{Y-H5-Bc}, and by the closing sentence of \Cref{Y-H5-Bc}; it is real by \Cref{Y-I-LB} and \Cref{Y-H5-Bc}\,(B8). As Steps~4 and~5 of the proof of \Cref{Y-J-ThA} put $\mathfrak n_{k_0+1}=F_{\ell_{k_0}}(u^\flat)$ and $\mathfrak n_{k+1}=F_{\ell_k}(\mathfrak n_k)$, they hold word for word, and $\mathfrak n_k=\bar{\mathfrak n}_k$. In this paragraph $\mathfrak n_{k_0+1}$, $\mathfrak n_{k+1}$ and $\mathfrak n_k$, written without an argument, denote the real system of \Cref{Y-J-ThA}\,(iii) that Steps~4 and~5 of its proof build, not the family $\mathfrak n_k(z)$ of \textup{(D1)} and not its representative. Hence $e_k\le3\vartheta^{k-k_0}$ for $k_0<k\le K$, by \Cref{Y-J-ThA}\,(iv), at the rate $\kappa_1$; Step~5 of the proof of that theorem proves this bound. The properties obtained in this way (reality; membership in $M_1$; the strong cap; and, for each of these $k$, the existence of a real number $\mathsf V$ with $\rho^{[K']}_{m+k}=e^{\mathsf V}\rho[\bar{\mathfrak n}_k;\ell_k]$ at every field, the positivity of this density at every field, and the bound $\mathsf D_{x_k}(\bar u_k,\bar{\mathfrak n}_k)\le3\vartheta^{k-k_0}$) are \textup{(D0)} at $z=0$.

We prove \textup{(D1)}\,(d). At $k_0$ the clause comes from \Cref{Y-H6-J}\,(J2), applied to the run $K'$ at its level $m+k_0$, where the blocks are the $\tilde S_i(k_0)$: the members of $\mathbb B_D(\tilde u)=(\tilde E^a+E^\Delta[D],\tilde E^h,\tilde we^{Q_D})$ at an index set $X$ are computed from those of $\tilde u$ at the same $X$, the partner levels have the same cubes, the same $\check r$ and the same $\mathrm{sep}$, and \textup{(Sep)}, proved below, holds. For the passage from $k$ to $k+1$, assume that $\mathfrak n_k^{\mathrm{pol}}(z)$ can depend on $z_i$ at $X$ only when $X^{+\check r(x_k)}$ meets $\tilde S_i(k)$. \Cref{Y-H1-U1}\,(f6) then makes the mixed derivative over $i\in S$ of $F_{\ell_k}(\mathfrak n_k^{\mathrm{pol}})(Y)$ equal to zero unless $Y^{+\hat\rho'}$ meets each set ${\rm cl}(\tilde S_i(k))=\tilde S_i(k+1)$, see \textup{(o3)} in item~2 of \S\ref{Y-H6-1} for the run $K'$; and $\hat\rho'=\lceil\check r(x_k)/L\rceil+1\le\check r(x_{k+1})$ by Step~1\,(R) of the proof of \Cref{Y-H6-J}.

We prove \textup{(Sep)} and measure lengths in the units of $T_l$. In these units the site $y$ of $\mathbb T_k$ lies at $L^{k-K}y$, and every bare site of either run that belongs to the block of $y$ is at distance at most $(L^{k-K}-\mathrm{sp})/2$ from it in each coordinate; here $\mathrm{sp}=a_K$ or $a_{K'}$ is the spacing of the lattice of that run, and we use that the blocks are centred (item~1 of \S\ref{Y-H2-1}). The origins of the two grids agree (Step~2\,(a) of the proof of \Cref{Y-J-ThA}). Consequently a cube of $\mathbb T_k$ is, for both runs, one and the same box of side $\mathrm{sd}:=ML^{k-K}$, and its bare sites are at distance at least $\mathrm{sp}/2$ from each of its walls.

Now let $\mathrm{cb}$ be a cube in $S_i(k)\cup\tilde S_i(k)$ and $\mathrm{cb}'$ a cube in $S_{i'}(k)\cup\tilde S_{i'}(k)$, where $i\ne i'$, at the distance $\mathrm{dc}$ in cubes; all distances are periodic sup distances. Let $\xi,\xi'$ be base sites in $\mathrm{cb}$ and $\mathrm{cb}'$ of plaquettes $p,p'$ with $f_i(x(p))\ne0$ and $f_{i'}(x(p'))\ne0$, and let $\mathrm{sp},\mathrm{sp}'$ be their spacings. In each coordinate $|\xi-\xi'|\le(\mathrm{dc}+1)\mathrm{sd}-\mathrm{sp}/2-\mathrm{sp}'/2$, while $x(p)-\xi$ is $0$ or $\mathrm{sp}/2$ and $x(p')-\xi'$ is $0$ or $\mathrm{sp}'/2$. Since the sources are separated by $d$, we get $d\le|x(p)-x(p')|\le(\mathrm{dc}+1)\mathrm{sd}-\min\{\mathrm{sp},\mathrm{sp}'\}/2<(\mathrm{dc}+1)\mathrm{sd}$, that is, $\mathrm{dc}\ge\lfloor dL^{K-k}/M\rfloor$. Passing to the neighbourhoods costs $2\,\check r(x_k)$, so the enlarged sets are at distance at least $\mathrm{sep}_k+2$, as in \textup{(G4)} in item~3 of \S\ref{Y-H6-1}.

\emph{Count.} In each coordinate, the base sites belonging to $f_1$, in the two runs taken together, are contained in a single interval of length $2R_f+a_K$. Such an interval meets at most $(2R_f+a_K)L^{K-k-1}/M+2\le2R_fL^{K-k-1}/M+3$ cells of $\mathbb T_{k+1}$, and the neighbourhood adds $2\,\check r(x_{k+1})$; this gives the bound $N_1(k)$. For the lower bound in \textup{(D1)}\,(d) on $\#X$ or $|X|_c$ of a set $X$ that meets $|S|\ge2$ of the enlarged sets, Step~1\,(C) of the proof of \Cref{Y-H6-J}, applied with $\sigma=\mathrm{sep}_k+2$, shows that any connected set of cubes meeting $|S|\ge2$ of these sets contains at least $|S|\,\mathrm{sep}_k/2$ cubes. If $X\in\mathcal P$, then $|X|_c=\#Z$ with $Z$ a connected set of cubes containing $X$ (item~2 of \S\ref{Y-J-1}), and $Z$ meets the sets that $X$ meets.

Finally we bound the majorant at the level $k_0$. As $\bar u_{k_0}\in M_1$ and $\bar{\mathfrak n}_{k_0}\in M_2$, the base points satisfy $\|\bar u_{k_0}-\bar{\mathfrak n}_{k_0}\|\le1+2=3$. On $\mathrm{Pd}(\mathfrak r_{k_0}/2)$, \Cref{Y-H6-J}\,(J5) for the run $K$ gives $\|\sum z^S\delta_S(k_0)\|^{\mathrm{st},(\kappa_{k_0})}\le0.455$. For the jets of $\mathfrak n_{k_0}$ we have $\|\sum z^S\tilde\delta_S(k_0)\|^{(\kappa_{k_0})}\le0.455$ in the weak norm, because $|\tilde w_Se^Q|_w^{\mathrm{wk}}\le|\tilde w_S|_w^{\mathrm{st}}$ by \Cref{Y-H5-Bc}\,(B5). The three bounds add up to $3.91$, which is the number $\Delta_{k_0}$.

\emph{Step 3: the singles in the $a$-channel.} The singles are not carried by the contraction of \S\ref{Y-J-1}: on the radii $\mathfrak r_k$ (item~1.3; \textup{(o5)} in item~2 of \S\ref{Y-H6-1}), for which \textup{(t3)} allows $\mathfrak r_{k+1}/\mathfrak r_k$ up to $L^{1/2}$, the factor $\vartheta[2]$ of \Cref{Y-H5-Af}\,(ii) times this growth is $\vartheta[2]L^{1/2}>1$, and the factor $\vartheta_a$ of \Cref{Y-H1-U1}\,(h3) times it is $\vartheta_aL^{1/2}=0.2431L^{1/4}>1$.

We put $\hat a_k:=a^D_k\,\mathsf b/(u_3(x_k)K_5)$ and show that $\hat a_{k_0}\le4+\Omega_0$ and that $\hat a_{k+1}\le\nu_a\hat a_k+c_ee_k$ for $k_0\le k<k_d$, where $c_e=2\nu_a/14+1/3\le0.36807$ at $k=k_0$ and $c_e=2\nu_a/15+1/3\le0.36575$ for the later $k$. These bounds yield \textup{(D2)}\,(i). Writing $\Omega_k:=\nu_a^{k-k_0}\check\omega_3(x_k)\Omega_0$, we show further that $\Omega_k\le C_5^{\#}K_5^{-1}K\,\nu_a^{k-k_0}L^{-(K-k)/4}\le\Omega^{\#}$ and that $n\,\mathfrak r_ka^D_k\le0.7\,\check\omega_3(x_k)\hat a_k\le0.0054\,\vartheta^{k-k_0}+0.7\,\Omega^{\#}$.

We fix $i$ and write $\mathbb D_k=\mathbb D_k^{[K]}[f_i]$ for the forms of the short run; both runs are decomposed along these forms. Let $\mathrm{rem}_k:=\delta_i^a(k)-s_{x_k}[\mathbb D_k/\mathsf b]$, which is the $\mathrm{rem}_{k,i}$ of \Cref{Y-H6-J}\,(J3), and $\widetilde{\mathrm{rem}}_k:=\tilde\delta_i^a(k)-s_{x_k}[\mathbb D_k/\mathsf b]$. Then $a^D_k=\max_i|\mathrm{rem}_k-\widetilde{\mathrm{rem}}_k|_a$. All these elements lie in $T^a(x_k)$: the jets by \Cref{Y-H5-LS}, and $s_x[D]=-E^\Delta[D]$ by \Cref{Y-H5-Bc}.

We derive the recursion. By \Cref{Y-H1-U1}\,(f1), $E'^{\rm new}$ depends on no member, and $A=E'^{\rm old}$ depends on $E^a$ alone. Taking the coefficient of $z_i$ we find $\delta_i^a(k+1)=DA_{\ell_k}(\bar u_k)[\delta_i^a(k)]$ and $\tilde\delta_i^a(k+1)=DA_{\ell_k}(\bar{\mathfrak n}_k)[\tilde\delta_i^a(k)]$. \Cref{Y-H3-U3} applies with \textup{(SD)}~$=(x_k,\ell_k)$, as $x_k\ge x_3$, at $u=\bar u_k$ and at $u=\bar{\mathfrak n}_k$, which both lie in $M_2$ and so inside $M_8$, and with $D=\mathbb D_k/\mathsf b$, which is symmetric and of range $\le\varrho(x_k)$ by \Cref{Y-H2-U2}. It yields \[\FitDisplay{\dimexpr\linewidth-1pt\relax}{DA_{\ell_k}(u)[s_{x_k}[\mathbb D_k/\mathsf b]]=s_{x_{k+1}}[\mathbb D_{k+1}/\mathsf b]+n[u;\mathbb D_k/\mathsf b];}\] here $D'=\mathrm{cut}_{\varrho(x_{k+1})}(T_{\ell_k}\mathbb D_k)/\mathsf b=\mathbb D_{k+1}/\mathsf b$ is determined by the rows $\ell_k$ alone and is independent of $u$. After the switch the two runs use one row family, so the same $s_{x_{k+1}}[\mathbb D_{k+1}/\mathsf b]$ is subtracted from both. Since $DA_{\ell_k}(u)$ is linear, it follows that $\mathrm{rem}_{k+1}=DA_{\ell_k}(\bar u_k)[\mathrm{rem}_k]+n[\bar u_k;\mathbb D_k/\mathsf b]$ and $\widetilde{\mathrm{rem}}_{k+1}=DA_{\ell_k}(\bar{\mathfrak n}_k)[\widetilde{\mathrm{rem}}_k]+n[\bar{\mathfrak n}_k;\mathbb D_k/\mathsf b]$, and subtraction gives \[\FitDisplay{\dimexpr\linewidth-1pt\relax}{\mathrm{rem}_{k+1}-\widetilde{\mathrm{rem}}_{k+1}=DA_{\ell_k}(\bar{\mathfrak n}_k)[\mathrm{rem}_k-\widetilde{\mathrm{rem}}_k]+\bigl(DA_{\ell_k}(\bar u_k)-DA_{\ell_k}(\bar{\mathfrak n}_k)\bigr)[\mathrm{rem}_k]+\bigl(n[\bar u_k;\mathbb D_k/\mathsf b]-n[\bar{\mathfrak n}_k;\mathbb D_k/\mathsf b]\bigr).}\]

We now bound the three terms at the rate $\kappa_1$, with the caps at $x_k$ on the input side and the caps at $x_{k+1}$ on the output side; here $B=2$ at $k=k_0$ and $B=1$ for the later $k$. A sum over cells dominates the supremum over $Y$. By \Cref{Y-H1-U1}\,(h3) at $u=\bar{\mathfrak n}_k$, the first term is at most $\vartheta_aa^D_k$.

For the second term we fix $g\in T^a(x_k)$ and consider $\Phi(u):=DA_{\ell_k}(u)[g]$, with values in the space that carries the cell-sum seminorm. By \textup{(h3)}, $\Phi$ is bounded there by $\vartheta_a|g|_a$ caps on $\{\|u\|<16\}$. It is holomorphic along affine lines because, by \Cref{Y-H1-U1}\,(h), $DA(u)[g]$ is holomorphic along families. On the open ball this can also be seen directly, since there $DA_{\ell_k}(u)[g]$ is, by \textup{(h)}, the $t$-derivative of $A(u+tg)$ at $0$: the map $(s,t)\mapsto A(u+sg'+tg)$ is holomorphic by \Cref{Y-H1-U1}\,(c+), hence the $t$-derivative at $0$ depends holomorphically on $s$. \Cref{Y-H5-LH}\,(b) and \Cref{Y-H6-J}\,(J3) now bound the second term by \[\vartheta_a|\mathrm{rem}_k|_a\|\bar u_k-\bar{\mathfrak n}_k\|/(16-B)\le\vartheta_a\,(2u_3(x_k)K_5/\mathsf b)\,e_k/(16-B).\]

In the third term the part $n^{\mathrm{fix}}$ depends on no member, so it cancels. By \Cref{Y-H3-U3}\,(d) on $M_2$, the difference of the parts $n^{\mathrm{mem}}$ is at most $1/3$ of their bound in \textup{(a)} times the $a$-part of $\mathsf D^{\mathrm{in}}$, which does not exceed $e_k$. Call that bound $B_{\mathrm{mem}}$; it is a part of the sum that \textup{(a)} bounds by $u_3(x_{k+1})\|D\|$. Indeed, \textup{(Z.4)}, in item~1 of \S\ref{Y-H3-1}, implies $18.2L^3c_a(x)\le e^{\kappa+3}L^3x^{-1/4}$, hence $B_{\mathrm{mem}}$ is at most the term $1001e^{\tilde\kappa+3}L^3x^{-1/4}\sigma_D$ of the last sentence of \Cref{Y-H3-U3}\,(a). As $\|\mathbb D_k/\mathsf b\|\le K_5/\mathsf b$, the third term is at most $(1/3)u_3(x_{k+1})(K_5/\mathsf b)e_k$. Dividing the three bounds by $u_3(x_{k+1})K_5/\mathsf b$ and using \textup{(t1)}, we arrive at $\hat a_{k+1}\le\nu_a\hat a_k+(2\nu_a/(16-B)+1/3)e_k$, which is the recursion of the claim.

\emph{Start.} At the start, the $a$-jets pass through $\mathbb B_D$ unchanged, by \textup{(D1)}\,(b). \Cref{Y-H6-J}\,(J3) for the run $K'$ gives $\tilde\delta_i^a(k_0)=s_x[\tilde{\mathbb D}_{m+k_0}/\mathsf b]+\widetilde{\mathrm{rem}}^{\mathrm{own}}$ with $|\widetilde{\mathrm{rem}}^{\mathrm{own}}|_a\le2u_3K_5/\mathsf b$, where $u_3$ and $\mathrm{amp}$ are taken at $x=x_{k_0}$. Therefore $\mathrm{rem}_{k_0}-\widetilde{\mathrm{rem}}_{k_0}=\mathrm{rem}_{k_0}-\widetilde{\mathrm{rem}}^{\mathrm{own}}-s_x[(\tilde{\mathbb D}_{m+k_0}-\mathbb D_{k_0})/\mathsf b]$. \Cref{Y-H5-Bc}\,(B2) gives $|s_x[D]|_a\le c_Q^{\rm su}p_0^2\|D\|/c_a(x)=(\mathrm{amp}(x)\mathsf b/K_5)\|D\|$; since this bound is homogeneous and $s_x$ is linear in $D$, it holds also when $\delta_\Delta>1$. The terms $\mathrm{rem}_{k_0}$ and $\widetilde{\mathrm{rem}}^{\mathrm{own}}$ contribute $2u_3K_5/\mathsf b$ each, by \textup{(J3)} for the two runs, and we obtain $a^D_{k_0}\le(\mathrm{amp}/K_5)\|\tilde{\mathbb D}_{m+k_0}-\mathbb D_{k_0}\|+4u_3K_5/\mathsf b$; by \textup{(t2)} this is $\hat a_{k_0}\le\Omega_0+4$.

\emph{Solution.} Recall that $e_{k_0+i}\le3\vartheta^i$ for $i\ge0$. As both values of $c_e$ are at most $0.36807$, for $j\ge1$ we have \[\hat a_{k_0+j}\le\nu_a^j(4+\Omega_0)+3\cdot0.36807\sum_{i<j}\nu_a^{j-1-i}\vartheta^i\le\nu_a^j\Omega_0+\bigl[4(\nu_a/\vartheta)^j+1.10421\,\vartheta^{-1}/(1-\nu_a/\vartheta)\bigr]\vartheta^j.\] Here $\nu_a/\vartheta\le0.54410$, $1/(1-\nu_a/\vartheta)\le2.19343$ and $1.10421\cdot2.19343/0.4468\le5.4208$, while $4\cdot0.54410=2.1764$; the sum $7.5972\le7.6$ bounds the bracket.

In units of the radius, \textup{(t2)} and \textup{(t3)} give $n\,\mathfrak r_ka^D_k=3.5\,(\mathrm{amp}(x_k)/\mathrm{amp}(Lx_k))\,\check\omega_3(x_k)\hat a_k\le0.7\,\check\omega_3(x_k)\hat a_k$. By \textup{(W3)}, the part with $\vartheta$ is at most $0.7\cdot10^{-3}\cdot7.6=0.00532$ times $\vartheta^{k-k_0}$, and the form part is $0.7\,\Omega_k$. By \textup{(t4)}, by \Cref{Y-H2-U5}\,(S3c) at $k_0$ (with $\|f_i\|_\sharp\le1$, $k_0\le K$ and $\vartheta_5\le L^{-1/5}$) and by Step~1\,(d), \[\Omega_k=\nu_a^{k-k_0}(\check\omega_3(x_k)/\check\omega_3(x_{k_0}))\|\tilde{\mathbb D}-\mathbb D\|/K_5\le C_5^{\#}K_5^{-1}K\,\nu_a^{k-k_0}L^{(k-k_0)/4-k_0/5},\] and $(k-k_0)/4-k_0/5=(K-k_0)/4-k_0/5-(K-k)/4\le-(K-k)/4$. Since $\nu_a\le0.3683$ and $L^{-1/4}\le10^{-4}\le0.3683$, it follows that $\Omega_k\le\Omega^{\#}$, also at $k=k_0$.

\emph{Step 4: the singles in the $h$- and $w$-channels.} Let $\hat h_k:=n\,\mathfrak r_kh^D_k$. We show that $\hat h_{k_0+1}\le0.2158+0.25\,\Omega^{\#}$ and that, for $k_0<k<k_d$, \[\hat h_{k+1}\le(1/15)\bigl(3.5\,\check\omega_3(x_k)\hat a_k+L^{-1/2}\hat h_k\bigr)+(e_k/56.25)\bigl(3.5(1+2\check\omega_3(x_k))+0.42L^{-1/2}\bigr);\] consequently $\hat h_k\le0.49\,\vartheta^{k-k_0}+0.25\,\Omega^{\#}$ for $k_0<k\le k_d$.

Taking the coefficient of $z_i$ in the $h$- and $w$-members, we get $\delta_i^{hw}(k+1)=DF_{\ell_k}^{hw}(\bar u_k)[\delta_i(k)]$ and $\tilde\delta_i^{hw}(k+1)=DF_{\ell_k}^{hw}(\bar{\mathfrak n}_k)[\tilde\delta_i(k)]$. Their difference splits as \[DF_{\ell_k}^{hw}(\bar{\mathfrak n}_k)[\delta_i-\tilde\delta_i]+\int_0^1D^2F_{\ell_k}^{hw}(u_\tau)[\bar u_k-\bar{\mathfrak n}_k,\delta_i]\,d\tau,\qquad u_\tau=\bar{\mathfrak n}_k+\tau(\bar u_k-\bar{\mathfrak n}_k),\] and $u_\tau$ lies in the convex set $M_B$, with $B$ as in Step~3.

In the mixed entry $D^2F_{\ell_k}^{hw}(u_\tau)[g_1,g_2]$ the first argument $g_1=\bar u_k-\bar{\mathfrak n}_k$ is a whole triple. \Cref{Y-H5-Af}\,(iv) covers this case. In that clause an argument $g$ that is a mixed triple is measured by $\|g\|$ at $p$ if it is left whole, and by the sum $|g^a|_a+\eta_h|g^h|_h+|g^{w\circ}|_w^{\mathrm{wk}}$ if it is split, where $g^{w\circ}:=g^w1\{X\text{ holds a plaquette of }D(V)\}$. We keep $g_1$ whole, with $\|g_1\|\le e_k$, and split only $g_2=\delta_i$. If $g_2=(g^a,0,0)$, we use \Cref{Y-H5-LH}\,(a) with $j=2$ and $A_{\#}=16$ in the norm with the pair $p=(\kappa_1,c_h)$, where $\mathsf m$ is one cap of $E'^h$ for the output $E'^h$ and the strong cap of $w'$ for the output $w'$ (\Cref{Y-H1-U1}\,(U); \Cref{Y-H-C9}\,(b)); this gives the bound $4\|g_1\||g^a|_a/(16-B)^2$. If $g_2=(0,0,g^w)$, then \Cref{Y-H1-U1}\,(f4) allows us to replace $g^w$ by $g^{w\circ}$, and the same lemma applies, now with $|g^{w\circ}|_w^{\mathrm{wk}}\le\eta_w|g^w|_w^{\mathrm{st}}$ by \Cref{Y-H5-Af}\,(iv). If $g_2=(0,g^h,0)$, we use the same lemma in the norm with $(\kappa_1,c_h/2)$, which is a pair of $\mathbb P$. At this pair neither the norm of $g_1$ nor that of $u_\tau$ grows, since the unit of $h$ is larger; the size of $g^h$ is $e^{-(c_h/2)p_0(x)^2}|g^h|_h$; and to bring $E'^h$ back to the unit $c_h(x')$ costs $e^{(c_h/2)p_0(x')^2}$. These factors together give $\eta_h$, while for $w'$ the bound holds as it stands: nothing needs to be brought back to that unit, in which only $E'^h$ is measured. Summing the three cases, \[|D^2F_{\ell_k}^{hw}(u_\tau)[g_1,g_2]|\le(4/(16-B)^2)\|g_1\|\bigl(|g_2^a|_a+\eta_h|g_2^h|_h+\eta_w|g_2^w|_w^{\mathrm{st}}\bigr).\]

First let $k>k_0$. We use the channels of \Cref{Y-H5-Af}\,(iv) with $B=1$ and the rate $\kappa_1$, with strong units on the output side. The first term is then at most $(a^D_k+(\eta_h+\eta_w)h^D_k)/15$, and the second term at most $(4/225)e_k\bigl(|\delta_i^a|_a+(\eta_h+\eta_w)\max\{|\delta_i^h|_h,|\delta_i^w|_w^{\mathrm{st}}\}\bigr)$. Here $|\delta_i^a(k)|_a\le\mathrm{amp}(x_k)(1+2\check\omega_3(x_k))$, by \Cref{Y-H6-J}\,(J3), \Cref{Y-H5-Bc}\,(B2), $\|\mathbb D_k\|\le K_5$ and \textup{(t2)}; and the maximum is at most $0.12\,\mathrm{amp}(Lx_k)$ by \Cref{Y-H6-J}\,(J4). We multiply by $n\,\mathfrak r_{k+1}=3.5/\mathrm{amp}(x_k)$. Then $n\,\mathfrak r_{k+1}a^D_k=3.5\,\check\omega_3(x_k)\hat a_k$, $n\,\mathfrak r_{k+1}h^D_k=(\mathfrak r_{k+1}/\mathfrak r_k)\hat h_k\le L^{1/2}\hat h_k$, and $n\,\mathfrak r_{k+1}\,0.12\,\mathrm{amp}(Lx_k)\le0.42L^{1/2}$. Finally \textup{(W2)} at $x_k$ turns $L^{1/2}(\eta_h+\eta_w)$ into $L^{-1/2}$, and $225/4=56.25$; this is the recursion of the claim.

Now let $k=k_0$, so that $B=2$ and $e_{k_0}=3$; we write $x=x_{k_0}$. For the $w$-jets of $\mathfrak n_{k_0}$ our bound is in the weak norm, so we use the channel $w\to h,w$ of \Cref{Y-H5-Af}\,(iv) in its first form, $\le|g^{w\circ}|_w^{\mathrm{wk}}/(16-B)$. The same clause gives $|\delta_i^{w\circ}|_w^{\mathrm{wk}}\le\eta_w|\delta_i^w|_w^{\mathrm{st}}$, and \Cref{Y-H5-Bc}\,(B5) gives $|(\tilde\delta^{\mathrm{own},w}_ie^Q)^{\circ}|_w^{\mathrm{wk}}\le e^{1/2}\eta_w|\tilde\delta^{\mathrm{own},w}_i|_w^{\mathrm{st}}$. Hence the $w$-input contributes at most $(1+e^{1/2})\eta_w\,0.12\,\mathrm{amp}(Lx)/14$, and the $h$-input at most $\eta_h\,0.24\,\mathrm{amp}(Lx)/14$ (\Cref{Y-H5-Af}\,(iv)). Altogether this is at most $0.24\,\mathrm{amp}(Lx)(\eta_h+e^{1/2}\eta_w)/14\le0.24\,\mathrm{amp}(Lx)/(14L)$. The $a$-input contributes $a^D_{k_0}/14$, and $3.5\,\check\omega_3(x_{k_0})\hat a_{k_0}\le3.5(4\cdot10^{-3}+\Omega_{k_0})\le0.014+3.5\,\Omega^{\#}$. The second term is $(4/196)\cdot3$ times the same bracket as for $k>k_0$. Therefore \[\FitDisplay{\dimexpr\linewidth-1pt\relax}{\hat h_{k_0+1}\le(1/14)(0.014+3.5\,\Omega^{\#}+0.84L^{-1/2})+(12/196)(3.5\cdot1.002+0.42L^{-1/2})\le0.001+0.25\,\Omega^{\#}+0.21472\le0.2158+0.25\,\Omega^{\#},}\] and $0.2158\le0.48300\,\vartheta\le0.49\,\vartheta$.

For the induction, let $\hat h_k\le0.49\,\vartheta^j+0.25\,\Omega^{\#}$ with $j=k-k_0\ge1$. By Step~3, $\check\omega_3(x_k)\hat a_k\le7.6\cdot10^{-3}\vartheta^j+\Omega_k$. In the recursion the $\vartheta$-part is at most \[(3.5/15)\,7.6\cdot10^{-3}+3\cdot3.507/56.25+10^{-8}\le0.00178+0.18704+10^{-8}\le0.18883\le0.42263\,\vartheta\] times $\vartheta^j$, and the $\Omega$-part is at most $(3.5/15)\Omega^{\#}+10^{-8}\Omega^{\#}\le0.2334\,\Omega^{\#}\le0.25\,\Omega^{\#}$; so the bound passes from $k$ to $k+1$.

The recursion closes because the system is triangular. A single has an $a$-part smaller than $1/\mathfrak r_k$ by the factor $L^{-1/2}$, because the radius $\mathfrak r_k$ is defined in terms of $\mathrm{amp}(Lx_k)$; the factor $1/L$ is carried by the channels $h,w\to h,w$; and by \Cref{Y-H1-U1}\,(f1) the channels $h,w\to a$ vanish.

\emph{Step 5: the sets with $|S|\ge2$ on the growing radii.} Fix $k$ with $k_0\le k<k_d$, and suppose that $u_k^{\mathrm{pol}}$ and $\mathfrak n_k^{\mathrm{pol}}$ are families in $M_B(x_k;\kappa_k,c_h)$, where $B\le2.6$, on the polydisc $\mathrm{Pd}(\mathfrak r_k/2)$. We claim that \[\sum_{|S|\ge2}(\mathfrak r_{k+1}/2)^{|S|}\|\delta_S(k+1)-\tilde\delta_S(k+1)\|^{\mathrm{st},(\kappa_{k+1})}\le c_B\Delta_k,\qquad c_B=0.0436895\,\vartheta[B];\] in particular $c_2\le0.012383$ and $c_{2.455}\le0.012799$.

By the triangle inequality, \[\sup_{\mathrm{Pd}(\mathfrak r_k/2)}\mathsf D^{\mathrm{in},(\kappa_k)}(u_k^{\mathrm{pol}},\mathfrak n_k^{\mathrm{pol}})\le e_k+\sum_S(\mathfrak r_k/2)^{|S|}\|\delta_S-\tilde\delta_S\|^{(\kappa_k)}\le\Delta_k;\] here the $n$ singles contribute $(n\,\mathfrak r_k/2)\max\{a^D_k,h^D_k\}$, norms at the rate $\kappa_1$ dominate those at $\kappa_k$, and at $k_0$ the bound is that of Step~2. For $k>k_0$, the same bound with strong norms is the first claim of \textup{(D2)}. \Cref{Y-H5-Af}\,(ii) with the pair $p=(\kappa_k,c_h)\in\mathbb P$ gives $\sup\mathsf D^{(\kappa_k)}_{x_{k+1}}(F_{\ell_k}u_k^{\mathrm{pol}},F_{\ell_k}\mathfrak n_k^{\mathrm{pol}})\le\vartheta[B]\Delta_k$. \Cref{Y-H5-LS} in the norm $\|\cdot\|^{\mathrm{st},(\kappa_k)}$ at $x_{k+1}$, together with \Cref{Y-H6-J}\,(J1) and the definition of $\tilde\delta_S(k+1)$, then yields $\|\delta_S(k+1)-\tilde\delta_S(k+1)\|^{\mathrm{st},(\kappa_k)}\le\vartheta[B]\Delta_k(\mathfrak r_k/2)^{-|S|}$.

\Cref{Y-H6-J}\,(J2) and \textup{(D1)}\,(d) show that the $a$- and $h$-parts of both jets are zero at $Y$ unless $\#Y$ is at least $|S|\,\mathrm{sep}_{k+1}/2$, and their $w$-parts unless $|Y|_c$ is at least this number; therefore the passage to the rate $\kappa_{k+1}=\kappa_k-\check\kappa_k$ gains the factor $e^{-\check\kappa_k(|S|\mathrm{sep}_{k+1}/2-1)}$. With \textup{(t3)}, \[q:=(\mathfrak r_{k+1}/\mathfrak r_k)e^{-\check\kappa_k\mathrm{sep}_{k+1}/2}\le L^{1/2}/(4nL^{1/2})=1/(4n),\] the factor $1/2$ of the radius cancelling in the quotient; and $e^{\check\kappa_k}\le e^{1/4}$ by Step~1\,(e). Multiplying the above bound on $\|\delta_S(k+1)-\tilde\delta_S(k+1)\|^{\mathrm{st},(\kappa_k)}$ by $(\mathfrak r_{k+1}/2)^{|S|}$ and by the gain, and summing over the sets $S$ with $|S|\ge2$, we find that the sum of the claim is at most $\vartheta[B]\Delta_ke^{1/4}[(1+q)^n-1-nq]$.

Since $q\le1/(4n)$, the bracket is at most $(1+1/(4n))^n-5/4$, which increases with $n$ towards $e^{1/4}-5/4\le0.0340255$. The product of the latter, $e^{1/4}(e^{1/4}-5/4)$, has the decimal expansion $0.0436894998\dots$ and is therefore a number $\le0.0436895$. Finally, $0.0436895\,\vartheta[2]\le0.012383$ and $0.0436895\,\vartheta[2.455]\le0.012799$.

\emph{Step 6: the induction closes; \textup{(D1)}\,(e) and \textup{(D2)}.} We claim: when $k_0<k\le k_d$, the two polynomials belong to $M_2(x_k;\kappa_k,c_h)$ on $\mathrm{Pd}(\mathfrak r_k/2)$, $\Delta_k\le3.36\,\vartheta^{k-k_0}+0.36\,\Omega^{\#}$, and $\Delta_k-e_k\le0.195$.

Let $\check\Sigma_k:=(1/2)\max\{n\,\mathfrak r_ka^D_k,\hat h_k\}$. Neither Step~3 nor Step~4 uses the membership of a polynomial $u_k^{\mathrm{pol}}$ or $\mathfrak n_k^{\mathrm{pol}}$, so the bounds of these two steps can be used in the induction below. From Steps~3 and~4 we obtain $\check\Sigma_k\le\max\{0.0027\,\vartheta^j+0.35\,\Omega^{\#},\,0.245\,\vartheta^j+0.125\,\Omega^{\#}\}\le0.245\,\vartheta^j+0.35\,\Omega^{\#}$, where $j=k-k_0\ge1$; and $e_k\le3\vartheta^j$ by Step~2.

For the first step, Step~2 gives $u_{k_0}^{\mathrm{pol}}\in M_{1.455}$, $\mathfrak n_{k_0}\in M_{2.455}$ and $\Delta_{k_0}=3.91$. Step~5 with $B=2.455$ bounds the sets with $|S|\ge2$ at $k_0+1$ by $0.012799\cdot3.91=0.05004409$; this is at most $0.112006\,\vartheta=0.05004428\dots$, and at most $0.050045$. Hence $\Delta_{k_0+1}\le(3+0.245+0.112006)\vartheta+0.35\,\Omega^{\#}=3.357006\,\vartheta+0.35\,\Omega^{\#}$.

Membership at a level $k>k_0$ follows from the bound at $k$: \[\|\mathfrak n_k^{\mathrm{pol}}(z)\|^{\mathrm{st},(\kappa_k)}\le\|\bar{\mathfrak n}_k\|^{\mathrm{st}}+\|\sum z^S\delta_S(k)\|^{\mathrm{st},(\kappa_k)}+\|\sum z^S(\tilde\delta_S-\delta_S)(k)\|^{\mathrm{st},(\kappa_k)}\le1+0.455+(\Delta_k-e_k);\] the three bounds come from \textup{(D1)}\,(c) (\Cref{Y-J-ThA}\,(iii)), from \Cref{Y-H6-J}\,(J5) and from the definition of $\Delta_k$. Here $\Delta_k-e_k=\check\Sigma_k+(\text{the sets }|S|\ge2)\le0.245\cdot0.4468+0.35\cdot0.1+0.050045\le0.19452$, for the last summand is largest at $k_0+1$: at the later levels it is at most $0.012383(3.36\,\vartheta+0.036)\le0.01904$. Hence the norm is at most $1.65\le2$; for $u_k^{\mathrm{pol}}$ it is at most $1.455$.

For the step from $k$ to $k+1$, with $k_0<k<k_d$, Step~5 with $B=2$ gives \[\FitDisplay{\dimexpr\linewidth-1pt\relax}{\begin{aligned}
&\Delta_{k+1}\le e_{k+1}+\check\Sigma_{k+1}+0.012383\,\Delta_k
\\
&\quad{}\le[3.245+0.012383\cdot3.36/0.4468]\vartheta^{j+1}
\\
&\quad{}+[0.35+0.012383\cdot0.36]\Omega^{\#}
\\
&\quad{}\le3.338122\,\vartheta^{j+1}+0.354458\,\Omega^{\#}.
\end{aligned}}\] Both coefficients are below $3.36$ and $0.36$, respectively. The first bound of \textup{(D2)}\,(iii), $\Delta_k\le3.36\,\vartheta^{k-k_0}+0.36\,\Omega^{\#}$, therefore holds; it implies the second bound of \textup{(D2)}\,(iii), because $3.36\le3.45$ and $0.36\le1.36$. The last clause of \textup{(D2)}\,(iii) follows from $\Delta_k-e_k\le0.195\le0.34$. For \textup{(D1)}\,(e), let $k_0<k\le k_d$. On $\mathrm{Pd}(\mathfrak r_k/2)$ the difference $\mathfrak n_k^{\mathrm{pol}}(z)-\bar{\mathfrak n}_k$ is the sum of the two quantities whose norms are the second and third summands in the display above that gives membership, so that $\|\mathfrak n_k^{\mathrm{pol}}(z)-\bar{\mathfrak n}_k\|^{\mathrm{st},(\kappa_k)}\le0.455+(\Delta_k-e_k)\le0.455+0.195\le0.455+0.34$; there, moreover, $\|u_k^{\mathrm{pol}}(z)-\bar u_k\|^{\mathrm{st},(\kappa_k)}\le0.7/2+0.21/4\le0.455$ by \Cref{Y-H6-J}\,(J5). Both polynomials belong to $M_2(x_k;\kappa_k,c_h)$, as shown after that display.

\emph{Step 7: the late levels; \textup{(D1)}\,(f), \textup{(D3)} and the depth $s$.} We prove \textup{(D1)}\,(f), the bound $\Delta_k\le\vartheta^{k-k_d}\Delta_{k_d}$ for $k_d\le k\le K$, and \textup{(D3)} both at $s=0$ and at each depth $s$ for which $K-s>k_0$. The rate at $k_d$ equals $\kappa_{k_d}=\kappa$, and by Step~6, as $k_d>k_0$, the two polynomials are families in $M_2(x_{k_d};\kappa,c_h)$ on $\mathrm{Pd}(\mathfrak r_*/2)$. The hypothesis \textup{(SD)} holds at every level $k<K$, and $k_d\le K-1$. \Cref{Y-H5-Af}\,(i) shows that the images belong to $M_1$, as families with the strong cap, at every level up to $K$, each of them being an output of $F$. \Cref{Y-H5-Af}\,(ii), first with $B=2$ and then with $B=1$, gives at each level a factor $\vartheta[2]\le0.28343\le\vartheta$, with the weak distance on the input side and the strong distance on the output side; hence $\Delta_k\le\vartheta^{k-k_d}\Delta_{k_d}$. The same clause, applied to the pairs $(u_k(z),\bar u_k)$ and $(\mathfrak n_k(z),\bar{\mathfrak n}_k)$, whose second entries are constant families, and started at $k_d$ from $0.455$ and from $0.455+0.34$, gives \textup{(D1)}\,(f). The jets are the true ones: for $\mathfrak n_k(z)$ by Step~2 and \textup{(D1)}\,(c), and for $u_k(z)$ by \Cref{Y-H6-J}\,(J6).

At the last level this gives \[\sup\mathsf D_{\mathsf b}^{(\kappa)}\le\vartheta^{K-k_d}(3.45\,\vartheta^{k_d-k_0}+1.36\,\Omega^{\#})\le(3.45+1.36\,C_5^{\#}K_5^{-1}K)\vartheta^{K-k_0},\] because $0.3683\le\vartheta$. By Step~5 of the proof of \Cref{Y-J-ThA}, $3\vartheta^{K-k_0}\le C_{\mathsf b}\theta^K$; and $3.45/3=1.15$ and $1.36/3\le0.46$.

Now let $k=K-s>k_0$. In both cases, namely for $k\le k_d$ by \textup{(D2)} and for $k>k_d$ by the contraction $\Delta_k\le\vartheta^{k-k_d}\Delta_{k_d}$ together with \textup{(D2)} at $k_d$, as at the last level, we have $\Delta_k\le3.45\,\vartheta^{k-k_0}+1.36\,\Omega^{\#}$, and $\Omega^{\#}\le C_5^{\#}K_5^{-1}K\,\vartheta^{K-k_0}\le C_5^{\#}K_5^{-1}K\,\vartheta^{k-k_0}$. Step~6\,(a) of the proof of \Cref{Y-J-ThA} shows that $K-s>k_0$ is equivalent to $K-s\ge K_0(L^s\mathsf b)$, in which case $3\vartheta^{K-s-k_0}\le C_{L^s\mathsf b}\theta^{K-s}$. Finally $k\le k_d$ if and only if $s\ge j^*$, and in that case $\mathfrak r_k=\mathfrak r(L^s\mathsf b)$.

\emph{Step 8: the identity \textup{(D6)}.} All derivatives in \textup{(D6)} are taken on the representatives of \textup{(D1)}. \emph{Normalisation.} Since the kernels of the step are probability kernels (item~2 of \S\ref{Y-J-1}), Fubini's theorem and linearity give $\int\rho'\,dV'=\int\rho\,dU$ for every bounded complex Borel function $\rho$. It follows that $\int\rho_k^z\,dV=Z_K(z)$ at every level.

We work on the interior of $\mathrm{Pd}_0$ of the run $K'$. For the run $K$, \Cref{Y-H4-U4}\,(h) gives \[Z_K(z)=\exp\{\sum_{0\le k<K}\mathsf V'_{\ell_k}[u_k(z)]\}\,I[u_K(z)].\] For the run $K'$ we apply \Cref{Y-H4-U4}\,(h) as far as the level $m+k_0$, next \Cref{Y-H5-Bc}\,(B1), and finally \textup{(D0)}. Since \textup{(B1)} is the identity $\rho[\tilde u;\ell;D]=\rho[\mathbb B_D\tilde u;\ell]$ at every field, without a prefactor $e^{\mathsf V}$, the switch adds nothing to the exponent of $Z_{K'}(z)$, and we obtain \[Z_{K'}(z)=\exp\{\sum_{0\le j<m+k_0}\mathsf V'_{\tilde\ell_j}[\tilde u_j(z)]+\sum_{k_0\le k<K}\mathsf V'_{\ell_k}[\mathfrak n^{\circ}_k(z)]\}\,I[\mathfrak n^{\circ}_K(z)].\] At $z=0$ all factors are $>0$; we take the logarithm near $0$ and apply $\partial_1\cdots\partial_n$ at $0$. For the run $K$ we have $\rho_0^z=\rho_0\exp\{\sum z_i\mathcal O_K(f_i)\}$, and for the run $K'$ we have $\rho_0^z=\rho_0\exp\{\sum z_i\mathcal O_{K'}(f_i)\}$ with the bare weight $\rho_0$ of that run; in both cases $\rho_0^z$ is the weight of \Cref{Y-H4-U4} taken at $\nu_p=\mathsf b^{-1}\sum z_if_i(x(p))$, where $p$ is a plaquette of the bare lattice of the run. Hence the truncated function $S^T_n$ of item~1 of \S\ref{Y-H6-1}, at the same label $\mathsf b$ and on the same torus, is $S^T_{n;K}=\partial_1\cdots\partial_n\log Z_K(0)$ for the run $K$ and $S^T_{n;K'}=\partial_1\cdots\partial_n\log Z_{K'}(0)$ for the run $K'$; the left side of \textup{(D6)} consists of these two terms.

\emph{Representatives.} Consider a function $\Phi$ that is holomorphic along affine lines in a ball about $u(0)$ and bounded on that ball. Since \Cref{Y-H5-LJ} holds for every such $\Phi$, it gives \[\partial_1\cdots\partial_n(\Phi\circ u)(0)=\sum_{\pi}D^{|\pi|}\Phi(u(0))[u_{S_1},\dots,u_{S_{|\pi|}}],\] where $\pi$ runs over the partitions of $\{1,\dots,n\}$ and $S_1,\dots,S_{|\pi|}$ are the blocks of $\pi$, as in \Cref{Y-H5-LJ}; since every variable is differentiated exactly once, the right side involves only $u(0)$ and the multi-affine jets. We apply this to $\mathsf V'_\ell[\cdot]$, to every $v(Y;\cdot)$, and, near a base point at which $I>0$, to $\log I$. For $I$ the reason is the following. For every $V$, $\rho[u;\ell_K](V)$ equals a polynomial in finitely many values of $w$ times the exponential of a linear form in finitely many values of $E^a$ and $E^h$, and it is bounded on $M_{16}$ uniformly in $V$, the torus being finite; therefore $I$ is bounded and holomorphic along affine lines, by the theorems of Fubini and Morera. Moreover $I[\bar u_K]>0$ and $I[\bar{\mathfrak n}_K]>0$, by \Cref{Y-J-ThA}\,(ii), (iii) taken as in \textup{(D1)}\,(c): both densities are $>0$ at every field.

\emph{The ball.} We take the norm at $p=(\kappa,c_h)$, the pair of the last level. The reason given for $I$ holds at every $p$ (\Cref{Y-H5-LJ}, its part on $I[u]$); in this norm both the honest families and the representatives at the level $K$ remain bounded (\Cref{Y-H4-U4}\,(h), \textup{(D0)}, \textup{(D1)}\,(f)); and each of the two base points $u_0$ has $\|u_0\|\le1$, since the norm at the rate $\kappa_1$ dominates. Let $\mathsf m$ be a bound of $|I|$ on $\{\|u\|<16\}$. For $\|g\|\le1$, the map $s\mapsto I[u_0+sg]$ is holomorphic and bounded by $\mathsf m$ on $|s|<15$. Applying Cauchy's formula on circles of radius $R<14$ centred at the points with $|s|\le1$ and letting $R\to14$, we see that at these points the derivative has modulus at most $\mathsf m/14$. Taking $s=\|u-u_0\|>0$ and $g=(u-u_0)/s$, we get $|I[u]-I[u_0]|\le\mathsf m\|u-u_0\|/14$ for $\|u-u_0\|\le1$; this is \Cref{Y-H5-LH}\,(b) with $A_{\#}=16$, $B=2$ for the single function $I$, carried out on a line.

Put $t:=\min\{1,7I[u_0]/\mathsf m\}$. For $\|u-u_0\|<t$ the quotient $I[u]/I[u_0]$ then lies in $\{|w-1|\le1/2\}$. So $\mathrm{Log}(I[u]/I[u_0])$, where $\mathrm{Log}$ is the principal branch of the complex logarithm on $|w-1|\le1/2$, is defined on the ball $\{\|u-u_0\|<t\}$, bounded there by $\log2$ and holomorphic along lines, and $\Phi[u]:=\log I[u_0]+\mathrm{Log}(I[u]/I[u_0])$ satisfies the hypothesis of \Cref{Y-H5-LJ}.

To conclude, \Cref{Y-H6-J}\,(J1), (J6) and \textup{(D1)}\,(c) allow us to replace the honest families by $u_k^{\mathrm{pol}}$, $\mathfrak n_k^{\mathrm{pol}}$ for $k\le k_d$ and by $u_k(z)$, $\mathfrak n_k(z)$ for $k\ge k_d$. Subtracting the two identities, we obtain \textup{(D6)}.

\emph{Step 9: \textup{(D4)} and \textup{(D5)}; the numbers $\mathsf V'$ and the two series.} We prove \textup{(D4)} and \textup{(D5)}, and that, once $j>j^*$, consecutive terms of $\mathrm{Tail}$, and likewise of $\mathrm{Ser}_2$, have quotient at most $1/2$.

For \textup{(D4)}, fix $k$ with $k_0\le k<K$ and a radius $r<\mathfrak r_{\min(k,k_d)}/2$, and let $u(z)$, $\mathfrak n(z)$ be the two representatives; they are families in $M_B(x_k;\kappa_k,c_h)$ with $B\le2.455$. By \Cref{Y-H1-U1}\,(g1), (f1), $\mathsf V'[\mathfrak n(z)]-\mathsf V'[u(z)]=\sum_Y[v(Y;\mathfrak n(z))-v(Y;u(z))]$, since $V_0^+$ depends on no member. We determine which $Y$ contribute. For $u$, the dependence on $z_i$ at $X$ requires that $X^{+\check r(x_k)}$ meet $S_i(k)$, and for $\mathfrak n$ that it meet $\tilde S_i(k)$; for the polynomials this is \Cref{Y-H6-J}\,(J2) and \textup{(D1)}\,(d), and after $k_d$ it follows level by level from \Cref{Y-H1-U1}\,(f5). By \Cref{Y-H1-U1}\,(f6), $\partial_1\cdots\partial_nv(Y;\cdot)(0)=0$ unless $Y^{+\check r(x_{k+1})}$ meets, for every $i$, the set $S_i(k+1)$ for $v(Y;u(z))$ and the set $\tilde S_i(k+1)$ for $v(Y;\mathfrak n(z))$. Every such $Y$ contains one of the $N_1(k)$ cells that lie within $\check r(x_{k+1})$ of $S_1(k+1)\cup\tilde S_1(k+1)$, by the part ``Count'' of \textup{(Sep)} in Step~2; and, for $k<k_d$, $\#Y\ge m_k$ follows from \textup{(Sep)} at $k+1$ and Step~1\,(C) of the proof of \Cref{Y-H6-J}.

For these $Y$ we write the mixed derivative at $0$ of $v(Y;\mathfrak n(z))-v(Y;u(z))$ by Cauchy's formula: \[\partial_1\cdots\partial_n[v(Y;\mathfrak n(z))-v(Y;u(z))](0)=(2\pi i)^{-n}\int_{|z_i|=r}[v(Y;\mathfrak n(z))-v(Y;u(z))]\prod dz_i/z_i^2.\] Moving the moduli inside the integral and summing over $Y$, we apply at each $z$ and at each anchor cell \Cref{Y-H5-Af}\,(v) with the rate $\kappa_k\ge\kappa$; it gives $\sum_{Y\text{ holding the cell},\,\#Y\ge m_k}|v(Y;\mathfrak n(z))-v(Y;u(z))|\le e^{-\kappa(m_k-1)}16C_gc_a(x_k)\mathsf D^{\mathrm{in}}/(16-B)$, where $16/13.545\le1.18125\le1.2$ and $\mathsf D^{\mathrm{in}}\le\Delta_k$. Then we let $r$ tend to the radius.

For the sum over the levels, $\Delta_{k_0}=3.91\le4$, and for $k>k_0$ we have $\Delta_k\le3.45\,\vartheta^{k-k_0}+1.36\,\Omega^{\#}$ by Steps~6 and~7, with $\Omega^{\#}\le C_5^{\#}K_5^{-1}K\,\vartheta^{K-k_0}$. Consequently, for every $k$, $\Delta_k\le(4+1.36\,C_5^{\#}K_5^{-1}K)\vartheta^{K-k_0}\vartheta^{-(K-k)}$. With $j=K-k\ge1$, $x_k=L^j\mathsf b$ and $\mathfrak r_{\min(k,k_d)}=\mathfrak r(L^{\max(j,j^*)}\mathsf b)$, the sum is at most $\mathrm{Ser}_2$ times the bracket times $\vartheta^{K-k_0}$.

For \textup{(D5)}, \Cref{Y-H6-J}\,(J7), (J8) for either run gives \[|\partial_1\cdots\partial_n\mathsf V'_{\ell_k}[u_k](0)|=|\sum_Y\partial_1\cdots\partial_nv_k(Y;u_k(z))(0)|\le\mathfrak r_{\min(k,k_d)}^{-n}N_1(k)\,2C_gc_a(x_k)e^{-\kappa(m_k-1)},\] and the right side is the term $j=K-k$ of $\mathrm{Tail}$. For the run $K$, $0\le k<k_0$ means $K\ge j>K-k_0$; for the run $K'$, $0\le k'<m+k_0$ means $K'\ge K'-k'>K-k_0$. The level $0$ is covered by \textup{(J8)}.

Let $j>j^*$. By \textup{(t3)}, Step~1\,(f) and \textup{(t5)}, $\mathfrak r(L^j\mathsf b)/\mathfrak r(L^{j+1}\mathsf b)\le L^{1/2}$, $N_1^{[j+1]}\le L^3N_1^{[j]}$, and $c_a(L^{j+1}\mathsf b)\le c_a(L^j\mathsf b)$. Moreover $m^{[j+1]}-m^{[j]}\ge n(\mathrm{sep}^{[j]}-\mathrm{sep}^{[j-1]})/2-1\ge n(L-1)\mathrm{sep}_*/2-1$ by Step~1\,(e), and $\mathrm{sep}_*\ge4\log L$. The quotient is therefore at most $\vartheta^{-1}L^{n/2+3}e^{\kappa}L^{-2\kappa n(L-1)}\le1/2$. Both series converge; $\mathrm{Ser}_2\le\sum_{j\le j^*}(\text{its terms})+2\,(\text{its term }j^*+1)$, and $\mathrm{Tail}(J)\le2\,(\text{its term }J+1)$ for $J\ge j^*$. In the latter term $m^{[J+1]}\ge n\,\mathrm{sep}^{[J]}/2\ge4nL^{J+1-j^*}\log(4nL^{1/2})$.

\emph{Step 10: \textup{(D7)}; on what the numbers depend, and local boundedness in $\mathsf b$.} The dependence follows from the definitions in item~1.3. In $c_Q^{\rm su}=306M^3\mathsf g(\lambda_\flat)^2$ only $M$ enters, besides a pure number. The constants $K_5$, $C_5$ and $C_1$ are pure numbers, by \textup{(O0a)} in item~(1) of \S\ref{Y-H0-1} and by \textup{(O0)} in item~1 of \S\ref{Y-J-1}; $C_5^{\#}$, $C_g$ and $\varepsilon_L$ involve $L$; and $c_4=15c_Q^{\rm su}C_1^2/E_a$. The functions $\mathrm{amp}$ and $c_a$ contain $p_0(x)=A_0(\log x)^{\mathfrak p_0}$ (\textup{(O8)} in item~1 of \S\ref{Y-J-1}), and $\check r(x)=r(x)+3$, with $r(x)=\lceil20c_R\log x/M\rceil$ (item~2 of \S\ref{Y-J-1}), contains $c_R$ and $M$. In the remainder of this step $p_0$ enters only through being continuous and $>0$ on $[x_s,\infty)$, and through \textup{(t3)}. Finally $\mathfrak r_*=\mathfrak r_{k_d}=\mathfrak r(L^{j^*}\mathsf b)$, because $K>j^*$.

We turn to monotonicity in $\mathsf b$. As $\check r(L^i\mathsf b)$ does not decrease in $\mathsf b$, $\mathrm{sep}^{[i]}$ does not increase and $j^*$ does not decrease; on $[\mathsf b_1,\mathsf b_2]$, therefore, $j^*(\mathsf b)\le j^*(\mathsf b_2)<\infty$ (Step~1\,(e)). The numbers $\Lambda$ and $K_0$ are continuous and increasing in $\mathsf b$. The left side of \textup{(K1a)} decreases in $\Lambda$ and its right side increases in $j^*$, and the left side of \textup{(K1b)} increases in $\Lambda$; hence $K_T(\mathsf b)$ is at most the $K_T$ obtained with $\Lambda(\mathsf b_2)$ and $j^*(\mathsf b_2)$. For every $j$ the function $\mathsf b\mapsto\mathfrak r(L^j\mathsf b)$ is $>0$ and continuous; since $j^*$ has only finitely many values on $[\mathsf b_1,\mathsf b_2]$, $\mathfrak r_*$ is bounded below there. The function $C_{\mathsf b}$ is continuous as well.

On $[\mathsf b_1,\mathsf b_2]$, each term of the series is bounded by the term defined as follows: $\mathrm{amp}$ and $c_a$ at their largest values, which exist since they are continuous functions on a compact set (in $\mathfrak r(L^{\max(j,j^*)}\mathsf b)$ we put $j^*(\mathsf b_2)$ for $j^*(\mathsf b)$, which is possible because $\mathrm{amp}(L^{i+1}\mathsf b)\ge\mathrm{amp}(L^i\mathsf b)$ by \textup{(t3)}); $N_1^{[j]}$ at $\mathsf b_2$; and $e^{-\kappa(m^{[j]}-1)}$ replaced by $1$ for $j\le j^*(\mathsf b_2)$ and by its value at $\mathsf b_2$ for $j>j^*(\mathsf b_2)$, since $m^{[j]}$ does not increase in $\mathsf b$ there. This majorant converges by Step~9.

Concerning the numerals: Steps~3 to~7 rest on $L^{1/4}\ge10^4$ and on \textup{(W2)}, \textup{(W3)}, \textup{(Y.0)(0b)}.\end{proof}

\subsection{Normalization at the last scale}\label{Y-H8-1}

We compare the densities $\rho[u_2;\ell]$ and $\rho[u_3;\ell]$, whose members may be complex, with one real positive first member: we bound their difference in the supremum and in $L^1$, and, when the condition \textup{(S)} holds, we bound quotients and logarithms. Along a pointwise holomorphic family $u(z)$ in the sources the integral $I[u(z)]$ is holomorphic, and Cauchy bounds control $\partial_S\log I$ and, in $L^1$, the jets of the laws; only the multi-affine jets count. For the sourced runs from two cutoffs the datum \textup{(J)} implies a formula for $S^T_{n;K}$, namely $\partial_1\cdots\partial_n\mathfrak V(0)$ plus a term of modulus at most $\log2\;\tau_7^n/\prod r_i$; a bound for what remains of the difference $S^T_{n;K'}-S^T_{n;K}$ when $\partial_1\cdots\partial_n\mathfrak V'(0)-\partial_1\cdots\partial_n\mathfrak V(0)$ is subtracted from it; and bounds in $L^1$ for the jets $\pi^{[K]}_{S,[s]}$ of the laws that the drawn fields follow.

\subsubsection*{Item 0: what the section does}

\begin{remark}[What passes from two real members to complex ones] \Cref{Y-J-ThB} concerns two points of one class, where both systems of members are real and both densities are positive. For the jets in the sources the further members are complex. \Cref{Y-J-ThB}\,(a) and (b) carry over to this case, (b) with an exponent that depends on the sizes. The first inequality of \Cref{Y-J-ThB}\,(c) depends on the first member alone and holds here unchanged. For a complex second member the argument for the second inequality of \Cref{Y-J-ThB}\,(c), an $L^1$ bound with $C_{\rm vol}$, does not carry over, since Step~9\,(6) of the proof of \Cref{Y-J-ThB} uses $\|g\|_1=\int g$, which requires $g\ge0$; instead, when the condition \textup{(S)} holds, \textup{(c2)} gives the bound with $2C_{\rm vol}$ in the place of $C_{\rm vol}$. \end{remark}

\subsubsection*{Item 1: letters, signs and definitions}

We use the letters of \S\ref{Y-H0-1} and of items 1 and 2 of \S\ref{Y-J-1}. The chain is run at the rate $\kappa_1=\kappa+1$, with $\kappa\ge6$: the sets $M_A^{\rm su}(x)$, the strong cap, the class $C^{\rm su}_\flat(x;\ell)$, the norm $\|\cdot\|_x$ and the distances $\mathsf D^{\rm in}_x$ and $\mathsf D_x$ are all defined with $\kappa_1$. The number of sources is $n$, and $q_N=(N/M)^3$ is the number of cubes.

For a subset $S$ of $\{1,\dots,n\}$ we write $z^S=\prod_{i\in S}z_i$, $\partial_S=\prod_{i\in S}\partial/\partial z_i$ and $r^S=\prod_{i\in S}r_i$. We write $\mathrm{Pd}(r)=\{z\in\mathbb C^n:|z_i|<r_i\}$, and $\mathrm{Pd}(r/t)$ for the polydisc with the radii $r_i/t$.

The letter $s$ denotes the depth when it has no argument and the radius $s(x)=8r(x)$ of item~2 of \S\ref{Y-J-1} when it has one, and both senses occur in $2M\,s(L^s\mathsf b)$, whereas in $x_s$ the index is a mark and does not denote the depth $s$. The symbol $z(ax)$ refers not to the sources $z$ but to the normalisation of the kernel $k_{ax}$, as defined in item~2 of \S\ref{Y-J-1}. The rate $\kappa_1$ is always $\kappa+1$, but $\kappa_K$, $\kappa_{K-s}$ and $\kappa_k$, whose indices are letters that denote levels, are the rates of \textup{(o4)} in item~2 of \S\ref{Y-H6-1}, and for $K=1$ we have $\kappa_K=\kappa$, not $\kappa+1$. The caps are $c_h(x)$ and $c_a(x)$, while in the phrase ``with the cap $c_h$'' and in $M_1(\mathsf b;\kappa_1,c_h)$ the letter $c_h$, without an argument, is the number of \textup{(O11)} in item~1 of \S\ref{Y-J-1}, and such a phrase says that the cap is $c_h(x)=e^{-c_hp_0(x)^2}$ itself.

\begin{definition}[Triple without a size]\label{Y-H8-Triple} A \emph{triple without a size} at $x$ is a triple $u=(E^a,E^h,w)$ that satisfies all clauses of \textup{(ca)}, \textup{(ch)} and \textup{(c$\Xi$)} of item~2 of \S\ref{Y-J-1} except the inequality. These triples form a complex vector space. The proofs use of such a triple only that $1_{X^\approx}(V)E^a(X,V)$, $1_{X^\approx}(V)E^h(X,V)$ and $w(X,V)$ are complex Borel functions of the $G$-valued field $V$. \end{definition}

\begin{definition}[Rate and norms at a rate]\label{Y-H8-Norms} A \emph{rate} is a number $\tilde\kappa\in[\kappa,\kappa_1]$, and the cap of $E^h$ is $c_h(x)$ itself. The norms below take their values in $[0,\infty]$. We let $\|u\|^\sim_E$ be the least $A$ such that $\sup_{\mathcal U_X(x)}|E^a(X,\cdot)|\le A\,c_a(x)e^{-\tilde\kappa(\#X-1)}$ and $\sup|E^h(X,\cdot)|\le A\,c_h(x)e^{-\tilde\kappa(\#X-1)}$ for all $X\in\mathcal D$. We let $\|u\|^\sim_w$ be the least $A$ such that $|w(X,V)|\le A\,e^{-\tilde\kappa(|X|_c-1)}$ for all $X\in\mathcal P$ and all $V$, that is, at the weak cap. We let $\|u\|^\sim$ be the larger of the two. Since $\tilde\kappa\le\kappa_1$, we have $\|u\|^\sim\le\|u\|_x$; hence $\|u_1-u_2\|^\sim\le\mathsf D^{\rm in}_x(u_1,u_2)\le\mathsf D_x(u_1,u_2)$, and $\|u\|^\sim\le1$ on $M_1^{\rm su}(x)$. If a norm measures $w$ against the strong cap at $\tilde\kappa$, it dominates $\|\cdot\|^\sim$. An inequality whose right side is infinite says nothing. \end{definition}

\begin{definition}[The last integral and the sum of the members]\label{Y-H8-Int} For $\|u\|^\sim$ finite we put $I[u]:=\int\rho[u;\ell](V)\,dV$. For the member $u_i$ we put $\mathrm{Sm}_i(V):=\sum_{X\in\mathcal D}1_{X^\approx}(V)\,[E_i^a+E_i^h](X,V)$. \end{definition}

\begin{definition}[The volume constant]\label{Y-H8-Cvol} We put \[C_{\rm vol}=C_{\rm vol}(N,x):=C^{\rm su}_{\rm vol}(N,x)=4q_Ne^{220q_N}e^{(3/2)\mathsf N_Zp_0(x)^2N^3}(384/\varepsilon(x)^3)^{3N^3},\] which is \textup{(O14)} in item~1 of \S\ref{Y-J-1}. If $\varepsilon(x)\le1$, then $C_{\rm vol}\ge4q_Ne^{220q_N}\ge400q_N$. At fixed $N$, $C_{\rm vol}$ is a continuous function of $x$ on $[e^2,\infty)$. \end{definition}

\begin{definition}[Pointwise holomorphic family]\label{Y-H8-Family} Let $\Omega\subset\mathbb C^n$ be open. A \emph{pointwise holomorphic family} on $\Omega$, in the sense of \Cref{Y-H1-U1}\,(c+), assigns to every $z\in\Omega$ a triple without a size $u(z)$ such that every value $E^a(X,\mathsf W;z)$ with $\mathsf W\in\mathcal U_X(x)$, every value $E^h(X,V;z)$ and every value $w(X,V;z)$ is holomorphic in $z$. \end{definition}

\begin{definition}[Congruence modulo the squares of the sources]\label{Y-H8-Congruence} For germs $f$ and $g$ of holomorphic functions at $0\in\mathbb C^n$ we write $f\equiv g$ if $\partial_Sf(0)=\partial_Sg(0)$ for all $S$, the empty set included, and we call $f$ and $g$ congruent modulo $z_i^2$. Two families are congruent if all their values are. \end{definition}



\subsubsection*{Item 2: the theorem}

\begin{theorem}[Two complex members against one positive member]\label{Y-H8-ThBC} Let the following data \textup{(B)} be given: $x\ge e^2$; $N$ even, $M\mid N$ and $N\ge2M\,s(x)$; $\varepsilon(x)\le1$; $(c_a+c_h)(x)\le1/2$; rows $\ell$ on $\mathbb Z^3/N\mathbb Z^3$ of range $\le\varrho(x)$; and $\tilde\kappa\in[\kappa,\kappa_1]$. As the first member, let $\rho_1=e^{\mathsf V_1}\rho[u_1;\ell]$ be a point of $C^{\rm su}_\flat(x;\ell)$ and $u_1$ a system of members of it; thus $u_1$ is real, lies in $M_1^{\rm su}(x)$ and has the strong cap, $\mathsf V_1$ is real, and $\rho_1>0$. As further members, let $u_2,u_3$ be triples without a size and $\mathsf V_2,\mathsf V_3\in\mathbb C$, and put $\rho_i:=e^{\mathsf V_i}\rho[u_i;\ell]$. Put $\mathsf D_{ij}:=\|u_i-u_j\|^\sim$. Whenever $\|u_i\|^\sim<\infty$, write $I_i:=I[u_i]$. Then the following hold.

\begin{enumerate} \item[\textup{(a$^C$)}] At every $V$ for which $D(V)$ is empty, $\rho_i(V)=\exp\{\mathsf V_i-x\mathrm{En}_\ell(V)+\mathrm{Sm}_i(V)\}$, which is not $0$, and $|\mathrm{Sm}_i-\mathrm{Sm}_j|(V)\le1.22\,q_N(c_a+c_h)(x)\,\mathsf D_{ij}$.

\item[\textup{(b$^C$)}] Let $a_E\ge\|u_2\|^\sim_E,\|u_3\|^\sim_E$ and $a_w\ge\max\{1,\|u_2\|^\sim_w,\|u_3\|^\sim_w\}$. Then \[\sup_V|e^{-\mathsf V_2}\rho_2-e^{-\mathsf V_3}\rho_3|\le1.61\,q_Ne^{(0.61a_E+218a_w)q_N}\mathsf D_{23}.\] In three cases the bound reads as follows.

\begin{enumerate} \item[\textup{(i)}] If $400q_N\max\{\mathsf D_{12},\mathsf D_{13}\}\le1$, the bound is $\le2q_Ne^{219q_N}\mathsf D_{23}$; for $u_3=u_1$ it is the bound of \Cref{Y-J-ThB}\,(b).

\item[\textup{(ii)}] If $a_w=1$ and $a_E\le2$, the bound is $\le1.61\,q_Ne^{219.22q_N}\mathsf D_{23}$. This case covers every output $F_\ell(u)$ of \Cref{Y-H-C9} with $A\le6$ and with $p=(\tilde\kappa,c_h)$ as the pair of parameters, because \Cref{Y-H-C9}\,(b), which holds for this $p$ by \Cref{Y-H1-U1}\,(U), bounds the $E^a$ of such an output by $0.35+0.248\cdot6=1.838$ times the cap $c_a$ and its $E^h$ by the cap $c_h$, and gives its $w$ the strong cap; hence $a_E\le2$ and $a_w=1$.

\item[\textup{(iii)}] If $\|u_2\|^\sim,\|u_3\|^\sim\le A$ with $A\ge1$, the bound is $\le2q_Ne^{219Aq_N}\mathsf D_{23}$. \end{enumerate}

\item[\textup{(c$^C$)}] Let moreover the rows obey $|\ell_y(p)|\le K_*e^{-\lambda_*|y-p|}$. Then \[I_1\ge e^{-0.61q_N}e^{-(3/2)\mathsf N_Zp_0^2N^3}(\varepsilon^3/384)^{3N^3};\] this is the first inequality of \Cref{Y-J-ThB}\,(c), and it depends on the first member alone. In the cases \textup{(i)} and \textup{(ii)}, $\int|\rho[u_2;\ell]-\rho[u_3;\ell]|\,dV\le\frac12C_{\rm vol}\mathsf D_{23}I_1$; in the case \textup{(iii)} the same holds with $C_{\rm vol}e^{219(A-1)q_N}$ in the place of $C_{\rm vol}$. If, in addition, $C_{\rm vol}\max\{\mathsf D_{12},\mathsf D_{13}\}\le1$, which is the condition \textup{(S)} and implies the case \textup{(i)}, then we have the following.

\begin{enumerate} \item[\textup{(c1)}] We have $|I_2-I_3|\le\frac12C_{\rm vol}\mathsf D_{23}I_1$ and $|I_i/I_1-1|\le\frac12C_{\rm vol}\mathsf D_{1i}\le\frac12$; moreover $I_i\ne0$ and $\int|\rho[u_i;\ell]|\,dV\le\frac32I_1\le3|I_i|$.

\item[\textup{(c2)}] We have \[\int|\rho_1/\!\int\!\rho_1-\rho_2/\!\int\!\rho_2|\,dV\le C_{\rm vol}\mathsf D_{12}/(1-C_{\rm vol}\mathsf D_{12}/2)\le2C_{\rm vol}\mathsf D_{12}.\]

\item[\textup{(c3)}] We have $\int|\rho_2/\!\int\!\rho_2-\rho_3/\!\int\!\rho_3|\,dV\le4C_{\rm vol}\mathsf D_{23}$.

\item[\textup{(c4)}] We have $|\mathrm{Log}(I_2/I_1)-\mathrm{Log}(I_3/I_1)|\le C_{\rm vol}\mathsf D_{23}$, where $\mathrm{Log}$ is the principal branch on $|w-1|\le1/2$. \end{enumerate}

\item[\textup{(d$^C$)}] Let $\eta_d:=1.22\,q_N(c_a+c_h)\mathsf D_{12}<\log2$. Since $c_a+c_h>0$ and $\|u_1\|^\sim\le1$, this implies $\|u_2\|^\sim\le1+\mathsf D_{12}<\infty$. Let $I_2\ne0$, let $\mathrm{Tm}:=\{D(V)\text{ empty}\}$, let $\mathrm{ms}_1$ be the mass of the complement of $\mathrm{Tm}$ under $\rho_1/\!\int\!\rho_1$, and let $\mathrm{ms}_2:=\int_{{\rm off}\,\mathrm{Tm}}|\rho_2|\,dV/|\!\int\!\rho_2|$. Then \[\int|\rho_1/\!\int\!\rho_1-\rho_2/\!\int\!\rho_2|\,dV\le2\,\mathrm{ms}_1+\mathrm{ms}_2\bigl(1+e^{\eta_d}/(2-e^{\eta_d})\bigr)+2(e^{\eta_d}-1)/(2-e^{\eta_d}).\] \end{enumerate} \end{theorem}

\subsubsection*{Item 3: the first corollary}

\begin{corollary}[Families in the sources and their jets]\label{Y-H8-CoF} Let the data \textup{(B)}, the bound on the rows of \textup{(c$^C$)} and the datum \textup{(F0)} be given; then \textup{(F1)} to \textup{(F5)} hold.

\begin{enumerate} \item[\textup{(F0)}] The datum is a pointwise holomorphic family $u(z)$ on $\mathrm{Pd}(r)$ with $u(0)=u_1$, the first member, together with a finite number $B_0\ge\sup_{\mathrm{Pd}(r)}\|u(z)-u_1\|^\sim$. We put $\tau_7:=\max\{1,B_0C_{\rm vol}\}$.

\item[\textup{(F1)}] The map $z\mapsto I[u(z)]$ is holomorphic on $\mathrm{Pd}(r)$, with $\partial_SI[u(z)]=\int\partial_S\rho[u(z);\ell](V)\,dV$, and $\|u(z)-u_1\|^\sim\le B_0\max_i|z_i|/r_i$.

\item[\textup{(F2)}] On $\mathrm{Pd}(r/\tau_7)$ we have $|I[u(z)]/I_1-1|\le1/2$, and for every non-empty $S$ we have $|\partial_S\log I[u(z)]|_{z=0}\le\log2\;\tau_7^{|S|}/r^S$.

\item[\textup{(F3)}] Let $\mathfrak n(z)$ be a pointwise holomorphic family on $\mathrm{Pd}(r/2\tau_7)$, with $\mathfrak n(0)$ arbitrary, let $\Delta:=\sup_{\mathrm{Pd}(r/2\tau_7)}\|u(z)-\mathfrak n(z)\|^\sim$, and let $S$ be non-empty; if $C_{\rm vol}\Delta\le1/2$, then $I[\mathfrak n(0)]\ne0$ and \[|\partial_S\log I[u(z)]-\partial_S\log I[\mathfrak n(z)]|_{z=0}\le C_{\rm vol}\Delta\,(2\tau_7)^{|S|}/r^S.\] If $\mathfrak n(z)$ is a pointwise holomorphic family on all of $\mathrm{Pd}(r)$, if $\mathfrak n(0)$ is a system of members of a point of $C^{\rm su}_\flat(x;\ell)$, hence real, in $M_1^{\rm su}(x)$, with the strong cap and with $\rho[\mathfrak n(0);\ell]>0$, and if $\sup_{\mathrm{Pd}(r)}\|\mathfrak n(z)-\mathfrak n(0)\|^\sim\le B_0$, then the left side is $\le3C_{\rm vol}\Delta\,(2\tau_7)^{|S|}/r^S$ for every value of $\Delta$.

\item[\textup{(F4)}] \emph{A lemma on jets.}

\begin{enumerate} \item[\textup{(i)}] If $\delta_S$, for non-empty $S$, are triples without a size with $\sum_Sr^S\|\delta_S\|^\sim\le B_0$, then $u^{\rm pol}(z):=u_1+\sum_Sz^S\delta_S$ obeys \textup{(F0)}.

\item[\textup{(ii)}] If $u(z)$ and $\hat u(z)$ are congruent pointwise holomorphic families near $0$, then $\rho[u(z);\ell](V)\equiv\rho[\hat u(z);\ell](V)$ at every $V$. If, moreover, the norms $\|\cdot\|^\sim$ of both families stay bounded near $0$ and $I[u(0)]\ne0$, then $\partial_S\log I[u(z)]=\partial_S\log I[\hat u(z)]$ at $0$ for every non-empty $S$.

\item[\textup{(iii)}] Let $u(z)$ be a pointwise holomorphic family whose norm $\|u(z)\|^\sim$ stays bounded near $0$, with $I[u(0)]\ne0$, and let $\mathfrak V(z)$ be holomorphic near $0$. Let $F(z,V)$ be a function that is Borel in $V$, holomorphic in $z$ near $0$ for each fixed $V$, and bounded uniformly in $(z,V)$, such that $F(\cdot,V)\equiv e^{\mathfrak V(z)}\rho[u(z);\ell](V)$ holds for each $V$. Then $\partial_S\log\int F\,dV=\partial_S\mathfrak V+\partial_S\log I[u(z)]$ at $z=0$ for each non-empty $S$, and for each $V$ we have $F(z,V)/\!\int\!F\,dV\equiv\rho[u(z);\ell](V)/I[u(z)]$. \end{enumerate}

Hence, under the hypotheses of the second sentence of \textup{(ii)}, the congruence $\rho[u(z);\ell](V)/I[u(z)]\equiv\rho[\hat u(z);\ell](V)/I[\hat u(z)]$ holds at every $V$.

\item[\textup{(F5)}] For the laws $\mu_u^z:=\rho[u(z);\ell]/I[u(z)]$ and for non-empty $S$, the jets obey $\int|\partial_S\mu_u^z|_{z=0}\,dV\le3\tau_7^{|S|}/r^S$, and, with $\mathfrak n$ and $\Delta$ as in \textup{(F3)}, if $C_{\rm vol}\Delta\le1/2$, then \[\int|\partial_S\mu_u^z-\partial_S\mu_{\mathfrak n}^z|_{z=0}\,dV\le4C_{\rm vol}\Delta\,(2\tau_7)^{|S|}/r^S.\] Under the hypotheses of the second sentence of \textup{(F3)}, this bound holds for every value of $\Delta$. \end{enumerate} \end{corollary}

\subsubsection*{Item 4: the second corollary}

\begin{corollary}[Two sourced runs and the jets of their laws]\label{Y-H8-CoR7} \emph{Setting.} The letters are those of \S\ref{Y-H0-1}; $\mathsf b\ge x_s$; $N=2L^{m_l}$ with $m_l\ge m'$; the condition \textup{(Tor)} $N\ge2M\,s(L\mathsf b)$ holds; $K'\ge K\ge1$ and $m=K'-K$; and the runs, tori, grids, rows $\ell_k=\ell_k^{[K]}$ and numbers $x_k=L^{K-k}\mathsf b$ are those of the definition ``Runs'' in item~2 of \S\ref{Y-J-1}. Let $f_1,\dots,f_n$ be complex bounded functions on $T_l$, and let $\mathcal O_K(f)=a_K^{-1}\sum_pf(x(p))\,a_p(U)$, where $a_K=L^{-K}$, $x(p)$ is the barycentre of $p$ and $a_p=1-\cos\phi_p$. We put $\rho^{[K],z}_0:=\rho^{[K]}_0\exp\{\sum_iz_i\mathcal O_K(f_i)\}$ and $\rho^{[K],z}_{k+1}:=(\rho^{[K],z}_k)'$, where the prime denotes the unsourced step with the noise $a x_k$. We put $Z_K(z):=\int\rho^{[K],z}_0$ and $S^T_{n;K}=\partial_1\cdots\partial_n\log Z_K(0)$.

\emph{Datum.} We call the clauses \textup{(J.a)} to \textup{(J.c)} below the datum \textup{(J)} at the depth $s\ge0$; it is taken at the level $k=K-s\ge0$, at $x=x_k=L^s\mathsf b$, on the torus $\mathbb Z^3/NL^s\mathbb Z^3$ and with the rows $\ell_k$, and it is the hypothesis of each clause below that names it. At $k=0$ the rows $\ell_0$ are empty, and the data \textup{(B)}, the bound on the rows of \textup{(c$^C$)}, with $\ell_y=0$ for empty rows, and Step~9\,(1) to (5) of the proof of \Cref{Y-J-ThB} admit empty rows as well.

\begin{enumerate} \item[\textup{(J.a)}] $\bar u_k$ is a system of members of the point $\rho^{[K]}_k$ of $C^{\rm su}_\flat(x_k;\ell_k)$.

\item[\textup{(J.b)}] $u(z)$ is a pointwise holomorphic family on $\mathrm{Pd}(r)$ with $u(0)=\bar u_k$ and $\sup_{\mathrm{Pd}(r)}\|u(z)-\bar u_k\|^\sim\le B_0$, $\mathfrak V(z)$ is holomorphic near $0$, and $\rho^{[K],z}_k(V)\equiv e^{\mathfrak V(z)}\rho[u(z);\ell_k](V)$ at every $V$.

\item[\textup{(J.c)}] $\mathfrak n(z)$ is a pointwise holomorphic family on $\mathrm{Pd}(r/2\tau_7)$, bounded in $\|\cdot\|^\sim$, $\mathfrak V'(z)$ is holomorphic near $0$, and $\rho^{[K'],z}_{m+k}(V)\equiv e^{\mathfrak V'(z)}\rho[\mathfrak n(z);\ell_k](V)$ at every $V$, with the rows $\ell_k$ of the short run; $\Delta$ is as in \textup{(F3)}. \end{enumerate}

Then, with $C_{\rm vol}=C_{\rm vol}(NL^s,L^s\mathsf b)$ and $\tau_7=\max\{1,B_0C_{\rm vol}\}$, the following hold.

\begin{enumerate} \item[\textup{(R0)}] For every level $k$ and every $z\in\mathbb C^n$ we have $\int\rho^{[K],z}_k\,dV=Z_K(z)$.

\item[\textup{(R1)}] If \textup{(J.a)} and \textup{(J.b)} hold, then $S^T_{n;K}=\partial_1\cdots\partial_n\mathfrak V(0)+\partial_1\cdots\partial_n\log I[u(z)](0)$, and the last term has modulus at most $\log2\;\tau_7^n/\prod r_i$.

\item[\textup{(R2)}] If \textup{(J.a)} to \textup{(J.c)} hold and $C_{\rm vol}\Delta\le1/2$, then \[|S^T_{n;K'}-S^T_{n;K}-(\partial_1\cdots\partial_n\mathfrak V'(0)-\partial_1\cdots\partial_n\mathfrak V(0))|\le C_{\rm vol}\Delta\,(2\tau_7)^n/\prod r_i,\] and under the hypotheses of the second sentence of \textup{(F3)} the same left side is $\le3C_{\rm vol}\Delta\,(2\tau_7)^n/\prod r_i$ for every $\Delta$. For polynomials $u=\bar u_k+\sum z^S\delta_S$ and $\mathfrak n=\bar{\mathfrak n}_k+\sum z^S\tilde\delta_S$ we have \[\Delta\le\|\bar u_k-\bar{\mathfrak n}_k\|^\sim+\sum_S(r/2\tau_7)^S\|\delta_S-\tilde\delta_S\|^\sim.\]

\item[\textup{(R3)}] For non-empty $S$ and $0\le s\le K$ we put \[\pi^{[K]}_{S,[s]}:=\partial_S\bigl(\rho^{[K],z}_{K-s}/Z_K(z)\bigr)\big|_{z=0},\] written $\pi_S$ for short; it is a bounded Borel function on the fields of $\mathbb Z^3/NL^s\mathbb Z^3$.

\begin{enumerate} \item[\textup{(i)}] For every bounded Borel function $\mathrm{obs}$, the integral $\int\mathrm{obs}\,\pi^{[K]}_{S,[s]}\,dV$ is the joint truncated function of the $\mathcal O_K(f_i)$, $i\in S$, and of $\mathrm{obs}(V_{K-s})$ under the law of the run. Moreover $\int\pi^{[K]}_{S,[s]}\,dV=0$, and $\pi_S$ is gauge invariant, and real for real $f_i$.

\item[\textup{(ii)}] For $1\le s\le K$ the drawn step with the noise $aL^s\mathsf b$ maps $\pi^{[K]}_{S,[s]}$ to $\pi^{[K]}_{S,[s-1]}$; this step does not increase the $L^1$ norm of a complex function, and it does not depend on $K$.

\item[\textup{(iii)}] If \textup{(J)} holds, then $\int|\pi^{[K]}_{S,[s]}|\,dV\le3\tau_7^{|S|}/r^S$, and if moreover $C_{\rm vol}\Delta\le1/2$, then \[\int|\pi^{[K']}_{S,[s]}-\pi^{[K]}_{S,[s]}|\,dV\le4C_{\rm vol}\Delta\,(2\tau_7)^{|S|}/r^S.\] Under the hypotheses of the second sentence of \textup{(F3)}, that is, if $\mathfrak n$ is a family on all of $\mathrm{Pd}(r)$, $\mathfrak n(0)$ is a system of members of a point of the class and $\sup_{\mathrm{Pd}(r)}\|\mathfrak n(z)-\mathfrak n(0)\|^\sim\le B_0$, the same bound holds for every value of $\Delta$; the condition $C_{\rm vol}(NL^s,L^s\mathsf b)\Delta\le1/2$ depends on $N$ and $s$. With the datum of item~5, $K$ has to satisfy $K\ge K_T(\mathsf b,n,d)$ and $K-s\ge K_0(L^s\mathsf b)$, and these conditions depend on $\mathsf b,n,d,s$ and on no $N$.

\item[\textup{(iv)}] Let the datum hold at the depths $s$ and $s-1$ for all $K'\ge K\ge K_\#$, with $r$ and $B_0$ free of $K,K'$ and with $\Delta\le\check\Delta_K$ for some numbers $\check\Delta_K\to0$; with the datum of item~5, $\check\Delta_K$ is the bound of \Cref{Y-H7-Df}\,(D3) at the depth $s$. Then, as $K\to\infty$, each of the sequences $\pi^{[K]}_{S,[s]}$ and $\pi^{[K]}_{S,[s-1]}$ converges in $L^1(dV)$, without passing to a subsequence. The $L^1(dV)$ distance of $\pi^{[K]}_{S,[s]}$ from its limit satisfies the second bound of \textup{(iii)} with that limit in the place of $\pi^{[K']}_{S,[s]}$ and with $\check\Delta_K$ in the place of $\Delta$, for the $K\ge K_\#$ with $C_{\rm vol}\check\Delta_K\le1/2$, and for all $K\ge K_\#$ under the hypotheses of the second sentence of \textup{(F3)}. At the depth $s-1$ the same is true of $\pi^{[K]}_{S,[s-1]}$, with the constants of that depth. The limits have the properties of \textup{(i)}, namely integral $0$, invariance and reality, and they obey the first bound of \textup{(iii)}; the step of \textup{(ii)} maps the limit at $s$ to the limit at $s-1$. Since their integral is $0$, none of them is $\ge0$ unless it vanishes almost everywhere.

\item[\textup{(v)}] Let $K-s\ge1$, and let $\hat f_i(y):=f_i(L^{-s}y)$ on $T_{NL^s}$. Then $\pi^{[K]}_{S,[s]}$, defined with $(K,\mathsf b,N,f)$, equals $L^{s|S|}$ times $\pi^{[K-s]}_{S,[0]}$, defined with $(K-s,L^s\mathsf b,NL^s,\hat f)$. Hence, for the run in which $(K-s,L^s\mathsf b,NL^s,\hat f)$ replaces $(K,\mathsf b,N,f)$, the datum \textup{(J)} may be taken at the depth $0$. The triple $(L^s\mathsf b,NL^s,\hat f)$ meets the conditions of the ``Setting'' of this corollary. If, moreover, the $f_i$ are in the setting of item~3 of \S\ref{Y-H7-1}, then any two of the sets $\{\hat f_i\ne0\}$ have sup distance $\ge L^sd$, each of them lies within sup distance $L^sR_f$ of a point, $\|\hat f_i\|_\infty=\|f_i\|_\infty$, $\mathrm{Lip}\,\hat f_i=L^{-s}\mathrm{Lip}\,f_i$ and $\|\hat f_i\|_\sharp=\max\{\|f_i\|_\infty,40M\,L^{-s}\mathrm{Lip}\,f_i\}\le1$, so that the $\hat f_i$ lie in that setting as well, with $(L^s\mathsf b,NL^s,L^sd,L^sR_f)$ instead of $(\mathsf b,N,d,R_f)$. In that case the datum \textup{(J)} holds at the depth $0$ for this run whenever $K'-s\ge K-s\ge K_T(L^s\mathsf b,n,L^sd)$, as follows from the parts ``First'' and ``Second'' of item~5 when $(K-s,K'-s,L^s\mathsf b,NL^s,\hat f_i,L^sd,L^sR_f)$ replaces $(K,K',\mathsf b,N,f_i,d,R_f)$ in them; there $r_i=\mathfrak r_*/2$ and $B_0=1$, by \Cref{Y-H7-Df}\,(D1)\,(f), and $\Delta$ is at most the bound of \Cref{Y-H7-Df}\,(D3). Alternatively the datum is taken at the depth $s$ itself, as in the part ``Datum at the depth $s\ge0$'' of item~5, for $K\ge K_T(\mathsf b,n,d)$ and $K-s\ge K_0(L^s\mathsf b)$; this threshold is the lower one, since $\max\{K_T(\mathsf b,n,d),\,s+K_0(L^s\mathsf b)\}\le s+K_T(L^s\mathsf b,n,L^sd)$ by the part ``The two thresholds'' of Step~P10\,(v). \end{enumerate} \end{enumerate} \end{corollary}

\subsubsection*{Item 5: the datum}

\begin{remark}[Which statements supply the datum]\label{Y-H8-Datum} The setting of item~4 admits any $n$ and any complex bounded $f_i$. The clauses supplying the datum are stated in narrower settings, and the parts ``First'' and ``Second'' below assume these settings. \Cref{Y-H6-J}\,(J6), (J7) hold in the setting of item~1 of \S\ref{Y-H6-1}: $n\ge2$; $d,R_f>0$; the $f_i$ are real and bounded with $\|f_i\|_\infty\le1$; any two of the sets $\{f_i\ne0\}$ have sup distance $\ge d$, and each of them lies within sup distance $R_f$ of a point; and $K\ge1$. The clauses of \Cref{Y-H7-Df} are formulated in the setting of item~3 of \S\ref{Y-H7-1}, which asks besides that the $f_i$ be Lipschitz with $\|f_i\|_\sharp\le1$ and that $K'\ge K\ge K_T(\mathsf b,n,d)$. Since $S^T_n$ and $\pi_S$ are linear in each $f_i$, the conditions $\|f_i\|_\infty\le1$ and $\|f_i\|_\sharp\le1$ are no restriction: the general case follows by scaling each $f_i$.

\emph{First.} We supply \textup{(F0)}, \textup{(J.a)} and \textup{(J.b)} at $s=0$. The last sentences of \Cref{Y-H4-U4}\,(h) give \textup{(J.a)} for every $K\ge1$. \Cref{Y-H6-J}\,(J6), (J7) give the family $u_K(z)$ on $\mathrm{Pd}(\mathfrak r_*)$, with $r_i=\mathfrak r_*$, $B_0=0.91$ and $\tilde\kappa=\kappa_K=\kappa$. This family is an output of \Cref{Y-H-C9} at the rate $\kappa<\kappa_1$, taken for a pair $p$ for which the clauses of that theorem hold by \Cref{Y-H1-U1}\,(U), and it falls under the case \textup{(ii)} of \Cref{Y-H8-ThBC}\,(b$^C$); on $\mathrm{Pd}(r/\tau_7)$, where \textup{(F2)} applies, the condition \textup{(S)} holds and implies the case \textup{(i)}, which suffices there. The family is pointwise holomorphic by \Cref{Y-H1-U1}\,(c+), and it has the true multi-affine jets. Then $\tau_7=\max\{1,0.91\,C_{\rm vol}(N,\mathsf b)\}$.

\emph{Second.} We supply the hypotheses of the second sentence of \textup{(F3)}, and \textup{(J.c)}, at $s=0$. \Cref{Y-H7-Df}\,(D1), (D3) give $\mathfrak n_K(z)$ on $\mathrm{Pd}(\mathfrak r_*/2)$, with $r_i=\mathfrak r_*/2$, with $B_0=1$, and with $\Delta$ at most the bound of (D3). That $\mathfrak n_K(0)$ is a system of members of a point of the class follows from \Cref{Y-J-ThA}\,(iii) together with \Cref{Y-H7-Df}\,(D1)\,(c). The value $B_0=1$ comes from \Cref{Y-H7-Df}\,(D1)\,(f).

\emph{Datum at the depth $s\ge0$.} We work in the setting of item~3 of \S\ref{Y-H7-1}, that is, with $n\ge2$, with real Lipschitz $f_i$ such that $\|f_i\|_\sharp\le1$, and with any two of the sets $\{f_i\ne0\}$ at sup distance $\ge d$, each within sup distance $R_f$ of a point; we take $K\ge K_T(\mathsf b,n,d)$ and $k=K-s>k_0$, that is, $K-s\ge K_0(L^s\mathsf b)$, and $\tilde\kappa=\kappa_{K-s}$. The clauses of the datum are then supplied as follows.

\begin{enumerate} \item[\textup{(J.a)}] \Cref{Y-H4-U4}\,(h) gives this clause.

\item[\textup{(J.b)}] We take $u=u_k^{\rm pol}$ if $s\ge j^*$ and $u=u_k(z)$ if $s<j^*$, on $\mathrm{Pd}(r)$ with $r_i=\mathfrak r(L^{\max(s,j^*)}\mathsf b)/2$, and $B_0=1$, since by \Cref{Y-H7-Df}\,(D1)\,(e), (f) the family $u(z)$ lies within $0.455$ of $\bar u_k$ on $\mathrm{Pd}(r)$. The congruence follows from \Cref{Y-H4-U4}\,(h), which holds at every field, from the definition of the jets in \textup{(o2)} of item~2 of \S\ref{Y-H6-1}, from \Cref{Y-H6-J}\,(J6) and from \textup{(F4)\,(ii)}.

\item[\textup{(J.c)}] For this clause and for the hypotheses of the second sentence of \textup{(F3)} we take $\mathfrak n=\mathfrak n_k^{\rm pol}$ if $s\ge j^*$ and $\mathfrak n=\mathfrak n_k(z)$ if $s<j^*$, on the same $\mathrm{Pd}(r)$; it lies within $0.795$ of $\bar{\mathfrak n}_k$, a system of members of a point of the class, by \Cref{Y-H7-Df}\,(D1)\,(c), (e), (f) and \Cref{Y-J-ThA}\,(iii). The congruence follows from \Cref{Y-H7-Df}\,(D0), (D1)\,(c) and \textup{(F4)\,(ii)}, as in Step~P11, and $\Delta$ is at most the bound of \Cref{Y-H7-Df}\,(D3) at the depth $s$. \end{enumerate}

Here $r$ and $B_0$ depend on no $K$, $K'$, $N$, and $\tau_7=C_{\rm vol}(NL^s,L^s\mathsf b)$. \end{remark}

\subsubsection*{Item 6: the place of this section in the order of the letters}

\begin{remark}[The letters and the floors of this section]\label{Y-H8-Order} In this section we impose no further condition on the letters \textup{(O0)} to \textup{(O12)}, and the floors in force here are those of \textup{(O13)}. Indeed, at $x=L^s\mathsf b$, where $\mathsf b\ge x_s$, the data \textup{(B)} follow from $x_s\ge e^2$, \textup{(X16)}, \textup{(X17)} and Step~2\,\textup{(c)}, \textup{(e)} of the proof of \Cref{Y-J-ThA}. Here we use that, with the rate $\kappa_1$, the number $x_s$ satisfies all of \textup{(O13)} of item~1 of \S\ref{Y-J-1} (\textup{(O13)} in item~\textup{(1)} of \S\ref{Y-H0-1}), so that the inequality $\mathsf b\ge x_s$ supplies what the floors \textup{(X16)}, \textup{(X17)}, \textup{(X12)} and the arguments for these two items of Step~2 require of $x_\flat$.

Step~2 belongs to the proof of \Cref{Y-J-ThA}, among whose data is the condition $K\ge K_0(\mathsf b)$, while the setting of item~4 of this section admits every $K\ge1$. The arguments for the items \textup{(c)} and \textup{(e)} of that step are free of $K_0$:

\begin{enumerate}

\item[\textup{(c)}] The argument for \textup{(c)} rests on $\mathsf b\ge x_\flat$, on \textup{(X12)}, on the geometric series over $1\le j\le K$, and on \Cref{Y-A-Th29}\,\textup{(b)}, applied under its hypotheses $L\ge L_A$ odd, $aL\ge3/c_L$ and $N$ even; this clause yields $\ell_k\in\mathcal R(b_*,K_*,\lambda_*,\frac19)$ for $0\le k\le K$, for every tower and all radii subject to \textup{(Cut)}, and its statement is free of $K_0$.

\item[\textup{(e)}] The argument for \textup{(e)} rests on $N\ge2M\,s(L\mathsf b)\ge2M\,s(\mathsf b)$ and on a clause of \textup{(RB)} in item~3 of \S\ref{Y-J-1}, namely \Cref{Y-GB-P63}: if $x\ge e^2$, $L\ge3$ and $N\ge2M\,s(x)$, then $NL^j\ge2M\,s(L^jx)$ for all $j\ge0$. We apply this clause with $x=\mathsf b\ge e^2$.

\end{enumerate}
Hence both items hold for every $K\ge1$.

The item \textup{(a)} of Step~2 enters Step~P9\,\textup{(R2)} for all $K'\ge K\ge1$. It asserts that the level $m+k$ of the run of length $K'$ has the same torus, the same value $x_k$, the same grid, the same numbers listed in item \textup{(a)} and the same step as the level $k$ of the run of length $K$. The reasons are that $K'-(m+k)=K-k$ is the exponent of $L$ in $x_k$ and in the side of the torus, that the origins of the grids are fixed backwards from the last level of a run, that each of these numbers is a function of $x_k$ and of the letters, and that the step depends on $x_k$, $a$ and $L$ and is independent of the rows, of the grid and of the members. Its argument, too, is free of $K_0$.

The bound $|\ell_{k,y}(p)|\le K_*e^{-\lambda_*|y-p|}$ is the clause on decay in the definition of $\mathcal R(b_*,K_*,\lambda_*,\cdot)$ (\Cref{Y-A-SetsOfRows}). The number $\tau_7$ is an entry of \textup{(O14)}, in item~\textup{(1)} of \S\ref{Y-H0-1}, and every floor is free of it. Finally, \textup{(S)} and $C_{\rm vol}\Delta\le1/2$ are conditions on distances, and no letter depends on them. \end{remark}

\subsubsection*{The proofs, in Steps P1 to P11}

\begin{proof}[Proof of \Cref{Y-H8-ThBC}] \emph{Step P1: the bound in the supremum norm.} We show that \[\sup_V|\rho[u_2;\ell]-\rho[u_3;\ell]|\le1.61\,q_Ne^{(0.61a_E+218a_w)q_N}\mathsf D_{23},\] together with the cases \textup{(i)}, \textup{(ii)}, \textup{(iii)} and with \textup{(a$^C$)}.

We write $\rho[u_i;\ell]=e^{-x\mathrm{En}_\ell}e^{\mathrm{Sm}_i}\Xi_i$, where $\Xi_i:=\Xi_{w_i}$. The energy $\mathrm{En}_\ell$, the indicators $1_{X^\approx}$ and the set of admitted families are the same for every $i$: they are determined by $V$, $\varepsilon$, $s(x)$ and the rows, and by nothing else. Moreover $0<e^{-x\mathrm{En}_\ell}\le1$, because $S_\varepsilon\ge0$, $m_y\in\{0,1\}$ and all remaining terms are squares.

\begin{enumerate}

\item[\textup{(1)}] \emph{The sums.} Let $1_{X^\approx}(V)=1$. Then the restriction of $V$ to the bonds of $X^\approx$ takes values in $G$ and is tame there; since $x>e$, it lies in $\mathcal U_X(x)$, by the definition ``Domain'' of item~2 of \S\ref{Y-J-1}. For such $V$ we therefore have \[|E_i^a+E_i^h|(X,V)\le\|u_i\|^\sim_E(c_a+c_h)e^{-\tilde\kappa(\#X-1)}\] and \[|(E_i^a+E_i^h)-(E_j^a+E_j^h)|(X,V)\le\mathsf D_{ij}(c_a+c_h)e^{-\tilde\kappa(\#X-1)}.\]

Every $X$ holds a cube, so that $\sum_X\le\sum_{\rm cube}\sum_{X\text{ holding it}}$. By the count $(26e)^{j-1}$ of item~2 of \S\ref{Y-J-1}, and since $\tilde\kappa\ge6$, the inner sum is at most $\sum_{j\ge1}(26e^{-5})^{j-1}=1/(1-26e^{-5})\le1.22$.

Summing over the $q_N$ cubes, we get $|\mathrm{Sm}_i-\mathrm{Sm}_j|\le1.22\,q_N(c_a+c_h)\mathsf D_{ij}$; as $c_a+c_h\le1/2$, it follows that $|\mathrm{Sm}_i|\le0.61\,q_Na_E$ for $i=2,3$ and that $|\mathrm{Sm}_2-\mathrm{Sm}_3|\le0.61\,q_N\mathsf D_{23}$.

When $D(V)$ is empty, the only admitted family is the empty one, because each member of an admitted family has to hold a plaquette of $D(V)$; the empty family gives $\Xi_i=1$. This is \textup{(a$^C$)}.

\item[\textup{(2)}] \emph{The exponentials.} For complex $\sigma,\sigma'$ we have $e^\sigma-e^{\sigma'}=(\sigma-\sigma')\int_0^1e^{\sigma'+t(\sigma-\sigma')}dt$, and therefore $|e^\sigma-e^{\sigma'}|\le e^{\max(\operatorname{Re}\sigma,\operatorname{Re}\sigma')}|\sigma-\sigma'|$. Together with \textup{(1)} this gives $|e^{\mathrm{Sm}_2}-e^{\mathrm{Sm}_3}|\le e^{0.61a_Eq_N}\,0.61\,q_N\mathsf D_{23}$ and $|e^{\mathrm{Sm}_3}|\le e^{0.61a_Eq_N}$.

\item[\textup{(3)}] \emph{A majorant for $\Xi$.} Two members of one family lie at a distance $>s(x)\ge8$ from each other, hence they are disjoint. We fix an order of the cubes and assign to each member its first cube in that order; distinct members receive distinct first cubes, and so no family has more than $q_N$ members.

Let $\mathrm{wt}(X):=e^{-\tilde\kappa(|X|_c-1)}$. Then \[\sum_{\rm families}\prod_ja_w\mathrm{wt}(X_j)\le\prod_{\rm cube}\bigl(1+a_w\sum_{X\text{ holding it}}\mathrm{wt}(X)\bigr)\le\exp\{a_wq_N\sum_{j\ge1}(52e)^je^{-6(j-1)}\},\] and the sum in the exponent equals $52e/(1-52e^{-5})\le217.588\le218$. Consequently $|\Xi_2|,|\Xi_3|\le e^{218a_wq_N}$.

\item[\textup{(4)}] \emph{Telescoping.} For finite products, \[\prod_jy_j-\prod_jy'_j=\sum_j(\prod_{i<j}y_i)(y_j-y'_j)(\prod_{i>j}y'_i).\] We combine this with $|w_2-w_3|(X,V)\le\mathsf D_{23}\mathrm{wt}(X)$ and $|w_i|(X,V)\le a_w\mathrm{wt}(X)$. For a family with $\#\mathrm{fam}$ members we find \[|\prod w_2-\prod w_3|\le\mathsf D_{23}\,\#\mathrm{fam}\;a_w^{\#\mathrm{fam}-1}\prod\mathrm{wt}(X_j)\le\mathsf D_{23}\,q_N\prod a_w\mathrm{wt}(X_j),\] where the second inequality holds as $\#\mathrm{fam}\le q_N$ and $a_w\ge1$; for the empty family the difference is $0$. Summing over the families and using \textup{(3)}, we arrive at $|\Xi_2-\Xi_3|\le\mathsf D_{23}\,q_Ne^{218a_wq_N}$.

\item[\textup{(5)}] \emph{The sum.} From \textup{(2)}, \textup{(3)} and \textup{(4)}, \[|e^{\mathrm{Sm}_2}\Xi_2-e^{\mathrm{Sm}_3}\Xi_3|\le|e^{\mathrm{Sm}_2}-e^{\mathrm{Sm}_3}||\Xi_2|+|e^{\mathrm{Sm}_3}||\Xi_2-\Xi_3|\le(0.61+1)\,q_N\mathsf D_{23}\,e^{(0.61a_E+218a_w)q_N}.\]

\end{enumerate}

We now turn to the three cases.

\begin{enumerate}

\item[\textup{(i)}] Here $\|u_1\|^\sim\le1$, so that $a_E=a_w=1+\mathsf D_{\max}$, with $\mathsf D_{\max}:=\max\{\mathsf D_{12},\mathsf D_{13}\}$, are admitted. The exponent is then $218.61\,q_N+218.61\,q_N\mathsf D_{\max}\le218.61\,q_N+218.61/400$, and $1.61\,e^{0.546525}\le2.78086<2e^{0.39}\le2e^{0.39q_N}$.

\item[\textup{(ii)}] In this case the coefficient of $q_N$ in the exponent is $0.61\cdot2+218=219.22$.

\item[\textup{(iii)}] Here $a_E=a_w=A$, and we use $218.61A\le219A$ and $1.61\le2$.

\end{enumerate}

\emph{Step P2: quotients and logarithms.} We first prove an inequality. Let $h>0$, and let $f,g$ be complex integrable functions with $\eta_f:=\|f-h\|_1/\!\int\!h<1$ and $\eta_g:=\|g-h\|_1/\!\int\!h<1$. Then \[\|f/\!\int\!f-g/\!\int\!g\|_1\le2\|f-g\|_1/\bigl(\int\!h\,(1-\eta_f)(1-\eta_g)\bigr),\] and when $f=h$ the left side is $\le2\eta_g/(1-\eta_g)$. From this inequality we then derive \textup{(c$^C$)}, with its clauses \textup{(c1)} to \textup{(c4)}.

For complex $f,g$ whose integrals are different from $0$ one has the identity \[f/\!\int\!f-g/\!\int\!g=(f-g)/\!\int\!f+g\,(\int\!g-\int\!f)/(\int\!f\int\!g).\] Taking the $L^1$ norm of both sides, we see that the left side has $L^1$ norm $\le(\|f-g\|_1/|\!\int\!f|)(1+\|g\|_1/|\!\int\!g|)$. At this point Step~9\,\textup{(6)} of the proof of \Cref{Y-J-ThB} uses $\|g\|_1/\!\int\!g=1$. Here we use instead $\|g\|_1\le(1+\eta_g)\int\!h$, $|\!\int\!g|\ge(1-\eta_g)\int\!h$ and $|\!\int\!f|\ge(1-\eta_f)\int\!h$, together with $1+(1+\eta_g)/(1-\eta_g)=2/(1-\eta_g)$.

\emph{The bound in $L^1$.} Since $dV$ is the product of the Haar probability measures and so has total mass $1$, the integral $\int|\rho[u_2;\ell]-\rho[u_3;\ell]|\,dV$ is at most the supremum bound of Step~P1. For the lower bound on $I_1$ we take the first inequality of \Cref{Y-J-ThB}\,\textup{(c)}, applied with $\rho_2:=\rho_1$ and with the rate $\kappa_1\ge6$. Its proof, which is Step~9\,\textup{(1)} to \textup{(5)} of the proof of \Cref{Y-J-ThB}, rests on $\mathrm{Sm}_1\ge-0.61\,q_N$, on $\Xi_1=1$ on $\{\operatorname{ang}\,V(b)\le\varepsilon/4\text{ at all bonds }b\}$, on the bound on the rows and on $\rho_1>0$. Dividing the bound in $L^1$ by this lower bound, we obtain in the three cases:

\begin{enumerate}

\item[\textup{(i)}] $2q_Ne^{219q_N}e^{0.61q_N}=2q_Ne^{219.61q_N}\le\frac12\,4q_Ne^{220q_N}$;

\item[\textup{(ii)}] $1.61\,q_Ne^{219.83q_N}\le2q_Ne^{220q_N}$;

\item[\textup{(iii)}] $2q_Ne^{219Aq_N+0.61q_N}\le2q_Ne^{220q_N}e^{219(A-1)q_N}$.

\end{enumerate}

The other two factors of $C_{\rm vol}$ are those of the lower bound.

When \textup{(S)} holds, the inequality $C_{\rm vol}\ge400q_N$ places us in case \textup{(i)}, and it does so for the pairs $(u_i,u_1)$ as well, where $\mathsf D_{11}=0$. With $h=\rho[u_1;\ell]$ we then have $\eta_i:=\|\rho[u_i;\ell]-h\|_1/I_1\le\frac12C_{\rm vol}\mathsf D_{1i}\le\frac12$.

\begin{enumerate}

\item[\textup{(c1)}] We have $|I_2-I_3|\le\|\rho[u_2;\ell]-\rho[u_3;\ell]\|_1$ and $|I_i/I_1-1|\le\eta_i$; moreover $\int|\rho[u_i;\ell]|\le(1+\eta_i)I_1\le\frac32I_1$, and $I_1\le2|I_i|$.

\item[\textup{(c2)}] Here $\rho_i/\!\int\!\rho_i=\rho[u_i;\ell]/I_i$, and we use the case $f=h$ of the inequality. The function $t\mapsto2t/(1-t)$ is increasing, and its value at $t=C_{\rm vol}\mathsf D_{12}/2$ is $C_{\rm vol}\mathsf D_{12}/(1-C_{\rm vol}\mathsf D_{12}/2)\le2C_{\rm vol}\mathsf D_{12}$.

\item[\textup{(c3)}] The inequality, with the bounds just obtained, gives $2\cdot(C_{\rm vol}\mathsf D_{23}/2)\cdot2\cdot2=4C_{\rm vol}\mathsf D_{23}$.

\item[\textup{(c4)}] Both quotients lie in the convex set $\{|w-1|\le1/2\}$, on which $|d\,\mathrm{Log}\,w/dw|=1/|w|\le2$; hence the difference of their logarithms is $\le2|I_2-I_3|/I_1\le C_{\rm vol}\mathsf D_{23}$.

\end{enumerate}

\emph{Step P3: the laws, with a complex second member.} With $\mu_i:=\rho_i/\!\int\!\rho_i$, we prove that for $\eta_d<\log2$ \[\int|\mu_1-\mu_2|\,dV\le2\,\mathrm{ms}_1+\mathrm{ms}_2\bigl(1+e^{\eta_d}/(2-e^{\eta_d})\bigr)+2(e^{\eta_d}-1)/(2-e^{\eta_d}).\]

First, $\mathrm{ms}_1<1$. Indeed, the set $\mathrm{Tm}$ contains $\{\operatorname{ang}\,V(b)\le\varepsilon/4\text{ at all bonds }b\}$, whose Haar mass is positive, as Step~9\,\textup{(1)}, \textup{(4)} of the proof of \Cref{Y-J-ThB} shows; moreover $\rho_1>0$.

On $\mathrm{Tm}$, \textup{(a$^C$)} gives $\mu_2=c'e^\sigma\mu_1$, where $\sigma:=\mathrm{Sm}_2-\mathrm{Sm}_1$ satisfies $|\sigma|\le\eta_d$ and where $c':=e^{\mathsf V_2-\mathsf V_1}\int\!\rho_1/\!\int\!\rho_2$. Splitting the integral of $\mu_2$ into its parts on and off $\mathrm{Tm}$, we get $1=\int\mu_2=c'(1-\mathrm{ms}_1)(1+t)+\mathrm{mo}$; here $t:=\int_{\mathrm{Tm}}(e^\sigma-1)\mu_1\,dV/(1-\mathrm{ms}_1)$ satisfies $|t|\le e^{\eta_d}-1<1$, and $\mathrm{mo}:=\int_{{\rm off}\,\mathrm{Tm}}\mu_2$ satisfies $|\mathrm{mo}|\le\mathrm{ms}_2$. It follows that $(1-\mathrm{ms}_1)c'=(1-\mathrm{mo})/(1+t)$, and that \[\FitDisplay{\dimexpr\linewidth-1pt\relax}{\int_{\mathrm{Tm}}|\mu_2-\mu_1|\le(1-\mathrm{ms}_1)\sup|c'e^\sigma-1|=\sup|(1-\mathrm{mo})e^\sigma/(1+t)-(1-\mathrm{ms}_1)|\le\sup|e^\sigma/(1+t)-1|+\mathrm{ms}_2e^{\eta_d}/(2-e^{\eta_d})+\mathrm{ms}_1.}\] Further, $|e^\sigma/(1+t)-1|\le(|e^\sigma-1|+|t|)/(1-|t|)\le2(e^{\eta_d}-1)/(2-e^{\eta_d})$.

Off $\mathrm{Tm}$, the integral of $|\mu_1-\mu_2|$ is at most $\mathrm{ms}_1+\mathrm{ms}_2$. \end{proof}

\begin{proof}[Proof of \Cref{Y-H8-CoF}] \emph{Step P4: the integral along a family.} We prove \textup{(F1)}: $I[u(z)]$ is holomorphic on $\mathrm{Pd}(r)$, the derivatives $\partial_S$ may be taken under the integral, and $\|u(z)-u_1\|^\sim\le B_0\max_i|z_i|/r_i$.

\begin{enumerate}

\item[\textup{(1)}] \emph{At a fixed field.} As the torus is finite, the sets $\mathcal D$ and $\mathcal P$ are finite, and so is the set of families. Thus $\rho[u(z);\ell](V)$ is the product of $e^{-x\mathrm{En}_\ell(V)}$, of the exponential of a finite sum, and of a polynomial in finitely many values, each of these values being holomorphic in $z$; hence it is holomorphic on $\mathrm{Pd}(r)$. For fixed $z$ it is a Borel function of $V$. By Step~P1\,\textup{(1)} and \textup{(3)}, applied with $a_E=a_w=1+B_0$, its modulus on $\mathrm{Pd}(r)$ is at most $\mathrm{Bd}:=e^{218.61(1+B_0)q_N}$.

\item[\textup{(2)}] \emph{The integral.} Fix radii $0<r'<r''<r$, the inequalities holding entry by entry. Take $z$ in the closure of $\mathrm{Pd}(r')$ and change $z_i$ by an amount of modulus $\le r''_i-r'_i$. The difference quotient in $z_i$ is then at most the supremum of $|\partial_i\rho|$ over the segment joining the two points, and Cauchy's estimate bounds that supremum by $\mathrm{Bd}/(r_i-r''_i)$. By dominated convergence, $I$ is therefore continuous and complex differentiable in each $z_i$, with $\partial_iI=\int\partial_i\rho\,dV$. Here $\partial_i\rho$ is a Borel function of $V$, being a limit of Borel functions, and it is holomorphic in $z$ and bounded on $\mathrm{Pd}(r'')$. Osgood's lemma states that a continuous function which is holomorphic in each variable separately is holomorphic; so $I$ is holomorphic, and by iteration $\partial_SI=\int\partial_S\rho\,dV$.

Only three properties of the integrand enter \textup{(2)}: it is holomorphic in $z$ at each $V$, Borel in $V$ at each $z$, and bounded. For $\rho[u(z);\ell]$ they are supplied by \textup{(1)}. If $u$ is any pointwise holomorphic family with $\sup\|u(z)\|^\sim\le A$ on a polydisc, then Step~P1\,\textup{(1)} and \textup{(3)}, applied with $a_E=A$ and $a_w=\max\{1,A\}$, bound the modulus of $\rho[u(z);\ell]$ there by $e^{(0.61A+218\max\{1,A\})q_N}$, which gives the third property. The other two properties follow for such a family from the first three sentences of \textup{(1)}, which use only that the family is pointwise holomorphic. The argument of \textup{(2)} therefore applies to every pointwise holomorphic family that is bounded in $\|\cdot\|^\sim$, and to every function $F(z,V)$ that is holomorphic in $z$, Borel in $V$ and bounded.

\item[\textup{(3)}] \emph{Schwarz's lemma for each value.} Fix $X$ and $\mathsf W$, and let $\mathrm{val}(z):=(E^a(X,\mathsf W;z)-E^a(X,\mathsf W;0))\,e^{\tilde\kappa(\#X-1)}/c_a(x)$. This function is holomorphic on $\mathrm{Pd}(r)$, with $|\mathrm{val}|\le B_0$ and $\mathrm{val}(0)=0$. For $z\ne0$ let $t:=\max_i|z_i|/r_i<1$. The map $\zeta\mapsto\mathrm{val}(\zeta z/t)$ is holomorphic on the unit disc, bounded by $B_0$ and equal to $0$ at $0$, so Schwarz's lemma at $\zeta=t$ gives $|\mathrm{val}(z)|\le B_0t$. The values of $E^h$ and of $w$ are treated in the same way, each multiplied by its own factor in $\|\cdot\|^\sim$ in place of $e^{\tilde\kappa(\#X-1)}/c_a(x)$. Taking suprema over all values, we obtain the bound on $\|u(z)-u_1\|^\sim$.

\end{enumerate}

\emph{Step P5: the logarithm of the integral.} We prove \textup{(F2)}: on $\mathrm{Pd}(r/\tau_7)$ we have $|I[u(z)]/I_1-1|\le1/2$, and $|\partial_S\log I[u(z)]|_0\le\log2\;\tau_7^{|S|}/r^S$.

On $\mathrm{Pd}(r/\tau_7)$, \textup{(F1)} gives $\|u(z)-u_1\|^\sim\le B_0/\tau_7\le1/C_{\rm vol}$. Hence \textup{(S)} holds for $u_2=u_3=u(z)$, and \textup{(c1)} yields $|\breve\rho_u(z)-1|\le1/2$, where $\breve\rho_u(z):=I[u(z)]/I_1$. So $\mathrm{Log}\,\breve\rho_u$ is holomorphic on this polydisc; and $|\mathrm{Log}(1+v)|\le\sum_j|v|^j/j=-\log(1-|v|)\le\log2$ whenever $|v|\le1/2$. As $I_1>0$ is a constant, for non-empty $S$ we have $\partial_S\log I[u(z)]=\partial_S\mathrm{Log}\,\breve\rho_u(z)$ near $0$, whichever branch of the logarithm is taken.

We use Cauchy's estimate in the variables of $S$, on the circles $|z_i|=t\,r_i/\tau_7$, with the other variables set to $0$ and one derivative in each variable of $S$; letting $t\to1$, we get $|\partial_S\mathrm{Log}\,\breve\rho_u(0)|\le\log2\prod_{i\in S}\tau_7/r_i$.

\emph{Step P6: two families against one positive reference.} We prove the two sentences of \textup{(F3)}.

\emph{The first sentence.} On $\mathrm{Pd}(r/2\tau_7)$, \textup{(F1)} gives $\|u(z)-u_1\|^\sim\le B_0/(2\tau_7)\le1/(2C_{\rm vol})$, and then $\|\mathfrak n(z)-u_1\|^\sim\le\Delta+1/(2C_{\rm vol})\le1/C_{\rm vol}$. So \textup{(S)} holds for $u_2=u(z)$ and $u_3=\mathfrak n(z)$, with $u_1$ as first member. Since $\|\mathfrak n(z)\|^\sim\le2$, Step~P4\,\textup{(2)} shows that $I[\mathfrak n(z)]$ is holomorphic. By \textup{(c1)}, both $\breve\rho_u(z)$ and $\breve\rho_{\mathfrak n}(z):=I[\mathfrak n(z)]/I_1$ lie in $\{|w-1|\le1/2\}$; in particular $I[\mathfrak n(0)]\ne0$.

By \textup{(c4)}, the function $\check\delta(z):=\mathrm{Log}\,\breve\rho_u(z)-\mathrm{Log}\,\breve\rho_{\mathfrak n}(z)$ is holomorphic, with $|\check\delta|\le C_{\rm vol}\Delta$. We apply Cauchy's estimate to it on $\mathrm{Pd}(r/2\tau_7)$, in the manner of Step~P5: in the variables of $S$, on circles whose radii are $t$ times those of $\mathrm{Pd}(r/2\tau_7)$, with the other variables at $0$ and one derivative in each variable of $S$, and then we let $t\to1$. For non-empty $S$ the constant $\log I_1$ is annihilated by $\partial_S$.

\emph{The second sentence.} If $C_{\rm vol}\Delta\le1/2$, the first sentence gives a stronger bound. If $C_{\rm vol}\Delta>1/2$, we apply \textup{(F2)} to $u$, and \textup{(F2)} to $\mathfrak n$ with $\mathfrak n(0)$ as its own first member. Here $C_{\rm vol}$ depends on $N$ and $x$ alone, and the number $\max\{1,C_{\rm vol}\sup_{\mathrm{Pd}(r)}\|\mathfrak n(z)-\mathfrak n(0)\|^\sim\}$, which \textup{(F0)} assigns to $\mathfrak n$ with the first member $\mathfrak n(0)$, is $\le\tau_7$, since $\sup_{\mathrm{Pd}(r)}\|\mathfrak n(z)-\mathfrak n(0)\|^\sim\le B_0$ by a hypothesis of the second sentence of \textup{(F3)}. The left side is therefore \[\le2\log2\;\tau_7^{|S|}/r^S<2\log2\cdot2C_{\rm vol}\Delta\,\tau_7^{|S|}/r^S\le3C_{\rm vol}\Delta\,(2\tau_7)^{|S|}/r^S.\]

\emph{Step P7: the lemma on jets.} We prove the clauses \textup{(i)}, \textup{(ii)}, \textup{(iii)} of \textup{(F4)}.

\begin{enumerate}

\item[\textup{(0)}] \emph{The ideal.} Let $\mathfrak a$ be the set of the germs $c$ with $c\equiv0$; we show that $\mathfrak a$ is an ideal in the ring of germs. When two monomials are multiplied, the result is $z^S$ only if they are $z^{S'}$ and $z^{S''}$ for disjoint $S'$, $S''$ with $S'\cup S''=S$. Hence, if $c\in\mathfrak a$ and $f$ is a germ, the coefficient of $z^S$ in $cf$ equals $\sum_{S'\subset S}c_{S'}f_{S\setminus S'}=0$. Consequently $\equiv$ is compatible with finite sums and finite products, and every $c\in\mathfrak a$ satisfies $c(0)=0$.

Moreover $e^c-1=c\cdot\mathrm{pw}_1(c)$ and $\log(1+c)=c\cdot\mathrm{pw}_2(c)$, where $\mathrm{pw}_1$ is entire and $\mathrm{pw}_2$ is holomorphic near $0$; so both lie in $\mathfrak a$. It follows that $f\equiv g$ implies $e^f-e^g=e^g(e^{f-g}-1)\in\mathfrak a$; and if $f\equiv g$ and $g(0)\ne0$, then $1/f-1/g=(g-f)/(fg)\in\mathfrak a$, and $\log f-\log g=\log(1+(f-g)/g)\in\mathfrak a$ for branches that agree at $0$.

\item[\textup{(i)}] Since the triples without a size form a complex vector space (\Cref{Y-H8-Triple}), $u^{\rm pol}(z)$ is such a triple for every $z$. Its values are polynomials, $u^{\rm pol}(0)=u_1$, and on $\mathrm{Pd}(r)$ we have $\|u^{\rm pol}(z)-u_1\|^\sim\le\sum_S|z^S|\,\|\delta_S\|^\sim\le B_0$.

\item[\textup{(ii)}] At a fixed $V$, $\rho[\cdot;\ell](V)$ is computed from finitely many values by sums, products and a single exponential, so \textup{(0)} applies to it. Step~P4\,\textup{(2)} lets us take $\partial_S$ under the integral for each of the two families; hence $I[u(z)]\equiv I[\hat u(z)]$. As $I[u(0)]\ne0$, \textup{(0)} applies to the logarithm, and $\partial_S$ at $0$ picks out the coefficient of $z^S$.

\item[\textup{(iii)}] By Step~P4\,\textup{(2)}, applied to $F$ and to $\rho[u(z);\ell]$, we have for all $S$ \[\partial_S\!\int\!F\,dV|_0=\int\partial_SF(0,V)\,dV=\int\partial_S\bigl(e^{\mathfrak V(z)}\rho[u(z);\ell](V)\bigr)|_0\,dV=\partial_S\bigl(e^{\mathfrak V(z)}I[u(z)]\bigr)|_0.\] So $\int F\,dV\equiv e^{\mathfrak V}I[u(z)]$, whose value at $0$ is $e^{\mathfrak V(0)}I[u(0)]\ne0$. By \textup{(0)}, $\log\int F\,dV\equiv\mathfrak V+\log I[u(z)]$ up to an additive constant in $2\pi i\mathbb Z$, and $F/\!\int\!F\,dV\equiv e^{\mathfrak V}\rho[u(z);\ell](V)/(e^{\mathfrak V}I[u(z)])$.

\end{enumerate}

\emph{Step P8: the jets of the laws in $L^1$.} We prove \textup{(F5)}: \[\int|\partial_S\mu_u^z|_0\,dV\le3\tau_7^{|S|}/r^S,\] and \[\int|\partial_S\mu_u^z-\partial_S\mu_{\mathfrak n}^z|_0\,dV\le4C_{\rm vol}\Delta\,(2\tau_7)^{|S|}/r^S\] if $C_{\rm vol}\Delta\le1/2$, and also for every $\Delta$ under the hypotheses of the second sentence of \textup{(F3)}.

On $\mathrm{Pd}(r/\tau_7)$ the condition \textup{(S)} holds, by Step~P5, and \textup{(c1)} gives $\|\mu_u^z\|_1=\int|\rho[u(z);\ell]|/|I[u(z)]|\le3$. On $\mathrm{Pd}(r/2\tau_7)$ the condition \textup{(S)} holds for the pair $u(z)$, $\mathfrak n(z)$, by Step~P6, and \textup{(c3)} gives $\|\mu_u^z-\mu_{\mathfrak n}^z\|_1\le4C_{\rm vol}\Delta$.

Let $m^z(V)$ be holomorphic in $z$ at each fixed $V$. Then $\partial_Sm^z(V)|_0$ equals $(2\pi i)^{-|S|}$ times the integral of $m^z(V)\prod_{i\in S}dz_i/z_i^2$ over the circles $|z_i|=\mathrm{rd}_i$, $i\in S$, the other variables being $0$. Since $m^z(V)$ is continuous in $z$ and Borel in $V$, it is jointly Borel, and Tonelli's theorem gives $\int|\partial_Sm^z|_0\,dV\le\sup_z\|m^z\|_1/\prod_{i\in S}\mathrm{rd}_i$. We let $\mathrm{rd}_i$ tend to $r_i/\tau_7$ for the first bound and to $r_i/2\tau_7$ for the second.

\emph{The second sentence.} If $C_{\rm vol}\Delta\le1/2$, it is contained in the first. So let $C_{\rm vol}\Delta>1/2$. On $\mathrm{Pd}(r/\tau_7)$, \textup{(S)} holds for $u_2=u(z)$, with $u_1$ as first member, by Step~P5, and \textup{(c2)} gives $\|\mu_u^z-\mu_u^0\|_1\le2C_{\rm vol}\|u(z)-u_1\|^\sim\le2$. As $\mu_u^0$ is free of $z$, $\partial_S\mu_u^z=\partial_S(\mu_u^z-\mu_u^0)$ for non-empty $S$. The Cauchy formula and Tonelli's theorem, now applied to $\mu_u^z-\mu_u^0$, whose $L^1$ norm is at most $2$, with $\mathrm{rd}_i$ tending to $r_i/\tau_7$, give $\int|\partial_S\mu_u^z|_0\,dV\le2\tau_7^{|S|}/r^S$.

The same bound holds for $\mathfrak n$, with $\mathfrak n(0)$ as its own first member; here the number that \textup{(F0)} assigns to $\mathfrak n$ is $\max\{1,C_{\rm vol}\sup_{\mathrm{Pd}(r)}\|\mathfrak n(z)-\mathfrak n(0)\|^\sim\}\le\tau_7$, because $\sup_{\mathrm{Pd}(r)}\|\mathfrak n(z)-\mathfrak n(0)\|^\sim\le B_0$ is one of the hypotheses of the second sentence of \textup{(F3)}. Adding the two bounds, we get $4\tau_7^{|S|}/r^S<8C_{\rm vol}\Delta\,\tau_7^{|S|}/r^S\le4C_{\rm vol}\Delta\,(2\tau_7)^{|S|}/r^S$, because $2^{|S|}\ge2$. \end{proof}

\begin{proof}[Proof of \Cref{Y-H8-CoR7}] \emph{Step P9: the integrals of the runs and the two formulas.} We prove $\int\rho^{[K],z}_k=Z_K(z)$ and the two formulas for $S^T$.

\emph{Bounds at a fixed $K$.} We have $0\le a_p\le2$, and the bare lattice has $3(NL^K)^3$ plaquettes, so that $|\mathcal O_K(f)|\le6L^K(NL^K)^3\|f\|_\infty$. It follows that $\rho^{[K],z}_0$ is continuous in $U$ and entire in $z$, and that it is bounded on compact sets of $z$, uniformly in $U$. The kernel $\prod_ck_{ax}(V'(c),\mathsf M(U)(c))$, where $\mathsf M$ is the average in the definition ``The step'' of item~2 of \S\ref{Y-J-1}, is a Borel function of $(V',U)$ with values $>0$ and $\le z(ax)^{-\#c}$, $\#c$ being the number of coarse bonds; at every $U$ it is a probability density in $V'$. By induction on $k$, with the help of Step~P4\,\textup{(2)}, every $\rho^{[K],z}_k$ is Borel in $V$ and entire in $z$, and it is bounded on compact sets of $z$, uniformly in $V$.

\begin{enumerate}

\item[\textup{(R0)}] Fubini's theorem for a bounded Borel complex function gives $\int\rho'(V')\,\allowbreak dV'=\int\rho(U)\,\allowbreak [\int\prod_ck\,\allowbreak dV']\,\allowbreak dU=\int\rho(U)\,\allowbreak dU$.

\item[\textup{(R1)}] At $x=L^s\mathsf b$, on $\mathbb Z^3/NL^s\mathbb Z^3$, both the data \textup{(B)} and the bound on the rows are satisfied. Indeed, $\mathsf b\ge x_s\ge e^2$; $NL^s$ is even, a multiple of $M$, and $\ge2M\,s(L^s\mathsf b)$, by Step~2\,\textup{(e)} of the proof of \Cref{Y-J-ThA}; $\varepsilon\le1/40$ by \textup{(X16)}; $c_a+c_h\le1/32$ by \textup{(X17)}; and $\ell_k$ has range $\le\varrho(x_k)$ and satisfies $|\ell_{k,y}(p)|\le K_*e^{-\lambda_*|y-p|}$, by Step~2\,\textup{(c)} of the proof of \Cref{Y-J-ThA} and by Step~11\,\textup{(a)} of the proof of \Cref{Y-J-CorC}. We apply \textup{(F4)}\,\textup{(iii)} with $F=\rho^{[K],z}_k$ and $S=\{1,\dots,n\}$; together with \textup{(R0)} it gives $S^T_{n;K}=\partial_1\cdots\partial_n\log\int F\,dV|_0$. We then apply \textup{(F2)}.

\item[\textup{(R2)}] We apply \textup{(F4)}\,\textup{(iii)} to both runs. At the level $m+k$ of the run of length $K'$ and at the level $k$ of the run of length $K$, the two runs share the torus, the value $x_k$, the grid and the step, by Step~2\,\textup{(a)} of the proof of \Cref{Y-J-ThA}, and both densities are defined with the rows $\ell_k$. Here $I[\mathfrak n(0)]\ne0$, by the first sentence of \textup{(F3)}, or since the density of which $\mathfrak n(0)$ is a system of members is positive. We subtract the two formulas and bound the difference by \textup{(F3)}.

For polynomial families, $\|u(z)-\mathfrak n(z)\|^\sim\le\|\bar u_k-\bar{\mathfrak n}_k\|^\sim+\sum_S|z^S|\,\|\delta_S-\tilde\delta_S\|^\sim$.

\end{enumerate}

\emph{Step P10: the jets of the laws of the drawn fields.} We prove the clauses \textup{(i)} to \textup{(v)} of \textup{(R3)}.

As $Z_K(0)>0$, the quotient $\rho^{[K],z}_{K-s}(V)/Z_K(z)$ is holomorphic near $0$ and bounded uniformly in $V$, by Step~P9. Therefore $\pi_S$ is bounded, by Cauchy's estimate, and Borel, as a limit of Borel functions.

\begin{enumerate}

\item[\textup{(i)}] By Step~P4\,\textup{(2)}, $\int\mathrm{obs}\,\pi_S\,dV=\partial_S\bigl[\int\mathrm{obs}\,\rho^{[K],z}_{K-s}\,dV/Z_K(z)\bigr]_0$. Under the joint law of the bare field $U$ and of the drawn fields we have $E[\mathrm{obs}(V_{K-s})\,e^{\sum z_i\mathcal O_K(f_i)}]=\int\mathrm{obs}\,\rho^{[K],z}_{K-s}\,dV/Z_K(0)$, since no kernel depends on a source. Hence the bracket equals $\frac{d}{dt}\log E[e^{\sum z_i\mathcal O_K(f_i)+t\,\mathrm{obs}}]$ at $t=0$, and its derivative $\partial_S$ at $z=0$ is the joint cumulant, that is, the truncated function. For $\mathrm{obs}=1$ we get $\partial_S1=0$.

The function $\mathcal O_K(f)$ is gauge invariant, and the step, which is linear in $\rho$, carries invariant functions into invariant functions (the definition ``The step'' of item~2 of \S\ref{Y-J-1}). Hence $\rho^{[K],z}_k$ and $\pi_S$ are gauge invariant as well. When the $f_i$ are real, the quotient takes real values at real $z$.

\item[\textup{(ii)}] Let $\mathrm{Kn}\,g(V'):=\int g(U)\prod_ck_{aL^s\mathsf b}(V'(c),\mathsf M(U)(c))\,dU$. Since $|\mathrm{Kn}\,g|\le\mathrm{Kn}|g|$ and $\int\mathrm{Kn}|g|=\int|g|$, we have $\|\mathrm{Kn}\,g\|_1\le\|g\|_1$. Further, $\mathrm{Kn}\,\rho^{[K],z}_{K-s}=\rho^{[K],z}_{K-s+1}$, the noise being $a x_{K-s}=aL^s\mathsf b$; $Z_K(z)$ does not depend on the field and is the same number in $\pi^{[K]}_{S,[s]}$ as in $\pi^{[K]}_{S,[s-1]}$; and $\partial_S$ passes under the integral over $U$, by Step~P4\,\textup{(2)}.

\item[\textup{(iii)}] By \textup{(R0)}, $\int F\,dV=Z_K(z)$ for $F=\rho^{[K],z}_{K-s}$, and $\int F\,dV=Z_{K'}(z)$ for $F=\rho^{[K'],z}_{K'-s}$. The second claim of \textup{(F4)}\,\textup{(iii)} yields $\pi^{[K]}_{S,[s]}=\partial_S\mu_u^z|_0$ for the first $F$ and $\pi^{[K']}_{S,[s]}=\partial_S\mu_{\mathfrak n}^z|_0$ for the second; here $I[\mathfrak n(0)]\ne0$, either by the first sentence of \textup{(F3)} or because $\mathfrak n(0)$ is a system of members of a density that is positive. We then apply \textup{(F5)}, together with its second sentence.

\item[\textup{(iv)}] Clause \textup{(iii)} shows that the sequence is a Cauchy sequence in the complete space $L^1(dV)$; we let $K'\to\infty$ in \textup{(iii)}. The integral, the invariance (the Haar measure being invariant), the reality and the first bound are preserved under limits in $L^1$. We let $K\to\infty$ on both sides of \textup{(ii)}, using that $\mathrm{Kn}$ is continuous in $L^1$ and independent of $K$.

\item[\textup{(v)}] The run $(K-s,L^s\mathsf b,NL^s)$ starts from $e^{-L^{K-s}L^s\mathsf b\sum(1-\cos\phi_p)}$ on $\mathbb Z^3/(NL^s)L^{K-s}\mathbb Z^3$, and its noises are $aL^{K-s-k}L^s\mathsf b=a x_k$. Hence at the levels $0\le k\le K-s$ it has the same densities as the run $(K,\mathsf b,N)$, and $\pi_S$ is independent of the grid, of the rows and of the members.

With $\mathrm{bar}(p)$ the barycentre of $p$ in lattice units, \[\FitDisplay{\dimexpr\linewidth-1pt\relax}{\mathcal O_{K-s}(\hat f)=L^{K-s}\sum_pf(L^{-s}L^{-(K-s)}\mathrm{bar}(p))\,\allowbreak a_p=L^{-s}\mathcal O_K(f)}\]. Thus $\sum z_i\mathcal O_K(f_i)=\sum(L^sz_i)\mathcal O_{K-s}(\hat f_i)$, and every derivative produces a factor $L^s$.

With $(L^s\mathsf b,NL^s)$ in the place of $(\mathsf b,N)$, all conditions of the setting of item~4 of this section are met: $L^s\mathsf b\ge x_s$; $NL^s=2L^{m_l+s}$; and \textup{(Tor)} follows from \textup{(RB)} in item~3 of \S\ref{Y-J-1}, by which $N\ge2M\,s(x)$ implies $NL^j\ge2M\,s(L^jx)$, applied with $x=L\mathsf b$.

\emph{The two thresholds.} Let $\mathrm{Kt}:=K_T(L^s\mathsf b,n,L^sd)$. We claim that $\mathrm{Kt}+s\ge K_T(\mathsf b,n,d)$ and $\mathrm{Kt}\ge K_0(L^s\mathsf b)$; then every $K$ with $K-s\ge\mathrm{Kt}$ satisfies $K\ge K_T(\mathsf b,n,d)$ and $K-s\ge K_0(L^s\mathsf b)$. The second claim is part of the definition of $K_T$.

For the first claim, note that $\Lambda(L^s\mathsf b)=\Lambda+s/2$, whence $K_0(L^s\mathsf b)=K_0(\mathsf b)+(5/4)s$; that the $\mathrm{sep}^{[i]}$ computed from $(L^s\mathsf b,L^sd)$ is the $\mathrm{sep}^{[i+s]}$ computed from $(\mathsf b,d)$, while $\mathrm{sep}_*$ depends on $n$ and $L$ alone, so that the definition of $j^*$, with $(L^s\mathsf b,L^sd)$ in the place of $(\mathsf b,d)$, gives $\max\{j^*-s,1\}$; and that $C_5^\#$ and $K_5$ are independent of $\mathsf b$. Let $\mathrm{Ex}_t:=(4/9)\mathrm{Kt}-(10/9)(\Lambda+s/2)-1$. This is the left side of \textup{(K1a)} and the exponent of $0.3683$ in \textup{(K1b)} when $\mathrm{Kt}$ replaces $K$ and $L^s\mathsf b$ replaces $\mathsf b$, and hence $\Lambda+s/2$ replaces $\Lambda$. Let $K:=\mathrm{Kt}+s$. Then $(4/9)K-(10/9)\Lambda-1=\mathrm{Ex}_t+s$.

\begin{enumerate}

\item[\textup{(K1a)}] This condition holds for the data $(K,\mathsf b,d)$, since $\mathrm{Ex}_t+s\ge\max\{j^*-s,1\}+1+s\ge j^*+1$.

\item[\textup{(K1b)}] This condition holds there as well, since $C_5^\#K_5^{-1}K\,0.3683^{\mathrm{Ex}_t+s}\le0.1\,(1+s/\mathrm{Kt})\,0.3683^s\le0.1\,(1+s/3)\,0.3683^s\le0.1$. Here an integer $\mathrm{Kt}\ge K_0(L^s\mathsf b)>9/4$ is $\ge3$, and $(1+s/3)\,0.3683^s\le0.4911$ for $s\ge1$: the value at $s=1$ is $0.49107$, and the quotient of the value at $s+1$ by the value at $s$ is $\le(5/4)\,0.3683<1$.

\end{enumerate}

Finally $K\ge K_0(L^s\mathsf b)+s\ge K_0(\mathsf b)$. So $K$ is an integer $\ge K_0(\mathsf b)$ with \textup{(K1a)} and \textup{(K1b)}, and $K_T(\mathsf b,n,d)$ is the least such integer; this proves the first claim.

\end{enumerate}\end{proof}

\begin{proof}[Proof of \Cref{Y-H8-Datum}] \emph{Step P11: the origin of the families and of the datum.} We show that the families of \Cref{Y-H6-J}\,\textup{(J6)} and of \Cref{Y-H7-Df}\,\textup{(D1)} at the last level have these properties: they are pointwise holomorphic; they are measured in a norm which dominates $\|\cdot\|^\sim$ at a rate in $[\kappa,\kappa_1]$, with the cap $c_h$; they are defined on polydiscs, the family of \textup{(J6)} on one containing $\mathrm{Pd}(\mathfrak r(L^{j^*}\mathsf b))$ and the family of \textup{(D1)} on one containing $\mathrm{Pd}(\mathfrak r(L^{j^*}\mathsf b)/2)$, with the radius $\mathfrak r(L^{j^*}\mathsf b)$ independent of $K$; and they start at systems of members of points of the class.

The assertion on the radius of the polydiscs follows from \textup{(G3)} in item~3 of \S\ref{Y-H6-1}, by which $\mathfrak r_*\ge\mathfrak r(L^{j^*}\mathsf b)$ for every $K\ge1$, with equality if $K\ge j^*$, while for $K<j^*$ the radius $\mathfrak r_*=\mathfrak r(L^K\mathsf b)$ depends on $K$. The restriction of a family to a smaller polydisc is a family as well, with the same $B_0$.

\Cref{Y-H6-J}\,\textup{(J6)} provides, for $k_d<k\le K$, the family $u_k(z)$ on $\mathrm{Pd}(\mathfrak r_*)$, in $M_1$ with the strong cap, and for $k_d\le k\le K$ the bound $\|u_k(z)-\bar u_k\|^{{\rm st},(\kappa)}\le0.91\cdot0.4^{k-k_d}\le0.91$, in the norm with the rate $\kappa_k$, where $\kappa_k=\kappa$ for $k\ge k_d$. At $k=k_d$ the family is $u_{k_d}^{\rm pol}$, which lies in $M_{1.91}$ by \textup{(J5)}; we use the case $k=K>k_d$ alone. The norm $\|\cdot\|^{{\rm st},(\kappa)}$ measures $w$ in strong units, so it dominates $\|\cdot\|^\sim$ at $\tilde\kappa=\kappa$. Since $k_d\le K-1$, we have $u_K(z)=F_{\ell_{K-1}}(u_{K-1}(z))$; for this output the clauses of \Cref{Y-H-C9} hold for the pair $(\kappa,c_h)$, whose rate is $\kappa<\kappa_1$, by \Cref{Y-H1-U1}\,\textup{(U)}, and \Cref{Y-H1-U1}\,\textup{(c+)} shows that it is pointwise holomorphic.

As to the true multi-affine jets: $u_K(z)\equiv\hat u_K(z)$, where $\hat u_K(z)$ is the family of \Cref{Y-H4-U4}\,\textup{(h)}, for which $\rho^{[K],z}_K=e^{\mathfrak V(z)}\rho[\hat u_K(z);\ell_K]$ at every field. This congruence is shown at the end of \textup{(3)} below, and with \textup{(F4)}\,\textup{(ii)} it gives \textup{(J.b)}. Moreover \Cref{Y-H6-J}\,\textup{(J7)} states the formula of \textup{(R1)} at depth $0$.

We now supply what \textup{(F3)} requires for $s=0$: the second family, the bound on $\Delta$, the polydisc and the base point. On $\mathrm{Pd}(\mathfrak r_*/2)$, \Cref{Y-H7-Df}\,\textup{(D1)} provides $\mathfrak n_K(z)$, and \textup{(D3)} bounds $\sup\mathsf D_{\mathsf b}^{(\kappa)}(u_K(z),\mathfrak n_K(z))$, the rate of this distance being $\kappa_K=\kappa$; the distance $\mathsf D_{\mathsf b}^{(\kappa)}$ dominates $\|\cdot\|^\sim$; and at $r_i=\mathfrak r_*/2$ the polydisc $\mathrm{Pd}(r/2\tau_7)$ lies in $\mathrm{Pd}(\mathfrak r_*/2)$. Further, $\mathfrak n_K(0)=\bar{\mathfrak n}_K$ is a real system as in \Cref{Y-J-ThA}\,\textup{(iii)}. That clause asserts that such a system exists, without singling one out, and the system in question is the one of \Cref{Y-H7-Df}\,\textup{(D1)}\,\textup{(c)}. What we use of it is \Cref{Y-H7-Df}\,\textup{(D0)} at $z=0$: the system is real, in $M_1(\mathsf b;\kappa_1,c_h)$, with the strong cap, and $\rho[\bar{\mathfrak n}_K;\ell_K]>0$, this density being a positive multiple of $\rho^{[K']}_{K'}$.

We use three further facts: \textup{(1)} for \textup{(J.a)}, for every $K\ge1$; \textup{(2)} for the value $B_0=1$ in the hypotheses of the second sentence of \textup{(F3)}; and \textup{(3)} for \textup{(J.c)} and for the congruence $u_K(z)\equiv\hat u_K(z)$ of the short run.

\begin{enumerate}

\item[\textup{(1)}] For every $K\ge1$, $\bar u_K$ is a system of members of a point of the class, its reality and positivity coming from \Cref{Y-H-C9}\,\textup{(d)} from level $0$ on. This is what the last sentences of \Cref{Y-H4-U4}\,\textup{(h)} say: at $z=0$, where the sources disappear from the weight, whatever the $f_i$ are, all members are real and $\rho[u_k(z);\ell_k]>0$; the $f_i$ of the setting may be complex, and for real $f_i$ the same holds at real $z$. Hence $\bar u_k$ is a system of members of $\rho_k$ for $0\le k\le K$: it is real, in $M_1^{\rm su}(x_k)$, with the strong cap.

\item[\textup{(2)}] We have $\sup_{\mathrm{Pd}(\mathfrak r_*/2)}\|\mathfrak n_K(z)-\mathfrak n_K(0)\|^\sim\le1$, the supremum being at most $0.455+0.34$. This is \Cref{Y-H7-Df}\,\textup{(D1)}\,\textup{(f)}, which gives $\le0.795\,\vartheta[2]^{K-k_d}$ in the strong norm with the rate $\kappa$, and that norm dominates $\|\cdot\|^\sim$; at a depth $s\ge j^*$ it is \textup{(D1)}\,\textup{(e)}.

\item[\textup{(3)}] The pointwise congruence $\rho^{[K'],z}_{K'}(V)\equiv e^{\mathfrak V'(z)}\rho[\mathfrak n_K(z);\ell_K](V)$ is \textup{(J.c)} at $s=0$. The families of \textup{(D1)} are constructed so that this congruence holds level by level, by \Cref{Y-H5-Bc}\,\textup{(B1)} at the switch and by \Cref{Y-H-C9}\,\textup{(a)} and \Cref{Y-H5-LJ} at every later step. Two clauses of \Cref{Y-H7-Df} yield it.

By \textup{(D0)}, on the interior of the polydisc $\mathrm{Pd}_0$ named there for the run of length $K'$ we have, at every field, \[\rho^{[K'],z}_{m+k}=\exp\{\tilde{\mathsf V}_0(z)+\sum_{k_0\le j<k}\mathsf V'_{\ell_j}[\mathfrak n^\circ_j(z)]\}\,\rho[\mathfrak n^\circ_k(z);\ell_k],\] where the numbers in the exponent are holomorphic; we write $\mathfrak V'(z)$ for the exponent. By \textup{(D1)}\,\textup{(c)} for $k\le k_d$, and by \textup{(D1)}\,\textup{(c)} with \textup{(f)} for $k_d<k\le K$, the family $\mathfrak n_k(z)$, which for $k\le k_d$ is $\mathfrak n_k^{\rm pol}(z)$, and the family $\mathfrak n^\circ_k(z)$ have the same base point and the same multi-affine jets; so they are congruent. The first claim of \textup{(F4)}\,\textup{(ii)} then gives $\rho[\mathfrak n_k(z);\ell_k](V)\equiv\rho[\mathfrak n^\circ_k(z);\ell_k](V)$ at every $V$, and multiplication by $e^{\mathfrak V'(z)}$ preserves $\equiv$, by Step~P7\,\textup{(0)}. This proves \textup{(J.c)}.

For the short run the argument is the same, with \Cref{Y-H4-U4}\,\textup{(h)}, \textup{(o2)} in item~2 of \S\ref{Y-H6-1} and \Cref{Y-H6-J}\,\textup{(J6)} in place of the two clauses of \Cref{Y-H7-Df}; this proves the congruence $u_K(z)\equiv\hat u_K(z)$ used above for \textup{(J.b)}.

\end{enumerate}\end{proof}

\section{Correlation bounds and limits}\label{sec:correlation-limits}

\subsection{Fixed-torus convergence}\label{Y-H9-1}

On a fixed torus we bound the truncated functions of the action density uniformly in the cutoff, show that they converge, with the explicit bound $C'_nK\theta^K$ of \textup{(T2)} and \textup{(T3)} on the speed of the convergence in $K$, and derive the properties of their limit. In Steps 0 to 3 of the proof of \Cref{Y-H9-T}, whose clauses are \textup{(T1)} to \textup{(T7)} and \textup{(T7+)}, we introduce the generating function, prove \Cref{Y-H9-LN} on the numbers of \textup{(O14)}, deduce \textup{(T1)} from \Cref{Y-H6-J} and \Cref{Y-H8-ThBC}, and deduce \textup{(T2)} from \Cref{Y-H7-Df} and \Cref{Y-H8-ThBC}. Steps 4 to 11 construct the limit, prove the dilation identities, the continuity in $\mathsf b$, the invariance, the reflection positivity and the regularity, pass to every $\mathsf b>0$ and every side, and treat the laws together with their jets.

\emph{Item 1: letters and numbers.} We retain the parameter table and floors of \S\ref{Y-H0-1}, with $\kappa_1=\kappa+1$. In particular $M=L^{m'}$ with $m'\ge1$, $p_0(x)=A_0(\log x)^{\mathfrak p_0}$, $\mathfrak p_0>2$, $\varepsilon(x)=p_0(x)x^{-1/2}$, $c_a(x)=E_ap_0(x)^3x^{-1/2}$ and $c_h(x)=e^{-c_hp_0(x)^2}$. The floor $x_s$ satisfies \textup{(O13)} there, including all the unsourced floors at $\kappa_1$, $Lx_s\ge\max\{x_9^U,x_3,x_4\}$, \textup{(X5)}, and the final conditions \textup{(W2)}, \textup{(W3)}. Thus $p_0\ge1$, $\varepsilon\le1/40$ and $16(c_a+c_h)\le1/2$ for $x\ge x_s$.

\emph{Radii.} We use $r(x)=\lceil20c_R\log x/M\rceil$, $s(x)=8r(x)$, $\varrho(x)=\lceil c_R\log x\rceil$ and $\check r(x)=r(x)+3$.

The quantities $\Lambda$, $K_0$, $k_0$, $\vartheta$, $\theta$, $C_{\mathsf b}$, $\mathrm{amp}$, $\mathfrak r$, $\mathrm{sep}^{[i]}$, $\mathrm{sep}_*$, $j^*$, $N_1^{[j]}$, $m^{[j]}$, $K_T$, $\mathrm{Tail}$ and $\mathrm{Ser}_2$ are those of items~1.1 and~1.3 of \S\ref{Y-H7-1}, for the data $(\mathsf b,n,d,R_f)$. The volume constant is \Cref{Y-H8-Cvol}. No new condition is imposed on these quantities. We additionally put
\[
\mathrm{Ser}_1:=\mathrm{Tail}(0),\qquad
P_*:=\mathfrak r(L^{j^*}\mathsf b)^{-n}N_1^{[j^*]}(C_g/16)e^\kappa.
\]
The definitions give the identity
\[
C_{L^s\mathsf b}\theta^{K-s}=3\vartheta^{(4/9)K-s-(10/9)\Lambda-1}.
\]
In a run of length $K$, $x_k=L^{K-k}\mathsf b$, the depth is $K-k$, $k_d=\max\{K-j^*,0\}$, $\mathfrak r_k=\mathfrak r(x_k)$ and $\mathfrak r_*=\mathfrak r_{k_d}\ge\mathfrak r(L^{j^*}\mathsf b)$, by \textup{(G3)} of item~3 of \S\ref{Y-H6-1}. We write $\mathrm{Pd}(r)=\{z\in\mathbb C^n:|z_i|<r\}$.

\emph{Distances.} The separation $d$, the radius $R_f$ and the balls refer to the sup distance of $T_l$; a tuple that is of the type $(d;R)$ in \Cref{support:150-3}, where the distance is Euclidean, is of the type $(d/\sqrt3;R)$ here. The Lipschitz constant $\mathrm{Lip}f$, by contrast, refers to the Euclidean distance, and $\|f\|_\sharp=\max\{\|f\|_\infty,40M\,\mathrm{Lip}f\}$, with $M$ the letter of \textup{(O7)}.

We use the following conventions for test functions and isometries.

\begin{definition}[Separated tuples and functionals on them]
\label{support:150-3} Let $\mathbb E=T_l$. For $O$ open in $\mathbb E$ and $n\ge0$, the set $\mathfrak S_n(O)$ consists of the $n$-tuples $X=(f_1,\dots,f_n)$ of real $C^1$ functions whose supports are compact, contained in $O$ and pairwise disjoint; the entry $0$ is allowed. The set $\mathfrak S_0(O)$ consists of the empty tuple alone, $\mathfrak S(O)=\bigcup_{n\ge0}\mathfrak S_n(O)$, and $\mathfrak S_n^\infty(\mathbb E)$ is the set of the tuples in $\mathfrak S_n(\mathbb E)$ all of whose entries are $C^\infty$.

For a set $J$ of indices we write $X_J=(f_i)_{i\in J}$, and $X\cdot Y$ is the tuple $X$ followed by the tuple $Y$; for an isometry $r$ we put $rf=f\circ r^{-1}$ and $rX=(rf_1,\dots,rf_n)$; and $\tau_v:x\mapsto x+v$. The cubic group $W_3$ consists of the $48$ signed permutation matrices of $\mathbb R^3$ . Each of them maps $l\mathbb Z^3$ and $a_K\mathbb Z^3$ onto themselves, so that $W_3$ acts on $T_l$ and on every lattice of spacing $a_K$; the reflection $y_\mu\mapsto-y_\mu$ is an element of $W_3$.

A tuple is of the \emph{type} $(d;R)$ if the supports of any two different entries are at distance $\ge d$ from each other and each support lies in a ball of radius $R$. We write $\mathrm{Part}_{\ge2}(n)$ for the set of the partitions of $\{1,\dots,n\}$ into blocks of size $\ge2$; the set $\mathrm{Part}_{\ge2}(0)$ consists of the empty partition, and $\mathrm{Part}_{\ge2}(1)$ is empty.

A map on $\mathfrak S_n(O)$ is \emph{multilinear} if it is $n$-linear on $C^1_c(O_1)\times\cdots\times C^1_c(O_n)$ for all open sets $O_i\subset O$ with pairwise disjoint closures, and \emph{symmetric} if every permutation of the entries leaves it unchanged. A \emph{functional on separated tuples} is a multilinear map $\sigma:\mathfrak S(\mathbb E)\to\mathbb R$ that takes the value $1$ at the empty tuple. A \emph{family of connected functions} is a sequence $S=(S_n)_{n\ge2}$ of multilinear maps $S_n:\mathfrak S_n(\mathbb E)\to\mathbb R$; its \emph{centred functional} \eqref{support:9-14} is \[\sigma_S(X)=\sum_{\pi\in\mathrm{Part}_{\ge2}(n)}\prod_{J\in\pi}S_{|J|}(X_J)\] for tuples $X$ of length $n\ge0$. A family $S$ is said to have a property of functionals when $\sigma_S$ has it, with one exception: $S$ is called symmetric when every $S_n$ is symmetric. A map is \emph{invariant} under $r$ if its value at $rX$ equals its value at $X$.

A \emph{mirror} of $T_l$ is a pair $\Sigma=(\mu,c)$ with $\mu\in\{1,2,3\}$ and $c\in\mathbb R/l\mathbb Z$; its reflection $\theta_\Sigma$ is induced by $y\mapsto y-2(y_\mu-c)e_\mu$, and its open side is $\mathbb E_+(\Sigma)=\{c<y_\mu<c+l/2\}$. A functional $\sigma$ is \emph{reflection positive at} $\Sigma$, in the sense of \eqref{support:9-28}, if \[\sum_{i,k}\bar\zeta_i\zeta_k\,\sigma(\theta_\Sigma X_i\cdot X_k)\ge0\] for every $r$, all $X_1,\dots,X_r\in\mathfrak S(\mathbb E_+(\Sigma))$ and all complex numbers $\zeta_i$.

Let $\mathbb E^n_{\ne}$ be the set of the points of $\mathbb E^n$ whose entries are pairwise different. A \emph{distribution} on $\mathbb E^n_{\ne}$ is a complex linear map $D$ on $C_c^\infty(\mathbb E^n_{\ne};\mathbb C)$ such that for each compact set there are $k$ and $C$ with $|D(F)|\le C\max_{|\alpha|\le k}\sup|\partial^\alpha F|$ for all $F$ carried by that set; it is \emph{of order $\le k$ on compact sets} if a single $k$ suffices for all compact sets. A functional $\sigma$ is \emph{regular} if for every $n$ there is a distribution $\sigma_n$ on $\mathbb E^n_{\ne}$ with $\sigma_n(f_1\otimes\cdots\otimes f_n)=\sigma(f_1,\dots,f_n)$ on $\mathfrak S_n^\infty(\mathbb E)$.
\end{definition}

On the lattice (\ref{support:150-3}) the functional is $\sigma_{K,l,\beta}(X)=\langle\prod_i(\mathcal O_K(f_i)-\langle \mathcal O_K(f_i)\rangle)\rangle$, and $\sigma_{K,l,\beta}=\sigma_S$ at $S=(S^T_{n;K,l,\beta})_{n\ge2}$; this is the classical expression of centred moments through cumulants.

In \textup{(T6)} and in Steps 5, 6 and 10, which are free of sources, $z\in\mathbb Z$ is the integer of the dilation; everywhere else $z$ stands for the sources in $\mathbb C^n$. The letter $s$ is the depth in \textup{(T7+)}, in $C_{L^s\mathsf b}$, in $C^{(s)}$ and in Step~11; $s(\cdot)$ denotes the radius function $s(x)=8r(x)$, also where $s$ without an argument is said to be non-decreasing; in the kernel $k_s(g,h)$ and in its formula, $s$ is the noise; and in $x_s$ it is part of the name of the floor, not the depth, whereas $x_k$, $x_{\check k}$ and $x_{K-1}$, with a level as index, are the couplings of these levels. In Steps 4 and 6, $e$ is a tolerance, and the upper index $e$ of $f^e$ is the same tolerance. With a plaquette as argument, $\partial U(p)$ is the product of $U$ along the boundary of $p$ (\ref{support:150-2}); in Step 8, $\theta_0:=\theta_{(\mu,0)}$ and $\theta_c:=\theta_{(\mu,c)}$ are the reflections of the mirrors $(\mu,0)$ and $(\mu,c)$ of $T_l$; and in Step 10, $n_i$ and $n_k$ are the lengths of $X_i$ and $X_k$.

In this section the tags \textup{(a$^C$)} to \textup{(d$^C$)}, \textup{(S)} and \textup{(c1)} to \textup{(c4)} are those of \Cref{Y-H8-ThBC}; \textup{(F0)} to \textup{(F5)} are those of \Cref{Y-H8-CoF}; \textup{(R0)} to \textup{(R3)} those of \Cref{Y-H8-CoR7}; \textup{(D0)} to \textup{(D7)} those of \Cref{Y-H7-Df}; \textup{(J0)} to \textup{(J8)} those of \Cref{Y-H6-J}; and \textup{(G1)} to \textup{(G4)} those of item~3 of \S\ref{Y-H6-1}.

\emph{Item 2: the statement.}

\begin{theorem}[Truncated functions of the action density: bounds, rate and limit]\label{Y-H9-T}
\emph{Data.} The group is $G=\operatorname{SU}(2)$, the dimension is $3$, and the letters are those of item 1. The variables are $\mathsf b\ge x_s$, the side $l=N=2L^{m_l}$ with $m_l\ge m'$, so that $N$ is even and $M\mid N$, subject to \textup{(Tor)} $N\ge2Ms(L\mathsf b)$, and $K\ge1$; we put $\beta=L^K\mathsf b$ and $a_K=L^{-K}$. The torus is $T_l=\mathbb R^3/l\mathbb Z^3$, and the lattice $a_K\mathbb Z^3/l\mathbb Z^3\subset T_l$ of \ref{support:150-2} has spacing $a_K$; as a graph it is $\mathbb Z^3/NL^K\mathbb Z^3$, with the site $y$ of the graph placed at the point $a_Ky$. So $0$ is a site, $a_K\mathbb Z^3\subset a_{K'}\mathbb Z^3$ for $K'\ge K$, and $l/2=L^{K+m_l}a_K$. The plaquette based at $y$ and spanned by $e_\mu,e_{\mu'}$ has the barycentre $x(p)=a_K(y+(e_\mu+e_{\mu'})/2)$, as in item~1 of \S\ref{Y-H2-1}.

Let $a_p(U)=1-\frac12\operatorname{tr}\partial U(p)$, $\mu_{K,l,\beta}=Z^{-1}e^{-\beta\sum_pa_p}dU$ and $\mathcal O_K(f)=a_K^{-1}\sum_pf(x(p))a_p(U)$, where $x(p)$ is the barycentre, and let $S^T_{n;K,l,\beta}(f_1,\dots,f_n)=\langle \mathcal O_K(f_1);\dots;\mathcal O_K(f_n)\rangle$ denote the joint cumulant. These coincide with the objects of \Cref{support:150-2}, whose formulas are taken at the $L$ and $M$ of the table (\textup{(O4)}, \textup{(O7)}) rather than at the $L$ and $M$ in the base Wilson model. Beyond the table and the data, the theorem has no hypothesis. In this setting the following hold.
\begin{enumerate}
\item[\textup{(T1)}] For $n\ge2$, $d>0$ and $R_f>0$, the bound \[|S^T_{n;K,l,L^K\mathsf b}(f_1,\dots,f_n)|\le C_n\prod\|f_i\|_\infty\] holds for every $K\ge1$ and for all real bounded functions $f_i$ on $T_l$ whose supports are pairwise at distance $\ge d$, each support lying in a ball of radius $R_f$. Here $C_n:=\mathrm{Ser}_1(\mathsf b,n,d,R_f)+\log2\,(\tau_7/\mathfrak r(L^{j^*}\mathsf b))^n$ with $\tau_7=\max\{1,0.91C_{\mathrm{vol}}(N,\mathsf b)\}$; this constant is free of $K$ and locally bounded in $\mathsf b$.
\item[\textup{(T2)}] For real Lipschitz functions $f_i$ as in \textup{(T1)} and for $K'\ge K\ge K_T(\mathsf b,n,d)$, \[\FitDisplay{\dimexpr\linewidth-1pt\relax}{|S^T_{n;K',l,L^{K'}\mathsf b}(f)-S^T_{n;K,l,L^K\mathsf b}(f)|\le C'_nK\theta^K\prod\|f_i\|_\sharp,}\] where $\theta=\vartheta^{4/9}\le0.69903$ and \[\FitDisplay{\dimexpr\linewidth-1pt\relax}{C'_n:=4P_*\vartheta^{-j^*-(10/9)\Lambda}+(C_{\mathsf b}/3)\mathrm{Ser}_2(4+1.36C_5^\#/K_5)+3C_{\mathrm{vol}}(N,\mathsf b)C_{\mathsf b}(1.15+0.46C_5^\#/K_5)(4C_{\mathrm{vol}}(N,\mathsf b)/\mathfrak r(L^{j^*}\mathsf b))^n.}\] Since $K_T$ and $C'_n$ are locally bounded in $\mathsf b$, the bound is uniform on $I_N:=\{\mathsf b\ge x_s:\text{(Tor)}\}$, which is either empty or one compact interval $[x_s,\mathsf b_N]$. The threshold $K_T$ does not depend on $N$.
\item[\textup{(T3)}] For real Lipschitz tuples with pairwise disjoint supports, the limit $S^{(l,\mathsf b)}_n(f):=\lim_KS^T_{n;K,l,L^K\mathsf b}(f)$ exists as $K\to\infty$ through all integers $K$, and $|S^{(l,\mathsf b)}_n-S^T_{n;K}|\le C'_nK\theta^K\prod\|f_i\|_\sharp$ for $K\ge K_T$. The limit is real and symmetric, it is $n$-linear on every product of spaces of functions supported in pairwise disjoint closed sets, one set for each factor, and it satisfies \textup{(T1)}. For continuous entries with pairwise disjoint supports the limit exists as well, and it coincides there with the extension given by \textup{(T1)}.
\item[\textup{(T4)}] For tuples $f$ as in \textup{(T3)}, Lipschitz or continuous, the map $\mathsf b\mapsto S^{(l,\mathsf b)}_n(f)$ is continuous. If $K_j\to\infty$ and $a_{K_j}\beta_j\to \mathsf b$, then $S^T_{n;K_j,l,\beta_j}(f)\to S^{(l,\mathsf b)}_n(f)$. Together with \textup{(T6)}, both statements hold for every $\mathsf b>0$; and if $(K_j,\beta_j)$ is a sequence of data with $K_j\to\infty$ whose labels $a_{K_j}\beta_j$ remain in a compact subset of $(0,\infty)$, one can pass to a subsequence along which $S^T_{n;K_j,l,\beta_j}(X)$ converges simultaneously for all $n\ge2$ and all $X$.
\item[\textup{(T5)}] The family $S=(S^{(l,\mathsf b)}_n)_{n\ge2}$, taken on the separated tuples $\mathfrak S_n(T_l)$ of \Cref{support:150-3}, whose entries are $C^1$ with pairwise disjoint supports, forms a family of connected functions as defined in \ref{support:150-3} and written out in item 1 of this section; moreover:
\begin{enumerate}
\item[\textup{(i)}] $S$ is invariant under all translations of $T_l$ and under $W_3$;
\item[\textup{(ii)}] the centred functional $\sigma_S(X)=\sum_{\pi\in\mathrm{Part}_{\ge2}(n)}\prod_{J\in\pi}S_{|J|}(X_J)$ of \eqref{support:9-14} is, at every mirror $(\mu,c)$ of $T_l$, reflection positive as in \eqref{support:9-28};
\item[\textup{(iii)}] each $S_n$ is symmetric;
\item[\textup{(iv)}] there is exactly one distribution $\tilde S_n$ on $(T_l)^n$ off the diagonals such that $\tilde S_n(f_1\otimes\cdots\otimes f_n)=S_n(f_1,\dots,f_n)$ on $\mathfrak S_n^\infty(T_l)$; its order on compact sets is $\le3n+1$, and $\sigma_S$ is regular.
\end{enumerate}
\item[\textup{(T6)}] By \eqref{support:150-42} and \eqref{support:150-43}, the clauses \textup{(T3)} to \textup{(T5)} hold for every $\mathsf b>0$ and every side $l=2L^m$, $m\in\mathbb Z$. Moreover, for every $z$ with $L^z\mathsf b\ge x_s$, $m+z\ge m'$ and $L^zl\ge2Ms(L^{z+1}\mathsf b)$, \[\FitDisplay{\dimexpr\linewidth-1pt\relax}{|S^{(l,\mathsf b)}_n(f)|\le L^{zn}C_n(L^z\mathsf b,n,L^zd,L^zR_f,L^zl)\prod\|f_i\|_\infty,}\] and, for real Lipschitz $f_i$ and $K-z\ge K_T(L^z\mathsf b,n,L^zd)$, \[\FitDisplay{\dimexpr\linewidth-1pt\relax}{|S^{(l,\mathsf b)}_n(f)-S^T_{n;K,l,L^K\mathsf b}(f)|\le L^{zn}C'_n(L^z\mathsf b,n,L^zd,L^zR_f,L^zl)(K-z)\theta^{K-z}\prod_i\max\{\|f_i\|_\infty,40ML^{-z}\mathrm{Lip}f_i\};}\] this is the rate of \textup{(T3)}, with its constants taken at the dilated data.
\item[\textup{(T7)}] \Cref{Y-J-CorC} holds with the rate $\kappa_1$ for the $\mathsf b$ and $N$ of the data.
\item[\textup{(T7+)}] \emph{The jets of the laws.} In the setting of \textup{(T2)}, let $\|f_i\|_\sharp\le1$, let $S$ be a non-empty subset of $\{1,\dots,n\}$, and let $0\le s\le K$. Let $\rho^{[K],z}_k$ be the run that starts from $\rho^{[K]}_0\exp\{\sum z_i\mathcal O_K(f_i)\}$, and $\pi^{[K]}_{S,[s]}:=\partial_S[\rho^{[K],z}_{K-s}/Z_K(z)]_{z=0}$. Then $\pi^{[K]}_{S,[s]}$ is a real function; it lies in $L^1(dV)$ on the fields of $\mathbb Z^3/NL^s\mathbb Z^3$; and, for every bounded Borel $F$, $\int F\pi\,dV$ is the joint cumulant of $F$, evaluated at the drawn field of depth $s$, with the $\mathcal O_K(f_i)$, $i\in S$. For $K'\ge K\ge K_T$ with $K-s-1\ge K_0(L^{s+1}\mathsf b)$, \[\|\pi^{[K']}_{S,[s]}-\pi^{[K]}_{S,[s]}\|_{L^1}\le8C^{(s)}(1.15+0.46C_5^\# K_5^{-1}K)C_{L^s\mathsf b}\theta^{K-s}(4C^{(s)}/\mathfrak r(L^{\max(s+1,j^*)}\mathsf b))^{|S|},\] where $C^{(s)}=C_{\mathrm{vol}}(NL^s,L^s\mathsf b)$. Hence $\pi^{[K]}_{S,[s]}$ converges in $L^1$ as $K\to\infty$ through all integers $K$, and, for $s\ge1$, the kernel of the drawn step with the noise $aL^s\mathsf b$ carries the limit at depth $s$ to the limit at depth $s-1$.
\end{enumerate}
\end{theorem}

\begin{remark}
The proof uses the blocking factor and cell scale of the data. Equation~\ref{support:150-42} transports the estimates to other sides and labels.\end{remark}

\begin{proof}[Proof of \Cref{Y-H9-T}]

\emph{Step 0: normalisation and the generating function.} We claim three things. First, it suffices to prove \textup{(T1)} for $\|f_i\|_\infty\le1$ and \textup{(T2)} for $\|f_i\|_\sharp\le1$. Second, for such $f$ and for the function $Z_K$ defined below, we have $S^T_{n;K,l,\beta}(f)=\partial_1\cdots\partial_n\log Z_K(0)$, and $Z_K(z)$ is the integral of the weight of \Cref{Y-H4-U4}: $Z_K(z)=\int\exp\{-\beta\sum_p(1-\nu_p)s_p\}dU$, with $s_p=a_p$ and $\nu_p:=\mathsf b^{-1}\sum_iz_if_i(x(p))$. Third, the data of \Cref{Y-H6-J} are those of \textup{(T1)}, and the data of \Cref{Y-H7-Df} are those of \textup{(T2)}.

For fixed $K$ the joint cumulant is $n$-linear in $(f_1,\dots,f_n)$, and the right sides of \textup{(T1)} and \textup{(T2)} change in the same way when one entry is multiplied by a number. If the norm of an entry vanishes, that entry vanishes at all barycentres, and both sides are $0$. We may therefore divide each $f_i$ by its norm.

The function $Z_K(z):=\int\exp\{-\beta\sum_pa_p+\sum_iz_i\mathcal O_K(f_i)\}dU$ is entire in $z\in\mathbb C^n$, because its integrand is continuous and the group is compact; moreover $Z_K(0)>0$, and $\partial_1\cdots\partial_n\log Z_K(0)$ is the joint cumulant, by the classical link between cumulants and the logarithm of the generating function. As $\beta=L^K\mathsf b$ and $a_K^{-1}=L^K$, we get $\beta\nu_pa_p=L^K\sum_iz_if_i(x(p))a_p$, and the sum over $p$ gives $\sum_p\beta\nu_pa_p=\sum_iz_i\mathcal O_K(f_i)$; so the weight is $\exp\{-\beta\sum_p(1-\nu_p)s_p\}$, the one of \Cref{Y-H4-U4}\,\textup{(a)}. Since the supports are pairwise disjoint, at each $p$ at most one $f_i(x(p))\ne0$, and hence $|\nu|_\infty\le\max_i|z_i|/\mathsf b$.

The variables $\mathsf b\ge x_s$, $N=2L^{m_l}$ with $m_l\ge m'$, \textup{(Tor)} and $K\ge1$ make up the data of \Cref{Y-H4-U4}, of \Cref{Y-H6-J} and of \Cref{Y-H7-Df}. The kernels of the step are probability kernels (item~2 of \S\ref{Y-J-1}, in its part ``The step''), and the step is linear in the density; hence, by Fubini's theorem, $\int\rho_k^zdV=Z_K(z)$ at every level, for complex $z$ as well as for real $z$.

\emph{Step 1: the numbers of \textup{(O14)}.} We first prove a lemma on these numbers.

\begin{lemma}[Growth of the separations and convergence of the tail series]\label{Y-H9-LN}
In the setting of item 1, the following hold for all data $(\mathsf b,n,d,R_f)$ with $\mathsf b\ge x_s$, $n\ge2$, $d>0$ and $R_f>0$.
\begin{enumerate}
\item[\textup{(N1)}] $r(Lx)\le r(x)+\log L/960+1$ for $x>1$.
\item[\textup{(N2)}] $\mathrm{sep}^{[i+1]}\ge L\,\mathrm{sep}^{[i]}$ for all $i\ge0$; the number $j^*$ exists and is the least $j\ge1$ with $\mathrm{sep}^{[j-1]}\ge\mathrm{sep}_*$; $\mathrm{sep}^{[j^*-1+t]}\ge L^t\mathrm{sep}_*$ for $t\ge0$; and $\mathsf b\mapsto j^*$ is non-decreasing.
\item[\textup{(N3)}] For $t\ge1$, the term $j=j^*+t$ of $\mathrm{Tail}$ is $\le P_*\exp\{-(\kappa n\,\mathrm{sep}_*/4)L^t\}$.
\item[\textup{(N4)}] $\mathrm{Tail}(J)\le2P_*\exp\{-(\kappa n\,\mathrm{sep}_*/4)L^{J+1-j^*}\}$ for $J\ge j^*$, and $\mathrm{Ser}_1<\infty$; the numbers $\mathrm{Ser}_1$, $\mathrm{Tail}$ and $P_*$ depend neither on $K$ nor on $N$.
\item[\textup{(N5)}] $\mathrm{Ser}_1$, $P_*$ and $j^*$ are bounded on every compact set of $\mathsf b\in[x_s,\infty)$.
\end{enumerate}
\end{lemma}

\begin{proof}[Proof of \Cref{Y-H9-LN}]
For \textup{(N1)} we apply $\lceil u+v\rceil\le\lceil u\rceil+v+1$ to $r(Lx)$, with $v=20c_R\log L/M$; by \textup{(M2)} and $\kappa_1+4\ge11$, this $v$ satisfies $v\le\log L/(96\cdot11)\le\log L/960$.

For \textup{(N2)} we put $t_i=dL^i/M$ and $r_i=r(L^i\mathsf b)$, so that $\mathrm{sep}^{[i]}=\lfloor t_i\rfloor-2r_i-8$. Then \[\mathrm{sep}^{[i+1]}-L\,\mathrm{sep}^{[i]}=[\lfloor Lt_i\rfloor-L\lfloor t_i\rfloor]+2(Lr_i-r_{i+1})+8(L-1).\] The bracket is $\ge0$, since $L$ is an integer, and \textup{(N1)} gives $r_{i+1}-Lr_i\le\log L/960+1\le4(L-1)$; so the difference is $\ge0$. As $t_i$ grows geometrically while $r_i$ grows at most linearly in $i$, we have $\mathrm{sep}^{[i]}\to\infty$; hence some $\mathrm{sep}^{[j-1]}\ge\mathrm{sep}_*>0$, and from that index on the sequence grows by the factor $L$. Finally, $\check r(L^i\mathsf b)$ is non-decreasing in $\mathsf b$, so that $\mathrm{sep}^{[i]}$ is non-increasing and $j^*$ is non-decreasing in $\mathsf b$.

For \textup{(N3)} let $j=j^*+t$ with $t\ge1$. We bound the factors of the term $j$ one by one.
\begin{enumerate}
\item[\textup{(i)}] Since $p_0$ is increasing, $\mathrm{amp}(L^{j+1}\mathsf b)/\mathrm{amp}(L^{j^*+1}\mathsf b)=L^{t/2}p_0(L^{j^*+1}\mathsf b)/p_0(L^{j+1}\mathsf b)\le L^{t/2}$, and therefore $\mathfrak r(L^j\mathsf b)^{-n}\le L^{nt/2}\mathfrak r(L^{j^*}\mathsf b)^{-n}$.
\item[\textup{(ii)}] By \textup{(N1)}, $\check r(L^{j-1}\mathsf b)\le\check r(L^{j^*-1}\mathsf b)+t(\log L/960+1)$. Combined with $2\,\check r+3\ge9$, $L^t-1\ge t(L-1)$ and $9(L-1)\ge2(\log L/960+1)$, this yields $2\,\check r(L^{j-1}\mathsf b)+3\le L^t(2\,\check r(L^{j^*-1}\mathsf b)+3)$, and so $N_1^{[j]}\le L^{3t}N_1^{[j^*]}$.
\item[\textup{(iii)}] By \textup{(X17)}, $2C_gc_a\le C_g/16$.
\item[\textup{(iv)}] $m^{[j]}\ge n\,\mathrm{sep}^{[j^*-1+t]}/2\ge(n/2)L^t\mathrm{sep}_*$.
\end{enumerate}

Consequently the term is $\le P_*L^{(n/2+3)t}\exp\{-(\kappa n\,\mathrm{sep}_*/2)L^t\}$. One half of the exponent absorbs the power, because \[(\kappa n\,\mathrm{sep}_*/4)L^t=2\kappa n\log(4nL^{1/2})L^t\ge6nL^t\log L\ge2nt\log L\ge(n/2+3)t\log L,\] where the last inequality uses $n\ge2$.

For \textup{(N4)} put $c:=\kappa n\,\mathrm{sep}_*/4\ge24\log(8L^{1/2})>49$. The quotient of consecutive bounds is $\exp\{-cL^t(L-1)\}\le1/2$, whence $\sum_{t\ge T}\exp\{-cL^t\}\le2\exp\{-cL^T\}$; here we take $T=J+1-j^*$. Moreover, $\mathrm{Ser}_1$ consists of the $j^*$ terms with $j\le j^*$ together with $\mathrm{Tail}(j^*)$, and none of the entries of the terms involves $K$ or $N$.

For \textup{(N5)} we put $\hat\jmath:=\max j^*$ on a given compact set, see \textup{(N2)}. The terms with $j\le\hat\jmath$ are products of the following factors: $c_a(L^j\mathsf b)$, continuous in $\mathsf b$; $\mathfrak r(L^i\mathsf b)^{-n}$ at $i=\max(j,j^*(\mathsf b))\le\hat\jmath$, which is one of finitely many continuous functions of $\mathsf b$ (and is $\le\mathfrak r(L^{\hat\jmath}\mathsf b)^{-n}$ by \textup{(G3)} of item~3 of \S\ref{Y-H6-1}); the non-decreasing $N_1^{[j]}$; and a factor $\le1$. So these terms are bounded. The number $P_*$ is the value at $j'=j^*\le\hat\jmath$ of a product of this kind, and it is bounded as well. For $j>\hat\jmath$ the term is $\le P_*\exp\{-cL^{j-j^*}\}\le P_*\exp\{-cL^{j-\hat\jmath}\}$.
\end{proof}

\emph{Step 2: the bound \textup{(T1)}.} We claim that, for $\|f_i\|_\infty\le1$ and all $K\ge1$, \[|S^T_{n;K,l,L^K\mathsf b}(f)|\le\mathrm{Ser}_1+\log2\,(\tau_7/\mathfrak r(L^{j^*}\mathsf b))^n,\] with $\tau_7=\max\{1,0.91C_{\mathrm{vol}}(N,\mathsf b)\}$, and that this bound is free of $K$ and locally bounded in $\mathsf b$. By Step 0 and \Cref{Y-H6-J}\,\textup{(J7)}, \[S^T_n=\sum_{0\le k<K}\sum_Y\partial_1\cdots\partial_nv_k(Y;u_k(z))(0)+\partial_1\cdots\partial_n\log I[u_K(z)](0),\] where $I[u]:=\int\rho[u;\ell_K]dV$. We treat the levels first and the last integral afterwards.

\emph{The levels.} By \Cref{Y-H6-J}\,\textup{(J8)}, the term $k$ is at most $\mathfrak r_{\min(k,k_d)}^{-n}N_1(k)\,2C_gc_a(x_k)e^{-\kappa(m_k-1)}$; in terms of the depth $j=K-k$ of the level $k$, we have $x_k=L^j\mathsf b$ and $N_1(k)=N_1^{[j]}$. If $K\ge j^*$, then $k_d=K-j^*$ and $\mathfrak r_{\min(k,k_d)}=\mathfrak r(L^{\max(j,j^*)}\mathsf b)$; moreover $k<k_d$ exactly when $j>j^*$, so that $m_k=m^{[j]}$. If $K<j^*$, then $k_d=0$, every $m_k=1=m^{[j]}$ because $j\le K<j^*$, and $\mathfrak r_{\min(k,0)}=\mathfrak r_0=\mathfrak r_*\ge\mathfrak r(L^{j^*}\mathsf b)=\mathfrak r(L^{\max(j,j^*)}\mathsf b)$ by \textup{(G3)} of item~3 of \S\ref{Y-H6-1}. In both cases the term $k$ is bounded by the term $j$ of $\mathrm{Tail}$, where $j$ runs over $1,\dots,K$; therefore the sum over the levels is at most $\mathrm{Tail}(0)=\mathrm{Ser}_1$, which is finite by \textup{(N4)}.

\emph{The last integral.} For the last integral we use \Cref{Y-H8-ThBC} on $\mathbb Z^3/N\mathbb Z^3$, applied with $x=\mathsf b$ and with the pair $(\kappa,c_h)$, that is, with the rate $\kappa$ and the cap $c_h(x)$. Its data are met: $\mathsf b\ge x_s\ge e^2$; $N$ is even, $M\mid N$, and $N\ge2Ms(L\mathsf b)\ge2Ms(\mathsf b)$ because $s$ is non-decreasing; $\varepsilon(\mathsf b)\le1$ and $(c_a+c_h)(\mathsf b)\le1/2$ by \textup{(X16)}, \textup{(X17)}; and $\ell_K$ has range $\le\varrho(\mathsf b)$, with $|\ell_{K,y}(p)|\le K_*e^{-\lambda_*|y-p|}$, by Step~2\,\textup{(c)} of the proof of \Cref{Y-J-ThA}.

These data are available at every $K\ge1$. Indeed, Step~2 of the proof of \Cref{Y-J-ThA} is carried out under the data of that theorem, which include $K\ge K_0(\mathsf b)$, but the reasoning in its items \textup{(c)} and \textup{(e)} makes no use of $K_0$:
\begin{enumerate}
\item[\textup{(c)}] uses $\mathsf b\ge x_\flat$ and \textup{(X12)}, the geometric series over $1\le j\le K$ and \Cref{Y-A-Th29}\,\textup{(b)}, with $L\ge L_A$ odd, $aL\ge3/c_L$ and $N$ even, and it gives $\ell_k\in\mathcal R(b_*,K_*,\lambda_*,1/9)$ for $0\le k\le K$, with range $\le\varrho(x_k)$, since the definition of the rows $\ell_k$ cuts them at the radius $\varrho(x_k)$; here, by \textup{(O13)} of item~(1) of \S\ref{Y-H0-1}, $x_s$ serves as the $x_\flat$ of \S\ref{Y-J-1}, with \textup{(X12)}, so that $\mathsf b\ge x_s$ gives $\mathsf b\ge x_\flat$;
\item[\textup{(e)}] uses $N\ge2Ms(L\mathsf b)\ge2Ms(\mathsf b)$ and \Cref{Y-GB-P63}, applied with $x=\mathsf b\ge e^2$; this proposition is stated with its proof in item~2 of Step~15 of the proof of \Cref{Y-GB-ThB}.
\end{enumerate}

So both items hold for every $K\ge1$. The pointwise bound above is the decay clause in the definition of $\mathcal R$ (the second condition in the display of \Cref{Y-A-SetsOfRows}).

The family is $u_K(z)$ on $\mathrm{Pd}(\mathfrak r_*)$, as in \textup{(J6)}. Since $j^*\ge1$, we have $k_d\le K-1$, so that $u_K(z)=F_{\ell_{K-1}}(u_{K-1}(z))$ is an output of \Cref{Y-H-C9} as extended by \Cref{Y-H1-U1}\,\textup{(U)}; here the rate is $\kappa_{K-1}=\kappa$, while \Cref{Y-H-C9} itself has a single rate, namely $\kappa_1$. This is case~\textup{(ii)} of \textup{(b$^C$)}. The argument does not depend on this case: \textup{(F2)} of \Cref{Y-H8-CoF} is asserted and derived on the polydisc $\mathrm{Pd}(\mathfrak r_*/\tau_7)$, where \textup{(S)} holds, and \textup{(S)} implies case~\textup{(i)} by \Cref{Y-H8-ThBC}\,\textup{(c$^C$)}.

We have $u_K(0)=\bar u_K$, a system of members of the point $\rho^{[K]}_K$ of the class, by \Cref{Y-H4-U4}\,\textup{(h)} at $z=0$ and \Cref{Y-H-C9}\,\textup{(d)}. The norm $\|u_K(z)-\bar u_K\|^\sim$ is at most the norm $\|\cdot\|^{\mathrm{st}}$, with the rate $\kappa_K=\kappa$, of the same difference, and the latter is $\le0.91\cdot0.4^{K-k_d}\le0.91=:B_0$. Now \textup{(F2)}, applied with $r_i=\mathfrak r_*$ and $S=\{1,\dots,n\}$, gives \[|\partial_1\cdots\partial_n\log I[u_K(z)](0)|\le\log2\,(\tau_7/\mathfrak r_*)^n\le\log2\,(\tau_7/\mathfrak r(L^{j^*}\mathsf b))^n.\]

\emph{Local boundedness.} Finally, the right side of the claim is locally bounded in $\mathsf b$, by \textup{(N5)}, by the continuity of $\mathsf b\mapsto\mathfrak r(L^j\mathsf b)$ and of $\mathsf b\mapsto C_{\mathrm{vol}}(N,\mathsf b)$, and since $j^*$ is bounded on compact sets.

\emph{Step 3: the clause \textup{(T2)}.} We show that, for $\|f_i\|_\sharp\le1$ and $K'\ge K\ge K_T(\mathsf b,n,d)$, one has $|S^T_{n;K'}-S^T_{n;K}|\le C'_nK\theta^K$, where $C'_n$ is the constant of \Cref{Y-H9-T}; that $K_T$ and $C'_n$ are locally bounded in $\mathsf b$; and that $I_N$ is empty or a compact interval.

Since $K\ge K_T\ge K_0(\mathsf b)$, \Cref{Y-J-ThA}\,(i) gives $\Lambda>0$ and $1\le k_0\le K-1$. From $k_0<(5K+10\Lambda)/9+1$ we obtain $J:=K-k_0>\frac49K-\frac{10}9\Lambda-1\ge j^*+1$, the second inequality being \textup{(K1a)}. As $J$ is an integer, it is therefore at least $j^*+2$, and \Cref{Y-H7-Df}\,(D1)\,(a) then yields $k_0+2\le k_d$.

In \textup{(D6)}, the early sum of each run has the step $0\to1$ of that run as its term $0$. With $m=K'-K$, the identity \textup{(D6)}, written in its four groups of terms, reads \[\FitDisplay{\dimexpr\linewidth-1pt\relax}{S^T_{n;K'}-S^T_{n;K}=[\text{early levels }0\le k<m+k_0\text{ of }K']-[\text{early levels }0\le k<k_0\text{ of }K]+[\text{pairs }k_0\le k<K]+[\text{last}].}\]

\emph{The early levels.} By \textup{(D5)}, each of the two brackets of early levels is at most $\mathrm{Tail}(J)$, because a level and its partner have the same depth and $j^*$ depends on $(\mathsf b,n,d)$ alone. We put $T:=J+1-j^*\ge3$ and $c:=\kappa n\,\mathrm{sep}_*/4>49$; then \textup{(N4)} gives $\mathrm{Tail}(J)\le2P_*e^{-cL^T}$. Since $\log(1/\vartheta)\le0.8057$, we have $cL^T\ge T\log(1/\vartheta)$, and hence $e^{-cL^T}\le\vartheta^T$. On the other hand $\frac49K<J+1+\frac{10}9\Lambda=T+j^*+\frac{10}9\Lambda$, whence $\theta^K=\vartheta^{(4/9)K}\ge\vartheta^{T+j^*+(10/9)\Lambda}$. The two brackets together are thus at most $4P_*\vartheta^{-j^*-(10/9)\Lambda}\theta^K$.

\emph{The pairs.} By \textup{(D4)}, the bracket of the pairs is at most $\mathrm{Ser}_2(4+1.36C_5^\# K_5^{-1}K)\vartheta^{K-k_0}$. By \Cref{Y-J-ThA}\,(iv), $3\vartheta^{K-k_0}\le3\vartheta^{(4/9)K-(10/9)\Lambda-1}=C_{\mathsf b}\theta^K$. Since $K\ge1$, the bracket of the pairs is therefore bounded by $(C_{\mathsf b}/3)\mathrm{Ser}_2(4+1.36C_5^\#/K_5)K\theta^K$.

\emph{The last bracket.} Here we apply the second sentence of \Cref{Y-H8-CoF}\,(F3) at $x=\mathsf b$, with the rate $\kappa$, the cap $c_h(\mathsf b)$ and $r_i=\mathfrak r_*/2$, and with the data of Step~2. The first family is $u_K(z)$; by \Cref{Y-H6-J}\,(J6) it stays within $0.364$ of $\bar u_K$. The second family is the family $\mathfrak n_K(z)$ of \textup{(D1)} on $\mathrm{Pd}(\mathfrak r_*/2)$. As $K-1\ge k_d>k_0$, we have $\mathfrak n_{K-1}(z)\in M_2$, since by \textup{(D1)} the member at the level $k_d$ is the polynomial and every later member is an output of \Cref{Y-H-C9} lying in $M_1$; and $\mathfrak n_{K-1}(0)\in M_1$, by \textup{(D0)} at $z=0$ together with \textup{(D1)}\,\textup{(c)}, and by \Cref{Y-J-ThA}\,(iii).

We use \Cref{Y-H5-Af}\,(ii) with $B=2$ and $p=(\kappa,c_h)$ at $x=x_{K-1}=L\mathsf b\ge Lx_s\ge x_9^U$, where \textup{(SD)} holds at $(x_{K-1};\ell_{K-1})$ by \Cref{Y-H4-U4}\,(b), and we take the second family constant. This gives \[\|\mathfrak n_K(z)-\mathfrak n_K(0)\|^\sim\le\vartheta[2](2+1)\le0.28343\cdot3\le0.8503\le1,\] and \Cref{Y-H7-Df}\,(D1)\,(f) provides the bound $0.795$ besides. Hence $B_0=1$ and $\tau_7=\max\{1,C_{\mathrm{vol}}\}=C_{\mathrm{vol}}$.

The member $\mathfrak n_K(0)$ is real and lies in $M_1$ with the strong cap, and $\rho[\mathfrak n_K(0);\ell_K]$ is a positive multiple of $\rho^{[K']}_{K'}$; so $\rho^{[K']}_{K'}$ is a point of $C^{\rm su}_\flat(\mathsf b;\ell_K)$, and $\mathfrak n_K(0)$ is a system of members of it, as the second sentence of \textup{(F3)} requires. This follows from \textup{(D0)} at $z=0$ with \textup{(D1)}\,\textup{(c)} and from \Cref{Y-J-ThA}\,(iii). With the rate $\kappa$ and the cap $c_h(\mathsf b)$, the norm $\|\cdot\|^\sim$ is bounded by the distance $\mathsf D_{\mathsf b}$ with the rate $\kappa_K=\kappa$, and therefore \textup{(D3)} yields $\Delta\le(1.15+0.46C_5^\# K_5^{-1}K)C_{\mathsf b}\theta^K$, where $3.45/3=1.15$ and $1.36/3\le0.4534$. Now \textup{(F3)} gives \[|\text{last}|\le3C_{\mathrm{vol}}\Delta(2\tau_7)^n(2/\mathfrak r_*)^n\le3C_{\mathrm{vol}}C_{\mathsf b}(1.15+0.46C_5^\#/K_5)(4C_{\mathrm{vol}}/\mathfrak r(L^{j^*}\mathsf b))^nK\theta^K.\] No threshold is imposed on $\Delta$; consequently $K_T$ does not depend on $N$.

\emph{The sum.} Adding the bounds for the early levels, for the pairs and for the last bracket, we arrive at the bound $C'_nK\theta^K$.

\emph{Uniformity.} The quantities $\Lambda$, $C_{\mathsf b}$, $\mathfrak r(L^j\mathsf b)$ and $C_{\mathrm{vol}}(N,\mathsf b)$ depend continuously on $\mathsf b$; $j^*$ and $P_*$ are bounded on compact sets by \textup{(N5)}; and $\mathrm{Ser}_2$ and $K_T$ are bounded on compact sets by \textup{(D7)}. The function $s(x)=8\lceil20c_R\log x/M\rceil$ is a non-decreasing step function and is continuous from the left. Hence $\{\mathsf b:2Ms(L\mathsf b)\le N\}$ is a half line closed on the right, and either $I_N=[x_s,\mathsf b_N]$ or $I_N$ is empty.

\emph{Step 4: the clause \textup{(T3)}.} On real Lipschitz tuples with pairwise disjoint supports, the limit exists as $K\to\infty$ through all integers $K$, and the convergence has the rate of \textup{(T2)}. This limit is real and symmetric, it is $n$-linear on every product of spaces of functions whose carriers are closed and pairwise disjoint, and it obeys \textup{(T1)}. For continuous entries the limit exists as well.

Supports are closed and $T_l$ is compact, so the least distance between two of the pairwise disjoint supports is a number $d>0$; each support lies in the ball of radius $R_f=N/2$, which is $T_l$ itself. By \textup{(T2)}, the sequence $(S^T_{n;K})_K$ is Cauchy, because $K\theta^K\to0$; and letting $K'\to\infty$ in \textup{(T2)} gives the rate.

At every $K$ the cumulant is real, symmetric and $n$-linear. Let $A_1,\dots,A_n$ be closed sets that are pairwise at distance $\ge d$. On the product of the spaces of Lipschitz functions carried by these sets, all tuples are of the same type $(d;N/2)$; therefore these properties, and \textup{(T1)} as well, pass to the limit.

\emph{Continuous entries.} Let the $f_i$ be continuous and carried by $A_i$. We choose Lipschitz functions $f_i^e$ with $\|f_i-f_i^e\|_\infty\le e$ that are carried by the closed $(d/4)$-neighbourhood of $A_i$; such functions are obtained by mollification, and these neighbourhoods are $\ge d/2$ apart. A telescoping sum in the entries, $n$-linearity at finite $K$ and \textup{(T1)} for tuples of the type $(d/2;N/2)$ give, uniformly in $K$, \[|S^T_{n;K}(f)-S^T_{n;K}(f^e)|\le nC_n(d/2)\,e\max_i(\|f_i\|_\infty+e)^{n-1}.\] Consequently $(S^T_{n;K}(f))_K$ is a Cauchy sequence. Its limit is the extension of $S_n$ by continuity with respect to the sup norms, and it retains all the properties above.

\emph{Step 5: dilation.} \Cref{lem:dilation} gives the finite-cutoff identity \eqref{support:150-42} for every $\beta>0$ and $z\in\mathbb Z$. With $\beta=L^K\mathsf b=L^{K-z}(L^z\mathsf b)$ it also gives \eqref{support:150-43} and the equivalence of existence of the two limits.

\emph{Step 6: the clause \textup{(T4)}.} We show that $\mathsf b\mapsto S^{(l,\mathsf b)}_n(f)$ is continuous at every $\mathsf b>0$; that $S^T_{n;K_j,l,\beta_j}(f)$ converges whenever $K_j\to\infty$ and $a_{K_j}\beta_j\to \mathsf b$; and that every sequence $(K_j,\beta_j)$ with $K_j\to\infty$ and with labels $a_{K_j}\beta_j$ in a compact subset of $(0,\infty)$ has a subsequence along which $S^T_{n;K_j,l,\beta_j}(f)$ converges.

\begin{enumerate}
\item[\textup{(i)}] \emph{On $I_N$.} For fixed $K$ the function $\mathsf b\mapsto S^T_{n;K,l,L^K\mathsf b}(f)$ is continuous: it is a polynomial in the moments, and every moment is the quotient of two integrals, over a compact group, of integrands that depend continuously on $\mathsf b$, the denominator being positive. By \textup{(T2)} and \textup{(T3)} we have, for $K\ge\sup_{I_N}K_T$, \[\sup_{I_N}|S^{(l,\mathsf b)}_n-S^T_{n;K}|\le\sup_{I_N}C'_nK\theta^K\prod\|f_i\|_\sharp\to0.\] Being a uniform limit of continuous functions, $\mathsf b\mapsto S^{(l,\mathsf b)}_n(f)$ is continuous on $I_N$.

\item[\textup{(ii)}] \emph{A compact interval about any $\mathsf b>0$, for any side $l=2L^m$.} We choose $z$ with $L^z\mathsf b/2\ge x_s$, $m+z\ge m'$ and $L^zl\ge2Ms(2L^{z+1}\mathsf b)$. Such a $z$ exists, because $s(2L^{z+1}\mathsf b)=8\lceil20c_R((z+1)\log L+\log2\mathsf b)/M\rceil$ grows linearly in $z$. Then $[L^z\mathsf b/2,2L^z\mathsf b]\subset I_{L^zl}$.

By Step~5, $S^{(l,\mathsf b')}_n(f)=L^{zn}S^{(L^zl,L^z\mathsf b')}_n(f^{[z]})$ for every $\mathsf b'$ in the interval $[\mathsf b/2,2\mathsf b]$, and $f^{[z]}$ is Lipschitz, of the type $(L^zd;L^zR_f)$. Applied on the torus of side $L^zl$, part \textup{(i)} now gives the continuity at $\mathsf b$. In particular, the end points of $I_N$ and all $\mathsf b<x_s$ are covered.

\item[\textup{(iii)}] Let $K_j\to\infty$ and $\mathsf b_j:=a_{K_j}\beta_j\to \mathsf b$. For all large $j$ we have $\mathsf b_j\in[\mathsf b/2,2\mathsf b]$ and $K_j-z\ge\sup K_T$. By Step~5, \[S^T_{n;K_j,l,\beta_j}(f)=L^{zn}S^T_{n;K_j-z,L^zl,\beta_j}(f^{[z]}),\] where the label on the right is $L^z\mathsf b_j$. Therefore \[|S^T_{n;K_j,l,\beta_j}(f)-S^{(l,\mathsf b)}_n(f)|\le L^{zn}\sup C'_n(K_j-z)\theta^{K_j-z}\prod\|f_i^{[z]}\|_\sharp+|S^{(l,\mathsf b_j)}_n(f)-S^{(l,\mathsf b)}_n(f)|\to0,\] by \textup{(ii)} and by \textup{(T3)}, applied on the torus of side $L^zl$ uniformly for the labels in $[L^z\mathsf b/2,2L^z\mathsf b]$, a subset of $I_{L^zl}$, that is, for $\mathsf b_j\in[\mathsf b/2,2\mathsf b]$. The suprema of $K_T$ and of $C'_n$ are taken over this dilated interval.

\item[\textup{(iv)}] A sequence of labels in a compact subset of $(0,\infty)$ has a convergent subsequence $\mathsf b_j\to \mathsf b>0$. Along this one subsequence, \textup{(iii)} holds for all $n$ and all tuples at once: the subsequence does not depend on $n$ or on the tuple.

\item[\textup{(v)}] \emph{Continuous entries.} Let the $f_i$ be continuous with pairwise disjoint supports, let $f^e$ be as in Step~4, and take the interval $[\mathsf b/2,2\mathsf b]$ and the $z$ of \textup{(ii)}. Combining Step~5 with the $n$-linearity at finite $K$ and with \textup{(T1)} on the torus of side $L^zl$, we obtain \[|S^T_{n;K,l,L^K\mathsf b'}(f)-S^T_{n;K,l,L^K\mathsf b'}(f^e)|\le n\sup_{\mathsf b'\in[\mathsf b/2,2\mathsf b]}L^{zn}C_n(L^z\mathsf b',n,L^zd/2,L^zl/2,L^zl)\,e\max_i(\|f_i\|_\infty+e)^{n-1}\] for every $K\ge z+1$ and every $\mathsf b'\in[\mathsf b/2,2\mathsf b]$. The supremum is finite, since $C_n$ is locally bounded in $\mathsf b$, and the same bound holds in the limit $K\to\infty$. Hence \textup{(ii)}, \textup{(iii)} and \textup{(iv)} pass from $f^e$ to $f$ by a uniform limit.
\end{enumerate}

\emph{Step 7: the clauses \textup{(T5)\,(i)} and \textup{(iii)}: a family of connected functions; translations, $W_3$ and symmetry.} The family $S=(S^{(l,\mathsf b)}_n)_{n\ge2}$ on $\mathfrak S_n(T_l)$ is symmetric and, in the sense of \ref{support:150-3}, multilinear. Further, every translation $r$ of $T_l$ and every $r\in W_3$ satisfy $S_n(rX)=S_n(X)$, and consequently $\sigma_S$ is invariant.

A $C^1$ function on $T_l$ is Lipschitz, so \textup{(T3)} applies on $\mathfrak S_n(T_l)$. The open sets $O_1,\dots,O_n\subset T_l$ in the definition of multilinearity have pairwise disjoint closures, and every entry in $C^1_c(O_i)$ is supported in the closure of $O_i$. Applied to these $n$ closed sets, one for each factor, \textup{(T3)} shows that $S_n$ is $n$-linear on $C^1_c(O_1)\times\cdots\times C^1_c(O_n)$; hence $S_n$ is multilinear as \ref{support:150-3} requires, and it is symmetric by \textup{(T3)}.

\emph{At finite $K$.} Let $r$ be the map $x\mapsto r^\circ x+v$ with $r^\circ\in W_3$ and $v\in a_K\mathbb Z^3$. It maps the lattice onto itself, each plaquette onto a plaquette and each barycentre to a barycentre. It does not, however, map base sites to base sites, and this is why the barycentre is required here. With $(rU)(b)=U(r^{-1}b)$ we have $a_p(rU)=a_{r^{-1}p}(U)$, since $\operatorname{tr}g=\operatorname{tr}g^{-1}$ on $\operatorname{SU}(2)$ and the trace around a plaquette depends neither on the corner at which it starts nor on its orientation. As the Haar measure and the action are invariant, $\mathcal O_K(f)(rU)=\mathcal O_K(f\circ r)(U)$, and $S^T_{n;K}(rX)=S^T_{n;K}(X)$.

\emph{The limit.} For $K'\ge K$ we have $a_K\mathbb Z^3\subset a_{K'}\mathbb Z^3$. So the identity holds for all large $K$ whenever $r^\circ\in W_3$ and $v$ lies in the union $A$ of the lattices $a_K\mathbb Z^3$; by \textup{(T3)} it then holds for $S_n$. The set $A$ is dense in $\mathbb R^3$.

\emph{All $v$.} Let $X$ be of the type $(d;R_f)$ and let $v$ be any vector of $\mathbb R^3$, in $A$ or not; let $v_j\in A$ with $|v-v_j|\le d/4$ and $v_j\to v$. We write $\tau_vX-\tau_{v_j}X$ as a telescoping sum. The entries of its $i$-th term are $\tau_{v_j}f_1,\dots,\tau_vf_i-\tau_{v_j}f_i,\dots,\tau_vf_n$; the entry in the $k$-th place is carried by the closed $(d/4)$-neighbourhood of $\tau_v\,\mathrm{supp}f_k$, and these neighbourhoods are $\ge d/2$ apart. By $n$-linearity and \textup{(T1)} for tuples of the type $(d/2;R_f+d/4)$, \[|S_n(\tau_vX)-S_n(\tau_{v_j}X)|\le nC_n\sqrt3|v-v_j|\max_i[\mathrm{Lip}f_i\prod_{k\ne i}\|f_k\|_\infty]\to0,\] while $S_n(\tau_{v_j}X)=S_n(X)$. Finally, $\sigma_S$ is a polynomial in the $S_{|J|}(X_J)$, and $(rX)_J=r(X_J)$.

\emph{Step 8: the clause \textup{(T5)\,(ii)}, reflection positivity of the centred functional.} For every mirror $\Sigma=(\mu,c)$ of $T_l$, with $\mathbb E_+(\Sigma)=\{c<y_\mu<c+l/2\}$, for all $X_1,\dots,X_r\in\mathfrak S(\mathbb E_+(\Sigma))$ and for all $\zeta\in\mathbb C^r$, we show that \[\sum_{i,k}\bar\zeta_i\zeta_k\,\sigma_S(\theta_\Sigma X_i\cdot X_k)\ge0.\]

\begin{enumerate}
\item[\textup{(a)}] \emph{Finite $K$ and $c=0$.} Both planes $y_\mu=0$ and $y_\mu=l/2=L^{K+m_l}a_K$ are planes of sites, and the reflection $\theta:y_\mu\mapsto-y_\mu$ maps the lattice onto itself and leaves each of the two planes pointwise fixed. We divide the bonds into three groups: $U_0$ consists of the bonds lying in the two planes, $U_+$ of the remaining bonds in $\{0\le y_\mu\le l/2\}$, and $U_-$ of the mirror images of the latter. Every plaquette lies either in one of the two planes or in exactly one of the two closed halves, so that \[\sum_pa_p=A_0(U_0)+A_+(U_0,U_+)+A_+(U_0,\theta U_-).\]

If the barycentre of a plaquette lies in the open side, then the plaquette lies in the closed half. Indeed, the $\mu$-coordinate of the barycentre is then either a multiple of $a_K$ strictly between the two planes, for a plaquette spanned by the two directions other than $\mu$ and lying in a plane $y_\mu=\mathrm{const}$ of sites, or an odd multiple of $a_K/2$, for a plaquette lying between two neighbouring planes of sites. Hence, for $f$ carried by $\mathbb E_+$, $\mathcal O_K(f)$ is a function of $(U_0,U_+)$.

We put $G:=\sum_k\zeta_k\prod_{f\in X_k}(\mathcal O_K(f)-\langle \mathcal O_K(f)\rangle)$. By Step~7, $\mathcal O_K(f)(\theta U)=\mathcal O_K(\theta f)(U)$ and $\langle \mathcal O_K(\theta f)\rangle=\langle \mathcal O_K(f)\rangle$, both being real. Fubini's theorem therefore gives \[\sum\bar\zeta_i\zeta_k\,\sigma_{K,l,\beta}(\theta X_i\cdot X_k)=\langle\overline{G\circ\theta}\,G\rangle=Z^{-1}\int dU_0\,e^{-\beta A_0}\bigl|\int dU_+\,e^{-\beta A_+}G\bigr|^2\ge0.\] In this computation $\theta$ reverses every bond of direction $\mu$: at every $y$, \[(\theta U)(\langle y,y+a_Ke_\mu\rangle)=U(\langle\theta y-a_Ke_\mu,\theta y\rangle)^{-1};\] since the Haar measure is invariant under inversion, the substitution $U_-\mapsto\theta U_-$ carries $dU_-$ to $dU_+$; and $A_0$, $A_+$ and the $\mathcal O_K(f)$ are real. Here $\sigma_{K,l,\beta}(X)=\langle\prod(\mathcal O_K(f_i)-\langle \mathcal O_K(f_i)\rangle)\rangle$.

\item[\textup{(b)}] The tuple $\theta X_i\cdot X_k$ is separated: the entries of $X_k$ are pairwise disjoint in $\mathbb E_+$, the entries of $\theta X_i$ are pairwise disjoint in $\theta\mathbb E_+$, and $\theta\mathbb E_+$ does not meet $\mathbb E_+$. By \ref{support:150-3}, $\sigma_{K,l,\beta}=\sigma_S$ at $S=(S^T_{n;K,l,\beta})_{n\ge2}$: centering makes the singleton cumulants zero, and all other cumulants are evaluated at separated tuples.

\item[\textup{(c)}] \emph{The limit.} The number $\sigma_S(Y)$ is a finite polynomial in the $S_{|J|}(Y_J)$, and each factor converges by \textup{(T3)}. A limit of non-negative reals is non-negative.

\item[\textup{(d)}] \emph{All $c$.} Let $\tau_vy=y+v$ with $v=ce_\mu$. Then $\theta_{(\mu,c)}=\tau_v\theta_{(\mu,0)}\tau_{-v}$, since $\tau_v\theta_0(y-v)=y-2(y_\mu-c)e_\mu$, and $\tau_{-v}$ maps $\mathbb E_+(\mu,c)$ onto $\mathbb E_+(\mu,0)$. By Step~7, $\sigma_S(\theta_cX_i\cdot X_k)=\sigma_S(\theta_0X'_i\cdot X'_k)$ with $X'=\tau_{-v}X\in\mathfrak S(\mathbb E_+(\mu,0))$, and so the claim follows from \textup{(c)}.
\end{enumerate}

\emph{Step 9: the clause \textup{(T5)\,(iv)}, regularity.} Let $n\ge2$ and $\mathbb E=T_l$, and let $\mathbb E^n_{\ne}$ be the set of points of $\mathbb E^n$ whose entries are pairwise different. We show that $\mathbb E^n_{\ne}$ carries exactly one distribution $\tilde S_n$ with $\tilde S_n(f_1\otimes\cdots\otimes f_n)=S_n(f_1,\dots,f_n)$ on the $C^\infty$ separated tuples, that this distribution is of order $\le3n+1$ on compact sets, and that $\sigma_S$ is regular in the sense of \ref{support:150-3}.

\emph{The local construction.} Let $\xi\in\mathbb E^n_{\ne}$, let $\delta:=\min_{i\ne k}|\xi_i-\xi_k|/4$, the distance being the sup distance, let $B_i$ be the open ball with centre $\xi_i$ and radius $\delta$, and let $Q=\prod B_i$. We fix $\chi_i\in C^\infty(T_l)$ with $0\le\chi_i\le1$, equal to $1$ on $B_i$ and carried by the closed ball of radius $3\delta/2$ about $\xi_i$. These closed balls are $\ge\delta$ apart. Let $S_n^{\mathbb C}$ be the extension of $S_n$ that is $n$-linear over $\mathbb C$ and is defined on entries that are continuous, complex valued and carried by these balls. Splitting every entry into its real and imaginary parts, we get from \textup{(T1)} and \textup{(T3)} \[|S_n^{\mathbb C}(g_1,\dots,g_n)|\le2^nC_n(\delta)\prod\|g_i\|_\infty.\]

For $F\in C_c^\infty(Q;\mathbb C)$ let $\hat F(k)$, $k\in(\mathbb Z^3)^n$, be its Fourier coefficients on $(T_l)^n$, where $e_k(x)=e^{2\pi ik\cdot x/l}$. For $k\ne0$ we integrate by parts $3n+1$ times in a coordinate with $|k_a|=|k|_\infty$ and obtain $|\hat F(k)|\le(l/2\pi|k|_\infty)^{3n+1}\sup|\partial_a^{3n+1}F|$; for $k=0$ we have $|\hat F(0)|\le\sup|F|$, and this term is covered by the maximum below. The number of indices with $|k|_\infty=t$ is at most $6n(2t+1)^{3n-1}$, and $\sum_tt^{-3n-1}(2t+1)^{3n-1}\le3^{3n-1}\sum_tt^{-2}$. Hence \[\sum_k|\hat F(k)|\le c(n,l)\max_{|\alpha|\le3n+1}\sup|\partial^\alpha F|.\] As $F=(\chi_1\otimes\cdots\otimes\chi_n)F$, we define \[D_Q(F):=\sum_k\hat F(k)S_n^{\mathbb C}(\chi_1e_{k_1},\dots,\chi_ne_{k_n}).\] Then $D_Q$ is linear, and $|D_Q(F)|\le2^nC_n(\delta)c(n,l)\max_{|\alpha|\le3n+1}\sup|\partial^\alpha F|$.

\emph{On products.} Let $f_i\in C_c^\infty(B_i)$, and let $F$ be their tensor product. Then $\hat F(k)=\prod\hat f_i(k_i)$, and the series $\sum_k\hat f_i(k)\chi_ie_k$ converges uniformly to $\chi_if_i=f_i$. Using $n$-linearity and the continuity in the sup norms, we conclude that $D_Q(f_1\otimes\cdots\otimes f_n)=S_n(f_1,\dots,f_n)$.

\emph{Uniqueness on $Q$.} Given $F\in C_c^\infty(Q)$, we take $\chi'_i\in C_c^\infty(B_i)$ equal to $1$ on the $i$-th projection of $\mathrm{supp}F$. The sums $(\chi'_1\otimes\cdots\otimes\chi'_n)\sum_{|k|\le R}\hat F(k)e_k$ are finite linear combinations of products of functions in $C_c^\infty(B_i;\mathbb C)$; they are all carried by a single compact subset of $Q$, and they converge to $F$ in every $C^m$. On that compact set every distribution has finite order, and so it is determined by its values on products.

\emph{Gluing.} The sets $Q$ cover $\mathbb E^n_{\ne}$. Each intersection $Q\cap Q'$ is again a product of open sets whose closures are pairwise disjoint; the argument for uniqueness works on it, so that $D_Q=D_{Q'}$ there. We now use the gluing theorem (L.~H\"ormander, \emph{The Analysis of Linear Partial Differential Operators I}, Theorem~2.2.4): let open sets $X_i$ cover $X$, and for each $i$ let $u_i$ be a distribution on $X_i$, such that $u_i=u_j$ on $X_i\cap X_j$ for all $i,j$; then $X$ carries a unique distribution $u$ that restricts to $u_i$ on each $X_i$. Applied in the flat charts of $(T_l)^n$, it yields exactly one distribution $\tilde S_n$ with $\tilde S_n=D_Q$ on every $Q$.

Any compact set is covered by finitely many sets $Q$, and with a partition of unity we get the order $\le3n+1$, with a constant that depends on the compact set. Now let a $C^\infty$ separated tuple be given, with separation $d$. A partition of unity splits each $f_i$ into pieces carried by small balls, namely open balls of sup radius $d/4$ about points $\xi_i$ of $\mathrm{supp}f_i$. At such a point $\xi$ we have $\delta=\min|\xi_i-\xi_k|/4\ge d/4$, so that every product of pieces is carried by a single $Q$. Since both sides are $n$-linear, $\tilde S_n(f_1\otimes\cdots\otimes f_n)=S_n(f)$. If another distribution has this property, then on every $Q$ it coincides with $D_Q$ on products; so it is $\tilde S_n$.

\emph{The functional $\sigma_S$.} We put $\sigma_n:=\sum_{\pi\in\mathrm{Part}_{\ge2}(n)}\bigotimes_{J\in\pi}\tilde S_{|J|}$, each tensor product being restricted to the open subset $\mathbb E^n_{\ne}$ of $\prod_{J\in\pi}\mathbb E^J_{\ne}$; and $\sigma_1=0$, because $\mathrm{Part}_{\ge2}(1)$ is empty. Each $\sigma_n$ is thus a distribution on $\mathbb E^n_{\ne}$. For $X=(f_1,\dots,f_n)\in\mathfrak S_n^\infty(\mathbb E)$ and each $\pi$, the tensor product over $J\in\pi$ takes at the product function $f_1\otimes\cdots\otimes f_n$ the value $\prod_{J\in\pi}\tilde S_{|J|}(\otimes_{i\in J}f_i)=\prod_{J\in\pi}S_{|J|}(X_J)$; hence \[\sigma_n(f_1\otimes\cdots\otimes f_n)=\sum_{\pi\in\mathrm{Part}_{\ge2}(n)}\prod_{J\in\pi}S_{|J|}(X_J)=\sigma_S(X),\] and $\sigma_S$ is regular in the sense of \ref{support:150-3}.

\emph{Step 10: the clause \textup{(T6)}.} For every $\mathsf b>0$ and every side $l=2L^m$, $m\in\mathbb Z$, the clauses \textup{(T3)}, \textup{(T4)} and \textup{(T5)} are true, with the bound and the rate that \Cref{Y-H9-T}\,(T6) states.

We take $z$ with $L^z\mathsf b\ge x_s$, $m+z\ge m'$ and $L^zl\ge2Ms(L^{z+1}\mathsf b)=16M\lceil20c_R\log(L^{z+1}\mathsf b)/M\rceil$; every large $z$ will do, as the right side is linear in $z$. Then $(L^zl,L^z\mathsf b)$ satisfies the conditions imposed on the side and on $\mathsf b$ in the paragraph ``Data'' of \Cref{Y-H9-T}, and by Step~5 every clause is carried back; the map $f\mapsto f^{[z]}$ is a bijection of the tuples.

\begin{enumerate}
\item[\textup{(T3)}] The map $f\mapsto f^{[z]}$ preserves Lipschitz, $C^1$ and continuous entries as well as disjoint supports; $\|f^{[z]}\|_\infty=\|f\|_\infty$; and the type $(d;R_f)$ turns into $(L^zd;L^zR_f)$, which gives the bound. As to the rate, Step~5 at $\beta=L^K\mathsf b=L^{K-z}(L^z\mathsf b)$ shows that $S^T_{n;K,l,L^K\mathsf b}(f)$ and $S^{(l,\mathsf b)}_n(f)$ equal $L^{zn}$ times the corresponding terms at $(K-z,L^zl,L^z\mathsf b,f^{[z]})$; there \textup{(T3)} holds for $K-z\ge K_T(L^z\mathsf b,n,L^zd)$, and $\mathrm{Lip}f^{[z]}=L^{-z}\mathrm{Lip}f$. Two admissible values of $z$ lead to the same number, as \eqref{support:150-43} is an identity.

\item[\textup{(T4)}] This is the content of parts \textup{(ii)} to \textup{(v)} of Step~6, which are formulated there for every $\mathsf b>0$ and every side.

\item[\textup{(T5)}] We have $(\tau_vf)^{[z]}=\tau_{L^zv}f^{[z]}$, and $W_3$ commutes with the dilation; under the dilation, the mirror $(\mu,c)$ is mapped onto the mirror $(\mu,L^zc)$ of the torus of side $L^zl$, and its open side onto the open side of the image, whose extent in the direction $\mu$ is $L^zl/2$, half the side of that torus. In $\sigma_S$ every partition of $n$ contributes $\prod_JL^{z|J|}=L^{zn}$; hence $\sigma_{S^{(l,\mathsf b)}}(X)=L^{zn}\sigma_{S^{(L^zl,L^z\mathsf b)}}(X^{[z]})$, and for the form of Step~8 we absorb the factor $L^{z(n_i+n_k)}$ into $\zeta_iL^{zn_i}$, a positive multiple. For the distributions, let $\Psi_z:(T_l)^n_{\ne}\to(T_{L^zl})^n_{\ne}$ be the componentwise dilation. With $D'$ the distribution constructed in Step~9 at the dilated data, define
\[
 D(F):=L^{zn}D'(F\circ\Psi_z^{-1}),
 \qquad F\in C_c^\infty((T_l)^n_{\ne};\mathbb C).
\]
Composition with $\Psi_z^{-1}$ sends compact off-diagonal supports to compact off-diagonal supports and satisfies $\partial^\alpha(F\circ\Psi_z^{-1})=L^{-z|\alpha|}(\partial^\alpha F)\circ\Psi_z^{-1}$. Hence $D$ has order at most $3n+1$ on compact sets. On tensor products this formula is exactly the dilation identity for $S_n$; the uniqueness argument of Step~9 therefore identifies $D$ with $\tilde S_n$ at $(l,\mathsf b)$.
\end{enumerate}

\emph{Step 11: the clauses \textup{(T7)} and \textup{(T7+)}.} \Cref{Y-J-CorC} holds for the $\mathsf b$ and the $N$ of the data of \Cref{Y-H9-T}, and \textup{(T7+)} holds.

\begin{enumerate}
\item[\textup{(T7)}] The number $x_s$ satisfies all of \textup{(O13)} of item~1 of \S\ref{Y-J-1}, evaluated at $\kappa_1$; and the conditions $\mathsf b\ge x_s$, $N$ even, $M\mid N$ and \textup{(Tor)} are the hypotheses of \Cref{Y-J-ThA} on $\mathsf b$ and $N$. To the conditions just named, the table of item~\textup{(1)} of \S\ref{Y-H0-1} only adds further ones. Therefore \Cref{Y-J-CorC} holds.

\item[\textup{(T7+)}] We put $k=K-s$ and $\check k=k-1$. By \Cref{Y-J-ThA}\,(iv), the condition $K-(s+1)\ge K_0(L^{s+1}\mathsf b)$ means that $\check k>k_0$.

\begin{enumerate}
\item[\textup{(0)}] \emph{The type of each object.} This part holds for every $K\ge1$, every $0\le s\le K$ and any finite number of real bounded $f_i$. Here $k$ ranges over all the levels $0\le k\le K$, whereas from part \textup{(1)} on we take $k=K-s$ again. We have $\rho^{[K]}_0=e^{-\beta\sum_pa_p}\le1$ and $|\mathcal O_K(f_i)|\le6L^K(NL^K)^3\|f_i\|_\infty$, because $0\le a_p\le2$ and there are $3(NL^K)^3$ plaquettes. Hence $\rho_0^z=\rho_0\exp\{\sum z_i\mathcal O_K(f_i)\}$ is continuous in $U$, entire in $z$ at every $U$, and bounded for $z$ in a compact set uniformly in $U$.

Next, the step from the level $k$ has the kernel $\prod_ck_{ax_k}(V'(c),\mathsf M(U)(c))$. It is Borel in $(V',U)$, since $k_s$ is continuous on $G\times G$ and $\mathsf M$ is a Borel function of $U$ (item~2 of \S\ref{Y-J-1}, in its part ``The step''). Its values lie between $0$ and a bound that depends only on $ax_k$ and on the number of coarse bonds, because $k_s(g,h)\le1/\int_Ge^{-s(1-\cos\operatorname{ang}\,g)}dg$.

By induction on $k$ it follows that $\rho_k^z(V)$ is Borel in $V$, entire in $z$ at every $V$, and bounded for $z$ in a compact set uniformly in $V$; and, by Cauchy's formula and Fubini's theorem, $\partial_S$ may be taken under the integrals of the steps and under $\int F\,\cdot\,dV$. The function $Z_K$ is entire, and $Z_K(0)>0$ by Step~0; hence $\pi^{[K]}_{S,[s]}$ is a bounded Borel function on a compact group, so it lies in $L^1(dV)$. For this we use Cauchy's estimate on a polydisc on which $|Z_K(z)|\ge Z_K(0)/2$, and the fact that a derivative is a limit of Borel functions. At real $z$ the functions $\rho_k^z$ and the number $Z_K(z)$ are real, as the $f_i$ are real; so are their derivatives at $0$.

The kernels do not depend on the sources. Hence, for every bounded Borel function $F$, \[\int F\rho^z_{K-s}dV/Z_K(z)=\langle F(V_{K-s})e^{\sum z_i\mathcal O_K(f_i)}\rangle/\langle e^{\sum z_i\mathcal O_K(f_i)}\rangle,\] where $\langle\ \rangle$ denotes the expectation under the joint law of the bare field $V_0=U$, whose law is $\mu_{K,l,\beta}$, and of the fields $V_1,\dots,V_K$ drawn in the run at $z=0$. This is the derivative $d/dt$ at $t=0$ of $\log\langle e^{tF(V_{K-s})+\sum z_i\mathcal O_K(f_i)}\rangle$, and differentiating it by $\partial_S$ at $z=0$ yields the joint cumulant of $F(V_{K-s})$ and the $\mathcal O_K(f_i)$, $i\in S$.

\item[\textup{(1)}] \emph{Families at the level $k$.} If $\check k\le k_d$, let $u^*_{\check k}$ and $\mathfrak n^*_{\check k}$ be the polynomials $u^{\mathrm{pol}}_{\check k}$ and $\mathfrak n^{\mathrm{pol}}_{\check k}$; if $\check k>k_d$, let them be the families $u_{\check k}(z)$ and $\mathfrak n_{\check k}(z)$ of \Cref{Y-H6-J}\,(J6) and \Cref{Y-H7-Df}\,(D1). We consider them on $\mathrm{Pd}(r)$, where $r:=\mathfrak r_{\min(\check k,k_d)}/2=\mathfrak r(L^{\max(s+1,j^*)}\mathsf b)/2$; note that $K\ge K_T$ gives $K>j^*$. By \textup{(J5)}, \textup{(J6)} and \textup{(D1)}, both lie in $M_2(x_{\check k};\kappa_{\check k},c_h)$. We put $\mathbb U(z):=F_{\ell_{\check k}}(u^*_{\check k}(z))$ and $\check N(z):=F_{\ell_{\check k}}(\mathfrak n^*_{\check k}(z))$. Since $x_{\check k}=L^{s+1}\mathsf b\ge Lx_s\ge x_9^U$, and since \textup{(SD)} holds at $(x_{\check k};\ell_{\check k})$ by \Cref{Y-H4-U4}\,(b), \Cref{Y-H5-Af}\,(i) shows that these are families in $M_1(x_k)$ with the strong cap, and outputs of \Cref{Y-H-C9}.

\item[\textup{(2)}] \emph{The distance.} For $\check k\le k_d$, \textup{(D2)} gives $\sup\mathsf D^{\mathrm{in}}(u^*,\mathfrak n^*)\le\Delta_{\check k}\le3.45\vartheta^{\check k-k_0}+1.36\Omega^{\#}$. For $\check k>k_d$, \Cref{Y-H5-Af}\,(ii), applied level by level with $B=2$, gives the bound $\vartheta[2]^{\check k-k_d}\Delta_{k_d}$, where $\vartheta[2]\le\vartheta$. Moreover $\Omega^{\#}=C_5^\# K_5^{-1}K\,0.3683^{K-k_0}\le C_5^\# K_5^{-1}K\vartheta^{K-k_0}$. In both cases, therefore, $\sup\mathsf D^{\mathrm{in}}\le(3.45+1.36C_5^\# K_5^{-1}K)\vartheta^{\check k-k_0}$. A further application of \Cref{Y-H5-Af}\,(ii) gives \[\Delta:=\sup_{\mathrm{Pd}(r)}\mathsf D_{x_k}(\mathbb U,\check N)\le(3.45+1.36C_5^\# K_5^{-1}K)\vartheta^{K-s-k_0}\le(1.15+0.46C_5^\# K_5^{-1}K)C_{L^s\mathsf b}\theta^{K-s},\] since $3\vartheta^{K-k_0-s}\le3\vartheta^{(4/9)K-s-(10/9)\Lambda-1}$.

\item[\textup{(3)}] \emph{The choice $B_0=1$.} By \textup{(J5)}, where $0.7/2+0.21/4=0.4025$, and by \textup{(J6)}, we have $\sup\|\mathbb U(z)-\mathbb U(0)\|\le\vartheta[2]\cdot0.455$; and $\sup\|\check N(z)-\check N(0)\|\le\vartheta[2](2+1)\le0.8503$.

\item[\textup{(4)}] \emph{The true jets.} On $\mathrm{Pd}_0$ we have $\rho^{[K],z}_k=\exp\{\sum_{j\le k}V_j(z)\}\rho[u_k(z);\ell_k]$ with the honest family, by \Cref{Y-H4-U4}\,(h). In this part and in the following ones, $I[u]=\int\rho[u;\ell_k]\,dV$, with the rows $\ell_k$ of the level $k$. By Step~0, $\int\rho_k^z\,dV=Z_K(z)$ at every level; hence, near $z=0$, where $Z_K(z)\ne0$ because $Z_K(0)>0$, we have $Z_K(z)=\exp\{\sum_{j\le k}V_j(z)\}I[u_k(z)]$, so that $I[u_k(z)]\ne0$ as well, and $\rho^z_k/Z_K(z)=\rho[u_k(z);\ell_k]/I[u_k(z)]$. By \textup{(J1)} if $\check k\le k_d$, by \textup{(J6)} if $\check k>k_d$, and by \Cref{Y-H5-LJ}, the values of $\mathbb U(z)$ and those of the honest $u_k(z)$ agree modulo the ideal $(z_i^2)$. For fixed $V$, $\rho[u;\ell](V)$ is an entire function of a finite number of values of $u$; so the two densities agree modulo the ideal. By bounded convergence, each coefficient of $I$ is the integral of the corresponding coefficient of $\rho$; a quotient whose denominator does not vanish at $z=0$ preserves congruences; and $\partial_S$ at $0$ annihilates the ideal. Therefore \[\pi^{[K]}_{S,[s]}=\partial_S[\rho[\mathbb U(z);\ell_k]/I[\mathbb U(z)]](0).\] For $K'$ the same identity holds with $\check N$ in place of $\mathbb U$, by \textup{(D1)}: the coefficients of $\check N$ are jets of a system of members of the long run with sources, taken over the rows $\ell_k^{[K]}$.

\item[\textup{(5)}] \emph{\Cref{Y-H8-ThBC} at $x=L^s\mathsf b$ on $\mathbb Z^3/NL^s\mathbb Z^3$, with the rate $\kappa_{\check k}$ and the cap $c_h(L^s\mathsf b)$.} The inequality $NL^s\ge2Ms(L^s\mathsf b)$ and the bound of the rows hold by Step~2\,(e), (c) of the proof of \Cref{Y-J-ThA}, and the floors hold by \textup{(X16)} and \textup{(X17)}. The first members $\mathbb U(0)=\bar u_k$ and $\check N(0)=\mathfrak n_k(0)$ are systems of members of points of one and the same class: for $\bar u_k$ by \Cref{Y-H4-U4}\,(h), and for $\mathfrak n_k(0)$ by \textup{(D0)} at $z=0$ and \textup{(D1)}\,\textup{(c)} together with \Cref{Y-J-ThA}\,(ii), (iii). Here $\tau_7=C^{(s)}$. We denote the two normalised densities by $p_2$ and $p_3$.

By parts \textup{(1)}, \textup{(3)} and the first paragraph of \textup{(5)}, the two families satisfy \textup{(F0)} and the hypotheses of the second sentence of \textup{(F3)}, with $B_0=1$ and $\tau_7=C^{(s)}$. Their integrals and normalized densities are therefore holomorphic on the smaller polydiscs specified in \Cref{Y-H8-CoF}.

\item[\textup{(6)}] \emph{The estimate for the jets.} Part \textup{(4)} identifies the true jets with those of the normalized densities. Apply \Cref{Y-H8-CoF}\,\textup{(F5)}, including its assertion for every $\Delta$, to obtain
\[
\|\pi^{[K']}_{S,[s]}-\pi^{[K]}_{S,[s]}\|_1
\le4C^{(s)}\Delta(2\tau_7/r)^{|S|}.
\]
Since $2\tau_7/r=4C^{(s)}/\mathfrak r(L^{\max(s+1,j^*)}\mathsf b)$, part \textup{(2)} gives the claimed bound of \textup{(T7+)} (with $4$ in place of its $8$).

\item[\textup{(7)}] \emph{The limit and consistency.} The bound tends to zero with $K$, uniformly for $K'\ge K$, so completeness of $L^1$ gives the limit. By \Cref{Y-H8-CoR7}\,\textup{(R3)(ii)}, the step at noise $aL^s\mathsf b$ maps $\pi^{[K]}_{S,[s]}$ to $\pi^{[K]}_{S,[s-1]}$, is independent of $K$ and is a contraction on $L^1$. Passing to the limits at the two fixed depths proves consistency, as in \textup{(R3)(iv)} there.

\end{enumerate}
\end{enumerate}
\end{proof}

\subsection{Estimates for separated sources}\label{Y-P1-1}

For $Q$ separated sources, the radii and decay rates can be chosen independently of $Q$. Linearity bounds the combined single jets by the sup norm of the combined source. A mixed jet can contribute to a polymer only through the sources within its reach, and a polymer meeting many separated source blocks must be large. These two facts give the estimates in \Cref{Y-P1-Jo}. Its vacuum bounds and the last-field integral estimate use the volume factor only once in \Cref{Y-P1-Cor}.

\begin{remark}[Item B: the range of depths] Each clause of \Cref{Y-P1-Jo} that is applied in \S\ref{Y-P6-1} holds without an upper bound on the depth $K-k$, on $x_k$ or on $\varrho$. \Cref{Y-P1-Jo}\,(J5$^\circ$)\,(d) and~(J8$^\circ$)\,(a$'$) assume $1\le k\le K$, $\varrho\ge1$ and a single lower bound, $\mathrm{sep}_k\ge8\log(4\varrho L^{1/2})$. \Cref{Y-P1-Jo}\,(lin), (J2$^\circ$) to~(J4$^\circ$), (J5$^\circ$)\,(a$'$) and~(J7$^\circ$)\,(iii) hold for all $0\le k\le K$, while \textup{(J1$^\circ$)} and \textup{(J7$^\circ$)}\,(ii), which concern the step from $k$ to $k+1$, hold for all $0\le k<K$, and \Cref{Y-P1-Jo}\,(L0) holds for every $K\ge1$. So they can be used at every depth $j\ge j_*(\mathsf b)$, at the level $k=K-j\ge1$, where $j_*(\mathsf b)$ is the depth of \S\ref{Y-P6-1}, set among its definitions: $j_*(\mathsf b)$ is the least $j\ge1$ with $L^j\mathsf b\ge\max\{x^\#,y(\mathsf b)^{4/3}\}$, where $x^\#$ is that of \S\ref{Y-P6-1} and $y(\mathsf b)=\mathsf b\max\{c_{\mathrm{vol}}(\mathsf b)^{1/3},L^{1-m'}\}$.\end{remark}

\emph{Item C: the data and the strong units.}

\begin{definition}[One run with $Q$ separated sources]\label{Y-P1-DATAo} We take the data of item~1 of \S\ref{Y-H6-1}: the group $G=\operatorname{SU}(2)$ in dimension $3$; $\mathsf b\ge x_s$; $N=2L^{m_l}$ with $m_l\ge m'$; the condition \textup{(Tor)}, $N\ge2Ms(L\mathsf b)$; and $K\ge1$. As there, $\beta=x_0=L^K\mathsf b$, $x_k=L^{K-k}\mathsf b$ and $a_K=L^{-K}$, with $\mathbb T_k$, the fitted grids and the rows $\ell_k$. In the letters of \S\ref{Y-00-1}, $Q\ge1$ is any number, $\check\delta>0$ and $\check\rho>0$. The sources $g_1,\dots,g_Q$ are bounded functions on $T_l=\mathbb R^3/N\mathbb Z^3$, real or complex, with $\|g_t\|_\infty\le1$. The sets $\{g_t\ne0\}$ lie pairwise at sup distance $\ge\check\delta$, and each of them lies within sup distance $\check\rho$ of one point.

The letters $n$, $f_i$, $d$, $R_f$ of \S\ref{Y-H6-1} become $Q$, $g_t$, $\check\delta$, $\check\rho$ here. We let $z\in\mathbb C^Q$, and $S$ stands for a non-empty subset of $\{1,\dots,Q\}$. If no $g_t$ is $0$ and $\check\delta\le N$, then $Q\le(N/\check\delta)^3$; a number of sources proportional to the volume is thus the general case. The objects $\mathcal O_K$, $Z_K(z)$, $\nu(z)$, the honest family $u_k^{\mathrm{hon}}(z)$ of \textup{(o1)} at the level $k$, $\delta_S(k)$, $u_k^{\mathrm{pol}}$, $S_t(k)$, $\check r$, $\mathrm{amp}$, $\mathrm{amp}_W$, the norms, $v_k$, $V_0^+$ and $I[u]$ are those of item~1 and of \textup{(o1)} to \textup{(o8)} in item~2 of \S\ref{Y-H6-1}.\end{definition}

\begin{definition}[Strong units at a rate]\label{Y-P1-UNITS} Let $\varkappa$ be a real number. The strong unit at $\varkappa$ at $(X,V)$ is $c_a(x_k)e^{-\varkappa(\#X-1)}$ for the $a$-member and $c_h(x_k)e^{-\varkappa(\#X-1)}$ for the $h$-member, and for the $w$-member it is \[\exp\{-c_\Xi x_k\varepsilon\sum_{p\in D(V)\cap X}\phi_p-\varkappa(|X|_c-1)\}.\]\end{definition}

\emph{Item D: numbers.}

\begin{definition}[Numbers that depend on $\mathsf b$, $\check\delta$ and $\check\rho$ alone]\label{Y-P1-D1} \textup{(D1)} The numbers below belong to the line \textup{(O14)} and are defined after the table and after $\mathsf b$. No floor depends on them, and none of them depends on $Q$, $N$, $K$ or the places of the sources.

We set $\mathrm{sep}^\circ:=8\log(8L^{1/2})$; this is the one numeral among them, and it is $\mathrm{sep}_*$ of \textup{(o4)} at $n=2$. Next, $\mathrm{sep}^{[i]}:=\lfloor\check\delta\,L^i/M\rfloor-2\,\check r(L^i\mathsf b)-2$, and $j^\circ=j^\circ(\mathsf b,\check\delta)$ is the least $j\ge1$ with $\mathrm{sep}^{[j-1]}\ge\mathrm{sep}^\circ$, which is $j^*$ of \textup{(o4)} at $n=2$ and $d=\check\delta$. We put $\mathfrak r^\circ(x):=3.5/\mathrm{amp}(Lx)$ and $r^\circ:=\mathfrak r^\circ(L^{j^\circ}\mathsf b)$, and \[n_{\mathrm{loc}}:=(\lfloor(2\,\check r(\mathsf b)+1)M/\check\delta\rfloor+1)^3,\qquad N_1^{[j]}:=(2\,\check\rho\,L^{j-1}/M+2\,\check r(L^{j-1}\mathsf b)+3)^3.\] For $j>j^\circ$ we put \[E^{[j]}:=N_1^{[j]}\,2C_gc_a(L^j\mathsf b)\,e^{1/(4L)}\exp\{-(\kappa-1/(4L))(\mathrm{sep}^{[j-1]}+1)\}.\] With these, \[C_W^\circ:=\tfrac12\sum_{j>j^\circ}E^{[j]}+6\sum_{1\le j\le j^\circ}N_1^{[j]}C_gc_a(L^j\mathsf b),\] and the letter of \S\ref{Y-00-1} is \[C_W(\mathsf b,\check\delta,\check\rho):=\sum_{j>j^\circ}2e\,N_1^{[j]}C_gc_a(L^j\mathsf b)e^{-(\kappa-1)(\mathrm{sep}^{[j-1]}-1)}+\sum_{1\le j\le j^\circ}6N_1^{[j]}C_gc_a(L^j\mathsf b).\] Term by term, $C_W^\circ\le C_W$. Finally, $e_h:=0.45\,N_1^{[1]}c_h(\mathsf b)$ and $G_1(\mathsf b,\check\delta,\check\rho):=\exp\{C_W+e_h\}/r^\circ$, and \[c_{\mathrm{vol}}(\mathsf b):=225/M^3+(3/2)\mathsf N_Zp_0(\mathsf b)^2+3\log(384/\varepsilon(\mathsf b)^3).\]\end{definition}

\begin{definition}[Quantities of a level]\label{Y-P1-D2} \textup{(D2)} The quantities of a level, defined below, depend on $K$ only through the depth $K-k$, with two exceptions: $k^\circ$ and $\mathfrak r^\circ_*$ depend on $K$ through $\min(K,j^\circ)$, as $k^\circ=K-\min(K,j^\circ)$ and $\mathfrak r^\circ_*=\mathfrak r^\circ(L^{\min(K,j^\circ)}\mathsf b)$. The sets $I_k(X)$ and the numbers $j_k(X)$ depend on the places of the sources, being defined on $\mathbb T_k$; the other quantities depend neither on $Q$ nor on $N$.

We set $\mathrm{sep}_k:=\mathrm{sep}^{[K-k]}$ and $k^\circ:=\max\{K-j^\circ,0\}$. We set $\check\kappa^\circ_k:=2\log(8L^{1/2})/\mathrm{sep}_{k+1}$ for $k<k^\circ$, and $\kappa^\circ_k:=\kappa+\sum_{k\le j<k^\circ}\check\kappa^\circ_j$. Next, $\mathfrak r^\circ_k:=\mathfrak r^\circ(x_k)$ and $\mathfrak r^\circ_*:=\mathfrak r^\circ_{k^\circ}$. Further, $I_k(X):=\{t:\ X^{+\check r(x_k)}\text{ meets }S_t(k)\}$, $j_k(X):=\#I_k(X)$ and $\mathfrak h_k:=\lceil(\mathrm{sep}_k+2)/2\rceil$. Finally, $N_1(k):=N_1^{[K-k]}$.\end{definition}

\emph{Item E: the statements.}

\begin{lemma}[The multi-affine part of a bounded holomorphic function]\label{Y-P1-LM} Let $r=(r_1,\dots,r_n)$ be a polyradius, and let $g$ be holomorphic on $\mathrm{Pd}(r)\subset\mathbb C^n$ with $|g|\le m$. Put $Mg(z):=\sum_Sz^S\partial_Sg(0)$, where the sum runs over all subsets $S$, the empty one included. Then $|Mg|\le m$ on $\{|z_i|\le r_i/2\}$, and no constant in this bound depends on $n$.\end{lemma}

\begin{proposition}[Jets of one run for any number of sources]\label{Y-P1-Jo} In the setting of Defs.~\ref{Y-P1-DATAo}, \ref{Y-P1-UNITS}, \ref{Y-P1-D1} and~\ref{Y-P1-D2}, the following hold.

\begin{enumerate} \item[\textup{(G$^\circ$)}] The clauses \textup{(G1)} to \textup{(G4)} of item~3 of \S\ref{Y-H6-1} hold in the following form. \begin{enumerate} \item[\textup{(G1), (G2)}] These two hold with $(\mathrm{sep}^\circ,j^\circ,k^\circ,\check\kappa^\circ,\kappa^\circ)$ in place of $(\mathrm{sep}_*,j^*,k_d,\check\kappa,\kappa_k)$: we have $\mathrm{sep}^{[i+1]}\ge L\,\mathrm{sep}^{[i]}$ and $k^\circ\le K-1$; for $k<k^\circ$ we have $\mathrm{sep}_{k+1}\ge L\,\mathrm{sep}^\circ$ and $\check\kappa^\circ_k\le1/(4L)$; further, $\kappa^\circ_{k^\circ}=\kappa$ and $\kappa\le\kappa^\circ_k\le\kappa+1/(4(L-1))$; and $j^\circ$ is finite, depends on $\mathsf b$, $\check\delta$, $L$, $M$ and $c_R$ alone, and does not decrease as $\mathsf b$ grows.

\item[\textup{(G3)}] We have $\mathfrak r^\circ_k\le\mathfrak r^\circ_{k+1}\le L^{1/2}\mathfrak r^\circ_k$, and $\mathfrak r^\circ_*\ge r^\circ$, with equality if $K\ge j^\circ$.

\item[\textup{(G4)}] The sets $S_t(k)^{+\check r(x_k)}$ lie pairwise at distance $\ge\mathrm{sep}_k+2$ in cubes.\end{enumerate}

\item[\textup{(G5$^\circ$)}] \begin{enumerate} \item[\textup{(a)}] If $\mathrm{sep}_k\ge-1$ and $j_k(X)\ge2$, then $\#X$ (for connected $X$), or $|X|_c$ (for $X\in\mathcal P$), is $\ge j_k(X)\mathfrak h_k\ge j_k(X)(\mathrm{sep}_k/2+1)$.

\item[\textup{(b)}] In every case, $j_k(X)\le n_{\mathrm{loc}}\#X$.\end{enumerate}

\item[\textup{(L0)}] \emph{Level $0$ in terms of $\nu$.} Let $z\mapsto\nu(z)$ be complex linear from $\mathbb C^m$ to the maps from the plaquettes of $\mathbb T_0$ to $\mathbb C$; let $R$ be a polyradius with $\sup_{\mathrm{Pd}(R)}|\nu(z)|_\infty\le\mathfrak r^\circ_0/\mathsf b$; and let $u_0^{\mathrm{pol}}(z)$ be the multi-affine part at $0$, taken value by value, of $u_0(\nu(z))$. Then $\sup_{\mathrm{Pd}(R)}\|u_0^{\mathrm{pol}}(z)-\bar u_0\|^{\mathrm{st},(\kappa_1)}\le0.014$ for every $K\ge1$ and every $m$, and this bound depends on no separation and on no disjointness of supports. The value of $u_0^{\mathrm{pol}}(z)$ at $X$ does not depend on $z_i$ unless some $\nu_p(\cdot)$ with $p\in X^{+r(x_0)}$ depends on $z_i$. If, besides, each $\nu_p(\cdot)$ depends on at most one variable, as it does in the setting of \Cref{Y-P1-DATAo}, then the value at $X$ depends on $z$ only through the $\nu_p(z)$ with $p\in X^{+r(x_0)}$. In the setting of \Cref{Y-P1-DATAo} we have $\nu(z)=\mathsf b^{-1}\sum_tz_tg_t(x(\cdot))$, and the polyradius with $R_t=\mathfrak r^\circ_0$ is one such $R$.

\item[\textup{(lin)}] This clause, \textup{(J0$^\circ$)}, \textup{(J3$^\circ$)}, \textup{(J4$^\circ$)} and the bound stated after \textup{(J4$^\circ$)} concern the single jets, for one complex bounded source of any support and for $0\le k\le K$; \textup{(J1$^\circ$)} and \textup{(J2$^\circ$)} hold in the setting of \Cref{Y-P1-DATAo}, for its $Q$ sources and every $S$. For a complex bounded $f$ on $T_l$ and $0\le k\le K$, let $\delta[f](k)$ be the derivative $\frac{d}{dz}\big|_{z=0}$, at the level $k$, of the honest family of \textup{(o1)} for the run with the single source $f$, in the variable $z\in\mathbb C$. Then $f\mapsto\delta[f](k)$ is complex linear, and $\delta_t(k)=\delta[g_t](k)$. Moreover, $\sum_tz_t\delta_t(k)=\delta[g_z](k)$ with $g_z:=\sum_tz_tg_t$, and $\|g_z\|_\infty\le\max_t|z_t|$.

\item[\textup{(J0$^\circ$)}] The quadratic $a$-member $s_{x_k}[\mathbb D_k[f]/\mathsf b]$ lies in $T^a(x_k)$, lives on single cubes, and has $|\cdot|_a\le\mathrm{amp}(x_k)\|f\|_\infty$ at every rate.

\item[\textup{(J3$^\circ$)}] We have $\delta^a[f](k)=s_{x_k}[\mathbb D_k[f]/\mathsf b]+\mathrm{rem}_k[f]$ with $\mathrm{rem}_k[f]\in T^a(x_k)$, and \[|\mathrm{rem}_k[f]|_a^{(\kappa_1)}\le2u_3(x_k)K_5\|f\|_\infty/\mathsf b\le0.002\,\mathrm{amp}(x_k)\|f\|_\infty.\]

\item[\textup{(J4$^\circ$)}] We have $|\delta^h[f](k)|_h^{(\kappa_1)},\ |\delta^w[f](k)|_w^{\mathrm{st},(\kappa_1)}\le0.12\,\mathrm{amp}(Lx_k)\|f\|_\infty$.

Hence $\|\delta[f](k)\|^{\mathrm{st},(\kappa_1)}\le0.2\,\mathrm{amp}(Lx_k)\|f\|_\infty$ for $0\le k\le K$.

\item[\textup{(J1$^\circ$)}] For $k<K$ and every $S$, the quantities $\delta_S(k+1)$, $\partial_SV_{k+1}(0)$ and $\partial_Sv_k(Y;u_k^{\mathrm{hon}}(z))(0)$ are the coefficients of $z^S$ of $F_{\ell_k}(u_k^{\mathrm{pol}}(z))$, $\mathsf V'_{\ell_k}[u_k^{\mathrm{pol}}(z)]$ and $v_k(Y;u_k^{\mathrm{pol}}(z))$, respectively.

\item[\textup{(J2$^\circ$)}] For $0\le k\le K$, $\delta_S(k)(X)=0$ unless $S\subset I_k(X)$.

\item[\textup{(AD)}] \begin{definition}[Admissible radii and rates]\label{Y-P1-AD} Admissible data consist of an integer $0\le k_1\le K$; radii $R_0,\dots,R_{k_1}>0$; and rates $\kappa_1\ge\varkappa_0\ge\dots\ge\varkappa_{k_1-1}\ge\kappa$ and a real $\varkappa_{k_1}\le\varkappa_{k_1-1}$ (for $k_1=0$: $\varkappa_0\le\kappa_1$), where we write $\check\Delta_k:=\varkappa_k-\varkappa_{k+1}$; they are subject to \textup{(i)} and \textup{(ii)}. \begin{enumerate} \item[\textup{(i)}] $R_k\le\mathfrak r^\circ_k$ for $k<k_1$; no condition is placed on $R_{k_1}$.

\item[\textup{(ii)}] For $k<k_1$: $\mathrm{sep}_{k+1}\ge-1$; $q_k:=(1+R_{k+1}/R_k)e^{-\check\Delta_k\mathfrak h_{k+1}}\le1$; and $0.364\,e^{\check\Delta_k}q_k^2\le0.03$.\end{enumerate}

The standard data are $k_1=k^\circ$, $R_k=\mathfrak r^\circ_k$ and $\varkappa_k=\kappa^\circ_k$.\end{definition}

The standard data are admissible, with $0.364\,e^{\check\Delta_k}q_k^2\le0.0057$.

\item[\textup{(J5$^\circ$)}] \begin{enumerate} \item[\textup{(a)}] For $0\le k\le K$ and $R>0$, \[\sup_{\mathrm{Pd}(R)}\|\sum_tz_t\delta_t(k)\|^{\mathrm{st},(\kappa_1)}\le0.2\,\mathrm{amp}(Lx_k)R=0.7R/\mathfrak r^\circ_k.\]

\item[\textup{(a$'$)}] \emph{Value by value.} At every $(X,V)$ and for each member, $\sum_t|\delta_t(k)(X,V)|\le0.2\,\mathrm{amp}(Lx_k)$ strong units at $\kappa_1$.

\item[\textup{(b)}] Let admissible data be given and let $1\le k\le k_1$. Then $\sup_{(X,\allowbreak V)}\sum_{|S|\ge2}R_k^{|S|}|\delta_S(k)(X,\allowbreak V)|\le0.03$ strong units at $\varkappa_k$, separately for every one of the three members.

\item[\textup{(c)}] For $k<k_1$, and for $k=k_1$ if $R_{k_1}\le\mathfrak r^\circ_{k_1}$ and $\varkappa_{k_1}\ge\kappa$, the map $u_k^{\mathrm{pol}}$ is a family on $\mathrm{Pd}(R_k)\subset\mathbb C^Q$ in $M_{1.91}(x_k;\varkappa_k,c_h)$, and $\sup\|u_k^{\mathrm{pol}}(z)-\bar u_k\|^{\mathrm{st},(\varkappa_k)}\le0.73$; for $k=0$ the bound is $\le0.014$, at $\kappa_1$.

\item[\textup{(d)}] \emph{Free last ratio.} Let $1\le k\le K$ and $\varrho\ge1$, with $\mathrm{sep}_k\ge8\log(4\varrho L^{1/2})$. Put $\check\kappa_i:=2\log(4L^{1/2})/\mathrm{sep}_{i+1}$ for $i<k-1$, and $\check\kappa_{\mathrm{last}}:=2\log(4\varrho L^{1/2})/\mathrm{sep}_k\le1/4$. Then the data $k_1=k$; $R_{k'}=\mathfrak r^\circ_{k'}$ for $k'<k$; $R_k=\varrho\,\mathfrak r^\circ_k$; $\varkappa_{k'}=\kappa+1/4+\check\kappa_{\mathrm{last}}+\sum_{k'\le i<k-1}\check\kappa_i$ for $k'<k$; and $\varkappa_k=\kappa+1/4$ are admissible, and all their rates lie in $[\kappa+1/4,\ \kappa+1/2+1/(4(L-1))]$. Consequently, \[\sup_{(X,V)}\sum_{|S|\ge2}(\varrho\,\mathfrak r^\circ_k)^{|S|}|\delta_S(k)(X,V)|\le0.03\] strong units at $\kappa+1/4$. This clause depends neither on $j^\circ$ nor on $k^\circ$.\end{enumerate}

\item[\textup{(J6$^\circ$)}] \emph{The representatives.} Put $u_k(z):=u_k^{\mathrm{pol}}(z)$ on $\mathrm{Pd}(\mathfrak r^\circ_k)$ for $k\le k^\circ$, and $u_{k+1}(z):=F_{\ell_k}(u_k(z))$ on $\mathrm{Pd}(\mathfrak r^\circ_*)$ for $k^\circ\le k<K$. The following holds for every $K\ge1$; if $K<j^\circ$, then $k^\circ=0$. \begin{enumerate} \item[\textup{(a)}] For $k^\circ<k\le K$, the $u_k$ are families in $M_1(x_k;\kappa,c_h)$ with the strong cap.

\item[\textup{(b)}] For $k^\circ\le k\le K$, $\sup_{\mathrm{Pd}(\mathfrak r^\circ_*)}\|u_k(z)-\bar u_k\|^{\mathrm{st},(\kappa)}\le0.91\cdot0.4^{k-k^\circ}$.

\item[\textup{(c)}] We have $u_k(0)=\bar u_k$ and $\partial_Su_k(0)=\delta_S(k)$ for every $S$.

\item[\textup{(d)}] Every representative has the property $(P_k)$ of Step~3 of the proof of \Cref{Y-H6-J}, and $u_K(z)$ is an output of $F_{\ell_{K-1}}$.

\item[\textup{(e)}] We have $|\delta_t^h(K)|_h^{\Sigma}\le0.0646/\mathfrak r^\circ_*$ at $(\kappa,c_h)$; so $\sum_X\sup|\delta_t^h(K)(X,\allowbreak \cdot)|\le0.0646\,\allowbreak N_1^{[1]}c_h(\mathsf b)/r^\circ$, and $e_t:=\sum_X\delta_t^h(K)(X,1)$ satisfies $r^\circ|e_t|\le0.065\,\allowbreak N_1^{[1]}c_h(\mathsf b)\le e_h$.\end{enumerate}

\item[\textup{(J7$^\circ$)}] Let $\equiv$ be the congruence of germs at $0$ in $\mathbb C^Q$ modulo $z_t^2$, as in \Cref{Y-H8-Congruence}. By any representatives we mean a choice, at every level $k'$, of a map $z\mapsto u_{k'}(z)$ defined on some neighbourhood of $0$ and mapping into $\mathfrak M(x_{k'})$, of norm $<16$, all of whose values, in the sense of item~0\,(4) of \S\ref{Y-H5-1}, are holomorphic in $z$, with $u_{k'}(0)=\bar u_{k'}$ and $\partial_Su_{k'}(0)=\delta_S(k')$ for all $S$. Among them are the honest family, $u_{k'}^{\mathrm{pol}}$ at every $k'\le K$, and those of \textup{(J6$^\circ$)}. We put $\gamma_t^{(k)}:=\sum_{k'<k}\partial_tV_{k'+1}(0)$ and $\gamma_t:=\gamma_t^{(K)}$, and \[W_{k'}(z):=\sum_Y\sum_{|S|\ge2}z^S\partial_Sv_{k'}(Y;u_{k'})(0);\] this is a polynomial, and it does not depend on the representatives. \begin{enumerate} \item[\textup{(i)}] For $0\le k\le K$, at every field $V$ of $\mathbb T_k$ and with any representatives, \[\rho_k^z(V)\equiv\exp\{\sum_{j\le k}V_j(0)+\sum_{k'<k}\sum_Y[v_{k'}(Y;u_{k'}(z))-v_{k'}(Y;\bar u_{k'})]\}\,\rho[u_k(z);\ell_k](V).\]

\item[\textup{(ii)}] We have $\partial_Sv_k(Y;u_k)(0)=0$ unless $S\subset I_{k+1}(Y)$.

\item[\textup{(iii)}] Moreover, \[\rho_k^z(V)\equiv\exp\{\sum_{j\le k}V_j(0)+\sum_tz_t\gamma_t^{(k)}+\sum_{k'<k}W_{k'}(z)\}\,\rho[u_k^{\mathrm{pol}}(z);\ell_k](V).\]

\item[\textup{(iv)}] With the representatives of \textup{(J6$^\circ$)}, the congruence \textup{(iii)} holds with $W^{(k)}$ in place of $\sum W_{k'}$ and $u_k(z)$ in place of $u_k^{\mathrm{pol}}(z)$, where \[W^{(k)}:=\sum_{k'<\min(k,k^\circ)}W_{k'}+\sum_{k^\circ\le k'<k}[w_{k'}-\sum_tz_t\partial_tw_{k'}(0)]\] and $w_{k'}(z):=\sum_Y[v_{k'}(Y;u_{k'}(z))-v_{k'}(Y;\bar u_{k'})]$, which is holomorphic on $\mathrm{Pd}(\mathfrak r^\circ_*)$. We put $W:=W^{(K)}$.

\item[\textup{(v)}] We have $Z_K(z)/Z_K(0)\equiv\exp\{\sum_tz_t\gamma_t+W(z)\}\,I[u_K(z)]/I[\bar u_K]$, and $I[\bar u_K]>0$.

\item[\textup{(vi)}] For every $S$, $\partial_S\log Z_K(0)=\sum_{k<K}\sum_Y\partial_Sv_k(Y;u_k)(0)+\partial_S\log I[u_K(z)](0)$.

\item[\textup{(vii)}] The number $\hat c_K(g_t):=\partial/\partial z_t\,\log\rho_K^z(1)$ at $z=0$ equals $\gamma_t+e_t$.\end{enumerate}

\item[\textup{(J8$^\circ$)}] \begin{enumerate} \item[\textup{(a)}] \emph{Free outer radius.} For admissible data, for $k<k_1$, for every $t$ and for every $R>0$ with $\lambda:=\log(1+R/R_k)/\mathfrak h_{k+1}\le\varkappa_k$, \[\sum_{S\ni t,\,|S|\ge2}R^{|S|}\sum_Y|\partial_Sv_k(Y;u_k^{\mathrm{pol}})(0)|\le N_1(k)\,2C_gc_a(x_k)\,e^{\lambda}\exp\{-(\varkappa_k-\lambda)(\mathrm{sep}_{k+1}+1)\}.\]

\item[\textup{(a$^\circ$)}] For the standard data, with $R=\mathfrak r^\circ_*$, $k<k^\circ$ and the depth $j=K-k>j^\circ$, we have $\lambda\le1/(4L)$, the right side of the bound in \textup{(a)} is $\le E^{[j]}$, and $\sup_{\mathrm{Pd}(\mathfrak r^\circ_*)}|W_k|\le(Q/2)E^{[j]}$.

\item[\textup{(a$'$)}] In the setting of \textup{(J5$^\circ$)}\,(d), for every $t$, \[\FitDisplay{\dimexpr\linewidth-1pt\relax}{\begin{aligned}
&\sum_{S\ni t,\,|S|\ge2}(\varrho\,\mathfrak r^\circ_k)^{|S|}\sum_{k'<k}\sum_Y|\partial_Sv_{k'}(Y;u_{k'}^{\mathrm{pol}})(0)|
\\
&\quad{}\le\sum_{1\le m\le k}N_1(k-m)\,2C_gc_a(x_{k-m})\,e^{1/4}
\\
&\quad{}\exp\{-(\kappa-1/4)(L^{m-1}\mathrm{sep}_k+1)\}
\\
&\quad{}\le E_{\mathrm{vac}},
\end{aligned}}\] where \[\FitDisplay{\dimexpr\linewidth-1pt\relax}{E_{\mathrm{vac}}:=\sum_{m\ge1}N_1^{[K-k+m]}\,2C_gc_a(L^{K-k+m}\mathsf b)\,e^{1/4}\exp\{-(\kappa-1/4)(L^{m-1}\mathrm{sep}_k+1)\}.}\] Of $K$ and $k$, $E_{\mathrm{vac}}$ depends on the depth $K-k$ alone.

\item[\textup{(b)}] \emph{Late levels.} For $k^\circ\le k<K$, on $\mathrm{Pd}(\mathfrak r^\circ_*)$, \[\sum_Y|v_k(Y;u_k(z))-v_k(Y;\bar u_k)|\le3QN_1(k)C_gc_a(x_k),\qquad|\partial_tw_k(0)|\le3N_1(k)C_gc_a(x_k)/\mathfrak r^\circ_*.\]

\item[\textup{(c)}] \emph{Late levels, anchored.} At these levels, $\sum_{S\ni t}(\mathfrak r^\circ_*/n_{\mathrm{loc}})^{|S|}\sum_Y|\partial_Sv_k(Y;u_k)(0)|\le2e\,N_1(k)C_gc_a(x_k)$.

\item[\textup{(d)}] We have $\sup_{\mathrm{Pd}(\mathfrak r^\circ_*)}|W(z)|\le QC_W^\circ\le QC_W$.\end{enumerate}\end{enumerate}\end{proposition}

\begin{lemma}[The tail of the vacuum numbers, and finiteness]\label{Y-P1-LNo} For $t\ge1$, \[E^{[j^\circ+t]}\le N_1^{[j^\circ]}(C_g/16)\,e^{1/(4L)}\exp\{-20L^t\log L\}.\] The numbers $C_W^\circ$ and $C_W$ are finite, depend neither on $K$ nor on $N$ nor on $Q$, and are bounded on compact sets of $\mathsf b\in[x_s,\infty)$.\end{lemma}

\begin{corollary}[The moment of $Q$ sources, with the volume constant to the first power]\label{Y-P1-Cor} The following holds for all $K\ge1$, all $Q\ge1$ and all complex $e_t$ with $r^\circ|e_t|\le e_h$, hence in particular for the $e_t$ of \Cref{Y-P1-Jo}\,(J6$^\circ$)\,(e). \begin{enumerate} \item[\textup{(a)}] We have \[|\langle\prod_{t\le Q}(\mathcal O_K(g_t)-\gamma_t-e_t)\rangle|\le G_1^Q(1+0.182\,C_{\mathrm{vol}}(N,\mathsf b))\le G_1^Qe^{c_{\mathrm{vol}}(\mathsf b)N^3}.\] Together with \Cref{Y-P1-Jo}\,(J7$^\circ$)\,(vii) this gives \[|\langle\prod_t(\mathcal O_K(g_t)-\hat c_K(g_t))\rangle|\le G_1^Q(1+0.182\,C_{\mathrm{vol}}(N,\mathsf b))\prod_t\|g_t\|_\infty\] for bounded $g_t$ of any norms. Here $C_{\mathrm{vol}}$ appears to the first power, and every sub-collection of the sources also falls under the setting of \Cref{Y-P1-DATAo}.

\item[\textup{(b)}] Let $F$ be a bounded Borel function on the fields of $\mathbb T_K$. In \[e^{-\sum_tz_t(\gamma_t+e_t)}\int F(V)\rho_K^z(V)\,dV/Z_K(0)\] the coefficient of $z_1\cdots z_Q$ is then, in modulus, $\le G_1^Q(1+0.182\,C_{\mathrm{vol}}(N,\mathsf b))\sup|F|$.\end{enumerate}\end{corollary}

\begin{remark}[Constants] Every numeral of this section is free of $K$, $k$, $N$, $\mathsf b$, $Q$, $\check\delta$, $\check\rho$ and $g$. The numerals $0.03$ and $0.0057$ depend on $L>1.2\cdot10^{13}$ of the table only through $(1+L^{-1/2})^2\le1.001$. The number $\mathfrak r^\circ$ depends on $\mathsf b$ and the depth; $\mathrm{sep}$, $j^\circ$, $\check\kappa^\circ$, $\kappa^\circ$ and $n_{\mathrm{loc}}$ depend on $(\mathsf b,\check\delta)$ and the depth; and $N_1$, $E^{[j]}$, $C_W^\circ$, $C_W$, $e_h$ and $G_1$ depend on $(\mathsf b,\check\delta,\check\rho)$. The number $Q$ occurs only as the dimension of $\mathrm{Pd}$, as the factor $Q/2$ in \Cref{Y-P1-Jo}\,(J8$^\circ$)\,(a$^\circ$), as the factor $Q$ in \textup{(J8$^\circ$)}\,(b) and~(d), and as the exponent in \Cref{Y-P1-Cor}.\end{remark}

\begin{remark}[Short distances, and the volume once] The proofs of all clauses of \Cref{Y-P1-Jo}, and that of \Cref{Y-P1-LNo}, work at short distances only. \Cref{Y-P1-LM}, on the multi-affine part on polyradii, is proved by function theory in one variable (Schwarz--Pick), one variable after the other.

The proof of \Cref{Y-P1-Cor} uses the volume once, through the $L^1$ bound of \Cref{Y-H8-ThBC}\,(c$^C$), case (ii), and it does not depend on any control at large distances; the step to a bound free of $N$, by the chessboard estimate, is taken in \S\ref{Y-P2-1}.\end{remark}

\begin{proof}
\emph{Proof 1: \Cref{Y-P1-LM} on the multi-affine part.} We show: if $r$ is a polyradius and $g$ is holomorphic on $\mathrm{Pd}(r)$ with $|g|\le m$, then $|Mg|\le m$ on $\{|z_i|\le r_i/2\}$, where $Mg=\sum_Sz^S\partial_Sg(0)$.

\emph{One variable.} Let $\phi$ be holomorphic on $|z|<r$ with $|\phi|\le1$, and put $a=|\phi(0)|$. By the Schwarz--Pick lemma, $|\phi'(0)|\le(1-a^2)/r$. Consequently, whenever $|z|\le r/2$, \[|\phi(0)+z\phi'(0)|\le a+(1-a^2)/2=1-(1-a)^2/2\le1.\]

\emph{Many variables.} Define $(M_ig)(z):=g|_{z_i=0}+z_i(\partial_ig)|_{z_i=0}$. On power series, $M_i$ deletes exactly the monomials $z^\mu$ with $\mu_i\ge2$ and leaves all others in place; therefore the operators $M_i$ commute with one another, and $M=M_n\cdots M_1$. The result is immediate for $m=0$. For $m>0$ put $g_0=g/m$ and $g_i=M_ig_{i-1}$.

We claim that $g_i$ is holomorphic, of modulus at most $1$, on \[A_i:=\{|z_j|\le r_j/2,\ j\le i\}\times\{|z_j|<r_j,\ j>i\};\] since it is affine in each of the first $i$ variables, it is holomorphic on a neighbourhood in these variables. For $i=0$ the claim is the hypothesis on $g$, as $A_0=\mathrm{Pd}(r)$.

To pass from $i-1$ to $i$, freeze all variables except $z_i$ at a point of $A_{i-1}$. Then $z_i\mapsto g_{i-1}$ is holomorphic on $|z_i|<r_i$ with modulus at most $1$, and $g_i$ arises from it as $\phi(0)+z_i\phi'(0)$ arises from $\phi$; the bound in one variable gives $|g_i|\le1$ for $|z_i|\le r_i/2$. At the end, $g_n=Mg/m$ on $A_n=\{|z_i|\le r_i/2\}$, which proves the lemma.

\emph{Proof 2: the clauses \textup{(G$^\circ$)} and \textup{(G5$^\circ$)}.} We show that \textup{(G1)} to \textup{(G4)} of item~3 of \S\ref{Y-H6-1} hold with the numbers of \textup{(G$^\circ$)}, and that $\#X\ge j_k(X)\mathfrak h_k$ and $j_k(X)\le n_{\mathrm{loc}}\#X$.
\begin{enumerate}
\item[\textup{(G$^\circ$)}] The numbers $\mathrm{sep}^\circ,j^\circ,k^\circ,\check\kappa^\circ_k,\kappa^\circ_k$ are the numbers $\mathrm{sep}_*,j^*,k_d,\check\kappa_k,\kappa_k$ of \textup{(o4)} in item~2 of \S\ref{Y-H6-1}, taken at $n=2$ and $d=\check\delta$.
\begin{enumerate}
\item[\textup{(G1), (G2)}] These two conditions are proved in Step~1\,(G1), (G2) of the proof of \Cref{Y-H6-J}. That argument involves $n$ only through $\mathrm{sep}_*>0$ and through $\check\kappa=(\mathrm{sep}_*/4)/\mathrm{sep}_{k+1}$, so it holds with the numbers of \textup{(G$^\circ$)}.

\item[\textup{(G3)}] By \textup{(G3)} of item~3 of \S\ref{Y-H6-1} and Step~0\,(b) of the proof of \Cref{Y-H6-J}, both taken at $x=x_k$, we have $\mathfrak r^\circ_{k+1}/\mathfrak r^\circ_k=\mathrm{amp}(Lx_k)/\mathrm{amp}(x_k)\in[5.01,L^{1/2}]$. The condition \textup{(G3)} is stated for $x\ge x_s$, and $x_k\ge\mathsf b\ge x_s$. If $K\ge j^\circ$, then $x_{k^\circ}=L^{j^\circ}\mathsf b$, and hence $\mathfrak r^\circ_*=\mathfrak r^\circ(x_{k^\circ})=r^\circ$, which is the case of equality in \textup{(G3)}. If $K<j^\circ$, then $k^\circ=0$ and $\mathfrak r^\circ_*=\mathfrak r^\circ(L^K\mathsf b)\ge r^\circ$, because the map $i\mapsto\mathfrak r^\circ(L^i\mathsf b)$ is decreasing.

\item[\textup{(G4)}] The argument of Step~1\,(G4) of the proof of \Cref{Y-H6-J} concerns one pair, and so it gives \textup{(G4)} pair by pair.
\end{enumerate}

\item[\textup{(G5$^\circ$)}]
\begin{enumerate}
\item[\textup{(a)}] The set $X^{+\check r(x_k)}$ meets $S_t(k)$ if and only if $X$ meets $S_t(k)^{+\check r(x_k)}$. Let $j=j_k(X)\ge2$. Then $X$, and likewise ${\rm hull}(X)$, meets the $j$ sets $S_t(k)^{+\check r(x_k)}$ with $t\in I_k(X)$; here ${\rm hull}(X)$ is connected, contains $X$ and has $\#{\rm hull}(X)=|X|_c$, by Step~1\,(C) of the proof of \Cref{Y-H6-J}. By \textup{(G4)}, these $j$ sets lie pairwise at distance at least $\sigma:=\mathrm{sep}_k+2\ge1$. Step~1\,(C) of the proof of \Cref{Y-H6-J} then bounds $\#X$ (for connected $X$) or $|X|_c$ (by way of ${\rm hull}(X)$) from below by $j\lceil\sigma/2\rceil=j\mathfrak h_k$.

\item[\textup{(b)}] Since $I_k(X)=\bigcup_cI_k(\{c\})$, the union running over the cubes $c$ of $X$, it suffices to treat $X=\{c\}$. For each $t\in I_k(\{c\})$ there is a plaquette $p_t$ with $g_t(x(p_t))\ne0$ that is based in the bare box of a cube of $c^{+\check r}$. The lifts of these base sites lie in a box whose side has $(2\check r(x_k)+1)ML^k$ sites. As $x(p)=a_K(z+(e_\mu+e_\nu)/2)$, the lifted barycentres spread, in each coordinate, over at most $a_K((2\check r+1)ML^k-1/2)<s_k:=(2\check r(x_k)+1)ML^{-(K-k)}$, so after a translation they lie in $[0,s_k)^3$. The sup distance of two lifts is at least their periodic distance, hence at least $\check\delta$. We cut each coordinate interval $[0,s_k)$ into $\lfloor s_k/\check\delta\rfloor+1$ half-open intervals of length $\check\delta$; two points in one cell of the resulting grid are closer than $\check\delta$, so each cell contains at most one of the lifted barycentres $x(p_t)$, and $\#I_k(\{c\})\le(\lfloor s_k/\check\delta\rfloor+1)^3$. Finally, Step~1\,(r) of the proof of \Cref{Y-H6-J} gives \[(2\check r(Lx)+1)/L\le(2\check r(x)+3+\log L/480)/L\le2\check r(x)+1,\] and therefore $s_k$ is largest at depth $0$.
\end{enumerate}
\end{enumerate}

\emph{Proof 3: the run in terms of $\nu$, and the clauses \textup{(lin)}, \textup{(J1$^\circ$)} and \textup{(J2$^\circ$)}.} We show that $U_k(\nu)$ is holomorphic on $|\nu|_\infty<15.99c_\nu\varepsilon(\beta)$, that $\delta[f](k)$ is linear in $f$ with $\sum z_t\delta_t(k)=\delta[g_z](k)$, and that \textup{(J1$^\circ$)} and \textup{(J2$^\circ$)} hold for every $Q\ge1$.
\begin{enumerate}
\item[\textup{1.}] \emph{Plaquette sources.} For a plaquette $p$ of $\mathbb T_0$, let $1_p$ be the function on $T_l$ equal to $1$ at $x(p)$ and to $0$ elsewhere. The point $x(p)=a_K(z+(e_\mu+e_\nu)/2)$ determines $(z;\mu\nu)$, so $\mathrm{supp}_01_p=\{p\}$. Thus the $1_p$ are complex bounded functions with pairwise disjoint $\mathrm{supp}_0$, and the part ``Sources'' of item~3 of \S\ref{Y-H4-1} admits them. The variables of these sources are $\zeta_p$, and we write $\nu_p=\zeta_p/\mathsf b$; then \[\mathrm{Pd}_0^o=\{|\zeta_p|<15.99/\mathrm{amp}_W\}=\{|\nu|_\infty<15.99c_\nu\varepsilon(\beta)\}=:B_\nu.\] By \Cref{Y-H4-U4}\,(h), the maps $U_0(\nu)=u_0(\nu)$, $U_{k+1}=F_{\ell_k}(U_k)$ and $\mathsf V_{k+1}=\mathsf V'[U_k]$ are defined on $B_\nu$, of norm at most $16$ in the strong cap, with all their values holomorphic in the finitely many $\nu_p$; moreover, writing $\rho_k^\nu$ for the $\rho_k^z$ of that clause for the plaquette sources $1_p$, with their variables $\zeta_p$ in the place of $z$, we have $\rho_k^\nu=\exp\{\sum_{j\le k}\mathsf V_j(\nu)\}\rho[U_k(\nu);\ell_k]$ at every field of $\mathbb T_k$.

\item[\textup{2.}] For a complex linear map $z\mapsto\nu(z)$ we put $u_k(z):=U_k(\nu(z))$ wherever $\nu(z)\in B_\nu$. In the setting of \Cref{Y-P1-DATAo}, this is the honest family of \textup{(o1)} in item~2 of \S\ref{Y-H6-1}: both arise by one and the same recursion from the same $u_0(\nu(z))$.

\item[\textup{3.}] Value by value, the chain rule gives \[\delta[f](k)=\mathsf b^{-1}\sum_pf(x(p))\,(\partial U_k/\partial\nu_p)(0),\] which is complex linear in $f$. The same rule gives $\delta_t(k)=\delta[g_t](k)$, and linearity then gives $\sum z_t\delta_t(k)=\delta[g_z](k)$. Since at each point at most one $g_t$ differs from $0$, and $\|g_t\|_\infty\le1$, we have $\|g_z\|_\infty\le\max|z_t|$.

\item[\textup{4.}] We prove \textup{(J1$^\circ$)}. By Step~2 of the proof of \Cref{Y-H6-J}, for $k<K$ the condition \textup{(SD)} of item~2 of \S\ref{Y-H-1} holds at $(x_k;\ell_k)$, through \Cref{Y-H4-U4}\,(b), and $x_k\ge Lx_s\ge\max\{x_9^U,x_3,x_4\}$; for $0\le k\le K$, moreover, $\bar u_k$ is real and lies in $M_1(x_k;\kappa_1,c_h)$ in the strong cap. \Cref{Y-H5-LS}\,(S1) and (S3), applied on $\mathrm{Pd}(1.99/\mathrm{amp}_W)$, put every $\delta_S(k)$, $0\le k\le K$, in $\mathfrak M(x_k)$, with finite strong norm at $\kappa_1$. Hence $u_k^{\mathrm{pol}}$ is a map defined near $0$, of norm less than $16$, that has the jets of the honest family. \Cref{Y-H5-LJ}, with $\Phi$ a value of $F_{\ell_k}$, of $\mathsf V'_{\ell_k}$ or of $v_k(Y;\cdot)$, now yields \textup{(J1$^\circ$)}; here $Q$ enters only as the dimension.

\item[\textup{5.}] We prove \textup{(J2$^\circ$)}. Let $k=0$. By \Cref{Y-H4-U4}\,(a4), $u_0(\nu)(X)$ depends only on the $\nu_p$ with $p\in X^{+r(x_0)}$. Now $\nu_p(z)$ depends on $z_t$ only if $g_t(x(p))\ne0$, that is, only if the cube that holds $p$ lies in $S_t(0)$; and $r\le\check r$. Hence $u_0(z)$, and with it $u_0^{\mathrm{pol}}$, has the property $(P_0)$ of the proof of \Cref{Y-H6-J}. For the passage from $k$ to $k+1$ with $k<K$ we use the claim of Step~3 of the proof of \Cref{Y-H6-J}, which comes from \Cref{Y-H1-U1}\,(f5) at $S=S_t(k)$ and $\hat\rho=\check r(x_k)$, applied one variable at a time, together with Step~1\,(R) of the same proof: $F_{\ell_k}(u_k^{\mathrm{pol}}(z))(Y)$ does not depend on $z_t$ if $Y^{+\check r(x_{k+1})}$ misses $S_t(k+1)$. By \textup{(J1$^\circ$)} its coefficient of $z^S$ is $\delta_S(k+1)(Y)$, which therefore vanishes whenever $t\in S$.
\end{enumerate}

\emph{Proof 4: the single jets.} Apply the single complex source estimate at the end of Step~6 of the proof of \Cref{Y-H6-J}. Restriction to one source variable gives the same first jet as the run with that source alone. The splitting and bounds on all levels give \textup{(J0$^\circ$)}, \textup{(J3$^\circ$)} and \textup{(J4$^\circ$)}.

\emph{Proof 5: the clause \textup{(L0)} at level $0$, for every $K\ge1$ and any number of variables.} We show that \[\sup_{\mathrm{Pd}(R)}\|u_0^{\mathrm{pol}}(z)-\bar u_0\|^{\mathrm{st},(\kappa_1)}\le0.014\] whenever $\nu(\cdot)$ is linear with $|\nu(z)|_\infty\le\mathfrak r^\circ_0/\mathsf b$ on $\mathrm{Pd}(R)$, and we prove the two sentences of \textup{(L0)} that say on what the value at $X$ depends.

Let $\theta:=\mathrm{amp}_W\mathfrak r^\circ_0=3.5\,\mathrm{amp}_W/\mathrm{amp}(Lx_0)$; by Step~0\,(b) of the proof of \Cref{Y-H6-J} and \textup{(G3)} of item~3 of \S\ref{Y-H6-1}, $\theta\le3.5/(51\cdot5.01)\le0.0136981$. Let $c:=15.99/\theta$; then $c\ge1167.3$. By linearity, $|\nu(z)|_\infty\le c\,\mathfrak r^\circ_0/\mathsf b=15.99/(\mathsf b\,\mathrm{amp}_W)=15.99c_\nu\varepsilon(\beta)$ on $\mathrm{Pd}(cR)$. There \Cref{Y-H4-U4}\,(a) to (a4) hold, and $u_0(\nu(z))$ has norm at most $16$.

\emph{The $a$-part.} By \Cref{Y-H4-U4}\,(a4), $E^a_\nu(\{c'\},V)=\sum_{p\in c'}[q_p+\beta\nu_ps_p]$ is affine in $\nu$, hence of degree at most $1$ in $z$, and so it coincides with its own multi-affine part. The $a$-part of $u_0^{\mathrm{pol}}(z)-\bar u_0$ is therefore $\beta\sum_{p\in c'}\nu_p(z)s_p$. Fix $z$ and let $f$ take the value $\mathsf b\nu_p(z)$ at $x(p)$ and the value $0$ off the barycentres. This is a single complex bounded source, with $\|f\|_\infty=\mathsf b|\nu(z)|_\infty$, and $\delta^a[f](0)(\{c'\})=(\beta/\mathsf b)\sum f(x(p))s_p$ equals the expression above. The $a$-part of \Cref{Y-H4-U4}\,(f) bounds it by $\mathrm{amp}_W\|f\|_\infty$ caps, and on $\mathrm{Pd}(R)$ we have $\mathrm{amp}_W\|f\|_\infty\le\mathrm{amp}_W\mathfrak r^\circ_0=\theta\le0.0137$.

\emph{The $h$-part} of $u_0^{\mathrm{pol}}(z)-\bar u_0$ is $0$.

\emph{The $w$-part.} Fix $(X,V)$, and let $\mathrm{cap}$ be the strong unit at $\kappa_1$ at $(X,V)$. The function $g(z):=w_{\nu(z)}(X,V)$ is entire by \Cref{Y-H4-U4}\,(a4), and $|g|\le\mathrm{cap}$ on $\mathrm{Pd}(cR)$ by \Cref{Y-H4-U4}\,(a2). The value of $u_0^{\mathrm{pol}}$ at $(X,V)$ is $Mg$, that of $\bar u_0$ is $g(0)$, and \Cref{Y-P1-LM} gives $|Mg|\le\mathrm{cap}$ on $\mathrm{Pd}(cR/2)$. Thus $Mg-g(0)$ vanishes at $0$ and is at most $2\,\mathrm{cap}$ on $\mathrm{Pd}(cR/2)$, and \Cref{Y-H5-LS}\,(S2) gives on $\mathrm{Pd}(R)$ \[|Mg-g(0)|\le2\,\mathrm{cap}\cdot\max|z_i|/(cR_i/2)\le(4/c)\,\mathrm{cap}\le0.003427\,\mathrm{cap}.\]

The strong norm at $\kappa_1$ of $u_0^{\mathrm{pol}}(z)-\bar u_0$ is the largest of the norms of its three parts, hence at most $0.0137\le0.014$. The coefficients of $u_0^{\mathrm{pol}}$ lie in $\mathfrak M(x_0)$, by \Cref{Y-H5-LS}\,(S3).

\emph{On what the value depends.} Fix $(X,V)$ and one of the members, and let $g(z)$ be the value of that member of $u_0(\nu(z))$ at $(X,V)$; the corresponding value of $u_0^{\mathrm{pol}}$ is $Mg$. By \Cref{Y-H4-U4}\,(a4), $g$ is a function of the $\nu_p(z)$ with $p\in X^{+r(x_0)}$. If none of these $\nu_p(\cdot)$ depends on $z_i$, then $g$ does not depend on $z_i$, so $\partial_Sg(0)=0$ for every $S\ni i$, and $Mg$ does not depend on $z_i$: this is the first sentence of \textup{(L0)}. Suppose besides that each $\nu_p(\cdot)$ depends on at most one variable, and let $T$ be the set of those $i$ on which some $\nu_p(\cdot)$ with $p\in X^{+r(x_0)}$ depends. For $i\in T$ there is such a $p$ with $\nu_p(z)=a_pz_i$, $a_p\ne0$, so that $z_i=\nu_p(z)/a_p$. By the first sentence, $Mg$ depends on $(z_i)_{i\in T}$ alone, and so it is a function of the $\nu_p(z)$ with $p\in X^{+r(x_0)}$: this is the second sentence. In the setting of \Cref{Y-P1-DATAo}, at most one $t$ has $g_t(x(p))\ne0$; for that $t$ we have $\nu_p(z)=\mathsf b^{-1}z_tg_t(x(p))$, and if there is none, $\nu_p(z)=0$.

\emph{Proof 6: the estimate \textup{(J5$^\circ$)}\,\textup{(a)} to \textup{(d)}.} We show: for admissible data the mixed jets are at most $0.03$, as a supremum over $(X,V)$ of a sum of moduli at $\varkappa_k$, and $u_k^{\mathrm{pol}}$ lies within $0.73$ of $\bar u_k$ on $\mathrm{Pd}(R_k)$; moreover, the standard data and the data of \textup{(d)} are admissible.
\begin{enumerate}
\item[\textup{(a)}] We have $\sum z_t\delta_t(k)=\delta[g_z](k)$ with $\|g_z\|_\infty\le\max|z_t|<R$, so Proof~4, applied to the source $g_z$, gives \textup{(a)}.

\item[\textup{(a$'$)}] Fix $(X,V)$. By a choice of the phases of the $z_t$, $\sup_{\mathrm{Pd}(R)}|\sum z_t\delta_t(k)(X,V)|=R\sum_t|\delta_t(k)(X,V)|$; dividing by $R$ gives \textup{(a$'$)}.

\item[\textup{(b), (c)}] We argue by induction on $k$. For $k=0$ we have $R_0\le\mathfrak r^\circ_0$ (if $k_1=0$, this is a hypothesis) and $|\nu(z)|_\infty\le\max|z_t|/\mathsf b$, so \textup{(L0)} yields the bound $0.014$ for the norm at the rate $\kappa_1$, and therefore, as $\kappa_1\ge\varkappa_0$, for the norm at the rate $\varkappa_0$ as well. Now let $k<k_1$; we pass from $k$ to $k+1$.
\begin{enumerate}
\item[\textup{1.}] By \textup{(c)} at $k$ and by $\|\bar u_k\|\le1$ at $\kappa_1$ (Step~0\,(c) of the proof of \Cref{Y-H6-J}), $u_k^{\mathrm{pol}}(z)$ and $\bar u_k$ lie in $M_2(x_k;\varkappa_k,c_h)$ for $z\in\mathrm{Pd}(R_k)$. Moreover $(\varkappa_k,c_h)\in\mathbb P$, and $k<K$, so \textup{(SD)} of item~2 of \S\ref{Y-H-1} holds and $x_k\ge x_9^U$.

\item[\textup{2.}] \Cref{Y-H-C9}\,(e) at $p=(\varkappa_k,c_h)$ (\Cref{Y-H1-U1}\,(U)) gives \[\|F_{\ell_k}(u_k^{\mathrm{pol}}(z))-\bar u_{k+1}\|^{\mathrm{st},(\varkappa_k)}\le0.4\cdot0.91=0.364\] on $\mathrm{Pd}(R_k)$.

\item[\textup{3.}] Fix $(Y,V')$ and one of the members. The value of that member of $F_{\ell_k}(u_k^{\mathrm{pol}}(z))$ at $(Y,V')$ is holomorphic on $\mathrm{Pd}(R_k)$ by \Cref{Y-H1-U1}\,(c+) at $\Omega=\mathrm{Pd}(R_k)$; the hypothesis of that clause, $u_z\in M_{16}(x_k;p)$ on all of $\Omega$ with holomorphic values, is supplied by item~1, the values being polynomials. By \textup{(J1$^\circ$)}, the coefficient of $z^S$ in this value is $\delta_S(k+1)(Y,V')$. \Cref{Y-H5-LS}\,(S1), applied to the seminorm given by this one weighted value, yields $|\delta_S(k+1)(Y,V')|\le0.364\,R_k^{-|S|}$ units at $\varkappa_k$.

\item[\textup{4.}] By \textup{(J2$^\circ$)}, $\delta_S(k+1)(Y,V')=0$ unless $S\subset I_{k+1}(Y)$. Since a mixed jet has $|S|\ge2$, only those $Y$ with $j:=j_{k+1}(Y)\ge2$ matter, and with $I=I_{k+1}(Y)$ \[\sum_{S\subset I,|S|\ge2}(R_{k+1}/R_k)^{|S|}\le(1+R_{k+1}/R_k)^j.\]

\item[\textup{5.}] A unit at $\varkappa_k$ equals $e^{-\check\Delta_k(\#Y-1)}$ units at $\varkappa_{k+1}$, with $|Y|_c$ in the place of $\#Y$ for the $w$-member. By \textup{(G5$^\circ$)}\,(a), $\#Y$ for the $a$- and $h$-members, or $|Y|_c$ in the case of the $w$-member, is at least $j\mathfrak h_{k+1}$.

\item[\textup{6.}] Collecting items 3 to 5, the sum over $|S|\ge2$ of $|\delta_S(k+1)(Y,V')|R_{k+1}^{|S|}$, in units at $\varkappa_{k+1}$, is at most \[0.364\,e^{\check\Delta_k}[(1+R_{k+1}/R_k)e^{-\check\Delta_k\mathfrak h_{k+1}}]^j=0.364\,e^{\check\Delta_k}q_k^j\le0.364\,e^{\check\Delta_k}q_k^2\le0.03,\] which is \textup{(b)} at $k+1$. Here $R_{k+1}$ enters only through the ratio $R_{k+1}/R_k$, and no estimate is applied at the rate $\varkappa_{k+1}$; so at $k+1=k_1$ no bound on $R_{k_1}$ enters, and $\varkappa_{k_1}$ may lie outside $[\kappa,\kappa_1]$.

\item[\textup{7.}] If $R_{k+1}\le\mathfrak r^\circ_{k+1}$ and $\varkappa_{k+1}\ge\kappa$, then \textup{(c)} holds at $k+1$: by \textup{(a)} the single jets contribute at most $0.2\,\mathrm{amp}(Lx_{k+1})\mathfrak r^\circ_{k+1}=0.7$ at $\kappa_1$, and the mixed jets at most $0.03$, together $0.73$. The coefficients lie in $\mathfrak M$ by \Cref{Y-H5-LS}\,(S3), the values are polynomials in $z$, and $\|\bar u_{k+1}\|\le1$; so the family lies in $M_{1.73}$, hence in $M_{1.91}$.
\end{enumerate}

The decay factor is raised to the power $j(Y)$, one factor for each of the sources near $Y$, and not to the power $|S|$, one for each element of $S$: with the power $|S|$, it would leave the sum $(1+q')^j-1-jq'$, which grows with $j$. Here $q'=(R_{k+1}/R_k)e^{-\check\Delta_k\mathfrak h_{k+1}}$.

\item[] \emph{Standard data.} For these, condition (i) of \textup{(AD)} holds with equality. Let $k<k^\circ$. Then $\mathrm{sep}_{k+1}\ge L\,\mathrm{sep}^\circ$, $R_{k+1}/R_k\le L^{1/2}$ and $\mathfrak h_{k+1}\ge\mathrm{sep}_{k+1}/2+1$; as $\check\kappa^\circ_k=(\mathrm{sep}^\circ/4)/\mathrm{sep}_{k+1}$ and $\mathrm{sep}^\circ=8\log(8L^{1/2})$, this gives $\check\kappa^\circ_k\mathfrak h_{k+1}\ge\log(8L^{1/2})+\check\kappa^\circ_k$, that is, $e^{-\check\kappa^\circ_k\mathfrak h_{k+1}}\le e^{-\check\kappa^\circ_k}/(8L^{1/2})$. Hence $q_k\le e^{-\check\kappa^\circ_k}(1+L^{-1/2})/8<1$, and \[0.364\,e^{\check\Delta}q^2\le0.364((1+L^{-1/2})/8)^2\le0.0057.\]

\item[\textup{(d)}] By \textup{(G1)} of item~3 of \S\ref{Y-H6-1}, $\mathrm{sep}_{i+1}\ge L^{k-i-1}\mathrm{sep}_k>0$ for $i<k$. Since $\mathrm{sep}_k\ge8\log(4\varrho L^{1/2})\ge8\log(4L^{1/2})$, we have $\check\kappa_{\mathrm{last}}\le1/4$ and $\sum_{i<k-1}\check\kappa_i\le\frac14\sum_{m\ge1}L^{-m}=1/(4(L-1))$. So the rates lie in $[\kappa+1/4,\kappa+1/2+1/(4(L-1))]$, inside $[\kappa,\kappa_1]$. For $i<k-1$, $q_i\le e^{-\check\kappa_i}(1+L^{-1/2})/4$ and $0.364\,e^{\check\Delta}q^2\le0.364((1+L^{-1/2})/4)^2\le0.02276$. At the last step, $R_k/R_{k-1}\le\varrho L^{1/2}$ and $e^{-\check\kappa_{\mathrm{last}}\mathfrak h_k}\le e^{-\check\kappa_{\mathrm{last}}}/(4\varrho L^{1/2})$, so $q\le e^{-\check\kappa_{\mathrm{last}}}(1+1/(\varrho L^{1/2}))/4$, and there too $0.364\,e^{\check\Delta}q^2\le0.02276$.
\end{enumerate}

\emph{Remark.} The standard data of \textup{(AD)} remain admissible when $\mathrm{sep}^\circ$ is replaced by an arbitrary numeral $\widehat{\mathrm{sep}}\ge8\log(4L^{1/2})$. Moreover, if two such numerals lie within a factor $L$ of one another, their values of $j^\circ$ differ by at most $1$. We put any numeral $\widehat{\mathrm{sep}}\ge8\log(4L^{1/2})$ in the place of $\mathrm{sep}^\circ$, and we define with it the number $j^\circ$ of \textup{(D1)} as well as the number $k^\circ$ and the rates $\kappa^\circ_k$ of \textup{(D2)}: $j^\circ$ is the least $j\ge1$ with $\mathrm{sep}^{[j-1]}\ge\widehat{\mathrm{sep}}$, $k^\circ=\max\{K-j^\circ,0\}$, $\check\kappa_k:=(\widehat{\mathrm{sep}}/4)/\mathrm{sep}_{k+1}$ for $k<k^\circ$, and the rates are $\kappa+\sum_{k\le j<k^\circ}\check\kappa_j$. In the two estimates below, $\check\kappa$ without an index denotes $\check\kappa_k$. The data so obtained, namely the standard data of \textup{(AD)} changed in this one respect, with $k_1=k^\circ$, $R_k=\mathfrak r^\circ_k$ and these rates as $\varkappa_k$, are the data of the first assertion. For them, the two estimates given above for the standard data go through with these changes: for $k<k^\circ$, with this $k^\circ$, we have $\mathrm{sep}_{k+1}\ge L\,\widehat{\mathrm{sep}}$, and $\check\kappa\,\mathfrak h_{k+1}\ge\widehat{\mathrm{sep}}/8+\check\kappa\ge\log(4L^{1/2})+\check\kappa$, so that $e^{-\check\kappa\,\mathfrak h_{k+1}}\le e^{-\check\kappa}/(4L^{1/2})$; hence $q_k\le e^{-\check\kappa}(1+L^{-1/2})/4<1$ and $0.364\,e^{\check\Delta}q^2\le0.364((1+L^{-1/2})/4)^2\le0.02276\le0.03$, while condition (i) of \textup{(AD)} again holds with equality. For $k<k^\circ$ we have $\mathrm{sep}_{k+1}\ge L^{K-k-j^\circ}\,\widehat{\mathrm{sep}}$ with $K-k-j^\circ\ge1$: this follows from \textup{(G1)} of item~3 of \S\ref{Y-H6-1}, applied from the depth $j^\circ-1$ on, by the reasoning of Step~1\,(G2) of the proof of \Cref{Y-H6-J}, where $\widehat{\mathrm{sep}}$ takes the place of $\mathrm{sep}_*$. We therefore get $\sum_{k<k^\circ}\check\kappa_k\le\frac14\sum_{m\ge1}L^{-m}=1/(4(L-1))$. So the rates $\varkappa_k=\kappa+\sum_{k\le j<k^\circ}\check\kappa_j$, $0\le k\le k^\circ$, are non-increasing in $k$, since each $\check\kappa_j>0$, and they lie in $[\kappa,\kappa+1/(4(L-1))]$, inside $[\kappa,\kappa_1]$; they thus form the chain of rates that \textup{(AD)} requires. By \textup{(G1)} of item~3 of \S\ref{Y-H6-1}, the second assertion of the remark follows.

These estimates use the $L$ of the table, $L>1.2\cdot10^{13}$ (\textup{(O4)} in item~1 of \S\ref{Y-H-1}, as cited in Step~0\,(b) of the proof of \Cref{Y-H6-J}); the size of $L$ matters for the numeral $0.364((1+L^{-1/2})/4)^2$, which arises when $4L^{1/2}$ is under the logarithm, rather than the $8L^{1/2}$ of the standard data: at $L=13$, for instance, it would be $0.0371>0.03$.

\emph{Proof 7: the clause \textup{(J6$^\circ$)}, on the representatives.} We prove \textup{(a)} to \textup{(e)} of \textup{(J6$^\circ$)} for every $K\ge1$.

The induction starts at $k=k^\circ$. There \textup{(J5$^\circ$)\,(c)}, taken for the standard data, for which $\kappa^\circ_{k^\circ}=\kappa$, shows that $u_{k^\circ}^{\mathrm{pol}}$ is a family on $\mathrm{Pd}(\mathfrak r^\circ_*)$ within $B=1.91$ caps at $p=(\kappa,c_h)$, whose distance to $\bar u_{k^\circ}$ is $\le0.73\le0.91$. If $K<j^\circ$, then $k^\circ=0$, and the start is \textup{(L0)}, in which no separation enters.

The passage from $k$ to $k+1$, for $k^\circ\le k<K$, is Step~8 of the proof of \Cref{Y-H6-J}, carried out with the following data. We apply \Cref{Y-H5-Af}\,(i), (ii) with $p=(\kappa,c_h)$, on $\Omega=\mathrm{Pd}(\mathfrak r^\circ_*)\subset\mathbb C^Q$ (item~0\,(5) of \S\ref{Y-H5-1} allows any $n\ge1$), to the pair consisting of $u_k(z)$ and the constant family $\bar u_k$, with $B=1.91$ at $k=k^\circ$ and $B=1$ at the later levels. By \Cref{Y-H5-Af}\,(i), the image is a family in $M_1(x_{k+1};\kappa,c_h)$ that obeys the strong cap, and its $a$-member satisfies $|E'^a|_a^{\Sigma}\le0.35+0.248\cdot1.91=0.82368\le1$. The strong distance of the images is at most $\vartheta[B]$ times the weak distance of the pair, where $\vartheta[1.91]=3.968/14.09\le0.2817$ and $\vartheta[1]\le0.2646$, both $\le0.4$. Since $F_{\ell_k}(\bar u_k)=\bar u_{k+1}$, this gives \textup{(a)} and \textup{(b)}.

\begin{enumerate}
\item[\textup{(c)}] If at level $k$ the family $u_k(z)$ and the honest family have the same $\partial_{S'}$ at $0$ for every $S'$, then \Cref{Y-H5-LJ}, applied value by value, shows that so do their images. At $k^\circ$ this holds by the form of $u^{\mathrm{pol}}$.

\item[\textup{(d)}] The property $(P_k)$ of the proof of \Cref{Y-H6-J} holds for $k\le k^\circ$ by \textup{(J2$^\circ$)}, and it passes from each level to the next by the claim of Step~3 of that proof. Here $k^\circ\le K-1$, by \textup{(G2)} of item~3 of \S\ref{Y-H6-1}, as restated in \textup{(G$^\circ$)}.

\item[\textup{(e)}] The last step is $k=K-1\ge k^\circ$. There \Cref{Y-H5-Af}\,(ii), in its part on the factors of the three members, bounds the $h$-member in the cell sum $|\cdot|_h^\Sigma$, with the factor $1/(16-B)$: $|E'^h[u_1]-E'^h[u_2]|_h^\Sigma\le \mathsf D^{\mathrm{in}}/(16-B)\le0.91/14.09\le0.0646$ on $\mathrm{Pd}(\mathfrak r^\circ_*)$, in caps $c_h(\mathsf b)$, at $(\kappa,c_h)$. The cell sum is an admissible seminorm (item~0\,(4) of \S\ref{Y-H5-1}), so \Cref{Y-H5-LS}\,(S1), applied with $\mu=1_{\{t\}}$, yields $|\delta_t^h(K)|_h^\Sigma\le0.0646/\mathfrak r^\circ_*$.

By \textup{(J2$^\circ$)} for $k=K$ and $S=\{t\}$, $\delta_t^h(K)(X)=0$ unless $X$ contains a cell of $S_t(K)^{+\check r(\mathsf b)}$, and there are at most $N_1(K-1)=N_1^{[1]}$ such cells (Step~1\,(A) of the proof of \Cref{Y-H6-J}). We add up the cell sums over these cells and use $\mathfrak r^\circ_*\ge r^\circ$. The field $1$, at which $e_t$ is evaluated, is $G$-valued and tame.
\end{enumerate}

\emph{Proof 8: the clause \textup{(J7$^\circ$)}, on congruences and vacuum functions.} We prove \textup{(i)} to \textup{(vii)} of \textup{(J7$^\circ$)}.

We use the squarefree congruence of \Cref{Y-H8-Congruence}. Its compatibility with sums, products and the exponential was proved in Step~P7 of the proof of \Cref{Y-H8-CoF}.

\begin{enumerate}
\item[\textup{(i)}] \Cref{Y-H4-U4}\,(h), applied on $\mathrm{Pd}_0^o$ to the honest family, gives $\rho_k^z=\exp\{\sum_{j\le k}V_j(z)\}\rho[u_k^{\mathrm{hon}}(z);\allowbreak \ell_k]$ at every field of $\mathbb T_k$. By \Cref{Y-H1-U1}\,(g1), (f1), $V_{k'+1}(z)-V_{k'+1}(0)=\sum_Y[v_{k'}(Y;\allowbreak u_{k'}^{\mathrm{hon}}(z))-v_{k'}(Y;\allowbreak \bar u_{k'})]$, a finite sum; the term $V_0^+$ involves no member. \Cref{Y-H5-LJ}, applied with $\Phi=v_{k'}(Y;\cdot)$, shows that the coefficient of $z^S$ depends only on $u(0)$ and the $u_{S'}$; therefore $v_{k'}(Y;\allowbreak u_{k'}^{\mathrm{hon}}(z))\equiv v_{k'}(Y;\allowbreak u_{k'}(z))$.

\Cref{Y-H5-LJ} applies as well with $\Phi=\rho[\cdot;\ell_k](V)$, with $V$ fixed, by its last part, the one on $I[u]$: at each $V$ this function is $e^{-x\mathrm{En}_\ell(V)}$, in which no member enters, times the exponential of a finite sum of values of $E^a,E^h$, times a polynomial in finitely many values of $w$; it is entire along lines and bounded on $\{\|u\|<R\}$. The lattices $\mathbb T_k$ with $k<K$ are those of item~2 of \S\ref{Y-I-1} at $x_k$, and $\mathbb T_K$ is the one named in \Cref{Y-H5-LJ}. Thus $\rho[u_k^{\mathrm{hon}}(z);\ell_k](V)\equiv\rho[u_k(z);\ell_k](V)$, and $\equiv$ then passes through the sum, the exponential and the product.

\item[\textup{(ii)}] By \textup{(J2$^\circ$)}, $u_k^{\mathrm{pol}}$ has the property $(P_k)$. The claim of Step~3 of the proof of \Cref{Y-H6-J} then shows that $v_k(Y;u_k^{\mathrm{pol}}(z))$ does not depend on $z_t$ if $Y^{+\check r(x_{k+1})}$ misses $S_t(k+1)$, so that its coefficient of $z^S$ vanishes whenever $t\in S$. By \Cref{Y-H5-LJ}, this coefficient is the same for all representatives.

\item[\textup{(iii)}] We take $u^{\mathrm{pol}}$ at all levels and replace each $v_{k'}(Y;u_{k'}^{\mathrm{pol}}(z))$ by its multi-affine part $v_{k'}(Y;\bar u_{k'})+\sum_Sz^S\partial_Sv_{k'}(Y;u_{k'}^{\mathrm{pol}})(0)$. The terms with $|S|=1$ add up to $\sum z_t\gamma_t^{(k)}$, because $\partial_tV_{k'+1}(0)=\sum_Y\partial_tv_{k'}(Y;u_{k'}^{\mathrm{pol}})(0)$.

\item[\textup{(iv)}] For $k'<k^\circ$ we replace, exactly as in \textup{(iii)}, each $v_{k'}(Y;u_{k'}^{\mathrm{pol}}(z))$ by its multi-affine part. For $k'\ge k^\circ$ we write $w_{k'}=\sum z_t\partial_tw_{k'}(0)+[w_{k'}-\sum z_t\partial_tw_{k'}(0)]$, where $\partial_tw_{k'}(0)=\partial_tV_{k'+1}(0)$. Each $w_{k'}$ is a finite sum of functions holomorphic on $\mathrm{Pd}(\mathfrak r^\circ_*)$, by \Cref{Y-H1-U1}\,(c+).

\item[\textup{(v)}] By \Cref{Y-H4-U4}\,(h), $Z_K(z)=\exp\{\sum V_j(z)\}I[u_K^{\mathrm{hon}}(z)]$, and \Cref{Y-H5-LJ} applies with $\Phi=I$. As $\rho[\bar u_K;\ell_K]>0$ (\Cref{Y-H4-U4}\,(h)), we have $I[\bar u_K]>0$, and we may divide by $Z_K(0)=\exp\{\sum V_j(0)\}I[\bar u_K]$.

\item[\textup{(vi)}] Near $0$ we have $\log Z_K=\sum V_j(z)+\log I[u_K^{\mathrm{hon}}(z)]$. The argument of Step~9 of the proof of \Cref{Y-H6-J}, with $S$ in place of $\{1,\dots,n\}$, shows that $\partial_S\log I$ is a polynomial in the $\partial_{S'}I/I(0)$, $S'\subset S$.

\item[\textup{(vii)}] Since $\rho_K^0(1)=\exp\{\sum V_j(0)\}\rho[\bar u_K;\ell_K](1)>0$, the logarithm of $\rho_K^z(1)$ exists near $0$. \Cref{Y-H4-U4}\,(h) at $V=1$ gives $\partial_t\log\rho_K^z(1)(0)=\gamma_t+\partial_t\log\rho[u_K^{\mathrm{hon}}(z);\ell_K](1)(0)$.

The set $D(1)$ is empty, since $\phi_p=0$ at the field $1$. Every polymer in a family contributing to $\Xi_w(V)$ contains a plaquette of $D(V)$, by item~3 of \S\ref{Y-H-1}. At $V=1$ only the empty family contributes, so $\Xi_w(1)=1$. Thus $\rho[u;\ell_K](1)=\exp\{-\mathsf b\mathrm{En}_\ell(1)+\sum_X[E^a+E^h](X,1)\}$, with $\mathrm{En}_\ell$ independent of the members. Moreover $E^a(X,1)=0$, by the clause ``$0$ at pure gauges'', a clause without a size which every member of the family satisfies. As $E^a$ is $0$ at the field $1$ for every $z$, we obtain $\partial_t\log\rho[u_K^{\mathrm{hon}}(z);\ell_K](1)(0)=\sum_X\partial_t[E^a+E^h](X,1)(0)=\sum_X\delta_t^h(K)(X,1)=e_t$.

Setting $z_s=0$ for $s\ne t$ turns the run with $Q$ sources into the run with the one source $g_t$. Hence $\partial_t\log\rho_K^z(1)(0)$ is $\hat c_K(g_t)$, which is $c^{(0)}_K(g_t)$ of \S\ref{Y-P2-1}, the subtraction of its item~2 and of \Cref{Y-P2-ThR}\,(R0).
\end{enumerate}

\emph{Proof 9: the clauses \textup{(J8$^\circ$)\,(a)}, \textup{(a$^\circ$)}, \textup{(a$'$)}, on the early vacuum numbers per source.} For \textup{(a)} we show that \[\sum_{S\ni t,|S|\ge2}R^{|S|}\sum_Y|\partial_Sv_k(Y;u_k^{\mathrm{pol}})(0)|\le N_1(k)2C_gc_a(x_k)e^{\lambda}\exp\{-(\varkappa_k-\lambda)(\mathrm{sep}_{k+1}+1)\},\] where $\lambda=\log(1+R/R_k)/\mathfrak h_{k+1}\le\varkappa_k$.

Fix $k<k_1$ and $r<R_k$, and put $\phi_Y(z):=v_k(Y;u_k^{\mathrm{pol}}(z))$; by \Cref{Y-H1-U1}\,(c+), this function is holomorphic on $\mathrm{Pd}(R_k)$.

\begin{enumerate}
\item[\textup{1.}] One mean suffices for every $S$: $\partial_S\phi_Y(0)=r^{-|S|}\mathrm{Mean}[\phi_Y(re^{i\theta})e^{-i\sum_{s\in S}\theta_s}]$, where the mean is taken over the whole $Q$-torus, because the mean of $z^\mu e^{-i\sum_S\theta}$ vanishes unless $\mu=1_S$. Hence $R^{|S|}|\partial_S\phi_Y(0)|\le(R/r)^{|S|}\mathrm{Mean}|\phi_Y|$, with the same mean for all $S$.

\item[\textup{2.}] By \textup{(J7$^\circ$)\,(ii)}, only sets $S\subset I_{k+1}(Y)$ occur. Over the $S\subset I_{k+1}(Y)$ with $t\in S$ and $|S|\ge2$, the factors $(R/r)^{|S|}$ add up to at most $(1+R/r)^{j(Y)}$. This sum is empty unless $Y\in\check Y_t:=\{Y$ connected: $t\in I_{k+1}(Y),\ j(Y)\ge2\}$, where $j(Y)=j_{k+1}(Y)$.

\item[\textup{3.}] The sums over $S$ and over $Y$ are finite, so we may take them inside the mean. Thus the left side of the bound is $\le\sup_{|z_s|=r}\sum_{Y\in\check Y_t}(1+R/r)^{j(Y)}|v_k(Y;u_k^{\mathrm{pol}}(z))|$.

\item[\textup{4.}] By \textup{(G5$^\circ$)\,(a)} for $k+1$, $\#Y\ge j(Y)\mathfrak h$ with $\mathfrak h=\mathfrak h_{k+1}$; therefore $(1+R/r)^{j(Y)}\le e^{\lambda_r\#Y}$, where $\lambda_r=\log(1+R/r)/\mathfrak h$. Also $\#Y\ge2\mathfrak h\ge\mathrm{sep}_{k+1}+2$.

\item[\textup{5.}] Each $Y\in\check Y_t$ contains a cell of $S_t(k+1)^{+\check r(x_{k+1})}$, and for any $t$ there are $\le N_1(k)$ such cells (Step~1\,(A) of the proof of \Cref{Y-H6-J}).

\item[\textup{6.}] We use \Cref{Y-H1-U1}\,(g2) with $A=2$ and $p=(\varkappa_k,c_h)$; it applies because $u_k^{\mathrm{pol}}(z)\in M_2$ by \textup{(J5$^\circ$)\,(c)}. For a cell, provided that $\lambda_r\le\varkappa_k$, \[\sum_{Y\ni\mathrm{cell},\#Y\ge2\mathfrak h}e^{\lambda_r\#Y}|v|=\sum e^{\lambda_r-(\varkappa_k-\lambda_r)(\#Y-1)}e^{\varkappa_k(\#Y-1)}|v|\le e^{\lambda_r}e^{-(\varkappa_k-\lambda_r)(2\mathfrak h-1)}2C_gc_a(x_k).\]
\end{enumerate}

Letting $r\to R_k$ gives the bound of \textup{(a)}; in the case $\lambda=\varkappa_k$ it follows by continuity of both sides in $R$.

\begin{enumerate}
\item[\textup{(a$^\circ$)}] Let $t':=k^\circ-k\ge1$. By \textup{(G3)} of item~3 of \S\ref{Y-H6-1}, as restated in \textup{(G$^\circ$)}, $\mathfrak r^\circ_*/\mathfrak r^\circ_k\le L^{t'/2}$, and $\log(1+L^{t'/2})\le t'\log(2L^{1/2})\le t'\mathrm{sep}^\circ/8$. The depth of $k+1$ is $j^\circ-1+t'$, so by \textup{(G1)} of the same item $\mathfrak h>\mathrm{sep}_{k+1}/2\ge L^{t'}\mathrm{sep}^\circ/2$, and hence $\lambda\le t'/(4L^{t'})\le1/(4L)$. Also $\varkappa_k\ge\kappa$. The sum over $t$ counts each $S$ exactly $|S|$ times, and $|S|\ge2$; therefore $\sup_{\mathrm{Pd}(\mathfrak r^\circ_*)}|W_k|\le(Q/2)E^{[j]}$.

\item[\textup{(a$'$)}] Put $k'=k-m$, $R=\varrho\,\mathfrak r^\circ_k$ and $R_{k'}=\mathfrak r^\circ_{k'}$. Then $R/R_{k'}\le\varrho L^{m/2}$, and $\log(1+\varrho L^{m/2})\le m\log(2\varrho L^{1/2})\le m\,\mathrm{sep}_k/8$. Since $\mathfrak h_{k'+1}>\mathrm{sep}_{k'+1}/2\ge L^{m-1}\mathrm{sep}_k/2$, we obtain $\lambda\le m/(4L^{m-1})\le1/4$. Moreover $\varkappa_{k'}-\lambda\ge\kappa\ge\kappa-1/4$ and $\mathrm{sep}_{k'+1}+1\ge L^{m-1}\mathrm{sep}_k+1$, and we sum over $1\le m\le k$. The numbers $N_1(k-m)=N_1^{[K-k+m]}$ and $x_{k-m}=L^{K-k+m}\mathsf b$ are defined for every $m\ge1$, and all terms are $\ge0$; so the finite sum is $\le E_{\mathrm{vac}}$.

The series $E_{\mathrm{vac}}$ is finite. Item~(ii) of the proof of \Cref{Y-H9-LN}\,(N3) uses \textup{(N1)} alone among the clauses of that lemma; carried out with $K-k+1$ in place of $j^*$, with $m-1$ in place of the offset $t$ of that item and with $\check\rho$ in place of $R_f$, it gives $N_1^{[K-k+m]}\le L^{3(m-1)}N_1^{[K-k+1]}$, with equality at $m=1$; moreover $2C_gc_a\le C_g/16$ by item~(iii) of that proof. Since $\mathrm{sep}_k\ge8\log(4L^{1/2})\ge4\log L$ and $\kappa-1/4\ge5.75$, \[(\kappa-1/4)L^{m-1}\mathrm{sep}_k\ge23L^{m-1}\log L\ge3(m-1)\log L+20L^{m-1}\log L.\] So $E_{\mathrm{vac}}\le N_1^{[K-k+1]}(C_g/8)e^{1/4}L^{-20}$.
\end{enumerate}

\emph{Proof 10: the clauses \textup{(J8$^\circ$)\,(b)}, \textup{(c)}, \textup{(d)}, and \Cref{Y-P1-LNo}.} We bound the late levels without shrinking the radius, prove $|W|\le QC_W^\circ\le QC_W$, and show that the series are finite and locally bounded.

\begin{enumerate}
\item[\textup{(b)}] Let $k^\circ\le k<K$ and $z\in\mathrm{Pd}(\mathfrak r^\circ_*)$. The representative has the property $(P_k)$, by \textup{(J6$^\circ$)\,(d)}. If $I_{k+1}(Y)$ is empty, then by the claim of Step~3 of the proof of \Cref{Y-H6-J}, $v_k(Y;u_k(z))$ does not depend on any $z_t$, so it equals its value at $0$, and the term of $Y$ is $0$. Every other $Y$ contains one of the $\le QN_1(k)$ cells of $\bigcup_tS_t(k+1)^{+\check r(x_{k+1})}$. \Cref{Y-H1-U1}\,(g2), applied with $p=(\kappa,c_h)$, with $A=1.91$ (at $k=k^\circ$) or $1$ for $u_k(z)$ and with $A=1$ for $\bar u_k$, bounds the contribution of each cell by $2.91\,C_gc_a(x_k)$.

For $\partial_tw_k(0)$ we set $z_s=0$ for $s\ne t$. Then every $Y$ with $t\notin I_{k+1}(Y)$ gives $0$, and each of the others contains one of $N_1(k)$ cells; hence $|w_k|\le3N_1(k)C_gc_a$ on $|z_t|<\mathfrak r^\circ_*$, and Cauchy's estimate applies. So $|w_k-\sum z_t\partial_tw_k(0)|\le6QN_1(k)C_gc_a(x_k)$ on $\mathrm{Pd}(\mathfrak r^\circ_*)$.

\item[\textup{(c)}] We repeat items~1 to 3 and 5 of Proof~9 with these changes: throughout, $R=\mathfrak r^\circ_*/n_{\mathrm{loc}}$; in item~1 the function is $v_k(Y;u_k(z))$ of \textup{(b)}, and the radius is any $r<\mathfrak r^\circ_*$; in items~2 and 3 the sums run over all $S\ni t$; in item~5 each connected $Y$ with $t\in I_{k+1}(Y)$, without the condition $j(Y)\ge2$, contains one of the $\le N_1(k)$ cells of $S_t(k+1)^{+\check r(x_{k+1})}$. Letting $r\to\mathfrak r^\circ_*$, we get $(1+R/r)^{j(Y)}\to(1+1/n_{\mathrm{loc}})^{j(Y)}\le(1+1/n_{\mathrm{loc}})^{n_{\mathrm{loc}}\#Y}\le e^{\#Y}$, by \textup{(G5$^\circ$)\,(b)} for $k+1$. For a cell, $\sum_{Y\ni\mathrm{cell}}e^{\#Y}|v|\le e\sum e^{\kappa(\#Y-1)}|v|\le e\cdot1.91\,C_gc_a$; the first inequality holds because $\kappa\ge1$, and the second follows from \Cref{Y-H1-U1}\,(g2), applied with $p=(\kappa,c_h)$, with $A=1.91$ (at $k=k^\circ$) or $1$. Summing over the $\le N_1(k)$ cells and using $1.91e\le2e$, we obtain the bound $2e\,N_1(k)C_gc_a(x_k)$ of \textup{(J8$^\circ$)\,(c)}.

\item[\textup{(d)}] We pass to the depth $j=K-k$, so that $x_k=L^j\mathsf b$ and $N_1(k)=N_1^{[j]}$. Then $k<k^\circ$ holds if and only if $j>j^\circ$, in which case $K>j^\circ$; so the late levels are the depths $1,\dots,\min(K,j^\circ)$. Combining \textup{(a$^\circ$)} and \textup{(b)} gives \[|W|\le Q[\frac12\sum_{j>j^\circ}E^{[j]}+6\sum_{j\le j^\circ}N_1^{[j]}C_gc_a(L^j\mathsf b)].\] For $C_W^\circ\le C_W$ we put $s=\mathrm{sep}^{[j-1]}>0$; then $e^{1/(4L)}\le2e$ and $(\kappa-1/(4L))(s+1)\ge(\kappa-1)(s-1)$.
\end{enumerate}

We turn to \Cref{Y-P1-LNo}. Let $j=j^\circ+t$ with $t\ge1$. Item~(ii) of the proof of \Cref{Y-H9-LN}\,(N3) uses \textup{(N1)} alone among the clauses of that lemma; carried out with $j^\circ$ in place of $j^*$ and with $\check\rho$ in place of $R_f$, it gives $N_1^{[j]}\le L^{3t}N_1^{[j^\circ]}$. By item~(iii) of the same proof, $2C_gc_a\le C_g/16$, and by \textup{(G1)} of item~3 of \S\ref{Y-H6-1}, as restated in \textup{(G$^\circ$)}, $\mathrm{sep}^{[j-1]}\ge L^t\mathrm{sep}^\circ$. Moreover $(\kappa-1/(4L))\mathrm{sep}^\circ\ge5.9\cdot8\log(8L^{1/2})\ge23.6\log L$ and $L^{3t}\le e^{3L^t\log L}$. Therefore $E^{[j]}\le N_1^{[j^\circ]}(C_g/16)e^{1/(4L)}\exp\{-20L^t\log L\}$, and these terms are summable.

For the early terms of $C_W$, again with $j=j^\circ+t$, we have $\mathrm{sep}^\circ=8\log8+4\log L$, $\kappa-1\ge5$ and $L^t\mathrm{sep}^\circ-1>0$. Hence \[(\kappa-1)(\mathrm{sep}^{[j-1]}-1)\ge5(L^t\mathrm{sep}^\circ-1)=20L^t\log L+40\log8\cdot L^t-5\ge20L^t\log L+83L^t-5.\] The term $83L^t$ absorbs the factor $L^{3t}$ in $N_1^{[j]}\le L^{3t}N_1^{[j^\circ]}$, as $L^{3t}=e^{3t\log L}$ and $L^t\ge tL\ge t\log L$. Next, $2ec_a\le2e/32<1$ by \textup{(X17)} of \textup{(O13)}, item~1 of \S\ref{Y-J-1}, which is the inequality $16(c_a+c_h)(x)\le\frac12$, that is, $c_a+c_h\le1/32$. So \[2eN_1^{[j]}C_gc_a(L^j\mathsf b)e^{-(\kappa-1)(\mathrm{sep}^{[j-1]}-1)}\le N_1^{[j^\circ]}C_g\exp\{-20L^t\log L+5\}.\]

Local boundedness follows the proof of \Cref{Y-H9-LN}\,(N5). Let $\mathsf b$ range over a compact set, and put $\hat\jmath:=\max j^\circ$, where the maximum is taken over this set. Here $j^\circ$ is non-decreasing in $\mathsf b$. The terms with $j\le\hat\jmath$, in either form, are products of the continuous $c_a(L^j\mathsf b)$, of the non-decreasing $N_1^{[j]}$ and of factors $\le6$, and so they are bounded: in the late form the factor is $6$, and in the early form it is $2e$ times an exponential $\le1$, because $\mathrm{sep}^{[j-1]}\ge L\,\mathrm{sep}^\circ\ge1$ there. For $j>\hat\jmath$ the term is $\le\hat P\exp\{-20L^{j-\hat\jmath}\log L+5\}$, where $\hat P$ is an upper bound for $N_1^{[j^\circ]}C_g$ on the compact set. No term depends on $K$, $N$ or $Q$.

\emph{Proof 11: \Cref{Y-P1-Cor}.} We prove \textup{(a)}, the bound on the moment of $Q$ sources with $C_{\mathrm{vol}}$ to the first power, for all $K\ge1$, and \textup{(b)}, the same bound with a bounded Borel function of the last field.

The data \textup{(B)} of \Cref{Y-H8-ThBC} with $x=\mathsf b$ on $\mathbb Z^3/N\mathbb Z^3$, with $\tilde\kappa=\kappa$, are met for every $K\ge1$: this is shown in Step~2 of the proof of \Cref{Y-H9-T}, in its part on the last integral, by an argument into which no source enters. These data are: $\mathsf b\ge x_s\ge e^2$; $N$ even, $M|N$, $N\ge2Ms(\mathsf b)$; $\varepsilon(\mathsf b)\le1$; $(c_a+c_h)(\mathsf b)\le1/2$; and $\ell_K$ of range $\le\rho(\mathsf b)$ with $|\ell_{K,y}(p)|\le K_*e^{-\lambda_*|y-p|}$.

The first member is $u_1=\bar u_K$, which by \Cref{Y-H4-U4}\,(h) is a system of members of a point of the class; we put $u_3:=u_1$. For the second we put $u_2:=u_K(z)$ with $z\in\mathrm{Pd}(\mathfrak r^\circ_*)$. By \textup{(J6$^\circ$)\,(a)} it lies in $M_1(\mathsf b;\kappa,c_h)$ with the strong cap, so $\|u_2\|^{\sim}_E\le1$ and $\|u_2\|^{\sim}_w\le1$; thus $a_w=1$ and $a_E\le2$, which is case \textup{(ii)} of \Cref{Y-H8-ThBC}\,(b$^C$). The pair is $(\kappa,c_h)$, since at $c_h/2$ the cap of $E^h$ would not give $a_E\le2$. Then $\mathsf D_{23}=\|u_K(z)-\bar u_K\|^{\sim}\le\|\cdot\|^{\mathrm{st},(\kappa)}\le0.91\cdot0.4^{K-k^\circ}\le0.364$, since $k^\circ\le K-1$.

In its part before the condition \textup{(S)}, \Cref{Y-H8-ThBC}\,(c$^C$) gives the $L^1$ bound $\int|\rho[u_2;\ell]-\rho[u_3;\ell]|\,dV\le\frac12C_{\mathrm{vol}}\mathsf D_{23}I_1$ in the cases \textup{(i)}, \textup{(ii)}. In case \textup{(ii)}, with $\mathsf D_{23}\le0.364$, this gives $\int|\rho[u_K(z);\ell_K]-\rho[\bar u_K;\ell_K]|dV\le0.182\,C_{\mathrm{vol}}I_1$. Hence, on $\mathrm{Pd}(\mathfrak r^\circ_*)$, \begin{equation}\tag{$*$}\int|\rho[u_K(z);\ell_K]|dV\le(1+0.182\,C_{\mathrm{vol}})I_1.\end{equation}

\begin{enumerate}
\item[\textup{(a)}] The moment $\langle\prod_t(\mathcal O_K(g_t)-a_t)\rangle$ is the coefficient of $z_1\cdots z_Q$ in $e^{-\sum a_tz_t}Z_K(z)/Z_K(0)$, a squarefree coefficient. With $a_t=\gamma_t+e_t$ and \textup{(J7$^\circ$)\,(v)}, it equals the same coefficient of $G(z):=\exp\{W(z)-\sum e_tz_t\}I[u_K(z)]/I_1$. The function $G$ is holomorphic on $\mathrm{Pd}(\mathfrak r^\circ_*)\supset\mathrm{Pd}(r^\circ)$, by \textup{(J7$^\circ$)\,(iv)} and by \Cref{Y-H8-CoF}\,(F0), (F1) with its datum $B_0=0.91$. On $\mathrm{Pd}(r^\circ)$ we have $|W|\le QC_W$ by \textup{(J8$^\circ$)\,(d)}, $|\sum e_tz_t|\le Qe_h$, and $|I/I_1|\le1+0.182\,C_{\mathrm{vol}}$ by \textup{($*$)}. By Cauchy's estimate the coefficient is $\le(r^\circ)^{-Q}e^{Q(C_W+e_h)}(1+0.182\,C_{\mathrm{vol}})=G_1^Q(1+0.182\,C_{\mathrm{vol}})$.

With $q=(N/M)^3\ge1$, \[1+0.182\,C_{\mathrm{vol}}\le[1+0.728\,qe^{220q}]e^{(c_{\mathrm{vol}}(\mathsf b)-225/M^3)N^3}\le e^{220.728q}e^{(c_{\mathrm{vol}}(\mathsf b)-225/M^3)N^3}\le e^{c_{\mathrm{vol}}(\mathsf b)N^3},\] where $c_{\mathrm{vol}}(\mathsf b)-225/M^3=\frac32\mathsf N_Zp_0(\mathsf b)^2+3\log(384/\varepsilon(\mathsf b)^3)$ and $225q=225N^3/M^3$; here we use the formulas that define $C_{\mathrm{vol}}$ in item~1 of \S\ref{Y-H8-1} and $c_{\mathrm{vol}}$ in \textup{(D1)} of item~D of \S\ref{Y-P1-1}. These steps rest on $1+aqe^{220q}\le e^{220q}(1+aq)\le e^{(220+a)q}$ at $a=0.728$, and on $220.728\le225$.

The moment is linear in each $g_t$ because $\gamma_t,e_t,\hat c_K$ are linear in $g_t$. Scaling each $g_t$ gives the factor $\prod_t\|g_t\|_\infty$. The moment is zero if any $g_t$ vanishes at all barycentres. A sub-collection of the sources has the same $\check\delta$ and $\check\rho$.

\item[\textup{(b)}] By \textup{(J7$^\circ$)\,(iv)} for $k=K$, $\rho_K^z(V)/Z_K(0)\equiv\exp\{\sum z_t\gamma_t+W(z)\}\rho[u_K(z);\ell_K](V)/I_1$ at every $V$. Both sides are holomorphic near $0$, bounded uniformly in $V$ (by the last part of \Cref{Y-H5-LJ}, $\mathbb T_K$ being finite), and Borel in $V$; so $\partial_S$ at $0$ is a Cauchy integral, which commutes with $\int F(V)dV$ by Fubini's theorem. The coefficient is therefore that of $\exp\{W-\sum e_tz_t\}\int F\rho[u_K(z);\ell_K]dV/I_1$. This function is holomorphic on $\mathrm{Pd}(\mathfrak r^\circ_*)$, by the theorems of Fubini and Morera, and through \textup{($*$)} it is bounded on $\mathrm{Pd}(r^\circ)$ by $e^{Q(C_W+e_h)}(1+0.182\,C_{\mathrm{vol}})\sup|F|$; Cauchy's estimate then gives the bound in \textup{(b)}. Here $\int\rho_K^zdV=Z_K(z)$, by Step~0 of the proof of \Cref{Y-H9-T}.
\end{enumerate}
\end{proof}

\subsection{Uniform bounds and limit procedures}\label{Y-P2-1}

This section proves one theorem, \Cref{Y-P2-ThR}, with the clauses \textup{(R0)} to \textup{(R8)}. They provide a subtraction $c^{(s)}_K$; a bound on moments per unit of volume, valid for every number $Q$ of sources; estimates in each cell; estimates of $\sigma$ and of $S^T_n$ that do not depend on $K$, on $N$ or on where the sources lie on the torus; and the extension to every $\mathsf b>0$ through a dilation. They also give the limits taken continuum first, volume first and jointly; the properties of every limit, namely invariance, reflection positivity, the distribution and the dilation identity; and the Hilbert space with $\Omega$, $H\ge0$ and the momenta. Finally they treat the block fields, with locality, reflection positivity without a collar, the laws in the limit and mixed bounds. The proofs occupy Lems.~\ref{Y-P2-L1} to~\ref{Y-P2-L15}.

\emph{Item 0.1: what is used of Sections~\ref{Y-H1-1} to~\ref{Y-H9-1}.} Of \Cref{Y-H1-U1} we use the clause \textup{(g1)}. Of item~1 of \S\ref{Y-H9-1} we use the table of letters, \textup{(O0)} to \textup{(O13)}, with the entries \textup{(O4a)} to \textup{(O13)} in the form in which that item lists them; the radii $r(x)$, $s(x)$, $\varrho(x)$, $\check r(x)$; the entry \textup{(O14)}; the numbers $\mathrm{sep}^{[i]}$ and $N_1^{[j]}$, which the item sets at the data $(\mathsf b,n,d,R_f)$; and $x_k$ in a run of length $K$. From the same item we take the paragraph on the sup distance of $T_l$, with the separation $d$, the radius $R_f$ and the balls, and the definition in that item of separated tuples and functionals on them, together with the letters and names of \Cref{support:150-3} that this definition states: $\mathfrak S_n(O)$, the type $(d;R)$, $\mathrm{Part}_{\ge2}$, $\sigma_S$, the mirror, reflection positivity, distributions and regularity. Of \Cref{Y-H9-T} we use the data and the clauses \textup{(T2)}, \textup{(T3)}, \textup{(T6)}, \textup{(T7)}. Further, we use the following steps of the proof of \Cref{Y-H9-T}.

\begin{enumerate}

\item[$\bullet$] Step~0, both its statement and its argument.

\item[$\bullet$] Step~1, which is \Cref{Y-H9-LN}; we use \Cref{Y-H9-LN}\,(N2), and item~(ii) of the proof of \Cref{Y-H9-LN}\,(N3).

\item[$\bullet$] Step~2, its part ``The last integral'': the data of \Cref{Y-H8-ThBC}, and the sentences there stating that items \textup{(c)} and \textup{(e)} of Step~2 of the proof of \Cref{Y-J-ThA} hold for every $K\ge1$.

\item[$\bullet$] Step~5, its statement and its argument.

\item[$\bullet$] Step~7, its parts ``At finite $K$'' and ``All $v$''.

\item[$\bullet$] Step~8\,(a), (b) and (d).

\item[$\bullet$] Step~9, with its parts ``The local construction'', ``On products'', ``Uniqueness on $Q$'', ``Gluing'' and ``The functional $\sigma_S$''.

\item[$\bullet$] Step~11\,(0), and the first sentence of Step~11\,(5).

\end{enumerate}

Of \S\ref{Y-H4-1} we use the part ``Sources'' of item~3 and \Cref{Y-H4-U4}\,(a), (h), (i). We use \Cref{Y-H5-LJ} together with its sentences on $I[u]$. Of \S\ref{Y-H6-1} we use \textup{(o3)}, \textup{(o6)} and \textup{(o7)} of item~2; the sentence on $\bar u_k$ in \Cref{Y-H6-J}\,(J6); the number $N_1(k)$ of \textup{(J8)}; and, from the proof of that proposition, Step~1\,(A) and the property $(P_k)$ of Step~3. Of item~1 of \S\ref{Y-H8-1} we use $q_N$, \Cref{Y-H8-Norms}, the number $C_{\rm vol}$ and \Cref{Y-H8-Congruence}, and of \Cref{Y-H8-ThBC} its data \textup{(B)}. We use \Cref{Y-H8-ThBC}\,(b$^C$) with its cases \textup{(ii)} and \textup{(iii)}. Of \Cref{Y-H8-ThBC}\,(c$^C$) we use the two assertions that come before \textup{(S)}: the lower bound on $I_1$, and the $L^1$ bound in the cases \textup{(i)}, \textup{(ii)} and \textup{(iii)}. We use \Cref{Y-H8-CoF}\,(F0) and (F1). We use the entries \textup{(X16)} and \textup{(X17)} of \textup{(O13)}, item~1 of \S\ref{Y-J-1}, in the form in which they enter Step~P9 of the proof of \Cref{Y-H8-CoR7}\,(R1) and the part ``The last integral'' of Step~2 of the proof of \Cref{Y-H9-T}.

\emph{Item 0.2: what is used of \S\ref{Y-J-1}, its items~2 and~3.} Of item~2 we use the part ``Lattice, cubes'', with $\mathcal D$ and the count $(26e)^{j-1}$ that is proved there; the set $X^{\approx}$ of the part ``At $x$''; the part ``Group, fields'' and the part ``Energy''; and the clauses \textup{(ca)}, \textup{(ch)}, \textup{(c$\Xi$)}. From the same item we use $\rho[u;\ell]$ and $\Xi_w$, the part ``The step'' and the part ``Runs''. Of item~3 we use \textup{(9a)}, \textup{(9b)} and \textup{(9d)}, which are clauses of \Cref{Y-H-C9}. We use \Cref{Y-J-CorC}\,(b) and (c).

\emph{Item 0.3. Wilson-measure estimates.}
We use the model of \S\ref{sec:wilson-conventions}, the board and its reflected copies from \Cref{support:4-1}, coordinate reflection positivity from \Cref{support:4-8}, and the chessboard seminorm and product estimate from \Cref{support:4-13} and \Cref{support:4-15}. \Cref{support:27-11} gives the reflection Hilbert space, and \Cref{support:14-9} describes the Gibbs laws and the torus-limit laws.

\emph{Item 0.4: what is used of \S\ref{Y-P1-1}, namely of \Cref{Y-P1-Jo}.} We use \Cref{Y-P1-Jo} with the single numeral $\mathrm{sep}^\circ:=8\log(8L^{1/2})$, for any number $Q\ge1$ of sources and for every $K\ge1$. Its data are those of \Cref{Y-P1-DATAo}: the functions $g_t$ are real and bounded with $\|g_t\|_\infty\le1$, and the sets $\{g_t\ne0\}$ lie pairwise at sup distance $\ge\check\delta$, each of them within $\check\rho$ of a point. We use four of its clauses, under the names \textup{(P1.0)} to \textup{(P1.3)}.

\begin{enumerate}

\item[\textup{(P1.0)}] This is \textup{(G$^\circ$)}: the number $j^\circ(\mathsf b,\check\delta)$ exists, and $k^\circ:=\max\{K-j^\circ,0\}\le K-1$; further $\mathfrak r^\circ_k\le\mathfrak r^\circ_{k+1}\le L^{1/2}\mathfrak r^\circ_k$, and $\mathfrak r^\circ_*:=\mathfrak r^\circ_{k^\circ}\ge r^\circ$; finally $\kappa^\circ_k\in[\kappa,\kappa+1/(4(L-1))]$.

\item[\textup{(P1.1)}] This is \textup{(J5$^\circ$)}\,(c) at the standard data together with \textup{(J6$^\circ$)}\,(a) to (d), on the representatives $u_k(z)$, $0\le k\le K$:

\begin{enumerate}

\item[$\bullet$] for $k\le k^\circ$, $u_k=u_k^{\mathrm{pol}}$ is a family on $\mathrm{Pd}(\mathfrak r^\circ_k)$ in $M_{1.91}(x_k;\kappa^\circ_k,c_h)$, at distance at most $0.91$ from $\bar u_k$ in $\|\cdot\|^{\mathrm{st},(\kappa^\circ_k)}$;

\item[$\bullet$] for $k^\circ<k\le K$, $u_k(z)=F_{\ell_{k-1}}(u_{k-1}(z))$ on $\mathrm{Pd}(\mathfrak r^\circ_*)$; this family lies in $M_1(x_k;\kappa,c_h)$, satisfies the strong cap, and stays within $0.91\cdot0.4^{k-k^\circ}$ of $\bar u_k$ in $\|\cdot\|^{\mathrm{st},(\kappa)}$;

\item[$\bullet$] for all $k$, $u_k(0)=\bar u_k$, and $\partial_Su_k(0)=\delta_S(k)$ for every $S$, where $\delta_S(k)$ are the jets of the honest family of \Cref{Y-H4-U4}\,(h); and $(P_k)$ holds.

\end{enumerate}

\item[\textup{(P1.2)}] This clause collects \textup{(J7$^\circ$)}\,(i), (iv) and (v); it is the identity at a level $k$, which holds pointwise in $V$. Thus, for $0\le k\le K$, at each field $V$ of $\mathbb T_k$ and as germs modulo $z_t^2$:

\begin{enumerate}

\item[\textup{(i)}] \[\rho_k^z(V)\equiv\exp\{\sum_{j\le k}V_j(0)+\sum_{k'<k}\sum_Y[v_{k'}(Y;u_{k'}(z))-v_{k'}(Y;\bar u_{k'})]\}\,\rho[u_k(z);\ell_k](V);\]

\item[\textup{(iv)}] \[\rho_k^z(V)\equiv\exp\{\sum_{j\le k}V_j(0)+\sum_tz_t\gamma_t^{(k)}+W^{(k)}(z)\}\,\rho[u_k(z);\ell_k](V),\] where $\gamma_t^{(k)}:=\sum_{k'<k}\partial_tV_{k'+1}(0)$ and \[W^{(k)}:=\sum_{k'<\min(k,k^\circ)}W_{k'}+\sum_{k^\circ\le k'<k}[w_{k'}-\sum_tz_t\partial_tw_{k'}(0)],\] with $W_{k'}$ and $w_{k'}$ as in \textup{(P1.3)};

\item[\textup{(v)}] \[Z_K(z)/Z_K(0)\equiv\exp\{\sum_tz_t\gamma_t^{(K)}+W^{(K)}(z)\}\,I[u_K(z)]/I[\bar u_K].\]

\end{enumerate}

Here $V_0:=0$, as in \Cref{Y-H4-U4}\,(h); hence $\gamma_t^{(k)}$ is the number $\partial_t\sum_{j\le k}V_j(0)$ of \Cref{Y-P2-L3}\,(c).

\item[\textup{(P1.3)}] This is \textup{(J8$^\circ$)}\,(a$^\circ$), (b) and (d), level by level, on $\mathrm{Pd}(\mathfrak r^\circ_*)$, where $j=K-k$ is the depth:

\begin{enumerate}

\item[$\bullet$] for $k<k^\circ$, the polynomial $W_k(z)=\sum_Y\sum_{|S|\ge2}z^S\partial_Sv_k(Y;u_k)(0)$ obeys \[|W_k|\le Q\,N_1^{[j]}C_gc_a(L^j\mathsf b)\,e^{1/(4L)}\exp\{-(\kappa-1/(4L))(\mathrm{sep}^{[j-1]}+1)\};\]

\item[$\bullet$] for $k^\circ\le k<K$, the function $w_k(z)=\sum_Y[v_k(Y;u_k(z))-v_k(Y;\bar u_k)]$ obeys $|w_k|\le3Q\,N_1^{[j]}C_gc_a(L^j\mathsf b)$ and $|\partial_tw_k(0)|\le3N_1^{[j]}C_gc_a(L^j\mathsf b)/\mathfrak r^\circ_*$;

\item[\textup{(d)}] $|W^{(K)}|\le Q\,C_W$.

\end{enumerate}

\end{enumerate}

We use \Cref{Y-P1-Cor}\,(a) and (b). We use \Cref{Y-P1-Jo}\,(J7$^\circ$)\,(vii) and (J6$^\circ$)\,(e). For $K<j^\circ$ we use \textup{(L0)}, \textup{(J5$^\circ$)}\,(c) at $k=k_1=0$, and \textup{(J6$^\circ$)} at $k^\circ=0$. \Cref{Y-P1-Jo}\,(lin) enters as a definition only: for $0\le k\le K$, $\delta[f](k)$ is the derivative $d/dz$ at $0$ of the honest run with the single source $f$, and its $h$-member $\delta^h[f](k)$ is the letter that enters the display of \textup{(R0)}. The numbers of \textup{(D1)}, item~D of \S\ref{Y-P1-1}, that this section uses are $\mathrm{sep}^\circ$, $\mathrm{sep}^{[i]}$, $j^\circ$, $\mathfrak r^\circ$, $r^\circ$, $N_1^{[j]}$, $C_W$, $e_h$, $G_1$ and $c_{\mathrm{vol}}$, and item~2 below defines each of them exactly as \textup{(D1)} does.

\emph{Item 0.5: local letters.} The letter $s$ is the depth, and $s(x)=8r(x)$; in $k_s$ it stands for the noise of the kernel, in the sense of ``The step'', item~2 of \S\ref{Y-J-1}. The letter $\theta$ is a reflection; with an exponent, it is the numeral of \Cref{Y-H9-T}\,(T2), (T3). The letter $z$ denotes the variables of the sources; it is also the integer of a dilation, a bound variable, with $z(\mathsf b)$ as in the setting \textup{(S7)} of \S\ref{Y-00-1}; and $z(t)$ normalises the kernel.

\emph{Item 1: data.} The group is $G=\operatorname{SU}(2)$, and the dimension is $3$. The table \textup{(O0)} to \textup{(O13)} and its floors are exactly those of item~(1) of \S\ref{Y-H0-1}, which is item~1 of \S\ref{Y-H9-1}. The blocking factor $L$ is odd, and $M=L^{m'}$.

We recall the data of the setting \textup{(S2)} of \S\ref{Y-00-1}: a label $\mathsf b\ge x_s$; a side $l=N=2L^{m_l}$ with $m_l\ge m'$, subject to \textup{(Tor)}, that is, $N\ge2M\,s(L\mathsf b)$; an integer $K\ge1$; and $\beta=L^K\mathsf b$, $a_K=L^{-K}$. We write $\mu_{K,N,\beta}$ for Wilson's measure on $a_K\mathbb Z^3/N\mathbb Z^3$ and set $\mathcal O_K(f)=a_K^{-1}\sum_pf(x(p))\,a_p(U)$, where $a_p=1-\tfrac12\operatorname{tr}\partial U(p)$ and $x(p)$ is the barycentre of $p$. The joint cumulant is $S^T_n$, and $\sigma_{K,N,\beta}(X)=\langle\prod_i(\mathcal O_Kf_i-\langle \mathcal O_Kf_i\rangle)\rangle$.

The level $k$ of the run is the lattice $L^{-(K-k)}\mathbb Z^3/N\mathbb Z^3$, whose depth is $K-k$. The point $0$ is a site of every level, and the last lattice has spacing $1$. The fields $(V_0=U,V_1,\dots,V_K)$ carry the joint law of Step~11\,(0) of the proof of \Cref{Y-H9-T}. Without an argument, $V_k$ is the field drawn at the level $k$, with $V_0=U$, whereas $V_j(z)$, $V_j(0)$ and the value $V_0:=0$ are numbers of \Cref{Y-H4-U4}\,(h), not these fields; and $V_0^+$ is the number of \Cref{Y-H1-U1}\,(g1). All distances are sup distances, and the support of a bounded function $g$ means the set $\{g\ne0\}$.

\emph{Boards.} For $-m_l\le i\le K-1$ and $o\in a_K\mathbb Z^3$, the board $B_i(o)$ has its walls at $o_\mu+L^{-i}(\mathbb Z+\tfrac12)$ and its cells of side $L^{-i}$; we put $B_i:=B_i(0)$. In the letters with which \S\ref{sec:chessboard-torus} writes a board, $B_i(o)=(a_K,\ \lambda=L^{K-i},\ \mathsf n_\mu=2L^{m_l+i},\ o)$.

No hypothesis enters beyond the data and the table.

\emph{Item 2: numbers.} The numbers are fixed in this order: first the table, then $\mathsf b$, then \textup{(O14)}. No floor depends on the numbers below, and all of these numbers are locally bounded in $\mathsf b$. None of them depends on $N$, on $K$, on the number of the sources or on their places, except as follows: $C_n^{(2)}$ and its $\check\delta_n$ depend on $n$, the number of the entries, but on none of $N$, $K$ and the places; $\mathfrak r^\circ_k=\mathfrak r^\circ(x_k)$ depends on $K$ and on $k$, but only through $x_k$; and the subtractions $c^{(s)}_K$ and $\hat c_K$ at the end of the item depend on $K$ and on $N$.

We set $\mathrm{sep}^\circ=8\log(8L^{1/2})$ and $\mathrm{sep}^{[i]}=\lfloor\check\delta\,L^i/M\rfloor-2\,\check r(L^i\mathsf b)-2$, and we let $j^\circ(\mathsf b,\check\delta)$ be the least $j\ge1$ with $\mathrm{sep}^{[j-1]}\ge\mathrm{sep}^\circ$. Next, $\mathfrak r^\circ(x)=3.5/\mathrm{amp}(Lx)$, $\mathfrak r^\circ_k=\mathfrak r^\circ(x_k)$, $r^\circ=\mathfrak r^\circ(L^{j^\circ}\mathsf b)$ and $r^\circ_s=\mathfrak r^\circ(L^{\max(s,j^\circ)}\mathsf b)$, so that $r^\circ_0=r^\circ$. We put $N_1^{[j]}=(2\,\check\rho\,L^{j-1}/M+2\,\check r(L^{j-1}\mathsf b)+3)^3$ and \[C_W(\mathsf b,\check\delta,\check\rho)=\sum_{j>j^\circ}2e\,N_1^{[j]}C_gc_a(L^j\mathsf b)\,e^{-(\kappa-1)(\mathrm{sep}^{[j-1]}-1)}+\sum_{1\le j\le j^\circ}6N_1^{[j]}C_gc_a(L^j\mathsf b);\] since, for the terms with $j>j^\circ$, we have $\mathrm{sep}^{[j-1]}\ge0$ and the inequalities $e^{1/(4L)}\le2e$ and $(\kappa-1/(4L))(\mathrm{sep}^{[j-1]}+1)\ge(\kappa-1)(\mathrm{sep}^{[j-1]}-1)$, this number dominates the bounds of \textup{(P1.3)} term by term.

Further, $e_h=e_h^{(0)}=0.45\,N_1^{[1]}c_h(\mathsf b)$, and $e_h^{(s)}=1.11\,N_1^{[s+1]}c_h(L^s\mathsf b)$ for $s\ge1$. We set $G_1(\mathsf b,\check\delta,\check\rho)=e^{C_W+e_h}/r^\circ$ and $G_{1,s}=e^{C_W+e_h^{(s)}}/r^\circ_s$. The letter $c_{\mathrm{vol}}$ stands for $c_{\mathrm{vol}}(\mathsf b)=225/M^3+\tfrac32\mathsf N_Zp_0(\mathsf b)^2+3\log(384/\varepsilon(\mathsf b)^3)$; hence $1+0.182\,C_{\rm vol}\le1+C_{\rm vol}\le2C_{\rm vol}\le e^{c_{\mathrm{vol}}N^3}$. For $s\ge1$ we set $c_{\mathrm{vol}}^{(s)}(\mathsf b)=L^{3s}(c_{\mathrm{vol}}(L^s\mathsf b)+200/M^3)$, and $c_{\mathrm{vol}}^{(0)}=c_{\mathrm{vol}}$.

We set $B_0(\mathsf b,\check\delta)=6L^{j^\circ}(\check\delta\,L^{j^\circ}/2+2)^3$; this is a number, whereas the board $B_0$ carries no argument. Then $G(\mathsf b,\check\delta)=\max\{G_1(\mathsf b,\check\delta,\check\delta/4)\,e^{c_{\mathrm{vol}}(\mathsf b)},\,B_0(\mathsf b,\check\delta)\}$. With $\check\delta_n:=\min\{d,1/(4n)\}$ we define $C_n^{(2)}(\mathsf b;d,R_f):=1.45^{n+1}\,n!\,(4R_f/\check\delta_n+2)^{3n}\,G(\mathsf b,\check\delta_n)^n$.

\emph{Subtractions.} For $0\le s\le K-1$ we set $c^{(s)}_K(f):=\frac{d}{dz}\log\rho_{K-s}^{[zf]}(1)\big|_{z=0}$, where $\rho_k^{[zf]}$ is the density drawn $k$ times from Wilson's weight multiplied by $e^{z\mathcal O_K(f)}$, taken at the trivial field of $\mathbb T_k$. We write $\hat c_K:=c^{(0)}_K$.

\emph{Item 3: the statement.}

\begin{theorem}[Uniform bounds, limits in three orders, reconstruction and block fields]\label{Y-P2-ThR} In every clause the objects, and the letters that denote them, are those of item~1, and the numbers are those of item~2. The clauses \textup{(R0)} to \textup{(R3)} and \textup{(R8)} assume the data of item~1, with $\mathsf b\ge x_s$, $m_l\ge m'$ and \textup{(Tor)}; the clauses \textup{(R4)} to \textup{(R7)} carry their own data, every $\mathsf b>0$ and, in \textup{(R4)}, every side $l=2L^m$. Under these assumptions the following clauses hold.

\begin{enumerate}

\item[\textup{(R0)}] \emph{Subtraction.}

\begin{enumerate}

\item[$\bullet$] The number $c^{(s)}_K(f)$ is the mean of $\mathcal O_K(f)$ under Wilson's measure weighted by the kernel density at $V_{K-s}=1$; it is real for real $f$, and it is linear in $f$.

\item[$\bullet$] $c^{(s)}_K(f\circ R)=c^{(s)}_K(f)$ for every map $R:y_\mu\mapsto c-y_\mu$ with $c\in L^{-s}\mathbb Z$, and for every translation by a vector of $L^{-s}\mathbb Z^3$. For $\hat c_K$ the invariance holds at every integer $c$, even or odd; the mirrors with odd $c$ are those in the walls of $B_0$.

\item[$\bullet$] For the data of the setting \textup{(S2)} and $0\le s\le K-1$, \[c^{(s)}_K(f)=\frac{d}{dz}\Big|_{z=0}\sum_{1\le j\le K-s}V_j(z)+\sum_X\delta^h[f](K-s)(X,1),\] where $X$ runs over $\mathcal D$ of the level $K-s$, and $V_j(z)$ are the numbers of \Cref{Y-H4-U4}\,(h) along the run with the single source $f$ and the variable $z$. If $f$ is of the type of one source of \textup{(R1)}, then, at any $\check\delta>0$, the modulus of the last sum, multiplied by $r^\circ_s$, is $\le e_h^{(s)}\|f\|_\infty$.

\end{enumerate}

\item[\textup{(R1)}] \emph{Per unit of volume.} Let $K\ge1$, let $Q\ge1$ be arbitrary, and let $g_t$ be real bounded functions on $T_N=\mathbb R^3/N\mathbb Z^3$ whose supports lie pairwise at distance $\ge\check\delta$, each of them inside a ball of radius $\check\rho$. Then \[|\langle\prod_t(\mathcal O_Kg_t-\hat c_Kg_t)\rangle|\le G_1^Q\,(1+0.182\,C_{\rm vol}(N,\mathsf b))\prod_t\|g_t\|_\infty.\] Here $C_{\rm vol}$ enters with the exponent $1$, whatever $Q$ is.

\item[\textup{(R2)}] \emph{Cells.} Let $B_i(o)$ be a board with $-m_l\le i\le0$ and $o\in\mathbb Z^3$, let $h=L^{-i}$, and take the data of \textup{(R1)}; suppose in addition that every support lies inside one open cell, at distance $\ge\check\delta/2$ from the walls, and let $q$ cells be occupied. Then the moment of \textup{(R1)} is $\le G_1^Q\,e^{q\,c_{\mathrm{vol}}(\mathsf b)h^3}\prod_t\|g_t\|_\infty$.

\begin{enumerate}

\item[\textup{(i)}] On $B_0$, whose walls are $\mathbb Z+\tfrac12$ and where $h=1$, the factor is $e^{q\,c_{\mathrm{vol}}(\mathsf b)}$. At $Q=1$ this gives $|\langle \mathcal O_Kg\rangle-\hat c_Kg|\le G_1e^{c_{\mathrm{vol}}}\|g\|_\infty$.

\item[\textup{(ii)}] \emph{Placed anywhere.} Let $n\ge1$ sources of the type $(d;R_f)$ be placed anywhere, and let $h=L^{-i}>n(2R_f+d+1)$ and $N\ge2h$. Then some $o\in\mathbb Z^3$ meets the condition above at $\check\delta=d$, $\check\rho=R_f$; hence the moment is $\le G_1(\mathsf b,d,R_f)^n\,e^{n\,c_{\mathrm{vol}}h^3}\prod_i\|f_i\|_\infty$.

\end{enumerate}

\item[\textup{(R3)}] \emph{Uniform bounds.} Let $n\ge2$, $d>0$ and $R_f>0$, and let $f_i$ be real bounded functions whose supports lie pairwise at distance $\ge d$, each of them inside a ball of radius $R_f$, wherever these balls lie. Recall from item~2 that $\check\delta_n=\min\{d,1/(4n)\}$. Then, for all $K\ge1$ and all $N$ of the data, \[|\sigma_{K,N,\beta}(f)|\le(4R_f/\check\delta_n+2)^{3n}\,(2G(\mathsf b,\check\delta_n))^n\prod_i\|f_i\|_\infty,\] and $|S^T_n(f)|\le C_n^{(2)}\prod_i\|f_i\|_\infty$, where $C_n^{(2)}(\mathsf b;d,R_f)=1.45^{n+1}\,n!\,(4R_f/\check\delta_n+2)^{3n}\,G(\mathsf b,\check\delta_n)^n$ is the number of item~2. The limits of \Cref{Y-H9-T}\,(T3) obey the same bounds.

\item[\textup{(R4)}] \emph{Every $\mathsf b>0$, every side $l=2L^m$.} Let $n\ge2$, $d>0$ and $R_f>0$, and let $f_i$ be real bounded functions on $T_l=\mathbb R^3/l\mathbb Z^3$ of the type $(d;R_f)$, as in \textup{(R3)}: their supports lie pairwise at distance $\ge d$, each of them inside a ball of radius $R_f$, wherever these balls lie. Let $z\in\mathbb Z$ be such that $L^z\mathsf b\ge x_s$, $K-z\ge1$, $m+z\ge m'$ and $L^zl\ge2M\,s(L^{z+1}\mathsf b)$. Then \[|S^T_{n;K,l,L^K\mathsf b}(f)|\le L^{zn}\,C_n^{(2)}(L^z\mathsf b;L^zd,L^zR_f)\prod_i\|f_i\|_\infty,\] and the same bound holds for $\sigma$.

\item[\textup{(R5)}] \emph{Limits, for every $\mathsf b>0$.} Let the sides $l_j=2L^{m_j}$ tend to $\infty$. In \textup{(a)} to \textup{(d)}, the cumulants and their limits are evaluated at tuples $f$ whose entries $f_1,\dots,f_n$ are continuous functions on $\mathbb R^3$ with compact and pairwise disjoint supports, periodised to $T_l$ as soon as $l$ exceeds twice the diameter of the union of the supports.

\begin{enumerate}

\item[\textup{(a)}] \emph{Continuum first.} At each side, $\lim_K$ exists as $K\to\infty$ through all integers $K$, by \Cref{Y-H9-T}\,(T3), (T6). From every sequence of sides one can extract a subsequence on which $S^{(l_j,\mathsf b)}_n(f)$ converges simultaneously for all $n\ge2$ and all such tuples. The same is true for labels $\mathsf b_j\to \mathsf b>0$.

\item[\textup{(b)}] \emph{Volume first.}

\begin{enumerate}

\item[$\bullet$] One subsequence of sides suffices for all $K$, and along it $\mu_{K,l_j,L^K\mathsf b}$ converges weakly to a state $\nu_K$ on $G^{\text{bonds of }a_K\mathbb Z^3}$.

\item[$\bullet$] $\nu_K$ is a Gibbs state of Wilson's action, and it is invariant under $a_K\mathbb Z^3$ and under $W_3$; so it lies in $\mathbb G_K(\mathsf b)$, the set of the setting \textup{(S6)} of \S\ref{Y-00-1}.

\item[$\bullet$] $\nu_K$ is reflection positive both with respect to planes of sites and with respect to the planes that lie midway between sites.

\item[$\bullet$] The cumulants of $\nu_K$ are the limits as $j\to\infty$, along the subsequence of the first bullet, of the cumulants $S^T_{n;K,l_j,L^K\mathsf b}(f)$ of $\mu_{K,l_j,L^K\mathsf b}$, and they obey \textup{(R3)} and \textup{(R4)}; these cumulants of $\nu_K$ then converge along a subsequence of the cutoffs $K$.

\end{enumerate}

\item[\textup{(c)}] \emph{Jointly.} Every sequence $(K_j,l_j)$ in which both entries tend to $\infty$ has a convergent subsequence; the same is true for labels $\mathsf b_j\to \mathsf b>0$.

\item[\textup{(d)}] If $K_j/l_j^3\to\infty$, the limit points of \textup{(c)} coincide with those of \textup{(a)} along the same sides $l_j$ and the same fixed or convergent labels.

\end{enumerate}

\item[\textup{(R6)}] \emph{Properties of every limit.} In \textup{(R6)} and \textup{(R7)}, the definition in item~1 of \S\ref{Y-H9-1} is taken with $\mathbb R^3$ in the place of $T_l$: a mirror is $(\mu,c)$ with $c$ real, and its open side is $\{y_\mu>c\}$. Every limit $S=(S_n)_{n\ge2}$ of \textup{(R5)} has the following properties.

\begin{enumerate}

\item[\textup{(i)}] $S_n$ is real and symmetric, and it is $n$-linear on every product of $n$ spaces in which each factor is the space of the continuous functions that vanish off one compact set, the $n$ sets being pairwise disjoint.

\item[\textup{(ii)}] $S$ obeys \textup{(R3)} and \textup{(R4)}, uniformly with respect to where the supports of the entries lie. For labels $\mathsf b_j\to \mathsf b$ the constant is the supremum over $U(\mathsf b)$, at $z=z(\mathsf b)$, with $U(\mathsf b)$ and $z(\mathsf b)$ as in the setting \textup{(S7)} of \S\ref{Y-00-1}.

\item[\textup{(iii)}] $S$ is invariant under all translations of $\mathbb R^3$ and under $W_3$.

\item[\textup{(iv)}] $\sigma_S$ is reflection positive in the sense of \eqref{support:9-28} at every coordinate mirror $(\mu,c)$.

\item[\textup{(v)}] On the open set of the points of $(\mathbb R^3)^n$ whose entries are pairwise different there is exactly one distribution whose value at $f_1\otimes\cdots\otimes f_n$ equals $S_n(f_1,\dots,f_n)$ whenever $f_1,\dots,f_n$ are smooth functions with pairwise disjoint compact supports, and on every compact subset of that open set this distribution has order $\le3n+1$. Moreover $\sigma_S$ is regular.

\item[\textup{(vi)}] The identity $S^{\mathsf b}_n(f)=L^{zn}S^{L^z\mathsf b}_n(f^{[z]})$ holds along matched subsequences, that is, if the limit for the label $\mathsf b$ is taken through the sides $l_j$ and the cutoffs $K$, and the limit for the label $L^z\mathsf b$ through the sides $L^zl_j$ and the cutoffs $K-z$, with the same order of the limits in both cases; when the labels move, it holds for each member of the sequence, by Step~5 of the proof of \Cref{Y-H9-T}.

\item[\textup{(vii)}] The group $\{R\in \operatorname{O}(3):\sigma_S(RX)=\sigma_S(X)\text{ for all }X\}$ is closed, and it is $W_3$ or $\operatorname{O}(3)$.

\end{enumerate}

\item[\textup{(R7)}] \emph{Reconstruction, in each coordinate direction $\mu$.} Let $D$ consist of the finite linear combinations, with complex coefficients, of the tuples of $\mathfrak S(\mathbb E_+)$, where $\mathbb E_+=\{y_\mu>0\}$, and equip it with $\langle F,F'\rangle=\sum\bar\zeta_i\zeta'_k\,\sigma_S(\theta X_i\cdot X'_k)$. The completion $\mathcal H(S)$ of $D/\text{null}$ is separable, and $\Omega=[\text{empty tuple}]$ has $\|\Omega\|=1$. The maps $\mathcal T_t[F]=[\tau_{te_\mu}F]$ form a strongly continuous semigroup of self-adjoint contractions $e^{-tH}$, with $H\ge0$ and $H\Omega=0$. The translations of the mirror plane act by strongly continuous unitary operators $U(v)$, with $U(v)\Omega=\Omega$ for all $v$ in that plane; their momenta $P_1,P_2$ satisfy $P_j\Omega=0$, and $H,P_1,P_2$ commute strongly; the joint spectrum lies in $[0,\infty)\times\mathbb R^2$, and $W_2$ acts by a unitary representation, where $W_2:=\{w\in W_3:\ we_\mu=e_\mu\}$. The class $[X]$ is defined also for tuples whose entries are continuous with compact disjoint supports in $\mathbb E_+$, with $\langle[X],[Y]\rangle=\sigma_S(\theta X\cdot Y)$, and $\langle\Omega,[f]\rangle=0$.

\item[\textup{(R8)}] \emph{Block fields.} In the setting of item~1, with $\mathsf b\ge x_s$, let $s\ge0$ be the depth and let $K\ge s+1$.

\begin{enumerate}

\item[\textup{(a)}] \emph{Locality.} Let $O_1,\dots,O_m$ be pairwise disjoint open sets of $T_N$ bounded by planes $y_\mu\in L^{-s}(\mathbb Z+\tfrac12)$, and let $F_i$ be a bounded Borel function of $V_{K-s}$ that depends only on the bonds with both ends in $O_i$. Then $E[F_i|U]$ is bounded and Borel, and it depends on $U$ only through the bonds with both ends in $O_i$. Moreover $E[\prod_iF_i|U]=\prod_iE[F_i|U]$, and $E[RF_i|U]=R\,E[F_i|U]$ for the maps $R$ of \textup{(R0)} at depth $s$.

\item[\textup{(b)}] \emph{Reflection positivity without a collar.} Let $c\in L^{-s}(\mathbb Z+\tfrac12)$ and $T_+=\{c<y_\mu<c+N/2\}$, and let $F$ be a bounded Borel function of the values of $V_{K-s}$ on the bonds that have both ends in $T_+$. Then $E[\overline{\theta F}\,F]\ge0$.

\item[\textup{(c)}] \emph{Limits.} Along a further subsequence of sides, and for every $s\ge0$, the laws $\mu^\infty_{[s]}$ of \Cref{Y-J-CorC}\,(b) converge, locally in total variation, to a law of the fields on $L^{-s}\mathbb Z^3$. This limit has the following properties.

\begin{enumerate}

\item[$\bullet$] It is gauge invariant.

\item[$\bullet$] It is invariant under $L^{-s}\mathbb Z^3$ and under $y_\mu\mapsto c-y_\mu$, $c\in L^{-s}\mathbb Z$.

\item[$\bullet$] For every finite set $B$ of bonds of $L^{-s}\mathbb Z^3$, its marginal on $B$ has a density, and that density lies between $(z^{-1}e^{-2t})^{\#B}$ and $z^{-\#B}$, where $t=aL^{s+1}\mathsf b$ and $z=z(t)$.

\item[$\bullet$] It is consistent under the drawn step with the noise $aL^s\mathsf b$.

\item[$\bullet$] It is reflection positive as in \textup{(b)}, with $T_+=\{y_\mu>c\}$.

\end{enumerate}

From every sequence $(K_j,N_j)$ with $N_j\to\infty$ one can extract a subsequence along which the same holds for all $s\le\liminf K_j-1$.

For $\mathsf b<x_s$, the dilation by $L^z$ gives the lattices $L^{-(s+z)}\mathbb Z^3$ only.

\item[\textup{(d)}] \emph{Mixed bounds.}

\begin{enumerate}

\item[\textup{(i)}] \emph{Without cells.} Let $F$ be a bounded Borel function of $V_{K-s}$, and take the data of \textup{(R1)}. Then \[|\langle F(V_{K-s})\prod_t(\mathcal O_Kg_t-c^{(s)}_Kg_t)\rangle|\le\sup|F|\;G_{1,s}^Q\,e^{c_{\mathrm{vol}}^{(s)}N^3}\prod_t\|g_t\|_\infty.\] At $s=0$ the factor $1+0.182\,C_{\rm vol}$ replaces the exponential.

\item[\textup{(ii)}] \emph{With cells.} Let $B_i(o)$ be a board with $-m_l\le i\le s$ and $o\in L^{-s}\mathbb Z^3$, let $h=L^{-i}$, and take sources as in \textup{(R2)}. Let the factors $F_{\mathrm{cell}}$ be functions of $V_{K-s}$ on those bonds that have both ends in the open cell, and let $q$ cells be occupied, each by a source or by a factor $\ne1$. Then the moment with these factors is bounded by $\prod\sup|F_{\mathrm{cell}}|\;G_{1,s}^Q\,e^{q\,c_{\mathrm{vol}}^{(s)}h^3}\prod_t\|g_t\|_\infty$. At $s=0$ an inner coarse bond exists only for $h\ge L$, while at $s\ge1$ there is one already in unit cells.

\end{enumerate}

\end{enumerate}

\end{enumerate}

\end{theorem}


The proof first estimates the subtracted moments at fixed volume.
Chessboard bounds then replace the volume factor by a product over
occupied cells. The resulting uniform bounds give compactness of the
correlations and allow reflection positivity to pass to the limits.

\emph{Statements and proofs of other sections, and how they are used.} The statements named below are used as stated, and any change in their use is given in the paragraph that names them. The data of \Cref{Y-H9-T} and Step~0 of its proof, for $Q$ sources: $Z_K(z)$ is entire; the weight is that of \Cref{Y-H4-U4} at $\nu_p=\mathsf b^{-1}\sum_tz_tg_t(x(p))$; since the supports are disjoint, $|\nu|_\infty\le\max_t|z_t|\,\|g_t\|_\infty/\mathsf b$; and $\int\rho_k^zdV=Z_K(z)$ at every level. The first three of these assertions are \Cref{Y-P2-L1}; the fourth is the last paragraph of Step~0 of the proof of \Cref{Y-H9-T}.

The bound \textup{(R3)} replaces \Cref{Y-H9-T}\,(T1): where \textup{(T1)} has $C_n=\mathrm{Ser}_1(\mathsf b,n,d,R_f)+\log2\,(\tau_7/\mathfrak r(L^{j^*}\mathsf b))^n$, $\tau_7=\max\{1,0.91\,C_{\rm vol}(N,\mathsf b)\}$, the bound \textup{(R3)} has the constant $C_n^{(2)}$, for all $N$ and all $K\ge1$. Of Step~2 of the proof of \Cref{Y-H9-T} we use only the verification of the data \textup{(B)} of \Cref{Y-H8-ThBC} at $x=\mathsf b$ for every $K\ge1$, namely the part ``The last integral'' with its sentences on every $K\ge1$; no line of this verification depends on the sources.

\emph{\Cref{Y-H9-T}\,(T2).} In the constant $C'_n$ of \textup{(T2)}, only $3C_{\rm vol}(N,\allowbreak \mathsf b)C_{\mathsf b}(1.15+0.46C_5^\#/K_5)(4C_{\rm vol}(N,\allowbreak \mathsf b)/\mathfrak r(L^{j^*}\mathsf b))^n$ depends on $N$, so that $\log C'_n\le(n+1)c_{\mathrm{vol}}N^3+O(1)$; and, as \textup{(T2)} states, $K_T$ does not depend on $N$.

\emph{\Cref{Y-H9-T}\,(T3) and (T6).} At a fixed side, the limit as $K\to\infty$ through all integers $K$ exists as those clauses state, also for continuous entries. Moreover this limit obeys \textup{(R3)} and \textup{(R4)}, being a limit of numbers that obey these bounds.

\emph{Step~7 of the proof of \Cref{Y-H9-T}.} Its part ``At finite $K$'' gives $\mathcal O_K(f)(rU)=\mathcal O_K(f\circ r)(U)$ for every bounded $f$. Its part ``All $v$'' holds with $C_n^{(2)}$ at $(d/2;R_f+d/4)$. Here $C_n^{(2)}$ at $(d;R)$ denotes the constant of \textup{(R6)\,(ii)}, namely $C_n^{(2)}(\mathsf b;d,R)$ at a fixed $\mathsf b\ge x_s$, $L^{zn}C_n^{(2)}(L^z\mathsf b;L^zd,L^zR)$ at a fixed $0<\mathsf b<x_s$, and the supremum over $U(\mathsf b)$ at $z(\mathsf b)$ for moving labels. This constant is free of $K$ and of $l$, and so the part ``All $v$'' applies to all three orders.

\emph{Step~8\,(a), (b), (d) of the proof of \Cref{Y-H9-T}.} Part \textup{(a)} is taken at finite $(K,N)$ and $c=0$; it enters the proof of \textup{(R6)}\,(iv) in each of the three orders, since every compact support in $\{y_\mu>0\}$ is contained in $\{0<y_\mu<l/2\}$ once $l$ is large; and from \textup{(R6)}\,(iii) we then obtain \textup{(d)}.

\emph{Step~9 of the proof of \Cref{Y-H9-T}.} In its part ``The local construction'' the Fourier series is taken with the period $4\delta$ about $\xi_i$, in the place of the period $l$. The nearest copy of $\mathrm{supp}\,F$ begins at the distance $3\delta$, and $\chi_i$ ends at $3\delta/2$; hence $F=(\chi\times\dots\times\chi)F^{\mathrm{per}}$. The constant $c(n,l)$ becomes $c(n,4\delta)$, and $C_n(\delta)$ becomes $C_n^{(2)}$ at $(\delta;3\delta/2)$, the constant of \textup{(R6)\,(ii)}. The parts ``On products'', ``Uniqueness on $Q$'' and ``Gluing'' are taken word for word.

Step~5 of the proof of \Cref{Y-H9-T}: the identity $S^T_{n;K,l,\beta}(f)=L^{zn}S^T_{n;K-z,L^zl,\beta}(f^{[z]})$ carries \textup{(R3)} to every $\mathsf b>0$. The side drops out, since $C_n^{(2)}$ depends on no side, and the conditions on $z$ are those of \Cref{Y-H9-T}\,(T6).

\emph{\Cref{Y-H9-T}\,(T7), \Cref{Y-J-CorC}\,(b), (c), and Step~11\,(0), (5) of the proof of \Cref{Y-H9-T}.} The clauses \textup{(b)} and \textup{(c)} of that corollary apply at each $N$. The rate $C_{\rm vol}(NL^s,\dots)\theta^{K-s}$ is lost as $N$ grows, and local densities that are equi-Lipschitz take its place; they come from the kernel alone (\Cref{Y-P2-L15}). Step~11\,(0) gives the joint law and the identity $\int F\rho^z_{K-s}dV/Z_K(z)=\langle Fe^{\sum_tz_t\mathcal O_K(g_t)}\rangle/\langle e^{\sum_tz_t\mathcal O_K(g_t)}\rangle$, and the first sentence of Step~11\,(5) gives the data \textup{(B)} at $(NL^s,L^s\mathsf b)$.

\emph{The part ``Sources'' of item~3 of \S\ref{Y-H4-1} and \Cref{Y-H4-U4}\,(a), (h), (i).} Clause \textup{(h)} states that $\rho_k^z=\exp\{\sum_{j\le k}V_j(z)\}\rho[u_k(z);\ell_k]$ at every field of $\mathbb T_k$; it states also that $\bar u_k$ is a system of members and that $\rho[\bar u_k]>0$. In \textup{(i)} the numerals are free of $n$.

\emph{\Cref{Y-H1-U1}\,(g1), the convention $V_0:=0$ of \Cref{Y-H4-U4}\,(h), and \textup{(X16)}, \textup{(X17)}.} By \textup{(g1)}, $\mathsf V'(u)=V_0^++\sum_Yv(Y;u)$; so $V_{k'+1}(0)=V_0^++\sum_Yv_{k'}(Y;\bar u_{k'})$, and together with $V_0:=0$ this is how the exponent of \textup{(P1.2)}\,(i) is computed. By \textup{(X16)} and \textup{(X17)}, Step~P9 of the proof of \Cref{Y-H8-CoR7}\,(R1) has $\varepsilon\le1/40$ and $c_a+c_h\le1/32$ at $x=L^s\mathsf b$, $\mathsf b\ge x_s$.

\emph{Of \S\ref{Y-H6-1}, \textup{(o3)}, \textup{(o6)}, \textup{(o7)}, the sentence of \Cref{Y-H6-J}\,(J6) on $\bar u_k$, the number $N_1(k)$, Step~1\,(A) and the property $(P_k)$ of Step~3.} By \textup{(o6)}, the $h$-norm is a supremum over $X$ and the field, in units $c_h(x)e^{-\tilde\kappa(\#X-1)}$. By \textup{(A)}, at most $N_1(k-1)$ cubes of $\mathbb T_k$ lie in $S_t(k)^{+\check r(x_k)}$. By $(P_k)$, $u(z)(X)$ does not depend on $z_t$ if $X^{+\check r(x_k)}$ misses $S_t(k)$.

In the place of \Cref{Y-H6-J}\,(J5) to (J8) we use \Cref{Y-P1-Jo}, for any number $Q\ge1$ of sources and clause by clause, as \textup{(P1.0)} to \textup{(P1.3)}, at $\mathrm{sep}^\circ=8\log(8L^{1/2})$. Its radii are free of $Q$; it provides representatives at all levels; it gives the congruence at a level $k$, pointwise in $V$; and it treats the vacuum sums per source, level by level.

\emph{\Cref{Y-H8-ThBC} with its data \textup{(B)}, \textup{(b$^C$)}\,(ii), (iii), and the two assertions of \textup{(c$^C$)} that come before \textup{(S)}.} The clause that carries the volume is this: in the cases \textup{(i)} and \textup{(ii)}, $\int|\rho[u_2;\ell]-\rho[u_3;\ell]|dV\le\tfrac12C_{\rm vol}\mathsf D_{23}I_1$, and in the case \textup{(iii)} the same holds with $C_{\rm vol}e^{219(A-1)q_N}$. It is used at $u_3=u_1=\bar u_k$, $u_2=u_k(z)$, and at $x=\mathsf b$ as well as at $x=L^s\mathsf b$ on $\mathbb Z^3/NL^s\mathbb Z^3$. The case \textup{(ii)} of \textup{(b$^C$)} is applied with the rate $\kappa$ and the cap $c_h$, the case \textup{(iii)} with the rate $\kappa^\circ_k$ and the cap $c_h$.

\emph{Defs.~\ref{Y-H8-Norms} and~\ref{Y-H8-Congruence}, \Cref{Y-H8-CoF}\,(F0), (F1) and \Cref{Y-H5-LJ}.} Besides $I[u]$, of which these statements speak, we use the functional $\Phi_F[u]=\int F\rho[u;\ell]dV$, with $F$ bounded Borel. This functional is holomorphic along a family by the proof of \Cref{Y-H8-CoF}\,(F1), which is Step~P4 of the proof of that corollary: pointwise holomorphy, together with the sup bound of \Cref{Y-H8-ThBC}\,(b$^C$).

\emph{The part ``Energy'' of item~2 of \S\ref{Y-J-1}, the clauses \textup{(ca)}, \textup{(ch)}, \textup{(c$\Xi$)}, and $\rho[u;\ell]$, $\Xi_w$, at the field $1$.} There $\phi_p=0$, so that $S_\varepsilon(0)=0$, $\zeta=0$ and $\mathrm{En}_\ell(1)=0$. The set $D(1)$ is empty, so that only the empty family is admitted and $\Xi_w(1)=1$. Further, $1_{X^{\approx}}(1)=1$, and the clause \textup{(ca)} includes the condition that the value at pure gauges is $0$.

\emph{The parts ``The step'' and ``Runs'' of item~2 of \S\ref{Y-J-1}.} This section proves two further properties of the drawn step. The first, \Cref{Y-P2-L2}, gives $\mathsf M(RU)=R\,\mathsf M(U)$ and the invariance of the kernel, for integer mirrors and coarse translations. The second, \Cref{Y-P2-L6}\,A, states that $\mathsf M(U)(c)$ depends on the bonds of $B(y)\cup B(y')$ alone. The proof given in ``The step'' yields gauge covariance only. ``The step'' states also that the step depends on $x$, $a$ and $L$, and on no rows, no grid of cubes, no members and no $K$; hence the covariance iterates. By ``Runs'', the noise of the step $k\to k+1$ is $a\,x_k$.

\emph{The clauses \textup{(9a)}, \textup{(9b)}, \textup{(9d)} of \Cref{Y-H-C9}, in item~3 of \S\ref{Y-J-1}.} Clause \textup{(9a)} holds at every $V'$. By \textup{(9b)}, the outputs keep all the clauses that carry no size, so that $E^a(X,1)=0$; and an input in $M_2$ gives an output in $M_1$ with the strong cap. Clause \textup{(9d)} gives positivity.

\emph{\Cref{support:4-1}\,(2) and the paragraphs of \S\ref{sec:chessboard-torus} that follow it and that item~0.3 names.} $B_i(o)$ is a board: $\lambda=L^{K-i}$ is odd and $\ge3$ if and only if $i\le K-1$; $\mathsf n_\mu=2L^{m_l+i}$ is an even integer if and only if $i\ge-m_l$; and $o\in a_K\mathbb Z^3$. The measure, $a_p$, $x(p)$, $\mathcal O(f)$, $(rU)(b)=U(r^{-1}b)$ for a bond $b$, $(rF)(U)=F(r^{-1}U)$ and $rf=f\circ r^{-1}$ coincide with the objects so named in the data of \Cref{Y-H9-T} and in Step~7 of its proof.

\Cref{support:4-8} is stated for all bounded Borel functions of the bonds that have both ends in $T_+$, with $c\in a(\mathbb Z+\tfrac12)$.

\Cref{support:4-13}, \Cref{support:4-15} and \Cref{support:4-17}\,(f) are applied on $B_i(o)$ and over the occupied cells. With $N_c=(N/h)^3$, the number of the cells of $B_i(o)$, the root $1/N_c$ turns $e^{c_{\mathrm{vol}}N^3}$ into $e^{c_{\mathrm{vol}}h^3}$.

\emph{\Cref{support:27-11}.} Its hypotheses are a semigroup, commuting maps, the fixing of $\mathsf 1$, and \textup{(1)} to \textup{(3)}. The conclusions \textup{(a)} to \textup{(d)} of \Cref{support:27-11} are \textup{(R7)}.

\begin{lemma}[The field and the generating identity]\label{Y-P2-L1} Let $Q\ge1$, and let $g_t$ be real bounded functions with pairwise disjoint supports. Then $Z_K(z):=\int\exp\{-\beta\sum_pa_p+\sum_tz_t\mathcal O_K(g_t)\}dU$ is entire, and $Z_K(0)>0$. For any numbers $m_t$, and with $m\cdot z=\sum_tm_tz_t$, the coefficient of $z_1\cdots z_Q$ in $e^{-m\cdot z}Z_K(z)/Z_K(0)$ is $\langle\prod_t(\mathcal O_Kg_t-m_t)\rangle$. Finally, the weight is that of \Cref{Y-H4-U4} at $\nu(z)_p=\mathsf b^{-1}\sum_tz_tg_t(x(p))$, and $|\nu(z)|_\infty\le\max_t|z_t|\,\|g_t\|_\infty/\mathsf b$.\end{lemma}

\begin{proof} The argument of Step~0 of the proof of \Cref{Y-H9-T} applies word for word, with the one change that $Q$ replaces $n$: the integrand is continuous on a compact group, and it is entire in $z$. Next, $e^{-m\cdot z}Z_K(z)/Z_K(0)=\langle\exp\sum_tz_t(\mathcal O_Kg_t-m_t)\rangle$, and the coefficient of $z_1\cdots z_Q$ in it is the mean of the product $\prod_t(\mathcal O_Kg_t-m_t)$. Since $\beta=L^K\mathsf b$, we have $\beta\nu_pa_p=L^K\sum_tz_tg_t(x(p))a_p$. At each $p$ at most one of the numbers $g_t(x(p))$ is $\ne0$, by the part ``Sources'' of item~3 of \S\ref{Y-H4-1}. Therefore the domain $\mathrm{Pd}_0$ of \Cref{Y-H4-U4}\,(h) depends on $\max_t|z_t|$ alone, and not on $Q$.\end{proof}

\begin{remark}[The field] The factor $\exp\{\sum_tz_t\mathcal O_K(g_t)\}$ turns Wilson's weight into the weight at the local label $\mathsf b-\sum_tz_tg_t(x)$. Hence $S_n$ is the $n$-th derivative of $\log Z$ with respect to minus the local label, taken at different points. These are the truncated functions of the action density. They concern one real gauge invariant scalar, which is given modulo a constant, and they are taken for $n\ge2$, off the diagonals.\end{remark}

\begin{lemma}[Covariance of the drawn step under integer mirrors and coarse translations]\label{Y-P2-L2} On the fine lattice $\mathbb Z^3/N'L\mathbb Z^3$, consider the drawn step that is defined in ``The step'', item~2 of \S\ref{Y-J-1}. Let $R$ be either the map $x_\mu\mapsto Lc-x_\mu$, in one direction $\mu$ and with $c\in\mathbb Z$, or a translation by a vector of $L\mathbb Z^3$, and let $(RU)(\text{bond})=U(R^{-1}\text{bond})$, where $U(\text{reversed bond})=U(\text{bond})^{-1}$. Then $\mathsf M(RU)(c)=\mathsf M(U)(R^{-1}c)$ for every coarse bond $c$, and \[\prod_ck_s((RV')(c),\mathsf M(RU)(c))=\prod_ck_s(V'(c),\mathsf M(U)(c)).\] Consequently, consider, in the units of the last lattice, a reflection $R:y_\mu\mapsto c-y_\mu$ with $c\in L^{-s}\mathbb Z$, or a translation $R$ by a vector of $L^{-s}\mathbb Z^3$. For every $k\le K-s$, the kernel density $\mathbb K_k(U;V)$, which is the density of $V_k$ conditional on $U$, satisfies $\mathbb K_k(RU;RV)=\mathbb K_k(U;V)$.\end{lemma}

\begin{proof} For translations the claim is evident. If $R$ is a reflection, then $R$ is an involution.

\emph{Step (1): geometry.} The map $R$ carries $L\mathbb Z^3$ onto itself, and it carries the centred block $B(y)=y+\{-(L-1)/2,\dots,(L-1)/2\}^3$ onto $B(Ry)$. Hence $y(Rx)=R\,y(x)$.

The map $R$ carries the coordinate path $\Gamma(y,x)$, which runs first along $e_1$, then along $e_2$, then along $e_3$, each time the short way, onto the coordinate path $\Gamma(Ry,Rx)$: only the steps in the direction $\mu$ change their sign, and the order of the axes is kept. For a path $\pi$ we have $(RU)(\pi)=U(R\pi)$; applied to the path $\Gamma(y(x),x)$, together with $y(Rx)=Ry(x)$, this gives $\mathsf t_{RU}(x)=\mathsf t_U(Rx)$.

\emph{Step (2): the bond $c=(y,y+Le_\nu)$ with $\nu\ne\mu$.} For the straight path $[x,x']$ along $e_\nu$ we have $R[x,x']=[Rx,Rx+Le_\nu]$, and therefore $h_{c,x}(RU)=h_{Rc,Rx}(U)$. The map $x\mapsto Rx$ is a bijection of $B(y)$ onto $B(Ry)$, and it sends $y$ to $Ry$. Hence $\bar\tau_c(RU)=\bar\tau_{Rc}(U)$ and $\mathsf M(RU)(c)=\mathsf M(U)(Rc)$.

\emph{Step (3): the bond $c=(y,y')$ with $y'=y+Le_\mu$.} Now $Rc=(Ry,Ry-Le_\mu)$ is the reverse of the bond $\tilde c:=(Ry',Ry)$. For $x\in B(y)$ we put $\tilde x:=Rx'=Rx-Le_\mu$; then $\tilde x\in B(Ry')$ and $\tilde x+Le_\mu=Rx$. As $[Rx,Rx']$ is the path $[\tilde x,Rx]$ run backwards, \[\FitDisplay{\dimexpr\linewidth-1pt\relax}{h_{c,x}(RU)=\mathsf t_U(Rx)U([Rx,Rx'])\mathsf t_U(Rx')^{-1}=[\mathsf t_U(\tilde x)U([\tilde x,Rx])\mathsf t_U(Rx)^{-1}]^{-1}=g_x^{-1},\qquad g_x:=h_{\tilde c,\tilde x}(U).}\]

Next, $g_x^{-1}g_y=g_y^{-1}(g_xg_y^{-1})^{-1}g_y$. The logarithm satisfies $\mathsf L(ugu^{-1})=\operatorname{Ad}_u\mathsf L(g)$ and $\mathsf L(g^{-1})=-\mathsf L(g)$; at $g=-1$ both sides of each identity are $0$. Hence $\mathsf L(g_x^{-1}g_y)=-\operatorname{Ad}_{g_y^{-1}}\mathsf L(g_xg_y^{-1})$. The map $x\mapsto\tilde x$ is a bijection of $B(y)$ onto $B(Ry')$ which sends $y$ to the base site $Ry'$ of $\tilde c$; summing over $x\in B(y)$ we obtain $\bar\tau_c(RU)=-\operatorname{Ad}_{g_y^{-1}}\bar\tau_{\tilde c}(U)$. Finally, \[\mathsf M(RU)(c)=\exp\{i\bar\tau_c(RU)\}g_y^{-1}=g_y^{-1}e^{-i\bar\tau_{\tilde c}(U)}=(e^{i\bar\tau_{\tilde c}(U)}g_y)^{-1}=\mathsf M(U)(\tilde c)^{-1}=\mathsf M(U)(Rc);\] the second equality holds because the exponential turns $\operatorname{Ad}_{g_y^{-1}}$ into conjugation by $g_y^{-1}$, and the last because $Rc$ is $\tilde c$ reversed.

\emph{Step (4) Kernel: both entries inverted.} We have $(g^{-1})(h^{-1})^{-1}=g^{-1}h$, and this is conjugate to $hg^{-1}=(gh^{-1})^{-1}$. Since $\operatorname{ang}$ is a class function and $\operatorname{ang}(u^{-1})=\operatorname{ang}(u)$, it follows that $k_s(g^{-1},h^{-1})=k_s(g,h)$. The map $R$ permutes the unoriented coarse bonds; at a bond whose orientation it reverses, both entries of the kernel are inverted, and the factor of that bond keeps its value.

\emph{Step (5) The last assertion: the levels down to depth $s$.} In the step from $\mathbb T_k$ to $\mathbb T_{k+1}$ the coarse site $y$ is located at the fine site $Ly$; so all levels are the lattices $L^{-j}\mathbb Z^3$ in one torus $\mathbb R^3/N\mathbb Z^3$, and $0$ is a site of each. In the units of depth $j$ the reflection $y_\mu\mapsto c-y_\mu$ of the last assertion of \Cref{Y-P2-L2} becomes the map $x\mapsto L^jc-x$, where $L^jc=L(L^{j-1}c)$ and $L^{j-1}c\in\mathbb Z$ for $j-1\ge s$. In the same units a translation by a vector $v\in L^{-s}\mathbb Z^3$ becomes the translation by $L^jv=L(L^{j-1}v)$, where $L^{j-1}v\in\mathbb Z^3$ for $j-1\ge s$; so it is a translation by a vector of $L\mathbb Z^3$, which is one of the maps $R$ of the first assertion of the lemma. Hence \textup{(1)} to \textup{(4)} hold at every step down to depth $s$.

The step depends on $x$, $a$ and $L$ alone: it depends on no rows, on no grid of cubes, on no members, and not on $K$; this is the last sentence of item~2 (``The step'') of \S\ref{Y-J-1}. The density $\mathbb K_k$ is the integral over $V_1,\dots,V_{k-1}$ of the product of the kernels of the steps. In this integral we change the variables of integration, $V_j\to RV_j$; the Haar measure is invariant under permutations of the bonds and under inversion, and \textup{(1)} to \textup{(4)} hold at each of these steps; so $\mathbb K_k(RU;RV)=\mathbb K_k(U;V)$ for $k\le K-s$. \end{proof}

\begin{lemma}[The subtraction as a mirror invariant weighted mean]\label{Y-P2-L3} Let $0\le s\le K-1$ and $k=K-s$.

\begin{enumerate}\item[\textup{(a)}] We have \[c^{(s)}_K(f)=\int \mathcal O_K(f)e^{-\beta A}\mathbb K_k(U;1)\,dU\big/\int e^{-\beta A}\mathbb K_k(U;1)\,dU;\] this number is real for real $f$, and it is linear in $f$. Here $A=\sum_pa_p$, with $a_p=1-\tfrac12\operatorname{tr}\partial U(p)$ (item~1 of \S\ref{Y-P2-1}).

\item[\textup{(b)}] For the maps $R$ of the last assertion of \Cref{Y-P2-L2} (in the units of the last lattice, the reflections $y_\mu\mapsto c-y_\mu$ with $c\in L^{-s}\mathbb Z$ and the translations by vectors of $L^{-s}\mathbb Z^3$), $c^{(s)}_K(f\circ R)=c^{(s)}_K(f)$.

\item[\textup{(c)}] Let the $g_t$ be the $Q\ge1$ sources of the data of \textup{(R1)} of \Cref{Y-P2-ThR}: real bounded functions whose supports are pairwise at distance $\ge\check\delta$, each in a ball of radius $\check\rho$. Then $c^{(s)}_K(g_t)=\gamma_t^{(k)}+e_t^{(s)}$, where $\gamma_t^{(k)}:=\partial_t\sum_{j\le k}V_j(0)$ and $e_t^{(s)}:=\sum_X\delta_t^h(k)(X,1)$. If $\|g_t\|_\infty\le1$, then $r^\circ_s|e_t^{(s)}|\le e_h^{(s)}$.\end{enumerate}\end{lemma}

At $s=0$, clause \textup{(c)} is \Cref{Y-P1-Jo}\,\textup{(J7$^\circ$)\,(vii)} together with \textup{(J6$^\circ$)\,(e)}, where the constant $0.065$ appears in the place of $0.45$.

\begin{proof} \emph{Step (a): the weighted mean.} By the definition of the run, $\rho_k^{[zf]}(1)=\int e^{-\beta A+z\mathcal O_K(f)}\mathbb K_k(U;\allowbreak 1)\,\allowbreak dU$. By Step~11 \textup{(0)} of the proof of \Cref{Y-H9-T}, the integrand is Borel, the integral is entire in $z$, and $d/dz$ may be taken under the integrals. As $\mathbb K_k>0$, we have $\rho_k(1)>0$. Differentiating the logarithm at $z=0$ we obtain the formula of \textup{(a)}; it shows that $c^{(s)}_K(f)$ is linear in $f$, and real for real $f$.

\emph{Step (b): mirror invariance.} In both integrals we change the variable of integration, $U\to RU$. The Haar measure and $A$ are invariant under this change. By the last assertion of \Cref{Y-P2-L2}, $\mathbb K_k(RU;RV)=\mathbb K_k(U;V)$ for $k\le K-s$, and so $\mathbb K_k(RU;1)=\mathbb K_k(RU;R1)=\mathbb K_k(U;1)$. By Step~7 (``At finite $K$'') of the proof of \Cref{Y-H9-T}, $\mathcal O_K(f)(RU)=\mathcal O_K(f\circ R)(U)$. So the numerator for $f$ equals the numerator for $f\circ R$, the denominator stays as it is, and \textup{(b)} follows.

\emph{Step (c) The formula: the value at the trivial field.} By \Cref{Y-H4-U4}\,\textup{(h)}, on $\mathrm{Pd}_0$ we have \[\rho_k^z(V)=\exp\{\sum_{j\le k}V_j(z)\}\,\rho[u_k^{\mathrm{hon}}(z);\ell_k](V)\] at \emph{every} field $V$ of $\mathbb T_k$. Here $u_k^{\mathrm{hon}}(z)$ is the honest family, the family $u_k(z)$ of \Cref{Y-H4-U4}\,\textup{(h)}. We write $u_k(z)$, without the superscript $\mathrm{hon}$, for the representative of \textup{(P1.1)}. By \textup{(P1.1)} the jets at $0$ of both families are the $\delta_S(k)$. Thus $E^h_k(X,1;z)$ below belongs to $u_k^{\mathrm{hon}}(z)$, while $E^h_k(X,\cdot\,;z)$ in the next step belongs to the representative. We evaluate the right side at the trivial field $V=1$, in four parts.

First, $\phi_p=0$, so that $S_\varepsilon(0)=0$, $\zeta_p=0$ and $\mathrm{En}_\ell(1)=0$ (item~2, ``Energy'', of \S\ref{Y-J-1}). Second, $D(1)$ is empty; as every $X_i$ of a family must hold a plaquette of $D(V)$, while the empty family is admitted, we get $\Xi_w(1)=1$ (item~2 of \S\ref{Y-J-1}, the definition of $\rho[u;\ell]$ and of $\Xi_w$). Third, the field $1$ is tame, so $1_{X^\approx}(1)=1$. Fourth, $E^a(X,1)=0$ by the clause ``$0$ at pure gauges'' of \textup{(ca)} in item~2 of \S\ref{Y-J-1}; this clause is not one of the size bounds, and the outputs of \Cref{Y-H-C9} obey every clause of \textup{(ca)}, \textup{(ch)} and \textup{(c$\Xi$)} other than the size bounds (\textup{(9b)} of item~3 of \S\ref{Y-J-1}; here $k\ge1$), as well as the size bounds of \Cref{Y-H-C9}\,\textup{(b)}.

Hence \[\log\rho_k^z(1)=\sum_{j\le k}V_j(z)+\sum_XE^h_k(X,1;z),\] a finite sum of functions holomorphic near $0$. The derivative $\partial_t$ at $0$, with the other variables set to $0$, is the derivative in the run with the single source $g_t$; this gives the formula of \textup{(c)}.

\emph{Step (c) The size: the bound on $e_t^{(s)}$.} By \textup{(P1.1)} in item~0.4 of \S\ref{Y-P2-1}, $\delta_t^h(k)$ is $\partial_tE^h$ of the representative. First let $k>k^\circ$. Then $s<j^\circ$, and $r^\circ_s=r^\circ\le\mathfrak r^\circ_*$. On the disc $|z_t|<r^\circ_s$, with the other variables $0$, the part of \textup{(P1.1)} for $k^\circ<k\le K$, taken in the norm $|\cdot|_h^{(\kappa)}$ of \textup{(o6)} (item~2 of \S\ref{Y-H6-1}), gives \[|E^h_k(X,\cdot\,;z)-\bar E^h_k(X,\cdot)|\le0.91\cdot0.4\,c_h(x_k)e^{-\kappa(\#X-1)},\] and Cauchy's estimate gives $|\delta_t^h(k)(X,1)|\le0.364\,c_h(L^s\mathsf b)e^{-\kappa(\#X-1)}/r^\circ_s$.

Now let $k\le k^\circ$. Then $s\ge j^\circ$, and $\mathfrak r^\circ_k=\mathfrak r^\circ(L^s\mathsf b)=r^\circ_s$. On the same disc the part of \textup{(P1.1)} for $k\le k^\circ$, taken in the norm $|\cdot|_h^{(\kappa^\circ_k)}$ of \textup{(o6)}, gives the bound of the previous case with $0.91$ in the place of $0.91\cdot0.4$ and with the rate $\kappa^\circ_k$ in the place of $\kappa$; since $\kappa^\circ_k\ge\kappa$, Cauchy's estimate now gives $|\delta_t^h(k)(X,1)|\le0.91\,c_h(L^s\mathsf b)e^{-\kappa(\#X-1)}/r^\circ_s$.

In both cases only the $X$ that hold a cube of $S_t(k)^{+\check r(x_k)}$ count: for the others $\delta_t(k)(X)=0$, by $(P_k)$ (Step~3 of the proof of \Cref{Y-H6-J}; \textup{(P1.1)}). The number of these cubes is at most $N_1(k-1)=N_1^{[s+1]}$ (Step~1 \textup{(A)} of the proof of \Cref{Y-H6-J}). The connected $X$ of $m$ cubes that hold a given cube are at most $(26e)^{m-1}$ in number (item~2, ``Lattice, cubes'', of \S\ref{Y-J-1}). As $\kappa\ge6$, the geometric series gives $\sum_m(26e\,e^{-6})^{m-1}\le1.2124$. So $r^\circ_s|e_t^{(s)}|$ is at most $N_1^{[s+1]}c_h(L^s\mathsf b)$ times $0.364\cdot1.2124$ if $k>k^\circ$, and times $0.91\cdot1.2124$ if $k\le k^\circ$. Now $0.364\cdot1.2124\le0.45$; this settles the case $s=0$, which falls under $k>k^\circ$ because $j^\circ\ge1$. For $s\ge1$ the numerical factor is at most $1.11$ in either case: if $k>k^\circ$, that is, if $1\le s<j^\circ$, it is at most $0.45$, and if $k\le k^\circ$, it is $0.91\cdot1.2124\le1.11$. As $e_h^{(0)}=0.45\,N_1^{[1]}c_h(\mathsf b)$ and $e_h^{(s)}=1.11\,N_1^{[s+1]}c_h(L^s\mathsf b)$ for $s\ge1$ (item~2 of \S\ref{Y-P2-1}), we obtain $r^\circ_s|e_t^{(s)}|\le e_h^{(s)}$ for every $0\le s\le K-1$, which is the bound of \textup{(c)}. \end{proof}

\begin{remark}[The constant of the subtraction] The constant of the subtraction is $c^{(s)}_K$ rather than the vacuum numbers $\gamma_t$ alone, because the grids of cubes and the base sites are not mirror covariant.\end{remark}

\begin{lemma}[The last integral bounded as a moment]\label{Y-P2-L4} Assume the data of item~0.4 of \S\ref{Y-P2-1}, let $u_k(z)$ be the representative of \textup{(P1.1)} for $k=K-s$ with $0\le s\le K-1$, let $z\in\mathrm{Pd}(r^\circ_s)$, and put $I_1:=I[\bar u_k]>0$. Then $\int|\rho[u_k(z);\ell_k]|\,dV\le\mathfrak V_sI_1$, where \[\mathfrak V_0=1+0.182\,C_{\mathrm{vol}}(N,\mathsf b)\le e^{c_{\mathrm{vol}}(\mathsf b)N^3}\] and, for $s\ge1$, \[\mathfrak V_s=1+0.455\,C_{\mathrm{vol}}(NL^s,L^s\mathsf b)e^{200q_N}\le e^{c_{\mathrm{vol}}^{(s)}(\mathsf b)N^3},\qquad q_N=(NL^s/M)^3.\] Moreover, for every bounded Borel function $F$ on the fields of $\mathbb T_k$, the map $z\mapsto\int F\rho[u_k(z);\ell_k]\,dV$ is holomorphic on $\mathrm{Pd}(r^\circ_s)$.\end{lemma}

\begin{proof} \emph{Step (Data \textup{(B)}): at $x=L^s\mathsf b$ the data of \Cref{Y-H8-ThBC} are met.} We take the data \textup{(B)} of \Cref{Y-H8-ThBC} at $x=L^s\mathsf b$ on $\mathbb Z^3/NL^s\mathbb Z^3$, with the rate $\tilde\kappa=\kappa$ if $k>k^\circ$ and $\tilde\kappa=\kappa^\circ_k$ if $k\le k^\circ$; both rates lie in $[\kappa,\kappa_1]$. These data are met. First, $x\ge x_s\ge e^2$. Next, $NL^s=2L^{m_l+s}$ is even, and $M$ divides it. Further, $NL^s\ge2Ms(L^s\mathsf b)$: for $s=0$ by \textup{(Tor)}, since $s(\cdot)$ is non-decreasing, and for $s\ge1$ by Step~11 \textup{(5)} of the proof of \Cref{Y-H9-T}.

Moreover, $\varepsilon\le1$ and $c_a+c_h\le1/2$ at $x\ge x_s$, by \textup{(X16)} and \textup{(X17)} (entry \textup{(O13)} of item~1 of \S\ref{Y-J-1}). The rows $\ell_k$ have range $\le\varrho(x_k)$ and obey the decay clause, for $0\le k\le K$ and every $K\ge1$ (this is Step~2, ``The last integral'', of the proof of \Cref{Y-H9-T}). As the first member we take $u_1=\bar u_k$; by \Cref{Y-H6-J}\,\textup{(J6)} and \Cref{Y-H4-U4}\,\textup{(h)}, it is a system of members of a point of $C^{\rm su}_\flat(x_k;\ell_k)$. We put $u_3:=u_1$ and $u_2:=u_k(z)$.

\emph{Step (Case $k>k^\circ$): case \textup{(ii)} of \textup{(b$^C$)}.} Here $u_2=F_{\ell_{k-1}}(u_{k-1}(z))$, with its input in $M_{1.91}\subset M_2$. So $u_2$ is an output of \Cref{Y-H-C9} applied with $A\le6$. For this output, and with $u_3\in M_1$, we may take $a_w=1$ and $a_E\le2$ in \Cref{Y-H8-ThBC}\,\textup{(b$^C$)}: this is case \textup{(ii)} of that clause, used with the pair $(\kappa,c_h)$. Further, $\mathsf D_{23}=\|u_2-u_1\|^\sim\le\|u_2-u_1\|^{\mathrm{st},(\kappa)}\le0.364$, since the strong norm of the same difference dominates (\Cref{Y-H8-Norms}). The $L^1$ bound of \Cref{Y-H8-ThBC}\,\textup{(c$^C$)} for the cases \textup{(i)}, \textup{(ii)} and \textup{(iii)} then gives in the case \textup{(ii)} \[\int|\rho[u_2]-\rho[u_1]|\,dV\le\frac12C_{\mathrm{vol}}\cdot0.364\,I_1.\]

\emph{Step (Case $k\le k^\circ$): case \textup{(iii)} of \textup{(b$^C$)}.} Here $\|u_2\|^\sim\le1.91$ and $\|u_3\|^\sim\le1$: this is case \textup{(iii)} of \textup{(b$^C$)}, applied with $A=1.91$ and $\mathsf D_{23}\le0.91$. The same integral is now at most $\frac12C_{\mathrm{vol}}e^{219\cdot0.91q_N}\,0.91\,I_1$, and $219\cdot0.91=199.29\le200$.

\emph{Step (Add): the first member.} We add $\int\rho[u_1]\,dV=I_1$ to the bounds of the two cases; this gives the first claim.

\emph{Step (Arithmetic): the two exponentials.} Let $X$ be the part of $C_{\mathrm{vol}}$ that is a power with exponent $N^3$. For $\mathfrak V_0$ we use $1+0.728\,q_Ne^{220q_N}X\le e^{220.728q_N}X$, and for $\mathfrak V_s$ we use $1+1.82\,q_Ne^{420q_N}X\le e^{421.82q_N}X\le e^{425q_N}X$; here $q_N=L^{3s}N^3/M^3$.

\emph{Step (Holomorphy): the integrals against $F$.} For each field $V$, $\rho[u;\ell](V)$ is an entire function of finitely many values of $u$ (the clause of \Cref{Y-H5-LJ} on $I[u]$); hence $\rho[u_k(z);\ell_k](V)$ is holomorphic in $z$. By \textup{(b$^C$)} of \Cref{Y-H8-ThBC}, $\sup_V|\rho[u_2]-\rho[u_1]|$ is bounded on the polydisc, and $\rho[u_1]$ is bounded. The function $F(V)\rho[u_k(z);\ell_k](V)$ is holomorphic in $z$, Borel in $V$ and bounded on $\mathrm{Pd}(r^\circ_s)$; so Step~P4\,\textup{(2)} of the proof of \Cref{Y-H8-CoF} applies to it, and it gives both the second claim and the identity $\partial_S\int F\rho[u_k(z);\ell_k]\,dV=\int F\,\partial_S\rho[u_k(z);\ell_k]\,dV$. \end{proof}

\begin{remark}[On the last integral] The lower bound on $I_1$ is crude, and this crudeness costs only the size of $c_{\mathrm{vol}}$.\end{remark}

\begin{lemma}[The master bound for centred sources]\label{Y-P2-L5} Assume the data of \textup{(R1)} of \Cref{Y-P2-ThR}, let $0\le s\le K-1$ and $k=K-s$, and let $F$ be a bounded Borel function of the fields of $\mathbb T_k$. Then \[\bigl|\bigl\langle F(V_k)\prod_t(\mathcal O_Kg_t-c^{(s)}_Kg_t)\bigr\rangle\bigr|\le\sup|F|\;G_{1,s}^Q\,\mathfrak V_s\prod\|g_t\|_\infty.\]\end{lemma}

Clause \textup{(R1)} of \Cref{Y-P2-ThR} is the case $s=0$, $F=1$. At $s=0$ the lemma is \Cref{Y-P1-Cor}\,\textup{(a)}, \textup{(b)} of \S\ref{Y-P1-1}.

\begin{proof} By homogeneity we may take $\|g_t\|_\infty=1$ for every $t$.

\emph{Step (1) The ideal: congruence of germs.} The germs at $0$ in $\mathbb C^Q$ whose squarefree coefficients all vanish form the ideal $(z_t^2)$: every monomial of such a germ is divisible by some $z_t^2$, and we group the monomials by the least such $t$. Hence $f\equiv g$ implies $fh\equiv gh$. It also implies $e^f\equiv e^g$, since $e^f-e^g=e^g(e^{f-g}-1)$.

\emph{Step (2) The exponent: its higher part.} By \textup{(P1.2)}, in the form \textup{(J7$^\circ$)\,(iv)} of \Cref{Y-P1-Jo}, the exponent at the level $k$ is $\sum_{j\le k}V_j(0)+\sum_tz_t\gamma_t^{(k)}+W^{(k)}(z)$, with \[W^{(k)}:=\sum_{k'<\min(k,k^\circ)}W_{k'}+\sum_{k^\circ\le k'<k}[w_{k'}-\text{linear part}].\] Bounding term by term with \textup{(P1.3)} on $\mathrm{Pd}(\mathfrak r^\circ_*)\supset\mathrm{Pd}(r^\circ_s)$, we get $|W^{(k)}|\le QC_W$. For a late term, of depth $j=K-k'$, we have $|w_{k'}|\le3QN_1^{[j]}C_gc_a(L^j\mathsf b)$ and $|\text{linear part}|\le Q\,\mathfrak r^\circ_*\cdot3N_1^{[j]}C_gc_a(L^j\mathsf b)/\mathfrak r^\circ_*$, which together give $6QN_1^{[j]}C_gc_a(L^j\mathsf b)$.

\emph{Step (3) The identity: the generating function.} By Step~11 \textup{(0)} of the proof of \Cref{Y-H9-T}, $\langle F(V_k)e^{\sum z_t\mathcal O_Kg_t}\rangle=\int F\rho_k^z\,dV/Z_K(0)$. We take \textup{(P1.2)} at the level $k$, pointwise in $V$; the derivatives $\partial_S$ may be taken under $\int F\,\cdot\,dV$ on both sides, on the left by Step~11 \textup{(0)} and on the right by Step (Holomorphy) of the proof of \Cref{Y-P2-L4}. Moreover $Z_K(0)=\exp\{\sum_{j\le k}V_j(0)\}I[\bar u_k]$ (by \Cref{Y-H4-U4}\,\textup{(h)} and Step~0 of the proof of \Cref{Y-H9-T}). With Steps~(1) and~(2) and \Cref{Y-P2-L3}\,\textup{(c)} we arrive at \[e^{-\sum z_tc^{(s)}_Kg_t}\langle F(V_k)e^{\sum z_t\mathcal O_Kg_t}\rangle\equiv G_F(z):=\exp\{W^{(k)}(z)-\sum_tz_te_t^{(s)}\}\int F\rho[u_k(z);\ell_k]\,dV\big/I[\bar u_k].\]

\emph{Step (4) Cauchy: the coefficient of $z_1\cdots z_Q$.} The function $G_F$ is holomorphic on $\mathrm{Pd}(r^\circ_s)$, and there $|G_F|\le e^{QC_W}e^{Qe_h^{(s)}}\sup|F|\,\mathfrak V_s$, by Lems.~\ref{Y-P2-L3} and~\ref{Y-P2-L4}. The coefficient of $z_1\cdots z_Q$ on the left is the moment, as in \Cref{Y-P2-L1}; by Step~(1) it is also the coefficient of $G_F$, and Cauchy's estimate bounds it by $(r^\circ_s)^{-Q}\sup|G_F|$. As $G_{1,s}=e^{C_W+e_h^{(s)}}/r^\circ_s$, this is the claim. \end{proof}

\begin{remark}[On the master bound] In \Cref{Y-P2-L5} the volume factor enters the bound with the exponent \emph{one}, whatever $Q$ is. All the ultraviolet input of the lemma lies in \textup{(P1.1)} to \textup{(P1.3)}, in \Cref{Y-H-C9}, in \Cref{Y-H1-U1}\,\textup{(g2)} and in \Cref{Y-H8-ThBC}\,\textup{(b$^C$)}.\end{remark}

\begin{lemma}[The joint law: locality, conditional independence, covariance]\label{Y-P2-L6} Consider the joint law \[d\mu(U)\prod_{k<K}\prod_ck_{ax_k}(V_{k+1}(c),\mathsf M(V_k)(c))\,dV_{k+1};\] it is the law of Step~11 \textup{(0)} of the proof of \Cref{Y-H9-T}. Let $0\le s\le K-1$. The following assertions hold for this law.

\begin{enumerate}\item[\textup{A}] For a coarse bond $c=(y,y')$, $\mathsf M(U)(c)$ is a function of $U$ on the bonds with both ends in $B(y)\cup B(y')$.

\item[\textup{B}] Let $O$ be an open set bounded by planes $y_\mu\in L^{-s}(\mathbb Z+\frac12)$, and let $F$ be a bounded Borel function of the values of $V_{K-s}$ at finitely many bonds, each with both ends in $O$. Then the function $E[F(V_{K-s})|U]$ of $U$ is bounded and Borel, and it depends only on the bonds with both ends in $O$.

\item[\textup{C}] Let the sets $O_i$ be pairwise disjoint, and let each $F_i$ be as in B for $O_i$. Then $E[\prod F_i|U]=\prod E[F_i|U]$.

\item[\textup{D}] For a map $R$ of the last assertion of \Cref{Y-P2-L2} and a bounded Borel function $F$ of $V_{K-s}$ we put $(RF)(V)=F(R^{-1}V)$, $\Phi=E[F|U]$ and $(R\Phi)(U):=\Phi(R^{-1}U)$. Then $E[(RF)(V_{K-s})|U]=(R\Phi)(U)$.\end{enumerate}\end{lemma}

\begin{proof} \emph{Step A. Reach of $\mathsf M$: the bonds on which the average depends.} Let $c=(y,y')$ be as in A. The path $\Gamma(y,x)$ leads from the centre of the box $B(y)$ to $x\in B(y)$ and stays in this box; the same is true of $\Gamma(y',x')$; and the straight path $[x,x+Le_\mu]$ stays in the two boxes. These three facts give claim A.

\emph{Step B. Locality: functions of an open set.} Let $O$ be as in B. As $L$ is odd, $L^{-s}(\mathbb Z+\frac12)\subset L^{-j}(\mathbb Z+\frac12)$ for all $j\ge s$. The block of a site $y$ of depth $j-1\ge s$ consists of the sites of depth $j$ in $y+L^{-(j-1)}(-\frac12,\frac12)^3$; so at every step down to depth $s$ the planes are faces of blocks, and no site lies on them. By A, the average at a bond of depth $j-1$ with both ends in $O$ depends only on bonds of depth $j$ with both ends in $O$.

Let $F$ be a bounded Borel function of $V_{k+1}$ on a finite set $B_F$ of such bonds. Then \[E[F(V_{k+1})|V_k]=\int F(V')\prod_{c\in B_F}k_{ax_k}(V'(c),\mathsf M(V_k)(c))\,dV',\] because the kernels at the other bonds integrate to $1$. This function is bounded; it is Borel, as $k_s$ is continuous and $\mathsf M$ is Borel; and it is a function of $V_k$ on the bonds with both ends in $O$. Iterating along the Markov chain we obtain claim B.

\emph{Step C. Conditional independence: disjoint open sets.} Let $O_i$ and $F_i$ be as in C. Then the sets $B_{F_i}$ are disjoint, and since the kernel is a \emph{product} over the coarse bonds, $E[\prod F_i(V_{k+1})|V_k]=\prod_i\Phi_i(V_k)$, where $\Phi_i(V_k):=E[F_i(V_{k+1})|V_k]$ is the function of Step B for $F_i$. Induction on the level gives claim C.

\emph{Step D. Covariance: the maps $R$.} In the integral against $\mathbb K$ we change the variable of integration, $V\to RV$; as the last assertion of \Cref{Y-P2-L2} gives $\mathbb K_k(RU;RV)=\mathbb K_k(U;V)$ for $k\le K-s$, we obtain claim D. \end{proof}

\begin{lemma}[Boards, functions of a cell, mirror maps, folds]\label{Y-P2-L7} Let $0\le s\le K-1$, $-m_l\le i\le\min(s,K-1)$, $o\in L^{-s}\mathbb Z^3$, $h=L^{-i}$ and $\mathsf n_\mu=N/h$. The cells of $B_i(o)$ are indexed by $\varkappa,\varkappa'\in\prod_\mu\mathbb Z/\mathsf n_\mu\mathbb Z$, and the cell $\varkappa$ is centred at $o+h\varkappa$.

\begin{enumerate}\item[\textup{(a)}] $B_i(o)$ is a board in the sense of \Cref{support:4-1}\,\textup{(2)}, with $N_c=(N/h)^3$ cells, and its walls lie in $L^{-s}(\mathbb Z+\frac12)\subset a_K(\mathbb Z+\frac12)$.

\item[\textup{(b)}] If $\{g\ne0\}$ lies in an open cell, then $\mathcal O_K(g)$ is a function of that cell.

\item[\textup{(c)}] Every mirror map $r_{\varkappa'\varkappa}$ is, coordinate by coordinate, a map of \textup{(R0)} of \Cref{Y-P2-ThR} at depth $s$.

\item[\textup{(d)}] Let $Y=\prod_{t\in I}(\mathcal O_Kg_t-c^{(s)}_Kg_t)$, with supports in the open cell $\varkappa$, and let $\Phi=E[F(V_{K-s})|U]$, where $F$ depends only on the bonds that have both ends in that open cell. Then \[\langle\operatorname{Fold}(\Phi Y)\rangle=\bigl\langle\prod_{\varkappa'}F^{[\varkappa']}(V_{K-s})\prod_{\varkappa'}\prod_{t\in I}(\mathcal O_K(rg_t)-c^{(s)}_K(rg_t))\bigr\rangle,\] where $r=r_{\varkappa'\varkappa}$.\end{enumerate}\end{lemma}

\begin{proof} \emph{Step (a): the board.} For the board $B_i(o)$, $\lambda=L^{K-i}$ is odd, and $\lambda\ge3$ because $i\le K-1$; $\mathsf n_\mu=N/h=2L^{m_l+i}$ is an even integer $\ge2$; and $o\in a_K\mathbb Z^3$ because $s\le K$. For the walls, $L^{-i}(m+\frac12)=L^{-s}L^{s-i}(2m+1)/2$, where $L^{s-i}(2m+1)$ is odd; adding $o_\mu\in L^{-s}\mathbb Z$ keeps the wall in $L^{-s}(\mathbb Z+\frac12)$.

\emph{Step (b): functions of a cell.} A plaquette $p$ enters $\mathcal O_K(g)$ only if its barycentre $x(p)$ lies in $\{g\ne0\}$, hence in the open cell. In the two directions of $p$, the coordinate of $x(p)$ lies in $a_K(\mathbb Z+\frac12)$, as the walls do; so it is at distance $\ge a_K$ from them, and the corners, at $\pm a_K/2$, lie inside. In the third direction the coordinate is that of a site, which lies off the walls. Hence the four corners of $p$ lie in the cell, and its four bonds are inner bonds; so $\mathcal O_K(g)$ is a function of the cell. Here The inner bonds are those of \Cref{support:4-1}. Cell functions are described in \eqref{support:2-8} and \Cref{support:2-9}(d).

\emph{Step (c): mirror maps.} By \eqref{support:2-18}, $r_{\varkappa'\varkappa}$ is the map $x\mapsto o+h\varkappa'+\sigma(x-o-h\varkappa)$, where $\sigma=\sigma(\varkappa',\varkappa)=(\sigma_\mu)_\mu$, with $\sigma_\mu=(-1)^{\varkappa'_\mu-\varkappa_\mu}$, acts coordinate by coordinate: $\sigma u=(\sigma_\mu u_\mu)_\mu$. In a coordinate $\mu$ with $\varkappa'_\mu-\varkappa_\mu$ even, it is the translation by $h(\varkappa'_\mu-\varkappa_\mu)\in L^{-s}\mathbb Z$. In a coordinate with $\varkappa'_\mu-\varkappa_\mu$ odd, it is $y\mapsto2o_\mu+h(\varkappa'_\mu+\varkappa_\mu)-y$, with the constant in $L^{-s}\mathbb Z$. Thus $r_{\varkappa'\varkappa}$ is a composition of maps in the different coordinates, and it is covered by \textup{(R0)} applied to one coordinate after the other.

\emph{Step (d): the fold.} By \Cref{Y-P2-L6}\,B, $\Phi$ is a function of the cell. By \eqref{support:2-26}, $(\Phi Y)^{[\varkappa']}$ is $r(\Phi Y)$ or its complex conjugate, and $Y$ is real. By Step~7 of the proof of \Cref{Y-H9-T}, $r\mathcal O_K(g)=\mathcal O_K(g)(r^{-1}\cdot)=\mathcal O_K(g\circ r^{-1})=\mathcal O_K(rg)$, and by \Cref{Y-P2-L3}\,\textup{(b)}, $c^{(s)}_K(g)=c^{(s)}_K(rg)$; so the copy of $Y$ is $\prod_{t\in I}(\mathcal O_K(rg_t)-c^{(s)}_K(rg_t))$. By \Cref{Y-P2-L6}\,D, $r\Phi=E[rF|U]$; complex conjugation passes through $E$, and so $\Phi^{[\varkappa']}=E[F^{[\varkappa']}(V_{K-s})|U]$. As the open cells are disjoint, \Cref{Y-P2-L6}\,C gives $\prod_{\varkappa'}E[F^{[\varkappa']}|U]=E[\prod F^{[\varkappa']}|U]$. The factors that depend on $U$ are measurable in $U$; taking the mean we obtain \textup{(d)}. \end{proof}

\begin{lemma}[Pieces, placing and mirror copies]\label{Y-P2-L8}\begin{enumerate}\item[\textup{(a)}] Let $f_1,\dots,f_n$ be of the type $(d;R_f)$, and let $0<\check\delta\le d$. Each $f_i$ is a sum of at most $(4R_f/\check\delta+2)^3$ pieces $g=f_i1_{\mathrm{cube}}$, over half-open cubes of side $\check\delta/2$. In a tuple with one piece of each $f_i$, the supports are pairwise at distance $\ge\check\delta$, and each of them lies in a ball of radius $\check\delta/4$.

\item[\textup{(b)}] Let, besides, $\check\delta\le1/(4n)$ and $a_K<5/(8n)$. Then there is $v\in a_K\mathbb Z^3$ such that every cube of such a tuple, translated by $v$, is at distance $\ge\check\delta/2$ from all the planes $y_\mu\in\mathbb Z+\frac12$, and so lies in one open cell of $B_0$. The condition on $a_K$ holds for $K\ge j^\circ(\mathsf b,\check\delta)$.

\item[\textup{(b$'$)}] Let $f_1,\dots,f_n$ be of the type $(d;R_f)$, let $h=L^{-i}>n(2R_f+d+1)$, and let $N\ge2h$. Then there is $o\in\mathbb Z^3$ such that each of the $n$ balls of radius $R_f$ that contain the supports of $f_1,\dots,f_n$ in the definition of the type is at distance $\ge d/2$ from the walls of $B_i(o)$.

\item[\textup{(c)}] On any board, let the sources of one cell be pairwise at distance $\ge\check\delta$, each in a ball of radius $\check\rho$ and at distance $\ge\check\delta/2$ from the walls. Then their mirror copies in all the cells are pairwise at distance $\ge\check\delta$, each lies in a ball of radius $\check\rho$, and they have the same sup norms.\end{enumerate}\end{lemma}

\begin{proof} \emph{Step (a): the pieces.} If $2R_f<N$, the ball lifts to a cube of side $2R_f$ in $\mathbb R^3$, which we cut by the grid $(\check\delta/2)\mathbb Z^3$; an interval of length $2R_f$ meets at most $4R_f/\check\delta+2$ of the grid intervals. If $2R_f\ge N$, we cut the fundamental cube $[0,N)^3$, the last row being smaller; per axis there are at most $2N/\check\delta+1\le4R_f/\check\delta+2$ intervals. Pieces of different $f_i$ are at distance $\ge d\ge\check\delta$, and a half-open cube of side $\check\delta/2$ lies in a closed sup ball of radius $\check\delta/4$.

\emph{Step (b): the translation.} We fix $\mu$; the $i$-th cube projects to $[\alpha_i,\alpha_i+\check\delta/2]$. The values of $v_\mu$ modulo $1$ for which $[\alpha_i+v_\mu,\alpha_i+\check\delta/2+v_\mu]$ comes closer than $\check\delta/2$ to $\mathbb Z+\frac12$ form an \emph{open} arc of length $\check\delta/2+\check\delta=1.5\,\check\delta$. Altogether these arcs have length $\le1.5\,n\check\delta\le3/8$. Their complement in $\mathbb R/\mathbb Z$ consists of at most $n$ \emph{closed} arcs of total length $\ge5/8$, so one of them has length $\ge5/(8n)>a_K$. That arc contains a point of $a_K\mathbb Z/\mathbb Z$, which is a grid of spacing $a_K$ because $1/a_K\in\mathbb Z$; we take $v_\mu$ there. An interval at positive distance from all the walls contains none of them; hence each translated cube lies in one open cell. Finally, if $K\ge j^\circ$, then $\mathrm{sep}^{[j^\circ-1]}\ge\mathrm{sep}^\circ$ forces $\check\delta L^{j^\circ-1}/M\ge\mathrm{sep}^\circ$, and so $a_K\le L^{-j^\circ}\le\check\delta/(LM\,\mathrm{sep}^\circ)<5/(8n)$.

\emph{Step \textup{(b$'$)}: the point $o$.} We fix $\mu$ and call $o_\mu\in\{0,\dots,h-1\}$ bad if some wall $o_\mu+h(\mathbb Z+\frac12)$ lies at distance less than $d/2$ from the projection of a ball, which is an interval of length $\le2R_f$. For one ball the bad $o_\mu$ are points of spacing one in an open interval of length $2R_f+d$, so there are at most $2R_f+d+1$ of them. As $n(2R_f+d+1)<h$, a good $o_\mu$ remains. The condition $N\ge2h$ is the condition $N/h=2L^{m_l+i}\ge2$.

\emph{Step (c): the mirror copies.} Inside one cell the copies are the image of the sources under one isometry, which preserves distances and sup norms and maps each ball to a ball of the same radius. If two cells are different, there is a coordinate $\mu$ in which they differ. Then, on the circle $\mathbb R/N\mathbb Z$, each arc that joins the two projections meets a wall; as each projection keeps a distance $\ge\check\delta/2$ from that wall, the distance is $\ge\check\delta$. \end{proof}

\begin{lemma}[The bound per occupied cell]\label{Y-P2-L9} In the setting of \Cref{Y-P2-L7}, let $B_i(o)$ be its board, assume the data of \textup{(R1)} of \Cref{Y-P2-ThR}, and let each source lie in one open cell, at distance $\ge\check\delta/2$ from the walls; let the factors $F_{\mathrm{cell}}$ be as in \Cref{Y-P2-L7}\,\textup{(d)}. We call a cell occupied if a source lies in it or if $F_{\mathrm{cell}}\ne1$, and we write $q$ for the number of occupied cells. Then \[\bigl|\bigl\langle\prod_{\mathrm{cells}}\bigl[F_{\mathrm{cell}}(V_{K-s})\prod_{t\in I_{\mathrm{cell}}}(\mathcal O_Kg_t-c^{(s)}_Kg_t)\bigr]\bigr\rangle\bigr|\le\prod\sup|F_{\mathrm{cell}}|\;G_{1,s}^Q\,e^{q\,c_{\mathrm{vol}}^{(s)}h^3}\prod\|g_t\|_\infty.\]\end{lemma}

Clause \textup{(R2)} of \Cref{Y-P2-ThR} is the case $s=0$, $F=1$, $o\in\mathbb Z^3$.

\begin{proof} By \Cref{Y-P2-L6}\,C and~B, the mean is $\langle\prod_{\mathrm{cells}}\Phi_{\mathrm{cell}}Y_{\mathrm{cell}}\rangle$, where $\Phi_{\mathrm{cell}}=E[F_{\mathrm{cell}}|U]$ and $Y_{\mathrm{cell}}$ is the product of the centred sources of the cell; each $\Phi_{\mathrm{cell}}Y_{\mathrm{cell}}$ is a function of its cell (\Cref{Y-P2-L7}\,\textup{(b)} and \Cref{Y-P2-L6}\,B). We apply \Cref{support:4-17}\,\textup{(f)} over the $q$ occupied cells: the mean is at most $\prod_{\mathrm{cells}}\|\Phi_{\mathrm{cell}}Y_{\mathrm{cell}}\|_{\mathrm{ch}}$, where $\|\cdot\|_{\mathrm{ch}}^{N_c}=\langle\operatorname{Fold}(\cdot)\rangle$ (\Cref{support:4-13}).

We fix an occupied cell and denote by $n_{\mathrm{cell}}$ the number of its sources. First let $n_{\mathrm{cell}}\ge1$. By \Cref{Y-P2-L7}\,\textup{(d)}, $\langle\operatorname{Fold}(\Phi Y)\rangle$ is a mean of the form of \Cref{Y-P2-L5}: its observable is $\prod_{\varkappa'}F^{[\varkappa']}(V_{K-s})$, whose sup is at most $(\sup|F|)^{N_c}$, and it has $n_{\mathrm{cell}}N_c$ sources, which by \Cref{Y-P2-L8}\,\textup{(c)} satisfy the data $(\check\delta;\check\rho)$ and have the sup norms of the sources of the cell. \Cref{Y-P2-L5} bounds it by $(\sup|F|)^{N_c}G_{1,s}^{n_{\mathrm{cell}}N_c}e^{c_{\mathrm{vol}}^{(s)}N^3}\prod\|g_t\|^{N_c}$. As $N^3/N_c=h^3$, the $N_c$-th root is $\sup|F|\,G_{1,s}^{n_{\mathrm{cell}}}e^{c_{\mathrm{vol}}^{(s)}h^3}\prod_{t\in I_{\mathrm{cell}}}\|g_t\|$. If $n_{\mathrm{cell}}=0$, then $Y=1$, and \Cref{Y-P2-L7}\,\textup{(d)} gives $\langle\operatorname{Fold}\Phi\rangle=\langle\prod_{\varkappa'}F^{[\varkappa']}(V_{K-s})\rangle\le(\sup|F|)^{N_c}$, which is below the bound claimed: its $N_c$-th root is at most $\sup|F|\le\sup|F|\,e^{c_{\mathrm{vol}}^{(s)}h^3}$. The product of these bounds over the occupied cells gives the claim. For \textup{(R2)\,(ii)} we use \Cref{Y-P2-L8}\,\textup{(b$'$)}. \end{proof}

\begin{remark}[Why one cell] If the support of a source meets $r>1$ cells, then $\mathcal O_K(g)$ is a \emph{sum} over cells, and the moment is not a product of functions of cells.\end{remark}

\begin{remark}[The placing of the sources] For the cumulants and for $\sigma$ we use unit cells, we translate the \emph{sources} by $a_K\mathbb Z^3$, and we cut into pieces of side $\check\delta_n/2$. This is allowed, since $S^T_n$ and $\sigma$ are invariant under $a_K\mathbb Z^3$. It is not allowed for moments centred at $\hat c$, which are invariant under $\mathbb Z^3$ alone. For moments centred at $\hat c$, with the sources placed anywhere and $n=1$ included, we move the board instead, by $o\in\mathbb Z^3$, with $h>n(2R_f+d+1)$. Both forms hold: they are \textup{(R2)\,(i)} and \textup{(R2)\,(ii)} of \Cref{Y-P2-ThR}.\end{remark}

\begin{lemma}[Uniform bounds, short runs and local boundedness in $\mathsf b$]\label{Y-P2-L10} \Cref{Y-P2-ThR}\,\textup{(R3)} holds: the bounds on $\sigma_{K,N,\beta}$ and on $S^T_n$ hold for all $K\ge1$, for all $N$ of the data and wherever the balls lie, and they hold for the limits of \textup{(T3)} of \Cref{Y-H9-T} too. Moreover, $\sum_{\pi\in\mathrm{Part}(n)}(|\pi|-1)!\le n!\,1.45^{n+1}$ for every $n\ge2$, and $C_n^{(2)}$ is locally bounded in $\mathsf b$.\end{lemma}

\begin{proof} \emph{Step (Case $K\ge j^\circ(\mathsf b,\check\delta_n)$): pieces on the unit board.} Let $K$ be finite. Since $S^T_n$ and $\sigma$ are $n$-linear in bounded entries, each of them is a sum of its values at the tuples of pieces of \Cref{Y-P2-L8}\,\textup{(a)} applied with $\check\delta=\check\delta_n$, and there are at most $(4R_f/\check\delta_n+2)^{3n}$ such tuples. We fix one tuple and take $v$ as in \Cref{Y-P2-L8}\,\textup{(b)}. By Step~7 (``At finite $K$'') of the proof of \Cref{Y-H9-T}, which holds for bounded entries, $S^T_n$ and $\sigma$ are unchanged under $\tau_v$: both are computed from means, and these are invariant under $a_K\mathbb Z^3$.

For the translated pieces we put $Y_i=\mathcal O_K(g_i)-\hat c_K(g_i)$. Cumulants of order $\ge2$ do not depend on shifts by constants; hence \[S^T_n=\sum_\pi(-1)^{|\pi|-1}(|\pi|-1)!\prod_{J\in\pi}\langle\prod_{i\in J}Y_i\rangle.\] We bound each factor by \Cref{Y-P2-ThR}\,\textup{(R2)\,(i)}, on the \emph{same} board $B_0$ for all factors, applied with the data $(\check\delta_n;\check\delta_n/4)$ and with $q_J\le|J|$: each factor is at most $(G_1e^{c_{\mathrm{vol}}})^{|J|}\prod_J\|g_i\|$. For $\sigma$ we write, for each of the translated pieces $g=g_i$ and with $Y=Y_i$, $\mathcal O_K(g)-\langle\mathcal O_K(g)\rangle=Y-\eta$, where $|\eta|=|\langle Y\rangle|\le G_1e^{c_{\mathrm{vol}}}\|g\|$ by \textup{(R2)} applied with $Q=1$; expanding the product gives $2^n$ terms, each at most $(G_1e^{c_{\mathrm{vol}}})^n\prod\|g_i\|$.

\emph{Step (Case $K<j^\circ$): short runs.} The barycentres of the plaquettes of one orientation form a translate of $a_K\mathbb Z^3$; a half-open cube of side $\check\delta/2$ contains at most $(\check\delta L^K/2+1)^3$ of them; there are three orientations, and $0\le a_p\le2$. Hence $|\mathcal O_K(g)|\le6L^K(\check\delta L^K/2+1)^3\|g\|\le B_0(\mathsf b,\check\delta)\|g\|$. The expansion of $S^T_n$ over the partitions, now with $Y=\mathcal O_K(g)$, gives the bound on $S^T_n$. For $\sigma$ we bound the product of the $n$ factors directly, each factor by $|\mathcal O_K(g)-\langle\mathcal O_K(g)\rangle|\le2B_0(\mathsf b,\check\delta)\|g\|$; this gives $|\sigma|\le\prod_i\sup|\mathcal O_K(g_i)-\langle\mathcal O_K(g_i)\rangle|\le(2B_0(\mathsf b,\check\delta))^n\prod\|g_i\|$.

In neither case does $N$ enter; if supports are far apart, the one effect is that they lie in different cells, and no factor depends on their distance. So the bounds of the two cases hold for all $N$ of the data and wherever the balls lie.

\emph{Step (The sum): the combinatorial bound.} Let $a_n:=\sum_\pi(|\pi|-1)!$. Its exponential generating function is $-\log(2-e^x)$, with derivative $2/(2-e^x)-1$; hence $a_{n+1}=2f_n$ for $n\ge1$, where $f_n=\sum_{m\ge0}m^n2^{-m-1}$. As $t^ne^{-t\log2}$ is unimodal, the sum is at most the integral plus the maximum: \[f_n\le\frac12[n!/(\log2)^{n+1}+(n/(e\log2))^n]\le1.2214\,n!\,(1/\log2)^n.\] Therefore $a_n\le2.443\,(n-1)!\,1.4427^{n-1}\le n!\,1.45^{n+1}$ for $n\ge2$.

\emph{Step (Local boundedness in $\mathsf b$): the constants.} The integer $j^\circ(\mathsf b,\check\delta)$ is non-decreasing in $\mathsf b$ and finite (\Cref{Y-H9-LN}\,\textup{(N2)}; \textup{(P1.0)}, that is, \Cref{Y-P1-Jo}\,\textup{(G$^\circ$)}), hence bounded on compact sets. The functions $\mathrm{amp}$, $c_a$, $c_h$, $p_0$ and $\varepsilon$ are continuous, and $\check r$ is a bounded step function. The tail of $C_W$ is bounded uniformly: $N_1^{[j^\circ+t]}\le L^{3t}N_1^{[j^\circ]}$ (item~\textup{(ii)} of the proof of \Cref{Y-H9-LN}\,\textup{(N3)}) is compensated by $e^{-(\kappa-1)(L^t\mathrm{sep}^\circ-1)}$ (\Cref{Y-H9-LN}\,\textup{(N2)}), with $c_a\le1/32$. In this step \textup{(N2)} of \Cref{Y-H9-LN} and item~\textup{(ii)} of the proof of its \textup{(N3)} are taken with $(\check\delta,\check\rho,\mathrm{sep}^\circ,j^\circ)$ in the place of $(d,R_f,\mathrm{sep}_*,j^*)$, and their proofs hold as they stand: $\mathrm{sep}^{[i]}$ and $N_1^{[j]}$ are the same functions, the proof of \textup{(N2)} uses no property of $\mathrm{sep}_*$ other than $\mathrm{sep}_*>0$, and item~\textup{(ii)} uses no property of $j^*$.

\emph{Step (The limits of \textup{(T3)}): the bounds pass to the limits.} The limits of \textup{(T3)} are limits of numbers bounded in this way, and so they obey the same bounds; for continuous entries the bounds hold for the extension given by \textup{(T1)}. \end{proof}

\begin{remark}[How $\mathrm{sep}^\circ$ enters the closed form of \S\ref{Y-P3-1}] Let $j^\circ\ge2$, and put $t=\check\delta L^{j^\circ-2}/M$. By the choice of $j^\circ$, $\mathrm{sep}^{[j^\circ-2]}<\mathrm{sep}^\circ$; with $\check r(x)\le\log x/1056+4$ this gives $t<\mathrm{sep}^\circ+11+\log(L^{j^\circ-2}\mathsf b)/528$. Hence $t(1-1/528)<\mathrm{sep}^\circ+11+[\log\mathsf b+\log(M/\check\delta)]/528$. Thus $\mathrm{sep}^\circ$ enters $L^{j^\circ}$, $r^\circ$, $N_1^{[j]}$ for $j\le j^\circ$, and $B_0(\mathsf b,\check\delta)$ as a letter.\end{remark}

\begin{lemma}[Every positive label, by dilation]\label{Y-P2-L11} \Cref{Y-P2-ThR}\,\textup{(R4)} holds, with the integer $z$ subject to the conditions $L^z\mathsf b\ge x_s$, $K-z\ge1$, $m+z\ge m'$ and $L^zl\ge2Ms(L^{z+1}\mathsf b)$. \end{lemma}

\begin{proof} We take the dilation $x\mapsto L^zx$ of Step~5 of the proof of \Cref{Y-H9-T}. It maps the lattice of $(K,l)$ onto the lattice of $(K-z,L^zl)$, sending plaquettes to plaquettes and barycentres to barycentres. The two graphs thus correspond to each other, and $\beta$ is the same for both; as Wilson's measure is determined by the graph and by $\beta$, the dilation carries the one measure onto the other. Under this correspondence $\mathcal O_{K-z}(f^{[z]})=L^{-z}\mathcal O_K(f)$, and a cumulant of order $n$, like a centred moment of order $n$, is homogeneous of degree $n$ in its entries. It follows that \[S^T_{n;K,l,\beta}(f)=L^{zn}S^T_{n;K-z,L^zl,\beta}(f^{[z]}),\] where $\beta=L^{K-z}(L^z\mathsf b)$. In the same way the centred moment of order $n$ obeys $\sigma_{K,l,\beta}(f)=L^{zn}\sigma_{K-z,L^zl,\beta}(f^{[z]})$. Under the dilation the type changes as $(d;R_f)\to(L^zd;L^zR_f)$, while the sup norms stay the same.

By the four conditions, $(L^z\mathsf b,L^zl,K-z)$ are data: $L^z\mathsf b\ge x_s$ is the condition on the label, $K-z\ge1$ the condition on the number of steps, $m+z\ge m'$ the condition on the exponent of the side $L^zl$, and $L^zl\ge2Ms(L^{z+1}\mathsf b)$ is \textup{(Tor)} for the label $L^z\mathsf b$. So both bounds of \textup{(R3)} apply to the dilated entries for these data. With $\check\delta_n=\min\{L^zd,1/(4n)\}$, the bound for $\sigma$ reads \[|\sigma_{K-z,L^zl,\beta}(f^{[z]})|\le(4L^zR_f/\check\delta_n+2)^{3n}(2G(L^z\mathsf b,\check\delta_n))^n\prod_i\|f_i\|_\infty,\] and, since $2^n\le1.45^{n+1}n!$ for $n\ge2$, the factor in front of the product obeys \[(4L^zR_f/\check\delta_n+2)^{3n}(2G(L^z\mathsf b,\check\delta_n))^n\le C_n^{(2)}(L^z\mathsf b;L^zd,L^zR_f).\] Together with the two identities above, the two bounds of \textup{(R3)} thus give $|S^T_{n;K,l,\beta}(f)|\le L^{zn}C_n^{(2)}(L^z\mathsf b;L^zd,L^zR_f)\prod_i\|f_i\|_\infty$ and the same bound for $|\sigma_{K,l,\beta}(f)|$. One integer suffices for all labels $\mathsf b'\in U(\mathsf b)$, namely $z=z(\mathsf b)$ of \Cref{Y-00-S7}, once the sides are large enough. The constant does not depend on the side. \end{proof}

\begin{lemma}[Compactness; the orders of the limits]\label{Y-P2-L12} \Cref{Y-P2-ThR}\,\textup{(R5)} holds in each of its parts: \begin{enumerate} \item[\textup{(a)}] the continuum limit taken first; \item[\textup{(b)}] the volume limit taken first; \item[\textup{(c)}] the two limits taken jointly; \item[\textup{(d)}] the joint limits with $K_j/l_j^3\to\infty$. \end{enumerate} \end{lemma}

\begin{proof} \emph{Step (Carrying): entries on the torus.} For a function $f$ on $\mathbb R^3$ with compact support, $f^{\langle l\rangle}(x)$ denotes the sum of the values of $f$ at all the lifts of $x$. If $l$ is larger than twice the diameter of the union of the supports, then on that union the covering map is an isometry for the sup distance, and so the periodised tuple has the type of the tuple.

\emph{Step (Frames): countably many frames.} A frame is an $n$-tuple $(A_1,\dots,A_n)$ of sets, each a finite union of closed cubes with rational corners, the sets being pairwise disjoint; there are countably many frames. Every tuple of compact supports with the separation $d$ is carried by a frame: for $A_i$ we take the union of those closed cubes of a rational grid of side $<d/4$ that meet $\operatorname{supp}f_i$, and then any two of the $A_i$ are $\ge d/2$ apart. For each $i$, the space $C(A_i)$ of the continuous functions vanishing off $A_i$ is separable.

Let a frame be of the type $(d_A;R_A)$. On the entries that it carries, all the forms $S^T_{n;K,l,L^K\mathsf b'}$ and $S^{(l,\mathsf b')}_n$ are $n$-linear and bounded by $L^{zn}C_n^{(2)}(L^z\mathsf b';L^zd_A,L^zR_A)\prod\|\cdot\|_\infty$. At $z=z(\mathsf b)$ this bound is uniform in $(K,l)$ with $K\ge z+1$ and $l$ large, and in $\mathsf b'\in U(\mathsf b)$, by \textup{(R3)}, \textup{(R4)}, \Cref{Y-00-S7} and \Cref{Y-P2-L10}: the last assertion of the lemma, the local boundedness of $C_n^{(2)}$ in the label, is used on the set $L^{z(\mathsf b)}U(\mathsf b)$, which by the definition is compact and contained in $[x_s,\infty)$. So the forms are equicontinuous on each frame. Taking a diagonal sequence over the frames and over countable dense subsets of the spaces $C(A_i)$, we obtain convergence on these subsets, and the equicontinuity extends it to all entries.

\emph{Step \textup{(a)}: the continuum limit first.} We apply Step (Frames) to the forms $S^{(l_j,\mathsf b_j)}_n$, which exist by \Cref{Y-H9-T}\,\textup{(T3)}, \textup{(T6)}, for continuous entries as well.

\emph{Step \textup{(c)}: the joint limits.} We apply Step (Frames) to the forms $S^T_{n;K_j,l_j,L^{K_j}\mathsf b_j}$.


\emph{Step \textup{(b)}. Volume limit first.}
At fixed $K$ and $\beta$, periodize the torus configurations on
$a_K\mathbb Z^3$. \Cref{support:14-16}(c) gives a weak
limit along a subsequence of sides, and the limit $\nu_K$ is a
translation-invariant Gibbs law. Choose the same subsequence for
all $K$ by diagonal extraction. \Cref{support:54-17}
gives its cubic invariance and reflection positivity at coordinate
site and bond mirrors. Every fixed lattice observable $\mathcal O_K(f)$
with compactly supported $f$ is continuous and local. Thus its
moments and cumulants converge on this subsequence, and the uniform
bounds pass to $\nu_K$. The frame argument then extracts the limit
in $K$ on all separated tuples.

\emph{Step \textup{(d)}: comparison along the same sides and labels.}
Suppose $\mathsf b_j\to\mathsf b>0$ and put $z=z(\mathsf b)$ and
$N_j=L^zl_j$. For large $j$ the dilated labels $L^z\mathsf b_j$ belong
to the compact set $L^zU(\mathsf b)\subset[x_s,\infty)$ and the sides
satisfy \textup{(Tor)} uniformly on this set. For Lipschitz entries,
\textup{(T2)}, \textup{(T3)} and \textup{(T6)} of
\Cref{Y-H9-T} give
\[
\big|S^T_{n;K_j,l_j,L^{K_j}\mathsf b_j}(f)
      -S^{(l_j,\mathsf b_j)}_n(f)\big|
\le L^{zn} C'_{n,j}(K_j-z)\theta^{K_j-z}
    \prod_i\max\{\|f_i\|_\infty,40ML^{-z}\mathrm{Lip}f_i\}.
\]
The cutoff threshold is independent of the side and locally bounded
in the label. The displayed expression for $C'_n$ in \textup{(T2)}
gives $\log C'_{n,j}\le C N_j^3+C'$ with $C,C'$ uniform on the compact
set of labels. Since $N_j=L^zl_j$ with fixed $z$, the hypothesis $K_j/l_j^3\to\infty$ implies $K_j/N_j^3\to\infty$, and the right side tends to zero.
The two sequences thus have the same limit points along the given
sides and labels.

For continuous entries we approximate uniformly inside a frame and invoke the equicontinuity. \end{proof}

\begin{lemma}[What every limit inherits]\label{Y-P2-L13} \Cref{Y-P2-ThR}\,\textup{(R6)} holds in its parts \textup{(i)} to \textup{(vii)}. \end{lemma}

\begin{proof} \emph{Step \textup{(i)}, \textup{(ii)}: reality, symmetry and the bound.} The forms $S_n$ are limits of real, symmetric, $n$-linear forms that all obey one bound; hence they are real, symmetric and $n$-linear themselves, and they obey the same bound.

We fix a convention for the rest of this proof and for \Cref{Y-P2-L14}: $C_n^{(2)}$ at $(d;R)$ denotes the constant of \textup{(ii)}. For a fixed label $\mathsf b\ge x_s$ this is the number $C_n^{(2)}(\mathsf b;d,R)$ itself. For a fixed label $0<\mathsf b<x_s$ it is $L^{zn}C_n^{(2)}(L^z\mathsf b;L^zd,L^zR)$, with $z$ as \textup{(R4)} admits. For moving labels it is the supremum of the latter over $\mathsf b'\in U(\mathsf b)$, taken at $z=z(\mathsf b)$. The arguments below depend on this constant only through the fact that it is independent of $K$, of $l$ and of the places of the entries.

\emph{Step \textup{(iii)}: invariance.} At finite $(K,l)$ the identity $S^T_n(rX)=S^T_n(X)$ is exact for $r^\circ\in W_3$, $v\in a_K\mathbb Z^3$ ($r$: $x\mapsto r^\circ x+v$), and these sets of maps grow with $K$. In the order \textup{(a)}, the identity at finite $K$ passes to the limit in $K$ of \Cref{Y-H9-T}\,\textup{(T3)}, \textup{(T6)}. Thus $S_n(rX)=S_n(X)$ for $r^\circ\in W_3$ and for every $v$ in the union $A$ of the lattices $a_K\mathbb Z^3$, a dense set.

For all $v$ and for $C^1$ entries, the argument is that of the part ``All $v$'' of Step~7 of the proof of \Cref{Y-H9-T}, changed in one point: where that step uses $C_n$, we use $C_n^{(2)}$ at $(d/2;R_f+d/4)$. For a sequence $v_j\in A$ with $v_j\to v$ the outcome is \[|S_n(\tau_vX)-S_n(\tau_{v_j}X)|\le nC_n^{(2)}\sqrt3\,|v-v_j|\max_i[\mathrm{Lip}f_i\prod_{k\ne i}\|f_k\|_\infty]\to0,\] and since $S_n(\tau_{v_j}X)=S_n(X)$, the identity holds at $v$. For continuous entries we conclude by the uniform continuity of the entries together with \textup{(ii)}.

\emph{Step \textup{(iv)}: reflection positivity.} Let $X_1,\dots,X_r\in\mathfrak S(\{y_\mu>0\})$; when $l$ is large, their supports lie in $\{0<y_\mu<l/2\}$. At finite $(K,l)$, Step~8\,\textup{(a)}, \textup{(b)} of the proof of \Cref{Y-H9-T} yields $\sum\bar\zeta_i\zeta_k\,\sigma_{K,l,\beta}(\theta X_i\cdot X_k)\ge0$, where $\sigma_{K,l,\beta}$ is $\sigma_S$ taken at $S^T$. Now $\sigma$ is a polynomial in the $S_{|J|}$, each factor of which converges, and a limit of non-negative numbers is non-negative, whether the limit is single, iterated or joint. For an arbitrary $c$ we write $\theta_c=\tau_v\theta_0\tau_{-v}$ with $v=ce_\mu$, and \textup{(iii)} carries the inequality from $\theta_0$ over to $\theta_c$.

\emph{Step \textup{(v)}: Step~9 of the proof of \Cref{Y-H9-T} on $\mathbb R^3$.} With $\mathbb R^3$ in place of $T_l$, that step changes as follows. Let $\xi$ be a point whose entries are pairwise different, let $\delta=\min|\xi_i-\xi_k|/4$, let $B_i$ be the open ball with centre $\xi_i$ and radius $\delta$, and let $Q=\prod B_i$. We choose $\chi_i$ smooth, with $0\le\chi_i\le1$, equal to $1$ on $B_i$ and carried by the closed ball with centre $\xi_i$ and radius $3\delta/2$; any two of these closed balls are at least $\delta$ apart. For $F\in C_c^\infty(Q)$ let $F^{\mathrm{per}}$ be the function that repeats $F$ with the period $4\delta$ in every one of the $3n$ coordinates, let $\hat F(k)$ be its Fourier coefficients, and let $e_k(x)=e^{2\pi ik\cdot x/(4\delta)}$.

In the $i$-th variable the nearest copy of $\operatorname{supp}F$ in $F^{\mathrm{per}}$ other than $\operatorname{supp}F$ itself, namely $\operatorname{supp}F$ shifted by $4\delta$ in one coordinate of that variable, begins at the distance $3\delta$ from $\xi_i$, while the support of $\chi_i$ ends at $3\delta/2$; therefore $F=(\chi_1\otimes\cdots\otimes\chi_n)F^{\mathrm{per}}$. We define \[D_Q(F):=\sum_k\hat F(k)\,S_n^{\mathbb C}(\chi_1e_{k_1},\dots,\chi_ne_{k_n}).\] The bounds of the part ``The local construction'' of Step~9, which hold there with the period $l$, the constant $c(n,l)$ and the bound $C_n(\delta)$, give here, with $C_n^{(2)}$ at $(\delta;3\delta/2)$, \[|D_Q(F)|\le2^nC_n^{(2)}\,c(n,4\delta)\max_{|\alpha|\le3n+1}\sup|\partial^\alpha F|.\] The parts ``On products'', ``Uniqueness on $Q$'', ``Gluing'' and ``The functional $\sigma_S$'' of Step~9 then carry over word for word.

\emph{Step \textup{(vi)}: dilation.} At finite $K$ we use Step~5 of the proof of \Cref{Y-H9-T}, matching the subsequence $(m_j)$ belonging to the label $\mathsf b$ with the subsequence $(m_j+z)$ belonging to the label $L^z\mathsf b$.

\emph{Step \textup{(vii)}: the group of symmetries.} The group in \textup{(vii)} is closed: if $R_j\to R$, then $\|f\circ R_j^{-1}-f\circ R^{-1}\|_\infty\to0$ on carriers that stay disjoint, and \textup{(ii)} carries the invariance to the limit. The group contains $W_3$, by \textup{(iii)}. Write $H$ for its rotation subgroup and $O=W_3\cap\operatorname{SO}(3)$, so $O\subset H$ and $|O|=24$. We give the subgroup argument explicitly.

If $H$ is infinite, compactness gives nonidentity rotations $h_j\to1$. Write $h_j=\exp(\theta_jJ_{v_j})$, where $\theta_j\downarrow0$ after taking a subsequence, $|v_j|=1$ and $J_vx=v\times x$. Passing to a further subsequence gives $v_j\to v$. For each real $t$, choose integers $m_j$ with $m_j\theta_j\to t$. Then $h_j^{m_j}\to\exp(tJ_v)$, which belongs to $H$ by closedness. Conjugating by $O$ gives the circles with axes $wv$, $w\in O$. Their axes span $\mathbb R^3$, because the action of $O$ is irreducible: averaging a vector with its half-turn about a coordinate axis projects onto that axis, and cyclically permuting the axes then generates all three from any nonzero invariant subspace. Choose three independent axes $v_1,v_2,v_3$ from this orbit. The map
\[
 (t_1,t_2,t_3)\longmapsto
       e^{t_1J_{v_1}}e^{t_2J_{v_2}}e^{t_3J_{v_3}}
\]
has invertible derivative at zero. The inverse function theorem therefore puts a neighborhood of $1$ in $H$. As $\operatorname{SO}(3)$ is connected, $H=\operatorname{SO}(3)$.

Suppose instead that $H$ is finite, of order $q\ge24$. Let $P\subset S^2$ be the set of poles of its nonidentity rotations, and let $n_1,\ldots,n_r\ge2$ be the stabilizer orders for its orbits in $P$. Each stabilizer is a finite group of rotations about one axis, hence cyclic. Counting pairs $(h,x)$ with $h\ne1$ and $hx=x$ gives
\[
 \sum_{i=1}^r\frac q{n_i}(n_i-1)=2(q-1),
 \qquad \sum_{i=1}^r\left(1-\frac1{n_i}\right)=2-\frac2q.
\]
Thus $2\le r\le3$. If $r=2$, then $1/n_1+1/n_2=2/q$ and $n_i\le q$ force $n_1=n_2=q$, giving an invariant axis, contrary to the irreducibility of $O$. For $r=3$, order the $n_i$. The inequality $\sum_i1/n_i>1$ gives $n_1=2$ and $n_2\le3$. The case $(2,2,n_3)$ gives $n_3=q/2\ge12$. Its cyclic stabilizer has index two, hence is normal, so its axis is invariant under $H$, again impossible. The remaining triples are $(2,3,3)$, $(2,3,4)$ and $(2,3,5)$, giving $q=12,24,60$. Since $24$ divides $q$, only $q=24$ remains, and $H=O$.

Finally, $-1\in W_3$, so the whole symmetry group is $H\cup(-1)H$. It is therefore $W_3$ or $\operatorname{O}(3)$. \end{proof}

\begin{lemma}[Hilbert space, vacuum, energy and momenta]\label{Y-P2-L14} The data $(D,\langle\cdot,\cdot\rangle,\mathsf 1=\text{empty tuple},\mathcal T^\circ_t=\tau_{te_\mu},U^\circ(v)=\tau_v)$, with $D$ and the form of \Cref{Y-P2-ThR}\,\textup{(R7)} and with $v$ in the mirror plane, satisfy the hypotheses of \Cref{support:27-11} applied with $N=2$, the dimension of the mirror plane. The Hilbert space that they define is separable, and the construction extends to the class of continuous entries. Consequently, by \Cref{support:27-11}\,\textup{(a)} to \textup{(d)}, \Cref{Y-P2-ThR}\,\textup{(R7)} holds. \end{lemma}

\begin{proof} \emph{Step (The form): Hermitian, positive, normalised.} The form is Hermitian: $\sigma_S$ is real, and $\sigma_S(\theta X\cdot X')=\sigma_S(X\cdot \theta X')=\sigma_S(\theta X'\cdot X)$, the first equality by the invariance under $\theta\in W_3$ and the second by the symmetry. The form is positive by \textup{(R6)}\,\textup{(iv)}. Moreover $\sigma_S(\emptyset)=1$.

\emph{Step (The maps): translations of tuples.} The maps $\tau_{te_\mu}$, for $t\ge0$, and $\tau_v$ send $\mathfrak S(\mathbb E_+)$ into itself and act linearly on $D$; the first family is a semigroup and the second a group; the two commute, and both fix the empty tuple.

\emph{Step \textup{(1)}: symmetry and isometry.} From $\theta\tau_{te_\mu}=\tau_{-te_\mu}\theta$ and \textup{(R6)}\,\textup{(iii)} we get $\sigma_S(\theta X\cdot \tau_tX')=\sigma_S(\tau_{-t}\theta X\cdot X')=\sigma_S(\theta\tau_tX\cdot X')$. From $\theta\tau_v=\tau_v\theta$ and \textup{(R6)}\,\textup{(iii)} again, $\sigma_S(\theta\tau_vX\cdot \tau_vX')=\sigma_S(\theta X\cdot X')$, so that $\tau_v$ is an isometry.

\emph{Step \textup{(2)}: bounded matrix elements.} Let $F=\sum\zeta_iX_i$, where $X_i$ has $n_i$ entries and the separation $d_i$. For all $t\ge0$ the tuple $\theta X_i\cdot \tau_tX_k$ has $n_i+n_k$ entries and is of the type $(\min\{d_i,d_k,2\operatorname{dist}(\text{supports},\text{mirror})\};R)$, where $R$ is a radius such that each entry of $X_i$ and of $X_k$ has its support in a ball of radius $R$. The bound that \textup{(R3)}, \textup{(R4)} give for $\sigma$ does not depend on the places. Hence the number $c_F:=\sup_{t\ge0}|\langle F,\mathcal T^\circ_tF\rangle|$ of hypothesis \textup{(2)} of \Cref{support:27-11} obeys $c_F\le\sum|\zeta_i||\zeta_k|C(i,k)<\infty$, where $C(i,k)$ is that bound for the pair $(i,k)$. What matters here is that the walls of the unit board $B_0$ in the proof of \textup{(R3)} enclose each piece of \Cref{Y-P2-L8}\,\textup{(a)}, and that no cell has to contain the whole tuple. Indeed, \Cref{Y-P2-L8}\,\textup{(b)} places each piece in one open cell of $B_0$.

\emph{Step \textup{(3)}: continuity.} By the definition of $\mathfrak S$, the entries of the tuples of $\mathfrak S(\mathbb E_+)$ are $C^1$. We pass from $\tau_tX_k$ to $X_k$ by changing one entry at a time, which is the telescoping of Step~7 of the proof of \Cref{Y-H9-T}; each term then has one entry of the form $\tau_{te_\mu}f-f$, and for small $t$ the carriers remain $d'/2$ apart, $d'$ being the separation of $\theta X_i\cdot X_k$. This yields \[|\sigma_S(\theta X_i\cdot \tau_tX_k)-\sigma_S(\theta X_i\cdot X_k)|\le C\sqrt3\,t\max\mathrm{Lip}\to0.\] For $\tau_v$, the telescoping in the entries gives $|\sigma_S(\theta X_i\cdot \tau_vX_k)-\sigma_S(\theta X_i\cdot X_k)|\le C\sqrt3\,|v|\max\mathrm{Lip}$, which tends to $0$ with $v$.

\emph{Step (Conclusions): the operators.} \Cref{support:27-11}\,\textup{(a)} to \textup{(c)} now provides $\|\mathcal T^\circ_tF\|\le\|F\|$, the operators $\mathcal T_t$ and $U(v)$, the identity $\mathcal T_t=e^{-tH}$ with $H\ge0$ and $H\Omega=0$, the momenta $P_1$ and $P_2$, and a joint spectral measure carried by $[0,\infty)\times\mathbb R^2$. For \Cref{support:27-11}\,\textup{(d)}, we use the group $W_2$ of \textup{(R7)}, whose elements, acting on the mirror plane, are the $8$ signed permutation matrices of that plane. Let $W^\circ$ be the action of an element $w$ of $W_2$. It is a linear bijection of $D$ that fixes the empty tuple $\mathsf 1$; it commutes with $\theta$ and with $\tau_{te_\mu}$, it maps $\mathbb E_+$ onto itself, and $W^\circ\tau_v=\tau_{wv}W^\circ$. By $W^\circ\theta=\theta W^\circ$ and \textup{(R6)}\,\textup{(iii)} it leaves the form unchanged, $\langle W^\circ F,W^\circ F'\rangle=\langle F,F'\rangle$; so we obtain a unitary representation of $W_2$. Finally, \Cref{support:27-11}\,\textup{(a)} gives $|\langle F,\mathcal T_tF\rangle|\le\|F\|^2$, and therefore $c_F=\langle F,F\rangle$.

\emph{Step (Separable): a countable dense set.} By \textup{(R6)}\,\textup{(ii)}, for tuples inside one frame (Step (Frames) of the proof of \Cref{Y-P2-L12}), $\|[X]-[X']\|^2$ is bounded in terms of the sup norms of the differences of the entries. There are countably many frames, and every $C(A_i)$ is separable; hence the Hilbert space is separable.

\emph{Step (Continuous entries): extension by continuity.} Let $f$ be continuous with compact support in $\mathbb E_+$. We choose $C^1$ functions $f^e\to f$ uniformly, all carried by one compact subset of $\mathbb E_+$. Then \[\|[f^e]-[f^{e'}]\|^2=S_2(\theta(f^e-f^{e'}),f^e-f^{e'})\le C_2^{(2)}\|f^e-f^{e'}\|_\infty^2,\] and for tuples we change the entries one at a time. The classes therefore form a Cauchy sequence, the limit does not depend on the approximating functions, and the form extends by continuity. Moreover $\langle\Omega,[f]\rangle=\sigma_S(f)=0$, because $\mathrm{Part}_{\ge2}(1)$ is empty.

Every mirror $c$ can be used equally well, by \textup{(R6)}\,\textup{(iii)}. An element of $W_3$ that carries $e_\mu$ to $e_\nu$ intertwines the three constructions, one for each direction. \end{proof}

\begin{lemma}[Reflection positivity without a collar; the limit laws on $\mathbb R^3$]\label{Y-P2-L15} \Cref{Y-P2-ThR}\,\textup{(R8)} holds in its parts \textup{(b)} and \textup{(c)}. \end{lemma}

\begin{proof} \emph{Step \textup{(b)}: positivity without a collar.} Let $c\in L^{-s}(\mathbb Z+\frac12)\subset a_K(\mathbb Z+\frac12)$. Since $N/2=L^{m_l}$ is a multiple of $L^{-s}$, we also have $c+N/2\in L^{-s}(\mathbb Z+\frac12)$. Hence $T_+$ and $\theta T_+$ are disjoint open sets whose boundaries are planes $y_\mu\in L^{-s}(\mathbb Z+\frac12)$. For $F$ as in \textup{(b)} let $\Phi:=E[F(V_{K-s})|U]$. By \Cref{Y-P2-L6}\,\textup{B}, $\Phi$ is bounded and Borel and depends only on the bonds whose two ends lie in $T_+$; it therefore lies in the class $\mathfrak F_+$ of \Cref{support:4-8}. The reflection $\theta$ is $y_\mu\mapsto2c-y_\mu$, and $2c\in L^{-s}\mathbb Z$; hence $E[\theta F|U]=\theta\Phi$, by \Cref{Y-P2-L6}\,\textup{D}, and $E[\overline{\theta F}F|U]=\overline{\theta\Phi}\Phi$, by \Cref{Y-P2-L6}\,\textup{C}. Consequently $E[\overline{\theta F}F]=\langle\overline{\theta\Phi}\Phi\rangle\ge0$, by \Cref{support:4-8}\,\textup{(a)}.

\emph{Step \textup{(c)} Densities: bounds on the marginals.} Let $K\ge s+1$, and let $B$ be a finite set of bonds of depth $s$. The marginal of $V_{K-s}$ on $B$ has the density \[p_B(V)=E[\prod_{c\in B}k_t(V(c),\mathsf M(V_{K-s-1})(c))],\] where $t=ax_{K-s-1}=aL^{s+1}\mathsf b$ (the definition ``Runs'' in item~2 of \S\ref{Y-J-1}). Because $1-\cos\operatorname{ang}(g)=1-\frac12\operatorname{tr}g\in[0,2]$, we have $z(t)^{-1}e^{-2t}\le k_t\le z(t)^{-1}$. Moreover $|k_t(g,h)-k_t(g',h)|\le z^{-1}t\,|\frac12\operatorname{tr}((g-g')h^{-1})|\le z^{-1}t\,\|g-g'\|$, uniformly in $h$. Hence $p_B$ is bounded from above and from below and is equi-Lipschitz, uniformly in $K$ and in $N$.

\emph{Step \textup{(c)} $K\to\infty$ at fixed $N$.} The law $\mu^{[K]}_{[s]}$ of $V_{K-s}$ converges in total variation by \Cref{Y-J-CorC}\,\textup{(b)}, which holds for $\mathsf b$ and $N$ of the data by \Cref{Y-H9-T}\,\textup{(T7)}. So the marginal densities converge in $L^1$; by Arzel\`a--Ascoli they converge uniformly along subsequences, and the limit has a density that obeys the same bounds.

\emph{Step \textup{(c)} $N_j\to\infty$.} The pairs $(s,B)$ form a countable set. Arzel\`a--Ascoli and a diagonal sequence give the uniform convergence of all the densities $p_B^{(N_j)}$, which is the same as convergence in total variation on every finite set of bonds. Marginalisation passes to the limit, and Kolmogorov's theorem yields the law on $G^{\text{bonds of }L^{-s}\mathbb Z^3}$.

\emph{Step \textup{(c)} Consistency: from depth $s$ to depth $s-1$.} Let $F$ be a bounded Borel function on a finite set $B'$ of bonds of depth $s-1$. For each side $N$, \Cref{Y-J-CorC}\,\textup{(c)} gives $\int F\,d\mu^\infty_{[s-1]}=\int(KF)\,d\mu^\infty_{[s]}$, with \[KF(V)=\int F(V')\prod_{c\in B'}k_{aL^s\mathsf b}(V'(c),\mathsf M(V)(c))\,dV';\] by \Cref{Y-P2-L6}\,\textup{A}, this function is bounded, Borel and local. The map $\mathsf M$ is discontinuous, so that weak convergence would not carry the identity to the limit. Convergence in total variation on finite sets of bonds, by contrast, implies that both sides of the identity converge.

\emph{Step \textup{(c)} Invariances: identities in the limit.} Each of the following is an identity or an inequality between integrals of bounded local Borel functions: the invariance under gauge transformations of finite support; the invariance at depth $s$ under the maps of \textup{(R0)}, which comes from \Cref{Y-P2-L2} and the invariance of $\mu_{K,N,\beta}$; and \textup{(b)} for local $F$. All of them are valid for $(K,N)$ with $N$ large, and they remain valid for the limits in total variation, taken first in $K$ and then in $N$.

\emph{Step (Other sequences).} For $N_j\to\infty$ the argument requires of $(K_j)$ only that $K_j\ge s+1$ from some $j$ on, that is $s\le\liminf K_j-1$. It provides a subsequence of $(K_j,N_j)$, and along this subsequence the conclusions of \textup{(c)} hold for these depths $s$. At finite $K$ the consistency is the definition of the run (the definition ``Runs'' in item~2 of \S\ref{Y-J-1}), which then takes the place of \Cref{Y-J-CorC}\,\textup{(c)}.

\emph{Step ($\mathsf b<x_s$): small labels.} By \Cref{Y-P2-L11}, the run with the data $(K,l,\mathsf b)$ is the run with the data $(K-z,L^zl,L^z\mathsf b)$. In the units of the data $(K,l,\mathsf b)$, which are the units before the dilation, the depth $s$ of this run is the lattice $L^{-(s+z)}\mathbb Z^3$. \end{proof}

\section{Growth on flat test functions}\label{Y-P3-1}

In this section we bound the numbers $j^\circ,r^\circ,N_1^{[j]},C_W,e_h,B_0(\mathsf b,\check\delta)$ of \S\ref{Y-P2-1} in closed form by powers of $1/\check\delta$ on $0<\check\delta\le1$; this is clause \textup{(G1)} of \Cref{Y-P3-ThG}, whose clauses are \textup{(G1)} to \textup{(G6+)}. A chessboard estimate, carried out cell by cell on the unit board together with \Cref{Y-P2-ThR}\,\textup{(R0)} and \textup{(R1)}, gives bounds for moments that contain no factorial, \textup{(G3)}; from these we obtain the product form $A_5^n(n!)^5\max\{1,d^{-5n}\}(1+R_f)^{3n}$, which is \textup{(G2)}. This bound holds uniformly in the cutoff $K$, in the positions of the supports and in the length $N$ of the sides of the torus; by a dilation it extends from $\mathsf b\ge x_s$ to all $\mathsf b>0$, as stated in \textup{(G2)}\,\textup{(iv)}. We then pass to $\mathcal S_{\ne}(\mathbb R^{3n})$, where the growth condition of \Cref{support:26-33} holds with $(13,1,C_E(\mathsf b),13)$, \textup{(G4)}, together with reality, symmetry, invariance and positivity, \textup{(G5)}; and a single field set between two vectors obeys a bound free of $n$, \textup{(G6)} and \textup{(G6+)}.

\emph{0.1. Inputs and conventions.} \textup{(A)} We use the table, radii and numbers of item~1 of \S\ref{Y-H9-1}, and the lattice data of \Cref{Y-H9-T}. In particular, $\kappa\ge6$, $\kappa_1=\kappa+1$, $M=L^{m'}$ with $m'\ge1$, $M\ge1920(\kappa_1+4)c_R$, and $L^{1/4}\ge10^4\mathsf N_{\rm c}^2$. The floors give $x_s\ge e^2$, $Lx_s\ge x_9^U$, $p_0\ge1$ and $\varepsilon\le1/40$ on $x\ge x_s$. We take $\mathrm{sep}^{[i]}$ and $N_1^{[j]}$ at $d=\check\delta$ and $R_f=\check\rho$; item~1 below recalls the numbers needed here.

The separated tuples, their centred functionals and reflection positivity are those of item~1 of \S\ref{Y-H9-1}. We use its convention $rf=f\circ r^{-1}$, with $\tau_v:x\mapsto x+v$, and $\|f\|_\sharp=\max\{\|f\|_\infty,40M\operatorname{Lip}f\}$. A Euclidean type $(d;R)$ is a sup type $(d/\sqrt3;R)$, as in \textup{(S3)} of \S\ref{Y-00-1}.

By \Cref{Y-H9-LN}\,\textup{(N2)}, $\mathrm{sep}^{[i+1]}\ge L\,\mathrm{sep}^{[i]}$ for $i\ge0$ and $d>0$. Its proof also gives $\mathrm{sep}^{[i]}\to\infty$: $dL^i/M$ grows geometrically, whereas $r(L^i\mathsf b)$ grows at most linearly in $i$.

We use the bound $c_a(x)\le x^{-3/8}/4000$ on the whole range $x\ge Lx_s$. Indeed, \textup{(O13)}\,\textup{(ii)} of item~0 of \S\ref{Y-H1-1} imposes \textup{(Y.3)} with $\kappa_1$ and $e^{-(c_h/2)p_0^2}$ for every $x\ge x_9^U$. Its non-negative $c_a$ term and the factors at least $1$ give $x^{1/4}\cdot50\cdot5\cdot16c_a(x)\le x^{-1/8}$; now use $Lx_s\ge x_9^U$.

Step~5 of the proof of \Cref{Y-H9-T} supplies dilation for $l=2L^m$, $m\in\mathbb Z$, $K\ge1-m$, $\beta>0$ and $z\in\mathbb Z$: with $f^{[z]}(x)=f(L^{-z}x)$, the data $(K,l,\beta)$ and $(K-z,L^zl,\beta)$ describe the same graph and measure, and $\mathcal O_{K-z}(f^{[z]})=L^{-z}\mathcal O_K(f)$. Thus the cumulants scale by $L^{zn}$; the corresponding limits exist together and obey the same identity along matched subsequences, also at moving labels. Step~7 of that proof supplies the finite-cutoff covariance $\mathcal O_K(f)(rU)=\mathcal O_K(f\circ r)(U)$ and invariance of the measure for $r:x\mapsto r^\circ x+v$, $r^\circ\in W_3$, $v\in a_K\mathbb Z^3$.

\textup{(B)} We use the boards, mirrors, copies and folds of \Cref{support:4-1}\,(2), with the action $(rF)(U)=F(r^{-1}U)$, and the chessboard estimate of \Cref{support:4-13} and \Cref{support:4-15}. For flat test functions we use \Cref{support:26-2}, the norm \eqref{support:26-4}, $\mathsf C_{\ne}=10^{53}$ of \eqref{support:26-10}, \Cref{support:26-23}, \Cref{support:26-24}, \Cref{support:26-33} and \Cref{support:26-41}\,(a)--(c). The table is admissible for these results: \Cref{support:26-1} admits odd $L\ge13$ and $M=L^{m'}$ with $m'\ge1$.

The conventions for $d(X)$, $R(X)$ and $\mathbb C\mathfrak S(O)$ are those of \Cref{support:9-5}, with Euclidean distance. Defs.~\textup{\ref{support:25-28}}\,(4), \textup{\ref{support:9-12}}\,(2), (3) and \textup{\ref{support:9-25}}\,(1), (3) are used as definitions only, at $\mathbb E=\mathbb R^3$; item~1 of \S\ref{Y-H9-1} gives the corresponding notions at $T_l$.

\textup{(C)} From \S\ref{Y-P2-1} we take the locally bounded numbers of item~2 and \Cref{Y-P2-ThR}\,\textup{(R0)}, \textup{(R1)}, \textup{(R5)}--\textup{(R7)}. We use \textup{(R0)} at $s=0$, and \textup{(R1)} only at $K\ge j^\circ(\mathsf b,\check\delta)$, for real bounded $g_t$ whose nonzero sets are pairwise at sup distance $\ge\check\delta$ and lie each in a sup-ball of radius $\check\rho>0$. The limits of \textup{(R5)}, including moving labels in parts (a), (c), are taken on continuous entries with compact, pairwise disjoint supports. Their reality, symmetry, invariance, reflection positivity and dilation follow from \textup{(R6)}\,(i), (iii), (iv), (vi); \textup{(R7)} supplies $\mathcal H(S)$, its unit vacuum and the contraction semigroup $e^{-tH}[F]=[\tau_{te_3}F]$.

When all cells carry the same pair $(\check\delta,\check\rho)$, \textup{(G3)}\,(a) is \Cref{Y-P2-ThR}\,\textup{(R2)}\,(i) restricted to $K\ge j^\circ(\mathsf b,\check\delta)$. The range $K<j^\circ$ of \textup{(R2)} lies outside \textup{(G3)}.

\emph{0.2. Local letters and tags.}

\begin{remark}[Letters set here]

The letter $\theta$ is the map $x_3\mapsto-x_3$.

The letter $T_N$ is the torus.

The labels \textup{(E0$'$)} to \textup{(E4)} denote the Euclidean axioms of \Cref{support:28-35}.

In the first sentence of \textup{(G2)}\,\textup{(iv)} and in \textup{(G4)}\,\textup{(e)}, $z=z_0(\mathsf b)$, in the setting \textup{(S7)} of \S\ref{Y-00-1}.\end{remark}

\emph{1. Data and numbers.}

\begin{definition}[Lattice data, centred moments and the unit board] We recall the data of \Cref{Y-H9-T}. Let $G=\operatorname{SU}(2)$ and $d=3$. Let $\mathsf b\ge x_s$, and let $N=2L^{m_l}$ with $m_l\ge m'$ and \textup{(Tor)} $N\ge2Ms(L\mathsf b)$. Let $K\ge1$, $\beta=L^K\mathsf b$ and $a_K=L^{-K}$; we write $\mu_{K,N,\beta}$ for the measure, and $\mathcal O_K(f)=a_K^{-1}\sum_pf(x(p))a_p$.

We put $\sigma_{K,N}(f_1,\dots,f_n)=\langle\prod_i(\mathcal O_Kf_i-\langle \mathcal O_Kf_i\rangle)\rangle$. Distances are sup distances, except where they are called Euclidean.

The board $B_0$, written without arguments, is $(a_K,L^K,(N,N,N),0)$ in the notation $(a,\lambda,\underline{\mathsf n},o)$ of \Cref{support:4-1}\,\textup{(2)} for a board. Its cells are $\varkappa+[-1/2,1/2]^3$, and its walls are $y_\mu\in\mathbb Z+1/2$, through the midpoints of bonds. The number $B_0(\mathsf b,\check\delta)$ always carries its arguments. By $S$ we denote any limit of \textup{(R5)}.\end{definition}

\begin{definition}[The numbers of the section] In the order in which the constants are fixed, the table \textup{(O0)} to \textup{(O13)} comes first, then $\mathsf b$, and then \textup{(O14)}, to which all numbers below belong. No floor depends on them, and none of them depends on $n$, $N$, $K$ or on where the sets $\{f_i\ne0\}$ lie on $T_N$; all are locally bounded in $\mathsf b$, by \textup{(Lab1)}.

From item~2 of \S\ref{Y-P2-1} we recall the following numbers: $\mathrm{sep}^\circ:=8\log(8L^{1/2})$, the one numeral among them; $\mathrm{sep}^{[i]}=\lfloor\check\delta L^i/M\rfloor-2\check r(L^i\mathsf b)-2$; $j^\circ(\mathsf b,\check\delta)$, the least $j\ge1$ with $\mathrm{sep}^{[j-1]}\ge\mathrm{sep}^\circ$; $\mathfrak r^\circ(x)=3.5/\mathrm{amp}(Lx)$ and $r^\circ=\mathfrak r^\circ(L^{j^\circ}\mathsf b)$; $N_1^{[j]}=(2\check\rho L^{j-1}/M+2\check r(L^{j-1}\mathsf b)+3)^3$; \[C_W(\mathsf b,\check\delta,\check\rho)=\sum_{j>j^\circ}2eN_1^{[j]}C_gc_a(L^j\mathsf b)e^{-(\kappa-1)(\mathrm{sep}^{[j-1]}-1)}+\sum_{1\le j\le j^\circ}6N_1^{[j]}C_gc_a(L^j\mathsf b);\] $e_h=0.45N_1^{[1]}c_h(\mathsf b)$; $G_1=e^{C_W+e_h}/r^\circ$; $B_0(\mathsf b,\check\delta)=6L^{j^\circ}(\check\delta L^{j^\circ}/2+2)^3$; and $c_{\rm vol}(\mathsf b):=225/M^3+(3/2)\mathsf N_Zp_0(\mathsf b)^2+3\log(384/\varepsilon(\mathsf b)^3)$.

We also set $\mathfrak l(\mathsf b,\check\delta):=\mathrm{sep}^\circ+11+[\log \mathsf b+\log^+(M/(264\check\delta))]/528$, $\mathfrak l_1(\mathsf b):=\mathrm{sep}^\circ+11+\log(\mathsf bM)/528$ and $\check\lambda(x):=\mathrm{sep}^\circ+11+\log x/528$.

We set $A_2:=c_Q^{\rm su}K_5L^{3/2}(2M)^{1/2}/(3.5E_a)$ and $T_L:=\sum_{t\ge1}L^{3t}\exp\{-(\kappa-1)(L^t\mathrm{sep}^\circ-1)\}$.

We set \[C_W^*(\mathsf b):=C_gL^3\{6\sum_{j\ge1}\check\lambda(L^j\mathsf b)^3c_a(L^j\mathsf b)+2eT_L\sup_{j\ge1}\check\lambda(L^j\mathsf b)^3c_a(L^j\mathsf b)\}\] and $e_h^*(\mathsf b):=0.45(2\check r(\mathsf b)+4)^3c_h(\mathsf b)$.

We set $\check\Gamma(\mathsf b,\check\delta,\check\rho):=G_1(\mathsf b,\check\delta,\check\rho)(1+e^{c_{\rm vol}(\mathsf b)})$ and $\check\Gamma(\mathsf b,\check\delta):=\check\Gamma(\mathsf b,\check\delta,\check\delta/4)$.

We set $A_4(\mathsf b):=\mathfrak l_1^4\max\{A_2\mathsf b^{-1/2}e^{C_W^*+e_h^*}(1+e^{c_{\rm vol}})e^{c_{\rm vol}},\ 24L^2M(L^2M+2)^3\}$, $A_5:=64(4e)^5A_4$, $A_6:=3^{5/2}A_5$ and $C_E(\mathsf b):=\mathsf C_{\ne}^6M8A_6$.

We set $\check\Gamma_*(\mathsf b,\eta):=2A_2e^{C_W^*+e_h^*}(1+e^{c_{\rm vol}})e^{c_{\rm vol}}(\mathfrak l(\mathsf b,\eta/4)/(\mathsf b\eta))^{1/2}$.

For $\eta\in(0,1]$ let $p':=\lceil1/(2\eta)\rceil$ and $\eta':=1/(2p')\in[\eta/2,\eta]$; equality on the left holds at $\eta=1$ alone. For $y\in\mathbb R^3$ let $|\!|\!|h|\!|\!|_{\eta,y}:=\sum_{c\in y+(\eta'/8)\mathbb Z^3}\sup\{|h(x)|:|x-c|_\infty\le\eta'/8\}$, and $|\!|\!|h|\!|\!|_\eta:=|\!|\!|h|\!|\!|_{\eta,0}$.\end{definition}

\emph{Labels.} As functions of the label $\mathsf b$, the numbers obey the following.\begin{enumerate}

\item[\textup{(Lab1)}] Every number above is locally bounded in $\mathsf b$ on $[x_s,\infty)$, and so is $j^\circ(\mathsf b,\check\delta)$ at fixed $\check\delta$. The numbers $A_5(\mathsf b)$, $A_6(\mathsf b)$, $C_E(\mathsf b)$ of \textup{(G2)}\,\textup{(iv)} and \textup{(G4)}\,\textup{(e)} are locally bounded on $(0,\infty)$.

\item[\textup{(Lab2)}] Let $e_h^{*c}(\mathsf b):=0.45(\log \mathsf b/528+12)^3c_h(\mathsf b)$; then $e_h^*\le e_h^{*c}$. For $X$ one of $A_4$, $A_5$, $A_6$, $C_E$, $\check\Gamma_*(\cdot,\eta)$ and $C_n^\infty$ of \textup{(G2)}\,\textup{(iii)}, let $X^c$ be $X$ with $e_h^{*c}$ in the place of $e_h^*$. Then $X\le X^c$, and $X^c$ is continuous in $\mathsf b$ on $[x_s,\infty)$.

\item[\textup{(Lab3)}] Let $S$ be a limit of \textup{(R5)}\,\textup{(a)} or \textup{(c)} along labels $\mathsf b_j\to \mathsf b$. Then every number $X(\mathsf b)$ in \textup{(G2)}\,\textup{(iii)}, \textup{(G4)}\,\textup{(a)} to \textup{(c)}, \textup{(G6)} and \textup{(G6+)} means $\limsup_jX(\mathsf b_j)$; this limsup may be taken along any subsequence, the limit $S$ staying the same.

That number is at most $\sup X$ over any set that contains all $\mathsf b_j$ of the subsequence from some $j$ on, for instance over $U(\mathsf b)$ of the setting \textup{(S7)} of \S\ref{Y-00-1} when $\mathsf b>x_s$. If $\mathsf b_j\ge x_s$ for infinitely many $j$, then along these $j$ it is at most $X^c(\mathsf b)$, a number of $\mathsf b$ alone.

If $\mathsf b_j<x_s$ for infinitely many $j$, we take a single integer $z\ge1$ with $L^z\mathsf b_j\ge x_s$ from some $j$ on, and we use \textup{(G2)}\,\textup{(iv)} at that $z$: the number is then at most $L^{4z}A_5^c(L^z\mathsf b)$, and likewise for $A_6$ and $C_E$.

As to \textup{(G6)}, its main clause and its first three bullets are asserted only when $\mathsf b_j\ge x_s$ for infinitely many $j$, and along these $j$. If $\mathsf b_j<x_s$ from some $j$ on, \textup{(G6)} is asserted in its last bullet alone, at such a $z$, with a number at most $L^z\check\Gamma_*^c(L^z\mathsf b,L^z\eta_z)$.\end{enumerate}

\emph{2. The theorem.}

\begin{theorem}[Growth condition and bounds in closed form]\label{Y-P3-ThG} In every clause the objects, letters and numbers are those of item~1. The data as recalled in item~1 are $\mathsf b\ge x_s$, $N=2L^{m_l}$ with $m_l\ge m'$, \textup{(Tor)} and $K\ge1$. They are assumed in full in the clauses at finite $(K,N)$: \textup{(G1)}, \textup{(G2)}\,\textup{(i)} and \textup{(ii)}, \textup{(G3)} and \textup{(G4)}\,\textup{(d)}. The clauses on a limit $S$ with $\mathsf b\ge x_s$ are \textup{(G2)}\,\textup{(iii)}, \textup{(G4)}\,\textup{(a)} to \textup{(c)}, \textup{(G5)}, the main clause and the first three bullets of \textup{(G6)}, and \textup{(G6+)}\,\textup{(m)}, which states its own condition on the $\mathsf b_j$, the same as that of the restriction in \textup{(G6)}. What these clauses assume of the data is $\mathsf b\ge x_s$ for the label $\mathsf b$ of the limit, while $N$, \textup{(Tor)} and $K$ concern the clauses at finite $(K,N)$. The labels $\mathsf b_j$ of the members are subject only to the conditions of \textup{(Lab3)} and of the restriction in \textup{(G6)}: $\mathsf b_j<x_s$ is admitted in \textup{(G2)}\,\textup{(iii)}, \textup{(G4)}\,\textup{(a)} to \textup{(c)} and \textup{(G5)}, with the numbers taken in the sense of \textup{(Lab3)}. In \textup{(G2)}\,\textup{(iv)} and \textup{(G4)}\,\textup{(e)}, the label, the sides and $K$ are those that the clause itself states, in the place of those of item~1. The last bullet of \textup{(G6)} and \textup{(G6+)}\,\textup{(n)} restate the label alone. The following hold.\begin{enumerate}

\item[\textup{(G1)}] \emph{The numbers.} For $\mathsf b\ge x_s$ and $0<\check\delta\le1$ the following lines hold.\begin{enumerate}

\item[\textup{0.}] $1+0.182\,C_{\rm vol}(N,\mathsf b)\le e^{c_{\rm vol}(\mathsf b)N^3}$ for every $N$ of the data; this line involves no $\check\delta$.

\item[\textup{1.}] $j^\circ$ exists and is non-increasing in $\check\delta$; $L^{j^\circ}\le2L^2(M/\check\delta)\mathfrak l$; and $a_K\le\check\delta/(LM(\mathrm{sep}^\circ+8))$ for $K\ge j^\circ$.

\item[\textup{2.}] $1/r^\circ\le A_2(\mathfrak l/(\mathsf b\check\delta))^{1/2}$.

\item[\textup{3.}] At $\check\rho=\check\delta/4$, $N_1^{[j]}\le(L\,\check\lambda(L^j\mathsf b))^3$ for $1\le j\le j^\circ$, and $N_1^{[j^\circ+t]}\le(L^{t+1}\check\lambda(L^{j^\circ+t}\mathsf b))^3$ for $t\ge1$.

\item[\textup{4.}] $C_W(\mathsf b,\check\delta,\check\rho)\le C_W(\mathsf b,\check\delta,\check\delta/4)\le C_W^*(\mathsf b)$ for $\check\rho\le\check\delta/4$, and $C_W^*(\mathsf b)\le0.27L^6((\mathrm{sep}^\circ+11)^3+1)(12+2eT_L)(L\mathsf b)^{-1/4}$; so $C_W^*(\mathsf b)$ is finite, free of $\check\delta$, and bounded on $[x_s,\infty)$.

\item[\textup{5.}] $e_h\le e_h^*(\mathsf b)$ for $\check\rho\le\check\delta/4$.

\item[\textup{6.}] $|\mathcal O_K(g)|\le B_0(\mathsf b,\check\delta)\|g\|_\infty$ for $K\le j^\circ$ whenever $\{g\ne0\}$ lies in a sup-ball of radius $\check\delta/4$; and $B_0(\mathsf b,\check\delta)\le12L^2M(L^2M+2)^3\mathfrak l^4/\check\delta$.

\item[\textup{7.}] $\mathfrak l\le\mathfrak l_1(1+\log(1/\check\delta)/528)$, and $\mathfrak l^4\le\mathfrak l_1^4/\check\delta$.

\item[\textup{8.}] $\max\{\check\Gamma(\mathsf b,\check\delta)e^{c_{\rm vol}},2B_0(\mathsf b,\check\delta)\}\le A_4(\mathsf b)\check\delta^{-2}$.\end{enumerate}

\item[\textup{(G2)}] \emph{Product form.} Let $K\ge1$ be arbitrary, let $N$ be any side of the data, $n\ge2$, $d>0$ and $R_f>0$, and let the $f_i$ be real bounded functions on $T_N$ whose sets $\{f_i\ne0\}$ are pairwise at distance $\ge d$ and lie each in a sup-ball of radius $R_f$, wherever these balls lie.\begin{enumerate}

\item[\textup{(i)}] $|\sigma_{K,N}(f)|\le A_5(\mathsf b)^n(n!)^5\max\{1,d^{-5n}\}(1+R_f)^{3n}\prod\|f_i\|_\infty$.

\item[\textup{(ii)}] $|S^T_{n;K,N,L^K\mathsf b}(f)|$ is at most $1.45^{n+1}n!$ times the right side of \textup{(i)}.

\item[\textup{(iii)}] Every limit $S$ of \textup{(R5)} obeys \textup{(i)} and \textup{(ii)} on continuous entries, at moving labels in the sense of \textup{(Lab3)}; for tuples of the Euclidean type $(d;R)$ it obeys them with $A_6$ in the place of $A_5$. The number $C_n^\infty(d;R):=1.45^{n+1}n!\,A_6^n(n!)^5\max\{1,d^{-5n}\}(1+R)^{3n}$ is non-increasing in $d$ and non-decreasing in $R$.

\item[\textup{(iv)}] For $0<\mathsf b<x_s$ and $z=z_0(\mathsf b)$, in the setting \textup{(S7)} of \S\ref{Y-00-1}, \textup{(i)} to \textup{(iii)} hold with $A_5(\mathsf b):=L^{4z}A_5(L^z\mathsf b)$, for $K\ge z+1$ and for the sides $l=2L^m$ with $m+z\ge m'$ and $L^zl\ge2Ms(L^{z+1}\mathsf b)$. The same holds for every $\mathsf b>0$ with any integer $z\ge0$ such that $L^z\mathsf b\ge x_s$, and with the number $L^{4z}A_5(L^z\mathsf b)$.\end{enumerate}

\item[\textup{(G3)}] \emph{Cell by cell on $B_0$, free of factorials.} Let $g_1,\dots,g_Q$ be real bounded functions on $T_N$, each set $\{g_i\ne0\}$ lying in one cell. For each occupied cell $t$ let $\check\delta_t\in(0,1]$ and $\check\rho_t>0$ be such that the $n_t$ sources of $t$ are pairwise at distance $\ge\check\delta_t$, at distance $\ge\check\delta_t/2$ from all walls, and each in a sup-ball of radius $\check\rho_t$. Let $K\ge\max_tj^\circ(\mathsf b,\check\delta_t)$.\begin{enumerate}

\item[\textup{(a)}] $|\langle\prod_i(\mathcal O_Kg_i-\hat c_Kg_i)\rangle|\le\prod_t[G_1(\mathsf b,\check\delta_t,\check\rho_t)^{n_t}e^{c_{\rm vol}(\mathsf b)}]\prod\|g_i\|_\infty$; at $Q=1$ this reads $|\langle \mathcal O_Kg\rangle-\hat c_Kg|\le G_1e^{c_{\rm vol}}\|g\|_\infty$.

\item[\textup{(b)}] $|\sigma_{K,N}(g_1,\dots,g_Q)|\le\prod_t[\check\Gamma(\mathsf b,\check\delta_t,\check\rho_t)^{n_t}e^{c_{\rm vol}(\mathsf b)}]\prod\|g_i\|_\infty$.\end{enumerate}

If $\check\rho_t\le\check\delta_t/4$, the condition that a source lies in one cell follows from its distance to the walls.

\item[\textup{(G4)}] \emph{The growth condition.}\begin{enumerate}

\item[\textup{(a)}] $\sigma_S$ satisfies the growth condition of \Cref{support:26-33} with $(k_{\rm E},\mathsf c_{\rm E},\mathsf C_{\rm E},\varkappa_{\rm E})=(13,1,C_E(\mathsf b),13)$: for every $n$ there is exactly one complex linear map $S^{\sigma_S}_n$ on $\mathcal S_{\ne}(\mathbb R^{3n})$ which is bounded for some $|\cdot|_k$ and takes the value $\sigma_S(X)$ at the tensors of $\mathfrak S_n^\infty(\mathbb R^3)$, and $|S^{\sigma_S}_n(\varphi)|\le C_E(\mathsf b)^n(n!)^{13}|\varphi|_{13n-1}$. Moreover $S_1=0$.

\item[\textup{(b)}] Independently of \Cref{support:26-24}, the sharper bound \[\FitDisplay{\dimexpr\linewidth-1pt\relax}{|S^{\sigma_S}_n(\varphi)|\le(10^{36}A_5(\mathsf b))^n(n!)^5|\varphi|_{12n+1}}\] holds. Since $10^{36}A_5\le C_E$ and $12n+1\le13n$, \textup{(b)} gives the quadruple of \textup{(a)} in the form \eqref{support:26-34}.

\item[\textup{(c)}] The extension $\tilde S_n$ of the truncated $S_n$ obeys \textup{(a)} and \textup{(b)} with the further factor $1.45^{n+1}n!$.

\item[\textup{(d)}] At finite $(K,N)$, for $\varphi$ smooth on $\mathbb R^{3n}$, flat on the diagonals and carried by $(-N/4+1,N/4-1)^{3n}$, \[|a_K^{-n}\sum_{p_1,\dots,p_n}\varphi(x(p_1),\dots,x(p_n))\langle\prod_i(a_{p_i}-\langle a_{p_i}\rangle)\rangle|\le(10^{36}A_5)^n(n!)^5|\varphi|_{12n+1},\] a bound free of $K$ and $N$.

\item[\textup{(e)}] For $0<\mathsf b<x_s$, \textup{(a)} to \textup{(d)} and \textup{(G5)} hold with $A_5(\mathsf b)$ of \textup{(G2)}\,\textup{(iv)}, with $A_6(\mathsf b):=3^{5/2}A_5(\mathsf b)$ and with $C_E(\mathsf b):=\mathsf C_{\ne}^6M8A_6(\mathsf b)$; in \textup{(d)} they hold for $K\ge z+1$, with the sides $l$ of \textup{(G2)}\,\textup{(iv)} in the place of $N$.\end{enumerate}

\item[\textup{(G5)}] On $\mathcal S_{\ne}(\mathbb R^{3n})$, with the axioms named as in \Cref{support:28-35}, the functionals $S^{\sigma_S}_n$ of \textup{(G4)} are real, have the symmetry \textup{(E3)}, and are invariant under all translations and under $W_3$, which is \textup{(E1)} for that group; and \textup{(E2)} holds in the form \eqref{support:26-42} at every coordinate mirror, for Schwartz functions flat outside $V_+^{(n)}$.

\item[\textup{(G6)}] \emph{A field between two vectors, free of $n$.} For $\mathsf b\ge x_s$, $\eta\in(0,1]$, $F,F'\in\mathbb C\mathfrak S(\{x_3>\eta\})$, $h$ real in $C^1_c(\{|x_3|<\eta/4\})$ and every $y\in\mathbb R^3$, \[|\sigma_S(\theta\overline F\cdot(h)\cdot F')|\le\check\Gamma_*(\mathsf b,\eta)\,|\!|\!|h|\!|\!|_{\eta,y}\,\|[\tau_{-\eta e_3}F]\|\,\|[\tau_{-\eta e_3}F']\|.\] At $y=0$ the clause holds with $|\!|\!|h|\!|\!|_\eta$.

\emph{Restriction.} At moving labels, the main clause and the first three bullets below hold when $\mathsf b_j\ge x_s$ for infinitely many $j$, with $\check\Gamma_*$ in the sense of \textup{(Lab3)} along these $j$. When $\mathsf b>x_s$, every sequence $\mathsf b_j\to \mathsf b$ is covered. When $\mathsf b=x_s$ and $\mathsf b_j<x_s$ from some $j$ on, the last bullet holds, and so does \textup{(G6+)}\,\textup{(n)}.\begin{enumerate}

\item[$\bullet$] $\check\Gamma_*(\mathsf b,\eta)\le4A_2e^{C_W^*+e_h^*}(1+e^{c_{\rm vol}})e^{c_{\rm vol}}(\mathfrak l_1/\mathsf b)^{1/2}/\eta$. For $h\in C^3_c$ and every $y$, $|\!|\!|h|\!|\!|_{\eta,y}\le4096\,\eta^{-3}\sum_{\alpha\in\{0,1\}^3}\eta^{|\alpha|}\|\partial^\alpha h\|_{L^1}$.

\item[$\bullet$] In time order: let the $h_i$ be real in $C^1_c$ with $\operatorname{supp}h_i\subset\{|x_3-t_i|<\eta_i/4\}$ and $t_{i+1}-t_i\ge\eta_i+\eta_{i+1}$; then $|\sigma_S(h_1,\dots,h_n)|\le\prod_i\check\Gamma_*(\mathsf b,\eta_i)|\!|\!|h_i|\!|\!|_{\eta_i}$.

\item[$\bullet$] The same holds in each coordinate direction.

\item[$\bullet$] Let $\mathsf b>0$, and let $z\ge1$ be an integer with $L^z\mathsf b_j\ge x_s$ from some $j$ on (at a fixed label: $L^z\mathsf b\ge x_s$). Let $\eta_z:=\min\{\eta,L^{-z}\}$, let $h$ be real in $C^1_c(\{|x_3|<\eta_z/4\})$, and let $F,F'$ be as above. Then the left side is at most $L^z\check\Gamma_*(L^z\mathsf b,L^z\eta_z)\,|\!|\!|h^{[z]}|\!|\!|_{L^z\eta_z,y}\,\|[\tau_{-\eta e_3}F]\|\,\|[\tau_{-\eta e_3}F']\|$, with $\check\Gamma_*(L^z\mathsf b,\cdot)$ in the sense of \textup{(Lab3)} along the labels $L^z\mathsf b_j$. Such a $z$ exists for every sequence $\mathsf b_j\to \mathsf b>0$.\end{enumerate}

\item[\textup{(G6+)}] \emph{The whole slab.}\begin{enumerate}

\item[\textup{(m)}] Let the label be fixed with $\mathsf b\ge x_s$, or let $\mathsf b_j\ge x_s$ for infinitely many $j$, $\check\Gamma_*$ then being taken in the sense of \textup{(Lab3)} along these $j$. For any $\eta>0$ put $\hat\eta:=\min\{1,3\eta/4\}$. Then for $F,F'\in\mathbb C\mathfrak S(\{x_3>\eta\})$, $h$ real in $C^1_c(\{|x_3|<\eta/4\})$ and every $y$, \[|\sigma_S(\theta\overline F\cdot(h)\cdot F')|\le\check\Gamma_*(\mathsf b,\hat\eta)\,|\!|\!|h|\!|\!|_{\hat\eta,y}\,\|[\tau_{-\eta e_3}F]\|\,\|[\tau_{-\eta e_3}F']\|.\]

\item[\textup{(n)}] Let $\mathsf b>0$, let $z\ge1$ be an integer with $L^z\mathsf b_j\ge x_s$ from some $j$ on, let $\eta\in(0,1]$, and put $e_z:=\min\{1,3L^z\eta/4\}$. Then for $F,F'\in\mathbb C\mathfrak S(\{x_3>\eta\})$, $h$ real in $C^1_c(\{|x_3|<\eta/4\})$ and every $y\in\mathbb R^3$, $|\sigma_S(\theta\overline F\cdot(h)\cdot F')|$ is at most $L^z\check\Gamma_*(L^z\mathsf b,e_z)\,|\!|\!|h^{[z]}|\!|\!|_{e_z,y}\,\|[\tau_{-\eta e_3}F]\|\,\|[\tau_{-\eta e_3}F']\|$. Here the slab $\{|x_3|<\eta/4\}$ is the whole slab, not the thin slab $\{|x_3|<\eta_z/4\}$ of the last bullet of \textup{(G6)}. In the case $\mathsf b=x_s$, $\mathsf b_j<x_s$ from some $j$ on, of the restriction in \textup{(G6)}, this holds with $z=1$.\end{enumerate}\end{enumerate}\end{theorem}


The useful dependence on separation comes from summing the decaying
analytic bounds over the late scales before exponentiating. The
resulting powers of the separation are combined with a dyadic
partition off the diagonals. Flatness of the tests makes the resulting
Schwartz bound summable over that partition.

\emph{Step 1: clause \textup{(G1)}, with \textup{(Lab1)} and \textup{(Lab2)}.}

\begin{lemma}[Closed forms of the numbers, and their dependence on the label]\label{Y-P3-L1} In the setting of item~1, for $\mathsf b\ge x_s$ and $0<\check\delta\le1$, the nine lines~0 to~8 of \textup{(G1)} hold, and so do \textup{(Lab1)} and \textup{(Lab2)}.\end{lemma}

\begin{proof}\emph{Step (0): radii.} By \textup{(M2)} and $\kappa_1+4\ge11$ we have $20c_R/M\le1/1056$. Hence $r(x)\le\log x/1056+1$, and as $\check r=r+3$ we get $2\check r(x)+2\le\log x/528+10$, then $2\check r(x)+3\le\log x/528+11\le\check\lambda(x)$, and $2\check r(x)+4\le\log x/528+12$.

\emph{Step (i): line 1.} The argument for \textup{(N2)} in the proof of \Cref{Y-H9-LN}, as recalled in item~0.1~\textup{(A)}, gives $\mathrm{sep}^{[i]}\to\infty$ at $d=\check\delta$, and by \textup{(N2)} the sequence $\mathrm{sep}^{[i]}$ grows by the factor $L$ from its first positive value on. Hence $j^\circ$ exists, and $\mathrm{sep}^{[j-1]}<\mathrm{sep}^\circ$ for $1\le j<j^\circ$. The number $\mathrm{sep}^{[i]}$ is non-decreasing in $\check\delta$; so $j^\circ$ is non-increasing in $\check\delta$.

Let $j^\circ\ge2$, and put $Y=L^{j^\circ-2}$ and $t=\check\delta Y/M$. The inequality $\mathrm{sep}^{[j^\circ-2]}<\mathrm{sep}^\circ$ reads $\lfloor t\rfloor<\mathrm{sep}^\circ+2\check r(Y\mathsf b)+2$, and with step \textup{(0)} it gives $t<\lfloor t\rfloor+1<\mathrm{sep}^\circ+11+\log(Y\mathsf b)/528$. Here $\log(Y\mathsf b)=\log \mathsf b+\log(M/(264\check\delta))+\log(264t)$, and $\log(264t)<264t$. Since $\log\le\log^+$, we obtain $t<\mathfrak l+t/2$, hence $t<2\mathfrak l$ and $L^{j^\circ}=L^2tM/\check\delta<2L^2(M/\check\delta)\mathfrak l$. If $j^\circ=1$, then $L\le2L^2(M/\check\delta)\mathfrak l$ holds as $\check\delta\le1$ and $\mathfrak l\ge11$.

For the last clause of line~1, $\mathrm{sep}^{[j^\circ-1]}\ge\mathrm{sep}^\circ$ gives $\check\delta L^{j^\circ-1}/M\ge\mathrm{sep}^\circ+2\check r+2\ge\mathrm{sep}^\circ+8$, so $a_K\le L^{-j^\circ}\le\check\delta/(LM(\mathrm{sep}^\circ+8))$ for $K\ge j^\circ$.

\emph{Step (ii): line 2.} By the definitions and \textup{(X14)}, \[1/r^\circ=\mathrm{amp}(L^{j^\circ+1}\mathsf b)/3.5=c_Q^{\rm su}K_5(L^{j^\circ+1}\mathsf b)^{1/2}/(3.5\,\mathsf bE_ap_0)\le[c_Q^{\rm su}K_5L^{1/2}/(3.5E_a)](L^{j^\circ}/\mathsf b)^{1/2}.\] We insert the bound of step \textup{(i)} for $L^{j^\circ}$; by the definition of $A_2$ this is line~2.

\emph{Step (iii): line 3, at $\check\rho=\check\delta/4$.} At $\check\rho=\check\delta/4$ we have $N_1^{[j]}=(\check\delta L^{j-1}/(2M)+2\check r(L^{j-1}\mathsf b)+3)^3$. Let $2\le j\le j^\circ$. Then $\mathrm{sep}^{[j-2]}<\mathrm{sep}^\circ$ gives $\check\delta L^{j-2}/M<\check\lambda(L^{j-2}\mathsf b)\le\check\lambda(L^j\mathsf b)$, hence $\check\delta L^{j-1}/(2M)<(L/2)\check\lambda(L^j\mathsf b)$, and with step \textup{(0)} we find $N_1^{1/3}\le(L/2+1)\check\lambda\le L\,\check\lambda$. For $j=1$ the bound requires $1/2+\log \mathsf b/528+11\le\check\lambda(L\mathsf b)$, which says $1/2\le\mathrm{sep}^\circ+\log L/528$ and holds by the definition of $\mathrm{sep}^\circ$. For $j=j^\circ+t$ we have $\check\delta L^{j-1}/(2M)=L^t\check\delta L^{j^\circ-1}/(2M)\le L^t(L/2)\check\lambda(L^j\mathsf b)$, where for $j^\circ=1$ we use $\check\delta/(2M)\le1/2$; so $N_1^{1/3}\le(L^{t+1}/2+1)\check\lambda\le L^{t+1}\check\lambda(L^j\mathsf b)$.

\emph{Step (iv): line 4.} Since $N_1$ is non-decreasing in $\check\rho$, it suffices to treat $\check\rho=\check\delta/4$, where line~3 applies. For the late levels, \[\sum_{j\le j^\circ}6N_1C_gc_a\le6C_gL^3\sum_{j\ge1}\check\lambda(L^j\mathsf b)^3c_a(L^j\mathsf b);\] this is the whole series, and it does not depend on $j^\circ$. For the early levels $j=j^\circ+t$, \textup{(N2)} gives $\mathrm{sep}^{[j-1]}\ge L^t\mathrm{sep}^\circ$, so the term is at most $2eC_gL^3L^{3t}e^{-(\kappa-1)(L^t\mathrm{sep}^\circ-1)}\check\lambda^3c_a$, and the sum of these terms is at most $2eC_gL^3T_L\sup\check\lambda^3c_a$. Hence $C_W(\mathsf b,\check\delta,\check\rho)\le C_W(\mathsf b,\check\delta,\check\delta/4)\le C_W^*(\mathsf b)$ for $\check\rho\le\check\delta/4$.

The number $C_W^*(\mathsf b)$ is finite. For $x=L^j\mathsf b\ge Lx_s$, the bound recalled from Step~0\,\textup{(b)} of the proof of \Cref{Y-H6-J} gives $c_a\le x^{-3/8}/4000$. We also have $\check\lambda^3\le4((\mathrm{sep}^\circ+11)^3+(\log x/528)^3)$, and $(\log x)^3\le(24/e)^3x^{1/8}\le688.3\,x^{1/8}$ with $688.3/528^3<1$. Hence $\check\lambda^3c_a\le((\mathrm{sep}^\circ+11)^3+1)x^{-1/4}/1000$. By \textup{(L3su)}, $\sum_j(L^j\mathsf b)^{-1/4}\le2(L\mathsf b)^{-1/4}$. With $C_g=263L^3$ this yields $C_W^*\le0.263L^6((\mathrm{sep}^\circ+11)^3+1)(12+2eT_L)(L\mathsf b)^{-1/4}$, which gives line~4, since $T_L<\infty$ by the double exponential decay of its terms.

We sum the whole series for the following reason. With $c_a\le1/32$ alone, each of the $j^\circ\sim\log(1/\check\delta)$ late terms would be of the size $\check\lambda^3\sim\log^3$, and $G_1$ would contain $\exp\{c\log^4(1/\check\delta)\}$, which is not a power of $\check\delta$.

\emph{Step (v): line 5.} For $\check\rho\le\check\delta/4$ we have $N_1^{[1]}\le(\check\delta/(2M)+2\check r(\mathsf b)+3)^3\le(2\check r(\mathsf b)+4)^3$, so $e_h\le e_h^*(\mathsf b)$.

\emph{Step (vi): line 6.} We have $0\le a_p\le2$. The barycentres of the plaquettes of one orientation form a translate of $a_K\mathbb Z^3$, so a closed cube of side $\check\delta/2$ contains at most $(\check\delta L^K/2+1)^3$ of them; and there are three orientations. A sup-ball of radius $\check\delta/4$ is such a cube, hence $|\mathcal O_Kg|\le L^K\cdot2\cdot3(\check\delta L^K/2+1)^3\|g\|\le B_0(\mathsf b,\check\delta)\|g\|$ for $K\le j^\circ$. With step \textup{(i)}, $B_0(\mathsf b,\check\delta)\le12L^2M(\mathfrak l/\check\delta)(L^2M\mathfrak l+2)^3\le12L^2M(L^2M+2)^3\mathfrak l^4/\check\delta$, as $\mathfrak l\ge1$.

\emph{Step (vii): line 7.} For $\check\delta\le1$, $\log^+(M/(264\check\delta))\le\log M+\log(1/\check\delta)$. Hence, with $u=\log(1/\check\delta)$, $\mathfrak l\le\mathfrak l_1+u/528\le\mathfrak l_1(1+u/528)$; and $(1+u/528)^4\le e^u$, which gives $\mathfrak l^4\le\mathfrak l_1^4/\check\delta$.

\emph{Step (viii): line 8.} By lines~6 and~7, $2B_0(\mathsf b,\check\delta)\le24L^2M(L^2M+2)^3\mathfrak l_1^4\check\delta^{-2}$. By lines~2, 4 and~5, $\check\Gamma\,e^{c_{\rm vol}}\le A_2\mathsf b^{-1/2}e^{C_W^*+e_h^*}(1+e^{c_{\rm vol}})e^{c_{\rm vol}}(\mathfrak l/\check\delta)^{1/2}$, and by step \textup{(vii)}, $(\mathfrak l/\check\delta)^{1/2}\le\mathfrak l_1^{1/2}\check\delta^{-5/8}\le\mathfrak l_1^4\check\delta^{-2}$. Line~8 follows by the definition of $A_4$.

\emph{Step (ix): line 0, the volume factor.} We write $C_{\rm vol}=4qe^{220q}E$, with $q=(N/M)^3\ge8$, since $N/M=2L^{m_l-m'}\ge2$, and with $E=e^{(3/2)\mathsf N_Zp_0^2N^3}(384/\varepsilon^3)^{3N^3}\ge1$, since $\varepsilon\le1/40$. Then $1+0.182\,C_{\rm vol}\le(1+0.728\,qe^{220q})E$. On $q\ge8$ we have $\log(1+0.728\,qe^{220q})\le220.23\,q$: the quotient is largest at $q=8$, where its value is at most $220.2203$. So this logarithm is at most $225N^3/M^3$, and line~0 follows from the definition of $c_{\rm vol}(\mathsf b)$.

\emph{Step (x): the labels.} For \textup{(Lab1)}: by \textup{(O8)} of the table, $p_0$, $\varepsilon$, $c_a$, $c_h$ are explicit positive functions, continuous on $x>1$; so $c_{\rm vol}$, $\mathfrak l$, $\mathfrak l_1$ and $\check\lambda$ are continuous in $\mathsf b$. The radius $\check r$ is a step function of $\log \mathsf b$, bounded on compact sets; $C_W^*$ is bounded by line~4; and $A_2$, $T_L$ do not depend on $\mathsf b$. The numbers $A_4$, $A_5$, $A_6$, $C_E$, $\check\Gamma_*$ are continuous expressions in these numbers, and $j^\circ(\mathsf b,\check\delta)$ is locally bounded by line~1, as $\mathfrak l$ is continuous. For $G_1$, $\check\Gamma(\mathsf b,\check\delta,\check\rho)$ and $B_0(\mathsf b,\check\delta)$ we use item~2 of \S\ref{Y-P2-1}. On $(0,x_s)$ the least $z$ is a step function of $\mathsf b$, bounded on compact sets.

For \textup{(Lab2)}: by step \textup{(0)}, $(2\check r(\mathsf b)+4)^3\le(\log \mathsf b/528+12)^3$, so $e_h^*\le e_h^{*c}$, and $e_h^{*c}$ is continuous. As to $C_W^*$, the terms $u_j(\mathsf b)=\check\lambda(L^j\mathsf b)^3c_a(L^j\mathsf b)$ are continuous and, by step \textup{(iv)}, at most $((\mathrm{sep}^\circ+11)^3+1)(L^jx_s)^{-1/4}/1000$ on $[x_s,\infty)$. So the series converges uniformly, and $\sup_ju_j$ is the uniform limit of $\max_{j\le J}u_j$; both are continuous. Each $X$ of \textup{(Lab2)} is a continuous expression which does not decrease as $e_h^*$ grows; hence $X\le X^c$, and $X^c$ is continuous.\end{proof}

\emph{Step 2: the cut, the placing on $B_0$, and the copies.}

\begin{lemma}[Cutting, placing and copying sources]\label{Y-P3-L2} In the setting of item~1 the following hold.\begin{enumerate}

\item[\textup{(a)}] Let $f$ be real and bounded on $T_N$, with $\{f\ne0\}$ in a sup-ball of radius $R_f$, and let $\check\delta>0$. Then $f$ is a sum of at most $(4R_f/\check\delta+2)^3$ pieces $g$ with $\|g\|\le\|f\|$ and $\{g\ne0\}\subset\{f\ne0\}$, each set $\{g\ne0\}$ lying in a sup-ball of radius $\check\delta/4$.

\item[\textup{(b)}] Let $n$ sets lie each in a sup-ball of radius $\check\delta/4$, where $\check\delta\le1/(4n)$, and let $a_K\le5/(8n)$. Then there is $v\in a_K\mathbb Z^3$ such that, after the translation by $v$, all of them keep a distance $\ge\check\delta/2$ from all walls $y_\mu\in\mathbb Z+1/2$.

\item[\textup{(c)}] Let the sources of one cell be pairwise at distance $\ge\check\delta$, and at distance $\ge\check\delta/2$ from the walls. Then their $N^3$ families of mirror copies are pairwise at distance $\ge\check\delta$.\end{enumerate}\end{lemma}

\begin{proof}\emph{Step (a): cutting.} The sup-ball is the image under the projection $\mathrm{pr}$ of a cube $C$ of $\mathbb R^3$ of side $\min\{2R_f,N\}$, half-open if the side is $N$, and $\mathrm{pr}$ is injective on $C$. We cut $\mathbb R^3$ into the half-open cubes $(\check\delta/2)(k+[0,1)^3)$. An interval of length $2R_f$ meets at most $4R_f/\check\delta+2$ of the intervals of length $\check\delta/2$, so at most $(4R_f/\check\delta+2)^3$ of these cubes meet $C$. For each of them we put $g:=f\,1_{\mathrm{pr}(C\cap\text{cube})}$; these pieces add up to $f$ and have the properties stated in \textup{(a)}.

\emph{Step (b): placing.} We fix $\mu$. The projection of the set $i$ to the coordinate $\mu$ lies in an interval of length $\check\delta/2$. The values of $v_\mu$ mod $1$ for which this interval, shifted by $v_\mu$, lies at distance $<\check\delta/2$ from $\mathbb Z+1/2$ form an open arc of length $\check\delta/2+\check\delta$. All these arcs together have length at most $1.5\,n\check\delta\le3/8$. Their complement is closed, consists of at most $n$ arcs, and has length $\ge5/8$; so one of its arcs is closed of length $\ge5/(8n)\ge a_K$, and it contains a point of $a_K\mathbb Z$ mod $1$, as $1/a_K$ is an integer. The same holds for $\mu=1,2,3$, and the three choices give $v$.

\emph{Step (c): copies.} Inside one cell the copies form an isometric image of the sources, so they have the same pairwise distances as the sources. Copies in two different cells differ in some coordinate $\mu$; in that coordinate the two supports lie in different unit intervals of $\mathbb R/N\mathbb Z$ cut out by $\mathbb Z+1/2$, each at distance $\ge\check\delta/2$ from every wall, and every path between them in that coordinate crosses a wall. So their distance is $\ge\check\delta$.

\emph{Step ``in one cell'': the last sentence of (G3).} Two points of a set which lie on the two sides of a wall, each at distance $\ge\check\delta/2$ from it, are at distance $\ge\check\delta$ from each other; but a sup-ball of radius $\check\delta/4$ has diameter $\check\delta/2$. So a source in such a ball, at distance $\ge\check\delta/2$ from the walls, lies in one cell.\end{proof}

\emph{Step 3: clause \textup{(G3)}, the chessboard estimate cell by cell, and centring at the true mean.}

\begin{lemma}[The chessboard estimate cell by cell, and centring at the true mean]\label{Y-P3-L3} In the setting of item~1, \textup{(G3)}\,\textup{(a)} and \textup{(b)} hold.\end{lemma}

\begin{proof}\emph{Step (1): cell functions.} For constants $k_i$ we put $F_t:=\prod_{i\in I_t}(\mathcal O_Kg_i-k_i)$, where $I_t$ indexes the $n_t$ sources of the cell $t$. Let $p$ be a plaquette with $g_i(x(p))\ne0$. Its four sites lie at sup distance $a_K/2$ from $x(p)$, the point $x(p)$ lies at distance $\ge\check\delta_t/2$ from the walls, and $a_K<\check\delta_t$ by line~1 of \textup{(G1)}. So the sites lie in the cell and off the walls, which contain no site, and the four bonds of $p$ are inner bonds. Hence $F_t$ is a real bounded continuous function of the cell $t$.

\emph{Step (2): the chessboard estimate.} We apply \Cref{support:4-15} on the board $B_0$, together with \Cref{support:4-13}. A cell of the board that holds no source receives the constant function $1$, and its fold is $\operatorname{Fold}1=1$. The chessboard estimate then reads \[|\langle\prod_tF_t\rangle|\le\prod_t\langle\operatorname{Fold}F_t\rangle^{1/N^3},\] where both products run over all cells $t$ of the board.

\emph{Step (3): the fold.} Let $\varkappa$ and $\varkappa'$ be cell indices, and let $\sigma=\sigma(\varkappa',\varkappa)$ be their signs. The mirror map $r=r_{\varkappa'\varkappa}:x\mapsto\varkappa'+\sigma(x-\varkappa)$ is of the form $r^\circ x+v$, where the linear part $r^\circ$ is a diagonal element of $W_3$ and $v\in\mathbb Z^3$. By Step~7 of the proof of \Cref{Y-H9-T}, in its paragraph on finite $K$, the copy of $\mathcal O_Kg_i-k_i$ under $r$ is $\mathcal O_K(g_i\circ r^{-1})-k_i$. Coordinate by coordinate, $r$ is either a translation by an even integer or a reflection $y\mapsto(\varkappa'_\mu+\varkappa_\mu)-y$, in which $\varkappa'_\mu+\varkappa_\mu$ is an odd integer. The parity of $\varkappa'-\varkappa$ is well defined, because $N$ is even.

\emph{Step (4): proof of \textup{(a)}.} In case \textup{(a)} of clause \textup{(G3)} of \Cref{Y-P3-ThG}, the constants are $k_i=\hat c_Kg_i$. By \textup{(R0)} of \Cref{Y-P2-ThR}, also $k_i=\hat c_K(g_i\circ r^{-1})$ for every mirror map $r$. Hence $\operatorname{Fold}F_t$ is a product of factors $(\mathcal O_Kg'-\hat c_Kg')$, one for each of $n_tN^3$ sources $g'$. Any two of these sources lie at distance at least $\check\delta_t$ from each other, by \Cref{Y-P3-L2}\,(c); each of them lies in a ball of radius $\check\rho_t$; and each has the sup norm of the source of which it is a copy. We apply \textup{(R1)} with $Q=n_tN^3$ and $K\ge j^\circ(\mathsf b,\check\delta_t)$ and obtain \[\langle\operatorname{Fold}F_t\rangle\le G_1^{n_tN^3}(1+0.182C_{\mathrm{vol}})\prod\|g_i\|_\infty^{N^3},\] where the product runs over $i\in I_t$. We take the $N^3$-th root and use line~0 of \textup{(G1)}; this bounds $\langle\operatorname{Fold}F_t\rangle^{1/N^3}$ by $G_1^{n_t}e^{c_{\mathrm{vol}}}\prod\|g_i\|_\infty$.

\emph{Step (5): proof of \textup{(b)}.} In case \textup{(b)} the constants are $k_i=\langle \mathcal O_Kg_i\rangle$; as the measure is invariant under the mirror maps, $k_i=\langle \mathcal O_Kg_i\rangle=\langle \mathcal O_K(g_i\circ r^{-1})\rangle$. For each copy $g'$ we split \[\mathcal O_Kg'-\langle \mathcal O_Kg'\rangle=(\mathcal O_Kg'-\hat c_Kg')-\mathfrak d(g'),\qquad\mathfrak d(g'):=\langle \mathcal O_Kg'\rangle-\hat c_Kg'.\] Case \textup{(a)}, applied with $Q=1$ to $g'$ as the only source of one cell with the same $\check\delta_t$ and $\check\rho_t$, gives $|\mathfrak d(g')|\le G_1(\mathsf b,\check\delta_t,\check\rho_t)e^{c_{\mathrm{vol}}}\|g'\|_\infty$. We expand the product over the $n_tN^3$ copies into a sum over the subsets $A$ of the copies: each copy in $A$ contributes the factor $\mathcal O_Kg'-\hat c_Kg'$, and each copy outside $A$ contributes the number $-\mathfrak d(g')$. To every sub-family $A$ we apply \textup{(R1)}, which admits every $Q\ge1$; for empty $A$ the moment is $1$. Summing over $A$, we find \[\langle\operatorname{Fold}F_t\rangle\le(1+0.182C_{\mathrm{vol}})\prod_{\text{copies}}(G_1+G_1e^{c_{\mathrm{vol}}})\|g'\|_\infty.\] The $N^3$-th root of this bound, taken with line~0 of \textup{(G1)}, is at most $\check\Gamma^{n_t}e^{c_{\mathrm{vol}}}\prod\|g_i\|_\infty$, with $\check\Gamma=\check\Gamma(\mathsf b,\check\delta_t,\check\rho_t)$.\end{proof}

\begin{remark}[On the proof of \Cref{Y-P3-L3}] The proof of \Cref{Y-P3-L3} involves no sum over partitions, and so its bounds carry no factorial. The fold $\operatorname{Fold}F_t$ of the function $F_t$ of a cell $t$ is made of copies of the sources of $t$ and of nothing else; for this reason the bounds on $\langle\operatorname{Fold}F_t\rangle^{1/N^3}$ at the end of Steps (4) and (5) involve, besides $\mathsf b$, $n_t$ and the sup norms of the sources of $t$, the pair $(\check\delta_t,\check\rho_t)$ of $t$ itself, and no pair of another cell. The copy $r_{\varkappa'\varkappa}F$ of \eqref{support:2-26} is $(r_{\varkappa'\varkappa}F)(U)=F(r_{\varkappa'\varkappa}^{-1}U)$, that is, $F$ composed with the action on bond fields of the inverse of the mirror map $r_{\varkappa'\varkappa}$ of Step (3). This follows from the notation recalled at the opening of \S\ref{sec:chessboard-torus}. In that opening, the paragraph ``The lattice, the action and the integrals'' provides the action of a diagonal isometry on functions and on bond fields, while the paragraph ``Mirror maps, copies and folds'' provides \eqref{support:2-26}, the fact that a mirror map is a diagonal isometry, and the definition \eqref{support:2-28} of the bond field $u^\sigma$.\end{remark}

\emph{Step 4: clause \textup{(G2)}, with \textup{(Lab3)}.}

\begin{lemma}[The bounds of clause \textup{(G2)}]\label{Y-P3-L4} In the setting of \Cref{Y-P3-ThG}, the bounds \textup{(i)} to \textup{(iv)} of its clause \textup{(G2)} hold, and \textup{(Lab3)} of item~1 of \S\ref{Y-P3-1} holds for \textup{(G2)\,(iii)}.\end{lemma}

\begin{proof} Throughout the proof we put $\check\delta:=\check\delta_n:=\min\{d,1/(4n)\}\le1/8$.

\emph{Step 4\,(a): the cut.} The functional $\sigma_{K,N}$ is $n$-linear on real bounded functions. By \Cref{Y-P3-L2}\,(a) and this $n$-linearity, $\sigma_{K,N}(f)$ is the sum of the numbers $\sigma_{K,N}(g_1,\dots,g_n)$, one for each way of picking a single piece $g_i$ of every $f_i$; this sum has at most $(4R_f/\check\delta+2)^{3n}$ terms. Pieces of different $f_i$ are at distance at least $d\ge\check\delta$ from each other, so every tuple of picked pieces is of type $(\check\delta;\check\delta/4)$.

\emph{Step 4\,(b): the case $K\ge j^\circ(\mathsf b,\check\delta)$.} Line~1 of \textup{(G1)} of \Cref{Y-P3-ThG} gives $a_K<\check\delta\le1/(4n)<5/(8n)$. Let $v$ be the vector of \Cref{Y-P3-L2}\,(b). By Step~7 of the proof of \Cref{Y-H9-T}, in its paragraph on finite $K$, the tuple $X$ of picked pieces satisfies $\sigma_{K,N}(\tau_vX)=\sigma_{K,N}(X)$. We use \textup{(G3)\,(b)} with $\check\delta_t=\check\delta$ and $\check\rho_t=\check\delta/4$; at most $n$ cells are occupied, and we obtain \[|\sigma_{K,N}(g_1,\dots,g_n)|\le(\check\Gamma(\mathsf b,\check\delta)e^{c_{\mathrm{vol}}})^n\prod\|g_i\|_\infty.\]

\emph{Step 4\,(c): the case $K<j^\circ$.} Here line~6 of \textup{(G1)} gives $|\mathcal O_Kg-\langle \mathcal O_Kg\rangle|\le2B_0(\mathsf b,\check\delta)\|g\|_\infty$ for each piece $g$.

\emph{Step 4\,(d): proof of \textup{(i)}.} In both cases, line~8 of \textup{(G1)} bounds each term of the sum of Step~4\,(a) by $(A_4\check\delta^{-2})^n\prod\|f_i\|_\infty$. As $4R_f/\check\delta+2\le4(1+R_f)/\check\delta$, the number of terms is at most $(4(1+R_f)/\check\delta)^{3n}$. So $|\sigma_{K,N}(f)|\le[64A_4(1+R_f)^3\check\delta^{-5}]^n\prod\|f_i\|_\infty$. Finally, $1/\check\delta=\max\{1/d,4n\}\le4n\max\{1,1/d\}$ and $n^n\le e^nn!$; inserted into the last bound, these two inequalities give \textup{(i)}.

\emph{Step 4\,(e): cumulants, and proof of \textup{(ii)}.} For $n\ge2$ the joint cumulant does not change when constants are added to its entries, and it is expressed through the centred moments: \[S^T_n=\sum_\pi(-1)^{|\pi|-1}(|\pi|-1)!\prod_{J\in\pi}\sigma_{K,N}(f_J),\] where $\pi$ runs over the partitions of $\{1,\dots,n\}$. A partition with a block of size $1$ contributes $0$, since the centred moment of a single function vanishes. We apply \textup{(i)} to each block with $|J|\ge2$ and use $\prod|J|!\le n!$; then every product in the sum is at most the right side of \textup{(i)}. Hence $|S^T_n|$ is at most $a_n$ times the right side of \textup{(i)}, where $a_n:=\sum_\pi(|\pi|-1)!$. The exponential generating function of the numbers $a_n$ is $-\log(2-e^x)$, whose derivative is $2/(2-e^x)-1$. Expanding this derivative into a series, one finds $a_n=2f_{n-1}$ for $n\ge2$, with $f_m=\sum_{j\ge0}j^m2^{-j-1}$. Since $t\mapsto t^m2^{-t}$ is unimodal on $[0,\infty)$, the sum $\sum_jj^m2^{-j}$ is at most the integral of this function plus its maximum; these are $m!(\log2)^{-m-1}$ and $(m/(e\log2))^m$, and $m!(\log2)^{-m-1}+(m/(e\log2))^m\le m!(\log2)^{-m-1}(1+\log2)$. It follows that $a_n\le1.694\,(n-1)!\,1.4427^n\le1.45^{n+1}n!$, which gives \textup{(ii)}.

\emph{Step 4\,(f): proof of \textup{(iii)}.} The bounds obtained so far hold, for every $(K,N)$, for the entries carried to the torus $T_N$, and the carried entries keep their type as soon as $N\ge2(\mathrm{diam}+d)$, where $\mathrm{diam}$ is the diameter, for the sup distance, of the union of the supports of the entries. A limit of numbers obeys every bound that the numbers obey, and $\sigma_S$ is a polynomial in the $S_{|J|}$. By item~1 of \S\ref{Y-H9-1}, in its part on distances, a tuple of Euclidean type $(d;R)$ is of sup type $(d/\sqrt3;R)$; this costs the factor $3^{5n/2}$.

\emph{Step 4\,(f$'$): proof of \textup{(Lab3)}.} Let the labels move, $\mathsf b_j\to \mathsf b$. The sides tend to $\infty$, while $s=8r$ is bounded on compact sets by item~1 of \S\ref{Y-H9-1}, in its part on radii; hence, for every large $j$, \textup{(Tor)} holds for the label $\mathsf b_j$. The $j$-th member of the sequence obeys \textup{(i)} and \textup{(ii)} with $A_5(\mathsf b_j)$. The limit of the values is at most the $\limsup$ of the bounds, and at fixed $n,d,R_f$ this is the bound taken at $\limsup_jA_5(\mathsf b_j)$. Passing to a subsequence does not change the limit of the values; so the $\limsup$ may be computed along a subsequence. The value of a $\limsup$ is not affected by finitely many members; hence it is at most the supremum over any set that contains a tail of the sequence. If $\mathsf b_j\ge x_s$ for infinitely many $j$, then along these $j$ we have $A_5(\mathsf b_j)\le A_5^c(\mathsf b_j)\to A_5^c(\mathsf b)$, by \textup{(Lab2)} of item~1 of \S\ref{Y-P3-1}. Otherwise the set $\{\mathsf b_j\}\cup\{\mathsf b\}$ is compact in $(0,\infty)$, and one integer $z\ge1$ satisfies $L^z\mathsf b_j\ge x_s$ for all $j$. For every large $j$ the ranges of \textup{(G2)\,(iv)} hold for this $z$, because the sides tend to $\infty$, the cutoff tends to $\infty$ (at each side in \textup{(R5)\,(a)} and along the sequence in \textup{(R5)\,(c)}), and $s$ is bounded on the compact set $\{L^{z+1}\mathsf b_j\}\cup\{L^{z+1}\mathsf b\}$. Then \textup{(iv)}, applied with that $z$, gives the bound with $L^{4z}A_5(L^z\mathsf b_j)$, whose $\limsup$ is at most $L^{4z}A_5^c(L^z\mathsf b)$.

\emph{Step 4\,(g): proof of \textup{(iv)}.} The argument of Step~5 of the proof of \Cref{Y-H9-T} gives \[\sigma_{K,l,L^K\mathsf b}(f)=L^{zn}\sigma_{K-z,L^zl,L^{K-z}(L^z\mathsf b)}(f^{[z]}).\] In addition, $(1+L^zR_f)^3\le L^{3z}(1+R_f)^3$, and $\max\{1,(L^zd)^{-5n}\}\le\max\{1,d^{-5n}\}$ for $z\ge0$. The argument applies to every integer $z\ge0$ with $L^z\mathsf b\ge x_s$, and it does not use that $z$ is the least such integer.\end{proof}

\begin{remark}[Where $n$ enters] The integer $n$ enters the bounds of \textup{(G2)} at the following places. The factor $(4n)^{3n}$ comes from the number of the pieces; this number remains as a factor because the bound of Step~4\,(d) uses the sup norm of each piece and not the volume of its support. The factor $(4n)^{2n}$ comes from $A_4\check\delta^{-2}$. The factor $n!$ comes from Step~4\,(e). Everything else is of the form $C^n$, because the exponent in \textup{(R1)} is linear in $Q$ and $C_{\mathrm{vol}}$ occurs only once.\end{remark}

\emph{Step 5: the extension lemma, and \textup{(G4)\,(b)}, \textup{(d)}.}

\begin{lemma}[Extension to test functions flat on the diagonals]\label{Y-P3-L5} Let $n\ge1$, $a\in\mathbb N_0$ and $\varGamma\ge0$. Let $T$ be a real $n$-linear map on the $n$-tuples of continuous functions on $\mathbb R^3$ with compact and pairwise disjoint supports. Suppose that $|T(f)|\le\varGamma\max\{1,\check\delta^{-a}\}\prod\|f_i\|_\infty$ for every $\check\delta>0$ and every such tuple whose supports are pairwise at sup distance at least $\check\delta$ and each lie in a sup-ball of radius $1$, wherever in $\mathbb R^3$ these balls are placed. Then there is a complex linear functional $\tilde T$ on $\mathcal S_{\ne}(\mathbb R^{3n})$ such that $\tilde T(\bigotimes X)=T(X)$ for every $X$ in $\mathfrak S_n^\infty(\mathbb R^3)$ and $|\tilde T(\varphi)|\le\varGamma\,C_4^n\,39^{a+3n+1}|\varphi|_{a+7n+1}$, where $C_4\le2.3\cdot10^{20}$.\end{lemma}

Steps (h) and (i) of the proof of \Cref{Y-P3-L5} yield \textup{(G4)\,(b)} and \textup{(d)} of \Cref{Y-P3-ThG}.

\begin{proof} \emph{Step (a): a partition of unity off the diagonals.} Let $\rho(t)=c_\rho(1-16t^2)^5$ for $|t|\le\frac14$ and $\rho(t)=0$ for $|t|>\frac14$; with $c_\rho=693/128$, its integral is $1$. The function $\phi_1=1_{[-1/2,1/2]}*\rho$ is of class $C^5$, is carried by $[-\frac34,\frac34]$, satisfies $0\le\phi_1\le1$, and has $\sum_k\phi_1(t-k)=1$, where $k$ runs over $\mathbb Z$. Put $\chi(y)=\prod_\mu\phi_1(y_\mu)$ for $y\in\mathbb R^3$. Then $X_4:=\max_{|\alpha|\le4}\sup|\partial^\alpha\chi|\le15131.7$: the derivatives of $\phi_1$ are pieced together from explicit polynomials, and $\sup|\phi_1^{(j)}|\le1,\ 5.42,\ 45.1,\ 866.3,\ 15131.7$ for $j=0,\dots,4$.

For $j\ge1$, $\phi_1^{(j)}(t)=\rho^{(j-1)}(t+\frac12)-\rho^{(j-1)}(t-\frac12)$, and the two terms have disjoint carriers; hence $\sup|\phi_1^{(j)}|=\sup|\rho^{(j-1)}|$. In the variable $u=4t$ one has $\rho=c_\rho(1-u^2)^5$. The first derivative $\rho'=-40c_\rho u(1-u^2)^4$ has its largest absolute value at $u=\frac13$, namely $40c_\rho\cdot4096/19683\le45.07$. The second derivative $\rho''=16c_\rho(1-u^2)^3(90u^2-10)$ has its largest absolute value at $u=0$, namely $160c_\rho=866.25$. The third derivative is $\rho'''=64c_\rho\cdot240u(1-u^2)^2(1-3u^2)$; the polynomial $u(1-u^2)^2(1-3u^2)$ has its critical points at $u^2=1$, at $u^2=(14-\sqrt{112})/42$ and at $u^2=(14+\sqrt{112})/42$, where its absolute values are $0$, $\le0.1819584$ and $\le0.0995$ respectively, and so $|\rho'''|$ is at most $15360c_\rho\cdot0.1819584\le15131.7$. For mixed derivatives of $\chi$ of total order at most $4$, the products of these suprema are smaller ($5.42\cdot866.3$, $45.1^2$, $5.42^2\cdot45.1$).

For $m\ge0$ let $h_m=2^{-m-1}$, and for $c\in(\mathbb Z^3)^n$ let $\Phi_{m,c}(\xi)=\prod_i\chi(\xi_i/h_m-c_i)$, where $\xi=(\xi_1,\dots,\xi_n)$ with $\xi_i\in\mathbb R^3$. We call $c$ admissible when $|c_i-c_k|_\infty\ge6$ for $i\ne k$. Let $\mathsf A_m=\sum_{c\ \mathrm{admissible}}\Phi_{m,c}$, and $\mathsf A_{-1}=0$. Suppose that $\mathsf A_{m-1}(\xi)>0$. Then $|\xi_i-\xi_k|\ge4.5h_{m-1}=9h_m$ for $i\ne k$, so that every $c$ with $\Phi_{m,c}(\xi)>0$ has $|c_i-c_k|\ge7.5$; all these $c$ are admissible, and $\mathsf A_m(\xi)=1$. It follows that $\mathsf A_m(1-\mathsf A_{m-1})=\mathsf A_m-\mathsf A_{m-1}=:\Upsilon_m$, and that $\sum_{m\le m_1}\Upsilon_m=\mathsf A_{m_1}$, which equals $1$ wherever $\min|\xi_i-\xi_k|\ge7.5h_{m_1}$. Thus $1=\sum_m\Upsilon_m$ off the diagonals, and the sum is locally finite; the $\Upsilon_m$ are obtained by sums and products alone, without any quotient. Finally, $1-\mathsf A_{m-1}$ is the sum of the $\Phi_{m-1,c'}$ over the non-admissible $c'$; at every point of its support, some pair satisfies $|\xi_i-\xi_k|_\infty\le6.5h_{m-1}=13h_m$.

\emph{Step (b): one piece.} Let $\varphi\in\mathcal S_{\ne}(\mathbb R^{3n})$. Fix $m$ and an admissible $c$, put $h=h_m$, and let $G=\Phi_{m,c}(1-\mathsf A_{m-1})\varphi$. Then $G$ is carried by $\prod_i\{|x-hc_i|\le3h/4\}$. It has continuous derivatives $\partial^\gamma$ for each $\gamma$ of order at most $4$ in the three coordinates of every entry, and so of total order at most $4n$: the cut-offs are finite sums of products of $C^5$ functions, each of a single coordinate, and $\varphi$ is smooth. Only these derivatives of $G$ enter the argument. For each $i$ we choose a continuous $\tilde\chi_i$ with $0\le\tilde\chi_i\le1$, equal to $1$ on $|x-hc_i|\le3h/4$ and carried by $|x-hc_i|\le5h/4$; the supports of the $\tilde\chi_i$ are at least $3.5h$ apart, and each lies in a ball of radius $1.25h<1$. A Fourier series of period $4h$ in each coordinate gives \[G=(\prod_i\tilde\chi_i)\sum_k\hat G(k)\prod_ie_{k_i},\qquad e_k(x)=\exp\{2\pi ik\cdot(x-hc_i)/(4h)\}.\] With $\lambda=4h/(2\pi)$ we have $(1-\lambda^2\Delta_i)e_{k_i}=(1+|k_i|^2)e_{k_i}$, and integration by parts yields \[|\hat G(k)|\le\prod_i(1+|k_i|^2)^{-2}\sup|\prod_i(1-\lambda^2\Delta_i)^2G|.\] We set \[\tilde T(G):=\sum_k\hat G(k)\,T^{\mathbb C}(\tilde\chi_1e_{k_1},\dots,\tilde\chi_ne_{k_n}),\] where $T^{\mathbb C}$ denotes the complex $n$-linear extension of $T$. The complex entries cost a factor $2^n$, and $\max\{1,(3.5h)^{-a}\}\le h^{-a}$. Further, $\sum_{\mathbb Z^3}(1+|k|^2)^{-2}\le10.1$, by Poisson's formula, which gives the value $\pi^2\sum_me^{-2\pi|m|}\le9.9986$ for this sum.

\emph{Step (c): derivatives.} Written out, $(1-\lambda^2\Delta)^2$ consists of $1+6+9=16$ monomials $\lambda^{|\gamma|}\partial^\gamma$, $|\gamma|\le4$. The Leibniz rule over the three factors of $G$ produces at most $3^{4n}$ terms. The first factor obeys $|\partial^\gamma\Phi_{m,c}|\le X_4^nh^{-|\gamma|}$. At any point at most $2^{3n}$ of the $c'$ contribute to $1-\mathsf A_{m-1}$, so $|\partial^\gamma(1-\mathsf A_{m-1})|\le8^nX_4^nh^{-|\gamma|}$. Since $\lambda/h=2/\pi<0.64<1$ and $\lambda<1$, no power of $\lambda$ or of $\lambda/h$ exceeds $1$. Collecting these factors and those of Step (b), we obtain \[|\tilde T(G)|\le\varGamma h^{-a}(2\cdot10.1\cdot16\cdot81\cdot8\cdot X_4^2)^n\max_{|\gamma|\le4n}\sup_{\mathrm{supp}\,G}|\partial^\gamma\varphi|.\]

\emph{Step (d): flatness, for $m\ge1$.} Let $m\ge1$ and $\xi\in\mathrm{supp}\,G$, and take the pair $i\ne k$ of Step (a), with $|\xi_i-\xi_k|_\infty\le13h$. Let $\xi^*$ be $\xi$ with $\xi_i$ replaced by $\xi_k$; this is a point of a diagonal, where every derivative of $\varphi$ vanishes. Put $v=\xi_i-\xi_k$; then $|v|_1\le39h$. Taylor's formula with remainder on the segment from $\xi^*$ to $\xi$ gives, for each order $r$, \[|\partial^\gamma\varphi(\xi)|\le\frac{(39h)^r}{r!}\max_{|\beta|=r}\sup_{\mathrm{seg}}|\partial^{\gamma+\beta}\varphi|.\]

\emph{Step (e): weights.} Because $(1+|y|^2)^n\ge\prod_i(1+|y_i|^2)$, we have, for $k\ge4n$, the bound $|\partial^\alpha\varphi(y)|\le|\varphi|_k\prod_i(1+|y_i|^2)^{-2}$ whenever $|\alpha|\le k$. On the segment of Step (d), $|y_i-hc_i|_\infty\le13.75h\le6.875$. The inequality $1+|u|^2\le2(1+|u-y|^2)(1+|y|^2)$ then gives $(1+|y_i|^2)^{-2}\le8.2\cdot10^4(1+|hc_i|^2)^{-2}$, as $4(1+3\cdot6.875^2)^2\le81{,}564$. The same inequality, used on the cubes of side $h$, together with the value $\pi^2$ of the integral of $(1+|x|^2)^{-2}$ over $\mathbb R^3$, gives $\sum_{c\in\mathbb Z^3}(1+h^2|c|^2)^{-2}\le5.641\pi^2h^{-3}\le57h^{-3}$. It is at this point that the hypothesis of the lemma is used for supports placed anywhere in $\mathbb R^3$.

\emph{Step (f): the sum.} We put $C_4:=2\cdot10.1\cdot16\cdot81\cdot8\cdot X_4^2\cdot8.2\cdot10^4\cdot57\le2.3\cdot10^{20}$. Let $G_{m,c}$ be the function $G$ of Step (b) for the scale $m$ and the admissible $c$. By Steps (c) to (e), the contributions $|\tilde T(G_{m,c})|$ of a scale $m\ge1$ add up to at most $\varGamma C_4^nh^{-a-3n}\frac{(39h)^r}{r!}|\varphi|_{4n+r}$; the choice $r=a+3n+1$ makes this at most $\varGamma C_4^n39^rh_m|\varphi|_{a+7n+1}$, and $\sum_{m\ge1}h_m=\frac12$. At the scale $m=0$, without the flatness of Step (d), the contributions add up to at most $\varGamma C_4^n2^{a+3n}|\varphi|_{4n}$. Altogether, the contributions of all scales add up to at most $\varGamma C_4^n39^{a+3n+1}|\varphi|_{a+7n+1}$. Hence the series $\tilde T(\varphi):=\sum_{m,c}\tilde T(G_{m,c})$ converges absolutely, and $\tilde T(\varphi)$ is linear in $\varphi$.

\emph{Step (g): values at tensors.} Let $X=(f_1,\dots,f_n)\in\mathfrak S_n^\infty$ and $K_i=\mathrm{supp}f_i$. Restricted to the product of the spaces of continuous functions carried by $K_i$, the map $T$ induces a linear map $\Lambda$ on their algebraic tensor product, which embeds in the space of functions on $\prod K_i$. Let $G=\bigotimes g_i$ be a tensor with $g_i\in C^5$ carried by $|x-hc_i|\le3h/4$. Then $g_i=\tilde\chi_i\sum_k\hat g_i(k)e_k$ with uniform convergence; as $T$ is continuous for the sup norm in each entry on fixed separated carriers, $\tilde T(G)=T(g)$. At $m=0$, $G_{0,c}$ is the tensor of the functions
$\chi(\cdot/h_0-c_i)f_i$, because $\mathsf A_{-1}=0$.
For $m\ge1$ use the complementary coarse-grid sum. Then $G_{m,c}$
is a finite sum over $c'$ of tensors of
$\chi(\cdot/h_m-c_i)\chi(\cdot/h_{m-1}-c'_i)f_i$.
Only finitely many $(m,c)$ contribute, since $\Upsilon_m=0$ on the support $\prod K_i$ of $\bigotimes X$ once the sets $K_i$ are pairwise at sup distance at least $7.5h_{m-1}$. As functions, these finitely many tensors add up to $\bigotimes X$. Therefore $\tilde T(\bigotimes X)=\Lambda(\bigotimes X)=T(X)$.

\emph{Step (h): proof of \textup{(G4)\,(b)}.} We take $T=\sigma_S$. By \textup{(G2)\,(iii)} of \Cref{Y-P3-ThG}, applied with $R_f=1$, the hypothesis of the lemma holds with $a=5n$ and $\varGamma=(8A_5)^n(n!)^5$. As $39^{8n+1}\le39^{9n}$ and $8C_439^9\le3.9\cdot10^{35}$, the conclusion of the lemma becomes the bound of \textup{(G4)\,(b)}, with a norm of order $12n+1$.

\emph{Step (i): proof of \textup{(G4)\,(d)}.} Now let $T=\sigma_{K,N}$ on bounded functions carried by $(-N/4,N/4)^3$; there the distance on the torus is that of $\mathbb R^3$. The hypothesis of the lemma is \textup{(G2)\,(i)} for every $d>0$; the scales $h_m<a_K$ are covered by Step~4\,(c). We define the lattice sum \[\tilde\sigma(G):=a_K^{-n}\sum G(x(p_1),\dots)\langle\prod(a_p-\langle a_p\rangle)\rangle.\] It is linear in $G$, it depends on finitely many values of $G$, and $\tilde\sigma(\bigotimes f)=T(f)$. Distinct plaquettes have distinct barycentres, and $\varphi=0$ at the points where two of the barycentres coincide. At each point the series of Step (b) converges absolutely, so $\tilde\sigma(G_{m,c})=\tilde T(G_{m,c})$; and at the finitely many points that enter the lattice sum, only finitely many $(m,c)$ contribute. The carrier of $\tilde\chi_i$ stays within $1.25h+0.75h=2h\le2h_0=1$ of the $i$-th projection of $\mathrm{supp}\,\varphi$: it lies within $1.25h$ of $hc_i$, and $hc_i$ lies within $0.75h$ of a point of that projection, since otherwise $G_{m,c}=0$. The margin $1$ in the statement of \textup{(G4)\,(d)} is therefore exactly what is needed, and the box in that statement is open.\end{proof}

\emph{Step 6: \textup{(G4)\,(a)}, \textup{(c)}, \textup{(e)} by \Cref{support:26-24}, uniqueness, and \textup{(G5)}.}

\begin{lemma}[Bound, uniqueness and symmetries of the extension]\label{Y-P3-L6} In the setting of \Cref{Y-P3-ThG}, for every $n\ge1$ and every $\varphi\in\mathcal S_{\ne}(\mathbb R^{3n})$, \[|S^{\sigma_S}_n(\varphi)|\le C_E(\mathsf b)^n(n!)^{13}|\varphi|_{13n-1};\] for every $n$ there is exactly one complex linear map $S^{\sigma_S}_n$ on $\mathcal S_{\ne}(\mathbb R^{3n})$ that is bounded with respect to some norm $|\cdot|_k$ and whose value at the tensor of each $X$ in $\mathfrak S_n^\infty(\mathbb R^3)$ is $\sigma_S(X)$; and these maps satisfy \textup{(E2)} of \Cref{support:28-35} in the form \eqref{support:26-42} at every coordinate mirror, are real and have the symmetry \textup{(E3)}, and are invariant under all translations and under $W_3$, which is \textup{(E1)} for $\mathbb R^3\rtimes W_3$. The same is true for labels $0<\mathsf b<x_s$, with the numbers $A_5(\mathsf b)$, $A_6(\mathsf b)$ and $C_E(\mathsf b)$ of \textup{(G4)\,(e)} of \Cref{Y-P3-ThG}, and for a limit $S$ of \textup{(R5)\,(a)} or \textup{(c)} along labels $\mathsf b_j\to\mathsf b$, with $C_E(\mathsf b)$ taken as in \textup{(Lab3)} of item~1 of \S\ref{Y-P3-1}, namely as $\limsup_jC_E(\mathsf b_j)$ along any subsequence. Uniqueness and the symmetries hold in both cases, whatever the label.\end{lemma}

\begin{proof} \emph{Step 6\,(a): the bound.} Fix $n\ge2$. By \textup{(R6)\,(i)} of \Cref{Y-P2-ThR}, the functional $T:=\sigma_S$ on $\mathfrak S_n^\infty(\mathbb R^3)$ is multilinear in the sense of \Cref{support:25-28}. Take any real number $d>0$. Let $X$ satisfy $R(X)\le1$ and $d(X)\ge d$, both in the Euclidean sense. Since every Euclidean ball is contained in the sup-ball of equal radius, \textup{(G2)\,(iii)} of \Cref{Y-P3-ThG} gives \[|T(X)|\le A_6^n(n!)^5\max\{1,d^{-5n}\}2^{3n}\prod\|f_i\|_\infty\le \mathsf c_{\mathsf T}\max\{1,d^{-5n}\}\prod\|f_i\|_\sharp,\qquad \mathsf c_{\mathsf T}=(8A_6)^n(n!)^5.\] This is \eqref{support:26-25} with $\mathsf a=5n$, and it holds for every such $X$, since the bound of \textup{(G2)} does not depend on where the supports lie. Then \eqref{support:26-26} gives \[|\tilde T(\varphi)|\le \mathsf C_{\ne}^{6n}M^n(n!)^8\mathsf c_{\mathsf T}|\varphi|_{13n-1}=C_E^n(n!)^{13}|\varphi|_{13n-1}\le C_E^n(n!)^{13}|\varphi|_{13n}.\] For $n=1$ the set $\mathrm{Part}_{\ge2}(1)$ is empty, so $\sigma_S(f)=0$, and $S^{\sigma_S}_1=0$: the zero functional is bounded, its value at the tensors is $\sigma_S(f)$, and by \Cref{support:26-23} it is the only such functional. Hence $\sigma_S$ satisfies the growth condition of \Cref{support:26-33}, with the quadruple of constants $(13,1,C_E,13)$.

\emph{Step 6\,(b): one functional.} \Cref{support:26-23} states that two extensions which are bounded for some norms and coincide on the tensors are equal. Hence \Cref{Y-P3-L5} and \Cref{support:26-24} give one and the same functional, and \textup{(G4)\,(b)} bounds that same object. Since $10^{36}A_5\le10^{318}M\,8\cdot3^{5/2}A_5$ and $(n!)^5\le(n!)^{13}$, \Cref{Y-P3-L5} and \Cref{support:26-23} suffice by themselves for the growth condition of \Cref{support:26-33} with the quadruple $(13,1,C_E,13)$.

\emph{Step 6\,(c): the cumulants.} By \textup{(R6)\,(i)}, the truncated function $S_n$ of the limit $S$ is $n$-linear. For $T=S_n$ we carry out Step~6\,(a), with $A_6$, and Step (h) of the proof of \Cref{Y-P3-L5}, with $A_5$ and $R_f=1$; in both we use the assertion of \textup{(G2)\,(iii)} on bound \textup{(ii)}, for $S_n$, in place of its assertion on bound \textup{(i)}, for $\sigma_S$. By \Cref{support:26-23} the two steps give the same functional, and the factor $1.45^{n+1}n!$ enters both bounds.

\emph{Step 6\,(d): proof of \textup{(G5)}.} The functional $\sigma_S$ is a symmetric functional on separated tuples over $\mathbb E=\mathbb R^3$, as in \Cref{support:9-12}\,(3), and it satisfies the growth condition. We apply the parts of \Cref{support:26-41} one after the other.\begin{enumerate}\item[\textup{(a)}] \Cref{support:26-41}\,(a) yields reality and symmetry.

\item[\textup{(b)}] \Cref{support:26-41}\,(b) is applied isometry by isometry, to every translation and to every element of $W_3$; its hypothesis, the invariance of $\sigma_S$ under each of these isometries, is given by \textup{(R6)\,(iii)}.

\item[\textup{(c)}] \Cref{support:26-41}\,(c) yields \eqref{support:26-42} at $(e_3,0)$. The hypothesis of this part is \textup{(R6)\,(iv)}: $\sigma_S$ is reflection positive at this mirror in the sense of \Cref{support:9-25}\,(3). For any other coordinate mirror $(\mu,c)$, we take the element of $\mathbb R^3\rtimes W_3$ that maps it to $(e_3,0)$, conjugate by it, and use the invariance obtained in item~(b) of this list.\end{enumerate}

\emph{Step 6\,(e): labels $0<\mathsf b<x_s$.} Steps~6\,(a) to 6\,(d), and Steps (h) and (i) of the proof of \Cref{Y-P3-L5}, depend on the label only through \textup{(G2)\,(i)} to \textup{(iii)}. By \textup{(G2)\,(iv)}, the clauses \textup{(i)} to \textup{(iii)} hold for these labels, with the number $A_5(\mathsf b)=L^{4z}A_5(L^z\mathsf b)$. Moreover, \textup{(R5)} holds for every $\mathsf b>0$, and every limit of \textup{(R5)} satisfies \textup{(R6)\,(i)}, \textup{(iii)}, \textup{(iv)}. Hence these steps carry over to the labels $0<\mathsf b<x_s$, with the numbers of \textup{(G4)\,(e)}.

\emph{Step 6\,(f): moving labels.} Steps~6\,(a) to 6\,(c) depend on the label only through the single number $A_5$ (or $A_6$) in \textup{(G2)\,(iii)}. This number is taken as \textup{(Lab3)} of item~1 of \S\ref{Y-P3-1} prescribes, and $C_E$ is linear in it.\end{proof}

\begin{remark}[The exponent of $d$ and the order of the norm] Flatness of order $a+3n+1$ pays for the factor $h^{-a}$ and for the factor $h^{-3n}$ that comes from the sum over the grid points. Hence the exponent $a$ is allowed to be linear in $n$, provided the norm has an order that is linear in $n$ too; this is the form of the growth condition of \Cref{support:26-33}, with its norm of order $k_{\rm E}n$. A factor $d^{-cn^2}$ or $e^{c/d}$ would not do, and neither would a prefactor $e^{cn^4}$.\end{remark}

\emph{Step 7: the field between two vectors, and \textup{(G6)}, \textup{(G6+)}.}

\begin{lemma}[Norm of the operator of one small piece]\label{Y-P3-L7} In the setting of item~1 of \S\ref{Y-P3-1}, let $S$ be any limit of \textup{(R5)}, whose label is any $\mathsf b>0$, either fixed or approached by labels $\mathsf b_j\to\mathsf b$; let $\eta\in(0,1]$, and let $g$ be a real function of class $C^1_c$ that is carried by a cube of side $\eta'/4$ lying inside $\{|x_3|<\eta'/2\}$. Put $\mathbb E_+=\{x_3>0\}$. Then the operator given on representatives $F\in\mathbb C\mathfrak S(\mathbb E_+)$ by $\hat A_g[F]:=[(\tau_{\eta'e_3}g)\cdot\tau_{2\eta'e_3}F]$ is well defined on $D=[\mathbb C\mathfrak S(\mathbb E_+)]$. If the label is fixed and $\mathsf b\ge x_s$, then \[\|\hat A_g\|\le\check\Gamma(\mathsf b,\eta'/2)e^{2\eta'c_{\mathrm{vol}}}\|g\|_\infty.\] If $\mathsf b_j\to\mathsf b$, in the order \textup{(a)} or \textup{(c)} of \textup{(R5)}, and $\mathsf b_j\ge x_s$ for infinitely many $j$, then \[\|\hat A_g\|\le\limsup_j\check\Gamma(\mathsf b_j,\eta'/2)e^{2\eta'c_{\mathrm{vol}}(\mathsf b_j)}\|g\|_\infty,\] the upper limit being taken along these $j$.\end{lemma}

From \Cref{Y-P3-L7} we obtain the main clause of \Cref{Y-P3-ThG}\,\textup{(G6)}, every bullet of \textup{(G6)}, and also \textup{(G6+)}.

\begin{proof}[Proof of \Cref{Y-P3-L7}, of \Cref{Y-P3-ThG}\,\textup{(G6)} and of \textup{(G6+)}] \emph{Step (a): reduction to one piece.} As $\eta'\ge\eta/2$, we have $\eta/4\le\eta'/2$, so the test function $h$ of \Cref{Y-P3-ThG}\,\textup{(G6)} has $\mathrm{supp}\,h\subset\{|x_3|<\eta'/2\}$; we define $\hat A_h$ by the formula of the lemma. For $G,G'\in\mathbb C\mathfrak S(\mathbb E_+)$, the scalar product of $\mathcal H(S)$ gives \[\langle[G],\hat A_h[G']\rangle=\sigma_S(\theta\overline G\cdot\tau_{\eta'e_3}h\cdot\tau_{2\eta'e_3}G').\] After a translation of every entry by $-\eta'e_3$, and with $\tau_{-\eta'e_3}\theta=\theta\tau_{\eta'e_3}$, the right side equals $\sigma_S(\theta\overline{(\tau_{\eta'e_3}G)}\cdot h\cdot\tau_{\eta'e_3}G')$. Now let $F,F'$ be as in \textup{(G6)}, and choose $G=\tau_{-\eta'e_3}F$ and $G'=\tau_{-\eta'e_3}F'$. Then \[\sigma_S(\theta\overline F\cdot(h)\cdot F')=\langle\Psi,\hat A_h\Psi'\rangle\] with $\Psi'=[\tau_{-\eta'e_3}F']$ and \[\Psi=[\tau_{-\eta'e_3}F]=e^{-(\eta-\eta')H}[\tau_{-\eta e_3}F];\] as $e^{-(\eta-\eta')H}$ is a contraction, $\|\Psi\|\le\|[\tau_{-\eta e_3}F]\|$.

We now split $h$. Let $s=\eta'/8$, let $y\in\mathbb R^3$ be arbitrary, and for $c\in y+s\mathbb Z^3$ put $\chi_c(x)=\prod_\mu w((x_\mu-c_\mu)/s)$, where $w\in C^5$ is carried by $[-1,1]$ and satisfies $0\le w\le1$ and $\sum_kw(t-k)=1$; we take $w=1_{[-1/2,1/2]}*\rho_2$, with $\rho_2$ carried by $[-\frac12,\frac12]$. Then $h=\sum_ch\chi_c$ is a finite sum, and each piece $h\chi_c$ is of class $C^1$ and is carried by the intersection of its cube of side $\eta'/4$ with $\mathrm{supp}\,h$, hence by a cube of side $\eta'/4$ inside $\{|x_3|<\eta'/2\}$. Moreover $\|h\chi_c\|_\infty\le\sup_{|x-c|\le\eta'/8}|h|$. Since $\hat A$ is linear in $h$, \[\|\hat A_h\|\le\sum_c\|\hat A_{h\chi_c}\|,\] and the lemma bounds the term of $c$ by $\check\Gamma(\mathsf b,\eta'/2)e^{2\eta'c_{\mathrm{vol}}}\sup_{|x-c|\le\eta'/8}|h|$; summed over $c$, these suprema give $|\!|\!|h|\!|\!|_{\eta,y}$. The steps (b) to (g) concern one piece $g$, and they are free of the grid.

\emph{Step (b): the adjoint.} The reflection $\theta_{\eta'}:x_3\mapsto2\eta'-x_3$ leaves $\sigma_S$ invariant. We have $\theta_{\eta'}\theta=\tau_{2\eta'e_3}$, then $\theta_{\eta'}\tau_{\eta'e_3}g=\tau_{\eta'e_3}g^r$ with $g^r(x)=g(x_1,x_2,-x_3)$, and $\theta_{\eta'}\tau_{2\eta'e_3}F=\theta F$. Applying $\theta_{\eta'}$ to every entry of $\langle[F'],\hat A_g[F]\rangle=\sigma_S(\theta\overline{F'}\cdot\tau_{\eta'e_3}g\cdot\tau_{2\eta'e_3}F)$, we obtain $\sigma_S(\tau_{2\eta'e_3}\overline{F'}\cdot\tau_{\eta'e_3}g^r\cdot\theta F)$. Next we apply $\theta$ to every entry, as \Cref{Y-P2-ThR}\,\textup{(R6)\,(iii)} permits; with $\theta^2=1$ and the symmetry of $\sigma_S$, this gives $\sigma_S(\theta\tau_{\eta'e_3}g^r\cdot\theta\tau_{2\eta'e_3}\overline{F'}\cdot F)$, which is $\langle\hat A_{g^r}[F'],[F]\rangle$ by the definition, since $\overline{g^r}=g^r$ for real $g$. Thus \[\langle[F'],\hat A_g[F]\rangle=\langle\hat A_{g^r}[F'],[F]\rangle\] for all $F,F'\in\mathbb C\mathfrak S(\mathbb E_+)$.

\emph{Step (c): multiple reflections.} Put $B=\hat A_{g^r}\hat A_g$ and, for $\Psi\in D$ and $k\ge0$, $a_k=\langle\Psi,B^{2^k}\Psi\rangle$. By step (b), $a_0=\|\hat A_g\Psi\|^2$ and $B$ is symmetric, so the Cauchy--Schwarz inequality yields $a_k\le\|\Psi\|a_{k+1}^{1/2}$. Applying this for $0,1,\dots,k-1$ in turn, we get \[a_0\le\|\Psi\|^{2-2^{1-k}}a_k^{2^{-k}}.\] These inequalities hold as well for the seminorm on $\mathbb C\mathfrak S(\mathbb E_+)$, with $\hat A_g$ acting on representatives; taken there, they show that $\hat A_g$ maps null vectors to null vectors, so that $\hat A_g$ is well defined on $D$.

\emph{Step (d): the moment $a_k$.} Let $m=2^k$ and $\Psi=[F_0]$ with $F_0\in\mathbb C\mathfrak S(\mathbb E_+)$. Unfolding the $2m$ factors of $B^m$, we find \[a_k=\sigma_S(\theta\overline{F_0}\cdot g_1\cdots g_{2m}\cdot\tau_{4m\eta'e_3}F_0),\] where $g_j=\tau_{(2j-1)\eta'e_3}(g^r\text{ or }g)$. The copy $g_j$ lies in the slab $((2j-2)\eta',2j\eta')$, at distance at least $\eta'/2$ from its faces. So the $g_j$ are pairwise at least $\eta'$ apart, each is at least $\eta'/2$ away from the entries of $\theta F_0$ and of $\tau_{4m\eta'e_3}F_0$, and each lies in a cube of side $\eta'/4$.

\emph{Step (e): placing.} We put $\check\delta:=\eta'/2$, so that the copies $g_j$ are of the type $(\check\delta;\check\delta/4)$. We write $F_0=\sum\zeta_aX_a$ with tuples $X_a$ of lengths $\le n_F$, and we fix a term $(\theta X_a,X_{a'})$ of the resulting expansion of $a_k$. Let \[\FitDisplay{\dimexpr\linewidth-1pt\relax}{\check\delta_F:=\min\{\check\delta,\ 1/(32n_F),\ \text{the least distance between two different entries of the same tuple, }X_a\text{ or }X_{a'}\};}\] this number depends on $F_0$ and $\eta$, but not on $m$. With the partition of step (a), taken at $s=\check\delta_F/4$, we cut each entry of $X_a$ and of $X_{a'}$ into pieces carried by cubes of side $\check\delta_F/2$. A term of the resulting sum contains one piece of each entry, so that its pieces are pairwise at least $\check\delta_F$ apart; and the number of terms does not depend on $m$.

Next we translate the whole tuple by a real vector $v$. In coordinate $3$ we have $1=p'\cdot2\eta'$; hence, if one wall lies within $\eta'/4$ of a face of a slab, every wall does, and then the copies are at distance at least $\eta'/4=\check\delta/2$ from the walls. The $v_3$ for which this happens form a set of measure $p'\cdot\eta'/2=\frac14$ modulo $1$. Each of the at most $2n_F$ pieces forbids an arc of length $1.5\check\delta_F$, so together they forbid at most $3/32$. In the coordinates $1$ and $2$, the stack of copies forbids $\eta'/4+\eta'/2\le3/8$, and the pieces forbid at most $3/32$. Hence a translation $v$ that avoids every forbidden set exists.

\emph{Step (f): the bound on $a_k$.} To the translated term we apply \Cref{Y-P3-ThG}\,\textup{(G3)\,(b)} at finite $(K,N)$, with $N$ large and with $K\ge j^\circ(\mathsf b,\check\delta_F)\ge j^\circ(\mathsf b,\check\delta)$. By line~1 of \textup{(G1)}, both numbers $j^\circ(\mathsf b,\check\delta_F)$ and $j^\circ(\mathsf b,\check\delta)$ exist, and $j^\circ$ is non-increasing in $\check\delta$; since $\check\delta_F\le\check\delta$, this gives the second inequality. A pure cell, that is, one that holds copies of $g$ alone, holds at most $p'$ copies, and for it $(\check\delta_t,\check\rho_t)=(\check\delta,\check\delta/4)$; since the stack has length $2m/p'$, there are at most $2m/p'+2$ pure cells. The mixed cells, which hold pieces of entries, number at most $2n_F$; for them $(\check\delta_t,\check\rho_t)=(\check\delta_F,\check\delta/4)$ and $n_t\le2n_F+p'$. Each source lies in one cell, at distance at least $\check\delta_F/2$ from the walls. We conclude that \[|a_k|\le C[\check\Gamma(\mathsf b,\check\delta)e^{c_{\mathrm{vol}}/p'}\|g\|_\infty]^{2m},\] where $C=C(F_0,\eta,g)$ does not depend on $m$: it is made of the number of terms, $|\zeta|^2$, the sup norms of the entries, $e^{(2+2n_F)c_{\mathrm{vol}}}$, $\max\{1,\check\Gamma(\mathsf b,\check\delta_F,\check\delta/4)\}^{2n_F(2n_F+p')}$, $\max\{1,\|g\|_\infty\}^{2n_Fp'}$ and, for $g\ne0$, $\max\{1,(\check\Gamma(\mathsf b,\check\delta)\|g\|_\infty)^{-2n_Fp'}\}$. This bound passes to $S$. It is here that the cell-wise form of \textup{(G3)} is needed: with one global $\check\delta$ we would get $\check\Gamma(\mathsf b,\check\delta_F)^{2m}$, which depends on $F_0$.

\emph{Step (g): the norm of $\hat A_g$.} As $C$ is free of $m$, the bound of step (f) on $a_k^{1/m}$ tends to $[\check\Gamma e^{c_{\mathrm{vol}}/p'}\|g\|_\infty]^2$ as $k\to\infty$, where $\check\Gamma=\check\Gamma(\mathsf b,\check\delta)$. Together with step (c), and with $\check\delta=\eta'/2$ and $1/p'=2\eta'$, this gives \[\|\hat A_g\|\le\check\Gamma(\mathsf b,\eta'/2)e^{2\eta'c_{\mathrm{vol}}}\|g\|_\infty,\] which is the bound of the lemma. Next we use lines~2, 4 and~5 of \Cref{Y-P3-ThG}\,\textup{(G1)} at $\check\delta=\eta'/2\in[\eta/4,\frac14]$: since $\mathfrak l$ is non-increasing in $\check\delta$, $1/\check\delta\le4/\eta$ and $2\eta'\le1$, the right side is at most $\check\Gamma_*(\mathsf b,\eta)\|g\|_\infty$. Combined with step (a), this bound gives the main clause of \Cref{Y-P3-ThG}\,\textup{(G6)}, for every $y\in\mathbb R^3$. Indeed, by step (a), $|\sigma_S(\theta\overline F\cdot(h)\cdot F')|=|\langle\Psi,\hat A_h\Psi'\rangle|$, where $\|\Psi\|\le\|[\tau_{-\eta e_3}F]\|$ and $\Psi'$ obeys the same bound with $F'$ in place of $F$; and by step (a), with the bound just proved for each piece, $\|\hat A_h\|\le\sum_c\|\hat A_{h\chi_c}\|\le\check\Gamma_*(\mathsf b,\eta)\,|\!|\!|h|\!|\!|_{\eta,y}$. Hence \[|\sigma_S(\theta\overline F\cdot(h)\cdot F')|\le\check\Gamma_*(\mathsf b,\eta)\,|\!|\!|h|\!|\!|_{\eta,y}\,\|[\tau_{-\eta e_3}F]\|\,\|[\tau_{-\eta e_3}F']\|.\] Line~7 of \textup{(G1)} gives $\mathfrak l(\mathsf b,\eta/4)\le\mathfrak l_1\cdot4/\eta$, whence \[\check\Gamma_*\le4A_2e^{C_W^*+e_h^*}(1+e^{c_{\mathrm{vol}}})e^{c_{\mathrm{vol}}}(\mathfrak l_1/\mathsf b)^{1/2}/\eta.\]

\emph{Step (h): time-ordered products.} Let $g_i=\tau_{-t_ie_3}h_i\in C^1_c(\{|x_3|<\eta_i/4\})$, and let $A_i=\hat A_{g_i}$ be the operator defined with $\eta'_i$. Once translated into $\mathbb E_+$, the class of $(h_1,\dots,h_n)$ is \[e^{-s_0H}A_1e^{-s_1H}A_2\cdots A_n\Omega,\] where $s_i=t_{i+1}-t_i-\eta'_i-\eta'_{i+1}\ge0$, and $\sigma_S(h)$ is the scalar product of $\Omega$ with this vector. Step (a) with the grid at $y=-t_ie_3$, together with step (g), gives $\|A_i\|\le\check\Gamma_*(\mathsf b,\eta_i)|\!|\!|g_i|\!|\!|_{\eta_i,-t_ie_3}$. Moreover \[|\!|\!|\tau_{-te_3}h|\!|\!|_{\eta,-te_3}=\sum_{c_0\in(\eta'/8)\mathbb Z^3}\sup\{|h(x+te_3)|:|x-(c_0-te_3)|\le\eta'/8\}=|\!|\!|h|\!|\!|_{\eta,0}.\] Therefore the bound holds with the fixed grid $y=0$, that is, with the norms $|\!|\!|h_i|\!|\!|_{\eta_i}$, as stated in \Cref{Y-P3-ThG}\,\textup{(G6)}.

\emph{Step (i): the norm $|\!|\!|\cdot|\!|\!|$.} On a cube $Q$ of side $s$, \[\sup_Q|u|\le\sum_{\alpha\in\{0,1\}^3}s^{|\alpha|-3}\int_Q|\partial^\alpha u|;\] indeed, in one variable the supremum is at most the mean plus $\int|u'|$, and we iterate this over the three coordinates. We apply it on the cubes that enter $|\!|\!|h|\!|\!|_{\eta,y}$, whose side is $s=\eta'/4\ge\eta/8$, so that $s^{|\alpha|-3}\le512\,\eta^{|\alpha|-3}$; almost every point lies in $8$ of these cubes, on every translate of the grid. The two factors give the constant $8\cdot512=4096$.

\emph{Step (j): other directions.} For the other coordinate directions we pass to $e_3$ by an element of the group $W_3$.

\emph{Step (k): the last bullet of \Cref{Y-P3-ThG}\,\textup{(G6)}.} Let $\mathsf b>0$, let $\mathsf b_j\to \mathsf b$ (or $\mathsf b_j=\mathsf b$ for all $j$), and let $z\ge1$ be an integer with $L^z\mathsf b_j\ge x_s$ from some $j$ on; such a $z$ exists because $\{\mathsf b_j\}\cup\{\mathsf b\}$ is compact in $(0,\infty)$. If the label is fixed, \Cref{Y-P2-ThR}\,\textup{(R6)\,(vi)} gives $S_n(f)=L^{zn}S'_n(f^{[z]})$; here $S'$ is taken along the matched subsequence, at the label $L^z\mathsf b\ge x_s$, as a limit in the same order. For moving labels we rely on Step~5 of the proof of \Cref{Y-H9-T}; its statement is valid for all $\beta>0$ and all $z$, and by it the member with index $j$ of the sequence forming $S$ is $L^{zn}$ times the value at $f^{[z]}$ of a member with label $L^z\mathsf b_j$. In the order \textup{(c)} of \textup{(R5)} this reads \[S^T_{n;K_j,l_j,L^{K_j}\mathsf b_j}(f)=L^{zn}S^T_{n;K_j-z,L^zl_j,L^{K_j-z}(L^z\mathsf b_j)}(f^{[z]}),\] and in the order \textup{(a)}, by the last sentence of that statement, \[S^{(l_j,\mathsf b_j)}_n(f)=L^{zn}S^{(L^zl_j,L^z\mathsf b_j)}_n(f^{[z]}).\]

The dilation $f\mapsto f^{[z]}$ maps the set of tuples whose entries are continuous, with compact and pairwise disjoint supports, one to one onto itself; hence the dilated members converge on all such tuples. Their limit $S'$ is a limit of \textup{(R5)}, in the same order of limits, taken along sides $L^zl_j=2L^{m_j+z}\to\infty$ and labels $L^z\mathsf b_j\to L^z\mathsf b$; these labels are $\ge x_s$ for all large $j$, and $S_n(f)=L^{zn}S'_n(f^{[z]})$. Condition \textup{(Tor)} of item~1 of \S\ref{Y-P3-1} at $L^z\mathsf b_j$ holds for large $j$. Since $\sigma_S$ is a sum, over partitions, of products, we get $\sigma_S(X)=L^{zn}\sigma_{S'}(X^{[z]})$. Let $D_z$ be the linear map $X\mapsto L^{zn}X^{[z]}$. Dilation commutes with $\theta$, and $(\tau_vf)^{[z]}=\tau_{L^zv}f^{[z]}$. Consequently $\langle[F],[F']\rangle_S=\langle[D_zF],[D_zF']\rangle_{S'}$, and \[\sigma_S(\theta\overline F\cdot(h)\cdot F')=\sigma_{S'}(\theta\overline{D_zF}\cdot(L^zh^{[z]})\cdot D_zF').\]

Here every entry of every tuple of $D_zF$ has its support in $\{x_3>L^z\eta\}\subset\{x_3>L^z\eta_z\}$, $\mathrm{supp}\,h^{[z]}$ lies in $\{|x_3|<L^z\eta_z/4\}$, and $L^z\eta_z\le1$. Clause \textup{(G6)} for $S'$ at $(L^z\mathsf b,L^z\eta_z)$, which holds for moving labels by step (l), bounds the modulus of this quantity by \[L^z\check\Gamma_*(L^z\mathsf b,L^z\eta_z)|\!|\!|h^{[z]}|\!|\!|_{L^z\eta_z,y}\|[\tau_{-L^z\eta_ze_3}D_zF]\|_{S'}\|[\tau_{-L^z\eta_ze_3}D_zF']\|_{S'}.\] Finally \[\|[\tau_{-L^z\eta_ze_3}D_zF]\|_{S'}=\|[\tau_{-\eta_ze_3}F]\|_S=\|e^{-(\eta-\eta_z)H}[\tau_{-\eta e_3}F]\|\le\|[\tau_{-\eta e_3}F]\|,\] and the same chain holds with $F'$ in place of $F$.

\emph{Step (l): moving labels.} Suppose that $\mathsf b_j\ge x_s$ for infinitely many $j$; we pass to these $j$, and the subsequence has the same limit $S$. In step (f), the $j$-th member of the sequence carries the label $\mathsf b_j$. By \textup{(Lab1)} of item~1 of \S\ref{Y-P3-1}, $j^\circ(\mathsf b_j,\check\delta_F)$ is bounded in $j$, while $K\to\infty$ in \Cref{Y-P2-ThR}\,\textup{(R5)\,(a)}, \textup{(c)}; so the condition on $K$ is met from some $j$ on, and, as in Step~4\,(f$'$) of the proof of \Cref{Y-P3-L4}, condition \textup{(Tor)} at $\mathsf b_j$ holds for large $j$, because the sides go to infinity while $s(x)=8r(x)$ is bounded on compact sets. The numbers that make up $C$ stay bounded in $j$: $c_{\mathrm{vol}}$ is continuous; by item~2 of \S\ref{Y-P2-1} and \Cref{Y-P3-ThG}\,\textup{(G1)}, $\check\Gamma(\cdot,\check\delta_F,\check\delta/4)$ is locally bounded; and since $1/r^\circ=\mathrm{amp}(L^{j^\circ+1}\mathsf b)/3.5$, with $j^\circ(\cdot,\check\delta)$ locally bounded on $[x_s,\infty)$ by \textup{(Lab1)} and with $\mathrm{amp}$ continuous and positive jointly in $x$ and the label, we have $\check\Gamma(\cdot,\check\delta)\ge1/r^\circ>0$, where $1/r^\circ$ has a positive lower bound on every compact set of labels in $[x_s,\infty)$. Hence \[|a_k|\le C'[\limsup_j\check\Gamma(\mathsf b_j,\check\delta)e^{c_{\mathrm{vol}}(\mathsf b_j)/p'}\|g\|_\infty]^{2m},\] with $C'$ free of $m$. From this and step (c), with $\check\delta=\eta'/2$ and $1/p'=2\eta'$, we argue as in step (g): the limit $k\to\infty$ yields the first inequality in \[\|\hat A_g\|\le\limsup_j\check\Gamma(\mathsf b_j,\eta'/2)e^{2\eta'c_{\mathrm{vol}}(\mathsf b_j)}\|g\|_\infty\le\limsup_j\check\Gamma_*(\mathsf b_j,\eta)\|g\|_\infty,\] and lines~2, 4 and~5 of \textup{(G1)} at $\check\delta=\eta'/2$, applied at each $\mathsf b_j\ge x_s$, yield the second, as in \textup{(Lab3)}. In the remaining case, $\mathsf b_j<x_s$ from some $j$ on (so that $\mathsf b\le x_s$), and steps (e) to (g) do not apply to these members. Indeed, $\check\Gamma_*(\mathsf b_j,\cdot)$ and $j^\circ(\mathsf b_j,\cdot)$ are defined under the data, and \textup{(G1)} and \textup{(R1)} are stated under these data, whose label is $\ge x_s$. This case is treated by step (k), and by step (n).

\emph{Step (m): \Cref{Y-P3-ThG}\,\textup{(G6+)\,(m)}.} Let $\eta>0$ be arbitrary, let $\hat\eta=\min\{1,3\eta/4\}$ and $s=(\hat\eta)'/8$, and write $h=\sum_ch\chi_c$ on the grid $y+s\mathbb Z^3$, with the partition of step (a). We fix a piece different from $0$ and let $t$ be the point of $[-\eta/4,\eta/4]$ nearest to $c_3$. This piece lies in $\{|x_3-t|\le s\}$. Indeed, $t=c_3$ when $|c_3|\le\eta/4$; when $c_3>\eta/4$, the points of the piece satisfy $c_3-s\le x_3<\eta/4=t$, whence $0<t-x_3<s$; and when $c_3<-\eta/4$, they satisfy $-\eta/4=t<x_3\le c_3+s$, whence $0<x_3-t<s$. The function $g:=\tau_{-te_3}(h\chi_c)$ is real and of class $C^1$, and it is carried by a cube of side $2s=(\hat\eta)'/4$ inside $\{|x_3|\le s\}\subset\{|x_3|<(\hat\eta)'/2\}$.

We translate all entries by $-te_3$, as \Cref{Y-P2-ThR}\,\textup{(R6)\,(iii)} permits, and use $\tau_{-te_3}\theta=\theta\tau_{te_3}$; this gives \[\sigma_S(\theta\overline F\cdot(h\chi_c)\cdot F')=\sigma_S(\theta\overline{(\tau_{te_3}F)}\cdot(g)\cdot\tau_{-te_3}F').\] The tuples $\tau_{te_3}F$ and $\tau_{-te_3}F'$ are carried by $\{x_3>\eta-|t|\}\subset\{x_3>3\eta/4\}\subset\{x_3>\hat\eta\}$. Steps (a) and (g), for the single piece $g$ at $\hat\eta\in(0,1]$ (and step (l) at moving labels), bound the term by \[\check\Gamma_*(\mathsf b,\hat\eta)\|g\|_\infty\|[\tau_{(t-(\hat\eta)')e_3}F]\|\,\|[\tau_{(-t-(\hat\eta)')e_3}F']\|.\] Now $[\tau_{(t-(\hat\eta)')e_3}F]=e^{-(\eta+t-(\hat\eta)')H}[\tau_{-\eta e_3}F]$ with $\eta+t-(\hat\eta)'\ge\eta-\eta/4-3\eta/4=0$, so its norm is at most $\|[\tau_{-\eta e_3}F]\|$, the semigroup being a contraction by \textup{(R7)}; the same holds with $-t$ in place of $t$ and $F'$ in place of $F$. Finally $\|g\|_\infty\le\sup\{|h(x)|:|x-c|\le s\}$, and summing over $c$ we obtain $|\!|\!|h|\!|\!|_{\hat\eta,y}$.

\emph{Step (n): \Cref{Y-P3-ThG}\,\textup{(G6+)\,(n)}.} We take the argument of step (k) without change, from the choice of $\mathsf b_j$ and $z$ up to the identity \[\sigma_S(\theta\overline F\cdot(h)\cdot F')=\sigma_{S'}(\theta\overline{D_zF}\cdot(L^zh^{[z]})\cdot D_zF').\] Now $D_zF,D_zF'$ are carried by $\{x_3>L^z\eta\}$, and $h^{[z]}$ by $\{|x_3|<L^z\eta/4\}$. Clause \textup{(G6+)\,(m)}, which admits any positive number, applied to $S'$ with $L^z\eta$ in the place of $\eta$, gives the bound \[L^z\check\Gamma_*(L^z\mathsf b,e_z)|\!|\!|h^{[z]}|\!|\!|_{e_z,y}\|[\tau_{-L^z\eta e_3}D_zF]\|_{S'}\|[\tau_{-L^z\eta e_3}D_zF']\|_{S'},\] and $\|[\tau_{-L^z\eta e_3}D_zF]\|_{S'}=\|[\tau_{-\eta e_3}F]\|_S$, and the same identity holds with $F'$ in place of $F$.\end{proof}

\begin{remark}[The lower bound on $\check\Gamma(\mathsf b_j,\check\delta)$ in step~\textup{(l)}] In step~\textup{(l)} of the proof of \Cref{Y-P3-L7}, the lower bound of $\check\Gamma(\mathsf b_j,\check\delta)$ enters only through the factor $\max\{1,(\check\Gamma\|g\|_\infty)^{-2n_Fp'}\}$. By the definitions of item~1 of \S\ref{Y-P3-1}, we have $\check\Gamma(\mathsf b_j,\check\delta)=G_1(\mathsf b_j,\check\delta,\check\delta/4)(1+e^{c_{\mathrm{vol}}(\mathsf b_j)})\ge G_1$ and $G_1=e^{C_W+e_h}/r^\circ\ge1/r^\circ$.\end{remark}

\begin{remark}[\Cref{support:4-1} and \Cref{support:4-15}] The notions recalled here are defined in \Cref{support:4-1}\,\textup{(2)} and in the paragraphs ``Boards'' and ``Mirror maps, copies and folds'', which follow that definition in \S\ref{sec:chessboard-torus}. The quadruple $(a_K,L^K,(N,N,N),0)$ is a board: $L^K$ is odd and $\ge3$, $N=2L^{m_l}$ is even and $\ge2$, a cell has side $1$, and there are $N^3$ cells. The mirrors $o_\mu+\mathsf l(t+\frac12)$ pass through midpoints of bonds and hold no site \eqref{support:2-14}. A function of a cell is a bounded Borel function of its inner bonds. A mirror map has the form $x\mapsto o+\mathsf l\varkappa'+\sigma(\varkappa',\varkappa)(x-o-\mathsf l\varkappa)$ \eqref{support:2-18}. The fold $\operatorname{Fold}F$ is the product of the $N^3$ copies, conjugated for odd pairs, which changes nothing for real $F$ \eqref{support:2-26}, \eqref{support:2-27}. With $\|F\|_{\rm ch}=\langle\operatorname{Fold}F\rangle^{1/N^3}$ as in \Cref{support:4-13}, \eqref{support:4-14}, \Cref{support:4-15} states $|\langle\prod F_\varkappa\rangle|\le\prod\langle\operatorname{Fold}F_\varkappa\rangle^{1/N^3}$.\end{remark}

\begin{remark}[\S\ref{Y-P5-1}] Let $z(\mathsf b)$ be the least $z\ge0$ with $L^z\mathsf b>x_s$, as in \textup{(S7)} of \S\ref{Y-00-1}. In \Cref{Y-P5-ThE}\,\textup{(E5)}, the growth bound with its quadruple of constants holds with $\mathsf C_{\ne}^6ML^z(1+L^z)^3A_6(L^z\mathsf b)$, $z=z(\mathsf b)$, in the place of $C_E$; the factor $L^z(1+L^z)^3$ there is to be compared with the factor $8L^{4z}$ here. When $z(\mathsf b)=z_0(\mathsf b)$, the number $\bar c_*(\mathsf b)=(3/(4\pi))L^{4z}A_6(L^z\mathsf b)$ and the number $(3/(4\pi))A_6(\mathsf b)$, with $A_6(\mathsf b)=L^{4z}A_6(L^z\mathsf b)$ of \textup{(G4)\,(e)}, are given by one formula, $(3/(4\pi))L^{4z}A_6(L^z\mathsf b)$. In \textup{(E5)}, however, $A_6(L^z\mathsf b)$ is the supremum over $L^zU(\mathsf b)$, whereas here it is the value at the label, or the number of \textup{(Lab3)}; hence $\bar c_*(\mathsf b)\ge(3/(4\pi))A_6(\mathsf b)$.\end{remark}

\begin{remark}[What \textup{(G1)} uses of $\mathrm{sep}^\circ$] Clause \textup{(G1)} of \Cref{Y-P3-ThG} uses $\mathrm{sep}^\circ\ge0$ in Step~\textup{(0)} of the proof of \Cref{Y-P3-L1}, for $2\check r+3\le\check\lambda$; it uses $\mathrm{sep}^\circ>0$ in line~1 and in $T_L<\infty$; and it uses $\frac12\le\mathrm{sep}^\circ+\log L/528$ in line~3 at $j=1$. At the value $\mathrm{sep}^\circ=8\log(8L^{1/2})$ of item~2 of \S\ref{Y-00-1}, we have $\mathrm{sep}^\circ\ge8\log8>16.6$, and all three conditions hold.\end{remark}

\section{Symmetries of continuum limits}\label{Y-P4-1}

On one torus and at one label, we bound the chessboard seminorm of a product of centred sources that lie in one cell, with free centres, first on the unit board and then on boards with small cells. Domination at a fixed cutoff on the infinite lattice carries these bounds over to every Gibbs state that is invariant under $a_K\mathbb Z^3$; this yields one table of moments for all $K\ge1$, and compactness. For every limit $S\in\mathrm{Lim}_G(\mathsf b)$ we then show that $\sigma_S$ has the properties \textup{(R0)} to \textup{(R7)} of \Cref{support:32-25}: an insertion bound with the power $\eta^{-1}$, a growth of the plates with exponent $7/2<4$, and positivity at the coordinate and at the diagonal mirrors. We then obtain invariance under $\mathbb R^3\rtimes\operatorname{O}(3)$, under $\operatorname{O}(3)$ from \Cref{support:37-34} and under translations from \textup{(R1)}; positivity at every mirror from \Cref{support:28-32}\,(c), or from \Cref{support:9-31}\,(b), (c); the circular cone and the bound $\|\mathcal T(z)\|\le1$ from \Cref{support:28-34}; and the group $\operatorname{O}(2)$ from \Cref{support:28-30}. The dilation $x\mapsto L^zx$ carries all of this over to every label $\mathsf b>0$. Besides, we treat joint limits with large sides, and we show that for every limit, reflection positivity at one diagonal mirror is equivalent to invariance under $\operatorname{O}(3)$.

Of \Cref{Y-00-Y}, this section proves clause~\textup{(Y5)}, the third line of~\textup{(Y2)}, the last two sentences of~\textup{(Y3)} and the last paragraph of~\textup{(Y4)}. The letters, the board, the number $c_{\rm vol}$ and the ranges are those of \S\ref{Y-00-1}.

\begin{remark}[The letters of this section]
Without a reference, \textup{(R0)} to \textup{(R7)} always denote the properties of \Cref{support:32-25}, while a clause of \Cref{Y-P2-ThR} always appears together with that reference. The letter $N$ always stands for the side of the torus, and $\mathsf l$ always for the side of a cell; where the data of \Cref{Y-H9-T} write $l$ for the side of the torus, we write $N$, and $m_l$ is the exponent of the side in those data. The threshold in $K$ is $j^\circ_U(\mathsf b,\check\delta)$; neither the number $K_0(\mathsf b)$ of entry~\textup{(O14)} in item~1 of \S\ref{Y-H9-1} nor the index set of \eqref{support:55-5} carries a letter here. Without arguments, $B_0$ is the board; with its arguments, $B_0(\mathsf b,\check\delta)$ is the number of item~2 of \S\ref{Y-P2-1}. The number $c_{\rm I}$ is the one of \Cref{support:32-25}; no number of this section depends on the letter $c_I$ that item~(1) of \S\ref{Y-H0-1} adds to entry~\textup{(O0)}.
\end{remark}

\emph{0.1. Data and constants.} We take the data of \Cref{Y-H9-T} and of item~1 of \S\ref{Y-H9-1}: the group $G=\operatorname{SU}(2)$ in dimension~$3$, a label $\mathsf b\ge x_s$, the side $N=2L^{m_l}$ with $m_l\ge m'$ and the condition \textup{(Tor)} $N\ge2Ms(L\mathsf b)$, a number of steps $K\ge1$, and $\beta=L^K\mathsf b$. From the same place come Wilson's measure $\mu_{K,N,\beta}$ on $a_K\mathbb Z^3/N\mathbb Z^3$, where $a_K=L^{-K}$, and the source $\mathcal O_K(f)$. There, too, the separation $d$, the radii and the balls refer to the sup distance of the torus, and a Euclidean type $(d;R)$ counts as a sup type $(d/\sqrt3;R)$. The entries \textup{(O4a)} and \textup{(O7)} in item~(1) of \S\ref{Y-H0-1}, the table of item~1 of \S\ref{Y-H9-1} and the remark after \Cref{Y-H9-T} in item~2 of \S\ref{Y-H9-1} fix $L$ and $M=L^{m'}$: by that remark $L$ is odd with $L\ge L_*^{\rm su}\ge10^{2048.9}$, and by the table $m'\ge1$. These values lie in the ranges in which the sections~\ref{support:4-1}, \ref{support:18-2}, \ref{support:54-1} and~\ref{support:55-1} are stated, namely $L$ odd with $L\ge13$ and $m'\ge0$ (\Cref{support:55-1}\,(1)), and \Cref{Y-P4-Th4} is asserted for the $L$ and $M$ of the table alone. From item~1 of \S\ref{Y-H8-1} we take $q_N=(N/M)^3\ge1$ and \[C_{\rm vol}(N,x)=4q_Ne^{220q_N}e^{(3/2)\mathsf N_Zp_0(x)^2N^3}(384/\varepsilon(x)^3)^{3N^3}.\] The table and the radii of item~1 of \S\ref{Y-H9-1} provide $p_0(x)=A_0(\log x)^{\mathfrak p_0}$, $\varepsilon(x)=p_0(x)x^{-1/2}$ and $c_h(x)=e^{-c_hp_0(x)^2}$, the floor $x_s\ge e^2$, the bound $p_0\ge1$ on $[x_s,\infty)$, and $\check r(x)=\lceil20c_R\log x/M\rceil+3$.

\emph{0.2. Results of earlier sections.} \Cref{Y-P2-ThR}\,\textup{(R0)} gives the lattice symmetries of the subtraction. It says that the subtraction $\hat c_K(f)=\frac{d}{dz}\log\rho_K^{[zf]}(1)\big|_{z=0}$ is real, linear in $f$, and dependent on $K$ and $N$ (we write it $\hat c_{K,N}$ where the dependence on $N$ matters), and that $\hat c_K(f\circ r)=\hat c_K(f)$ for $r:y_\mu\mapsto c-y_\mu$ at every integer $c$; we use it at odd $c$. The subtracted moment estimate is \Cref{Y-P2-ThR}\,\textup{(R1)} and we use it for $K\ge j^\circ$. Its data are a label $\mathsf b\ge x_s$, a side $N$ allowed by the data of \Cref{Y-H9-T}, a number $K\ge j^\circ(\mathsf b,\check\delta)$, an arbitrary $Q\ge1$, and real bounded functions $g_t$ on $T_N$ whose sets $\{g_t\ne0\}$ are pairwise at sup distance at least $\check\delta$ and lie each in a sup-ball of radius $\check\rho$. For these it gives the bound \[\bigl|\langle\prod_t(\mathcal O_Kg_t-\hat c_Kg_t)\rangle_{K,N}\bigr|\le G_1(\mathsf b,\check\delta,\check\rho)^Q(1+0.182\,C_{\rm vol}(N,\mathsf b))\prod\|g_t\|_\infty.\] On the board $B_0$, \Cref{Y-P2-ThR}\,\textup{(R2)}\,(i), taken at $K\ge j^\circ(\mathsf b,\check\delta)$, is \Cref{Y-P4-Th4}\,\textup{(4.1)}\,(b), (c). In this section these two clauses also follow from \Cref{Y-P4-Th4}\,\textup{(4.1)}\,(a) (Step~1 of the proof of \Cref{Y-P4-Th4}). Of \Cref{Y-P3-ThG} we use clause~\textup{(G1)}, on $0<\check\delta\le1$, with the following lines:
\begin{enumerate}
\item[\textup{Line 1:}] $L^{j^\circ(\mathsf b,\check\delta)}\le2L^2(M/\check\delta)\,\mathfrak l(\mathsf b,\check\delta)$;
\item[\textup{Line 2:}] $1/r^\circ\le A_2(\mathfrak l/(\mathsf b\check\delta))^{1/2}$;
\item[\textup{Line 4:}] $C_W(\mathsf b,\check\delta,\check\delta/4)\le C_W^*(\mathsf b)$;
\item[\textup{Line 5:}] $e_h\le e_h^*(\mathsf b)$;
\item[\textup{Line 7:}] $\mathfrak l(\mathsf b,\check\delta)\le\mathfrak l_1(\mathsf b)(1+\log(1/\check\delta)/528)$;
\item[\textup{Line 8:}] $\max\{\check\Gamma(\mathsf b,\check\delta)e^{c_{\rm vol}},2B_0(\mathsf b,\check\delta)\}\le A_4(\mathsf b)\check\delta^{-2}$.
\end{enumerate}

We use \Cref{Y-P2-ThR}\,\textup{(R5)}, \Cref{Y-P2-ThR}\,\textup{(R6)}\,(i) to (iv) and (vi), \Cref{Y-P3-ThG}\,\textup{(G2)}\,(i) to (iii), of which (i) and (ii) are lines 1 and~2 of \Cref{Y-00-Y}\,\textup{(Y2)}, and the main clause and the first bullet of \Cref{Y-P3-ThG}\,\textup{(G6)}.

\begin{definition}[Item 1: labels, the class of states, limit sets, boards, normal form]
\emph{Labels.} A number $\mathsf b>0$ is the label of a limit, and $\mathsf b'$ is the label of an approximant; always $\beta=L^K\mathsf b'>0$.

\emph{The class.} We denote by $\mathbb G_K(\mathsf b')$ the set of the Gibbs states $\nu$ at $(a_K,L^K\mathsf b')$, in the sense of \Cref{support:14-9}\,(2), with $(\tau_v)_*\nu=\nu$ for all $v\in a_K\mathbb Z^3$. By \Cref{support:14-9}\,(3), with period number~$1$, we have $\mathbb G_K(\mathsf b')\subset\mathcal G^{\rm per}_{K,\beta}$; by \Cref{support:14-16}\,(c), $\mathbb G_K(\mathsf b')\supset\mathcal G^{\rm tor}_{K,\beta}$, so that $\mathbb G_K(\mathsf b')$ is not empty. For $\rho,\nu\in\mathbb G_K(\mathsf b')$ and a tuple $X=(f_1,\dots,f_n)$ we put $X^\nu:=\prod_i(\mathcal O_Kf_i-\nu(\mathcal O_Kf_i))$, give the empty tuple the value $1$, and extend linearly; then \[\sigma_\rho(X):=\rho(X^\rho)=\sum_{\pi\in\mathrm{Part}_{\ge2}(n)}\prod_{J\in\pi}S^T_{|J|}[\rho](X_J).\]

\emph{Limit sets.} We let $\mathrm{Lim}_G(\mathsf b)$ be the set of the families $S=(S_n)_{n\ge2}$ with $S_n(X)=\lim_jS^T_n[\rho_j](X)$ for all $n\ge2$ and all separated $C^1_c$ tuples $X$ over $\mathbb R^3$, where $K_j\to\infty$, $\rho_j\in\mathbb G_{K_j}(\mathsf b_j)$ and $\mathsf b_j\to \mathsf b$. The set $\mathrm{Lim}_{\rm tor}(\mathsf b)$ consists of those families for which $\rho_j\in\mathcal G^{\rm tor}$. The limits of \Cref{Y-P2-ThR}\,\textup{(R5)}\,(b), in which first $N\to\infty$ along a subsequence at fixed $K$ and then $K\to\infty$, lie in $\mathrm{Lim}_{\rm tor}(\mathsf b)$: by \Cref{Y-P2-ThR}\,\textup{(R5)}\,(b) the measures converge weakly along the subsequence of sides, so that their limit lies in $\mathcal G^{\rm tor}$ by \Cref{support:14-9}\,(4), and the cumulants are polynomials in means of continuous local functions. For every family $S=(S_n)_{n\ge2}$, $\sigma_S$ denotes the centred functional \eqref{support:9-14}, with $\sigma_S(\emptyset)=1$.

\emph{Boards.} We use one family $B_i$, the board $(i,0)$ of $\mathbb L_K$ for $-m_l\le i\le K-1$, whose cells are $\mathsf lk+[-\mathsf l/2,\mathsf l/2]^3$ with $\mathsf l=L^{-i}$ and whose walls are $\mathsf l(\mathbb Z+\frac12)$; besides, we use its translates $(i,o)$ with $o\in a_K\mathbb Z^3$. In this section only $i\ge0$ occurs.

\emph{Normal form.} The set $U(\mathsf b)$ is that of item~\textup{(S7)} of \S\ref{Y-00-1}; for $\mathsf b>x_s$, $U(\mathsf b)=[\mathsf b/2,2\mathsf b]\cap[x_s,\infty)$. The labels are in normal form if $\mathsf b>x_s$; then $\mathsf b_j\in U(\mathsf b)$ from some $j$ on.

\emph{Remark.} At $\mathsf b=x_s$, the clauses are given by \Cref{Y-P4-Th4}\,\textup{(4.8)} at $z=1$, for every sequence. For sequences with $\mathsf b_j\ge x_s$, the numbers defined by the supremum over $[x_s,2x_s]$ may be used in the clauses as well, because throughout each proof the label is a single fixed $\mathsf b'\ge x_s$.

Every number $X(\mathsf b)$ below means $\sup_{\mathsf b'\in U(\mathsf b)}X(\mathsf b')$. Each of them is thus one function of $\mathsf b$, which depends neither on the sequence $(\mathsf b_j)$ nor on the state. These suprema are finite, on the following grounds. The number $c_{\rm vol}$ is continuous on $[x_s,\infty)$, by the formulas for $p_0$ and $\varepsilon$ in the table of item~1 of \S\ref{Y-H9-1}; there $p_0\ge1$, so that $\varepsilon>0$. The number $C_W^*$ is bounded on $[x_s,\infty)$, by line~4 of \Cref{Y-P3-ThG}\,\textup{(G1)}. On $U(\mathsf b)$ we have $e_h^*(\mathsf b')=0.45(2\check r(\mathsf b')+4)^3c_h(\mathsf b')\le0.45(2\check r(2\mathsf b)+4)^3$, since $c_h(\cdot)\le1$ and $\check r$ does not decrease (the table and the radii of item~1 of \S\ref{Y-H9-1}). The function $\mathfrak l_1$ is continuous. The numbers $A_4$, $\hat G$, $A_6$, $C_{\rm I}$ and $C_\varphi$ are continuous expressions in these four numbers and in $1/\mathsf b'$. Finally, $\check\Gamma(\mathsf b',\check\delta)\le\hat G(\mathsf b')/\check\delta$ and $\Gamma(\mathsf b',\mathsf l)\le3\hat G(\mathsf b')e^{c_{\rm vol}(\mathsf b')}/(2\mathsf l)$.

We put $j^\circ_U(\mathsf b,\check\delta):=\sup_{\mathsf b'\in U(\mathsf b)}j^\circ(\mathsf b',\check\delta)$, which is finite by line~1 of \Cref{Y-P3-ThG}\,\textup{(G1)}, as $\mathfrak l$ is continuous in $\mathsf b'$. Clauses \textup{(4.1)} and \textup{(4.2)} of \Cref{Y-P4-Th4} are statements on one torus at one label, and no supremum is taken in them. In these two clauses, and throughout the proofs of clauses \textup{(4.3)} to \textup{(4.5)}, a single label $\mathsf b'$ is kept fixed, and the supremum is taken at the end. Every other label $\mathsf b>0$ is treated in clause~\textup{(4.8)}.
\end{definition}

\begin{definition}[Item 2: the numbers of the theorem]
All numbers of this item are placed at \textup{(O14)}. No floor depends on them, and none of them depends on $N$, $K$, $n$, the places of the sources or the state, with one exception: $\check\delta_n(d)$, which depends on $n$ and on $d$. We set $\mathrm{sep}^\circ:=8\log(8L^{1/2})$; it enters here only as a letter, inside $j^\circ$, $\mathfrak l$, $\mathfrak l_1$, $\check\lambda$ and $T_L$. We set \[c_{\rm vol}(\mathsf b):=225/M^3+\frac32\mathsf N_Zp_0(\mathsf b)^2+3\log(384/\varepsilon(\mathsf b)^3),\] so that $1+0.182\,C_{\rm vol}(N,\mathsf b)\le e^{c_{\rm vol}(\mathsf b)N^3}$. We define $\check\Gamma(\mathsf b,\check\delta):=G_1(\mathsf b,\check\delta,\check\delta/4)(1+e^{c_{\rm vol}(\mathsf b)})$ and \[\hat G(\mathsf b):=A_2e^{C_W^*(\mathsf b)+e_h^*(\mathsf b)}(1+e^{c_{\rm vol}(\mathsf b)})(\mathfrak l_1(\mathsf b)/\mathsf b)^{1/2}.\] For $0<\check\delta\le1$, \[\check\Gamma(\mathsf b,\check\delta)\le\hat G(\mathsf b)\check\delta^{-1/2}(1+\log(1/\check\delta))^{1/2}\le\hat G(\mathsf b)/\check\delta;\] the first inequality comes from lines 2, 4, 5 and~7 of \Cref{Y-P3-ThG}\,\textup{(G1)}, and the second from $x(1+\log(1/x))\le1$ on $(0,1]$, taken at $x=\check\delta$. We put $\Gamma(\mathsf b,\mathsf l):=\check\Gamma(\mathsf b,2\mathsf l/3)e^{c_{\rm vol}(\mathsf b)\mathsf l^3}$, $\check\delta_n(d):=\min\{d/\sqrt3,1/(4n)\}$, where $d=\infty$ is allowed, and $A_6(\mathsf b):=3^{5/2}\cdot64(4e)^5A_4(\mathsf b)$. In the order of \Cref{support:32-25}, $(c_{\rm I},C_{\rm I},\mathsf p_{\rm I}):=(1/(48L),\ 24Le^{c_{\rm vol}(\mathsf b)}\hat G(\mathsf b),\ 1)$. Moreover $\mathsf a_\Gamma:=7/2$ and \[C_\varphi:=288\cdot512^2(16\lambda_\varphi+1)^2\|\varphi\|_\infty^2e^{2c_{\rm vol}(\mathsf b)}\hat G(\mathsf b)^2.\]

For an integer $z\ge1$ with $L^z\mathsf b>x_s$, we put $(c_{\rm I}^{(z)},C_{\rm I}^{(z)},\mathsf p_{\rm I}):=(c_{\rm I}L^{-z}/2,\ 2L^zC_{\rm I}(L^z\mathsf b),\ 1)$. We also put \[C_\varphi^{(z)}:=L^{2z}\max\{L^{-7z/2}C_{\varphi^{[z]}}(L^z\mathsf b),\ 9\|\varphi\|_\infty^2C^\sigma_2(2;(2+\lambda_\varphi)L^z)\},\] where the table $C^\sigma$ of \Cref{Y-P4-Th4}\,\textup{(4.6)} is taken at the label $L^z\mathsf b$, and where $\varphi^{[z]}(x)=\varphi(L^{-z}x)$ has $\lambda_{\varphi^{[z]}}=L^z\lambda_\varphi$ and the sup norm of $\varphi$.
\end{definition}

\begin{theorem}[Invariant Gibbs states; rotations and mirrors of the limits taken volume first]\label{Y-P4-Th4}
In clauses \textup{(4.3)} to \textup{(4.7)} the labels are in normal form; clauses \textup{(4.1)} and \textup{(4.2)} hold at a label $\mathsf b\ge x_s$, in the setting of item~\textup{(S2)} of \S\ref{Y-00-1}.
\begin{enumerate}
\item[\textup{(4.1)}] \emph{The unit board $B_0$, with walls $\mathbb Z+\frac12$.} Let $\mathsf b\ge x_s$, let $N=2L^{m_l}$ with $m_l\ge m'$ satisfy \textup{(Tor)}, and let $K\ge j^\circ(\mathsf b,\check\delta)$.
\begin{enumerate}
\item[\textup{(a)}] \emph{One cell, free centres.} Let $g_1,\dots,g_r$ be real bounded functions that vanish outside one cell, such that the sets $\{g_i\ne0\}$ are pairwise at distance at least $\check\delta$, lie at distance at least $\check\delta/2$ from the faces, and lie each in a sup-ball of radius $\check\rho$, and let the $c_i$ be real numbers. Then \[\bigl\|\prod_i(\mathcal O_Kg_i-c_i)\bigr\|^{(0,0)}_{{\rm ch},N}\le e^{c_{\rm vol}(\mathsf b)}\prod_i\bigl(G_1(\mathsf b,\check\delta,\check\rho)\|g_i\|_\infty+|c_i-\hat c_{K,N}(g_i)|\bigr).\]
\item[\textup{(b)}] This clause is the first sentence of \Cref{Y-P2-ThR}\,\textup{(R2)\,(i)}, taken at $K\ge j^\circ(\mathsf b,\check\delta)$. Let $Q\ge1$ sources with these data be given, each in one cell, and let $q$ cells be occupied. Then \[\bigl|\langle\prod_t(\mathcal O_Kg_t-\hat c_Kg_t)\rangle_{K,N}\bigr|\le G_1^Qe^{qc_{\rm vol}(\mathsf b)}\prod\|g_t\|_\infty.\]
\item[\textup{(c)}] For one source $g$ with the data of clause~\textup{(a)} and $m_N(g):=\langle \mathcal O_Kg\rangle_{K,N}$ we have $|m_N(g)-\hat c_{K,N}(g)|\le G_1e^{c_{\rm vol}(\mathsf b)}\|g\|_\infty$. (This clause is the second sentence of \Cref{Y-P2-ThR}\,\textup{(R2)\,(i)}, taken at $K\ge j^\circ(\mathsf b,\check\delta)$.)
\end{enumerate}
\item[\textup{(4.2)}] \emph{Small cells, centred at the torus mean.} Let $i\ge0$, $\mathsf l=L^{-i}$ and $K\ge\max\{i+1,j^\circ(\mathsf b,\check\delta)\}$, and let $(i,o)$ be any board of $\mathbb L_K$. Let $g^1,\dots,g^r$ be real bounded functions that vanish outside one cell, such that the sets where they do not vanish are pairwise at distance at least $\check\delta$, lie at distance at least $\check\delta/2$ from the faces, and lie each in a sup-ball of radius $\check\delta/4$, and let the $c^s$ be real numbers. Then \[\bigl\|\prod_s(\mathcal O_Kg^s-c^s)\bigr\|^{(i,o)}_{{\rm ch},N}\le e^{c_{\rm vol}(\mathsf b)\mathsf l^3}\prod_s\bigl(\check\Gamma(\mathsf b,\check\delta)\|g^s\|_\infty+|c^s-m_N(g^s)|\bigr).\]
\item[\textup{(4.3)}] \emph{All of $\mathbb G_K(\mathsf b')$.} Let $\mathsf b'\in U(\mathsf b)$ and $\rho,\nu\in\mathbb G_K(\mathsf b')$.
\begin{enumerate}
\item[\textup{(a)}] \emph{Cell by cell.} Let sources on $B_0$ be given, with the data of clause~\textup{(4.1)\,(a)} at $(\check\delta_t;\check\delta_t/4)$ in the cell $t$, and let $K\ge\max_tj^\circ(\mathsf b',\check\delta_t)$, for which $K\ge\max_tj^\circ_U(\mathsf b,\check\delta_t)$ suffices. Then \[\bigl|\rho\bigl(\prod_i(\mathcal O_Kg_i-\nu(\mathcal O_Kg_i))\bigr)\bigr|\le\prod_{t\ {\rm occupied}}\bigl[e^{c_{\rm vol}}\check\Gamma(\mathsf b,\check\delta_t)^{n_t}\bigr]\cdot\prod\|g_i\|_\infty.\] On a board $(i,o)$ with $i\ge0$, with the data of clause~\textup{(4.2)} at the label $\mathsf b'$, the same holds with the factor $e^{c_{\rm vol}\mathsf l^3}\bigl[\check\Gamma(\mathsf b,\check\delta_t)(1+e^{c_{\rm vol}\mathsf l^3})\bigr]^{n_t}$ for each occupied cell.
\item[\textup{(b)}] \emph{The table.} This clause is the third line of \Cref{Y-00-Y}\,\textup{(Y2)}. For all $K\ge1$, all $n\ge1$ and all real bounded $f_i$ of Euclidean type $(d;R)$, wherever they lie, \[\bigl|\rho\bigl(\prod(\mathcal O_Kf_i-\nu(\mathcal O_Kf_i))\bigr)\bigr|\le A_6(\mathsf b')^n(n!)^5\max\{1,d^{-5n}\}(1+R)^{3n}\prod\|f_i\|_\infty.\] Here the plain value $A_6(\mathsf b')$ is obtained by evaluating the formula of $A_6$ in item~2 at $\mathsf b'$ itself, without a supremum over a set $U(\cdot)$ (item~\textup{(S8)\,(n1)} of \S\ref{Y-00-1}; the paragraph ``The label'' of Step~3 of the proof of \Cref{Y-P4-Th4}), and $A_6(\mathsf b')\le A_6(\mathsf b)$. Since every $\mathsf b'\ge x_s$ lies in $U(\mathsf b)$ for some $\mathsf b>x_s$, this bound holds at every $\mathsf b'\ge x_s$. For $n\ge2$, moreover, $|S^T_n[\rho](f)|\le1.45^{n+1}n!$ times the right side of the display of this clause.
\item[\textup{(c)}] Every sequence $(K_j,\rho_j,\mathsf b_j)$ with $K_j\to\infty$ and $\mathsf b_j\to \mathsf b$ has a subsequence that converges on all separated tuples. Hence $\mathrm{Lim}_{\rm tor}(\mathsf b)$ and $\mathrm{Lim}_G(\mathsf b)$ are not empty.
\end{enumerate}
\item[\textup{(4.4)}] \emph{Insertion on the lattice.} Let $(i,o)$ be a board with $i\ge0$, let $\{x_3=c\}$ be one of its walls, and let $\square$ be a cell in $\{c\le x_3\le c+\mathsf l\}$. Let $f$ be real and bounded and vanish outside $\square$, with $\{f\ne0\}$ at distance at least $\mathsf l/3$ from the faces and in a sup-ball of radius $\mathsf l/6$. Let $\mathsf b'\in U(\mathsf b)$, let $K\ge\max\{i+1,j^\circ(\mathsf b',2\mathsf l/3)\}$, for which $K\ge\max\{i+1,j^\circ_U(\mathsf b,2\mathsf l/3)\}$ suffices, let $\rho\in\mathbb G_K(\mathsf b')$, and let $F,F'\in\mathcal C_{\rm loc}(\mathbb H^+_c)$. Then \[\bigl|\rho\bigl(\overline{\theta_cF'}\,[\mathcal O_Kf-\rho(\mathcal O_Kf)]\,\tau_{\mathsf le_3}F\bigr)\bigr|\le2\Gamma(\mathsf b',\mathsf l)\|f\|_\infty\|F'\|_{\rho,c}\|F\|_{\rho,c},\] and $\Gamma(\mathsf b',\mathsf l)\le\Gamma(\mathsf b,\mathsf l)$.
\item[\textup{(4.5)}] \emph{Plates.} Let $\mathsf b'\in U(\mathsf b)$, $\rho\in\mathbb G_K(\mathsf b')$, $0<t\le1$ and $\alpha\in\mathsf{Prof}_t$, let $\varphi$ be a plate label and $f=\alpha\otimes\varphi$, and let $K\ge j^\circ_U(\mathsf b,t/8)$. Then $|S^T_2[\rho](\theta f,f)|\le C_\varphi t^{-7/2}$. More precisely, the left side is at most of the order $t^{-3}(1+\log(8/t))$, as we show in Step~6 of the proof by summing over the pairs of pieces.
\item[\textup{(4.6)}] \emph{The table of hypotheses.} For every $S\in\mathrm{Lim}_G(\mathsf b)$, $\sigma_S$ is a functional on separated tuples over $\mathbb R^3$ with the properties \textup{(R0)} to \textup{(R7)} of \Cref{support:32-25}:
\begin{enumerate}
\item[\textup{(R0)}] holds by \eqref{support:9-14};
\item[\textup{(R1)}] holds by Proposition~\textup{\ref{support:54-17}\,(b)} and clause~\textup{(4.3)\,(b)};
\item[\textup{(R2)}] holds with $C^\sigma_n(d;R):=A_6^n(n!)^5\max\{1,d^{-5n}\}(1+R)^{3n}$, a table that is non-increasing in $d$ by its form and whose bound is in the sup norm, where $\|\cdot\|_\infty\le\|\cdot\|_\sharp$;
\item[\textup{(R3)}] holds by Proposition~\textup{\ref{support:54-17}\,(a)} at $(e_3,0)$;
\item[\textup{(R4)}] holds by \Cref{support:28-21} with Remark~\textup{\ref{support:32-37}\,(1), (2)};
\item[\textup{(R5)}] holds with $(c_{\rm I},C_{\rm I},\mathsf p_{\rm I})$, by clause~\textup{(4.4)};
\item[\textup{(R6)}] holds by Proposition~\textup{\ref{support:54-17}\,(a)} at $(e_1+e_3,0)$, and that clause gives the full positivity at $\Sigma$;
\item[\textup{(R7)}] holds with $\mathsf a_\Gamma=7/2<4$, by clause~\textup{(4.5)}.
\end{enumerate}
\item[\textup{(4.7)}] \emph{Rotations and mirrors.} For every $S\in\mathrm{Lim}_G(\mathsf b)$ we have $\sigma_S(rX)=\sigma_S(X)$ and $S_n(rX)=S_n(X)$ for all $r\in\mathbb R^3\rtimes\operatorname{O}(3)$, and $\sigma_S$ is reflection positive at every mirror $(u,c)$ of $\mathbb R^3$. Moreover $H\ge|P_k|$ for $k=1,2$, the joint spectrum lies in $\{E\ge|q|\}$, and $\|\mathcal T(z)\|\le1$ for $\operatorname{Re}z_3\ge|\operatorname{Im}(z_1,z_2)|$. The group $\operatorname{O}(2)$ acts on $\mathcal H(\sigma_S)$ by a strongly continuous unitary representation, which fixes $\Omega$, commutes with $H$ and extends the representation of $W_2$. Even without \Cref{support:37-34}, the reflection positivity of $\sigma_S$ holds at every coordinate plane, and also at each plane that belongs to one of the six diagonal families.
\item[\textup{(4.8)}] \emph{Every label $\mathsf b>0$.} Take an integer $z\ge0$ with $L^z\mathsf b>x_s$. For $K\ge z+1$, the relabelling $x\mapsto L^zx$ carries $\mathbb G_K(\mathsf b')$ onto $\mathbb G_{K-z}(L^z\mathsf b')$, and $\mathcal O_K(f)=L^z\mathcal O_{K-z}(f^{[z]})$ with $f^{[z]}(x)=f(L^{-z}x)$. Hence every $S\in\mathrm{Lim}_G(\mathsf b)$ is of the form $S_n(X)=L^{zn}S'_n(X^{[z]})$ with $S'\in\mathrm{Lim}_G(L^z\mathsf b)$. The map $D_z:x\mapsto L^zx$ carries $\tilde\mu_{K,2L^m,\beta}$ to $\tilde\mu_{K-z,2L^{m+z},\beta}$; since limit points depend on tails only (Definition~\textup{\ref{support:14-9}\,(4)}), $D_z$ carries $\mathcal G^{\rm tor}_{K,\beta}$ onto $\mathcal G^{\rm tor}_{K-z,\beta}$, and $\mathrm{Lim}_{\rm tor}(\mathsf b)$ corresponds to $\mathrm{Lim}_{\rm tor}(L^z\mathsf b)$. Clause~\textup{(4.3)\,(c)} holds at $\mathsf b$, as one sees by applying it to $(K_j-z,D_z\rho_j,L^z\mathsf b_j)$; so $\mathrm{Lim}_{\rm tor}(\mathsf b)$ and $\mathrm{Lim}_G(\mathsf b)$ are not empty. Clause~\textup{(4.7)} holds for $S$ word for word. The properties \textup{(R0)} to \textup{(R4)} and \textup{(R6)} hold for $S$, with the table $L^{zn}C^\sigma_n(L^zd;L^zR)$ taken at the label $L^z\mathsf b$. For $z=0$, properties \textup{(R5)} and \textup{(R7)} have the constants of \textup{(4.6)}. For $z\ge1$, they hold with $(c_{\rm I}^{(z)},C_{\rm I}^{(z)},1)$ and $(7/2,C_\varphi^{(z)})$, by \Cref{Y-P4-LZ}.
\item[\textup{(4.9)}] \emph{Further limits.} Let $\mathsf b>0$ be any label, and let $\mathsf b_j\to \mathsf b$.
\begin{enumerate}
\item[\textup{(i)}] There is a threshold $N^*(K,\mathsf b')$, which is not explicit, with the following property: whenever a joint sequence $\sigma_{K_j,N_j}$, taken at labels $\mathsf b_j\to \mathsf b$ and with $N_j\ge N^*(K_j,\mathsf b_j)$, has a limit point, that limit point equals the functional of some member of $\mathrm{Lim}_{\rm tor}(\mathsf b)$.
\item[\textup{(ii)}] This clause covers every limit of \Cref{Y-P2-ThR}\,\textup{(R5)}: those of \textup{(a)} and \textup{(c)} also at labels $\mathsf b_j\to \mathsf b$, and those of \textup{(b)} at a fixed label; a limit in the order of \textup{(b)} at moving labels is a member of $\mathrm{Lim}_{\rm tor}(\mathsf b)$, and for it clauses \textup{(4.6)} to \textup{(4.8)} apply. Every such limit has the properties \textup{(R0)} to \textup{(R3)}, \textup{(R5)} and \textup{(R7)}, with these numbers: in normal form, the table $C^\sigma_n$ of clause~\textup{(4.6)}, the triple $(c_{\rm I},C_{\rm I},\mathsf p_{\rm I})$ and the pair $(\mathsf a_\Gamma,C_\varphi)$; otherwise, with $z\ge1$ and $L^z\mathsf b>x_s$, the table $L^{zn}C^\sigma_n(L^zd;L^zR)$ taken at $L^z\mathsf b$, the triple $(c_{\rm I}^{(z)},C_{\rm I}^{(z)},1)$ and the pair $(7/2,C_\varphi^{(z)})$. For such a limit the following are equivalent: invariance under $\operatorname{O}(3)$; reflection positivity at the one mirror $\Sigma=((e_1+e_3)/\sqrt2,0)$; and \textup{(R4)} together with \textup{(R6)}.
\end{enumerate}
\end{enumerate}
\end{theorem}

The per-cell estimate transfers to periodic Gibbs states by chessboard
domination. We use the transferred estimates in two ways. A one-cell
insertion bounds the local form in the reflection Hilbert space. A sum
over pairs of small pieces bounds the plate covariance by
$O(t^{-3}(1+\log(8/t)))$. This is the input for the rotation argument in
Appendix~\ref{app:continuum-rotations}.


The volume factor used below satisfies
\[
 1+0.182\,C_{\rm vol}(N,\mathsf b)
       \le e^{c_{\rm vol}(\mathsf b)N^3}.
\]
Indeed, write $q=(N/M)^3\ge1$ and
$C_{\rm vol}=4qe^{220q}Y$, where $Y\ge1$ is the remaining
factor in the formula of paragraph~0.1. Then
\[
 1+0.182\,C_{\rm vol}
 \le e^{220q}(1+q)Y
 \le e^{221q}Y
 \le e^{225q}Y=e^{c_{\rm vol}N^3}.
\]

\begin{proof}[Proof of \Cref{Y-P4-Th4}]
\emph{Step 1: the unit board $B_0$ with walls $\mathbb Z+\frac12$, with free centres.} We prove \textup{(4.1)\,(a)}, \textup{(b)} and \textup{(c)}.

\emph{Board.} For $K\ge1$ the pair $B_0=(0,0)$ is a board of $\mathbb L_K$, as $\mathsf l/a_K=L^K$ is odd and at least $3$. Its cells are the cubes $k+[-\frac12,\frac12]^3$, and its walls $\mathbb Z+\frac12$ lie in $a_K(\mathbb Z+\frac12)$. The side of the torus satisfies $N=2L^{m_l}\ge2\mathsf l$.

\emph{Functions of a cell.} By \Cref{support:18-7}\,(b), each $\mathcal O_K(g_i)$ is a real continuous element of $\mathcal A(\square)$; so is $F:=\prod_i(\mathcal O_K(g_i)-c_i)$.

\emph{Fold.} By \eqref{support:18-17} in \Cref{support:18-16}, \[\|F\|_{{\rm ch},N}^{N^3}=\tilde\mu_{K,N,\beta}\bigl(\prod_{k'}F^{[k']}\bigr),\] the product running over one period box of $N^3$ cells; we call this mean the fold. As $F$ is real, the conjugations drop out of the formula (\Cref{support:18-7}\,(e)). The map $\hat r_{k'\leftarrow k}$ is a composite of wall reflections $y_\mu\mapsto(2t+1)-y_\mu$, and it carries plaquettes to plaquettes and barycentres to barycentres. The action $a_p$ of a plaquette does not depend on its orientation, because $\operatorname{tr}g^{-1}=\operatorname{tr}g$ in $\operatorname{SU}(2)$. Hence $\hat r\mathcal O_K(g)=\mathcal O_K(g\circ\hat r^{-1})=:\mathcal O_K(g^{k'})$, and $F^{[k']}=\prod_i(\mathcal O_K(g_i^{k'})-c_i)$. After composition with $\mathsf{per}_N$, the fold is the mean under $\mu_{K,N,\beta}$ of a product over the sources $g^{k'}$ on $T_N$.

\emph{Data of the $rN^3$ copies on $T_N$.} Inside one cell the copies are isometric images of the given sources; so they are pairwise at least $\check\delta$ apart, and each of them lies in a sup-ball of radius $\check\rho$. Now take two points of copies that lie in different cells of the board on the torus. In some coordinate $\mu$ their cell indices differ modulo $N$, and on the circle $\mathbb R/N\mathbb Z$ both arcs joining their $\mu$-coordinates leave the interval of one cell and enter that of the other; hence each arc has length at least $\check\delta/2+\check\delta/2$. Since $N$ is even, the pattern of the copies closes up on the torus.

\emph{Subtraction.} We apply \textup{(Y0)} with $c=2t+1$ to each factor of $\hat r$ and obtain $\hat c_{K,N}(g_i^{k'})=\hat c_{K,N}(g_i)$. We set $\varepsilon_i:=\hat c_{K,N}(g_i)-c_i$.

\emph{Expansion.} Then \[\prod_{k',i}\bigl[(\mathcal O_K(g_i^{k'})-\hat c_{K,N}(g_i^{k'}))+\varepsilon_i\bigr]=\sum_T\prod_T(\mathcal O_K-\hat c_{K,N})\cdot\prod_{T^c}\varepsilon,\] where $T$ runs over the subsets of the $rN^3$ copies. To each sub-collection $T$ we apply \textup{(Y1)}, with the same data; here $|T|\ge1$ is arbitrary, while $T=\emptyset$ contributes $1$. This gives \[\bigl|\langle\prod_T(\mathcal O_K-\hat c_{K,N})\rangle_{K,N}\bigr|\le(1+0.182C_{\rm vol})\prod_TG_1\|g\|_\infty,\] and the sum over $T$ yields \[|\text{fold}|\le(1+0.182C_{\rm vol}(N,\mathsf b))\prod_i(G_1\|g_i\|_\infty+|\varepsilon_i|)^{N^3}.\]

\emph{Root.} Since $(1+0.182C_{\rm vol})^{1/N^3}\le e^{c_{\rm vol}(\mathsf b)}$, the $N^3$-th root of this bound is \textup{(a)}.

\emph{Proof of \textup{(b)}.} Clause \textup{(b)} is \Cref{Y-P2-ThR}\,(R2)\,(i); we derive it from \textup{(a)}. Let $F_t:=\prod_{i\in I_t}(\mathcal O_K(g_i)-\hat c_{K,N}(g_i))$, a function of the cell $t$, since every source lies in one cell. \Cref{support:18-18}\,(b), applied to the set of the occupied cells, gives $|\langle\prod_tF_t\rangle_{K,N}|\le\prod_t\|F_t\|_{{\rm ch},N}$, and \textup{(a)}, applied with $c_i=\hat c_{K,N}(g_i)$, bounds each factor.

\emph{Proof of \textup{(c)}.} Clause \textup{(c)} is \Cref{Y-P2-ThR}\,(R2)\,(i) in the case $Q=1$, and it too follows from \textup{(a)}: by \Cref{support:18-18}\,(c), $|\tilde\mu(F)|\le\|F\|_{{\rm ch},N}$, and we apply \textup{(a)} with $r=1$ and $c=\hat c_{K,N}(g)$.

\emph{Step 2: the chessboard seminorm of sources on a board with small cells.} We prove \textup{(4.2)}.

\emph{\textup{(i)} Offset.} By \eqref{support:18-27} applied with $r=\tau_{-o}$, $\|F\|^{(i,o)}_{{\rm ch},N}=\|\tau_{-o}F\|^{(i,0)}_{{\rm ch},N}$. Moreover $\tau_{-o}\mathcal O_K(g)=\mathcal O_K(\tau_{-o}g)$, and $m_N(\tau_{-o}g)=m_N(g)$, because $\mu_{K,N,\beta}$ is invariant under $a_K\mathbb Z^3$. We may therefore take $o=0$; then the walls are $\mathsf l(\mathbb Z+\frac12)$. As $L^i$ is odd, $k+\frac12=\mathsf l\bigl(L^ik+(L^i-1)/2+\frac12\bigr)$, so every plane of $\mathbb Z+\frac12$ is a wall of $B_i$.

\emph{\textup{(ii)} Fold.} We have \[\|\cdot\|_{{\rm ch},N}^{Q_c}=\bigl\langle\prod_\kappa\prod_s(\mathcal O_K(g^s_\kappa)-c^s)\bigr\rangle_{K,N},\] with $Q_c=(N/\mathsf l)^3$. The reasoning of Step~1 carries over: the functions are real, so that no conjugation occurs, and the wall reflections carry plaquettes to plaquettes and barycentres to barycentres, so that the mirror copies are again sources. On $T_N$ the $rQ_c$ copies are pairwise at least $\check\delta$ apart, and each lies in a sup-ball of radius $\check\delta/4$; here $N/\mathsf l=2L^{m_l+i}$ is even.

\emph{\textup{(iii)} One mean.} A wall reflection $y_\mu\mapsto2w-y_\mu$ has $2w\in a_K\mathbb Z$. It is therefore an isometry of $a_K\mathbb Z^3/N\mathbb Z^3$, and $\mu_{K,N,\beta}$ is invariant under it. Hence $\langle\mathcal O_K(g^s_\kappa)\rangle_{K,N}=m_N(g^s)=:m^s$ for every $\kappa$.

\emph{\textup{(iv)} The subtraction.} For $i\ge1$, \textup{(Y0)} yields no invariance of $\hat c_{K,N}$ under the small mirrors. We set $\delta^s_\kappa:=\hat c_{K,N}(g^s_\kappa)-m^s$. By \textup{(i)}, each copy stays at least $\check\delta/2$ away from the planes $\mathbb Z+\frac12$ and lies in a sup-ball of radius $\check\delta/4$; it therefore lies in one unit cell, and \textup{(4.1)\,(c)} gives $|\delta^s_\kappa|\le G_1(\mathsf b,\check\delta,\check\delta/4)e^{c_{\rm vol}}\|g^s\|_\infty$.

\emph{\textup{(v)} Expansion.} We write \[\mathcal O_K(g^s_\kappa)-c^s=\bigl(\mathcal O_K(g^s_\kappa)-\hat c_{K,N}(g^s_\kappa)\bigr)+(\delta^s_\kappa+m^s-c^s)\] and expand the product of these sums over the subsets $T$ of the copies: the term of $T$ is the mean of $\prod_T(\mathcal O_K-\hat c_{K,N})$, multiplied by the constants $\delta^s_\kappa+m^s-c^s$ of the copies outside $T$. For each $T$, \textup{(Y1)} bounds that mean by $(1+0.182C_{\rm vol})\prod_TG_1\|g^s\|_\infty$, and $T=\emptyset$ contributes $1$. With the bound of \textup{(iv)}, the sum over $T$ gives \[|\text{fold}|\le(1+0.182C_{\rm vol})\prod_{\kappa,s}\bigl[G_1\|g^s\|_\infty+G_1e^{c_{\rm vol}}\|g^s\|_\infty+|c^s-m^s|\bigr].\]

\emph{\textup{(vi)} Root.} Finally $(1+0.182C_{\rm vol}(N,\mathsf b))^{(\mathsf l/N)^3}\le e^{c_{\rm vol}(\mathsf b)\mathsf l^3}$, and the $Q_c$-th root of the last bound gives \textup{(4.2)}.

\emph{Remark.} We use \eqref{support:18-27} and \Cref{support:18-7} in the form recalled in \Cref{support:55-1}\,(2)(e).

\emph{Step 3: all of $\mathbb G_K(\mathsf b')$, a single table of moment bounds for all $K\ge1$, compactness.} We prove \textup{(4.3)\,(a)}, \textup{(b)} and \textup{(c)}.

\emph{The label.} We prove \textup{(a)} and \textup{(b)} at one label $\mathsf b'\in U(\mathsf b)$. The data, \textup{(Y1)}, \textup{(4.1)}, \textup{(4.2)} and \textup{(G1)} are taken at $\mathsf b'$: each number is the plain value $X(\mathsf b')$, and each threshold is $j^\circ(\mathsf b',\cdot)$. Only at the end do we use $X(\mathsf b')\le X(\mathsf b)$, where $X(\mathsf b)$ is the supremum over $U(\mathsf b)$.

\emph{\textup{(a)} on $B_0$.} At fixed $K$ we have $|m_N(g)|\le2a_K^{-1}\#\{p:x(p)\in\operatorname{supp}g\}\|g\|_\infty$; together with \textup{(4.1)\,(c)}, it follows that $\hat c_{K,N}(g)$ is bounded in $N$ as well. There is therefore a sequence $N_\bullet=(N_k)\to\infty$ in $\{2L^m\}$, all of whose members are sides of the data, along which $\hat c_{K,N_k}(g_i)\to c_i^\infty$ for each of the finitely many $i$. By \Cref{support:55-40}\,(a), $\nu$ has \textup{(Dom; $N_\bullet$, $\mathcal C_{\rm loc}$)}. First take the one cell of $g_i$ and $F=\mathcal O_K(g_i)-c_i^\infty$, which does not depend on $N_k$; by \eqref{support:55-5} and \textup{(4.1)\,(a)}, \[|\nu(\mathcal O_K(g_i))-c_i^\infty|\le\liminf e^{c_{\rm vol}}\bigl(G_1\|g_i\|_\infty+|c_i^\infty-\hat c_{K,N_k}(g_i)|\bigr)=e^{c_{\rm vol}}G_1\|g_i\|_\infty.\] Next let $F_t:=\prod_{I_t}(\mathcal O_K(g_i)-\nu(\mathcal O_K(g_i)))$; it belongs to $\mathcal A(\square_t)\cap\mathcal C_{\rm loc}$ and does not depend on $N_k$. By \eqref{support:55-5} for $\rho$ and \textup{(4.1)\,(a)} applied with $c_i=\nu(\mathcal O_K(g_i))$, \[|\rho(\prod_tF_t)|\le\liminf\prod_te^{c_{\rm vol}}\prod_{I_t}\bigl(G_1\|g_i\|_\infty+|\nu(\mathcal O_K(g_i))-\hat c_{K,N_k}(g_i)|\bigr).\] This limit exists, and by the bound for one cell it is at most $\prod_te^{c_{\rm vol}}\prod_{I_t}G_1(1+e^{c_{\rm vol}})\|g_i\|_\infty$, where $G_1=G_1(\mathsf b',\check\delta_t,\check\delta_t/4)$. Only the copies of the cell $t$ occur in its fold; so the data may differ from cell to cell.

\emph{\textup{(a)} on a board $(i,o)$.} We repeat the argument for $B_0$, with the following changes. The means are bounded in $N$ by the first inequality above, and we take a sequence $N_\bullet$ of sides of the data with $m_{N_k}(g_i)\to m_i^*$. Each use of \textup{(4.1)\,(a)} is replaced by a use of \textup{(4.2)}, so that the factor $e^{c_{\rm vol}}$ of a cell becomes $e^{c_{\rm vol}\mathsf l^3}$. For one cell, with $F=\mathcal O_K(g_i)-m_i^*$, the bound becomes $|\nu(\mathcal O_K(g_i))-m_i^*|\le e^{c_{\rm vol}\mathsf l^3}\check\Gamma\|g_i\|_\infty$. For the product over the cells of the board $(i,o)$, each factor of the final bound becomes $\check\Gamma\|g_i\|_\infty+e^{c_{\rm vol}\mathsf l^3}\check\Gamma\|g_i\|_\infty$.

\emph{\textup{(b)}: the table.} Let $\check\delta:=\check\delta_n(d)\le1/(4n)$. A Euclidean distance of at least $d$ gives a sup distance of at least $d/\sqrt3\ge\check\delta$.

\emph{\textup{(b)}: cut.} We cut each $f_i$ sharply along the half-open cubes of side $\check\delta/2$ of one grid. A Euclidean ball of radius $R$ lies in a cube of side $2R$; so $f_i$ falls into at most $(4R/\check\delta+2)^3$ pieces, each lying in a sup-ball of radius $\check\delta/4$. Then $\mathcal O_K(f_i)$ is exactly the sum of the $\mathcal O_K$ of its pieces, and by multilinearity the moment is a sum over the choices of one piece from each $f_i$. Pieces of different $f_i$ are at least $\check\delta$ apart, and two pieces of one and the same $f_i$ never occur in one term.

The argument also applies to tuples of sup type $(d;R)$: their
separation is at least $d\ge\check\delta$, and each sup-ball of
radius $R$ is a cube of side $2R$. Thus the number of pieces and
every subsequent factor are unchanged.

\emph{\textup{(b)}, case $K\ge j^\circ(\mathsf b',\check\delta)$: placing.} In each coordinate, a piece forbids an arc of length $\check\delta/2+\check\delta$ on the circle $\mathbb R/\mathbb Z$ of shifts. The $n$ pieces together forbid at most $3n\,\check\delta/2\le3/8$, and the complement consists of at most $n$ arcs of total length at least $5/8$; so one of these arcs has length at least $5/(8n)\ge(5/2)\check\delta$. From $\mathrm{sep}^{[j^\circ-1]}\ge\mathrm{sep}^\circ$ we obtain $\check\delta L^{j^\circ-1}/M\ge\mathrm{sep}^\circ+2$; hence $a_K\le L^{-j^\circ}<\check\delta$, and the arc contains a point of $a_K\mathbb Z$. Thus some $v\in a_K\mathbb Z^3$ moves all $n$ pieces to distance at least $\check\delta/2$ from $\mathbb Z+\frac12$, each into one cell. As $\rho$ and $\nu$ are both invariant under $\tau_v$, the moment of the translated pieces is the same, and \textup{(a)}, with $q\le n$ cells, bounds it by $(e^{c_{\rm vol}}\check\Gamma(\mathsf b',\check\delta))^n$ times the product of the sup norms of the $n$ pieces.

\emph{\textup{(b)}, case $K<j^\circ(\mathsf b',\check\delta)$.} Here $0\le a_p\le2$, and a half-open cube of side $\check\delta/2$ contains at most $3(\check\delta L^K/2+1)^3$ barycentres. So, for a piece $g$, $|\mathcal O_K(g)|\le6L^K(\check\delta L^K/2+2)^3\|g\|_\infty\le B_0(\mathsf b',\check\delta)\|g\|_\infty$, and $|\mathcal O_K(g)-\nu(\mathcal O_K(g))|\le2B_0(\mathsf b',\check\delta)\|g\|_\infty$ under any law.

\emph{\textup{(b)}: sum.} By line~8 of \textup{(G1)}, in both cases each choice of pieces contributes at most $(A_4\check\delta^{-2})^n$ times the product of the sup norms. Moreover $(4R/\check\delta+2)^3\le64(1+R)^3\check\delta^{-3}$, $\check\delta^{-5n}\le3^{5n/2}(4n)^{5n}\max\{1,d^{-5n}\}$ and $n^n\le e^nn!$. Combining these, we bound the moment by $A_6^n(n!)^5\max\{1,d^{-5n}\}(1+R)^{3n}$ times the product of the sup norms of the $f_i$.

\emph{\textup{(b)}: cumulants, $n\ge2$.} By \eqref{support:9-13} with $\sigma_\rho((f))=0$, \[|S^T_n|\le\sum_\pi(|\pi|-1)!\prod_J\text{bound}(|J|),\] where $\text{bound}(|J|)$ denotes the bound just found, for the $|J|$ entries of a block $J$. Here $\prod_J(|J|!)^5\le(n!)^5$, and \[\sum_\pi(|\pi|-1)!\le F_n:=\sum_kk!\,S(n,k)\le0.52\cdot1.45^{n+1}n!,\] with $S(n,k)$ the Stirling numbers of the second kind. The last inequality follows from $F_n=(n!/2)\sum_{k\in\mathbb Z}(\log2+2\pi ik)^{-n-1}$. In $F_n/n!$ the term $k=0$ is $\frac12(1/\log2)^{n+1}\le\frac12\cdot1.45^{n+1}$, as $1/\log2\le1.4427$. For $n\ge1$ the moduli of the terms $k\ne0$ add up to \[\frac12\sum_{k\ne0}|\log2+2\pi ik|^{-n-1}\le\sum_{k\ge1}(2\pi k)^{-n-1}\le(\pi^2/6)(2\pi)^{-n-1}\le0.02\cdot1.45^{n+1},\] because $\sum_{k\ge1}k^{-n-1}\le\pi^2/6$. Hence $F_n/n!\le0.52\cdot1.45^{n+1}$.

\emph{\textup{(c)}: compactness.} Let $\mathcal U$ denote the countable set of all tuples $(O_1,\dots,O_n)$ whose members are finite unions of open rational cubes and have pairwise disjoint closures. Each $C^1_c(O_i)$ is separable in the sup norm, and we fix countable dense sets $D(O_i)$. The values $\sigma_{\rho_j}(X)$ with $X\in\prod D(O_i)$ are bounded by \textup{(b)}, the numbers being suprema over $U(\mathsf b)$; hence a diagonal subsequence converges at all of them. A separated tuple $X$ has $\operatorname{supp}f_i\subset O_i$ for some member of $\mathcal U$. On $\prod C^1_c(O_i)$ the type $(d;R)$ is fixed, and multilinearity together with \textup{(b)} gives, uniformly in $j$, \[|\sigma_{\rho_j}(X)-\sigma_{\rho_j}(X')|\le C\sum_i\|f_i-f'_i\|_\infty\prod_{k\ne i}\max\{\|f_k\|_\infty,\|f'_k\|_\infty\}.\] Approximating $X$ by tuples $X'$ from $\prod D(O_i)$, we see that $\sigma_{\rho_j}(X)$ is a Cauchy sequence. By \eqref{support:9-13} and \eqref{support:9-14}, the convergence of all $\sigma$ amounts to the convergence of all $S^T_n$ with $n\ge2$. Finally, $\mathbb G_K(\mathsf b_j)\supset\mathcal G^{\rm tor}$ is not empty (\Cref{support:14-16}\,(c)).

\emph{Remark.} Let $S\in\mathrm{Lim}_G(\mathsf b)$, with labels $\mathsf b_j$ in normal form, and let $X$ be a separated $C^1_c$ tuple of type $(d;R)$ with $n$ entries $f_i$. For all $j$ from some $j$ on, namely those with $\mathsf b_j\in U(\mathsf b)$, \textup{(4.3)\,(b)} applied with $\nu=\rho=\rho_j$ bounds $|\sigma_{\rho_j}(X)|$ by the table $C^\sigma_n(d;R)=A_6^n(n!)^5\max\{1,d^{-5n}\}(1+R)^{3n}$ times the product of the sup norms of the $f_i$, because $A_6(\mathsf b_j)\le A_6(\mathsf b)$ for $\mathsf b_j\in U(\mathsf b)$. As $\sigma_S(X)=\lim_j\sigma_{\rho_j}(X)$ on separated $C^1_c$ tuples (Step~4), the bound passes to $\sigma_S(X)$ in the limit, and since $\|\cdot\|_\infty\le\|\cdot\|_\sharp$ it holds also with the norm $\|\cdot\|_\sharp$. By its form the table does not increase in $d$; and for $n=1$ we have $\sigma_S((f))=0$, which is \textup{(R0)}. Altogether, these facts give \textup{(R2)} of \textup{(4.6)}.

\emph{Step 4: symmetries and mirrors.} We prove \textup{(R0)}, \textup{(R1)}, \textup{(R3)}, \textup{(R6)} and \textup{(R4)} of \textup{(4.6)}. Thus, for $S\in\mathrm{Lim}_G(\mathsf b)$, we show: $\sigma_S$ is a functional, centred and symmetric; it is invariant under $\mathbb R^3\rtimes W_3$; it is reflection positive (\Cref{support:9-25}\,(3)) in each coordinate plane and in each plane of the six diagonal families; and $H\ge|P_1|$.

\emph{Functional; \textup{(R0)}.} Each $S_n$ is a limit of real, symmetric and multilinear maps; hence $\sigma_S$ is a functional (\Cref{support:9-12}\,(3), (4)), and on every separated tuple $\sigma_S=\lim\sigma_{\rho_j}$. As $\mathrm{Part}_{\ge2}$ of a set with one element is empty, $\sigma_S((f))=0$.

\emph{\textup{(R1)}: $W_3$.} Each $\rho_j$ is invariant under $\mathsf{Iso}(\mathbb L_{K_j})=a_{K_j}\mathbb Z^3\rtimes W_3$ (\Cref{support:54-17}\,(b)), and $r\mathcal O_K(f)=\mathcal O_K(f\circ r^{-1})$ (Step~1). Hence $r(X^\rho)=(rX)^\rho$, and $\sigma_{\rho_j}(rX)=\sigma_{\rho_j}(X)$ holds exactly, for $r\in\mathsf{Iso}(\mathbb L_{K_j})$.

\emph{\textup{(R1)}: translations (equicontinuity).} For real $v$ we pick $v_j\in a_{K_j}\mathbb Z^3$ with $v_j\to v$, and we replace $f_i(\cdot-v)$ by $f_i(\cdot-v_j)$, one entry at a time. The difference has sup norm at most $\operatorname{Lip}f_i\,|v-v_j|$, and for large $j$ the tuples are of the type $(d/2;2R)$. So \textup{(4.3)\,(b)} gives \[|\sigma_{\rho_j}(\tau_vX)-\sigma_{\rho_j}(\tau_{v_j}X)|\le C|v-v_j|\sum_i\operatorname{Lip}f_i\prod_{k\ne i}\|f_k\|_\infty\to0,\] while $\sigma_{\rho_j}(\tau_{v_j}X)=\sigma_{\rho_j}(X)$; in the limit $\sigma_S(\tau_vX)=\sigma_S(X)$. The same estimate holds if only some of the entries are moved, provided the tuples stay separated by a fixed amount.

\emph{Mirrors.} Let $n=e_3$ or $n=e_1+e_3$, and $c=0\in\operatorname{Adm}(n)$; here $\mathbb E_+$ is the open side of the mirror $(n,0)$ at hand. Let $X_1,\dots,X_r\in\mathfrak S(\mathbb E_+)$, let $\zeta_a$ be complex numbers, and put $F_j:=\sum_a\zeta_aX_a^{\rho_j}$. The supports are compact and lie in the open side, so $n\cdot x\ge\delta>0$ on them. A point of a plaquette differs from the barycentre by at most $a_K/2$ in each of two coordinates; hence, for large $j$, $n\cdot x\ge\delta-a_K\ge0$ on its bonds, and $F_j\in\mathcal F_{\rm loc}(\mathbb H^+(n,0))$. \Cref{support:54-17}\,(a) gives $\rho_j\bigl(\overline{\theta_{n,0}F_j}\,F_j\bigr)\ge0$. Here $\theta_{e_1+e_3,0}x=x-(x_1+x_3)(e_1+e_3)=(-x_3,x_2,-x_1)=\theta_\Sigma x$ (\Cref{support:32-25}\,\textup{(R6)}), and as $\theta_{n,0}\in\mathsf{Iso}(\mathbb L_K)$, we have $\theta(X^\rho)=(\theta X)^\rho$. The left side is therefore \[\sum\bar\zeta_a\zeta_b\,\sigma_{\rho_j}(\theta X_a\cdot X_b)\to\sum\bar\zeta_a\zeta_b\,\sigma_S(\theta X_a\cdot X_b),\] so that the limit is $\ge0$. This is \textup{(R3)} and the full positivity at $\Sigma$; \textup{(R6)} is the case of tuples with one entry. With \textup{(R1)} and \Cref{support:9-31}\,(d), (a), we reach every coordinate mirror and every diagonal mirror, on both sides.

\emph{\textup{(R4)}.} By \Cref{support:32-37}\,(1), \textup{(R1)} and \textup{(R2)} place $\sigma_S$ in some $\mathsf{Fct}(\underline C,\bar{\mathfrak c})$. \Cref{support:28-21} applied with $k=1$, $\varepsilon=1$, with $e_+=(e_3+e_1)/\sqrt2$, gives $H_\sigma\ge|P_{1,\sigma}|$, and by \Cref{support:32-37}\,(2) this is \textup{(R4)}.

\emph{Step 5: the insertion bound.} On the lattice we prove \textup{(4.4)}. On the continuum we prove \textup{(R5)} in the following form. Let $\eta\in(0,1]$; let $h$ be a real $C^1_c$ function with $\operatorname{supp}h\subset B(y,\eta/(48L))$, where $|y_3|\le\eta/8$; and let $F,G\in\mathbb C\mathfrak S(\{x_3>\eta/4\})$. Then \[|\sigma_S(\theta\bar F\cdot(h)\cdot G)|\le C_{\rm I}\eta^{-1}\|h\|_\sharp\,\|[\tau_{-(\eta/4)e_3}F]\|\,\|[\tau_{-(\eta/4)e_3}G]\|.\]

\emph{Lattice.} We fix a label $\mathsf b'\in U(\mathsf b)$, and we take \textup{(4.2)} and the data at $\mathsf b'$; at the end we use $X(\mathsf b')\le X(\mathsf b)$, the supremum over $U(\mathsf b)$. The state $\rho$ has \textup{(Inv)} and \textup{(Pos)} by \Cref{support:54-17}(a), (b), since $c\in a_K(\mathbb Z+\frac12)\subset\operatorname{Adm}(e_3)$, and it has \textup{(Dom)} by \Cref{support:55-40}(a). We choose a sequence $N_\bullet$ of sides of the data with $m_{N_k}(f)\to m^*$, and we put $D:=\mathcal O_K(f)-\rho(\mathcal O_K(f))$, which lies in $\mathcal A(\square)\cap\mathcal C_{\rm loc}$ (\Cref{support:18-7}\,(b)). The single source $f$ has the data of \textup{(4.2)} with $r=1$ and $\check\delta=2\mathsf l/3$: the set $\{f\ne0\}$ is at distance at least $\mathsf l/3=\check\delta/2$ from the faces, and it lies in a sup-ball of radius $\mathsf l/6=\check\delta/4$. By \Cref{support:18-18}\,(c), \[\|D\|_{{\rm ch},N_k}\le\|\mathcal O_K(f)-m^*\|_{{\rm ch},N_k}+|m^*-\rho(\mathcal O_K(f))|,\] and by \textup{(4.2)}, \[\|\mathcal O_K(f)-m^*\|_{{\rm ch},N_k}\le e^{c_{\rm vol}(\mathsf b')\mathsf l^3}\bigl(\check\Gamma(\mathsf b',2\mathsf l/3)\|f\|_\infty+|m^*-m_{N_k}(f)|\bigr)\to\Gamma(\mathsf b',\mathsf l)\|f\|_\infty.\] Also, \eqref{support:55-5} for the one cell gives $|\rho(\mathcal O_K(f))-m^*|\le\Gamma(\mathsf b',\mathsf l)\|f\|_\infty$. So $\|D\|_{{\rm ch},N_\bullet}\le2\,\Gamma(\mathsf b',\mathsf l)\,\|f\|_\infty$, and \textup{(4.4)} follows by \Cref{support:55-9}, with the number $2\Gamma(\mathsf b',\mathsf l)$. At the end of the step we pass to $\Gamma(\mathsf b,\mathsf l)$, since $\Gamma(\mathsf b',\mathsf l)\le\Gamma(\mathsf b,\mathsf l)$ by $X(\mathsf b')\le X(\mathsf b)$.

\emph{Geometry.} The half-open intervals $(x/L,x]$ with $x=L^k\eta/8$, $k\in\mathbb Z$, tile $(0,\infty)$, and $1$ lies in exactly one of them; so exactly one $i$ has $\mathsf l=L^{-i}\in(\eta/(8L),\eta/8]$, and $i\ge1$, as $\mathsf l\le1/8$. Then $c_{\rm I}\eta=\eta/(48L)<\mathsf l/6$, strictly. We take $j$ so large that $K=K_j\ge\max\{i+1,j^\circ_U(\mathsf b,2\mathsf l/3)\}$, $a_K/2\le\mathsf l/6-c_{\rm I}\eta$ and $a_K<\eta/8$. We pick $o\in a_K\mathbb Z^3$ with $|o_\mu-y_\mu|\le a_K/2$, and we let $\square=o+[-\mathsf l/2,\mathsf l/2]^3$ and $c:=o_3-\mathsf l/2$; then $c\in a_K(\mathbb Z+\frac12)$, because $L^{K-i}$ is odd. As a Euclidean ball lies in the sup-ball of the same radius, $\operatorname{supp}h$ is at distance at least $\mathsf l/2-a_K/2-c_{\rm I}\eta\ge\mathsf l/3$ from the faces and lies in a sup-ball of radius $c_{\rm I}\eta\le\mathsf l/6$. Finally $c\ge-3\eta/16-a_K/2>-\eta/4$ and $c+\mathsf l\le3\eta/16+a_K/2<\eta/4$.

\emph{The two sides, on the lattice.} By the rules of \Cref{support:55-1}\,(2)(d), $\theta_c\theta_0=\tau_{2ce_3}$, and $\theta_c\tau_v=\tau_{-v}\theta_c$ for $v$ parallel to $e_3$; moreover $2c$, $\mathsf l$ and $2c+2\mathsf l$ lie in $a_K\mathbb Z$. Let $F^\rho$ be the linear extension of $X\mapsto X^\rho$. It commutes with $\mathsf{Iso}(\mathbb L_K)$, and $\overline{F^\rho}=(\bar F)^\rho$, the $X^\rho$ being real. With $D$ defined as above from the source $h$, \[\sigma_\rho(\theta\bar F\cdot(h)\cdot G)=\rho\bigl(\overline{\theta_0F^\rho}\,D\,G^\rho\bigr).\] Put $F':=\tau_{2ce_3}F^\rho$ and $F'':=\tau_{-\mathsf l e_3}G^\rho$; then $\theta_cF'=\theta_0F^\rho$ and $\tau_{\mathsf l e_3}F''=G^\rho$. The entries of $F'$ lie in $x_3>\eta/4+2c>c$, and those of $F''$ in $x_3>\eta/4-\mathsf l>c$. All bonds of a plaquette lie in $\mathbb H^+_c$ as soon as its barycentre lies strictly above the bond mirror $c$: the height $x_3$ of the barycentre either lies in $a_K\mathbb Z$, and is then at least $c+a_K/2$, the plaquette being horizontal; or it lies in $a_K(\mathbb Z+\frac12)$, and is then at least $c+a_K$. Hence $F',F''\in\mathcal C_{\rm loc}(\mathbb H^+_c)$, and \textup{(4.4)} applies. Its two norms are $\|F'\|^2_{\rho,c}=\rho\bigl(\overline{\theta_0F^\rho}\,\tau_{2ce_3}F^\rho\bigr)=\sigma_\rho(\theta\bar F\cdot\tau_{2ce_3}F)$ and, since $\theta_c\tau_{-\mathsf l e_3}=\tau_{(2c+\mathsf l)e_3}\theta_0$, $\|F''\|^2_{\rho,c}=\sigma_\rho(\tau_{(2c+\mathsf l)e_3}\theta\bar G\cdot\tau_{-\mathsf l e_3}G)$. A translation by the lattice vector $-(2c+\mathsf l)e_3$ turns the latter into $\sigma_\rho(\theta\bar G\cdot\tau_{-2(c+\mathsf l)e_3}G)$. The gaps of these two tuples are $\eta/2+2c$ and $\eta/2-2(c+\mathsf l)$, and both are at least $\eta/8-a_K>0$. We stress that no translation by $(\eta/4)e_3$ is made on the lattice: the vector $(\eta/4)e_3$ need not lie in $a_K\mathbb Z^3$, and $\sigma_\rho$ is invariant under $a_K\mathbb Z^3$ alone.

\emph{Limit.} The index $i$ is fixed by $\eta$, and $c_j=o_{3,j}-\mathsf l/2\to c_\infty:=y_3-\mathsf l/2$, because $|o_{3,j}-y_3|\le a_{K_j}/2$. From $|y_3|\le\eta/8$ and $\mathsf l\le\eta/8$ we get $-3\eta/16\le c_\infty$ and $c_\infty+\mathsf l\le3\eta/16$. The left side tends to $\sigma_S(\theta\bar F\cdot(h)\cdot G)$, since the tuples are fixed. By the equicontinuity of Step~4, with a common positive separation given by the minimum of $\eta/16$ and the internal separations of the finitely many tuples occurring in $F$ and $G$, we have $\sigma_{\rho_j}(\theta\bar F\cdot\tau_{2c_je_3}F)\to\sigma_S(\theta\bar F\cdot\tau_{2c_\infty e_3}F)$; in the same way $\sigma_{\rho_j}(\theta\bar G\cdot\tau_{-2(c_j+\mathsf l)e_3}G)\to\sigma_S(\theta\bar G\cdot\tau_{-2(c_\infty+\mathsf l)e_3}G)$.

\emph{After the limit, by \textup{(R1)} for $\sigma_S$.} Let $F_0:=\tau_{-(\eta/4)e_3}F\in\mathbb C\mathfrak S(\mathbb E_+)$; then $\theta\bar F=\tau_{-(\eta/4)e_3}\theta\bar F_0$. A translation by $+(\eta/4)e_3$ gives \[\sigma_S(\theta\bar F\cdot\tau_{2c_\infty e_3}F)=\sigma_S(\theta\bar F_0\cdot\tau_{(2c_\infty+\eta/2)e_3}F_0)=\langle[F_0],\mathcal T_{2c_\infty+\eta/2}[F_0]\rangle\le\|[F_0]\|^2.\] In the last step we use $2c_\infty+\eta/2\ge\eta/8\ge0$ and $0\le\mathcal T_s\le1$ (\Cref{support:32-16}\,(b), which applies as $\sigma_S$ has \textup{(R1)} to \textup{(R3)}). In the same manner, with $G_0:=\tau_{-(\eta/4)e_3}G$, the translation by $+(\eta/4)e_3$ gives \[\sigma_S(\theta\bar G\cdot\tau_{-2(c_\infty+\mathsf l)e_3}G)=\langle[G_0],\mathcal T_{\eta/2-2(c_\infty+\mathsf l)}[G_0]\rangle\le\|[G_0]\|^2,\] where now $\eta/2-2(c_\infty+\mathsf l)\ge\eta/8\ge0$.

\emph{Power.} Keep the plain label $\mathsf b'\in U(\mathsf b)$ in the lattice estimate until the complete coefficient has been bounded. Pointwise in that label,
\[
 2\Gamma(\mathsf b',\mathsf l)
 \le\frac{3e^{c_{\rm vol}(\mathsf b')}\hat G(\mathsf b')}{\mathsf l}
 <\frac{24L e^{c_{\rm vol}(\mathsf b')}\hat G(\mathsf b')}{\eta}
 \le\frac{C_{\rm I}(\mathsf b)}{\eta},
\]
where $C_{\rm I}(\mathsf b)$ is the supremum over $U(\mathsf b)$ of the whole plain-label expression $24L e^{c_{\rm vol}(\mathsf b')}\hat G(\mathsf b')$. Thus this last coefficient bounds every member of the sequence before taking the two limits above. Since $\|h\|_\infty\le\|h\|_\sharp$, \textup{(4.4)} and those limits give the bound stated at the beginning of the step.

\emph{Step 6: plates, their growth exponent $7/2<4$.} We prove \textup{(4.5)} and \textup{(R7)}. Here $\mathsf b'$, $\rho$, $t$, $\alpha$, $\varphi$, $f$ and $K$ are as in clause \textup{(4.5)}. From \textup{(4.5)} it follows that $\sigma_S(\theta f,f)=S_2(\theta f,f)\le C_\varphi t^{-7/2}$, which is \textup{(R7)}.

\emph{One choice for all $t$.} We take $\check\delta:=t/8\le1/8$, on $B_0$.

\emph{Pieces.} The function $f$ vanishes outside the box $\{x_3\in[t,2t],\ x_1\in[-t,t],\ x_2\in$ the interval of $\varphi\}$, and since it is continuous it also vanishes on the boundary of the box. We cut $f$ sharply by half-open cubes of side $\check\delta/2=t/16$, aligned with the box. This yields $16\cdot32\cdot\lceil16\lambda_\varphi/t\rceil\le512(16\lambda_\varphi+1)/t$ pieces $g_p$, with $\|g_p\|_\infty\le3\|\varphi\|_\infty$, each in a sup-ball of radius $\check\delta/4$. Bilinearity, at the lattice level, gives $S^T_2[\rho](\theta f,f)=\sum_{p,q}\sigma_\rho(\theta g_p,g_q)$. Along $x_3$, $\theta g_p$ and $g_q$ are at least $2t\ge\check\delta$ apart; between two pieces on the same side no condition is required.

\emph{Placing.} In each coordinate two pieces forbid at most $3\,\check\delta\le3/8$, so one free arc has length at least $5/16>a_K$. A point of $a_K\mathbb Z$ on the free arc of each coordinate gives, as in Step~3, a common $v\in a_K\mathbb Z^3$ which moves both pieces to distance at least $\check\delta/2$ from $\mathbb Z+\frac12$, each into one cell; and $\sigma_\rho$ is invariant under $\tau_v$.

\emph{Bound.} Retain the plain label $\mathsf b'$ throughout the estimate. The fixed-label proof of \textup{(4.3)\,(a)}, applied with $\nu=\rho$, covers both cases, that of one cell with two separated sources and that of two cells:
\[
 |\sigma_\rho(\theta g_p,g_q)|
 \le e^{2c_{\rm vol}(\mathsf b')}\check\Gamma(\mathsf b',t/8)^2
       \|g_p\|_\infty\|g_q\|_\infty.
\]
Here the numbers at $\mathsf b'$ are their plain values. By \textup{(G1)}, $\check\Gamma(\mathsf b',t/8)^2\le\hat G(\mathsf b')^2(8/t)(1+\log(8/t))$.

\emph{Sum.} Summing over the pairs of pieces gives
\[
 72\cdot512^2(16\lambda_\varphi+1)^2\|\varphi\|_\infty^2
 e^{2c_{\rm vol}(\mathsf b')}\hat G(\mathsf b')^2
 t^{-3}(1+\log(8/t)).
\]
The function $t^{1/2}(1+\log(8/t))$ increases on $(0,1]$, its maximum lying at $t=8/e>1$; so it is at most $1+\log8\le3.0795\le4$. Writing $t^{-3}=t^{-7/2}t^{1/2}$, this is at most $C_\varphi(\mathsf b')t^{-7/2}$, where $C_\varphi(\mathsf b')$ is the plain expression in item~2 with the factor $288=72\cdot4$. Now take the supremum of this whole expression over $\mathsf b'\in U(\mathsf b)$ to obtain $C_\varphi(\mathsf b)t^{-7/2}$.

\emph{Limit.} The pair $(\theta f,f)$ is a separated $C^1_c$ tuple (\Cref{support:32-30}); hence $S^T_2[\rho_j](\theta f,f)\to S_2(\theta f,f)=\sigma_S(\theta f,f)$.

\emph{Remark: the margin.} The argument needs $G_1\le C\,\check\delta^{-p}$ with $2+2p<4$, that is, $p<1$. Line~2 of \textup{(G1)} yields the exponent $p=1/2$, with a logarithm.

\emph{Remark: why \textup{(4.5)} has a threshold in $K$.} The square $B_0(\mathsf b,\check\delta)^2$ is of the order of $\check\delta^{-2}$, and the closed form $d^{-5n}$ gives $t^{-10}$. These bounds are of use for $K<j^\circ$ alone, and not for \textup{(R7)}; hence the threshold in $K$ in \textup{(4.5)}.

\emph{Step 7: \Cref{support:37-34} applied; every mirror; the cone; $\operatorname{O}(2)$; every $\mathsf b>0$.} We prove \textup{(4.7)} and \textup{(4.8)}.

\emph{\textup{(4.7)}: labels in normal form.} By Steps~3 to~6, $\sigma_S$ is a functional with \textup{(R0)} to \textup{(R7)}, where $c_{\rm I}=1/(48L)\in(0,1/16]$, $C_{\rm I}>0$, $\mathsf p_{\rm I}=1\ge0$ and $\mathsf a_\Gamma=7/2<4$. \Cref{support:37-34} gives $\sigma_S(rX)=\sigma_S(X)$ for $r\in\operatorname{O}(3)$, and with \textup{(R1)} we obtain invariance under all of $\mathbb R^3\rtimes\operatorname{O}(3)$. By \eqref{support:9-13} with $\sigma_S((f))=0$, $S_n=S^{\sigma_S}_n$, and that formula commutes with $r$.

\emph{Every mirror.} This is \Cref{support:28-32}\,(c). By \Cref{support:32-37}\,(1), $\sigma_S$ lies in some $\mathsf{Fct}(\underline C,\bar{\mathfrak c})$, and \Cref{support:28-34} then gives $H\ge|P_k|$ for $k=1,2$, \eqref{support:28-33}, $\|\mathcal T(z)\|\le1$, and all of \Cref{support:28-30}.

\emph{\textup{(4.8)}.} The dilation $D_z$, $x\mapsto L^zx$, maps $a_K\mathbb Z^3$ onto $a_{K-z}\mathbb Z^3$ for $K\ge z+1$, and it maps bonds to bonds and plaquettes to plaquettes. As $H_B=-\beta\sum a_p$ does not depend on the spacing (\eqref{support:14-7}), $D_z$ carries the Gibbs states at $(a_K,\beta)$ onto those at $(a_{K-z},\beta)$, with $\beta=L^{K-z}(L^z\mathsf b')$, and it carries $a_K\mathbb Z^3$-invariance to $a_{K-z}\mathbb Z^3$-invariance. Further, \[\mathcal O_K(f)=a_K^{-1}\sum f(x(p))a_p=L^za_{K-z}^{-1}\sum f^{[z]}(L^zx(p))a_p,\] and $X\mapsto X^{[z]}$ is a bijection of the separated $C^1_c$ tuples. Therefore $S'_n(Y):=\lim S^T_n[D_z\rho_j](Y)$ exists on all tuples, $S'\in\mathrm{Lim}_G(L^z\mathsf b)$, and $S_n(X)=L^{zn}S'_n(X^{[z]})$. The dilation $D_z$ carries $\tilde\mu_{K,2L^m,\beta}$ to $\tilde\mu_{K-z,2L^{m+z},\beta}$; since limit points depend on the tails only (\Cref{support:14-9}\,(4)), it carries $\mathcal G^{\rm tor}_{K,\beta}$ onto $\mathcal G^{\rm tor}_{K-z,\beta}$, and $\mathrm{Lim}_{\rm tor}(\mathsf b)$ corresponds to $\mathrm{Lim}_{\rm tor}(L^z\mathsf b)$. Since $L^z\mathsf b_j\to L^z\mathsf b>x_s$, we have $L^z\mathsf b_j\in U(L^z\mathsf b)$ from some $j$ on, so the labels are in normal form. The invariance follows from $(rX)^{[z]}=r(X^{[z]})$ for $r\in\operatorname{O}(3)$ and $(\tau_vX)^{[z]}=\tau_{L^zv}X^{[z]}$.

\emph{Mirrors.} We have $(\theta_{(u,c)}X)^{[z]}=\theta_{(u,L^zc)}X^{[z]}$. Hence, with $\theta'=\theta_{(u,L^zc)}$ and with $n_a$ the number of entries of $X_a$, \[\sum\bar\zeta_a\zeta_b\,\sigma_S(\theta X_a\cdot X_b)=\sum\overline{\zeta_aL^{zn_a}}\,(\zeta_bL^{zn_b})\,\sigma_{S'}(\theta'X_a^{[z]}\cdot X_b^{[z]})\ge0.\] For \textup{(R2)} of $\sigma_S$, the type $(d;R)$ becomes $(L^zd;L^zR)$ and the sup norm of each entry does not change under $f\mapsto f^{[z]}$; the table $L^{zn}C^\sigma_n(L^zd;L^zR)$ does not increase in $d$. Hence \Cref{support:32-37}\,(1) and \Cref{support:28-34} are applicable to $\sigma_S$ itself; \textup{(R4)} is contained in this, and \textup{(R0)}, \textup{(R1)}, \textup{(R3)}, \textup{(R6)} are contained in the invariance and the mirrors. For $z=0$, \textup{(R5)} and \textup{(R7)} are already part of \textup{(4.6)}. For $z\ge1$, apply \Cref{Y-P4-LZ} (Step~8) to $\sigma'=\sigma_{S'}$, which has \textup{(R0)} to \textup{(R7)} by \textup{(4.6)}.

\emph{Step 8: joint limits along large sides, the diagonal mirror $\Sigma$ at which reflection positivity is equivalent to invariance under $\operatorname{O}(3)$, and the dilation lemma, for all $\mathsf b>0$.} We prove \textup{(4.9)\,(i)} and \textup{(ii)}, at every $\mathsf b>0$. First we state the lemma on dilations.

\begin{lemma}[Dilation by $L^z$ carries the hypotheses]\label{Y-P4-LZ}
Let $z\ge1$ be an integer, and put $D_zX:=L^{z|X|}X^{[z]}$. Let $\sigma'$ be a functional with \textup{(R0)} to \textup{(R3)}, for a table $C'_n(d;R)$ that gives the bound of \textup{(R2)} in terms of the sup norms $\|f_i\|_\infty$ of the entries and does not decrease in $R$; with \textup{(R5)}, for the triple $(c_{\rm I},C'_{\rm I},1)$; and with \textup{(R7)}, for the exponent $7/2$ and numbers $C'_\varphi$. Then $\sigma:=\sigma'\circ D_z$ is a functional with \textup{(R0)} to \textup{(R3)}, for the table $L^{zn}C'_n(L^zd;L^zR)$; with \textup{(R5)}, for the triple $(c_{\rm I}L^{-z}/2,\ 2L^zC'_{\rm I},\ 1)$; and with \textup{(R7)}, for the exponent $7/2$ and the numbers $L^{2z}\max\{L^{-7z/2}C'_{\varphi^{[z]}},\ 9\|\varphi\|_\infty^2C'_2(2;(2+\lambda_\varphi)L^z)\}$.
\end{lemma}

\emph{Proof of \textup{(4.9)\,(i)}, at fixed $(K,\mathsf b')$, for any $\mathsf b'>0$.} Let $X^1,X^2,\dots$ be the countable family of Step~3\,(c), and fix $K$ and $\mathsf b'$. We claim that there is an $N^*$ such that for every side $N\ge N^*$ in $\{2L^m\}$ some $\nu\in\mathcal G^{\rm tor}_{K,L^K\mathsf b'}$ satisfies $|\sigma_{K,N}(X^q)-\sigma_\nu(X^q)|\le1/K$ for all $q\le K$. Otherwise some $N_k\to\infty$ violates this for every $\nu$. Since $\mathcal V(\mathbb L_K)$ is compact, a subsequence of $\tilde\mu_{K,N_k,\beta}$ converges weakly to some $\nu\in\mathcal G^{\rm tor}$ (\Cref{support:14-9}\,(4)). The finitely many $\sigma(X^q)$ are polynomials in means of continuous local functions, so they converge along this subsequence, a contradiction. Besides, let $z_0(\mathsf b')$ be the least $z\ge0$ with $L^z\mathsf b'\ge x_s$ (item~\textup{(S7)} of \S\ref{Y-00-1}). We choose $N^*(K,\mathsf b')$ large enough that, whenever $N\ge N^*$, the side $L^zN$ belongs to the data at the label $L^z\mathsf b'$, both for $z=z_0(\mathsf b')$ and for $z=z_0(\mathsf b')+1$. The conditions are $m+z\ge m'$ and $L^zN\ge2Ms(L^{z+1}\mathsf b')$, and they hold for all large $N$.

\emph{\textup{(i)}: the limit, in normal form.} Along sides $N_j\ge N^*(K_j,\mathsf b_j)$ we choose such states $\nu_j$; then $\sigma_{K_j,N_j}(X^q)-\sigma_{\nu_j}(X^q)\to0$ for every $q$. Both families obey sup-norm bounds uniform in $j$: on the tori by \Cref{Y-P3-ThG}\,\textup{(G2)(i)}, and for the $\nu_j$ by \textup{(4.3)\,(b)}. The argument of Step~3\,(c) therefore applies to both: on each product $\prod C^1_c(O_i)$ with $(O_1,\dots,O_n)\in\mathcal U$ the type is fixed, and multilinearity with these bounds yields the estimate of Step~3\,(c) for both families, uniformly in $j$. As the differences tend to $0$ on the countable family, they tend to $0$ on every separated tuple; so the two families have the same limit points on all tuples.

\emph{\textup{(i)}: the limit, outside the normal form.} Let $z=z(\mathsf b)$ (item~\textup{(S7)} of \S\ref{Y-00-1}); for $\mathsf b\le x_s$ it is the least $z\ge1$ with $L^z\mathsf b>x_s$. For large $j$ we have $L^z\mathsf b_j>x_s$ and $L^{z-2}\mathsf b_j<x_s$; thus $z$ is $z_0(\mathsf b_j)$ or $z_0(\mathsf b_j)+1$, and $L^zN_j$ is a side of the data at $L^z\mathsf b_j$. The map $x\mapsto L^zx$ carries $a_K\mathbb Z^3/N\mathbb Z^3$ onto $a_{K-z}\mathbb Z^3/L^zN\mathbb Z^3$, and Wilson's weight does not depend on the spacing. Hence, by the formula for $\mathcal O_K(f)$ in the proof of \textup{(4.8)}, $\sigma_{K,N}(X)$ at the label $\mathsf b'$ equals $L^{z|X|}\sigma_{K-z,L^zN}(X^{[z]})$ at the label $L^z\mathsf b'$. So \Cref{Y-P3-ThG}\,\textup{(G2)(i)}, applied to the torus with cutoff, side and label $(K_j-z,L^zN_j,L^z\mathsf b_j)$, bounds the torus family by $L^{zn}$ times the table at $(L^zd;L^zR)$, with the numbers as suprema over $U(L^z\mathsf b)$; and \textup{(4.8)} bounds $\sigma_{\nu_j}$ in the same way. We now follow the argument of the normal form, with one change: these bounds take the place of \Cref{Y-P3-ThG}\,\textup{(G2)(i)} and of \textup{(4.3)\,(b)}. The choice of the $\nu_j$ along $N_j\ge N^*(K_j,\mathsf b_j)$ stays as it is, the estimate of Step~3\,(c) holds for both families uniformly in $j$, and the two families have the same limit points on all tuples.

\begin{proof}[Proof of \Cref{Y-P4-LZ}]
\emph{Functional.} The map $X\mapsto X^{[z]}$ is a bijection of the separated $C^1_c$ tuples, and it is linear in each entry; hence $\sigma$ is a functional, and $\sigma(\emptyset)=1$.

\emph{Intertwining.} We have $(rX)^{[z]}=r(X^{[z]})$ for $r\in\operatorname{O}(3)$, $(\tau_vX)^{[z]}=\tau_{L^zv}X^{[z]}$, $(\theta X)^{[z]}=\theta X^{[z]}$ and $D_z(X\cdot Y)=D_zX\cdot D_zY$; and we extend $D_z$ linearly to $\mathbb C\mathfrak S$. Properties \textup{(R0)}, \textup{(R1)} and \textup{(R3)} follow from these relations as in the proof of \textup{(4.8)} in Step~7: $\sigma((f))=L^z\sigma'((f^{[z]}))=0$, and $\sigma$ is symmetric, since $D_z$ acts entry by entry; for $r\in W_3$ and real $v$, $\sigma(rX)=L^{z|X|}\sigma'(rX^{[z]})=\sigma(X)$ and $\sigma(\tau_vX)=L^{z|X|}\sigma'(\tau_{L^zv}X^{[z]})=\sigma(X)$; and, with $n_a$ the number of entries of $X_a$, $\sum\bar\zeta_a\zeta_b\,\sigma(\theta X_a\cdot X_b)=\sum\overline{\zeta_aL^{zn_a}}\,(\zeta_bL^{zn_b})\,\sigma'(\theta X_a^{[z]}\cdot X_b^{[z]})\ge0$. For \textup{(R2)}, the type $(d;R)$ becomes $(L^zd;L^zR)$, while the sup norm of each entry does not change under $f\mapsto f^{[z]}$.

\emph{Norms.} For $F\in\mathbb C\mathfrak S(\mathbb E_+)$ we have $\|[F]\|^2_\sigma=\sigma'(\theta\overline{D_zF}\cdot D_zF)=\|[D_zF]\|^2_{\sigma'}$, and $D_z\tau_{-(\eta/4)e_3}F=\tau_{-(L^z\eta/4)e_3}D_zF$.

\emph{\textup{(R5)} at $\eta\le\eta_1:=L^{-z}$, for $\operatorname{supp}h\subset B(y,c_{\rm I}\eta)$.} Put $\tilde\eta:=L^z\eta\le1$ and $\tilde h:=L^zh^{[z]}$. Then $\operatorname{supp}\tilde h\subset B(L^zy,c_{\rm I}\tilde\eta)$, $|L^zy_3|\le\tilde\eta/8$, and $D_zF,D_zG\in\mathbb C\mathfrak S(\{x_3>\tilde\eta/4\})$. Moreover $\sigma(\theta\bar F\cdot(h)\cdot G)=\sigma'(\theta\overline{D_zF}\cdot(\tilde h)\cdot D_zG)$, and \[\|\tilde h\|_\sharp=L^z\max\{\|h\|_\infty,40ML^{-z}\operatorname{Lip}h\}\le L^z\|h\|_\sharp.\] Property \textup{(R5)} of $\sigma'$ at $\tilde\eta$ bounds the modulus by $C'_{\rm I}L^{-z}\eta^{-1}L^z\|h\|_\sharp$ times the two norms; by the relations under Norms, this equals \[C'_{\rm I}\eta^{-1}\|h\|_\sharp\,\|[\tau_{-(\eta/4)e_3}F]\|\,\|[\tau_{-(\eta/4)e_3}G]\|.\]

\emph{\textup{(R5)} at $\eta\in(\eta_1,1]$, for $\operatorname{supp}h\subset B(y,c_{\rm I}\eta_1\eta/2)$.} Put $\eta_0:=\min\{\eta/2,\eta_1\}$; then $c_{\rm I}\eta_1\eta/2\le c_{\rm I}\eta_0$. We translate by $-y_3e_3$, using $\tau_{-y_3e_3}\theta=\theta\tau_{y_3e_3}$: the functions $F_1:=\tau_{y_3e_3}F$ and $G_1:=\tau_{-y_3e_3}G$ live in $\{x_3>\eta/8\}\subset\{x_3>\eta_0/4\}$, and $h_1:=\tau_{-y_3e_3}h$ has its centre at height $0$. The case above, at $\eta_0$, gives the bound \[C'_{\rm I}\eta_0^{-1}\|h\|_\sharp\,\|[\tau_{-(\eta_0/4)e_3}F_1]\|\,\|[\tau_{-(\eta_0/4)e_3}G_1]\|.\] Here $[\tau_{-(\eta_0/4)e_3}F_1]=\mathcal T_s[\tau_{-(\eta/4)e_3}F]$ with $s=\eta/4+y_3-\eta_0/4\ge\eta/8-\eta_0/4\ge0$, and $\mathcal T_s$ is a contraction (\Cref{support:32-16}\,(b), which applies as $\sigma$ has \textup{(R1)} to \textup{(R3)}); the same holds for $G_1$, with $-y_3$ in the place of $y_3$. Moreover $\eta_0^{-1}\le2L^z/\eta$. At $\eta\le\eta_1$ the smaller ball lies in $B(y,c_{\rm I}\eta)$, and $C'_{\rm I}\le2L^zC'_{\rm I}$; so one triple covers all $\eta\in(0,1]$, and $c_{\rm I}L^{-z}/2\in(0,1/16]$.

\emph{\textup{(R7)}; $f=\alpha\otimes\varphi$, $\alpha\in\mathsf{Prof}_t$, $0<t\le1$.} We have $f^{[z]}=\alpha^{[z]}\otimes\varphi^{[z]}$. The function $\alpha^{[z]}$ has its support in $[L^zt,2L^zt]\times[-L^zt,L^zt]$, satisfies $|\alpha^{[z]}|\le3$, and has $\operatorname{Lip}\alpha^{[z]}\le12/(L^zt)$; so $\alpha^{[z]}\in\mathsf{Prof}_{L^zt}$, and $\varphi^{[z]}$ is a plate label. Moreover $\sigma(\theta f,f)=L^{2z}\sigma'(\theta f^{[z]},f^{[z]})$. For $L^zt\le1$, \textup{(R7)} of $\sigma'$ bounds this by $L^{2z}C'_{\varphi^{[z]}}(L^zt)^{-7/2}$. For $1<L^zt\le L^z$, the pair is of the type $(2L^zt;R)$, where $2L^zt>2$ and $R$ is half the diagonal of the box: $R=(L^z/2)(5t^2+\lambda_\varphi^2)^{1/2}\le(2+\lambda_\varphi)L^z$. Then \textup{(R2)} of $\sigma'$ for $n=2$, whose table does not rise with $d$, gives the bound $L^{2z}C'_2(2;(2+\lambda_\varphi)L^z)\,9\|\varphi\|_\infty^2$, which is a constant; and $1\le t^{-7/2}$.
\end{proof}

\emph{Proof of \textup{(4.9)\,(ii)}, in normal form.} Properties \textup{(R0)} to \textup{(R3)} come from \Cref{Y-P2-ThR}\,(R6)\,(i) to (iv). The table comes from \Cref{Y-P3-ThG}\,(G2)\,(i), (iii) for $n\ge2$; for the Euclidean type, this is the $C^\sigma_n$ of \textup{(4.6)}, which is monotone in $d$ by its form. For $n=1$, $\sigma_S((f))=0$.

For \textup{(R7)} we follow the proof of \textup{(4.5)} in Step~6, now on the tori, with the following changes. In Placing, the invariance of $\mu_{K,N,\beta}$ under $a_K\mathbb Z^3$ is used. In Bound, \Cref{support:18-18}\,(b) takes the place of \eqref{support:55-5}, on which \textup{(4.3)\,(a)} rests; and \textup{(4.1)\,(a)} applied with $c_i=m_N(g_i)$, together with \textup{(4.1)\,(c)}, gives $|\sigma_{K,N}(\theta g_p,g_q)|\le e^{2c_{\rm vol}}\check\Gamma^2\|g_p\|_\infty\|g_q\|_\infty$ for $K\ge j^\circ(\mathsf b',t/8)$ at the label $\mathsf b'$ of the torus, hence for $K\ge j^\circ_U(\mathsf b,t/8)$. In Limit, the bound passes to every limit.

\emph{\textup{(ii)}, in normal form: \textup{(R5)} from \Cref{Y-P3-ThG}\,(G6).} Let $\eta_0,h,y,F,G$ be as in \textup{(R5)}, with $c_{\rm I}=1/(48L)\le1/32$. We translate by $-y_3e_3$, using $\tau_{-y_3e_3}\theta=\theta\tau_{y_3e_3}$, and obtain $\sigma(\theta\bar F\cdot(h)\cdot G)=\sigma(\theta\bar F_1\cdot(h_1)\cdot G_1)$, where $F_1=\tau_{y_3e_3}F$ and $G_1=\tau_{-y_3e_3}G$ live in $\{x_3>\eta_0/8\}$, and $h_1=\tau_{-y_3e_3}h$ lives in $\{|x_3|<\eta_0/32\}$. Now \textup{(G6)} at $\eta=\eta_0/8$ gives the bound \[\check\Gamma_*(\mathsf b,\eta)\,|\!|\!|h_1|\!|\!|_\eta\,\|[\tau_{-\eta e_3}F_1]\|\,\|[\tau_{-\eta e_3}G_1]\|.\] Here $[\tau_{-\eta e_3}F_1]=\mathcal T_{y_3+\eta_0/8}[\tau_{-(\eta_0/4)e_3}F]$ with $y_3+\eta_0/8\ge0$, and the operator $\mathcal T_{y_3+\eta_0/8}$ is a contraction; the same holds for $G$, with $-y_3$. Let $\eta'\in[\eta/2,\eta]$ be as in item~1 of \S\ref{Y-P3-1}, taken at $\eta=\eta_0/8$. In each coordinate, at most $256c_{\rm I}+3<4$ points of the grid $(\eta'/8)\mathbb Z$ lie at distance at most $\eta'/8$ from the image of the support of $h_1$ under the projection onto that coordinate axis. The sum that defines $|\!|\!|h_1|\!|\!|_\eta$ runs over $(\eta'/8)\mathbb Z^3$. Its term at a point vanishes unless each coordinate of that point is one of these grid points; hence $|\!|\!|h_1|\!|\!|_\eta\le4^3\|h\|_\infty$. At each plain label $\mathsf b'\in U(\mathsf b)$, the first bullet of \textup{(G6)} gives
\[
 \check\Gamma_*(\mathsf b',\eta)
 \le\frac{32e^{c_{\rm vol}(\mathsf b')}\hat G(\mathsf b')}{\eta_0}.
\]
For moving labels use the upper limit prescribed by \textup{(Lab3)}, and then bound it by the supremum over $U(\mathsf b)$ of this whole expression. Combining with the factor $4^3$ just obtained gives
\[
 \frac{2048}{24L}\,C_{\rm I}(\mathsf b)\eta_0^{-1}
       \|h\|_\infty
\]
times the two norms. Since $2048\le24L$, the triple $(c_{\rm I},C_{\rm I},\mathsf p_{\rm I})$ of item~2 of \S\ref{Y-P4-1} suffices.

\emph{\textup{(ii)}, outside the normal form.} Let $z\ge1$ with $L^z\mathsf b>x_s$. \Cref{Y-P2-ThR}\,(R6)\,(vi), at moving labels, gives $S_n(X)=L^{zn}S'_n(X^{[z]})$. At finite $(K,N)$ the identity is the one found in \textup{(i)} above. Here $S'$ is a limit of \Cref{Y-P2-ThR}\,(R5) along the matched subsequence, and its two limits are taken in the same order as for $S$, namely in one of the orders \textup{(a)}, \textup{(b)} and \textup{(c)} of that clause. The labels of $S'$ are $L^z\mathsf b_j\to L^z\mathsf b$, and they are in normal form. Hence $\sigma_S=\sigma_{S'}\circ D_z$, by \eqref{support:9-14}. With the numbers of the normal form at $L^z\mathsf b$, \Cref{Y-P4-LZ} now gives \textup{(R0)} to \textup{(R3)}, \textup{(R5)} and \textup{(R7)} for $\sigma_S$, with the numbers stated in \textup{(4.9)\,(ii)}.

\emph{Equivalence, at every $\mathsf b>0$.} This part uses only that $\sigma_S$ has \textup{(R0)} to \textup{(R3)}, \textup{(R5)} and \textup{(R7)}. Positivity at $\Sigma$ gives \textup{(R6)}, which is this positivity for tuples with one entry. It also gives \textup{(R4)}, by \Cref{support:28-21} with \Cref{support:32-37}. Then \textup{(R4)} and \textup{(R6)}, together with \textup{(R0)} to \textup{(R3)}, \textup{(R5)} and \textup{(R7)}, give invariance of $\sigma_S$ under $\operatorname{O}(3)$ by \Cref{support:37-34}. Conversely, invariance under $\operatorname{O}(3)$ gives positivity at $\Sigma$ (\Cref{support:28-32}\,(c)).
\end{proof}

\section{A reflected test for nontriviality}\label{Y-P6-1}

In this section we prove clause \textup{(Y7)} of \Cref{Y-00-Y}, the bound $\mathsf b^2S^T_2(\theta f_*,f_*)\ge c_*>0$, in which $\theta$ is a reflection and $c_*$ a pure number: first on every torus of the data and for every cutoff $K\ge j+1$, then for the limits and for every label $\mathsf b>0$. The witness is a function $G_-$ of the drawn block field $V^{(k)}$, and it lives on one side of a mirror; by reflection positivity, $\mathrm{Cov}^2\le\mathrm{Var}\cdot S^T_2$. We bound the covariance from below by integration by parts in $\nu_k$ and by the free Green function, and the variance from above. All conditions on the depth $j$ are collected in one list, and we give a depth $j_*(\mathsf b)$ in closed form that is bounded on compact sets.

\emph{Item 0: what is used, letters.}

\emph{Item 0.1: what is used of Sections~\ref{Y-G0-1} to \ref{Y-H9-1}.} From item~1 of \S\ref{Y-G0-1}, ``Cubes'', we take the counts $(26e)^{n-1}$ and $(52e)^n$, together with the meaning that is given there to ``plaquettes of $A$'', to ``$p$ in $A$'' and to ``tame there''. From the same item, ``At $x$'', we take $p_0$, $\varepsilon$, $\rho(x)$ (see item~0.4), $r(x)$, $s(x)=8r(x)$, $\ell_w(x)$, $c_a$ and $c_h$, and from its parts ``Fields'' and ``Domains'' the conditions \textup{(u1)} to \textup{(u3)}. The members of its part ``Members'' are objects of the same kinds as those of item~3 of \S\ref{Y-H-1}, to which item~0 of \S\ref{Y-H1-1}, ``Pairs'', refers: they are the same triples, with the same functions, domains and invariance. From ``Members'' we take \textup{(ca)}, \textup{(ch)} and \textup{(c$\Xi$)}, the strong cap, the formula for $\rho[u;\ell]$, and $\Xi_w$ together with its covering clause, by which, in the sum that defines $\Xi_w$, each set $X_i$ holds a plaquette of $D(V)$ and $D(V)$ lies in the union of the $X_i$. Item~1 of \S\ref{Y-G0-1} also gives \textup{(SD)}, and item~3 of \S\ref{Y-G0-1}, \textup{(O3)}, gives $\mathsf N_Z=\mathsf N_{\mathsf B}=1+3K_*^2\mathsf g_*^2$.

From item~2 of \S\ref{Y-J-1} we take ``Energy'', ``The step'' and ``Runs''. From item~3 of \S\ref{Y-J-1}, \textup{(R9)}, we take \textup{(9a)} and \textup{(9e)}; they are \Cref{Y-H-C9}\,(a), (e).

Of the floor \textup{(W0)} of \Cref{Y-E-C6} we use one clause alone, $\rho_0\le\rho(x)+1$. Here it enters as the entry \textup{(j7) (vii)} of the list in \Cref{Y-P6-N}, rather than as a floor.

From \S\ref{Y-A-1} we take the following: the bounds $\mathsf g\le51\lambda^{-3}$ and $\mathsf g_1\le182\lambda^{-4}$, and the count of $24r^2+2$ points at distance $r$ (\Cref{Y-A-Sums}, after the display of $\mathsf g_1,\mathsf g_2$); the covariant rows (\Cref{Y-A-CovRows}); the defect $\tau(\ell)$ (\Cref{Y-A-Defect}); the sets $\mathcal R(b,K,\lambda,\tau)$ with their bounds (\Cref{Y-A-SetsOfRows}); the numbers $b_*$, $\lambda_*$, $K_*$, $\mathsf t_*(\rho)$ and $L_A$ (\Cref{Y-A-PureNumbers}); the statements of \Cref{Y-A-Th29}\,(a), (b); the number $b_{\mathsf B}=\|\mathsf B\|$ (\Cref{Y-A-Rows}); and the numerals (the order of the constants in \Cref{Y-A-Fr}, (0)).

We keep $\kappa\ge6$ as a letter, with $\kappa_1=\kappa+1$ (entry~1 of item~(0) of \S\ref{Y-H0-1}; item~0 (O5) of \S\ref{Y-H1-1}; item~1 (O5) of \S\ref{Y-H3-1}). The quadratic member $s_x[D]$ is that of entry~5 of item~(0) of \S\ref{Y-H0-1}. By \textup{(M2)}, $\sigma_M=160(\kappa_1+1)c_R/M\le1/12$ (item~(1) (O7a) of \S\ref{Y-H0-1}), and \textup{(M2)} itself is stated in the table of item~1 of \S\ref{Y-H9-1}. The function $\eta_w$ is that of item~(1) (O12b) of \S\ref{Y-H0-1}, and the floors and $\mathbb P$ are those of item~(1) (O13) there. By item~(1), ``No circle'', of \S\ref{Y-H0-1}, the exponent of $p_0$ is $>2$, and each floor is a family of one-sided inequalities in $x$.

Of \Cref{Y-H9-T} we use the paragraph ``Data'', \textup{(T3)} and \textup{(T6)}. We use the part ``Pairs'' of item~0 of \S\ref{Y-H1-1}, and \Cref{Y-H1-U1}\,(U), (c+), (f1), (f4), (g1), (g2).

From \S\ref{Y-H2-1} we take $e_R\ge256$ (item~1, ``Letters and data''); ``Rows'', ``Blocks'', ``Places'', the norm of a family and the charge (the definition ``Families, norm, charge, step and cut''), the source forms (the definition ``Source forms and Riemann means'') and $\mathrm{Mn}$ (item~1); the numerals (item~2), together with the letter $g_{1b}$ (item~1 (O0a) of \S\ref{Y-H3-1}); \Cref{Y-H2-U2}\,(U2.1) to (U2.3); and \Cref{Y-H2-U5}\,(S2). We use \Cref{Y-H4-U4}\,(b), (h); $C_g=263L^3$ (item~0\,(1) of \S\ref{Y-H5-1}); \Cref{Y-H5-LS}; and \Cref{Y-H5-Af}\,(i).

From \S\ref{Y-H6-1} we take the data (item~1) and \textup{(o2)} to \textup{(o7)} (item~2); \textup{(G1)} to \textup{(G4)} (item~3); \Cref{Y-H6-J}\,(J0) to (J4); the statement of item~5 that \textup{(J1)} to \textup{(J4)} hold for all $k\le K-1$, by the same proof; and \textup{(C)} and \textup{(A)} of Step~1 of the proof of \Cref{Y-H6-J}.

From \S\ref{Y-H8-1} we take $q_N$, the norms, $C_{\rm vol}$ and the congruence (item~1); and from \Cref{Y-H8-ThBC} we take the data \textup{(B)}, \textup{(b$^C$)} (ii), and \textup{(c$^C$)}, the sentence before \textup{(S)}.

From Step~2 of the proof of \Cref{Y-H9-T}, ``The last integral'', we take $\ell_k\in\mathcal R(b_*,K_*,\lambda_*,\frac19)$ for $0\le k\le K$, for every $K\ge1$. This inclusion is proved in Step~2\,\textup{(c)} of the proof of \Cref{Y-J-ThA}, where its data include $K\ge1$.

\emph{Item 0.2. The reflection and chessboard estimates.}
We use the board, copies and fold of \Cref{support:4-1}, the reflected
weight identities of \Cref{support:4-8}, and the seminorm and product
estimate of \Cref{support:4-13} and \Cref{support:4-15}.
The two special cases for a constant and for one cell follow from
\Cref{support:4-17}.

\emph{Item 0.3: what is used of earlier sections.} Of \Cref{Y-P1-Jo} we use the setting of \Cref{Y-P1-DATAo}, with any number $Q\ge1$ of sources, and then the following, cited by the tags of the proposition: \textup{(G$^\circ$)} and \textup{(G5$^\circ$)} (a); \textup{(lin)}, together with the line ``Hence $\|\delta[f](k)\|^{\mathrm{st},(\kappa_1)}\le0.2\,\mathrm{amp}(Lx_k)\|f\|_\infty$'' that follows \textup{(J4$^\circ$)}; \textup{(J1$^\circ$)}, \textup{(J2$^\circ$)}, \textup{(J3$^\circ$)} and \textup{(J4$^\circ$)}; of \textup{(L0)}, the bound $0.014$ in the norm with the rate $\kappa_1$, this bound being also \textup{(J5$^\circ$)} (c) at $k=0$; \textup{(J5$^\circ$)} (a$'$), (c) at $k=0$, and (d); \textup{(J7$^\circ$)} (ii), (iii); \textup{(J8$^\circ$)} (a$'$); and $\mathrm{sep}^\circ:=8\log(8L^{1/2})$, as it is set in item~2 of \S\ref{Y-00-1}.

Of \Cref{Y-P2-ThR} we use, in its own letters, \textup{(R4)}, \textup{(R5)} (a), (b), (c), \textup{(R6)} (i), (iii), (vi), and \textup{(R7)}; they are \textup{(Y3)}, \textup{(Y4)} and \textup{(Y8)} of \Cref{Y-00-Y}.

Of \S\ref{Y-00-1} we use \textup{(S9)}, the boards $B_i$ for $-m_l\le i\le K-1$, with walls $L^{-i}(\mathbb Z+\frac12)$; item~2, with $c_{\mathrm{vol}}(\mathsf b)=225/M^3+\frac32\mathsf N_Zp_0(\mathsf b)^2+3\log(384/\varepsilon(\mathsf b)^3)$; and \textup{(S5)}, by which $Q\ge1$ is arbitrary. Further, $\beta=L^K\mathsf b>0$ by the data.

The boards carry no offset. Every other place is reached by a translation as in \Cref{Y-P6-L1}\,(K), namely by a site of level $k$, in $L^{-j}\mathbb Z^3$; and if a board is shifted by an offset in $L^{-j}\mathbb Z^3$ and carries the cell function translated accordingly, its fold is the same.

\emph{Item 0.4: letters.} We put $a_u:=\langle u,C_Bu\rangle$, where $C_B$ is that of \Cref{Y-P6-L11}. The sign $\mathfrak h_k$ has the meaning given in \textup{(D2)} of \S\ref{Y-P1-1}, and $\mathrm{ex}:=x^{-8L}$.

Without an argument, $N_1$ is $(D_0/4M+2\check r(x)+3)^3$ (\Cref{Y-P6-L2}), whereas $N_1(k-m)$ is that of \textup{(D2)} of \S\ref{Y-P1-1} (\Cref{Y-P6-L6}). For $s\ge0$, $\tau_sf$ is the translate of $f$ by the distance $s$ along the normal to the mirror of $\theta$, away from that mirror.

The bare tags \textup{(f1)}, \textup{(f4)}, \textup{(g1)}, \textup{(g2)}, \textup{(c+)}, \textup{(U)}, \textup{(S2)}, \textup{(U2.1)} to \textup{(U2.3)}, \textup{(SD)}, \textup{(9a)}, \textup{(9e)}, \textup{(b$^C$)}, \textup{(c$^C$)}, \textup{(M2)}, \textup{(W0)}, \textup{(J0)} to \textup{(J4)} and \textup{(G1)} to \textup{(G4)} refer to the statements and items named in item~0.1. With an index, as in $R_{k'}$ and $R_k$, the letter $R$ denotes the radii of \textup{(AD)} of \Cref{Y-P1-Jo}. In Sections~\ref{Y-G0-1} to \ref{Y-H9-1}, the radius $\rho(x)=\lceil c_R\log x\rceil$ of \textup{(S1)} of \S\ref{Y-00-1} is denoted by $\varrho(x)$; in this section and in \S\ref{Y-P5-1}, $\varrho$ without an argument is $x^{1/2}p_0^6$.

\begin{theorem}[A lower bound on the two-point function, uniform in cutoff, side, depth and label]\label{Y-P6-N} The gauge group is $\operatorname{SU}(2)$, and $d=3$.

\textup{Data.} We work in the setting of the paragraph ``Data'' of \Cref{Y-H9-T}: $\mathsf b\ge x_s$; $N=2L^{m_l}$ with $m_l\ge m'$; \textup{(Tor)} $N\ge2Ms(L\mathsf b)$; $K\ge1$; and $\beta=L^K\mathsf b$. The observable is $\mathcal O_K(f)=a_K^{-1}\sum_pf(x(p))a_p$. Let $j\ge1$, $k:=K-j\ge1$ and $x:=x_k=L^j\mathsf b$; we write $\varepsilon=p_0(x)x^{-1/2}$, where $p_0=A_0(\log x)^{\mathfrak p_0}$ and $\mathfrak p_0>2$, and $\rho=\rho(x)$. Let $\nu_k$ be the law of $V=V^{(k)}$ on $\mathbb T_k=\mathbb Z^3/NL^j\mathbb Z^3$; its density is $\rho_k/\int\rho_k$, where $\rho_k=e^{\mathsf V}\rho[\bar u_k;\ell_k]>0$ at every field. The member $\bar u_k$ is real and lies in $M_1(x_k;\kappa_1,c_h)$, with the strong cap (\Cref{Y-H4-U4}\,(h)). The table \textup{(O0)} to \textup{(O13)} and all floors are those of \Cref{Y-H9-T}; none of them is chosen here, and the new numbers come after them, at \textup{(O14)}. Beyond these data, the table and the floors, we make no hypothesis.

\textup{Geometry.} We measure in units of level $k$, that is in units of $L^{-j}$, and $r=|y|$ is the sup distance. We centre the construction at the site $0$, which is the base site of $P_0$ and the centre of the boxes and of the cell $0$; other places are reached by the translations of \Cref{Y-P6-L1}\,(K). The mirror is $y_1=\frac12$, and it is a wall of $B_j$. Let $P_0$ be the 23-plaquette based at $0$, and let $u:=e_{P_0}-e_{P_0+e_2}$. Let $T_y:=\operatorname{Ad}_{V(\Gamma(0,y))}$, $\tilde F_p:=T_{\pi(p)}F_p$ and $u^\alpha(F):=\sum_pu_p\tilde F_p^\alpha$. The boxes are $B_R=\{r\le R\}$, $B'=\{r\le R+2\rho+3\}$ and $\mathrm{Box}=\{r\le r_B\}$, with $r_B:=R+h/8+1$.

Let $\chi$ be the product of $\chi_0(\phi_p/\varepsilon)$ over the plaquettes of $B'$, where $\chi_0$ is smooth, equal to $1$ on $[0,\frac12]$ and to $0$ on $[\frac34,\infty)$, with $|\chi_0'|\le8$ and $|\chi_0''|\le250$; such a $\chi_0$ exists, since for $-\chi_0'$ we can take a smooth trapezium of height $16/3$ with two ramps of width $1/16$. Let $\chi_-$ be the same product over those plaquettes all of whose sites lie in $y_1<\frac12$. We put $G_-:=\frac x2\sum_\alpha u^\alpha(F)^2\chi_-$; its values lie in $[0,\frac98p_0^2]$, and it depends only on the bonds with $\frac12-NL^j/2<y_1<\frac12$. By $G$ we denote the same function with $\chi$ in place of $\chi_-$.

Let $\chi_R(y):=\prod_\mu\chi_1(|y_\mu|/R)$, where $\chi_1=1$ on $[0,\frac12]$, $\chi_1=0$ on $[1,\infty)$, and $\chi_1(t)=1-3s^2+2s^3$ for $\frac12\le t\le1$, with $s=2t-1$; thus $|\chi_1'|\le3$ and $\chi_1'$ is $24$-Lipschitz. Let $y_c:=\lceil D_0/\sqrt2\rceil(e_1+e_2)$. Let $f_0\in C^\infty$ satisfy $0\le f_0\le1$, $f_0=1$ on $|y-y_c|\le D_0/16$ and $f_0=0$ off $|y-y_c|\le D_0/8$, and let $f_*(\xi):=f_0(L^j\xi)$. The boards are $B_{j-1}$, whose cells have side $L$ (\Cref{Y-P6-L3}), and $B_{j-i_h}$, whose cells have side $h$ (Lems.~\ref{Y-P6-L8} and \ref{Y-P6-L9}); in both, the cell $0$ is centred at $0$, and there is no offset.

\textup{New numbers.} All new numbers are placed at \textup{(O14)}, after $\mathfrak r^\circ$ and $\mathrm{sep}^\circ$. No floor depends on them; none of them depends on $N$ or on $K$; and each depends only on those that precede it, in the order below. The pure numbers are $D_0:=10^{140}$ and \[c_F:=\frac34\sum_{|y-y_c|\le D_0/16-1}|(\Pi u)(y)|^2\in[5.4\cdot10^{-5},4.0\cdot10^{-4}]D_0^{-5},\] where $\Pi$ is the lattice projection on exact 2-forms; then $c_*:=c_F^2/7\ge10^{-1410}$, $r_1:=10^{130}$ and $W_*:=10^3r_1^4$; the numbers $C_e$, $C_1$, $C_2$, $C_S$, $C_R:=60C_SC_1$, $C_6$ and $C_{\rm lat}$ of item~17 of \S\ref{Y-P6-1}, the table of the numerals; and $C_h:=64$ and $C_{\rm ext}:=32$.

The following depend on $x$ alone: $R:=\lceil C_Rp_0^5/c_F\rceil$; $h=L^{i_h}$, the least power of $L$ that is $\ge\max\{16R,64LM\log x\}$; $n_0:=\lfloor h/16M\rfloor$, $t:=x^{-2}$ and $r_\eta:=x^4$; $\varrho:=x^{1/2}p_0^6$ and $r_z/\mathsf b=3.5E_ap_0^6p_0(Lx)/(c_Q^{\rm su}K_5L^{1/2})$; $\tau_x:=3.82\cdot10^{56}x^{-256}/(1-L^{-256})$; $W_0(x):=78r_1^3(1+\log R(x))$; $\Theta(x):=\varepsilon(x)(2\ell_w(x)+4/\vartheta_0^\flat)W_0(x)$, in which $\vartheta_0^\flat$ is the number of \textup{(u1)}; and $C_a(x):=2.5(2r(x)+1)^3E_aW_*$. After these comes $r_z:=\varrho\,\mathfrak r^\circ(x)$, which depends on $\mathsf b$ besides $x$.

\textup{The list of the conditions on the depth.} The list below contains every condition on the depth. Its entries fall into two families:
\begin{enumerate}
\item[\textup{(X)}] the entries that depend on $x$ alone:
\begin{enumerate}
\item[\textup{(j1)}] each of \textup{(E1)} to \textup{(E6)} of \Cref{Y-P6-L13} is $\le c_F/12$; it follows that \textup{(E7)} $:=4[(\mathrm{E2})+(\mathrm{E3})+(\mathrm{E4})+(\mathrm{E6})]\le4c_F/3<1$, and the bound on \textup{(E2)} holds of itself;
\item[\textup{(j5)}] $C_{\rm ext}r_B^2\varepsilon\le1$;
\item[\textup{(j6)}] $x^{10}\eta_w(x)\le1$, $p_0^{10}\le x$ and $(r_z/\mathsf b)(4C_Sp_0^2+5)\le x$;
\item[\textup{(j7)}] the conditions
\begin{enumerate}
\item[\textup{(i)}] $(c_P/2)p_0^2\ge L^3(24\mathsf N_Z+4.5\log x+5+220M^{-3})$, where $c_P:=\min\{c_\Xi,\frac18\}$;
\item[\textup{(ii)}] $R/2-1\ge2D_0+M(r(x)+2)$;
\item[\textup{(iii)}] $8\rho(x)+8\le R$;
\item[\textup{(iv)}] $8M(r(x)+1)\le R$;
\item[\textup{(vi)}] $p_0(x)\ge4$, $(c_a+c_h)(x)\le\frac12$, $\log x\ge3$ and $\Theta(x)\le1$;
\item[\textup{(vii)}] $\rho_0\le\rho(x)+1$;
\end{enumerate}
\end{enumerate}
\item[\textup{(B)}] the entries that depend on $\mathsf b$ besides $x$:
\begin{enumerate}
\item[\textup{(j2)}] $c_{\mathrm{vol}}(\mathsf b)h^3L^{-3j}\le1$;
\item[\textup{(j3)}] $i_h\le j+m'-1$.
\end{enumerate}
\end{enumerate}

Once \textup{(X)} holds, the following hold of themselves (\Cref{Y-P6-L16}\,(b)). First, \textup{(E2)} $\le c_F/(12p_0)$ and \textup{(E7)} $<1$. Second, the two conditions
\begin{enumerate}
\item[\textup{(j4)}] in both of its lines, $\lfloor h/2M\rfloor-2\check r(x)-2\ge8\log(4\varrho L^{1/2})$ and $\lfloor h/2LM\rfloor-2\check r(x/L)-2\ge\mathrm{sep}^\circ$;
\item[\textup{(j7) (v)}] $\varrho L^{1/2}\ge75$.
\end{enumerate}

Third, $R\ge6r_1^3$, because $R\ge10^{1410}p_0^5$; moreover $R\ge4r_1$, $4t\sup|w|<\varepsilon/4$ for the 1-form $w:=\mathrm{d}^*(-\Delta)^{-1}\omega$ of \Cref{Y-P6-L11}, $x\ge L$ and $h<0.71x^{1/4}$.

\textup{Definitions.} Let $J(\mathsf b):=\{j\ge1$: \textup{(X)} holds at $x=L^j\mathsf b$, and \textup{(B)} holds$\}$. Let $x^\#:=1+\inf\{y\ge L$: \textup{(X)} holds at every $x\ge y\}$; this number depends on the table and on the pure numbers alone. Let $y(\mathsf b):=\mathsf b\max\{c_{\mathrm{vol}}(\mathsf b)^{1/3},L^{1-m'}\}$. For $\mathsf b\ge x_s$, let $j_*(\mathsf b)$ be the least integer $j\ge1$ with $L^j\mathsf b\ge\max\{x^\#,y(\mathsf b)^{4/3}\}$; for $0<\mathsf b<x_s$, let $j_*(\mathsf b):=z+j_*(L^z\mathsf b)$, where $z=z_0(\mathsf b)$ is that of \textup{(S7)} of \S\ref{Y-00-1}. For $\mathsf b\ge x_s$, let $j^+(\mathsf b)$ be the supremum of $j_*$ over $[\mathsf b/2,2\mathsf b]\cap[x_s,\infty)$. Finally, let $x^{\#\#}:=\max\{x^\#,\ L^3x_s,\ \sup_{\mathsf b'\in[x_s,L^2x_s]}y(\mathsf b')^{4/3}\}$.

\textup{Claims.} Let $\theta$ be the reflection $y_1\to1-y_1$.

\textup{(D)} The number $x^\#$ is finite; $[j_*(\mathsf b),\infty)\subset J(\mathsf b)$; and $j_*$ is bounded on compact sets of $[x_s,\infty)$, so that $j^+(\mathsf b)$ is finite.

For every $j\in J(\mathsf b)$, hence for all $j\ge j_*(\mathsf b)$, and for all $K\ge j+1$ and all $N$ of the data, the following clauses hold.
\begin{enumerate}
\item[\textup{(N0)}] $\mathrm{Cov}(G_-,\mathcal O_K(f))^2\le\mathrm{Var}(G_-)\,S^T_2(\theta f,f)$ for every real bounded $f$ such that $f(x(p))=0$ unless all sites of $p$ lie in $\frac12<y_1<\frac12+NL^j/2$. This clause does not depend on any entry of the list: it holds for every $j\ge1$ with $k\ge1$, once the coordinate $y_1$ is taken in the fundamental domain $\frac12-NL^j/2<y_1\le\frac12+NL^j/2$.
\item[\textup{(N1)}] $\mathsf b\,\mathrm{Cov}(G_-(V^{(k)}),\mathcal O_K(f_*))\ge c_F$.
\item[\textup{(N2)}] $\mathrm{Var}(G_-)\le\frac32a_u^2+1\le7$.
\item[\textup{(N3)}] $\mathsf b^2S^T_{2;K,N,L^K\mathsf b}(\theta f_*,f_*)\ge c_*$, where $c_*$ is pure, that is free of the cutoff, of the side, of $j$ and of $\mathsf b$.
\item[\textup{(N4)}] The same bound holds for $S_2^{(N,\mathsf b)}$ of \textup{(T3)}; for each limit of \textup{(R5)} (a), (b), (c) at the label $\mathsf b$; and, at every $j\ge j^+(\mathsf b)$, for each limit of \textup{(R5)} (a), (c) with labels $\mathsf b_i\to \mathsf b$, $\mathsf b_i\ge x_s$.
\item[\textup{(N5)}] For every limit $S$ of \textup{(R5)}, $S_2(\theta\tau_sf_*,\tau_sf_*)>0$ for every $s\ge0$, and $\dim \mathcal H(S)\ge2$; at moving labels this holds at the $j$ of \textup{(N4)} or at the $i$ of \textup{(N6$'$)}.
\item[\textup{(N6)}] The same holds for every $\mathsf b>0$ and every $j\ge j_*(\mathsf b)$, because $\mathsf b^2S_2$ is invariant under the dilation of \textup{(T6)}, \textup{(R4)}, \textup{(R6)} (vi). This clause does not include the clause of \textup{(N4)} on moving labels; for labels below $x_s$, that clause is carried by \textup{(N6$'$)} with $Lx^{\#\#}$.
\item[\textup{(N6$'$)}] One inequality holds for all labels: for every $\mathsf b>0$ and every integer $i$ with $L^i\mathsf b\ge x^{\#\#}$, $\mathsf b^2S_2(\theta_if_0(L^i\cdot),f_0(L^i\cdot))\ge c_*$, where $\theta_i$ is the mirror at $\xi_1=L^{-i}/2$. It is valid for each limit of \textup{(R5)} at the fixed label $\mathsf b$ and, provided $K\ge i+1$, on every torus whose side satisfies the conditions that \textup{(R4)} sets for it at some integer $z$ with $L^z\mathsf b\in[x_s,L^2x_s]$; with $Lx^{\#\#}$ in place of $x^{\#\#}$, it is valid also for labels $\mathsf b_n\to \mathsf b$.
\end{enumerate}
\end{theorem}

\begin{remark} The function $f_*$ depends on $\mathsf b$ through the depth $j\ge j_*(\mathsf b)$. The number $c_*$ is pure but tiny, and the depth $j_*$ is astronomically large; the numerals are generous, and none of them is sharp. For $x^\#$, $x^{\#\#}$ and $j_*$ we give closed forms only, since the table gives $L$, $M$, $A_0$, $\mathfrak p_0$, $c_\Xi$, $c_h$, $E_a$, $C_3$, $\vartheta_0^\flat$, $\rho_0$, $c_R$ and $\kappa$ as letters. The depth $j_*(\mathsf b)$ is defined by its closed form, and not as the least depth at which the list holds; for large $\mathsf b$ this closed form grows like $\frac13\log_L\mathsf b$.
\end{remark}

\emph{The parts of the proof.} The proof of \Cref{Y-P6-N} falls into the parts below.

The first part consists of two statements. One is the clause \textup{(N0)}, which is \Cref{Y-P6-L1}\,(W), the witness inequality; it follows from \textup{(L)}, \textup{(I)} and \textup{(K)} of that lemma (the joint law: locality, independence and the covariance of the kernel) and from \textup{\ref{support:4-8}}\,(c), (d). The other, \Cref{Y-P6-L2}, is the split $E[\mathcal O_K(f)|V]=c_f+A_f+W_f$; it is used for \textup{(N1)}. The inputs of this part are ``The step'' and ``Runs'' of item~2 of \S\ref{Y-J-1}, the clauses (c) and (d) of \textup{\ref{support:4-8}}, \Cref{Y-H4-U4}\,(h), \Cref{Y-H6-J}\,(J0) to (J4), and item~5 of \S\ref{Y-H6-1}.

\Cref{Y-P6-L3} is the Peierls bound at $\varepsilon/2$; it is free of $N$, $K$ and $\mathsf b$, and it is stated in the norm form \textup{(P$^\circ$)}. In its proof, the numerator and the denominator are treated cell by cell, and then the root of order equal to the number of cells is taken. In this way the volume drops out by soft means: by \textup{\ref{support:4-15}} on $B_{j-1}$ for the block law, by the covering clause together with the strong cap, and by a lower bound for the denominator, the integral of the density, which we obtain on a small ball of fields about $V=1$ (Step~7 of the proof of \Cref{Y-P6-L3}).

\Cref{Y-P6-L4} extends a field that is tame on the box $\{r\le r_B\}$ to a globally tame field on $\mathbb T_k$. \Cref{Y-P6-L5} compares two sums, the sum $\Lambda$ of the lemma and the sum $c_f+A_f$ of the split, with their images under the maps of the fold of $B_{j-i_h}$, that is under the mirrors and the translations of its cells; only these sums enter the comparison, not the single members, so that no covariance of a member is used. In both lemmas no member is reflected.

\Cref{Y-P6-L6} is \Cref{Y-P1-Jo} with a free last ratio. It is \textup{(J5$^\circ$)} (d) and \textup{(J8$^\circ$)} (a$'$) of \Cref{Y-P1-Jo}; its proof is an induction and uses \textup{(9e)}, \textup{(g2)}, \textup{(J1$^\circ$)} and \textup{(J2$^\circ$)}.

For any $\hat w\ge0$ that takes the place of the $w$-member of $\bar u_k$ and lies within the cap $2.6e^{-\kappa(|X|_c-1)}$, \Cref{Y-P6-L7} bounds the integral of the density by $(1+C_{\rm vol})\int\rho[\bar u_k;\ell_k]$, and it gives the root bound. Lems.~\ref{Y-P6-L8} and \ref{Y-P6-L9} build their majorants and use this lemma. \Cref{Y-P6-L8} bounds the part $W_f$ of the split, and \Cref{Y-P6-L9} the defect of the integration by parts caused by $\Xi_w$; both the part and the defect come from the weights, the $w$-member of $\bar u_k$. The last two lemmas are stated in the norm forms \textup{(W$^\circ$)} and \textup{(X$^\circ$)}. Their bounds are independent of $N$ and of $K$, whereas the first bound of \Cref{Y-P6-L7} contains the factor $1+C_{\rm vol}$, which depends on $N$. The proofs of Lems.~\ref{Y-P6-L8} and \ref{Y-P6-L9} take the fold over the cells of $B_{j-i_h}$ and then the root of order equal to the number of cells. The volume factors that arise in these proofs are removed by \textup{\ref{support:4-15}} on $B_{j-i_h}$ and by both bounds of \Cref{Y-P6-L7}. The proof of the first of these bounds, with its factor $1+C_{\rm vol}$, uses \textup{(9a)} with \textup{(f1)}, that is the exact cancellation of the vacuum factors along $k\le i<K$; \Cref{Y-H5-Af}\,(i); and \Cref{Y-H8-ThBC}\,(c$^C$). The last bound of \Cref{Y-P6-L7} follows from \textup{(j2)}: the root of $1+C_{\rm vol}$ of order equal to the number of cells of $B_{j-i_h}$ is $\le e$.

Lems.~\ref{Y-P6-L10}, \ref{Y-P6-L12} and \ref{Y-P6-L13} carry out the integration by parts in $\nu_k$, and give the key line, four identities and the list of errors. Every power in their estimates is explicit.

Lems.~\ref{Y-P6-L11} and \ref{Y-P6-L14} give the Green function of $\mathrm{d}^*\mathsf B_k\mathrm{d}$ uniformly in the rows, and $q\ge c_F$. Both concern the free model, and both have explicit constants.

The clauses \textup{(N1)} to \textup{(N6)} and \textup{(N6$'$)} are derived in \Cref{Y-P6-L15}. The derivation rests on \textup{(T3)} and \textup{(T6)}, and on \textup{(R4)}, \textup{(R5)}, \textup{(R6)} (i), (iii), (vi) and \textup{(R7)} of \Cref{Y-P2-ThR}.

\Cref{Y-P6-L16} treats the list of the conditions on $j$ and the depth $j_*(\mathsf b)$. It gives each entry together with its use and the reason why it holds; $x^\#$ is finite, $[j_*(\mathsf b),\infty)\subset J(\mathsf b)$, and $j_*$ is locally bounded.

The numerals are collected in one table. For $C_2$, $C_6$ and the numeral of \textup{(E4)}, the table gives generous closed forms, with the rows from which they come.

\emph{Statements and items of other sections, as used here.}

\emph{Item~1 of \S\ref{Y-G0-1} (``Cubes'', ``At $x$'', ``Fields'' and ``Domains'', ``Members'', \textup{(SD)}) and item~3 there, \textup{(O3)}.} These two items supply the formula for the score of $\nu_k$. On fields that are tame on $B'$, $S_\varepsilon(t)=t^2/2$, $\zeta=F$, and $m_y=1$ for every row that the flow moves.

\emph{Item~2 of \S\ref{Y-J-1}, ``Energy'', with \textup{(W0)} of \Cref{Y-E-C6}.} The mask depends on $D(V)$ within $\rho_0$ of $y$, and on the plaquettes with $\phi_p>1$ within $\rho(x)+1$ of $y$; the rows have range $\le\rho(x)$; and the paths go along $e_1$, then $e_2$, then $e_3$. Once $\rho_0\le\rho(x)+1$, which is \textup{(j7) (vii)}, this fixes the radius $R+2\rho+3$ of $B'$.

\emph{Item~2 of \S\ref{Y-J-1}, ``The step'', ``Runs''.} The average $\mathsf M(U)(c)$ depends on $B(y)$, $B(y')$ and straight paths. Since the kernel $k_{ax}$ of the step is a probability kernel, $\int\rho'=\int\rho$. The step is independent of the rows, of the grid and of the members. The site $y$ of $\mathbb T_k$ lies at $L^{-j}y$ (``Places'', item~1 of \S\ref{Y-H2-1}, which sets it at $a_KL^ky=L^{-j}y$).

\emph{Item~3 of \S\ref{Y-J-1}, \textup{(R9)}, with \textup{(9a)}, \textup{(9e)}.} We apply \textup{(9a)}, which holds at every $V'$, to a single real member $\hat u$ and to its iterates (\Cref{Y-P6-L7}); and in the range over which \Cref{Y-H1-U1}\,(U) extends the clauses (a) to (e) of \Cref{Y-H-C9}, that is at every pair $p\in\mathbb P$, we apply \textup{(9e)} to $u_1,u_2\in M_2$ (\Cref{Y-P6-L6}).

\emph{\S\ref{Y-A-1}, with the definitions and \Cref{Y-A-Th29} named in item~0.1.} The statement of \Cref{Y-A-Th29}\,(a) gives $\ell'^{\rm c}\in\mathcal R(b_*,K_*,\lambda_*,\tau+\mathsf t_*(\rho))$, and induction from $\tau(\emptyset)=0$ gives $\tau(\ell_k)\le\sum_{i\le k}\mathsf t_*(\rho(x_i))$. The definition of $\mathsf t_*$ (\Cref{Y-A-PureNumbers}), together with $\rho(x)\ge c_R\log x$ and $e_R\ge256$, then yields $\tau(\ell_k)\le\tau_x$. If only the stated bound $\tau\le\frac19$ were at hand, the term $C_\tau\tau^2r^{-4}$ of \Cref{Y-P6-L11}\,(b) would swamp the continuum form $\pi$ of \Cref{Y-P6-L14}\,(A), which is of size $r^{-4}$; the sharper bound is therefore needed. The bound $b_{\mathsf B}\le b_*$ holds at the momenta of the torus, and Step~0 of the proof of \Cref{Y-P6-L11} carries it to all $p$.

\emph{Item~0 of \S\ref{Y-H1-1}, ``Pairs'' ($\mathbb P$), and \Cref{Y-H1-U1}\,(U), (c+), (f1), (f4), (g1), (g2).} By \textup{(f1)}, neither $E'^a$ nor $\mathsf V'$ depends on $E^h$ or on $w$, so the vacuum factors of $\hat u$ and $\bar u$ cancel exactly. The function $w(X,V)$ of \textup{(f4)} enters only for sets $X$ that hold a plaquette of $D(V)$. We apply \textup{(g2)} with $A=2$ and with a pair $p\in\mathbb P$ whose first entry is a rate $\ge\kappa$, one of the rates of \Cref{Y-P6-L6}, and whose second entry is $c_h$.

\emph{\S\ref{Y-H2-1}, with ``Rows'', ``Blocks'', ``Places'', the charge (the definition ``Families, norm, charge, step and cut'' of item~1), $\mathrm{Mn}$, \textup{(U2.1)} to \textup{(U2.3)}, \textup{(S2)} and the numerals.} The barycentres of $B^k(y)$ lie within $\frac12$ of $y$, in units of level $k$. Hence $\mathrm{Mn}=1$ for $|y-y_c|\le D_0/16-\frac12$, $\mathrm{Mn}\ge0$ always, and $\mathbb D_{k,y}[f_*]=0$ for $|y-y_c|>D_0/8+\frac12$. For each site $y$, the error of \textup{(S2)} is bounded in operator norm, and these bounds are summed over the sites $y$ of the outer ball $|y-y_c|\le D_0/8+\frac12$.

\emph{\Cref{Y-H4-U4}\,(b), (h).} The condition \textup{(SD)} holds at $(x_i;\ell_i)$ for $0\le i<K$. At every field, $\rho_k^z=\exp\{\sum V_i(z)\}\rho[u_k(z);\ell_k]$.

\emph{\Cref{Y-H5-LS} and \Cref{Y-H5-Af}\,(i).} We apply \Cref{Y-H5-Af}\,(i) with $B=2.6$ and $p=(\kappa,c_h)$: the theorem assumes families within $B$ caps at $p$, with the weak cap, and its clause (i) yields $F_\ell(u_i)$ in $M_1(x';p)$, with the strong cap ($0.35+0.248\cdot2.6=0.9948$). \Cref{Y-H5-LS}\,(S1) is applied value by value.

\emph{\Cref{Y-H6-J}, with items 1, 2, 3, 5 of \S\ref{Y-H6-1} and with \textup{(C)}, \textup{(A)} of Step~1 of its proof.} We use \textup{(J2)}, \textup{(J3)} and \textup{(J4)} for one source. \Cref{Y-P1-Jo} contains them as \textup{(J2$^\circ$)}, \textup{(J3$^\circ$)} and \textup{(J4$^\circ$)}, for one complex source of any support and for $0\le k\le K$. \Cref{Y-H6-J} by itself yields them for $n=2$ as well, if the second source is a translate, shifted by the distance $(h/2)L^{-j}$, and is taken at $z_2=0$.

\emph{\Cref{Y-H8-ThBC}, with item~1 of \S\ref{Y-H8-1}.} The sentence $\int|\rho[u_2;\ell]-\rho[u_3;\ell]|\le\frac12C_{\rm vol}\mathsf D_{23}I_1$ of \textup{(c$^C$)}, case (ii), is applied in \Cref{Y-P6-L7}, with $x=\mathsf b$, with $u_1=u_3=\bar u_K$ and with $u_2=\hat u_K$, an output of \Cref{Y-H-C9} with $A\le2.6\le6$; and $\mathsf D_{23}\le2$.

\emph{\textup{\ref{support:4-1}}\,(2), \textup{\ref{support:4-2}}\,(1), \textup{\ref{support:4-8}}\,(c), (d), (e), \textup{\ref{support:4-13}} with \eqref{support:4-14}, \textup{\ref{support:4-15}}, \textup{\ref{support:4-17}}\,(b), (f).} The boards have $\lambda$ odd with $\lambda\ge3$, $\mathsf n_\mu$ even with $\mathsf n_\mu\ge2$, and walls at the midpoints of bonds. The functions are bounded Borel functions of those bonds whose two ends lie in the cell. Clause (c) of \textup{\ref{support:4-8}} is the Schwarz inequality. For $\mathfrak F_+$ we take a half torus, and clause (d) of \textup{\ref{support:4-8}} gives the other half. The copies are $F(u^\sigma)$; for real $F$, no conjugation enters.

\emph{\Cref{Y-P1-Jo}, with \textup{(G$^\circ$)}, \textup{(G5$^\circ$)} (a), \textup{(lin)}, \textup{(J1$^\circ$)} to \textup{(J4$^\circ$)}, \textup{(L0)}, \textup{(J5$^\circ$)} (a$'$), (c), (d), \textup{(J7$^\circ$)} (ii), (iii) and \textup{(J8$^\circ$)} (a$'$).} \Cref{Y-P6-L6}\,(A) is \textup{(J5$^\circ$)} (d), and \Cref{Y-P6-L6}\,(B) is \textup{(J8$^\circ$)} (a$'$). We take $R_{k'}=\mathfrak r^\circ_{k'}$ for $k'<k$, and $R_k=\varrho\,\mathfrak r^\circ_k$. The rates satisfy $\varkappa_{k'}\in[\kappa+\frac14,\kappa+0.51]$, and the last drop satisfies $2\log(4\varrho L^{1/2})/\mathrm{sep}_k\le\frac14$. The hypothesis \textup{(AD)} (ii) holds: for its number $q_{k'}$ we have $q_{k'}\le(1+L^{1/2})/(4\,L^{1/2})$ for $k'<k-1$ and $q_{k-1}\le(1+\varrho L^{1/2})/(4\,\varrho L^{1/2})$. Both right-hand sides are $\le(1+1/75)/4$, hence $\le1$, and since every drop $\varkappa_{k'}-\varkappa_{k'+1}$ is $\le\frac14$, it follows that $0.364e^{1/4}q_{k'}^2\le0.029996$. Here $L^{1/2}$ and $\varrho L^{1/2}$ are the bounds on $R_{k'+1}/R_{k'}$ that \textup{(G3)} gives.

\emph{\Cref{Y-P2-ThR}, with \textup{(R4)}, \textup{(R5)} (a) to (c), \textup{(R6)} (i), (iii), (vi) and \textup{(R7)}.} The dilation covariance is \textup{(R6)} (vi). Moving labels are treated in Step~4 of the proof of \Cref{Y-P6-L15} and in \textup{(N6$'$)}. The space $\mathcal H(S)$ and all real translations are at hand for the limits of \textup{(R5)} alone.

\emph{\textup{(S9)} of \S\ref{Y-00-1}, the boards.} The mirror $y_1=\frac12$ is a wall of $B_j$ ($\lambda=L^k\ge3$). The board of \Cref{Y-P6-L3} is $B_{j-1}$ ($\lambda=L^{k+1}$, $\mathsf n_\mu=2L^{m_l+j-1}$). The board of Lems.~\ref{Y-P6-L8} and \ref{Y-P6-L9} is $B_{j-i_h}$ ($\lambda=L^{k+i_h}$, $\mathsf n_\mu=2L^{m_l+j-i_h}\ge2L$ by \textup{(j3)}); its cells are centred at $h\mathbb Z^3$, and $h(\mathbb Z+\frac12)\subset\mathbb Z+\frac12$ because $h$ is odd.

\emph{Item~2 of \S\ref{Y-00-1}, with $c_{\mathrm{vol}}(\mathsf b)$ and the subtraction $c_f+A_f$.} We have \[\log(1+C_{\rm vol})\le\log(8q)+220q+[\frac32\mathsf N_Zp_0(\mathsf b)^2+3\log(384/\varepsilon(\mathsf b)^3)]N^3\le c_{\mathrm{vol}}(\mathsf b)N^3,\] with the numeral $225$ in $c_{\mathrm{vol}}(\mathsf b)$, because $q=q_N=(N/M)^3\ge1$ (item~1 of \S\ref{Y-H8-1}) and $\log(8q)\le2.08q$ for $q\ge1$. The subtraction $c_f+A_f$ is neither $\hat c_K$ nor the torus mean; its mirror defect is treated in \Cref{Y-P6-L5}.

\begin{lemma}[The joint law and the witness inequality]\label{Y-P6-L1}
Let $P$ be the joint law of $(U=V^{(0)},V^{(1)},\dots,V^{(K)})$, that is, Wilson's measure multiplied by the kernels $\prod_ck_{ax_i}(V^{(i+1)}(c),\allowbreak \mathsf M(V^{(i)})(c))$, $0\le i<K$, as they are defined in item~2 of \S\ref{Y-J-1}, ``The step'' and ``Runs''. As $\mathsf M$ is Borel and $k_s(\cdot,h)dg$ is a probability measure, $P$ is a probability measure; we denote by $\nu_k$ the law of the component $V^{(k)}$ under $P$. For a set $A$ of sites of $\mathbb T_k$ we write $A^{(i)}$ for the union of the $(k-i)$-fold centred blocks of the sites of $A$.
\begin{enumerate}
\item[\textup{(L)}] The kernels from which the values of $V^{(i+1)}$ on the bonds of $A^{(i+1)}$ are drawn depend on $V^{(i)}$ only through its values on the bonds of $A^{(i)}$.
\item[\textup{(I)}] Let $A_1,\dots,A_m$ be pairwise disjoint, and for every $s$ let $\Phi_s$ be a bounded Borel function of the values of $U$ on the bonds of $A_s^{(0)}$ together with the values of $V^{(k)}$ on the bonds of $A_s$. Then $E[\prod\Phi_s|U]=\prod E[\Phi_s|U]$, and each $E[\Phi_s|U]$ is bounded, Borel, and determined by the values of $U$ on the bonds of $A_s^{(0)}$ alone.
\item[\textup{(K)}] Let $\check\theta$ act in each coordinate either as $y_\mu\to y_\mu+n_\mu$ or as $y_\mu\to n_\mu-y_\mu$, with $n_\mu\in\mathbb Z$ measured in units of level $k$, and define $(\check\Theta V)(b):=V(\check\theta\,b)$ at every level $i\le k$, under the convention $V(\text{reversed }b)=V(b)^{-1}$. Then $\mathsf M(\check\Theta U)(c)=\mathsf M(U)(\check\theta\,c)$; acting on the levels $0,\dots,k$ simultaneously, $\check\Theta$ leaves the marginal of $P$ on these levels invariant; $E[\Phi\circ\check\Theta|U]=E[\Phi|\cdot](\check\Theta U)$; $\rho_k\circ\check\Theta=\rho_k$; and $\rho_k^{[zf]}(\check\Theta V)=\rho_k^{[zf\circ\check\theta^{-1}]}(V)$, where $\rho_k^{[zf]}(V)$ is the integral of Wilson's weight, times $e^{z\mathcal O_K(f)}$, times the $k$ kernels.
\item[\textup{(W)}] Let $G_-$ be a bounded real Borel function of those bonds of $V^{(k)}$ whose two ends lie in the region $\frac12-NL^j/2<y_1<\frac12$, and let $f$ be as in \textup{(N0)}. Then $\mathrm{Cov}(G_-,\mathcal O_K(f))^2\le\mathrm{Var}(G_-)S^T_2(\theta f,f)$.
\end{enumerate}
\end{lemma}

\begin{proof}
\emph{Step \textup{(L)}: locality.}
\begin{enumerate}
\item[\textup{1.}] Let $c=(y,y')$ with $y'=y+Le_\mu$. The formula for $\mathsf M$ (item~2 of \S\ref{Y-J-1}, ``The step'') depends on $U$ only through $\mathsf t_U(x)=U(\Gamma(y(x),x))$, through the straight paths $[x,x+Le_\mu]$ with $x\in B(y)$, and through $\mathsf t_U(x')$.
\item[\textup{2.}] A coordinate path running from the centre of a box to one of its points changes a single coordinate at a time, each of them monotonically; therefore it never leaves the box.
\item[\textup{3.}] The straight path is contained in $B(y)\cup B(y')$.
\item[\textup{4.}] Consequently $\mathsf M(U)(c)$ depends only on bonds with both ends in $B(y)\cup B(y')$, and \textup{(L)} follows by induction on the levels.
\end{enumerate}

\emph{Step \textup{(I)}: independence.}
\begin{enumerate}
\item[\textup{1.}] Conditionally on $U$, the law of $(V^{(1)},\dots,V^{(K)})$ has a density, namely a product over all bonds of the levels $1,\dots,K$ that contains one kernel factor for every bond.
\item[\textup{2.}] The factors of the levels above $k$ integrate to $1$, level by level, beginning with $K$ and descending.
\item[\textup{3.}] We call a bond of level $i$ a bond of $s$ if both of its ends lie in $A_s^{(i)}$, and free if it is a bond of no $s$.
\item[\textup{4.}] By \textup{(L)}, the factor of a bond of $s$ depends only on bonds of $s$ one level lower.
\item[\textup{5.}] We integrate over the free bonds, beginning at level $k$ and descending. A free bond of level $k$ occurs neither in any $\Phi_s$ nor in any factor besides its own, so its factor integrates to $1$.
\item[\textup{6.}] Suppose that the free bonds of all levels above $i$ are already integrated out. The factors that remain at level $i+1$ belong to bonds of some $s$, and by item~4 they depend only on bonds of $s$ at level $i$; hence a free bond of level $i$ occurs in no remaining factor except its own, and its factor integrates to $1$. We apply this for $i=k-1,k-2,\dots,1$ in turn.
\item[\textup{7.}] What is left is a product over $s$ of integrals, and the $s$-th integral depends only on the values of $U$ on the bonds of $A_s^{(0)}$.
\end{enumerate}

\emph{Step \textup{(K)}: covariance.}
\begin{enumerate}
\item[\textup{1.}] At the transition from level $i$ to level $i+1$, $i<k$, the offsets of $\check\theta$ in fine units are $L^{k-i}n_\mu$, hence multiples of $L$. Apply \Cref{Y-P2-L2} to the coordinate reflections and translations composing $\check\theta^{-1}$; its convention $(RV)(b)=V(R^{-1}b)$ then agrees with $\check\Theta$. This gives the covariance of $\mathsf M$ and the invariance of each transition kernel, including reversed bonds and the value $\mathsf L(-1)=0$.
\item[\textup{2.}] Wilson's weight and the Haar measure are invariant, as $\operatorname{tr}g^{-1}=\operatorname{tr}g$.
\item[\textup{3.}] Barycentres are mapped to barycentres; hence $\mathcal O_K(f)(\check\Theta U)=a_K^{-1}\sum_pf(x(p))a_{\check\theta p}(U)=\mathcal O_K(f\circ\check\theta^{-1})(U)$.
\item[\textup{4.}] The levels above $k$ integrate to $1$ (item~2 of Step~\textup{(I)}). The change of variables $V^{(i)}\to\check\Theta V^{(i)}$ at the levels $0\le i\le k$ then yields the five assertions of \textup{(K)}.
\end{enumerate}

\emph{Step \textup{(W)}: the witness inequality.}
\begin{enumerate}
\item[\textup{1.}] At level $k$ the planes $y_1=\frac12$ and $y_1=\frac12+NL^j/2$ lie in $\mathbb Z+\frac12$, because $NL^j/2=L^{m_l+j}$.
\item[\textup{2.}] At level $i$, the plane $m+\frac12$ of level $i+1$ becomes the plane $L(m+\frac12)$. Since $L$ is odd, this plane lies in $\mathbb Z+\frac12$; moreover it is a face of the centred blocks. In bare units, therefore, both planes lie in $\mathbb Z+\frac12$.
\item[\textup{3.}] Let $G_+:=G_-\circ\check\Theta$ with $\check\theta=\theta$ ($n_1=1$). Applied with $A$ the right half, \textup{(I)} shows that $\breve\Gamma:=E[G_+|U]-EG_+$ belongs to $\mathfrak F_+$ of \textup{\ref{support:4-8}}; and \textup{(K)} gives $E[G_-|U]-EG_-=\theta\,\breve\Gamma$.
\item[\textup{4.}] By the hypothesis on $f$, also $F:=\mathcal O_K(f)-\langle \mathcal O_K(f)\rangle$ belongs to $\mathfrak F_+$.
\item[\textup{5.}] All functions here are real, and $\mathrm{Cov}(G_-,\mathcal O_K(f))=E[(E[G_-|U]-EG_-)F]=\langle\theta\,\breve\Gamma\,F\rangle$.
\item[\textup{6.}] By \textup{\ref{support:4-8}}\,(c), $\langle\theta\,\breve\Gamma\,F\rangle^2\le\langle\theta\,\breve\Gamma\,\breve\Gamma\rangle\langle\theta F\,F\rangle$.
\item[\textup{7.}] By the Schwarz inequality and the invariance of the measure, $\langle\theta\,\breve\Gamma\,\breve\Gamma\rangle\le\langle\breve\Gamma^2\rangle$; furthermore $\langle\breve\Gamma^2\rangle=\mathrm{Var}\,E[G_+|U]\le\mathrm{Var}\,G_+=\mathrm{Var}\,G_-$.
\item[\textup{8.}] Item~3 of Step~\textup{(K)} gives $\langle\theta F\,F\rangle=S^T_2(\theta f,f)$. Combining items 5 to 8 we arrive at $\mathrm{Cov}^2\le\mathrm{Var}(G_-)S^T_2(\theta f,f)$.
\end{enumerate}
\end{proof}

\begin{lemma}[A decomposition of the conditional mean]\label{Y-P6-L2}
At every $V$ of $\mathbb T_k$ put $T_f(V):=\frac d{dz}\log\rho_k^{[zf]}(V)$ at $z=0$. Then $T_f(V)$ is a version of $E[\mathcal O_K(f)|V^{(k)}=V]$ that is defined at every $V$, and $T_f=c_f+A_f(V)+W_f(V)$. Here $c_f=\sum_{i\le k}\frac{d}{dz}\big|_{z=0}V_i(z)$; next, $A_f=\sum_{X\in\mathcal D}1_{X^{\approx}}(\delta^a+\delta^h)(X,V)=S_f/\mathsf b+\mathrm{Rem}$, where $S_f$ denotes the sum over all cubes of $1_{\mathrm{cube}^{\approx}}s_x[\mathbb D_k[f]]$; and $W_f(V)=\big[\frac{d}{dz}\big|_{z=0}\Xi_{w(z)}(V)\big]/\Xi_{\bar w}(V)$. In these formulas $V_i(z)$ and $u_k(z)=(\cdot,\cdot,w(z))$ are those of \Cref{Y-H4-U4}\,(h) along the run with the single source $f$ and the variable $z$, and $\bar w=w(0)$. For $f=f_*$ we have $\sup|S_f|\le C_Sp_0^2$ and \[\mathsf b\sup|\mathrm{Rem}|\le1.2124N_1[2C_3K_5E_ap_0^3x^{-1/4}+0.12c_Q^{\rm su}K_5(Lx)^{1/2}c_h(x)/(E_ap_0(Lx))],\] with $N_1=(D_0/4M+2\check r(x)+3)^3$; hence $\mathsf b\sup|A_f|\le C_Sp_0^2+1$ under the condition \textup{(j1)}.
\end{lemma}

\begin{proof}
\emph{Step 1: the conditional mean as a logarithmic derivative.} The density $\rho_k^{[zf]}(V)$ is the integral of Wilson's weight, times $e^{z\mathcal O_K(f)}$, times the $k$ kernels. Since $\mathcal O_K(f)$ is bounded and $\rho_k>0$ at every $V$ (\Cref{Y-H4-U4}\,(h)), we have $E[e^{z\mathcal O_K(f)}|V]=\rho_k^{[zf]}(V)/\rho_k(V)$, and differentiation at $0$ is legitimate.

\emph{Step 2: the density in terms of members.} According to \Cref{Y-H4-U4}\,(h), whose density $\rho_k^z$ is $\rho_k^{[zf]}$ for the single source $f$, we have $\rho_k^{[zf]}=\exp\{\sum V_i(z)\}\rho[u_k(z);\ell_k]$ at every field, and the values are holomorphic.

\emph{Step 3: the three parts.} The formula for $\rho[u;\ell]$ (item~1 of \S\ref{Y-G0-1}, ``Members'') produces the three parts of $T_f$. Here $\Xi_{\bar w}(V)>0$, as $\rho[\bar u_k]>0$.

\emph{Step 4: the $a$-part.} From \Cref{Y-H6-J}\,(J3) we have $\delta^a=s_x[\mathbb D_k[f]/\mathsf b]+\mathrm{rem}$, where \[|\mathrm{rem}(X,\cdot)|\le2u_3(x)K_5\mathsf b^{-1}c_a(x)e^{-\kappa_1(\#X-1)}=2C_3K_5E_ap_0^3x^{-1/4}\mathsf b^{-1}e^{-\kappa_1(\#X-1)}.\]

\emph{Step 5: the $h$-part.} From \Cref{Y-H6-J}\,(J4) we have \[\FitDisplay{\dimexpr\linewidth-1pt\relax}{|\delta^h(X,\allowbreak \cdot)|\le0.12\,\allowbreak \mathrm{amp}(Lx)c_h(x)e^{-\kappa_1(\#X-1)}}\], and \[\FitDisplay{\dimexpr\linewidth-1pt\relax}{\mathsf b\,\allowbreak \mathrm{amp}(Lx)=c_Q^{\rm su}K_5(Lx)^{1/2}/(E_ap_0(Lx))}\] by item~2 of \S\ref{Y-H6-1}, \textup{(o5)}.

\emph{Step 6: the sets that occur.} \Cref{Y-H6-J}\,(J2) shows that an $X$ occurs only when $X^{+\check r}$ meets $S(k)$.

\emph{Step 7: the anchor cubes.} The support of $f_*$ is a sup-ball of radius $D_0/8$. It meets no more than $(D_0/4M+2)^3$ cubes, and there are therefore at most $N_1$ anchor cubes.

\emph{Step 8: connected sets through a cube.} For a given cube, the sum of $e^{-\kappa(\#X-1)}$ over the connected $X$ that contain it is at most $1/(1-26e\,e^{-6})<1.2124$ (item~1 of \S\ref{Y-G0-1}, ``Cubes''; $\kappa\ge6$).

\emph{Step 9: the quadratic member.} The member $s_x[D]$ is given by the formula of entry~5 of item~(0) of \S\ref{Y-H0-1}. On fields tame on $\mathrm{cube}^{\approx}$, each $F_p$ entering it satisfies $|F_p|\le\varepsilon$, because the range does not exceed $\rho(x)\le Mr(x)-2$.

\emph{Step 10: the bound on $S_f$.} We use \textup{(U2.1)}, the three orientations at a site, and \textup{(U2.3)}, by which only sites $y\in Y_{\rm out}=\{|y-y_c|\le D_0/8+\frac12\}$ occur; this gives $|S_f|\le\frac x2\#Y_{\rm out}K_5(3\mathsf g_\flat)^2\varepsilon^2=C_Sp_0^2$.

\emph{Step 11: a single source.} Steps~4 to~6 use \textup{(J3)}, \textup{(J4)} and \textup{(J2)} for the single source $f$ at the level $k$. \Cref{Y-P1-Jo} contains them, with the same bounds, as \textup{(J3$^\circ$)}, \textup{(J4$^\circ$)} and \textup{(J2$^\circ$)} for one complex source of any support and for $0\le k\le K$.
\end{proof}

\begin{lemma}[A Peierls bound at the threshold $\varepsilon/2$]\label{Y-P6-L3}
\textup{(P$^\circ$)} In norm form: at $j\in J(\mathsf b)$, on the board $B_{j-1}$, and for $p$ based at $0$, \[\|E[1\{\phi_p\ge\varepsilon/2\}|U]\|_{\rm ch}\le\exp\{-c_Pp_0^2+L^3(24\mathsf N_Z+4.5\log x+5+220M^{-3})\}\le e^{-(c_P/2)p_0^2}.\] Consequently, for every plaquette $p$ of $\mathbb T_k$, \[\nu_k(\phi_p\ge\varepsilon/2)\le\exp\{-c_Pp_0^2+L^3(24\mathsf N_Z+4.5\log x+5+220M^{-3})\}\le e^{-(c_P/2)p_0^2},\] where $c_P=\min\{c_\Xi,\frac18\}$. This bound does not depend on $N$, $K$ or $\mathsf b$. Hence \[\nu_k(\text{some }\phi_p\ge\varepsilon/2\text{ in }B')\le3(2R+4\rho+7)^3e^{-(c_P/2)p_0^2}=:\nu_P.\]
\end{lemma}

\begin{proof}
\emph{Step 1: translation to the origin.} Using \Cref{Y-P6-L1}\,(K) for translations, we may take $p$ based at $0$.

\emph{Step 2: the board.} The board $B_{j-1}$ has cells of side $L$ in units of level $k$, with $\lambda=L^{k+1}$ and $\mathsf n_\mu=2L^{m_l+j-1}$. The cell $0$ is centred at $0$; as $L\ge13$, the plaquette $p$ is therefore inner.

\emph{Step 3: a function of one cell.} According to \Cref{Y-P6-L1}\,(I), $F:=E[1\{\phi_p\ge\varepsilon/2\}|U]$ is a function of the cell $0$.

\emph{Step 4: the chessboard estimate.} By \textup{\ref{support:4-15}}, with the function $1$ in all other cells (\textup{\ref{support:4-17}}\,(b), (f)), we get $\nu_k(\phi_p\ge\varepsilon/2)=\langle F\rangle\le\langle\operatorname{Fold}F\rangle^{1/\#\mathrm{cells}}=\|F\|_{\rm ch}$. Steps 5 to 8 bound the fold and thereby prove \textup{(P$^\circ$)}.

\emph{Step 5: the copies.} The copies are $F\circ\check\Theta_c$, with $\check\Theta_c$ as in the data of \Cref{Y-P6-L5}, but with the side $L$ in place of $h$; the map $\check\theta_c^{-1}$ is of the kind allowed in \textup{(K)}. From \textup{(K)} and \textup{(I)} we get $\langle\operatorname{Fold}F\rangle=\nu_k(\phi_{\check\theta_c^{-1}p}\ge\varepsilon/2\text{ in every cell }c)$. The event is the same for either orientation of the mirrored plaquette $\check\theta_c^{-1}p$, since $\operatorname{ang}\,g^{-1}=\operatorname{ang}\,g$.

\emph{Step 6: the numerator, pointwise.} To begin with, $\mathrm{En}_\ell\ge\sum_pS_\varepsilon(\phi_p)$ (item~2 of \S\ref{Y-J-1}, ``Energy''). Next, writing $q$ for the number of cubes, $|\sum1(E^a+E^h)|\le1.2124q(c_a+c_h)\le0.61q$. Further, since $\bar w$ has the strong cap, and since the members of any family are pairwise disjoint and together cover $D(V)$, we obtain \[|\Xi_w|\le\exp\{-c_\Xi x\varepsilon\sum_{D(V)}\phi_p\}\prod_{\mathrm{cubes}}(1+\sum_{X\text{ holding the cube}}e^{-\kappa(|X|_c-1)})\le\exp\{-c_\Xi x\varepsilon\sum_{D(V)}\phi_p\}e^{217.6q},\] where $\sum_{n\ge1}(52e)^ne^{-6(n-1)}=52e/(1-52e\,e^{-6})<217.6$. Lastly, the marked plaquette of a cell contributes the factor $e^{-x(\varepsilon/2)^2/2}=e^{-p_0^2/8}$ if $\phi\le\varepsilon$, and the factor $e^{-c_\Xi x\varepsilon\phi}\le e^{-c_\Xi p_0^2}$ if $\phi>\varepsilon$.

\emph{Step 7: the denominator.} We restrict the integral to the ball $\{\operatorname{ang}\,V(b)\le x^{-1/2}\text{ for all }b\}$. On this ball $\phi_p\le4x^{-1/2}\le\varepsilon$ by \textup{(j7)}\,(vi), so that $D(V)$ is empty and $\Xi=1$. The total energy $x\mathrm{En}_\ell=\frac x2\langle F,\mathsf B^VF\rangle$ is at most $\frac x2\mathsf N_Z\cdot3\cdot16/x=24\mathsf N_Z$ per site (\Cref{Y-A-SetsOfRows}, second paragraph; item~3 of \S\ref{Y-G0-1}, \textup{(O3)}). The members are bounded below by $-0.61q$. Finally, $\mathrm{Haar}\{\operatorname{ang}\le\check\delta\}=\frac2\pi\int_0^{\check\delta}\sin^2\ge\check\delta^3/5$ for $\check\delta^2\le0.28$; since a site carries three bonds, the restriction costs $3(1.5\log x+\log5)\le4.5\log x+4.83$ per site.

\emph{Step 8: conclusion.} The factor $e^{\mathsf V}$ cancels. A cell contains $L^3$ sites, and $218.9q$ does not exceed $220M^{-3}$ per site. This proves the first inequality of \textup{(P$^\circ$)}; the second one is \textup{(j7)}\,(i).
\end{proof}

\begin{lemma}[Extension of a field that is tame on a box]\label{Y-P6-L4}
Let $V$ be a field with values in $G$, defined on the bonds of the box $\{r\le r_B\}$ and tame there. Suppose that $r_B\ge7$, $C_{\rm ext}r_B^2\varepsilon\le1$, and the side of $\mathbb T_k$ is at least $10r_B+2$. Then $V$ is the restriction of a field $\tilde V$ on $\mathbb T_k$ that is globally tame.
\end{lemma}

\begin{proof}
\emph{Step 1: the comb gauge.} Let $g(y)=V(\Gamma(0,y))$ and $V^g(b)=g(y)V(b)g(y')^{-1}=:e^{iA_b}$. Then $V^g_b$ is the holonomy of the loop $\Gamma(0,y).b.\Gamma(0,y')^{-1}$. When $b$ points along $e_3$, this loop is trivial; when $b$ points along $e_2$, the loop bounds a strip of $|y_3|$ plaquettes; and when $b$ points along $e_1$, the loop encloses two strips, one of $|y_2|$ and one of $|y_3|$ plaquettes. In each case no more than $2r_B$ plaquettes of the box are involved. This holonomy is a product of conjugates of the holonomies of plaquettes; as $\operatorname{ang}(gh)\le\operatorname{ang}\,g+\operatorname{ang}\,h$ and $\operatorname{ang}$ is a class function, it follows that $|A_b|\le a:=2r_B\varepsilon$. Since $C_{\rm ext}=32$ and $r_B\ge7$, the assumed bound gives $a\le1/(16r_B)<0.01$. The prescribed box radius satisfies $r_B\ge7$.

\emph{Step 2: the continuation.} Let $\mathrm{pr}$ be the map that replaces each coordinate of a site by the nearest value that this coordinate takes in the box; for a site $y$ let $n(y)$ be its sup distance from the box; and let $s(n):=\max(0,1-n/(4r_B))$. If $b$ is a bond outside the box, its image $\mathrm{pr}\,b$ is a bond or a single point; the former happens when $\mathrm{pr}$ leaves the coordinate along $b$ unchanged, and then $b$ has the same $n$ at both ends. We put $\tilde V_b:=\exp(is(n)A_{\mathrm{pr}\,b})$ in the first case and $\tilde V_b:=1$ in the second. All bonds with $n\ge4r_B$ then carry the value $1$, and the side, being at least $2(5r_B)+2$, leaves room for this on the torus.

\emph{Step 3: plaquettes outside the box.} The image of such a plaquette under $\mathrm{pr}$ is a plaquette, a bond or a point.
\begin{enumerate}
\item[\textup{(i)}] A plaquette. Here $n$ is constant. Write its four oriented link logarithms as $A_1,\ldots,A_4$, with $|A_i|\le a$, and put
\[
 F(sA)=-i\log\bigl(e^{isA_1}\cdots e^{isA_4}\bigr)
                  =\sum_{q\ge1}s^qF_q.
\]
We verify $|d^2F(sA)/ds^2|\le14a^2$ on $0\le s\le1$.
The six commutators in degree two give $|F_2|\le6a^2$.
For completeness, after $m$ factors write the first three
coefficients as $S_1,S_2,S_3$. Multiplication by $e^{isA}$,
followed by expansion of the logarithm through degree three,
gives the bounds
\[
 |S_2^{\rm new}|\le|S_2|+ma^2,\qquad
 |S_3^{\rm new}|\le|S_3|+a|S_2|
                             +\tfrac13(m^2+m)a^3,
\]
using $|S_1|\le ma$ and the bracket constant $2$.
The successive bounds for $|S_2|/a^2$ are $1,3,6$, and those for $|S_3|/a^3$ are $2/3,11/3,32/3$. Thus $|F_3|\le(32/3)a^3$.
More generally, put $S=\sum_i|A_i|\le4a$. The coefficients of
the logarithm are bounded by those of $H(Ss)$, where
$H(u)=-\log(2-e^u)$: the product minus $1$ is majorized
coefficient by coefficient by $e^{Ss}-1$, and then the
logarithm series is absolutely majorized by $H(Ss)$.
In particular, the remaining coefficients are bounded by
those of $H(4as)$.
Since all coefficients of $H$ are nonnegative,
\[
 \left|\frac{d^2}{ds^2}F(sA)\right|
 \le12a^2+64a^3+
       16a^2\bigl[H''(4a)-2-24a\bigr]
 <13.01a^2<14a^2.
\]
Here $H''(u)=2e^u/(2-e^u)^2$, and for $a\le0.01$ the
coefficient on the right is at most
$12+0.64+16[H''(0.04)-2-0.24]<13.01$.
Now let $f(s)=F(sA)-sF(A)$. It has $f(0)=f(1)=0$;
the Green function for the second derivative on $[0,1]$
therefore gives $|f(s)|\le7s(1-s)a^2$. Consequently
\[
 |F(sA)|\le s\varepsilon+(1-s)28r_B^2\varepsilon^2
                                                    \le\varepsilon.
\]
\item[\textup{(ii)}] A bond $b'$. Here the holonomy is $\exp(i(s-s')A_{b'})$ with $|n-n'|\le1$, and its angle is at most $a/(4r_B)=\varepsilon/2$.
\item[\textup{(iii)}] A point. Here the holonomy is $1$.
\end{enumerate}

\emph{Step 4: back to the original gauge.} We undo the comb gauge: the transformation is $g^{-1}$ at the sites of the box and $1$ at all other sites. Tameness is gauge invariant, and on the bonds of the box the resulting field restricts to $V$.
\end{proof}

\begin{lemma}[Mirrors and a comparison of sums]\label{Y-P6-L5}
At $j\in J(\mathsf b)$, on the board $B_{j-i_h}$, let $\check\theta_c$ be the map of the fold that takes the cell $c$ onto the cell $0$; in each coordinate it is given by $y_\mu\to y_\mu-c_\mu h$ when $c_\mu$ is even and by $y_\mu\to c_\mu h-y_\mu$ when $c_\mu$ is odd. Let $\check\Theta_cV:=V\circ\check\theta_c^{-1}$, which is the map $\check\Theta$ of \Cref{Y-P6-L1}\,(K) applied with $\check\theta=\check\theta_c^{-1}$. For a bounded function $f$ let $f_c:=f\circ\check\theta_c$, and let $\mathrm{Box}_c:=\check\theta_c^{-1}\mathrm{Box}$, with $B'_c$ and $B_{R,c}$ defined in the same manner. Accordingly, the copy $\Phi\circ\check\Theta_c$ of a function $\Phi$ of the cell $0$ is a function of the cell $c$, and $\mathrm{Box}_c$ is contained in the cell $c$. If the support of $f$ lies in the cell $0$, as it does for $f=f_*$, then the support of $f_c$ lies in the cell $c$. Let $\Lambda:=-x\mathrm{En}_\ell+\sum1(\bar E^a+\bar E^h)$ and $\Lambda_c:=\Lambda\circ\check\Theta_c$. With an upper index ${\rm tr}$, a sum over $X$ is restricted to the $X$ with $\#X\le n_0$, and the index ${\rm long}$ refers to all other $X$. Finally, let $\tau_t$ be the flow of \Cref{Y-P6-L10} and $\tau^c:=\check\Theta_c^{-1}\tau\,\check\Theta_c$.
\begin{enumerate}
\item[\textup{(a)}] $\rho_k\circ\check\Theta_c=\rho_k$, and $T_f(\check\Theta_cV)=T_{f_c}(V)$.
\item[\textup{(b)}] At globally tame $V$, $\Lambda_c=\Lambda$ and $c_f+A_f(\check\Theta_cV)=c_{f_c}+A_{f_c}(V)$.
\item[\textup{(c)}] Let $f=f_*$. Put $\Delta_c(V):=c_{f_c}+A^{\rm tr}_{f_c}(V)-c_f-A^{\rm tr}_f(\check\Theta_cV)$ and $d_c(V):=(\Lambda^{\rm tr}\circ\tau^c_t-\Lambda^{\rm tr})-(\Lambda_c^{\rm tr}\circ\tau^c_t-\Lambda_c^{\rm tr})$. Each of them depends on the bonds of $\mathrm{Box}_c$ alone. On fields that are tame on $\mathrm{Box}_c$, and in the case of $d_c$ satisfy moreover $\phi_p\le3\varepsilon/4+4t$ on $B'_c$, we have $\mathsf b|\Delta_c|,\ |d_c|\le\mathrm{ex}$. At every field, $\mathsf b|\Delta_c|\le4(C_Sp_0^2+1)$. On fields with $\phi_p\le3\varepsilon/4+4t$ on $B'_c$, $|d_c|\le1$.
\end{enumerate}
\end{lemma}

\begin{proof}
\emph{Step \textup{(a)}: the mirrors as maps of the kind \textup{(K)}.} The inverse map acts by $y_\mu\to y_\mu+c_\mu h$ if $c_\mu$ is even and by $y_\mu\to c_\mu h-y_\mu$ if $c_\mu$ is odd; so $\check\theta_c^{-1}$ is of the kind \textup{(K)}, with $n_\mu=c_\mu h$. \Cref{Y-P6-L1}\,(K), applied with $\check\theta=\check\theta_c^{-1}$, gives \textup{(a)}, because $f\circ\check\theta^{-1}=f\circ\check\theta_c=f_c$.

\emph{Step \textup{(b)}: globally tame fields.}
\begin{enumerate}
\item[\textup{1.}] If $D(V)$ is empty, then the empty family is the only family in $\Xi_{w(z)}$ (item~1 of \S\ref{Y-G0-1}, ``Members''). Hence $\Xi_{w(z)}(V)=1$ for all $z$, and $W_f(V)=0$.
\item[\textup{2.}] It follows that $\Lambda=\log\rho_k-\mathsf V$ and $c_f+A_f=T_f$ at such fields.
\item[\textup{3.}] The field $\check\Theta_cV$ is globally tame as well, and \textup{(a)} now gives \textup{(b)}.
\end{enumerate}

\emph{Step \textup{(c)}, dependence: the bonds that enter.}
\begin{enumerate}
\item[\textup{1.}] When a short $X$ contributes to $A^{\rm tr}_f$, the set $X^{+\check r}$ meets $\mathrm{supp}f$; and this support lies in $r\le0.84D_0+2\le R/2$.
\item[\textup{2.}] When it contributes to $\Lambda^{\rm tr}\circ\tau-\Lambda^{\rm tr}$, the set $X^{\approx}$ meets $B_R^{+1}$.
\item[\textup{3.}] The term depends on $X^{\approx}$, and every point of $X^{\approx}$ has $r$ at most $R+1+M(n_0+2r(x)+4)\le R+1+h/16+R\le r_B$; the three inequalities used are $Mn_0\le h/16$, $M(2r+4)\le R$ from \textup{(j7)}\,(ii), and $R\le h/16$.
\item[\textup{4.}] The moved terms of $\mathrm{En}_\ell$ depend on $B'$, which lies in $\mathrm{Box}$ since $2\rho+3\le R$.
\item[\textup{5.}] Under a mirror, the least corner of a cube or a plaquette moves by one unit; this is where the $+1$ in $r_B$ comes from.
\end{enumerate}

\emph{Step \textup{(c)}, smallness: fields tame on the box.}
\begin{enumerate}
\item[\textup{1.}] By \Cref{Y-P6-L4} the restriction of $V$ to $\mathrm{Box}_c$ extends to a field $\tilde V$; the side obeys $NL^j\ge2Lh\ge10r_B+2$ by \textup{(j3)}.
\item[\textup{2.}] Then $\Delta_c(V)=\Delta_c(\tilde V)$, and by \textup{(b)} this equals minus the long parts.
\item[\textup{3.}] Because $s_x$ lives on single cubes, the long parts consist of $\mathrm{rem}$ and $\delta^h$ only.
\item[\textup{4.}] They are bounded by \[\FitDisplay{\dimexpr\linewidth-1pt\relax}{2N_1\sum_{n>n_0}(26e)^{n-1}e^{-\kappa_1(n-1)}[2C_3K_5E_ap_0^3x^{-1/4}+0.12c_Q^{\rm su}K_5(Lx)^{1/2}c_h(x)/(E_ap_0(Lx))]/\mathsf b}\], where the bracket is that of the display of \Cref{Y-P6-L2}.
\item[\textup{5.}] As $\kappa_1-\log(26e)\ge2.7419$ and $n_0\ge4L\log x-1$, the bound is $x^{-10.9L}e^{2.75}$ times $2N_1[\,\cdot\,]/\mathsf b$, with the same bracket, and this is at most $\mathrm{ex}/\mathsf b$: the number $2N_1[\,\cdot\,]$ is at most $1$ by \textup{(E1)}~$\le c_F/12$.
\item[\textup{6.}] Turning to $d_c$: the flow $\tau^c_t$ moves only plaquettes touching $B_{R,c}$, each by at most $4t\sup|w|\le4t$, where $w$ is the form fixed for \Cref{Y-P6-L10} and not the third member. All sites of a moved plaquette lie in $B'_c$, where $\phi_p(\tilde V)=\phi_p(V)\le3\varepsilon/4+4t$ by the hypothesis of \textup{(c)}, since $B'_c\subset\mathrm{Box}_c$; at $\tau^c_t\tilde V$ such a plaquette therefore obeys $\phi_p\le3\varepsilon/4+8t\le\varepsilon$. The last inequality holds because $32x^{-2}\le p_0x^{-1/2}$, by $t=x^{-2}$ and \textup{(j7)}\,(vi), which gives $p_0\ge4$ and $\log x\ge3$. A plaquette that is not moved keeps $\phi_p(\tilde V)\le\varepsilon$, so $\tau^c_t\tilde V$ is globally tame.
\item[\textup{7.}] Applying \textup{(b)} both at $\tilde V$ and at $\tau^c_t\tilde V$, we see that the untruncated differences cancel. There remain four long sums, each bounded by $(2R/M+2r+3)^3(c_a+c_h)x^{-10.9L}e^{2.75}$. Since $2R/M+2r+3\le h<x^{1/4}$ (\Cref{Y-P6-L16}\,(b)), the four together are at most $\mathrm{ex}$.
\end{enumerate}

\emph{Step \textup{(c)}, rough bounds: all fields.}
\begin{enumerate}
\item[\textup{1.}] At $V=1$, \textup{(b)} gives $|c_{f_c}-c_f|\le2\sup|A|$, and so $\mathsf b|\Delta_c|\le4\mathsf b\sup|A|$.
\item[\textup{2.}] For $d_c$, let $\Lambda^a:=-x\mathrm{En}_\ell+\sum1\bar E^a$, so that $\Lambda=\Lambda^a+\sum1\bar E^h$. We split both differences, each into the $a$-part coming from $\Lambda^a$ and the $h$-part coming from $\sum1\bar E^h$. The index ${\rm tr}$ acts on the sums over $X$ alone.
\item[\textup{3.}] We estimate the two $a$-parts along the flow segments. Their moving plaquettes have angles at most $3\varepsilon/4+8t<\varepsilon$: indeed $32x^{-2}<p_0x^{-1/2}$ follows from $p_0\ge4$ and $\log x\ge3$. The masks of every varying energy row are therefore $1$, as in Step~2 of the proof of \Cref{Y-P6-L10}.

Write $D=D_w^\alpha$ in the cell's translated coordinates, and use the sums $W_0,W_1,W_2$ of \Cref{Y-P6-L10}. Steps~1--3 of the proof of \Cref{Y-P6-L12}, which do not use the Green-function identity, give
\[
 |xD\mathrm{En}_\ell|
       \le x\varepsilon N_{B0}W_1+C_2p_0^2R.
\]
The leading term is bounded by $\varepsilon N_{B0}W_1$ because the absolute row sums of $B$ are at most $N_{B0}$. The remainder includes both the curvature error and the derivatives of the row transports.

For the other difference, use the reflected oriented bonds and plaquettes, with the root translated to the origin. Reversing a bond turns a left-invariant derivative into a right-invariant one; it preserves its speed $|w_b|$ and the signed plaquette sum $\mathrm{d}w$. We bound the resulting conjugations directly. In the comb gauge each link of a plaquette based at radius $r$ has angle at most $2(r+2)\varepsilon$. A plaquette prefix uses at most four links, and conversion of a right derivative to a left derivative adds at most one endpoint link. Thus each derivative axis differs from $e_\alpha$ by at most $20(r+2)\varepsilon$. The logarithm differential of Step~1 of the proof of \Cref{Y-P6-L12}, with $\varepsilon\le0.01$, gives the safe bound
\[
 |DF_p-(\mathrm{d}w)_pe_\alpha|
       \le\varepsilon\{40(r+2)|w|_p+|\mathrm{d}w|_p\}.
\]
The scalar row matrix is a signed reindexing of $B$, so its absolute row-sum bound remains $N_{B0}$. The displacement sums use the pulled-back oriented plaquettes and their distinguished endpoints, so their path lengths and displacement counts are unchanged. A varying curvature plaquette lies within radius $R+1$, so a retained row has center at radius at most $R+\rho+1$; its paths, plaquettes and mask lie within $R+2\rho+3$, hence in $B'$. If the least base site changes under reflection, its sup radius changes by at most one. The added unit in $r+2$ and this possible change are both covered by
\[
 W_2\ \hbox{replaced by }W_2+8W_0,
 \qquad
 V_1\ \hbox{replaced by }V_1+W_1,
 \qquad V_1=\sum_p(r_p+1)|(\mathrm{d}w_R)_p|.
\]
Here $\sum_p|w|_p=4W_0$. In the notation of Step~3 there, the row remainder is thus at most
\begin{align*}
 p_0^2\bigl[&N_{B0}\{40(W_2+8W_0)+W_1\}+18N_{B1}W_0\\
       &+24N_{B1}(V_1+W_1)+72N_{B2}W_1\bigr]
 \le400(N_{B0}+N_{B1}+N_{B2})r_1^4p_0^2R.
\end{align*}
The bare plaquette remainder is at most $12000r_1^3p_0^2R$. Indeed $W_0\le r_1^3R$ for $R\ge10^6$, $W_2\le289r_1^3R$, and $W_1\le400r_1^4\le r_1^3R$. The two remainders together have coefficient at most
\[
 12000r_1^3+400(N_{B0}+N_{B1}+N_{B2})r_1^4
        <10^{665}<C_2=10^{670}.
\]
Thus the same energy bound holds for the reflected orientation.

The first Cauchy estimate of \Cref{Y-P6-L10} also holds on these slightly longer segments. In Step~7.1 of its proof, replace the starting plaquette bound by $\varepsilon$; the domain estimate becomes
$\varepsilon+e^{1/2}\varepsilon/4<2\varepsilon$. Its other domain estimates are unchanged. For the reflected orientation use
$\delta_p=|\mathrm{d}w|_p+40(r_p+2)\varepsilon|w|_p$ and the same definition of $s_X$. The frames remain constant along the conjugated flow, as follows by conjugating Step~7.0. The enlarged weighted correction is less than $9$, since
\[
 W_0\le78r_1^3(1+\log R)\le r_1^3R,\qquad
 \varepsilon\le(288R^2)^{-1},\qquad R\ge6r_1^3.
\]
Indeed $5\varepsilon W_2\le1$ gives $40\varepsilon W_2\le8$, while $320\varepsilon W_0\le320/1728<1$. Hence
$W_1+40\varepsilon(W_2+8W_0)\le400r_1^4+9\le401r_1^4$, retaining the constants $1605r_1^4$ and $C_a(x)$ of that proof. Restricting the sum to short $X$ only decreases its positive majorant.

For either $a$-part, we therefore have
\[
 \sup|D\Lambda^{a,{\rm tr}}|
 \le x\varepsilon N_{B0}W_1+(C_2R+C_a(x))p_0^2
 \le2N_{B0}W_*p_0x^{1/2}.
\]
Indeed $W_1\le0.4W_*$, while \textup{(E3)}$\le c_F/12$ gives
$(C_2R+C_a)p_0^2\le c_Fx^{1/2}/(60C_Sp_0^3)$.
Multiplying by $t$ bounds the corresponding difference.
\item[\textup{4.}] The $h$-part is at most $2.5(2R/M+2r+3)^3c_h$.
\item[\textup{5.}] By \textup{(E6)}~$\le c_F/12$, both parts are at most $\frac14$, because $C_6\ge N_{B0}W_*$.
\end{enumerate}
\end{proof}

\begin{lemma}[Jets of many sources with a free last ratio]\label{Y-P6-L6}
In the setting of \Cref{Y-P1-DATAo} with $Q$ sources, denoted there by $g_t$ and here by $f_i$, each with sup norm $\|f_i\|_\infty\le1$, let $1\le k\le K-1$, let $\varrho$ be any number $\ge1$ (the lemma is used with $\varrho=x^{1/2}p_0^6$), and let $\mathrm{sep}_k\ge8\log(4\varrho L^{1/2})$. Here $\varrho L^{1/2}\ge75$, because $L>1.2\cdot10^{13}$.
\begin{enumerate}
\item[\textup{(A)}] In the strong unit of each of the three members at the rate $\kappa+\frac14$, \[\sup_{(X,V)}\sum_{|S|\ge2}(\varrho\,\mathfrak r^\circ_k)^{|S|}|\delta_S(k)(X,V)|\le0.03.\]
\item[\textup{(B)}] For each $i$, \[\sum_{S\ni i,|S|\ge2}(\varrho\,\mathfrak r^\circ_k)^{|S|}\sum_{k'<k}\sum_Y|\partial_Sv_{k'}(Y;u_{k'}^{\mathrm{pol}})(0)|\le E_{\mathrm{vac}},\] where $E_{\mathrm{vac}}:=\sum_{m\ge1}N_1(k-m)2C_gc_a(x_{k-m})e^{1/4}\exp\{-(\kappa-\frac14)(L^{m-1}\mathrm{sep}_k+1)\}$.
\end{enumerate}
Neither of the two bounds depends on $Q$, $N$ or $K$.
\end{lemma}

\begin{proof}
Apply \Cref{Y-P1-Jo}\,\textup{(J5$^\circ$)}\,(d) with
$R_{k'}=\mathfrak r^\circ_{k'}$ for $k'<k$ and
$R_k=\varrho\mathfrak r^\circ_k$.
Its rates are
\[
 \varkappa_{k'}=\kappa+\tfrac14+
 \frac{2\log(4\varrho L^{1/2})}{\mathrm{sep}_k}
 +\sum_{k'\le i<k-1}
       \frac{2\log(4L^{1/2})}{\mathrm{sep}_{i+1}}
 \quad(k'<k),\qquad
 \varkappa_k=\kappa+\tfrac14.
\]
The assumed separation is precisely the hypothesis of that clause,
which proves admissibility of these data and gives \textup{(A)}.
The additional $1/4$ in the rates is retained: in
\textup{(J8$^\circ$)}\,(a$'$), the loss
$\lambda\le1/4$ leaves $\varkappa_{k'}-\lambda\ge\kappa$.
That clause therefore gives \textup{(B)} with the stated
$E_{\rm vac}$, since
$N_1(k-m)=N_1^{[K-k+m]}$ and $x_{k-m}=L^mx_k$.
Neither result imposes an upper bound on the depth or on $\varrho$.

\emph{Step: the size of $E_{\mathrm{vac}}$.} Throughout this step we take the sources $f_c$ of \Cref{Y-P6-L8}, $\check\delta=(h/2)L^{-j}$, $\check\rho=(D_0/8)L^{-j}$, $k=K-j$ and $x=x_k$, and \textup{(X)} is in force. We have $N_1(k-m)\le((D_0/4M)L^{m-1}+3+2\check r(L^{m-1}x))^3$ and $\mathrm{sep}_k\ge\lfloor h/2M\rfloor-2\check r(x)-2\ge31L\log x$. The latter holds because \textup{(M2)} gives $r(x)\le1+\log x/768$, hence $2\check r(x)+2\le10+\log x/384$, while $h/2M\ge32L\log x$ (\Cref{Y-P6-L16}\,(b)). Moreover $C_g=263L^3$ (item~0\,(1) of \S\ref{Y-H5-1}), and $x^{8L}\ge L^{8L}$. As \textup{(M2)} is a condition on the letters alone, it gives $r(y)\le1+\log y/768$ at every $y>1$; at $y=L^{m-1}x$ this yields $3+2\check r(L^{m-1}x)\le11+(\log x+(m-1)\log L)/384$. We have $(\kappa-\frac14)\cdot31\ge178.25$, and $c_a(x_{k-m})=c_a(L^mx)\le1/32$ by \textup{(X17)} (an entry of \textup{(O13)}, item~1 of \S\ref{Y-J-1}). Hence the $m$-th term of $E_{\mathrm{vac}}$ is at most \[((D_0/4M)L^{m-1}+11+(\log x+(m-1)\log L)/384)^3\cdot526L^3e^{1/4}x^{-178.25L^m}.\] By $x^{8L}\ge L^{8L}$, the term with $m=1$ is at most $\frac12x^{-170L}$, and the terms with $m\ge2$ form a smaller geometric tail. Summing over $m\ge1$, we obtain $E_{\mathrm{vac}}\le x^{-170L}\le\mathrm{ex}$.
\end{proof}

\begin{lemma}[A positive majorant and its integral]\label{Y-P6-L7}
Let $\hat w(X,V)\ge0$ be Borel and gauge invariant, depending on $V$ only through its values on $X^{\approx\approx}$. Suppose that $\hat w=0$ unless $X$ holds a plaquette of $D(V)$, and that $\hat w\le2.6e^{-\kappa(|X|_c-1)}$. Let $\hat u:=(\bar E_k^a,\bar E_k^h,\hat w)$. Then \[0\le\int\rho[\hat u;\ell_k]dV\le(1+C_{\rm vol}(N,\mathsf b))\int\rho[\bar u_k;\ell_k]dV,\] and $(1+C_{\rm vol})^{1/Q}\le\exp\{c_{\mathrm{vol}}(\mathsf b)h^3L^{-3j}\}\le e$ for $Q=(NL^j/h)^3$, at $j\in J(\mathsf b)$.
\end{lemma}

\begin{proof}
\emph{Step 1: the class of $\hat u$.} The triple $\hat u$ is real and belongs to $M_{2.6}(x;\kappa,c_h)$: each of its $a$- and $h$-members has a single cap at $\kappa_1$, and therefore one at $\kappa$ as well (item~0 of \S\ref{Y-H1-1}, ``Pairs''). As $\hat w\ge0$, also $\rho[\hat u]\ge0$. Membership in $M_A$ does not depend on any translation covariance of $w$ (item~1 of \S\ref{Y-G0-1}, ``Members'').

\emph{Step 2: the data along the run.} For $k\le i<K$, \Cref{Y-H4-U4}\,(b) gives \textup{(SD)} at $(x_i;\ell_i)$; moreover $x_i\ge L\mathsf b\ge Lx_s\ge x_9^U$ (item~(1)~(O13) of \S\ref{Y-H0-1}).

\emph{Step 3: the chain of members.} Let $\hat u_k:=\hat u$, and define $\hat u_{i+1}:=F_{\ell_i}(\hat u_i)$ recursively. From \textup{(9a)}, applied at every $V'$, together with $\int\rho'=\int\rho$ (item~3 of \S\ref{Y-J-1}, \textup{(R9)}, and item~2 of the same section, ``The step''), we get $\int\rho[\hat u_i;\ell_i]=e^{\mathsf V'(\hat u_i)}\int\rho[\hat u_{i+1};\ell_{i+1}]$.

\emph{Step 4: the classes along the chain.} \Cref{Y-H5-Af}\,(i), applied with $B=2.6$ and $p=(\kappa,c_h)$, gives $\hat u_{k+1}\in M_1(x_{k+1};p)$ with the strong cap. At each later level $i$ the same theorem, applied to $\hat u_i$, gives $\hat u_{i+1}\in M_1(x_{i+1};p)$ with the strong cap.

\emph{Step 5: the vacuum numbers.} According to \textup{(f1)} (\Cref{Y-H1-U1}), neither $E'^a$ nor $\mathsf V'$ depends on $E^h$ or on $w$. The $a$-members of $\hat u_k$ and $\bar u_k$ coincide; by induction they coincide at every $i$, and so $\mathsf V'(\hat u_i)=\mathsf V'(\bar u_i)$.

\emph{Step 6: exact cancellation.} Since $\bar u_{i+1}=F_{\ell_i}(\bar u_i)$ (\Cref{Y-H4-U4}\,(h)), the vacuum factors cancel exactly, and \[\int\rho[\hat u;\ell_k]/\int\rho[\bar u_k;\ell_k]=\int\rho[\hat u_K;\ell_K]/\int\rho[\bar u_K;\ell_K].\]

\emph{Step 7: the last level.} For every $K\ge1$ the data \textup{(B)} of \Cref{Y-H8-ThBC} are satisfied for $x=\mathsf b$ and the pair $(\kappa,c_h)$, by Step~2 of the proof of \Cref{Y-H9-T}, and $|\ell_{K,y}(p)|\le K_*e^{-\lambda_*|y-p|}$ by the same step, ``The last integral''. Here $u_1=u_3=\bar u_K$, while $u_2=\hat u_K$ arises as an output of \Cref{Y-H-C9} at $A\le2.6\le6$: since $j\ge1$, we have $k=K-j<K$, so the chain of Step~3 has at least one step, and $\hat u_K=F_{\ell_{K-1}}(\hat u_{K-1})$ with $\hat u_{K-1}\in M_A$ and $A\le2.6$ (Steps 3 and 4). We are in case \textup{(ii)}, with $a_w=1$ and $a_E\le2$ (\Cref{Y-H8-ThBC}\,(b$^C$)\,(ii)). Moreover $\mathsf D_{23}\le2$. The sentence of \textup{(c$^C$)} before \textup{(S)} in \Cref{Y-H8-ThBC} then yields $\int|\rho[\hat u_K]-\rho[\bar u_K]|\le\frac12C_{\rm vol}\cdot2\cdot I_1$.

\emph{Step 8: the root.} We have \[\log(1+C_{\rm vol})\le\log(2C_{\rm vol})=\log(8q_N)+220q_N+[\frac32\mathsf N_Zp_0(\mathsf b)^2+3\log(384/\varepsilon(\mathsf b)^3)]N^3\le c_{\mathrm{vol}}(\mathsf b)N^3,\] since $q_N=(N/M)^3$ (item~1 of \S\ref{Y-H8-1}: $q_N$; $C_{\rm vol}$) and $\log(8q)\le2.08q\le5q$ for $q\ge1$. It remains to divide by $Q$ and to apply \textup{(j2)}.
\end{proof}

\begin{lemma}[The part of the source carried by the weights]\label{Y-P6-L8} In the setting of \Cref{Y-P6-N}, at $j\in J(\mathsf b)$, the following hold. \textup{(W$^\circ$)} In norm form: on the board $B_{j-i_h}$ and for every constant $g\in[0,(9/8)p_0^2]$, \[ \|E[(G_--g)(\mathcal O_K(f_*)-c_f-A_f^{\rm tr}(V))\mid U]\|_{\rm ch}\le(9/8)p_0^2\cdot2e^{2.2}/r_z. \]

As a consequence, for every such $g$, and so in particular for $g=EG_-$, \[ \mathsf b\,|E[(G_--g)W_{f_*}]|\le(9/8)p_0^2\cdot2e^{2.2}\mathsf b/r_z+(9/4)p_0^2\,\mathrm{ex}\le\textup{(E5)}, \] where we define \[ \textup{(E5)}:=5.81\,c_Q^{\rm su}K_5L^{1/2}/(E_ap_0^4p_0(Lx))+(9/4)p_0^2\,\mathrm{ex}. \] Here $\mathsf b/r_z=c_Q^{\rm su}K_5(Lx)^{1/2}/(3.5\varrho E_ap_0(Lx))$. The bound is independent of $N$ and of $K$.\end{lemma}

\begin{proof} We write $f$ for $f_*$.

\emph{Step 1: the cell function.} We work on the board $B_{j-i_h}$, which has $Q=(NL^j/h)^3$ cells. The box $\mathrm{Box}$ lies in the cell $0$ and is centred in it ($r_B\le(h-1)/2$). We set $\Phi:=(G_--g)(\mathcal O_K(f)-c_f-A_f^{\rm tr}(V))$; it is a bounded Borel function of the bonds of $U$ and of $V^{(k)}$ that belong to this cell. \Cref{Y-P6-L2} gives $E\Phi=E[(G_--g)(W_f+A_f^{\rm long})]$, because by that lemma $T_f=c_f+A_f+W_f$ is a version of $E[\mathcal O_K(f)|V]$, and because $G_-$ is a function of $V$; the indices ${\rm tr}$ and ${\rm long}$ are those of the statement of \Cref{Y-P6-L5}. The bound $\mathsf b\sup|A^{\rm long}|\le\mathrm{ex}$ follows from the sizes in \Cref{Y-P6-L2} and from items 3 to 5 of Step~\textup{(c)}, smallness, of the proof of \Cref{Y-P6-L5}. Indeed, $A^{\rm long}$ is the long part of $\mathrm{Rem}$; as $s_x$ lives on single cubes, only $\mathrm{rem}$ and $\delta^h$ have long parts, and \[ \mathsf b\sup|A^{\rm long}|\le N_1\sum_{n>n_0}(26e)^{n-1}e^{-\kappa_1(n-1)}\,\mathsf b\,(\text{sizes of \Cref{Y-P6-L2}})\le\mathrm{ex}, \] where $N_1$ counts the anchors, $n_0\ge4L\log x-1$, and the prefactor is at most $1$ by $\textup{(E1)}\le c_F/12$.

\emph{Step 2: the chessboard estimate.} By \Cref{Y-P6-L1}\,\textup{(I)}, $E[\Phi|U]$ is a function of the cell. \Cref{support:4-15}, with the function $1$ in all other cells, yields \[ |E\Phi|^Q\le\langle\operatorname{Fold}E[\Phi|U]\rangle=\|E[\Phi|U]\|_{\rm ch}^Q. \] Steps 3 to 8 bound the fold, and in doing so they prove \textup{(W$^\circ$)}. By \textup{(K)} and \textup{(I)} of \Cref{Y-P6-L1}, the fold equals $E\prod_c\Phi\circ\check\Theta_c$.

\emph{Step 3: the copies.} For every cell $c$ we have $\Phi\circ\check\Theta_c=(G_c-g)(\mathcal O_K(f_c)-a_c(V))$, where $G_c=G_-\circ\check\Theta_c$ and \[ a_c=c_f+A_f^{\rm tr}(\check\Theta_cV)=c_{f_c}+A^{\rm tr}_{f_c}(V)-\Delta_c(V). \] The functions $f_c$ are themselves Wilson sources, with $\sup|f_c|\le1$, and they enter as $\mathcal O_K(f_c)$; each of them is supported within $D_0$ of a point, and the supports of two distinct ones are at least $h-2D_0\ge h/2$ apart. Hence the sources $f_c$ satisfy the conditions of the setting of \Cref{Y-P1-DATAo}, with $\check\delta=(h/2)L^{-j}$ and $\check\rho=(D_0/8)L^{-j}$, and $\mathrm{sep}_k\ge\lfloor h/2M\rfloor-2\check r(x)-2$. The first line of \textup{(j4)} then gives the hypothesis of \Cref{Y-P6-L6}.

\emph{Step 4: conditioning.} Each $G_c$ is a function of $V$. Moreover \[ E[\prod_c(\mathcal O_K(f_c)-a_c)\mid V]\,\rho_k(V)=\partial_1\cdots\partial_Q[e^{-\sum z_ca_c(V)}\rho_k^z(V)](0). \] We apply \Cref{Y-P1-Jo}\,\textup{(J7$^\circ$)}\,\textup{(iii)} at the level $k$ and at this $V$, with $\gamma_c^{(k)}=c_{f_c}$, modulo $z_c^2$ (the congruence of item~1 of \S\ref{Y-H8-1}); it shows that the bracket equals \[ e^{\mathsf V+\Lambda(V)}\exp\{W_2(z)+\mathrm{Mix}(z,V)\}\prod_c(1+z_c(\Delta_c+A^{\rm long}_{f_c})(V))\,\Xi_{w^{\rm pol}(z)}(V). \] Here $W_2$ gathers the vacuum jets with $|S|\ge2$ of the levels $k'<k$, and $\mathrm{Mix}=\sum_X1_{X^\approx}\sum_{|S|\ge2}z^S(\delta_S^a+\delta_S^h)$. The terms of a single source, $c_{f_c}+A^{\rm tr}_{f_c}$, cancel identically against $a_c$.

\emph{Step 5: the majorant.} The function $\Xi_w$ is a polynomial in the values of $w$, and its coefficients are non-negative. Hence the modulus of the coefficient of $z_1\cdots z_Q$ is at most $r_z^{-Q}$ times the value at $z_c=r_z$ of the majorant series, in which every coefficient is replaced by its modulus. For the majorants in the exponent, we have $\hat W_2\le(Q/2)E_{\rm vac}\le Q\,\mathrm{ex}$: the first inequality is \Cref{Y-P6-L6}\,\textup{(B)}, and the second follows from the step on the size of $E_{\mathrm{vac}}$ in the proof of \Cref{Y-P6-L6}, which gives $E_{\mathrm{vac}}\le x^{-170L}\le\mathrm{ex}$. Next, \Cref{Y-P6-L6}\,\textup{(A)} with \Cref{Y-P1-Jo}\,\textup{(J2$^\circ$)} gives \[ \widehat{\mathrm{Mix}}\le QN_1\,0.03(c_a+c_h)\sum_{n\ge\mathrm{sep}_k}(26e)^{n-1}e^{-(\kappa+1/4)(n-1)}\le Q\,\mathrm{ex}. \] We put \[ \hat w:=[\,|\bar w|+\sum_Sr_z^{|S|}|\delta_S^w|\,]\,1\{X\text{ holds a plaquette of }D(V)\}. \] \Cref{Y-P1-Jo}\,\textup{(J5$^\circ$)}\,\textup{(a$'$)}, taken value by value, gives $\sum_cr_z|\delta_c^w|\le0.2\,\mathrm{amp}(Lx)r_z=0.7\varrho$ strong units at $\kappa_1$. So $\hat w\le(1.03+0.7\varrho)$ strong units at $\kappa+1/4$.

\emph{Step 6: the factors.} Let $\mathsf y:=(r_z/\mathsf b)(4C_Sp_0^2+5)$; then $\mathsf y\le x$, by \textup{(j6)}. By \Cref{Y-P6-L5}\,\textup{(c)}, on fields tame on $\mathrm{Box}_c$ we have $1+r_z|\Delta_c+A^{\rm long}|\le1+(r_z/\mathsf b)2\,\mathrm{ex}\le1+2x^{1-8L}$, the last step by \textup{(j6)}, while on the other fields this quantity is at most $1+\mathsf y\,1\{D(V)$ meets $\mathrm{Box}_c\}$. A field fails to be tame on $\mathrm{Box}_c$ exactly when some plaquette of $\mathrm{Box}_c$ lies in $D(V)$ (item~1 of \S\ref{Y-G0-1}, on cubes). Let $S$ be a set of cells. On the event that $D(V)$ meets $\mathrm{Box}_c$ for all $c\in S$, the covering clause shows that in every family there is, for each $c\in S$, a member holding a site of $\mathrm{Box}_c$. Consequently $\prod_{c\in S}1\{\ldots\}\,\Xi_{\hat w}\le x^{-|S|}\Xi_{\hat w'}$, where $\hat w':=\hat w\,x^{m(X)}$ and $m(X):=\#\{c:X$ meets $\mathrm{Box}_c\}$. Finally, $\sum_S(\mathsf y/x)^{|S|}\le2^Q$.

\emph{Step 7: the cap of $\hat w'$.} At a set $X$ holding a rough plaquette, a strong unit at a rate is at most $\eta_w(x)e^{-{\rm rate}\,(|X|_c-1)}$. If $m\le1$, then $(1.03+0.7\varrho)x\,\eta_w\le x^{10}\eta_w\le1$, by \textup{(j6)}. If $m\ge2$, the boxes are more than $h-2r_B>h/2$ apart, so that their $h/4$-neighbourhoods are disjoint; the connected hull of $X$ leaves each of these neighbourhoods, and therefore $|X|_c\ge mh/(4M)$. The margin $1/4$ in the rate then gives $e^{1/4}(xe^{-h/(16M)})^m\le1$, because $h\ge64LM\log x$. Thus $\hat w'\le2.6e^{-\kappa(|X|_c-1)}$; moreover $\hat w'$ is Borel, gauge invariant and non-negative, and it is a function of $V$ on $X^{\approx\approx}$.

\emph{Step 8: use of \Cref{Y-P6-L7}.} \Cref{Y-P6-L7} now gives \[ \langle\operatorname{Fold}E[\Phi|U]\rangle=\|E[\Phi|U]\|_{\rm ch}^Q\le((9/8)p_0^2)^Qr_z^{-Q}2^Qe^{Q(2\mathrm{ex}+2x^{1-8L})}(1+C_{\rm vol}), \] and by Step~2 the right side also bounds $|E\Phi|^Q$. Taking the $Q$-th root we get $\|E[\Phi|U]\|_{\rm ch}\le(9/8)p_0^2r_z^{-1}2e^{1.01}$, and the same bound for $|E\Phi|$. As $e^{1.01}\le e^{2.2}$, the bound on $\|E[\Phi|U]\|_{\rm ch}$ implies \textup{(W$^\circ$)} of the statement; and as $(9/8)\cdot2e^{2.2}/3.5<5.802\le5.81$, the value of $\mathsf b/r_z$ given in the statement yields the first term of \textup{(E5)}.\end{proof}

\begin{remark}[One kind of source in each cell] The witness $G_c$ carries no source. It is taken out as a supremum after the conditioning, and this is legitimate because the bound of Step~5 of the proof of \Cref{Y-P6-L8} holds pointwise in $V$ and is positive.\end{remark}

\begin{remark}[The role of the subtracted constant] The proof of \Cref{Y-P6-L8} depends on $g$ only through the bound $\sup|G_c-g|\le(9/8)p_0^2$, which enters in Step~8. The identity of Step~1 holds for every constant $g$.\end{remark}

For \Cref{Y-P6-L9} we fix the following data. Let $w:=w_R$ be the form of \Cref{Y-P6-L11}, and let $\tau_t$, $0<t\le x^{-2}$, be the flow of $D_w^\alpha$ of \Cref{Y-P6-L10}. In Lems.~\ref{Y-P6-L9} and~\ref{Y-P6-L10} and their proofs, $w$ denotes the $w$-member $\bar w$ of $\bar u_k$ where it appears as an index of $\Xi$ or carries the arguments $(X,V)$. The lemma depends on $w$ only through these properties: $w$ is real, it is carried by the bonds of $B_R$, $|w|\le1$, and the two bounds on $d_c$ in \Cref{Y-P6-L5}\,\textup{(c)} hold. Let $\Phi$ be a bounded Borel function that depends on the bonds of $B'$ and is carried by $\{\phi_p\le3\varepsilon/4$ on $B'\}$. We split $\Lambda=\Lambda^a+H$, with $H=\sum1\bar E^h$, and we put $\Psi:=(\Phi e^{\Lambda^a})\circ\tau_t\,e^{-\Lambda^a}$.

\begin{lemma}[The defect of the integration by parts through the weights]\label{Y-P6-L9} In the setting of \Cref{Y-P6-N}, at $j\in J(\mathsf b)$, and with the data fixed before the lemma, the following hold. \textup{(X$^\circ$)} In norm form: on the board $B_{j-i_h}$, \[ \|E[\Phi^{\rm tr}_{\rm cell}\mid U]\|_{\rm ch}\le\sup|\Psi|\cdot2e^{2.2}/(r_\eta-1), \] where $\Phi^{\rm tr}_{\rm cell}:=\Psi^{\rm tr}\circ\tau_{-t}\,e^{\Lambda^{\rm tr}\circ\tau_{-t}-\Lambda^{\rm tr}}-\Psi^{\rm tr}$, and $\Psi^{\rm tr}$ is defined as $\Psi$ is, but with $\Lambda^{a,{\rm tr}}$. \textup{Claim.} With $r_\eta=x^4$, we have \[ |\int\Psi e^{\Lambda}[\Xi_w\circ\tau_t-\Xi_w]\,dV|\le\sup|\Psi|\,(2e^{2.2}/(r_\eta-1)+\mathrm{ex})\int e^{\Lambda}\Xi_w, \] and the bound is independent of $N$ and of $K$.\end{lemma}

\begin{proof} \emph{Step 1: a mean.} The flow $\tau_t$ preserves Haar measure (Step~1 of the proof of \Cref{Y-P6-L10}). Hence the integral on the left side of the ``Claim'' equals $E_\nu[\Phi_{\rm cell}]\int e^\Lambda\Xi_w$, with $\Phi_{\rm cell}=\Psi\circ\tau_{-t}\,e^{\Lambda\circ\tau_{-t}-\Lambda}-\Psi$. The function $\Psi$ is carried by $\{\phi_p\le3\varepsilon/4+4t$ on $B'\}$. We replace $\Lambda$ and $\Lambda^a$ by $\Lambda^{\rm tr}$ and $\Lambda^{a,{\rm tr}}$ once, before the chessboard estimate. We justify its pointwise cost and the comparison of the two suprema. Put $H=\sum1\bar E^h$ and let $\delta_sA=A\circ\tau_s-A$. The long-polymer estimate in Step~\textup{(c)} of the proof of \Cref{Y-P6-L5} gives, for $s=\pm t$ on the supports at issue,
\[
 |\delta_s\Lambda^{a,{\rm long}}|
       +|\delta_sH^{\rm long}|
 \le T:=2(2R/M+2r+3)^3(c_a+c_h)e^{2.75}x^{-10.9L}
 \le\mathrm{ex}/100.
\]
For the last inequality use $2R/M+2r+3\le h<x^{1/4}$,
$c_a+c_h\le1/2$, $L\ge13$, $\log x\ge3$, and
$\mathrm{ex}=x^{-8L}$. In fact
$T/\mathrm{ex}\le e^{2.75}x^{3/4-2.9L}
\le e^{2.75+3(3/4-2.9\cdot13)}=e^{-108.1}<1/100$.
The energy itself has no long part.
The definitions give the exact identities
\[
 \Phi_{\rm cell}=\Phi e^{\delta_{-t}H}-\Psi,\qquad
 \Phi^{\rm tr}_{\rm cell}=\Phi e^{\delta_{-t}H^{\rm tr}}-\Psi^{\rm tr},
 \qquad
 \Psi^{\rm tr}=\Psi e^{-\delta_t\Lambda^{a,{\rm long}}}.
\]
The flow bound in the data of \Cref{Y-P6-L10} gives
$\sup|\Phi|\le e^{1/2}\sup|\Psi|$, and the rough $h$-part
bound of \Cref{Y-P6-L5} gives
$|\delta_{-t}H^{\rm tr}|\le1/4$ on $\operatorname{supp}\Phi$.
Consequently
\[
 |\Phi_{\rm cell}-\Phi^{\rm tr}_{\rm cell}|
 \le(e^{3/4}+1)(e^T-1)\sup|\Psi|
 \le\mathrm{ex}\sup|\Psi|,
 \qquad
 \sup|\Psi^{\rm tr}|\le(1+\mathrm{ex})\sup|\Psi|.
\]
Thus there is one pointwise truncation error, and the latter
comparison supplies the factor used in Step~7. The function $\Phi^{\rm tr}_{\rm cell}$ is then a function of the field on $\mathrm{Box}$ alone, by items 2 to 4 of Step~\textup{(c)}, dependence, of the proof of \Cref{Y-P6-L5}: there $\Lambda^{\rm tr}\circ\tau-\Lambda^{\rm tr}$ and $\Lambda^{a,{\rm tr}}\circ\tau-\Lambda^{a,{\rm tr}}$ depend on bonds with $r\le r_B$, and the moved terms of $\mathrm{En}_\ell$ depend on $B'$, which lies in $\mathrm{Box}$.

\emph{Step 2: the chessboard estimate on $B_{j-i_h}$.} By \textup{(I)} and \textup{(K)} of \Cref{Y-P6-L1} and by \Cref{support:4-15}, \[ |E\Phi^{\rm tr}_{\rm cell}|^Q\le\|E[\Phi^{\rm tr}_{\rm cell}|U]\|_{\rm ch}^Q=E\prod_c[\Psi_c\circ\tau^c_{-t}e^{a_c}-\Psi_c]. \] Steps 3 to 7 bound this fold, and so they prove \textup{(X$^\circ$)}. Here $\Psi_c=\Psi^{\rm tr}\circ\check\Theta_c$ and $a_c=\Lambda_c^{\rm tr}\circ\tau^c_{-t}-\Lambda_c^{\rm tr}$. The function $\Psi_c$ is bounded, and no regularity of it is at hand.

\emph{Step 3: the expansion.} We expand the product over the subsets $S'$ of the set of cells, and in the term of $S'$ we change the variable, $V\to\tau_t^{S'}V$. The flows commute, since each of them depends only on the bonds of its own $B_{R,c}$ and moves only these. The function $\Psi_c$ is independent of the bonds moved by $\tau^{c'}$ for $c'\ne c$. Hence \[ E\prod_c[\Psi_c\circ\tau^c_{-t}e^{a_c}-\Psi_c]=\int\prod\Psi_c\,e^\Lambda\sum_{S'}(-1)^{Q-|S'|}e^{\mathsf E_{S'}}\,\Xi_w\circ\tau_t^{S'}\big/\int e^\Lambda\Xi_w, \] with $\mathsf E_{S'}=\Lambda\circ\tau_t^{S'}-\Lambda-\sum_{c\in S'}(\Lambda_c^{\rm tr}\circ\tau_t^c-\Lambda_c^{\rm tr})$.

\emph{Step 4: the switch.} On the support of $\prod\Psi_c$, the plaquettes of $B'_c$ have angles at most $3\varepsilon/4+4t$. Along $\tau_s^{S'}$, $0\le s\le t$, their angles are at most $3\varepsilon/4+8t<\varepsilon$, by \textup{(j7)}\,\textup{(vi)}. Thus $D(V)$ and the angles of its plaquettes stay fixed, as do all the indicators and masks. The clip equals the varying curvature $F_p$ on each moved plaquette and is unchanged on the others. The difference $\mathrm{En}_\ell\circ\tau_t^{S'}-\mathrm{En}_\ell$ is additive over $c$, because each energy term that moves under $\tau^c$ depends on $B'_c$, and the other flows leave those bonds fixed. A short $X$ has $X^\approx$ of diameter at most $h/8$, so it touches at most one of the sets $B_{R,c}$. For a function $g$ of $V$ we put \[ g(\eta):=\sum_{S''}\prod_{c\in S''}\eta_c\prod_{c\notin S''}(1-\eta_c)\,g(\tau_t^{S''}V), \] with one variable $\eta_c$ for each cell $c$ for which $g$ depends on the bonds of $B_{R,c}$. Then $\mathsf E_{S'}$ is the value at $\eta=1_{S'}$ of $\sum_c\eta_cd_c+\mathsf L(\eta)$, where $\mathsf L(\eta):=\sum_{X\,{\rm long}}([1E(X)](\eta)-1E(X))$. We write $\mathsf L$ for $\mathsf L(\eta)$. We put \[ \mathsf G(\eta):=\int\prod\Psi_c\,e^\Lambda\prod_c(1+\eta_c(e^{d_c}-1))\,e^{\mathsf L(\eta)}\,\Xi_{w(\eta)}\,dV. \] This function is entire, it takes the right values at the vertices, and the families of $\Xi$ are the same at all vertices. Therefore $E\prod_c[\Psi_c\circ\tau^c_{-t}e^{a_c}-\Psi_c]=[\prod_c(\mathrm{ev}_1-\mathrm{ev}_0)\mathsf G]\big/\int e^\Lambda\Xi_w$.

\emph{Step 5: bounds for $|\eta_c|\le r_\eta$.} In this range, $|g(\eta)|\le(2r_\eta+1)^{\text{number of variables}}\sup|g|$. Let $\tau(X)$ be the set of cells $c$ such that $X^{\approx\approx}$ meets $B_{R,c}$. Then \textup{(j7)}\,\textup{(iv)} gives $\tau(X)\subset\{c:X$ meets $\mathrm{Box}_c\}$, hence $|\tau(X)|\le m(X)$. Next, \[ |\mathsf L|\le Q(2R/M+2r+3)^3\sum_{n>n_0}(26e)^{n-1}2(2r_\eta+1)^{\max(1,4Mn/h)}(c_a+c_h)e^{-\kappa_1(n-1)}\le Q\,\mathrm{ex}. \] Further, $|w(\eta)(X,V)|$ is at most $(2r_\eta+1)^{|\tau(X)|}$ times the strong cap at $\kappa_1$ at $V$. By \Cref{Y-P6-L5}\,\textup{(c)}, $|1+\eta_c(e^{d_c}-1)|\le1+2r_\eta\mathrm{ex}$ on fields tame on $\mathrm{Box}_c$, and $|1+\eta_c(e^{d_c}-1)|\le1+2r_\eta1\{D(V)$ meets $\mathrm{Box}_c\}$ on the others.

We now bound the product of these factors with $\Xi_{w(\eta)}$ by the argument of Step~6 of the proof of \Cref{Y-P6-L8}, with the following changes. The factor $2r_\eta$ takes the place of $\mathsf y$, and on tame fields the error $\mathrm{ex}':=2r_\eta\mathrm{ex}=2x^{4-8L}$ takes the place of $2x^{1-8L}$. The covering clause enters unchanged: on the event that $D(V)$ meets $\mathrm{Box}_c$ for all $c$ in a set $S$ of cells, every family has, for each $c\in S$, a member holding a site of $\mathrm{Box}_c$. The weight receives the factor $(2r_\eta)^{m(X)}$ in place of $x^{m(X)}$; this factor absorbs $(2r_\eta)^{|S|}\prod_{c\in S}1\{\ldots\}$, so that the sum over $S$ gives $2^Q$ in place of $\sum_S(\mathsf y/x)^{|S|}$. In this way the product is at most $e^{Q\mathrm{ex}'}2^Q\Xi_{\hat w'}$, with \[ \hat w':=(2r_\eta+1)^{|\tau(X)|}(2r_\eta)^{m(X)}\exp\{-c_\Xi x\varepsilon\sum_{D\cap X}\phi_p-\kappa_1(|X|_c-1)\}\,1\{X\text{ holds a plaquette of }D(V)\}. \] The weight $\hat w'$ is an explicit function of $V$ on $X$; so the bound is unaffected by the dependence of $w(X,\tau V)$ on comb paths outside $X^{\approx\approx}$. For the cap of $\hat w'$ there are two cases. If $m\le1$, then $(2x^4+1)^2\eta_w\le x^{10}\eta_w\le1$. If $m\ge2$, then $|X|_c\ge mh/(4M)$, and the margin $\kappa_1-\kappa=1$ absorbs $(3x^4)^{2m}$, because $h/(4M)\ge16L\log x$.

\emph{Step 6: a contour integral.} In each variable, \[ g(1)-g(0)=(1/2\pi i)\oint_{|\zeta|=r_\eta}g(\zeta)/(\zeta(\zeta-1))\,d\zeta, \] and the modulus of this is at most $\sup|g|/(r_\eta-1)$.

\emph{Step 7: the conclusion.} \Cref{Y-P6-L7}, followed by the $Q$-th root, gives $\|E[\Phi^{\rm tr}_{\rm cell}|U]\|_{\rm ch}\le\sup|\Psi|(1+\mathrm{ex})2e^{1.01}/(r_\eta-1)$, and by Step~2 the same estimate holds for $|E\Phi^{\rm tr}_{\rm cell}|$. As $(1+\mathrm{ex})e^{1.01}\le e^{2.2}$, the bound on the norm implies \textup{(X$^\circ$)}; and the bound on $|E\Phi^{\rm tr}_{\rm cell}|$, together with the single pointwise error $\sup|\Psi|\,\mathrm{ex}$ of Step~1, yields the ``Claim''.\end{proof}

For \Cref{Y-P6-L10} we fix the following data. Let $w$ be real on the bonds of $B_R$, with $|w|\le1$, and let $t\sup|D_w^\alpha\Lambda^a|\le1/2$ hold on $\{\phi_p\le3\varepsilon/4+4t$ on $B'\}$. The integration-by-parts lemma is applied to $w=w_R$ of \Cref{Y-P6-L11}. For this form, Step~0 of the proof of \Cref{Y-P6-L13} gives $\sup|D\Lambda^a|\le2p_0x^{1/2}$. Its proof also gives the same bound for $\Lambda^{a,{\rm tr}}$, since the estimate on the analytic corrections is a sum of absolute values. As \textup{(j6)} gives $p_0^{10}\le x$ and \textup{(j7)}\,\textup{(vi)} gives $\log x\ge3$, we have
\[
 t\sup|D\Lambda^a|\le2p_0x^{-3/2}
       \le2x^{-1.4}<\tfrac12.
\]
In the reflected copies used in Lems.~\ref{Y-P6-L5} and~\ref{Y-P6-L9}, the flow and the weight are conjugated together:
\[
 \tau_s^c=\check\Theta_c^{-1}\tau_s\check\Theta_c,
 \qquad \Lambda_c^a=\Lambda^a\circ\check\Theta_c,
 \qquad
 \frac{d}{ds}\Lambda_c^a(\tau_s^cV)\Big|_{s=0}
       =(D\Lambda^a)(\check\Theta_cV).
\]
Thus the same $2p_0x^{1/2}$ bound applies to these conjugated weights, and to their truncations, on the corresponding reflected tame sets.

The comparison of $\Lambda^{a,{\rm tr}}$ itself along $\tau_s^c$ in \Cref{Y-P6-L5} uses the separate bound $2N_{B0}W_*p_0x^{1/2}$ proved in its rough-bounds step. Its product with $t$ is also less than $1/2$: from \textup{(E6)}$\le c_F/12$ and $C_6\ge N_{B0}W_*$,
\[
 2N_{B0}W_*p_0x^{-3/2}
       \le\frac{c_F}{36p_0^4}<\tfrac12.
\]
This comparison does not apply the Green-function identity to a reflected form against the original row matrix. There is no circularity: Lems.~\ref{Y-P6-L11} and~\ref{Y-P6-L12}, and the part ``Besides'' of \Cref{Y-P6-L10}, use nothing of the ``Claim'' of \Cref{Y-P6-L10} or of \Cref{Y-P6-L9}. We put $D_w^\alpha:=\sum_bw_b\cdot($derivative of $V_b\to e^{isX_b}V_b$ at $s=0)$, where $X_b:=T_{\pi(b)}^{-1}e_\alpha$. Here the $e_\alpha$, $\alpha=1,2,3$, form an orthonormal basis of $(\mathfrak g,\langle Y,Y'\rangle=\frac12\operatorname{tr}YY')$, and we write $Y^\alpha=\langle Y,e_\alpha\rangle$; further, $\pi(b)$ is the base site of the bond $b$, namely its least end, so that $\pi(b)=z$ for $b=(z,z+e_\mu)$. Finally, $\Phi$ is $C^2$ along the flow, depends on the bonds of $B'$, and is carried by $\{\phi_p\le3\varepsilon/4$ on $B'\}$.

\begin{lemma}[Integration by parts for members that are only Borel]\label{Y-P6-L10} In the setting of \Cref{Y-P6-N}, at $j\in J(\mathsf b)$, the following hold. \textup{Claim.} With the data fixed before the lemma and at $t=x^{-2}$, the bound \[ |E_\nu[D_w^\alpha\Phi+\Phi D_w^\alpha\Lambda^a]|\le tB_2+t^{-1}e\sup|\Phi|\,(2.5(2R/M+2r(x)+3)^3c_h(x)+6e^{2.2}x^{-4}+2\mathrm{ex}) \] holds, where $B_2=\sup|e^{-\Lambda^a}D_w^2(\Phi e^{\Lambda^a})|$. \textup{Besides.} At $w=w_R$ and on the set $\{\phi_p\le3\varepsilon/4+4t$ on $B'\}$ that occurs in the data fixed before the lemma, two bounds of Cauchy type hold: first, $|D_w\sum1\bar E^a|\le C_a(x)p_0^2$, with $C_a(x)=2.5(2r(x)+1)^3E_aW_*$; second, $|D_w^2\sum1\bar E^a|\le12.5E_ap_0x^{1/2}(2r(x)+1)^3W_*^2$. \textup{The sums.} Let $W_0=\sum_b|w_b|$, $W_{00}=\sum_bw_b^2$, $W_1=\sum_p|\mathrm{d}w|_p$ and $W_2=\sum_p(r_p+1)|w|_p$, where $|w|_p$ denotes the sum over the four bonds of $p$. At $w=w_R$ these sums satisfy $W_0\le W_0(x)$, $W_{00}\le60r_1^3$, $W_1\le400r_1^4$, $W_2\le289r_1^3R$, and $5\varepsilon W_2\le1$.\end{lemma}

The available bound on $W_0$ grows logarithmically with $R$; the argument uses it only through the product $\Theta(x)$.

\begin{proof} \emph{Step 1: the divergence.} The simple path $\Gamma(0,\pi(b))$ can contain the bond $b$ only as its last step, run backwards. In that case $V(\Gamma)=V(\Gamma')V_b^{-1}$, hence $X_bV_b=V_bX'_b$ with $X'_b$ free of $V_b$; in the other case $X_b$ itself is free of $V_b$. In both cases the $b$-component is an invariant vector field in the variable $V_b$. So its divergence is $0$, and the flow $\tau_t$, smooth on a compact manifold, preserves Haar measure.

\emph{Step 2. Regularity along the relevant flow segments.}
For a configuration $V$, consider $\tau_sV$ for $0\le s\le t$.
If the segment misses $\operatorname{supp}\Phi$, the product
$\Phi(\tau_sV)e^{\Lambda^a(\tau_sV)-\Lambda^a(V)}$ vanishes there
with its derivatives. Otherwise some point of the segment lies in
$\operatorname{supp}\Phi$. The displacement bound $4t$ then puts the
whole segment in $\{\phi_p\le3\varepsilon/4+4t\text{ on }B'\}$.
Since $4t<\varepsilon/4$, all these plaquettes remain tame. Every $1_{X^\approx}$ is constant, because the plaquettes off $B'$ stay fixed. A row depending on a plaquette that moves, or having a path $\Gamma(y,\pi p)$ that crosses a moving bond, has $r(y)\le R+\rho+1$. By \textup{(j7)}\,\textup{(vii)}, the mask of such a row depends on the plaquettes within $\max\{\rho_0,\rho(x)+1\}=\rho(x)+1$ of $y$; recall that the mask has two radii, $\rho_0$ and $\rho(x)+1$ (item~2 of \S\ref{Y-J-1}, on the energy). These plaquettes belong to $B'$, so $m_y=1$ and $\zeta=F$. There $\mathrm{En}_\ell$ is a polynomial in $F$ and in the $T(y,\cdot)$, and the $E^a$ are holomorphic.

\emph{Step 3. Taylor's formula.} On each segment meeting the support,
apply Taylor's formula to the product from Step~2. Its second
$\tau$-derivative has modulus at most
$B_2e^{t\sup|D_w\Lambda^a|}$ by the bounds on the enlarged tame set.
On the remaining segments the product vanishes. Integrating this
pointwise estimate for $J(t)=E_\nu[\Phi\circ\tau_t\,
 e^{\Lambda^a\circ\tau_t-\Lambda^a}]$ gives \[ |J(t)-J(0)-tJ'(0)|\le(t^2/2)B_2e^{t\sup|D\Lambda^a|}\le t^2B_2, \] by the data, as $e^{1/2}<2$.

\emph{Step 4: a split.} We put $Z:=\int e^\Lambda\Xi_w\,dV$. By Step~1, $J(0)=\int(\Phi e^{\Lambda^a}e^H\Xi_w)\circ\tau_t/Z$. So \[ J(t)-J(0)=\int(\Phi e^{\Lambda^a})\circ\tau_t\,[e^H\Xi_w-(e^H\Xi_w)\circ\tau_t]/Z, \] and this is \[ -\int\Phi\circ\tau_t(1-e^{H-H\circ\tau_t})\,(e^\Lambda\Xi_w)\circ\tau_t/Z-\int\Psi e^\Lambda[\Xi_w\circ\tau_t-\Xi_w]/Z. \]

\emph{Step 5: the first term.} Here $|H-H\circ\tau_t|\le2\cdot1.2124(2R/M+2r(x)+3)^3c_h$. This error acts multiplicatively on $(e^\Lambda\Xi_w)\circ\tau_t>0$, whose integral is $Z$; and $|1-e^y|\le e|y|$.

\emph{Step 6: the second term.} Here we apply \Cref{Y-P6-L9}, with $\sup|\Psi|\le\sup|\Phi|e^{t\sup|D\Lambda^a|}\le e^{1/2}\sup|\Phi|$ by the data.

\emph{Step 7: the Cauchy estimates for $E^a$. Step 7.0: the frames do not move.} We have $D_w^\alpha X_b=0$ exactly, and so $\tau_sV_b=e^{isw_bX_b(V)}V_b$. Let $g_z:=V(\Gamma(0,z))$. We claim that $D_wg_z=i\,\check\sigma(z)e_\alpha g_z$, where $\check\sigma(z)$ is the signed sum of the $w_b$ along $\Gamma(0,z)$. This follows by induction along the comb path, whose initial segments are the comb paths of the sites it passes through ($e_1$, then $e_2$, then $e_3$: item~2 of \S\ref{Y-J-1}, on the energy). At a forward step $b=(z,z+e_\mu)$, with $\pi(b)=z$, we have $DV_b=iw_bX_bV_b$ and $X_b=g_z^{-1}e_\alpha g_z$, so that $g_zDV_b=iw_be_\alpha g_zV_b$. At a backward step $b=(z-e_\mu,z)$, with $\pi(b)=z-e_\mu$, we set $g:=g_zV_b^{-1}$; then $X_b=g^{-1}e_\alpha g$ and $D(V_b^{-1})=-iw_bV_b^{-1}X_b$, so that $g_zD(V_b^{-1})=-iw_be_\alpha g$. Hence $DX_b=g^{-1}[e_\alpha,i\,\check\sigma\,e_\alpha]g=0$: the field is constant along its own flow.

\emph{Step 7.1: the curve and its domain.} For complex $s$ let $V_b(s)=e^{isw_bX_b}V_b$. By Step~7.0 this curve is the flow itself for real $s$, so that $d^n/ds^n$ at $0$ is $D_w^n$, for every $n$. Carrying all four factors of a plaquette to its base site, we find $\partial V(s)(p)=\prod_be^{\pm isw_bY_b}\,\partial V(p)$, where $Y_b=T_{\pi p}^{-1}\operatorname{Ad}_{\mathrm{hol}_b}e_\alpha$ does not depend on $s$ and, with $Y:=T_{\pi p}^{-1}e_\alpha$, satisfies $|Y_b-Y|\le5(r_p+1)\varepsilon$ (Step~1 of the proof of \Cref{Y-P6-L12}). By telescoping, $|\prod_be^{\pm isw_bY_b}-1|\le|s|(|\mathrm{d}w|_p+5(r_p+1)\varepsilon|w|_p)e^{|s||w|_p}$. We put $\delta_p:=|\mathrm{d}w|_p+5(r_p+1)\varepsilon|w|_p$, we let $\delta_X$ be the supremum of $\delta_p$ over the plaquettes of $X^\approx$, and $w_X$ the supremum of $|w_b|$ over $X^\approx$. We define $s_X$ by $1/s_X:=4\delta_X/\varepsilon+(2\ell_w(x)+4/\vartheta_0^\flat)w_X$. For $|s|\le s_X$ the following hold.

\begin{enumerate}\item[\textup{(u3)}] $|\partial V(s)(p)-1|\le3\varepsilon/4+4t+e^{1/2}\varepsilon/4<2\varepsilon$.

\item[\textup{(u2)}] $|V(s)({\rm path})|\le e^{|s|\ell_ww_X}\le e^{1/2}<2-1/\log x$.

\item[\textup{(u1)}] $|V_b(s)|\le e^{\vartheta_0^\flat/4}\le1+\vartheta_0^\flat/2<1+\vartheta_0^\flat(1-1/\log x)$.\end{enumerate}

Here \textup{(u1)}, \textup{(u2)} and \textup{(u3)} are the three conditions of the domain $\mathcal U_X(x)$ (item~1 of \S\ref{Y-G0-1}, on the domains). In \textup{(u3)} the last inequality follows from the bound $4t<\varepsilon/4$ of Step~2, and in \textup{(u1)} and \textup{(u2)} it rests on $\log x\ge3$, which is \textup{(j7)}\,\textup{(vi)} of the list of \Cref{Y-P6-N}.

\emph{Step 7.2: the first order.} The Cauchy estimate gives $|D_wE^a(X,V)|\le c_ae^{-\kappa_1(\#X-1)}/s_X$, and $1/s_X\le\sum_{{\rm cubes\ in\ }X}1/s_{\rm cube}$. We have $\sum_{\rm cubes}\sup_{p\in{\rm cube}^\approx}\delta_p\le(2r+1)^3\sum_p\delta_p$, where $\sum_p$ is the sum over all plaquettes. Likewise, $\sum_{\rm cubes}\sup_{b\in{\rm cube}^\approx}|w_b|\le(2r+1)^3\sum_b|w_b|$, with $\sum_b$ running over all bonds. We bound the sums with $|w|\le\min\{1,r_1^3r^{-3}\}$ and $|\mathrm{d}w_R|\le3r_1^4r^{-4}$ (\Cref{Y-P6-L11}\,\textup{(a0)}, \textup{(a)}), and with the count of $3(24r^2+2)\le78r^2$ bonds, and of as many plaquettes, at distance $r\ge1$ (\Cref{Y-A-Sums}). This gives $W_0\le3+78\sum_{1\le r\le r_1}r^2+78r_1^3\log(R/r_1)\le W_0(x)$ and $W_1\le234r_1^4\pi^2/6+($terms at $0)\le400r_1^4$. For $W_2$, a plaquette based at radius $r$ has its bond bases at
radius at least $r-1$, and it cannot meet a moved bond if $r>R+1$.
For $r\le r_1+1$ use $|w|_p\le4$. The sum of
$4(72r^2+6)(r+1)$ over those radii, including the origin, is at most
$73r_1^4$. For $r\ge r_1+2$ use $|w|_p\le4r_1^3/(r-1)^3$.
The corresponding summand is at most $288.5r_1^3$, since
$4(72r^2+6)(r+1)/(r-1)^3\le288.5$ for $r\ge r_1+2$.
There are at most $R-r_1$ such radii. Hence $W_2\le289r_1^3R$. As $r_B\ge R+h/8\ge3R$, \textup{(j5)} gives $\varepsilon\le1/(288R^2)$; hence $5\varepsilon W_2\le5.02r_1^3/R\le1$, by $R\ge6r_1^3$, one of the entries of the list of \Cref{Y-P6-N} that hold of themselves. Also $W_{00}\le3+\sum78r^2\min\{1,r_1^6r^{-6}\}\le60r_1^3$. So \[ \sum_{\rm cubes}1/s_{\rm cube}\le(2r+1)^3[4(W_1+5\varepsilon W_2)+\Theta(x)]/\varepsilon\le(2r+1)^3\cdot1605r_1^4/\varepsilon, \] by $\Theta\le1$. Since $c_a/\varepsilon=E_ap_0^2$ and $1.2124\cdot1605r_1^4\le2.5W_*$, this yields $C_a(x)$.

\emph{Step 7.3: the second order.} The Cauchy estimate gives $|D_w^2E^a(X,V)|\le2c_ae^{-\kappa_1(\#X-1)}/s_X^2$. As $\delta_X$ and $w_X$ are maxima over the cubes of $X$, \[ 1/s_X^2\le\sum_{{\rm cubes\ in\ }X}[32\delta_{\rm cube}^2/\varepsilon^2+2(2\ell_w+4/\vartheta_0^\flat)^2w_{\rm cube}^2]. \] We sum over the cubes in the manner of Step~7.2, that is, with the factor $(2r+1)^3$ times the sum over all plaquettes or over all bonds; the result is at most \[ 2\cdot1.2124c_a(2r+1)^3\varepsilon^{-2}[32\sum_p\delta_p^2+2(\varepsilon(2\ell_w+4/\vartheta_0^\flat))^2W_{00}]. \] Here $\sum\delta_p^2\le(W_1+5\varepsilon W_2)^2\le(401r_1^4)^2\le0.161W_*^2$, and $\varepsilon(2\ell_w+4/\vartheta_0^\flat)\le1$, by $\Theta\le1$ and $W_0(x)\ge1$. With $c_a\varepsilon^{-2}=E_ap_0x^{1/2}$ the total is $2.425(5.146W_*^2+120r_1^3)\le12.5W_*^2$ (the coefficient is $12.48$).

\end{proof}

\begin{lemma}[The Green function, uniformly in the rows]\label{Y-P6-L11} Let $\ell=\ell_k^{[K]}$, $1\le k\le K$, be unfolded to $\mathbb Z^3$. We put $B=1+\ell^{\mathsf T}\ell$, $M=\mathrm{d}^*B\mathrm{d}$, $C_B=\mathrm{d}M^+\mathrm{d}^*$, $\omega:=C_Bu$, $w:=\mathrm{d}^*(-\Delta)^{-1}\omega$, $w_R:=\chi_Rw$ and $e:=\mathrm{d}^*u-Mw_R$. Then the following hold.

\begin{enumerate}\item[\textup{(a0)}] At every $y$ we have $|w|\le1$, $|\omega|\le\sqrt2$, $\mathrm{d}w=\omega$ and $Mw=\mathrm{d}^*u$.

\item[\textup{(a)}] For $r\ge r_1=10^{130}$: $|w|\le C_wr^{-3}$, $|\nabla w|\le C_{gw}r^{-4}$, $|\omega|\le C_\omega r^{-4}$ and $|\nabla\omega|\le C_{g\omega}r^{-5}$. For all $r\ge1$: $|w|\le r_1^3r^{-3}$ and $|\mathrm{d}w_R|\le3r_1^4r^{-4}$.

\item[\textup{(b)}] For $r\ge r_1$: $|\omega-\Pi u|\le C_Gr^{-6}+C_\tau\tau_x^2r^{-4}$; and, for the gradient, $|\nabla(\omega-\Pi u)|\le C_G'r^{-7}+1.9\cdot10^{36}\tau_x^2r^{-5}+A_{\rm out}r^{-16}$.

\item[\textup{(c)}] $0\le a_u=\langle u,\omega\rangle\le2$.

\item[\textup{(d)}] If $8\rho+8\le R$ and $R\ge4r_1$, then $e$ is carried by $R/2-2\rho-2\le r\le R+2\rho+2$, $|e|\le C_eR^{-5}$, and $\|e\|_1\le81C_eR^{-2}$.\end{enumerate}

None of the constants depends on $\rho(x)$, $k$, $K$ or $N$; the values of $C_w$, $C_{gw}$, $C_\omega$, $C_{g\omega}$, $C_G$, $C_G'$, $C_\tau$, $A_{\rm out}$ and $C_e$ are given in item~17 of \S\ref{Y-P6-1}.\end{lemma}

\begin{proof} \emph{Step 0: unfolding.} The rows have range at most $\rho(x)$, and the side is $S=NL^j>4\rho+2$ (by \textup{(Tor)} in \textup{(SD)}, and \Cref{Y-H4-U4}\,\textup{(b)}). Hence, for $|y-p|\le\rho$, the formula $\ell_y(p):=\ell_{[y]}([p])$ gives covariant rows on $\mathbb Z^3$; their forms $\mathsf E_y$ are unchanged, and so are the bound $K_*e^{-\lambda_*|y-p|}$ (Step~2 of the proof of \Cref{Y-H9-T}, in the part headed ``The last integral'') and the defect $\tau$. Identify the target spaces $\mathsf H_y$ orthogonally with
$\mathsf H_0$ using translation covariance. The unfolded row has
finite support, so its symbol $\hat\ell(q)$ is a matrix of finite
Laurent polynomials and
$\hat B(q)=1+\hat\ell(-q)^{\mathsf T}\hat\ell(q)$.
For real $p$, $\hat B(p)=1+\hat\ell(p)^*\hat\ell(p)\ge1$.
The inequality $S>4\rho+2$ ensures that periodizing the unfolded
kernel, whose range is at most $2\rho$, recovers the torus kernel.
At $p\in(2\pi/S)\mathbb Z^3$ the symbol $\hat B(p)$ is therefore
that of the torus, and it is at most $b_*$. At any other $p$ we choose $p'\in(2\pi/S)\mathbb Z^3$ with $|p-p'|_\infty\le\pi/S$; then \[ \|\hat B(p)-\hat B(p')\|\le9K_*^2\mathsf g(\lambda_*/2)\mathsf g_1(\lambda_*/2)|p-p'|_\infty\le4.81\cdot10^{127}|p-p'|_\infty<1, \] because $S\ge2L\ge2L_A\ge10^{149}$ (by \Cref{Y-A-PureNumbers}, with $\mathsf g_1(\lambda)\ge2.2\lambda^{-4}$). So $1\le\hat B\le b^\circ:=b_*+1$ at all $p$. In the application to \Cref{Y-P6-L12}, $j\in J(\mathsf b)$.
Condition \textup{(j3)} gives $S\ge2Lh$, while $h\ge16R$ and
\textup{(j7)(iii)} gives $8\rho+8\le R$.
Thus $B'$ and every range-$2\rho$ interaction entering $Mw_R$
fit inside one fundamental domain. The torus and $\mathbb Z^3$
therefore give the same $Mw_R$ there; no torus Green function
is identified with the infinite-lattice Green function.

We now prove the bound $\tau(\ell_k)\le\tau_x$. \Cref{Y-A-Th29}\,\textup{(a)} says that $\ell\in \mathcal R(b_*,K_*,\lambda_*,\tau)$, $\tau\le1/9$ and $\mathsf t_*(\rho)\le1/9$ imply $\ell'^{\rm c}\in \mathcal R(b_*,K_*,\lambda_*,\tau+\mathsf t_*(\rho))$. Along the runs (item~2 of \S\ref{Y-J-1}; the rows of item~1 of \S\ref{Y-H2-1}) we have $\ell_i=(\ell_{i-1}')^{\rm c}$, with the cut at $\rho(x_i)$, and $\tau(\emptyset)=0$. All $\ell_i$ lie in $\mathcal R(b_*,K_*,\lambda_*,1/9)$, by Step~2 of the proof of \Cref{Y-H9-T} (for $0\le k\le K$ and every $K\ge1$). Therefore $\tau(\ell_k)\le\sum_{i\le k}\mathsf t_*(\rho(x_i))$. The definition of $\mathsf t_*$ (\Cref{Y-A-PureNumbers}), together with $\rho(x)\ge c_R\log x$ and $e_R=\lambda_*c_R/2\ge256$ (item~1 of \S\ref{Y-H2-1}), gives $\mathsf t_*(\rho(x))\le3K_*\mathsf g(\lambda_*/2)x^{-e_R}$. As $x_i=L^{k-i}x$, the sum is a geometric series and is at most $3K_*\mathsf g(\lambda_*/2)x^{-e_R}/(1-L^{-e_R})\le\tau_x$.

\emph{Step 1: symbols.} In Steps 1--5, Fourier-vector norms are Euclidean; the spatial radius remains $r=|y|_\infty$. We write $k_\mu=e^{ip_\mu}-1$ and $k^\#_\mu=e^{-ip_\mu}-1$; then $(2/\pi)|p|\le|k|\le|p|$. The range of $\hat {\mathrm{d}}$ is $\{F:k^{\mathsf T}F=0\}$. With $\mathsf R=B^{-1}$ and $s=k^{\mathsf T}\mathsf Rk^\#$ we have $\hat C_B=\mathsf R-\mathsf Rk^\#k^{\mathsf T}\mathsf R/s$: both sides annihilate $k^\#$, map into $\{k^{\mathsf T}F=0\}$, and invert $\hat B$ on that space modulo $k^\#$. Hence $\hat B\hat C_B=1-k^\#k^{\mathsf T}\mathsf R/s$, and $\widehat{\mathrm{d}^*}k^\#=0$; so $\mathrm{d}^*BC_Bu=\mathrm{d}^*u$, that is, $Mw=\mathrm{d}^*u$. On 2-forms $\mathrm{d}\mathrm{d}^*+\mathrm{d}_2^*\mathrm{d}_2=-\Delta$, and $k^{\mathsf T}\hat\omega=0$; so $\mathrm{d}w=\omega$. Further, $\hat u=e_1(1-e^{-ip_2})$. Let $P:=\hat\Pi=1-k^\#k^{\mathsf T}/(k^{\mathsf T}k^\#)$ be the symbol of the projection $\Pi$ of \Cref{Y-P6-N}, continued to the discs. We have $\hat C_B\hat BP=P$, hence $\hat C_B-P=-\hat C_B(\hat B-1)P$. The symbol $(\hat B-1)(q)$ is holomorphic for $|\operatorname{Im}q_\mu|\le\lambda_*/6$. Finally, $|\hat\ell(q)e_i-\hat\ell(q')e_i|\le(\mu_1/3)|q-q'|_\infty$, where $\mu_1:=9K_*\mathsf g_1(\lambda_*/2)\le3.78\cdot10^{71}$, and $|\hat\ell(0)e_i|\le\tau_x$.

\emph{Step 2: inner discs.} Let $|p|\le p_1:=1/(20\mu_1)$, and let $q$ differ from $p$ in one coordinate only, with $|q_\mu-p_\mu|\le|p|/20$. Then $|\sum q^2|\ge(1-0.1-0.0025)|p|^2=0.8975|p|^2$ and $\sum|q|^2\le1.05^2|p|^2=1.1025|p|^2$, and $\|\hat B(q)-1\|\le3(\tau_x+1.05\mu_1|p|)^2\le0.0086$. Hence $|s|\ge0.80|p|^2$, $\|\hat C_B\|\le2.5107$, $\|P\|\le2.442$ and $|\hat u|\le1.051|p|$. The gradients below are full tensors, with their Euclidean norms: their symbols are $k\otimes\hat w$, $k\otimes\hat\omega$ and $k\otimes((C_B-\Pi)u)^\wedge$. The pure numbers give $p_1<10^{-12}$; thus $|q|\le1.05|p|$ and the exponential series imply $|k(q)|,|k^\#(q)|\le1.0501|p|$. Six symbols get majorants of the form $A'|p|^a$: for $\hat w$ it is $3.502$, with $a=0$; for $k\otimes\hat w$, $3.68|p|$; for $\hat\omega$, $2.64|p|$; for $k\otimes\hat\omega$, $2.78|p|^2$; for $((C_B-\Pi)u)^\wedge$, $42.7\mu_1^2|p|^3+38.7\tau_x^2|p|$; and for its full gradient, $1.051|p|$ times that. Indeed $3.502\cdot1.0501<3.68$ and $2.64\cdot1.0501<2.78$. The Cauchy estimate for these vector- or tensor-valued symbols gives $|\partial_\mu^im|\le i!(20/|p|)^iA'|p|^a$. The symbol $P$ is discontinuous at $0$, and we take no derivative of it there.

\emph{Step 3: outer discs.} Now let $|p|>p_1$, and take discs of radius $d_{\rm out}:=\min\{p_1/(200b^\circ),\allowbreak 1/(80b^\circ\check\Lambda)\}\ge4.2\cdot10^{-90}$, where $\check\Lambda=3\mu_1b^{\circ1/2}$. Write $b=b^\circ$, $d=d_{\rm out}$ and $K=|k(p)|$. The row-symbol bound of Step~1 gives $\|\hat\ell(q)-\hat\ell(p)\|\le\mu_1d/\sqrt3$ and the same bound at $-q,-p$. Since $\|\hat\ell(p)\|\le\sqrt b$, it follows that
\[
 \|\hat B(q)-\hat B(p)\|
 \le2\sqrt{b/3}\,\mu_1d+\mu_1^2d^2/3
 \le\check\Lambda d\le\frac1{80b}.
\]
Thus $\|\mathsf R(q)\|\le80/79<1.013$ and $\|\mathsf R(q)-\mathsf R(p)\|\le1/(79b)<0.013/b$. Also $K\ge2|p|/\pi$ and $|k(q)-k(p)|,|k^\#(q)-k^\#(p)|\le2d\le0.016K/b$. The real symbol has $s(p)\ge K^2/b$. Expanding its three factors gives
\[
 |s(q)-s(p)|
 \le\frac{K^2}{b}
 [0.016\cdot1.013\cdot1.016+0.013\cdot1.016+0.016]
 <0.046K^2/b,
\]
so $|s(q)|>0.95K^2/b$. Likewise $|k(q)\cdot k^\#(q)|\ge(1-2\cdot0.016-0.016^2)K^2>0.967K^2$. The formulas of Step~1 now give $\|\hat C_B(q)\|<3b$ and $\|P(q)\|<3$. Moreover $|k(q)|,|k^\#(q)|<3.5$, $|\hat u(q)|<2.1$ and
\[
 |\hat w(q)|\le
 \frac{|k^\#(q)|^2}{|k(q)\cdot k^\#(q)|}\,\|\hat C_B(q)\|
 <4b.
\]
The six full symbols of Step~2 are therefore bounded respectively by $4b$, $14b$, $6.3b$, $22.05b$, $12.6b$ and $44.1b$, all less than $70b$. Consequently $|\partial_\mu^{16}m|\le16!\,\allowbreak d_{\rm out}^{-16}70b^\circ=:A_{\rm out}\le10^{1455.5}$.

\emph{Step 4: a bound for multipliers.} We first state it. Let $m$ be periodic and $C^n$ off $0$, with $a\ge0$ and $n\ge a+4$; let $|\partial_\mu^im|\le A|p|^{a-i}$ for $|p|\le p_1$ and $i\le n$, and let $|\partial_\mu^nm|\le A_{\rm out}$ for $|p|>p_1$. Then, for $r=|y|_\infty=|y_\mu|\ge\sqrt3/p_1$, \[ |m^\vee(y)|\le c_{a,n}Ar^{-3-a}+A_{\rm out}r^{-n}, \] with $c_{a,n}=(8\cdot3^{a/2}(n+1)+4\pi/(n-a-3))/(2\pi)^3$.

We now prove this bound. On $Q=\{|p|_\infty\le1/r\}$ we have $\int|m|\le8r^{-3}A(\sqrt3/r)^a$. Off $Q$ we integrate by parts $n$ times in $p_\mu$, and the boundary terms at the ends $p_\mu=\pm\pi$ cancel by periodicity. Only the lines with $|p_\perp|_\infty\le1/r$ have further ends, namely at $\pm1/r$, and the boundary terms there are at most $\sum_{i<n}r^{-i-1}\cdot2\cdot(4/r^2)A\cdot3^{a/2}r^{i-a}$. What remains is at most $r^{-n}[A\cdot4\pi\int_{1/r}^\infty s^{a-n+2}ds+(2\pi)^3A_{\rm out}]$. This proof is linear in the majorant.

\emph{Step 5: clauses \textup{(a)} and \textup{(b)}.}
Apply Step~4 with $n=16$ and $A=16!\,20^{16}A'$ separately to
the powers in each majorant of Step~2. The outer remainder is
$A_{\rm out}r^{-16}$, with $A_{\rm out}\le10^{1455.5}$.
At $r_1=10^{130}$ its ratio to $c_{a,16}Ar^{-3-a}$ is at most
$10^{-9.29}$ for the four symbols
$\hat w,k\otimes\hat w,\hat\omega,k\otimes\hat\omega$ and
for the term $42.7\mu_1^2|p|^3$ of $((C_B-\Pi)u)^\wedge$.
These ratios decrease for $r\ge r_1$. The numerical constants
$C_w,C_{gw},C_\omega,C_{g\omega},C_G$ in item~17 exceed
$(1+10^{-9.29})c_{a,16}A$ for the corresponding majorants.
This proves (a) and absorbs the outer remainder in $C_Gr^{-6}$
in the first estimate of (b).

The terms containing $\tau_x^2$ are handled by linearity of the
integral estimate in Step~4. Their coefficients are bounded by
$C_\tau$ and $1.9\cdot10^{36}$. For the gradient of the
$\mu_1^2|p|^3$ term, use $a=4$ and the constant $C_G'$.
At $r=r_1$ the outer remainder cannot be absorbed into this
$r^{-7}$ bound, so it remains as the explicit
$A_{\rm out}r^{-16}$ term in (b).

For \textup{(a0)}: $|\hat w|\le|k||\hat\omega|/|k|^2\le|k_2|/|k|\le1$, and $|\omega(y)|\le\|C_Bu\|_2\le\sqrt2$. For $r<r_1$ we have $1\le r_1^3r^{-3}$. Finally, $\mathrm{d}w_R=\chi_R\omega+\eta_1$. Summing the squares of the three antisymmetric components in the formula for $\eta_1$ below gives
\[
 |\eta_1(y)|^2\le\frac{54}{R^2}
                  \max_\alpha|w(y+e_\alpha)|^2,
 \qquad |\mathrm{d}w_R(y)|\le\sqrt2+8/R.
\]
Here and in this display the norms include all form components. The prescribed $R$ satisfies $R\ge4r_1$. Thus $\mathrm{d}w_R=\omega$ for $r<r_1$. If $\eta_1(y)\ne0$, then $R/2-1\le r\le R$, and the bound on $w$ gives
\[
 |\eta_1(y)|\le\frac{8r_1^3}{R(r-1)^3}
                  \le9r_1^3r^{-4}.
\]
For $r\ge r_1$, therefore, $|\mathrm{d}w_R|\le(C_\omega+9r_1^3)r^{-4}<3r_1^4r^{-4}$; for smaller $r$ the same conclusion follows from $|\omega|\le\sqrt2$.

\emph{Step 6: clause \textup{(c)}.} The operator $B^{1/2}C_BB^{1/2}$ is an orthogonal projection. So $0\le C_B\le B^{-1}\le1$; and $|u|^2=2$.

\emph{Step 7: clause \textup{(d)}, an exact identity and the bounds.} At a cell, $\chi$ is evaluated at the base site of the cell. Then $\mathrm{d}(\chi w)=\chi\omega+\eta_1$, where \[ \eta_1(y;\mu\nu)=(\partial_\mu\chi)(y)w_\nu(y+e_\mu)-(\partial_\nu\chi)(y)w_\mu(y+e_\nu); \] next, $B(\chi\omega)=\chi B\omega+\eta_3$, where \[ \eta_3(p)=\sum_{p'}(B-1)(p,p')(\chi(p')-\chi(p))\omega_{p'}; \] and $\mathrm{d}^*(\chi G)=\chi \mathrm{d}^*G+\eta_2$. We take $G=B\omega$, for which $\mathrm{d}^*G=\mathrm{d}^*u$; since $\chi=1$ on ${\rm supp}\,\mathrm{d}^*u$, we obtain the exact identity $e=-\eta_2-\mathrm{d}^*\eta_3-\mathrm{d}^*B\eta_1$.

We now turn to the bounds of clause \textup{(d)}. For the product cut-off, $|\partial_\alpha\chi|\le3/R$ in each direction, $|\chi(y')-\chi(y)|\le(9/R)|y-y'|_\infty$, and the second differences are at most $24/R^2$. We write $S_\cdot$ for a supremum over $r\ge R/2-2\rho-2\ge R/4\ge r_1$. Then $|\eta_2|\le(6/R)N_{B0}S_\omega$ and $|\mathrm{d}^*\eta_3|\le4(9/R)N_{B1}S_\omega$. The form $\eta_1$ is only of order $R^{-1}R^{-3}$. However, $\mathrm{d}^*B$ annihilates constant forms for any covariant rows, because $Bc$ is then constant. Let $D_1(\eta_1)$ be the supremum of $|\partial_\alpha\eta_1(p)|$ over directions and plaquettes of all orientations in this region. Subtract the constant form $c_i=\eta_1(y_b;i)$. Each component difference is a sum along a coordinate path of length at most $3(|z|_\infty+1)$; the four terms in $\mathrm{d}^*$ therefore give $|\mathrm{d}^*B\eta_1(b)|\le12(N_{B0}+N_{B1})D_1(\eta_1)$. The product rule in each direction gives $D_1(\eta_1)\le2[(24/R^2)S_w+(3/R)S_{gw}]$. Here $N_{B0}=1+3K_*^2\mathsf g(\lambda_*/2)^2\le4.9\cdot10^{112}$ and $N_{B1}=3K_*^2\mathsf g(\lambda_*/2)\mathsf g_1(\lambda_*/2)\le1.61\cdot10^{127}$; moreover $S_w\le64C_wR^{-3}$, $S_\omega\le256C_\omega R^{-4}$ and $S_{gw}\le256C_{gw}R^{-4}$. Altogether, this gives \[ C_e=256(6N_{B0}+36N_{B1})C_\omega+24(N_{B0}+N_{B1})(1536C_w+768C_{gw})\le7.1\cdot10^{166}. \] The number of bonds that carry $e$ is at most $3(2R+4\rho+5)^3\le81R^3$.\end{proof}

\begin{lemma}[The key line with its non-abelian remainders]\label{Y-P6-L12} In the setting of \Cref{Y-P6-N}, at $j\in J(\mathsf b)$, on fields with $\phi_p<\varepsilon$ for every plaquette of $B'$ and for $\alpha=1,2,3$, \[ |x\,u^\alpha(F)-x\,D^\alpha_{w_R}\mathrm{En}_\ell(V)|\le C_1x\varepsilon/R+C_2p_0^2R, \] where $C_1=203C_e$. In its applications the fields satisfy the stronger bound $\phi_p\le3\varepsilon/4+4t<\varepsilon$ on $B'$.\end{lemma}

\begin{proof} \emph{Step 0: the gauge.} Under a gauge transformation $g$ with $g(0)=1$, the quantities $\tilde F$ and $\mathrm{En}_\ell$ are invariant, and the vector field is covariant ($X_b\to\operatorname{Ad}_{g(\pi b)}X_b$). We work in the comb gauge on $B'$, as in Step~1 of the proof of \Cref{Y-P6-L4}. There, at the point $V$, $T_y=1$ and $X_b=e_\alpha$; also $\tilde F=F$, and $V=e^{iA}$, where $|A_b|\le2r\varepsilon$ by the count of plaquettes in that step: for a bond at the site $y$, it gives at most $|y_2|+|y_3|\le2r$ plaquettes.

\emph{Step 1: the variation of a plaquette, in any gauge.} The product $\delta(\partial V(p))\,\partial V(p)^{-1}$ is $i$ times the sum, over the four bonds, of $\pm w_b\operatorname{Ad}_{\rm path}X_b$. Each term equals $T_{\pi p}^{-1}\operatorname{Ad}_{\mathrm{hol}_b}e_\alpha$, where $\mathrm{hol}_b$ is the holonomy of a loop composed of the comb paths to two neighbouring sites and closed along $p$; its area is at most $2r+2$. Since $\operatorname{Ad}_h$ is a rotation by $2\,\operatorname{ang}\,h$, we get $|\operatorname{Ad}_{\mathrm{hol}_b}e_\alpha-e_\alpha|\le2(2r+2)\varepsilon$. For the left-trivialized variation $\delta g\,g^{-1}=iY$, write $g=\cos\phi\,1+i\sin\phi\,n\cdot\boldsymbol\sigma$ in an orthonormal Pauli basis, with $\phi=\operatorname{ang}g$ and $|n|=1$, and decompose $Y=Y_\parallel+Y_\perp$ along $n$. Differentiating this formula gives
\[
 d\mathsf L_g(Y)=Y_\parallel+\phi\cot\phi\,Y_\perp
                         +\phi\,n\times Y_\perp.
\]
For $0<\phi\le0.01$, the inequalities $\tan\phi\ge\phi$, $\sin\phi\le\phi$ and $\cos\phi\ge1-\phi^2/2$ give $0\le1-\phi\cot\phi\le\phi^2/2$. Orthogonality of $Y_\perp$ and $n\times Y_\perp$ yields $|d\mathsf L_g(Y)-Y|\le\phi\sqrt{1+\phi^2/4}|Y|\le1.01\phi|Y|$; the value at $\phi=0$ follows by continuity. Here $\phi\le\varepsilon$, as the plaquette $p$ is tame, and $\varepsilon\le2r_B\varepsilon\le0.01$ by Step~1 of the proof of \Cref{Y-P6-L4}. Since $|\mathrm{d}w|_p\le|w|_p$, the excess over $\phi|Y|$ fits into the room between $4(r+1)$ and $5(r+1)$. Hence $\delta F_p=T_{\pi p}^{-1}[(\mathrm{d}w)_pe_\alpha]+\mathrm{rem}_p$, with $|\mathrm{rem}_p|\le\varepsilon(5(r+1)|w|_p+|\mathrm{d}w|_p)$.

\emph{Step 2: the plaquette part.} We have \[ x\sum_pF_p\cdot\delta F_p=x\langle\tilde F^\alpha,\mathrm{d}w_R\rangle+O(x\varepsilon^2\sum_p(5(r+1)|w|_p+|\mathrm{d}w|_p)). \] With $|w|\le r_1^3r^{-3}$ and $24r^2+2$ sites at distance $r$, the sum is at most $1600r_1^3R$. So the error is $1600r_1^3p_0^2R$.

\emph{Step 3: the row part.} Write $\mathrm{En}_{\rm row}=\frac12\sum_y m_y(V)|\ell_y^V\zeta|^2$. Every tame indicator is locally constant along the flow: each plaquette that can move lies in $B'$ and has angle strictly below $\varepsilon$. Retain the rows whose curvature arguments or transport paths meet a bond moved by $w_R$. Their centers have radius at most $R+\rho+1$, and their plaquettes, transport paths and masks lie in $B'$, using $\rho_0\le\rho+1$. On these rows the tame indicators equal $1$ and the clipped curvature equals $F$. The other rows contribute neither to the derivative nor to its proposed leading term. We may therefore use $|F_p|\le\varepsilon$ on every retained row and extend all ensuing positive bounds to the full lattice.

Put $r_p=|\pi(p)|_\infty$, $d=|y-\pi(p)|_\infty$, $d'=|y-\pi(p')|_\infty$ and
\[
 A_{y,p}=\operatorname{Ad}_{V(\Gamma(y,\pi(p)))},\qquad
 a_{y,p}=\sum_{b\in\Gamma(y,\pi(p))}|w_b|.
\]
In the gauge of Step~0, differentiation and symmetry give the row derivative, before multiplication by $x$, as
\[
 \sum_{y,p,p'}\langle\ell_y(p),\ell_y(p')\rangle
       \langle D(A_{y,p}F_p),A_{y,p'}F_{p'}\rangle.
\]
The leading term is $\langle F^\alpha,(B-1)\mathrm{d}w_R\rangle$. We bound its three errors separately.

The curvature remainder of Step~1 contributes at most
$\varepsilon^2N_{B0}(5W_2+W_1)$, with the sums from \Cref{Y-P6-L10}. Differentiating the transport gives $\|DA_{y,p}\|\le2a_{y,p}$. For a fixed displacement $v=y-\pi(p)$ and fixed plaquette orientation $\iota$, summing over translates of its coordinate path gives
\[
 \sum_{z\in\mathbb Z^3}a_{z+v,(z;\iota)}\le3|v|_\infty W_0.
\]
Indeed each of its at most $3d$ bond positions contributes at most $W_0$. Summing both plaquette orientations and using the exponential row bound therefore makes this error at most
\[
 54\varepsilon^2K_*^2\mathsf g(\lambda_*)
                   \mathsf g_1(\lambda_*)W_0
       \le18\varepsilon^2N_{B1}W_0.
\]

Finally, each bond on $\Gamma(y,\pi(p))$ has radius at most $|y|_\infty+d$ and angle at most $2(|y|_\infty+d)\varepsilon$. The adjoint bound and telescoping along the path yield
\[
 \|A_{y,p}-1\|\le12d(|y|_\infty+d)\varepsilon.
\]
Since $|y|_\infty\le r_p+d$, the discrepancy between the two transports in the leading term is bounded by
\[
 \|A_{y,p'}^{-1}A_{y,p}-1\|
 \le12\varepsilon\bigl[(d+d')(r_p+1)+3d^2+3d'^2\bigr].
\]
For each fixed $p$, the row weights summed over $y,p'$ have zeroth, first and second displacement moments bounded by $N_{B0}$, $N_{B1}$ and $N_{B2}$, respectively, for either $d$ or $d'$. The first two constants are those of Step~7 of the proof of \Cref{Y-P6-L11}; the last is
\[
 N_{B2}=3K_*^2\mathsf g(\lambda_*/2)
                  \,600(\lambda_*/2)^{-5}\le4.9\cdot10^{141}.
\]
These bounds follow by summing $K_*^2e^{-\lambda_*(d+d')}$ over the three orientations of $p'$ and then using the shell sums of \Cref{Y-A-Sums}. Thus the transport discrepancy contributes at most
\[
 \varepsilon^2(24N_{B1}V_1+72N_{B2}W_1),\qquad
 V_1:=\sum_p(r_p+1)|(\mathrm{d}w_R)_p|.
\]

To bound the new sum, use $|\mathrm{d}w_R|\le\min\{2,3r_1^4r^{-4}\}$ and $72r^2+6$ plaquettes at radius $r\ge1$. The radii $r\le r_1$ contribute less than $40r_1^4$; at $r>r_1\ge1000$ each shell contributes at most $217r_1^4/r$. Since $\mathrm{d}w_R$ is supported at radii at most $R+1$,
\[
 V_1\le500r_1^4(1+\log(R+1))\le r_1^4R.
\]
The last inequality holds for $R\ge10^6$, as is the case here. Likewise the existing bounds give $W_0,W_1,W_2\le r_1^4R$, since $r_1\ge289$ and $R\ge6r_1^3$. Combining the three errors gives
\begin{align*}
 &\left|xD\mathrm{En}_{\rm row}
              -x\langle F^\alpha,(B-1)\mathrm{d}w_R\rangle\right|\\
 &\quad\le p_0^2\bigl[N_{B0}(5W_2+W_1)+18N_{B1}W_0
                        +24N_{B1}V_1+72N_{B2}W_1\bigr]\\
 &\quad\le200(N_{B0}+N_{B1}+N_{B2})r_1^4p_0^2R.
\end{align*}

\emph{Step 4: the two parts together.} So $xD\mathrm{En}_\ell=x\langle\tilde F^\alpha,B\,\mathrm{d}w_R\rangle+O(p_0^2R)$, with the undressed scalar $B$.

\emph{Step 5: from $F$ to $A$.} We have $F_p=(\mathrm{d}A)_p+O(28r^2\varepsilon^2)$. Here $r$ is the distance of $p$. At second order, the four factors $e^{\pm iA_b}$ of $\partial V(p)$ give the term $\frac12\sum_{\rm pairs}[A_b,A_{b'}]$, a sum over six pairs. Of the bonds of $p$, the one at the neighbouring site has $|A_b|\le(2r+1)\varepsilon$, two have $|A_b|$ at most $r\varepsilon$ or $0$, and the fourth has $|A_b|\le2r\varepsilon$. With the bracket constant $2$, that is, with $|[Y,Y']|\le2|Y||Y'|$, the term of second order is therefore at most $\frac12\cdot2\sum_{\rm pairs}|A_b||A_{b'}|\le(13r^2+4r)\varepsilon^2\le24r^2\varepsilon^2$ for $r\ge1$. For the higher terms put $S=\sum_b|A_b|\le(6r+1)\varepsilon\le7r\varepsilon\le0.035$, since $2r\varepsilon\le2r_B\varepsilon\le0.01$. The logarithm coefficient bound proved in Step~3\,(i) of \Cref{Y-P6-L4}, with $H(s)=-\log(2-e^s)$, bounds the terms of degree at least three by $H(S)-S-S^2$. The series converges absolutely since $e^S-1<1$. Since
\[
 H'''(s)=\frac{2e^s(2+e^s)}{(2-e^s)^3},\qquad
 H'''(0.035)/6<1.18,
\]
Taylor's theorem gives $H(S)-S-S^2\le1.18S^3\le2.024r^2\varepsilon^2$, using $r\varepsilon\le0.005$. For real $A_b$, the logarithm after multiplication by $-i$, its homogeneous terms and their remainder are traceless Hermitian matrices; their operator norm is the norm $|\cdot|$ used here. The total error is consequently less than $(24+2.024)r^2\varepsilon^2<28r^2\varepsilon^2$. At $r=0$ three of the four $A_b$ vanish, so that $F_p=(\mathrm{d}A)_p$ exactly. Since $r_p^2\le2r_{p'}^2+2|p-p'|^2$, and with $|\mathrm{d}w_R|\le\min\{\sqrt2+8/R,\,3r_1^4r^{-4}\}$ for the factor $370$, \[ \sum_pr^2|B\,\mathrm{d}w_R|_p\le2N_{B0}\cdot225r_1^4R+2N_{B2}\cdot370r_1^4. \] The factor $370$ is unchanged: $R\ge4r_1$ gives $\sqrt2+8/R<1.415$, and the three plaquette types give
\[
 \sum_p|\mathrm{d}w_R|_p
 \le r_1^4\left[81\cdot1.415
       +216\sum_{m\ge2}m^{-2}+18\sum_{m\ge2}m^{-4}\right]
 <337r_1^4<370r_1^4,
\]
using the integral bounds $\sum_{m\ge2}m^{-2}\le1$ and $\sum_{m\ge2}m^{-4}\le1/3$. The error is $28$ times the preceding weighted sum, times $p_0^2$.

\emph{Step 6: summation by parts inside $B'$.} The form $B\,\mathrm{d}w_R$ lives in $r\le R+2\rho+1$, and \[ \langle \mathrm{d}A^\alpha,B\,\mathrm{d}w_R\rangle=\langle A^\alpha,Mw_R\rangle=\langle A^\alpha,\mathrm{d}^*u\rangle-e(A^\alpha)=u^\alpha(F)+O(100\varepsilon^2)-e(A^\alpha). \]

\emph{Step 7: the shell.} By \Cref{Y-P6-L11}\,\textup{(d)} and \textup{(j7)}\,\textup{(iii)}, \[ |e(A)|\le\|e\|_1\cdot2(R+2\rho+2)\varepsilon\le81C_eR^{-2}\cdot2.5R\varepsilon\le203C_e\varepsilon/R. \] The unknown infrared part of the field enters the proof only in this step; as $u$ annihilates constants, it is bounded by a constant times $\varepsilon/R$, where $\varepsilon$ is the bound of a tame field.

\emph{Step 8: the constant $C_2$.} The constant $C_2:=10^{670}$ covers Steps 2, 3, 5 and 6; the largest of their contributions is $28\cdot740N_{B2}r_1^4\le1.1\cdot10^{666}$. The contributions are collected in the rows for $C_2$ in item~17 of \S\ref{Y-P6-1}.\end{proof}

\begin{lemma}[Three bounds for the covariance and the variance; the list of errors]\label{Y-P6-L13} In the setting of \Cref{Y-P6-N}, put $q:=\langle \mathrm{d}w,\mathbb D_k[f_*]\mathrm{d}w\rangle=\sum_y\mathbb D_{k,y}[f_*](\omega,\omega)$. Then the following three bounds hold at $j\in J(\mathsf b)$.

\begin{enumerate} \item[\textup{(i)}] $|\mathrm{Cov}_\nu(G,S_f)-(3/2)q|\le$ \textup{(E2)}$+$\textup{(E3)}$+$\textup{(E4)}$+$\textup{(E6)}.

\item[\textup{(ii)}] $|\mathsf b\,\mathrm{Cov}(G_-,\mathcal O_K(f_*))-(3/2)q|\le$ \textup{(E1)}$+\ldots+$\textup{(E6)}.

\item[\textup{(iii)}] $|\mathrm{Var}(G_-)-(3/2)a_u^2|\le$ \textup{(E7)}. \end{enumerate}

Here the errors are the following quantities.

\begin{enumerate} \item[\textup{(E1)}] is the number \[(9/4)p_0^2\cdot1.2124N_1[2C_3K_5E_ap_0^3x^{-1/4}+0.12c_Q^{\rm su}K_5(Lx)^{1/2}c_h(x)/(E_ap_0(Lx))];\]

\item[\textup{(E2)}] is the number $5C_SC_1p_0^4/R$, and this number is at most $c_F/(12p_0)$;

\item[\textup{(E3)}] is the number $5C_S(C_2R+C_a(x))p_0^5x^{-1/2}$;

\item[\textup{(E4)}] is the number $10^5W_*C_Sp_0^4\nu_P$, where $\nu_P=3(2R+4\rho+7)^3e^{-(c_P/2)p_0^2}$;

\item[\textup{(E5)}] is the bound given by \Cref{Y-P6-L8};

\item[\textup{(E6)}] is the number \[6[C_6p_0^5x^{-3/2}+1.5eC_Sp_0^3x^{3/2}(2.5(2R/M+2r(x)+3)^3c_h(x)+6e^{2.2}x^{-4}+2\mathrm{ex})];\]

\item[\textup{(E7)}] is the number $4[$\textup{(E2)}$+$\textup{(E3)}$+$\textup{(E4)}$+$\textup{(E6)}$]$. \end{enumerate} \end{lemma}

\begin{proof} Throughout the proof we write $D=D^\alpha_{w_R}$.

\emph{Step 0: the tool.} On ${\rm supp}\,\chi$ we have $D\Lambda^a=-xu^\alpha+\mathsf k$, where $|\mathsf k|\le C_1x\varepsilon/R+C_2p_0^2R+C_a(x)p_0^2$: \Cref{Y-P6-L12} is used for $-xD\mathrm{En}_\ell$, and the first bound of the part ``Besides'' of \Cref{Y-P6-L10}, namely $|D_w\sum1\bar E^a|\le C_a(x)p_0^2$, accounts for the term $C_a(x)p_0^2$ of $\mathsf k$. Neither of the two depends on the ``Claim'' of \Cref{Y-P6-L10}, and that first bound does not depend on \Cref{Y-P6-L9} either. This remainder satisfies $|\mathsf k|\le0.5p_0x^{1/2}$: the choice of $R$ yields $C_1/R\le c_F/(60C_Sp_0^5)$, and the inequality \textup{(E3)} $\le c_F/12$ yields $(C_2R+C_a)p_0^2\le c_Fx^{1/2}/(60C_Sp_0^3)$. The identity and this bound remain valid on the set $\{\phi_p\le3\varepsilon/4+4t$ on $B'\}$, on which $|xu^\alpha|\le1.5p_0x^{1/2}+8/x$. The two bounds on $C_1/R$ and on $(C_2R+C_a)p_0^2$ give $|\mathsf k|\le c_Fx^{1/2}/(30C_Sp_0^3)\le x^{1/2}$, where the second inequality follows from $c_F<1\le C_S$. Since $8/x\le x^{1/2}$ by $x\ge L$ and $0.5p_0x^{1/2}\ge2x^{1/2}$ by $p_0\ge4$, we obtain $|\mathsf k|+8/x\le0.5p_0x^{1/2}$. Hence we have $\sup|D\Lambda^a|\le2p_0x^{1/2}$ on that set, and this is the datum that Lems.~\ref{Y-P6-L9} and~\ref{Y-P6-L10} take as input. With it, \Cref{Y-P6-L10} states that $E[xu^\alpha\Phi^\alpha]=E[D\Phi^\alpha]+O(\sup|\Phi|\sup|\mathsf k|)+\mathrm{def}(\Phi)$, where the defect $\mathrm{def}(\Phi)$ is bounded by the right side of the ``Claim'' of \Cref{Y-P6-L10} at $t=x^{-2}$: \[|\mathrm{def}(\Phi)|\le tB_2+t^{-1}e\sup|\Phi|\,(2.5(2R/M+2r(x)+3)^3c_h(x)+6e^{2.2}x^{-4}+2\mathrm{ex}),\] with $B_2=\sup|e^{-\Lambda^a}D^2(\Phi e^{\Lambda^a})|$.

On ${\rm supp}\,\chi$ one has $|u(F)|\le2(3\varepsilon/4)$, and so $G\le(9/8)p_0^2$. For the transported field strength we have \[D^\alpha\tilde F^\beta_p=(\mathrm{d}w_R)_p\delta_{\alpha\beta}+O(\varepsilon(5(r+1)|w|_p+|\mathrm{d}w|_p+2\sum_{\Gamma(0,\pi p)}|w|)).\] For the two plaquettes defining $u$, the numerical error can be bounded independently of their radii. Work in the comb gauge. In each of the 23-plaquettes based at $0$ and $e_2$, three edges belong to the comb, so their link values are $1$. If $g$ is its plaquette holonomy, Step~7.0 of the proof of \Cref{Y-P6-L10} gives
\[
 g(s)=e^{isae_\alpha}g e^{-isbe_\alpha},\qquad
 |a|,|b|\le2,\qquad a-b=(\mathrm{d}w_R)_p.
\]
Here $a$ and $b$ are the sums of the weights on the two pairs of edges. The logarithm differential from Step~1 of the proof of \Cref{Y-P6-L12} obeys $\|d\mathsf L_g\|\le1.01$: its perpendicular norm is $\phi/\sin\phi\le(1-\phi^2/6)^{-1}<1.01$ for $\phi\le0.01$. Since $\|\operatorname{Ad}_g-1\|\le2\phi$, the left-trivialized generator $ae_\alpha-b\operatorname{Ad}_g e_\alpha$ gives
\[
 |D^\alpha F_p-(\mathrm{d}w_R)_p e_\alpha|
 \le4\cdot1.01\varepsilon+2\cdot1.01\cdot2\varepsilon
 =8.08\varepsilon.
\]
The derivative of the transport to the origin vanishes at the first base. At the base $e_2$ its signed comb sum has modulus at most $1$, and its contribution is at most $2\varepsilon$. Thus the error for their difference is at most $18.16\varepsilon<50\varepsilon$. As $\chi_R=1$ near these plaquettes and $\langle u,\mathrm{d}w\rangle=a_u$, we obtain $D^\alpha u^\beta=a_u\delta_{\alpha\beta}+O(50\varepsilon)$.

\emph{Step 1: the first identity.} Let $\Phi^\alpha=u^\alpha\chi(S_f-\bar s)/2$, with $\bar s=ES_f$; then $\sup|\Phi|\le1.5C_S\varepsilon p_0^2$. On the left, $\sum_\alpha xu^\alpha\Phi^\alpha=G(S_f-\bar s)$, and the mean of this product is $\mathrm{Cov}_\nu(G,S_f)$. On the right, $\sum_\alpha D\Phi^\alpha$ is the sum of the three parts \textup{(a)}, \textup{(b)}, \textup{(c)} below.

\begin{enumerate} \item[\textup{(a)}] The first part equals $(3/2)a_u\chi(S_f-\bar s)+O(\varepsilon C_Sp_0^2)$. In the mean of its main term the factor $E[\chi(S_f-\bar s)]=E[(\chi-1)(S_f-\bar s)]\le2C_Sp_0^2\nu_P$ appears.

\item[\textup{(b)}] The second part is the one in which $D\chi$ occurs. Here $|D\chi|\le8(W_1+5\varepsilon W_2)/\varepsilon$, and $D\chi$ vanishes outside the event $\{$some $\phi_p\ge\varepsilon/2$ in $B'\}$; the mean of this part is therefore at most $(3/2)(3\varepsilon/2)(16W_*/\varepsilon)(2C_Sp_0^2)\nu_P$.

\item[\textup{(c)}] The third part is $(1/2)\sum_\alpha u^\alpha\chi DS_f$. \end{enumerate}

The remainder $\mathsf k$ produces two further errors. Its shell term $C_1x\varepsilon/R$ costs $3\cdot1.5C_S\varepsilon p_0^2\cdot C_1x\varepsilon/R=4.5C_SC_1p_0^4/R$, and the remaining part of $\mathsf k$ costs $4.5C_Sp_0^3x^{-1/2}(C_2p_0^2R+C_ap_0^2)$.

\emph{Step 2: the derivative of $S_f$.} We prove
\[
 DS_f=x\langle\mathbb D_k\mathrm{d}w,\tilde F^\alpha\rangle
             +O(C_Dp_0^2),\qquad C_D:=10^{200}C_S,
\]
where $C_D$ is a constant of this step. By \textup{(U2.3)}, only the centers $y\in Y_{\rm out}$ occur, and $|y|\le D_0$. A plaquette in the range of such a center obeys $|\pi(p)|\le D_0+\rho(x)$, which may exceed $2D_0$. Nevertheless, \textup{(j7)}\,\textup{(ii)} and $\rho(x)\le Mr(x)-2$ put all these plaquettes in the region where $\chi_R=1$. Thus $\mathrm{d}w_R=\mathrm{d}w$ there. The expanded source cubes lie in $B'$, so every indicator in $S_f$ is identically $1$ in a neighborhood of ${\rm supp}\,\chi$.

Work at the field in the root comb gauge and set
\[
 d=|y-\pi(p)|_\infty,\qquad d'=|y-\pi(p')|_\infty,
 \qquad A_{y,p}=\operatorname{Ad}_{V(\Gamma(y,\pi(p)))}.
\]
There are at most $3d$ bonds on this path, each based at radius at most $D_0+d$. The comb-gauge bound of Step~0 of the proof of \Cref{Y-P6-L12} and the adjoint bound of \Cref{Y-F-L1}\,(b) therefore give
\[
 \|A_{y,p}-1\|\le12d(D_0+d)\varepsilon,
 \qquad \|D A_{y,p}\|\le6d.
\]
The second estimate follows by differentiating the path product: each bond has speed at most $|w_b|\le1$, and the adjoint action has differential norm at most twice that speed. All these paths and their combs lie in $B'$. By \Cref{Y-P6-L11}\,(a0), $|\mathrm{d}w_R|\le\sqrt2+8/R<2$, while $|w|_p\le4$. Step~1 of the proof of \Cref{Y-P6-L12} consequently yields
\[
 |D F_p-(\mathrm{d}w)_p e_\alpha|
       \le\varepsilon[20(D_0+d)+22].
\]
Differentiate the quadratic expression of entry~5 of item~(0) of \S\ref{Y-H0-1}. Symmetry of $\mathbb D_{k,y}$ combines the two differentiated factors. Since $|F_p|,|F_{p'}|\le\varepsilon$, the preceding three estimates bound the error for a pair $(p,p')$, after removing $x|\mathbb D_{k,y}(p,p')|\varepsilon^2$, by
\[
 20D_0+26d+22+24D_0(d+d')+24(d^2+d'^2)
       \le100(D_0+1)(1+d+d')^2.
\]
Here the last two terms account for replacing both $A_{y,p}$ and $A_{y,p'}$ by $1$ in the leading term; at the field in this gauge, $\tilde F=F$.

The displacement factors are summable uniformly in the range. Put $\lambda=\lambda_\flat$. For $s\ge0$,
\[
 (1+s)^2e^{-\lambda s}
       \le(1+4/(e\lambda))^2e^{-\lambda s/2}.
\]
By \textup{(U2.1)} and the three plaquette orientations, the error is at most
\[
 100(D_0+1)p_0^2\#Y_{\rm out}K_5
       (1+4/(e\lambda))^2[3\mathsf g(\lambda/2)]^2.
\]
For completeness, writing $z=2v+e$ with $e\in\{0,1\}^3$ gives
$\mathsf g(\lambda/2)\le8e^{\lambda/2}\mathsf g(\lambda)$.
Also $\lambda=\lambda_*/16>10^{-15}$ by the pure-number intervals in \Cref{Y-A-Fr}, and $\lambda\le1$. Comparing with
$C_S=\frac12\#Y_{\rm out}K_5[3\mathsf g(\lambda)]^2$, the coefficient is at most
\[
 200(D_0+1)(1+4/(e\lambda))^2
       \left[\frac{\mathsf g(\lambda/2)}{\mathsf g(\lambda)}\right]^2C_S
       <10^{176}C_S\le C_D.
\]
This proves the derivative estimate without an upper bound on $\rho(x)$ or on the depth.

\emph{Step 3: the second identity.} Let
$\Phi'^\alpha=\chi\langle\mathbb D_k\mathrm{d}w,\tilde F^\alpha\rangle/2$; its supremum is at most $(\varepsilon/2)\cdot2\sqrt2C_S$. The sum $\sum_\alpha E[xu^\alpha\Phi'^\alpha]$ is the mean of the leading term in part~\textup{(c)} of Step~1. By Step~2, the error in that substitution is at most
$\frac94 C_D\varepsilon p_0^2$, since $|u^\alpha|\le3\varepsilon/2$ and there are three values of $\alpha$.

The derivative of the root transport at $p$ contributes at most $6|\pi(p)|\varepsilon$. Together with Step~1 of the proof of \Cref{Y-P6-L12}, this gives
\[
 |D^\alpha\tilde F_p-(\mathrm{d}w)_p e_\alpha|
                \le\varepsilon(26|\pi(p)|+22).
\]
The same displacement-moment bound as in Step~2, using
$|\pi(p)|\le D_0+|y-\pi(p)|_\infty$ and $|\mathrm{d}w|\le2$,
therefore shows that
\[
 \sum_\alpha D^\alpha\Phi'^\alpha
       =\tfrac32q\chi+O(3C_D\varepsilon)
          +\tfrac12\sum_\alpha(D^\alpha\chi)
                  \langle\mathbb D_k\mathrm{d}w,\tilde F^\alpha\rangle.
\]
The last term and the replacement of $q\chi$ by $q$ are included in \textup{(E4)}, as are the cutoff terms of Step~1. The derivative error in part~\textup{(a)} of Step~1 is at most $150C_S\varepsilon p_0^2$, by the $50\varepsilon$ estimate of Step~0. Thus all three non-cutoff derivative errors just identified total at most
\[
 150C_S\varepsilon p_0^2+\tfrac94C_D\varepsilon p_0^2
                       +3C_D\varepsilon
       \le10^{202}C_Sp_0^3x^{-1/2}.
\]

We record the budget in \textup{(E2)} and \textup{(E3)}. The two shell costs are at most
$C_SC_1p_0^4R^{-1}(4.5+4.3/p_0^2)<5C_SC_1p_0^4/R$.
The two remaining costs from $\mathsf k$ are at most
\[
 C_S(C_2R+C_a(x))p_0^5x^{-1/2}
                       (4.5+4.3/p_0^2).
\]
Since $p_0\ge4$, $R\ge1$ and $C_2=10^{670}$,
$10^{202}\le0.01C_2Rp_0^2$. The displayed derivative errors therefore add at most $0.01$ to this last coefficient, and
$4.5+4.3/16+0.01<5$. All these errors fit within \textup{(E2)} and \textup{(E3)}; the integration-by-parts defects are those of \textup{(E6)}. This proves \textup{(i)}. 

\emph{Step 4: the bound \textup{(ii)}.} Lems.~\ref{Y-P6-L1} and~\ref{Y-P6-L2} give \[\mathsf b\,\mathrm{Cov}(G_-,\mathcal O_K(f))=\mathrm{Cov}_\nu(G_-,S_f)+\mathsf b\,\mathrm{Cov}_\nu(G_-,\mathrm{Rem})+\mathsf b\,E[(G_--g)W_f].\] Here $f=f_*$ and $g:=EG_-$, so that $E(G_--g)=0$. The second term is at most $2(9/8)p_0^2\mathsf b\sup|\mathrm{Rem}|\le$ \textup{(E1)}, and the third is bounded by \textup{(E5)}. In the first term we compare $G_-$ with $G$: one has $G_--G\ne0$ only where some $\phi_p>\varepsilon/2$ in $B'$, so $|\mathrm{Cov}(G_--G,S_f)|\le(9/2)C_Sp_0^4\nu_P$, an amount that is included in \textup{(E4)}.

\emph{Step 5: the bound \textup{(iii)}.} The third identity uses $\Phi=u^\alpha\chi/2$ and gives $EG=(3/2)a_u+\mathrm{err}_1$, the shell part of $\mathrm{err}_1$ being $(9/4)C_1p_0^2/R$. The fourth identity uses $\Phi=(x/4)u^\alpha|u|^2\chi^2$, whose supremum is at most $(x/4)(3\varepsilon/2)^3=(27/32)\varepsilon p_0^2$, smaller than $1.5C_S\varepsilon p_0^2$. For this choice $\sum_\alpha xu^\alpha\Phi=G^2$, and $\sum_\alpha D(u^\alpha|u|^2)=(3+2)a_u|u|^2+O(\varepsilon^3)$, so that $\sum_\alpha D\Phi=(5/2)a_uG\chi+\ldots$, up to terms that enter $\mathrm{err}_2$. Hence $EG^2=(15/4)a_u^2+\mathrm{err}_2$, and so $\mathrm{Var}(G)=(3/2)a_u^2+\mathrm{err}_2-3a_u\mathrm{err}_1-\mathrm{err}_1^2$. Here $\mathrm{err}_2$ contains, besides the errors of the fourth identity itself, the term $(5/2)a_u\,\mathrm{err}_1$. Each entry of $\mathrm{err}_1$ and of $\mathrm{err}_2$ is at most $C_S^{-1}$ times the error that arises in the same way, namely \textup{(E2)} for the shell term, \textup{(E3)} for the remaining part of $\mathsf k$ and, as in Step~3, for the terms $O(\varepsilon)$, \textup{(E4)} for the rows with $\nu_P$ and \textup{(E6)} for the defect: the suprema $3\varepsilon/4$ and $(27/32)\varepsilon p_0^2$ take the place of $1.5C_S\varepsilon p_0^2$, the rows with $\nu_P$ have no factor $C_S$, and for both functions $B_2\le10^3N_{B0}W_*^2p_0^5x^{1/2}$. Since $[$\textup{(E2)}$+$\textup{(E3)}$+$\textup{(E4)}$+$\textup{(E6)}$]<1$, we obtain $|\mathrm{err}_2|+6|\mathrm{err}_1|+\mathrm{err}_1^2\le8C_S^{-1}[$\textup{(E2)}$+$\textup{(E3)}$+$\textup{(E4)}$+$\textup{(E6)}$]$, which lies far below \textup{(E7)}. The passage from $G$ to $G_-$ costs $(81/16)p_0^4\nu_P$.

\emph{Step 6: the defect.} Work on the tame segments of Step~2 of
the proof of \Cref{Y-P6-L10}. The moving plaquettes and the
plaquettes entering a varying row or its mask lie in $B'$. On
these segments the corresponding energy terms are bare curvature
squares and dressed row squares.

For the plaquette product in Step~7.1 of that proof, set
$\delta_p=|\mathrm{d}w|_p+5(r_p+1)\varepsilon|w|_p$. Its first derivative has
norm at most $\delta_p$. Its second derivative is the square of
the first-order generator plus the pairwise commutators, followed
by right multiplication by $\partial V(p)$. Each commutator of
unit generators has norm at most $20(r_p+1)\varepsilon$, by
comparison with their common generator. The absolute pair weights
sum to at most $|w|_p^2/2$. On $\|g-1\|\le\varepsilon\le1/2$ the
first and second derivatives of the matrix logarithm have norms
at most $2$ and $4$, by its power series. It follows that
\[
 |DF_p|\le2\delta_p,\qquad
 |D^2F_p|\le6\delta_p^2+20(r_p+1)\varepsilon|w|_p^2 .
\]
Put $A=\sum_p\delta_p$ and $B=\sum_p|D^2F_p|$. By the sums in
\Cref{Y-P6-L10} and the weighted square sum
$\sum_p(r_p+1)|w|_p^2\le10^3r_1^4$ in the numerical bounds below,
$A\le0.401W_*$ and $B\le W_*^2$. The bare curvature squares
contribute at most $4A^2+\varepsilon B<2W_*^2$.

For a path $\Gamma$ put $a_\Gamma=\sum_{b\in\Gamma}|w_b|$.
Differentiation of its adjoint product gives
$\|DT\|\le2a_\Gamma$ and $\|D^2T\|\le4a_\Gamma^2$.
At displacement $v$, the path has length at most $3|v|_\infty$,
so $a_\Gamma\le3|v|_\infty$ and its $\ell^2$ norm over path
translates is at most $3|v|_\infty W_{00}^{1/2}$. For
$G_y=\ell_y^VF$, apply the row bound $K_*e^{-\lambda_*|v|_\infty}$
and sum over the three plaquette types. Schur's bound on the
term with $DF$ and the triangle inequality on the transport terms
give
\begin{align*}
 \|DG\|_2
 &\le 2\sqrt{N_{B0}-1}\,A
       +18K_*\mathsf g_1(\lambda_*)\varepsilon W_{00}^{1/2},\\
 \sum_y|D^2G_y|
 &\le K_*\mathsf g(\lambda_*)B+24K_*\mathsf g_1(\lambda_*)A
       +108K_*\mathsf g_2(\lambda_*)\varepsilon W_{00}.
\end{align*}
In the second line, sum the terms differentiating $F$ first over
the row centers. For the term with $D^2T$, sum $a_\Gamma^2$ over
the translates before summing the displacements.
On the varying rows, $|G_y|\le3K_*\mathsf g(\lambda_*)\varepsilon$.
The bounds in \Cref{Y-A-Sums}, together with
$\mathsf g_2(\lambda)\le(1+4/(e\lambda))^2\mathsf g(\lambda/2)$,
give $K_*\mathsf g_1(\lambda_*)<10^{72}$ and
$K_*\mathsf g_2(\lambda_*)<10^{88}$.
With $W_{00}\le60r_1^3$ and $W_*=10^3r_1^4$ we get
$\|DG\|_2<0.81\sqrt{N_{B0}}W_*$ and
$\sum_y|G_y||D^2G_y|<0.51N_{B0}W_*^2$.
Including the bare part gives $|D^2\mathrm{En}_\ell|\le5N_{B0}W_*^2$.

Where a derivative of $\chi_0$ is nonzero, $\phi_p\ge\varepsilon/2$.
The group distance gives $|D\phi_p|\le\delta_p$. Differentiating
$|F_p|$ then gives $|D^2\phi_p|\le |D^2F_p|+8\delta_p^2/\varepsilon$.
The product rule yields
\[
 |D\chi|\le8A/\varepsilon,\qquad
 |D^2\chi|\le\bigl[(64+250+64)A^2+8\varepsilon B\bigr]/\varepsilon^2
              \le100W_*^2/\varepsilon^2 .
\]
The transports in $S_f$, $u$ and the linear form of Step~3 have
displacement at most $R$, with $\varepsilon R^2\le1/288$. Thus
$|D(TF)|\le 2A+6R\varepsilon\le W_*$ and
\[
 \varepsilon|D^2(TF)|
 \le\varepsilon B+24R\varepsilon A+36R^2\varepsilon^2\le W_*^2 .
\]
For $\widehat D=\varepsilon D/W_*$, the bounds on $TF$ and its
first two scaled derivatives are therefore all $\varepsilon$.
The sum of absolute quadratic coefficients of $2S_f/x$ is at
most $2C_S$, by the definition of $C_S$. The sum of absolute
coefficients of $\langle\mathbb D_k\mathrm{d}w,\tilde F^\alpha\rangle$
is at most $8C_S$, since $|\mathrm{d}w|_p\le4$. The product rule gives
$|D^i\Phi|\le10^3C_S\varepsilon p_0^2(W_*/\varepsilon)^i$ for
$i=1,2$ in the four tests of Steps~1, 3 and~5. For the last two
tests one can replace $C_S$ by $1$. For example, the cubic
$u^\alpha|u|^2$ and its first two scaled derivatives are bounded
by $8\varepsilon^3$, $24\varepsilon^3$ and $72\varepsilon^3$.
Those of $\chi^2$ are bounded by $1$, $8$ and $232$. The factor
$x/4$ in the last test gives the stated bound.

Finally,
\[
 e^{-\Lambda^a}D^2(\Phi e^{\Lambda^a})
 =D^2\Phi+2D\Phi D\Lambda^a+
              \Phi\bigl(D^2\Lambda^a+(D\Lambda^a)^2\bigr).
\]
Use $|\Phi|\le1.5C_S\varepsilon p_0^2$,
$|D\Lambda^a|\le2p_0x^{1/2}$ and the energy estimate above.
The Cauchy term in $|\Phi D^2\Lambda^a|$ is at most
$48C_SE_a(2r+1)^3W_*^2p_0^4$ by Step~7.3 of the proof of
\Cref{Y-P6-L10}. This is $19.2C_SW_*p_0^4C_a(x)$ and hence
at most $0.32W_*c_Fx^{1/2}/p_0$ by \textup{(E3)} $\le c_F/12$.
Adding the estimates gives
$B_2\le10^3C_SN_{B0}W_*^2p_0^5x^{1/2}\le C_6p_0^5x^{1/2}$.
For the two tests of Step~5 the same bound holds with
$10^3N_{B0}W_*^2$ in place of $C_6$.
Apply \Cref{Y-P6-L10} six times for (i) and six times for (iii).
The resulting defect in (i) is bounded by \textup{(E6)}, and
the defect in (iii) is at most $C_S^{-1}$ times that bound.
\end{proof}

\begin{remark}[The factor $3/2$] For $Y\cdot Y'=(1/2)\operatorname{tr}YY'$ there are three colours of unit length. For Gaussian colours with covariance $C/x$ one finds \[\mathrm{Cov}((x/2)|u\cdot\phi|^2,(x/2)\langle\phi,D\phi\rangle)=3\cdot(1/4)\cdot2\langle Cu,DCu\rangle.\] \end{remark}



\begin{lemma}[The lower bound for $q$]\label{Y-P6-L14} In the setting of \Cref{Y-P6-N} the following hold.

\begin{enumerate} \item[\textup{(A)}] For $y$ off ${\rm supp}\,u$ with $r=|y|_\infty\ge1.2\cdot10^5$, one has $|(\Pi u)(y)-\pi(y)|\le C_{\rm lat}r^{-5}$, where $C_{\rm lat}=10^{11}$ and $\pi_i:=\partial_i(3y_1y_2/(4\pi|y|_2^5))$.

\item[\textup{(B)}] At $D_0=10^{140}$, for all rows of the runs, all $1\le k\le K$ and all $x\ge L$, one has $q\ge c_F$, and $5.4\cdot10^{-5}D_0^{-5}\le c_F\le4.0\cdot10^{-4}D_0^{-5}$. \end{enumerate} \end{lemma}

\begin{proof}[Proof of \textup{(A)}] \emph{Step 1: the multiplier.} We pass to the dual picture, in which a $2$-form becomes a vector field through the normal directions of the plaquettes. Then $u=e_1(\delta_0-\delta_{e_2})$ and $\Pi u=u-\nabla\Delta^{-1}\operatorname{div}u$. Off ${\rm supp}\,u$ this gives $(\Pi u)_i=[m_i]^\vee$ with $m_i=k^\#_ik_1k^\#_2/(k\cdot k^\#)$. On $\mathbb R^3$ the multiplier $m_{0,i}=-ip_ip_1p_2/|p|^2$ has the transform $\pi_i$, because $\partial_1\partial_2(1/r)=3y_1y_2/r^5$.

\emph{Step 2: bounds on small discs.}
All Fourier-vector norms in this proof are Euclidean. For each
coordinate $\mu$, take $q=p+ze_\mu$ with $|z|\le P/20$,
where $0<P=|p|_2\le1/2$. Then $|q|_2\le1.05P$ and
$|\sum_\nu q_\nu^2|\ge0.8975P^2$. Write
$k_\nu=iq_\nu(1+e_\nu)$ and
$k_\nu^\#=-iq_\nu(1+e_\nu^\#)$, using the continuous values
at $q_\nu=0$. The exponential series gives
$|e_\nu|,|e_\nu^\#|\le0.6304P$. The cosine series gives
\[
 \left|k\cdot k^\#-\sum_\nu q_\nu^2\right|
 \le0.0842\sum_\nu|q_\nu|^4
 \le0.0842(1.05P)^4.
\]
Thus $k\cdot k^\#=(\sum_\nu q_\nu^2)(1-\eta)$ with
$|\eta|\le0.1141P^2<0.0793P$. Moreover,
\[
 |m_0(q)|_2
 \le\frac{|q|_2^3}{2|\sum_\nu q_\nu^2|}
 \le\frac{1.05^3}{2\cdot0.8975}P<0.757P.
\]
The three numerator factors and the denominator therefore give
\[
 |m(q)-m_0(q)|_2
 \le0.757P\left(\frac{(1+0.6304P)^3}{1-0.0793P}-1\right)
 \le2.08P^2.
\]
The quotient of the parenthesis by $P$ increases on $(0,1/2]$;
its endpoint gives the last bound. These one-coordinate discs
suffice for every Cauchy derivative estimate used below.

\emph{Step 3: the difference as an integral.} We have \[\Pi u-\pi=(2\pi)^{-3}[\int_{T^3}(m-m_0)e^{ipy}-\int_{\mathbb R^3\setminus T^3}m_0e^{ipy}].\] To bound the right side we follow Step~4 of the proof of \Cref{Y-P6-L11}, where the inverse transform of a multiplier is estimated, and we take $n=6$. The homogeneous multiplier $m_0$ has degree one. Interpret its
inverse transform as a tempered distribution. To justify the
integration by parts in the high-frequency contribution, insert a
smooth cutoff that equals one on $|p|\le R_\infty$ and vanishes for
$|p|\ge 2R_\infty$. Six integrations by parts produce terms with
derivatives of this cutoff. A term with $j\ge1$ cutoff derivatives
is $O(R_\infty^{-j}R_\infty^{1-(6-j)})$ on the annulus. Its integral
is $O(R_\infty^{-2})$ and tends to zero. The sixth derivative of
$m_0$ is integrable at infinity. This gives the inverse transform
away from the origin with the bounds below. The boundary terms of
$m$ at opposite faces of the torus cancel by periodicity and
$y\in\mathbb Z^3$. For $m_0$ the inner and outer face terms cancel
before the cutoff is removed.

\emph{Step 4: the numbers.} With the bound of Step~2, the argument yields the constant $c_{2,6}\cdot6!\cdot20^6\cdot2.08<7.0\cdot10^{10}$. The multiplier $m_0$ is homogeneous of degree $1$, so the bound $|m_0|\le0.757|p|$ of that step holds on the discs $|q-p|\le|p|/20$ for every $p$; hence $m_0$ on $|p|>1/2$, on both sides of $\partial T^3$, contributes $3.6\cdot10^9$, a coefficient of $r^{-6}$. For $m$ on $|p|_2>1/2$, use the one-coordinate discs
$q=p+ze_\mu$, $|z|\le1/100$, and put $K=|k(p)|_2\ge0.318$.
Writing $D(q)=k(q)\cdot k^\#(q)$, the cosine formula and
$|\sin p_\mu|\le|k_\mu(p)|$ give
\[
 |D(q)-D(p)|
 \le2K\sinh(0.01)+2(\cosh(0.01)-1)<0.064K^2.
\]
Thus $|D(q)|\ge0.935K^2$. Also
$|k_\nu(q)-k_\nu(p)|,|k_\nu^\#(q)-k_\nu^\#(p)|\le0.0102$,
and only one component changes. Hence
\[
 |k_1(q)k_2^\#(q)|\le\tfrac12K^2+0.0102K<0.547K^2,
 \qquad |k^\#(q)|_2\le\sqrt{12}+0.0102.
\]
Consequently $|m(q)|_2\le(\sqrt{12}+0.0102)\,0.547/0.935<2.1$,
and Cauchy's estimate gives
$|\partial_\mu^6m|_2\le6!\cdot100^6\cdot2.1<3.3\cdot10^{15}$. In sum, $7.0\cdot10^{10}r^{-5}+3.4\cdot10^{15}r^{-6}\le10^{11}r^{-5}$. \end{proof}

\begin{proof}[Proof of \textup{(B)}] We use $r_c:=|y_c|_2\in[D_0,D_0+\sqrt2]$ and the two sets $Y_{\rm out}=\{|y-y_c|\le D_0/8+1/2\}$ and $Y_{\rm in}=\{|y-y_c|\le D_0/16-1\}$.

\emph{Step 1: localisation.} We write $\mathbb D_y:=\mathbb D_{k,y}[f_*]$ and $Q_y:=\mathbb D_{k,y}[f_*](\omega,\omega)$. In units of level $k$, the barycentres of the plaquettes based in $B^k(y)$ lie within $1/2$ of $y$; the places of the sites and the barycentres $x(p)$ are those of the definition ``Places'' in item~1 of \S\ref{Y-H2-1}. With \Cref{Y-H2-U2}\,\textup{(U2.3)} and ${\rm supp}\,f_0\subset\{|y-y_c|\le D_0/8\}$, this gives $q=\sum_{Y_{\rm out}}Q_y$. By \Cref{Y-H2-U2}\,\textup{(U2.1)}, $|\mathbb D_y(p,p')|\le K_5e^{-\lambda_\flat(|y-p|+|y-p'|)}$.

\emph{Step 2: the region.} The set $Z=\{|p-y_c|\le0.1875D_0+1/2\}$ contains every $p$ with $|y-p|\le D_0/16$ for some $y\in Y_{\rm out}$, and on $Z$ we have $|p|_\infty\ge0.5195r_c\ge r_1$ and $|p|_2\ge0.7347r_c$. A closed form is available: with $Y=3y_1y_2$, \[|\nabla(Y/r^5)|^2=|\nabla Y|^2/r^{10}+(25-20)Y^2/r^{12}=9(y_1^2+y_2^2)/r^{10}+45y_1^2y_2^2/r^{12}\le20.25/r^8.\] Therefore $|\pi|\le(20.25/157.9)^{1/2}/0.7347^4\,r_c^{-4}<1.23r_c^{-4}$ on $Z$. To bound the Jacobian, put $n=y/r$, where $r=|y|_2$.
Direct differentiation gives
\[
\begin{split}
 r^5\partial_i\partial_j\partial_1\partial_2(1/r)
 ={}&105n_in_jn_1n_2\\
 &-15(\delta_{ij}n_1n_2+\delta_{i1}n_jn_2
       +\delta_{i2}n_jn_1+\delta_{j1}n_in_2
       +\delta_{j2}n_in_1)\\
 &+3(\delta_{i1}\delta_{j2}+\delta_{i2}\delta_{j1}).
\end{split}
\]
For $i=j\in\{1,2\}$ the right side has modulus at most $30$;
for $\{i,j\}=\{1,2\}$ it is at most $177/4$; the remaining
cases are at most $45$. These bounds use
$|n_an_b|\le1/2$ for $a\ne b$ and $\sum_a n_a^2=1$.
In particular every component is bounded by $105$.
Since $(\mathrm{Jac}\,\pi)_{ij}=(4\pi)^{-1}
\partial_i\partial_j\partial_1\partial_2(1/r)$, the nine entries
and the lower bound $r\ge0.7347r_c$ give
$|\mathrm{Jac}\,\pi|_F\le G_4r_c^{-5}$ with
$G_4=3\cdot105/(4\pi\cdot0.7347^5)<117.1$. On $Y_{\rm in}$ one has $y_1,y_2\ge0.6446r_c$ and $r\le1.0902r_c$; it follows that $|\pi|^2r_c^8\ge0.0374$ there, and also $|\pi|^2r_c^8\le0.2689$.

\emph{Step 3: $\omega$ against $\pi$ on $Z$.} \Cref{Y-P6-L11}\,\textup{(b)} gives \[|\omega-\Pi u|\le[C_G\cdot50.9/r_c^2+13.8C_\tau\tau_x^2]r_c^{-4}\le(10^{-98.6}+3\cdot10^{150}x^{-512})r_c^{-4},\] and \textup{(A)} gives $|\Pi u-\pi|\le26.5C_{\rm lat}r_c^{-5}$.

\emph{Step 4: the split.} For fixed $y$ let $\bar\omega_y$ be the constant form that has the components of $\omega$ at $y$, and set $d_p:=\omega_p-\bar\omega_{y,\iota(p)}$. For $|y-p|\le D_0/16$, \[|d_p|\le\sqrt3|y-p|G_4r_c^{-5}+2C_{\rm lat}(0.5195r_c)^{-5}+3|y-p|[C_G'(0.5195r_c)^{-7}+\ldots];\] here the gradient is taken only of the smooth function $\pi$ and of $\omega-\Pi u$. Put $\mathsf t:=r_c^4\sum_pe^{-\lambda_\flat|y-p|}|d_p|$. Then $\mathsf t\le3\sqrt3g_{1b}G_4/D_0+159\mathsf g_\flat C_{\rm lat}/D_0+\ldots+$ the far part of the sum, and this bound is at most $3.3178\cdot10^{-76}+3.3\cdot10^{-81}+6\cdot10^{-158}+8.5\,r_c^4\,\mathsf g'e^{-\lambda_\flat D_0/32}$, where $\lambda_\flat D_0/32\ge4\cdot10^{123}$. Here $6\cdot10^{-158}$ bounds the contribution of all three terms of the bracket in the bound on $|d_p|$. The term $3|y-p|C_G'(0.5195r_c)^{-7}$ contributes $9g_{1b}C_G'\,0.5195^{-7}D_0^{-3}\le4.2\cdot10^{-176}$, summed over the three orientations; the other two terms, denoted by the dots in the bracket, are those of \Cref{Y-P6-L11}\,\textup{(b)} at $0.5195r_c$, namely $1.9\cdot10^{36}\tau_x^2(0.5195r_c)^{-5}$ and $A_{\rm out}(0.5195r_c)^{-16}$, and the largest of the three contributions is $9g_{1b}A_{\rm out}0.5195^{-16}D_0^{-12}\le5.6\cdot10^{-158}$.

\emph{Step 5: the error from replacing $\omega$ by $\bar\omega$.} With $\log_{10}K_5\le24.776$, $\log_{10}\mathsf g_\flat\le46.315$, the factor $1.23$ and $\mathsf t\le3.318\cdot10^{-76}$, \[|\mathbb D_y(\omega,\omega)-\mathbb D_y(\bar\omega,\bar\omega)|\le K_5(6\mathsf g_\flat\cdot1.23\mathsf t+\mathsf t^2)r_c^{-8}\le3.02\cdot10^{-4}r_c^{-8}.\]

\emph{Step 6: the charge.} For the constant form, $\mathbb D_y(\bar\omega,\bar\omega)=c_y(\bar\omega,\bar\omega)$. By \Cref{Y-H2-U5}\,\textup{(S2)}, $c_y$ coincides with ${\rm diag}\,\mathrm{Mn}$ up to $0.0102/L$ in operator norm. Moreover $\mathrm{Mn}\ge0$, and $\mathrm{Mn}=1$ on $Y_{\rm in}$.

\emph{Step 7: the number of points.} We have $\#Y_{\rm out}/\#Y_{\rm in}\le((D_0/4+2)/(D_0/8-3))^3\le8.01$. Hence \[\FitDisplay{\dimexpr\linewidth-1pt\relax}{q\ge\sum_{Y_{\rm in}}|\bar\omega|^2-\#Y_{\rm out}(3.02\cdot10^{-4}+0.0155/L)r_c^{-8}\ge\sum_{Y_{\rm in}}|\Pi u|^2[1-10^{-97}-8.01(3.02\cdot10^{-4}+0.0155/L)/0.037],}\] where the term with $1/L$ costs $3.4/L<10^{-148}$. The bracket is at least $0.93\ge3/4$, and so $q\ge c_F$.

\emph{Step 8. Bounds on $c_F$.}
On $Y_{\rm in}$, Step~2 gives
$0.0374\le r_c^8|\pi(y)|^2\le0.2689$.
By (A) the difference $r_c^4|\Pi u-\pi|$ is at most
$10^{11}(0.6446)^{-5}/r_c<10^{-127}$ there. In particular,
$0.037\le r_c^8|\Pi u|^2\le0.2690$.
Counting sites in the sup cube gives
$(D_0/8-3)^3\le\#Y_{\rm in}\le(D_0/8+1)^3$.
With $D_0\le r_c\le D_0+\sqrt2$, we obtain
\[
 5.4\cdot10^{-5}D_0^{-5}\le
 \tfrac34\sum_{y\in Y_{\rm in}}|\Pi u(y)|^2
 \le4.0\cdot10^{-4}D_0^{-5}.
\] \end{proof}



\begin{remark}[The inner numerals of Step~2 of the proof of \textup{(A)}] In Step~2 of the proof of \textup{(A)}, the quotient $(e^t-1-t)/t$ at $t=0.525$ is less than $0.6304\cdot0.5$; the sharper bound $|\eta|\le0.069|p|$ lies below $0.0793|p|$; and $0.757[(1+0.6304|p|)^3/(1-0.0793|p|)-1]\le2.073|p|$ for $|p|\le1/2$, which lies below $2.08$. \end{remark}

\begin{lemma}[Assembly of the clauses]\label{Y-P6-L15} In the setting of \Cref{Y-P6-N}, the clauses \textup{(N1)} to \textup{(N6)} and \textup{(N6$'$)} of that theorem hold at each $j\in J(\mathsf b)$. \end{lemma}

\begin{proof} \emph{Step 1: the clause \textup{(N1)}.} By \Cref{Y-P6-L13}\,\textup{(ii)}, \Cref{Y-P6-L14} and \textup{(j1)}, \[\mathsf b\,\mathrm{Cov}\ge(3/2)q-6c_F/12\ge(3/2)c_F-c_F/2=c_F.\]

\emph{Step 2: the clause \textup{(N2)}.} \Cref{Y-P6-L13}\,\textup{(iii)}, together with $a_u\le2$ and \textup{(E7)} $<1$, gives $\mathrm{Var}\le(3/2)\cdot4+1=7$.

\emph{Step 3: the clause \textup{(N3)}.} Clause \textup{(W)} of \Cref{Y-P6-L1} applies. Indeed, the support of $f_*$ lies in $y_1\ge(0.707-0.125)D_0\ge0.58D_0$ and in $y_1\le D_0<NL^j/2$, while $G_-$ depends on $P_0$, on $P_0+e_2$ and on the one bond between them, all of which have $y_1=0$, and on the plaquettes of $B'$ on the left. By \textup{(N1)}, $\mathrm{Cov}\ge c_F/\mathsf b>0$. Therefore \textup{(N0)} forces $\mathrm{Var}(G_-)>0$, and $\mathsf b^2S^T_2(\theta f_*,f_*)\ge c_F^2/\mathrm{Var}(G_-)\ge c_F^2/7$.

\emph{Step 4: the clause \textup{(N4)}.} The function $f_*$ is smooth, and $\theta f_*$ and $f_*$ have disjoint supports. Hence \Cref{Y-H9-T}\,\textup{(T3)} applies: it provides the limit through all cutoffs $K\ge j+1$ at a fixed side, and the closed inequality persists in that limit. For the limits \textup{(a)} and \textup{(c)} of \Cref{Y-P2-ThR}\,\textup{(R5)} at the fixed label $\mathsf b$, one has $K_i\ge j+1$ for every $i$ from some index on, and once more the closed inequality persists.

Consider next the limits \textup{(a)} and \textup{(c)} of \textup{(R5)} along labels $\mathsf b_i\to \mathsf b$ with $\mathsf b_i\ge x_s$. We have $\mathsf b_i\in[\mathsf b/2,2\mathsf b]$ for all but finitely many $i$. We take $j\ge j^+(\mathsf b)$, which is finite by \Cref{Y-P6-L16}\,\textup{(d)}; then $j\ge j_*(\mathsf b_i)$, and so \textup{(N3)} holds at $(j,\mathsf b_i)$. The function $f_*$ and the mirror depend on $j$ alone, whereas $G_-$, $R$ and $h$ may change with $i$. Further, $\mathsf b_i^2\to \mathsf b^2$, and for all sufficiently large $i$ the sides satisfy \textup{(Tor)} at $\mathsf b_i\le2\mathsf b$. In the case \textup{(c)} we have in addition $K_i\ge j+1$.

Labels $\mathsf b_i<x_s$ can occur only when $\mathsf b=x_s$, and they are handled by \textup{(N6$'$)}. For the limit \textup{(b)} of \textup{(R5)} we note that $S^T_2$ is a polynomial in means of bounded continuous local functions; it therefore passes to the weak limit in the side at each $K\ge j+1$, and then along $K$.

\emph{Step 5: the clause \textup{(N5)}.} Let $S$ be a limit of \textup{(R5)}. For such limits, and only for them, \Cref{Y-P2-ThR}\,\textup{(R6)}\,\textup{(iii)} and \textup{(R7)} supply all real translations, the space $\mathcal H(S)$ and the operators $e^{-sH}$. By \textup{(R6)}\,\textup{(iii)} we move the mirror to that of \textup{(R7)}. In $\mathcal H(S)$, let $v$ be the class of the tuple with the single entry $f_*$. Then $\langle\Omega,v\rangle=\sigma_S(f_*)=0$, because the centred functional is a sum over partitions into blocks of size $\ge2$. Moreover $\|v\|^2=S_2(\theta f_*,f_*)\ge c_*/\mathsf b^2>0$, so that $\dim \mathcal H(S)\ge2$. Finally $S_2(\theta\tau_sf_*,\tau_sf_*)=\|e^{-sH}v\|^2>0$, since $e^{-sH}$ is injective, $H$ being self-adjoint and $\ge0$. Thus, for $s>0$, the argument yields $S_2(\theta\tau_sf_*,\tau_sf_*)>0$, but no lower bound for this quantity.

\emph{Step 6: the clause \textup{(N6)}.} Let $\mathsf b>0$, and choose $z$ with $L^z\mathsf b\ge x_s$. We regard the lattice $\mathbb Z^3/NL^K\mathbb Z^3$ at $\beta=L^K\mathsf b$ as the lattice with the data $(K-z,NL^z,L^z\mathsf b)$; then $a_{K-z}=L^za_K$. With $f^{[z]}(\xi):=f(L^{-z}\xi)$ we have $\mathcal O_{K-z}(f^{[z]})=L^{-z}\mathcal O_K(f)$, and hence $\mathsf b^2S_2^{(\mathsf b)}(f,f')=(L^z\mathsf b)^2S_2^{(L^z\mathsf b)}(f^{[z]},f'^{[z]})$. We apply \textup{(N3)} to \textup{(N5)} at the dilated data, for $K\ge z+j+1$ and for the sides of \textup{(R4)}, with the test function $f_*^{[-z]}$; for this we cite \Cref{Y-H9-T}\,\textup{(T6)} and \Cref{Y-P2-ThR}\,\textup{(R4)}, \textup{(R6)}\,\textup{(vi)}. In the label $\mathsf b$ the depth is $j+z$: indeed $f_*^{[-z]}(\xi)=f_0(L^{j+z}\xi)$, and $x=L^{j+z}\mathsf b$ is the same number. With $z=z_0(\mathsf b)$ this is the rule $j_*(\mathsf b)=z+j_*(L^z\mathsf b)$ of the definitions of \Cref{Y-P6-N}.

\emph{Step 7: the clause \textup{(N6$'$)}.} Let $\mathsf b>0$ with $L^i\mathsf b\ge x^{\#\#}$. Choose $z\in\mathbb Z$ with $\mathsf b':=L^z\mathsf b\in[x_s,L^2x_s]$, and put $j:=i-z$. Then $x=L^j\mathsf b'=L^i\mathsf b\ge L^3x_s\ge L\mathsf b'$, so $j\ge1$; moreover $x\ge\max\{x^\#,y(\mathsf b')^{4/3}\}$, so $j\ge j_*(\mathsf b')$. By Step~6, $\mathsf b^2S_2^{(\mathsf b)}(f,f')=\mathsf b'^2S_2^{(\mathsf b')}(f^{[z]},f'^{[z]})$. Here $f^{[z]}=f_0(L^j\cdot)$ if and only if $f=f_0(L^i\cdot)$, and the mirror $y_1=1/2$ of depth $j$ lies at $\xi_1=L^{-z}L^{-j}/2=L^{-i}/2$. We apply \textup{(N3)} to \textup{(N5)} at $(K-z,NL^z,\mathsf b')$. This requires $K-z\ge j+1$, that is $K\ge i+1$, and the sides of \textup{(R4)}, namely $m+z\ge m'$ and $L^zl\ge2Ms(L^{z+1}\mathsf b)$. For $\mathsf b>L^2x_s$ the integer $z$ is negative, and only large sides are reached; for the limits of \textup{(R5)} this loses nothing. For moving labels $\mathsf b_n\to \mathsf b$ we choose $z$ with $L^z\mathsf b\in[L^{1/2}x_s,L^{3/2}x_s]$ and require $L^i\mathsf b\ge Lx^{\#\#}$; then the same $z$ and the same $i$ can be used for all large $n$. \end{proof}

\begin{lemma}[The conditions on the depth]\label{Y-P6-L16} In the setting of \Cref{Y-P6-N}, every condition on $x=L^j\mathsf b$, and on $\mathsf b$ and $m'$, that Lems.~\ref{Y-P6-L1} to~\ref{Y-P6-L15} use is an entry of the list of \Cref{Y-P6-N}, or follows from that list. These conditions affect no floor. Moreover the following hold.

\begin{enumerate} \item[\textup{(a)}] Each entry of \textup{(X)} holds for all sufficiently large $x$. The number $x^\#$ is finite and depends on the table alone, and \textup{(X)} holds at every $x\ge x^\#$.

\item[\textup{(b)}] In the setting \textup{(X)}, $h<0.71x^{1/4}$, and the entries that \Cref{Y-P6-N} lists as holding of themselves do hold.

\item[\textup{(c)}] $[j_*(\mathsf b),\infty)\subset J(\mathsf b)$.

\item[\textup{(d)}] $\sup_{[\mathsf b_1,\mathsf b_2]}j_*\le\max\{1,\lceil\log_L(\max\{x^\#,\sup y^{4/3}\}/\mathsf b_1)\rceil\}$.

\item[\textup{(e)}] The clauses \textup{(N1)} to \textup{(N3)} at the depth $j$ follow from \textup{(X)} and \textup{(B)} taken at the single value $x=L^j\mathsf b$, that is, from $j\in J(\mathsf b)$ in the sense of the definitions of \Cref{Y-P6-N}; the list enters at this one value of $x$ and at no other. \end{enumerate} \end{lemma}

\begin{proof} \emph{The growth of the quantities.} By \textup{(O8)} of item~1 of \S\ref{Y-H9-1}, $p_0=A_0(\log x)^{\mathfrak p_0}$ with $\mathfrak p_0>2$, and $\varepsilon=p_0x^{-1/2}$. The quantities $\eta_w=e^{-c_\Xi p_0^2}$ and $c_h(x)=e^{-c_hp_0^2}$ both decay faster than every power of $x$, since $p_0^2\ge A_0^2(\log x)^4$. By the definition ``At $x$'' in item~1 of \S\ref{Y-G0-1}, and by item~2 of \S\ref{Y-H6-1}, \textup{(o4)}, we have $\rho=\lceil c_R\log x\rceil$, $r=\lceil20c_R\log x/M\rceil$, $\ell_w=6Mr+12$ and $\check r=r+3$, so that $\rho$, $Mr$ and $\ell_w$ are linear in $\log x$. From $\sigma_M\le1/12$ (item~(1) of \S\ref{Y-H0-1}, by \textup{(M2)}) and $\kappa_1\ge7$ we obtain $20c_R/M\le1/768$, and hence $r(x)\le1+\log x/768$. Finally $R$ is of the order $(\log x)^{5\mathfrak p_0}$, where $5\mathfrak p_0>10$, and $h\in[m,Lm)$ with $m=\max\{16R,64LM\log x\}$.

\emph{The entries of \textup{(X)}.} Every entry is a lower bound on $x$: the quotient of its left side by its right side tends to $0$ as $x$ grows. We take the entries in turn.

\begin{enumerate} \item[\textup{(j1)}] The errors \textup{(E1)} to \textup{(E6)} are each at most $c_F/12$. This entry is used in Steps~1 and~2 of the proof of \Cref{Y-P6-L15}, in \Cref{Y-P6-L5}\,\textup{(c)} for the smallness and for the rough bound, in \Cref{Y-P6-L2} for $\mathsf b\sup|\mathrm{Rem}|\le1$. As to the growth: \textup{(E1)}, \textup{(E3)} and the first term of \textup{(E6)} are at most a constant times a power of $\log x$ times $x^{-1/4}$, $x^{-1/2}$ and $x^{-3/2}$, respectively; \textup{(E2)} is at most $c_F/(12p_0)$ by the choice of $R$, so it imposes no condition; \textup{(E4)} and the $c_h$-term of \textup{(E6)} are a power of $x$ times $e^{-cp_0^2}$; and \textup{(E5)}, of the order $p_0^{-5}$, decays the slowest.

\item[\textup{(j5)}] $C_{\rm ext}r_B^2\varepsilon\le1$. It is used in \Cref{Y-P6-L4}, for $a\le0.01$, for $5\varepsilon W_2\le1$ and for $h<0.71x^{1/4}$. It holds for large $x$ because its left side is at most a constant times a power of $\log x$ times $x^{-1/2}$.

\item[\textup{(j6)}] This entry consists of three inequalities. The first, $x^{10}\eta_w\le1$, is used for the caps in Lems.~\ref{Y-P6-L8} and~\ref{Y-P6-L9}. The second, $p_0^{10}\le x$, is used for $\varrho x\le x^{2.1}$ and in \textup{(j4)}. The third, $(r_z/\mathsf b)(4C_Sp_0^2+5)\le x$, is used in Step~6 of the proof of \Cref{Y-P6-L8}; here $r_z/\mathsf b=3.5E_ap_0^6p_0(Lx)/(c_Q^{\rm su}K_5L^{1/2})$ is at most a constant times a power of $\log x$ and is free of $\mathsf b$ (item~2 of \S\ref{Y-H6-1}, \textup{(o5)}).

\item[\textup{(j7)\,(i)}] $(c_P/2)p_0^2\ge L^3(24\mathsf N_Z+4.5\log x+5+220M^{-3})$. This entry is used for the second inequality of \Cref{Y-P6-L3}. It holds for large $x$ because $p_0^2$ outgrows $L^3\log x$, as $2\mathfrak p_0>4$.

\item[\textup{(j7)\,(ii)}] $R/2-1\ge2D_0+M(r(x)+2)$. It is used to make the indicators of $S_f$ equal to $1$ and to have $\mathrm{d}w_R=\mathrm{d}w$ (Step~2 of the proof of \Cref{Y-P6-L13}), and for the dependence on $\mathrm{Box}$ alone (\Cref{Y-P6-L5}\,\textup{(c)}).

\item[\textup{(j7)\,(iii)}] $8\rho+8\le R$. It is used in \Cref{Y-P6-L11}\,\textup{(d)}, in Step~7 of the proof of \Cref{Y-P6-L12}, and for the inclusion of $B'$ in $\mathrm{Box}$.

\item[\textup{(j7)\,(iv)}] $8M(r(x)+1)\le R$. It is used to compare $\tau(X)$ with $m(X)$ (\Cref{Y-P6-L9}).

The entries \textup{(ii)} to \textup{(iv)} hold for large $x$ for one common reason: $R\ge10^{1410}A_0^5(\log x)^{5\mathfrak p_0}$, to be compared with $D_0$, with $Mr(x)\le20c_R\log x+M$ and with $\rho\le c_R\log x+1$. They are genuine conditions, since $M=L^{m'}$ may exceed $10^{1411}$; what carries them is the growth of $p_0^5$ against $\log x$, and not the numeral.

\item[\textup{(j7)\,(vi)}] $p_0\ge4$, $(c_a+c_h)(x)\le1/2$ and $\log x\ge3$. They are used in Steps~6 and~7 of the proof of \Cref{Y-P6-L3}, in \textup{(E2)}, and in \textup{(u1)}, \textup{(u2)}.

The fourth entry of \textup{(j7)\,(vi)} is $\Theta(x)\le1$. It is used for $C_a(x)$ and for the second-order Cauchy estimate, in Steps~7.2 and~7.3 of the proof of \Cref{Y-P6-L10}. It holds for large $x$ because \[\Theta\le p_0x^{-1/2}(12M(1+\log x/768)+24+4/\vartheta_0^\flat)\cdot78r_1^3(3252+5\log p_0),\] which is at most a constant times a power of $\log x$ times $\log\log x$ times $x^{-1/2}$; here $\log R\le3251+5\log p_0$.

\item[\textup{(j7)\,(vii)}] $\rho_0\le\rho(x)+1$. It is used for the mask in Step~2 of the proof of \Cref{Y-P6-L10} and for the radius $R+2\rho+3$ of $B'$. It holds for large $x$ because $\rho(x)$ is nondecreasing and tends to infinity. \end{enumerate}

\emph{The entries of \textup{(B)}.} They are used as follows.

\begin{enumerate} \item[\textup{(j2)}] $c_{\rm vol}(\mathsf b)h^3L^{-3j}\le1$. It is used for the root in \Cref{Y-P6-L7}.

\item[\textup{(j3)}] $i_h\le j+m'-1$. It is used for the fact that $\mathsf n_\mu=2L^{m_l+j-i_h}$ is an even integer $\ge2L$, and for the side $\ge10r_B+2$ (Lems.~\ref{Y-P6-L4} and~\ref{Y-P6-L5}). \end{enumerate}

\emph{Step \textup{(a)}: $x^\#$ is finite.} The list \textup{(X)} has finitely many entries. By the growth of the quantities, each entry has the form $A(x)\le1$ with $A(x)\to0$, so that each holds for all sufficiently large $x$. Among the definitions of \Cref{Y-P6-N}, that of $x^\#$ requires of a number $y$ that \textup{(X)} hold at every $x\ge y$, and hence no monotonicity of the $A(x)$ in $x$ is needed. The quantities $R$, $h$, $r$, $\rho$ are ceilings, and an entry such as \textup{(j7)\,(ii)} may become true just after a jump of $R$. The set of such $y$ is an upper half-line, whose endpoint may be excluded. With $x^\#:=1+\inf$, \textup{(X)} holds at every $x\ge x^\#$. The number $x^\#$ depends on the letters of the table and on pure numbers alone; it is free of $\mathsf b$, $j$, $K$ and $N$.

\emph{Step \textup{(b)}: in the setting \textup{(X)}.} First, $h<0.71x^{1/4}$: as $r_B>h/8$, \textup{(j5)} gives $32(h^2/64)p_0x^{-1/2}<1$, hence $h^2<2x^{1/2}/p_0\le x^{1/2}/2$, by $p_0\ge4$. The first inequality of \textup{(j4)} is used for $\mathrm{sep}_k\ge8\log(4\varrho L^{1/2})$ in \Cref{Y-P6-L6}. Its left side is at least $32L\log x-1-2(\log x/768+4)-2\ge31L\log x$, and its right side is $8\log4+4\log x+48\log p_0+4\log L\le11.1+12.8\log x$, by $p_0^{10}\le x$ and $x\ge L$. The second inequality of \textup{(j4)} has its left side at least $31.9\log x-11$ and its right side $\mathrm{sep}^\circ=8\log8+4\log L\le16.64+4\log x$.

The entry \textup{(j7)\,(v)} holds since $\varrho L^{1/2}\ge L^{1/2}\ge75$, because $L>1.2\cdot10^{13}$; it is used for the numeral $0.029996$ of \Cref{Y-P6-L6}. Next, $R\ge6r_1^3\ge4r_1$, because $R\ge10^{1410}p_0^5$; the bound $R\ge6r_1^3$ is used for $5\varepsilon W_2\le1$ (Step~7.2 of the proof of \Cref{Y-P6-L10}). Further, $4t\sup|w|<\varepsilon/4$, because $|w|\le1$ and $16x^{-2}<p_0x^{-1/2}$; and $x\ge L$, because $j\ge1$ and $\mathsf b\ge x_s\ge1$. Finally, \textup{(E7)} $<1$.

\emph{Step \textup{(c)}: $[j_*(\mathsf b),\infty)\subset J(\mathsf b)$.} As $L^j=x/\mathsf b$, the entry \textup{(j2)} holds if and only if $x/h\ge \mathsf b\,c_{\rm vol}(\mathsf b)^{1/3}$. The entry \textup{(j3)} says that $h\le(x/\mathsf b)L^{m'-1}$, that is, $x/h\ge \mathsf bL^{1-m'}$. By $h^2<x^{1/2}/2$ of Step~\textup{(b)}, $x/h>1.41x^{3/4}\ge y(\mathsf b)$ as soon as $x\ge y(\mathsf b)^{4/3}$. Every $j\ge j_*(\mathsf b)$ has $L^j\mathsf b\ge\max\{x^\#,y(\mathsf b)^{4/3}\}$, so \textup{(X)} and \textup{(B)} hold at $j$.

\emph{Step \textup{(d)}: local boundedness.} The function \[c_{\rm vol}(\mathsf b)=225/M^3+(3/2)\mathsf N_Zp_0(\mathsf b)^2+3\log(384\mathsf b^{3/2}/p_0(\mathsf b)^3)\] is continuous on $[x_s,\infty)$, and so is $y$. Hence $\sup_{[\mathsf b_1,\mathsf b_2]}j_*\le\max\{1,\lceil\log_L(\max\{x^\#,\sup y^{4/3}\}/\mathsf b_1)\rceil\}$. On the whole of $[x_s,\infty)$, however, $j_*$ is unbounded.

\emph{Step \textup{(e)}: one $x$ at a time.} We show that, at the depth $j$, the clauses \textup{(N1)} to \textup{(N3)} require the list only at the one value $x=L^j\mathsf b$. Finer levels $k-m$, with $x_{k-m}=L^mx$, appear in two places in which a bound is needed at $L^mx$, $m\ge1$. In \Cref{Y-P6-L6} they occur through \textup{(G1)} (item~3 of \S\ref{Y-H6-1}) and through $N_1(k-m)c_a(L^mx)$; this product is multiplied by the factor $\exp\{-5.75L^{m-1}\mathrm{sep}_k\}$, where $\mathrm{sep}_k\ge31L\log x$. Step~0 of the proof of \Cref{Y-P6-L11} uses them through $\mathsf t_*(\rho(L^mx))$. In both places the argument depends on no entry of the list at $L^mx$. Coarser levels $i>k$, with $x_i\ge L\mathsf b$, occur in \Cref{Y-P6-L7} alone, through clauses of Sections~\ref{Y-A-1} to~\ref{Y-H9-1} under their own floors. Thus \textup{(N1)} to \textup{(N3)} hold at every $j\in J(\mathsf b)$.

Every condition is a condition on $j$; none of them is a condition on $x_s$, $L$, $M$, $\kappa$ or $c_R$. Hence these conditions affect no floor. \end{proof}



\emph{Item 17: the numerals in one table.} Each constant of the section is stated below by its numerical value or, where letters of the table \textup{(O0)} to \textup{(O13)} occur, by a closed expression in those letters.

\emph{Input.}
Use the exact definitions in \Cref{Y-A-PureNumbers} and the
sums of \Cref{Y-A-Sums}. For $\lambda>0$, writing $z=e^{-\lambda}$
gives the useful closed forms
\[
 \mathsf g(\lambda)=1+\frac{24z(1+z)}{(1-z)^3}+\frac{2z}{1-z},
 \qquad
 \mathsf g_1(\lambda)=
       \frac{24z(1+4z+z^2)}{(1-z)^4}+\frac{2z}{(1-z)^2}.
\]
They follow by differentiating the geometric series and counting
the $24r^2+2$ sites in the shell of radius $r\ge1$.
The bounds used below are
$\log_{10}\mathsf g_\flat\le46.315$,
$\log_{10}g_{1b}\le61.7366$ and $\log_{10}K_5\le24.776$.
Here $g_{1b}=\mathsf g_1(\lambda_\flat)$, and $e_R\ge256$.

Since $\kappa\ge6$, the connected-set sums are bounded by
their values at $\kappa=6$. In particular,
$(1-26e^{-5})^{-1}<1.2124$,
$52e/(1-52e^{-5})<217.59$ and
$\kappa_1-\log(26e)>2.7419$.

\emph{The numerals of \Cref{Y-P6-L11}.}
The defining formulas and the sums in \Cref{Y-A-Sums} give
\[
 \mu_1\le3.78\cdot10^{71},\quad
 \check\Lambda\le1.56\cdot10^{77},\quad
 p_1\ge1.32\cdot10^{-73},\quad
 d_{\rm out}\ge4.2\cdot10^{-90},\quad
 A_{\rm out}\le10^{1455.5}.
\]
We choose
\[
 C_w=5.3\cdot10^{34},\quad C_{gw}=9.7\cdot10^{34},\quad
 C_\omega=7.0\cdot10^{34},\quad C_{g\omega}=1.3\cdot10^{35},
\]
and
\[
 C_G=4.8\cdot10^{179},\quad C_G'=8.7\cdot10^{179},\quad
 C_\tau=1.1\cdot10^{36}.
\]
For example, $16!\,20^{16}=1.371195958099968\cdot10^{34}$.
The coefficients $c_{a,16}16!\,20^{16}A'$ for the first four
majorants in Step~2 are respectively below
$2.66\cdot10^{34}$, $4.82\cdot10^{34}$,
$3.46\cdot10^{34}$ and $6.29\cdot10^{34}$.
The remaining coefficients are below
$2.39\cdot10^{179}$, $4.35\cdot10^{179}$,
$5.07\cdot10^{35}$ and $9.21\cdot10^{35}$.
These bounds verify the choices above.
The exact formulas for $N_{B0},N_{B1},N_{B2}$ give the bounds
$4.9\cdot10^{112}$, $1.61\cdot10^{127}$ and $4.9\cdot10^{141}$.
The formula in Step~7 then gives $C_e\le7.1\cdot10^{166}$.

\emph{The numerals of \Cref{Y-P6-L12}.} Here $C_1=203C_e\le1.5\cdot10^{169}$, and $C_2:=10^{670}$; its largest contribution is at most $1.1\cdot10^{666}$.

\emph{The numerals of \Cref{Y-P6-L14}.} Here $C_{\rm lat}=10^{11}$, because $7.0\cdot10^{10}+(3.6\cdot10^9+3.3\cdot10^{15})/r\le9.8\cdot10^{10}$ for $r\ge1.2\cdot10^5$, and $D_0=10^{140}$.

\emph{The numerals of Lems.~\ref{Y-P6-L2} and~\ref{Y-P6-L13}.} We have $C_S=(9/2)K_5\mathsf g_\flat^2\#Y_{\rm out}$ and $C_R=60C_SC_1$. The closed forms above and
$\#Y_{\rm out}=(D_0/4+1)^3$ give
$1.58\cdot10^{536}<C_S<1.59\cdot10^{536}$,
$10^{707}<C_R<1.02\cdot10^{707}$ and
$\log_{10}(C_R/c_F)<1411.3$. The constant $W_*=10^3r_1^4=10^{523}$ bounds $W_1+5\varepsilon W_2+1\le402r_1^4$ alone and is no bound for $W_0$. The sums of $w_R$ (Step~7.2 of the proof of \Cref{Y-P6-L10}) are as follows: $W_0\le W_0(x)=78r_1^3(1+\log R(x))$; $W_{00}\le60r_1^3$; $W_1\le400r_1^4$; $W_2\le289r_1^3R$; the sum of $|w|$ along a comb path is at most $5r_1$; and $\sum_p(r_p+1)|w|_p^2\le10^3r_1^4$.
For the last sum, a shell of radius $r$ has at most $72r^2+6$
plaquettes, and $|w|_p\le4$ on each plaquette. Summing the shells
through $r_1+1$ gives at most $300r_1^4$. For $r\ge r_1+2$ we have
$|w|_p\le4r_1^3/(r-1)^3$, so the shell contributes at most
$1156r_1^6 r^{-3}$. The tail is bounded by $578r_1^4$.
The numerical bounds used above include $1.2124\cdot1605<2500$ and $2\cdot1.2124\cdot32\cdot0.401^2<12.5$. We set $C_6:=10^{1700}$, which covers $10^3C_SN_{B0}W_*^2<10^{1699}$. The numeral of \textup{(E4)} gives $10^5W_*C_S<1.59\cdot10^{1064}$.

\emph{Input for the rows behind $C_2$, the $10^5$ of \textup{(E4)} and $C_6$.} The rows use $|w|\le\min\{1,r_1^3r^{-3}\}$ and $|\mathrm{d}w_R|\le3r_1^4r^{-4}$; there are $72r^2+6$ plaquettes at distance $r$, and as many bonds; $N_{B0}\le4.9\cdot10^{112}$, $N_{B1}\le1.61\cdot10^{127}$, $N_{B2}\le4.9\cdot10^{141}$; and $W_*=10^{523}$ and $C_S<1.59\cdot10^{536}$.

\emph{The rows for $C_2$.} The steps named in these rows are those of the proof of \Cref{Y-P6-L12}, and every row is a multiple of $p_0^2=x\varepsilon^2$.

\begin{enumerate} \item[1.] Step~2: $5W_2+W_1\le1445r_1^3R+400r_1^4\le1600r_1^3R$, for $R\ge2.6r_1$.

\item[2.] Step~3: the curvature remainder, the derivative of the transports and their discrepancy contribute, together, at most
\[
 N_{B0}(5W_2+W_1)+18N_{B1}W_0+24N_{B1}V_1+72N_{B2}W_1
       \le200(N_{B0}+N_{B1}+N_{B2})r_1^4R.
\]
The displacement moments and the estimate $V_1\le r_1^4R$ are proved in that step.

\item[3.] Step~5: $\sum r^2|\mathrm{d}w_R|\le216r_1^4R+30r_1^4\le225r_1^4R$, which gives $28[2N_{B0}\cdot225r_1^4R+2N_{B2}\cdot370r_1^4]$; the largest piece is $28\cdot740N_{B2}r_1^4\le1.1\cdot10^{666}$.

\item[4.] Step~6: $100$, from $100\varepsilon^2$ times $x$. \end{enumerate}

Summing the rows, we obtain a total $\le4\cdot10^{665}R+1.1\cdot10^{666}+$ smaller terms $\le10^{670}R$.

\emph{The rows for the $10^5$ of \textup{(E4)}.} Every row is a multiple of $C_S\nu_P$.

\begin{enumerate} \item[\textup{(a)}] The part \textup{(a)} of Step~1 of the proof of \Cref{Y-P6-L13} gives $6p_0^2$.

\item[\textup{(b)}] The part \textup{(b)} gives $72W_*p_0^2$, from $|D\chi|\le8\cdot401r_1^4/\varepsilon=3.208W_*/\varepsilon\le16W_*/\varepsilon$. \end{enumerate}

The part with $D\chi$ of the second identity gives $3\cdot3.208W_*\cdot\sqrt2<13.7W_*$; the term $(3/2)q\nu_P$ enters with $q\le4C_S$; and the passage from $G$ to $G_-$ gives $(9/2)p_0^4$. In \textup{(iii)} the rows are $(81/16)p_0^4/C_S$, and the parts with $D\chi$ are at most $40W_*p_0^2/C_S$. Altogether the sum is at most $10^3W_*C_Sp_0^4\nu_P$, which does not exceed \textup{(E4)}.

\begin{remark}[The size of the constants] The choice of the constants $C_2$, $C_6$ and the numeral of \textup{(E4)} does not affect the clauses: no floor depends on \textup{(O14)}, and a larger constant only raises $x^\#$. \end{remark}

\section{Components and reconstruction}\label{Y-P5-1}

In this section we show that the Gibbs states at cutoff $K$ invariant under $a_K\mathbb Z^3$ make up one class, with bounds that depend neither on the cutoff nor on the state nor on the places; from this we obtain the frame and the canonical decomposition of $\sigma_S$ for $S\in\mathrm{Lim}_G(\mathsf b)$, and each component has full Euclidean invariance, reflection positivity at every mirror, a unique vacuum, clustering, the axioms of \Cref{support:28-35}, point functions and a Wightman theory (\Cref{Y-P5-ThE}). We then carry the lower bound of \S\ref{Y-P6-1} over from the tori to all periodic Gibbs states (\Cref{Y-P5-ThP}), and further to each component at each $\xi$ (\Cref{Y-P5-ThC}). In consequence (\Cref{Y-P5-Cor}), for every $\mathsf b>0$ each component of each $S\in\mathrm{Lim}_G(\mathsf b)$ is, up to a shift, the functional of some member of $\mathrm{Lim}_G(\mathsf b)$ whose canonical component measure is concentrated on that component, and it satisfies all the axioms together with $S_2\ne0$.

\begin{remark}[The premises of the results] Every result below is a consequence of the clauses named in items A to D.

\end{remark}

\emph{Item A: the model, the step and the runs.} From the data of \Cref{Y-H9-T}, in item~2 of \S\ref{Y-H9-1}, we take the model, together with $\mu_{K,l,\beta}$, $\mathcal O_K(f)$ and $S^T$, which are those of \Cref{support:150-2}. The same data give the sides $N=2L^{m_l}$ with $m_l\ge m'$, and the condition \textup{(Tor)}. By the part of item~1 of \S\ref{Y-H9-1} on distances, a tuple of type $(d;R)$ in the sense of \Cref{support:150-3}, where distances are Euclidean, is of type $(d/\sqrt3;R)$ in the sup distance. From the step of item~2 of \S\ref{Y-J-1} we take the formula of $\mathsf M(U)(c)$, which is a Borel function of $U$ built from $U$, $L$, the blocks and the paths alone. By the same item, the step, that is, the kernel $\prod_ck_{ax}(V'(c),\mathsf M(U)(c))$, is a probability kernel at every $U$. The step depends on $x$, $a$ and $L$, but on no rows, on no grid of cubes, on no members, and not on $K$. The runs of item~2 of \S\ref{Y-J-1} supply the following: $\rho_0$ is Wilson's weight at $L^K\mathsf b$; the noises are $a\,x_k$ with $x_k=L^{K-k}\mathsf b$; and the rows form the chain described in \Cref{Y-A-Towers}.

\emph{Item B. Gibbs and continuum inputs.}
For the Gibbs-state estimates we use the chessboard domination
\ref{support:18-38} and its finite-side form \ref{support:18-43}.
The symmetry and conditioning statements are
\ref{support:54-17} and \ref{support:54-18}.
For the limit we use the frame and component construction
\ref{support:35-2}--\ref{support:35-36} and the conclusions of
\ref{support:77-15}. The point-insertion bound supplies the
reconstruction hypotheses through \ref{support:77-12}.

\emph{Item C: the clauses of other sections that this section uses.} From \S\ref{Y-P2-1} we take \Cref{Y-P2-ThR}\,(R0) for every integer $c$ and for all translations by $\mathbb Z^3$. We use \Cref{Y-P2-ThR}\,(R1) for every $Q\ge1$, with $C_{\rm vol}$ to the first power, and the inequality $1+0.182\,C_{\rm vol}(N,\mathsf b)\le e^{c_{\rm vol}(\mathsf b)N^3}$. The numbers $\mathrm{sep}^\circ$, $j^\circ$, $G_1$, $B_0(\mathsf b,\check\delta)$, $c_{\rm vol}$ and $\hat c_K$ of item~2 of \S\ref{Y-P2-1} enter as well; each of them is locally bounded in $\mathsf b$.

\S\ref{Y-P3-1} contributes a single inequality of \Cref{Y-P3-ThG}\,(G1), namely $\max\{\check\Gamma(\mathsf b,\allowbreak \check\delta)e^{c_{\rm vol}},\allowbreak 2B_0(\mathsf b,\allowbreak \check\delta)\}\le A_4(\mathsf b)\check\delta^{-2}$ on $0<\check\delta\le1$. Besides this inequality, we take the letters of item~1 of \S\ref{Y-P3-1} and \textup{(Lab1)}.

\Cref{Y-P4-Th4}\,(4.3)\,(a), (b) is the moment bound of \textup{(E0)}. By \Cref{Y-P4-Th4}\,(4.6), every $S\in\mathrm{Lim}_G(\mathsf b)$ has \textup{(R0)} to \textup{(R7)} of \Cref{support:32-25}, and \textup{(R5)} holds with $(c_{\rm I},C_{\rm I},\mathsf p_{\rm I})=(1/(48L),\,24Le^{c_{\rm vol}}\hat G(\mathsf b),\,1)$. We use what \Cref{Y-P4-Th4}\,(4.7) states for all isometries, for every mirror, for the circular cone and for $\operatorname{O}(2)$. From \Cref{Y-P4-Th4}\,(4.8) we take the dilation, and from \Cref{Y-P4-Th4}\,(4.9)\,(i) the joint limits.

From \S\ref{Y-P6-1} we take the following. We use its geometry and, of its numbers, $D_0$, $c_F$ and $c_*=c_F^2/7$, and also $R$, $h=L^{i_h}$, $n_0$, $t$, $r_\eta$, $\varrho$, $r_z$ and $a_u$. The depth $j_*(\mathsf b)$ of the data enters with \Cref{Y-P6-N}\,(D) and \Cref{Y-P6-L16}\,(a), (c): the number $x^\#$ is finite and depends on the table alone, and every entry of \textup{(j1)} to \textup{(j7)} and of \textup{(N0)} to \textup{(N6)} is valid for all $j\ge j_*(\mathsf b)$. We also use \Cref{Y-P6-L16}\,(d), by which $j_*$ is bounded on compact sets. We need \Cref{Y-P6-L1}\,(L), (I), (K), that is, locality, independence, and the covariance of the $k$-fold kernel under $y_\mu\to n_\mu-y_\mu$ and under $y_\mu\to y_\mu+n_\mu$, for $n_\mu\in\mathbb Z$ of level $k$. The statement of \Cref{Y-P6-L2}, together with Steps~4 to~8 of its proof, where the remainder is bounded term by term, is used as well, and so are two parts of the proof of \Cref{Y-P6-L5}: item~1 of Step~(b), which gives $W_f=0$ at globally tame fields, and Step~(c), on dependence, that is, on the bonds that enter. Lems.~\ref{Y-P6-L3}, \ref{Y-P6-L8} and \ref{Y-P6-L9} are used in their norm forms \textup{(P$^\circ$)}, \textup{(W$^\circ$)} and \textup{(X$^\circ$)}, with \textup{(W$^\circ$)} for every constant $g\in[0,\frac98p_0^2]$. We rely, moreover, on Steps 2, 5, 6 and 7 of the proof of \Cref{Y-P6-L10} together with the part ``Besides'' of that lemma, the data of \Cref{Y-P6-L11} and its clause (c), \Cref{Y-P6-L12}, Steps 0 to 6 of the proof of \Cref{Y-P6-L13} with the list of errors \textup{(E1)} to \textup{(E7)} of \Cref{Y-P6-L13}, which in this section always appear together with that lemma, and \Cref{Y-P6-L14}\,(B), which gives $q\ge c_F$. Finally we take \Cref{Y-P6-N}\,(N3), (N6).

\emph{Item D: classical inputs.} We use the following textbook facts, each by its name: the spectral theorem, dominated convergence, normal families, the identity theorem, the reflection principle, and the mean ergodic theorem. The boundary-value estimate is \Cref{support:76-23}.

\begin{definition}[Data] We recall the settings \textup{(S1)}, \textup{(S2)}, \textup{(S3)}, \textup{(S6)}, \textup{(S7)}, \textup{(S8)} and \textup{(S9)} of \S\ref{Y-00-1}. The group is $G=\operatorname{SU}(2)$ and the dimension is 3. The table \textup{(O0)} to \textup{(O13)} and the floors are those of item~(1) of \S\ref{Y-H0-1} and of item~1 of \S\ref{Y-H9-1}; $L$ and $M=L^{m'}$ are the $L$ and the $M$ of the table, and $L$ is odd. We write $a_K=L^{-K}$; the number $\mathsf b>0$ is the label of the limit, $\mathsf b'$ is the label of an approximant, and always $\beta=L^K\mathsf b'>0$. We set \[\mathcal O_K(f)=a_K^{-1}\sum_pf(x(p))\,a_p(U),\] where $a_p=1-\frac12\operatorname{tr}\partial U(p)\in[0,2]$ and $x(p)$ is the barycentre of $p$; and $\|f\|_\sharp=\max\{\|f\|_\infty,40M\,\mathrm{Lip}f\}$.

\emph{The class.} The class $\mathbb G_K(\mathsf b')$ is the set of the $\nu\in\mathcal G_{a_K,L^K\mathsf b'}$, the Gibbs states of \Cref{support:14-9}\,(2), with $(\tau_v)_*\nu=\nu$ for all $v\in a_K\mathbb Z^3$. We have $\mathbb G_K(\mathsf b')\subset\mathcal G^{\rm per}$, with period number 1 in the sense of \Cref{support:14-9}\,(3), and $\mathbb G_K(\mathsf b')\supset\mathcal G^{\rm tor}$ by \Cref{support:14-16}\,(c). We remark that $\mathbb G_K(\mathsf b')$ is in fact all of $\mathcal G^{\rm per}$, by \Cref{support:54-17}\,(c).

\emph{The boards.} For $i\le K-1$ the board $B_i$ has its walls at $L^{-i}(\mathbb Z+\frac12)$ and its cells of side $L^{-i}$, the cell 0 being centred at 0; these are the boards $(i,0)$ of \Cref{support:18-2}, and $B_0$ is the unit board. Written without arguments, $B_0$ is this board; written with them, $B_0(\mathsf b,\check\delta)$ is the number. The boards have no offset; other places are reached by translating the functions by vectors of $a_K\mathbb Z^3$, as in \Cref{Y-P5-L1}, or of $L^{-j}\mathbb Z^3$, as in \Cref{Y-P5-ThP}.

We put $X^\nu_K=\prod_i(\mathcal O_Kf_i-\nu(\mathcal O_Kf_i))$, and $S^T_n[\rho]$ stands for the joint cumulants under $\rho$. The set $\mathrm{Lim}_G(\mathsf b)$ consists of the $S=(S_n)_{n\ge2}$ such that $S_n(X)=\lim_jS^T_n[\nu_j](X)$ for all separated $C^1_c$ tuples, for some $K_j\to\infty$, $\nu_j\in\mathbb G_{K_j}(\mathsf b_j)$ and $\mathsf b_j\to \mathsf b$; and $\mathrm{Lim}_{\rm tor}(\mathsf b)$ consists of those obtained with $\nu_j\in\mathcal G^{\rm tor}$. The functional $\sigma_S$ is the centred functional of \eqref{support:9-14}. The components of $S$ are the $\sigma_\xi$ of a canonical decomposition $(\mathsf X,\lambda,(\sigma_\xi))$ of $\sigma_S$, that is, of one with the properties of \Cref{support:35-34}. Further, $\theta$ is the reflection at $(e_3,0)$, $\mathbb E_+=\{x_3>0\}$, and $\mathrm{sh}_c$ is the algebra map $(f)\mapsto(f)+c\int f$.

\emph{Labels.} We let $z(\mathsf b)$ be the least integer $z\ge0$ with $L^z\mathsf b>x_s$; it is a $z$ as in \Cref{Y-P4-Th4}\,(4.8). The set $U(\mathsf b):=[\mathsf b/2,2\mathsf b]\cap[L^{-z(\mathsf b)}x_s,\infty)$ is a compact neighbourhood of $\mathsf b$ in $(0,\infty)$ that is determined by $\mathsf b$ alone and satisfies $L^{z(\mathsf b)}U(\mathsf b)\subset[x_s,\infty)$; for $\mathsf b>x_s$ we have $z(\mathsf b)=0$ and $U(\mathsf b)=[\mathsf b/2,2\mathsf b]\cap[x_s,\infty)$. Of a sequence $\mathsf b_j\to \mathsf b$ only finitely many labels lie outside $U(\mathsf b)$, and we drop these $j$: by \Cref{support:35-2}\,(5), a subsequence of a class system is again a class system, with the same bound system, and the limit does not change.

\emph{Normal form.} The normal form is the case $\mathsf b>x_s$. There every number ``at $\mathsf b$'' is defined as its supremum over $U(\mathsf b)$; it is thus one function of $\mathsf b$, which depends neither on the sequence $(\mathsf b_j)$ nor on the state. Numbers at the label $\mathsf b'$ of an approximant are plain values, whereas numbers at the limit label $\mathsf b$ are suprema over $U(\mathsf b)$. All the numbers are locally bounded in $\mathsf b$: see item~2 of \S\ref{Y-P2-1} and \textup{(Lab1)} of \S\ref{Y-P3-1}; further, $c_{\rm vol}$, $p_0$ and $\varepsilon$ are continuous on $[x_s,\infty)$, and for $j_*$ we refer to the paragraph ``The depth'' below. The labels outside the normal form, those with $\mathsf b\le x_s$, are the subject of \textup{(E5)}.

\emph{The depth.} We recall the depth from the definitions of \Cref{Y-P6-N}, from \Cref{Y-P6-N}\,(D) and from \Cref{Y-P6-L16}\,(a), (c). The number $x^\#$ is the one of those definitions, namely $1+\inf\{y\ge L:$ each entry of the list of \S\ref{Y-P6-1} depending on $x$ alone is valid for all $x\ge y\}$. Next, $y(\mathsf b):=\mathsf b\max\{c_{\rm vol}(\mathsf b)^{1/3},L^{1-m'}\}$, and for $\mathsf b\ge x_s$ the number $j_*(\mathsf b)$ is the least $j\ge1$ with $L^j\mathsf b\ge\max\{x^\#,y(\mathsf b)^{4/3}\}$. If $j\ge j_*(\mathsf b)$, then also $L^{j'}\mathsf b\ge\max\{x^\#,y(\mathsf b)^{4/3}\}$ for every $j'\ge j$. Moreover, as in \Cref{Y-P6-L16}\,(d), $j_*(\mathsf b')\le\max\{1,\lceil\log_L(\max\{x^\#,y(\mathsf b')^{4/3}\}/\mathsf b')\rceil\}$, and since $y$ is continuous on $[x_s,\infty)$, $j_*$ is bounded on compact sets.

\emph{The geometry of \S\ref{Y-P6-1}, recalled.} We measure in units of level $k=K-j$, that is, in units of $L^{-j}$; $r=|y|$ is the sup distance; and $x=L^j\mathsf b'$, $\varepsilon=p_0(x)x^{-1/2}$, $\rho=\rho(x)$. The mirror is the plane $y_1=\frac12$, and $u$, $T_y$, $\tilde F_p$ and $u^\alpha(F)$ are those of the geometry of \Cref{Y-P6-N}. We set $B_R=\{r\le R\}$, $B'=\{r\le R+2\rho+3\}$ and $\mathrm{Box}=\{r\le r_B\}$, where $r_B:=R+h/8+1\le(h-1)/2$. Let $\chi$ be the product over the plaquettes of $B'$ of $\chi_0(\phi_p/\varepsilon)$, where $\chi_0$ is smooth, equal to $1$ on $[0,\frac12]$ and to $0$ on $[\frac34,\infty)$, and let $\chi_-$ be the analogous product taken over those plaquettes all of whose sites lie in $y_1<\frac12$. We set $G_-:=\frac x2\sum_\alpha u^\alpha(F)^2\chi_-\in[0,\frac98p_0^2]$; it is a function of those drawn bonds of $B'$ that lie in $y_1<\frac12$, and $G$ is defined in the same way with $\chi$. Let $y_c:=\lceil D_0/\sqrt2\rceil(e_1+e_2)$, and let $f_0\in C^\infty$ with $0\le f_0\le1$, equal to $1$ on $|y-y_c|\le D_0/16$ and to $0$ off $|y-y_c|\le D_0/8$; we set $f_*:=f_0(L^j\,\cdot\,)$, and it keeps a distance $\ge0.58\,D_0L^{-j}$ from the mirror.

The pure numbers are $D_0=10^{140}$, $c_F\in[5.4\cdot10^{-5},4.0\cdot10^{-4}]\,D_0^{-5}$ and $c_*=c_F^2/7$. The following are functions of $x$ alone: $R=\lceil C_Rp_0^5/c_F\rceil$; $h=L^{i_h}$, the least power of $L$ that is $\ge\max\{16R,64LM\log x\}$; $n_0=\lfloor h/16M\rfloor$; $t=x^{-2}$; $r_\eta=x^4$; $\varrho=x^{1/2}p_0^6$; and $r_z/\mathsf b'$. Here $r_z:=\varrho\,\mathfrak r^\circ(x)$, which is not itself a member of the list, depends on $\mathsf b'$ as well as on $x$. The boards are $B_{j-1}$, with cells of side $L$, and $B_{j-i_h}$, with cells of side $h$.

\end{definition}

\begin{definition}[New numbers] All the numbers below are defined at \textup{(O14)}, after the table and $\mathsf b$. No floor depends on them, and they are independent of $K$, $N$, the state, the sequence of labels, $\xi$ and the places. From item~2 of \S\ref{Y-P2-1} we recall $\mathrm{sep}^\circ:=8\log(8L^{1/2})$, the only numeral among the numbers of that item, and, from the same item, \[c_{\rm vol}(\mathsf b):=225/M^3+\frac32\mathsf N_Zp_0(\mathsf b)^2+3\log(384/\varepsilon(\mathsf b)^3).\] We recall $\check\Gamma(\mathsf b,\check\delta)=G_1(\mathsf b,\check\delta,\check\delta/4)(1+e^{c_{\rm vol}})$; and $A_4$ and $\hat G(\mathsf b)$ are those of item~1 of \S\ref{Y-P3-1} and of item~2 of \S\ref{Y-P4-1}. From item~1 of \S\ref{Y-P3-1} we recall $A_6(\mathsf b):=3^{5/2}\cdot64(4e)^5A_4(\mathsf b)$, and we set $C_1(R):=A_6(1+R)^3$. From item~2 of \S\ref{Y-00-1} we recall \[C_n^\infty(d;R):=1.45^{n+1}n!\,A_6^n(n!)^5\max\{1,d^{-5n}\}(1+R)^{3n},\] the value $d=\infty$ included, and we set $C^{\mathsf b}:=(C_1,(C_n^\infty))$. We let $\bar c(\mathsf b)$ be the value of the formula \eqref{support:35-14} at $C^{\mathsf b}$. We set $\bar c_*(\mathsf b):=\frac3{4\pi}A_6(\mathsf b)$.

From \Cref{Y-P3-ThG}\,(G4)\,(a) we recall the quadruple $(13,1,C_E,13)$, with $C_E=\mathsf C_{\ne}^6\,M\,8A_6(\mathsf b)$ and $\mathsf C_{\ne}=10^{53}$. From item~2 of \S\ref{Y-P4-1} we recall $(c_{\rm I},C_{\rm I},\mathsf p_{\rm I})=(1/(48L),\,24Le^{c_{\rm vol}}\hat G(\mathsf b),\,1)$. The formulas \eqref{support:34-6} at $(c_{\rm I},C_{\rm I},\mathsf p_{\rm I})$ give $\wp=\mathsf p_{\rm I}+4=5$ and \[C_{\rm pt}=320\cdot8^4\,\mathsf{mv}_3\,c_{\rm I}^{-4}C_{\rm I}M,\] where $\mathsf{mv}_3=m_3$ is defined in \Cref{support:34-3}; we set $C_\wp=32C_{\rm pt}$, $c_0=\frac12$ and $C_*=C_\wp+\bar c_*$. For $\mathsf b\ge x_s$ we recall from the definitions of \Cref{Y-P6-N} that $j^+(\mathsf b)$ is the supremum of $j_*$ over $[\mathsf b/2,2\mathsf b]\cap[x_s,\infty)$; for $\mathsf b>x_s$ this set is $U(\mathsf b)$, and the supremum is finite, since $j_*$ is bounded on compact sets by the last sentence of the paragraph on the depth in the data. We set $c_{**}:=(4/L)^6c_*$ and, in the normal form, $r_*(\mathsf b):=D_0L^{-j^+(\mathsf b)}>0$; in general $r_*(\mathsf b)$ is given in \textup{(C4)}. We write $D_W$ for the domain $\mathcal D_0$ of \Cref{support:76-59}.

\end{definition}

\begin{theorem}[The class, the limits and their components]\label{Y-P5-ThE} The following hold.

\begin{enumerate}

\item[\textup{(E0)}] \emph{Class and bounds.} Let $\mathsf b'\ge x_s$ and $K\ge1$. Then $\mathbb G_K(\mathsf b')$ is a class, in the sense of \textup{(Cls1)}, \textup{(Cls2)} of \Cref{support:35-2}. For $\rho,\nu\in\mathbb G_K(\mathsf b')$, for $n\ge1$ and for real bounded functions $f_i$ of the Euclidean type $(d;R)$, wherever the balls lie, \[|\rho(X^\nu_K)|\le A_6(\mathsf b')^n(n!)^5\max\{1,d^{-5n}\}(1+R)^{3n}\prod_i\|f_i\|_\infty,\] and for $n\ge2$ also $|S^T_n[\rho](X)|\le C_n^\infty(d;R)\prod_i\|f_i\|_\infty$. These two bounds are \Cref{Y-P4-Th4}\,(4.3)\,(b). For $K_j\to\infty$, $\mathsf b>x_s$ and $\mathsf b_j\in U(\mathsf b)$, the sequence $(K_j,\mathbb G_{K_j}(\mathsf b_j))$ is a class system that obeys $C^{\mathsf b}$ in the sense of \textup{(Bd1)}, \textup{(Bd2)}.

\item[\textup{(E1)}] \emph{The limit.} Let $K_j\to\infty$, $\mathsf b_j\to\mathsf b>x_s$ and
$\nu_j\in\mathbb G_{K_j}(\mathsf b_j)$. After discarding finitely many
indices we have $\mathsf b_j\in U(\mathsf b)$. Every such frame has a
convergent subsequence. By \Cref{support:35-5}, the sequence $(K_j,\nu_j)$ is a frame in the class system $(K_j,\mathbb G_{K_j}(\mathsf b_j))$ of \textup{(E0)}. Along $(K_j,\nu_j)$ an element $S\in\mathrm{Lim}_G(\mathsf b)$ is obtained if and only if the frame converges in the sense of \eqref{support:35-6}, and then its limit $\sigma_{\rm frame}$, the functional of \eqref{support:35-6}, equals $\sigma_S$. The functional $\sigma_S$ has the properties of \Cref{support:77-15}\,(1)(a) to (e); it is invariant under $\mathbb R^3\rtimes\operatorname{O}(3)$ and reflection positive at every plane; it satisfies the growth condition \eqref{support:26-34} with the quadruple; and it has data \textup{(Rec1)} to \textup{(Rec4)} with $(C_\wp,5,\frac12)$, where \textup{(Rec1)} to \textup{(Rec5)} are as recalled in \ref{support:77-2} before \Cref{support:77-12}.

\item[\textup{(E2)}] \emph{The decomposition.} A canonical decomposition of $\sigma_S$ exists, and it is unique up to exactly one isomorphism. The map $P_0[F]\mapsto(\xi\mapsto\sigma_\xi(F))$ is a unitary map of $\ker H$ onto $L^2(\mathsf X,\lambda)$ that carries $\Omega$ to $1$. Moreover $\#\mathsf X=\dim\ker H$ when $\mathsf X$ is finite, and $\ker H$ is infinite-dimensional when $\mathsf X$ is infinite, and every isometry acts on $\mathsf X$ as the identity.

\item[\textup{(E3)}] \emph{Every limit in the frame.} Let $\sigma'$ be any limit in the frame; in particular, $\sigma'$ may be any component $\sigma_\xi$.

\begin{enumerate}

\item[\textup{(i)}] $(S_n^{\sigma'})_{n\ge2}\in\mathrm{Lim}_G(\mathsf b)$, and we write $\sigma'^{\rm c}$ for its centred functional; moreover $\sigma'=\sigma'^{\rm c}\circ\mathrm{sh}_c$ with $c=c_{\sigma'}$. The functional $\sigma'$ is invariant under $\mathbb R^3\rtimes\operatorname{O}(3)$, reflection positive at every plane, and regular of order $\le4n+1$. When $\sigma'$ is a component $\sigma_\xi$, it has approximants $\rho_j=\nu_j(\,\cdot\mid E_j)$ with $E_j$ invariant and $\nu_j(E_j)>0$.

\item[\textup{(ii)}] $S_1=c_{\sigma'}\int f$ with $|c_{\sigma'}|\le\bar c_*$; the number $c_\xi$ is continuous in $\xi$, and $\int c_\xi\,d\lambda=0$. The moment bound of \textup{(E0)} holds for $\sigma'$, and so does the growth condition \eqref{support:26-34} of \textup{(E1)}, with the same quadruple.

\item[\textup{(iii)}] $\mathcal H(\sigma')$ is separable and $H\ge0$; the joint spectrum lies in $E\ge|q|$, and $\operatorname{O}(2)$ acts. When $\sigma'=\sigma_\xi$, we have $\ker H_\xi=\mathbb C\Omega$.

\item[\textup{(iv)}] $\sigma_\xi(X\cdot\tau_vY)\to\sigma_\xi(X)\sigma_\xi(Y)$ as $|v|\to\infty$, uniformly in $\xi$; \Cref{support:36-17} gives, in addition, the error $C(X,Y)/|v|$.

\end{enumerate}

\item[\textup{(E4)}] \emph{Reconstruction.} The following hold.

\begin{enumerate}

\item[\textup{(a)}] The Schwinger functions of every $\sigma_\xi$ satisfy \textup{(E0$'$)}, \textup{(E1)} for $\mathbb R^3\rtimes\operatorname{O}(3)$, \textup{(E2)}, \textup{(E3)} and \textup{(E4)} of \Cref{support:28-35}.

\item[\textup{(b1)}] Every limit $\sigma'$ in the frame has \textup{(PtB)} of \Cref{support:77-2} with $(C_\wp,5,\frac12)$.

\item[\textup{(b2)}] Every $\sigma_\xi$ has data \textup{(Rec1)} to \textup{(Rec5)} with $(C_*,5,\frac12)$.

\item[\textup{(b3)}] \emph{Point functions.} For each $n$ there is one real-analytic, symmetric and isometry-invariant function $\mathsf S_n^\xi$ on $\{x_i\ne x_k\}$ such that $\sigma_\xi(f_1,\dots,f_n)=\int\prod_if_i(x_i)\,\mathsf S_n^\xi(x)\,dx$ for all separated tuples; and \[|\mathsf S_n^\xi(x)|\le[C_*\max\{1,(2n^2/d(x))^5\}]^n,\] where $d(x)$ is the least distance between distinct entries, with $d(x)=+\infty$ for $n=1$.

\item[\textup{(b4)}] \emph{Wightman.} For each $n$ there is a function $\mathsf W_n$, holomorphic on $T^{n-1}$, where $T=\{(z^0,\mathbf z)\in\mathbb C\times\mathbb C^2:\ \operatorname{Re}z^0>|\operatorname{Im}\mathbf z|\}$, such that $|\mathsf W_n|\le(C_*\varepsilon_T^{-5})^n$ with $\varepsilon_T=\min\{1,\frac12\min_k(\operatorname{Re}z_k^0-|\operatorname{Im}\mathbf z_k|)\}$, with $\varepsilon_T=1$ when $n=1$. At points of increasing height, $\mathsf S_n(x)=\mathsf W_n(x_2-x_1,\dots,x_n-x_{n-1})$. The quintuple $(\mathcal H(\sigma_\xi),U,\Omega,D_W,\Phi)$ is a Wightman theory of one Hermitian scalar field on $\mathbb R^{1+2}$, with the domain, covariance, locality, spectrum and continuation properties of \Cref{support:76-59}. For $\sigma_S$ all of this holds too, except the uniqueness of the vacuum.

\end{enumerate}

\item[\textup{(E5)}] \emph{Every $\mathsf b>0$.} For $\mathsf b>x_s$, use \textup{(E0)} to \textup{(E4)} with their stated constants. For $0<\mathsf b\le x_s$, put $z=z(\mathsf b)\ge1$, take $K\ge z+1$, and let every number ``at $L^z\mathsf b$'' be its supremum over $L^zU(\mathsf b)$. Then \textup{(E0)} to \textup{(E4)} hold with the changes listed below. In both cases \textup{(E7)} holds as well, since its proof uses only \textup{(E2)} to \textup{(E4)}. For a label $\mathsf b'<x_s$ of $U(\mathsf b)$, \Cref{Y-P5-ThP} is to be understood with $(K-z,L^z\mathsf b')$ in place of $(K,\mathsf b')$. The changes are these: $A_6(\mathsf b)(1+R)^3$ becomes $L^zA_6(L^z\mathsf b)(1+L^zR)^3$; the growth condition \eqref{support:26-34} of \textup{(E1)} holds with the quadruple in which $C_E$ is replaced by $\mathsf C_{\ne}^6\,M\,L^z(1+L^z)^3A_6(L^z\mathsf b)$; $\bar c_*(\mathsf b)=\frac3{4\pi}L^{4z}A_6(L^z\mathsf b)$; $(c_{\rm I},C_{\rm I})$ becomes $(c_{\rm I}L^{-z}/2,\,2L^zC_{\rm I}(L^z\mathsf b))$; and $C_{\rm pt}$ and $C_*$ are computed at these numbers, so that $C_{\rm pt}(\mathsf b)=32L^{5z}C_{\rm pt}(L^z\mathsf b)$.

\item[\textup{(E6)}] \emph{Joint limits.} For any $\mathsf b>0$, the limit points of joint sequences at labels $\mathsf b_j\to \mathsf b$ with $N_j\ge N^*(K_j,\mathsf b_j)$ are functionals of members of $\mathrm{Lim}_{\rm tor}(\mathsf b)$; this is \Cref{Y-P4-Th4}\,(4.9)\,(i).

\item[\textup{(E7)}] \emph{The relation between $S_2$ and the components.} Let $g_f(\xi):=S_2^{\sigma_\xi}(\theta f,f)$ and $V:=\int c_\xi^2\,d\lambda$.

\begin{enumerate}

\item[\textup{(a)}] For $f\in C^1_c(\mathbb E_+)$ we have $S_2(\theta f,f)=\int g_f\,d\lambda+V(\int f)^2$; here $g_f\ge0$ is continuous, $\int g_f\,d\lambda=\|(1-P_0)[(f)]\|^2$, and $V(\int f)^2=\|P_0[(f)]\|^2\le\bar c_*^2(\int f)^2$.

\item[\textup{(b)}] For a single $\xi$ the following are equivalent: $\dim \mathcal H(\sigma_\xi)=1$; $S_2^{\sigma_\xi}=0$; $g_f(\xi)=0$ for one $f\ge0$ with $f\ne0$; and $\sigma_\xi(f_1,\dots,f_n)=\prod_ic_\xi\int f_i$. The centred two-point function of $\sigma_\xi$ is $s_\xi(|x-y|)$, where $s_\xi\ge0$ is non-increasing.

\item[\textup{(c)}] The conditions $S\ne0$ and $\dim \mathcal H(S)\ge2$ by themselves decide nothing about the components; \Cref{ex:constant-mixture} gives an instance.

\end{enumerate}

\end{enumerate}

\end{theorem}

\begin{theorem}[A lower bound in every periodic Gibbs state]\label{Y-P5-ThP} Let $\mathsf b'\ge x_s$, $j\ge j_*(\mathsf b')$ and $K\ge j+1$ be arbitrary, put $k=K-j$ and $x=L^j\mathsf b'$, and let $\rho\in\mathbb G_K(\mathsf b')$ be arbitrary. Let $P_\rho$ be $\rho$ times the $k$-fold kernel of the step of item~2 of \S\ref{Y-J-1}, on the infinite lattices, and write $V=V^{(k)}$. Let $\theta_*$ be the reflection at $x_1=cL^{-j}$, where $c\in\mathbb Z+\frac12$; for $c=\frac12$, $G_-$ and $f_*$ are those of the geometry in the data, and for other $c$ they are their translates by $(c-\frac12)L^{-j}e_1$. Then the following hold.

\begin{enumerate}

\item[\textup{(P0)}] $\mathrm{Cov}_{P_\rho}(G_-,\mathcal O_Kf)^2\le\mathrm{Var}_{P_\rho}(G_-)\,S^T_2[\rho](\theta_*f,f)$ for every real bounded local $f$ carried by the open right half.

\item[\textup{(P1)}] $\mathsf b'\,\mathrm{Cov}_{P_\rho}(G_-,\mathcal O_Kf_*)\ge c_F$.

\item[\textup{(P2)}] $\mathrm{Var}_{P_\rho}(G_-)\le7$.

\item[\textup{(P3)}] $\mathsf b'^2S^T_2[\rho](\theta_*f_*,f_*)\ge c_*$.

\item[\textup{(P4)}] For $\nu\in\mathbb G_K(\mathsf b')$, $\nu$-almost surely, $\mathsf b'^2\,\mathrm{Cov}_\nu(\mathcal O_K\theta_*f_*,\mathcal O_Kf_*\mid\mathcal I)\ge c_*$, where $\mathcal I$ is the one of \Cref{support:14-9}\,(5) at $a=a_K$.

\end{enumerate}

\end{theorem}

\Cref{Y-P5-ThP} involves no torus and no side, and each of its clauses is a statement at one $j$.

\begin{theorem}[A lower bound for every component]\label{Y-P5-ThC} Let $S\in\mathrm{Lim}_G(\mathsf b)$ and $\xi\in\mathsf X$, where every point of $\mathsf X$ is admitted as $\xi$. Clauses \textup{(C1)} to \textup{(C3)} below are stated in normal form, $\mathsf b>x_s$. Clause \textup{(C4)} gives the corresponding bounds for every $\mathsf b>0$, including $\mathsf b=x_s$.

\begin{enumerate}

\item[\textup{(C1)}] For every $j\ge j^+(\mathsf b)$, with the mirror carried to $(e_3,0)$, we have $\mathsf b^2g_{f_*}(\xi)\ge c_*$; hence every component is non-trivial in the sense of \textup{(E7)}\,(b).

\item[\textup{(C2)}] $s_\xi(r)\ge c_{**}\mathsf b^{-2}r^{-6}$ for $0<r\le r_*(\mathsf b)$. Moreover, $H_\xi$ is unbounded and $\dim \mathcal H(\sigma_\xi)$ is infinite.

\item[\textup{(C3)}] $\|(1-P_0)[(f_*)]\|^2\ge c_*/\mathsf b^2$; hence $H\ne0$ on $\mathcal H(\sigma_S)$.

\item[\textup{(C4)}] The same holds for every $\mathsf b>0$, with $z=z(\mathsf b)$, $f_*=f_0(L^{j+z}\,\cdot\,)$ and $j\ge j^+(L^z\mathsf b)$, where $j^+(L^z\mathsf b)$ is the supremum of $j_*$ over $L^zU(\mathsf b)=U(L^z\mathsf b)$. The quantities $\mathsf b^2S_2$ and $\mathsf b^2r^6s(r)$ are invariant under the dilation, and $r_*(\mathsf b)=L^{-z}D_0L^{-j^+(L^z\mathsf b)}$.

\end{enumerate}

\end{theorem}

The clauses of \Cref{Y-P5-ThC} differ in the values of $j$ they require. \textup{(C1)}, \textup{(C3)} and the non-triviality of the components require one $j\ge j^+(\mathsf b)$; the lower bound on $s_\xi$ in \textup{(C2)} requires all $j\ge j^+(\mathsf b)$; and the unboundedness of $H_\xi$ together with the infinite dimension in \textup{(C2)} requires infinitely many. \Cref{Y-P6-N}\,(D) and \Cref{Y-P6-L16}\,(a), (c) provide the whole half line of such $j$.

\begin{corollary}[Each component is a theory in its own right, satisfying all axioms, with $S_2^{\sigma_\xi}\ne0$]\label{Y-P5-Cor} Let $\mathsf b>0$. Then $\mathrm{Lim}_G(\mathsf b)\ne\emptyset$. For every $S\in\mathrm{Lim}_G(\mathsf b)$, every component $\sigma_\xi$ of $S$ is, up to the shift, the functional of a member of $\mathrm{Lim}_G(\mathsf b)$ whose only component is itself: $\sigma_\xi=\sigma_\xi^{\rm c}\circ\mathrm{sh}_{c_\xi}$, where $\sigma_\xi^{\rm c}$ is the centred functional of $S^\xi=(S_n^{\sigma_\xi})_{n\ge2}\in\mathrm{Lim}_G(\mathsf b)$. Each such component has the following properties. It satisfies \textup{(E0$'$)} to \textup{(E4)} of \Cref{support:28-35}, with full Euclidean invariance. It is reflection positive at every plane. It has a unique vacuum, and it clusters. It carries the point functions of clause \textup{(E4)}(b3) of \Cref{Y-P5-ThE} and the Wightman theory of its clause \textup{(E4)}(b4). Finally, $S_2^{\sigma_\xi}\ne0$, with the lower bound \textup{(C2)}.

\end{corollary}

\begin{remark}[Two words of this section] In this section, an \emph{ergodic component} is understood as in \Cref{support:35-34}. The word \emph{non-trivial} is meant in the sense of \textup{(E7)(b)}.\end{remark}


The moment estimates give the class and bound system needed for
canonical decomposition. The centering law is held fixed while the
integrating law is conditioned. The point-insertion bound then
applies to all limits of the resulting frame. To prove nontriviality
in each component, we first transfer the reflected test estimate to
all periodic Gibbs states and then pass to the conditioned limits.

The letter $P_0$ denotes the projection on $\ker H$. We also write $U$ for the unitary representation $v\mapsto U(v)$ of $\mathbb R^2$ on $\mathcal H(\sigma')$ with $U(v)[F]=[\tau_{(v,0)}F]$, which is the $\mathcal U$ of \textup{(Rec2)} and of \Cref{support:27-22}.

\begin{lemma}[Step 1: the class, the moment bound and the bound system]\label{Y-P5-L1} Let $\mathsf b'\ge x_s$, which in this step we write $\mathsf b$; let $K\ge1$, and let $\rho,\nu\in\mathbb G_K(\mathsf b)$. Every number of this step, that is, $j^\circ$, $G_1$, $\check\Gamma$, $B_0(\cdot,\cdot)$, $c_{\rm vol}$ and $A_4$, is evaluated at the single label $\mathsf b$ itself, not as a supremum over $U(\cdot)$.

\begin{enumerate}\item[\textup{(a)}] \emph{One cell.} Let $g_1,\dots,g_r$ be real bounded functions carried by one cell of $B_0$, at sup distance $\ge\check\delta$ from one another and at distance $\ge\check\delta/2$ from the walls $\mathbb Z+\frac12$, each in a sup-ball of radius $\check\delta/4$, where $0<\check\delta\le1$ and $K\ge j^\circ(\mathsf b,\check\delta)$; let the $c_i$ be real. Then, for every $N$ of the data, \[\|\prod_i(\mathcal O_Kg_i-c_i)\|_{{\rm ch},N}\le e^{c_{\rm vol}}\prod_i\bigl(G_1\|g_i\|_\infty+|c_i-\hat c_{K,N}(g_i)|\bigr),\] where $G_1=G_1(\mathsf b,\check\delta,\check\delta/4)$.

\item[\textup{(b)}] \emph{Cell by cell.} Let finitely many pieces be placed so that, in each cell $t$, the $n_t$ pieces lying in $t$ satisfy the conditions of \textup{(a)} with a number $\check\delta_t$ in the place of $\check\delta$. Then \[|\rho(\prod_i(\mathcal O_Kg_i-\nu(\mathcal O_Kg_i)))|\le\prod_{t\ \rm occupied}\bigl[\check\Gamma(\mathsf b,\check\delta_t)^{n_t}e^{c_{\rm vol}}\bigr]\cdot\prod_i\|g_i\|_\infty.\]

\item[\textup{(c)}] The moment bound of \textup{(E0)}, the bound on $|\rho(X^\nu_K)|$, holds for every $n\ge1$, for real bounded $f_i$ of the Euclidean type $(d;R)$ and for all $K\ge1$.

\item[\textup{(d)}] The bound of \textup{(E0)} on $|S^T_n[\rho](X)|$ holds for $n\ge2$, and so do \textup{(Cls1)}, \textup{(Cls2)}, \textup{(Bd1)} and \textup{(Bd2)}.\end{enumerate}\end{lemma}


\begin{proof}
Clause (a) is \Cref{Y-P4-Th4}(4.1)(a) at
$\check\rho=\check\delta/4$. The fixed-label estimate in the proof of
(4.3)(a) gives (b) before the supremum over a label neighborhood is
taken. Its application to each occupied cell uses the same sequence
of sides for the integrating law and the centering law. This gives
(c) by (4.3)(b), which also gives the cumulant bound in (d). The
one-point bound (Bd2) is (c) at $n=1$, and (Bd1) follows from the
cumulant estimate since $\|f\|_\infty\le\|f\|_\sharp$.
\Cref{support:54-17}(a), (b) gives (Cls1). For (Cls2),
\Cref{support:54-18}(a) gives the Gibbs equations after
invariant conditioning, and the conditioned law remains invariant
under $a_K\mathbb Z^3$.
\end{proof}

\begin{proof}[Proof of \Cref{Y-P5-ThE}] \emph{Step 3: from \textup{(E0)} to \textup{(E1)}, \textup{(E2)}, \textup{(E3)} and \textup{(E4)}(a).} In this step we derive \textup{(E1)} to \textup{(E4)}(a) from three inputs: \textup{(E0)}; \Cref{Y-P4-Th4}\,(4.6), (4.7); and the statements proved in the supporting sections that item~B of \S\ref{Y-P5-1} cites.

\textup{(1)} \emph{The data of \Cref{support:77-1}.} Once the $j$ whose labels lie outside $U(\mathsf b)$ are dropped, $(K_j,\mathbb G_{K_j}(\mathsf b_j))$ is a class system obeying $C^{\mathsf b}$, and $(K_j,\nu_j)$ is a frame. We have $\sigma_{\nu_j}(X)=\sum_{\mathrm{Part}_{\ge2}}\prod_JS^T_{|J|}[\nu_j](X_J)$; conversely, the cumulants are fixed polynomials in the moments (\eqref{support:9-13}). Hence $S_n=\lim S^T_n[\nu_j]$ exists on all tuples if and only if the frame converges, and then $\sigma_{\rm frame}=\sigma_S$. Converging subsequences exist by \Cref{support:35-39}, and a decomposition exists by \Cref{support:35-36}.

\emph{\textup{(2)} The hypotheses of \Cref{support:77-15}.} We substitute the data of Step~3\,(1) for the letters of \Cref{support:77-15}; then the hypotheses of that theorem are met. Its clauses \textup{(1)(a)} to \textup{(e)} give the assertion of \textup{(E1)} that $\sigma_S$ has the properties of these clauses, while the other three assertions of \textup{(E1)} on $\sigma_S$ come from Step~3\,(3), which gives the invariance under $\mathbb R^3\rtimes\operatorname{O}(3)$ and the reflection positivity at every mirror, from Step~3\,(4), which gives the growth condition \eqref{support:26-34} with the quadruple stated there, and from Step~4 in the case $c=0$, which gives the data \textup{(Rec1)} to \textup{(Rec4)} with $(C_\wp,5,\frac12)$. Its clauses \textup{(2)(a)} and \textup{(b)}, which are the content of \Cref{support:36-9}\,(a) and of \Cref{support:36-25} by way of \Cref{support:36-27}, give \textup{(E2)} except for the assertion that every isometry acts on $\mathsf X$ as the identity; this assertion is given by clause \textup{(2)(c)} with $h_r=\mathrm{id}$, which the condition \textup{(Rot)} of Step~3\,(3) yields. Its clauses \textup{(3)(i)}, \textup{(ii)(a)}, the first half of \textup{(ii)(b)}, \textup{(iii)} and \textup{(iv)} give \textup{(E3)} for the components $\sigma'=\sigma_\xi$, apart from its part on rotations. For a general limit $\sigma'$ in the frame, \Cref{support:36-12}\,(a), (b), (c), (e) shows that $\sigma'$ is symmetric, invariant under translations and $W_3$, and reflection positive at the coordinate and the diagonal mirrors, that $\sigma'\in\mathsf{Fct}$, and that $\sigma'$ is regular of order $\le4n+1$; \Cref{support:36-13}\,(a) gives $S_1=c_{\sigma'}\int f$. The datum \textup{(Rec2)} of \Cref{Y-P5-hs4}\,(b2) shows that $\mathcal H(\sigma')$ is separable and that $H\ge0$, since, by Step~3\,(3), every limit in the frame is invariant under $\operatorname{O}(3)$, and this invariance is the hypothesis of clause \textup{(b2)}. As $\sigma'\in\mathsf{Fct}$ by \Cref{support:36-12}\,(e), and as Step~3\,(3) gives the invariance of $\sigma'$ under $\mathbb R^3\rtimes\operatorname{O}(3)$ and its reflection positivity at $(e_3,0)$, \Cref{support:28-34} applies to $\sigma'$ itself: it gives the representation of $\operatorname{O}(2)$ of \Cref{support:28-30}, and it shows that the joint spectrum lies in the cone $E\ge|q|$, which is the datum \textup{(Rec3)} of \Cref{Y-P5-hs4}\,(b2).

\emph{\textup{(3)} The condition \textup{(Rot)} for every limit in the frame.} Let $\sigma'$ be a limit in the frame, with approximants $\rho_j\in\mathbb G_{K_j}(\mathsf b_j)$. \Cref{support:36-12}\,(d) gives $S_n^{\sigma'}=\lim_jS^T_n[\rho_j]$ for $n\ge2$, and $\sigma'^{\rm c}=\lim\sigma_{\rho_j}$. By the definition of $\mathrm{Lim}_G(\mathsf b)$, the sequence $(S_n^{\sigma'})_{n\ge2}$ is therefore a member of $\mathrm{Lim}_G(\mathsf b)$, and its centred functional is $\sigma'^{\rm c}$. \Cref{Y-P4-Th4}\,(4.7) then shows that $\sigma'^{\rm c}$ is invariant under $\mathbb R^3\rtimes \operatorname{O}(3)$ and reflection positive at every mirror. Next, \Cref{support:36-13}\,(a) gives $\sigma'((f))=c\int f$, and from \eqref{support:9-21} we get $\sigma'=\sigma'^{\rm c}\circ\mathrm{sh}_c$. The map $\mathrm{sh}_c$ commutes with every isometry $r$, because $\int f\circ r=\int f$; it commutes with conjugation, because $c$ is real; and it commutes with concatenation. Hence $\sigma'(rX)=\sigma'(X)$, and for every mirror $m$ \[\sigma'(\theta_m\bar F\cdot F)=\sigma'^{\rm c}(\theta_m\overline{\mathrm{sh}_cF}\cdot\mathrm{sh}_cF)\ge0.\] This is the condition \textup{(Rot)} of \Cref{support:77-2}. It yields the clauses \textup{(1)(g)} and \textup{(2)(c)} of \Cref{support:77-15}, with $h_r=\mathrm{id}$, and those halves of \textup{(3)(i)} and \textup{(iii)} that are conditional on \textup{(Rot)}: the circular cone, by \Cref{support:28-34}, and the group $\operatorname{O}(2)$, by \Cref{support:28-30}.

\emph{\textup{(4)} Growth.} By \textup{(E0)}, $\rho(X^{\nu_j})\in[-B,B]$ for all $\rho\in\mathbb G_{K_j}(\mathsf b_j)$, where the interval is closed and $B$ is the bound of \textup{(E0)}. \Cref{support:36-12}\,(f), applied with $k=1$, therefore gives $|\sigma'(X)|\le B$. For $R\le1$ this becomes \[|\sigma'(X)|\le(8A_6)^n(n!)^5\max\{1,d^{-5n}\}\prod\|f_i\|_\sharp.\] We now apply \Cref{support:26-24} to $\mathsf T=\sigma'$, with $\mathsf a=5n$ and $\mathsf c_{\mathsf T}=(8A_6)^n(n!)^5$. It yields the complex linear functional $\mathsf S_n^{\sigma'}$ of \Cref{support:26-33} on $\mathcal S_{\ne}(\mathbb R^{3n})$, which satisfies $\mathsf S_n^{\sigma'}(\bigotimes X)=\sigma'(X)$ on tensors, together with the bound \[|\mathsf S_n^{\sigma'}(\varphi)|\le\mathsf C_{\ne}^{\,6n}M^n(n!)^{13}(8A_6)^n|\varphi|_{13n-1},\] the case $n=1$ being included. This is the bound \eqref{support:26-34} with the quadruple $(13,1,C_E,13)$.

\emph{\textup{(5)} The axioms.} We verify the hypotheses of \Cref{support:28-48} for the components $\sigma_\xi$. Membership in $\mathsf{Fct}(C,\bar c)$ is supplied by \Cref{support:36-12}\,(e), the growth condition by Step~3\,(4), the invariance by Step~3\,(3), reflection positivity at $(e_3,0)$ by \Cref{support:36-12}\,(c), and clustering by \Cref{support:36-17}.

\emph{\textup{(6)} The constant $\bar c_*$.} For $n=1$ the bound of Step~3\,(4) reads $|c_{\sigma'}|\,|\int f|\le A_6(1+R)^3\|f\|_\infty$. We take $f\in C^1_c$ with $0\le f\le1$ which equals $1$ on a ball of radius $R-1$ and is carried by the ball of radius $R$, and we let $R\to\infty$. Since $(1+R)^3/(\frac{4\pi}3(R-1)^3)\to\frac3{4\pi}$, this yields $|c_{\sigma'}|\le\frac3{4\pi}A_6$.

\emph{Step 4: the clauses \textup{(E4)(b1)} and \textup{(b2)}.} They make up the following proposition.

\begin{proposition}[Point bounds and reconstruction data]\label{Y-P5-hs4} \begin{enumerate}\item[\textup{(b1)}] The point bound \eqref{support:77-5} holds in the form \[\FitDisplay{\dimexpr\linewidth-1pt\relax}{|\sigma'^{\rm c}(\theta\bar F\cdot(h)\cdot F')|\le32C_{\rm pt}\,\eta^{-5}\|h\|_{L^1}\,\|[\tau_{-\eta}F]\|\,\|[\tau_{-\eta}F']\|}\] for every limit $\sigma'$ in the frame, every $\eta\in(0,1]$, all $F,F'\in\mathbb C\mathfrak S(\{x_3>\eta\})$ and every $h\in C^1_c(\{|x_3|<\eta/2\})$.

\item[\textup{(b2)}] Every limit $\sigma'$ in the frame that is invariant under $\operatorname{O}(3)$ carries the data \textup{(Rec1)} to \textup{(Rec4)} with $(C_\wp+|c_{\sigma'}|,5,\frac12)$, and the components $\sigma_\xi$ satisfy \textup{(Rec5)} as well.\end{enumerate}\end{proposition}

\begin{proof}[Proof of \Cref{Y-P5-hs4}] \emph{The clause \textup{(b1)}.} By Step~3\,(3), $\sigma'^{\rm c}$ is the centred functional of a member of $\mathrm{Lim}_G(\mathsf b)$. \Cref{Y-P4-Th4}\,(4.6) therefore gives it the properties \textup{(R1)} to \textup{(R4)}, and \textup{(R5)} with the triple $(c_{\rm I},C_{\rm I},\mathsf p_{\rm I})$ of item~2 of \S\ref{Y-P4-1}, where $C_{\rm I}$ is the supremum of the whole plain-label expression; these numbers do not depend on $\sigma'$. They are exactly the hypotheses of \Cref{support:34-42}, with $\wp=\mathsf p_{\rm I}+4=5$ and with $C_{\rm pt}$ as in \eqref{support:34-6}. Clause \textup{(b)} of that corollary, applied with $d=\eta$, gives $\sigma'^{\rm c}(\theta\bar F\cdot(h)\cdot F')=\int h(x)\mathsf v_\eta(x)\,dx$ with a function $\mathsf v_\eta$ that obeys, by \eqref{support:34-46}, \[|\mathsf v_\eta(x)|\le C_{\rm pt}\min\{1,\eta-|x_3|\}^{-5}\|[\tau_{-\eta}F]\|\,\|[\tau_{-\eta}F']\|.\] On the support of $h$ we have $\eta-|x_3|>\eta/2$; as also $\eta\le1$, the factor with the minimum is $\le2^5\eta^{-5}$ there, and \textup{(b1)} follows.


For \textup{(b2)}, Euclidean invariance of the frame limits was proved
in Step~3(3). Use (b1) as the point-bound hypothesis of
\Cref{support:77-12}. Its scalar shift gives
$C_\wp+|c_{\sigma'}|$, and the one-dimensional zero-energy space for
a component gives \textup{(Rec5)}.
\end{proof}

\begin{lemma}[One Hilbert space for a limit and its centred functional; the point bound carries over]\label{Y-P5-LShift} Let $\sigma'$ be a limit in the frame, and let $c=c_{\sigma'}$, so that $\sigma'=\sigma'^{\rm c}\circ\mathrm{sh}_c$. Then \[\langle[F],[F']\rangle_{\sigma'}=\sigma'^{\rm c}(\theta\overline{\mathrm{sh}_cF}\cdot\mathrm{sh}_cF')=\langle[\mathrm{sh}_cF],[\mathrm{sh}_cF']\rangle_{\sigma'^{\rm c}}.\] The map $[F]\mapsto[\mathrm{sh}_cF]$ is isometric and onto, it intertwines $e^{-tH}$ and $U(v)$, and it fixes $\Omega$. Moreover, $\sigma'$ satisfies \textup{(Rec4)} with the constant $C_\wp+|c|$, and $|c|\le\bar c_*$.\end{lemma}


\begin{proof}
\Cref{support:28-12} gives the unitary and its intertwining
properties, with $\sigma=\sigma'^{\rm c}$ and shift $c$.
\Cref{Y-P5-hs4}(b1) supplies the centered bound required
in \Cref{support:75-50}. That lemma gives \textup{(Rec4)} with
constant $C_\wp+|c|$. The one-point estimate in Step~3(6) bounds
$|c|$ by $\bar c_*$. For $\sigma_S$ the constant $c$ is zero.
\end{proof}


\emph{The combined constant.} We work in normal form, $\mathsf b>x_s$. The convention of \textup{(S8)} takes the supremum of the whole plain-label expression $C_\wp+\bar c_*$, so we keep these two terms at one label. Let $\sigma'$ be a limit in the frame, represented by $\rho_j\in\mathbb G_{K_j}(\mathsf b_j)$ with the original centring laws $\nu_j$. By local boundedness we may pass to a subsequence, without changing $\sigma'$, on which the plain values satisfy
\[
 A_6(\mathsf b_j)\longrightarrow A,
 \qquad C_{\rm I}(\mathsf b_j)\longrightarrow I.
\]
The fixed-label lattice estimate in Step~5 of the proof of \Cref{Y-P4-Th4} has coefficient $C_{\rm I}(\mathsf b_j)\eta^{-1}$. Its limiting argument therefore gives \textup{(R5)} for $\sigma'^{\rm c}$ with $I+\epsilon$ in place of $C_{\rm I}$, for every $\epsilon>0$. Applying \Cref{support:34-42} and letting $\epsilon\downarrow0$ yields the centered point bound with coefficient
\[
 P:=32\cdot320\cdot8^4\mathsf{mv}_3c_{\rm I}^{-4}MI
   =\lim_j C_\wp(\mathsf b_j).
\]
The plain-label estimate of \Cref{Y-P4-Th4}\,\textup{(4.3)(b)}, used with $n=1$ and the centring law $\nu_j$ and passed to the same subsequence, gives
\[
 |c_{\sigma'}|\,\left|\int f\right|
       \le A(1+R)^3\|f\|_\infty.
\]
The ball tests of Step~3(6) imply $|c_{\sigma'}|\le3A/(4\pi)$. The shift lemma \ref{support:75-50} now gives \textup{(Rec4)} for $\sigma'$ with coefficient at most
\[
 P+\frac{3A}{4\pi}
 =\lim_j\bigl(C_\wp(\mathsf b_j)+\bar c_*(\mathsf b_j)\bigr)
 \le\sup_{\mathsf b'\in U(\mathsf b)}
       \bigl(C_\wp(\mathsf b')+\bar c_*(\mathsf b')\bigr)
 =C_*(\mathsf b).
\]
All quantities inside this supremum are plain values. The same argument applies after the fixed dilation in \textup{(E5)}, using its modified plain coefficients before taking their sum.

\emph{Point kernels and reconstruction.}
The preceding argument supplies the insertion bound for the
component functionals with constant at most $C_*$. Apply the ordered
kernel construction \eqref{support:75-66}--\eqref{support:75-76} and the
extension \eqref{support:75-79} to each component. If $n\ge2$, choose
a direction with successive projection gaps at least $d(x)/n^2$ by
\eqref{support:75-53}. The bound \eqref{support:75-85} with
$\eta=\min\{1,d(x)/(2n^2)\}$ is the estimate in \textup{(E4)(b3)}.
For $n=1$ use $\mathbf W_1=\langle\Omega,B(1,1)\Omega\rangle$.

\Cref{support:76-59} gives \textup{(E4)(b4)}, including locality
and tempered matrix elements. The complex time variable in that
statement is the third Euclidean coordinate. The norm estimate
\eqref{support:75-72} at $\eta=\varepsilon_T$ gives the claimed bound on
$\mathsf W_n$. The same construction applies to $\sigma_S$ with its
vacuum space $L^2(\mathsf X,\lambda)$. The one-dimensional vacuum
conclusion uses \textup{(Rec5)} for a component.

\begin{proposition}[Every label $\mathsf b>0$ and labels $\mathsf b_j\to \mathsf b$; joint limits at every $\mathsf b$]\label{Y-P5-hs8} The clauses \textup{(E0)} to \textup{(E4)} hold at labels $\mathsf b\le x_s$, with the constants of \textup{(E5)} and with $z=z(\mathsf b)$. In particular \textup{(R5)} holds for every $\eta\in(0,1]$, with $(c_{\rm I}L^{-z}/2,\ 2L^zC_{\rm I}(L^z\mathsf b),\ 1)$. The clause \textup{(E6)} holds at every $\mathsf b>0$.\end{proposition}

\begin{proof} \emph{\textup{(1)} Dilation.} The map $x\mapsto L^zx$ carries $a_K\mathbb Z^3$ to $a_{K-z}\mathbb Z^3$, and its pushforward is a bijection from $\mathbb G_K(\mathsf b')$ to $\mathbb G_{K-z}(L^z\mathsf b')$, since $\beta=L^{K-z}(L^z\mathsf b')$. Moreover $\mathcal O_K(f)=L^z\mathcal O_{K-z}(f^{[z]})$ with $f^{[z]}(y)=f(L^{-z}y)$. The type $(d;R)$ of a tuple becomes $(L^zd;L^zR)$, and $f^{[z]}$ has the same sup norm as $f$. Hence \textup{(E0)} holds for $K\ge z+1$ with $L^zA_6(L^z\mathsf b)(1+L^zR)^3$ for each entry, and $\max\{1,(L^zd)^{-5n}\}\le\max\{1,d^{-5n}\}$. The one number $z=z(\mathsf b)$ suffices for all labels of $U(\mathsf b)$, because $L^zU(\mathsf b)\subset[x_s,\infty)$; we discard the finitely many $j$ with $K_j\le z$ or $\mathsf b_j\notin U(\mathsf b)$. At $\mathsf b=x_s$ the definition of $z(\mathsf b)$ gives $z=1$ and $U(\mathsf b)=[x_s/2,2x_s]$, so that the same numbers also cover the sequences that tend to $x_s$ from below. With this bound system, the frame and point estimates of Steps~3 and~4, scalar reconstruction in Appendix~\ref{app:scalar-reconstruction}, and the argument of Step~9 apply at label $\mathsf b$. Parts~(2) to~(5) below verify the corresponding constants and hypotheses.

\emph{\textup{(2)} The numbers.} For $R\le1$ we have $\mathsf c_{\mathsf T}=(L^z(1+L^z)^3A_6(L^z\mathsf b))^n(n!)^5$; with this $\mathsf c_{\mathsf T}$ in place of $(8A_6)^n(n!)^5$, \Cref{support:26-24} gives, as in Step~3\,(4), the number that takes the place of $C_E$. From $|c|\,|\int f|\le L^zA_6(L^z\mathsf b)(1+L^zR)^3\|f\|_\infty$, with the functions $f$ of Step~3\,(6) and $R\to\infty$, we obtain $\bar c_*(\mathsf b)=\frac3{4\pi}L^{4z}A_6(L^z\mathsf b)$.

\emph{\textup{(3)} The properties \textup{(R1)} to \textup{(R4)}.} For $\sigma'^{\rm c}$ the properties \textup{(R1)} to \textup{(R4)} do not depend on the label. The properties \textup{(R1)}, \textup{(R2)}, \textup{(R3)} come from \Cref{support:36-12}\,(a), (b), (c), (e); the property \textup{(R4)} comes from \Cref{support:28-21}, by way of the diagonal mirrors; and the invariance under $\operatorname{O}(3)$ comes from \Cref{Y-P4-Th4}\,(4.8).

\emph{\textup{(4)} The property \textup{(R5)}.}
Put $D_zX=L^{z|X|}X^{[z]}$. The centered functionals satisfy
$\sigma'^{\rm c}=\sigma''^{\rm c}\circ D_z$, where $\sigma''$
is the corresponding frame limit at $(K_j-z,L^z\mathsf b_j)$.
By \textup{(E3)(i)} in normal form, $\sigma''^{\rm c}$ is the
centered functional of a member of $\mathrm{Lim}_G(L^z\mathsf b)$.
It therefore satisfies the hypotheses of \Cref{Y-P4-LZ}, by
\Cref{Y-P4-Th4}\,\textup{(4.6)}. That lemma gives \textup{(R5)}
for every $\eta\in(0,1]$, with
$(c_{\rm I}L^{-z}/2,\,2L^zC_{\rm I}(L^z\mathsf b),\,1)$.

\emph{\textup{(5)} The conditions \textup{(PtB)} and \textup{(Rec)}.} They follow by the proof of \Cref{Y-P5-hs4} in Step~4, with these changes: \textup{(R1)} to \textup{(R4)} are taken from Step~8\,(3), and \textup{(R5)} from Step~8\,(4); where that proof applies \Cref{support:34-42}, $C_{\rm pt}$ is taken at $(c_{\rm I}L^{-z}/2,\,2L^zC_{\rm I}(L^z\mathsf b))$, which is $2^4L^{4z}\cdot2L^z=32L^{5z}$ times $C_{\rm pt}(L^z\mathsf b)$; and in \Cref{Y-P5-LShift} the bound on $|c|$ is the one of Step~8\,(2).

\emph{\textup{(6)} The clause \textup{(E6)}.} This clause is \Cref{Y-P4-Th4}\,(4.9)\,(i).\end{proof}

\emph{Step 9: the exact relation between $S_2$ and the components.} In this step we prove the clauses \textup{(E7)(a)}, \textup{(b)} and \textup{(c)} of \Cref{Y-P5-ThE}.

\emph{\textup{(a)} The clause \textup{(E7)(a)}.} By \eqref{eq:component-plateau-moments}, applied with $q=1$, $\sigma_S(\theta f,f)=\int\sigma_\xi(\theta f,f)\,d\lambda$; moreover $\sigma_S(\theta f,f)=S_2(\theta f,f)$, because $S_1[\sigma_S]=0$. As $\int\theta f=\int f$, we have $\sigma_\xi(\theta f,f)=S_2^\xi(\theta f,f)+c_\xi^2(\int f)^2$. In $\mathcal H(\sigma_\xi)$ we have $g_f(\xi)=\|[(f)]-c_\xi\int f\,\Omega\|^2\ge0$, and $g_f$ is continuous in $\xi$ by \Cref{support:35-34}. \Cref{support:36-20}\,(a) and \eqref{eq:component-plateau-moments}, applied with $q=2$, give \[\|P_0[(f)]\|^2=\lim_t\sigma_S(\theta f\cdot\tau_tf)=\int\sigma_\xi((\theta f))\sigma_\xi((f))\,d\lambda=V(\int f)^2.\] Subtracting, we obtain $\int g_f\,d\lambda=\|(1-P_0)[(f)]\|^2$. Finally $V\le\bar c_*^2$ by \textup{(E3)(ii)}.

\emph{\textup{(b)} The clause \textup{(E7)(b)}.} By \textup{(b3)}, $\mathsf S_2^\xi(x,y)-c_\xi^2=s(|x-y|)$.
Write $B_{\mathsf e}=B(\mathsf e,\mathsf e)$ as in
\eqref{support:75-34}. It is self-adjoint, and
$c_\xi=\langle\Omega,B_{\mathsf e}\Omega\rangle$ by the one-point case
of \eqref{support:75-66}. We place two points at the distance $t$ along $e_3$ and put $\psi_{\mathsf e}:=B_{\mathsf e}\Omega-c_\xi\Omega$, a vector orthogonal to $\Omega$. Then \[s(t)=\langle\psi_{\mathsf e},e^{-(t-2\mathsf e)H}\psi_{\mathsf e}\rangle=\int e^{-(t-2\mathsf e)E}\,d\mu_{\mathsf e}(E)\] for $0<\mathsf e\le t/2$, with $\mu_{\mathsf e}\ge0$. Hence $s\ge0$, and $s$ is non-increasing. If $s(t_0)=0$, then $e^{-(t_0-2\mathsf e)H/2}\psi_{\mathsf e}=0$, and so $\psi_{\mathsf e}=0$, since $e^{-sH}$ is injective. As $\psi_{\mathsf e'}=e^{-(\mathsf e'-\mathsf e)H}\psi_{\mathsf e}$ for $\mathsf e'\ge\mathsf e$, it follows that $\psi_{\mathsf e}=0$ for every $\mathsf e$. Suppose now that $g_f(\xi)=\int\!\!\int f(\theta x)f(y)s(|x-y|)\,dx\,dy=0$ for a continuous $f\ge0$ with $f\ne0$. Then $s$ vanishes somewhere, and hence $B_{\mathsf e}\Omega=c_\xi\Omega$ for all $\mathsf e$. In \eqref{support:75-68} the last insertion then gives
$c_\xi\Omega$. The translations fix $\Omega$, so induction gives
$\Psi_n(x)=c_\xi^n\Omega$. The ordered tuple classes are total by the
approximation in \S\ref{sec:ordered-density}, and therefore
$\dim\mathcal H=1$; moreover $\mathsf S_n=c_\xi^n$, which is the product formula. Conversely, if $\dim\mathcal H=1$, then $[(f)]=\sigma_\xi((f))\Omega$, so that $g_f=0$ for all $f$, and polarisation together with translations and $\operatorname{O}(3)$ gives $S_2^\xi=0$.

\emph{\textup{(c)} The clause \textup{(E7)(c)}.} As an instance we take $\frac12(\sigma_0+\sigma_0\circ\mathrm{sh}_{c'})$, centred, where $\sigma_0$ is the functional equal to $1$ at the empty tuple and to $0$ elsewhere. It has $S_2=(c'^2/4)\int f\int f'$, two trivial components and $\dim\mathcal H(S)=2$, by \Cref{ex:constant-mixture}. Hence \Cref{Y-P6-N}\,(N5) by itself decides nothing about the components; non-triviality of every component follows from \Cref{Y-P5-ThC}.\end{proof}

\begin{proof}[Proof of \Cref{Y-P5-ThP}] \emph{Step 10: domination at one side; cell functions of the pair $(U,V^{(k)})$.} The first of the three steps towards \Cref{Y-P5-ThP} consists of two lemmas: \Cref{Y-P5-LD} gives domination at one side for a cell function that is allowed to depend on the side, and \Cref{Y-P5-LK} treats the cell functions of the pair $(U,V^{(k)})$, with the three bounds of \S\ref{Y-P6-1} as seminorms.

\begin{lemma}[Domination at one side for a cell function allowed to depend on the side]\label{Y-P5-LD} Let $\rho\in\mathbb G_K(\mathsf b')$, and let $B_i$ be a board of $\mathbb L_K$ with cells of side $l_c=a_i$. Let $N\in\{2L^m\}$ with $N\ge2l_c$, and let $F\in\mathcal A(\square_0)$ with $\sup|F|\le\mathsf m$. Put $\epsilon_N:=36\beta(N/a_K)^2\lfloor N/2l_c\rfloor^{-3}$. If $\epsilon_N\le\frac14$, then $|\rho(F)|\le\|F\|_{{\rm ch},N}+2\mathsf m\,\epsilon_N^{1/2}$. Moreover $\epsilon_N\le2304\,\beta a_K^{-2}l_c^3/N\to0$ as $N\to\infty$. As a consequence, if $F_N\in\mathcal A(\square_0)$, $\sup_N\sup|F_N|<\infty$ and $\|F_N\|_{{\rm ch},N}\le A$ for all large $N$, then $\limsup_N|\rho(F_N)|\le A$.\end{lemma}

\begin{lemma}[Cell functions of the pair $(U,V^{(k)})$, and three bounds as seminorms]\label{Y-P5-LK} Let $\Phi$ be a bounded real Borel function of the inner bonds of a cell and the drawn bonds of level $k$ whose two $k$-blocks both lie in that cell, where the walls of the board lie in $\mathbb Z+\frac12$ at level $k$.

\begin{enumerate}

\item[\textup{(a)}] $F:=E[\Phi\mid U]\in\mathcal A(\square)$, and $P_\rho(\Phi)=\rho(F)$.

\item[\textup{(b)}] $\|F\|_{{\rm ch},N}^Q=E_{{\rm torus}\,N}\prod_c\Phi\circ\check\Theta_c$, where $Q$ is the number of cells and $\check\theta_c$, $\check\Theta_c$ are the maps of the fold, those of the data of \Cref{Y-P6-L5} with $h$ replaced by the side of the cells of the board, counted in units of level $k$.

\item[\textup{(c)}] For every $N$ of the data, that is, every $N=2L^{m_l}$ with $m_l\ge m'$ and with \textup{(Tor)} holding at $\mathsf b'$, for all $K\ge j+1$ and all $j\ge j_*(\mathsf b')$, and on boards without offset, the three bounds below hold.

\begin{enumerate}

\item[\textup{(P$^\circ$)}] (\Cref{Y-P6-L3}) On $B_{j-1}$ and for $p$ based at $0$, $\|E[1\{\phi_p\ge\varepsilon/2\}\mid U]\|_{{\rm ch},N}\le e^{-(c_P/2)p_0^2}$.

\item[\textup{(W$^\circ$)}] (\Cref{Y-P6-L8}) On $B_{j-i_h}$ and for every constant $g\in[0,\frac98p_0^2]$, \[\|E[(G_--g)(\mathcal O_Kf_*-c_f^N-A_f^{{\rm tr},N}(V))\mid U]\|_{{\rm ch},N}\le\frac98p_0^2\cdot2e^{2.2}/r_z.\]

\item[\textup{(X$^\circ$)}] (\Cref{Y-P6-L9}) On the same board and at $w=w_R$, \[\|E[\Psi\circ\tau_{-t}\,e^{\Lambda^{\rm tr}\circ\tau_{-t}-\Lambda^{\rm tr}}-\Psi\mid U]\|_{{\rm ch},N}\le\sup|\Psi^\circ|\cdot2e^{2.2}/(r_\eta-1).\] Here $\Psi$ is the function of Step~11, that is, the $\Psi^{\rm tr}$ of \Cref{Y-P6-L9}, and $\Psi^\circ:=(\Phi e^{\Lambda^a})\circ\tau_t\,e^{-\Lambda^a}$, with $t=x^{-2}$ the number $t$ of \S\ref{Y-P6-1}, is the $\Psi$ of \Cref{Y-P6-L9}. In \textup{(X$^\circ$)}, $\Phi$ is not the function of the opening sentence of this lemma but, as in the data of \Cref{Y-P6-L9}, a bounded Borel function of the bonds of $B'$, carried by $\{\phi_p\le3\varepsilon/4\text{ on }B'\}$.

\end{enumerate}

\end{enumerate}
\end{lemma}

\begin{proof}[Proof of \Cref{Y-P5-LD}]

\begin{enumerate}

\item[\textup{1.}] Since $\rho$ is invariant under $a_K\mathbb Z^3$, the number $P=1$ is a period number in the sense of \Cref{support:14-9}\,\textup{(3)}.

\item[\textup{2.}] As data of \Cref{support:18-32} we take $\mathcal K_0=\{0\}$, $F_0=F$, $\mathsf q=2$, which is even and a multiple of $P$, and $\mathbf k_0=0$. With these data the letter $\mathsf P$ of \Cref{support:18-32} equals $2l_c$, and the number $\epsilon_\ell$ of \eqref{support:18-33} is $\epsilon_N$.

\item[\textup{3.}] \Cref{support:18-43}, applied at this one side, is the claim. In that corollary $\mathsf m$ is $\sup|F|$, and a larger choice of $\mathsf m$ would only make the bound weaker.

\item[\textup{4.}] Because $\lfloor y\rfloor\ge y/2$ for $y\ge1$, we have $\lfloor N/2l_c\rfloor^{-3}\le(4l_c/N)^3$, and hence $\epsilon_N\le36\cdot64\,\beta a_K^{-2}l_c^3/N$.

\item[\textup{5.}] \emph{A cell function that may depend on the side.} \Cref{support:18-38}\,\textup{(a)} and \Cref{support:18-43} assume that the data of a domination are given and that $\ell$ belongs to the set of sides with $\ell\ge\mathsf P$; the corollary assumes in addition that $\epsilon_\ell\le\frac14$. For each $N_0$ we give the data with $F=F_{N_0}$, and we use the conclusion at $\ell=N_0$ alone. No hypothesis asks that $F$ be the same at two sides.

\end{enumerate}
\end{proof}

\begin{proof}[Proof of \Cref{Y-P5-LK}]

\begin{enumerate}

\item[\textup{(a)}] This clause follows as well from \Cref{Y-P6-L1}\,\textup{(L)} and \textup{(I)}; we give the argument. By item~2 of \S\ref{Y-J-1}, which describes the step, the average $\mathsf M(U)(c)$ is determined by $\mathsf t_U$ on $B(y)$ and on $B(y')$ together with the straight paths $[x,x+Le_\mu]$, all of which lie in $B(y)\cup B(y')$. This average is a Borel function of $U$, and $k_s$ is a probability kernel. We follow the dependence level by level, from level $k$ to level $0$: a drawn bond of level $k$ depends on the drawn bonds of level $k-1$ in its two blocks, and so in the end on the bare bonds in the union of its two centred $k$-blocks. Since $L$ is odd, every plane in $\mathbb Z+\frac12$ of level $k$ is a face of the centred blocks of each finer level, and lies in $a_K(\mathbb Z+\frac12)$.

Therefore $F$ is a bounded Borel function of the inner bonds. The law $P_\rho$ is $\rho$ times a product kernel, and $\Phi$ depends on finitely many of its draws; it follows that $P_\rho(\Phi)=\rho(F)$. The formula of the kernel is local and does not depend on the side, so $F\circ\mathsf{per}_N$ is the function that the same formula gives on the torus.

\item[\textup{(b)}] The identity $\check\Theta_cE[\Phi\mid U]=E[\Phi\circ\check\Theta_c\mid U]$ is the third claim of \Cref{Y-P6-L1}\,\textup{(K)}. Its proof uses two pointwise facts, in which no torus enters: $\mathsf M(\check\Theta U)(c)=\mathsf M(U)(\check\theta\,c)$ and $k_s(g^{-1},h^{-1})=k_s(g,h)$. Here $\check\theta$ and $\check\Theta$ are the maps of \Cref{Y-P6-L1}\,\textup{(K)} for any choice, in each coordinate, of a translation $y_\mu\to y_\mu+n_\mu$ or a reflection $y_\mu\to n_\mu-y_\mu$. The draws used in different cells, at all levels $\le k$, form disjoint families, and by \Cref{Y-P6-L1}\,\textup{(I)} these families are independent given $U$. All functions here are real.

By \Cref{support:18-14}\,\textup{(2)}, $\|F\|_{{\rm ch},N}=\|F\circ\mathsf{per}_N\|_{\rm ch}$, the seminorm on the right being that of \eqref{support:4-14}, whose $Q$-th power is the mean of the product of the copies. By the identity and the independence above, this mean equals $E_{{\rm torus}\,N}\prod_c\Phi\circ\check\Theta_c$.

\item[\textup{(c)}] The three clauses come from \S\ref{Y-P6-1}. There Lems.~\ref{Y-P6-L3}, \ref{Y-P6-L8} and \ref{Y-P6-L9} are formulated as bounds on $\|E[\Phi\mid U]\|_{\rm ch}$: they give \textup{(P$^\circ$)}, \textup{(W$^\circ$)} for every constant $g\in[0,\frac98p_0^2]$, and \textup{(X$^\circ$)}. By \textup{(b)}, the fold $E\prod_c\Phi\circ\check\Theta_c$ is $\|F\|_{{\rm ch},N}^Q$. In \textup{(X$^\circ$)} the truncation is part of the definition of $\mathsf\Lambda$ and of $\Lambda^{\rm tr}$ in Step~11, and these two are the $\Lambda^{a,{\rm tr}}$ and $\Lambda^{\rm tr}$ of \S\ref{Y-P6-1}.

\end{enumerate}
\end{proof}

\emph{The boards.} The cells of the two boards have the sides $a_{j-1}$ and $a_{j-i_h}$, where $h=L^{i_h}$; both indices are $\le j\le K-1$. There is no offset: as in \S\ref{Y-P6-1}, the cell $0$ is centred at $0$. All sites of a plaquette $p$ based at $0$ lie in the cell of side $L$, and all sites of $\mathrm{Box}$ lie in the cell of side $h$, because $r_B\le(h-1)/2$. Hence every bond of $p$ or of $\mathrm{Box}$ is a drawn bond of level $k$ with both of its $k$-blocks in the respective cell, as in the opening sentence of \Cref{Y-P5-LK}. Every other place is reached by a translation in $L^{-j}\mathbb Z^3$, the set of sites of level $k$, which lies in $a_K\mathbb Z^3$. Since $\rho$ is invariant under these translations and the kernel is covariant under them by \Cref{Y-P6-L1}\,\textup{(K)}, the law $P_\rho$ is invariant under them. In each direction there are $2L^{m_l+j-1}$ and $2L^{m_l+j-i_h}\ge2L$ cells respectively, by entry \textup{(j3)} of the list of conditions on the depth in \S\ref{Y-P6-1}.

\emph{Sup bounds free of $N$, for \Cref{Y-P5-LD}.}
Keep $K$, $j$ and $\mathsf b'$ fixed. Only plaquettes whose
barycentres meet $\operatorname{supp}f_*$ contribute to
$\mathcal O_Kf_*$. This support lies in a cube of side
$(D_0/4)L^{-j}$. For each orientation its barycentres form a
translate of $a_K\mathbb Z^3$, so
\[
 \sup|\mathcal O_Kf_*|
 \le 6L^K\bigl((D_0/4)L^{K-j}+2\bigr)^3=:C_{K,j}.
\]
This number is independent of $N$. The termwise estimates of
\Cref{Y-P6-L2} give
\[
 \max\{\sup|A_f|,\sup|A_f^{{\rm tr},N}|\}
 \le (C_Sp_0^2+1)/\mathsf b'=:A_*.
\]
At $V=1$, the field is globally tame, so $W_f(1)=0$ by
Step~\textup{(b)} of the proof of \Cref{Y-P6-L5}.
Consequently \Cref{Y-P6-L2} gives
$|c_f^N|\le C_{K,j}+A_*$.
For every $g\in[0,9p_0^2/8]$, the cell function in
\textup{(W$^\circ$)} therefore has supremum at most
$(9/4)p_0^2(C_{K,j}+A_*)$.
The cell function in \textup{(X$^\circ$)} has supremum at most
$2e\sup|\Phi|$; for the four tests used below,
$\sup|\Phi|\le1.5C_S\varepsilon p_0^2$.
Conditional expectation preserves these supremum bounds.
Thus each application of \Cref{Y-P5-LD} has an error bounded by
a constant independent of $N$ times $\epsilon_N^{1/2}$.

\emph{Step 11: two lemmas at one side.} We fix a side $N$ of the data, that is, $N=2L^{m_l}$ with $m_l\ge m'$ and with \textup{(Tor)} holding at $\mathsf b'$, such that $\epsilon_N\le\frac14$ on $B_{j-i_h}$; both lemmas below are stated at this side. The rows, the members, $w_R$, $S_f$ and $\mathbb D_k$ are those of the torus of side $N$, and what follows holds however they depend on $N$. For the function $\Phi$ of \Cref{Y-P5-L10per} below we put $\delta_N:=4e\sup|\Phi|\,\epsilon_N^{1/2}$.

\begin{lemma}[Approximate integration by parts under every periodic Gibbs state]\label{Y-P5-L10per} Let $\rho\in\mathbb G_K(\mathsf b')$, and put $\mathsf\Lambda:=\Lambda^{a,{\rm tr}}$. Let $t=x^{-2}$ and $w=w_R$; here $w_R$ is a real function on the bonds of $B_R$ with $|w_R|\le1$. Suppose that the datum $t\sup|D_w^\alpha\mathsf\Lambda|\le\frac12$ holds on $\{\phi_p\le3\varepsilon/4+4t\text{ on }B'\}$. Let $\Phi$ be as in the data of \Cref{Y-P6-L10}, that is, of class $C^2$ along the flow, a function of $B'$, and carried by $\{\phi_p\le3\varepsilon/4\text{ on }B'\}$, and on the support of $\Phi$ put $\varkappa:=D_w\mathsf\Lambda+xu^\alpha$. Then \[\FitDisplay{\dimexpr\linewidth-1pt\relax}{|P_\rho[D_w^\alpha\Phi-xu^\alpha\Phi]|\le tB_2+t^{-1}e\sup|\Phi|\,\bigl(2.5(2R/M+2r(x)+3)^3c_h(x)+6e^{2.2}x^{-4}\bigr)+\sup|\Phi|\sup|\varkappa|+\delta_N/t.}\]
\end{lemma}

At this $w$ the datum is satisfied. Indeed, Step~0 of the proof of \Cref{Y-P6-L13} proves $\sup|D\Lambda^a|\le2p_0x^{1/2}$ by bounding the terms one by one, so the same bound holds for $\mathsf\Lambda=\Lambda^{a,{\rm tr}}$, a truncated sum being a part of the sum, and therefore $t\sup|D_w\mathsf\Lambda|\le2p_0x^{-3/2}$. Moreover the data of \Cref{Y-P6-L10}, taken at the same $x$, include the inequality $2p_0x^{-3/2}<\frac12$ between the numbers $x$ and $p_0$ alone; so the datum holds for the truncated sum $\mathsf\Lambda$. In the claim, $B_2$ is defined with $\mathsf\Lambda$.

\begin{lemma}[Covariance and variance, centred at the means of the state]\label{Y-P5-L13per} Let $\rho\in\mathbb G_K(\mathsf b')$, and let $N$ be the side fixed in Step~11. Let $a_{u,N}=\langle u,C_Bu\rangle$ and $q_N=\langle \mathrm{d}w,\mathbb D_k[f_*]\mathrm{d}w\rangle$. Then there is a multiple $\Delta_N$ of $\epsilon_N^{1/2}$, with a factor that depends on $K$, $j$ and $\mathsf b'$ but not on $N$, such that the following hold.

\begin{enumerate}

\item[\textup{(i)}] $|\mathrm{Cov}_{P_\rho}(G,S_f)-\frac32q_N|\le$ \textup{(E2)} $+$ \textup{(E3)} $+$ \textup{(E4)} $+$ \textup{(E6)} $+\Delta_N$, where the four errors are those of \Cref{Y-P6-L13}.

\item[\textup{(ii)}] $|\mathsf b'\mathrm{Cov}_{P_\rho}(G_-,\mathcal O_Kf_*)-\frac32q_N|\le$ \textup{(E1)} $+\dots+$ \textup{(E6)} $+\Delta_N$, where the errors are those of \Cref{Y-P6-L13}.

\item[\textup{(iii)}] $|\mathrm{Var}_{P_\rho}(G_-)-\frac32a_{u,N}^2|\le$ \textup{(E7)} $+\Delta_N$, where the error is that of \Cref{Y-P6-L13}.

\end{enumerate}
\end{lemma}

By \Cref{Y-P6-L14}\,\textup{(B)}, $q_N\ge c_F$ for all rows of the runs, and by \Cref{Y-P6-L11}\,\textup{(c)}, $0\le a_{u,N}\le2$.

\begin{proof}[Proof of \Cref{Y-P5-L10per}] Throughout, all members are truncated at $\#X\le n_0$ and taken on $\mathrm{Box}$.

\begin{enumerate}

\item[\textup{1.}] With $\mathsf\Lambda$ as in the statement, $\Lambda^{\rm tr}=\mathsf\Lambda+H^{\rm tr}$, where $H^{\rm tr}$ is the truncated sum of the $E^h$. We define $\Psi:=(\Phi e^{\mathsf\Lambda})\circ\tau_t\,e^{-\mathsf\Lambda}$ and \[\Phi_{\rm cell}:=\Psi\circ\tau_{-t}\,e^{\Lambda^{\rm tr}\circ\tau_{-t}-\Lambda^{\rm tr}}-\Psi.\] The argument depends on $\mathsf\Lambda$ and $H^{\rm tr}$ only through the differences $\mathsf\Lambda\circ\tau-\mathsf\Lambda$ and $H^{\rm tr}\circ\tau-H^{\rm tr}$. These are finite sums over the terms that depend on a moved bond, so they are functions of $\mathrm{Box}$, as Step~\textup{(c)}, dependence, of the proof of \Cref{Y-P6-L5} shows.

\item[\textup{2.}] \emph{Pointwise algebra.} We have $\Psi\circ\tau_{-t}=\Phi e^{\mathsf\Lambda-\mathsf\Lambda\circ\tau_{-t}}$. Hence the first term of $\Phi_{\rm cell}$ is \[\Phi\exp\{\mathsf\Lambda-\mathsf\Lambda\circ\tau_{-t}+\mathsf\Lambda\circ\tau_{-t}+H^{\rm tr}\circ\tau_{-t}-\mathsf\Lambda-H^{\rm tr}\}=\Phi e^{dH},\] with $dH=H^{\rm tr}\circ\tau_{-t}-H^{\rm tr}$. The second term of $\Phi_{\rm cell}$ is $\Phi\circ\tau_t\,e^{\mathsf\Lambda\circ\tau_t-\mathsf\Lambda}$.

\item[\textup{3.}] \emph{Taylor's formula along the flow.}
Use the support argument in Step~2 of the proof of \Cref{Y-P6-L10}.
A segment meeting the support stays in the enlarged tame set, with
fixed indicators and smooth $\mathsf\Lambda$. On the other segments
all terms involving $\Phi$ and its derivatives vanish. Thus the
Taylor identity holds on every configuration: \[\Phi\circ\tau_t\,e^{\mathsf\Lambda\circ\tau_t-\mathsf\Lambda}=\Phi+t(D_w\Phi+\Phi D_w\mathsf\Lambda)+\mathrm{rem},\] where, by the datum, $|\mathrm{rem}|\le(t^2/2)B_2e^{t\sup|D_w\mathsf\Lambda|}\le t^2B_2$.

\item[\textup{4.}] Pointwise, $|dH|\le2\cdot1.2124(2R/M+2r(x)+3)^3c_h\le1$. The first inequality follows from Step~5 of the proof of \Cref{Y-P6-L10}, since a truncated sum is a part of the sum. For the second, entry \textup{(j1)} of the list of \S\ref{Y-P6-1} bounds the error \textup{(E6)} of \Cref{Y-P6-L13} by $c_F/12$, and this gives $2.5(2R/M+2r(x)+3)^3c_h\le\frac14$.

\item[\textup{5.}] It follows that, pointwise, \[t(D_w\Phi+\Phi D_w\mathsf\Lambda)=-\Phi_{\rm cell}+\Phi(e^{dH}-1)-\mathrm{rem}.\]

\item[\textup{6.}] Applying $P_\rho$, we use \Cref{Y-P5-LK}\,\textup{(a)}, then \Cref{Y-P5-LD} at the side $N$, then \textup{(X$^\circ$)}, and obtain \[|P_\rho(\Phi_{\rm cell})|\le\sup|\Psi^\circ|\,2e^{2.2}/(r_\eta-1)+\delta_N.\] Here $\sup|\Psi^\circ|\le e^{1/2}\sup|\Phi|$ by Step~6 of the proof of \Cref{Y-P6-L10}, and the datum of \Cref{Y-P6-L10} is met at $w_R$ because of Step~0 in the proof of \Cref{Y-P6-L13}; further, $\sup|\Psi|\le e^{1/2}\sup|\Phi|$ by the datum, and $\sup|\Phi_{\rm cell}|\le2e\sup|\Phi|$.

\item[\textup{7.}] On the support of $\Phi$ we have $D_w\mathsf\Lambda=-xu^\alpha+\varkappa$ with $|\varkappa|\le C_1x\varepsilon/R+C_2p_0^2R+C_a(x)p_0^2$, where $C_1$ and $C_2$ are the numbers of \Cref{Y-P6-L12}. This rests on \Cref{Y-P6-L12}, which is pointwise and deterministic, and on the Cauchy bound for the truncated sum of the $E^a$, given by the part ``Besides'' of \Cref{Y-P6-L10}, by Step~7 of its proof and by Step~0 of the proof of \Cref{Y-P6-L13}; the caps in that bound do not depend on $N$. The number $C_a(x)$ is the function of $x$ alone from \S\ref{Y-P6-1}, a fixed multiple of $(2r(x)+1)^3$.

\item[\textup{8.}] Items 5 and 7 give \[P_\rho[D_w^\alpha\Phi-xu^\alpha\Phi]=t^{-1}P_\rho[-\Phi_{\rm cell}+\Phi(e^{dH}-1)-\mathrm{rem}]-P_\rho[\Phi\varkappa].\] To bound the right side we use item 4 together with $2\cdot1.2124\le2.5$ and $|e^y-1|\le e|y|$ for $|y|\le1$, then item 3 in the form $|\mathrm{rem}|/t\le tB_2$, and item 6 with $r_\eta=x^4$ and $e^{1/2}\cdot2/(x^4-1)\le6e\,x^{-4}$. This gives the claim, an inequality at the one side $N$.

\end{enumerate}
\end{proof}

\begin{remark} On the torus, $E_\nu\Phi_{\rm cell}$ becomes an integral against $e^{\Lambda}[\Xi_w\circ\tau_t-\Xi_w]$ by changing variables in the density: the proof of \Cref{Y-P6-L10} reaches this integral in its Step~4, after the divergence, the regularity and Taylor's formula, and Step~1 of the proof of \Cref{Y-P6-L9} identifies it with $E_\nu[\Phi_{\rm cell}]\int e^\Lambda\Xi_w$. A periodic Gibbs state has no density of this kind. Under $\rho$ the argument does not depend on the change of variables in the density, because $\Phi_{\rm cell}$ is a bounded cell function given explicitly, whose fold is a quantity on the torus. No term $2\,\mathrm{ex}$ arises, since the truncation is part of the definition; consequently \textup{(E6)}, in the form in which \Cref{Y-P6-L13} states it, is the larger bound.
\end{remark}

\begin{proof}[Proof of \Cref{Y-P5-L13per}] This is the proof of \Cref{Y-P6-L13} with the changes in items 1 to 8 below.

\begin{enumerate}

\item[\textup{1.}] \emph{The split, without a conditional expectation.} Let $g:=P_\rho(G_-)\in[0,\frac98p_0^2]$. As $g$ is the mean, plain algebra gives \[\mathrm{Cov}_{P_\rho}(G_-,\mathcal O_Kf_*)=P_\rho[(G_--g)(\mathcal O_Kf_*-c_f^N-A_f^{{\rm tr},N}(V))]+\mathrm{Cov}_{P_\rho}(G_-,A_f^{{\rm tr},N}).\] By \textup{(W$^\circ$)}, \Cref{Y-P5-LK} and \Cref{Y-P5-LD}, the first term is at most $\frac98p_0^2\cdot2e^{2.2}/r_z$ plus the term of \Cref{Y-P5-LD}, which goes into $\Delta_N$. For the error \textup{(E5)} of \Cref{Y-P6-L13} we have $\frac98p_0^2\cdot2e^{2.2}/r_z\le$ \textup{(E5)}$/\mathsf b'$. Since single cubes are short, $A^{\rm tr}=S_f/\mathsf b'+\mathrm{Rem}^{\rm tr}$, and $|\mathrm{Cov}_{P_\rho}(G_-,\mathrm{Rem}^{\rm tr})|\le2\cdot\frac98p_0^2\sup|\mathrm{Rem}^{\rm tr}|\le$ \textup{(E1)}$/\mathsf b'$. For the second inequality we use that the bound of \Cref{Y-P6-L2} on $\mathsf b'\sup|\mathrm{Rem}|$ comes, in its proof, from estimating each term separately, with the sizes of the terms from Steps~4 and~5 of that proof, the number of anchors from Steps~6 and~7 of that proof and the constant $1.2124$ from Step~8 of that proof. As a truncated sum is a part of the sum, the same bound holds for $\mathsf b'\sup|\mathrm{Rem}^{\rm tr}|$, and $2\cdot\frac98p_0^2$ times this bound equals the error \textup{(E1)} of \Cref{Y-P6-L13}.

\item[\textup{2.}] \emph{Peierls bound.} A translation in $L^{-j}\mathbb Z^3$ carries $p$ to a plaquette based at $0$, and $P_\rho$ is invariant under it by the paragraph on the boards in Step~10. By \Cref{Y-P5-LK}(a), this probability is $\rho(E[1\{\phi_p\ge\varepsilon/2\}\mid U])$. The conditional expectation is a cell function independent of the torus side. Apply \Cref{support:18-38}(b) to this function and use \textup{(P$^\circ$)} to obtain $P_\rho(\phi_p\ge\varepsilon/2)\le e^{-(c_P/2)p_0^2}$. So $P_\rho(\chi<1)\le3(2R+4\rho(x)+7)^3e^{-(c_P/2)p_0^2}=\nu_P$. Since $G_--G\ne0$ only where some $\phi_p>\varepsilon/2$ in $B'$, we get $|\mathrm{Cov}_{P_\rho}(G_--G,S_f)|\le\frac92C_Sp_0^4\nu_P$.

\item[\textup{3.}] \emph{First identity.} We apply \Cref{Y-P5-L10per} with $\Phi^\alpha=u^\alpha\chi(S_f-\bar s)/2$ as its function $\Phi$, where $\bar s:=P_\rho(S_f)$. The sup bound of this function is the one of \Cref{Y-P6-L13}, because $|S_f-\bar s|\le2\sup|S_f|$. Now $\sum_\alpha xu^\alpha\Phi^\alpha=G(S_f-\bar s)$, and its $P_\rho$-mean is $\mathrm{Cov}_{P_\rho}(G,S_f)$. From $Du^\alpha=a_u+O(\varepsilon)$, and as $\bar s$ is the mean of $S_f$, we get \[\frac32a_uP_\rho[\chi(S_f-\bar s)]=-\frac32a_uP_\rho[(1-\chi)(S_f-\bar s)]=O(\sup|S_f|\,P_\rho(\chi<1)).\] The term with $D\chi$ is carried by $\{\text{some }\phi_p\ge\varepsilon/2\text{ in }B'\}$. There remains $\frac12\sum_\alpha P_\rho[u^\alpha\chi DS_f]$.

\item[\textup{4.}] Pointwise, $DS_f=x\langle\mathbb D_k\mathrm{d}w,\tilde F^\alpha\rangle+O(p_0^2C_D)$, by Step~2 of the proof of \Cref{Y-P6-L13}.

\item[\textup{5.}] \emph{Second identity.} The choice $\Phi'^\alpha=\chi\langle\mathbb D_k\mathrm{d}w,\tilde F^\alpha\rangle/2$ yields $\frac32q_NP_\rho(\chi)$ plus errors.

\item[\textup{6.}] \emph{Variance.} The choice $\Phi=u^\alpha\chi/2$ yields $P_\rho G=\frac32a_u+\mathrm{err}$. The function $\Phi=\frac x4u^\alpha|u|^2\chi^2$ satisfies $\sup\le0.844\,p_0^2\varepsilon$ and yields $P_\rho G^2=\frac{15}4a_u^2+\mathrm{err}$. Recall that in \Cref{Y-P6-L13} the error \textup{(E7)} is $4[$\textup{(E2)} $+$ \textup{(E3)} $+$ \textup{(E4)} $+$ \textup{(E6)}$]$, with the errors of that lemma on the right.

\item[\textup{7.}] All the other lines of the proof of \Cref{Y-P6-L13} are pointwise. The list of errors is the one of \Cref{Y-P6-L13}, with $\Delta_N$ added. In $\Delta_N$ we collect the contribution of the single application of \textup{(W$^\circ$)} and of the twelve applications of \textup{(X$^\circ$)}, which arise from three colours and four functions; each enters with a factor that does not depend on $N$, namely $\mathsf b'$, $1/t$ and the sup bounds of Step~10.

\item[\textup{8.}] \emph{What does not depend on the side.} Neither $G_-$, nor $f_*$, nor $\rho$, nor the kernel depends on the side; hence neither $\mathrm{Cov}_{P_\rho}(G_-,\mathcal O_Kf_*)$ nor $\mathrm{Var}_{P_\rho}(G_-)$ does. The argument holds whether or not $q$, $a_u$, $S_f$, $\mathbb D_k$ and $w_R$ depend on the side.

\end{enumerate}
\end{proof}

\emph{Step 12: the witness inequality \textup{(P0)} under $\rho$, and the clauses \textup{(P1)} to \textup{(P4)}.} It is enough to take $c=\frac12$, since the remaining mirrors are reached by the translations in $L^{-j}\mathbb Z^3$ from the paragraph on the boards in Step~10.

\begin{enumerate}

\item[\textup{(P0)}]

\begin{enumerate}

\item[\textup{1.}] The plane $x_1=cL^{-j}$ lies in $a_K(\mathbb Z+\frac12)$, and so in $(a_K/2)\mathbb Z=\operatorname{Adm}(e_1)$; it is therefore a lattice mirror.

\item[\textup{2.}] Let $G_+:=G_-\circ\theta_*$, $\breve\Gamma:=E[G_+\mid U]-\rho(E[G_+\mid U])$, $F:=\mathcal O_Kf-\rho(\mathcal O_Kf)$. Unlike $G_+$, which is a function of the drawn field $V^{(k)}$, the two functions $\breve\Gamma$ and $F$ are bounded Borel functions of finitely many bonds of $U$ in $\mathbb H^+$, by the argument of \Cref{Y-P5-LK}\,\textup{(a)} with a half space instead of a cell: a drawn bond whose two ends lie on one side depends on that side only. If the barycentre of a plaquette lies strictly to the right of a mid-bond plane, then every site of the plaquette lies to the right of that plane.

\item[\textup{3.}] The identity $E[G_-\mid U]=\theta_*E[G_+\mid U]$ is the third claim of \Cref{Y-P6-L1}\,\textup{(K)} for the map $\check\theta$ given by the reflection $y_1\mapsto n_1-y_1$, with $n_1=2c$ an odd integer. The state $\rho$ is $\theta_*$-invariant by \Cref{support:54-17}\,\textup{(b)}; so the two means agree, and $\rho(\mathcal O_K\theta_*f)=\rho(\mathcal O_Kf)$.

\item[\textup{4.}] Since $\mathcal O_Kf$ is a function of $U$, $\mathrm{Cov}_{P_\rho}(G_-,\mathcal O_Kf)=\rho(\theta_*\breve\Gamma\,F)$.

\item[\textup{5.}] By \Cref{support:54-17}\,\textup{(a)}, the form $(A,B)\mapsto\rho(\theta_*A\,B)$ is positive when restricted to real functions in $\mathcal F_{\rm loc}(\mathbb H^+)$, and it is symmetric by \textup{(b)} of that theorem and $\theta_*^2=1$. The Schwarz inequality then gives $\rho(\theta_*\breve\Gamma\,F)^2\le\rho(\theta_*\breve\Gamma\,\breve\Gamma)\,\rho(\theta_*F\,F)$.

\item[\textup{6.}] By Cauchy's inequality and invariance, $\rho(\theta_*\breve\Gamma\,\breve\Gamma)\le\rho(\breve\Gamma^2)$. Moreover $\rho(\breve\Gamma^2)=\mathrm{Var}_{P_\rho} E[G_+\mid U]\le\mathrm{Var}_{P_\rho}G_+=\mathrm{Var}_{P_\rho}G_-$.

\item[\textup{7.}] Finally, $\rho(\theta_*F\,F)=S^T_2[\rho](\theta_*f,f)$.

\end{enumerate}

Items 4 to 7 together give \textup{(P0)}, in product form, so that no quotient occurs.

\item[\textup{(P1)}] We combine \Cref{Y-P5-L13per}\,\textup{(ii)}, taken at the side $N$, with $q_N\ge c_F$ and with entry \textup{(j1)} of the list of \S\ref{Y-P6-1}; this gives \[\mathsf b'\mathrm{Cov}_{P_\rho}(G_-,\mathcal O_Kf_*)\ge\frac32c_F-6c_F/12-\Delta_N=c_F-\Delta_N.\] The left side does not depend on $N$, and letting $N\to\infty$ we obtain \textup{(P1)}.

\item[\textup{(P2)}] Clause \textup{(iii)}, together with $a_{u,N}\le2$ and entry \textup{(j1)}, yields $\mathrm{Var}\le\frac32\cdot4+$ \textup{(E7)} $+\Delta_N\le6+4c_F/3+\Delta_N$, where the error is that of \Cref{Y-P6-L13}. As $N\to\infty$, this becomes $\mathrm{Var}\le7$.

\item[\textup{(P3)}] \textup{(P1)} forces $\mathrm{Var}_{P_\rho}(G_-)>0$. Then \textup{(P0)}, \textup{(P1)} and \textup{(P2)} give $\mathsf b'^2S^T_2\ge c_F^2/7=c_*$.

\item[\textup{(P4)}] The last sentence of \Cref{support:54-25}\,\textup{(2)} gives $\nu^U\in\mathcal G^{\rm per}$ for $\nu$-almost every $U$. This state is ergodic, hence invariant under $a_K\mathbb Z^3$, and therefore $\nu^U\in\mathbb G_K(\mathsf b')$. Moreover $\mathrm{Cov}_\nu(\cdot,\cdot\mid\mathcal I)(U)$ is the covariance under $\nu^U$. We apply \textup{(P3)} to $\nu^U$.

\end{enumerate}
\end{proof}

\begin{remark} For non-triviality alone, \textup{(P2)} can be replaced by the bound $\mathrm{Var}\le(\frac98)^2p_0^4$, which holds because $G_-$ takes its values in $[0,\frac98p_0^2]$.
\end{remark}

\begin{proof}[Proof of \Cref{Y-P5-ThC} and of \Cref{Y-P5-Cor}] \emph{Step 13: the depth, the clauses \textup{(C1)} to \textup{(C4)}, and the corollary.} We first work in normal form, $\mathsf b>x_s$. Item \textup{(C0)} shows that $j^+(\mathsf b)<\infty$ and $r_*(\mathsf b)>0$, so that the ranges $j\ge j^+(\mathsf b)$ and $0<r\le r_*(\mathsf b)$ are not empty. Item \textup{(C4)} treats every $\mathsf b>0$.

\begin{enumerate}

\item[\textup{(C0)}] \emph{The depth.} By \Cref{Y-P6-N}\,\textup{(D)} and \Cref{Y-P6-L16}\,\textup{(a)}, \textup{(c)}, the conditions on the depth listed in \S\ref{Y-P6-1} and the clauses \textup{(N0)} to \textup{(N6)} of \Cref{Y-P6-N} are satisfied for every $j\ge j_*(\mathsf b')$. By \Cref{Y-P6-L16}\,\textup{(d)}, $j_*$ is bounded on the compact set $U(\mathsf b)\subset[x_s,\infty)$. Hence $j^+(\mathsf b)$ is an integer, $r_*(\mathsf b)=D_0L^{-j^+(\mathsf b)}>0$, and every $j\ge j^+(\mathsf b)$ is $\ge j_*(\mathsf b')$ for every $\mathsf b'\in U(\mathsf b)$.

\item[\textup{(C1)}] We take $\mathsf b$ in normal form. Let $S\in\mathrm{Lim}_G(\mathsf b)$ with the frame $(K_i,\nu_i)$, $\nu_i\in\mathbb G_{K_i}(\mathsf b_i)$, $\mathsf b_i\to \mathsf b$, all $\mathsf b_i\in U(\mathsf b)$, and let $\xi\in\mathsf X$. By \Cref{support:35-36}\,\textup{(iii)} there are $\rho_i=\nu_i(\cdot\mid E_i)\in\mathbb G_{K_i}(\mathsf b_i)$ which are approximants of $\sigma_\xi$, and this for every $\xi\in\mathsf X$. By \Cref{support:36-12}\,\textup{(d)}, $S_2^{\sigma_\xi}=\lim_iS^T_2[\rho_i]$ on $\mathfrak S_2(\mathbb R^3)$.

Fix $j\ge j^+(\mathsf b)$. For every $i$ we then have $j\ge j_*(\mathsf b_i)$, and once $i$ is large, $K_i\ge j+1$. Only $D_0$ and $j$ enter $f_*=f_0(L^j\,\cdot)$ and the mirror, whereas $G_-$, $R$ and $h$ depend on $x_i=L^j\mathsf b_i$; these are allowed to change with $i$, since they are used only as witnesses. As $f_0$ is $C^\infty$ and the supports of $\theta_*f_*$ and $f_*$ are $\ge D_0L^{-j}$ apart, $(\theta_*f_*,f_*)$ is a separated $C^1_c$ tuple, so it lies in $\mathfrak S_2(\mathbb R^3)$, the set on which the limit above holds. By \textup{(P3)}, $S^T_2[\rho_i](\theta_*f_*,f_*)\ge c_*/\mathsf b_i^2\to c_*/\mathsf b^2$. By \textup{(E3)(i)} the mirror can be carried to $(e_3,0)$. By \textup{(E7)(b)}, non-triviality amounts to $g_f(\xi)>0$ for one $f\ge0$.

\item[\textup{(C2)}] We have $g_{f_*}(\xi)=\int\!\!\int f_*(\theta_*x)f_*(y)s_\xi(|x-y|)$, where $s_\xi\ge0$ is non-increasing. The set $\operatorname{supp}f_*$ keeps a distance $\ge(\lceil D_0/\sqrt2\rceil-\frac12-D_0/8)L^{-j}\ge0.58D_0L^{-j}$ from the mirror, as $1/\sqrt2-1/8=0.5821\dots$; hence $|\theta_*x-y|\ge D_0L^{-j}$ for $x$ and $y$ in that set. Since $0\le f_*\le1$ and $f_*$ is carried by a sup-ball of radius $(D_0/8)L^{-j}$, we have $\int f_*\le(D_0/4)^3L^{-3j}$. As $s_\xi$ is non-increasing and $|\theta_*x-y|\ge D_0L^{-j}$, it follows that $g_{f_*}(\xi)\le s_\xi(D_0L^{-j})(\int f_*)^2\le s_\xi(D_0L^{-j})(D_0/4)^6L^{-6j}$. By \textup{(C1)}, $c_*/\mathsf b^2\le g_{f_*}(\xi)$, and so $c_*/\mathsf b^2\le s_\xi(D_0L^{-j})(D_0/4)^6L^{-6j}$. For $D_0L^{-j-1}<r\le D_0L^{-j}$ this gives \[s_\xi(r)\ge s_\xi(D_0L^{-j})\ge c_*\mathsf b^{-2}(4/D_0)^6L^{6j}>c_*(4/L)^6\mathsf b^{-2}r^{-6},\] because $L^j>D_0/(Lr)$. Taken together, all $j\ge j^+$ cover $(0,r_*]$.

Suppose that $H_\xi$ were bounded, $H_\xi\le\check\Lambda$. For $t<t_0$ and one $\mathsf e\le t/2$ we would then have $s(t_0)=\int e^{-(t_0-2\mathsf e)E}\,d\mu_{\mathsf e}\ge e^{-(t_0-t)\check\Lambda}s(t)$, so that $s$ would stay bounded as $t\to0$, which contradicts the lower bound with $r^{-6}$. In finite dimension $H$ is necessarily bounded; hence $\dim\mathcal H(\sigma_\xi)$ is infinite.

\item[\textup{(C3)}] This follows from \textup{(E7)(a)} and \textup{(C1)}.

\item[\textup{(C4)}] Put $z=z(\mathsf b)$. Step~8\,\textup{(1)} gives $S^T_2[\rho;K](f,f')=L^{2z}S^T_2[\rho;K-z](f^{[z]},f'^{[z]})$; moreover \Cref{Y-P5-ThP}, taken at $(K-z,L^z\mathsf b')$, concerns the same measures. We take $f=f_0(L^{j+z}\,\cdot)$, so that $f^{[z]}=f_0(L^j\,\cdot)$, with $j\ge j^+(L^z\mathsf b)$. Then \[\mathsf b_i^2S^T_2[\rho_i;K_i](\theta_*f,f)=(L^z\mathsf b_i)^2S^T_2[\rho_i;K_i-z](\theta_*f^{[z]},f^{[z]})\ge c_*.\] The steps of \textup{(C1)} and \textup{(C2)} go through with $L^{-j-z}$ in the place of $L^{-j}$ in each of their formulas. In the limit, $S_n^{(\mathsf b)}(X)=L^{zn}S_n^{(L^z\mathsf b)}(X^{[z]})$, and hence $\mathsf S_n^{(\mathsf b)}(x)=L^{4zn}\mathsf S_n^{(L^z\mathsf b)}(L^zx)$ and $s^{(\mathsf b)}(r)=L^{8z}s^{(L^z\mathsf b)}(L^zr)$. The quantities $\mathsf b^2S_2$ and $\mathsf b^2r^6s(r)$ are therefore invariant, as $2+6=8$.

\end{enumerate}

\emph{The corollary.} By \textup{(E1)}, $\mathrm{Lim}_G(\mathsf b)$ is not empty, since $\mathcal G^{\rm tor}$ is not empty by compactness. For $\xi\in\mathsf X$ we have $S^\xi:=(S_n^{\sigma_\xi})_{n\ge2}\in\mathrm{Lim}_G(\mathsf b)$, with centred functional $\sigma_\xi^{\rm c}$, and $\sigma_\xi=\sigma_\xi^{\rm c}\circ\mathrm{sh}_{c_\xi}$ by \textup{(E3)(i)}. Thus, up to the shift, $\sigma_\xi$ is the functional of a member, and it coincides with that functional exactly when $c_\xi=0$. By \Cref{support:36-17}, $\sigma_\xi$ clusters, and clustering passes through $\mathrm{sh}_{-c}$; so the one-point space is a canonical decomposition of $\sigma_\xi^{\rm c}$, and by \Cref{support:36-9}\,\textup{(a)} it is the only one. The list of the corollary follows from \textup{(E3)}, \textup{(E4)} and \Cref{Y-P5-ThC}. This list is stated for $\sigma_\xi$; \Cref{Y-P5-LShift} carries it over to $\sigma_\xi^{\rm c}$, and the two have the same $S_2$.

\end{proof}

\appendix
\section{Reflection positivity and chessboard estimates}\label{app:reflection-and-chessboard}

\subsection{Wilson reflection positivity}
\label{sec:plaquette-reflection}

To estimate the source derivatives uniformly in the volume, we first prove the reflection inequality for the Wilson measure. We use rectangular tori because the chessboard proof treats one coordinate direction at a time.

Let $a>0$, let $N_\mu\ge3$ be integers, and put
\[
 T=\prod_{\mu=1}^3\mathbb R/(aN_\mu\mathbb Z),\qquad
 \Lambda=\prod_{\mu=1}^3 a\mathbb Z/(aN_\mu\mathbb Z).
\]
Positive bonds and plaquettes have the usual coordinate orientations.
For reversed bonds set $U(\bar b)=U(b)^{-1}$. Let $\partial U(p)$ denote the ordered product around $p$ and put $a_p(U)=1-\tfrac12\operatorname{tr}\partial U(p)$.
For $\beta\ge0$ define
\begin{equation}\label{support:1-9}
 A(U)=\sum_p a_p(U),\qquad
 \mathcal O_a(f)=a^{-1}\sum_p f(x(p))a_p(U),\qquad
 d\mu_{a,T,\beta}=Z_{a,T,\beta}^{-1}e^{-\beta A(U)}\,dU.
\end{equation}
Here $dU$ is product Haar probability measure and $Z_{a,T,\beta}$
is the integral of $e^{-\beta A}$. Since $0\le a_p\le2$,
this normalization is strictly positive, and $\mathcal O_a(f)$ is
bounded for every function $f$ on the finite set of plaquette centres.
For a bounded Borel function $F$ of the configuration, write
\[
 \mathcal Z_{a,T,\beta}(F)=\int F(U)e^{-\beta A(U)}\,dU,
 \qquad \langle F\rangle_{a,T,\beta}
 =Z_{a,T,\beta}^{-1}\mathcal Z_{a,T,\beta}(F).
\]
Thus, on the cubic lattice with $a=L^{-K}$,
$Z_{K,\ell,\beta}[w]=\mathcal Z_{a,T_\ell,\beta}
(e^{\mathcal O_a(w)})$.

For $n\ge1$ the joint cumulant is
\begin{equation}\label{support:1-23}
 S^T_{n;a,T,\beta}(f_1,\ldots,f_n)
 =\left.\partial_{z_1}\cdots\partial_{z_n}
 \log\left\langle e^{\sum_i z_i\mathcal O_a(f_i)}
 \right\rangle_{a,T,\beta}\right|_{z=0}.
\end{equation}
The expectation is entire and equals one at zero, so the logarithm
in this formula exists near zero. If $\mathcal P_n$ denotes the set
of partitions of $\{1,\ldots,n\}$, differentiating the logarithm gives
\begin{equation}\label{support:1-25}
 S^T_n(f_1,\ldots,f_n)=
 \sum_{\pi\in\mathcal P_n}(-1)^{|\pi|-1}(|\pi|-1)!
 \prod_{B\in\pi}\left\langle\prod_{i\in B}\mathcal O_a(f_i)
 \right\rangle.
\end{equation}
Conversely, the joint moment is the sum over $\pi\in\mathcal P_n$
of the products of the cumulants indexed by its blocks. Both identities
follow by comparing the coefficient of $z_1\cdots z_n$ in the
exponential identity relating the expectation to its logarithm.
In particular, cumulants are symmetric and multilinear. For $n\ge2$
the cumulant of the random variables $\mathcal O_a(f_i)$ is unchanged
if a scalar constant is added to any of these random variables:
such a change adds a linear function of $z$ to the logarithm.

Translations and signed coordinate permutations transport oriented bonds
and preserve product Haar measure, with inversion at reversed bonds.
They carry a plaquette holonomy to a conjugate of the original holonomy
or its inverse, both with the same trace in $\operatorname{SU}(2)$.
The Wilson law is therefore invariant under translations and under
signed permutations that preserve the rectangular period lattice.
Other signed permutations carry it to the Wilson law on the permuted
torus. In either case $\mathcal O_a(f)$ transforms by pullback of $f$.

\subsubsection{The operator across one layer}

Fix a coordinate direction, write $N$ for the number of sites in
that direction, and let $\Sigma$ be the two-dimensional cross-section.
Let $\mathsf B(\Sigma)$ be its positive bonds and set
$\mathcal X=\operatorname{SU}(2)^{\mathsf B(\Sigma)}$.
A configuration consists of the transverse fields $V^t\in\mathcal X$
and the bonds $S^t\in\operatorname{SU}(2)^\Sigma$ joining layers $t$ and
$t+1$, with $t$ taken modulo $N$. Write $A_\Sigma(V)$ for the sum
of the plaquette actions within a transverse layer and put
\begin{equation}\label{support:3-6}
 A_\times(V,S,W)=\sum_{b\in\mathsf B(\Sigma)}
 \left(1-\tfrac12\operatorname{tr}
       V_bS_{b_+}W_b^{-1}S_{b_-}^{-1}\right).
\end{equation}
The two types of plaquettes partition all plaquettes, giving
\[
 A(U)=\sum_{t\bmod N}
       \bigl[A_\Sigma(V^t)+A_\times(V^t,S^t,V^{t+1})\bigr].
\]
On $\mathcal H=L^2(\mathcal X,dV)$ define the integral operator
\[
 \mathsf K(W,V)=\int e^{-\beta A_\times(V,S,W)}\,dS.
\]
The kernel is continuous, symmetric and bounded by one. Symmetry follows
by replacing $S$ by $S^{-1}$ and using invariance of the trace under
cyclic permutations and inversion. Thus $\mathsf K$ is self-adjoint
and Hilbert--Schmidt. Its positivity needs a separate argument.

\begin{proposition}\label{support:3-16}
The crossing operator satisfies $\mathsf K\ge0$.
\end{proposition}
\begin{proof}
Put $\mathsf k(u)=e^{-\beta(1-\operatorname{tr}u/2)}$ and
$\mathsf K_0(V,W)=\prod_b\mathsf k(V_bW_b^{-1})$.
For $R\in\operatorname{SU}(2)^\Sigma$ let
$(V^R)_b=R_{b_-}V_bR_{b_+}^{-1}$.
The crossing weight is $\mathsf K_0(V,W^S)$, and
$\mathsf K_0(V^R,W^R)=\mathsf K_0(V,W)$. Haar invariance therefore gives
\begin{equation}\label{support:3-18}
 \mathsf K(W,V)=\int dR\,dS\,\mathsf K_0(V^R,W^S).
\end{equation}
For $v,w\in\operatorname{SU}(2)$,
$\operatorname{tr}(vw^{-1})=\sum_{i,j}v_{ij}\overline{w_{ij}}$.
Expanding the exponential at every bond consequently yields
\[
 \mathsf K_0(V,W)=\sum_\alpha c_\alpha
      \phi_\alpha(V)\overline{\phi_\alpha(W)},\qquad c_\alpha\ge0.
\]
Here $\alpha$ specifies a nonnegative integer $n_b$ and $n_b$ pairs
of matrix indices at each bond,
\[
 c_\alpha=e^{-\beta|\mathsf B(\Sigma)|}
       \prod_b\frac{(\beta/2)^{n_b}}{n_b!},\qquad
 \phi_\alpha(V)=\prod_b\prod_{s=1}^{n_b}(V_b)_{i_{b,s}j_{b,s}}.
\]
Every matrix entry has modulus at most one. There are $4^{n_b}$
choices of indices at bond $b$, so
\[
 \sum_\alpha c_\alpha
 =e^{-\beta|\mathsf B(\Sigma)|}
       \prod_b\sum_{n\ge0}\frac{(2\beta)^n}{n!}
 =e^{\beta|\mathsf B(\Sigma)|}<\infty.
\]
This proves absolute uniform convergence. Put
$\widehat\phi_\alpha(V)=\int\phi_\alpha(V^R)\,dR$. Substituting the series in \eqref{support:3-18} gives
\[
 \mathsf K(W,V)=\sum_\alpha c_\alpha
       \widehat\phi_\alpha(V)\overline{\widehat\phi_\alpha(W)}.
\]
Since $dV$ is a probability measure, every $\psi\in\mathcal H$
is integrable. Absolute convergence permits termwise integration, giving
\[
 \langle\psi,\mathsf K\psi\rangle
   =\sum_\alpha c_\alpha
       \left|\int\widehat\phi_\alpha(V)\psi(V)\,dV\right|^2\ge0.
\]
\end{proof}

\subsubsection{Slabs and the trace formula}

For an integer $q\ge2$, a slab of $q$ layers has variables
$(V^1,\ldots,V^q;S^1,\ldots,S^{q-1})$ and weight
\[
 w_q=\exp\left[-\beta\sum_{t=1}^q A_\Sigma(V^t)
       -\beta\sum_{t=1}^{q-1}A_\times(V^t,S^t,V^{t+1})\right].
\]
For a bounded Borel function $\Phi$ of these variables define
\begin{equation}\label{support:3-10}
 \mathfrak S[\Phi](W,V)=
 \int\prod_{1<t<q}dV^t\prod_{1\le t<q}dS^t\,
       (\Phi w_q)\big|_{V^1=V,V^q=W},\qquad
 \mathfrak Y[\Phi]=\mathsf K^{1/2}\mathfrak S[\Phi]\mathsf K^{1/2}.
\end{equation}
These are Hilbert--Schmidt operators, since the first kernel is bounded
by $\|\Phi\|_\infty$ and $\mathsf K^{1/2}$ is bounded.
Reverse the order of the layers, invert the joining bonds, and conjugate
the value of $\Phi$; denote the resulting function by $\Phi^\Theta$.
This change of variables preserves Haar measure and $w_q$, and exchanges
the two endpoint fields. Consequently
\begin{equation}\label{support:3-25}
 \mathfrak S[\Phi^\Theta]=\mathfrak S[\Phi]^*,\qquad
 \mathfrak Y[\Phi^\Theta]=\mathfrak Y[\Phi]^*.
\end{equation}

\begin{proposition}[Trace formula]\label{support:3-26}
Partition the $N$ layers cyclically into $r\ge2$ consecutive slabs
of lengths $q_s\ge2$. Let $\Phi_s$ be a bounded Borel function
of the inner bonds of slab $s$. Then
\begin{equation}\label{support:3-27}
 \mathcal Z_{a,T,\beta}\left(\prod_{s=1}^r\Phi_s\right)
 =\operatorname{Tr}\bigl(\mathfrak S[\Phi_r]\mathsf K\cdots
                  \mathfrak S[\Phi_1]\mathsf K\bigr)
 =\operatorname{Tr}\bigl(\mathfrak Y[\Phi_r]\cdots
                  \mathfrak Y[\Phi_1]\bigr).
\end{equation}
All the operators attached to a slab depend on its length, the
cross-section and $\beta$, but not on the total length $N$.
\end{proposition}
\begin{proof}
Keep the first and last transverse fields $V_s,W_s$ of each slab.
Every plaquette either lies inside one slab or crosses from its last
layer to the first layer of the next slab. Integrating all other
variables gives
\[
 \int\prod_{s=1}^r dV_s\,dW_s\,
       \prod_{s=1}^r\mathfrak S[\Phi_s](W_s,V_s)
                         \mathsf K(V_{s+1},W_s),\qquad V_{r+1}=V_1.
\]
The integrands are bounded and all measures are finite, so Fubini
applies. The kernel formula for a product, followed by the
Hilbert--Schmidt identity
$\operatorname{Tr}(CD)=\iint C(y,x)D(x,y)\,dx\,dy$,
gives the first trace. There are at least two Hilbert--Schmidt factors,
so the product is trace class. Splitting each $\mathsf K$ into two
square roots and moving the last square root around the trace gives
the second identity. The remaining factors form a trace-class operator,
which justifies this cyclic permutation. The definitions contain no
total layer number $N$, proving the last assertion.
\end{proof}

\begin{corollary}[Reflection at bond midpoints]\label{support:4-8}
Suppose that $N$ is even. Let $\theta$ be reflection in a coordinate plane at a half-integer multiple of $a$, and let $T_+$ be one open half-torus bounded by this plane and its opposite parallel plane. Write $\mathfrak F_+$ for the bounded Borel functions depending only on bonds whose two endpoints lie in $T_+$.
\begin{enumerate}
\item[(a)] The form $\langle\overline{\theta F}H\rangle_{a,T,\beta}$ on $\mathfrak F_+$ is positive semidefinite.
\item[(b)] It is Hermitian and sesquilinear.
\item[(c)] For $F,H\in\mathfrak F_+$,
\begin{equation}\label{support:4-10}
 |\langle\overline{\theta F}H\rangle|^2
 \le\langle\overline{\theta F}F\rangle
       \langle\overline{\theta H}H\rangle.
\end{equation}
\item[(d)] The assertions hold on the other open half-torus as well.
\item[(e)] When the plane is a wall of a board in \Cref{support:4-1}, the product of functions of cells on one side is a function of that open half-torus.
\end{enumerate}
\end{corollary}
\begin{proof}
The two halves contain $q=N/2\ge2$ layers each. Their inner bond
variables identify $F,H$ with slab functions $\Phi,\Psi$.
Reflection sends the first slab to the second and reverses its layer
order, so $\overline{\theta F}$ is the placed function $\Phi^\Theta$.
By \eqref{support:3-25} and \ref{support:3-26},
\[
 \mathcal Z_{a,T,\beta}(\overline{\theta F}H)
 =\operatorname{Tr}\bigl(\mathfrak Y[\Phi]^*\mathfrak Y[\Psi]\bigr).
\]
This is the Hilbert--Schmidt inner product. Its positivity, Hermitian symmetry and Cauchy--Schwarz inequality survive division by $Z_{a,T,\beta}>0$. The same trace proof starts with either half-torus. For (e), the two walls bounding a half-torus have half-integer coordinates in units of $a$. They contain no sites. The cells of that half-torus have their sites strictly between the two walls, so all their inner bonds have both endpoints on that side.
\end{proof}

\subsection{The chessboard estimate on a torus}\label{sec:chessboard-torus}
\label{sec:chessboard}

The trace formula reduces the chessboard estimate in one coordinate
direction to an inequality for the slab operators. Applying this
bound in all three directions gives the estimate on the full board.

\subsubsection{Cells and reflected copies}

\begin{definition}[A board on a torus]\label{support:4-1}
\begin{enumerate}
\item[(1)] We use the Wilson measure of \eqref{support:1-9} and its action on functions by lattice isometries.
\item[(2)] Fix $a>0$, an odd integer $\lambda\ge3$, even integers
$\mathsf n_\mu\ge2$, and an offset $o\in a\mathbb Z^3$.
We call these data a board. Its torus has sides
$\ell_\mu=\mathsf n_\mu\mathsf l$, where $\mathsf l=\lambda a$.
Write
\begin{equation}\label{support:2-3}
 \mathsf Z=\prod_{\mu=1}^3\mathbb Z/\mathsf n_\mu\mathbb Z,
 \qquad \mathsf N=\prod_\mu\mathsf n_\mu,
 \qquad h_\lambda=(\lambda-1)/2.
\end{equation}
The cell $\square_\kappa$ is the image in the torus of
$o+\mathsf l\tilde\kappa+[-\mathsf l/2,\mathsf l/2]^3$, with
$\tilde\kappa$ any integer representative of $\kappa\in\mathsf Z$.
Its sites are
\begin{equation}\label{support:2-8}
 o+\mathsf l\tilde\kappa+an,\qquad n\in\mathsf C_\lambda
 :=\{-h_\lambda,\ldots,h_\lambda\}^3.
\end{equation}
The walls normal to $e_\mu$ have coordinates
\begin{equation}\label{support:2-14}
 x_\mu=o_\mu+\mathsf l(t+\tfrac12),\qquad t\in\mathbb Z/\mathsf n_\mu\mathbb Z.
\end{equation}
The cell faces contain no sites since their varying coordinate is a half-integer multiple of $a$. A plaquette whose centre lies in a cell interior has all its corners in that cell. Otherwise its centre lies on a cell face.
\end{enumerate}
\end{definition}

\begin{lemma}[Cell coordinates]\label{support:2-9}
\begin{enumerate}
\item[(a)] The sites of each cell have the representatives \eqref{support:2-8}.
\item[(b)] The number of cells is $\mathsf N$.
\item[(c)] Every site belongs to exactly one cell.
\item[(d)] A function of a cell is a bounded complex Borel function of its inner bond variables, where an inner bond has both endpoints in the cell. Transporting the positive bonds of $\mathsf C_\lambda$
to $\square_\kappa$ identifies such a function with a unique model
function $F_0$ on the finite product
$\operatorname{SU}(2)^{\mathsf B^+(\mathsf C_\lambda)}$; its placement
in $\square_\kappa$ is denoted by $F_0^{(\kappa)}$.
\end{enumerate}
\end{lemma}
\begin{proof}
The site formula gives (a). The cardinality in (b) follows by multiplying the three index ranges. In each coordinate the sets of $\lambda$ consecutive sites of the cells are disjoint and exhaust the sites in that coordinate. This gives (c). For (d), arbitrary values of the inner bond variables extend to a torus configuration, so placement preserves the supremum norm and determines the model function uniquely.
\end{proof}

\begin{definition}[Copies and folds]\label{support:4-2}
\begin{enumerate}
\item[(1)] For $\kappa,\kappa'\in\mathsf Z$ set
$\epsilon_\mu=(-1)^{\tilde\kappa'_\mu-\tilde\kappa_\mu}$ and define
the diagonal isometry
\begin{equation}\label{support:2-18}
 \mathsf r_{\kappa'\kappa}(x)
 =o+\mathsf l\tilde\kappa'
    +\epsilon(x-o-\mathsf l\tilde\kappa).
\end{equation}
It carries $\square_\kappa$ to $\square_{\kappa'}$.
The evenness of each $\mathsf n_\mu$ makes the signs independent
of the representatives; changing a representative changes the
formula only by a torus period. Transport configurations and
observables by the action used in \ref{sec:plaquette-reflection},
including inversion when a bond orientation is reversed.
For a function $F$ of $\square_\kappa$, define
\begin{equation}\label{support:2-26}
 F^{[\kappa']}=
 \begin{cases}
  \mathsf r_{\kappa'\kappa}F,&\prod_\mu\epsilon_\mu=1,\\
  \overline{\mathsf r_{\kappa'\kappa}F},&\prod_\mu\epsilon_\mu=-1,
 \end{cases}
\end{equation}
\item[(2)] The fold is
\begin{equation}\label{support:2-27}
 \operatorname{Fold}F=\prod_{\kappa'\in\mathsf Z}F^{[\kappa']}.
\end{equation}
\item[(3)] For a sign vector $\epsilon$ and a model bond variable $u$, put
\begin{equation}\label{support:2-28}
 u^\epsilon(\langle n,n+e_\mu\rangle)=
 \begin{cases}
 u(\langle\epsilon n,\epsilon n+e_\mu\rangle),&\epsilon_\mu=1,\\
 u(\langle\epsilon n-e_\mu,\epsilon n\rangle)^{-1},&\epsilon_\mu=-1.
 \end{cases}
\end{equation}
The sign map preserves the model cube, so both positive bond arguments lie in the model. Formula \eqref{support:2-26} is therefore the placement of $F_0(u^\epsilon)$, with conjugation for a negative product of the signs.
\end{enumerate}
\end{definition}
\begin{lemma}[Operations on copies]\label{alt:copy-operations}
Copies preserve products and sums and commute with conjugation.
For all cell indices,
\begin{equation}\label{support:2-30}
 (F^{[\kappa']})^{[\kappa'']}=F^{[\kappa'']},
 \qquad \operatorname{Fold}(F^{[\kappa']})
 =\operatorname{Fold}F.
\end{equation}
\end{lemma}
\begin{proof}
Pullback and conjugation preserve products and sums and commute with
conjugation. The mirror maps compose since their sign matrices
multiply to the sign matrix for the composite map. The determinants
multiply as well, so the composite copy conjugates exactly when the
composite mirror map has negative determinant. This proves the first
identity. Multiplying over $\kappa''$ gives the second.
\end{proof}
Exactly half the cells have each parity, since changing one cell
index by one interchanges the parities. Consequently
\begin{equation}\label{support:2-31}
 \operatorname{Fold}(cF)=|c|^{\mathsf N}\operatorname{Fold}F,
 \qquad
 \operatorname{Fold}\bar F=\overline{\operatorname{Fold}F},
 \qquad
 |\operatorname{Fold}F|\le\|F\|_\infty^{\mathsf N}.
\end{equation}

\subsubsection{Multiple reflections}

We first give the trace inequality needed for an arbitrary even
number of slabs. For a Hilbert--Schmidt operator $A$ put
$\|A\|_{2k}=(\operatorname{Tr}(A^*A)^k)^{1/(2k)}$.
This is finite. Cyclicity of the trace gives $\|A^*\|_{2k}=\|A\|_{2k}$. The spectral theorem for $A^*A$ shows that this number vanishes only when $A=0$.

\begin{lemma}[Trace H\"older inequality]\label{support:3-35}
For $k\ge1$ and Hilbert--Schmidt operators $A_1,\ldots,A_{2k}$,
the product is trace class and
\begin{equation}\label{support:3-36}
 |\operatorname{Tr}(A_1\cdots A_{2k})|
 \le\prod_{s=1}^{2k}\|A_s\|_{2k}.
\end{equation}
\end{lemma}
\begin{proof}
The trace-class property follows from the presence of two
Hilbert--Schmidt factors. Unless one factor is zero, normalize each
$A_s$ to have $\|A_s\|_{2k}=1$. Let $\mathcal A$ be the finite
set consisting of these normalized operators and their adjoints,
and maximize $|\operatorname{Tr}(B_1\cdots B_{2k})|$ over
$\mathcal A^{2k}$; denote the maximum by $m$.
An alternating tuple $(C,C^*,\ldots,C,C^*)$ has trace one,
so $m\ge1$.

For $X=B_1\cdots B_k$ and $Y=B_{k+1}\cdots B_{2k}$,
Hilbert--Schmidt Cauchy--Schwarz gives
\[
 |\operatorname{Tr}(XY)|^2
 \le\operatorname{Tr}(XX^*)\operatorname{Tr}(Y^*Y).
\]
Both traces on the right are nonnegative and correspond to
tuples in $\mathcal A^{2k}$. If the original tuple maximizes
the modulus, both reflected tuples also maximize it.
Among all maximizing tuples choose one with a longest string
of consecutive entries alternating between an operator and its
adjoint, counting positions cyclically. Let its length be $j$.
If $j\le k$, cyclically shift it to end at position $k$;
if $k<j<2k$, shift it to start at position one. Reflect the first
$k$ entries to obtain
$(B_1,\ldots,B_k,B_k^*,\ldots,B_1^*)$.
This is again a maximizer and has an alternating string of length
$2\min\{j,k\}>j$, a contradiction. Thus $j=2k$ and the
maximizing trace is one. Rescaling proves the inequality.
\end{proof}

Let $\mathbf F=(F_\kappa)_{\kappa\in\mathsf Z}$ be a family
with $F_\kappa$ a function of $\square_\kappa$, and abbreviate
\[
 Z[\mathbf F]=\mathcal Z_{a,T,\beta}\left(\prod_\kappa F_\kappa\right).
\]
For a direction $\mu$ and an index $t\bmod\mathsf n_\mu$,
write $\kappa|t$ for $\kappa$ with its $\mu$th component
replaced by $t$. Define the reflected family by
\begin{equation}\label{support:4-3}
 (R_{\mu,t}\mathbf F)_\kappa=(F_{\kappa|t})^{[\kappa]}.
\end{equation}
The cells with $\kappa_\mu=t$ lie in a slab of $\lambda$
layers, starting just after layer
$\tau_0(t)=o_\mu/a+\lambda t-h_\lambda-1$.
Their product is a bounded Borel function $\Phi^{(t)}$ of the
inner slab bonds. Here $2\le\lambda\le\lambda\mathsf n_\mu-1$,
so these are slabs to which \ref{support:3-26} applies.

\begin{proposition}\label{support:4-5}
For every board, every $\beta\ge0$, and every such family,
\begin{align}
 Z[R_{\mu,t}\mathbf F]
 &=\operatorname{Tr}
   \bigl(\mathfrak Y[\Phi^{(t)}]^*
              \mathfrak Y[\Phi^{(t)}]\bigr)^{\mathsf n_\mu/2}
 =\|\mathfrak Y[\Phi^{(t)}]\|_{\mathsf n_\mu}^{\mathsf n_\mu}
 \ge0,\label{support:4-6}\\
 |Z[\mathbf F]|&\le
 \prod_{t\bmod\mathsf n_\mu}
             Z[R_{\mu,t}\mathbf F]^{1/\mathsf n_\mu}.
 \label{support:4-7}
\end{align}
\end{proposition}
\begin{proof}
The trace formula expresses $Z[\mathbf F]$ as the trace of
the product of the $\mathsf n_\mu$ slab operators, in reverse
cyclic order. In $R_{\mu,t}\mathbf F$, an even displacement
from slab $t$ is a translation of its slab function; an odd
displacement is a translation of the reflected, conjugated function
$(\Phi^{(t)})^\Theta$. Indeed, in the varying coordinate the
map between two slabs is respectively a translation by
$(t'-t)\mathsf l$ or reflection followed by translation, with
reflection formula $x_\mu\mapsto2o_\mu+(t+t')\mathsf l-x_\mu$.
The other coordinates are unchanged. Thus the operators in this
trace alternate between $Y=\mathfrak Y[\Phi^{(t)}]$ and $Y^*$,
by \eqref{support:3-25}. Since $\mathsf n_\mu$ is even, their trace
is $\operatorname{Tr}(Y^*Y)^{\mathsf n_\mu/2}$, independently
of the initial parity. This proves \eqref{support:4-6}.
Apply \ref{support:3-35} to the original slab product and use
\eqref{support:4-6} for each factor to obtain \eqref{support:4-7}.
\end{proof}

Composition of copies in \eqref{support:2-30} implies that
$R_{3,t_3}R_{2,t_2}R_{1,t_1}\mathbf F$ is the family of all
copies of $F_{(t_1,t_2,t_3)}$. In particular its product is
$\operatorname{Fold}F_{(t_1,t_2,t_3)}$.
For the family of all copies of one $F$, applying $R_{1,t}$
does not change the family. Equation \eqref{support:4-6} therefore
shows that $\mathcal Z_{a,T,\beta}(\operatorname{Fold}F)\ge0$.
\begin{definition}[Chessboard seminorm]\label{support:4-13}
Define
\begin{equation}\label{support:4-14}
 \|F\|_{\mathrm{ch}}
   =\langle\operatorname{Fold}F\rangle_{a,T,\beta}^{1/\mathsf N}.
\end{equation}
\end{definition}

\begin{proposition}[Chessboard estimate]\label{support:4-15}
For every board, every $\beta\ge0$, and every family of
bounded Borel functions $F_\kappa$ of their respective cells,
\begin{equation}\label{support:4-16}
 \left|\left\langle\prod_{\kappa\in\mathsf Z}F_\kappa
                  \right\rangle_{a,T,\beta}\right|
 \le\prod_{\kappa\in\mathsf Z}\|F_\kappa\|_{\mathrm{ch}}.
\end{equation}
\end{proposition}
\begin{proof}
Apply \eqref{support:4-7} in directions $1$, $2$, and $3$, in
that order. The first reflected integrals are nonnegative, so the
inequality can be applied successively under their positive roots.
The composition of copies just proved gives
\[
 |Z[\mathbf F]|
 \le\prod_{t_1,t_2,t_3}
 \mathcal Z_{a,T,\beta}
       (\operatorname{Fold}F_{(t_1,t_2,t_3)})^{1/\mathsf N}.
\]
Divide by $Z_{a,T,\beta}>0$, distributing its $\mathsf N$th
root among the $\mathsf N$ factors.
\end{proof}

\begin{proposition}[Properties of the chessboard seminorm]\label{support:4-17}
\begin{enumerate}
\item[(a)] $\|cF\|_{\mathrm{ch}}=|c|\|F\|_{\mathrm{ch}}$.
\item[(b)] $\|1\|_{\mathrm{ch}}=1$.
\item[(c)] $\|\bar F\|_{\mathrm{ch}}=\|F\|_{\mathrm{ch}}$.
\item[(d)] A mirror copy has the same fold and the same seminorm.
\item[(e)] $\|F\|_{\mathrm{ch}}\le\|F\|_\infty$.
\item[(f)] $|\langle F\rangle_{a,T,\beta}|\le\|F\|_{\mathrm{ch}}$. More generally, for functions $F_k$ of distinct cells, $|\langle\prod_kF_k\rangle_{a,T,\beta}|\le\prod_k\|F_k\|_{\mathrm{ch}}$.
\item[(g)] $\|F+H\|_{\mathrm{ch}}\le\|F\|_{\mathrm{ch}}+\|H\|_{\mathrm{ch}}$ for functions of the same cell.
\item[(h)] For a fixed model function the seminorm is independent of the chosen cell and of the offset of the board.
\end{enumerate}
\end{proposition}
\begin{proof}
The homogeneity, conjugation, and supremum assertions follow
from \eqref{support:2-30}--\eqref{support:2-31} and the nonnegativity
of the folded integral. Set the other cell functions equal to
one in \Cref{support:4-15} to get (f). A constant-one function has fold one, which gives (b). For the triangle inequality, expand
\[
 \operatorname{Fold}(F+H)
  =\sum_{J\subset\mathsf Z}
        \prod_{\kappa\in J}F^{[\kappa]}
        \prod_{\kappa\notin J}H^{[\kappa]}.
\]
The chessboard estimate and invariance under copies bound the
modulus of the mean of the $J$th term by
$\|F\|_{\mathrm{ch}}^{|J|}
 \|H\|_{\mathrm{ch}}^{\mathsf N-|J|}$.
Summing and taking the $\mathsf N$th root gives the triangle
inequality. Finally, changing $(o,\kappa)$ to $(o',\kappa')$
for the same model function translates its entire fold by
$o'-o+\mathsf l(\tilde\kappa'-\tilde\kappa)\in a\mathbb Z^3$.
The Wilson integral is translation invariant.
\end{proof}
\section{Periodic Gibbs states}\label{app:periodic-gibbs}

\subsection{Infinite-volume lattice states}\label{sec:infinite-volume-states}

We first take the volume limit at a fixed spacing. The configuration
space is compact, but the finite periodic measures do not themselves
satisfy the infinite-lattice Gibbs equations. We therefore check those
equations on each fixed set of bonds before passing to a limit.

Let $a>0$, let $\mathbb L_a=a\mathbb Z^3$, and let $\mathsf{Bd}$
be its positively oriented nearest-neighbour bonds. A configuration
is an element of
\[
 \mathcal V(\mathbb L_a)=\operatorname{SU}(2)^{\mathsf{Bd}},
 \qquad U(\bar b)=U(b)^{-1}.
\]
We use the product topology. This is a compact metrizable space.
For $B\subset\mathsf{Bd}$, write $\mathcal F(B)$ for the bounded
Borel functions depending only on $U_B$, and $\mathcal C_{\rm loc}$
for the continuous functions depending on finitely many bonds.
For a spatial set $Q$, $\mathsf{Bd}(Q)$ consists of bonds whose
closed segments lie in $Q$. The notation
$\mathcal C_{\rm loc}(Q)$ requires a finite bond support in
$\mathsf{Bd}(Q)$.

Enumerating the bonds gives a product metric with summable coordinate
weights. Replacing all but the first $N$ coordinates by the identity
approximates the identity map uniformly in this metric. Uniform
continuity therefore shows that $\mathcal C_{\rm loc}$ is uniformly
dense in $C(\mathcal V(\mathbb L_a))$. In particular, its integrals
determine a Borel probability measure. We call such a measure a law.
The laws form a compact metrizable space in the weak topology.

For a plaquette $p$, retain the notation
$a_p(U)=1-\frac12\operatorname{tr}\partial U(p)\in[0,2]$.
For bounded compactly supported $f$ on $\mathbb R^3$, put
\begin{equation}\label{support:13-6}
 \mathcal O_a(f)=a^{-1}\sum_{p\in\mathsf{Pl}}f(x(p))a_p.
\end{equation}
The sum is finite, so $\mathcal O_a(f)\in\mathcal C_{\rm loc}$
even when $f$ is merely bounded. For a law $\nu$ define
\begin{equation}\label{support:13-23}
 S_n^T[\nu](f_1,\ldots,f_n)
 =\langle\mathcal O_a(f_1);\ldots;\mathcal O_a(f_n)\rangle_\nu.
\end{equation}
If $X=(f_1,\ldots,f_n)$ and $\rho,\nu$ are two laws, write
\begin{equation}\label{support:13-25}
 X_a^\nu=\prod_{q=1}^n
       \bigl(\mathcal O_a(f_q)-\nu(\mathcal O_a(f_q))\bigr),
 \qquad \sigma_{\rho,\nu}(X)=\rho(X_a^\nu),
 \qquad \sigma_\nu=\sigma_{\nu,\nu}.
\end{equation}
The empty product is one. The first index specifies the integrating
law, the second the centring law. These expressions are symmetric
and multilinear in the entries. Lattice isometries
$\mathsf r\in a\mathbb Z^3\rtimes W_3$ act by pullback on
configurations and functions, with inversion when a bond orientation
is reversed. Plaquette traces are unchanged by cyclic permutation
and inversion, hence
\[
 \mathsf r\mathcal O_a(f)=\mathcal O_a(\mathsf rf),\qquad
 \sigma_{\mathsf r_*\rho,\mathsf r_*\nu}(\mathsf rX)
 =\sigma_{\rho,\nu}(X).
\]
The same covariance holds for the cumulants.

\subsubsection{Gibbs equations and periodic limits}

A period lattice is a finite-index subgroup $\Gamma\subset a\mathbb Z^3$
with $\operatorname{sys}(\Gamma):=\min_{0\ne v\in\Gamma}|v|_\infty
\ge3a$. Choose one representative of every positive bond class
modulo $\Gamma$. Independent Haar variables at these representatives,
extended periodically, define $dU_\Gamma$ on the closed space
$\mathcal V_\Gamma$ of periodic configurations. Changing representatives
only permutes the variables. If $D$ contains one site of each class,
then the plaquettes based at sites of $D$ represent the plaquette
classes once. Thus
\begin{equation}\label{support:13-30}
 S_\Gamma=\sum_{p\in\mathsf{Pl}/\Gamma}a_p,
 \quad Z_\Gamma[\Phi]=\int e^{-\beta S_\Gamma+\Phi}\,dU_\Gamma,
 \quad d\mu_\Gamma^\Phi=Z_\Gamma[\Phi]^{-1}
                        e^{-\beta S_\Gamma+\Phi}\,dU_\Gamma
\end{equation}
are independent of representatives; here $\beta\ge0$ and
$\Phi\in C(\mathcal V_\Gamma)$ is real. Write $\mu_{\Gamma,\beta}$
when $\Phi=0$. The action and Haar measure show that
$\mathsf r_*\mu_\Gamma^\Phi=\mu_\Gamma^{\mathsf r\Phi}$ whenever
$\mathsf r^\circ\Gamma=\Gamma$.

For finite $B\subset\mathsf{Bd}$, let $\mathsf{Pl}[B]$ be the
plaquettes containing a bond of $B$, and set
\begin{equation}\label{support:14-7}
 H_B=-\beta\sum_{p\in\mathsf{Pl}[B]}a_p.
\end{equation}
Define the kernel by
\begin{equation}\label{alt:eq:dlr-kernel}
 (\gamma_BF)(U)=
 \frac{\int F(V\vee U_{B^c})e^{H_B(V\vee U_{B^c})}\,dV_B}
      {\int e^{H_B(V\vee U_{B^c})}\,dV_B}.
\end{equation}
Every bond is in four plaquettes, so $\#\mathsf{Pl}[B]\le4\#B$.
The denominator lies in $[e^{-2\beta\#\mathsf{Pl}[B]},1]$.
The kernel is positive, preserves one and contracts the supremum
norm. Integration over the finite set $B$ preserves continuity;
if $F$ is local, the remaining support lies in the union of its
support and the bonds of $\mathsf{Pl}[B]$, with $B$ removed.
Thus $\gamma_B\mathcal C_{\rm loc}\subset\mathcal C_{\rm loc}$.

\begin{definition}[Gibbs states and invariance]\label{support:14-9}
\begin{enumerate}
\item[(1)] The Gibbs equation on a finite bond set $B$ is
\begin{equation}\label{alt:eq:dlr}
 \nu(\gamma_BF)=\nu(F)\qquad(F\in\mathcal C_{\rm loc}).
\end{equation}
It then holds for every bounded Borel $F$, since the two probability measures agree on continuous local tests.
\item[(2)] A Gibbs state satisfies this equation for every finite $B$. Write $\mathcal G_{a,\beta}$ for all such states.
\item[(3)] A state is periodic if it is invariant under a finite-index subgroup of $a\mathbb Z^3$. Equivalently it is invariant under $Pa\mathbb Z^3$ for some period number $P\in\mathbb N$, because the orders of the coordinate classes in the finite quotient have a common multiple. The class of periodic Gibbs states is $\mathcal G^{\rm per}_{a,\beta}$.
\item[(4)] For $a=a_K$, $\mathcal G^{\rm tor}_{K,\beta}$ consists of the weak limits of periodically extended cubic torus laws along subsequences of sides $2L^m\to\infty$.
\item[(5)] For a translation-invariant state let $\mathcal I$ be the sigma field of translation-invariant events modulo null sets. Such a state is ergodic when $\mathcal I$ is trivial.
\end{enumerate}
\end{definition}

We also need locally finite perturbations. Let
$\mathfrak A=(A_i)_{i\in I}$ be real local continuous functions,
each with a nonempty finite support $B_i$, such that each bond
lies in only finitely many $B_i$, counting indices. Replace $H_B$
in \eqref{support:14-7} by
\[
 H_B^{\mathfrak A}=H_B+\sum_{i:B_i\cap B\ne\emptyset}A_i
\]
to define $\gamma_B^{\mathfrak A}$ and the corresponding Gibbs
equations.

\begin{proposition}\label{support:14-14}
The Gibbs states form a convex weakly compact set. The subset invariant
under any specified subgroup of lattice translations is also convex
and weakly compact. Lattice isometries carry Gibbs states to Gibbs
states and preserve period numbers. If $t_j\to t$ and $\nu_j\to\nu$
weakly, with $\nu_j$ Gibbs for $t_j\mathfrak A$, then $\nu$ is Gibbs
for $t\mathfrak A$.
\end{proposition}
\begin{proof}
Each Gibbs equation is a closed affine condition on laws, since
$\gamma_BF-F$ is continuous. Invariance is likewise expressed by
$\nu(F\circ\mathsf r-F)=0$. Intersect these conditions in the compact
space of laws. Relabelling bonds, and inverting the reversed ones,
preserves product Haar measure and gives
$\gamma_{\mathsf rB}(\mathsf rF)=\mathsf r(\gamma_BF)$.

For the last assertion put
$A_B=\sum_{i:B_i\cap B\ne\emptyset}A_i$ and
$\delta_j=|t_j-t|\|A_B\|_\infty$.
The ratio of the two normalized kernel densities lies between
$e^{-2\delta_j}$ and $e^{2\delta_j}$. Consequently
\[
 \|\gamma_B^{t_j\mathfrak A}F-
       \gamma_B^{t\mathfrak A}F\|_\infty
 \le(e^{2\delta_j}-1)\|F\|_\infty\longrightarrow0.
\]
The fixed kernel applied to $F\in\mathcal C_{\rm loc}$ is continuous.
Pass to the limit in the equations for $\nu_j$.
\end{proof}

\begin{proposition}[Periodic approximations]\label{support:14-16}
\begin{enumerate}
\item[(a)] Let $\Gamma_j$ be period lattices with
$\operatorname{sys}(\Gamma_j)\to\infty$. Suppose that
$\Gamma_j$ acts on the indices of $\mathfrak A$, with
$A_{v\cdot i}=\tau_vA_i$ and $B_{v\cdot i}=B_i+v$.
Put $\Phi_j=\sum_{[i]\in I/\Gamma_j}A_i$ on $\mathcal V_{\Gamma_j}$.
Every weak limit of $\mu_{\Gamma_j}^{\Phi_j}$ is Gibbs for
$\mathfrak A$.
\item[(b)] Limits of $\mu_{\Gamma_j,\beta}$ are translation-invariant Gibbs states.
\item[(c)] At $a=a_K$, let $\widetilde\mu_{K,\ell,\beta}$ be the periodic
lift of the cubic torus measure and define
\[
 \mathcal G^{\rm tor}_{K,\beta}
 =\bigcap_{m_0\ge1}
   \overline{\{\widetilde\mu_{K,2L^m,\beta}:m\ge m_0\}}.
\]
This is a nonempty compact subset of $\mathcal G^{\rm per}_{K,\beta}$.
Every member is invariant under $a_K\mathbb Z^3\rtimes W_3$ and
has period number one.
\end{enumerate}
\end{proposition}
\begin{proof}
The index action is free: $B_i+v=B_i$ with $B_i$ finite nonempty
forces $v=0$. Every orbit has an index whose support meets a fixed
bond transversal, so local finiteness makes $I/\Gamma_j$ finite.
Periodicity makes $\Phi_j$ independent of the representatives.

Fix a finite $B$ and a local continuous $F$. Choose a closed cube
$Q$ containing its bond support, all plaquettes of $\mathsf{Pl}[B]$,
and all $B_i$ meeting $B$. For large $j$, its side is smaller than
$\operatorname{sys}(\Gamma_j)$. Distinct bonds and plaquettes in
$Q$ are then incongruent modulo $\Gamma_j$. Choose a bond
transversal containing $\mathsf{Bd}(Q)$ and write its variables as
$(V_B,W)$. In these variables,
\[
 -\beta S_{\Gamma_j}+\Phi_j
       =H_B^{\mathfrak A}+H',
\]
where $H'$ is independent of $V_B$. Indeed, a remaining plaquette
class, or perturbation orbit, depending on a bond of $B$ would
have a representative already included in $H_B^{\mathfrak A}$.
For functions supported in $Q$, replacing $V_B$ in a periodic
configuration agrees on $Q$ with the nonperiodic replacement in
the definition of $\gamma_B^{\mathfrak A}$. Integrating first over
$V_B$ therefore gives
\[
 \mu_{\Gamma_j}^{\Phi_j}(\gamma_B^{\mathfrak A}F)
       =\mu_{\Gamma_j}^{\Phi_j}(F).
\]
Both functions are continuous and independent of $j$. The equality
passes to every weak limit.

For the empty perturbation family, periodic Haar measure and the
action are invariant under every lattice translation. Invariance
passes to weak limits. For the cubic sequence, the same is true
of every signed coordinate permutation. Its nested tail closures
are nonempty compact sets in the space of laws, so their intersection
is nonempty and compact. Metrizability identifies this intersection
with the subsequential limits. These limits are Gibbs by the first
part, proving all assertions.
\end{proof}

\begin{lemma}[Gauge invariance]\label{support:14-13}
Every Gibbs state is invariant under every gauge transformation.
Two Gibbs states agreeing on gauge-invariant local continuous
functions are equal.
\end{lemma}
\begin{proof}
First let a gauge transformation differ from the identity only at
one site $x$. In the kernel for the six incident bonds, it is a
left or right translation of each Haar variable and preserves
every plaquette trace. Thus $\gamma_{B_x}(F\circ u)=\gamma_{B_x}F$,
and the Gibbs equation gives invariance. Compose these changes
for finitely many sites. On a local function, an arbitrary gauge
transformation uses only its values at the endpoints of a finite
bond support; hence invariance holds there and therefore for the law.

Averaging a local continuous $F$ over the gauge variables at those
endpoints gives a gauge-invariant local continuous $F^{\rm inv}$
with the same integral in every Gibbs state. Agreement on such
averages implies agreement on $\mathcal C_{\rm loc}$, hence equality
of the laws.
\end{proof}

\subsubsection{Reflection positivity for adapted periods}

The diagonal reflections require period lattices that split in the normal direction.
\begin{definition}[Lattice mirrors]\label{support:54-1}
\begin{enumerate}
\item[(1)] Fix $a>0$ and $\beta\ge0$.
\item[(2)] Let
\[
 \mathfrak N=\{e_\mu:1\le\mu\le3\}
   \cup\{e_\mu\pm e_\nu:\mu<\nu\},
 \qquad \Sigma(n,c)=\{x:n\cdot x=c\}.
\]
The admissible heights are $\operatorname{Adm}(n)=\frac a2\mathbb Z$ for coordinate normals and $\operatorname{Adm}(n)=a\mathbb Z$ for diagonal normals. These are exactly the heights in the respective families at which $\theta_{n,c}x=x-2(n\cdot x-c)n/|n|^2$ preserves the lattice: the linear reflection is a signed permutation and its translation part is $2cn/|n|^2$. Opposite normals give the same mirrors upon negating the height. The coordinate mirrors are site mirrors when $c\in a\mathbb Z$ and bond mirrors otherwise.
\end{enumerate}
\end{definition}

A period lattice is adapted to $n$ if, for some $h>0$,
\begin{equation}\label{support:15-23}
 \Gamma=\Gamma_\perp\oplus\mathbb Zv_\Gamma,
 \qquad \Gamma_\perp=\Gamma\cap n^\perp,
 \qquad v_\Gamma=h n/|n|^2.
\end{equation}
Then $n\cdot\Gamma=h\mathbb Z$, and the linear reflection preserves
$\Gamma$. If $n$ is coordinate, $h\in a\mathbb N$ and $h\ge3a$;
if $n$ is diagonal, $h\in2a\mathbb N$ and $h/2\ge3a$.
Thus $c+h/2$ is also admissible. The reflections in $c$ and $c+h/2$
differ by translation through $v_\Gamma$ and act identically on
$\mathcal V_\Gamma$.

\begin{proposition}[Two-mirror positivity]\label{support:16-26}
Let $a>0$, $\beta\ge0$, $n\in\mathfrak N$, and let $\Gamma$
be a period lattice adapted to $n$, with height period $h$.
For admissible $c$, put $\theta=\theta_{n,c}$ and
$\mathbb S^+=\{c\le n\cdot x\le c+h/2\}$. Then
\begin{equation}\label{support:16-27}
 \mu_{\Gamma,\beta}(\overline{\theta F}F)\ge0
 \qquad\bigl(F\in\mathcal F(\mathsf{Bd}(\mathbb S^+))\bigr).
\end{equation}
The law is $\theta$-invariant. Both assertions hold also for
$\mu_\Gamma^\Phi$ if
$\Phi=(\Phi_++\theta\Phi_+)|_{\mathcal V_\Gamma}$ with
$\Phi_+$ real continuous and dependent only on the bonds of
$\mathbb S^+$. No gauge-invariance assumption on $F$ is imposed;
bonds lying in either mirror are included.
\end{proposition}
\begin{proof}
We retain the boundary variables instead of integrating them into
either half. Write $t(x)=n\cdot x-c$. Modulo $\Gamma_\perp$, choose
representatives $R^0$ of the bonds lying in $t=0$ or $t=h/2$,
$R^+$ of the other bonds contained in $0\le t\le h/2$, and
$R^\times$ of the bonds crossing either boundary. The remaining
representatives are the reflected bonds $(\theta b)^+$,
$b\in R^+$. For a bond or plaquette, use its midpoint or
barycentre to measure height. Translation by $jv_\Gamma$ changes
this height by $jh$. Every object has a representative in the
two half-slabs, their boundaries, or the crossing sets, with height
in $(-h/2,h/2]$. Two such representatives congruent modulo $\Gamma$
differ by an element of $\Gamma_\perp$, since their height difference
has modulus less than $h$. This proves that the chosen representatives
are disjoint and exhaustive. There are finitely many of them.

Consequently the coordinates
\[
 U_0=U_{R^0},\quad U_+=U_{R^+},\quad
 U'_+=(\theta U)_{R^+},\quad U_\times=U_{R^\times}
\]
are independent Haar variables. The necessary inversions in
$U'_+$ preserve Haar measure. On the boundary bonds $\theta$ fixes
the coordinates modulo $\Gamma$, and it interchanges $U_+$ and
$U'_+$. Thus $F(U)=\widehat F(U_0,U_+)$ and
$(\theta F)(U)=\widehat F(U_0,U'_+)$.

Choose plaquette representatives in the same way. Those lying in
the boundaries contribute a weight $W_0(U_0)>0$. Those in the
positive half but not in its boundaries contribute $W_+(U_0,U_+)>0$,
and their reflected copies contribute $W_+(U_0,U'_+)$.
The remaining plaquettes meet both sides of a mirror. Their form
follows by listing the corner heights
\[
 t(x),\quad t(x)+an_\mu,\quad
 t(x)+a(n_\mu+n_\nu),\quad t(x)+an_\nu.
\]
A site mirror has none. At a diagonal mirror the only possibility
has two opposite corners $z,z'$ in the plane and two reflected
corners off it, at heights $\pm a$ relative to the plane.
With $y$ the corner in $\mathbb S^+$, set
$v_p(U)=U(\langle z,y\rangle)U(\langle y,z'\rangle)$.
These two bonds belong to the positive half and not to its boundary;
the inequality $a<h/2$ ensures that they do not reach the other mirror.
Tracing the plaquette boundary gives
\[
 e^{-\beta a_p(U)}=k_\beta(v_p(U_+)v_p(U'_+)^{-1}),
 \qquad k_\beta(g)=e^{-\beta}e^{(\beta/2)\operatorname{tr}g}.
\]

At a bond mirror, let $x,y$ be neighbouring sites in the adjacent
layer of $\mathbb S^+$. Put
$u_b=U(\langle x,y\rangle)$,
$u'_b=(\theta U)(\langle x,y\rangle)$ and
$w_x=U(\langle x,\theta^{(k)}x\rangle)$, where $\theta^{(k)}$
is reflection in the boundary under consideration. The four-bond
loop gives
\[
 e^{-\beta a_p(U)}
   =k_\beta(u_bw_y(u'_b)^{-1}w_x^{-1}).
\]
The crossing variables at the two mirrors are disjoint: their
midpoint heights differ by $h/2$, not by a multiple of $h$.
Their Haar integral at each boundary is therefore
\[
 K(u,u')=\int\prod_{b=(x,y)}
             k_\beta\bigl(u_b(w\cdot u')_b^{-1}\bigr)\,dw,
 \qquad (w\cdot u')_b=w_xu'_bw_y^{-1}.
\]

Each resulting kernel has a uniformly absolutely convergent expansion
$\sum_\alpha c_\alpha\Phi_\alpha(V)
\overline{\Phi_\alpha(V')}$ with
$c_\alpha\ge0$ and
$\sum_\alpha c_\alpha\|\Phi_\alpha\|_\infty^2<\infty$.
Indeed, the matrix-entry expansion used in \ref{support:3-16} gives
\[
 k_\beta(gg'^{-1})
 =e^{-\beta}\sum_{N\ge0}\frac{(\beta/2)^N}{N!}
   \sum_{I,J\in\{1,2\}^N}
        g_{IJ}\overline{g'_{IJ}},
 \qquad \sum_{N,I,J}e^{-\beta}\frac{(\beta/2)^N}{N!}=e^\beta.
\]
Products and pullbacks preserve this property. For a bond boundary,
simultaneous gauge invariance and Haar invariance replace the single
average in $K$ by independent averages on its two entries. The
expansion then has features
$\Phi_\alpha^{\rm av}(u)=\int\Phi_\alpha(w\cdot u)\,dw$;
their supremum norms do not increase. This proves the assertion
also for the product $\mathcal K$ of the two boundary kernels.

After integrating $U_\times$, the full Wilson weight is
\[
 W_0(U_0)W_+(U_0,U_+)W_+(U_0,U'_+)\mathcal K(U_+,U'_+).
\]
Termwise integration is justified by the summable supremum bounds
and gives the unnormalized quadratic form
\[
 \int dU_0\,W_0(U_0)
   \sum_\alpha c_\alpha
    \left|\int \widehat F(U_0,U_+)W_+(U_0,U_+)
                    \Phi_\alpha(U_+)\,dU_+\right|^2\ge0.
\]
Division by $Z_\Gamma[0]>0$ proves \eqref{support:16-27}.
For the perturbed measure apply this calculation to $Fe^{\Phi_+}$.
Finally $\theta$ preserves Haar measure and the action, and the
stated perturbation is $\theta$-invariant, proving invariance.
\end{proof}

\subsubsection{The torus-limit states}

Periodization $f^{\langle\ell\rangle}(x)=
\sum_{v\in\ell\mathbb Z^3}f(\widetilde x+v)$ is well defined
for compactly supported $f$. Grouping the plaquettes by period class
gives the exact identity
\begin{equation}\label{support:17-14}
 \mathcal O_K(f)(\operatorname{per}_\ell U)
       =\mathcal O_{a_K}(f^{\langle\ell\rangle})(U).
\end{equation}
If the supports lie in $(-\ell/4,\ell/4)^3$, periodization preserves
their stated separation and per-entry radius bounds, and their
$\|\cdot\|_\sharp$ norms. To check the distance assertion, each
coordinate difference of two points in this cube has modulus
less than $\ell/2$, so it is already the shortest representative
of its class modulo $\ell$. The same argument and local lifting
give the supremum and Lipschitz norm identities. Periodization
commutes with the lattice isometries whenever the source and its
image are considered on the same quotient.

\begin{proposition}\label{support:17-45}
Let $K\ge1$, $\beta\ge0$ and $\nu\in\mathcal G^{\rm tor}_{K,\beta}$.
There are sides $\ell_j\to\infty$ with
$\widetilde\mu_{K,\ell_j,\beta}\to\nu$ weakly.
Along every such sequence, for bounded compactly supported test
functions and $w=\sum_qz_qf_q$,
\begin{gather}
 \nu(e^{\mathcal O_K(w)})
   =\lim_j\frac{Z_{K,\ell_j,\beta}[w^{\langle\ell_j\rangle}]}
                         {Z_{K,\ell_j,\beta}[0]},\label{support:17-46}\\
 S_n^T[\nu](X)
   =\lim_j S^T_{n;K,\ell_j,\beta}(X^{\langle\ell_j\rangle}),
 \qquad
 \sigma_\nu(X)
   =\lim_j\sigma_{K,\ell_j,\beta}(X^{\langle\ell_j\rangle}).
 \label{support:17-47}
\end{gather}
Any bound on the right uniform for tuples in
$(-\ell_j/4,\ell_j/4)^3$ with prescribed separation and per-entry
radius passes to the left for tuples anywhere in $\mathbb R^3$.
The state $\nu$ is reflection positive at every coordinate site
and bond mirror, for bounded local functions of the closed positive
half-space.
\end{proposition}
\begin{proof}
The sequence exists by the tail-closure description in
\ref{support:14-16}. Every product of the observables and their
exponentials is local continuous. Weak convergence and
\eqref{support:17-14} therefore give convergence of all moments and
the first formula. Cumulants and centred moments are fixed finite
polynomials in these moments, proving \eqref{support:17-47}.
A fixed compact union of supports eventually lies in the central
cube, so each uniform bound passes to the limit.

The cubic period lattice is adapted to every coordinate normal,
with height period $\ell_j$. A fixed finite bond support in
$\{x_\mu\ge c\}$ eventually lies in
$\{c\le x_\mu\le c+\ell_j/2\}$. Apply \ref{support:16-26}
and pass to the weak limit for continuous local $F$.
To extend the resulting inequality to bounded Borel $F$ with the
same finite support, approximate its function of the finite bond
coordinates in $L^2$ of the corresponding marginal of $\nu$ by
continuous functions. Reflection invariance gives the same
$L^2$ convergence for their reflected functions. Cauchy--Schwarz
then passes the quadratic form to the limit. Reversing the side
follows by reflection invariance.
\end{proof}

The last assertion does not supply diagonal positivity. For
$n=e_\mu\pm e_\nu$, the heights of $\ell\mathbb Z^3$ are
$\ell\mathbb Z$, whereas its shortest nonzero pure-normal vector
is $\ell n$, of height $2\ell$. Thus the cubic period lattice
does not satisfy \eqref{support:15-23} for a diagonal normal.


\subsubsection{Periodic arrays and their reflection}

The comparison of Gibbs states below uses two finite-volume
inequalities. The first compares a periodic array with its expectation
in a Gibbs state; the second increases its partition function by
reflecting the array across successive slabs.

For a site transversal $D$ of $\Gamma$, write
$\mathsf{Bd}_D$ for the positive bonds starting in $D$, and let
\[
 \partial_rD=\{x\in D:x+w\notin D\text{ for some }w\in a\mathbb Z^3,
                                      |w|_\infty\le ra\}.
\]
If $\Gamma\subset\Gamma'\subset a\mathbb Z^3$ have finite index,
the array of a local function $A$ is
$\mathbb A=\sum_{v\in\Gamma'/\Gamma}\tau_vA$ on
$\mathcal V_\Gamma$, with $|\mathbb A|=[\Gamma':\Gamma]$.

\begin{lemma}[Lower tangent]\label{support:14-20}
Let $\Gamma$ be a period lattice, $\Gamma'\supset\Gamma$ a
finite-index subgroup, $D$ a site transversal of $\Gamma$, and
$\nu\in\mathcal G_{a,\beta}$ invariant under $\Gamma'$.
Let $A\in\mathcal C_{\rm loc}$ be real and $t\in\mathbb R$.
If $N_{\rm in}$ is the number of classes in $\Gamma'/\Gamma$
having a representative $v$ for which $\tau_vA$ has a support
in $\mathsf{Bd}_D$, then
\begin{equation}\label{support:14-21}
 \log\frac{Z_\Gamma[t\mathbb A]}{Z_\Gamma[0]}
 \ge t|\mathbb A|\nu(A)-24\beta\#\partial_1D
       -2|t|\|A\|_\infty(|\mathbb A|-N_{\rm in}).
\end{equation}
If $A$ has a bond support in $x_0+[0,ra]^3$, with
$x_0\in a\mathbb Z^3$ and integer $r\ge1$, then
\[
 |\mathbb A|-N_{\rm in}\le\#\partial_rD
          \le(2r-1)^3\#\partial_1D.
\]
\end{lemma}
\begin{proof}
Put $B=\mathsf{Bd}_D$. Plaquettes whose four bonds lie in $B$
give identical action terms for $V\vee U_{B^c}$ and the periodic
extension $V^{\rm per}$. They represent distinct period classes.
Let $n_1$ count the remaining plaquettes meeting $B$, and $n_2$
the period classes not represented by an interior plaquette.
Then
\[
 \sum_{p\in\mathsf{Pl}[B]}a_p(V\vee U_{B^c})
        -S_\Gamma(V^{\rm per})=\Sigma_1-\Sigma_2,
 \qquad 0\le\Sigma_i\le2n_i.
\]
If a positive bond starts at $x\notin\partial_1D$, all plaquettes
containing it have their bonds in $B$. There are at most nine
plaquettes meeting the three positive bonds starting at a specified
$x$, and three plaquettes based at $x$. Hence
$n_1\le9\#\partial_1D$ and $n_2\le3\#\partial_1D$.

Choose one good representative for each class counted by $N_{\rm in}$,
and set $\Phi_B=t\sum_{\rm good}\tau_vA$ and
$e=|t|\|A\|_\infty(|\mathbb A|-N_{\rm in})$.
The preceding action comparison, applied to numerator and denominator,
gives for every exterior configuration
\[
 \frac{Z_\Gamma[\Phi_B]}{Z_\Gamma[0]}
       \ge e^{-2\beta(n_1+n_2)}\gamma_B(e^{\Phi_B}).
\]
Jensen's inequality for the normalized kernel, followed by the Gibbs
equation and $\Gamma'$-invariance, yields
\[
 \log\frac{Z_\Gamma[\Phi_B]}{Z_\Gamma[0]}
       \ge tN_{\rm in}\nu(A)-2\beta(n_1+n_2).
\]
On periodic configurations $|t\mathbb A-\Phi_B|\le e$;
also $|t(|\mathbb A|-N_{\rm in})\nu(A)|\le e$.
The two errors give \eqref{support:14-21}.

For the last estimate, assign each array class the unique point of
$x_0+v+\Gamma$ in $D$. If that point is outside $\partial_rD$,
the translated cube and all its bonds lie in $\mathsf{Bd}_D$.
Different bad classes therefore give different points of
$\partial_rD$. Finally, from $x\in\partial_rD$ follow a lattice path
to a point outside $D$, moving at most $a$ in each coordinate at each
step and using at most $r$ steps. The point immediately before its first exit from $D$
lies in $\partial_1D$ and differs from $x$ by a vector of sup norm
at most $(r-1)a$. Thus $\partial_rD$ is contained in the union of
$(2r-1)^3$ translates of $\partial_1D$.
\end{proof}

Fix admissible $c_0$, $h>0$ with $v_h=hn/|n|^2\in a\mathbb Z^3$,
and put
\[
 \Sigma_j=\Sigma(n,c_0+jh),\quad\theta_j=\theta_{n,c_0+jh},
 \quad\mathbb S_j=\{(j-1)h\le n\cdot x-c_0\le jh\}.
\]
The maps
$\mathsf w_{2k+1}=\tau_{2kv_h}$ and $\mathsf w_{2k}=\theta_k$
carry $\mathbb S_1$ to $\mathbb S_j$. Define
$f^{[j]}=\mathsf w_jf$ for odd $j$ and
$f^{[j]}=\overline{\mathsf w_jf}$ for even $j$.
Writing $\vartheta=\theta_{n,c_0+h/2}$, the reflection formulas give
\begin{equation}\label{support:15-31}
 \theta_m\mathsf w_j=\mathsf w_{2m+1-j},\quad
 \mathsf w_{j+2N}=\tau_{2Nv_h}\mathsf w_j,\quad
 \tau_{v_h}\mathsf w_j=\mathsf w_{j+1}\vartheta.
\end{equation}
All these maps preserve the lattice. The group generated by the
$\theta_j$ consists of them and the translations $\tau_{2kv_h}$.

\begin{proposition}[Reflection of arrays]\label{support:16-31}
Suppose $\Gamma$ is adapted to $n$ with height period $2Nh$,
$N\ge1$. If a law $\mu$ on $\mathcal V_\Gamma$ is reflection
positive at every pair $\Sigma_j,\Sigma_{j+N}$, then
\begin{equation}\label{support:16-33}
 \left|\mu\left(\prod_{j=1}^{2N}f_j^{[j]}\right)\right|
 \le\prod_{j=1}^{2N}
       \mu\left(\prod_{k=1}^{2N}f_j^{[k]}\right)^{1/(2N)}
\end{equation}
for bounded Borel functions of the bonds in $\mathbb S_1$;
each quantity under a root is nonnegative.

Suppose additionally that $\Gamma'\supset\Gamma$ has finite index
and the splitting \eqref{support:15-23} with height period $h$.
No lower bound on its systole is required. For real
$A\in\mathcal C_{\rm loc}(\mathbb S_1)$,
let $T$ represent $(\Gamma'\cap n^\perp)/(\Gamma\cap n^\perp)$,
and set
\[
 A_T=\sum_{v\in T}\tau_vA,\qquad
 \mathbb A=\sum_{j=1}^{2N}\tau_{(j-1)v_h}A_T,\qquad
 \mathbb A^{\rm f}=\sum_{j=1}^{2N}\mathsf w_jA_T.
\]
These define functions on $\mathcal V_\Gamma$ independent of
representatives, and $|\mathbb A|=2N\#T$. For every real $t$,
\begin{equation}\label{support:16-34}
 Z_\Gamma[t\mathbb A]\le Z_\Gamma[t\mathbb A^{\rm f}].
\end{equation}
The law $\mu_\Gamma^{t\mathbb A^{\rm f}}$ is reflection positive
at each pair $\Sigma_j,\Sigma_{j+N}$, invariant under the wall
group and under $\Gamma'\cap n^\perp$, and satisfies
\begin{equation}\label{support:16-35}
 \mu_\Gamma^{t\mathbb A^{\rm f}}(\mathbb A^{\rm f})
       =|\mathbb A|\mu_\Gamma^{t\mathbb A^{\rm f}}(A).
\end{equation}
\end{proposition}
\begin{proof}
Let $\mathrm{cy}(f_1,\ldots,f_{2N})$ denote the integral on the
left of \eqref{support:16-33}. The copies are cyclic modulo $2N$ on
$\mathcal V_\Gamma$. Reflection in $\Sigma_m$ changes index $j$
to $2m+1-j$ and conjugates the copy. Cauchy--Schwarz for the
positive form at that pair therefore bounds $|\mathrm{cy}|^2$
by the product of the two integrals obtained by reflecting either
half of the cyclic array into the other half. Those two integrals
are nonnegative. In particular constant arrays have nonnegative
integrals.

Assume first that the constant-array roots $\|f_j\|$ are positive.
Among the finitely many arrays with entries chosen from
$f_1,\ldots,f_{2N}$ maximize
$|\mathrm{cy}(\underline f)|/\prod_j\|f_j\|$.
The maximum is at least one and is unchanged by either half-reflection:
the product bound and maximality force both reflected arrays to
attain it. Choose a maximizing array with a longest constant cyclic
string, of length $k$. Reflect across its terminal boundary.
If $k\le N$, this doubles the string; if $N<k<2N$, it makes the
whole array constant. Maximality of $k$ therefore forces $k=2N$,
so the maximum is one. If a root vanishes, replace every $f_j$ by
$f_j+\varepsilon$ with real $\varepsilon$. Each constant-array
integral is a nonnegative polynomial of degree $2N$ with leading
coefficient one. Choose $\varepsilon\to0$ avoiding their finitely
many zeros and pass to the limit. This proves \eqref{support:16-33}.

For the array comparison take $f=e^{tA_T}$. The splitting of both
period groups gives the displayed expression for $\mathbb A$.
For odd $j$, $\tau_{(j-1)v_h}=\mathsf w_j$; for even $j$ it is
$\mathsf w_j\vartheta$. Thus
\[
 Z_\Gamma[t\mathbb A]/Z_\Gamma[0]
       =\mathrm{cy}(f,\vartheta f,\ldots,f,\vartheta f),
 \quad
 Z_\Gamma[t\mathbb A^{\rm f}]/Z_\Gamma[0]
       =\mathrm{cy}(f,\ldots,f).
\]
The unperturbed measure has the required positivity by
\ref{support:16-26}. Translation by $v_h$ takes the fold of $f$ to the
fold of $\vartheta f$, by \eqref{support:15-31}; translation invariance
therefore gives $\|\vartheta f\|=\|f\|$. Inequality
\eqref{support:16-33} proves \eqref{support:16-34}.

For any $m$, the folded perturbation is $\Phi_++\theta_m\Phi_+$,
where $\Phi_+=t\sum_{j=m+1}^{m+N}\mathsf w_jA_T$.
Apply the perturbed assertion of \ref{support:16-26}. This gives
positivity and invariance under every wall reflection.
Translations in $\Gamma'\cap n^\perp$ permute the transverse
classes and commute with the $\mathsf w_j$, so they preserve the
perturbation and the law. All $2N\#T$ summands in
$\mathbb A^{\rm f}$ consequently have the same expectation,
which proves \eqref{support:16-35}.
\end{proof}

\subsubsection{Domination of periodic Gibbs states}

\begin{definition}[Infinite-lattice boards]\label{support:18-2}
Fix $K\ge1$, $a=a_K=L^{-K}$ and $\beta>0$. For each integer $i\le K-1$ and offset $o\in a\mathbb Z^3$, the board $(i,o)$ has cells
\[
 \square_k=o+l k+[-l/2,l/2]^3,
 \qquad l=a_i=L^{-i},\qquad k\in\mathbb Z^3.
\]
The space $\mathcal A(\square_k)$ consists of the bounded Borel functions of the inner bonds of the cell. Mirror copies are defined by \eqref{support:2-18} and \eqref{support:2-26} with indices in $\mathbb Z^3$. The reflections in the walls generate the wall group of the board.
\end{definition}

\begin{lemma}[Cell properties]\label{support:18-7}
\begin{enumerate}
\item[(a)] The walls of $(i,o)$ are bond-midpoint mirrors. Each site belongs to exactly one cell and lies in its interior. The inner bond sets of distinct cells are disjoint.
\item[(b)] A plaquette whose barycentre is in the interior of $\square_k$ has all its bonds in $\square_k$. If $f$ vanishes outside the interior of $\square_k$, the observable $\mathcal O_K(f)$ is continuous and belongs to $\mathcal A(\square_k)$.
\item[(c)] The site chart identifies the inner bond coordinates with those of the integer cube $\mathsf C_{l/a}$ from \eqref{support:2-8}. It is $n\mapsto o+lk+an$.
\item[(d)] The mirror maps \eqref{support:2-18} are lattice isometries and carry the bond coordinates of one cell to those of the other, using inverse variables at reversed bonds.
\item[(e)] The copies preserve products, sums and sup norms and commute with conjugation. Copies compose as in \eqref{support:2-30}. An even copy is the pullback by its mirror map, while an odd copy also conjugates the function.
\end{enumerate}
\end{lemma}
\begin{proof}
The side $l/a=L^{K-i}$ is odd, so $o_\mu+l(t+1/2)$ is a half-integer multiple of $a$. Consecutive wall positions in each coordinate enclose exactly $l/a$ sites. A plaquette can cross a wall only along one of its varying coordinates, in which case its barycentre in that coordinate is the wall midpoint. These facts give (a)--(c). The signs in \eqref{support:2-18} preserve $a\mathbb Z^3$ and map the corresponding site charts onto one another. In each coordinate an odd index displacement gives reflection in the wall halfway between the two centres. An even index displacement gives translation by a multiple of $2l$, which is a product of two wall reflections. Thus the full map lies in the wall group and (d) holds. The composition and parity calculation in \Cref{alt:copy-operations} applies also with integer indices and proves (e).
\end{proof}

\begin{definition}[Period boxes and seminorms]\label{support:18-14}
Write $\mathfrak L=\{2L^m:m\in\mathbb Z\}$ for the sides. Let $(i,o)$ be a board and $\ell\in\mathfrak L$ with $\ell\ge2l$.
\begin{enumerate}
\item[(1)] A period box of cell indices is $k_0+\{0,\ldots,\ell/l-1\}^3$, where $k_0\in\mathbb Z^3$.
\item[(2)] For a function $F$ of one cell, $\|F\|^{(i,o)}_{{\rm ch},\ell}$ is the chessboard seminorm of its restriction to periodically lifted configurations, as a function of the corresponding torus cell. Write $\|F\|_{{\rm ch},\ell}$ when the board is understood.
\end{enumerate}
\end{definition}

\begin{lemma}[Period-box formula]\label{support:18-16}
For a period box $\mathcal K_\ell$ and a cell function $F$,
\begin{equation}\label{support:18-17}
 \|F\|_{{\rm ch},\ell}^{(\ell/l)^3}
 =\widetilde\mu_{K,\ell,\beta}\left(\prod_{k'\in\mathcal K_\ell}F^{[k']}\right).
\end{equation}
\end{lemma}
\begin{proof}
Since $\ell/l$ is an even integer, taking a copy commutes with periodic projection and its parity is unchanged by changing the lift of an index. Every torus cell has exactly one index in the box. The restriction of the product to periodic configurations is therefore the torus fold. Apply \eqref{support:4-14}.
\end{proof}

\begin{corollary}[Chessboard estimates for the lifted laws]\label{support:18-18}
\begin{enumerate}
\item[(a)] The lifted measure $\widetilde\mu_{K,\ell,\beta}$ has reflection positivity on local functions in each open half-period strip at a coordinate bond mirror.
\item[(b)] For cell functions $F_k$ in any subset of a single period box,
\[
 \left|\widetilde\mu_{K,\ell,\beta}\left(\prod_k F_k\right)\right|
 \le\prod_k\|F_k\|_{{\rm ch},\ell}.
\]
\item[(c)] The properties of \Cref{support:4-17} hold for the lifted seminorms, with the same normalization, bound by the sup norm, and invariance under mirror copies.
\end{enumerate}
\end{corollary}
\begin{proof}
Periodic projection carries the functions in (a) into the half-torus algebra. Apply \Cref{support:4-8}. In (b) distinct indices in a period box project to distinct cells. Apply \Cref{support:4-15}, with all unused factors equal to one. Since restriction to the periodic configurations preserves products, sums and conjugation and commutes with copies, \Cref{support:4-17} gives (c).
\end{proof}

\begin{lemma}[Moving a board]\label{support:18-26}
Let $r$ be a lattice isometry with linear part $r^\circ$ and $(i,o)$ a board.
\begin{enumerate}
\item[(a)] It carries cell $k$ of $(i,o)$ to cell $r^\circ k$ of $(i,ro)$ and carries their inner bond algebras onto one another.
\item[(b)] Conjugation by $r$ carries the mirror map from $k$ to $k'$ to the mirror map from $r^\circ k$ to $r^\circ k'$ and preserves its parity.
\item[(c)] For each $\ell\in\mathfrak L$ with $\ell\ge2l$,
\begin{equation}\label{support:18-27}
 \|rF\|^{(i,ro)}_{{\rm ch},\ell}=\|F\|^{(i,o)}_{{\rm ch},\ell}.
\end{equation}
\end{enumerate}
\end{lemma}
\begin{proof}
A signed permutation preserves the cube $[-l/2,l/2]^3$. Substitution in the mirror formula \eqref{support:2-18} gives the conjugate map, with the diagonal signs permuted and determinant unchanged. This proves (a) and (b). The image of a period box is a period box, and the lifted cubic Wilson measure is invariant under $r$. Formula \eqref{support:18-17} therefore proves (c).
\end{proof}

The following surface estimate removes the boundary condition.
\begin{proposition}\label{support:18-28}
Let $\ell\in\mathfrak L$ with $\ell/a$ an integer at least $3$.
For $y\in a(\mathbb Z+1/2)^3$, put
$D=(y+[0,\ell]^3)\cap a\mathbb Z^3$ and let $W$ be the set of
positive bonds whose initial point lies in $D$. For every
nonnegative bounded Borel function $\Psi$ of $W$ and every exterior
configuration $\xi$,
\begin{equation}\label{support:18-29}
 e^{-36\beta(\ell/a)^2}\widetilde\mu_{K,\ell,\beta}(\Psi)
 \le\gamma_W\Psi(\xi)
 \le e^{36\beta(\ell/a)^2}\widetilde\mu_{K,\ell,\beta}(\Psi).
\end{equation}
The same inequalities hold with $\gamma_W\Psi(\xi)$ replaced by
$\nu(\Psi)$ for any Gibbs state $\nu$. For every real bounded
function $H$ of $W$,
\begin{equation}\label{support:18-30}
 \nu(H)\le\log\widetilde\mu_{K,\ell,\beta}(e^H)
                          +36\beta(\ell/a)^2.
\end{equation}
\end{proposition}
\begin{proof}
Use the bonds of $W$ as Haar coordinates for both the Gibbs kernel
and the torus measure, and put $n=\ell/a$. In each of the three
plaquette orientations, $n(n-1)^2$ plaquettes have all their bonds
in $W$ and occur with the same action in both measures. There are
$n^3$ torus plaquettes and $n^3+2n^2$ plaquettes touching $W$ in
that orientation. Thus, after cancellation, the Gibbs action has
$N_1=12n^2-3n$ terms and the torus action has $N_2=6n^2-3n$ terms.
Since $0\le a_p\le2$, their difference is between $-2N_2$ and
$2N_1$. Comparing both numerators and both normalizing denominators
gives a density ratio between
$e^{-2\beta(N_1+N_2)}$ and $e^{2\beta(N_1+N_2)}$;
$2(N_1+N_2)=36n^2-12n\le36n^2$ proves \eqref{support:18-29}.
Integrate the kernel inequality against $\nu$ and use its Gibbs
equation. Finally Jensen's inequality in the kernel gives
$\gamma_WH\le\log\gamma_We^H$; apply \eqref{support:18-29} and integrate
once more to obtain \eqref{support:18-30}.
\end{proof}

\begin{definition}[Domination data]\label{support:18-32}
Let $\nu$ be a periodic Gibbs state with period number $P$, let $(i,o)$ be a board with cell side $l$, and let $F_k\in\mathcal A(\square_k)$ for $k$ in a nonempty finite set $\mathcal K_0$. Choose an even positive multiple $q$ of $P$ and $k_0$ with $\mathcal K_0\subset k_0+\{0,\ldots,q-1\}^3$. Such a choice exists by choosing coordinatewise minima for $k_0$ and increasing $q$. Set $\mathsf P=ql$ and $m=\prod_k\|F_k\|_\infty$. For $\ell\in\mathfrak L$, $\ell\ge\mathsf P$, set
\begin{equation}\label{support:18-33}
 N_\ell=\lfloor\ell/\mathsf P\rfloor^3,\qquad
 \epsilon_\ell=36\beta(\ell/a)^2/N_\ell.
\end{equation}
\end{definition}

\begin{theorem}[Chessboard domination]\label{support:18-38}
\begin{enumerate}
\item[(a)] For data of \Cref{support:18-32}, $\ell\ge\mathsf P$ and $t>0$ with $2tm\le1$,
\begin{equation}\label{support:18-39}
 \left|\nu\left(\prod_kF_k\right)\right|
 \le\prod_k\|F_k\|_{{\rm ch},\ell}+tm^2+\epsilon_\ell/t.
\end{equation}
\item[(b)] For every such state and family and every sequence $\ell_j\in\mathfrak L$ tending to infinity,
\begin{equation}\label{support:18-40}
 \left|\nu\left(\prod_{k\in\mathcal K_0}F_k\right)\right|
 \le\liminf_{j\to\infty}\prod_{k\in\mathcal K_0}
                  \|F_k\|_{{\rm ch},\ell_j}
 \le\prod_{k\in\mathcal K_0}\limsup_{j\to\infty}
                  \|F_k\|_{{\rm ch},\ell_j}.
\end{equation}
\item[(c)] For a real bounded function $H$ of one cell, the same packing, with $\ell\ge\mathsf P$, gives
\begin{equation}\label{support:18-41}
 \nu(H)\le\log\|e^H\|_{{\rm ch},\ell}+\epsilon_\ell.
\end{equation}
\item[(d)] If $\nu$ is invariant under every reflection in the walls of the board, then for every $\ell\in\mathfrak L$ with $\ell\ge2l$,
\begin{equation}\label{support:18-42}
 \nu(H)\le\log\|e^H\|_{{\rm ch},\ell}
                              +36\beta l^3a^{-2}\ell^{-1}.
\end{equation}
\end{enumerate}
\end{theorem}
\begin{proof}
The vectors
\[
 V_\ell=\{v\in\mathsf P\mathbb Z^3:
                    0\le v_\mu\le\ell-\mathsf P\}
\]
are periods of $\nu$, since $l/a$ is an integer and $P\mid q$.
They belong to $2l\mathbb Z^3$, so $\tau_vF_k$ is the mirror copy
in cell $k+v/l$. All these cell indices are distinct and fit in
one $\ell$-period box: division with remainder by $q$ separates
the translated blocks. Their inner bonds lie in the set $W$ of
\ref{support:18-28}, with corner
$y=o+l(k_0-(1/2,1/2,1/2))$. Applying \eqref{support:18-30} to the sum
of the $N_\ell$ translates of a real local function $\Phi$ gives
\begin{equation}\label{support:18-37}
 N_\ell\nu(\Phi)\le36\beta(\ell/a)^2+
 \log\widetilde\mu_{K,\ell,\beta}
                         \left(\prod_{v\in V_\ell}e^{\tau_v\Phi}\right).
\end{equation}
Here $\#V_\ell=N_\ell$, and
$\lfloor\ell/\mathsf P\rfloor\ge\ell/(2\mathsf P)$ proves the
bound $\epsilon_\ell\le288\beta\mathsf P^3a^{-2}\ell^{-1}$.

Write $F=\prod_kF_k$, and choose $|\zeta|=1$ so that
$\zeta\nu(F)=|\nu(F)|$. If $m=0$ there is nothing to prove.
Otherwise put $G=\operatorname{Re}(\zeta F)$ and
$\Phi=\log(1+tG)$. Since $|tG|\le1/2$ and
$\log(1+x)\ge x-x^2$ on $[-1/2,1/2]$,
\[
 \nu(\Phi)\ge t|\nu(F)|-t^2m^2.
\]
Expand $\prod_{v\in V_\ell}(1+t\tau_vG)$ over subsets
$Y\subset V_\ell$. Each product $\prod_{v\in Y}\tau_vG$ is
$2^{-|Y|}$ times a sum of $2^{|Y|}$ products of mirror copies of
the $F_k$ or $\overline{F_k}$, with coefficients of modulus one.
The finite chessboard estimate and invariance of its seminorm
under copying and conjugation bound each such mean by
$\Pi_\ell^{|Y|}$, where $\Pi_\ell=\prod_k\|F_k\|_{{\rm ch},\ell}$.
Consequently
\[
 \widetilde\mu_{K,\ell,\beta}
      \left(\prod_{v\in V_\ell}e^{\tau_v\Phi}\right)
 \le(1+t\Pi_\ell)^{N_\ell}\le e^{N_\ell t\Pi_\ell}.
\]
Insert this in \eqref{support:18-37} and divide by $N_\ell t$ to
obtain \eqref{support:18-39}. First let $\ell=\ell_j\to\infty$ at
fixed $t$, and then let $t\downarrow0$, proving
the first inequality of \eqref{support:18-40}. The second follows by bounding each factor eventually by its upper limit plus an arbitrary positive error, since the family is finite and all the factors are bounded. For $H$, choose
the corner $k_0$ so that the packing box contains its cell, retaining
the same $q$. Apply \eqref{support:18-37} and bound the torus mean
by $\|e^H\|_{{\rm ch},\ell}^{N_\ell}$ to obtain
\eqref{support:18-41}. The seminorm is positive because its
defining integrand is bounded below by a positive number.

For \eqref{support:18-42}, sum all the mirror copies of $H$ in an
$\ell$-period box. The mean of this sum is $(\ell/l)^3\nu(H)$,
while the torus exponential mean is exactly
$\|e^H\|_{{\rm ch},\ell}^{(\ell/l)^3}$. Use
\eqref{support:18-30} and divide by $(\ell/l)^3$.
\end{proof}

\begin{corollary}[Finite-side error]\label{support:18-43}
With the data of \Cref{support:18-32}, if $\ell\ge\mathsf P$ and $\epsilon_\ell\le1/4$, then
\[
 \left|\nu\left(\prod_kF_k\right)\right|
 \le \prod_k\|F_k\|_{{\rm ch},\ell}+2m\sqrt{\epsilon_\ell}.
\]
For $\ell\ge\mathsf P$, the error condition follows from $\ell\ge1152\beta\mathsf P^3a^{-2}$.
\end{corollary}
\begin{proof}
The case $m=0$ is immediate. If $\epsilon_\ell=0$, let $t\downarrow0$ in \eqref{support:18-39}. Otherwise take $t=\sqrt{\epsilon_\ell}/m$ there. The estimate for $\epsilon_\ell$ in the proof of \Cref{support:18-38} gives the threshold.
\end{proof}
\subsection{Two-layer exchange and comparison states}
\label{sec:gibbs-two-layer}

The comparison argument for diagonal reflections requires more than
positivity of a finite-volume measure.  We first prove an exchange
identity for two neighbouring layers of every periodic Gibbs state.
The proof inserts a unitary operator in the transfer matrix and compares
the resulting partition function with a Gibbs expectation.  Strict
positivity of the Fourier coefficients of the plaquette weight then
turns this operator identity into an identity for arbitrary continuous
functions of the two layers.

\subsubsection{The plaquette kernel and its square root}

Throughout this subsection $G=\operatorname{SU}(2)$ has normalized Haar
measure and $\beta>0$.  Set
\[
 k_\beta(g)=\exp\{-\beta(1-\tfrac12\operatorname{tr}g)\},\qquad
 k_S(u,u')=\prod_{b\in S}k_\beta(u(b)u'(b)^{-1})
\]
for a finite set $S$.  We write $K_S$ for the integral operator with
kernel $k_S$ on $L^2(G^S)$.  Inner products are antilinear in the first
argument.

For $l\in\frac12\mathbb N_0$, let $D^l$ be the representation on the
homogeneous polynomials of degree $2l$ in two variables, acting by
$\phi(z)\mapsto\phi(zg)$, with an invariant inner product obtained by
averaging any positive inner product over $G$.  Choose an orthonormal
basis and write $D^l_{mn}$ for its matrix coefficients and
$\chi_l=\operatorname{tr}D^l$ for its character.  Put
$\chi_{-1/2}=0$.  The normalized products
\[
 e_\alpha(u)=\prod_{b\in S}(2l_b+1)^{1/2}D^{l_b}_{m_bn_b}(u(b))
\]
are indexed by $\alpha=(\underline l,\underline m,\underline n)$.
Their finite linear combinations form $\mathcal T(G^S)$.

\begin{lemma}\label{support:29-4}
The coefficients satisfy
\begin{equation}\label{support:29-5}
 \int_G D^l_{mn}(g)\overline{D^{l'}_{m'n'}(g)}\,dg
 =\frac{\delta_{ll'}\delta_{mm'}\delta_{nn'}}{2l+1}.
\end{equation}
The functions $e_\alpha$ are an orthonormal basis of $L^2(G^S)$.
The algebra $\mathcal T(G^S)$ is invariant under conjugation and
uniformly dense in $C(G^S)$.  Each fixed-spin coefficient space is
invariant under conjugation.  Moreover, $\chi_l$ is real,
$\chi_l(g^{-1})=\chi_l(g)$, $|\chi_l|\le2l+1$, and
\begin{equation}\label{support:29-7}
 (\operatorname{tr}g)^N=\sum_l m_{N,l}\chi_l(g),\quad
 m_{N,l}\ge0,\quad m_{2l,l}=1,\quad
 \sum_l(2l+1)m_{N,l}=2^N,
\end{equation}
where $m_{N,l}=0$ for $2l>N$.
\end{lemma}

\begin{proof}
In the monomial basis $z_1^jz_2^{2l-j}$, a diagonal group element with
entries $e^{i\omega},e^{-i\omega}$ acts with weights $2j-2l$.
Taking $\omega=\pi/(2l+1)$ makes these eigenvalues distinct.  An
invariant subspace therefore contains every nonzero monomial component
of each of its elements, by polynomial spectral projections.  The
matrix
$2^{-1/2}\left(\begin{smallmatrix}1&1\\-1&1\end{smallmatrix}\right)$
sends any monomial to a polynomial with a nonzero $z_1^{2l}$
coefficient, and sends $z_1^{2l}$ to a polynomial with every monomial
coefficient nonzero.  Thus the representation is irreducible.
An intertwiner has invariant kernel and image; it is zero between
different dimensions, and is scalar in one irreducible representation
by subtracting an eigenvalue.  Averaging
$D^l(g)BD^{l'}(g)^{-1}$ and taking $B$ to be a matrix unit proves
\eqref{support:29-5}.

Let $\mathcal P_N$ be the restrictions to $G$ of polynomials of degree
at most $N$ in its four matrix entries.  The coefficients with
$2l\le N$ span a subspace of dimension
$\sum_{j=0}^N(j+1)^2=(N+1)(N+2)(2N+3)/6$.
The identity $g_{11}g_{22}-g_{12}g_{21}=1$ expresses every monomial of
degree at most $N$ as a sum of monomials of degrees $N$ and $N-1$.
The dimensions of these two homogeneous polynomial spaces add to the
same number.  Consequently $\mathcal P_N$ is exactly the orthogonal
sum of these coefficient spaces.  Conjugating an entry of an
$\operatorname{SU}(2)$ matrix gives another entry, up to sign.  Hence
$\mathcal P_N$ is conjugation invariant, as is its orthogonal
complement to $\mathcal P_{N-1}$, the spin-$N/2$ space.
Products over $S$ now give a self-adjoint algebra which contains the
constants and separates points.  Stone--Weierstrass and
\eqref{support:29-5} prove the asserted density and completeness.

Every $g\in G$ is conjugate in $G$ to
$\operatorname{diag}(e^{i\omega},e^{-i\omega})$, so
$\chi_l(g)=\sum_{j=0}^{2l}e^{i(2j-2l)\omega}$.
This proves the reality and bounds and gives
$\chi_{1/2}\chi_l=\chi_{l+1/2}+\chi_{l-1/2}$.
Starting with $m_{0,l}=\delta_{0l}$, the recursion
$m_{N+1,l}=m_{N,l-1/2}+m_{N,l+1/2}$ proves
\eqref{support:29-7}; evaluation at the identity gives its last equality.
\end{proof}

\begin{lemma}\label{support:29-11}
Define
\[
 \lambda_l=\frac1{2l+1}\int_Gk_\beta(g)\chi_l(g)\,dg,
 \qquad \lambda_\alpha=\prod_{b\in S}\lambda_{l_b}.
\]
Then
\begin{equation}\label{support:29-13}
 \frac{e^{-\beta}(\beta/2)^{2l}}{(2l)!(2l+1)}
 \le\lambda_l\le\frac{\beta^{2l}}{(2l+1)^2(2l)!},
 \qquad \sum_l(2l+1)^2\lambda_l=1.
\end{equation}
The operator $K_S$ is positive, injective and trace class, with
$K_Se_\alpha=\lambda_\alpha e_\alpha$ and
$\operatorname{Tr}K_S=1$.  Its positive square root has the continuous
real symmetric kernel
\[
 q_S(u,u')=\prod_{b\in S}q_\beta(u(b)u'(b)^{-1}),\qquad
 q_\beta=\sum_l(2l+1)\sqrt{\lambda_l}\chi_l.
\]
Writing $q_u=q_S(u,\cdot)$, we have
\begin{equation}\label{support:29-18}
 \langle q_u,q_{u'}\rangle=k_S(u,u'),\qquad \|q_u\|=1.
\end{equation}
Both $K_S$ and its square root map each fixed-spin space to itself
by a strictly positive scalar.  In particular, $K_S$ maps
$\mathcal T(G^S)$ onto itself.  If a finite complex measure $\sigma$
on $G^S\times Y$, with $Y$ compact metrizable, annihilates
$(K_S\varphi)(u)\psi(y)$ for every $\varphi\in\mathcal T(G^S)$
and a uniformly dense collection of $\psi\in C(Y)$, then $\sigma=0$.
\end{lemma}

\begin{proof}
Expanding the exponential and using \eqref{support:29-7} gives
\[
 (2l+1)\lambda_l
 =e^{-\beta}\sum_{N\ge2l}\frac{(\beta/2)^N}{N!}m_{N,l}.
\]
The term $N=2l$ gives the lower bound.  For the upper bound use
$m_{N,l}\le2^N/(2l+1)$ and
$\sum_{N\ge j}\beta^N/N!\le\beta^j e^\beta/j!$.
The character expansion
$k_\beta=\sum_l(2l+1)\lambda_l\chi_l$ converges absolutely
uniformly, since the exponential expansion has a summable majorant.
At the identity it gives the sum rule in \eqref{support:29-13}.
Schur orthogonality gives the stated eigenvalues of $K_S$; their sum
with multiplicities is $1$.  Positivity, injectivity and the trace
assertions follow.

The upper bound in \eqref{support:29-13} gives
$\sum_l(2l+1)^2\sqrt{\lambda_l}<\infty$, by the ratio test.
Thus the series defining $q_\beta$ converges uniformly.  Applying
orthogonality once more shows that its product kernel has eigenvalues
$\sqrt{\lambda_\alpha}$, and convolution of two copies gives $k_S$.
This proves \eqref{support:29-18}.  Finally, division by the nonzero
eigenvalues is legitimate on every finite coefficient sum.  The
assertion about measures follows from this surjectivity, density of
products in $C(G^S\times Y)$, and uniqueness of representing measures.
\end{proof}

For a finite bond scheme with vertex set $X$ and bond set $S$, put
$(w\cdot u)(b)=w(b_-)u(b)w(b_+)^{-1}$ and
\[
 (R_w\varphi)(u)=\varphi(w^{-1}\cdot u),\qquad
 P\varphi=\int R_w\varphi\,dw.
\]
Haar invariance shows that $R_w$ is unitary and $P$ is the orthogonal
projection onto gauge invariant functions.  Conjugation invariance
of $q_\beta$ gives $q_S(w\cdot u,w\cdot u')=q_S(u,u')$.
Consequently $K_S$ and $K_S^{1/2}$ commute with $P$ and every $R_w$.
If $C$ has continuous kernel $c$, then $CP$ has kernel
$\int c(u,w\cdot u')\,dw$.

\subsubsection{Operator insertions and the finite transfer formula}

Use the third direction as time.  A configuration is specified by
horizontal layers $u_t$ on the positive bonds of $a\mathbb Z^2$ and
vertical variables $w_t(y)$ between heights $ta$ and $(t+1)a$.
The vertical plaquette over a horizontal bond $b$ has weight
\begin{equation}\label{support:29-23}
 k_\beta\bigl(u_t(b)[(w_t\cdot u_{t+1})(b)]^{-1}\bigr).
\end{equation}
Indeed, its holonomy is
$u_t(b)w_t(b_+)u_{t+1}(b)^{-1}w_t(b_-)^{-1}$.
Horizontal plaquette actions are unchanged by a planar gauge
transformation.

For a bounded operator $X$ on $L^2(G^S)$ define
\begin{equation}\label{support:29-25}
 \kappa_X(u,u')=\langle q_u,Xq_{u'}\rangle,
 \qquad A_X(u,u')=\frac{\kappa_X(u,u')}{k_S(u,u')},
 \qquad
 A_X^{(t)}(U)=A_X(u_t|_S,(w_t\cdot u_{t+1})|_S).
\end{equation}
The last expression uses the horizontal copies of $S$ at heights $t$
and $t+1$ and the vertical bonds at their endpoints.  In particular,
it is a continuous local function supported in that closed two-layer
slab.

\begin{proposition}\label{support:29-28}
The map $X\mapsto A_X$ is linear, $A_1=1$, and
\[
 \|A_X\|_\infty\le e^{2\beta\#S}\|X\|.
\]
It preserves reality, tensor products on disjoint bond sets, and
relabeling of bonds.  The function $\kappa_X$ is a continuous kernel
of $K_S^{1/2}XK_S^{1/2}$.

Let $n_1,n_2,s\ge3$ be integers and
$\Gamma=a(n_1\mathbb Z\times n_2\mathbb Z\times s\mathbb Z)$.
Let $S_0$ and $X_0$ be the horizontal bonds and vertices modulo the
first two periods.  On $L^2(G^{S_0})$, let $E$ multiply by
$e^{-\beta V}$, where $V$ is the sum of the horizontal plaquette
actions in one layer, and set
\[
 Q=K_{S_0}^{1/2},\qquad M=QPEQ.
\]
Then $M\ge0$, $M$ is trace class and $Z_\Gamma[0]=\operatorname{Tr}M^s$.
At each height $t$, insert operators $X_{t,j}$ on finitely many
horizontal bond sets $S_{t,j}$. Require that the map
$(j,b)\mapsto[b]$ from the disjoint union of the $S_{t,j}$ to $S_0$
is injective. Relabel each operator by the resulting bijection onto
its image in $S_0$, and let $X_t$ be their tensor product with the
identity on the remaining bonds. Then
\begin{equation}\label{support:29-30}
 \mu_{\Gamma,\beta}\left(\prod_{t=1}^s\prod_j A_{X_{t,j}}^{(t)}\right)
 =\frac{\operatorname{Tr}(X_1M\cdots X_sM)}{\operatorname{Tr}M^s},
 \qquad
 \left|\mu_{\Gamma,\beta}\left(\prod_{t,j}A_{X_{t,j}}^{(t)}\right)\right|
 \le\prod_{t,j}\|X_{t,j}\|.
\end{equation}
Heights in this formula are understood modulo $s$.
\end{proposition}

\begin{proof}
The map $u\mapsto q_u$ is continuous in $L^2$, by uniform continuity
of the kernel.  Hence $\kappa_X$ is continuous and
$|\kappa_X|\le\|X\|$.  Since $k_S\ge e^{-2\beta\#S}$, the norm
bound follows.  Applying the root kernel on both sides of $X$ gives
the kernel assertion, first against continuous test functions and
then in $L^2$.  Formula \eqref{support:29-18} gives $A_1=1$.
Product Haar measure and the product formula for $q_S$ prove the tensor
and relabeling assertions.  They also prove preservation of reality.

The torus Haar measure is the product of the measures $du_t\,dw_t$,
and \eqref{support:29-23} gives
\[
 e^{-\beta S_\Gamma(U)}
 =\prod_{t=1}^s e^{-\beta V(u_t)}k_{S_0}(u_t,w_t\cdot u_{t+1}),
 \qquad u_{s+1}=u_1.
\]
Multiplication by the insertions replaces the $t$th kernel by
$\kappa_{X_t}$.  Integrating $w_t$ gives the continuous kernel of
$F_t=EQX_tQP$.  Integrating the remaining cyclic product of kernels
therefore gives $\operatorname{Tr}(F_1\cdots F_s)$.
Each factorization needed for cyclicity contains two Hilbert--Schmidt
factors $Q$; moving $EQ$ from the beginning to the end yields
\[
 Z_\Gamma[0]\mu_{\Gamma,\beta}\left(\prod_{t,j}A_{X_{t,j}}^{(t)}\right)
 =\operatorname{Tr}(X_1M\cdots X_sM).
\]
Since $V$ is gauge invariant, $P$ commutes with $E^{1/2}$ and
$M=(E^{1/2}PQ)^*(E^{1/2}PQ)$ is positive trace class.  Empty
insertions give $Z_\Gamma[0]=\operatorname{Tr}M^s>0$.

For the norm bound, let $L\ge0$ be trace class and let
$B_1,\ldots,B_s$ be bounded. The trace H\"older inequality gives
\begin{equation}\label{support:29-20}
 |\operatorname{Tr}(B_1L\cdots B_sL)|
 \le\prod_t\|B_t\|\operatorname{Tr}L^s.
\end{equation}
Indeed, apply \Cref{support:3-35} to the $2s$ factors
$B_tL^{1/2},L^{1/2}$.  For positive trace class $A\le B$,
$\operatorname{Tr}A^s\le\operatorname{Tr}B^s$: in an eigenbasis
of $A$, bound each eigenvalue by $\langle e,Be\rangle$, and then
use scalar Jensen in an eigenbasis of $B$.
Applying this to
$L^{1/2}B_t^*B_tL^{1/2}\le\|B_t\|^2L$ bounds the Schatten
norm of $B_tL^{1/2}$ by
$\|B_t\|(\operatorname{Tr}L^s)^{1/(2s)}$.
This proves \eqref{support:29-20}, and hence the second part of
\eqref{support:29-30}.
\end{proof}

\subsubsection{Exchange of two neighbouring layers}

\begin{proposition}\label{support:29-36}
Let $a>0$, $\beta>0$, and let $\nu$ be a periodic Gibbs state.
For a finite horizontal bond set $S$, an integer $t$, and every
$H\in C(G^S\times G^S)$,
\[
 \nu H(u_t|_S,(w_t\cdot u_{t+1})|_S)
 =\nu H((w_t\cdot u_{t+1})|_S,u_t|_S).
\]
More precisely, if $D$ is a bounded real operator on $L^2(G^S)$
with transpose $D^{\mathsf T}=-D$, then $\nu(A_D^{(t)})=0$.
\end{proposition}

\begin{proof}
For a real operator, $D^{\mathsf T}=D^*$, so
$X_\varepsilon=e^{\varepsilon D}$ is real unitary for real
$\varepsilon$.  Its remainder after $1+\varepsilon D$ has norm at
most $\frac12\varepsilon^2\|D\|^2e^{|\varepsilon|\|D\|}$.
By \Cref{support:29-28}, for sufficiently small
$|\varepsilon|$ the continuous local insertion
$A_{X_\varepsilon}^{(t)}$ lies in $[1/2,3/2]$.  Set
$A=\log A_{X_\varepsilon}^{(t)}$.  Then
\begin{equation}\label{support:29-37}
 \|A\|_\infty\le\log2,
 \qquad \|A-\varepsilon A_D^{(t)}\|_\infty\le C_D\varepsilon^2.
\end{equation}
Choose a period number $N_0$ of $\nu$ and a multiple $N'\ge3$ of
$N_0$ large enough that distinct translates of $S$ by $N'a\mathbb Z^2$
are disjoint.  For each $i\ge1$, take
$\Gamma' =N'a\mathbb Z^3$, $\Gamma=iN'a\mathbb Z^3$ and the
array $\mathbb A$ of $A$ over $\Gamma'/\Gamma$.
On the torus, $e^{\mathbb A}$ is a product of unitary insertions:
at heights $t+N'k_3$ use the $i^2$ translates by
$N'a(k_1,k_2)$, $0\le k_1,k_2<i$.  Their bond sets are pairwise
incongruent modulo the horizontal periods: a congruence would give
a nonzero $N'a\mathbb Z^2$ translation intersecting $S$, or equal
representatives.  Thus \eqref{support:29-30} gives
\[
 \log\frac{Z_\Gamma[\mathbb A]}{Z_\Gamma[0]}\le0.
\]
Enclose the support of $A$ in a lattice cube of side $ra$.  The lower
tangent inequality~\ref{support:14-20}, applied on the coordinate domain
of side $iN'a$, gives
\[
 0\ge i^3\nu(A)
 -\bigl(24\beta+2\log2(2r-1)^3\bigr)\#\partial_1D_i.
\]
Since $\#\partial_1D_i\le6i^2N'^2$, division by $i^3$ and
$i\to\infty$ yield $\nu(A)\le0$.
Using \eqref{support:29-37} for both signs of $\varepsilon$ and then
letting $\varepsilon\to0$ proves $\nu(A_D^{(t)})=0$.

It remains to recover arbitrary functions of the layers.  Let $\Phi$
be a real antisymmetric trigonometric polynomial on $G^S\times G^S$,
and write its finite expansion as
\[
 \Phi(u,u')=\sum_{\alpha,\alpha'}c_{\alpha\alpha'}
 e_\alpha(u)\overline{e_{\alpha'}(u')}.
\]
Divide its coefficients by
$\sqrt{\lambda_\alpha\lambda_{\alpha'}}$ to obtain $\Phi'$.
This operation preserves reality and antisymmetry: each fixed pair
of spin spaces is conjugation invariant, and swapping the two spaces
does not change the positive divisor.  The finite-rank operator $D$
with kernel $\Phi'$ is therefore real and skew-adjoint.  Applying
$K_S^{1/2}$ in both variables gives $\kappa_D=\Phi$, and hence
\[
 \nu\left[\frac{\Phi}{k_S}
 (u_t|_S,(w_t\cdot u_{t+1})|_S)\right]=0.
\]
For real continuous $H$, the function
$k_S(u,u')[H(u,u')-H(u',u)]$ is real and antisymmetric.
Approximate it uniformly by trigonometric polynomials, then take real
parts and antisymmetrize.  Division by $k_S\ge e^{-2\beta\#S}$
preserves uniform convergence, proving the desired exchange identity
for real $H$.  Complex linearity finishes the proof.
\end{proof}

\subsubsection{Comparison with a state positive at a prescribed mirror}

\begin{proposition}\label{support:29-40}
Let $a>0$, $\beta\ge0$, $n\in\mathfrak N$ and
$c\in\operatorname{Adm}(n)$.  For every periodic Gibbs state $\nu$
and every real $A\in\mathcal C_{\mathrm{loc}}(\mathbb H^+(n,c))$
there exists a periodic Gibbs state $\nu_0$, reflection positive at
$\Sigma(n,c)$, such that $\nu(A)\le\nu_0(A)$.
\end{proposition}

\begin{proof}
Choose a period number $N_0$ of $\nu$, a nonempty finite bond support
$B_A\subset\mathsf{Bd}(\mathbb H^+(n,c))$, and a lattice cube
$Q=x_0+[0,ra]^3$ containing its bonds.  Choose
$h\in2N_0a\mathbb N$ so that $B_A$ lies in
$\mathbb S_1=\{c\le n\cdot x\le c+h\}$, and put
$v_h=hn/|n|^2$.  If $n=e_\mu$, let
$\mathfrak T=N_0a(\mathbb Ze_{\mu'}\oplus\mathbb Ze_{\mu''})$.
If $n=e_\mu+s e_\nu$, $s=\pm1$, let
$\mathfrak T=N_0a\mathbb Z(e_\mu-s e_\nu)\oplus N_0a\mathbb Ze_\kappa$,
where $\kappa$ is the remaining direction.  Set
\[
 \Gamma'=\mathfrak T\oplus\mathbb Zv_h,\qquad
 \iota=[a\mathbb Z^3:\Gamma'],\qquad
 \Gamma_i=i\mathfrak T\oplus\mathbb Z(2iv_h),\quad i\ge3.
\]
These groups are adapted to $n$, with respective height periods $h$
and $2ih$.  Moreover $\Gamma'\subset N_0a\mathbb Z^3$, so $\nu$
is $\Gamma'$ invariant, and $\operatorname{sys}(\Gamma_i)\ge ia$.

In the coordinate case choose the domain
\[
 D_i=\{x\in a\mathbb Z^3:0\le x_\mu<2ih,
             \ 0\le x_{\mu'},x_{\mu''}<iN_0a\}.
\]
In the diagonal case choose
\[
 D_i=\{x\in a\mathbb Z^3:
 0\le x_\mu-sx_\nu<2iN_0a,\quad
 0\le x_\kappa<iN_0a,\quad
 0\le x_\mu+sx_\nu<2ih\}.
\]
The two diagonal coordinates divided by $a$ have the same parity.
Reducing them modulo $2iN_0$ and $2ih/a$ preserves this parity and
gives one representative per class. Thus these sets are domains of
$\Gamma_i$. In the definition of $\partial_1D_i$, the displacements
$w$ satisfy $|w|_\infty\le a$. Each ordinary coordinate changes by
at most $a$, while $x_\mu-sx_\nu$ and $x_\mu+sx_\nu$ change by
at most $2a$. Count one lattice layer at each end of an ordinary
coordinate interval and two at each end of a diagonal coordinate
interval. The parity count above gives the same number of sites
in each layer. The union bound therefore gives
\[
 \frac{\#\partial_1D_i}{\#D_i}
 \le\frac2i\left(\frac2{N_0}+\frac ah\right).
\]

Let $\mathbb A$ be the array of $A$ over $\Gamma'/\Gamma_i$, and
$\mathbb A^{\mathrm f}$ its folded array from
\S\ref{sec:infinite-volume-states}.
Its number of terms is $|\mathbb A|=\#D_i/\iota$.
For $t>0$, the lower tangent inequality, folding inequality, and
Jensen's inequality in the folded measure give, in that order,
\[
\begin{split}
 t|\mathbb A|\nu(A)
 -\bigl(24\beta+2t\|A\|_\infty(2r-1)^3\bigr)\#\partial_1D_i
 &\le \log\frac{Z_{\Gamma_i}[t\mathbb A]}{Z_{\Gamma_i}[0]}\\
 &\le \log\frac{Z_{\Gamma_i}[t\mathbb A^{\mathrm f}]}{Z_{\Gamma_i}[0]}\\
 &\le t\mu_{\Gamma_i}^{t\mathbb A^{\mathrm f}}(\mathbb A^{\mathrm f})
 =t|\mathbb A|\mu_{\Gamma_i}^{t\mathbb A^{\mathrm f}}(A).
\end{split}
\]
The last equality is the wall symmetry of the folded measure.  Hence
\begin{equation}\label{support:29-41}
 \nu(A)\le\mu_{\Gamma_i}^{t\mathbb A^{\mathrm f}}(A)
 +\iota\left(\frac{24\beta}{t}
       +2\|A\|_\infty(2r-1)^3\right)
       \frac{\#\partial_1D_i}{\#D_i}.
\end{equation}

We keep $t$ fixed while taking the volume limit.  To identify its
Gibbs specification, write the infinite folded family with its indices:
\begin{equation}\label{support:29-38}
 I=\mathbb Z\times\mathfrak T,\qquad
 A_{(j,v)}=t w_j\tau_vA,\qquad
 B_{(j,v)}=w_j(B_A+v).
\end{equation}
Here $w_j$ are the alternating translations and reflections which
send the first slab to the $j$th slab.  Every bond belongs to at most
$2\#B_A$ of these supports.  For
$g=g_\perp+2ikv_h\in\Gamma_i$, define
$g\cdot(j,v)=(j+2ik,v+g_\perp)$.
The relations $w_{j+2ik}=\tau_{2ikv_h}w_j$ and
$\tau_{g_\perp}w_j=w_j\tau_{g_\perp}$ show that this free index
action translates both the functions and their supports.  Representatives
of $I/\Gamma_i$ are $\{1,\ldots,2i\}\times R_i$, where $R_i$
represents $\mathfrak T/i\mathfrak T$, and their sum on the torus is
$t\mathbb A^{\mathrm f}$.  Thus the finite-volume DLR limit applies
to one fixed locally finite perturbation family.

By compactness, a subsequence of the folded torus measures converges
to a Gibbs state $\nu_t$ for this family.  The finite measures, and
hence their limit, are invariant under
$\mathfrak T\oplus\mathbb Z(2v_h)$.  A fixed continuous local
$F$ in $\mathbb H^+(n,c)$ lies in
$\{c\le n\cdot x\le c+ih\}$ for all sufficiently large $i$.
The two-mirror positivity of the folded torus measure then gives
$\mu_{\Gamma_i}^{t\mathbb A^{\mathrm f}}(\overline{\theta F}F)\ge0$,
and this inequality passes to $\nu_t$.
The error in \eqref{support:29-41} tends to zero, so $\nu(A)\le\nu_t(A)$.

Finally take a weak limit of a subsequence of $\nu_{1/k}$.
Continuity of the local DLR kernels as the strength of the fixed
family tends to zero makes its limit $\nu_0$ an unperturbed Gibbs
state.  The common period group, reflection positivity, and the
expectation inequality all pass to this limit.  The period group has
finite index, so $\nu_0$ is periodic, as required.
\end{proof}

\subsection{Symmetries of periodic Gibbs states}
\label{sec:gibbs-symmetries}

We first use the two-layer identity to obtain a common translation
group for all periodic Gibbs states.  This makes their class compact
and convex.  The comparison states can then test membership in its
reflection-positive subclass by separation.

\subsubsection{Temporal gauge and determination by two layers}

Fix an integer $t_0$.  Given the vertical variables $w_t$, define
\[
 \hat w_{t_0}=1,\qquad \hat w_{t+1}=\hat w_tw_t,
 \qquad \tilde u_t=\hat w_t\cdot u_t.
\]
This defines a continuous map
$\mathsf{tg}:\mathcal V(\mathbb L_a)\to(\mathcal V^{(2)})^{\mathbb Z}$:
each of its coordinates uses finitely many vertical variables and one
horizontal bond.  The gauge transformation $\hat w$ makes all vertical
bonds equal to $1$.  Write $U[\tilde u]$ for the resulting configuration,
with horizontal layers $\tilde u$ and vertical variables $1$.
For a gauge invariant local continuous function $F$,
\[
 F(U)=\tilde F(\mathsf{tg}(U)),\qquad
 \tilde F(\tilde u)=F(U[\tilde u]).
\]

For finite horizontal $S$, let $V_S(u)$ be the sum of the planar
plaquette actions which meet $S$.  It contains at most $2\#S$
terms, so $0\le V_S\le4\#S$.  Define the conditional layer kernel
\begin{equation}\label{support:54-6}
 \tilde k_{S,t}(v;\tilde u)
 =k_S(\tilde u_{t-1}|_S,v)
   e^{-\beta V_S(v\vee\tilde u_t|_{S^c})}
   k_S(v,\tilde u_{t+1}|_S).
\end{equation}
The kernel $\tilde\gamma_{S,t}$ replaces $\tilde u_t|_S$ by $v$
and integrates with the normalized density
$\tilde k_{S,t}(v;\tilde u)\,dv$.  This density and its normalizer
are bounded below by $e^{-8\beta\#S}$ before normalization.
They depend continuously on finitely many coordinates.  A layer law
obeys the chain equations if it is fixed by every $\tilde\gamma_{S,t}$
on continuous local functions.

\begin{lemma}\label{support:54-9}
The temporal-gauge image of every Gibbs state obeys the chain
equations.  The equations are invariant under each index reflection
$(\tilde\theta_m\tilde u)_t=\tilde u_{m-t}$.
If $\mathsf{tg}_*\nu$ is invariant under $\tilde\theta_{2t_0+1}$,
then $\nu$ is invariant under the bond reflection
$\theta_{e_3,(t_0+1/2)a}$.
\end{lemma}

\begin{proof}
The vertical plaquette formula becomes
$1-\frac12\operatorname{tr}(\tilde u_t(b)\tilde u_{t+1}(b)^{-1})$,
and the horizontal plaquette action is unchanged by $\hat w_t$.
If only the horizontal bonds $B=\{b^{(t)}:b\in S\}$ are integrated
in the DLR kernel, the variables $\hat w_s$ do not change.
The substitution $v=\hat w_t\cdot U_B$ preserves product Haar measure.
The plaquettes meeting $B$ consist of the planar plaquettes meeting
$S$, and one vertical plaquette above and one below each bond of $S$.
Their weight is precisely \eqref{support:54-6}.  Thus
\[
 \gamma_B(\psi\circ\mathsf{tg})
 =(\tilde\gamma_{S,t}\psi)\circ\mathsf{tg},
\]
which proves the chain equations.

Symmetry of $k_S$ exchanges the two neighbouring layers in
\eqref{support:54-6}; hence
$(\tilde\gamma_{S,m-t}\psi)\circ\tilde\theta_m
=\tilde\gamma_{S,t}(\psi\circ\tilde\theta_m)$.
This proves covariance of the equations.
With $\theta=\theta_{e_3,(t_0+1/2)a}$ one has
$\theta U[\tilde u]=U[\tilde\theta_{2t_0+1}\tilde u]$.
Consequently invariance of the layer law implies that $\nu$ and
$\theta_*\nu$ agree on all gauge invariant local continuous
functions.  Both measures are Gibbs states, so the gauge averaging
lemma~\ref{support:14-13} makes them equal.
\end{proof}

\begin{lemma}\label{support:54-13}
For $\beta>0$, two laws obeying the chain equations and agreeing on
two or more consecutive entire layers agree on all layers.
\end{lemma}

\begin{proof}
Suppose that the two laws agree on layers $m,\ldots,n$, with $m<n$.
We extend agreement to $n+1$.  Fix a nonempty finite horizontal set
$S$, and use the coordinates
\[
 u'=\tilde u_{n+1}|_S,\qquad v=\tilde u_n|_S,\qquad
 c=(\tilde u_n|_{S^c},\tilde u_{n-1},\ldots,\tilde u_m)\in Y.
\]
Here $Y$ is a compact countable product of copies of $G$.  Put
\[
 \Phi(v;c)=k_S(\tilde u_{n-1}|_S,v)
 e^{-\beta V_S(v\vee\tilde u_n|_{S^c})},\qquad
 N(u',c)=\int k_S(v,u')\Phi(v;c)\,dv.
\]
The assumption $m<n$ ensures that $\tilde u_{n-1}$ is among the
coordinates of $c$.  Both functions are continuous and local, with
$\Phi\ge e^{-6\beta\#S}$ and $N\ge e^{-8\beta\#S}$.
Let $\rho_i'$ be the image of the $i$th law on $(u',c)$ and set
$\sigma_i=N^{-1}\rho_i'$.  The chain equation at $(S,n)$ reads
\begin{equation}\label{support:54-14}
 \rho_i(\psi)=\int\sigma_i(du',dc)\int
 k_S(v,u')\Phi(v;c)\psi(u',v,c)\,dv
\end{equation}
for every continuous local $\psi$.
Take $\psi=\varphi(v)\chi(c)/\Phi(v;c)$, with
$\varphi\in\mathcal T(G^S)$ and $\chi$ continuous local.
This is a function of the already equal layers.  Subtracting the two
identities gives
\[
 \int(K_S\varphi)(u')\chi(c)\,d(\sigma_1-\sigma_2)=0.
\]
The measure injectivity in \Cref{support:29-11} yields
$\sigma_1=\sigma_2$.  Equation~\eqref{support:54-14} then gives agreement
for all continuous local functions of $(u',v,c)$.  Enlarging $S$ to
contain the finitely many coordinates of any test on layer $n+1$
proves agreement on layers $m,\ldots,n+1$.

Apply the same argument after replacing $t$ by $m+n-t$, using the
covariance in \Cref{support:54-9}, to extend agreement to
layer $m-1$.  Iterating both
extensions gives agreement on every finite block of layers.
Continuous local functions determine laws on this compact product,
so the two laws are equal.
\end{proof}

\begin{proposition}\label{support:54-16}
Every periodic Gibbs state is invariant under every coordinate bond
reflection, and consequently under translations by $2a\mathbb Z^3$.
The periodic Gibbs states are exactly the Gibbs states invariant under
$2a\mathbb Z^3$; their set is nonempty, compact and convex.
\end{proposition}

\begin{proof}
For $\beta>0$, take the temporal gauge based at $t_0$.
Its layers $t_0,t_0+1$ are $u_{t_0}$ and
$w_{t_0}\cdot u_{t_0+1}$.  \Cref{support:29-36} says that their
joint law is invariant under exchange.  The layer law and its
$\tilde\theta_{2t_0+1}$ reflection both obey the chain equations and
agree on these two layers.  \Cref{support:54-13} makes them equal,
and \Cref{support:54-9} gives the bond reflection invariance of $\nu$.
The same reasoning in each coordinate direction follows by applying
it to the coordinate permutation of $\nu$, which is again a periodic
Gibbs state.

At $\beta=0$, the DLR equation on the support of a local function
sets its expectation equal to its product Haar integral.  This
determines a unique Gibbs state, which is invariant under all lattice
isometries.

The product of the reflections at $x_\mu=a/2$ and $x_\mu=3a/2$
is translation by $2ae_\mu$.  Thus every periodic state is invariant
under $2a\mathbb Z^3$.  Conversely, this group has finite index,
so a Gibbs state invariant under it is periodic.  Compactness and
convexity follow from the DLR equations and the fixed common period
group; nonemptiness follows from the periodic volume construction.
\end{proof}

\subsubsection{All lattice reflections}

\begin{proposition}\label{support:54-17}
For every $a>0$ and $\beta\ge0$ the following hold.
\begin{enumerate}
\item[(a)] Every periodic Gibbs state is reflection positive at all coordinate site, coordinate bond and diagonal lattice mirrors, on bounded local Borel functions of either closed half-space.
\item[(b)] Every periodic Gibbs state is invariant under $a\mathbb Z^3\rtimes W_3$.
\item[(c)] The periodic Gibbs states are exactly the translation-invariant Gibbs states. Their class is nonempty, compact and convex and contains every cubic torus-limit state.
\end{enumerate}
\end{proposition}

\begin{proof}
Fix a mirror $\Sigma(n,c)$, $n\in\mathfrak N$, and write
$\theta=\theta_{n,c}$.  Let $\mathcal R$ be the periodic Gibbs
states reflection positive at this mirror.  It is nonempty by
\Cref{support:29-40} with $A=0$.  It is compact and convex:
\Cref{support:54-16} gives those properties for the ambient class,
and each condition
$\rho(\overline{\theta F}F)\in[0,\infty)$, with $F$ continuous
local in the positive half-space, is closed and affine.

Suppose $\nu\notin\mathcal R$.  There is a real continuous $f$ and
$\delta>0$ such that
$\nu(f)\ge\rho(f)+3\delta$ for every $\rho\in\mathcal R$.
Here separation can be obtained in a finite-dimensional space.
Indeed, if every finite family of continuous test integrals of $\nu$
were attained by an element of $\mathcal R$, compactness and the
finite intersection property would give an element with all the same
integrals, hence $\nu\in\mathcal R$.  Otherwise the finite vector
of integrals of $\nu$ lies outside the corresponding compact convex
set.  Its nearest point separates it by the vector joining that point
to the given vector, which gives $f$ and a positive gap.

Choose real local $A_0$ with $\|f-A_0\|_\infty\le\delta$.
Then $\nu(A_0)\ge\rho(A_0)+\delta$ for every $\rho\in\mathcal R$.
Translate $A_0$ by $2kan$, with $k$ so large that its support lies in
$\mathbb H^+(n,c)$, and call the result $A$.
All the states in question are invariant under this translation by
\Cref{support:54-16}, so the same strict inequality holds for $A$.
\Cref{support:29-40} now supplies $\nu_0\in\mathcal R$ with
$\nu(A)\le\nu_0(A)$, a contradiction.  This proves positivity for
continuous local functions.

The same positivity holds for bounded Borel local functions.  On
their finite compact coordinate space, approximate such a function
in $L^2$ for the sum of the marginal laws of $\nu$ and $\theta_*\nu$
by continuous functions.  Cauchy--Schwarz bounds the change in
$\nu(\overline{\theta F}F)$, so positivity passes to the limit.
Replacing $F$ by the conjugate of its reflection also gives positivity
on the negative half-space and for the opposite normal.

To obtain reflection invariance, positivity makes the form
$\nu(\overline{\theta F}F')$ Hermitian.  Taking one entry equal to
$1$ gives $\nu(F)=\nu(\theta F)$ for local $F$ in the positive
half-space.  Any local $F_0$ can be translated into that half-space
by $v=2kan$.  Since both $v$ and $-v$ are periods of $\nu$ and
$\theta\tau_v=\tau_{-v}\theta$, the same equality holds for $F_0$.
Thus $\theta_*\nu=\nu$ on all continuous local tests and hence on
all Borel sets.  The coordinate and diagonal lattice reflections
generate $a\mathbb Z^3\rtimes W_3$: adjacent site and bond mirrors
give unit translations, coordinate mirrors give sign changes, and
$x_\mu=x_\nu$ gives coordinate transpositions.
This proves the symmetry assertion.  The description and compactness
of the state class then follow from the translation-invariant DLR
equations, and the torus-limit inclusion is already known.
\end{proof}

\subsubsection{Invariant conditioning}

\begin{proposition}[Invariant reweighting]\label{support:54-18}
\begin{enumerate}
\item[(a)] Let $\nu$ be a periodic Gibbs state. Suppose $h\ge0$ is bounded, is $\mathcal I$-measurable and satisfies $\nu(h)=1$. Then $h\nu$ is a periodic Gibbs state.
\item[(b)] The extreme points of the periodic Gibbs class are exactly its ergodic states.
\end{enumerate}
\end{proposition}
\begin{proof}
By \Cref{support:54-17} the state $\nu$ is invariant under all lattice translations. Fix a finite bond set $B$. Approximate $h$ in $L^1(\nu)$ by uniformly bounded continuous local functions, and translate each approximant so that its support is outside $B$. The error is unchanged because both $h$ and $\nu$ are translation invariant. The DLR equation on $B$ implies $\nu(h_j\gamma_BF)=\nu(h_jF)$ for each translated approximant $h_j$. Pass to the limit to get the same equality for $h$. The law $h\nu$ is translation invariant by the invariance of $h$ and $\nu$, so this proves (a).

If an invariant event has probability strictly between zero and one, (a) decomposes $\nu$ into the two conditioned Gibbs laws. Conversely, in any decomposition into invariant states with positive weights, each summand has a bounded Radon--Nikodym density with respect to $\nu$. Uniqueness of that density makes it invariant. Ergodicity makes each density constant, so the decomposition is trivial.
\end{proof}

\begin{proposition}[Ergodic conditional Gibbs laws]\label{support:54-25}
\begin{enumerate}
\item[(1)] In the periodic Gibbs class, ergodicity is equivalent to extremality.
\item[(2)] Let $\nu$ be a periodic Gibbs state and let $\nu^U$ denote its regular conditional law given $\mathcal I$. Then for $\nu$-almost every $U$, the law $\nu^U$ belongs to $\mathcal G^{\rm per}_{a,\beta}$ and is ergodic.
\end{enumerate}
\end{proposition}
\begin{proof}
Part (1) is \Cref{support:54-18}(b). For (2), the DLR and translation-invariance equalities hold after integration over any invariant event by \Cref{support:54-18}(a). Thus their conditional expectations given $\mathcal I$ vanish. Use a countable uniformly dense family of continuous local tests, and the countability of finite bond sets and translations, to obtain all the equalities for $\nu^U$ simultaneously outside one null set. The equality on the dense family extends to continuous tests by contraction of the kernels in supremum norm. This proves the Gibbs and invariance assertions.

For a continuous local $F$, the cubic averages $A_nF$ of its translates converge in $L^2(\nu)$ to $\nu(F\mid\mathcal I)$. Indeed the commuting coordinate translations are unitary. Their one-direction Ces\`aro averages converge strongly to their invariant projections, by the spectral theorem, and the product converges to the common invariant projection. Choose a common subsequence with summable squared errors for every member of the countable dense family. By disintegration, the same subsequence converges in $L^2(\nu^U)$ for almost every $U$, with constant limit $\nu^U(F)$. Here the invariant random variable $\nu(F\mid\mathcal I)$ equals its value at $U$, $\nu^U$-almost surely. Applying the same Hilbert-space mean argument under $\nu^U$ identifies its invariant projection of each test as that constant. Density proves that the invariant subspace consists of constants, so $\nu^U$ is ergodic.
\end{proof}

\subsection{Insertion estimates for periodic Gibbs states}
\label{sec:gibbs-insertions}

\begin{definition}[Insertion setting]\label{support:55-1}
\begin{enumerate}
\item[(1)] Fix the odd integer $L$ and the cube parameter $M$ of the renormalisation argument. The insertion estimates only require $L\ge3$ odd and do not restrict $M$.
\item[(2)] At fixed cutoff $K$ and coupling $\beta>0$, use (a) the lattice $a_K\mathbb Z^3$, (b) the Wilson laws and their periodic lifts, (c) the Gibbs equations of \Cref{support:14-9}, (d) the mirrors and admissible heights of \Cref{support:54-1}, and (e) the boards and their seminorms from Defs.~\ref{support:18-2} and \ref{support:18-14}.
\end{enumerate}
\end{definition}

\begin{definition}[Invariance, positivity and domination]\label{support:55-2}
For a law $\rho$ on $a_K\mathbb Z^3$ and a sequence $\ell_\bullet$ of admissible sides tending to infinity, use the following conditions.
\begin{enumerate}
\item[(1)] \textup{(Inv)} is invariance under all lattice isometries.
\item[(2)] \textup{(Pos)} is reflection positivity at every bond mirror $x_3=c$ for bounded local functions in the positive half-space.
\item[(3)] \textup{(Dom)} for a class of functions $\mathfrak C$ means that on every board and for every finite family of cell functions $F_k\in\mathfrak C$,
\begin{equation}\label{support:55-5}
 \left|\rho\left(\prod_kF_k\right)\right|
 \le \liminf_{j\to\infty}\prod_k\|F_k\|_{{\rm ch},\ell_j}.
\end{equation}
Here $\mathfrak C$ is the bounded local Borel class or its continuous subclass.
\item[(4)] For a cell function $D$, write $\|D\|_{{\rm ch},\ell_\bullet}=\liminf_j\|D\|_{{\rm ch},\ell_j}$.
\item[(5)] At a bond mirror $c$, put $\langle F',F\rangle_{\rho,c}=\rho(\overline{\theta_c F'}F)$ and, under \textup{(Pos)}, $\|F\|_{\rho,c}=\langle F,F\rangle_{\rho,c}^{1/2}$.
\end{enumerate}
\end{definition}

\begin{corollary}[Domination and positivity of Gibbs laws]\label{support:55-40}
\begin{enumerate}
\item[(a)] Every periodic Gibbs state satisfies \textup{(Dom)} for every sequence of sides tending to infinity and for both function classes.
\item[(b)] Each limit of cubic torus measures belongs to the periodic Gibbs class and satisfies \textup{(Inv)} and \textup{(Pos)} as well.
\end{enumerate}
\end{corollary}
\begin{proof}
\Cref{support:18-38}(b) gives (a). \Cref{support:14-16}(c) gives the Gibbs inclusion in (b), and \Cref{support:54-17} gives the invariance and positivity. In fact the same two properties hold for every periodic Gibbs state.
\end{proof}

\subsubsection{Multiple reflections with an insertion}

We need the following consequence of the finite chessboard estimate.
Call $(i',o')$ coarser than $(i,o)$ if $i'\le i$ and
$o'-o\in l\mathbb Z^3$.  Then $d=a_{i'}/a_i=L^{i-i'}$ is odd,
and each coarse cell is the union of $d^3$ fine cells.

\begin{lemma}\label{support:18-24}
For a coarse cell $C'$ and functions $D_k$ of its fine cells,
\begin{equation}\label{support:18-25}
 \left\|\prod_{C_k\subset C'}D_k\right\|_{\mathrm{ch},\ell}^{\rm coarse}
 \le\prod_{C_k\subset C'}\|D_k\|_{\mathrm{ch},\ell}^{\rm fine},
 \qquad \ell\ge2a_{i'}.
\end{equation}
\end{lemma}

\begin{proof}
Writing $z=(o'-o)/l$, the fine cells of coarse cell $j$ have indices
\[
 z+dj+\{-(d-1)/2,\ldots,(d-1)/2\}^3.
\]
Every coarse wall is a fine wall.  Thus a coarse mirror copy sends
each fine cell to a fine cell by the corresponding fine mirror map;
its parity is unchanged, since parity is the determinant of that map.
The disseminated coarse product is consequently a product over all
fine cells in a torus period box.  Apply the fine chessboard inequality
to this product.  Every original factor has one copy for each coarse
cell, and its fine seminorm is unchanged by copying.  Taking the root
whose order is the number of coarse cells proves \eqref{support:18-25}.
\end{proof}

\begin{proposition}\label{support:55-9}
Let $\rho$ satisfy \textup{(Inv)}, \textup{(Pos)} and \textup{(Dom)} along $\ell_\bullet$ for the class $\mathfrak C$ of \Cref{support:55-2}. Suppose a cell $C$ of side $l$ lies between consecutive board walls $x_3=c$ and $x_3=c+l$. For a function $D\in\mathfrak C$ of this cell and functions $F,F'\in\mathfrak C$ in $\{x_3\ge c\}$,
\begin{equation}\label{support:55-10}
 \left|\rho(\overline{\theta_cF'}D\tau_{le_3}F)\right|
 \le\|D\|_{\mathrm{ch},\ell_\bullet}
       \|F'\|_{\rho,c}\|F\|_{\rho,c}.
\end{equation}
\end{proposition}

\begin{proof}
Use the reflected form of \Cref{support:55-2}(5), and
write $\|\cdot\|$ for its seminorm. Put $\widehat DF=D\tau_{le_3}F$
and $D^\sharp=\overline{\theta_{c+l/2}D}$, where the middle
reflection preserves the cell. Translation invariance and
$\theta_c\theta_{c+l/2}=\tau_{-le_3}$ give
\[
 \langle F',\widehat DF\rangle_{\rho,c}
 =\langle\widehat{D^\sharp}F',F\rangle_{\rho,c}.
\]
Let $F_n=(\widehat{D^\sharp}\widehat D)^nF$ and
$s_n=\langle F,F_n\rangle_{\rho,c}$.
Moving adjoints gives $s_{2p}=\|F_p\|^2$ and
$s_{2p+1}=\|\widehat DF_p\|^2$. The odd-indexed moment is also
$\langle F_p,F_{p+1}\rangle_{\rho,c}$, and the even-indexed
moment $s_{2p}$, for $p\ge1$, is
$\langle\widehat DF_{p-1},\widehat DF_p\rangle_{\rho,c}$.
Cauchy--Schwarz in these two identities gives
$s_n^2\le s_{n-1}s_{n+1}$ for $n\ge1$. Therefore, for $n\ge2$,
\begin{equation}\label{support:55-13}
 s_1^n\le s_0^{n-1}s_n.
\end{equation}
If $s_1>0$, log convexity first gives $s_n>0$ for all $n\ge0$ and
then makes the ratios $s_{n+1}/s_n$ increasing. If $s_1=0$,
the inequality holds directly.

Expanding the iterates yields
\[
 s_n=\rho\left(\overline{\theta_cF}
       \prod_{r=0}^{2n-1}D_r\,\tau_{2nle_3}F\right),\qquad
 D_r=\begin{cases}\tau_{rle_3}D^\sharp,&r\text{ even},\\
                     \tau_{rle_3}D,&r\text{ odd}.
       \end{cases}
\]
These are the successive mirror copies of $D^\sharp$.
Reflection in the cell centre preserves the board, so \Cref{support:18-26}(c)
and invariance under conjugation give
$\|D^\sharp\|_{\mathrm{ch},\ell}=\|D\|_{\mathrm{ch},\ell}$.
Choose an odd coarsening factor $d=L^{r_0}$ large enough that $F$
is supported in a coarse cell directly above the wall $c$.
Such a board has side $l'=dl$ and center obtained from that of $C$
by adding $(d-1)le_3/2$; its bottom wall remains $c$.
For $n=n'd$, group the $2n$ functions $D_r$ into $2n'$ consecutive
coarse cells, with $F$ and its reflected copy in the two end cells.
At each side $\ell_j\ge2l'$, the product of coarse seminorms is
at most $\|F\|_\infty^2\|D\|_{\mathrm{ch},\ell_j}^{2n}$
by \eqref{support:18-25} and the supremum bound on the end
factors. Chessboard domination and the lower limit give
\[
 s_n\le\|F\|_\infty^2
       \|D\|_{\mathrm{ch},\ell_\bullet}^{2n}.
\]
Combine this estimate with \eqref{support:55-13}, take $n$th roots
and let $n'\to\infty$. This gives
$\|\widehat DF\|\le\|D\|_{\mathrm{ch},\ell_\bullet}\|F\|$
when $\|F\|>0$. If $\|F\|=0$, \eqref{support:55-13} gives
$s_1=0$ directly.
Cauchy--Schwarz now proves \eqref{support:55-10}.
\end{proof}
\section{Separated tests and growth}\label{app:tests-and-growth}

\subsection{Separated functions and compactness}\label{app:separated}

Uniform estimates at separated supports give one subsequence for all
orders and all test tuples. The support restriction matters in this
argument: approximating the entries must preserve their separation.
We give the compactness statement on both a fixed torus and
$\mathbb R^3$.

\begin{definition}[Separated tuples]\label{support:9-5}
Let $\mathbb E=T_\ell$ or $\mathbb R^3$. For an open set
$O\subset\mathbb E$, let $\mathfrak S_n(O)$ be the tuples
$(f_1,\ldots,f_n)$ of real functions in $C_c^1(O)$ with pairwise
disjoint supports. Zero entries are allowed. Put
$\mathfrak S_0(O)=\{\varnothing\}$ and
$\mathfrak S(O)=\bigcup_{n\ge0}\mathfrak S_n(O)$.
A tuple has type $(d;R)$ if distinct supports are at distance at
least $d>0$ and each support lies in a ball of radius $R>0$;
the centres of these balls may differ. Every separated tuple has
such a type, since its supports are compact.
For a tuple $X=(f_1,\ldots,f_n)$, let $d(X)$ be the least distance between distinct supports, with value $\infty$ when fewer than two supports are nonempty. Let $R(X)$ be the largest of the least radii of closed balls containing the individual supports, with $R(\varnothing)=0$. These radii are attained by compactness of the support and of the set of candidate centres within any fixed radius of a support point. We write $|X|=n$, $\operatorname{supp}X=\bigcup_i\operatorname{supp}f_i$, and $\|X\|_\sharp=\prod_i\|f_i\|_\sharp$. For $J\subseteq\{1,\ldots,n\}$ the tuple $X_J$ retains the entries indexed by $J$ in order. The tuples with smooth entries form $\mathfrak S_n^\infty(O)$.

An isometry $r$ acts on tests by $rf=f\circ r^{-1}$ and entrywise on tuples. The notation $X\cdot Y$ denotes concatenation when the combined tuple is separated. We write $\mathbb C\mathfrak S(O)$ for the free complex vector space with basis $\mathfrak S(O)$. Actions on tuples extend linearly, and concatenation extends bilinearly whenever every pair of basis tuples with nonzero coefficients can be concatenated. Conjugation $F\mapsto\bar F$ conjugates the coefficients of a formal sum. A tuple of complex tests is evaluated by expansion in real and imaginary parts. This expansion and the usual linear relations in each entry hold after applying a multilinear functional.
\end{definition}

A map on $\mathfrak S_n(O)$ is called multilinear if its restriction
to $\prod_i C_c^1(O_i)$ is multilinear whenever the open sets
$O_i\subset O$ have pairwise disjoint closures. We use
\begin{equation}\label{support:9-3}
 \|f\|_\sharp=\max(\|f\|_\infty,40M\operatorname{Lip} f)
\end{equation}
on both spaces. With
$\|f\|_{C^1}=\max(\|f\|_\infty,\|\nabla f\|_\infty)$,
integration along a shortest line segment, lifted on the torus,
gives $\|f\|_\sharp\le40M\|f\|_{C^1}$.

Fix $n_0\ge1$ and finite constants $C_n(d;R)\ge0$ for
$n\ge n_0$ and $d,R>0$. A family $\Phi=(\Phi_n)_{n\ge n_0}$
obeys these bounds if
\begin{equation}\label{support:10-9}
 |\Phi_n(f_1,\ldots,f_n)|
 \le C_n(d;R)\prod_i\|f_i\|_\sharp
\end{equation}
on every tuple of type $(d;R)$. We impose no bound as $d$ tends
to zero.

\begin{proposition}[Simultaneous subsequence extraction]\label{support:10-25}
Let $O_j\subset\mathbb E$ be open sets such that every compact
subset of $\mathbb E$ is contained in $O_j$ for all sufficiently
large $j$. Suppose that the real multilinear families
$\Phi^{(j)}$ on $\mathfrak S(O_j)$ satisfy
\eqref{support:10-9} for $n\ge n_0$, with the same constants,
once $j\ge j_0(n,R)$. Then there is a subsequence along which
$\Phi^{(j)}_n(X)$ converges for every $n\ge n_0$ and every
$X\in\mathfrak S_n(\mathbb E)$. The subsequence can also be
chosen so that any prescribed countable collection of bounded
real sequences converges along it. The limiting family is
multilinear and satisfies the same bounds. Symmetry is preserved.
\end{proposition}

\begin{proof}
\emph{A countable family with controlled supports.}
Let $\mathfrak O$ consist of finite unions of rational open boxes,
using their projections on the torus. These sets have compact
closures and form a countable basis under finite unions.
For each $W\in\mathfrak O$ choose a countable $C^1$-dense subset
$\mathfrak f(W)\subset C_c^1(W)$.

Here is the separability argument with the support fixed.
For $f\in C_c^1(W)$ consider
$g_f=(f,\partial_1f,\partial_2f,\partial_3f)$ on $\overline W$.
The $C^1$ distance of two functions is the uniform distance of
these maps for the norm
$|(v_0,v')|_* =\max(|v_0|,|v'|)$.
For each integer $p\ge1$ choose a finite $1/p$-net in
$\overline W$. Classify the functions by a common modulus bound
$|g_f(x)-g_f(y)|_*\le1/q$ for $|x-y|\le1/p$, and by rational
approximations within $1/q$ to their values at the points of the
net. There are countably many classes. Each has diameter at most
$4/q$, by comparison through a nearest net point. Uniform continuity
of $g_f$ on $\overline W$ places every $f$ in such a class for each
$q$, with $p$ sufficiently large. Choosing one member of each
nonempty class therefore gives the required dense subset. For
$W=\varnothing$ take $\mathfrak f(W)=\{0\}$.

Take all tuples of all lengths with entries in
$\mathfrak f(W_i)$, where the $W_i$ have pairwise disjoint closures.
This is a countable family; enumerate it by $(X^k)_{k\ge1}$.
Given a tuple $X=(f_1,\ldots,f_n)$ of type $(d;R)$, compactness
of each support permits a choice of such sets with
\begin{equation}\label{support:10-16}
 \operatorname{supp}f_i\subset W_i
 \subset\{x:\operatorname{dist}(x,\operatorname{supp}f_i)<d/3\}.
\end{equation}
Their mutual distances are at least $d/3$, and each lies in a ball
of radius $R+d/3$. The union of their closures is compact.
For a zero entry choose $W_i=\varnothing$.

\emph{Uniform approximation of evaluations.}
If $f_i,h_i$ have supports in $W_i$, multilinearity gives
\[
 \Phi_n(f_1,\ldots,f_n)-\Phi_n(h_1,\ldots,h_n)
 =\sum_{i=1}^n\Phi_n(h_1,\ldots,h_{i-1},f_i-h_i,
                                  f_{i+1},\ldots,f_n).
\]
Every tuple on the right retains separation $d/3$.
Choose $h_i\in\mathfrak f(W_i)$ with
$\|f_i-h_i\|_{C^1}\le\epsilon\le1$. For every family obeying
\eqref{support:10-9} at $(n,R+d/3)$, the preceding identity gives
\begin{equation}\label{support:10-18}
 |\Phi_n(X)-\Phi_n(h_1,\ldots,h_n)|
 \le C_n(d/3;R+d/3)\,n(40M)^n\epsilon
            (1+\max_i\|f_i\|_{C^1})^{n-1}.
\end{equation}
The constant does not depend on the approximating family or its
domain, provided the latter contains all $W_i$.

\emph{Extraction and extension.}
For each $X^k$ of length at least $n_0$, its evaluation by $\Phi^{(j)}$ is defined and
bounded for all sufficiently large $j$. Assign arbitrary values
to the finitely many undefined entries. Apply the diagonal
subsequence argument to these countably many bounded real sequences
and to the additional sequences in the statement. Each selected
evaluation then converges. For a fixed arbitrary $X$, use
\eqref{support:10-18} to approximate it by one $X^k$, with an error
uniform in all sufficiently large $j$. The triangle inequality
shows that the evaluations at $X$ are Cauchy. This proves convergence
on every separated tuple along the same subsequence.

Each multilinearity or symmetry identity involves only finitely
many evaluations, all eventually defined, so it passes to the limit.
The same is true of \eqref{support:10-9} for every fixed tuple and type.
\end{proof}

\begin{proposition}\label{support:10-28}
The set of symmetric multilinear families on $\mathfrak S(\mathbb E)$
satisfying \eqref{support:10-9} is compact and metrizable for pointwise
convergence on separated tuples.
\end{proposition}
\begin{proof}
Use the determining family $X^k$ in the preceding proof and let $n_k$
be the length of $X^k$. A metric is
\begin{equation}\label{support:10-27}
 d(\Phi,\Psi)=\sum_{k:n_k\ge n_0}2^{-k}
        \min\{1,|\Phi_{n_k}(X^k)-\Psi_{n_k}(X^k)|\}.
\end{equation}
The triangle inequality follows term by term. If this expression
vanishes, \eqref{support:10-18} makes every evaluation of the two families
equal, so it is a metric. That same estimate shows that every
evaluation is continuous for the metric: approximate its tuple by
one $X^k$, uniformly over all the families. Conversely, the tail
of the series is uniformly bounded by $\sum_{k>N}2^{-k}$, and its
first $N$ terms are continuous for pointwise convergence. The two
topologies thus agree. By \cref{support:10-25}, every sequence has a
convergent subsequence whose limit still belongs to the set. Sequential
compactness in this metric proves compactness.
\end{proof}


\begin{definition}[Functionals and cumulants]\label{support:9-12}
In (1), for a finite set $J$, let $\operatorname{Part}(J)$ be its set of partitions and $\operatorname{Part}_{\ge2}(J)$ the partitions with all blocks of size at least two. For $J=\{1,\ldots,n\}$ write $\operatorname{Part}(n)$ and $\operatorname{Part}_{\ge2}(n)$. The empty set has one partition, with no blocks. In (2), multilinearity is on products of test spaces with separated closures and symmetry means permutation invariance. In (3), a functional $\sigma$ is real and multilinear on separated tuples, with $\sigma(\varnothing)=1$. For $n\ge1$, its connected functions are
\begin{equation}\label{support:9-13}
S_n^\sigma(X)=\sum_{\pi\in\operatorname{Part}(n)}(-1)^{|\pi|-1}(|\pi|-1)!\prod_{J\in\pi}\sigma(X_J).
\end{equation}
In (4), for a family $S=(S_n)_{n\ge2}$, the centred functional $\sigma_S$ is defined by \eqref{support:9-14}. For a functional $\sigma$, we write $\sigma^{\rm c}=\sigma_{(S_n^\sigma)_{n\ge2}}$. The notation $\mathsf{Fct}(\underline C,\bar c)$ denotes symmetric functionals whose one-point function is $c\int f$ with $|c|\le\bar c$ and whose connected functions obey the table $\underline C=(C_n(d;R))_{n\ge2}$ of \eqref{support:10-9}. The corresponding symmetric connected families form $\mathcal K_{\mathbb E}(\underline C)$.
\end{definition}

\begin{lemma}[Moment inversion and centring]
Every symmetric functional obeys
\begin{equation}\label{support:9-21}
\sigma(X)=\sum_{A\subseteq\{1,\ldots,n\}}\prod_{i\in A}\sigma((f_i))\,\sigma^{\rm c}(X_{A^c}).
\end{equation}
In particular, if $\sigma((f))=c\int f$, the map $\operatorname{sh}_c$ replaces the entry $(f)$ by $(f)+c(\int f)\varnothing$ in each position and $\sigma=\sigma^{\rm c}\circ\operatorname{sh}_c$. It is invertible with inverse $\operatorname{sh}_{-c}$. The shift commutes with isometries and with concatenation on separated arguments.
\end{lemma}
\begin{proof}
In the commutative algebra on variables $z_1,\ldots,z_n$ subject to $z_i^2=0$, let $z_J=\prod_{i\in J}z_i$ and $M=\sum_J\sigma(X_J)z_J$. The coefficient of $z_J$ in $\log M$ is the cumulant of the subtuple $X_J$, since an unordered partition with $k$ blocks appears $k!$ times in $(M-1)^k$. Exponentiating gives the partition expansion for moments. Grouping the partitions by their singleton blocks gives \eqref{support:9-21}. The shift identities follow by expanding the product of the replacements. Isometries preserve $\int f$.
\end{proof}

\begin{definition}[Mirrors and reflection positivity]\label{support:9-25}
In (1), a mirror in $\mathbb R^3$ $(u,c)$, $|u|=1$, has reflection $\theta x=x-2(u\cdot x-c)u$ and open side $\{x:u\cdot x>c\}$. Its opposite mirror is $(-u,-c)$. The normals of coordinate and diagonal mirrors are $\pm e_\mu$ and $(\pm e_\mu\pm e_\nu)/\sqrt2$, $\mu\ne\nu$, respectively. In (2), a mirror on $T_\ell$ $(\mu,c)$ is the projection of the two parallel planes with $x_\mu=c$ and $x_\mu=c+\ell/2$. Its reflection is induced by the coordinate reflection about $x_\mu=c$ and its open side is the projection of $c<x_\mu<c+\ell/2$. The opposite mirror is $(\mu,c+\ell/2)$. All torus formulas descend since a change of lift shifts coordinates by periods.

In (3), a functional is reflection positive at the mirror if for all $F\in\mathbb C\mathfrak S(\mathbb E_+)$ on its open side,
\begin{equation}\label{support:9-28}
\sigma(\theta\bar F\cdot F)\ge0 .
\end{equation}
The inequality means a real nonnegative number. Each concatenation is separated since the reflection maps the open side to the opposite open side.
\end{definition}
\begin{proposition}[Transport of reflection positivity]\label{support:9-31}
For a real symmetric functional, reflection positivity passes to (a) the opposite side, (b) every translate of the mirror by a translation preserving the functional, and (c) every image under a linear isometry preserving the functional. On the torus the linear isometry must preserve the period lattice. Consequently, (d) invariance under translations and $W_3$ gives positivity at every mirror of the same kind as a positive mirror. Full Euclidean invariance and positivity at one mirror give positivity at every plane.
\end{proposition}
\begin{proof}
In the matrix of the reflected form on the opposite side substitute $X_i=\theta Y_i$ with $Y_i$ on the original side. The entries become $\sigma(Y_i\cdot\theta Y_j)=\sigma(\theta Y_j\cdot Y_i)$ by symmetry. The resulting matrix is the transpose of the original positive semidefinite matrix and is positive semidefinite. For an isometry $r$ carrying one mirror to another, substitute $X_i=rY_i$ and use $\theta_{r\Sigma}r=r\theta_\Sigma$. The transports in (b) and (c) follow by invariance. The cubic group acts transitively on the coordinate normals and on the diagonal normals, and $\operatorname{O}(3)$ acts transitively on all unit normals. This proves the last assertions.
\end{proof}
\subsection{Distributional extensions and growth}
\label{sec:distribution-extension}
\begin{remark}[Admissible geometric constants]\label{support:26-1}
The estimates on separated tests and their flat Schwartz extension below hold for $M\ge1$. In particular, they hold for every odd $L\ge13$ and $M=L^{m'}$ with $m'\ge1$.
\end{remark}
\begin{definition}[Regularity on separated tests]\label{support:25-28}
Write $E^n_{\ne}$ for the configurations of $n$ distinct points in $E$. A functional is regular of order at most $k$ on compact subsets of this set if its values on separated smooth tensors have a distributional extension there of order at most $k$ on each compact set. Multilinearity in this definition is multilinearity on products of spaces of tests with mutually separated supports as above.
\end{definition}

The bounds on separated tests define distributions away from
coincident points. Their power dependence on the separation also
controls Schwartz functions flat on the diagonals. We prove both
facts by one local Fourier construction. The dyadic partition used
below keeps the supports of the individual Fourier factors separated;
flatness then makes its fine-scale sum converge.

Write $E^n_{\ne}=\{(x_1,\ldots,x_n):x_i\ne x_j\ (i\ne j)\}$,
where $E=\mathbb R^3$ or $T_\ell$, and set
$m(x)=\min_{i\ne j}|x_i-x_j|_\infty$, with $m=\infty$ for $n=1$.
We complexify multilinear forms by complex linearity in each entry.
Thus a bound by $c\prod_i\|f_i\|_\sharp$ on real entries gives
the bound $c\prod_i(\|\operatorname{Re}f_i\|_\sharp+
\|\operatorname{Im}f_i\|_\sharp)$ on complex entries.

\subsubsection{A partition adapted to the diagonals}

Put $v(t)=e^{-1/t}$ for $t>0$ and $v(t)=0$ for $t\le0$,
$\vartheta(t)=v(1-16t^2)$, and $c=\int\vartheta$.
The two cutoffs
\begin{equation}\label{support:25-9}
 \chi(t)=c^{-1}\int_{t-1/2}^{t+1/2}\vartheta(s)\,ds,
 \qquad
 \widetilde\chi(t)=c^{-1}\int_{t-1}^{t+1}\vartheta(s)\,ds
\end{equation}
have values in $[0,1]$. The support of $\chi$ is $[-3/4,3/4]$,
$\sum_{z\in\mathbb Z}\chi(t-z)=1$, and at most two summands
are nonzero. The second cutoff equals one on $[-3/4,3/4]$,
vanishes outside $[-5/4,5/4]$, and satisfies
$|\widetilde\chi'|\le16$. We shall use
\begin{equation}\label{support:25-5}
 |v^{(k)}|\le8^k(k!)^2,\qquad
 |\vartheta^{(k)}|\le200^k(k!)^2,
 \qquad |\chi^{(k)}|\le200^k(k!)^2.
\end{equation}
Here is the derivative estimate. On $t>0$,
$v^{(k)}(t)=P_k(1/t)e^{-1/t}$, where
$P_{k+1}(s)=s^2(P_k(s)-P_k'(s))$ and $P_0=1$.
The weighted coefficient sum $A_k=\sum_j|[s^j]P_k|j!$
satisfies $A_{k+1}\le2(2k+1)^2A_k\le8(k+1)^2A_k$;
also $s^je^{-s}\le j!$. All derivatives tend to zero at
$t=0$, so this proves the first bound globally. For the second,
the chain rule gives
\[
 \vartheta^{(k)}(t)=
 \sum_{0\le2j\le k}\frac{k!}{j!(k-2j)!}
       (-32t)^{k-2j}(-16)^jv^{(k-j)}(1-16t^2).
\]
On $|t|<1/4$, use $|32t|\le8$,
$k!/(j!(k-2j)!)\le3^k j!$, and
$j!((k-j)!)^2\le(k!)^2$. The resulting sum is bounded by
$(32/31)192^k(k!)^2\le200^k(k!)^2$ for $k\ge1$.
Outside this interval it vanishes. Finally
$c\ge\tfrac14e^{-4/3}>1/16$, and differentiating
\eqref{support:25-9} proves the last bound and the stated properties
of the cutoffs.

Let $h_q=2^{-q}/16$, $q\ge0$. On $\mathbb R^3$ use the grid
$h_q\mathbb Z^3$; on $T_\ell$ use its quotient and periodize
the cutoffs. Here $\ell=2L^m$, $m\ge0$, so $\ell/h_q$ is
an integer and the cutoff supports embed in the torus. Define
\[
 \chi_{q,z}(y)=\prod_{\mu=1}^3\chi(y_\mu/h_q-z_\mu),
 \qquad
 \widetilde\chi_{q,z}(y)=
             \prod_{\mu=1}^3\widetilde\chi(y_\mu/h_q-z_\mu),
\]
with the periodized interpretation on the torus. Let $Q_{q,z}$
and $\widetilde Q_{q,z}$ be the closed cubes of sup radii
$3h_q/4$ and $5h_q/4$ about $h_qz$.
An index tuple $\mathbf z=(z_1,\ldots,z_n)$ is good when
$D(\mathbf z):=\min_{i\ne j}|z_i-z_j|_\infty\ge6$.
Write $\chi_{q,\mathbf z}=\bigotimes_i\chi_{q,z_i}$ and
\begin{equation}\label{support:25-18}
 G_q=\sum_{\mathbf z\text{ good}}\chi_{q,\mathbf z},
 \qquad G_{-1}=0,
 \qquad \Upsilon_q=G_q(1-G_{q-1}).
\end{equation}
The enlarged cubes in a good tuple have mutual distance at least
$(D(\mathbf z)-5/2)h_q\ge7h_q/2$ and each lies in a ball of
radius less than one. The grid identities give
\begin{gather}
 G_q=0\quad(m\le9h_q/2),\qquad
 G_q=1\quad(m\ge15h_q/2),\label{support:25-22}\\
 G_qG_{q-1}=G_{q-1},\qquad
 \Upsilon_q=G_q-G_{q-1}\ge0,
 \qquad \sum_{q\ge0}\Upsilon_q=1\quad\text{on }E^n_{\ne}.
 \label{support:25-23}
\end{gather}
Indeed, a nonzero grid factor has
$|x_i-h_qz_i|_\infty<3h_q/4$, so the distance of two
centres differs from that of their points by less than $3h_q/2$.
This proves the two plateaux. If $G_{q-1}\ne0$, then
$m>9h_q$, which proves the product identity. In particular,
for $q\ge1$,
\[
 \operatorname{supp}\Upsilon_q\subset
       \{9h_q/2\le m\le15h_q\}.
\]
The partition is locally finite on $E^n_{\ne}$.

For a pure derivative in any coordinate, the product rule and
\eqref{support:25-5} yield
\begin{equation}\label{support:25-25}
 |\partial_{i\mu}^k
       (\chi_{q,\mathbf z}(1-G_{q-1})F)|
 \le2\cdot600^k(k!)^2
       \max_{0\le j\le k}h_q^{j-k}|\partial_{i\mu}^jF|.
\end{equation}
For the derivatives of $G_q$, only two terms in the
differentiated coordinate can contribute; the sums in every
other coordinate are one. Thus
$|\partial_{i\mu}^kG_q|\le2h_q^{-k}200^k(k!)^2$.
For mixed derivatives of a sum of any subfamily of the grid
tensors, the corresponding bound is
$2^{3n}h_q^{-|\alpha|}200^{|\alpha|}(|\alpha|!)^2$.

One count will supply summability over positions. For $q\ge1$
and each unit box $\mathbf w+[0,1)^{3n}$, there are at most
\begin{equation}\label{support:25-27}
 \binom n2\,33^3h_q^{-(3n-3)}
\end{equation}
good $\mathbf z$ with $h_q\mathbf z$ in that box and
$Q_{q,\mathbf z}$ meeting $\operatorname{supp}\Upsilon_q$.
At an intersection point some pair has distance at most $15h_q$;
the two centres therefore differ by at most $16h_q$ in each
coordinate. Choose the pair, the other $n-1$ centres, and one
of $33^3$ differences. At $q=0$ the unrestricted count is
$16^{3n}$.

\subsubsection{The local Fourier construction}

\begin{proposition}[Kernel on separated points]\label{support:25-55}
Let $T$ be a real multilinear form on separated smooth compactly
supported $n$-tuples over $E$. Suppose that for every $d>0$
there is $c_T(d)<\infty$ such that
\begin{equation}\label{support:25-31}
 |T(f_1,\ldots,f_n)|\le c_T(d)\prod_i\|f_i\|_\sharp
\end{equation}
whenever the supports have mutual distance at least $d$ and
each lies in a ball of radius one. There is a unique distribution
$\widetilde T$ on $E^n_{\ne}$ such that
$\widetilde T(\bigotimes_i f_i)=T(f_1,\ldots,f_n)$.
It has order at most $4n+1$ on compact sets, with a bound
depending only on that compact set, $M$, and the function $c_T$.
Symmetries of $T$ under permutations or isometries are inherited
by $\widetilde T$.
\end{proposition}
\begin{proof}
Fix a good box at scale $q$ and put $h=h_q$, $\lambda=5h/2$.
Lift its enlarged cubes to cubes $C_i\subset\mathbb R^3$ of
side $\lambda$. For $\omega\in\mathbb Z^3$ let
$e_\omega(t)=e^{2\pi i\omega\cdot t/\lambda}$ and let
$E_{i,\omega}$ be $\widetilde\chi_{q,z_i}e_\omega$, interpreted
through the lift and periodized if necessary. For a smooth $F$
supported in $Q_{q,\mathbf z}$, extend its lift by zero and set
\[
 \widehat F(\omega)=\lambda^{-3n}
       \int_{\prod_iC_i}F(t)\prod_i\overline{e_{\omega_i}(t_i)}\,dt,
 \qquad
 [F]_{q,k}=\max\left\{\|F\|_\infty,
       (\lambda/2\pi)^k\max_{i,\mu}
                          \|\partial_{i\mu}^kF\|_\infty\right\}.
\]
Integration by parts in a coordinate attaining
$|\omega|_\infty$ gives
\[
 |\widehat F(\omega)|\le
       2^k(1+|\omega|_\infty)^{-k}[F]_{q,k}.
\]
The real and imaginary parts of $E_{i,\omega}$ have sharp norm
at most $1109Mh^{-1}(1+|\omega|_\infty)$: the cutoff gradient
is at most $16\sqrt3/h$, and differentiating the exponential
adds at most $16\sqrt3|\omega|_\infty/h$.
Also, grouping by $|\omega|_\infty$ gives
$\sum_{\omega\in\mathbb Z^{3n}}(1+|\omega|_\infty)^{-3n-1}
\le3n8^n$. With $k=4n+1$ these bounds prove absolute
convergence of
\begin{equation}\label{support:25-41}
 T_{q,\mathbf z}(F)=
 \sum_{\omega\in(\mathbb Z^3)^n}
       \widehat F(\omega)T(E_{1,\omega_1},\ldots,E_{n,\omega_n})
\end{equation}
and the estimate
\begin{equation}\label{support:25-49}
 |T_{q,\mathbf z}(F)|\le
 (2\cdot10^6)^nM^nc_\square h^{-n}[F]_{q,4n+1}.
\end{equation}
Here $c_\square$ can be either $c_T(7h/2)$ or
$c_T((D(\mathbf z)-5/2)h)$, whenever the argument is finite;
for $n=1$ any valid uniform one-entry bound serves.
The numerical factor before $M^nc_\square h^{-n}[F]$ is
$6n\,283904^n\le(2\cdot10^6)^n$.
Changing a lift multiplies the entries by phases and the Fourier
coefficient by their inverse product; hence \eqref{support:25-41}
is independent of the lifts.

The Fourier partial sums, multiplied by
$\bigotimes_i\widetilde\chi_{q,z_i}$, converge to $F$ with
all derivatives and have support in $\widetilde Q_{q,\mathbf z}$.
For completeness, integration by parts makes the Fourier series
of every derivative absolutely summable. Its sum is that
derivative: the product Fej\'er kernels are nonnegative approximate
identities, so their convolution converges uniformly to a continuous
periodic function and also to the sum of an absolutely convergent
Fourier series. The cutoff equals one on the support of $F$,
which proves the claimed convergence after multiplication.
If $F=\bigotimes_i f_i$ with
$\operatorname{supp}f_i\subset Q_{q,z_i}$, its coefficients
factor. The partial sums therefore give $T$ of the individual
Fourier approximations; their convergence in the sharp norm and
the telescoping estimate prove
$T_{q,\mathbf z}(F)=T(f_1,\ldots,f_n)$.

Define
\begin{equation}\label{support:25-56}
 \widetilde T(F)=\sum_{q\ge0}\sum_{\mathbf z\text{ good}}
 T_{q,\mathbf z}\bigl(\chi_{q,\mathbf z}(1-G_{q-1})F\bigr).
\end{equation}
For $F$ in a fixed compact subset of $E^n_{\ne}$, only
finitely many scales and positions occur. Combining
\eqref{support:25-25} with \eqref{support:25-49} gives the asserted
distribution bound of order $4n+1$.
For a tensor $F=\bigotimes_i f_i$, each piece in
\eqref{support:25-56} is the finite sum
\[
 \bigotimes_i(\chi_{q,z_i}f_i)
 -\sum_{\mathbf z'\text{ good at }q-1}
       \bigotimes_i(\chi_{q,z_i}\chi_{q-1,z_i'}f_i),
\]
where the second sum is absent at $q=0$. All entries stay
inside the original separated supports. Multilinearity and
\eqref{support:25-23} therefore give
$\widetilde T(F)=T(f_1,\ldots,f_n)$. In making this computation,
the linear form induced on finite sums of tensors is well defined:
on the finite-dimensional spans of their factors, point evaluations
span the duals, so a tensor representing the zero function has all
its coefficients zero.

Finally, the same partition and the local Fourier approximation
approximate every test in $C_c^\infty(E^n_{\ne})$ by finite
sums of separated tensors, in all derivatives with supports in
one compact subset of $E^n_{\ne}$. This proves uniqueness.
Applying uniqueness to the pullback by a symmetry proves the
last assertion.
\end{proof}

If the form is originally defined on separated $C_c^1$ tuples
with the same bounds, these values are determined by its kernel.
Indeed, cut each entry by the $q=0$ partition, mollify the
finitely many pieces, and let the mollifier radius tend to zero.
Their supports retain positive separation, their sharp norms
converge, and the telescoping bound recovers the original value.

\subsubsection{Flat tests and the global bound}

For $\varphi\in\mathcal S(\mathbb R^{3n})$ put
\begin{equation}\label{support:26-4}
 |\varphi|_k=\sup_{x,\,|\alpha|\le k}
                  (1+|x|^2)^{k/2}|\partial^\alpha\varphi(x)|.
\end{equation}
\begin{definition}[Flat Schwartz test spaces]\label{support:26-2}
Let $\mathcal S_{\ne}(\mathbb R^{3n})$ consist of Schwartz
functions all of whose derivatives vanish whenever two entries
coincide. Let $\mathcal S_+$ be the subspace additionally flat
outside $V_+^{(n)}=\{x\in(\mathbb R^3)^n_{\ne}:x_{i3}>0\text{ for all }i\}$.
In degree zero both spaces are $\mathbb C$.
\end{definition}

\begin{lemma}\label{support:26-17}
Finite sums of tensors of separated smooth compactly supported
tests are dense in $\mathcal S_{\ne}$ for every norm
$|\cdot|_k$. In $\mathcal S_+$ the approximating entries may
all be supported in the open positive half-space.
For $k\ge4$ and $A\subset\{1,\ldots,n\}$, partial integration
$\psi_A(y)=\int\varphi(y,x_A)\,dx_A$ maps
$\mathcal S_{\ne}(\mathbb R^{3n})$ into
$\mathcal S_{\ne}(\mathbb R^{3(n-|A|)})$ and satisfies
\begin{equation}\label{support:26-19}
 |\psi_A|_{k(n-|A|)}\le C_{\int}(k)^{|A|}|\varphi|_{kn},
 \qquad C_{\int}(k)=\int_{\mathbb R^3}(1+|x|^2)^{-k/2}\,dx.
\end{equation}
In particular $C_{\int}(4)=\pi^2$ and
$C_{\int}(10)=5\pi^2/64$. Also
$|\varphi\otimes\psi|_k\le|\varphi|_k|\psi|_k$.
If both factors belong to $\mathcal S_{\ne}$ and are flat outside
the negative and positive half-spaces respectively, their tensor
is flat on every diagonal.
\end{lemma}
\begin{proof}
Flatness gives the following Taylor estimate. If $\varphi$ is
flat at $x^*$, $|x-x^*|\le r$ with $r\ge1$, and
$N\ge|\alpha|+s$, then
\begin{equation}\label{support:26-15}
 |\partial^\alpha\varphi(x)|\le
 \frac{|x-x^*|_1^s}{s!}(1+2r^2)^{N/2}
       |\varphi|_N(1+|x|^2)^{-N/2}.
\end{equation}
Apply Taylor's formula with integral remainder to
$t\mapsto\partial^\alpha\varphi(x^*+t(x-x^*))$; along this
segment $1+|x|^2\le(1+2r^2)(1+|x^*+t(x-x^*)|^2)$.
When $m(x)$ is small, replace a closest pair by its midpoint.
This gives a diagonal point $x^*$ with
$|x-x^*|_1\le3m(x)$ and $|x-x^*|\le\sqrt{3/2}m(x)$.
In the half-space case one may instead replace a smallest third
coordinate by its minimum with zero.

For the half-space partition restrict the good indices further
by $z_{i3}\ge6$, and denote the resulting sum by $G_q^+$.
The plateau and derivative estimates above hold with
$m^+(x)=\min\{m(x),x_{13},\ldots,x_{n3}\}$.
Taylor's estimate with $s=k+1$, followed by the mixed derivative
bound for the cutoffs, gives
\[
 |(1-G_Q)\varphi|_k\le C(n,k)h_Q|\varphi|_{2k+1},
\]
and the same bound with $G_Q^+$ for $\varphi\in\mathcal S_+$.
Indeed, the derivatives of the cutoff cost at most $h_Q^{-k}$,
whereas Taylor's remainder contributes $h_Q^{k+1}$.
Fix $Q$ to make this error small. Among the remaining finitely
many scales, discard grid centres outside a large box. Their
derivatives are bounded with $Q$ fixed, while one extra Schwartz
weight makes their $|\cdot|_k$ contribution tend to zero.
Only finitely many boxes remain. Approximate each by its local
Fourier tensors as in the preceding proof. Their enlarged
supports are separated, and in the half-space case have third
coordinate at least $19h_q/4$. This proves both density assertions.

For partial integrals, use
\[
 (1+|x|^2)^{kn/2}\ge
 (1+|y|^2)^{k(n-|A|)/2}
                    \prod_{i\in A}(1+|x_i|^2)^{k/2}.
\]
This permits differentiation under the integral, proves
\eqref{support:26-19}, and preserves flatness on the remaining
diagonals. Polar coordinates give the two stated constants.
The tensor bound follows from
$1+|x|^2+|y|^2\le(1+|x|^2)(1+|y|^2)$.
At a diagonal joining the negative and positive groups, one
factor and all its derivatives vanish; diagonals inside a group
are already covered by its flatness.
\end{proof}

\begin{proposition}\label{support:26-24}
Suppose that $T$ satisfies \eqref{support:25-31} on $\mathbb R^3$
with $a\in\mathbb N_0$ and the bound
\begin{equation}\label{support:26-25}
 |T(f_1,\ldots,f_n)|\le c_T\max\{1,d^{-a}\}\prod_i\|f_i\|_\sharp .
\end{equation}
Its kernel has a unique extension to $\mathcal S_{\ne}$
bounded for a Schwartz norm. Use the constant
\begin{equation}\label{support:26-10}
 C_{\ne}=10^{53}.
\end{equation}
Then
\begin{equation}\label{support:26-26}
 |\widetilde T(\varphi)|\le
 C_{\ne}^{n+a}M^n(n!)^8c_T|\varphi|_{8n+a-1}.
\end{equation}
\end{proposition}
\begin{proof}
Use the series \eqref{support:25-56}, now summed over all scales
and positions. Set $k=4n+1$, $s=a+4n-2$, and $N=k+s$.
For $F_{q,\mathbf z}=\chi_{q,\mathbf z}(1-G_{q-1})\varphi$,
\eqref{support:25-25} and $\lambda/(2\pi)<2h_q/5$ give
\begin{equation}\label{support:25-46}
 [F_{q,\mathbf z}]_{q,k}\le
 2\cdot240^k(k!)^2
 \sup_{\operatorname{supp}F_{q,\mathbf z}}
       \max_{i,\mu,\,0\le j\le k}h_q^j|\partial_{i\mu}^j\varphi|.
\end{equation}
For $q\ge1$, this support has $m\le15h_q$. The midpoint
choice in \eqref{support:26-15} has $|x-x^*|_1\le45h_q$ and
$|x-x^*|<1$. Since $N\ge4n$, it follows that
\[
 |\partial_{i\mu}^j\varphi(x)|\le
 \frac{(45h_q)^s}{s!}3^{N/2}|\varphi|_N
                  \prod_{l=1}^n(1+|x_l|^2)^{-2},\qquad j\le k.
\]
To sum the positions, assign $h_q\mathbf z$ to its unit box
$\mathbf w+[0,1)^{3n}$. Its support lies in
$\mathbf w+[-1/20,21/20]^{3n}$, and
\[
 \sum_{\mathbf w}\sup_{x\in\mathbf w+[-1/20,21/20]^{3n}}
          \prod_l(1+|x_l|^2)^{-2}\le(81\pi^2)^n.
\]
For one factor, compare a point in the enlarged unit box with
a point in its unit cube through the common centre: the squared
weights differ by at most $81$. Integrate over that cube and
sum; then take the product in the $n$ entries.
Combining this estimate, \eqref{support:25-27},
\eqref{support:25-49}, and \eqref{support:25-46} bounds the sum at
scale $q\ge1$ by
\begin{align}
 \Sigma_q\le{}&(2\cdot10^6)^nM^nc_T(2/7)^a
 2\cdot240^k(k!)^2\frac{45^s}{s!}3^{N/2}
 \binom n2\,33^3(81\pi^2)^n
 h_q^{s-a-n-(3n-3)}|\varphi|_N .\label{support:26-31}
\end{align}
The exponent of $h_q$ is one. Thus the fine-scale sum converges.

At $q=0$, split the centres into $D(\mathbf z)\ge19$ and
$6\le D(\mathbf z)\le18$. In the first group the enlarged
supports are separated by more than one, so $c_\square=c_T$;
use the Schwartz bound without Taylor's formula and the count
$16^{3n}$. In the second group a pair in the box has distance
at most $39/32$ in the sup norm. Its midpoint has
$|x-x^*|_1\le117/32$ and $|x-x^*|\le3/2$; hence
\eqref{support:26-15} supplies the factor
$(117/32)^s(11/2)^{N/2}/s!$, and
$c_\square=c_T(32/7)^a$.
These give, respectively,
\begin{align*}
 \Sigma_0^{\mathrm A}&\le
 (2\cdot10^6)^nM^nc_T16^{4n}
 2\cdot240^k(k!)^2(81\pi^2)^n|\varphi|_N,\\
 \Sigma_0^{\mathrm B}&\le
 (2\cdot10^6)^nM^nc_T16^{4n}
 2\cdot240^k(k!)^2(81\pi^2)^n
 (32/7)^a\frac{(117/32)^s}{s!}(11/2)^{N/2}|\varphi|_N.
\end{align*}
For $n=1$ the second group and every $q\ge1$ contribution
are empty. In general use
$(k!)^2\le25^n4^{8n}(n!)^8$,
$(2/7)^a3^{N/2}\le81^n$, and $45^s/s!\le3^{45}$
in the fine-scale bound, with $\sum_{q\ge1}h_q=1/16$.
For the second coarse group use
$(32/7)^a(11/2)^{a/2}\le11^a\le11^s$ and
$(11\cdot117/32)^s/s!\le3^{41}$.
The three sums, divided by $M^n(n!)^8c_T|\varphi|_N$,
are at most $10^{28}(10^{28})^n$,
$10^3(10^{30})^n$, and $10^{23}(10^{33})^n$.
For $n\ge2$ their sum is at most $3(10^{45})^n$;
for $n=1$ the sole coarse sum is at most $10^{33}$.
Both are bounded by $C_{\ne}^n$.

The series is therefore absolutely convergent, linear, and
satisfies the stronger bound with $C_{\ne}^n$ in place of
$C_{\ne}^{n+a}$. On compactly supported tests away from the
diagonals it is the previous finite sum. Uniqueness follows from
the density in \cref{support:26-17}.
\end{proof}

\begin{proposition}[Growth of the Schwinger functions]\label{support:26-37}
Suppose a symmetric family $S=(S_n)_{n\ge2}$ on separated
$C_c^1$ tuples satisfies
\[
 |S_n(f_1,\ldots,f_n)|\le
 A\max\{1,d^{-a}\}B^n(n!)^\varkappa\prod_i\|f_i\|_\sharp
\]
when the separation is at least $d$ and each support lies in a
ball of radius one, where $a\in\mathbb N$, $A,B\ge1$, and
$\varkappa\ge0$. Its centred moments $\sigma_S$ extend uniquely
to functionals $\mathsf S_n^{\sigma_S}$ on $\mathcal S_{\ne}$,
with $\mathsf S_0^{\sigma_S}=1$, $\mathsf S_1^{\sigma_S}=0$,
and
\begin{equation}\label{support:26-39}
 |\mathsf S_n^{\sigma_S}(\varphi)|\le
 C_c^n(n!)^{\varkappa+9}|\varphi|_{kn},\qquad
 k=8+\lceil a/2\rceil,\qquad
 C_c=C_{\ne}^{1+a/2}MB\sqrt A.
\end{equation}
If a functional $\sigma$ has these connected functions and
one-point function $c\int f$, with $|c|\le\bar c$, the same
conclusion holds with $C_c$ replaced by
$C=C_c+\bar c\,C_{\int}(k)$.
\end{proposition}
\begin{proof}
In the moment partition sum, there are at most $\lfloor n/2\rfloor$
blocks and no singleton blocks. Also
$\prod_{J\in\pi}|J|!\le n!$, and the number of partitions
is at most $n!$: each partition is the cycle partition of a
permutation, choosing increasing cycles. Thus
\[
 |\sigma_S(f_1,\ldots,f_n)|\le
 A^{\lfloor n/2\rfloor}
 \max\{1,d^{-a\lfloor n/2\rfloor}\}
 B^n(n!)^{\varkappa+1}\prod_i\|f_i\|_\sharp.
\]
Apply \cref{support:26-24} with singularity order
$a\lfloor n/2\rfloor$. The derivative order is
$8n+a\lfloor n/2\rfloor-1\le kn$, and its bound is
\eqref{support:26-39}. For $n=1$ the functional is zero.
For an uncentred functional use the finite expansion
\begin{equation}\label{support:26-40}
 \mathsf S_n^\sigma(\varphi)=
 \sum_{A_0\subset\{1,\ldots,n\}}c^{|A_0|}
             \mathsf S_{n-|A_0|}^{\sigma_S}(\psi_{A_0}).
\end{equation}
It agrees with the original moments on separated tensors.
By \eqref{support:26-19}, each integrated entry costs at most
$C_{\int}(k)$, so the binomial formula gives $C^n$.
Uniqueness again follows from density.
\end{proof}


\begin{proposition}\label{support:26-41}
Suppose the moments of a real symmetric functional $\sigma$
extend to continuous linear functionals $\mathsf S_n$ on
$\mathcal S_{\ne}(\mathbb R^{3n})$, agreeing with $\sigma$
on separated smooth tensors and satisfying
\[
 |\mathsf S_n(\varphi)|\le c_E C_E^n(n!)^\kappa
                           |\varphi|_{k_E n}
\]
for fixed $k_E\in\mathbb N$ and $c_E,C_E,\kappa\ge0$.
In (a) these extensions preserve reality and symmetry. In (b), every Euclidean isometry under which $\sigma$ is invariant also leaves them invariant. In (c), if $\sigma$ is reflection positive, then for
$\varphi_n\in\mathcal S_+(\mathbb R^{3n})$, $0\le n\le n_0$,
with $\mathcal S_+(\mathbb R^0)=\mathbb C$,
\begin{equation}\label{support:26-42}
 \sum_{n,n'=0}^{n_0}
 \mathsf S_{n+n'}(\Theta\bar\varphi_n\otimes\varphi_{n'})
 \ge0.
\end{equation}
Here $\Theta$ reflects every time coordinate.
\end{proposition}
\begin{proof}
Conjugation and permutation preserve the Schwartz seminorms
and flatness on diagonals. The respective identities hold
on separated smooth tensors and extend by
\cref{support:26-17}. For an isometry $rx=\rho x+v$, its
pullback preserves the flatness condition and satisfies
\begin{equation}\label{support:26-43}
 |\varphi\circ(r^{-1})^{\times n}|_k
 \le[6(1+n|v|^2)]^{k/2}|\varphi|_k.
\end{equation}
Indeed each row of $\rho$ has sum of absolute values at
most $\sqrt3$, while
$1+|x|^2\le2(1+n|v|^2)(1+|(r^{-1})^{\times n}x|^2)$.
The chain rule gives the estimate and shows that the
pullback is Schwartz and flat on the diagonals. Density
now transfers the isometry identity.

For positivity, approximate each $\varphi_n$ by finite
sums $\varphi_n^{(j)}$ of separated positive-time smooth
tensors in the seminorm of order $k=2k_{\rm E}n_0$;
leave $\varphi_0$ unchanged. Reflection and conjugation
preserve this seminorm. The tensor estimate and the
flatness assertion in \cref{support:26-17} give
\[
 |\Theta\bar\varphi_n\otimes\varphi_{n'}
  -\Theta\bar\varphi_n^{(j)}\otimes\varphi_{n'}^{(j)}|_k
 \le |\varphi_n-\varphi_n^{(j)}|_k|\varphi_{n'}|_k
   +|\varphi_n^{(j)}|_k|\varphi_{n'}-\varphi_{n'}^{(j)}|_k
 \longrightarrow0.
\]
Since $k_{\rm E}(n+n')\le k$, continuity permits this
limit in every term of \eqref{support:26-42}. For each $j$,
expansion of the finite tensor sums expresses the result
as one reflection-positive quadratic form for $\sigma$,
including the empty tuple. Its limit is nonnegative.
\end{proof}

\begin{corollary}[Uniqueness on the flat test space]\label{support:26-23}
Two linear functionals on $\mathcal S_{\ne}(\mathbb R^{3n})$, each bounded for some Schwartz seminorm, agree if their values on separated smooth compactly supported tensors agree.
\end{corollary}
\begin{proof}
Take the larger of the two orders of the seminorms and use the density statement in \Cref{support:26-17}.
\end{proof}
\begin{definition}[Growth condition]\label{support:26-33}
The Schwinger functionals $\mathsf S_n$ of a symmetric functional satisfy the growth condition with $(k_{\rm E},\mathsf c_{\rm E},\mathsf C_{\rm E},\varkappa_{\rm E})$ if their restrictions to separated smooth tensors are its moments and, for each $n\ge1$ and $\varphi\in\mathcal S_{\ne}(\mathbb R^{3n})$,
\begin{equation}\label{support:26-34}
|\mathsf S_n(\varphi)|\le \mathsf c_{\rm E}\,\mathsf C_{\rm E}^{n}(n!)^{\varkappa_{\rm E}}|\varphi|_{k_{\rm E}n} .
\end{equation}
Here $k_{\rm E}\ge1$ is an integer and the other entries are finite nonnegative numbers. In degree zero $\mathsf S_0(z)=z$. \Cref{support:26-23} gives uniqueness for each degree.
\end{definition}
\section{Continuum reflection and rotations}\label{app:continuum-rotations}

\subsection{Reflection positivity and the translation spectrum}
\label{sec:reflection-hilbert}
\begin{lemma}[Reconstruction from a positive form]\label{support:27-11}
Let $\mathcal D$ carry a positive semidefinite Hermitian form, antilinear in its first argument, and a unit vector $\mathsf 1$. Suppose $T_t^\circ$ and $U^\circ(v)$ are commuting linear representations of $[0,\infty)$ and $\mathbb R^N$ on $\mathcal D$, with both identities represented by the identity and both representations fixing $\mathsf 1$. Assume
\begin{enumerate}
\item[(1)] $T_t^\circ$ is symmetric for the form and $U^\circ(v)$ preserves it.
\item[(2)] For every $F$, $c_F=\sup_{t\ge0}|\langle F,T_t^\circ F\rangle|<\infty$.
\item[(3)] $\langle F,T_t^\circ F\rangle$ and $\langle F,U^\circ(v)F\rangle$ tend to $\|F\|^2$ as $t\downarrow0$ and $v\to0$, respectively.
\end{enumerate}
In the completion $\mathcal H$ of the null-space quotient, (a) the induced $T_t$ are a strongly continuous semigroup of positive self-adjoint contractions, and $U(v)$ a strongly continuous unitary group. (b) $T_t=e^{-tH}$ for a unique $H\ge0$, with $H[\mathsf 1]=0$. (c) There is a joint projection-valued measure of $H$ and the momentum generators, supported in $[0,\infty)\times\mathbb R^N$. (d) Any group of form-preserving maps on $\mathcal D$ fixing $\mathsf 1$ induces a unitary action. When the original maps commute with time or intertwine spatial translations, the induced action has the same property.
\end{lemma}
\begin{proof}
Cauchy--Schwarz gives
\begin{equation}\label{support:27-12}
 \|T_t^\circ F\|^2=\langle F,T_{2t}^\circ F\rangle
       \le\|F\|\,\|T_{2t}^\circ F\|.
\end{equation}
Iterating $j$ times gives $\|T_t^\circ F\|\le \|F\|^{1-2^{-j}}c_F^{2^{-j-1}}$. For non-null $F$ send $j$ to infinity, and for null $F$ the first inequality already shows the image is null. The $U^\circ$ maps preserve the null space and have isometric inverses. The identity $T_t=T_{t/2}^2$ gives positivity. At zero, the norm-square of $(T_t-1)[F]$ is bounded by $2\|F\|^2-2\operatorname{Re}\langle F,T_t^\circ F\rangle$, and the analogous expression for $U(v)$ is exact. This proves strong continuity on a dense subspace and hence on the completion. The spectral theorem for a strongly continuous self-adjoint contraction semigroup gives $H=-\log T_1\ge0$, and the joint spectral theorem with the commuting unitary translation group gives (c). The fixed-vector assertion follows by the spectral theorem from $T_t[\mathsf 1]=[\mathsf 1]$. The final assertion follows by passing each form-preserving invertible map and each intertwining identity to the quotient and then to the completion.
\end{proof}

The bounds on separated correlations suffice to construct the
Hilbert space of a reflection. The additional positivity at a
diagonal plane then restricts the joint energy--momentum spectrum.
This construction uses neither coincident-point moments nor a
probability measure on continuum fields.

Let $\sigma$ be a real, symmetric, normalized functional on
$\mathfrak S(\mathbb R^3)$, multilinear on products of test spaces
with disjoint closures. Assume that its one-point function is
$\sigma(f)=c\int f$ and that its connected functions satisfy a
bound table $C_n(d;R)$ as in \cref{app:separated}.
Thus, on tuples of separation at least $d$ whose individual
supports lie in balls of radius $R$,
\begin{equation}\label{support:27-21}
 \sigma(X)=\sum_{A\subset\{1,\ldots,n\}}
       \prod_{i\in A}\left(c\int f_i\right)
       \sum_{\pi\in\operatorname{Part}_{\ge2}(A^c)}
                         \prod_{J\in\pi}S^\sigma_{|J|}(X_J).
\end{equation}
In particular the absolute value is bounded by the same finite
sum with $|c|\|f_i\|_1$ in place of $c\int f_i$ and
$C_{|J|}(d;R)\prod_{i\in J}\|f_i\|_\sharp$ in place of each
connected function. Multilinear telescoping shows that $\sigma$
is continuous when the entries converge in the sharp norm and
their supports remain in fixed separated compact sets.

For a unit vector $u$, put $E_+(u)=\{x:u\cdot x>0\}$,
let $\theta_u x=x-2(u\cdot x)u$, and let
$\mathcal D_u=\mathbb C\mathfrak S(E_+(u))$ be the free
complex vector space on the separated tuples in this half-space.
Concatenation is used only when the resulting tuple is separated.
The product $\theta_u\bar F\cdot G$ is always defined for
$F,G\in\mathcal D_u$. Reflection positivity means
\[
 \sigma(\theta_u\bar F\cdot F)\ge0\qquad(F\in\mathcal D_u).
\]
It makes
\begin{equation}\label{support:27-19}
 \langle F,G\rangle_{\sigma,u}
       =\sigma(\theta_u\bar F\cdot G)
\end{equation}
a positive semidefinite Hermitian form, antilinear in its first
argument. Quotient by its null space and complete to obtain
$\mathcal H(\sigma,u)$. Write $[F]$ for a class and
$\Omega=[\varnothing]$, so that $\|\Omega\|=1$.

\begin{proposition}[The translation operators]\label{support:27-22}
Suppose that $\sigma$ is translation invariant and reflection
positive at $u\cdot x=0$, with the bound table above. In (a), the reflection form and the uncompleted translation maps satisfy the hypotheses of \Cref{support:27-11}, with $N=2$. In (b), $\mathcal H(\sigma,u)$ is separable. In (c), the maps
\[
 T_t[F]=[\tau_{tu}F]\quad(t\ge0),\qquad
 U(v)[F]=[\tau_vF]\quad(v\in u^\perp)
\]
extend to a strongly continuous semigroup of positive
self-adjoint contractions and a strongly continuous unitary
group. They commute and fix $\Omega$. There are strongly
commuting self-adjoint operators $H\ge0,P_1,P_2$ such that
\begin{equation}\label{support:27-13}
\begin{aligned}
 T_t&=e^{-tH},& U(v)&=e^{i(v_1P_1+v_2P_2)},\\
 e^{i\tau H}U(v)&=\int e^{i\tau E+iv\cdot q}\,d\mathcal P(E,q),&
 T_tU(v)&=\int e^{-tE+iv\cdot q}\,d\mathcal P(E,q),
\end{aligned}
\end{equation}
where coordinates are taken in any fixed orthonormal basis of
$u^\perp$, and $\mathcal P$ is carried by
$[0,\infty)\times\mathbb R^2$. Moreover $\ker T_t=0$,
$H\Omega=P_1\Omega=P_2\Omega=0$, and
$T_t\to P_{\ker H}$ strongly as $t\to\infty$.

In (d), if an orthogonal transformation $r$ fixes $u$ and leaves
$\sigma$ invariant, $[F]\mapsto[rF]$ extends to a unitary
$U(r)$ fixing $\Omega$. It commutes with $H$ and satisfies
$U(r)U(v)U(r)^*=U(rv)$. The joint spectrum is invariant under
$(E,q)\mapsto(E,rq)$.
\end{proposition}

\begin{proof}
On $\mathcal D_u$ write $T_t^0=\tau_{tu}$ and
$U^0(v)=\tau_v$. The identities
$\theta_u\tau_{tu}=\tau_{-tu}\theta_u$ and
$\theta_u\tau_v=\tau_v\theta_u$, together with translation
invariance, give
\[
 \langle F,T_t^0G\rangle=\langle T_t^0F,G\rangle,
 \qquad
 \langle U^0(v)F,U^0(v)G\rangle=\langle F,G\rangle.
\]
For each $F=\sum_i z_iX_i$ the number
\[
 c_F=\sup_{t\ge0}|\langle F,T_t^0F\rangle|
\]
is finite. Choose $\delta>0$ below the Euclidean separations
of the finitely many tuples $\theta_uX_i\cdot X_j$, and choose
$R>0$ bounding the radii of their individual supports.
Put $d=\delta/\sqrt3$. Within each half-space the Euclidean
separations remain at least $\delta$: the negative entries stay
fixed and all positive entries receive the same translation.
Across the mirror,
\[
 |y+tu-x|^2=|y-x|^2+2t\,u\cdot(y-x)+t^2\ge|y-x|^2
 \quad(u\cdot x<0<u\cdot y).
\]
Thus $|y+tu-x|_\infty\ge |y+tu-x|/\sqrt3\ge d$ for all
$t\ge0$. Every translated tuple has sup separation at least
$d$. The sharp and $L^1$ norms and the radii of individual
supports are unchanged. The bound following
\eqref{support:27-21}, summed with coefficients $|z_i z_j|$,
therefore bounds $c_F$.

For small translations, enlarge the finitely many supports by
less than $d/4$. They remain separated by $d/2$. Translations
of a compactly supported $C^1$ function converge in the sharp
norm, by uniform continuity of the function and its first
derivatives. Hence the continuity following
\eqref{support:27-21} gives
$\langle F,T_t^0F\rangle\to\|F\|^2$ and
$\langle F,U^0(v)F\rangle\to\|F\|^2$ at the origin.
Thus all three hypotheses of \Cref{support:27-11} hold.
That lemma gives the strongly continuous translations and their
generators. The joint representation
$e^{isH}U(v)$ has spectral measure \eqref{support:27-13}.
The fixed vector $\Omega$ has spectral measure at $(0,0)$.
The spectral theorem gives $\ker T_t=0$ for finite $t$ and
$T_t\to P_{\ker H}$ strongly as $t\to\infty$.

For separability use the countable family of separated test
tuples constructed in \cref{app:separated}.
Given $X\in\mathfrak S(E_+(u))$, choose $d>0$ smaller than
all sup separations in $\theta_uX\cdot X$. Rational neighborhoods
of its positive supports of sup width less than $d/(4\sqrt3)$
still lie in $E_+(u)$: on those supports $u\cdot x\ge d/2$,
and $|u\cdot v|\le\sqrt3\,|v|_\infty$.
Approximate each entry in $C^1$ inside these neighborhoods.
The four terms in
\[
 \|[X]-[Y]\|^2
 =\sigma(\theta_uX\cdot X)-\sigma(\theta_uX\cdot Y)
  -\sigma(\theta_uY\cdot X)+\sigma(\theta_uY\cdot Y)
\]
then converge to the same value. Rational complex linear
combinations of the resulting countable set of classes are
dense. Finally $r$ commutes with $\theta_u$ and preserves
the form. Its unitary extension has the asserted covariance
on the dense set of classes, hence everywhere. Spectral
uniqueness gives the covariance of $\mathcal P$.
\end{proof}

\subsubsection{Changing the one-point constant}

For real $b$ define, on a separated tuple $X=(f_1,\ldots,f_n)$,
\begin{equation}\label{support:28-13}
 \mathsf{sh}_bX=
   \sum_{A\subset\{1,\ldots,n\}}\prod_{i\in A}
                         \left(b\int f_i\right)X_{A^c}.
\end{equation}
The map extends linearly, preserves every open-set tuple space,
commutes with isometries and conjugation, and respects every
defined concatenation. Expanding the subsets gives
$\mathsf{sh}_b\mathsf{sh}_{b'}=\mathsf{sh}_{b+b'}$.
Thus it is invertible with inverse $\mathsf{sh}_{-b}$.
The moment--cumulant partition formula shows that
$\sigma' =\sigma\circ\mathsf{sh}_b$ has one-point constant
$c+b$ and the same connected functions of order at least two.

\begin{proposition}[Recentering and reflection]\label{support:28-12}
The functionals $\sigma$ and $\sigma'$ have the same isometry
invariances and the same reflection positivities. They cluster
simultaneously. In the setting of \cref{support:27-22}, the map
\[
 [F]_{\sigma',u}\longmapsto
                   [\mathsf{sh}_bF]_{\sigma,u}
\]
extends to a unitary which fixes the vacuum and intertwines
the translation operators and their joint spectral measures.
\end{proposition}
\begin{proof}
The algebraic identities just established imply
\begin{equation}\label{support:28-15}
 \sigma'(\theta_u\bar F\cdot G)
 =\sigma(\theta_u\overline{\mathsf{sh}_bF}
                                  \cdot\mathsf{sh}_bG).
\end{equation}
They prove reflection positivity in both directions, using
$\mathsf{sh}_{-b}$ for the converse, and give a surjective
isometry on the quotients. The same commutation identities
intertwine translations and isometries. Uniqueness of the
spectral measure proves the last assertion. Invariance follows
also directly from the subset expansion. Each shifted clustering
difference is a finite sum of unshifted clustering differences
of subtuples, with constant integral factors; the converse uses
the inverse shift.
\end{proof}

\subsubsection{Positivity at a diagonal plane}

The following measure estimate is the analytic step in comparing
the Hilbert spaces of two reflections.

\begin{lemma}\label{support:28-17}
Let $\mu$ be a finite positive measure on
$[0,\infty)\times\mathbb R$, and set $w=(E-iq)/\sqrt2$.
Suppose that, on a Hilbert space $\mathcal K$, there are
vectors $a,b$ and self-adjoint operators $H_+\ge0,Q$ with
\begin{equation}\label{support:28-18}
 \int e^{-sw-s'\bar w}\,d\mu(E,q)
       =\langle a,e^{-sH_+}e^{-is'Q}b\rangle
       \qquad(s,s'\ge0).
\end{equation}
Then $\mu(\{|q|>E\})=0$. No commutation of $H_+$ and $Q$
is required.
\end{lemma}
\begin{proof}
For $c\in L^1(0,\infty)$ put
$p(z)=\int_0^\infty c(s)e^{-sz}\,ds$ on $\operatorname{Re}z\ge0$.
Multiply \eqref{support:28-18} by $c(s)\overline{c(s')}$ and
integrate. Fubini is justified by the bound
$|e^{-sw-s'\bar w}|\le1$, the finiteness of $\mu$, and the
contraction and unitary bounds on the right. Spectral calculus
in the indicated order gives
\begin{equation}\label{support:28-19}
 \int|p(w)|^2\,d\mu
 =\langle a,p(H_+)p^\sharp(Q)b\rangle
 \le\|a\|\|b\|\sup_{x\ge0}|p(x)|
                         \sup_{y\in\mathbb R}|p(iy)|,
 \qquad p^\sharp(y)=\overline{p(-iy)}.
\end{equation}
Take $z_0=\varepsilon+i\kappa$, where $\varepsilon>0$ and
$\kappa\in\mathbb R$, and
$c(s)=\Gamma(z_0)^{-1}s^{z_0-1}e^{-\eta s}$ with $\eta>0$.
The Gamma integral gives $p(z)=(z+\eta)^{-z_0}$, with the
principal logarithm. Consequently
\[
 \int |w+\eta|^{-2\varepsilon}
         e^{2\kappa\arg(w+\eta)}\,d\mu
 \le\|a\|\|b\|\eta^{-2\varepsilon}e^{|\kappa|\pi/2}.
\]
Fatou's lemma first as $\varepsilon\downarrow0$ and then
as $\eta\downarrow0$ yields
\begin{equation}\label{support:28-20}
 \int_{w\ne0}e^{2\kappa\arg w}\,d\mu
                      \le\|a\|\|b\|e^{|\kappa|\pi/2}.
\end{equation}
For every $\delta>0$, the measure of
$\{\arg w\ge\pi/4+\delta\}$ is at most
$\|a\|\|b\|e^{-2\kappa\delta}$ for all $\kappa>0$,
and hence is zero. Taking $\kappa<0$ proves the corresponding
assertion for $\arg w\le-\pi/4-\delta$. Their countable
union is exactly $\{|q|>E\}$.
\end{proof}

\begin{proposition}[The spectral bound from two reflections]
\label{support:28-21}
Let $\sigma$ satisfy the bounds and translation invariance of
\cref{support:27-22}. Fix $k\in\{1,2\}$ and $\epsilon\in\{1,-1\}$.
If $\sigma$ is reflection positive for the normals $e_3$ and
$e_+=(e_3+\epsilon e_k)/\sqrt2$, then its spectral measure
in $\mathcal H(\sigma,e_3)$ satisfies
\[
 \mathcal P\{(E,q):E<|q_k|\}=0.
\]
We denote this spectral-support assertion by $H\ge|P_k|$.
\end{proposition}
\begin{proof}
Put $e_-=(e_3-\epsilon e_k)/\sqrt2$ and
$W=\{x:x_3>|x_k|\}$. For $F\in\mathbb C\mathfrak S(W)$,
both $F$ and $\theta_{e_+}\theta_{e_3}F$ lie in the positive
half-space for $e_+$. Let $\mu$ be the image of the positive
spectral measure $\langle[F],\mathcal P(\cdot)[F]\rangle$
under $(E,q)\mapsto(E,\epsilon q_k)$. The same correlation
has two representations:
\begin{equation}\label{support:28-23}
 \begin{split}
 \int e^{-sw-s'\bar w}\,d\mu
 &=\sigma\bigl(\theta_{e_3}\bar F\cdot
                         \tau_{se_++s'e_-}F\bigr)\\
 &=\bigl\langle[\theta_{e_+}\theta_{e_3}F]_+,
                  T_s^+U^+(s'e_-)[F]_+\bigr\rangle_+.
 \end{split}
\end{equation}
Indeed $se_++s'e_-=(s+s')e_3/\sqrt2+
\epsilon(s-s')e_k/\sqrt2$, which gives the first equality
by \eqref{support:27-13}; the second is the defining reflection
form in $\mathcal H(\sigma,e_+)$. There
$T_s^+=e^{-sH_+}$ and $U^+(s'e_-)=e^{-is'Q}$ for the
negative of the momentum generator along $e_-$. Hence
\cref{support:28-17} shows that the spectral projection of
$\{E<|q_k|\}$ annihilates $[F]$.

For arbitrary $F\in\mathcal D_{e_3}$, a sufficiently large
positive time translation puts all its finitely many supports
in $W$. The same projection therefore annihilates $T_r[F]$.
It commutes with $T_r$, whose kernel is zero by
\cref{support:27-22}, and so annihilates $[F]$. Density proves
that the projection is zero.
\end{proof}

\begin{proposition}[Rotations fixing the time axis]\label{support:28-30}
If a functional as in \Cref{support:27-22} is invariant under the isometries $\widehat R$ acting by $R\in\operatorname{O}(2)$ on $e_3^\perp$ and fixing $e_3$, then $\mathcal U(R)[F]=[\widehat RF]$ defines a strongly continuous unitary representation of $\operatorname{O}(2)$. It fixes $\Omega$, commutes with time translations, and satisfies $\mathcal U(R)U(v)\mathcal U(R)^{-1}=U(Rv)$.
\end{proposition}
\begin{proof}
The reflection commutes with $\widehat R$, so invariance of the functional preserves the reflection form and the null space. The inverse action is that of $R^{-1}$. The commutation and covariance identities hold on tuple classes by the same identities for their supports and tests. For $R_j\to R$, the moved tests converge in the $C^1$ norm and for large $j$ have supports in fixed separated compact sets on the positive side. The reflection form therefore gives norm convergence on tuple classes, hence on their span and, by unitarity, the full Hilbert space.
\end{proof}

\begin{proposition}[The circular spectral cone]\label{support:28-32}
Suppose the hypotheses of \Cref{support:28-30} hold. If $H\ge|P_1|$, then the joint spectral measure satisfies
\begin{equation}\label{support:28-33}
\mathcal P\bigl(\{(E,q):E<|q|\}\bigr)=0.
\end{equation}
For $z\in\mathbb C^3$ with $\operatorname{Re}z_3\ge|\operatorname{Im}(z_1,z_2)|$, the integral
\[
\mathcal T(z)=\int e^{-z_3 E+i(z_1q_1+z_2q_2)}\,d\mathcal P(E,q)
\]
defines a contraction. In addition, (c) a Euclidean-invariant functional that is reflection positive at $x_3=0$ is reflection positive at every mirror.
\end{proposition}
\begin{proof}
Conjugating the spectral support inequality by $\mathcal U(R)$ gives $E\ge|q\cdot Re_1|$ for each $R\in\operatorname{O}(2)$. Use a countable dense set of $Re_1$ and take the supremum over it to obtain \eqref{support:28-33}. The norm bound follows from $-E\operatorname{Re}z_3-q\cdot\operatorname{Im}(z_1,z_2)\le0$ on this cone. Clause (c) is \Cref{support:9-31}.
\end{proof}
\begin{corollary}[Euclidean invariance and the circular cone]\label{support:28-34}
A Euclidean-invariant functional as in \Cref{support:27-22} which is reflection positive at $x_3=0$ is reflection positive at every mirror and has the circular spectral cone, the complex translation contractions, and the unitary $\operatorname{O}(2)$ action just described.
\end{corollary}
\begin{proof}
Transport reflection positivity to the diagonal plane with normal $(e_1+e_3)/\sqrt2$. \Cref{support:28-21} gives $H\ge|P_1|$, so \Cref{support:28-32} applies.
\end{proof}
\subsection{Slab operators and analytic translations}
\label{sec:slab-operators}

We now isolate the hypotheses used to recover rotations from
the cubic symmetries. The insertion estimate constructs bounded
operators across a slab. Together with the spectral bound, these
operators will give holomorphic correlation kernels. A separate
two-point growth estimate will exclude the nonconstant angular
harmonics.

Throughout this section, complex test functions are expanded
multilinearly into real and imaginary parts. On them we use
$\|h\|_\sharp=\max\{\|h\|_\infty,40M\operatorname{Lip}h\}$;
each of the real and imaginary parts has norm at most
$\|h\|_\sharp$. Such multilinear identities hold also for
Hilbert-space classes: pairing their difference with every
$[G]$ reduces the assertion to the multilinearity of $\sigma$,
and density makes the difference zero.

\begin{definition}[Plate geometry]\label{support:32-19}
Write $y=(x_3,x_1)$, leaving $x_2$ as the third coordinate.
Let $R_t$ rotate this plane by
$(y^0,y^1)\mapsto(y^0\cos t-y^1\sin t,y^0\sin t+y^1\cos t)$.
Put $n_\upsilon=(\cos\upsilon,\sin\upsilon)$, so
$n_0=(1,0)$ corresponds to $e_3$.
A plate is $f(x)=\alpha(y)\varphi(x_2)$, where
$\varphi\in C_c^1(\mathbb R)$ is real and
$\alpha\in C_c^1(\mathbb R^2;\mathbb C)$. Its label is
$\varphi$; denote the length of the smallest interval containing
$\operatorname{supp}\varphi$ by $\lambda_\varphi$, with
$\lambda_0=0$. Set
\begin{equation}\label{support:32-20}
 \operatorname{Prof}_t=
 \{\alpha\in C_c^1(\mathbb R^2):
   \operatorname{supp}\alpha\subset[t,2t]\times[-t,t],\
   |\alpha|\le3,\ \operatorname{Lip}\alpha\le12/t\}.
\end{equation}

\end{definition}

\begin{definition}[Hypotheses for the rotation argument]
\label{support:32-25}
For a real normalized functional $\sigma$ on separated tuples,
the hypotheses are as follows.
\begin{enumerate}[label=(R\arabic*),start=0]
\item $\sigma(f)=0$ for every $f\in C_c^1(\mathbb R^3)$.
\item $\sigma$ is symmetric and invariant under translations and
the cubic group $W_3$ of signed coordinate permutations.
\item For each $n\ge1$ and $d,R>0$ there is a finite
$C_n^\sigma(d;R)$, non-increasing in $d$, such that
\[
 |\sigma(f_1,\ldots,f_n)|\le C_n^\sigma(d;R)
                                      \prod_i\|f_i\|_\sharp
\]
for every tuple with separation at least $d$ and each support
in some ball of radius $R$; the centers are arbitrary.
\item $\sigma$ is reflection positive at $x_3=0$.
\item In the corresponding Hilbert space, $H\ge|P_1|$ in the
sense of spectral support.
\item There are $0<c_{\rm I}\le1/16$, $C_{\rm I}>0$ and
$p_{\rm I}\ge0$ such that, for $0<\eta\le1$, every real $h$
supported in $B(z,c_{\rm I}\eta)$ with $|z_3|\le\eta/8$, and
$F,G\in\mathbb C\mathfrak S(\{x_3>\eta/4\})$,
\begin{equation}\label{support:32-27}
 |\sigma(\theta\bar F\cdot(h)\cdot G)|
 \le C_{\rm I}\eta^{-p_{\rm I}}\|h\|_\sharp
       \|[\tau_{-(\eta/4)e_3}F]\|
       \|[\tau_{-(\eta/4)e_3}G]\|.
\end{equation}
\item For the diagonal reflection with normal
$(e_1+e_3)/\sqrt2$, the two-point form is positive:
\[
 \sum_{a,b}\bar z_a z_b\sigma(\theta_\Sigma f_a,f_b)\ge0
 \quad\bigl(f_a\in C_c^1(\{x_1+x_3>0\})\bigr).
\]
\item There is an exponent $a_\Gamma<4$ such that for every
label $\varphi$ there is $C_\varphi<\infty$ with
\begin{equation}\label{support:32-29}
 \sigma(\theta f,f)\le C_\varphi t^{-a_\Gamma}
 \quad(0<t\le1,\ f=\alpha\otimes\varphi,
                         \ \alpha\in\operatorname{Prof}_t).
\end{equation}
\end{enumerate}
\end{definition}

The Hilbert construction of \cref{support:27-22} applies under
(R1)--(R3). Its proof uses only the uniform moment bound in
(R2); alternatively the partition formula bounds each connected
function by a finite sum of products of those bounds.

\begin{lemma}[Coordinate reflection pictures]\label{support:32-16}
Under (R1)--(R3), (a) the three coordinate reflection Hilbert spaces are unitarily identified by the signed coordinate permutations that send one positive normal to another. (b) Their time translations satisfy $0\le\mathcal T_t\le1$ for $t\ge0$.
\end{lemma}
\begin{proof}
The substitution $F\mapsto wF$ in the reflection form for $w\in W_3$ preserves the form by (R1), since $w\theta_uw^{-1}=\theta_{wu}$. It also intertwines translations. The contraction and positivity are the conclusions of \Cref{support:27-22}, whose bounds follow from (R2) by the partition formula for cumulants.
\end{proof}
The coordinate identifications transport (R4) to every ordered pair
$\mu\ne\nu$:
\begin{equation}\label{support:32-33}
 H^{(\mu)}\ge|P_\nu^{(\mu)}|.
\end{equation}

\begin{remark}[Conventions in the rotation hypotheses]\label{support:32-30}
In (R5) the ball for the insertion test is open. In (R7) the supports of the plate and its reflection are separated by at least $2t$ in the $x_3$ coordinate when $\alpha\in\operatorname{Prof}_t$.
\end{remark}
\begin{remark}[Moment bounds and the energy--momentum inequality]\label{support:32-37}
(1) If (R1) and (R2) hold, translation invariance and the one-point bound give $\sigma((f))=c\int f$. In fact a continuous translation-invariant linear functional on compactly supported tests is a multiple of Lebesgue measure: it vanishes on each derivative of a smooth compactly supported test by differentiating translation invariance, and then agrees with $c\int f$ by convolution with a fixed unit-integral test. The cumulant partition formula converts the moment table to a connected-function table, taking the sum of products of the finitely many bounds. Thus $\sigma\in\mathsf{Fct}(\underline C,|c|)$ for some table $\underline C$.

(2) The spectral measure of the $e_3$ picture is exactly the joint measure of the reflected time and spatial translations of \Cref{support:27-22}. In particular, the spectral assertion of \Cref{support:28-21} for $k=1$ is (R4). Under Euclidean invariance its square-cone integral agrees with $\mathcal T(z)$ on the smaller tube where both are defined, by \eqref{support:28-33}.
\end{remark}
\subsubsection{The holomorphic translation semigroup}

For $z\in\mathbb C^3$ and $\xi\in\mathbb C^2$ put
\begin{equation}\label{support:32-24}
 d_3(z)=\operatorname{Re}z_3-|\operatorname{Im}z_1|
                              -|\operatorname{Im}z_2|,
 \qquad \mathsf{Tb}_3=\{d_3>0\},
 \qquad z^*=(-\bar z_1,-\bar z_2,\bar z_3),
\end{equation}
and
\begin{equation}\label{support:32-22}
 d_{\mathsf{Tb}}(\xi)=\operatorname{Re}\xi^0
                                -|\operatorname{Im}\xi^1|,
 \quad\mathsf{Tb}=\{d_{\mathsf{Tb}}>0\},
 \quad\xi^*=(\bar\xi^0,-\bar\xi^1),
 \quad\xi^\pm=\xi^0\pm i\xi^1.
\end{equation}
Complex rotations are defined by
$(R_\zeta\xi)^\pm=e^{\pm i\zeta}\xi^\pm$.
Both tubes are convex cones, stable under their respective
involutions. The map $\xi\mapsto(\xi^1,0,\xi^0)$ embeds
$\mathsf{Tb}$ in $\mathsf{Tb}_3$.

\begin{lemma}\label{support:32-34}
Under (R1)--(R4),
\begin{equation}\label{support:32-35}
 \mathsf T_3(z)=\int e^{-z_3E+i(z_1q_1+z_2q_2)}\,d\mathcal P(E,q)
 \quad(z\in\overline{\mathsf{Tb}_3})
\end{equation}
is a contraction semigroup, strongly continuous on the closed
tube and holomorphic in operator norm on its interior.
It satisfies $\mathsf T_3(z)^*=\mathsf T_3(z^*)$,
$\mathsf T_3(z)\Omega=\Omega$, and
$\mathsf T_3(x)[F]=[\tau_xF]$ for real $x$ with $x_3\ge0$.
Its restriction is denoted
$\mathsf T(\xi)=\mathsf T_3(\xi^1,0,\xi^0)$.
\end{lemma}
\begin{proof}
On the support of $\mathcal P$, \eqref{support:32-33} gives
$|q_1|,|q_2|\le E$, so the integrand has modulus at most
$e^{-d_3(z)E}\le1$. Spectral calculus gives the semigroup,
adjoint, vacuum and real-translation identities. Dominated
convergence in each scalar spectral measure proves strong
continuity. On a compact subset of the open tube there is
$\delta>0$ with $d_3\ge\delta$. Every derivative of degree
$m$ in the integrand is bounded by $E^m e^{-\delta E}$.
Taylor's formula on a smaller neighborhood, with this uniform
bound, permits differentiation in operator norm. This proves
holomorphy.
\end{proof}

\subsubsection{Insertion across a slab}

\begin{proposition}[Slab operators]\label{support:32-38}
Under (R1)--(R3) and (R5), for every $l>0$ and
$h\in C_c^1(\{0<x_3<l\};\mathbb C)$ there is a unique bounded
operator $B_l(h)$ such that
\begin{equation}\label{support:32-39}
 B_l(h)[F]=[(h)\cdot\tau_{le_3}F].
\end{equation}
It is complex linear in $h$. For every compact subset $K$ of
the slab there is $C(K,l)<\infty$ with
$\|B_l(h)\|\le C(K,l)\|h\|_\sharp$ when
$\operatorname{supp}h\subset K$. It has the identities
\begin{align}
 B_l(h)&=T_cB_{c'-c}(\tau_{-ce_3}h)T_{l-c'},
 &&\operatorname{supp}h\subset\{c<x_3<c'\},\label{support:32-41}\\
 B_l(h)^*&=B_l(\overline{\theta_{l/2}h}),
 & U(v)B_l(h)U(-v)&=B_l(\tau_vh),\label{support:32-42}
\end{align}
where $0\le c<c'\le l$ and $v\perp e_3$.
For real $h$ supported in $B(z,c_{\rm I}\eta)$ with
$\eta/8\le z_3\le3\eta/8$, $0<\eta\le1$,
\begin{equation}\label{support:32-44}
 \|B_{\eta/2}(h)\|\le C_{\rm I}\eta^{-p_{\rm I}}
                                             \|h\|_\sharp.
\end{equation}
If $0=c_0<\cdots<c_n$ and
$\operatorname{supp}h_j\subset\{c_{j-1}<x_3<c_j\}$, then
\begin{equation}\label{support:32-45}
 [(h_1,\ldots,h_n)]
 =\prod_{j=1}^n B_{c_j-c_{j-1}}
                   (\tau_{-c_{j-1}e_3}h_j)\Omega,
\end{equation}
with increasing $j$ from left to right.
\end{proposition}
\begin{proof}
For a small piece as in \eqref{support:32-44}, put
$h_0=\tau_{-(\eta/4)e_3}h$. Apply (R5) to $h_0$ and
$\tau_{(\eta/4)e_3}F,\tau_{(\eta/4)e_3}G$. Translation
invariance changes the resulting form into
\[
 \sigma(\theta\bar F\cdot(h)\cdot\tau_{(\eta/2)e_3}G),
\]
whose modulus is at most
$C_{\rm I}\eta^{-p_{\rm I}}\|h\|_\sharp\|[F]\|\|[G]\|$.
It descends to the quotient and extends continuously to the
Hilbert space. The bounded-form representation theorem gives
an operator with this bound. Pairing with every class identifies
its value on $[G]$ with $[(h)\cdot\tau_{(\eta/2)e_3}G]$.

For general compact $K$ in the slab, set
$\delta=\min_{x\in K}\min(x_3,l-x_3)>0$,
$\eta=\min(1,4\delta)$ and $b=c_{\rm I}\eta/2$.
Use the $C^1$ partition
\[
 \chi_{b,k}(x)=\prod_{\mu=1}^3\chi(x_\mu/b-k_\mu),
 \quad
 \chi(s)=\begin{cases}\cos^2(\pi s/2),&|s|\le1,\\0,&|s|>1.
 \end{cases}
\]
The identity $\chi(s)+\chi(s-1)=1$ for $0\le s\le1$
gives $\sum_k\chi_{b,k}=1$. Each support is a cube of
radius $b$, and $\operatorname{Lip}\chi_{b,k}\le\sqrt3\pi/(2b)$.
Only finitely many of these cubes meet $K$. Write
$h_k=\chi_{b,k}h$ and let $c_k$ be the projection of
$bk_3-\eta/4$ onto $[0,l-\eta/2]$. Because
$b\le\eta/32$ and $\eta/4\le\delta$, the center of
$\tau_{-c_ke_3}h_k$ has third coordinate between
$7\eta/32$ and $9\eta/32$. Its support lies in a ball
of radius $\sqrt3 b<c_{\rm I}\eta$. The small-piece
construction therefore applies, and
\[
 B_l(h)=\sum_k T_{c_k}B_{\eta/2}(\tau_{-c_ke_3}h_k)
                                      T_{l-c_k-\eta/2}
\]
has the required value on classes. The product estimate
\[
 \|h_k\|_\sharp\le
 \left(1+\frac{40\sqrt3\pi M}{c_{\rm I}\eta}\right)
                                               \|h\|_\sharp
\]
gives the compact-support bound, with an extra factor two for
the real and imaginary parts. The construction is independent
of the partition because its values on the dense set of classes
are prescribed.

Translation of the tuple in \eqref{support:32-39} proves
\eqref{support:32-41} and tangential covariance. Reflection in
$x_3=l/2$ interchanges its left and right entries; symmetry
and conjugation give the adjoint formula. Finally, iterating
\eqref{support:32-39} through the ordered slabs gives
\eqref{support:32-45}.
\end{proof}

\subsubsection{Removing coincident planar projections}

\begin{proposition}\label{support:32-48}
Assume (R1)--(R4). Fix $n\ge2$ and a real angle $t$.
If $\sigma(R_tX)=\sigma(X)$ for every $n$-tuple whose
supports have pairwise disjoint projections onto the
$(x_3,x_1)$ plane, the identity holds for all separated
$n$-tuples.
\end{proposition}
\begin{proof}
First suppose that two groups of entries lie respectively
in $x_2<c-\delta$ and $x_2>c+\delta$, where $\delta>0$.
In the Hilbert space with normal $e_2$, their correlation
after a relative translation $su$, $u\perp e_2$, $|u|=1$,
is a matrix element of $T_{2\delta}^{(2)}U^{(2)}(su)$.
For its finite complex spectral measure $m$ this is
\[
 \Phi(s)=\int e^{-2\delta E+is(u_1q_1+u_3q_3)}\,dm(E,q).
\]
Since $|q_1|,|q_3|\le E$ and $|u_1|+|u_3|\le\sqrt2$,
this is holomorphic for $|\operatorname{Im}s|<\sqrt2\delta$.
Thus the relative planar translation is analytic whenever
there is such an $x_2$ gap.

Let $A_j$ be the planar projections of the supports and $I_j$
their $x_2$ interval hulls. Suppose first that every pair of
intersecting $A_j$ has a level separating its two intervals
and lying outside all the $I_j$. Induct on the number of
intersecting pairs. Choose a separating level for one pair,
and translate all entries above it by $se_1$. For all sufficiently
large $s$, the planar projections on opposite sides no longer
meet, and the number of intersecting pairs strictly decreases.
The interval hulls and their separating levels are unchanged.
The induction hypothesis gives rotation invariance for these
large $s$. Both correlations are holomorphic on a common strip:
after rotation, the relative translation is $sR_te_1$, still
perpendicular to $e_2$. The identity theorem gives equality
also at $s=0$.

For a general tuple of separation $d>0$, choose
$0<\varepsilon<d/n$ and partition each entry in the $x_2$
coordinate with $\chi(2x_2/\varepsilon-k)$.
Multilinearity gives a finite sum of tuples whose interval
hulls have length at most $\varepsilon$. If two planar
projections of a nonzero summand meet, points over their common
projection are at $x_2$ distance at least $d$. The gap between
their two interval hulls has length at least
$d-2\varepsilon>(n-2)\varepsilon$. The other $n-2$ interval
hulls cannot cover this gap, so it contains a level outside all
hulls. The preceding induction applies to each summand.
The partition commutes with $R_t$, and summing proves the
assertion. Tuples with a zero entry vanish by multilinearity.
\end{proof}

\subsection{The two-point rotation argument}
\label{sec:two-point-rotation}

For a plate label $\varphi$, reflection positivity gives a
positive spectral measure on the energy--momentum half-plane.
Its transform represents the two-point function. The plate
growth bound leaves only the angular frequencies $0,4,-4$.
Positivity at the coordinate and diagonal planes gives opposite
signs for the coefficient at frequency four.

\subsubsection{The spectral measure of a label}

Put $Y=[0,\infty)\times\mathbb R$, $V=\{(E,q):E\ge|q|\}$,
and, for a profile supported in $y^0>0$, set
\begin{equation}\label{support:33-9}
 \widetilde\alpha(E,q)=\int\alpha(y)e^{-y^0E+iy^1q}\,dy.
\end{equation}
This transform is continuous and bounded by
$\|\alpha\|_1e^{-\varepsilon E}$ if the support has
$y^0\ge\varepsilon>0$. Convolution of profiles multiplies
their transforms.

We first work with any symmetric bilinear two-point form $s$
on separated supports which is translation invariant, satisfies
the order-two bound in (R2), and is reflection positive at
$y^0=0$. The form $s(\theta\bar f,h)$ on
$C_c^1(\{x_3>0\};\mathbb C)$ defines a one-field Hilbert
space. The proof of \cref{support:27-22}, with single test functions
instead of tuples and with no distinguished vacuum, gives
positive time translations, tangential unitaries and a joint
spectral measure $\mathcal P^s$ on $Y$.

\begin{proposition}\label{support:33-21}
For every real label $\varphi$ there is a unique positive
Borel measure $\rho[s,\varphi]$ on $Y$, finite on compact
sets, such that all profile transforms are in $L^2(\rho)$ and
\begin{equation}\label{support:33-22}
 s\bigl(\theta(\bar\alpha\otimes\varphi),
                       \alpha'\otimes\varphi\bigr)
       =\int\overline{\widetilde\alpha}\,
                         \widetilde\alpha'\,d\rho.
\end{equation}
More precisely, for every Borel set $B\subset Y$,
\begin{equation}\label{support:33-23}
 \langle[\alpha\otimes\varphi],\mathcal P^s(B)
                              [\alpha'\otimes\varphi]\rangle
 =\int_B\overline{\widetilde\alpha}\,
                         \widetilde\alpha'\,d\rho.
\end{equation}
\end{proposition}
\begin{proof}
Write $\Psi_\alpha=[\alpha\otimes\varphi]$ and
$\mu_{\alpha,\alpha'}(B)=
\langle\Psi_\alpha,\mathcal P^s(B)\Psi_{\alpha'}\rangle$.
Riemann sums for convolution converge in $C^1$, with supports
in a fixed compact subset of the positive half-plane. The
two-point bound therefore gives
\begin{equation}\label{support:33-24}
 \Psi_{\alpha'*\alpha}
       =\widetilde\alpha'(\mathcal P^s)\Psi_\alpha.
\end{equation}
On the right, the same Riemann sums converge weakly by the
spectral representation of translations and dominated
convergence; this identifies their norm limit on the left.
Commutativity of convolution implies
\begin{equation}\label{support:33-25}
 |\widetilde\alpha'|^2\mu_{\alpha,\alpha}
       =|\widetilde\alpha|^2\mu_{\alpha',\alpha'}.
\end{equation}

On $O_\alpha=\{\widetilde\alpha\ne0\}$ define
$d\rho_\alpha=|\widetilde\alpha|^{-2}
d\mu_{\alpha,\alpha}$. Equation \eqref{support:33-25} makes
these local measures agree on overlaps. The sets $O_\alpha$
cover $Y$: at $(E_0,q_0)$ take
$\alpha(y)=\chi(y^0-2)\chi(y^1)e^{-iy^1q_0}$, whose
transform there is strictly positive. Compact exhaustion of
$Y$ gives a countable subcover. Partitioning it into disjoint
Borel pieces defines a positive measure $\rho$ with the
prescribed restrictions. It is locally finite, since near a
point one transform is bounded away from zero.

Equation \eqref{support:33-25} also shows that
$\mu_{\alpha,\alpha}$ vanishes on the zeros of
$\widetilde\alpha$: intersect that zero set with each member
of the countable cover and use the nonzero transform there.
Hence $d\mu_{\alpha,\alpha}=|\widetilde\alpha|^2d\rho$
on all of $Y$. Polarization proves \eqref{support:33-23}, and
taking $B=Y$ gives \eqref{support:33-22}.

For uniqueness, a finite complex measure $m$ on $Y$ is
determined by its Laplace--Fourier transform. Indeed, vanishing
of $\int e^{-tE+ivq}dm$ at all integers $t\ge0$ and all
$v\in\mathbb R$ first gives vanishing against
$p(e^{-E})e^{ivq}$ for polynomials $p$. Uniform polynomial
approximation on $[0,1]$, followed by Fourier uniqueness in
$q$, gives vanishing against $g(e^{-E})\mathbf1_B(q)$ for
continuous $g$ and Borel $B$. Approximation of indicators of
closed subsets of $[0,\infty)$ and the monotone class theorem
then give $m=0$. Thus vanishing against every
$\widetilde\alpha$ also implies $m=0$, by testing the
continuous Laplace--Fourier transform against profiles.
If two proposed measures satisfy \eqref{support:33-22}, replace
$\alpha'$ by $\alpha''*\alpha$. The preceding uniqueness
applies to the finite measure
$|\widetilde\alpha|^2(\rho'-\rho'')$ and makes it zero.
Division on $O_\alpha$ and the countable cover give
$\rho'=\rho''$.
\end{proof}

Apply this to $s=\sigma|_{\mathfrak S_2}$ under (R1)--(R4)
and write $\rho_\varphi=\rho[s,\varphi]$. The one-field
space embeds isometrically into $\mathcal H(\sigma)$ through
$[f]\mapsto[(f)]$. It intertwines translations, so
$|\widetilde\alpha|^2\rho_\varphi$ is the image of the
full spectral measure of $[(\alpha\otimes\varphi)]$ under
$(E,q_1,q_2)\mapsto(E,q_1)$. Consequently $\rho_\varphi$
is carried by $V$. Reflection of $x_1$ and uniqueness in
\cref{support:33-21} make it invariant under $q\mapsto-q$.

\subsubsection{From plate bounds to a kernel bound}

Let $a_0(s)=1+\cos(2\pi(s-3/2))$ on $[1,2]$, zero
elsewhere, and let
$b_1=(3/2)\mathbf1_{[-1/3,1/3]}$, $b=b_1*b_1*b_1$.
For $1/2\le\lambda\le1$ define
$\alpha_{t,\lambda}(y)=a_0(y^0/t)b(y^1/(\lambda t))$.
These profiles belong to $\operatorname{Prof}_t$:
$\int a_0=\int b=1$, $0\le a_0\le2$, $|a_0'|\le2\pi$,
$0\le b\le9/8$, and $|b'|\le9/4$. The last bounds follow
by writing $b=(3/2)\int_{u-1/3}^{u+1/3}T(v)dv$ with
$T(v)=(9/4)(2/3-|v|)_+$.

There is an absolute $C_b<\infty$ such that every positive
locally finite measure $\rho$ carried by $V$ satisfies
\begin{equation}\label{support:33-40}
 \Gamma[\rho](5t):=\int e^{-5tE}\,d\rho
 \le C_b t^{-4}\sup_{1/2\le\lambda\le1}
                    \int|\widetilde\alpha_{t,\lambda}|^2d\rho.
\end{equation}
Here the inequality is valid also for infinite integrals.
To see it, the transform of $b$ is
$\widehat b(q)=(\sin(q/3)/(q/3))^3$, and
$\widehat a_0(s)\ge e^{-2s}$ for $s\ge0$. Thus
\[
 |\widetilde\alpha_{t,\lambda}(E,q)|^2
      \ge(t^4/4)e^{-4tE}\widehat b(\lambda tq)^2.
\]
Let $c_6>0$ be the infimum of the means of $\sin^6$ over
intervals of length at least $1/2$. It is positive by
periodicity, compactness for lengths between $1/2$ and $\pi$,
and counting full periods for larger lengths. Set
$c_b=\min(\sin^61,3^6c_6)$. For $|s|\le3$, concavity of
sine gives $\widehat b(\lambda s)^2\ge\sin^61$; for
$|s|\ge3$, integration over $\lambda$ gives the mean of
$\sin^6$ on an interval of length $|s|/6$. Therefore
\[
 \int_{1/2}^1 2\widehat b(\lambda s)^2d\lambda
          \ge c_b(1+|s|)^{-6}.
\]
On $V$ this is at least $c_b6^{-6}e^5e^{-tE}$ when
$s=tq$, since $(1+x)^{-6}\ge6^{-6}e^5e^{-x}$.
Tonelli proves \eqref{support:33-40} with
$C_b=4\cdot6^6e^{-5}/c_b$.

For $\rho_\varphi$, the supremum in \eqref{support:33-40} is
finite at each fixed $t$ by (R2), so $\Gamma_\varphi(t)$
is finite for all $t>0$. Under (R7), \eqref{support:33-40} gives
$\Gamma_\varphi(5t)\le C_bC_\varphi t^{-4-a_\Gamma}$ for
$0<t\le1$. If $a_\Gamma<-4$, letting $t\downarrow0$ and using
monotonicity forces $\rho_\varphi=0$. If $a_\Gamma\ge-4$,
apply this bound at $\min\{t/5,1\}$ and use monotonicity again.
In both cases,
\begin{equation}\label{support:33-54}
 \Gamma_\varphi(t)\le
 C'_\varphi\max(1,t^{-1})^{4+a_\Gamma},
 \qquad C'_\varphi=5^{4+a_\Gamma}C_bC_\varphi.
\end{equation}
Define
\begin{equation}\label{support:33-29}
 k_\varphi(\xi)=\int_V e^{-\xi^0E+i\xi^1q}\,d\rho_\varphi.
\end{equation}
It is holomorphic on $\mathsf{Tb}$ and bounded by
$\Gamma_\varphi(d_{\mathsf{Tb}}(\xi))$. The integrable
majorant follows from $|q|\le E$; on smaller tubes it also
bounds all derivatives. It is even in $\xi^1$ and obeys
$\overline{k_\varphi(\xi)}=k_\varphi(\xi^*)$.

\begin{proposition}[The two-point kernels]\label{support:33-47}
Under (R1)--(R4), $k_\varphi$ extends by quarter turns to a
real analytic function $\widetilde k_\varphi$ on
$\mathbb R^2\setminus\{0\}$. For disjointly supported profiles,
\begin{equation}\label{support:33-50}
 \sigma(\alpha\otimes\varphi,\alpha'\otimes\varphi)
   =\iint\alpha(y)\alpha'(y')
                         \widetilde k_\varphi(y'-y)\,dy\,dy'.
\end{equation}
For two labels set
$k_{\varphi,\psi}=(k_{\varphi+\psi}-k_{\varphi-\psi})/4$,
and likewise for $\widetilde k$. Then
\begin{equation}\label{support:33-51}
 \sigma(\alpha\otimes\varphi,\alpha'\otimes\psi)
 =\sigma(\alpha\otimes\psi,\alpha'\otimes\varphi)
 =\iint\alpha(y)\alpha'(y')
                    \widetilde k_{\varphi,\psi}(y'-y)\,dy\,dy'.
\qedhere
\end{equation}
\end{proposition}
\begin{proof}
If the profiles lie on opposite sides of $y^0=c$, translate
to $c=0$ and apply \eqref{support:33-22}. Fubini gives
\eqref{support:33-50} with $k_\varphi$ on the positive half-plane.
For $v^0,v^1>0$, take profiles in small neighborhoods of
$0,v$ and rotate both through $-\pi/2$. Cubic invariance
equates their correlations, and the integral representation
gives $k_\varphi(v)=k_\varphi(R_{-\pi/2}v)$ by shrinking
the neighborhoods. This equality makes the quarter-turned
definitions agree on overlaps. They cover the punctured plane
and give $\widetilde k_\varphi$. Reality follows from the
evenness in $q$.

For arbitrary disjoint profile supports of distance $d>0$,
partition each into squares of diameter less than $d/4$ using
the functions $\chi$ above. Any pair of nonzero pieces is
ordered along one of the signed coordinate axes, since one
coordinate of a vector between them has modulus at least
$d/\sqrt2$. The ordered formula applies after a quarter turn,
and its sum is \eqref{support:33-50}.

Polarization in the labels gives the sum of the two mixed
correlations in \eqref{support:33-51}. To see that they are equal,
use an even approximate identity $j_\delta$ on the plane.
Translation invariance and Riemann sums give
\[
 \sigma((\alpha*j_\delta)\otimes\varphi,
                      (\alpha'*j_\delta)\otimes\psi)
 =\iint\alpha(u)\alpha'(u')h_\delta(u'-u)\,du\,du',
\]
where $h_\delta(v)=\sigma(j_\delta\otimes\varphi,
(\tau_vj_\delta)\otimes\psi)$ and $\delta$ is smaller
than a quarter of the separation. Apply the same identity
to $\alpha'(-u),\alpha(-u)$ and change both integration
variables to their negatives. Its right side is unchanged.
Rotation by $\pi$ and permutation of the two entries turn
its left side into the expression with the labels exchanged.
As $\delta\downarrow0$, convolution converges in $C^1$ on
fixed separated neighborhoods. The order-two bound permits
the limit and proves the equality of the mixed correlations.
\end{proof}

\subsubsection{Eliminating the fourth harmonic}

\begin{proposition}\label{support:33-53}
Under (R1)--(R4) and (R7), for each $r>0$ the function
$k_{\varphi,r}(\zeta)=k_\varphi(R_\zeta(r,0))$, initially
defined for $|\operatorname{Re}\zeta|<\pi/2$, extends to
an entire function with period $\pi/2$. There are real
$c_0(r),c_1(r)$ such that
\begin{equation}\label{support:33-55}
 k_{\varphi,r}(\zeta)=c_0(r)+2c_1(r)\cos4\zeta,
 \qquad
 \widetilde k_\varphi(r\cos t,r\sin t)
                      =c_0(r)+2c_1(r)\cos4t.
\end{equation}
\end{proposition}
\begin{proof}
For $\zeta=t+is$ in the initial strip,
$d_{\mathsf{Tb}}(R_\zeta(r,0))=re^{-|s|}\cos t$.
The quarter-turn identity of \cref{support:33-47} agrees on
overlapping strips by holomorphic uniqueness, so translates
by $j\pi/2$ give one entire periodic function. Move $t$
modulo $\pi/2$ into $[-\pi/4,\pi/4]$. Then
\[
 |k_{\varphi,r}(t+is)|\le
       \Gamma_\varphi(re^{-|s|}/\sqrt2)
       \le C_r e^{A|s|},\qquad A=\max(0,4+a_\Gamma)<8.
\]
For an entire $\pi/2$-periodic function with this bound,
shifting the contour of its Fourier coefficient from the
real axis to height $s$ shows that the coefficient of
$e^{4im\zeta}$ vanishes if $4|m|>A$. Only $m=0,1,-1$
remain. Evenness in $\xi^1$ equates the last two coefficients,
and reality on real points makes both remaining coefficients
real. This proves the formulas.
\end{proof}

\begin{proposition}[Rotation invariance at two plates]\label{support:33-57}
Under (R1)--(R4), (R6) and (R7), $c_1(r)=0$ for all
labels and all $r>0$. Thus $\widetilde k_{\varphi,\psi}(v)$
depends only on $|v|$, and correlations of two plates with
disjoint planar supports are invariant under every $R_t$.
Moreover, for $v^0>0$ and real $s$,
\begin{equation}\label{support:33-59}
 k_{\varphi,\psi}(R_{is}v)=k_{\varphi,\psi}(v).
\end{equation}
\end{proposition}
\begin{proof}
For imaginary angles the integral in \eqref{support:33-29} is
\[
 k_{\varphi,r}(is)
 =\int_V e^{-r(E\cosh s+q\sinh s)}\,d\rho_\varphi\ge0.
\]
Dividing $c_0(r)+2c_1(r)\cosh4s$ by $\cosh4s$ and
letting $s\to\infty$ gives $c_1(r)\ge0$.

For the opposite sign, turn the two-point form through
$\pi/4$: $s^{\rm d}(f,h)=\sigma(R_{\pi/4}f,R_{\pi/4}h)$.
It is translation invariant. The rotation preserves the sharp and
$L^1$ norms. Since it and its inverse have sup-norm operator norm
$\sqrt2$, it sends a tuple of sup type $(d;R)$ to one of type
$(d/\sqrt2;\sqrt2 R)$. Thus its order-two moment bound has table
$C_2^\sigma(d/\sqrt2;\sqrt2 R)$, and (R6) gives its coordinate
reflection positivity. Hence
\cref{support:33-21} constructs a positive measure $\rho_{\rm d}$
on $Y$. At this point no cone support for $\rho_{\rm d}$
is assumed.

Let $b_s(E,q)=(E\cosh s+q\sinh s,q\cosh s+E\sinh s)$.
It preserves $V$, has inverse $b_{-s}$, and its first
coordinate is at least $e^{-|s|}E$ on $V$. Thus
$(b_s)_*\rho_\varphi$ has finite Laplace transforms, and
its kernel at real $v^0>0$ is $k_\varphi(R_{is}v)$.
The three harmonics in \eqref{support:33-55} give
\[
 k_\varphi(R_iv)+k_\varphi(R_{-i}v)
 = (\cosh4+1)k_\varphi(v)
       -(\cosh4-1)\widetilde k_\varphi(R_{\pi/4}v).
\]
Multiply by $\overline{\alpha(u)}\alpha'(u')$ at
$v=u'-\theta u$ and integrate. The last term represents
the turned form by \eqref{support:33-50}; the other terms are
the spectral representations for the boosted measures.
Uniqueness in \cref{support:33-21} therefore gives the identity
of positive measures
\begin{equation}\label{support:33-63}
 (\cosh4-1)\rho_{\rm d}
       +(b_1)_*\rho_\varphi+(b_{-1})_*\rho_\varphi
                   =(\cosh4+1)\rho_\varphi.
\end{equation}
In particular
$0\le\rho_{\rm d}\le(\cosh4+1)(\cosh4-1)^{-1}\rho_\varphi$.
This proves both its cone support and finiteness of its
Laplace transforms. Its kernel equals
$\widetilde k_\varphi(R_{\pi/4}v)$ on real $v^0>0$,
by the integral identity against profiles. Analytic
continuation in the angle thus gives
\[
 k[\rho_{\rm d}](R_{is}(r,0))
                   =c_0(r)-2c_1(r)\cosh4s\ge0.
\]
The same limit proves $c_1(r)\le0$. Hence it vanishes.
Polarization in $\varphi\pm\psi$ gives the radiality of
the mixed kernel. Formula \eqref{support:33-51} proves invariance
of the two-plate correlations. Holomorphic uniqueness in the
initial strip gives \eqref{support:33-59}.
\end{proof}

\subsection{Point fields and plate kernels}
\label{sec:point-kernels}

The insertion estimate first gives operators smeared inside a slab.
Holomorphy of the translation semigroup improves this to a pointwise
operator bound. The essential step is a norm Cauchy estimate for
mollified insertions; weak compactness alone would not give the
products needed below. Throughout this section we assume (R1)--(R5).

Write $B_d^{\rm c}(h)=B_{2d}(\tau_{de_3}h)$.
\begin{definition}[Point-bound constants]\label{support:34-3}
Put
\begin{equation}\label{support:34-6}
 \wp=p_{\rm I}+4,\qquad
 C_{\rm plt}=480\cdot8^3m_2c_{\rm I}^{-4}C_{\rm I}M,
 \qquad C_{\rm pt}=320\cdot8^4m_3c_{\rm I}^{-4}C_{\rm I}M.
\end{equation}
The absolute constants are obtained from the radial weight
\[
 m(s)=\begin{cases}
 c_m,&0\le s\le1/3,\\
 c_m\cos^2(3\pi(s-1/3)/4),&1/3\le s\le1,\\
 0,&s\ge1,
 \end{cases}\qquad c_m=\frac{9\pi}{5\pi^2-8},
\]
by setting
$m_0=c_m$, $m_1=3\pi c_m/4$,
$m_2=\sqrt2m_0m_1+2m_0^2$ and
$m_3=m_0^3+\sqrt3m_0^2m_1$.
The function $m$ is $C^1$, nonnegative and constant near zero;
direct integration gives $2\pi\int_0^1m(s)s\,ds=1$.
Let $j_\delta^{(N)}$, $N=2,3$, be $m(2|x|/\delta)$ divided
by its integral. These are nonnegative radial $C^1$ mollifiers
supported in $|x|\le\delta/2$. We also write $\mathsf{mv}_3=m_3$.
\end{definition}

For $r>0$, the product weight
\[
 \alpha_{r,t}^{(N)}(x)=\prod_{i=1}^N
 r^{-2}m\bigl(r^{-1}\sqrt{x_i^2+t_i^2}\bigr)
\]
has supremum at most $m_0^Nr^{-2N}$ and Lipschitz constant at
most $\sqrt N m_0^{N-1}m_1r^{-2N-1}$. Iterating the one-variable
mean-value formula in polar coordinates gives
\begin{equation}\label{support:34-16}
 u(0)=\int_{(-r,r)^N}\!dt\int\!dx\,
             \alpha_{r,t}^{(N)}(x)u(x+it)
\end{equation}
whenever $u$ is holomorphic on a neighborhood of the closed
polydisc of radius $r$. All operator integrals below converge in
norm; their integrands are continuous with compact support.

\subsubsection{Removing the smearing}

\begin{proposition}[Point plates]\label{support:34-17}
For every real label $\varphi\in C_c^1(\mathbb R)$ there is a
norm-holomorphic family $B_\varphi(\xi,\xi')$, linear in
$\varphi$, on $\mathsf{Tb}\times\mathsf{Tb}$ such that
\begin{equation}\label{support:34-18}
 \|B_\varphi(\xi,\xi')\|\le
 C_{\rm plt}(\lambda_\varphi+1)\|\varphi\|_{C^1}
 \min\{1,d_{\mathsf{Tb}}(\xi),d_{\mathsf{Tb}}(\xi')\}^{-\wp}.
\end{equation}
For $\kappa,\kappa'\in\overline{\mathsf{Tb}}$ it satisfies
\begin{align}
 \mathsf T(\kappa)B_\varphi(\xi,\xi')\mathsf T(\kappa')
 &=B_\varphi(\kappa+\xi,\xi'+\kappa'),\label{support:34-19}\\
 B_\varphi(\xi,\xi')^*&=B_\varphi(\xi'^*,\xi^*).
 \label{support:34-20}
\end{align}
If $\operatorname{supp}\alpha\subset\{|y^0|<d\}$, then
\begin{equation}\label{support:34-21}
 B_d^{\rm c}(\alpha\otimes\varphi)
 =\int\alpha(y)B_\varphi(dn_0+y,dn_0-y)\,dy.
\end{equation}
These properties determine the family already among weakly
continuous families satisfying the shift law and the integral
identity for real profiles.
\end{proposition}
\begin{proof}
Put $j_\delta=j_\delta^{(2)}$ and
\[
 A_\delta=B_\delta^{\rm c}(j_\delta\otimes\varphi),\qquad
 B^\delta(\xi,\xi')=
 \mathsf T(\xi-\delta n_0)A_\delta
 \mathsf T(\xi'-\delta n_0).
\]
This is defined when both tube distances are at least $\delta$,
and holomorphic when they are larger. Radiality of the mollifier
and the slab adjoint formula give $A_\delta^*=A_\delta$.
The shift and adjoint laws follow at once. Slab localization and
tangential covariance give
\[
 B^\delta(dn_0+y,dn_0-y)
       =B_d^{\rm c}((\tau_yj_\delta)\otimes\varphi).
\]
Consequently, for continuous $\alpha$ supported in
$|y^0|<d-\delta$,
\begin{equation}\label{support:34-25}
 \int\alpha(y)B^\delta(dn_0+y,dn_0-y)\,dy
        =B_d^{\rm c}((\alpha*j_\delta)\otimes\varphi).
\end{equation}
Indeed, convolution Riemann sums converge in $C^1$ with one
compact support, and the slab bound permits passage to the limit.

We first bound $B^\delta$ independently of $\delta$. Set
$\eta=\min(1,d_{\mathsf{Tb}}(\xi),d_{\mathsf{Tb}}(\xi'))$,
$r=c_{\rm I}\eta/8$ and $0<\delta\le r$. For unit vectors
$v,v'$, apply \eqref{support:34-16} to
$u(z)=\langle v,B^\delta(\xi+z,\xi'-z)v'\rangle$.
The polydisc of radius $2r$ lies in its domain because
$\eta\ge128r$. Move the imaginary shift into the vectors by
the two contractions with arguments
$\xi-\eta n_0/4+it$ and $\xi'-\eta n_0/4-it$.
The remaining inner integral is a matrix element of
\[
 B_{\eta/4}^{\rm c}(\gamma_t\otimes\varphi),
 \qquad \gamma_t=\alpha_{r,t}^{(2)}*j_\delta.
\]
Its profile is supported in $[-2r,2r]^2$, with supremum
$m_0^2r^{-4}$ and Lipschitz constant
$\sqrt2m_0m_1r^{-5}$.

Partition the label as
$\varphi_k(s)=\chi(s/(4r)-k)\varphi(s)$, using the partition
in the proof of \cref{support:32-38}. At most
$3(\lambda_\varphi+1)/(c_{\rm I}\eta)$ terms are nonzero.
Each $\gamma_t\otimes\varphi_k$ is supported in
$B(4rke_2,c_{\rm I}\eta)$ and has sharp norm at most
$40Mm_2r^{-5}\|\varphi\|_{C^1}$.
The centered version of \eqref{support:32-44}, followed by integration
over $(-r,r)^2$, therefore gives
\[
 \|B^\delta(\xi,\xi')\|\le
 (2r)^2\frac{3(\lambda_\varphi+1)}{c_{\rm I}\eta}
 C_{\rm I}\eta^{-p_{\rm I}}40Mm_2r^{-5}\|\varphi\|_{C^1}.
\]
This is \eqref{support:34-18} with $B^\delta$ in place of
$B_\varphi$.

To obtain norm convergence, fix $\epsilon>0$ and set
\[
 \delta_\epsilon=\frac{c_{\rm I}}{16}\min(1,\epsilon),\qquad
 b_\epsilon=C_{\rm plt}(\lambda_\varphi+1)\|\varphi\|_{C^1}
                  (\tfrac12\min(1,\epsilon))^{-\wp}.
\]
For $\delta\le\delta_\epsilon$, the holomorphic map
$z\mapsto B^\delta(\epsilon n_0+z,\epsilon n_0-z)$ is
bounded by $b_\epsilon$ on the polydisc of radius
$\epsilon/4$. Cauchy's derivative estimate gives the bound
$8b_\epsilon|y|/\epsilon$ for its change from $z=0$ to a
real $y$ with $|y|\le\epsilon/8$. Integrate against
$j_{\delta'}$ and use \eqref{support:34-25}; interchange
$\delta,\delta'$ and use commutativity of convolution. The result is
\[
 \|B^\delta(\epsilon n_0,\epsilon n_0)
       -B^{\delta'}(\epsilon n_0,\epsilon n_0)\|
 \le\frac{4b_\epsilon}{\epsilon}(\delta+\delta').
\]
Thus $B_{[\epsilon]}=\lim_{\delta\downarrow0}
B^\delta(\epsilon n_0,\epsilon n_0)$ exists in norm.
For $0<\epsilon\le\min(d_{\mathsf{Tb}}(\xi),d_{\mathsf{Tb}}(\xi'))$,
define
\begin{equation}\label{support:34-27}
 B_\varphi(\xi,\xi')=
 \mathsf T(\xi-\epsilon n_0)B_{[\epsilon]}
 \mathsf T(\xi'-\epsilon n_0).
\end{equation}
This is the limit of $B^\delta(\xi,\xi')$, uniformly when
both tube distances exceed $\epsilon$, so it does not depend
on the chosen $\epsilon$. Choosing a smaller $\epsilon$ locally
proves norm holomorphy directly from this formula.
The bound, shifts, adjoints and linearity pass to the limit.
In \eqref{support:34-25}, uniform norm convergence on the support
and $\alpha*j_\delta\to\alpha$ in $C^1$ prove
\eqref{support:34-21}.

For uniqueness, the integral identity against all real profiles
and weak continuity identify the matrix elements at
$(dn_0+y,dn_0-y)$, in particular at $(dn_0,dn_0)$.
The shift law then identifies every pair of tube points.
\end{proof}

\begin{proposition}[Point fields]\label{support:34-30}
There is a unique weakly continuous family
$B(z,z')$, $(z,z')\in\mathsf{Tb}_3^2$, satisfying the
three-dimensional shift law and the following slab integral
formula. It is norm holomorphic and obeys
\begin{align}
 \|B(z,z')\|&\le C_{\rm pt}
 \min\{1,d_3(z),d_3(z')\}^{-\wp},
 \label{support:34-31}\\
 \mathsf T_3(\kappa)B(z,z')\mathsf T_3(\kappa')
 &=B(\kappa+z,z'+\kappa'),\label{support:34-32}\\
 B(z,z')^*&=B(z'^*,z^*),\label{support:34-33}\\
 B_d^{\rm c}(h)&=\int h(x)B(de_3+x,de_3-x)\,dx.
 \label{support:34-34}
\end{align}
Here $\kappa,\kappa'\in\overline{\mathsf{Tb}_3}$, and in
the last identity $h\in C_c^1(\{|x_3|<d\};\mathbb C)$.
\end{proposition}
\begin{proof}
Use $j_\delta^{(3)}$ in the preceding construction, with
$A_\delta=B_\delta^{\rm c}(j_\delta^{(3)})$ and
$\mathsf T_3$ in place of $\mathsf T$. The convolution
identity and the shift and adjoint laws have the same proofs.
For completeness, the change in the uniform estimate is as follows.
Put $\eta=\min(1,d_3(z),d_3(z'))$
and $r=c_{\rm I}\eta/8$. The polydisc of radius $2r$
loses at most $6r$ in tube distance. The three-variable
mean-value formula, after moving the imaginary shifts into
contractions, gives the centered insertion of
$\gamma_t=\alpha_{r,t}^{(3)}*j_\delta^{(3)}$.
Its support lies in $[-2r,2r]^3\subset B(0,c_{\rm I}\eta)$,
and $\|\gamma_t\|_\sharp\le40Mm_3r^{-7}$.
No label partition is needed. Hence the bound is
\[
 (2r)^3C_{\rm I}\eta^{-p_{\rm I}}40Mm_3r^{-7}
       =C_{\rm pt}\eta^{-\wp}.
\]
At $(\epsilon e_3,\epsilon e_3)$ use a polydisc of radius
$\epsilon/6$. With
$b_\epsilon=C_{\rm pt}(\tfrac12\min(1,\epsilon))^{-\wp}$,
Cauchy's estimate and the symmetric convolution argument give
\[
 \|B^\delta(\epsilon e_3,\epsilon e_3)
        -B^{\delta'}(\epsilon e_3,\epsilon e_3)\|
 \le6b_\epsilon(\delta+\delta')/\epsilon.
\]
Norm convergence, extension to arbitrary tube pairs, holomorphy,
the integral formula and uniqueness now follow exactly by the
limit and shift arguments above.
\end{proof}

\begin{corollary}[Slab operators in $L^1$]\label{support:34-42}
Let $\sigma$ satisfy (R1)--(R5), with $\wp$ and $C_{\rm pt}$ as in \eqref{support:34-6}. In (a), for $0\le d'<d$ and $\operatorname{supp}h\subset\{|x_3|\le d'\}$,
\begin{equation}\label{support:34-43}
 \|B_d^{\rm c}(h)\|\le
 C_{\rm pt}\min(1,d-d')^{-\wp}\|h\|_{L^1}.
\end{equation}
In (b), if $F,F'\in\mathbb C\mathfrak S(\{x_3>d\})$, the function
\begin{equation}\label{support:34-44}
 v_d(x)=\langle[\tau_{-de_3}F],
 B(de_3+x,de_3-x)[\tau_{-de_3}F']\rangle
\end{equation}
is real analytic for $|x_3|<d$ and satisfies
\begin{align}
 \sigma(\theta\bar F\cdot(h)\cdot F')
       &=\int h(x)v_d(x)\,dx,\label{support:34-45}\\
 |v_d(x)|&\le C_{\rm pt}\min(1,d-|x_3|)^{-\wp}
       \|[\tau_{-de_3}F]\|\|[\tau_{-de_3}F']\|.
       \label{support:34-46}
\end{align}
In (c), identifying $\xi=(\xi^0,\xi^1)$ with
$(\xi^1,0,\xi^0)$,
\begin{equation}\label{support:34-47}
 B_\varphi(\xi,\xi')
       =\int\varphi(s)B(\xi+se_2,\xi'-se_2)\,ds.
\end{equation}
\end{corollary}
\begin{proof}
Integrate \eqref{support:34-31} to obtain \eqref{support:34-43}.
For every $h\in C_c^1(\{|x_3|<d\})$, the slab form gives \eqref{support:34-45}; its density is the
restriction of a holomorphic matrix element on
$|\operatorname{Re}z_3|+|\operatorname{Im}z_1|
+|\operatorname{Im}z_2|<d$, and the same point bound gives
\eqref{support:34-46}. The right side of \eqref{support:34-47} is
weakly continuous, satisfies the plate shift law, and, by
Fubini and \eqref{support:34-34}, represents the slab operator of
$\alpha\otimes\varphi$. Uniqueness in \cref{support:34-17}
proves the identity.
\end{proof}

\subsubsection{Products and their analytic continuation}

At an interior junction of two point plates, only the sum of
the two adjoining arguments matters. More precisely, for
$\kappa_1,\kappa_2\in\mathsf{Tb}$,
$\kappa_0\in\overline{\mathsf{Tb}}$ and
$0<\epsilon\le\min(d_{\mathsf{Tb}}(\kappa_1),
d_{\mathsf{Tb}}(\kappa_2))$,
\begin{equation}\label{support:34-50}
 \begin{split}
 B_\varphi(\xi,\kappa_1)\mathsf T(\kappa_0)
 B_\psi(\kappa_2,\xi')
 =B_\varphi(\xi,\epsilon n_0)
 \mathsf T(\kappa_1+\kappa_0+\kappa_2-2\epsilon n_0)
 B_\psi(\epsilon n_0,\xi').
 \end{split}
\end{equation}
This follows by extracting the two semigroup factors with
\eqref{support:34-19}. Thus one may replace the adjoining arguments
by half their sum. An outer argument next to $\Omega$ may be
changed arbitrarily, since $\mathsf T(\kappa)\Omega=\Omega$.
The same statements hold for point fields and $\mathsf T_3$.

For $\xi_1,\ldots,\xi_{n-1}\in\mathsf{Tb}$, define
\begin{equation}\label{support:34-56}
 K_n(\xi;\underline\varphi)=
 \left\langle\Omega,
 B_{\varphi_1}(\kappa,\xi_1/2)
 B_{\varphi_2}(\xi_1/2,\xi_2/2)\cdots
 B_{\varphi_n}(\xi_{n-1}/2,\kappa)\Omega\right\rangle.
\end{equation}
Here $\kappa\in\mathsf{Tb}$ is arbitrary; the preceding outer
argument identity makes the answer independent of it. Set
$K_0=1$ and $K_1(\varphi)=
\langle\Omega,B_\varphi(\kappa,\kappa)\Omega\rangle$.
For
\[
 \mathsf{Tb}^{(n)}=\{x:x_1\in\mathsf{Tb},\
                   x_{j+1}-x_j\in\mathsf{Tb}\},
\]
let $\Psi(x;\underline\varphi)$ be the operator product in
\eqref{support:34-56}, applied to $\Omega$, with successive
differences $\xi_j=x_{j+1}-x_j$ and first outer argument
$x_1$ in place of $\kappa$; set $\Psi(\varnothing)=\Omega$.

\begin{proposition}\label{support:34-59}
The kernels $K_n$ are holomorphic, and $\Psi$ is norm holomorphic
on $\mathsf{Tb}^{(n)}$. They are multilinear in the labels, and for $n\ge2$,
\begin{equation}\label{support:34-60}
 |K_n(\xi;\underline\varphi)|\le
 \prod_{a=1}^n\bigl(2^\wp C_{\rm plt}
       (\lambda_{\varphi_a}+1)\|\varphi_a\|_{C^1}\bigr)
 \min\{1,\min_a d_{\mathsf{Tb}}(\xi_a)\}^{-\wp n}.
\end{equation}
For profiles whose supports are strictly ordered in time,
\begin{equation}\label{support:34-61}
 \sigma(\alpha_1\otimes\varphi_1,\ldots,\alpha_n\otimes\varphi_n)
 =\int\prod_a\alpha_a(y_a)
 K_n(y_2-y_1,\ldots,y_n-y_{n-1};\underline\varphi)\,dY.
\end{equation}
If the first support also lies in $y^0>0$, the class of this
tuple is
\begin{equation}\label{support:34-62}
 [(\alpha_1\otimes\varphi_1,\ldots,\alpha_n\otimes\varphi_n)]
       =\int\prod_a\alpha_a(y_a)\Psi(Y;\underline\varphi)\,dY.
\end{equation}
The vectors satisfy
\begin{equation}\label{support:34-64}
 \begin{split}
 \langle\Psi(x;\underline\varphi),\Psi(z;\underline\psi)\rangle
 =K_{n+n'}\bigl(&(x_n-x_{n-1})^*,\ldots,(x_2-x_1)^*,\\
 &x_1^*+z_1,z_2-z_1,\ldots,z_{n'}-z_{n'-1};
 \varphi_n,\ldots,\varphi_1,\psi_1,\ldots,\psi_{n'}\bigr).
 \end{split}
\end{equation}
When one side is empty, its formula is obtained by removing
the first outer argument of the other vector; when both sides
are empty, the value is one. In addition,
\begin{equation}\label{support:34-66}
 \overline{K_n(\xi;\underline\varphi)}
       =K_n(\bar\xi;\underline\varphi),\qquad
 K_2(\xi;\varphi,\psi)=k_{\varphi,\psi}(\xi).
\end{equation}
Under (R0), $K_1=0$.
\end{proposition}
\begin{proof}
Holomorphy follows by taking products and matrix elements of
the norm-holomorphic point plates. Choose $\kappa=n_0$ and
apply \eqref{support:34-18} to each factor to get
\eqref{support:34-60}. Taking adjoints reverses the first product;
at its junction with the second, \eqref{support:34-50} replaces
$x_1^*,z_1$ by half their sum. This gives the Gram formula,
including its empty-vector variants.

For the integral formulas choose separating times
$0=c_0<c_1<\cdots<c_n$ for positive-time profiles.
The slab product \eqref{support:32-45} and
\eqref{support:34-21} express the class as the integral of
\[
 \prod_a\alpha_a(y_a)
 \prod_a B_{\varphi_a}(y_a-c_{a-1}n_0,c_an_0-y_a)\Omega.
\]
At each junction the arguments sum to $y_{a+1}-y_a$.
Replacing them by half this sum gives \eqref{support:34-62}.
Taking the inner product with $\Omega$ gives
\eqref{support:34-61}. A common translation reduces arbitrary
time-ordered supports to this case.

Reality of $\sigma$ on real tests shows that the imaginary
part of the continuous kernel vanishes on real time-ordered
configurations: test it against arbitrary products of real
profiles in small disjoint neighborhoods. The holomorphic
identity theorem gives the first identity in
\eqref{support:34-66}. Applying the same argument to the two
integral representations of the two-point function identifies
$K_2$ with $k_{\varphi,\psi}$ first on $\mathbb R^2_+$ and
then on $\mathsf{Tb}$. Finally a profile of integral one
identifies $K_1(\varphi)$ with the one-point value, which
vanishes under (R0).
\end{proof}

For a permutation $\mathfrak s$ and a frame
$\upsilon\in\frac\pi2\mathbb Z$, set
\begin{equation}\label{support:34-69}
 \begin{split}
 \mathcal Z(\mathfrak s,\upsilon)
 &=\{Z:z_{\mathfrak s(a+1)}-z_{\mathfrak s(a)}
                         \in R_\upsilon\mathsf{Tb},\ a<n\},\\
 K_\upsilon^{\mathfrak s}(Z;\underline\varphi)
 &=K_n\bigl(R_{-\upsilon}\partial(\mathfrak sZ);
                              \mathfrak s\underline\varphi\bigr).
 \end{split}
\end{equation}
Here $\partial Z=(z_{a+1}-z_a)_{a<n}$, and permutations act
on points and labels in the same order. Let $\mathcal Z_n$
be the union of these tubes and
$\mathcal O_n=\mathcal Z_n\cap(\mathbb R^2)^n$.

\begin{lemma}\label{support:34-72}
The functions $K_\upsilon^{\mathfrak s}$ agree on overlaps
and define one holomorphic function
$\mathcal K_n(Z;\underline\varphi)$ on the connected open set
$\mathcal Z_n$. This set is stable under common complex
translations, permutations, quarter turns, conjugation, and
$R_{is}$ for real $s$. The kernel is invariant under the first
three operations, with labels permuted together with points,
and obeys
$\overline{\mathcal K_n(Z;\underline\varphi)}
=\mathcal K_n(\bar Z;\underline\varphi)$.
Its real domain is
\[
 \mathcal O_n=\{Y:\text{all }y_a^0\text{ are distinct, or all }
                         y_a^1\text{ are distinct}\}.
\]
For profiles whose product support lies in one tube, integration
of $K_\upsilon^{\mathfrak s}$ against their product gives
$\sigma(\alpha_1\otimes\varphi_1,\ldots,\alpha_n\otimes\varphi_n)$.
\end{lemma}
\begin{proof}
Each tube is a nonempty open convex cone, the preimage of
$(R_\upsilon\mathsf{Tb})^{n-1}$ under a surjective linear
difference map. A real configuration belongs to it precisely
when $n_\upsilon\cdot y_{\mathfrak s(a)}$ is strictly
increasing. Permuting the entries and making the quarter turn
therefore reduces its integral formula to
\eqref{support:34-61}, by cubic invariance of $\sigma$.

Every nonempty intersection of two tubes is convex and
contains the real part of each of its points. On its real
open subset, the two kernels have equal integrals against
all products of small supported profiles. Their difference
therefore vanishes there. The holomorphic identity theorem
on the connected intersection proves agreement everywhere.
This proves that the kernels glue without a choice of
continuation path.

The stated stability follows from the difference inequalities
defining the tubes; in light-cone coordinates $R_{is}$ multiplies
the two coordinates by positive real numbers. The kernel
symmetries follow from its definition and
\eqref{support:34-66}. There are at most $4n!$ tubes.
The complement of $\mathcal O_n$ is the finite union
\[
 \{y_a^0=y_b^0,\ y_k^1=y_l^1\},\qquad a\ne b,\ k\ne l.
\]
Each subspace has codimension two, so its finite complement
is connected. Every tube is connected and meets this real
set. Their union $\mathcal Z_n$ is therefore connected.
\end{proof}

\subsection{Rotation invariance of all correlations}
\label{sec:all-rotations}

We now assume (R0)--(R5). The polynomial growth of the point
plates implies that their kernels have only finitely many angular
frequencies. Positivity imposes an arithmetic restriction on the
highest frequency; this reduces the rotation problem to the
two-point function.

\subsubsection{Angular frequencies and their phases}

\begin{lemma}\label{support:37-2}
Put $\mathcal M_n=\{m\in\mathbb Z:4|m|\le\wp n\}$.
For each tuple of $n\ge2$ labels there are unique holomorphic
functions $c_m$ on $\mathcal Z_n$, zero for $m\notin\mathcal M_n$,
such that
\begin{equation}\label{support:37-4}
 \mathcal K_n(R_{is}Z;\underline\varphi)
       =\sum_{m\in\mathcal M_n}c_m(Z;\underline\varphi)e^{-4ms}.
\end{equation}
They inherit permutation and translation invariance and satisfy
\[
 c_m(R_{is}Z)=e^{-4ms}c_m(Z),\qquad
 \overline{c_m(X)}=c_{-m}(X)\quad(X\in\mathcal O_n).
\]
They extend uniquely to real-analytic functions $\widetilde c_m$
on configurations of distinct planar points, with
\begin{equation}\label{support:37-14}
 \widetilde c_m(v+R_t(X-v);\underline\varphi)
              =e^{4imt}\widetilde c_m(X;\underline\varphi).
\end{equation}
Here a common vector $v$ is subtracted from every point.
\end{lemma}
\begin{proof}
First let the consecutive differences of a real configuration
$X$ make angles of modulus less than $\pi/8$ with $n_0$.
For every complex $\zeta=t+is$, choose a quarter-turn frame
$\upsilon$ with $|t-\upsilon|\le\pi/4$.
All the differences of $R_\zeta X$ lie in
$R_\upsilon\mathsf{Tb}$, with tube distances at least
\[
 \min_a|x_{a+1}-x_a|\cos(3\pi/8)e^{-|s|}.
\]
Thus $h_X(\zeta)=\mathcal K_n(R_\zeta X)$ is entire,
$\pi/2$-periodic, and bounded by $C_Xe^{\wp n|s|}$.
Shifting the Fourier contour vertically shows that its
coefficient at $e^{4im\zeta}$ vanishes if $4|m|>\wp n$.
Fourier uniqueness on the real circle then gives a finite
expansion, valid throughout the complex plane.

To define the coefficients at every $Z$, let
$\mu=\#\mathcal M_n$ and invert the Vandermonde matrix
$V_{lm}=e^{-4ml}$, $0\le l<\mu$. If $W=V^{-1}$, put
\[
 c_m(Z)=\sum_{l=0}^{\mu-1}W_{ml}\mathcal K_n(R_{il}Z).
\]
These functions are holomorphic. On each tube, permuted
quarter-turned configurations of the preceding narrow angular
type form a nonempty real open set. On that set the formula
recovers the Fourier coefficients. The holomorphic identity
theorem on the tube proves \eqref{support:37-4} everywhere.
The same Vandermonde comparison proves uniqueness and all
the stated symmetry identities.

We give the phase argument explicitly because $\mathcal O_n$
need not contain an entire real rotation orbit. If $X$ is a
real point of a fixed tube, write its successive differences
in that frame as $r_an_{\vartheta_a}$,
$|\vartheta_a|<\pi/2$. For
\[
 t\in I_X=\bigcap_a(-\pi/2-\vartheta_a,\pi/2-\vartheta_a),
\]
the full vertical strip $R_{t+is}X$ remains in this tube.
Analytically continuing \eqref{support:37-4} in the parameter and
comparing coefficients gives
$c_m(R_tX)=e^{4imt}c_m(X)$ on $I_X$.

Consider next configurations with no coincident points and no
perpendicular pair of distinct difference vectors. They form
a rotation-invariant dense open set contained in $\mathcal O_n$.
For density, the product of all squared differences and all
such scalar products is a polynomial nonzero at
$((a,a^2))_{a=1}^n$; on the line segment from any configuration
to this one it has only finitely many zeros.
For these configurations the preceding local phase identity
holds along the entire orbit, so
$t\mapsto e^{-4imt}c_m(R_tX)$ is locally constant on $\mathbb R$,
hence constant. Continuity extends the identity to every
$X,R_tX\in\mathcal O_n$.

For any configuration of distinct points, all but finitely
many angles modulo $\pi$ give distinct time coordinates.
Choose such a $t$ and define
\[
 \widetilde c_m(X)=e^{-4imt}c_m(R_tX).
\]
The preceding identity makes this independent of $t$.
Near each configuration one may keep $t$ fixed, giving a
holomorphic local extension and hence real analyticity.
Translation invariance and the same definition give
\eqref{support:37-14}, including rotations about arbitrary centers.
\end{proof}

Set $\bar m(0)=\bar m(1)=0$. For $n\ge2$, let $\bar m(n)$
be the largest $|m|$ for which $c_m$ is nonzero for some labels,
or zero if all kernels vanish. Then
$0\le\bar m(n)\le\wp n/4$. A nonzero harmonic is nonzero at
some strictly time-ordered real configuration, after permuting
labels: otherwise the identity theorem on a tube would make
it zero. Conjugation shows that, if $\bar m(n)>0$, the
coefficient $c_{-\bar m(n)}$ is nonzero at such a configuration.
Write $\widetilde K_n=\sum_m\widetilde c_m$.

\subsubsection{Positivity of the highest harmonic}

\begin{lemma}\label{support:37-18}
For all $a,b\ge0$,
\begin{equation}\label{support:37-19}
 2\bar m(a+b)\le\bar m(2a)+\bar m(2b).
\end{equation}
For $a\ge1$, let
$t_a=\widetilde c^{(2a)}_{-\bar m(2a)}$ on labelled
configurations, where a labelled configuration is the unordered
set of its distinct points and their labels. On configurations
of $a$ points in $\mathbb R^2_+$,
\[
 G_a(S,T)=t_a(\theta S\cup T)
\]
is a positive Hermitian kernel. If $\mathsf L$ is any line,
and $S,T$ lie strictly on one side, then
\begin{equation}\label{support:37-25}
 |t_a(\theta_{\mathsf L}S\cup T)|^2
 \le |t_a(\theta_{\mathsf L}S\cup S)|
       |t_a(\theta_{\mathsf L}T\cup T)|.
\end{equation}
\end{lemma}
\begin{proof}
For strictly time-ordered positive configurations $x,z$,
the Gram formula and $(R_{is}x)^*=R_{is}x^*$ give
\[
 \langle\Psi(R_{is}x),\Psi(R_{is}z)\rangle
 =\mathcal K_{a+b}(R_{is}(\theta x_a,\ldots,\theta x_1,
                                    z_1,\ldots,z_b)).
\]
An empty configuration gives $\Omega$. If $\bar m(a+b)=0$,
\eqref{support:37-19} follows from nonnegativity. Otherwise, by symmetry,
take $b\ge1$. Choose a time-ordered configuration where
$c_{-\bar m(a+b)}\ne0$ and translate a separating time between
its first $a$ and last $b$ points to zero; when $a=0$, choose
this time below the first point. Cauchy--Schwarz for the two resulting vectors bounds
the squared expansion with leading exponent
$8\bar m(a+b)s$ by a product whose maximal exponent is
$4(\bar m(2a)+\bar m(2b))s$. Letting $s\to+\infty$ proves
\eqref{support:37-19}.

For a finite family of $a$-point positive configurations,
multiply its Gram matrix at $R_{is}$ by
$e^{-4\bar m(2a)s}$ and let $s\to+\infty$.
Its limit is the matrix of $G_a$, hence positive Hermitian.
If some times agree, perturb the $k$th time by $k\epsilon$,
avoiding the finitely many exceptional $\epsilon$, and use
continuity of $\widetilde c_m$.

If the positive normal to $\mathsf L$ has angle $\upsilon$
and $v\in\mathsf L$, the map $\iota y=R_{-\upsilon}(y-v)$
takes its selected side to $\mathbb R^2_+$. The phase law gives
\[
 t_a(\theta_{\mathsf L}S\cup T)
       =e^{-4i\bar m(2a)\upsilon}G_a(\iota S,\iota T).
\]
Cauchy--Schwarz for this positive kernel proves
\eqref{support:37-25}; the common phase has modulus one.
\end{proof}

\begin{proposition}\label{support:37-27}
One has
\begin{equation}\label{support:37-29}
 \bar m(2a)=a\bar m(2),\qquad
 \bar m(n)\le\frac n2\bar m(2).
\end{equation}
\end{proposition}
\begin{proof}
We first prove $a\mid2\bar m(2a)$. Let $\mathcal N_a(S)$
mean that $t_a$ vanishes on every $2a$-point labelled
configuration containing $S$.

We need a halving observation. Suppose $S'$ has $j'\le a$
points on one side of a line $\mathsf L$, and
$S\subset\theta_{\mathsf L}S'\cup S'$ satisfies
$\mathcal N_a(S)$. Fix the points and labels of
$\theta_{\mathsf L}S'$ and arbitrary labels on another
$2a-j'$ points. The resulting value of $t_a$ is real
analytic in these free points on the complement of their
coincidence sets. This domain is connected, since the excluded
affine subspaces have codimension two. On the nonempty open
subset where $a-j'$ free points lie on the negative side and
$a$ on the positive side, write the configuration as
$\theta_{\mathsf L}T\cup T'$ with $S'\subset T$.
The diagonal value $t_a(\theta_{\mathsf L}T\cup T)$ vanishes
by $\mathcal N_a(S)$; \eqref{support:37-25} makes the mixed value
zero. The real-analytic identity theorem therefore gives
$\mathcal N_a(\theta_{\mathsf L}S')$ on all configurations.
Exchanging the two sides proves $\mathcal N_a(S')$ as well.

Suppose now $a\nmid2\bar m(2a)$, and fix any label $\varphi$.
Give all vertices of a regular $2a$-gon this label. Rotation
through $\pi/a$ preserves the labelled configuration, whereas
the phase law multiplies its top harmonic by
$e^{-4i\bar m(2a)\pi/a}\ne1$. Thus $\mathcal N_a$ holds for
every full polygon. We descend from arcs of length $j\ge2$
to arcs of length $j'=\lceil j/2\rceil$. Write the vertices
as $P_k=p+\rho n_{\omega+k\pi/a}$. Reflection in the line
through $p$ of direction
$\omega+(k_0+1/2)\pi/a$ sends $P_{k_0+k}$ to
$P_{k_0+1-k}$. The arc $P_{k_0+1},\ldots,P_{k_0+j'}$
lies strictly on one side, and its reflected union is an
arc of length $2j'$ containing an arc of length $j$.
The halving observation proves $\mathcal N_a$ for the shorter
arc. Iteration reaches a singleton at any prescribed point.
Since the label was arbitrary, $t_a$ then vanishes on every
configuration. This contradicts the definition of the
nonzero top harmonic. Hence the divisibility holds.

Put $\nu(a)=\bar m(2a)$. By \eqref{support:37-19}, $\nu$ is a
convex integer sequence with $\nu(0)=0$ and
$\nu(a)\le\wp a/2$. Its increments are increasing bounded
integers, hence eventually equal to an integer $d$.
Thus $\nu(a)=da-c$ eventually, where
$c=\sum_k(d-(\nu(k)-\nu(k-1)))\ge0$ is a finite integer.
Since $a\mid2\nu(a)$, every sufficiently large $a$ divides
$2c$, forcing $c=0$. All increments equal $d=\nu(1)$.
This proves the even identity. Apply \eqref{support:37-19} with
$a=\lfloor n/2\rfloor$ and $b=n-a$ to obtain the second one.
\end{proof}

\subsubsection{From kernels to separated tests}

\begin{proposition}\label{support:37-30}
For arbitrary profiles with pairwise disjoint supports,
\begin{equation}\label{support:37-31}
 \sigma(\alpha_1\otimes\varphi_1,\ldots,\alpha_n\otimes\varphi_n)
       =\int\prod_a\alpha_a(y_a)
                 \widetilde K_n(Y;\underline\varphi)\,dY.
\end{equation}
If $\bar m(n)=0$, the functional is invariant under $R_t$
on all $n$-tuples with pairwise disjoint planar projections.
\end{proposition}
\begin{proof}
Subtract the right side of \eqref{support:37-31} from the left
and call the difference $D(\underline\alpha)$.
It is multilinear when each profile is restricted to its
own separated compact support. We first use real profiles;
complexification then gives the general statement.

Let $A_a^0,A_a^1$ be the time and space interval hulls of
the supports. Suppose each pair whose $A_a^1$ intersect
has a separating time level outside every $A_e^0$.
Induct on the number of such pairs. When there are none,
all supports lie in one permuted quarter-turned tube, so
\cref{support:34-72} gives $D=0$.
For the induction step choose a level $c$ and split the
profiles below and above it. Translate the upper group
by $sn_{\pi/2}$. Both terms of $D$ are real analytic in
$s$. For the correlation this follows from the matrix
element of $\mathsf T(2\delta,s)$, where $2\delta>0$ is
a time gap about $c$; it is holomorphic on
$|\operatorname{Im}s|<2\delta$. For the kernel integral,
the compact product support remains in the distinct-point
domain for every real $s$. Local holomorphic extensions
on a finite cover of this compact set justify integration
and give real analyticity. For large $s$ all cross-group
spatial intervals are disjoint, so the number of
intersecting pairs decreases. The induction hypothesis
and the identity theorem imply $D=0$ also at $s=0$.

For general profiles with separation $\delta$, partition
the plane into the $C^1$ squares obtained from
$\chi(2y^0/\epsilon-j^0)\chi(2y^1/\epsilon-j^1)$, with
$0<\epsilon<\delta/(n+1)$. Each nonzero tuple of pieces
has interval lengths at most $\epsilon$. If two spatial
intervals overlap, their time intervals have a gap at
least $\sqrt{\delta^2-4\epsilon^2}>(n-2)\epsilon$.
The other $n-2$ time intervals cannot cover this gap,
so there is a level outside all intervals. The preceding
argument applies to every tuple of pieces; multilinearity
proves \eqref{support:37-31}.

If $\bar m(n)=0$, the phase law makes
$\widetilde K_n=\widetilde c_0$ invariant under all real
rotations. Changing variables in \eqref{support:37-31} proves
invariance for plates. For a general tuple with disjoint
planar projections, convolve each entry with the product
mollifier
\[
 J_\epsilon(x)=\prod_{\mu=1}^3
             \frac2\epsilon\chi(2x_\mu/\epsilon).
\]
For fixed sufficiently small $\epsilon$, the convolution
Riemann sums are finite sums of translates of $J_\epsilon$,
each of which is a plate. Their planar supports stay in
fixed mutually disjoint neighborhoods. Multilinearity
proves rotation invariance for these sums. They converge
in $C^1$ to the convolutions, which converge in $C^1$ to
the original entries. The telescoping consequence of (R2),
\[
 |\sigma(g_1,\ldots,g_n)-\sigma(g'_1,\ldots,g'_n)|
 \le C_n^\sigma(d;R)\sum_a\|g_a-g'_a\|_\sharp
             \prod_{b\ne a}\max(\|g_b\|_\sharp,\|g'_b\|_\sharp),
\]
on fixed separated compact supports permits both limits.
\end{proof}

\begin{theorem}[Rotation invariance]\label{support:37-34}
Every functional satisfying (R0)--(R7) is invariant under
$\operatorname{O}(3)$.
\end{theorem}
\begin{proof}
The two-point result \cref{support:33-57} and the identification
$K_2=k$ show that $K_2(R_{is}\xi)$ is independent of $s$
for real $\xi\in\mathbb R^2_+$. Comparing the coefficients
in \eqref{support:37-4}, and then applying the identity theorem,
gives $\bar m(2)=0$. By \eqref{support:37-29}, all $\bar m(n)$
vanish. \Cref{support:37-30} gives invariance under
$R_t$ on tuples with disjoint planar projections, and
\cref{support:32-48} gives it on every separated tuple.
The empty and one-point values are invariant by normalization
and (R0).

Cubic conjugation gives rotations about the other coordinate
axes. These generate $\operatorname{SO}(3)$: rotate an arbitrary
unit vector to $e_3$ by two coordinate-plane rotations,
then the remaining transformation fixes $e_3$ and is a
rotation in its perpendicular plane. A coordinate reflection
already belongs to the cubic group, so adjoining it gives
all of $\operatorname{O}(3)$.
\end{proof}

\section{Clustering and decomposition}\label{app:clustering-decomposition}

\subsection{Uniform clustering on a lattice}
\label{sec:lattice-clustering}

The estimate needed for decomposition into continuum components must
be uniform as the spacing tends to zero.  A decay estimate involving
the supremum of a lattice polynomial would not suffice.  We instead
bound correlations by the reflection norms of their local factors.
The argument uses coordinate transfer operators, the spectral
restriction imposed by a diagonal reflection, and a Fourier cutoff
which removes zero energy.

Throughout this section, $\nu$ is a translation-invariant probability
law on the configurations over $a\mathbb Z^3$, $a>0$.  We assume RP
at the six coordinate mirrors
$\Sigma(e_\mu,0)$ and $\Sigma(e_\mu,a/2)$, and at one diagonal
mirror $\Sigma(e_\mu+\epsilon e_\lambda,c)$ in each coordinate
plane, where $\epsilon\in\{1,-1\}$ and $c\in a\mathbb Z$.
Translations move these diagonal mirrors to height zero and move
the coordinate mirrors to every height in $(a/2)\mathbb Z$.
In particular these hypotheses hold for every periodic Gibbs state
by \Cref{support:54-17}.

\subsubsection{Transfer operators and the lattice spectral restriction}

For a mirror $\Sigma(n,c)$ and a positive normal lattice vector $v$,
suppose RP also holds at $\Sigma(n,c+n\cdot v/2)$.  Complete the
bounded local functions supported in the upper half-space in the form
$\langle[F],[G]\rangle=\nu(\overline{\theta_{n,c}F}G)$, after
quotienting its null space.  This is the reflection Hilbert space at
the mirror.  Coordinate pictures use $n=e_\mu$ and $v=ae_\mu$;
diagonal pictures use $n=e_\mu+\epsilon e_\lambda$ and $v=an$.

\begin{lemma}\label{support:30-5}
In these pictures the normal translation and the tangential
translations induce commuting operators
\[
 T[F]=[\tau_vF],\qquad V(w)[F]=[\tau_wF]\quad(n\cdot w=0),
 \qquad 0\le T\le1,
\]
where the $V(w)$ are unitary.  Moreover
\begin{equation}\label{support:30-9}
 \langle[F],T^kV(w)[G]\rangle
 =\nu(\overline{\theta_{n,c}F}\,\tau_{kv+w}G),
 \qquad k\ge0,
\end{equation}
and $\|[F]\|^2\le\nu(|F|^2)\le\|F\|_\infty^2$.
In a coordinate picture there is a joint spectral measure
$\mathfrak E$ on $[0,1]\times(S^1)^2$ such that
\begin{equation}\label{support:30-17}
 T^kV(a y)=\int t^k\zeta^y\,d\mathfrak E(t,\zeta),
 \qquad y_\mu=0.
\end{equation}
\end{lemma}

\begin{proof}
Polarization of RP, followed by a translation of an arbitrary local
function into a half-space, gives invariance under the mirror.
Ordinary Cauchy--Schwarz therefore bounds the reflection norm by
the $L^2(\nu)$ norm.  Normal translation is symmetric in the
reflection form.  Reflection Cauchy--Schwarz gives
\[
 \|[\tau_vF]\|^2
 \le\|[F]\|\,\|[\tau_{2v}F]\|.
\]
Iteration and the uniform bound by $\|F\|_\infty$ show
$\|[\tau_vF]\|\le\|[F]\|$.  Thus translation descends through
the null space to a self-adjoint contraction.  The identity
$\theta_{n,c}=\theta_{n,c+n\cdot v/2}\tau_v$ and RP at the
second mirror show $\langle[F],T[F]\rangle\ge0$.
Tangential translations preserve the form and have inverse
translations, hence induce unitaries.  Their commutation and
\eqref{support:30-9} follow on local classes and extend by continuity.
The joint spectral theorem for the commuting normal operators
$T,V(ae_\lambda),V(ae_{\lambda'})$ gives
\eqref{support:30-17}.
\end{proof}

The following elementary complex-variable estimate isolates the
spectral consequence of comparing a coordinate picture with a
diagonal picture.

\begin{lemma}\label{support:30-21}
Suppose a finite positive measure $\eta$ on the closed unit disc
satisfies, for every polynomial $q$,
\begin{equation}\label{support:30-22}
 \int|q|^2\,d\eta\le C\sup_{[0,1]}|q|\sup_{S^1}|q|.
\end{equation}
Then $\eta\{|1+z|^2<4|z|\}=0$.
\end{lemma}

\begin{proof}
Uniform Taylor approximation extends \eqref{support:30-22} to functions
holomorphic on a disc of radius greater than one.  On the unit disc
put
\[
 \Lambda(\xi)=\log\frac{1+\xi}{1-\xi},\qquad
 g_s(z)=\cos\bigl(s\Lambda(\sqrt z)\bigr),\qquad s>0.
\]
The principal logarithm is analytic because its argument is in the
right half-plane.  Since $\Lambda$ is odd and cosine is even,
$g_s$ is a single-valued holomorphic function of $z$.
Write $A(z)=|\operatorname{Im}\Lambda(\sqrt z)|$.  Then
\[
 \sinh^2(sA(z))\le|g_s(z)|^2\le\cosh^2(sA(z)),
 \qquad 0\le A(z)<\pi/2,
\]
and $A=0$ on $[0,1)$.  Applying \eqref{support:30-22} to
$g_s(rz)$, $0<r<1$, gives
\[
 \int\sinh^2(sA(rz))\,d\eta(z)\le C\cosh(s\pi/2).
\]
Extend $A$ to the unit circle by $A=\pi/2$ away from $1$ and
$A(1)=0$; these are its radial limits.  Bounded convergence allows
$r\uparrow1$.  For every $\delta>0$ the last inequality implies
\[
 \eta\{A\ge\pi/4+\delta\}\le16C e^{-2s\delta}
\]
for all sufficiently large $s$, so this set is null.
Inside the disc,
\[
 \tan A(z)=\frac{\sqrt{2(|z|-\operatorname{Re}z)}}{1-|z|}.
\]
Consequently $A(z)>\pi/4$ is equivalent to
$|1+z|^2<4|z|$, with the same equivalence at the boundary by the
radial values.  A countable union over $\delta>0$ proves the claim.
\end{proof}

\begin{proposition}\label{support:30-32}
In every coordinate picture and for each tangential direction
$\lambda\ne\mu$,
\begin{equation}\label{support:30-33}
 (1-T)^2\ge T\bigl(2-V(ae_\lambda)-V(ae_\lambda)^*\bigr).
\end{equation}
Equivalently, its joint spectral measure is supported where
$(1-t)^2\ge t|1-\zeta_\lambda|^2$.  In particular, with
\begin{equation}\label{support:30-42}
 \mathsf k=\frac{2}{\pi}\log(1+\sqrt2)>0,
\end{equation}
the spectral measure is supported where
\begin{equation}\label{support:30-45}
 t\le\exp\{-\mathsf k|\arg\zeta_\lambda|\}.
\end{equation}
\end{proposition}

\begin{proof}
Use the available diagonal normal
$n=e_\mu+\epsilon e_\lambda$ and write
$n^\perp=e_\mu-\epsilon e_\lambda$.
For a local upper-half-space function $F$, let $\varsigma$ be the
image of its spectral measure
$\langle[F],\mathfrak E(\cdot)[F]\rangle$ under
$(t,\zeta)\mapsto z=t\zeta_\lambda^\epsilon$.
Expansion of a polynomial $q(z)=\sum q_kz^k$ gives
\begin{equation}\label{support:30-35}
 \int|q(z)|^2|z|^{2N}\,d\varsigma
 =\sum_{k,l}q_k\overline{q_l}\,
 \nu\bigl(\overline{\theta F}\,
       \tau_{kan+lan^\perp}\tau_{2Nae_\mu}F\bigr).
\end{equation}
Choose a diagonal mirror $\Sigma(n,c)$ beyond every endpoint of
the bonds supporting $\theta F$, then choose $N$ so that every
endpoint supporting $\tau_{2Nae_\mu}F$ lies above that mirror.
Both choices depend only on $F$.  Set
$F_1=\theta_{n,c}\theta F$ and $F_2=\tau_{2Nae_\mu}F$.
Both functions lie in its upper half-space.  In the diagonal
picture, the right side of \eqref{support:30-35} equals
\[
 \langle[F_1],q(T_\parallel)\overline q(V_\perp)[F_2]\rangle,
\]
where $\overline q$ denotes coefficient conjugation.
Its modulus is at most
$\|F\|_\infty^2\sup_{[0,1]}|q|\sup_{S^1}|q|$.
\Cref{support:30-21} applied to $|z|^{2N}d\varsigma$ shows
that $\varsigma$ gives no mass to $|1+z|^2<4|z|$: the weight
is strictly positive there, since $z=0$ is outside this set.
The identity
\[
 |1+t\zeta_\lambda^\epsilon|^2-4t
 =(1-t)^2-t|1-\zeta_\lambda|^2
\]
proves the spectral restriction for every local vector, hence for
the whole spectral measure by density.  Functional calculus gives
\eqref{support:30-33}.

For $t=e^{-x}>0$ and $\zeta_\lambda=e^{iy}$, the restriction is
$\sinh(x/2)\ge|\sin(y/2)|$.  The function
$2\operatorname{arcsinh}(\sin(y/2))$ is concave on $[0,\pi]$:
its derivative is
$\cos(y/2)(1+\sin^2(y/2))^{-1/2}$, a product of nonnegative
decreasing functions.  It therefore lies above the chord
$\mathsf k y$.  This proves \eqref{support:30-45}; the case $t=0$
is immediate.
\end{proof}

\subsubsection{The invariant part of a correlation}

Let $\mathcal I$ be the sigma-algebra of translation-invariant
Borel events, completed modulo $\nu$-null sets.  In a coordinate
picture let $\Pi_0$ be the orthogonal projection onto
$\ker(1-T)$.  It describes the nondecaying part of a correlation.

\begin{proposition}\label{support:31-12}
One has $T^n\to\Pi_0$ strongly and $\Pi_0V(w)=\Pi_0$ for
tangential translations.  Moreover
\begin{equation}\label{support:31-14}
 \langle[F],\Pi_0[G]\rangle
 =\nu\bigl(\overline{E_\nu[F\mid\mathcal I]}
                   E_\nu[G\mid\mathcal I]\bigr).
\end{equation}
Conditional expectation onto $\mathcal I$ is unchanged by
translations or coordinate reflections.  In particular,
$\nu$ is ergodic precisely when $\ker(1-T)=\mathbb C[1]$.
\end{proposition}

\begin{proof}
By the spectral theorem,
$\Pi_0=\mathfrak E(\{t=1\})$ and $T^n\to\Pi_0$ strongly.
\Cref{support:30-32} in both tangential directions says that
$t=1$ forces $\zeta=(1,1)$, proving $\Pi_0V(w)=\Pi_0$.
Let $P_\mu$ be the $L^2(\nu)$ projection onto the functions fixed
by $\tau_{ae_\mu}$.  The mean ergodic theorem and
\eqref{support:30-9} give, for upper-half-space functions,
\begin{equation}\label{support:31-16}
 \langle[F],\Pi_0[G]\rangle
 =\langle\theta F,P_\mu G\rangle_{L^2(\nu)}.
\end{equation}
Since $\Pi_0V(w)=\Pi_0$, the right side is unchanged on replacing
$G$ by $\tau_wG$.  For arbitrary local $F_1,F_2$, choose
$u=kae_\mu$ large enough that
$F=\theta\tau_{-u}F_1$ and $G=\tau_uF_2$ are in the upper
half-space.  Substitution in this identity, using
$P_\mu\tau_{\pm u}=P_\mu$, gives
\[
 \langle F_1,P_\mu\tau_wF_2\rangle_{L^2(\nu)}
 =\langle F_1,P_\mu F_2\rangle_{L^2(\nu)}.
\]
Local functions are dense in $L^2(\nu)$, so
$P_\mu\tau_w=P_\mu$ and, by adjoints,
$\tau_wP_\mu=P_\mu$.  Thus a function fixed by the one normal
translation is fixed by every lattice translation.  The invariant
subspace is $L^2(\mathcal I,\nu)$, so
$P_\mu=E_\nu[\,\cdot\mid\mathcal I]$.

Reflection preserves that subspace and its orthogonal complement,
hence commutes with $P_\mu$.  Positivity of the left side of
\eqref{support:31-16} at $G=F$ implies
$\langle\theta P_\mu F,P_\mu F\rangle_{L^2(\nu)}\ge0$.
The vectors $P_\mu F$ with $F$ in the upper half-space are dense
in the invariant subspace: translate each local function there
without changing its projection.  Reflection is consequently a
positive involution on that subspace and must be the identity;
its negative eigenspace is zero.  Formula~\eqref{support:31-16} now
becomes \eqref{support:31-14}.
If $\nu$ is ergodic, the right side is
$\overline{\nu(F)}\nu(G)$, so $\Pi_0$ has range $\mathbb C[1]$.
Conversely that range identity implies
$\|E_\nu[F\mid\mathcal I]\|_2=|\nu(F)|$, first for local
$F$ and then for all $L^2$ functions.  For $F=1_E$, $E\in\mathcal I$,
it says $\nu(E)^{1/2}=\nu(E)$, proving ergodicity.
\end{proof}

\subsubsection{A bounded density after removing zero energy}

Write $Q_r=[-ra,ra]^3$, where $r\ge1$ is an integer, and
$\theta_\mu=\theta_{e_\mu,0}$.  For a bounded Borel function $A$
of the bonds in $Q_r$, use the site picture to set
\[
 \psi_A^\mu=[\tau_{rae_\mu}A]_\mu.
\]
If $\mathsf S_8$ is the group of coordinate sign changes, define
\begin{align}
 \mathcal N_r(A)&=\max_{\mu,\,s\in\mathsf S_8}
                      \|\psi_{sA}^\mu\|^2,
 \label{support:31-32}\\
 \mathcal N'_{r,\mu}(A)&=\max\left\{
   \|\psi_A^\mu\|^2,
   \max_{\lambda\ne\mu}
      \|\psi_{\theta_\lambda A}^\lambda\|
      \|\psi_A^\lambda\|\right\}.
 \label{support:31-33}
\end{align}
Then $\mathcal N'_{r,\mu}(A)\le\mathcal N_r(A)
\le\nu(|A|^2)\le\|A\|_\infty^2$.
The square root of $\mathcal N_r$ is a seminorm, invariant under
$\mathsf S_8$.  Consequently the parity parts
$A_\pm=(A\pm\theta_\mu A)/2$ satisfy
$\mathcal N_r(A_\pm)\le\mathcal N_r(A)$.
Translation in the reflection form also gives
\begin{equation}\label{support:31-38}
 \nu(\overline A\tau_xB)
 =\langle\psi_{\theta_\mu A}^\mu,
       T_\mu^{x_\mu/a-2r}V_\mu(x_\perp)\psi_B^\mu\rangle,
 \qquad x_\mu\ge2ra.
\end{equation}
The correlation at $-x$ is the conjugate of
$\nu(\overline B\tau_xA)$; coordinate reflection changes
$x_\mu$ to $-x_\mu$ and reflects both functions.  These identities
allow the same correlation to be estimated in any coordinate picture.

We require a scalar weight which vanishes at $t=1$ but controls
$1-t$ from below.  Let $z_m$ be the even, period-four sequence
$(z_0,z_1,z_2,z_3)=(0,1,-2,1)$, and let
\[
 B(u)=8(1-4|u|/\pi)_+\quad(-\pi\le u\le\pi)
\]
be extended with period $2\pi$.  Its Fourier coefficients are
\[
 b_0=1,\qquad b_m=\left(\frac{8\sin(m\pi/8)}{m\pi}\right)^2
 \quad(m\ne0).
\]
Integration by parts gives this formula; in particular the Fourier
series converges absolutely.  Put
\begin{equation}\label{support:31-22}
 B_z(u)=\sum_{m\in\mathbb Z}z_mb_me^{imu},\qquad
 \Psi(t)=2\sum_{m\ge1}z_mb_mt^{m-1},\quad 0\le t\le1.
\end{equation}
The convention $0^0=1$ is used here.

\begin{lemma}\label{support:31-23}
The above series converge uniformly.  One has
\[
 \sum_m|z_m|b_m=8,\qquad |B_z|\le8,\qquad
 B_z(u)=0\quad\text{if }\operatorname{dist}(u,2\pi\mathbb Z)\le\pi/4.
\]
Moreover $0\le\Psi\le8$, $\Psi(1)=0$, $\Psi(0)=2b_1>0$, and
\begin{equation}\label{support:31-26}
 \Psi(e^{-x})\ge c_\Psi\tanh x\quad(x>0),\qquad
 c_\Psi=(1+\sqrt2)^{-1}.
\end{equation}
\end{lemma}

\begin{proof}
The identities $|z_m|=1-\cos(m\pi/2)$ and
$z_m=\cos(m\pi/2)-\cos(m\pi)$ give
\[
 \sum_m|z_m|b_m=B(0)-B(\pi/2)=8,
 \qquad B_z(u)=\tfrac12B(u+\pi/2)+\tfrac12B(u-\pi/2)-B(u+\pi).
\]
The three translated triangles have disjoint interiors of their
supports and vanish for $\operatorname{dist}(u,2\pi\mathbb Z)
\le\pi/4$.  This proves the stated bounds and, at $u=0$,
$\Psi(1)=0$.  Absolute summability gives the uniform convergence.
For $0<t<1$, the Poisson kernel
$P_t(u)=(1-t^2)/(1-2t\cos u+t^2)$ yields
\[
 t\Psi(t)=\frac1{2\pi}\int_{-\pi/4}^{\pi/4}B(v)
 \left[\tfrac12P_t(v-\pi/2)+\tfrac12P_t(v+\pi/2)-P_t(v-\pi)\right]dv.
\]
Writing $s=1+t^2$, the bracket is at least
\[
 (1-t^2)\left[\frac1s-\frac1{s+\sqrt2t}\right]
 =\frac{(1-t^2)\sqrt2t}{s(s+\sqrt2t)}:
\]
the average of $(s\mp2t\sin v)^{-1}$ is at least $1/s$, whereas
$\cos v\ge1/\sqrt2$.  Since $(2\pi)^{-1}\int B=1$, division
by $t$ gives
\[
 \Psi(t)\ge\frac{1-t^2}{1+t^2}
                  \frac{\sqrt2}{1+t^2+\sqrt2t}
 \ge c_\Psi\frac{1-t^2}{1+t^2}.
\]
This is \eqref{support:31-26}.  Together with $\Psi(0)>0$ and
$\Psi(1)=0$ it proves positivity on $[0,1)$.
\end{proof}

Set
\begin{equation}\label{support:31-4}
 C_{\mathrm{dn}}=\frac1{4\pi^2}
 \left[128+64\left(\frac{16}{\pi\mathsf k}+2+
                              \frac{64}{(\pi\mathsf k)^2}\right)\right].
\end{equation}

\begin{lemma}\label{support:31-41}
If $\Phi$ is supported in $Q_r$ and $\theta_\mu\Phi=\epsilon\Phi$,
$\epsilon\in\{1,-1\}$, let
$\varrho_\psi=\langle\psi,\mathfrak E_\mu(\cdot)\psi\rangle$,
where $\psi=\psi_\Phi^\mu$.  The image of
$\Psi(t^{2r})d\varrho_\psi(t,\zeta)$ on $(S^1)^2$ has a
continuous density $h$ relative to unnormalized angular Lebesgue
measure $d\alpha$ on $(-\pi,\pi]^2$, and
\begin{equation}\label{support:31-44}
 0\le h\le C_{\mathrm{dn}}r^2\mathcal N'_{r,\mu}(\Phi).
\end{equation}
\end{lemma}

\begin{proof}
Write $u(v)=\nu(\overline{\theta_\mu\Phi}\tau_v\Phi)
=\epsilon\nu(\overline\Phi\tau_v\Phi)$ and
\[
 K(y)=\sum_{m\in\mathbb Z}z_mb_m\,
         u\bigl(a(2rme_\mu+y)\bigr),\qquad y_\mu=0.
\]
The series is absolutely convergent.  Reflection parity makes
$u(a(-ke_\mu+y))=u(a(ke_\mu+y))$.  By
\eqref{support:31-38}, for $k\ge2r$,
\[
 u(a(ke_\mu+y))=\langle\psi,T_\mu^{k-2r}V_\mu(ay)\psi\rangle.
\]
Pairing $m$ with $-m$, and observing $z_0=0$, therefore gives
\begin{equation}\label{support:31-48}
 K(y)=\int\Psi(t^{2r})\zeta^y\,d\varrho_\psi(t,\zeta),
 \qquad |K(y)|\le8\|\psi\|^2.
\end{equation}
Thus $K$ is the sequence of Fourier coefficients of the weighted
momentum measure.

To bound large $y$, choose a tangential coordinate $\lambda$ with
$|y_\lambda|=|y|_\infty$.  Since $K(-y)=\overline{K(y)}$, assume
$y_\lambda\ge2r$.  Applying \eqref{support:31-38} in the
$\lambda$ picture and summing its absolutely convergent series gives
\[
 K(y)=\epsilon\langle\psi_{\theta_\lambda\Phi}^\lambda,
  X\psi_\Phi^\lambda\rangle,
 \quad
 X=\int t^{y_\lambda-2r}
        B_z(2r\arg\zeta_\mu)\zeta_{\lambda'}^{y_{\lambda'}}
                  \,d\mathfrak E_\lambda.
\]
Where $B_z(2r\arg\zeta_\mu)$ is nonzero,
$|\arg\zeta_\mu|>\pi/(8r)$.  The lattice spectral restriction
\eqref{support:30-45} then gives $t\le e^{-\pi\mathsf k/(8r)}$.
Consequently
\begin{equation}\label{support:31-49}
 |K(y)|\le8e^{-\pi\mathsf k(|y|_\infty-2r)/(8r)}
                         \mathcal N'_{r,\mu}(\Phi),
 \qquad |y|_\infty\ge2r.
\end{equation}
There are at most $16r^2$ points with $|y|_\infty<2r$ and
$8j$ points with $|y|_\infty=j$.  For $c=\pi\mathsf k/(8r)$,
\[
 \sum_{j\ge2r}64j e^{-c(j-2r)}
 =64\left[\frac{2r}{1-e^{-c}}+
                 \frac{e^{-c}}{(1-e^{-c})^2}\right]
 \le64\left[2r+\frac{2r}{c}+\frac1{c^2}\right].
\]
Here $(1-e^{-c})^{-1}\le1+c^{-1}$ and
$e^{-c}(1-e^{-c})^{-2}\le c^{-2}$.
Together with \eqref{support:31-48}, this proves
\[
 \sum_y|K(y)|\le4\pi^2C_{\mathrm{dn}}r^2\mathcal N'_{r,\mu}(\Phi).
\]
The Fourier series
$h(\zeta)=(2\pi)^{-2}\sum_yK(y)\zeta^{-y}$ is therefore
uniformly convergent and bounded by the right side of
\eqref{support:31-44}.  Its angular measure has the same Fourier
coefficients as the weighted momentum measure by
\eqref{support:31-48}.  Density of trigonometric polynomials gives
equality of the measures.  Positivity of the weighted measure and
continuity of $h$ imply $h\ge0$ everywhere.
\end{proof}

\subsubsection{The clustering estimate}

\begin{proposition}\label{support:31-51}
Under the nine-mirror hypotheses of this section, put
\begin{equation}\label{support:31-5}
 C_{\mathrm{dc}}=16c_\Psi^{-1}\mathsf k^{-2}C_{\mathrm{dn}}.
\end{equation}
For every $A,B$ supported in $Q_r$ and every integer $n\ge1$,
\begin{equation}\label{support:31-52}
 0\le\langle\psi_A^\mu,T_\mu^n(1-\Pi_0)\psi_A^\mu\rangle
 \le C_{\mathrm{dc}}\mathcal N_r(A)
                  \left(\frac rn+2\frac{r^2}{n^2}\right).
\end{equation}
If $x\in a\mathbb Z^3$ and $n=|x|_\infty/a-2r\ge1$, then
\begin{multline}\label{support:31-53}
 \left|\nu(\overline A\tau_xB)
  -\nu\bigl(\overline{E_\nu[A\mid\mathcal I]}
                        E_\nu[B\mid\mathcal I]\bigr)\right|\\
 \le C_{\mathrm{dc}}\mathcal N_r(A)^{1/2}\mathcal N_r(B)^{1/2}
                  \left(\frac rn+2\frac{r^2}{n^2}\right).
\end{multline}
In particular, for $|x|\ge4\sqrt3\,ra$,
\begin{multline}\label{support:31-55}
 \left|\nu(\overline A\tau_xB)
  -\nu\bigl(\overline{E_\nu[A\mid\mathcal I]}
                        E_\nu[B\mid\mathcal I]\bigr)\right|\\
 \le4\sqrt3 C_{\mathrm{dc}}\,
       \mathcal N_r(A)^{1/2}\mathcal N_r(B)^{1/2}\frac{ra}{|x|}.
\end{multline}
The constants are independent of the spacing, the law, and the
observables.  For an ergodic law the subtracted term is
$\overline{\nu(A)}\nu(B)$.
\end{proposition}

\begin{proof}
Set $D_n=T_\mu^n(1-\Pi_0)=\int t^n1_{t<1}\,d\mathfrak E_\mu$.
It is positive and commutes with tangential translations.
First suppose $\theta_\mu\Phi=\pm\Phi$ and let
$\psi=\psi_\Phi^\mu$.  By \Cref{support:31-41}, the weighted
momentum measure has a bounded density.  It gives no mass to
$\zeta=(1,1)$; since $\Psi(t^{2r})>0$ for $t<1$, the same is
true of $\varrho_\psi$ on $\{t<1,\zeta=(1,1)\}$.
For $0<t<1$ write $s=-\log t$.  The two tangential versions of
\eqref{support:30-45} imply $s\ge\mathsf k|\alpha|_\infty$,
where $\zeta=e^{i\alpha}$.  Since
$s\mapsto e^{-ns}/\tanh(2rs)$ is decreasing and
$\Psi(t^{2r})\ge c_\Psi\tanh(2rs)$,
\[
 t^n\le c_\Psi^{-1}
 \frac{e^{-n\mathsf k|\alpha|_\infty}}
      {\tanh(2r\mathsf k|\alpha|_\infty)}\,
 \Psi(t^{2r})
\]
away from the null point $\alpha=0$.  At $t=0$, the integrand
$t^n$ is zero.  Integrating against the bounded density and using
$1/\tanh u\le1+1/u$ gives
\begin{align*}
 \langle\psi,D_n\psi\rangle
 &\le c_\Psi^{-1}C_{\mathrm{dn}}r^2\mathcal N'_{r,\mu}(\Phi)
    \int_{(-\pi,\pi]^2}
       \frac{e^{-n\mathsf k|\alpha|_\infty}}
            {\tanh(2r\mathsf k|\alpha|_\infty)}\,d\alpha\\
 &\le c_\Psi^{-1}C_{\mathrm{dn}}\mathsf k^{-2}
       \mathcal N'_{r,\mu}(\Phi)
          \left(\frac{8r^2}{n^2}+\frac{4r}{n}\right).
\end{align*}
Indeed the square integral equals
$8\int_0^\pi u e^{-n\mathsf k u}/\tanh(2r\mathsf k u)\,du$;
extend the upper limit to infinity to obtain the displayed bound.
For arbitrary $A=A_++A_-$, positivity yields
$\langle\psi_A,D_n\psi_A\rangle\le
2\langle\psi_{A_+},D_n\psi_{A_+}\rangle+
2\langle\psi_{A_-},D_n\psi_{A_-}\rangle$.
Use $\mathcal N'_{r,\mu}(A_\pm)\le\mathcal N_r(A)$ to obtain
\eqref{support:31-52} with \eqref{support:31-5}.

For \eqref{support:31-53}, choose $\mu$ with
$|x_\mu|=|x|_\infty$ and first suppose $x_\mu>0$.
Insert $T_\mu^n=\Pi_0+D_n$ into \eqref{support:31-38}.
\Cref{support:31-12} identifies the $\Pi_0$ term with the
subtracted conditional expectation.  Cauchy--Schwarz for $D_n$,
its commutation with $V_\mu(x_\perp)$, and
\eqref{support:31-52} bound the remainder as claimed.
For $x_\mu<0$, conjugate and interchange $A,B$.
If $|x|_\infty\ge4ra$, then
$n\ge|x|_\infty/(2a)\ge2r$, and
$r/n+2r^2/n^2\le4ra/|x|_\infty$.
Finally $|x|_\infty\ge|x|/\sqrt3$ proves
\eqref{support:31-55}.  The ergodic assertion follows from
\Cref{support:31-12}.
\end{proof}

\subsection{Decomposition by invariant conditional expectations}
\label{sec:conditional-decomposition}

A continuum limit of invariant states need not cluster.  We construct
its components from conditional expectations on the lattice.  The
essential estimate is uniform in the cutoff: two local products
approach the product of their invariant conditional expectations at
rate inverse distance.  This identifies all moments of the component
measure before a component is selected.

\begin{definition}[Classes, class systems and bound systems]\label{support:35-2}
At cutoff $K$, let $\mathcal I$ be the Borel sigma-field of events invariant under lattice translations.
\begin{enumerate}
\item[(1)] A class $\mathcal G$ at cutoff $K$ is a nonempty set of lattice laws satisfying the following conditions.
\begin{enumerate}
\item[(Cls1)] The laws are invariant under lattice isometries and reflection positive at every lattice mirror.
\item[(Cls2)] If $\nu\in\mathcal G$ and $h$ is bounded, nonnegative and $\mathcal I$-measurable with $\nu(h)>0$, then $\nu_h=h\nu/\nu(h)$ belongs to $\mathcal G$.
\end{enumerate}
\item[(2)] A class system is a sequence $(K_j,\mathcal G_j)$ with $K_j\to\infty$ and $\mathcal G_j$ a class at cutoff $K_j$.
\item[(3)] A bound system consists of finite positive numbers $C_1(R)$ and $C_n^\infty(d;R)$ for $n\ge2$, $d\in(0,\infty]$ and $R>0$. All are nondecreasing in $R$, and the latter are nonincreasing in $d$.
\item[(4)] The class system obeys this bound system if, in every class,
\begin{align}
 |S_n^T[\rho](X)|
 &\le C_n^\infty(d;R)\|X\|_\sharp
       &&\textup{(Bd1)}\quad(n\ge2,\ X\text{ of type }(d;R)),\label{support:35-3}\\
 |\rho(\mathcal O_K(f))-\nu(\mathcal O_K(f))|
 &\le C_1(R)\|f\|_\sharp
       &&\textup{(Bd2)}\quad(\rho,\nu\in\mathcal G).\label{support:35-4}
\end{align}
In the second line $\operatorname{supp}f$ lies in a ball of radius $R$.
\item[(5)] The same table applies to every class in the system and to every subsequence of the system, by restriction. We write $\underline C=(C_n^\infty)_{n\ge2}$ for its connected-function table.
\end{enumerate}
\end{definition}

For the remainder of this section, fix a class system obeying one
bound system in the sense of \textup{(Bd1)}--\textup{(Bd2)}.

\begin{definition}[Frames and limits in a frame]\label{support:35-5}
A \emph{frame} is a choice $\nu_j\in\mathcal G_j$. It converges if
\begin{equation}\label{support:35-6}
 \sigma(X)=\lim_j\nu_j(A_j(X)),\qquad
 A_j(X):=X^{\nu_j}_{K_j},
\end{equation}
exists for every separated tuple.  A limit in this frame is, more
generally, $\lim_j\rho_j(A_j(X))$ with $\rho_j\in\mathcal G_j$.
The centring law remains $\nu_j$ when the integrating law changes. Every such limit is normalized, real, symmetric and multilinear on separated supports. Its connected functions are the cumulant polynomials of \eqref{support:9-13}.
\end{definition}

Write $d(X)$ for the separation and $R(X)$ for the largest of the
minimal radii of balls containing individual supports.  Set
$d(X)=\infty$ for a tuple of length at most one.  Define
\begin{equation}\label{support:35-9}
 c_n(d;R)=\sum_{\pi\in\operatorname{Part}(n)}
                 \prod_{J\in\pi}b_{|J|},\qquad
 b_1=C_1(R),\quad b_k=C_k^\infty(d;R)\ (k\ge2),\quad c_0=1 .
\end{equation}
The moment bound and the bound used in reflections are
\begin{align}
 C_{\rm mom}(X)&=c_{|X|}(d(X);R(X))\|X\|_\sharp,\label{support:35-10}\\
 \delta(X;\bar R)&=
 \operatorname{dist}_\infty(\operatorname{supp}X,
                  \mathbb R^3\setminus(-\bar R,\bar R)^3),\label{support:35-11}\\
 N(X;\bar R)&=
 c_{2|X|}(\min(d(X),2\delta(X;\bar R));R(X))
                         \|X\|_\sharp^2.\label{support:35-12}
\end{align}
Here the support is contained in the indicated open cube.  Both
bounds are zero if a tuple has a zero entry, and both are one on
the empty tuple.  These conventions avoid evaluating $c_n$ at
$R=0$.  Extend $C_{\rm mom}$ to finite formal sums by the triangle
inequality.  Put
\begin{equation}\label{support:35-13}
 k_1=4\sqrt3,\qquad k_0=4\sqrt3\,C_{\rm dc},
\end{equation}
where $C_{\rm dc}$ is the absolute constant in the lattice
clustering estimate of \cref{support:31-51}.

\begin{proposition}[Uniform conditional bounds and clustering]
\label{support:35-17}
For $\rho,\nu\in\mathcal G$ one has
$|\rho(X^\nu_K)|\le C_{\rm mom}(X)$.  If
$\operatorname{supp}f_i\cup\operatorname{supp}f'_i\subset Q_i$,
the sets $Q_i$ lie in balls of radius $R$ and are mutually at
distance at least $d$, then
\begin{equation}\label{support:35-19}
 |\rho(X^\nu_K-X'^\nu_K)|
 \le c_n(d;R)\sum_i\|f_i-f'_i\|_\sharp
                    \prod_{l\ne i}\max(\|f_l\|_\sharp,\|f'_l\|_\sharp).
\end{equation}
The same bounds hold almost surely for conditional expectations
given $\mathcal I$.

Suppose $\bar R\in a_K\mathbb N$, both tuples are supported in
$(-\bar R,\bar R)^3$, and
$a_K\le2\min(\delta(X;\bar R),\delta(Y;\bar R))$.  For
$v\in a_K\mathbb Z^3$, $|v|\ge k_1\bar R$,
\begin{equation}\label{support:35-24}
 \left|\rho((X\cdot\tau_vY)^\nu_K)
 -\rho\!\left(E_\rho[X^\nu_K\mid\mathcal I]\,
             E_\rho[Y^\nu_K\mid\mathcal I]\right)\right|
 \le k_0\sqrt{N(X;\bar R)N(Y;\bar R)}\,\frac{\bar R}{|v|}.
\end{equation}
All constants are independent of the cutoff and of the two laws.
\end{proposition}

\begin{proof}
The moment--cumulant identity for two centring laws is
\begin{equation}\label{support:35-18}
 \rho(X^\nu_K)=\sum_{\pi\in\operatorname{Part}(n)}
                       \prod_{J\in\pi}\kappa_J,\qquad
 \kappa_{\{i\}}=\rho(\mathcal O_K(f_i))-\nu(\mathcal O_K(f_i)),
 \quad \kappa_J=S_{|J|}^T[\rho](X_J)\ (|J|\ge2).
\end{equation}
Applying \eqref{support:35-3}--\eqref{support:35-4} block by block gives
$c_n(d;R)\|X\|_\sharp$ on any tuple of type $(d;R)$.
Telescope the product defining $X^\nu_K-X'^\nu_K$ one factor
at a time.  Each resulting tuple remains supported in the same
sets $Q_i$; the preceding bound gives \eqref{support:35-19}.

For a bounded Borel function $A$ and an invariant density $h$,
every version $e=E_\rho[A\mid\mathcal I]$ remains a version under
$\rho_h$, since
\[
 \rho_h(kA)=\rho(hkA)/\rho(h)=\rho(hke)/\rho(h)
\]
for every bounded invariant $k$.  If $\rho'(A)\le c$ for every
$\rho'\in\mathcal G$, condition on
$\{E_\rho[A\mid\mathcal I]>c+\varepsilon\}$.  A positive
probability for this event would contradict that bound, by
\Cref{support:35-2}.  Applying this observation to both signs of
$X^\nu_K$ and of its difference from $X'^\nu_K$ proves the
conditional estimates.  A version can be changed to zero on
its exceptional null set to make the absolute bound hold
everywhere.

Put $r=\bar R/a_K$.  A plaquette lies within sup-norm distance
$a_K/2$ of its barycentre.  The margin assumption therefore
makes $A=X^\nu_K$ and $B=Y^\nu_K$ functions of the bonds in
$[-ra_K,ra_K]^3$.  Lattice covariance of $\mathcal O_K$ and
invariance of $\nu$ give
\[
 (\mathsf rX)^\nu_K=\mathsf r(X^\nu_K),\qquad
 (X\cdot\tau_vY)^\nu_K=A\,\tau_vB .
\]
The concatenation is separated because
$|v|_\infty\ge4\bar R$.

The reflection size of the real function $A$ is the maximum,
over coordinate directions $\mu$ and coordinate sign changes
$\mathsf g$, of
\[
 \rho\bigl((\theta_\mu\mathsf gA)
                   (\tau_{2\bar Re_\mu}\mathsf gA)\bigr).
\]
Within either reflected tuple the separation is $d(X)$.
Between the two tuples it is at least $2\delta(X;\bar R)$,
since their $\mu$-coordinates lie respectively below
$\bar R-\delta$ and above $\bar R+\delta$.
The moment bound thus makes this reflection size at most
$N(X;\bar R)$, and the same argument applies to $Y$.
The coordinate and diagonal mirror hypotheses of
\cref{support:31-51} hold for $\rho$.  Its estimate, with
$r a_K=\bar R$, is precisely \eqref{support:35-24}.
\end{proof}

The estimates also give continuity under small translations of
individual entries, uniformly in $K,\rho,\nu$.  Indeed,
$\|\tau_u f-f\|_\sharp\to0$ as $u\to0$, because $\nabla f$
is uniformly continuous.  The analogous fact holds for
$\|\mathsf Rf-f\|_\sharp$ as $\mathsf R\to1$ in $\operatorname{O}(3)$:
on a fixed ball the chain rule gives
$\nabla(\mathsf Rf)(x)=\mathsf R\nabla f(\mathsf R^{-1}x)$.
Enlarge each support by one quarter of
$\min(1,d(X))$ and apply \eqref{support:35-19}; the enlarged
supports have at least half the original separation.
This supplies a common local modulus of continuity for all
the translations used below.

\subsubsection{Limits at spatial infinity}

For a converging frame put
\[
 e_j(X)=E_{\nu_j}[A_j(X)\mid\mathcal I],\qquad
 m_j(X^1,\ldots,X^q)=\nu_j\!\left(\prod_{i=1}^q e_j(X^i)\right).
\]
Choose real versions with $|e_j(X)|\le C_{\rm mom}(X)$.
Versions do not affect $m_j$.  Limits below are iterated: the
shift $v_2$ tends to infinity first, with all later shifts fixed.

\begin{proposition}[Iterated limits]\label{support:35-29}
For any separated tuples $X^1,\ldots,X^q$, not necessarily
separated from one another, $m_j(X^1,\ldots,X^q)$ converges and
\begin{equation}\label{support:35-35}
 \lim_j m_j(X^1,\ldots,X^q)
 =\lim_{|v_q|\to\infty}\cdots\lim_{|v_2|\to\infty}
       \sigma(X^1\cdot\tau_{v_2}X^2\cdots\tau_{v_q}X^q).
\end{equation}
Every intermediate limit is continuous in the remaining shifts
on the open set where the remaining concatenation is separated.

More precisely, put
$X'_p=X^1\cdot\tau_{v_{p+1}}X^{p+1}\cdots\tau_{v_q}X^q$
and denote the intermediate limits by $G^{(p)}$, so that
$G^{(1)}$ is the unmodified correlation.  If
$\bar R'>0$ belongs to $\bigcup_Ka_K\mathbb Z$ and
$\operatorname{supp}X'_p\cup\operatorname{supp}X^p
 \subset(-\bar R',\bar R')^3$, then
\begin{equation}\label{support:35-30}
 |G^{(p-1)}(v_p,\ldots,v_q)-G^{(p)}(v_{p+1},\ldots,v_q)|
 \le k_0\prod_{i=2}^{p-1}C_{\rm mom}(X^i)
       \sqrt{N(X'_p;\bar R')N(X^p;\bar R')}
                         \frac{\bar R'}{|v_p|}
\end{equation}
for $|v_p|\ge k_1\bar R'$.
\end{proposition}

\begin{proof}
Define
\[
 G_j^{(p)}=\nu_j\!\left(\prod_{i=2}^p e_j(X^i)\,A_j(X'_p)\right).
\]
For $p=1$ these converge by the definition of the frame.
For $p=q$ conditional expectation gives
$G_j^{(q)}=m_j(X^1,\ldots,X^q)$.

A real invariant weight $w$ with $|w|\le\bar w$ gives
\begin{equation}\label{support:35-32}
 \nu_j(wF)=c_+\rho_+(F)-c_-\rho_-(F),\qquad
 c_\pm=\nu_j(w_\pm),\quad c_++c_-\le\bar w .
\end{equation}
Terms with $c_\pm=0$ are omitted; otherwise
$\rho_\pm=(\nu_j)_{w_\pm}\in\mathcal G_j$.
The conditional expectations $e_j(X)$ remain versions for
$\rho_\pm$.  Apply \eqref{support:35-24} to these two laws, centred
at $\nu_j$, with $w=\prod_{i=2}^{p-1}e_j(X^i)$.
It bounds $G_j^{(p-1)}-G_j^{(p)}$ by the right side of
\eqref{support:35-30}, whenever the shifts and the cube side
belong to the lattice and $j$ is sufficiently large.
Using \eqref{support:35-19} in the same signed decomposition
also gives a local modulus of continuity of $G_j^{(p)}$
in the remaining shifts, independent of $j$.

Assume $G_j^{(p-1)}$ converges. First fix lattice-compatible
later shifts and a lattice-compatible $\bar R'$. These shifts
belong to every sufficiently fine lattice. For each sufficiently
large lattice-compatible $v_p$, the difference between the
upper and lower limits of $G_j^{(p)}$ is at most twice the
right side of \eqref{support:35-30}, since $G_j^{(p-1)}$
converges there. Letting $v_p=ze_1$, $z\in\mathbb N$, tend
to infinity proves that $G_j^{(p)}$ converges at the chosen
later shifts.

At arbitrary later shifts with $X'_p$ separated, approximate the
shifts by lattice-compatible ones in the same separated neighborhood.
The common local modulus bounds the difference between the values of
$G_j^{(p)}$ at the original and approximating shifts uniformly in $j$.
The Cauchy criterion now gives convergence at the original shifts,
and the same modulus gives continuity of the limit.

For the bound at arbitrary shifts, first take $|v_p|>k_1\bar R'$.
Round all remaining shifts to the lattice at cutoff $K_j$.
For large $j$, the rounded shifts still satisfy the support and
distance hypotheses. Apply \eqref{support:31-55} in the two laws
of \eqref{support:35-32} before bounding the reflection norms.
Each term in the maximum \eqref{support:31-32} for the rounded
$X'_p$ is a reflected moment. The original and rounded reflected
tuples have fixed separated neighborhoods for their entries.
The uniform bound \eqref{support:35-19} therefore compares their
moments with error $o(1)$ in both laws. The original moment is
bounded by $N(X'_p;\bar R')$ by \eqref{support:35-12}, since the
moment bound applies to tests at arbitrary shifts.
Thus the reflection norm squared for the rounded tuple is at
most $N(X'_p;\bar R')+o(1)$. The norm for $X^p$ is bounded by
$N(X^p;\bar R')$ as before. The common local modulus also
allows the rounded shifts in $G_j^{(p-1)},G_j^{(p)}$ to be
replaced by the original shifts with error $o(1)$. Passing to
the limit gives \eqref{support:35-30}. The equality case
$|v_p|=k_1\bar R'$ follows by continuity in $v_p$ with the
later shifts fixed. The right side tends to zero as
$|v_p|\to\infty$, identifying the limit as $G^{(p)}$.
Induction gives \eqref{support:35-35} at $p=q$.
\end{proof}

The constants in \eqref{support:35-30} depend on the later shifts
through $X'_p$ and $\bar R'$.  No joint limit with all shifts
tending to infinity is asserted.

\subsubsection{Construction of the component measure}

\begin{definition}[Canonical decomposition]\label{support:35-34}
A decomposition of the following kind is determined by the
values at spatial infinity.  Its space $\mathsf X$ is compact
and metrizable, its probability measure $\lambda$ has full
support, and each component $\sigma_\xi$ is a normalized,
real, symmetric, multilinear functional on separated tuples.
Require that each evaluation $\xi\mapsto\sigma_\xi(X)$ be
continuous, that distinct points give distinct functionals,
and that
\begin{equation}\label{eq:component-plateau-moments}
 \int_{\mathsf X}\prod_{i=1}^q\sigma_\xi(X^i)\,\lambda(d\xi)
 =\lim_{|v_q|\to\infty}\cdots\lim_{|v_2|\to\infty}
       \sigma(X^1\cdot\tau_{v_2}X^2\cdots\tau_{v_q}X^q).
\end{equation}
For $q=1$ this is the decomposition of $\sigma$ itself.
\end{definition}

\begin{theorem}[Existence of the component measure]\label{support:35-36}
Every converging frame in classes satisfying \Cref{support:35-2} has a canonical decomposition as above, with the following properties.
\begin{enumerate}
\item[(i)] Each evaluation $\xi\mapsto\sigma_\xi(X)$ is continuous, and the components separate the points of $\mathsf X$.
\item[(ii)] The mixed moments of the components are the iterated limits in \eqref{eq:component-plateau-moments}.
\item[(iii)] Each component is a limit in the same frame. More precisely, for every $\xi\in\mathsf X$ there are invariant events $E_j$ with $\nu_j(E_j)>0$ such that $\rho_j=(\nu_j)_{1_{E_j}}$ obeys
\begin{equation}\label{support:35-37}
 \|E_{\rho_j}[A_j(X)\mid\mathcal I]-\sigma_\xi(X)\|_{L^\infty(\rho_j)}\longrightarrow0
 \quad\text{for every separated }X.
\end{equation}
\item[(iv)] Every component satisfies $|\sigma_\xi(X)|\le C_{\rm mom}(X)$ and the difference estimate \eqref{support:35-19}.
\item[(v)] $\#\mathsf X\le\liminf_j\dim L^2(\mathcal V(\mathbb L_{K_j}),\mathcal I,\nu_j)$, with infinity allowed on either side.
\end{enumerate}
\end{theorem}

\begin{proof}
Let $\mathfrak F=\{X_1,X_2,\ldots\}$ be the countable family
with controlled supports constructed in
\cref{app:separated}.  For every nonempty tuple
$X=(f_1,\ldots,f_n)$ and $0<\varepsilon\le1$, choose
$X^{[\varepsilon]}\in\mathfrak F$ whose entries differ from
those of $X$ by at most $\varepsilon$ in $\|\cdot\|_\sharp$.
The construction permits both supports of each entry to lie
in a set contained in a ball of radius $R(X)+1/4$, with
different sets at distance at least
$\frac12\min(1,d(X))$.  Indeed, take the open neighbourhoods
within distance $\frac14\min(1,d(X))$ of the original
supports and use their countable dense test sets.
The estimate \eqref{support:35-19}, also for conditional
expectations, then gives the error
\begin{equation}\label{support:35-38}
 \varepsilon_X(\varepsilon)=
 c_n\!\left(\tfrac12\min(1,d(X));R(X)+\tfrac14\right)
       n\varepsilon\prod_i(\|f_i\|_\sharp+1)
       \longrightarrow0 .
\end{equation}
Choose the approximation using the fixed enumeration, so
that it is independent of the law and of $\xi$.

Consider the compact metrizable product
\[
 \Pi=\prod_{X\in\mathfrak F}[-C_{\rm mom}(X),C_{\rm mom}(X)].
\]
The map $y_j=(e_j(X))_{X\in\mathfrak F}$ is
$\mathcal I$-measurable, and
$\lambda_j=(y_j)_*\nu_j$ is a probability measure on $\Pi$.
By \cref{support:35-29}, the integral against $\lambda_j$ of
every monomial in the coordinates converges.  Coordinate
polynomials contain the constants and separate points;
the real Stone--Weierstrass theorem makes them uniformly
dense in $C(\Pi)$.  Consequently $\lambda_j(\phi)$ converges
for every $\phi\in C(\Pi)$: approximate $\phi$ by a polynomial
and use $|\lambda_j(\phi)|\le\|\phi\|_\infty$.
The limiting positive functional, equal to one on constants,
is represented by a probability measure $\lambda$.
Let $\mathsf X=\operatorname{supp}\lambda$.
The countable base of $\Pi$ shows $\lambda(\mathsf X)=1$,
so this is a compact metrizable space with full support.

For every open $V\subset\Pi$,
$\liminf_j\lambda_j(V)\ge\lambda(V)$.  One may verify this
directly by increasing the continuous functions
$\min(1,s\operatorname{dist}(\,\cdot\,,\Pi\setminus V))$
to $1_V$.  Fix $\xi\in\mathsf X$ and set
\[
 V_s=\{\zeta:|\zeta_{X_k}-\xi_{X_k}|<1/s,\quad 1\le k\le s\}.
\]
Then $\lambda(V_s)>0$, and hence $\lambda_j(V_s)>0$
for all sufficiently large $j$.  Choose increasing
$j_s$ with this property for every $j\ge j_s$.
For $j_s\le j<j_{s+1}$ put $E_j=y_j^{-1}(V_s)$;
for the finitely many earlier indices use the whole space.
These invariant events have positive measure, and
$\rho_j=(\nu_j)_{1_{E_j}}\in\mathcal G_j$.
The unchanged versions of conditional expectation give
\[
 \|E_{\rho_j}[A_j(X_k)\mid\mathcal I]-\xi_{X_k}\|_\infty
 \le1/s\qquad(k\le s,\ j_s\le j<j_{s+1}).
\]

For an arbitrary nonempty $X$, comparison with
$X^{[\varepsilon]}$ now yields
\[
 \limsup_j\|E_{\rho_j}[A_j(X)\mid\mathcal I]
                    -\xi_{X^{[\varepsilon]}}\|_\infty
 \le\varepsilon_X(\varepsilon).
\]
In particular the numbers $\xi_{X^{[\varepsilon]}}$
are Cauchy as $\varepsilon\downarrow0$, with the difference
between two of them at most the sum of the two errors.
Define their limit to be $\sigma_\xi(X)$ and set
$\sigma_\xi(\varnothing)=1$.  Then
\[
 |\sigma_\xi(X)-\xi_{X^{[\varepsilon]}}|
 \le\varepsilon_X(\varepsilon).
\]
This bound and the preceding limsup prove \eqref{support:35-37}.
Integration shows $\rho_j(A_j(X))\to\sigma_\xi(X)$.
It follows that $\sigma_\xi$ is a functional of the asserted
kind and inherits the moment and difference estimates.
For $X\in\mathfrak F$ it equals the coordinate $\xi_X$.

The last approximation is uniform in $\xi$.
Thus every evaluation $\sigma_\xi(X)$ is continuous, as a
uniform limit of coordinates.  Equality of all evaluations
implies equality of all coordinates, so the component map
is injective.

For tuples in $\mathfrak F$, convergence of coordinate
monomials proves \eqref{eq:component-plateau-moments}.
For general tuples replace each $X^i$ by
$(X^i)^{[\varepsilon]}$.  Telescoping the product bounds
both the component integral error and the error in $m_j$
by
\[
 \sum_i\varepsilon_{X^i}(\varepsilon)
       \prod_{l\ne i}
        (C_{\rm mom}(X^l)+\varepsilon_{X^l}(\varepsilon)).
\]
This tends to zero independently of $j$.
Empty tuples contribute a factor one.  Passing first in
$j$ and then in $\varepsilon$ proves the required identity.

Finally, if the displayed liminf dimension is a finite
integer $N$, it equals $N$ on infinitely many indices.
If $\mathsf X$ had $N+1$ distinct points, choose disjoint
open neighbourhoods of them in $\Pi$.  Each has positive
$\lambda$-measure and therefore positive $\lambda_j$-measure
for all sufficiently large $j$.  Their inverse images under
$y_j$ give $N+1$ nonzero orthogonal invariant indicator
functions, contradicting dimension $N$ on those indices.
\end{proof}

\begin{lemma}[Subsequences of frames]\label{support:35-39}
Every frame in classes obeying the moment and cumulant bounds in \Cref{support:35-2} has a convergent subsequence.
\end{lemma}
\begin{proof}
The bounds in \eqref{support:35-10} apply to all orders of centred moments. They give the common separated-support estimates of \Cref{support:10-25} on $\mathbb R^3$. The same subsequence then works for cumulants, by the finite polynomial \eqref{support:9-13}, and the moments of the limit are recovered by \eqref{support:9-21}.
\end{proof}
\subsection{Clustering components and the vacuum space}
\label{sec:component-properties}

We first show that the preceding decomposition is intrinsic to the
limit functional.  Reflection positivity will then force the
symmetries of each component, and the uniform estimate
\eqref{support:35-24} will give clustering for every component.

\begin{proposition}[Uniqueness of the decomposition]\label{support:36-9}
In (a), two decompositions satisfying
\eqref{eq:component-plateau-moments} and the stated compactness,
continuity, injectivity and full-support conditions are related
by a unique measure-preserving homeomorphism which identifies
their component functionals.

In (b), if an isometry $\mathsf r$ preserves $\sigma$, it induces a unique
measure-preserving homeomorphism $h_{\mathsf r}$ of $\mathsf X$ with
\begin{equation}\label{support:36-10}
 \sigma_{h_{\mathsf r}(\xi)}=\sigma_\xi\circ\mathsf r^{-1},
 \qquad h_{\mathsf r\mathsf r'}=h_{\mathsf r}h_{\mathsf r'}.
\end{equation}
In (c), $h_{\mathsf r}$ is the identity if and only if every component is invariant under $\mathsf r$, by the identity in (b) and injectivity. For $s>0$ define
\begin{equation}\label{support:36-4}
 (\operatorname{dil}_s\sigma)(f_1,\ldots,f_n)
   =s^{-n}\sigma(f_1(s\,\cdot),\ldots,f_n(s\,\cdot)).
\end{equation}
The dilated components on the same probability space decompose
$\operatorname{dil}_s\sigma$.
\end{proposition}

\begin{proof}
For any finite collection of tuples, the two component measures
have the same moments in their evaluation coordinates, by
\eqref{eq:component-plateau-moments}, allowing repeated tuples.
Both coordinate images lie in a compact box.  Polynomial
approximation and uniqueness of the representing measure show
that their finite-dimensional distributions agree.

Embed the two component spaces in
$\mathbb R^{\mathfrak S(\mathbb R^3)}$ by all evaluations.
They are compact subsets of this Hausdorff product, and their
embedding maps are homeomorphisms onto their images.  Every
basic neighbourhood of a point in the first image is determined
by finitely many evaluations.  Its inverse image is a nonempty
open set and hence has positive measure.  Equality of the
finite-dimensional distributions makes it meet the second
image.  Since that image is closed, it contains the point.
Interchanging the two spaces proves equality of the images.
Their identification is the required homeomorphism; it is
unique by injectivity of the evaluations.  Polynomials in
these evaluations are uniformly dense on either component
space, so equality of their integrals also proves preservation
of the measures.

For an isometry, the functions
$\sigma_\xi(\mathsf r^{-1}X)$ satisfy the decomposition identity:
conjugating translations changes $v$ by the orthogonal linear
part of $\mathsf r$, and therefore preserves the iterated
limits $|v|\to\infty$.  Apply uniqueness to obtain
$h_{\mathsf r}$.  Composing its defining identities proves the
group law and $h_{\mathsf r^{-1}}=h_{\mathsf r}^{-1}$.
For dilation, $f(s\,\cdot)$ is translated by $v/s$ when $f$
is translated by $v$.  Each iterated limit is thus rescaled
by $s^{-\sum_i|X^i|}$.  This is exactly the product of the
factors in \eqref{support:36-4}.  Continuity and injectivity are
preserved because dilation of test tuples is invertible.
\end{proof}

\begin{proposition}[Reflection invariance of the components]
\label{support:36-29}
Suppose a real symmetric functional $\sigma$ is translation
invariant, is reflection positive at a mirror $\Sigma=(u,c)$,
and has a decomposition as above.  Then $\sigma$ and every
component are invariant under $\theta_\Sigma$.  Each component
is also invariant under
$\mathsf g\theta_\Sigma\mathsf g^{-1}$ whenever the isometry
$\mathsf g$ preserves $\sigma$.

Consequently a Euclidean-invariant $\sigma$ which is reflection
positive at one mirror has Euclidean-invariant components.
\end{proposition}

\begin{proof}
On the positive half-space, positivity for the two formal
vectors $\varnothing,X$ implies
$\sigma(\theta_\Sigma X)=\sigma(X)$: use coefficients $1,i$
and equate imaginary parts.  Move any tuple into this
half-space by a translation parallel to $u$ and use
$\theta_\Sigma\tau_{tu}=\tau_{-tu}\theta_\Sigma$.
Translation invariance proves the same identity for all tuples.

Let $F$ be a finite formal sum of tuples and put
$F_t=\tau_{tu/2}F$.  For large $t$, $F_t$ lies in the positive
half-space.  Translation invariance and the case $q=2$ of
the decomposition identity give
\begin{equation}\label{support:36-30}
 0\le\lim_{t\to\infty}\sigma(\theta_\Sigma\overline{F_t}\cdot F_t)
   =\int_{\mathsf X}
       \overline{\sigma_\xi(\theta_\Sigma F)}
                         \sigma_\xi(F)\,\lambda(d\xi).
\end{equation}
Indeed, term by term the expression before the limit is
$\sigma(\theta_\Sigma\bar F\cdot\tau_{tu}F)$, whose large
translation limit is prescribed by
\eqref{eq:component-plateau-moments}.
Replace $F$ by $F-\theta_\Sigma F$.  The last integral is
then
\[
 -\int_{\mathsf X}
       |\sigma_\xi(F-\theta_\Sigma F)|^2\,\lambda(d\xi).
\]
It must vanish.  The integrand is continuous, and $\lambda$
has full support; hence every component is fixed by
$\theta_\Sigma$.

Equivalently $h_{\theta_\Sigma}$ in \eqref{support:36-10} is the
identity.  Its conjugate by $h_{\mathsf g}$ is also the
identity, proving the assertion for the conjugated mirror.
If $\sigma$ is Euclidean invariant, every mirror is such a
conjugate.  Products of two parallel reflections give
translations.  Every orthogonal map is a product of at most
three reflections: reflect its image of a basis vector back
to that vector, preserving the previously fixed vectors, and
repeat.  Thus every component is Euclidean invariant.
\end{proof}

\begin{corollary}[Symmetries of every component]\label{support:36-31}
For a translation-invariant functional with a canonical decomposition, (i) cubic invariance and reflection positivity at one coordinate and one diagonal mirror imply that every component is translation and cubic invariant. In (ii), full Euclidean invariance and positivity at one mirror imply that every component is Euclidean invariant.
\end{corollary}
\begin{proof}
\Cref{support:9-31} transports the reflection positivity to all mirrors of the relevant kinds. \Cref{support:36-29} then gives invariance of each component under their reflections. Coordinate mirrors generate translations by products of parallel reflections, and their reflections through the origin include the coordinate sign flips. Reflections in $x_i=x_j$ are the transpositions of the coordinate axes, giving (i). In (ii), all reflections are available and generate the Euclidean group.
\end{proof}

\subsubsection{Properties inherited from the frame}

The following conclusions apply to every component, including
points of measure zero in the component space.  Their lattice
approximants are the conditioned laws in \eqref{support:35-37}.

\begin{proposition}[Properties of a limit in a frame]\label{support:36-12}
Let $\sigma'$ be a limit in a frame satisfying \Cref{support:35-2}, with integrating laws $\rho_j$ and centring laws $\nu_j$.
\begin{enumerate}
\item[(a)] It is a symmetric functional, has the moment bound \eqref{support:35-10} and the continuity bound \eqref{support:35-19}, and is continuous under small translations of individual entries.
\item[(b)] It is invariant under translations and cubic isometries.
\item[(c)] It is reflection positive at every coordinate and diagonal mirror.
\item[(d)] Its cumulants of order at least two are the limits of those of $\rho_j$ and obey \eqref{support:35-3}. Its one-point function is the limit of $\rho_j(\mathcal O_{K_j}(f))-\nu_j(\mathcal O_{K_j}(f))$. The frame with integrating and centring laws both $\rho_j$ converges to $\sigma'^{\rm c}$.
\item[(e)] Its one-point function is $c_{\sigma'}\int f$ with $|c_{\sigma'}|\le\bar c$ of \eqref{support:35-14}. It lies in $\mathsf{Fct}(\underline C,\bar c)$, its off-diagonal distributions have order at most $4n+1$ on compact sets, and the coordinate reflection Hilbert spaces and their translation semigroups are defined.
\item[(f)] If a finite vector of tuple moments $\bigl(\rho(A_j(X^i))\bigr)_{i=1}^k$ belongs to a fixed closed subset of $\mathbb R^k$ for all sufficiently large $j$ and every $\rho\in\mathcal G_j$, then $\bigl(\sigma'(X^i)\bigr)_{i=1}^k$ belongs to that set.
\end{enumerate}
\end{proposition}
\begin{proof}
The bounds pass to the limit by \Cref{support:35-17}. The continuity estimate applies to translated entries since $C_c^1$ tests vary continuously in the sharp norm under translations. Translations on the lattice-compatible dense subgroup and cubic symmetries pass directly to the limit. Continuity extends translation invariance to every translation. This proves (a) and (b).

For (c), the laws $\rho_j$ are reflection positive on the lattice in the coordinate and diagonal planes through the origin. The compact supports of finitely many test tuples in an open half-space are a positive distance from its plane. At sufficiently small spacing all plaquettes entering the associated smeared observables lie on the same side. Reflection covariance of the lattice observable and invariance of the centring laws $\nu_j$ give entrywise convergence of the positive lattice reflection matrices to the continuum reflection matrices. Transport to all heights by \Cref{support:9-31}.

Cumulants are polynomials in moments, and centring shifts only the first cumulant. Inverting the cumulants gives the last assertion of (d). Lattice symmetry makes $\rho_j(a_p)-\nu_j(a_p)$ independent of the plaquette. The barycentre Riemann sum therefore gives
\[
\rho_j(\mathcal O_{K_j}(f))-\nu_j(\mathcal O_{K_j}(f))
 =3a_{K_j}^{-4}\bigl(\rho_j(a_p)-\nu_j(a_p)\bigr)r_j(f),
\qquad r_j(f)=\frac{a_{K_j}^3}{3}\sum_p f(x(p))\longrightarrow\int f.
\]
A unit-integral test gives convergence of the coefficient, and the bound of \eqref{support:35-4} gives (e) by taking the infimum over those tests in the unit ball. The distribution and Hilbert-space conclusions follow from Props.~\ref{support:25-55} and~\ref{support:27-22}. For (f), take $\rho=\rho_j$ in the hypothesis and pass to the limit in the closed set.
\end{proof}

\begin{proposition}[One-point functions of frame limits]\label{support:36-13}
\begin{enumerate}
\item[(a)] For a limit $\sigma'$ with integrating laws $\rho_j$,
\begin{equation}\label{support:36-14}
 \sigma'(f)=c_{\sigma'}\int f,\qquad
 c_{\sigma'}=\lim_j3a_{K_j}^{-4}\bigl(\rho_j(a_p)-\nu_j(a_p)\bigr),\qquad |c_{\sigma'}|\le\bar c,
\end{equation}
where
\begin{equation}\label{support:35-14}
 \bar c=C_1(1)\inf\{\|f\|_\sharp:f\in C_c^1(B(0,1)),\ \int f=1\}.
\end{equation}
\item[(b)] In the decomposition of a converging frame, each component has the approximants of \eqref{support:35-37}, $\xi\mapsto c_\xi$ is continuous, $|c_\xi|\le\bar c$, and
\begin{equation}\label{support:36-15}
 \sigma(f)=0,\qquad\int_{\mathsf X}c_\xi\,\lambda(d\xi)=0.
\end{equation}
If $\kappa=\int c_\xi^2\,d\lambda$, then for disjoint $f,f'$,
\begin{equation}\label{support:36-16}
 S_2^\sigma(f,f')=\int_{\mathsf X}S_2^{\sigma_\xi}(f,f')\,\lambda(d\xi)+\kappa\int f\int f'.
\end{equation}
\end{enumerate}
\end{proposition}
\begin{proof}
The Riemann-sum argument in \Cref{support:36-12} proves (a). For (b), \Cref{support:35-36} supplies the approximants. For a unit-integral test $f_0$, $c_\xi=\sigma_\xi(f_0)$ is a continuous evaluation. Every one-point moment of the frame is zero by its centring, and the decomposition identity gives \eqref{support:36-15}. Integrate $S_2^{\sigma_\xi}(f,f')=\sigma_\xi(f,f')-c_\xi^2\int f\int f'$ to obtain \eqref{support:36-16}.
\end{proof}

\begin{proposition}[Uniform clustering of the components]\label{support:36-17}
For the decomposition of a converging frame, let
$\bar R>0$ belong to $\bigcup_Ka_K\mathbb Z$ and contain
$\operatorname{supp}X\cup\operatorname{supp}Y$ in its open
cube.  Then, for every $\xi\in\mathsf X$ and
$|v|\ge k_1\bar R$,
\begin{equation}\label{support:36-18}
 |\sigma_\xi(X\cdot\tau_vY)-\sigma_\xi(X)\sigma_\xi(Y)|
 \le k_0\sqrt{N(X;\bar R)N(Y;\bar R)}\,\frac{\bar R}{|v|}.
\end{equation}
The bound is uniform over all components.
\end{proposition}

\begin{proof}
Use the conditioned approximants $\rho_j$ for the given
component.  For a lattice-compatible $v$, all the hypotheses
of \eqref{support:35-24} hold eventually.  Its first term converges
to $\sigma_\xi(X\cdot\tau_vY)$.
By \eqref{support:35-37}, the conditional expectations of
$A_j(X)$ and $A_j(Y)$ converge in their respective
$L^\infty(\rho_j)$ norms to the constants
$\sigma_\xi(X)$ and $\sigma_\xi(Y)$.  If their errors are
$\varepsilon_j(X)$ and $\varepsilon_j(Y)$, the error in
the integral of their product is at most
\[
 \varepsilon_j(X)(|\sigma_\xi(Y)|+\varepsilon_j(Y))
       +|\sigma_\xi(X)|\varepsilon_j(Y)\longrightarrow0 .
\]
Thus \eqref{support:36-18} follows on the dense subgroup.
Approximate any remaining $v$ by such vectors with
$|v_j|\ge k_1\bar R$; the continuity under translations
already proved gives the assertion.  The right side
contains no component-dependent quantity.
\end{proof}

\subsubsection{Zero energy and the component space}
\begin{proposition}[Zero energy]\label{support:36-20}
Let $\sigma$ be a translation-invariant functional with the bounds of \S\ref{sec:reflection-hilbert}, reflection positive at the mirror with unit normal $u$. In its reflection Hilbert space put $P_0=1_{\{0\}}(H)$.
\begin{enumerate}
\item[(a)] For positive-half-space vectors $F,F'$,
\begin{equation}\label{support:36-21}
 \lim_{t\to\infty}\sigma(\theta\bar F\cdot\tau_{tu}F')=\langle[F],P_0[F']\rangle .
\end{equation}
\item[(b)] The equality $\ker H=\mathbb C\Omega$ is equivalent to
\begin{equation}\label{support:36-22}
 \lim_{t\to\infty}\sigma(\theta\bar F\cdot\tau_{tu}F')=\overline{\sigma(F)}\sigma(F')
\end{equation}
for all such $F,F'$.
\item[(c)] For each component of a converging frame obeying \Cref{support:35-2}, the coordinate reflection Hamiltonians have one-dimensional kernels.
\end{enumerate}
\end{proposition}
\begin{proof}
By the spectral theorem $T_t=e^{-tH}\to P_0$ strongly. This is (a). Since $\langle\Omega,[F]\rangle=\sigma(F)$, the identity in (b) is exactly the equality of matrix elements of $P_0$ and $|\Omega\rangle\langle\Omega|$ on a dense subspace. Each component satisfies the limit in (b) by the bound \eqref{support:36-18} and reflection invariance, which proves (c).

\end{proof}

\begin{theorem}[The vacuum space of the mixture]\label{support:36-25}
Let $\sigma$ satisfy the hypotheses of \Cref{support:36-20}
and have a canonical decomposition as in \Cref{support:35-34}.
Assume its components
are translation invariant and invariant under the reflection
$\theta$ defining $\mathcal H(\sigma,u)$, and that
\[
 \sup_\xi|\sigma_\xi(X\cdot\tau_{tu}Y)
                     -\sigma_\xi(X)\sigma_\xi(Y)|\longrightarrow0
 \quad(t\to\infty)
\]
for positive-half-space tuples $X,Y$.  Then
\begin{equation}\label{support:36-26}
 \langle P_0[F],P_0[F']\rangle
    =\int_{\mathsf X}\overline{\sigma_\xi(F)}
                             \sigma_\xi(F')\,\lambda(d\xi).
\end{equation}
The assignment
$P_0[F]\mapsto(\xi\mapsto\sigma_\xi(F))$ extends to a unitary
map $\ker H\to L^2(\mathsf X,\lambda)$, carrying $\Omega$
to $1$.  The number of components equals $\dim\ker H$
when either is finite; if $\mathsf X$ is infinite, the
kernel is a separable infinite-dimensional Hilbert space.

These conclusions apply to every coordinate reflection
Hilbert space of the limit of a converging frame.
\end{theorem}

\begin{proof}
For positive-half-space tuples $X,Y$, apply
\eqref{eq:component-plateau-moments} with first tuple $\theta X$,
second tuple $Y$ and shift $tu$. Reflection invariance of each
component gives
\[
 \lim_{t\to\infty}\sigma(\theta X\cdot\tau_{tu}Y)
      =\int_{\mathsf X}\sigma_\xi(X)\sigma_\xi(Y)\,d\lambda .
\]
Combine this with \eqref{support:36-21} and extend sesquilinearly
to obtain \eqref{support:36-26}.
It makes the proposed map well defined and isometric.
The vectors $P_0[F]$ are dense in $\ker H$, so it extends
to an isometry with closed range, and sends $\Omega$ to $1$.

To prove surjectivity, put
$\psi_F(\xi)=\sigma_\xi(F)$ and take the uniform closure of
their span in $C(\mathsf X)$, initially for positive-half-space
$F$.  It contains constants and is closed under conjugation.
The uniform clustering assumption gives
\[
 \psi_{X\cdot\tau_{tu}Y}\longrightarrow\psi_X\psi_Y
          \quad\text{uniformly on }\mathsf X .
\]
Thus the closure is an algebra.  It separates points:
any tuple distinguishing two components may be translated
into the positive half-space without changing its values.
Stone--Weierstrass makes the closure $C(\mathsf X)$.
Since continuous functions are dense in $L^2(\lambda)$,
the isometry has dense as well as closed range.

Full support implies that disjoint nonempty open sets give
nonzero orthogonal indicator functions in $L^2(\lambda)$.
Any $n$ distinct points have $n$ such neighbourhoods.
Hence a finite component space has exactly that dimension,
and an infinite component space has infinite dimension.
Separability follows from that of the reflection Hilbert
space.  For frame limits, \cref{support:36-12,support:36-17} supply
both invariance and uniform clustering required above.
\end{proof}


\begin{corollary}[Vacuum space of a frame limit]\label{support:36-27}
\Cref{support:36-25} applies to every canonical decomposition of a converging frame satisfying the stated moment and class bounds.
\end{corollary}
\begin{proof}
For the canonical decomposition constructed by \Cref{support:35-36}, reflection invariance and translation invariance follow from \Cref{support:36-12}, and the uniform limit follows from \Cref{support:36-17}. Every other canonical decomposition agrees through the unique measure-preserving homeomorphism of \Cref{support:36-9}.
\end{proof}
\begin{remark}[A two-component constant field]\label{ex:constant-mixture}
Let $c_1\ne c_2$ and $0<p<1$. The functional given by the mixture of the two deterministic constant fields with values $c_1,c_2$ has a two-dimensional reflection Hilbert space and $H=0$. Indeed, the reflection form is the Euclidean inner product on the evaluation vectors at $c_1,c_2$ with weights $p,1-p$. The vacuum and a one-point test of integral one already span both coordinates. Every translation acts trivially. The connected two-point function is $p(1-p)(c_1-c_2)^2\int f\int h$, whereas each component's connected two-point function vanishes.
\end{remark}
\section{Scalar reconstruction}\label{app:scalar-reconstruction}

\subsection{Point insertions and energy bounds for reconstruction}
\label{sec:reconstruction-energy}

The reconstruction uses more than reflection positivity and bounds
on separated moments.  Its additional input is an $L^1$ estimate
for one insertion between two positive-time vectors.  We first
construct point forms with polynomial energy bounds, then use
them to define smeared fields on a common domain.

Let $\sigma$ be a real, normalized, symmetric multilinear
functional on separated tuples.  Throughout this section assume:
\begin{enumerate}[label=(Rec\arabic*)]
\item $\sigma$ is invariant under all Euclidean isometries.
\item It is reflection positive at $x_3=0$.  Its reflection
Hilbert space $\mathcal H$ is separable, and its translations
are represented by $e^{-tH}$ and $\mathcal U(v)$, with $H\ge0$,
the stated strong continuity, and joint spectral measure
$\mathcal P$:
\begin{equation}\label{support:75-12}
 e^{-tH}[F]=[\tau_{te_3}F],\qquad
 \mathcal U(v)[F]=[\tau_{(v,0)}F],\qquad
 e^{i\tau H}\mathcal U(v)
   =\int e^{i(\tau E+v\cdot q)}\,d\mathcal P(E,q).
\end{equation}
\item $\mathcal P$ is carried by $\{(E,q):E\ge|q|\}$.
\item For some $C_\wp,\wp\ge0$ and $0<c_0\le1/2$,
put $\mathrm{Sl}_\eta=\{|x_3|<c_0\eta\}$.  For
$0<\eta\le1$, $h\in C_c^1(\mathrm{Sl}_\eta)$ and
$F,F'\in\mathbb C\mathfrak S(\{x_3>\eta\})$, assume
\begin{equation}\label{support:75-14}
 |\sigma(\theta\bar F\cdot(h)\cdot F')|
 \le C_\wp\eta^{-\wp}\|h\|_1
       \|[\tau_{-\eta e_3}F]\|
       \|[\tau_{-\eta e_3}F']\|.
\end{equation}
\end{enumerate}
We refer to these four conditions as (Rec).
Vacuum uniqueness, $\ker H=\mathbb C\Omega$, will be a separate
assumption when needed.  The numerical triple in (Rec4) is
not part of the definition of the resulting field.

Write $\vec z=(z_1,z_2)$ and define
\begin{equation}\label{support:75-5}
 d_\circ(z)=\operatorname{Re}z_3-|\operatorname{Im}\vec z|,
 \qquad \mathrm{Tb}_\circ=\{d_\circ>0\},\qquad
 \mathcal T(z)=\int e^{-z_3E+i\vec z\cdot q}\,d\mathcal P .
\end{equation}
The last integral is defined on the closed tube.
The cone bound gives
$|e^{-z_3E+i\vec z\cdot q}|\le e^{-d_\circ(z)E}$.
Consequently $\mathcal T$ consists of contractions, is strongly
continuous on the closed tube and norm holomorphic in its
interior, and satisfies
\[
 \mathcal T(z)\mathcal T(w)=\mathcal T(z+w),\qquad
 \mathcal T(z)^*=\mathcal T(z^*),\qquad
 z^*=(-\bar z_1,-\bar z_2,\bar z_3).
\]
For norm holomorphy, expand the exponential around a point
with positive margin; the factor $e^{-d_\circ(z)E}$ bounds
every coefficient and the quadratic remainder uniformly
on a smaller ball.  Dominated convergence against scalar
spectral measures proves strong continuity at the boundary.
In the interior, $E^ke^{-d_\circ(z)E}$ is bounded for every
$k$, so the range of $\mathcal T(z)$ lies in
\[
 \mathcal D_\infty=\bigcap_{k\ge1}\operatorname{dom}H^k .
\]
The vacuum is fixed by $\mathcal T(z)$ and its adjoint.
For Minkowski coordinates $x=(t,\vec x)$ put
\begin{equation}\label{support:75-8}
 \iota(t,\vec x)=(\vec x,-it).
\end{equation}
This identifies the forward tube
$\mathbb R^{1+2}+i\{y:y^0>|\vec y|\}$ with
$\mathrm{Tb}_\circ$, and
$\mathcal T(\iota x)=e^{itH}\mathcal U(\vec x)$ is unitary.

\subsubsection{Densities of bounded insertions}

\begin{lemma}\label{support:75-22}
For $0<\eta\le1$ there is a unique complex linear map
$h\mapsto B_\eta^{\rm in}(h)$ from $L^1(\mathrm{Sl}_\eta)$
to bounded operators such that
\begin{align}
 \|B_\eta^{\rm in}(h)\|&\le
                   C_\wp\eta^{-\wp}\|h\|_1,\label{support:75-23}\\
 B_\eta^{\rm in}(h)[F]&=
       [(\tau_{\eta e_3}h)\cdot\tau_{2\eta e_3}F]
       \quad(h\in C_c^1(\mathrm{Sl}_\eta)).\label{support:75-24}
\end{align}
For real $h$, its adjoint is $B_\eta^{\rm in}(\theta h)$.
Each scalar matrix element has an $L^\infty$ density of
norm at most $C_\wp\eta^{-\wp}\|\psi\|\|\psi'\|$.
\end{lemma}

\begin{proof}
Apply \eqref{support:75-14} after translating both formal vectors
by $\eta e_3$.  The bounded sesquilinear form
\begin{equation}\label{support:75-25}
 \langle[F],B_\eta^{\rm in}(h)[F']\rangle
   =\sigma(\theta\overline{\tau_{\eta e_3}F}
                         \cdot(h)\cdot\tau_{\eta e_3}F')
\end{equation}
vanishes on null classes and therefore defines a bounded
operator by Riesz representation.  Translate the whole tuple
by $\eta e_3$ to obtain \eqref{support:75-24}.
The support inequalities $c_0\le1/2$ make that tuple
separated.  Real linearity and the norm bound extend this
map from real $C_c^1$ functions to real $L^1$ functions.

For fixed vectors the real and imaginary parts of the
matrix element are bounded real functionals on $L^1$, hence
have a complex density $b$.  The same constant bounds
$|b|$: apply the norm estimate to normalized indicators
of small balls and use Lebesgue differentiation.  Thus
$\|b\|_\infty\le C_\wp\eta^{-\wp}\|\psi\|\|\psi'\|$.
Complex linear extension now preserves this constant,
since $|\int hb|\le\|h\|_1\|b\|_\infty$.
Reflection invariance, symmetry and reality in
\eqref{support:75-25} give the adjoint relation on $C_c^1$,
then on $L^1$ by approximation.  The class identity,
density of classes, and $L^1$ continuity prove uniqueness.
\end{proof}

\begin{proposition}[Point insertions]\label{support:75-27}
There is a unique family of bounded operators $B(s,s')$,
$s,s'>0$, with the shift law
\begin{equation}\label{support:75-28}
 T_tB(s,s')T_{t'}=B(s+t,s'+t')\qquad(t,t'\ge0)
\end{equation}
and the density identity
\begin{equation}\label{support:75-29}
 \langle\psi,B_\eta^{\rm in}(h)\psi'\rangle
 =\int h(x)\langle\psi,\mathcal U(\vec x)
        B(\eta+x_3,\eta-x_3)\mathcal U(-\vec x)\psi'\rangle\,dx .
\end{equation}
It satisfies
\begin{equation}\label{support:75-30}
 \|B(s,s')\|\le C_\wp\min(s,s',1)^{-\wp},
 \qquad B(s,s')^*=B(s',s).
\end{equation}
It is norm continuous and commutes with spatial orthogonal
transformations acting on $\mathcal H$.  The operator density
\[
 b_\eta(x)=\mathcal U(\vec x)
                  B(\eta+x_3,\eta-x_3)\mathcal U(-\vec x)
\]
is norm continuous on $\mathrm{Sl}_\eta$,
has norm at most $C_\wp\eta^{-\wp}$ there, and has
real-analytic scalar matrix elements.
\end{proposition}

\begin{proof}
For $0<\eta'<\eta\le1$, $h\in L^1(\mathrm{Sl}_{\eta'})$
and real $w$ with $|w_3|<c_0(\eta-\eta')$, translation
covariance of \eqref{support:75-25} gives
\begin{equation}\label{support:75-32}
 B_\eta^{\rm in}(\tau_wh)
 =\mathcal T((\eta-\eta')e_3+w)\,
      B_{\eta'}^{\rm in}(h)\,
        \mathcal T((\eta-\eta')e_3-w).
\end{equation}
The two semigroup arguments have positive third coordinate,
and the translated support stays in $\mathrm{Sl}_\eta$.
It suffices to prove the identity on $C_c^1$ tests and
dense classes; $L^1$ continuity then proves the assertion
as written.

Fix $\psi,\psi'$ and let $b$ be their scalar density at
margin $\eta$.  Take nonnegative continuous mollifiers
$j_\varepsilon$ of integral one, supported in
$B(0,\varepsilon)$, with $\varepsilon<c_0\eta'$.
The functions
\[
 b_\varepsilon(w)=
 \langle\psi,\mathcal T((\eta-\eta')e_3+w)
 B_{\eta'}^{\rm in}(j_\varepsilon)
 \mathcal T((\eta-\eta')e_3-w)\psi'\rangle
\]
are holomorphic and bounded by
$C_\wp\eta'^{-\wp}\|\psi\|\|\psi'\|$ on
\[
 \{w\in\mathbb C^3:
       |\operatorname{Re}w_3|+|\operatorname{Im}\vec w|
                                      <\eta-\eta'\}.
\]
On the smaller real slab
$|w_3|<c_0(\eta-\eta')$, \eqref{support:75-32}
identifies them with $\int j_\varepsilon(x-w)b(x)\,dx$.
Montel's theorem gives a locally uniformly convergent
subsequence with holomorphic limit.  Mollification also
converges to $b$ in $L^1$ on compact real subsets, so this
limit agrees with $b$ almost everywhere.
As $\eta'\downarrow0$ the smaller real slabs exhaust
$\mathrm{Sl}_\eta$.  Their continuous representatives
agree on overlaps and give one real-analytic representative
of $b$ on the whole slab.

Linearity in $\psi'$ and antilinearity in $\psi$ hold
everywhere for these representatives: they hold almost
everywhere for the original densities, and both sides are
continuous.  The original $L^\infty$ bound likewise holds
at every point.  Riesz representation thus defines bounded
operators $b_\eta(x)$ with
$\|b_\eta(x)\|\le C_\wp\eta^{-\wp}$.
This constructs the operator density from scalar densities.

Integrate \eqref{support:75-32} against arbitrary tests.
Uniqueness of the continuous scalar densities gives
\begin{equation}\label{support:75-33}
 b_\eta(x+w)=\mathcal T((\eta-\eta')e_3+w)\,
                 b_{\eta'}(x)\,
                     \mathcal T((\eta-\eta')e_3-w).
\end{equation}
For $0<\eta\le\min(s,s',1)$ define
\begin{equation}\label{support:75-34}
 B(s,s')=T_{s-\eta}\,b_\eta(0)\,T_{s'-\eta}.
\end{equation}
The case $x=w=0$ of \eqref{support:75-33} proves independence
of $\eta$.  It also proves the shift law; choosing
$\eta=\min(s,s',1)$ gives the norm bound.
For $x\in\mathrm{Sl}_\eta$, choose $\eta'>0$ with
$|x_3|<c_0(\eta-\eta')$ and apply \eqref{support:75-33} at
$(0,x)$.  Factoring spatial translations from the semigroup
gives the stated formula for $b_\eta(x)$ and hence
\eqref{support:75-29}.

In \eqref{support:75-33} choose a fixed smaller margin locally.
Its semigroup factors have positive margins and are norm
holomorphic, proving norm continuity of $b_\eta$.
The same reasoning in \eqref{support:75-34} proves norm
continuity of $B$.  The insertion adjoint identity makes
$b_\eta(x)^*=b_\eta(\theta x)$, so $b_\eta(0)$ is
self-adjoint.  This proves the adjoint relation for $B$.
Spatial rotations in \eqref{support:75-25} give covariance of
$b_\eta$ and invariance of $b_\eta(0)$; they commute with
$T_t$, proving the remaining assertion.

For uniqueness, any bounded-operator family with the
shift law has continuous scalar densities in
\eqref{support:75-29}: factor it through $B(\eta/2,\eta/2)$.
The density identity then determines its value at $x=0$,
hence $B(\eta,\eta)$.  The shift law determines all other
values.  The same argument shows that changing the
admissible triple $C_\wp,\wp,c_0$ does not change the
operators, by first restricting insertions to the common
smaller slab.
\end{proof}

\subsubsection{Resolvent regularization of the point form}

Put $R_H=(1+H)^{-1}$.  For real $l>\wp$ define
\begin{equation}\label{support:75-40}
 A_l=\Gamma(l)^{-2}\int_0^\infty\!\!\int_0^\infty
            (ss')^{l-1}e^{-s-s'}B(s,s')\,ds\,ds'.
\end{equation}

\begin{lemma}\label{support:75-41}
The integral is norm integrable and defines a bounded
self-adjoint operator with
\begin{equation}\label{support:75-42}
 \|A_l\|\le C_\wp\left(1+\frac{2\Gamma(l-\wp)}{\Gamma(l)}\right).
\end{equation}
For $a>0$ and $s,s'>0$,
\begin{equation}\label{support:75-43}
 R_H^aA_lR_H^a=A_{l+a},\qquad
 T_sA_lT_{s'}=R_H^lB(s,s')R_H^l.
\end{equation}
Consequently
\begin{equation}\label{support:75-49}
 q_0(\psi,\psi')=
       \langle(1+H)^l\psi,A_l(1+H)^l\psi'\rangle
\end{equation}
defines a Hermitian form on
$\bigcup_{l>\wp}\operatorname{dom}H^l$, independent of the
admissible exponent used for a given pair of vectors.
It satisfies the energy bound
\[
 |q_0(\psi,\psi')|\le
       \|A_l\|\,\|(1+H)^l\psi\|\,\|(1+H)^l\psi'\|
\]
on $\operatorname{dom}H^l$, and
\begin{equation}\label{support:75-45}
 q_0(T_s\psi,T_{s'}\psi')=\langle\psi,B(s,s')\psi'\rangle.
\end{equation}
\end{lemma}

\begin{proof}
By \eqref{support:75-30},
$\|B(s,s')\|\le C_\wp(1+s^{-\wp}+s'^{-\wp})$.
The two Gamma integrals therefore converge for $l>\wp$
and give \eqref{support:75-42}.  Norm continuity makes the
operator integrand measurable.  Self-adjointness follows
by interchanging $s,s'$ and using $B(s,s')^*=B(s',s)$.

Functional calculus and the scalar Gamma integral give
\[
 R_H^a=\Gamma(a)^{-1}\int_0^\infty t^{a-1}e^{-t}T_t\,dt
\]
weakly for $a>0$.  Insert this on both sides of $A_l$
and use the shift law.  Fubini is justified by the
integrable bound just proved and contractivity.
Set $s=r+t$ and $s'=r'+t'$.  The two inner integrals are
\[
 \int_0^s t^{a-1}(s-t)^{l-1}\,dt
      =s^{a+l-1}\frac{\Gamma(a)\Gamma(l)}{\Gamma(a+l)}.
\]
This proves the first identity in \eqref{support:75-43}.
For the second, shift $B(r,r')$ by $(s,s')$, then write
$B(r+s,r'+s')=T_rB(s,s')T_{r'}$ and use the same
resolvent integral.

If $l'>l>\wp$ and both vectors lie in
$\operatorname{dom}H^{l'}$, the first identity and
$R_H^{l'-l}(1+H)^{l'}=(1+H)^l$ prove equality of the
two definitions of $q_0$.  The domains are nested, so a
common admissible exponent exists for every pair in the
stated union.  Hermiticity and the bound follow from those
of $A_l$.  Finally $(1+H)^lT_s$ is bounded for $s>0$.
Apply the second identity at $(s/2,s'/2)$ and cancel
the resolvent factors against these bounded regularized
powers.  The shift law gives \eqref{support:75-45}.
\end{proof}

\subsubsection{Recovering an insertion bound after centering}

The hypotheses (Rec) are also available for a functional with
nonzero one-point function.

\begin{lemma}\label{support:75-50}
Suppose $\sigma$ satisfies (Rec1)--(Rec2), its spectrum obeys
$E\ge|q_1|$, and $\sigma(f)=c\int f$.  If its centered
functional satisfies
\begin{equation}\label{support:75-51}
 |\sigma^{\rm c}(\theta\bar F\cdot(h)\cdot F')|
 \le C_\wp\eta^{-\wp}\|h\|_1
 \|[\tau_{-\eta e_3}F]\|_{\sigma^{\rm c}}
 \|[\tau_{-\eta e_3}F']\|_{\sigma^{\rm c}},
\end{equation}
with the supports in (Rec4), then $\sigma$ satisfies (Rec)
with constants $(C_\wp+|c|,\wp,c_0)$.
\end{lemma}

\begin{proof}
Spatial rotations conjugate the joint spectral measure, so
$E\ge|n\cdot q|$ for every unit spatial vector $n$.
A countable dense set of such vectors gives $E\ge|q|$.
The shift identities of \S\ref{sec:reflection-hilbert}
give
$\sigma=\sigma^{\rm c}\circ\mathsf{sh}_c$ and an isometry
$[F]_\sigma\mapsto[\mathsf{sh}_cF]_{\sigma^{\rm c}}$.
The shift commutes with translations, conjugation and
reflection, preserves supports, and is multiplicative.
Writing $F_c=\mathsf{sh}_cF$, therefore,
\begin{equation}\label{support:75-52}
 \sigma(\theta\bar F\cdot(h)\cdot F')
 =\sigma^{\rm c}(\theta\bar F_c\cdot(h)\cdot F'_c)
       +c\int h\,\sigma^{\rm c}(\theta\bar F_c\cdot F'_c).
\end{equation}
The first term has the bound \eqref{support:75-51} with the
$\sigma$ norms.  Cauchy--Schwarz and contractivity of
$T_\eta$ bound the second by
$|c|\|h\|_1\|[\tau_{-\eta e_3}F]\|
\|[\tau_{-\eta e_3}F']\|$.
Since $\eta\le1$ and $\wp\ge0$, their sum has the required
bound.  Centering preserves the other stated symmetries
and reflection positivity by the same shift identities.
\end{proof}

\subsection{Analytic continuation of the Euclidean kernels}
\label{sec:reconstruction-kernels}

The point-insertion bound gives a product estimate without
a factorial.  Rotation of the ordering direction then covers
all distinct Euclidean points, including configurations
with equal third coordinates.

Call $(h_1,\ldots,h_n)$ \emph{$\eta$-ordered along $u$}
at heights $t_1,\ldots,t_n$ if $u\in S^2$, $0<\eta\le1$,
$t_{k+1}-t_k\ge2\eta$, and
\[
 \operatorname{supp}h_k
       \subset\{|u\cdot x-t_k|<c_0\eta\}.
\]
Such tuples are separated by at least $\eta$.
For $u=e_3$, set $g_k=\tau_{-t_ke_3}h_k$ and
$s_k=t_{k+1}-t_k-2\eta\ge0$.

\begin{proposition}\label{support:75-61}
For an $\eta$-ordered tuple along $e_3$,
\begin{equation}\label{support:75-62}
 \sigma(h_1,\ldots,h_n)=
 \langle\Omega,B_\eta^{\rm in}(g_1)T_{s_1}
       B_\eta^{\rm in}(g_2)\cdots
       T_{s_{n-1}}B_\eta^{\rm in}(g_n)\Omega\rangle.
\end{equation}
If $t_1\ge\eta$, then
\begin{equation}\label{support:75-63}
 [(h_1,\ldots,h_n)]=T_{t_1-\eta}
 B_\eta^{\rm in}(g_1)T_{s_1}\cdots
 B_\eta^{\rm in}(g_n)\Omega,\qquad
 \|[(h_1,\ldots,h_n)]\|
       \le\prod_k C_\wp\eta^{-\wp}\|h_k\|_1.
\end{equation}
\end{proposition}

\begin{proof}
Starting at the right, \eqref{support:75-24} proves by downward
induction that the product beginning with
$B_\eta^{\rm in}(g_k)$ represents
$[\tau_{(\eta-t_k)e_3}(h_k,\ldots,h_n)]$.
Indeed the intervening translation has height
$s_k+\eta-t_{k+1}=-t_k-\eta$, and the insertion
adds $2\eta$ to the remaining tuple.
All these classes have supports in
$\{x_3>(1-c_0)\eta\}$.
Taking the vacuum component and using translation
invariance proves \eqref{support:75-62}.
If $t_1\ge\eta$, translate the class forward by
$t_1-\eta$ to obtain \eqref{support:75-63};
contractivity and \eqref{support:75-23} give its norm bound.
\end{proof}

For later use, the scalar density formula can be iterated:
for $g_k\in L^1(\mathrm{Sl}_\eta)$ and bounded $A_0,\ldots,A_n$,
\begin{equation}\label{support:75-60}
 \begin{split}
 &\langle\psi,A_0B_\eta^{\rm in}(g_1)A_1\cdots
                 B_\eta^{\rm in}(g_n)A_n\psi'\rangle\\
 &\qquad=\int\prod_k g_k(y_k)
       \langle\psi,A_0b_\eta(y_1)A_1\cdots
                          b_\eta(y_n)A_n\psi'\rangle\,dy .
 \end{split}
\end{equation}
The second factor is continuous and bounded by
$(C_\wp\eta^{-\wp})^n\|\psi\|\|\psi'\|\prod_j\|A_j\|$.
Thus the integrand is absolutely integrable.  Apply
\eqref{support:75-29} to the first insertion, then to the
remaining ones; Fubini proves the formula.

\subsubsection{Ordered complex configurations}

For $\zeta_1,\ldots,\zeta_{n-1}\in\mathrm{Tb}_\circ$ define
\begin{equation}\label{support:75-66}
 \mathbf W_n(\zeta_1,\ldots,\zeta_{n-1})
 =\langle\Omega,B(\eta,\eta)\mathcal T(\zeta_1-2\eta e_3)
          B(\eta,\eta)\cdots
          \mathcal T(\zeta_{n-1}-2\eta e_3)
          B(\eta,\eta)\Omega\rangle,
\end{equation}
where $0<\eta\le\min(1,\frac12\min_k d_\circ(\zeta_k))$.
For $n=1$ this is the real number
$\mathbf W_1=\langle\Omega,B(1,1)\Omega\rangle$.
On
\[
 \mathrm{Tb}^{(n)}_\circ
 =\{z:z_1\in\mathrm{Tb}_\circ,\
               z_{k+1}-z_k\in\mathrm{Tb}_\circ\}
\]
put
\begin{align}
 \eta(z)&=\min\left(1,d_\circ(z_1),
                     \tfrac12\min_{k<n}d_\circ(z_{k+1}-z_k)\right),
                                                        \label{support:75-67}\\
 \Psi_n(z)&=\mathcal T(z_1-\eta e_3)B(\eta,\eta)
       \mathcal T(z_2-z_1-2\eta e_3)B(\eta,\eta)\cdots
       \mathcal T(z_n-z_{n-1}-2\eta e_3)B(\eta,\eta)\Omega,
                                                        \label{support:75-68}
\end{align}
using any $0<\eta\le\eta(z)$; the inner minimum is omitted
when $n=1$.

These definitions are independent of $\eta$.
To compare $\eta'<\eta$, substitute
\[
 B(\eta,\eta)=T_{\eta-\eta'}B(\eta',\eta')T_{\eta-\eta'}
\]
at every insertion.  The semigroup law absorbs the adjacent
factors into the differences, and $T_t\Omega=\Omega$
removes the final factor.  For the scalar product it removes
the initial factor as well.

\begin{proposition}\label{support:75-71}
The function $\mathbf W_n$ is holomorphic and $\Psi_n$ is
holomorphic in norm.  They satisfy
\begin{align}
 |\mathbf W_n(\zeta)|&\le(C_\wp\eta^{-\wp})^n
       &&\text{if }d_\circ(\zeta_k)\ge2\eta,\ 0<\eta\le1,
                                                        \label{support:75-72}\\
 \|\Psi_n(z)\|&\le(C_\wp\eta(z)^{-\wp})^n,                 \label{support:75-73}\\
 \mathcal T(w)\Psi_n(z)&=\Psi_n(z_1+w,\ldots,z_n+w),\notag\\
 \langle\Omega,\Psi_n(z)\rangle
       &=\mathbf W_n(z_2-z_1,\ldots,z_n-z_{n-1})           \label{support:75-74}
\end{align}
for $w\in\overline{\mathrm{Tb}_\circ}$.
For ordered tuples,
\begin{equation}\label{support:75-76}
 \sigma(h_1,\ldots,h_n)
 =\int\mathbf W_n(x_2-x_1,\ldots,x_n-x_{n-1})
                              \prod_k h_k(x_k)\,dx.
\end{equation}
If also $t_1\ge\eta$, then
\begin{equation}\label{support:75-77}
 \langle\psi,[(h_1,\ldots,h_n)]\rangle
       =\int\langle\psi,\Psi_n(x)\rangle\prod_k h_k(x_k)\,dx.
\end{equation}
\end{proposition}

\begin{proof}
Contractivity and \eqref{support:75-30} give the bounds.
Locally choose a fixed $\eta$ strictly smaller than every
allowed margin.  All semigroup factors are then norm
holomorphic; the independence of $\eta$ proves holomorphy
of the defined functions.  The semigroup law and the
vacuum identities give \eqref{support:75-74}.

The density can be written with half the margin:
\begin{equation}\label{support:75-75}
 b_\eta(y)=\mathcal T(y+\tfrac\eta2e_3)
        B(\tfrac\eta2,\tfrac\eta2)
                  \mathcal T(-y+\tfrac\eta2e_3).
\end{equation}
Both arguments have positive margin, since
$|y_3|<c_0\eta\le\eta/2$.  This follows directly from
the shift law and the formula for $b_\eta$.
Insert \eqref{support:75-75} in \eqref{support:75-60}, applied to
the products \eqref{support:75-62}--\eqref{support:75-63}.
With $x_k=y_k+t_ke_3$, consecutive semigroup factors
combine into
$\mathcal T(x_{k+1}-x_k-\eta e_3)$.
At the right the factor fixes $\Omega$.
At the left it either disappears in the vacuum component
or becomes $\mathcal T(x_1-\eta e_3/2)$.
These are precisely \eqref{support:75-66} and
\eqref{support:75-68} with margin $\eta/2$, proving the
two integral formulas.  The supports are compact and
stay inside the respective open tubes, so their
continuous integrands are integrable.
\end{proof}

\subsubsection{Changing the ordering direction}

For $u\in S^2$ put
\[
 d_\circ^u(z)=\operatorname{Re}(u\cdot z)
       -|\operatorname{Im}(z-(u\cdot z)u)|,\qquad
 \mathrm{Tb}^{u,\pi}_n
 =\{z:d_\circ^u(z_{\pi(k+1)}-z_{\pi(k)})>0,\ k<n\}.
\]
If $R\in\operatorname{SO}(3)$ satisfies $Re_3=u$, then
$d_\circ^u(z)=d_\circ(R^{-1}z)$.
Each tube is open and convex; it contains $\operatorname{Re}z$
whenever it contains $z$.  For $n=1$ it is all of $\mathbb C^3$.
Define on it
\begin{equation}\label{support:75-69}
 \mathbf S_n^{R,\pi}(z)
 =\mathbf W_n\bigl(R^{-1}(z_{\pi(2)}-z_{\pi(1)}),\ldots,
                   R^{-1}(z_{\pi(n)}-z_{\pi(n-1)})\bigr).
\end{equation}

Any $n$ distinct real points admit a separating direction.
More quantitatively, if their mutual distances are at least
$\delta>0$, there is $u\in S^2$ with
\begin{equation}\label{support:75-53}
 |u\cdot(x_i-x_j)|\ge\delta/n^2\quad(i\ne j).
\end{equation}
Indeed the spherical zone
$|u\cdot (x_i-x_j)/|x_i-x_j||<n^{-2}$ has area $4\pi/n^2$.
The union over the $\binom n2$ pairs has area less than
$2\pi$, so cannot cover $S^2$.
Ordering the projections shows that the real parts of the
tubes cover $(\mathbb R^3)^n_{\ne}$.

\begin{proposition}\label{support:75-78}
There is a unique continuous kernel $\mathsf s_n$ on
$(\mathbb R^3)^n_{\ne}$ such that
\begin{equation}\label{support:75-79}
 \sigma(f_1,\ldots,f_n)=\int\mathsf s_n(x)\prod_i f_i(x_i)\,dx
\end{equation}
for all separated tuples.  It is real analytic, symmetric
and invariant under Euclidean isometries.
It extends uniquely to a holomorphic function $\mathbf S_n$
on
$\mathrm{Tb}^{\rm all}_n=\bigcup_{u,\pi}\mathrm{Tb}^{u,\pi}_n$,
whose restriction to each tube is \eqref{support:75-69},
independently of the choice of $R$.  This extension is
symmetric and obeys
\begin{equation}\label{support:75-85}
 |\mathbf S_n(z)|\le(C_\wp\eta^{-\wp})^n
 \quad\text{if }d_\circ^u(z_{\pi(k+1)}-z_{\pi(k)})
                  \ge2\eta,\quad0<\eta\le1.
\end{equation}
In particular,
\begin{equation}\label{support:75-80}
 \mathsf s_n(x)=\mathbf W_n(x_2-x_1,\ldots,x_n-x_{n-1})
                  \quad(x_{13}<\cdots<x_{n3}),
 \qquad \mathsf s_1=\mathbf W_1 .
\end{equation}
\end{proposition}

\begin{proof}
For a tuple ordered along $u$, rotate it by $R^{-1}$,
permute its entries, and apply \eqref{support:75-76}.
Euclidean invariance and symmetry give the kernel formula
with $\mathbf S_n^{R,\pi}$.
If a real configuration belongs to two ordering regions,
choose sufficiently small balls about its entries so that
every product of tests in those balls is ordered both ways,
with one common positive margin.  The two kernel formulas
have the same value on every such product.  Continuous
functions are determined by these product tests, for example
by taking products of approximate identities.
Thus the kernels agree on each real overlap.

The real ordering regions cover all distinct points by
\eqref{support:75-53}.  They therefore define one real-analytic
function $\mathsf s_n$.  To obtain \eqref{support:75-79} for
arbitrary separated supports, let their least distance be
$d>0$ and choose
$\eta=\min(d/(4n^2),1)$.
Partition each test into finitely many pieces supported
in open Euclidean balls of radius $c_0\eta/2$.  Choose a point of the
original support in every nonzero piece.
Equation \eqref{support:75-53} orders these chosen points with
gaps at least $4\eta$; each piece is supported in a
$c_0\eta$ slab about its chosen point.
Each resulting product tuple is therefore ordered.
Apply the formula to each product and sum by multilinearity.
For $n=1$ the same partition uses $\eta=1$.
This proves the formula on every separated tuple.
Products of arbitrarily small supported tests also prove
uniqueness.  Applying uniqueness to permuted or isometrically
transported tests gives the claimed real symmetries.

A nonempty complex overlap of two tubes is convex, and
contains the real part of each of its points.
It consequently has a nonempty open real part.
The two holomorphic pieces agree there by the preceding
argument.  The identity theorem on the connected overlap
makes them equal throughout it.  They therefore glue to
one holomorphic function on the union.
The same identity theorem on each tube proves uniqueness
and independence of $R$; it also extends the permutation
symmetry.  The bound is \eqref{support:75-72} after rotation.
All constructions depend on $\sigma$, not on the chosen
admissible numerical triple.
\end{proof}

\subsubsection{Growth on tests flat at the diagonals}

The extension to the full space of tests used for Schwinger
functions requires a bound near every diagonal, not only
for products with fixed separated supports.

\begin{proposition}\label{support:75-87}
Put $k_{\rm E}=\lceil\wp\rceil+4$ and
$C_{\rm E}=C_\wp2^\wp e^{2\wp}3^{k_{\rm E}/2}\pi^2$.
Then $|\mathsf s_1|\le C_\wp$, and for $n\ge2$,
\begin{equation}\label{support:75-88}
 |\mathsf s_n(x)|\le
 C_\wp^n(2n^2)^{\wp n}
            \min(1,\min_{i\ne j}|x_i-x_j|)^{-\wp n}.
\end{equation}
For every $\varphi\in\mathcal S_{\ne}(\mathbb R^{3n})$,
\begin{equation}\label{support:75-89}
 \int|\mathsf s_n(x)\varphi(x)|\,dx
       \le2C_{\rm E}^n(n!)^{2\wp}|\varphi|_{k_{\rm E}n}.
\end{equation}
Consequently the Schwinger functions on this space are
\begin{equation}\label{support:75-92}
 \mathsf S_n^\sigma(\varphi)=\int\mathsf s_n(x)\varphi(x)\,dx.
\end{equation}
\end{proposition}

\begin{proof}
For a configuration with least Euclidean distance $\rho$,
choose the direction in \eqref{support:75-53} and use
\eqref{support:75-85} with $\eta=\min(\rho/(2n^2),1)$.
This proves \eqref{support:75-88}; the one-point bound follows
from $B(1,1)$.

Let $k=k_{\rm E}n$, $N=\lceil\wp n\rceil$.
Then $N\le k$ and $k\ge4n$.
If $\rho\le1$ and $N\ge1$, move a closest pair to its
midpoint, obtaining a diagonal point $x^*$ with
$|x-x^*|\le1$ and $|x-x^*|_1\le\sqrt3\rho$.
Flatness at $x^*$ and Taylor's remainder, with the
weighted seminorm of \S\ref{sec:distribution-extension},
give
\[
 |\varphi(x)|\le
 \frac{(\sqrt3\rho)^N}{N!}\,3^{k/2}
       |\varphi|_k(1+|x|^2)^{-k/2}
 \le2\,3^{k/2}\rho^N
       |\varphi|_k(1+|x|^2)^{-k/2}.
\]
The factor $3^{k/2}$ accounts for moving at most unit
distance along the Taylor segment.
Since $N\ge\wp n$, this absorbs the singular factor.
When $\rho>1$ or $\wp=0$, the unmodified seminorm bound
suffices.  Finally
\[
 n^n\le e^nn!,\qquad
 \int_{\mathbb R^{3n}}(1+|x|^2)^{-k/2}\,dx
 \le\prod_{i=1}^n\int_{\mathbb R^3}(1+|x_i|^2)^{-2}\,dx_i
 =\pi^{2n}.
\]
Combining these estimates proves \eqref{support:75-89}.
The integral functional agrees with $\sigma$ on separated
smooth tensors by \eqref{support:75-79}.
The tensor-density and uniqueness result of
\S\ref{sec:distribution-extension} identifies it
with the Schwinger function.
\end{proof}

\subsection{Smeared fields and their boundary distributions}
\label{sec:reconstruction-fields}

We continue under (Rec).  Minkowski space has coordinates
$x=(t,\vec x)$, product $x\cdot y=t s-\vec x\cdot\vec y$,
and forward cone $V_+=\{t>|\vec x|\}$.
Write $U(x)=\mathcal T(\iota x)$ and fix an integer $l>\wp$.
For $f\in\mathcal S(\mathbb R^{1+2})$ put
\begin{equation}\label{support:76-5}
 c_N(f)=\|A_l\|\sum_{k=0}^N\binom Nk\|\partial_t^kf\|_1.
\end{equation}
The Schwartz seminorms used above give
$\|\partial^\alpha F\|_1\le\pi^{2n}|F|_{|\alpha|+4n}$
on $\mathbb R^{3n}$, by comparison with
$\prod_i(1+|x_i|^2)^{-2}$.  In particular
\begin{equation}\label{support:76-15}
 c_N(f)\le\|A_l\|2^N\pi^2|f|_{N+4}.
\end{equation}

For integers $j\ge0$, we use the graph seminorms $\|(1+H)^j\psi\|$ on
$\mathcal D_\infty$.
The spectral projections $P_N=\mathcal P(\{E\le N\})$
have ranges in $\mathcal D_\infty$ and converge strongly
to the identity.  Bounded joint spectral multipliers
preserve every $\operatorname{dom}H^j$ and commute there
with $(1+H)^j$.
We will also use the following domain criterion:
\begin{equation}\label{eq:iv-domain-criterion}
 |\langle(1+H)^j\psi,\chi\rangle|\le C\|\psi\|
 \quad(\psi\in\mathcal D_\infty)
 \ \Longrightarrow\
 \chi\in\operatorname{dom}H^j,\quad
 \|(1+H)^j\chi\|\le C.
\end{equation}
Indeed, insert $\psi=(1+H)^jP_N\chi$ to bound
$\int_{E\le N}(1+E)^{2j}\,dm_\chi$ by $C^2$,
then let $N\to\infty$.

\subsubsection{Construction on a common domain}

\begin{proposition}\label{support:76-25}
There is a unique complex linear map
$f\mapsto\phi(f)$ from $\mathcal S(\mathbb R^{1+2})$
to linear operators $\mathcal D_\infty\to\mathcal D_\infty$
with
\begin{equation}\label{support:76-26}
 \langle\psi,\phi(f)\psi'\rangle
   =\int f(x)\,q_0(U(-x)\psi,U(-x)\psi')\,dx.
\end{equation}
It is independent of $l$ and of the numerical triple in
(Rec).  For every integer $j\ge0$,
\begin{equation}\label{support:76-28}
 \|(1+H)^j\phi(f)\psi\|
     \le c_{l+j}(f)\|(1+H)^{2l+j}\psi\|.
\end{equation}
On $\mathcal D_\infty$ it satisfies
\begin{align}
 \langle\psi,\phi(f)\psi'\rangle
    &=\langle\phi(\bar f)\psi,\psi'\rangle,\nonumber\\
 U(a)\phi(f)U(-a)&=\phi(f(\cdot-a)),\qquad
 [H,\phi(f)]=i\phi(\partial_tf).                         \label{support:76-27}
\end{align}
\end{proposition}

\begin{proof}
The integrand in \eqref{support:76-26} equals
\[
 f(x)\langle U(-x)(1+H)^l\psi,
                    A_lU(-x)(1+H)^l\psi'\rangle.
\]
Its second factor is bounded and continuous.
Denote the resulting sesquilinear form by $q_f$.
It obeys
\begin{equation}\label{support:76-31}
 |q_f(\psi,\psi')|\le
 \|f\|_1\|A_l\|\|(1+H)^l\psi\|\|(1+H)^l\psi'\|.
\end{equation}
To turn this form into an operator, move energy powers
from the first vector to the second.
The map $t\mapsto e^{-itH}\xi$ is norm differentiable
for $\xi\in\operatorname{dom}H$, with derivative
$-ie^{-itH}H\xi$, by dominated convergence in the spectral
calculus.  Differentiating both vectors in the preceding
matrix element, then integrating by parts in $t$, gives
\[
 q_f(H\psi,\psi')-q_f(\psi,H\psi')
                   =i q_{\partial_tf}(\psi,\psi').
\]
The boundary terms vanish because $f$ is Schwartz and
the matrix element and its derivative are bounded.
Iteration yields
\begin{equation}\label{support:76-33}
 q_f((1+H)^N\chi,\psi')
 =\sum_{k=0}^N\binom Nk i^k
                 q_{\partial_t^kf}(\chi,(1+H)^{N-k}\psi').
\end{equation}
Choose $\chi=R_H^l\psi$ and $N=l+j$.
Using \eqref{support:76-31} and monotonicity of graph norms,
\begin{equation}\label{support:76-34}
 |q_f((1+H)^j\psi,\psi')|
 \le c_{l+j}(f)\|\psi\|\|(1+H)^{2l+j}\psi'\|.
\end{equation}
At $j=0$, Riesz representation gives a unique vector
$\phi(f)\psi'$.  For every $j$, the domain criterion
\eqref{eq:iv-domain-criterion} gives its membership in
$\operatorname{dom}H^j$ and the bound \eqref{support:76-28}.
Thus it belongs to $\mathcal D_\infty$.

Linearity, uniqueness and independence of the auxiliary
choices follow from the defining form.
Hermiticity follows from that of $q_0$.
Changing variables in \eqref{support:76-26} proves translation
covariance.  The identity obtained by integration by parts,
tested against dense vectors in $\mathcal D_\infty$,
proves the commutator formula.
\end{proof}

\subsubsection{Regularized products and their continuity}

\begin{lemma}\label{support:76-35}

Let $w,w'\in\mathrm{Tb}_\circ$ and
$0<s\le\min(1,d_\circ(w),d_\circ(w'))$.
Then
\begin{equation}\label{support:76-36}
 \mathcal T(w)\phi(f)\mathcal T(w')
 =\int f(x)\mathcal T(w-se_3+\iota x)
             B(s,s)\mathcal T(w'-se_3-\iota x)\,dx
\end{equation}
in matrix elements on $\mathcal H$.
The integrand after removal of $f$ is strongly continuous
and bounded by $C_\wp s^{-\wp}$.
For $w_1,\ldots,w_n\in\mathrm{Tb}_\circ$, set
$z_k(x)=\sum_{r\le k}w_r+\iota x_k$.  Then
\begin{equation}\label{support:76-37}
 \begin{split}
 &\langle\psi,\mathcal T(w_1)\phi(f_1)\cdots
                    \mathcal T(w_n)\phi(f_n)\Omega\rangle\\
 &\qquad=\int\prod_k f_k(x_k)\langle\psi,\Psi_n(z(x))\rangle\,dx .
 \end{split}
\end{equation}
\end{lemma}

\begin{proof}
The positive margins put
$\mathcal T(w)^*\psi,\mathcal T(w')\psi'$ in
$\mathcal D_\infty$.
Apply \eqref{support:76-26} and split $T_s$ from each
semigroup.  Equation \eqref{support:75-45} changes the form
between these two factors into $B(s,s)$; commuting the
remaining spectral multipliers gives \eqref{support:76-36}.
The margin is unchanged by addition of $\iota x$, so all
remaining factors are contractions.  Strong continuity
and the point bound give the stated properties.

For the product choose
$s=\min(1,d_\circ(w_1),\frac12\min_{k\ge2}d_\circ(w_k))$.
The inner minimum is omitted when $n=1$.
Split each intermediate semigroup into equal halves, and
append $T_s$ at the right, where it fixes $\Omega$.
This writes the product as a product of bounded
regularized fields, each covered by \eqref{support:76-36}.
Successively insert those formulas.
The uniform point bounds and $\prod_k|f_k|\in L^1$
justify Fubini.  Consecutive semigroup arguments combine
to $z_{k+1}(x)-z_k(x)-2se_3$; the first is
$z_1(x)-se_3$, and the final spatial-time translation
fixes $\Omega$.  The resulting product is
$\Psi_n(z(x))$ with margin $s$, proving
\eqref{support:76-37}.
\end{proof}

\begin{lemma}\label{support:76-38}
Suppose $c_N(f_k^\epsilon-f_k^0)\to0$ for every $N$,
and $C_k^\epsilon$ are contractive joint spectral
multipliers converging strongly to $C_k^0$.
Then
\[
 C_1^\epsilon\phi(f_1^\epsilon)\cdots
 C_n^\epsilon\phi(f_n^\epsilon)\Omega
 \longrightarrow
 C_1^0\phi(f_1^0)\cdots C_n^0\phi(f_n^0)\Omega
\]
in every graph norm of $H$.
\end{lemma}

\begin{proof}
Work from right to left.
If $\chi^\epsilon\to\chi^0$ in all graph norms, then
\[
 \begin{split}
 &\|(1+H)^j(\phi(f^\epsilon)\chi^\epsilon
                                      -\phi(f^0)\chi^0)\|\\
 &\quad\le c_{l+j}(f^\epsilon-f^0)
                       \|(1+H)^{2l+j}\chi^\epsilon\|
       +c_{l+j}(f^0)
                       \|(1+H)^{2l+j}(\chi^\epsilon-\chi^0)\|
       \longrightarrow0.
 \end{split}
\]
The spectral multipliers commute with the graph powers.
Their strong convergence and uniform norm bound preserve
the just-proved convergence.  Induction starts at
$\chi^\epsilon=\Omega$.
\end{proof}

\subsubsection{Tempered boundary values}

Here we give the boundary argument explicitly; pointwise
polynomial growth near the tube boundary alone does not
justify passing a limit through a test integral.

\begin{lemma}\label{support:76-23}
Let $\mathbf F$ be holomorphic on an open set containing
\[
 z(x;\eta)=(\iota x_j+\eta h_je_3)_{j=1}^n,
 \qquad x\in K,\quad0<\eta\le\eta_0\le1,
\]
where $K$ is compact, and suppose that for real $b,p\ge0$,
$|\mathbf F(z(x;\eta))|\le b\eta^{-p}$.
For $g\in C_c^\infty$ supported in $K$, put
$I_g(\eta)=\int g(x)\mathbf F(z(x;\eta))\,dx$ and
$D_h=\sum_jh_j\partial_{t_j}$.
If the integer $N\ge p+1/2$, then
\begin{equation}\label{eq:iv-boundary-estimate}
 \sup_{0<\eta\le\eta_0}|I_g(\eta)|
 \le 2^N\eta_0^{-(N-1/2)}
                          b\sum_{r=0}^N\|D_h^rg\|_1.
\end{equation}
The limit at $\eta=0$ exists.
\end{lemma}

\begin{proof}
On each compact interval of positive heights,
holomorphy permits differentiation under the integral.
The Cauchy--Riemann equations and integration by parts give
$I_g'=I_{-iD_hg}$.
Set $J_r=I_{(-iD_h)^rg}$.
For $0\le r\le N$, these functions obey
$J_r'=J_{r+1}$ when $r<N$, and
$|J_r(\eta)|\le b\|D_h^rg\|_1\eta^{-(N-1/2)}$.
Integrate successively from $\eta_0$ toward $\eta$.
Every integration of $\eta^{-a}$ with $a\ge3/2$
lowers the singular exponent by one at a cost at most
$2$; after $N-1$ integrations the exponent is $1/2$.
Its final integral is bounded by $2$.
The values at $\eta_0$ contribute the same bound with
the corresponding $\|D_h^rg\|_1$, giving
\eqref{eq:iv-boundary-estimate}.
Apply that estimate once more to $-iD_hg$.
Then $I_g'$ is bounded uniformly down to zero,
so $I_g$ is Lipschitz there and has a limit.
\end{proof}

\begin{proposition}\label{prop:iv-wightman-distributions}
There is a unique tempered distribution
$\mathcal W_n^\sigma$ on $\mathbb R^{3n}$ satisfying
\begin{equation}\label{support:76-29}
 \mathcal W_n^\sigma(f_1\otimes\cdots\otimes f_n)
 =\langle\Omega,\phi(f_1)\cdots\phi(f_n)\Omega\rangle .
\end{equation}
For arbitrary $\alpha_1,\ldots,\alpha_{n-1}\in V_+$,
\begin{equation}\label{support:76-30}
 \begin{split}
 &\mathcal W_n^\sigma(f_1\otimes\cdots\otimes f_n)\\
 &\quad=\lim_{\epsilon\downarrow0}\int\prod_k f_k(x_k)
 \mathbf W_n\bigl(\iota(x_2-x_1+i\epsilon\alpha_1),\ldots,
             \iota(x_n-x_{n-1}+i\epsilon\alpha_{n-1})\bigr)\,dx .
 \end{split}
\end{equation}
For $n=1$, $\mathcal W_1^\sigma(f)=\mathbf W_1\int f$.
\end{proposition}

\begin{proof}
First take the vacuum component of \eqref{support:76-37},
with $w_1=e_3$ and $w_{k+1}=\iota(i\epsilon\alpha_k)$.
It identifies the integral on the right of
\eqref{support:76-30} with
\[
 \langle\Omega,\phi(f_1)\mathcal T(w_2)
                  \phi(f_2)\cdots\mathcal T(w_n)\phi(f_n)\Omega\rangle.
\]
Its integrand is bounded by
$(C_\wp\eta^{-\wp})^n\prod_k|f_k(x_k)|$, with
$2\eta=\min(2,\epsilon\min_k(\alpha_k^0-|\vec\alpha_k|))$.
The spectral contractions tend strongly to the identity.
\Cref{support:76-38} therefore proves the boundary formula
on tensors of Schwartz functions.

To obtain a distribution on general tests, apply
\Cref{support:76-23} to
$\mathbf F(z)=\mathbf W_n(z_2-z_1,\ldots,z_n-z_{n-1})$
with $h_j=2j$, $\eta_0=1$, $b=C_\wp^n$,
$p=\wp n$, and $N=\lceil\wp n\rceil+1$.
Every consecutive margin is $2\eta$, so the hypothesis
follows from \eqref{support:75-72}.
The resulting limit functional on $C_c^\infty$
has bound
\[
 |\mathcal L(F)|\le
 2^N(N+1)C_\wp^n(n(n+1))^N\pi^{2n}|F|_{N+4n}.
\]
Here $\sum_jh_j=n(n+1)$ bounds the coefficients in
$D_h^r$, and the $L^1$ derivative estimate gives the
Schwartz norm.

For a smooth cutoff $\chi$ equal to one near the origin,
$\chi(\cdot/R)F\to F$ in every Schwartz seminorm.
Thus the bound extends $\mathcal L$ continuously to
$\mathcal S(\mathbb R^{3n})$.
It agrees with \eqref{support:76-29} first on compactly
supported tensors by the boundary identity, then on
Schwartz tensors by \Cref{support:76-38}.
For uniqueness, finite sums of compactly supported smooth
tensors are dense in each fixed Schwartz seminorm:
cut off the test, convolve with a product mollifier,
and approximate the resulting compact integral by
Riemann sums of translates of that product mollifier.
All derivatives up to the fixed order converge uniformly
on a common compact set.  A tempered functional vanishing
on those tensors must therefore vanish.
\end{proof}

\subsubsection{Cyclicity}\label{sec:ordered-density}

Let $\mathcal D_0$ be the span of $\Omega$ and all
$\phi(f_1)\cdots\phi(f_n)\Omega$.
It lies in $\mathcal D_\infty$, is invariant under the
fields and translations, and is dense.

To prove the last assertion, let $\psi\perp\mathcal D_0$.
Translation covariance first gives
\[
 \langle\psi,U(a_1)\phi(f_1)\cdots
                       U(a_n)\phi(f_n)\Omega\rangle=0 .
\]
Replace the translations successively by
$\mathcal T(w_j)$, $w_j\in\mathrm{Tb}_\circ$.
For one replacement, move all preceding regularized
factors to the first vector using Hermiticity.
The first positive-margin semigroup puts that vector
in $\mathcal D_\infty$, so every move is legitimate.
The resulting expression is $\langle\chi,U(a)\chi'\rangle$.
Its vanishing for all $a$ says that the Fourier transform
of the finite joint spectral measure $m_{\chi,\chi'}$
vanishes.  Uniqueness of Fourier transforms of finite
measures gives $m_{\chi,\chi'}=0$, and hence
$\langle\chi,\mathcal T(w)\chi'\rangle=0$.
Induction proves vanishing for every fully regularized
product.

By \eqref{support:76-37} and product approximate identities,
$\langle\psi,\Psi_n(z)\rangle=0$ throughout
$\mathrm{Tb}^{(n)}_\circ$.
Equation \eqref{support:75-77} then makes $\psi$ orthogonal to
classes of time-ordered positive-time tuples.
We must still approximate arbitrary positive-time tuples;
their supports need not be time ordered.

Fix such a tuple $X=(f_1,\ldots,f_n)$.
Partition its entries with nonnegative functions
$\chi_{h,k}$ supported in cubes of side $2h$ about $hk$,
summing to one.  Call a product of pieces good when all
their center heights differ by at least $4h/c_0$.
For small $h$, every good product, after permutation,
is ordered with $\eta=2h/c_0$ and first height at least
$\eta$.  Its class is orthogonal to $\psi$.

Let $\xi_h$ be the sum of the remaining classes.
On the compact set consisting of the original positive
supports and their reflected negative supports,
the continuous kernel $\mathsf s_{2n}$ is bounded by
some $M_X$.  The kernel formula gives
\begin{equation}\label{support:76-43}
 \|\xi_h\|^2\le M_X
 \left(\int\prod_i|f_i(x_i)|
       \sum_{\mathbf k\ {\rm bad}}\prod_i\chi_{h,k_i}(x_i)\,dx\right)^2 .
\end{equation}
The sum is at most one and vanishes unless some two
heights differ by less than $(4/c_0+2)h$.
Its integral is $O_X(h)$: integrate one height first
over an interval of that length, bound that test by
its supremum norm on its fixed spatial support,
then integrate the remaining tests in $L^1$.
Thus $\|\xi_h\|\to0$ and
$\langle\psi,[X]\rangle=\langle\psi,\xi_h\rangle=0$.
All tuple classes span a dense subspace of the reflection
Hilbert space, so $\psi=0$ and $\mathcal D_0$ is dense.

The translation representation already has joint spectrum
$E\ge|q|$.  In the convention
$U(x)=\int e^{ip\cdot x}\,d\mathsf E(p)$ its momentum
measure is the image of $\mathcal P$ under
$(E,q)\mapsto(E,-q)$, and is carried by
$\overline{V_+}$.  If $\ker H=\mathbb C\Omega$, the only
translation-invariant vectors are multiples of $\Omega$:
invariance under $e^{itH}$ forces spectral support at
$E=0$.  No vacuum-uniqueness assertion is used in the
preceding construction.

\subsection{Lorentz covariance and locality}
\label{sec:reconstruction-covariance}

Let $\mathcal L=\operatorname{SO}^+(1,2)$.
Every element is a product of spatial rotations and the boosts
\[
 B(s)(t,x_1,x_2)
   =(t\cosh s+x_1\sinh s,\ x_1\cosh s+t\sinh s,\ x_2).
\]
For completeness, $\Lambda e_0$ is a future unit timelike
vector.  A spatial rotation and a boost send $e_0$ to it.
Their inverse followed by $\Lambda$ fixes $e_0$ and hence
acts by an element of $\operatorname{SO}(2)$ on its orthogonal
spatial plane.  This proves the factorization.
These matrices have determinant one and preserve $V_+$.

\begin{proposition}\label{support:76-44}
There is a unique strongly continuous unitary representation
$U(\Lambda)$ of $\mathcal L$ fixing $\Omega$ and satisfying
\begin{equation}\label{support:76-47}
 U(\Lambda)\mathcal D_0=\mathcal D_0,\qquad
 U(\Lambda)\phi(f)U(\Lambda)^{-1}\psi
             =\phi(f\circ\Lambda^{-1})\psi
                   \quad(\psi\in\mathcal D_0).
\end{equation}
Together with translations it defines the strongly continuous
representation
\begin{equation}\label{support:76-48}
 U(a,\Lambda)=U(a)U(\Lambda),\qquad
 U(a,\Lambda)U(a',\Lambda')
                  =U(a+\Lambda a',\Lambda\Lambda').
\end{equation}
Its spatial-rotation subgroup agrees with the representation
on the reflection Hilbert space.
\end{proposition}

\begin{proof}
We first prove invariance of the analytic functions:
\begin{equation}\label{support:76-45}
 \mathbf W_n(\iota\Lambda\zeta_1,\ldots,\iota\Lambda\zeta_{n-1})
          =\mathbf W_n(\iota\zeta_1,\ldots,\iota\zeta_{n-1})
 \quad(\zeta_k\in\mathbb R^{1+2}+iV_+).
\end{equation}
The complex linear map $\widehat\Lambda=\iota\Lambda\iota^{-1}$
preserves $\mathrm{Tb}_\circ$.  Both sides, viewed as
functions of $\iota\zeta_k$, are holomorphic on its
connected product.  It suffices to compare them at real
differences $\xi_k$ of positive third coordinate.
Choose $x_1=0$, $x_{k+1}=x_k+\xi_k$.
Then $\mathbf W_n(\xi)=\mathsf s_n(x)$.
Spatial-rotation invariance follows from that of the
Euclidean kernel.

For a boost, use the complex rotation
\begin{equation}\label{support:76-49}
 r_z(y_1,y_2,y_3)
       =(y_1\cos z+y_3\sin z,\ y_2,\ y_3\cos z-y_1\sin z).
\end{equation}
It satisfies $r_{is}=\widehat{B(s)}$.
Fix real positive-time differences $\xi_k$.
The set of $z$ for which every $r_z\xi_k$ lies in
$\mathrm{Tb}_\circ$ is open and contains the imaginary
axis, since real boosts preserve the forward tube.
A sufficiently small neighborhood of the compact segment
from $0$ to $is$ lies in this set and is connected.
For small real $z$, $r_z$ is an ordinary Euclidean
rotation preserving the ordering of these points.
Thus $\mathbf W_n(r_z\xi_1,\ldots,r_z\xi_{n-1})$
is constant there by Euclidean invariance.
The identity theorem continues this equality to $z=is$.
The identity theorem in all difference variables, then
the rotation--boost factorization, proves
\eqref{support:76-45}.

In \eqref{support:76-30} change variables $x_j=\Lambda y_j$
and use \eqref{support:76-45}.  The approach vectors become
$\Lambda^{-1}\alpha_k\in V_+$, which are allowed in
that boundary formula.  Hence
\begin{equation}\label{support:76-46}
 \mathcal W_n^\sigma(F\circ(\Lambda^{-1})^{\times n})
                         =\mathcal W_n^\sigma(F).
\end{equation}
This first holds on tensors; continuity under invertible
linear changes of Schwartz variables and tensor density
extend it to all tests.

On a polynomial vector define
\[
 U(\Lambda)\phi(f_1)\cdots\phi(f_n)\Omega
       =\phi(f_1\circ\Lambda^{-1})\cdots
                       \phi(f_n\circ\Lambda^{-1})\Omega .
\]
Inner products of two such vectors are Wightman values:
move the first string across using Hermiticity, reversing
its order and conjugating its tests.
Equation \eqref{support:76-46} preserves every such inner
product.  The rule therefore respects all linear
relations, is isometric, and has inverse given by
$\Lambda^{-1}$.  Density of $\mathcal D_0$ gives a unitary
extension; the group law and covariance hold on its
spanning vectors and hence as stated.
This also proves uniqueness.

As $\Lambda\to1$, the time derivatives of
$f\circ\Lambda^{-1}-f$ tend to zero in $L^1$.
The chain rule expresses them as finite sums of derivatives
of $f$ after the linear change; dominated convergence,
using Schwartz decay uniformly for $\Lambda$ near $1$,
proves the assertion.
\Cref{support:76-38} gives strong continuity first on
polynomial vectors, then on all of $\mathcal H$ by
unitarity and density.
On polynomial vectors the conjugation
$U(\Lambda)U(a)U(\Lambda)^{-1}$ translates every test
by $\Lambda a$.  This proves \eqref{support:76-48}
and continuity of the combined representation.

Finally, the original spatial rotations commute with
$H$ and $B(s,s')$, hence with $A_l$, and preserve $q_0$.
They conjugate $U(x)$ to $U((1\oplus R)x)$.
Changing variables in \eqref{support:76-26} therefore gives
the same field transformation as $U(1\oplus R)$.
Both fix the vacuum, so they agree on $\mathcal D_0$
and on its closure.
\end{proof}

\subsubsection{Interchanging spacelike-separated insertions}

\begin{proposition}\label{support:76-52}
If $f,h\in\mathcal S(\mathbb R^{1+2})$ have spacelike
separated supports, then
\begin{equation}\label{support:76-53}
 \phi(f)\phi(h)\psi=\phi(h)\phi(f)\psi
                         \qquad(\psi\in\mathcal D_0).
\end{equation}
\end{proposition}

\begin{proof}
First suppose all tests have compact support.
Testing against polynomial vectors and using Hermiticity
reduces the assertion to invariance of
$\mathcal W_n^\sigma(f_1\otimes\cdots\otimes f_n)$
under interchange of adjacent entries $k,k+1$ whose
supports are spacelike separated.
Set $K=\prod_j\operatorname{supp}f_j$; zero entries are
trivial.  On $K$, write
$\xi=\vec x_{k+1}-\vec x_k$, $\tau=t_{k+1}-t_k$, and put
\[
 \begin{gathered}
 \delta=\min_K(|\xi|-|\tau|)>0,\quad A=4/\delta,\quad
 \eta_0=\min((2A)^{-1},1/2),\\
 h_*=2A\max_{x\in K,\ j<n}
           (|\vec x_{j+1}-\vec x_j|+|t_{j+1}-t_j|)+6.
 \end{gathered}
\]
We deform the two adjacent heights past each other:
\begin{equation}\label{support:76-54}
 \begin{gathered}
 h_k^\lambda=-\lambda/2,\quad h_{k+1}^\lambda=\lambda/2,\\
 h_j^\lambda=\tfrac12+(j-k-1)h_*\quad(j>k+1),\qquad
 h_j^\lambda=-\tfrac12-(k-j)h_*\quad(j<k),\\
 z_j(x;\lambda,\eta)=\iota x_j+\eta h_j^\lambda e_3,
 \qquad -1\le\lambda\le1,\quad0<\eta\le\eta_0 .
 \end{gathered}
\end{equation}
Although the vertical ordering fails during this
deformation, the points remain in the union of tilted
tubes.  Indeed choose $\sin\theta=A\eta$ and
\[
 u_\theta=\cos\theta\,e_3+\sin\theta\,(\xi/|\xi|,0).
\]
Then $\cos\theta\ge\sqrt3/2$, and direct calculation gives
\begin{equation}\label{support:76-55}
 \begin{split}
 d_\circ^{u_\theta}(z_{j+1}-z_j)
 &=\eta(h_{j+1}^\lambda-h_j^\lambda)\cos\theta\\
 &\quad+A\eta\,\frac{\xi}{|\xi|}
                   \cdot(\vec x_{j+1}-\vec x_j)
             -A\eta|t_{j+1}-t_j|.
 \end{split}
\end{equation}
For $j=k$ this is at least $-\eta+A\eta\delta=3\eta$.
For $j\ne k$ the height gap is at least $h_*$;
the definition of $h_*$ makes the same lower bound hold.
Consequently
\begin{equation}\label{support:76-56}
 z(x;\lambda,\eta)\in\mathrm{Tb}^{\rm all}_n,\qquad
 |\mathbf S_n(z(x;\lambda,\eta))|
                      \le(C_\wp\eta^{-\wp})^n
\end{equation}
uniformly in $x,\lambda,\eta$ in these ranges.
The witnessing direction may depend on $x$ and $\eta$;
the holomorphic function $\mathbf S_n$ is the same on
all of these tubes.

For a smooth $g$ supported in $K$, let
$I_g(\lambda,\eta)=\int g(x)\mathbf S_n(z(x;\lambda,\eta))\,dx$.
The Cauchy--Riemann equations and integration by parts give
\begin{equation}\label{support:76-57}
 \partial_\eta I_g=I_{-iD_\lambda g},\qquad
 \partial_\lambda I_g=\eta I_{-iD'g},\qquad
 D_\lambda=\sum_jh_j^\lambda\partial_{t_j},\quad
 D'=\tfrac12(\partial_{t_{k+1}}-\partial_{t_k}).
\end{equation}
Compactness at every positive height justifies differentiation.
\Cref{support:76-23}, with $N=\lceil\wp n\rceil+1$,
gives the uniform bound
\begin{equation}\label{support:76-58}
 |I_g(\lambda,\eta)|\le Q(g):=
 2^N\eta_0^{-(N-1/2)}C_\wp^n
          \sum_{r=0}^N\sup_{|\lambda|\le1}\|D_\lambda^rg\|_1 .
\end{equation}
The coefficients of $D_\lambda$ are uniformly bounded,
so $Q(g)<\infty$.

Take $g=\bigotimes_jf_j$.
At $\lambda=1$ the heights increase in the original
order, and \eqref{support:76-30} identifies the limit
$I_g(1,\eta)$ with the original Wightman value.
At $\lambda=-1$ they increase in the order with $k,k+1$
interchanged.  The symmetry of the single function
$\mathbf S_n$ and the same boundary formula identify
$I_g(-1,\eta)$ with the interchanged value.
But \eqref{support:76-57}--\eqref{support:76-58} imply
\[
 |I_g(1,\eta)-I_g(-1,\eta)|\le2\eta Q(D'g)\longrightarrow0.
\]
This proves the adjacent interchange for compact tests.

For polynomial vectors with other Schwartz entries,
cut off those entries and use
\Cref{support:76-38} to pass to the limit in both
products.  Density of polynomial test vectors gives
\eqref{support:76-53} for compactly supported $f,h$.
Finally cut off $f,h$ themselves inside their original
supports.  Their spacelike separation is preserved,
and graph-norm continuity passes the equality to the
original Schwartz functions.
\end{proof}

\begin{theorem}[Scalar reconstruction]\label{support:76-59}
Under (Rec), the reflection Hilbert space carries a
Hermitian scalar field $\phi$ with the following properties:
\begin{enumerate}
\item The proper orthochronous Poincar\'e group acts by
a strongly continuous unitary representation $U$ fixing
the unit vector $\Omega$.  Its translation spectrum lies
in $\overline{V_+}$.
\item The dense polynomial domain $\mathcal D_0$ contains
$\Omega$ and is invariant under $U$ and every $\phi(f)$.
For $\psi,\psi'\in\mathcal D_0$,
$f\mapsto\langle\psi,\phi(f)\psi'\rangle$ is tempered,
and the Hermiticity relation in \eqref{support:76-27} holds.
\item The field is covariant:
\[
 U(a,\Lambda)\phi(f)U(a,\Lambda)^{-1}
      =\phi\bigl(f(\Lambda^{-1}(\,\cdot-a))\bigr)
             \quad\hbox{on }\mathcal D_0,
\]
and fields with spacelike separated test supports commute
there.
\item The vacuum is cyclic.  It is the unique
translation-invariant vector up to scalar multiples
if $\ker H=\mathbb C\Omega$.
\end{enumerate}
The vacuum distributions are $\mathcal W_n^\sigma$.
Their analytic functions are $\mathbf W_n$, with
boundary values \eqref{support:76-30}, and their Euclidean
continuations recover $\sigma$ on every separated
$C_c^1$ tuple by \eqref{support:75-79}.
On $\mathcal S_{\ne}(\mathbb R^{3n})$ the Schwinger
functions are the absolutely convergent integrals
\eqref{support:75-92}.
Thus, with the additional vacuum-uniqueness hypothesis,
these objects form a scalar Wightman theory.
\end{theorem}

\begin{proof}
The domain, spectrum, cyclicity and Hermiticity were proved
in \S\ref{sec:reconstruction-fields}.
Equation \eqref{support:76-28}, followed by
\eqref{support:76-15}, gives tempered matrix elements.
Props.~\ref{support:76-44} and \ref{support:76-52}
give covariance and locality on $\mathcal D_0$.
The boundary and Euclidean identities follow from
\eqref{support:76-30}, \eqref{support:75-79} and \eqref{support:75-80}.
On differences with third coordinate at least
$\epsilon>0$ and real spatial coordinates,
\eqref{support:75-72} bounds $\mathbf W_n$ by
$(C_\wp\min(1,\epsilon/2)^{-\wp})^n$;
in particular the Euclidean continuation has the required
polynomial growth.  Finally
\eqref{support:75-89} proves the asserted absolute convergence.
\end{proof}


\begin{definition}[Euclidean axioms on flat tests]\label{support:28-35}
The test spaces are those in \Cref{support:26-2}. Isometries act by simultaneous pullback by their inverses and permutations by rearranging arguments. The axioms are growth \textup{(E0$'$)} as in \eqref{support:26-34}, Euclidean invariance \textup{(E1)}, positivity \textup{(E2)} in \eqref{support:26-42}, permutation symmetry \textup{(E3)}, and clustering \textup{(E4)} in \eqref{support:28-40} on positive-time flat tests.
\end{definition}

The passage from clustering of separated tuples to clustering of
Schwartz tests uses the reflection Hilbert space.  Directly
translating a Schwartz approximation would give bounds depending
on the translation; the Hilbert-space argument keeps them uniform.

\begin{proposition}\label{support:28-48}
Suppose a symmetric, normalized separated-tuple functional
$\sigma$ has the bounds of
\S\ref{sec:reflection-hilbert}, is Euclidean invariant and
reflection positive, and has Schwinger functions satisfying
\[
 |\mathsf S_n(\varphi)|\le c_E C_E^n(n!)^\kappa
                  |\varphi|_{kn},\qquad
 \varphi\in\mathcal S_{\ne}(\mathbb R^{3n}),
\]
where $k\in\mathbb N$ and $c_E,C_E,\kappa\ge0$.
Assume also that
\[
 \sigma(X\cdot\tau_vY)\longrightarrow\sigma(X)\sigma(Y)
 \quad(|v|\to\infty)
\]
for all separated tuples $X,Y$, with the concatenation defined
for sufficiently large $|v|$.
Then the Schwinger functions satisfy Euclidean invariance,
permutation symmetry, reflection positivity
\eqref{support:26-42}, the displayed growth bound, and
\begin{equation}\label{support:28-40}
 \lim_{s\to\infty}
 \mathsf S_{n+n'}(\Theta\bar\varphi\otimes\tau_{sv}\varphi')
 =\mathsf S_n(\Theta\bar\varphi)\mathsf S_{n'}(\varphi')
\end{equation}
for $\varphi\in\mathcal S_+(\mathbb R^{3n})$,
$\varphi'\in\mathcal S_+(\mathbb R^{3n'})$ and
$0\ne v\in\mathbb R^2\times\{0\}$.  We refer to these five
properties as \textup{(E0$'$)--(E4)}, respectively growth,
Euclidean invariance, reflection positivity, symmetry and
clustering.  The test spaces are the flat Schwartz spaces defined
in \S\ref{sec:distribution-extension}.
\end{proposition}

\begin{proof}
\Cref{support:26-41} gives the first four properties.
For the last, form the algebraic direct sum
$\underline{\mathcal S}_+=\bigoplus_{n\ge0}
\mathcal S_+(\mathbb R^{3n})$, with the positive form
\[
 \langle\underline\varphi,\underline\psi\rangle_{\mathcal S}
 =\sum_{n,n'}\mathsf S_{n+n'}
                       (\Theta\bar\varphi_n\otimes\psi_{n'}).
\]
Let $\mathcal H^{\mathcal S}$ be its completion modulo the null
space.  If both sequences vanish above degree $n_0$, their
form is bounded in absolute value by
\begin{equation}\label{support:28-49}
 \max(1,c_E)\max(1,C_E)^{2n_0}((2n_0)!)^\kappa
 \left(\sum_{n\le n_0}|\varphi_n|_{2kn_0}\right)
 \left(\sum_{n\le n_0}|\psi_n|_{2kn_0}\right).
\end{equation}
This follows term by term from the tensor-product seminorm
bound; the degree-zero term is ordinary multiplication.

Send a smooth separated positive-time tuple $X$ to the sequence
$\iota X$ with $\bigotimes X$ in its degree and zero elsewhere.
For finite sums of such tuples,
\begin{equation}\label{support:28-50}
 \langle\iota F,\iota F'\rangle_{\mathcal S}
        =\sigma(\theta\bar F\cdot F').
\end{equation}
\Cref{support:26-17} and \eqref{support:28-49} show that their
classes are dense in $\mathcal H^{\mathcal S}$.
They are also dense in $\mathcal H(\sigma)$: mollify each
entry of a $C_c^1$ tuple with a common radius tending to zero,
smaller than its distance to the reflection plane and half
its least support separation.  The resulting tuples remain
separated in the positive half-space.  The moving-entry
continuity estimate gives convergence of the three terms in
\[
 \|[X]-[X_j]\|^2
 =\sigma(\theta X\cdot X)
  -2\operatorname{Re}\sigma(\theta X\cdot X_j)
  +\sigma(\theta X_j\cdot X_j).
\]
Consequently \eqref{support:28-50} extends to a unitary
$V:\mathcal H(\sigma)\to\mathcal H^{\mathcal S}$ carrying
$\Omega$ to the constant sequence.  For a fixed tangential
translation, the seminorm estimate \eqref{support:26-43} shows
that the same approximations converge after translation.
Thus
\[
 V\mathcal U(v')V^{-1}[\underline\varphi]
       =[\tau_{(v',0)}\underline\varphi].
\]

Clustering of tuples gives, first for tuple classes,
\[
 \langle\chi,\mathcal U(sv')\chi'\rangle
       \longrightarrow
       \langle\chi,\Omega\rangle\langle\Omega,\chi'\rangle,
       \qquad v'\ne0.
\]
The difference between the two sides has modulus at most
$2\|\chi\|\|\chi'\|$, uniformly in $s$, since
$\mathcal U(sv')$ is unitary.  Density extends this limit to
all vectors.  Apply it to the inverse images under $V$ of
the classes with $\varphi$ and $\varphi'$ in degrees $n,n'$.
The form identity gives exactly \eqref{support:28-40}.
\end{proof}

\subsection{Assembly of the continuum limits}
\begin{definition}[Frame data]\label{support:77-1}
The data are a converging frame $(K_j,\nu_j)$ in classes $\mathcal G_j$ obeying the hypotheses of \S\ref{sec:conditional-decomposition}. Write $\sigma$ for its limit and $\sigma'$ for an arbitrary limit in that frame, keeping the centring by $\nu_j$ when changing the integrating laws. The reconstruction hypotheses \textup{(Rec1)} to \textup{(Rec4)} are those of \S\ref{sec:reconstruction-energy} for the functional itself. We write \textup{(Rec5)} for the condition $\ker H=\mathbb C\Omega$.
\end{definition}
\begin{definition}[Optional frame hypotheses]\label{support:77-2}
The rotation condition \textup{(Rot)} is Euclidean invariance of all limits in the frame. The point-bound condition \textup{(PtB)} asks for common constants $C_\wp,\wp,c_0$ with $0<c_0\le1/2$ for which every limit satisfies
\begin{equation}\label{support:77-5}
|\sigma'^{\rm c}(\theta\bar F\cdot(h)\cdot G)|
\le C_\wp\eta^{-\wp}\|h\|_1\|[\tau_{-\eta e_3}F]\|_{\sigma'^{\rm c}}\|[\tau_{-\eta e_3}G]\|_{\sigma'^{\rm c}}
\end{equation}
for $0<\eta\le1$, $h$ supported in $\{|x_3|<c_0\eta\}$ and $F,G$ supported in $\{x_3>\eta\}$.
\end{definition}

\begin{proposition}[Point bound and reconstruction hypotheses]\label{support:77-12}
Let $\sigma'$ be a limit in a frame as above. Under \textup{(Rot)} and \textup{(PtB)}, it has \textup{(Rec1)} to \textup{(Rec4)}. For a component it has \textup{(Rec5)} as well.
\end{proposition}
\begin{proof}
The first condition is \textup{(Rot)}. The limit properties of \Cref{support:36-12} give its one-point function, bounds and reflection positivity, so \Cref{support:27-22} applies and gives \textup{(Rec2)}. \Cref{support:28-34} gives \textup{(Rec3)}. By \textup{(PtB)}, the centered functional obeys \eqref{support:75-51}.
Apply \Cref{support:75-50} with $c=c_{\sigma'}$ to obtain
\textup{(Rec4)} with constant $C_\wp+|c_{\sigma'}|$. For components the clustering estimate of \Cref{support:36-17} and the spectral limit in \eqref{support:36-22} give \textup{(Rec5)}.
\end{proof}

\begin{theorem}[The frame limit and its components]\label{support:77-15}
For a convergent frame, the following hold.
\begin{enumerate}
\item[(1)] The frame limit $\sigma$ (a) is symmetric and equals the centred functional of its connected functions $S_n=\lim_jS_n^T[\nu_j]$, $n\ge2$, with table $\underline C$ given by the frame bounds. It is (b) translation and cubic invariant and (c) reflection positive at coordinate and diagonal mirrors. In (d), its moments and connected functions have unique off-diagonal distributional extensions of order at most $4n+1$ on compact sets. In (e), its coordinate Hilbert space is separable, has the translation semigroups of \Cref{support:27-22} and a unitary $W_2$ action, and satisfies $H\ge|P_k|$ for $k=1,2$. Under \textup{(Rot)}, (g) the invariance is Euclidean, all mirrors are positive, the spectral cone is circular and $\operatorname{O}(2)$ acts unitarily on the coordinate Hilbert space as in \Cref{support:28-34}.
\item[(2)] (a) The decomposition of $\sigma$ is canonical and unique up to the unique measure-preserving homeomorphism matching the components. (b) The map $P_0[F]\mapsto(\xi\mapsto\sigma_\xi(F))$ gives the vacuum-space unitary of \Cref{support:36-25}, including the dimension statement. (c) An isometry preserving $\sigma$ acts by a measure-preserving homeomorphism on the component space. This homeomorphism is the identity for translations and cubic isometries, and for every Euclidean isometry under \textup{(Rot)}.
\item[(3)] Each component $\sigma_\xi$ has the following properties. (i) It is a limit in the original frame. It lies in $\mathsf{Fct}(\underline C,\bar c)$ with $\bar c$ given by \eqref{support:35-14}, is invariant under translations and cubic isometries, is reflection positive at coordinate and diagonal mirrors, and is regular of order at most $4n+1$ on compact off-diagonal sets. Under \textup{(Rot)} it is Euclidean invariant and reflection positive at every mirror. (ii)(a) Its mean is $c_\xi\int f$, with $|c_\xi|\le\bar c$, $c_\xi$ continuous in $\xi$, and $\int c_\xi\,d\lambda=0$. (b) Its connected functions obey the original bound table. (iii) Its reflection Hilbert space obeys (1)(e), and its zero-energy subspace is one-dimensional. Under \textup{(Rot)}, the circular spectral cone and the $\operatorname{O}(2)$ action also hold. (iv) For each pair of tuples $X,Y$, $\sigma_\xi(X\cdot\tau_vY)-\sigma_\xi(X)\sigma_\xi(Y)$ is bounded by $C(X,Y)/|v|$ for sufficiently large $|v|$, uniformly in $\xi$.
\end{enumerate}
\end{theorem}
\begin{proof}
Apply \Cref{support:36-12} to the frame limit and to every limit in the frame, using \Cref{support:35-36} for the conditioned approximants to each component. The frame limit has zero mean because the integrating and centring laws coincide, and \eqref{support:9-21} then gives (1)(a). The same proposition gives (b) and (c). \Cref{support:25-55} gives (d) for both the moments and the connected functions. In (e), \Cref{support:27-22} gives the representation and separability, the signed permutations fixing $e_3$ preserve its reflection form, and \Cref{support:28-21} gives each spectral inequality by (c). Under \textup{(Rot)}, \Cref{support:28-34} gives (g).

Existence in (2)(a) is \Cref{support:35-36}; uniqueness and the action on components in (c) follow from \Cref{support:36-9}. The action is the identity whenever each component is invariant, since different points give different functionals. The vacuum-space statement (b) is \Cref{support:36-25} using the uniform clustering estimate of \Cref{support:36-17}. For (3)(i) and (ii) apply \Cref{support:36-12}, \eqref{support:36-15}, and reflection transport under \textup{(Rot)}. The spectral and Hilbert-space statements follow by the same argument as (1)(e), (g) and the equivalence \eqref{support:36-22}. The uniform estimate (iv) is \Cref{support:36-17}.
\end{proof}

\begin{corollary}[Reconstruction of the components]\label{support:77-14}
Under \textup{(Rot)} and \textup{(PtB)}, every component has a scalar Wightman reconstruction with unique invariant vacuum. The frame limit has the scalar reconstruction with vacuum space given by \Cref{support:36-25}. If the component Schwinger functionals have the growth bound \eqref{support:26-34}, they satisfy the Euclidean axioms of \Cref{support:28-48}.
\end{corollary}
\begin{proof}
\Cref{support:77-12} verifies the hypotheses of \Cref{support:76-59} and establishes the vacuum condition for the components. \Cref{support:28-48} applies to their Schwinger functionals by the bound, reflection positivity, Euclidean invariance and uniform component clustering.
\end{proof}

\section{Index of notation and conventions}\label{app:notation}

\begin{center}
\begin{tabular}{@{}p{.22\linewidth}p{.53\linewidth}l@{}}
\textbf{Notation}&\textbf{Meaning}&\textbf{Definition}\\[2pt]
$G$, $U$, $a_p(U)$ & Gauge group, link field and plaquette action & \S\ref{sec:wilson-conventions}\\
$L$, $K$, $a_K$ & Blocking factor, cutoff and lattice spacing & \S\ref{Y-00-1}\\
$\ell$, $\mathsf b$, $\beta$ & Physical side, label and inverse lattice coupling & \S\ref{Y-00-1}\\
$\mathcal O_K(f)$, $S_n^T$ & Smeared action density and connected functions & Def.~\ref{support:150-2}\\
$V_k$, $x_k$ & Block field and reference coupling at step $k$ & Def.~\ref{Y-00-S4}\\
$k_{ax}$, $\mathsf M$ & Coarse-link transition density and its center & \S\ref{Y-J-1}\\
$\ell_k$, $\varrho$ & Gaussian rows and their truncation radius & \S\ref{Y-A-1}\\
$E^a,E^h,w$ & Local analytic term, remainder and large-field weight & \S\ref{Y-H-1}\\
$\rho[u;\ell]$ & Effective density with interaction data $u$ and rows $\ell$ & \S\ref{Y-J-1}\\
$c_K^{(s)}$, $\hat c_K$ & Source subtractions at depth $s$ and at the last scale & Def.~\ref{Y-00-S4}\\
$\mathbb G_K(\mathsf b)$ & Translation-invariant Wilson Gibbs states & Def.~\ref{Y-00-S6}\\
$\sigma_S$ & Centered functional of the connected family $S$ & Def.~\ref{support:9-12}\\
$\mathcal S_{\ne}$ & Schwartz tests flat at coincidence diagonals & \S\ref{sec:distribution-extension}\\
$\mathcal H(\sigma),\Omega,H$ & Reflection Hilbert space, vacuum and Hamiltonian & \S\ref{sec:reflection-hilbert}\\
$\mathsf X,\lambda,\sigma_\xi$ & Component space, measure and functionals & Def.~\ref{support:35-34}\\
$B(s,s')$, $\mathbf W_n$ & Point insertion operator and holomorphic vacuum functions & \S\ref{app:scalar-reconstruction}\\
\end{tabular}
\end{center}

The cutoff $K$ increases toward the continuum, the step number $k$
increases toward the coarser lattice, and depth is $K-k$. Unless a
Euclidean distance is specified, separation of supports is measured
with the sup distance. The Lipschitz constant in $\|f\|_\sharp$ uses the
Euclidean distance. On the periodic lattices, distances are quotient
distances. Haar measure on $G$ has mass one.

\begin{thebibliography}{9}
\bibitem{OpenAI215}
OpenAI (2026).
\emph{The canonical massive continuum limit of the two-dimensional
$\operatorname{O}(3)$ model}. Preprint, 4 October 2026.
\href{https://github.com/openai/math/blob/main/preprints/The-canonical-massive-continuum-limit-of-the-two-dimensional-O3-model-October-4-2026/massive-continuum-o3.pdf}{Paper}.
\href{https://github.com/openai/math/blob/main/overview.tex}{OpenAI mathematics result family 215}.
\end{thebibliography}
\end{document}
